% Punto de entrada aditivo del sucesor; el main.tex rector queda inmutable.
\RequirePackage{fix-cm}
\documentclass[11pt,a4paper,twoside,openany]{book}

\usepackage[T1]{fontenc}
\usepackage[utf8]{inputenc}
\DeclareUnicodeCharacter{03BC}{\ensuremath{\mu}}
\DeclareUnicodeCharacter{03C4}{\ensuremath{\tau}}
\DeclareUnicodeCharacter{03C3}{\ensuremath{\sigma}}
\DeclareUnicodeCharacter{03C0}{\ensuremath{\pi}}
\DeclareUnicodeCharacter{03B3}{\ensuremath{\gamma}}
\DeclareUnicodeCharacter{03A9}{\ensuremath{\Omega}}
\DeclareUnicodeCharacter{03C1}{\ensuremath{\rho}}
\DeclareUnicodeCharacter{03C6}{\ensuremath{\varphi}}
\DeclareUnicodeCharacter{03C8}{\ensuremath{\psi}}
\DeclareUnicodeCharacter{03A5}{\ensuremath{\Upsilon}}
\DeclareUnicodeCharacter{03B7}{\ensuremath{\eta}}
\DeclareUnicodeCharacter{03B1}{\ensuremath{\alpha}}
\DeclareUnicodeCharacter{03B5}{\ensuremath{\varepsilon}}
\DeclareUnicodeCharacter{03B6}{\ensuremath{\zeta}}
\DeclareUnicodeCharacter{03B8}{\ensuremath{\theta}}
\DeclareUnicodeCharacter{03BB}{\ensuremath{\lambda}}
\DeclareUnicodeCharacter{00B1}{\ensuremath{\pm}}
\DeclareUnicodeCharacter{2192}{\ensuremath{\longrightarrow}}
\DeclareUnicodeCharacter{2194}{\ensuremath{\leftrightarrow}}
\DeclareUnicodeCharacter{21A6}{\ensuremath{\longmapsto}}
\DeclareUnicodeCharacter{21D2}{\ensuremath{\Rightarrow}}
\DeclareUnicodeCharacter{2202}{\ensuremath{\partial}}
% lcdf-typetools glyphtounicode.tex, Version 2.95
% Contents: Glyph mapping information for pdftex, used for PDF searching
% Generated from:
% - glyphlist.txt, Version 2.0
% - texglyphlist.txt, Version 2.95
% - texglyphlist-g2u.txt, Version 2.95
\pdfglyphtounicode{A}{0041}
\pdfglyphtounicode{AE}{00C6}
\pdfglyphtounicode{AEacute}{01FC}
\pdfglyphtounicode{AEmacron}{01E2}
\pdfglyphtounicode{AEsmall}{00E6}
\pdfglyphtounicode{Aacute}{00C1}
\pdfglyphtounicode{Aacutesmall}{00E1}
\pdfglyphtounicode{Abreve}{0102}
\pdfglyphtounicode{Abreveacute}{1EAE}
\pdfglyphtounicode{Abrevecyrillic}{04D0}
\pdfglyphtounicode{Abrevedotbelow}{1EB6}
\pdfglyphtounicode{Abrevegrave}{1EB0}
\pdfglyphtounicode{Abrevehookabove}{1EB2}
\pdfglyphtounicode{Abrevetilde}{1EB4}
\pdfglyphtounicode{Acaron}{01CD}
\pdfglyphtounicode{Acircle}{24B6}
\pdfglyphtounicode{Acircumflex}{00C2}
\pdfglyphtounicode{Acircumflexacute}{1EA4}
\pdfglyphtounicode{Acircumflexdotbelow}{1EAC}
\pdfglyphtounicode{Acircumflexgrave}{1EA6}
\pdfglyphtounicode{Acircumflexhookabove}{1EA8}
\pdfglyphtounicode{Acircumflexsmall}{00E2}
\pdfglyphtounicode{Acircumflextilde}{1EAA}
\pdfglyphtounicode{Acute}{00B4}
\pdfglyphtounicode{Acutesmall}{00B4}
\pdfglyphtounicode{Acyrillic}{0410}
\pdfglyphtounicode{Adblgrave}{0200}
\pdfglyphtounicode{Adieresis}{00C4}
\pdfglyphtounicode{Adieresiscyrillic}{04D2}
\pdfglyphtounicode{Adieresismacron}{01DE}
\pdfglyphtounicode{Adieresissmall}{00E4}
\pdfglyphtounicode{Adotbelow}{1EA0}
\pdfglyphtounicode{Adotmacron}{01E0}
\pdfglyphtounicode{Agrave}{00C0}
\pdfglyphtounicode{Agravesmall}{00E0}
\pdfglyphtounicode{Ahookabove}{1EA2}
\pdfglyphtounicode{Aiecyrillic}{04D4}
\pdfglyphtounicode{Ainvertedbreve}{0202}
\pdfglyphtounicode{Alpha}{0391}
\pdfglyphtounicode{Alphatonos}{0386}
\pdfglyphtounicode{Amacron}{0100}
\pdfglyphtounicode{Amonospace}{FF21}
\pdfglyphtounicode{Aogonek}{0104}
\pdfglyphtounicode{Aring}{00C5}
\pdfglyphtounicode{Aringacute}{01FA}
\pdfglyphtounicode{Aringbelow}{1E00}
\pdfglyphtounicode{Aringsmall}{00E5}
\pdfglyphtounicode{Asmall}{0061}
\pdfglyphtounicode{Atilde}{00C3}
\pdfglyphtounicode{Atildesmall}{00E3}
\pdfglyphtounicode{Aybarmenian}{0531}
\pdfglyphtounicode{B}{0042}
\pdfglyphtounicode{Bcircle}{24B7}
\pdfglyphtounicode{Bdotaccent}{1E02}
\pdfglyphtounicode{Bdotbelow}{1E04}
\pdfglyphtounicode{Becyrillic}{0411}
\pdfglyphtounicode{Benarmenian}{0532}
\pdfglyphtounicode{Beta}{0392}
\pdfglyphtounicode{Bhook}{0181}
\pdfglyphtounicode{Blinebelow}{1E06}
\pdfglyphtounicode{Bmonospace}{FF22}
\pdfglyphtounicode{Brevesmall}{02D8}
\pdfglyphtounicode{Bsmall}{0062}
\pdfglyphtounicode{Btopbar}{0182}
\pdfglyphtounicode{C}{0043}
\pdfglyphtounicode{Caarmenian}{053E}
\pdfglyphtounicode{Cacute}{0106}
\pdfglyphtounicode{Caron}{02C7}
\pdfglyphtounicode{Caronsmall}{02C7}
\pdfglyphtounicode{Ccaron}{010C}
\pdfglyphtounicode{Ccedilla}{00C7}
\pdfglyphtounicode{Ccedillaacute}{1E08}
\pdfglyphtounicode{Ccedillasmall}{00E7}
\pdfglyphtounicode{Ccircle}{24B8}
\pdfglyphtounicode{Ccircumflex}{0108}
\pdfglyphtounicode{Cdot}{010A}
\pdfglyphtounicode{Cdotaccent}{010A}
\pdfglyphtounicode{Cedillasmall}{00B8}
\pdfglyphtounicode{Chaarmenian}{0549}
\pdfglyphtounicode{Cheabkhasiancyrillic}{04BC}
\pdfglyphtounicode{Checyrillic}{0427}
\pdfglyphtounicode{Chedescenderabkhasiancyrillic}{04BE}
\pdfglyphtounicode{Chedescendercyrillic}{04B6}
\pdfglyphtounicode{Chedieresiscyrillic}{04F4}
\pdfglyphtounicode{Cheharmenian}{0543}
\pdfglyphtounicode{Chekhakassiancyrillic}{04CB}
\pdfglyphtounicode{Cheverticalstrokecyrillic}{04B8}
\pdfglyphtounicode{Chi}{03A7}
\pdfglyphtounicode{Chook}{0187}
\pdfglyphtounicode{Circumflexsmall}{02C6}
\pdfglyphtounicode{Cmonospace}{FF23}
\pdfglyphtounicode{Coarmenian}{0551}
\pdfglyphtounicode{Csmall}{0063}
\pdfglyphtounicode{D}{0044}
\pdfglyphtounicode{DZ}{01F1}
\pdfglyphtounicode{DZcaron}{01C4}
\pdfglyphtounicode{Daarmenian}{0534}
\pdfglyphtounicode{Dafrican}{0189}
\pdfglyphtounicode{Dbar}{0110}
\pdfglyphtounicode{Dcaron}{010E}
\pdfglyphtounicode{Dcedilla}{1E10}
\pdfglyphtounicode{Dcircle}{24B9}
\pdfglyphtounicode{Dcircumflexbelow}{1E12}
\pdfglyphtounicode{Dcroat}{0110}
\pdfglyphtounicode{Ddotaccent}{1E0A}
\pdfglyphtounicode{Ddotbelow}{1E0C}
\pdfglyphtounicode{Decyrillic}{0414}
\pdfglyphtounicode{Deicoptic}{03EE}
\pdfglyphtounicode{Delta}{2206}
\pdfglyphtounicode{Deltagreek}{0394}
\pdfglyphtounicode{Dhook}{018A}
\pdfglyphtounicode{Dieresis}{00A8}
\pdfglyphtounicode{DieresisAcute}{F6CC}
\pdfglyphtounicode{DieresisGrave}{F6CD}
\pdfglyphtounicode{Dieresissmall}{00A8}
\pdfglyphtounicode{Digamma}{D875 DFCB}
\pdfglyphtounicode{Digammagreek}{03DC}
\pdfglyphtounicode{Djecyrillic}{0402}
\pdfglyphtounicode{Dlinebelow}{1E0E}
\pdfglyphtounicode{Dmonospace}{FF24}
\pdfglyphtounicode{Dotaccentsmall}{02D9}
\pdfglyphtounicode{Dslash}{0110}
\pdfglyphtounicode{Dsmall}{0064}
\pdfglyphtounicode{Dtopbar}{018B}
\pdfglyphtounicode{Dz}{01F2}
\pdfglyphtounicode{Dzcaron}{01C5}
\pdfglyphtounicode{Dzeabkhasiancyrillic}{04E0}
\pdfglyphtounicode{Dzecyrillic}{0405}
\pdfglyphtounicode{Dzhecyrillic}{040F}
\pdfglyphtounicode{E}{0045}
\pdfglyphtounicode{Eacute}{00C9}
\pdfglyphtounicode{Eacutesmall}{00E9}
\pdfglyphtounicode{Ebreve}{0114}
\pdfglyphtounicode{Ecaron}{011A}
\pdfglyphtounicode{Ecedillabreve}{1E1C}
\pdfglyphtounicode{Echarmenian}{0535}
\pdfglyphtounicode{Ecircle}{24BA}
\pdfglyphtounicode{Ecircumflex}{00CA}
\pdfglyphtounicode{Ecircumflexacute}{1EBE}
\pdfglyphtounicode{Ecircumflexbelow}{1E18}
\pdfglyphtounicode{Ecircumflexdotbelow}{1EC6}
\pdfglyphtounicode{Ecircumflexgrave}{1EC0}
\pdfglyphtounicode{Ecircumflexhookabove}{1EC2}
\pdfglyphtounicode{Ecircumflexsmall}{00EA}
\pdfglyphtounicode{Ecircumflextilde}{1EC4}
\pdfglyphtounicode{Ecyrillic}{0404}
\pdfglyphtounicode{Edblgrave}{0204}
\pdfglyphtounicode{Edieresis}{00CB}
\pdfglyphtounicode{Edieresissmall}{00EB}
\pdfglyphtounicode{Edot}{0116}
\pdfglyphtounicode{Edotaccent}{0116}
\pdfglyphtounicode{Edotbelow}{1EB8}
\pdfglyphtounicode{Efcyrillic}{0424}
\pdfglyphtounicode{Egrave}{00C8}
\pdfglyphtounicode{Egravesmall}{00E8}
\pdfglyphtounicode{Eharmenian}{0537}
\pdfglyphtounicode{Ehookabove}{1EBA}
\pdfglyphtounicode{Eightroman}{2167}
\pdfglyphtounicode{Einvertedbreve}{0206}
\pdfglyphtounicode{Eiotifiedcyrillic}{0464}
\pdfglyphtounicode{Elcyrillic}{041B}
\pdfglyphtounicode{Elevenroman}{216A}
\pdfglyphtounicode{Emacron}{0112}
\pdfglyphtounicode{Emacronacute}{1E16}
\pdfglyphtounicode{Emacrongrave}{1E14}
\pdfglyphtounicode{Emcyrillic}{041C}
\pdfglyphtounicode{Emonospace}{FF25}
\pdfglyphtounicode{Encyrillic}{041D}
\pdfglyphtounicode{Endescendercyrillic}{04A2}
\pdfglyphtounicode{Eng}{014A}
\pdfglyphtounicode{Enghecyrillic}{04A4}
\pdfglyphtounicode{Enhookcyrillic}{04C7}
\pdfglyphtounicode{Eogonek}{0118}
\pdfglyphtounicode{Eopen}{0190}
\pdfglyphtounicode{Epsilon}{0395}
\pdfglyphtounicode{Epsilontonos}{0388}
\pdfglyphtounicode{Ercyrillic}{0420}
\pdfglyphtounicode{Ereversed}{018E}
\pdfglyphtounicode{Ereversedcyrillic}{042D}
\pdfglyphtounicode{Escyrillic}{0421}
\pdfglyphtounicode{Esdescendercyrillic}{04AA}
\pdfglyphtounicode{Esh}{01A9}
\pdfglyphtounicode{Esmall}{0065}
\pdfglyphtounicode{Eta}{0397}
\pdfglyphtounicode{Etarmenian}{0538}
\pdfglyphtounicode{Etatonos}{0389}
\pdfglyphtounicode{Eth}{00D0}
\pdfglyphtounicode{Ethsmall}{00F0}
\pdfglyphtounicode{Etilde}{1EBC}
\pdfglyphtounicode{Etildebelow}{1E1A}
\pdfglyphtounicode{Euro}{20AC}
\pdfglyphtounicode{Ezh}{01B7}
\pdfglyphtounicode{Ezhcaron}{01EE}
\pdfglyphtounicode{Ezhreversed}{01B8}
\pdfglyphtounicode{F}{0046}
\pdfglyphtounicode{FFIsmall}{0066 0066 0069}
\pdfglyphtounicode{FFLsmall}{0066 0066 006C}
\pdfglyphtounicode{FFsmall}{0066 0066}
\pdfglyphtounicode{FIsmall}{0066 0069}
\pdfglyphtounicode{FLsmall}{0066 006C}
\pdfglyphtounicode{Fcircle}{24BB}
\pdfglyphtounicode{Fdotaccent}{1E1E}
\pdfglyphtounicode{Feharmenian}{0556}
\pdfglyphtounicode{Feicoptic}{03E4}
\pdfglyphtounicode{Fhook}{0191}
\pdfglyphtounicode{Finv}{2132}
\pdfglyphtounicode{Fitacyrillic}{0472}
\pdfglyphtounicode{Fiveroman}{2164}
\pdfglyphtounicode{Fmonospace}{FF26}
\pdfglyphtounicode{Fourroman}{2163}
\pdfglyphtounicode{Fsmall}{0066}
\pdfglyphtounicode{G}{0047}
\pdfglyphtounicode{GBsquare}{3387}
\pdfglyphtounicode{Gacute}{01F4}
\pdfglyphtounicode{Gamma}{0393}
\pdfglyphtounicode{Gammaafrican}{0194}
\pdfglyphtounicode{Gangiacoptic}{03EA}
\pdfglyphtounicode{Gbreve}{011E}
\pdfglyphtounicode{Gcaron}{01E6}
\pdfglyphtounicode{Gcedilla}{0122}
\pdfglyphtounicode{Gcircle}{24BC}
\pdfglyphtounicode{Gcircumflex}{011C}
\pdfglyphtounicode{Gcommaaccent}{0122}
\pdfglyphtounicode{Gdot}{0120}
\pdfglyphtounicode{Gdotaccent}{0120}
\pdfglyphtounicode{Gecyrillic}{0413}
\pdfglyphtounicode{Germandbls}{0053 0053}
\pdfglyphtounicode{Germandblssmall}{0073 0073}
\pdfglyphtounicode{Ghadarmenian}{0542}
\pdfglyphtounicode{Ghemiddlehookcyrillic}{0494}
\pdfglyphtounicode{Ghestrokecyrillic}{0492}
\pdfglyphtounicode{Gheupturncyrillic}{0490}
\pdfglyphtounicode{Ghook}{0193}
\pdfglyphtounicode{Gimarmenian}{0533}
\pdfglyphtounicode{Gjecyrillic}{0403}
\pdfglyphtounicode{Gmacron}{1E20}
\pdfglyphtounicode{Gmir}{2141}
\pdfglyphtounicode{Gmonospace}{FF27}
\pdfglyphtounicode{Grave}{0060}
\pdfglyphtounicode{Gravesmall}{0060}
\pdfglyphtounicode{Gsmall}{0067}
\pdfglyphtounicode{Gsmallhook}{029B}
\pdfglyphtounicode{Gstroke}{01E4}
\pdfglyphtounicode{H}{0048}
\pdfglyphtounicode{H18533}{25CF}
\pdfglyphtounicode{H18543}{25AA}
\pdfglyphtounicode{H18551}{25AB}
\pdfglyphtounicode{H22073}{25A1}
\pdfglyphtounicode{HPsquare}{33CB}
\pdfglyphtounicode{Haabkhasiancyrillic}{04A8}
\pdfglyphtounicode{Hadescendercyrillic}{04B2}
\pdfglyphtounicode{Hardsigncyrillic}{042A}
\pdfglyphtounicode{Hbar}{0126}
\pdfglyphtounicode{Hbrevebelow}{1E2A}
\pdfglyphtounicode{Hcedilla}{1E28}
\pdfglyphtounicode{Hcircle}{24BD}
\pdfglyphtounicode{Hcircumflex}{0124}
\pdfglyphtounicode{Hdieresis}{1E26}
\pdfglyphtounicode{Hdotaccent}{1E22}
\pdfglyphtounicode{Hdotbelow}{1E24}
\pdfglyphtounicode{Hmonospace}{FF28}
\pdfglyphtounicode{Hoarmenian}{0540}
\pdfglyphtounicode{Horicoptic}{03E8}
\pdfglyphtounicode{Hsmall}{0068}
\pdfglyphtounicode{Hungarumlaut}{02DD}
\pdfglyphtounicode{Hungarumlautsmall}{02DD}
\pdfglyphtounicode{Hzsquare}{3390}
\pdfglyphtounicode{I}{0049}
\pdfglyphtounicode{IAcyrillic}{042F}
\pdfglyphtounicode{IJ}{0132}
\pdfglyphtounicode{IUcyrillic}{042E}
\pdfglyphtounicode{Iacute}{00CD}
\pdfglyphtounicode{Iacutesmall}{00ED}
\pdfglyphtounicode{Ibreve}{012C}
\pdfglyphtounicode{Icaron}{01CF}
\pdfglyphtounicode{Icircle}{24BE}
\pdfglyphtounicode{Icircumflex}{00CE}
\pdfglyphtounicode{Icircumflexsmall}{00EE}
\pdfglyphtounicode{Icyrillic}{0406}
\pdfglyphtounicode{Idblgrave}{0208}
\pdfglyphtounicode{Idieresis}{00CF}
\pdfglyphtounicode{Idieresisacute}{1E2E}
\pdfglyphtounicode{Idieresiscyrillic}{04E4}
\pdfglyphtounicode{Idieresissmall}{00EF}
\pdfglyphtounicode{Idot}{0130}
\pdfglyphtounicode{Idotaccent}{0130}
\pdfglyphtounicode{Idotbelow}{1ECA}
\pdfglyphtounicode{Iebrevecyrillic}{04D6}
\pdfglyphtounicode{Iecyrillic}{0415}
\pdfglyphtounicode{Ifractur}{2111}
\pdfglyphtounicode{Ifraktur}{2111}
\pdfglyphtounicode{Igrave}{00CC}
\pdfglyphtounicode{Igravesmall}{00EC}
\pdfglyphtounicode{Ihookabove}{1EC8}
\pdfglyphtounicode{Iicyrillic}{0418}
\pdfglyphtounicode{Iinvertedbreve}{020A}
\pdfglyphtounicode{Iishortcyrillic}{0419}
\pdfglyphtounicode{Imacron}{012A}
\pdfglyphtounicode{Imacroncyrillic}{04E2}
\pdfglyphtounicode{Imonospace}{FF29}
\pdfglyphtounicode{Iniarmenian}{053B}
\pdfglyphtounicode{Iocyrillic}{0401}
\pdfglyphtounicode{Iogonek}{012E}
\pdfglyphtounicode{Iota}{0399}
\pdfglyphtounicode{Iotaafrican}{0196}
\pdfglyphtounicode{Iotadieresis}{03AA}
\pdfglyphtounicode{Iotatonos}{038A}
\pdfglyphtounicode{Ismall}{0069}
\pdfglyphtounicode{Istroke}{0197}
\pdfglyphtounicode{Itilde}{0128}
\pdfglyphtounicode{Itildebelow}{1E2C}
\pdfglyphtounicode{Izhitsacyrillic}{0474}
\pdfglyphtounicode{Izhitsadblgravecyrillic}{0476}
\pdfglyphtounicode{J}{004A}
\pdfglyphtounicode{Jaarmenian}{0541}
\pdfglyphtounicode{Jcircle}{24BF}
\pdfglyphtounicode{Jcircumflex}{0134}
\pdfglyphtounicode{Jecyrillic}{0408}
\pdfglyphtounicode{Jheharmenian}{054B}
\pdfglyphtounicode{Jmonospace}{FF2A}
\pdfglyphtounicode{Jsmall}{006A}
\pdfglyphtounicode{K}{004B}
\pdfglyphtounicode{KBsquare}{3385}
\pdfglyphtounicode{KKsquare}{33CD}
\pdfglyphtounicode{Kabashkircyrillic}{04A0}
\pdfglyphtounicode{Kacute}{1E30}
\pdfglyphtounicode{Kacyrillic}{041A}
\pdfglyphtounicode{Kadescendercyrillic}{049A}
\pdfglyphtounicode{Kahookcyrillic}{04C3}
\pdfglyphtounicode{Kappa}{039A}
\pdfglyphtounicode{Kastrokecyrillic}{049E}
\pdfglyphtounicode{Kaverticalstrokecyrillic}{049C}
\pdfglyphtounicode{Kcaron}{01E8}
\pdfglyphtounicode{Kcedilla}{0136}
\pdfglyphtounicode{Kcircle}{24C0}
\pdfglyphtounicode{Kcommaaccent}{0136}
\pdfglyphtounicode{Kdotbelow}{1E32}
\pdfglyphtounicode{Keharmenian}{0554}
\pdfglyphtounicode{Kenarmenian}{053F}
\pdfglyphtounicode{Khacyrillic}{0425}
\pdfglyphtounicode{Kheicoptic}{03E6}
\pdfglyphtounicode{Khook}{0198}
\pdfglyphtounicode{Kjecyrillic}{040C}
\pdfglyphtounicode{Klinebelow}{1E34}
\pdfglyphtounicode{Kmonospace}{FF2B}
\pdfglyphtounicode{Koppacyrillic}{0480}
\pdfglyphtounicode{Koppagreek}{03DE}
\pdfglyphtounicode{Ksicyrillic}{046E}
\pdfglyphtounicode{Ksmall}{006B}
\pdfglyphtounicode{L}{004C}
\pdfglyphtounicode{LJ}{01C7}
\pdfglyphtounicode{LL}{004C 004C}
\pdfglyphtounicode{Lacute}{0139}
\pdfglyphtounicode{Lambda}{039B}
\pdfglyphtounicode{Lcaron}{013D}
\pdfglyphtounicode{Lcedilla}{013B}
\pdfglyphtounicode{Lcircle}{24C1}
\pdfglyphtounicode{Lcircumflexbelow}{1E3C}
\pdfglyphtounicode{Lcommaaccent}{013B}
\pdfglyphtounicode{Ldot}{013F}
\pdfglyphtounicode{Ldotaccent}{013F}
\pdfglyphtounicode{Ldotbelow}{1E36}
\pdfglyphtounicode{Ldotbelowmacron}{1E38}
\pdfglyphtounicode{Liwnarmenian}{053C}
\pdfglyphtounicode{Lj}{01C8}
\pdfglyphtounicode{Ljecyrillic}{0409}
\pdfglyphtounicode{Llinebelow}{1E3A}
\pdfglyphtounicode{Lmonospace}{FF2C}
\pdfglyphtounicode{Lslash}{0141}
\pdfglyphtounicode{Lslashsmall}{0142}
\pdfglyphtounicode{Lsmall}{006C}
\pdfglyphtounicode{M}{004D}
\pdfglyphtounicode{MBsquare}{3386}
\pdfglyphtounicode{Macron}{00AF}
\pdfglyphtounicode{Macronsmall}{00AF}
\pdfglyphtounicode{Macute}{1E3E}
\pdfglyphtounicode{Mcircle}{24C2}
\pdfglyphtounicode{Mdotaccent}{1E40}
\pdfglyphtounicode{Mdotbelow}{1E42}
\pdfglyphtounicode{Menarmenian}{0544}
\pdfglyphtounicode{Mmonospace}{FF2D}
\pdfglyphtounicode{Msmall}{006D}
\pdfglyphtounicode{Mturned}{019C}
\pdfglyphtounicode{Mu}{039C}
\pdfglyphtounicode{N}{004E}
\pdfglyphtounicode{NJ}{01CA}
\pdfglyphtounicode{Nacute}{0143}
\pdfglyphtounicode{Ncaron}{0147}
\pdfglyphtounicode{Ncedilla}{0145}
\pdfglyphtounicode{Ncircle}{24C3}
\pdfglyphtounicode{Ncircumflexbelow}{1E4A}
\pdfglyphtounicode{Ncommaaccent}{0145}
\pdfglyphtounicode{Ndotaccent}{1E44}
\pdfglyphtounicode{Ndotbelow}{1E46}
\pdfglyphtounicode{Ng}{014A}
\pdfglyphtounicode{Nhookleft}{019D}
\pdfglyphtounicode{Nineroman}{2168}
\pdfglyphtounicode{Nj}{01CB}
\pdfglyphtounicode{Njecyrillic}{040A}
\pdfglyphtounicode{Nlinebelow}{1E48}
\pdfglyphtounicode{Nmonospace}{FF2E}
\pdfglyphtounicode{Nowarmenian}{0546}
\pdfglyphtounicode{Nsmall}{006E}
\pdfglyphtounicode{Ntilde}{00D1}
\pdfglyphtounicode{Ntildesmall}{00F1}
\pdfglyphtounicode{Nu}{039D}
\pdfglyphtounicode{O}{004F}
\pdfglyphtounicode{OE}{0152}
\pdfglyphtounicode{OEsmall}{0153}
\pdfglyphtounicode{Oacute}{00D3}
\pdfglyphtounicode{Oacutesmall}{00F3}
\pdfglyphtounicode{Obarredcyrillic}{04E8}
\pdfglyphtounicode{Obarreddieresiscyrillic}{04EA}
\pdfglyphtounicode{Obreve}{014E}
\pdfglyphtounicode{Ocaron}{01D1}
\pdfglyphtounicode{Ocenteredtilde}{019F}
\pdfglyphtounicode{Ocircle}{24C4}
\pdfglyphtounicode{Ocircumflex}{00D4}
\pdfglyphtounicode{Ocircumflexacute}{1ED0}
\pdfglyphtounicode{Ocircumflexdotbelow}{1ED8}
\pdfglyphtounicode{Ocircumflexgrave}{1ED2}
\pdfglyphtounicode{Ocircumflexhookabove}{1ED4}
\pdfglyphtounicode{Ocircumflexsmall}{00F4}
\pdfglyphtounicode{Ocircumflextilde}{1ED6}
\pdfglyphtounicode{Ocyrillic}{041E}
\pdfglyphtounicode{Odblacute}{0150}
\pdfglyphtounicode{Odblgrave}{020C}
\pdfglyphtounicode{Odieresis}{00D6}
\pdfglyphtounicode{Odieresiscyrillic}{04E6}
\pdfglyphtounicode{Odieresissmall}{00F6}
\pdfglyphtounicode{Odotbelow}{1ECC}
\pdfglyphtounicode{Ogoneksmall}{02DB}
\pdfglyphtounicode{Ograve}{00D2}
\pdfglyphtounicode{Ogravesmall}{00F2}
\pdfglyphtounicode{Oharmenian}{0555}
\pdfglyphtounicode{Ohm}{2126}
\pdfglyphtounicode{Ohookabove}{1ECE}
\pdfglyphtounicode{Ohorn}{01A0}
\pdfglyphtounicode{Ohornacute}{1EDA}
\pdfglyphtounicode{Ohorndotbelow}{1EE2}
\pdfglyphtounicode{Ohorngrave}{1EDC}
\pdfglyphtounicode{Ohornhookabove}{1EDE}
\pdfglyphtounicode{Ohorntilde}{1EE0}
\pdfglyphtounicode{Ohungarumlaut}{0150}
\pdfglyphtounicode{Oi}{01A2}
\pdfglyphtounicode{Oinvertedbreve}{020E}
\pdfglyphtounicode{Omacron}{014C}
\pdfglyphtounicode{Omacronacute}{1E52}
\pdfglyphtounicode{Omacrongrave}{1E50}
\pdfglyphtounicode{Omega}{2126}
\pdfglyphtounicode{Omegacyrillic}{0460}
\pdfglyphtounicode{Omegagreek}{03A9}
\pdfglyphtounicode{Omegainv}{2127}
\pdfglyphtounicode{Omegaroundcyrillic}{047A}
\pdfglyphtounicode{Omegatitlocyrillic}{047C}
\pdfglyphtounicode{Omegatonos}{038F}
\pdfglyphtounicode{Omicron}{039F}
\pdfglyphtounicode{Omicrontonos}{038C}
\pdfglyphtounicode{Omonospace}{FF2F}
\pdfglyphtounicode{Oneroman}{2160}
\pdfglyphtounicode{Oogonek}{01EA}
\pdfglyphtounicode{Oogonekmacron}{01EC}
\pdfglyphtounicode{Oopen}{0186}
\pdfglyphtounicode{Oslash}{00D8}
\pdfglyphtounicode{Oslashacute}{01FE}
\pdfglyphtounicode{Oslashsmall}{00F8}
\pdfglyphtounicode{Osmall}{006F}
\pdfglyphtounicode{Ostrokeacute}{01FE}
\pdfglyphtounicode{Otcyrillic}{047E}
\pdfglyphtounicode{Otilde}{00D5}
\pdfglyphtounicode{Otildeacute}{1E4C}
\pdfglyphtounicode{Otildedieresis}{1E4E}
\pdfglyphtounicode{Otildesmall}{00F5}
\pdfglyphtounicode{P}{0050}
\pdfglyphtounicode{Pacute}{1E54}
\pdfglyphtounicode{Pcircle}{24C5}
\pdfglyphtounicode{Pdotaccent}{1E56}
\pdfglyphtounicode{Pecyrillic}{041F}
\pdfglyphtounicode{Peharmenian}{054A}
\pdfglyphtounicode{Pemiddlehookcyrillic}{04A6}
\pdfglyphtounicode{Phi}{03A6}
\pdfglyphtounicode{Phook}{01A4}
\pdfglyphtounicode{Pi}{03A0}
\pdfglyphtounicode{Piwrarmenian}{0553}
\pdfglyphtounicode{Pmonospace}{FF30}
\pdfglyphtounicode{Psi}{03A8}
\pdfglyphtounicode{Psicyrillic}{0470}
\pdfglyphtounicode{Psmall}{0070}
\pdfglyphtounicode{Q}{0051}
\pdfglyphtounicode{Qcircle}{24C6}
\pdfglyphtounicode{Qmonospace}{FF31}
\pdfglyphtounicode{Qsmall}{0071}
\pdfglyphtounicode{R}{0052}
\pdfglyphtounicode{Raarmenian}{054C}
\pdfglyphtounicode{Racute}{0154}
\pdfglyphtounicode{Rcaron}{0158}
\pdfglyphtounicode{Rcedilla}{0156}
\pdfglyphtounicode{Rcircle}{24C7}
\pdfglyphtounicode{Rcommaaccent}{0156}
\pdfglyphtounicode{Rdblgrave}{0210}
\pdfglyphtounicode{Rdotaccent}{1E58}
\pdfglyphtounicode{Rdotbelow}{1E5A}
\pdfglyphtounicode{Rdotbelowmacron}{1E5C}
\pdfglyphtounicode{Reharmenian}{0550}
\pdfglyphtounicode{Rfractur}{211C}
\pdfglyphtounicode{Rfraktur}{211C}
\pdfglyphtounicode{Rho}{03A1}
\pdfglyphtounicode{Ringsmall}{02DA}
\pdfglyphtounicode{Rinvertedbreve}{0212}
\pdfglyphtounicode{Rlinebelow}{1E5E}
\pdfglyphtounicode{Rmonospace}{FF32}
\pdfglyphtounicode{Rsmall}{0072}
\pdfglyphtounicode{Rsmallinverted}{0281}
\pdfglyphtounicode{Rsmallinvertedsuperior}{02B6}
\pdfglyphtounicode{S}{0053}
\pdfglyphtounicode{SF010000}{250C}
\pdfglyphtounicode{SF020000}{2514}
\pdfglyphtounicode{SF030000}{2510}
\pdfglyphtounicode{SF040000}{2518}
\pdfglyphtounicode{SF050000}{253C}
\pdfglyphtounicode{SF060000}{252C}
\pdfglyphtounicode{SF070000}{2534}
\pdfglyphtounicode{SF080000}{251C}
\pdfglyphtounicode{SF090000}{2524}
\pdfglyphtounicode{SF100000}{2500}
\pdfglyphtounicode{SF110000}{2502}
\pdfglyphtounicode{SF190000}{2561}
\pdfglyphtounicode{SF200000}{2562}
\pdfglyphtounicode{SF210000}{2556}
\pdfglyphtounicode{SF220000}{2555}
\pdfglyphtounicode{SF230000}{2563}
\pdfglyphtounicode{SF240000}{2551}
\pdfglyphtounicode{SF250000}{2557}
\pdfglyphtounicode{SF260000}{255D}
\pdfglyphtounicode{SF270000}{255C}
\pdfglyphtounicode{SF280000}{255B}
\pdfglyphtounicode{SF360000}{255E}
\pdfglyphtounicode{SF370000}{255F}
\pdfglyphtounicode{SF380000}{255A}
\pdfglyphtounicode{SF390000}{2554}
\pdfglyphtounicode{SF400000}{2569}
\pdfglyphtounicode{SF410000}{2566}
\pdfglyphtounicode{SF420000}{2560}
\pdfglyphtounicode{SF430000}{2550}
\pdfglyphtounicode{SF440000}{256C}
\pdfglyphtounicode{SF450000}{2567}
\pdfglyphtounicode{SF460000}{2568}
\pdfglyphtounicode{SF470000}{2564}
\pdfglyphtounicode{SF480000}{2565}
\pdfglyphtounicode{SF490000}{2559}
\pdfglyphtounicode{SF500000}{2558}
\pdfglyphtounicode{SF510000}{2552}
\pdfglyphtounicode{SF520000}{2553}
\pdfglyphtounicode{SF530000}{256B}
\pdfglyphtounicode{SF540000}{256A}
\pdfglyphtounicode{SS}{0053 0053}
\pdfglyphtounicode{SSsmall}{0073 0073}
\pdfglyphtounicode{Sacute}{015A}
\pdfglyphtounicode{Sacutedotaccent}{1E64}
\pdfglyphtounicode{Sampigreek}{03E0}
\pdfglyphtounicode{Scaron}{0160}
\pdfglyphtounicode{Scarondotaccent}{1E66}
\pdfglyphtounicode{Scaronsmall}{0161}
\pdfglyphtounicode{Scedilla}{015E}
\pdfglyphtounicode{Schwa}{018F}
\pdfglyphtounicode{Schwacyrillic}{04D8}
\pdfglyphtounicode{Schwadieresiscyrillic}{04DA}
\pdfglyphtounicode{Scircle}{24C8}
\pdfglyphtounicode{Scircumflex}{015C}
\pdfglyphtounicode{Scommaaccent}{0218}
\pdfglyphtounicode{Sdotaccent}{1E60}
\pdfglyphtounicode{Sdotbelow}{1E62}
\pdfglyphtounicode{Sdotbelowdotaccent}{1E68}
\pdfglyphtounicode{Seharmenian}{054D}
\pdfglyphtounicode{Sevenroman}{2166}
\pdfglyphtounicode{Shaarmenian}{0547}
\pdfglyphtounicode{Shacyrillic}{0428}
\pdfglyphtounicode{Shchacyrillic}{0429}
\pdfglyphtounicode{Sheicoptic}{03E2}
\pdfglyphtounicode{Shhacyrillic}{04BA}
\pdfglyphtounicode{Shimacoptic}{03EC}
\pdfglyphtounicode{Sigma}{03A3}
\pdfglyphtounicode{Sixroman}{2165}
\pdfglyphtounicode{Smonospace}{FF33}
\pdfglyphtounicode{Softsigncyrillic}{042C}
\pdfglyphtounicode{Ssmall}{0073}
\pdfglyphtounicode{Stigmagreek}{03DA}
\pdfglyphtounicode{T}{0054}
\pdfglyphtounicode{Tau}{03A4}
\pdfglyphtounicode{Tbar}{0166}
\pdfglyphtounicode{Tcaron}{0164}
\pdfglyphtounicode{Tcedilla}{0162}
\pdfglyphtounicode{Tcircle}{24C9}
\pdfglyphtounicode{Tcircumflexbelow}{1E70}
\pdfglyphtounicode{Tcommaaccent}{0162}
\pdfglyphtounicode{Tdotaccent}{1E6A}
\pdfglyphtounicode{Tdotbelow}{1E6C}
\pdfglyphtounicode{Tecyrillic}{0422}
\pdfglyphtounicode{Tedescendercyrillic}{04AC}
\pdfglyphtounicode{Tenroman}{2169}
\pdfglyphtounicode{Tetsecyrillic}{04B4}
\pdfglyphtounicode{Theta}{0398}
\pdfglyphtounicode{Thook}{01AC}
\pdfglyphtounicode{Thorn}{00DE}
\pdfglyphtounicode{Thornsmall}{00FE}
\pdfglyphtounicode{Threeroman}{2162}
\pdfglyphtounicode{Tildesmall}{02DC}
\pdfglyphtounicode{Tiwnarmenian}{054F}
\pdfglyphtounicode{Tlinebelow}{1E6E}
\pdfglyphtounicode{Tmonospace}{FF34}
\pdfglyphtounicode{Toarmenian}{0539}
\pdfglyphtounicode{Tonefive}{01BC}
\pdfglyphtounicode{Tonesix}{0184}
\pdfglyphtounicode{Tonetwo}{01A7}
\pdfglyphtounicode{Tretroflexhook}{01AE}
\pdfglyphtounicode{Tsecyrillic}{0426}
\pdfglyphtounicode{Tshecyrillic}{040B}
\pdfglyphtounicode{Tsmall}{0074}
\pdfglyphtounicode{Twelveroman}{216B}
\pdfglyphtounicode{Tworoman}{2161}
\pdfglyphtounicode{U}{0055}
\pdfglyphtounicode{Uacute}{00DA}
\pdfglyphtounicode{Uacutesmall}{00FA}
\pdfglyphtounicode{Ubreve}{016C}
\pdfglyphtounicode{Ucaron}{01D3}
\pdfglyphtounicode{Ucircle}{24CA}
\pdfglyphtounicode{Ucircumflex}{00DB}
\pdfglyphtounicode{Ucircumflexbelow}{1E76}
\pdfglyphtounicode{Ucircumflexsmall}{00FB}
\pdfglyphtounicode{Ucyrillic}{0423}
\pdfglyphtounicode{Udblacute}{0170}
\pdfglyphtounicode{Udblgrave}{0214}
\pdfglyphtounicode{Udieresis}{00DC}
\pdfglyphtounicode{Udieresisacute}{01D7}
\pdfglyphtounicode{Udieresisbelow}{1E72}
\pdfglyphtounicode{Udieresiscaron}{01D9}
\pdfglyphtounicode{Udieresiscyrillic}{04F0}
\pdfglyphtounicode{Udieresisgrave}{01DB}
\pdfglyphtounicode{Udieresismacron}{01D5}
\pdfglyphtounicode{Udieresissmall}{00FC}
\pdfglyphtounicode{Udotbelow}{1EE4}
\pdfglyphtounicode{Ugrave}{00D9}
\pdfglyphtounicode{Ugravesmall}{00F9}
\pdfglyphtounicode{Uhookabove}{1EE6}
\pdfglyphtounicode{Uhorn}{01AF}
\pdfglyphtounicode{Uhornacute}{1EE8}
\pdfglyphtounicode{Uhorndotbelow}{1EF0}
\pdfglyphtounicode{Uhorngrave}{1EEA}
\pdfglyphtounicode{Uhornhookabove}{1EEC}
\pdfglyphtounicode{Uhorntilde}{1EEE}
\pdfglyphtounicode{Uhungarumlaut}{0170}
\pdfglyphtounicode{Uhungarumlautcyrillic}{04F2}
\pdfglyphtounicode{Uinvertedbreve}{0216}
\pdfglyphtounicode{Ukcyrillic}{0478}
\pdfglyphtounicode{Umacron}{016A}
\pdfglyphtounicode{Umacroncyrillic}{04EE}
\pdfglyphtounicode{Umacrondieresis}{1E7A}
\pdfglyphtounicode{Umonospace}{FF35}
\pdfglyphtounicode{Uogonek}{0172}
\pdfglyphtounicode{Upsilon}{03A5}
\pdfglyphtounicode{Upsilon1}{03D2}
\pdfglyphtounicode{Upsilonacutehooksymbolgreek}{03D3}
\pdfglyphtounicode{Upsilonafrican}{01B1}
\pdfglyphtounicode{Upsilondieresis}{03AB}
\pdfglyphtounicode{Upsilondieresishooksymbolgreek}{03D4}
\pdfglyphtounicode{Upsilonhooksymbol}{03D2}
\pdfglyphtounicode{Upsilontonos}{038E}
\pdfglyphtounicode{Uring}{016E}
\pdfglyphtounicode{Ushortcyrillic}{040E}
\pdfglyphtounicode{Usmall}{0075}
\pdfglyphtounicode{Ustraightcyrillic}{04AE}
\pdfglyphtounicode{Ustraightstrokecyrillic}{04B0}
\pdfglyphtounicode{Utilde}{0168}
\pdfglyphtounicode{Utildeacute}{1E78}
\pdfglyphtounicode{Utildebelow}{1E74}
\pdfglyphtounicode{V}{0056}
\pdfglyphtounicode{Vcircle}{24CB}
\pdfglyphtounicode{Vdotbelow}{1E7E}
\pdfglyphtounicode{Vecyrillic}{0412}
\pdfglyphtounicode{Vewarmenian}{054E}
\pdfglyphtounicode{Vhook}{01B2}
\pdfglyphtounicode{Vmonospace}{FF36}
\pdfglyphtounicode{Voarmenian}{0548}
\pdfglyphtounicode{Vsmall}{0076}
\pdfglyphtounicode{Vtilde}{1E7C}
\pdfglyphtounicode{W}{0057}
\pdfglyphtounicode{Wacute}{1E82}
\pdfglyphtounicode{Wcircle}{24CC}
\pdfglyphtounicode{Wcircumflex}{0174}
\pdfglyphtounicode{Wdieresis}{1E84}
\pdfglyphtounicode{Wdotaccent}{1E86}
\pdfglyphtounicode{Wdotbelow}{1E88}
\pdfglyphtounicode{Wgrave}{1E80}
\pdfglyphtounicode{Wmonospace}{FF37}
\pdfglyphtounicode{Wsmall}{0077}
\pdfglyphtounicode{X}{0058}
\pdfglyphtounicode{Xcircle}{24CD}
\pdfglyphtounicode{Xdieresis}{1E8C}
\pdfglyphtounicode{Xdotaccent}{1E8A}
\pdfglyphtounicode{Xeharmenian}{053D}
\pdfglyphtounicode{Xi}{039E}
\pdfglyphtounicode{Xmonospace}{FF38}
\pdfglyphtounicode{Xsmall}{0078}
\pdfglyphtounicode{Y}{0059}
\pdfglyphtounicode{Yacute}{00DD}
\pdfglyphtounicode{Yacutesmall}{00FD}
\pdfglyphtounicode{Yatcyrillic}{0462}
\pdfglyphtounicode{Ycircle}{24CE}
\pdfglyphtounicode{Ycircumflex}{0176}
\pdfglyphtounicode{Ydieresis}{0178}
\pdfglyphtounicode{Ydieresissmall}{00FF}
\pdfglyphtounicode{Ydotaccent}{1E8E}
\pdfglyphtounicode{Ydotbelow}{1EF4}
\pdfglyphtounicode{Yen}{00A5}
\pdfglyphtounicode{Yericyrillic}{042B}
\pdfglyphtounicode{Yerudieresiscyrillic}{04F8}
\pdfglyphtounicode{Ygrave}{1EF2}
\pdfglyphtounicode{Yhook}{01B3}
\pdfglyphtounicode{Yhookabove}{1EF6}
\pdfglyphtounicode{Yiarmenian}{0545}
\pdfglyphtounicode{Yicyrillic}{0407}
\pdfglyphtounicode{Yiwnarmenian}{0552}
\pdfglyphtounicode{Ymonospace}{FF39}
\pdfglyphtounicode{Ysmall}{0079}
\pdfglyphtounicode{Ytilde}{1EF8}
\pdfglyphtounicode{Yusbigcyrillic}{046A}
\pdfglyphtounicode{Yusbigiotifiedcyrillic}{046C}
\pdfglyphtounicode{Yuslittlecyrillic}{0466}
\pdfglyphtounicode{Yuslittleiotifiedcyrillic}{0468}
\pdfglyphtounicode{Z}{005A}
\pdfglyphtounicode{Zaarmenian}{0536}
\pdfglyphtounicode{Zacute}{0179}
\pdfglyphtounicode{Zcaron}{017D}
\pdfglyphtounicode{Zcaronsmall}{017E}
\pdfglyphtounicode{Zcircle}{24CF}
\pdfglyphtounicode{Zcircumflex}{1E90}
\pdfglyphtounicode{Zdot}{017B}
\pdfglyphtounicode{Zdotaccent}{017B}
\pdfglyphtounicode{Zdotbelow}{1E92}
\pdfglyphtounicode{Zecyrillic}{0417}
\pdfglyphtounicode{Zedescendercyrillic}{0498}
\pdfglyphtounicode{Zedieresiscyrillic}{04DE}
\pdfglyphtounicode{Zeta}{0396}
\pdfglyphtounicode{Zhearmenian}{053A}
\pdfglyphtounicode{Zhebrevecyrillic}{04C1}
\pdfglyphtounicode{Zhecyrillic}{0416}
\pdfglyphtounicode{Zhedescendercyrillic}{0496}
\pdfglyphtounicode{Zhedieresiscyrillic}{04DC}
\pdfglyphtounicode{Zlinebelow}{1E94}
\pdfglyphtounicode{Zmonospace}{FF3A}
\pdfglyphtounicode{Zsmall}{007A}
\pdfglyphtounicode{Zstroke}{01B5}
\pdfglyphtounicode{a}{0061}
\pdfglyphtounicode{aabengali}{0986}
\pdfglyphtounicode{aacute}{00E1}
\pdfglyphtounicode{aadeva}{0906}
\pdfglyphtounicode{aagujarati}{0A86}
\pdfglyphtounicode{aagurmukhi}{0A06}
\pdfglyphtounicode{aamatragurmukhi}{0A3E}
\pdfglyphtounicode{aarusquare}{3303}
\pdfglyphtounicode{aavowelsignbengali}{09BE}
\pdfglyphtounicode{aavowelsigndeva}{093E}
\pdfglyphtounicode{aavowelsigngujarati}{0ABE}
\pdfglyphtounicode{abbreviationmarkarmenian}{055F}
\pdfglyphtounicode{abbreviationsigndeva}{0970}
\pdfglyphtounicode{abengali}{0985}
\pdfglyphtounicode{abopomofo}{311A}
\pdfglyphtounicode{abreve}{0103}
\pdfglyphtounicode{abreveacute}{1EAF}
\pdfglyphtounicode{abrevecyrillic}{04D1}
\pdfglyphtounicode{abrevedotbelow}{1EB7}
\pdfglyphtounicode{abrevegrave}{1EB1}
\pdfglyphtounicode{abrevehookabove}{1EB3}
\pdfglyphtounicode{abrevetilde}{1EB5}
\pdfglyphtounicode{acaron}{01CE}
\pdfglyphtounicode{acircle}{24D0}
\pdfglyphtounicode{acircumflex}{00E2}
\pdfglyphtounicode{acircumflexacute}{1EA5}
\pdfglyphtounicode{acircumflexdotbelow}{1EAD}
\pdfglyphtounicode{acircumflexgrave}{1EA7}
\pdfglyphtounicode{acircumflexhookabove}{1EA9}
\pdfglyphtounicode{acircumflextilde}{1EAB}
\pdfglyphtounicode{acute}{00B4}
\pdfglyphtounicode{acutebelowcmb}{0317}
\pdfglyphtounicode{acutecmb}{0301}
\pdfglyphtounicode{acutecomb}{0301}
\pdfglyphtounicode{acutedeva}{0954}
\pdfglyphtounicode{acutelowmod}{02CF}
\pdfglyphtounicode{acutetonecmb}{0341}
\pdfglyphtounicode{acyrillic}{0430}
\pdfglyphtounicode{adblgrave}{0201}
\pdfglyphtounicode{addakgurmukhi}{0A71}
\pdfglyphtounicode{adeva}{0905}
\pdfglyphtounicode{adieresis}{00E4}
\pdfglyphtounicode{adieresiscyrillic}{04D3}
\pdfglyphtounicode{adieresismacron}{01DF}
\pdfglyphtounicode{adotbelow}{1EA1}
\pdfglyphtounicode{adotmacron}{01E1}
\pdfglyphtounicode{ae}{00E6}
\pdfglyphtounicode{aeacute}{01FD}
\pdfglyphtounicode{aekorean}{3150}
\pdfglyphtounicode{aemacron}{01E3}
\pdfglyphtounicode{afii00208}{2015}
\pdfglyphtounicode{afii08941}{20A4}
\pdfglyphtounicode{afii10017}{0410}
\pdfglyphtounicode{afii10018}{0411}
\pdfglyphtounicode{afii10019}{0412}
\pdfglyphtounicode{afii10020}{0413}
\pdfglyphtounicode{afii10021}{0414}
\pdfglyphtounicode{afii10022}{0415}
\pdfglyphtounicode{afii10023}{0401}
\pdfglyphtounicode{afii10024}{0416}
\pdfglyphtounicode{afii10025}{0417}
\pdfglyphtounicode{afii10026}{0418}
\pdfglyphtounicode{afii10027}{0419}
\pdfglyphtounicode{afii10028}{041A}
\pdfglyphtounicode{afii10029}{041B}
\pdfglyphtounicode{afii10030}{041C}
\pdfglyphtounicode{afii10031}{041D}
\pdfglyphtounicode{afii10032}{041E}
\pdfglyphtounicode{afii10033}{041F}
\pdfglyphtounicode{afii10034}{0420}
\pdfglyphtounicode{afii10035}{0421}
\pdfglyphtounicode{afii10036}{0422}
\pdfglyphtounicode{afii10037}{0423}
\pdfglyphtounicode{afii10038}{0424}
\pdfglyphtounicode{afii10039}{0425}
\pdfglyphtounicode{afii10040}{0426}
\pdfglyphtounicode{afii10041}{0427}
\pdfglyphtounicode{afii10042}{0428}
\pdfglyphtounicode{afii10043}{0429}
\pdfglyphtounicode{afii10044}{042A}
\pdfglyphtounicode{afii10045}{042B}
\pdfglyphtounicode{afii10046}{042C}
\pdfglyphtounicode{afii10047}{042D}
\pdfglyphtounicode{afii10048}{042E}
\pdfglyphtounicode{afii10049}{042F}
\pdfglyphtounicode{afii10050}{0490}
\pdfglyphtounicode{afii10051}{0402}
\pdfglyphtounicode{afii10052}{0403}
\pdfglyphtounicode{afii10053}{0404}
\pdfglyphtounicode{afii10054}{0405}
\pdfglyphtounicode{afii10055}{0406}
\pdfglyphtounicode{afii10056}{0407}
\pdfglyphtounicode{afii10057}{0408}
\pdfglyphtounicode{afii10058}{0409}
\pdfglyphtounicode{afii10059}{040A}
\pdfglyphtounicode{afii10060}{040B}
\pdfglyphtounicode{afii10061}{040C}
\pdfglyphtounicode{afii10062}{040E}
\pdfglyphtounicode{afii10063}{F6C4}
\pdfglyphtounicode{afii10064}{F6C5}
\pdfglyphtounicode{afii10065}{0430}
\pdfglyphtounicode{afii10066}{0431}
\pdfglyphtounicode{afii10067}{0432}
\pdfglyphtounicode{afii10068}{0433}
\pdfglyphtounicode{afii10069}{0434}
\pdfglyphtounicode{afii10070}{0435}
\pdfglyphtounicode{afii10071}{0451}
\pdfglyphtounicode{afii10072}{0436}
\pdfglyphtounicode{afii10073}{0437}
\pdfglyphtounicode{afii10074}{0438}
\pdfglyphtounicode{afii10075}{0439}
\pdfglyphtounicode{afii10076}{043A}
\pdfglyphtounicode{afii10077}{043B}
\pdfglyphtounicode{afii10078}{043C}
\pdfglyphtounicode{afii10079}{043D}
\pdfglyphtounicode{afii10080}{043E}
\pdfglyphtounicode{afii10081}{043F}
\pdfglyphtounicode{afii10082}{0440}
\pdfglyphtounicode{afii10083}{0441}
\pdfglyphtounicode{afii10084}{0442}
\pdfglyphtounicode{afii10085}{0443}
\pdfglyphtounicode{afii10086}{0444}
\pdfglyphtounicode{afii10087}{0445}
\pdfglyphtounicode{afii10088}{0446}
\pdfglyphtounicode{afii10089}{0447}
\pdfglyphtounicode{afii10090}{0448}
\pdfglyphtounicode{afii10091}{0449}
\pdfglyphtounicode{afii10092}{044A}
\pdfglyphtounicode{afii10093}{044B}
\pdfglyphtounicode{afii10094}{044C}
\pdfglyphtounicode{afii10095}{044D}
\pdfglyphtounicode{afii10096}{044E}
\pdfglyphtounicode{afii10097}{044F}
\pdfglyphtounicode{afii10098}{0491}
\pdfglyphtounicode{afii10099}{0452}
\pdfglyphtounicode{afii10100}{0453}
\pdfglyphtounicode{afii10101}{0454}
\pdfglyphtounicode{afii10102}{0455}
\pdfglyphtounicode{afii10103}{0456}
\pdfglyphtounicode{afii10104}{0457}
\pdfglyphtounicode{afii10105}{0458}
\pdfglyphtounicode{afii10106}{0459}
\pdfglyphtounicode{afii10107}{045A}
\pdfglyphtounicode{afii10108}{045B}
\pdfglyphtounicode{afii10109}{045C}
\pdfglyphtounicode{afii10110}{045E}
\pdfglyphtounicode{afii10145}{040F}
\pdfglyphtounicode{afii10146}{0462}
\pdfglyphtounicode{afii10147}{0472}
\pdfglyphtounicode{afii10148}{0474}
\pdfglyphtounicode{afii10192}{F6C6}
\pdfglyphtounicode{afii10193}{045F}
\pdfglyphtounicode{afii10194}{0463}
\pdfglyphtounicode{afii10195}{0473}
\pdfglyphtounicode{afii10196}{0475}
\pdfglyphtounicode{afii10831}{F6C7}
\pdfglyphtounicode{afii10832}{F6C8}
\pdfglyphtounicode{afii10846}{04D9}
\pdfglyphtounicode{afii299}{200E}
\pdfglyphtounicode{afii300}{200F}
\pdfglyphtounicode{afii301}{200D}
\pdfglyphtounicode{afii57381}{066A}
\pdfglyphtounicode{afii57388}{060C}
\pdfglyphtounicode{afii57392}{0660}
\pdfglyphtounicode{afii57393}{0661}
\pdfglyphtounicode{afii57394}{0662}
\pdfglyphtounicode{afii57395}{0663}
\pdfglyphtounicode{afii57396}{0664}
\pdfglyphtounicode{afii57397}{0665}
\pdfglyphtounicode{afii57398}{0666}
\pdfglyphtounicode{afii57399}{0667}
\pdfglyphtounicode{afii57400}{0668}
\pdfglyphtounicode{afii57401}{0669}
\pdfglyphtounicode{afii57403}{061B}
\pdfglyphtounicode{afii57407}{061F}
\pdfglyphtounicode{afii57409}{0621}
\pdfglyphtounicode{afii57410}{0622}
\pdfglyphtounicode{afii57411}{0623}
\pdfglyphtounicode{afii57412}{0624}
\pdfglyphtounicode{afii57413}{0625}
\pdfglyphtounicode{afii57414}{0626}
\pdfglyphtounicode{afii57415}{0627}
\pdfglyphtounicode{afii57416}{0628}
\pdfglyphtounicode{afii57417}{0629}
\pdfglyphtounicode{afii57418}{062A}
\pdfglyphtounicode{afii57419}{062B}
\pdfglyphtounicode{afii57420}{062C}
\pdfglyphtounicode{afii57421}{062D}
\pdfglyphtounicode{afii57422}{062E}
\pdfglyphtounicode{afii57423}{062F}
\pdfglyphtounicode{afii57424}{0630}
\pdfglyphtounicode{afii57425}{0631}
\pdfglyphtounicode{afii57426}{0632}
\pdfglyphtounicode{afii57427}{0633}
\pdfglyphtounicode{afii57428}{0634}
\pdfglyphtounicode{afii57429}{0635}
\pdfglyphtounicode{afii57430}{0636}
\pdfglyphtounicode{afii57431}{0637}
\pdfglyphtounicode{afii57432}{0638}
\pdfglyphtounicode{afii57433}{0639}
\pdfglyphtounicode{afii57434}{063A}
\pdfglyphtounicode{afii57440}{0640}
\pdfglyphtounicode{afii57441}{0641}
\pdfglyphtounicode{afii57442}{0642}
\pdfglyphtounicode{afii57443}{0643}
\pdfglyphtounicode{afii57444}{0644}
\pdfglyphtounicode{afii57445}{0645}
\pdfglyphtounicode{afii57446}{0646}
\pdfglyphtounicode{afii57448}{0648}
\pdfglyphtounicode{afii57449}{0649}
\pdfglyphtounicode{afii57450}{064A}
\pdfglyphtounicode{afii57451}{064B}
\pdfglyphtounicode{afii57452}{064C}
\pdfglyphtounicode{afii57453}{064D}
\pdfglyphtounicode{afii57454}{064E}
\pdfglyphtounicode{afii57455}{064F}
\pdfglyphtounicode{afii57456}{0650}
\pdfglyphtounicode{afii57457}{0651}
\pdfglyphtounicode{afii57458}{0652}
\pdfglyphtounicode{afii57470}{0647}
\pdfglyphtounicode{afii57505}{06A4}
\pdfglyphtounicode{afii57506}{067E}
\pdfglyphtounicode{afii57507}{0686}
\pdfglyphtounicode{afii57508}{0698}
\pdfglyphtounicode{afii57509}{06AF}
\pdfglyphtounicode{afii57511}{0679}
\pdfglyphtounicode{afii57512}{0688}
\pdfglyphtounicode{afii57513}{0691}
\pdfglyphtounicode{afii57514}{06BA}
\pdfglyphtounicode{afii57519}{06D2}
\pdfglyphtounicode{afii57534}{06D5}
\pdfglyphtounicode{afii57636}{20AA}
\pdfglyphtounicode{afii57645}{05BE}
\pdfglyphtounicode{afii57658}{05C3}
\pdfglyphtounicode{afii57664}{05D0}
\pdfglyphtounicode{afii57665}{05D1}
\pdfglyphtounicode{afii57666}{05D2}
\pdfglyphtounicode{afii57667}{05D3}
\pdfglyphtounicode{afii57668}{05D4}
\pdfglyphtounicode{afii57669}{05D5}
\pdfglyphtounicode{afii57670}{05D6}
\pdfglyphtounicode{afii57671}{05D7}
\pdfglyphtounicode{afii57672}{05D8}
\pdfglyphtounicode{afii57673}{05D9}
\pdfglyphtounicode{afii57674}{05DA}
\pdfglyphtounicode{afii57675}{05DB}
\pdfglyphtounicode{afii57676}{05DC}
\pdfglyphtounicode{afii57677}{05DD}
\pdfglyphtounicode{afii57678}{05DE}
\pdfglyphtounicode{afii57679}{05DF}
\pdfglyphtounicode{afii57680}{05E0}
\pdfglyphtounicode{afii57681}{05E1}
\pdfglyphtounicode{afii57682}{05E2}
\pdfglyphtounicode{afii57683}{05E3}
\pdfglyphtounicode{afii57684}{05E4}
\pdfglyphtounicode{afii57685}{05E5}
\pdfglyphtounicode{afii57686}{05E6}
\pdfglyphtounicode{afii57687}{05E7}
\pdfglyphtounicode{afii57688}{05E8}
\pdfglyphtounicode{afii57689}{05E9}
\pdfglyphtounicode{afii57690}{05EA}
\pdfglyphtounicode{afii57694}{FB2A}
\pdfglyphtounicode{afii57695}{FB2B}
\pdfglyphtounicode{afii57700}{FB4B}
\pdfglyphtounicode{afii57705}{FB1F}
\pdfglyphtounicode{afii57716}{05F0}
\pdfglyphtounicode{afii57717}{05F1}
\pdfglyphtounicode{afii57718}{05F2}
\pdfglyphtounicode{afii57723}{FB35}
\pdfglyphtounicode{afii57793}{05B4}
\pdfglyphtounicode{afii57794}{05B5}
\pdfglyphtounicode{afii57795}{05B6}
\pdfglyphtounicode{afii57796}{05BB}
\pdfglyphtounicode{afii57797}{05B8}
\pdfglyphtounicode{afii57798}{05B7}
\pdfglyphtounicode{afii57799}{05B0}
\pdfglyphtounicode{afii57800}{05B2}
\pdfglyphtounicode{afii57801}{05B1}
\pdfglyphtounicode{afii57802}{05B3}
\pdfglyphtounicode{afii57803}{05C2}
\pdfglyphtounicode{afii57804}{05C1}
\pdfglyphtounicode{afii57806}{05B9}
\pdfglyphtounicode{afii57807}{05BC}
\pdfglyphtounicode{afii57839}{05BD}
\pdfglyphtounicode{afii57841}{05BF}
\pdfglyphtounicode{afii57842}{05C0}
\pdfglyphtounicode{afii57929}{02BC}
\pdfglyphtounicode{afii61248}{2105}
\pdfglyphtounicode{afii61289}{2113}
\pdfglyphtounicode{afii61352}{2116}
\pdfglyphtounicode{afii61573}{202C}
\pdfglyphtounicode{afii61574}{202D}
\pdfglyphtounicode{afii61575}{202E}
\pdfglyphtounicode{afii61664}{200C}
\pdfglyphtounicode{afii63167}{066D}
\pdfglyphtounicode{afii64937}{02BD}
\pdfglyphtounicode{agrave}{00E0}
\pdfglyphtounicode{agujarati}{0A85}
\pdfglyphtounicode{agurmukhi}{0A05}
\pdfglyphtounicode{ahiragana}{3042}
\pdfglyphtounicode{ahookabove}{1EA3}
\pdfglyphtounicode{aibengali}{0990}
\pdfglyphtounicode{aibopomofo}{311E}
\pdfglyphtounicode{aideva}{0910}
\pdfglyphtounicode{aiecyrillic}{04D5}
\pdfglyphtounicode{aigujarati}{0A90}
\pdfglyphtounicode{aigurmukhi}{0A10}
\pdfglyphtounicode{aimatragurmukhi}{0A48}
\pdfglyphtounicode{ainarabic}{0639}
\pdfglyphtounicode{ainfinalarabic}{FECA}
\pdfglyphtounicode{aininitialarabic}{FECB}
\pdfglyphtounicode{ainmedialarabic}{FECC}
\pdfglyphtounicode{ainvertedbreve}{0203}
\pdfglyphtounicode{aivowelsignbengali}{09C8}
\pdfglyphtounicode{aivowelsigndeva}{0948}
\pdfglyphtounicode{aivowelsigngujarati}{0AC8}
\pdfglyphtounicode{akatakana}{30A2}
\pdfglyphtounicode{akatakanahalfwidth}{FF71}
\pdfglyphtounicode{akorean}{314F}
\pdfglyphtounicode{alef}{05D0}
\pdfglyphtounicode{alefarabic}{0627}
\pdfglyphtounicode{alefdageshhebrew}{FB30}
\pdfglyphtounicode{aleffinalarabic}{FE8E}
\pdfglyphtounicode{alefhamzaabovearabic}{0623}
\pdfglyphtounicode{alefhamzaabovefinalarabic}{FE84}
\pdfglyphtounicode{alefhamzabelowarabic}{0625}
\pdfglyphtounicode{alefhamzabelowfinalarabic}{FE88}
\pdfglyphtounicode{alefhebrew}{05D0}
\pdfglyphtounicode{aleflamedhebrew}{FB4F}
\pdfglyphtounicode{alefmaddaabovearabic}{0622}
\pdfglyphtounicode{alefmaddaabovefinalarabic}{FE82}
\pdfglyphtounicode{alefmaksuraarabic}{0649}
\pdfglyphtounicode{alefmaksurafinalarabic}{FEF0}
\pdfglyphtounicode{alefmaksurainitialarabic}{FEF3}
\pdfglyphtounicode{alefmaksuramedialarabic}{FEF4}
\pdfglyphtounicode{alefpatahhebrew}{FB2E}
\pdfglyphtounicode{alefqamatshebrew}{FB2F}
\pdfglyphtounicode{aleph}{2135}
\pdfglyphtounicode{allequal}{224C}
\pdfglyphtounicode{alpha}{03B1}
\pdfglyphtounicode{alphatonos}{03AC}
\pdfglyphtounicode{amacron}{0101}
\pdfglyphtounicode{amonospace}{FF41}
\pdfglyphtounicode{ampersand}{0026}
\pdfglyphtounicode{ampersandmonospace}{FF06}
\pdfglyphtounicode{ampersandsmall}{0026}
\pdfglyphtounicode{amsquare}{33C2}
\pdfglyphtounicode{anbopomofo}{3122}
\pdfglyphtounicode{angbopomofo}{3124}
\pdfglyphtounicode{angbracketleft}{27E8}
\pdfglyphtounicode{angbracketright}{27E9}
\pdfglyphtounicode{angkhankhuthai}{0E5A}
\pdfglyphtounicode{angle}{2220}
\pdfglyphtounicode{anglebracketleft}{3008}
\pdfglyphtounicode{anglebracketleftvertical}{FE3F}
\pdfglyphtounicode{anglebracketright}{3009}
\pdfglyphtounicode{anglebracketrightvertical}{FE40}
\pdfglyphtounicode{angleleft}{2329}
\pdfglyphtounicode{angleright}{232A}
\pdfglyphtounicode{angstrom}{212B}
\pdfglyphtounicode{anoteleia}{0387}
\pdfglyphtounicode{anticlockwise}{27F2}
\pdfglyphtounicode{anudattadeva}{0952}
\pdfglyphtounicode{anusvarabengali}{0982}
\pdfglyphtounicode{anusvaradeva}{0902}
\pdfglyphtounicode{anusvaragujarati}{0A82}
\pdfglyphtounicode{aogonek}{0105}
\pdfglyphtounicode{apaatosquare}{3300}
\pdfglyphtounicode{aparen}{249C}
\pdfglyphtounicode{apostrophearmenian}{055A}
\pdfglyphtounicode{apostrophemod}{02BC}
\pdfglyphtounicode{apple}{F8FF}
\pdfglyphtounicode{approaches}{2250}
\pdfglyphtounicode{approxequal}{2248}
\pdfglyphtounicode{approxequalorimage}{2252}
\pdfglyphtounicode{approximatelyequal}{2245}
\pdfglyphtounicode{approxorequal}{224A}
\pdfglyphtounicode{araeaekorean}{318E}
\pdfglyphtounicode{araeakorean}{318D}
\pdfglyphtounicode{arc}{2312}
\pdfglyphtounicode{archleftdown}{21B6}
\pdfglyphtounicode{archrightdown}{21B7}
\pdfglyphtounicode{arighthalfring}{1E9A}
\pdfglyphtounicode{aring}{00E5}
\pdfglyphtounicode{aringacute}{01FB}
\pdfglyphtounicode{aringbelow}{1E01}
\pdfglyphtounicode{arrowboth}{2194}
\pdfglyphtounicode{arrowbothv}{2195}
\pdfglyphtounicode{arrowdashdown}{21E3}
\pdfglyphtounicode{arrowdashleft}{21E0}
\pdfglyphtounicode{arrowdashright}{21E2}
\pdfglyphtounicode{arrowdashup}{21E1}
\pdfglyphtounicode{arrowdblboth}{21D4}
\pdfglyphtounicode{arrowdblbothv}{21D5}
\pdfglyphtounicode{arrowdbldown}{21D3}
\pdfglyphtounicode{arrowdblleft}{21D0}
\pdfglyphtounicode{arrowdblright}{21D2}
\pdfglyphtounicode{arrowdblup}{21D1}
\pdfglyphtounicode{arrowdown}{2193}
\pdfglyphtounicode{arrowdownleft}{2199}
\pdfglyphtounicode{arrowdownright}{2198}
\pdfglyphtounicode{arrowdownwhite}{21E9}
\pdfglyphtounicode{arrowheaddownmod}{02C5}
\pdfglyphtounicode{arrowheadleftmod}{02C2}
\pdfglyphtounicode{arrowheadrightmod}{02C3}
\pdfglyphtounicode{arrowheadupmod}{02C4}
\pdfglyphtounicode{arrowhorizex}{F8E7}
\pdfglyphtounicode{arrowleft}{2190}
\pdfglyphtounicode{arrowleftbothalf}{21BD}
\pdfglyphtounicode{arrowleftdbl}{21D0}
\pdfglyphtounicode{arrowleftdblstroke}{21CD}
\pdfglyphtounicode{arrowleftoverright}{21C6}
\pdfglyphtounicode{arrowlefttophalf}{21BC}
\pdfglyphtounicode{arrowleftwhite}{21E6}
\pdfglyphtounicode{arrownortheast}{2197}
\pdfglyphtounicode{arrownorthwest}{2196}
\pdfglyphtounicode{arrowparrleftright}{21C6}
\pdfglyphtounicode{arrowparrrightleft}{21C4}
\pdfglyphtounicode{arrowright}{2192}
\pdfglyphtounicode{arrowrightbothalf}{21C1}
\pdfglyphtounicode{arrowrightdblstroke}{21CF}
\pdfglyphtounicode{arrowrightheavy}{279E}
\pdfglyphtounicode{arrowrightoverleft}{21C4}
\pdfglyphtounicode{arrowrighttophalf}{21C0}
\pdfglyphtounicode{arrowrightwhite}{21E8}
\pdfglyphtounicode{arrowsoutheast}{2198}
\pdfglyphtounicode{arrowsouthwest}{2199}
\pdfglyphtounicode{arrowtableft}{21E4}
\pdfglyphtounicode{arrowtabright}{21E5}
\pdfglyphtounicode{arrowtailleft}{21A2}
\pdfglyphtounicode{arrowtailright}{21A3}
\pdfglyphtounicode{arrowtripleleft}{21DA}
\pdfglyphtounicode{arrowtripleright}{21DB}
\pdfglyphtounicode{arrowup}{2191}
\pdfglyphtounicode{arrowupdn}{2195}
\pdfglyphtounicode{arrowupdnbse}{21A8}
\pdfglyphtounicode{arrowupdownbase}{21A8}
\pdfglyphtounicode{arrowupleft}{2196}
\pdfglyphtounicode{arrowupleftofdown}{21C5}
\pdfglyphtounicode{arrowupright}{2197}
\pdfglyphtounicode{arrowupwhite}{21E7}
\pdfglyphtounicode{arrowvertex}{F8E6}
\pdfglyphtounicode{asciicircum}{005E}
\pdfglyphtounicode{asciicircummonospace}{FF3E}
\pdfglyphtounicode{asciitilde}{007E}
\pdfglyphtounicode{asciitildemonospace}{FF5E}
\pdfglyphtounicode{ascript}{0251}
\pdfglyphtounicode{ascriptturned}{0252}
\pdfglyphtounicode{asmallhiragana}{3041}
\pdfglyphtounicode{asmallkatakana}{30A1}
\pdfglyphtounicode{asmallkatakanahalfwidth}{FF67}
\pdfglyphtounicode{asterisk}{002A}
\pdfglyphtounicode{asteriskaltonearabic}{066D}
\pdfglyphtounicode{asteriskarabic}{066D}
\pdfglyphtounicode{asteriskcentered}{2217}
\pdfglyphtounicode{asteriskmath}{2217}
\pdfglyphtounicode{asteriskmonospace}{FF0A}
\pdfglyphtounicode{asterisksmall}{FE61}
\pdfglyphtounicode{asterism}{2042}
\pdfglyphtounicode{asuperior}{0061}
\pdfglyphtounicode{asymptoticallyequal}{2243}
\pdfglyphtounicode{at}{0040}
\pdfglyphtounicode{atilde}{00E3}
\pdfglyphtounicode{atmonospace}{FF20}
\pdfglyphtounicode{atsmall}{FE6B}
\pdfglyphtounicode{aturned}{0250}
\pdfglyphtounicode{aubengali}{0994}
\pdfglyphtounicode{aubopomofo}{3120}
\pdfglyphtounicode{audeva}{0914}
\pdfglyphtounicode{augujarati}{0A94}
\pdfglyphtounicode{augurmukhi}{0A14}
\pdfglyphtounicode{aulengthmarkbengali}{09D7}
\pdfglyphtounicode{aumatragurmukhi}{0A4C}
\pdfglyphtounicode{auvowelsignbengali}{09CC}
\pdfglyphtounicode{auvowelsigndeva}{094C}
\pdfglyphtounicode{auvowelsigngujarati}{0ACC}
\pdfglyphtounicode{avagrahadeva}{093D}
\pdfglyphtounicode{aybarmenian}{0561}
\pdfglyphtounicode{ayin}{05E2}
\pdfglyphtounicode{ayinaltonehebrew}{FB20}
\pdfglyphtounicode{ayinhebrew}{05E2}
\pdfglyphtounicode{b}{0062}
\pdfglyphtounicode{babengali}{09AC}
\pdfglyphtounicode{backslash}{005C}
\pdfglyphtounicode{backslashmonospace}{FF3C}
\pdfglyphtounicode{badeva}{092C}
\pdfglyphtounicode{bagujarati}{0AAC}
\pdfglyphtounicode{bagurmukhi}{0A2C}
\pdfglyphtounicode{bahiragana}{3070}
\pdfglyphtounicode{bahtthai}{0E3F}
\pdfglyphtounicode{bakatakana}{30D0}
\pdfglyphtounicode{bar}{007C}
\pdfglyphtounicode{bardbl}{2225}
\pdfglyphtounicode{barmonospace}{FF5C}
\pdfglyphtounicode{bbopomofo}{3105}
\pdfglyphtounicode{bcircle}{24D1}
\pdfglyphtounicode{bdotaccent}{1E03}
\pdfglyphtounicode{bdotbelow}{1E05}
\pdfglyphtounicode{beamedsixteenthnotes}{266C}
\pdfglyphtounicode{because}{2235}
\pdfglyphtounicode{becyrillic}{0431}
\pdfglyphtounicode{beharabic}{0628}
\pdfglyphtounicode{behfinalarabic}{FE90}
\pdfglyphtounicode{behinitialarabic}{FE91}
\pdfglyphtounicode{behiragana}{3079}
\pdfglyphtounicode{behmedialarabic}{FE92}
\pdfglyphtounicode{behmeeminitialarabic}{FC9F}
\pdfglyphtounicode{behmeemisolatedarabic}{FC08}
\pdfglyphtounicode{behnoonfinalarabic}{FC6D}
\pdfglyphtounicode{bekatakana}{30D9}
\pdfglyphtounicode{benarmenian}{0562}
\pdfglyphtounicode{bet}{05D1}
\pdfglyphtounicode{beta}{03B2}
\pdfglyphtounicode{betasymbolgreek}{03D0}
\pdfglyphtounicode{betdagesh}{FB31}
\pdfglyphtounicode{betdageshhebrew}{FB31}
\pdfglyphtounicode{beth}{2136}
\pdfglyphtounicode{bethebrew}{05D1}
\pdfglyphtounicode{betrafehebrew}{FB4C}
\pdfglyphtounicode{between}{226C}
\pdfglyphtounicode{bhabengali}{09AD}
\pdfglyphtounicode{bhadeva}{092D}
\pdfglyphtounicode{bhagujarati}{0AAD}
\pdfglyphtounicode{bhagurmukhi}{0A2D}
\pdfglyphtounicode{bhook}{0253}
\pdfglyphtounicode{bihiragana}{3073}
\pdfglyphtounicode{bikatakana}{30D3}
\pdfglyphtounicode{bilabialclick}{0298}
\pdfglyphtounicode{bindigurmukhi}{0A02}
\pdfglyphtounicode{birusquare}{3331}
\pdfglyphtounicode{blackcircle}{25CF}
\pdfglyphtounicode{blackdiamond}{25C6}
\pdfglyphtounicode{blackdownpointingtriangle}{25BC}
\pdfglyphtounicode{blackleftpointingpointer}{25C4}
\pdfglyphtounicode{blackleftpointingtriangle}{25C0}
\pdfglyphtounicode{blacklenticularbracketleft}{3010}
\pdfglyphtounicode{blacklenticularbracketleftvertical}{FE3B}
\pdfglyphtounicode{blacklenticularbracketright}{3011}
\pdfglyphtounicode{blacklenticularbracketrightvertical}{FE3C}
\pdfglyphtounicode{blacklowerlefttriangle}{25E3}
\pdfglyphtounicode{blacklowerrighttriangle}{25E2}
\pdfglyphtounicode{blackrectangle}{25AC}
\pdfglyphtounicode{blackrightpointingpointer}{25BA}
\pdfglyphtounicode{blackrightpointingtriangle}{25B6}
\pdfglyphtounicode{blacksmallsquare}{25AA}
\pdfglyphtounicode{blacksmilingface}{263B}
\pdfglyphtounicode{blacksquare}{25A0}
\pdfglyphtounicode{blackstar}{2605}
\pdfglyphtounicode{blackupperlefttriangle}{25E4}
\pdfglyphtounicode{blackupperrighttriangle}{25E5}
\pdfglyphtounicode{blackuppointingsmalltriangle}{25B4}
\pdfglyphtounicode{blackuppointingtriangle}{25B2}
\pdfglyphtounicode{blank}{2423}
\pdfglyphtounicode{blinebelow}{1E07}
\pdfglyphtounicode{block}{2588}
\pdfglyphtounicode{bmonospace}{FF42}
\pdfglyphtounicode{bobaimaithai}{0E1A}
\pdfglyphtounicode{bohiragana}{307C}
\pdfglyphtounicode{bokatakana}{30DC}
\pdfglyphtounicode{bparen}{249D}
\pdfglyphtounicode{bqsquare}{33C3}
\pdfglyphtounicode{braceex}{F8F4}
\pdfglyphtounicode{braceleft}{007B}
\pdfglyphtounicode{braceleftbt}{F8F3}
\pdfglyphtounicode{braceleftmid}{F8F2}
\pdfglyphtounicode{braceleftmonospace}{FF5B}
\pdfglyphtounicode{braceleftsmall}{FE5B}
\pdfglyphtounicode{bracelefttp}{F8F1}
\pdfglyphtounicode{braceleftvertical}{FE37}
\pdfglyphtounicode{braceright}{007D}
\pdfglyphtounicode{bracerightbt}{F8FE}
\pdfglyphtounicode{bracerightmid}{F8FD}
\pdfglyphtounicode{bracerightmonospace}{FF5D}
\pdfglyphtounicode{bracerightsmall}{FE5C}
\pdfglyphtounicode{bracerighttp}{F8FC}
\pdfglyphtounicode{bracerightvertical}{FE38}
\pdfglyphtounicode{bracketleft}{005B}
\pdfglyphtounicode{bracketleftbt}{F8F0}
\pdfglyphtounicode{bracketleftex}{F8EF}
\pdfglyphtounicode{bracketleftmonospace}{FF3B}
\pdfglyphtounicode{bracketlefttp}{F8EE}
\pdfglyphtounicode{bracketright}{005D}
\pdfglyphtounicode{bracketrightbt}{F8FB}
\pdfglyphtounicode{bracketrightex}{F8FA}
\pdfglyphtounicode{bracketrightmonospace}{FF3D}
\pdfglyphtounicode{bracketrighttp}{F8F9}
\pdfglyphtounicode{breve}{02D8}
\pdfglyphtounicode{brevebelowcmb}{032E}
\pdfglyphtounicode{brevecmb}{0306}
\pdfglyphtounicode{breveinvertedbelowcmb}{032F}
\pdfglyphtounicode{breveinvertedcmb}{0311}
\pdfglyphtounicode{breveinverteddoublecmb}{0361}
\pdfglyphtounicode{bridgebelowcmb}{032A}
\pdfglyphtounicode{bridgeinvertedbelowcmb}{033A}
\pdfglyphtounicode{brokenbar}{00A6}
\pdfglyphtounicode{bstroke}{0180}
\pdfglyphtounicode{bsuperior}{0062}
\pdfglyphtounicode{btopbar}{0183}
\pdfglyphtounicode{buhiragana}{3076}
\pdfglyphtounicode{bukatakana}{30D6}
\pdfglyphtounicode{bullet}{2022}
\pdfglyphtounicode{bulletinverse}{25D8}
\pdfglyphtounicode{bulletoperator}{2219}
\pdfglyphtounicode{bullseye}{25CE}
\pdfglyphtounicode{c}{0063}
\pdfglyphtounicode{caarmenian}{056E}
\pdfglyphtounicode{cabengali}{099A}
\pdfglyphtounicode{cacute}{0107}
\pdfglyphtounicode{cadeva}{091A}
\pdfglyphtounicode{cagujarati}{0A9A}
\pdfglyphtounicode{cagurmukhi}{0A1A}
\pdfglyphtounicode{calsquare}{3388}
\pdfglyphtounicode{candrabindubengali}{0981}
\pdfglyphtounicode{candrabinducmb}{0310}
\pdfglyphtounicode{candrabindudeva}{0901}
\pdfglyphtounicode{candrabindugujarati}{0A81}
\pdfglyphtounicode{capslock}{21EA}
\pdfglyphtounicode{careof}{2105}
\pdfglyphtounicode{caron}{02C7}
\pdfglyphtounicode{caronbelowcmb}{032C}
\pdfglyphtounicode{caroncmb}{030C}
\pdfglyphtounicode{carriagereturn}{21B5}
\pdfglyphtounicode{cbopomofo}{3118}
\pdfglyphtounicode{ccaron}{010D}
\pdfglyphtounicode{ccedilla}{00E7}
\pdfglyphtounicode{ccedillaacute}{1E09}
\pdfglyphtounicode{ccircle}{24D2}
\pdfglyphtounicode{ccircumflex}{0109}
\pdfglyphtounicode{ccurl}{0255}
\pdfglyphtounicode{cdot}{010B}
\pdfglyphtounicode{cdotaccent}{010B}
\pdfglyphtounicode{cdsquare}{33C5}
\pdfglyphtounicode{cedilla}{00B8}
\pdfglyphtounicode{cedillacmb}{0327}
\pdfglyphtounicode{ceilingleft}{2308}
\pdfglyphtounicode{ceilingright}{2309}
\pdfglyphtounicode{cent}{00A2}
\pdfglyphtounicode{centigrade}{2103}
\pdfglyphtounicode{centinferior}{00A2}
\pdfglyphtounicode{centmonospace}{FFE0}
\pdfglyphtounicode{centoldstyle}{00A2}
\pdfglyphtounicode{centsuperior}{00A2}
\pdfglyphtounicode{chaarmenian}{0579}
\pdfglyphtounicode{chabengali}{099B}
\pdfglyphtounicode{chadeva}{091B}
\pdfglyphtounicode{chagujarati}{0A9B}
\pdfglyphtounicode{chagurmukhi}{0A1B}
\pdfglyphtounicode{chbopomofo}{3114}
\pdfglyphtounicode{cheabkhasiancyrillic}{04BD}
\pdfglyphtounicode{check}{2713}
\pdfglyphtounicode{checkmark}{2713}
\pdfglyphtounicode{checyrillic}{0447}
\pdfglyphtounicode{chedescenderabkhasiancyrillic}{04BF}
\pdfglyphtounicode{chedescendercyrillic}{04B7}
\pdfglyphtounicode{chedieresiscyrillic}{04F5}
\pdfglyphtounicode{cheharmenian}{0573}
\pdfglyphtounicode{chekhakassiancyrillic}{04CC}
\pdfglyphtounicode{cheverticalstrokecyrillic}{04B9}
\pdfglyphtounicode{chi}{03C7}
\pdfglyphtounicode{chieuchacirclekorean}{3277}
\pdfglyphtounicode{chieuchaparenkorean}{3217}
\pdfglyphtounicode{chieuchcirclekorean}{3269}
\pdfglyphtounicode{chieuchkorean}{314A}
\pdfglyphtounicode{chieuchparenkorean}{3209}
\pdfglyphtounicode{chochangthai}{0E0A}
\pdfglyphtounicode{chochanthai}{0E08}
\pdfglyphtounicode{chochingthai}{0E09}
\pdfglyphtounicode{chochoethai}{0E0C}
\pdfglyphtounicode{chook}{0188}
\pdfglyphtounicode{cieucacirclekorean}{3276}
\pdfglyphtounicode{cieucaparenkorean}{3216}
\pdfglyphtounicode{cieuccirclekorean}{3268}
\pdfglyphtounicode{cieuckorean}{3148}
\pdfglyphtounicode{cieucparenkorean}{3208}
\pdfglyphtounicode{cieucuparenkorean}{321C}
\pdfglyphtounicode{circle}{25CB}
\pdfglyphtounicode{circleR}{00AE}
\pdfglyphtounicode{circleS}{24C8}
\pdfglyphtounicode{circleasterisk}{229B}
\pdfglyphtounicode{circlecopyrt}{20DD}
\pdfglyphtounicode{circledivide}{2298}
\pdfglyphtounicode{circledot}{2299}
\pdfglyphtounicode{circleequal}{229C}
\pdfglyphtounicode{circleminus}{2296}
\pdfglyphtounicode{circlemultiply}{2297}
\pdfglyphtounicode{circleot}{2299}
\pdfglyphtounicode{circleplus}{2295}
\pdfglyphtounicode{circlepostalmark}{3036}
\pdfglyphtounicode{circlering}{229A}
\pdfglyphtounicode{circlewithlefthalfblack}{25D0}
\pdfglyphtounicode{circlewithrighthalfblack}{25D1}
\pdfglyphtounicode{circumflex}{02C6}
\pdfglyphtounicode{circumflexbelowcmb}{032D}
\pdfglyphtounicode{circumflexcmb}{0302}
\pdfglyphtounicode{clear}{2327}
\pdfglyphtounicode{clickalveolar}{01C2}
\pdfglyphtounicode{clickdental}{01C0}
\pdfglyphtounicode{clicklateral}{01C1}
\pdfglyphtounicode{clickretroflex}{01C3}
\pdfglyphtounicode{clockwise}{27F3}
\pdfglyphtounicode{club}{2663}
\pdfglyphtounicode{clubsuitblack}{2663}
\pdfglyphtounicode{clubsuitwhite}{2667}
\pdfglyphtounicode{cmcubedsquare}{33A4}
\pdfglyphtounicode{cmonospace}{FF43}
\pdfglyphtounicode{cmsquaredsquare}{33A0}
\pdfglyphtounicode{coarmenian}{0581}
\pdfglyphtounicode{colon}{003A}
\pdfglyphtounicode{colonmonetary}{20A1}
\pdfglyphtounicode{colonmonospace}{FF1A}
\pdfglyphtounicode{colonsign}{20A1}
\pdfglyphtounicode{colonsmall}{FE55}
\pdfglyphtounicode{colontriangularhalfmod}{02D1}
\pdfglyphtounicode{colontriangularmod}{02D0}
\pdfglyphtounicode{comma}{002C}
\pdfglyphtounicode{commaabovecmb}{0313}
\pdfglyphtounicode{commaaboverightcmb}{0315}
\pdfglyphtounicode{commaaccent}{F6C3}
\pdfglyphtounicode{commaarabic}{060C}
\pdfglyphtounicode{commaarmenian}{055D}
\pdfglyphtounicode{commainferior}{002C}
\pdfglyphtounicode{commamonospace}{FF0C}
\pdfglyphtounicode{commareversedabovecmb}{0314}
\pdfglyphtounicode{commareversedmod}{02BD}
\pdfglyphtounicode{commasmall}{FE50}
\pdfglyphtounicode{commasuperior}{002C}
\pdfglyphtounicode{commaturnedabovecmb}{0312}
\pdfglyphtounicode{commaturnedmod}{02BB}
\pdfglyphtounicode{compass}{263C}
\pdfglyphtounicode{complement}{2201}
\pdfglyphtounicode{compwordmark}{200C}
\pdfglyphtounicode{congruent}{2245}
\pdfglyphtounicode{contourintegral}{222E}
\pdfglyphtounicode{control}{2303}
\pdfglyphtounicode{controlACK}{0006}
\pdfglyphtounicode{controlBEL}{0007}
\pdfglyphtounicode{controlBS}{0008}
\pdfglyphtounicode{controlCAN}{0018}
\pdfglyphtounicode{controlCR}{000D}
\pdfglyphtounicode{controlDC1}{0011}
\pdfglyphtounicode{controlDC2}{0012}
\pdfglyphtounicode{controlDC3}{0013}
\pdfglyphtounicode{controlDC4}{0014}
\pdfglyphtounicode{controlDEL}{007F}
\pdfglyphtounicode{controlDLE}{0010}
\pdfglyphtounicode{controlEM}{0019}
\pdfglyphtounicode{controlENQ}{0005}
\pdfglyphtounicode{controlEOT}{0004}
\pdfglyphtounicode{controlESC}{001B}
\pdfglyphtounicode{controlETB}{0017}
\pdfglyphtounicode{controlETX}{0003}
\pdfglyphtounicode{controlFF}{000C}
\pdfglyphtounicode{controlFS}{001C}
\pdfglyphtounicode{controlGS}{001D}
\pdfglyphtounicode{controlHT}{0009}
\pdfglyphtounicode{controlLF}{000A}
\pdfglyphtounicode{controlNAK}{0015}
\pdfglyphtounicode{controlRS}{001E}
\pdfglyphtounicode{controlSI}{000F}
\pdfglyphtounicode{controlSO}{000E}
\pdfglyphtounicode{controlSOT}{0002}
\pdfglyphtounicode{controlSTX}{0001}
\pdfglyphtounicode{controlSUB}{001A}
\pdfglyphtounicode{controlSYN}{0016}
\pdfglyphtounicode{controlUS}{001F}
\pdfglyphtounicode{controlVT}{000B}
\pdfglyphtounicode{coproduct}{2A3F}
\pdfglyphtounicode{copyright}{00A9}
\pdfglyphtounicode{copyrightsans}{00A9}
\pdfglyphtounicode{copyrightserif}{00A9}
\pdfglyphtounicode{cornerbracketleft}{300C}
\pdfglyphtounicode{cornerbracketlefthalfwidth}{FF62}
\pdfglyphtounicode{cornerbracketleftvertical}{FE41}
\pdfglyphtounicode{cornerbracketright}{300D}
\pdfglyphtounicode{cornerbracketrighthalfwidth}{FF63}
\pdfglyphtounicode{cornerbracketrightvertical}{FE42}
\pdfglyphtounicode{corporationsquare}{337F}
\pdfglyphtounicode{cosquare}{33C7}
\pdfglyphtounicode{coverkgsquare}{33C6}
\pdfglyphtounicode{cparen}{249E}
\pdfglyphtounicode{cruzeiro}{20A2}
\pdfglyphtounicode{cstretched}{0297}
\pdfglyphtounicode{ct}{0063 0074}
\pdfglyphtounicode{curlyand}{22CF}
\pdfglyphtounicode{curlyleft}{21AB}
\pdfglyphtounicode{curlyor}{22CE}
\pdfglyphtounicode{curlyright}{21AC}
\pdfglyphtounicode{currency}{00A4}
\pdfglyphtounicode{cwm}{200C}
\pdfglyphtounicode{cyrBreve}{02D8}
\pdfglyphtounicode{cyrFlex}{00A0 0311}
\pdfglyphtounicode{cyrbreve}{02D8}
\pdfglyphtounicode{cyrflex}{00A0 0311}
\pdfglyphtounicode{d}{0064}
\pdfglyphtounicode{daarmenian}{0564}
\pdfglyphtounicode{dabengali}{09A6}
\pdfglyphtounicode{dadarabic}{0636}
\pdfglyphtounicode{dadeva}{0926}
\pdfglyphtounicode{dadfinalarabic}{FEBE}
\pdfglyphtounicode{dadinitialarabic}{FEBF}
\pdfglyphtounicode{dadmedialarabic}{FEC0}
\pdfglyphtounicode{dagesh}{05BC}
\pdfglyphtounicode{dageshhebrew}{05BC}
\pdfglyphtounicode{dagger}{2020}
\pdfglyphtounicode{daggerdbl}{2021}
\pdfglyphtounicode{dagujarati}{0AA6}
\pdfglyphtounicode{dagurmukhi}{0A26}
\pdfglyphtounicode{dahiragana}{3060}
\pdfglyphtounicode{dakatakana}{30C0}
\pdfglyphtounicode{dalarabic}{062F}
\pdfglyphtounicode{dalet}{05D3}
\pdfglyphtounicode{daletdagesh}{FB33}
\pdfglyphtounicode{daletdageshhebrew}{FB33}
\pdfglyphtounicode{daleth}{2138}
\pdfglyphtounicode{dalethatafpatah}{05D3 05B2}
\pdfglyphtounicode{dalethatafpatahhebrew}{05D3 05B2}
\pdfglyphtounicode{dalethatafsegol}{05D3 05B1}
\pdfglyphtounicode{dalethatafsegolhebrew}{05D3 05B1}
\pdfglyphtounicode{dalethebrew}{05D3}
\pdfglyphtounicode{dalethiriq}{05D3 05B4}
\pdfglyphtounicode{dalethiriqhebrew}{05D3 05B4}
\pdfglyphtounicode{daletholam}{05D3 05B9}
\pdfglyphtounicode{daletholamhebrew}{05D3 05B9}
\pdfglyphtounicode{daletpatah}{05D3 05B7}
\pdfglyphtounicode{daletpatahhebrew}{05D3 05B7}
\pdfglyphtounicode{daletqamats}{05D3 05B8}
\pdfglyphtounicode{daletqamatshebrew}{05D3 05B8}
\pdfglyphtounicode{daletqubuts}{05D3 05BB}
\pdfglyphtounicode{daletqubutshebrew}{05D3 05BB}
\pdfglyphtounicode{daletsegol}{05D3 05B6}
\pdfglyphtounicode{daletsegolhebrew}{05D3 05B6}
\pdfglyphtounicode{daletsheva}{05D3 05B0}
\pdfglyphtounicode{daletshevahebrew}{05D3 05B0}
\pdfglyphtounicode{dalettsere}{05D3 05B5}
\pdfglyphtounicode{dalettserehebrew}{05D3 05B5}
\pdfglyphtounicode{dalfinalarabic}{FEAA}
\pdfglyphtounicode{dammaarabic}{064F}
\pdfglyphtounicode{dammalowarabic}{064F}
\pdfglyphtounicode{dammatanaltonearabic}{064C}
\pdfglyphtounicode{dammatanarabic}{064C}
\pdfglyphtounicode{danda}{0964}
\pdfglyphtounicode{dargahebrew}{05A7}
\pdfglyphtounicode{dargalefthebrew}{05A7}
\pdfglyphtounicode{dasiapneumatacyrilliccmb}{0485}
\pdfglyphtounicode{dbar}{0111}
\pdfglyphtounicode{dblGrave}{00A0 030F}
\pdfglyphtounicode{dblanglebracketleft}{300A}
\pdfglyphtounicode{dblanglebracketleftvertical}{FE3D}
\pdfglyphtounicode{dblanglebracketright}{300B}
\pdfglyphtounicode{dblanglebracketrightvertical}{FE3E}
\pdfglyphtounicode{dblarchinvertedbelowcmb}{032B}
\pdfglyphtounicode{dblarrowdwn}{21CA}
\pdfglyphtounicode{dblarrowheadleft}{219E}
\pdfglyphtounicode{dblarrowheadright}{21A0}
\pdfglyphtounicode{dblarrowleft}{21D4}
\pdfglyphtounicode{dblarrowright}{21D2}
\pdfglyphtounicode{dblarrowup}{21C8}
\pdfglyphtounicode{dblbracketleft}{27E6}
\pdfglyphtounicode{dblbracketright}{27E7}
\pdfglyphtounicode{dbldanda}{0965}
\pdfglyphtounicode{dblgrave}{00A0 030F}
\pdfglyphtounicode{dblgravecmb}{030F}
\pdfglyphtounicode{dblintegral}{222C}
\pdfglyphtounicode{dbllowline}{2017}
\pdfglyphtounicode{dbllowlinecmb}{0333}
\pdfglyphtounicode{dbloverlinecmb}{033F}
\pdfglyphtounicode{dblprimemod}{02BA}
\pdfglyphtounicode{dblverticalbar}{2016}
\pdfglyphtounicode{dblverticallineabovecmb}{030E}
\pdfglyphtounicode{dbopomofo}{3109}
\pdfglyphtounicode{dbsquare}{33C8}
\pdfglyphtounicode{dcaron}{010F}
\pdfglyphtounicode{dcedilla}{1E11}
\pdfglyphtounicode{dcircle}{24D3}
\pdfglyphtounicode{dcircumflexbelow}{1E13}
\pdfglyphtounicode{dcroat}{0111}
\pdfglyphtounicode{ddabengali}{09A1}
\pdfglyphtounicode{ddadeva}{0921}
\pdfglyphtounicode{ddagujarati}{0AA1}
\pdfglyphtounicode{ddagurmukhi}{0A21}
\pdfglyphtounicode{ddalarabic}{0688}
\pdfglyphtounicode{ddalfinalarabic}{FB89}
\pdfglyphtounicode{dddhadeva}{095C}
\pdfglyphtounicode{ddhabengali}{09A2}
\pdfglyphtounicode{ddhadeva}{0922}
\pdfglyphtounicode{ddhagujarati}{0AA2}
\pdfglyphtounicode{ddhagurmukhi}{0A22}
\pdfglyphtounicode{ddotaccent}{1E0B}
\pdfglyphtounicode{ddotbelow}{1E0D}
\pdfglyphtounicode{decimalseparatorarabic}{066B}
\pdfglyphtounicode{decimalseparatorpersian}{066B}
\pdfglyphtounicode{decyrillic}{0434}
\pdfglyphtounicode{defines}{225C}
\pdfglyphtounicode{degree}{00B0}
\pdfglyphtounicode{dehihebrew}{05AD}
\pdfglyphtounicode{dehiragana}{3067}
\pdfglyphtounicode{deicoptic}{03EF}
\pdfglyphtounicode{dekatakana}{30C7}
\pdfglyphtounicode{deleteleft}{232B}
\pdfglyphtounicode{deleteright}{2326}
\pdfglyphtounicode{delta}{03B4}
\pdfglyphtounicode{deltaturned}{018D}
\pdfglyphtounicode{denominatorminusonenumeratorbengali}{09F8}
\pdfglyphtounicode{dezh}{02A4}
\pdfglyphtounicode{dhabengali}{09A7}
\pdfglyphtounicode{dhadeva}{0927}
\pdfglyphtounicode{dhagujarati}{0AA7}
\pdfglyphtounicode{dhagurmukhi}{0A27}
\pdfglyphtounicode{dhook}{0257}
\pdfglyphtounicode{dialytikatonos}{0385}
\pdfglyphtounicode{dialytikatonoscmb}{0344}
\pdfglyphtounicode{diamond}{2666}
\pdfglyphtounicode{diamondmath}{22C4}
\pdfglyphtounicode{diamondsolid}{2666}
\pdfglyphtounicode{diamondsuitwhite}{2662}
\pdfglyphtounicode{dieresis}{00A8}
\pdfglyphtounicode{dieresisacute}{00A0 0308 0301}
\pdfglyphtounicode{dieresisbelowcmb}{0324}
\pdfglyphtounicode{dieresiscmb}{0308}
\pdfglyphtounicode{dieresisgrave}{00A0 0308 0300}
\pdfglyphtounicode{dieresistonos}{0385}
\pdfglyphtounicode{difference}{224F}
\pdfglyphtounicode{dihiragana}{3062}
\pdfglyphtounicode{dikatakana}{30C2}
\pdfglyphtounicode{dittomark}{3003}
\pdfglyphtounicode{divide}{00F7}
\pdfglyphtounicode{dividemultiply}{22C7}
\pdfglyphtounicode{divides}{2223}
\pdfglyphtounicode{divisionslash}{2215}
\pdfglyphtounicode{djecyrillic}{0452}
\pdfglyphtounicode{dkshade}{2593}
\pdfglyphtounicode{dlinebelow}{1E0F}
\pdfglyphtounicode{dlsquare}{3397}
\pdfglyphtounicode{dmacron}{0111}
\pdfglyphtounicode{dmonospace}{FF44}
\pdfglyphtounicode{dnblock}{2584}
\pdfglyphtounicode{dochadathai}{0E0E}
\pdfglyphtounicode{dodekthai}{0E14}
\pdfglyphtounicode{dohiragana}{3069}
\pdfglyphtounicode{dokatakana}{30C9}
\pdfglyphtounicode{dollar}{0024}
\pdfglyphtounicode{dollarinferior}{0024}
\pdfglyphtounicode{dollarmonospace}{FF04}
\pdfglyphtounicode{dollaroldstyle}{0024}
\pdfglyphtounicode{dollarsmall}{FE69}
\pdfglyphtounicode{dollarsuperior}{0024}
\pdfglyphtounicode{dong}{20AB}
\pdfglyphtounicode{dorusquare}{3326}
\pdfglyphtounicode{dotaccent}{02D9}
\pdfglyphtounicode{dotaccentcmb}{0307}
\pdfglyphtounicode{dotbelowcmb}{0323}
\pdfglyphtounicode{dotbelowcomb}{0323}
\pdfglyphtounicode{dotkatakana}{30FB}
\pdfglyphtounicode{dotlessi}{0131}
\pdfglyphtounicode{dotlessj}{0237}
\pdfglyphtounicode{dotlessjstrokehook}{0284}
\pdfglyphtounicode{dotmath}{22C5}
\pdfglyphtounicode{dotplus}{2214}
\pdfglyphtounicode{dottedcircle}{25CC}
\pdfglyphtounicode{doubleyodpatah}{FB1F}
\pdfglyphtounicode{doubleyodpatahhebrew}{FB1F}
\pdfglyphtounicode{downfall}{22CE}
\pdfglyphtounicode{downslope}{29F9}
\pdfglyphtounicode{downtackbelowcmb}{031E}
\pdfglyphtounicode{downtackmod}{02D5}
\pdfglyphtounicode{dparen}{249F}
\pdfglyphtounicode{dsuperior}{0064}
\pdfglyphtounicode{dtail}{0256}
\pdfglyphtounicode{dtopbar}{018C}
\pdfglyphtounicode{duhiragana}{3065}
\pdfglyphtounicode{dukatakana}{30C5}
\pdfglyphtounicode{dz}{01F3}
\pdfglyphtounicode{dzaltone}{02A3}
\pdfglyphtounicode{dzcaron}{01C6}
\pdfglyphtounicode{dzcurl}{02A5}
\pdfglyphtounicode{dzeabkhasiancyrillic}{04E1}
\pdfglyphtounicode{dzecyrillic}{0455}
\pdfglyphtounicode{dzhecyrillic}{045F}
\pdfglyphtounicode{e}{0065}
\pdfglyphtounicode{eacute}{00E9}
\pdfglyphtounicode{earth}{2641}
\pdfglyphtounicode{ebengali}{098F}
\pdfglyphtounicode{ebopomofo}{311C}
\pdfglyphtounicode{ebreve}{0115}
\pdfglyphtounicode{ecandradeva}{090D}
\pdfglyphtounicode{ecandragujarati}{0A8D}
\pdfglyphtounicode{ecandravowelsigndeva}{0945}
\pdfglyphtounicode{ecandravowelsigngujarati}{0AC5}
\pdfglyphtounicode{ecaron}{011B}
\pdfglyphtounicode{ecedillabreve}{1E1D}
\pdfglyphtounicode{echarmenian}{0565}
\pdfglyphtounicode{echyiwnarmenian}{0587}
\pdfglyphtounicode{ecircle}{24D4}
\pdfglyphtounicode{ecircumflex}{00EA}
\pdfglyphtounicode{ecircumflexacute}{1EBF}
\pdfglyphtounicode{ecircumflexbelow}{1E19}
\pdfglyphtounicode{ecircumflexdotbelow}{1EC7}
\pdfglyphtounicode{ecircumflexgrave}{1EC1}
\pdfglyphtounicode{ecircumflexhookabove}{1EC3}
\pdfglyphtounicode{ecircumflextilde}{1EC5}
\pdfglyphtounicode{ecyrillic}{0454}
\pdfglyphtounicode{edblgrave}{0205}
\pdfglyphtounicode{edeva}{090F}
\pdfglyphtounicode{edieresis}{00EB}
\pdfglyphtounicode{edot}{0117}
\pdfglyphtounicode{edotaccent}{0117}
\pdfglyphtounicode{edotbelow}{1EB9}
\pdfglyphtounicode{eegurmukhi}{0A0F}
\pdfglyphtounicode{eematragurmukhi}{0A47}
\pdfglyphtounicode{efcyrillic}{0444}
\pdfglyphtounicode{egrave}{00E8}
\pdfglyphtounicode{egujarati}{0A8F}
\pdfglyphtounicode{eharmenian}{0567}
\pdfglyphtounicode{ehbopomofo}{311D}
\pdfglyphtounicode{ehiragana}{3048}
\pdfglyphtounicode{ehookabove}{1EBB}
\pdfglyphtounicode{eibopomofo}{311F}
\pdfglyphtounicode{eight}{0038}
\pdfglyphtounicode{eightarabic}{0668}
\pdfglyphtounicode{eightbengali}{09EE}
\pdfglyphtounicode{eightcircle}{2467}
\pdfglyphtounicode{eightcircleinversesansserif}{2791}
\pdfglyphtounicode{eightdeva}{096E}
\pdfglyphtounicode{eighteencircle}{2471}
\pdfglyphtounicode{eighteenparen}{2485}
\pdfglyphtounicode{eighteenperiod}{2499}
\pdfglyphtounicode{eightgujarati}{0AEE}
\pdfglyphtounicode{eightgurmukhi}{0A6E}
\pdfglyphtounicode{eighthackarabic}{0668}
\pdfglyphtounicode{eighthangzhou}{3028}
\pdfglyphtounicode{eighthnotebeamed}{266B}
\pdfglyphtounicode{eightideographicparen}{3227}
\pdfglyphtounicode{eightinferior}{2088}
\pdfglyphtounicode{eightmonospace}{FF18}
\pdfglyphtounicode{eightoldstyle}{0038}
\pdfglyphtounicode{eightparen}{247B}
\pdfglyphtounicode{eightperiod}{248F}
\pdfglyphtounicode{eightpersian}{06F8}
\pdfglyphtounicode{eightroman}{2177}
\pdfglyphtounicode{eightsuperior}{2078}
\pdfglyphtounicode{eightthai}{0E58}
\pdfglyphtounicode{einvertedbreve}{0207}
\pdfglyphtounicode{eiotifiedcyrillic}{0465}
\pdfglyphtounicode{ekatakana}{30A8}
\pdfglyphtounicode{ekatakanahalfwidth}{FF74}
\pdfglyphtounicode{ekonkargurmukhi}{0A74}
\pdfglyphtounicode{ekorean}{3154}
\pdfglyphtounicode{elcyrillic}{043B}
\pdfglyphtounicode{element}{2208}
\pdfglyphtounicode{elevencircle}{246A}
\pdfglyphtounicode{elevenparen}{247E}
\pdfglyphtounicode{elevenperiod}{2492}
\pdfglyphtounicode{elevenroman}{217A}
\pdfglyphtounicode{ellipsis}{2026}
\pdfglyphtounicode{ellipsisvertical}{22EE}
\pdfglyphtounicode{emacron}{0113}
\pdfglyphtounicode{emacronacute}{1E17}
\pdfglyphtounicode{emacrongrave}{1E15}
\pdfglyphtounicode{emcyrillic}{043C}
\pdfglyphtounicode{emdash}{2014}
\pdfglyphtounicode{emdashvertical}{FE31}
\pdfglyphtounicode{emonospace}{FF45}
\pdfglyphtounicode{emphasismarkarmenian}{055B}
\pdfglyphtounicode{emptyset}{2205}
\pdfglyphtounicode{enbopomofo}{3123}
\pdfglyphtounicode{encyrillic}{043D}
\pdfglyphtounicode{endash}{2013}
\pdfglyphtounicode{endashvertical}{FE32}
\pdfglyphtounicode{endescendercyrillic}{04A3}
\pdfglyphtounicode{eng}{014B}
\pdfglyphtounicode{engbopomofo}{3125}
\pdfglyphtounicode{enghecyrillic}{04A5}
\pdfglyphtounicode{enhookcyrillic}{04C8}
\pdfglyphtounicode{enspace}{2002}
\pdfglyphtounicode{eogonek}{0119}
\pdfglyphtounicode{eokorean}{3153}
\pdfglyphtounicode{eopen}{025B}
\pdfglyphtounicode{eopenclosed}{029A}
\pdfglyphtounicode{eopenreversed}{025C}
\pdfglyphtounicode{eopenreversedclosed}{025E}
\pdfglyphtounicode{eopenreversedhook}{025D}
\pdfglyphtounicode{eparen}{24A0}
\pdfglyphtounicode{epsilon}{03B5}
\pdfglyphtounicode{epsilon1}{03F5}
\pdfglyphtounicode{epsiloninv}{03F6}
\pdfglyphtounicode{epsilontonos}{03AD}
\pdfglyphtounicode{equal}{003D}
\pdfglyphtounicode{equaldotleftright}{2252}
\pdfglyphtounicode{equaldotrightleft}{2253}
\pdfglyphtounicode{equalmonospace}{FF1D}
\pdfglyphtounicode{equalorfollows}{22DF}
\pdfglyphtounicode{equalorgreater}{2A96}
\pdfglyphtounicode{equalorless}{2A95}
\pdfglyphtounicode{equalorprecedes}{22DE}
\pdfglyphtounicode{equalorsimilar}{2242}
\pdfglyphtounicode{equalsdots}{2251}
\pdfglyphtounicode{equalsmall}{FE66}
\pdfglyphtounicode{equalsuperior}{207C}
\pdfglyphtounicode{equivalence}{2261}
\pdfglyphtounicode{equivasymptotic}{224D}
\pdfglyphtounicode{erbopomofo}{3126}
\pdfglyphtounicode{ercyrillic}{0440}
\pdfglyphtounicode{ereversed}{0258}
\pdfglyphtounicode{ereversedcyrillic}{044D}
\pdfglyphtounicode{escyrillic}{0441}
\pdfglyphtounicode{esdescendercyrillic}{04AB}
\pdfglyphtounicode{esh}{0283}
\pdfglyphtounicode{eshcurl}{0286}
\pdfglyphtounicode{eshortdeva}{090E}
\pdfglyphtounicode{eshortvowelsigndeva}{0946}
\pdfglyphtounicode{eshreversedloop}{01AA}
\pdfglyphtounicode{eshsquatreversed}{0285}
\pdfglyphtounicode{esmallhiragana}{3047}
\pdfglyphtounicode{esmallkatakana}{30A7}
\pdfglyphtounicode{esmallkatakanahalfwidth}{FF6A}
\pdfglyphtounicode{estimated}{212E}
\pdfglyphtounicode{esuperior}{0065}
\pdfglyphtounicode{eta}{03B7}
\pdfglyphtounicode{etarmenian}{0568}
\pdfglyphtounicode{etatonos}{03AE}
\pdfglyphtounicode{eth}{00F0}
\pdfglyphtounicode{etilde}{1EBD}
\pdfglyphtounicode{etildebelow}{1E1B}
\pdfglyphtounicode{etnahtafoukhhebrew}{0591}
\pdfglyphtounicode{etnahtafoukhlefthebrew}{0591}
\pdfglyphtounicode{etnahtahebrew}{0591}
\pdfglyphtounicode{etnahtalefthebrew}{0591}
\pdfglyphtounicode{eturned}{01DD}
\pdfglyphtounicode{eukorean}{3161}
\pdfglyphtounicode{euro}{20AC}
\pdfglyphtounicode{evowelsignbengali}{09C7}
\pdfglyphtounicode{evowelsigndeva}{0947}
\pdfglyphtounicode{evowelsigngujarati}{0AC7}
\pdfglyphtounicode{exclam}{0021}
\pdfglyphtounicode{exclamarmenian}{055C}
\pdfglyphtounicode{exclamdbl}{203C}
\pdfglyphtounicode{exclamdown}{00A1}
\pdfglyphtounicode{exclamdownsmall}{00A1}
\pdfglyphtounicode{exclammonospace}{FF01}
\pdfglyphtounicode{exclamsmall}{0021}
\pdfglyphtounicode{existential}{2203}
\pdfglyphtounicode{ezh}{0292}
\pdfglyphtounicode{ezhcaron}{01EF}
\pdfglyphtounicode{ezhcurl}{0293}
\pdfglyphtounicode{ezhreversed}{01B9}
\pdfglyphtounicode{ezhtail}{01BA}
\pdfglyphtounicode{f}{0066}
\pdfglyphtounicode{fadeva}{095E}
\pdfglyphtounicode{fagurmukhi}{0A5E}
\pdfglyphtounicode{fahrenheit}{2109}
\pdfglyphtounicode{fathaarabic}{064E}
\pdfglyphtounicode{fathalowarabic}{064E}
\pdfglyphtounicode{fathatanarabic}{064B}
\pdfglyphtounicode{fbopomofo}{3108}
\pdfglyphtounicode{fcircle}{24D5}
\pdfglyphtounicode{fdotaccent}{1E1F}
\pdfglyphtounicode{feharabic}{0641}
\pdfglyphtounicode{feharmenian}{0586}
\pdfglyphtounicode{fehfinalarabic}{FED2}
\pdfglyphtounicode{fehinitialarabic}{FED3}
\pdfglyphtounicode{fehmedialarabic}{FED4}
\pdfglyphtounicode{feicoptic}{03E5}
\pdfglyphtounicode{female}{2640}
\pdfglyphtounicode{ff}{0066 0066}
\pdfglyphtounicode{ffi}{0066 0066 0069}
\pdfglyphtounicode{ffl}{0066 0066 006C}
\pdfglyphtounicode{fi}{0066 0069}
\pdfglyphtounicode{fifteencircle}{246E}
\pdfglyphtounicode{fifteenparen}{2482}
\pdfglyphtounicode{fifteenperiod}{2496}
\pdfglyphtounicode{figuredash}{2012}
\pdfglyphtounicode{filledbox}{25A0}
\pdfglyphtounicode{filledrect}{25AC}
\pdfglyphtounicode{finalkaf}{05DA}
\pdfglyphtounicode{finalkafdagesh}{FB3A}
\pdfglyphtounicode{finalkafdageshhebrew}{FB3A}
\pdfglyphtounicode{finalkafhebrew}{05DA}
\pdfglyphtounicode{finalkafqamats}{05DA 05B8}
\pdfglyphtounicode{finalkafqamatshebrew}{05DA 05B8}
\pdfglyphtounicode{finalkafsheva}{05DA 05B0}
\pdfglyphtounicode{finalkafshevahebrew}{05DA 05B0}
\pdfglyphtounicode{finalmem}{05DD}
\pdfglyphtounicode{finalmemhebrew}{05DD}
\pdfglyphtounicode{finalnun}{05DF}
\pdfglyphtounicode{finalnunhebrew}{05DF}
\pdfglyphtounicode{finalpe}{05E3}
\pdfglyphtounicode{finalpehebrew}{05E3}
\pdfglyphtounicode{finaltsadi}{05E5}
\pdfglyphtounicode{finaltsadihebrew}{05E5}
\pdfglyphtounicode{firsttonechinese}{02C9}
\pdfglyphtounicode{fisheye}{25C9}
\pdfglyphtounicode{fitacyrillic}{0473}
\pdfglyphtounicode{five}{0035}
\pdfglyphtounicode{fivearabic}{0665}
\pdfglyphtounicode{fivebengali}{09EB}
\pdfglyphtounicode{fivecircle}{2464}
\pdfglyphtounicode{fivecircleinversesansserif}{278E}
\pdfglyphtounicode{fivedeva}{096B}
\pdfglyphtounicode{fiveeighths}{215D}
\pdfglyphtounicode{fivegujarati}{0AEB}
\pdfglyphtounicode{fivegurmukhi}{0A6B}
\pdfglyphtounicode{fivehackarabic}{0665}
\pdfglyphtounicode{fivehangzhou}{3025}
\pdfglyphtounicode{fiveideographicparen}{3224}
\pdfglyphtounicode{fiveinferior}{2085}
\pdfglyphtounicode{fivemonospace}{FF15}
\pdfglyphtounicode{fiveoldstyle}{0035}
\pdfglyphtounicode{fiveparen}{2478}
\pdfglyphtounicode{fiveperiod}{248C}
\pdfglyphtounicode{fivepersian}{06F5}
\pdfglyphtounicode{fiveroman}{2174}
\pdfglyphtounicode{fivesuperior}{2075}
\pdfglyphtounicode{fivethai}{0E55}
\pdfglyphtounicode{fl}{0066 006C}
\pdfglyphtounicode{flat}{266D}
\pdfglyphtounicode{floorleft}{230A}
\pdfglyphtounicode{floorright}{230B}
\pdfglyphtounicode{florin}{0192}
\pdfglyphtounicode{fmonospace}{FF46}
\pdfglyphtounicode{fmsquare}{3399}
\pdfglyphtounicode{fofanthai}{0E1F}
\pdfglyphtounicode{fofathai}{0E1D}
\pdfglyphtounicode{follownotdbleqv}{2ABA}
\pdfglyphtounicode{follownotslnteql}{2AB6}
\pdfglyphtounicode{followornoteqvlnt}{22E9}
\pdfglyphtounicode{follows}{227B}
\pdfglyphtounicode{followsequal}{2AB0}
\pdfglyphtounicode{followsorcurly}{227D}
\pdfglyphtounicode{followsorequal}{227F}
\pdfglyphtounicode{fongmanthai}{0E4F}
\pdfglyphtounicode{forall}{2200}
\pdfglyphtounicode{forces}{22A9}
\pdfglyphtounicode{forcesbar}{22AA}
\pdfglyphtounicode{fork}{22D4}
\pdfglyphtounicode{four}{0034}
\pdfglyphtounicode{fourarabic}{0664}
\pdfglyphtounicode{fourbengali}{09EA}
\pdfglyphtounicode{fourcircle}{2463}
\pdfglyphtounicode{fourcircleinversesansserif}{278D}
\pdfglyphtounicode{fourdeva}{096A}
\pdfglyphtounicode{fourgujarati}{0AEA}
\pdfglyphtounicode{fourgurmukhi}{0A6A}
\pdfglyphtounicode{fourhackarabic}{0664}
\pdfglyphtounicode{fourhangzhou}{3024}
\pdfglyphtounicode{fourideographicparen}{3223}
\pdfglyphtounicode{fourinferior}{2084}
\pdfglyphtounicode{fourmonospace}{FF14}
\pdfglyphtounicode{fournumeratorbengali}{09F7}
\pdfglyphtounicode{fouroldstyle}{0034}
\pdfglyphtounicode{fourparen}{2477}
\pdfglyphtounicode{fourperiod}{248B}
\pdfglyphtounicode{fourpersian}{06F4}
\pdfglyphtounicode{fourroman}{2173}
\pdfglyphtounicode{foursuperior}{2074}
\pdfglyphtounicode{fourteencircle}{246D}
\pdfglyphtounicode{fourteenparen}{2481}
\pdfglyphtounicode{fourteenperiod}{2495}
\pdfglyphtounicode{fourthai}{0E54}
\pdfglyphtounicode{fourthtonechinese}{02CB}
\pdfglyphtounicode{fparen}{24A1}
\pdfglyphtounicode{fraction}{2044}
\pdfglyphtounicode{franc}{20A3}
\pdfglyphtounicode{frown}{2322}
\pdfglyphtounicode{g}{0067}
\pdfglyphtounicode{gabengali}{0997}
\pdfglyphtounicode{gacute}{01F5}
\pdfglyphtounicode{gadeva}{0917}
\pdfglyphtounicode{gafarabic}{06AF}
\pdfglyphtounicode{gaffinalarabic}{FB93}
\pdfglyphtounicode{gafinitialarabic}{FB94}
\pdfglyphtounicode{gafmedialarabic}{FB95}
\pdfglyphtounicode{gagujarati}{0A97}
\pdfglyphtounicode{gagurmukhi}{0A17}
\pdfglyphtounicode{gahiragana}{304C}
\pdfglyphtounicode{gakatakana}{30AC}
\pdfglyphtounicode{gamma}{03B3}
\pdfglyphtounicode{gammalatinsmall}{0263}
\pdfglyphtounicode{gammasuperior}{02E0}
\pdfglyphtounicode{gangiacoptic}{03EB}
\pdfglyphtounicode{gbopomofo}{310D}
\pdfglyphtounicode{gbreve}{011F}
\pdfglyphtounicode{gcaron}{01E7}
\pdfglyphtounicode{gcedilla}{0123}
\pdfglyphtounicode{gcircle}{24D6}
\pdfglyphtounicode{gcircumflex}{011D}
\pdfglyphtounicode{gcommaaccent}{0123}
\pdfglyphtounicode{gdot}{0121}
\pdfglyphtounicode{gdotaccent}{0121}
\pdfglyphtounicode{gecyrillic}{0433}
\pdfglyphtounicode{gehiragana}{3052}
\pdfglyphtounicode{gekatakana}{30B2}
\pdfglyphtounicode{geomequivalent}{224E}
\pdfglyphtounicode{geometricallyequal}{2251}
\pdfglyphtounicode{gereshaccenthebrew}{059C}
\pdfglyphtounicode{gereshhebrew}{05F3}
\pdfglyphtounicode{gereshmuqdamhebrew}{059D}
\pdfglyphtounicode{germandbls}{00DF}
\pdfglyphtounicode{gershayimaccenthebrew}{059E}
\pdfglyphtounicode{gershayimhebrew}{05F4}
\pdfglyphtounicode{getamark}{3013}
\pdfglyphtounicode{ghabengali}{0998}
\pdfglyphtounicode{ghadarmenian}{0572}
\pdfglyphtounicode{ghadeva}{0918}
\pdfglyphtounicode{ghagujarati}{0A98}
\pdfglyphtounicode{ghagurmukhi}{0A18}
\pdfglyphtounicode{ghainarabic}{063A}
\pdfglyphtounicode{ghainfinalarabic}{FECE}
\pdfglyphtounicode{ghaininitialarabic}{FECF}
\pdfglyphtounicode{ghainmedialarabic}{FED0}
\pdfglyphtounicode{ghemiddlehookcyrillic}{0495}
\pdfglyphtounicode{ghestrokecyrillic}{0493}
\pdfglyphtounicode{gheupturncyrillic}{0491}
\pdfglyphtounicode{ghhadeva}{095A}
\pdfglyphtounicode{ghhagurmukhi}{0A5A}
\pdfglyphtounicode{ghook}{0260}
\pdfglyphtounicode{ghzsquare}{3393}
\pdfglyphtounicode{gihiragana}{304E}
\pdfglyphtounicode{gikatakana}{30AE}
\pdfglyphtounicode{gimarmenian}{0563}
\pdfglyphtounicode{gimel}{05D2}
\pdfglyphtounicode{gimeldagesh}{FB32}
\pdfglyphtounicode{gimeldageshhebrew}{FB32}
\pdfglyphtounicode{gimelhebrew}{05D2}
\pdfglyphtounicode{gjecyrillic}{0453}
\pdfglyphtounicode{glottalinvertedstroke}{01BE}
\pdfglyphtounicode{glottalstop}{0294}
\pdfglyphtounicode{glottalstopinverted}{0296}
\pdfglyphtounicode{glottalstopmod}{02C0}
\pdfglyphtounicode{glottalstopreversed}{0295}
\pdfglyphtounicode{glottalstopreversedmod}{02C1}
\pdfglyphtounicode{glottalstopreversedsuperior}{02E4}
\pdfglyphtounicode{glottalstopstroke}{02A1}
\pdfglyphtounicode{glottalstopstrokereversed}{02A2}
\pdfglyphtounicode{gmacron}{1E21}
\pdfglyphtounicode{gmonospace}{FF47}
\pdfglyphtounicode{gohiragana}{3054}
\pdfglyphtounicode{gokatakana}{30B4}
\pdfglyphtounicode{gparen}{24A2}
\pdfglyphtounicode{gpasquare}{33AC}
\pdfglyphtounicode{gradient}{2207}
\pdfglyphtounicode{grave}{0060}
\pdfglyphtounicode{gravebelowcmb}{0316}
\pdfglyphtounicode{gravecmb}{0300}
\pdfglyphtounicode{gravecomb}{0300}
\pdfglyphtounicode{gravedeva}{0953}
\pdfglyphtounicode{gravelowmod}{02CE}
\pdfglyphtounicode{gravemonospace}{FF40}
\pdfglyphtounicode{gravetonecmb}{0340}
\pdfglyphtounicode{greater}{003E}
\pdfglyphtounicode{greaterdbleqlless}{2A8C}
\pdfglyphtounicode{greaterdblequal}{2267}
\pdfglyphtounicode{greaterdot}{22D7}
\pdfglyphtounicode{greaterequal}{2265}
\pdfglyphtounicode{greaterequalorless}{22DB}
\pdfglyphtounicode{greaterlessequal}{22DB}
\pdfglyphtounicode{greatermonospace}{FF1E}
\pdfglyphtounicode{greatermuch}{226B}
\pdfglyphtounicode{greaternotdblequal}{2A8A}
\pdfglyphtounicode{greaternotequal}{2A88}
\pdfglyphtounicode{greaterorapproxeql}{2A86}
\pdfglyphtounicode{greaterorequalslant}{2A7E}
\pdfglyphtounicode{greaterorequivalent}{2273}
\pdfglyphtounicode{greaterorless}{2277}
\pdfglyphtounicode{greaterornotdbleql}{2269}
\pdfglyphtounicode{greaterornotequal}{2269}
\pdfglyphtounicode{greaterorsimilar}{2273}
\pdfglyphtounicode{greateroverequal}{2267}
\pdfglyphtounicode{greatersmall}{FE65}
\pdfglyphtounicode{gscript}{0261}
\pdfglyphtounicode{gstroke}{01E5}
\pdfglyphtounicode{guhiragana}{3050}
\pdfglyphtounicode{guillemotleft}{00AB}
\pdfglyphtounicode{guillemotright}{00BB}
\pdfglyphtounicode{guilsinglleft}{2039}
\pdfglyphtounicode{guilsinglright}{203A}
\pdfglyphtounicode{gukatakana}{30B0}
\pdfglyphtounicode{guramusquare}{3318}
\pdfglyphtounicode{gysquare}{33C9}
\pdfglyphtounicode{h}{0068}
\pdfglyphtounicode{haabkhasiancyrillic}{04A9}
\pdfglyphtounicode{haaltonearabic}{06C1}
\pdfglyphtounicode{habengali}{09B9}
\pdfglyphtounicode{hadescendercyrillic}{04B3}
\pdfglyphtounicode{hadeva}{0939}
\pdfglyphtounicode{hagujarati}{0AB9}
\pdfglyphtounicode{hagurmukhi}{0A39}
\pdfglyphtounicode{haharabic}{062D}
\pdfglyphtounicode{hahfinalarabic}{FEA2}
\pdfglyphtounicode{hahinitialarabic}{FEA3}
\pdfglyphtounicode{hahiragana}{306F}
\pdfglyphtounicode{hahmedialarabic}{FEA4}
\pdfglyphtounicode{haitusquare}{332A}
\pdfglyphtounicode{hakatakana}{30CF}
\pdfglyphtounicode{hakatakanahalfwidth}{FF8A}
\pdfglyphtounicode{halantgurmukhi}{0A4D}
\pdfglyphtounicode{hamzaarabic}{0621}
\pdfglyphtounicode{hamzadammaarabic}{0621 064F}
\pdfglyphtounicode{hamzadammatanarabic}{0621 064C}
\pdfglyphtounicode{hamzafathaarabic}{0621 064E}
\pdfglyphtounicode{hamzafathatanarabic}{0621 064B}
\pdfglyphtounicode{hamzalowarabic}{0621}
\pdfglyphtounicode{hamzalowkasraarabic}{0621 0650}
\pdfglyphtounicode{hamzalowkasratanarabic}{0621 064D}
\pdfglyphtounicode{hamzasukunarabic}{0621 0652}
\pdfglyphtounicode{hangulfiller}{3164}
\pdfglyphtounicode{hardsigncyrillic}{044A}
\pdfglyphtounicode{harpoondownleft}{21C3}
\pdfglyphtounicode{harpoondownright}{21C2}
\pdfglyphtounicode{harpoonleftbarbup}{21BC}
\pdfglyphtounicode{harpoonleftright}{21CC}
\pdfglyphtounicode{harpoonrightbarbup}{21C0}
\pdfglyphtounicode{harpoonrightleft}{21CB}
\pdfglyphtounicode{harpoonupleft}{21BF}
\pdfglyphtounicode{harpoonupright}{21BE}
\pdfglyphtounicode{hasquare}{33CA}
\pdfglyphtounicode{hatafpatah}{05B2}
\pdfglyphtounicode{hatafpatah16}{05B2}
\pdfglyphtounicode{hatafpatah23}{05B2}
\pdfglyphtounicode{hatafpatah2f}{05B2}
\pdfglyphtounicode{hatafpatahhebrew}{05B2}
\pdfglyphtounicode{hatafpatahnarrowhebrew}{05B2}
\pdfglyphtounicode{hatafpatahquarterhebrew}{05B2}
\pdfglyphtounicode{hatafpatahwidehebrew}{05B2}
\pdfglyphtounicode{hatafqamats}{05B3}
\pdfglyphtounicode{hatafqamats1b}{05B3}
\pdfglyphtounicode{hatafqamats28}{05B3}
\pdfglyphtounicode{hatafqamats34}{05B3}
\pdfglyphtounicode{hatafqamatshebrew}{05B3}
\pdfglyphtounicode{hatafqamatsnarrowhebrew}{05B3}
\pdfglyphtounicode{hatafqamatsquarterhebrew}{05B3}
\pdfglyphtounicode{hatafqamatswidehebrew}{05B3}
\pdfglyphtounicode{hatafsegol}{05B1}
\pdfglyphtounicode{hatafsegol17}{05B1}
\pdfglyphtounicode{hatafsegol24}{05B1}
\pdfglyphtounicode{hatafsegol30}{05B1}
\pdfglyphtounicode{hatafsegolhebrew}{05B1}
\pdfglyphtounicode{hatafsegolnarrowhebrew}{05B1}
\pdfglyphtounicode{hatafsegolquarterhebrew}{05B1}
\pdfglyphtounicode{hatafsegolwidehebrew}{05B1}
\pdfglyphtounicode{hbar}{0127}
\pdfglyphtounicode{hbopomofo}{310F}
\pdfglyphtounicode{hbrevebelow}{1E2B}
\pdfglyphtounicode{hcedilla}{1E29}
\pdfglyphtounicode{hcircle}{24D7}
\pdfglyphtounicode{hcircumflex}{0125}
\pdfglyphtounicode{hdieresis}{1E27}
\pdfglyphtounicode{hdotaccent}{1E23}
\pdfglyphtounicode{hdotbelow}{1E25}
\pdfglyphtounicode{he}{05D4}
\pdfglyphtounicode{heart}{2665}
\pdfglyphtounicode{heartsuitblack}{2665}
\pdfglyphtounicode{heartsuitwhite}{2661}
\pdfglyphtounicode{hedagesh}{FB34}
\pdfglyphtounicode{hedageshhebrew}{FB34}
\pdfglyphtounicode{hehaltonearabic}{06C1}
\pdfglyphtounicode{heharabic}{0647}
\pdfglyphtounicode{hehebrew}{05D4}
\pdfglyphtounicode{hehfinalaltonearabic}{FBA7}
\pdfglyphtounicode{hehfinalalttwoarabic}{FEEA}
\pdfglyphtounicode{hehfinalarabic}{FEEA}
\pdfglyphtounicode{hehhamzaabovefinalarabic}{FBA5}
\pdfglyphtounicode{hehhamzaaboveisolatedarabic}{FBA4}
\pdfglyphtounicode{hehinitialaltonearabic}{FBA8}
\pdfglyphtounicode{hehinitialarabic}{FEEB}
\pdfglyphtounicode{hehiragana}{3078}
\pdfglyphtounicode{hehmedialaltonearabic}{FBA9}
\pdfglyphtounicode{hehmedialarabic}{FEEC}
\pdfglyphtounicode{heiseierasquare}{337B}
\pdfglyphtounicode{hekatakana}{30D8}
\pdfglyphtounicode{hekatakanahalfwidth}{FF8D}
\pdfglyphtounicode{hekutaarusquare}{3336}
\pdfglyphtounicode{henghook}{0267}
\pdfglyphtounicode{herutusquare}{3339}
\pdfglyphtounicode{het}{05D7}
\pdfglyphtounicode{hethebrew}{05D7}
\pdfglyphtounicode{hhook}{0266}
\pdfglyphtounicode{hhooksuperior}{02B1}
\pdfglyphtounicode{hieuhacirclekorean}{327B}
\pdfglyphtounicode{hieuhaparenkorean}{321B}
\pdfglyphtounicode{hieuhcirclekorean}{326D}
\pdfglyphtounicode{hieuhkorean}{314E}
\pdfglyphtounicode{hieuhparenkorean}{320D}
\pdfglyphtounicode{hihiragana}{3072}
\pdfglyphtounicode{hikatakana}{30D2}
\pdfglyphtounicode{hikatakanahalfwidth}{FF8B}
\pdfglyphtounicode{hiriq}{05B4}
\pdfglyphtounicode{hiriq14}{05B4}
\pdfglyphtounicode{hiriq21}{05B4}
\pdfglyphtounicode{hiriq2d}{05B4}
\pdfglyphtounicode{hiriqhebrew}{05B4}
\pdfglyphtounicode{hiriqnarrowhebrew}{05B4}
\pdfglyphtounicode{hiriqquarterhebrew}{05B4}
\pdfglyphtounicode{hiriqwidehebrew}{05B4}
\pdfglyphtounicode{hlinebelow}{1E96}
\pdfglyphtounicode{hmonospace}{FF48}
\pdfglyphtounicode{hoarmenian}{0570}
\pdfglyphtounicode{hohipthai}{0E2B}
\pdfglyphtounicode{hohiragana}{307B}
\pdfglyphtounicode{hokatakana}{30DB}
\pdfglyphtounicode{hokatakanahalfwidth}{FF8E}
\pdfglyphtounicode{holam}{05B9}
\pdfglyphtounicode{holam19}{05B9}
\pdfglyphtounicode{holam26}{05B9}
\pdfglyphtounicode{holam32}{05B9}
\pdfglyphtounicode{holamhebrew}{05B9}
\pdfglyphtounicode{holamnarrowhebrew}{05B9}
\pdfglyphtounicode{holamquarterhebrew}{05B9}
\pdfglyphtounicode{holamwidehebrew}{05B9}
\pdfglyphtounicode{honokhukthai}{0E2E}
\pdfglyphtounicode{hookabovecomb}{0309}
\pdfglyphtounicode{hookcmb}{0309}
\pdfglyphtounicode{hookpalatalizedbelowcmb}{0321}
\pdfglyphtounicode{hookretroflexbelowcmb}{0322}
\pdfglyphtounicode{hoonsquare}{3342}
\pdfglyphtounicode{horicoptic}{03E9}
\pdfglyphtounicode{horizontalbar}{2015}
\pdfglyphtounicode{horncmb}{031B}
\pdfglyphtounicode{hotsprings}{2668}
\pdfglyphtounicode{house}{2302}
\pdfglyphtounicode{hparen}{24A3}
\pdfglyphtounicode{hsuperior}{02B0}
\pdfglyphtounicode{hturned}{0265}
\pdfglyphtounicode{huhiragana}{3075}
\pdfglyphtounicode{huiitosquare}{3333}
\pdfglyphtounicode{hukatakana}{30D5}
\pdfglyphtounicode{hukatakanahalfwidth}{FF8C}
\pdfglyphtounicode{hungarumlaut}{02DD}
\pdfglyphtounicode{hungarumlautcmb}{030B}
\pdfglyphtounicode{hv}{0195}
\pdfglyphtounicode{hyphen}{002D}
\pdfglyphtounicode{hyphenchar}{002D}
\pdfglyphtounicode{hypheninferior}{002D}
\pdfglyphtounicode{hyphenmonospace}{FF0D}
\pdfglyphtounicode{hyphensmall}{FE63}
\pdfglyphtounicode{hyphensuperior}{002D}
\pdfglyphtounicode{hyphentwo}{2010}
\pdfglyphtounicode{i}{0069}
\pdfglyphtounicode{iacute}{00ED}
\pdfglyphtounicode{iacyrillic}{044F}
\pdfglyphtounicode{ibengali}{0987}
\pdfglyphtounicode{ibopomofo}{3127}
\pdfglyphtounicode{ibreve}{012D}
\pdfglyphtounicode{icaron}{01D0}
\pdfglyphtounicode{icircle}{24D8}
\pdfglyphtounicode{icircumflex}{00EE}
\pdfglyphtounicode{icyrillic}{0456}
\pdfglyphtounicode{idblgrave}{0209}
\pdfglyphtounicode{ideographearthcircle}{328F}
\pdfglyphtounicode{ideographfirecircle}{328B}
\pdfglyphtounicode{ideographicallianceparen}{323F}
\pdfglyphtounicode{ideographiccallparen}{323A}
\pdfglyphtounicode{ideographiccentrecircle}{32A5}
\pdfglyphtounicode{ideographicclose}{3006}
\pdfglyphtounicode{ideographiccomma}{3001}
\pdfglyphtounicode{ideographiccommaleft}{FF64}
\pdfglyphtounicode{ideographiccongratulationparen}{3237}
\pdfglyphtounicode{ideographiccorrectcircle}{32A3}
\pdfglyphtounicode{ideographicearthparen}{322F}
\pdfglyphtounicode{ideographicenterpriseparen}{323D}
\pdfglyphtounicode{ideographicexcellentcircle}{329D}
\pdfglyphtounicode{ideographicfestivalparen}{3240}
\pdfglyphtounicode{ideographicfinancialcircle}{3296}
\pdfglyphtounicode{ideographicfinancialparen}{3236}
\pdfglyphtounicode{ideographicfireparen}{322B}
\pdfglyphtounicode{ideographichaveparen}{3232}
\pdfglyphtounicode{ideographichighcircle}{32A4}
\pdfglyphtounicode{ideographiciterationmark}{3005}
\pdfglyphtounicode{ideographiclaborcircle}{3298}
\pdfglyphtounicode{ideographiclaborparen}{3238}
\pdfglyphtounicode{ideographicleftcircle}{32A7}
\pdfglyphtounicode{ideographiclowcircle}{32A6}
\pdfglyphtounicode{ideographicmedicinecircle}{32A9}
\pdfglyphtounicode{ideographicmetalparen}{322E}
\pdfglyphtounicode{ideographicmoonparen}{322A}
\pdfglyphtounicode{ideographicnameparen}{3234}
\pdfglyphtounicode{ideographicperiod}{3002}
\pdfglyphtounicode{ideographicprintcircle}{329E}
\pdfglyphtounicode{ideographicreachparen}{3243}
\pdfglyphtounicode{ideographicrepresentparen}{3239}
\pdfglyphtounicode{ideographicresourceparen}{323E}
\pdfglyphtounicode{ideographicrightcircle}{32A8}
\pdfglyphtounicode{ideographicsecretcircle}{3299}
\pdfglyphtounicode{ideographicselfparen}{3242}
\pdfglyphtounicode{ideographicsocietyparen}{3233}
\pdfglyphtounicode{ideographicspace}{3000}
\pdfglyphtounicode{ideographicspecialparen}{3235}
\pdfglyphtounicode{ideographicstockparen}{3231}
\pdfglyphtounicode{ideographicstudyparen}{323B}
\pdfglyphtounicode{ideographicsunparen}{3230}
\pdfglyphtounicode{ideographicsuperviseparen}{323C}
\pdfglyphtounicode{ideographicwaterparen}{322C}
\pdfglyphtounicode{ideographicwoodparen}{322D}
\pdfglyphtounicode{ideographiczero}{3007}
\pdfglyphtounicode{ideographmetalcircle}{328E}
\pdfglyphtounicode{ideographmooncircle}{328A}
\pdfglyphtounicode{ideographnamecircle}{3294}
\pdfglyphtounicode{ideographsuncircle}{3290}
\pdfglyphtounicode{ideographwatercircle}{328C}
\pdfglyphtounicode{ideographwoodcircle}{328D}
\pdfglyphtounicode{ideva}{0907}
\pdfglyphtounicode{idieresis}{00EF}
\pdfglyphtounicode{idieresisacute}{1E2F}
\pdfglyphtounicode{idieresiscyrillic}{04E5}
\pdfglyphtounicode{idotbelow}{1ECB}
\pdfglyphtounicode{iebrevecyrillic}{04D7}
\pdfglyphtounicode{iecyrillic}{0435}
\pdfglyphtounicode{ieungacirclekorean}{3275}
\pdfglyphtounicode{ieungaparenkorean}{3215}
\pdfglyphtounicode{ieungcirclekorean}{3267}
\pdfglyphtounicode{ieungkorean}{3147}
\pdfglyphtounicode{ieungparenkorean}{3207}
\pdfglyphtounicode{igrave}{00EC}
\pdfglyphtounicode{igujarati}{0A87}
\pdfglyphtounicode{igurmukhi}{0A07}
\pdfglyphtounicode{ihiragana}{3044}
\pdfglyphtounicode{ihookabove}{1EC9}
\pdfglyphtounicode{iibengali}{0988}
\pdfglyphtounicode{iicyrillic}{0438}
\pdfglyphtounicode{iideva}{0908}
\pdfglyphtounicode{iigujarati}{0A88}
\pdfglyphtounicode{iigurmukhi}{0A08}
\pdfglyphtounicode{iimatragurmukhi}{0A40}
\pdfglyphtounicode{iinvertedbreve}{020B}
\pdfglyphtounicode{iishortcyrillic}{0439}
\pdfglyphtounicode{iivowelsignbengali}{09C0}
\pdfglyphtounicode{iivowelsigndeva}{0940}
\pdfglyphtounicode{iivowelsigngujarati}{0AC0}
\pdfglyphtounicode{ij}{0133}
\pdfglyphtounicode{ikatakana}{30A4}
\pdfglyphtounicode{ikatakanahalfwidth}{FF72}
\pdfglyphtounicode{ikorean}{3163}
\pdfglyphtounicode{ilde}{02DC}
\pdfglyphtounicode{iluyhebrew}{05AC}
\pdfglyphtounicode{imacron}{012B}
\pdfglyphtounicode{imacroncyrillic}{04E3}
\pdfglyphtounicode{imageorapproximatelyequal}{2253}
\pdfglyphtounicode{imatragurmukhi}{0A3F}
\pdfglyphtounicode{imonospace}{FF49}
\pdfglyphtounicode{increment}{2206}
\pdfglyphtounicode{infinity}{221E}
\pdfglyphtounicode{iniarmenian}{056B}
\pdfglyphtounicode{integerdivide}{2216}
\pdfglyphtounicode{integral}{222B}
\pdfglyphtounicode{integralbottom}{2321}
\pdfglyphtounicode{integralbt}{2321}
\pdfglyphtounicode{integralex}{F8F5}
\pdfglyphtounicode{integraltop}{2320}
\pdfglyphtounicode{integraltp}{2320}
\pdfglyphtounicode{intercal}{22BA}
\pdfglyphtounicode{interrobang}{203D}
\pdfglyphtounicode{interrobangdown}{2E18}
\pdfglyphtounicode{intersection}{2229}
\pdfglyphtounicode{intersectiondbl}{22D2}
\pdfglyphtounicode{intersectionsq}{2293}
\pdfglyphtounicode{intisquare}{3305}
\pdfglyphtounicode{invbullet}{25D8}
\pdfglyphtounicode{invcircle}{25D9}
\pdfglyphtounicode{invsmileface}{263B}
\pdfglyphtounicode{iocyrillic}{0451}
\pdfglyphtounicode{iogonek}{012F}
\pdfglyphtounicode{iota}{03B9}
\pdfglyphtounicode{iotadieresis}{03CA}
\pdfglyphtounicode{iotadieresistonos}{0390}
\pdfglyphtounicode{iotalatin}{0269}
\pdfglyphtounicode{iotatonos}{03AF}
\pdfglyphtounicode{iparen}{24A4}
\pdfglyphtounicode{irigurmukhi}{0A72}
\pdfglyphtounicode{ismallhiragana}{3043}
\pdfglyphtounicode{ismallkatakana}{30A3}
\pdfglyphtounicode{ismallkatakanahalfwidth}{FF68}
\pdfglyphtounicode{issharbengali}{09FA}
\pdfglyphtounicode{istroke}{0268}
\pdfglyphtounicode{isuperior}{0069}
\pdfglyphtounicode{iterationhiragana}{309D}
\pdfglyphtounicode{iterationkatakana}{30FD}
\pdfglyphtounicode{itilde}{0129}
\pdfglyphtounicode{itildebelow}{1E2D}
\pdfglyphtounicode{iubopomofo}{3129}
\pdfglyphtounicode{iucyrillic}{044E}
\pdfglyphtounicode{ivowelsignbengali}{09BF}
\pdfglyphtounicode{ivowelsigndeva}{093F}
\pdfglyphtounicode{ivowelsigngujarati}{0ABF}
\pdfglyphtounicode{izhitsacyrillic}{0475}
\pdfglyphtounicode{izhitsadblgravecyrillic}{0477}
\pdfglyphtounicode{j}{006A}
\pdfglyphtounicode{jaarmenian}{0571}
\pdfglyphtounicode{jabengali}{099C}
\pdfglyphtounicode{jadeva}{091C}
\pdfglyphtounicode{jagujarati}{0A9C}
\pdfglyphtounicode{jagurmukhi}{0A1C}
\pdfglyphtounicode{jbopomofo}{3110}
\pdfglyphtounicode{jcaron}{01F0}
\pdfglyphtounicode{jcircle}{24D9}
\pdfglyphtounicode{jcircumflex}{0135}
\pdfglyphtounicode{jcrossedtail}{029D}
\pdfglyphtounicode{jdotlessstroke}{025F}
\pdfglyphtounicode{jecyrillic}{0458}
\pdfglyphtounicode{jeemarabic}{062C}
\pdfglyphtounicode{jeemfinalarabic}{FE9E}
\pdfglyphtounicode{jeeminitialarabic}{FE9F}
\pdfglyphtounicode{jeemmedialarabic}{FEA0}
\pdfglyphtounicode{jeharabic}{0698}
\pdfglyphtounicode{jehfinalarabic}{FB8B}
\pdfglyphtounicode{jhabengali}{099D}
\pdfglyphtounicode{jhadeva}{091D}
\pdfglyphtounicode{jhagujarati}{0A9D}
\pdfglyphtounicode{jhagurmukhi}{0A1D}
\pdfglyphtounicode{jheharmenian}{057B}
\pdfglyphtounicode{jis}{3004}
\pdfglyphtounicode{jmonospace}{FF4A}
\pdfglyphtounicode{jparen}{24A5}
\pdfglyphtounicode{jsuperior}{02B2}
\pdfglyphtounicode{k}{006B}
\pdfglyphtounicode{kabashkircyrillic}{04A1}
\pdfglyphtounicode{kabengali}{0995}
\pdfglyphtounicode{kacute}{1E31}
\pdfglyphtounicode{kacyrillic}{043A}
\pdfglyphtounicode{kadescendercyrillic}{049B}
\pdfglyphtounicode{kadeva}{0915}
\pdfglyphtounicode{kaf}{05DB}
\pdfglyphtounicode{kafarabic}{0643}
\pdfglyphtounicode{kafdagesh}{FB3B}
\pdfglyphtounicode{kafdageshhebrew}{FB3B}
\pdfglyphtounicode{kaffinalarabic}{FEDA}
\pdfglyphtounicode{kafhebrew}{05DB}
\pdfglyphtounicode{kafinitialarabic}{FEDB}
\pdfglyphtounicode{kafmedialarabic}{FEDC}
\pdfglyphtounicode{kafrafehebrew}{FB4D}
\pdfglyphtounicode{kagujarati}{0A95}
\pdfglyphtounicode{kagurmukhi}{0A15}
\pdfglyphtounicode{kahiragana}{304B}
\pdfglyphtounicode{kahookcyrillic}{04C4}
\pdfglyphtounicode{kakatakana}{30AB}
\pdfglyphtounicode{kakatakanahalfwidth}{FF76}
\pdfglyphtounicode{kappa}{03BA}
\pdfglyphtounicode{kappasymbolgreek}{03F0}
\pdfglyphtounicode{kapyeounmieumkorean}{3171}
\pdfglyphtounicode{kapyeounphieuphkorean}{3184}
\pdfglyphtounicode{kapyeounpieupkorean}{3178}
\pdfglyphtounicode{kapyeounssangpieupkorean}{3179}
\pdfglyphtounicode{karoriisquare}{330D}
\pdfglyphtounicode{kashidaautoarabic}{0640}
\pdfglyphtounicode{kashidaautonosidebearingarabic}{0640}
\pdfglyphtounicode{kasmallkatakana}{30F5}
\pdfglyphtounicode{kasquare}{3384}
\pdfglyphtounicode{kasraarabic}{0650}
\pdfglyphtounicode{kasratanarabic}{064D}
\pdfglyphtounicode{kastrokecyrillic}{049F}
\pdfglyphtounicode{katahiraprolongmarkhalfwidth}{FF70}
\pdfglyphtounicode{kaverticalstrokecyrillic}{049D}
\pdfglyphtounicode{kbopomofo}{310E}
\pdfglyphtounicode{kcalsquare}{3389}
\pdfglyphtounicode{kcaron}{01E9}
\pdfglyphtounicode{kcedilla}{0137}
\pdfglyphtounicode{kcircle}{24DA}
\pdfglyphtounicode{kcommaaccent}{0137}
\pdfglyphtounicode{kdotbelow}{1E33}
\pdfglyphtounicode{keharmenian}{0584}
\pdfglyphtounicode{kehiragana}{3051}
\pdfglyphtounicode{kekatakana}{30B1}
\pdfglyphtounicode{kekatakanahalfwidth}{FF79}
\pdfglyphtounicode{kenarmenian}{056F}
\pdfglyphtounicode{kesmallkatakana}{30F6}
\pdfglyphtounicode{kgreenlandic}{0138}
\pdfglyphtounicode{khabengali}{0996}
\pdfglyphtounicode{khacyrillic}{0445}
\pdfglyphtounicode{khadeva}{0916}
\pdfglyphtounicode{khagujarati}{0A96}
\pdfglyphtounicode{khagurmukhi}{0A16}
\pdfglyphtounicode{khaharabic}{062E}
\pdfglyphtounicode{khahfinalarabic}{FEA6}
\pdfglyphtounicode{khahinitialarabic}{FEA7}
\pdfglyphtounicode{khahmedialarabic}{FEA8}
\pdfglyphtounicode{kheicoptic}{03E7}
\pdfglyphtounicode{khhadeva}{0959}
\pdfglyphtounicode{khhagurmukhi}{0A59}
\pdfglyphtounicode{khieukhacirclekorean}{3278}
\pdfglyphtounicode{khieukhaparenkorean}{3218}
\pdfglyphtounicode{khieukhcirclekorean}{326A}
\pdfglyphtounicode{khieukhkorean}{314B}
\pdfglyphtounicode{khieukhparenkorean}{320A}
\pdfglyphtounicode{khokhaithai}{0E02}
\pdfglyphtounicode{khokhonthai}{0E05}
\pdfglyphtounicode{khokhuatthai}{0E03}
\pdfglyphtounicode{khokhwaithai}{0E04}
\pdfglyphtounicode{khomutthai}{0E5B}
\pdfglyphtounicode{khook}{0199}
\pdfglyphtounicode{khorakhangthai}{0E06}
\pdfglyphtounicode{khzsquare}{3391}
\pdfglyphtounicode{kihiragana}{304D}
\pdfglyphtounicode{kikatakana}{30AD}
\pdfglyphtounicode{kikatakanahalfwidth}{FF77}
\pdfglyphtounicode{kiroguramusquare}{3315}
\pdfglyphtounicode{kiromeetorusquare}{3316}
\pdfglyphtounicode{kirosquare}{3314}
\pdfglyphtounicode{kiyeokacirclekorean}{326E}
\pdfglyphtounicode{kiyeokaparenkorean}{320E}
\pdfglyphtounicode{kiyeokcirclekorean}{3260}
\pdfglyphtounicode{kiyeokkorean}{3131}
\pdfglyphtounicode{kiyeokparenkorean}{3200}
\pdfglyphtounicode{kiyeoksioskorean}{3133}
\pdfglyphtounicode{kjecyrillic}{045C}
\pdfglyphtounicode{klinebelow}{1E35}
\pdfglyphtounicode{klsquare}{3398}
\pdfglyphtounicode{kmcubedsquare}{33A6}
\pdfglyphtounicode{kmonospace}{FF4B}
\pdfglyphtounicode{kmsquaredsquare}{33A2}
\pdfglyphtounicode{kohiragana}{3053}
\pdfglyphtounicode{kohmsquare}{33C0}
\pdfglyphtounicode{kokaithai}{0E01}
\pdfglyphtounicode{kokatakana}{30B3}
\pdfglyphtounicode{kokatakanahalfwidth}{FF7A}
\pdfglyphtounicode{kooposquare}{331E}
\pdfglyphtounicode{koppacyrillic}{0481}
\pdfglyphtounicode{koreanstandardsymbol}{327F}
\pdfglyphtounicode{koroniscmb}{0343}
\pdfglyphtounicode{kparen}{24A6}
\pdfglyphtounicode{kpasquare}{33AA}
\pdfglyphtounicode{ksicyrillic}{046F}
\pdfglyphtounicode{ktsquare}{33CF}
\pdfglyphtounicode{kturned}{029E}
\pdfglyphtounicode{kuhiragana}{304F}
\pdfglyphtounicode{kukatakana}{30AF}
\pdfglyphtounicode{kukatakanahalfwidth}{FF78}
\pdfglyphtounicode{kvsquare}{33B8}
\pdfglyphtounicode{kwsquare}{33BE}
\pdfglyphtounicode{l}{006C}
\pdfglyphtounicode{labengali}{09B2}
\pdfglyphtounicode{lacute}{013A}
\pdfglyphtounicode{ladeva}{0932}
\pdfglyphtounicode{lagujarati}{0AB2}
\pdfglyphtounicode{lagurmukhi}{0A32}
\pdfglyphtounicode{lakkhangyaothai}{0E45}
\pdfglyphtounicode{lamaleffinalarabic}{FEFC}
\pdfglyphtounicode{lamalefhamzaabovefinalarabic}{FEF8}
\pdfglyphtounicode{lamalefhamzaaboveisolatedarabic}{FEF7}
\pdfglyphtounicode{lamalefhamzabelowfinalarabic}{FEFA}
\pdfglyphtounicode{lamalefhamzabelowisolatedarabic}{FEF9}
\pdfglyphtounicode{lamalefisolatedarabic}{FEFB}
\pdfglyphtounicode{lamalefmaddaabovefinalarabic}{FEF6}
\pdfglyphtounicode{lamalefmaddaaboveisolatedarabic}{FEF5}
\pdfglyphtounicode{lamarabic}{0644}
\pdfglyphtounicode{lambda}{03BB}
\pdfglyphtounicode{lambdastroke}{019B}
\pdfglyphtounicode{lamed}{05DC}
\pdfglyphtounicode{lameddagesh}{FB3C}
\pdfglyphtounicode{lameddageshhebrew}{FB3C}
\pdfglyphtounicode{lamedhebrew}{05DC}
\pdfglyphtounicode{lamedholam}{05DC 05B9}
\pdfglyphtounicode{lamedholamdagesh}{05DC 05B9 05BC}
\pdfglyphtounicode{lamedholamdageshhebrew}{05DC 05B9 05BC}
\pdfglyphtounicode{lamedholamhebrew}{05DC 05B9}
\pdfglyphtounicode{lamfinalarabic}{FEDE}
\pdfglyphtounicode{lamhahinitialarabic}{FCCA}
\pdfglyphtounicode{laminitialarabic}{FEDF}
\pdfglyphtounicode{lamjeeminitialarabic}{FCC9}
\pdfglyphtounicode{lamkhahinitialarabic}{FCCB}
\pdfglyphtounicode{lamlamhehisolatedarabic}{FDF2}
\pdfglyphtounicode{lammedialarabic}{FEE0}
\pdfglyphtounicode{lammeemhahinitialarabic}{FD88}
\pdfglyphtounicode{lammeeminitialarabic}{FCCC}
\pdfglyphtounicode{lammeemjeeminitialarabic}{FEDF FEE4 FEA0}
\pdfglyphtounicode{lammeemkhahinitialarabic}{FEDF FEE4 FEA8}
\pdfglyphtounicode{largecircle}{25EF}
\pdfglyphtounicode{latticetop}{22A4}
\pdfglyphtounicode{lbar}{019A}
\pdfglyphtounicode{lbelt}{026C}
\pdfglyphtounicode{lbopomofo}{310C}
\pdfglyphtounicode{lcaron}{013E}
\pdfglyphtounicode{lcedilla}{013C}
\pdfglyphtounicode{lcircle}{24DB}
\pdfglyphtounicode{lcircumflexbelow}{1E3D}
\pdfglyphtounicode{lcommaaccent}{013C}
\pdfglyphtounicode{ldot}{0140}
\pdfglyphtounicode{ldotaccent}{0140}
\pdfglyphtounicode{ldotbelow}{1E37}
\pdfglyphtounicode{ldotbelowmacron}{1E39}
\pdfglyphtounicode{leftangleabovecmb}{031A}
\pdfglyphtounicode{lefttackbelowcmb}{0318}
\pdfglyphtounicode{less}{003C}
\pdfglyphtounicode{lessdbleqlgreater}{2A8B}
\pdfglyphtounicode{lessdblequal}{2266}
\pdfglyphtounicode{lessdot}{22D6}
\pdfglyphtounicode{lessequal}{2264}
\pdfglyphtounicode{lessequalgreater}{22DA}
\pdfglyphtounicode{lessequalorgreater}{22DA}
\pdfglyphtounicode{lessmonospace}{FF1C}
\pdfglyphtounicode{lessmuch}{226A}
\pdfglyphtounicode{lessnotdblequal}{2A89}
\pdfglyphtounicode{lessnotequal}{2A87}
\pdfglyphtounicode{lessorapproxeql}{2A85}
\pdfglyphtounicode{lessorequalslant}{2A7D}
\pdfglyphtounicode{lessorequivalent}{2272}
\pdfglyphtounicode{lessorgreater}{2276}
\pdfglyphtounicode{lessornotdbleql}{2268}
\pdfglyphtounicode{lessornotequal}{2268}
\pdfglyphtounicode{lessorsimilar}{2272}
\pdfglyphtounicode{lessoverequal}{2266}
\pdfglyphtounicode{lesssmall}{FE64}
\pdfglyphtounicode{lezh}{026E}
\pdfglyphtounicode{lfblock}{258C}
\pdfglyphtounicode{lhookretroflex}{026D}
\pdfglyphtounicode{lira}{20A4}
\pdfglyphtounicode{liwnarmenian}{056C}
\pdfglyphtounicode{lj}{01C9}
\pdfglyphtounicode{ljecyrillic}{0459}
\pdfglyphtounicode{ll}{006C 006C}
\pdfglyphtounicode{lladeva}{0933}
\pdfglyphtounicode{llagujarati}{0AB3}
\pdfglyphtounicode{llinebelow}{1E3B}
\pdfglyphtounicode{llladeva}{0934}
\pdfglyphtounicode{llvocalicbengali}{09E1}
\pdfglyphtounicode{llvocalicdeva}{0961}
\pdfglyphtounicode{llvocalicvowelsignbengali}{09E3}
\pdfglyphtounicode{llvocalicvowelsigndeva}{0963}
\pdfglyphtounicode{lmiddletilde}{026B}
\pdfglyphtounicode{lmonospace}{FF4C}
\pdfglyphtounicode{lmsquare}{33D0}
\pdfglyphtounicode{lochulathai}{0E2C}
\pdfglyphtounicode{logicaland}{2227}
\pdfglyphtounicode{logicalnot}{00AC}
\pdfglyphtounicode{logicalnotreversed}{2310}
\pdfglyphtounicode{logicalor}{2228}
\pdfglyphtounicode{lolingthai}{0E25}
\pdfglyphtounicode{longdbls}{017F 017F}
\pdfglyphtounicode{longs}{017F}
\pdfglyphtounicode{longsh}{017F 0068}
\pdfglyphtounicode{longsi}{017F 0069}
\pdfglyphtounicode{longsl}{017F 006C}
\pdfglyphtounicode{longst}{017F 0074}
\pdfglyphtounicode{lowlinecenterline}{FE4E}
\pdfglyphtounicode{lowlinecmb}{0332}
\pdfglyphtounicode{lowlinedashed}{FE4D}
\pdfglyphtounicode{lozenge}{25CA}
\pdfglyphtounicode{lparen}{24A7}
\pdfglyphtounicode{lscript}{2113}
\pdfglyphtounicode{lslash}{0142}
\pdfglyphtounicode{lsquare}{2113}
\pdfglyphtounicode{lsuperior}{006C}
\pdfglyphtounicode{ltshade}{2591}
\pdfglyphtounicode{luthai}{0E26}
\pdfglyphtounicode{lvocalicbengali}{098C}
\pdfglyphtounicode{lvocalicdeva}{090C}
\pdfglyphtounicode{lvocalicvowelsignbengali}{09E2}
\pdfglyphtounicode{lvocalicvowelsigndeva}{0962}
\pdfglyphtounicode{lxsquare}{33D3}
\pdfglyphtounicode{m}{006D}
\pdfglyphtounicode{mabengali}{09AE}
\pdfglyphtounicode{macron}{00AF}
\pdfglyphtounicode{macronbelowcmb}{0331}
\pdfglyphtounicode{macroncmb}{0304}
\pdfglyphtounicode{macronlowmod}{02CD}
\pdfglyphtounicode{macronmonospace}{FFE3}
\pdfglyphtounicode{macute}{1E3F}
\pdfglyphtounicode{madeva}{092E}
\pdfglyphtounicode{magujarati}{0AAE}
\pdfglyphtounicode{magurmukhi}{0A2E}
\pdfglyphtounicode{mahapakhhebrew}{05A4}
\pdfglyphtounicode{mahapakhlefthebrew}{05A4}
\pdfglyphtounicode{mahiragana}{307E}
\pdfglyphtounicode{maichattawalowleftthai}{F895}
\pdfglyphtounicode{maichattawalowrightthai}{F894}
\pdfglyphtounicode{maichattawathai}{0E4B}
\pdfglyphtounicode{maichattawaupperleftthai}{F893}
\pdfglyphtounicode{maieklowleftthai}{F88C}
\pdfglyphtounicode{maieklowrightthai}{F88B}
\pdfglyphtounicode{maiekthai}{0E48}
\pdfglyphtounicode{maiekupperleftthai}{F88A}
\pdfglyphtounicode{maihanakatleftthai}{F884}
\pdfglyphtounicode{maihanakatthai}{0E31}
\pdfglyphtounicode{maitaikhuleftthai}{F889}
\pdfglyphtounicode{maitaikhuthai}{0E47}
\pdfglyphtounicode{maitholowleftthai}{F88F}
\pdfglyphtounicode{maitholowrightthai}{F88E}
\pdfglyphtounicode{maithothai}{0E49}
\pdfglyphtounicode{maithoupperleftthai}{F88D}
\pdfglyphtounicode{maitrilowleftthai}{F892}
\pdfglyphtounicode{maitrilowrightthai}{F891}
\pdfglyphtounicode{maitrithai}{0E4A}
\pdfglyphtounicode{maitriupperleftthai}{F890}
\pdfglyphtounicode{maiyamokthai}{0E46}
\pdfglyphtounicode{makatakana}{30DE}
\pdfglyphtounicode{makatakanahalfwidth}{FF8F}
\pdfglyphtounicode{male}{2642}
\pdfglyphtounicode{maltesecross}{2720}
\pdfglyphtounicode{mansyonsquare}{3347}
\pdfglyphtounicode{maqafhebrew}{05BE}
\pdfglyphtounicode{mars}{2642}
\pdfglyphtounicode{masoracirclehebrew}{05AF}
\pdfglyphtounicode{masquare}{3383}
\pdfglyphtounicode{mbopomofo}{3107}
\pdfglyphtounicode{mbsquare}{33D4}
\pdfglyphtounicode{mcircle}{24DC}
\pdfglyphtounicode{mcubedsquare}{33A5}
\pdfglyphtounicode{mdotaccent}{1E41}
\pdfglyphtounicode{mdotbelow}{1E43}
\pdfglyphtounicode{measuredangle}{2221}
\pdfglyphtounicode{meemarabic}{0645}
\pdfglyphtounicode{meemfinalarabic}{FEE2}
\pdfglyphtounicode{meeminitialarabic}{FEE3}
\pdfglyphtounicode{meemmedialarabic}{FEE4}
\pdfglyphtounicode{meemmeeminitialarabic}{FCD1}
\pdfglyphtounicode{meemmeemisolatedarabic}{FC48}
\pdfglyphtounicode{meetorusquare}{334D}
\pdfglyphtounicode{mehiragana}{3081}
\pdfglyphtounicode{meizierasquare}{337E}
\pdfglyphtounicode{mekatakana}{30E1}
\pdfglyphtounicode{mekatakanahalfwidth}{FF92}
\pdfglyphtounicode{mem}{05DE}
\pdfglyphtounicode{memdagesh}{FB3E}
\pdfglyphtounicode{memdageshhebrew}{FB3E}
\pdfglyphtounicode{memhebrew}{05DE}
\pdfglyphtounicode{menarmenian}{0574}
\pdfglyphtounicode{merkhahebrew}{05A5}
\pdfglyphtounicode{merkhakefulahebrew}{05A6}
\pdfglyphtounicode{merkhakefulalefthebrew}{05A6}
\pdfglyphtounicode{merkhalefthebrew}{05A5}
\pdfglyphtounicode{mhook}{0271}
\pdfglyphtounicode{mhzsquare}{3392}
\pdfglyphtounicode{middledotkatakanahalfwidth}{FF65}
\pdfglyphtounicode{middot}{00B7}
\pdfglyphtounicode{mieumacirclekorean}{3272}
\pdfglyphtounicode{mieumaparenkorean}{3212}
\pdfglyphtounicode{mieumcirclekorean}{3264}
\pdfglyphtounicode{mieumkorean}{3141}
\pdfglyphtounicode{mieumpansioskorean}{3170}
\pdfglyphtounicode{mieumparenkorean}{3204}
\pdfglyphtounicode{mieumpieupkorean}{316E}
\pdfglyphtounicode{mieumsioskorean}{316F}
\pdfglyphtounicode{mihiragana}{307F}
\pdfglyphtounicode{mikatakana}{30DF}
\pdfglyphtounicode{mikatakanahalfwidth}{FF90}
\pdfglyphtounicode{minus}{2212}
\pdfglyphtounicode{minusbelowcmb}{0320}
\pdfglyphtounicode{minuscircle}{2296}
\pdfglyphtounicode{minusmod}{02D7}
\pdfglyphtounicode{minusplus}{2213}
\pdfglyphtounicode{minute}{2032}
\pdfglyphtounicode{miribaarusquare}{334A}
\pdfglyphtounicode{mirisquare}{3349}
\pdfglyphtounicode{mlonglegturned}{0270}
\pdfglyphtounicode{mlsquare}{3396}
\pdfglyphtounicode{mmcubedsquare}{33A3}
\pdfglyphtounicode{mmonospace}{FF4D}
\pdfglyphtounicode{mmsquaredsquare}{339F}
\pdfglyphtounicode{mohiragana}{3082}
\pdfglyphtounicode{mohmsquare}{33C1}
\pdfglyphtounicode{mokatakana}{30E2}
\pdfglyphtounicode{mokatakanahalfwidth}{FF93}
\pdfglyphtounicode{molsquare}{33D6}
\pdfglyphtounicode{momathai}{0E21}
\pdfglyphtounicode{moverssquare}{33A7}
\pdfglyphtounicode{moverssquaredsquare}{33A8}
\pdfglyphtounicode{mparen}{24A8}
\pdfglyphtounicode{mpasquare}{33AB}
\pdfglyphtounicode{mssquare}{33B3}
\pdfglyphtounicode{msuperior}{006D}
\pdfglyphtounicode{mturned}{026F}
\pdfglyphtounicode{mu}{00B5}
\pdfglyphtounicode{mu1}{00B5}
\pdfglyphtounicode{muasquare}{3382}
\pdfglyphtounicode{muchgreater}{226B}
\pdfglyphtounicode{muchless}{226A}
\pdfglyphtounicode{mufsquare}{338C}
\pdfglyphtounicode{mugreek}{03BC}
\pdfglyphtounicode{mugsquare}{338D}
\pdfglyphtounicode{muhiragana}{3080}
\pdfglyphtounicode{mukatakana}{30E0}
\pdfglyphtounicode{mukatakanahalfwidth}{FF91}
\pdfglyphtounicode{mulsquare}{3395}
\pdfglyphtounicode{multicloseleft}{22C9}
\pdfglyphtounicode{multicloseright}{22CA}
\pdfglyphtounicode{multimap}{22B8}
\pdfglyphtounicode{multiopenleft}{22CB}
\pdfglyphtounicode{multiopenright}{22CC}
\pdfglyphtounicode{multiply}{00D7}
\pdfglyphtounicode{mumsquare}{339B}
\pdfglyphtounicode{munahhebrew}{05A3}
\pdfglyphtounicode{munahlefthebrew}{05A3}
\pdfglyphtounicode{musicalnote}{266A}
\pdfglyphtounicode{musicalnotedbl}{266B}
\pdfglyphtounicode{musicflatsign}{266D}
\pdfglyphtounicode{musicsharpsign}{266F}
\pdfglyphtounicode{mussquare}{33B2}
\pdfglyphtounicode{muvsquare}{33B6}
\pdfglyphtounicode{muwsquare}{33BC}
\pdfglyphtounicode{mvmegasquare}{33B9}
\pdfglyphtounicode{mvsquare}{33B7}
\pdfglyphtounicode{mwmegasquare}{33BF}
\pdfglyphtounicode{mwsquare}{33BD}
\pdfglyphtounicode{n}{006E}
\pdfglyphtounicode{nabengali}{09A8}
\pdfglyphtounicode{nabla}{2207}
\pdfglyphtounicode{nacute}{0144}
\pdfglyphtounicode{nadeva}{0928}
\pdfglyphtounicode{nagujarati}{0AA8}
\pdfglyphtounicode{nagurmukhi}{0A28}
\pdfglyphtounicode{nahiragana}{306A}
\pdfglyphtounicode{nakatakana}{30CA}
\pdfglyphtounicode{nakatakanahalfwidth}{FF85}
\pdfglyphtounicode{nand}{22BC}
\pdfglyphtounicode{napostrophe}{0149}
\pdfglyphtounicode{nasquare}{3381}
\pdfglyphtounicode{natural}{266E}
\pdfglyphtounicode{nbopomofo}{310B}
\pdfglyphtounicode{nbspace}{00A0}
\pdfglyphtounicode{ncaron}{0148}
\pdfglyphtounicode{ncedilla}{0146}
\pdfglyphtounicode{ncircle}{24DD}
\pdfglyphtounicode{ncircumflexbelow}{1E4B}
\pdfglyphtounicode{ncommaaccent}{0146}
\pdfglyphtounicode{ndotaccent}{1E45}
\pdfglyphtounicode{ndotbelow}{1E47}
\pdfglyphtounicode{negationslash}{0338}
\pdfglyphtounicode{nehiragana}{306D}
\pdfglyphtounicode{nekatakana}{30CD}
\pdfglyphtounicode{nekatakanahalfwidth}{FF88}
\pdfglyphtounicode{newsheqelsign}{20AA}
\pdfglyphtounicode{nfsquare}{338B}
\pdfglyphtounicode{ng}{014B}
\pdfglyphtounicode{ngabengali}{0999}
\pdfglyphtounicode{ngadeva}{0919}
\pdfglyphtounicode{ngagujarati}{0A99}
\pdfglyphtounicode{ngagurmukhi}{0A19}
\pdfglyphtounicode{ngonguthai}{0E07}
\pdfglyphtounicode{nhiragana}{3093}
\pdfglyphtounicode{nhookleft}{0272}
\pdfglyphtounicode{nhookretroflex}{0273}
\pdfglyphtounicode{nieunacirclekorean}{326F}
\pdfglyphtounicode{nieunaparenkorean}{320F}
\pdfglyphtounicode{nieuncieuckorean}{3135}
\pdfglyphtounicode{nieuncirclekorean}{3261}
\pdfglyphtounicode{nieunhieuhkorean}{3136}
\pdfglyphtounicode{nieunkorean}{3134}
\pdfglyphtounicode{nieunpansioskorean}{3168}
\pdfglyphtounicode{nieunparenkorean}{3201}
\pdfglyphtounicode{nieunsioskorean}{3167}
\pdfglyphtounicode{nieuntikeutkorean}{3166}
\pdfglyphtounicode{nihiragana}{306B}
\pdfglyphtounicode{nikatakana}{30CB}
\pdfglyphtounicode{nikatakanahalfwidth}{FF86}
\pdfglyphtounicode{nikhahitleftthai}{F899}
\pdfglyphtounicode{nikhahitthai}{0E4D}
\pdfglyphtounicode{nine}{0039}
\pdfglyphtounicode{ninearabic}{0669}
\pdfglyphtounicode{ninebengali}{09EF}
\pdfglyphtounicode{ninecircle}{2468}
\pdfglyphtounicode{ninecircleinversesansserif}{2792}
\pdfglyphtounicode{ninedeva}{096F}
\pdfglyphtounicode{ninegujarati}{0AEF}
\pdfglyphtounicode{ninegurmukhi}{0A6F}
\pdfglyphtounicode{ninehackarabic}{0669}
\pdfglyphtounicode{ninehangzhou}{3029}
\pdfglyphtounicode{nineideographicparen}{3228}
\pdfglyphtounicode{nineinferior}{2089}
\pdfglyphtounicode{ninemonospace}{FF19}
\pdfglyphtounicode{nineoldstyle}{0039}
\pdfglyphtounicode{nineparen}{247C}
\pdfglyphtounicode{nineperiod}{2490}
\pdfglyphtounicode{ninepersian}{06F9}
\pdfglyphtounicode{nineroman}{2178}
\pdfglyphtounicode{ninesuperior}{2079}
\pdfglyphtounicode{nineteencircle}{2472}
\pdfglyphtounicode{nineteenparen}{2486}
\pdfglyphtounicode{nineteenperiod}{249A}
\pdfglyphtounicode{ninethai}{0E59}
\pdfglyphtounicode{nj}{01CC}
\pdfglyphtounicode{njecyrillic}{045A}
\pdfglyphtounicode{nkatakana}{30F3}
\pdfglyphtounicode{nkatakanahalfwidth}{FF9D}
\pdfglyphtounicode{nlegrightlong}{019E}
\pdfglyphtounicode{nlinebelow}{1E49}
\pdfglyphtounicode{nmonospace}{FF4E}
\pdfglyphtounicode{nmsquare}{339A}
\pdfglyphtounicode{nnabengali}{09A3}
\pdfglyphtounicode{nnadeva}{0923}
\pdfglyphtounicode{nnagujarati}{0AA3}
\pdfglyphtounicode{nnagurmukhi}{0A23}
\pdfglyphtounicode{nnnadeva}{0929}
\pdfglyphtounicode{nohiragana}{306E}
\pdfglyphtounicode{nokatakana}{30CE}
\pdfglyphtounicode{nokatakanahalfwidth}{FF89}
\pdfglyphtounicode{nonbreakingspace}{00A0}
\pdfglyphtounicode{nonenthai}{0E13}
\pdfglyphtounicode{nonuthai}{0E19}
\pdfglyphtounicode{noonarabic}{0646}
\pdfglyphtounicode{noonfinalarabic}{FEE6}
\pdfglyphtounicode{noonghunnaarabic}{06BA}
\pdfglyphtounicode{noonghunnafinalarabic}{FB9F}
\pdfglyphtounicode{noonhehinitialarabic}{FEE7 FEEC}
\pdfglyphtounicode{nooninitialarabic}{FEE7}
\pdfglyphtounicode{noonjeeminitialarabic}{FCD2}
\pdfglyphtounicode{noonjeemisolatedarabic}{FC4B}
\pdfglyphtounicode{noonmedialarabic}{FEE8}
\pdfglyphtounicode{noonmeeminitialarabic}{FCD5}
\pdfglyphtounicode{noonmeemisolatedarabic}{FC4E}
\pdfglyphtounicode{noonnoonfinalarabic}{FC8D}
\pdfglyphtounicode{notapproxequal}{2247}
\pdfglyphtounicode{notarrowboth}{21AE}
\pdfglyphtounicode{notarrowleft}{219A}
\pdfglyphtounicode{notarrowright}{219B}
\pdfglyphtounicode{notbar}{2224}
\pdfglyphtounicode{notcontains}{220C}
\pdfglyphtounicode{notdblarrowboth}{21CE}
\pdfglyphtounicode{notdblarrowleft}{21CD}
\pdfglyphtounicode{notdblarrowright}{21CF}
\pdfglyphtounicode{notelement}{2209}
\pdfglyphtounicode{notelementof}{2209}
\pdfglyphtounicode{notequal}{2260}
\pdfglyphtounicode{notexistential}{2204}
\pdfglyphtounicode{notfollows}{2281}
\pdfglyphtounicode{notfollowsoreql}{2AB0 0338}
\pdfglyphtounicode{notforces}{22AE}
\pdfglyphtounicode{notforcesextra}{22AF}
\pdfglyphtounicode{notgreater}{226F}
\pdfglyphtounicode{notgreaterdblequal}{2267 0338}
\pdfglyphtounicode{notgreaterequal}{2271}
\pdfglyphtounicode{notgreaternorequal}{2271}
\pdfglyphtounicode{notgreaternorless}{2279}
\pdfglyphtounicode{notgreaterorslnteql}{2A7E 0338}
\pdfglyphtounicode{notidentical}{2262}
\pdfglyphtounicode{notless}{226E}
\pdfglyphtounicode{notlessdblequal}{2266 0338}
\pdfglyphtounicode{notlessequal}{2270}
\pdfglyphtounicode{notlessnorequal}{2270}
\pdfglyphtounicode{notlessorslnteql}{2A7D 0338}
\pdfglyphtounicode{notparallel}{2226}
\pdfglyphtounicode{notprecedes}{2280}
\pdfglyphtounicode{notprecedesoreql}{2AAF 0338}
\pdfglyphtounicode{notsatisfies}{22AD}
\pdfglyphtounicode{notsimilar}{2241}
\pdfglyphtounicode{notsubset}{2284}
\pdfglyphtounicode{notsubseteql}{2288}
\pdfglyphtounicode{notsubsetordbleql}{2AC5 0338}
\pdfglyphtounicode{notsubsetoreql}{228A}
\pdfglyphtounicode{notsucceeds}{2281}
\pdfglyphtounicode{notsuperset}{2285}
\pdfglyphtounicode{notsuperseteql}{2289}
\pdfglyphtounicode{notsupersetordbleql}{2AC6 0338}
\pdfglyphtounicode{notsupersetoreql}{228B}
\pdfglyphtounicode{nottriangeqlleft}{22EC}
\pdfglyphtounicode{nottriangeqlright}{22ED}
\pdfglyphtounicode{nottriangleleft}{22EA}
\pdfglyphtounicode{nottriangleright}{22EB}
\pdfglyphtounicode{notturnstile}{22AC}
\pdfglyphtounicode{nowarmenian}{0576}
\pdfglyphtounicode{nparen}{24A9}
\pdfglyphtounicode{nssquare}{33B1}
\pdfglyphtounicode{nsuperior}{207F}
\pdfglyphtounicode{ntilde}{00F1}
\pdfglyphtounicode{nu}{03BD}
\pdfglyphtounicode{nuhiragana}{306C}
\pdfglyphtounicode{nukatakana}{30CC}
\pdfglyphtounicode{nukatakanahalfwidth}{FF87}
\pdfglyphtounicode{nuktabengali}{09BC}
\pdfglyphtounicode{nuktadeva}{093C}
\pdfglyphtounicode{nuktagujarati}{0ABC}
\pdfglyphtounicode{nuktagurmukhi}{0A3C}
\pdfglyphtounicode{numbersign}{0023}
\pdfglyphtounicode{numbersignmonospace}{FF03}
\pdfglyphtounicode{numbersignsmall}{FE5F}
\pdfglyphtounicode{numeralsigngreek}{0374}
\pdfglyphtounicode{numeralsignlowergreek}{0375}
\pdfglyphtounicode{numero}{2116}
\pdfglyphtounicode{nun}{05E0}
\pdfglyphtounicode{nundagesh}{FB40}
\pdfglyphtounicode{nundageshhebrew}{FB40}
\pdfglyphtounicode{nunhebrew}{05E0}
\pdfglyphtounicode{nvsquare}{33B5}
\pdfglyphtounicode{nwsquare}{33BB}
\pdfglyphtounicode{nyabengali}{099E}
\pdfglyphtounicode{nyadeva}{091E}
\pdfglyphtounicode{nyagujarati}{0A9E}
\pdfglyphtounicode{nyagurmukhi}{0A1E}
\pdfglyphtounicode{o}{006F}
\pdfglyphtounicode{oacute}{00F3}
\pdfglyphtounicode{oangthai}{0E2D}
\pdfglyphtounicode{obarred}{0275}
\pdfglyphtounicode{obarredcyrillic}{04E9}
\pdfglyphtounicode{obarreddieresiscyrillic}{04EB}
\pdfglyphtounicode{obengali}{0993}
\pdfglyphtounicode{obopomofo}{311B}
\pdfglyphtounicode{obreve}{014F}
\pdfglyphtounicode{ocandradeva}{0911}
\pdfglyphtounicode{ocandragujarati}{0A91}
\pdfglyphtounicode{ocandravowelsigndeva}{0949}
\pdfglyphtounicode{ocandravowelsigngujarati}{0AC9}
\pdfglyphtounicode{ocaron}{01D2}
\pdfglyphtounicode{ocircle}{24DE}
\pdfglyphtounicode{ocircumflex}{00F4}
\pdfglyphtounicode{ocircumflexacute}{1ED1}
\pdfglyphtounicode{ocircumflexdotbelow}{1ED9}
\pdfglyphtounicode{ocircumflexgrave}{1ED3}
\pdfglyphtounicode{ocircumflexhookabove}{1ED5}
\pdfglyphtounicode{ocircumflextilde}{1ED7}
\pdfglyphtounicode{ocyrillic}{043E}
\pdfglyphtounicode{odblacute}{0151}
\pdfglyphtounicode{odblgrave}{020D}
\pdfglyphtounicode{odeva}{0913}
\pdfglyphtounicode{odieresis}{00F6}
\pdfglyphtounicode{odieresiscyrillic}{04E7}
\pdfglyphtounicode{odotbelow}{1ECD}
\pdfglyphtounicode{oe}{0153}
\pdfglyphtounicode{oekorean}{315A}
\pdfglyphtounicode{ogonek}{02DB}
\pdfglyphtounicode{ogonekcmb}{0328}
\pdfglyphtounicode{ograve}{00F2}
\pdfglyphtounicode{ogujarati}{0A93}
\pdfglyphtounicode{oharmenian}{0585}
\pdfglyphtounicode{ohiragana}{304A}
\pdfglyphtounicode{ohookabove}{1ECF}
\pdfglyphtounicode{ohorn}{01A1}
\pdfglyphtounicode{ohornacute}{1EDB}
\pdfglyphtounicode{ohorndotbelow}{1EE3}
\pdfglyphtounicode{ohorngrave}{1EDD}
\pdfglyphtounicode{ohornhookabove}{1EDF}
\pdfglyphtounicode{ohorntilde}{1EE1}
\pdfglyphtounicode{ohungarumlaut}{0151}
\pdfglyphtounicode{oi}{01A3}
\pdfglyphtounicode{oinvertedbreve}{020F}
\pdfglyphtounicode{okatakana}{30AA}
\pdfglyphtounicode{okatakanahalfwidth}{FF75}
\pdfglyphtounicode{okorean}{3157}
\pdfglyphtounicode{olehebrew}{05AB}
\pdfglyphtounicode{omacron}{014D}
\pdfglyphtounicode{omacronacute}{1E53}
\pdfglyphtounicode{omacrongrave}{1E51}
\pdfglyphtounicode{omdeva}{0950}
\pdfglyphtounicode{omega}{03C9}
\pdfglyphtounicode{omega1}{03D6}
\pdfglyphtounicode{omegacyrillic}{0461}
\pdfglyphtounicode{omegalatinclosed}{0277}
\pdfglyphtounicode{omegaroundcyrillic}{047B}
\pdfglyphtounicode{omegatitlocyrillic}{047D}
\pdfglyphtounicode{omegatonos}{03CE}
\pdfglyphtounicode{omgujarati}{0AD0}
\pdfglyphtounicode{omicron}{03BF}
\pdfglyphtounicode{omicrontonos}{03CC}
\pdfglyphtounicode{omonospace}{FF4F}
\pdfglyphtounicode{one}{0031}
\pdfglyphtounicode{onearabic}{0661}
\pdfglyphtounicode{onebengali}{09E7}
\pdfglyphtounicode{onecircle}{2460}
\pdfglyphtounicode{onecircleinversesansserif}{278A}
\pdfglyphtounicode{onedeva}{0967}
\pdfglyphtounicode{onedotenleader}{2024}
\pdfglyphtounicode{oneeighth}{215B}
\pdfglyphtounicode{onefitted}{0031}
\pdfglyphtounicode{onegujarati}{0AE7}
\pdfglyphtounicode{onegurmukhi}{0A67}
\pdfglyphtounicode{onehackarabic}{0661}
\pdfglyphtounicode{onehalf}{00BD}
\pdfglyphtounicode{onehangzhou}{3021}
\pdfglyphtounicode{oneideographicparen}{3220}
\pdfglyphtounicode{oneinferior}{2081}
\pdfglyphtounicode{onemonospace}{FF11}
\pdfglyphtounicode{onenumeratorbengali}{09F4}
\pdfglyphtounicode{oneoldstyle}{0031}
\pdfglyphtounicode{oneparen}{2474}
\pdfglyphtounicode{oneperiod}{2488}
\pdfglyphtounicode{onepersian}{06F1}
\pdfglyphtounicode{onequarter}{00BC}
\pdfglyphtounicode{oneroman}{2170}
\pdfglyphtounicode{onesuperior}{00B9}
\pdfglyphtounicode{onethai}{0E51}
\pdfglyphtounicode{onethird}{2153}
\pdfglyphtounicode{oogonek}{01EB}
\pdfglyphtounicode{oogonekmacron}{01ED}
\pdfglyphtounicode{oogurmukhi}{0A13}
\pdfglyphtounicode{oomatragurmukhi}{0A4B}
\pdfglyphtounicode{oopen}{0254}
\pdfglyphtounicode{oparen}{24AA}
\pdfglyphtounicode{openbullet}{25E6}
\pdfglyphtounicode{option}{2325}
\pdfglyphtounicode{ordfeminine}{00AA}
\pdfglyphtounicode{ordmasculine}{00BA}
\pdfglyphtounicode{orthogonal}{221F}
\pdfglyphtounicode{orunderscore}{22BB}
\pdfglyphtounicode{oshortdeva}{0912}
\pdfglyphtounicode{oshortvowelsigndeva}{094A}
\pdfglyphtounicode{oslash}{00F8}
\pdfglyphtounicode{oslashacute}{01FF}
\pdfglyphtounicode{osmallhiragana}{3049}
\pdfglyphtounicode{osmallkatakana}{30A9}
\pdfglyphtounicode{osmallkatakanahalfwidth}{FF6B}
\pdfglyphtounicode{ostrokeacute}{01FF}
\pdfglyphtounicode{osuperior}{006F}
\pdfglyphtounicode{otcyrillic}{047F}
\pdfglyphtounicode{otilde}{00F5}
\pdfglyphtounicode{otildeacute}{1E4D}
\pdfglyphtounicode{otildedieresis}{1E4F}
\pdfglyphtounicode{oubopomofo}{3121}
\pdfglyphtounicode{overline}{203E}
\pdfglyphtounicode{overlinecenterline}{FE4A}
\pdfglyphtounicode{overlinecmb}{0305}
\pdfglyphtounicode{overlinedashed}{FE49}
\pdfglyphtounicode{overlinedblwavy}{FE4C}
\pdfglyphtounicode{overlinewavy}{FE4B}
\pdfglyphtounicode{overscore}{00AF}
\pdfglyphtounicode{ovowelsignbengali}{09CB}
\pdfglyphtounicode{ovowelsigndeva}{094B}
\pdfglyphtounicode{ovowelsigngujarati}{0ACB}
\pdfglyphtounicode{owner}{220B}
\pdfglyphtounicode{p}{0070}
\pdfglyphtounicode{paampssquare}{3380}
\pdfglyphtounicode{paasentosquare}{332B}
\pdfglyphtounicode{pabengali}{09AA}
\pdfglyphtounicode{pacute}{1E55}
\pdfglyphtounicode{padeva}{092A}
\pdfglyphtounicode{pagedown}{21DF}
\pdfglyphtounicode{pageup}{21DE}
\pdfglyphtounicode{pagujarati}{0AAA}
\pdfglyphtounicode{pagurmukhi}{0A2A}
\pdfglyphtounicode{pahiragana}{3071}
\pdfglyphtounicode{paiyannoithai}{0E2F}
\pdfglyphtounicode{pakatakana}{30D1}
\pdfglyphtounicode{palatalizationcyrilliccmb}{0484}
\pdfglyphtounicode{palochkacyrillic}{04C0}
\pdfglyphtounicode{pansioskorean}{317F}
\pdfglyphtounicode{paragraph}{00B6}
\pdfglyphtounicode{parallel}{2225}
\pdfglyphtounicode{parenleft}{0028}
\pdfglyphtounicode{parenleftaltonearabic}{FD3E}
\pdfglyphtounicode{parenleftbt}{F8ED}
\pdfglyphtounicode{parenleftex}{F8EC}
\pdfglyphtounicode{parenleftinferior}{208D}
\pdfglyphtounicode{parenleftmonospace}{FF08}
\pdfglyphtounicode{parenleftsmall}{FE59}
\pdfglyphtounicode{parenleftsuperior}{207D}
\pdfglyphtounicode{parenlefttp}{F8EB}
\pdfglyphtounicode{parenleftvertical}{FE35}
\pdfglyphtounicode{parenright}{0029}
\pdfglyphtounicode{parenrightaltonearabic}{FD3F}
\pdfglyphtounicode{parenrightbt}{F8F8}
\pdfglyphtounicode{parenrightex}{F8F7}
\pdfglyphtounicode{parenrightinferior}{208E}
\pdfglyphtounicode{parenrightmonospace}{FF09}
\pdfglyphtounicode{parenrightsmall}{FE5A}
\pdfglyphtounicode{parenrightsuperior}{207E}
\pdfglyphtounicode{parenrighttp}{F8F6}
\pdfglyphtounicode{parenrightvertical}{FE36}
\pdfglyphtounicode{partialdiff}{2202}
\pdfglyphtounicode{paseqhebrew}{05C0}
\pdfglyphtounicode{pashtahebrew}{0599}
\pdfglyphtounicode{pasquare}{33A9}
\pdfglyphtounicode{patah}{05B7}
\pdfglyphtounicode{patah11}{05B7}
\pdfglyphtounicode{patah1d}{05B7}
\pdfglyphtounicode{patah2a}{05B7}
\pdfglyphtounicode{patahhebrew}{05B7}
\pdfglyphtounicode{patahnarrowhebrew}{05B7}
\pdfglyphtounicode{patahquarterhebrew}{05B7}
\pdfglyphtounicode{patahwidehebrew}{05B7}
\pdfglyphtounicode{pazerhebrew}{05A1}
\pdfglyphtounicode{pbopomofo}{3106}
\pdfglyphtounicode{pcircle}{24DF}
\pdfglyphtounicode{pdotaccent}{1E57}
\pdfglyphtounicode{pe}{05E4}
\pdfglyphtounicode{pecyrillic}{043F}
\pdfglyphtounicode{pedagesh}{FB44}
\pdfglyphtounicode{pedageshhebrew}{FB44}
\pdfglyphtounicode{peezisquare}{333B}
\pdfglyphtounicode{pefinaldageshhebrew}{FB43}
\pdfglyphtounicode{peharabic}{067E}
\pdfglyphtounicode{peharmenian}{057A}
\pdfglyphtounicode{pehebrew}{05E4}
\pdfglyphtounicode{pehfinalarabic}{FB57}
\pdfglyphtounicode{pehinitialarabic}{FB58}
\pdfglyphtounicode{pehiragana}{307A}
\pdfglyphtounicode{pehmedialarabic}{FB59}
\pdfglyphtounicode{pekatakana}{30DA}
\pdfglyphtounicode{pemiddlehookcyrillic}{04A7}
\pdfglyphtounicode{perafehebrew}{FB4E}
\pdfglyphtounicode{percent}{0025}
\pdfglyphtounicode{percentarabic}{066A}
\pdfglyphtounicode{percentmonospace}{FF05}
\pdfglyphtounicode{percentsmall}{FE6A}
\pdfglyphtounicode{period}{002E}
\pdfglyphtounicode{periodarmenian}{0589}
\pdfglyphtounicode{periodcentered}{00B7}
\pdfglyphtounicode{periodhalfwidth}{FF61}
\pdfglyphtounicode{periodinferior}{002E}
\pdfglyphtounicode{periodmonospace}{FF0E}
\pdfglyphtounicode{periodsmall}{FE52}
\pdfglyphtounicode{periodsuperior}{002E}
\pdfglyphtounicode{perispomenigreekcmb}{0342}
\pdfglyphtounicode{perpcorrespond}{2A5E}
\pdfglyphtounicode{perpendicular}{22A5}
\pdfglyphtounicode{pertenthousand}{2031}
\pdfglyphtounicode{perthousand}{2030}
\pdfglyphtounicode{peseta}{20A7}
\pdfglyphtounicode{pfsquare}{338A}
\pdfglyphtounicode{phabengali}{09AB}
\pdfglyphtounicode{phadeva}{092B}
\pdfglyphtounicode{phagujarati}{0AAB}
\pdfglyphtounicode{phagurmukhi}{0A2B}
\pdfglyphtounicode{phi}{03C6}
\pdfglyphtounicode{phi1}{03D5}
\pdfglyphtounicode{phieuphacirclekorean}{327A}
\pdfglyphtounicode{phieuphaparenkorean}{321A}
\pdfglyphtounicode{phieuphcirclekorean}{326C}
\pdfglyphtounicode{phieuphkorean}{314D}
\pdfglyphtounicode{phieuphparenkorean}{320C}
\pdfglyphtounicode{philatin}{0278}
\pdfglyphtounicode{phinthuthai}{0E3A}
\pdfglyphtounicode{phisymbolgreek}{03D5}
\pdfglyphtounicode{phook}{01A5}
\pdfglyphtounicode{phophanthai}{0E1E}
\pdfglyphtounicode{phophungthai}{0E1C}
\pdfglyphtounicode{phosamphaothai}{0E20}
\pdfglyphtounicode{pi}{03C0}
\pdfglyphtounicode{pi1}{03D6}
\pdfglyphtounicode{pieupacirclekorean}{3273}
\pdfglyphtounicode{pieupaparenkorean}{3213}
\pdfglyphtounicode{pieupcieuckorean}{3176}
\pdfglyphtounicode{pieupcirclekorean}{3265}
\pdfglyphtounicode{pieupkiyeokkorean}{3172}
\pdfglyphtounicode{pieupkorean}{3142}
\pdfglyphtounicode{pieupparenkorean}{3205}
\pdfglyphtounicode{pieupsioskiyeokkorean}{3174}
\pdfglyphtounicode{pieupsioskorean}{3144}
\pdfglyphtounicode{pieupsiostikeutkorean}{3175}
\pdfglyphtounicode{pieupthieuthkorean}{3177}
\pdfglyphtounicode{pieuptikeutkorean}{3173}
\pdfglyphtounicode{pihiragana}{3074}
\pdfglyphtounicode{pikatakana}{30D4}
\pdfglyphtounicode{pisymbolgreek}{03D6}
\pdfglyphtounicode{piwrarmenian}{0583}
\pdfglyphtounicode{planckover2pi}{210F}
\pdfglyphtounicode{planckover2pi1}{210F}
\pdfglyphtounicode{plus}{002B}
\pdfglyphtounicode{plusbelowcmb}{031F}
\pdfglyphtounicode{pluscircle}{2295}
\pdfglyphtounicode{plusminus}{00B1}
\pdfglyphtounicode{plusmod}{02D6}
\pdfglyphtounicode{plusmonospace}{FF0B}
\pdfglyphtounicode{plussmall}{FE62}
\pdfglyphtounicode{plussuperior}{207A}
\pdfglyphtounicode{pmonospace}{FF50}
\pdfglyphtounicode{pmsquare}{33D8}
\pdfglyphtounicode{pohiragana}{307D}
\pdfglyphtounicode{pointingindexdownwhite}{261F}
\pdfglyphtounicode{pointingindexleftwhite}{261C}
\pdfglyphtounicode{pointingindexrightwhite}{261E}
\pdfglyphtounicode{pointingindexupwhite}{261D}
\pdfglyphtounicode{pokatakana}{30DD}
\pdfglyphtounicode{poplathai}{0E1B}
\pdfglyphtounicode{postalmark}{3012}
\pdfglyphtounicode{postalmarkface}{3020}
\pdfglyphtounicode{pparen}{24AB}
\pdfglyphtounicode{precedenotdbleqv}{2AB9}
\pdfglyphtounicode{precedenotslnteql}{2AB5}
\pdfglyphtounicode{precedeornoteqvlnt}{22E8}
\pdfglyphtounicode{precedes}{227A}
\pdfglyphtounicode{precedesequal}{2AAF}
\pdfglyphtounicode{precedesorcurly}{227C}
\pdfglyphtounicode{precedesorequal}{227E}
\pdfglyphtounicode{prescription}{211E}
\pdfglyphtounicode{prime}{2032}
\pdfglyphtounicode{primemod}{02B9}
\pdfglyphtounicode{primereverse}{2035}
\pdfglyphtounicode{primereversed}{2035}
\pdfglyphtounicode{product}{220F}
\pdfglyphtounicode{projective}{2305}
\pdfglyphtounicode{prolongedkana}{30FC}
\pdfglyphtounicode{propellor}{2318}
\pdfglyphtounicode{propersubset}{2282}
\pdfglyphtounicode{propersuperset}{2283}
\pdfglyphtounicode{proportion}{2237}
\pdfglyphtounicode{proportional}{221D}
\pdfglyphtounicode{psi}{03C8}
\pdfglyphtounicode{psicyrillic}{0471}
\pdfglyphtounicode{psilipneumatacyrilliccmb}{0486}
\pdfglyphtounicode{pssquare}{33B0}
\pdfglyphtounicode{puhiragana}{3077}
\pdfglyphtounicode{pukatakana}{30D7}
\pdfglyphtounicode{punctdash}{2014}
\pdfglyphtounicode{pvsquare}{33B4}
\pdfglyphtounicode{pwsquare}{33BA}
\pdfglyphtounicode{q}{0071}
\pdfglyphtounicode{qadeva}{0958}
\pdfglyphtounicode{qadmahebrew}{05A8}
\pdfglyphtounicode{qafarabic}{0642}
\pdfglyphtounicode{qaffinalarabic}{FED6}
\pdfglyphtounicode{qafinitialarabic}{FED7}
\pdfglyphtounicode{qafmedialarabic}{FED8}
\pdfglyphtounicode{qamats}{05B8}
\pdfglyphtounicode{qamats10}{05B8}
\pdfglyphtounicode{qamats1a}{05B8}
\pdfglyphtounicode{qamats1c}{05B8}
\pdfglyphtounicode{qamats27}{05B8}
\pdfglyphtounicode{qamats29}{05B8}
\pdfglyphtounicode{qamats33}{05B8}
\pdfglyphtounicode{qamatsde}{05B8}
\pdfglyphtounicode{qamatshebrew}{05B8}
\pdfglyphtounicode{qamatsnarrowhebrew}{05B8}
\pdfglyphtounicode{qamatsqatanhebrew}{05B8}
\pdfglyphtounicode{qamatsqatannarrowhebrew}{05B8}
\pdfglyphtounicode{qamatsqatanquarterhebrew}{05B8}
\pdfglyphtounicode{qamatsqatanwidehebrew}{05B8}
\pdfglyphtounicode{qamatsquarterhebrew}{05B8}
\pdfglyphtounicode{qamatswidehebrew}{05B8}
\pdfglyphtounicode{qarneyparahebrew}{059F}
\pdfglyphtounicode{qbopomofo}{3111}
\pdfglyphtounicode{qcircle}{24E0}
\pdfglyphtounicode{qhook}{02A0}
\pdfglyphtounicode{qmonospace}{FF51}
\pdfglyphtounicode{qof}{05E7}
\pdfglyphtounicode{qofdagesh}{FB47}
\pdfglyphtounicode{qofdageshhebrew}{FB47}
\pdfglyphtounicode{qofhatafpatah}{05E7 05B2}
\pdfglyphtounicode{qofhatafpatahhebrew}{05E7 05B2}
\pdfglyphtounicode{qofhatafsegol}{05E7 05B1}
\pdfglyphtounicode{qofhatafsegolhebrew}{05E7 05B1}
\pdfglyphtounicode{qofhebrew}{05E7}
\pdfglyphtounicode{qofhiriq}{05E7 05B4}
\pdfglyphtounicode{qofhiriqhebrew}{05E7 05B4}
\pdfglyphtounicode{qofholam}{05E7 05B9}
\pdfglyphtounicode{qofholamhebrew}{05E7 05B9}
\pdfglyphtounicode{qofpatah}{05E7 05B7}
\pdfglyphtounicode{qofpatahhebrew}{05E7 05B7}
\pdfglyphtounicode{qofqamats}{05E7 05B8}
\pdfglyphtounicode{qofqamatshebrew}{05E7 05B8}
\pdfglyphtounicode{qofqubuts}{05E7 05BB}
\pdfglyphtounicode{qofqubutshebrew}{05E7 05BB}
\pdfglyphtounicode{qofsegol}{05E7 05B6}
\pdfglyphtounicode{qofsegolhebrew}{05E7 05B6}
\pdfglyphtounicode{qofsheva}{05E7 05B0}
\pdfglyphtounicode{qofshevahebrew}{05E7 05B0}
\pdfglyphtounicode{qoftsere}{05E7 05B5}
\pdfglyphtounicode{qoftserehebrew}{05E7 05B5}
\pdfglyphtounicode{qparen}{24AC}
\pdfglyphtounicode{quarternote}{2669}
\pdfglyphtounicode{qubuts}{05BB}
\pdfglyphtounicode{qubuts18}{05BB}
\pdfglyphtounicode{qubuts25}{05BB}
\pdfglyphtounicode{qubuts31}{05BB}
\pdfglyphtounicode{qubutshebrew}{05BB}
\pdfglyphtounicode{qubutsnarrowhebrew}{05BB}
\pdfglyphtounicode{qubutsquarterhebrew}{05BB}
\pdfglyphtounicode{qubutswidehebrew}{05BB}
\pdfglyphtounicode{question}{003F}
\pdfglyphtounicode{questionarabic}{061F}
\pdfglyphtounicode{questionarmenian}{055E}
\pdfglyphtounicode{questiondown}{00BF}
\pdfglyphtounicode{questiondownsmall}{00BF}
\pdfglyphtounicode{questiongreek}{037E}
\pdfglyphtounicode{questionmonospace}{FF1F}
\pdfglyphtounicode{questionsmall}{003F}
\pdfglyphtounicode{quotedbl}{0022}
\pdfglyphtounicode{quotedblbase}{201E}
\pdfglyphtounicode{quotedblleft}{201C}
\pdfglyphtounicode{quotedblmonospace}{FF02}
\pdfglyphtounicode{quotedblprime}{301E}
\pdfglyphtounicode{quotedblprimereversed}{301D}
\pdfglyphtounicode{quotedblright}{201D}
\pdfglyphtounicode{quoteleft}{2018}
\pdfglyphtounicode{quoteleftreversed}{201B}
\pdfglyphtounicode{quotereversed}{201B}
\pdfglyphtounicode{quoteright}{2019}
\pdfglyphtounicode{quoterightn}{0149}
\pdfglyphtounicode{quotesinglbase}{201A}
\pdfglyphtounicode{quotesingle}{0027}
\pdfglyphtounicode{quotesinglemonospace}{FF07}
\pdfglyphtounicode{r}{0072}
\pdfglyphtounicode{raarmenian}{057C}
\pdfglyphtounicode{rabengali}{09B0}
\pdfglyphtounicode{racute}{0155}
\pdfglyphtounicode{radeva}{0930}
\pdfglyphtounicode{radical}{221A}
\pdfglyphtounicode{radicalex}{F8E5}
\pdfglyphtounicode{radoverssquare}{33AE}
\pdfglyphtounicode{radoverssquaredsquare}{33AF}
\pdfglyphtounicode{radsquare}{33AD}
\pdfglyphtounicode{rafe}{05BF}
\pdfglyphtounicode{rafehebrew}{05BF}
\pdfglyphtounicode{ragujarati}{0AB0}
\pdfglyphtounicode{ragurmukhi}{0A30}
\pdfglyphtounicode{rahiragana}{3089}
\pdfglyphtounicode{rakatakana}{30E9}
\pdfglyphtounicode{rakatakanahalfwidth}{FF97}
\pdfglyphtounicode{ralowerdiagonalbengali}{09F1}
\pdfglyphtounicode{ramiddlediagonalbengali}{09F0}
\pdfglyphtounicode{ramshorn}{0264}
\pdfglyphtounicode{rangedash}{2013}
\pdfglyphtounicode{ratio}{2236}
\pdfglyphtounicode{rbopomofo}{3116}
\pdfglyphtounicode{rcaron}{0159}
\pdfglyphtounicode{rcedilla}{0157}
\pdfglyphtounicode{rcircle}{24E1}
\pdfglyphtounicode{rcommaaccent}{0157}
\pdfglyphtounicode{rdblgrave}{0211}
\pdfglyphtounicode{rdotaccent}{1E59}
\pdfglyphtounicode{rdotbelow}{1E5B}
\pdfglyphtounicode{rdotbelowmacron}{1E5D}
\pdfglyphtounicode{referencemark}{203B}
\pdfglyphtounicode{reflexsubset}{2286}
\pdfglyphtounicode{reflexsuperset}{2287}
\pdfglyphtounicode{registered}{00AE}
\pdfglyphtounicode{registersans}{00AE}
\pdfglyphtounicode{registerserif}{00AE}
\pdfglyphtounicode{reharabic}{0631}
\pdfglyphtounicode{reharmenian}{0580}
\pdfglyphtounicode{rehfinalarabic}{FEAE}
\pdfglyphtounicode{rehiragana}{308C}
\pdfglyphtounicode{rehyehaleflamarabic}{0631 FEF3 FE8E 0644}
\pdfglyphtounicode{rekatakana}{30EC}
\pdfglyphtounicode{rekatakanahalfwidth}{FF9A}
\pdfglyphtounicode{resh}{05E8}
\pdfglyphtounicode{reshdageshhebrew}{FB48}
\pdfglyphtounicode{reshhatafpatah}{05E8 05B2}
\pdfglyphtounicode{reshhatafpatahhebrew}{05E8 05B2}
\pdfglyphtounicode{reshhatafsegol}{05E8 05B1}
\pdfglyphtounicode{reshhatafsegolhebrew}{05E8 05B1}
\pdfglyphtounicode{reshhebrew}{05E8}
\pdfglyphtounicode{reshhiriq}{05E8 05B4}
\pdfglyphtounicode{reshhiriqhebrew}{05E8 05B4}
\pdfglyphtounicode{reshholam}{05E8 05B9}
\pdfglyphtounicode{reshholamhebrew}{05E8 05B9}
\pdfglyphtounicode{reshpatah}{05E8 05B7}
\pdfglyphtounicode{reshpatahhebrew}{05E8 05B7}
\pdfglyphtounicode{reshqamats}{05E8 05B8}
\pdfglyphtounicode{reshqamatshebrew}{05E8 05B8}
\pdfglyphtounicode{reshqubuts}{05E8 05BB}
\pdfglyphtounicode{reshqubutshebrew}{05E8 05BB}
\pdfglyphtounicode{reshsegol}{05E8 05B6}
\pdfglyphtounicode{reshsegolhebrew}{05E8 05B6}
\pdfglyphtounicode{reshsheva}{05E8 05B0}
\pdfglyphtounicode{reshshevahebrew}{05E8 05B0}
\pdfglyphtounicode{reshtsere}{05E8 05B5}
\pdfglyphtounicode{reshtserehebrew}{05E8 05B5}
\pdfglyphtounicode{revasymptequal}{22CD}
\pdfglyphtounicode{reversedtilde}{223D}
\pdfglyphtounicode{reviahebrew}{0597}
\pdfglyphtounicode{reviamugrashhebrew}{0597}
\pdfglyphtounicode{revlogicalnot}{2310}
\pdfglyphtounicode{revsimilar}{223D}
\pdfglyphtounicode{rfishhook}{027E}
\pdfglyphtounicode{rfishhookreversed}{027F}
\pdfglyphtounicode{rhabengali}{09DD}
\pdfglyphtounicode{rhadeva}{095D}
\pdfglyphtounicode{rho}{03C1}
\pdfglyphtounicode{rho1}{03F1}
\pdfglyphtounicode{rhook}{027D}
\pdfglyphtounicode{rhookturned}{027B}
\pdfglyphtounicode{rhookturnedsuperior}{02B5}
\pdfglyphtounicode{rhosymbolgreek}{03F1}
\pdfglyphtounicode{rhotichookmod}{02DE}
\pdfglyphtounicode{rieulacirclekorean}{3271}
\pdfglyphtounicode{rieulaparenkorean}{3211}
\pdfglyphtounicode{rieulcirclekorean}{3263}
\pdfglyphtounicode{rieulhieuhkorean}{3140}
\pdfglyphtounicode{rieulkiyeokkorean}{313A}
\pdfglyphtounicode{rieulkiyeoksioskorean}{3169}
\pdfglyphtounicode{rieulkorean}{3139}
\pdfglyphtounicode{rieulmieumkorean}{313B}
\pdfglyphtounicode{rieulpansioskorean}{316C}
\pdfglyphtounicode{rieulparenkorean}{3203}
\pdfglyphtounicode{rieulphieuphkorean}{313F}
\pdfglyphtounicode{rieulpieupkorean}{313C}
\pdfglyphtounicode{rieulpieupsioskorean}{316B}
\pdfglyphtounicode{rieulsioskorean}{313D}
\pdfglyphtounicode{rieulthieuthkorean}{313E}
\pdfglyphtounicode{rieultikeutkorean}{316A}
\pdfglyphtounicode{rieulyeorinhieuhkorean}{316D}
\pdfglyphtounicode{rightangle}{221F}
\pdfglyphtounicode{rightanglene}{231D}
\pdfglyphtounicode{rightanglenw}{231C}
\pdfglyphtounicode{rightanglese}{231F}
\pdfglyphtounicode{rightanglesw}{231E}
\pdfglyphtounicode{righttackbelowcmb}{0319}
\pdfglyphtounicode{righttriangle}{22BF}
\pdfglyphtounicode{rihiragana}{308A}
\pdfglyphtounicode{rikatakana}{30EA}
\pdfglyphtounicode{rikatakanahalfwidth}{FF98}
\pdfglyphtounicode{ring}{02DA}
\pdfglyphtounicode{ringbelowcmb}{0325}
\pdfglyphtounicode{ringcmb}{030A}
\pdfglyphtounicode{ringhalfleft}{02BF}
\pdfglyphtounicode{ringhalfleftarmenian}{0559}
\pdfglyphtounicode{ringhalfleftbelowcmb}{031C}
\pdfglyphtounicode{ringhalfleftcentered}{02D3}
\pdfglyphtounicode{ringhalfright}{02BE}
\pdfglyphtounicode{ringhalfrightbelowcmb}{0339}
\pdfglyphtounicode{ringhalfrightcentered}{02D2}
\pdfglyphtounicode{ringinequal}{2256}
\pdfglyphtounicode{rinvertedbreve}{0213}
\pdfglyphtounicode{rittorusquare}{3351}
\pdfglyphtounicode{rlinebelow}{1E5F}
\pdfglyphtounicode{rlongleg}{027C}
\pdfglyphtounicode{rlonglegturned}{027A}
\pdfglyphtounicode{rmonospace}{FF52}
\pdfglyphtounicode{rohiragana}{308D}
\pdfglyphtounicode{rokatakana}{30ED}
\pdfglyphtounicode{rokatakanahalfwidth}{FF9B}
\pdfglyphtounicode{roruathai}{0E23}
\pdfglyphtounicode{rparen}{24AD}
\pdfglyphtounicode{rrabengali}{09DC}
\pdfglyphtounicode{rradeva}{0931}
\pdfglyphtounicode{rragurmukhi}{0A5C}
\pdfglyphtounicode{rreharabic}{0691}
\pdfglyphtounicode{rrehfinalarabic}{FB8D}
\pdfglyphtounicode{rrvocalicbengali}{09E0}
\pdfglyphtounicode{rrvocalicdeva}{0960}
\pdfglyphtounicode{rrvocalicgujarati}{0AE0}
\pdfglyphtounicode{rrvocalicvowelsignbengali}{09C4}
\pdfglyphtounicode{rrvocalicvowelsigndeva}{0944}
\pdfglyphtounicode{rrvocalicvowelsigngujarati}{0AC4}
\pdfglyphtounicode{rsuperior}{0072}
\pdfglyphtounicode{rtblock}{2590}
\pdfglyphtounicode{rturned}{0279}
\pdfglyphtounicode{rturnedsuperior}{02B4}
\pdfglyphtounicode{ruhiragana}{308B}
\pdfglyphtounicode{rukatakana}{30EB}
\pdfglyphtounicode{rukatakanahalfwidth}{FF99}
\pdfglyphtounicode{rupeemarkbengali}{09F2}
\pdfglyphtounicode{rupeesignbengali}{09F3}
\pdfglyphtounicode{rupiah}{20A8}
\pdfglyphtounicode{ruthai}{0E24}
\pdfglyphtounicode{rvocalicbengali}{098B}
\pdfglyphtounicode{rvocalicdeva}{090B}
\pdfglyphtounicode{rvocalicgujarati}{0A8B}
\pdfglyphtounicode{rvocalicvowelsignbengali}{09C3}
\pdfglyphtounicode{rvocalicvowelsigndeva}{0943}
\pdfglyphtounicode{rvocalicvowelsigngujarati}{0AC3}
\pdfglyphtounicode{s}{0073}
\pdfglyphtounicode{sabengali}{09B8}
\pdfglyphtounicode{sacute}{015B}
\pdfglyphtounicode{sacutedotaccent}{1E65}
\pdfglyphtounicode{sadarabic}{0635}
\pdfglyphtounicode{sadeva}{0938}
\pdfglyphtounicode{sadfinalarabic}{FEBA}
\pdfglyphtounicode{sadinitialarabic}{FEBB}
\pdfglyphtounicode{sadmedialarabic}{FEBC}
\pdfglyphtounicode{sagujarati}{0AB8}
\pdfglyphtounicode{sagurmukhi}{0A38}
\pdfglyphtounicode{sahiragana}{3055}
\pdfglyphtounicode{sakatakana}{30B5}
\pdfglyphtounicode{sakatakanahalfwidth}{FF7B}
\pdfglyphtounicode{sallallahoualayhewasallamarabic}{FDFA}
\pdfglyphtounicode{samekh}{05E1}
\pdfglyphtounicode{samekhdagesh}{FB41}
\pdfglyphtounicode{samekhdageshhebrew}{FB41}
\pdfglyphtounicode{samekhhebrew}{05E1}
\pdfglyphtounicode{saraaathai}{0E32}
\pdfglyphtounicode{saraaethai}{0E41}
\pdfglyphtounicode{saraaimaimalaithai}{0E44}
\pdfglyphtounicode{saraaimaimuanthai}{0E43}
\pdfglyphtounicode{saraamthai}{0E33}
\pdfglyphtounicode{saraathai}{0E30}
\pdfglyphtounicode{saraethai}{0E40}
\pdfglyphtounicode{saraiileftthai}{F886}
\pdfglyphtounicode{saraiithai}{0E35}
\pdfglyphtounicode{saraileftthai}{F885}
\pdfglyphtounicode{saraithai}{0E34}
\pdfglyphtounicode{saraothai}{0E42}
\pdfglyphtounicode{saraueeleftthai}{F888}
\pdfglyphtounicode{saraueethai}{0E37}
\pdfglyphtounicode{saraueleftthai}{F887}
\pdfglyphtounicode{sarauethai}{0E36}
\pdfglyphtounicode{sarauthai}{0E38}
\pdfglyphtounicode{sarauuthai}{0E39}
\pdfglyphtounicode{satisfies}{22A8}
\pdfglyphtounicode{sbopomofo}{3119}
\pdfglyphtounicode{scaron}{0161}
\pdfglyphtounicode{scarondotaccent}{1E67}
\pdfglyphtounicode{scedilla}{015F}
\pdfglyphtounicode{schwa}{0259}
\pdfglyphtounicode{schwacyrillic}{04D9}
\pdfglyphtounicode{schwadieresiscyrillic}{04DB}
\pdfglyphtounicode{schwahook}{025A}
\pdfglyphtounicode{scircle}{24E2}
\pdfglyphtounicode{scircumflex}{015D}
\pdfglyphtounicode{scommaaccent}{0219}
\pdfglyphtounicode{sdotaccent}{1E61}
\pdfglyphtounicode{sdotbelow}{1E63}
\pdfglyphtounicode{sdotbelowdotaccent}{1E69}
\pdfglyphtounicode{seagullbelowcmb}{033C}
\pdfglyphtounicode{second}{2033}
\pdfglyphtounicode{secondtonechinese}{02CA}
\pdfglyphtounicode{section}{00A7}
\pdfglyphtounicode{seenarabic}{0633}
\pdfglyphtounicode{seenfinalarabic}{FEB2}
\pdfglyphtounicode{seeninitialarabic}{FEB3}
\pdfglyphtounicode{seenmedialarabic}{FEB4}
\pdfglyphtounicode{segol}{05B6}
\pdfglyphtounicode{segol13}{05B6}
\pdfglyphtounicode{segol1f}{05B6}
\pdfglyphtounicode{segol2c}{05B6}
\pdfglyphtounicode{segolhebrew}{05B6}
\pdfglyphtounicode{segolnarrowhebrew}{05B6}
\pdfglyphtounicode{segolquarterhebrew}{05B6}
\pdfglyphtounicode{segoltahebrew}{0592}
\pdfglyphtounicode{segolwidehebrew}{05B6}
\pdfglyphtounicode{seharmenian}{057D}
\pdfglyphtounicode{sehiragana}{305B}
\pdfglyphtounicode{sekatakana}{30BB}
\pdfglyphtounicode{sekatakanahalfwidth}{FF7E}
\pdfglyphtounicode{semicolon}{003B}
\pdfglyphtounicode{semicolonarabic}{061B}
\pdfglyphtounicode{semicolonmonospace}{FF1B}
\pdfglyphtounicode{semicolonsmall}{FE54}
\pdfglyphtounicode{semivoicedmarkkana}{309C}
\pdfglyphtounicode{semivoicedmarkkanahalfwidth}{FF9F}
\pdfglyphtounicode{sentisquare}{3322}
\pdfglyphtounicode{sentosquare}{3323}
\pdfglyphtounicode{seven}{0037}
\pdfglyphtounicode{sevenarabic}{0667}
\pdfglyphtounicode{sevenbengali}{09ED}
\pdfglyphtounicode{sevencircle}{2466}
\pdfglyphtounicode{sevencircleinversesansserif}{2790}
\pdfglyphtounicode{sevendeva}{096D}
\pdfglyphtounicode{seveneighths}{215E}
\pdfglyphtounicode{sevengujarati}{0AED}
\pdfglyphtounicode{sevengurmukhi}{0A6D}
\pdfglyphtounicode{sevenhackarabic}{0667}
\pdfglyphtounicode{sevenhangzhou}{3027}
\pdfglyphtounicode{sevenideographicparen}{3226}
\pdfglyphtounicode{seveninferior}{2087}
\pdfglyphtounicode{sevenmonospace}{FF17}
\pdfglyphtounicode{sevenoldstyle}{0037}
\pdfglyphtounicode{sevenparen}{247A}
\pdfglyphtounicode{sevenperiod}{248E}
\pdfglyphtounicode{sevenpersian}{06F7}
\pdfglyphtounicode{sevenroman}{2176}
\pdfglyphtounicode{sevensuperior}{2077}
\pdfglyphtounicode{seventeencircle}{2470}
\pdfglyphtounicode{seventeenparen}{2484}
\pdfglyphtounicode{seventeenperiod}{2498}
\pdfglyphtounicode{seventhai}{0E57}
\pdfglyphtounicode{sfthyphen}{00AD}
\pdfglyphtounicode{shaarmenian}{0577}
\pdfglyphtounicode{shabengali}{09B6}
\pdfglyphtounicode{shacyrillic}{0448}
\pdfglyphtounicode{shaddaarabic}{0651}
\pdfglyphtounicode{shaddadammaarabic}{FC61}
\pdfglyphtounicode{shaddadammatanarabic}{FC5E}
\pdfglyphtounicode{shaddafathaarabic}{FC60}
\pdfglyphtounicode{shaddafathatanarabic}{0651 064B}
\pdfglyphtounicode{shaddakasraarabic}{FC62}
\pdfglyphtounicode{shaddakasratanarabic}{FC5F}
\pdfglyphtounicode{shade}{2592}
\pdfglyphtounicode{shadedark}{2593}
\pdfglyphtounicode{shadelight}{2591}
\pdfglyphtounicode{shademedium}{2592}
\pdfglyphtounicode{shadeva}{0936}
\pdfglyphtounicode{shagujarati}{0AB6}
\pdfglyphtounicode{shagurmukhi}{0A36}
\pdfglyphtounicode{shalshelethebrew}{0593}
\pdfglyphtounicode{sharp}{266F}
\pdfglyphtounicode{shbopomofo}{3115}
\pdfglyphtounicode{shchacyrillic}{0449}
\pdfglyphtounicode{sheenarabic}{0634}
\pdfglyphtounicode{sheenfinalarabic}{FEB6}
\pdfglyphtounicode{sheeninitialarabic}{FEB7}
\pdfglyphtounicode{sheenmedialarabic}{FEB8}
\pdfglyphtounicode{sheicoptic}{03E3}
\pdfglyphtounicode{sheqel}{20AA}
\pdfglyphtounicode{sheqelhebrew}{20AA}
\pdfglyphtounicode{sheva}{05B0}
\pdfglyphtounicode{sheva115}{05B0}
\pdfglyphtounicode{sheva15}{05B0}
\pdfglyphtounicode{sheva22}{05B0}
\pdfglyphtounicode{sheva2e}{05B0}
\pdfglyphtounicode{shevahebrew}{05B0}
\pdfglyphtounicode{shevanarrowhebrew}{05B0}
\pdfglyphtounicode{shevaquarterhebrew}{05B0}
\pdfglyphtounicode{shevawidehebrew}{05B0}
\pdfglyphtounicode{shhacyrillic}{04BB}
\pdfglyphtounicode{shiftleft}{21B0}
\pdfglyphtounicode{shiftright}{21B1}
\pdfglyphtounicode{shimacoptic}{03ED}
\pdfglyphtounicode{shin}{05E9}
\pdfglyphtounicode{shindagesh}{FB49}
\pdfglyphtounicode{shindageshhebrew}{FB49}
\pdfglyphtounicode{shindageshshindot}{FB2C}
\pdfglyphtounicode{shindageshshindothebrew}{FB2C}
\pdfglyphtounicode{shindageshsindot}{FB2D}
\pdfglyphtounicode{shindageshsindothebrew}{FB2D}
\pdfglyphtounicode{shindothebrew}{05C1}
\pdfglyphtounicode{shinhebrew}{05E9}
\pdfglyphtounicode{shinshindot}{FB2A}
\pdfglyphtounicode{shinshindothebrew}{FB2A}
\pdfglyphtounicode{shinsindot}{FB2B}
\pdfglyphtounicode{shinsindothebrew}{FB2B}
\pdfglyphtounicode{shook}{0282}
\pdfglyphtounicode{sigma}{03C3}
\pdfglyphtounicode{sigma1}{03C2}
\pdfglyphtounicode{sigmafinal}{03C2}
\pdfglyphtounicode{sigmalunatesymbolgreek}{03F2}
\pdfglyphtounicode{sihiragana}{3057}
\pdfglyphtounicode{sikatakana}{30B7}
\pdfglyphtounicode{sikatakanahalfwidth}{FF7C}
\pdfglyphtounicode{siluqhebrew}{05BD}
\pdfglyphtounicode{siluqlefthebrew}{05BD}
\pdfglyphtounicode{similar}{223C}
\pdfglyphtounicode{similarequal}{2243}
\pdfglyphtounicode{sindothebrew}{05C2}
\pdfglyphtounicode{siosacirclekorean}{3274}
\pdfglyphtounicode{siosaparenkorean}{3214}
\pdfglyphtounicode{sioscieuckorean}{317E}
\pdfglyphtounicode{sioscirclekorean}{3266}
\pdfglyphtounicode{sioskiyeokkorean}{317A}
\pdfglyphtounicode{sioskorean}{3145}
\pdfglyphtounicode{siosnieunkorean}{317B}
\pdfglyphtounicode{siosparenkorean}{3206}
\pdfglyphtounicode{siospieupkorean}{317D}
\pdfglyphtounicode{siostikeutkorean}{317C}
\pdfglyphtounicode{six}{0036}
\pdfglyphtounicode{sixarabic}{0666}
\pdfglyphtounicode{sixbengali}{09EC}
\pdfglyphtounicode{sixcircle}{2465}
\pdfglyphtounicode{sixcircleinversesansserif}{278F}
\pdfglyphtounicode{sixdeva}{096C}
\pdfglyphtounicode{sixgujarati}{0AEC}
\pdfglyphtounicode{sixgurmukhi}{0A6C}
\pdfglyphtounicode{sixhackarabic}{0666}
\pdfglyphtounicode{sixhangzhou}{3026}
\pdfglyphtounicode{sixideographicparen}{3225}
\pdfglyphtounicode{sixinferior}{2086}
\pdfglyphtounicode{sixmonospace}{FF16}
\pdfglyphtounicode{sixoldstyle}{0036}
\pdfglyphtounicode{sixparen}{2479}
\pdfglyphtounicode{sixperiod}{248D}
\pdfglyphtounicode{sixpersian}{06F6}
\pdfglyphtounicode{sixroman}{2175}
\pdfglyphtounicode{sixsuperior}{2076}
\pdfglyphtounicode{sixteencircle}{246F}
\pdfglyphtounicode{sixteencurrencydenominatorbengali}{09F9}
\pdfglyphtounicode{sixteenparen}{2483}
\pdfglyphtounicode{sixteenperiod}{2497}
\pdfglyphtounicode{sixthai}{0E56}
\pdfglyphtounicode{slash}{002F}
\pdfglyphtounicode{slashmonospace}{FF0F}
\pdfglyphtounicode{slong}{017F}
\pdfglyphtounicode{slongdotaccent}{1E9B}
\pdfglyphtounicode{slurabove}{2322}
\pdfglyphtounicode{slurbelow}{2323}
\pdfglyphtounicode{smile}{2323}
\pdfglyphtounicode{smileface}{263A}
\pdfglyphtounicode{smonospace}{FF53}
\pdfglyphtounicode{sofpasuqhebrew}{05C3}
\pdfglyphtounicode{softhyphen}{00AD}
\pdfglyphtounicode{softsigncyrillic}{044C}
\pdfglyphtounicode{sohiragana}{305D}
\pdfglyphtounicode{sokatakana}{30BD}
\pdfglyphtounicode{sokatakanahalfwidth}{FF7F}
\pdfglyphtounicode{soliduslongoverlaycmb}{0338}
\pdfglyphtounicode{solidusshortoverlaycmb}{0337}
\pdfglyphtounicode{sorusithai}{0E29}
\pdfglyphtounicode{sosalathai}{0E28}
\pdfglyphtounicode{sosothai}{0E0B}
\pdfglyphtounicode{sosuathai}{0E2A}
\pdfglyphtounicode{space}{0020}
\pdfglyphtounicode{spacehackarabic}{0020}
\pdfglyphtounicode{spade}{2660}
\pdfglyphtounicode{spadesuitblack}{2660}
\pdfglyphtounicode{spadesuitwhite}{2664}
\pdfglyphtounicode{sparen}{24AE}
\pdfglyphtounicode{sphericalangle}{2222}
\pdfglyphtounicode{square}{25A1}
\pdfglyphtounicode{squarebelowcmb}{033B}
\pdfglyphtounicode{squarecc}{33C4}
\pdfglyphtounicode{squarecm}{339D}
\pdfglyphtounicode{squarediagonalcrosshatchfill}{25A9}
\pdfglyphtounicode{squaredot}{22A1}
\pdfglyphtounicode{squarehorizontalfill}{25A4}
\pdfglyphtounicode{squareimage}{228F}
\pdfglyphtounicode{squarekg}{338F}
\pdfglyphtounicode{squarekm}{339E}
\pdfglyphtounicode{squarekmcapital}{33CE}
\pdfglyphtounicode{squareln}{33D1}
\pdfglyphtounicode{squarelog}{33D2}
\pdfglyphtounicode{squaremg}{338E}
\pdfglyphtounicode{squaremil}{33D5}
\pdfglyphtounicode{squareminus}{229F}
\pdfglyphtounicode{squaremm}{339C}
\pdfglyphtounicode{squaremsquared}{33A1}
\pdfglyphtounicode{squaremultiply}{22A0}
\pdfglyphtounicode{squareoriginal}{2290}
\pdfglyphtounicode{squareorthogonalcrosshatchfill}{25A6}
\pdfglyphtounicode{squareplus}{229E}
\pdfglyphtounicode{squaresolid}{25A0}
\pdfglyphtounicode{squareupperlefttolowerrightfill}{25A7}
\pdfglyphtounicode{squareupperrighttolowerleftfill}{25A8}
\pdfglyphtounicode{squareverticalfill}{25A5}
\pdfglyphtounicode{squarewhitewithsmallblack}{25A3}
\pdfglyphtounicode{squiggleleftright}{21AD}
\pdfglyphtounicode{squiggleright}{21DD}
\pdfglyphtounicode{srsquare}{33DB}
\pdfglyphtounicode{ssabengali}{09B7}
\pdfglyphtounicode{ssadeva}{0937}
\pdfglyphtounicode{ssagujarati}{0AB7}
\pdfglyphtounicode{ssangcieuckorean}{3149}
\pdfglyphtounicode{ssanghieuhkorean}{3185}
\pdfglyphtounicode{ssangieungkorean}{3180}
\pdfglyphtounicode{ssangkiyeokkorean}{3132}
\pdfglyphtounicode{ssangnieunkorean}{3165}
\pdfglyphtounicode{ssangpieupkorean}{3143}
\pdfglyphtounicode{ssangsioskorean}{3146}
\pdfglyphtounicode{ssangtikeutkorean}{3138}
\pdfglyphtounicode{ssuperior}{0073}
\pdfglyphtounicode{st}{0073 0074}
\pdfglyphtounicode{star}{22C6}
\pdfglyphtounicode{sterling}{00A3}
\pdfglyphtounicode{sterlingmonospace}{FFE1}
\pdfglyphtounicode{strokelongoverlaycmb}{0336}
\pdfglyphtounicode{strokeshortoverlaycmb}{0335}
\pdfglyphtounicode{subset}{2282}
\pdfglyphtounicode{subsetdbl}{22D0}
\pdfglyphtounicode{subsetdblequal}{2AC5}
\pdfglyphtounicode{subsetnoteql}{228A}
\pdfglyphtounicode{subsetnotequal}{228A}
\pdfglyphtounicode{subsetorequal}{2286}
\pdfglyphtounicode{subsetornotdbleql}{2ACB}
\pdfglyphtounicode{subsetsqequal}{2291}
\pdfglyphtounicode{succeeds}{227B}
\pdfglyphtounicode{suchthat}{220B}
\pdfglyphtounicode{suhiragana}{3059}
\pdfglyphtounicode{sukatakana}{30B9}
\pdfglyphtounicode{sukatakanahalfwidth}{FF7D}
\pdfglyphtounicode{sukunarabic}{0652}
\pdfglyphtounicode{summation}{2211}
\pdfglyphtounicode{sun}{263C}
\pdfglyphtounicode{superset}{2283}
\pdfglyphtounicode{supersetdbl}{22D1}
\pdfglyphtounicode{supersetdblequal}{2AC6}
\pdfglyphtounicode{supersetnoteql}{228B}
\pdfglyphtounicode{supersetnotequal}{228B}
\pdfglyphtounicode{supersetorequal}{2287}
\pdfglyphtounicode{supersetornotdbleql}{2ACC}
\pdfglyphtounicode{supersetsqequal}{2292}
\pdfglyphtounicode{svsquare}{33DC}
\pdfglyphtounicode{syouwaerasquare}{337C}
\pdfglyphtounicode{t}{0074}
\pdfglyphtounicode{tabengali}{09A4}
\pdfglyphtounicode{tackdown}{22A4}
\pdfglyphtounicode{tackleft}{22A3}
\pdfglyphtounicode{tadeva}{0924}
\pdfglyphtounicode{tagujarati}{0AA4}
\pdfglyphtounicode{tagurmukhi}{0A24}
\pdfglyphtounicode{taharabic}{0637}
\pdfglyphtounicode{tahfinalarabic}{FEC2}
\pdfglyphtounicode{tahinitialarabic}{FEC3}
\pdfglyphtounicode{tahiragana}{305F}
\pdfglyphtounicode{tahmedialarabic}{FEC4}
\pdfglyphtounicode{taisyouerasquare}{337D}
\pdfglyphtounicode{takatakana}{30BF}
\pdfglyphtounicode{takatakanahalfwidth}{FF80}
\pdfglyphtounicode{tatweelarabic}{0640}
\pdfglyphtounicode{tau}{03C4}
\pdfglyphtounicode{tav}{05EA}
\pdfglyphtounicode{tavdages}{FB4A}
\pdfglyphtounicode{tavdagesh}{FB4A}
\pdfglyphtounicode{tavdageshhebrew}{FB4A}
\pdfglyphtounicode{tavhebrew}{05EA}
\pdfglyphtounicode{tbar}{0167}
\pdfglyphtounicode{tbopomofo}{310A}
\pdfglyphtounicode{tcaron}{0165}
\pdfglyphtounicode{tccurl}{02A8}
\pdfglyphtounicode{tcedilla}{0163}
\pdfglyphtounicode{tcheharabic}{0686}
\pdfglyphtounicode{tchehfinalarabic}{FB7B}
\pdfglyphtounicode{tchehinitialarabic}{FB7C}
\pdfglyphtounicode{tchehmedialarabic}{FB7D}
\pdfglyphtounicode{tchehmeeminitialarabic}{FB7C FEE4}
\pdfglyphtounicode{tcircle}{24E3}
\pdfglyphtounicode{tcircumflexbelow}{1E71}
\pdfglyphtounicode{tcommaaccent}{0163}
\pdfglyphtounicode{tdieresis}{1E97}
\pdfglyphtounicode{tdotaccent}{1E6B}
\pdfglyphtounicode{tdotbelow}{1E6D}
\pdfglyphtounicode{tecyrillic}{0442}
\pdfglyphtounicode{tedescendercyrillic}{04AD}
\pdfglyphtounicode{teharabic}{062A}
\pdfglyphtounicode{tehfinalarabic}{FE96}
\pdfglyphtounicode{tehhahinitialarabic}{FCA2}
\pdfglyphtounicode{tehhahisolatedarabic}{FC0C}
\pdfglyphtounicode{tehinitialarabic}{FE97}
\pdfglyphtounicode{tehiragana}{3066}
\pdfglyphtounicode{tehjeeminitialarabic}{FCA1}
\pdfglyphtounicode{tehjeemisolatedarabic}{FC0B}
\pdfglyphtounicode{tehmarbutaarabic}{0629}
\pdfglyphtounicode{tehmarbutafinalarabic}{FE94}
\pdfglyphtounicode{tehmedialarabic}{FE98}
\pdfglyphtounicode{tehmeeminitialarabic}{FCA4}
\pdfglyphtounicode{tehmeemisolatedarabic}{FC0E}
\pdfglyphtounicode{tehnoonfinalarabic}{FC73}
\pdfglyphtounicode{tekatakana}{30C6}
\pdfglyphtounicode{tekatakanahalfwidth}{FF83}
\pdfglyphtounicode{telephone}{2121}
\pdfglyphtounicode{telephoneblack}{260E}
\pdfglyphtounicode{telishagedolahebrew}{05A0}
\pdfglyphtounicode{telishaqetanahebrew}{05A9}
\pdfglyphtounicode{tencircle}{2469}
\pdfglyphtounicode{tenideographicparen}{3229}
\pdfglyphtounicode{tenparen}{247D}
\pdfglyphtounicode{tenperiod}{2491}
\pdfglyphtounicode{tenroman}{2179}
\pdfglyphtounicode{tesh}{02A7}
\pdfglyphtounicode{tet}{05D8}
\pdfglyphtounicode{tetdagesh}{FB38}
\pdfglyphtounicode{tetdageshhebrew}{FB38}
\pdfglyphtounicode{tethebrew}{05D8}
\pdfglyphtounicode{tetsecyrillic}{04B5}
\pdfglyphtounicode{tevirhebrew}{059B}
\pdfglyphtounicode{tevirlefthebrew}{059B}
\pdfglyphtounicode{tfm:cmbsy10/diamond}{2662}
\pdfglyphtounicode{tfm:cmbsy10/heart}{2661}
\pdfglyphtounicode{tfm:cmbsy5/diamond}{2662}
\pdfglyphtounicode{tfm:cmbsy5/heart}{2661}
\pdfglyphtounicode{tfm:cmbsy6/diamond}{2662}
\pdfglyphtounicode{tfm:cmbsy6/heart}{2661}
\pdfglyphtounicode{tfm:cmbsy7/diamond}{2662}
\pdfglyphtounicode{tfm:cmbsy7/heart}{2661}
\pdfglyphtounicode{tfm:cmbsy8/diamond}{2662}
\pdfglyphtounicode{tfm:cmbsy8/heart}{2661}
\pdfglyphtounicode{tfm:cmbsy9/diamond}{2662}
\pdfglyphtounicode{tfm:cmbsy9/heart}{2661}
\pdfglyphtounicode{tfm:cmmi10/phi}{03D5}
\pdfglyphtounicode{tfm:cmmi10/phi1}{03C6}
\pdfglyphtounicode{tfm:cmmi12/phi}{03D5}
\pdfglyphtounicode{tfm:cmmi12/phi1}{03C6}
\pdfglyphtounicode{tfm:cmmi5/phi}{03D5}
\pdfglyphtounicode{tfm:cmmi5/phi1}{03C6}
\pdfglyphtounicode{tfm:cmmi6/phi}{03D5}
\pdfglyphtounicode{tfm:cmmi6/phi1}{03C6}
\pdfglyphtounicode{tfm:cmmi7/phi}{03D5}
\pdfglyphtounicode{tfm:cmmi7/phi1}{03C6}
\pdfglyphtounicode{tfm:cmmi8/phi}{03D5}
\pdfglyphtounicode{tfm:cmmi8/phi1}{03C6}
\pdfglyphtounicode{tfm:cmmi9/phi}{03D5}
\pdfglyphtounicode{tfm:cmmi9/phi1}{03C6}
\pdfglyphtounicode{tfm:cmmib10/phi}{03D5}
\pdfglyphtounicode{tfm:cmmib10/phi1}{03C6}
\pdfglyphtounicode{tfm:cmmib5/phi}{03D5}
\pdfglyphtounicode{tfm:cmmib5/phi1}{03C6}
\pdfglyphtounicode{tfm:cmmib6/phi}{03D5}
\pdfglyphtounicode{tfm:cmmib6/phi1}{03C6}
\pdfglyphtounicode{tfm:cmmib7/phi}{03D5}
\pdfglyphtounicode{tfm:cmmib7/phi1}{03C6}
\pdfglyphtounicode{tfm:cmmib8/phi}{03D5}
\pdfglyphtounicode{tfm:cmmib8/phi1}{03C6}
\pdfglyphtounicode{tfm:cmmib9/phi}{03D5}
\pdfglyphtounicode{tfm:cmmib9/phi1}{03C6}
\pdfglyphtounicode{tfm:cmsy10/diamond}{2662}
\pdfglyphtounicode{tfm:cmsy10/heart}{2661}
\pdfglyphtounicode{tfm:cmsy5/heart}{2661}
\pdfglyphtounicode{tfm:cmsy6/diamond}{2662}
\pdfglyphtounicode{tfm:cmsy6/heart}{2661}
\pdfglyphtounicode{tfm:cmsy7/diamond}{2662}
\pdfglyphtounicode{tfm:cmsy7/heart}{2661}
\pdfglyphtounicode{tfm:cmsy8/diamond}{2662}
\pdfglyphtounicode{tfm:cmsy8/heart}{2661}
\pdfglyphtounicode{tfm:cmsy9/diamond}{2662}
\pdfglyphtounicode{tfm:cmsy9/heart}{2661}
\pdfglyphtounicode{tfm:eurb10/phi}{03D5}
\pdfglyphtounicode{tfm:eurb10/phi1}{03C6}
\pdfglyphtounicode{tfm:eurb5/phi}{03D5}
\pdfglyphtounicode{tfm:eurb5/phi1}{03C6}
\pdfglyphtounicode{tfm:eurb6/phi}{03D5}
\pdfglyphtounicode{tfm:eurb6/phi1}{03C6}
\pdfglyphtounicode{tfm:eurb7/phi}{03D5}
\pdfglyphtounicode{tfm:eurb7/phi1}{03C6}
\pdfglyphtounicode{tfm:eurb8/phi}{03D5}
\pdfglyphtounicode{tfm:eurb8/phi1}{03C6}
\pdfglyphtounicode{tfm:eurb9/phi}{03D5}
\pdfglyphtounicode{tfm:eurb9/phi1}{03C6}
\pdfglyphtounicode{tfm:eurm10/phi}{03D5}
\pdfglyphtounicode{tfm:eurm10/phi1}{03C6}
\pdfglyphtounicode{tfm:eurm5/phi}{03D5}
\pdfglyphtounicode{tfm:eurm5/phi1}{03C6}
\pdfglyphtounicode{tfm:eurm6/phi}{03D5}
\pdfglyphtounicode{tfm:eurm6/phi1}{03C6}
\pdfglyphtounicode{tfm:eurm7/phi}{03D5}
\pdfglyphtounicode{tfm:eurm7/phi1}{03C6}
\pdfglyphtounicode{tfm:eurm8/phi}{03D5}
\pdfglyphtounicode{tfm:eurm8/phi1}{03C6}
\pdfglyphtounicode{tfm:eurm9/phi}{03D5}
\pdfglyphtounicode{tfm:eurm9/phi1}{03C6}
\pdfglyphtounicode{tfm:fplmbi/phi}{03D5}
\pdfglyphtounicode{tfm:fplmbi/phi1}{03C6}
\pdfglyphtounicode{tfm:fplmri/phi}{03D5}
\pdfglyphtounicode{tfm:fplmri/phi1}{03C6}
\pdfglyphtounicode{tfm:lmbsy10/diamond}{2662}
\pdfglyphtounicode{tfm:lmbsy10/heart}{2661}
\pdfglyphtounicode{tfm:lmbsy5/diamond}{2662}
\pdfglyphtounicode{tfm:lmbsy5/heart}{2661}
\pdfglyphtounicode{tfm:lmbsy7/diamond}{2662}
\pdfglyphtounicode{tfm:lmbsy7/heart}{2661}
\pdfglyphtounicode{tfm:lmmi10/phi}{03D5}
\pdfglyphtounicode{tfm:lmmi10/phi1}{03C6}
\pdfglyphtounicode{tfm:lmmi12/phi}{03D5}
\pdfglyphtounicode{tfm:lmmi12/phi1}{03C6}
\pdfglyphtounicode{tfm:lmmi5/phi}{03D5}
\pdfglyphtounicode{tfm:lmmi5/phi1}{03C6}
\pdfglyphtounicode{tfm:lmmi6/phi}{03D5}
\pdfglyphtounicode{tfm:lmmi6/phi1}{03C6}
\pdfglyphtounicode{tfm:lmmi7/phi}{03D5}
\pdfglyphtounicode{tfm:lmmi7/phi1}{03C6}
\pdfglyphtounicode{tfm:lmmi8/phi}{03D5}
\pdfglyphtounicode{tfm:lmmi8/phi1}{03C6}
\pdfglyphtounicode{tfm:lmmi9/phi}{03D5}
\pdfglyphtounicode{tfm:lmmi9/phi1}{03C6}
\pdfglyphtounicode{tfm:lmmib10/phi}{03D5}
\pdfglyphtounicode{tfm:lmmib10/phi1}{03C6}
\pdfglyphtounicode{tfm:lmmib5/phi}{03D5}
\pdfglyphtounicode{tfm:lmmib5/phi1}{03C6}
\pdfglyphtounicode{tfm:lmmib7/phi}{03D5}
\pdfglyphtounicode{tfm:lmmib7/phi1}{03C6}
\pdfglyphtounicode{tfm:lmsy10/diamond}{2662}
\pdfglyphtounicode{tfm:lmsy10/heart}{2661}
\pdfglyphtounicode{tfm:lmsy5/diamond}{2662}
\pdfglyphtounicode{tfm:lmsy5/heart}{2661}
\pdfglyphtounicode{tfm:lmsy6/diamond}{2662}
\pdfglyphtounicode{tfm:lmsy6/heart}{2661}
\pdfglyphtounicode{tfm:lmsy7/diamond}{2662}
\pdfglyphtounicode{tfm:lmsy7/heart}{2661}
\pdfglyphtounicode{tfm:lmsy8/diamond}{2662}
\pdfglyphtounicode{tfm:lmsy8/heart}{2661}
\pdfglyphtounicode{tfm:lmsy9/diamond}{2662}
\pdfglyphtounicode{tfm:lmsy9/heart}{2661}
\pdfglyphtounicode{tfm:msam10/diamond}{2662}
\pdfglyphtounicode{tfm:msam5/diamond}{2662}
\pdfglyphtounicode{tfm:msam6/diamond}{2662}
\pdfglyphtounicode{tfm:msam7/diamond}{2662}
\pdfglyphtounicode{tfm:msam8/diamond}{2662}
\pdfglyphtounicode{tfm:msam9/diamond}{2662}
\pdfglyphtounicode{tfm:pxbmia/phi}{03D5}
\pdfglyphtounicode{tfm:pxbmia/phi1}{03C6}
\pdfglyphtounicode{tfm:pxbsy/diamond}{2662}
\pdfglyphtounicode{tfm:pxbsy/heart}{2661}
\pdfglyphtounicode{tfm:pxbsya/diamond}{2662}
\pdfglyphtounicode{tfm:pxmia/phi}{03D5}
\pdfglyphtounicode{tfm:pxmia/phi1}{03C6}
\pdfglyphtounicode{tfm:pxsy/diamond}{2662}
\pdfglyphtounicode{tfm:pxsy/heart}{2661}
\pdfglyphtounicode{tfm:pxsya/diamond}{2662}
\pdfglyphtounicode{tfm:pzdr/a1}{2701}
\pdfglyphtounicode{tfm:pzdr/a10}{2721}
\pdfglyphtounicode{tfm:pzdr/a100}{275E}
\pdfglyphtounicode{tfm:pzdr/a101}{2761}
\pdfglyphtounicode{tfm:pzdr/a102}{2762}
\pdfglyphtounicode{tfm:pzdr/a103}{2763}
\pdfglyphtounicode{tfm:pzdr/a104}{2764}
\pdfglyphtounicode{tfm:pzdr/a105}{2710}
\pdfglyphtounicode{tfm:pzdr/a106}{2765}
\pdfglyphtounicode{tfm:pzdr/a107}{2766}
\pdfglyphtounicode{tfm:pzdr/a108}{2767}
\pdfglyphtounicode{tfm:pzdr/a109}{2660}
\pdfglyphtounicode{tfm:pzdr/a11}{261B}
\pdfglyphtounicode{tfm:pzdr/a110}{2665}
\pdfglyphtounicode{tfm:pzdr/a111}{2666}
\pdfglyphtounicode{tfm:pzdr/a112}{2663}
\pdfglyphtounicode{tfm:pzdr/a117}{2709}
\pdfglyphtounicode{tfm:pzdr/a118}{2708}
\pdfglyphtounicode{tfm:pzdr/a119}{2707}
\pdfglyphtounicode{tfm:pzdr/a12}{261E}
\pdfglyphtounicode{tfm:pzdr/a120}{2460}
\pdfglyphtounicode{tfm:pzdr/a121}{2461}
\pdfglyphtounicode{tfm:pzdr/a122}{2462}
\pdfglyphtounicode{tfm:pzdr/a123}{2463}
\pdfglyphtounicode{tfm:pzdr/a124}{2464}
\pdfglyphtounicode{tfm:pzdr/a125}{2465}
\pdfglyphtounicode{tfm:pzdr/a126}{2466}
\pdfglyphtounicode{tfm:pzdr/a127}{2467}
\pdfglyphtounicode{tfm:pzdr/a128}{2468}
\pdfglyphtounicode{tfm:pzdr/a129}{2469}
\pdfglyphtounicode{tfm:pzdr/a13}{270C}
\pdfglyphtounicode{tfm:pzdr/a130}{2776}
\pdfglyphtounicode{tfm:pzdr/a131}{2777}
\pdfglyphtounicode{tfm:pzdr/a132}{2778}
\pdfglyphtounicode{tfm:pzdr/a133}{2779}
\pdfglyphtounicode{tfm:pzdr/a134}{277A}
\pdfglyphtounicode{tfm:pzdr/a135}{277B}
\pdfglyphtounicode{tfm:pzdr/a136}{277C}
\pdfglyphtounicode{tfm:pzdr/a137}{277D}
\pdfglyphtounicode{tfm:pzdr/a138}{277E}
\pdfglyphtounicode{tfm:pzdr/a139}{277F}
\pdfglyphtounicode{tfm:pzdr/a14}{270D}
\pdfglyphtounicode{tfm:pzdr/a140}{2780}
\pdfglyphtounicode{tfm:pzdr/a141}{2781}
\pdfglyphtounicode{tfm:pzdr/a142}{2782}
\pdfglyphtounicode{tfm:pzdr/a143}{2783}
\pdfglyphtounicode{tfm:pzdr/a144}{2784}
\pdfglyphtounicode{tfm:pzdr/a145}{2785}
\pdfglyphtounicode{tfm:pzdr/a146}{2786}
\pdfglyphtounicode{tfm:pzdr/a147}{2787}
\pdfglyphtounicode{tfm:pzdr/a148}{2788}
\pdfglyphtounicode{tfm:pzdr/a149}{2789}
\pdfglyphtounicode{tfm:pzdr/a15}{270E}
\pdfglyphtounicode{tfm:pzdr/a150}{278A}
\pdfglyphtounicode{tfm:pzdr/a151}{278B}
\pdfglyphtounicode{tfm:pzdr/a152}{278C}
\pdfglyphtounicode{tfm:pzdr/a153}{278D}
\pdfglyphtounicode{tfm:pzdr/a154}{278E}
\pdfglyphtounicode{tfm:pzdr/a155}{278F}
\pdfglyphtounicode{tfm:pzdr/a156}{2790}
\pdfglyphtounicode{tfm:pzdr/a157}{2791}
\pdfglyphtounicode{tfm:pzdr/a158}{2792}
\pdfglyphtounicode{tfm:pzdr/a159}{2793}
\pdfglyphtounicode{tfm:pzdr/a16}{270F}
\pdfglyphtounicode{tfm:pzdr/a160}{2794}
\pdfglyphtounicode{tfm:pzdr/a161}{2192}
\pdfglyphtounicode{tfm:pzdr/a162}{27A3}
\pdfglyphtounicode{tfm:pzdr/a163}{2194}
\pdfglyphtounicode{tfm:pzdr/a164}{2195}
\pdfglyphtounicode{tfm:pzdr/a165}{2799}
\pdfglyphtounicode{tfm:pzdr/a166}{279B}
\pdfglyphtounicode{tfm:pzdr/a167}{279C}
\pdfglyphtounicode{tfm:pzdr/a168}{279D}
\pdfglyphtounicode{tfm:pzdr/a169}{279E}
\pdfglyphtounicode{tfm:pzdr/a17}{2711}
\pdfglyphtounicode{tfm:pzdr/a170}{279F}
\pdfglyphtounicode{tfm:pzdr/a171}{27A0}
\pdfglyphtounicode{tfm:pzdr/a172}{27A1}
\pdfglyphtounicode{tfm:pzdr/a173}{27A2}
\pdfglyphtounicode{tfm:pzdr/a174}{27A4}
\pdfglyphtounicode{tfm:pzdr/a175}{27A5}
\pdfglyphtounicode{tfm:pzdr/a176}{27A6}
\pdfglyphtounicode{tfm:pzdr/a177}{27A7}
\pdfglyphtounicode{tfm:pzdr/a178}{27A8}
\pdfglyphtounicode{tfm:pzdr/a179}{27A9}
\pdfglyphtounicode{tfm:pzdr/a18}{2712}
\pdfglyphtounicode{tfm:pzdr/a180}{27AB}
\pdfglyphtounicode{tfm:pzdr/a181}{27AD}
\pdfglyphtounicode{tfm:pzdr/a182}{27AF}
\pdfglyphtounicode{tfm:pzdr/a183}{27B2}
\pdfglyphtounicode{tfm:pzdr/a184}{27B3}
\pdfglyphtounicode{tfm:pzdr/a185}{27B5}
\pdfglyphtounicode{tfm:pzdr/a186}{27B8}
\pdfglyphtounicode{tfm:pzdr/a187}{27BA}
\pdfglyphtounicode{tfm:pzdr/a188}{27BB}
\pdfglyphtounicode{tfm:pzdr/a189}{27BC}
\pdfglyphtounicode{tfm:pzdr/a19}{2713}
\pdfglyphtounicode{tfm:pzdr/a190}{27BD}
\pdfglyphtounicode{tfm:pzdr/a191}{27BE}
\pdfglyphtounicode{tfm:pzdr/a192}{279A}
\pdfglyphtounicode{tfm:pzdr/a193}{27AA}
\pdfglyphtounicode{tfm:pzdr/a194}{27B6}
\pdfglyphtounicode{tfm:pzdr/a195}{27B9}
\pdfglyphtounicode{tfm:pzdr/a196}{2798}
\pdfglyphtounicode{tfm:pzdr/a197}{27B4}
\pdfglyphtounicode{tfm:pzdr/a198}{27B7}
\pdfglyphtounicode{tfm:pzdr/a199}{27AC}
\pdfglyphtounicode{tfm:pzdr/a2}{2702}
\pdfglyphtounicode{tfm:pzdr/a20}{2714}
\pdfglyphtounicode{tfm:pzdr/a200}{27AE}
\pdfglyphtounicode{tfm:pzdr/a201}{27B1}
\pdfglyphtounicode{tfm:pzdr/a202}{2703}
\pdfglyphtounicode{tfm:pzdr/a203}{2750}
\pdfglyphtounicode{tfm:pzdr/a204}{2752}
\pdfglyphtounicode{tfm:pzdr/a205}{276E}
\pdfglyphtounicode{tfm:pzdr/a206}{2770}
\pdfglyphtounicode{tfm:pzdr/a21}{2715}
\pdfglyphtounicode{tfm:pzdr/a22}{2716}
\pdfglyphtounicode{tfm:pzdr/a23}{2717}
\pdfglyphtounicode{tfm:pzdr/a24}{2718}
\pdfglyphtounicode{tfm:pzdr/a25}{2719}
\pdfglyphtounicode{tfm:pzdr/a26}{271A}
\pdfglyphtounicode{tfm:pzdr/a27}{271B}
\pdfglyphtounicode{tfm:pzdr/a28}{271C}
\pdfglyphtounicode{tfm:pzdr/a29}{2722}
\pdfglyphtounicode{tfm:pzdr/a3}{2704}
\pdfglyphtounicode{tfm:pzdr/a30}{2723}
\pdfglyphtounicode{tfm:pzdr/a31}{2724}
\pdfglyphtounicode{tfm:pzdr/a32}{2725}
\pdfglyphtounicode{tfm:pzdr/a33}{2726}
\pdfglyphtounicode{tfm:pzdr/a34}{2727}
\pdfglyphtounicode{tfm:pzdr/a35}{2605}
\pdfglyphtounicode{tfm:pzdr/a36}{2729}
\pdfglyphtounicode{tfm:pzdr/a37}{272A}
\pdfglyphtounicode{tfm:pzdr/a38}{272B}
\pdfglyphtounicode{tfm:pzdr/a39}{272C}
\pdfglyphtounicode{tfm:pzdr/a4}{260E}
\pdfglyphtounicode{tfm:pzdr/a40}{272D}
\pdfglyphtounicode{tfm:pzdr/a41}{272E}
\pdfglyphtounicode{tfm:pzdr/a42}{272F}
\pdfglyphtounicode{tfm:pzdr/a43}{2730}
\pdfglyphtounicode{tfm:pzdr/a44}{2731}
\pdfglyphtounicode{tfm:pzdr/a45}{2732}
\pdfglyphtounicode{tfm:pzdr/a46}{2733}
\pdfglyphtounicode{tfm:pzdr/a47}{2734}
\pdfglyphtounicode{tfm:pzdr/a48}{2735}
\pdfglyphtounicode{tfm:pzdr/a49}{2736}
\pdfglyphtounicode{tfm:pzdr/a5}{2706}
\pdfglyphtounicode{tfm:pzdr/a50}{2737}
\pdfglyphtounicode{tfm:pzdr/a51}{2738}
\pdfglyphtounicode{tfm:pzdr/a52}{2739}
\pdfglyphtounicode{tfm:pzdr/a53}{273A}
\pdfglyphtounicode{tfm:pzdr/a54}{273B}
\pdfglyphtounicode{tfm:pzdr/a55}{273C}
\pdfglyphtounicode{tfm:pzdr/a56}{273D}
\pdfglyphtounicode{tfm:pzdr/a57}{273E}
\pdfglyphtounicode{tfm:pzdr/a58}{273F}
\pdfglyphtounicode{tfm:pzdr/a59}{2740}
\pdfglyphtounicode{tfm:pzdr/a6}{271D}
\pdfglyphtounicode{tfm:pzdr/a60}{2741}
\pdfglyphtounicode{tfm:pzdr/a61}{2742}
\pdfglyphtounicode{tfm:pzdr/a62}{2743}
\pdfglyphtounicode{tfm:pzdr/a63}{2744}
\pdfglyphtounicode{tfm:pzdr/a64}{2745}
\pdfglyphtounicode{tfm:pzdr/a65}{2746}
\pdfglyphtounicode{tfm:pzdr/a66}{2747}
\pdfglyphtounicode{tfm:pzdr/a67}{2748}
\pdfglyphtounicode{tfm:pzdr/a68}{2749}
\pdfglyphtounicode{tfm:pzdr/a69}{274A}
\pdfglyphtounicode{tfm:pzdr/a7}{271E}
\pdfglyphtounicode{tfm:pzdr/a70}{274B}
\pdfglyphtounicode{tfm:pzdr/a71}{25CF}
\pdfglyphtounicode{tfm:pzdr/a72}{274D}
\pdfglyphtounicode{tfm:pzdr/a73}{25A0}
\pdfglyphtounicode{tfm:pzdr/a74}{274F}
\pdfglyphtounicode{tfm:pzdr/a75}{2751}
\pdfglyphtounicode{tfm:pzdr/a76}{25B2}
\pdfglyphtounicode{tfm:pzdr/a77}{25BC}
\pdfglyphtounicode{tfm:pzdr/a78}{25C6}
\pdfglyphtounicode{tfm:pzdr/a79}{2756}
\pdfglyphtounicode{tfm:pzdr/a8}{271F}
\pdfglyphtounicode{tfm:pzdr/a81}{25D7}
\pdfglyphtounicode{tfm:pzdr/a82}{2758}
\pdfglyphtounicode{tfm:pzdr/a83}{2759}
\pdfglyphtounicode{tfm:pzdr/a84}{275A}
\pdfglyphtounicode{tfm:pzdr/a85}{276F}
\pdfglyphtounicode{tfm:pzdr/a86}{2771}
\pdfglyphtounicode{tfm:pzdr/a87}{2772}
\pdfglyphtounicode{tfm:pzdr/a88}{2773}
\pdfglyphtounicode{tfm:pzdr/a89}{2768}
\pdfglyphtounicode{tfm:pzdr/a9}{2720}
\pdfglyphtounicode{tfm:pzdr/a90}{2769}
\pdfglyphtounicode{tfm:pzdr/a91}{276C}
\pdfglyphtounicode{tfm:pzdr/a92}{276D}
\pdfglyphtounicode{tfm:pzdr/a93}{276A}
\pdfglyphtounicode{tfm:pzdr/a94}{276B}
\pdfglyphtounicode{tfm:pzdr/a95}{2774}
\pdfglyphtounicode{tfm:pzdr/a96}{2775}
\pdfglyphtounicode{tfm:pzdr/a97}{275B}
\pdfglyphtounicode{tfm:pzdr/a98}{275C}
\pdfglyphtounicode{tfm:pzdr/a99}{275D}
\pdfglyphtounicode{tfm:rpxbmi/phi}{03D5}
\pdfglyphtounicode{tfm:rpxbmi/phi1}{03C6}
\pdfglyphtounicode{tfm:rpxmi/phi}{03D5}
\pdfglyphtounicode{tfm:rpxmi/phi1}{03C6}
\pdfglyphtounicode{tfm:rpzdr/a1}{2701}
\pdfglyphtounicode{tfm:rpzdr/a10}{2721}
\pdfglyphtounicode{tfm:rpzdr/a100}{275E}
\pdfglyphtounicode{tfm:rpzdr/a101}{2761}
\pdfglyphtounicode{tfm:rpzdr/a102}{2762}
\pdfglyphtounicode{tfm:rpzdr/a103}{2763}
\pdfglyphtounicode{tfm:rpzdr/a104}{2764}
\pdfglyphtounicode{tfm:rpzdr/a105}{2710}
\pdfglyphtounicode{tfm:rpzdr/a106}{2765}
\pdfglyphtounicode{tfm:rpzdr/a107}{2766}
\pdfglyphtounicode{tfm:rpzdr/a108}{2767}
\pdfglyphtounicode{tfm:rpzdr/a109}{2660}
\pdfglyphtounicode{tfm:rpzdr/a11}{261B}
\pdfglyphtounicode{tfm:rpzdr/a110}{2665}
\pdfglyphtounicode{tfm:rpzdr/a111}{2666}
\pdfglyphtounicode{tfm:rpzdr/a112}{2663}
\pdfglyphtounicode{tfm:rpzdr/a117}{2709}
\pdfglyphtounicode{tfm:rpzdr/a118}{2708}
\pdfglyphtounicode{tfm:rpzdr/a119}{2707}
\pdfglyphtounicode{tfm:rpzdr/a12}{261E}
\pdfglyphtounicode{tfm:rpzdr/a120}{2460}
\pdfglyphtounicode{tfm:rpzdr/a121}{2461}
\pdfglyphtounicode{tfm:rpzdr/a122}{2462}
\pdfglyphtounicode{tfm:rpzdr/a123}{2463}
\pdfglyphtounicode{tfm:rpzdr/a124}{2464}
\pdfglyphtounicode{tfm:rpzdr/a125}{2465}
\pdfglyphtounicode{tfm:rpzdr/a126}{2466}
\pdfglyphtounicode{tfm:rpzdr/a127}{2467}
\pdfglyphtounicode{tfm:rpzdr/a128}{2468}
\pdfglyphtounicode{tfm:rpzdr/a129}{2469}
\pdfglyphtounicode{tfm:rpzdr/a13}{270C}
\pdfglyphtounicode{tfm:rpzdr/a130}{2776}
\pdfglyphtounicode{tfm:rpzdr/a131}{2777}
\pdfglyphtounicode{tfm:rpzdr/a132}{2778}
\pdfglyphtounicode{tfm:rpzdr/a133}{2779}
\pdfglyphtounicode{tfm:rpzdr/a134}{277A}
\pdfglyphtounicode{tfm:rpzdr/a135}{277B}
\pdfglyphtounicode{tfm:rpzdr/a136}{277C}
\pdfglyphtounicode{tfm:rpzdr/a137}{277D}
\pdfglyphtounicode{tfm:rpzdr/a138}{277E}
\pdfglyphtounicode{tfm:rpzdr/a139}{277F}
\pdfglyphtounicode{tfm:rpzdr/a14}{270D}
\pdfglyphtounicode{tfm:rpzdr/a140}{2780}
\pdfglyphtounicode{tfm:rpzdr/a141}{2781}
\pdfglyphtounicode{tfm:rpzdr/a142}{2782}
\pdfglyphtounicode{tfm:rpzdr/a143}{2783}
\pdfglyphtounicode{tfm:rpzdr/a144}{2784}
\pdfglyphtounicode{tfm:rpzdr/a145}{2785}
\pdfglyphtounicode{tfm:rpzdr/a146}{2786}
\pdfglyphtounicode{tfm:rpzdr/a147}{2787}
\pdfglyphtounicode{tfm:rpzdr/a148}{2788}
\pdfglyphtounicode{tfm:rpzdr/a149}{2789}
\pdfglyphtounicode{tfm:rpzdr/a15}{270E}
\pdfglyphtounicode{tfm:rpzdr/a150}{278A}
\pdfglyphtounicode{tfm:rpzdr/a151}{278B}
\pdfglyphtounicode{tfm:rpzdr/a152}{278C}
\pdfglyphtounicode{tfm:rpzdr/a153}{278D}
\pdfglyphtounicode{tfm:rpzdr/a154}{278E}
\pdfglyphtounicode{tfm:rpzdr/a155}{278F}
\pdfglyphtounicode{tfm:rpzdr/a156}{2790}
\pdfglyphtounicode{tfm:rpzdr/a157}{2791}
\pdfglyphtounicode{tfm:rpzdr/a158}{2792}
\pdfglyphtounicode{tfm:rpzdr/a159}{2793}
\pdfglyphtounicode{tfm:rpzdr/a16}{270F}
\pdfglyphtounicode{tfm:rpzdr/a160}{2794}
\pdfglyphtounicode{tfm:rpzdr/a161}{2192}
\pdfglyphtounicode{tfm:rpzdr/a162}{27A3}
\pdfglyphtounicode{tfm:rpzdr/a163}{2194}
\pdfglyphtounicode{tfm:rpzdr/a164}{2195}
\pdfglyphtounicode{tfm:rpzdr/a165}{2799}
\pdfglyphtounicode{tfm:rpzdr/a166}{279B}
\pdfglyphtounicode{tfm:rpzdr/a167}{279C}
\pdfglyphtounicode{tfm:rpzdr/a168}{279D}
\pdfglyphtounicode{tfm:rpzdr/a169}{279E}
\pdfglyphtounicode{tfm:rpzdr/a17}{2711}
\pdfglyphtounicode{tfm:rpzdr/a170}{279F}
\pdfglyphtounicode{tfm:rpzdr/a171}{27A0}
\pdfglyphtounicode{tfm:rpzdr/a172}{27A1}
\pdfglyphtounicode{tfm:rpzdr/a173}{27A2}
\pdfglyphtounicode{tfm:rpzdr/a174}{27A4}
\pdfglyphtounicode{tfm:rpzdr/a175}{27A5}
\pdfglyphtounicode{tfm:rpzdr/a176}{27A6}
\pdfglyphtounicode{tfm:rpzdr/a177}{27A7}
\pdfglyphtounicode{tfm:rpzdr/a178}{27A8}
\pdfglyphtounicode{tfm:rpzdr/a179}{27A9}
\pdfglyphtounicode{tfm:rpzdr/a18}{2712}
\pdfglyphtounicode{tfm:rpzdr/a180}{27AB}
\pdfglyphtounicode{tfm:rpzdr/a181}{27AD}
\pdfglyphtounicode{tfm:rpzdr/a182}{27AF}
\pdfglyphtounicode{tfm:rpzdr/a183}{27B2}
\pdfglyphtounicode{tfm:rpzdr/a184}{27B3}
\pdfglyphtounicode{tfm:rpzdr/a185}{27B5}
\pdfglyphtounicode{tfm:rpzdr/a186}{27B8}
\pdfglyphtounicode{tfm:rpzdr/a187}{27BA}
\pdfglyphtounicode{tfm:rpzdr/a188}{27BB}
\pdfglyphtounicode{tfm:rpzdr/a189}{27BC}
\pdfglyphtounicode{tfm:rpzdr/a19}{2713}
\pdfglyphtounicode{tfm:rpzdr/a190}{27BD}
\pdfglyphtounicode{tfm:rpzdr/a191}{27BE}
\pdfglyphtounicode{tfm:rpzdr/a192}{279A}
\pdfglyphtounicode{tfm:rpzdr/a193}{27AA}
\pdfglyphtounicode{tfm:rpzdr/a194}{27B6}
\pdfglyphtounicode{tfm:rpzdr/a195}{27B9}
\pdfglyphtounicode{tfm:rpzdr/a196}{2798}
\pdfglyphtounicode{tfm:rpzdr/a197}{27B4}
\pdfglyphtounicode{tfm:rpzdr/a198}{27B7}
\pdfglyphtounicode{tfm:rpzdr/a199}{27AC}
\pdfglyphtounicode{tfm:rpzdr/a2}{2702}
\pdfglyphtounicode{tfm:rpzdr/a20}{2714}
\pdfglyphtounicode{tfm:rpzdr/a200}{27AE}
\pdfglyphtounicode{tfm:rpzdr/a201}{27B1}
\pdfglyphtounicode{tfm:rpzdr/a202}{2703}
\pdfglyphtounicode{tfm:rpzdr/a203}{2750}
\pdfglyphtounicode{tfm:rpzdr/a204}{2752}
\pdfglyphtounicode{tfm:rpzdr/a205}{276E}
\pdfglyphtounicode{tfm:rpzdr/a206}{2770}
\pdfglyphtounicode{tfm:rpzdr/a21}{2715}
\pdfglyphtounicode{tfm:rpzdr/a22}{2716}
\pdfglyphtounicode{tfm:rpzdr/a23}{2717}
\pdfglyphtounicode{tfm:rpzdr/a24}{2718}
\pdfglyphtounicode{tfm:rpzdr/a25}{2719}
\pdfglyphtounicode{tfm:rpzdr/a26}{271A}
\pdfglyphtounicode{tfm:rpzdr/a27}{271B}
\pdfglyphtounicode{tfm:rpzdr/a28}{271C}
\pdfglyphtounicode{tfm:rpzdr/a29}{2722}
\pdfglyphtounicode{tfm:rpzdr/a3}{2704}
\pdfglyphtounicode{tfm:rpzdr/a30}{2723}
\pdfglyphtounicode{tfm:rpzdr/a31}{2724}
\pdfglyphtounicode{tfm:rpzdr/a32}{2725}
\pdfglyphtounicode{tfm:rpzdr/a33}{2726}
\pdfglyphtounicode{tfm:rpzdr/a34}{2727}
\pdfglyphtounicode{tfm:rpzdr/a35}{2605}
\pdfglyphtounicode{tfm:rpzdr/a36}{2729}
\pdfglyphtounicode{tfm:rpzdr/a37}{272A}
\pdfglyphtounicode{tfm:rpzdr/a38}{272B}
\pdfglyphtounicode{tfm:rpzdr/a39}{272C}
\pdfglyphtounicode{tfm:rpzdr/a4}{260E}
\pdfglyphtounicode{tfm:rpzdr/a40}{272D}
\pdfglyphtounicode{tfm:rpzdr/a41}{272E}
\pdfglyphtounicode{tfm:rpzdr/a42}{272F}
\pdfglyphtounicode{tfm:rpzdr/a43}{2730}
\pdfglyphtounicode{tfm:rpzdr/a44}{2731}
\pdfglyphtounicode{tfm:rpzdr/a45}{2732}
\pdfglyphtounicode{tfm:rpzdr/a46}{2733}
\pdfglyphtounicode{tfm:rpzdr/a47}{2734}
\pdfglyphtounicode{tfm:rpzdr/a48}{2735}
\pdfglyphtounicode{tfm:rpzdr/a49}{2736}
\pdfglyphtounicode{tfm:rpzdr/a5}{2706}
\pdfglyphtounicode{tfm:rpzdr/a50}{2737}
\pdfglyphtounicode{tfm:rpzdr/a51}{2738}
\pdfglyphtounicode{tfm:rpzdr/a52}{2739}
\pdfglyphtounicode{tfm:rpzdr/a53}{273A}
\pdfglyphtounicode{tfm:rpzdr/a54}{273B}
\pdfglyphtounicode{tfm:rpzdr/a55}{273C}
\pdfglyphtounicode{tfm:rpzdr/a56}{273D}
\pdfglyphtounicode{tfm:rpzdr/a57}{273E}
\pdfglyphtounicode{tfm:rpzdr/a58}{273F}
\pdfglyphtounicode{tfm:rpzdr/a59}{2740}
\pdfglyphtounicode{tfm:rpzdr/a6}{271D}
\pdfglyphtounicode{tfm:rpzdr/a60}{2741}
\pdfglyphtounicode{tfm:rpzdr/a61}{2742}
\pdfglyphtounicode{tfm:rpzdr/a62}{2743}
\pdfglyphtounicode{tfm:rpzdr/a63}{2744}
\pdfglyphtounicode{tfm:rpzdr/a64}{2745}
\pdfglyphtounicode{tfm:rpzdr/a65}{2746}
\pdfglyphtounicode{tfm:rpzdr/a66}{2747}
\pdfglyphtounicode{tfm:rpzdr/a67}{2748}
\pdfglyphtounicode{tfm:rpzdr/a68}{2749}
\pdfglyphtounicode{tfm:rpzdr/a69}{274A}
\pdfglyphtounicode{tfm:rpzdr/a7}{271E}
\pdfglyphtounicode{tfm:rpzdr/a70}{274B}
\pdfglyphtounicode{tfm:rpzdr/a71}{25CF}
\pdfglyphtounicode{tfm:rpzdr/a72}{274D}
\pdfglyphtounicode{tfm:rpzdr/a73}{25A0}
\pdfglyphtounicode{tfm:rpzdr/a74}{274F}
\pdfglyphtounicode{tfm:rpzdr/a75}{2751}
\pdfglyphtounicode{tfm:rpzdr/a76}{25B2}
\pdfglyphtounicode{tfm:rpzdr/a77}{25BC}
\pdfglyphtounicode{tfm:rpzdr/a78}{25C6}
\pdfglyphtounicode{tfm:rpzdr/a79}{2756}
\pdfglyphtounicode{tfm:rpzdr/a8}{271F}
\pdfglyphtounicode{tfm:rpzdr/a81}{25D7}
\pdfglyphtounicode{tfm:rpzdr/a82}{2758}
\pdfglyphtounicode{tfm:rpzdr/a83}{2759}
\pdfglyphtounicode{tfm:rpzdr/a84}{275A}
\pdfglyphtounicode{tfm:rpzdr/a85}{276F}
\pdfglyphtounicode{tfm:rpzdr/a86}{2771}
\pdfglyphtounicode{tfm:rpzdr/a87}{2772}
\pdfglyphtounicode{tfm:rpzdr/a88}{2773}
\pdfglyphtounicode{tfm:rpzdr/a89}{2768}
\pdfglyphtounicode{tfm:rpzdr/a9}{2720}
\pdfglyphtounicode{tfm:rpzdr/a90}{2769}
\pdfglyphtounicode{tfm:rpzdr/a91}{276C}
\pdfglyphtounicode{tfm:rpzdr/a92}{276D}
\pdfglyphtounicode{tfm:rpzdr/a93}{276A}
\pdfglyphtounicode{tfm:rpzdr/a94}{276B}
\pdfglyphtounicode{tfm:rpzdr/a95}{2774}
\pdfglyphtounicode{tfm:rpzdr/a96}{2775}
\pdfglyphtounicode{tfm:rpzdr/a97}{275B}
\pdfglyphtounicode{tfm:rpzdr/a98}{275C}
\pdfglyphtounicode{tfm:rpzdr/a99}{275D}
\pdfglyphtounicode{tfm:rtxbmi/phi}{03D5}
\pdfglyphtounicode{tfm:rtxbmi/phi1}{03C6}
\pdfglyphtounicode{tfm:rtxmi/phi}{03D5}
\pdfglyphtounicode{tfm:rtxmi/phi1}{03C6}
\pdfglyphtounicode{tfm:txbmia/phi}{03D5}
\pdfglyphtounicode{tfm:txbmia/phi1}{03C6}
\pdfglyphtounicode{tfm:txbsy/diamond}{2662}
\pdfglyphtounicode{tfm:txbsy/heart}{2661}
\pdfglyphtounicode{tfm:txbsya/diamond}{2662}
\pdfglyphtounicode{tfm:txmia/phi}{03D5}
\pdfglyphtounicode{tfm:txmia/phi1}{03C6}
\pdfglyphtounicode{tfm:txsy/diamond}{2662}
\pdfglyphtounicode{tfm:txsy/heart}{2661}
\pdfglyphtounicode{tfm:txsya/diamond}{2662}
\pdfglyphtounicode{tfm:zd/a1}{2701}
\pdfglyphtounicode{tfm:zd/a10}{2721}
\pdfglyphtounicode{tfm:zd/a100}{275E}
\pdfglyphtounicode{tfm:zd/a101}{2761}
\pdfglyphtounicode{tfm:zd/a102}{2762}
\pdfglyphtounicode{tfm:zd/a103}{2763}
\pdfglyphtounicode{tfm:zd/a104}{2764}
\pdfglyphtounicode{tfm:zd/a105}{2710}
\pdfglyphtounicode{tfm:zd/a106}{2765}
\pdfglyphtounicode{tfm:zd/a107}{2766}
\pdfglyphtounicode{tfm:zd/a108}{2767}
\pdfglyphtounicode{tfm:zd/a109}{2660}
\pdfglyphtounicode{tfm:zd/a11}{261B}
\pdfglyphtounicode{tfm:zd/a110}{2665}
\pdfglyphtounicode{tfm:zd/a111}{2666}
\pdfglyphtounicode{tfm:zd/a112}{2663}
\pdfglyphtounicode{tfm:zd/a117}{2709}
\pdfglyphtounicode{tfm:zd/a118}{2708}
\pdfglyphtounicode{tfm:zd/a119}{2707}
\pdfglyphtounicode{tfm:zd/a12}{261E}
\pdfglyphtounicode{tfm:zd/a120}{2460}
\pdfglyphtounicode{tfm:zd/a121}{2461}
\pdfglyphtounicode{tfm:zd/a122}{2462}
\pdfglyphtounicode{tfm:zd/a123}{2463}
\pdfglyphtounicode{tfm:zd/a124}{2464}
\pdfglyphtounicode{tfm:zd/a125}{2465}
\pdfglyphtounicode{tfm:zd/a126}{2466}
\pdfglyphtounicode{tfm:zd/a127}{2467}
\pdfglyphtounicode{tfm:zd/a128}{2468}
\pdfglyphtounicode{tfm:zd/a129}{2469}
\pdfglyphtounicode{tfm:zd/a13}{270C}
\pdfglyphtounicode{tfm:zd/a130}{2776}
\pdfglyphtounicode{tfm:zd/a131}{2777}
\pdfglyphtounicode{tfm:zd/a132}{2778}
\pdfglyphtounicode{tfm:zd/a133}{2779}
\pdfglyphtounicode{tfm:zd/a134}{277A}
\pdfglyphtounicode{tfm:zd/a135}{277B}
\pdfglyphtounicode{tfm:zd/a136}{277C}
\pdfglyphtounicode{tfm:zd/a137}{277D}
\pdfglyphtounicode{tfm:zd/a138}{277E}
\pdfglyphtounicode{tfm:zd/a139}{277F}
\pdfglyphtounicode{tfm:zd/a14}{270D}
\pdfglyphtounicode{tfm:zd/a140}{2780}
\pdfglyphtounicode{tfm:zd/a141}{2781}
\pdfglyphtounicode{tfm:zd/a142}{2782}
\pdfglyphtounicode{tfm:zd/a143}{2783}
\pdfglyphtounicode{tfm:zd/a144}{2784}
\pdfglyphtounicode{tfm:zd/a145}{2785}
\pdfglyphtounicode{tfm:zd/a146}{2786}
\pdfglyphtounicode{tfm:zd/a147}{2787}
\pdfglyphtounicode{tfm:zd/a148}{2788}
\pdfglyphtounicode{tfm:zd/a149}{2789}
\pdfglyphtounicode{tfm:zd/a15}{270E}
\pdfglyphtounicode{tfm:zd/a150}{278A}
\pdfglyphtounicode{tfm:zd/a151}{278B}
\pdfglyphtounicode{tfm:zd/a152}{278C}
\pdfglyphtounicode{tfm:zd/a153}{278D}
\pdfglyphtounicode{tfm:zd/a154}{278E}
\pdfglyphtounicode{tfm:zd/a155}{278F}
\pdfglyphtounicode{tfm:zd/a156}{2790}
\pdfglyphtounicode{tfm:zd/a157}{2791}
\pdfglyphtounicode{tfm:zd/a158}{2792}
\pdfglyphtounicode{tfm:zd/a159}{2793}
\pdfglyphtounicode{tfm:zd/a16}{270F}
\pdfglyphtounicode{tfm:zd/a160}{2794}
\pdfglyphtounicode{tfm:zd/a161}{2192}
\pdfglyphtounicode{tfm:zd/a162}{27A3}
\pdfglyphtounicode{tfm:zd/a163}{2194}
\pdfglyphtounicode{tfm:zd/a164}{2195}
\pdfglyphtounicode{tfm:zd/a165}{2799}
\pdfglyphtounicode{tfm:zd/a166}{279B}
\pdfglyphtounicode{tfm:zd/a167}{279C}
\pdfglyphtounicode{tfm:zd/a168}{279D}
\pdfglyphtounicode{tfm:zd/a169}{279E}
\pdfglyphtounicode{tfm:zd/a17}{2711}
\pdfglyphtounicode{tfm:zd/a170}{279F}
\pdfglyphtounicode{tfm:zd/a171}{27A0}
\pdfglyphtounicode{tfm:zd/a172}{27A1}
\pdfglyphtounicode{tfm:zd/a173}{27A2}
\pdfglyphtounicode{tfm:zd/a174}{27A4}
\pdfglyphtounicode{tfm:zd/a175}{27A5}
\pdfglyphtounicode{tfm:zd/a176}{27A6}
\pdfglyphtounicode{tfm:zd/a177}{27A7}
\pdfglyphtounicode{tfm:zd/a178}{27A8}
\pdfglyphtounicode{tfm:zd/a179}{27A9}
\pdfglyphtounicode{tfm:zd/a18}{2712}
\pdfglyphtounicode{tfm:zd/a180}{27AB}
\pdfglyphtounicode{tfm:zd/a181}{27AD}
\pdfglyphtounicode{tfm:zd/a182}{27AF}
\pdfglyphtounicode{tfm:zd/a183}{27B2}
\pdfglyphtounicode{tfm:zd/a184}{27B3}
\pdfglyphtounicode{tfm:zd/a185}{27B5}
\pdfglyphtounicode{tfm:zd/a186}{27B8}
\pdfglyphtounicode{tfm:zd/a187}{27BA}
\pdfglyphtounicode{tfm:zd/a188}{27BB}
\pdfglyphtounicode{tfm:zd/a189}{27BC}
\pdfglyphtounicode{tfm:zd/a19}{2713}
\pdfglyphtounicode{tfm:zd/a190}{27BD}
\pdfglyphtounicode{tfm:zd/a191}{27BE}
\pdfglyphtounicode{tfm:zd/a192}{279A}
\pdfglyphtounicode{tfm:zd/a193}{27AA}
\pdfglyphtounicode{tfm:zd/a194}{27B6}
\pdfglyphtounicode{tfm:zd/a195}{27B9}
\pdfglyphtounicode{tfm:zd/a196}{2798}
\pdfglyphtounicode{tfm:zd/a197}{27B4}
\pdfglyphtounicode{tfm:zd/a198}{27B7}
\pdfglyphtounicode{tfm:zd/a199}{27AC}
\pdfglyphtounicode{tfm:zd/a2}{2702}
\pdfglyphtounicode{tfm:zd/a20}{2714}
\pdfglyphtounicode{tfm:zd/a200}{27AE}
\pdfglyphtounicode{tfm:zd/a201}{27B1}
\pdfglyphtounicode{tfm:zd/a202}{2703}
\pdfglyphtounicode{tfm:zd/a203}{2750}
\pdfglyphtounicode{tfm:zd/a204}{2752}
\pdfglyphtounicode{tfm:zd/a205}{276E}
\pdfglyphtounicode{tfm:zd/a206}{2770}
\pdfglyphtounicode{tfm:zd/a21}{2715}
\pdfglyphtounicode{tfm:zd/a22}{2716}
\pdfglyphtounicode{tfm:zd/a23}{2717}
\pdfglyphtounicode{tfm:zd/a24}{2718}
\pdfglyphtounicode{tfm:zd/a25}{2719}
\pdfglyphtounicode{tfm:zd/a26}{271A}
\pdfglyphtounicode{tfm:zd/a27}{271B}
\pdfglyphtounicode{tfm:zd/a28}{271C}
\pdfglyphtounicode{tfm:zd/a29}{2722}
\pdfglyphtounicode{tfm:zd/a3}{2704}
\pdfglyphtounicode{tfm:zd/a30}{2723}
\pdfglyphtounicode{tfm:zd/a31}{2724}
\pdfglyphtounicode{tfm:zd/a32}{2725}
\pdfglyphtounicode{tfm:zd/a33}{2726}
\pdfglyphtounicode{tfm:zd/a34}{2727}
\pdfglyphtounicode{tfm:zd/a35}{2605}
\pdfglyphtounicode{tfm:zd/a36}{2729}
\pdfglyphtounicode{tfm:zd/a37}{272A}
\pdfglyphtounicode{tfm:zd/a38}{272B}
\pdfglyphtounicode{tfm:zd/a39}{272C}
\pdfglyphtounicode{tfm:zd/a4}{260E}
\pdfglyphtounicode{tfm:zd/a40}{272D}
\pdfglyphtounicode{tfm:zd/a41}{272E}
\pdfglyphtounicode{tfm:zd/a42}{272F}
\pdfglyphtounicode{tfm:zd/a43}{2730}
\pdfglyphtounicode{tfm:zd/a44}{2731}
\pdfglyphtounicode{tfm:zd/a45}{2732}
\pdfglyphtounicode{tfm:zd/a46}{2733}
\pdfglyphtounicode{tfm:zd/a47}{2734}
\pdfglyphtounicode{tfm:zd/a48}{2735}
\pdfglyphtounicode{tfm:zd/a49}{2736}
\pdfglyphtounicode{tfm:zd/a5}{2706}
\pdfglyphtounicode{tfm:zd/a50}{2737}
\pdfglyphtounicode{tfm:zd/a51}{2738}
\pdfglyphtounicode{tfm:zd/a52}{2739}
\pdfglyphtounicode{tfm:zd/a53}{273A}
\pdfglyphtounicode{tfm:zd/a54}{273B}
\pdfglyphtounicode{tfm:zd/a55}{273C}
\pdfglyphtounicode{tfm:zd/a56}{273D}
\pdfglyphtounicode{tfm:zd/a57}{273E}
\pdfglyphtounicode{tfm:zd/a58}{273F}
\pdfglyphtounicode{tfm:zd/a59}{2740}
\pdfglyphtounicode{tfm:zd/a6}{271D}
\pdfglyphtounicode{tfm:zd/a60}{2741}
\pdfglyphtounicode{tfm:zd/a61}{2742}
\pdfglyphtounicode{tfm:zd/a62}{2743}
\pdfglyphtounicode{tfm:zd/a63}{2744}
\pdfglyphtounicode{tfm:zd/a64}{2745}
\pdfglyphtounicode{tfm:zd/a65}{2746}
\pdfglyphtounicode{tfm:zd/a66}{2747}
\pdfglyphtounicode{tfm:zd/a67}{2748}
\pdfglyphtounicode{tfm:zd/a68}{2749}
\pdfglyphtounicode{tfm:zd/a69}{274A}
\pdfglyphtounicode{tfm:zd/a7}{271E}
\pdfglyphtounicode{tfm:zd/a70}{274B}
\pdfglyphtounicode{tfm:zd/a71}{25CF}
\pdfglyphtounicode{tfm:zd/a72}{274D}
\pdfglyphtounicode{tfm:zd/a73}{25A0}
\pdfglyphtounicode{tfm:zd/a74}{274F}
\pdfglyphtounicode{tfm:zd/a75}{2751}
\pdfglyphtounicode{tfm:zd/a76}{25B2}
\pdfglyphtounicode{tfm:zd/a77}{25BC}
\pdfglyphtounicode{tfm:zd/a78}{25C6}
\pdfglyphtounicode{tfm:zd/a79}{2756}
\pdfglyphtounicode{tfm:zd/a8}{271F}
\pdfglyphtounicode{tfm:zd/a81}{25D7}
\pdfglyphtounicode{tfm:zd/a82}{2758}
\pdfglyphtounicode{tfm:zd/a83}{2759}
\pdfglyphtounicode{tfm:zd/a84}{275A}
\pdfglyphtounicode{tfm:zd/a85}{276F}
\pdfglyphtounicode{tfm:zd/a86}{2771}
\pdfglyphtounicode{tfm:zd/a87}{2772}
\pdfglyphtounicode{tfm:zd/a88}{2773}
\pdfglyphtounicode{tfm:zd/a89}{2768}
\pdfglyphtounicode{tfm:zd/a9}{2720}
\pdfglyphtounicode{tfm:zd/a90}{2769}
\pdfglyphtounicode{tfm:zd/a91}{276C}
\pdfglyphtounicode{tfm:zd/a92}{276D}
\pdfglyphtounicode{tfm:zd/a93}{276A}
\pdfglyphtounicode{tfm:zd/a94}{276B}
\pdfglyphtounicode{tfm:zd/a95}{2774}
\pdfglyphtounicode{tfm:zd/a96}{2775}
\pdfglyphtounicode{tfm:zd/a97}{275B}
\pdfglyphtounicode{tfm:zd/a98}{275C}
\pdfglyphtounicode{tfm:zd/a99}{275D}
\pdfglyphtounicode{tfm:zpzdr-reversed/a1}{2701}
\pdfglyphtounicode{tfm:zpzdr-reversed/a10}{2721}
\pdfglyphtounicode{tfm:zpzdr-reversed/a100}{275E}
\pdfglyphtounicode{tfm:zpzdr-reversed/a101}{2761}
\pdfglyphtounicode{tfm:zpzdr-reversed/a102}{2762}
\pdfglyphtounicode{tfm:zpzdr-reversed/a103}{2763}
\pdfglyphtounicode{tfm:zpzdr-reversed/a104}{2764}
\pdfglyphtounicode{tfm:zpzdr-reversed/a105}{2710}
\pdfglyphtounicode{tfm:zpzdr-reversed/a106}{2765}
\pdfglyphtounicode{tfm:zpzdr-reversed/a107}{2766}
\pdfglyphtounicode{tfm:zpzdr-reversed/a108}{2767}
\pdfglyphtounicode{tfm:zpzdr-reversed/a109}{2660}
\pdfglyphtounicode{tfm:zpzdr-reversed/a11}{261B}
\pdfglyphtounicode{tfm:zpzdr-reversed/a110}{2665}
\pdfglyphtounicode{tfm:zpzdr-reversed/a111}{2666}
\pdfglyphtounicode{tfm:zpzdr-reversed/a112}{2663}
\pdfglyphtounicode{tfm:zpzdr-reversed/a117}{2709}
\pdfglyphtounicode{tfm:zpzdr-reversed/a118}{2708}
\pdfglyphtounicode{tfm:zpzdr-reversed/a119}{2707}
\pdfglyphtounicode{tfm:zpzdr-reversed/a12}{261E}
\pdfglyphtounicode{tfm:zpzdr-reversed/a120}{2460}
\pdfglyphtounicode{tfm:zpzdr-reversed/a121}{2461}
\pdfglyphtounicode{tfm:zpzdr-reversed/a122}{2462}
\pdfglyphtounicode{tfm:zpzdr-reversed/a123}{2463}
\pdfglyphtounicode{tfm:zpzdr-reversed/a124}{2464}
\pdfglyphtounicode{tfm:zpzdr-reversed/a125}{2465}
\pdfglyphtounicode{tfm:zpzdr-reversed/a126}{2466}
\pdfglyphtounicode{tfm:zpzdr-reversed/a127}{2467}
\pdfglyphtounicode{tfm:zpzdr-reversed/a128}{2468}
\pdfglyphtounicode{tfm:zpzdr-reversed/a129}{2469}
\pdfglyphtounicode{tfm:zpzdr-reversed/a13}{270C}
\pdfglyphtounicode{tfm:zpzdr-reversed/a130}{2776}
\pdfglyphtounicode{tfm:zpzdr-reversed/a131}{2777}
\pdfglyphtounicode{tfm:zpzdr-reversed/a132}{2778}
\pdfglyphtounicode{tfm:zpzdr-reversed/a133}{2779}
\pdfglyphtounicode{tfm:zpzdr-reversed/a134}{277A}
\pdfglyphtounicode{tfm:zpzdr-reversed/a135}{277B}
\pdfglyphtounicode{tfm:zpzdr-reversed/a136}{277C}
\pdfglyphtounicode{tfm:zpzdr-reversed/a137}{277D}
\pdfglyphtounicode{tfm:zpzdr-reversed/a138}{277E}
\pdfglyphtounicode{tfm:zpzdr-reversed/a139}{277F}
\pdfglyphtounicode{tfm:zpzdr-reversed/a14}{270D}
\pdfglyphtounicode{tfm:zpzdr-reversed/a140}{2780}
\pdfglyphtounicode{tfm:zpzdr-reversed/a141}{2781}
\pdfglyphtounicode{tfm:zpzdr-reversed/a142}{2782}
\pdfglyphtounicode{tfm:zpzdr-reversed/a143}{2783}
\pdfglyphtounicode{tfm:zpzdr-reversed/a144}{2784}
\pdfglyphtounicode{tfm:zpzdr-reversed/a145}{2785}
\pdfglyphtounicode{tfm:zpzdr-reversed/a146}{2786}
\pdfglyphtounicode{tfm:zpzdr-reversed/a147}{2787}
\pdfglyphtounicode{tfm:zpzdr-reversed/a148}{2788}
\pdfglyphtounicode{tfm:zpzdr-reversed/a149}{2789}
\pdfglyphtounicode{tfm:zpzdr-reversed/a15}{270E}
\pdfglyphtounicode{tfm:zpzdr-reversed/a150}{278A}
\pdfglyphtounicode{tfm:zpzdr-reversed/a151}{278B}
\pdfglyphtounicode{tfm:zpzdr-reversed/a152}{278C}
\pdfglyphtounicode{tfm:zpzdr-reversed/a153}{278D}
\pdfglyphtounicode{tfm:zpzdr-reversed/a154}{278E}
\pdfglyphtounicode{tfm:zpzdr-reversed/a155}{278F}
\pdfglyphtounicode{tfm:zpzdr-reversed/a156}{2790}
\pdfglyphtounicode{tfm:zpzdr-reversed/a157}{2791}
\pdfglyphtounicode{tfm:zpzdr-reversed/a158}{2792}
\pdfglyphtounicode{tfm:zpzdr-reversed/a159}{2793}
\pdfglyphtounicode{tfm:zpzdr-reversed/a16}{270F}
\pdfglyphtounicode{tfm:zpzdr-reversed/a160}{2794}
\pdfglyphtounicode{tfm:zpzdr-reversed/a161}{2192}
\pdfglyphtounicode{tfm:zpzdr-reversed/a162}{27A3}
\pdfglyphtounicode{tfm:zpzdr-reversed/a163}{2194}
\pdfglyphtounicode{tfm:zpzdr-reversed/a164}{2195}
\pdfglyphtounicode{tfm:zpzdr-reversed/a165}{2799}
\pdfglyphtounicode{tfm:zpzdr-reversed/a166}{279B}
\pdfglyphtounicode{tfm:zpzdr-reversed/a167}{279C}
\pdfglyphtounicode{tfm:zpzdr-reversed/a168}{279D}
\pdfglyphtounicode{tfm:zpzdr-reversed/a169}{279E}
\pdfglyphtounicode{tfm:zpzdr-reversed/a17}{2711}
\pdfglyphtounicode{tfm:zpzdr-reversed/a170}{279F}
\pdfglyphtounicode{tfm:zpzdr-reversed/a171}{27A0}
\pdfglyphtounicode{tfm:zpzdr-reversed/a172}{27A1}
\pdfglyphtounicode{tfm:zpzdr-reversed/a173}{27A2}
\pdfglyphtounicode{tfm:zpzdr-reversed/a174}{27A4}
\pdfglyphtounicode{tfm:zpzdr-reversed/a175}{27A5}
\pdfglyphtounicode{tfm:zpzdr-reversed/a176}{27A6}
\pdfglyphtounicode{tfm:zpzdr-reversed/a177}{27A7}
\pdfglyphtounicode{tfm:zpzdr-reversed/a178}{27A8}
\pdfglyphtounicode{tfm:zpzdr-reversed/a179}{27A9}
\pdfglyphtounicode{tfm:zpzdr-reversed/a18}{2712}
\pdfglyphtounicode{tfm:zpzdr-reversed/a180}{27AB}
\pdfglyphtounicode{tfm:zpzdr-reversed/a181}{27AD}
\pdfglyphtounicode{tfm:zpzdr-reversed/a182}{27AF}
\pdfglyphtounicode{tfm:zpzdr-reversed/a183}{27B2}
\pdfglyphtounicode{tfm:zpzdr-reversed/a184}{27B3}
\pdfglyphtounicode{tfm:zpzdr-reversed/a185}{27B5}
\pdfglyphtounicode{tfm:zpzdr-reversed/a186}{27B8}
\pdfglyphtounicode{tfm:zpzdr-reversed/a187}{27BA}
\pdfglyphtounicode{tfm:zpzdr-reversed/a188}{27BB}
\pdfglyphtounicode{tfm:zpzdr-reversed/a189}{27BC}
\pdfglyphtounicode{tfm:zpzdr-reversed/a19}{2713}
\pdfglyphtounicode{tfm:zpzdr-reversed/a190}{27BD}
\pdfglyphtounicode{tfm:zpzdr-reversed/a191}{27BE}
\pdfglyphtounicode{tfm:zpzdr-reversed/a192}{279A}
\pdfglyphtounicode{tfm:zpzdr-reversed/a193}{27AA}
\pdfglyphtounicode{tfm:zpzdr-reversed/a194}{27B6}
\pdfglyphtounicode{tfm:zpzdr-reversed/a195}{27B9}
\pdfglyphtounicode{tfm:zpzdr-reversed/a196}{2798}
\pdfglyphtounicode{tfm:zpzdr-reversed/a197}{27B4}
\pdfglyphtounicode{tfm:zpzdr-reversed/a198}{27B7}
\pdfglyphtounicode{tfm:zpzdr-reversed/a199}{27AC}
\pdfglyphtounicode{tfm:zpzdr-reversed/a2}{2702}
\pdfglyphtounicode{tfm:zpzdr-reversed/a20}{2714}
\pdfglyphtounicode{tfm:zpzdr-reversed/a200}{27AE}
\pdfglyphtounicode{tfm:zpzdr-reversed/a201}{27B1}
\pdfglyphtounicode{tfm:zpzdr-reversed/a202}{2703}
\pdfglyphtounicode{tfm:zpzdr-reversed/a203}{2750}
\pdfglyphtounicode{tfm:zpzdr-reversed/a204}{2752}
\pdfglyphtounicode{tfm:zpzdr-reversed/a205}{276E}
\pdfglyphtounicode{tfm:zpzdr-reversed/a206}{2770}
\pdfglyphtounicode{tfm:zpzdr-reversed/a21}{2715}
\pdfglyphtounicode{tfm:zpzdr-reversed/a22}{2716}
\pdfglyphtounicode{tfm:zpzdr-reversed/a23}{2717}
\pdfglyphtounicode{tfm:zpzdr-reversed/a24}{2718}
\pdfglyphtounicode{tfm:zpzdr-reversed/a25}{2719}
\pdfglyphtounicode{tfm:zpzdr-reversed/a26}{271A}
\pdfglyphtounicode{tfm:zpzdr-reversed/a27}{271B}
\pdfglyphtounicode{tfm:zpzdr-reversed/a28}{271C}
\pdfglyphtounicode{tfm:zpzdr-reversed/a29}{2722}
\pdfglyphtounicode{tfm:zpzdr-reversed/a3}{2704}
\pdfglyphtounicode{tfm:zpzdr-reversed/a30}{2723}
\pdfglyphtounicode{tfm:zpzdr-reversed/a31}{2724}
\pdfglyphtounicode{tfm:zpzdr-reversed/a32}{2725}
\pdfglyphtounicode{tfm:zpzdr-reversed/a33}{2726}
\pdfglyphtounicode{tfm:zpzdr-reversed/a34}{2727}
\pdfglyphtounicode{tfm:zpzdr-reversed/a35}{2605}
\pdfglyphtounicode{tfm:zpzdr-reversed/a36}{2729}
\pdfglyphtounicode{tfm:zpzdr-reversed/a37}{272A}
\pdfglyphtounicode{tfm:zpzdr-reversed/a38}{272B}
\pdfglyphtounicode{tfm:zpzdr-reversed/a39}{272C}
\pdfglyphtounicode{tfm:zpzdr-reversed/a4}{260E}
\pdfglyphtounicode{tfm:zpzdr-reversed/a40}{272D}
\pdfglyphtounicode{tfm:zpzdr-reversed/a41}{272E}
\pdfglyphtounicode{tfm:zpzdr-reversed/a42}{272F}
\pdfglyphtounicode{tfm:zpzdr-reversed/a43}{2730}
\pdfglyphtounicode{tfm:zpzdr-reversed/a44}{2731}
\pdfglyphtounicode{tfm:zpzdr-reversed/a45}{2732}
\pdfglyphtounicode{tfm:zpzdr-reversed/a46}{2733}
\pdfglyphtounicode{tfm:zpzdr-reversed/a47}{2734}
\pdfglyphtounicode{tfm:zpzdr-reversed/a48}{2735}
\pdfglyphtounicode{tfm:zpzdr-reversed/a49}{2736}
\pdfglyphtounicode{tfm:zpzdr-reversed/a5}{2706}
\pdfglyphtounicode{tfm:zpzdr-reversed/a50}{2737}
\pdfglyphtounicode{tfm:zpzdr-reversed/a51}{2738}
\pdfglyphtounicode{tfm:zpzdr-reversed/a52}{2739}
\pdfglyphtounicode{tfm:zpzdr-reversed/a53}{273A}
\pdfglyphtounicode{tfm:zpzdr-reversed/a54}{273B}
\pdfglyphtounicode{tfm:zpzdr-reversed/a55}{273C}
\pdfglyphtounicode{tfm:zpzdr-reversed/a56}{273D}
\pdfglyphtounicode{tfm:zpzdr-reversed/a57}{273E}
\pdfglyphtounicode{tfm:zpzdr-reversed/a58}{273F}
\pdfglyphtounicode{tfm:zpzdr-reversed/a59}{2740}
\pdfglyphtounicode{tfm:zpzdr-reversed/a6}{271D}
\pdfglyphtounicode{tfm:zpzdr-reversed/a60}{2741}
\pdfglyphtounicode{tfm:zpzdr-reversed/a61}{2742}
\pdfglyphtounicode{tfm:zpzdr-reversed/a62}{2743}
\pdfglyphtounicode{tfm:zpzdr-reversed/a63}{2744}
\pdfglyphtounicode{tfm:zpzdr-reversed/a64}{2745}
\pdfglyphtounicode{tfm:zpzdr-reversed/a65}{2746}
\pdfglyphtounicode{tfm:zpzdr-reversed/a66}{2747}
\pdfglyphtounicode{tfm:zpzdr-reversed/a67}{2748}
\pdfglyphtounicode{tfm:zpzdr-reversed/a68}{2749}
\pdfglyphtounicode{tfm:zpzdr-reversed/a69}{274A}
\pdfglyphtounicode{tfm:zpzdr-reversed/a7}{271E}
\pdfglyphtounicode{tfm:zpzdr-reversed/a70}{274B}
\pdfglyphtounicode{tfm:zpzdr-reversed/a71}{25CF}
\pdfglyphtounicode{tfm:zpzdr-reversed/a72}{274D}
\pdfglyphtounicode{tfm:zpzdr-reversed/a73}{25A0}
\pdfglyphtounicode{tfm:zpzdr-reversed/a74}{274F}
\pdfglyphtounicode{tfm:zpzdr-reversed/a75}{2751}
\pdfglyphtounicode{tfm:zpzdr-reversed/a76}{25B2}
\pdfglyphtounicode{tfm:zpzdr-reversed/a77}{25BC}
\pdfglyphtounicode{tfm:zpzdr-reversed/a78}{25C6}
\pdfglyphtounicode{tfm:zpzdr-reversed/a79}{2756}
\pdfglyphtounicode{tfm:zpzdr-reversed/a8}{271F}
\pdfglyphtounicode{tfm:zpzdr-reversed/a81}{25D7}
\pdfglyphtounicode{tfm:zpzdr-reversed/a82}{2758}
\pdfglyphtounicode{tfm:zpzdr-reversed/a83}{2759}
\pdfglyphtounicode{tfm:zpzdr-reversed/a84}{275A}
\pdfglyphtounicode{tfm:zpzdr-reversed/a85}{276F}
\pdfglyphtounicode{tfm:zpzdr-reversed/a86}{2771}
\pdfglyphtounicode{tfm:zpzdr-reversed/a87}{2772}
\pdfglyphtounicode{tfm:zpzdr-reversed/a88}{2773}
\pdfglyphtounicode{tfm:zpzdr-reversed/a89}{2768}
\pdfglyphtounicode{tfm:zpzdr-reversed/a9}{2720}
\pdfglyphtounicode{tfm:zpzdr-reversed/a90}{2769}
\pdfglyphtounicode{tfm:zpzdr-reversed/a91}{276C}
\pdfglyphtounicode{tfm:zpzdr-reversed/a92}{276D}
\pdfglyphtounicode{tfm:zpzdr-reversed/a93}{276A}
\pdfglyphtounicode{tfm:zpzdr-reversed/a94}{276B}
\pdfglyphtounicode{tfm:zpzdr-reversed/a95}{2774}
\pdfglyphtounicode{tfm:zpzdr-reversed/a96}{2775}
\pdfglyphtounicode{tfm:zpzdr-reversed/a97}{275B}
\pdfglyphtounicode{tfm:zpzdr-reversed/a98}{275C}
\pdfglyphtounicode{tfm:zpzdr-reversed/a99}{275D}
\pdfglyphtounicode{thabengali}{09A5}
\pdfglyphtounicode{thadeva}{0925}
\pdfglyphtounicode{thagujarati}{0AA5}
\pdfglyphtounicode{thagurmukhi}{0A25}
\pdfglyphtounicode{thalarabic}{0630}
\pdfglyphtounicode{thalfinalarabic}{FEAC}
\pdfglyphtounicode{thanthakhatlowleftthai}{F898}
\pdfglyphtounicode{thanthakhatlowrightthai}{F897}
\pdfglyphtounicode{thanthakhatthai}{0E4C}
\pdfglyphtounicode{thanthakhatupperleftthai}{F896}
\pdfglyphtounicode{theharabic}{062B}
\pdfglyphtounicode{thehfinalarabic}{FE9A}
\pdfglyphtounicode{thehinitialarabic}{FE9B}
\pdfglyphtounicode{thehmedialarabic}{FE9C}
\pdfglyphtounicode{thereexists}{2203}
\pdfglyphtounicode{therefore}{2234}
\pdfglyphtounicode{theta}{03B8}
\pdfglyphtounicode{theta1}{03D1}
\pdfglyphtounicode{thetasymbolgreek}{03D1}
\pdfglyphtounicode{thieuthacirclekorean}{3279}
\pdfglyphtounicode{thieuthaparenkorean}{3219}
\pdfglyphtounicode{thieuthcirclekorean}{326B}
\pdfglyphtounicode{thieuthkorean}{314C}
\pdfglyphtounicode{thieuthparenkorean}{320B}
\pdfglyphtounicode{thirteencircle}{246C}
\pdfglyphtounicode{thirteenparen}{2480}
\pdfglyphtounicode{thirteenperiod}{2494}
\pdfglyphtounicode{thonangmonthothai}{0E11}
\pdfglyphtounicode{thook}{01AD}
\pdfglyphtounicode{thophuthaothai}{0E12}
\pdfglyphtounicode{thorn}{00FE}
\pdfglyphtounicode{thothahanthai}{0E17}
\pdfglyphtounicode{thothanthai}{0E10}
\pdfglyphtounicode{thothongthai}{0E18}
\pdfglyphtounicode{thothungthai}{0E16}
\pdfglyphtounicode{thousandcyrillic}{0482}
\pdfglyphtounicode{thousandsseparatorarabic}{066C}
\pdfglyphtounicode{thousandsseparatorpersian}{066C}
\pdfglyphtounicode{three}{0033}
\pdfglyphtounicode{threearabic}{0663}
\pdfglyphtounicode{threebengali}{09E9}
\pdfglyphtounicode{threecircle}{2462}
\pdfglyphtounicode{threecircleinversesansserif}{278C}
\pdfglyphtounicode{threedeva}{0969}
\pdfglyphtounicode{threeeighths}{215C}
\pdfglyphtounicode{threegujarati}{0AE9}
\pdfglyphtounicode{threegurmukhi}{0A69}
\pdfglyphtounicode{threehackarabic}{0663}
\pdfglyphtounicode{threehangzhou}{3023}
\pdfglyphtounicode{threeideographicparen}{3222}
\pdfglyphtounicode{threeinferior}{2083}
\pdfglyphtounicode{threemonospace}{FF13}
\pdfglyphtounicode{threenumeratorbengali}{09F6}
\pdfglyphtounicode{threeoldstyle}{0033}
\pdfglyphtounicode{threeparen}{2476}
\pdfglyphtounicode{threeperiod}{248A}
\pdfglyphtounicode{threepersian}{06F3}
\pdfglyphtounicode{threequarters}{00BE}
\pdfglyphtounicode{threequartersemdash}{F6DE}
\pdfglyphtounicode{threeroman}{2172}
\pdfglyphtounicode{threesuperior}{00B3}
\pdfglyphtounicode{threethai}{0E53}
\pdfglyphtounicode{thzsquare}{3394}
\pdfglyphtounicode{tihiragana}{3061}
\pdfglyphtounicode{tikatakana}{30C1}
\pdfglyphtounicode{tikatakanahalfwidth}{FF81}
\pdfglyphtounicode{tikeutacirclekorean}{3270}
\pdfglyphtounicode{tikeutaparenkorean}{3210}
\pdfglyphtounicode{tikeutcirclekorean}{3262}
\pdfglyphtounicode{tikeutkorean}{3137}
\pdfglyphtounicode{tikeutparenkorean}{3202}
\pdfglyphtounicode{tilde}{02DC}
\pdfglyphtounicode{tildebelowcmb}{0330}
\pdfglyphtounicode{tildecmb}{0303}
\pdfglyphtounicode{tildecomb}{0303}
\pdfglyphtounicode{tildedoublecmb}{0360}
\pdfglyphtounicode{tildeoperator}{223C}
\pdfglyphtounicode{tildeoverlaycmb}{0334}
\pdfglyphtounicode{tildeverticalcmb}{033E}
\pdfglyphtounicode{timescircle}{2297}
\pdfglyphtounicode{tipehahebrew}{0596}
\pdfglyphtounicode{tipehalefthebrew}{0596}
\pdfglyphtounicode{tippigurmukhi}{0A70}
\pdfglyphtounicode{titlocyrilliccmb}{0483}
\pdfglyphtounicode{tiwnarmenian}{057F}
\pdfglyphtounicode{tlinebelow}{1E6F}
\pdfglyphtounicode{tmonospace}{FF54}
\pdfglyphtounicode{toarmenian}{0569}
\pdfglyphtounicode{tohiragana}{3068}
\pdfglyphtounicode{tokatakana}{30C8}
\pdfglyphtounicode{tokatakanahalfwidth}{FF84}
\pdfglyphtounicode{tonebarextrahighmod}{02E5}
\pdfglyphtounicode{tonebarextralowmod}{02E9}
\pdfglyphtounicode{tonebarhighmod}{02E6}
\pdfglyphtounicode{tonebarlowmod}{02E8}
\pdfglyphtounicode{tonebarmidmod}{02E7}
\pdfglyphtounicode{tonefive}{01BD}
\pdfglyphtounicode{tonesix}{0185}
\pdfglyphtounicode{tonetwo}{01A8}
\pdfglyphtounicode{tonos}{0384}
\pdfglyphtounicode{tonsquare}{3327}
\pdfglyphtounicode{topatakthai}{0E0F}
\pdfglyphtounicode{tortoiseshellbracketleft}{3014}
\pdfglyphtounicode{tortoiseshellbracketleftsmall}{FE5D}
\pdfglyphtounicode{tortoiseshellbracketleftvertical}{FE39}
\pdfglyphtounicode{tortoiseshellbracketright}{3015}
\pdfglyphtounicode{tortoiseshellbracketrightsmall}{FE5E}
\pdfglyphtounicode{tortoiseshellbracketrightvertical}{FE3A}
\pdfglyphtounicode{totaothai}{0E15}
\pdfglyphtounicode{tpalatalhook}{01AB}
\pdfglyphtounicode{tparen}{24AF}
\pdfglyphtounicode{trademark}{2122}
\pdfglyphtounicode{trademarksans}{2122}
\pdfglyphtounicode{trademarkserif}{2122}
\pdfglyphtounicode{tretroflexhook}{0288}
\pdfglyphtounicode{triagdn}{25BC}
\pdfglyphtounicode{triaglf}{25C4}
\pdfglyphtounicode{triagrt}{25BA}
\pdfglyphtounicode{triagup}{25B2}
\pdfglyphtounicode{triangle}{25B3}
\pdfglyphtounicode{triangledownsld}{25BC}
\pdfglyphtounicode{triangleinv}{25BD}
\pdfglyphtounicode{triangleleft}{25C1}
\pdfglyphtounicode{triangleleftequal}{22B4}
\pdfglyphtounicode{triangleleftsld}{25C0}
\pdfglyphtounicode{triangleright}{25B7}
\pdfglyphtounicode{trianglerightequal}{22B5}
\pdfglyphtounicode{trianglerightsld}{25B6}
\pdfglyphtounicode{trianglesolid}{25B2}
\pdfglyphtounicode{ts}{02A6}
\pdfglyphtounicode{tsadi}{05E6}
\pdfglyphtounicode{tsadidagesh}{FB46}
\pdfglyphtounicode{tsadidageshhebrew}{FB46}
\pdfglyphtounicode{tsadihebrew}{05E6}
\pdfglyphtounicode{tsecyrillic}{0446}
\pdfglyphtounicode{tsere}{05B5}
\pdfglyphtounicode{tsere12}{05B5}
\pdfglyphtounicode{tsere1e}{05B5}
\pdfglyphtounicode{tsere2b}{05B5}
\pdfglyphtounicode{tserehebrew}{05B5}
\pdfglyphtounicode{tserenarrowhebrew}{05B5}
\pdfglyphtounicode{tserequarterhebrew}{05B5}
\pdfglyphtounicode{tserewidehebrew}{05B5}
\pdfglyphtounicode{tshecyrillic}{045B}
\pdfglyphtounicode{tsuperior}{0074}
\pdfglyphtounicode{ttabengali}{099F}
\pdfglyphtounicode{ttadeva}{091F}
\pdfglyphtounicode{ttagujarati}{0A9F}
\pdfglyphtounicode{ttagurmukhi}{0A1F}
\pdfglyphtounicode{tteharabic}{0679}
\pdfglyphtounicode{ttehfinalarabic}{FB67}
\pdfglyphtounicode{ttehinitialarabic}{FB68}
\pdfglyphtounicode{ttehmedialarabic}{FB69}
\pdfglyphtounicode{tthabengali}{09A0}
\pdfglyphtounicode{tthadeva}{0920}
\pdfglyphtounicode{tthagujarati}{0AA0}
\pdfglyphtounicode{tthagurmukhi}{0A20}
\pdfglyphtounicode{tturned}{0287}
\pdfglyphtounicode{tuhiragana}{3064}
\pdfglyphtounicode{tukatakana}{30C4}
\pdfglyphtounicode{tukatakanahalfwidth}{FF82}
\pdfglyphtounicode{turnstileleft}{22A2}
\pdfglyphtounicode{turnstileright}{22A3}
\pdfglyphtounicode{tusmallhiragana}{3063}
\pdfglyphtounicode{tusmallkatakana}{30C3}
\pdfglyphtounicode{tusmallkatakanahalfwidth}{FF6F}
\pdfglyphtounicode{twelvecircle}{246B}
\pdfglyphtounicode{twelveparen}{247F}
\pdfglyphtounicode{twelveperiod}{2493}
\pdfglyphtounicode{twelveroman}{217B}
\pdfglyphtounicode{twentycircle}{2473}
\pdfglyphtounicode{twentyhangzhou}{5344}
\pdfglyphtounicode{twentyparen}{2487}
\pdfglyphtounicode{twentyperiod}{249B}
\pdfglyphtounicode{two}{0032}
\pdfglyphtounicode{twoarabic}{0662}
\pdfglyphtounicode{twobengali}{09E8}
\pdfglyphtounicode{twocircle}{2461}
\pdfglyphtounicode{twocircleinversesansserif}{278B}
\pdfglyphtounicode{twodeva}{0968}
\pdfglyphtounicode{twodotenleader}{2025}
\pdfglyphtounicode{twodotleader}{2025}
\pdfglyphtounicode{twodotleadervertical}{FE30}
\pdfglyphtounicode{twogujarati}{0AE8}
\pdfglyphtounicode{twogurmukhi}{0A68}
\pdfglyphtounicode{twohackarabic}{0662}
\pdfglyphtounicode{twohangzhou}{3022}
\pdfglyphtounicode{twoideographicparen}{3221}
\pdfglyphtounicode{twoinferior}{2082}
\pdfglyphtounicode{twomonospace}{FF12}
\pdfglyphtounicode{twonumeratorbengali}{09F5}
\pdfglyphtounicode{twooldstyle}{0032}
\pdfglyphtounicode{twoparen}{2475}
\pdfglyphtounicode{twoperiod}{2489}
\pdfglyphtounicode{twopersian}{06F2}
\pdfglyphtounicode{tworoman}{2171}
\pdfglyphtounicode{twostroke}{01BB}
\pdfglyphtounicode{twosuperior}{00B2}
\pdfglyphtounicode{twothai}{0E52}
\pdfglyphtounicode{twothirds}{2154}
\pdfglyphtounicode{u}{0075}
\pdfglyphtounicode{uacute}{00FA}
\pdfglyphtounicode{ubar}{0289}
\pdfglyphtounicode{ubengali}{0989}
\pdfglyphtounicode{ubopomofo}{3128}
\pdfglyphtounicode{ubreve}{016D}
\pdfglyphtounicode{ucaron}{01D4}
\pdfglyphtounicode{ucircle}{24E4}
\pdfglyphtounicode{ucircumflex}{00FB}
\pdfglyphtounicode{ucircumflexbelow}{1E77}
\pdfglyphtounicode{ucyrillic}{0443}
\pdfglyphtounicode{udattadeva}{0951}
\pdfglyphtounicode{udblacute}{0171}
\pdfglyphtounicode{udblgrave}{0215}
\pdfglyphtounicode{udeva}{0909}
\pdfglyphtounicode{udieresis}{00FC}
\pdfglyphtounicode{udieresisacute}{01D8}
\pdfglyphtounicode{udieresisbelow}{1E73}
\pdfglyphtounicode{udieresiscaron}{01DA}
\pdfglyphtounicode{udieresiscyrillic}{04F1}
\pdfglyphtounicode{udieresisgrave}{01DC}
\pdfglyphtounicode{udieresismacron}{01D6}
\pdfglyphtounicode{udotbelow}{1EE5}
\pdfglyphtounicode{ugrave}{00F9}
\pdfglyphtounicode{ugujarati}{0A89}
\pdfglyphtounicode{ugurmukhi}{0A09}
\pdfglyphtounicode{uhiragana}{3046}
\pdfglyphtounicode{uhookabove}{1EE7}
\pdfglyphtounicode{uhorn}{01B0}
\pdfglyphtounicode{uhornacute}{1EE9}
\pdfglyphtounicode{uhorndotbelow}{1EF1}
\pdfglyphtounicode{uhorngrave}{1EEB}
\pdfglyphtounicode{uhornhookabove}{1EED}
\pdfglyphtounicode{uhorntilde}{1EEF}
\pdfglyphtounicode{uhungarumlaut}{0171}
\pdfglyphtounicode{uhungarumlautcyrillic}{04F3}
\pdfglyphtounicode{uinvertedbreve}{0217}
\pdfglyphtounicode{ukatakana}{30A6}
\pdfglyphtounicode{ukatakanahalfwidth}{FF73}
\pdfglyphtounicode{ukcyrillic}{0479}
\pdfglyphtounicode{ukorean}{315C}
\pdfglyphtounicode{umacron}{016B}
\pdfglyphtounicode{umacroncyrillic}{04EF}
\pdfglyphtounicode{umacrondieresis}{1E7B}
\pdfglyphtounicode{umatragurmukhi}{0A41}
\pdfglyphtounicode{umonospace}{FF55}
\pdfglyphtounicode{underscore}{005F}
\pdfglyphtounicode{underscoredbl}{2017}
\pdfglyphtounicode{underscoremonospace}{FF3F}
\pdfglyphtounicode{underscorevertical}{FE33}
\pdfglyphtounicode{underscorewavy}{FE4F}
\pdfglyphtounicode{union}{222A}
\pdfglyphtounicode{uniondbl}{22D3}
\pdfglyphtounicode{unionmulti}{228E}
\pdfglyphtounicode{unionsq}{2294}
\pdfglyphtounicode{universal}{2200}
\pdfglyphtounicode{uogonek}{0173}
\pdfglyphtounicode{uparen}{24B0}
\pdfglyphtounicode{upblock}{2580}
\pdfglyphtounicode{upperdothebrew}{05C4}
\pdfglyphtounicode{uprise}{22CF}
\pdfglyphtounicode{upsilon}{03C5}
\pdfglyphtounicode{upsilondieresis}{03CB}
\pdfglyphtounicode{upsilondieresistonos}{03B0}
\pdfglyphtounicode{upsilonlatin}{028A}
\pdfglyphtounicode{upsilontonos}{03CD}
\pdfglyphtounicode{upslope}{29F8}
\pdfglyphtounicode{uptackbelowcmb}{031D}
\pdfglyphtounicode{uptackmod}{02D4}
\pdfglyphtounicode{uragurmukhi}{0A73}
\pdfglyphtounicode{uring}{016F}
\pdfglyphtounicode{ushortcyrillic}{045E}
\pdfglyphtounicode{usmallhiragana}{3045}
\pdfglyphtounicode{usmallkatakana}{30A5}
\pdfglyphtounicode{usmallkatakanahalfwidth}{FF69}
\pdfglyphtounicode{ustraightcyrillic}{04AF}
\pdfglyphtounicode{ustraightstrokecyrillic}{04B1}
\pdfglyphtounicode{utilde}{0169}
\pdfglyphtounicode{utildeacute}{1E79}
\pdfglyphtounicode{utildebelow}{1E75}
\pdfglyphtounicode{uubengali}{098A}
\pdfglyphtounicode{uudeva}{090A}
\pdfglyphtounicode{uugujarati}{0A8A}
\pdfglyphtounicode{uugurmukhi}{0A0A}
\pdfglyphtounicode{uumatragurmukhi}{0A42}
\pdfglyphtounicode{uuvowelsignbengali}{09C2}
\pdfglyphtounicode{uuvowelsigndeva}{0942}
\pdfglyphtounicode{uuvowelsigngujarati}{0AC2}
\pdfglyphtounicode{uvowelsignbengali}{09C1}
\pdfglyphtounicode{uvowelsigndeva}{0941}
\pdfglyphtounicode{uvowelsigngujarati}{0AC1}
\pdfglyphtounicode{v}{0076}
\pdfglyphtounicode{vadeva}{0935}
\pdfglyphtounicode{vagujarati}{0AB5}
\pdfglyphtounicode{vagurmukhi}{0A35}
\pdfglyphtounicode{vakatakana}{30F7}
\pdfglyphtounicode{vav}{05D5}
\pdfglyphtounicode{vavdagesh}{FB35}
\pdfglyphtounicode{vavdagesh65}{FB35}
\pdfglyphtounicode{vavdageshhebrew}{FB35}
\pdfglyphtounicode{vavhebrew}{05D5}
\pdfglyphtounicode{vavholam}{FB4B}
\pdfglyphtounicode{vavholamhebrew}{FB4B}
\pdfglyphtounicode{vavvavhebrew}{05F0}
\pdfglyphtounicode{vavyodhebrew}{05F1}
\pdfglyphtounicode{vcircle}{24E5}
\pdfglyphtounicode{vdotbelow}{1E7F}
\pdfglyphtounicode{vector}{20D7}
\pdfglyphtounicode{vecyrillic}{0432}
\pdfglyphtounicode{veharabic}{06A4}
\pdfglyphtounicode{vehfinalarabic}{FB6B}
\pdfglyphtounicode{vehinitialarabic}{FB6C}
\pdfglyphtounicode{vehmedialarabic}{FB6D}
\pdfglyphtounicode{vekatakana}{30F9}
\pdfglyphtounicode{venus}{2640}
\pdfglyphtounicode{verticalbar}{007C}
\pdfglyphtounicode{verticallineabovecmb}{030D}
\pdfglyphtounicode{verticallinebelowcmb}{0329}
\pdfglyphtounicode{verticallinelowmod}{02CC}
\pdfglyphtounicode{verticallinemod}{02C8}
\pdfglyphtounicode{vewarmenian}{057E}
\pdfglyphtounicode{vhook}{028B}
\pdfglyphtounicode{vikatakana}{30F8}
\pdfglyphtounicode{viramabengali}{09CD}
\pdfglyphtounicode{viramadeva}{094D}
\pdfglyphtounicode{viramagujarati}{0ACD}
\pdfglyphtounicode{visargabengali}{0983}
\pdfglyphtounicode{visargadeva}{0903}
\pdfglyphtounicode{visargagujarati}{0A83}
\pdfglyphtounicode{visiblespace}{2423}
\pdfglyphtounicode{visualspace}{2423}
\pdfglyphtounicode{vmonospace}{FF56}
\pdfglyphtounicode{voarmenian}{0578}
\pdfglyphtounicode{voicediterationhiragana}{309E}
\pdfglyphtounicode{voicediterationkatakana}{30FE}
\pdfglyphtounicode{voicedmarkkana}{309B}
\pdfglyphtounicode{voicedmarkkanahalfwidth}{FF9E}
\pdfglyphtounicode{vokatakana}{30FA}
\pdfglyphtounicode{vparen}{24B1}
\pdfglyphtounicode{vtilde}{1E7D}
\pdfglyphtounicode{vturned}{028C}
\pdfglyphtounicode{vuhiragana}{3094}
\pdfglyphtounicode{vukatakana}{30F4}
\pdfglyphtounicode{w}{0077}
\pdfglyphtounicode{wacute}{1E83}
\pdfglyphtounicode{waekorean}{3159}
\pdfglyphtounicode{wahiragana}{308F}
\pdfglyphtounicode{wakatakana}{30EF}
\pdfglyphtounicode{wakatakanahalfwidth}{FF9C}
\pdfglyphtounicode{wakorean}{3158}
\pdfglyphtounicode{wasmallhiragana}{308E}
\pdfglyphtounicode{wasmallkatakana}{30EE}
\pdfglyphtounicode{wattosquare}{3357}
\pdfglyphtounicode{wavedash}{301C}
\pdfglyphtounicode{wavyunderscorevertical}{FE34}
\pdfglyphtounicode{wawarabic}{0648}
\pdfglyphtounicode{wawfinalarabic}{FEEE}
\pdfglyphtounicode{wawhamzaabovearabic}{0624}
\pdfglyphtounicode{wawhamzaabovefinalarabic}{FE86}
\pdfglyphtounicode{wbsquare}{33DD}
\pdfglyphtounicode{wcircle}{24E6}
\pdfglyphtounicode{wcircumflex}{0175}
\pdfglyphtounicode{wdieresis}{1E85}
\pdfglyphtounicode{wdotaccent}{1E87}
\pdfglyphtounicode{wdotbelow}{1E89}
\pdfglyphtounicode{wehiragana}{3091}
\pdfglyphtounicode{weierstrass}{2118}
\pdfglyphtounicode{wekatakana}{30F1}
\pdfglyphtounicode{wekorean}{315E}
\pdfglyphtounicode{weokorean}{315D}
\pdfglyphtounicode{wgrave}{1E81}
\pdfglyphtounicode{whitebullet}{25E6}
\pdfglyphtounicode{whitecircle}{25CB}
\pdfglyphtounicode{whitecircleinverse}{25D9}
\pdfglyphtounicode{whitecornerbracketleft}{300E}
\pdfglyphtounicode{whitecornerbracketleftvertical}{FE43}
\pdfglyphtounicode{whitecornerbracketright}{300F}
\pdfglyphtounicode{whitecornerbracketrightvertical}{FE44}
\pdfglyphtounicode{whitediamond}{25C7}
\pdfglyphtounicode{whitediamondcontainingblacksmalldiamond}{25C8}
\pdfglyphtounicode{whitedownpointingsmalltriangle}{25BF}
\pdfglyphtounicode{whitedownpointingtriangle}{25BD}
\pdfglyphtounicode{whiteleftpointingsmalltriangle}{25C3}
\pdfglyphtounicode{whiteleftpointingtriangle}{25C1}
\pdfglyphtounicode{whitelenticularbracketleft}{3016}
\pdfglyphtounicode{whitelenticularbracketright}{3017}
\pdfglyphtounicode{whiterightpointingsmalltriangle}{25B9}
\pdfglyphtounicode{whiterightpointingtriangle}{25B7}
\pdfglyphtounicode{whitesmallsquare}{25AB}
\pdfglyphtounicode{whitesmilingface}{263A}
\pdfglyphtounicode{whitesquare}{25A1}
\pdfglyphtounicode{whitestar}{2606}
\pdfglyphtounicode{whitetelephone}{260F}
\pdfglyphtounicode{whitetortoiseshellbracketleft}{3018}
\pdfglyphtounicode{whitetortoiseshellbracketright}{3019}
\pdfglyphtounicode{whiteuppointingsmalltriangle}{25B5}
\pdfglyphtounicode{whiteuppointingtriangle}{25B3}
\pdfglyphtounicode{wihiragana}{3090}
\pdfglyphtounicode{wikatakana}{30F0}
\pdfglyphtounicode{wikorean}{315F}
\pdfglyphtounicode{wmonospace}{FF57}
\pdfglyphtounicode{wohiragana}{3092}
\pdfglyphtounicode{wokatakana}{30F2}
\pdfglyphtounicode{wokatakanahalfwidth}{FF66}
\pdfglyphtounicode{won}{20A9}
\pdfglyphtounicode{wonmonospace}{FFE6}
\pdfglyphtounicode{wowaenthai}{0E27}
\pdfglyphtounicode{wparen}{24B2}
\pdfglyphtounicode{wreathproduct}{2240}
\pdfglyphtounicode{wring}{1E98}
\pdfglyphtounicode{wsuperior}{02B7}
\pdfglyphtounicode{wturned}{028D}
\pdfglyphtounicode{wynn}{01BF}
\pdfglyphtounicode{x}{0078}
\pdfglyphtounicode{xabovecmb}{033D}
\pdfglyphtounicode{xbopomofo}{3112}
\pdfglyphtounicode{xcircle}{24E7}
\pdfglyphtounicode{xdieresis}{1E8D}
\pdfglyphtounicode{xdotaccent}{1E8B}
\pdfglyphtounicode{xeharmenian}{056D}
\pdfglyphtounicode{xi}{03BE}
\pdfglyphtounicode{xmonospace}{FF58}
\pdfglyphtounicode{xparen}{24B3}
\pdfglyphtounicode{xsuperior}{02E3}
\pdfglyphtounicode{y}{0079}
\pdfglyphtounicode{yaadosquare}{334E}
\pdfglyphtounicode{yabengali}{09AF}
\pdfglyphtounicode{yacute}{00FD}
\pdfglyphtounicode{yadeva}{092F}
\pdfglyphtounicode{yaekorean}{3152}
\pdfglyphtounicode{yagujarati}{0AAF}
\pdfglyphtounicode{yagurmukhi}{0A2F}
\pdfglyphtounicode{yahiragana}{3084}
\pdfglyphtounicode{yakatakana}{30E4}
\pdfglyphtounicode{yakatakanahalfwidth}{FF94}
\pdfglyphtounicode{yakorean}{3151}
\pdfglyphtounicode{yamakkanthai}{0E4E}
\pdfglyphtounicode{yasmallhiragana}{3083}
\pdfglyphtounicode{yasmallkatakana}{30E3}
\pdfglyphtounicode{yasmallkatakanahalfwidth}{FF6C}
\pdfglyphtounicode{yatcyrillic}{0463}
\pdfglyphtounicode{ycircle}{24E8}
\pdfglyphtounicode{ycircumflex}{0177}
\pdfglyphtounicode{ydieresis}{00FF}
\pdfglyphtounicode{ydotaccent}{1E8F}
\pdfglyphtounicode{ydotbelow}{1EF5}
\pdfglyphtounicode{yeharabic}{064A}
\pdfglyphtounicode{yehbarreearabic}{06D2}
\pdfglyphtounicode{yehbarreefinalarabic}{FBAF}
\pdfglyphtounicode{yehfinalarabic}{FEF2}
\pdfglyphtounicode{yehhamzaabovearabic}{0626}
\pdfglyphtounicode{yehhamzaabovefinalarabic}{FE8A}
\pdfglyphtounicode{yehhamzaaboveinitialarabic}{FE8B}
\pdfglyphtounicode{yehhamzaabovemedialarabic}{FE8C}
\pdfglyphtounicode{yehinitialarabic}{FEF3}
\pdfglyphtounicode{yehmedialarabic}{FEF4}
\pdfglyphtounicode{yehmeeminitialarabic}{FCDD}
\pdfglyphtounicode{yehmeemisolatedarabic}{FC58}
\pdfglyphtounicode{yehnoonfinalarabic}{FC94}
\pdfglyphtounicode{yehthreedotsbelowarabic}{06D1}
\pdfglyphtounicode{yekorean}{3156}
\pdfglyphtounicode{yen}{00A5}
\pdfglyphtounicode{yenmonospace}{FFE5}
\pdfglyphtounicode{yeokorean}{3155}
\pdfglyphtounicode{yeorinhieuhkorean}{3186}
\pdfglyphtounicode{yerahbenyomohebrew}{05AA}
\pdfglyphtounicode{yerahbenyomolefthebrew}{05AA}
\pdfglyphtounicode{yericyrillic}{044B}
\pdfglyphtounicode{yerudieresiscyrillic}{04F9}
\pdfglyphtounicode{yesieungkorean}{3181}
\pdfglyphtounicode{yesieungpansioskorean}{3183}
\pdfglyphtounicode{yesieungsioskorean}{3182}
\pdfglyphtounicode{yetivhebrew}{059A}
\pdfglyphtounicode{ygrave}{1EF3}
\pdfglyphtounicode{yhook}{01B4}
\pdfglyphtounicode{yhookabove}{1EF7}
\pdfglyphtounicode{yiarmenian}{0575}
\pdfglyphtounicode{yicyrillic}{0457}
\pdfglyphtounicode{yikorean}{3162}
\pdfglyphtounicode{yinyang}{262F}
\pdfglyphtounicode{yiwnarmenian}{0582}
\pdfglyphtounicode{ymonospace}{FF59}
\pdfglyphtounicode{yod}{05D9}
\pdfglyphtounicode{yoddagesh}{FB39}
\pdfglyphtounicode{yoddageshhebrew}{FB39}
\pdfglyphtounicode{yodhebrew}{05D9}
\pdfglyphtounicode{yodyodhebrew}{05F2}
\pdfglyphtounicode{yodyodpatahhebrew}{FB1F}
\pdfglyphtounicode{yohiragana}{3088}
\pdfglyphtounicode{yoikorean}{3189}
\pdfglyphtounicode{yokatakana}{30E8}
\pdfglyphtounicode{yokatakanahalfwidth}{FF96}
\pdfglyphtounicode{yokorean}{315B}
\pdfglyphtounicode{yosmallhiragana}{3087}
\pdfglyphtounicode{yosmallkatakana}{30E7}
\pdfglyphtounicode{yosmallkatakanahalfwidth}{FF6E}
\pdfglyphtounicode{yotgreek}{03F3}
\pdfglyphtounicode{yoyaekorean}{3188}
\pdfglyphtounicode{yoyakorean}{3187}
\pdfglyphtounicode{yoyakthai}{0E22}
\pdfglyphtounicode{yoyingthai}{0E0D}
\pdfglyphtounicode{yparen}{24B4}
\pdfglyphtounicode{ypogegrammeni}{037A}
\pdfglyphtounicode{ypogegrammenigreekcmb}{0345}
\pdfglyphtounicode{yr}{01A6}
\pdfglyphtounicode{yring}{1E99}
\pdfglyphtounicode{ysuperior}{02B8}
\pdfglyphtounicode{ytilde}{1EF9}
\pdfglyphtounicode{yturned}{028E}
\pdfglyphtounicode{yuhiragana}{3086}
\pdfglyphtounicode{yuikorean}{318C}
\pdfglyphtounicode{yukatakana}{30E6}
\pdfglyphtounicode{yukatakanahalfwidth}{FF95}
\pdfglyphtounicode{yukorean}{3160}
\pdfglyphtounicode{yusbigcyrillic}{046B}
\pdfglyphtounicode{yusbigiotifiedcyrillic}{046D}
\pdfglyphtounicode{yuslittlecyrillic}{0467}
\pdfglyphtounicode{yuslittleiotifiedcyrillic}{0469}
\pdfglyphtounicode{yusmallhiragana}{3085}
\pdfglyphtounicode{yusmallkatakana}{30E5}
\pdfglyphtounicode{yusmallkatakanahalfwidth}{FF6D}
\pdfglyphtounicode{yuyekorean}{318B}
\pdfglyphtounicode{yuyeokorean}{318A}
\pdfglyphtounicode{yyabengali}{09DF}
\pdfglyphtounicode{yyadeva}{095F}
\pdfglyphtounicode{z}{007A}
\pdfglyphtounicode{zaarmenian}{0566}
\pdfglyphtounicode{zacute}{017A}
\pdfglyphtounicode{zadeva}{095B}
\pdfglyphtounicode{zagurmukhi}{0A5B}
\pdfglyphtounicode{zaharabic}{0638}
\pdfglyphtounicode{zahfinalarabic}{FEC6}
\pdfglyphtounicode{zahinitialarabic}{FEC7}
\pdfglyphtounicode{zahiragana}{3056}
\pdfglyphtounicode{zahmedialarabic}{FEC8}
\pdfglyphtounicode{zainarabic}{0632}
\pdfglyphtounicode{zainfinalarabic}{FEB0}
\pdfglyphtounicode{zakatakana}{30B6}
\pdfglyphtounicode{zaqefgadolhebrew}{0595}
\pdfglyphtounicode{zaqefqatanhebrew}{0594}
\pdfglyphtounicode{zarqahebrew}{0598}
\pdfglyphtounicode{zayin}{05D6}
\pdfglyphtounicode{zayindagesh}{FB36}
\pdfglyphtounicode{zayindageshhebrew}{FB36}
\pdfglyphtounicode{zayinhebrew}{05D6}
\pdfglyphtounicode{zbopomofo}{3117}
\pdfglyphtounicode{zcaron}{017E}
\pdfglyphtounicode{zcircle}{24E9}
\pdfglyphtounicode{zcircumflex}{1E91}
\pdfglyphtounicode{zcurl}{0291}
\pdfglyphtounicode{zdot}{017C}
\pdfglyphtounicode{zdotaccent}{017C}
\pdfglyphtounicode{zdotbelow}{1E93}
\pdfglyphtounicode{zecyrillic}{0437}
\pdfglyphtounicode{zedescendercyrillic}{0499}
\pdfglyphtounicode{zedieresiscyrillic}{04DF}
\pdfglyphtounicode{zehiragana}{305C}
\pdfglyphtounicode{zekatakana}{30BC}
\pdfglyphtounicode{zero}{0030}
\pdfglyphtounicode{zeroarabic}{0660}
\pdfglyphtounicode{zerobengali}{09E6}
\pdfglyphtounicode{zerodeva}{0966}
\pdfglyphtounicode{zerogujarati}{0AE6}
\pdfglyphtounicode{zerogurmukhi}{0A66}
\pdfglyphtounicode{zerohackarabic}{0660}
\pdfglyphtounicode{zeroinferior}{2080}
\pdfglyphtounicode{zeromonospace}{FF10}
\pdfglyphtounicode{zerooldstyle}{0030}
\pdfglyphtounicode{zeropersian}{06F0}
\pdfglyphtounicode{zerosuperior}{2070}
\pdfglyphtounicode{zerothai}{0E50}
\pdfglyphtounicode{zerowidthjoiner}{FEFF}
\pdfglyphtounicode{zerowidthnonjoiner}{200C}
\pdfglyphtounicode{zerowidthspace}{200B}
\pdfglyphtounicode{zeta}{03B6}
\pdfglyphtounicode{zhbopomofo}{3113}
\pdfglyphtounicode{zhearmenian}{056A}
\pdfglyphtounicode{zhebrevecyrillic}{04C2}
\pdfglyphtounicode{zhecyrillic}{0436}
\pdfglyphtounicode{zhedescendercyrillic}{0497}
\pdfglyphtounicode{zhedieresiscyrillic}{04DD}
\pdfglyphtounicode{zihiragana}{3058}
\pdfglyphtounicode{zikatakana}{30B8}
\pdfglyphtounicode{zinorhebrew}{05AE}
\pdfglyphtounicode{zlinebelow}{1E95}
\pdfglyphtounicode{zmonospace}{FF5A}
\pdfglyphtounicode{zohiragana}{305E}
\pdfglyphtounicode{zokatakana}{30BE}
\pdfglyphtounicode{zparen}{24B5}
\pdfglyphtounicode{zretroflexhook}{0290}
\pdfglyphtounicode{zstroke}{01B6}
\pdfglyphtounicode{zuhiragana}{305A}
\pdfglyphtounicode{zukatakana}{30BA}
 \pdfgentounicode=1
\usepackage[english]{babel}
\usepackage{newtxtext}
\usepackage[protrusion=true,expansion=true,tracking=true]{microtype}
\usepackage[a4paper,inner=2.72cm,outer=2.34cm,top=2.42cm,bottom=2.68cm,
  headheight=31pt,headsep=18pt,footskip=30pt]{geometry}
\usepackage{amsmath,amsthm,amscd,mathtools,mathrsfs,bm}
\usepackage{newtxmath}
% Las jerarquías de títulos, el índice y las tablas emplean escalas
% intermedias. Se mantiene la familia RSFS y se habilita su escalado continuo,
% de modo que la notación matemática conserve el cuerpo tipográfico exacto
% sin sustituciones silenciosas de tamaño.
\DeclareFontFamily{U}{rsfs}{\skewchar\font127}
\DeclareFontShape{U}{rsfs}{m}{n}{%
  <-6> s*[1.0] rsfs5
  <6-8> s*[1.0] rsfs7
  <8-> s*[1.0] rsfs10
}{}
\usepackage{array,booktabs,longtable,tabularx,multirow}
\usepackage{enumitem,graphicx,float,pdfpages,pdflscape}
\usepackage{subcaption}
\usepackage[table]{xcolor}
\pdfinclusioncopyfonts=1
\usepackage{tikz}
\usepackage{tikz-cd}
\usetikzlibrary{arrows.meta,positioning,fit,calc,matrix,shapes.geometric,
  decorations.pathreplacing,decorations.markings,backgrounds,patterns}
\usepackage[most]{tcolorbox}
\usepackage{setspace,csquotes,xurl,fancyhdr,fancyvrb,listings,tocloft,titlesec,needspace,etoolbox}
\usepackage{xstring}
\usepackage{hyperref}
% Los marcadores PDF son parte de la jerarquía científica de la obra.  Esta
% tabla conserva en texto Unicode la notación que aparece en los títulos, en
% lugar de dejar que hyperref suprima silenciosamente flechas y constantes.
\pdfstringdefDisableCommands{%
  \def\({}\def\){}\def\[{}\def\]{}%
  \def\ensuremath#1{#1}%
  \def\mathrm#1{#1}\def\mathbf#1{#1}\def\mathsf#1{#1}%
  \def\mathbb#1{#1}\def\mathcal#1{#1}\def\mathfrak#1{#1}%
  \def\operatorname#1{#1}\def\text#1{#1}\def\bm#1{#1}%
  \def\allowbreak{}%
  \def\frac#1#2{#1/#2}\def\sqrt#1{root(#1)}%
  \def\to{→}\def\longrightarrow{→}\def\mapsto{↦}%
  \def\leftrightarrow{↔}\def\Rightarrow{⇒}\def\Longrightarrow{⇒}%
  \def\pi{π}\def\varphi{φ}\def\phi{φ}\def\alpha{α}%
  \def\varepsilon{ε}\def\epsilon{ε}\def\mu{μ}\def\lambda{λ}%
  \def\Omega{Ω}\def\zeta{ζ}\def\partial{∂}%
  \def\ast{*}\def\cdot{·}\def\natural{natural}\def\log{log}%
}
\usepackage[nameinlink,noabbrev]{cleveref}
\crefname{appendix}{appendix}{appendices}
\Crefname{appendix}{Appendix}{Appendices}
% Referencias descriptivas a resultados del tratado que se transducen en esta
% monografía bajo una numeración propia.  La sustitución evita remitir al lector
% a la numeración editorial de otra obra sin duplicar los resultados.
\makeatletter
\newcommand{\HMTDeclareDescriptor}[2]{%
  \expandafter\gdef\csname hmt@xref@#1\endcsname{#2}}
\newcommand{\HMTIfDescriptorTF}[3]{%
  \ifcsname hmt@xref@#1\endcsname #2\else #3\fi}

\HMTDeclareDescriptor{app:catalogo-semillas-app-rev9}{chapter devoted to the APP seed domain and the canonical emission catalogue}
\HMTDeclareDescriptor{chap:83-08}{chapter devoted to the genealogical decomposition of the five oriented regions of \(\pi\)}
\HMTDeclareDescriptor{chap:alpha}{chapter devoted to the dodecaphasic closure of \(\alpha\)}
\HMTDeclareDescriptor{chap:banda-particula-virtual}{chapter devoted to the particle band, anti-section and virtuality}
\HMTDeclareDescriptor{chap:bandera-excepcional-acumulativa}{chapter devoted to dodecaphasic incidence and the Witt--Leech--Monster corridor}
\HMTDeclareDescriptor{chap:biografia-pi-autonoma}{chapter devoted to the construction of the closure character and the structural biography of \(\pi\)}
\HMTDeclareDescriptor{chap:cinco-ciclos-pi}{chapter devoted to the minimality of the closures of the five regions of \(\pi\)}
\HMTDeclareDescriptor{chap:funtor-dimensional-hmt}{chapter devoted to magnitude lines and the dimensional realization of the kernel}
\HMTDeclareDescriptor{chap:integral-genealogica}{chapter devoted to the genealogical path algebra and its integral representation}
\HMTDeclareDescriptor{chap:numero-medicion}{chapter devoted to cylindrical systems, survival and the construction of number}
\HMTDeclareDescriptor{chap:orden-determinantal-accion}{chapter devoted to the determinantal order of the action line}
\HMTDeclareDescriptor{chap:pantallas-dualidad}{chapter devoted to subspaces, quotients, rank reductions and duality}
\HMTDeclareDescriptor{chap:semillas-emisor-54}{chapter devoted to the index-\(54\) emitter and its oriented seeds}
\HMTDeclareDescriptor{chap:sistema-inverso-continuidad-genealogica}{chapter devoted to the inverse system and genealogical continuity}
\HMTDeclareDescriptor{chap:tpk-definicion}{chapter devoted to the transport system and phase evolution}
\HMTDeclareDescriptor{def:banda-dodecafasica-ruta}{definition of the dodecaphasic route band}
\HMTDeclareDescriptor{def:odometro-9-12}{definition of the nonadic--dodecaphasic odometer}
\HMTDeclareDescriptor{eq:G-interno}{internal constitutive identity of gravitation}
\HMTDeclareDescriptor{eq:consistencia-proyectiva-medidas-u024}{projective consistency identity of cylindrical measures}
\HMTDeclareDescriptor{eq:continuo-lector-arquimediano}{coverage identity of the Archimedean reader}
\HMTDeclareDescriptor{eq:descomposiciones-cinco-regiones-pi}{decomposition of the five regions of \(\pi\)}
\HMTDeclareDescriptor{eq:frontal-doble-proyeccion}{frontal diagram of the double projection}
\HMTDeclareDescriptor{eq:modular-finite-first-law}{first finite modular identity with defect term}
\HMTDeclareDescriptor{eq:modular-inverse-channels}{joint reconstruction identity of the area channels}
\HMTDeclareDescriptor{eq:realizacion-circular-ct108}{circular realization identity of the \(108\)-phase calendar}
\HMTDeclareDescriptor{espiral:chap:realizacion-em}{chapter devoted to the restricted electromagnetic realization}
\HMTDeclareDescriptor{espiral:eq:cociclos-radiales}{identity of the radial cocycles of the tritic lift}
\HMTDeclareDescriptor{hap:def:entropia-bicapa-rev6}{definition of oriented bilayer entropy}
\HMTDeclareDescriptor{hmtctx:p12:chap:orbitas-elevacion-r36}{chapter devoted to orbits and lifting to the \(R_{36}\) gauge}
\HMTDeclareDescriptor{mdg:chap:torres}{chapter devoted to the global hierarchy of dimensional towers}
\HMTDeclareDescriptor{mdg:sec:cierre-electron-two-way-rev11}{section devoted to bidirectional electronic closure}
\HMTDeclareDescriptor{mdg:sec:ley-general-masa-accion-autonoma}{section devoted to the universal action operator}
\HMTDeclareDescriptor{mdg:sec:selector-gravitatorio-rev11}{section devoted to the natural gravitational selector}
\HMTDeclareDescriptor{p12-83-c19-s1:hmtch:eq-celula-trit-completa}{identity of the complete tritic cell}
\HMTDeclareDescriptor{prop:cierre-antiseccion-ct108}{proposition on closure of the forms of the \(108\)-phase calendar}
\HMTDeclareDescriptor{prop:proyeccion-fase-longitud}{phase--length projection proposition}
\HMTDeclareDescriptor{rm:ch:cosmo}{chapter devoted to the cosmological scales of the \(108\)-phase calendar}
\HMTDeclareDescriptor{rm:ch:metrology}{chapter devoted to the quantum--electrical constants, the thermal scale and the Planck system}
\HMTDeclareDescriptor{sec:tres-generadores-concretos-trit}{section devoted to the three concrete generators of the tritic regimes}
\HMTDeclareDescriptor{sec:w12-w24-elevacion}{section devoted to the dodecaphasic residual design and the lifting of hexads}
\HMTDeclareDescriptor{thm:censo-atlas-81}{classification theorem for the families and generations of the nonadic atlas}
\HMTDeclareDescriptor{thm:corredor-hmt-moonshine}{theorem on the corridor from Leech to the Moonshine algebra}
\HMTDeclareDescriptor{thm:dos-realizaciones-fuente-cuadratica-rev3}{theorem on the two realizations of an integral quadratic presentation}
\HMTDeclareDescriptor{thm:generacion-monodromica-conjunta-rev7}{joint filtration theorem at finite depth}
\HMTDeclareDescriptor{thm:homogeneidad-pch}{homogeneity theorem for the Palatini--Cartan--Holst action}
\HMTDeclareDescriptor{thm:principio-conjugacion-masa-banda}{mass and band character conjugation theorem}
\HMTDeclareDescriptor{thm:tpk-cociente-predictivo-36}{theorem on the minimal predictive-memory quotient}

% Copias semánticas, no simples alias: \cref y \Cref son órdenes robustas
% gestionadas por xparse y un \let conservaría el despachador mutable,
% provocando recursión después de \RenewDocumentCommand.
\NewCommandCopy{\HMTOriginalCref}{\cref}
\NewCommandCopy{\HMTOriginalCrefCapital}{\Cref}
\NewCommandCopy{\HMTOriginalEqref}{\eqref}
% Los capítulos transducidos conservan centenares de referencias semánticas
% cuya numeración pertenecía al tratado integral.  Una referencia que ya ha
% recibido etiqueta local conserva su numeración propia; una referencia que
% sólo existe en el volumen rector se expresa por su clase matemática.  De
% este modo el texto sigue siendo legible y autosuficiente, sin signos ``??''
% ni remisiones ficticias a una numeración ausente.
\newcommand{\HMTFallbackReference}[1]{%
  \IfSubStr{#1}{eq:}{the corresponding identity}{%
  \IfSubStr{#1}{thm:}{the corresponding theorem}{%
  \IfSubStr{#1}{theorem:}{the corresponding theorem}{%
  \IfSubStr{#1}{prop:}{the corresponding proposition}{%
  \IfSubStr{#1}{lem:}{the corresponding lemma}{%
  \IfSubStr{#1}{cor:}{the corresponding corollary}{%
  \IfSubStr{#1}{def:}{the corresponding definition}{%
  \IfSubStr{#1}{fig:}{the corresponding figure}{%
  \IfSubStr{#1}{tab:}{the corresponding table}{%
  \IfSubStr{#1}{chap:}{the corresponding chapter}{%
  \IfSubStr{#1}{ch:}{the corresponding chapter}{%
  \IfSubStr{#1}{sec:}{the corresponding section}{%
  \IfSubStr{#1}{crit:}{the corresponding criterion}{%
  the corresponding result}}}}}}}}}}}}}}
\newcommand{\HMTFallbackReferenceCapital}[1]{%
  \IfSubStr{#1}{eq:}{The corresponding identity}{%
  \IfSubStr{#1}{thm:}{The corresponding theorem}{%
  \IfSubStr{#1}{theorem:}{The corresponding theorem}{%
  \IfSubStr{#1}{prop:}{The corresponding proposition}{%
  \IfSubStr{#1}{lem:}{The corresponding lemma}{%
  \IfSubStr{#1}{cor:}{The corresponding corollary}{%
  \IfSubStr{#1}{def:}{The corresponding definition}{%
  \IfSubStr{#1}{fig:}{The corresponding figure}{%
  \IfSubStr{#1}{tab:}{The corresponding table}{%
  \IfSubStr{#1}{chap:}{The corresponding chapter}{%
  \IfSubStr{#1}{ch:}{The corresponding chapter}{%
  \IfSubStr{#1}{sec:}{The corresponding section}{%
  \IfSubStr{#1}{crit:}{The corresponding criterion}{%
  The corresponding result}}}}}}}}}}}}}}
\newcommand{\HMTIfLocalLabelTF}[3]{%
  \ifcsname r@#1\endcsname #2\else #3\fi}
\RenewDocumentCommand{\cref}{s m}{%
  \HMTIfDescriptorTF{#2}{\csname hmt@xref@#2\endcsname}{%
    \HMTIfLocalLabelTF{#2}{%
      \IfBooleanTF{#1}{\HMTOriginalCref*{#2}}{\HMTOriginalCref{#2}}%
    }{\HMTFallbackReference{#2}}}}
\RenewDocumentCommand{\Cref}{s m}{%
  \HMTIfDescriptorTF{#2}{\csname hmt@xref@#2\endcsname}{%
    \HMTIfLocalLabelTF{#2}{%
      \IfBooleanTF{#1}{\HMTOriginalCrefCapital*{#2}}{\HMTOriginalCrefCapital{#2}}%
    }{\HMTFallbackReferenceCapital{#2}}}}
\RenewDocumentCommand{\eqref}{m}{%
  \HMTIfDescriptorTF{#1}{\textup{(\csname hmt@xref@#1\endcsname)}}{%
    \HMTIfLocalLabelTF{#1}{\HMTOriginalEqref{#1}}{\textup{(corresponding identity)}}}}
\makeatother
 \usepackage[noautomatic,quiet]{imakeidx}
\makeindex
\allowdisplaybreaks

\definecolor{hmtblue}{RGB}{24,57,89}
\definecolor{hmtgold}{RGB}{145,101,26}
\definecolor{hmtgreen}{RGB}{29,104,78}
\definecolor{hmtred}{RGB}{151,54,45}
\definecolor{hmtviolet}{RGB}{92,64,125}
\definecolor{hmtgray}{RGB}{87,95,102}
\definecolor{hmtpaper}{RGB}{248,247,243}
\definecolor{hmtlightblue}{RGB}{232,240,247}
\definecolor{hmtlightgold}{RGB}{249,242,225}
\definecolor{hmtlightgreen}{RGB}{232,244,238}
\definecolor{hmtlightred}{RGB}{250,235,233}
\definecolor{hmtlightviolet}{RGB}{241,236,248}

% Paleta pública común para los cuadros, el atlas y los inventarios
% integrales de Mecánica Dimensional. Los nombres en versales son alias
% tipográficos; no introducen una segunda identidad visual.
\colorlet{HMTNavy}{hmtblue}
\colorlet{HMTBlue}{hmtblue}
\colorlet{HMTGold}{hmtgold}
\colorlet{HMTLight}{hmtlightblue}
\colorlet{HMTGray}{hmtgray}
\colorlet{HMTLine}{hmtlightblue}
\providecommand{\HMTtablehead}[1]{\color{white}\sffamily\bfseries #1}
\providecommand{\RaggedRight}{\raggedright}

\newcommand{\HMT}{\ensuremath{\mathrm{HMT}}}
\providecommand{\gammaH}{\gamma_{\HMT}}
\providecommand{\ii}{\mathrm i}
\providecommand{\C}{\mathbb C}
\providecommand{\Z}{\mathbb Z}
\providecommand{\U}{\mathcal U}
\providecommand{\arsinh}{\operatorname{arsinh}}
\providecommand{\arcosh}{\operatorname{arcosh}}
\providecommand{\HHM}{K_{\mathrm{HHM}}}
\providecommand{\Dmod}{\Delta_{4}^{\mathrm{mod}}}
\providecommand{\Val}{\operatorname{Val}}
\providecommand{\Ext}{\operatorname{Ext}}
\providecommand{\ev}{\operatorname{ev}}
\providecommand{\End}{\operatorname{End}}
\providecommand{\diam}{\operatorname{diam}}
\providecommand{\im}{\operatorname{im}}
\providecommand{\supp}{\operatorname{supp}}
\providecommand{\tr}{\operatorname{tr}}
\providecommand{\Gnine}{\mathcal G_9}
\providecommand{\NTHW}{\mathfrak N_{\mathrm{HW}}}
\providecommand{\readerquestion}[1]{%
  \begin{quote}\small\itshape\color{hmtblue}#1\end{quote}}
\newcommand{\TRIT}{\ensuremath{\widehat{\mathbb T}}}
\newcommand{\AAPP}{\mathcal A_9}
\newcommand{\AAPPlift}{\mathcal A_9^{\#}}
\newcommand{\GraphAPP}{\Gamma_{\mathrm{APP},9}}
\newcommand{\MonNine}{\mathscr M_9}
\newcommand{\XTPK}{\mathcal X_{\mathrm{TPK}}}
\newcommand{\CellRQ}{\mathcal C_{9}^{\mathrm{rq}}}
\newcommand{\WordMonoid}{\mathfrak W_{9}^{\times}}
\newcommand{\Cdoce}{C_{12}}
\newcommand{\Kseal}{K}
\newcommand{\Usgn}{U_{12}^{\mathrm{sgn}}}
\newcommand{\Etrans}{\mathcal E_{3,4}}
\newcommand{\Rpi}{\mathscr R_{\pi}}
\newcommand{\Reu}{\mathscr R_{e}}
\newcommand{\Rphi}{\mathscr R_{\varphi}}
\newcommand{\RR}{\mathbb R}
\newcommand{\ZZ}{\mathbb Z}
\newcommand{\QQ}{\mathbb Q}
\newcommand{\FF}{\mathbb F}
\newcommand{\CC}{\mathbb C}
\newcommand{\dd}{\,\mathrm d}
\newcommand{\ee}{\mathrm e}
\newcommand{\Aut}{\operatorname{Aut}}
\newcommand{\Hol}{\operatorname{Hol}}
\newcommand{\dr}{\operatorname{dr}_9}
\newcommand{\Surv}{\operatorname{Surv}}
\newcommand{\val}{\operatorname{val}}
\newcommand{\Dnine}{D_9}
\newcommand{\Wnine}{\mathscr W_9}
\newcommand{\Krad}{\mathcal K_{\#}^{+}}
\newcommand{\Hplus}{\mathcal H_{+}}
\newcommand{\Dom}{\operatorname{Dom}}
\newcommand{\Tr}{\operatorname{Tr}}
\newcommand{\Spec}{\operatorname{Spec}}
\providecommand{\Sha}{\operatorname{Sha}}
\newcommand{\Irr}{\mathfrak I}
\newcommand{\NN}{\mathbb N}
\newcommand{\N}{\mathbb N}
\newcommand{\R}{\mathbb R}
\providecommand{\Adeg}{A_{\rm deg}}
\providecommand{\Arad}{A_{\rm rad}}
\providecommand{\Cdeg}{C^{*}_{\rm deg}}
\providecommand{\Crad}{C^{*}_{\rm rad}}
\providecommand{\etaRetrad}{\eta_{\rm ret}^{(\rm rad)}}
\providecommand{\Gtr}{\ensuremath{\Gamma_{\rm tr}}}
\providecommand{\HTr}{\ensuremath{\mathsf{Tr}_{q}}}
\providecommand{\HAdd}{\ensuremath{\mathsf{Add}}}
\providecommand{\HDef}{\ensuremath{\mathsf{Def}}}
\providecommand{\HRes}{\ensuremath{\mathsf{Res}}}
\providecommand{\HLoc}{\ensuremath{\mathsf{Loc}}}

\theoremstyle{plain}
\newtheorem{theorem}{Theorem}[chapter]
\newtheorem{lemma}[theorem]{Lemma}
\newtheorem{proposition}[theorem]{Proposition}
\newtheorem{corollary}[theorem]{Corollary}
\theoremstyle{definition}
\newtheorem{definition}[theorem]{Definition}
\newtheorem{construction}[theorem]{Construction}
\newtheorem{algorithm}[theorem]{Algorithm}
\newtheorem{ablation}[theorem]{Structural Necessity Proposition}
\newtheorem{criterion}[theorem]{Test criterion}
\theoremstyle{remark}
\newtheorem{remark}[theorem]{Remark}

\tcbset{enhanced,breakable,boxrule=.45pt,arc=0mm,left=2.5mm,right=2.5mm,top=1.8mm,bottom=1.8mm}
\newtcolorbox{exactbox}[1][]{colback=white,colframe=black,title={Exact result},fonttitle=\bfseries,#1}

\newcolumntype{Y}{>{\raggedright\arraybackslash}X}
\newcolumntype{R}{>{\raggedleft\arraybackslash}X}
\newcolumntype{L}[1]{>{\raggedright\arraybackslash}p{#1}}
\newcolumntype{C}{>{\centering\arraybackslash}p{1.65cm}}

\setlength{\emergencystretch}{3em}
 % Emisión de afirmaciones vivas durante la compilación de REV.9.
% Una marca sólo cuenta si la rama TeX que la contiene se ejecuta.
\newwrite\HMTClaimsStream
\immediate\openout\HMTClaimsStream=\jobname.hmtclaims
\newcommand{\HMTClaim}[1]{%
  \immediate\write\HMTClaimsStream{HMT-CLAIM:#1}%
}
\AtEndDocument{\immediate\closeout\HMTClaimsStream}
 % Compatibilidad tipada de las tres aportaciones incorporadas en la edición
% sucesora de 2026-08-12. Este archivo no redefine objetos de REV.3: aporta
% únicamente nombres, entornos y operadores que los cuadernos autónomos
% necesitan para residir en el grafo único del tratado.

\RequirePackage{xspace}

\definecolor{hmtteal}{RGB}{10,107,112}

\providecommand{\Rhyper}{\mathcal R_{\mathrm{hyp}}}
\providecommand{\OmegaInf}{\Omega_{\infty}}
\providecommand{\Clop}{\operatorname{Clop}}
\providecommand{\Closed}{\operatorname{Closed}}
\providecommand{\Proj}{\operatorname{Proj}}
\providecommand{\id}{\operatorname{id}}
\providecommand{\coker}{\operatorname{coker}}
\providecommand{\ord}{\operatorname{ord}}
\providecommand{\HALT}{\mathsf{HALT}}
\providecommand{\Prob}{\mathcal P}
\providecommand{\status}[1]{\emph{#1}}
\providecommand{\citep}{\cite}
\providecommand{\QMS}{\ensuremath{\mathrm{QMS}}\xspace}
\providecommand{\QLS}{\ensuremath{\mathrm{QLS}}\xspace}
\providecommand{\QPM}{\ensuremath{\mathrm{QPM}}\xspace}
\providecommand{\UCP}{\ensuremath{\mathrm{UCP}}\xspace}
\providecommand{\MUB}{\ensuremath{\mathrm{MUB}}\xspace}
\providecommand{\PVM}{\ensuremath{\mathrm{PVM}}\xspace}
\providecommand{\POVM}{\ensuremath{\mathrm{POVM}}\xspace}
\providecommand{\Mat}{\operatorname{Mat}}
\providecommand{\Her}{\operatorname{Her}}
\providecommand{\mconv}{\operatorname{mconv}}
\providecommand{\Span}{\operatorname{span}}
\providecommand{\Ad}{\operatorname{Ad}}
\providecommand{\diag}{\operatorname{diag}}

\theoremstyle{definition}
\newtheorem{example}[theorem]{Example}

\newtcolorbox{resultbox}[1][]{%
  colback=hmtlightgreen,colframe=hmtgreen,
  title={Result},fonttitle=\bfseries,#1}
\newtcolorbox{controlbox}[1][]{%
  colback=white,colframe=hmtgold!75!black,
  title={Structural obstruction},fonttitle=\bfseries,#1}
\newtcolorbox{sourcebox}[1][]{%
  colback=white,colframe=hmtblue!55,
  title={Premises and derivation},fonttitle=\bfseries,#1}
\newtcolorbox{intuitionbox}[1][]{%
  colback=hmtlightgold,colframe=hmtgold,
  title={Scientific reading},fonttitle=\bfseries,#1}
\newtcolorbox{workbox}[1][]{%
  colback=white,colframe=hmtgray,
  title={Mathematical extensions},fonttitle=\bfseries,#1}
\newtcolorbox{importbox}[1][]{%
  colback=white,colframe=hmtblue!55,
  title={Antecedent theorem},fonttitle=\bfseries,#1}
\newtcolorbox{bridgebox}[1][]{%
  colback=hmtlightgold,colframe=hmtgold,
  title={Result of the matrix bridge},fonttitle=\bfseries,#1}
 % Estado material V3 comentario-only: conserva los 219 resultados de V5 y
% registra las dos correcciones editoriales autorizadas del puente matricial.
% Los sidecars 218/219 previos quedan como testigos inactivos e inmutables.
% HMT--MD REV.4: sidecar comentario-only de continuidad material V3/219.
% Encadena V2, conserva 219 IDs y no publica teoremas, títulos ni contenido paginado.
% Las correcciones registrales V3 no alteran enunciados ni fuerzas probatorias.
% HMT_RESULT_BEGIN:R3-001:residence
% {"material_control":"REV3_BASE_GRANULAR+PASS_NO_PERDIDA_REV3","material_state":"CONSERVADO_POR_BASE_GRANULAR_REV3","residences":[{"path":"/Users/ruben/Documents/HMT2/CONTINUIDAD_EDITORIAL/INTEGRACION_REV3_VON_NEUMANN_MATRICIAL_REALIZACION_2026-08-12/REV3_BASE_GRANULAR.json","sha256":"d89983feac6e064f0ce77d6506d60b94dd1eaf00ed174352d3e32bfc5a9b88d4"}],"result_id":"R3-001","statement_sha256":"ec22d0aed84b1a5ccd6d8675f494a2d54bb746ae05ca67325087c42e7697061d"}
% HMT_RESULT_END:R3-001:residence
% HMT_RESULT_BEGIN:R3-002:residence
% {"material_control":"REV3_BASE_GRANULAR+PASS_NO_PERDIDA_REV3","material_state":"CONSERVADO_POR_BASE_GRANULAR_REV3","residences":[{"path":"/Users/ruben/Documents/HMT2/CONTINUIDAD_EDITORIAL/INTEGRACION_REV3_VON_NEUMANN_MATRICIAL_REALIZACION_2026-08-12/REV3_BASE_GRANULAR.json","sha256":"d89983feac6e064f0ce77d6506d60b94dd1eaf00ed174352d3e32bfc5a9b88d4"}],"result_id":"R3-002","statement_sha256":"94b5e013944e5d31216adcca582bf4469c088ecd478d1b65dccc1a3343942490"}
% HMT_RESULT_END:R3-002:residence
% HMT_RESULT_BEGIN:R3-003:residence
% {"material_control":"REV3_BASE_GRANULAR+PASS_NO_PERDIDA_REV3","material_state":"CONSERVADO_POR_BASE_GRANULAR_REV3","residences":[{"path":"/Users/ruben/Documents/HMT2/CONTINUIDAD_EDITORIAL/INTEGRACION_REV3_VON_NEUMANN_MATRICIAL_REALIZACION_2026-08-12/REV3_BASE_GRANULAR.json","sha256":"d89983feac6e064f0ce77d6506d60b94dd1eaf00ed174352d3e32bfc5a9b88d4"}],"result_id":"R3-003","statement_sha256":"1c2e8dc0a9e7ba35279bfc5f48df699c614e4d96af09dcdaf1b76a665374c7cc"}
% HMT_RESULT_END:R3-003:residence
% HMT_RESULT_BEGIN:R3-004:residence
% {"material_control":"REV3_BASE_GRANULAR+PASS_NO_PERDIDA_REV3","material_state":"CONSERVADO_POR_BASE_GRANULAR_REV3","residences":[{"path":"/Users/ruben/Documents/HMT2/CONTINUIDAD_EDITORIAL/INTEGRACION_REV3_VON_NEUMANN_MATRICIAL_REALIZACION_2026-08-12/REV3_BASE_GRANULAR.json","sha256":"d89983feac6e064f0ce77d6506d60b94dd1eaf00ed174352d3e32bfc5a9b88d4"}],"result_id":"R3-004","statement_sha256":"067caf5027a3ca629eaa3f299767eebe36ab8da36d4a594f0490a6fcba2959a3"}
% HMT_RESULT_END:R3-004:residence
% HMT_RESULT_BEGIN:R3-005:residence
% {"material_control":"REV3_BASE_GRANULAR+PASS_NO_PERDIDA_REV3","material_state":"CONSERVADO_POR_BASE_GRANULAR_REV3","residences":[{"path":"/Users/ruben/Documents/HMT2/CONTINUIDAD_EDITORIAL/INTEGRACION_REV3_VON_NEUMANN_MATRICIAL_REALIZACION_2026-08-12/REV3_BASE_GRANULAR.json","sha256":"d89983feac6e064f0ce77d6506d60b94dd1eaf00ed174352d3e32bfc5a9b88d4"}],"result_id":"R3-005","statement_sha256":"7f015f268d0e7772da1a33238cd19903733ee7cfa8947c6372241210ebcd0f01"}
% HMT_RESULT_END:R3-005:residence
% HMT_RESULT_BEGIN:R3-006:residence
% {"material_control":"REV3_BASE_GRANULAR+PASS_NO_PERDIDA_REV3","material_state":"CONSERVADO_POR_BASE_GRANULAR_REV3","residences":[{"path":"/Users/ruben/Documents/HMT2/CONTINUIDAD_EDITORIAL/INTEGRACION_REV3_VON_NEUMANN_MATRICIAL_REALIZACION_2026-08-12/REV3_BASE_GRANULAR.json","sha256":"d89983feac6e064f0ce77d6506d60b94dd1eaf00ed174352d3e32bfc5a9b88d4"}],"result_id":"R3-006","statement_sha256":"9e3f5af6f313c9c4c256d7e6a1f22a0bc0d112d95a7909a4b61b35638a89ea73"}
% HMT_RESULT_END:R3-006:residence
% HMT_RESULT_BEGIN:R3-007:residence
% {"material_control":"REV3_BASE_GRANULAR+PASS_NO_PERDIDA_REV3","material_state":"CONSERVADO_POR_BASE_GRANULAR_REV3","residences":[{"path":"/Users/ruben/Documents/HMT2/CONTINUIDAD_EDITORIAL/INTEGRACION_REV3_VON_NEUMANN_MATRICIAL_REALIZACION_2026-08-12/REV3_BASE_GRANULAR.json","sha256":"d89983feac6e064f0ce77d6506d60b94dd1eaf00ed174352d3e32bfc5a9b88d4"}],"result_id":"R3-007","statement_sha256":"afe6a7f442b4b9b73f09e57dcf993019f691bf05fabf3cf8a866c0cb30e0b20e"}
% HMT_RESULT_END:R3-007:residence
% HMT_RESULT_BEGIN:R3-008:residence
% {"material_control":"REV3_BASE_GRANULAR+PASS_NO_PERDIDA_REV3","material_state":"CONSERVADO_POR_BASE_GRANULAR_REV3","residences":[{"path":"/Users/ruben/Documents/HMT2/CONTINUIDAD_EDITORIAL/INTEGRACION_REV3_VON_NEUMANN_MATRICIAL_REALIZACION_2026-08-12/REV3_BASE_GRANULAR.json","sha256":"d89983feac6e064f0ce77d6506d60b94dd1eaf00ed174352d3e32bfc5a9b88d4"}],"result_id":"R3-008","statement_sha256":"3f0aaf6fa196a66818755e0fe8d0a900ec8170dc90de754fc9b133303d4627b7"}
% HMT_RESULT_END:R3-008:residence
% HMT_RESULT_BEGIN:R3-009:residence
% {"material_control":"REV3_BASE_GRANULAR+PASS_NO_PERDIDA_REV3","material_state":"CONSERVADO_POR_BASE_GRANULAR_REV3","residences":[{"path":"/Users/ruben/Documents/HMT2/CONTINUIDAD_EDITORIAL/INTEGRACION_REV3_VON_NEUMANN_MATRICIAL_REALIZACION_2026-08-12/REV3_BASE_GRANULAR.json","sha256":"d89983feac6e064f0ce77d6506d60b94dd1eaf00ed174352d3e32bfc5a9b88d4"}],"result_id":"R3-009","statement_sha256":"84c9171db979cfcda487379afcdf640f4d78458cf4f7d5697a3ec0b30a2f3096"}
% HMT_RESULT_END:R3-009:residence
% HMT_RESULT_BEGIN:R3-010:residence
% {"material_control":"REV3_BASE_GRANULAR+PASS_NO_PERDIDA_REV3","material_state":"CONSERVADO_POR_BASE_GRANULAR_REV3","residences":[{"path":"/Users/ruben/Documents/HMT2/CONTINUIDAD_EDITORIAL/INTEGRACION_REV3_VON_NEUMANN_MATRICIAL_REALIZACION_2026-08-12/REV3_BASE_GRANULAR.json","sha256":"d89983feac6e064f0ce77d6506d60b94dd1eaf00ed174352d3e32bfc5a9b88d4"}],"result_id":"R3-010","statement_sha256":"3ce9bef7730788d82ad408181bd8942d9cc48d02b7e2b99b36b5f763b3a20968"}
% HMT_RESULT_END:R3-010:residence
% HMT_RESULT_BEGIN:R3-011:residence
% {"material_control":"REV3_BASE_GRANULAR+PASS_NO_PERDIDA_REV3","material_state":"CONSERVADO_POR_BASE_GRANULAR_REV3","residences":[{"path":"/Users/ruben/Documents/HMT2/CONTINUIDAD_EDITORIAL/INTEGRACION_REV3_VON_NEUMANN_MATRICIAL_REALIZACION_2026-08-12/REV3_BASE_GRANULAR.json","sha256":"d89983feac6e064f0ce77d6506d60b94dd1eaf00ed174352d3e32bfc5a9b88d4"}],"result_id":"R3-011","statement_sha256":"5de78d3622665d2d9585dcb93d9f3d762ca07c9959ab042c4c085b8ef5943597"}
% HMT_RESULT_END:R3-011:residence
% HMT_RESULT_BEGIN:R3-012:residence
% {"material_control":"REV3_BASE_GRANULAR+PASS_NO_PERDIDA_REV3","material_state":"CONSERVADO_POR_BASE_GRANULAR_REV3","residences":[{"path":"/Users/ruben/Documents/HMT2/CONTINUIDAD_EDITORIAL/INTEGRACION_REV3_VON_NEUMANN_MATRICIAL_REALIZACION_2026-08-12/REV3_BASE_GRANULAR.json","sha256":"d89983feac6e064f0ce77d6506d60b94dd1eaf00ed174352d3e32bfc5a9b88d4"}],"result_id":"R3-012","statement_sha256":"a873e4e36ef58361dfd662e6f6b08f0484e978fb07f397dab890d355b8ae0df8"}
% HMT_RESULT_END:R3-012:residence
% HMT_RESULT_BEGIN:R3-013:residence
% {"material_control":"REV3_BASE_GRANULAR+PASS_NO_PERDIDA_REV3","material_state":"CONSERVADO_POR_BASE_GRANULAR_REV3","residences":[{"path":"/Users/ruben/Documents/HMT2/CONTINUIDAD_EDITORIAL/INTEGRACION_REV3_VON_NEUMANN_MATRICIAL_REALIZACION_2026-08-12/REV3_BASE_GRANULAR.json","sha256":"d89983feac6e064f0ce77d6506d60b94dd1eaf00ed174352d3e32bfc5a9b88d4"}],"result_id":"R3-013","statement_sha256":"892af66958bb697bdbc9d87aa439e0dcebdecff82281fa756e35aac1088a3c74"}
% HMT_RESULT_END:R3-013:residence
% HMT_RESULT_BEGIN:R3-014:residence
% {"material_control":"REV3_BASE_GRANULAR+PASS_NO_PERDIDA_REV3","material_state":"CONSERVADO_POR_BASE_GRANULAR_REV3","residences":[{"path":"/Users/ruben/Documents/HMT2/CONTINUIDAD_EDITORIAL/INTEGRACION_REV3_VON_NEUMANN_MATRICIAL_REALIZACION_2026-08-12/REV3_BASE_GRANULAR.json","sha256":"d89983feac6e064f0ce77d6506d60b94dd1eaf00ed174352d3e32bfc5a9b88d4"}],"result_id":"R3-014","statement_sha256":"3b135fbe4cbe5d0e292ffcee032f81cef6c993d833f85dde740957bd50df821f"}
% HMT_RESULT_END:R3-014:residence
% HMT_RESULT_BEGIN:R3-015:residence
% {"material_control":"REV3_BASE_GRANULAR+PASS_NO_PERDIDA_REV3","material_state":"CONSERVADO_POR_BASE_GRANULAR_REV3","residences":[{"path":"/Users/ruben/Documents/HMT2/CONTINUIDAD_EDITORIAL/INTEGRACION_REV3_VON_NEUMANN_MATRICIAL_REALIZACION_2026-08-12/REV3_BASE_GRANULAR.json","sha256":"d89983feac6e064f0ce77d6506d60b94dd1eaf00ed174352d3e32bfc5a9b88d4"}],"result_id":"R3-015","statement_sha256":"f94da8f805b5238f940f359590ae265b58be35f6bf4257fac8fa16cf8933c5f8"}
% HMT_RESULT_END:R3-015:residence
% HMT_RESULT_BEGIN:R3-016:residence
% {"material_control":"REV3_BASE_GRANULAR+PASS_NO_PERDIDA_REV3","material_state":"CONSERVADO_POR_BASE_GRANULAR_REV3","residences":[{"path":"/Users/ruben/Documents/HMT2/CONTINUIDAD_EDITORIAL/INTEGRACION_REV3_VON_NEUMANN_MATRICIAL_REALIZACION_2026-08-12/REV3_BASE_GRANULAR.json","sha256":"d89983feac6e064f0ce77d6506d60b94dd1eaf00ed174352d3e32bfc5a9b88d4"}],"result_id":"R3-016","statement_sha256":"660c7b55c80036946d6df58772d5bebbfded71d302316af3ab877a9e3cd1f3b7"}
% HMT_RESULT_END:R3-016:residence
% HMT_RESULT_BEGIN:R3-017:residence
% {"material_control":"REV3_BASE_GRANULAR+PASS_NO_PERDIDA_REV3","material_state":"CONSERVADO_POR_BASE_GRANULAR_REV3","residences":[{"path":"/Users/ruben/Documents/HMT2/CONTINUIDAD_EDITORIAL/INTEGRACION_REV3_VON_NEUMANN_MATRICIAL_REALIZACION_2026-08-12/REV3_BASE_GRANULAR.json","sha256":"d89983feac6e064f0ce77d6506d60b94dd1eaf00ed174352d3e32bfc5a9b88d4"}],"result_id":"R3-017","statement_sha256":"0834862c1f2fd11986bcab78a6a9ad03cd89152459ab8f8a6b83e9d1d4a0d0a8"}
% HMT_RESULT_END:R3-017:residence
% HMT_RESULT_BEGIN:R3-018:residence
% {"material_control":"REV3_BASE_GRANULAR+PASS_NO_PERDIDA_REV3","material_state":"CONSERVADO_POR_BASE_GRANULAR_REV3","residences":[{"path":"/Users/ruben/Documents/HMT2/CONTINUIDAD_EDITORIAL/INTEGRACION_REV3_VON_NEUMANN_MATRICIAL_REALIZACION_2026-08-12/REV3_BASE_GRANULAR.json","sha256":"d89983feac6e064f0ce77d6506d60b94dd1eaf00ed174352d3e32bfc5a9b88d4"}],"result_id":"R3-018","statement_sha256":"7a6ecf5a52d2b3bb92ca8b82bb5e3852c1a1f729a12ca8d39ad7353a61b2259b"}
% HMT_RESULT_END:R3-018:residence
% HMT_RESULT_BEGIN:R3-019:residence
% {"material_control":"REV3_BASE_GRANULAR+PASS_NO_PERDIDA_REV3","material_state":"CONSERVADO_POR_BASE_GRANULAR_REV3","residences":[{"path":"/Users/ruben/Documents/HMT2/CONTINUIDAD_EDITORIAL/INTEGRACION_REV3_VON_NEUMANN_MATRICIAL_REALIZACION_2026-08-12/REV3_BASE_GRANULAR.json","sha256":"d89983feac6e064f0ce77d6506d60b94dd1eaf00ed174352d3e32bfc5a9b88d4"}],"result_id":"R3-019","statement_sha256":"949668449b4331519738227ea1e782aefbfbe4067e371147c0821de5c4e2314c"}
% HMT_RESULT_END:R3-019:residence
% HMT_RESULT_BEGIN:R3-020:residence
% {"material_control":"REV3_BASE_GRANULAR+PASS_NO_PERDIDA_REV3","material_state":"CONSERVADO_POR_BASE_GRANULAR_REV3","residences":[{"path":"/Users/ruben/Documents/HMT2/CONTINUIDAD_EDITORIAL/INTEGRACION_REV3_VON_NEUMANN_MATRICIAL_REALIZACION_2026-08-12/REV3_BASE_GRANULAR.json","sha256":"d89983feac6e064f0ce77d6506d60b94dd1eaf00ed174352d3e32bfc5a9b88d4"}],"result_id":"R3-020","statement_sha256":"e6097d0d922f2e8e30dd23b38e77066a77f46464e19a875d95d0e9c21b9acd23"}
% HMT_RESULT_END:R3-020:residence
% HMT_RESULT_BEGIN:C-001:residence
% {"material_control":"PASS_MATRIZ_ATOMICA_AUTOCONTENIDA_REV3","material_state":"CONTROL_ACTIVO_EN_PLAN_Y_FUENTE","residences":[{"path":"/Users/ruben/Documents/HMT2/CONTINUIDAD_EDITORIAL/INTEGRACION_REV3_VON_NEUMANN_MATRICIAL_REALIZACION_2026-08-12/verificar_matriz_atomica_autocontenida.py","sha256":"a6598c0e1da8e8a8c61acc4150033a66348c039fb2ebf241791b0517dc458472"}],"result_id":"C-001","statement_sha256":"beae79d27d25d820609da41c6347a8561c6d41d85b1ec0525092421ad8721cf3"}
% HMT_RESULT_END:C-001:residence
% HMT_RESULT_BEGIN:C-002:residence
% {"material_control":"PASS_MATRIZ_ATOMICA_AUTOCONTENIDA_REV3","material_state":"CONTROL_ACTIVO_EN_PLAN_Y_FUENTE","residences":[{"path":"/Users/ruben/Documents/HMT2/CONTINUIDAD_EDITORIAL/INTEGRACION_REV3_VON_NEUMANN_MATRICIAL_REALIZACION_2026-08-12/verificar_matriz_atomica_autocontenida.py","sha256":"a6598c0e1da8e8a8c61acc4150033a66348c039fb2ebf241791b0517dc458472"}],"result_id":"C-002","statement_sha256":"5cf2ae9f5be634f470caa6670cf89d72015dc631a63870dc323d018d9f03d0e0"}
% HMT_RESULT_END:C-002:residence
% HMT_RESULT_BEGIN:C-003:residence
% {"material_control":"PASS_MATRIZ_ATOMICA_AUTOCONTENIDA_REV3","material_state":"CONTROL_ACTIVO_EN_PLAN_Y_FUENTE","residences":[{"path":"/Users/ruben/Documents/HMT2/CONTINUIDAD_EDITORIAL/INTEGRACION_REV3_VON_NEUMANN_MATRICIAL_REALIZACION_2026-08-12/verificar_matriz_atomica_autocontenida.py","sha256":"a6598c0e1da8e8a8c61acc4150033a66348c039fb2ebf241791b0517dc458472"}],"result_id":"C-003","statement_sha256":"2cabb42c3d350ac50184c7358426ca161c539e6dfd45b760168861d7461247ad"}
% HMT_RESULT_END:C-003:residence
% HMT_RESULT_BEGIN:C-004:residence
% {"material_control":"PASS_MATRIZ_ATOMICA_AUTOCONTENIDA_REV3","material_state":"CONTROL_ACTIVO_EN_PLAN_Y_FUENTE","residences":[{"path":"/Users/ruben/Documents/HMT2/CONTINUIDAD_EDITORIAL/INTEGRACION_REV3_VON_NEUMANN_MATRICIAL_REALIZACION_2026-08-12/verificar_matriz_atomica_autocontenida.py","sha256":"a6598c0e1da8e8a8c61acc4150033a66348c039fb2ebf241791b0517dc458472"}],"result_id":"C-004","statement_sha256":"9ff96e268873d82ee4ac8616df2a572ca23330c63f928a3e7a67ea10e2d46084"}
% HMT_RESULT_END:C-004:residence
% HMT_RESULT_BEGIN:C-005:residence
% {"material_control":"PASS_MATRIZ_ATOMICA_AUTOCONTENIDA_REV3","material_state":"CONTROL_ACTIVO_EN_PLAN_Y_FUENTE","residences":[{"path":"/Users/ruben/Documents/HMT2/CONTINUIDAD_EDITORIAL/INTEGRACION_REV3_VON_NEUMANN_MATRICIAL_REALIZACION_2026-08-12/verificar_matriz_atomica_autocontenida.py","sha256":"a6598c0e1da8e8a8c61acc4150033a66348c039fb2ebf241791b0517dc458472"}],"result_id":"C-005","statement_sha256":"f28286ac6df6ee110b8bdd8f2de9acd882a13dafee92dd7bc0645accfe4384d5"}
% HMT_RESULT_END:C-005:residence
% HMT_RESULT_BEGIN:C-006:residence
% {"material_control":"PASS_MATRIZ_ATOMICA_AUTOCONTENIDA_REV3","material_state":"CONTROL_ACTIVO_EN_PLAN_Y_FUENTE","residences":[{"path":"/Users/ruben/Documents/HMT2/CONTINUIDAD_EDITORIAL/INTEGRACION_REV3_VON_NEUMANN_MATRICIAL_REALIZACION_2026-08-12/verificar_matriz_atomica_autocontenida.py","sha256":"a6598c0e1da8e8a8c61acc4150033a66348c039fb2ebf241791b0517dc458472"}],"result_id":"C-006","statement_sha256":"61914b1116661f7db384d4e45bba72ec85e7ff331800fc50353e8e93f850dd76"}
% HMT_RESULT_END:C-006:residence
% HMT_RESULT_BEGIN:C-007:residence
% {"material_control":"PASS_MATRIZ_ATOMICA_AUTOCONTENIDA_REV3","material_state":"CONTROL_ACTIVO_EN_PLAN_Y_FUENTE","residences":[{"path":"/Users/ruben/Documents/HMT2/CONTINUIDAD_EDITORIAL/INTEGRACION_REV3_VON_NEUMANN_MATRICIAL_REALIZACION_2026-08-12/verificar_matriz_atomica_autocontenida.py","sha256":"a6598c0e1da8e8a8c61acc4150033a66348c039fb2ebf241791b0517dc458472"}],"result_id":"C-007","statement_sha256":"12e1672ff9d4988000a45bcd1a7b89ef0b97692c2ca5f6756ca06ad9670ede54"}
% HMT_RESULT_END:C-007:residence
% HMT_RESULT_BEGIN:MAS-RECTOR-001:residence
% {"material_control":"PASS_PAQUETE_RECTOR_MASAS_HMT_MD_2026_08_06","material_state":"CONSERVADO_CON_APTO_FOCAL","residences":[{"path":"/Users/ruben/Documents/New project/PUBLICACION_HMT/PAQUETE_RECTOR_LEY_GENERAL_MASAS_HMT_MD_2026-08-06/MATRIZ_RECTORA.tsv","sha256":"d4d58d9e5394804b83f48a2996ef884b275669b0f9468a28c9ed7c0a58affdc5"}],"result_id":"MAS-RECTOR-001","statement_sha256":"0a768984b144ee78bd8e25d36a2945dd57a6f83a5071c4fcfc7ba1a0b68f8d7f"}
% HMT_RESULT_END:MAS-RECTOR-001:residence
% HMT_RESULT_BEGIN:MAS-RECTOR-002:residence
% {"material_control":"PASS_PAQUETE_RECTOR_MASAS_HMT_MD_2026_08_06","material_state":"CONSERVADO_CON_APTO_FOCAL","residences":[{"path":"/Users/ruben/Documents/New project/PUBLICACION_HMT/PAQUETE_RECTOR_LEY_GENERAL_MASAS_HMT_MD_2026-08-06/MATRIZ_RECTORA.tsv","sha256":"d4d58d9e5394804b83f48a2996ef884b275669b0f9468a28c9ed7c0a58affdc5"}],"result_id":"MAS-RECTOR-002","statement_sha256":"89fa9d9d6538d06602255e139245ecc085a236981ab5f973f21ada8ca0f83cd3"}
% HMT_RESULT_END:MAS-RECTOR-002:residence
% HMT_RESULT_BEGIN:MAS-RECTOR-003:residence
% {"material_control":"PASS_PAQUETE_RECTOR_MASAS_HMT_MD_2026_08_06","material_state":"CONSERVADO_CON_APTO_FOCAL","residences":[{"path":"/Users/ruben/Documents/New project/PUBLICACION_HMT/PAQUETE_RECTOR_LEY_GENERAL_MASAS_HMT_MD_2026-08-06/MATRIZ_RECTORA.tsv","sha256":"d4d58d9e5394804b83f48a2996ef884b275669b0f9468a28c9ed7c0a58affdc5"}],"result_id":"MAS-RECTOR-003","statement_sha256":"a859bf6639b8355be0ff5ebcf8ba32630249b5d4842f80ebc5272733f0715109"}
% HMT_RESULT_END:MAS-RECTOR-003:residence
% HMT_RESULT_BEGIN:MAS-RECTOR-004:residence
% {"material_control":"PASS_PAQUETE_RECTOR_MASAS_HMT_MD_2026_08_06","material_state":"CONSERVADO_CON_APTO_FOCAL","residences":[{"path":"/Users/ruben/Documents/New project/PUBLICACION_HMT/PAQUETE_RECTOR_LEY_GENERAL_MASAS_HMT_MD_2026-08-06/MATRIZ_RECTORA.tsv","sha256":"d4d58d9e5394804b83f48a2996ef884b275669b0f9468a28c9ed7c0a58affdc5"}],"result_id":"MAS-RECTOR-004","statement_sha256":"5478fc19a7db005a9c0a02175258a722aa9289c8aae2b40c0f8be1a5539331e6"}
% HMT_RESULT_END:MAS-RECTOR-004:residence
% HMT_RESULT_BEGIN:MAS-RECTOR-005:residence
% {"material_control":"PASS_PAQUETE_RECTOR_MASAS_HMT_MD_2026_08_06","material_state":"CONSERVADO_CON_APTO_FOCAL","residences":[{"path":"/Users/ruben/Documents/New project/PUBLICACION_HMT/PAQUETE_RECTOR_LEY_GENERAL_MASAS_HMT_MD_2026-08-06/MATRIZ_RECTORA.tsv","sha256":"d4d58d9e5394804b83f48a2996ef884b275669b0f9468a28c9ed7c0a58affdc5"}],"result_id":"MAS-RECTOR-005","statement_sha256":"9ac8216130e6ddb43e8032b3eb72c4dd18ab87c05f238ee8fa65f06fdea782d2"}
% HMT_RESULT_END:MAS-RECTOR-005:residence
% HMT_RESULT_BEGIN:MAS-RECTOR-006:residence
% {"material_control":"PASS_PAQUETE_RECTOR_MASAS_HMT_MD_2026_08_06","material_state":"CONSERVADO_CON_APTO_FOCAL","residences":[{"path":"/Users/ruben/Documents/New project/PUBLICACION_HMT/PAQUETE_RECTOR_LEY_GENERAL_MASAS_HMT_MD_2026-08-06/MATRIZ_RECTORA.tsv","sha256":"d4d58d9e5394804b83f48a2996ef884b275669b0f9468a28c9ed7c0a58affdc5"}],"result_id":"MAS-RECTOR-006","statement_sha256":"f64d301daf63969b4ae52cb1259e81411210bb18d9c5e297f815a8f7ea6698e3"}
% HMT_RESULT_END:MAS-RECTOR-006:residence
% HMT_RESULT_BEGIN:MAS-RECTOR-007:residence
% {"material_control":"PASS_PAQUETE_RECTOR_MASAS_HMT_MD_2026_08_06","material_state":"CONSERVADO_CON_APTO_FOCAL","residences":[{"path":"/Users/ruben/Documents/New project/PUBLICACION_HMT/PAQUETE_RECTOR_LEY_GENERAL_MASAS_HMT_MD_2026-08-06/MATRIZ_RECTORA.tsv","sha256":"d4d58d9e5394804b83f48a2996ef884b275669b0f9468a28c9ed7c0a58affdc5"}],"result_id":"MAS-RECTOR-007","statement_sha256":"537017eed6d0ecd896c8a40aed8e3197aeef62bb6895d356b95a7bbc2948643e"}
% HMT_RESULT_END:MAS-RECTOR-007:residence
% HMT_RESULT_BEGIN:MAS-RECTOR-008:residence
% {"material_control":"PASS_PAQUETE_RECTOR_MASAS_HMT_MD_2026_08_06","material_state":"CONSERVADO_CON_APTO_FOCAL","residences":[{"path":"/Users/ruben/Documents/New project/PUBLICACION_HMT/PAQUETE_RECTOR_LEY_GENERAL_MASAS_HMT_MD_2026-08-06/MATRIZ_RECTORA.tsv","sha256":"d4d58d9e5394804b83f48a2996ef884b275669b0f9468a28c9ed7c0a58affdc5"}],"result_id":"MAS-RECTOR-008","statement_sha256":"9563efbb323a0406d27d86e100b82f2b9ca50dd08d2197ca5d02418e4656c614"}
% HMT_RESULT_END:MAS-RECTOR-008:residence
% HMT_RESULT_BEGIN:MAS-RECTOR-009:residence
% {"material_control":"PASS_PAQUETE_RECTOR_MASAS_HMT_MD_2026_08_06","material_state":"CONSERVADO_CON_APTO_FOCAL","residences":[{"path":"/Users/ruben/Documents/New project/PUBLICACION_HMT/PAQUETE_RECTOR_LEY_GENERAL_MASAS_HMT_MD_2026-08-06/MATRIZ_RECTORA.tsv","sha256":"d4d58d9e5394804b83f48a2996ef884b275669b0f9468a28c9ed7c0a58affdc5"}],"result_id":"MAS-RECTOR-009","statement_sha256":"cda0251b6312a403ff360c8c40e29d102a3395cef7e4b3f6ea6b2b72ef982e76"}
% HMT_RESULT_END:MAS-RECTOR-009:residence
% HMT_RESULT_BEGIN:MAS-RECTOR-010:residence
% {"material_control":"PASS_PAQUETE_RECTOR_MASAS_HMT_MD_2026_08_06","material_state":"CONSERVADO_CON_APTO_FOCAL","residences":[{"path":"/Users/ruben/Documents/New project/PUBLICACION_HMT/PAQUETE_RECTOR_LEY_GENERAL_MASAS_HMT_MD_2026-08-06/MATRIZ_RECTORA.tsv","sha256":"d4d58d9e5394804b83f48a2996ef884b275669b0f9468a28c9ed7c0a58affdc5"}],"result_id":"MAS-RECTOR-010","statement_sha256":"5913a73e04b4c752216af319b77e125d31914639200b39b45d1a6b28b536508b"}
% HMT_RESULT_END:MAS-RECTOR-010:residence
% HMT_RESULT_BEGIN:MAS-RECTOR-011:residence
% {"material_control":"PASS_PAQUETE_RECTOR_MASAS_HMT_MD_2026_08_06","material_state":"CONSERVADO_CON_APTO_FOCAL","residences":[{"path":"/Users/ruben/Documents/New project/PUBLICACION_HMT/PAQUETE_RECTOR_LEY_GENERAL_MASAS_HMT_MD_2026-08-06/MATRIZ_RECTORA.tsv","sha256":"d4d58d9e5394804b83f48a2996ef884b275669b0f9468a28c9ed7c0a58affdc5"}],"result_id":"MAS-RECTOR-011","statement_sha256":"b56b1d18f5b984455801b89eace1fbf9143fa965b645cd443b7bbef759b09251"}
% HMT_RESULT_END:MAS-RECTOR-011:residence
% HMT_RESULT_BEGIN:MAS-RECTOR-012:residence
% {"material_control":"PASS_PAQUETE_RECTOR_MASAS_HMT_MD_2026_08_06","material_state":"CONSERVADO_CON_APTO_FOCAL","residences":[{"path":"/Users/ruben/Documents/New project/PUBLICACION_HMT/PAQUETE_RECTOR_LEY_GENERAL_MASAS_HMT_MD_2026-08-06/MATRIZ_RECTORA.tsv","sha256":"d4d58d9e5394804b83f48a2996ef884b275669b0f9468a28c9ed7c0a58affdc5"}],"result_id":"MAS-RECTOR-012","statement_sha256":"aa319b29fb1a3ad47b2d521716fea82da5660479225223eac504798c36b42f24"}
% HMT_RESULT_END:MAS-RECTOR-012:residence
% HMT_RESULT_BEGIN:MAS-RECTOR-013:residence
% {"material_control":"PASS_PAQUETE_RECTOR_MASAS_HMT_MD_2026_08_06","material_state":"CONSERVADO_CON_APTO_FOCAL","residences":[{"path":"/Users/ruben/Documents/New project/PUBLICACION_HMT/PAQUETE_RECTOR_LEY_GENERAL_MASAS_HMT_MD_2026-08-06/MATRIZ_RECTORA.tsv","sha256":"d4d58d9e5394804b83f48a2996ef884b275669b0f9468a28c9ed7c0a58affdc5"}],"result_id":"MAS-RECTOR-013","statement_sha256":"c9f012d4f1b9fa1d9a642670582f65c700f433930615ddc0500e4beef203d5e9"}
% HMT_RESULT_END:MAS-RECTOR-013:residence
% HMT_RESULT_BEGIN:MAS-RECTOR-014:residence
% {"material_control":"PASS_PAQUETE_RECTOR_MASAS_HMT_MD_2026_08_06","material_state":"CONSERVADO_CON_APTO_FOCAL","residences":[{"path":"/Users/ruben/Documents/New project/PUBLICACION_HMT/PAQUETE_RECTOR_LEY_GENERAL_MASAS_HMT_MD_2026-08-06/MATRIZ_RECTORA.tsv","sha256":"d4d58d9e5394804b83f48a2996ef884b275669b0f9468a28c9ed7c0a58affdc5"}],"result_id":"MAS-RECTOR-014","statement_sha256":"12aadab1bbdd24c3cf826e52a333c34f70ced2a37f655b2919eaaaf3281fef04"}
% HMT_RESULT_END:MAS-RECTOR-014:residence
% HMT_RESULT_BEGIN:MAS-RECTOR-015:residence
% {"material_control":"PASS_PAQUETE_RECTOR_MASAS_HMT_MD_2026_08_06","material_state":"CONSERVADO_CON_APTO_FOCAL","residences":[{"path":"/Users/ruben/Documents/New project/PUBLICACION_HMT/PAQUETE_RECTOR_LEY_GENERAL_MASAS_HMT_MD_2026-08-06/MATRIZ_RECTORA.tsv","sha256":"d4d58d9e5394804b83f48a2996ef884b275669b0f9468a28c9ed7c0a58affdc5"}],"result_id":"MAS-RECTOR-015","statement_sha256":"ae01c9c511c63968eed7516b99a2cf7e2a65d0d1c6f56e756638a3d524bb5f9f"}
% HMT_RESULT_END:MAS-RECTOR-015:residence
% HMT_RESULT_BEGIN:MAS-RECTOR-016:residence
% {"material_control":"PASS_PAQUETE_RECTOR_MASAS_HMT_MD_2026_08_06","material_state":"CONSERVADO_CON_APTO_FOCAL","residences":[{"path":"/Users/ruben/Documents/New project/PUBLICACION_HMT/PAQUETE_RECTOR_LEY_GENERAL_MASAS_HMT_MD_2026-08-06/MATRIZ_RECTORA.tsv","sha256":"d4d58d9e5394804b83f48a2996ef884b275669b0f9468a28c9ed7c0a58affdc5"}],"result_id":"MAS-RECTOR-016","statement_sha256":"685df6bdb27ac04026246ca3f09b212f9aabebf22929fb63ddfe1222da2a6e76"}
% HMT_RESULT_END:MAS-RECTOR-016:residence
% HMT_RESULT_BEGIN:MAS-RECTOR-017:residence
% {"material_control":"PASS_PAQUETE_RECTOR_MASAS_HMT_MD_2026_08_06","material_state":"CONSERVADO_CON_APTO_FOCAL","residences":[{"path":"/Users/ruben/Documents/New project/PUBLICACION_HMT/PAQUETE_RECTOR_LEY_GENERAL_MASAS_HMT_MD_2026-08-06/MATRIZ_RECTORA.tsv","sha256":"d4d58d9e5394804b83f48a2996ef884b275669b0f9468a28c9ed7c0a58affdc5"}],"result_id":"MAS-RECTOR-017","statement_sha256":"2d6f3806c75ac84503e6b8ce4c6fe49dbc650e525915991cfef4ecc29e12bb55"}
% HMT_RESULT_END:MAS-RECTOR-017:residence
% HMT_RESULT_BEGIN:MAS-RECTOR-018:residence
% {"material_control":"PASS_PAQUETE_RECTOR_MASAS_HMT_MD_2026_08_06","material_state":"CONSERVADO_CON_APTO_FOCAL","residences":[{"path":"/Users/ruben/Documents/New project/PUBLICACION_HMT/PAQUETE_RECTOR_LEY_GENERAL_MASAS_HMT_MD_2026-08-06/MATRIZ_RECTORA.tsv","sha256":"d4d58d9e5394804b83f48a2996ef884b275669b0f9468a28c9ed7c0a58affdc5"}],"result_id":"MAS-RECTOR-018","statement_sha256":"1fc9dcfc0f77f0e3dcf950b1b3ffb6ee7ee7314e9e8407bb22f75fa7efd5300c"}
% HMT_RESULT_END:MAS-RECTOR-018:residence
% HMT_RESULT_BEGIN:MAS-RECTOR-019:residence
% {"material_control":"PASS_PAQUETE_RECTOR_MASAS_HMT_MD_2026_08_06","material_state":"CONSERVADO_CON_APTO_FOCAL","residences":[{"path":"/Users/ruben/Documents/New project/PUBLICACION_HMT/PAQUETE_RECTOR_LEY_GENERAL_MASAS_HMT_MD_2026-08-06/MATRIZ_RECTORA.tsv","sha256":"d4d58d9e5394804b83f48a2996ef884b275669b0f9468a28c9ed7c0a58affdc5"}],"result_id":"MAS-RECTOR-019","statement_sha256":"ebbe49ae2f4afbc32a86801a2996b2c2d2bf71c0ff4f86d9dc3e58cf6366ae60"}
% HMT_RESULT_END:MAS-RECTOR-019:residence
% HMT_RESULT_BEGIN:MAS-RECTOR-020:residence
% {"material_control":"PASS_PAQUETE_RECTOR_MASAS_HMT_MD_2026_08_06","material_state":"CONSERVADO_CON_APTO_FOCAL","residences":[{"path":"/Users/ruben/Documents/New project/PUBLICACION_HMT/PAQUETE_RECTOR_LEY_GENERAL_MASAS_HMT_MD_2026-08-06/MATRIZ_RECTORA.tsv","sha256":"d4d58d9e5394804b83f48a2996ef884b275669b0f9468a28c9ed7c0a58affdc5"}],"result_id":"MAS-RECTOR-020","statement_sha256":"ea7797a9dc2c4c52d16299bf176fa0b0178158922903c442a8cb1aa85af17879"}
% HMT_RESULT_END:MAS-RECTOR-020:residence
% HMT_RESULT_BEGIN:MAS-RECTOR-021:residence
% {"material_control":"PASS_PAQUETE_RECTOR_MASAS_HMT_MD_2026_08_06","material_state":"CONSERVADO_CON_APTO_FOCAL","residences":[{"path":"/Users/ruben/Documents/New project/PUBLICACION_HMT/PAQUETE_RECTOR_LEY_GENERAL_MASAS_HMT_MD_2026-08-06/MATRIZ_RECTORA.tsv","sha256":"d4d58d9e5394804b83f48a2996ef884b275669b0f9468a28c9ed7c0a58affdc5"}],"result_id":"MAS-RECTOR-021","statement_sha256":"7f875aaae09118d70a3022037a232a45a3e2fe74a8c92aae01e5486aa2ef7f95"}
% HMT_RESULT_END:MAS-RECTOR-021:residence
% HMT_RESULT_BEGIN:MAS-RECTOR-022:residence
% {"material_control":"PASS_PAQUETE_RECTOR_MASAS_HMT_MD_2026_08_06","material_state":"CONSERVADO_CON_APTO_FOCAL","residences":[{"path":"/Users/ruben/Documents/New project/PUBLICACION_HMT/PAQUETE_RECTOR_LEY_GENERAL_MASAS_HMT_MD_2026-08-06/MATRIZ_RECTORA.tsv","sha256":"d4d58d9e5394804b83f48a2996ef884b275669b0f9468a28c9ed7c0a58affdc5"}],"result_id":"MAS-RECTOR-022","statement_sha256":"c6e9560103ad03e63b2460672442d06e18c09c1f866c7d34ee4946f543255d51"}
% HMT_RESULT_END:MAS-RECTOR-022:residence
% HMT_RESULT_BEGIN:MAS-RECTOR-023:residence
% {"material_control":"PASS_PAQUETE_RECTOR_MASAS_HMT_MD_2026_08_06","material_state":"CONSERVADO_CON_APTO_FOCAL","residences":[{"path":"/Users/ruben/Documents/New project/PUBLICACION_HMT/PAQUETE_RECTOR_LEY_GENERAL_MASAS_HMT_MD_2026-08-06/MATRIZ_RECTORA.tsv","sha256":"d4d58d9e5394804b83f48a2996ef884b275669b0f9468a28c9ed7c0a58affdc5"}],"result_id":"MAS-RECTOR-023","statement_sha256":"dd2590f85503d95b6ed6280bb6540f982aadade680f17a9ca0298c3abb7fe3c6"}
% HMT_RESULT_END:MAS-RECTOR-023:residence
% HMT_RESULT_BEGIN:MAS-RECTOR-024:residence
% {"material_control":"PASS_PAQUETE_RECTOR_MASAS_HMT_MD_2026_08_06","material_state":"CONSERVADO_CON_APTO_FOCAL","residences":[{"path":"/Users/ruben/Documents/New project/PUBLICACION_HMT/PAQUETE_RECTOR_LEY_GENERAL_MASAS_HMT_MD_2026-08-06/MATRIZ_RECTORA.tsv","sha256":"d4d58d9e5394804b83f48a2996ef884b275669b0f9468a28c9ed7c0a58affdc5"}],"result_id":"MAS-RECTOR-024","statement_sha256":"4d3e6032d279cb845f8bb3fe2e3bcd153b9585996d56fd93c0e721793c358eae"}
% HMT_RESULT_END:MAS-RECTOR-024:residence
% HMT_RESULT_BEGIN:MAS-RECTOR-025:residence
% {"material_control":"PASS_PAQUETE_RECTOR_MASAS_HMT_MD_2026_08_06","material_state":"CONSERVADO_CON_APTO_FOCAL","residences":[{"path":"/Users/ruben/Documents/New project/PUBLICACION_HMT/PAQUETE_RECTOR_LEY_GENERAL_MASAS_HMT_MD_2026-08-06/MATRIZ_RECTORA.tsv","sha256":"d4d58d9e5394804b83f48a2996ef884b275669b0f9468a28c9ed7c0a58affdc5"}],"result_id":"MAS-RECTOR-025","statement_sha256":"6e3bb65bdf6f87c71551bf92d9f6b8e7c60ba9785857316070c5d88c9a59fb9a"}
% HMT_RESULT_END:MAS-RECTOR-025:residence
% HMT_RESULT_BEGIN:MAS-RECTOR-026:residence
% {"material_control":"PASS_PAQUETE_RECTOR_MASAS_HMT_MD_2026_08_06","material_state":"CONSERVADO_CON_APTO_FOCAL","residences":[{"path":"/Users/ruben/Documents/New project/PUBLICACION_HMT/PAQUETE_RECTOR_LEY_GENERAL_MASAS_HMT_MD_2026-08-06/MATRIZ_RECTORA.tsv","sha256":"d4d58d9e5394804b83f48a2996ef884b275669b0f9468a28c9ed7c0a58affdc5"}],"result_id":"MAS-RECTOR-026","statement_sha256":"33fe0ca572be6ea55887ec4b42e78250c112e42699de77e13eaf5ddfdf608e12"}
% HMT_RESULT_END:MAS-RECTOR-026:residence
% HMT_RESULT_BEGIN:MAS-RECTOR-027:residence
% {"material_control":"PASS_PAQUETE_RECTOR_MASAS_HMT_MD_2026_08_06","material_state":"CONSERVADO_CON_APTO_FOCAL","residences":[{"path":"/Users/ruben/Documents/New project/PUBLICACION_HMT/PAQUETE_RECTOR_LEY_GENERAL_MASAS_HMT_MD_2026-08-06/MATRIZ_RECTORA.tsv","sha256":"d4d58d9e5394804b83f48a2996ef884b275669b0f9468a28c9ed7c0a58affdc5"}],"result_id":"MAS-RECTOR-027","statement_sha256":"6168ae29b3e71bf7885c166f21fc815528ad39b20fa81a628be34700e902949c"}
% HMT_RESULT_END:MAS-RECTOR-027:residence
% HMT_RESULT_BEGIN:MAS-RECTOR-028:residence
% {"material_control":"PASS_PAQUETE_RECTOR_MASAS_HMT_MD_2026_08_06","material_state":"CONSERVADO_CON_APTO_FOCAL","residences":[{"path":"/Users/ruben/Documents/New project/PUBLICACION_HMT/PAQUETE_RECTOR_LEY_GENERAL_MASAS_HMT_MD_2026-08-06/MATRIZ_RECTORA.tsv","sha256":"d4d58d9e5394804b83f48a2996ef884b275669b0f9468a28c9ed7c0a58affdc5"}],"result_id":"MAS-RECTOR-028","statement_sha256":"52f3251b7aec29d108b1b6c06d3cc064189cfe5cad74e31c311b1eab0be4b3d1"}
% HMT_RESULT_END:MAS-RECTOR-028:residence
% HMT_RESULT_BEGIN:MAS-PUBLIC-023:residence
% {"material_control":"PASS_PAQUETE_RECTOR_MASAS_HMT_MD_2026_08_06","material_state":"CONSERVADO_CON_APTO_FOCAL","residences":[{"path":"/Users/ruben/Documents/New project/PUBLICACION_HMT/PAQUETE_RECTOR_LEY_GENERAL_MASAS_HMT_MD_2026-08-06/MATRIZ_RECTORA.tsv","sha256":"d4d58d9e5394804b83f48a2996ef884b275669b0f9468a28c9ed7c0a58affdc5"}],"result_id":"MAS-PUBLIC-023","statement_sha256":"2490d5f9a177b2e5c480be8158dc8b1bccbb5f0785cb1e610be8e85f3a6b8918"}
% HMT_RESULT_END:MAS-PUBLIC-023:residence
% HMT_RESULT_BEGIN:MAS-COND-017:residence
% {"material_control":"PASS_PAQUETE_RECTOR_MASAS_HMT_MD_2026_08_06","material_state":"CONSERVADO_CON_APTO_FOCAL","residences":[{"path":"/Users/ruben/Documents/New project/PUBLICACION_HMT/PAQUETE_RECTOR_LEY_GENERAL_MASAS_HMT_MD_2026-08-06/MATRIZ_RECTORA.tsv","sha256":"d4d58d9e5394804b83f48a2996ef884b275669b0f9468a28c9ed7c0a58affdc5"}],"result_id":"MAS-COND-017","statement_sha256":"d9b0c82973bb442a515af30c9bfe351d9b9ab183eb81db8250f32c615d5f50ba"}
% HMT_RESULT_END:MAS-COND-017:residence
% HMT_RESULT_BEGIN:INF-001:residence
% {"material_control":"GATE_MATERIAL_64+PASS_PUBLICACION_INTEGRAL","material_state":"INTEGRADO_O_REMITIDO_SIN_PERDIDA","residences":[{"path":"manuscrito/sections/hmt/05_tpk_numero_medicion.tex","sha256":"51deffafc13a6d83bc07e15f77892d1ec32ba05f54cc2b01d6417ea13e3599a2"}],"result_id":"INF-001","statement_sha256":"d2d4fc4dc0a639778129fc6046b7f974c7a2cc39dae5b6aede5bd2ce96d91b1d"}
% HMT_RESULT_END:INF-001:residence
% HMT_RESULT_BEGIN:INF-002:residence
% {"material_control":"GATE_MATERIAL_64+PASS_PUBLICACION_INTEGRAL","material_state":"INTEGRADO_O_REMITIDO_SIN_PERDIDA","residences":[{"path":"manuscrito/sections/hmt/05_tpk_numero_medicion.tex","sha256":"51deffafc13a6d83bc07e15f77892d1ec32ba05f54cc2b01d6417ea13e3599a2"}],"result_id":"INF-002","statement_sha256":"8e32b4ab88538fa1ca2b52a628278c2f084dc54f423cf78093a3270700d157a5"}
% HMT_RESULT_END:INF-002:residence
% HMT_RESULT_BEGIN:INF-003:residence
% {"material_control":"GATE_MATERIAL_64+PASS_PUBLICACION_INTEGRAL","material_state":"INTEGRADO_O_REMITIDO_SIN_PERDIDA","residences":[{"path":"manuscrito/sections/hmt/05_tpk_numero_medicion.tex","sha256":"51deffafc13a6d83bc07e15f77892d1ec32ba05f54cc2b01d6417ea13e3599a2"}],"result_id":"INF-003","statement_sha256":"667c50d6578b21e1abcff2bc241db42bbe6e270c7268fef0d64172e29d096836"}
% HMT_RESULT_END:INF-003:residence
% HMT_RESULT_BEGIN:N9-001:residence
% {"material_control":"GATE_MATERIAL_64+PASS_PUBLICACION_INTEGRAL","material_state":"INTEGRADO_O_REMITIDO_SIN_PERDIDA","residences":[{"path":"manuscrito/sections/hmt/06aa_genealogia_filtracion_g9_autonoma.tex","sha256":"ba4816ac627ae2d7eb4141421a19a240bf66022f8728f4ce8bf873bd439efb39"}],"result_id":"N9-001","statement_sha256":"f38e038d91253c1364af870bc92f1e5a243d0a5e5c69c5f85bc4d05cc2ef6fb9"}
% HMT_RESULT_END:N9-001:residence
% HMT_RESULT_BEGIN:N9-002:residence
% {"material_control":"GATE_MATERIAL_64+PASS_PUBLICACION_INTEGRAL","material_state":"INTEGRADO_O_REMITIDO_SIN_PERDIDA","residences":[{"path":"manuscrito/sections/hmt/06aa_genealogia_filtracion_g9_autonoma.tex","sha256":"ba4816ac627ae2d7eb4141421a19a240bf66022f8728f4ce8bf873bd439efb39"}],"result_id":"N9-002","statement_sha256":"e4c29231470c3fe69fe80b927a6249121b8a79287a407ca4318e1889a8ed9fb6"}
% HMT_RESULT_END:N9-002:residence
% HMT_RESULT_BEGIN:N9-003:residence
% {"material_control":"GATE_MATERIAL_64+PASS_PUBLICACION_INTEGRAL","material_state":"INTEGRADO_O_REMITIDO_SIN_PERDIDA","residences":[{"path":"manuscrito/sections/hmt/06b_cinco_ciclos_pi.tex","sha256":"83f8a2be69639ac01847a0e25720d3e58b193ae689abfe5fd567a552fb301012"},{"path":"manuscrito/sections/hmt/06d_biografias_numeros_rev6.tex","sha256":"a01607a1a6f5ee4fadbd482daa384b09b192d315da15048ff661d2baf7d197f6"},{"path":"manuscrito/sections/hmt/06e_certificado_generacion_monodromica_rev7.tex","sha256":"2a0cf8028e12605a97923c53726fcb62c1798c458af8e9f78f12c1129b0e1b94"}],"result_id":"N9-003","statement_sha256":"a3b1e710d9bbef1129d8ddb4da929c36fe0b2eaf7ea7e20bbc91312eff72cb57"}
% HMT_RESULT_END:N9-003:residence
% HMT_RESULT_BEGIN:N9-004:residence
% {"material_control":"GATE_MATERIAL_64+PASS_PUBLICACION_INTEGRAL","material_state":"INTEGRADO_O_REMITIDO_SIN_PERDIDA","residences":[{"path":"manuscrito/sections/hmt/06e_certificado_generacion_monodromica_rev7.tex","sha256":"2a0cf8028e12605a97923c53726fcb62c1798c458af8e9f78f12c1129b0e1b94"}],"result_id":"N9-004","statement_sha256":"5288a67a9b48b60a9a1896b482933707612c1ae1e4b4862ba2b23f9ff393e77f"}
% HMT_RESULT_END:N9-004:residence
% HMT_RESULT_BEGIN:HEL-001:residence
% {"material_control":"GATE_MATERIAL_64+PASS_PUBLICACION_INTEGRAL","material_state":"INTEGRADO_O_REMITIDO_SIN_PERDIDA","residences":[{"path":"manuscrito/sections/hmt/03l_elevacion_cohomologica_dimension_rev2.tex","sha256":"1b0f6b86dc275a3947b64e319b0b19789d62c83bcc3714d5c0b8a2cea2243a8c"}],"result_id":"HEL-001","statement_sha256":"3ed30cd71d3b16fcdaf015298295033782765571ae3a2c7d60f65aa8b0bee52d"}
% HMT_RESULT_END:HEL-001:residence
% HMT_RESULT_BEGIN:OPS-001:residence
% {"material_control":"GATE_MATERIAL_64+PASS_PUBLICACION_INTEGRAL","material_state":"INTEGRADO_O_REMITIDO_SIN_PERDIDA","residences":[{"path":"manuscrito/sections/hmt/03k_genealogia_operaciones_rev2.tex","sha256":"fae7aca94c6fb20754431ed79a5aa2f005a7cf2f5681736fc5826da4c694fb98"}],"result_id":"OPS-001","statement_sha256":"61ed21a147f32626c42baccde2ffe7594b14cf6c1c83912b96d5f3f8776405e5"}
% HMT_RESULT_END:OPS-001:residence
% HMT_RESULT_BEGIN:ZF-001:residence
% {"material_control":"GATE_MATERIAL_64+PASS_PUBLICACION_INTEGRAL","material_state":"INTEGRADO_O_REMITIDO_SIN_PERDIDA","residences":[{"path":"manuscrito/sections/integracion_20260812/von_neumann_64/04_zfc_fundacion_ordinales.tex","sha256":"5d81fd31bf31a950f9d068f278e76ef0ebfb37cb8cf6ee5483575ab87f5fddb0"}],"result_id":"ZF-001","statement_sha256":"1d6985e3c4c3690eb43465daa6bcc3d7fb39cb3d0ae1fe31a3cd6c70db621314"}
% HMT_RESULT_END:ZF-001:residence
% HMT_RESULT_BEGIN:ZF-002:residence
% {"material_control":"GATE_MATERIAL_64+PASS_PUBLICACION_INTEGRAL","material_state":"INTEGRADO_O_REMITIDO_SIN_PERDIDA","residences":[{"path":"manuscrito/sections/integracion_20260812/von_neumann_64/04_zfc_fundacion_ordinales.tex","sha256":"5d81fd31bf31a950f9d068f278e76ef0ebfb37cb8cf6ee5483575ab87f5fddb0"}],"result_id":"ZF-002","statement_sha256":"e78afa0785eb7978d61364af148ab43c40092242ad1edf80db0e5c19f11915cd"}
% HMT_RESULT_END:ZF-002:residence
% HMT_RESULT_BEGIN:EXT-001:residence
% {"material_control":"GATE_MATERIAL_64+PASS_PUBLICACION_INTEGRAL","material_state":"INTEGRADO_O_REMITIDO_SIN_PERDIDA","residences":[{"path":"manuscrito/sections/hmt/14d_ontologia_numeros_rev11.tex","sha256":"c96b3c94af4668f15c6a061260d1c1900c52b69b8a4a5d2442268f6ff1d0abcc"}],"result_id":"EXT-001","statement_sha256":"664eeafa4a1a0cfad63aa9414daaccdf8f0ea8a15f92851a674ad7cd5d26c9a8"}
% HMT_RESULT_END:EXT-001:residence
% HMT_RESULT_BEGIN:FND-001:residence
% {"material_control":"GATE_MATERIAL_64+PASS_PUBLICACION_INTEGRAL","material_state":"INTEGRADO_O_REMITIDO_SIN_PERDIDA","residences":[{"path":"manuscrito/sections/integracion_20260812/von_neumann_64/04_zfc_fundacion_ordinales.tex","sha256":"5d81fd31bf31a950f9d068f278e76ef0ebfb37cb8cf6ee5483575ab87f5fddb0"},{"path":"manuscrito/sections/hmt/00a_continuo_numero_medida_hmt_revfinal.tex","sha256":"1be7648180a45eb337db43d6d18a63624876b9f568ea29998c80b00d13b1a819"}],"result_id":"FND-001","statement_sha256":"6b0867dcae5d6924d26cbfdfb829827b9b60d7461148fab933762bd8ab425ac7"}
% HMT_RESULT_END:FND-001:residence
% HMT_RESULT_BEGIN:ORD-001:residence
% {"material_control":"GATE_MATERIAL_64+PASS_PUBLICACION_INTEGRAL","material_state":"INTEGRADO_O_REMITIDO_SIN_PERDIDA","residences":[{"path":"manuscrito/sections/integracion_20260812/von_neumann_64/04_zfc_fundacion_ordinales.tex","sha256":"5d81fd31bf31a950f9d068f278e76ef0ebfb37cb8cf6ee5483575ab87f5fddb0"},{"path":"manuscrito/sections/hmt/14f_cierre_genealogico_numero.tex","sha256":"4522230d8297a1420ccb586d0d6c1fefcb08b5e8e2deeda85b11cc73afd8ded5"}],"result_id":"ORD-001","statement_sha256":"1c7c37b7877465ee0663efb162f12aa0614e44999754b59066edcfa489f7f6ae"}
% HMT_RESULT_END:ORD-001:residence
% HMT_RESULT_BEGIN:ORD-002:residence
% {"material_control":"GATE_MATERIAL_64+PASS_PUBLICACION_INTEGRAL","material_state":"INTEGRADO_O_REMITIDO_SIN_PERDIDA","residences":[{"path":"manuscrito/sections/integracion_20260812/von_neumann_64/04_zfc_fundacion_ordinales.tex","sha256":"5d81fd31bf31a950f9d068f278e76ef0ebfb37cb8cf6ee5483575ab87f5fddb0"}],"result_id":"ORD-002","statement_sha256":"c8188c7c87197d0287f9d7bde9efd04e7246384d5271159053cefc1b62278bf1"}
% HMT_RESULT_END:ORD-002:residence
% HMT_RESULT_BEGIN:PWR-001:residence
% {"material_control":"GATE_MATERIAL_64+PASS_PUBLICACION_INTEGRAL","material_state":"INTEGRADO_O_REMITIDO_SIN_PERDIDA","residences":[{"path":"manuscrito/sections/integracion_20260812/von_neumann_64/05_potencia_eleccion.tex","sha256":"cb412bdee19cc8e65181f75c0666a966dc20111f1dafcaee563bf1ee0e3d2684"}],"result_id":"PWR-001","statement_sha256":"15e7dd2441fdb8f717dbe8bbfb825a15a1165399b2158782c52216b71f1cb8e4"}
% HMT_RESULT_END:PWR-001:residence
% HMT_RESULT_BEGIN:PWR-002:residence
% {"material_control":"GATE_MATERIAL_64+PASS_PUBLICACION_INTEGRAL","material_state":"INTEGRADO_O_REMITIDO_SIN_PERDIDA","residences":[{"path":"manuscrito/sections/integracion_20260812/von_neumann_64/05_potencia_eleccion.tex","sha256":"cb412bdee19cc8e65181f75c0666a966dc20111f1dafcaee563bf1ee0e3d2684"}],"result_id":"PWR-002","statement_sha256":"667fc09531d0de7134ded48a16fb499d56212ba6d7734d3cbc1e04ec1144ad52"}
% HMT_RESULT_END:PWR-002:residence
% HMT_RESULT_BEGIN:CHO-001:residence
% {"material_control":"GATE_MATERIAL_64+PASS_PUBLICACION_INTEGRAL","material_state":"INTEGRADO_O_REMITIDO_SIN_PERDIDA","residences":[{"path":"manuscrito/sections/integracion_20260812/von_neumann_64/05_potencia_eleccion.tex","sha256":"cb412bdee19cc8e65181f75c0666a966dc20111f1dafcaee563bf1ee0e3d2684"},{"path":"manuscrito/sections/hmt/14f_cierre_genealogico_numero.tex","sha256":"4522230d8297a1420ccb586d0d6c1fefcb08b5e8e2deeda85b11cc73afd8ded5"}],"result_id":"CHO-001","statement_sha256":"8b2d6fc14cefc2d7b601d295b4441c23d96827e08dd639c47877a67cb5b36929"}
% HMT_RESULT_END:CHO-001:residence
% HMT_RESULT_BEGIN:CHO-002:residence
% {"material_control":"GATE_MATERIAL_64+PASS_PUBLICACION_INTEGRAL","material_state":"INTEGRADO_O_REMITIDO_SIN_PERDIDA","residences":[{"path":"manuscrito/sections/integracion_20260812/von_neumann_64/05_potencia_eleccion.tex","sha256":"cb412bdee19cc8e65181f75c0666a966dc20111f1dafcaee563bf1ee0e3d2684"}],"result_id":"CHO-002","statement_sha256":"26fee3b2fa8f6cd1bc1073d172a9682b085c58a85cae5e2a6864e89d45f7ddea"}
% HMT_RESULT_END:CHO-002:residence
% HMT_RESULT_BEGIN:CAT-001:residence
% {"material_control":"GATE_MATERIAL_64+PASS_PUBLICACION_INTEGRAL","material_state":"INTEGRADO_O_REMITIDO_SIN_PERDIDA","residences":[{"path":"manuscrito/sections/integracion_20260812/von_neumann_64/04_zfc_fundacion_ordinales.tex","sha256":"5d81fd31bf31a950f9d068f278e76ef0ebfb37cb8cf6ee5483575ab87f5fddb0"},{"path":"manuscrito/sections/hmt/00b_complejo_discreto_medida_rev2.tex","sha256":"c4bd51e17232daf45c4fe95b5863bfa0f945f2a926d22ec6c9c8eccdfa5f1dd8"},{"path":"manuscrito/sections/hmt/03e_funtor_fase_longitud.tex","sha256":"b1a874c3ecb9a434f01fa3ed05c0cef286ac57a85a9269945a994c08df70e78b"},{"path":"manuscrito/sections/hmt/14f_cierre_genealogico_numero.tex","sha256":"4522230d8297a1420ccb586d0d6c1fefcb08b5e8e2deeda85b11cc73afd8ded5"}],"result_id":"CAT-001","statement_sha256":"cd69268cc61050d66556a6a33bd6fb31ab1fc444f678229ff48d09dfeeb7ec1b"}
% HMT_RESULT_END:CAT-001:residence
% HMT_RESULT_BEGIN:TOP-001:residence
% {"material_control":"GATE_MATERIAL_64+PASS_PUBLICACION_INTEGRAL","material_state":"INTEGRADO_O_REMITIDO_SIN_PERDIDA","residences":[{"path":"manuscrito/sections/integracion_20260812/von_neumann_64/04_zfc_fundacion_ordinales.tex","sha256":"5d81fd31bf31a950f9d068f278e76ef0ebfb37cb8cf6ee5483575ab87f5fddb0"},{"path":"manuscrito/sections/hmt/05_tpk_numero_medicion.tex","sha256":"51deffafc13a6d83bc07e15f77892d1ec32ba05f54cc2b01d6417ea13e3599a2"}],"result_id":"TOP-001","statement_sha256":"cf2c0f501167aac25e8c173e2c9c6cca8aa4cc123f668b06c5a7ca8c42f21eb3"}
% HMT_RESULT_END:TOP-001:residence
% HMT_RESULT_BEGIN:HIL-001:residence
% {"material_control":"GATE_MATERIAL_64+PASS_PUBLICACION_INTEGRAL","material_state":"INTEGRADO_O_REMITIDO_SIN_PERDIDA","residences":[{"path":"manuscrito/sections/integracion_20260812/von_neumann_64/06_hilbert_heisenberg.tex","sha256":"6f35e2e9b56af3cbf33d20c28ae1b5bfac2b92f4a0827f444be39774046e130b"},{"path":"manuscrito/sections/hmt/03l_elevacion_cohomologica_dimension_rev2.tex","sha256":"1b0f6b86dc275a3947b64e319b0b19789d62c83bcc3714d5c0b8a2cea2243a8c"}],"result_id":"HIL-001","statement_sha256":"7529a7dd47eb815a976287f21bd184c40d9cc3dfb0dcb1d2c3f1a2b4b4eeaaa7"}
% HMT_RESULT_END:HIL-001:residence
% HMT_RESULT_BEGIN:UHF-001:residence
% {"material_control":"GATE_MATERIAL_64+PASS_PUBLICACION_INTEGRAL","material_state":"INTEGRADO_O_REMITIDO_SIN_PERDIDA","residences":[{"path":"manuscrito/sections/integracion_20260812/von_neumann_64/07_uhf_factor_ii1.tex","sha256":"a3f413168ac1514605dc7e9bf8a7866e73420380c5a93390fbd5f5d55fa25765"},{"path":"manuscrito/sections/hmt/03l_elevacion_cohomologica_dimension_rev2.tex","sha256":"1b0f6b86dc275a3947b64e319b0b19789d62c83bcc3714d5c0b8a2cea2243a8c"}],"result_id":"UHF-001","statement_sha256":"0cf650bcfd502563a4ff8cab59c8c8c4ae997565ebed1f058034f5cce7e5b378"}
% HMT_RESULT_END:UHF-001:residence
% HMT_RESULT_BEGIN:II1-001:residence
% {"material_control":"GATE_MATERIAL_64+PASS_PUBLICACION_INTEGRAL","material_state":"INTEGRADO_O_REMITIDO_SIN_PERDIDA","residences":[{"path":"manuscrito/sections/integracion_20260812/von_neumann_64/07_uhf_factor_ii1.tex","sha256":"a3f413168ac1514605dc7e9bf8a7866e73420380c5a93390fbd5f5d55fa25765"},{"path":"manuscrito/sections/hmt/03l_elevacion_cohomologica_dimension_rev2.tex","sha256":"1b0f6b86dc275a3947b64e319b0b19789d62c83bcc3714d5c0b8a2cea2243a8c"}],"result_id":"II1-001","statement_sha256":"dc01c84423637eb3704db747bb7a96106332b3cc6f8f36258e1af6cf02c8ead6"}
% HMT_RESULT_END:II1-001:residence
% HMT_RESULT_BEGIN:DIM-001:residence
% {"material_control":"GATE_MATERIAL_64+PASS_PUBLICACION_INTEGRAL","material_state":"INTEGRADO_O_REMITIDO_SIN_PERDIDA","residences":[{"path":"manuscrito/sections/integracion_20260812/von_neumann_64/07_uhf_factor_ii1.tex","sha256":"a3f413168ac1514605dc7e9bf8a7866e73420380c5a93390fbd5f5d55fa25765"},{"path":"manuscrito/sections/hmt/03l_elevacion_cohomologica_dimension_rev2.tex","sha256":"1b0f6b86dc275a3947b64e319b0b19789d62c83bcc3714d5c0b8a2cea2243a8c"}],"result_id":"DIM-001","statement_sha256":"9b8936bc39190cf70a463a94ca6fb311c2d8d8e66262fb2a3d44d831c5653c4b"}
% HMT_RESULT_END:DIM-001:residence
% HMT_RESULT_BEGIN:CLK-001:residence
% {"material_control":"GATE_MATERIAL_64+PASS_PUBLICACION_INTEGRAL","material_state":"INTEGRADO_O_REMITIDO_SIN_PERDIDA","residences":[{"path":"manuscrito/sections/integracion_20260812/von_neumann_64/08_reloj_logica_ergodicidad.tex","sha256":"aeda0eee8e4376ae44e5f6505cf4414ba767a64705a111c66ef2917d7411a0d9"},{"path":"manuscrito/sections/hmt/19_rotaciones_vacancias_criticidad.tex","sha256":"2d46e077e597811a4b0545b2a6ded98b5d56d0d66c65d71f416a1d06c282c363"}],"result_id":"CLK-001","statement_sha256":"b89435d53375a0c540e471c24439fbb5b83ca598038eaa6e60ad19138fe1d106"}
% HMT_RESULT_END:CLK-001:residence
% HMT_RESULT_BEGIN:INT-001:residence
% {"material_control":"GATE_MATERIAL_64+PASS_PUBLICACION_INTEGRAL","material_state":"INTEGRADO_O_REMITIDO_SIN_PERDIDA","residences":[{"path":"manuscrito/sections/integracion_20260812/von_neumann_64/08_reloj_logica_ergodicidad.tex","sha256":"aeda0eee8e4376ae44e5f6505cf4414ba767a64705a111c66ef2917d7411a0d9"}],"result_id":"INT-001","statement_sha256":"aa1ffba490fcc7437280d3a98fe5f9a543b85b037100711af707edb8fdc41527"}
% HMT_RESULT_END:INT-001:residence
% HMT_RESULT_BEGIN:QLG-001:residence
% {"material_control":"GATE_MATERIAL_64+PASS_PUBLICACION_INTEGRAL","material_state":"INTEGRADO_O_REMITIDO_SIN_PERDIDA","residences":[{"path":"manuscrito/sections/integracion_20260812/von_neumann_64/08_reloj_logica_ergodicidad.tex","sha256":"aeda0eee8e4376ae44e5f6505cf4414ba767a64705a111c66ef2917d7411a0d9"},{"path":"manuscrito/sections/integracion_20260812/von_neumann_64/06_hilbert_heisenberg.tex","sha256":"6f35e2e9b56af3cbf33d20c28ae1b5bfac2b92f4a0827f444be39774046e130b"},{"path":"manuscrito/sections/hmt/03_app_ortograma.tex","sha256":"aaf0fd1dec2c270d544cd19562b4461c57ebb009437e051d973be43395428390"}],"result_id":"QLG-001","statement_sha256":"5af886b5ee8e9ef572b0db2f8860b44c9a5685f03b042ad7934083b6157af5d3"}
% HMT_RESULT_END:QLG-001:residence
% HMT_RESULT_BEGIN:ERG-001:residence
% {"material_control":"GATE_MATERIAL_64+PASS_PUBLICACION_INTEGRAL","material_state":"INTEGRADO_O_REMITIDO_SIN_PERDIDA","residences":[{"path":"manuscrito/sections/integracion_20260812/von_neumann_64/08_reloj_logica_ergodicidad.tex","sha256":"aeda0eee8e4376ae44e5f6505cf4414ba767a64705a111c66ef2917d7411a0d9"},{"path":"manuscrito/sections/hmt/19_rotaciones_vacancias_criticidad.tex","sha256":"2d46e077e597811a4b0545b2a6ded98b5d56d0d66c65d71f416a1d06c282c363"}],"result_id":"ERG-001","statement_sha256":"0e5e1bd65b68a7b812b621df955227e4bd7880b25b43edc0e1a9db24e3a80b8e"}
% HMT_RESULT_END:ERG-001:residence
% HMT_RESULT_BEGIN:DEN-001:residence
% {"material_control":"GATE_MATERIAL_64+PASS_PUBLICACION_INTEGRAL","material_state":"INTEGRADO_O_REMITIDO_SIN_PERDIDA","residences":[{"path":"manuscrito/sections/integracion_20260812/von_neumann_64/08_reloj_logica_ergodicidad.tex","sha256":"aeda0eee8e4376ae44e5f6505cf4414ba767a64705a111c66ef2917d7411a0d9"},{"path":"manuscrito/sections/md/04ab_postulados_cuanticos_genealogicos_rev2.tex","sha256":"1895f84bd75b104e873a3d73c5b0e8ae1b5df606bd6feefeb4f9666a92966804"}],"result_id":"DEN-001","statement_sha256":"d3f54d28f9a40e95ef70680ed00ece354f13bef3fa9616bf65771ff984899f55"}
% HMT_RESULT_END:DEN-001:residence
% HMT_RESULT_BEGIN:LAN-001:residence
% {"material_control":"GATE_MATERIAL_64+PASS_PUBLICACION_INTEGRAL","material_state":"INTEGRADO_O_REMITIDO_SIN_PERDIDA","residences":[{"path":"manuscrito/sections/md/04g_informacion_reversible_landauer_rev7.tex","sha256":"3f60a3da983f28f197fe487896eda3667d44247d0ad30f6f3d1a598b368f0e8e"}],"result_id":"LAN-001","statement_sha256":"197b4b23802d24d8e911bcd3787ba07daf3a5694505b103e27acf6a540021b23"}
% HMT_RESULT_END:LAN-001:residence
% HMT_RESULT_BEGIN:LAN-002:residence
% {"material_control":"GATE_MATERIAL_64+PASS_PUBLICACION_INTEGRAL","material_state":"INTEGRADO_O_REMITIDO_SIN_PERDIDA","residences":[{"path":"manuscrito/sections/integracion_20260812/bloque_64_landauer.tex","sha256":"ef8b27ff2b888ef3c5aa3228093c84ed139d63a8e8a504a4b6da8b56f256cbe6"}],"result_id":"LAN-002","statement_sha256":"aa495d4ee76ccc4ec19bbfee8fca847aa39dfb3e7785b0654264e883afa8e9f3"}
% HMT_RESULT_END:LAN-002:residence
% HMT_RESULT_BEGIN:GDL-001:residence
% {"material_control":"GATE_MATERIAL_64+PASS_PUBLICACION_INTEGRAL","material_state":"INTEGRADO_O_REMITIDO_SIN_PERDIDA","residences":[{"path":"manuscrito/sections/integracion_20260812/von_neumann_64/10_godel_turing.tex","sha256":"c17b938966df53455f1927258ae4e4fd2603927ea1827a08620dfdbb386e7863"}],"result_id":"GDL-001","statement_sha256":"e65a071d6f27cf044d0d665fa43c1c77e8b1c800787b0879ee63f693a954b4eb"}
% HMT_RESULT_END:GDL-001:residence
% HMT_RESULT_BEGIN:GDL-002:residence
% {"material_control":"GATE_MATERIAL_64+PASS_PUBLICACION_INTEGRAL","material_state":"INTEGRADO_O_REMITIDO_SIN_PERDIDA","residences":[{"path":"manuscrito/sections/integracion_20260812/von_neumann_64/10_godel_turing.tex","sha256":"c17b938966df53455f1927258ae4e4fd2603927ea1827a08620dfdbb386e7863"}],"result_id":"GDL-002","statement_sha256":"30af7308ca4cbed2f47ea0f0aaf875baaf611eda5be904a33538de7e5ee41864"}
% HMT_RESULT_END:GDL-002:residence
% HMT_RESULT_BEGIN:IND-001:residence
% {"material_control":"GATE_MATERIAL_64+PASS_PUBLICACION_INTEGRAL","material_state":"INTEGRADO_O_REMITIDO_SIN_PERDIDA","residences":[{"path":"manuscrito/sections/integracion_20260812/von_neumann_64/11_forcing_borel_minimax.tex","sha256":"c4b14714d216d63c6db3b930602a3cd7953d40dd61d50a093bd50d840ccfc33c"}],"result_id":"IND-001","statement_sha256":"5e22be0c9938a9dceab95188d043eebbbffab47eced16a62dc2ded6048b7891d"}
% HMT_RESULT_END:IND-001:residence
% HMT_RESULT_BEGIN:FOR-001:residence
% {"material_control":"GATE_MATERIAL_64+PASS_PUBLICACION_INTEGRAL","material_state":"INTEGRADO_O_REMITIDO_SIN_PERDIDA","residences":[{"path":"manuscrito/sections/integracion_20260812/von_neumann_64/11_forcing_borel_minimax.tex","sha256":"c4b14714d216d63c6db3b930602a3cd7953d40dd61d50a093bd50d840ccfc33c"},{"path":"manuscrito/sections/hmt/05_tpk_numero_medicion.tex","sha256":"51deffafc13a6d83bc07e15f77892d1ec32ba05f54cc2b01d6417ea13e3599a2"}],"result_id":"FOR-001","statement_sha256":"1a8d72a91afdb3a1411a9d0431a099066dffdd6a4603d5d3bc5cb570330fa919"}
% HMT_RESULT_END:FOR-001:residence
% HMT_RESULT_BEGIN:FOR-002:residence
% {"material_control":"GATE_MATERIAL_64+PASS_PUBLICACION_INTEGRAL","material_state":"INTEGRADO_O_REMITIDO_SIN_PERDIDA","residences":[{"path":"manuscrito/sections/integracion_20260812/von_neumann_64/11_forcing_borel_minimax.tex","sha256":"c4b14714d216d63c6db3b930602a3cd7953d40dd61d50a093bd50d840ccfc33c"}],"result_id":"FOR-002","statement_sha256":"59f4f58f84a7735d60c6fd2454ba7965e566bd09c58b3d45392f878cb3dec369"}
% HMT_RESULT_END:FOR-002:residence
% HMT_RESULT_BEGIN:BOR-001:residence
% {"material_control":"GATE_MATERIAL_64+PASS_PUBLICACION_INTEGRAL","material_state":"INTEGRADO_O_REMITIDO_SIN_PERDIDA","residences":[{"path":"manuscrito/sections/integracion_20260812/von_neumann_64/11_forcing_borel_minimax.tex","sha256":"c4b14714d216d63c6db3b930602a3cd7953d40dd61d50a093bd50d840ccfc33c"}],"result_id":"BOR-001","statement_sha256":"ac8e973306e90a8e1806a33402b6ef6e087583edf2b49e475fcf64fc7bffc827"}
% HMT_RESULT_END:BOR-001:residence
% HMT_RESULT_BEGIN:MIN-001:residence
% {"material_control":"GATE_MATERIAL_64+PASS_PUBLICACION_INTEGRAL","material_state":"INTEGRADO_O_REMITIDO_SIN_PERDIDA","residences":[{"path":"manuscrito/sections/integracion_20260812/von_neumann_64/11_forcing_borel_minimax.tex","sha256":"c4b14714d216d63c6db3b930602a3cd7953d40dd61d50a093bd50d840ccfc33c"}],"result_id":"MIN-001","statement_sha256":"ed4856dacc07059c530b3404157bcc15541d967d4e205eaba615ed2e7f06bd2e"}
% HMT_RESULT_END:MIN-001:residence
% HMT_RESULT_BEGIN:III-001:residence
% {"material_control":"GATE_MATERIAL_64+PASS_PUBLICACION_INTEGRAL","material_state":"INTEGRADO_O_REMITIDO_SIN_PERDIDA","residences":[{"path":"manuscrito/sections/integracion_20260812/von_neumann_64/08_reloj_logica_ergodicidad.tex","sha256":"aeda0eee8e4376ae44e5f6505cf4414ba767a64705a111c66ef2917d7411a0d9"},{"path":"manuscrito/sections/md/04e_termodinamica_rutas.tex","sha256":"ab1429f9145eee89e5287095b94f327291fa4aa247ae0abcc95e3f5611de87c1"}],"result_id":"III-001","statement_sha256":"ed27269e518d94d2dd09835c4f6d5c6be473aa08906e8ed2d633e2686f5bb4a7"}
% HMT_RESULT_END:III-001:residence
% HMT_RESULT_BEGIN:TRI-001:residence
% {"material_control":"GATE_MATERIAL_64+PASS_PUBLICACION_INTEGRAL","material_state":"INTEGRADO_O_REMITIDO_SIN_PERDIDA","residences":[{"path":"manuscrito/sections/integracion_20260812/von_neumann_64/12_cierre_triple_doble_proyeccion.tex","sha256":"9773623f121e57b194d95202c9b01d24af38a7f1e391fd219beb76ce9dbf2a86"},{"path":"manuscrito/sections/md/10c_puente_retorno_holonomia_cartan.tex","sha256":"b206dc7147f5eb4295dace2d1960e61cb5f6c4b40a0f6828535d1ba4024a7ba9"}],"result_id":"TRI-001","statement_sha256":"7717c12ab00203b9ef7e04dfea02661344e850d6889ac9e04b4906e296f8b383"}
% HMT_RESULT_END:TRI-001:residence
% HMT_RESULT_BEGIN:EXC-001:residence
% {"material_control":"GATE_MATERIAL_64+PASS_PUBLICACION_INTEGRAL","material_state":"INTEGRADO_O_REMITIDO_SIN_PERDIDA","residences":[{"path":"manuscrito/sections/integracion_20260812/von_neumann_64/12_cierre_triple_doble_proyeccion.tex","sha256":"9773623f121e57b194d95202c9b01d24af38a7f1e391fd219beb76ce9dbf2a86"},{"path":"manuscrito/sections/hmt/09_doble_proyeccion.tex","sha256":"ff538add4f3afedca7f4cbd5cea553e3b8132e05117117c40be30e55d33dcdba"},{"path":"manuscrito/sections/hmt/15_bandera_y_rama_excepcional.tex","sha256":"6c42a7012881b036232a5e12c6e2ce382d48964ef9ed12cc0a4d5256126154e5"},{"path":"manuscrito/sections/hmt/15c_cierre_micro_macro_nivel_tres.tex","sha256":"0a5cf2476d9825b9005494bfab1b937deddbba26d5c69f76a7822077be3797c5"},{"path":"manuscrito/sections/hmt/15d_voa_moonshine_orden_dos_rev11.tex","sha256":"6bc88137eb614a0bba834212c9f8f95be09e48e8f570f6c8126ac33dbe007e86"}],"result_id":"EXC-001","statement_sha256":"b9c1e5194c3ddd0cd113377935663265ca76fce3648e3ea5576c0e8b41aa6747"}
% HMT_RESULT_END:EXC-001:residence
% HMT_RESULT_BEGIN:MD-001:residence
% {"material_control":"GATE_MATERIAL_64+PASS_PUBLICACION_INTEGRAL","material_state":"INTEGRADO_O_REMITIDO_SIN_PERDIDA","residences":[{"path":"manuscrito/sections/integracion_20260812/von_neumann_64/12_cierre_triple_doble_proyeccion.tex","sha256":"9773623f121e57b194d95202c9b01d24af38a7f1e391fd219beb76ce9dbf2a86"},{"path":"manuscrito/sections/md/05f_operadores_masa_two_way_rev11.tex","sha256":"e0c576ff54f405c41adea7093ca6b117134e9322749c969f1fc4cca6e8615d35"},{"path":"manuscrito/sections/md/06h_atlas_familias_rev6.tex","sha256":"f6cf1c90f4d42d24598380ca0bff5c9ee66c3269aad819c08d85d32a31066367"}],"result_id":"MD-001","statement_sha256":"3def61d573286af3af73231496757ed80209391421e2f5b759252b1e4f75a582"}
% HMT_RESULT_END:MD-001:residence
% HMT_RESULT_BEGIN:PSC-001:residence
% {"material_control":"GATE_MATERIAL_64+PASS_PUBLICACION_INTEGRAL","material_state":"INTEGRADO_O_REMITIDO_SIN_PERDIDA","residences":[{"path":"manuscrito/sections/integracion_20260812/von_neumann_64/11_forcing_borel_minimax.tex","sha256":"c4b14714d216d63c6db3b930602a3cd7953d40dd61d50a093bd50d840ccfc33c"}],"result_id":"PSC-001","statement_sha256":"aa11f37512474aa438b0ea7a0c6719d2337b671cab88d78f0bfbda69d2d029d4"}
% HMT_RESULT_END:PSC-001:residence
% HMT_RESULT_BEGIN:COV-001:residence
% {"material_control":"GATE_MATERIAL_64+PASS_PUBLICACION_INTEGRAL","material_state":"INTEGRADO_O_REMITIDO_SIN_PERDIDA","residences":[{"path":"/Users/ruben/Documents/HMT2/CONTINUIDAD_EDITORIAL/INTEGRACION_REV3_VON_NEUMANN_MATRICIAL_REALIZACION_2026-08-12/MATRIZ_ATOMICA_AUTOCONTENIDA_REV3_V4.tsv","sha256":"a89bbb749d42ca8e123b0549913409dadda6196187891b4bf7f83b993d2ffa9b"}],"result_id":"COV-001","statement_sha256":"2d2a3767b8f662d7a68bd271a690c4b557d2bc3f3a256077632ca30e518bd4d7"}
% HMT_RESULT_END:COV-001:residence
% HMT_RESULT_BEGIN:IINF-001:residence
% {"material_control":"GATE_MATERIAL_64+PASS_PUBLICACION_INTEGRAL","material_state":"INTEGRADO_O_REMITIDO_SIN_PERDIDA","residences":[{"path":"manuscrito/sections/integracion_20260812/von_neumann_64/07_uhf_factor_ii1.tex","sha256":"a3f413168ac1514605dc7e9bf8a7866e73420380c5a93390fbd5f5d55fa25765"}],"result_id":"IINF-001","statement_sha256":"217587fdc70a6738c237837099cbbba8d41f6d59c40619c866bf48c95e0198d7"}
% HMT_RESULT_END:IINF-001:residence
% HMT_RESULT_BEGIN:DEN-003:residence
% {"material_control":"GATE_MATERIAL_64+PASS_PUBLICACION_INTEGRAL","material_state":"INTEGRADO_O_REMITIDO_SIN_PERDIDA","residences":[{"path":"manuscrito/sections/integracion_20260812/von_neumann_64/08_reloj_logica_ergodicidad.tex","sha256":"aeda0eee8e4376ae44e5f6505cf4414ba767a64705a111c66ef2917d7411a0d9"}],"result_id":"DEN-003","statement_sha256":"921dc3fb0d30a2764747bddc32233eb05b1589bbcbb17f4a98f3aae0ea462a13"}
% HMT_RESULT_END:DEN-003:residence
% HMT_RESULT_BEGIN:CARD-001:residence
% {"material_control":"GATE_MATERIAL_64+PASS_PUBLICACION_INTEGRAL","material_state":"INTEGRADO_O_REMITIDO_SIN_PERDIDA","residences":[{"path":"manuscrito/sections/integracion_20260812/von_neumann_64/04_zfc_fundacion_ordinales.tex","sha256":"5d81fd31bf31a950f9d068f278e76ef0ebfb37cb8cf6ee5483575ab87f5fddb0"},{"path":"manuscrito/sections/hmt/00_construccion_app_gnomon_rev11.tex","sha256":"f027dd8cb64274e9f984a180e01da910cdb1fc6235d856ad6dc736f7477381b9"}],"result_id":"CARD-001","statement_sha256":"657e91ba07505f78793eeac979c673b8bc6ec88f1adc9067345b9c06ef806c89"}
% HMT_RESULT_END:CARD-001:residence
% HMT_RESULT_BEGIN:PWR-003:residence
% {"material_control":"GATE_MATERIAL_64+PASS_PUBLICACION_INTEGRAL","material_state":"INTEGRADO_O_REMITIDO_SIN_PERDIDA","residences":[{"path":"manuscrito/sections/integracion_20260812/von_neumann_64/05_potencia_eleccion.tex","sha256":"cb412bdee19cc8e65181f75c0666a966dc20111f1dafcaee563bf1ee0e3d2684"},{"path":"manuscrito/sections/hmt/09_doble_proyeccion.tex","sha256":"ff538add4f3afedca7f4cbd5cea553e3b8132e05117117c40be30e55d33dcdba"}],"result_id":"PWR-003","statement_sha256":"eaefebdd49a0763208853c604924034d7d4c705c6371b2990f7fec1b29b0032a"}
% HMT_RESULT_END:PWR-003:residence
% HMT_RESULT_BEGIN:WEYL-001:residence
% {"material_control":"GATE_MATERIAL_64+PASS_PUBLICACION_INTEGRAL","material_state":"INTEGRADO_O_REMITIDO_SIN_PERDIDA","residences":[{"path":"manuscrito/sections/integracion_20260812/von_neumann_64/06_hilbert_heisenberg.tex","sha256":"6f35e2e9b56af3cbf33d20c28ae1b5bfac2b92f4a0827f444be39774046e130b"},{"path":"manuscrito/sections/hmt/14_numero_heisenberg_y_d108.tex","sha256":"36c7812d22940e488c5a5434ffb85399d9430a12d32f6a376138c58474b3247b"}],"result_id":"WEYL-001","statement_sha256":"3c9c006467a020b44e1f152943e542f7efd63f1fbafa62c1459d799fb5a7fb8d"}
% HMT_RESULT_END:WEYL-001:residence
% HMT_RESULT_BEGIN:DSC-001:residence
% {"material_control":"GATE_MATERIAL_64+PASS_PUBLICACION_INTEGRAL","material_state":"INTEGRADO_O_REMITIDO_SIN_PERDIDA","residences":[{"path":"manuscrito/sections/md/02_pantalla_accion.tex","sha256":"40796ae0bdab2ff0e8c2cc2f07181c0f5efc77c14aa4cca2d5334e899dc0f1ef"}],"result_id":"DSC-001","statement_sha256":"67da6399420386a7c52ecce6942df44aef3e17b49f2b6c2ea0005b3f95cae4ca"}
% HMT_RESULT_END:DSC-001:residence
% HMT_RESULT_BEGIN:REC-001:residence
% {"material_control":"GATE_MATERIAL_64+PASS_PUBLICACION_INTEGRAL","material_state":"INTEGRADO_O_REMITIDO_SIN_PERDIDA","residences":[{"path":"manuscrito/sections/md/04h_entropias_tipadas_rev2.tex","sha256":"f47cec558bb3f375af069dfde73f82b69bf48d9f168a635630f0e9a25bab3d91"}],"result_id":"REC-001","statement_sha256":"0bd633094af48b759b62a734d6cefa9feb253f605c7052145c392284c4b6a32f"}
% HMT_RESULT_END:REC-001:residence
% HMT_RESULT_BEGIN:LAN-003:residence
% {"material_control":"GATE_MATERIAL_64+PASS_PUBLICACION_INTEGRAL","material_state":"INTEGRADO_O_REMITIDO_SIN_PERDIDA","residences":[{"path":"manuscrito/sections/md/04g_informacion_reversible_landauer_rev7.tex","sha256":"3f60a3da983f28f197fe487896eda3667d44247d0ad30f6f3d1a598b368f0e8e"}],"result_id":"LAN-003","statement_sha256":"2a5b4ef14486d6da04b11ef915d58b57331c845e9d106cf402e883204d561861"}
% HMT_RESULT_END:LAN-003:residence
% HMT_RESULT_BEGIN:FIN-001:residence
% {"material_control":"GATE_MATERIAL_64+PASS_PUBLICACION_INTEGRAL","material_state":"INTEGRADO_O_REMITIDO_SIN_PERDIDA","residences":[{"path":"manuscrito/sections/integracion_20260812/von_neumann_64/10_godel_turing.tex","sha256":"c17b938966df53455f1927258ae4e4fd2603927ea1827a08620dfdbb386e7863"}],"result_id":"FIN-001","statement_sha256":"a73d6ef20cd13e14f35306811ecee9eb3057af5c97ddf492c7ba1f86dd04f1ce"}
% HMT_RESULT_END:FIN-001:residence
% HMT_RESULT_BEGIN:DEN-002:residence
% {"material_control":"GATE_MATERIAL_64+PASS_PUBLICACION_INTEGRAL","material_state":"INTEGRADO_O_REMITIDO_SIN_PERDIDA","residences":[{"path":"manuscrito/sections/integracion_20260812/von_neumann_64/08_reloj_logica_ergodicidad.tex","sha256":"aeda0eee8e4376ae44e5f6505cf4414ba767a64705a111c66ef2917d7411a0d9"}],"result_id":"DEN-002","statement_sha256":"35b2ba3c1b4020b9163f1ef3d443b6bcb565b09c8a8622d543cf8e4f57289896"}
% HMT_RESULT_END:DEN-002:residence
% HMT_RESULT_BEGIN:UNI-001:residence
% {"material_control":"GATE_MATERIAL_64+PASS_PUBLICACION_INTEGRAL","material_state":"INTEGRADO_O_REMITIDO_SIN_PERDIDA","residences":[{"path":"manuscrito/sections/integracion_20260812/von_neumann_64/12_cierre_triple_doble_proyeccion.tex","sha256":"9773623f121e57b194d95202c9b01d24af38a7f1e391fd219beb76ce9dbf2a86"}],"result_id":"UNI-001","statement_sha256":"c937c36037e7f572961ab452004513b8c1c29a3ba0d578a7b0da6590a849a422"}
% HMT_RESULT_END:UNI-001:residence
% HMT_RESULT_BEGIN:UNI-002:residence
% {"material_control":"GATE_MATERIAL_64+PASS_PUBLICACION_INTEGRAL","material_state":"INTEGRADO_O_REMITIDO_SIN_PERDIDA","residences":[{"path":"manuscrito/sections/integracion_20260812/von_neumann_64/12_cierre_triple_doble_proyeccion.tex","sha256":"9773623f121e57b194d95202c9b01d24af38a7f1e391fd219beb76ce9dbf2a86"}],"result_id":"UNI-002","statement_sha256":"6c0f67e5980d3e7df62bb7613076467f3e8c1e71d7bd85ae9b8e5a5a6e6eefac"}
% HMT_RESULT_END:UNI-002:residence
% HMT_RESULT_BEGIN:P01:residence
% {"material_control":"PASS_PUENTE_GENEALOGICO_MATRICIAL","material_state":"INTEGRADO_EN_GRAFO_ACTIVO","residences":[{"path":"manuscrito/sections/integracion_20260812/matricial_52/02_hmt_en_profundidad.tex","sha256":"fe6c23eac09e147199090a091fe2564e5587b7b1155ea85518d2bf47aade4cfc"}],"result_id":"P01","statement_sha256":"53d8dc07ede69f64f0d718b277459101f4824f60da865637ded87662edab5245"}
% HMT_RESULT_END:P01:residence
% HMT_RESULT_BEGIN:P02:residence
% {"material_control":"PASS_PUENTE_GENEALOGICO_MATRICIAL","material_state":"INTEGRADO_EN_GRAFO_ACTIVO","residences":[{"path":"manuscrito/sections/integracion_20260812/matricial_52/02_hmt_en_profundidad.tex","sha256":"fe6c23eac09e147199090a091fe2564e5587b7b1155ea85518d2bf47aade4cfc"}],"result_id":"P02","statement_sha256":"69c00b394120ef47751db68d34eb4b918078692555e4fa5e663b2f0297a8548c"}
% HMT_RESULT_END:P02:residence
% HMT_RESULT_BEGIN:P03:residence
% {"material_control":"PASS_PUENTE_GENEALOGICO_MATRICIAL","material_state":"INTEGRADO_EN_GRAFO_ACTIVO","residences":[{"path":"manuscrito/sections/integracion_20260812/matricial_52/02_hmt_en_profundidad.tex","sha256":"fe6c23eac09e147199090a091fe2564e5587b7b1155ea85518d2bf47aade4cfc"}],"result_id":"P03","statement_sha256":"2fe6041370eb6e955a3d0ba0f2e07811bc0aebd09c6793150785e46a510b9359"}
% HMT_RESULT_END:P03:residence
% HMT_RESULT_BEGIN:P04:residence
% {"material_control":"PASS_PUENTE_GENEALOGICO_MATRICIAL","material_state":"INTEGRADO_EN_GRAFO_ACTIVO","residences":[{"path":"manuscrito/sections/integracion_20260812/matricial_52/02_hmt_en_profundidad.tex","sha256":"fe6c23eac09e147199090a091fe2564e5587b7b1155ea85518d2bf47aade4cfc"}],"result_id":"P04","statement_sha256":"c9c80bd3c7cc1697f5e94b0d7401b0649ab52d4f48fd42943fe4bd715aa71c05"}
% HMT_RESULT_END:P04:residence
% HMT_RESULT_BEGIN:P05:residence
% {"material_control":"PASS_PUENTE_GENEALOGICO_MATRICIAL","material_state":"INTEGRADO_EN_GRAFO_ACTIVO","residences":[{"path":"manuscrito/sections/integracion_20260812/matricial_52/02_hmt_en_profundidad.tex","sha256":"fe6c23eac09e147199090a091fe2564e5587b7b1155ea85518d2bf47aade4cfc"}],"result_id":"P05","statement_sha256":"87558450ea6ada90154488b513914bbcbe56a44e45cb40411baa8054fa0a9b11"}
% HMT_RESULT_END:P05:residence
% HMT_RESULT_BEGIN:P06:residence
% {"material_control":"PASS_PUENTE_GENEALOGICO_MATRICIAL","material_state":"INTEGRADO_EN_GRAFO_ACTIVO","residences":[{"path":"manuscrito/sections/integracion_20260812/matricial_52/02_hmt_en_profundidad.tex","sha256":"fe6c23eac09e147199090a091fe2564e5587b7b1155ea85518d2bf47aade4cfc"}],"result_id":"P06","statement_sha256":"b3609e5998fdaf08a9020b1ff80a65442e473516d2ecf24b5b68474c1f6c94c0"}
% HMT_RESULT_END:P06:residence
% HMT_RESULT_BEGIN:P07:residence
% {"material_control":"PASS_PUENTE_GENEALOGICO_MATRICIAL","material_state":"INTEGRADO_EN_GRAFO_ACTIVO","residences":[{"path":"manuscrito/sections/integracion_20260812/matricial_52/02_hmt_en_profundidad.tex","sha256":"fe6c23eac09e147199090a091fe2564e5587b7b1155ea85518d2bf47aade4cfc"},{"path":"manuscrito/sections/md/04ab_postulados_cuanticos_genealogicos_rev2.tex","sha256":"1895f84bd75b104e873a3d73c5b0e8ae1b5df606bd6feefeb4f9666a92966804"}],"result_id":"P07","statement_sha256":"f8147d53ff5923fc1c43835f30a9f07bf581e5e699eb63dd74ac1414533a629e"}
% HMT_RESULT_END:P07:residence
% HMT_RESULT_BEGIN:P08:residence
% {"material_control":"PASS_PUENTE_GENEALOGICO_MATRICIAL","material_state":"INTEGRADO_EN_GRAFO_ACTIVO","residences":[{"path":"manuscrito/sections/integracion_20260812/matricial_52/02_hmt_en_profundidad.tex","sha256":"fe6c23eac09e147199090a091fe2564e5587b7b1155ea85518d2bf47aade4cfc"}],"result_id":"P08","statement_sha256":"f9a8d4c5ec3b2afef25f422ad348af58e28287d1d33e38fd8b67b7529c84166d"}
% HMT_RESULT_END:P08:residence
% HMT_RESULT_BEGIN:P09:residence
% {"material_control":"PASS_PUENTE_GENEALOGICO_MATRICIAL","material_state":"INTEGRADO_EN_GRAFO_ACTIVO","residences":[{"path":"manuscrito/sections/integracion_20260812/matricial_52/01_objeto_y_diccionario.tex","sha256":"95d254edbb6dd7184d3a5db5ad83006aee1adbbdd4ee0a5f631e89d6300648da"}],"result_id":"P09","statement_sha256":"cb5502d8e34b02758b04dda42755bf0a91a83e8ba03dbc45f8e4453942b6ce77"}
% HMT_RESULT_END:P09:residence
% HMT_RESULT_BEGIN:P10:residence
% {"material_control":"PASS_PUENTE_GENEALOGICO_MATRICIAL","material_state":"INTEGRADO_EN_GRAFO_ACTIVO","residences":[{"path":"manuscrito/sections/integracion_20260812/matricial_52/01_objeto_y_diccionario.tex","sha256":"95d254edbb6dd7184d3a5db5ad83006aee1adbbdd4ee0a5f631e89d6300648da"}],"result_id":"P10","statement_sha256":"42f48b2ff4254875c1d787ae3c1d25e5f5e086049b6b5f8e963f462ceab4b4d9"}
% HMT_RESULT_END:P10:residence
% HMT_RESULT_BEGIN:P11:residence
% {"material_control":"PASS_PUENTE_GENEALOGICO_MATRICIAL","material_state":"INTEGRADO_EN_GRAFO_ACTIVO","residences":[{"path":"manuscrito/sections/integracion_20260812/matricial_52/01_objeto_y_diccionario.tex","sha256":"95d254edbb6dd7184d3a5db5ad83006aee1adbbdd4ee0a5f631e89d6300648da"}],"result_id":"P11","statement_sha256":"bcd1e3fbe8b15b11e4f83caac1d5b8879ef6f7e0ff3881fdd1ffba891a6f32ce"}
% HMT_RESULT_END:P11:residence
% HMT_RESULT_BEGIN:P12:residence
% {"material_control":"PASS_PUENTE_GENEALOGICO_MATRICIAL","material_state":"INTEGRADO_EN_GRAFO_ACTIVO","residences":[{"path":"manuscrito/sections/integracion_20260812/matricial_52/04_cuadrados_cuanticos.tex","sha256":"7131342598b20a77a08f6d9d98e12d41ec47e331372ea6c19754df21a58ab585"}],"result_id":"P12","statement_sha256":"3eca57321c7a0f5f618f0ffad0376d39eb8c8ab5968847fb959fdbb09fc734ab"}
% HMT_RESULT_END:P12:residence
% HMT_RESULT_BEGIN:P13:residence
% {"material_control":"PASS_PUENTE_GENEALOGICO_MATRICIAL","material_state":"INTEGRADO_EN_GRAFO_ACTIVO","residences":[{"path":"manuscrito/sections/integracion_20260812/matricial_52/04_cuadrados_cuanticos.tex","sha256":"7131342598b20a77a08f6d9d98e12d41ec47e331372ea6c19754df21a58ab585"}],"result_id":"P13","statement_sha256":"55db6d54634c091c6187c36369faa76bdf9f508b14a4a3b18f5a348d2c9cf011"}
% HMT_RESULT_END:P13:residence
% HMT_RESULT_BEGIN:P14:residence
% {"material_control":"PASS_PUENTE_GENEALOGICO_MATRICIAL","material_state":"INTEGRADO_EN_GRAFO_ACTIVO","residences":[{"path":"manuscrito/sections/integracion_20260812/matricial_52/03_bigraduacion_ucp.tex","sha256":"c054ec7cb105bc2fda9c3228bfaee808f01db65e38e126a8865a3420ed37864d"}],"result_id":"P14","statement_sha256":"57d9a1d7ea29cb176674b4f73e0033f34f82419195809fa7e8ff1f235d95e410"}
% HMT_RESULT_END:P14:residence
% HMT_RESULT_BEGIN:P15:residence
% {"material_control":"PASS_PUENTE_GENEALOGICO_MATRICIAL","material_state":"INTEGRADO_EN_GRAFO_ACTIVO","residences":[{"path":"manuscrito/sections/integracion_20260812/matricial_52/03_bigraduacion_ucp.tex","sha256":"c054ec7cb105bc2fda9c3228bfaee808f01db65e38e126a8865a3420ed37864d"}],"result_id":"P15","statement_sha256":"83ae0271d5db4482e0f2f1c4ec1fc0a1b3e5bbde57a18fc5a75a65cf03960e6f"}
% HMT_RESULT_END:P15:residence
% HMT_RESULT_BEGIN:P16:residence
% {"material_control":"PASS_PUENTE_GENEALOGICO_MATRICIAL","material_state":"INTEGRADO_EN_GRAFO_ACTIVO","residences":[{"path":"manuscrito/sections/integracion_20260812/matricial_52/03_bigraduacion_ucp.tex","sha256":"c054ec7cb105bc2fda9c3228bfaee808f01db65e38e126a8865a3420ed37864d"}],"result_id":"P16","statement_sha256":"bdf8eff661630a7ba8edf744bf36bded26cbd2417e7ef7e7259b9c6c7505a21e"}
% HMT_RESULT_END:P16:residence
% HMT_RESULT_BEGIN:P17:residence
% {"material_control":"PASS_PUENTE_GENEALOGICO_MATRICIAL","material_state":"INTEGRADO_EN_GRAFO_ACTIVO","residences":[{"path":"manuscrito/sections/integracion_20260812/matricial_52/03_bigraduacion_ucp.tex","sha256":"c054ec7cb105bc2fda9c3228bfaee808f01db65e38e126a8865a3420ed37864d"}],"result_id":"P17","statement_sha256":"b0cce98edf7fa782d1064e6b7ac175d850f4ad264c1b6d8e62288f2568ad45f5"}
% HMT_RESULT_END:P17:residence
% HMT_RESULT_BEGIN:P18:residence
% {"material_control":"PASS_PUENTE_GENEALOGICO_MATRICIAL","material_state":"INTEGRADO_EN_GRAFO_ACTIVO","residences":[{"path":"manuscrito/sections/integracion_20260812/matricial_52/04_cuadrados_cuanticos.tex","sha256":"7131342598b20a77a08f6d9d98e12d41ec47e331372ea6c19754df21a58ab585"}],"result_id":"P18","statement_sha256":"1473f4abe925a8e28298eb8bf1dfd39fb9af561154966420712347199f61b2ce"}
% HMT_RESULT_END:P18:residence
% HMT_RESULT_BEGIN:P19:residence
% {"material_control":"PASS_PUENTE_GENEALOGICO_MATRICIAL","material_state":"INTEGRADO_EN_GRAFO_ACTIVO","residences":[{"path":"manuscrito/sections/integracion_20260812/matricial_52/04_cuadrados_cuanticos.tex","sha256":"7131342598b20a77a08f6d9d98e12d41ec47e331372ea6c19754df21a58ab585"}],"result_id":"P19","statement_sha256":"fd3e9565106ad07b0ffeff64653cdd8f4e25c01bd8125c56b954e4ad58015bce"}
% HMT_RESULT_END:P19:residence
% HMT_RESULT_BEGIN:P20:residence
% {"material_control":"PASS_PUENTE_GENEALOGICO_MATRICIAL","material_state":"INTEGRADO_EN_GRAFO_ACTIVO","residences":[{"path":"manuscrito/sections/integracion_20260812/matricial_52/04_cuadrados_cuanticos.tex","sha256":"7131342598b20a77a08f6d9d98e12d41ec47e331372ea6c19754df21a58ab585"}],"result_id":"P20","statement_sha256":"f0301b81ceca519949c82bef3e2c9c9d0415d1943aadd55e64bafd83dc369e25"}
% HMT_RESULT_END:P20:residence
% HMT_RESULT_BEGIN:P21:residence
% {"material_control":"PASS_PUENTE_GENEALOGICO_MATRICIAL","material_state":"INTEGRADO_EN_GRAFO_ACTIVO","residences":[{"path":"manuscrito/sections/integracion_20260812/matricial_52/08_sistemas_conicos_y_universalidad.tex","sha256":"2b94acee5d5df46c4c6e2c4b9048962de8961703fc585df45d08a42f9f4ae0d2"}],"result_id":"P21","statement_sha256":"9677a2892d35d3020f115f04121dc2d8f34f6bcda52cdb28ac4ad708ca6fcfe2"}
% HMT_RESULT_END:P21:residence
% HMT_RESULT_BEGIN:P22:residence
% {"material_control":"PASS_PUENTE_GENEALOGICO_MATRICIAL","material_state":"INTEGRADO_EN_GRAFO_ACTIVO","residences":[{"path":"manuscrito/sections/integracion_20260812/matricial_52/05_geometria_simplectica_y_lectores.tex","sha256":"087482ebc4a880dba1096a9364dab1a05b0f0bddd3495b6f5b967fb9ccb79ccd"}],"result_id":"P22","statement_sha256":"f79d8ca5659786ae2dfdcb69f4a77431e62f85c8e159a0cb789c16dd900f0e08"}
% HMT_RESULT_END:P22:residence
% HMT_RESULT_BEGIN:P23:residence
% {"material_control":"PASS_PUENTE_GENEALOGICO_MATRICIAL","material_state":"INTEGRADO_EN_GRAFO_ACTIVO","residences":[{"path":"manuscrito/sections/integracion_20260812/matricial_52/05_geometria_simplectica_y_lectores.tex","sha256":"087482ebc4a880dba1096a9364dab1a05b0f0bddd3495b6f5b967fb9ccb79ccd"}],"result_id":"P23","statement_sha256":"4ebed4b0f0e733875915bad691b12c053780dfaa50e1f0f4898837314d3edf9a"}
% HMT_RESULT_END:P23:residence
% HMT_RESULT_BEGIN:P24:residence
% {"material_control":"PASS_PUENTE_GENEALOGICO_MATRICIAL","material_state":"INTEGRADO_EN_GRAFO_ACTIVO","residences":[{"path":"manuscrito/sections/integracion_20260812/matricial_52/06_operatorios_doble_proyeccion.tex","sha256":"0f183ed734248c42a405ba89f43a6ac6085d705a188f911c44ed76e8eb4bf745"}],"result_id":"P24","statement_sha256":"6196ff6b2fa346db4f4fd64dc27219eba3734ae04559393e65a5ad14e3c297e0"}
% HMT_RESULT_END:P24:residence
% HMT_RESULT_BEGIN:P25:residence
% {"material_control":"PASS_PUENTE_GENEALOGICO_MATRICIAL","material_state":"INTEGRADO_EN_GRAFO_ACTIVO","residences":[{"path":"manuscrito/sections/integracion_20260812/matricial_52/06_operatorios_doble_proyeccion.tex","sha256":"0f183ed734248c42a405ba89f43a6ac6085d705a188f911c44ed76e8eb4bf745"}],"result_id":"P25","statement_sha256":"339d4b25f9c31d5aa693c11a2377a405ec7ba754affe3143c351475a2d50885d"}
% HMT_RESULT_END:P25:residence
% HMT_RESULT_BEGIN:P26:residence
% {"material_control":"PASS_PUENTE_GENEALOGICO_MATRICIAL","material_state":"INTEGRADO_EN_GRAFO_ACTIVO","residences":[{"path":"manuscrito/sections/integracion_20260812/matricial_52/06_operatorios_doble_proyeccion.tex","sha256":"0f183ed734248c42a405ba89f43a6ac6085d705a188f911c44ed76e8eb4bf745"}],"result_id":"P26","statement_sha256":"1f0955a242b45acf69644be639ccca7ceac39f2d83009c4aafcdb1172802a601"}
% HMT_RESULT_END:P26:residence
% HMT_RESULT_BEGIN:P27:residence
% {"material_control":"PASS_PUENTE_GENEALOGICO_MATRICIAL","material_state":"INTEGRADO_EN_GRAFO_ACTIVO","residences":[{"path":"manuscrito/sections/integracion_20260812/matricial_52/06_operatorios_doble_proyeccion.tex","sha256":"0f183ed734248c42a405ba89f43a6ac6085d705a188f911c44ed76e8eb4bf745"}],"result_id":"P27","statement_sha256":"60ff5f6df5943b80338a53b54928755aa4cae85b90764d63a153d6ad50dfed24"}
% HMT_RESULT_END:P27:residence
% HMT_RESULT_BEGIN:P28:residence
% {"material_control":"PASS_PUENTE_GENEALOGICO_MATRICIAL","material_state":"INTEGRADO_EN_GRAFO_ACTIVO","residences":[{"path":"manuscrito/sections/integracion_20260812/matricial_52/07_continuo_tensores_y_memoria.tex","sha256":"3df0b8a366a6458f7acf2199859a25f5b40379b3f9039ccaa0b959e2f25e36a5"}],"result_id":"P28","statement_sha256":"de45d3b80b483962a79b544ce428c4298718a36362562637d064386d09251da4"}
% HMT_RESULT_END:P28:residence
% HMT_RESULT_BEGIN:P29:residence
% {"material_control":"PASS_PUENTE_GENEALOGICO_MATRICIAL","material_state":"INTEGRADO_EN_GRAFO_ACTIVO","residences":[{"path":"manuscrito/sections/integracion_20260812/matricial_52/07_continuo_tensores_y_memoria.tex","sha256":"3df0b8a366a6458f7acf2199859a25f5b40379b3f9039ccaa0b959e2f25e36a5"}],"result_id":"P29","statement_sha256":"eed86be2eeabf3a9de1c3150b3fb89bdee7fc9a0e51ea1e9052cfab634322586"}
% HMT_RESULT_END:P29:residence
% HMT_RESULT_BEGIN:P30:residence
% {"material_control":"PASS_PUENTE_GENEALOGICO_MATRICIAL","material_state":"INTEGRADO_EN_GRAFO_ACTIVO","residences":[{"path":"manuscrito/sections/integracion_20260812/matricial_52/09_metrologia_y_programa.tex","sha256":"d5cf28b96e5cde3577601d038089754378d0eb0e6118b134ed2b27e98004a1e7"}],"result_id":"P30","statement_sha256":"16dae690fc342fd1d4bbb3ee5c7067af92476cfdcc61842908df31b55a844797"}
% HMT_RESULT_END:P30:residence
% HMT_RESULT_BEGIN:P31:residence
% {"material_control":"PASS_PUENTE_GENEALOGICO_MATRICIAL","material_state":"INTEGRADO_EN_GRAFO_ACTIVO","residences":[{"path":"manuscrito/sections/integracion_20260812/matricial_52/08_sistemas_conicos_y_universalidad.tex","sha256":"2b94acee5d5df46c4c6e2c4b9048962de8961703fc585df45d08a42f9f4ae0d2"}],"result_id":"P31","statement_sha256":"cabf80d0abd369646b4013a8ac2aad89b7f4c6cd7e0b0092d5f4e07d76e74471"}
% HMT_RESULT_END:P31:residence
% HMT_RESULT_BEGIN:P32:residence
% {"material_control":"PASS_PUENTE_GENEALOGICO_MATRICIAL","material_state":"INTEGRADO_EN_GRAFO_ACTIVO","residences":[{"path":"manuscrito/sections/integracion_20260812/matricial_52/08_sistemas_conicos_y_universalidad.tex","sha256":"2b94acee5d5df46c4c6e2c4b9048962de8961703fc585df45d08a42f9f4ae0d2"}],"result_id":"P32","statement_sha256":"d808a7bd125136593c39a42df20427716bf90ad4b1ac4a033227a3c4e370c64a"}
% HMT_RESULT_END:P32:residence
% HMT_RESULT_BEGIN:QPM-001:residence
% {"material_control":"PASS_PUENTE_GENEALOGICO_MATRICIAL","material_state":"INTEGRADO_EN_GRAFO_ACTIVO","residences":[{"path":"manuscrito/sections/integracion_20260812/matricial_52/01_objeto_y_diccionario.tex","sha256":"95d254edbb6dd7184d3a5db5ad83006aee1adbbdd4ee0a5f631e89d6300648da"}],"result_id":"QPM-001","statement_sha256":"876638c74d1df52e66cf8b186c110f2cbb879465ca4aed51526e0ca3b191405f"}
% HMT_RESULT_END:QPM-001:residence
% HMT_RESULT_BEGIN:QLS-001:residence
% {"material_control":"PASS_PUENTE_GENEALOGICO_MATRICIAL","material_state":"INTEGRADO_EN_GRAFO_ACTIVO","residences":[{"path":"manuscrito/sections/integracion_20260812/matricial_52/01_objeto_y_diccionario.tex","sha256":"95d254edbb6dd7184d3a5db5ad83006aee1adbbdd4ee0a5f631e89d6300648da"}],"result_id":"QLS-001","statement_sha256":"60abad49376bb9a2df1df2851bfc068d7fd0e59b90f09c5812f8ac93138d5308"}
% HMT_RESULT_END:QLS-001:residence
% HMT_RESULT_BEGIN:QLSI-001:residence
% {"material_control":"PASS_PUENTE_GENEALOGICO_MATRICIAL","material_state":"INTEGRADO_EN_GRAFO_ACTIVO","residences":[{"path":"manuscrito/sections/integracion_20260812/matricial_52/01_objeto_y_diccionario.tex","sha256":"95d254edbb6dd7184d3a5db5ad83006aee1adbbdd4ee0a5f631e89d6300648da"}],"result_id":"QLSI-001","statement_sha256":"3cc12a3dcfe1865a34b892c489519c5c02e989e8075c94e9782f1e01e67b2c1d"}
% HMT_RESULT_END:QLSI-001:residence
% HMT_RESULT_BEGIN:QMSI-001:residence
% {"material_control":"PASS_PUENTE_GENEALOGICO_MATRICIAL","material_state":"INTEGRADO_EN_GRAFO_ACTIVO","residences":[{"path":"manuscrito/sections/integracion_20260812/matricial_52/01_objeto_y_diccionario.tex","sha256":"95d254edbb6dd7184d3a5db5ad83006aee1adbbdd4ee0a5f631e89d6300648da"}],"result_id":"QMSI-001","statement_sha256":"641653814fe7d90a81405d016aa63d2840bdf02215f16c828afe4ece7989c37b"}
% HMT_RESULT_END:QMSI-001:residence
% HMT_RESULT_BEGIN:QMSI-002:residence
% {"material_control":"PASS_PUENTE_GENEALOGICO_MATRICIAL","material_state":"INTEGRADO_EN_GRAFO_ACTIVO","residences":[{"path":"manuscrito/sections/integracion_20260812/matricial_52/01_objeto_y_diccionario.tex","sha256":"95d254edbb6dd7184d3a5db5ad83006aee1adbbdd4ee0a5f631e89d6300648da"}],"result_id":"QMSI-002","statement_sha256":"79afee550d7e55837c8ed6a9067de90d22cfee54bbff7c98636ded289843016f"}
% HMT_RESULT_END:QMSI-002:residence
% HMT_RESULT_BEGIN:QMSI-003:residence
% {"material_control":"PASS_PUENTE_GENEALOGICO_MATRICIAL","material_state":"INTEGRADO_EN_GRAFO_ACTIVO","residences":[{"path":"manuscrito/sections/integracion_20260812/matricial_52/01_objeto_y_diccionario.tex","sha256":"95d254edbb6dd7184d3a5db5ad83006aee1adbbdd4ee0a5f631e89d6300648da"}],"result_id":"QMSI-003","statement_sha256":"1bbe978d6a9ba0ec29aa916bb4b2a72df5bf3936b89bf14a66b2916e1d4fd090"}
% HMT_RESULT_END:QMSI-003:residence
% HMT_RESULT_BEGIN:QBR-001:residence
% {"material_control":"PASS_PUENTE_GENEALOGICO_MATRICIAL","material_state":"INTEGRADO_EN_GRAFO_ACTIVO","residences":[{"path":"manuscrito/sections/integracion_20260812/matricial_52/02_hmt_en_profundidad.tex","sha256":"fe6c23eac09e147199090a091fe2564e5587b7b1155ea85518d2bf47aade4cfc"},{"path":"manuscrito/sections/md/04ab_postulados_cuanticos_genealogicos_rev2.tex","sha256":"1895f84bd75b104e873a3d73c5b0e8ae1b5df606bd6feefeb4f9666a92966804"}],"result_id":"QBR-001","statement_sha256":"38ad3305cd9c643e9067d133a541df231165278ab3b01e906f494ec45be28104"}
% HMT_RESULT_END:QBR-001:residence
% HMT_RESULT_BEGIN:QDY-001:residence
% {"material_control":"PASS_PUENTE_GENEALOGICO_MATRICIAL","material_state":"INTEGRADO_EN_GRAFO_ACTIVO","residences":[{"path":"manuscrito/sections/integracion_20260812/matricial_52/02_hmt_en_profundidad.tex","sha256":"fe6c23eac09e147199090a091fe2564e5587b7b1155ea85518d2bf47aade4cfc"},{"path":"manuscrito/sections/md/04ab_postulados_cuanticos_genealogicos_rev2.tex","sha256":"1895f84bd75b104e873a3d73c5b0e8ae1b5df606bd6feefeb4f9666a92966804"}],"result_id":"QDY-001","statement_sha256":"bffa8460b3d4625a598b73816f5f9e8a9218458a12b69fd3cc36c8ba18d69a89"}
% HMT_RESULT_END:QDY-001:residence
% HMT_RESULT_BEGIN:QKR-001:residence
% {"material_control":"PASS_PUENTE_GENEALOGICO_MATRICIAL","material_state":"INTEGRADO_EN_GRAFO_ACTIVO","residences":[{"path":"manuscrito/sections/integracion_20260812/matricial_52/02_hmt_en_profundidad.tex","sha256":"fe6c23eac09e147199090a091fe2564e5587b7b1155ea85518d2bf47aade4cfc"},{"path":"manuscrito/sections/md/01a_medicion_cuantica_genealogica_rev2.tex","sha256":"725bbfb625e4bd74c3f5845dc2d3ad117b173fac63f11d300f7f6b8344fb505a"}],"result_id":"QKR-001","statement_sha256":"e458e06f5341341d7cbbeee4dbcc1fa333c7c1d293266aec038128fe9c334858"}
% HMT_RESULT_END:QKR-001:residence
% HMT_RESULT_BEGIN:QMO-001:residence
% {"material_control":"PASS_PUENTE_GENEALOGICO_MATRICIAL","material_state":"INTEGRADO_EN_GRAFO_ACTIVO","residences":[{"path":"manuscrito/sections/integracion_20260812/matricial_52/02_hmt_en_profundidad.tex","sha256":"fe6c23eac09e147199090a091fe2564e5587b7b1155ea85518d2bf47aade4cfc"},{"path":"manuscrito/sections/md/04ab_postulados_cuanticos_genealogicos_rev2.tex","sha256":"1895f84bd75b104e873a3d73c5b0e8ae1b5df606bd6feefeb4f9666a92966804"}],"result_id":"QMO-001","statement_sha256":"14d863637e09c435b127f3b5fc5dccc9f9d873d837809884a8efc4b2da2953be"}
% HMT_RESULT_END:QMO-001:residence
% HMT_RESULT_BEGIN:DIC-001:residence
% {"material_control":"PASS_PUENTE_GENEALOGICO_MATRICIAL","material_state":"INTEGRADO_EN_GRAFO_ACTIVO","residences":[{"path":"manuscrito/sections/integracion_20260812/matricial_52/01_objeto_y_diccionario.tex","sha256":"95d254edbb6dd7184d3a5db5ad83006aee1adbbdd4ee0a5f631e89d6300648da"}],"result_id":"DIC-001","statement_sha256":"51d602fff6bbd77d754ba0ebe0732f3998ac22c81a82514440408ff0f6545d30"}
% HMT_RESULT_END:DIC-001:residence
% HMT_RESULT_BEGIN:QMC-001:residence
% {"material_control":"PASS_PUENTE_GENEALOGICO_MATRICIAL","material_state":"INTEGRADO_EN_GRAFO_ACTIVO","residences":[{"path":"manuscrito/sections/integracion_20260812/matricial_52/04_cuadrados_cuanticos.tex","sha256":"7131342598b20a77a08f6d9d98e12d41ec47e331372ea6c19754df21a58ab585"}],"result_id":"QMC-001","statement_sha256":"900c10fc899cb62a7793541419f3c6ce103296ab9690d816954062e980960504"}
% HMT_RESULT_END:QMC-001:residence
% HMT_RESULT_BEGIN:QDT-001:residence
% {"material_control":"PASS_PUENTE_GENEALOGICO_MATRICIAL","material_state":"INTEGRADO_EN_GRAFO_ACTIVO","residences":[{"path":"manuscrito/sections/integracion_20260812/matricial_52/05_geometria_simplectica_y_lectores.tex","sha256":"087482ebc4a880dba1096a9364dab1a05b0f0bddd3495b6f5b967fb9ccb79ccd"}],"result_id":"QDT-001","statement_sha256":"a42d2bae56ce5410c45794269f3181c68a38ea3fb2d5861cebb013294580287d"}
% HMT_RESULT_END:QDT-001:residence
% HMT_RESULT_BEGIN:ENV-001:residence
% {"material_control":"PASS_PUENTE_GENEALOGICO_MATRICIAL","material_state":"INTEGRADO_EN_GRAFO_ACTIVO","residences":[{"path":"manuscrito/sections/integracion_20260812/matricial_52/03_bigraduacion_ucp.tex","sha256":"c054ec7cb105bc2fda9c3228bfaee808f01db65e38e126a8865a3420ed37864d"}],"result_id":"ENV-001","statement_sha256":"a2c7c6251ce55a5d5d21e2825a8eecad64f3da571d4130d8119999126371726c"}
% HMT_RESULT_END:ENV-001:residence
% HMT_RESULT_BEGIN:RNG-001:residence
% {"material_control":"PASS_PUENTE_GENEALOGICO_MATRICIAL","material_state":"INTEGRADO_EN_GRAFO_ACTIVO","residences":[{"path":"manuscrito/sections/integracion_20260812/matricial_52/06_operatorios_doble_proyeccion.tex","sha256":"0f183ed734248c42a405ba89f43a6ac6085d705a188f911c44ed76e8eb4bf745"}],"result_id":"RNG-001","statement_sha256":"6f564d18bb3910cde9dd86c3bcc20cf5be7d388abb45c8acebdc31553910900c"}
% HMT_RESULT_END:RNG-001:residence
% HMT_RESULT_BEGIN:DYN-001:residence
% {"material_control":"PASS_PUENTE_GENEALOGICO_MATRICIAL","material_state":"INTEGRADO_EN_GRAFO_ACTIVO","residences":[{"path":"manuscrito/sections/integracion_20260812/matricial_52/07_continuo_tensores_y_memoria.tex","sha256":"3df0b8a366a6458f7acf2199859a25f5b40379b3f9039ccaa0b959e2f25e36a5"}],"result_id":"DYN-001","statement_sha256":"931716fef6fe78eab6afb720f6b8621e7cf08908dfd2b425ac48605bb361219a"}
% HMT_RESULT_END:DYN-001:residence
% HMT_RESULT_BEGIN:TEN-001:residence
% {"material_control":"PASS_PUENTE_GENEALOGICO_MATRICIAL","material_state":"INTEGRADO_EN_GRAFO_ACTIVO","residences":[{"path":"manuscrito/sections/integracion_20260812/matricial_52/07_continuo_tensores_y_memoria.tex","sha256":"3df0b8a366a6458f7acf2199859a25f5b40379b3f9039ccaa0b959e2f25e36a5"}],"result_id":"TEN-001","statement_sha256":"458096f924295967751d528c1f5b73f3ae28b5e1cec81ab6ac3036eda48c54fd"}
% HMT_RESULT_END:TEN-001:residence
% HMT_RESULT_BEGIN:CHP-001:residence
% {"material_control":"PASS_PUENTE_GENEALOGICO_MATRICIAL","material_state":"INTEGRADO_EN_GRAFO_ACTIVO","residences":[{"path":"manuscrito/sections/integracion_20260812/matricial_52/08_sistemas_conicos_y_universalidad.tex","sha256":"2b94acee5d5df46c4c6e2c4b9048962de8961703fc585df45d08a42f9f4ae0d2"}],"result_id":"CHP-001","statement_sha256":"8fce51f0c85e0120694a493a3cd73f2edb220a7f02de5bc0d820e4697456ff1f"}
% HMT_RESULT_END:CHP-001:residence
% HMT_RESULT_BEGIN:SDP-001:residence
% {"material_control":"PASS_PUENTE_GENEALOGICO_MATRICIAL","material_state":"INTEGRADO_EN_GRAFO_ACTIVO","residences":[{"path":"manuscrito/sections/integracion_20260812/matricial_52/08_sistemas_conicos_y_universalidad.tex","sha256":"2b94acee5d5df46c4c6e2c4b9048962de8961703fc585df45d08a42f9f4ae0d2"}],"result_id":"SDP-001","statement_sha256":"0fa05be883df7a0e3a59ceb8d9386d4775fc4456de540e0b4fa29e7d7716aa46"}
% HMT_RESULT_END:SDP-001:residence
% HMT_RESULT_BEGIN:EQV-001:residence
% {"material_control":"PASS_PUENTE_GENEALOGICO_MATRICIAL","material_state":"INTEGRADO_EN_GRAFO_ACTIVO","residences":[{"path":"manuscrito/sections/integracion_20260812/matricial_52/05_geometria_simplectica_y_lectores.tex","sha256":"087482ebc4a880dba1096a9364dab1a05b0f0bddd3495b6f5b967fb9ccb79ccd"}],"result_id":"EQV-001","statement_sha256":"0ef0dde22fd08d35555aa0a88a1eb4fa608e3b298286a47559fb223af824f79f"}
% HMT_RESULT_END:EQV-001:residence
% HMT_RESULT_BEGIN:QMT-001:residence
% {"material_control":"PASS_PUENTE_GENEALOGICO_MATRICIAL","material_state":"INTEGRADO_EN_GRAFO_ACTIVO","residences":[{"path":"manuscrito/sections/integracion_20260812/matricial_52/04_cuadrados_cuanticos.tex","sha256":"7131342598b20a77a08f6d9d98e12d41ec47e331372ea6c19754df21a58ab585"}],"result_id":"QMT-001","statement_sha256":"e40aff576b435e19d95904fbf6c160fedf5d8f5bcbe4634eab9591b8c6960630"}
% HMT_RESULT_END:QMT-001:residence
% HMT_RESULT_BEGIN:TRN-001:residence
% {"material_control":"PASS_PUENTE_GENEALOGICO_MATRICIAL","material_state":"INTEGRADO_EN_GRAFO_ACTIVO","residences":[{"path":"manuscrito/sections/integracion_20260812/matricial_52/09_metrologia_y_programa.tex","sha256":"d5cf28b96e5cde3577601d038089754378d0eb0e6118b134ed2b27e98004a1e7"}],"result_id":"TRN-001","statement_sha256":"076c0e63737e5420b981d0d4223a441dd2be4af3e62fe56af2608779d7c60bb0"}
% HMT_RESULT_END:TRN-001:residence
% HMT_RESULT_BEGIN:MPSH-001:residence
% {"material_control":"PASS_PUENTE_GENEALOGICO_MATRICIAL","material_state":"INTEGRADO_EN_GRAFO_ACTIVO","residences":[{"path":"manuscrito/sections/integracion_20260812/matricial_52/07_continuo_tensores_y_memoria.tex","sha256":"3df0b8a366a6458f7acf2199859a25f5b40379b3f9039ccaa0b959e2f25e36a5"}],"result_id":"MPSH-001","statement_sha256":"e8ee9870050a98d998dcfa91b8f581a5626e54e41b516743f9666140829b1f9e"}
% HMT_RESULT_END:MPSH-001:residence
% HMT_RESULT_BEGIN:UNIIMP-001:residence
% {"material_control":"PASS_PUENTE_GENEALOGICO_MATRICIAL","material_state":"INTEGRADO_EN_GRAFO_ACTIVO","residences":[{"path":"manuscrito/sections/integracion_20260812/matricial_52/08_sistemas_conicos_y_universalidad.tex","sha256":"2b94acee5d5df46c4c6e2c4b9048962de8961703fc585df45d08a42f9f4ae0d2"}],"result_id":"UNIIMP-001","statement_sha256":"342e5b5e49f7e76a154e4376612d523325fc967212ee3b2f4156d578ed7818be"}
% HMT_RESULT_END:UNIIMP-001:residence
% HMT_RESULT_BEGIN:PRG-001:residence
% {"material_control":"PASS_PUENTE_GENEALOGICO_MATRICIAL","material_state":"INTEGRADO_EN_GRAFO_ACTIVO","residences":[{"path":"manuscrito/sections/integracion_20260812/matricial_52/09_metrologia_y_programa.tex","sha256":"d5cf28b96e5cde3577601d038089754378d0eb0e6118b134ed2b27e98004a1e7"}],"result_id":"PRG-001","statement_sha256":"42475f38a74d5f5c373f586a656ec2ec049a0f9fd5b8ed896dfb8bd11bf3aec5"}
% HMT_RESULT_END:PRG-001:residence
% HMT_RESULT_BEGIN:QCT-001:residence
% {"material_control":"HASH_REGENERACION+NON_RESULTS_6+SCOPE_LIMITS_5","material_state":"CONTROL_FOCAL_ACTIVO","residences":[{"path":"manuscrito/sections/integracion_20260812/matricial_52/A_pruebas_y_certificados.tex","sha256":"ba316b48d94c537b548e76a95774f361aab946f725468bb6acc0700fbe9b0cdb"},{"path":"pruebas/integracion_20260812/verificar_puente_genealogico_matricial.py","sha256":"4661187f03768de0b87fd454620287e9bb407a9402fa6d7e43e6d2969dd2dc69"},{"path":"certificados/integracion_20260812/puente_genealogico_matricial.json","sha256":"fd45cc4e1e1d12225874e214e27fb648e2cea586a5162ba49afdabb981216683"}],"result_id":"QCT-001","statement_sha256":"6a04ea7ae52a0c3f2571d9787272c10ddc18456acf94ddfef7adef2764393d8c"}
% HMT_RESULT_END:QCT-001:residence
% HMT_RESULT_BEGIN:PRV-001:residence
% {"material_control":"HASH_REGENERACION+NON_RESULTS_6+SCOPE_LIMITS_5","material_state":"CONTROL_FOCAL_ACTIVO","residences":[{"path":"manuscrito/sections/integracion_20260812/matricial_52/C_fuentes_y_bibliografia.tex","sha256":"ba54c19f8729fa9e16033a263c3409a80a6fbb6788dfddcb037a657c14f598ff"},{"path":"pruebas/integracion_20260812/verificar_puente_genealogico_matricial.py","sha256":"4661187f03768de0b87fd454620287e9bb407a9402fa6d7e43e6d2969dd2dc69"},{"path":"certificados/integracion_20260812/puente_genealogico_matricial.json","sha256":"fd45cc4e1e1d12225874e214e27fb648e2cea586a5162ba49afdabb981216683"}],"result_id":"PRV-001","statement_sha256":"3898b9fa7cd9ca405abd8eb49767227624e28b2704db43c64ed70e7a0061bf68"}
% HMT_RESULT_END:PRV-001:residence
% HMT_RESULT_BEGIN:MOR-001:residence
% {"material_control":"GATE_MATERIAL_68_22_22+AUDITORIA_ESTATICA","material_state":"INTEGRADO_O_REMITIDO_SIN_PERDIDA","residences":[{"path":"manuscrito/sections/hmt/14f_cierre_genealogico_numero.tex","sha256":"4522230d8297a1420ccb586d0d6c1fefcb08b5e8e2deeda85b11cc73afd8ded5"}],"result_id":"MOR-001","statement_sha256":"81ec51bdba281285df6bfdc1ea9bb82c7f59300a8a07d463fa438037c958c27a"}
% HMT_RESULT_END:MOR-001:residence
% HMT_RESULT_BEGIN:MOR-002:residence
% {"material_control":"GATE_MATERIAL_68_22_22+AUDITORIA_ESTATICA","material_state":"INTEGRADO_O_REMITIDO_SIN_PERDIDA","residences":[{"path":"manuscrito/sections/hmt/03_app_ortograma.tex","sha256":"aaf0fd1dec2c270d544cd19562b4461c57ebb009437e051d973be43395428390"}],"result_id":"MOR-002","statement_sha256":"4d2f31444c015ef52d258636bae153f77c798e12f034cb624fdf8d1721eda040"}
% HMT_RESULT_END:MOR-002:residence
% HMT_RESULT_BEGIN:MOR-003:residence
% {"material_control":"GATE_MATERIAL_68_22_22+AUDITORIA_ESTATICA","material_state":"INTEGRADO_O_REMITIDO_SIN_PERDIDA","residences":[{"path":"manuscrito/sections/hmt/02_trit_geometrias.tex","sha256":"aefb5c1ec6732534c810d0021834eff92a662f59072b97c35c175a6ce7952fa6"}],"result_id":"MOR-003","statement_sha256":"56c13700ad07f59cb91c4278006ecbd0864a2c64b2bc08c14219693a1d5c1961"}
% HMT_RESULT_END:MOR-003:residence
% HMT_RESULT_BEGIN:MOR-004:residence
% {"material_control":"GATE_MATERIAL_68_22_22+AUDITORIA_ESTATICA","material_state":"INTEGRADO_O_REMITIDO_SIN_PERDIDA","residences":[{"path":"manuscrito/sections/hmt/03k_genealogia_operaciones_rev2.tex","sha256":"fae7aca94c6fb20754431ed79a5aa2f005a7cf2f5681736fc5826da4c694fb98"}],"result_id":"MOR-004","statement_sha256":"e709d75fe6d60011165caa4daee14d1c7b07aa4b57f2395a2ef27267672e3448"}
% HMT_RESULT_END:MOR-004:residence
% HMT_RESULT_BEGIN:MOR-005:residence
% {"material_control":"GATE_MATERIAL_68_22_22+AUDITORIA_ESTATICA","material_state":"INTEGRADO_O_REMITIDO_SIN_PERDIDA","residences":[{"path":"manuscrito/sections/integracion_20260812/realizacion_68/01_cilindros_kronecker_dirac.tex","sha256":"782450c6fc1e432641df8b3cc605425c604109f1db5830873b0b778c59ae6946"}],"result_id":"MOR-005","statement_sha256":"39a6a2f053bee151e082142609deae0baa2067c95df7e2cc57eabe0e0cdb0b35"}
% HMT_RESULT_END:MOR-005:residence
% HMT_RESULT_BEGIN:MOR-006:residence
% {"material_control":"GATE_MATERIAL_68_22_22+AUDITORIA_ESTATICA","material_state":"INTEGRADO_O_REMITIDO_SIN_PERDIDA","residences":[{"path":"manuscrito/sections/integracion_20260812/realizacion_68/01_cilindros_kronecker_dirac.tex","sha256":"782450c6fc1e432641df8b3cc605425c604109f1db5830873b0b778c59ae6946"}],"result_id":"MOR-006","statement_sha256":"969223c89123c8150a4a16dde30b45f11ccb2c9ba9a773533e000cbb936e1d0c"}
% HMT_RESULT_END:MOR-006:residence
% HMT_RESULT_BEGIN:MOR-007:residence
% {"material_control":"GATE_MATERIAL_68_22_22+AUDITORIA_ESTATICA","material_state":"INTEGRADO_O_REMITIDO_SIN_PERDIDA","residences":[{"path":"manuscrito/sections/integracion_20260812/bloque_64_operatorial.tex","sha256":"46e39285e0cee5f56142e4c31f8aece125754649847a8864878db03c39927c0f"}],"result_id":"MOR-007","statement_sha256":"cd14cec597ee88698599f74945cfe83dee11ed410ee4f5f9c55eea289b84c15d"}
% HMT_RESULT_END:MOR-007:residence
% HMT_RESULT_BEGIN:MOR-008:residence
% {"material_control":"GATE_MATERIAL_68_22_22+AUDITORIA_ESTATICA","material_state":"INTEGRADO_O_REMITIDO_SIN_PERDIDA","residences":[{"path":"manuscrito/sections/md/01a_medicion_cuantica_genealogica_rev2.tex","sha256":"725bbfb625e4bd74c3f5845dc2d3ad117b173fac63f11d300f7f6b8344fb505a"},{"path":"manuscrito/sections/md/04ab_postulados_cuanticos_genealogicos_rev2.tex","sha256":"1895f84bd75b104e873a3d73c5b0e8ae1b5df606bd6feefeb4f9666a92966804"}],"result_id":"MOR-008","statement_sha256":"a220ba7c55fda4a44314abe4540421d4c8d4b4e75f7ffa6f6bc273734217fc40"}
% HMT_RESULT_END:MOR-008:residence
% HMT_RESULT_BEGIN:MOR-009:residence
% {"material_control":"GATE_MATERIAL_68_22_22+AUDITORIA_ESTATICA","material_state":"INTEGRADO_O_REMITIDO_SIN_PERDIDA","residences":[{"path":"manuscrito/sections/md/04ab_postulados_cuanticos_genealogicos_rev2.tex","sha256":"1895f84bd75b104e873a3d73c5b0e8ae1b5df606bd6feefeb4f9666a92966804"},{"path":"manuscrito/sections/md/00b_realizacion_espectral_numero_particula.tex","sha256":"741b8542e90d89c9973f78712ccfb0625689f72a5af27ecedcc81d193f38d8f6"}],"result_id":"MOR-009","statement_sha256":"3b85ec28f733e941f275615780570ab9c2a3d81fc638a62c0e6728f5ae1f5925"}
% HMT_RESULT_END:MOR-009:residence
% HMT_RESULT_BEGIN:MOR-010:residence
% {"material_control":"GATE_MATERIAL_68_22_22+AUDITORIA_ESTATICA","material_state":"INTEGRADO_O_REMITIDO_SIN_PERDIDA","residences":[{"path":"manuscrito/sections/integracion_20260812/realizacion_68/02_realizacion_fisica_discreta_corregida.tex","sha256":"219a1a44dc7a0979b3728cebdc4308038cfbef74342bb0fcba50fe754c30292e"}],"result_id":"MOR-010","statement_sha256":"f2c89228291df78e200ef2085a4f4d478ac1be2e6684f9a715f658123de31342"}
% HMT_RESULT_END:MOR-010:residence
% HMT_RESULT_BEGIN:MOR-011:residence
% {"material_control":"GATE_MATERIAL_68_22_22+AUDITORIA_ESTATICA","material_state":"INTEGRADO_O_REMITIDO_SIN_PERDIDA","residences":[{"path":"manuscrito/sections/integracion_20260812/realizacion_68/05_involucion_constitutiva_bidireccional.tex","sha256":"c5cc64e9d5b8014a44b0c30ef8b2885e64b0379a107a073c6026b15381f76a7c"}],"result_id":"MOR-011","statement_sha256":"1ff2ade4f9f1721656ba26a95ce77ff15b5def7ccc20e9c135257ef0a7b97f0b"}
% HMT_RESULT_END:MOR-011:residence
% HMT_RESULT_BEGIN:MOR-012:residence
% {"material_control":"GATE_MATERIAL_68_22_22+AUDITORIA_ESTATICA","material_state":"INTEGRADO_O_REMITIDO_SIN_PERDIDA","residences":[{"path":"manuscrito/sections/integracion_20260812/realizacion_68/05_involucion_constitutiva_bidireccional.tex","sha256":"c5cc64e9d5b8014a44b0c30ef8b2885e64b0379a107a073c6026b15381f76a7c"}],"result_id":"MOR-012","statement_sha256":"c71dfcd2e0330c46692b8833f1bc5d41d7878bfc9db5a82f4cbfc84d7549ab7c"}
% HMT_RESULT_END:MOR-012:residence
% HMT_RESULT_BEGIN:MOR-013:residence
% {"material_control":"GATE_MATERIAL_68_22_22+AUDITORIA_ESTATICA","material_state":"INTEGRADO_O_REMITIDO_SIN_PERDIDA","residences":[{"path":"manuscrito/sections/md/04c_realizaciones_nula_material.tex","sha256":"d4bc1ef3e11a544359a35157497479312fce619dec624b4e01d33b6df16a84d9"}],"result_id":"MOR-013","statement_sha256":"ffe00c7b9b7d0edf7f35a609adfde95e7d3498e2638a94dfd78b707d810b9890"}
% HMT_RESULT_END:MOR-013:residence
% HMT_RESULT_BEGIN:MOR-014:residence
% {"material_control":"GATE_MATERIAL_68_22_22+AUDITORIA_ESTATICA","material_state":"INTEGRADO_O_REMITIDO_SIN_PERDIDA","residences":[{"path":"manuscrito/sections/md/04c_realizaciones_nula_material.tex","sha256":"d4bc1ef3e11a544359a35157497479312fce619dec624b4e01d33b6df16a84d9"},{"path":"manuscrito/sections/md/06ka_clausura_nulo_material_electron_rev3.tex","sha256":"4e83aee8d1ef22ee3bc6e71b10989c62a7a01c4f41153c30efc054909337c201"}],"result_id":"MOR-014","statement_sha256":"892c56b47c56b0f46d3c0955bfbb8bd5ef3f1c07286a797739e61c565e0dc2a2"}
% HMT_RESULT_END:MOR-014:residence
% HMT_RESULT_BEGIN:MOR-015:residence
% {"material_control":"GATE_MATERIAL_68_22_22+AUDITORIA_ESTATICA","material_state":"INTEGRADO_O_REMITIDO_SIN_PERDIDA","residences":[{"path":"manuscrito/sections/integracion_20260812/realizacion_68/03_hopf_villarceau_memoria_5_2.tex","sha256":"5ff5c6ec5855f74338db3440d9a7ddeaa0d93fec54b1ca5bb5094d6a14bd0fff"}],"result_id":"MOR-015","statement_sha256":"49efabbeafb0ec3ca4300a5684772f9b39fbc9836d768664adc1d443c2be762b"}
% HMT_RESULT_END:MOR-015:residence
% HMT_RESULT_BEGIN:MOR-016:residence
% {"material_control":"GATE_MATERIAL_68_22_22+AUDITORIA_ESTATICA","material_state":"INTEGRADO_O_REMITIDO_SIN_PERDIDA","residences":[{"path":"manuscrito/sections/md/04b_banda_particula_virtual.tex","sha256":"d9b8c3db3bc26d37b5a78251a5fbcf67076eaacf35c194b87bf6552d1edfac08"},{"path":"manuscrito/sections/md/21_persistencia_ruta_mobius_familias_actualizada.tex","sha256":"375c7e897d21777f05f9a506bb2429e5de76351b55f10e15f275913d9c306999"}],"result_id":"MOR-016","statement_sha256":"d549af55e36e03afb5ec1ca58bf41f7c6c0c5bd39e17a4a0320aff3b421fec3b"}
% HMT_RESULT_END:MOR-016:residence
% HMT_RESULT_BEGIN:MOR-017:residence
% {"material_control":"GATE_MATERIAL_68_22_22+AUDITORIA_ESTATICA","material_state":"INTEGRADO_O_REMITIDO_SIN_PERDIDA","residences":[{"path":"manuscrito/sections/integracion_20260812/realizacion_68/04_control_cglmp_exacto_68.tex","sha256":"278ced97b864e5c8a06650fca5a102c1a213716af14d3148e8206da54f91fdd7"}],"result_id":"MOR-017","statement_sha256":"8aceaf43834801dfe28d27e802dea64160a4e53dd788c2579d1ea1828dd1aba1"}
% HMT_RESULT_END:MOR-017:residence
% HMT_RESULT_BEGIN:MOR-018:residence
% {"material_control":"GATE_MATERIAL_68_22_22+AUDITORIA_ESTATICA","material_state":"INTEGRADO_O_REMITIDO_SIN_PERDIDA","residences":[{"path":"manuscrito/sections/md/10_constitutivas_si.tex","sha256":"852439fe31f4004a04c10217e89f3414c049bc8803c5b86e3d5b2097c04da0a1"}],"result_id":"MOR-018","statement_sha256":"1c4bd3e069ce1c87d0a969b52dbb80e067cb34adc73c877aa272e4018e841ed0"}
% HMT_RESULT_END:MOR-018:residence
% HMT_RESULT_BEGIN:MOR-019:residence
% {"material_control":"GATE_MATERIAL_68_22_22+AUDITORIA_ESTATICA","material_state":"INTEGRADO_O_REMITIDO_SIN_PERDIDA","residences":[{"path":"manuscrito/sections/md/03b_escala_accion_determinantal.tex","sha256":"73e012418f7e982e277e3e978a0fb03d76d0aebacc8fe085341c70121b20d832"},{"path":"manuscrito/sections/md/04e_termodinamica_rutas.tex","sha256":"ab1429f9145eee89e5287095b94f327291fa4aa247ae0abcc95e3f5611de87c1"}],"result_id":"MOR-019","statement_sha256":"12f6cc8a7f2131682e716fb00c3f2546dbb1f63249994d5d8c1c14d83510c994"}
% HMT_RESULT_END:MOR-019:residence
% HMT_RESULT_BEGIN:MOR-020:residence
% {"material_control":"GATE_MATERIAL_68_22_22+AUDITORIA_ESTATICA","material_state":"INTEGRADO_O_REMITIDO_SIN_PERDIDA","residences":[{"path":"manuscrito/sections/md/05aa_ley_general_masa_accion_autonoma.tex","sha256":"43dbef95faacb5db35585fa447b16406038f05b5296852f259687292bc0313a2"},{"path":"manuscrito/sections/md/05f_operadores_masa_two_way_rev11.tex","sha256":"e0c576ff54f405c41adea7093ca6b117134e9322749c969f1fc4cca6e8615d35"},{"path":"manuscrito/sections/md/06m_tabla_completa_masas_hmt_sm_revfinal.tex","sha256":"1dd8ca6908a9b94af9d5533d57785a2222d455f88e7ac70aa9acf34667ed35b1"}],"result_id":"MOR-020","statement_sha256":"ef9647c7358991a2b5046f1b253d9497443b3f0c31a95b7f6a9906100d45378e"}
% HMT_RESULT_END:MOR-020:residence
% HMT_RESULT_BEGIN:MOR-021:residence
% {"material_control":"GATE_MATERIAL_68_22_22+AUDITORIA_ESTATICA","material_state":"INTEGRADO_O_REMITIDO_SIN_PERDIDA","residences":[{"path":"manuscrito/sections/md/06k_electron_two_way_rev11.tex","sha256":"e9857301411fd54a3f9d9ad99f56466fc2d7aa72b59d7a4a12f14076fc151e0b"},{"path":"manuscrito/sections/md/06ka_clausura_nulo_material_electron_rev3.tex","sha256":"4e83aee8d1ef22ee3bc6e71b10989c62a7a01c4f41153c30efc054909337c201"}],"result_id":"MOR-021","statement_sha256":"79bbb961eec13af6b1fb10be55c0f3a17f871661f753bc37ab8c5e6ce19c0a6a"}
% HMT_RESULT_END:MOR-021:residence
% HMT_RESULT_BEGIN:MOR-022:residence
% {"material_control":"GATE_MATERIAL_68_22_22+AUDITORIA_ESTATICA","material_state":"INTEGRADO_O_REMITIDO_SIN_PERDIDA","residences":[{"path":"manuscrito/sections/md/03d_identidad_autodual_planck_ct108_rev3.tex","sha256":"522a1ce5c537fb26513227fc0766eadefd1b6fd898382137129677aaa69cfaed"}],"result_id":"MOR-022","statement_sha256":"e741dca3aef92e05581285a7ecd9802c495fc2b37445a4604229e67b619c98da"}
% HMT_RESULT_END:MOR-022:residence
% HMT_RESULT_BEGIN:MOR-023:residence
% {"material_control":"GATE_MATERIAL_68_22_22+AUDITORIA_ESTATICA","material_state":"INTEGRADO_O_REMITIDO_SIN_PERDIDA","residences":[{"path":"manuscrito/sections/integracion_20260812/realizacion_68/02_realizacion_fisica_discreta_corregida.tex","sha256":"219a1a44dc7a0979b3728cebdc4308038cfbef74342bb0fcba50fe754c30292e"}],"result_id":"MOR-023","statement_sha256":"bfaa0b7d864dd992d466327275ed93894d90946d897887db756d6d92268920b5"}
% HMT_RESULT_END:MOR-023:residence
% HMT_RESULT_BEGIN:MOR-024:residence
% {"material_control":"GATE_MATERIAL_68_22_22+AUDITORIA_ESTATICA","material_state":"INTEGRADO_O_REMITIDO_SIN_PERDIDA","residences":[{"path":"manuscrito/sections/md/11c_gravedad_holonomia_rev6.tex","sha256":"87d99cb98c92f5e7a771a9ab750690274e370d49a630e461b9d524e02cf6cf71"},{"path":"manuscrito/sections/md/11n_selector_gravitatorio_rev11.tex","sha256":"14b45e691dfd01916515b4a15b910eba1ff698031f794e68ba13d72da9973c9d"}],"result_id":"MOR-024","statement_sha256":"7736b062aa565d44a6b82fd9985d0d4dd7dd7bf8810791d818212081641834d4"}
% HMT_RESULT_END:MOR-024:residence
% HMT_RESULT_BEGIN:MOR-025:residence
% {"material_control":"GATE_MATERIAL_68_22_22+AUDITORIA_ESTATICA","material_state":"INTEGRADO_O_REMITIDO_SIN_PERDIDA","residences":[{"path":"manuscrito/sections/integracion_20260812/realizacion_68/06_programas_mach_cosmologia_68.tex","sha256":"5cca10348cd22a33e86ac0edaaffbd486e433e7e3d80bdf4c97a9368ab868e89"}],"result_id":"MOR-025","statement_sha256":"560fed3a547431a0a0bb145c63be86a48bd71bbf9d20c9ce42e8660b1aa64996"}
% HMT_RESULT_END:MOR-025:residence
% HMT_RESULT_BEGIN:MOR-026:residence
% {"material_control":"GATE_MATERIAL_68_22_22+AUDITORIA_ESTATICA","material_state":"INTEGRADO_O_REMITIDO_SIN_PERDIDA","residences":[{"path":"manuscrito/sections/hmt/14f_cierre_genealogico_numero.tex","sha256":"4522230d8297a1420ccb586d0d6c1fefcb08b5e8e2deeda85b11cc73afd8ded5"},{"path":"manuscrito/sections/hmt/19b_zeta_vacancias_primos_rev11.tex","sha256":"37b1f8fff4663d03826091250843b7a33c8a83e3fd08659c5fa6613a4c346534"}],"result_id":"MOR-026","statement_sha256":"e27ece362c8e14e5d61ce86444197f785ee77f3c6b2a79b2057a895ebb30db26"}
% HMT_RESULT_END:MOR-026:residence
% HMT_RESULT_BEGIN:MOR-027:residence
% {"material_control":"GATE_MATERIAL_68_22_22+AUDITORIA_ESTATICA","material_state":"INTEGRADO_O_REMITIDO_SIN_PERDIDA","residences":[{"path":"manuscrito/sections/hmt/14d_ontologia_numeros_rev11.tex","sha256":"c96b3c94af4668f15c6a061260d1c1900c52b69b8a4a5d2442268f6ff1d0abcc"},{"path":"manuscrito/sections/hmt/14f_cierre_genealogico_numero.tex","sha256":"4522230d8297a1420ccb586d0d6c1fefcb08b5e8e2deeda85b11cc73afd8ded5"}],"result_id":"MOR-027","statement_sha256":"5b813e0066c4fb64f0e94a22e81f2741e55dc4dec333c9aea2d7d02f31523a5d"}
% HMT_RESULT_END:MOR-027:residence
% HMT_RESULT_BEGIN:MOR-028:residence
% {"material_control":"GATE_MATERIAL_68_22_22+AUDITORIA_ESTATICA","material_state":"INTEGRADO_O_REMITIDO_SIN_PERDIDA","residences":[{"path":"manuscrito/sections/hmt/03f_puente_central_app_emrh_rev2.tex","sha256":"f4088242d198ece37ba686e3314432777582aa93e3b7960e39986a5e9a54358f"}],"result_id":"MOR-028","statement_sha256":"3efd7a28586cf53974d753f4d22e4e7b42a73517f5454c21aefb37550f6b1eae"}
% HMT_RESULT_END:MOR-028:residence
% HMT_RESULT_BEGIN:MOR-029:residence
% {"material_control":"GATE_MATERIAL_68_22_22+AUDITORIA_ESTATICA","material_state":"INTEGRADO_O_REMITIDO_SIN_PERDIDA","residences":[{"path":"manuscrito/sections/hmt/06e_certificado_generacion_monodromica_rev7.tex","sha256":"2a0cf8028e12605a97923c53726fcb62c1798c458af8e9f78f12c1129b0e1b94"}],"result_id":"MOR-029","statement_sha256":"91fbdaa9566d579a9f7f0fc50c56eae1187bb6ce719915dc9f11c1031024480f"}
% HMT_RESULT_END:MOR-029:residence
% HMT_RESULT_BEGIN:MOR-030:residence
% {"material_control":"GATE_MATERIAL_68_22_22+AUDITORIA_ESTATICA","material_state":"INTEGRADO_O_REMITIDO_SIN_PERDIDA","residences":[{"path":"manuscrito/sections/integracion_20260812/realizacion_68/08_interfaces_falsadores_realizacion_68.tex","sha256":"323f27caa4ce2d4e3e96795394eb2843c9549f0c7137bde4a8702bf79561e05e"}],"result_id":"MOR-030","statement_sha256":"ae3db0712fc04e6f7197110ed9ac53af8a4e6d1ce41056c0f1849cbe71519bc2"}
% HMT_RESULT_END:MOR-030:residence
% HMT_RESULT_BEGIN:PLS-001:residence
% {"material_control":"GATE_MATERIAL_68_22_22+AUDITORIA_ESTATICA","material_state":"INTEGRADO_O_REMITIDO_SIN_PERDIDA","residences":[{"path":"manuscrito/sections/md/03d_identidad_autodual_planck_ct108_rev3.tex","sha256":"522a1ce5c537fb26513227fc0766eadefd1b6fd898382137129677aaa69cfaed"}],"result_id":"PLS-001","statement_sha256":"dd9b9e630ef17a042a9d2eb03dfb5bfb4f12afef20d488d3eab6c9d3236bc69a"}
% HMT_RESULT_END:PLS-001:residence
% HMT_RESULT_BEGIN:MAD-001:residence
% {"material_control":"GATE_MATERIAL_68_22_22+AUDITORIA_ESTATICA","material_state":"INTEGRADO_O_REMITIDO_SIN_PERDIDA","residences":[{"path":"manuscrito/sections/md/03d_identidad_autodual_planck_ct108_rev3.tex","sha256":"522a1ce5c537fb26513227fc0766eadefd1b6fd898382137129677aaa69cfaed"}],"result_id":"MAD-001","statement_sha256":"5bce3c37e8568b7428605807933327d2bddbd9668451a4eac13bc22f055af271"}
% HMT_RESULT_END:MAD-001:residence
% HMT_RESULT_BEGIN:AGR-001:residence
% {"material_control":"GATE_MATERIAL_68_22_22+AUDITORIA_ESTATICA","material_state":"INTEGRADO_O_REMITIDO_SIN_PERDIDA","residences":[{"path":"manuscrito/sections/md/03d_identidad_autodual_planck_ct108_rev3.tex","sha256":"522a1ce5c537fb26513227fc0766eadefd1b6fd898382137129677aaa69cfaed"}],"result_id":"AGR-001","statement_sha256":"4fee4acf8c84a7dbbaf710eff510d2cb37e76574c593f4912526198dafaa1747"}
% HMT_RESULT_END:AGR-001:residence
% HMT_RESULT_BEGIN:PLE-001:residence
% {"material_control":"GATE_MATERIAL_68_22_22+AUDITORIA_ESTATICA","material_state":"INTEGRADO_O_REMITIDO_SIN_PERDIDA","residences":[{"path":"manuscrito/sections/md/03d_identidad_autodual_planck_ct108_rev3.tex","sha256":"522a1ce5c537fb26513227fc0766eadefd1b6fd898382137129677aaa69cfaed"},{"path":"manuscrito/sections/md/06ka_clausura_nulo_material_electron_rev3.tex","sha256":"4e83aee8d1ef22ee3bc6e71b10989c62a7a01c4f41153c30efc054909337c201"}],"result_id":"PLE-001","statement_sha256":"354993392866e30255c6503239114f91f02cba03ce48f1bd36427071e5939fc8"}
% HMT_RESULT_END:PLE-001:residence
% HMT_RESULT_BEGIN:ANG-001:residence
% {"material_control":"GATE_MATERIAL_68_22_22+AUDITORIA_ESTATICA","material_state":"INTEGRADO_O_REMITIDO_SIN_PERDIDA","residences":[{"path":"manuscrito/sections/integracion_20260812/realizacion_68/03_hopf_villarceau_memoria_5_2.tex","sha256":"5ff5c6ec5855f74338db3440d9a7ddeaa0d93fec54b1ca5bb5094d6a14bd0fff"}],"result_id":"ANG-001","statement_sha256":"f8f0e5ef3962f4ee1a30115f54f35751e6d96da9a4b54f83017c293e67683683"}
% HMT_RESULT_END:ANG-001:residence
% HMT_RESULT_BEGIN:BDR-001:residence
% {"material_control":"GATE_MATERIAL_68_22_22+AUDITORIA_ESTATICA","material_state":"INTEGRADO_O_REMITIDO_SIN_PERDIDA","residences":[{"path":"manuscrito/sections/integracion_20260812/realizacion_68/05_involucion_constitutiva_bidireccional.tex","sha256":"c5cc64e9d5b8014a44b0c30ef8b2885e64b0379a107a073c6026b15381f76a7c"}],"result_id":"BDR-001","statement_sha256":"6a0929cf24c9401330b5cc9ca4baa6fbfe43f7d9efa70a75883aa7f5168a4682"}
% HMT_RESULT_END:BDR-001:residence
% HMT_RESULT_BEGIN:OCC-001:residence
% {"material_control":"GATE_MATERIAL_68_22_22+AUDITORIA_ESTATICA","material_state":"INTEGRADO_O_REMITIDO_SIN_PERDIDA","residences":[{"path":"manuscrito/sections/md/04e_termodinamica_rutas.tex","sha256":"ab1429f9145eee89e5287095b94f327291fa4aa247ae0abcc95e3f5611de87c1"}],"result_id":"OCC-001","statement_sha256":"ee25ce148dd36c42aeaace51d97662de4acb35aa1ed8f59a27a40f937515c612"}
% HMT_RESULT_END:OCC-001:residence
% HMT_RESULT_BEGIN:BOL-001:residence
% {"material_control":"GATE_MATERIAL_68_22_22+AUDITORIA_ESTATICA","material_state":"INTEGRADO_O_REMITIDO_SIN_PERDIDA","residences":[{"path":"manuscrito/sections/md/04e_termodinamica_rutas.tex","sha256":"ab1429f9145eee89e5287095b94f327291fa4aa247ae0abcc95e3f5611de87c1"}],"result_id":"BOL-001","statement_sha256":"b336e9229e7eae82bc749151dce330d09f3ec616c75b361fac0728e7c87f6c18"}
% HMT_RESULT_END:BOL-001:residence
% HMT_RESULT_BEGIN:COS-001:residence
% {"material_control":"GATE_MATERIAL_68_22_22+AUDITORIA_ESTATICA","material_state":"INTEGRADO_O_REMITIDO_SIN_PERDIDA","residences":[{"path":"manuscrito/sections/integracion_20260812/realizacion_68/06_programas_mach_cosmologia_68.tex","sha256":"5cca10348cd22a33e86ac0edaaffbd486e433e7e3d80bdf4c97a9368ab868e89"}],"result_id":"COS-001","statement_sha256":"bd40c86eaf9fa174e90b2113b1e1001dc7db1a73fb69eb230c228131867c4969"}
% HMT_RESULT_END:COS-001:residence
% HMT_RESULT_BEGIN:GWV-001:residence
% {"material_control":"GATE_MATERIAL_68_22_22+AUDITORIA_ESTATICA","material_state":"INTEGRADO_O_REMITIDO_SIN_PERDIDA","residences":[{"path":"manuscrito/sections/integracion_20260812/realizacion_68/06_programas_mach_cosmologia_68.tex","sha256":"5cca10348cd22a33e86ac0edaaffbd486e433e7e3d80bdf4c97a9368ab868e89"}],"result_id":"GWV-001","statement_sha256":"3cf172b69a92d50f0b499502ca1a16ea845a1c3adb4ce9e5c07cfda16f9c38fa"}
% HMT_RESULT_END:GWV-001:residence
% HMT_RESULT_BEGIN:DRK-001:residence
% {"material_control":"GATE_MATERIAL_68_22_22+AUDITORIA_ESTATICA","material_state":"INTEGRADO_O_REMITIDO_SIN_PERDIDA","residences":[{"path":"manuscrito/sections/integracion_20260812/realizacion_68/06_programas_mach_cosmologia_68.tex","sha256":"5cca10348cd22a33e86ac0edaaffbd486e433e7e3d80bdf4c97a9368ab868e89"}],"result_id":"DRK-001","statement_sha256":"cfbb50fa8ff9e04eb1292280e0f1cac17fb57e837209733564e0026a2be31639"}
% HMT_RESULT_END:DRK-001:residence
% HMT_RESULT_BEGIN:CMK-001:residence
% {"material_control":"GATE_MATERIAL_68_22_22+AUDITORIA_ESTATICA","material_state":"INTEGRADO_O_REMITIDO_SIN_PERDIDA","residences":[{"path":"manuscrito/sections/integracion_20260812/realizacion_68/07_cmk_coestructura_dinamica.tex","sha256":"66c38ef3af1739a961696c1d7c23c719df6530be5afced7b53e2c3de4b79d1fd"}],"result_id":"CMK-001","statement_sha256":"d549072e45a1534a427abeaf657bff1597774b49ed6b4c93ac49b3a4a6c836a6"}
% HMT_RESULT_END:CMK-001:residence
% HMT_RESULT_BEGIN:MRC-001:residence
% {"material_control":"GATE_MATERIAL_68_22_22+AUDITORIA_ESTATICA","material_state":"INTEGRADO_O_REMITIDO_SIN_PERDIDA","residences":[{"path":"manuscrito/sections/integracion_20260812/realizacion_68/09_certificado_y_disposicion_68.tex","sha256":"87dc8390235dd0f1c606ca180c5b56639642fa4818d7400f4ce633c8df913d0c"}],"result_id":"MRC-001","statement_sha256":"80a7a36147cfc1163735bc538f698ff947ea95661565393b3489da4e24a85dd7"}
% HMT_RESULT_END:MRC-001:residence
% HMT_RESULT_BEGIN:PRG68-001:residence
% {"material_control":"GATE_MATERIAL_68_22_22+AUDITORIA_ESTATICA","material_state":"INTEGRADO_O_REMITIDO_SIN_PERDIDA","residences":[{"path":"manuscrito/sections/integracion_20260812/realizacion_68/08_interfaces_falsadores_realizacion_68.tex","sha256":"323f27caa4ce2d4e3e96795394eb2843c9549f0c7137bde4a8702bf79561e05e"}],"result_id":"PRG68-001","statement_sha256":"cf2a64c0deb4ed89400145c71e75b98276573e3fd427b0a93e0ecf04d4114f26"}
% HMT_RESULT_END:PRG68-001:residence
% HMT_RESULT_BEGIN:GTR-001:residence
% {"material_control":"GATE_MATERIAL_68_22_22+AUDITORIA_ESTATICA","material_state":"INTEGRADO_O_REMITIDO_SIN_PERDIDA","residences":[{"path":"manuscrito/sections/md/11c_gravedad_holonomia_rev6.tex","sha256":"87d99cb98c92f5e7a771a9ab750690274e370d49a630e461b9d524e02cf6cf71"},{"path":"manuscrito/sections/md/11n_selector_gravitatorio_rev11.tex","sha256":"14b45e691dfd01916515b4a15b910eba1ff698031f794e68ba13d72da9973c9d"}],"result_id":"GTR-001","statement_sha256":"ddfba4182cba6421426c961c29f8e8fc1053105316f55429bf5db6218eae8a25"}
% HMT_RESULT_END:GTR-001:residence
% HMT_RESULT_BEGIN:AMB-001:residence
% {"material_control":"PASS_AMBIVALENCE_AREAL_VOLUMETRIC+GATE_MATERIAL_V5_219","material_state":"CONSERVADO_Y_ATOMIZADO_EN_DELTA_SUCESOR","residences":[{"path":"manuscrito/sections/hmt/05b_extensiones_dinamica_cociente.tex","sha256":"2c562fa974ec9637a3fe09b0594b57cf0652f1f6813e6eefcae2a14929248882"}],"result_id":"AMB-001","statement_sha256":"a33b07c534544852bd3e4e64ece4fe0c7c8af500d5eb8b8addbb97736170835e"}
% HMT_RESULT_END:AMB-001:residence
 % Sidecar operativo sucesor: conserva los 219 resultados atómicos de REV.4 y
% añade los 29 resultados CHI, separado del estado material testigo.
% Sidecar atómico de la edición integral reforzada: 261 IDs estables.
% PREPARED: no es recibo y no autoriza por sí solo la compilación.
% HMT_RESULT_BEGIN:R3-001:residence
% HMT_ATOMIC_RESULT {"material_state":"CONSERVADO_POR_BASE_GRANULAR_REV3","proof_strength":"EXACTO_INTERNO","provenance":"RESULTADO_RECUPERADO","reinforcement_state":"DEPLOYED_WITHOUT_NEW_RESULT_ID","residence":"portada, resumen, introducción, síntesis","result_id":"R3-001","statement":"Título y subtítulo como tesis: cierre holográfico, estructura discreta del continuo y origen de constantes","status":"PRESENT_IN_SUCCESSOR_SOURCE"}
% HMT_RESULT_END:R3-001:residence
% HMT_RESULT_BEGIN:R3-002:residence
% HMT_ATOMIC_RESULT {"material_state":"CONSERVADO_POR_BASE_GRANULAR_REV3","proof_strength":"EXACTO_INTERNO","provenance":"RESULTADO_RECUPERADO","reinforcement_state":"DEPLOYED_WITHOUT_NEW_RESULT_ID","residence":"`hmt/00_construccion_app_gnomon_rev11`, `03_app_ortograma`, `03f`","result_id":"R3-002","statement":"APP completo: suma, producto, residuo, cociente, gnomón, fibras, acarreo y bloque central","status":"PRESENT_IN_SUCCESSOR_SOURCE"}
% HMT_RESULT_END:R3-002:residence
% HMT_RESULT_BEGIN:R3-003:residence
% HMT_ATOMIC_RESULT {"material_state":"CONSERVADO_POR_BASE_GRANULAR_REV3","proof_strength":"EXACTO_INTERNO","provenance":"RESULTADO_RECUPERADO","reinforcement_state":"DEPLOYED_WITHOUT_NEW_RESULT_ID","residence":"`hmt/02_trit_geometrias`, `03g`","result_id":"R3-003","statement":"TRIT orientado y temporal, con presente torsional y memoria","status":"PRESENT_IN_SUCCESSOR_SOURCE"}
% HMT_RESULT_END:R3-003:residence
% HMT_RESULT_BEGIN:R3-004:residence
% HMT_ATOMIC_RESULT {"material_state":"CONSERVADO_POR_BASE_GRANULAR_REV3","proof_strength":"EXACTO_INTERNO","provenance":"RESULTADO_RECUPERADO","reinforcement_state":"DEPLOYED_WITHOUT_NEW_RESULT_ID","residence":"/Users/ruben/Documents/HMT2/EDICION_AUTONOMA_HMT_MD_2026-08-09_REV_3_EN_TRABAJO/01_FUENTE_REVISADA/manuscrito/sections/hmt/03h_arbol_operatorio_tpk_rev11.tex; /Users/ruben/Documents/HMT2/EDICION_AUTONOMA_HMT_MD_2026-08-09_REV_3_EN_TRABAJO/01_FUENTE_REVISADA/manuscrito/sections/hmt/03i_inventario_operatorio_tpk_rev11.tex; /Users/ruben/Documents/HMT2/EDICION_AUTONOMA_HMT_MD_2026-08-09_REV_3_EN_TRABAJO/01_FUENTE_REVISADA/manuscrito/sections/hmt/03k_genealogia_operaciones_rev2.tex; /Users/ruben/Documents/HMT2/EDICION_AUTONOMA_HMT_MD_2026-08-09_REV_3_EN_TRABAJO/01_FUENTE_REVISADA/manuscrito/sections/hmt/03l_elevacion_cohomologica_dimension_rev2.tex; /Users/ruben/Documents/HMT2/EDICION_AUTONOMA_HMT_MD_2026-08-09_REV_3_EN_TRABAJO/01_FUENTE_REVISADA/manuscrito/sections/hmt/20_arquitectura_operatoria_tpk_actualizada.tex","result_id":"R3-004","statement":"TPK-full: objetos, rutas, memoria, frontera, prolongación, operadores y cálculo genealógico","status":"PRESENT_IN_SUCCESSOR_SOURCE"}
% HMT_RESULT_END:R3-004:residence
% HMT_RESULT_BEGIN:R3-005:residence
% HMT_ATOMIC_RESULT {"material_state":"CONSERVADO_POR_BASE_GRANULAR_REV3","proof_strength":"EXACTO_INTERNO","provenance":"RESULTADO_RECUPERADO","reinforcement_state":"DEPLOYED_WITHOUT_NEW_RESULT_ID","residence":"`hmt/00a`, `00b`, `14f`","result_id":"R3-005","statement":"Cuenta→medida, continuo como secciones compatibles y recta como evaluación","status":"PRESENT_IN_SUCCESSOR_SOURCE"}
% HMT_RESULT_END:R3-005:residence
% HMT_RESULT_BEGIN:R3-006:residence
% HMT_ATOMIC_RESULT {"material_state":"CONSERVADO_POR_BASE_GRANULAR_REV3","proof_strength":"EXACTO_INTERNO","provenance":"RESULTADO_RECUPERADO","reinforcement_state":"DEPLOYED_WITHOUT_NEW_RESULT_ID","residence":"`hmt/06aa`, `06b--06f`","result_id":"R3-006","statement":"Monodromía \\(9\\to1\\) con memoria, generación coinductiva de \\(\\pi,e,\\varphi\\), 396/729/43/387 y cinco regiones de \\(\\pi\\)","status":"PRESENT_IN_SUCCESSOR_SOURCE"}
% HMT_RESULT_END:R3-006:residence
% HMT_RESULT_BEGIN:R3-007:residence
% HMT_ATOMIC_RESULT {"material_state":"CONSERVADO_POR_BASE_GRANULAR_REV3","proof_strength":"EXACTO_INTERNO","provenance":"RESULTADO_RECUPERADO","reinforcement_state":"DEPLOYED_WITHOUT_NEW_RESULT_ID","residence":"capítulos 8, Planck, coordenadas conjugadas y Barbero","result_id":"R3-007","statement":"\\(\\alpha_{\\rm HMT}\\), Planck, Barbero y constantes en su orden causal","status":"PRESENT_IN_SUCCESSOR_SOURCE"}
% HMT_RESULT_END:R3-007:residence
% HMT_RESULT_BEGIN:R3-008:residence
% HMT_ATOMIC_RESULT {"material_state":"CONSERVADO_POR_BASE_GRANULAR_REV3","proof_strength":"EXACTO_INTERNO","provenance":"RESULTADO_RECUPERADO","reinforcement_state":"DEPLOYED_WITHOUT_NEW_RESULT_ID","residence":"`hmt/09`, `09g`, `09e`","result_id":"R3-008","statement":"Doble proyección desde el mismo estado holonómico integral","status":"PRESENT_IN_SUCCESSOR_SOURCE"}
% HMT_RESULT_END:R3-008:residence
% HMT_RESULT_BEGIN:R3-009:residence
% HMT_ATOMIC_RESULT {"material_state":"CONSERVADO_POR_BASE_GRANULAR_REV3","proof_strength":"EXACTO_INTERNO","provenance":"RESULTADO_RECUPERADO","reinforcement_state":"DEPLOYED_WITHOUT_NEW_RESULT_ID","residence":"`hmt/09c`, `15`, `15c`, `15d`","result_id":"R3-009","statement":"Corredor Paley--Witt--Golay--Leech--\\(V^\\natural\\)--Monstruo--Moonshine","status":"PRESENT_IN_SUCCESSOR_SOURCE"}
% HMT_RESULT_END:R3-009:residence
% HMT_RESULT_BEGIN:R3-010:residence
% HMT_ATOMIC_RESULT {"material_state":"CONSERVADO_POR_BASE_GRANULAR_REV3","proof_strength":"EXACTO_INTERNO","provenance":"RESULTADO_RECUPERADO","reinforcement_state":"DEPLOYED_WITHOUT_NEW_RESULT_ID","residence":"`md/00b_realizacion_espectral_numero_particula`","result_id":"R3-010","statement":"Número→partícula: sección, ruta, espectro, persistencia y frontera MD","status":"PRESENT_IN_SUCCESSOR_SOURCE"}
% HMT_RESULT_END:R3-010:residence
% HMT_RESULT_BEGIN:R3-011:residence
% HMT_ATOMIC_RESULT {"material_state":"CONSERVADO_POR_BASE_GRANULAR_REV3","proof_strength":"EXACTO_INTERNO","provenance":"RESULTADO_RECUPERADO","reinforcement_state":"DEPLOYED_WITHOUT_NEW_RESULT_ID","residence":"`md/04ab_postulados_cuanticos_genealogicos_rev2`","result_id":"R3-011","statement":"Postulados cuánticos HMT: Hilbert, rayos/densidades, observables, PVM, Born, Stone--Schrödinger y tensor","status":"PRESENT_IN_SUCCESSOR_SOURCE"}
% HMT_RESULT_END:R3-011:residence
% HMT_RESULT_BEGIN:R3-012:residence
% HMT_ATOMIC_RESULT {"material_state":"CONSERVADO_POR_BASE_GRANULAR_REV3","proof_strength":"EXACTO_INTERNO","provenance":"RESULTADO_RECUPERADO","reinforcement_state":"DEPLOYED_WITHOUT_NEW_RESULT_ID","residence":"`md/01a_medicion_cuantica_genealogica_rev2`","result_id":"R3-012","statement":"Medición: fibra, Lüders, Kraus, CPTP, Stinespring y PVM común","status":"PRESENT_IN_SUCCESSOR_SOURCE"}
% HMT_RESULT_END:R3-012:residence
% HMT_RESULT_BEGIN:R3-013:residence
% HMT_ATOMIC_RESULT {"material_state":"CONSERVADO_POR_BASE_GRANULAR_REV3","proof_strength":"EXACTO_INTERNO","provenance":"RESULTADO_RECUPERADO","reinforcement_state":"DEPLOYED_WITHOUT_NEW_RESULT_ID","residence":"`md/04aa`, capítulos cuánticos","result_id":"R3-013","statement":"Caminos, operador unitario \\(27\\times27\\), límite cilíndrico y reconstrucción cuántica","status":"PRESENT_IN_SUCCESSOR_SOURCE"}
% HMT_RESULT_END:R3-013:residence
% HMT_RESULT_BEGIN:R3-014:residence
% HMT_ATOMIC_RESULT {"material_state":"CONSERVADO_POR_BASE_GRANULAR_REV3","proof_strength":"EXACTO_INTERNO","provenance":"RESULTADO_RECUPERADO","reinforcement_state":"DEPLOYED_WITHOUT_NEW_RESULT_ID","residence":"`md/04d`, `04e`, `04g`, `04h`","result_id":"R3-014","statement":"Landauer finito/informacional y estadística cuántica tipada","status":"PRESENT_IN_SUCCESSOR_SOURCE"}
% HMT_RESULT_END:R3-014:residence
% HMT_RESULT_BEGIN:R3-015:residence
% HMT_ATOMIC_RESULT {"material_state":"CONSERVADO_POR_BASE_GRANULAR_REV3","proof_strength":"EXACTO_INTERNO","provenance":"RESULTADO_RECUPERADO","reinforcement_state":"DEPLOYED_WITHOUT_NEW_RESULT_ID","residence":"`md/03d`, `06k`, `06ka`","result_id":"R3-015","statement":"Planck--CT108--Compton, recta autodual y clausura electrónica de rango uno","status":"PRESENT_IN_SUCCESSOR_SOURCE"}
% HMT_RESULT_END:R3-015:residence
% HMT_RESULT_BEGIN:R3-016:residence
% HMT_ATOMIC_RESULT {"material_state":"CONSERVADO_POR_BASE_GRANULAR_REV3","proof_strength":"EXACTO_INTERNO","provenance":"RESULTADO_RECUPERADO","reinforcement_state":"DEPLOYED_WITHOUT_NEW_RESULT_ID","residence":"`md/05aa`, `05f`, `06m`, `06i`","result_id":"R3-016","statement":"Ley operatorial de masas: \\(10^{-34}\\), \\(K_D/K_\\Omega\\), 81/56/13/324, 23+23 y 17 condicionadas","status":"PRESENT_IN_SUCCESSOR_SOURCE"}
% HMT_RESULT_END:R3-016:residence
% HMT_RESULT_BEGIN:R3-017:residence
% HMT_ATOMIC_RESULT {"material_state":"CONSERVADO_POR_BASE_GRANULAR_REV3","proof_strength":"EXACTO_INTERNO","provenance":"RESULTADO_RECUPERADO","reinforcement_state":"DEPLOYED_WITHOUT_NEW_RESULT_ID","residence":"`md/10c`, `11`, `11c`, `11m`, `11p`","result_id":"R3-017","statement":"Cartan--Holst, gravitación, autoinercia y sus fronteras declaradas","status":"PRESENT_IN_SUCCESSOR_SOURCE"}
% HMT_RESULT_END:R3-017:residence
% HMT_RESULT_BEGIN:R3-018:residence
% HMT_ATOMIC_RESULT {"material_state":"CONSERVADO_POR_BASE_GRANULAR_REV3","proof_strength":"EXACTO_INTERNO","provenance":"RESULTADO_RECUPERADO","reinforcement_state":"DEPLOYED_WITHOUT_NEW_RESULT_ID","residence":"`hmt/17_sintesis`, `md/12_sintesis`","result_id":"R3-018","statement":"Síntesis causal APP→TRIT→TPK→MD, incluyendo defectos 4/2/6/54","status":"PRESENT_IN_SUCCESSOR_SOURCE"}
% HMT_RESULT_END:R3-018:residence
% HMT_RESULT_BEGIN:R3-019:residence
% HMT_ATOMIC_RESULT {"material_state":"CONSERVADO_POR_BASE_GRANULAR_REV3","proof_strength":"EXACTO_INTERNO","provenance":"RESULTADO_RECUPERADO","reinforcement_state":"DEPLOYED_WITHOUT_NEW_RESULT_ID","residence":"activos de REV.3","result_id":"R3-019","statement":"Portada, referencias, tablas y QA ya aptos","status":"PRESENT_IN_SUCCESSOR_SOURCE"}
% HMT_RESULT_END:R3-019:residence
% HMT_RESULT_BEGIN:R3-020:residence
% HMT_ATOMIC_RESULT {"material_state":"CONSERVADO_POR_BASE_GRANULAR_REV3","proof_strength":"EXACTO_INTERNO","provenance":"RESULTADO_RECUPERADO","reinforcement_state":"DEPLOYED_WITHOUT_NEW_RESULT_ID","residence":"grafo completo","result_id":"R3-020","statement":"Todo archivo y toda afirmación de REV.3 no afectados por el delta","status":"PRESENT_IN_SUCCESSOR_SOURCE"}
% HMT_RESULT_END:R3-020:residence
% HMT_RESULT_BEGIN:C-001:residence
% HMT_ATOMIC_RESULT {"material_state":"CONTROL_ACTIVO_EN_PLAN_Y_FUENTE","proof_strength":"EXACTO_INTERNO","provenance":"FORMALIZACION_NUEVA","reinforcement_state":"DEPLOYED_WITHOUT_NEW_RESULT_ID","residence":"gobierno transversal","result_id":"C-001","statement":"La monografía de 68 pp abrevia KMS como \\(F_{A,B}(t+i\\beta)=F_{B,A}(t)\\) bajo una definición que no lo permite","status":"CONTROLLED"}
% HMT_RESULT_END:C-001:residence
% HMT_RESULT_BEGIN:C-002:residence
% HMT_ATOMIC_RESULT {"material_state":"CONTROL_ACTIVO_EN_PLAN_Y_FUENTE","proof_strength":"EXACTO_INTERNO","provenance":"FORMALIZACION_NUEVA","reinforcement_state":"DEPLOYED_WITHOUT_NEW_RESULT_ID","residence":"gobierno transversal","result_id":"C-002","statement":"MOR-010 enumera siete piezas, mientras teorema y certificado tienen cuatro componentes","status":"CONTROLLED"}
% HMT_RESULT_END:C-002:residence
% HMT_RESULT_BEGIN:C-003:residence
% HMT_ATOMIC_RESULT {"material_state":"CONTROL_ACTIVO_EN_PLAN_Y_FUENTE","proof_strength":"EXACTO_INTERNO","provenance":"FORMALIZACION_NUEVA","reinforcement_state":"DEPLOYED_WITHOUT_NEW_RESULT_ID","residence":"gobierno transversal","result_id":"C-003","statement":"El verificador de 68 pp declara parte de la funtorialidad mediante `True`","status":"CONTROLLED"}
% HMT_RESULT_END:C-003:residence
% HMT_RESULT_BEGIN:C-004:residence
% HMT_ATOMIC_RESULT {"material_state":"CONTROL_ACTIVO_EN_PLAN_Y_FUENTE","proof_strength":"EXACTO_INTERNO","provenance":"FORMALIZACION_NUEVA","reinforcement_state":"DEPLOYED_WITHOUT_NEW_RESULT_ID","residence":"gobierno transversal","result_id":"C-004","statement":"QMS no conmutativos del puente son semiclásicos","status":"CONTROLLED"}
% HMT_RESULT_END:C-004:residence
% HMT_RESULT_BEGIN:C-005:residence
% HMT_ATOMIC_RESULT {"material_state":"CONTROL_ACTIVO_EN_PLAN_Y_FUENTE","proof_strength":"EXACTO_INTERNO","provenance":"FORMALIZACION_NUEVA","reinforcement_state":"DEPLOYED_WITHOUT_NEW_RESULT_ID","residence":"gobierno transversal","result_id":"C-005","statement":"El límite inverso HMT y el continuo MPS son objetos diferentes","status":"CONTROLLED"}
% HMT_RESULT_END:C-005:residence
% HMT_RESULT_BEGIN:C-006:residence
% HMT_ATOMIC_RESULT {"material_state":"CONTROL_ACTIVO_EN_PLAN_Y_FUENTE","proof_strength":"EXACTO_INTERNO","provenance":"FORMALIZACION_NUEVA","reinforcement_state":"DEPLOYED_WITHOUT_NEW_RESULT_ID","residence":"gobierno transversal","result_id":"C-006","statement":"`06_teorema_morfismo.tex` presupone un dominio común y hace preservar inversas a magnitudes positivas; además sugerir que la realización positiva se recupera de la orientada pierde variación total","status":"CONTROLLED"}
% HMT_RESULT_END:C-006:residence
% HMT_RESULT_BEGIN:C-007:residence
% HMT_ATOMIC_RESULT {"material_state":"CONTROL_ACTIVO_EN_PLAN_Y_FUENTE","proof_strength":"EXACTO_INTERNO","provenance":"FORMALIZACION_NUEVA","reinforcement_state":"DEPLOYED_WITHOUT_NEW_RESULT_ID","residence":"gobierno transversal","result_id":"C-007","statement":"El apéndice del puente matricial afirma que sus seis no-certificaciones se incluyen en el JSON, pero `scope_limits` contiene cinco controles distintos","status":"CONTROLLED"}
% HMT_RESULT_END:C-007:residence
% HMT_RESULT_BEGIN:MAS-RECTOR-001:residence
% HMT_ATOMIC_RESULT {"material_state":"CONSERVADO_CON_APTO_FOCAL","proof_strength":"EXACTO_INTERNO","provenance":"RESULTADO_RECUPERADO","reinforcement_state":"DEPLOYED_WITHOUT_NEW_RESULT_ID","residence":"/Users/ruben/Documents/HMT2/EDICION_AUTONOMA_HMT_MD_2026-08-09_REV_3_EN_TRABAJO/01_FUENTE_REVISADA/manuscrito/sections/md/06h_atlas_familias_rev6.tex; /Users/ruben/Documents/HMT2/EDICION_AUTONOMA_HMT_MD_2026-08-09_REV_3_EN_TRABAJO/01_FUENTE_REVISADA/manuscrito/sections/md/06i_inventario_humano_23_sectores_rev8.tex","result_id":"MAS-RECTOR-001","statement":"81 celdas; 56 familias; 13 multisecciones; 324 rutas orientadas; sector nulo separado","status":"PRESENT_IN_SUCCESSOR_SOURCE"}
% HMT_RESULT_END:MAS-RECTOR-001:residence
% HMT_RESULT_BEGIN:MAS-RECTOR-002:residence
% HMT_ATOMIC_RESULT {"material_state":"CONSERVADO_CON_APTO_FOCAL","proof_strength":"EXACTO_INTERNO","provenance":"ARQUITECTURA_AUTORAL_PREEXISTENTE","reinforcement_state":"DEPLOYED_WITHOUT_NEW_RESULT_ID","residence":"/Users/ruben/Documents/HMT2/EDICION_AUTONOMA_HMT_MD_2026-08-09_REV_3_EN_TRABAJO/01_FUENTE_REVISADA/manuscrito/sections/md/00b_realizacion_espectral_numero_particula.tex; /Users/ruben/Documents/HMT2/EDICION_AUTONOMA_HMT_MD_2026-08-09_REV_3_EN_TRABAJO/01_FUENTE_REVISADA/manuscrito/sections/md/05aa_ley_general_masa_accion_autonoma.tex","result_id":"MAS-RECTOR-002","statement":"extensión superviviente→ruta→ledger→firma Z5→carácter→torres→recta de masa","status":"PRESENT_IN_SUCCESSOR_SOURCE"}
% HMT_RESULT_END:MAS-RECTOR-002:residence
% HMT_RESULT_BEGIN:MAS-RECTOR-003:residence
% HMT_ATOMIC_RESULT {"material_state":"CONSERVADO_CON_APTO_FOCAL","proof_strength":"EXACTO_INTERNO","provenance":"RESULTADO_RECUPERADO","reinforcement_state":"DEPLOYED_WITHOUT_NEW_RESULT_ID","residence":"/Users/ruben/Documents/HMT2/EDICION_AUTONOMA_HMT_MD_2026-08-09_REV_3_EN_TRABAJO/01_FUENTE_REVISADA/manuscrito/sections/hmt/03_app_ortograma.tex","result_id":"MAS-RECTOR-003","statement":"Bc=9P_++5P_-; salto primordial=4; defecto global posterior 2→6→54","status":"PRESENT_IN_SUCCESSOR_SOURCE"}
% HMT_RESULT_END:MAS-RECTOR-003:residence
% HMT_RESULT_BEGIN:MAS-RECTOR-004:residence
% HMT_ATOMIC_RESULT {"material_state":"CONSERVADO_CON_APTO_FOCAL","proof_strength":"EXACTO_INTERNO","provenance":"RESULTADO_RECUPERADO","reinforcement_state":"DEPLOYED_WITHOUT_NEW_RESULT_ID","residence":"/Users/ruben/Documents/HMT2/EDICION_AUTONOMA_HMT_MD_2026-08-09_REV_3_EN_TRABAJO/01_FUENTE_REVISADA/manuscrito/sections/md/03b_escala_accion_determinantal.tex","result_id":"MAS-RECTOR-004","statement":"SNF(H4)=(1,2,2,4); |coker|=16; Sigma_H=phi alpha^16; ord10=-34","status":"PRESENT_IN_SUCCESSOR_SOURCE"}
% HMT_RESULT_END:MAS-RECTOR-004:residence
% HMT_RESULT_BEGIN:MAS-RECTOR-005:residence
% HMT_ATOMIC_RESULT {"material_state":"CONSERVADO_CON_APTO_FOCAL","proof_strength":"DEMOSTRADO_CON_ESTRUCTURA_DE_PARTIDA_EXPLICITA","provenance":"RESULTADO_RECUPERADO","reinforcement_state":"DEPLOYED_WITHOUT_NEW_RESULT_ID","residence":"/Users/ruben/Documents/HMT2/EDICION_AUTONOMA_HMT_MD_2026-08-09_REV_3_EN_TRABAJO/01_FUENTE_REVISADA/manuscrito/sections/md/03b_escala_accion_determinantal.tex; /Users/ruben/Documents/HMT2/EDICION_AUTONOMA_HMT_MD_2026-08-09_REV_3_EN_TRABAJO/01_FUENTE_REVISADA/manuscrito/sections/md/03c_elipse_accion_holonomia_rev8.tex","result_id":"MAS-RECTOR-005","statement":"delta_2w=H5-etaRet=(90/pi)Cpi(alpha); R_act=H5/etaRet","status":"PRESENT_IN_SUCCESSOR_SOURCE"}
% HMT_RESULT_END:MAS-RECTOR-005:residence
% HMT_RESULT_BEGIN:MAS-RECTOR-006:residence
% HMT_ATOMIC_RESULT {"material_state":"CONSERVADO_CON_APTO_FOCAL","proof_strength":"DEMOSTRADO_CON_ESTRUCTURA_DE_PARTIDA_EXPLICITA","provenance":"FORMALIZACION_NUEVA","reinforcement_state":"DEPLOYED_WITHOUT_NEW_RESULT_ID","residence":"/Users/ruben/Documents/HMT2/EDICION_AUTONOMA_HMT_MD_2026-08-09_REV_3_EN_TRABAJO/01_FUENTE_REVISADA/manuscrito/sections/md/03d_identidad_autodual_planck_ct108_rev3.tex; /Users/ruben/Documents/HMT2/EDICION_AUTONOMA_HMT_MD_2026-08-09_REV_3_EN_TRABAJO/01_FUENTE_REVISADA/manuscrito/sections/md/05aa_ley_general_masa_accion_autonoma.tex","result_id":"MAS-RECTOR-006","statement":"hbar_ret=etaRet 10^-34 U_S; omega0=2pi/(108t0); m_clk=hbar_ret omega0 c^-2","status":"PRESENT_IN_SUCCESSOR_SOURCE"}
% HMT_RESULT_END:MAS-RECTOR-006:residence
% HMT_RESULT_BEGIN:MAS-RECTOR-007:residence
% HMT_ATOMIC_RESULT {"material_state":"CONSERVADO_CON_APTO_FOCAL","proof_strength":"EXACTO_INTERNO","provenance":"RESULTADO_RECUPERADO","reinforcement_state":"DEPLOYED_WITHOUT_NEW_RESULT_ID","residence":"/Users/ruben/Documents/HMT2/EDICION_AUTONOMA_HMT_MD_2026-08-09_REV_3_EN_TRABAJO/01_FUENTE_REVISADA/manuscrito/sections/md/05f_operadores_masa_two_way_rev11.tex","result_id":"MAS-RECTOR-007","statement":"K_D=(D27+D36,D1,(D4-D3)/2); 14 perfiles","status":"PRESENT_IN_SUCCESSOR_SOURCE"}
% HMT_RESULT_END:MAS-RECTOR-007:residence
% HMT_RESULT_BEGIN:MAS-RECTOR-008:residence
% HMT_ATOMIC_RESULT {"material_state":"CONSERVADO_CON_APTO_FOCAL","proof_strength":"EXACTO_INTERNO","provenance":"RESULTADO_RECUPERADO","reinforcement_state":"DEPLOYED_WITHOUT_NEW_RESULT_ID","residence":"/Users/ruben/Documents/HMT2/EDICION_AUTONOMA_HMT_MD_2026-08-09_REV_3_EN_TRABAJO/01_FUENTE_REVISADA/manuscrito/sections/md/05f_operadores_masa_two_way_rev11.tex","result_id":"MAS-RECTOR-008","statement":"K_Omega=(Omega,54,P6); 9 perfiles","status":"PRESENT_IN_SUCCESSOR_SOURCE"}
% HMT_RESULT_END:MAS-RECTOR-008:residence
% HMT_RESULT_BEGIN:MAS-RECTOR-009:residence
% HMT_ATOMIC_RESULT {"material_state":"CONSERVADO_CON_APTO_FOCAL","proof_strength":"EXACTO_INTERNO","provenance":"CERTIFICADO_NUEVO","reinforcement_state":"DEPLOYED_WITHOUT_NEW_RESULT_ID","residence":"/Users/ruben/Documents/HMT2/EDICION_AUTONOMA_HMT_MD_2026-08-09_REV_3_EN_TRABAJO/01_FUENTE_REVISADA/manuscrito/sections/md/05f_operadores_masa_two_way_rev11.tex","result_id":"MAS-RECTOR-009","statement":"40 pares; 18 coincidencias; valor común único (80,54,6)","status":"PRESENT_IN_SUCCESSOR_SOURCE"}
% HMT_RESULT_END:MAS-RECTOR-009:residence
% HMT_RESULT_BEGIN:MAS-RECTOR-010:residence
% HMT_ATOMIC_RESULT {"material_state":"CONSERVADO_CON_APTO_FOCAL","proof_strength":"EXACTO_INTERNO","provenance":"RESULTADO_RECUPERADO","reinforcement_state":"DEPLOYED_WITHOUT_NEW_RESULT_ID","residence":"/Users/ruben/Documents/HMT2/EDICION_AUTONOMA_HMT_MD_2026-08-09_REV_3_EN_TRABAJO/01_FUENTE_REVISADA/manuscrito/sections/md/06f_biografia_electron_rev6.tex","result_id":"MAS-RECTOR-010","statement":"(5,5,++); tau12=+++---+++---; D=(54,56,68,48,32)","status":"PRESENT_IN_SUCCESSOR_SOURCE"}
% HMT_RESULT_END:MAS-RECTOR-010:residence
% HMT_RESULT_BEGIN:MAS-RECTOR-011:residence
% HMT_ATOMIC_RESULT {"material_state":"CONSERVADO_CON_APTO_FOCAL","proof_strength":"DEMOSTRADO_CON_ESTRUCTURA_DE_PARTIDA_EXPLICITA","provenance":"RESULTADO_RECUPERADO","reinforcement_state":"DEPLOYED_WITHOUT_NEW_RESULT_ID","residence":"/Users/ruben/Documents/HMT2/EDICION_AUTONOMA_HMT_MD_2026-08-09_REV_3_EN_TRABAJO/01_FUENTE_REVISADA/manuscrito/sections/md/06f_biografia_electron_rev6.tex; /Users/ruben/Documents/HMT2/EDICION_AUTONOMA_HMT_MD_2026-08-09_REV_3_EN_TRABAJO/01_FUENTE_REVISADA/manuscrito/sections/md/06ka_clausura_nulo_material_electron_rev3.tex","result_id":"MAS-RECTOR-011","statement":"e0=(sqrt3/4)[exp(phi/pi^2)-22alpha^3](1+15Delta4)=0.51099892248806607250987087719134943763; etapa basal, no salida final","status":"PRESENT_IN_SUCCESSOR_SOURCE"}
% HMT_RESULT_END:MAS-RECTOR-011:residence
% HMT_RESULT_BEGIN:MAS-RECTOR-012:residence
% HMT_ATOMIC_RESULT {"material_state":"CONSERVADO_CON_APTO_FOCAL","proof_strength":"EXACTO_INTERNO","provenance":"RESULTADO_RECUPERADO","reinforcement_state":"DEPLOYED_WITHOUT_NEW_RESULT_ID","residence":"/Users/ruben/Documents/HMT2/EDICION_AUTONOMA_HMT_MD_2026-08-09_REV_3_EN_TRABAJO/01_FUENTE_REVISADA/manuscrito/sections/md/06k_electron_two_way_rev11.tex; /Users/ruben/Documents/HMT2/EDICION_AUTONOMA_HMT_MD_2026-08-09_REV_3_EN_TRABAJO/01_FUENTE_REVISADA/manuscrito/sections/md/06ka_clausura_nulo_material_electron_rev3.tex","result_id":"MAS-RECTOR-012","statement":"kappa_e=80-54alpha+6Delta4; m_e^beta=e0 R_act^kappa_e=0.51099895069000080860825716483792339562463486771513350353936201542059113716425185","status":"PRESENT_IN_SUCCESSOR_SOURCE"}
% HMT_RESULT_END:MAS-RECTOR-012:residence
% HMT_RESULT_BEGIN:MAS-RECTOR-013:residence
% HMT_ATOMIC_RESULT {"material_state":"CONSERVADO_CON_APTO_FOCAL","proof_strength":"EXACTO_INTERNO","provenance":"RESULTADO_RECUPERADO","reinforcement_state":"DEPLOYED_WITHOUT_NEW_RESULT_ID","residence":"/Users/ruben/Documents/HMT2/EDICION_AUTONOMA_HMT_MD_2026-08-09_REV_3_EN_TRABAJO/01_FUENTE_REVISADA/manuscrito/sections/md/06f_biografia_electron_rev6.tex; /Users/ruben/Documents/HMT2/EDICION_AUTONOMA_HMT_MD_2026-08-09_REV_3_EN_TRABAJO/01_FUENTE_REVISADA/manuscrito/sections/md/06k_electron_two_way_rev11.tex","result_id":"MAS-RECTOR-013","statement":"firma cruda; firma par CT108; eventos; v_e=0 no se fusionan","status":"PRESENT_IN_SUCCESSOR_SOURCE"}
% HMT_RESULT_END:MAS-RECTOR-013:residence
% HMT_RESULT_BEGIN:MAS-RECTOR-014:residence
% HMT_ATOMIC_RESULT {"material_state":"CONSERVADO_CON_APTO_FOCAL","proof_strength":"EXACTO_INTERNO","provenance":"RESULTADO_RECUPERADO","reinforcement_state":"DEPLOYED_WITHOUT_NEW_RESULT_ID","residence":"/Users/ruben/Documents/HMT2/EDICION_AUTONOMA_HMT_MD_2026-08-09_REV_3_EN_TRABAJO/01_FUENTE_REVISADA/manuscrito/sections/md/05a_jerarquia_espectral_infinita_rev7.tex","result_id":"MAS-RECTOR-014","statement":"s_n=q_-^n-q_+^n; c_n=q_-^n+q_+^n; S=<90,120>={90,120}∪{30r:r≥6}","status":"PRESENT_IN_SUCCESSOR_SOURCE"}
% HMT_RESULT_END:MAS-RECTOR-014:residence
% HMT_RESULT_BEGIN:MAS-RECTOR-015:residence
% HMT_ATOMIC_RESULT {"material_state":"CONSERVADO_CON_APTO_FOCAL","proof_strength":"EXACTO_INTERNO","provenance":"FORMALIZACION_NUEVA","reinforcement_state":"DEPLOYED_WITHOUT_NEW_RESULT_ID","residence":"/Users/ruben/Documents/HMT2/EDICION_AUTONOMA_HMT_MD_2026-08-09_REV_3_EN_TRABAJO/01_FUENTE_REVISADA/manuscrito/sections/md/05a_jerarquia_espectral_infinita_rev7.tex","result_id":"MAS-RECTOR-015","statement":"k[S0]=k[X90,Y120]/(X90^4-Y120^3); dos genealogías a 360 antes del cociente","status":"PRESENT_IN_SUCCESSOR_SOURCE"}
% HMT_RESULT_END:MAS-RECTOR-015:residence
% HMT_RESULT_BEGIN:MAS-RECTOR-016:residence
% HMT_ATOMIC_RESULT {"material_state":"CONSERVADO_CON_APTO_FOCAL","proof_strength":"EXACTO_INTERNO","provenance":"RESULTADO_RECUPERADO","reinforcement_state":"DEPLOYED_WITHOUT_NEW_RESULT_ID","residence":"/Users/ruben/Documents/HMT2/EDICION_AUTONOMA_HMT_MD_2026-08-09_REV_3_EN_TRABAJO/01_FUENTE_REVISADA/manuscrito/sections/md/05f_operadores_masa_two_way_rev11.tex","result_id":"MAS-RECTOR-016","statement":"B_M^D y B_M^Omega diagonales exactos con espectros completos","status":"PRESENT_IN_SUCCESSOR_SOURCE"}
% HMT_RESULT_END:MAS-RECTOR-016:residence
% HMT_RESULT_BEGIN:MAS-RECTOR-017:residence
% HMT_ATOMIC_RESULT {"material_state":"CONSERVADO_CON_APTO_FOCAL","proof_strength":"EXACTO_INTERNO","provenance":"FORMALIZACION_NUEVA","reinforcement_state":"DEPLOYED_WITHOUT_NEW_RESULT_ID","residence":"/Users/ruben/Documents/HMT2/EDICION_AUTONOMA_HMT_MD_2026-08-09_REV_3_EN_TRABAJO/01_FUENTE_REVISADA/manuscrito/sections/md/05f_operadores_masa_two_way_rev11.tex","result_id":"MAS-RECTOR-017","statement":"M_M^{beta,r}=R_act^{B_M^r/2} M_M^(0) R_act^{B_M^r/2}","status":"PRESENT_IN_SUCCESSOR_SOURCE"}
% HMT_RESULT_END:MAS-RECTOR-017:residence
% HMT_RESULT_BEGIN:MAS-RECTOR-018:residence
% HMT_ATOMIC_RESULT {"material_state":"CONSERVADO_CON_APTO_FOCAL","proof_strength":"EXACTO_INTERNO","provenance":"CERTIFICADO_NUEVO","reinforcement_state":"DEPLOYED_WITHOUT_NEW_RESULT_ID","residence":"/Users/ruben/Documents/HMT2/EDICION_AUTONOMA_HMT_MD_2026-08-09_REV_3_EN_TRABAJO/01_FUENTE_REVISADA/manuscrito/sections/md/06m_tabla_completa_masas_hmt_sm_revfinal.tex","result_id":"MAS-RECTOR-018","statement":"Toda clase conserva residencia, operador/salida y estatuto; electrón beta y nulos cerrados; la publicación mantiene 23 filas de salida interna y 23 filas de comparabilidad física, una por clase","status":"PRESENT_IN_SUCCESSOR_SOURCE"}
% HMT_RESULT_END:MAS-RECTOR-018:residence
% HMT_RESULT_BEGIN:MAS-RECTOR-019:residence
% HMT_ATOMIC_RESULT {"material_state":"CONSERVADO_CON_APTO_FOCAL","proof_strength":"EXACTO_INTERNO","provenance":"RESULTADO_RECUPERADO","reinforcement_state":"DEPLOYED_WITHOUT_NEW_RESULT_ID","residence":"/Users/ruben/Documents/HMT2/EDICION_AUTONOMA_HMT_MD_2026-08-09_REV_3_EN_TRABAJO/01_FUENTE_REVISADA/manuscrito/sections/md/06m_tabla_completa_masas_hmt_sm_revfinal.tex","result_id":"MAS-RECTOR-019","statement":"m=0 en carta nula; kappa no aplicable","status":"PRESENT_IN_SUCCESSOR_SOURCE"}
% HMT_RESULT_END:MAS-RECTOR-019:residence
% HMT_RESULT_BEGIN:MAS-RECTOR-020:residence
% HMT_ATOMIC_RESULT {"material_state":"CONSERVADO_CON_APTO_FOCAL","proof_strength":"INTERFAZ_FISICA","provenance":"RESULTADO_RECUPERADO","reinforcement_state":"DEPLOYED_WITHOUT_NEW_RESULT_ID","residence":"/Users/ruben/Documents/HMT2/EDICION_AUTONOMA_HMT_MD_2026-08-09_REV_3_EN_TRABAJO/01_FUENTE_REVISADA/manuscrito/sections/md/06i_inventario_humano_23_sectores_rev8.tex; /Users/ruben/Documents/HMT2/EDICION_AUTONOMA_HMT_MD_2026-08-09_REV_3_EN_TRABAJO/01_FUENTE_REVISADA/manuscrito/sections/md/06m_tabla_completa_masas_hmt_sm_revfinal.tex","result_id":"MAS-RECTOR-020","statement":"Rama up/down, generación, color y carácter fino; scalarización exige join individual y esquema/escala","status":"PRESENT_IN_SUCCESSOR_SOURCE"}
% HMT_RESULT_END:MAS-RECTOR-020:residence
% HMT_RESULT_BEGIN:MAS-RECTOR-021:residence
% HMT_ATOMIC_RESULT {"material_state":"CONSERVADO_CON_APTO_FOCAL","proof_strength":"INTERFAZ_FISICA","provenance":"RESULTADO_RECUPERADO","reinforcement_state":"DEPLOYED_WITHOUT_NEW_RESULT_ID","residence":"/Users/ruben/Documents/HMT2/EDICION_AUTONOMA_HMT_MD_2026-08-09_REV_3_EN_TRABAJO/01_FUENTE_REVISADA/manuscrito/sections/md/06i_inventario_humano_23_sectores_rev8.tex; /Users/ruben/Documents/HMT2/EDICION_AUTONOMA_HMT_MD_2026-08-09_REV_3_EN_TRABAJO/01_FUENTE_REVISADA/manuscrito/sections/md/08b_transporte_ckm_pmns_rev11.tex; /Users/ruben/Documents/HMT2/EDICION_AUTONOMA_HMT_MD_2026-08-09_REV_3_EN_TRABAJO/01_FUENTE_REVISADA/manuscrito/sections/md/05a1_generaciones_quarks_neutrinos_rev2.tex","result_id":"MAS-RECTOR-021","statement":"Tres clases de sabor neutrínico sobre el par operatorial común B_F^D/B_F^Omega; los tres sabores comparten H_nu=directsum_{j=1}^4 H_{nu,G_j} de dimensión 20; f_e,f_mu,f_tau es base de interacción y G_j profundidad estructural; PMNS transporta bases y no crea autovalores, escala ni masa de sabor","status":"PRESENT_IN_SUCCESSOR_SOURCE"}
% HMT_RESULT_END:MAS-RECTOR-021:residence
% HMT_RESULT_BEGIN:MAS-RECTOR-022:residence
% HMT_ATOMIC_RESULT {"material_state":"CONSERVADO_CON_APTO_FOCAL","proof_strength":"RECONOCIMIENTO_EXTERNO","provenance":"RESULTADO_RECUPERADO","reinforcement_state":"DEPLOYED_WITHOUT_NEW_RESULT_ID","residence":"/Users/ruben/Documents/HMT2/EDICION_AUTONOMA_HMT_MD_2026-08-09_REV_3_EN_TRABAJO/01_FUENTE_REVISADA/manuscrito/sections/md/08_ckm_cp.tex; /Users/ruben/Documents/HMT2/EDICION_AUTONOMA_HMT_MD_2026-08-09_REV_3_EN_TRABAJO/01_FUENTE_REVISADA/manuscrito/sections/md/08b_transporte_ckm_pmns_rev11.tex","result_id":"MAS-RECTOR-022","statement":"Cuatro fórmulas angulares, matriz unitaria y J; contraste posterior","status":"PRESENT_IN_SUCCESSOR_SOURCE"}
% HMT_RESULT_END:MAS-RECTOR-022:residence
% HMT_RESULT_BEGIN:MAS-RECTOR-023:residence
% HMT_ATOMIC_RESULT {"material_state":"CONSERVADO_CON_APTO_FOCAL","proof_strength":"PROGRAMA_ABIERTO","provenance":"RESULTADO_RECUPERADO","reinforcement_state":"DEPLOYED_WITHOUT_NEW_RESULT_ID","residence":"/Users/ruben/Documents/HMT2/EDICION_AUTONOMA_HMT_MD_2026-08-09_REV_3_EN_TRABAJO/01_FUENTE_REVISADA/manuscrito/sections/md/06g_composicion_especies_rev6.tex","result_id":"MAS-RECTOR-023","statement":"Rutas compuestas con frontera/enlace; 432 y 1728 expresiones no son especies","status":"PRESENT_IN_SUCCESSOR_SOURCE"}
% HMT_RESULT_END:MAS-RECTOR-023:residence
% HMT_RESULT_BEGIN:MAS-RECTOR-024:residence
% HMT_ATOMIC_RESULT {"material_state":"CONSERVADO_CON_APTO_FOCAL","proof_strength":"EXACTO_INTERNO","provenance":"RESULTADO_RECUPERADO","reinforcement_state":"DEPLOYED_WITHOUT_NEW_RESULT_ID","residence":"/Users/ruben/Documents/HMT2/EDICION_AUTONOMA_HMT_MD_2026-08-09_REV_3_EN_TRABAJO/01_FUENTE_REVISADA/manuscrito/sections/md/06c_sectores_topologicos_no_omision.tex","result_id":"MAS-RECTOR-024","statement":"Anillos, dimensiones y trenzado son codominios propios, no masas escalares automáticas","status":"PRESENT_IN_SUCCESSOR_SOURCE"}
% HMT_RESULT_END:MAS-RECTOR-024:residence
% HMT_RESULT_BEGIN:MAS-RECTOR-025:residence
% HMT_ATOMIC_RESULT {"material_state":"CONSERVADO_CON_APTO_FOCAL","proof_strength":"EXACTO_INTERNO","provenance":"RESULTADO_RECUPERADO","reinforcement_state":"DEPLOYED_WITHOUT_NEW_RESULT_ID","residence":"/Users/ruben/Documents/HMT2/EDICION_AUTONOMA_HMT_MD_2026-08-09_REV_3_EN_TRABAJO/01_FUENTE_REVISADA/manuscrito/sections/md/04b_banda_particula_virtual.tex","result_id":"MAS-RECTOR-025","statement":"Sección de prolongación finita con ledger hasta L; no especie estable ni masa imaginaria","status":"PRESENT_IN_SUCCESSOR_SOURCE"}
% HMT_RESULT_END:MAS-RECTOR-025:residence
% HMT_RESULT_BEGIN:MAS-RECTOR-026:residence
% HMT_ATOMIC_RESULT {"material_state":"CONSERVADO_CON_APTO_FOCAL","proof_strength":"EXACTO_INTERNO","provenance":"CERTIFICADO_NUEVO","reinforcement_state":"DEPLOYED_WITHOUT_NEW_RESULT_ID","residence":"/Users/ruben/Documents/HMT2/EDICION_AUTONOMA_HMT_MD_2026-08-09_REV_3_EN_TRABAJO/01_FUENTE_REVISADA/manuscrito/sections/md/06m_tabla_completa_masas_hmt_sm_revfinal.tex","result_id":"MAS-RECTOR-026","statement":"Salida interna congelada antes de abrir el catálogo externo","status":"PRESENT_IN_SUCCESSOR_SOURCE"}
% HMT_RESULT_END:MAS-RECTOR-026:residence
% HMT_RESULT_BEGIN:MAS-RECTOR-027:residence
% HMT_ATOMIC_RESULT {"material_state":"CONSERVADO_CON_APTO_FOCAL","proof_strength":"EXACTO_INTERNO","provenance":"RESULTADO_RECUPERADO","reinforcement_state":"DEPLOYED_WITHOUT_NEW_RESULT_ID","residence":"/Users/ruben/Documents/HMT2/EDICION_AUTONOMA_HMT_MD_2026-08-09_REV_3_EN_TRABAJO/01_FUENTE_REVISADA/manuscrito/sections/md/06m_tabla_completa_masas_hmt_sm_revfinal.tex","result_id":"MAS-RECTOR-027","statement":"Sólo trazabilidad histórica; no residencia narrativa pública ni dominio","status":"PRESENT_IN_SUCCESSOR_SOURCE"}
% HMT_RESULT_END:MAS-RECTOR-027:residence
% HMT_RESULT_BEGIN:MAS-RECTOR-028:residence
% HMT_ATOMIC_RESULT {"material_state":"CONSERVADO_CON_APTO_FOCAL","proof_strength":"DEMOSTRADO_CON_ESTRUCTURA_DE_PARTIDA_EXPLICITA","provenance":"RESULTADO_RECUPERADO","reinforcement_state":"DEPLOYED_WITHOUT_NEW_RESULT_ID","residence":"/Users/ruben/Documents/HMT2/EDICION_AUTONOMA_HMT_MD_2026-08-09_REV_3_EN_TRABAJO/01_FUENTE_REVISADA/manuscrito/sections/md/06l_vacancias_corriente_rev11.tex","result_id":"MAS-RECTOR-028","statement":"Matriz [[21,22],[6,7]] con det=15 acompaña (-22,+15); uso local, no global","status":"PRESENT_IN_SUCCESSOR_SOURCE"}
% HMT_RESULT_END:MAS-RECTOR-028:residence
% HMT_RESULT_BEGIN:MAS-PUBLIC-023:residence
% HMT_ATOMIC_RESULT {"material_state":"CONSERVADO_CON_APTO_FOCAL","proof_strength":"EXACTO_INTERNO","provenance":"CERTIFICADO_NUEVO","reinforcement_state":"DEPLOYED_WITHOUT_NEW_RESULT_ID","residence":"/Users/ruben/Documents/HMT2/EDICION_AUTONOMA_HMT_MD_2026-08-09_REV_3_EN_TRABAJO/01_FUENTE_REVISADA/manuscrito/sections/md/06m_tabla_completa_masas_hmt_sm_revfinal.tex","result_id":"MAS-PUBLIC-023","statement":"las 23 claves internas reaparecen una vez y en el mismo orden en la capa pública de comparabilidad; las seis clases de codominio no escalar se comparan en su codominio nativo y no reciben una masa universal inventada","status":"PRESENT_IN_SUCCESSOR_SOURCE"}
% HMT_RESULT_END:MAS-PUBLIC-023:residence
% HMT_RESULT_BEGIN:MAS-COND-017:residence
% HMT_ATOMIC_RESULT {"material_state":"CONSERVADO_CON_APTO_FOCAL","proof_strength":"INTERFAZ_FISICA","provenance":"RESULTADO_RECUPERADO","reinforcement_state":"DEPLOYED_WITHOUT_NEW_RESULT_ID","residence":"/Users/ruben/Documents/HMT2/EDICION_AUTONOMA_HMT_MD_2026-08-09_REV_3_EN_TRABAJO/01_FUENTE_REVISADA/manuscrito/sections/md/06m_tabla_completa_masas_hmt_sm_revfinal.tex","result_id":"MAS-COND-017","statement":"17 evaluaciones exactas, 11 elementales y 6 compuestas, residen en una tabla separada y rotulada como no prospectiva; no son salidas cerradas, no seleccionan ruta y sólo admiten residual cuando los escalares, unidad, esquema y escala son comparables","status":"PRESENT_IN_SUCCESSOR_SOURCE"}
% HMT_RESULT_END:MAS-COND-017:residence
% HMT_RESULT_BEGIN:INF-001:residence
% HMT_ATOMIC_RESULT {"material_state":"INTEGRADO_O_REMITIDO_SIN_PERDIDA","proof_strength":"EXACTO_INTERNO","provenance":"RESULTADO_RECUPERADO","reinforcement_state":"DEPLOYED_WITHOUT_NEW_RESULT_ID","residence":"apertura del capítulo del continuo, tras `hmt/00b`","result_id":"INF-001","statement":"El continuo HMT no se forma con valores desnudos: cada casilla transporta bloque, firma, acarreo, fase, etiqueta y orientación. \\(\\xi_m=(B_m,\\sigma_m,c_m,\\omega_m)\\).","status":"PRESENT_IN_SUCCESSOR_SOURCE"}
% HMT_RESULT_END:INF-001:residence
% HMT_RESULT_BEGIN:INF-002:residence
% HMT_ATOMIC_RESULT {"material_state":"INTEGRADO_O_REMITIDO_SIN_PERDIDA","proof_strength":"EXACTO_INTERNO","provenance":"RESULTADO_RECUPERADO","reinforcement_state":"DEPLOYED_WITHOUT_NEW_RESULT_ID","residence":"bloque fundacional, § Estado, sección y valor, tras INF-001","result_id":"INF-002","statement":"Un número HMT es una sección compatible, no su expansión ni un conjunto acabado sin genealogía. \\(\\Omega_\\infty=\\varprojlim_m\\mathcal S_m\\).","status":"PRESENT_IN_SUCCESSOR_SOURCE"}
% HMT_RESULT_END:INF-002:residence
% HMT_RESULT_BEGIN:INF-003:residence
% HMT_ATOMIC_RESULT {"material_state":"INTEGRADO_O_REMITIDO_SIN_PERDIDA","proof_strength":"DEMOSTRADO_CON_ESTRUCTURA_DE_PARTIDA_EXPLICITA","provenance":"RESULTADO_RECUPERADO","reinforcement_state":"DEPLOYED_WITHOUT_NEW_RESULT_ID","residence":"bloque fundacional, § Existencia y unicidad coinductivas, tras INF-002","result_id":"INF-003","statement":"Finitud, no vaciedad y sobreyectividad dan existencia por König; la unicidad exige además fibras unitarias compatibles a toda profundidad.","status":"PRESENT_IN_SUCCESSOR_SOURCE"}
% HMT_RESULT_END:INF-003:residence
% HMT_RESULT_BEGIN:N9-001:residence
% HMT_ATOMIC_RESULT {"material_state":"INTEGRADO_O_REMITIDO_SIN_PERDIDA","proof_strength":"DEMOSTRADO_CON_ESTRUCTURA_DE_PARTIDA_EXPLICITA","provenance":"ARQUITECTURA_AUTORAL_PREEXISTENTE","reinforcement_state":"DEPLOYED_WITHOUT_NEW_RESULT_ID","residence":"`hmt/06aa` y transición inicial","result_id":"N9-001","statement":"La cadena APP--TRIT--TPK abre la frontera coinductiva; el retorno de fase no es retorno de estado.","status":"PRESENT_IN_SUCCESSOR_SOURCE"}
% HMT_RESULT_END:N9-001:residence
% HMT_RESULT_BEGIN:N9-002:residence
% HMT_ATOMIC_RESULT {"material_state":"INTEGRADO_O_REMITIDO_SIN_PERDIDA","proof_strength":"EXACTO_INTERNO","provenance":"RESULTADO_RECUPERADO","reinforcement_state":"DEPLOYED_WITHOUT_NEW_RESULT_ID","residence":"`hmt/06aa`","result_id":"N9-002","statement":"Las tres evaluaciones son lecturas coordinadas de un mismo estado; \\(\\varphi\\) fija el eje y \\(\\pi/e\\) constituyen la fibra especular orientada.","status":"PRESENT_IN_SUCCESSOR_SOURCE"}
% HMT_RESULT_END:N9-002:residence
% HMT_RESULT_BEGIN:N9-003:residence
% HMT_ATOMIC_RESULT {"material_state":"INTEGRADO_O_REMITIDO_SIN_PERDIDA","proof_strength":"EXACTO_INTERNO","provenance":"RESULTADO_RECUPERADO","reinforcement_state":"DEPLOYED_WITHOUT_NEW_RESULT_ID","residence":"`hmt/06b--06e`","result_id":"N9-003","statement":"Una salida visible conserva cinco genealogías regionales. \\(5=3_{++}+1_{--}+1_\\perp\\), \\(432+4\\cdot144=1008\\).","status":"PRESENT_IN_SUCCESSOR_SOURCE"}
% HMT_RESULT_END:N9-003:residence
% HMT_RESULT_BEGIN:N9-004:residence
% HMT_ATOMIC_RESULT {"material_state":"INTEGRADO_O_REMITIDO_SIN_PERDIDA","proof_strength":"EXACTO_INTERNO","provenance":"CERTIFICADO_NUEVO","reinforcement_state":"DEPLOYED_WITHOUT_NEW_RESULT_ID","residence":"`hmt/06e--06f`","result_id":"N9-004","statement":"Las 1.000 cifras son una restricción reproducible, no el techo ni por sí solas la prueba lógica del límite.","status":"PRESENT_IN_SUCCESSOR_SOURCE"}
% HMT_RESULT_END:N9-004:residence
% HMT_RESULT_BEGIN:HEL-001:residence
% HMT_ATOMIC_RESULT {"material_state":"INTEGRADO_O_REMITIDO_SIN_PERDIDA","proof_strength":"EXACTO_INTERNO","provenance":"RESULTADO_RECUPERADO","reinforcement_state":"DEPLOYED_WITHOUT_NEW_RESULT_ID","residence":"`hmt/03l`","result_id":"HEL-001","statement":"«Diez vuelve a uno» significa misma fase y memoria incrementada.","status":"PRESENT_IN_SUCCESSOR_SOURCE"}
% HMT_RESULT_END:HEL-001:residence
% HMT_RESULT_BEGIN:OPS-001:residence
% HMT_ATOMIC_RESULT {"material_state":"INTEGRADO_O_REMITIDO_SIN_PERDIDA","proof_strength":"EXACTO_INTERNO","provenance":"RESULTADO_RECUPERADO","reinforcement_state":"DEPLOYED_WITHOUT_NEW_RESULT_ID","residence":"`hmt/03k`","result_id":"OPS-001","statement":"Suma, producto, potencia, raíz, logaritmo, derivación e integración deben elevarse con memoria de fibra.","status":"PRESENT_IN_SUCCESSOR_SOURCE"}
% HMT_RESULT_END:OPS-001:residence
% HMT_RESULT_BEGIN:ZF-001:residence
% HMT_ATOMIC_RESULT {"material_state":"INTEGRADO_O_REMITIDO_SIN_PERDIDA","proof_strength":"EXACTO_INTERNO","provenance":"RESULTADO_RECUPERADO","reinforcement_state":"DEPLOYED_WITHOUT_NEW_RESULT_ID","residence":"nuevo bloque «Genealogía, extensionalidad y teoría de conjuntos» después de completar APP--TRIT--TPK y `03l`","result_id":"ZF-001","statement":"HMT no refuta ZFC: construye estructuras genealógicas codificables en ZFC y muestra qué información pierde la decategorización conjuntista.","status":"PRESENT_IN_SUCCESSOR_SOURCE"}
% HMT_RESULT_END:ZF-001:residence
% HMT_RESULT_BEGIN:ZF-002:residence
% HMT_ATOMIC_RESULT {"material_state":"INTEGRADO_O_REMITIDO_SIN_PERDIDA","proof_strength":"PROGRAMA_ABIERTO","provenance":"RESULTADO_NUEVO","reinforcement_state":"DEPLOYED_WITHOUT_NEW_RESULT_ID","residence":"bloque lógico, § Alcance axiomático, tras ZF-001","result_id":"ZF-002","statement":"No se afirmará que HMT sea ya un modelo interno completo de ZF, IZF o NBG: Separación, Reemplazo, Potencia plena e Infinito transfinito no se han verificado conjuntamente.","status":"PRESENT_IN_SUCCESSOR_SOURCE"}
% HMT_RESULT_END:ZF-002:residence
% HMT_RESULT_BEGIN:EXT-001:residence
% HMT_ATOMIC_RESULT {"material_state":"INTEGRADO_O_REMITIDO_SIN_PERDIDA","proof_strength":"EXACTO_INTERNO","provenance":"RESULTADO_RECUPERADO","reinforcement_state":"DEPLOYED_WITHOUT_NEW_RESULT_ID","residence":"bloque fundacional, § Extensionalidad tipada, tras DSC-001","result_id":"EXT-001","statement":"Dos secciones son idénticas si coinciden a toda profundidad; igualdad de valor es un cociente salvo lector fiel.","status":"PRESENT_IN_SUCCESSOR_SOURCE"}
% HMT_RESULT_END:EXT-001:residence
% HMT_RESULT_BEGIN:FND-001:residence
% HMT_ATOMIC_RESULT {"material_state":"INTEGRADO_O_REMITIDO_SIN_PERDIDA","proof_strength":"EXACTO_INTERNO","provenance":"RESULTADO_RECUPERADO","reinforcement_state":"DEPLOYED_WITHOUT_NEW_RESULT_ID","residence":"bloque lógico, § Fundación y coinducción, tras TOP-001","result_id":"FND-001","statement":"Los prefijos son bien fundados por rango natural; una sección infinita es coinductiva y no introduce ciclos de pertenencia.","status":"PRESENT_IN_SUCCESSOR_SOURCE"}
% HMT_RESULT_END:FND-001:residence
% HMT_RESULT_BEGIN:ORD-001:residence
% HMT_ATOMIC_RESULT {"material_state":"INTEGRADO_O_REMITIDO_SIN_PERDIDA","proof_strength":"EXACTO_INTERNO","provenance":"FORMALIZACION_NUEVA","reinforcement_state":"DEPLOYED_WITHOUT_NEW_RESULT_ID","residence":"subsección «Ordinal, cardinal y ruta»","result_id":"ORD-001","statement":"El ordinal de von Neumann se eleva conservando el estado recorrido; la proyección recupera el ordinal desnudo.","status":"PRESENT_IN_SUCCESSOR_SOURCE"}
% HMT_RESULT_END:ORD-001:residence
% HMT_RESULT_BEGIN:ORD-002:residence
% HMT_ATOMIC_RESULT {"material_state":"INTEGRADO_O_REMITIDO_SIN_PERDIDA","proof_strength":"PROGRAMA_ABIERTO","provenance":"FORMALIZACION_NUEVA","reinforcement_state":"DEPLOYED_WITHOUT_NEW_RESULT_ID","residence":"bloque lógico, § Jerarquía transfinita, tras ORD-001","result_id":"ORD-002","statement":"El corpus cierra profundidades finitas y secciones de longitud \\(\\omega\\); la extensión a todo ordinal es programa.","status":"PRESENT_IN_SUCCESSOR_SOURCE"}
% HMT_RESULT_END:ORD-002:residence
% HMT_RESULT_BEGIN:PWR-001:residence
% HMT_ATOMIC_RESULT {"material_state":"INTEGRADO_O_REMITIDO_SIN_PERDIDA","proof_strength":"EXACTO_INTERNO","provenance":"FORMALIZACION_NUEVA","reinforcement_state":"DEPLOYED_WITHOUT_NEW_RESULT_ID","residence":"subsección «Potencia genealógica»","result_id":"PWR-001","statement":"Los predicados decidibles a profundidad finita forman los clopen del límite.","status":"PRESENT_IN_SUCCESSOR_SOURCE"}
% HMT_RESULT_END:PWR-001:residence
% HMT_RESULT_BEGIN:PWR-002:residence
% HMT_ATOMIC_RESULT {"material_state":"INTEGRADO_O_REMITIDO_SIN_PERDIDA","proof_strength":"EXACTO_INTERNO","provenance":"FORMALIZACION_NUEVA","reinforcement_state":"DEPLOYED_WITHOUT_NEW_RESULT_ID","residence":"bloque lógico, § Potencia cerrada, tras PWR-001","result_id":"PWR-002","statement":"Familias compatibles de subconjuntos finitos reconstruyen cerrados, no todos los subconjuntos.","status":"PRESENT_IN_SUCCESSOR_SOURCE"}
% HMT_RESULT_END:PWR-002:residence
% HMT_RESULT_BEGIN:CHO-001:residence
% HMT_ATOMIC_RESULT {"material_state":"INTEGRADO_O_REMITIDO_SIN_PERDIDA","proof_strength":"EXACTO_INTERNO","provenance":"RESULTADO_RECUPERADO","reinforcement_state":"DEPLOYED_WITHOUT_NEW_RESULT_ID","residence":"subsección «Elección tipada»","result_id":"CHO-001","statement":"Se separan elección conjuntista, natural, algebraica y efectiva. Toda sección natural, algebraica o efectiva es en particular una sección conjuntista, pero naturalidad, algebraicidad y efectividad son requisitos adicionales generalmente incomparables entre sí.","status":"PRESENT_IN_SUCCESSOR_SOURCE"}
% HMT_RESULT_END:CHO-001:residence
% HMT_RESULT_BEGIN:CHO-002:residence
% HMT_ATOMIC_RESULT {"material_state":"INTEGRADO_O_REMITIDO_SIN_PERDIDA","proof_strength":"EXACTO_INTERNO","provenance":"RESULTADO_RECUPERADO","reinforcement_state":"DEPLOYED_WITHOUT_NEW_RESULT_ID","residence":"bloque lógico, § Elección tipada y obstrucción homomorfa, tras CHO-001","result_id":"CHO-002","statement":"Hay representantes conjuntistas pero no sección homomorfa; eso mide memoria estructural, no fracaso de AC.","status":"PRESENT_IN_SUCCESSOR_SOURCE"}
% HMT_RESULT_END:CHO-002:residence
% HMT_RESULT_BEGIN:CAT-001:residence
% HMT_ATOMIC_RESULT {"material_state":"INTEGRADO_O_REMITIDO_SIN_PERDIDA","proof_strength":"EXACTO_INTERNO","provenance":"RESULTADO_RECUPERADO","reinforcement_state":"DEPLOYED_WITHOUT_NEW_RESULT_ID","residence":"cierre categorial del nuevo bloque","result_id":"CAT-001","statement":"Objetos como torres tipadas; morfismos conmutan restricciones; las secciones son el funtor límite.","status":"PRESENT_IN_SUCCESSOR_SOURCE"}
% HMT_RESULT_END:CAT-001:residence
% HMT_RESULT_BEGIN:TOP-001:residence
% HMT_ATOMIC_RESULT {"material_state":"INTEGRADO_O_REMITIDO_SIN_PERDIDA","proof_strength":"EXACTO_INTERNO","provenance":"FORMALIZACION_NUEVA","reinforcement_state":"DEPLOYED_WITHOUT_NEW_RESULT_ID","residence":"nota técnica subordinada a CAT-001","result_id":"TOP-001","statement":"La categoría pequeña de caminos TPK admite una envolvente de presheaves; no se le atribuyen elección o lógica clásica sin prueba.","status":"PRESENT_IN_SUCCESSOR_SOURCE"}
% HMT_RESULT_END:TOP-001:residence
% HMT_RESULT_BEGIN:HIL-001:residence
% HMT_ATOMIC_RESULT {"material_state":"INTEGRADO_O_REMITIDO_SIN_PERDIDA","proof_strength":"EXACTO_INTERNO","provenance":"RESULTADO_RECUPERADO","reinforcement_state":"DEPLOYED_WITHOUT_NEW_RESULT_ID","residence":"`hmt/03l` y `md/04ab`","result_id":"HIL-001","statement":"Cada profundidad posee una realización de Hilbert y una álgebra matricial.","status":"PRESENT_IN_SUCCESSOR_SOURCE"}
% HMT_RESULT_END:HIL-001:residence
% HMT_RESULT_BEGIN:UHF-001:residence
% HMT_ATOMIC_RESULT {"material_state":"INTEGRADO_O_REMITIDO_SIN_PERDIDA","proof_strength":"EXACTO_INTERNO","provenance":"FORMALIZACION_NUEVA","reinforcement_state":"DEPLOYED_WITHOUT_NEW_RESULT_ID","residence":"nuevo bloque «Torre nonádica de observables» tras `hmt/03l`","result_id":"UHF-001","statement":"Conservar simultáneamente las nueve fibras preserva unidad y traza y genera la rama UHF--II_1; una esquina marcada no lo hace y su corredor de tipo I_infty se conserva separadamente en IINF-001.","status":"PRESENT_IN_SUCCESSOR_SOURCE"}
% HMT_RESULT_END:UHF-001:residence
% HMT_RESULT_BEGIN:II1-001:residence
% HMT_ATOMIC_RESULT {"material_state":"INTEGRADO_O_REMITIDO_SIN_PERDIDA","proof_strength":"EXACTO_INTERNO","provenance":"FORMALIZACION_NUEVA","reinforcement_state":"DEPLOYED_WITHOUT_NEW_RESULT_ID","residence":"torre operatorial, § Cierre tracial hiperinfinito II_1, tras IINF-001","result_id":"II1-001","statement":"El límite unital de todas las puertas produce UHF nonádico y su cierre tracial es el factor hiperinfinito \\(II_1\\).","status":"PRESENT_IN_SUCCESSOR_SOURCE"}
% HMT_RESULT_END:II1-001:residence
% HMT_RESULT_BEGIN:DIM-001:residence
% HMT_ATOMIC_RESULT {"material_state":"INTEGRADO_O_REMITIDO_SIN_PERDIDA","proof_strength":"EXACTO_INTERNO","provenance":"ARQUITECTURA_AUTORAL_PREEXISTENTE","reinforcement_state":"DEPLOYED_WITHOUT_NEW_RESULT_ID","residence":"torre operatorial, § Dimensión continua de proyectores, tras II1-001","result_id":"DIM-001","statement":"Para todo t∈[0,1], m_d=floor(9^d t) permite elegir proyectores anidados cuyo límite fuerte P_t satisface tau(P_t)=t; en el factor finito, P~Q iff tau(P)=tau(Q).","status":"PRESENT_IN_SUCCESSOR_SOURCE"}
% HMT_RESULT_END:DIM-001:residence
% HMT_RESULT_BEGIN:CLK-001:residence
% HMT_ATOMIC_RESULT {"material_state":"INTEGRADO_O_REMITIDO_SIN_PERDIDA","proof_strength":"EXACTO_INTERNO","provenance":"RESULTADO_RECUPERADO","reinforcement_state":"DEPLOYED_WITHOUT_NEW_RESULT_ID","residence":"subsección «Realización por reloj profinito»","result_id":"CLK-001","statement":"Una segunda ruta llega abstractamente al mismo factor mediante producto cruzado.","status":"PRESENT_IN_SUCCESSOR_SOURCE"}
% HMT_RESULT_END:CLK-001:residence
% HMT_RESULT_BEGIN:INT-001:residence
% HMT_ATOMIC_RESULT {"material_state":"INTEGRADO_O_REMITIDO_SIN_PERDIDA","proof_strength":"PROGRAMA_ABIERTO","provenance":"FORMALIZACION_NUEVA","reinforcement_state":"DEPLOYED_WITHOUT_NEW_RESULT_ID","residence":"torre operatorial, § Entrelazador APP--reloj, tras CLK-001","result_id":"INT-001","statement":"La isomorfía abstracta de factores no identifica genealogías ni los antecedentes \\(C^*\\).","status":"PRESENT_IN_SUCCESSOR_SOURCE"}
% HMT_RESULT_END:INT-001:residence
% HMT_RESULT_BEGIN:QLG-001:residence
% HMT_ATOMIC_RESULT {"material_state":"INTEGRADO_O_REMITIDO_SIN_PERDIDA","proof_strength":"EXACTO_INTERNO","provenance":"RESULTADO_RECUPERADO","reinforcement_state":"DEPLOYED_WITHOUT_NEW_RESULT_ID","residence":"`md/04ab` después de PVM/contextos","result_id":"QLG-001","statement":"\\(\\operatorname{Proj}(R_{\\rm hyp})\\) es un retículo ortomodular completo, continuo y no distributivo; cada contexto conmutativo es booleano.","status":"PRESENT_IN_SUCCESSOR_SOURCE"}
% HMT_RESULT_END:QLG-001:residence
% HMT_RESULT_BEGIN:ERG-001:residence
% HMT_ATOMIC_RESULT {"material_state":"INTEGRADO_O_REMITIDO_SIN_PERDIDA","proof_strength":"EXACTO_INTERNO","provenance":"RESULTADO_RECUPERADO","reinforcement_state":"DEPLOYED_WITHOUT_NEW_RESULT_ID","residence":"capítulo de torre/reloj","result_id":"ERG-001","statement":"Los promedios temporales proyectan sobre invariantes; para el reloj ergódico, sobre la media de Haar.","status":"PRESENT_IN_SUCCESSOR_SOURCE"}
% HMT_RESULT_END:ERG-001:residence
% HMT_RESULT_BEGIN:DEN-001:residence
% HMT_ATOMIC_RESULT {"material_state":"INTEGRADO_O_REMITIDO_SIN_PERDIDA","proof_strength":"EXACTO_INTERNO","provenance":"RESULTADO_RECUPERADO","reinforcement_state":"DEPLOYED_WITHOUT_NEW_RESULT_ID","residence":"`md/04ab` y torre \\(II_1\\)","result_id":"DEN-001","statement":"La incertidumbre sobre rutas se representa por estado positivo; no obliga a borrar la genealogía subyacente.","status":"PRESENT_IN_SUCCESSOR_SOURCE"}
% HMT_RESULT_END:DEN-001:residence
% HMT_RESULT_BEGIN:LAN-001:residence
% HMT_ATOMIC_RESULT {"material_state":"INTEGRADO_O_REMITIDO_SIN_PERDIDA","proof_strength":"EXACTO_INTERNO","provenance":"RESULTADO_RECUPERADO","reinforcement_state":"DEPLOYED_WITHOUT_NEW_RESULT_ID","residence":"`md/04g`","result_id":"LAN-001","statement":"Una actualización biyectiva U del estado completo satisface H(X|U(X))=0 y, para el término ideal de borrado lógico, Q_L^{min}[U]=0.","status":"PRESENT_IN_SUCCESSOR_SOURCE"}
% HMT_RESULT_END:LAN-001:residence
% HMT_RESULT_BEGIN:LAN-002:residence
% HMT_ATOMIC_RESULT {"material_state":"INTEGRADO_O_REMITIDO_SIN_PERDIDA","proof_strength":"EXACTO_INTERNO","provenance":"FORMALIZACION_NUEVA","reinforcement_state":"DEPLOYED_WITHOUT_NEW_RESULT_ID","residence":"`md/04g`, después del caso finito","result_id":"LAN-002","statement":"Automorfismos traza-preservantes conservan entropía; una esperanza condicional incrementa la entropía exactamente por entropía relativa.","status":"PRESENT_IN_SUCCESSOR_SOURCE"}
% HMT_RESULT_END:LAN-002:residence
% HMT_RESULT_BEGIN:GDL-001:residence
% HMT_ATOMIC_RESULT {"material_state":"INTEGRADO_O_REMITIDO_SIN_PERDIDA","proof_strength":"EXACTO_INTERNO","provenance":"RESULTADO_RECUPERADO","reinforcement_state":"DEPLOYED_WITHOUT_NEW_RESULT_ID","residence":"nuevo bloque «Límite efectivo de la clausura»","result_id":"GDL-001","statement":"Cada nivel puede calcularse mientras la unicidad de la fibra global codifica HALT.","status":"PRESENT_IN_SUCCESSOR_SOURCE"}
% HMT_RESULT_END:GDL-001:residence
% HMT_RESULT_BEGIN:GDL-002:residence
% HMT_ATOMIC_RESULT {"material_state":"INTEGRADO_O_REMITIDO_SIN_PERDIDA","proof_strength":"EXACTO_INTERNO","provenance":"RESULTADO_RECUPERADO","reinforcement_state":"DEPLOYED_WITHOUT_NEW_RESULT_ID","residence":"bloque metalógico, § Incompletitud semántica Gödel--Turing, tras GDL-001","result_id":"GDL-002","statement":"Toda teoría recursivamente axiomatizable y correcta para la familia \\(U_e\\) deja alguna sentencia sin decidir; no implica inconsistencia de ZFC.","status":"PRESENT_IN_SUCCESSOR_SOURCE"}
% HMT_RESULT_END:GDL-002:residence
% HMT_RESULT_BEGIN:IND-001:residence
% HMT_ATOMIC_RESULT {"material_state":"INTEGRADO_O_REMITIDO_SIN_PERDIDA","proof_strength":"EXACTO_INTERNO","provenance":"FORMALIZACION_NUEVA","reinforcement_state":"DEPLOYED_WITHOUT_NEW_RESULT_ID","residence":"bloque metalógico, § Frontera Gödel--Cohen, tras GDL-002","result_id":"IND-001","statement":"Bajo las hipótesis de consistencia correspondientes, CH es independiente de ZFC y AC es independiente de ZF; HMT enriquece objetos, no decide por sí sola esos enunciados.","status":"PRESENT_IN_SUCCESSOR_SOURCE"}
% HMT_RESULT_END:IND-001:residence
% HMT_RESULT_BEGIN:FOR-001:residence
% HMT_ATOMIC_RESULT {"material_state":"INTEGRADO_O_REMITIDO_SIN_PERDIDA","proof_strength":"EXACTO_INTERNO","provenance":"FORMALIZACION_NUEVA","reinforcement_state":"DEPLOYED_WITHOUT_NEW_RESULT_ID","residence":"subsección de filtros de prefijos","result_id":"FOR-001","statement":"Una sección determina un filtro que atraviesa todos los densos de profundidad; eso no basta para genericidad sobre un modelo.","status":"PRESENT_IN_SUCCESSOR_SOURCE"}
% HMT_RESULT_END:FOR-001:residence
% HMT_RESULT_BEGIN:FOR-002:residence
% HMT_ATOMIC_RESULT {"material_state":"INTEGRADO_O_REMITIDO_SIN_PERDIDA","proof_strength":"INTERFAZ_FISICA","provenance":"FORMALIZACION_NUEVA","reinforcement_state":"DEPLOYED_WITHOUT_NEW_RESULT_ID","residence":"bloque metalógico, § Forcing cilíndrico y filtro genérico, tras FOR-001","result_id":"FOR-002","statement":"Árbol perfecto, numerable y sin átomos da forcing cilíndrico Cohen-equivalente; rama única da forcing trivial.","status":"PRESENT_IN_SUCCESSOR_SOURCE"}
% HMT_RESULT_END:FOR-002:residence
% HMT_RESULT_BEGIN:BOR-001:residence
% HMT_ATOMIC_RESULT {"material_state":"INTEGRADO_O_REMITIDO_SIN_PERDIDA","proof_strength":"EXACTO_INTERNO","provenance":"FORMALIZACION_NUEVA","reinforcement_state":"DEPLOYED_WITHOUT_NEW_RESULT_ID","residence":"subsección final del bloque","result_id":"BOR-001","statement":"Los cilindros generan una clase Borel estable por complemento, unión numerable e imágenes inversas continuas; sus juegos usan determinación boreliana clásica.","status":"PRESENT_IN_SUCCESSOR_SOURCE"}
% HMT_RESULT_END:BOR-001:residence
% HMT_RESULT_BEGIN:MIN-001:residence
% HMT_ATOMIC_RESULT {"material_state":"INTEGRADO_O_REMITIDO_SIN_PERDIDA","proof_strength":"INTERFAZ_FISICA","provenance":"FORMALIZACION_NUEVA","reinforcement_state":"DEPLOYED_WITHOUT_NEW_RESULT_ID","residence":"bloque metalógico, § Minimax finito y proyectivo, tras BOR-001","result_id":"MIN-001","statement":"En cada profundidad finita, von Neumann da valor mixto; un límite global requiere compactitud, continuidad y consistencia proyectiva, y no selecciona por sí solo rama pura.","status":"PRESENT_IN_SUCCESSOR_SOURCE"}
% HMT_RESULT_END:MIN-001:residence
% HMT_RESULT_BEGIN:III-001:residence
% HMT_ATOMIC_RESULT {"material_state":"INTEGRADO_O_REMITIDO_SIN_PERDIDA","proof_strength":"PROGRAMA_ABIERTO","provenance":"RESULTADO_NUEVO","reinforcement_state":"DEPLOYED_WITHOUT_NEW_RESULT_ID","residence":"frontera tras \\(II_1\\)","result_id":"III-001","statement":"La construcción actual produce naturalmente \\(II_1\\); tipo \\(III\\) exige GNS no tracial/KMS o cociclo de Radon--Nikodym no trivial.","status":"PRESENT_IN_SUCCESSOR_SOURCE"}
% HMT_RESULT_END:III-001:residence
% HMT_RESULT_BEGIN:TRI-001:residence
% HMT_ATOMIC_RESULT {"material_state":"INTEGRADO_O_REMITIDO_SIN_PERDIDA","proof_strength":"INTERFAZ_FISICA","provenance":"FORMALIZACION_NUEVA","reinforcement_state":"DEPLOYED_WITHOUT_NEW_RESULT_ID","residence":"después de exponer íntegramente la doble proyección y completar el corredor Paley--Witt--Golay--Leech--\\(V^\\natural\\)--Monstruo--Moonshine","result_id":"TRI-001","statement":"Punto, dimensión y genealogía requieren tres capas coordinadas.","status":"PRESENT_IN_SUCCESSOR_SOURCE"}
% HMT_RESULT_END:TRI-001:residence
% HMT_RESULT_BEGIN:EXC-001:residence
% HMT_ATOMIC_RESULT {"material_state":"INTEGRADO_O_REMITIDO_SIN_PERDIDA","proof_strength":"DEMOSTRADO_CON_ESTRUCTURA_DE_PARTIDA_EXPLICITA","provenance":"RESULTADO_RECUPERADO","reinforcement_state":"DEPLOYED_WITHOUT_NEW_RESULT_ID","residence":"`hmt/09g`, `15`","result_id":"EXC-001","statement":"El mismo estado superviviente alimenta lectura real e incidencia excepcional hasta Leech--\\(V^\\natural\\)--Monstruo/Moonshine; el Monstruo no actúa literalmente sobre dígitos.","status":"PRESENT_IN_SUCCESSOR_SOURCE"}
% HMT_RESULT_END:EXC-001:residence
% HMT_RESULT_BEGIN:MD-001:residence
% HMT_ATOMIC_RESULT {"material_state":"INTEGRADO_O_REMITIDO_SIN_PERDIDA","proof_strength":"PROGRAMA_ABIERTO","provenance":"RESULTADO_RECUPERADO","reinforcement_state":"DEPLOYED_WITHOUT_NEW_RESULT_ID","residence":"`md/00b`, antes de Planck/electrón","result_id":"MD-001","statement":"Una sección compatible no es todavía una partícula: la ruta es su historia observable y la partícula una sección persistente enriquecida antes del morfismo de realización física.","status":"PRESENT_IN_SUCCESSOR_SOURCE"}
% HMT_RESULT_END:MD-001:residence
% HMT_RESULT_BEGIN:PSC-001:residence
% HMT_ATOMIC_RESULT {"material_state":"INTEGRADO_O_REMITIDO_SIN_PERDIDA","proof_strength":"EXACTO_INTERNO","provenance":"ARQUITECTURA_AUTORAL_PREEXISTENTE","reinforcement_state":"DEPLOYED_WITHOUT_NEW_RESULT_ID","residence":"apéndice crítico, fuera del tronco probatorio","result_id":"PSC-001","statement":"Las ramas AD/MM, CH y «cierre fuerte» conservan prioridad intuitiva, no fuerza probatoria vigente.","status":"PRESENT_IN_SUCCESSOR_SOURCE"}
% HMT_RESULT_END:PSC-001:residence
% HMT_RESULT_BEGIN:COV-001:residence
% HMT_ATOMIC_RESULT {"material_state":"INTEGRADO_O_REMITIDO_SIN_PERDIDA","proof_strength":"EXACTO_INTERNO","provenance":"RESULTADO_RECUPERADO","reinforcement_state":"DEPLOYED_WITHOUT_NEW_RESULT_ID","residence":"gate de cobertura final","result_id":"COV-001","statement":"Constantes, electrón, masas, Barbero, CKM/CP, gravedad y corredor excepcional se preservan por remisión, sin versión degradada.","status":"PRESENT_IN_SUCCESSOR_SOURCE"}
% HMT_RESULT_END:COV-001:residence
% HMT_RESULT_BEGIN:IINF-001:residence
% HMT_ATOMIC_RESULT {"material_state":"INTEGRADO_O_REMITIDO_SIN_PERDIDA","proof_strength":"EXACTO_INTERNO","provenance":"RESULTADO_RECUPERADO","reinforcement_state":"DEPLOYED_WITHOUT_NEW_RESULT_ID","residence":"torre operatorial, § Rama marcada de tipo I_infty, tras UHF-001","result_id":"IINF-001","statement":"Las esquinas j_{d,j}(a)=a⊗p_j, implementadas por V_{d,j}aV_{d,j}^*, generan en el límite los operadores de rango finito; su clausura normativa es K(H_infty) y su clausura débil B(H_infty), factor de tipo I_infty.","status":"PRESENT_IN_SUCCESSOR_SOURCE"}
% HMT_RESULT_END:IINF-001:residence
% HMT_RESULT_BEGIN:DEN-003:residence
% HMT_ATOMIC_RESULT {"material_state":"INTEGRADO_O_REMITIDO_SIN_PERDIDA","proof_strength":"EXACTO_INTERNO","provenance":"RESULTADO_RECUPERADO","reinforcement_state":"DEPLOYED_WITHOUT_NEW_RESULT_ID","residence":"torre \\(II_1\\), inmediatamente después de DEN-002","result_id":"DEN-003","statement":"Si h=9^d rho, entonces log h=d log9 I+log rho y S_tau(h)=S_vN(rho)-d log9; bajo rho_{d+1}=rho_d⊗I9/9, S_vN aumenta en log9 y S_tau permanece invariante.","status":"PRESENT_IN_SUCCESSOR_SOURCE"}
% HMT_RESULT_END:DEN-003:residence
% HMT_RESULT_BEGIN:CARD-001:residence
% HMT_ATOMIC_RESULT {"material_state":"INTEGRADO_O_REMITIDO_SIN_PERDIDA","proof_strength":"EXACTO_INTERNO","provenance":"RESULTADO_RECUPERADO","reinforcement_state":"DEPLOYED_WITHOUT_NEW_RESULT_ID","residence":"bloque lógico, § Cardinalidad como funtor no fiel, tras ORD-002","result_id":"CARD-001","statement":"Contar es decategorizar: el perfil (|X_m|)_m no determina restricciones, acarreo, orientación o memoria.","status":"PRESENT_IN_SUCCESSOR_SOURCE"}
% HMT_RESULT_END:CARD-001:residence
% HMT_RESULT_BEGIN:PWR-003:residence
% HMT_ATOMIC_RESULT {"material_state":"INTEGRADO_O_REMITIDO_SIN_PERDIDA","proof_strength":"EXACTO_INTERNO","provenance":"FORMALIZACION_NUEVA","reinforcement_state":"DEPLOYED_WITHOUT_NEW_RESULT_ID","residence":"bloque lógico, § Tres escalas de potencia, tras PWR-002","result_id":"PWR-003","statement":"En un Cantor HMT, clopen, cerrados y potencia plena tienen tamaños distintos.","status":"PRESENT_IN_SUCCESSOR_SOURCE"}
% HMT_RESULT_END:PWR-003:residence
% HMT_RESULT_BEGIN:WEYL-001:residence
% HMT_ATOMIC_RESULT {"material_state":"INTEGRADO_O_REMITIDO_SIN_PERDIDA","proof_strength":"EXACTO_INTERNO","provenance":"RESULTADO_RECUPERADO","reinforcement_state":"DEPLOYED_WITHOUT_NEW_RESULT_ID","residence":"`hmt/14_numero_heisenberg_y_d108.tex`, `md/04ab`","result_id":"WEYL-001","statement":"Fase y desplazamiento forman un par no conmutativo; fijado el carácter central primitivo, la representación irreducible es única hasta unitario.","status":"PRESENT_IN_SUCCESSOR_SOURCE"}
% HMT_RESULT_END:WEYL-001:residence
% HMT_RESULT_BEGIN:DSC-001:residence
% HMT_ATOMIC_RESULT {"material_state":"INTEGRADO_O_REMITIDO_SIN_PERDIDA","proof_strength":"EXACTO_INTERNO","provenance":"RESULTADO_RECUPERADO","reinforcement_state":"DEPLOYED_WITHOUT_NEW_RESULT_ID","residence":"bloque fundacional, § Presentación y descenso por fibras, tras INF-003","result_id":"DSC-001","statement":"Existe un único operador descendido O_bar:V→Y con O=O_bar q si y sólo si O es constante en cada fibra de q.","status":"PRESENT_IN_SUCCESSOR_SOURCE"}
% HMT_RESULT_END:DSC-001:residence
% HMT_RESULT_BEGIN:REC-001:residence
% HMT_ATOMIC_RESULT {"material_state":"INTEGRADO_O_REMITIDO_SIN_PERDIDA","proof_strength":"EXACTO_INTERNO","provenance":"RESULTADO_RECUPERADO","reinforcement_state":"DEPLOYED_WITHOUT_NEW_RESULT_ID","residence":"`md/04g`, antes del teorema de Landauer lógico","result_id":"REC-001","statement":"Si q_hat=(q,M) es inyectivo sobre supp(X), entonces H(X|q(X),M(X))=0; H(X|q(X)) mide la fibra olvidada por el lector visible.","status":"PRESENT_IN_SUCCESSOR_SOURCE"}
% HMT_RESULT_END:REC-001:residence
% HMT_RESULT_BEGIN:LAN-003:residence
% HMT_ATOMIC_RESULT {"material_state":"INTEGRADO_O_REMITIDO_SIN_PERDIDA","proof_strength":"EXACTO_INTERNO","provenance":"RESULTADO_RECUPERADO","reinforcement_state":"DEPLOYED_WITHOUT_NEW_RESULT_ID","residence":"`md/04g`, después de LAN-001","result_id":"LAN-003","statement":"Para un reset no inyectivo de tres estados equiprobables, H(X)=log 3 y el término ideal satisface Q_reset≥k_B Θ log 3.","status":"PRESENT_IN_SUCCESSOR_SOURCE"}
% HMT_RESULT_END:LAN-003:residence
% HMT_RESULT_BEGIN:FIN-001:residence
% HMT_ATOMIC_RESULT {"material_state":"INTEGRADO_O_REMITIDO_SIN_PERDIDA","proof_strength":"EXACTO_INTERNO","provenance":"FORMALIZACION_NUEVA","reinforcement_state":"DEPLOYED_WITHOUT_NEW_RESULT_ID","residence":"nuevo bloque «Límite efectivo de la clausura», antes de GDL-001","result_id":"FIN-001","statement":"Fijado N, con dominios finitos y observables atómicos computables, toda fórmula del lenguaje finito L_N es semánticamente decidible.","status":"PRESENT_IN_SUCCESSOR_SOURCE"}
% HMT_RESULT_END:FIN-001:residence
% HMT_RESULT_BEGIN:DEN-002:residence
% HMT_ATOMIC_RESULT {"material_state":"INTEGRADO_O_REMITIDO_SIN_PERDIDA","proof_strength":"EXACTO_INTERNO","provenance":"RESULTADO_RECUPERADO","reinforcement_state":"DEPLOYED_WITHOUT_NEW_RESULT_ID","residence":"torre \\(II_1\\), después de DEN-001","result_id":"DEN-002","statement":"Para h=9^d rho, S_tau(h)=S_vN(rho)-d log9; la extensión rho_{d+1}=rho_d⊗I9/9 aumenta S_vN en log9 y deja invariante S_tau.","status":"PRESENT_IN_SUCCESSOR_SOURCE"}
% HMT_RESULT_END:DEN-002:residence
% HMT_RESULT_BEGIN:UNI-001:residence
% HMT_ATOMIC_RESULT {"material_state":"INTEGRADO_O_REMITIDO_SIN_PERDIDA","proof_strength":"EXACTO_INTERNO","provenance":"RESULTADO_RECUPERADO","reinforcement_state":"DEPLOYED_WITHOUT_NEW_RESULT_ID","residence":"cierre APP--EMRH","result_id":"UNI-001","statement":"La frontera distingue 999=729+243+27 de 1000=729+243+27+1=729+271; el acarreo c_{m+1}=1 conserva la unidad visible y registra una vuelta adicional.","status":"PRESENT_IN_SUCCESSOR_SOURCE"}
% HMT_RESULT_END:UNI-001:residence
% HMT_RESULT_BEGIN:UNI-002:residence
% HMT_ATOMIC_RESULT {"material_state":"INTEGRADO_O_REMITIDO_SIN_PERDIDA","proof_strength":"PROGRAMA_ABIERTO","provenance":"ARQUITECTURA_AUTORAL_PREEXISTENTE","reinforcement_state":"DEPLOYED_WITHOUT_NEW_RESULT_ID","residence":"transición APP--EMRH hacia monodromía e inclusión unital nonádica","result_id":"UNI-002","statement":"La compatibilidad entre c_{m+1}=1, la monodromía 9→1 y a↦a⊗I9 es una obligación de entrelazamiento: debe construirse un diagrama que transporte la misma unidad con memoria entre los tres dominios.","status":"PRESENT_IN_SUCCESSOR_SOURCE"}
% HMT_RESULT_END:UNI-002:residence
% HMT_RESULT_BEGIN:P01:residence
% HMT_ATOMIC_RESULT {"material_state":"INTEGRADO_EN_GRAFO_ACTIVO","proof_strength":"EXACTO_INTERNO","provenance":"FORMALIZACION_NUEVA","reinforcement_state":"DEPLOYED_WITHOUT_NEW_RESULT_ID","residence":"APP vigente","result_id":"P01","statement":"APP tiene soporte $X_9$, 81 vértices y todos los caminos finitos.","status":"PRESENT_IN_SUCCESSOR_SOURCE"}
% HMT_RESULT_END:P01:residence
% HMT_RESULT_BEGIN:P02:residence
% HMT_ATOMIC_RESULT {"material_state":"INTEGRADO_EN_GRAFO_ACTIVO","proof_strength":"EXACTO_INTERNO","provenance":"FORMALIZACION_NUEVA","reinforcement_state":"DEPLOYED_WITHOUT_NEW_RESULT_ID","residence":"APP vigente","result_id":"P02","statement":"$\\Sigma$ es latina; $\\Pi$ sólo permuta en unidades.","status":"PRESENT_IN_SUCCESSOR_SOURCE"}
% HMT_RESULT_END:P02:residence
% HMT_RESULT_BEGIN:P03:residence
% HMT_ATOMIC_RESULT {"material_state":"INTEGRADO_EN_GRAFO_ACTIVO","proof_strength":"EXACTO_INTERNO","provenance":"FORMALIZACION_NUEVA","reinforcement_state":"DEPLOYED_WITHOUT_NEW_RESULT_ID","residence":"APP/TRIT","result_id":"P03","statement":"La elevación residuo--cociente recompone $1215=54+9\\cdot129$.","status":"PRESENT_IN_SUCCESSOR_SOURCE"}
% HMT_RESULT_END:P03:residence
% HMT_RESULT_BEGIN:P04:residence
% HMT_ATOMIC_RESULT {"material_state":"INTEGRADO_EN_GRAFO_ACTIVO","proof_strength":"EXACTO_INTERNO","provenance":"FORMALIZACION_NUEVA","reinforcement_state":"DEPLOYED_WITHOUT_NEW_RESULT_ID","residence":"`hmt/03f`","result_id":"P04","statement":"$B_c=9P_++5P_-$ y el salto es cuatro.","status":"PRESENT_IN_SUCCESSOR_SOURCE"}
% HMT_RESULT_END:P04:residence
% HMT_RESULT_BEGIN:P05:residence
% HMT_ATOMIC_RESULT {"material_state":"INTEGRADO_EN_GRAFO_ACTIVO","proof_strength":"EXACTO_INTERNO","provenance":"FORMALIZACION_NUEVA","reinforcement_state":"DEPLOYED_WITHOUT_NEW_RESULT_ID","residence":"TRIT","result_id":"P05","statement":"TRIT visible no determina el estado levantado.","status":"PRESENT_IN_SUCCESSOR_SOURCE"}
% HMT_RESULT_END:P05:residence
% HMT_RESULT_BEGIN:P06:residence
% HMT_ATOMIC_RESULT {"material_state":"INTEGRADO_EN_GRAFO_ACTIVO","proof_strength":"EXACTO_INTERNO","provenance":"FORMALIZACION_NUEVA","reinforcement_state":"DEPLOYED_WITHOUT_NEW_RESULT_ID","residence":"TPK","result_id":"P06","statement":"TPK compone rutas con hoja, acarreo y memoria.","status":"PRESENT_IN_SUCCESSOR_SOURCE"}
% HMT_RESULT_END:P06:residence
% HMT_RESULT_BEGIN:P07:residence
% HMT_ATOMIC_RESULT {"material_state":"INTEGRADO_EN_GRAFO_ACTIVO","proof_strength":"EXACTO_INTERNO","provenance":"RESULTADO_RECUPERADO","reinforcement_state":"DEPLOYED_WITHOUT_NEW_RESULT_ID","residence":"`md/04ab`","result_id":"P07","statement":"Los operadores de traslación X_u y fase Z_v satisfacen Z_vX_u=omega^{u·v}X_uZ_v, generan M_{9^d}(C) y, fijado el carácter central primitivo, dan la representación irreducible finita única salvo equivalencia unitaria.","status":"PRESENT_IN_SUCCESSOR_SOURCE"}
% HMT_RESULT_END:P07:residence
% HMT_RESULT_BEGIN:P08:residence
% HMT_ATOMIC_RESULT {"material_state":"INTEGRADO_EN_GRAFO_ACTIVO","proof_strength":"EXACTO_INTERNO","provenance":"RESULTADO_RECUPERADO","reinforcement_state":"DEPLOYED_WITHOUT_NEW_RESULT_ID","residence":"bloque UHF--\\(II_1\\)","result_id":"P08","statement":"El límite inductivo de a↦a⊗I_9 es UHF(9^infty), canónicamente isomorfa a UHF(3^infty); la representación GNS de su traza única cierra en el factor hiperinfinito II_1.","status":"PRESENT_IN_SUCCESSOR_SOURCE"}
% HMT_RESULT_END:P08:residence
% HMT_RESULT_BEGIN:P09:residence
% HMT_ATOMIC_RESULT {"material_state":"INTEGRADO_EN_GRAFO_ACTIVO","proof_strength":"EXACTO_INTERNO","provenance":"FORMALIZACION_NUEVA","reinforcement_state":"DEPLOYED_WITHOUT_NEW_RESULT_ID","residence":"apertura del corredor matricial","result_id":"P09","statement":"Un QMS es una cuadrícula positiva con sumas $I_s$.","status":"PRESENT_IN_SUCCESSOR_SOURCE"}
% HMT_RESULT_END:P09:residence
% HMT_RESULT_BEGIN:P10:residence
% HMT_ATOMIC_RESULT {"material_state":"INTEGRADO_EN_GRAFO_ACTIVO","proof_strength":"EXACTO_INTERNO","provenance":"RESULTADO_RECUPERADO","reinforcement_state":"DEPLOYED_WITHOUT_NEW_RESULT_ID","residence":"corredor matricial, § Semiclasicidad como factorización conmutativa, tras P09/QPM-001","result_id":"P10","statement":"A es semiclásico si y sólo si existe un mapa positivo unital Phi:C(S_n)→M_s tal que A_ij=Phi(f_ij), donde f_ij(sigma)=1_{sigma(i)=j}; por ser conmutativo el dominio, Phi es completamente positivo.","status":"PRESENT_IN_SUCCESSOR_SOURCE"}
% HMT_RESULT_END:P10:residence
% HMT_RESULT_BEGIN:P11:residence
% HMT_ATOMIC_RESULT {"material_state":"INTEGRADO_EN_GRAFO_ACTIVO","proof_strength":"EXACTO_INTERNO","provenance":"RESULTADO_RECUPERADO","reinforcement_state":"DEPLOYED_WITHOUT_NEW_RESULT_ID","residence":"corredor matricial, § QLS fáciles y semiclásicos, tras QLS-001","result_id":"P11","statement":"Un cuadrado latino cuántico es semiclásico si y sólo si es fácil, es decir, procede de un cuadrado latino clásico y una base ortonormal fija.","status":"PRESENT_IN_SUCCESSOR_SOURCE"}
% HMT_RESULT_END:P11:residence
% HMT_RESULT_BEGIN:P12:residence
% HMT_ATOMIC_RESULT {"material_state":"INTEGRADO_EN_GRAFO_ACTIVO","proof_strength":"EXACTO_INTERNO","provenance":"FORMALIZACION_NUEVA","reinforcement_state":"DEPLOYED_WITHOUT_NEW_RESULT_ID","residence":"capítulo qutrit","result_id":"P12","statement":"Cuatro contextos HMT producen cuatro QLS fáciles.","status":"PRESENT_IN_SUCCESSOR_SOURCE"}
% HMT_RESULT_END:P12:residence
% HMT_RESULT_BEGIN:P13:residence
% HMT_ATOMIC_RESULT {"material_state":"INTEGRADO_EN_GRAFO_ACTIVO","proof_strength":"EXACTO_INTERNO","provenance":"RESULTADO_NUEVO","reinforcement_state":"DEPLOYED_WITHOUT_NEW_RESULT_ID","residence":"cuadrados operatoriales, § No-go de QLS cíclicos, tras P12","result_id":"P13","statement":"Dos filas MUB no pueden apilarse en una QLS.","status":"PRESENT_IN_SUCCESSOR_SOURCE"}
% HMT_RESULT_END:P13:residence
% HMT_RESULT_BEGIN:P14:residence
% HMT_ATOMIC_RESULT {"material_state":"INTEGRADO_EN_GRAFO_ACTIVO","proof_strength":"EXACTO_INTERNO","provenance":"FORMALIZACION_NUEVA","reinforcement_state":"DEPLOYED_WITHOUT_NEW_RESULT_ID","residence":"tras TPK y torre UHF","result_id":"P14","statement":"Para n fijo, profundidad genealógica m y nivel matricial s definen una bigraduación; si n varía, la familia es trigraduada por (m,n,s).","status":"PRESENT_IN_SUCCESSOR_SOURCE"}
% HMT_RESULT_END:P14:residence
% HMT_RESULT_BEGIN:P15:residence
% HMT_ATOMIC_RESULT {"material_state":"INTEGRADO_EN_GRAFO_ACTIVO","proof_strength":"EXACTO_INTERNO","provenance":"FORMALIZACION_NUEVA","reinforcement_state":"DEPLOYED_WITHOUT_NEW_RESULT_ID","residence":"sistema inverso UCP, § Compresiones y truncación, tras P14","result_id":"P15","statement":"Si Phi_r∈K_{m+1,s_r} y V_r:C^s→C^{s_r} satisfacen sum_r V_r^*V_r=I_s, entonces la columna V=(V_1,\\ldots,V_k)^T es isométrica y la restricción genealógica conmuta con la combinación matricial: rho_m^s(sum_r V_r^*Phi_r V_r)=sum_r V_r^*(rho_m^{s_r}Phi_r)V_r.","status":"PRESENT_IN_SUCCESSOR_SOURCE"}
% HMT_RESULT_END:P15:residence
% HMT_RESULT_BEGIN:P16:residence
% HMT_ATOMIC_RESULT {"material_state":"INTEGRADO_EN_GRAFO_ACTIVO","proof_strength":"EXACTO_INTERNO","provenance":"FORMALIZACION_NUEVA","reinforcement_state":"DEPLOYED_WITHOUT_NEW_RESULT_ID","residence":"sistema inverso UCP, § Dilatación cilíndrica compatible y PVM común, tras P15","result_id":"P16","statement":"Existe una única Phi_infty:C(X_infty)→M_s que extiende todas las cartas cilíndricas y una dilatación común Phi_infty(f)=V*∫f dP V, con E_x^(m)=V*P([x]_m)V; la dilatación mínima es única hasta unitario y, fijada pi, la PVM P es única.","status":"PRESENT_IN_SUCCESSOR_SOURCE"}
% HMT_RESULT_END:P16:residence
% HMT_RESULT_BEGIN:P17:residence
% HMT_ATOMIC_RESULT {"material_state":"INTEGRADO_EN_GRAFO_ACTIVO","proof_strength":"EXACTO_INTERNO","provenance":"FORMALIZACION_NUEVA","reinforcement_state":"DEPLOYED_WITHOUT_NEW_RESULT_ID","residence":"continuación UHF","result_id":"P17","statement":"Si Psi_m=Psi_{m+1}∘iota_m, existe una única UCP Psi_infty:A_infty→M_s y una dilatación común Psi_m(a)=V* pi(iota_{m,infty}(a))V; pi puede hacerse fiel.","status":"PRESENT_IN_SUCCESSOR_SOURCE"}
% HMT_RESULT_END:P17:residence
% HMT_RESULT_BEGIN:P18:residence
% HMT_ATOMIC_RESULT {"material_state":"INTEGRADO_EN_GRAFO_ACTIVO","proof_strength":"EXACTO_INTERNO","provenance":"FORMALIZACION_NUEVA","reinforcement_state":"DEPLOYED_WITHOUT_NEW_RESULT_ID","residence":"capítulo «Cuadrados cuánticos de APP»","result_id":"P18","statement":"Existe QMS aditivo $(9,3)$.","status":"PRESENT_IN_SUCCESSOR_SOURCE"}
% HMT_RESULT_END:P18:residence
% HMT_RESULT_BEGIN:P19:residence
% HMT_ATOMIC_RESULT {"material_state":"INTEGRADO_EN_GRAFO_ACTIVO","proof_strength":"EXACTO_INTERNO","provenance":"RESULTADO_NUEVO","reinforcement_state":"DEPLOYED_WITHOUT_NEW_RESULT_ID","residence":"cuadrados APP, § Construcción A^Sigma, tras P18","result_id":"P19","statement":"Tres contextos producen QMS $(9,3)$ no conmutativo y semiclásico.","status":"PRESENT_IN_SUCCESSOR_SOURCE"}
% HMT_RESULT_END:P19:residence
% HMT_RESULT_BEGIN:P20:residence
% HMT_ATOMIC_RESULT {"material_state":"INTEGRADO_EN_GRAFO_ACTIVO","proof_strength":"EXACTO_INTERNO","provenance":"RESULTADO_NUEVO","reinforcement_state":"DEPLOYED_WITHOUT_NEW_RESULT_ID","residence":"cuadrados APP, § Construcción A^(9), tras P19","result_id":"P20","statement":"Las doce direcciones producen QMS $(12,3)$ semiclásico.","status":"PRESENT_IN_SUCCESSOR_SOURCE"}
% HMT_RESULT_END:P20:residence
% HMT_RESULT_BEGIN:P21:residence
% HMT_ATOMIC_RESULT {"material_state":"INTEGRADO_EN_GRAFO_ACTIVO","proof_strength":"EXACTO_INTERNO","provenance":"RESULTADO_NUEVO","reinforcement_state":"DEPLOYED_WITHOUT_NEW_RESULT_ID","residence":"cuadrados APP, § Construcción B(12,3) y cubo de tres índices, tras P20","result_id":"P21","statement":"Las mismas \\POVM producen cubos de órdenes 9 y 12.","status":"PRESENT_IN_SUCCESSOR_SOURCE"}
% HMT_RESULT_END:P21:residence
% HMT_RESULT_BEGIN:P22:residence
% HMT_ATOMIC_RESULT {"material_state":"INTEGRADO_EN_GRAFO_ACTIVO","proof_strength":"EXACTO_INTERNO","provenance":"FORMALIZACION_NUEVA","reinforcement_state":"DEPLOYED_WITHOUT_NEW_RESULT_ID","residence":"subsección de fibración finita","result_id":"P22","statement":"P^1(Z/9Z) tiene doce puntos y r tiene cuatro fibras de tres; la fibración es canónica, mientras el orden de cada fibra y la bijección Xi hacia resultados qutrit constituyen un calibre.","status":"PRESENT_IN_SUCCESSOR_SOURCE"}
% HMT_RESULT_END:P22:residence
% HMT_RESULT_BEGIN:P23:residence
% HMT_ATOMIC_RESULT {"material_state":"INTEGRADO_EN_GRAFO_ACTIVO","proof_strength":"EXACTO_INTERNO","provenance":"FORMALIZACION_NUEVA","reinforcement_state":"DEPLOYED_WITHOUT_NEW_RESULT_ID","residence":"geometría qutrit, § Fibración 12→4×3 y POVM calibrada, tras P22","result_id":"P23","statement":"F_k^(pi)=(3/pi)P_k para k=0,1,2 y F_r^(pi)=((1-3/pi)/6)I_3 para r=3,…,8 forman una POVM; para un proyector Q de otro contexto MUB, Tr(QF_k^(pi))=(3/pi)(1/3)=1/pi","status":"PRESENT_IN_SUCCESSOR_SOURCE"}
% HMT_RESULT_END:P23:residence
% HMT_RESULT_BEGIN:P24:residence
% HMT_ATOMIC_RESULT {"material_state":"INTEGRADO_EN_GRAFO_ACTIVO","proof_strength":"EXACTO_INTERNO","provenance":"RESULTADO_NUEVO","reinforcement_state":"DEPLOYED_WITHOUT_NEW_RESULT_ID","residence":"después de separación espectral APP","result_id":"P24","statement":"El polinomio de CC* es t^5(t-1)(t-5/7)(t-1/3)^2 y la descomposición de Halmos da C*(P_{Sigma,2},P_{Pi,2})≅C^4⊕M_2(C)⊕M_2(C), de dimensión 12.","status":"PRESENT_IN_SUCCESSOR_SOURCE"}
% HMT_RESULT_END:P24:residence
% HMT_RESULT_BEGIN:P25:residence
% HMT_ATOMIC_RESULT {"material_state":"INTEGRADO_EN_GRAFO_ACTIVO","proof_strength":"EXACTO_INTERNO","provenance":"RESULTADO_NUEVO","reinforcement_state":"DEPLOYED_WITHOUT_NEW_RESULT_ID","residence":"doble proyección operatorial, § Intervalo libre del bloque B_c, tras P24","result_id":"P25","statement":"Para todo nivel s, W_s(B_c)={Y=Y*:5I_s≼Y≼9I_s}; 7 es centro, 2 radio operatorial y 4 diámetro espectral.","status":"PRESENT_IN_SUCCESSOR_SOURCE"}
% HMT_RESULT_END:P25:residence
% HMT_RESULT_BEGIN:P26:residence
% HMT_ATOMIC_RESULT {"material_state":"INTEGRADO_EN_GRAFO_ACTIVO","proof_strength":"EXACTO_INTERNO","provenance":"FORMALIZACION_NUEVA","reinforcement_state":"DEPLOYED_WITHOUT_NEW_RESULT_ID","residence":"doble proyección, antes de TRI-001","result_id":"P26","statement":"Sus PVM diagonales conmutan y admiten el refinamiento común P_{a,e}=1_{R_m=a}1_{Q_m=e}; si los lectores respetan truncaciones, el refinamiento pasa a una PVM cilíndrica común.","status":"PRESENT_IN_SUCCESSOR_SOURCE"}
% HMT_RESULT_END:P26:residence
% HMT_RESULT_BEGIN:P27:residence
% HMT_ATOMIC_RESULT {"material_state":"INTEGRADO_EN_GRAFO_ACTIVO","proof_strength":"EXACTO_INTERNO","provenance":"FORMALIZACION_NUEVA","reinforcement_state":"DEPLOYED_WITHOUT_NEW_RESULT_ID","residence":"doble proyección operatorial, § Límite de dilatación MUB, tras P26","result_id":"P27","statement":"No admiten una PVM conjunta aguda: sus proyectores tienen solapamiento 1/3 y no conmutan; una realización conjunta requiere ruido, POVM o un ambiente mayor.","status":"PRESENT_IN_SUCCESSOR_SOURCE"}
% HMT_RESULT_END:P27:residence
% HMT_RESULT_BEGIN:P28:residence
% HMT_ATOMIC_RESULT {"material_state":"INTEGRADO_EN_GRAFO_ACTIVO","proof_strength":"EXACTO_INTERNO","provenance":"RESULTADO_RECUPERADO","reinforcement_state":"DEPLOYED_WITHOUT_NEW_RESULT_ID","residence":"frontera del capítulo de continuo operatorial","result_id":"P28","statement":"El teorema importado afirma que la MPS posee límite continuo si y sólo si E_A es infinitamente divisible, equivalentemente E_A=P exp(L), P^2=P y PLP=PL; una potencia E_A^r basta para el límite continuo grueso.","status":"PRESENT_IN_SUCCESSOR_SOURCE"}
% HMT_RESULT_END:P28:residence
% HMT_RESULT_BEGIN:P29:residence
% HMT_ATOMIC_RESULT {"material_state":"INTEGRADO_EN_GRAFO_ACTIVO","proof_strength":"EXACTO_INTERNO","provenance":"FORMALIZACION_NUEVA","reinforcement_state":"DEPLOYED_WITHOUT_NEW_RESULT_ID","residence":"medición/expectativa tracial","result_id":"P29","statement":"E_m es UCP, satisface E_m(a⊗I_9)=a y posee Kraus W_a psi=(1/3)psi⊗e_a y una dilatación cuyo ancilla registra el dígito promediado; en la diagonal, la retracción de rama kappa_{m,a}(f)(x)=f(x,a) es multiplicativa y no equivale a promediar. Esa memoria no se identifica aún con TPK.","status":"PRESENT_IN_SUCCESSOR_SOURCE"}
% HMT_RESULT_END:P29:residence
% HMT_RESULT_BEGIN:P30:residence
% HMT_ATOMIC_RESULT {"material_state":"INTEGRADO_EN_GRAFO_ACTIVO","proof_strength":"EXACTO_INTERNO","provenance":"FORMALIZACION_NUEVA","reinforcement_state":"DEPLOYED_WITHOUT_NEW_RESULT_ID","residence":"APP matricial","result_id":"P30","statement":"El doble centrado inserta una señal APP en un QMS.","status":"PRESENT_IN_SUCCESSOR_SOURCE"}
% HMT_RESULT_END:P30:residence
% HMT_RESULT_BEGIN:P31:residence
% HMT_ATOMIC_RESULT {"material_state":"INTEGRADO_EN_GRAFO_ACTIVO","proof_strength":"PROGRAMA_ABIERTO","provenance":"ARQUITECTURA_AUTORAL_PREEXISTENTE","reinforcement_state":"DEPLOYED_WITHOUT_NEW_RESULT_ID","residence":"perspectivas operatoriales","result_id":"P31","statement":"El programa asigna a cada tipo c un cono de lectores positivos A(c), con transformaciones naturales por niveles; requiere funtor de conos, objeto dualizante, evaluaciones, coevaluaciones y compatibilidad de la composición TPK. La profundidad de separación es el menor m donde los sistemas mínimo y máximo difieren; el nivel matricial de detección es el menor s que exhibe esa diferencia; y el defecto de memoria es la mínima información de ruta necesaria para levantar un punto del sistema máximo al mínimo.","status":"PRESENT_IN_SUCCESSOR_SOURCE"}
% HMT_RESULT_END:P31:residence
% HMT_RESULT_BEGIN:P32:residence
% HMT_ATOMIC_RESULT {"material_state":"INTEGRADO_EN_GRAFO_ACTIVO","proof_strength":"PROGRAMA_ABIERTO","provenance":"ARQUITECTURA_AUTORAL_PREEXISTENTE","reinforcement_state":"DEPLOYED_WITHOUT_NEW_RESULT_ID","residence":"perspectivas, § Universalidad genealógicamente fiel, tras UNIIMP-001","result_id":"P32","statement":"La versión HMT exige un funtor entre categorías de rutas que preserve composición, retorno y memoria, además de los datos del marco target--context importado.","status":"PRESENT_IN_SUCCESSOR_SOURCE"}
% HMT_RESULT_END:P32:residence
% HMT_RESULT_BEGIN:QPM-001:residence
% HMT_ATOMIC_RESULT {"material_state":"INTEGRADO_EN_GRAFO_ACTIVO","proof_strength":"EXACTO_INTERNO","provenance":"FORMALIZACION_NUEVA","reinforcement_state":"DEPLOYED_WITHOUT_NEW_RESULT_ID","residence":"apertura del corredor matricial, tras P09","result_id":"QPM-001","statement":"Una QPM es un QMS cuyas celdas son proyecciones; es conmutativa sólo si todas esas proyecciones conmutan conjuntamente. La definición no se obtiene dilatando cada fila por separado.","status":"PRESENT_IN_SUCCESSOR_SOURCE"}
% HMT_RESULT_END:QPM-001:residence
% HMT_RESULT_BEGIN:QLS-001:residence
% HMT_ATOMIC_RESULT {"material_state":"INTEGRADO_EN_GRAFO_ACTIVO","proof_strength":"EXACTO_INTERNO","provenance":"FORMALIZACION_NUEVA","reinforcement_state":"DEPLOYED_WITHOUT_NEW_RESULT_ID","residence":"apertura del corredor matricial, antes de QLSI-001","result_id":"QLS-001","statement":"Un QLS es una cuadrícula de vectores unitarios v_ij en C^n cuyas filas y columnas son bases ortonormales; las proyecciones P_ij=|v_ij><v_ij| forman un QMS proyectivo. Un QLS fácil se obtiene de un cuadrado latino clásico y una base ortonormal fija.","status":"PRESENT_IN_SUCCESSOR_SOURCE"}
% HMT_RESULT_END:QLS-001:residence
% HMT_RESULT_BEGIN:QLSI-001:residence
% HMT_ATOMIC_RESULT {"material_state":"INTEGRADO_EN_GRAFO_ACTIVO","proof_strength":"EXACTO_INTERNO","provenance":"RESULTADO_RECUPERADO","reinforcement_state":"DEPLOYED_WITHOUT_NEW_RESULT_ID","residence":"apertura del corredor matricial, tras P11","result_id":"QLSI-001","statement":"La envolvente matricial convexa de los QLS contiene estrictamente más que la clase semiclásica.","status":"PRESENT_IN_SUCCESSOR_SOURCE"}
% HMT_RESULT_END:QLSI-001:residence
% HMT_RESULT_BEGIN:QMSI-001:residence
% HMT_ATOMIC_RESULT {"material_state":"INTEGRADO_EN_GRAFO_ACTIVO","proof_strength":"EXACTO_INTERNO","provenance":"RESULTADO_RECUPERADO","reinforcement_state":"DEPLOYED_WITHOUT_NEW_RESULT_ID","residence":"apertura del corredor matricial, tras P10","result_id":"QMSI-001","statement":"Un QMS es semiclásico si y sólo si admite una dilatación simultánea a una QPM conmutativa mediante una única compresión común.","status":"PRESENT_IN_SUCCESSOR_SOURCE"}
% HMT_RESULT_END:QMSI-001:residence
% HMT_RESULT_BEGIN:QMSI-002:residence
% HMT_ATOMIC_RESULT {"material_state":"INTEGRADO_EN_GRAFO_ACTIVO","proof_strength":"EXACTO_INTERNO","provenance":"RESULTADO_RECUPERADO","reinforcement_state":"DEPLOYED_WITHOUT_NEW_RESULT_ID","residence":"apertura matricial, § Existencia de QMS no semiclásicos, tras QMSI-001","result_id":"QMSI-002","statement":"Para n≥3 y s≥2 existen QMS no semiclásicos.","status":"PRESENT_IN_SUCCESSOR_SOURCE"}
% HMT_RESULT_END:QMSI-002:residence
% HMT_RESULT_BEGIN:QMSI-003:residence
% HMT_ATOMIC_RESULT {"material_state":"INTEGRADO_EN_GRAFO_ACTIVO","proof_strength":"EXACTO_INTERNO","provenance":"RESULTADO_RECUPERADO","reinforcement_state":"DEPLOYED_WITHOUT_NEW_RESULT_ID","residence":"apertura del corredor matricial, tras QMSI-002","result_id":"QMSI-003","statement":"Para todo n≥3 existen QMS sobre M_2(C) que no admiten dilatación a ninguna QPM.","status":"PRESENT_IN_SUCCESSOR_SOURCE"}
% HMT_RESULT_END:QMSI-003:residence
% HMT_RESULT_BEGIN:QBR-001:residence
% HMT_ATOMIC_RESULT {"material_state":"INTEGRADO_EN_GRAFO_ACTIVO","proof_strength":"EXACTO_INTERNO","provenance":"RESULTADO_RECUPERADO","reinforcement_state":"DEPLOYED_WITHOUT_NEW_RESULT_ID","residence":"`md/04ab`, después de la realización hilbertiana finita","result_id":"QBR-001","statement":"Los observables se representan por operadores autoadjuntos y sus PVM; para un suceso cilíndrico P y estado normalizado Psi, la probabilidad es ||P Psi||^2.","status":"PRESENT_IN_SUCCESSOR_SOURCE"}
% HMT_RESULT_END:QBR-001:residence
% HMT_RESULT_BEGIN:QDY-001:residence
% HMT_ATOMIC_RESULT {"material_state":"INTEGRADO_EN_GRAFO_ACTIVO","proof_strength":"EXACTO_INTERNO","provenance":"RESULTADO_RECUPERADO","reinforcement_state":"DEPLOYED_WITHOUT_NEW_RESULT_ID","residence":"`md/04ab`, dinámica cuántica","result_id":"QDY-001","statement":"Un grupo unitario fuertemente continuo posee un generador autoadjunto H y satisface la ecuación de Schrödinger en su dominio.","status":"PRESENT_IN_SUCCESSOR_SOURCE"}
% HMT_RESULT_END:QDY-001:residence
% HMT_RESULT_BEGIN:QKR-001:residence
% HMT_ATOMIC_RESULT {"material_state":"INTEGRADO_EN_GRAFO_ACTIVO","proof_strength":"EXACTO_INTERNO","provenance":"RESULTADO_RECUPERADO","reinforcement_state":"DEPLOYED_WITHOUT_NEW_RESULT_ID","residence":"`md/01a`, instrumentos y canales","result_id":"QKR-001","statement":"Una familia de operadores de Kraus con suma K_a^*K_a=I define un canal completamente positivo y traza-preservante; cada resultado conserva su efecto e instrumento tipados.","status":"PRESENT_IN_SUCCESSOR_SOURCE"}
% HMT_RESULT_END:QKR-001:residence
% HMT_RESULT_BEGIN:QMO-001:residence
% HMT_ATOMIC_RESULT {"material_state":"INTEGRADO_EN_GRAFO_ACTIVO","proof_strength":"EXACTO_INTERNO","provenance":"RESULTADO_RECUPERADO","reinforcement_state":"DEPLOYED_WITHOUT_NEW_RESULT_ID","residence":"`md/04ab`, composición de registros","result_id":"QMO-001","statement":"Una realización monoidal fuerte transporta la composición independiente de registros al producto tensorial, con asociatividad, unidad y coherencias declaradas.","status":"PRESENT_IN_SUCCESSOR_SOURCE"}
% HMT_RESULT_END:QMO-001:residence
% HMT_RESULT_BEGIN:DIC-001:residence
% HMT_ATOMIC_RESULT {"material_state":"INTEGRADO_EN_GRAFO_ACTIVO","proof_strength":"EXACTO_INTERNO","provenance":"FORMALIZACION_NUEVA","reinforcement_state":"DEPLOYED_WITHOUT_NEW_RESULT_ID","residence":"apertura del corredor matricial","result_id":"DIC-001","statement":"El diccionario conserva exactamente: tamaño exterior n↔cuadrícula publicada; tamaño interior s↔carta operatorial; profundidad m↔genealogía; POVM↔familia de lectores positivos; PVM↔contexto espectral; UCP↔lector matricial; compresión↔sombra; dilatación↔antecedente operatorio; semiclasicidad↔POVM primaria común reordenada; no semiclasicidad↔inexistencia de una POVM primaria común que reproduzca todas las celdas; sistema de operadores↔espacio de lectores; rango matricial↔sombras a todos los s; límite inductivo↔cierre de refinamientos operatorios; y límite inverso↔sección infinita compatible. La composición genuinamente bidimensional es la interpretación HMT propuesta, no la definición importada.","status":"PRESENT_IN_SUCCESSOR_SOURCE"}
% HMT_RESULT_END:DIC-001:residence
% HMT_RESULT_BEGIN:QMC-001:residence
% HMT_ATOMIC_RESULT {"material_state":"INTEGRADO_EN_GRAFO_ACTIVO","proof_strength":"EXACTO_INTERNO","provenance":"RESULTADO_RECUPERADO","reinforcement_state":"DEPLOYED_WITHOUT_NEW_RESULT_ID","residence":"capítulo de cuadrados cuánticos, antes de las instancias \\((9,3)\\)","result_id":"QMC-001","statement":"A es un QMS (n,s), admite la descomposición semiclásica A=sum_t P_{sigma_t}⊗E_t y toda dilatación Naimark común E_t=V*Q_tV induce un QPM conmutativo simultáneo.","status":"PRESENT_IN_SUCCESSOR_SOURCE"}
% HMT_RESULT_END:QMC-001:residence
% HMT_RESULT_BEGIN:QDT-001:residence
% HMT_ATOMIC_RESULT {"material_state":"INTEGRADO_EN_GRAFO_ACTIVO","proof_strength":"EXACTO_INTERNO","provenance":"FORMALIZACION_NUEVA","reinforcement_state":"DEPLOYED_WITHOUT_NEW_RESULT_ID","residence":"geometría qutrit, antes de la fibración \\(12\\to4\\times3\\)","result_id":"QDT-001","statement":"M_3(C)=CI⊕direct_sum_{a∈P^1(F_3)}D_a^0 ortogonalmente para Hilbert--Schmidt y dim=9=1+4·2.","status":"PRESENT_IN_SUCCESSOR_SOURCE"}
% HMT_RESULT_END:QDT-001:residence
% HMT_RESULT_BEGIN:ENV-001:residence
% HMT_ATOMIC_RESULT {"material_state":"INTEGRADO_EN_GRAFO_ACTIVO","proof_strength":"EXACTO_INTERNO","provenance":"FORMALIZACION_NUEVA","reinforcement_state":"DEPLOYED_WITHOUT_NEW_RESULT_ID","residence":"dilatación cilíndrica, después de P16","result_id":"ENV-001","statement":"Sumar una representación fiel pi_f a una dilatación conserva todas las compresiones finitas y hace fiel la representación ambiental; esto no prueba todavía fidelidad dinámica TPK.","status":"PRESENT_IN_SUCCESSOR_SOURCE"}
% HMT_RESULT_END:ENV-001:residence
% HMT_RESULT_BEGIN:RNG-001:residence
% HMT_ATOMIC_RESULT {"material_state":"INTEGRADO_EN_GRAFO_ACTIVO","proof_strength":"EXACTO_INTERNO","provenance":"RESULTADO_NUEVO","reinforcement_state":"DEPLOYED_WITHOUT_NEW_RESULT_ID","residence":"rango matricial APP, después de P24","result_id":"RNG-001","statement":"(X,Y) pertenece a W_s(APP_2) exactamente cuando existen cuatro efectos escalares E_ab y dos mapas CP Psi_{5/7},Psi_{1/3}:M_2→M_s que satisfacen las ecuaciones completas de normalización y coordenadas.","status":"PRESENT_IN_SUCCESSOR_SOURCE"}
% HMT_RESULT_END:RNG-001:residence
% HMT_RESULT_BEGIN:DYN-001:residence
% HMT_ATOMIC_RESULT {"material_state":"INTEGRADO_EN_GRAFO_ACTIVO","proof_strength":"EXACTO_INTERNO","provenance":"FORMALIZACION_NUEVA","reinforcement_state":"DEPLOYED_WITHOUT_NEW_RESULT_ID","residence":"continuo operatorial, después de Stinespring","result_id":"DYN-001","statement":"La compatibilidad UCP a toda profundidad no implica ancilla finito uniforme; la fidelidad dinámica exige J_{m+1}T_m^TPK=tilde T_m J_m y naturalidad de las compresiones.","status":"PRESENT_IN_SUCCESSOR_SOURCE"}
% HMT_RESULT_END:DYN-001:residence
% HMT_RESULT_BEGIN:TEN-001:residence
% HMT_ATOMIC_RESULT {"material_state":"INTEGRADO_EN_GRAFO_ACTIVO","proof_strength":"PROGRAMA_ABIERTO","provenance":"ARQUITECTURA_AUTORAL_PREEXISTENTE","reinforcement_state":"DEPLOYED_WITHOUT_NEW_RESULT_ID","residence":"perspectivas del continuo tensorial","result_id":"TEN-001","statement":"El programa exige declarar tensor local, grupo de simetría, régimen exacto/aproximado/asintótico y estabilidad de rango; propone una torre coherente de raíces E_{kℓ}^ℓ=E_k equipada con registros TPK.","status":"PRESENT_IN_SUCCESSOR_SOURCE"}
% HMT_RESULT_END:TEN-001:residence
% HMT_RESULT_BEGIN:CHP-001:residence
% HMT_ATOMIC_RESULT {"material_state":"INTEGRADO_EN_GRAFO_ACTIVO","proof_strength":"PROGRAMA_ABIERTO","provenance":"FORMALIZACION_NUEVA","reinforcement_state":"DEPLOYED_WITHOUT_NEW_RESULT_ID","residence":"sistemas cónicos, antes de universalidad","result_id":"CHP-001","statement":"Las condiciones sum_j J(Phi_ij)=J(id) y sum_i J(Phi_ij)=J(id), con J(Phi_ij) positivo, linealizan la normalización mágica; la fidelidad de hoja y ruta es una obligación adicional.","status":"PRESENT_IN_SUCCESSOR_SOURCE"}
% HMT_RESULT_END:CHP-001:residence
% HMT_RESULT_BEGIN:SDP-001:residence
% HMT_ATOMIC_RESULT {"material_state":"INTEGRADO_EN_GRAFO_ACTIVO","proof_strength":"EXACTO_INTERNO","provenance":"FORMALIZACION_NUEVA","reinforcement_state":"DEPLOYED_WITHOUT_NEW_RESULT_ID","residence":"cierre del capítulo qutrit","result_id":"SDP-001","statement":"Se buscan Q_sigma positivos con sum_sigma Q_sigma=I y A_ij=sum_{sigma: sigma(i)=j}Q_sigma; factibilidad certifica semiclasicidad e inviabilidad requiere testigo dual.","status":"PRESENT_IN_SUCCESSOR_SOURCE"}
% HMT_RESULT_END:SDP-001:residence
% HMT_RESULT_BEGIN:EQV-001:residence
% HMT_ATOMIC_RESULT {"material_state":"INTEGRADO_EN_GRAFO_ACTIVO","proof_strength":"EXACTO_INTERNO","provenance":"FORMALIZACION_NUEVA","reinforcement_state":"DEPLOYED_WITHOUT_NEW_RESULT_ID","residence":"geometría simpléctica y lectores, después del calibre \\(3/\\pi\\)","result_id":"EQV-001","statement":"El calibre preliminar falla un test antiunitario y no es aún entrelazador; la promoción exige Xi(g·ell)=U_g Xi(ell)U_g* en todas las fibras, formulada como naturalidad con origen y destino cuando g es una flecha.","status":"PRESENT_IN_SUCCESSOR_SOURCE"}
% HMT_RESULT_END:EQV-001:residence
% HMT_RESULT_BEGIN:QMT-001:residence
% HMT_ATOMIC_RESULT {"material_state":"INTEGRADO_EN_GRAFO_ACTIVO","proof_strength":"PROGRAMA_ABIERTO","provenance":"FORMALIZACION_NUEVA","reinforcement_state":"DEPLOYED_WITHOUT_NEW_RESULT_ID","residence":"cierre del capítulo de cuadrados cuánticos APP","result_id":"QMT-001","statement":"La tabla conserva tres filas distintas: A^Sigma (9,3), que cierra hoja suma, positividad, normalización y semiclasicidad pero no hoja producto, cociente ni rutas; A^(9) (9,3), que cierra tres contextos, no conmutatividad y dilatación común pero no cuarto contexto, incidencia Hensel, gnomon ni memoria TPK; y B (12,3), que cierra el atlas 12→4×3 y su dilatación común pero no equivarianza, dinámica ni monodromía.","status":"PRESENT_IN_SUCCESSOR_SOURCE"}
% HMT_RESULT_END:QMT-001:residence
% HMT_RESULT_BEGIN:TRN-001:residence
% HMT_ATOMIC_RESULT {"material_state":"INTEGRADO_EN_GRAFO_ACTIVO","proof_strength":"PROGRAMA_ABIERTO","provenance":"ARQUITECTURA_AUTORAL_PREEXISTENTE","reinforcement_state":"DEPLOYED_WITHOUT_NEW_RESULT_ID","residence":"cierre metrológico del corredor matricial","result_id":"TRN-001","statement":"Conserva ocho filas completas: (1) monodromía 9→1→retorno del observable con estado ambiental distinto→clasificar dilataciones que preservan holonomía; (2) pi,e,phi,alpha→cuatro lectores compatibles de una genealogía→existencia y minimalidad de un sistema UCP común; (3) doble proyección→dos marginales de un antecedente cilíndrico→antecedente no conmutativo y separación conjunta de fibras; (4) B_c y salto 4→intervalo matricial libre [5I,9I]→extremos y transporte de la jerarquía 4;2;6;54; (5) Planck y CT108→punto autodual y lectores conjugados→realización como simetría de un sistema cónico con memoria; (6) ley de masas→familia positiva sobre fibras y rutas→rango matricial, extremos y realizaciones por sector; (7) CKM/CP→transporte de rango tres y fase no trivial→canal covariante que preserve espectro y orientación; (8) simetría excepcional→segundo codominio de la doble proyección→entrelazador entre atlas nonádico, código y representación.","status":"PRESENT_IN_SUCCESSOR_SOURCE"}
% HMT_RESULT_END:TRN-001:residence
% HMT_RESULT_BEGIN:MPSH-001:residence
% HMT_ATOMIC_RESULT {"material_state":"INTEGRADO_EN_GRAFO_ACTIVO","proof_strength":"PROGRAMA_ABIERTO","provenance":"ARQUITECTURA_AUTORAL_PREEXISTENTE","reinforcement_state":"DEPLOYED_WITHOUT_NEW_RESULT_ID","residence":"inmediatamente después de P28","result_id":"MPSH-001","statement":"Aplicar el teorema MPS a HMT exige construir una celda local uniforme, tensor A, canal E_A, forma irreducible, divisibilidad infinita propia o de una potencia y transporte al límite de hoja, acarreo y monodromía.","status":"PRESENT_IN_SUCCESSOR_SOURCE"}
% HMT_RESULT_END:MPSH-001:residence
% HMT_RESULT_BEGIN:UNIIMP-001:residence
% HMT_ATOMIC_RESULT {"material_state":"INTEGRADO_EN_GRAFO_ACTIVO","proof_strength":"EXACTO_INTERNO","provenance":"RESULTADO_RECUPERADO","reinforcement_state":"DEPLOYED_WITHOUT_NEW_RESULT_ID","residence":"perspectivas operatoriales, antes de P32","result_id":"UNIIMP-001","statement":"Un simulador s:P⊗C→T⊗C es universal cuando existe la reducción laxa prescrita al simulador trivial; un monótono compatible cuyo supremo sobre la imagen funcional quede por debajo de algún target constituye un no-go.","status":"PRESENT_IN_SUCCESSOR_SOURCE"}
% HMT_RESULT_END:UNIIMP-001:residence
% HMT_RESULT_BEGIN:PRG-001:residence
% HMT_ATOMIC_RESULT {"material_state":"INTEGRADO_EN_GRAFO_ACTIVO","proof_strength":"PROGRAMA_ABIERTO","provenance":"ARQUITECTURA_AUTORAL_PREEXISTENTE","reinforcement_state":"DEPLOYED_WITHOUT_NEW_RESULT_ID","residence":"cierre del corredor matricial","result_id":"PRG-001","statement":"Los cierres son: sistema APP filtrado, QMS APP--TPK fiel, frontera semiclásica, canal de continuo, sistema cónico intrínseco y realización metrológica UCP conjunta; ninguno se promueve sin dominio, prueba/certificado, positividad, compatibilidad, falsadores y trazabilidad.","status":"PRESENT_IN_SUCCESSOR_SOURCE"}
% HMT_RESULT_END:PRG-001:residence
% HMT_RESULT_BEGIN:QCT-001:residence
% HMT_ATOMIC_RESULT {"material_state":"CONTROL_FOCAL_ACTIVO","proof_strength":"EXACTO_INTERNO","provenance":"CERTIFICADO_NUEVO","reinforcement_state":"DEPLOYED_WITHOUT_NEW_RESULT_ID","residence":"apéndice digital y gate del puente","result_id":"QCT-001","statement":"Registra exactamente dieciocho controles: PVM qutrit; idempotencia; MUB; no-go QLS; QMS aditivo (9,3); QMS de tres contextos; no conmutatividad; QMS atlas (12,3); descomposición semiclásica; cubos de órdenes 9 y 12; P^1(Z/9)→P^1(F3); truncación--compresión; lector 1/pi; calibre 3/pi; polinomio de Gram; intersecciones de Halmos; y C*-álgebra APP d=2.","status":"CONTROLLED"}
% HMT_RESULT_END:QCT-001:residence
% HMT_RESULT_BEGIN:PRV-001:residence
% HMT_ATOMIC_RESULT {"material_state":"CONTROL_FOCAL_ACTIVO","proof_strength":"EXACTO_INTERNO","provenance":"FORMALIZACION_NUEVA","reinforcement_state":"DEPLOYED_WITHOUT_NEW_RESULT_ID","residence":"registro bibliográfico del corredor matricial","result_id":"PRV-001","statement":"Conserva autores, firmas e identificadores exactamente: De las Cuevas--Drescher--Netzer, Quantum magic squares: dilations and their limitations—DOI 10.1063/5.0022344—QMS/QPM/dilatación; Magic squares: Latin, Semiclassical and Quantum, De las Cuevas--Netzer--Valentiner-Branth—DOI 10.1063/5.0127393—QLS/envolventes; Quantum information theory and free semialgebraic geometry, De las Cuevas--Netzer—arXiv:2102.04240—torres/Choi; Beyond Operator Systems, De les Coves--Mirte van der Eyden--Netzer—DOI 10.1016/j.jmaa.2025.130102—sistemas cónicos; Continuum limits of Matrix Product States, De las Cuevas--Schuch--Pérez-García--Cirac—DOI 10.1103/PhysRevB.98.174303—divisibilidad; Tensor decompositions on simplicial complexes with invariance, De las Cuevas--Hoogsteder Riera--Netzer—DOI 10.1016/j.jsc.2024.102299—simetría; Approximate tensor decompositions, De las Cuevas--Klingler--Netzer—DOI 10.1063/5.0033876—aproximación; A Framework for Universality in Physics, Computer Science, and Beyond, Gonda--Reinhart--Stengele--De les Coves—DOI 10.46298/compositionality-6-3—target--context. Separa literatura simpléctica vecina y fija De las Cuevas/De les Coves.","status":"CONTROLLED"}
% HMT_RESULT_END:PRV-001:residence
% HMT_RESULT_BEGIN:MOR-001:residence
% HMT_ATOMIC_RESULT {"material_state":"INTEGRADO_O_REMITIDO_SIN_PERDIDA","proof_strength":"EXACTO_INTERNO","provenance":"ARQUITECTURA_AUTORAL_PREEXISTENTE","reinforcement_state":"DEPLOYED_WITHOUT_NEW_RESULT_ID","residence":"apertura continuo y `md/00b`","result_id":"MOR-001","statement":"Una cifra es una coordenada finita publicada; un valor es la imagen arquimediana, espectral o física de una sección; un número HMT es la sección completa con profundidad, hoja, orientación, acarreo, ruta, frontera, memoria y supervivencia; una simetría HMT conserva la identidad enriquecida y transporta su memoria.","status":"PRESENT_IN_SUCCESSOR_SOURCE"}
% HMT_RESULT_END:MOR-001:residence
% HMT_RESULT_BEGIN:MOR-002:residence
% HMT_ATOMIC_RESULT {"material_state":"INTEGRADO_O_REMITIDO_SIN_PERDIDA","proof_strength":"EXACTO_INTERNO","provenance":"RESULTADO_RECUPERADO","reinforcement_state":"DEPLOYED_WITHOUT_NEW_RESULT_ID","residence":"APP","result_id":"MOR-002","statement":"suma/producto y memoria de fibras","status":"PRESENT_IN_SUCCESSOR_SOURCE"}
% HMT_RESULT_END:MOR-002:residence
% HMT_RESULT_BEGIN:MOR-003:residence
% HMT_ATOMIC_RESULT {"material_state":"INTEGRADO_O_REMITIDO_SIN_PERDIDA","proof_strength":"EXACTO_INTERNO","provenance":"ARQUITECTURA_AUTORAL_PREEXISTENTE","reinforcement_state":"DEPLOYED_WITHOUT_NEW_RESULT_ID","residence":"TRIT","result_id":"MOR-003","statement":"orientación temporal y acarreo","status":"PRESENT_IN_SUCCESSOR_SOURCE"}
% HMT_RESULT_END:MOR-003:residence
% HMT_RESULT_BEGIN:MOR-004:residence
% HMT_ATOMIC_RESULT {"material_state":"INTEGRADO_O_REMITIDO_SIN_PERDIDA","proof_strength":"EXACTO_INTERNO","provenance":"FORMALIZACION_NUEVA","reinforcement_state":"DEPLOYED_WITHOUT_NEW_RESULT_ID","residence":"TPK","result_id":"MOR-004","statement":"composición con memoria","status":"PRESENT_IN_SUCCESSOR_SOURCE"}
% HMT_RESULT_END:MOR-004:residence
% HMT_RESULT_BEGIN:MOR-005:residence
% HMT_ATOMIC_RESULT {"material_state":"INTEGRADO_O_REMITIDO_SIN_PERDIDA","proof_strength":"EXACTO_INTERNO","provenance":"RESULTADO_RECUPERADO","reinforcement_state":"DEPLOYED_WITHOUT_NEW_RESULT_ID","residence":"continuo, tras `hmt/00b`","result_id":"MOR-005","statement":"9^n 1_C mu -> delta_x","status":"PRESENT_IN_SUCCESSOR_SOURCE"}
% HMT_RESULT_END:MOR-005:residence
% HMT_RESULT_BEGIN:MOR-006:residence
% HMT_ATOMIC_RESULT {"material_state":"INTEGRADO_O_REMITIDO_SIN_PERDIDA","proof_strength":"EXACTO_INTERNO","provenance":"FORMALIZACION_NUEVA","reinforcement_state":"DEPLOYED_WITHOUT_NEW_RESULT_ID","residence":"bisagra continuo--cuántica, § Kronecker--cilindro--delta de Dirac, tras MOR-005","result_id":"MOR-006","statement":"proyectores condicionales y núcleo","status":"PRESENT_IN_SUCCESSOR_SOURCE"}
% HMT_RESULT_END:MOR-006:residence
% HMT_RESULT_BEGIN:MOR-007:residence
% HMT_ATOMIC_RESULT {"material_state":"INTEGRADO_O_REMITIDO_SIN_PERDIDA","proof_strength":"EXACTO_INTERNO","provenance":"RESULTADO_RECUPERADO","reinforcement_state":"DEPLOYED_WITHOUT_NEW_RESULT_ID","residence":"UHF--\\(II_1\\) y `md/04ab`","result_id":"MOR-007","statement":"La inclusión unital nonádica genera UHF(9^infty), canónicamente isomorfa a UHF(3^infty); la traza GNS conduce al factor hiperinfinito II_1. El cuaderno de 64 páginas conserva la prueba propietaria completa.","status":"PRESENT_IN_SUCCESSOR_SOURCE"}
% HMT_RESULT_END:MOR-007:residence
% HMT_RESULT_BEGIN:MOR-008:residence
% HMT_ATOMIC_RESULT {"material_state":"INTEGRADO_O_REMITIDO_SIN_PERDIDA","proof_strength":"EXACTO_INTERNO","provenance":"RESULTADO_RECUPERADO","reinforcement_state":"DEPLOYED_WITHOUT_NEW_RESULT_ID","residence":"`md/01a`, `04ab`","result_id":"MOR-008","statement":"Born--Lüders--Kraus","status":"PRESENT_IN_SUCCESSOR_SOURCE"}
% HMT_RESULT_END:MOR-008:residence
% HMT_RESULT_BEGIN:MOR-009:residence
% HMT_ATOMIC_RESULT {"material_state":"INTEGRADO_O_REMITIDO_SIN_PERDIDA","proof_strength":"INTERFAZ_FISICA","provenance":"RESULTADO_RECUPERADO","reinforcement_state":"DEPLOYED_WITHOUT_NEW_RESULT_ID","residence":"`md/04ab`, `00b`","result_id":"MOR-009","statement":"Q:P_TPK→Hilb_{R,J}; la extensión a flechas inversas se afirma sólo en el subdominio reversible donde el transporte unitario fuente teoremática la demuestra","status":"PRESENT_IN_SUCCESSOR_SOURCE"}
% HMT_RESULT_END:MOR-009:residence
% HMT_RESULT_BEGIN:MOR-010:residence
% HMT_ATOMIC_RESULT {"material_state":"INTEGRADO_O_REMITIDO_SIN_PERDIDA","proof_strength":"PROGRAMA_ABIERTO","provenance":"FORMALIZACION_NUEVA","reinforcement_state":"DEPLOYED_WITHOUT_NEW_RESULT_ID","residence":"después de `md/00b`, antes de la ley dimensional/Planck","result_id":"MOR-010","statement":"R_MD^{disc,+} es el candidato de producto con blanco Hilbert general y sólo queda construido tras demostrar identidad y composición componente a componente; U_reg existe sólo en el subdominio unitario; R_MD^{disc,or} existe únicamente cuando espectro e Hilbert son invertibles, las relaciones localizadas se respetan y la dimensional toma valores en Gr(M_D^+); positiva y orientada son lectoras hermanas; no se presupone ni exige una factorización horizontal entre ellas. El monoide de coeficientes satisface (n,a,b,lambda)(n',a',b',lambda')=(n+n',a+a',b+b',lambda lambda') y sólo después se aplica la completación de Grothendieck.","status":"PRESENT_IN_SUCCESSOR_SOURCE"}
% HMT_RESULT_END:MOR-010:residence
% HMT_RESULT_BEGIN:MOR-011:residence
% HMT_ATOMIC_RESULT {"material_state":"INTEGRADO_O_REMITIDO_SIN_PERDIDA","proof_strength":"EXACTO_INTERNO","provenance":"RESULTADO_RECUPERADO","reinforcement_state":"DEPLOYED_WITHOUT_NEW_RESULT_ID","residence":"`md/11m`","result_id":"MOR-011","statement":"C U(gamma) C^-1=U(gamma^-1)","status":"PRESENT_IN_SUCCESSOR_SOURCE"}
% HMT_RESULT_END:MOR-011:residence
% HMT_RESULT_BEGIN:MOR-012:residence
% HMT_ATOMIC_RESULT {"material_state":"INTEGRADO_O_REMITIDO_SIN_PERDIDA","proof_strength":"INTERFAZ_FISICA","provenance":"ARQUITECTURA_AUTORAL_PREEXISTENTE","reinforcement_state":"DEPLOYED_WITHOUT_NEW_RESULT_ID","residence":"`md/11m`","result_id":"MOR-012","statement":"retardo/adelanto y borde","status":"PRESENT_IN_SUCCESSOR_SOURCE"}
% HMT_RESULT_END:MOR-012:residence
% HMT_RESULT_BEGIN:MOR-013:residence
% HMT_ATOMIC_RESULT {"material_state":"INTEGRADO_O_REMITIDO_SIN_PERDIDA","proof_strength":"EXACTO_INTERNO","provenance":"RESULTADO_RECUPERADO","reinforcement_state":"DEPLOYED_WITHOUT_NEW_RESULT_ID","residence":"`md/04c`","result_id":"MOR-013","statement":"N_sp(uP_++vP_-)=uv","status":"PRESENT_IN_SUCCESSOR_SOURCE"}
% HMT_RESULT_END:MOR-013:residence
% HMT_RESULT_BEGIN:MOR-014:residence
% HMT_ATOMIC_RESULT {"material_state":"INTEGRADO_O_REMITIDO_SIN_PERDIDA","proof_strength":"INTERFAZ_FISICA","provenance":"RESULTADO_RECUPERADO","reinforcement_state":"DEPLOYED_WITHOUT_NEW_RESULT_ID","residence":"`md/04c`, `06ka`","result_id":"MOR-014","statement":"producto positivo/interior","status":"PRESENT_IN_SUCCESSOR_SOURCE"}
% HMT_RESULT_END:MOR-014:residence
% HMT_RESULT_BEGIN:MOR-015:residence
% HMT_ATOMIC_RESULT {"material_state":"INTEGRADO_O_REMITIDO_SIN_PERDIDA","proof_strength":"PROGRAMA_ABIERTO","provenance":"RESULTADO_NUEVO","reinforcement_state":"DEPLOYED_WITHOUT_NEW_RESULT_ID","residence":"`md/04ac` y persistencia","result_id":"MOR-015","statement":"gamma_+ es fibra de Hopf y gamma_- se proyecta con grado dos; sus clases (1,1) y (1,-1) dan las familias de Villarceau. La conjugación compleja total preserva cada familia no orientada y revierte su orientación: no intercambia ambas familias.","status":"PRESENT_IN_SUCCESSOR_SOURCE"}
% HMT_RESULT_END:MOR-015:residence
% HMT_RESULT_BEGIN:MOR-016:residence
% HMT_ATOMIC_RESULT {"material_state":"INTEGRADO_O_REMITIDO_SIN_PERDIDA","proof_strength":"PROGRAMA_ABIERTO","provenance":"RESULTADO_RECUPERADO","reinforcement_state":"DEPLOYED_WITHOUT_NEW_RESULT_ID","residence":"/Users/ruben/Documents/HMT2/EDICION_AUTONOMA_HMT_MD_2026-08-09_REV_3_EN_TRABAJO/01_FUENTE_REVISADA/manuscrito/sections/md/21_persistencia_ruta_mobius_familias_actualizada.tex","result_id":"MOR-016","statement":"El doble retorno 2pi/4pi y el fibrado real de línea L_chi con chi(1)=-1 son exactos; un círculo aislado no es Möbius. Ascender el signo C_M a la holonomía del lazo Villarceau requiere un único entrelazador físico.","status":"PRESENT_IN_SUCCESSOR_SOURCE"}
% HMT_RESULT_END:MOR-016:residence
% HMT_RESULT_BEGIN:MOR-017:residence
% HMT_ATOMIC_RESULT {"material_state":"INTEGRADO_O_REMITIDO_SIN_PERDIDA","proof_strength":"EXACTO_INTERNO","provenance":"FORMALIZACION_NUEVA","reinforcement_state":"DEPLOYED_WITHOUT_NEW_RESULT_ID","residence":"`md/04fa`","result_id":"MOR-017","statement":"En el escenario qutrit declarado, las POVM completas producen una distribución no señalizante y la cota local es I_3≤2. Para el estado máximamente entrelazado, I_3=4(3+2sqrt(3))/9=2.872934…; para |Psi_gamma>=(|00>+gamma|11>+|22>)/sqrt(2+gamma^2), gamma=(sqrt(11)-sqrt(3))/2, se alcanza I_3=1+sqrt(11/3)=2.914854…. Por tanto, máxima entropía no equivale a máxima violación. Derivar el estado y los instrumentos íntegramente desde el registro HMT sigue siendo una obligación separada.","status":"PRESENT_IN_SUCCESSOR_SOURCE"}
% HMT_RESULT_END:MOR-017:residence
% HMT_RESULT_BEGIN:MOR-018:residence
% HMT_ATOMIC_RESULT {"material_state":"INTEGRADO_O_REMITIDO_SIN_PERDIDA","proof_strength":"INTERFAZ_FISICA","provenance":"RESULTADO_RECUPERADO","reinforcement_state":"DEPLOYED_WITHOUT_NEW_RESULT_ID","residence":"`md/10_constitutivas_si`","result_id":"MOR-018","statement":"mu,epsilon,Z,c","status":"PRESENT_IN_SUCCESSOR_SOURCE"}
% HMT_RESULT_END:MOR-018:residence
% HMT_RESULT_BEGIN:MOR-019:residence
% HMT_ATOMIC_RESULT {"material_state":"INTEGRADO_O_REMITIDO_SIN_PERDIDA","proof_strength":"INTERFAZ_FISICA","provenance":"RESULTADO_RECUPERADO","reinforcement_state":"DEPLOYED_WITHOUT_NEW_RESULT_ID","residence":"Planck y paquete de masas","result_id":"MOR-019","statement":"Sigma_H y 10^-34","status":"PRESENT_IN_SUCCESSOR_SOURCE"}
% HMT_RESULT_END:MOR-019:residence
% HMT_RESULT_BEGIN:MOR-020:residence
% HMT_ATOMIC_RESULT {"material_state":"INTEGRADO_O_REMITIDO_SIN_PERDIDA","proof_strength":"INTERFAZ_FISICA","provenance":"RESULTADO_RECUPERADO","reinforcement_state":"DEPLOYED_WITHOUT_NEW_RESULT_ID","residence":"`md/05aa`, `05f`, `06m`","result_id":"MOR-020","statement":"operadores K_D/K_Omega y torres","status":"PRESENT_IN_SUCCESSOR_SOURCE"}
% HMT_RESULT_END:MOR-020:residence
% HMT_RESULT_BEGIN:MOR-021:residence
% HMT_ATOMIC_RESULT {"material_state":"INTEGRADO_O_REMITIDO_SIN_PERDIDA","proof_strength":"INTERFAZ_FISICA","provenance":"RESULTADO_RECUPERADO","reinforcement_state":"DEPLOYED_WITHOUT_NEW_RESULT_ID","residence":"`md/06k`, `06ka`, `06m`","result_id":"MOR-021","statement":"basal->beta (80,54,6)","status":"PRESENT_IN_SUCCESSOR_SOURCE"}
% HMT_RESULT_END:MOR-021:residence
% HMT_RESULT_BEGIN:MOR-022:residence
% HMT_ATOMIC_RESULT {"material_state":"INTEGRADO_O_REMITIDO_SIN_PERDIDA","proof_strength":"EXACTO_INTERNO","provenance":"RESULTADO_RECUPERADO","reinforcement_state":"DEPLOYED_WITHOUT_NEW_RESULT_ID","residence":"`md/03d`, con remisiones a `06k` y `06ka`","result_id":"MOR-022","statement":"Con T_108=108t_0, omega_108=2pi/T_108, R_108=c/omega_108 y m_108=hbar omega_108/c^2, se cumple hbar/(m_108 c)=R_108 y h/(m_108 c)=2pi R_108=C_108.","status":"PRESENT_IN_SUCCESSOR_SOURCE"}
% HMT_RESULT_END:MOR-022:residence
% HMT_RESULT_BEGIN:MOR-023:residence
% HMT_ATOMIC_RESULT {"material_state":"INTEGRADO_O_REMITIDO_SIN_PERDIDA","proof_strength":"INTERFAZ_FISICA","provenance":"RESULTADO_RECUPERADO","reinforcement_state":"DEPLOYED_WITHOUT_NEW_RESULT_ID","residence":"nuevo teorema producto + `md/10c`","result_id":"MOR-023","statement":"holonomía, coframe, conexión","status":"PRESENT_IN_SUCCESSOR_SOURCE"}
% HMT_RESULT_END:MOR-023:residence
% HMT_RESULT_BEGIN:MOR-024:residence
% HMT_ATOMIC_RESULT {"material_state":"INTEGRADO_O_REMITIDO_SIN_PERDIDA","proof_strength":"INTERFAZ_FISICA","provenance":"RESULTADO_RECUPERADO","reinforcement_state":"DEPLOYED_WITHOUT_NEW_RESULT_ID","residence":"`md/11`, `11c`, `11p`","result_id":"MOR-024","statement":"A=(omega,e), F=(R,T)","status":"PRESENT_IN_SUCCESSOR_SOURCE"}
% HMT_RESULT_END:MOR-024:residence
% HMT_RESULT_BEGIN:MOR-025:residence
% HMT_ATOMIC_RESULT {"material_state":"INTEGRADO_O_REMITIDO_SIN_PERDIDA","proof_strength":"INTERFAZ_FISICA","provenance":"RESULTADO_RECUPERADO","reinforcement_state":"DEPLOYED_WITHOUT_NEW_RESULT_ID","residence":"`md/11m`","result_id":"MOR-025","statement":"Si b es el lector de base, todo observable constante en las fibras de b factoriza de manera única a través de b. Una semilla s_0 se transporta consistentemente alrededor de un lazo exactamente cuando s_0 pertenece a Fix(Hol). El dato de fibra asciende al estado TPK completo sólo si la realización empleada es inyectiva.","status":"PRESENT_IN_SUCCESSOR_SOURCE"}
% HMT_RESULT_END:MOR-025:residence
% HMT_RESULT_BEGIN:MOR-026:residence
% HMT_ATOMIC_RESULT {"material_state":"INTEGRADO_O_REMITIDO_SIN_PERDIDA","proof_strength":"EXACTO_INTERNO","provenance":"RESULTADO_RECUPERADO","reinforcement_state":"DEPLOYED_WITHOUT_NEW_RESULT_ID","residence":"Euler/09b/14f","result_id":"MOR-026","statement":"J^2=-I y z->iz","status":"PRESENT_IN_SUCCESSOR_SOURCE"}
% HMT_RESULT_END:MOR-026:residence
% HMT_RESULT_BEGIN:MOR-027:residence
% HMT_ATOMIC_RESULT {"material_state":"INTEGRADO_O_REMITIDO_SIN_PERDIDA","proof_strength":"EXACTO_INTERNO","provenance":"FORMALIZACION_NUEVA","reinforcement_state":"DEPLOYED_WITHOUT_NEW_RESULT_ID","residence":"/Users/ruben/Documents/HMT2/EDICION_AUTONOMA_HMT_MD_2026-08-09_REV_3_EN_TRABAJO/01_FUENTE_REVISADA/manuscrito/sections/hmt/14d_ontologia_numeros_rev11.tex","result_id":"MOR-027","statement":"valor regularizado de zeta(-1)","status":"PRESENT_IN_SUCCESSOR_SOURCE"}
% HMT_RESULT_END:MOR-027:residence
% HMT_RESULT_BEGIN:MOR-028:residence
% HMT_ATOMIC_RESULT {"material_state":"INTEGRADO_O_REMITIDO_SIN_PERDIDA","proof_strength":"EXACTO_INTERNO","provenance":"ARQUITECTURA_AUTORAL_PREEXISTENTE","reinforcement_state":"DEPLOYED_WITHOUT_NEW_RESULT_ID","residence":"`hmt/03f`, `04_cubo_emrh`","result_id":"MOR-028","statement":"doble evaluación del antecedente","status":"PRESENT_IN_SUCCESSOR_SOURCE"}
% HMT_RESULT_END:MOR-028:residence
% HMT_RESULT_BEGIN:MOR-029:residence
% HMT_ATOMIC_RESULT {"material_state":"INTEGRADO_O_REMITIDO_SIN_PERDIDA","proof_strength":"EXACTO_INTERNO","provenance":"RESULTADO_RECUPERADO","reinforcement_state":"DEPLOYED_WITHOUT_NEW_RESULT_ID","residence":"`hmt/06aa--06f`","result_id":"MOR-029","statement":"Para todo K, la misma regla particiona exhaustivamente las 729 prolongaciones, selecciona una única superviviente y conmuta con las restricciones; los cilindros compatibles anidados determinan una sección coinductiva única.","status":"PRESENT_IN_SUCCESSOR_SOURCE"}
% HMT_RESULT_END:MOR-029:residence
% HMT_RESULT_BEGIN:MOR-030:residence
% HMT_ATOMIC_RESULT {"material_state":"INTEGRADO_O_REMITIDO_SIN_PERDIDA","proof_strength":"EXACTO_INTERNO","provenance":"FORMALIZACION_NUEVA","reinforcement_state":"DEPLOYED_WITHOUT_NEW_RESULT_ID","residence":"junto al teorema producto y síntesis final","result_id":"MOR-030","statement":"Los falsadores exactos son: continuo, pérdida de masa unitaria o de convergencia débil; Hilbert, ausencia de PVM compatible o dinámica no unitaria; morfismo discreto, fallo de composición; conjugación, C^2≠I, CJC^{-1}≠-J o CHC^{-1}≠H; absorbedor, falta de operador de onda o de condiciones causales; constitutivas, inversión que no intercambia mu/epsilon o cambia c sin extensión; masa, selección por PDG o promoción condicionada; electrón, imposibilidad de obtener carga, espín o estabilidad del mismo objeto; Planck, otra solución positiva autodual; Cartan, dependencia de subdivisión o confusión con Delta_4; Bell, señalización o instrumentos escogidos tras el blanco; filtración nonádica, consulta de prefijos objetivo o pérdida de una región de pi; Euler/zeta, regularización tratada como suma ordinaria; cosmología, incapacidad de reproducir simultáneamente expansión, lentes y estructura.","status":"PRESENT_IN_SUCCESSOR_SOURCE"}
% HMT_RESULT_END:MOR-030:residence
% HMT_RESULT_BEGIN:PLS-001:residence
% HMT_ATOMIC_RESULT {"material_state":"INTEGRADO_O_REMITIDO_SIN_PERDIDA","proof_strength":"EXACTO_INTERNO","provenance":"RESULTADO_RECUPERADO","reinforcement_state":"DEPLOYED_WITHOUT_NEW_RESULT_ID","residence":"`md/03d`, subsección CT108--Compton","result_id":"PLS-001","statement":"Para G(theta)=theta c_int^5 t_0^2/hbar_int, la clausura circular selecciona de manera única theta_*=(54/pi)^2 y hace coincidir el radio CT108 con las escalas de Compton, gravitatoria y de Planck en la convención declarada. Se distinguen r_g=Gm/c^2 y r_s=2r_g.","status":"PRESENT_IN_SUCCESSOR_SOURCE"}
% HMT_RESULT_END:PLS-001:residence
% HMT_RESULT_BEGIN:MAD-001:residence
% HMT_ATOMIC_RESULT {"material_state":"INTEGRADO_O_REMITIDO_SIN_PERDIDA","proof_strength":"EXACTO_INTERNO","provenance":"RESULTADO_RECUPERADO","reinforcement_state":"DEPLOYED_WITHOUT_NEW_RESULT_ID","residence":"`md/03d`, recta autodual de masas","result_id":"MAD-001","statement":"Con P_±=(I±j)/2, Z(m)=(m_P/m)P_++(m/m_P)P_- satisface N_sp(Z(m))=1; la conjugación implementa m↦m_P^2/m, el punto fijo positivo es m=m_P y x=log(m/m_P) es la rapidez de la recta.","status":"PRESENT_IN_SUCCESSOR_SOURCE"}
% HMT_RESULT_END:MAD-001:residence
% HMT_RESULT_BEGIN:AGR-001:residence
% HMT_ATOMIC_RESULT {"material_state":"INTEGRADO_O_REMITIDO_SIN_PERDIDA","proof_strength":"DEMOSTRADO_CON_ESTRUCTURA_DE_PARTIDA_EXPLICITA","provenance":"RESULTADO_RECUPERADO","reinforcement_state":"DEPLOYED_WITHOUT_NEW_RESULT_ID","residence":"`md/03d`, covariancia pre/retorno","result_id":"AGR-001","statement":"Las ramas pre y retorno transforman recíprocamente acción y acoplamiento gravitatorio, mantienen invariante su producto geométrico y conservan la longitud característica construida.","status":"PRESENT_IN_SUCCESSOR_SOURCE"}
% HMT_RESULT_END:AGR-001:residence
% HMT_RESULT_BEGIN:PLE-001:residence
% HMT_ATOMIC_RESULT {"material_state":"INTEGRADO_O_REMITIDO_SIN_PERDIDA","proof_strength":"INTERFAZ_FISICA","provenance":"RESULTADO_RECUPERADO","reinforcement_state":"DEPLOYED_WITHOUT_NEW_RESULT_ID","residence":"`md/03d`, `06k`, `06ka`","result_id":"PLE-001","statement":"El punto autodual de Planck fija el origen de la coordenada; la genealogía electrónica determina un desplazamiento posterior y la clausura escalar de rango uno. Este cierre discreto no demuestra por sí solo carga, espín, estabilidad, factor de forma, momento magnético ni un Hamiltoniano electromagnético ligado.","status":"PRESENT_IN_SUCCESSOR_SOURCE"}
% HMT_RESULT_END:PLE-001:residence
% HMT_RESULT_BEGIN:ANG-001:residence
% HMT_ATOMIC_RESULT {"material_state":"INTEGRADO_O_REMITIDO_SIN_PERDIDA","proof_strength":"INTERFAZ_FISICA","provenance":"ARQUITECTURA_AUTORAL_PREEXISTENTE","reinforcement_state":"DEPLOYED_WITHOUT_NEW_RESULT_ID","residence":"Hopf/Villarceau, después de MOR-016","result_id":"ANG-001","statement":"Para epsilon=+1, Spec(B_mem^+)={7,3}; resolver R+r=7ell y R-r=3ell da exactamente (R,r)=(5ell,2ell) y sin(beta)=2/5. La lectura como bordes de toro es interfaz declarada; el bloque total conserva por separado la razón 7:2.","status":"PRESENT_IN_SUCCESSOR_SOURCE"}
% HMT_RESULT_END:ANG-001:residence
% HMT_RESULT_BEGIN:BDR-001:residence
% HMT_ATOMIC_RESULT {"material_state":"INTEGRADO_O_REMITIDO_SIN_PERDIDA","proof_strength":"INTERFAZ_FISICA","provenance":"RESULTADO_RECUPERADO","reinforcement_state":"DEPLOYED_WITHOUT_NEW_RESULT_ID","residence":"constitutivas, con remisión desde bidireccionalidad","result_id":"BDR-001","statement":"La inversión conserva c_hat=exp(-E), invierte Z_hat=exp(B) e intercambia mu_hat y epsilon_hat; la realización vigente tiene velocidad común en ambos sentidos e impedancia orientada. Velocidades distintas exigen una extensión anisótropa adicional.","status":"PRESENT_IN_SUCCESSOR_SOURCE"}
% HMT_RESULT_END:BDR-001:residence
% HMT_RESULT_BEGIN:OCC-001:residence
% HMT_ATOMIC_RESULT {"material_state":"INTEGRADO_O_REMITIDO_SIN_PERDIDA","proof_strength":"EXACTO_INTERNO","provenance":"RESULTADO_RECUPERADO","reinforcement_state":"DEPLOYED_WITHOUT_NEW_RESULT_ID","residence":"Gibbs--KMS y ocupación","result_id":"OCC-001","statement":"Gobierna la identidad tipada de REV.3 F_{A,B}(t+i beta)=Tr(rho alpha_t(B)A), no la abreviación errónea del PDF de 68 pp; para 0<C<I, K_1=log((I-C)C^{-1}) y C=(I+e^{K_1})^{-1}, con tipos CAR/CCR separados.","status":"PRESENT_IN_SUCCESSOR_SOURCE"}
% HMT_RESULT_END:OCC-001:residence
% HMT_RESULT_BEGIN:BOL-001:residence
% HMT_ATOMIC_RESULT {"material_state":"INTEGRADO_O_REMITIDO_SIN_PERDIDA","proof_strength":"EXACTO_INTERNO","provenance":"RESULTADO_RECUPERADO","reinforcement_state":"DEPLOYED_WITHOUT_NEW_RESULT_ID","residence":"`md/04e`, `04g` y síntesis, con remisión CT108","result_id":"BOL-001","statement":"Boltzmann actúa como k_B:L_Theta→L_E y, para el reloj CT108, satisface k_B Theta_clk log(3)=hbar_int omega_0 con omega_0=2pi/(108t_0). La identidad fija el producto acción--información; insertar la coordenada SI declarada de k_B sólo evalúa Theta_clk y no constituye una derivación autónoma de ese número SI.","status":"PRESENT_IN_SUCCESSOR_SOURCE"}
% HMT_RESULT_END:BOL-001:residence
% HMT_RESULT_BEGIN:COS-001:residence
% HMT_ATOMIC_RESULT {"material_state":"INTEGRADO_O_REMITIDO_SIN_PERDIDA","proof_strength":"PROGRAMA_ABIERTO","provenance":"ARQUITECTURA_AUTORAL_PREEXISTENTE","reinforcement_state":"DEPLOYED_WITHOUT_NEW_RESULT_ID","residence":"perspectivas cosmológicas","result_id":"COS-001","statement":"La coexistencia de cartas permite un modelo estratificado, pero la aceleración sólo queda explicada tras derivar una ecuación efectiva con signo observado y límites de curvatura.","status":"PRESENT_IN_SUCCESSOR_SOURCE"}
% HMT_RESULT_END:COS-001:residence
% HMT_RESULT_BEGIN:GWV-001:residence
% HMT_ATOMIC_RESULT {"material_state":"INTEGRADO_O_REMITIDO_SIN_PERDIDA","proof_strength":"PROGRAMA_ABIERTO","provenance":"RESULTADO_RECUPERADO","reinforcement_state":"DEPLOYED_WITHOUT_NEW_RESULT_ID","residence":"gravedad sin gravitón y espectro","result_id":"GWV-001","statement":"Linealizar y reducir por calibre puede producir dos grados transversos de una onda gravitatoria colectiva; esto no obliga ni excluye un gravitón elemental, cuya existencia es una pregunta espectral adicional.","status":"PRESENT_IN_SUCCESSOR_SOURCE"}
% HMT_RESULT_END:GWV-001:residence
% HMT_RESULT_BEGIN:DRK-001:residence
% HMT_ATOMIC_RESULT {"material_state":"INTEGRADO_O_REMITIDO_SIN_PERDIDA","proof_strength":"PROGRAMA_ABIERTO","provenance":"ARQUITECTURA_AUTORAL_PREEXISTENTE","reinforcement_state":"DEPLOYED_WITHOUT_NEW_RESULT_ID","residence":"perspectivas cosmológicas","result_id":"DRK-001","statement":"Una alternativa estructural debe reproducir curvas de rotación, lentes, crecimiento, CMB y expansión con los mismos parámetros; Landauer nulo reversible no genera por sí solo un ciclo, que requiere dinámica de frontera acción--información.","status":"PRESENT_IN_SUCCESSOR_SOURCE"}
% HMT_RESULT_END:DRK-001:residence
% HMT_RESULT_BEGIN:CMK-001:residence
% HMT_ATOMIC_RESULT {"material_state":"INTEGRADO_O_REMITIDO_SIN_PERDIDA","proof_strength":"PROGRAMA_ABIERTO","provenance":"ARQUITECTURA_AUTORAL_PREEXISTENTE","reinforcement_state":"DEPLOYED_WITHOUT_NEW_RESULT_ID","residence":"cierre Euler/TPK","result_id":"CMK-001","statement":"El programa debe construir un diagrama de módulos y transformaciones naturales que tipen pares adjuntos o involutivos; una enumeración de operadores no constituye la coestructura.","status":"PRESENT_IN_SUCCESSOR_SOURCE"}
% HMT_RESULT_END:CMK-001:residence
% HMT_RESULT_BEGIN:MRC-001:residence
% HMT_ATOMIC_RESULT {"material_state":"INTEGRADO_O_REMITIDO_SIN_PERDIDA","proof_strength":"EXACTO_INTERNO","provenance":"CERTIFICADO_NUEVO","reinforcement_state":"DEPLOYED_WITHOUT_NEW_RESULT_ID","residence":"apéndice digital y gate de la monografía de realización","result_id":"MRC-001","statement":"Conserva diez familias: cilindros 0--1--infinito; espectros APP/memoria; identidad Perron--Fibonacci; inversión antiunitaria; producto semidirecto de Poincaré; constitutivas; norma partida y clausura material; Planck/autodualidad; CGLMP; e invariantes de filtración nonádica y masas.","status":"PRESENT_IN_SUCCESSOR_SOURCE"}
% HMT_RESULT_END:MRC-001:residence
% HMT_RESULT_BEGIN:PRG68-001:residence
% HMT_ATOMIC_RESULT {"material_state":"INTEGRADO_O_REMITIDO_SIN_PERDIDA","proof_strength":"PROGRAMA_ABIERTO","provenance":"ARQUITECTURA_AUTORAL_PREEXISTENTE","reinforcement_state":"DEPLOYED_WITHOUT_NEW_RESULT_ID","residence":"cierre del capítulo de realización física","result_id":"PRG68-001","statement":"Los seis cierres distintos son: límite suave Poincaré--Cartan independiente de subdivisión; un único operador C global que separe C, P y T; electrón ligado con carga, espín, estabilidad, factor de forma y momento magnético; mapa desde Delta_4 y la clausura electrónica hasta S_matter y la corriente de espín; estado e instrumentos HMT que deriven Bell/CGLMP sin usar el blanco; y selector general theta_G del transductor gravitatorio sin introducir la coordenada externa.","status":"PRESENT_IN_SUCCESSOR_SOURCE"}
% HMT_RESULT_END:PRG68-001:residence
% HMT_RESULT_BEGIN:GTR-001:residence
% HMT_ATOMIC_RESULT {"material_state":"INTEGRADO_O_REMITIDO_SIN_PERDIDA","proof_strength":"EXACTO_INTERNO","provenance":"RESULTADO_RECUPERADO","reinforcement_state":"DEPLOYED_WITHOUT_NEW_RESULT_ID","residence":"`md/11c`, `11n` y síntesis","result_id":"GTR-001","statement":"El mapa G_int(theta_G)=theta_G c_int^5 t_0^2/hbar_int tiene tipo dimensional exacto. En la clausura circular CT108 se especializa con theta_*=(54/pi)^2; una carta declarada evalúa 6.67430×10^-11. El selector general theta_G sigue siendo una obligación distinta y no se obtiene introduciendo esa coordenada externa.","status":"PRESENT_IN_SUCCESSOR_SOURCE"}
% HMT_RESULT_END:GTR-001:residence
% HMT_RESULT_BEGIN:AMB-001:residence
% HMT_ATOMIC_RESULT {"material_state":"CONSERVADO_Y_ATOMIZADO_EN_DELTA_SUCESOR","proof_strength":"DEMOSTRADO_CON_ESTRUCTURA_DE_PARTIDA_EXPLICITA","provenance":"RESULTADO_RECUPERADO","reinforcement_state":"DEPLOYED_WITHOUT_NEW_RESULT_ID","residence":"manuscrito/sections/hmt/05b_extensiones_dinamica_cociente.tex","result_id":"AMB-001","statement":"El cociente de modos del TPK induce la incidencia primitiva F_av; la medida de Parry del envolvente de memoria uno satisface mu_ar/mu_vol=phi_HMT^-2. El lector coinductivo marca una sola vez el retorno de frontera y produce pi_HMT/phi_HMT^2. Para R(x,y)=x/y^2 y J(x,y)=(y,x), las orientaciones conjugadas son pi_HMT/phi_HMT^2 y phi_HMT/pi_HMT^2; son funcionales coordinados y no intercambiables del mismo estado holonómico integral con memoria, sin factorización horizontal.","status":"PRESENT_IN_SUCCESSOR_SOURCE"}
% HMT_RESULT_END:AMB-001:residence
% HMT_RESULT_BEGIN:CHI-II-HDG-01:residence
% HMT_ATOMIC_RESULT {"material_state":"AUDITADO_PENDIENTE_INTEGRACION","proof_strength":"DEMOSTRADO_CON_ESTRUCTURA_DE_PARTIDA_EXPLICITA","provenance":"FORMALIZACION_NUEVA","reinforcement_state":"DEPLOYED_WITHOUT_NEW_RESULT_ID","residence":"manuscrito/sections/sucesora/02_estructura_discreta_continuo_20260817.tex:Hodge_macroobjetos","result_id":"CHI-II-HDG-01","statement":"La realización tipada R_Hdg:Xχ→Perf conserva la historia enriquecida del carácter χ y publica cada estado de Xχ como complejo perfecto, sin reconstruir el antecedente desde su sombra cohomológica.","status":"PRESENT_IN_SUCCESSOR_SOURCE"}
% HMT_RESULT_END:CHI-II-HDG-01:residence
% HMT_RESULT_BEGIN:CHI-II-HDG-02:residence
% HMT_ATOMIC_RESULT {"material_state":"AUDITADO_PENDIENTE_INTEGRACION","proof_strength":"DEMOSTRADO_CON_ESTRUCTURA_DE_PARTIDA_EXPLICITA","provenance":"FORMALIZACION_NUEVA","reinforcement_state":"DEPLOYED_WITHOUT_NEW_RESULT_ID","residence":"manuscrito/sections/sucesora/02_estructura_discreta_continuo_20260817.tex:Hodge_identidad_realizacion","result_id":"CHI-II-HDG-02","statement":"La identidad cl∘chloc∘R_Hdg=ev∞ conmuta exactamente en Xχ: la realización perfecta, su carácter de Chern localizado y la clase de ciclo publican la misma evaluación terminal.","status":"PRESENT_IN_SUCCESSOR_SOURCE"}
% HMT_RESULT_END:CHI-II-HDG-02:residence
% HMT_RESULT_BEGIN:CHI-II-HDG-03:residence
% HMT_ATOMIC_RESULT {"material_state":"INTEGRADO_EN_SUCESORA_PENDIENTE_DE_SELLO","proof_strength":"DEMOSTRADO_CON_ESTRUCTURA_DE_PARTIDA_EXPLICITA","provenance":"FORMALIZACION_NUEVA","reinforcement_state":"DEPLOYED_WITHOUT_NEW_RESULT_ID","residence":"manuscrito/sections/sucesora/02_estructura_discreta_continuo_20260817.tex:eq:continuo-hodge-completitud; manuscrito/sections/sucesora/04_hodge_bsd_fragment_20260817.tex:thm:hodge-suc-generacion","result_id":"CHI-II-HDG-03","statement":"Xχ cubre el dominio Hdg y, para toda clase α de ese dominio, existe ξ_α∈Xχ con β_α=chloc(R_Hdg(ξ_α)); por la identidad de realización, cl(β_α)=α.","status":"PRESENT_IN_SUCCESSOR_SOURCE"}
% HMT_RESULT_END:CHI-II-HDG-03:residence
% HMT_RESULT_BEGIN:CHI-IV-HDG-01:residence
% HMT_ATOMIC_RESULT {"material_state":"AUDITADO_PENDIENTE_INTEGRACION","proof_strength":"DEMOSTRADO_CON_ESTRUCTURA_DE_PARTIDA_EXPLICITA","provenance":"RESULTADO_NUEVO","reinforcement_state":"DEPLOYED_WITHOUT_NEW_RESULT_ID","residence":"manuscrito/sections/sucesora/04_cinco_realizaciones_fundamentales_20260817.tex:Hodge_APP_MV_coend","result_id":"CHI-IV-HDG-01","statement":"La celda APP–Mayer–Vietoris κ materializa suma, producto, diferencia y carry en Perf; el coend sobre historias compatibles totaliza las nueve fases coordinadas y conserva como terminal Cone(I−U_Γ9), con memoria visible.","status":"PRESENT_IN_SUCCESSOR_SOURCE"}
% HMT_RESULT_END:CHI-IV-HDG-01:residence
% HMT_RESULT_BEGIN:CHI-IV-HDG-C01:residence
% HMT_ATOMIC_RESULT {"material_state":"AUDITADO_PENDIENTE_INTEGRACION","proof_strength":"EXACTO_INTERNO","provenance":"RESULTADO_NUEVO","reinforcement_state":"DEPLOYED_WITHOUT_NEW_RESULT_ID","residence":"manuscrito/sections/sucesora/04_cinco_realizaciones_fundamentales_20260817.tex:Hodge_control_no_go_eliptico","result_id":"CHI-IV-HDG-C01","statement":"En una curva elíptica compleja E no existe un frame racional finito que cubra H²(E,Q(1)) y sea invariante bajo todas las traslaciones, ni un selector natural no basado de la clase de un punto; la obstrucción obliga a conservar la rigidificación y su descenso.","status":"PRESENT_IN_SUCCESSOR_SOURCE"}
% HMT_RESULT_END:CHI-IV-HDG-C01:residence
% HMT_RESULT_BEGIN:CHI-II-BSD-01:residence
% HMT_ATOMIC_RESULT {"material_state":"AUDITADO_PENDIENTE_INTEGRACION","proof_strength":"DEMOSTRADO_CON_ESTRUCTURA_DE_PARTIDA_EXPLICITA","provenance":"FORMALIZACION_NUEVA","reinforcement_state":"DEPLOYED_WITHOUT_NEW_RESULT_ID","residence":"manuscrito/sections/sucesora/02_estructura_discreta_continuo_20260817.tex:BSD_Gmn_AW","result_id":"CHI-II-BSD-01","statement":"El complejo G_mn de cadenas normalizadas del nervio de historias TPK supervivientes porta la diagonal Alexander–Whitney; ésta conserva la unidad genealógica y hace group-like cada vértice de estado completo.","status":"PRESENT_IN_SUCCESSOR_SOURCE"}
% HMT_RESULT_END:CHI-II-BSD-01:residence
% HMT_RESULT_BEGIN:CHI-II-BSD-02:residence
% HMT_ATOMIC_RESULT {"material_state":"AUDITADO_PENDIENTE_INTEGRACION","proof_strength":"DEMOSTRADO_CON_ESTRUCTURA_DE_PARTIDA_EXPLICITA","provenance":"RESULTADO_NUEVO","reinforcement_state":"DEPLOYED_WITHOUT_NEW_RESULT_ID","residence":"manuscrito/sections/sucesora/02_estructura_discreta_continuo_20260817.tex:BSD_Manin_producto_fibrado","result_id":"CHI-II-BSD-02","statement":"El blanco conservativo BSD es Manin×_A G: el producto fibrado acopla el complejo Manin basado con G sobre el aumento A y conserva simultáneamente la proyección modular y la historia profunda.","status":"PRESENT_IN_SUCCESSOR_SOURCE"}
% HMT_RESULT_END:CHI-II-BSD-02:residence
% HMT_RESULT_BEGIN:CHI-II-BSD-03:residence
% HMT_ATOMIC_RESULT {"material_state":"AUDITADO_PENDIENTE_INTEGRACION","proof_strength":"DEMOSTRADO_CON_ESTRUCTURA_DE_PARTIDA_EXPLICITA","provenance":"RESULTADO_NUEVO","reinforcement_state":"DEPLOYED_WITHOUT_NEW_RESULT_ID","residence":"manuscrito/sections/sucesora/02_estructura_discreta_continuo_20260817.tex:BSD_publicacion_acoplada_K_PT","result_id":"CHI-II-BSD-03","statement":"Las publicaciones Kato y Poitou–Tate se construyen sobre Manin×_A G y publican la misma unidad genealógica de G; la comparación conserva raíz y complemento sin factorizar por el mero aumento.","status":"PRESENT_IN_SUCCESSOR_SOURCE"}
% HMT_RESULT_END:CHI-II-BSD-03:residence
% HMT_RESULT_BEGIN:CHI-IV-BSD-01:residence
% HMT_ATOMIC_RESULT {"material_state":"AUDITADO_PENDIENTE_INTEGRACION","proof_strength":"DEMOSTRADO_CON_ESTRUCTURA_DE_PARTIDA_EXPLICITA","provenance":"RESULTADO_NUEVO","reinforcement_state":"DEPLOYED_WITHOUT_NEW_RESULT_ID","residence":"manuscrito/sections/sucesora/04_cinco_realizaciones_fundamentales_20260817.tex:BSD_raiz_dualidad_HFS","result_id":"CHI-IV-BSD-01","statement":"La reducción ortogonal y la dualidad chain-level fijan J_red; la orientación λ elimina el twist y fuerza p_root zK=zPT, mientras H_FS construye h* para contraer el complemento finita–singular sin alterar la raíz.","status":"PRESENT_IN_SUCCESSOR_SOURCE"}
% HMT_RESULT_END:CHI-IV-BSD-01:residence
% HMT_RESULT_BEGIN:CHI-IV-BSD-02:residence
% HMT_ATOMIC_RESULT {"material_state":"INTEGRADO_EN_SUCESORA_PENDIENTE_DE_SELLO","proof_strength":"DEMOSTRADO_CON_ESTRUCTURA_DE_PARTIDA_EXPLICITA","provenance":"FORMALIZACION_NUEVA","reinforcement_state":"DEPLOYED_WITHOUT_NEW_RESULT_ID","residence":"manuscrito/sections/sucesora/04_hodge_bsd_fragment_20260817.tex:thm:bsd-suc-bockstein,thm:bsd-suc-local-global,eq:bsd-suc-formula-final","result_id":"CHI-IV-BSD-02","statement":"El regulador derivado R_Boc^der coincide como elemento basado con el término inicial LT_T(S); unido a la igualdad raíz Kato/PT, la altura, Sha, Tamagawa, periodo y torsión ya tipados, publica la fórmula BSD final sin definir un lado por el otro.","status":"PRESENT_IN_SUCCESSOR_SOURCE"}
% HMT_RESULT_END:CHI-IV-BSD-02:residence
% HMT_RESULT_BEGIN:CHI-IV-BSD-C01:residence
% HMT_ATOMIC_RESULT {"material_state":"AUDITADO_PENDIENTE_INTEGRACION","proof_strength":"EXACTO_INTERNO","provenance":"RESULTADO_NUEVO","reinforcement_state":"DEPLOYED_WITHOUT_NEW_RESULT_ID","residence":"manuscrito/sections/sucesora/04_cinco_realizaciones_fundamentales_20260817.tex:BSD_control_twist4","result_id":"CHI-IV-BSD-C01","statement":"La ablación del pullback genealógico o la inserción del twist raíz 4 conserva sombras modulares parciales pero destruye la unidad basada: zK^(4)=4r+3e+b pasa módulo 3 y falla desde módulo 9, por lo que no puede satisfacer p_root zK=zPT.","status":"PRESENT_IN_SUCCESSOR_SOURCE"}
% HMT_RESULT_END:CHI-IV-BSD-C01:residence
% HMT_RESULT_BEGIN:CHI-II-RH-01:residence
% HMT_ATOMIC_RESULT {"material_state":"AUDITADO_PENDIENTE_INTEGRACION","proof_strength":"EXACTO_INTERNO","provenance":"FORMALIZACION_NUEVA","reinforcement_state":"DEPLOYED_WITHOUT_NEW_RESULT_ID","residence":"manuscrito/sections/sucesora/02_estructura_discreta_continuo_20260817.tex:RH_macroobjeto","result_id":"CHI-II-RH-01","statement":"La Parte II construye M_RH=(Xhat_infinito,Gamma9,Rhat=(R,K),Carr_R,S0,M,O9,Xi,P,D_R_cof): Gamma9 conserva la fase y eleva la profundidad; Carr_R compone como cociclo; S0=sum_{q=0}^{N-1}M^{*q}O9^*O9M^q+M^{*N}S0M^N conserva el terminal; y B_W=Xi^*(I-P)Xi publica el complemento de nivel/fibra sin borrar memoria.","status":"PRESENT_IN_SUCCESSOR_SOURCE"}
% HMT_RESULT_END:CHI-II-RH-01:residence
% HMT_RESULT_BEGIN:CHI-IV-RH-01:residence
% HMT_ATOMIC_RESULT {"material_state":"AUDITADO_PENDIENTE_INTEGRACION","proof_strength":"DEMOSTRADO_CON_ESTRUCTURA_DE_PARTIDA_EXPLICITA","provenance":"RESULTADO_RECUPERADO","reinforcement_state":"DEPLOYED_WITHOUT_NEW_RESULT_ID","residence":"manuscrito/sections/sucesora/04_cinco_realizaciones_fundamentales_20260817.tex:RH_teorema","result_id":"CHI-IV-RH-01","statement":"Con q=5/9, M9=Phi A8···A0=qV9 y la identidad finita de memoria dan B_W,R=L_R^*L_R≥0 al extinguirse su terminal; la realización común produce B_W=Xi^*(I-P)Xi≥0 y, sobre D_R_cof con los canales completos, el criterio reversible publica Re(rho)=1/2 para todo cero no trivial.","status":"PRESENT_IN_SUCCESSOR_SOURCE"}
% HMT_RESULT_END:CHI-IV-RH-01:residence
% HMT_RESULT_BEGIN:CHI-IV-RH-C01:residence
% HMT_ATOMIC_RESULT {"material_state":"AUDITADO_PENDIENTE_INTEGRACION","proof_strength":"EXACTO_INTERNO","provenance":"RESULTADO_RECUPERADO","reinforcement_state":"DEPLOYED_WITHOUT_NEW_RESULT_ID","residence":"manuscrito/sections/sucesora/04_cinco_realizaciones_fundamentales_20260817.tex:RH_control_ablacion","result_id":"CHI-IV-RH-C01","statement":"El control retira por separado ortogonalidad, antecedente común, canal primo/gamma/polar, cofinalidad, carry, memoria o terminal; cada ablación rompe una igualdad anterior al criterio y no puede presentarse como una variante del teorema.","status":"PRESENT_IN_SUCCESSOR_SOURCE"}
% HMT_RESULT_END:CHI-IV-RH-C01:residence
% HMT_RESULT_BEGIN:CHI-II-NS-01:residence
% HMT_ATOMIC_RESULT {"material_state":"AUDITADO_PENDIENTE_INTEGRACION","proof_strength":"EXACTO_INTERNO","provenance":"FORMALIZACION_NUEVA","reinforcement_state":"DEPLOYED_WITHOUT_NEW_RESULT_ID","residence":"manuscrito/sections/sucesora/02_estructura_discreta_continuo_20260817.tex:NS_macroobjeto","result_id":"CHI-II-NS-01","statement":"La Parte II construye M_NS,m=(Xi_NS,m,C_{m←n},M_m^+,M_m^-,D_m,H_m,E9,J9,P9,Q_micro,R_m,E_m,ell_m): C_{m←n}=P_mB(u_n,u_n)-B_m(P_m u_n,P_m u_n) conserva el carry; E9J9=I, P9=J9E9 y Q_micro=I-P9 tipan la fibra; E9||Xi_{m+1}||²+E9 ell_m=||Xi_m||² conserva el terminal y archiva memoria positiva.","status":"PRESENT_IN_SUCCESSOR_SOURCE"}
% HMT_RESULT_END:CHI-II-NS-01:residence
% HMT_RESULT_BEGIN:CHI-IV-NS-01:residence
% HMT_ATOMIC_RESULT {"material_state":"AUDITADO_PENDIENTE_INTEGRACION","proof_strength":"EXACTO_INTERNO","provenance":"RESULTADO_RECUPERADO","reinforcement_state":"DEPLOYED_WITHOUT_NEW_RESULT_ID","residence":"manuscrito/sections/sucesora/04_cinco_realizaciones_fundamentales_20260817.tex:NS_teorema","result_id":"CHI-IV-NS-01","statement":"El puerto del macroobjeto satisface [B_{s,m}+(1-eta)nu D_{s,m}]^-≤C_NS ell_m+b_{parcial,m}; al componerlo con la telescopía y la energía se obtiene sup_m||u_m||_{L∞H^s}²+sup_m||u_m||_{L²H^{s+1}}²≤C_T(u0,f,Xi0) para s>5/2, y se publican existencia global, unicidad y suavidad para t>0 en el dominio declarado.","status":"PRESENT_IN_SUCCESSOR_SOURCE"}
% HMT_RESULT_END:CHI-IV-NS-01:residence
% HMT_RESULT_BEGIN:CHI-IV-NS-C01:residence
% HMT_ATOMIC_RESULT {"material_state":"AUDITADO_PENDIENTE_INTEGRACION","proof_strength":"EXACTO_INTERNO","provenance":"RESULTADO_RECUPERADO","reinforcement_state":"DEPLOYED_WITHOUT_NEW_RESULT_ID","residence":"manuscrito/sections/sucesora/04_cinco_realizaciones_fundamentales_20260817.tex:NS_control_ablacion","result_id":"CHI-IV-NS-C01","statement":"El control elimina por separado retorno conjunto, carry de Galerkin, memoria positiva, Q_micro, terminal o uniformidad del puerto; una rama de norma 3^J, C_NS dependiente del corte o b_parcial,m no sumable impiden la telescopía y rechazan la publicación global.","status":"PRESENT_IN_SUCCESSOR_SOURCE"}
% HMT_RESULT_END:CHI-IV-NS-C01:residence
% HMT_RESULT_BEGIN:CHI-II-YM-01:residence
% HMT_ATOMIC_RESULT {"material_state":"AUDITADO_PENDIENTE_INTEGRACION","proof_strength":"EXACTO_INTERNO","provenance":"FORMALIZACION_NUEVA","reinforcement_state":"DEPLOYED_WITHOUT_NEW_RESULT_ID","residence":"manuscrito/sections/sucesora/02_estructura_discreta_continuo_20260817.tex:YM_macroobjeto","result_id":"CHI-II-YM-01","statement":"La Parte II construye M_YM,m=(Lambda_m,mu_m,T_m,J_m,Theta_{nu,m},K_t,D_m^YM,P_Omega_m,Q_m^phys): Lambda_m=a_m Z^4 y K_{a_{m+1}}^9=K_{a_m}; (T_{m+1})^9J_m=J_mT_m y Theta_{nu,m+1}J_m=J_mTheta_{nu,m} conservan nivel/fibra, carry de fusión y memoria; P_Omega_m es el terminal raíz.","status":"PRESENT_IN_SUCCESSOR_SOURCE"}
% HMT_RESULT_END:CHI-II-YM-01:residence
% HMT_RESULT_BEGIN:CHI-IV-YM-01:residence
% HMT_ATOMIC_RESULT {"material_state":"AUDITADO_PENDIENTE_INTEGRACION","proof_strength":"EXACTO_INTERNO","provenance":"RESULTADO_RECUPERADO","reinforcement_state":"DEPLOYED_WITHOUT_NEW_RESULT_ID","residence":"manuscrito/sections/sucesora/04_cinco_realizaciones_fundamentales_20260817.tex:YM_teorema","result_id":"CHI-IV-YM-01","statement":"Sobre el macroobjeto común, omega_m(overline{F(Theta_nu,m U)}F(U))=||B_nu,mF||²≥0 para nu=1,…,4; el generador D_m^YM satisface ker D_m^YM=C Omega_m y D_m^YM≥3kappa_av(I-P_Omega_m); q_{m+1}^9=q_m y -a_m^{-1}log q_m=3kappa_av publican una teoría local, vacío simple y brecha Delta_YM^HMT≥3kappa_av>0 en el sector colectivo gauge-invariante.","status":"PRESENT_IN_SUCCESSOR_SOURCE"}
% HMT_RESULT_END:CHI-IV-YM-01:residence
% HMT_RESULT_BEGIN:CHI-IV-YM-C01:residence
% HMT_ATOMIC_RESULT {"material_state":"AUDITADO_PENDIENTE_INTEGRACION","proof_strength":"EXACTO_INTERNO","provenance":"RESULTADO_RECUPERADO","reinforcement_state":"DEPLOYED_WITHOUT_NEW_RESULT_ID","residence":"manuscrito/sections/sucesora/04_cinco_realizaciones_fundamentales_20260817.tex:YM_control_ablacion","result_id":"CHI-IV-YM-C01","statement":"El control elimina por separado semigrupo común, medida compartida, localidad uniforme, entrelazamiento de refinamiento, carry, memoria, vacío terminal o escala física; cada ablación rompe una identidad previa a la brecha y no puede repararse usando Wilson como selector retrospectivo.","status":"PRESENT_IN_SUCCESSOR_SOURCE"}
% HMT_RESULT_END:CHI-IV-YM-C01:residence
% HMT_RESULT_BEGIN:CHI-I-TPK-01:residence
% HMT_ATOMIC_RESULT {"material_state":"INTEGRADO_EN_LIENZO_SUCESOR_NO_ACTIVADO","proof_strength":"EXACTO_INTERNO","provenance":"FORMALIZACION_NUEVA","reinforcement_state":"DEPLOYED_WITHOUT_NEW_RESULT_ID","residence":"manuscrito/sections/sucesora/01_nucleo_formal_hmt_20260817.tex:conexion_nonadica_conjunta","result_id":"CHI-I-TPK-01","statement":"La holonomía única Γ9=Hol_C9(∇_TPK) conserva el retorno de fase con evolución del estado; en una realización isométrica satisface U_Γ9 J=JT+η, J*η=0 e I−T*T=η*η, y la iteración anti-especular conserva la suma de memoria junto con el terminal finito positivo.","status":"PRESENT_IN_SUCCESSOR_SOURCE"}
% HMT_RESULT_END:CHI-I-TPK-01:residence
% HMT_RESULT_BEGIN:CHI-I-EMRH-01:residence
% HMT_ATOMIC_RESULT {"material_state":"INTEGRADO_EN_LIENZO_SUCESOR_NO_ACTIVADO","proof_strength":"EXACTO_INTERNO","provenance":"RESULTADO_RECUPERADO","reinforcement_state":"DEPLOYED_WITHOUT_NEW_RESULT_ID","residence":"manuscrito/sections/sucesora/01_nucleo_formal_hmt_20260817.tex:completacion_cubica_conservacion_proyectiva","result_id":"CHI-I-EMRH-01","statement":"Para C̄9^(d)=(E9⊔{⋆})^d, el estrato con r coordenadas de frontera tiene cardinal B_r(d)=binom(d,r)9^(d−r), la proyección que olvida la última coordenada satisface (π_d)_*μ_d=μ_{d−1} y la masa total permanece uno; en dimensión tres, 1000=729+243+27+1=729+271.","status":"PRESENT_IN_SUCCESSOR_SOURCE"}
% HMT_RESULT_END:CHI-I-EMRH-01:residence
% HMT_RESULT_BEGIN:CHI-I-FAV-01:residence
% HMT_ATOMIC_RESULT {"material_state":"INTEGRADO_EN_LIENZO_SUCESOR_NO_ACTIVADO","proof_strength":"DEMOSTRADO_CON_ESTRUCTURA_DE_PARTIDA_EXPLICITA","provenance":"FORMALIZACION_NUEVA","reinforcement_state":"DEPLOYED_WITHOUT_NEW_RESULT_ID","residence":"manuscrito/sections/sucesora/01_nucleo_formal_hmt_20260817.tex:cociente_areal_volumetrico","result_id":"CHI-I-FAV-01","statement":"El bloque primitivo (+1,−1,0) induce la órbita de modos (Σ,Π,Π); el lector tipado de hojas transporta su incidencia a F_av=[[0,1],[1,1]], y la repetición del calendario produce N9=3F_av, N27=9F_av y N108=36F_av.","status":"PRESENT_IN_SUCCESSOR_SOURCE"}
% HMT_RESULT_END:CHI-I-FAV-01:residence
% HMT_RESULT_BEGIN:CHI-III-VAC-01:residence
% HMT_ATOMIC_RESULT {"material_state":"AUDITADO_PENDIENTE_INTEGRACION","proof_strength":"EXACTO_INTERNO","provenance":"FORMALIZACION_NUEVA","reinforcement_state":"DEPLOYED_WITHOUT_NEW_RESULT_ID","residence":"manuscrito/sections/sucesora/03_genealogia_constantes_20260817.tex:descomposicion_Fitting_vacancias","result_id":"CHI-III-VAC-01","statement":"En T9=(Z/9Z)^2, M=[[3,4],[6,7]] satisface M²−M=3I; E=M³ es idempotente y N=M−E cumple N²=0 y EN=NE, de modo que T9=im(E)⊕ker(E) y la filtración de imágenes tiene cardinales 81→27→9→9.","status":"PRESENT_IN_SUCCESSOR_SOURCE"}
% HMT_RESULT_END:CHI-III-VAC-01:residence
% HMT_RESULT_BEGIN:CHI-III-MASS-01:residence
% HMT_ATOMIC_RESULT {"material_state":"AUDITADO_PENDIENTE_INTEGRACION","proof_strength":"DEMOSTRADO_CON_ESTRUCTURA_DE_PARTIDA_EXPLICITA","provenance":"FORMALIZACION_NUEVA","reinforcement_state":"DEPLOYED_WITHOUT_NEW_RESULT_ID","residence":"manuscrito/sections/sucesora/03_genealogia_constantes_20260817.tex:dominio_completo_ley_general_masas","result_id":"CHI-III-MASS-01","statement":"El dominio interno vinculante de la Ley General de Masas comprende 81 celdas, 56 familias de ruta, 13 multisecciones, 324 rutas orientadas y un sector nulo separado; cada extensión superviviente publica ruta, registro, firma, carácter y operador de fibra antes de toda carta metrológica. Las antiguas ventanas de diez o doce masas son sólo testigos históricos.","status":"PRESENT_IN_SUCCESSOR_SOURCE"}
% HMT_RESULT_END:CHI-III-MASS-01:residence
% HMT_RESULT_BEGIN:CHI-III-ELECTRON-01:residence
% HMT_ATOMIC_RESULT {"material_state":"AUDITADO_PENDIENTE_INTEGRACION","proof_strength":"DEMOSTRADO_CON_ESTRUCTURA_DE_PARTIDA_EXPLICITA","provenance":"FORMALIZACION_NUEVA","reinforcement_state":"DEPLOYED_WITHOUT_NEW_RESULT_ID","residence":"manuscrito/sections/sucesora/03_genealogia_constantes_20260817.tex:electron_central_two_way","result_id":"CHI-III-ELECTRON-01","statement":"La fibra electrónica ocupa el único centro (5,5) con ruta Gamma_e=(5,5,++); sus lectores independientes coinciden en K_D(Gamma_e)=K_Omega(Gamma_e)=(80,54,6), fijan kappa_e=80−54alpha_HMT+6Delta_4 y, después del basal y de la recta de acción, publican m_e^beta=0.5109989506900008086082571648… MeV. El basal permanece como etapa, no como salida rectora.","status":"PRESENT_IN_SUCCESSOR_SOURCE"}
% HMT_RESULT_END:CHI-III-ELECTRON-01:residence
% HMT_RESULT_BEGIN:CHI-III-BOLTZ-01:residence
% HMT_ATOMIC_RESULT {"material_state":"AUDITADO_PENDIENTE_INTEGRACION","proof_strength":"EXACTO_INTERNO","provenance":"FORMALIZACION_NUEVA","reinforcement_state":"DEPLOYED_WITHOUT_NEW_RESULT_ID","residence":"manuscrito/sections/sucesora/03_genealogia_constantes_20260817.tex:boltzmann_accion_informacion","result_id":"CHI-III-BOLTZ-01","statement":"La constante de Boltzmann se tipa como el isomorfismo positivo k_B:L_Theta→L_E; una decisión TRIT uniforme aporta I_TRIT=log 3 y la periodicidad de orden 108 fija k_B Theta_clk log 3=hbar_int 2pi/(108 t0). La construcción determina internamente la energía por unidad de información; la carta kelvin–joule publica después k_B=1.380649×10^−23 J K^−1.","status":"PRESENT_IN_SUCCESSOR_SOURCE"}
% HMT_RESULT_END:CHI-III-BOLTZ-01:residence
% HMT_RESULT_BEGIN:CHI-III-PLANCK108-01:residence
% HMT_ATOMIC_RESULT {"material_state":"AUDITADO_PENDIENTE_INTEGRACION","proof_strength":"DEMOSTRADO_CON_ESTRUCTURA_DE_PARTIDA_EXPLICITA","provenance":"FORMALIZACION_NUEVA","reinforcement_state":"DEPLOYED_WITHOUT_NEW_RESULT_ID","residence":"manuscrito/sections/sucesora/03_genealogia_constantes_20260817.tex:planck_periodicidad_108","result_id":"CHI-III-PLANCK108-01","statement":"Las cartas de acción hbar_pre y hbar_ret admiten una realización circular de período T_108=108t0, con omega_108=2pi/T_108, R_108=c_int/omega_108 y C_108=108 ell0; para m_108,a=hbar_a omega_108/c_int² se cumplen barlambda_C=R_108 y lambda_C=C_108. La condición de radio gravitatorio reducido selecciona theta_108=(54/pi)² y coordina las unidades internas de Planck en ambas cartas.","status":"PRESENT_IN_SUCCESSOR_SOURCE"}
% HMT_RESULT_END:CHI-III-PLANCK108-01:residence
% HMT_RESULT_BEGIN:CHI-VI-GRAV-ALG-01:residence
% HMT_ATOMIC_RESULT {"material_state":"INTEGRADO_EN_LIENZO_SUCESOR_NO_ACTIVADO","proof_strength":"EXACTO_INTERNO","provenance":"FORMALIZACION_NUEVA","reinforcement_state":"DEPLOYED_WITHOUT_NEW_RESULT_ID","residence":"manuscrito/sections/sucesora/06_mecanica_dimensional_20260817.tex:sector_tensorial_icosaedrico_y_proyeccion_TT","result_id":"CHI-VI-GRAV-ALG-01","statement":"La representación icosaédrica de A5 sobre V12 posee el proyector ortogonal P5 de rango cinco; la igualdad de caracteres identifica V5=im(P5) con Sym²_0(R³) mediante una isometría A5-equivariante única hasta signo, y el proyector transversal sin traza Lambda(n) tiene rango dos. La cadena exacta es 12→5→5→2 y no se confunde con una realización gravitatoria dinámica.","status":"PRESENT_IN_SUCCESSOR_SOURCE"}
% HMT_RESULT_END:CHI-VI-GRAV-ALG-01:residence
% HMT_RESULT_BEGIN:AMP-I-RECTA-01:residence
% HMT_ATOMIC_RESULT {"material_state":"AMPLIFICACION_AUTORIZADA_2026_08_17","proof_strength":"EXACTO_INTERNO","provenance":"FORMALIZACION_NUEVA","reinforcement_state":"DEPLOYED_WITHOUT_NEW_RESULT_ID","residence":"manuscrito/sections/ampliacion_20260817/00_recta_horizonte_completacion.tex","result_id":"AMP-I-RECTA-01","statement":"La sucesión de cartas [0,10]→[0,1000]→[0,10^m] construye un sistema compatible de marcas, intervalos y fronteras cuya evaluación final vive en R; el quinto postulado euclídeo se trata como antecedente histórico de cierre geométrico y la completación se formula mediante mapas y límites, no como metáfora.","status":"PRESENT_IN_AMPLIFIED_SOURCE"}
% HMT_RESULT_END:AMP-I-RECTA-01:residence
% HMT_RESULT_BEGIN:AMP-I-UNIDAD-01:residence
% HMT_ATOMIC_RESULT {"material_state":"AMPLIFICACION_AUTORIZADA_2026_08_17","proof_strength":"EXACTO_INTERNO","provenance":"FORMALIZACION_NUEVA","reinforcement_state":"DEPLOYED_WITHOUT_NEW_RESULT_ID","residence":"manuscrito/sections/ampliacion_20260817/00_recta_horizonte_completacion.tex","result_id":"AMP-I-UNIDAD-01","statement":"La unidad normalizada se conserva bajo refinamiento, elevación dimensional y doble publicación; en dimensión d la estratificación binomial de (E9⊔{★})^d suma 10^d y los push-forward proyectivos conservan masa uno.","status":"PRESENT_IN_AMPLIFIED_SOURCE"}
% HMT_RESULT_END:AMP-I-UNIDAD-01:residence
% HMT_RESULT_BEGIN:AMP-II-EXC-01:residence
% HMT_ATOMIC_RESULT {"material_state":"AMPLIFICACION_AUTORIZADA_2026_08_17","proof_strength":"DEMOSTRADO_CON_ESTRUCTURA_DE_PARTIDA_EXPLICITA","provenance":"RESULTADO_RECUPERADO","reinforcement_state":"DEPLOYED_WITHOUT_NEW_RESULT_ID","residence":"manuscrito/sections/sucesora/02_estructura_discreta_continuo_20260817.tex; manuscrito/sections/sucesora/06_mecanica_dimensional_20260817.tex","result_id":"AMP-II-EXC-01","statement":"La rama excepcional se bifurca tipadamente: C_W→W12→W24 por diseños, y C_W→N(A2^12)→Λ24→V_Λ→V^natural→Monster→Moonshine por retículo y orbifold; no existe una serialización Witt→Golay→A2^12.","status":"PRESENT_IN_AMPLIFIED_SOURCE"}
% HMT_RESULT_END:AMP-II-EXC-01:residence
% HMT_RESULT_BEGIN:AMP-III-TIPOS-01:residence
% HMT_ATOMIC_RESULT {"material_state":"AMPLIFICACION_AUTORIZADA_2026_08_17","proof_strength":"DEMOSTRADO_CON_ESTRUCTURA_DE_PARTIDA_EXPLICITA","provenance":"RESULTADO_RECUPERADO","reinforcement_state":"DEPLOYED_WITHOUT_NEW_RESULT_ID","residence":"manuscrito/sections/sucesora/03_genealogia_constantes_20260817.tex; manuscrito/sections/sucesora/06_mecanica_dimensional_20260817.tex","result_id":"AMP-III-TIPOS-01","statement":"G se expone en tres capas —transductor exacto, selector circular de orden 108 y criterio holonómico general—; Boltzmann deriva internamente k_B Θ_clk log3 y sólo después abre la carta SI; raíz de -1 permanece en Parte II y -1/12 en Parte III.","status":"PRESENT_IN_AMPLIFIED_SOURCE"}
% HMT_RESULT_END:AMP-III-TIPOS-01:residence
% HMT_RESULT_BEGIN:AMP-IV-REAL-01:residence
% HMT_ATOMIC_RESULT {"material_state":"AMPLIFICACION_AUTORIZADA_2026_08_17","proof_strength":"DEMOSTRADO_CON_ESTRUCTURA_DE_PARTIDA_EXPLICITA","provenance":"FORMALIZACION_NUEVA","reinforcement_state":"DEPLOYED_WITHOUT_NEW_RESULT_ID","residence":"manuscrito/sections/ampliacion_20260817/04_principio_realizacion_y_supersesion.tex","result_id":"AMP-IV-REAL-01","statement":"Una familia F_N:S_N^surv→C_N compatible con truncación, orientación, transporte, acarreo, unidad y terminal induce F_∞ por la propiedad universal del límite inverso y transporta conjuntamente observable y memoria.","status":"PRESENT_IN_AMPLIFIED_SOURCE"}
% HMT_RESULT_END:AMP-IV-REAL-01:residence
% HMT_RESULT_BEGIN:AMP-IV-RHNS-01:residence
% HMT_ATOMIC_RESULT {"material_state":"AMPLIFICACION_AUTORIZADA_2026_08_17","proof_strength":"DEMOSTRADO_CON_ESTRUCTURA_DE_PARTIDA_EXPLICITA","provenance":"RESULTADO_RECUPERADO","reinforcement_state":"DEPLOYED_WITHOUT_NEW_RESULT_ID","residence":"manuscrito/sections/ampliacion_20260817/04a_refuerzo_riemann_weil.tex; manuscrito/sections/ampliacion_20260817/04b_refuerzo_navier_stokes.tex","result_id":"AMP-IV-RHNS-01","statement":"RH incorpora el residual finito 729=81+648 y R_m,R=0 antes del criterio; NS incorpora cota conjunta, ∂_t u en L2_loc H^{s-1}, pegado por unicidad y presión normalizada, sin sustituir sus demostraciones vigentes.","status":"PRESENT_IN_AMPLIFIED_SOURCE"}
% HMT_RESULT_END:AMP-IV-RHNS-01:residence
% HMT_RESULT_BEGIN:AMP-IV-YM-01:residence
% HMT_ATOMIC_RESULT {"material_state":"AMPLIFICACION_AUTORIZADA_2026_08_17","proof_strength":"DEMOSTRADO_CON_ESTRUCTURA_DE_PARTIDA_EXPLICITA","provenance":"RESULTADO_RECUPERADO","reinforcement_state":"DEPLOYED_WITHOUT_NEW_RESULT_ID","residence":"manuscrito/sections/ampliacion_20260817/04c_refuerzo_yang_mills.tex","result_id":"AMP-IV-YM-01","statement":"La familia proyectiva común produce cuatro formas OS, un Hamiltoniano común, acción fuerte de E(4), funcionales de Schwinger, publicación en H^{-s}_loc(R4), s>2, lectores clover y F², vacío simple y brecha exacta 3κ_av.","status":"PRESENT_IN_AMPLIFIED_SOURCE"}
% HMT_RESULT_END:AMP-IV-YM-01:residence
% HMT_RESULT_BEGIN:AMP-IV-HDG-01:residence
% HMT_ATOMIC_RESULT {"material_state":"AMPLIFICACION_AUTORIZADA_2026_08_17","proof_strength":"DEMOSTRADO_CON_ESTRUCTURA_DE_PARTIDA_EXPLICITA","provenance":"FORMALIZACION_NUEVA","reinforcement_state":"DEPLOYED_WITHOUT_NEW_RESULT_ID","residence":"manuscrito/sections/ampliacion_20260817/04d_hodge_completitud_desplegada.tex","result_id":"AMP-IV-HDG-01","statement":"La completitud de Xχ(X,p) se despliega como teorema coinductivo de realizaciones finitas compatibles; produce primero ξ_α y sólo después β_α=ch_p^loc(R_Hdg(ξ_α)), preservando el cuadrado de publicación.","status":"PRESENT_IN_AMPLIFIED_SOURCE"}
% HMT_RESULT_END:AMP-IV-HDG-01:residence
% HMT_RESULT_BEGIN:AMP-IV-BSD-01:residence
% HMT_ATOMIC_RESULT {"material_state":"AMPLIFICACION_AUTORIZADA_2026_08_17","proof_strength":"DEMOSTRADO_CON_ESTRUCTURA_DE_PARTIDA_EXPLICITA","provenance":"FORMALIZACION_NUEVA","reinforcement_state":"DEPLOYED_WITHOUT_NEW_RESULT_ID","residence":"manuscrito/sections/ampliacion_20260817/04e_bsd_comparadores_desplegados.tex","result_id":"AMP-IV-BSD-01","statement":"Los comparadores Kato–FERS/Iwasawa, Poitou–Tate/Poincaré–Néron–Tate, Smith–Tamagawa–Shafarevich, torsión y período se construyen y componen sobre una única línea orientada; la igualdad raíz elimina la unidad relativa y publica rango, finitud de Sha y fórmula final.","status":"PRESENT_IN_AMPLIFIED_SOURCE"}
% HMT_RESULT_END:AMP-IV-BSD-01:residence
% HMT_RESULT_BEGIN:AMP-IV-SUP-01:residence
% HMT_ATOMIC_RESULT {"material_state":"AMPLIFICACION_AUTORIZADA_2026_08_17","proof_strength":"EXACTO_INTERNO","provenance":"RESULTADO_RECUPERADO","reinforcement_state":"DEPLOYED_WITHOUT_NEW_RESULT_ID","residence":"manuscrito/sections/ampliacion_20260817/04_principio_realizacion_y_supersesion.tex; manuscrito/appendices_acumulativa_revfinal.tex","result_id":"AMP-IV-SUP-01","statement":"Los rótulos históricos que declaraban abiertos RH, NS, YM, Hodge y BSD quedan supersedidos en esta edición por los cinco teoremas de Parte IV; se conservan únicamente como trazabilidad de formulaciones anteriores y controles tipados, sin autoridad sobre el corte vigente.","status":"PRESENT_IN_AMPLIFIED_SOURCE"}
% HMT_RESULT_END:AMP-IV-SUP-01:residence
% HMT_RESULT_BEGIN:AMP-V-GRAV5-01:residence
% HMT_ATOMIC_RESULT {"material_state":"AMPLIFICACION_AUTORIZADA_2026_08_17","proof_strength":"EXACTO_INTERNO","provenance":"FORMALIZACION_NUEVA","reinforcement_state":"DEPLOYED_WITHOUT_NEW_RESULT_ID","residence":"manuscrito/sections/ampliacion_20260817/05_sector_tensorial_gravitatorio_cinco_componentes.tex","result_id":"AMP-V-GRAV5-01","statement":"La acción icosaédrica dodecafásica induce V12≅V1⊕V3⊕V3'⊕V5; el proyector central P5, el entrelazador explícito Γ5 hacia Sym0²(R3), la base STF y la reducción TT construyen la cadena 12→5→5→2 y separan portador quíntuple, espín dos masivo, GR y clases E(2).","status":"PRESENT_IN_AMPLIFIED_SOURCE"}
% HMT_RESULT_END:AMP-V-GRAV5-01:residence
% HMT_RESULT_BEGIN:AMP-V-GRAVPHY-01:residence
% HMT_ATOMIC_RESULT {"material_state":"AMPLIFICACION_AUTORIZADA_2026_08_17","proof_strength":"INTERFAZ_FISICA","provenance":"FORMALIZACION_NUEVA","reinforcement_state":"DEPLOYED_WITHOUT_NEW_RESULT_ID","residence":"manuscrito/sections/ampliacion_20260817/05_sector_tensorial_gravitatorio_cinco_componentes.tex","result_id":"AMP-V-GRAVPHY-01","statement":"El transductor G, la realización circular de orden 108 y el sector tensorial algebraico conservan sus conclusiones exactas; la realización PCH se expone como cadena tipada que determina el espacio físico mediante ecuaciones, restricciones y calibre, sin atribuir cinco polarizaciones a Einstein, Feynman ni GR.","status":"PRESENT_IN_AMPLIFIED_SOURCE"}
% HMT_RESULT_END:AMP-V-GRAVPHY-01:residence
% HMT_RESULT_BEGIN:AMP-EPI-R-01:residence
% HMT_ATOMIC_RESULT {"material_state":"AMPLIFICACION_AUTORIZADA_2026_08_17","proof_strength":"DEMOSTRADO_CON_ESTRUCTURA_DE_PARTIDA_EXPLICITA","provenance":"FORMALIZACION_NUEVA","reinforcement_state":"DEPLOYED_WITHOUT_NEW_RESULT_ID","residence":"manuscrito/sections/ampliacion_20260817/99_epilogo_retorno_recta_real.tex","result_id":"AMP-EPI-R-01","statement":"El epílogo retorna a R y demuestra que la recta real es la publicación arquimediana de una sección compatible, mientras la publicación excepcional conserva incidencia; ambas proceden del mismo estado y ninguna reconstruye por sí sola su genealogía.","status":"PRESENT_IN_AMPLIFIED_SOURCE"}
% HMT_RESULT_END:AMP-EPI-R-01:residence
 % Sidecar sucesor de la reparación semántica aprobada: 32 residencias
% rectoras, comentario-only y enlazadas al nuevo recibo de precompilación.
% Sidecar probatorio de la reparación semántica 2026-08-19.
% No produce salida tipográfica; cada intervalo liga una residencia rectora.
% HMT_REPAIR_RESULT_BEGIN::HMT-CORE-APP-TRIT-TPK
% RESULT_ID: HMT-CORE-APP-TRIT-TPK
% STATUS: DEMONSTRATED
% RESIDENCE: manuscrito/sections/sucesora/01_nucleo_formal_hmt_20260817.tex#eq:nucleo-app-dos-hojas
% STATEMENT_SHA256: 6657284d9bd2918464c9d964efdaae3608a94c55cd238f3de467fe85c00ad0ac
% STATEMENT: APP construye X_9; TRIT tipa sus dos hojas; después se forman los dominios de semilla, las emisiones y las palabras visibles, y sólo entonces se produce el estado holonómico integral que conserva hoja, orientación, residuo, cociente, acarreo, firma, cilindro, frontera, ruta, memoria, multisección y supervivencia.
% HMT_REPAIR_RESULT_END::HMT-CORE-APP-TRIT-TPK
% HMT_REPAIR_RESULT_BEGIN::HMT-CORE-NONADIC-JOINT-CONTINUITY
% RESULT_ID: HMT-CORE-NONADIC-JOINT-CONTINUITY
% STATUS: DEMONSTRATED
% RESIDENCE: manuscrito/sections/sucesora/01_nucleo_formal_hmt_20260817.tex#thm:nucleo-continuidad-conexion-nonadica-conjunta
% STATEMENT_SHA256: bea2a7fd24c115638fea7eb3682c82f9b684160dc0a4e855abb04024c2e484c4
% STATEMENT: Γ_9 devuelve la fase tras nueve pasos sin reinicializar el estado; conserva memoria, acarreo y terminal, y sus potencias son composiciones relacionales compatibles con truncamientos.
% HMT_REPAIR_RESULT_END::HMT-CORE-NONADIC-JOINT-CONTINUITY
% HMT_REPAIR_RESULT_BEGIN::HMT-RECTOR-CONTINUUM-CLOSURE-CONSTANTS
% RESULT_ID: HMT-RECTOR-CONTINUUM-CLOSURE-CONSTANTS
% STATUS: DEMONSTRATED
% RESIDENCE: manuscrito/sections/sucesora/02_estructura_discreta_continuo_20260817.tex#eq:continuo-teorema-fundamental-discreto
% STATEMENT_SHA256: 34c506d5e3311116e4a5f5d4baf1febf44fb1fcbd22e17a6540197f17479fbd7
% STATEMENT: HMT cierra el problema del continuo en el dominio genealógico: las publicaciones continuas son lectores compatibles de secciones discretas completas; las constantes emergen del mismo antecedente, y los cinco problemas del Milenio son especializaciones concretas y subordinadas del cierre rector.
% HMT_REPAIR_RESULT_END::HMT-RECTOR-CONTINUUM-CLOSURE-CONSTANTS
% HMT_REPAIR_RESULT_BEGIN::HMT-CORE-NO-ABLATION
% RESULT_ID: HMT-CORE-NO-ABLATION
% STATUS: DEMONSTRATED
% RESIDENCE: gestion/reparacion_semantica_20260818/MAPA_RECONSTRUCCION_CAUSAL_FUENTE_20260818.md §7.1
% STATEMENT_SHA256: 6a8421d9f55a806faa959c19578d710ae27e286555775c9a1f34551268c7e592
% STATEMENT: El mapa de olvido U retiene una sombra visible, conjuntista, ordinal, escalar o convencional y elimina genealogía, memoria, orientación, cociente/residuo, acarreo, frontera/ruta, supervivencia o multisección; ninguna objeción a esa sombra se transfiere al objeto completo sin probar que U conserva y refleja la propiedad examinada.
% HMT_REPAIR_RESULT_END::HMT-CORE-NO-ABLATION
% HMT_REPAIR_RESULT_BEGIN::HMT-ZFC-VON-NEUMANN-NONFAITHFUL-CH
% RESULT_ID: HMT-ZFC-VON-NEUMANN-NONFAITHFUL-CH
% STATUS: DEMONSTRATED
% RESIDENCE: manuscrito/sections/integracion_20260812/von_neumann_64/04_zfc_fundacion_ordinales.tex#thm:zfc-perdida-genealogia-olvido
% STATEMENT_SHA256: 9fcf0f78fc08d6066f7bf7f9f8aeb86d0a5e87348f2efb1affef488e061e0b9b
% STATEMENT: Tres resultados distintos quedan tipados: U_60 colapsa dos objetos no isomorfos con la misma Sig_60 y fibras Ext_81 de cardinales 0 y 12, por lo que no refleja isomorfías ni clasifica; V no es fiel porque id_ρ y Γ_9 tienen la misma sombra unipuntual pero grados de acarreo/memoria (0,0) y (1,1); los ordinales finitos con memorias k y k+1 proyectan al mismo ordinal de von Neumann. CH es un blanco distinto.
% HMT_REPAIR_RESULT_END::HMT-ZFC-VON-NEUMANN-NONFAITHFUL-CH
% HMT_REPAIR_RESULT_BEGIN::HMT-CORE-DOUBLE-PUBLICATION
% RESULT_ID: HMT-CORE-DOUBLE-PUBLICATION
% STATUS: DEMONSTRATED
% RESIDENCE: manuscrito/sections/sucesora/02_estructura_discreta_continuo_20260817.tex#eq:continuo-doble-publicacion
% STATEMENT_SHA256: 3051702377a6942a887f59b9323a35081f9d1ba17e9b3e8a72326b7039ab0e43
% STATEMENT: El lector arquimediano contrae cilindros compatibles y el lector incidencial conserva la genealogía de frontera; sus codominios son distintos. La rama excepcional es bifurcada: Hadamard organiza U_12^sgn en paralelo y Paley produce C_W, desde donde parten las ramas Witt–Golay–Mathieu y Niemeier–Leech–V^natural–Monstruo–Moonshine.
% HMT_REPAIR_RESULT_END::HMT-CORE-DOUBLE-PUBLICATION
% HMT_REPAIR_RESULT_BEGIN::HMT-MILLENNIUM-RH-WEIL-CLOSED
% RESULT_ID: HMT-MILLENNIUM-RH-WEIL-CLOSED
% STATUS: DEMONSTRATED
% RESIDENCE: manuscrito/sections/sucesora/04_cinco_realizaciones_fundamentales_20260817.tex#thm:realizaciones-rh
% STATEMENT_SHA256: a4f088786afdbee666309a8fc5034b3722fade65c9f6a259d05da2b5732626b5
% STATEMENT: En cada región R se fija q=5/9 y se construye E_R=j_-Q_RXi_R con AE_R=qE_R y E_R*E_R=B_{W,R}; la telescopía finita conserva el terminal q^(2N)B_{W,R}, cuya extinción sólo se toma al pasar N→∞, y L_R satisface L_R*L_R=B_{W,R}⪰0. El descenso cofinal publica la forma de Weil global como cuadrado positivo y el criterio de Weil reconoce la hipótesis de Riemann.
% HMT_REPAIR_RESULT_END::HMT-MILLENNIUM-RH-WEIL-CLOSED
% HMT_REPAIR_RESULT_BEGIN::HMT-MILLENNIUM-NS-CLOSED
% RESULT_ID: HMT-MILLENNIUM-NS-CLOSED
% STATUS: DEMONSTRATED
% RESIDENCE: manuscrito/sections/sucesora/04_cinco_realizaciones_fundamentales_20260817.tex#thm:realizaciones-ns
% STATEMENT_SHA256: f8e920750843d201de3d006e4d249eb58a18849ab02f1db10372401a049235db
% STATEMENT: La rama propietaria PC–Fock–KYP construye prospectivamente el Gram del terminal y las pérdidas futuras, deriva la identidad Stein y la telescopía con terminal, factoriza directamente el bloque KYP, establece balance, observabilidad y registro de torre, y produce el presupuesto raíz uniforme y la cota H^s global. El comparador G_{M,j}=P_aux*U^term_{M,j}P_aux es sólo un control no gobernante.
% HMT_REPAIR_RESULT_END::HMT-MILLENNIUM-NS-CLOSED
% HMT_REPAIR_RESULT_BEGIN::HMT-MILLENNIUM-YM-CLOSED
% RESULT_ID: HMT-MILLENNIUM-YM-CLOSED
% STATUS: DEMONSTRATED
% RESIDENCE: manuscrito/sections/sucesora/04_cinco_realizaciones_fundamentales_20260817.tex#thm:realizaciones-ym
% STATEMENT_SHA256: 7bae1f46e683874449e24e2e5e76323c00b6953291777330a09a1c2324fc5167
% STATEMENT: La misma familia proyectiva construye medida, red local, generador D_m^YM, semigrupo K_{m,t}=e^(−tD_m^YM), vacío, cuatro formas OS, acción fuerte E(4) y brecha exacta. Se cumplen K_{m,t}P_{Ω_m}=P_{Ω_m} y ||K_{m,t}(I−P_{Ω_m})||≤e^(−3κ_av t); el testigo W_c alcanza la cota. El lector clover/F² reconoce después la curvatura del mismo proceso.
% HMT_REPAIR_RESULT_END::HMT-MILLENNIUM-YM-CLOSED
% HMT_REPAIR_RESULT_BEGIN::HMT-MILLENNIUM-HODGE-CLOSED
% RESULT_ID: HMT-MILLENNIUM-HODGE-CLOSED
% STATUS: DEMONSTRATED
% RESIDENCE: manuscrito/sections/sucesora/04_hodge_bsd_fragment_20260817.tex#thm:hodge-suc-generacion
% STATEMENT_SHA256: 871872c44580d0e280556a6d45a7a549870a5b324223da5310a7a61a0ee0f72b
% STATEMENT: La realización gobernante parte de X_χ=lim_K Surv_K(X,p), conmuta con el carácter local y la publicación de clases y, para cada α, construye ξ_α y β_α con cl(β_α)=α. Los núcleos relativos y la presentación coend son controles de autonomía expositiva, no la cadena gobernante.
% HMT_REPAIR_RESULT_END::HMT-MILLENNIUM-HODGE-CLOSED
% HMT_REPAIR_RESULT_BEGIN::HMT-MILLENNIUM-BSD-CLOSED
% RESULT_ID: HMT-MILLENNIUM-BSD-CLOSED
% STATUS: DEMONSTRATED
% RESIDENCE: manuscrito/sections/sucesora/04_hodge_bsd_fragment_20260817.tex#eq:bsd-suc-formula-final
% STATEMENT_SHA256: 58f17d0b3189c35e80a258c376c281f796aee2740fea34308df90316cbe709c7
% STATEMENT: La coálgebra Alexander–Whitney y su counidad construyen G_{m,n} y el producto fibrado P_E^gen; una misma unidad genera las publicaciones Kato/PT, la igualdad de raíz, el relleno h_*, la bandera de Bockstein, los comparadores y la igualdad basada que cierra BSD. La unipotencia generador a generador es sólo control de despliegue.
% HMT_REPAIR_RESULT_END::HMT-MILLENNIUM-BSD-CLOSED
% HMT_REPAIR_RESULT_BEGIN::HMT-EXCEPTIONAL-CORRIDOR
% RESULT_ID: HMT-EXCEPTIONAL-CORRIDOR
% STATUS: DEMONSTRATED
% RESIDENCE: manuscrito/sections/hmt/09c_genealogia_excepcional_completa_autonoma.tex#sec:genealogia-excepcional-completa-autonoma
% STATEMENT_SHA256: 9e1f0885079d7e94e3bbf4142ac1fef08314e4d1d537ffc3403e23bd7d3bb08b
% STATEMENT: Hadamard organiza U_12^sgn sin convertirse en antecedente serial de Paley; la cadena C_P→A_W→C_W bifurca desde C_W hacia Witt–Golay–Mathieu y hacia Niemeier–Leech–V_Λ–V^natural; después Aut(V^natural)=Monstruo y J=T_{1A} precede a la familia {T_g} de Moonshine.
% HMT_REPAIR_RESULT_END::HMT-EXCEPTIONAL-CORRIDOR
% HMT_REPAIR_RESULT_BEGIN::HMT-INTERNAL-I-QUARTER-TURN
% RESULT_ID: HMT-INTERNAL-I-QUARTER-TURN
% STATUS: DEMONSTRATED
% RESIDENCE: manuscrito/sections/hmt/09a_paley_codigo_witt.tex#eq:AW-raiz-menos-uno
% STATEMENT_SHA256: 89af0b10c8bc5da185fb92a05117661adee2496d45d963a0922f2c33d9cf247e
% STATEMENT: La matriz incidencial A_W satisface A_W²=−I en el sector Paley–Witt; ésta es la raíz finita interna. La estructura J_E²=−I y su lector exponencial publican después la realización compleja arquimediana, sin identificar ambos operadores por nombre.
% HMT_REPAIR_RESULT_END::HMT-INTERNAL-I-QUARTER-TURN
% HMT_REPAIR_RESULT_BEGIN::HMT-INTERNAL-MINUS-ONE-TWELFTH
% RESULT_ID: HMT-INTERNAL-MINUS-ONE-TWELFTH
% STATUS: DEMONSTRATED
% RESIDENCE: manuscrito/sections/sucesora/03_genealogia_constantes_20260817.tex#eq:constantes-menos-un-doce-proyectores
% STATEMENT_SHA256: d04beccf9787ff880f26f73a4ff543adb5aab378144657d59b5d156e5fd65284
% STATEMENT: El valor −1/12 emerge en seis dominios tipados: incidencia 90/120, diferencia de periodos, Gram positivo con pseudoinversa, extractor de escala, exponente eta y traza de proyectores del corrimiento cíclico S_12:Q^12→Q^12, con S_12^12=I, donde D_r=I−S_12^r y los proyectores sobre ker D_3 y ker D_4 tienen rangos 3 y 4. La igualdad ζ(−1)=−1/12 es un reconocimiento analítico separado, no el origen de esas realizaciones.
% HMT_REPAIR_RESULT_END::HMT-INTERNAL-MINUS-ONE-TWELFTH
% HMT_REPAIR_RESULT_BEGIN::HMT-LORENTZ-MINKOWSKI
% RESULT_ID: HMT-LORENTZ-MINKOWSKI
% STATUS: DEMONSTRATED
% RESIDENCE: manuscrito/sections/hmt/03_app_ortograma.tex#prop:app-trit-espectral-markov-lorentz-rev8
% STATEMENT_SHA256: 7da5585d1b7752e0cd76fc580631d618b0aedef2c4bc51e326b4cc8bae1649e5
% STATEMENT: El bloque central B_c de APP produce M_c=B_c/√45∈SO^+(1,1) y satisface M_c^Tη_{1,1}M_c=η_{1,1}; ésta es una realización exacta de Lorentz y del plano de Minkowski 1+1. Toda elevación 3+1 requiere un mapa dimensional tipado y no demueve el resultado 1+1.
% HMT_REPAIR_RESULT_END::HMT-LORENTZ-MINKOWSKI
% HMT_REPAIR_RESULT_BEGIN::HMT-HILBERT-REALIZATION
% RESULT_ID: HMT-HILBERT-REALIZATION
% STATUS: DEMONSTRATED
% RESIDENCE: manuscrito/sections/integracion_20260812/von_neumann_64/06_hilbert_heisenberg.tex
% STATEMENT_SHA256: 5323eb3e082071f29482590ea4ede54b91c92314e2f7d6f865ddf5cde5fea80d
% STATEMENT: El espacio de Hilbert matricial, el límite UHF y su cierre GNS, la dinámica compleja y Born se construyen por ramas tipadas que no se implican automáticamente: C_9^(d) produce H_d,A_d,UHF(9^∞), su traza τ produce la representación GNS (π_τ,H_τ,Ω_τ) y π_τ(A_(9^∞))ʺ=R_hyp; J_E,H∈B(H_τ), con J_E²=−I y H=H*, producen U(t); el lector Born GNS usa P=P*=P²∈A_UHF actuando por π_τ(P) y p_τ(P)=<Ω_τ,π_τ(P)Ω_τ>; una segunda realización finita conserva Pr_ρ(a)=Tr(ρE_a) desde una PVM y un estado normalizado.
% HMT_REPAIR_RESULT_END::HMT-HILBERT-REALIZATION
% HMT_REPAIR_RESULT_BEGIN::HMT-T-DUALITY-M-SCREENS
% RESULT_ID: HMT-T-DUALITY-M-SCREENS
% STATUS: DEMONSTRATED
% RESIDENCE: manuscrito/sections/hmt/16b_pantallas_dualidad_teoria_m.tex#thm:dualidad-visible-memoria
% STATEMENT_SHA256: a6c3111b2029fd7c5e93e357ea6b65a0cd583f22589de276bcbf2a3215fc5d31
% STATEMENT: En D_dual=Z²×R_(>0), T(m,w,R_0)=(w,m,R_0^(−1)) satisface T²=id y preserva E_(R_0)(m,w)=m²/R_0²+w²R_0² mediante E_(R_0)(m,w)=E_(R_0^(−1))(w,m); su realización unitaria convive con las reducciones 11→10 y 26→24. La interfaz física separada R_M:X_HMT^enr→X_11^phys debe producir y compatibilizar g_(MN)∈Γ(Sym² T*X_11), C_(MNP)∈Ω³(X_11), ψ_M∈Γ(S_11⊗T*X_11), su acción y transformaciones, sin contaminar el teorema interno.
% HMT_REPAIR_RESULT_END::HMT-T-DUALITY-M-SCREENS
% HMT_REPAIR_RESULT_BEGIN::HMT-CONSTANT-PI
% RESULT_ID: HMT-CONSTANT-PI
% STATUS: DEMONSTRATED
% RESIDENCE: manuscrito/sections/ampliacion_20260818/02_generacion_infinita_pi_e_phi.tex#eq:generacion-infinita-gaussiana-pi-20260818
% STATEMENT_SHA256: 26b11a14165eacab3aa7d06fcb90f31ff0f2eb0dd328b3f7c11cf183fee1293f
% STATEMENT: La órbita y el carácter 5/4/239 se seleccionan antes de abrir intervalos reales; la identidad de Machin reconoce posteriormente el lector ya construido.
% HMT_REPAIR_RESULT_END::HMT-CONSTANT-PI
% HMT_REPAIR_RESULT_BEGIN::HMT-CONSTANT-E
% RESULT_ID: HMT-CONSTANT-E
% STATUS: DEMONSTRATED
% RESIDENCE: manuscrito/sections/ampliacion_20260818/02_generacion_infinita_pi_e_phi.tex#eq:generacion-infinita-recurrencia-e-20260818
% STATEMENT_SHA256: 2c21447e726bae153a60dc20b39c536b01f8b4dba8836c1176eeb252af969ec0
% STATEMENT: La ley multiplicativa del canal fuerza la recurrencia a_n=1/n! antes de reconocer la serie exponencial y su valor e.
% HMT_REPAIR_RESULT_END::HMT-CONSTANT-E
% HMT_REPAIR_RESULT_BEGIN::HMT-CONSTANT-PHI
% RESULT_ID: HMT-CONSTANT-PHI
% STATUS: DEMONSTRATED
% RESIDENCE: capítulo autónomo de la constante de autoescala, ecuación de prolongación 20260818
% STATEMENT_SHA256: 8151c9f0c4164daffee20672e7f6751ed9452b96c233d2ecd5e81aedca948338
% STATEMENT: La incidencia areal–volumétrica induce F_av y su rayo de Perron fuerza x²−x−1=0; la raíz positiva se reconoce después como φ.
% HMT_REPAIR_RESULT_END::HMT-CONSTANT-PHI
% HMT_REPAIR_RESULT_BEGIN::HMT-CONSTANT-ALPHA
% RESULT_ID: HMT-CONSTANT-ALPHA
% STATUS: DEMONSTRATED
% RESIDENCE: manuscrito/sections/sucesora/03_genealogia_constantes_20260817.tex#thm:constantes-alpha-independencia-principal
% STATEMENT_SHA256: 8eb80629b88f75eaf8f624a0e6992e43812a576d8f34b530137934ff6bf0d7a6
% STATEMENT: En el perfil HMT completo, B_term y U_12^sgn son proyecciones paralelas del mismo estado terminal; R_sgn extrae U, H_12 invertible da K, y el acarreo y la recurrencia producen α sin α, CODATA ni masa como entrada.
% HMT_REPAIR_RESULT_END::HMT-CONSTANT-ALPHA
% HMT_REPAIR_RESULT_BEGIN::HMT-EULER-HOLOGRAPHIC
% RESULT_ID: HMT-EULER-HOLOGRAPHIC
% STATUS: DEMONSTRATED
% RESIDENCE: manuscrito/sections/ampliacion_20260818/03_extension_holografica_euler.tex#thm:euler-holografica-tres-cierres-20260818
% STATEMENT_SHA256: f54720a289461a9177197a7184922777c9bb186d4617980d744cae299400c025
% STATEMENT: A_W²=−I proporciona la raíz finita incidencial; en una rama distinta, J_E²=−I y el lector exponencial producen la identidad arquimediana de Euler, extendida por J_τ a tres ramas. El contraángulo no nace de Euler: sólo después de α_HMT, la microfibra 169, D_A, H_5 y la recurrencia regional generan C*, cuya identidad inversa se reconoce al final.
% HMT_REPAIR_RESULT_END::HMT-EULER-HOLOGRAPHIC
% HMT_REPAIR_RESULT_BEGIN::HMT-EMRH-CUBIC-UNIT
% RESULT_ID: HMT-EMRH-CUBIC-UNIT
% STATUS: DEMONSTRATED
% RESIDENCE: manuscrito/sections/ampliacion_20260818/01_emrh_desarrollo_integral.tex#eq:emrh-cierre-apice-20260818
% STATEMENT_SHA256: 6f630c70e8d9bf2ba536e9a19b4f59ce7f4f90709e8c1f3410e2cf49cfab5724
% STATEMENT: La completación correcta es 10^3=729+271, con 271 como corona E y 729 como bloque alfa; además (1/271+1/729)^−1=271·729/1000 expresa exactamente la razón inversa asociada a la completación volumétrica.
% HMT_REPAIR_RESULT_END::HMT-EMRH-CUBIC-UNIT
% HMT_REPAIR_RESULT_BEGIN::HMT-CONSTANTS-FOUR-STRATA
% RESULT_ID: HMT-CONSTANTS-FOUR-STRATA
% STATUS: DEMONSTRATED
% RESIDENCE: manuscrito/sections/ampliacion_20260818/05_constantes_fisicas_genealogia.tex#thm:amp18-constantes-cuatro-estratos
% STATEMENT_SHA256: f67dfab70770a2b7aed78a5925e7f56252e00166de0b0e1ce805a0be3520e083
% STATEMENT: Toda constante se presenta en el orden constructor discreto/genealógico, caracterización analítica, transducción dimensional equivariante y comparación metrológica; la carta SI no construye el invariante.
% HMT_REPAIR_RESULT_END::HMT-CONSTANTS-FOUR-STRATA
% HMT_REPAIR_RESULT_BEGIN::HMT-CONSTANT-PLANCK
% RESULT_ID: HMT-CONSTANT-PLANCK
% STATUS: DEMONSTRATED
% RESIDENCE: manuscrito/sections/md/03b_escala_accion_determinantal.tex#thm:orden-accion
% STATEMENT_SHA256: 0bc899b4d06add106a366bf912d47203626d8fde1c6911f1073da88cc17eeef0
% STATEMENT: El lector determinantal Σ_H=φ_HMT α_HMT^16 y su cokernel local de orden 16 fijan la recta de acción y la década −34; la mantisa y la carta J·s se distinguen, y cualquier residual permanece visible.
% HMT_REPAIR_RESULT_END::HMT-CONSTANT-PLANCK
% HMT_REPAIR_RESULT_BEGIN::HMT-CONSTANT-BOLTZMANN
% RESULT_ID: HMT-CONSTANT-BOLTZMANN
% STATUS: DEMONSTRATED
% RESIDENCE: manuscrito/sections/ampliacion_20260818/05_constantes_fisicas_genealogia.tex#prop:amp18-constantes-covariancia-termica
% STATEMENT_SHA256: c61e51204132bcfb86f480ea2592cef625c3aa9bdbbce4632ba8c5995cc47828
% STATEMENT: El invariante interno fija k_B Θ_clk log 3=ħ_int 2π/(108t_0); la covariancia Θ→λΘ, k_B→k_B/λ demuestra que el kelvin fija la carta, no el objeto.
% HMT_REPAIR_RESULT_END::HMT-CONSTANT-BOLTZMANN
% HMT_REPAIR_RESULT_BEGIN::HMT-CONSTITUTIVES-MU0-EPSILON0
% RESULT_ID: HMT-CONSTITUTIVES-MU0-EPSILON0
% STATUS: DEMONSTRATED
% RESIDENCE: manuscrito/sections/ampliacion_20260818/05_constantes_fisicas_genealogia.tex#eq:amp18-constantes-identidades-vacio
% STATEMENT_SHA256: 5a67e3f2c22019847d22853c8038fa4258f414e3c63a0d1568a2e08908c456fd
% STATEMENT: A partir de α_HMT y una única carta global de e,h,c se construye la impedancia Z_0 y, mediante Z_0=μ_0c y Z_0c=1/ε_0, se publican μ_0 y ε_0; sus valores SI son comparación posterior.
% HMT_REPAIR_RESULT_END::HMT-CONSTITUTIVES-MU0-EPSILON0
% HMT_REPAIR_RESULT_BEGIN::HMT-CONSTANT-G
% RESULT_ID: HMT-CONSTANT-G
% STATUS: DEMONSTRATED
% RESIDENCE: manuscrito/sections/ampliacion_20260818/05_constantes_fisicas_genealogia.tex#eq:amp18-constantes-selector-108
% STATEMENT_SHA256: 4e8377ac898d655608bc820b018d8760487159e3f6349a41d5aef00743e4fae5
% STATEMENT: El selector circular fija θ_108=(54/π_HMT)^2 y el transductor exacto define G_int(θ_108). La vía holonómica independiente da la plantilla homogénea G_hol=(c_int^3/ħ_int)(δ_θΛ_27)^2, cuya cifra exige construir y normalizar Λ_27; la carta decimal SI se mantiene separada y no construye ninguna de las dos vías.
% HMT_REPAIR_RESULT_END::HMT-CONSTANT-G
% HMT_REPAIR_RESULT_BEGIN::HMT-MASS-LAW-ATLAS
% RESULT_ID: HMT-MASS-LAW-ATLAS
% STATUS: DEMONSTRATED
% RESIDENCE: manuscrito/sections/ampliacion_20260818/04_electron_y_ley_general_masas.tex#eq:amp18-masas-ley-operatorial
% STATEMENT_SHA256: 4666f47d67c5e3e1b827ca0cb3802e083fbee5770e215ceea03a3ab43eefdc84
% STATEMENT: La partícula es sección persistente y la masa es carácter positivo de esa persistencia. El electrón central único (5,5) construye primero el basal adimensional \widetilde{\mathfrak e}_0=(√3/4)[exp(φ_HMT/π_HMT²)−22α_HMT³](1+15Δ_4), después \widetilde{\mathfrak e}_e^β=\widetilde{\mathfrak e}_0R_act^(80−54α_HMT+6Δ_4), y sólo entonces la carta común T_E:=\varsigma_{u_E}:R_{>0}^{adim}→R_{>0}u_E, T_E(x)=xu_E, u_E=1 MeV publica \mathfrak e_e^β=T_E(\widetilde{\mathfrak e}_e^β)=m_e^βc²=0.5109989506900008086082571648… MeV y m_e^β=\mathfrak e_e^β/c²=0.5109989506900008086082571648… MeV/c². La misma carta común gobierna el atlas 81/56/13/324; los 23 sectores son sólo un inventario humano separado.
% HMT_REPAIR_RESULT_END::HMT-MASS-LAW-ATLAS
% HMT_REPAIR_RESULT_BEGIN::HMT-BARBERO-IMMIRZI
% RESULT_ID: HMT-BARBERO-IMMIRZI
% STATUS: DEMONSTRATED
% RESIDENCE: manuscrito/sections/md/07_hexada_octada_barbero.tex#chap:barbero
% STATEMENT_SHA256: 360316071f48bb431d8e0b9a737194c74511243d354eafeaf5e398f932b048e9
% STATEMENT: Las incidencias 90/120, las series de los dos canales, el residuo 1/24 y la minimalidad diofántica fijan los pesos (12,1) y el funcional Γ_6→8; la identificación γ=Γ ocurre después y no define retroactivamente A o C*.
% HMT_REPAIR_RESULT_END::HMT-BARBERO-IMMIRZI
% HMT_REPAIR_RESULT_BEGIN::HMT-CKM-CP
% RESULT_ID: HMT-CKM-CP
% STATUS: DEMONSTRATED
% RESIDENCE: manuscrito/sections/ampliacion_20260818/05_constantes_fisicas_genealogia.tex#eq:amp18-constantes-ckm-cuatro-ecuaciones
% STATEMENT_SHA256: a245142e9d5486d59a0cd73afac1263a4916dd3b49d3437964457a148a0d3878
% STATEMENT: El registro de rutas y el único par angular construyen cuatro ángulos, la matriz unitaria CKM y un invariante de Jarlskog no nulo; la comparación PDG es posterior y CP no se confunde con CPT.
% HMT_REPAIR_RESULT_END::HMT-CKM-CP
% HMT_REPAIR_RESULT_BEGIN::HMT-GOVERNANCE-NO-REGRESSION
% RESULT_ID: HMT-GOVERNANCE-NO-REGRESSION
% STATUS: DEMONSTRATED
% RESIDENCE: gestion/reparacion_semantica_20260818/MAPA_RECONSTRUCCION_CAUSAL_FUENTE_20260818.md §7
% STATEMENT_SHA256: c59a160a4fe22784718235a339c7a3a3b5df6305829c5ef7e443d93c2ca38b7a
% STATEMENT: Todo artefacto obsoleto que invierta APP→TRIT→estado del TPK, haga externas las semillas, convierta los cinco cierres en tareas pendientes o subordine el continuo al marco clásico se retira de la gobernanza o queda como testigo inactivo; no se arrastra a la edición sucesora.
% HMT_REPAIR_RESULT_END::HMT-GOVERNANCE-NO-REGRESSION
 % Entorno editorial para incorporar cuerpos científicos completos sin
% crear capitulos nuevos. El grupo conserva intactos los archivos de origen:
% solo traduce su jerarquia visible un nivel hacia abajo durante el input.
% La procedencia material de cada uso se registra en comentarios del destino
% y en gestion/RESTAURACION_16_CLAVES_SIN_DESTINO_20260824.tsv.
\newenvironment{HMTScientificOwnerSubordinated}{%
  \begingroup
  \let\HMTScientificOwnerSection\section
  \let\HMTScientificOwnerSubsection\subsection
  \let\HMTScientificOwnerSubsubsection\subsubsection
  \let\HMTScientificOwnerParagraph\paragraph
  \let\HMTScientificOwnerSubparagraph\subparagraph
  \let\chapter\HMTScientificOwnerSection
  \let\section\HMTScientificOwnerSubsection
  \let\subsection\HMTScientificOwnerSubsubsection
  \let\subsubsection\HMTScientificOwnerParagraph
  \let\paragraph\HMTScientificOwnerSubparagraph
}{%
  \par
  \endgroup
}
 \colorlet{hmtlight}{hmtpaper}
\theoremstyle{definition}
\newtheorem{principle}[theorem]{Principle}

% Jerarquía tipográfica de monografía científica. Los títulos de capítulo
% comparten línea con su numeración, se componen con cuerpo moderado y dejan
% espacio suficiente para el primer párrafo o la primera ecuación.
\titleformat{\part}[display]
  {\centering\normalfont\bfseries\color{hmtblue}}
  {\Large\scshape\partname\ \thepart}
  {.45ex}
  {\Huge}
\titlespacing*{\part}{0pt}{1.25cm}{1.0cm}
\titleformat{\chapter}[hang]
  {\normalfont\bfseries\color{hmtblue}\LARGE}
  {\normalsize\scshape\chaptertitlename\ \thechapter}
  {1.0em}
  {\raggedright}
\titlespacing*{\chapter}{0pt}{-0.35cm}{0.75cm}
\titleformat{name=\chapter,numberless}[block]
  {\normalfont\bfseries\color{hmtblue}\LARGE}
  {}{0pt}{\raggedright}
\titlespacing*{name=\chapter,numberless}{0pt}{-0.35cm}{.75cm}
\titleformat{\section}
  {\normalfont\bfseries\Large\color{hmtblue}}
  {\thesection}{.65em}{}
\titlespacing*{\section}{0pt}{2.0ex plus .45ex minus .2ex}{.65ex}
\titleformat{\subsection}
  {\normalfont\bfseries\large}
  {\thesubsection}{.65em}{}
\titlespacing*{\subsection}{0pt}{1.65ex plus .4ex minus .2ex}{.55ex}
\titleformat{\subsubsection}
  {\normalfont\itshape\normalsize}
  {\thesubsubsection}{.65em}{}
\titlespacing*{\subsubsection}{0pt}{1.35ex plus .3ex minus .2ex}{.4ex}

% Jerarquía editorial de la edición revisada. Los archivos heredados
% conservan sus rótulos y etiquetas, pero sus capítulos se presentan como
% secciones dentro de capítulos macroscópicos.
\let\HMTBookChapter\chapter
\let\HMTBookSection\section
\let\HMTBookSubsection\subsection
\let\HMTBookSubsubsection\subsubsection
\let\HMTBookParagraph\paragraph
\let\HMTBookSubparagraph\subparagraph
\newcommand{\HMTDemoteHeadings}{%
  \let\chapter\HMTBookSection
  \let\section\HMTBookSubsection
  \let\subsection\HMTBookSubsubsection
  \let\subsubsection\HMTBookParagraph
  \let\paragraph\HMTBookSubparagraph}
\newcommand{\HMTRestoreHeadings}{%
  \let\chapter\HMTBookChapter
  \let\section\HMTBookSection
  \let\subsection\HMTBookSubsection
  \let\subsubsection\HMTBookSubsubsection
  \let\paragraph\HMTBookParagraph
  \let\subparagraph\HMTBookSubparagraph}

% Notación de la frontera con Mecánica Dimensional.
\providecommand{\APP}{\ensuremath{\mathrm{APP}}}
\providecommand{\TPK}{\ensuremath{\mathrm{TPK}}}
\providecommand{\R}{\mathbb R}
\providecommand{\Z}{\mathbb Z}
\providecommand{\N}{\mathbb N}
\providecommand{\F}{\mathbb F}
\providecommand{\M}{\mathbb M}
\providecommand{\one}{\mathbf 1}
\providecommand{\Halg}{\mathfrak H}
\providecommand{\cA}{\mathcal A}
\providecommand{\cC}{\mathcal C}
\providecommand{\cR}{\mathcal R}
\providecommand{\cX}{\mathcal X}
\providecommand{\dd}{\mathrm d}
\providecommand{\e}{\mathrm e}
\providecommand{\ii}{\mathrm i}
\providecommand{\Znine}{\mathbb Z/9\mathbb Z}
\providecommand{\Cstar}{C^{*}}
\providecommand{\Cstarrad}{C^{*}_{\rm rad}}
\providecommand{\hbarHMT}{\hbar_{\mathrm{HMT}}}
\providecommand{\etaRet}{\eta_{\mathrm{ret}}^{(\mathrm{deg})}}
\providecommand{\etaRetRad}{\eta_{\mathrm{ret}}^{(\mathrm{rad})}}
\providecommand{\alphaAngle}{\alpha_{\angle,\mathrm{deg}}}
\providecommand{\alphaHMT}{\alpha}
\providecommand{\piHMT}{\pi}
\providecommand{\phiHMT}{\varphi}
\providecommand{\eexp}{\mathrm e}
\providecommand{\Dfour}{\Delta_{4}}
\providecommand{\zCorr}{\zeta_{270}^{\mathrm{corr}}}
\providecommand{\sSpec}{s_{270}^{\mathrm{spec}}}
\providecommand{\Kfield}{\mathbb K}
\providecommand{\rank}{\operatorname{rank}}
\providecommand{\nint}{\operatorname{nint}}
\providecommand{\MD}{\ensuremath{\mathrm{MD}}}
\providecommand{\CT}{\ensuremath{\mathrm{CT}_{108}}}
\providecommand{\RR}{\mathbb R}
\providecommand{\ZZ}{\mathbb Z}
\providecommand{\PP}{\mathbb P}
\providecommand{\hbarpre}{\hbar_{\mathrm{pre}}}
\providecommand{\hbarret}{\hbar_{\mathrm{ret}}}
\providecommand{\Ract}{R_{\mathrm{act}}}
\providecommand{\ellP}{\ell_{\mathrm P}}
\providecommand{\mP}{m_{\mathrm P}}
\providecommand{\tP}{t_{\mathrm P}}
\providecommand{\EP}{E_{\mathrm P}}
\providecommand{\rg}{r_{\mathrm g}}
\providecommand{\rs}{r_{\mathrm s}}
\providecommand{\lambdabarC}{\bar\lambda_{\mathrm C}}
\providecommand{\lambdaC}{\lambda_{\mathrm C}}
\providecommand{\Nsp}{N_{\mathrm{sp}}}
\providecommand{\thetaStar}{\theta_{\!*}}
\providecommand{\Gammae}{\Gamma_{e}}
\providecommand{\mebeta}{m_{e}^{\beta}}
% En esta monografía las letras universales no designan duplicados de las
% constantes: son las mismas coordenadas producidas por la construcción y
% reconocidas después en la realización arquimediana.
\renewcommand{\alphaHMT}{\alpha}
\renewcommand{\piHMT}{\pi}
\renewcommand{\phiHMT}{\varphi}
\providecommand{\epigraphtext}[2]{%
  \begin{flushright}%
  \begin{minipage}{0.76\textwidth}\itshape
  #1\par\smallskip
  \raggedleft\normalfont--- #2
  \end{minipage}%
  \end{flushright}%
}

\newtcolorbox{historybox}[1][]{%
  colback=hmtlightblue,colframe=hmtblue,
  title={Documented historical fact},fonttitle=\bfseries,#1}
\newtcolorbox{readingbox}[1][]{%
  colback=hmtlightgold,colframe=hmtgold,
  title={Lectura HMT},fonttitle=\bfseries,#1}
\newtcolorbox{statusbox}[1][]{%
  colback=white,colframe=hmtblue!55,
  title={Hypothesis, domain, and conclusion},fonttitle=\bfseries,#1}
\newtcolorbox{warningbox}[1][]{%
  colback=white,colframe=hmtgold!75!black,
  title={Compatibility condition},fonttitle=\bfseries,#1}
\newtcolorbox{thesisbox}[1][]{%
  colback=white,colframe=hmtviolet,
  title={Tesis rectora},fonttitle=\bfseries,#1}
\newtcolorbox{proofstatusbox}[1][]{%
  colback=white,colframe=hmtgreen,
  title={Estatuto probatorio},fonttitle=\bfseries,#1}
\newtcolorbox{falsifierbox}[1][]{%
  colback=white,colframe=hmtred,
  title={Falsification criterion},fonttitle=\bfseries,#1}
\newtcolorbox{falsadorBox}[1][]{%
  colback=hmtlightred,colframe=hmtred,
  title={Control negativo},fonttitle=\bfseries,#1}
\newtcolorbox{fronteraBox}[1][]{%
  colback=hmtlightgold,colframe=hmtgold,
  title={Boundary of validity},fonttitle=\bfseries,#1}
\newtcolorbox{programbox}[1][]{%
  colback=hmtlightgold,colframe=hmtgold,
  title={Hypotheses and test program},fonttitle=\bfseries,#1}
\newtcolorbox{takeawaybox}[1][]{%
  colback=white,colframe=hmtblue,
  title={Result of the section},fonttitle=\bfseries,#1}
% Entornos académicos conservados por los fuentes teoremáticas científicos extensos.
% Se tipan aquí con la jerarquía gráfica única del tratado para evitar que cada
% incorporación restablezca una tipografía o un preámbulo autónomos.
\newtcolorbox{resultado}[1][]{%
  colback=white,colframe=hmtgreen,
  title={Result},fonttitle=\bfseries,#1}
\newtcolorbox{definicionnarrativa}[1][]{%
  colback=hmtlightblue,colframe=hmtblue,
  title={Definition},fonttitle=\bfseries,#1}
\newtcolorbox{lecturafisica}[1][]{%
  colback=hmtlightgold,colframe=hmtgold,
  title={Physical realization},fonttitle=\bfseries,#1}
\newtcolorbox{frontera}[1][]{%
  colback=hmtlightgold,colframe=hmtgold,
  title={Boundary of validity},fonttitle=\bfseries,#1}
\newtheorem{falsifier}[theorem]{Falsification criterion}
\newtheorem{test}[theorem]{Testing program}
\newenvironment{keyresult}
  {\begin{quote}\color{hmtblue}\bfseries}
  {\end{quote}}

\lstdefinestyle{hmtcode}{
  basicstyle=\ttfamily\scriptsize,
  backgroundcolor=\color{hmtpaper},
  frame=single,
  rulecolor=\color{hmtgray!45},
  breaklines=true,
  showstringspaces=false,
  columns=fullflexible,
  keepspaces=true
}

\hypersetup{
  colorlinks=true,
  bookmarksdepth=3,
  pdflang={es-ES},
  linkcolor=hmtblue,
  citecolor=hmtblue,
  urlcolor=hmtgold,
  pdftitle={Double holographic projection of an aperiodic time crystal},
  pdfsubject={Construction of the continuum, generation of constants, and exceptional stabilizers of boundary incidence},
  pdfauthor={Oumar Haidara Fall; Rubén Ramos Balsa},
  pdfkeywords={nonadic arithmetic, ternary orientation, holonomy,
    aperiodic time crystal, continuum, fundamental constants,
    exceptional symmetries, action, information, gravity}
}

\setstretch{1.035}
\setlength{\parskip}{0.16em}
\setlength{\parindent}{1.25em}
% El índice científico de la edición definitiva hace visibles capítulos y
% secciones rectoras. Los niveles inferiores permanecen en el cuerpo para
% conservar la demostración sin convertir el índice en un inventario.
\setcounter{tocdepth}{3}
\setcounter{secnumdepth}{3}
\renewcommand{\contentsname}{Index}
\renewcommand{\cftchappagefont}{\fontsize{11.5}{13.4}\selectfont\bfseries}
\renewcommand{\cftchapfont}{\fontsize{11.5}{13.4}\selectfont\bfseries}
\renewcommand{\cftsecfont}{\fontsize{10.75}{12.8}\selectfont\bfseries}
\renewcommand{\cftsecpagefont}{\fontsize{10.75}{12.8}\selectfont\bfseries}
\renewcommand{\cftsubsecfont}{\fontsize{10.25}{12.2}\selectfont}
\renewcommand{\cftsubsecpagefont}{\fontsize{10.25}{12.2}\selectfont}
\renewcommand{\cftsubsubsecfont}{\fontsize{9.5}{11.4}\selectfont\itshape}
\renewcommand{\cftsubsubsecpagefont}{\fontsize{9.5}{11.4}\selectfont\itshape}
\setlength{\cftbeforechapskip}{.38em}
\setlength{\cftbeforesecskip}{.05em}
\setlength{\cftchapnumwidth}{3.25em}
\setlength{\cftsecindent}{3.25em}
\setlength{\cftsecnumwidth}{4.15em}
\setlength{\cftsubsecnumwidth}{5.4em}
\setlist[itemize]{leftmargin=1.65em,itemsep=.25em,topsep=.35em}
\setlist[enumerate]{leftmargin=2em,itemsep=.25em,topsep=.35em}
\emergencystretch=3em
\widowpenalty=10000
\clubpenalty=10000
\displaywidowpenalty=10000
\brokenpenalty=10000
\predisplaypenalty=500
\postdisplaypenalty=500
\raggedbottom
% Cada epígrafe debe permanecer unido a un desarrollo mínimo. Se protegen
% tanto los mandos públicos como sus copias editoriales, pues los capítulos
% heredados se degradan localmente mediante \HMTDemoteHeadings.
\pretocmd{\section}{\Needspace{9\baselineskip}}{}{}
\pretocmd{\subsection}{\Needspace{7\baselineskip}}{}{}
\pretocmd{\subsubsection}{\Needspace{5\baselineskip}}{}{}
\pretocmd{\HMTBookSection}{\Needspace{9\baselineskip}}{}{}
\pretocmd{\HMTBookSubsection}{\Needspace{7\baselineskip}}{}{}
\pretocmd{\HMTBookSubsubsection}{\Needspace{5\baselineskip}}{}{}
\pretocmd{\HMTBookParagraph}{\Needspace{9\baselineskip}}{}{}

\crefname{chapter}{chapter}{chapters}
\Crefname{chapter}{Chapter}{Chapters}
\crefname{section}{section}{sections}
\Crefname{section}{Section}{Sections}
\crefname{subsection}{section}{sections}
\Crefname{subsection}{Section}{Sections}
\crefname{subsubsection}{section}{sections}
\Crefname{subsubsection}{Section}{Sections}
\crefname{figure}{Figure}{Figures}
\Crefname{figure}{Figure}{Figures}
\crefname{table}{table}{tables}
\Crefname{table}{Table}{Tables}
\crefname{equation}{equation}{equations}
\Crefname{equation}{Equation}{Equations}
\crefname{theorem}{theorem}{theorems}
\Crefname{theorem}{Theorem}{Theorems}
\crefname{lemma}{lemma}{lemmas}
\Crefname{lemma}{Lemma}{Lemmas}
\crefname{proposition}{proposition}{propositions}
\Crefname{proposition}{Proposition}{Propositions}
\crefname{corollary}{corollary}{corollaries}
\Crefname{corollary}{Corollary}{Corollaries}
\crefname{definition}{definition}{definitions}
\Crefname{definition}{Definition}{Definitions}
\crefname{construction}{construction}{constructions}
\Crefname{construction}{Construction}{Constructions}
\crefname{algorithm}{algorithm}{algorithms}
\Crefname{algorithm}{Algorithm}{Algorithms}
\crefname{ablation}{proposition of structural necessity}
  {propositions of structural necessity}
\Crefname{ablation}{Structural Necessity Proposition}
  {Structural Necessity Propositions}
\crefname{remark}{remark}{observations}
\Crefname{remark}{Remark}{Observations}
\crefname{principle}{principle}{principles}
\Crefname{principle}{Principle}{Principles}
\crefname{criterion}{test criterion}{criteria of contrast}
\Crefname{criterion}{Test criterion}{Test criteria}
\AtBeginDocument{%
  \renewcommand{\crefrangeconjunction}{ to\nobreakspace}%
  \renewcommand{\crefpairconjunction}{ and\nobreakspace}%
  \renewcommand{\crefmiddleconjunction}{, }%
  \renewcommand{\creflastconjunction}{ and\nobreakspace}%
  \renewcommand{\crefpairgroupconjunction}{ and\nobreakspace}%
  \renewcommand{\crefmiddlegroupconjunction}{, }%
  \renewcommand{\creflastgroupconjunction}{ and\nobreakspace}%
}

\pagestyle{fancy}
\fancyhf{}
\fancyhead[LE,RO]{\small\thepage}
% Las cajas de marca se reservan dentro del 84 % central del encabezado. Así
% los títulos extensos pueden ocupar dos líneas sin invadir el folio exterior.
\fancyhead[LO]{\parbox[b]{.84\headwidth}{\raggedright\footnotesize\nouppercase{\leftmark}}}
\fancyhead[RE]{\parbox[b]{.84\headwidth}{\raggedleft\footnotesize\nouppercase{\leftmark}}}
\renewcommand{\headrulewidth}{0.25pt}
\fancypagestyle{plain}{%
  \fancyhf{}\fancyfoot[C]{\thepage}\renewcommand{\headrulewidth}{0pt}}

% Jerarquía óptica moderada aprobada para el sucesor y mapa declarativo de
% rótulos públicos. El adaptador del radión aporta sólo sus macros locales;
% no importa el estilo autónomo ni modifica la tipografía general.
% Escala tipográfica aprobada por Rubén el 28 de agosto de 2026.
% Este archivo sustituye exclusivamente la escala de encabezamientos; no
% altera los cuerpos científicos ni sus relaciones de numeración.
\titleformat{\part}[display]
  {\centering\normalfont\color{hmtblue}}
  {\fontsize{12.4}{14.2}\selectfont\mdseries\scshape\partname\ \thepart}
  {2.5pt}
  {\fontsize{16.2}{18.4}\selectfont\bfseries}
\titlespacing*{\part}{0pt}{1.1cm}{.82cm}

\titleformat{\chapter}[display]
  {\normalfont\color{hmtblue}}
  {\fontsize{9.4}{11.2}\selectfont\mdseries\scshape\chaptertitlename\ \thechapter}
  {2.4pt}
  {\fontsize{13.7}{15.7}\selectfont\bfseries\raggedright}
\titlespacing*{\chapter}{0pt}{-.22cm}{.62cm}

\titleformat{name=\chapter,numberless}[block]
  {\normalfont\fontsize{13.7}{15.7}\selectfont\bfseries\color{hmtblue}}
  {}{0pt}{\raggedright}
\titlespacing*{name=\chapter,numberless}{0pt}{-.22cm}{.62cm}

\titleformat{\section}
  {\normalfont\fontsize{11.4}{13.4}\selectfont\bfseries\color{hmtblue}}
  {\thesection}{.65em}{}
\titlespacing*{\section}{0pt}{13pt plus 2pt minus 1pt}{3.5pt}

\titleformat{\subsection}
  {\normalfont\fontsize{10.4}{12.3}\selectfont\bfseries}
  {\thesubsection}{.65em}{}
\titlespacing*{\subsection}{0pt}{8pt plus 1.5pt minus 1pt}{3pt}

\titleformat{\subsubsection}
  {\normalfont\fontsize{9.7}{11.6}\selectfont\itshape}
  {\thesubsubsection}{.65em}{}
\titlespacing*{\subsubsection}{0pt}{6pt plus 1.2pt minus .8pt}{2.5pt}

% El índice reproduce la misma jerarquía descendente y evita que un nivel
% subordinado aparezca ópticamente mayor que su antecedente.
\renewcommand{\cftchappagefont}{\fontsize{10.5}{12.4}\selectfont\bfseries}
\renewcommand{\cftchapfont}{\fontsize{10.5}{12.4}\selectfont\bfseries}
\renewcommand{\cftsecfont}{\fontsize{9.9}{11.8}\selectfont\bfseries}
\renewcommand{\cftsecpagefont}{\fontsize{9.9}{11.8}\selectfont\bfseries}
\renewcommand{\cftsubsecfont}{\fontsize{9.35}{11.2}\selectfont}
\renewcommand{\cftsubsecpagefont}{\fontsize{9.35}{11.2}\selectfont}
\renewcommand{\cftsubsubsecfont}{\fontsize{8.9}{10.8}\selectfont\itshape}
\renewcommand{\cftsubsubsecpagefont}{\fontsize{8.9}{10.8}\selectfont\itshape}
% Los folios de cuatro cifras requieren una caja propia: el margen derecho
% evita que el último término del rótulo colisione con la paginación.
\cftsetpnumwidth{3em}
\cftsetrmarg{3.8em}
 % Mapa declarativo de retitulado aprobado. No altera los cuerpos científicos.
% Cada entrada procede de una coincidencia estática uno-a-uno y conserva el nivel.
\makeatletter
\DeclareUnicodeCharacter{00B2}{\ensuremath{^{2}}}
\DeclareUnicodeCharacter{00B3}{\ensuremath{^{3}}}
\DeclareUnicodeCharacter{00B7}{\ensuremath{\cdot}}
\DeclareUnicodeCharacter{00D7}{\ensuremath{\times}}
\DeclareUnicodeCharacter{0393}{\ensuremath{\Gamma}}
\DeclareUnicodeCharacter{03B1}{\ensuremath{\alpha}}
\DeclareUnicodeCharacter{03B6}{\ensuremath{\zeta}}
\DeclareUnicodeCharacter{03BE}{\ensuremath{\xi}}
\DeclareUnicodeCharacter{2013}{--}
\DeclareUnicodeCharacter{2074}{\ensuremath{^{4}}}
\DeclareUnicodeCharacter{2076}{\ensuremath{^{6}}}
\DeclareUnicodeCharacter{207B}{\ensuremath{^{-}}}
\DeclareUnicodeCharacter{2080}{\ensuremath{_{0}}}
\DeclareUnicodeCharacter{2081}{\ensuremath{_{1}}}
\DeclareUnicodeCharacter{2082}{\ensuremath{_{2}}}
\DeclareUnicodeCharacter{2083}{\ensuremath{_{3}}}
\DeclareUnicodeCharacter{2085}{\ensuremath{_{5}}}
\DeclareUnicodeCharacter{2087}{\ensuremath{_{7}}}
\DeclareUnicodeCharacter{2088}{\ensuremath{_{8}}}
\DeclareUnicodeCharacter{2089}{\ensuremath{_{9}}}
\DeclareUnicodeCharacter{210F}{\ensuremath{\hbar}}
\DeclareUnicodeCharacter{2192}{\ensuremath{\to}}
\DeclareUnicodeCharacter{21A6}{\ensuremath{\mapsto}}
\DeclareUnicodeCharacter{2212}{\ensuremath{-}}
\DeclareUnicodeCharacter{2218}{\ensuremath{\circ}}
\DeclareUnicodeCharacter{221A}{\ensuremath{\surd}}
\DeclareUnicodeCharacter{2260}{\ensuremath{\ne}}
\DeclareUnicodeCharacter{2261}{\ensuremath{\equiv}}
\DeclareUnicodeCharacter{27E8}{\ensuremath{\langle}}
\DeclareUnicodeCharacter{27E9}{\ensuremath{\rangle}}
\newcommand{\HMTTitleMapEntry}[2]{%
  \edef\HMTTitleProbe{\detokenize{#1}}%
  \ifx\HMTTitleKey\HMTTitleProbe\def\HMTResolvedTitle{#2}\fi}
\newcommand{\HMTResolvePublicTitle}[1]{%
  \def\HMTResolvedTitle{#1}%
  \edef\HMTTitleKey{\detokenize{#1}}%
  % El prefacio fue aprobado como texto autoral literal; su subtítulo no se
  % retitula ni se normaliza.
  \HMTTitleMapEntry{Complex interaction between two simple machines}{Complex interaction between two simple machines}% ADP102-0001-LITERAL
  \HMTTitleMapEntry{General introduction}{Causal architecture of the continuum problem and mathematical-physical scope of HMT–MD}% ADP102-0002
  \HMTTitleMapEntry{Initial object and causal order: finite scale, observation horizon, and Archimedean publication}{Projective horizon: dependence between scale, boundary, and measure}% ADP102-0003
  \HMTTitleMapEntry{Oriented interval, marks, cells, and closing edge}{Construction of the finite line through order, scale, and conservation of the unit}% ADP102-0004
  \HMTTitleMapEntry{Refinement, survival and thresholds of compatible prolongation}{Limit relation between infinitesimal, infinite prolongation, and projective threshold}% ADP102-0005
  \HMTTitleMapEntry{Archimedean evaluation of compatible sections and reconstruction of \(\mathbb R\)}{Compatible limit of finite horizons and realization of the real line}% ADP102-0006
  \HMTTitleMapEntry{Formal kernel of Triadic Modular Holography}{Construction and operatorial architecture of the APP–TRIT–TPK formal kernel: arithmetic support, ternary orientation, and genealogical dynamics}% ADP102-0007
  \HMTTitleMapEntry{Lifts and carry cocycle}{Residue-quotient extensions and nonadic carry cocycle}% ADP102-0008
  \HMTTitleMapEntry{Central block of \(2\) positions}{Bivariate central block of APP and decomposition into 2 spectral modes}% ADP102-0009
  \HMTTitleMapEntry{Mixed curvature of the \(2\) laws}{Discrete curvature of the interaction between additive and multiplicative leaves}% ADP102-0010
  \HMTTitleMapEntry{Correspondence between \(2\) charts of the same residue class}{Conjugation between the positive and residual charts of the same class modulo 9}% ADP102-0011
  \HMTTitleMapEntry{Interior kernel and gnomonic complement}{Decomposition of the APP orthogram into inner core and gnomonic complement}% ADP102-0012
  \HMTTitleMapEntry{Marked segments, interstices and the nonadic cell}{Construction of the nonadic cell from marked segments and interstices}% ADP102-0013
  \HMTTitleMapEntry{Finite positional chart and decimal residual projection}{Decimal residual projection of the finite positional chart of APP}% ADP102-0014
  \HMTTitleMapEntry{Extraction of the cell of \(9\) phases}{Construction of the nonadic cell of 9 phases via residual projection}% ADP102-0015
  \HMTTitleMapEntry{Position, height and Euclidean division}{Position and height coordinates induced by the Euclidean division}% ADP102-0016
  \HMTTitleMapEntry{Inaugural domain of the nonadic cell and causal exclusion of APP, TRIT, TPK, and analytic constants}{Autonomous arithmetic domain of the nonadic cell and causal precedence with respect to APP–TRIT–TPK and the constants analiticas}% ADP102-0017
  \HMTTitleMapEntry{Defect law at fixed modulus and variable length}{Cyclotomic Defect Invariance at Fixed Modulus Under Length Variation}% ADP102-0018
  \HMTTitleMapEntry{Primary cyclotomic operator}{Cyclotomic operator of APP and action on the 12 A₂-type fibers}% ADP102-0019
  \HMTTitleMapEntry{Localization of the aggregate \texorpdfstring{\(54\)}{54}}{Derivation of the sum-product defect 459−405=54 in the central block of APP}% ADP102-0020
  \HMTTitleMapEntry{Resonance of the incidence operator in rank \(6\)}{Range 6 of the incidence operator and multiplicity of its resonant mode}% ADP102-0021
  \HMTTitleMapEntry{TRIT as the minimal unit of topological information: ternary orientation, carry memory, and temporal order}{TRIT as a ternary oriented unit: internal geometries, temporal order, and dynamic conditions of time crystal}% ADP102-0022
  \HMTTitleMapEntry{Composition and carry cocycle}{Carry cocycle of the oriented ternary composition}% ADP102-0023
  \HMTTitleMapEntry{Universal family of quadratic algebras}{Unified classification of the 3 quadratic algebras of the TRIT}% ADP102-0024
  \HMTTitleMapEntry{Topological Phase Kernel: observable base and integral holonomic state}{TPK: observable base, integral holonomic state, and typed separation of its projections}% ADP102-0025
  \HMTTitleMapEntry{The double base of APP paths}{Bivariant category of APP paths as the base of the integral holonomic TPK state}% ADP102-0026
  \HMTTitleMapEntry{Nonadic phase, dodecaphasic sector and calendar}{Non-split extension C₁₂→C₁₀₈→C₉ and dodecaphase record of the nonadic phase}% ADP102-0027
  \HMTTitleMapEntry{Cylinder and boundary}{Cylindrical compatibility and boundary signature of the integral holonomic state}% ADP102-0028
  \HMTTitleMapEntry{Fiber transport}{Cardinal operator \(T_d\) of transport between enriched fibers with conservation of path and memory}% ADP102-0029
  \HMTTitleMapEntry{Extension relations}{Typed prolongation relation \(\operatorname{Ext}_r\) between enriched fibers}% ADP102-0030
  \HMTTitleMapEntry{Projections and typed losses}{Observable quotients of the integral holonomic state and preserved genealogical fields}% ADP102-0031
  \HMTTitleMapEntry{Operator tree of the container}{Insertion of the 6 fibered families of the integral holonomic state into the TPK operator hierarchy}% ADP102-0032
  \HMTTitleMapEntry{Operator algebra of the Topological Phase Kernel and causal update dynamics of the integral holonomic state}{Typed hierarchy of the 34 heterogeneous canonical constructions and causal dynamics of TPK updating}% ADP102-0033
  \HMTTitleMapEntry{Classification of the integral holonomic state into \(8\) invariant structural classes}{Typed L₀–L₇ hierarchy of the 34 heterogeneous canonical TPK constructions}% ADP102-0034
  \HMTTitleMapEntry{Functional decomposition into \(14\) groupings compatible with the
APP–TRIT–TPK operators}{Functional grouping of the 34 canonical constructions in 14 families of composition APP–TRIT–TPK}% ADP102-0035
  \HMTTitleMapEntry{Domain and typing of the \(34\) entries of the update operator}{Mathematical classes, domains, codomains, and causal dependencies of the 34 canonical constructions}% ADP102-0036
  \HMTTitleMapEntry{Causal composition of the update of the integral holonomic state}{Update of the integral holonomic state through \(U_t=\operatorname{Upd}_t\circ\operatorname{Tra}_t\circ\operatorname{Sel}_t\)}% ADP102-0037
  \HMTTitleMapEntry{Typed separation between transport, memory, calendar, extraction and return operators}{Typed separation between operators, states, projections, and publications of the TPK}% ADP102-0038
  \HMTTitleMapEntry{Operator decomposition of incidence channels \texorpdfstring{\(90/120\)}{90/120}: temporal selector, nilpotent generator and angular component}{Tritic incidence 90/120 and typed separation of the temporal, nilpotent, and angular operators}% ADP102-0039
  \HMTTitleMapEntry{Distribution of the indices \(90/120\) among the \(6\) objects of the incidence system}{Typed residences of the indices 90/120 in 6 non-equivalent operator objects}% ADP102-0040
  \HMTTitleMapEntry{Normalized incidence defect}{Normalized defect 1−90/120=1/4 of the tritic incidence}% ADP102-0041
  \HMTTitleMapEntry{Typing invariants}{Domain invariants, orientation, and memory of the 90/120 realizations}% ADP102-0042
  \HMTTitleMapEntry{Local state and sheet feedback}{Oriented transition local operator and APP sheet update}% ADP102-0043
  \HMTTitleMapEntry{Cylinders of the \(2\) readers}{Cylindrical systems of the 2 arithmetic readers and truncation compatibility}% ADP102-0044
  \HMTTitleMapEntry{Decomposition of the Topological Phase Kernel into advance, rotation, carry, and memory operators}{Factorization of the TPK transition as \(U_t=\operatorname{Upd}_t\circ\operatorname{Tra}_t\circ\operatorname{Sel}_t\)}% ADP102-0045
  \HMTTitleMapEntry{Ordered support and cyclic coordinate}{Cyclic coordinate on the time-oriented 12-sector support}% ADP102-0046
  \HMTTitleMapEntry{Nonsplit cyclic extension}{nonsplit extension 0→C₁₂→C₁₀₈→C₉→0 and carry cocycle}% ADP102-0047
  \HMTTitleMapEntry{Signed event extractor}{Extraction of the signed dodecaphase state from the enriched time record}% ADP102-0048
  \HMTTitleMapEntry{\(2\) dodecaphasic words of different provenance}{Genealogical distinction between 2 dodecaphasic words of equal period}% ADP102-0049
  \HMTTitleMapEntry{Minimal refinement by microinstants}{Minimal temporal subdivision compatible with semiconjugation}% ADP102-0050
  \HMTTitleMapEntry{\(3\) noninterchangeable quaternary operations}{Commutators of 3 quaternary operations and separation of their actions}% ADP102-0051
  \HMTTitleMapEntry{Tritic reader of the oriented state: temporal order and carry transport}{Triticale reader with quotient and carry transport}% ADP102-0052
  \HMTTitleMapEntry{Minimum predictive-memory quotient}{Memory quotient of order 1 and predictive sufficiency of the modal projection}% ADP102-0053
  \HMTTitleMapEntry{Unbounded vertical memory}{Unbounded Growth of Memory Depth Under Phase Returns}% ADP102-0054
  \HMTTitleMapEntry{Geometric realizations of TRIT: elliptic, parabolic and hyperbolic regimes and double temporal orientation}{TRIT as an oriented ternary state: carry memory, temporal order, and 3 quadratic regimes}% ADP102-0055
  \HMTTitleMapEntry{APP, TRIT and TPK: support, orientation and transport}{APP–TRIT–TPK Composition: arithmetic support, ternary orientation, and memory-based transport}% ADP102-0056
  \HMTTitleMapEntry{TPK as a transport system}{TPK as the dynamics of selection, transport, memory, return, and deployment.}% ADP102-0057
  \HMTTitleMapEntry{Temporal periodicity of (108) phases and return laws at (27), (54), and (108) iterations}{Temporal periodicity of 108 phases and returns at 27, 54, and 108 iterations}% ADP102-0058
  \HMTTitleMapEntry{APP, TRIT and TPK: governing definitions}{Generative composition APP→TRIT→TPK and construction of the integral holonomic state}% ADP102-0059
  \HMTTitleMapEntry{Consequences and publications of the constructed continuum}{Joint nonadic connection Γ₉ as a compatible correspondence of the inverse system}% ADP102-0060
  \HMTTitleMapEntry{Conditioning as pruning of futures}{Conditioning of the measure through restriction of the future extension tree}% ADP102-0061
  \HMTTitleMapEntry{Nonadic holonomy \texorpdfstring{\(\Gamma_9\)}{Γ₉} of the Topological Phase Kernel: fiber \texorpdfstring{\(\mathbb F_3^6\)}{F₃⁶}, coinductive prolongation, and phase return with state memory}{Nonadic holonomy Γ₉ of the Topological Phase Kernel: transport of the integral holonomic state, coinductive prolongation, and return 9→1 with memory advance}% ADP102-0062
  \HMTTitleMapEntry{Joint nonadic connection, compression, and terminal memory}{Joint Holonomy Γ₉: observable compression T, orthogonal component η and terminal memory conservation}% ADP102-0063
  \HMTTitleMapEntry{Local relations and holonomy}{Local transition relations and construction of the nonadic holonomy Γ₉}% ADP102-0064
  \HMTTitleMapEntry{Obstructions of the holonomy \(\Gamma_9\): return \(9\to1\), composition indices and preservation of phase, carry, boundary, route and memory}{Obstructions of the holonomy Γ₉: 9→1 return, composition indices and phase conservation, carry, boundary, path and memory}% ADP102-0065
  \HMTTitleMapEntry{Decomposition of the closure microfiber of \texorpdfstring{\(\pi\)}{π} into \texorpdfstring{\(5\)}{5} oriented regions: multiplicity \texorpdfstring{\(432+4\cdot144=1008\)}{432+4·144=1008}, \texorpdfstring{\(\varphi\)}{φ} axis, and mirror involution \texorpdfstring{\(\pi/e\)}{π/e}}{Oriented classification of the closure microfiber of π: 5 regions, multiplicity 432+4·144=1008, φ-axis, and π/e involution}% ADP102-0066
  \HMTTitleMapEntry{Coinductive generation of \texorpdfstring{\(\pi\)}{π}, \texorpdfstring{\(e\)}{e}, and \texorpdfstring{\(\varphi\)}{φ} through compatible nonadic refinement}{Coinductive Generation of π, e and φ via Compatible Nonadic Refinement}% ADP102-0067
  \HMTTitleMapEntry{Visible lift and chain of blocks}{Coinductive prolongation of the exponential emission via compatible blocks}% ADP102-0068
  \HMTTitleMapEntry{Decimal signature and positional square}{Construction of the decimal signature via the ternary–decimal positional square}% ADP102-0069
  \HMTTitleMapEntry{Visible return and memory rupture}{Phase return with non-trivial memory state update}% ADP102-0070
  \HMTTitleMapEntry{Adversarial certificate for the generation of \(e\): orbital uniqueness, factorial recurrence, rational brackets and cylindrical prolongation}{Orbital uniqueness, factorial recurrence, and cylindrical convergence of electron generation}% ADP102-0071
  \HMTTitleMapEntry{Blockwise lift and 1st coinductive boundary}{Blockwise extension and construction of the 1st coinductive boundary}% ADP102-0072
  \HMTTitleMapEntry{From the tritic calendar to the incidence multigraph}{Construction of the incidence multigraph from the triticale calendar}% ADP102-0073
  \HMTTitleMapEntry{Characterization, brackets and prolongation}{Characterization of self-similarity via rational forks and coinductive prolongation}% ADP102-0075
  \HMTTitleMapEntry{Minimal diagonal selection}{Minimal diagonal selection of compatible cylinders and convergence to arbitrary precision}% ADP102-0076
  \HMTTitleMapEntry{Markov transfer compatible with the entire inverse limit}{Markov dynamics as a modal quotient of order-1 memory of the holonomy Γ₉, compatible with the inverse limit}% ADP102-0077
  \HMTTitleMapEntry{2nd coordinate system: step differences \(3\) and \(4\)}{Differential coordinate system induced by step 3 and 4 operators}% ADP102-0078
  \HMTTitleMapEntry{Obstructions of the affine prolongation of \(\alpha\): telescoping, inversion of \(K\), boundary \(662/663\) and isolation of target data}{Obstructions to affine prolongation of α: telescoping, inversion of K, boundary 662/663, and isolation of target data}% ADP102-0079
  \HMTTitleMapEntry{Construction of \texorpdfstring{\(\sqrt{-1}\)}{√(-1)} by means of \texorpdfstring{\(J^2=-I\)}{J²=-I} and coordination of elliptic, parabolic, and hyperbolic geometries}{Construction of √(−1) and the tritic exponential family of Euler: operators J_τ and coordination of the elliptic, parabolic, and hyperbolic regimes}% ADP102-0080
  \HMTTitleMapEntry{The Hensel--Witt Transduction Kernel}{Transduction kernel Hensel–Witt: domain, operation and genealogical preservation}% ADP102-0081
  \HMTTitleMapEntry{Hensel–Witt Transduction Kernel: microregions of \(\pi\), holonomy \(\Gamma_9\), dodecaphasic register and hexad–octad transport}{Hensel–Witt Transduction Kernel: π microregions, holonomy Γ₉, dodecaphase record, and hexad–octad transport}% ADP102-0082
  \HMTTitleMapEntry{From the ternary code to the Witt system}{Translation of the ternary code into the Witt system}% ADP102-0083
  \HMTTitleMapEntry{Hexad, octad and the corridor \texorpdfstring{\(90/120\)}{90/120}}{Hexagonal–Octagonal Inclusion and Incidence Transport through the 90/120 Channels}% ADP102-0084
  \HMTTitleMapEntry{Triadic scale extractor}{Triadic Scale Extraction Operator and Its Rational Image}% ADP102-0085
  \HMTTitleMapEntry{Prime-power defect and number operator}{Representation of Primary Power Defects via the Number Operator}% ADP102-0086
  \HMTTitleMapEntry{Ablations of the selector and continuation}{Uniqueness conditions for the regional selector and the analytic prolongation}% ADP102-0087
  \HMTTitleMapEntry{Generation chain \(\mathcal F_\pi\to b_\pi\to\rho_\pi\to D_A\to C_*^{\rm HMT}\to\Gamma_{6\to8}\): regional selection, determinantal prolongation and oriented coordinate}{Generation chain \(\mathcal F_\pi\to b_\pi\to\rho_\pi\to D_A\to C_*^{\mathrm{HMT}}\to\Gamma_{6\to8}\): regional selection, determinantal prolongation and oriented coordinate}% ADP102-0088
  \HMTTitleMapEntry{Hexad--octad incidence decomposition in the sectors \texorpdfstring{\(q_{90}\)}{q₉₀} and \texorpdfstring{\(q_{120}\)}{q₁₂₀}: intertwiner \texorpdfstring{\(6\to8\)}{6→8} and precursor structure of the area spectrum}{Incidence channels 90/120 of the Topological Phase Kernel: tritic semicrowns, sheet factorization, and hexad--octad realization of the morphism 6→8}% ADP102-0089
  \HMTTitleMapEntry{Dodecaphasic geometry of mixing: CKM--CP angular antecedent and precursor
reader of the electron}{Dodecaphase Mixing Operator: CKM–CP Phase and Construction of the Pre-Dimensional Electron Reader}% ADP102-0090
  \HMTTitleMapEntry{Archimedean and exceptional morphisms from a common enriched domain}{Principle of Ambivalence: involution of the areal and volumetric readers on the integral holonomic state}% ADP102-0091
  \HMTTitleMapEntry{The local operation on the 729 words}{Common enriched domain and inequivalent codomains of the 2 projections}% ADP102-0092
  \HMTTitleMapEntry{Naturality and compatibility at every depth}{Joint separation of histories by the Archimedean and exceptional projections}% ADP102-0093
  \HMTTitleMapEntry{Postgenerative coupling with Euler's operator closure}{Common incidence invariant between the integral holonomic state and the Witt system}% ADP102-0094
  \HMTTitleMapEntry{Natural globalization of the Archimedean and exceptional projections over a common inverse system}{Decomposition of the microfiber of π into 5 oriented regions: 4 peripheral components and 1 fixed axis}% ADP102-0095
  \HMTTitleMapEntry{Compatible systems at finite depth}{Reticular prolongation of incidence and construction of dimensional screens}% ADP102-0096
  \HMTTitleMapEntry{Gauge of the exceptional publication}{Factorization of the incidence reader through the integral holonomic state}% ADP102-0097
  \HMTTitleMapEntry{Joint naturality of the publications of the genealogical state}{Double publication of the integral holonomic state: Archimedean evaluation and incidence preservation}% ADP102-0098
  \HMTTitleMapEntry{Joint separation of value and exceptional incidence}{Uniqueness of the Initial Caliber of the Exceptional Projection}% ADP102-0099
  \HMTTitleMapEntry{Common noncommutative antecedent of the double projection}{Noncommutative Algebra, Common and Factorization of the 2 Projections}% ADP102-0100
  \HMTTitleMapEntry{From the space of histories to the real line}{Quotient of the space of compatible histories and realization of the real line}% ADP102-0101
  \HMTTitleMapEntry{Archimedean mean and double projection of the same state}{Compatibility of the mean Archimedean evaluation with the double projection of the integral holonomic state}% ADP102-0102
  \HMTTitleMapEntry{Unitary calendar, modular memory and stable mode}{Stable Mode of the Unitary Calendar under Modular Memory Update}% ADP102-0103
  \HMTTitleMapEntry{Analytic publications of HMT histories}{Analytic realizations of HMT histories via the Archimedean and bidirectional charts}% ADP102-0104
  \HMTTitleMapEntry{Consequence for a positive action}{Positivity of the Action Induced by the Dodecaphasic Braider–Witt}% ADP102-0105
  \HMTTitleMapEntry{From the spectral jump \texorpdfstring{\(4\)}{4} to the Leech lattice: APP unfolding \texorpdfstring{\(4\to2\to6\to54\)}{4→2→6→54} and the Paley--Witt--Mathieu--Golay--Niemeier corridor}{From spectral jump 4 to the Leech lattice: APP unfolding 4→2→6→54 and the Paley–Witt–Mathieu–Golay–Niemeier corridor}% ADP102-0106
  \HMTTitleMapEntry{Nature of the structure obtained}{Algebraic classification of the self-dual ternary structure obtained}% ADP102-0107
  \HMTTitleMapEntry{Incidence law \texorpdfstring{\(4+1\)}{4+1} on a tetrad}{4+1 Decomposition of the incidence of a tetrad into peripheral and axial components}% ADP102-0108
  \HMTTitleMapEntry{Discriminant module and form}{Decomposition of the incidence of a tetrad into 4 peripheral components and 1 axial component}% ADP102-0109
  \HMTTitleMapEntry{\(12\) factors and gluing map}{Correspondence between incidence strata and free orbit of the positive triple}% ADP102-0110
  \HMTTitleMapEntry{Character and congruence kernel}{Discriminant gluing of 12 A₂ factors and construction of the integral extension}% ADP102-0111
  \HMTTitleMapEntry{Mixed APP curvatures}{Ordered system of incidences and radial camera orientation}% ADP102-0112
  \HMTTitleMapEntry{Weyl--Heisenberg representation and central phase}{Construction of the radial character and determination of its congruence kernel}% ADP102-0113
  \HMTTitleMapEntry{Exact status of the hierarchy \(4\to2\to6\to54\) and of the lattice
identification.}{Exact status of the hierarchy 4→2→6→54 and of the lattice identification.}% ADP102-0114
  \HMTTitleMapEntry{Graded trace of the ternary generator}{Graded trace of the action of the ternary generator on \(V^{\natural}\)}% ADP102-0115
  \HMTTitleMapEntry{Full-support word and fixed-point-free isometry}{Isometry Without Fixed Points Induced by the Ternary Word of Full Support}% ADP102-0116
  \HMTTitleMapEntry{Construction of Hilbert space from the nonadic connection and its phase operators}{Projective realization of Bloch from the internal complex structure}% ADP102-0117
  \HMTTitleMapEntry{The marked branch changes type}{Corner immersion of a marked phase: compact ideal \(\mathcal K(\mathcal H_\infty)\) and factor \(\mathcal B(\mathcal H_\infty)\) of type \(\mathrm I_\infty\)}% ADP102-0118
  \HMTTitleMapEntry{Visible residue and boundary quotient}{Operational extension of APP: visible modes and memory modes}% ADP102-0119
  \HMTTitleMapEntry{T-duality involution}{Completion 1000=729+243+27+1 and stratification of the unit}% ADP102-0120
  \HMTTitleMapEntry{Stratified completion of the unit}{Decomposition 243=1+22+90+120+10 and typing of its components}% ADP102-0121
  \HMTTitleMapEntry{The block \texorpdfstring{\(243\)}{243} as boundary and perfect ternary ball}{Derivation of the constant channels from the decomposition of block 243}% ADP102-0122
  \HMTTitleMapEntry{Internal decomposition of the block \texorpdfstring{\(243\)}{243}}{Construction of dimension 11 via compatible residual readers}% ADP102-0123
  \HMTTitleMapEntry{Dimension \texorpdfstring{\(11\)}{11} and compatible readings}{Closure of the genealogical register under the T-duality involution}% ADP102-0124
  \HMTTitleMapEntry{Verifiable identities of the construction}{Exact identities of involution, modularity, and residual conservation}% ADP102-0125
  \HMTTitleMapEntry{Level, linear mass exchange, and modularity: \(3\) nonfusion controls}{Level separation, linear mass exchange, and modularity of the 3 realizations}% ADP102-0126
  \HMTTitleMapEntry{\(2\) descendants of the extended ternary code}{Dimensional reduction 26→24→11→10 and coordination with M-theory screens}% ADP102-0127
  \HMTTitleMapEntry{The icosahedral representation of dimension \(12\)}{Coordination of exact screens with the physical realization of M theory}% ADP102-0128
  \HMTTitleMapEntry{Isotropic reduction of rank \(26\to24\)}{Operatorial representation of APP and supersymmetric spectral triple in dimension 11}% ADP102-0129
  \HMTTitleMapEntry{Coordination of the Moonshine corridor with the M-theory screens: grading,
Lorentzian lattice, and dual realization of the HMT state}{Compatibility theorem between the Moonshine corridor and M-theory screens through grading and lattice transport}% ADP102-0130
  \HMTTitleMapEntry{Common antecedent of readers: spectral refinement and joint-dilation limit}{Common spectral refinement and limit of the joint dilation of matrix readers}% ADP102-0131
  \HMTTitleMapEntry{Substructures, global object, and projections}{Construction of 4 cyclic quantum Latin squares in an elementary operator class}% ADP102-0132
  \HMTTitleMapEntry{Direct and inverse correspondences between the formalisms}{Realization of 3 operator contexts in a quantum square (9,3)}% ADP102-0133
  \HMTTitleMapEntry{Extension of the bigrading to 3 indices}{Representation of the complete genealogical atlas in a quantum square (12,3)}% ADP102-0134
  \HMTTitleMapEntry{4 cyclic quantum Latin squares of easy operator class}{Cyclic classification of 4 quantum Latin squares into an elementary operator class}% ADP102-0135
  \HMTTitleMapEntry{Cyclic lemma for quantum magic squares associated with a positive operator-valued measure}{Bigraduation of the genealogical-matrix family}% ADP102-0136
  \HMTTitleMapEntry{Matrix cones, ranks, and completely positive extensions}{Matrix representation of the genealogical composition APP–TRIT–TPK}% ADP102-0137
  \HMTTitleMapEntry{Explicit construction of quantum magic cubes and compatibility of their margins}{TPK as a category of routes with selection, transport, memory, and return.}% ADP102-0138
  \HMTTitleMapEntry{Magic combinations of CP maps}{Completely positive compositions compatible with the genealogical-matricial bidegree}% ADP102-0139
  \HMTTitleMapEntry{Spectral hierarchy \(4\to2\to6\to54\) of the central APP block}{Subsequent metrological realization of the double projection of the integral holonomic state}% ADP102-0140
  \HMTTitleMapEntry{Nonadic refinement through conditional expectation}{Symplectic decomposition 12→4×3 of the phase coordinates}% ADP102-0141
  \HMTTitleMapEntry{Functional taxonomy of \(114\) numerical expressions and \(106\) operator readings}{Operator classification of 114 numerical expressions by domain, generator, invariant, and codomain}% ADP102-0142
  \HMTTitleMapEntry{Phase transfer on the multigraph of complete blocks}{Closure microfiber of π: 5 oriented regions, multiplicities \(432+4\cdot144=1008\), and common projection \(w_\pi\)}% ADP102-0143
  \HMTTitleMapEntry{\(5\) oriented closure cycles of the microfiber}{Publication of a reduced memory of a complete nonad block via the phase transfer multigraph}% ADP102-0144
  \HMTTitleMapEntry{Incidence operator and oriented half-turn}{Oriented holonomies of the closure of the 5 regions of π}% ADP102-0145
  \HMTTitleMapEntry{Rational composition and compensator \texorpdfstring{\(239\)}{239}}{Determination of the slope 1/5 by the multiplicity of the regional microfiber of π}% ADP102-0146
  \HMTTitleMapEntry{Gaussian identity and \(1^{\mathrm a}\) oriented half-turn}{Rational composition of the half-turn and determination of the compensating term 239}% ADP102-0147
  \HMTTitleMapEntry{Coinductive prolongation of the section of \(\pi\) through the inverse system of surviving histories}{Filtration of F₃⁶ by regional admissibility and selection of compatible subcylinders}% ADP102-0148
  \HMTTitleMapEntry{Inverse system of compatible histories}{Uniqueness of the Minimal Diagonal Section in the Rational Fork System}% ADP102-0149
  \HMTTitleMapEntry{Decimal positional signature and compatibility square}{Exponential Section Prolongation via the Visible Block Chain}% ADP102-0150
  \HMTTitleMapEntry{Phase return with a change of state and conservation of traversed memory}{Construction of the decimal signature and the positional square of the exponential section}% ADP102-0151
  \HMTTitleMapEntry{Factorial semigroup of the exponential character: recurrence \(a_n=1/n!\)}{Phase return with inequality of integral holonomic states: \(g(K+9)=g(K)\) and \(B_{K+9}\ne B_K\)}% ADP102-0152
  \HMTTitleMapEntry{Internal certificate of orbital uniqueness, factorial recurrence, convergence, and cylindrical prolongation}{Inductive prolongation of e via the nonadic holonomy Γ₉ of the TPK}% ADP102-0153
  \HMTTitleMapEntry{Invariant positive ray and uniqueness of self-scaling}{Construction of the modal incidence multigraph from the triticale calendar}% ADP102-0154
  \HMTTitleMapEntry{Internal certificate of uniqueness, compatibility and convergence of Perron self-scaling}{Characterization of φ by the Perron ray, fork embedding, and coinductive prolongation}% ADP102-0156
  \HMTTitleMapEntry{Internal derivations of the fine-structure constant \(\alpha\): reversible integer dodecaphasic closure and unique analytic root of the kernel \(E_{21/22}\)}{dodecaphasic closure of \(\alpha\) and analytic characterization by means of \(E_{21/22}\): reversible generation, simple root and nonadic prolongation of order 9}% ADP102-0157
  \HMTTitleMapEntry{Common terminal state and 2 coordinated internal readings}{Terminal and dodecaphasic evaluations of the common integral holonomic state: blocks \(B_{\mathrm{term}}\) and signed chart \(U_{12}^{\mathrm{sgn}}\)}% ADP102-0158
  \HMTTitleMapEntry{Path A: dodecaphasic integer closure}{Reversible generation of \(\alpha\) through the integral dodecaphasic closure}% ADP102-0159
  \HMTTitleMapEntry{Hensel--Witt domain, Teichmüller maps, and rational target \(-1/12\)}{Hensel–Witt transduction operator between microregions of π, holonomy Γ₉, the dodecaphasic register and hexad–octad incidence}% ADP102-0160
  \HMTTitleMapEntry{Teichmüller map from the ternary code to the Witt ring}{Construction of the Witt system from the compatible ternary coding}% ADP102-0161
  \HMTTitleMapEntry{Hexad--octad incidence and modular channels \(90/120\)}{Hexad--octad realization of the 90/120 channels previously derived from the tritic incidence of the TPK}% ADP102-0162
  \HMTTitleMapEntry{4 compatible routes to \texorpdfstring{\(-1/12\)}{-1/12}}{4 operational realizations of −1/12 and separation of their domains}% ADP102-0163
  \HMTTitleMapEntry{Determinantal derivation of the action scale: Smith normal form, cokernel of order \(16\) and valuation \(10^{-34}\)}{Determinantal derivation of the absolute action scale: Smith normal form, cokernel of order 16, and valuation 10⁻³⁴}% ADP102-0164
  \HMTTitleMapEntry{Quotient measure and positive character of action: canonical Hilbertization, regional--dodecaphasic defect and functorial law on paths}{Construction of the positive action measure over the quotient of emissions: Hilbertization, regional--dodecaphasic defect, and functoriality of paths}% ADP102-0165
  \HMTTitleMapEntry{Self-dual derivation of the Planck constant: temporal periodicity of \(108\) phases, action and Compton wavelength}{Self-dual derivation of the Planck constant and its action line: temporal periodicity of 108 phases, Compton closure and bidirectional covariance}% ADP102-0166
  \HMTTitleMapEntry{Equivalence of normalizations on the self-dual identity}{Commutativity of the angular, temporal, and gravitational normalizations of the self-dual identity}% ADP102-0167
  \HMTTitleMapEntry{Predimensional angular operator: decomposition into \(27\) sectors, spectral radius and effective composition of the double projection}{Phase operator prior to dimensional transduction: decomposition into 27 sectors, spectral radius, and factorization of the double projection}% ADP102-0168
  \HMTTitleMapEntry{Quadratic characters and spectral constants: Catalan, Apéry and the Euler--Stieltjes tower}{quadratic characters and spectral constants: Catalan, Apery, Euler–Mascheroni and Stieltjes hierarchy}% ADP102-0169
  \HMTTitleMapEntry{Unimodal criticality of the dodecaphasic family: exact analytic profile, Feigenbaum approximants and reduction of universality to the hyperbolic certificate}{Unimodal criticality of the dodecaphase family: exact analytic profile, Feigenbaum approximants, and hyperbolic renormalization certificate}% ADP102-0170
  \HMTTitleMapEntry{Arithmetic of mode \(30\): Rogers--Ramanujan identities, normalizer \(899\) and nonadic Pell lift}{Arithmetic of the dodecaphasic modular family of period 30: Rogers--Ramanujan identities, normalizer 899, and nonadic Pell lift}% ADP102-0171
  \HMTTitleMapEntry{Governing theorem of the radical lattice and the power--vacancy law}{Classification of the radical network 3→6→9 and derivation of the power–vacancy law}% ADP102-0172
  \HMTTitleMapEntry{Cyclic orbit \(369\to693\to936\to369\)}{Cyclic orbit 369→693→936→369}% ADP102-0173
  \HMTTitleMapEntry{Kernel \(7227\), spectral decomposition and link to the radical cross}{Determination of the arithmetic core 7227 by incidence with the radical cross of APP}% ADP102-0174
  \HMTTitleMapEntry{Mechanical coding of vacancies through nonadic powers}{Sturmian coding of vacancies induced by nonadic powers}% ADP102-0175
  \HMTTitleMapEntry{Internal root of \(-1\), exponential structure, and elliptic, parabolic, and hyperbolic geometries}{Triadic exponential family \(J_\tau^2=-\tau I\): elliptic, parabolic, and hyperbolic realizations of Euler's identity}% ADP102-0176
  \HMTTitleMapEntry{Continuous realization of the induced quadratic type}{Continuous extension of the common quadratic operator and conservation of its spectral type}% ADP102-0177
  \HMTTitleMapEntry{Common operator support of \(3\) geometric readings}{Conjugation of the arithmetic, angular, and torsional realizations via the common quadratic operator}% ADP102-0178
  \HMTTitleMapEntry{Constitutive chain of the gravitational sector: chronological scale \(108\), tensor reduction, holonomy and metrological evaluation of \(G\)}{Causal composition of the phase, action, and length invariants in the derivation of G}% ADP102-0179
  \HMTTitleMapEntry{Incidence cardinalities 90 and 120}{Tritic derivation of the cardinalities 90/120 and subsequent realization in hexad--octad incidence}% ADP102-0180
  \HMTTitleMapEntry{Causal separation between the area spectrum, modular memory and the Cartan--Holst realization}{Typed distinction between Barbero–area, area–modular Hamiltonian, and Cartan–Holst realization}% ADP102-0182
  \HMTTitleMapEntry{Unfiltered Hensel transducer and full shift on the ternary fibers}{Hensel's Integral Transducer as a Full Shift Operator}% ADP102-0183
  \HMTTitleMapEntry{Surjectivity of the unfiltered Archimedean chart onto \(\mathbb R\)}{Exact coverage of the preceding real chart via compatible sections}% ADP102-0184
  \HMTTitleMapEntry{2 classification axes: visible value and transport genealogy}{Bifactorial classification of the number by Archimedean value and transport genealogy}% ADP102-0185
  \HMTTitleMapEntry{Position of \texorpdfstring{\(\pi,e,\varphi\) and \(\alpha\)}{pi, e, phi, and alpha} in the classification}{Position of π, e, φ and α in the classification}% ADP102-0186
  \HMTTitleMapEntry{Unfiltered Hensel chart, survival, and global boundary}{Previous chart, inductive survival, and construction of the global boundary}% ADP102-0187
  \HMTTitleMapEntry{The Euler operator family as a closure grammar}{Factorization of the genealogical closure of the number through Euler's operational system}% ADP102-0188
  \HMTTitleMapEntry{Mathematical dependencies.}{Causal dependence between inverse section, Archimedean evaluation, and exceptional conservation of incidence}% ADP102-0189
  \HMTTitleMapEntry{\(4\) choice requirements in the projective system}{Necessary choice conditions for the global prolongation of compatible sections}% ADP102-0190
  \HMTTitleMapEntry{Effective decidability of each finite window of projective multiplicities}{Algorithmic decidability of each finite level of the prefix system}% ADP102-0191
  \HMTTitleMapEntry{Exact relationship with the \(9\) phases of the Topological Phase Kernel}{Restriction of the uniform obstruction to the nonadic temporal system of 9 phases}% ADP102-0192
  \HMTTitleMapEntry{From constructible consistency to independence by forcing}{Relative consistency in L and independence via forcing}% ADP102-0193
  \HMTTitleMapEntry{Minimax optimization over finite windows of compatible prefixes}{Minimax theorem at each finite level of the poset of compatible prefixes}% ADP102-0194
  \HMTTitleMapEntry{The chain that the bridge receives as proved}{Theoretical Background of the Genealogical-Metrical Transduction Functor}% ADP102-0195
  \HMTTitleMapEntry{Bidirectional derivation of the electronic exponent from the enriched route}{Electron realization band and separation between base mass and beta closure}% ADP102-0196
  \HMTTitleMapEntry{\(2\) causal branches and their confluence}{Confluence of the determinantal and dodecaphasic branches of the action scale}% ADP102-0197
  \HMTTitleMapEntry{Subsequent metrological chart and dimensional evaluation of the generated invariants}{Systems of transduction between genealogical number and physical magnitude}% ADP102-0198
  \HMTTitleMapEntry{Ambivalence principle: involution of the areal and volumetric readers,
local equivalence, and conjugate temporal orientation}{Ambivalence Principle in Dimensional Mechanics: areal--volumetric involution, autoinertial systems, and conjugate temporal orientation}% ADP102-0199
  \HMTTitleMapEntry{Conditions for a prospective evaluation}{Conditions for prospective evaluation of the Principle of Ambivalence on the integral holonomic state}% ADP102-0200
  \HMTTitleMapEntry{Finite HMT time crystal: drive, subharmonic response, and spectral stability}{Finite-time Temporal Crystal HMT: Periodically forced Hamiltonian, Subharmonic response and Spectral stability}% ADP102-0201
  \HMTTitleMapEntry{Drive constructed from the 108-state clock}{Construction of the periodic forcing from the 108-state temporal system}% ADP102-0202
  \HMTTitleMapEntry{Finite spectral stability}{Nilpotent jet of the triadic evolution and compatible prolongation}% ADP102-0203
  \HMTTitleMapEntry{HMT–MD confluence theorem: electrogravitational form, Einstein–Schrödinger system and coordination of action, matter, fields, geometry and information}{HMT–MD confluence theorem: Einstein–Schrodinger system and electrogravitational coordination of action, matter, fields, geometry, and information}% ADP102-0204
  \HMTTitleMapEntry{Horizontal independence and vertical compatibility}{Naturality of the realizations and vertical preservation of the invariants of the integral holonomic state}% ADP102-0205
  \HMTTitleMapEntry{1st confluence: coupling, vacuum and the electroweak sector}{Constitutive coupling between the vacuum and the electroweak sector}% ADP102-0206
  \HMTTitleMapEntry{2nd Confluence: Action, Gravity, and Electrogravitational Form}{Reciprocity between action and gravity in the electrogravitational form}% ADP102-0207
  \HMTTitleMapEntry{3rd confluence: geometry, information and mixing}{Confluence between geometry, information and mixing}% ADP102-0208
  \HMTTitleMapEntry{4th Confluence: Screen, Torsion, and Helical Memory}{Transduction of the helical memory in the torsional realization}% ADP102-0209
  \HMTTitleMapEntry{From the electroweak determinant to the positive constitutive vacuum operator and its spectral completion \(3\to4\)}{From the electroweak determinant to the positive vacuum constitutive operator and its spectral completion 3→4}% ADP102-0210
  \HMTTitleMapEntry{Published value and fiber memory}{Spectral decomposition of the path measure and conditioning of futures}% ADP102-0211
  \HMTTitleMapEntry{Algebraic foundations of phase sewing and self-measurement}{Joint observer–observable state, reversible instrument, and self-measurement}% ADP102-0212
  \HMTTitleMapEntry{Even information and odd geometric transport in the bilayer}{Odd-even decomposition of the bilayer state: interleaf information and oriented geometric transport}% ADP102-0214
  \HMTTitleMapEntry{\texorpdfstring{\(3\)}{3} realizations of genealogical memory: dodecaphase transport, non-local Green function with a local generator and cosmological expansion}{Dynamic, propagatory, and cosmological realizations of genealogical memory: dodecaphasic transport, Green function, and spatial expansion}% ADP102-0215
  \HMTTitleMapEntry{Constitutive chain of the gravitational sector: chronological scale \texorpdfstring{\(108\)}{108}, tensor reduction, holonomy and metrological evaluation of \texorpdfstring{\(G\)}{G}}{Icosahedral descent A₅/C₅ and gravitational carrier selection}% ADP102-0216
  \HMTTitleMapEntry{Intertwiner with symmetric traceless tensors and transverse reduction \texorpdfstring{\(12\to5\to5\to2\)}{12→5→5→2}}{Intertwiner with traceless symmetric tensors and transverse reduction 12→5→5→2}% ADP102-0217
  \HMTTitleMapEntry{Structural criteria for the gravitational selector: dimensional transducer, enriched holonomy, and separation between torsional and circular selectors}{Gravitational holonomic selector: dimensional transduction, naturality, refinement, and separation of the torsional and circular regimes}% ADP102-0218
  \HMTTitleMapEntry{Declared decimal chart and metrological comparison}{Posterior metrological recognition of G and causal independence of the derivation}% ADP102-0219
  \HMTTitleMapEntry{Tetrad, connection, and \(2\) defects}{Tetrad, spin connection, and torsion of the Cartan–Holst realization}% ADP102-0220
  \HMTTitleMapEntry{1st-order action}{Palatini–Holst Action and Variational Equations}% ADP102-0221
  \HMTTitleMapEntry{Horizon–area–information–entropy–action confluence}{Theorem of Confluence Between Horizon Energy, Area, Information, Entropy, and Action}% ADP102-0222
  \HMTTitleMapEntry{\(2\) orientations of the same propagation}{Advanced and Retarded Propagators on Oriented Paths}% ADP102-0223
  \HMTTitleMapEntry{Elliptical orbit, action and \(3\) J invariants}{Temporal Involution Invariants and Particle–Antiparticle Conjugation}% ADP102-0224
  \HMTTitleMapEntry{Cosmology of periodicity \texorpdfstring{\(108\)}{108}: sheet inversion, horizon, Big Bang, torsional bounce and dark sector}{Torsional Cosmology of Periodicity 108: Expansion, Memory-Return, and Einstein–Cartan Bounce}% ADP102-0225
  \HMTTitleMapEntry{Typological separation of the \(2\) cosmological realizations}{Separation between de Sitter acceleration and Einstein–Cartan torsional bounce}% ADP102-0226
  \HMTTitleMapEntry{\(3\) exact foliations of de Sitter}{Cosmological foliations compatible with sheet inversion}% ADP102-0227
  \HMTTitleMapEntry{Universe with boundary}{Boundary state of the cosmological horizon and return memory}% ADP102-0228
  \HMTTitleMapEntry{Spectral data of a supersymmetric realization: \((\mathcal H_{\rm phys},H,\mathcal Q)\), supercharge and compactification in dimension \(11\)}{Realization of M-theory through compactification, strings, branes, and 11 dimensions}% ADP102-0229
  \HMTTitleMapEntry{General Mass Law on the atlas \texorpdfstring{\(81/56/13/324\)}{81/56/13/324}: trajectories, register, signature, character, towers \texorpdfstring{\(90/120\)}{90/120} and null sector}{General Mass Law over the 81/56/13/324 atlas: trajectories, record, signature, character, harmonic semigroup ⟨90,120⟩, and null sector}% ADP102-0230
  \HMTTitleMapEntry{Universal material-realization algorithm and spectral equalizer of the readers \(K_D\) and \(K_\Omega\)}{Functorial construction of mass families over the atlas}% ADP102-0231
  \HMTTitleMapEntry{Null bifurcation and positivity}{Bifurcation of the previous null sector from the positive mass branch}% ADP102-0232
  \HMTTitleMapEntry{Gateway to the exhaustive domain}{Exhaustive construction of the atlas 81/56/13/324 and conservation of the null sector}% ADP102-0233
  \HMTTitleMapEntry{Internal domain and realization problem}{Theorem of material domain exhaustiveness and null sector stability under refinement}% ADP102-0234
  \HMTTitleMapEntry{Fiber product between descriptors and routes}{Fiber product between species invariants and genealogical paths}% ADP102-0235
  \HMTTitleMapEntry{Bidirectional readers \texorpdfstring{\(K_D\)}{K D} and \texorpdfstring{\(K_\Omega\)}{K Omega}: profiles, coincidences, \texorpdfstring{\(90/120\)}{90/120} towers, and operators by multisection}{Bidirectional readers \(K_D\) and \(K_\Omega\): profiles, coincidences, harmonic semigroup \(\langle90,120\rangle\), and operators by multisection}% ADP102-0236
  \HMTTitleMapEntry{Readers \(K_D\) and \(K_\Omega\): dodecaphasic profiles, path domain and spectral codomain}{Readers \(K_D\) and \(K_\Omega\): dodecaphasic profiles, path domain and spectral codomain}% ADP102-0237
  \HMTTitleMapEntry{Spectral, scalar and null realizations of the General Mass Law: material families, CKM--PMNS transport and photonic and gluonic sectors}{Spectral, scalar, and null realizations of the General Mass Law: material families, CKM quark mixing, PMNS neutrino mixing, and photonic and gluonic null sectors}% ADP102-0238
  \HMTTitleMapEntry{Color, quark flavors, and CKM mixing}{Fiber, orbit, and multisection mass operators of the atlas}% ADP102-0239
  \HMTTitleMapEntry{Composite states and binding law}{Correspondence by domain and codomain between the material inventory and the realizations of the atlas}% ADP102-0240
  \HMTTitleMapEntry{Classifier of material species by fibre, representation, charge, spin, flavour and generation}{Classification of material species by fiber, representation, charge, spin, flavor, and generation}% ADP102-0241
  \HMTTitleMapEntry{Evidentiary typing of persistence, atlas and families.}{Persistence and correspondence conditions between atlases, families, and representations}% ADP102-0242
  \HMTTitleMapEntry{Material ontology, charge, spin, statistics, color, and flavor}{Structural definitions of particle, mass, charge, spin, statistics, color, flavor and generation}% ADP102-0243
  \HMTTitleMapEntry{Typed closure theorem for the mass realization}{Theorem of the Closure of the Material Realization via Persistent Sections and Finite Representations}% ADP102-0244
  \HMTTitleMapEntry{Metrological comparison and exhaustive archive of realizations}{Posterior metrological recognition and exhaustive inventory of the material realizations}% ADP102-0245
  \HMTTitleMapEntry{Common natural architecture of the terminal realizations of the genealogical continuum}{Joint factorization theorem of the \(5\) terminal realizations through the discrete structure of the continuum}% ADP102-0246
  \HMTTitleMapEntry{Terminal publication of the unique continuum through \(5\) non-permutable readers}{Terminal publication through 5 non-permutable readers and common factor of the integral holonomic state}% ADP102-0247
  \HMTTitleMapEntry{Dynamic closures of the continuum: solenoidal transport and gauge quotient}{Dynamics and curvature of the ancestral lineage continuum preceding classical realizations}% ADP102-0248
  \HMTTitleMapEntry{Dynamics and curvature of the genealogical continuum}{Joint Factorization of the Terminal Realizations from the Same State of the Continuum}% ADP102-0249
  \HMTTitleMapEntry{Algebraic closures of the continuum and limits of extensional reconstruction}{Common Factorization, Extensional Consistency, and Limits of Terminal Reconstruction}% ADP102-0250
  \HMTTitleMapEntry{The example of von Neumann ordinals}{Ordinal von Neumann Representation as a Model of Extensional Consistency}% ADP102-0251
  \HMTTitleMapEntry{Domain, maps, and compatibilities of the common terminal object}{Internal view \(G_T^{\mathrm{int}}\) of the continuum and map \(\operatorname{App}_{T,d_T}\): preservation of fiber, carry, memory and terminal}% ADP102-0252
  \HMTTitleMapEntry{Internal unfolding of the correlative realizations before their identification with classical problems}{Correlative emergence of the 5 terminal actions over the same state of the continuum}% ADP102-0253
  \HMTTitleMapEntry{\(5\) non-permutable terminal readers of the genealogical continuum}{Non-permutability Theorem of the 5 terminal readers and classification of their codomains}% ADP102-0254
  \HMTTitleMapEntry{Obstructions to the ablative projection of the continuum: loss of phase, orientation, or genealogical memory}{Positive Projectors of the Guinand–Weil Form and Spectral Localization}% ADP102-0255
  \HMTTitleMapEntry{Fixing the critical line through the functional involution \(s\mapsto1-\bar s\)}{Fixing of the critical line by the functional involution \texorpdfstring{\(s\mapsto 1-\bar{s}\)}{s -> 1 - conjugate(s)}}% ADP102-0256
  \HMTTitleMapEntry{Typed separation of the zeros of \(\zeta\), \(\xi\), and the spectral determinant}{Typed separation of the zeros of ζ, ξ and the spectral determinant}% ADP102-0257
  \HMTTitleMapEntry{Positive representative of the null residual class: \(n\equiv0\pmod 9\) and \(\rho_9(n)=9\)}{Positive representative of the null residual class: \(n\equiv0\pmod 9\) and \(\rho_9(n)=9\)}% ADP102-0258
  \HMTTitleMapEntry{Twisted transport isometry \(\widetilde U_k\) on cylindrical states, projective weights, and Weil memory}{Twisted transport isometry \(\widetilde U_k\) on cylindrical states, projective weights, and Weil memory}% ADP102-0259
  \HMTTitleMapEntry{Graded invariant of the spectral corridor \(4\to2\to6\to54\)}{Graded invariant of the spectral runner 4→2→6→54}% ADP102-0260
  \HMTTitleMapEntry{Coupled trace of the Weil spectrum and the Moonshine grading}{Independent construction of the spectral field over the double circle}% ADP102-0261
  \HMTTitleMapEntry{Hydrodynamic domain and Galerkin publication}{Solenoidal domain and conservation of zero divergence}% ADP102-0262
  \HMTTitleMapEntry{Unique starting datum and dense cores}{Initial data, boundary compatibility and regular propagation}% ADP102-0263
  \HMTTitleMapEntry{Finite realization of the Kalman–Yakubovich–Popov inequality and auxiliary
function in energy closure}{KYP Inequality and Coercive Control of the Solenoidal Evolution}% ADP102-0264
  \HMTTitleMapEntry{Material extraction of the Sobolev bound}{Sobolev Estimation and Exclusion of Singular Concentration}% ADP102-0265
  \HMTTitleMapEntry{The incidence fibre is \(1\) a physical coordinate of the complete state}{Fiber curvature and coercivity of the gauge action}% ADP102-0266
  \HMTTitleMapEntry{Simple vacuum and a witness attaining the gap}{Uniqueness of the gauge-invariant vacuum and gap saturation \(3\kappa_{\mathrm{av}}\) by the eigenvector \(W_c\)}% ADP102-0267
  \HMTTitleMapEntry{\(4\) Osterwalder–Schrader reconstructions of common dynamics}{Osterwalder–Schrader Positivity and Quantum Reconstruction}% ADP102-0268
  \HMTTitleMapEntry{Euclidean gauge theory and Yang--Mills gap}{Spectral gap of the gauge sector and projective limit}% ADP102-0269
  \HMTTitleMapEntry{Generator-by-generator presentation of the new chart}{Cohomological Generators and Compatible Projectors}% ADP102-0270
  \HMTTitleMapEntry{Initial fiber, Mittag--Leffler and affirmative publication}{Lefschetz–Murre and algebraic descent of Hodge classes}% ADP102-0271
  \HMTTitleMapEntry{Bocksteins on all pages}{Bockstein homomorphisms of the spectral sequence: compatibility between pages and determination of the order in T of the universal complex}% ADP102-0272
  \HMTTitleMapEntry{Complete Kato--FERS table before the determinant}{Kato–Fermat Determinant, Rank, and Leading Term of BSD}% ADP102-0273
  \HMTTitleMapEntry{Separation between structure and application}{Internal view \(G_T\) and terminal reader \(P_T\): logical independence between structural existence and classical realization}% ADP102-0274
  \HMTTitleMapEntry{Obstructions common to the \(5\) terminal publications: loss of domain, naturality, certificate, or compatible limit}{Common Obstructions of Terminal Readers: Loss of Domain, Naturalness, or Compatible Limit}% ADP102-0275
  \HMTTitleMapEntry{Assisted formalization of the algebraic kernel}{Formal Verified Formalization of the Algebraic Core and Certification of Finite Identities}% ADP102-0276
  \HMTTitleMapEntry{Correspondence between scientific statements, proofs, certificates, and theorem sources}{Correspondence between statements, proofs, certificates, and primary scientific sources}% ADP102-0277
  \HMTTitleMapEntry{Typing and data representation conventions}{Mathematical Conventions and Data Representation}% ADP102-0278
  \HMTTitleMapEntry{Typed heterogeneity and catalog coverage}{Heterogeneity of Physical Domains and Exhaustiveness of the Catalogue}% ADP102-0279
  \HMTTitleMapEntry{Individual records of the oriented paths}{Complete unfolding of the 324 oriented routes: correspondence among 53 canonical fields, 57 enriched fields, and 49 invariants published per section}% ADP102-0280
  \HMTTitleMapEntry{Universal typed inventory of 471 masses and 384 widths}{Exhaustive inventory of 471 masses and 384 widths: identification, provenance, and correspondence with the atlas}% ADP102-0281
  \HMTTitleMapEntry{Source, deduplication, and type limits}{Data provenance, record uniqueness, and validity domains}% ADP102-0282
  \HMTTitleMapEntry{Object code and HMT--MD reader}{Identification of objects and application of the HMT--MD reader}% ADP102-0283
  % Correspondencias espirales: sustituciones científicas unívocas aprobadas
  % por cotejo focal. Se modifican sólo los rótulos; los cuerpos permanecen
  % íntegros en sus residencias distribuidas.
  \HMTTitleMapEntry{TPK: selection, transport, update, and return}{Dynamic composition of the Topological Phase Kernel on the integral holonomic state: tritic selection, cardinal transport, memory update, and typed return}% ADP102-0284
  \HMTTitleMapEntry{The typed transition}{Definition and factorization of the enriched transition operator \(U_t\)}% ADP102-0285
  \HMTTitleMapEntry{Composition of carry}{Bimodular cocycle law of the carry under path concatenation}% ADP102-0286
  \HMTTitleMapEntry{Dodecaphase calendar and non-splitting}{Nonsplit cyclic extension \(0\to C_{12}\to C_{108}\to C_9\to0\) and internal chronology of the TPK}% ADP102-0287
  \HMTTitleMapEntry{Nonadic helix of phase and depth}{Helical covering of the nonadic phase and integral coordinate of memory depth}% ADP102-0288
  \HMTTitleMapEntry{Normalized geometric immersion}{Immersion \(H_9\) of the nonadic quotient--residue into \(\mathbb R^3\) and calculation of the Frenet invariants}% ADP102-0289
  \HMTTitleMapEntry{Nested circles and cylinders}{Circular sections and cylindrical foliation of the phase--memory covering}% ADP102-0290
  \HMTTitleMapEntry{Orbital sheets and seam of the suspension space}{Orbits of the skew product and identification relation of the suspension space}% ADP102-0291
  \HMTTitleMapEntry{Stationary affine case}{Iteration of the stationary affine endomorphism and classification of its fixed points}% ADP102-0292
  \HMTTitleMapEntry{Route inversion}{Antipodal of the reversible path groupoid and conjugation of the affine holonomy}% ADP102-0293
  \HMTTitleMapEntry{Mirror of constants and fixed axis}{Involution \(\pi/e\) of the fiber of constants, invariant \(u^\pi+u^e\) and fixed axis \(\varphi\)}% ADP102-0294
  \HMTTitleMapEntry{Visible geometry as publication of the integral holonomic state}{Archimedean evaluation of the integral holonomic state and genealogical fidelity defect}% ADP102-0295
  \HMTTitleMapEntry{The inversion of the explanatory order}{Causal Factorization from the Enriched History to the Archimedean Evaluation}% ADP102-0296
  \HMTTitleMapEntry{State, history and publication}{Space of integral holonomic states, category of TPK histories, and Archimedean evaluation}% ADP102-0297
  \HMTTitleMapEntry{Intrinsic circle and projected circle}{Distinction between intrinsic radial orbit and circular projection of a helical trajectory}% ADP102-0298
  \HMTTitleMapEntry{The nonadic cell and the two-dimensional birth}{Construction of the bidimensional nonadic support \(X_9=(\mathbb Z/9\mathbb Z)^2\)}% ADP102-0299
  \HMTTitleMapEntry{APP antecedents of radial classification}{APP invariants that induce the accumulated radial character}% ADP102-0300
  \HMTTitleMapEntry{Lifted ternary state}{Residue and carry enriched TRIT state: multiplicity of threshold representatives}% ADP102-0301
  \HMTTitleMapEntry{Quadratic algebras}{Family \(\mathcal A_\tau=\mathbb R[J]/(J^2+\tau I)\) and elliptic, parabolic, and hyperbolic classification}% ADP102-0302
  \HMTTitleMapEntry{Double projection}{Joint factorization of the Archimedean and exceptional publications of the integral holonomic state}% ADP102-0303
  \HMTTitleMapEntry{Local representation of a transition}{Affine representation \(e\mapsto h_e\) of the generators of the route category}% ADP102-0304
  \HMTTitleMapEntry{Prefix holonomies}{Ordered affine products \texorpdfstring{\(H_j=h_{e_j}\cdots h_{e_1}\)}{H_j = h_ej ... h_e1} and preservation of prefixes}% ADP102-0305
  \HMTTitleMapEntry{Lifted category and elementary characters}{Category of enriched paths by 2 TRIT orientations and definition of the characters \(p\) and \(a\)}% ADP102-0306
  \HMTTitleMapEntry{Balanced decomposition}{Unique decomposition \(p=r_3(p)+3c_3(p)\) into balanced residue and carry}% ADP102-0307
  \HMTTitleMapEntry{Involution}{Action of \((\tau_+,\tau_-)\mapsto(-\tau_-,-\tau_+)\) over the characters \(p\) and \(a\)}% ADP102-0308
  \HMTTitleMapEntry{Definition and domain}{Definition of \(L_p\) on \(\mathbb R_{>0}\) and domain of the inverse \(\exp_p\)}% ADP102-0309
  \HMTTitleMapEntry{Deformed addition}{Composition law \(L_p(xy)=L_p(x)+L_p(y)+pL_p(x)L_p(y)\)}% ADP102-0310
  \HMTTitleMapEntry{Non-exhaustiveness and exact scope}{Classificatory domain of \(p\in\{2,1,0,-1,-2\}\) and general affine extension}% ADP102-0311
  \HMTTitleMapEntry{Obstruction to the reduced reader}{Defecto of naturality of the reader reduced \((p,a,\Lambda,C,m,r,\theta)\)}% ADP102-0312
  \HMTTitleMapEntry{Definition of the complete reader}{Enriched radial reader \(\mathcal P^{\mathrm{enr}}_{\mathrm{rad},N}\): domain, codomain, and preservation of prefixes}% ADP102-0313
  \HMTTitleMapEntry{Conservation as traceability}{Dynamic conservation of carry, boundary, and memory by the radial reader}% ADP102-0314
  \HMTTitleMapEntry{Quadruple subdivision}{Naturality of the reader under the quadruple refinement \(\operatorname{sd}_4\)}% ADP102-0315
  \HMTTitleMapEntry{Genealogical classification and non-faithfulness of publication}{Classification by the full radial signature and failure of fidelity of the Archimedean evaluation}% ADP102-0316
  \HMTTitleMapEntry{Classification by 3 visible increments}{Kinematic classification through the increments \((\omega,\beta,\mu)\)}% ADP102-0317
  \HMTTitleMapEntry{3 fibers of the same visible evaluation}{\(3\) genealogical classes with a same Archimedean image}% ADP102-0318
  \HMTTitleMapEntry{Screw law}{Spectral identity \(V^{-1}\mathscr S V=(2+\sqrt3)\mathrm i\mathscr S\) and equation of the logarithmic suspension}% ADP102-0319
  \HMTTitleMapEntry{Common tritic mode}{Common singular mode of the APP sheets in the canonical chart \(1,\ldots,9\)}% ADP102-0320
  \HMTTitleMapEntry{Ledger material}{Genealogical record of the material route: surviving extension, full signature, and dimensionless character}% ADP102-0321
  \HMTTitleMapEntry{Subsequent dimensional publication}{Realization of the constitutive vacuum operator in the fields \((E_\sigma,B_\sigma)\)}% ADP102-0322
  \HMTTitleMapEntry{Planar polar curve}{Curvature and arc length of a polar curve \(r=r(\theta)\)}% ADP102-0323
  \HMTTitleMapEntry{Specializations}{Evaluation of \(\kappa_{\mathrm F}\) in the logarithmic, Archimedean, Fermat, hyperbolic, and lituus families}% ADP102-0324
  \HMTTitleMapEntry{Spatial suspension}{Spatial immersion \(\Gamma(\theta)\) and Frenet--Lancret differential criteria}% ADP102-0325
  \HMTTitleMapEntry{Geometric trichromy}{Coordination of the regimes \(J^2=-I,0,+I\) with the Frenet geometry}% ADP102-0326
  \HMTTitleMapEntry{Conclusion}{Theorem of closure of genealogical spiral correspondences}% ADP102-0327
  \HMTTitleMapEntry{Affine semidirect product and route composition}{Affine monoid \(\operatorname{End}(U_{\mathrm{APP}})\ltimes U_{\mathrm{APP}}\) and ordered composition of routes}% ADP102-0328
  \HMTTitleMapEntry{Exhaustive table of the local chart of 2 orientations}{Exhaustive enumeration of the local chart \((\tau_+,\tau_-)\mapsto(p,a,r_3,c_3)\)}% ADP102-0329
  \HMTTitleMapEntry{5 regions of \texorpdfstring{\(\pi\)}{pi} and non-faithfulness of evaluation}{Decomposition of the microfiber of \(\pi\) into 5 regions and multiplicity of Archimedean evaluation}% ADP102-0330
  \HMTTitleMapEntry{Power--logarithmic algebra and 5 canonical publications}{Algebra of \(L_p\) and classification of the 5 canonical families \(p\in\{2,1,0,-1,-2\}\)}% ADP102-0331
  \HMTTitleMapEntry{Derivation of the 5 canonical curves}{Evaluation of \(\exp_p(u_0+\beta m)\) in the 5 canonical families}% ADP102-0332
  \HMTTitleMapEntry{Sheet inversion}{Conjugation \(p\leftrightarrow-p\) of the power--logarithmic sheets}% ADP102-0333
  \HMTTitleMapEntry{General affine sector}{Classification of the affine sector by the pair \((p,[\Lambda,C])\)}% ADP102-0334
  \HMTTitleMapEntry{Naturality, equivariance and non-faithfulness}{Naturality under truncation, nonadic equivariance, and Archimedean fidelity defect}% ADP102-0335
  \HMTTitleMapEntry{The prefix reader}{Affine prefix reader and enriched record conservation}% ADP102-0336
  \HMTTitleMapEntry{Non-trivial refinements}{Compatibility criterion for general refinements of paths}% ADP102-0337
  \HMTTitleMapEntry{Equivariance, not invariance}{Equivariance of the radial cocycle under the holonomy \(\Gamma_9\)}% ADP102-0338
  \HMTTitleMapEntry{3 proofs of non-faithfulness}{\(3\) explicit obstructions to the faithfulness of Archimedean evaluation}% ADP102-0339
  \HMTTitleMapEntry{Same endpoint, different affine order}{Final product coincidence and dependence on the order of affine factors}% ADP102-0340
  \HMTTitleMapEntry{Same curve, different genealogical region}{Coincidence of the Archimedean image between distinct genealogical regions}% ADP102-0341
  \HMTTitleMapEntry{Classification by mathematical level}{Domain stratification: history, reader, Archimedean evaluation, and physical realization}% ADP102-0342
  \HMTTitleMapEntry{Classification by rotation, radius and memory}{Classification by the increments \((\omega,\beta,\mu)\)}% ADP102-0343
  \HMTTitleMapEntry{Subsequent realizations and type checks}{Electromagnetic and Material Realizations with Domain and Codomain Separation}% ADP102-0344
  \HMTTitleMapEntry{Electronic centre and subsequent reader}{Compatibility of the radial reader with the electronic fixed point \((5,5)\) and atlas \(81/56/13/324\)}% ADP102-0345
  \HMTTitleMapEntry{Closed matrix of negative controls}{Matrix of validity conditions by causal level}% ADP102-0346
  \HMTTitleMapEntry{First prolongation and coinductive boundary}{Initial coinductive prolongation from \(w_6\) until the compatibility boundary of order \(1\)}% ADP102-0347
  \HMTTitleMapEntry{Path B: vacancy kernel \texorpdfstring{\(E_{21/22}\)}{E21/22}}{Characterization of \(\alpha\) as a simple root of the analytic kernel \(E_{21/22}\)}% ADP102-0348
  % Purismo residual: rótulos activos localizados en el índice del sucesor.
  % Las sustituciones siguientes no alteran el nivel jerárquico ni el cuerpo.
  \HMTTitleMapEntry{Observable chronology of the two cursors}{Chronological evolution of the cursor pair and determination of the observable state}% ADP102-0349
  \HMTTitleMapEntry{Finite two-cursor emitter}{Finite emitter associated with the cursor pair and generation of admissible states}% ADP102-0350
  \HMTTitleMapEntry{From one cell to two coordinates: birth of APP}{Construction of APP via decomposition of 1 cell into 2 coordinates}% ADP102-0351
  \HMTTitleMapEntry{The branch of \(20\) blocks}{Prolongation of the surviving section via 20 compatible blocks}% ADP102-0352
  \HMTTitleMapEntry{From the prefix to a real point}{Archimedean Evaluation of the Compatible Prefix and Determination of the Real Point}% ADP102-0353
  \HMTTitleMapEntry{A reproducible emission of twelve triads}{Reproducible emission of 12 triads and reconstruction of their integral holonomic state}% ADP102-0354
  \HMTTitleMapEntry{Construction of the closure character and structural genealogical trajectories of \(\pi\)}{Construction of the closure character and of the genealogical trajectories of \(\pi\)}% ADP102-0355
  \HMTTitleMapEntry{The slope 1/5 forced by the \(5\) regions}{Determination of the slope 1/5 via the balance of the 5 oriented regions}% ADP102-0356
  \HMTTitleMapEntry{Dodecaphasic window of the signed state and reversible closure of the cycle}{Finite restriction of the signed dodecaphase state and reversible closure of the cycle}% ADP102-0357
  \HMTTitleMapEntry{Finite dodecaphasic window}{Finite domain of the dodecaphonic closure}% ADP102-0358
  \HMTTitleMapEntry{Remote prolongation \texorpdfstring{\(\mathsf T_{\alpha,\infty}\)}{T alpha infinity} of the integer path}{Prolongation coinductive \(\mathsf T_{\alpha,\infty}\) of the entire dodecaphasic closure}% ADP102-0359
  \HMTTitleMapEntry{Theorem of the four integer channels}{Classification Theorem of the 4 Integer Channels}% ADP102-0360
  \HMTTitleMapEntry{First intrinsic encoder and its limit}{Initial intrinsic encoder and existence of its limit}% ADP102-0361
  \HMTTitleMapEntry{Typing hypotheses and negative controls}{Typing conditions and incompatibility criteria of the nonadic vacancy theory}% ADP102-0362
  \HMTTitleMapEntry{Chronological path of length, action, and spectrum}{Chronological composition of length, action, and spectrum in the determination of \(G\)}% ADP102-0363
  \HMTTitleMapEntry{Twisted operator and deterministic negative control}{Operator with torsion and deterministic obstruction to multiplicative continuation}% ADP102-0364
  \HMTTitleMapEntry{Mathematical stratification of the Majorana object into four levels}{Mathematical stratification of the Majorana object in 4 levels}% ADP102-0365
  \HMTTitleMapEntry{Six characters and four braid twists}{Classification of 6-character by 4 braid torsion classes}% ADP102-0366
  \HMTTitleMapEntry{specific obstructions and non-degeneracy conditions}{Specific Obstructions and Non-Degeneracy Conditions}% ADP102-0367
  \HMTTitleMapEntry{Up quark and up antiquark}{Quark sectors \(u\) and antiquark sectors \(\bar u\)}% ADP102-0368
  \HMTTitleMapEntry{Down quark and down antiquark}{Quark sectors \(d\) and antiquark sectors \(\bar d\)}% ADP102-0369
  \HMTTitleMapEntry{Refutation controls.}{Validity conditions and incompatibility of the equivariant realization of species}% ADP102-0370
  \HMTTitleMapEntry{The local square of \texorpdfstring{\(e\)}{e} and the memory break}{Local positional square of \(e\) and genealogical differentiation of states with the same visible block}% ADP102-0371
  \HMTTitleMapEntry{The general law before its species-specific realizations}{Causal precedence of the General Mass Law with respect to species-specific realizations}% ADP102-0372
  \HMTTitleMapEntry{Unified foundation of mass}{Extension–route factorization and genealogical record conservation}% ADP102-0373
  \HMTTitleMapEntry{The route band}{Definition of the route band through the orbit of a persistent section}% ADP102-0374
  \HMTTitleMapEntry{The incidence channel is local curvature, not an auxiliary factor}{Interpretation of the incidence channel as the local curvature of the gauge system}% ADP102-0375
  % Publicaciones terminales del capítulo 95: se conserva cada desarrollo
  % completo y se sustituye la enumeración narrativa por objeto, operación y resultado.
  \HMTTitleMapEntry{1st resolution: Riemann and positivity as a constructed complement}{Spectral–polar publication of the continuum via the positive complement of Guinand–Weil}% ADP102-0376
  \HMTTitleMapEntry{Global existence, regularity and uniqueness for Navier–Stokes through a uniform bound and compactness}{Solenoidal publication of the continuum: coercive bound, compactness and global regularity}% ADP102-0377
  \HMTTitleMapEntry{3rd solution: Yang–Mills and the curvature sector gap}{Publication-grade continuum: positivity by reflection and curvature sector gap}% ADP102-0378
  \HMTTitleMapEntry{4th resolution: Hodge and the effective construction of the algebraic preimage}{Cohomological publication of the continuum: compatible lifting and effective construction of the algebraic preimage}% ADP102-0379
  \HMTTitleMapEntry{5th resolution: BSD and the coincidence of the analytic and arithmetic publications}{Determinantal publication of the continuum: Kato–Poitou–Tate comparison and identity of the leading term of Birch–Swinnerton–Dyer}% ADP102-0380
  \HMTTitleMapEntry{The scattering identity that preserves cross power}{Scattering identity and cross-power conservation}% ADP102-0381
  \HMTTitleMapEntry{Persistence of the terminal object under passage to the limit}{Persistence of the terminal solenoidal state under passage to the limit}% ADP102-0382
  \HMTTitleMapEntry{Global existence, regularity, and uniqueness of the Navier–Stokes solution obtained by a uniform bound and compactness}{Global Solenoidal Closure Theorem via Uniform Bound and Compactness}% ADP102-0383
  \HMTTitleMapEntry{The universal extension operator}{Universal Operator for the Extension of Compatible Cohomology Class Extensions}% ADP102-0384
  \HMTTitleMapEntry{Factorization of the terminal RH--NS--YM--Hodge--BSD intertwiners through the common genealogical continuum}{Information losses specific to the 5 classical publications and preservation in the integral holonomic state}% ADP102-0385
  \HMTTitleMapEntry{Factorization of RH–NS–YM–Hodge–BSD as terminal publications of the common genealogical continuum}{Terminal factorization theorem of RH–NS–YM–Hodge–BSD to traves of the common genealogical continuum}% ADP102-0386
  \HMTTitleMapEntry{Correlative emergence of \texorpdfstring{\(5\)}{5} non-exhaustive intrinsic actions}{Expository decomposition of the correlative action into 5 intrinsic views of the same state}% ADP102-0387
  \HMTTitleMapEntry{One and the same boundary: interior flow and gauge tension}{Common factorization of the boundary operator in the solenoidal and gauge realizations}% ADP102-0388
  \HMTTitleMapEntry{The common boundary state}{Integral holonomic boundary state and parity decomposition of transport}% ADP102-0389
  \HMTTitleMapEntry{The minimal operator and its spectrum}{Minimal Boundary Complex and Spectral Equivalence of the 0 and 1 Degree Laplacians}% ADP102-0390
  \HMTTitleMapEntry{Hydrodynamic reading}{Solenoidal realization of the boundary operator: flow contraction and energy conservation}% ADP102-0391
  \HMTTitleMapEntry{Yang--Mills reading}{Boundary operator calibration: coercivity, gauge quotient, and spectral gap}% ADP102-0392
  % Promoción jerárquica de los certificados causales de \(\Gamma_9\) y de
  % la generación trítica de los canales enteros.
  \HMTTitleMapEntry{Monodromic filtration and calculation at prescribed depth}{Finite certification of Γ₉: 396 levels, 2.376 trits, 43 complete revolutions, 387 pairs \((K,K+9)\) per channel, and an ambient fiber of 729 elements}% ADP102-0393
  \HMTTitleMapEntry{Finite census of transitions and decimal generation}{Exhaustive census of 396 transitions, 2.376 trits, 43 turns and 387 equal-phase pairs}% ADP102-0394
  \HMTTitleMapEntry{Integer structure of the constant readers}{Arithmetic classification of integer channels and compatibility with the triticale semicrown}% ADP102-0395
}
\let\HMTNativeChapter\chapter
\let\HMTNativeSection\section
\let\HMTNativeSubsection\subsection
\let\HMTNativeSubsubsection\subsubsection
\let\HMTNativeParagraph\paragraph
\let\HMTNativeSubparagraph\subparagraph
\RenewDocumentCommand{\chapter}{s o m}{%
  \HMTResolvePublicTitle{#3}%
  \IfBooleanTF{#1}{\HMTNativeChapter*{\HMTResolvedTitle}}{%
    \IfNoValueTF{#2}{\HMTNativeChapter{\HMTResolvedTitle}}{\HMTNativeChapter[\HMTResolvedTitle]{\HMTResolvedTitle}}}}
\RenewDocumentCommand{\section}{s o m}{%
  \HMTResolvePublicTitle{#3}%
  \IfBooleanTF{#1}{\HMTNativeSection*{\HMTResolvedTitle}}{%
    \IfNoValueTF{#2}{\HMTNativeSection{\HMTResolvedTitle}}{\HMTNativeSection[\HMTResolvedTitle]{\HMTResolvedTitle}}}}
\RenewDocumentCommand{\subsection}{s o m}{%
  \HMTResolvePublicTitle{#3}%
  \IfBooleanTF{#1}{\HMTNativeSubsection*{\HMTResolvedTitle}}{%
    \IfNoValueTF{#2}{\HMTNativeSubsection{\HMTResolvedTitle}}{\HMTNativeSubsection[\HMTResolvedTitle]{\HMTResolvedTitle}}}}
\RenewDocumentCommand{\subsubsection}{s o m}{%
  \HMTResolvePublicTitle{#3}%
  \IfBooleanTF{#1}{\HMTNativeSubsubsection*{\HMTResolvedTitle}}{%
    \IfNoValueTF{#2}{\HMTNativeSubsubsection{\HMTResolvedTitle}}{\HMTNativeSubsubsection[\HMTResolvedTitle]{\HMTResolvedTitle}}}}
\RenewDocumentCommand{\paragraph}{s o m}{%
  \HMTResolvePublicTitle{#3}%
  \IfBooleanTF{#1}{\HMTNativeParagraph*{\HMTResolvedTitle}}{%
    \IfNoValueTF{#2}{\HMTNativeParagraph{\HMTResolvedTitle}}{\HMTNativeParagraph[\HMTResolvedTitle]{\HMTResolvedTitle}}}}
\RenewDocumentCommand{\subparagraph}{s o m}{%
  \HMTResolvePublicTitle{#3}%
  \IfBooleanTF{#1}{\HMTNativeSubparagraph*{\HMTResolvedTitle}}{%
    \IfNoValueTF{#2}{\HMTNativeSubparagraph{\HMTResolvedTitle}}{\HMTNativeSubparagraph[\HMTResolvedTitle]{\HMTResolvedTitle}}}}
% Las utilidades de degradación posteriores deben conservar el mapa de títulos.
\let\HMTBookChapter\chapter
\let\HMTBookSection\section
\let\HMTBookSubsection\subsection
\let\HMTBookSubsubsection\subsubsection
\let\HMTBookParagraph\paragraph
\let\HMTBookSubparagraph\subparagraph
\makeatother
 % Adaptador de compatibilidad del fuente teoremática íntegro del radión.
%
% El tratado rector ya carga tipografía, geometría, hipervínculos, teoremas,
% tablas, TikZ y tcolorbox.  Este archivo aporta únicamente los nombres y
% entornos exclusivos del fuente teoremática, sin importar su preámbulo autónomo ni
% alterar la jerarquía tipográfica general.

\makeatletter
\@ifundefined{color@RadNavy}{\definecolor{RadNavy}{HTML}{102A43}}{}
\@ifundefined{color@RadBlue}{\definecolor{RadBlue}{HTML}{1F5E7A}}{}
\@ifundefined{color@RadTeal}{\definecolor{RadTeal}{HTML}{168C8C}}{}
\@ifundefined{color@RadCopper}{\definecolor{RadCopper}{HTML}{B66A3C}}{}
\@ifundefined{color@RadGold}{\definecolor{RadGold}{HTML}{D4A73A}}{}
\@ifundefined{color@RadInk}{\definecolor{RadInk}{HTML}{1A202C}}{}
\@ifundefined{color@RadMist}{\definecolor{RadMist}{HTML}{EDF4F7}}{}
\@ifundefined{color@RadCream}{\definecolor{RadCream}{HTML}{FAF7F0}}{}
\@ifundefined{color@RadRose}{\definecolor{RadRose}{HTML}{8D3F55}}{}
\@ifundefined{color@RadGreen}{\definecolor{RadGreen}{HTML}{2F7D5D}}{}
\@ifundefined{color@RadGray}{\definecolor{RadGray}{HTML}{66788A}}{}

\providecommand{\TRIT}{\ensuremath{\mathrm{TRIT}}}
\providecommand{\HMT}{\ensuremath{\mathrm{HMT}}}
\providecommand{\Rstate}{\mathcal R_{12}}
\providecommand{\epsR}{\varepsilon_R}
\providecommand{\epsone}{\varepsilon_R^{(1)}}
\providecommand{\pih}{\pi}
\providecommand{\alphah}{\alpha}
\providecommand{\hret}{\hbar_{\mathrm{ret}}}
\providecommand{\hfull}{h_{\mathrm{ret}}}
\providecommand{\diff}{\mathop{}\!\mathrm d}
\providecommand{\e}{\mathrm e}
\providecommand{\R}{\mathbb R}
\providecommand{\C}{\mathbb C}
\providecommand{\F}{\mathbb F}
\providecommand{\Z}{\mathbb Z}
\providecommand{\dd}{\,\mathrm d}
\providecommand{\status}[1]{\textsf{\bfseries #1}}
\providecommand{\sourcepath}[1]{\nolinkurl{#1}}

\@ifundefined{typebox}{%
  \newtcolorbox{typebox}[1][]{enhanced,breakable,colback=white,
    colframe=RadBlue,boxrule=.7pt,arc=1.5mm,left=4mm,right=4mm,
    top=3mm,bottom=3mm,title={Type receipt},fonttitle=\bfseries,
    coltitle=RadNavy,#1}%
}{}
\@ifundefined{narrativebox}{%
  \newtcolorbox{narrativebox}[1][]{enhanced,breakable,colback=white,
    colframe=RadGold,borderline west={2.4pt}{0pt}{RadGold},
    boxrule=.25pt,arc=0mm,left=5mm,right=4mm,top=3mm,bottom=3mm,#1}%
}{}

% Las once portadillas internas del fuente teoremática autónomo se convierten en
% divisiones científicas ordinarias dentro del capítulo 52.  Se evita así
% introducir partes autónomas, gigantismo tipográfico o páginas vacías.
\providecommand{\partopener}[3]{\section{#2}\noindent\emph{#3}\par\medskip}

\providecommand{\BigRadionEquation}{%
  \begin{equation*}
  \boxed{\displaystyle
  \frac{180^{\circ}}{\pih}
  =50^{\circ}+A^{\circ}-\frac{20}{81}\,\Delta_4^{\circ}+\epsR^{\circ}}
  \end{equation*}}
\makeatother
 % Adaptador mínimo para las correspondencias espirales genealógicas.
% El preámbulo autónomo no se importa: sólo se registran símbolos y colores
% que los cuerpos científicos usan de manera efectiva.
\makeatletter
\@ifundefined{color@HMTNavy}{\definecolor{HMTNavy}{HTML}{17324D}}{}
\@ifundefined{color@HMTBlue}{\definecolor{HMTBlue}{HTML}{315E7D}}{}
\@ifundefined{color@HMTRed}{\definecolor{HMTRed}{HTML}{8A3D36}}{}
\@ifundefined{color@HMTGold}{\definecolor{HMTGold}{HTML}{9B7B3E}}{}
\@ifundefined{color@HMTGray}{\definecolor{HMTGray}{HTML}{5F6870}}{}
\@ifundefined{color@HMTLight}{\definecolor{HMTLight}{HTML}{F2F4F5}}{}
\@ifundefined{color@HMTLine}{\definecolor{HMTLine}{HTML}{C9CFD3}}{}
\providecommand{\enr}{\mathrm{enr}}
\providecommand{\AffEnd}{\operatorname{AffEnd}}
\providecommand{\Hist}{\operatorname{Hist}}
\providecommand{\Led}{\operatorname{Led}}
\providecommand{\Carr}{\operatorname{Carr}}
\providecommand{\ev}{\operatorname{ev}}
\providecommand{\sd}{\operatorname{sd}}
\providecommand{\id}{\operatorname{id}}
\providecommand{\im}{\operatorname{im}}
\providecommand{\Spec}{\operatorname{Spec}}
\providecommand{\arcosh}{\operatorname{arcosh}}
\providecommand{\N}{\mathbb N}
\providecommand{\T}{\mathbb T}
\providecommand{\ii}{\mathrm i}
\providecommand{\twoway}{\textit{bidirectional}}
\providecommand{\ph}{\mathrm{ph}}
\providecommand{\mem}{\mathrm{mem}}
\providecommand{\rad}{\mathrm{rad}}
\providecommand{\arq}{\mathrm{arq}}
\providecommand{\e}{\mathrm e}
\providecommand{\F}{\mathbb F}
\providecommand{\Z}{\mathbb Z}
\providecommand{\R}{\mathbb R}
\providecommand{\dd}{\,\mathrm d}

% Los fuentes teoremáticas de correspondencias espirales se distribuyen por
% residencia causal. Sus etiquetas conservan un espacio de nombres único y
% sus divisiones autónomas se subordinan a la jerarquía del capítulo receptor.
\newenvironment{HMTSpiralIntegratedBody}{%
  \begingroup
  % Las copias públicas ya llevan el prefijo material ``espiral:'' en cada
  % etiqueta y referencia. Esto incluye los labels diferidos de amsmath, que
  % no podían asegurarse mediante una redefinición local de \label.
  \let\part\HMTBookSection
  \let\chapter\HMTBookSubsection
  \let\section\HMTBookSubsubsection
  \let\subsection\HMTBookParagraph
  \let\subsubsection\HMTBookSubparagraph
}{\endgroup}

\newenvironment{HMTSpiralChapterBody}{%
  \begingroup
  \let\chapter\HMTBookSection
  \let\section\HMTBookSubsection
  \let\subsection\HMTBookSubsubsection
  \let\subsubsection\HMTBookParagraph
}{\endgroup}

\newenvironment{HMTSpiralAppendixBody}{%
  \begingroup
  \let\part\HMTBookChapter
  \let\chapter\HMTBookSection
  \let\section\HMTBookSubsection
  \let\subsection\HMTBookSubsubsection
  \let\subsubsection\HMTBookParagraph
}{\endgroup}
\makeatother
 
% Identificador estable para builds reproducibles.
\ifdefined\pdfinfoomitdate
  \pdfinfoomitdate=1
\fi
\ifdefined\pdfsuppressptexinfo
  \pdfsuppressptexinfo=15
\fi
\ifdefined\pdftrailerid
  \pdftrailerid{<554E49444144484F4C4F47524146494341>
                <44494D454E53494F4E414C484D54323026>}
\fi

\makeatletter\begin{document}
\frontmatter
\hypersetup{pageanchor=false}
\begin{titlepage}
\thispagestyle{empty}
\centering
\vspace*{1.8cm}
{\color{hmtgold}\rule{0.68\textwidth}{0.8pt}\par}
\vspace{1.35cm}
{\fontsize{31}{37}\selectfont\color{hmtblue}\bfseries
Double holographic projection\par
of an aperiodic time crystal\par}
\vspace{1.0cm}
{\Large\bfseries
Construction of the continuum as an inverse limit of compatible cylinders,\par
coinductive generation of \(\pi,\varphi,e,\alpha\) and exceptional\par
stabilizers of boundary incidence\par}
\vspace{1.1cm}
{\color{hmtgold}\rule{0.68\textwidth}{0.8pt}\par}
\vfill
{\large Oumar Haidara Fall\par}
{\large Rubén Ramos Balsa\par}
\vspace*{1.8cm}
\end{titlepage}
\hypersetup{pageanchor=true}

\chapter*{Abstract}
\addcontentsline{toc}{chapter}{Abstract}

An arithmetic of oriented histories is constructed in which the unit is
preserved under refinement while resolution, holonomic height,
and memory increase. The finite support of all possible paths
simultaneously maintains the additive and multiplicative evaluations, together with
their residues, quotients, and carries; a tritic fiber distinguishes the
elliptic, parabolic, and hyperbolic regimes; and a topological phase kernel selects,
transports, and updates each state without confusing phase return with
state return. The composition of these three structures determines a
noncommutative associative algebra, graded by memory, and a system of
compatible refinements whose history space is an inverse limit.
Cylindrical observables extend covariantly and constitute
the dual direct limit. This double variance provides the operator
definition of holography used in the article.

The dynamics combines a finite nonadic phase, a mechanical coding of
slope $\delta=\log_{10}(10/9)$, and a carry cocycle. The result is
an aperiodic topological-temporal order of linear complexity and zero
topological entropy, endowed with a holonomy that raises memory at each nonadic
return. Its Weyl relations rigorously separate local periodicity
from global aperiodicity. On this support, coinductive
prolongation correlatively generates $\pi$, $\varphi$, $e$, and $\alpha$: the
first three coordinates realize, respectively, elliptic closure,
hyperbolic autoscaling, and exponential transport; the fine-structure
constant is the compatibility invariant that brings together, through two
internal derivations, dodecaphasic closure and aperiodic prolongation.

The same structure produces a unit multiplicative-carry vacancy
in an invariant additive fiber. Its local algebra is
$\mathrm{Cl}_{1,1}\cong M_2(\R)$ and simultaneously contains elliptic
and hyperbolic generators and null Witt directions. This section determines
the electronic kernel and extends the genealogy toward action, the Planck
and gravity scales, the Barbero--Immirzi parameter, the area spectrum, and
boundary information. The second projection preserves incidence and
leads, through the Paley and Witt configurations, the Golay code, the
Leech lattice, and the vertex algebra $V^\natural$, to the exceptional
stabilizers and the Monster group. Thus, the Archimedean realization and the
exceptional realization do not constitute independent sources: they are two
representations of a unique discrete structure of the continuum.

The uniform dodecaphasic reader also determines a sequence of first roots of critical return. Its separation ratio is proved to converge to the Feigenbaum constant, preserving the original selection, orientation and memory of the construction.

\vspace{4mm}
\noindent\textbf{Keywords:} triadic modular holography; inverse
limit; nonadic holonomy; Sturmian word; aperiodic time crystal;
path algebra; fine-structure constant; split Clifford;
Barbero--Immirzi; Leech lattice; Monster group.

\clearpage

\chapter*{Abstract}
\addcontentsline{toc}{chapter}{Abstract}

We construct an arithmetic of oriented histories in which the unit is
preserved under refinement while resolution, holonomic height and memory
increase.  The finite support of all possible paths retains the additive and
multiplicative evaluations simultaneously, together with their remainders,
quotients and carries; a tritic fibre distinguishes the elliptic, parabolic
and hyperbolic regimes; and a topological phase kernel selects, transports
and updates each state without identifying phase return with state return.
Their composition determines a memory-graded, associative and
noncommutative algebra and a compatible refinement system whose history
space is an inverse limit.  Cylindrical observables extend covariantly and
form the dual direct limit.  This variance duality is the operative notion of
holography used throughout the article.

The dynamics combines a finite nonadic phase, a mechanical coding of slope
$\delta=\log_{10}(10/9)$ and a carry cocycle.  It yields an aperiodic
topological--temporal order with linear complexity and zero topological
entropy, endowed with a holonomy that raises memory at every nonadic return.
On this support the coinductive prolongation generates $\pi$, $\varphi$,
$e$ and $\alpha$ as correlated outputs.  The same structure produces a
unit multiplicative-carry vacancy on an invariant additive fibre, whose
local algebra is $\mathrm{Cl}_{1,1}\cong M_2(\R)$.  The Archimedean
realization and the exceptional boundary-incidence realization are thereby
shown to be two representations of one discrete structure of the continuum.

The uniform dodecaphasic reader also determines a sequence of first roots of critical return. Its separation ratio is proved to converge to the Feigenbaum constant, preserving the original selection, orientation and memory of the construction.

\vspace{4mm}
\noindent\textbf{Keywords:} triadic modular holography; inverse limit;
nonadic holonomy; Sturmian word; aperiodic time crystal; path algebra;
fine-structure constant; split Clifford algebra; Barbero--Immirzi parameter;
Leech lattice; Monster group.
\clearpage
\tableofcontents
\clearpage

\let\cleardoublepage\clearpage
\mainmatter
\part{Two-sheeted arithmetic, trigeometric orientation, and phase dynamics}

\chapter{The unit as a compatible section of an arithmetic of histories}
\begingroup
\renewcommand{\chapter}[2][]{}%
\chapter{Thesis and mathematical architecture}
\label{sec:tesis-inaugural-helicoidal-hmt-md}

The uniform dodecaphasic reader also determines a sequence of first roots of critical return. Its separation ratio is proved to converge to the Feigenbaum constant, preserving the original selection, orientation and memory of the construction.


\section{The unit as a refinement invariant}

The central question is not to obtain a collection of scalar
values from a finite table. It is to determine the minimal
structure that preserves a unit through unlimited refinements,
distinguishes three classes of curvature, transports the memory of its paths, and
generates, through typed representations, numbers, constants, sections,
fields, and symmetries. The real line is not presupposed as the universal domain of
this construction. It appears at the end of a compatible contraction of
cylinders; it preserves the scalar value of the history, but not necessarily
the entirety of its genealogy.

For each depth $n\geq0$, let $X_n$ be the finite set of admissible
states. A state is not an isolated digit, but a structured fiber

\begin{equation}
x=\bigl(a,b;\Sigma,\Pi;R^+,q^+,R^\times,q^\times;
\kappa;\tau,o;\theta_9,\theta_{12};h;\partial;\mu;\rho;s\bigr),
\label{eq:estado-integral}
\end{equation}

where $(a,b)$ locates the cell; $\Sigma$ and $\Pi$ are the additive and
multiplicative sheets; $(R^+,q^+)$ and $(R^\times,q^\times)$ preserve residue and
quotient; $\kappa$ is the carry; $\tau\in\{+1,0,-1\}$ types the curvature;
$o$ fixes the orientation; $\theta_9$ and $\theta_{12}$ are the nonadic and
dodecaphasic phases; $h$ measures the helical height; $\partial$ records boundary
incidence; $\mu$ is memory; $\rho$ preserves the path and $s$ the survival
condition. These coordinates do not form a free product: they satisfy
the relations of path arithmetic, trigeometric typing, and the
transport dynamics to be constructed in the following chapters.

The maps

\begin{equation}
r_{m,n}:X_m\longrightarrow X_n,\qquad m>n,
\label{eq:restriccion}
\end{equation}

forget the last layer of resolution, but preserve every relation visible at
depth $n$. The space of infinitely refinable histories is

\begin{equation}
X_\infty=\varprojlim_n X_n.
\label{eq:limite-inverso}
\end{equation}

If $e_x$ denotes the characteristic idempotent of the cylinder determined by
$x\in X_n$, the prolongation of observables satisfies

\begin{equation}
r_{m,n}^{*}e_x=\sum_{y:\,r_{m,n}(y)=x}e_y.
\label{eq:elevacion-idempotente}
\end{equation}

Consequently,

\begin{equation}
r_{m,n}^{*}\!\left(\sum_{x\in X_n}e_x\right)
=\sum_{y\in X_m}e_y,
\label{eq:conservacion-unidad}
\end{equation}

so that the unit is preserved even as the number of extensions grows. Evolution
does not preserve a stationary configuration: it preserves the augmentation, the
relations, and the possibility of reconstructing each antecedent from the complete
history.

\section{Double variance and double projection}

The depth of states and the extension of observables advance in
opposite categorical directions. The restrictions of
\eqref{eq:restriccion} correspond to prolongation morphisms

\begin{equation}
r_{m,n}^{*}:\Halg_n\longrightarrow\Halg_m,
\qquad f\longmapsto f\circ r_{m,n},
\end{equation}

and the algebra of cylindrical observables is obtained as a direct limit:

\begin{equation}
\Halg_{\mathrm{cyl}}=\varinjlim_n(\Halg_n,r_{n+1,n}^{*}).
\label{eq:limite-directo}
\end{equation}

The pair formed by \eqref{eq:limite-inverso} and
\eqref{eq:limite-directo} provides the fundamental holographic structure:
the state gains resolution through contravariant compatibility and the language
of observation gains scope through covariant prolongation. This is not a
figurative similarity between scales, but the naturality of the
evaluation

\begin{equation}
\langle r_{m,n}^{*}a,y\rangle_m
=\langle a,r_{m,n}y\rangle_n.
\label{eq:naturalidad-evaluacion}
\end{equation}

The double projection in the title designates two representations of this same
pair. The first contracts compatible cylinders toward an Archimedean
realization; the second retains boundary incidence and extends it
toward exceptional structures:

\begin{equation}
\rho_{\mathrm{arq}}:\Halg_{\mathrm{cyl}}\longrightarrow\cA_{\R},
\qquad
\rho_{\mathrm{exc}}:\Halg_{\mathrm{cyl}}\longrightarrow\cA_{\mathrm{exc}}.
\label{eq:doble-proyeccion}
\end{equation}

Both have the same source and preserve different invariants. The first
retains the scalar coordinate of a compatible section; the second retains
the incidence patterns that determine its stabilizers. The word
\emph{double} does not duplicate holography: it distinguishes two non-equivalent
codomains of a single refinement and reading mechanism.

\section{Central statement}

\begin{theorem}[Conservative evolution and production of symmetries]
There exists a graded algebra of oriented histories
\begin{equation}
\Halg=\left[
\varinjlim_n
\frac{\R\langle e_x,J_x,T_u\rangle}
{\cR_{\mathrm{APP\text{-}TRIT\text{-}TPK}}}
\right]\rtimes_{\Gamma_9}\N,
\label{eq:algebra-rectora}
\end{equation}
that acts on $X_\infty$ and satisfies the following properties:
\begin{enumerate}[label=\textup{(\roman*)},leftmargin=9mm]
\item the unit and augmentation are preserved under every compatible refinement;
\item the nonadic return closes the phase and raises memory;
\item the elliptic, parabolic, and hyperbolic regimes belong to a single
trigeometric functional calculus;
\item the symbolic factor has linear complexity, zero topological entropy,
and long-range aperiodic order;
\item the discrete structure of the continuum arises jointly from the space of
histories and the algebra of observables;
\item $\pi$, $\varphi$, $e$, and $\alpha$ are correlated outputs of the
same coinductive prolongation;
\item a symmetry of a realization is the stabilizer of the relations and
boundary incidence that this realization preserves.
\end{enumerate}
\end{theorem}

The proof will be distributed throughout the article. No known
constant will be used to select its own genealogy, nor a symmetry
group to retrospectively construct the object on which it acts. The
logical order will always be the same: two-sheeted arithmetic, trigeometric typing,
holonomic transport, compatible refinement, continuum, constants,
physical sections, and boundary stabilizers.

\section{Irreducibility of the generating kernel}

The structure has three functionally irreducible components. If the
two arithmetic sheets are identified, the potential
$\Delta_{\Sigma\Pi}$ and the commutator of their projectors disappear. If the
three tritic regimes are identified, the spectrum disappears

\begin{equation}
J_+^2=-I,\qquad J_0^2=0,\qquad J_-^2=I.
\end{equation}

If the transport cocycle is trivialized, the increment

\begin{equation}
[D_M,\Gamma_9]=\Gamma_9
\label{eq:conmutador-memoria-intro}
\end{equation}

vanishes and the phase turn becomes indistinguishable from state return.
Each forgetful quotient destroys an invariant of a different type that the other
two factors do not reconstruct. This irreducibility explains why the
structure sought cannot be reduced to an arithmetic table, an aperiodic
word, a Clifford algebra, or an isolated matrix.

The conceptual consequence is precise. A number is a scalar coordinate
of a representation; a constant is a stable return or autoscale
invariant; a particle is a persistent section; a field transports
sections; a symmetry stabilizes the relations of its boundary. These five
classes come from the same algebra, but do not share a codomain and cannot
be confused with one another.
 \endgroup

\chapter{Arithmetic support of all possible paths}
\label{chap:83-01}

% RESULTADO_RECUPERADO. Owner integro: sections/hmt/03_app_ortograma.tex
% Destino causal: capitulo 1; la jerarquia se subordina sin alterar el owner.
% Los dos bloques científicos del owner se activan en una única residencia.
\begingroup
\def\HMTAPPFormal{1}
\def\HMTAPPDevelopments{1}
\begin{HMTScientificOwnerSubordinated}
\ifdefined\HMTAPPFormal
\chapter{APP: support, operations, and paths}
\label{chap:app}
\index{APP}

\noindent The acronym APP derives from \emph{All Possible Paths}. It designates a single path structure on a finite arithmetic support: the underlying graph has eighty-one vertices, whereas the set of finite paths, taken over all lengths, is countable. Its principal arithmetic matrix is the nonadic multiplication table. Addition is the internal law whose accumulation generates each row and each column of that table, and it also acts as a comparison observable on the same support. APP thus brings together addition, multiplication, and path composition before the ternary projection or the Archimedean evaluation is performed. The following definition exhaustively fixes the universe of paths under consideration.

\section{Domain, generators and preformal relations of the APP support}

The support of APP is a finite arithmetic torus: each of its two
coordinates ranges over the nine residue classes and returns to its origin after
nine unit translations. Its eighty-one points are therefore
pairs of classes modulo nine. The associated Cayley graph preserves this
cyclic structure: a vertex represents a position, an oriented edge
represents a unit cardinal displacement, and a finite sequence of
edges represents a path. Thus, \emph{all possible paths} denotes the set of finite cardinal words with an
initial point on the torus.

Each vertex simultaneously receives two arithmetic evaluations. The \emph{additive sheet} publishes the sum of its two coordinates, and the \emph{multiplicative sheet} publishes their product. The term \emph{sheet} here designates an evaluation of the same point and the same path; APP preserves a single support and two functions upon it.

The operative difference between the two sheets is exact. In a row or column,
addition acts by translation and repeatedly accumulates the same increment.
Multiplication by a unit also permutes the classes; multiplication
by a zero divisor folds several classes onto one and the same
residue. To preserve the information that this folding conceals, APP lifts
each evaluation to its residue--quotient coordinates:
\[
 i+j=R^+_{ij}+9q^+_{ij},
 \qquad
 ij=R^\times_{ij}+9q^\times_{ij}.
\]
The pairs \((R^+,q^+)\) and \((R^\times,q^\times)\) recover the integer values
prior to reduction. When blocks are concatenated, additive carry
\(\kappa_9\) records the discrepancy between adding integer representatives and
first adding their residues; together with both quotients, it makes recovery
compositional. Thus, additive accumulation constructs visible trajectories,
whereas the residue--quotient lift unfolds the fibers created by
multiplicative folding.

For a path \(\gamma=(x_0,\ldots,x_r)\), the genealogy visible is the
sucesion of positions and pairs of evaluations
\[
 \mathcal G_{\rm vis}(\gamma)
 =\bigl(\gamma;(\Sigma(x_k),\Pi(x_k))_{k=0}^{r}\bigr).
\]
At each step, its enriched genealogy adjoins the quotients, carry, and
transported orientation. The enriched--visible projection preserves the
path and residues and forgets precisely those memory coordinates. This
distinction will be the input to the complete state of TRIT and of the Topological Phase Kernel.
The complete domain of (104\,976) double-cursor initial conditions and
the (468) emission fibers they produce are published in full in the
\cref{app:catalogo-semillas-app-rev9}; there each seed is distinguished from the
emission to which it belongs, and the multiplicity of the fiber is preserved.

\section{Aritmetic support, translations and observables}

Let
\[
\dr(n)=1+((n-1)\bmod9)\in\{1,\ldots,9\}.
\]
This map chooses the representative \(9\) for the null class. We define
\(X_9=(\ZZ/9\ZZ)^2\), on which acts the group of translations
\((\ZZ/9\ZZ)^2\), and fix the symmetric set of generators
\[
 \mathscr D=\{(1,0),(0,1),(-1,0),(0,-1)\}.
\]
The Cayley graph
\[
 \GraphAPP=\operatorname{Cay}(X_9,\mathscr D)
\]
has \(81\) vertices and four oriented edges originating at each vertex.
An oriented path of length \(r\) is a sequence
\(\gamma=(x_0,\ldots,x_r)\) in \(X_9\) such that
\(x_{k+1}-x_k\in\mathscr D\) for every \(k<r\). We denote by
\(\operatorname{Path}(\GraphAPP)\) the set of all finite paths
so defined.

\begin{definition}[APP Structure]
The arithmetic structure of possible paths is the tuple
\[
 \AAPP=(X_9,\GraphAPP,\Sigma,\Pi,
        \operatorname{Path}(\GraphAPP)),
\]
where the two vertex observables are
\[
\Sigma(i,j)=\dr(i+j),
\qquad
\Pi(i,j)=\dr(ij),
\qquad1\le i,j\le9.
\]
\end{definition}

\subsection*{Marked cardinal and path words}

To make the orientation of the torus explicit, we use coordinates
\((\text{row},\text{column})\), with the row increasing southward and the
column increasing eastward, and write
\[
 \delta(N)=(-1,0),\qquad \delta(E)=(0,1),\qquad
 \delta(S)=(1,0),\qquad \delta(O)=(0,-1).
\]
Let \(\mathscr D_{\rm card}=\{E,N,O,S\}\), where \(O\) denotes west. Given an initial point
\(x_0\in X_9\) and a word \(w=d_1\cdots d_r\in
\mathscr D_{\rm card}^{r}\), its trajectory is
\[
 \gamma(x_0;w)_k
 =x_0+\sum_{j=1}^{k}\delta(d_j)\pmod 9,
 \qquad 0\leq k\leq r.
\]
The correspondence \((x_0,w)\mapsto\gamma(x_0;w)\) is bijective between
\(X_9\times\mathscr D_{\rm card}^{r}\) and the paths oriented of
length \(r\). In particular, existen exactly \(81\cdot4^r\) paths
of that length when it allowed returns and backtracks.

The word \(w\) closes in the torus precisely when
\[
\#S(w)-\#N(w)\equiv0\pmod9,
\qquad
\#E(w)-\#O(w)\equiv0\pmod9.
\]
For a closed word, the integer displacement computed before reduction
modulo nine defines the winding vector
\[
 \operatorname{wind}(w)=\frac1{9}
 \bigl(\#S(w)-\#N(w),\#E(w)-\#O(w)\bigr)\in\ZZ^2.
\]
Reversal of orientation reverses the order of the word, exchanges
\(E\leftrightarrow O\) and \(N\leftrightarrow S\), and transforms \(\operatorname{wind}(w)\) into its
opposite. This cardinal framework is subsequently transported in the phase
kernel and forms part of the definition of each trajectory.

The expressions are written here in the chart of representatives
\(\{1,\ldots,9\}\). The formula
\(\widetilde\Sigma(i,j)=\dr(i+j-1)\) is a translated chart. It represents the
same torus only when the translation is also applied to indices, initial
conditions, positional trajectories, and labels. All the following proofs employ
\(\Sigma(i,j)=\dr(i+j)\) and do not alternate both charts.

\begin{figure}[H]
\centering
\resizebox{\textwidth}{!}{\begin{tikzpicture}[font=\scriptsize,cell/.style={minimum width=4.2mm,minimum height=4.2mm,inner sep=0pt,draw=hmtgray!35}]
  \begin{scope}[xshift=0cm]
    \node[font=\bfseries\scriptsize,hmtblue,align=center] at (1.72,.92)
      {APP arithmetic matrix\\[-.5mm]\(\Pi(i,j)=\dr(ij)\)};
    \foreach \i in {1,...,9}{
      \node[font=\bfseries] at (-.35,{-(\i-1)*.43}) {\i};
      \foreach \j in {1,...,9}{
        \pgfmathtruncatemacro{\v}{mod(\i*\j-1,9)+1}
        \pgfmathtruncatemacro{\border}{(\i==9)||(\j==9)}
        \pgfmathtruncatemacro{\rad}{(mod(\i,3)==0)||(mod(\j,3)==0)}
        \ifnum\border=1
          \node[cell,fill=hmtlightblue,text=hmtblue,font=\bfseries] at ({(\j-1)*.43},{-(\i-1)*.43}) {\v};
        \else
          \ifnum\rad=1
            \node[cell,fill=hmtlightviolet] at ({(\j-1)*.43},{-(\i-1)*.43}) {\v};
          \else
            \node[cell,fill=white] at ({(\j-1)*.43},{-(\i-1)*.43}) {\v};
          \fi
        \fi
      }
    }
    \foreach \j in {1,...,9}{\node[font=\bfseries] at ({(\j-1)*.43},.38) {\j};}
    \draw[hmtgold,very thick,rounded corners=1pt] (1.075,-1.075) rectangle (1.935,-1.935);
    \draw[hmtblue,very thick] (3.72,.22)--(3.72,-3.72)--(-.22,-3.72);
    \node[hmtblue,anchor=north] at (1.75,-4.02) {subconjunto coordenado de la clase nula};
  \end{scope}

  \begin{scope}[xshift=5.25cm]
    \node[draw=hmtblue,rounded corners=2pt,fill=hmtlightblue,
      text width=6.25cm,align=left,inner sep=8pt] at (2.9,-1.35) {%
      \textbf{Ley interna de acumulación}\par\smallskip
      \[
        \Pi(i,j)=\dr(\underbrace{i+\cdots+i}_{j\ {\rm sumandos}})
        =\dr(ij).
      \]
      Addition generates the multiplicative rows and, in the transverse
      direction, the columns. The observable
      \(\Sigma(i,j)=\dr(i+j)\) preserves that elementary law in order to compare
      both readings at the same point.};
    \node[draw=hmtgold,rounded corners=2pt,fill=hmtlightgold,
      text width=6.25cm,align=center,inner sep=8pt] at (2.9,-3.65) {%
      \(\displaystyle
      \AAPP=(X_9,\GraphAPP,\Pi;\Sigma,
      \operatorname{Path}(\GraphAPP))\)\par
      a single APP, one principal matrix, one internal additive law, and four
      oriented directions \(N,E,S,O\).};
  \end{scope}
\end{tikzpicture}
 }
\caption{The unique APP matrix and its internal law. The nonadic multiplication table occupies the main body; each of its rows is generated by additive accumulation. The violet shading indicates the nonunitary rows or columns; the blue line identifies the subset determined by the representative \(9\) of the null class; the gold frame highlights the central block \(2\times2\).}
\label{fig:app-ortograma}
\end{figure}

Every point of \(X_9\) simultaneously receives the product and sum values,
but it does not thereby belong to two distinct APP matrices. The paths of
\(\operatorname{Path}(\GraphAPP)\) traverse a single support; \(\Pi\)
publishes its main table and \(\Sigma\) conserves the internal operation that
generates the accumulations and allows its incompatibility to be measured.

\begin{proposition}[Structure of Rows and Columns]
\label{prop:app-latina-unidades}
The matrix of \(\Sigma\) is a Latin square of order nine. The matrix of
\(\Pi\) is not; nevertheless, for each
\(u\in(\ZZ/9\ZZ)^\times=\{1,2,4,5,7,8\}\), the row and column
corresponding to \(u\) are permutations of the nine residue classes.
\end{proposition}

\begin{proof}
With \(i\) fixed, the map \(j\mapsto i+j\) is a translation of \(\ZZ/9\ZZ\); the same holds when \(j\) is fixed. Hence each representative occurs exactly once in every row and column of \(\Sigma\). For \(\Pi\), multiplication by \(u\) permutes the classes if and only if \(u\) is a unit. The rows corresponding to \(3,6,9\) contain repetitions, so the Latin property does not extend to the entire multiplicative matrix.
\end{proof}

\section{Multiplicative structure of the local ring}

Each row of the multiplication table is a repeated sum. Row \(i\)
accumulates \(i\) along one direction; column \(j\) accumulates \(j\) along
the perpendicular direction. Cell \((i,j)\) is the intersection
of both accumulations and returns \(ij=ji\). The table thus represents product
and commutativity through two transverse accumulations on the same
lattice.

The second row is
\[
2,4,6,8,1,3,5,7,9.
\]
Before the modular return the even ones appear and after the odd ones. The third
row is
\[
3,6,9,3,6,9,3,6,9,
\]
and makes the radical visible. Let
\[
 R:=\ZZ/9\ZZ,
 \qquad
 \mathfrak m:=(3)=\{\overline0,\overline3,\overline6\}.
\]
Multiplication by three defines
\[
 M_3:R\longrightarrow R,
 \qquad x\longmapsto3x,
 \qquad
 \ker M_3=\operatorname{im}M_3=\mathfrak m.
\]
By the first isomorphism theorem induces a linear identification
\begin{equation}
 \overline M_3:R/\mathfrak m\xrightarrow{\ \sim\ }\mathfrak m,
 \qquad \overline x\longmapsto3x,
 \label{eq:multiplicador-tres-radical-rev3}
\end{equation}
of vector spaces on \(\mathbb F_3\). No is an isomorphism of
rings: \(\mathfrak m^2=0\), whereas the quotient is a field.

In positive representatives, the three phases of
\eqref{eq:multiplicador-tres-radical-rev3} are
\begin{equation}
 369\longmapsto693\longmapsto936\longmapsto369,
 \label{eq:orbita-radical-369-rev3}
\end{equation}
produced by the phase shift \(j\mapsto j+1\). The row
corresponding to \(6\equiv-3\pmod9\) reverses the orientation of the same
cycle. This radical orbit must not be conflated with the integer
\(369\,693\,100\), which will subsequently appear as the digit count of
\(9^{9^9}\).

In direct reduction modulo three,
\[
3,6,9\longmapsto0,0,0.
\]
If we first extract the factor three and retain its internal coordinate,
\[
\eta_{\mathfrak m}(3a)=a\bmod3,
\qquad
3,6,9\longmapsto1,2,0.
\]
The second map preserves the internal coordinate of the ideal generated by
three. It must not be confused with the oriented ternary projection to be defined
in the section devoted to the TRIT\@.

\section{Coordinate subset of the null class}

In the graph of the multiplicative observable,
\[
\Pi(9,j)=\Pi(i,9)=9.
\]
The last row and the last column form a union of seventeen positions,
all with value nine. They do not constitute an intrinsic boundary of the arithmetic
torus: they are the subset of the chart in which at least one coordinate
represents the null class by means of \(9\). There are also four interior null representatives in
\[
(3,3),(3,6),(6,3),(6,6).
\]
The figure thus distinguishes this subset from the internal zeros of the
nonunital ideal. The additional symbol used later in the cubic
completion belongs to an enlarged space and is not identified with any of
them.

\section{Residue--quotient extension}

Reduction modulo nine preserves the residue and eliminates the quotient of the
Euclidean division. To study both components, we define, at each point,
\[
i+j=9q^+_{ij}+R^+_{ij},
\qquad
ij=9q^\times_{ij}+R^\times_{ij},
\]
where \(R^+=\Sigma\), \(R^\times=\Pi\) and the \(q\) preserve the quotient.
The residue--quotient extension of APP it is denoted by
\[
\AAPPlift=(R^+,q^+;R^\times,q^\times).
\]

Let \(s_9:\ZZ/9\ZZ\to\{1,\ldots,9\}\) be the section of representatives fixed
by \(\dr\). The additive carry associated with this section is
\[
 \kappa_9(a,b)
 =\frac{s_9(a)+s_9(b)-s_9(a+b)}9\in\ZZ.
\]

\begin{lemma}[Cocycle Identity of the Carry]
\label{lem:app-cociclo-acarreo}
For arbitrary \(a,b,c\in\ZZ/9\ZZ\),
\[
 \kappa_9(a,b)+\kappa_9(a+b,c)
 =\kappa_9(b,c)+\kappa_9(a,b+c).
\]
\end{lemma}

\begin{proof}
Both sides are equal to
\[
 \frac{s_9(a)+s_9(b)+s_9(c)-s_9(a+b+c)}9.\qedhere
\]
\end{proof}

This cocycle records the discrepancy between summing integer representatives and
first summing in the quotient. It is the compositional form of the
quotient coordinate that will reappear in base \(1000\) concatenation.

The gross totals are
\[
\sum_{i,j}(i+j)
=9\sum_i i+9\sum_j j
=2\cdot9\cdot45
=810,
\]
and
\[
\sum_{i,j}ij
=\left(\sum_i i\right)\left(\sum_j j\right)
=45^2
=2025.
\]
The visible sums are obtained from the tables:
\[
\sum R^+=405,
\qquad
\sum R^\times=459.
\]
Therefore, the total sums of the quotient coordinates are
\[
\sum q^+=\frac{810-405}{9}=45,
\qquad
\sum q^\times=\frac{2025-459}{9}=174.
\]

\begin{exactbox}[title={Decomposition exacta of the difference integrada}]
\[
\boxed{
2025-810
=1215
=(459-405)+9(174-45)
=54+9\cdot129.}
\]
The difference between the sums of residues is \(54\), whereas the
difference between the sums of quotients is \(129\). The equality recomposes the
total difference between the integer observables and proves that the quotient
component contains arithmetic information not determined by the residual
class.
\end{exactbox}

\fi

\ifdefined\HMTAPPDevelopments
\chapter{Arithmetic Operators and Local Geometry of the Path Space}
\label{chap:app-operadores}

\section{Localization of the excess in the nilpotent radical}

The group of units of \(\ZZ/9\ZZ\) is
\[
(\ZZ/9\ZZ)^\times=\{1,2,4,5,7,8\},
\]
and the maximal ideal is
\[
\mathfrak m=(3)=\{0,3,6\}=\{9,3,6\}_{\rm visual},
\qquad\mathfrak m^2=0.
\]
Each unit row of \(\Pi\) is a permutation of \(1,\ldots,9\) and sums to
\(45\). Rows \(3\) and \(6\) sum to \(54\); row \(9\) sums to \(81\).
Hence
\[
459=6\cdot45+2\cdot54+81
\]
and
\[
459-405=2(54-45)+(81-45)=54.
\]
All the excess is located in the nonunitary sector.

The union of rows and columns \(3,6,9\) contains \(45\) cells and sums to
\(297\). Its central square \(3\times3\) sums to \(81=3\cdot27\); the outer
arms, to \(216=3\cdot72\). Thus,
\[
297=216+81=3(72+27)=3\cdot99.
\]
The interpretation of \(72\) and \(27\) is obtained later from the
uniquely selected central block.

\section{Conditional expectation with respect to the projection \(\mathbb Z/9\mathbb Z\to\mathbb Z/3\mathbb Z\)}

The canonical reduction
\(\ZZ/9\ZZ\to\ZZ/3\ZZ\) determines the partition
\[
 C_1=\{1,4,7\},\qquad
 C_2=\{2,5,8\},\qquad
 C_0=\{3,6,9\}.
\]
The uniform conditional expectation averages each observable on the nine points of each block
\(C_a\times C_b\). The aggregate matrices are
\[
 M_{\Pi}=
 \begin{pmatrix}4&5&6\\5&4&6\\6&6&9\end{pmatrix},
 \qquad
 M_{\Sigma}=
 \begin{pmatrix}5&6&4\\6&4&5\\4&5&6\end{pmatrix}.
\]
Since each block contains nine positions,
\[
 9\sum_{a,b}(M_\Pi)_{ab}=459,\qquad
 9\sum_{a,b}(M_\Sigma)_{ab}=405.
\]
Thus, the averaging preserves the totals of both observables and relocates their difference:
\[
 \sum M_\Pi-\sum M_\Sigma=51-45=6,\qquad 9\cdot6=54.
\]

\begin{theorem}[Spectrum of the aggregated multiplicative operator]
\label{thm:app-espectro-agregado}
The matrix operator \(M_\Pi\) satisfies
\[
 \det M_\Pi=-9
\]
and its real spectrum is
\[
 \operatorname{Spec}(M_\Pi)
 =\{-1,\,9-6\sqrt2,\,9+6\sqrt2\}.
\]
In particular, the associated quadratic form has signature \((2,1)\).
\end{theorem}

\begin{proof}
The expansion of the determinant gives \(-9\). The characteristic polynomial
factors as
\[
 (\lambda+1)\bigl((\lambda-9)^2-72\bigr).
\]
Its roots are those indicated and \(9-6\sqrt2>0\); therefore, there exist two
positive directions and one negative direction.
\end{proof}

The spectral decomposition admits an integral form over an invariant sublattice. For
\[
 v_-=(1,-1,0),\qquad v_+=(1,1,0),\qquad v_0=(0,0,1),
\]
one has \(M_\Pi v_-=-v_-\), whereas the restriction of \(M_\Pi\) to the
submodule generated by \(v_+,v_0\) is represented by
\[
 \begin{pmatrix}9&6\\12&9\end{pmatrix}
 =3H,\qquad
 H=\begin{pmatrix}3&2\\4&3\end{pmatrix}\in SL(2,\ZZ).
\]
The eigenvalues of \(H\) are
\[
 3\pm2\sqrt2=(1\pm\sqrt2)^2.
\]
The three vectors \(v_-,v_+,v_0\) generate a sublattice of index two in
\(\ZZ^3\), since the determinant of the matrix having them as columns equals
\(-2\). Therefore, an integral direct sum of all of
\(\ZZ^3\) is not asserted here: on that invariant sublattice, the aggregate operator separates an
eigendirection of eigenvalue \(-1\) and a hyperbolic submodule governed by the
quadratic unit \(3+2\sqrt2\); over \(\RR^3\), the spectral decomposition
is complete.

\section{Commutator of the arithmetic projectors and separation of their images}

The difference between \(\Sigma\) and \(\Pi\) is not reducible to their global sums.
For \(d\ge2\), let \(\mathscr A_d=\{1,\ldots,9\}^d\), and define
\[
 \Sigma_d(x)=\dr(x_1+\cdots+x_d),\qquad
 \Pi_d(x)=\dr(x_1\cdots x_d).
\]
In the Hilbert space \(\ell^2(\mathscr A_d)\), endowed with the
uniform measure, let \(P_{\Sigma,d}\) and \(P_{\Pi,d}\) denote the orthogonal
conditional expectations onto the functions measurable with respect to
\(\Sigma_d\) and \(\Pi_d\), respectively. The normalized joint table is
\[
 C^{(d)}_{ab}
 =\frac{N^{(d)}_{ab}}{\sqrt{n_a m_b}},
\quad
 N^{(d)}_{ab}
 =\#\{x:\Sigma_d(x)=a,\ \Pi_d(x)=b\},
\]
where \(n_a=\sum_bN^{(d)}_{ab}\) and
\(m_b=\sum_aN^{(d)}_{ab}\).

On the original two-dimensional support, the exact joint correlation is
\[
N^{(2)}=
\begin{pmatrix}
0&0&2&0&0&2&3&0&2\\
3&0&2&0&0&2&0&0&2\\
0&2&0&0&2&0&0&2&3\\
0&0&2&3&0&2&0&0&2\\
0&0&2&3&0&2&0&0&2\\
0&2&0&0&2&0&0&2&3\\
3&0&2&0&0&2&0&0&2\\
0&0&2&0&0&2&3&0&2\\
0&2&0&0&2&0&0&2&3
\end{pmatrix}.
\]
The rows, ordered by the value of \(\Sigma\), sum to nine; the columns,
ordered by the value of \(\Pi\), sum to \((6,6,12,6,6,12,6,6,21)\). Complete
enumeration thus determines the joint correlation of the two observables
before any subsequent evaluation map is introduced.

\begin{proposition}[Switches in two and three dimensions]
\label{prop:app-no-conmutatividad}
The two pairs of projectors satisfy
\[
 \|[P_{\Sigma,2},P_{\Pi,2}]\|=\frac{\sqrt2}{3},
 \qquad
 \|[P_{\Sigma,3},P_{\Pi,3}]\|=\frac{\sqrt3}{4}.
\]
In particular, the additive and multiplicative observables do not induce simultaneously compatible decompositions.
\end{proposition}

\begin{proof}
The singular values \(s\) of \(C^{(d)}\) are the cosines of the principal angles
between the two subspaces. In the corresponding two-dimensional
block, the norm of the commutator is \(s\sqrt{1-s^2}\). The exact
enumeration of the \(9^d\) words produces, for \(d=2\), the nontrivial eigenvalues
\[
 s^2\in\left\{\frac57,\frac13,\frac13\right\}.
\]
The maximum of \(s^2(1-s^2)\) is \(2/9\), whence
\(\sqrt2/3\). For \(d=3\), the same construction gives the maximum
\(s^2(1-s^2)=3/16\), attained at \(s^2=1/4\), and therefore the norm is
\(\sqrt3/4\). In both cases, the commutator is nonzero.
\end{proof}

\begin{proposition}[Intersection of the observable subspaces]
\label{prop:app-interseccion-proyectores}
The exact enumeration of the words of \(\mathscr A_d\), for
\(2\leq d\leq15\), verifies
\[
 \operatorname{Ran}P_{\Sigma,d}\cap
 \operatorname{Ran}P_{\Pi,d}
 =\CC\mathbf 1.
\]
In that interval of dimensions, the constant functions constitute the
only subspace common to the two arithmetic decompositions.
\end{proposition}

\begin{proof}
For two orthogonal projections, the dimension of the intersection of their ranges coincides with the multiplicity of the singular value \(1\) in the overlap matrix \(C^{(d)}\). The integer computation of the tables \(N^{(d)}\), followed by the exact computation of that multiplicity, gives the value one for each \(d=2,\ldots,15\). The constant vector belongs to both ranges; it therefore generates the entire intersection.
\end{proof}

The sequence of norms is not extrapolated from the first two cases.
In particular, for \(d=4\) the same calculation gives
\[
 \|[P_{\Sigma,4},P_{\Pi,4}]\|
 =\sqrt{\frac{2618939}{22240656}}
 \simeq0.3431538652,
\]
which does not coincide with the elementary prolongation \(2/(d+1)\). This case
constitutes a counterexample to any formula inferred solely from
\(d=2,3\).

\section{Central block and spectral decomposition}

The characterization of the central block is obtained by means of a
congruence condition. For two consecutive representatives \(k,k+1\), consider
\[
B_k=
\begin{pmatrix}
k^2&k(k+1)\\
k(k+1)&(k+1)^2
\end{pmatrix}\pmod9.
\]

\begin{theorem}[Uniqueness of the Adjacent Diagonal Block]
\label{thm:app-bloque-central}
There is a unique block \(B_k\) whose diagonal entries coincide. It is
determined by \(k=4\) and is
\[
B_c=
\begin{pmatrix}7&2\\2&7\end{pmatrix}.
\]
\end{theorem}

\begin{proof}
The equality of the diagonals is equivalent to
\[
k^2\equiv(k+1)^2\pmod9,
\]
that is, \(2k+1\equiv0\pmod9\). Since \(2\) is a unit of
\(\ZZ/9\ZZ\), the congruence has the unique solution \(k\equiv4\pmod9\).
\end{proof}

Therefore, the positions \(4,5\) of the multiplicative observable form
\[
B_c=
\begin{pmatrix}7&2\\2&7\end{pmatrix}
=7I+2R,
\qquad
R=\begin{pmatrix}0&1\\1&0\end{pmatrix},
\qquad R^2=I.
\]
With \(P_\pm=(I\pm R)/2\),
\[
B_c=9P_++5P_-.
\]
\begin{theorem}[Typed hierarchy of the sum-product separation]
\label{thm:app-jerarquia-cuatro-dos-seis-cincuentaycuatro}
The central block has a local spectral gap (9-5=4). Its
off-diagonal coefficient is (2). Compression by the three residue classes modulo three
produces the aggregate difference (51-45=6), and unfolding across the nine
positions produces the global defect
\[
 9\cdot6=459-405=54.
\]
The integers (4), (2), (6), and (54) belong, respectively, to the
spectral jump, the local coupling, the modular compression, and the
two-dimensional integration.
\end{theorem}

\begin{proof}
The decomposition \(B_c=7I+2R=9P_++5P_-\) gives the jump and the
local coefficient. The aggregated matrices \(M_\Pi,M_\Sigma\) constructed
above satisfy
\(\sum M_\Pi=51\) and \(\sum M_\Sigma=45\). Each aggregated entry represents
nine positions, so unfolding multiplies the difference by
nine and recovers the totals (459) and (405).
\end{proof}
Their ordered concatenations produce the blocks \(72\) and \(27\):
\[
72+27=99,
\qquad
72-27=45=\det B_c.
\]
The concatenation of the diagonal and antidiagonal entries produces,
respectively, \(77\) and \(22\). Therefore,
\[
 72+27=77+22=99,\qquad
 77-22=55=54+1.
\]
These equalities express two exact partitions of the same block; they do not
by themselves identify an external constant.

\begin{proposition}[Local--global relationship of the radical sector]
\label{prop:app-local-global-radical}
The partition by rows of the central block reproduces, with scale factor
three, the partition of the radical sector:
\[
 3\cdot72=216,\qquad
 3\cdot27=81,\qquad
 3\cdot99=297.
\]
\end{proposition}

\begin{proof}
The block \(C_0\times C_0\) contains nine null representatives and sums to
\(81\). The remaining four arms of the rows and columns of \(C_0\)
sum to \(216\). The identities follow from \(72+27=99\).
\end{proof}

Moreover,
\[
M_c:=\frac{B_c}{\sqrt{45}}
=\exp\left(\frac12\log\frac95\,R\right)
\in SO^+(1,1).
\]
If \(\eta=\operatorname{diag}(1,-1)\), the normalized matrix \(M_c\) satisfies
\(M_c^{\mathsf T}\eta M_c=\eta\), \(\det M_c=1\), and belongs to
\(SO^+(1,1)\). This finite realization is confined to the constructed block \(2\times2\);
its space-time interpretation requires the dimensional bridge
developed in the treatise on Dimensional Mechanics.

\begin{proposition}[Spectral family APP--TRIT and reading
Markov--Lorentz local]
\label{prop:app-trit-espectral-markov-lorentz-rev8}
For \(\tau\in\{-1,0,+1\}\), let
\[
 B_\tau=(4+3\tau)I+(5-3\tau)R.
\]
Then
\[
 \operatorname{Spec}(B_\tau)=\{9,\,6\tau-1\},
 \qquad
 B_\tau=9P_++(6\tau-1)P_-,
\]
and the differential mode recovers the orientation through
\[
 \tau=\frac{\lambda_-+1}{6}.
\]
In particular, \(B_{+1}=B_c\). Its stochastic normalization
\[
 P_B:=\frac{1}{9}B_c=P_++\frac59P_-
\]
satisfies, for all \(n\geq0\),
\[
 P_B^n=P_++\left(\frac59\right)^n P_-.
\]
Therefore, the contraction of the antisymmetric mode is
\(\kappa_B=5/9\) and the spectral gap of this local chain is \(4/9\).
If
\[
 \eta_c:=\frac12\log\frac95,
\]
the lorentzian normalization of the same block satisfies
\[
 \frac{B_c}{\sqrt{45}}=\exp(\eta_c R),
 \qquad
 \boxed{\frac59=e^{-2\eta_c}}.
\]
\end{proposition}

\begin{proof}
Since \(R\) acts with eigenvalues \(+1\) and \(-1\) on
\(\operatorname{im}P_+\) and \(\operatorname{im}P_-\), respectively, the
two eigenvalues of \(B_\tau\) are
\[
 (4+3\tau)+(5-3\tau)=9,\qquad
 (4+3\tau)-(5-3\tau)=6\tau-1.
\]
The formula for the powers follows from
\(P_\pm^2=P_\pm\) and \(P_+P_-=0\). Finally,
\(e^{-2\eta_c}=e^{-\log(9/5)}=5/9\).
\end{proof}

This coincidence identifies two exact readings of the local block: the
differential memory decreases by the same quotient that encodes the speed
in the decomposed realization. It does not turn \(B_\tau\) into the global
Markov kernel of TPK, or \(5/9\) into a physical constant. Nor does it by
itself extend the realization \(SO^+(1,1)\) to a Minkowski space
\(3+1\): those steps require additional dimensional maps.

\section{Discrete curvature induced by the interaction of the sum and product leaves and its commutation defect}

On the torus we write, now in residues,
\[
S(x,y)=x+y,
\qquad
P(x,y)=xy\pmod9.
\]
For every function \(T:X_9\to\ZZ/9\ZZ\), we fix the forward differences
\[
 \Delta_xT(x,y)=T(x+1,y)-T(x,y),\qquad
 \Delta_yT(x,y)=T(x,y+1)-T(x,y),
\]
where the coordinates are taken modulo nine. This convention also fixes the
positive orientation of each plaquette: first the direction
\(x\) is traversed, then the direction \(y\), and one returns along the opposite edges.
It is important to type these two objects correctly. \(S\) and \(P\) are
vertex potentials; the link variables are their discrete
differences
\[
A_x^T=\Delta_xT,
\qquad
A_y^T=\Delta_yT.
\]
For an ordered mixture of potentials \((T_x,T_y)\), we define the
plaquette curvature by
\[
F(T_x,T_y)
=\Delta_x A_y^{T_y}-\Delta_y A_x^{T_x}
=\Delta_x\Delta_yT_y-\Delta_y\Delta_xT_x.
\]
Since
\[
\Delta_xS=\Delta_yS=1,
\qquad
\Delta_xP=y,
\qquad
\Delta_yP=x,
\]
is obtained in the \(81\) plaquettes:
\[
F(S,S)=F(P,P)=0,
\qquad
F(S,P)=+1,
\qquad
F(P,S)=-1\pmod9.
\]
The connections associated with a single observable are flat; the mixed composition of the two observables produces curvature of opposite sign according to the order. For a coordinate rectangle \(R_{a,b}\) contained in the chart, we define the additive boundary Wilson observable as the oriented sum of the link variables. Cancellation of the interior edges gives
\[
 W_{\partial R_{a,b}}(S,P)=ab,\qquad
 W_{\partial R_{a,b}}(P,S)=-ab\pmod9,
\]
which is the finite Stokes identity for this convention.

The distinction between potential and linkage is algebraically necessary. If
\(A_x,A_y\) were directly substituted by \(S,P\), terms depending on
\(x,y\) would appear and the preceding constant values would not be obtained. The typed
formulation is what produces the mixed cocycle used later in the finite
Heisenberg closure.

\begin{remark}
The integrated difference \(54\) and the mixed curvature \(\pm1\) are distinct invariants. The former compares global sums; the latter depends on the order of the two potentials and is localized on each plaquette. The formulation in terms of self-adjoint operators and commutators will be developed after the corresponding state space has been constructed.
\end{remark}

\section{Path Transition Operator Expansion}

The definition of APP includes the space
\(\operatorname{Path}(\GraphAPP)\). Once a transition operator
\(U\) has been fixed on the vertices of \(\GraphAPP\), we assume that
\(U_{yx}=0\) whenever \(x\to y\) is not an edge of the graph. Under this
support hypothesis, its powers admit the expansion
\begin{proposition}[Finite Path Expansion]
\label{prop:app-suma-caminos}
For each \(R\ge1\) and each pair of vertices \(i,f\),
\[
(U^R)_{fi}
=\sum_{\substack{\gamma:i\to f\\|\gamma|=R}}
\prod_{k=0}^{R-1}U_{x_{k+1},x_k},
\qquad
\gamma=(x_0=i,x_1,\ldots,x_R=f).
\]
\end{proposition}

\begin{proof}
For \(R=1\) the identity is the definition of the matrix elements of
\(U\). If it holds for \(R\), then
\[
 (U^{R+1})_{fi}
 =\sum_j U_{fj}(U^R)_{ji}.
\]
Substituting the induction hypothesis concatenates to every path of length \(R\)
from \(i\) to \(j\) the final edge \(j\to f\); thus, every path of
length~\(R+1\) appears exactly once.\qedhere
\end{proof}

Consequently, APP brings together in a single structure the finite arithmetic torus,
the sum and product observables, the residue--quotient decomposition, the
nilpotent radical, and an oriented mixed curvature. The representative \(9\) of
the null class fixes a visible chart of the quotient; replacing it with \(0\) without
simultaneously transporting the chart alters the registers used by the
subsequent dynamics.
\fi
 \end{HMTScientificOwnerSubordinated}
\endgroup

% Unidad U003. Redacción autónoma; tablas y censos reconstruidos desde 1063/live.

\section{Aritmetic Matrix of APP and Local Ring Structure}
\label{p12-83-c1-s1:chap:app-hojas-algebra-local}

The \cref{p12-83-c1-s2:chap:app-completacion-gnomonica} constructed a single support
\(X_9\). APP has a single principal arithmetic matrix: nonadic multiplication
published by the positive digital root. Addition does not constitute a
second APP table; it is the internal accumulation law that generates the rows and
columns of the multiplicative matrix and provides an auxiliary observable for
comparison. This distinction will be maintained throughout the work.

\subsection{Nonadic multiplicative matrix and digital root}

For \(a,b\in\mathcal D_9=\{1,\ldots,9\}\), the APP entry is obtained
directly through
\begin{equation}
 \boxed{\Pi(a,b)=\operatorname{dr}_9(ab)
 =1+((ab-1)\bmod9).}
 \label{p12-83-c1-s1:eq:app-matriz-multiplicativa-directa}
\end{equation}
The only APP matrix is, therefore,
\[
\begin{array}{c|ccccccccc}
\Pi&1&2&3&4&5&6&7&8&9\\\hline
1&1&2&3&4&5&6&7&8&9\\
2&2&4&6&8&1&3&5&7&9\\
3&3&6&9&3&6&9&3&6&9\\
4&4&8&3&7&2&6&1&5&9\\
5&5&1&6&2&7&3&8&4&9\\
6&6&3&9&6&3&9&6&3&9\\
7&7&5&3&1&8&6&4&2&9\\
8&8&7&6&5&4&3&2&1&9\\
9&9&9&9&9&9&9&9&9&9
\end{array}
\]
Each row can also be reconstructed by reiterated addition of its index,
but that recurrence describes the internal law of the matrix and not a second
sheet. The rows and columns of indices \(1,2,4,5,7,8\) traverse the nine
classes; those of indices \(3,6,9\) are exactly those that contract the support.

In the residual chart, the set of units is
\[
 R_9^\times=(\mathbb Z/9\mathbb Z)^\times
   =\{[1],[2],[4],[5],[7],[8]\}.
\]
The maximal ideal is
\[
 \mathfrak m=(3)=\{[0],[3],[6]\}.
\]

\begin{theorem}[Classification of multiplicative rows]
\label{p12-83-c1-s1:thm:app-clasificacion-filas-producto}
For \(a\in\mathcal D_9\), the map
\[
 m_a:\mathcal D_9\longrightarrow\mathcal D_9,
 \qquad b\longmapsto\Pi(a,b)
\]
is bijective if and only if \([a]_9\in R_9^\times\). If
\([a]_9\in\mathfrak m\setminus\{0\}\), its image has three elements. If
\([a]_9=[0]_9\), its image has a single element.
\end{theorem}

\begin{proof}
The residual reduction of \(m_a\) is multiplication by \([a]_9\). In a
finite ring, that multiplication is bijective exactly when \([a]_9\)
is a unit. If \([a]_9=[3]\) or \([6]\), its image is the ideal
\(\mathfrak m\), of cardinality three. If \([a]_9=[0]\), every entry is sent
to the null class, published as \(9\).
\end{proof}

\begin{proposition}[Internality of the additive law]
\label{p12-83-c1-s1:prop:app-latinidad-aditiva}
The auxiliary observable
\[
 \Sigma(a,b)=\operatorname{dr}_9(a+b)
\]
acts by translations: for each \(a,c\in\mathcal D_9\) there exists a unique
\(b\in\mathcal D_9\) such that \(\Sigma(a,b)=c\), and analogously in the other
coordinate.
\end{proposition}

\begin{proof}
In the residual chart, \([a+b]_9=[c]_9\) has the unique solution
\([b]_9=[c-a]_9\). The positive section transports that class to a unique
element of \(\mathcal D_9\).
\end{proof}

The difference of geometry of fibers remains thus typed without duplicar the matrix
APP: the law additive internal only translates, whereas the matrix multiplicative
permutes on the units and contracts on the radical.

\begin{proposition}[Linear structure of the radical]
\label{p12-83-c1-s1:prop:app-radical-lineal}
The map
\[
 M_3:R_9\longrightarrow R_9,
 \qquad x\longmapsto3x
\]
satisfies
\[
 \ker M_3=\operatorname{im}M_3=\mathfrak m.
\]
The induced map
\[
 R_9/\mathfrak m\longrightarrow\mathfrak m,
 \qquad [x]\longmapsto3x
\]
is an isomorphism of vector spaces over \(\mathbb F_3\), but not a
ring isomorphism.
\end{proposition}

\begin{proof}
The congruence \(3x\equiv0\pmod9\) is equivalent to \(x\equiv0\pmod3\), whence
the kernel \(\mathfrak m\). The image is
\(\{[0],[3],[6]\}=\mathfrak m\). The first isomorphism theorem gives the
linear isomorphism. It is not a ring isomorphism because
\(\mathfrak m^2=0\) in \(R_9\), whereas the quotient
\(R_9/\mathfrak m\cong\mathbb F_3\) has a nonzero unit.
\end{proof}

The positive section publishes the non-trivial radical cycle as
\[
 369\longmapsto693\longmapsto936\longmapsto369,
\]
according to the position from which the triple is read. This notation preserves the
order of the coordinates; it does not turn the ideal into three independent
integers.

\subsection{Lifts and carry cocycle}

Let \(s_0:R_9\to\{0,\ldots,8\}\) be the normalized residual section. The
binary carry of the sum is
\[
 \kappa_0(a,b)
 =\frac{s_0(a)+s_0(b)-s_0(a+b)}9.
\]

\begin{lemma}[Cocyclic Identity]
\label{p12-83-c1-s1:lem:app-identidad-cociclo-acarreo}
For \(a,b,c\in R_9\),
\[
 \kappa_0(a,b)+\kappa_0(a+b,c)
 =\kappa_0(b,c)+\kappa_0(a,b+c).
\]
\end{lemma}

\begin{proof}
Multiply by nine and expand either of the two sides produces
\[
 s_0(a)+s_0(b)+s_0(c)-s_0(a+b+c).
\]
Therefore, both are equal.
\end{proof}

Define \(\chi_0(a)=1\) if \(a=[0]_9\) and \(0\) otherwise. Since
\[
 s_+(a)=s_0(a)+9\chi_0(a),
\]
the carryovers of the two sections are related by
\[
 \kappa_+(a,b)=
 \kappa_0(a,b)+\chi_0(a)+\chi_0(b)-\chi_0(a+b).
\]
Changing the representative modifies the cocycle by a coboundary. This is the
exact form for transporting memory when the null class is published as
\(9\) instead of \(0\).

\subsection{Integrated censuses of residues and quotients}

Summing over all pairs \((i,j)\in\mathcal D_9^2\), one obtains
\[
 \sum_{i,j}(i+j)=810,
 \qquad
 \sum_{i,j}ij=2025.
\]
The positive residues satisfy
\[
 \sum_{i,j}R^+_{i,j}=405,
 \qquad
 \sum_{i,j}R^\times_{i,j}=459.
\]
By the raised identities,
\[
 \sum_{i,j}Q^+_{i,j}=\frac{810-405}{9}=45,
 \qquad
 \sum_{i,j}Q^\times_{i,j}=\frac{2025-459}{9}=174.
\]

\begin{theorem}[Exact separation of the sum-product excess]
\label{p12-83-c1-s1:thm:app-separacion-exceso}
The integrated difference between product and sum decomposes uniquely
into a published part and a quotient part:
\[
 \boxed{
  2025-810
  =(459-405)+9(174-45)
  =54+9\cdot129.}
\]
\end{theorem}

\begin{proof}
We subtract, term to term, the identities
\(ij=R^\times_{i,j}+9Q^\times_{i,j}\) and
\(i+j=R^+_{i,j}+9Q^+_{i,j}\), and sum on the eighty and a pairs. The
censuses preceding dan the equality. The uniqueness proceeds of the division
Euclidean in each pair.
\end{proof}

The value \(54\) is not the complete difference between two supposed matrices.
It is the visible defect between the multiplicative matrix and the additive accumulator
in the positive chart. The term \(9\cdot129\) is the contribution
that remains in the quotients. Omitting it would confuse publication with
genealogy.

\subsection{Central block of \(2\) positions}

For \(k\in\mathbb Z/9\mathbb Z\), consideramos the block symmetric
\[
 B_k=
 \begin{pmatrix}
  k^2&k(k+1)\\
  k(k+1)&(k+1)^2
 \end{pmatrix}\pmod9.
\]
The diagonals coincide exactly when
\[
 k^2\equiv(k+1)^2\pmod9
 \quad\Longleftrightarrow\quad
 2k+1\equiv0\pmod9.
\]
Since \(2^{-1}=5\pmod9\), the unique solution is \(k=4\). Thus appears the
block
\[
 B_c=
 \begin{pmatrix}7&2\\2&7\end{pmatrix}.
\]
Let
\[
 R=\begin{pmatrix}0&1\\1&0\end{pmatrix},
 \qquad
 P_\pm=\frac12(I\pm R).
\]
Then
\[
 B_c=7I+2R=9P_++5P_-.
\]

\begin{theorem}[Uniqueness and spectrum of the central block]
\label{p12-83-c1-s1:thm:app-bloque-central}
The block \(B_c\) is the unique adjacent block in the family \(B_k\) with
equal diagonals. Its eigenvalues on the symmetric and
antisymmetric subspaces are, respectively, \(9\) and \(5\).
\end{theorem}

\begin{proof}
The congruence before proves the uniqueness. Since that \(R\) acts by
\(+1\) on \(\operatorname{im}P_+\) and by \(-1\) on
\(\operatorname{im}P_-\), the matrix \(7I+2R\) acts by \(9\) and \(5\),
respectively.
\end{proof}

The two ordered readings of its digits give
\[
 72+27=99,
 \qquad
 72-27=45=\det B_c,
\]
and the two diagonals give
\[
 77+22=99,
 \qquad
 77-22=55=54+1.
\]
These identities are internal controls of the block; they do not by themselves define the global defect.

\subsection{Mixed curvature of the \(2\) laws}

On the residual chart we write
\[
 S(x,y)=x+y,
 \qquad
 P(x,y)=xy.
\]
For a function \(T:X_9\to R_9\), let
\[
 \Delta_xT(x,y)=T(x+1,y)-T(x,y),
 \qquad
 \Delta_yT(x,y)=T(x,y+1)-T(x,y).
\]
We define the ordered curvature of a pair of laws by
\[
 \mathcal F(T_x,T_y)
 :=\Delta_x\Delta_yT_y-\Delta_y\Delta_xT_x.
\]

\begin{proposition}[Orientations of the mixed composition]
\label{p12-83-c1-s1:prop:app-curvatura-mixta}
The following hold
\[
 \mathcal F(S,S)=0,
 \qquad
 \mathcal F(P,P)=0,
\]
and, for mixed compositions,
\[
 \mathcal F(S,P)=+1,
 \qquad
 \mathcal F(P,S)=-1
 \quad\text{in }R_9.
\]
\end{proposition}

\begin{proof}
The mixed second differences of the linear law \(S\) vanish. For
\(P(x,y)=xy\), both pure mixed differences equal one and cancel when
compared in the same order. Combining a linear law and a bilinear one
leaves a single unit whose sign changes when the order of the sheets is reversed.
\end{proof}

The curvature not creates another matrix APP. Is recorded that the order
\(\Sigma\to\Pi\) and the order \(\Pi\to\Sigma\) are orientations opposite,
whereas the sectors pure remain flat. This triad
\(-1,0,+1\) will be typed in the unit TRIT.

\subsection{Obstructions of the arithmetic matrix of APP: translations, local radical, central block, mixed curvature, and conservation of carry}

\begin{criterion}[Refutation criteria for the APP matrix]
The unit is refuted if:
\begin{enumerate}
 \item the internal additive law ceases to act by translations;
 \item a nonunitary multiplicative row is declared bijective;
 \item \(R_9/\mathfrak m\) is identified with \(\mathfrak m\) as a ring;
 \item the exceso \(1215\) it replaces by \(54\) without conservar quotients;
 \item the central block is chosen without resolving \(2k+1=0\pmod9\);
 \item reversal of the mixed order does not change the sign of the curvature;
 \item the carry it eliminates to the cambiar between the sections \(s_0\) and
       \(s_+\).
\end{enumerate}
\end{criterion}

APP is now constructed as support, category of trajectories, nonadic
multiplicative matrix, internal additive law, local radical, carry cocycle,
and mixed orientation. The following unit
will not summarize these structures: it will take the triad of orientations
that has just emerged and construct the complete TRIT state with its three
quadratic regimes.


% Unidad U002. Redacción autónoma a partir del generador integral 1063/live.

\section{Gnomonic completion and arithmetic support of APP}
\label{p12-83-c1-s2:chap:app-completacion-gnomonica}

This unit constructs the APP support from the
nonadic cell of the \cref{p12-83-c1-s4:chap:segmentos-marcas-intersticios}. It does not presuppose a
continuous plane, an infinite lattice, or a numerical constant. The starting
object is finite: two ordered copies of the positive chart of
\(\mathbb Z/9\mathbb Z\). Completion of its interior interstices
produces exactly eighty-one states, and the residual translation law
turns that square support into a discrete torus.

Throughout the unit, APP denotes a mathematical structure and not an image.
Henceforth, the acronym APP names this structure; the properties that
follow depend on its objects and morphisms, not on a verbal expansion of the
acronym.

\subsection{Correspondence between \(2\) charts of the same residue class}

We define
\[
 \mathcal D_9:=\{1,2,\ldots,9\},
 \qquad
 \mathcal D_8:=\{1,2,\ldots,8\}.
\]
The residual publication application is
\[
 \lambda_9:\mathcal D_9\longrightarrow\mathbb Z/9\mathbb Z,
 \qquad
 \lambda_9(a)=[a]_9.
\]
Therefore, \(\lambda_9(9)=[0]_9\). Its inverse as a set map is the positive section
\[
 s_+([a]_9)=
 \begin{cases}
  a,&[a]_9\neq[0]_9,\quad 1\leq a\leq8,\\
  9,&[a]_9=[0]_9.
 \end{cases}
\]
The two expressions \(9\) and \([0]_9\) are not two phases: they are two
representatives of the same element in different charts.

For two coordinates we will write
\[
 \lambda_9^{\times2}:\mathcal D_9^2
       \longrightarrow X_9:=(\mathbb Z/9\mathbb Z)^2,
 \qquad
 (a,b)\longmapsto([a]_9,[b]_9).
\]
This map is a bijection of sets. The chart
\(\mathcal D_9^2\) makes the positive marks visible; the group \(X_9\)
makes the residual sum and the translations visible.

\begin{lemma}[State-preserving chart change]
\label{p12-83-c1-s2:lem:app-cambio-carta-sin-perdida}
The maps \(\lambda_9^{\times2}\) and
\(s_+^{\times2}\) are inverses. In particular, the positive representation and the
residual representation contain exactly the same eighty-one states.
\end{lemma}

\begin{proof}
In each coordinate,
\(\lambda_9\circ s_+=\operatorname{id}_{\mathbb Z/9\mathbb Z}\) and
\(s_+\circ\lambda_9=\operatorname{id}_{\mathcal D_9}\) hold. The Cartesian product of these identities gives the two required identities.
\end{proof}

\subsection{Interior kernel and gnomonic complement}

The inner core is
\[
 \mathcal N_8:=\mathcal D_8\times\mathcal D_8,
 \qquad |\mathcal N_8|=8^2=64.
\]
In the plane of indices, \(\mathcal N_8\) is the square that not uses the
mark of closure in no coordinate. For publish also the null class
it must completar the ninth row and the ninth column. We define the
complement gnomonic disjoint by
\[
 \mathcal G_9:=
  (\{9\}\times\mathcal D_9)
  \sqcup
  (\mathcal D_8\times\{9\}).
\]
The corner \((9,9)\) belongs to the first piece; for that reason the second uses
\(\mathcal D_8\) and not \(\mathcal D_9\). Consequently,
\[
 |\mathcal G_9|=9+8=17,
 \qquad
 \mathcal D_9^2=\mathcal N_8\sqcup\mathcal G_9,
 \qquad
 81=64+17.
\]

\begin{proposition}[Minimality of the Completion]
\label{p12-83-c1-s2:prop:app-minimalidad-completacion}
Let \(Y\) be a subset of \(\mathcal D_9^2\) that contains \(\mathcal N_8\) and whose projection onto each coordinate is all of \(\mathcal D_9\). If, in addition, \(Y\) contains all pairs required to translate one coordinate while the other remains arbitrary, then \(\mathcal G_9\subseteq Y\). Therefore, the addition of seventeen states is the minimal completion compatible with the two translations.
\end{proposition}

\begin{proof}
To transport the first coordinate through the closure class while
the second remains at an arbitrary value \(b\in\mathcal D_9\), all
pairs \((9,b)\) are necessary. Nine states are obtained. To
transport the second coordinate through closure while the first assumes
an arbitrary value \(a\in\mathcal D_8\), the remaining eight pairs
\((a,9)\) are necessary. The corner \((9,9)\) was already counted in the first
family. In total, \(9+8=17\) states are necessary, exactly those of
\(\mathcal G_9\).
\end{proof}

The hypothesis of compatibility with the two translations is essential. It would suffice
to add a single mark per coordinate for both projections to be
surjective, but that set would not contain the boundary states
needed to move one coordinate while holding the other fixed. APP requires the
complete product, not two disconnected axes.

\subsection{Exact realization of the interstitial interval via APP cells and boundary coordinates}

Let \(P_9\) be the ordered path with nine marks and eight interior edges
\[
 E(P_9)=\{e_1,\ldots,e_8\}.
\]
Closing the endpoints incorporates an edge \(e_9\) and produces the cycle
\(C_9\):
\[
 E(C_9)=E(P_9)\sqcup\{e_9\}.
\]
The pairs of interior interstices form \(E(P_9)^2\). The pairs of
interstices of the complete cycle form \(E(C_9)^2\). The difference is
\begin{align}
 E(C_9)^2\setminus E(P_9)^2
   &=(\{e_9\}\times E(C_9))
     \sqcup(E(P_9)\times\{e_9\}).
\label{p12-83-c1-s2:eq:app-complemento-intersticial}
\end{align}
The decomposition is disjoint because the corner \((e_9,e_9)\) is reserved for
the first piece.

\begin{theorem}[Equivalence of the mark model and the interstice model]
\label{p12-83-c1-s2:thm:app-equivalencia-modelos}
The map
\[
 \beta:E(C_9)^2\longrightarrow\mathcal D_9^2,
 \qquad
 \beta(e_i,e_j)=(i,j),
\]
is bijective. Its restrictions satisfy
\[
 \beta(E(P_9)^2)=\mathcal N_8,
 \qquad
 \beta(E(C_9)^2\setminus E(P_9)^2)=\mathcal G_9.
\]
Therefore, the identity \(81=64+17\) simultaneously counts pairs of
published marks and pairs of closed interstices.
\end{theorem}

\begin{proof}
Each edge of \(C_9\) has a unique index in \(\mathcal D_9\). The map \((i,j)\mapsto(e_i,e_j)\) is the explicit inverse of \(\beta\). If \(i,j\leq8\), the pair belongs to \(E(P_9)^2\) and its image belongs to \(\mathcal D_8^2\). If at least one of the indices is nine, the pair belongs to the complement of \cref{p12-83-c1-s2:eq:app-complemento-intersticial}; separating the case \(i=9\) from the case \(i\leq8,j=9\) produces exactly the two pieces of \(\mathcal G_9\).
\end{proof}

\subsection{Toroidal closure and translations}

Let
\[
 \mathbf e_1=([1]_9,[0]_9),
 \qquad
 \mathbf e_2=([0]_9,[1]_9).
\]
The set of elementary directions is
\[
 \mathscr D_{\rm APP}:=
 \{+\mathbf e_1,-\mathbf e_1,+\mathbf e_2,-\mathbf e_2\}.
\]
We define the oriented Cayley graph.
\[
 \Gamma_{\rm APP}:=\operatorname{Cay}(X_9,\mathscr D_{\rm APP}).
\]
Its vertices are the eighty-one elements of \(X_9\). For every
\(x\in X_9\) and \(d\in\mathscr D_{\rm APP}\), there is an oriented edge
\(x\xrightarrow{d}x+d\). The inverse edge bears the label \(-d\).

In the positive chart, unit translation is written by means of the published
successor and its carry:
\[
 u_+(a)=
 \begin{cases}
  a+1,&a<9,\\
  1,&a=9,
 \end{cases}
 \qquad
 \kappa_+(a)=
 \begin{cases}
  0,&a<9,\\
  1,&a=9.
 \end{cases}
\]
The raised identity is
\begin{equation}
 a+1=u_+(a)+9\kappa_+(a).
\label{p12-83-c1-s2:eq:app-sucesor-publicado-acarreo}
\end{equation}
Upon applying \(\lambda_9\), the carry term vanishes, leaving
\[
 \lambda_9(u_+(a))=\lambda_9(a)+[1]_9.
\]

\begin{lemma}[Compatibility between closure, residue, and elevation]
\label{p12-83-c1-s2:lem:app-compatibilidad-cierre-residuo}
The addition of \(e_9\), the identification \(9\sim[0]_9\), and the translation
\(x\mapsto x+\mathbf e_i\) describe three compatible operations in
distinct domains. The residual quotient publishes the return; the carry of
\cref{p12-83-c1-s2:eq:app-sucesor-publicado-acarreo} preserves the elevation that the
quotient omits.
\end{lemma}

\begin{proof}
Topological closure incorporates the edge joining the endpoints of the path. The
residual chart identifies the published endpoint with the null class. If
\(a<9\), the successor has zero carry; if \(a=9\), the lifted equality
is \(10=1+9\). Both cases descend to the same residual sum, but the
value of \(\kappa_+\) distinguishes them.
\end{proof}

\begin{corollary}[Toroidal support]
\label{p12-83-c1-s2:cor:app-soporte-toroidal}
Identifying opposite sides of the chart \(\mathcal D_9^2\) yields a
geometric realization of the group \(X_9\). Each elementary translation is
defined at every vertex and has an inverse.
\end{corollary}

\subsection{APP Path Category: Objects, Prolongation Morphisms and Compatible Composition}

\begin{definition}[Oriented trajectory]
A trajectory of length \(n\geq0\) is a pair
\[
 \gamma=(x_0;d_1,\ldots,d_n),
 \qquad
 x_0\in X_9,\qquad d_k\in\mathscr D_{\rm APP}.
\]
Its vertices are
\[
 x_k=x_0+\sum_{r=1}^k d_r,
 \qquad 0\leq k\leq n.
\]
The source is \(s(\gamma)=x_0\) and the target is
\(t(\gamma)=x_n\). For \(n=0\) the identity trajectory in
\(x_0\) is obtained.
\end{definition}

If \(t(\gamma)=s(\delta)\), its concatenation is
\[
 \delta\circ\gamma
  =(x_0;d_1,\ldots,d_n,e_1,\ldots,e_m),
\]
where \(e_1,\ldots,e_m\) are the directions of \(\delta\). Concatenation is associative, and paths of zero length are identities. We thus obtain the free category
\[
 \mathsf{Path}(\Gamma_{\rm APP}).
\]

\begin{proposition}[Finiteness of Support and Countability of Trajectories]
\label{p12-83-c1-s2:prop:app-soporte-finito-trayectorias-numerables}
The APP state set is finite, of cardinality \(81\). For a fixed source,
there are \(4^n\) direction words of length \(n\). The union
of all finite trajectories is countable.
\end{proposition}

\begin{proof}
The first assertion follows from \(|X_9|=9^2\). At each of the \(n\)
steps there are four independent choices, whence \(4^n\). Finally,
\(\bigcup_{n\geq0}X_9\times\mathscr D_{\rm APP}^n\) is a countable union
of finite sets.
\end{proof}

This proposition separates two levels that will not be identified: APP has a
finite support, but admits trajectories of arbitrary length. The
prolongation of a route does not create new vertices of the support; it creates new
genealogy over them.

\subsection{Multiplicative matrix and internal additive law on a single support}

For \(n\in\mathbb Z_{>0}\), the root digital positive is
\[
 \operatorname{dr}_9(n):=1+((n-1)\bmod9).
\]
On the positive chart we define the principal matrix and the internal law.
\[
 \Sigma(a,b):=\operatorname{dr}_9(a+b),
 \qquad
 \Pi(a,b):=\operatorname{dr}_9(ab).
\]
Both maps have the same domain \(\mathcal D_9^2\) and the same
codomain \(\mathcal D_9\), but they do not have the same status: \(\Pi\) publishes
the unique APP matrix, and \(\Sigma\) generates its accumulations. Their residual versions
are
\[
 \bar\Sigma([a]_9,[b]_9)=[a+b]_9,
 \qquad
 \bar\Pi([a]_9,[b]_9)=[ab]_9.
\]
The change of chart is compatible with both operations:
\begin{align*}
 \lambda_9\circ\Sigma
   &=\bar\Sigma\circ\lambda_9^{\times2},\\
 \lambda_9\circ\Pi
   &=\bar\Pi\circ\lambda_9^{\times2}.
\end{align*}

To preserve the integer quotient, lifts are introduced.
\begin{align}
 a+b &= R^+_{a,b}+9Q^+_{a,b},
 &1\leq R^+_{a,b}\leq9,
\label{p12-83-c1-s2:eq:app-lift-suma}\\
 ab &= R^\times_{a,b}+9Q^\times_{a,b},
 &1\leq R^\times_{a,b}\leq9.
\label{p12-83-c1-s2:eq:app-lift-producto}
\end{align}
Here
\[
 R^+_{a,b}=\Sigma(a,b),
 \qquad
 R^\times_{a,b}=\Pi(a,b),
\]
and the quotients are uniquely determined by
\[
 Q^+_{a,b}=\frac{a+b-R^+_{a,b}}9,
 \qquad
 Q^\times_{a,b}=\frac{ab-R^\times_{a,b}}9.
\]

\begin{definition}[APP genealogical]
\label{p12-83-c1-s2:def:app-genealogica-autonoma}
The APP structure at this stage is the tuple
\[
 \operatorname{APP}:=
 \bigl(
  X_9,
  \Gamma_{\rm APP},
  \mathsf{Path}(\Gamma_{\rm APP}),
  \Sigma,\Pi,
  R^+,Q^+,R^\times,Q^\times
 \bigr).
\]
The genealogy of a trajectory preserves, for each vertex visited, the
position, incoming direction, multiplicative entry, additive accumulation,
and their respective residues and quotients.
\end{definition}

\begin{theorem}[Causal construction of the APP support]
\label{p12-83-c1-s2:thm:app-cadena-constructiva-autonoma}
The chain
\[
 \mathcal D_9
 \longrightarrow
 \mathcal D_8^2
 \longrightarrow
 \mathcal D_8^2\sqcup\mathcal G_9
 \longrightarrow
 \mathcal D_9^2
 \xrightarrow{\lambda_9^{\times2}}
 X_9
 \longrightarrow
 \Gamma_{\rm APP}
 \longrightarrow
 \mathsf{Path}(\Gamma_{\rm APP})
 \longrightarrow
 (\Sigma,\Pi)
\]
constructs, without loss of cardinality, the support, motion, paths, APP multiplicative matrix and its internal additive law.
\end{theorem}

\begin{proof}
The \cref{p12-83-c1-s2:prop:app-minimalidad-completacion} constructs the complete chart from
the kernel. The \cref{p12-83-c1-s2:thm:app-equivalencia-modelos} identifies its
interstitial model. The \cref{p12-83-c1-s2:lem:app-cambio-carta-sin-perdida} transports
the eighty-one states to \(X_9\). The Cayley graph incorporates the four
reversible translations without changing the vertices. The free category
incorporates all finite concatenations. Finally, \(\Sigma\) and
\(\Pi\) are distinct total maps on the same domain; the
equations \cref{p12-83-c1-s2:eq:app-lift-suma,p12-83-c1-s2:eq:app-lift-producto} preserve the part
that residual reduction does not publish.
\end{proof}

\subsection{Decomposition of the APP orthogram into kernel \(8\times8\), gnomon of \(17\) states, and chart \(9\times9\) with opposite sides identified}

The \cref{fig:app-ortograma} must be read in three layers: the inner square
of sixty-four states, the gnomon of seventeen states, and the
identification of opposite sides. The nine-by-nine chart thereby completes the
eight-by-eight core by means of the closure gnomon. The multiplicative
matrix and the internal additive law are evaluated at every vertex of the same
state, preserve distinct residue--quotient lifts, and do not constitute
two APP squares.

\subsection{Obstructions to the fidelity of the APP support: cardinals \(81/17\), reversible translations, and separation of the sum and product sheets}

\begin{criterion}[Refutation criteria for the APP support]
The construction is refuted if any of the following facts
occurs:
\begin{enumerate}
 \item the complete support has cardinality other than \(81\);
 \item the gnomonic complement does not have cardinal \(17\);
 \item a state of \(X_9\) lacks any of the four translations;
 \item it identify \(\Sigma\) and \(\Pi\) because coincide in a value;
 \item \(Q^+\) or \(Q^\times\) is eliminated upon returning through a boundary;
 \item trajectories of different lengths are presented as new
       support vertices;
 \item it interpreta \(9\) and \([0]_9\) as phases independent.
\end{enumerate}
\end{criterion}

The unit has constructed APP up to its complete genealogical evaluation. It has not yet
oriented the difference between the matrix and its accumulator, has not
introduced the TRIT,
and has not defined the Topological Phase Kernel. The following unit will study
the algebraic structure of rows and columns, the sector of units, and the
carry cocycle that prepares that orientation.


% Unidad U005. Redacción autónoma desde 03c/03h/20 de la autoridad 1063/live.

\long\def\HMTDeferredTPKBase{%
\section{Topological Phase Kernel: observable base and integral holonomic state}
\label{p12-83-c1-s3:chap:tpk-categoria-estado-enriquecido}

APP has supplied a support, a multiplicative matrix, and an internal
additive law; TRIT has
supplied a lifted orientation. The Topological Phase Kernel, hereinafter
TPK, specifies how those data are transported jointly, which
coordinates are preserved when a route is closed, and which readers can be published
without identifying the publication with the state. Its mathematical object is a
category over a base of paths, endowed with an explicit family of
fibers and prolongation relations.

\subsection{The double base of APP paths}

Let \(\mathsf P_{\rm APP}\) be the category libre constructed in the
\cref{p12-83-c1-s2:chap:app-completacion-gnomonica}. We define
\[
 \mathsf P_{\rm APP}^{(2)}
 :=\mathsf P_{\rm APP}\times\mathsf P_{\rm APP}.
\]
An object is a pair of positions
\((p^\Sigma,p^\Pi)\in X_9^2\). A morphism is a pair of trajectories; one
of them may be the identity. The first component is evaluated by
\(\Sigma\), the second by \(\Pi\). They are not two lattices: they are two typed
cursors over the same APP support.

The set of cardinal directions is
\[
 \mathscr D_{\rm card}:=\{N,E,S,O\},
\]
with shifts
\[
 \delta(N)=(-1,0),\quad
 \delta(E)=(0,1),\quad
 \delta(S)=(1,0),\quad
 \delta(O)=(0,-1)
 \quad\text{in }X_9.
\]
The orientation involution is
\[
 \overline N=S,
 \quad \overline S=N,
 \quad \overline E=O,
 \quad \overline O=E.
\]

\subsection{Nonadic phase, dodecaphasic sector and calendar}

The complete clock of this stage is
\[
 \Theta_{108}:=\mathbb Z/108\mathbb Z.
\]
For \(\theta\in\Theta_{108}\), let
\(\widetilde\theta\in\{0,\ldots,107\}\) be its normalized representative.
We define three distinct coordinates:
\begin{align*}
 g_0(\theta)&=[\widetilde\theta+1]_9\in C_9,\\
 r(\theta)&=1+(\widetilde\theta\bmod9)\in\mathcal D_9,\\
 \beta(\theta)&=
 \left[\left\lfloor\frac{\widetilde\theta}{9}\right\rfloor\right]_{12}
 \in C_{12}.
\end{align*}
The first is the phase class, the second its positive label, and the third
the dodecaphasic sector. The map \(\beta\) is a block map;
it is not the canonical reduction \(\theta\bmod12\).

The signature and the active sheet are
\[
 \tau(\theta):=\varepsilon_9(r(\theta))\in\mathsf{TRIT},
\]
and
\[
 \mu(\theta)=
 \begin{cases}
  \Sigma,&\tau(\theta)=+1,\\
  \Pi,&\tau(\theta)=-1,\\
  \varnothing,&\tau(\theta)=0.
 \end{cases}
\]
The symbol \(\varnothing\) denotes a threshold step in which neither cursor
moves; it does not denote the absence of a state.

\begin{definition}[Observable configuration]
\label{p12-83-c1-s3:def:tpk-configuracion-observable}
The observable base is
\[
 \mathcal Q_{\rm obs}
 :=X_9^2\times\mathscr D_{\rm card}^2\times\Theta_{108}.
\]
An element is written
\[
 q=(p^\Sigma,p^\Pi,d^\Sigma,d^\Pi,\theta).
\]
It contains two positions, two cardinal orientations, and a calendar index.
\end{definition}

\subsection{Sheet selector \(M_{\rm ph}:\mathcal D_9\to\mathsf{TRIT}\) and joint evolution of position, orientation, and calendar}

The function
\[
 M_{\rm ph}:\mathcal D_9\longrightarrow\mathsf{TRIT},
 \qquad M_{\rm ph}(r)=\varepsilon_9(r)
\]
is the internal sheet selector. It selects which evaluation acts; it does not by itself
displace any position.

We define the orientation inversion indicator by
\[
 \eta_{27}(\theta)=
 \begin{cases}
  -1,&\widetilde\theta+1\in\{27,54,81,108\},\\
  +1,&\text{otherwise}.
 \end{cases}
\]
The new directions are
\[
 d^{\nu,+}=
 \begin{cases}
  \overline{d^\nu},&\eta_{27}(\theta)=-1,\\
  d^\nu,&\eta_{27}(\theta)=+1,
 \end{cases}
 \qquad \nu\in\{\Sigma,\Pi\}.
\]
The positions evolve via
\[
 (p^{\Sigma,+},p^{\Pi,+})=
 \begin{cases}
 (p^\Sigma+\delta(d^\Sigma),p^\Pi),&\tau(\theta)=+1,\\
 (p^\Sigma,p^\Pi+\delta(d^\Pi)),&\tau(\theta)=-1,\\
 (p^\Sigma,p^\Pi),&\tau(\theta)=0.
 \end{cases}
\]

\begin{definition}[Observable evolution]
\label{p12-83-c1-s3:def:tpk-evolucion-observable}
The map
\[
 F_{\rm obs}:\mathcal Q_{\rm obs}\longrightarrow\mathcal Q_{\rm obs}
\]
is defined by
\[
 F_{\rm obs}(q)=
 (p^{\Sigma,+},p^{\Pi,+},d^{\Sigma,+},d^{\Pi,+},[\theta+1]_{108}).
\]
The category libre of the graph \(q\to F_{\rm obs}(q)\) it is denoted by
\(\mathsf P_{\rm obs}\).
\end{definition}

Each arrow of \(\mathsf P_{\rm obs}\) induces a pair of trajectories in
APP, so there exists a functor
\[
 B:\mathsf P_{\rm obs}\longrightarrow
 \mathsf P_{\rm APP}^{(2)}.
\]
If \(\tau=+1\), the second component of \(B\) is the identity; if
\(\tau=-1\), the first is the identity; if \(\tau=0\), both are the identity.

\begin{proposition}[Separation between selection, motion, and remark]
\label{p12-83-c1-s3:prop:tpk-separacion-seleccion-movimiento}
The selector \(M_{\rm ph}\), the translations
\(T_d(p)=p+\delta(d)\), and any temporal sampling have distinct domains and
codomains. They are not interchangeable operators.
\end{proposition}

\begin{proof}
\(M_{\rm ph}\) takes a phase label and returns a tritic orientation.
\(T_d\) is an automorphism of \(X_9\). A sampling acts on a
sequence or a category of arrows. Identifying two of them would violate
typing. \(F_{\rm obs}\) composes them in a declared order without erasing that
distinction.
\end{proof}

\subsection{Decomposition of the integral holonomic state into \(6\) fibered families}

Over each \(q\in\mathcal Q_{\rm obs}\) we fix the total fiber
\[
 \mathcal F_q=
 \mathsf O_q\times
 \mathsf H_q\times
 \mathsf C_q\times
 \mathsf S_q\times
 \mathsf B_q\times
 \Lambda_q.
\]
Each factor has its own type and function.

\subsubsection{Joint determination of the orientation and sheet of the integral holonomic state}

\[
 \mathsf O_q\subseteq
 \mathscr D_{\rm card}^2\times\{+1,-1\}.
\]
The two directions coincide with those stored in \(q\); the last
coordinate records the sheet orientation when the realization requires it.
The sheet label does not replace the cardinal directions.

\subsubsection{Euclidean decomposition of the state into compatible residue and quotient}

The first lift fiber is
\[
 \mathsf H^{(1)}=\mathcal D_9\times\mathbb Z,
 \qquad
 \operatorname{rec}_9(r,u)=r+9u.
\]
For each integer \(n\) there exist unique
\[
 r=\operatorname{dr}_9(n),
 \qquad
 u=\frac{n-r}{9}.
\]
A greater depth consists of a declared tower of these fibers and their
truncation maps. The term ``Hensel'' does not replace the definition of
that tower.

\subsubsection{Carry operator and joint update of residue, quotient, and memory}

\(\mathsf C_q\) is an abelian group. In the readers of decimal blocks
one takes \(\mathsf C_q=\mathbb Z\). The increment is not an independent
coordinate, but a cochain over the arrows:
\[
 \kappa(r)\in\mathsf C_{q'}
 \quad\text{for }r:q\to q'.
\]

\subsubsection{Construction of the full signature from orientation, carry, and boundary}

The signature is derived via a typed application.
\[
 \operatorname{Sig}_q:
 \mathsf H_q\times\mathsf C_q\times\Lambda_q
 \longrightarrow\mathsf S_q.
\]
In the nonadic boundary realization, individual conservation is maintained,
\[
 q_\partial,a_\partial,c_\partial\in\mathbb F_3^6,
 \qquad
 \operatorname{colw}\in\mathbb Z_{\geq0}^{6},
 \qquad
 \rho_W,r_\partial\in\mathbb F_3^3.
\]
These coordinates are margin, axial combination, cubic condition, column
weights, Witt residual class, and ternary return. They do not condense into a
single word if a subsequent proof uses any of them.

\subsubsection{Cylinder and boundary}

\[
 \mathsf B_q:=\mathsf{Cyl}_q\times\mathsf{Fr}_q.
\]
The cylinder localizes a family of prefixes; the boundary determines which
extensions are compatible. Survival will be a relation between
states endowed with both data, not a number added a posteriori.

\subsubsection{Genealogical transport record}

\(\Lambda_q\) contains the transport records whose endpoint is
\(q\). A prolongation \(r:q\to q'\) acts by concatenation:
\[
 \lambda\longmapsto r\circ\lambda.
\]
The genealogical register preserves the route; its observable endpoint does not determine it.

\subsection{Fiber transport}

For each arrow \(r:q\to q'\) maps are declared
\[
 \mathsf H(r):\mathsf H_q\to\mathsf H_{q'},
 \quad
 \mathsf C(r):\mathsf C_q\to\mathsf C_{q'},
 \quad
 \Lambda(r):\Lambda_q\to\Lambda_{q'}.
\]
They respect identities and composition, and \(\mathsf C(r)\) is a homomorphism.
The carry satisfies
\begin{equation}
 \kappa(r_2\circ r_1)
 =\kappa(r_2)+\mathsf C(r_2)\bigl(\kappa(r_1)\bigr).
\label{p12-83-c1-s3:eq:tpk-cocadena-transporte}
\end{equation}
Therefore, the transport of the determined coordinates is
\begin{align*}
 h'&=\mathsf H(r)(h),\\
 c'&=\mathsf C(r)(c)+\kappa(r),\\
 \lambda'&=\Lambda(r)(\lambda),\\
 \sigma'&=\operatorname{Sig}_{q'}(h',c',\lambda').
\end{align*}

\begin{lemma}[Associativity of the transported carry]
\label{p12-83-c1-s3:lem:tpk-asociatividad-acarreo}
The equation \cref{p12-83-c1-s3:eq:tpk-cocadena-transporte} guarantees that transporting a
state by \(r_1\) and then by \(r_2\) produces the same carry as
transporting it by \(r_2\circ r_1\).
\end{lemma}

\begin{proof}
The successive transport gives
\[
 \mathsf C(r_2)\bigl(\mathsf C(r_1)c+\kappa(r_1)\bigr)+\kappa(r_2).
\]
By functoriality of \(\mathsf C\), the first term is
\(\mathsf C(r_2r_1)c\); the sum of the other two is
\(\kappa(r_2r_1)\) by \cref{p12-83-c1-s3:eq:tpk-cocadena-transporte}.
\end{proof}

\subsection{Extension relations}

Orientation and boundary are not imposed by means of a universal function.
For each arrow \(r:q\to q'\), a relation is declared
\[
 \operatorname{Ext}_r\subseteq\mathcal F_q\times\mathcal F_{q'}.
\]
The two conditions are required.
\[
 \Delta_q\subseteq\operatorname{Ext}_{\operatorname{id}_q},
 \qquad
 \operatorname{Ext}_{r_2}\circ\operatorname{Ext}_{r_1}
 \subseteq\operatorname{Ext}_{r_2\circ r_1}.
\]
The first preserves identities; the second preserves composition. A
relation may admit several local extensions and leave only one after a
larger horizon.

\begin{definition}[The Topological Phase Kernel]
\label{p12-83-c1-s3:def:tpk-categoria-autonoma}
An object of the TPK is a tuple
\[
 x=(q,o,h,c,\sigma,C,b,\lambda)
\]
such that
\[
 q\in\mathcal Q_{\rm obs},\quad
 o\in\mathsf O_q,\quad
 h\in\mathsf H_q,\quad
 c\in\mathsf C_q,\quad
 \lambda\in\Lambda_q,
\]
\[
 \sigma=\operatorname{Sig}_q(h,c,\lambda),
 \qquad
 (C,b)\in\mathsf{Cyl}_q\times\mathsf{Fr}_q.
\]
We associate to each object its fiber component.
\[
 \mathbf f(x):=(o,h,c,\sigma,(C,b),\lambda)\in\mathcal F_q.
\]
A morphism \(x\to x'\) is a triple \((r,x,x')\) in which
\(r:q\to q'\) belongs to \(\mathsf P_{\rm obs}\), the transportable
coordinates satisfy the preceding laws, and
\((\mathbf f(x),\mathbf f(x'))\in
\operatorname{Ext}_r\). Composition is
\[
 (r_2,x',x'')\circ(r_1,x,x')=(r_2r_1,x,x''),
\]
and the identity of \(x\) is \((\operatorname{id}_q,x,x)\).
\end{definition}

\begin{theorem}[Categorical structure of the TPK]
\label{p12-83-c1-s3:thm:tpk-es-categoria}
The \cref{p12-83-c1-s3:def:tpk-categoria-autonoma} defines a category and the map
\[
 p:\mathsf{TPK}\longrightarrow\mathsf P_{\rm obs},
 \qquad p(x)=q,\qquad p(r,x,x')=r,
\]
is a functor.
\end{theorem}

\begin{proof}
Route composition is associative. The
\cref{p12-83-c1-s3:lem:tpk-asociatividad-acarreo} proves compatibility of the
carry coordinate; the declared functoriality of \(\mathsf H\),
\(\mathsf C\), and \(\Lambda\) proves compatibility of the other determined
coordinates. The stability of \(\operatorname{Ext}\) guarantees that the
composite remains admissible, and its diagonal condition provides the
identities. The map \(p\) preserves identities and composition by
construction.
\end{proof}

\subsection{Projections and typed losses}

The elementary projections are
\[
 \phi(x)=g_0(\theta),
 \qquad
 \tau(x)=\varepsilon_9(r(\theta)),
 \qquad
 \mu(x)=\mu(\theta).
\]
The projection to the local TRIT state is
\[
 \pi_{\rm TRIT}(x)
 =\bigl(\tau,\mu,o,h,c,\sigma,\lambda,\phi\bigr).
\]
It forgets the cylinder, the boundary, and part of \(q\); it does not define another container.
Two states may have the same image under \(\pi_{\rm TRIT}\) and differ
in the forgotten information.

\subsection{Criterion of route survival via conservation of the projected boundary}

Let \(\mathcal X_{\rm TPK}^{(n)}\) be the family of states at depth
\(n\). For \(N\geq n\), we define
\[
 \operatorname{Surv}_n^{(N)}=
 \left\{x_n:
 \begin{array}{l}
 \text{there exist }x_{n+1},\ldots,x_N\text{ with}\\
 (x_j,x_{j+1})\in\operatorname{Ext}_{r_j}
 \text{ for }n\leq j<N
 \end{array}\right\}.
\]

\begin{proposition}[Survival filtration]
\label{p12-83-c1-s3:prop:tpk-filtracion-supervivencia}
For \(N\geq n\),
\[
 \operatorname{Surv}_n^{(N+1)}
 \subseteq
 \operatorname{Surv}_n^{(N)}.
\]
If the scheme is defined to all depth, the states
indefinitely prolongable are
\[
 \operatorname{Surv}_n
 =\bigcap_{N\geq n}\operatorname{Surv}_n^{(N)}.
\]
\end{proposition}

\begin{proof}
A prolongation to \(N+1\) truncates to a prolongation to \(N\),
which proves the inclusion. The intersection expresses exactly
prolongability for every finite horizon.
\end{proof}

\subsection{Operator tree of the container}

\begin{figure}[htbp]
 \centering
 \begin{tikzpicture}[
  node distance=7mm and 9mm,
  every node/.style={font=\small,align=center},
  root/.style={draw=black,very thick,rounded corners,inner sep=5pt},
  op/.style={draw=black,rounded corners,inner sep=4pt},
  edge/.style={-{Latex[length=2mm]},semithick}
 ]
  \node[root] (tpk) {TPK\\$\mathcal X_{\rm TPK}^{\rm enr}$};
  \node[op,below left=of tpk] (sel) {selection\\$M_{\rm ph}$};
  \node[op,below=of tpk] (mov) {transport\\$T_d,F_{\rm obs}$};
  \node[op,below right=of tpk] (fib) {fibers\\$\mathsf H,\mathsf C,\mathsf S,
    \mathsf B,\Lambda$};
  \node[op,below=17mm of mov] (cal) {calendar\\$C_{12}\hookrightarrow
    C_{108}\twoheadrightarrow C_9$};
  \node[op,below=of cal] (ext) {extension, cylinders and survival};
  \draw[edge] (tpk) -- (sel);
  \draw[edge] (tpk) -- (mov);
  \draw[edge] (tpk) -- (fib);
  \draw[edge] (sel) |- (cal);
  \draw[edge] (mov) -- (cal);
  \draw[edge] (fib) |- (cal);
  \draw[edge] (cal) -- (ext);
 \end{tikzpicture}
 \caption{Operator containment of the Topological Phase Kernel. The selector,
 transport, readers, and fibers are not parallel containers: they are
 maps or coordinates of the same integral holonomic state.}
 \label{p12-83-c1-s3:fig:tpk-arbol-operatorio-autonomo}
\end{figure}

\subsection{Obstructions of the integral holonomic TPK state: selector typing, carry cochain, compatible extension, and genealogical preservation}

\begin{criterion}[Refutation criteria for the integral holonomic state]
The definition is refuted if:
\begin{enumerate}
 \item the selector \(M_{\rm ph}\) it usa as translation;
 \item \(\theta\), \(g_0(\theta)\), \(r(\theta)\), and
       \(\beta(\theta)\) are identified;
 \item a signature is allowed to be chosen independently of
       \((h,c,\lambda)\);
 \item the composition of carries incumple
       \cref{p12-83-c1-s3:eq:tpk-cocadena-transporte};
 \item \(\operatorname{Ext}_r\) is replaced by an unproved functional choice;
 \item a projection that has forgotten the cylinder, boundary,
       or genealogical record is called a complete state;
 \item the phase return erases any of the coordinates of fiber.
\end{enumerate}
\end{criterion}

This unit has defined the container. The next will enumerate its readers,
operators, odometer, and calendars, and will prove why the returns of
nine, twenty-seven, and one hundred eight steps are not restarts of the state.
}%


% Unidad U001. Redacción autónoma; no procede del borrador condensado.

\section{Marked segments, interstices and the nonadic cell}
\label{p12-83-c1-s4:chap:segmentos-marcas-intersticios}

The construction begins before every constant and before every choice of
physical unit. The inaugural datum is not an already constituted real number, but an
ordered segment, a family of marks, and the adjacency relation between
consecutive marks. This separation is necessary because a mark, an
interval, the length of that interval, and the numeral that publishes that
length are objects of different types.

\subsection{Finite positional chart and decimal residual projection}

\begin{definition}[Marked segment of order ten]
Let
\[
 \mathsf M_{10}:=\{m_0,m_1,\ldots,m_{10}\}
\]
a set totally ordered by
\(m_0<m_1<\cdots<m_{10}\). Its adjacency graph is the path
\[
 \mathsf P_{10}:\quad
 m_0\longrightarrow m_1\longrightarrow\cdots\longrightarrow m_{10}.
\]
We shall denote by
\[
 \mathsf I_k:=(m_k,m_{k+1}),\qquad 0\leq k\leq9,
\]
the abstract interstice between two consecutive marks. The
family of interstices is
\(\mathsf E_{10}:=\{\mathsf I_0,\ldots,\mathsf I_9\}\).
\end{definition}

The open notation \((m_k,m_{k+1})\) does not yet presuppose the real line. It
indicates only the ordered pair of endpoints and the edge joining them.
When the Archimedean chart is chosen
\[
 \chi_{10}(m_k)=k,
\]
the abstract interstice \(\mathsf I_k\) is published by the semiopen interval
\([k,k+1)\). Therefore, the index \(k\), the edge
\(\mathsf I_k\), and the interval \([k,k+1)\) will not be identified in the text:
the first is a label, the second is a combinatorial element, and the
third is its realization in a chart.

\begin{lemma}[Census of Marks and Gaps]
\label{p12-83-c1-s4:lem:censo-marcas-intersticios}
The segment \(\mathsf P_{10}\) contains eleven marks and ten interstices. More generally, a path with \(n+1\) marks has exactly \(n\) edges.
\end{lemma}

\begin{proof}
The map
\[
 \{0,\ldots,n-1\}\longrightarrow E(\mathsf P_n),
 \qquad k\longmapsto(m_k,m_{k+1}),
\]
is bijective. The particular assertion is obtained with \(n=10\).
\end{proof}

The endpoint \(m_0\) fixes the origin of the chart, and \(m_{10}\) fixes its threshold. Neither by itself constitutes an additional interstice. This elementary remark prevents three subsequent confusions: counting endpoints as cells, confusing closure with a new phase state, and treating zero as a tenth phase independent of the nonadic cycle.

\subsection{Extraction of the cell of \(9\) phases}

The decimal chart distinguishes the anchor \([0,1)\) of the periodic cell. The
nine positive marks
\[
 \mathsf D_9:=\{1,2,3,4,5,6,7,8,9\}
\]
it utilizaran as representantes publicos of
\(\mathbb Z/9\mathbb Z\), haciendo corresponder the mark \(9\) with the class
residual \(\overline 0\). We define
\[
 \lambda_9:\mathsf D_9\longrightarrow\mathbb Z/9\mathbb Z,
 \qquad \lambda_9(a)=\overline a.
\]

\begin{proposition}[Positive chart of the nonadic classes]
\label{p12-83-c1-s4:prop:carta-positiva-nonadica}
The map \(\lambda_9\) is bijective. Its inverse is the positive section
\[
 s_+(\overline a)=
 \begin{cases}
 a,&1\leq a\leq8,\\
 9,&\overline a=\overline0.
 \end{cases}
\]
In particular, the public representation \(9\) and the residual representation \(0\)
represent the same class, but belong to different charts.
\end{proposition}

\begin{proof}
The nine classes of \(\mathbb Z/9\mathbb Z\) have unique representatives in
\(\{0,1,\ldots,8\}\). Replacing the representative \(0\) by \(9\) produces
another complete transversal, so \(\lambda_9\) is bijective and \(s_+\)
is its inverse.
\end{proof}

The phase cell is not obtained by arbitrarily erasing the origin. The segment
\([0,1)\) fixes the decimal chart from which the cycle is observed; the nine
phases are the nine classes published as \(1,\ldots,9\). In the lifted
chart, the step that closes the cell is represented by \([9,10)\); in the
residual chart, that same step returns from \(\overline0\) to \(\overline1\).
Closure and change of representative are not the same operation.

\subsection{Position, height and Euclidean division}

\begin{definition}[Nonadic Coordinates]
For each integer \(n\geq1\), we define
\[
 \rho_9(n):=1+((n-1)\bmod9),
 \qquad
 q_9(n):=\left\lfloor\frac{n-1}{9}\right\rfloor.
\]
The first coordinate is called phase position; the second, turn height or memory quotient.
\end{definition}

\begin{theorem}[Unique nonadic decomposition]
\label{p12-83-c1-s4:thm:descomposicion-nonadica-unica}
Every integer \(n\geq1\) admits a unique expression
\[
 n=9q+\rho,
 \qquad q\in\mathbb Z_{\geq0},\quad \rho\in\mathsf D_9.
\]
In esa escritura \(q=q_9(n)\) and \(\rho=\rho_9(n)\).
\end{theorem}

\begin{proof}
We apply Euclidean division to \(n-1\): there exist unique
\(q\geq0\) and \(r\in\{0,\ldots,8\}\) such that
\(n-1=9q+r\). Taking \(\rho=r+1\) gives
\(n=9q+\rho\), with \(1\leq\rho\leq9\). The uniqueness of \(q,r\) implies that
of \(q,\rho\).
\end{proof}

\begin{corollary}[Phase return with memory increment]
\label{p12-83-c1-s4:cor:retorno-fase-incremento-memoria}
For every \(n\geq1\),
\[
 \rho_9(n+9)=\rho_9(n),
 \qquad
 q_9(n+9)=q_9(n)+1.
\]
Consequently, a full rotation preserves the visible phase and changes the raised state.
\end{corollary}

\begin{proof}
If \(n=9q+\rho\), then \(n+9=9(q+1)+\rho\); uniqueness of
\cref{p12-83-c1-s4:thm:descomposicion-nonadica-unica} gives the two identities.
\end{proof}

The one-step evolution is written unambiguously as
\[
 (q,\rho)\longmapsto
 \begin{cases}
 (q,\rho+1),&1\leq\rho<9,\\
 (q+1,1),&\rho=9.
 \end{cases}
\]
The second line is not a restart. It recovers the initial phase label and
increases the quotient that records how many turns have been performed.

\subsection{Circular projection periodicity and injectivity of the nonadic helical lifting}

Let
\[
 \vartheta_9(\rho)=\frac{2\pi(\rho-1)}9.
\]
The phase projection onto the unit circle is
\[
 \mathcal R_9(n)=
 \bigl(\cos\vartheta_9(\rho_9(n)),
       \sin\vartheta_9(\rho_9(n))\bigr).
\]
To simultaneously preserve position and height we introduce the helical immersion.
\[
 \mathcal H_9(n)=
 \left(
  \cos\vartheta_9(\rho_9(n)),
  \sin\vartheta_9(\rho_9(n)),
  q_9(n)+\frac{\rho_9(n)-1}{9}
 \right).
\]

\begin{proposition}[The circular projection does not determine the state]
\label{p12-83-c1-s4:prop:rueda-no-determina-estado}
The map \(\mathcal R_9\) is periodic with period nine. The map
\(\mathcal H_9\) is injective. In particular,
\[
 \mathcal R_9(n+9)=\mathcal R_9(n),
 \qquad
 \mathcal H_9(n+9)\neq\mathcal H_9(n).
\]
\end{proposition}

\begin{proof}
The first equality follows from
\cref{p12-83-c1-s4:cor:retorno-fase-incremento-memoria}. The third coordinate of
\(\mathcal H_9(n+9)\) exceeds that of \(\mathcal H_9(n)\) by one unit, hence
both points are distinct. Finally, \(n\) is recovered from the third coordinate
by multiplying by nine and adding one to the phase part;
therefore, \(\mathcal H_9\) is injective.
\end{proof}

This proposition is the minimal model of the principle that will run through the entire
work:
\[
 \boxed{\text{phase return}\neq\text{state return}.}
\]
The nonadic connection of the Topological Phase Kernel will preserve more fields than
the elementary helix ---sheet, orientation, carry, boundary, terminal, and
route memory---, but it will not alter this inaugural distinction.

\subsection{Correspondence between Discrete Cardinality and Projective Measure}

\begin{definition}[Counting]
The count of a finite set \(X\) is its cardinality \(|X|\). It requires
neither an orientation, a scale, nor an additional reader.
\end{definition}

\begin{definition}[Measured data]
A discrete measurement datum is a triple
\[
 (x,\mathcal L,\mathcal C),
\]
where \(x\) is a state, \(\mathcal L\) is a typed reader, and
\(\mathcal C\) is a publication chart compatible with the codomain of
\(\mathcal L\). Its visible output is
\[
 \mathcal C(\mathcal L(x)).
\]
\end{definition}

The same set may have a single cardinality and admit multiple measures,
because the local cochain, the orientation of integration, or the
publication chart may vary. Conversely, two distinct states may publish the
same coordinate. The work will therefore always preserve two levels:
\[
 \text{state with genealogy}
 \quad\longrightarrow\quad
 \text{published coordinate}.
\]

\begin{criterion}[Refutation criteria for the inaugural unit]
The construction is refuted if any of the following
facts occurs:
\begin{enumerate}
 \item ten independent phases are counted in the residue cell;
 \item the mark \(9\) is identified with the integer \(0\) without declaring the
 change of chart;
 \item the passage \(9\to1\) forces \(q\mapsto0\);
 \item two distinct integers have the same pair \((q_9,\rho_9)\);
 \item a physical unit is attributed to the cell before defining the
 dimensional transducer.
\end{enumerate}
\end{criterion}

\subsection{Inaugural domain of the nonadic cell and causal exclusion of APP, TRIT, TPK, and analytic constants}

The unit constructed so far contains only:
\[
 \mathsf P_{10},\quad
 \mathsf D_9,\quad
 \lambda_9,\quad
 (q_9,\rho_9),\quad
 \mathcal R_9,\quad
 \mathcal H_9.
\]
It does not yet contain APP, TRIT, the Topological Phase Kernel, or an analytic
constant. The following unit will take two copies of the nonadic cell,
construct its gnomonic completion, and prove how the bidimensional support appears
over which the multiplicative matrix of APP and its internal additive law act.


% Unidad U003B. Esqueleto ciclotómico de APP y completación binomial.
% Redacción autónoma a partir de los objetos, no de la arquitectura del PDF abreviado.

\section{Cyclotomic skeleton of APP and binomial completion of the unit}
\label{p12-83-c1-s5:chap:app-esqueleto-completacion}

The defect between the APP matrix and its internal accumulator is not an isolated
number. Its law in
Cartesian length, its decomposition into orbits of order three, and the
completion of the nonadic cell produce three different objects that must
remain typed:
\[
 \text{scalar defect}\quad D_d,
 \qquad
 \text{cyclotomic module}\quad W_{\rm APP},
 \qquad
 \text{completion}\quad\overline{\mathcal C}_9.
\]
The first counts; the second preserves orientation and rank; the third adds
boundary strata. Equalities between their cardinals do not authorize identifying
their elements.

\subsection{Defect law at fixed modulus and variable length}

For each integer \(d\geq1\), let
\[
 B_d:=\{1,\ldots,9\}^d.
\]
The nonary digital root is represented by
\[
 \operatorname{dr}_9(n)=
 \begin{cases}
 n\bmod9,&n\not\equiv0\pmod9,\\
 9,&n\equiv0\pmod9.
 \end{cases}
\]
We define the additive accumulation and the multiplicative matrix evaluation on the same domain:
\[
 A_d(x_1,\ldots,x_d)
 =\operatorname{dr}_9\!\left(\sum_{j=1}^d x_j\right),
\]
\[
 M_d(x_1,\ldots,x_d)
 =\operatorname{dr}_9\!\left(\prod_{j=1}^d x_j\right).
\]
Their visible masses and their defect are
\[
 S_\Sigma(d):=\sum_{x\in B_d}A_d(x),
 \qquad
 S_\Pi(d):=\sum_{x\in B_d}M_d(x),
\]
\[
 D_d:=S_\Pi(d)-S_\Sigma(d).
\]
Here \(d\) counts Cartesian positions. The module \(9\) does not vary, and no
physical dimension has yet been introduced.

\begin{lemma}[Mass of the additive accumulator]
\label{p12-83-c1-s5:lem:masa-hoja-aditiva-longitud}
For every \(d\geq1\),
\[
 S_\Sigma(d)=45\,9^{d-1}.
\]
\end{lemma}

\begin{proof}
Once the first \(d-1\) coordinates and a published residue
\(r\in\{1,\ldots,9\}\) have been fixed, there exists a unique last coordinate that makes the
sum have residue \(r\). Each visible value thus has \(9^{d-1}\)
preimages. Summing the weights \(1+\cdots+9=45\) gives the formula.
\end{proof}

\begin{lemma}[Census of the multiplicative matrix]
\label{p12-83-c1-s5:lem:censo-hoja-multiplicativa-longitud}
For each of the six unit residues there are \(6^{d-1}\) preimages for each
fixed product value. The residues \(3\) and \(6\) have
\(d6^{d-1}\) preimages each. The zero residue, published as
\(9\), has
\[
 9^d-6^d-2d6^{d-1}
\]
preimages. Consequently,
\[
 S_\Pi(d)=9^{d+1}-9(d+3)6^{d-1}.
\]
\end{lemma}

\begin{proof}
The unit group of \(\mathbb Z/9\mathbb Z\) has six elements. Upon
fixing \(d-1\) units, the last is determined by the prescribed
unit product, giving \(6^{d-1}\) solutions for each of the
six classes.

To obtain a product of ternary valuation exactly one, one chooses which
of the \(d\) coordinates belongs to one of the two classes \(3,6\), while
the others are units. Once the product \(3\) or \(6\) is fixed, the last unit
is determined; hence \(d6^{d-1}\) for each class. All remaining tuples
have product zero modulo \(9\). Subtracting the two preceding
counts from \(9^d\) gives their multiplicity. The sum weighted by the
visible representatives yields the final identity.
\end{proof}

\begin{theorem}[Internal dimensional law of APP]
\label{p12-83-c1-s5:thm:ley-dimensional-app-autonoma}
For every \(d\geq1\),
\[
 \boxed{D_d=4\,9^d-9(d+3)6^{d-1}.}
\]
In particular,
\[
 D_1=0,
 \qquad
 D_2=54,
\]
and
\[
 \sum_{d\geq1}D_dz^d
 =\frac{54z^2(1-3z)}{(1-9z)(1-6z)^2}.
\]
Moreover,
\[
 \frac{D_d}{9^d}
 =4-(d+3)\left(\frac23\right)^{d-1}
 \longrightarrow4.
\]
\end{theorem}

\begin{proof}
The formulas from the two lemmas are subtracted. For the generating function, one uses
\[
 \sum_{d\geq1}9^dz^d=\frac{9z}{1-9z},
 \qquad
 \sum_{d\geq1}(d+3)6^{d-1}z^d
 =\frac{z(4-18z)}{(1-6z)^2},
\]
and the common numerator is reduced. The normalized formula follows by dividing
by \(9^d\); the exponential term multiplied by \(d+3\) converges to
zero.
\end{proof}

The equality \(D_2=54\) reproduce the census \(459-405\), but the formula
before and the law of modulo variable
\(\Delta_{3^m}^{\rm pl}=m3^{2m-1}\) are enunciados distintos: the first
varia \(d\) manteniendo the modulo \(9\); the second varia the modulo.

\subsection{Primary cyclotomic operator}

Let
\[
 G_{A_2}=\begin{pmatrix}2&-1\\-1&2\end{pmatrix},
 \qquad
 C=\begin{pmatrix}0&-1\\1&-1\end{pmatrix}.
\]

\begin{proposition}[Coxeter of type \(A_2\)]
\label{p12-83-c1-s5:prop:coxeter-a2-autonomo}
They are verified.
\[
 C^3=I_2,
 \qquad
 C^{\mathsf T}G_{A_2}C=G_{A_2},
 \qquad
 \det(XI_2-C)=X^2+X+1.
\]
Therefore \(C\) preserves the form \(A_2\), has order exactly three, and
has no nonzero fixed vector.
\end{proposition}

\begin{proof}
The multiplication gives
\[
 C^2=\begin{pmatrix}-1&1\\-1&0\end{pmatrix},
 \qquad C^3=I_2.
\]
Direct substitution into \(C^{\mathsf T}G_{A_2}C\) yields
\(G_{A_2}\). Finally, \(1\) is not a root of \(X^2+X+1\), hence
\(\ker(C-I_2)=0\).
\end{proof}

\subsection{Decomposition of APP into \(12\) fibers of type \(A_2\)}

Consider the action
\[
 a:\mathbb Z/9\mathbb Z\longrightarrow\mathbb Z/9\mathbb Z,
 \qquad a(r)=4r.
\]
Since \(4^3\equiv1\pmod9\), its restriction to the units has two free
orbits:
\[
 \mathcal O_1=(1,4,7),
 \qquad
 \mathcal O_2=(2,8,5).
\]
For an orbit \(\mathcal O=(r,4r,7r)\), we define the augmentation module
\[
 A(\mathcal O):=
 \ker\!\left(\mathbb Z[\mathcal O]
 \xrightarrow{\,\sum\,}\mathbb Z\right).
\]
At the base
\[
 b_1=e_r-e_{4r},
 \qquad b_2=e_{4r}-e_{7r},
\]
the restricted Euclidean form has Gram matrix \(G_{A_2}\), and the cycle
\(r\mapsto4r\mapsto7r\mapsto r\) acts by \(C\).

For \(x=(x_1,x_2,x_3)\in(\mathbb Z/9\mathbb Z)^3\), the three unordered pairs are
\[
 \{1,2\},\quad\{2,3\},\quad\{3,1\}.
\]
Over each pair, two typed laws are evaluated:
\[
 \ell_{ij}^{\Sigma}(x)=x_i+x_j,
 \qquad
 \ell_{ij}^{\Pi}(x)=x_ix_j.
\]
Since
\[
 \ell_{ij}^{\Sigma}(ax)=a\ell_{ij}^{\Sigma}(x),
 \qquad
 \ell_{ij}^{\Pi}(ax)=a^{-1}\ell_{ij}^{\Pi}(x),
\]
the sum carries the orientation \(C\) and the product the inverse orientation
\(C^{-1}\).

\begin{definition}[Cyclotomic module of APP]
\label{p12-83-c1-s5:def:modulo-ciclotomico-app-autonomo}
It is defined
\[
 \boxed{
 W_{\rm APP}:=
 \bigoplus_{\{i,j\}}
 \bigoplus_{\mu\in\{\Sigma,\Pi\}}
 \bigoplus_{\mathcal O\in\{\mathcal O_1,\mathcal O_2\}}
 A(\mathcal O).}
\]
There are \(3\cdot2\cdot2=12\) summands, each of rank two; therefore
\[
 W_{\rm APP}\cong A_2^{12},
 \qquad \operatorname{rank}W_{\rm APP}=24.
\]
The par--law--orbit labels are part of the object.
\end{definition}

For \(u\in\mathbb Z/9\mathbb Z\), we write
\[
 \beta_{\mathcal O}(u)=
 \begin{cases}
 (1,0),&u=r,\\
 (0,1),&u=4r,\\
 (-1,-1),&u=7r,\\
 (0,0),&u\notin\mathcal O,
 \end{cases}
\]
in the basis \((b_1,b_2)\). The state map is
\[
 \Psi(x)_{ij,\mu,\mathcal O}
 :=\beta_{\mathcal O}(\ell_{ij}^{\mu}(x)).
\]

\begin{theorem}[Equivariance and rank of the state skeleton]
\label{p12-83-c1-s5:thm:esqueleto-estatal-app-autonomo}
Let \(\rho(a)\) be the action using \(C\) on the six additive
fibers and \(C^{-1}\) on the six multiplicative fibers. Then
\[
 \boxed{\Psi(ax)=\rho(a)\Psi(x)}
\]
for the \(729\) states of \((\mathbb Z/9\mathbb Z)^3\). The rational submodule
generated by the images has rank \(24\).
\end{theorem}

\begin{proof}
Under the additive law, multiplying all coordinates by \(4\) advances one
position in each orbit; under the multiplicative law, it multiplies the value by
\(4^2\equiv7\equiv4^{-1}\pmod9\) and therefore reverses the cycle. This proves
componentwise equivariance.

In order for the rank not to remain hidden in a computational statement, we fix the coordinate order.
\[
 (12,\Sigma,\mathcal O_1),(12,\Sigma,\mathcal O_2),
 (12,\Pi,\mathcal O_1),(12,\Pi,\mathcal O_2),
\]
followed by the same order for the pairs \(23\) and \(31\), and, within each
fiber, the two coordinates \((b_1,b_2)\). Consider the states
\[
\begin{gathered}
001,002,004,005,010,011,012,013,014,015,016,020,\\
022,023,040,050,101,102,104,105,110,120,140,150.
\end{gathered}
\]
The integer matrix \(M_\Psi\) whose rows are its twenty-four images has
entries in ({-1,0,1}). Bareiss's fraction-free elimination yields the following
successive leading principal minors
\[
1,1,1,-1,-1,-1,-1,1,-1,1,1,1,1,-2,-1,1,1,1,2,1,4,-6,12,9.
\]
In particular,
\[
 \det M_\Psi=9\neq0.
\]
These twenty-four images are linearly independent. Since the
codomain already has rank \(24\), they generate \(W_{\rm APP}\otimes\mathbb Q\).
\end{proof}

The proof exhibits a finite reproducible certificate; it does not deduce full rank
from the mere surjectivity of each component projection.

\subsection{Localization of the aggregate \texorpdfstring{\(54\)}{54}}

Let
\[
 \mathcal M:=\mathbb Z[\mathbb Z/9\mathbb Z]
\]
and, for \(\mu\in\{\Sigma,\Pi\}\),
\[
 \nu_\mu:=\sum_{x,y\in\mathbb Z/9\mathbb Z}
 e_{\ell^\mu(x,y)}.
\]
The census of the two tables gives
\[
 \delta_{\Pi\Sigma}:=\nu_\Pi-\nu_\Sigma
 =12e_0+3e_3+3e_6
 -3\sum_{u\in(\mathbb Z/9\mathbb Z)^\times}e_u.
\]
Its coefficients are constant on the orbits of \(a\); therefore
\[
 (I-a)\delta_{\Pi\Sigma}=0.
\]
We define the increase projection
\[
 T:\mathcal M\longrightarrow
 A(\mathcal O_1)\oplus A(\mathcal O_2),
\]
\[
 T(\nu)=
 \left(((I-a)\nu)|_{\mathcal O_1},
       ((I-a)\nu)|_{\mathcal O_2}\right).
\]
Then
\[
 T(\delta_{\Pi\Sigma})=0.
\]
However, the visible reader
\[
 w(e_0)=9,
 \qquad w(e_r)=r\quad(1\leq r\leq8)
\]
satisfies
\[
 w(\delta_{\Pi\Sigma})=54.
\]

\begin{corollary}[Isotopic Obstruction]
\label{p12-83-c1-s5:cor:obstruccion-isotipica-app-autonoma}
The scalar \(54\) belongs to the \(C_3\)-invariant sector and cannot
be recovered by means of a linear functional that factors exclusively through
the augmentation modules. In particular,
\[
 \sum_x\Psi(x)=0
\]
and, for every linear functional \(L\) on \(W_{\rm APP}\otimes\mathbb Q\),
\[
 L\!\left(\sum_x\Psi(x)\right)=0\neq54.
\]
\end{corollary}

This obstruction prevents presenting \(54\) as a retrospective linear sum
of the twelve fibers. The subsequent APP--Fock coupling preserves the
orientation and uses an additional object; it does not contradict the corollary.

\subsection{Cubic completion of the nonad cell}

Let \(E_9\) be the underlying set of \(\mathbb Z/9\mathbb Z\), and let
\(\star\notin E_9\) be a boundary state. We define
\[
 \mathcal C_9:=E_9^3,
 \qquad
 \overline{\mathcal C}_9:=(E_9\sqcup\{\star\})^3.
\]
The symbol \(\star\) is not the zero class: it indicates that a coordinate has
reached the closure stratum.

\begin{theorem}[Binomial stratification of the unit cube]
\label{p12-83-c1-s5:thm:emrh-binomial-autonoma}
The stratum with exactly \(r\) coordinates equal to \(\star\) contains
\[
 \binom3r9^{3-r}
\]
elements. Therefore,
\[
 \boxed{1000=729+243+27+1=729+271.}
\]
Upon removing only the apex \((\star,\star,\star)\), there remains
\[
 999=729+243+27=729+270.
\]
\end{theorem}

\begin{proof}
The \(r\) boundary coordinates are chosen in \(\binom3r\) ways; each
remaining coordinate admits nine values. The four strata are disjoint
and exhaust the Cartesian product. The sum is the expansion of
\((9+1)^3\).
\end{proof}

\begin{figure}[htbp]
\centering
\resizebox{0.96\linewidth}{!}{\begin{tikzpicture}[font=\small]
  \begin{scope}[x={(.95cm,0cm)},y={(.45cm,.33cm)},z={(0cm,.95cm)}]
  \coordinate (O) at (0,0,0);
  \coordinate (X) at (3.4,0,0);
  \coordinate (Y) at (0,3.4,0);
  \coordinate (XY) at (3.4,3.4,0);
  \coordinate (Z) at (0,0,3.4);
  \coordinate (XZ) at (3.4,0,3.4);
  \coordinate (YZ) at (0,3.4,3.4);
  \coordinate (XYZ) at (3.4,3.4,3.4);
  \draw[hmtgray,dashed,semithick] (O)--(Y)--(XY);
  \draw[hmtgray,dashed,semithick] (Y)--(YZ);
  \path[fill=hmtlightblue,fill opacity=.78]
    (O)--(X)--(XZ)--(Z)--cycle;
  \path[fill=hmtlightgold,fill opacity=.62]
    (X)--(XY)--(XYZ)--(XZ)--cycle;
  \path[fill=hmtlightgreen,fill opacity=.72]
    (Z)--(XZ)--(XYZ)--(YZ)--cycle;
  \draw[hmtblue,thick] (O)--(X)--(XZ)--(Z)--cycle;
  \draw[hmtgold,thick] (X)--(XY)--(XYZ)--(XZ);
  \draw[hmtgreen,thick] (Z)--(YZ)--(XYZ);
  \node[font=\bfseries\Large,hmtgreen,fill=white,fill opacity=.78,
    text opacity=1,inner sep=2pt] at (1.7,.55,1.72) {\(9^3=729\)};
  \node[font=\scriptsize,hmtgray,fill=white,fill opacity=.78,
    text opacity=1,inner sep=1.5pt] at (1.7,.55,1.22)
    {interior nonadic};
  \draw[-{Latex[length=2mm]},hmtblue,semithick]
    (O) -- (1.18,0,0)
    node[pos=.78,below=5pt,font=\scriptsize] {\(x=(x_0,x_1)\)};
  \draw[-{Latex[length=2mm]},hmtgreen,semithick]
    (O) -- (0,1.18,0)
    node[pos=1,right=5pt,font=\scriptsize] {\(y=(y_0,y_1)\)};
  \draw[-{Latex[length=2mm]},hmtblue,semithick]
    (O) -- (0,0,1.18)
    node[pos=.78,left=2pt,font=\scriptsize] {\(z=(z_0,z_1)\)};
  \end{scope}
  \node[align=center,text width=5.2cm,font=\scriptsize,hmtgray,
    anchor=north] at (2.55,-.72)
    {paired ternary coordinates in three independent directions};
  \draw[decorate,decoration={brace,amplitude=5pt},hmtgold,thick]
    (5.20,.65)--(5.20,4.80);
  \node[anchor=west,align=left,text width=1.25cm,font=\scriptsize,hmtgold]
    at (5.40,3.82) {of side \(9\) to \(10\)\\crown \(271\)};
  \begin{scope}[xshift=7cm,yshift=.45cm]
    \node[anchor=west,font=\bfseries\large,hmtblue] at (0,4.05)
      {Cubic completion};
    \draw[fill=hmtlightgreen,draw=hmtgreen] (0,3.15) rectangle (6.1,3.75);
    \node at (3.05,3.45) {internal: \(729\)};
    \draw[fill=hmtlightblue,draw=hmtblue] (0,2.35) rectangle (2.9,2.95);
    \node[align=center,text width=2.75cm,font=\scriptsize] at (1.45,2.65)
      {caras nuevas\\\(243=3\cdot9^2\)};
    \draw[fill=hmtlightgold,draw=hmtgold] (2.9,2.35) rectangle (4.9,2.95);
    \node[align=center,text width=1.8cm,font=\scriptsize] at (3.9,2.65)
      {edges\\\(27=3\cdot9\)};
    \draw[fill=hmtlightred,draw=hmtred] (4.9,2.35) rectangle (6.1,2.95);
    \node[align=center,font=\scriptsize] at (5.5,2.65)
      {vertex\\\(1\)};
    \node[font=\bfseries\Large,hmtgold] at (3.05,1.65)
      {\(1000=729+243+27+1=729+271\)};
    \node[align=center,hmtgray] at (3.05,.65)
      {boundary internal of the cubo of lado \(9\): \(386\) puntos\\
       corona of the completion \(9^3\subset10^3\): \(271\) states};
  \end{scope}
\end{tikzpicture}
 }
\caption{Cubic completion of the nonadic cell. The interior of
cardinality \(729\) and the boundary strata of cardinalities \(243\), \(27\), and
\(1\) are disjoint parts of a single object.}
\label{p12-83-c1-s5:fig:cubo-emrh-autonomo}
\end{figure}

This stratification receives in HMT the name of extreme and mean
holographic ratio. Mathematically it is the typed binomial completion of the nonadic cell;
the name does not replace its definition.

For every \(d\geq1\), let
\[
 \overline{\mathcal C}^{(d)}_9=(E_9\sqcup\{\star\})^d.
\]
Asignemos peso \(1/10\) to each symbol of a coordinate and tomemos the measure
product \(\mu_d\).

\begin{theorem}[Projective compatibility of the measure]
\label{p12-83-c1-s5:thm:medida-proyectiva-completacion}
The projection that forgets the last coordinate satisfies
\[
 (\pi_d)_*\mu_d=\mu_{d-1}.
\]
Moreover,
\[
 10^d=\sum_{r=0}^{d}\binom dr9^{d-r},
 \qquad
 \mu_d(\overline{\mathcal C}^{(d)}_9)=1.
\]
The number of states increases with \(d\), while the normalized total mass
is conserved.
\end{theorem}

\begin{proof}
For every cylinder \(A\times(E_9\sqcup\{\star\})\),
\[
 \mu_d(A\times(E_9\sqcup\{\star\}))
 =\mu_{d-1}(A)\sum_{u\in E_9\sqcup\{\star\}}\frac1{10}
 =\mu_{d-1}(A).
\]
The cilindros generate the algebra finite of subconjuntos, what that determines the
measure image. The identity cardinal again becomes the binomio of Newton.
\end{proof}

\subsection{Morphism from the central block of APP to its Euclidean realization}

The already constructed central block is
\[
 B_c=\begin{pmatrix}7&2\\2&7\end{pmatrix}.
\]
Row reading and diagonal reading give
\[
 72+27=77+22=99,
\]
\[
 72-27=45=\det B_c,
 \qquad
 77-22=55=54+1.
\]
The relation with the completion is not obtained by concatenating symbols, but by the Euclidean division.
\[
 729=10\cdot72+9,
 \qquad
 271=10\cdot27+1.
\]

\begin{theorem}[Uniqueness of the Decimal Lift]
\label{p12-83-c1-s5:thm:levantamiento-decimal-emrh-autonomo}
For \(L_d(n)=10n+d\), with \(d\in\{0,\ldots,9\}\), there exists a single pair
\((d,d')\) such that
\[
 L_d(72)=729,
 \qquad
 L_d(72)+L_{d'}(27)=1000,
 \qquad
 d+d'=10.
\]
It is the pair
\[
 (d,d')=(9,1).
\]
\end{theorem}

\begin{proof}
The first equality is \(720+d=729\), hence \(d=9\). The second becomes
\(729+270+d'=1000\), whence \(d'=1\). The third is verified, and the
uniqueness of quotient and remainder rules out any other solution.
\end{proof}

Thus the causal diagram is obtained.
\[
 B_c\longrightarrow(72,27)
 \longleftarrow((72,9),(27,1))
 \longleftarrow(729,271).
\]
The block and its reading belong to APP; the dividends and remainders belong
to the completion. The diagram shares quotients; it does not invert
causality.

\subsection{Incidence of the cube and separation with respect to sextets and octets}

Let \(M_d\) be the incidence matrix of the \(2d\) codimension-one faces
of a hypercube on its \(2^d\) vertices. Each face contains
\(2^{d-1}\) vertices; two opposite faces are disjoint, and two
transversal faces share \(2^{d-2}\) vertices.

\begin{theorem}[Hypercube incidence spectrum]
\label{p12-83-c1-s5:thm:espectro-incidencia-cubo-autonomo}
\[
 \operatorname{Spec}(M_dM_d^{\mathsf T})
 =\left\{
 d2^{d-1}\,^{(1)},
 2^{d-1}\,^{(d)},
 0^{(d-1)}
 \right\}.
\]
Para \(d=3\),
\[
 \operatorname{Spec}(M_3M_3^{\mathsf T})
 =\{12,4,4,4,0,0\}.
\]
\end{theorem}

\begin{proof}
The row space decomposes into: the constant line; the \(d\)
differences of each pair of opposite faces; and the (d-1) combinations of
pair sums whose total sum is zero. The intersections described above
cause the Gram matrix to act by \(d2^{d-1}\), \(2^{d-1}\), and \(0\),
respectively.
\end{proof}

The three-dimensional cube has six faces, eight vertices, and twelve edges. Those
cardinals anticipate an incidence pattern (6/8), but a face is not
a hexad and the set of vertices is not an octad. The exceptional
identification requires a subsequent transporter that preserves supports and
gauge.

\subsection{Resonance of the incidence operator in rank \(6\)}

Define the strata that preserve exactly \(k\) nonadic
coordinates in rank \(d\):
\[
 \mathcal I_k(d):=\binom dk9^k.
\]
The defects of APP are
\[
 D_{\rm vis}=54,
 \qquad
 D_{\rm tot}=1215,
 \qquad
 D_{\rm coc}=129,
 \qquad
 1215=54+9\cdot129.
\]

\begin{theorem}[Uniqueness of the Binomial Resonance]
\label{p12-83-c1-s5:thm:resonancia-binomial-autonoma}
The unique positive integer \(d\) that simultaneously satisfies
\[
 \mathcal I_1(d)=54,
 \qquad
 \mathcal I_2(d)=1215
\]
is \(d=6\). In that range,
\[
 \frac{\mathcal I_2(6)-\mathcal I_1(6)}9=129.
\]
\end{theorem}

\begin{proof}
The first equality is \(9d=54\), which forces \(d=6\). Then
\[
 \mathcal I_2(6)=81\binom62=81\cdot15=1215,
\]
and
\[
 \frac{1215-54}{9}=129.
\]
There remains no other possible rank because the first equality already determines \(d\).
\end{proof}

\subsection{Obstructions of the cyclotomic skeleton of APP: fibers \(A_2^{12}\), completion \(729+271\), rank resonance \(6\), and cubic incidence}

\begin{criterion}[Typed falsification of the unit]
The construction is refuted if any of the following facts occurs:
\begin{enumerate}
 \item the law of \(D_d\) not reproduce \(D_1=0\) and \(D_2=54\);
 \item variation of length is fused with variation of modulus;
 \item \(C\) ceases to preserve \(G_{A_2}\) or acquires a fixed vector;
 \item the labels pair--sheet--orbit not produce twelve fibers;
 \item the smallest certificate of \(\Psi\) has zero determinant;
 \item a functional that factors through the augmentation module recovers \(54\);
 \item the symbol of boundary \(\star\) it identifica with the zero residue;
 \item the four strata do not sum to \(1000\);
 \item a cubic face is identified with a hexad without an incidence
       transporter.
\end{enumerate}
\end{criterion}

The unit has constructed, without recognizing them retrospectively, the module
\(A_2^{12}\), the completion \(729+271\), the projective measure, and the
rank-six resonance. These structures will feed, in distinct ways,
the generation of constants and the exceptional branches.


% Unidad U006F. Inventario operatorio completo del Topological Phase Kernel.

\long\def\HMTDeferredTPKDevelopment{%
\section{Operator algebra of the Topological Phase Kernel and causal update dynamics of the integral holonomic state}
\label{p12-83-c1-s6:chap:inventario-operatorio-completo-tpk}

The Topological Phase Kernel is neither a reduction modulo three nor a clock with
one hundred and eight positions. It is a category of integral holonomic states equipped
with selection, transport, reading, memory, periodicity,
prolongation, and publication operators. To prevent an abbreviation from replacing the
object, each operator is recorded by its type and by the information it
preserves.

\subsection{Typed schema of the 34 canonical TPK constructions: domain, codomain, action, invariants, and memory}

\begin{definition}[Operator specification]
\label{p12-83-c1-s6:def:u006f-ficha-operatoria}
A canonical realization of the TPK is a sextuple
\begin{equation}
 \mathcal O=
 (D_{\mathcal O},C_{\mathcal O},f_{\mathcal O},
 I_{\mathcal O},L_{\mathcal O},M_{\mathcal O}),
 \label{p12-83-c1-s6:eq:u006f-ficha}
\end{equation}
where:
\begin{enumerate}
 \item \(f_{\mathcal O}:D_{\mathcal O}\to C_{\mathcal O}\) is its rule;
 \item \(I_{\mathcal O}\) is the family of conserved relations;
 \item \(L_{\mathcal O}\) declares the information published or lost;
 \item \(M_{\mathcal O}\) enumerates the memory remaining in the integral
 holonomic state.
\end{enumerate}
\end{definition}

Two realizations represent the same operator only when there exists a
typed equivalence that intertwines rules and invariants. Sharing a period,
cardinality, or visible output is not sufficient.

\subsection{Classification of the integral holonomic state into \(8\) invariant structural classes}

The inventory is partitioned into:
\begin{enumerate}
 \item[\textbf{C1.}] supports and states;
 \item[\textbf{C2.}] internal selection;
 \item[\textbf{C3.}] transporte;
 \item[\textbf{C4.}] readers and quotients;
 \item[\textbf{C5.}] memory and boundary;
 \item[\textbf{C6.}] periodicity and return;
 \item[\textbf{C7.}] prolongation;
 \item[\textbf{C8.}] salidas.
\end{enumerate}
The class is determined by domain, codomain, and function, not by the name of a number associated with it.

\Needspace{.92\textheight}
\subsection{Functional decomposition into \(14\) groupings compatible with the
APP–TRIT–TPK operators}

\begingroup
\small
\renewcommand{\arraystretch}{1.12}
\begin{longtable}{@{}>{\raggedright\arraybackslash}p{.20\textwidth}>{\raggedright\arraybackslash}p{.29\textwidth}>{\raggedright\arraybackslash}p{.43\textwidth}@{}}
\toprule
Grouping&Operators or spaces&Function and invariants\\
\midrule
\endfirsthead
\toprule
Grouping&Operators or spaces&Function and invariants\\
\midrule
\endhead
State and route memory&
\(\mathcal X_{\rm TPK}^{\rm enr}\), \(\mathcal F_q\)&
Preserve position, sheet, orientation, residue, quotient, carry, signature,
boundary, cylinder, prefix, provenance and memory.\\

Internal selection&
\(M_{\rm ph}:\{1,\ldots,9\}\to\{-1,0,+1\}\)&
Activate the additive, multiplicative evaluation or the threshold step; does not move
the state.\\

Transporte espacial&
\(T_N,T_E,T_S,T_O,F_{\rm loc},F_{\rm obs}\)&
Perform cardinal translations, sheet feedback, and timeline of two cursors.\\

Lectores modulares&
\(\rho_9,q_9,\rho_3^{\rm or},c_3,q_{1000},r_{1000}\)&
Publish phase, orientation and decimal block preserving the lift quotients.\\

Block and memory&
\(\mathcal K_{27}=F_{\rm obs}^{27}\), \(\mathcal N_{27}\)&
The former composes twenty-seven transports; the latter filters memory
events. An equal index does not imply an equal operator.\\

Twelve-phase state&
\(\mathcal R_{12},C_{12},\Theta_{108},\Gamma_{108}\)&
Distinguish enriched record, sector, cyclic law, and ordered support.\\

Periodicity and depths&
\(C_{12}\hookrightarrow C_{108}\twoheadrightarrow C_9\),
\(w_{108},w_{1080}\)&
Preserves the non-splittable extent, words of different lengths and
the vertical memory.\\

Channels \(90/120\)&
\(\mathcal F_6,\mathcal F_8,\mathcal I_{6\to8},
S_{90},S_{120}\)&
Separate finite incidences, analytic tails, and their normalizations.\\

Bidirectionality&
\(P_\pm,\operatorname{rev},\operatorname{Ext}_\pm,
N_\pm,\beta,C_M\)&
Perform exchange, route inversion, future and past prolongation, valence, balance, and conjugation.\\

Emission and coinduction&
\(L_0,L_1,R_{36},G_9,\varprojlim\operatorname{Cyl}_m\)&
Lift the seed, open the boundary, and organize finite and compatible extensions.\\

Numeric Readers&
\(\mathscr C_\pi,\mathscr P_e,F_{\rm av}\)&
They construct rational forks, factorial propagation, and Fibonacci self-scaling.\\

Exceptional Projection&
\(\Gamma_{A_W}\), weight, projective support&
Transport the same emission to ternary incidence without selecting Archimedean numbers.\\

Atlas and realization&
route signature, atlas \(81/56/13/324\), \(C_M\)&
Classify stories and characters before applying a dimensional chart.\\

Finite validation&
censuses, structural necessity tests, integrity identities and
independence controls&
Checks coverage at the executed depth and separates construction,
independent verification, and comparison data.\\
\bottomrule
\end{longtable}
\endgroup

The eight classes describe mathematical nature; the fourteen groupings describe
execution function. They are two gradings of the same inventory.

\Needspace{.92\textheight}
\subsection{Domain and typing of the \(34\) entries of the update operator}

\begin{equation}
 \mathfrak O_{\rm TPK}:=(\mathfrak o_1,\ldots,\mathfrak o_{34}).
 \label{p12-83-c1-s6:eq:u006f-registro-operadores}
\end{equation}

\begingroup
\scriptsize
\renewcommand{\arraystretch}{1.12}
\begin{longtable}{@{}r p{.19\textwidth}p{.28\textwidth}p{.42\textwidth}@{}}
\toprule
\#&Symbol&Domain and codomain&Rule or function\\
\midrule
\endfirsthead
\toprule
\#&Symbol&Domain and codomain&Rule or function\\
\midrule
\endhead
1&\(\mathcal A_9\)&\((\mathbb Z/9)^2\), matrix and internal law&
Additive-multiplicative arithmetic cell.\\
2&\((\rho_9,q_9)\)&\(\mathbb Z_{>0}\to
\{1,\ldots,9\}\times\mathbb Z_{\geq0}\)&
Nonadic division with positive representative and quotient.\\
3&\(\rho_3^{\rm or},c_3\)&\(\mathbb Z\to
\{-1,0,+1\}\times\mathbb Z\)&
Lift ternary oriented \(n=\rho_3^{\rm or}(n)+3c_3(n)\).\\
4&\(\iota_{1000}:=(q_{1000},r_{1000})\)&\(\mathbb Z\to
\mathbb Z\times\{0,\ldots,999\}\)&
Euclidean division \(N=1000q_{1000}+r_{1000}\).\\
5&\(T_d\)&\(X_9\to X_9\), \(d\in\{N,E,S,O\}\)&
Cardinal oriented translation.\\
6&\(F_{\rm loc}\)&\(\mathcal Q_{\rm loc}\times
\mathscr D_{\rm card}\to\mathcal Q_{\rm loc}\)&
Local position feedback and sheet.\\
7&\(T_{\min}\)&\(\Omega_{\rm seed}\to\mathcal U_6\)&
Minimum emitter of the two cursors, including the coordinate signature.\\
8&\(G\)&\(\operatorname{im}s_\Sigma\times
\operatorname{im}s_\Pi\to\{0,\ldots,999\}^6\)&
Coordinate-by-coordinate decimal coupling.\\
9&\(\Psi_6\)&\(\mathcal U_6\to\mathbb F_3^6\)&
Tertiary reduction of the catalog.\\
10&\(\mathcal R_3\)&\(\mathcal Q_{\rm loc}\times
\mathscr D_{\rm card}^3\to\{-1,0,+1\}\)&
Local three-step observable factor.\\
11&\(\mathcal K_{27}\)&\(\mathcal Q_{\rm obs}\to\mathcal Q_{\rm obs}\)&
Composition \(F_{\rm obs}^{27}\).\\
12&\(\mathcal X_{\rm TPK}^{\rm enr}\)&category of integral holonomic states&
Preserve all the typed coordinates of the nucleus.\\
13&\(\mathcal H\)&relation on integral holonomic states&
Hensel Transition that retains residue and quotient.\\
14&\(\mathcal C_{3,1000}\)&crossed states
\((N,L,D,K,A,B)\)&
Exact update of ternary-decimal cylinders.\\
15&\(\Sigma_{\rm B}\)&boundary memories \(\to\) signature joint&
Publish the block orientation preserving its genealogical record.\\
16&\({\rm EF}\text{--}G_9\)&extension relation&
Filter compatible extensions and organize the nonadic genealogy.\\
17&\(\Gamma_9\)&enriched histories \(\longrightarrow\) histories&
Holonomy of a phase twist with boundary memory.\\
18&\(X_\chi\)&\(\varprojlim_m\operatorname{Surv}^{(\chi)}_m\)&
Global section when non-emptiness, compatibility, and uniqueness are demonstrated in full depth.\\
19&\(\mathcal P_{\rm APP}^{(2)}\)&bivariate path category&
Transports simultaneously the multiplicative evaluation and the additive accumulation.\\
20&\(F_{\rm obs}\)&\(\mathcal Q_{\rm obs}\to\mathcal Q_{\rm obs}\)&
Deterministic observable chronology with two cursors.\\
21&\(\mathcal A_{42}\)&alphabet of forty-two triads&
Extended decimal support of finite emission.\\
22&\(\operatorname{Ext}_r\)&
\(\operatorname{Ext}_r\subseteq\mathcal F_q\times\mathcal F_{q'}\)&
Typed relation of fiber prolongation.\\
23&\(\mathcal K_{\rm ph}\)&\(\mathcal U_6\times C_9\) about itself&
Transfer of phase-compatible explicit continuation.\\
24&\((P,H,Q)\)&\(\mathcal X_\times\to\mathcal X_\times\)&
Advance, half-turn, and quarter-turn of seeds.\\
25&\((c_{\rm mode},c_{\rm or})\)&caminos \(\to\mathbb Z^2\)&
Mode and orientation change cocycles.\\
26&\(E_{20}\)&finite states \(\to\) twenty blocks&
Finite triadic prolongation with sourcebook.\\
27&\(\mathcal R_{12}\)&history enriched \(\to\) record of twelve
sectors&
Timed system with route, signature, carry, and genealogical record.\\
28&\(R_{\rm sgn}\)&terminal state \(\to U_{12}^{\rm sgn}
\subseteq\mathbb Z^{12}\)&
Reader of the signed harmonic observable.\\
29&\(\mathcal H_{12}\)&\(\mathbb Z^{12}\to\mathbb Z^{12}\)&
Transformation by three orbits of step three; inverse
\(\mathcal H_{12}/4\) over the pertinent image.\\
30&\(\mathcal E_{\rm dod}=(D_3,D_4,Q)\)&
\(\mathbb Z^{12}\to\mathbb Z^{25}\)&
Transversal extractor set of arities three and four.\\
31&\(\mathcal C_\alpha\)&blocks and carries \(\to\) blocks&
Dodecaphasic closure recurrence in base thousand.\\
32&\((L_{\rm tick},\chi_{\rm mass})\)&enriched history \(\to\)
temporal character&
Posterior dimensional interface; does not participate in the numeric constructor of
this work.\\
33&\(\mathcal W_\pm\)&bidirectional lattice of index twelve&
Reading in both directions preserving orientation and origin.\\
34&\(\mathcal A_{\rm species}\)&extension families \(\to\) atlas&
Posterior realization interface, not a primitive fourth.\\
\bottomrule
\end{longtable}
\endgroup

\subsection{Causal composition of the update of the integral holonomic state}

\begin{definition}[Complete canonical state]
\label{p12-83-c1-s6:def:u006f-esquema-estado}
Let
\[
 \mathcal E:=
 \mathbb Z_{\rm quo}\times\mathbb Z_{\rm car}\times
 \mathsf{Mem}\times\mathsf{Term}\times\mathsf{Cyl}\times
 \mathsf{Fr}\times\mathsf{Pref}\times\Lambda .
\]
The canonical integral holonomic state at time \(t\) is
\begin{equation}
 x_t=(p_t,m_t,\sigma_t,b_t,r_t;
 q_t,c_t,\eta_t,s_t,C_t,\partial_t,N_t,\lambda_t)
 \in
 X_9\times\{\Sigma,\Pi\}\times\mathsf{TRIT}\times
 C_{12}\times C_9\times\mathcal E.
 \label{p12-83-c1-s6:eq:u006f-esquema-estado}
\end{equation}
The state abbreviations used in other units are projections of
this tuple. In particular, the quotient \(q_t\), carry
\(c_t\), terminal \(s_t\), cylinder \(C_t\), boundary
\(\partial_t\), prefix \(N_t\), and ledger \(\lambda_t\) are not omitted.
\end{definition}

Let \(\rho_t\) be the observable arrow determined by \(x_t\). The transport
laws of U005 and the transition \(\mathcal C_{3,1000}\) determine a memory
endomorphism
\begin{align}
 \mathbf E_t:\mathcal E&\longrightarrow\mathcal E,\notag\\
 (q,c,\eta,s,C,\partial,N,\lambda)&\longmapsto
 \left(
 \begin{aligned}
 &\mathsf H(\rho_t)q,
  \mathsf C(\rho_t)c+\kappa(\rho_t),
  \mathsf M(\rho_t)\eta,\\
 &\mathsf T(\rho_t)s,
  \mathsf{Cyl}(\rho_t)C,
  \mathsf{Fr}(\rho_t)\partial,\\
 &\mathsf{Pref}(\rho_t)N,
  \Lambda(\rho_t)\lambda
 \end{aligned}
 \right).
 \label{p12-83-c1-s6:eq:u006f-actualizacion-memoria}
\end{align}
Each component published in \eqref{p12-83-c1-s6:eq:u006f-actualizacion-memoria} has the
domain indicated by its fiber; \(\kappa(\rho_t)\) is the carry cochain.

We define three endomorphisms of the canonical product. If
\(\sigma'=M_{\rm ph}(g(t))\) and \(m'=\mu(\sigma',m)\), then
\begin{align}
 \mathsf{Sel}_t(p,m,\sigma,b,r;e)
  &:=(p,m',\sigma',b,r;e),\label{p12-83-c1-s6:eq:u006f-selector-tipado}\\
 \mathsf{Tra}_t(p,m,\sigma,b,r;e)
  &:=(T_{d(t)}p,m,\sigma,b,r;e),\label{p12-83-c1-s6:eq:u006f-transporte-tipado}\\
 \mathsf{Upd}_t(p,m,\sigma,b,r;e)
  &:=(p,m,\sigma,b',r';\mathbf E_t(e)),
 \label{p12-83-c1-s6:eq:u006f-memoria-tipada}
\end{align}
where \((b',r')\) is the destination of \(\rho_t:(b,r)\to(b',r')\). The causal update is a composition of maps, not the composition of a map with a tritic value:
\begin{equation}
 \boxed{\mathcal U_t:=
 \mathsf{Upd}_t\circ\mathsf{Tra}_t\circ\mathsf{Sel}_t.}
 \label{p12-83-c1-s6:eq:u006f-tuberia}
\end{equation}

\begin{theorem}[Closure of causal update]
\label{p12-83-c1-s6:thm:u006f-cierre-tuberia}
The map \(\mathcal U_t\) is an endomorphism of the complete canonical
state. Readers, prolongations, and projections are applied after
\(\mathcal U_t\) only when their domain has been produced.
\end{theorem}

\begin{proof}
\(\mathsf{Sel}_t\) changes only the mode and the TRIT;
\(\mathsf{Tra}_t\) changes only the APP position; and \(\mathsf{Upd}_t\) applies the
typed product \eqref{p12-83-c1-s6:eq:u006f-actualizacion-memoria} and the phase arrow. All three codomains are the
same product of \cref{p12-83-c1-s6:def:u006f-esquema-estado}; their composition is therefore defined and
preserves all state coordinates.
\end{proof}

\subsection{Typed separation between transport, memory, calendar, extraction and return operators}

The inventory imposes
\begin{equation}
 \mathcal S_4\neq D_4\neq\operatorname{sd}_4,
 \qquad
 \mathcal R_{12}\neq C_{12}\neq\Theta_{108}\neq\Gamma_{108},
 \qquad
 F_{\rm obs}\neq\mathcal K_{\rm ph},
 \label{p12-83-c1-s6:eq:u006f-no-identificacion-uno}
\end{equation}
and
\begin{equation}
 U_{12}^{\rm cat}\neq U_{12}^{\rm sgn},
 \qquad
 \mathcal E_{\rm dod}\neq D_{\rm raw}^{90,120}
 \neq\mathcal K_{6\to8}.
 \label{p12-83-c1-s6:eq:u006f-no-identificacion-dos}
\end{equation}
They are type incompatibilities, not mere value inequalities.

\begin{criterion}[Verification of operator completeness]
\label{p12-83-c1-s6:crit:u006f-completitud}
An execution is operator-complete only if:
\begin{enumerate}
 \item declares the inputs of \(\mathfrak O_{\rm TPK}\) used;
 \item is recorded the domain and the codomain of each composition;
 \item preserves the causal order of \eqref{p12-83-c1-s6:eq:u006f-tuberia};
 \item distinguishes dynamics, transfer, reading, prolongation, and output;
 \item permits reconstruction from \(\lambda_t\) of every block change and
 every carry;
 \item does not promote a finite check to an infinite section without
 proving non-emptiness, compatibility, and uniqueness at every depth.
\end{enumerate}
Any symbol used without a type, or any fusion prohibited by
\eqref{p12-83-c1-s6:eq:u006f-no-identificacion-uno} and
\eqref{p12-83-c1-s6:eq:u006f-no-identificacion-dos}, constitutes a verifiable failure.
\end{criterion}


% Unidad U006E. Separación tipada de los canales 90/120 y de los símbolos
% epsilon. La realización de incidencia se construye después de W12 y W24.

\section{Operator decomposition of incidence channels \texorpdfstring{\(90/120\)}{90/120}: temporal selector, nilpotent generator and angular component}
\label{p12-83-c1-s7:chap:canales-90-120-epsilon}

The integers \(90\) and \(120\) appear in several strata of the construction. Their
numerical coincidence does not establish an identity of objects. Likewise,
the letter epsilon denotes three constructions with incompatible domains. The
purpose of this unit is to fix their types before using them and to indicate
precisely which realizations require structures that have not yet been
constructed.

\subsection{Temporal, nilpotent, and angular operators: domains and codomains}

\subsubsection{Action of the tritic temporal selector over \(\{-1,0,+1\}\): opening, threshold, and return with memory}

The internal calendar selector is
\begin{equation}
 \varepsilon_9:\{1,\ldots,9\}
 \longrightarrow\{-1,0,+1\},
 \qquad
 \varepsilon_9(r)=
 \begin{cases}
  +1,&r\in\{1,4,7\},\\
  -1,&r\in\{2,5,8\},\\
  0,&r\in\{3,6,9\}.
 \end{cases}
 \label{p12-83-c1-s7:eq:u006e-epsilon-temporal}
\end{equation}
Its codomain is the oriented TRIT. It determines which sheet acts and does not itself shift any position.

\subsubsection{Nilpotent parabolic generator \(J_0^2=0\) and evolution at the temporal threshold}

The class
\begin{equation}
 \varepsilon_{\rm nil}
 \in
 \mathcal Q_{\rm n}:=
 \mathbb F_3[\varepsilon_{\rm nil}]
 /(\varepsilon_{\rm nil}^{\,2})
 \label{p12-83-c1-s7:eq:u006e-epsilon-nilpotente}
\end{equation}
satisfies
\begin{equation}
 \varepsilon_{\rm nil}\neq0,
 \qquad
 \varepsilon_{\rm nil}^{\,2}=0.
 \label{p12-83-c1-s7:eq:u006e-epsilon-nilpotente-relacion}
\end{equation}
It is an element of a degenerate quadratic algebra; it is not a temporal selector.

\subsubsection{Angular residue induced after temporal selection and orientation preservation}

The symbol
\begin{equation}
 \varepsilon_{\rm res}^{\circ}\in\mathbb R
 \label{p12-83-c1-s7:eq:u006e-epsilon-residual-tipo}
\end{equation}
designates, when its geometric parameters are introduced, the residue of an
angular identity. At this stage only its type remains reserved. It is not
obtained from \(\varepsilon_9\) or from
\(\varepsilon_{\rm nil}\), and it does not intervene in the construction of
\(\pi,e,\varphi\) or \(\alpha\).

\begin{proposition}[No identification of the epsilon symbols]
\label{p12-83-c1-s7:prop:u006e-no-identificacion-epsilon}
There is no well-typed equality between
\[
 \varepsilon_9,\qquad
 \varepsilon_{\rm nil},\qquad
 \varepsilon_{\rm res}^{\circ}.
\]
\end{proposition}

\begin{proof}
The first object is a function between finite sets; the second, a
nilpotent element of an algebra over \(\mathbb F_3\); the third, a
real scalar with angular unit. Their domains, codomains, and laws of
composition are distinct.
\end{proof}

\subsection{Distribution of the indices \(90/120\) among the \(6\) objects of the incidence system}

\begin{definition}[Dodecaphasic Extractor]
\label{p12-83-c1-s7:def:u006e-extractor-dodecafasico}
The integral extractor
\begin{equation}
 \mathcal E_{\rm dod}:=(D_3,D_4,Q):
 \mathbb Z^{12}\longrightarrow\mathbb Z^{25}
 \label{p12-83-c1-s7:eq:u006e-extractor-dodecafasico}
\end{equation}
reads differences of arities three and four and a sum coordinate. The
subscripts are not angular lengths.
\end{definition}

\begin{definition}[Raw difference of sequences]
\label{p12-83-c1-s7:def:u006e-diferencia-cruda}
For \(|q|<1\), define
\begin{equation}
 S_m^{\rm raw}(q):=\frac{q^m}{1-q^{3m}},
 \qquad
 D_{\rm raw}^{90,120}(q)
 :=S_{90}^{\rm raw}(q)-S_{120}^{\rm raw}(q).
 \label{p12-83-c1-s7:eq:u006e-series-crudas}
\end{equation}
This object is an analytic function of \(q\); it is not a finite census.
\end{definition}

\begin{definition}[Marked families of incidence]
\label{p12-83-c1-s7:def:u006e-familias-incidencia}
Once a hexad \(H\), its image
\(\ell(H)\) in a binary dodecad \(D\), and the lifted octad
\(O_H\) have been constructed, one will define
\begin{align}
 \mathcal F_6(H)
 &:=\ell(H)\times\binom{\ell(H)}2,
 \label{p12-83-c1-s7:eq:u006e-familia-seis}\\
 \mathcal F_8(H)
 &:=O_H\times\binom{\ell(H)}2.
 \label{p12-83-c1-s7:eq:u006e-familia-ocho}
\end{align}
By construction,
\begin{equation}
 |\mathcal F_6(H)|
 =6\binom62=90,
 \qquad
 |\mathcal F_8(H)|
 =8\binom62=120.
 \label{p12-83-c1-s7:eq:u006e-cardinales-incidencia}
\end{equation}
The inclusion \(\ell(H)\hookrightarrow O_H\) inducira
\begin{equation}
 \mathcal I_{6\to8}:
 \mathcal F_6(H)\hookrightarrow\mathcal F_8(H).
 \label{p12-83-c1-s7:eq:u006e-inclusion-incidencia}
\end{equation}
\end{definition}

The \cref{p12-83-c1-s7:eq:u006e-familia-seis,p12-83-c1-s7:eq:u006e-familia-ocho} fix the
type from now on. Their effective existence will be proved after constructing
\(C_W\), \(W_{12}\), the binary datum \(G_{24},D,\ell\), and \(W_{24}\).
That branch is not anticipated here.

\begin{definition}[Normalized Queue Combination]
\label{p12-83-c1-s7:def:u006e-combinacion-colas}
For parameters \(q_-,q_+\in(0,1)\), let
\begin{equation}
 S_m(q):=\frac{q^m}{1-q^{3m}},
 \qquad
 \Delta S_m:=S_m(q_-)-S_m(q_+).
 \label{p12-83-c1-s7:eq:u006e-colas-normalizadas}
\end{equation}
The functional
\begin{equation}
 \mathcal K_{6\to8}
 :=12\,\Delta S_{90}-\Delta S_{120}
 \label{p12-83-c1-s7:eq:u006e-funcional-seis-ocho}
\end{equation}
is an analytic combination of tails. Its coefficient twelve belongs to its
normalization and does not turn it into the record
\(\mathcal R_{12}\).
\end{definition}

\begin{definition}[Semigroup of Degrees and Moments]
\label{p12-83-c1-s7:def:u006e-semigrupo-momentos}
The additive semigroup generated by both degrees is
\begin{equation}
 \mathfrak S_{90,120}
 :=\langle90,120\rangle
 =\{90a+120b:a,b\in\mathbb N_0\}.
 \label{p12-83-c1-s7:eq:u006e-semigrupo}
\end{equation}
For \(n\in\mathfrak S_{90,120}\), the oriented moment is
\begin{equation}
 s_n:=q_-^n-q_+^n.
 \label{p12-83-c1-s7:eq:u006e-momento}
\end{equation}
\end{definition}

The moments are differential monomials indexed by a semigroup; the
series of \eqref{p12-83-c1-s7:eq:u006e-colas-normalizadas} are infinite sums of those
moments over specific progressions.

\subsection{Normalized incidence defect}

When the inclusion \(\mathcal I_{6\to8}\) has been constructed, its
difference of cardinalities will be
\[
 |\mathcal F_6(H)|-|\mathcal F_8(H)|=90-120=-30.
\]
If normalization is declared
\[
 24\binom62=360,
\]
we obtain
\begin{equation}
 \frac{90-120}{24\binom62}
 =\frac{6-8}{24}
 =-\frac1{12}.
 \label{p12-83-c1-s7:eq:u006e-defecto-normalizado}
\end{equation}
The inclusion determines \(-30\); the rational \(-1/12\) depends moreover of
the declared normalization. No it what must present as if the cardinality
by if alone fixed the denominator.

\subsection{Typing invariants}

\begin{theorem}[Separation of the six objects]
\label{p12-83-c1-s7:thm:u006e-separacion-seis-objetos}
The objects
\begin{equation}
 \mathcal E_{\rm dod},\quad
 D_{\rm raw}^{90,120},\quad
 \mathcal I_{6\to8},\quad
 \mathcal K_{6\to8},\quad
 \mathfrak S_{90,120},\quad
 \varepsilon_{\rm res}^{\circ}
 \label{p12-83-c1-s7:eq:u006e-seis-objetos}
\end{equation}
are mutually distinct by type: they are, respectively, an integer reader,
an analytic function, a finite morphism, a combination of series, an
index semigroup, and a real angular scalar.
\end{theorem}

\begin{proof}
The conclusion is obtained by comparing their domains and codomains:
\[
 \mathbb Z^{12}\to\mathbb Z^{25},\quad
 \mathbb D\to\mathbb C,\quad
 \mathcal F_6(H)\to\mathcal F_8(H),\quad
 (0,1)^2\to\mathbb R,\quad
 \mathbb N_0^2\to\mathbb N_0,\quad
 \{\ast\}\to\mathbb R.
\]
There are no two identical signatures.
\end{proof}

\begin{criterion}[Refutation criteria]
\label{p12-83-c1-s7:crit:u006e-refutacion}
The classification is refuted if any of its domains or
codomains is altered, if the census of
\eqref{p12-83-c1-s7:eq:u006e-cardinales-incidencia} does not reproduce \(90\) and \(120\), if
\(-1/12\) is attributed to the inclusion without declaring the normalization, or if
the angular residual is presented as an output of the constants generator of
this monograph.
\end{criterion}


% Unidad U006A. Factor local, observacion estroboscopica y transicion exacta
% de cilindros.

\section{Local transport factor and ternary-decimal transition}
\label{p12-83-c1-s8:chap:tpk-factor-local-transicion-cilindros}

This unit separates three operations involved in the transport of the
Topological Phase Kernel that do not share a domain: the cardinal realization
of a trit, the stroboscopic sampling of a sequence, and the exact update
of a pair of ternary--decimal cylinders. The first is a finite
dynamics; the second is a temporal reader; the third preserves the integers
that make it possible to decide compatibility without floating-point arithmetic.

\subsection{Local state and sheet feedback}

Let
\[
 X_9=(\mathbb Z/9\mathbb Z)^2,
 \qquad
 \mathscr D_{\rm card}=\{N,E,S,O\},
\]
with
\[
 \delta(N)=(-1,0),\quad \delta(E)=(0,1),\quad
 \delta(S)=(1,0),\quad \delta(O)=(0,-1).
\]
The visible evaluations of APP are
\[
 L_\Sigma(x,y)=\rho_9((x+1)+(y+1)),
 \qquad
 L_\Pi(x,y)=\rho_9((x+1)(y+1)),
\]
where \(\rho_9(n)=1+((n-1)\bmod9)\). We further define
\[
 \rho_3^{\rm or}([0])=0,
 \qquad
 \rho_3^{\rm or}([1])=+1,
 \qquad
 \rho_3^{\rm or}([2])=-1.
\]

\begin{definition}[Local factor of APP--TRIT]
\label{p12-83-c1-s8:def:u006a-factor-local}
The local state space is
\[
 \mathcal Q_{\rm loc}=X_9\times\{\Sigma,\Pi\},
 \qquad |\mathcal Q_{\rm loc}|=162.
\]
For \(q=(p,m)\), \(d\in\mathscr D_{\rm card}\), we set
\[
 p_d=p+\delta(d),
 \qquad
 a_d(q)=\rho_3^{\rm or}([L_m(p_d)]_3),
\]
and
\[
 \mu(a,m)=
 \begin{cases}
  \Sigma,&a=+1,\\
  \Pi,&a=-1,\\
  m,&a=0.
 \end{cases}
\]
The local transition is
\[
 F_d^{\rm loc}(p,m)=(p_d,\mu(a_d(p,m),m)).
\]
\end{definition}

The position it is modified before evaluating the sheet. The null trit preserves the
last active mode; not eliminates the sheet nor reinitializes the history.

For a cardinal word \(u=d_1d_2d_3\), let
\[
 q_j=F_{d_j}^{\rm loc}(q_{j-1}),\qquad q_0=q,
\]
and let us denote by
\[
 R_3(q,u)\in\{-1,0,+1\}
\]
the trit produced in the third displacement.

\begin{definition}[Remark on step three]
\label{p12-83-c1-s8:def:u006a-obs3}
On an oriented sequence \(a=(a_t)_{t\geq1}\) we define
\[
 \operatorname{Obs}_3(a)=(a_3,a_6,a_9,\ldots).
\]
The map \(R_3\) takes a state and a cardinal word of length three;
\(\operatorname{Obs}_3\) takes a sequence and publishes a subsequence. They are
operators of different types.
\end{definition}

\begin{theorem}[Exact surjectivity of the sampled factor]
\label{p12-83-c1-s8:thm:u006a-sobreyectividad-r3}
For every \(q\in\mathcal Q_{\rm loc}\) and every
\(b\in\{-1,0,+1\}\), the fiber
\[
 \mathcal R_3(q,b)
 =\{u\in\mathscr D_{\rm card}^3:R_3(q,u)=b\}
\]
is non-empty and satisfies
\[
 \boxed{8\leq |\mathcal R_3(q,b)|\leq40.}
\]
In particular, the \(162\cdot3=486\) state--output fibers are non-empty.
\end{theorem}

\begin{proof}
From each of the \(162\) states, the \(4^3=64\) cardinal words are
enumerated, and the recurrence of
\cref{p12-83-c1-s8:def:u006a-factor-local} is applied literally. The exhaustive census of the \(486\) pairs
\((q,b)\) has a minimum of eight, a maximum of forty, and no zero value. The
statement is therefore a finite identity on the defined system, not
a mixing hypothesis.
\end{proof}

\begin{corollary}[Realization of an arbitrary oriented sequence]
\label{p12-83-c1-s8:cor:u006a-realizacion-sucesion}
Remark. For every sequence
\((b_k)_{k\geq1}\in\{-1,0,+1\}^{\mathbb N}\) there exists a cardinal history
whose observation satisfies
\[
 (a_3,a_6,a_9,\ldots)=(b_1,b_2,b_3,\ldots).
\]
\end{corollary}

\begin{proof}
Fixing a total order on \(\mathscr D_{\rm card}^3\), at stage \(k\) the
first element of \(\mathcal R_3(q_{k-1},b_k)\) is taken. The terminal state
of the block again belongs to \(\mathcal Q_{\rm loc}\), so that the
choice may be iterated indefinitely.
\end{proof}

The corollary proves the realization capacity of the sampled factor. It does not
select the sections of \(\pi,e,\varphi\), does not fix the two intermediate
trits, and does not replace the compatibility of the integral holonomic state.

\subsection{Cylinders of the \(2\) readers}

A ternary word of length \(L\), interpreted as the integer \(N\),
determines
\[
 I_3(N,L)=\left[\frac{N}{3^L},\frac{N+1}{3^L}\right).
\]
A prefix of \(K\) decimal triads, interpreted as the integer \(D\),
determines
\[
 I_{1000}(D,K)=
 \left[\frac{D}{1000^K},\frac{D+1}{1000^K}\right).
\]
We introduce the cross residues
\[
 A=N1000^K-D3^L,
 \qquad
 B=(N+1)1000^K-D3^L=A+1000^K.
\]

\begin{lemma}[Integer compatibility criterion]
\label{p12-83-c1-s8:lem:u006a-criterio-entero}
The cylinders \(I_3(N,L)\) and \(I_{1000}(D,K)\) intersect if and only if
\[
 \boxed{A<3^L,\qquad B>0.}
\]
Equivalently,
\[
 N1000^K<(D+1)3^L,
 \qquad
 (N+1)1000^K>D3^L.
\]
\end{lemma}

\begin{proof}
Two half-open intervals \([a,b)\) and \([c,d)\) intersect
exactly when \(a<d\) and \(c<b\). Multiplication by
\(3^L1000^K>0\) produces the two inequalities. Subtracting \(D3^L\) and
\(N1000^K\), respectively, gives the form in \(A,B\).
\end{proof}

\subsection{Completion operator \(C_{3,1000}\): ternary extension and preservation of the unit in \(10^3\) positions}

A new block of six trits is represented by
\[
 u\in\{0,\ldots,3^6-1\}=\{0,\ldots,728\}.
\]
The ternary update is
\[
 N'=729N+u,
 \qquad L'=L+6.
\]
Decimal publication has two cases. If the refined cylinder does not yet
determine a new triad, then \(K'=K\), \(D'=D\). If it determines
\(\delta\in\{0,\ldots,999\}\), then
\[
 K'=K+1,
 \qquad D'=1000D+\delta.
\]

\begin{definition}[Crossed state and exact transition]
\label{p12-83-c1-s8:def:u006a-c31000}
The cylinder state is the sextuple
\[
 \mathfrak c=(N,L,D,K,A,B)
\]
subject to
\[
 A=N1000^K-D3^L,
 \qquad B=A+1000^K.
\]
The transition \(\mathcal C_{3,1000}\) adjoins \(u\) and, when corresponds,
\(\delta\), and updates
\[
 (N,L,D,K,A,B)\longmapsto
 (N',L',D',K',A',B')
\]
through means of
\[
 A'=
 \begin{cases}
  729A+u1000^K,&K'=K,\\[1mm]
  729000A+u1000^{K+1}-\delta\,729\,3^L,&K'=K+1,
 \end{cases}
 \qquad
 B'=A'+1000^{K'}.
\]
The prolongation is admissible exactly when
\[
 A'<3^{L'},\qquad B'>0.
\]
\end{definition}

\begin{theorem}[Accuracy and Sufficiency of the Cross State]
\label{p12-83-c1-s8:thm:u006a-exactitud-c31000}
The transition of the \cref{p12-83-c1-s8:def:u006a-c31000} exactly preserves the
identities
\[
 A'=N'1000^{K'}-D'3^{L'},
 \qquad
 B'=(N'+1)1000^{K'}-D'3^{L'}.
\]
Consequently, \((A',B',L',K')\) suffices to decide the compatibility of the
new pair of cylinders, while \((N',D')\) preserves its reconstructive
provenance.
\end{theorem}

\begin{proof}
If a decimal triplet does not appear,
\[
 N'1000^K-D3^{L+6}
 =(729N+u)1000^K-729D3^L
 =729A+u1000^K.
\]
Si aparece \(\delta\),
\begin{align*}
 N'1000^{K+1}-D'3^{L+6}
 &=(729N+u)1000^{K+1}-(1000D+\delta)7293^L\\
 &=729000A+u1000^{K+1}-\delta\,729\,3^L.
\end{align*}
In both cases, increasing \(N'\) by one unit adds \(1000^{K'}\), which
yields \(B'\). The criterion of the \cref{p12-83-c1-s8:lem:u006a-criterio-entero} concludes.
\end{proof}

\begin{proposition}[Compatibility with concatenation]
\label{p12-83-c1-s8:prop:u006a-concatenacion-c31000}
Successively applying \(\mathcal C_{3,1000}\) to two blocks produces the same pair
\((N'',D'')\) and the same cross-residues as appending at once the two
ternary words and the decimal triads published in the same order.
\end{proposition}

\begin{proof}
The recurrences \(N\mapsto729N+u\) and
\(D\mapsto1000D+\delta\) are the laws of positional concatenation in
bases \(729\) and \(1000\). The
\cref{p12-83-c1-s8:thm:u006a-exactitud-c31000} identifies \(A,B\) with functions of those
four integers; hence its stagewise update coincides with the
evaluation subsequent to concatenation.
\end{proof}

\subsection{Decomposition of the Topological Phase Kernel into advance, rotation, carry, and memory operators}

The correct sequence is
\[
 (q,u)\xmapsto{\ R_3\ }b,
 \qquad
 (a_t)_{t\geq1}
 \xmapsto{\ \operatorname{Obs}_3\ }
 (a_{3k})_{k\geq1},
 \qquad
 (\mathfrak c,u,\delta)
 \xmapsto{\ \mathcal C_{3,1000}\ }
 \mathfrak c'.
\]
\(R_3\) proves local realizability; \(\operatorname{Obs}_3\) publishes a
subsequence; \(\mathcal C_{3,1000}\) transports compatibility between two
bases. None of the three can replace the survival relation
of the full TPK.

\begin{criterion}[Refutation criteria]
\label{p12-83-c1-s8:crit:u006a-refutacion}
The construction is refuted in its domain if any of the following
facts occurs:
\begin{enumerate}
 \item any of the \(486\) fibers \(\mathcal R_3(q,b)\) is empty or remains
       outside of the interval \([8,40]\);
 \item the null trit changes the active mode;
 \item \(R_3\) is identified with \(\operatorname{Obs}_3\);
 \item the recurrence of \(A'\) does not reproduce its direct definition;
 \item a declared compatible prolongation violates \(A'<3^{L'}\) or
       \(B'>0\);
 \item the decision depends on a numerical tolerance absent from the
       integer inequalities.
\end{enumerate}
\end{criterion}


% Unidad U006D. Registro dodecafásico integral y descomposición 1+5+6.

\section{Record dodecafasico integral, extension ciclica and decomposition \texorpdfstring{\(1+5+6\)}{1+5+6}}
\label{p12-83-c1-s9:chap:registro-dodecafasico-integral}

The observable period of one hundred eight positions contains twelve sectors of
nine positions. This equality of cardinalities does not suffice to describe the
temporal state: the ordered support, cyclic law, sector index, and
enriched register of transports, orientations, carries, and incidences must be
distinguished. This unit constructs the four objects and proves
the integral reconstruction of the twelve memory coordinates.

\subsection{Ordered support and cyclic coordinate}

For \(m\in\{1,\ldots,12\}\), define
\begin{equation}
 E_m:=\{9(m-1)+1,\ldots,9m\}.
 \label{p12-83-c1-s9:eq:u006d-bloques-nonadicos}
\end{equation}
The ordered support and the cyclic group are
\begin{equation}
 \Gamma_{108}:=
 E_1\sqcup\cdots\sqcup E_{12},
 \qquad
 \Theta_{108}:=\mathbb Z/108\mathbb Z.
 \label{p12-83-c1-s9:eq:u006d-gamma-theta}
\end{equation}
\(\Gamma_{108}\) preserves position, order and pertenencia to block;
\(\Theta_{108}\) preserves the law of translation. No are the same object.

For \(\theta\in\Theta_{108}\), let
\(\widetilde\theta\in\{0,\ldots,107\}\) be its representative. We define
\begin{equation}
 \beta(\theta):=
 \left[
 \left\lfloor\frac{\widetilde\theta}{9}\right\rfloor
 \right]_{12},
 \qquad
 r(\theta):=1+(\widetilde\theta\bmod9).
 \label{p12-83-c1-s9:eq:u006d-coordenadas-bloque-residuo}
\end{equation}
The first coordinate locates the duodecaphase sector; the second publishes
the positive position \(1,\ldots,9\) within the sector.

\subsection{Nonsplit cyclic extension}

Let \(C_n=\mathbb Z/n\mathbb Z\). The applications
\begin{equation}
 \iota:C_{12}\longrightarrow C_{108},
 \quad \iota([b]_{12})=[9b]_{108},
 \qquad
 p:C_{108}\longrightarrow C_9,
 \quad p([\theta]_{108})=[\theta]_9
 \label{p12-83-c1-s9:eq:u006d-mapas-extension}
\end{equation}
form the sequence
\begin{equation}
 0\longrightarrow C_{12}
 \xrightarrow{\;\iota\;}C_{108}
 \xrightarrow{\;p\;}C_9
 \longrightarrow0.
 \label{p12-83-c1-s9:eq:u006d-extension}
\end{equation}
For the set-theoretic section
\[
 s([r]_9):=[\widetilde r]_{108},
 \qquad 0\leq\widetilde r<9,
\]
the defect of additivity is
\begin{equation}
 c(r,r'):=
 \left[
 \left\lfloor
 \frac{\widetilde r+\widetilde r'}9
 \right\rfloor
 \right]_{12},
 \qquad
 s(r)+s(r')
 =s(r+r')+\iota(c(r,r')).
 \label{p12-83-c1-s9:eq:u006d-cociclo-extension}
\end{equation}

\begin{theorem}[Accuracy, non-division, and cohomological class]
\label{p12-83-c1-s9:thm:u006d-extension-no-escindida}
The sucesion \eqref{p12-83-c1-s9:eq:u006d-extension} is exacta and not it escinde. For the
action trivial,
\begin{equation}
 H^2(C_9,C_{12})
 \simeq C_{12}/9C_{12}
 \simeq C_3,
 \label{p12-83-c1-s9:eq:u006d-cohomologia-extension}
\end{equation}
and the cocycle \eqref{p12-83-c1-s9:eq:u006d-cociclo-extension} represents its generating
class.
\end{theorem}

\begin{proof}
The image of \(\iota\) is the subgroup of multiples of nine, which coincides
with \(\ker p\); \(p\) is surjective. This proves exactness. If the
extension were split, \(C_{108}\) would be isomorphic to
\(C_{12}\times C_9\), but
\[
 \exp(C_{108})=108,
 \qquad
 \exp(C_{12}\times C_9)=\operatorname{mcm}(12,9)=36.
\]
The cocycle identity is the Euclidean division of \(\widetilde r+\widetilde r'\) by nine.
For a cyclic group with trivial action, \(H^2(C_9,C_{12})\simeq C_{12}/9C_{12}\); the unit carry
projects to \([1]\), which generates the quotient of order three.
\end{proof}

The non-splitting typifies the return:
\[
 g_0(\theta+9)=g_0(\theta),
 \qquad
 \beta(\theta+9)=\beta(\theta)+1\pmod{12}.
\]
Nine iterations close the visible phase and shift one sector; one hundred and
eight close both clock coordinates. The integral holonomic state also preserves
vertical memory, the cylinder, the boundary, and genealogy.

\subsection{Enhanced temporal registration: joint encoding of phase, sheet, boundary, and memory}

\begin{definition}[Oriented dodecaphonic register]
\label{p12-83-c1-s9:def:u006d-registro-dodecafasico}
For an enriched history, let \(\tau_m\) be the ordered subroute on
\(E_m\), \(\sigma_m\) its orientation, \(\kappa_m\) the carry crossing the
block boundary, and \(\Lambda_m\) its local incidence ledger. The
complete record is
\begin{equation}
 \mathcal R_{12}:=
 \bigl(
 (E_m,\tau_m,\sigma_m,\kappa_m,\Lambda_m)
 \bigr)_{m=1}^{12}.
 \label{p12-83-c1-s9:eq:u006d-registro-r12}
\end{equation}
\end{definition}

Thus four levels are distinguished:
\begin{equation}
 \boxed{
 \mathcal R_{12}
 \neq C_{12}
 \neq\Theta_{108}
 \neq\Gamma_{108}.}
 \label{p12-83-c1-s9:eq:u006d-cuatro-objetos}
\end{equation}
The first is an enriched register; the second, a sector index; the
third, a group of translations; the fourth, an ordered support. There exist
projections between them, but no identifications.

Let \(q_{\rm mem}\in\mathbb Z_{\geq0}\) be the unbounded number of nonadic
turns. The duodecaphase coordinate is only its residue
\[
 \beta=[q_{\rm mem}]_{12}.
\]
After one hundred and eight steps,
\[
 \beta\longmapsto\beta,
 \qquad
 q_{\rm mem}\longmapsto q_{\rm mem}+12.
\]
The calendar returns; the vertical memory is not reset.

\subsection{Signed event extractor}

Let \((d_t)_{t=1}^{108}\) be the digital record on
\(\Gamma_{108}\). We define
\begin{equation}
 \mathcal D_{\rm evt}:=\{2,4,6,8\},
 \qquad
 \Sigma_{12}:=(+,+,+,-,-,-,+,+,+,-,-,-).
 \label{p12-83-c1-s9:eq:u006d-datos-eventos}
\end{equation}
For \(m=0,\ldots,11\), let
\begin{equation}
 e_m:=
 \Sigma_{12}(m+1)\,
 \#\{t\in E_{m+1}:d_t\in\mathcal D_{\rm evt}\}
 \in\mathbb Z.
 \label{p12-83-c1-s9:eq:u006d-eventos-firmados}
\end{equation}
The application preserves the local census and the sectorial orientation as separate factors.

We pair the two semicrowns by means of
\begin{equation}
 p_i:=e_i+e_{i+6},
 \qquad
 a_i:=e_i-e_{i+6},
 \qquad 0\leq i\leq5.
 \label{p12-83-c1-s9:eq:u006d-pares-antipodales}
\end{equation}
We define the total census, the five differences with respect to the sixth pair,
and the six antisymmetries:
\begin{equation}
 s_C:=\sum_{i=0}^{5}p_i,
 \qquad
 M_i:=p_i-p_5\quad(0\leq i\leq4),
 \qquad
 A_6:=(a_0,\ldots,a_5).
 \label{p12-83-c1-s9:eq:u006d-coordenadas-uno-cinco-seis}
\end{equation}

\begin{definition}[Integral reader \(1+5+6\)]
\label{p12-83-c1-s9:def:u006d-lector-integral}
The linear reader is
\begin{equation}
 \mathcal L_{12}:\mathbb Z^{12}\longrightarrow\mathbb Z^{12},
 \qquad
 \mathcal L_{12}(e)
 :=
 (s_C,M_0,\ldots,M_4,a_0,\ldots,a_5).
 \label{p12-83-c1-s9:eq:u006d-lector-l12}
\end{equation}
The coordinate distribution is
\begin{equation}
 \underbrace{1}_{s_C}
 +\underbrace{5}_{M_0,\ldots,M_4}
 +\underbrace{6}_{a_0,\ldots,a_5}
 =12.
 \label{p12-83-c1-s9:eq:u006d-uno-cinco-seis}
\end{equation}
\end{definition}

\begin{theorem}[Exact reconstruction and integral image]
\label{p12-83-c1-s9:thm:u006d-inversa-l12}
The reader \(\mathcal L_{12}\) is injective. On its image,
\begin{align}
 p_5&=
 \frac{s_C-\sum_{i=0}^{4}M_i}{6},
 &
 p_i&=p_5+M_i\quad(0\leq i\leq4),
 \label{p12-83-c1-s9:eq:u006d-inversa-p}\\
 e_i&=\frac{p_i+a_i}{2},
 &
 e_{i+6}&=\frac{p_i-a_i}{2}
 \quad(0\leq i\leq5).
 \label{p12-83-c1-s9:eq:u006d-inversa-e}
\end{align}
An output tuple belongs to \(\operatorname{im}\mathcal L_{12}\) if and
only if
\begin{equation}
 s_C-\sum_{i=0}^{4}M_i\equiv0\pmod6
 \label{p12-83-c1-s9:eq:u006d-condicion-divisibilidad}
\end{equation}
and, for the reconstructed \(p_i\),
\begin{equation}
 p_i\equiv a_i\pmod2
 \qquad(0\leq i\leq5).
 \label{p12-83-c1-s9:eq:u006d-condicion-paridad}
\end{equation}
\end{theorem}

\begin{proof}
From \(M_i=p_i-p_5\) is obtained
\[
 s_C=\sum_{i=0}^{4}(p_5+M_i)+p_5
 =6p_5+\sum_{i=0}^{4}M_i,
\]
which gives the first equation of \eqref{p12-83-c1-s9:eq:u006d-inversa-p}; the others are
immediate. Adding and subtracting
\(p_i=e_i+e_{i+6}\) and \(a_i=e_i-e_{i+6}\) produces
\eqref{p12-83-c1-s9:eq:u006d-inversa-e}. Divisibility by six and the six
parity congruences are precisely the necessary and
sufficient conditions for all reconstructed coordinates to be integers.
\end{proof}

Division by six appears after the oriented census, not during it. The
integrality conditions for this unit are exactly
\eqref{p12-83-c1-s9:eq:u006d-condicion-divisibilidad} and
\eqref{p12-83-c1-s9:eq:u006d-condicion-paridad}; no Smith normal form is attributed here to a
matrix that has not previously been defined.

\subsection{\(2\) dodecaphasic words of different provenance}

The distinction must be preserved.
\begin{equation}
 U_{12}^{\rm cat}:=U_6\Vert U_6,
 \qquad
 U_{12}^{\rm sgn}:=\mathcal H_{12}(K).
 \label{p12-83-c1-s9:eq:u006d-dos-palabras-doce}
\end{equation}
The first concatenates two catalog emissions of length six. The second
is a signed coordinate reversibly linked to \(K\) by the operator
\(\mathcal H_{12}\). Their common length does not imply the same provenance,
semantics, or transport law.

\begin{criterion}[Falsification criteria]
\label{p12-83-c1-s9:crit:u006d-refutacion}
The construction is refuted if the sequence
\eqref{p12-83-c1-s9:eq:u006d-extension} is not exact, if its cohomology class is
trivial, if a tuple in the image violates
\eqref{p12-83-c1-s9:eq:u006d-condicion-divisibilidad} or
\eqref{p12-83-c1-s9:eq:u006d-condicion-paridad}, if the inverse formulas do not recover
the twelve initial integers, or if the return of \(\beta\) is interpreted as a
reset of \(q_{\rm mem}\).
\end{criterion}


% Unidad U006B. Funtor ciclotomico y elevacion modular de caminos.

\section{Cyclotomic phase functor--length and modular elevation of paths}
\label{p12-83-c1-s10:chap:funtor-ciclotomico-elevacion-caminos}

The odometric quotient \(\mathcal O_{9,12}\) simultaneously preserves the nonadic position, the dodecaphase register, and the integral length of each arrow. In this unit, its phase is realized over fibers \(A_2\), and the cardinal lift that replaces each edge with four oriented edges is proved. The cyclotomic representation and the absolute length will be retained as distinct coordinates.

\subsection{Coxeter representation of order \(3\) on the dodecaphasic fibers of type \(A_2\)}

Let
\begin{equation}
 G_{A_2}:=
 \begin{pmatrix}2&-1\\-1&2\end{pmatrix},
 \qquad
 C:=
 \begin{pmatrix}0&-1\\1&-1\end{pmatrix}.
 \label{p12-83-c1-s10:eq:u006b-dato-coxeter}
\end{equation}
By the \cref{p12-83-c1-s5:prop:coxeter-a2-autonomo},
\begin{equation}
 C^3=I_2,
 \qquad
 C^{\mathsf T}G_{A_2}C=G_{A_2},
 \qquad
 \det(XI_2-C)=X^2+X+1.
 \label{p12-83-c1-s10:eq:u006b-identidades-coxeter}
\end{equation}
Thus, \(C\) realizes the character cyclotomic of \(C_3\) on \(A_2\).

For each \(b\in\mathbb Z/12\mathbb Z\), let \(A_{2,b}\) be a copy of
\(A_2\), let
\[
 \iota_b:A_{2,b}\xrightarrow{\;\sim\;}A_2
\]
an integral identification, and let
\(P_b\in O(G_{A_2})\) be a frame that depends on \(b\), but not on the residue
\(r\). We define
\begin{equation}
 \rho_b:=
 \iota_b^{-1}P_bCP_b^{-1}\iota_b,
 \qquad
 \rho_b^3=\operatorname{id}_{A_{2,b}}.
 \label{p12-83-c1-s10:eq:u006b-rho-fibra}
\end{equation}
The independencia respect of \(r\) liga the character to the fiber
dodecaphasic; \(r\) preserves the position internal of the turn nonadic.

\subsection{Odometer category of the TPK: cylinders, truncations, and sequence morphisms}

Recall that \(\mathcal O_{9,12}\) has objects
\[
 (b,r)\in\mathbb Z/12\mathbb Z\times\{0,\ldots,8\}.
\]
For \(n\in\mathbb N\), the flecha unique of length \(n\) is
\begin{equation}
 \gamma_{(b,r),n}:
 (b,r)\longrightarrow
 \left(
 b+\kappa_r(n),[r+n]_9
 \right),
 \qquad
 \kappa_r(n):=
 \left\lfloor\frac{r+n}{9}\right\rfloor.
 \label{p12-83-c1-s10:eq:u006b-flecha-odometro}
\end{equation}
The composition relies on
\begin{equation}
 \kappa_r(n+m)
 =\kappa_r(n)+\kappa_{[r+n]_9}(m).
 \label{p12-83-c1-s10:eq:u006b-cociclo-odometro}
\end{equation}
The length \(\ell(\gamma_{(b,r),n})=n\) is additive.

\begin{definition}[Cyclotomic Phase--Length Functor]
\label{p12-83-c1-s10:def:u006b-funtor-ciclotomico}
Let \(\mathbf{Rep}_{\mathbb C}(C_3)\) be the category of complex
representations of \(C_3\), and let \(\mathbf B\mathbb N\) be the monoidal category
with a single object whose endomorphisms are \((\mathbb N,+)\). We define
\begin{equation}
 \mathfrak F_{\rm cyc}:
 \mathcal O_{9,12}
 \longrightarrow
 \mathbf{Rep}_{\mathbb C}(C_3)\times\mathbf B\mathbb N
 \label{p12-83-c1-s10:eq:u006b-signatura-funtor}
\end{equation}
by
\[
 \mathfrak F_{\rm cyc}(b,r)
 =(A_{2,b}\otimes_{\mathbb Z}\mathbb C,\rho_b)
\]
and
\begin{equation}
 \mathfrak F_{\rm cyc}(\gamma)
 :=
 \left(
 \iota_{b'}^{-1}P_{b'}C^{\ell(\gamma)}
 P_b^{-1}\iota_b,
 \ell(\gamma)
 \right)
 \label{p12-83-c1-s10:eq:u006b-flecha-funtor}
\end{equation}
for \(\gamma:(b,r)\to(b',r')\).
\end{definition}

\begin{theorem}[Functoriality and separation of the records]
\label{p12-83-c1-s10:thm:u006b-functorialidad}
The assignment of the \cref{p12-83-c1-s10:def:u006b-funtor-ciclotomico} is a functor. Its
first component preserves only the length class modulo three; the
second preserves the complete integer and, by reduction, determines that class.
The cyclotomic component does not allow the absolute length to be reconstructed.
\end{theorem}

\begin{proof}
If \(\gamma_1:(b_0,r_0)\to(b_1,r_1)\) and
\(\gamma_2:(b_1,r_1)\to(b_2,r_2)\) have lengths \(n_1,n_2\), the
intermediate trivialization cancels:
\begin{align}
 &\bigl(
 \iota_{b_2}^{-1}P_{b_2}C^{n_2}P_{b_1}^{-1}\iota_{b_1}
 \bigr)
 \bigl(
 \iota_{b_1}^{-1}P_{b_1}C^{n_1}P_{b_0}^{-1}\iota_{b_0}
 \bigr)\notag\\
 &\qquad=
 \iota_{b_2}^{-1}P_{b_2}C^{n_1+n_2}
 P_{b_0}^{-1}\iota_{b_0}.
 \label{p12-83-c1-s10:eq:u006b-composicion-funtor}
\end{align}
The second component composes by addition. Since \(C^n\) depends only on
\([n]_3\), whereas \(\ell=n\) preserves the integer, the two coordinates are not
identified.
\end{proof}

If
\[
 Q:\mathsf{TPK}\longrightarrow\mathcal O_{9,12}
\]
is the odometer functor of U006, the typed realization is
\begin{equation}
 \mathfrak F_{\rm cyc}\circ Q:
 \mathsf{TPK}\longrightarrow
 \mathbf{Rep}_{\mathbb C}(C_3)\times\mathbf B\mathbb N.
 \label{p12-83-c1-s10:eq:u006b-compuesto-tpk-ciclotomico}
\end{equation}

\subsection{Modular elevation of cardinal paths}

Let \(\mathcal A_{N,d}\) be the category libre of paths cardinal of
\((\mathbb Z/N\mathbb Z)^d\). Let
\[
 N\ge2,\quad u\ge1,\quad\gcd(u,N)=1,
 \quad k:=\operatorname{ord}_N(u),
 \quad q:=\frac{u^k-1}{N}.
\]

\begin{definition}[Cardinal Substitution by \(u\)]
\label{p12-83-c1-s10:def:u006b-sustitucion-cardinal}
The functor
\[
 \mathcal S_u:\mathcal A_{N,d}\longrightarrow\mathcal A_{N,d}
\]
sends an object \(p\) to \(up\) and each oriented edge to the ordered concatenation of \(u\) copies with the same direction.
\end{definition}

\begin{theorem}[Modular lifting to endomorphism of paths]
\label{p12-83-c1-s10:thm:u006b-elevacion-modular}
For every object \(p\), path \(\gamma\), and
\[
 \kappa_{N,r}(n):=
 \left\lfloor\frac{r+n}{N}\right\rfloor,
\]
we have
\begin{align}
 \mathcal S_u^k(p)&=p,
 \label{p12-83-c1-s10:eq:u006b-retorno-objetos}\\
 \ell(\mathcal S_u^k\gamma)&=u^k\ell(\gamma),
 \label{p12-83-c1-s10:eq:u006b-longitud-sustituida}\\
 \kappa_{N,r}(u^kn)-\kappa_{N,r}(n)&=qn.
 \label{p12-83-c1-s10:eq:u006b-diferencia-vueltas}
\end{align}
\end{theorem}

\begin{proof}
The category is free, so the images of objects and generators
determine a unique functor. Each map multiplies the length by
\(u\). Since \(u^k=1+qN\), the objects return modulo \(N\), and
\[
 \left\lfloor
 \frac{r+n+qNn}{N}
 \right\rfloor
 =
 \left\lfloor\frac{r+n}{N}\right\rfloor+qn.
\]
\end{proof}

In the nonadic case,
\begin{equation}
 (N,u,k,q)=(9,4,3,7),
 \qquad
 4^3=64=1+7\cdot9.
 \label{p12-83-c1-s10:eq:u006b-parametros-nonadicos}
\end{equation}
Therefore,
\begin{equation}
 \boxed{
 \mathcal S_4^3(e)
 \text{ preserves the endpoints and contains seven nonadic turns
 oriented per initial edge}.}
 \label{p12-83-c1-s10:eq:u006b-siete-vueltas}
\end{equation}
The endpoints and residual position return in three applications; the
integral holonomic state need not return. Without a proof of
invertibility, \(\mathcal S_4^3\) is a vertical path endomorphism,
not a monodromy.

\begin{remark}[Separation with respect to clock dilation]
\label{p12-83-c1-s10:rem:u006b-s4-no-d4}
\(\mathcal S_4\) replaces objects and edges in a path category.
The map \(D_4\) of U006C will multiply the coordinate
\(\Theta_{108}\) by four. They share an integer factor, but not a domain, codomain, or
composition law.
\end{remark}

\begin{criterion}[Criteria of refutation]
\label{p12-83-c1-s10:crit:u006b-refutacion}
The block remains refuted if fails any of the identities
\eqref{p12-83-c1-s10:eq:u006b-identidades-coxeter},
\eqref{p12-83-c1-s10:eq:u006b-cociclo-odometro},
\eqref{p12-83-c1-s10:eq:u006b-composicion-funtor} or
\eqref{p12-83-c1-s10:eq:u006b-parametros-nonadicos}; also remains badly typed if it
identify the phase \(C^n\), the length \(n\), the substitution
\(\mathcal S_4\) and the dilation \(D_4\).
\end{criterion}
}%


% Unidad U004. Redacción autónoma desde los generadors TRIT de 1063/live.

\section{TRIT as the minimal unit of topological information: ternary orientation, carry memory, and temporal order}
\label{p12-83-c1-s11:chap:trit-levantado-regimenes}

The mixed curvature of APP produced three oriented values: zero on the pure sheets and opposite signs when the sum--product order is exchanged. The TRIT types that triad and preserves, in addition to the visible residue, the quotient required to reconstruct the lifted integer. This unit explicitly separates four levels: ternary class, balanced representative, carry, and quadratic regime.

\begin{definition}[Topological Information Unit TRIT]
The \emph{TRIT} is the elementary oriented unit whose state space is
\(\mathsf{TRIT}=\{-1,0,+1\}\). Its native logarithmic capacity is
\(\log 3\), and its expression in binary units is \(\log_2 3\). Unlike
a bit, its three states do not reduce to a cardinal encoding: they respectively type
opening, threshold, and return, and preserve the carry quotient
required to reconstruct the transition that produced the visible
residue.
\end{definition}

\subsection{Oriented ternary system and balanced representation of TRIT states}

We define the oriented set
\[
 \mathsf{TRIT}:=\{-1,0,+1\}.
\]
The balanced chart is the bijection of sets.
\[
 \operatorname{bal}:\mathbb F_3\longrightarrow\mathsf{TRIT},
 \qquad
 [0]_3\mapsto0,\qquad[1]_3\mapsto+1,\qquad[2]_3\mapsto-1.
\]
It is not of an equality of types. \(\mathbb F_3\) carries sum and product
of field; \(\mathsf{TRIT}\) carries order and orientation. Every operation
transported between both must indicate the chart used.

The orientation of the nine phases is
\[
 \varepsilon_9(r)=
 \begin{cases}
 +1,&r\in\{1,4,7\},\\
 -1,&r\in\{2,5,8\},\\
 0,&r\in\{3,6,9\}.
 \end{cases}
\]
His complete table is
\[
\begin{array}{c|ccccccccc}
r&1&2&3&4&5&6&7&8&9\\\hline
\varepsilon_9(r)&+1&-1&0&+1&-1&0&+1&-1&0
\end{array}.
\]
Periodicity of three does not eliminate the nonadic phase: the phases
\(1,4,7\), for example, have the same orientation and different positions
in the circuit.

\subsection{Balanced ternary division algorithm and uniqueness of the quotient–residue pair}

\begin{theorem}[Local ternary lifting]
\label{p12-83-c1-s11:thm:trit-levantamiento-local}
For each \(n\in\{-4,-3,\ldots,4\}\) there exists a single pair
\((r,c)\in\mathsf{TRIT}^2\) such that
\[
 n=r+3c.
\]
\end{theorem}

\begin{proof}
The existence is verified in the table.
\[
\begin{array}{c|rrrrrrrrr}
n&-4&-3&-2&-1&0&1&2&3&4\\\hline
r&-1&0&+1&-1&0&+1&-1&0&+1\\
c&-1&-1&-1&0&0&0&+1&+1&+1
\end{array}.
\]
If \(r+3c=r'+3c'\), then \(r-r'=3(c'-c)\). Since
\(r-r'\in\{-2,-1,0,1,2\}\), the only possible multiple of three is zero;
therefore \(r=r'\) and \(c=c'\).
\end{proof}

The visible projection \((r,c)\mapsto r\) has three local states over the value zero:
\[
 0^-:=(0,-1),
 \qquad
 0^0:=(0,0),
 \qquad
 0^+:=(0,+1).
\]
They publish the same residue, but reconstruct \(-3,0,+3\). This triple fiber
is the minimal form of a distinction that will reappear in the integral holonomic state
of the Topological Phase Kernel.

\begin{corollary}[Balanced Iteration]
\label{p12-83-c1-s11:cor:trit-iteracion-balanceada}
Every integer \(N\) admits a finite balanced ternary expansion
\[
 N=r_0+3r_1+\cdots+3^mr_m,
 \qquad r_j\in\mathsf{TRIT}.
\]
It is obtained by iterating the division \(N_k=r_k+3N_{k+1}\), where \(r_k\) is the
balanced representative of \(N_k\bmod3\).
\end{corollary}

\begin{proof}
Euclidean division modulo three yields a residue in \(\{0,1,2\}\); replacing \(2\) with \(-1\) increases the quotient by one and produces \(r_k\in\mathsf{TRIT}\). The absolute value of the quotient decreases outside a finite set, so the process terminates. Uniqueness follows by comparing the first differing digit and reducing modulo three.
\end{proof}

\subsection{Composition and carry cocycle}

For \(r,s\in\mathsf{TRIT}\), we define \(r\oplus s\in\mathsf{TRIT}\) and
\(\kappa_3(r,s)\in\mathbb Z\) by the unique equality
\[
 r+s=(r\oplus s)+3\kappa_3(r,s).
\]
The table is
\[
\begin{array}{c|ccc}
\oplus&-1&0&+1\\\hline
-1&+1&-1&0\\
0&-1&0&+1\\
+1&0&+1&-1
\end{array},
\qquad
\begin{array}{c|ccc}
\kappa_3&-1&0&+1\\\hline
-1&-1&0&0\\
0&0&0&0\\
+1&0&0&+1
\end{array}.
\]

\begin{proposition}[Ternary cocyclic identity]
\label{p12-83-c1-s11:prop:trit-cociclo}
For \(r,s,t\in\mathsf{TRIT}\),
\[
 \kappa_3(r,s)+\kappa_3(r\oplus s,t)
 =\kappa_3(s,t)+\kappa_3(r,s\oplus t).
\]
\end{proposition}

\begin{proof}
Both sides are the total quotient that appears when writing
\(r+s+t\) as a balanced representative plus a multiple of three. The
uniqueness of the lift proves the equality.
\end{proof}

If \((r,c)\) and \((s,d)\) are states levantados, its sum is
\[
 (r,c)\boxplus(s,d)
 :=\bigl(r\oplus s,\ c+d+\kappa_3(r,s)\bigr).
\]
The \cref{p12-83-c1-s11:prop:trit-cociclo} proves that \(\boxplus\) is associative and that the
reconstruction \(\operatorname{rec}_3(r,c)=r+3c\) satisfies
\[
 \operatorname{rec}_3(x\boxplus y)
 =\operatorname{rec}_3(x)+\operatorname{rec}_3(y).
\]

\subsection{Orientation and Temporal Order Involution}

For to tritic word \(\tau=(\tau_1,\ldots,\tau_n)\), define
\[
 \mathcal C_{\rm rev}(\tau_1,\ldots,\tau_n)
 :=(-\tau_n,\ldots,-\tau_1).
\]

\begin{lemma}[Oriented reversal]
\label{p12-83-c1-s11:lem:trit-reversion}
The map \(\mathcal C_{\rm rev}\) is an involution. It interchanges the
sectors \(+1\) and \(-1\) and fixes the sector \(0\).
\end{lemma}

\begin{proof}
Applying it twice reverses the order twice and changes every sign twice.
The initial word is thereby recovered. The remaining assertions are
immediate in each component.
\end{proof}

When a time-oriented parameter is fixed, the TPK convention is
\[
 +1:\text{return with memory},
 \qquad
 0:\text{present threshold},
 \qquad
 -1:\text{prolongation opening}.
\]
This assignment does not add a fourth temporal state. The direction is preserved
in the sign and the exact position is preserved in the nonadic phase.

\subsection{Universal family of quadratic algebras}

For \(\tau\in\mathsf{TRIT}\), define
\[
 \mathbb A_\tau:=\mathbb R[J_\tau]/(J_\tau^2+\tau I).
\]
The three cases are:
\[
\begin{array}{c|c|c|c}
\tau&J_\tau^2&\text{regime}&\exp(tJ_\tau)\\\hline
+1&-I&\text{elliptic}&\cos t\,I+\sin t\,J_{+1}\\
0&0&\text{parabolic}&I+tJ_0\\
-1&+I&\text{hyperbolic}&\cosh t\,I+\sinh t\,J_{-1}
\end{array}.
\]

\begin{theorem}[Uniform quadratic exponential]
\label{p12-83-c1-s11:thm:trit-exponencial-uniforme}
If \(J_\tau^2=-\tau I\), then
\[
 \exp(tJ_\tau)
 =C_\tau(t)I+S_\tau(t)J_\tau,
\]
where
\[
\begin{array}{c|ccc}
\tau&+1&0&-1\\\hline
C_\tau(t)&\cos t&1&\cosh t\\
S_\tau(t)&\sin t&t&\sinh t
\end{array}.
\]
\end{theorem}

\begin{proof}
We separate the exponential series into even and odd powers. The relations
\(J_{+1}^{2m}=(-1)^mI\), \(J_0^2=0\), and
\(J_{-1}^{2m}=I\) respectively produce the cosine and sine series,
the linear truncation, and the hyperbolic series.
\end{proof}

On \(\mathbb F_3\) the analogues appear
\[
 \mathbb F_9=\mathbb F_3[\iota]/(\iota^2+1),
 \qquad
 \mathbb D_3=\mathbb F_3[\epsilon]/(\epsilon^2),
 \qquad
 \mathbb S_3=\mathbb F_3[j]/(j^2-1).
\]
The first is a field of nine elements; the second contains a nonzero
nilpotent; the third decomposes through the idempotents
\((1\pm j)/2\) and is isomorphic to \(\mathbb F_3\times\mathbb F_3\). The
coincidence of cardinality does not identify these three algebras.

The integral source of the elliptic regime is
\[
 \mathbb Q_{\mathbb Z}:=\mathbb Z[u]/(u^2+1).
\]
Reduction modulo three sends \(u\) to \(\iota\) and produces
\(\mathbb F_9\). The extension of scalars to \(\mathbb R\) and a
\(u\mapsto J_{+1}\) representation produce the real elliptic realizer.
There is therefore no untyped map \(\mathbb F_9\to\mathbb R\): both
realizations follow from a common integral source.

\subsection{Matrix representations of the TRIT selected by the sign of \(J_\tau^2\)}

Let
\[
 C=\begin{pmatrix}0&-1\\1&-1\end{pmatrix},
 \qquad C^3=I,
\]
and let us define
\[
 J_e:=\frac{2C+I}{\sqrt3}
 =\frac1{\sqrt3}\begin{pmatrix}1&-2\\2&-1\end{pmatrix}.
\]
A direct multiplication gives
\[
 (2C+I)^2=-3I,
 \qquad J_e^2=-I.
\]
The parabolic realizer is
\[
 J_p:=N=\begin{pmatrix}0&1\\0&0\end{pmatrix},
 \qquad N^2=0.
\]
For the hyperbolic regime we take
\[
 F_{\rm av}=\begin{pmatrix}0&1\\1&1\end{pmatrix},
 \qquad
 A:=F_{\rm av}^2=\begin{pmatrix}1&1\\1&2\end{pmatrix},
\]
and
\[
 J_h:=\frac{2A-3I}{\sqrt5}
 =\frac1{\sqrt5}\begin{pmatrix}-1&2\\2&1\end{pmatrix}.
\]
Then
\[
 (2A-3I)^2=5I,
 \qquad J_h^2=I.
\]

The bilinear form
\[
 G=\begin{pmatrix}2&1\\1&-2\end{pmatrix}
\]
satisfies
\[
 A^TGA=G.
\]
Moreover,
\[
 \cosh(2\log\varphi)=\frac32,
 \qquad
 \sinh(2\log\varphi)=\frac{\sqrt5}{2},
\]
for which
\[
 A=\exp(2\log\varphi\,J_h)
  =\cosh(2\log\varphi)I+\sinh(2\log\varphi)J_h.
\]
This identity will be recovered causally in the biography of
\(\varphi\); here only the matrix representation is verified.

The three exponential controls are
\[
 \exp(\pi J_e)+I=0,
 \qquad
 \exp(J_p)=I+J_p,
 \qquad
 \exp(2\log\varphi\,J_h)=F_{\rm av}^2.
\]
The presence of \(\pi\) and \(\varphi\) in these verification identities
does not authorize using them as constructors of those constants. Their
constructors will appear later from the nonadic genealogy.

\subsection{Local coupling APP--TRIT}

With the exchange
\[
 R=\begin{pmatrix}0&1\\1&0\end{pmatrix}
\]
we define, for \(\tau\in\mathsf{TRIT}\),
\[
 B_\tau=(4+3\tau)I+(5-3\tau)R.
\]
Since \(P_\pm=(I\pm R)/2\),
\[
 B_\tau=9P_++(6\tau-1)P_-.
\]

\begin{theorem}[Spectral Recovery of Orientation]
\label{p12-83-c1-s11:thm:app-trit-recuperacion-espectral}
The spectrum of \(B_\tau\) is
\[
 \operatorname{Spec}(B_\tau)=\{9,6\tau-1\}.
\]
The differential eigenvalue \(\lambda_-\) uniquely determines the TRIT:
\[
 \tau=\frac{\lambda_-+1}{6}.
\]
Moreover, the block central of APP is \(B_{+1}\).
\end{theorem}

\begin{proof}
The projectors \(P_+\) and \(P_-\) diagonalize \(R\) with eigenvalues
\(+1\) and \(-1\). The spectral formula is obtained by substitution. The second
identity solves for \(\tau\). For \(\tau=+1\),
\(B_{+1}=7I+2R=\begin{psmallmatrix}7&2\\2&7\end{psmallmatrix}\), which is
the unique block of \cref{p12-83-c1-s1:thm:app-bloque-central}.
\end{proof}

Therefore, the TRIT not it anade as retrospective label: can be recovered
of the differential mode of a family constructed on the internal exchange
of APP.

\subsection{Obstructions of TRIT: residue--representative--carry separation, involutivity, and spectral recovery of \(\tau\)}

\begin{criterion}[Refutation Criteria for the TRIT]
The unit is refuted if:
\begin{enumerate}
 \item it identify class residual, representative balanceado and carry;
 \item it is asserted that \(0^-\), \(0^0\), and \(0^+\) are the same state;
 \item the composition \(\boxplus\) ceases to reconstruct the integer sum;
 \item \(\mathcal C_{\rm rev}\) is not involutive;
 \item the three algebras are identified because they have the same cardinal;
 \item some realizer fails to satisfy \(J_\tau^2=-\tau I\);
 \item the value of \(\tau\) is chosen after observing the outcome and is not
       recovered from the spectrum of \(B_\tau\).
\end{enumerate}
\end{criterion}

APP provides the support, the multiplicative matrix, and the internal additive
law; TRIT provides orientation,
lift, and quadratic type. The next unit will construct the
Topological Phase Kernel as the transport category of these data, without
confusing the internal selector with the complete state.

\HMTDeferredTPKBase
\HMTDeferredTPKDevelopment


% Unidad U006C. Dilatación aritmética, cociclo de acarreo y
% semiconjugación cuaternaria.

\section{Quaternary dilation of the register, observable obstruction, and refined semiconjugacy}
\label{p12-83-c1-s12:chap:dilatacion-cuaternaria-semiconjugacion}

The cardinal substitution \(\mathcal S_4\) constructed in the preceding unit
acts on paths. In this unit, a second operation, arithmetic in
nature, is defined on the register
\(\mathbb Z/12\mathbb Z\times\mathbb Z/9\mathbb Z\), and a third,
categorical operation subdivides each observable transition into four
microtransitions. The aim is not to merge them, but to prove the exact
relationship between their realizations.

\subsection{Dilation on the product register}

Let
\[
 \mathcal Q_{12,9}:=
 \mathbb Z/12\mathbb Z\times\mathbb Z/9\mathbb Z.
\]
For \(r\in\mathbb Z/9\mathbb Z\), we denote its normalized representative by
\(\widetilde r\in\{0,\ldots,8\}\).

\begin{definition}[Quaternary dilation of the register]
\label{p12-83-c1-s12:def:u006c-dilatacion-registro}
We define
\begin{equation}
 D_4(b,r):=
 \left(
 \left[
  4b+\left\lfloor\frac{4\widetilde r}{9}\right\rfloor
 \right]_{12},
 [4r]_9
 \right).
 \label{p12-83-c1-s12:eq:u006c-d4-producto}
\end{equation}
The total coordinate is
\begin{equation}
 \Theta:\mathcal Q_{12,9}\xrightarrow{\;\sim\;}
 \mathbb Z/108\mathbb Z,
 \qquad
 \Theta(b,r):=[9b+\widetilde r]_{108}.
 \label{p12-83-c1-s12:eq:u006c-coordenada-total}
\end{equation}
\end{definition}

\begin{proposition}[Realization by modular multiplication]
\label{p12-83-c1-s12:prop:u006c-d4-total}
In the total coordinate,
\begin{equation}
 \boxed{
 \Theta\circ D_4\circ\Theta^{-1}(\theta)=[4\theta]_{108}.}
 \label{p12-83-c1-s12:eq:u006c-d4-total}
\end{equation}
In particular,
\begin{equation}
 \operatorname{im}(D_4)
 =4\mathbb Z/108\mathbb Z
 \simeq\mathbb Z/27\mathbb Z,
 \label{p12-83-c1-s12:eq:u006c-imagen-d4}
\end{equation}
and the induced restriction on the image has order nine.
\end{proposition}

\begin{proof}
For \(\theta=9b+\widetilde r\),
\[
 4\theta
 =9\left(
  4b+\left\lfloor\frac{4\widetilde r}{9}\right\rfloor
 \right)
 +(4\widetilde r\bmod9).
\]
This is exactly the quotient--residue decomposition of
\eqref{p12-83-c1-s12:eq:u006c-d4-producto}. The image consists of the classes that are
multiples of four. Dividing by four yields
\(\mathbb Z/27\mathbb Z\), and
\[
\begin{aligned}
 4^1&\equiv4,& 4^2&\equiv16,& 4^3&\equiv10,\\
 4^4&\equiv13,& 4^5&\equiv25,& 4^6&\equiv19,\\
 4^7&\equiv22,& 4^8&\equiv7,& 4^9&\equiv1\pmod{27}.
\end{aligned}
\]
No preceding power is one; therefore
\(\operatorname{ord}_{27}(4)=9\).
\end{proof}

\subsection{Exact carry identity}

For \(r\in\{0,\ldots,8\}\) and \(n\in\mathbb N\), define
\[
 \kappa_r(n):=\left\lfloor\frac{r+n}{9}\right\rfloor,
 \qquad
 r':=[r+n]_9.
\]

\begin{theorem}[Quaternary carry cocycle]
\label{p12-83-c1-s12:thm:u006c-cociclo-d4}
We have
\begin{equation}
 \boxed{
 \kappa_{[4r]_9}(4n)
 =4\kappa_r(n)
 +\left\lfloor\frac{4\widetilde r'}{9}\right\rfloor
 -\left\lfloor\frac{4\widetilde r}{9}\right\rfloor.}
 \label{p12-83-c1-s12:eq:u006c-cociclo-d4}
\end{equation}
The last difference is the coboundary produced by the choice of
representatives \(r\mapsto\widetilde r\).
\end{theorem}

\begin{proof}
Write
\[
 r+n=9\kappa_r(n)+\widetilde r'.
\]
When multiplied by four,
\[
 4r+4n=36\kappa_r(n)+4\widetilde r'.
\]
Separating quotient and residue modulo nine and subtracting the quotient already contained in
\(4\widetilde r\) yields
\eqref{p12-83-c1-s12:eq:u006c-cociclo-d4}.
\end{proof}

\begin{corollary}[Third iteration]
\label{p12-83-c1-s12:cor:u006c-tercera-iteracion}
For every \((b,r)\in\mathcal Q_{12,9}\),
\begin{equation}
 \boxed{D_4^3(b,r)=([4b+7\widetilde r]_{12},r).}
 \label{p12-83-c1-s12:eq:u006c-d4-cubo}
\end{equation}
\end{corollary}

\begin{proof}
By \cref{p12-83-c1-s12:prop:u006c-d4-total},
\[
 \Theta(D_4^3(b,r))
 =[4^3(9b+\widetilde r)]_{108}
 =[64(9b+\widetilde r)]_{108}.
\]
Since \(64=1+7\cdot9\), the residue modulo nine is again \(r\), whereas
the block coordinate becomes
\([4b+7\widetilde r]_{12}\).
\end{proof}

The return of \(r\) is not return of the complete register. The difference
\(7\widetilde r\) preserves the memory of the internal representative traversed.

\subsection{Obstruction in the observable quotient}

In the observable configuration of
\cref{p12-83-c1-s3:def:tpk-configuracion-observable}, we write
\[
 q=(p^\Sigma,p^\Pi,d^\Sigma,d^\Pi,\theta)
 \in\mathcal Q_{\rm obs}.
\]
The positional projection is
\begin{equation}
 \pi_{\rm pos}(q):=(p^\Sigma,p^\Pi)\in X_9^2.
 \label{p12-83-c1-s12:eq:u006c-proyeccion-posicional}
\end{equation}
The temporal signature is
\[
 \tau(q):=\varepsilon_9(r(\theta))\in\{-1,0,+1\}.
\]
We define the oriented observable shift
\begin{equation}
 a(q):=
 \begin{cases}
  (\delta(d^\Sigma),0),&\tau(q)=+1,\\
  (0,\delta(d^\Pi)),&\tau(q)=-1,\\
  (0,0),&\tau(q)=0.
 \end{cases}
 \label{p12-83-c1-s12:eq:u006c-desplazamiento-observable}
\end{equation}
Then
\[
 \pi_{\rm pos}(F_{\rm obs}q)=\pi_{\rm pos}(q)+a(q),
\]
with the address update and clock defined in U005.

\begin{theorem}[Obstruction to the ordinary observable endofunctor]
\label{p12-83-c1-s12:thm:u006c-obstruccion}
There is no endofunctor of the ordinary observable space that simultaneously satisfies:
\begin{enumerate}
 \item multiplying the geometric position by four;
 \item replace each transition by four genuine microtransitions;
 \item preserving sheet and orientation in each microtransition;
 \item sending a complete duodecaphase fiber to a single
 duodecaphase fiber.
\end{enumerate}
The obstruction persiste in a step of umbral with \(a(q)=0\).
\end{theorem}

\begin{proof}
Let
\[
 \mathcal B_b:=\{(b,r):0\leq r\leq8\}.
\]
Replace one elementary iteration with four force
\[
 h(\theta+1)=h(\theta)+4
\]
and, by induction,
\[
 h(\theta)=4\theta+c\pmod{108}.
\]
Take \(0\leq\widetilde c<108\) and write \(\widetilde c=9c_b+c_0\), with \(0\leq c_0<9\). If the whole
fiber \(\mathcal B_b\) were sent to a single fiber, the destination block index
would have to be independent of \(r\); however, it contains
\begin{equation}
 4b+c_b+
 \left\lfloor\frac{4r+c_0}{9}\right\rfloor
 \pmod{12},
 \label{p12-83-c1-s12:eq:u006c-bloque-obstruccion}
\end{equation}
which is not constant in \(r\). Moreover, the neutral transition operator
satisfies \(P_0^{\rm TPK}=I\), distinct from the cyclotomic
trivializations \(P_b\); the ordinary calendar \((+1,-1,0)\) causes
four elementary iterations to visit both arithmetic modes and the
neutral mode.
\end{proof}

\subsection{Minimal refinement by microinstants}

\begin{definition}[Subdivided observable state]
\label{p12-83-c1-s12:def:u006c-estado-subdividido}
Let
\begin{equation}
 \widetilde{\mathcal Q}_4
 :=\mathcal Q_{\rm obs}\times\{0,1,2,3\},
 \qquad
 \widetilde\pi(q,j)
 :=4\pi_{\rm pos}(q)+j\,a(q).
 \label{p12-83-c1-s12:eq:u006c-estado-refinado}
\end{equation}
We define
\begin{equation}
 \widetilde F(q,j):=
 \begin{cases}
  (q,j+1),&0\leq j<3,\\
  (F_{\rm obs}q,0),&j=3,
 \end{cases}
 \qquad
 \iota_4(q):=(q,0).
 \label{p12-83-c1-s12:eq:u006c-dinamica-refinada}
\end{equation}
\end{definition}

\begin{theorem}[Four-step semiconjugacy]
\label{p12-83-c1-s12:thm:u006c-semiconjugacion}
The inclusion \(\iota_4\) satisfies
\begin{equation}
 \boxed{
 \widetilde F^4\circ\iota_4
 =\iota_4\circ F_{\rm obs}.}
 \label{p12-83-c1-s12:eq:u006c-semiconjugacion}
\end{equation}
The intermediate positions are exactly
\[
 4\pi_{\rm pos}(q)+j\,a(q),
 \qquad j=0,1,2,3.
\]
\end{theorem}

\begin{proof}
Starting from \((q,0)\), the first three applications increment only the
microinstant index. The fourth applies \(F_{\rm obs}\) and returns the index to
zero. This gives
\[
 (q,0)\mapsto(q,1)\mapsto(q,2)\mapsto(q,3)
 \mapsto(F_{\rm obs}q,0).
\]
The formula for the positions is the definition of
\(\widetilde\pi\). The refinement makes visible a determined subdivision;
it does not add a dynamic decision.
\end{proof}

\subsection{Subdivided category and cyclotomic elevation}

Let \(\operatorname{sd}_4\mathsf P_{\rm obs}\) be the category obtained by
replacing each generator
\(\gamma:q\to F_{\rm obs}q\) with the chain
\begin{equation}
 (q,0)\to(q,1)\to(q,2)\to(q,3)\to(F_{\rm obs}q,0).
 \label{p12-83-c1-s12:eq:u006c-cadena-subdivision}
\end{equation}
The functor
\[
 \operatorname{sd}_4:
 \mathsf P_{\rm obs}\longrightarrow
 \operatorname{sd}_4\mathsf P_{\rm obs}
\]
sends each generator to that string. The collapse functor
\[
 \operatorname{col}_4:
 \operatorname{sd}_4\mathsf P_{\rm obs}
 \longrightarrow\mathsf P_{\rm obs}
\]
recompose the four microarrows and satisfies
\begin{equation}
 \operatorname{col}_4\circ\operatorname{sd}_4
 =\operatorname{id}_{\mathsf P_{\rm obs}}.
 \label{p12-83-c1-s12:eq:u006c-colapso}
\end{equation}

In a chain that remains in the fiber \(b\), the first three
microarrows act by \(\rho_b\); the fourth incorporates the interfiber transport
if a dodecaphase crossing occurs. By \(\rho_b^3=I\), the ordered product
of the four actions coincides with the cyclotomic morphism of the
original arrow.

\begin{proposition}[Cyclotomic Invariance under Subdivision]
\label{p12-83-c1-s12:prop:u006c-invariancia-ciclotomica}
There exists a typed lift
\[
 \widetilde{\mathfrak F}_{\rm cyc}:
 \operatorname{sd}_4\mathsf P_{\rm obs}
 \longrightarrow
 \mathbf{Rep}_{\mathbb C}(C_3)\times\mathbf B\mathbb N
\]
such that, for every arrow \(\gamma\) of
\(\mathsf P_{\rm obs}\),
\begin{equation}
 \boxed{
 \widetilde{\mathfrak F}_{\rm cyc}
 \bigl(\operatorname{sd}_4\gamma\bigr)
 =
 \bigl(\mathfrak F_{\rm cyc}\circ Q\bigr)(\gamma).}
 \label{p12-83-c1-s12:eq:u006c-invariancia-ciclotomica}
\end{equation}
\end{proposition}

\begin{proof}
The successive powers of
\(\rho_b\) are assigned to the internal microarrows, and transport between the frames of the
source and target fibers is assigned to the fourth. The composition of the first three factors is the
identity because \(\rho_b^3=I\). The remaining factor is exactly the
morphism that \(\mathfrak F_{\rm cyc}\circ Q\) assigns to the
collapsed arrow. The subdivided category is free on those microarrows, so
the assignment extends uniquely.
\end{proof}

\subsection{\(3\) noninterchangeable quaternary operations}

The architecture contains three distinct operations:
\begin{enumerate}
 \item \(\mathcal S_4\), cardinal substitution of objects and paths;
 \item \(D_4\), multiplication by four in
 \(\Theta_{108}\);
 \item \(\operatorname{sd}_4\), functorial subdivision of an observable
 transition.
\end{enumerate}
Their compatibilities are given, respectively, by
\[
 4^3=1+7\cdot9,\qquad
 \Theta D_4=4\Theta,\qquad
 \widetilde F^4\iota_4=\iota_4F_{\rm obs}.
\]
Sharing the integer four converts none into another.

\begin{criterion}[Refutation criteria]
\label{p12-83-c1-s12:crit:u006c-refutacion}
The construction is refuted if any of the identities
\eqref{p12-83-c1-s12:eq:u006c-d4-total},
\eqref{p12-83-c1-s12:eq:u006c-cociclo-d4},
\eqref{p12-83-c1-s12:eq:u006c-d4-cubo},
\eqref{p12-83-c1-s12:eq:u006c-semiconjugacion}, or
\eqref{p12-83-c1-s12:eq:u006c-invariancia-ciclotomica} fails; likewise if an observable
realization satisfies the four conditions of
\cref{p12-83-c1-s12:thm:u006c-obstruccion} without preserving a register equivalent to the
microinstant.
\end{criterion}


% Unidad U006. Lectores, odómetro y calendario; redacción autónoma.

\section{Modular readers, odometer, and dodecaphasic calendar}
\label{p12-83-c1-s13:chap:tpk-lectores-odometro-calendario}

The integral holonomic TPK state is not identified with any of its readings.
This unit defines the readers modulo nine, modulo three, and modulo one thousand, the
phase--sector odometer, the predictive quotient of thirty-six states, and the
cyclic extension of one hundred and eight positions. Each construction declares what it
publishes and what memory it preserves.

\subsection{Typed scheme of the modular readers of the TPK: domain, codomain, publication, and conserved memory}

\begin{definition}[Operator specification]
\label{p12-83-c1-s13:def:tpk-ficha-operatoria}
A canonical TPK operator is recorded by means of a sextuple
\[
 \mathcal O=(D_{\mathcal O},C_{\mathcal O},f_{\mathcal O},
 I_{\mathcal O},L_{\mathcal O},M_{\mathcal O}),
\]
where \(D_{\mathcal O}\) is the domain, \(C_{\mathcal O}\) the codomain,
\(f_{\mathcal O}\) the rule, \(I_{\mathcal O}\) the preserved invariants,
\(L_{\mathcal O}\) the published or forgotten information, and
\(M_{\mathcal O}\) the memory remaining in the integral holonomic state.
\end{definition}

Two symbols with equal cardinality or equal period do not represent the same
operator unless there exists an equivalence intertwining their domains,
rules, and invariants. This convention will prevent confusing the predictive quotient
\(C_{36}^{\rm pred}\), the boundary \(R_{36}\), and a
cardinal defect of value thirty-six.

\subsection{Nonadic reader of the integral holonomic state: phase, transition, and return with memory}

For \(n\geq1\), define
\[
 \rho_9(n)=1+((n-1)\bmod9),
 \qquad
 q_9(n)=\left\lfloor\frac{n-1}{9}\right\rfloor.
\]
Then
\[
 n=9q_9(n)+\rho_9(n).
\]
The card of \(\rho_9\) is
\[
 \bigl(
 \mathbb Z_{>0},\mathcal D_9,\rho_9,
 n\equiv\rho_9(n)\pmod9,
 \text{positive phase},
 q_9(n)
 \bigr).
\]
The positive digital root numerically coincides with \(\rho_9\), but its
semantics is distinct: one is presented as a section of the quotient; the other,
as an iterated sum of digits.

\subsection{Tritic reader of the oriented state: temporal order and carry transport}

The balanced division is written
\[
 n=r_3(n)+3c_3(n),
 \qquad r_3(n)\in\mathsf{TRIT},\quad c_3(n)\in\mathbb Z.
\]
Its card is
\[
 \bigl(
 \mathbb Z,\mathsf{TRIT}\times\mathbb Z,
 n\mapsto(r_3(n),c_3(n)),
 \operatorname{rec}_3(r,c)=r+3c,
 r_3,
 c_3
 \bigr).
\]
The publication may display only \(r_3\), but the TPK preserves
\(c_3\) if a subsequent composition requires reconstruction.

\subsection{Reader positional by blocks in base \(10^3\)}

The Euclidean division defines
\[
 N=1000q_{1000}(N)+r_{1000}(N),
 \qquad 0\leq r_{1000}(N)<1000.
\]
Its card is
\[
 \bigl(
 \mathbb Z,\mathbb Z\times\{0,\ldots,999\},
 N\mapsto(q_{1000},r_{1000}),
 \operatorname{rec}_{1000}(q,r)=1000q+r,
 r_{1000},
 q_{1000}
 \bigr).
\]

\begin{proposition}[Concatenation and residual reading]
\label{p12-83-c1-s13:prop:tpk-concatenacion-base-mil}
If \(u_1,\ldots,u_m\in\{0,\ldots,999\}\) and
\[
 D_m=1000D_{m-1}+u_m,
 \qquad D_0=0,
\]
then
\[
 D_m=\sum_{j=1}^m u_j1000^{m-j}.
\]
Moreover,
\[
 D_m\equiv\sum_{j=1}^m u_j\pmod9
\]
because \(1000\equiv1\pmod9\).
\end{proposition}

\begin{proof}
The first equality is obtained by induction on \(m\). The second follows
by reducing it modulo nine.
\end{proof}

The congruence does not reconstruct the order of the blocks. The genealogical register preserves the sequence \((u_1,\ldots,u_m)\); the residual reader publishes only its sum modulo nine.

\subsection{Odometer category of the TPK: cylindrical objects and composition by truncation}

\begin{definition}[Nonadic odometer--dodecaphasic]
\label{p12-83-c1-s13:def:tpk-odometro-9-12}
The category \(\mathcal O_{9,12}\) has objects
\[
 (b,r)\in C_{12}\times\{0,\ldots,8\}.
\]
An arrow of length \(n\in\mathbb N\) is
\[
 (b,r)\xrightarrow{\ n\ }
 \left(
  b+\left\lfloor\frac{r+n}{9}\right\rfloor,
  r+n\bmod9
 \right).
\]
The composition sums lengths.
\end{definition}

Let
\[
 \kappa_r(n)=\left\lfloor\frac{r+n}{9}\right\rfloor.
\]

\begin{lemma}[Cocycle of the odometer]
\label{p12-83-c1-s13:lem:tpk-cociclo-odometro}
If \(r_1=r+n_1\bmod9\), then
\[
 \kappa_r(n_1+n_2)
 =\kappa_r(n_1)+\kappa_{r_1}(n_2).
\]
\end{lemma}

\begin{proof}
We write \(r+n_1=9\kappa_r(n_1)+r_1\), with
\(0\leq r_1<9\). Adding \(n_2\) and again dividing by nine
yields the identity.
\end{proof}

For a state \(x\) whose base contains \(\theta\), we define
\[
 Q_0(x)=(\beta(\theta),\widetilde\theta\bmod9).
\]
A TPK arrow has length equal to the number of observable generators that
compose it.

\begin{theorem}[Phase--length functor]
\label{p12-83-c1-s13:thm:tpk-funtor-fase-longitud}
The assignment sending a state \(x\) to \(Q_0(x)\) and an arrow of
length \(n\) to the odometric arrow of length \(n\) defines a functor
\[
 Q:\mathsf{TPK}\longrightarrow\mathcal O_{9,12}.
\]
\end{theorem}

\begin{proof}
Each generator increases \(\theta\) by one unit. After \(n\) steps, the
phase residue increases by \(n\) modulo nine and the sector increases by
\(\lfloor(r+n)/9\rfloor\). The identity has zero length and the
\cref{p12-83-c1-s13:lem:tpk-cociclo-odometro} proves compatibility with composition.
\end{proof}

After nine steps,
\[
 (b,r)\longmapsto(b+1,r).
\]
After one hundred and eight,
\[
 (b,r)\longmapsto(b,r).
\]
These identities belong to the odometric quotient. They do not assert that the
cylinder, boundary, genealogical register, or vertical memory have
returned.

\subsection{Minimum predictive-memory quotient}

Let
\[
 X_{9,12}=\{(a,m):a\in\mathbb Z/9\mathbb Z, 0\leq m<12\}
\]
with successor
\[
 S(a,m)=
 \begin{cases}
  (a,m+1),&m<11,\\
  (a+1\bmod9,0),&m=11.
 \end{cases}
\]
We define the observables
\[
 \beta_{\rm pred}(a,m)=4a\pmod{12},
 \qquad
 d_{\rm pred}(a,m)=1+(m+\beta_{\rm pred}(a,m)\pmod{12}).
\]
For an observable \(q\), the inclusive future word of horizon \(h\)
is
\[
 Q_h^q(x)=(q(x),q(Sx),\ldots,q(S^hx)).
\]

\begin{theorem}[Minimum predictive quotient]
\label{p12-83-c1-s13:thm:tpk-cociente-predictivo-36}
The observables \(\beta_{\rm pred}\), \(d_{\rm pred}\), and the pair
\((\beta_{\rm pred},d_{\rm pred})\) produce the same stable quotient
\[
 \pi_{36}(a,m)=(a\bmod3,m).
\]
It has thirty-six states and fibers of cardinality three. The minimal
horizons are \(11\), \(8\), and \(0\), respectively.
\end{theorem}

\begin{proof}
For a finite system \((X,T)\), let
\[
 E_h=\operatorname{Eq}(Q_h^q).
\]
Then
\[
 E_{h+1}=E_h\cap(T\times T)^{-1}(E_h).
\]
The decreasing chain stabilizes and its limit is the largest
\(T\)-compatible equivalence contained in \(\operatorname{Eq}(q)\). In the declared
system, the exact enumeration produces the following number of classes:
\[
\begin{array}{c|rrrrrrrrrrrr}
h&0&1&2&3&4&5&6&7&8&9&10&11\\\hline
\#(X/E_h^{\beta})&3&6&9&12&15&18&21&24&27&30&33&36
\end{array},
\]
and
\[
\begin{array}{c|rrrrrrrrr}
h&0&1&2&3&4&5&6&7&8\\\hline
\#(X/E_h^{d})&12&15&18&21&24&27&30&33&36.
\end{array}
\]
The pair already distinguishes the thirty-six classes at horizon zero. In all
cases, the stable classes preserve \(m\) and \(a\bmod3\); the three
possibilities for \(a\) with the same residue modulo three form each fiber.
\end{proof}

We denote this object by \(C_{36}^{\rm pred}\). No is the boundary
\(R_{36}\), that will contain three ordered channels of six trits, nor a census
of thirty-six elements of another construction.

\subsection{Central extension of the calendar}

We define
\[
 \iota:C_{12}\longrightarrow C_{108},
 \qquad \iota([b]_{12})=[9b]_{108},
\]
and
\[
 p:C_{108}\longrightarrow C_9,
 \qquad p([\theta]_{108})=[\theta]_9.
\]
The exact sequence is obtained.
\begin{equation}
 0\longrightarrow C_{12}
 \xrightarrow{\ \iota\ }C_{108}
 \xrightarrow{\ p\ }C_9
 \longrightarrow0.
\label{p12-83-c1-s13:eq:tpk-extension-12-108-9}
\end{equation}
A section of sets \(s([r]_9)=[r]_{108}\),
\(0\leq r<9\), induces the cocycle
\[
 c(r,r')=
 \left[\left\lfloor\frac{r+r'}{9}\right\rfloor\right]_{12}.
\]

\begin{proposition}[The extension does not split]
\label{p12-83-c1-s13:prop:tpk-extension-no-escindida}
The sucesion \cref{p12-83-c1-s13:eq:tpk-extension-12-108-9} not is isomorfa to the extension
product \(C_{12}\times C_9\).
\end{proposition}

\begin{proof}
The exponent of \(C_{108}\) is \(108\). The exponent of
\(C_{12}\times C_9\) is
\(\operatorname{mcm}(12,9)=36\). Two isomorphic finite groups have the
same exponent; therefore, they are not isomorphic.
\end{proof}

\begin{theorem}[Necessity and exact period of one hundred eight positions]
\label{p12-83-c1-s13:thm:tpk-periodo-exacto-108}
The translation \(\theta\mapsto\theta+1\) on \(C_{108}\) has order
exactly \(108\). Consequently, an iteration \(n\) returns the
complete calendar if and only if \(108\mid n\).
\end{theorem}

\begin{proof}
The \(n\)-th power sends \([\theta]\) to \([\theta+n]_{108}\). It is
the identity for every \(\theta\) exactly when \([n]_{108}=0\), that is,
when \(108\mid n\).
\end{proof}

In particular, \(1008\) is not another return of the calendar:
\[
 1008\equiv36\pmod{108}.
\]
After \(1008\) steps, the phase modulo nine returns, but the coordinate of
\(C_{108}\) has shifted by thirty-six positions. Naturally,
\(216,324,\ldots\) also return the calendar, but they are repetitions of
the fundamental circuit of length \(108\).

\begingroup\fontsize{10.2}{11.7}\selectfont
\setlength{\parskip}{2pt}\setlength{\abovedisplayskip}{5pt}\setlength{\belowdisplayskip}{5pt}
\subsection{Unbounded vertical memory}

Let \(q_{\rm mem}\in\mathbb Z_{\geq0}\) be the unbounded number of nonadic
turns, and let
\[
 \beta=[q_{\rm mem}]_{12}.
\]
Then a nine-step turn satisfies
\[
 q_{\rm mem}\longmapsto q_{\rm mem}+1,
 \qquad
 \beta\longmapsto\beta+1\pmod{12}.
\]
After one hundred eight steps,
\[
 \beta\longmapsto\beta,
 \qquad
 q_{\rm mem}\longmapsto q_{\rm mem}+12.
\]
The calendar returns; the vertical memory does not.

\begin{corollary}[Hierarchy of returns]
\label{p12-83-c1-s13:cor:tpk-jerarquia-retornos}
In the phase quotient,
\(9\) steps constitute one turn. The composite morphisms
\[
 \mathcal K_{27}=F_{\rm obs}^{27},
 \qquad
 F_{\rm obs}^{54},
 \qquad
 F_{\rm obs}^{108}
\]
contain, respectively, three, six, and twelve nonadic turns. None
erases the genealogical record or the boundary by definition.
\end{corollary}

\subsection{Minimal type system of the integral holonomic state and compatibilities between its readers}

The following table answers explicitly which internal objects
already exist in the TPK before introducing seeds:
\[
\begin{array}{p{.23\textwidth}p{.27\textwidth}p{.38\textwidth}}
\toprule
Objeto u operador&Dominio y codominio&Función\\\midrule
\(X_9\)&\((\mathbb Z/9)^2\)&soporte APP\\
\(T_N,T_E,T_S,T_O\)&\(X_9\to X_9\)&traslaciones cardinales\\
\(M_{\rm ph}\)&\(\mathcal D_9\to\mathsf{TRIT}\)&selección de hoja\\
\(F_{\rm obs}\)&\(\mathcal Q_{\rm obs}\to\mathcal Q_{\rm obs}\)&sucesor observable\\
\((\rho_9,q_9)\)&\(\mathbb Z_{>0}\to\mathcal D_9\times\mathbb Z_{\ge0}\)&fase y altura\\
\((r_3,c_3)\)&\(\mathbb Z\to\mathsf{TRIT}\times\mathbb Z\)&orientación y carry\\
\((q_{1000},r_{1000})\)&\(\mathbb Z\to\mathbb Z\times\{0,\ldots,999\}\)&bloque decimal y cociente\\
\(Q\)&\(\mathsf{TPK}\to\mathcal O_{9,12}\)&fase y longitud\\
\(C_{36}^{\rm pred}\)&cociente de \(X_{9,12}\)&memoria predictiva mínima\\
\(C_{108}\)&calendario cíclico&extensión no escindida\\
\(q_{\rm mem}\)&\(\mathbb Z_{\ge0}\)&memoria vertical no acotada\\
\(\operatorname{Ext}_r\)&
\(\operatorname{Ext}_r\subseteq\mathcal F_q\times\mathcal F_{q'}\)&
prolongaciones compatibles\\
\bottomrule
\end{array}
\]

\subsection{Obstructions of the readers and the TPK calendar: Euclidean quotients, \(C_{36}^{\mathrm{pred}}\), extension \(C_{108}\), and vertical memory}

\begin{criterion}[Reader Refutation Criteria and Calendar]
The unit is refuted if:
\begin{enumerate}
 \item a residual reader is presented as a complete transition;
 \item the quotient of some subsequently used Euclidean division is omitted;
 \item \(C_{36}^{\rm pred}\) is identified with \(R_{36}\);
 \item it is asserted that the extension \(C_{108}\) is the product
       \(C_{12}\times C_9\);
 \item complete return is declared for a number not divisible by \(108\);
 \item the return of \(\beta\) fuerza \(q_{\rm mem}=0\);
 \item \(\mathcal K_{27}\) is presented as a distinct clock and not as
       a power of the transport.
\end{enumerate}
\end{criterion}

\section{Complete operatorial structure of APP and arithmetic-geometric hierarchy}

The operatorial development of APP fixes the hierarchy
\(4\to2\to6\to54\), the arithmetic projectors, the central block, and the
local geometry of paths. These consequences remain subordinated to the
already constructed formal kernel and do not introduce a second origin of APP.

The next unit will construct the quadruple subdivision, the
cyclotomic functor, and the internal channels of ninety and one hundred twenty positions.
Only afterward will the exhaustive space of seeds be introduced.
 

\par\endgroup
\chapter{Operator architecture of selection, transport, and update}
\section{Categorical structure of TPK: integral holonomic states, transport
morphisms, and operator classification}
\label{sec:arbol-operatorio-tpk-rev11}
\index{TPK!operational classification}
\index{phase selector!suboperator of the TPK}

The TPK is the categorical container that transports over APP the
oriented states of TRIT. Its definition precedes the inventory of realizations:
first objects, morphisms, composition, and operatorial classes are fixed; only
afterward are the operators that realize each class introduced.

\HMTClaim{HMT-R-TPK-ENRICHED-STATE}
\begin{definition}[TPK as a category over the observable base with fiber family]
\label{def:estado-enriquecido-tpk-rev11}
Let \(\mathsf P_{\rm obs}\) be the category of observable paths over the Cayley
graph of APP. For each configuration \(q\) the fiber is fixed
\[
 \mathcal F_q=
 \mathsf O_q\times\mathsf H_q\times\mathsf C_q\times
 \mathsf S_q\times\mathsf B_q\times\mathsf\Lambda_q,
\]
whose factors retain, respectively, orientation and sheet; residue and
quotient; carry; signature; cylinder and boundary; and transport register. A
TPK object is a pair \(x=(q,f)\), with \(f\in\mathcal F_q\). A morphism
\(\Gamma:x\to y\) is a route \(\gamma:q\to q'\) of
\(\mathsf P_{\rm obs}\) together with the fiber transport that preserves the
declared extension relations between \(f\) and \(f'\). The identity is the
empty route with identity transport, and composition is the concatenation of
routes with composition of the fiber transports. The projection
\[
 p:\mathcal X_{\rm TPK}^{\rm enr}\longrightarrow\mathsf P_{\rm obs}
\]
sends each state to its observable configuration and each morphism to the path that
transports. This definition fixes a category over the base and a family of
fibers; the existence is not presupposed the existence of Cartesian lifts for every
flecha of \(\mathsf P_{\rm obs}\).
\end{definition}

\begin{definition}[Ontological--operational classification of the TPK]
\label{def:categorias-operatorias-tpk-rev11}
The internal structure of the TPK is organized into eight classes distinguished by type of
domain, codomain, and function:
\begin{enumerate}
 \item \textbf{Support and state.} The
base is \(\mathsf P_{\rm obs}\) and the carrier is \(\mathcal X_{\rm TPK}^{\rm enr}\to\mathsf P_{\rm obs}\). This class contains
objects; the next seven contain maps between objects or readers of those
objects.
 \item \textbf{Internal selection.} The selector
 \[
 M_{\rm ph}:\{1,\ldots,9\}\longrightarrow\{+1,-1,0\}
 \]
 activates additive evaluation, multiplicative evaluation, or the threshold passage. It acts
 on the leaf coordinate and does not by itself modify the position in APP.
 \item \textbf{Transport.} For each cardinal direction \(d\),
 \(T_d:X_9\to X_9\), \(p\mapsto p+\delta(d)\), it transports the position. The
 evolution \(F_{\rm obs}:\mathcal Q_{\rm obs}\to\mathcal Q_{\rm obs}\)
 composes selection, translation, orientation, and phase advance.
 \item \textbf{Readers and quotients.} \(\rho_9\) publishes the nonadic class;
 \((r_3,c_3)\) conserves tritic residue and quotient; and
 \((q_{1000},r_{1000})\) separates decimal block and carry. They are
 reading maps, not state transitions.
 \item \textbf{Memory and boundary.} The transports of
 \(\mathsf H_q,\mathsf C_q,\mathsf S_q,\mathsf B_q\) and
 \(\mathsf\Lambda_q\) update quotient, carry, signature, cylinder, boundary,
 and register. Their visible projection may coincide even though the integral
 holonomic states differ.
 \item \textbf{Periodicity and return.} The extension
 \[
 0\longrightarrow C_{12}\longrightarrow C_{108}
 \longrightarrow C_9\longrightarrow0
 \]
 types sector, calendar, and visible phase. The nonadic return restores the
 phase; the periodicity of one hundred eight iterations restores the calendar;
 neither erases
 by definition the boundary memory. The loop
 \(\mathcal K_{27}=F_{\rm obs}^{27}\) is a composite morphism, not another
 clock.
 \item \textbf{Prolongation.} The relations
 \(\operatorname{Ext}_r\subseteq\mathcal F_q\times\mathcal F_{q'}\)
 select compatible continuations. The cylinders and their inverse system
 publish finite survival and, under the stated compatibility
 hypotheses, a global section.
 \item \textbf{Outputs.} The evaluation functors contract the same integral
 holonomic state toward the Archimedean output, the exceptional incidence
 output, or the atlas classifiers. These outputs are reading
 codomains; they do not constitute new containers alongside TPK.
\end{enumerate}
\end{definition}

Each operator appearing below is determined by the sextuple \((D,C,f,I,L,M)\): domain \(D\), codomain \(C\), rule \(f\), preserved invariants \(I\), published or lost information \(L\), and retained memory \(M\). Equality of cardinalities, periods, or historical names does not identify two operators with distinct sextuples. Once this classification has been fixed, the following diagram summarizes five functional groupings of realizations that refine the eight operator classes of TPK:
\[
\begin{tikzpicture}[
  node distance=7mm and 9mm,
  every node/.style={font=\small,align=center},
  root/.style={draw=hmtblue,very thick,rounded corners,fill=hmtlightblue,
    inner sep=5pt},
  op/.style={draw=hmtblue,rounded corners,fill=white,inner sep=4pt},
  sub/.style={draw=hmtgray,rounded corners,fill=hmtpaper,inner sep=3pt},
  edge/.style={-{Latex[length=2mm]},semithick,hmtblue}
]
\node[root] (tpk) {TPK\\$\mathcal X_{\rm TPK}^{\rm enr}$};
\node[op,below left=of tpk] (sel) {selección\\de lectura};
\node[op,below=of tpk] (trans) {transporte\\y retorno};
\node[op,below right=of tpk] (lift) {elevación\\y memoria};
\node[sub,below=of sel] (mov) {$M_{\rm ph}$};
\node[sub,below=of trans] (card) {$T_d,\ \mathcal K_{27},\ F_{\rm obs}$};
\node[sub,below=of lift] (fib) {$\mathsf H,\mathsf C,\mathsf S,
  \mathsf B,\mathsf\Lambda$};
\node[op,below=17mm of trans] (clock) {calendario y registro\\
  $C_{12}\hookrightarrow C_{108}\twoheadrightarrow C_9$, $\Theta_{108}$};
\node[op,below=of clock] (ext) {prolongación, cilindros, frontera\\
  y filtración de supervivencia};
\draw[edge] (tpk) -- (sel);
\draw[edge] (tpk) -- (trans);
\draw[edge] (tpk) -- (lift);
\draw[edge] (sel) -- (mov);
\draw[edge] (trans) -- (card);
\draw[edge] (lift) -- (fib);
\draw[edge] (mov) |- (clock);
\draw[edge] (card) -- (clock);
\draw[edge] (fib) |- (clock);
\draw[edge] (clock) -- (ext);
\end{tikzpicture}
\]
The diagram establishes a containment, not a succession of theories. In particular,
the internal selector, cardinal translations, stroboscopic
observation, and modular readers preserve the types fixed in
\cref{def:categorias-operatorias-tpk-rev11}. TPK coordinates their compositions
and preserves their provenance in the same integral holonomic state.

\HMTClaim{HMT-R-TPK-OPERATOR-SEPARATION}
\begin{proposition}[Separation of selection, motion, and remark]
\label{prop:tpk-separacion-operadores-rev11}
The selector \(M_{\rm ph}\), cardinal transport \(T_d\), and sampling
\(\operatorname{Obs}_3\) are noninterchangeable morphisms: the first chooses
the active evaluation; the second changes position and winding; the
third publishes a temporal subsequence. Their composition within the TPK
determines an observable route, while the fiber \(\mathcal F_q\) preserves
the data required to distinguish prolongations with the same endpoint.
\end{proposition}

\begin{proof}
The three morphisms have different codomains: \(M_{\rm ph}\) takes values
in the ternary activation set, \(T_d\) is an automorphism of the torus, and
\(\operatorname{Obs}_3\) acts on sequences. An identification between any two
of them would violate the domain or codomain typing. The extension by fibers
further incorporates residue--quotient, carry, boundary, and register, so
that projection to the endpoint of the path is an explicit forgetful map.
\end{proof}

\begin{statusbox}[title={Canonical nomenclature}]
The public name of the container is \emph{TPK}. When it is necessary
to name the transported data, \emph{integral holonomic TPK state} or
\(\mathcal X_{\rm TPK}^{\rm enr}\) will be used. The morphism \(M_{\rm ph}\) designates
exclusively internal phase selection.
\end{statusbox}
 \section{Fundamental operators of the TPK: domains, actions, invariants, and memory}
\label{sec:inventario-operatorio-tpk-rev11}
\index{TPK!operational inventory}

The \cref{def:categorias-operatorias-tpk-rev11} first fixed the eight
ontological--operatorial classes: support and state; selection; transport;
readers and quotients; memory and boundary; periodicity and return;
prolongation; and outputs. This section defines their realizations. For each
operator, the sextuple \((D,C,f,I,L,M)\) of domain, codomain, rule,
invariants, published or lost information, and preserved memory is declared. The subsequent
nominal census groups symbols already defined by these six data and does not
replace their domains or rules.

\begin{definition}[Typed Record for a TPK Operator]
A canonical operator is a sextuple
\[
 \mathcal O=(D_{\mathcal O},C_{\mathcal O},f_{\mathcal O},
 I_{\mathcal O},L_{\mathcal O},M_{\mathcal O}),
\]
where \(f_{\mathcal O}:D_{\mathcal O}\to C_{\mathcal O}\) is the action
rule, \(I_{\mathcal O}\) is the family of relations it preserves,
\(L_{\mathcal O}\) declares the published coordinates and the equivalence
relation induced by forgetting, and \(M_{\mathcal O}\) enumerates the data
remaining in the integral holonomic state after publication. Two
realizations represent the same operator only when there is a
typed equivalence intertwining their rules and invariants; sharing a
period, cardinality, or visible output is not enough.
\end{definition}

\subsection{Modular readers built}

The nonadic reduction is fixed by the published representative.
\begin{equation}
 \rho_9(n)=1+((n-1)\bmod 9)
 =n-9\left\lfloor\frac{n-1}{9}\right\rfloor,
 \label{eq:rho9-construida-tpk-rev11}
\end{equation}
so that the null class is written as \(9\). For \(n\geq1\), the
digital root satisfies
\[
 \operatorname{dr}_9(n)=\rho_9(n),
\]
but the two expressions preserve distinct semantics: the first is a
section of the quotient; the second publishes the iterated sum of digits. The tritic
lift writes uniquely
\begin{equation}
 n=r_3(n)+3c_3(n),
 \qquad r_3(n)\in\{-1,0,+1\},\quad c_3(n)\in\mathbb Z,
 \label{eq:lift-mod3-tpk-rev11}
\end{equation}
and simultaneously preserves oriented residue and carry quotient.
Finally, the decimal-block reader is the Euclidean division
\begin{equation}
 N=1000q_{1000}(N)+r_{1000}(N),
 \qquad 0\leq r_{1000}(N)<1000,
 \label{eq:lector-mod1000-tpk-rev11}
\end{equation}
therefore basis \(1000\), residue modulo \(1000\) and carry \(q_{1000}\)
are three data related and not a same object.

\subsection{Classification of Operator Realizations by Domain and Function}

The preceding definitions govern the reading of the table. Its fourteen
functional groupings refine the eight operator classes and summarize the
realizations that share a function; a class may therefore occupy more
than one row. The table does not operate as a replacement definition. Symbols that
share a cardinal remain separated by domain and codomain.

\begingroup
\footnotesize
\renewcommand{\arraystretch}{1.06}
\begin{longtable}{@{}p{.21\textwidth}p{.31\textwidth}p{.40\textwidth}@{}}
\toprule
Functional grouping & Operators and Spaces & Domain, Action and Invariants \\
\midrule
\endfirsthead
\toprule
Grouping functional & Operators and Spaces & Domain, Action and Invariants \\
\midrule
\endhead
State and route memory &
\(\mathcal X_{\rm TPK}^{\rm enr}\), fiber \(\mathcal F_q\), route memory &
They preserve position, leaf, orientation, residue, quotient, carry, signature,
boundary, cylinder, prefix, genealogy, and memory. \\

Sheet selection &
\(M_{\rm ph}:\{1,\ldots,9\}\to\{-1,0,+1\}\) &
It activates the additive evaluation, the multiplicative evaluation, or the threshold passage; it does not
by itself move the state. \\

Transporte espacial &
\(T_N,T_E,T_S,T_O\), reversible two-cursor engine,
\(F_{\rm obs}\) &
It realizes the translations of the Cayley graph, updates sheet and orientation and
preserves each transition of the path. \\

Lectores modulares &
\(\rho_9\), \((r_3,c_3)\), \((q_{1000},r_{1000})\) &
Publish nonadic phase, tritic orientation and decimal block without erasing the lift quotients. \\

Block and memory &
\(\mathcal K_{27}=F_{\rm obs}^{27}\), memory filter
\(\mathcal N_{27}\) &
The first composes twenty-seven transports; the second selects the
memory events. Equal cardinality does not imply an equal operator. \\

Twelve-phase state &
\(\mathcal R_{12}\), centers \(\theta_m=30m-10\) degrees,
refinement \(\mathcal R_{120}\) &
Organize the twelve phases, their opposition, and the angular refinement in increments of
three degrees. \\

Periodicity and depths &
\(C_{12}\hookrightarrow C_{108}\twoheadrightarrow C_9\),
\(\Pi_{108}^{\rm cur}\), \(w_{108}\), \(w_{1080}\),
\(\Gamma_0(1080)\) &
It distinguishes the total periodicity of order \(108\), the cursor projection, the
words of 108 and 1080 trits, and modular level 1080. \\

Channels (90/120) &
mask \(C_m\), \(S_{90},S_{120}\), orbital grammar,
supertrace, and invariants \(\nu_{120},\nu_{270}\) &
It separates the two genealogies, their geometric tails, their weights, and the
coordinates that they extract from the state. \\

Bidirectionality &
\(P_\pm\), route inversion, \(\operatorname{Ext}_\pm\),
\(N_\pm\), \(\beta\), \(C_M=-\operatorname{rev}\) &
It implements, respectively, spectral exchange, spatial outward passage and return,
future/past prolongation, valence, balance, and word conjugation. \\

Emission and coinduction &
\(L_0,L_1\), \(R_{36}\), \(G_9\), partition \(3^6=729\),
\(\varprojlim\operatorname{Cyl}_m\) &
It lifts the seed, opens the boundary, and exhaustively enumerates the
finite prolongations. The numerical limit and the enriched lift
belong to distinct types. \\

Numeric Readers &
\(\mathscr C_\pi\), \(\mathscr P_e\), \(F_{\rm av}\) &
Through rational forks, factorial recurrence, and Fibonacci
self-scaling, they construct the published sections of \(\pi,e,\varphi\). \\

Exceptional Projection &
Witt vector lifting, weight enumerator, and non-selective shadow &
Transports the same state to ternary incidence without intervening in the Archimedean selection of digits. \\

Atlas and realization &
route signature, atlas (81/56/13/324), involution \(C_M\), electronic
center, and fiber operators &
Convert the route state into classifiers, characters, and operators before
applying any dimensional chart. \\

Finite validation &
exhaustive censuses, ablations, and causal independence &
It checks coverage and uniqueness at the computed depth, and verifies that
no target prefix intervenes in the generating map. \\
\bottomrule
\end{longtable}
\endgroup

\Needspace{46\baselineskip}
\subsection{Functional Index of the Canonical Operators}

The preceding classification organizes the realizations by function. The following table
gathers the operators after the categorical definition and the constructed
readers. Each row declares domain, codomain, rule, and invariants; the
symbols make it possible to recognize a single definition throughout its
different realizations without identifying operators of different types.

\begingroup
\footnotesize
\renewcommand{\arraystretch}{1.06}
\enlargethispage{2\baselineskip}
\begin{longtable}{@{}r >{\raggedright\arraybackslash}p{.42\textwidth}
  >{\raggedright\arraybackslash}p{.22\textwidth}
  >{\raggedright\arraybackslash}p{.22\textwidth}@{}}
\toprule
\# & Canonical Operator & Symbol & Capa tipada \\
\midrule
\endfirsthead
\toprule
\# & Canonical Operator & Symbol & Capa tipada \\
\midrule
\endhead
1 & Additive--multiplicative arithmetic cell of order nine
  & \(\mathcal A_9\) & arithmetic support \\
2 & Residual decomposition and nonadic quotient
  & \((\rho_9,q_9)\) & reduction and elevation \\
3 & Ternary oriented reduction
  & \(\rho_3^{\rm or}\) & reduction and elevation \\
4 & Ternary-decimal positional reader of base thousand
  & \(\iota_{1000}\) & reduction and elevation \\
5 & Cardinal Action Oriented on the Arithmetic Torus
  & \(T_d,\ d\in\{N,E,S,O\}\) & transporte cardinal \\
6 & Additive--multiplicative local feedback
  & \(F_{\rm loc}\) & transporte cardinal \\
7 & Minimum emitter of two cursors
  & \(T_{\min}\) & finite emission \\
8 & Decimal Arithmetic Sign Coupler
  & \(G\) & finite emission \\
9 & Ternary projection of the catalog
  & \(\Psi_6\) & finite emission \\
10 & Local Surjectivity of Cardinal Transport
  & \(\mathcal R_3\) & transporte cardinal \\
11 & Transportation loop of twenty-seven iterations
  & \(\mathcal K_{27}\) & transporte cardinal \\
12 & Integral holonomic state of the topological phase kernel
  & \(\mathcal X_{\rm TPK}^{\rm enr}\) & integral holonomic state \\
13 & Hensel State Transition
  & \(\mathcal H\) & integral holonomic state \\
14 & Exact Transition of Ternary-Decimal Cylinders
  & \(\mathcal C_{3,1000}\) & integral holonomic state \\
15 & Joint boundary signature
  & \(\Sigma_{\rm B}\) & integral holonomic state \\
16 & Relationship of extension and nonadic survival
  & \({\rm EF}\text{--}G_9\) & prolongation and monodromy \\
17 & Nonadic monodromy with boundary holonomy
  & \(\mathcal M_9\) & prolongation and monodromy \\
18 & Global Section of the Inverse Survivor System
  & \(X_\chi=\varprojlim_m{\rm Surv}^{(\chi)}_m\)
  & prolongation and monodromy \\
19 & Bivariant Category of Arithmetic Paths
  & \(\mathcal P_{\rm APP}^{(2)}\) & transporte cardinal \\
20 & Observable chronology of two cursors
  & \(F_{\rm obs}\) & transporte cardinal \\
21 & Decimal Alphabet of Forty-Two Trinities
  & \(\mathcal A_{42}\) & finite emission \\
22 & Typed extension relation of fibers
  & \({\rm Ext}_r\) & integral holonomic state \\
23 & Nonadic transfer of compatible fibers
  & \(K_{\rm ph}\) & prolongation and monodromy \\
24 & Refinement of Seeds by Advancement and Rotation
  & \((P,H,Q)\) & transporte cardinal \\
25 & Mode and Orientation Cocycles
  & \((c_{\rm mode},c_{\rm or})\) & integral holonomic state \\
26 & Triadic prolongation of twenty blocks
  & \(E_{20}\) & prolongation and monodromy \\
27 & Dodecaphonic oriented system
  & \(\mathcal R_{12}\) & dodecaphonic reading \\
28 & Signed harmonic dodecaphonic observable
  & \(U_{12}^{\rm sgn}\) & dodecaphonic reading \\
29 & Harmonic transformation by three-step orbits
  & \(\mathcal H_{12}\) & dodecaphonic reading \\
30 & Cyclic difference transversal extractor
  & \(\mathcal E_{3,4}\) & dodecaphonic reading \\
31 & Dodecaphasic closure recurrence in base thousand
  & \(\mathcal C_\alpha\) & close in base thousand \\
32 & Enhanced history transducer to angular character
  & \((L_{\rm tick},\chi_{\rm mass})\) & interfaz dimensional \\
33 & Bidiirectional lattice of index twelve
  & \(\mathcal W_{\pm}\) & interfaz dimensional \\*
34 & Nonadic Atlas of Extension Families
  & \(\mathcal A_{\rm species}\) & interfaz dimensional \\
\bottomrule
\end{longtable}
\endgroup

The symbol \(L_{\rm tick}\) denotes the iteration and return reader. Its
codomain is the temporal memory of the state; it does not define an independent physical
time unit.

The internal selector \(M_{\rm ph}\) belongs to the TPK\@. Cardinal transport
\(T_d\) occupies another factor of the transition \(\mathcal U_t\) and effects the
motion on APP\@. This separation preserves a single carrier,
distinguishes phase selection from transport, and preserves the complete architecture
by types.

\Needspace{9\baselineskip}
\begin{resultbox}[title={Finite depth and enriched limit}]
The enumeration, censuses, and contraction of numerical cylinders are
exact results in their domains. The existence and uniqueness of an enriched
section jointly subject to all survival predicates
follow when the lifting fibers are nonempty, singletons, and
compatible at every depth. A calculation on a finite window
verifies the fibers in that window, but does not by itself imply the
quantifier over all depths.
\end{resultbox}

\begin{theorem}[Composition of the internal dynamics]
\label{thm:composicion-dinamica-tpk-rev11}
A transition observable of the TPK is the composition typed
\[
 \mathcal U_t
 =\operatorname{Rec}_t\circ\operatorname{Mem}_t
 \circ T_{d(t)}\circ M_{\rm ph}(g(t))
\]
on \(\mathcal X_{\rm TPK}^{\rm enr}\), followed, when appropriate, by
a modular reader, a cylindrical extension, or one of the two projections.
Each factor acts on a distinct component of the state; for that reason the
selector, the movement, the clock, the memory, and the publication cannot
collapse into a single modular reduction.
\end{theorem}

The preceding table enumerates the operator families. Equivalences among
local realizations must intertwine their domains, rules, and invariants.
Subsequent developments select compositions of these operations;
they introduce no fourth primitive object alongside APP, TRIT, and TPK.
 \chapter{Transport, selection, and prolongation in the Topological Phase Kernel}
\label{sec:arquitectura-tpk}

The acronym TPK designates the \emph{Topological Phase Kernel}, or the topological phase
kernel.\footnote{The denomination likewise preserves the
authorial dedication to Tan Peng Kian.} It does not denote a single 108-step algorithm, a table of
results, or the catalog of 468 emissions. It is the system of
operators that enables the additive and multiplicative readings to coexist
over the APP path category without identifying their
genealogies.

The word \emph{topological} has a precise content here. The support
of APP is the Cayley graph of the quotient group
\[
 X_9=(\mathbb Z/9\mathbb Z)^2
\]
with its cardinal edges. A closed route admits a winding class
after being lifted to cover \(\mathbb Z^2\). TPK transports that class,
the route orientation, and phase holonomy. Two paths ending
in the same cell may therefore represent distinct states. The phase kernel
is precisely the information that remains invisible when publishing
only the final cell or its residue.

The dynamical basis is the free category \(\mathsf P_{\rm APP}\) of the
cardinal graph of APP\@. Since the two arithmetic evaluations do not induce the
same partition, transport resides on the bivariate category
\[
 \mathsf P_{\rm APP}^{(2)}
 =\mathsf P_{\rm APP}\times\mathsf P_{\rm APP}.
\]
The first component is evaluated through the additive sheet \(\Sigma\) and the
second through the multiplicative sheet \(\Pi\). This separation is
constitutive, but it does not duplicate APP: it records two ordered readings of one
and the same matrix of paths. The multiplicative component is the principal table;
the additive component preserves the law of accumulation that generates its rows and makes it possible
to measure what information a visible projection eliminates.

\ifdefined\HMTTPKRepeatAPP
\section{Canonical chart and sum-product architecture}

The support is written in cursor coordinates.
\[
 I_9=\{0,\ldots,8\},\qquad X_9=I_9^2,
\]
but its arithmetic labels belong to \(\{1,\ldots,9\}\). The chart
that relates both levels is
\[
 \chi:I_9\longrightarrow\{1,\ldots,9\},\qquad \chi(i)=i+1.
\]
Let
\[
 \rho_9(n)=1+((n-1)\bmod9),\qquad
 q_9(n)=\frac{n-\rho_9(n)}9.
\]
In particular, \(9\) is the visible representative of the null class. The two
observables of the same vertex are
\[
 \Sigma(i,j)=\rho_9\bigl(\chi(i)+\chi(j)\bigr),\qquad
 \Pi(i,j)=\rho_9\bigl(\chi(i)\chi(j)\bigr).
\]
This convention prevents the zero cursor coordinate from being mixed with the visible
label nine. APP is the tuple
\[
 \mathcal A_9=
 \bigl(X_9,\operatorname{Cay}(X_9),\Sigma,\Pi,
       \operatorname{Path}(X_9)\bigr),
\]
and not the juxtaposition of two tables: each seed \((i,j)\) simultaneously bears
the two evaluations and each route preserves the order in which it consults them.

The elevation that retains residue and quotient is
\[
 \mathcal A_9^\#=(R^+,Q^+;R^\times,Q^\times),
\]
with, for \(a=\chi(i)\) and \(b=\chi(j)\),
\[
 a+b=R^+_{ij}+9Q^+_{ij},\qquad
 ab=R^\times_{ij}+9Q^\times_{ij},
\]
\[
 R^+_{ij}=\Sigma(i,j),\qquad R^\times_{ij}=\Pi(i,j).
\]
The exact totals are
\[
 \sum R^+=405,\quad \sum R^\times=459,\quad
 \sum Q^+=45,\quad \sum Q^\times=174.
\]
Therefore,
\[
 2025-810=1215=(459-405)+9(174-45)=54+9\cdot129.
\]
The equality does not create the TPK by simple counting: it identifies the datum that the
transport must preserve. The visible residue differs by \(54\) and the quotient
memory by \(129\).

The composition of the additive sheet has the carry cocycle
\[
 \kappa_9(a,b)=
 \frac{\rho_9(a)+\rho_9(b)-\rho_9(a+b)}9,
\]
which satisfies
\[
 \kappa_9(a,b)+\kappa_9(a+b,c)
 =\kappa_9(b,c)+\kappa_9(a,b+c).
\]
The multiplicative sector also possesses the radical
\[
 \mathfrak m=(3)=\{3,6,9\},\qquad \mathfrak m^2=0
 \quad\text{in }\mathbb Z/9\mathbb Z.
\]
The six unit rows sum to \(45\), the rows \(3,6\) sum to \(54\), and
the row \(9\) sums to \(81\); thus
\[
 459=6\cdot45+2\cdot54+81,
\]
and all the excess \(459-405=54\) is localized in the nonunital sector. This
localization, the carry cocycle, and the quotient difference are three
distinct data of the same sum--product architecture.
\fi

\section{Residual reductions and decimal representation}

The architecture uses three reductions, each with a distinct purpose:
\[
 \mathbb Z\xrightarrow{(\rho_9,q_9)}
 \{1,\ldots,9\}\times\mathbb Z,
 \qquad
 \rho_3^{\rm or}:\mathbb Z/3\mathbb Z\longrightarrow
 \{-1,0,+1\},
\]
\[
 (d_1,d_2,\ldots)\in\{0,\ldots,999\}^{\mathbb N}
 \xrightarrow{\iota_{1000}}
 \sum_{m\geq1}d_m1000^{-m}.
\]
The ternary orientation is
\[
 \rho_3^{\rm or}([0])=0,\qquad
 \rho_3^{\rm or}([1])=+1,\qquad
 \rho_3^{\rm or}([2])=-1.
\]
The first separates the visible residue from the nonadic quotient; the second assigns the oriented type of the transition; the third groups the Archimedean reading into decimal triads, whose register preserves the carry. None can replace another.
In particular, reduction modulo three does not reconstruct the quotient modulo nine, and a decimal triad is not a six-trit word.
Addition accumulates displacements; multiplication folds classes and contracts the noninvertible rows. The need for TPK arises because the two readings share a support but do not share a partition, period, or quotient.

\Needspace{36\baselineskip}
\section{Correspondence between TPK symbols, domains, and compositions}
\label{sec:diccionario-interno-tpk}

Historical names are fixed here by object, not by numerical similarity. This
table forms part of the public definition of TPK.

\begin{longtable}{@{}>{\raggedright\arraybackslash}p{.19\textwidth}
>{\raggedright\arraybackslash}p{.30\textwidth}
>{\raggedright\arraybackslash}p{.43\textwidth}@{}}
\toprule
Name&Object&Function within the core\\\midrule
\endhead
nonadic coordinate&
\(\rho_9(n)\) together with \(q_9(n)\)&
Separa the representative visible modulo \(9\) of the quotient that preserves the
memory of the cruce of boundary. No is a root digital.\\
ternary temporal selector&
\(M_{\rm ph}(t)=\varepsilon(t)\in\{+1,-1,0\}\)&
At \(1,4,7\) it activates additive accumulation; at \(2,5,8\) it activates
multiplicative accumulation; at \(3,6,9\) it maintains both cursors and records the threshold.
It does not denote the spatial directions.\\
transporte cardinal&
\(T_N,T_E,T_S,T_O\)&
It shifts the position on the APP torus and preserves the path word and its
winding. It is independent of the temporal selector.\\
ternary reader&
\(\rho_3^{\rm or}\)&
It converts the three residual classes into orientation
\(0,+1,-1\). It does not by itself reconstruct the nonadic quotient.\\
reader in base thousand&
\(\iota_{1000}\)&
It groups the Archimedean publication into decimal triads and transports their carries. It is a positional reader, not a congruence ``modulo \(1000\)''.\\
return of twenty-seven&
\(\mathcal K_{27}=F_{\rm obs}^{27}\)&
It composes twenty-seven positional and oriented updates. Its fourth iterate restores the observable phase, while the complete register preserves the accumulated memory.\\
dodecaphonic register&
\(\Lambda\in\mathcal R_{12}\)&
It preserves six positions, five differences, and one total. It is the support of the
reversible transport \(1+5+6=12\).\\
bidirectional reading&
\(q_+,q_-\) and its conjugate weights&
It simultaneously preserves the outward and return orientations. In Dimensional
Mechanics, a signature \((n_A,s_C)\) first produces
\(W_\pm=6n_A\pm s_C\); only afterward does
\(s_C/6\) appear in the angular character.\\
\bottomrule
\end{longtable}

The temporal selector, cardinal transport, and ternary orientation are therefore three distinct maps. Their separation preserves the difference between \emph{when} a reading is activated, \emph{in which direction} the route advances, and \emph{with which orientation} the step is recorded. TPK contains all three and preserves their relations.

\ifdefined\HMTTPKWordArithmetic
\section{From the local cell to the nonadic word monoid}

The sum--product architecture does not end at the local table. To prove that
the transported quotient generates global arithmetic, one must construct the
operation on words without introducing as a definition the integer multiplication
that is to be recovered. Let
\[
 \mathcal D_9=\{1,\ldots,9\},\qquad
 \mathcal W_9=\{\varnothing\}\sqcup
 \bigsqcup_{\ell\geq1}\mathcal D_9^\ell .
\]
The words are written from the least significant position and evaluated
by
\[
 \operatorname{ev}_9(u_0,\ldots,u_{\ell-1})
 =\sum_{j=0}^{\ell-1}u_j9^j,\qquad
 \operatorname{ev}_9(\varnothing)=0.
\]
The symbol \(9\) preserves the boundary crossing; it is not replaced by a zero.

\begin{theorem}[Normal form nonadic]
The evaluation \(\operatorname{ev}_9\) is a bijection between
\(\mathcal W_9\) and \(\mathbb N_0\).
\end{theorem}
\begin{proof}
The words of length \(\ell\) exactly cover the interval
\[
 \frac{9^\ell-1}{8}
 \leq\operatorname{ev}_9(u)
 \leq9\frac{9^\ell-1}{8}.
\]
The minimum of the next interval is one unit greater than the maximum of the
preceding interval. For \(n>0\), the division of \(n-1\) by nine determines uniquely
\[
 d=1+((n-1)\bmod9)\in\mathcal D_9
\]
and a strictly smaller quotient, to which the same rule is recursively
applied. The process terminates and is reversible.
\end{proof}
For example,
\[
 9\leftrightarrow(9),\quad
 10\leftrightarrow(1,1),\quad
 81\leftrightarrow(9,8),\quad
 1000\leftrightarrow(1,3,3,1).
\]
The last identity displays in a single normal form the two fundamental
partitions
\[
 1000=729+271=729+243+27+1;
\]
its interpretation by channels is introduced thereafter, not in the definition
of the numbering.

\section{Sum, Horner product and quotient preservation}

The native sum \(u\boxplus_\#v\) traverses the words by positions. At each
position it applies the additive cell to the two digits present and, if there exists
an incoming carry, applies a second additive cell to the provisional residue. The
residue is published and the sum of quotients passes to the next position. The
partial invariant is
\[
 \sum_{j<k}(u_j+v_j)9^j
 =\sum_{j<k}w_j9^j+c_k9^k,
\]
from which it is obtained
\[
 \operatorname{ev}_9(u\boxplus_\#v)
 =\operatorname{ev}_9(u)+\operatorname{ev}_9(v).
\]

For \(d\in\mathcal D_9\), the scalar transducer \(d\odot_\#u\) first uses
\[
 u_jd=9q_j+r_j,\qquad
 (r_j,q_j)=\bigl(R^\times(u_j,d),Q^\times(u_j,d)\bigr),
\]
and normalizes the collision between \(r_j\) and the incoming quotient by means of the
additive cell. Thus,
\[
 \operatorname{ev}_9(d\odot_\#u)
 =d\operatorname{ev}_9(u).
\]
If \(v=(v_0,\ldots,v_{m-1})\), it is defined, from right to left,
\[
 a_m=\varnothing,\qquad
 a_j=(9\odot_\#a_{j+1})\boxplus_\#(v_j\odot_\#u),
\]
and
\[
 u\boxtimes_\#v:=a_0.
\]
This is a Horner construction performed only with
the two local cells and their quotients.

This description fixes the operatorial dependence, but it does not constitute a
second residence of the result. The global product, its Horner proof, and
the identity
\(operatorname{ev}_9(u\boxtimes_\#v)
=\operatorname{ev}_9(u)\operatorname{ev}_9(v)\)
are established only once in the
\cref{eq:ev-multiplicative}, within the specialized chapter on native
product.
The shortest boundary example is
\[
 (9)\boxtimes_\#(9)=(9,8),\qquad 9+8\cdot9=81.
\]
Preserving only \(R^\times\) would destroy this theorem: \(2\cdot5=10\)
would be reduced to the visible residue \(1\). The quotient is not an auxiliary
annotation but the memory needed for the product to survive passage through
the boundary.

\section{Coproduct and Irreducibility}

The multiplicative opening is defined internally by
\[
 \Delta_\#e_w
 =\sum_{u\boxtimes_\#v=w}e_u\otimes e_v.
\]
Its reduced version eliminates the factors \((1)\). If
\(\widetilde\Delta_\#\) denotes that version and
\[
 A_\#=
 \bigl(\overline{\widetilde\Delta_\#}\bigr)^*
 \overline{\widetilde\Delta_\#},
\]
then \(A_\#e_w=d_\#(w)e_w\), where \(d_\#(w)\) counts the pairs ordenados
not triviales that produce \(w\). The selector
\[
 P_{\mathrm{Irr},\#}
 =(I-P_{(1)})\,\mathbf1_{\{0\}}(A_\#)
\]
chooses exactly the words distinct from the unit without nontrivial native
opening. Therefore, irreducibility is read in the product fiber itself;
it is not introduced by an external table of primes.

This result is exact in the monoid of normal forms. Its extension to
complete histories requires additional compatibility. If \(P_n\) projects
onto surviving histories of depth \(n\), a family of
injections \(\iota_n\) must satisfy
\[
 P_n\iota_n\mu_\#
 =P_n\mu_n(\iota_n\otimes\iota_n).
\]
The preceding theorems construct \(\mu_\#\) and fix this equation; by
themselves they do not assert the existence or uniqueness of the family
\((\iota_n)_n\). This separation prevents the global arithmetic of
words, already proved, from being confused with its monoidal prolongation to the deep
TPK\@ memory space.

\fi

\section{Cardinal Action and Arithmetic Feedback}

Let
\[
 \nu_N=(-1,0),\quad \nu_E=(0,1),\quad
 \nu_S=(1,0),\quad \nu_O=(0,-1)
\]
the four generators cardinal. For \(d\in\{N,E,S,O\}\),
\[
 T_d(i,j)=(i,j)+\nu_d\pmod9.
\]
The letter \(O\) means west. These directions are displacements over
APP and not names of the channels \(\pi,e,\varphi\).

A word \(w=d_1\cdots d_r\) determines
\[
 \gamma(x_0;w)_k=x_0+\sum_{j=1}^{k}\nu_{d_j}\pmod9.
\]
When it closes, its lift preserves
\[
 \operatorname{wind}(w)=\frac1{9}
 \bigl(\#S-\#N,\#E-\#O\bigr)\in\mathbb Z^2.
\]
The bidirectional reading preserves word, orientation, and sheet. Its
formal reversal inverts the order and exchanges
\(N\leftrightarrow S\), \(E\leftrightarrow O\). Its eventual reading as
advance--delay is a later physical interface, not part of this
route definition.

The locally observable state is
\[
 s_t=(i_t,j_t,m_t),
 \qquad m_t\in\{\Sigma,\Pi\}.
\]
The active sheet is evaluated,
\[
 a_t=\rho_3^{\rm or}\bigl(m_t(i_t,j_t)\bigr),
\]
and the mode is updated via
\[
 m_{t+1}=
 \begin{cases}
 \Sigma,&a_t=+1,\\
 \Pi,&a_t=-1,\\
 m_t,&a_t=0.
 \end{cases}
\]
Thus, each cardinal step modifies both the position and the future
arithmetic reading. The path is not superimposed on a previously constructed word: it is
the positional history that produces it.

\begin{theorem}[Surjectivity of the sampled three-step factor]
For each of the \(162=9^2\cdot2\) states \((i,j,m)\) and each target
oriented trit \(b\in\{-1,0,+1\}\), there exists a route
\(u\in\{N,E,S,O\}^3\) whose third step emits \(b\). In the \(486\) state--target
cases, the number of admissible routes is between
\(8\) and \(40\). Consequently, for every target succession
\((b_k)_{k\geq1}\), there exists a cardinal history whose output satisfies
\[
 (a_3,a_6,a_9,\ldots)=(b_1,b_2,b_3,\ldots).
\]
The two intermediate trits \(a_{3k-2}\) and \(a_{3k-1}\) of each block are not
prescribed by this assertion.
\end{theorem}
\begin{proof}
The \(4^3=64\) three-step routes from each state are enumerated. The included exhaustive certificate records, for each of the \(162\cdot3\) cases, the route executed, the final state, and the total number of witnesses; its indices \(0,1,2\) are transported by \(\rho_3^{\rm or}\) to \(0,+1,-1\). All counts are positive; the minimum is \(8\) and the maximum is \(40\). After each block, the same result is applied to the state reached.
The choice, for example, of the first witness in lexicographic order recursively produces a history for any target sequence.
\end{proof}

Surjectivity concerns the factor observed every three steps, not the
word emitted at unit time. This theorem establishes that sampled capacity
of the route support. The selection of the sections
\(\pi,e,\varphi\) belongs to the enriched EF--G9 relation defined
later: capacity for realization and selection that is genealogical are two
distinct operators.

\section{Observable chronology of the two cursors}

The calendar is formulated over
\[
 \Theta_{108}=\mathbb Z/108\mathbb Z,\qquad
 \beta(\theta)=
 \left[\left\lfloor\frac{\widetilde\theta}{9}\right\rfloor\right]_{12},
 \qquad
 r(\theta)=1+(\theta\bmod9),
\]
where \(\widetilde\theta\in\{0,\ldots,107\}\). The active signature is
\[
 \varepsilon(r)=
 \begin{cases}
 +1,&r\in\{1,4,7\},\\
 -1,&r\in\{2,5,8\},\\
 0,&r\in\{3,6,9\}.
 \end{cases}
\]
The observable space of two cursors is
\[
 \mathcal Q_{\rm obs}
 =X_9^2\times\{N,E,S,O\}^2\times\Theta_{108}.
\]
If \(\varepsilon=+1\), the additive cursor advances; if
\(\varepsilon=-1\), the multiplicative cursor advances; if \(\varepsilon=0\), neither advances.
The two orientations are reversed after steps
\(27,54,81,108\). This rule defines a permutation
\[
 F_{\rm obs}:\mathcal Q_{\rm obs}\longrightarrow\mathcal Q_{\rm obs}.
\]
Positions and orientations are restored after \(54\) steps and the complete
phase after \(108\):
\[
 F_{\rm obs}^{108}=I.
\]
Therefore, the chronology of a single route is reversible and does not by itself
select a distribution on the continuations; selection of all
compatible prolongations belongs to the fiber transfer operator
defined below.

\section{Finite two-cursor emitter}

A cursor seed is
\[
 \mathcal S=\mathbb T_9^2\times\{N,E,S,O\},
 \qquad |\mathcal S|=324.
\]
Minimal TPK receives an additive cursor \(P\in\mathcal S\) and a
multiplicative cursor \(M\in\mathcal S\). Their six-window signatures are
\(s(P)\) and \(e(M)\). In each coordinate one applies
\[
 G(s,e)=100s+10\bigl((s+e)\bmod10\bigr)
       +\bigl((3s+5e+1)\bmod10\bigr).
\]
The figure of hundreds recovers \(s\); known \(s\), the of tens recovers
\(e\). Therefore \(G\) is injective on the effective signatures. The emitter
factorizes as
\[
 \mathcal S^2\xrightarrow{s\times e}
 \operatorname{im}s\times\operatorname{im}e
 \xrightarrow{G}\mathcal U_6,
\]
with
\[
 |\mathcal S^2|=104\,976,\qquad
 |\operatorname{im}s|=18,\qquad
 |\operatorname{im}e|=26,\qquad
 |\mathcal U_6|=18\cdot26=468.
\]
The factorization \(468=18\cdot26\) is therefore a decomposition of the
object: each emission conserves an additive and a multiplicative signature.

The projection
\[
 \Psi_6:\mathcal U_6\longrightarrow\mathbb F_3^6,
 \qquad \Psi_6(U)=-U\pmod3,
\]
has an image of \(243\) words. The chain typed is
\[
 104\,976\ \text{seeds}
 \longrightarrow468\ \text{decimal emissions }U_6
 \longrightarrow243\ \text{ternary words }w_6.
\]
The ambient space \(\mathbb F_3^6\) contains \(729\) words; it is not identified
with the effective image of \(243\). Nor are \(U_6\) and \(w_6\) the same
type.

The emitter minimo repite, in the chart dodecaphasic of catalogo,
\[
 U_{12}^{\rm cat}=U_6\Vert U_6.
\]
Its visible period six explains why it cannot replace the signed state
\(U_{12}^{\rm sgn}\). This difference is not an insufficiency of TPK: it separates
the catalogue projection from the enriched history.

Each emitted decimal block belongs to an exact alphabet of
\(7\cdot6=42\) triads. If \(\overline{d_1d_2d_3}\) is one of them,
\[
 \boxed{d_3\equiv5d_2-2d_1+1\pmod{10}.}
\]
The sweep of the \(104\,976\) seeds produces no triad outside this
law. It is therefore useful to distinguish the complete chain of quotients:
\[
 \mathcal S^2
 \longrightarrow\Omega_{\rm seed}
 \longrightarrow\mathcal U_6
 \xrightarrow{\Psi_6}\operatorname{im}\Psi_6\subset\mathbb F_3^6.
\]
The microscopic presentations, the \(468\) decimal emissions, the
\(243\) ternary words, and the environment of \(729\) words are four
distinct objects.

\paragraph{Orbital Commutativity of Typed Morphisms.}
Let
\(W_{243}:=\operatorname{im}\Psi_6\subset\mathbb F_3^6\). The complete funnel and its dihedral quotient are represented without turning inclusions into selections:
\[
 \underbrace{\mathcal S^2}_{104\,976\ \mathrm{semillas}}
 \xrightarrow[\text{fiber census}]{\text{emission}}
 \underbrace{\mathcal U_6}_{468\ U_6}
 \xrightarrow[\text{visible projection}]{\Psi_6}
 \underbrace{W_{243}}_{243\ w_6}
 \hookrightarrow
 \underbrace{\mathbb F_3^6}_{729\ \mathrm{palabras}},
 \qquad
 W_{243}\xrightarrow{\,/D_3^{\rm cat}\,}
 \underbrace{\mathcal O_{43}}_{43\ \text{orbits}}.
 \tag{TPK-embudo}\label{eq:embudo-orbital-tipado-rev8}
\]
The first arrow groups initial conditions by the emission they produce;
the second forgets the decimal presentation; the third is an inclusion in
the algebraic ambient space; the last is a quotient by the action
\(D_3^{\rm cat}\) on \(W_{243}\). In particular, there is no implicit quotient
\(729\to43\). The exact census is
\[
 243=38\cdot6+5\cdot3,
\]
that is, thirty-eight orbits of size six and five of size three.
These \(43\) orbital classes are not the \(43\) complete nonadic turns
per channel of the thousand-digit G9 witness either: the first \(43\) belongs
to the finite quotient \(W_{243}/D_3^{\rm cat}\); the second is a depth
counter of a certified run.

\section{Configuration of the integral holonomic TPK state: sheet, orientation, residue, quotient, carry, boundary, path, and memory}

The complete dynamic unit is a history, not an isolated word. At
depth \(m\) we write
\[
 v_m=
 \bigl(
 B_m,\boldsymbol\delta_m,\boldsymbol x_m,
 \boldsymbol{\mathcal C}_m,\Sigma_m,
 \omega_m,\ell_m,g(m),\Lambda_m
 \bigr),
\]
where
\[
 B_m=(u_m^\pi,u_m^e,u_m^\varphi)\in(\mathbb F_3^6)^3.
\]
Their components have distinct functions:
\begin{enumerate}
 \item \(B_m\) is the joint visible block of the three channels;
 \item \(\boldsymbol\delta_m\) retains the new decimal triads and the
       associated carry;
 \item \(\boldsymbol x_m\) is the deep Hensel state;
 \item \(\boldsymbol{\mathcal C}_m\) contains the ternary and
       decimal cylinders and their integer residues;
 \item \(\Sigma_m\) is the signature joint of boundary;
 \item \(\omega_m\) and \(\ell_m\) conserve orientation and bidirectional sheet;
 \item \(g(m)\in\mathbb Z/9\mathbb Z\) is the nonadic phase;
 \item \(\Lambda_m\) is the dodecaphasic record oriented.
\end{enumerate}

\begin{definition}[Operatorial presentation and TPK category]
The operator inventory of the integral holonomic TPK state is
\[
 \mathfrak T_{\rm HMT}=
 \bigl(
 \mathcal A_9,\TRIT,T_{N,E,S,O},T_{\min},
 H,\mathcal C_{3,1000},\Sigma_B,EF,\mathcal G_9,
 \mathcal R_{12}
 \bigr)
\]
which transports histories \((v_0,\ldots,v_m)\) and their truncations. Its
formal carrier is the category \(\mathsf{TPK}\) over the observable base:
its objects are integral holonomic states and its morphisms are the prolongations
\(\operatorname{Ext}_r\) stable under identity and concatenation, defined
below. The inventory does not replace that category. The visible projection
eliminates components of \(v_m\); it does not define an equivalence with the
complete state.
\end{definition}

This definition contains exactly what disappears when
cursors are reset or only the residue is conserved: the path, the sheet, the quotient, the
carry, the orientation, and the boundary. TPK is, for that reason, the point at which
modular arithmetic ceases to be an operation without history.

\section{Ternary-decimal cylinders and state of extension}

Let \(N\) be the integer determined by a ternary prefix of length
\(L\), and let \(D\) be the integer determined by
\(K\) decimal triads. We define
\[
 A=N1000^K-D3^L,
 \qquad
 B=(N+1)1000^K-D3^L.
\]
The cylinders intersect exactly when
\[
 \boxed{A<3^L,\qquad B>0.}
\]
When incorporating a ternary block \(u\in\{0,\ldots,728\}\),
\[
 N'=729N+u.
\]
If no appears a new triad decimal,
\[
 A'=729A+u1000^K.
\]
Si aparece \(\delta\in\{0,\ldots,999\}\),
\[
 D'=1000D+\delta,
\]
\[
 A'=729000A+u1000^{K+1}-\delta\,729\,3^L.
\]
In both cases,
\[
 B'=A'+1000^{K+\Delta K}.
\]
These integer recurrences transport the decimal datum without turning it into an
external target of the update.

The Hensel state preserves the information that does not appear in the residue. The
active matrix is
\[
A_H=\begin{pmatrix}
2&2&2&1&2&1\\2&1&2&2&1&1\\1&0&0&1&0&2\\
0&0&1&1&1&0\\2&2&2&0&2&1\\2&1&2&1&0&0
\end{pmatrix},\qquad \det A_H=1,
\]
and the residue--quotient division is
\[
 y=xA_H,\qquad u=y\bmod3,\qquad x^+=\frac{y-u}{3}.
\]
For the three channels, the joint cocycles are additionally formed
\[
 q=(\pi+e+\varphi)\bmod3,
 \qquad
 c=\frac{\pi+e+\varphi-q}{3},
\]
\[
 a=(\pi+e-\varphi)\bmod3,
 \qquad
 h=\frac{\pi+e-\varphi-a}{3}.
\]
The second-order update corrects the next state with \(c\) and \(h\). Two branches with the same visible block are therefore not necessarily the same state: they may differ in their Hensel extension and in their cylinders. Since \(\det A_H=1\), the reduction of \(A_H\) belongs to \(\operatorname{GL}_6(\mathbb F_3)\). Hence, before cylinders, boundaries, or survival are imposed, the residual Hensel chart covers exactly \(\mathbb F_3^6\). This raw coverage does not imply coverage of the filtered system.

Formally, over each observable configuration \(q\) a fiber is fixed
\[
 \mathcal F_q=
 \mathsf O_q\times\mathsf H_q\times\mathsf C_q\times
 \mathsf S_q\times\mathsf B_q\times\mathsf\Lambda_q,
\]
corresponding, respectively, to orientation and sheet, residue and quotient,
carry, signature, cylinder and boundary, and transport record. For a
cardinal arrow \(r:q\to q'\), the admissible prolongation does not reduce to a
function on residues; it is the relation
\[
 \operatorname{Ext}_r\subseteq\mathcal F_q\times\mathcal F_{q'},
 \qquad
 \operatorname{Ext}_{r_2}\circ\operatorname{Ext}_{r_1}
 \subseteq\operatorname{Ext}_{r_2\circ r_1}.
\]
The carry increment is a cocycle:
\[
 \kappa(r_2\circ r_1)
 =\kappa(r_2)+\mathsf C(r_2)\bigl(\kappa(r_1)\bigr).
\]
Thus the forgetful map is typed: projecting an extension to its
observable route eliminates components of \(\mathcal F_q\), but does not authorize
identifying both objects.

\section{Compound returns and transport calendars}

The length of a trace does not determine its type. The academic presentation
separates the composite return, calendar, trace, and atlas according to their
respective preserved states.

The symbol \(\mathcal K_{27}\) is reserved for the composition of
twenty-seven updates of a positional and oriented transport loop.
The integer \(27\) is its composition length; this does not turn it into a
trace of 108 iterations or into an infinite emitter. Another construction, historical
in origin, groups four loops whose joint periods reach 108; that joint rotor
is not denoted by \(\mathcal K_{27}\) in isolation.

\begin{longtable}{@{}>{\raggedright\arraybackslash}p{.24\textwidth}>{\raggedright\arraybackslash}p{.31\textwidth}>{\raggedright\arraybackslash}p{.37\textwidth}@{}}
\toprule
Object&Estado conservado&Significado de 108\\\midrule
\endhead
Contador simplificado&position and velocity&its real period is 18; it is a
historical control.\\
Rotor set of four loops&four routes, reading, rotation, and motion&
the individual periods are \(54,54,36,36\); the joint state has
period 108.\\
Raw clocked trace&four positions, orientations, and phase&the dynamic content repeats at 54.\\
Trace of \(T_{\min}\)&two positions and two directions&the state returns to
54 and the output is \(U_6\Vert U_6\).\\
Reinitialized pseudocode&cursors reset each window&collapses to
\((222,222,222,333,333,333)^2\); it does not represent the integral holonomic TPK state.\\
Rutas condicionadas&path and prescribed word&108 is horizon, not period.\\
Sectorial odometer&\(\mathbb Z/9\mathbb Z\times\{0,\ldots,11\}\)&the
complete state has period 108; the sectorial observable, 36.\\
Projection raw&quotient by forgetfulness of \(X_{\rm TPK}^{\rm enr}\)&does not possess
by itself the signed reading.\\
Traza enriquecida&all the fields of \(v_m\) during 108 updates&is
the temporal restriction of the integral holonomic TPK state.\\
Atlas of Routes&81 cells/56 signatures or 324 oriented seeds&is a taxonomy,
not an orbit of 108 states.\\
\bottomrule
\end{longtable}

The loop \(\mathcal K_{27}\) is the finite transport operator of
twenty-seven updates. The windows, joint rotor, and calendrical projection
receive distinct symbols, so that the shared length
does not replace the domain or composition of each object.

\section{Fiber transfer and spin memory}

The factorization \(\mathcal U_6\simeq\mathcal S_+\times\mathcal S_\times\)
allows one to define, with respect to the uniform count, the conditional expectations
\[
 (E_+f)(s,e)=\frac1{18}\sum_{s'\in\mathcal S_+}f(s',e),
 \qquad
 (E_\times f)(s,e)=\frac1{26}
 \sum_{e'\in\mathcal S_\times}f(s,e').
\]
\(E_+\) renews the additive fiber and preserves the multiplicative one;
\(E_\times\) realizes the dual operation. They are the only Markov projectors
that completely marginalize the active coordinate, preserve the inactive one, and
preserve the uniform count.

Let \(\tau=(+1,-1,0)^3\) be the phase word on \(C_9\), and
\(P_{+1}=E_+\), \(P_{-1}=E_\times\), \(P_0=I\). The nonadic transfer is
\[
 \mathcal K_{\rm ph}\bigl((u,r),(v,r+1)\bigr)
 =P_{\tau(r)}(u,v).
\]
It is doubly stochastic, irreducible, of period \(9\), and has a unique
stationary distribution \(1/(9\cdot468)\). Moreover,
\[
 \mathcal K_{\rm ph}^{\,9}=E_+E_\times,\qquad
 (E_+E_\times)(u,v)=\frac1{468}.
\]
This equality expresses the simultaneous renewal of all compatible
continuations, not the deterministic successor of a single path.

Three geometric operations on the seeds preserve additional
information: \(P\), advance in the active orientation; \(H\), a half-turn; and
\(Q\), a quarter-turn. The pair \((P,H)\) refines the \(26\) product
signatures into \(28\) classes. In particular, the signature \(555555\) splits into three sheets
of eight representatives, and
\[
 18\cdot28=504=468+36.
\]
The term \(36\) counts additional sheets of the refinement and is not the
object \(R_{36}\). The pair \((P,Q)\) produces \(162\) classes of two
elements, transitively permuted by
\[
 J(i,j,d)=\bigl(7-i,\,7-j,\,\overline d\bigr)\pmod9.
\]

In a trace of \(108\) steps, there appear \(36\) signs of each type,
\(72\) effective changes between the two sheets, and four orientation
reversals; the total signed cocycle is zero. These counts are invariants
of the evolution record and link nonadic dynamics with the dodecaphase record.

\section{Tertiary prolongation and joint boundary}

Relative to the named seeds, boundaries, and readers of the published
profile, the initial visible words are
\[
 w_6^\pi=010211,\qquad w_6^e=201101,\qquad
 w_6^\varphi=121200.
\]
After five blocks, the prefixes are
\[
\begin{aligned}
w_{30}^{\pi}&=010211|012222|010211|002111|110221,\\
w_{30}^{e}&=201101|121221|102011|012222|102011,\\
w_{30}^{\varphi}&=121200|112202|121020|010210|010200.
\end{aligned}
\]
The first deep boundary is
\[
 R_{36}=
 \begin{pmatrix}222220\\021222\\102011\end{pmatrix}.
\]
For each channel the prefix
\(w_{36}^\chi:=w_{30}^\chi|u_5^\chi\) is defined; \(R_{36}\) is the ordered triple of
the three sixth blocks. Neither is an emitter nor a terminal state.

For \(B_k=(u_k^\pi,u_k^e,u_k^\varphi)\), the signatures
\[
 q_k=u_k^\pi+u_k^e+u_k^\varphi,\qquad
 a_k=u_k^\pi+u_k^e-u_k^\varphi
\]
determine
\[
 u_k^\varphi=2(q_k-a_k),\qquad
 u_k^\pi+u_k^e=q_k-u_k^\varphi
 \quad\text{in }\mathbb F_3^6.
\]
Thus, \(\varphi\) is the fixed axis and the residual freedom is the
internal distribution \(\pi/e\). In the ninth phase,
\[
 B_9=(100100|020112|010122)
\]
admits three local outputs with the same axis \(\varphi\) and the same aggregate
\(\pi+e\); all three survive one horizon and only
\[
 B_{10}=(101222|112221|112002)
\]
survives the second. This is the executable\(3\to1\) descent of EF--G9.
The published branch continues through twenty blocks per
channel: \(w_{30}\) and \(R_{36}\) are prefix and boundary, not
closure.

\section{Relation of extension and nonadic monodromy}
\label{sec:extension-monodromia-tpk}

Let \(\mathcal H_m\) be the set of admissible enriched histories of
depth \(m\), and let
\[
 p_{n,m}:\mathcal H_n\longrightarrow\mathcal H_m
\]
the truncation. For \(N\geq m\),
\[
 \operatorname{Surv}_m^{(N)}=p_{N,m}(\mathcal H_N),
 \qquad
 \operatorname{Surv}_m=\bigcap_{N\geq m}\operatorname{Surv}_m^{(N)}.
\]
The EF--G9 relation conserves the extensions that belong to this stable core.
The phase satisfies
\[
 g(m+9)=g(m),
\]
but the transported state satisfies, in general,
\[
 v_{m+9}\ne v_m.
\]
The return acts on the fiber of states of a same phase; does not reset the
history.

\begin{definition}[Number HMT]
\label{def:numero-hmt-tpk-completo}
An HMT number of channel \(\chi\) is solely a compatible global section
\[
 x=(x_m)_{m\geq0}\in
 X_\chi:=\varprojlim_m\operatorname{Surv}_m^{(\chi)}.
\]
The reader and the published value are not part of the section.
\end{definition}

A presentation or measurement of the number is a pair \((x,L)\), where
\(L\) is a declared reader whose domain contains \(x\). When \(L\) is the
Archimedean reader of cylinders, \(L(x)\) is the published value; \(x\)
conserves in addition orientation, quotient, carry, and boundary.

\begin{theorem}[Existence of a global prolongation]
\label{thm:prolongacion-global-tpk-completo}
Let \(h\in\mathcal H_m\). If the prolongation tree of \(h\) is finitely branching and has at least one node at every depth, then there exists an infinite history whose truncations prolong \(h\). If, in addition, each stable level is a singleton and the unique states are compatible, the global section is unique.
\end{theorem}
\begin{proof}
The non-empty finite levels form a finitely ramified tree. König's
lemma provides an infinite branch. The compatibility of the
truncations identifies that branch with an element of the inverse limit. If
each level contains a single state, there is no second compatible branch.
\end{proof}

The preceding theorem is abstract and conditional: by itself it asserts neither nonemptiness nor the stable uniqueness of a concrete channel. For \(\pi,e,\varphi\), those hypotheses are subsequently discharged through the executable rational filtration of \cref{thm:generacion-monodromica-conjunta-rev7}. This placement preserves here only the inverse-limit principle and does not anticipate its application.

\section{Composition of the fundamental operators}

The preceding operators lead to the terminal holonomic state without
flattening its coordinates:
\[
 \begin{aligned}
 \mathrm{APP}&\longrightarrow\mathrm{TRIT}
 \longrightarrow\text{oriented route}
 \longrightarrow\text{integral holonomic state}\\
 &\longrightarrow\text{cylinders and survival}
 \longrightarrow x_{\rm term}.
 \end{aligned}
\]
The dodecaphonic closure then receives two coordinated evaluations of the same
state:
\[
 x_{\rm term}
 \xrightarrow{(\pi_{\rm term},R_{\rm sgn})}
 \bigl(B_{\rm term},U_{12}^{\rm sgn}\bigr),
 \qquad
 U_{12}^{\rm sgn}\xleftrightarrow{\mathcal H_{12}}K,
\]
\[
 (P,E,\Phi,K,c_{13}=0)\longleftrightarrow\alpha_{12}.
\]
No arrow \(B_{\rm term}\to U_{12}^{\rm sgn}\) appears: both
coordinates follow from \(x_{\rm term}\). The subsequent passage to Dimensional
Mechanics likewise creates no other machine; it prolongs the preserved genealogy
through an enriched chronological record, signature, character, and towers.

\ifdefined\HMTInternalTraceability
\section{Academic record and traceability}

The file \path{datos/registro_operadores_tpk.json} forms part of the package.
For each operator it declares an academic name, a historical alias, domain,
codomain, transition, invariants, status, and local source. It also includes
the separation of the ten historical constructions associated with the number
108. The publication gate checks this register and blocks an output
that again reduces TPK to its catalog or to a raw trace.

In the remainder of the article, four already typed abbreviations will be used:
\[
 w_6\in\mathbb F_3^6,\qquad
 w_{30}\in\mathbb F_3^{30},\qquad
 R_{36}\in(\mathbb F_3^6)^3,\qquad
 U_{12}^{\rm sgn}\in\mathbb Z^{12}.
\]
The first is a seed word; the second, a five-block prefix; the third, the
joint entry boundary into EF--G9; the fourth, the declared signed coordinate
of the enriched dodecaphasic state. No cardinal equality
authorizes their identification.
\fi
  \part{Holonomic path algebra and local electronic structure}

\chapter{Universal presentation of the holonomic algebra}
\section{Generating data}

Fix a depth \(n\), and let \(X_n\) be the finite set of
APP--TRIT--TPK states preserving, at the corresponding level, sheet, regime,
orientation, residue, quotient, carry, boundary, height, path, memory, and
survival. Let \(\mathcal C_n\) be the small category generated by the steps

\[
U_t={\rm Upd}_t\circ{\rm Tra}_t\circ{\rm Sel}_t.
\]

The objects of \(\mathcal C_n\) are the states of \(X_n\); its arrows are the
admissible histories. Composition exists when the terminal state of the
first history coincides with the initial state of the second.

Each \(x\in X_n\) is associated with the tritic fiber

\[
A_x=A_{\tau(x)}
=\mathbb R[J_x]/(J_x^2+\tau(x)e_x),
\qquad \tau(x)\in\{+1,0,-1\}.
\]

The symbol \(e_x\) designates the local identity of the fiber. For each arrow
\(u:x\to y\), TPK transport induces an action

\[
\rho_u:A_x\longrightarrow A_y
\]

on the coefficients surviving the step. When a transition changes
regime, \(\rho_u\) is interpreted as the typed bimodule of that
transition, not as an identification of the two quadratic algebras.

The concatenation memory register defines a multiplier

\[
\Omega(v,u)\in Z(A_z)^\times,
\qquad
u:x\to y,\qquad v:y\to z,
\]

that records the carry, torsion, and boundary contribution produced
by composition. Its compatibility law is

\[
\rho_w(\Omega(v,u))\Omega(w,v\circ u)
=\Omega(w,v)\Omega(w\circ v,u).
\tag{1.1}
\]

Equation (1.1) is the cocycle identity that prevents the memory of a
history from depending on how its composition is parenthesized.

\section{Presentation by generators and relations}

Let \(\mathcal F_n\) be the free associative algebra generated by

\[
\{e_x,J_x:x\in X_n\}\cup\{T_u:u\in{\rm Mor}(\mathcal C_n)\}.
\]

We define the ideal \(\mathcal I_{\rm HMT}\) through the relations

\[
e_xe_y=\delta_{xy}e_x,
\qquad
\sum_{x\in X_n}e_x=1_n,
\tag{2.1}
\]

\[
J_x=e_xJ_xe_x,
\qquad
J_x^2=-\tau(x)e_x,
\tag{2.2}
\]

\[
T_u=e_{t(u)}T_ue_{s(u)},
\tag{2.3}
\]

\[
T_ua=\rho_u(a)T_u,
\qquad a\in A_{s(u)},
\tag{2.4}
\]

and

\[
T_vT_u=\Omega(v,u)T_{v\circ u}
\tag{2.5}
\]

when \(v\circ u\) exists; the product is zero when the arrows are not
composable. The finite holonomic algebra is

\[
\boxed{
\mathfrak H_n=\mathcal F_n/\mathcal I_{\rm HMT}.}
\tag{2.6}
\]

Addition in \(\mathfrak H_n\) superposes possible paths. Multiplication concatenates
histories. The factor \(\Omega\) prevents concatenation from erasing the carry.
Thus, addition and multiplication have acquired a meaning that is simultaneously
arithmetic and dynamic.

\section{Associativity theorem}

\textbf{Theorem 3.1.} If \(\rho\) respects local identities and \(\Omega\)
satisfies (1.1), the product defined by (2.5) is associative.

\textbf{Proof.} Let \(u:x\to y\), \(v:y\to z\), and \(w:z\to t\). For
coefficients \(a,b,c\) in the relevant terminal fibers,

\[
\begin{aligned}
((cT_w)(bT_v))(aT_u)
={}&c\rho_w(b)\Omega(w,v)
\rho_{w\circ v}(a)\Omega(w\circ v,u)T_{wvu},\\
(cT_w)((bT_v)(aT_u))
={}&c\rho_w(b\rho_v(a)\Omega(v,u))
\Omega(w,v\circ u)T_{wvu}.
\end{aligned}
\]

The covariance of \(\rho\) equates the factors containing \(a\) and \(b\).
Identity (1.1) equates the two products of \(\Omega\). Associativity
on monomials follows, and bilinearity extends it to all of
\(\mathfrak H_n\). \(\square\)

The proof clarifies the role of memory. Without \(\Omega\),
concatenation would preserve the bare path, but not the accumulated discrepancy
between its lifts. With \(\Omega\), two ways of grouping the same
history produce the same complete state.

\section{Universal representation property}

\textbf{Theorem 4.1.} Let \(\{E_x\}_{x\in X_n}\) be a family of modules. Suppose
local representations are given

\[
\pi_x:A_x\to\operatorname{End}(E_x)
\]

and operators \(V_u:E_x\to E_y\) satisfying

\[
V_u\pi_x(a)=\pi_y(\rho_u(a))V_u,
\tag{4.1}
\]

\[
V_vV_u=\pi_z(\Omega(v,u))V_{v\circ u}.
\tag{4.2}
\]

There exists a unique representation

\[
\Pi:\mathfrak H_n\to
\operatorname{End}\!\left(\bigoplus_{x\in X_n}E_x\right)
\]

that realizes \(e_x,J_x,T_u\) as the given projections, generators, and
transports.

\textbf{Proof.} The assignment of the generators defines a homomorphism
from the free algebra \(\mathcal F_n\). Equations (4.1)--(4.2) and the
local relations imply that \(\mathcal I_{\rm HMT}\) belongs to its
kernel. The homomorphism factors through the quotient (2.6). Uniqueness follows
because the classes of \(e_x,J_x,T_u\) generate the quotient. \(\square\)

The universal property turns the causal precedence of the generator into an algebraic
statement. A subsequent reader is legitimate when it represents the relations
of the generator; it does not need to reconstruct them from its codomain. The continuum,
the electron, gravity, and the exceptional branch are distinct realizations
of a single presentation.

\section{The elementary nonadic cocycle}

The canonical section \(s:C_9\to\{0,\ldots,8\}\subset C_{108}\) produces

\[
c_9(a,b)=\left\lfloor\frac{a+b}{9}\right\rfloor\in C_{12}.
\tag{5.1}
\]

The identity

\[
c_9(a,b)+c_9(a+b\bmod9,c)
=c_9(b,c)+c_9(a,b+c\bmod9)
\tag{5.2}
\]

is exact. If \(X\) is the lift operator and \(\lambda\) the image of the
memory generator, then

\[
X^9=\lambda I
\]

and

\[
X^aX^b
=\lambda^{c_9(a,b)}X^{a+b\bmod9}.
\tag{5.3}
\]

Equation (5.3) is the minimal form of the helix. The visible exponent
returns modulo nine; the factor \(\lambda\) preserves the turn. If
\(\lambda\ne1\), phase return and state return are distinct
operations.

\section{State refinement and reader prolongation}

The finite sets of states form an inverse system

\[
\cdots\xrightarrow{r_{n+2,n+1}}X_{n+1}
\xrightarrow{r_{n+1,n}}X_n.
\]

The deep state is

\[
X_\infty=\varprojlim_nX_n.
\tag{6.1}
\]

Functions and cylinder idempotents move in the opposite direction:

\[
r_{n+1,n}^*:f\longmapsto f\circ r_{n+1,n}.
\tag{6.2}
\]

If \(e_x\) is the indicator of the cylinder \(x\in X_n\),

\[
r^*e_x=\sum_{y:r(y)=x}e_y.
\tag{6.3}
\]

The family of algebras therefore forms a direct system:

\[
\mathfrak H_{\rm cyl}=\varinjlim_n\mathfrak H_n.
\tag{6.4}
\]

\textbf{Theorem 6.1.} The inclusions (6.2) are unital; cylinder idempotents
remain idempotent; and every compatible reader at finite
depth determines a unique reader on \(\mathfrak H_{\rm cyl}\).

\textbf{Proof.} The preimage of a partition is a partition. Hence
\(r^*(fg)=r^*f\,r^*g\), \(r^*1_n=1_{n+1}\), and the sum (6.3) is orthogonal.
The universal property of the direct limit extends every compatible family
of readers. \(\square\)

This theorem supplies the exact meaning of «fractalizing inward and
expanding outward». States are completed by restriction;
readers are completed by prolongation. Preservation of the unit is the
unitality of these inclusions.

\section{Holonomy and memory grading}

Let \(d_M\) be the memory degree. The holonomy satisfies

\[
d_M(\Gamma_9x)=d_M(x)+1.
\tag{7.1}
\]

In the reversible domain, \(\Gamma_9\) induces an automorphism of
\(\mathfrak H_{\rm cyl}\) and one can form

\[
\mathfrak H_{\rm rev}\rtimes_{\Gamma_9}\mathbb Z.
\tag{7.2}
\]

In the general domain, advance produces an endomorphism and the object is

\[
\mathfrak H_{\rm cyl}\rtimes_{\Gamma_9}\mathbb N.
\tag{7.3}
\]

The distinction between (7.2) and (7.3) prevents introducing an inversion where only
an advance has been constructed. The already typed helical inversion acts in
(7.2); the mirror \(\pi/e\) acts on the reader coordinates and constitutes
another involution.

\section{Total tritic algebra}

Let

\[
\mathbb A_{\rm tri}=A_+\oplus A_0\oplus A_-
\]

with central idempotents \(p_+,p_0,p_-\). The generators satisfy

\[
J_+^2=-p_+,
\qquad
J_0^2=0,
\qquad
J_-^2=p_-.
\tag{8.1}
\]

Let

\[
\mathbb J=p_+J_++p_0J_0+p_-J_-.
\]

The exponential defined by its internal series satisfies

\[
\boxed{
\begin{aligned}
\operatorname{Exp}_\star(t\mathbb J)
={}&p_+(\cos t+J_+\sin t)\\
&+p_0(1+tJ_0)\\
&+p_-(\cosh t+J_-\sinh t).
\end{aligned}}
\tag{8.2}
\]

\textbf{Proof.} In the elliptic fiber, the even and odd powers of
\(J_+\) alternate signs. In the parabolic fiber, all powers of
order at least two disappear. In the hyperbolic fiber, the even powers
are the identity and the odd powers are \(J_-\). Grouping the three series proves
(8.2). \(\square\)

The identity simultaneously preserves the three geometries. None of them
is obtained by informal continuation of another; all three are orthogonal
summands of the same object.

\section{Euler as elliptic projection}

Let \(\pi{}\) be the elliptic closure already produced by the nonadic
connection. Then

\[
\cos\pi{}=-1,
\qquad
\sin\pi{}=0
\]

in the circular realization, and (8.2) gives

\[
p_+\operatorname{Exp}_\star(\pi{}\mathbb J)+p_+=0.
\tag{9.1}
\]

Equation (9.1) recovers Euler's identity on the elliptic face. The
total element also preserves

\[
p_0\operatorname{Exp}_\star(\pi{}\mathbb J)
=p_0(1+\pi{}J_0),
\]

and

\[
p_-\operatorname{Exp}_\star(\pi{}\mathbb J)
=p_-(\cosh\pi{}+J_-\sinh\pi{}).
\]

The information omitted by the complex representation is precisely the threshold
jet and dilation. The tritic identity does not replace Euler's
identity: it identifies which projection produces it and what memory remains outside that
projection.

\section{The hyperbolic autoscale and \(\varphi\)}

Let

\[
F_{\rm av}=\begin{pmatrix}0&1\\1&1\end{pmatrix}.
\]

Its square has determinant one and trace three. We define

\[
H_\varphi=\frac{2F_{\rm av}^2-3I}{\sqrt5}.
\tag{10.1}
\]

A direct calculation gives

\[
H_\varphi^2=I.
\tag{10.2}
\]

Since

\[
\cosh(2\log\varphi)=\frac32,
\qquad
\sinh(2\log\varphi)=\frac{\sqrt5}{2},
\]

we obtain

\[
\boxed{
F_{\rm av}^2
=\operatorname{Exp}(2\log\varphi{}\,H_\varphi).}
\tag{10.3}
\]

The hyperbolic regime does not merely admit a quadratic unit: the
area--volume incidence matrix realizes that unit exactly. The ratio
\(\varphi^{-2}\) is the stable direction of a unimodular evolution.

\section{The structural role of \(e\)}

The series \(\operatorname{Exp}_\star\) transforms the additive generator into
multiplicative transport:

\[
\operatorname{Exp}_\star(a+b)
=\operatorname{Exp}_\star(a)\operatorname{Exp}_\star(b)
\]

when \(a\) and \(b\) commute. Its scalar realization on the identity
defines

\[
\operatorname{Exp}_\star(1)=e{}I.
\tag{11.1}
\]

Equation (11.1) does not replace the nonadic decimal generation of \(e\).
It identifies its algebraic role: \(e\) is the scale that converts infinitesimal
addition into a multiplicative transport law. This is the intrinsic
bridge between the two APP sheets in the exponential representation.

\section{Alpha as a dodecaphasic transgression functional}

The cocycle (5.1) represents the impossibility of choosing a homomorphic
section \(C_9\to C_{108}\). The visible turn produces an element of the
dodecaphasic kernel. The closure of \(\alpha\) receives that element along with the
signed word and determines it on a scalar line.

It is useful to distinguish the cocycle from its evaluation. Let

\[
\mathcal T_\alpha:
Z^2(C_9;C_{12})\times X_{12}^{\rm sgn}
\longrightarrow\mathcal L_\alpha
\tag{12.1}
\]

the functional realizing concatenation with memory. Then

\[
\alpha{}=\mathcal T_\alpha(c_9,K).
\tag{12.2}
\]

Equation (12.2) does not identify \(\alpha\) with a cohomology class.
It states that \(\alpha\) is the scalar realization of that class after
incorporating the dodecaphasic state. The second route determines the same section
through

\[
P_9(\alpha{})=\delta,
\qquad
\delta=\log_{10}\frac{10}{9},
\tag{12.3}
\]

and both meet in the resolution fiber product.

This algebraic type explains the singularity of \(\alpha\). \(\pi,\varphi,e\)
are closure, scale, and transport coordinates. \(\alpha\) is the functional
that measures how those coordinates survive the non-split lift.

\section{Compact matrix of the correlated system}

The four functions come together without collapsing into

\[
\boxed{
\begin{pmatrix}
\operatorname{Exp}_\star(\pi{}J_+) & 0&0&0\\
0&\operatorname{Exp}_\star(I)&0&0\\
0&0&\operatorname{Exp}_\star(2\log\varphi{}H_\varphi)&0\\
0&0&0&\operatorname{Exp}_\star(P_9(\alpha{})J_0)
\end{pmatrix}
=
\begin{pmatrix}
-I&0&0&0\\
0&e{}I&0&0\\
0&0&F_{\rm av}^2&0\\
0&0&0&I+\delta J_0
\end{pmatrix}.}
\tag{13.1}
\]

Matrix (13.1) is an identity of functions, not a table of values. Its
four blocks mean elliptic closure, exponential multiplication,
hyperbolic autoscaling, and parabolic seam. The dodecaphasic compatibility
of the last row preserves the joint genealogy.

\section{The electronic corner}

Let

\[
P=\begin{pmatrix}1&0\\0&0\end{pmatrix},
\qquad
Q=\begin{pmatrix}3/4&\sqrt3/4\\\sqrt3/4&1/4\end{pmatrix}.
\]

We define

\[
S=2P-I,
\qquad
J=\frac4{\sqrt3}[P,Q].
\]

Then

\[
S^2=I,
\qquad
J^2=-I,
\qquad
SJ=-JS.
\tag{14.1}
\]

Therefore,

\[
\boxed{
e_{\rm vac}\mathfrak H_ne_{\rm vac}\twoheadrightarrow
{\rm Cl}_{1,1}\cong M_2(\mathbb R).}
\tag{14.2}
\]

Arrow (14.2) expresses that the electronic block is a quotient or a
representation of a local corner of holonomic arithmetic.
Noncommutativity is not added from quantum mechanics: it already resides in the
order of the APP sheets. The internal root of \(-1\) is the elliptic generator;
the split involutions preserve vacancy and dilation.

\section{Readers and output types}

Let

\[
\operatorname{Read}(\mathfrak H)
=\coprod_B\operatorname{Rep}_{\rm comp}(\mathfrak H,B).
\tag{15.1}
\]

A constant is obtained as

\[
c=\sigma_B(\rho(x_\infty)),
\tag{15.2}
\]

where \(\rho\) is a compatible representation and \(\sigma_B\) a scalar
invariant specific to its codomain. This definition includes without confusing:

\[
\begin{array}{c|c|c}
\text{output}&B&\sigma_B\\ \hline
\pi&\text{elliptic module}&\text{minimum closure time}\\
\varphi&\text{hyperbolic module}&\text{stable spectral radius}\\
e&\text{exponential semigroup}&\text{unit evaluation}\\
\alpha&\text{dodecaphasic fiber product}&\text{scalar gluing}\\
G_{\rm Cat}&\text{character module}&\operatorname{Tr}(N^{-2}Q_4)\\
k_B&\text{thermal and energy lines}&\text{unit transduction}\\
\gamma_{\rm BI}&\text{area--volume module}&(12,-1)\text{ on }90/120\\
G&\text{gravitational module}&\text{contragredient character}
\end{array}
\]

A particle is not a digit in (15.2): it is a persistent section of a
module. A field is the transport of those sections. A symmetry is the
automorphism group preserving the relations and boundary of the module.

\section{Conservation and emergent symmetry theorem}

Let \(1_n=\sum_{x\in X_n}e_x\). The refinement inclusions satisfy

\[
r^*1_n=1_{n+1}.
\tag{16.1}
\]

An individual history also defines a compatible family of
cylinder idempotents. The first is the global unit of the algebra; the
second is the genealogical unit of a section. Distinguishing them avoids
confusing preservation of the space of possibilities with preservation of a
particular path.

For a representation \(\rho\) and a boundary \(b_\partial\), we define

\[
\operatorname{Sym}(\rho,b_\partial)
=\{g\in\operatorname{Aut}(\rho(\mathfrak H)):
g\rho(\mathcal R_{\rm HMT})=\rho(\mathcal R_{\rm HMT}),
\;gb_\partial=b_\partial\}.
\tag{16.2}
\]

Symmetry is, by definition, the stabilizer of an already constructed
conservation. In the exceptional branch, (16.2) leads to the
Witt--Golay--Leech--\(V^\natural\)--Monster incidence stabilizers. In the physical branch,
it leads to the groups preserving the electronic, gauge, or tensor
modules. No symmetry is imposed to select the state; one obtains
the symmetry leaving its projection invariant.

\section{Information, entropy, and quotients}

If \(U\) is bijective on the complete state,

\[
H(X\mid U(X))=0.
\tag{17.1}
\]

Let \(q\) be a quotient that forgets a coordinate. Then

\[
H(X)=H(qX)+H(X\mid qX).
\tag{17.2}
\]

A forgotten uniform trit contributes

\[
H(X\mid qX)=\log3.
\]

The thermal transducer determines

\[
Q_{\rm reset}\ge k_B\Theta\log3.
\tag{17.3}
\]

Equations (17.1)--(17.3) show that entropy does not measure the advance of
holonomy. It measures the information discarded by a reader. Sturmian
complexity \(p(m)=m+1\) allows infinite memory with zero topological entropy;
the thermal cost appears only when a projection is made irreversible.

\section{Core theorem}

\textbf{Theorem 18.1.} The object

\[
\boxed{
\mathfrak H_{\rm HMT}^{\rm alg}
=\left[\varinjlim_n
\bigl(\mathbb R\langle e_x,J_x,T_u\rangle/
\mathcal R_{\rm HMT}\bigr)\right]
\rtimes_{\Gamma_9}\mathbb N}
\tag{18.1}
\]

is a trigeometric holonomic algebra, graded by phase and memory, with an
associative product twisted by the carry. Its state space is the
inverse limit \(X_\infty\). Every compatible realization of APP, TRIT, and TPK
factors through (18.1). Its typed representations determine the joint discrete
structure of the continuum, the correlated system, the electron, action,
gravity, constants, and symmetries without turning their terminal
values into generators.

The notation in (18.1) is deliberately algebraic: it imposes no norm
before the codomain is chosen. If
\(\rho:\mathfrak H_{\rm HMT}^{\rm alg}\to B(\mathcal K)\) is a
compatible bounded representation, its completed realization is

\[
\mathfrak H_{\rm HMT}^{\rho}
=\overline{\rho(\mathfrak H_{\rm HMT}^{\rm alg})}^{\|\cdot\|}.
\tag{18.2}
\]

Thus each reader supplies the topology it effectively preserves;
no conventional completion is introduced as part of the generator.

\textbf{Proof.} Associativity follows from Theorem 3.1. Universality follows
from Theorem 4.1. Preservation of the unit and refinement duality follow
from Theorem 6.1. Memory elevation follows from the nonadic cocycle (5.3).
The coexistence of curvatures follows from (8.1). The realizations are obtained
through the universal property and the readers (15.1)--(15.2). \(\square\)

\section{Finite certification}

Exhaustive verification over the finite domains proves:

\begin{enumerate}
\item the \(9^3\) carry cocycle identities;
\item the matrix representation \(X^9=2I\) and its twisted law;
\item the three relations \(J_+^2=-I\), \(J_0^2=0\), \(J_-^2=I\);
\item the three branches of the exponential;
\item \(H_\varphi^2=I\) y \(\exp(2\log\varphi H_\varphi)=F_{\rm av}^2\);
\item the corner \({\rm Cl}_{1,1}\);
\item the contravariant/covariant unitality of refinement;
\item the synchrony \(A_2\) with diagonal advance \(-16\);
\item the difference between a reversible permutation and a tritic quotient;
\item the null Witt pair contained in the electronic corner;
\item the exact relations of the pair of phase observables and the memory commutator;
\item the irreducibility of APP, TRIT, and TPK through forgetful quotients, together with potential inversion in the sequence reader.
\end{enumerate}

These identities constitute a reproducible finite verification of the algebraic system.

\section{The null Witt pair within the electronic corner}

The electronic corner provides an even more complete local realization of
trigeometry. Let

\[
S=2P-I,\qquad
J=\frac{[P,Q]}{\sqrt3/4}.
\tag{20.1}
\]

The already proven relations are

\[
S^2=I,\qquad J^2=-I,\qquad SJ=-JS.
\tag{20.2}\label{eq:relaciones-cl11}
\]

Let us now define

\[
n_+=\frac{S+J}{2},\qquad
n_-=\frac{S-J}{2}.
\tag{20.3}
\]

Por (20.2),

\[
n_+^2=n_-^2=0,\qquad
n_+n_-+n_-n_+=I.
\tag{20.4}\label{eq:par-nulo-witt}
\]

The proof is immediate but structurally decisive:
\((S\pm J)^2=S^2+J^2\pm(SJ+JS)=0\), while the sum of the
two cross products equals
\(\tfrac12(S^2-J^2)=I\). Thus, the same block
\({\rm Cl}_{1,1}\cong M_2(\mathbb R)\) contains:

\[
\begin{array}{c|c|c}
\text{local generator}&\text{square}&\text{regime}\\ \hline
J&-I&\text{elliptic}\\
n_\pm&0&\text{parabolic}\\
S&I&\text{hyperbolic}.
\end{array}
\tag{20.5}
\]

This is an exact local Witt decomposition. The electronic vacancy
thus acquires a second equivalent characterization in its corner: in addition to
being the minimal loss of one multiplicative-carry unit under the
APP conservations, it is a parabolic null direction between the two
orientations \(n_+\) and \(n_-\). TRIT globally types these three classes and TPK
transports them with sheet, phase, and memory. The subsequent Paley--Witt
incidence structure preserves its own combinatorial domain; the
decomposition (20.3)--(20.5) supplies the local model from which its
realization can be formulated without equating the two objects.

\section{Crossed product of the nonadic and Sturmian phase observables}

Let

\[
\delta=\log_{10}\frac{10}{9},
\qquad
\zeta_9=\text{primitive character of }C_9,
\qquad
\omega_\delta=e^{2\pi i\delta},
\tag{21.1}
\]

and let \(T\) be the oriented step on the reversible lift. In the basis
\(\lvert t,m\rangle\), where \(t\) records the step and \(m\) the memory
height, we introduce the two diagonal characters

\[
V_9\lvert t,m\rangle=\zeta_9^t\lvert t,m\rangle,\qquad
W_\delta\lvert t,m\rangle=\omega_\delta^t\lvert t,m\rangle.
\tag{21.2}
\]

The covariance of the step produces the Weyl relations

\[
V_9T=\zeta_9\,TV_9,\qquad
W_\delta T=\omega_\delta\,TW_\delta.
\tag{21.3}
\]

For the nonadic turn \(\Gamma_9=T^9\), enriched by the TPK
update, one obtains

\[
V_9\Gamma_9=\Gamma_9V_9,\qquad
W_\delta\Gamma_9=\omega_\delta^9\Gamma_9W_\delta.
\tag{21.4}
\]

The nonadic observable closes; the rotation observable advances because
\(\delta\notin\mathbb Q\).
If \(D_M\lvert t,m\rangle=m\lvert t,m\rangle\), the law
\(\Gamma_9\lvert t,m\rangle=\lvert t+9,m+1\rangle\) is equivalent to

\[
[D_M,\Gamma_9]=\Gamma_9,
\qquad
\Gamma_9^{-1}D_M\Gamma_9=D_M+I
\tag{21.5}
\]

on the reversible sector. This is the operator equation of the helix:
phase closure is compatible with a unit increase in the memory
degree. The basal translation group remains abelian; the algebra
of observables and translations ceases to be so by (21.3). The Sturmian
word appears when the TRIT threshold is applied to the orbit of
\(W_\delta\); its aperiodicity therefore does not decorate the finite phase, but
constitutes the symbolic factor of the irrational rotation.

The operator presentation belongs to the crossed product

\[
\mathfrak W_{9,\delta}
=C\!\left(C_9\times\mathbb T\right)\rtimes_T\mathbb Z,
\tag{21.6}
\]

enriched by the grading \(D_M\). The spectrum of
\(C(C_9\times\mathbb T)\) simultaneously records the finite phase and the irrational
phase; the generator of \(\mathbb Z\) implements their joint translation;
and (21.5) prevents the graded representation from descending to a closed
orbit after forgetting memory. The subshift obtained through the tritic
partition of the second coordinate is the Sturmian hull. The helical suspension
is therefore the graded extension of that crossed product, not
the metaphorical juxtaposition of two clocks.

The characters \(\zeta_9\) and \(\omega_\delta\) belong to the holonomic
generator. Their subsequent elliptic representation recognizes

\[
\rho_+(\zeta_9)
=\operatorname{Exp}_+\!\left(\frac{2\pi{}}9J_+\right),
\qquad
\rho_+(\omega_\delta)
=\operatorname{Exp}_+\!\left(2\pi{}\delta J_+\right).
\tag{21.7}
\]

The conventional exponential thus appears as a realization of the already constructed
internal character; it neither selects the observables nor introduces a target value
into the generator.

Relations (21.3)--(21.5) separate exactly three facts:

\[
\text{phase period }9,\qquad
\text{absence of irrational period},\qquad
\text{integer memory advance}.
\tag{21.8}
\]

Their coexistence defines the Sturmian helical suspension over the nonadic
holonomic extension. The symbolic hull has long-range aperiodic
order and zero topological entropy, although its spectral module simultaneously
contains the rational nonadic frequency and the irrational frequency
\(\delta\). The physical name «time crystal» additionally requires fixing
the dynamic realization, its response, and its stability; the theorem formulated
here is the topological-symbolic object that precedes it.

\section{The minimal polynomial relation of the three curvatures}

The orthogonal sum of the three TRIT generators can be written as

\[
\mathbb J=p_+J_++p_0J_0+p_-J_-,
\qquad p_++p_0+p_-=I.
\tag{22.1}
\]

Their three squares imply a single relation:

\[
\boxed{\mathbb J^6=\mathbb J^2.}
\tag{22.2}
\]

Conversely, if \(K=\mathbb J^2\), then \(K^3=K\) and the three
regime idempotents are recovered polynomially:

\[
p_+=\frac{K^2-K}{2},\qquad
p_0=I-K^2,\qquad
p_-=\frac{K^2+K}{2}.
\tag{22.3}
\]

Thus, the elliptic, parabolic, and hyperbolic classification does not require three
algebras juxtaposed from outside. It is contained in the algebraic spectral
calculus of a single tritic generator. Substituting (22.3) into the
total exponential yields the compact identity

\[
\boxed{
\operatorname{Exp}_\star(t\mathbb J)=
p_+(\cos t+\mathbb J\sin t)+
p_0(I+t\mathbb J)+
p_-(\cosh t+\mathbb J\sinh t).
}
\tag{22.4}
\]

The Euler face is the elliptic projection of (22.4); parabolic
propagation occupies the threshold; the hyperbolic autoscale contains the realization
of \(\varphi\). Composing it with the holonomy,

\[
\mathcal E_{\rm HMT}(t)
=\Gamma_9\operatorname{Exp}_\star(t\mathbb J),
\qquad [D_M,\Gamma_9]=\Gamma_9,
\tag{22.5}
\]

trigeometry and the helical lift come together in two equations. APP
supplies the idempotents and the tension between sheets; TRIT supplies
\(\mathbb J\); TPK supplies \(\Gamma_9\), its cocycle, and the degree \(D_M\).
Numerical and physical outputs are typed representations of (22.5).
This presentation constitutes the minimal operator backbone of the standalone
article, while the universal algebra (18.1) preserves the full multiplicity
of states, paths, and refinements needed to reconstruct the integral
development.
 
\chapter{Carry vacancy, split Clifford algebra, and null Witt pair}
\begingroup
\renewcommand{\chapter}[2][]{}%
\chapter{Electronic vacancy, action, gravity, and boundary information}

\section{Vacancy as a minimal carry defect}

The level-ten additive fiber,

\begin{equation}
 \mathcal F_{10}=\{(i,j)\in D_9^2:i+j=10\},
 \label{eq:fibra-suma-diez}
\end{equation}

contains the center $Q_0=(5,5)$ and the oriented deformations
$Q_t=(5+t,5-t)$. Their product is $25-t^2$. Preserving the center's multiplicative
residue requires

\begin{equation}
 t^2\equiv0\pmod9.
\end{equation}

For $|t|\leq4$, the only solutions are $t=0,\pm3$. The two
noncentral deformations are therefore

\begin{equation}
 Q_{+3}=(8,2),\qquad Q_{-3}=(2,8),
 \label{eq:vacancias-orientadas}
\end{equation}

exchanged by cell inversion. In the APP chart,

\begin{equation}
 \mathcal A(5,5)=((1,1),(7,2)),
 \qquad
 \mathcal A(8,2)=((1,1),(7,1)).
 \label{eq:hojas-vacancia}
\end{equation}

The additive evaluation, its quotient, the multiplicative residue, and
the tritic regime are preserved; only
$q^\times:2\mapsto1$ changes. Vacancy is, by definition, this minimal loss of
one multiplicative-carry unit while preserving the remaining
invariants. It is not a zero added to the state.

The same fiber is reached from the nonadic reduction

\begin{equation}
 (369,126)\xmapsto{/9}(41,14)\longmapsto(5,5).
 \label{eq:369-centro}
\end{equation}

It has two distinct prolongations:

\begin{align}
 (41,14)&\xrightarrow{(+3,0)}(44,14)\longmapsto(8,5),
 \label{eq:rama-22-7}\\
 (41,14)&\xrightarrow{(+3,-3)}(44,11)\longmapsto(8,2).
 \label{eq:rama-vacancia-conservativa}
\end{align}

The first satisfies

\begin{equation}
 \frac{396}{126}=\frac{22}{7};
 \label{eq:aproximante-pi-vacancia}
\end{equation}

the second preserves

\begin{equation}
 41+14=44+11=55=54+1.
 \label{eq:unidad-borde-vacancia}
\end{equation}

The two branches separate the closure reader from the vacancy reader. The
former recognizes a period ratio; the latter preserves the boundary
unit while reducing the carry exactly.

\begingroup\fontsize{10.2}{11.7}\selectfont
\setlength{\parskip}{2pt}\setlength{\abovedisplayskip}{5pt}\setlength{\belowdisplayskip}{5pt}
\section{Matrix corner and electronic section}

The relations \eqref{eq:relaciones-cl11} and \eqref{eq:par-nulo-witt}
prove that the central fiber locally realizes the three geometries. To
specify its singular mode without normalization ambiguity, let
$\Sigma_3,\Pi_3:D_9^3\to D_9$ be the additive and multiplicative evaluations of
the triples and define their joint table by

\begin{equation}
 N^{(3)}_{ab}
 :=\#\{x\in D_9^3:\Sigma_3(x)=a,\ \Pi_3(x)=b\},
 \qquad
 C^{(3)}_{ab}:=\frac{N^{(3)}_{ab}}{\sqrt{n_a m_b^{\Pi}}},
 \label{eq:tabla-conjunta-tridimensional}
\end{equation}

where $n_a=\sum_bN^{(3)}_{ab}$ and
$m_b^{\Pi}=\sum_aN^{(3)}_{ab}$. Exhaustive enumeration of the
$9^3=729$ triples gives row marginals equal to $81$ and marginals equal
to $36$ in the six active columns. With

\begin{equation}
 m_{\mathrm{ph}}=(1,-1,0)^{\oplus3},
 \qquad
 N^{(3)}m_{\mathrm{ph}}
 =N^{(3)\mathsf T}m_{\mathrm{ph}}=-27m_{\mathrm{ph}},
 \label{eq:modo-entero-electronico}
\end{equation}

the normalization $\sqrt{81\cdot36}=54$ produces exactly

\begin{equation}
 C^{(3)}m_{\mathrm{ph}}
 =C^{(3)*}m_{\mathrm{ph}}=-\frac12m_{\mathrm{ph}},
 \label{eq:modo-singular-electronico}
\end{equation}

and the incompatibility of the sheet projectors equals

\begin{equation}
 \|[P_{\Sigma,3},P_{\Pi,3}]\|=\frac{\sqrt3}{4}.
 \label{eq:torsion-minima}
\end{equation}

The areal--volumetric incidence provides

\begin{equation}
 1-\left(\frac12\right)^2=\frac34=\frac{90}{120}.
 \label{eq:incidencia-electronica}
\end{equation}

The point $(5,5)$ is the unique fixed point of cell inversion. In this
way, the electronic section is not selected through a known mass:
it is singled out by the coincidence of centrality, minimal vacancy,
null direction, torsion \eqref{eq:torsion-minima}, and incidence
\eqref{eq:incidencia-electronica}. Its dimensionless character is

\begin{equation}
 \mathfrak e=
 \frac{\sqrt3}{4}
 \left(e^{\varphi/\pi^2}-22\alpha^3\right)(1+15\Delta_4),
 \label{eq:caracter-electronico}
\end{equation}

where each factor comes from an already constructed stage: minimal torsion, the
reciprocal orientation of the reader $R(x,y)=x/y^2$, the vacancy complement,
the compatibility constant, and the determinant of the fiber. The energy
chart is subsequently applied to this section.

\section{Polarity and constitutive equations}

Let $\mathsf C$ be sheet inversion. The common character $\mathcal E$ and the
oriented character $\mathcal B$ are defined by

\begin{equation}
 \mathcal E\circ\mathsf C=\mathcal E,
 \qquad
 \mathcal B\circ\mathsf C=-\mathcal B.
 \label{eq:paridad-constitutiva}
\end{equation}

The two combinations

\begin{equation}
 \widehat\mu=e^{\mathcal E+\mathcal B},
 \qquad
 \widehat\varepsilon=e^{\mathcal E-\mathcal B}
 \label{eq:mu-epsilon}
\end{equation}

satisfy

\begin{equation}
 \widehat\mu\,\widehat\varepsilon=e^{2\mathcal E},
 \qquad
 \frac{\widehat\mu}{\widehat\varepsilon}=e^{2\mathcal B}.
 \label{eq:constitutivas-comun-orientada}
\end{equation}

The product fixes the propagation cone and the quotient retains the impedance
orientation. The potential difference belongs to the boundary of the vacancy;
the transverse response belongs to its conjugate transport. The neutral
sector is $\ker\mathcal B$, and every section fixed by $\mathsf C$ belongs to
it. This structure makes polarity a property of vacancy
transport, not a label imposed afterward.

% La sección electrónica estructural concluye aquí. Las realizaciones de
% acción, gravitación, área e información se desarrollan posteriormente con
% sus fuentes rectoras completas y no se duplican en este capítulo.
 % 
 \endgroup
 \part{Aperiodic topological-temporal order and nonadic holonomy}


\par\endgroup
\chapter{Internal capacity calendar, phase clocks, and Diophantine renormalization}
\section{Aperiodicity as the internal calendar of nonadic capacity}

The Sturmian word and the nonadic helix are not superimposed
structures. Both are realizations of a single capacity arithmetic.
Let

\[
 \rho:=6\log_{1000}3=\log_{10}9,
 \qquad
 \delta:=1-\rho=\log_{10}\!\left(\frac{10}{9}\right).
\]

The decimal capacity index opened after ternary block number
\(t\) is

\[
 K_t:=\left\lfloor (t+1)\rho\right\rfloor,
 \qquad t\geq0,
\]

and the mechanical vacancy word is

\[
 z_t:=\lfloor(t+1)\delta\rfloor-\lfloor t\delta\rfloor,
 \qquad t\geq1.
\]

Since \(\delta\) is irrational, \((t+1)\delta\notin\mathbb Z\). Therefore,

\[
\begin{aligned}
 K_t
 &=\left\lfloor(t+1)(1-\delta)\right\rfloor\\
 &=t+1-\left\lceil(t+1)\delta\right\rceil\\
 &=t-\left\lfloor(t+1)\delta\right\rfloor.
\end{aligned}
\]

Subtracting two successive terms yields the link that remained to be made explicit:

\[
 \boxed{K_t-K_{t-1}=1-z_t.}
\]

This is not a similarity between two counts. It is an exact identity between the
capacity increment used by generation and the vacancy word
produced by the incompatibility \(729/1000\). Its two cases are:

\[
 \begin{array}{c|c|c}
 z_t & K_t-K_{t-1} & \text{visible operation}\\
 \hline
 0 & 1 & \text{the block opens }729\text{ subcells and a new decimal pruning enters},\\
 1 & 0 & \text{the block opens }729\text{ subcells without increasing decimal pruning}.
 \end{array}
\]

The second case is reopening or \emph{stutter}. The aperiodic word therefore decides
when capacity advances and when it remains fixed.
Aperiodicity does not govern a periodic generator from outside: it is the internal
calendar of its open--prune operation.

Telescoping gives a global form of the same theorem. For \(0\leq a<b\),

\[
 K_b-K_a=(b-a)-\sum_{t=a+1}^{b}z_t.
\]

Each interval is decomposed, without remainder, into capacity increments and
vacancies. In particular, every conservation formulated on the pair

\[
 \bigl(\text{acquired capacity},\text{recorded vacancy}\bigr)
\]

must preserve their sum, which is the length of the block interval. This is
the most elementary arithmetic form of the evolutionary conservation of the unit
at this level; it is not yet the projective measure of the continuum, but it is its
exact local model.

\section{Two times, two returns, and phase retention with memory advance}

The preceding identity requires distinguishing two clocks.

The \textbf{block time} \(t\) increases with each application of the dynamics. The
\textbf{capacity time} \(K_t\) counts the decimal triads actually
pruned. The determined nonadic phase is

\[
 g_t:=[K_t]_9\in C_9.
\]

From \(\Delta K_t=1-z_t\) it follows.

\[
 \boxed{g_t-g_{t-1}\equiv1-z_t\pmod9.}
\]

When \(z_t=0\), the capacity clock and phase advance. When \(z_t=1\),
the time block advances but \(K_t\) and \(g_t\) remain. This is an
exact event of \textbf{phase retention with state advance}: the new
block opens another fiber of \(729\) possibilities, changes the signature, and preserves
the survival history even though the phase residue is the same. The
statement «phase returns and memory changes» thus has two internal
realizations that must not be confused:

\begin{enumerate}
\item local retention \(K_t=K_{t-1}\), caused by a Sturmian vacancy;
\item holonomic return \(K\mapsto K+9\), which traverses a complete circuit of \(C_9\).
\end{enumerate}

In both cases the determined phase coincides and the integral holonomic state differs,
but the first mechanism is a capacity pause and the second is a complete nonadic
turn.

This allows the base system to be written without reducing TPK to a rotation. Let
\(\theta_t=\{t\delta\}\in\mathbb T\) and let
\(x_t\in\mathcal X_{\rm TPK}^{\rm enr}\). The effective step has the form

\[
 (g_{t-1},\theta_{t-1},x_{t-1})
 \longmapsto
 \bigl(g_{t-1}+1-z_t,\;\theta_{t-1}+\delta,\;
       U_t(x_{t-1})\bigr),
\]

with

\[
 U_t=\operatorname{Upd}_t\circ\operatorname{Tra}_t
     \circ\operatorname{Sel}_t.
\]

The coordinate \(z_t\) schedules capacity advance; it does not replace the action
of \(U_t\). The selector preserves sheet and regime, transport preserves path
and orientation, and the update incorporates carry, boundary, and register.
Therefore, equality of \(g\) never authorizes equality of \(x\).

\section{The inverse clock and the suspension with nonconstant roof}

The monotone inversion of capacity reveals a second mechanical word.
For \(k\geq1\), let

\[
 T(k):=\min\{t:K_t=k\}
      =\left\lceil\frac{k}{\rho}\right\rceil-1.
\]

Let us introduce

\[
 \sigma:=\frac{1}{\rho}-1=\frac{\delta}{\rho}.
\]

Since \(\rho\) is irrational, \(\sigma\) is also irrational, and

\[
 T(k)=k+\lceil k\sigma\rceil-1.
\]

The number of blocks needed to increase capacity by one unit is

\[
 r_k:=T(k+1)-T(k)
 =1+\lceil(k+1)\sigma\rceil-\lceil k\sigma\rceil
 \in\{1,2\}.
\]

The value \(r_k=2\) records exactly that a reopening occurred between capacities \(k\) and
\(k+1\). The exceptional positions begin

\[
 20,41,62,83,104,125,145,166,187,208,\ldots,
\]

and their separations are \(20\) or \(21\). This is the same dynamics viewed on
the inverse clock; it must not be confused with the separations \(21/22\) between
vacancies in block time.

For a capacity increment of length \(m\), the return time is

\[
\begin{aligned}
 R_m(k)
 &:=T(k+m)-T(k)\\
 &=m+\lceil(k+m)\sigma\rceil-\lceil k\sigma\rceil.
\end{aligned}
\]

By mechanical balance,

\[
 R_m(k)\in
 \{m+\lfloor m\sigma\rfloor,\;m+\lceil m\sigma\rceil\}.
\]

Two exact consequences follow:

\[
 \boxed{R_9(k)\in\{9,10\}},
 \qquad
 \boxed{R_{108}(k)\in\{113,114\}}.
\]

One nonadic phase turn requires nine or ten blocks; one turn of the
complete clock \(C_{108}\) requires one hundred thirteen or one hundred fourteen. The finite
support preserves its phases exactly, but the block time required
to traverse it is aperiodic. This is a more precise formulation of the
\textbf{Sturmian helical suspension}: the finite clock \(C_{108}\) is the base,
the irrational rotation provides the roof \(r_k\), and the TPK cocycle
transports the enriched fiber.

The pair \(6/7\) reappears here because

\[
 R_6(k)\in\{6,7\}.
\]

However, this \(6/7\) measures blocks consumed by six capacity
increments. It differs both from the number of vacancies in a window of
length \(132\) and from the separation between short gaps in the primary
word. All three phenomena descend from the same slope, but only an
explicit conjugacy would allow them to be identified as the same operator.

\section{Diophantine renormalization and unimodular conservation of the unit}

The continued fraction of the vacancy density, calculated after
\(\delta\) has been generated by \(729/1000\), begins with

\[
 \delta=
 [0;21,1,5,1,6,2,3,1,1,8,1,2,2,5,\ldots].
\]

Its first convergents are

\[
 \frac1{21},\quad
 \frac1{22},\quad
 \frac6{131},\quad
 \frac7{153},\quad
 \frac{48}{1049},\quad
 \frac{103}{2251},\quad
 \frac{357}{7802},\quad
 \frac{460}{10053},\quad
 \frac{817}{17855},\quad
 \frac{6996}{152893},\quad
 \frac{7813}{170748},\ldots
\]

This list explains, within a single hierarchy, the primary gaps and their
induced returns. The denominators \(21,22\) are the first approximation
times; the pairs \((6,131)\) and \((7,153)\) are the next times
with their event counts. In particular,

\[
 131=14\cdot9+5,
 \qquad
 153=17\cdot9,
\]

so that the return of secondary length \(6\) leaves a phase residue
\(+5\), while the one of length \(7\) closes the nonadic phase.

The integer \(5\) does not appear here through fitting. It is the partial quotient that
produces

\[
 (6,131)=5(1,22)+(1,21),
\]

and the next step gives

\[
 (7,153)=(6,131)+(1,22).
\]

An exact function of level \(5\) within vacancy
renormalization is thus located. Identifying this partial quotient with the
Ramanujan continued fraction of level \(5\) still requires a map between the two
objects; equality of the integer does not provide it on its own.

Let \(p_n/q_n\) be the \(n\)-th convergent. Two consecutive convergents
satisfy

\[
 p_{n+1}q_n-p_nq_{n+1}=(-1)^n.
\]

At the first two stages this is seen as

\[
 1\cdot22-1\cdot21=1,
 \qquad
 6\cdot153-7\cdot131=1.
\]

Therefore, the consecutive vectors \((p_n,q_n)\) form a primitive basis
of \(\mathbb Z^2\) and preserve unit lattice area. The orientation
alternates, but the modulus of the determinant remains equal to one. This identity
demonstrates the conservative evolution of the dimensional unit: the real slope is approximated at
increasing depth while each change of scale preserves a fundamental
cell and admits an integer inverse.

The recurrence is written using

\[
 A(a)=
 \begin{pmatrix}a&1\\1&0\end{pmatrix},
 \qquad \det A(a)=-1.
\]

The ordered products of the matrices \(A(a_n)\) generate the pairs
\((p_n,q_n)\). For \(a\ne b\),

\[
 A(a)A(b)=
 \begin{pmatrix}ab+1&a\\b&1\end{pmatrix}
 \ne
 \begin{pmatrix}ab+1&b\\a&1\end{pmatrix}
 =A(b)A(a).
\]

The rotation by \(\delta\) remains abelian. Noncommutativity
appears in the renormalization cocycle that remembers the order of the
partial quotients, in addition to that already demonstrated in the APP sheets and in the
affine cocycle. The crystal therefore possesses an abelian phase layer and
at least two ordered nonabelian memory layers.

The oriented errors

\[
 \varepsilon_n=q_n\delta-p_n
\]

alternate in sign and never vanish. Their sign defines a tritic renormalization
reader with values \(\pm1\); the regime \(0\) is reserved for the
exact threshold, excluded here by irrationality. This reader does not replace
the local TRIT of APP. It is a projection at another scale whose naturality with the
tritic phase selector must be proved on the enriched fiber.

\section{First Diophantine resonance of the block clock and separation of the capacity clock}

The sequence of convergents contains returns that close the
finite coordinate exactly and only approximate the irrational coordinate. Among the
convergents written above:

\[
 153\equiv0\pmod9,
 \qquad
 10053\equiv0\pmod9,
 \qquad
 170748\equiv0\pmod{108}.
\]

The last is the first convergent in this chain whose denominator closes the
block clock —the chronological coordinate θ indexed by the number of
iterations—:

\[
 170748=1581\cdot108.
\]

At the same time,

\[
 170748\,\delta-7813
 =-0.0000017458436865028488500088576\ldots
\]

The return is not periodic: the chronological coordinate
\([t]_{108}\in C_{108}^{(t)}\) returns exactly, but the irrational
coordinate remains shifted by a nonzero residue. It is indeed a
\textbf{privileged Diophantine quasi-turn}, because it comes from a convergent and
not from a decimal search.

This clock must be separated from the capacity quotient

\[
 C_{108}^{(K)}:\qquad g_{108}(t)=[K_t]_{108}.
\]

At the canonical origin, after \(170748\) blocks, the acquired capacity
is

\[
 \lfloor170748\rho\rfloor
 =170748-7813
 =162935
 \equiv71\pmod{108}.
\]

Therefore, that convergent does \textbf{not} close \(C_{108}^{(K)}\). The returns
\(R_{108}\in\{113,114\}\) of the previous section belong precisely to
this second clock: they count how many blocks are needed to acquire
\(108\) units of capacity. The coincidence of the name \(C_{108}\) does not
authorize merging the two indices. The observable TPK calendar is the chronological
clock; the capacity clock is a construction induced by the
compactification \(729\to1000\).

The system thus distinguishes four notions:

\begin{enumerate}
\item exact closure of the finite chronological clock;
\item nonclosure, in general, of the finite capacity clock;
\item optimal proximity of the Sturmian phase;
\item nontrivial transformation of the integral holonomic state by the TPK cocycle.
\end{enumerate}

A complete repetition would require simultaneous closure of both clocks,
vanishing of the irrational residue and return of the fiber. Here only the
first closes. This decoupling is the mathematical content of the
helical aperiodic crystal: exact return in a finite factor, drift in
the other, hierarchical near-return in the irrational factor and genuine
evolution in the fiber.

\section{Vacancy spectrum over the internal cardinalities of HMT}

The complementary identity allows any internal length
\(N\) to be evaluated without using a target constant. Every window of \(N\) steps has

\[
 V_N\in\{\lfloor N\delta\rfloor,\lceil N\delta\rceil\},
 \qquad
 P_N=N-V_N,
\]

where \(V_N\) counts vacancies and \(P_N\) capacity increments. For
some cardinalities already constructed by HMT one obtains:

\[
\begin{array}{c|c|c}
N&V_N&P_N\\
\hline
108&4\text{ or }5&104\text{ or }103\\
132&6\text{ or }7&126\text{ or }125\\
243&11\text{ or }12&232\text{ or }231\\
396&18\text{ or }19&378\text{ or }377\\
468&21\text{ or }22&447\text{ or }446\\
729&33\text{ or }34&696\text{ or }695\\
1000&45\text{ or }46&955\text{ or }954.
\end{array}
\]

Two rows require special attention.

At length \(132\), the branch with six vacancies produces exactly

\[
 132=6+126.
\]

The integer \(126\) reappears in the electronic census, but the two objects are not
yet identified: one is complementary capacity in a temporal
window and the other is the cardinality of an electronic class. The possible bridge
must send window, phase, orientation and memory to the census, not merely preserve
the number.

At length \(468\), which coincides with the cardinality of the catalog \(U_6\), the
number of vacancies is \(21\) or \(22\). Here the primary scale reappears
as the multiplicity of events within a previously generated cardinality.
The identity is exact; promoting it to self-similarity of the emitter requires
proving that the enumeration of the \(468\) sextets respects the order of the
mechanical rotation.

The complete decimal window gives an even more explicit closure:

\[
\begin{array}{c|c|c|c}
V_{1000}&P_{1000}&P_{1000}-900&V_{1000}+P_{1000}\\
\hline
45&955&55&1000\\
46&954&54&1000.
\end{array}
\]

Thus, the pairs \(45/46\) and \(55/54\) are exact complements after
removing the hundred nonadic turns \(900=100\cdot9\). The corridor
\(45\leftrightarrow54\leftrightarrow55\) is not added through a
coincidence: it appears as a balance of pruning and vacancy in the base window
\(1000\) itself. The subsequent interpretation of \(54\) in the modular corridor
preserves its own APP derivation; what has been constructed here is the
co-occurrence map that both constructions will have to make commute.

\section{The window \(396\), the boundary \(27\) and the radical factorization \(369\)}

The first interval of \(396=3\cdot132\) steps contains, in the canonical phase,
three consecutive windows with six vacancies each. Therefore,

\[
 396=18+378,
 \qquad
 18=3\cdot6,
 \qquad
 378=3\cdot126=42\cdot9.
\]

The radical and electronic censuses satisfy by another route

\[
 396=369+27,
 \qquad
 27=3\cdot9.
\]

The two decompositions come together arithmetically in

\[
 \boxed{369=378-9=3(132-6)-9}
\]

and

\[
 \boxed{27=18+9.}
\]

These two identities establish the cardinal factorization of the square: \(369\)
is the capacity obtained by separating one nonadic turn from the three lower windows
of \(132\), while \(27\) brings together the eighteen vacancies and that turn.
The assertion is strictly cardinal at this point; the realization on
carry memory is obtained later through the integer reader of the
specular transit matrix. The type distinction prevents confusing an
identity of cardinalities with a premature identification of fibers.

It also establishes the scope of the certified numbers \(396\), \(2376\),
\(387\) and \(43\). The certificate executes \(396\) events of six trits, whence

\[
 2376=396\cdot6.
\]

Comparing each level with the one nine positions later leaves
\(396-9=387\) pairs, grouped into

\[
 387=43\cdot9
\]

complete turns. The additional equality \(2376=22\cdot108\) is true,
but does not identify the trits of the certificate with twenty-two copies of the clock
\(C_{108}\); the former is a count of tritic positions and the latter a
nonsplit calendar. Only a map preserving sector, phase, carry and
memory could promote that factorization to an equivalence of objects.

\section{Synchronization, orientation and reopening of the three channels}

The census extended to \(t=120\) modifies the picture of a single closure at
the ninth gate. The synchronization of cardinalities reappears in clusters:

\[
\begin{aligned}
&\{8,9,10,12,13\},\quad
\{16,17,18,19,20\},\quad
\{29\},\\
&\{35,36,37,38,40,42\},\quad
\{49,51,52,54,55,56,57\},\\
&\{63\},\quad\{72,74\},\quad\{77,79,80\},
\quad\{83,84\},\\
&\{101\},\quad\{104,105,107,108\},
\quad\{116,118,119\}.
\end{aligned}
\]

The first reopenings are at

\[
 21,43,65,87,109,131,152,174,196,\ldots,
\]

with gaps \(21/22\). In the certified window \(0\leq t\leq120\), no
\emph{stutter} is a triple synchronization. The datum is finite, not an asymptotic
theorem, but reveals a breathing dynamics:

\[
 \text{wide opening}
 \longrightarrow\text{pruning}
 \longrightarrow\text{intermittent synchronization}
 \longrightarrow\text{aperiodic reopening}.
\]

The first breath is explicit. After the triple synchronization

\[
 (59,59,59)\to(43,43,43)\to(32,32,32),
\]

new small synchronizations appear

\[
 (17,17,17),\ (13,13,13),\ (6,6,6),\ (4,4,4),\
 (4,4,4),\ (3,3,3),\ (2,2,2),
\]

and the vacancy \(t=21\), with \(K_{21}=K_{20}=20\), reopens the fiber as

\[
 (611,527,723).
\]

The same capacity phase \([20]_9=2\) remains, but the cardinalities,
orientation and signature change drastically. This is direct material
evidence that the crystal is not an orbit contracted toward a fixed
point: it alternates contraction and reopening under an aperiodic calendar, while
survival maintains long-range coherence.

Equality of cardinalities does not erase orientation either. At \(t=8\) there are
three local orientations and at \(t=9\) two remain; future survival
selects one. At the orientation gates, \(\varphi\) remains fixed and
\(\pi+e\) remains constant, while the distribution \(\pi/e\) is reflected.
Synchronization must therefore be defined as equality of the size of the
channel fibers, not as identity of words, signatures or states.

\section{From triple synchronization to the dodecaphasic closure of \(\alpha\)}

The new structure clarifies the causal position of \(\alpha\). The channels
\(\pi\), \(e\) and \(\varphi\) do not need \(\alpha\) to synchronize their
cardinalities: the previous clusters arise during filtration and
survival. Equality of cardinalities is not sufficient to produce
\(\alpha\) either, because orientation, specular distribution, carry and
boundary condition still remain.

The dodecaphasic closure receives precisely that additional information. From
the same terminal holonomic state, two readings open up:

\[
 x_{\rm term}
 \longmapsto
 B_{\rm term},
 \qquad
 x_{\rm term}
 \longmapsto
 U_{12}^{\rm sgn}.
\]

The first determines the three surviving channels; the second preserves the
signed coordinate needed to propagate the carry. The integer route produces

\[
 P+E-\Phi-K+C^+-1000C=A,
\]

while the vacancy route evaluates the same genealogy through

\[
 \mathcal P(x)=\delta,
 \qquad
 \mathcal P(x)=
 2\pi x-\frac74x^2+\frac{x^3}{2\pi}
 +\frac{x^4}{20}-\frac{2x^5}{21}-\frac{x^6}{46}.
\]

The identity \(\Delta K=1-z\) explains why the second route belongs to the
same generator: \(\delta\) is not an external constant placed alongside the
three channels, but the density of the calendar that determines their capacity and their
reopenings. The synchronization of \((\pi,\varphi,e)\) and the realization of
\(\alpha\) are not two separate stories. They are two levels of the same state:

\[
 \text{fiber synchronization}
 \longrightarrow
 \text{oriented resolution and carry memory}
 \longrightarrow
 \text{dodecaphasic closure of }\alpha.
\]

This preserves the causal lock: the correlated system
\((\pi,\varphi,e,\alpha)\) is a correlated output of coinductive
generation. Within that common genealogy, the three channels are determined and
synchronized before the dodecaphasic seal determines the coordinate
\(\alpha\). No conventional value of the four constants enters
as an input.

\section{Multiplicative breathing cocycle and exactly zero logarithmic drift}

The complementarity between capacity and vacancy admits a multiplicative
realization that makes the author’s picture of cylinders contracting and
reopening more precise. We define the ambient scale of the cylinder after block \(t\) by

\[
 \Omega_t:=\frac{729^{t+1}}{1000^{K_t}},
 \qquad
 K_t=\lfloor(t+1)\rho\rfloor,
 \qquad
 \rho=\log_{1000}729=\log_{10}9.
\]

The equality \(729=1000^\rho\) immediately gives

\[
 \boxed{\Omega_t=1000^{\{(t+1)\rho\}}.}
\]

Therefore,

\[
 1\leq\Omega_t<1000
\]

for every \(t\), and the set \(\{\Omega_t:t\geq0\}\) is dense in the
half-open multiplicative interval \([1,1000)\). Density is not a
statistical hypothesis: it follows from the irrationality of \(\rho\). Indeed,
if \(\rho=a/b\in\mathbb Q\), then \(9^b=10^a\), contradicting the
uniqueness of prime factorization.

Let

\[
 c_t:=K_t-K_{t-1}=1-z_t\in\{0,1\}
\]

the compactification carry. The ratio between two consecutive scales
is exactly

\[
 \frac{\Omega_t}{\Omega_{t-1}}
 =\frac{729}{1000^{c_t}}
 =
 \begin{cases}
 729/1000,&z_t=0,\\[1mm]
 729,&z_t=1.
 \end{cases}
\]

The first branch contracts; the second reopens. These are not two rules added to the
calendar: they are the two charts of the same multiplication by \(729\) reduced modulo
the scale \(1000\). Before the first reopening one observes, for example,

\[
 \Omega_{20}=1.3100205086376203\ldots,
 \qquad
 \Omega_{21}=729\Omega_{20}
 =955.0049507968252\ldots.
\]

The accumulated contraction does not lead to an attractor because each vacancy
restores the scale through an exact jump. In logarithmic coordinates,

\[
 \ell_t:=\log_{1000}\frac{\Omega_t}{\Omega_{t-1}}
 =\rho-c_t=z_t-\delta,
 \qquad \delta=1-\rho.
\]

Moreover,

\[
 \sum_{j=1}^{n}\ell_j
 =\log_{1000}\Omega_n-\log_{1000}\Omega_0,
\]

so that all partial sums remain uniformly bounded. In
particular,

\[
 \boxed{\lim_{n\to\infty}\frac1n
 \log\frac{\Omega_n}{\Omega_0}=0.}
\]

This is a stronger assertion than mere equality of a mean
frequency: the logarithmic cocycle is a bounded coboundary. Breathing preserves
the scale in the long term with zero Lyapunov exponent, although each step is
a contraction or a very different reopening.

The natural geometric form of the result is the multiplicative circle

\[
 \mathbb A_{1000}:=\mathbb R_{>0}/1000^{\mathbb Z}.
\]

The map

\[
 R_{729}:[x]\longmapsto[729x]
\]

is conjugate, through \(u=\log_{1000}x\), to the irrational rotation
\(u\mapsto u+\rho\) of \(\mathbb R/\mathbb Z\). The integer \(c_t\) is the
index of the deck transformation that returns \(729\Omega_{t-1}\) to
the fundamental section \([1,1000)\); the vacancy \(z_t=1-c_t\) is,
exactly, the minimal absence of that carry. This formulation brings together in a
single object compactification, multiplicative elasticity, the torsion of
the covering and the memory of how many times its boundary was crossed.

The invariant measure induced on the fundamental section is logarithmic:

\[
 d\mu_{\Omega}(x)=\frac{dx}{x\log1000},
 \qquad 1\leq x<1000.
\]

Therefore, for \(s\ne0\),

\[
 \int x^s\,d\mu_{\Omega}(x)
 =\frac{1000^s-1}{s\log1000},
\]

and the geometric mean is \(1000^{1/2}=10\sqrt{10}\). These quantities
describe the mathematical distribution of the multiplicative radius; they are not
identified with a physical length or with the Planck length without the
corresponding dimensional transducer.

The zero exponent does not replace the Landauer theorem of the integral
holonomic state either. Here the absence of exponential drift in a scale
coordinate is demonstrated; zero Landauer requires recoverability of the complete
evolution. These are compatible and typologically distinct results.

\textbf{Refutation criterion.} A value of \(\Omega_t\) outside \([1,1000)\), a ratio
different from \(729/1000\) or \(729\), or an unbounded logarithmic sum.

\section{Phase and vacancy as coordinates of a single lifted rotation}

The finite phase and the aperiodic word do not form an arbitrary product. They
can be glued exactly as two coordinates of a single irrational
translation. Let

\[
 u_t:=\{(t+1)\rho\}\in[0,1),
 \qquad
 g_t^{(q)}:=[K_t]_q\in C_q,
\]

for \(q\geq2\). On

\[
 X_q=C_q\times[0,1)
\]

with the interval endpoints glued by the covering, we define

\[
 c(u):=\lfloor u+\rho\rfloor\in\{0,1\}
\]

and

\[
 F_q(g,u)=
 \bigl(g+c(u),\;u+\rho-c(u)\bigr).
\]

The second component remains in \([0,1)\); the first preserves the boundary
crossing. The map

\[
 \Psi_q:X_q\longrightarrow\mathbb R/q\mathbb Z,
 \qquad
 \Psi_q(g,u)=[g+u]_q,
\]

satisfies

\[
 \boxed{\Psi_q\circ F_q=R_\rho\circ\Psi_q,}
\qquad
R_\rho(y)=y+\rho\pmod q.
\]

On the canonical orbit,

\[
 \Psi_q(g_t^{(q)},u_t)=[(t+1)\rho]_q.
\]

Therefore, the phase \([K_t]_q\), the irrational coordinate \(u_t\) and the
carry \(c_t=1-z_t\) are not three juxtaposed decorations. They are the integer
part, the fractional part and the covering cocycle of a single point on
the circle of length \(q\).

This conjugacy provides a structural proof of the observed
complexities. The coding of \(R_\rho\) by the \(q\) unit intervals
has, at length \(m\), exactly \(qm\) words. Its boundaries are

\[
 j-k\rho\pmod q,
 \qquad j=0,\ldots,q-1,\quad k=0,\ldots,m-1,
\]

and they are all distinct because \(\rho\) is irrational. By adding \(z_t\), a
block of length \(m\) reconstructs the immediately preceding phase
symbol through

\[
 g_{t-1}^{(q)}=g_t^{(q)}-(1-z_t)\pmod q.
\]

There is therefore a bijection between decorated blocks of length
\(m\) and phase blocks of length \(m+1\). One obtains

\[
 \boxed{p_{g,q}(m)=qm,\qquad
 p_{(g,z),q}(m)=q(m+1).}
\]

For \(q=3,9,12,108\), these laws explain the certified linear
complexities without assuming independence between a periodic clock and a
Sturmian decoration. Growth remains linear and topological entropy
remains zero, but the finite factor exactly multiplies the
available language.

The inversion of the helix is now unambiguous: in the coordinate \(y\) it is the
rotation \(y\mapsto y-\rho\). The descending traversal does not invent another
word; it reads the same invertible system with the conjugate covering cocycle.
The phase can return, while \(u\) and the TPK fiber retain the history.

This formulation also requires distinguishing two objects of order \(108\):

\[
 \theta_t=[t]_{108}\in C_{108}^{(t)},
 \qquad
 g_t^{(108)}=[K_t]_{108}\in C_{108}^{(K)}.
\]

The denominator \(170748=1581\cdot108\) closes the first. By contrast,

\[
 \lfloor170748\rho\rfloor=162935\equiv71\pmod{108},
\]

so it does not close the second. The observable dodecaphasic TPK clock and
the induced capacity clock are compatible, but not identical. Their future
comparison requires specifying the index map, not reusing the symbol
\(C_{108}\) as though it erased the domain.

\textbf{Refutation criterion.} A collision between
boundaries \(j-k\rho\), a complexity different from \(qm\) or \(q(m+1)\), or
the assertion that \(170748\) closes the capacity phase.

\section{Induced vacancy clock and tritic typing of the returns \(6/7\)}

The secondary scale \(6/7\) does not merely count returns of short gaps. When
passing from the block clock to the event clock, it organizes exactly the
duration of the three regimes of the tritic selector.

For \(n\geq1\), let

\[
 p_n:=\left\lfloor\frac{n}{\delta}\right\rfloor
\]

the position of the \(n\)-th vacancy event and let

\[
 \kappa_n:=K_{p_n}
\]

the capacity retained at that event. Since

\[
 p_n\delta<n<(p_n+1)\delta,
\]

we obtain

\[
 \boxed{\kappa_n=p_n-n.}
\]

If

\[
 h_n:=p_{n+1}-p_n\in\{21,22\},
 \qquad
 s_n:=22-h_n\in\{0,1\},
\]

then

\[
 \kappa_{n+1}-\kappa_n=h_n-1=21-s_n.
\]

Reducing modulo \(3\),

\[
 \boxed{[\kappa_{n+1}]_3=[\kappa_n-s_n]_3.}
\]

A long gap \(22\) preserves the regime; a short gap \(21\) rotates it by one
unit. Since short gaps reappear after six or seven events,
the plateaus of the induced regime have lengths exactly \(6\) or \(7\).
The first twenty classes are

\[
 \underbrace{2,\ldots,2}_{6},
 \underbrace{1,\ldots,1}_{7},
 \underbrace{0,\ldots,0}_{7}.
\]

To relate them to the canonical selector, the positive class is determined

\[
 r_n=1+(\kappa_n\bmod9)
\]

and one applies

\[
 M_{\rm ph}(r)=
 \begin{cases}
 +1,&r\in\{1,4,7\},\\
 -1,&r\in\{2,5,8\},\\
 0,&r\in\{3,6,9\}.
 \end{cases}
\]

Thus, the first plateaus traverse

\[
 \boxed{0\longrightarrow-1\longrightarrow+1\longrightarrow0}
\]

and each transition is triggered by a short gap. This is an exact
formulation of TRIT as a unit of topological-temporal information: it does not replace
the binary vacancy word, but integrates its events and determines in which
regime the next transfer operates. Information is preserved because
the pair \((\kappa_n,s_n)\) reconstructs the previous increment;
orientation is preserved because the cyclic order of the classes is not flattened.

The result does not turn the three values into three separate theories either.
They are three internal regimes of a single selector, and the induced word
visits them under the same cocycle. Their subsequent identification with elliptic,
parabolic and hyperbolic curvatures uses the algebras
\(\mathbb R[J_\tau]/(J_\tau^2+\tau I)\) already constructed; realization as
geometric curvature or a physical field requires the corresponding transducer.

\textbf{Refutation criterion.} A plateau of length different from \(6/7\), a regime
transition during a gap \(22\), or a transition other than the cyclic
descent through a gap \(21\).

\section{Universal matrix of the two balanced branches}

Sturmian balance has a linear form that brings together conservation and
orientation. In a window of length \(N\), let

\[
 n_N=\lfloor N\delta\rfloor.
\]

The two counting possibilities are \(n_N\) and \(n_N+1\), and their complements
are \(N-n_N\) and \(N-n_N-1\). We form the branch matrix

\[
 W_N:=
 \begin{pmatrix}
 n_N&N-n_N\\
 n_N+1&N-n_N-1
 \end{pmatrix}.
\]

Its identities are universal:

\[
 W_N\binom11=N\binom11,
 \qquad
 (1,-1)W_N=-(1,-1),
\]

\[
 \operatorname{tr}W_N=N-1,
 \qquad
 \det W_N=-N,
\]

and, therefore,

\[
 \boxed{\chi_{W_N}(\lambda)=(\lambda-N)(\lambda+1).}
\]

The diagonal mode preserves the total length of the window. The oriented
functional that subtracts one branch from the other changes sign with eigenvalue
\(-1\). The oriented unit is not introduced from outside: it appears because the
two balanced words differ by a single event and a single complementary
carry.

For the internal lengths considered one obtains, among others,

\[
 W_{132}=\begin{pmatrix}6&126\\7&125\end{pmatrix},
 \quad
 W_{468}=\begin{pmatrix}21&447\\22&446\end{pmatrix},
 \quad
 W_{1000}=\begin{pmatrix}45&955\\46&954\end{pmatrix}.
\]

Each has the eigenvalues \(N\) and \(-1\). This law precisely separates
conservation of the total unit and the elementary orientation of
the vacancy. The third tritic value \(0\) corresponds to the threshold of the
oriented functional, not to a third counting branch added to the binary
word.

The window \(1000\) contains an additional internal reduction. After
separating \(900=100\cdot9\) units of capacity, what remains is

\[
 \widetilde W_{100}=
 \begin{pmatrix}
 45&55\\
 46&54
 \end{pmatrix},
\]

with

\[
 \operatorname{tr}\widetilde W_{100}=99,
 \qquad
 \det\widetilde W_{100}=-100,
 \qquad
 \operatorname{Spec}(\widetilde W_{100})=\{100,-1\}.
\]

The half-imbalances of its rows are

\[
 \frac{55-45}{2}=5,
 \qquad
 \frac{54-46}{2}=4.
\]

By another route, the central APP block

\[
 B_c=\begin{pmatrix}7&2\\2&7\end{pmatrix}
\]

has spectrum \(\{9,5\}\) and jump \(9-5=4\). Therefore, the two
scalar readers

\[
 \chi_{\rm vac}(\widetilde W_{100})=(5,4),
 \qquad
 \chi_{\rm APP}(B_c)=(\lambda_{\min},
 \lambda_{\max}-\lambda_{\min})=(5,4)
\]

coincide exactly. This constructs the common codomain of the
vacancies--APP square. It does not yet demonstrate that \(\widetilde W_{100}\) and \(B_c\)
are the same operator: for that promotion, the arrow transporting
the two window branches to the two APP eigenspaces while preserving phase,
sheet and carry is missing. The reservation is located at that arrow and does not weaken either of the
two derivations.

\textbf{Refutation criterion.} A
balanced window whose matrix does not have spectrum \(\{N,-1\}\), or a
purported identification with APP that does not transport the branches.

\section{Synchronization quotient \(A_2\) and projective direction of the first reopening}

The censuses of the three channels admit a canonical quotient that eliminates
common survival and preserves only their misalignment. For

\[
 N=(N_\pi,N_e,N_\varphi)\in\mathbb Z^3
\]

we define

\[
 \partial N=
 (N_\pi-N_e,\;N_e-N_\varphi,\;N_\varphi-N_\pi).
\]

Its image is contained in

\[
 A_2=\{(x,y,z)\in\mathbb Z^3:x+y+z=0\},
\]

and is all of \(A_2\). The kernel is the diagonal
\(\mathbb Z(1,1,1)\). Consequently,

\[
 \boxed{\mathbb Z^3/\mathbb Z(1,1,1)\simeq A_2.}
\]

The quadratic form

\[
 Q(N)=\frac12\|\partial N\|^2
\]

vanishes exactly at triple synchronizations. The unit defects
of the phases \(11,14,15\) satisfy \(Q=1\) and are minimal roots of
\(A_2\). Equalities of two channels occupy the three mirrors of the
Weyl group \(W(A_2)\simeq S_3\); the orientation on one side or the other of the mirror remains
in the sign of the root.

The initial corridor of mirrors is

\[
\begin{array}{c|c|c|c|c}
t& (N_\pi,N_e,N_\varphi)&\text{mirror}&\partial N&d_t\\
\hline
4&(207,207,122)&\pi=e&(0,85,-85)&85\\
5&(39,151,151)&e=\varphi&(-112,0,112)&112\\
6&(110,110,69)&\pi=e&(0,41,-41)&41\\
7&(81,29,81)&\pi=\varphi&(52,-52,0)&52\\
8&(59,59,59)&S_3&(0,0,0)&0\\
9&(43,43,43)&S_3&(0,0,0)&0.
\end{array}
\]

The four amplitudes preceding the first triple closure form the matrix

\[
 D_{\rm syn}:=
 \begin{pmatrix}85&112\\41&52\end{pmatrix}.
\]

Without introducing \(\alpha\), the electron or an experimental value, the identities
are verified

\[
 112-85=27,
 \qquad
 52-41=11,
\]

\[
 85+41=126,
 \qquad
 112+41=153,
 \qquad
 85+52=137,
\]

and

\[
 \boxed{\operatorname{tr}D_{\rm syn}=137,\qquad
 \det D_{\rm syn}=-172=-4\cdot43.}
\]

The same corridor thus contains the census \(126\), the return \(153\), the
integer direction \(137\), the boundary \(27\), the rank \(11\), the jump \(4\)
and the second synchronization \(43\). All these equalities are exact. The
matrix, considered in isolation, determines an integer direction of trace
\(137\), not the constant \(\alpha\). The derivation of \(\alpha\) additionally
requires the signed state, the dodecaphasic closure and the complete jet, which are
constructed in the corresponding sections.

The first reopening offers a more precise arrow. At \(t=20\),

\[
 N_{20}=(2,2,2),\qquad \partial N_{20}=0.
\]

At \(t=21\), the capacity remains \(K_{21}=K_{20}=20\), the determined phase
is the same and the selector is in the threshold class
\(r=3\), for which \(M_{\rm ph}=0\). However,

\[
 N_{21}=(611,527,723)
\]

and

\[
 \boxed{\partial N_{21}
 =(84,-196,112)=28(3,-7,4).}
\]

The primitive vector \((3,-7,4)\) belongs to \(A_2\). In the oriented projective
chart whose denominator is the component \(\varphi-\pi\),

\[
 \chi_{e\varphi\mid\varphi\pi}(\partial N)
 :=\frac{N_e-N_\varphi}{N_\varphi-N_\pi},
\]

we obtain

\[
 \boxed{\chi_{e\varphi\mid\varphi\pi}(\partial N_{21})
 =\frac{-196}{112}=-\frac74.}
\]

This is exactly the quadratic coefficient of the kernel
\(E_{21/22}\). It is no longer a matter of finding the integers \(7\) and
\(4\) separately: the complete oriented ratio emerges at the first reopening of the
three channels, after synchronization and before the resolution of
\(\alpha\). The identity is internal and does not use the target value of
\(\alpha\).

\section{Sturmian-triadic genesis of \(22/7\) and internal portrait of the jet \(E_{21/22}\)}

The first seven effective vacancy positions are

\[
 21,43,65,87,109,131,152.
\]

Between the first and the seventh, \(132\) steps elapse and six
intervals appear. In that first closed window,

\[
 132=6\cdot22,
 \qquad
 126=6\cdot21.
\]

Tritic amplification gives

\[
 396=3\cdot132,
\]

and, therefore,

\[
 \boxed{\frac{396}{126}
 =\frac{3\cdot6\cdot22}{6\cdot21}
 =\frac{22}{7}.}
\]

The rational seal \(22/7\) is not introduced as a historical approximation of
\(\pi\): it is obtained from the ratio between the tritic amplification of the
window and its survival on the branch of six long gaps. The subsequent comparison
with \(\pi\) preserves its status as recognition; it does not generate the
ratio.

The dynamics also provides five pairs of endpoints invariant under a change
of origin of the mechanical word:

\[
 \mathcal G_{\rm prim}=\{21,22\},
 \quad
 \mathcal R_{\rm sec}=\{6,7\},
 \quad
 \mathcal G_{\rm inv}=\{20,21\},
\]

\[
 \mathcal V_{1000}=\{45,46\},
 \qquad
 \mathcal H_{100}=\{4,5\}.
\]

With the orientation already fixed by the TPK, consider the extraction

\[
 r^+=7,\qquad h^-=4,\qquad
 g_{\rm inv}^-=20,\qquad g_{\rm prim}^-=21,\qquad
 v_{1000}^+=46.
\]

These five invariants reproduce exactly the nontranscendental
coefficients of the sixth-order jet:

\[
 -\frac{r^+}{h^-}=-\frac74,
 \qquad
 \frac1{g_{\rm inv}^-}=\frac1{20},
 \qquad
 -\frac2{g_{\rm prim}^-}=-\frac2{21},
 \qquad
 -\frac1{v_{1000}^+}=-\frac1{46}.
\]

Degrees one and three form the reciprocal pair

\[
 (2\pi)\,(2\pi)^{-1}=1.
\]

Thus, the declared kernel can be written as

\[
\begin{aligned}
 P_6(x)={}&
 2\pi{}x
 -\frac{r^+}{h^-}x^2
 +(2\pi)^{-1}x^3
 +\frac{x^4}{g_{\rm inv}^-}\\
 &-\frac{2x^5}{g_{\rm prim}^-}
 -\frac{x^6}{v_{1000}^+}.
\end{aligned}
\]

This expression does not use the root of \(P_6(x)=\delta\) to choose the
coefficients. All the numbers on the right-hand side have appeared earlier as
invariants of the compactification calendar. The coefficient \(-7/4\)
also has the second internal factorization from the preceding section as
a projective coordinate of the first defect \(A_2\).

The preceding equality is an \emph{exact genealogical factorization} of the
kernel. The selected endpoints, their degrees and their signs are brought together
by the graded reader \(\mathfrak E_9\), defined later on the invariant
signature of the integral holonomic state; this reader commutes with
transports that preserve origin, orientation, windows and memory.

The prolongation of degrees \(7,8,9\) already has another status. The reference constructions
demonstrate

\[
 T_{7:9}(x)
 =-\frac{x^7}{120}+\frac{x^8}{45}+\frac{2x^9}{495},
\]

with

\[
 120=|\mathcal F_8^{\rm ph}|,
 \qquad
 45=\det B_c,
 \qquad
 11=\operatorname{rank}(D_3,D_4),
 \qquad
 495=45\cdot11,
\]

and the nilpotent resolvent \((I-N)^{-1}\). The complete jet

\[
 P_9=P_6+T_{7:9}
\]

is thus stratified: degrees \(1{:}6\) come from the synchronization
signature and degrees \(7{:}9\) from phase incidence, the
boundary determinant and the dodecaphasic rank. The simple root of
\(E_{21/22}\) is obtained after that factorization and does not enter into the
selection of the coefficients.

\section{Refined causal chain: from compactification to the coordinate \(\alpha\)}

The results of the previous sections allow a verbal sequence to be replaced
by a diagram of domains and quotients:

\[
\begin{aligned}
&\text{APP--TRIT--TPK and internal ratio }729/1000\\
&\quad\longrightarrow
 (\mathbb A_{1000},R_{729},c_t,z_t)\\
&\quad\longrightarrow
 (g_t^{(9)},u_t,x_t^{\rm enr})\\
&\quad\longrightarrow
 (N_\pi,N_e,N_\varphi)\\
&\quad\xrightarrow{\ \partial\ }
 A_2\\
&\quad\longrightarrow
 \bigl(B_{\rm term},U_{12}^{\rm sgn}\bigr)\\
&\quad\longrightarrow
 \begin{cases}
 \text{integer dodecaphasic closure},\\
 \text{simple root of the jet }E_{21/22},
 \end{cases}\\
&\quad\longrightarrow\alpha.
\end{aligned}
\]

The first level constructs the multiplicative rotation and its carry
cocycle. The second simultaneously preserves phase, irrational part and state.
The third determines the three correlated channels. The quotient \(A_2\)
measures their misalignment without erasing common survival. The dodecaphasic
closure then receives the orientation and memory that the census
quotient does not contain.

Synchronization therefore does not by itself generate \(\alpha\). It does provide
two concrete internal precursors:

\begin{enumerate}
\item the rectangle of defects, whose trace is \(137\);
\item the first threshold reopening, whose projective coordinate is \(-7/4\).
\end{enumerate}

Both appear before solving the analytic equation and before any
experimental comparison. The first precursor locates an integer direction;
the second locates a coefficient of the jet. The remaining information resides
in the signed state, in the dodecaphasic carry and in the higher degrees
of the kernel.

This architecture specifies in what sense the aperiodic crystal can be the
nerve of the construction, without turning it into an all-encompassing phrase. The
crystal is the minimal system

\[
 \bigl(\mathbb A_{1000},R_{729},c_t;
 C_9,\Gamma_9;\mathcal X_{\rm TPK}^{\rm enr}\bigr)
\]

that brings together an irrational base of zero entropy, an integer carry, a
lifted finite holonomy and a fiber with memory. Its realizations
\(\pi,\varphi,e,\alpha\) are not inputs to the crystal; they are correlated
characters of the state it prolongs. The dimensional,
fractal and exceptional realizations are obtained afterward through their own transducers.
The graded reader \(\mathfrak E_9\) brings together the defect \(28(3,-7,4)\), the
matrix \(D_{\rm syn}\) and the vacancy endpoints in the coefficient signature
of \(P_9\), so that synchronization and the derivation of \(\alpha\)
are linked without introducing the value of the constant as a datum.

\section{The compactification slope is transcendental and its renormalization is nonstationary}

The aperiodic structure does not depend only on the vacancy word
having no period. Its slope has a stronger arithmetic property.
Recall that

\[
 \delta=\log_{10}\frac{10}{9},
 \qquad
 \rho=1-\delta=\log_{10}9.
\]

Both quantities arise from the internal ratio between the ternary opening of a
six-trit block, \(3^6=729\), and decimal compactification
\(10^3=1000\). They are not introduced as parameters of an externally chosen
rotation.

First, \(\delta\) is irrational. If
\(\delta=p/q\in\mathbb Q\), with \(q>0\), then

\[
 10^{p/q}=\frac{10}{9}
 \quad\Longrightarrow\quad
 9^q10^p=10^q.
\]

The equality contradicts the uniqueness of prime factorization: the left-hand
side contains the factor \(3\), while the right-hand side contains only
\(2\) and \(5\). Second, \(\delta\) cannot be algebraic irrational. In
that case, the Gelfond--Schneider theorem would make
\(10^\delta\) transcendental, but

\[
 10^\delta=\frac{10}{9}
\]

is algebraic. One concludes

\[
 \boxed{\delta\ \text{is transcendental}.}
\]

This classical recognition is applied after constructing
\(\delta\) from APP--TRIT--TPK; it does not act as a generator of the calendar.
Its structural consequences are immediate. The continued fraction of
\(\delta\) is not eventually periodic, because an eventually periodic continued
fraction represents a quadratic irrational. Therefore, the
renormalization

\[
 [0;a_1,a_2,a_3,\ldots]
 =[0;21,1,5,1,6,2,3,1,1,8,\ldots]
\]

does not close on a stationary Sturmian substitution. The word admits
an \(S\)-adic description directed by the partial quotients, but that
directive does not repeat periodically. The crystal is ordered and recurrent,
without being the fixed point of a single quadratic inflation.

The same conclusion can be formulated as the absence of a nontrivial modular
stabilizer. If

\[
 M=\begin{pmatrix}a&b\\c&d\end{pmatrix}\in
 \operatorname{GL}_2(\mathbb Z)
\]

fixed the slope through a fractional transformation,

\[
 \frac{a\delta+b}{c\delta+d}=\delta,
\]

then

\[
 c\delta^2+(d-a)\delta-b=0.
\]

The transcendence of \(\delta\) forces
\(c=b=0\) and \(d=a\). Since \(\det M=\pm1\), the resulting projective action
is the identity. Therefore,

\[
 \boxed{\operatorname{Stab}_{\operatorname{PGL}_2(\mathbb Z)}(\delta)
 =\{1\}.}
\]

Nonrepetition of renormalization does not eliminate long-range order.
The mechanical word

\[
 z_n=\mathbf 1_{[1-\delta,1)}
 \bigl(\{n\delta+\beta\}\bigr)
\]

is a regular model set: it comes from cutting a lattice orbit by
an interval window whose boundary has zero measure. Its hull is
minimal and uniquely ergodic, has complexity \(p(m)=m+1\), zero topological
entropy and pure dynamical spectrum. The expression
“aperiodic crystal” is thus made precise: diffractive order and uniform recurrence, together with
absence of translational periodicity and stationary renormalization.

\textbf{Refutation criterion.} A period of the
word, an eventually periodic continued fraction or a nontrivial modular
transformation fixing \(\delta\).

\section{The three channels as cuts of the same window over \(729\) extensions}

The synchronization of \(\pi\), \(e\) and \(\varphi\) can be described without
reducing it to equality of censuses. At level \(t\), let \(T_{c,t}\) be the integer
represented by the first \(6t\) ternary digits of channel
\(c\in\{\pi,e,\varphi\}\), and let \(D_{c,t}\) be the integer represented by its
first \(K_t\) decimal blocks. Each local extension is identified with
an integer

\[
 u\in\{0,1,\ldots,728\}.
\]

The corresponding ternary cell is

\[
 \left[
 \frac{729T_{c,t}+u}{3^{6(t+1)}},
 \frac{729T_{c,t}+u+1}{3^{6(t+1)}}
 \right),
\]

while the decimal cylinder is

\[
 \left[
 \frac{D_{c,t}}{1000^{K_t}},
 \frac{D_{c,t}+1}{1000^{K_t}}
 \right).
\]

Let us define the common scale and the internal displacement

\[
 \Omega_t=\frac{3^{6(t+1)}}{1000^{K_t}},
 \qquad
 \xi_{c,t}=D_{c,t}\Omega_t-729T_{c,t}.
\]

The two cylinders intersect exactly when

\[
 u<\xi_{c,t}+\Omega_t,
 \qquad
 u+1>\xi_{c,t}.
\]

Consequently, the set of surviving extensions is

\[
 I_{c,t}
 =\{0,\ldots,728\}\cap
 \bigl[\lfloor\xi_{c,t}\rfloor,
       \lceil\xi_{c,t}+\Omega_t\rceil-1\bigr]_{\mathbb Z},
\]

and the determined census satisfies

\[
 \boxed{N_c(t)=|I_{c,t}|.}
\]

The formula provides the geometric mechanism of synchronization. The
three channels receive the same width \(\Omega_t\); they differ only in the
translation \(\xi_{c,t}\) of the window. The equality
\(N_\pi=N_e=N_\varphi\) means that three translated intervals cut the
local domain of \(729\) cells with the same cardinality. It does not imply that
the intervals, words or their histories are equal.

In the terminal state immediately preceding the first vacancy,
one obtains the integer intervals

\[
\begin{array}{c|c|c}
c&t=20:\ \text{raw interval}&I_{c,20}\\
\hline
\pi&[407,408]&[407,408]\\
e&[172,173]&[172,173]\\
\varphi&[313,314]&[313,314].
\end{array}
\]

The three intervals occupy different positions, but contain two cells.
This is the exact form of the synchronized state
\((2,2,2)\): equality of discrete measure with distinct positional memory.

The next vacancy preserves \(K_{21}=K_{20}=20\) and multiplies the scale
by \(729\). The raw intervals become

\[
\begin{array}{c|c|c|c}
c&t=21:\ \text{raw interval}&I_{c,21}&N_c(21)\\
\hline
\pi&[118,1073]&[118,728]&611\\
e&[-429,526]&[0,526]&527\\
\varphi&[6,961]&[6,728]&723.
\end{array}
\]

The three raw intervals have the same integer width, \(956\), but
are clipped by opposite boundaries of the finite fiber. If

\[
 \operatorname{Def}_*
 :=729(1,1,1)-\bigl(N_\pi(21),N_e(21),N_\varphi(21)\bigr),
\]

then

\[
 \operatorname{Def}_*=(118,202,6)
 =6(1,1,1)+28(4,7,0).
\]

The diagonal term \(6(1,1,1)\) records the common loss of a six-trit
block. Taking the quotient by diagonal survival leaves the oriented
divisor

\[
 [\operatorname{Def}_*]_{A_2}=28(4,7,0).
\]

Equivalently,

\[
 \partial N_{21}
 =(84,-196,112)
 =28(3,-7,4).
\]

The slope

\[
 \frac{N_e-N_\varphi}{N_\varphi-N_\pi}
 =-\frac{196}{112}
 =-\frac74
\]

is therefore the ratio of two oriented clippings of a common window.
It is not a fraction added to the analytic kernel.

The full recalculation of the \(171\) levels available in the reference construction
shows two finite uniqueness properties: \(t=21\) is the only vacancy preceded by
a chain of five consecutive synchronizations, and it is the only level whose
projective defect in that chart is \(-7/4\). The assertion is exact in the
domain of definition; its coinductive extension to arbitrary depth
will have to explicitly preserve the marked state.

\textbf{Refutation criterion.} The censuses not
coinciding with the interval intersections, the three raw widths
at \(t=21\) differing, or another slope \(-7/4\) appearing in the declared
domain.

\section{Specular transit matrix and joint closure of \(21/22\), \(6/7\), \(137\) and the lattice \(369\)–\(126\)}

The four phases preceding the first triple synchronization should not
be read as four anecdotes. Their defect amplitudes form the matrix

\[
 D_{\rm syn}
 =\begin{pmatrix}
 85&112\\
 41&52
 \end{pmatrix}.
\]

The synchronized census of the ninth phase is \(N_9=43\). Therefore,

\[
 \det D_{\rm syn}=-172=-4\cdot43
\]

defines internally

\[
 a:=-\frac{\det D_{\rm syn}}{N_9}=4.
\]

Through an independent branch, inducing the vacancy word on
its short gaps produces the secondary spectrum
\(\mathcal R_{\rm sec}=\{6,7\}\). Let

\[
 b:=\max\mathcal R_{\rm sec}=7.
\]

The first reopening then satisfies

\[
 t_*=3b=21,
 \qquad
 K_*=3b-1=20,
\]

where \(3\) is the cardinality of the TRIT regimes. Moreover,

\[
 \left\lfloor1000\delta\right\rfloor=45,
 \qquad
 \left\lceil1000\delta\right\rceil=46
 =2(3b)+a.
\]

The clipping decomposition of the preceding section now takes the closed
form

\[
 \boxed{
 \operatorname{Def}_*
 =6(1,1,1)+ab(a,b,0).}
\]

Indeed,

\[
 6(1,1,1)+4\cdot7(4,7,0)
 =(118,202,6).
\]

This equality is a compatibility square between two independent
extractions: \(a=4\) comes from the synchronization determinant,
\(b=7\) comes from the induced clock and the common clipping finds them again
as boundary weights. The fraction \(-b/a=-7/4\) is thus determined before
solving any equation for \(\alpha\).

The same matrix brings together Diophantine renormalization. The two secondary
returns give

\[
 q_6=22\cdot6-1=131,
 \qquad
 q_7=22\cdot7-1=153,
\]

and satisfy

\[
 6q_7-7q_6=1.
\]

Within \(D_{\rm syn}\), those denominators reappear as

\[
 q_6=(85+41)+5=131,
 \qquad
 q_7=112+41=153,
\]

where \(5\) is the upper branch of the half-imbalance \(4/5\) of
\(\widetilde W_{100}\). Finally,

\[
 \boxed{\operatorname{tr}D_{\rm syn}=137=q_6+6.}
\]

Thus, \(137\) is not obtained here by rounding
\(\alpha^{-1}\). It is an integer direction of synchronization transit and,
simultaneously, the return denominator \(131\) prolonged by a block
of six trits. This does not replace the two constructions of \(\alpha\): it establishes
a common integer precursor that the dodecaphasic closure and the analytic kernel
must transport.

The radical--electronic lattice is also reconstructed from
\(D_{\rm syn}\). Let us define

\[
 A_{\rm surv}:=85+41=126,
 \qquad
 B_{\rm bord}:=112-85=27,
\]

and

\[
 T_{\rm tri}:=3(A_{\rm surv}+6)=396.
\]

The matrix reader

\[
 \mathcal R(D_{\rm syn})
 :=
 \begin{pmatrix}
 T_{\rm tri}-B_{\rm bord}&T_{\rm tri}\\
 T_{\rm tri}-A_{\rm surv}-B_{\rm bord}&
 T_{\rm tri}-A_{\rm surv}
 \end{pmatrix}
\]

produces exactly

\[
 \boxed{
 \mathcal R(D_{\rm syn})
 =\begin{pmatrix}369&396\\243&270\end{pmatrix}.}
\]

Consequently, the four numerical vertices of the parallelogram are no longer
collected as an independent set: they come from an integer reader of
channel synchronization, the six-trit block and tritic
amplification. The identification of each vertex with its radical or
electronic object nevertheless preserves the domain and the transducer of its
domains of origin; equality of the numerical matrix does not erase those fibers.

\textbf{Refutation criterion.} \(a\) not being an integer, the induced
word not producing \(b=7\), failure of the clipping law, or
\(\mathcal R(D_{\rm syn})\) not reproducing the four vertices.

\section{Internal factorization of all the coefficients of the jet \(P_9\)}

The preceding results make it possible to bring together, for the first time within this
development, the genealogy of all the coefficients of the jet without introducing
its root. Let

\[
 \omega:=2\pi,
 \qquad
 m:=3,
 \qquad
 a:=4,
 \qquad
 b:=7,
\]

and let us define

\[
 t_*=mb=21,
 \qquad
 K_*=mb-1=20,
\]

\[
 v_-:=\lfloor1000\delta\rfloor=45,
 \qquad
 v_+:=\lceil1000\delta\rceil=46.
\]

The lower branch simultaneously satisfies

\[
 v_-=45=\det
 \begin{pmatrix}7&2\\2&7\end{pmatrix}.
\]

The internal difference of the second row of \(D_{\rm syn}\) gives

\[
 r_{\rm dod}:=52-41=11,
\]

which coincides with the dodecaphasic transverse rank already demonstrated. The
phase half-crown contributes

\[
 f_{\rm ph}=8\binom62=120,
\]

and

\[
 v_-r_{\rm dod}=45\cdot11=495.
\]

On the invariant signature

\[
 \mathcal I_9
 =(\omega,m,a,b,K_*,t_*,v_-,v_+,f_{\rm ph},r_{\rm dod}),
\]

we define the graded reader

\[
\begin{aligned}
 \mathfrak E_9(\mathcal I_9)(x)
 :={}&
 \omega x-\frac ba x^2+\omega^{-1}x^3
 +\frac{x^4}{K_*}-\frac{2x^5}{t_*}
 -\frac{x^6}{v_+}\\
 &-\frac{x^7}{f_{\rm ph}}
 +\frac{x^8}{v_-}
 +\frac{2x^9}{v_-r_{\rm dod}}.
\end{aligned}
\]

Substitution of the invariants produces

\[
\boxed{
\begin{aligned}
\mathfrak E_9(\mathcal I_9)(x)
={}&2\pi{}x-\frac74x^2
+\frac{x^3}{2\pi}
+\frac{x^4}{20}-\frac{2x^5}{21}-\frac{x^6}{46}\\
&-\frac{x^7}{120}+\frac{x^8}{45}
+\frac{2x^9}{495}
=P_9(x).
\end{aligned}}
\]

The formula reveals a relation that remained hidden when
\(P_6\) and \(T_{7:9}\) were presented separately: the vacancy pair \(45/46\) crosses the boundary
between them. The upper branch appears in
\(-x^6/46\), while the lower branch reappears in \(+x^8/45\) and,
coupled to the rank \(11\), in \(2x^9/495\). The prolongation does not add a
second arithmetic; it continues the same balanced calendar through the
nilpotent resolvent.

The typed diagram is

\[
 \mathcal X_{\rm TPK}^{\rm enr,*}
 \xrightarrow{\ \mathcal A_*\ }
 \mathcal I_9
 \xrightarrow{\ \mathfrak E_9\ }
 \mathbb Q(\pi)[x]/(x^{10}),
\]

where the asterisk indicates that the object preserves the coinductive origin, the
order of the channels and the orientation of their sheets. The extractor
\(\mathcal A_*\) performs the demonstrated operations:

\[
 D_{\rm syn}\mapsto a,\qquad
 \mathcal R_{\rm sec}\mapsto b,\qquad
 \operatorname{Def}_*\mapsto(a,b),
\]

\[
 \delta\mapsto(v_-,v_+),\qquad
 B_c\mapsto v_-,\qquad
 (D_3,D_4)\mapsto r_{\rm dod}.
\]

Every morphism that preserves the marked state, the intersection windows,
the channel orientation, the capacity clock and the two-way involution
preserves \(\mathcal I_9\), and therefore makes the reader commute. What has been
closed is the internal factorization of the coefficient signature. The
stronger assertion that \(\mathfrak E_9\) is the only possible reader
in an unpointed category requires classifying all automorphisms that
can forget or permute the origin and orientation. That reservation is
concrete and reopens neither the simple root of \(E_{21/22}\), nor the nilpotent
resolvent, nor the genealogy of the coefficients.

\textbf{Refutation criterion.} Any coefficient requiring the value
of the root, two declared extractions producing different values, or
an allowed automorphism changing \(\mathcal I_9\).

\section{The factor \(28\), the exact projection \(36\times28\) and the partition into seven tetrads}

The primitive scale of the first defect is

\[
 \lambda_*:=\gcd(84,196,112)=28=ab.
\]

In the constructed catalog of the region of \(\pi\), the oriented microfiber
has multiplicity

\[
 M(w_\pi)=432+4\cdot144=1008
\]

and admits the factorizations

\[
 1008=36\cdot28
 =\binom92\binom82
 =7\cdot144.
\]

Therefore,

\[
 \boxed{M(w_\pi)=36\,\lambda_*}
\]

is an exact scalar compatibility between the clipping defect of the
first reopening and the multiplicity of the microfiber of \(\pi\). The
factor \(36\) also belongs to \(R_{36}\) and to the set of pairs of a
nonad; \(28\) is the number of pairs of a set of eight elements.

The initial reading of this equality as mere scalar compatibility was
incomplete. The reference construction already constructs on the microfiber

\[
 \mathcal R_{010211}
 =\{(P,T):w_6(U_6(P,T))=010211\}
\]

an explicit projection. For

\[
 P=(i,j,d_P),\qquad T=(k,l,d_T),
\]

the directional gauge is fixed

\[
 \beta(N)=1,\quad\beta(E)=2,\quad
 \beta(S)=3,\quad\beta(O)=4
\]

and one defines

\[
 x(P)=5(i-j)\pmod 9,
 \qquad
 h(T)=5\beta(d_T)\pmod9,
\]

\[
 \boxed{
 q_\beta(P,T)=\{x(P)-h(T),x(P)+h(T)\}.}
\]

The map

\[
 q_\beta:\mathcal R_{010211}\longrightarrow
 \binom{\mathbb Z/9\mathbb Z}{2}
\]

is surjective, has a base of cardinality \(36\) and all its fibers have
cardinality \(28\). Each fiber receives, in the order of the five regions of
\(\pi\), exactly

\[
 12+4+4+4+4=28
\]

paths. It is not obtained by dividing \(1008\) by \(36\): the certificate
reconstructs the \(104,976\) seed pairs, computes each image and proves
the uniformity of the fibers.

The additive translation of \(C_9\) and its reflection preserve \(U_6\) and make
\(q_\beta\) equivariant with respect to the dihedral action on pairs. On the
multiplicative cursor, two involutions generate a free \(V_4\) that preserves
the complete signature. The group

\[
 H_0=C_9\times V_4
\]

The ratio

\[
 \Omega=\mathcal R_{010211}/H_0,
 \qquad |\Omega|=28,
\]

has an additional intrinsic partition:

\[
 \Omega=
 \mathcal B_3\sqcup\mathcal B_6\sqcup\mathcal B_9
 \sqcup\mathcal L_1\sqcup\mathcal L_2
 \sqcup\mathcal L_3\sqcup\mathcal L_4,
\]

with seven blocks of four elements. The three blocks
\(\mathcal B_3,\mathcal B_6,\mathcal B_9\) are transverse and are
distinguished by the multiplicative residue \(3,6,9\); the four blocks
\(\mathcal L_j\) come from the four lateral regions. The reduced
census vector of this partition is

\[
 v_\Omega=(3,-7,4)\in A_2.
\]

The first reopening of the channels satisfies, exactly,

\[
 \boxed{
 \partial N_{21}=28v_\Omega.}
\]

Thus, the direction \((3,-7,4)\) does not merely contain the slope \(-7/4\): it is the
conservation relation between three transverse tetrads, seven total
tetrads and four lateral tetrads. Moreover,

\[
 28=7\cdot4,
 \qquad
 1008=36\cdot7\cdot4.
\]

The integers \(a=4\) and \(b=7\), obtained earlier from the synchronization
determinant and the induced return, now reappear as block size and
number of blocks of the microfiber.

This identification has decisive local rigidity. The integral group of
isometries of \(A_2\), represented by permutations of the three coordinates
and global sign change, has order \(12\). Since the coordinates
\(3,-7,4\) are distinct and the multiset \(\{3,-7,4\}\) does not coincide with
its negative, only the identity fixes \(v_\Omega\). Therefore,

\[
 \boxed{
 \operatorname{Stab}_{\operatorname{Aut}(A_2)}(v_\Omega)=\{1\}.}
\]

Once the first reopening is characterized, the defect itself recovers the
order of the three channels and their orientation: it needs no external mark.

The octadic boundary is typed by the seven tetrads. After choosing a
gauge, bijective charts

\[
 \Phi:\Omega\overset\sim\longrightarrow
 \binom{\mathbb F_2^3}{2}
\]

and with them an exact incidence \(J(8,2)\). However, those charts are not
selected by the source relations tested. The reflection stabilizing
each fiber acts freely with cycle type \(2^{14}\), whereas every
involution of \(\operatorname{Aut}J(8,2)\cong S_8\) fixes at least four
pairs. The proved no-go follows:

\[
 \boxed{
 \text{there is no family of charts }J(8,2)
 \text{ equivariant under every }D_9\times V_4.}
\]

The oriented datum \(v_\Omega\), whose stabilizer is trivial, provides the
reflection breaking required by this obstruction theorem. Consequently,
the Johnson charts on each fiber \(F_b=q_\beta^{-1}(b)\) are conjugate gauge
charts and not a canonical selection under the complete action
\(D_9\times V_4\). The projection \(36\times28\) remains intrinsic; the
choice of an octadic chart is not part of its definition.

A temporal ambiguity is also corrected. The integer \(1008\) is not
the period of the dodecaphasic clock:

\[
 1008\equiv36\pmod{108}.
\]

After \(1008\) steps the nonadic phase returns, but the clock \(C_{108}\)
advances by \(36\) positions. On this front, \(1008\) is a microfiber
multiplicity and a combinatorial incidence; \(108\) is the finite temporal
period. The projection \(q_\beta\) explains the product \(36\times28\); an
octadic gauge chart and the clock map are distinct operations.

\textbf{Refutation criterion.} The seed sweep changing a fiber or the partition
\(7\times4\), a nontrivial isometry of \(A_2\) fixing \(v_\Omega\),
the existence of a completely \(D_9\times V_4\)-equivariant Johnson chart, or
\(1008\) being used as the period of \(C_{108}\).

\section{Two clocks, an integrated memory and two spectral modules}

Let

\[
 Z_t:=\sum_{j=0}^{t}z_j
 =\lfloor(t+1)\delta\rfloor.
\]

Since \(\rho=1-\delta\) and
\((t+1)\delta\notin\mathbb Z\),

\[
\begin{aligned}
 K_t
 &=\lfloor(t+1)\rho\rfloor\\
 &=\lfloor(t+1)-(t+1)\delta\rfloor\\
 &=t-\lfloor(t+1)\delta\rfloor
 =t-Z_t.
\end{aligned}
\]

If

\[
 \theta_t=[t]_q
 \quad\text{and}\quad
 g_t=[K_t]_q,
\]

then

\[
 \boxed{g_t=\theta_t-[Z_t]_q.}
\]

The capacity phase is the uniform phase corrected by the integrated memory
of vacancies. This identity expresses exactly the difference between
phase return and state return: knowing \(\theta_t\) does not determine
\(g_t\) if \(Z_t\) has been forgotten.

The distinction produces two legitimate spectral modules. On the block
clock, the translation

\[
 S_q(\theta,u)=(\theta+1,u+\delta)
\]

on \(C_q\times\mathbb T\) has additive frequencies

\[
 \mathcal M_q^{\rm blk}
 =\frac1q\mathbb Z+\delta\mathbb Z
 =\frac1q\mathbb Z+\rho\mathbb Z
 \pmod1.
\]

In particular,

\[
 \mathcal M_9^{\rm blk}
 =\frac19\mathbb Z+\log_{10}9\,\mathbb Z,
\]

and

\[
 \mathcal M_{108}^{\rm blk}
 =\frac1{108}\mathbb Z+\log_{10}9\,\mathbb Z.
\]

This is the exact sense in which the internal decimal logarithm appears
as a spectral module: not as positive entropy, but as an irrational
generator of a dense group of pure frequencies.

On the glued capacity clock, the dynamics \(F_q\) of section 49 is
conjugate to

\[
 y\longmapsto y+\rho
 \quad\text{on}\quad
 \mathbb R/q\mathbb Z.
\]

Its frequencies are

\[
 \mathcal M_q^{\rm cap}
 =\frac{\rho}{q}\mathbb Z\pmod1.
\]

There is no contradiction between the two modules: they belong to two algebras of
observables of the same integral holonomic state. The first preserves the block
index; the second absorbs the accumulated carry into the finite coordinate.
The relation \(g_t=\theta_t-Z_t\) is the gauge change that compares them.

This also clarifies two returns mentioned earlier. The denominator
\(170748=1581\cdot108\) closes the block clock, but

\[
 K_{170748-1}\equiv71\pmod{108},
\]

so it does not close the capacity clock. Likewise, \(1008\) closes
\(C_9\), but shifts \(C_{108}\) by \(36\). The integral holonomic state retains
height and signature in both cases, so neither of those returns
equates to restarting the dynamics.

\textbf{Refutation criterion.} That
\(K_t\ne t-Z_t\), the same spectrum being attributed to two distinct domains of
observables, or a finite return erasing memory.
 
\chapter{Joint nonadic connection and biography of return with memory}
\label{chap:83-07}

The connection is constructed from the complete orbital classification and from
its first coinductive elevation. The following development fixes the three
initial histories, the oriented boundaries and the chain
\(w_6\to w_{30}\to R_{36}\) before introducing the nonadic holonomy.

\begin{HMTScientificOwnerSubordinated}
% Unidad U008. Clasificación orbital y elevación; redacción autónoma.

\chapter{Orbital classification, initial words and 1st boundary}
\label{chap:orbitas-elevacion-r36}

The ternary reduction of the emitter realizes 243 words within the ambient
fiber \(\mathbb F_3^6\). This unit does not choose three words by
decimal comparison. It first constructs the symmetry action of the catalogue,
classifies its orbits by integer invariants, and only afterwards orients
the three structural representatives.

\section{Dihedral action on \(6\) coordinates}

For \(u=(u_1,u_2,u_3,u_4,u_5,u_6)\in\mathbb F_3^6\), define
\begin{align*}
 r(u)&=(u_3,u_1,u_2,u_5,u_6,u_4),\\
 s(u)&=(u_4,u_5,u_6,u_1,u_2,u_3).
\end{align*}
We verify
\[
 r^3=s^2=1,
 \qquad srs=r^{-1}.
\]
Therefore, \(r,s\) generate an action of \(D_3\).

The same permutation acts on the six blocks of each emission
\(U=(U_1,\ldots,U_6)\). The projection \(\Psi\) satisfies
\[
 \Psi(gU)=g\Psi(U),
 \qquad g\in D_3.
\]

\begin{theorem}[Orbital decomposition of the catalog]
\label{thm:catalogo-descomposicion-d3}
The 243 realized words decompose into 43 orbits:
\[
 38\text{ orbits of cardinality }6,
 \qquad
 5\text{ orbits of cardinality }3.
\]
In particular,
\[
 38\cdot6+5\cdot3=243.
\]
\end{theorem}

\begin{proof}
The set is applied to each word.
\[
 \{u,ru,r^2u,su,sru,sr^2u\}
\]
and repetitions are removed. Exhaustive enumeration of the catalog produces
the indicated histogram. The equivariance of \(\Psi\) ensures that no
orbit leaves the realized image.
\end{proof}

\section{Fiber Invariants}

For a word realized \(w\), let
\[
 \mathcal F_w:=\{U\in\operatorname{im}T_{\min}:\Psi(U)=w\},
\]
\[
 n(w):=|\mathcal F_w|,
 \qquad
 M(w):=\sum_{U\in\mathcal F_w}\operatorname{mult}(U),
\]
and
\[
 c(w)=\bigl(\#\{j:w_j=0\},\#\{j:w_j=1\},\#\{j:w_j=2\}\bigr).
\]
The enriched catalogue further preserves, for each emission, a diagonal
coordinate \(\operatorname{diag}_c(U)\in\{0,\ldots,8\}\), a cardinal
orientation \(h_+(U)\), and the multiplicative signature \(E_{\rm sig}(U)\). All
of them are computed from the seed route; they are not read from a target
constant. Their componentwise algorithm will be printed alongside the complete
catalogue in the certification appendix and forms part of the input state
of the following predicates.

\begin{definition}[Closure orbit]
\label{def:orbita-clausura}
An orbit \(\mathcal O\) is of closure if it has cardinality six and, for
each \(w\in\mathcal O\),
\[
 n(w)=5,
 \qquad
 M(w)=1008,
\]
\[
 \{\operatorname{mult}(U):U\in\mathcal F_w\}
 =\{432,144,144,144,144\},
\]
and the common tritic census is \(c(w)=(2,3,1)\).
\end{definition}

\begin{definition}[Radical Minimal Orbits]
\label{def:orbitas-minimas-radicales}
An orbit \(\mathcal O\) is minimally radical if it has cardinality six, each
word possesses a unique emission of multiplicity \(144\), and
\[
 \{\operatorname{diag}_c(U):\Psi(U)\in\mathcal O\}=\{0,3,6\}.
\]
Among them, the orbit with census \((2,3,1)\) is called the propagation
orbit, and the orbit with census \((2,2,2)\) is called the self-scaling orbit.
\end{definition}

\begin{theorem}[Uniqueness of the Three Structural Orbits]
\label{thm:unicidad-tres-orbitas-estructurales}
In the catalogue of 468 emissions there exists a unique closure orbit, a
unique minimally radical orbit of propagation, and a unique minimally radical
orbit of self-scale.
\end{theorem}

\begin{proof}
The 43 orbits of \cref{thm:catalogo-descomposicion-d3} are enumerated. For each one, \(n(w)\), \(M(w)\), the multiset of multiplicities, the census \(c(w)\), and the diagonal support are computed. The predicates of \cref{def:orbita-clausura,def:orbitas-minimas-radicales} leave exactly one orbit in each class. The proof is exhaustive over a finite set and does not use decimal digits of the analytic characters.
\end{proof}

The three orbits are
\begin{align*}
\mathcal O_{\rm cl}
 &=\{001112,010211,100121,112001,121100,211010\},\\
\mathcal O_{\rm prop}
 &=\{011120,012110,101201,110012,120011,201101\},\\
\mathcal O_{\rm aut}
 &=\{002112,020211,112002,121200,200121,211020\}.
\end{align*}

\section{Internal Orientation of Each Orbit}

We fix the internal caliber.
\[
 \operatorname{diag}_c(U)=6,
 \qquad h_+(U)=ES.
\]
A word in an orbit is admissible in that gauge if some emission in
its fiber satisfies both conditions.

\begin{proposition}[Unique oriented representative]
\label{prop:representante-orientado-unico}
Each of the three structural orbits contains a unique word
admissible in the gauge \((6,ES)\). They are
\[
 \boxed{
 w_6^\pi=010211,
 \qquad
 w_6^e=201101,
 \qquad
 w_6^\varphi=121200.}
\]
\end{proposition}

\begin{proof}
The enumeration of the catalogue fibers counts the words containing
an emission with \(\operatorname{diag}_c=6\) and \(h_+=ES\). In each of
the three orbits, the number of coincidences is one. The superscripts
\(\pi,e,\varphi\) henceforth name the structural functions of
closure, propagation, and self-scaling; they did not participate in the selection.
\end{proof}

\section{Closure microfiber}

The fiber of the representative \(w_6^\pi\) consists of five emissions:
\[
\begin{array}{c|c|c|c|c}
U_6&E_{\rm sig}&\operatorname{mult}&\operatorname{diag}_c&\text{stratum}\\\hline
501|614|498|169|272|272&555555&432&4&\perp\\
810|923|870|169|272|272&339555&144&6&++\\
810|983|810|169|272|272&393555&144&6&++\\
870|923|810|169|272|272&933555&144&6&++\\
870|923|810|418|674|521&933717&144&5&--
\end{array}
\]
Thus,
\[
 432+4\cdot144=1008,
 \qquad
 5=1_\perp+3_{++}+1_{--}.
\]
The unipolar emissions oriented of the other two channels are
\[
 U_e=850|963|890|149|252|212,
\]
\[
 U_\varphi=830|943|830|109|252|252.
\]
The mass \(1008\) of this microfiber is not the period of the calendar
\(C_{108}\); the two numbers belong to different domains.

\section{Ternary lifting matrices}

On row vectors of \(\mathbb F_3^6\), we define
\[
L_0=
\begin{pmatrix}
2&2&2&1&2&1\\
2&1&2&2&1&1\\
1&0&0&1&0&2\\
0&0&1&1&1&0\\
2&2&2&0&2&1\\
2&1&2&1&0&0
\end{pmatrix},
\qquad
L_1=
\begin{pmatrix}
2&1&2&0&2&2\\
1&2&1&1&2&0\\
1&2&1&1&1&2\\
0&1&0&0&2&1\\
2&0&2&1&0&1\\
0&2&2&2&1&1
\end{pmatrix}.
\]
The elevation schedule is \((L_0,L_0,L_1,L_1)\).

\begin{definition}[Block recurrence and concatenation]
\label{def:recurrencia-bloques-l0-l1}
For each channel \(\chi\in\{\pi,e,\varphi\}\), let
\(b_0^\chi=w_6^\chi\) and
\[
 b_{k+1}^\chi=b_k^\chi L_{\epsilon_k}\pmod3,
 \qquad
 (\epsilon_0,\epsilon_1,\epsilon_2,\epsilon_3)=(0,0,1,1).
\]
The accumulated word is
\[
 w_{30}^\chi
 =b_0^\chi\Vert b_1^\chi\Vert b_2^\chi
  \Vert b_3^\chi\Vert b_4^\chi.
\]
\end{definition}

This formulation distinguishes the current block \(b_k\in\mathbb F_3^6\) from
the concatenated word \(w_{6(k+1)}\). A word of length \(6k\) is not
multiplied by a matrix \(6\times6\).

\begin{theorem}[Ternary biographies of thirty positions]
\label{thm:palabras-w30}
The recurrence of the \cref{def:recurrencia-bloques-l0-l1} produce
\begin{align*}
w_{30}^{\pi}
 &=010211|012222|010211|002111|110221,\\
w_{30}^{e}
 &=201101|121221|102011|012222|102011,\\
w_{30}^{\varphi}
 &=121200|112202|121020|010210|010200.
\end{align*}
\end{theorem}

\begin{proof}
The five row vectors are successively multiplied by the indicated
matrices, reducing each coordinate modulo three after each product.
The concatenation of the five computed blocks gives the three expressions.
Each equality can be verified component by component by means of 24
scalar products per channel.
\end{proof}

The environmental indices of the five blocks are
\[
\begin{array}{c|rrrrr}
&b_0&b_1&b_2&b_3&b_4\\\hline
\pi&103&161&103&67&349\\
e&523&457&301&161&301\\
\varphi&450&398&438&102&99
\end{array},
\]
where each word of six trits is interpreted as an integer in base three.

\section{1st common boundary}

The boundary is obtained through a finite relation that preserves axis, aggregate,
saturation, column occupancy, and dual neutrality. The incidence matrix
used in this relation is
\begin{equation}
 A_W=
 \begin{pmatrix}
 0&1&1&1&1&1\\
 1&0&1&2&2&1\\
 1&1&0&1&2&2\\
 1&2&1&0&1&2\\
 1&2&2&1&0&1\\
 1&1&2&2&1&0
 \end{pmatrix}\in M_6(\mathbb F_3).
 \label{eq:u008-aw-frontera}
\end{equation}
If \(u_4^\chi\) is the fifth block of \(w_{30}^\chi\), the axis and the
aggregate prior to separating the two lateral channels are
\begin{equation}
 u_5^\varphi=u_4^e=102011,
 \qquad
 u_5^\pi+u_5^e=210112,
 \qquad
 u_4^\varphi A_WL_1^3=111101.
 \label{eq:u008-eje-agregado-r36}
\end{equation}
From them, they calculate
\begin{equation}
 q=012120,\qquad a=111101,\qquad
 qA_W=012000,\qquad aA_W=120210.
 \label{eq:u008-q-a-r36}
\end{equation}
Depth six imposes the saturation \(c=111111\); the minimal compatible
occupancy is \(\operatorname{colw}=223232\); and dual neutrality requires
\begin{equation}
 (BA_W)\mathbf1=000.
 \label{eq:u008-neutralidad-dual-r36}
\end{equation}

\begin{proposition}[Enumeration of the first boundary]
\label{prop:u008-enumeracion-r36}
The conditions \eqref{eq:u008-eje-agregado-r36}--
\eqref{eq:u008-neutralidad-dual-r36} leave exactly two matrices:
\begin{equation}
 B_-=
 \begin{pmatrix}021222\\222220\\102011\end{pmatrix},
 \qquad
 B_+=
 \begin{pmatrix}222220\\021222\\102011\end{pmatrix}.
 \label{eq:u008-dos-fronteras}
\end{equation}
The mirror involution swaps its first two rows. Its internal cards are, respectively,
\begin{equation}
 B_-\longmapsto(789,048,894),
 \qquad
 B_+\longmapsto(793,045,894).
 \label{eq:u008-cartas-frontera}
\end{equation}
The orientation APP of the cylinder semiabierto selecciona \(B_+\).
\end{proposition}

\begin{proof}
The axis fixes the third row and the aggregate fixes the sum of the first two.
Saturation and occupation eliminate the distributions that do not fill the
six columns. The enumeration of the solutions of the six linear equations
of \eqref{eq:u008-neutralidad-dual-r36} leaves the two matrices of
\eqref{eq:u008-dos-fronteras}; the semi-open chart and the orientation
distinguish the second without consulting a target decimal expansion.
\end{proof}

We therefore define,
\begin{equation}
 R_{36}:=B_+=
 \begin{pmatrix}
 222220\\
 021222\\
 102011
 \end{pmatrix},
 \qquad
 \operatorname{ind}_3(R_{36})=(726,215,301).
 \label{eq:u008-r36-construida}
\end{equation}
The object \(R_{36}\) contains three ordered blocks of six trits. It is not the
predictive quotient \(C_{36}^{\rm pred}\) of
\cref{thm:tpk-cociente-predictivo-36} or the excess of 36 classes of the seed
gauge.

\section{Residue--quotient lifting}

Let \(A\in M_6(\mathbb Z)\) be the integer lift of \(L_0\):
\[
 A=
\begin{pmatrix}
2&2&2&1&2&1\\
2&1&2&2&1&1\\
1&0&0&1&0&2\\
0&0&1&1&1&0\\
2&2&2&0&2&1\\
2&1&2&1&0&0
\end{pmatrix},
\qquad \det A=1.
\]
For \(x_k\in\mathbb Z^6\), define
\[
 y_k=x_kA,
 \qquad
 u_k=y_k\bmod3\in\{0,1,2\}^6,
 \qquad
 x_{k+1}=\frac{y_k-u_k}{3}.
\]
Then
\begin{equation}
 x_kA=u_k+3x_{k+1}.
\label{eq:bockstein-hensel-local}
\end{equation}

\begin{theorem}[Reversibility of the lift]
\label{thm:reversibilidad-bockstein-hensel}
The pair \((u_k,x_{k+1})\) is uniquely determined by \(x_k\), and
\[
 x_k=(u_k+3x_{k+1})A^{-1}.
\]
\end{theorem}

\begin{proof}
Euclidean division by three in each coordinate of \(x_kA\) determines
residue and quotient. Since \(\det A=1\), the inverse \(A^{-1}\) has
integer coefficients. Multiplying \cref{eq:bockstein-hensel-local} by it
recovers \(x_k\).
\end{proof}

In channel \(e\), the visible block \(102011\) reappears twice, but
the prestates modulo 27 are
\[
 (25,11,7,17,8,16)
 \neq
 (7,8,4,23,26,7),
\]
and the subsequent quotients are
\[
 (6,13,9,2,13,1)
 \neq
 (15,0,3,19,23,25).
\]
This explicitly verifies that the visible return of a word is not
the return of the integral holonomic state.

\section{Obstructions to orbital selection and to the lifting \(w_6\to w_{30}\to R_{36}\): uniqueness, gauge, and Hensel reversibility}

\begin{criterion}[Orbital selection refutation criteria]
The unit is refuted if:
\begin{enumerate}
 \item the action \(D_3\) leaves the catalogue or changes multiplicities;
 \item any one of the three structural predicates leaves more than one orbit;
 \item the gauge \((6,ES)\) selects two words in an orbit;
 \item a target digit is used to define the predicates;
 \item the accumulated word \(w_{6k}\) is multiplied by \(L_0\) or
       \(L_1\) instead of the block \(b_k\);
 \item \(\det A\neq1\) or fails the inverse of the lift;
 \item two equal visible residues are declared to be equal complete states.
\end{enumerate}
\end{criterion}

The first joint boundary has already been constructed. The next unit
will classify the nine phase components, prove the nonadic return
without resetting, and construct the oriented mirror with its local resolution
\(3\to1\).
 \end{HMTScientificOwnerSubordinated}

% Unidad U009. Conexión nonádica conjunta; redacción autónoma.

\section{Joint nonadic connection, compression, and terminal memory}
\label{p12-83-c7-s1:chap:conexion-nonadica-conjunta}

The nine phases that appear after \(R_{36}\) are not nine independent
proofs. They are the nine local components of a single connection
on integral holonomic states. Their composition around one turn defines
a holonomy. The return of the phase label preserves the continuity of
prefix, carry, boundary, Witt shadow and memory.

\subsection{Transport of the integral holonomic state by nonadic holonomy and memory update}

At the level \(K\), we write
\[
 \widehat x_K=
 (B_K,C_K,p_K,c_K,\Sigma_K,W_K,g(K),q_{{\rm mem},K}),
\]
where:
\begin{enumerate}
 \item \(B_K=(u_K^\pi,u_K^e,u_K^\varphi)\) is the visible block of three
       words of six trits;
 \item \(C_K\) is the cylinder ternario compatible;
 \item \(p_K\) is the complete prefix;
 \item \(c_K\) is the accumulated carry;
 \item \(\Sigma_K\) is the signature joint of boundary;
 \item \(W_K\) is the Witt incidence shadow;
 \item \(g(K)=1+((K-1)\bmod9)\) is the public phase label;
 \item \(q_{{\rm mem},K}\in\mathbb Z_{\geq0}\) is the memory vertical not
       acotada.
\end{enumerate}
The dodecaphasic coordinate is derived:
\[
 \beta_K=[q_{{\rm mem},K}]_{12}.
\]
\(q_{\rm mem}\) is not replaced by \(\beta\), since reduction modulo
twelve would lose the number of circuits.

\subsection{Local relations and holonomy}

For each \(K\), let
\[
 \mathcal R_K\subseteq
 \widehat X_K\times\widehat X_{K+1}
\]
the relation of compatible extension between complete states. The relational composition of a loop is
\begin{equation}
 \Gamma_{9,K}
 :=\mathcal R_{K+8}\circ\cdots\circ\mathcal R_{K+1}\circ\mathcal R_K
 :\widehat X_K\longrightarrow\widehat X_{K+9}.
\label{p12-83-c7-s1:eq:gamma9-composicion-local}
\end{equation}

\begin{definition}[Nonadic holonomy]
\label{p12-83-c7-s1:def:holonomia-nonadica}
The holonomy of one turn is
\[
 \Gamma_9:=\operatorname{Hol}_{C_9}(\nabla_{\rm TPK}),
\]
whose realization at level \(K\) is
\(\Gamma_{9,K}\). The nine relations \(\mathcal R_K,\ldots,
\mathcal R_{K+8}\) are charts of this single connection.
\end{definition}

\begin{theorem}[Phase return without restart]
\label{p12-83-c7-s1:thm:retorno-nonadico-sin-reinicio}
If \((\widehat x_K,\widehat x_{K+9})\in\Gamma_{9,K}\), then
\[
 g(K+9)=g(K),
 \qquad
 q_{{\rm mem},K+9}=q_{{\rm mem},K}+1,
\]
and
\[
 \beta_{K+9}=\beta_K+1\pmod{12}.
\]
In general, \(\widehat x_{K+9}\neq\widehat x_K\).
\end{theorem}

\begin{proof}
The first identity follows from the definition of \(g\). A complete turn increments the unbounded winding counter by definition. Reducing that identity modulo twelve yields the third. The remaining coordinates are transported by the composition \cref{p12-83-c7-s1:eq:gamma9-composicion-local}; no law of the TPK resets them. Therefore, equality of phase does not imply equality of the complete state.
\end{proof}

\subsection{Exhaustive classification of the \(9\) phases of the Topological Phase Kernel}

The visible blocks and the local censuses of the first round are
\[
\begin{array}{c|rrrrr|ccc}
K&N_\pi&N_e&N_\varphi&N_{\rm loc}&N_{\rm ext}
 &u^\pi&u^e&u^\varphi\\\hline
1&378&422&349&3&2&012222&121221&112202\\
2&238&368&388&2&5&010211&102011&121020\\
3&184&284&173&5&4&002111&012222&010210\\
4&207&207&122&4&1&110221&102011&010200\\
5&39&151&151&1&1&222220&021222&102011\\
6&110&110&69&1&3&111201&201202&221021\\
7&81&29&81&3&3&212121&222210&200221\\
8&59&59&59&3&2&200121&212212&110110\\
9&43&43&43&2&1&100100&020112&010122
\end{array}.
\]

\begin{theorem}[Functions of the nine components]
\label{p12-83-c7-s1:thm:funciones-nueve-componentes}
In the preceding turn, the nine components perform the following
functions:
\begin{enumerate}
 \item local opening of three orientations and first resolution;
 \item binary split with a local maximum of five extensions;
 \item maximal local multiplicity of five states;
 \item contraction of four local states to an extension;
 \item first strong saturation, with a single continuation;
 \item unique local present and triple opening of the future;
 \item triple mirror fiber;
 \item synchronization of the three censuses in \(59\);
 \item oriented closure: two local input states and a relation that,
       after the additional horizon, conserves a prolongation.
\end{enumerate}
\end{theorem}

\begin{proof}
Each assertion is read from the censuses \(N_{\rm loc},N_{\rm ext}\), from the
equality or inequality of \(N_\pi,N_e,N_\varphi\), and from the explicit extension
relation in the ninth, tenth, and eleventh phases to be proved
in the following unit. The table enumerates the entire revolution; no subset
of phases has been chosen.
\end{proof}

The fifth component coincides with the first common boundary.
\[
 R_{36}=(222220|021222|102011)^T.
\]
The ninth component contains
\[
 B_9=(100100|020112|010122).
\]
The presence of two local states in the ninth phase and of three extensions
from the oriented state belongs to different layers; they are not contradictory
censuses.

Holonomy acts on the complete integral holonomic state. Before publishing
its modal quotient, one therefore constructs the involution interchanging
the fibers of \(\pi\) and \(e\), the fixed axis of \(\varphi\), the ternary
invariants and the local resolution \(3\to1\) determined by the second
horizon.

\begin{HMTScientificOwnerSubordinated}
% Unidad U010. Espejo orientado y selección 3->1; redacción autónoma.

\chapter{Mirror involution and local selection \(3\to1\) of the surviving extension}
\label{chap:espejo-seleccion-tres-uno}

The ninth phase contains two local input states. Once the oriented
state \(B_9\) has been fixed, the boundary relation produces three extensions
toward the first phase. The three share an axis of self-scale and a ternary
aggregate; they differ in the ordered distribution of that aggregate between the branches of
closure and propagation. A second horizon conserves exactly one.

\section{Symmetric and axial coordinates}

Let
\[
 \mathcal B:=(\mathbb F_3^6)^3.
\]
A visible block is
\[
 B=(u^\pi,u^e,u^\varphi)\in\mathcal B.
\]
We define
\[
 q(B):=u^\pi+u^e+u^\varphi,
 \qquad
 a(B):=u^\pi+u^e-u^\varphi.
\]
All operations are performed componentwise in
\(\mathbb F_3^6\).

\begin{theorem}[Reconstruction of the axis and the aggregate]
\label{thm:espejo-reconstruccion-eje-agregado}
The coordinates \(q,a\) determine
\[
 \boxed{u^\varphi=2(q-a),}
 \qquad
 \boxed{u^\pi+u^e=q-u^\varphi.}
\]
\end{theorem}

\begin{proof}
Subtracting the definitions gives
\(q-a=2u^\varphi\). In \(\mathbb F_3\), \(2^{-1}=2\); therefore,
\(u^\varphi=2(q-a)\). Substitution into the definition of \(q\) recovers the
remaining aggregate.
\end{proof}

The specular involution is
\[
 \mathfrak c(u^\pi,u^e,u^\varphi)
 =(u^e,u^\pi,u^\varphi).
\]

\begin{proposition}[Invariants of the Involution]
\label{prop:espejo-invariantes}
\(\mathfrak c^2=\operatorname{id}\), y
\[
 q(\mathfrak cB)=q(B),
 \qquad
 a(\mathfrak cB)=a(B).
\]
In particular, \(u^\varphi\) and \(u^\pi+u^e\) are invariantes.
\end{proposition}

\begin{proof}
Exchanging the first coordinates twice recovers \(B\). Both
\(q\) and \(a\) depend on them only through their sum, so that the
exchange preserves them. The \cref{thm:espejo-reconstruccion-eje-agregado}
transfers that invariance to the axis and the aggregate.
\end{proof}

The involution asserts neither that \(\pi=e\) nor that one is the quotient of the other.
It asserts that the two branches occupy conjugate positions with respect to a
fixed structure.

\section{Oriented state in phase \(9\) and associated boundary data}

The visible projection of the selected state is
\[
 B_9=(100100\mid020112\mid010122).
\]
The induced relation on visible blocks is denoted
\[
 \mathcal M_9^B\subseteq\mathcal B\times\mathcal B.
\]
The two local input states of the phase table precede this
relation. After selecting the complete state that projects to \(B_9\), its
output fiber has cardinal three.

\begin{theorem}[The three extensions of the ninth phase]
\label{thm:espejo-tres-extensiones}
The relational image of \(B_9\) contains exactly
\begin{align*}
B_{10}^{(0)}&=(101222\mid112221\mid112002),\\
B_{10}^{(1)}&=(102221\mid111222\mid112002),\\
B_{10}^{(2)}&=(111221\mid102222\mid112002).
\end{align*}
The three satisfy
\[
 q_{10}=022112,
 \qquad
 a_{10}=101111,
\]
and, consequently,
\[
 u_{10}^\varphi=112002,
 \qquad
 u_{10}^\pi+u_{10}^e=210110.
\]
\end{theorem}

\begin{proof}
The extension relation on complete states is enumerated in the fiber of
\(B_9\) and produces the three indicated rows. Adding and subtracting their three
components gives the same pair \((q_{10},a_{10})\) in all three cases. The
\cref{thm:espejo-reconstruccion-eje-agregado} produces the common axis and
aggregate. The first two components are distinct, so the three
rows represent different distributions of the same aggregate.
\end{proof}

All their publications in base three are
\[
 (279,352,365),
 \qquad
 (280,351,365),
 \qquad
 (282,350,365).
\]

\section{Exact compatibility of cylinders}

A ternary word of length \(L\) and index \(N\) determines the half-open
cylinder
\[
 I_3(N,L)=\left[\frac{N}{3^L},\frac{N+1}{3^L}\right).
\]
A prefix of \(b\) decimal blocks, with associated integer \(D\), determines
\[
 I_{1000}(D,b)=
 \left[\frac{D}{1000^b},\frac{D+1}{1000^b}\right).
\]

\begin{lemma}[Integer intersection test]
\label{lem:espejo-test-entero-cilindros}
The cylinders \(I_3(N,L)\) and \(I_{1000}(D,b)\) intersect if and only if
\[
 N1000^b<(D+1)3^L,
 \qquad
 (N+1)1000^b>D3^L.
\]
\end{lemma}

\begin{proof}
Two half-open intervals \([a,b)\) and \([c,d)\) intersect if and only if \(a<d\) and \(c<b\). We substitute the four endpoints and multiply by the positive integer \(3^L1000^b\).
\end{proof}

The test uses only integer arithmetic. It does not depend on a floating-point tolerance.

\section{2nd horizon and uniqueness}

The following block is
\[
 B_{11}=(022212\mid110001\mid100021),
\]
and the corresponding decimal boundary is
\[
 (502,662,638).
\]

\begin{theorem}[Local resolution \(3\to1\)]
\label{thm:espejo-seleccion-tres-uno}
The three extensions of \cref{thm:espejo-tres-extensiones} are compatible at depth one. When
\(B_{11}\) and the boundary \((502,662,638)\) are imposed simultaneously,
only \(B_{10}^{(0)}\) prolongs as the same complete state. Therefore,
\[
 \boxed{
 \#\operatorname{Surv}^{\rm cal}_1(B_9)=3,
 \qquad
 \#\operatorname{Surv}^{\rm cal}_2(B_9)=1.}
\]
\end{theorem}

\begin{proof}
We apply \cref{lem:espejo-test-entero-cilindros} to each of the three
channels. At the first horizon, the six inequalities of every candidate are
satisfied. After adjoining \(B_{11}\), the inequalities of the three channels
remain satisfied for \(B_{10}^{(0)}\). In \(B_{10}^{(1)}\) and
\(B_{10}^{(2)}\), compatibility of channel \(\varphi\) is preserved, but
the conditions of branches \(\pi\) and \(e\) fail. Enumeration therefore
leaves a unique state at depth two.
\end{proof}

This proof does not define the global connection from a local example. It is an
explicit witness of how a multivalued relation at one horizon can
become single-valued by preserving memory and requiring a longer prolongation.

\section{Phase return with block change}

The first phase of the initial loop is
\[
 B_1=(012222\mid121221\mid112202),
\]
while the first phase of the next turn is
\[
 B_{10}=(101222\mid112221\mid112002).
\]
Thus,
\[
 g(1)=g(10)=1,
 \qquad
 B_1\neq B_{10}.
\]
The phase returns, but the axis, the aggregate, the distribution, and the memory of
quotient have evolved.

\section{Diagram of the mirror involution \(\pi/e\), the fixed axis \(\varphi\), and the local selection \(3\to1\) by prolongability}

\begin{figure}[htbp]
\centering
\resizebox{0.98\textwidth}{!}{%
\begin{tikzpicture}[
 node distance=10mm and 15mm,
 every node/.style={font=\small,align=center},
 state/.style={draw=black,rounded corners,inner sep=5pt},
 inv/.style={draw=black,double,rounded corners,inner sep=5pt},
 arr/.style={-{Latex[length=2mm]},semithick}
]
 \node[state] (pi) {$u^\pi$\\closure branch};
 \node[state,right=70mm of pi] (e) {$u^e$\\propagation branch};
 \coordinate (mirroraxis) at ($(pi)!0.5!(e)$);
 \node[inv,above=9mm of mirroraxis] (phi) {$u^\varphi$\\eje fijo};
 \draw[arr,<->] (pi) -- node[below] {$\mathfrak c$} (e);
 \node[above=8mm of phi] (aggregate)
   {$u^\pi+u^e$ preserved in the fiber};
 \draw[dashed] (phi.north) -- (aggregate.south);
 \node[state,below=20mm of mirroraxis] (b9) {phase $9$: $B_9$\\two local inputs};
 \node[state,below left=18mm and 25mm of b9] (b100) {$B_{10}^{(0)}$};
 \node[state,below=18mm of b9] (b101) {$B_{10}^{(1)}$};
 \node[state,below right=18mm and 25mm of b9] (b102) {$B_{10}^{(2)}$};
 \node[state,below=18mm of b101] (b11) {$B_{11}$\\unique prolongation};
 \draw[arr] (b9) -- (b100);
 \draw[arr] (b9) -- (b101);
 \draw[arr] (b9) -- (b102);
 \draw[arr] (b100) -- (b11);
 \draw[arr,densely dotted] (b101) -- node[right] {obstructed} (b11);
\draw[arr,densely dotted] (b102) -- node[right] {obstructed} (b11);
\end{tikzpicture}
}
\caption{The mirror preserves the axis \(u^\varphi\) and the aggregate
\(u^\pi+u^e\). The return of the ninth phase opens three visible extensions;
the memory of the second horizon preserves one. The circuit does not reset the
state.}
\label{fig:espejo-tres-uno-arrastre}
\end{figure}

\section{Residue--quotient memory during the mirror}

The recurrence
\[
 x_kA=u_k+3x_{k+1}
\]
of the \cref{eq:bockstein-hensel-local} is applied to each of the three
channels. The part \(u_k\) publishes the visible block; \(x_{k+1}\) conserves
the quotient. Since \(A\in\operatorname{GL}_6(\mathbb Z)\), the pair recovers
the preceding state. The mirror can exchange the visible branches without
eliminating their provenance quotients.

\section{Obstructions to the mirror involution \(\pi/e\) and to the local selection \(3\to1\): invariants, depth-\(2\) horizon, and Hensel memory}

\begin{criterion}[Refutation Criteria for the Mirror]
The construction is refuted if:
\begin{enumerate}
 \item the involution modifies \(u^\varphi\) or \(u^\pi+u^e\);
 \item the two local input states are conflated with the three outputs;
 \item one of the three outputs has a different \((q_{10},a_{10})\);
 \item the cylinder test uses a floating tolerance;
 \item more than one output survives the declared second horizon;
 \item phase return is interpreted as equality \(B_1=B_{10}\);
 \item the visible exchange erases the Hensel quotients.
\end{enumerate}
\end{criterion}

The mirror and the local selection have now been constructed. The following
units will develop the Sturmian vacancies and then the three
complete analytic biographies.
 \end{HMTScientificOwnerSubordinated}

\subsection{Joint transfer of the fiber of 468 emissions}

The factorization of the catalog is
\[
 \mathcal U_6^{\rm cat}\cong
 \mathfrak S_+\times\mathfrak S_\times,
 \qquad
 |\mathfrak S_+|=18,\qquad |\mathfrak S_\times|=26.
\]
For a function \(f\) on the catalogue, let
\[
 (E_+f)(s,e)=\frac1{18}\sum_{s'\in\mathfrak S_+}f(s',e),
\]
and
\[
 (E_\times f)(s,e)=\frac1{26}
 \sum_{e'\in\mathfrak S_\times}f(s,e').
\]
The projectors \(E_+\) and \(E_\times\) commute.

Let
\[
 \tau=(+1,-1,0,+1,-1,0,+1,-1,0),
 \qquad
 P_{+1}=E_+,quad P_{-1}=E_\times,quad P_0=I.
\]
On \(\mathcal U_6^{\rm cat}\times C_9\) we define phase transfer
by
\begin{equation}
 \mathcal K_{\rm ph}\bigl((u,r),(v,[r+1]_9)\bigr)
 :=P_{\tau(r)}(u,v),
 \label{p12-83-c7-s1:eq:u009-kph-explicita}
\end{equation}
and we make it zero for the remaining pairs of phases. Thus, the cyclic permutation
and the operator applied in each of the nine steps are determined.

\begin{proposition}[Joint Renewal After a Turn]
\label{p12-83-c7-s1:prop:transferencia-conjunta-nueve}
\[
 \mathcal K_{\rm ph}^{\,9}=E_+E_\times.
\]
In particular, on a point mass,
\[
 (E_+E_\times)(u,v)=\frac1{18\cdot26}=\frac1{468}.
\]
\end{proposition}

\begin{proof}
The ordered product over a loop is
\[
 P_{\tau(8)}P_{\tau(7)}\cdots P_{\tau(0)}
 =I E_\times E_+ I E_\times E_+ I E_\times E_+
 =E_+^3E_\times^3I^3=E_+E_\times.
\]
The second equality uses commutativity and the final one idempotence. The
second result is the uniform distribution on the product of the two
signatures.
\end{proof}

This identity expresses a joint renewal, not nine independent filtrations.

\subsection{Decomposition of the holonomy \(U_{\Gamma_9}\) into visible compression \(T\) and orthogonal component \(\eta\): \(I-T^*T=\eta^*\eta\)}

Let \(\mathcal H_{\rm vis}\) and \(\mathcal H_{\rm full}\) be Hilbert spaces for the visible and complete realizations. Let
\[
 J:\mathcal H_{\rm vis}\longrightarrow\mathcal H_{\rm full}
\]
an isometric inclusion and
\(U_{\Gamma_9}\) an isometric realization of the holonomy. We define
\[
 T:=J^*U_{\Gamma_9}J,
 \qquad
 \eta:=(I-JJ^*)U_{\Gamma_9}J.
\]

\begin{theorem}[Compression-decomposition expression]
\label{p12-83-c7-s1:thm:conexion-compresion-expresion}
The following hold
\[
 \boxed{U_{\Gamma_9}J=JT+\eta,}
\]
\[
 J^*\eta=0,
 \qquad
 \boxed{I-T^*T=\eta^*\eta.}
\]
\end{theorem}

\begin{proof}
We decompose \(U_{\Gamma_9}J\) by means of the orthogonal projections
\(JJ^*\) and \(I-JJ^*\):
\[
 U_{\Gamma_9}J
 =JJ^*U_{\Gamma_9}J+(I-JJ^*)U_{\Gamma_9}J
 =JT+\eta.
\]
Orthogonality gives \(J^*\eta=0\). Since \(J\) and
\(U_{\Gamma_9}\) are isometries,
\[
 I=J^*U_{\Gamma_9}^*U_{\Gamma_9}J
 =T^*T+\eta^*\eta.
\]
Reorder the proof of the last identity.
\end{proof}

\(T\) is the visible compression; \(\eta\) is the complementary expression.
Neither should be called a strict contraction. For that function a
different operator is reserved.

\subsection{Strict contraction and terminal telescoping}

Let \(V_9\) be an isometry and let \(q\) be a parameter with \(0<q<1\). We define
\[
 M_9=qV_9,
 \qquad
 D_9=(I-M_9^*M_9)^{1/2}.
\]
Then \(\|M_9\|=q<1\) and
\[
 I-M_9^*M_9=D_9^*D_9.
\]

\begin{theorem}[Stein Identity with Terminal Term]
\label{p12-83-c7-s1:thm:stein-telescopia-terminal}
For every \(L\geq1\),
\[
 \boxed{
 I=
 \sum_{r=0}^{L-1}M_9^{*r}D_9^*D_9M_9^r
 +M_9^{*L}M_9^L.}
\]
More generally, if \(S_0\) satisfies
\(S_0-M_9^*S_0M_9=D_9^*D_9\), then
\[
 \boxed{
 S_0=
 \sum_{r=0}^{L-1}M_9^{*r}D_9^*D_9M_9^r
 +M_9^{*L}S_0M_9^L.}
\]
\end{theorem}

\begin{proof}
We multiply the defect identity on the left by \(M_9^{*r}\) and on the
right by \(M_9^r\). Upon summing over \(r=0,\ldots,L-1\), the interior terms
telescope, leaving the initial and terminal terms. The same proof
applies to the general Stein equation.
\end{proof}

The final term does not vanish at finite depth. Its disappearance in a
limit requires a convergence proof; it is not part of the algebraic
identity.

\subsection{Carry Cocycle and Relational Powers}

For composable arrows \(\gamma_1,\gamma_2\), the carry satisfies
\[
 \operatorname{Carr}(\gamma_2\gamma_1)
 =\operatorname{Carr}(\gamma_2)H(\gamma_1)
 +Y(\gamma_2)\operatorname{Carr}(\gamma_1).
\]
The formula preserves the contribution of the first stage after transporting it through the second.

The longer loops are written with displaced indices:
\begin{align*}
 \Gamma_{27,K}
 &=\Gamma_{9,K+18}\circ\Gamma_{9,K+9}\circ\Gamma_{9,K},\\
 \Gamma_{54,K}
 &=\Gamma_{9,K+45}\circ\cdots\circ\Gamma_{9,K},\\
 \Gamma_{108,K}
 &=\Gamma_{9,K+99}\circ\cdots\circ\Gamma_{9,K}.
\end{align*}
Only after declaring a global operator on the graded union may
this notation be abbreviated as \(\Gamma_9^3,\Gamma_9^6,\Gamma_9^{12}\).

\subsection{Complete nonadic state: phase, carry, boundary, route, and memory after the prolongation}

The one-turn connection is recorded by the quintuple
\[
 \mathfrak N_9=
 \bigl(\Gamma_9,T,\eta,\operatorname{Carr},S_{\rm term}\bigr),
\]
where \(S_{\rm term}=M_9^{*L}S_0M_9^L\) at a finite depth
\(L\). An application functor that uses the connection must declare its
compatibility with all five fields. Commuting only with the visible phase does not
constitute naturality with respect to the complete nonadic state.

\subsection{Obstructions of the holonomy \(\Gamma_9\): return \(9\to1\), composition indices and preservation of phase, carry, boundary, route and memory}

\begin{criterion}[Refutation Criteria of the Connection]
The construction is refuted if:
\begin{enumerate}
 \item the nine phases are presented as independent filters;
 \item the return \(9\to1\) resets \(q_{\rm mem}\), the prefix, or the carry;
 \item \(T\) is called a strict contraction;
 \item some identity of
       \cref{p12-83-c7-s1:thm:conexion-compresion-expresion} fails;
 \item the terminal term is removed at finite depth;
 \item \(\Gamma_{27,K}\) is composed without shifting the indices;
 \item an application functor ignores any of the five fields of
       \(\mathfrak N_9\).
\end{enumerate}
\end{criterion}

The complete developments that follow first reconstruct the finite emitter,
its factorization and its microfibers; they then unfold the specular involution,
the three extensions of \(B_9\), uniqueness at two horizons and the
Bockstein--Hensel inverse that preserves memory during phase return.

\begin{HMTScientificOwnerSubordinated}
\chapter[Finite arithmetic image and factorization \(468=18\cdot26\)]{Finite arithmetic image and factorization \(468=18\cdot26\)}
\label{chap:alfabeto-42}
\index{admissible decimal code}
\index{check digit}

The image in base \(1000\) of the finite map that assigns the decimal blocks is a
proper subset of the one thousand decimal blocks. The evolution law defined in the
\cref{chap:tpk-definicion} and the arithmetic of \(\mathbb Z/9\mathbb Z\)
determine an image of 42 triples, organized into six classes of seven. The
same separation of the additive and multiplicative observables subsequently factors the 468
outputs into 18 additive signatures and 26 multiplicative signatures.

\section{Primitives and recurrence of the finite sender}

\begin{definition}[Initial data of \(T_{\min}\)]
\label{pf84:E02-15}
Let \(D=\{N,E,S,O\}\) be given, with cardinal vectors
\[
 v(N)=(-1,0),\quad v(E)=(0,1),\quad
 v(S)=(1,0),\quad v(O)=(0,-1)
 \qquad\text{in }(\mathbb Z/9\mathbb Z)^2.
\]
The temporal signature and the multiplicative index are fixed by
\[
 \operatorname{sgn}_9(t)=
 \begin{cases}
 +1,&t\bmod9\in\{1,4,7\},\\
 -1,&t\bmod9\in\{2,5,8\},\\
 0,&t\bmod9\in\{0,3,6\},
 \end{cases}
\]
and
\[
 \operatorname{ind}_2(1,2,4,8,7,5)=(0,1,2,3,4,5),
\]
with value zero outside the units. An entry of \(T_{\min}\) is a pair
of seeds
\[
 I_9:=\{0,1,\ldots,8\},\qquad
 s_+=(i_+,j_+,d_+),\qquad s_\times=(i_\times,j_\times,d_\times)
 \in I_9^2\times D.
\]
These tables, the readers \(L_\Sigma,L_\Pi\), the six blocks of nine
instants, and the inversion of both vectors when \(27\mid t\) are all the
starting data. Neither \(\pi\), \(e\), \(\varphi\), nor any other
target value is introduced.
\end{definition}

\begin{algorithm}[54-step recurrence]
\label{pf84:E02-16}
At the outset, the two cursors occupy their seed positions and carry
the vectors \(v(d_+)\), \(v(d_\times)\). For
\(t=1,\ldots,54\), in this order:
\begin{enumerate}
  \item if \(\operatorname{sgn}_9(t)=+1\), \(L_\Sigma(i_+(t),j_+(t))\) is accumulated
  and the additive cursor is advanced;
  \item if \(\operatorname{sgn}_9(t)=-1\), then
  \(\operatorname{ind}_2(L_\Pi(i_\times(t),j_\times(t)))\) is accumulated and the
  multiplicative cursor is shifted;
  \item if \(\operatorname{sgn}_9(t)=0\), both cursors remain stationary;
  \item after that update, both vectors are reversed if
  \(27\mid t\).
\end{enumerate}
The accumulators are reset at the beginning of each block
\(9(r-1)+1\leq t\leq9r\). If \(S_r,E_r\) are their final values, then there is emitted
\[
 U_r=100\bar S_r+10(\bar S_r+\bar E_r)_{10}
       +(3S_r+5E_r+1)_{10},
 \qquad
 \bar S_r=S_r\bmod10,\quad \bar E_r=E_r\bmod10,
\]
y \(w_r=(-U_r)\bmod3\). Thus it is obtained.
\[
 T_{\min}(s_+,s_\times)=(U_1,\ldots,U_6)
 \in\mathcal U_6^{\rm cat},
 \qquad
 \Psi(U_6)=(-U_6)\bmod3=:w_6\in\mathbb F_3^6.
\]
\end{algorithm}

\begin{criterion}[Target-value-independent dual control] \label{pf84:E02-17} The recurrence admits two independent formulations: channel factorization and joint simulation. Exhaustive traversal of the \(104\,976\) inputs produces the same sextuple by both routes for every seed, without consulting a target constant. A discrepancy for a single seed would refute the equivalence.
\end{criterion}

\section{Bivariate dynamics and block observables}

We work on the torus \(\mathbb Z_9^2\). Under the preceding convention, a
complete seed is
\[
\omega=(p_+,p_\times,h_+,h_\times)
\in(\mathbb Z_9^2)^2\times\{N,E,S,O\}^2,
\]
so that
\[
|\Omega|=9^4\cdot4^2=104\,976.
\]
In each window of nine steps, the times \(1,4,7\) evaluate the additive observable,
the times \(2,5,8\) evaluate the multiplicative observable, and the times \(3,6,9\)
increment the neutral counter. If \(S_m,E_m,T_m\) are the accumulators
of the window \(m\), the decimal observable is
\[
d_1\equiv S_m,\qquad
d_2\equiv S_m+E_m,\qquad
d_3\equiv3S_m+5E_m+7T_m\pmod{10},
\]
and the decimal triple is \(U_m=100d_1+10d_2+d_3\).

\begin{lemma}[Neutral counter]
For every window, \(T_m=3\).
\end{lemma}
\begin{proof}
Each block of nine times contains exactly the three neutral times
\(3,6,9\). Therefore the counter is incremented three times.
\end{proof}

\section{Image of the additive accumulator: \(\operatorname{Im}S_m=\{6,9,12,15,18,21,24\}\)}

Let \(s=i+j\pmod9\). The visible additive observable is
\[
f(s)=\operatorname{dr}_9(s+2)=(2,3,4,5,6,7,8,9,1).
\]
A cardinal translation modifies \(s\) in \(+1\) or in \(-1\). In a window,
three consecutive terms are read; reversing the orientation only reverses their
order. The nine cyclic sums are
\[
9,12,15,18,21,24,18,12,6.
\]

\begin{lemma}[Image of the additive accumulator]
\[
S_m\in\{6,9,12,15,18,21,24\}.
\]
In particular, \(d_1=S_m\bmod10\) traverses precisely
\(\{1,2,4,5,6,8,9\}\).
\end{lemma}
\nopagebreak[4]
\begin{proof}
The preceding list enumerates all initial positions of the trajectory and the
two orientations. Their decimal residues are seven and distinct.\qedhere
\end{proof}

\section{Image of the multiplicative accumulator: \(\operatorname{Im}E_m=\{0,1,3,5,7,9\}\)}

The unit group \((\mathbb Z/9\mathbb Z)^\times\) is cyclic of order
six. We fix the multiplicative coordinate function
\[
\ell_\times(1,2,3,4,5,6,7,8,9)=(0,1,0,2,5,0,4,3,0).
\]
If the row or column traversed has a factor divisible by three, each
evaluation belongs to \(\{3,6,9\}\) and contributes zero. If the fixed factor is a
unit, the sums of three consecutive terms of
\(b\mapsto\ell_\times(\operatorname{dr}_9(ab))\) are:
\[
\begin{array}{c|c}
a&\text{possible sum}\\
\midrule
1&\{1,3,7,9\}\\
2&\{3,5,9\}\\
4&\{1,5,7\}\\
5&\{1,5,7\}\\
7&\{3,5,9\}\\
8&\{1,3,7,9\}.
\end{array}
\]

\begin{lemma}[Image of the multiplicative accumulator]
\[
E_m\in\{0,1,3,5,7,9\}.
\]
\end{lemma}
\begin{proof}
The noninvertible case contributes \(0\). The finite table of the six units
contains all possibilities in the invertible case, and their union is
\(\{1,3,5,7,9\}\).
\end{proof}

\section{Decimal congruence of the ternary alphabet: uniqueness of the 3rd digit}

\begin{theorem}[Control congruence]
Every terna \(\overline{d_1d_2d_3}\) emitted by the system satisfies
\[
\boxed{d_3\equiv5d_2-2d_1+1\pmod{10}.}
\]
\end{theorem}
\begin{proof}
From \(d_2\equiv S_m+E_m\) one obtains
\(E_m\equiv d_2-d_1\pmod{10}\). Since \(T_m=3\),
\[
d_3\equiv3d_1+5(d_2-d_1)+21
\equiv5d_2-2d_1+1\pmod{10}.\qedhere
\]
\end{proof}

\begin{theorem}[Exact image of the decimal application]
The system produces exactly \(7\cdot6=42\) decimal triples. Grouped by
\(E=d_2-d_1\pmod{10}\), they are:
\[
\begin{array}{c|rrrrrrr}
E&\multicolumn{7}{c}{\text{decimal triples}}\\
\midrule
0&114&227&443&556&669&885&998\\
1&129&232&458&561&674&890&903\\
3&149&252&478&581&694&810&923\\
5&169&272&498&501&614&830&943\\
7&189&292&418&521&634&850&963\\
9&109&212&438&541&654&870&983.
\end{array}
\]
\end{theorem}
\begin{proof}
There are seven possibilities for \(d_1\) and six for \(E\). For each pair,
\(d_2\equiv d_1+E\) and the control congruence fix a unique \(d_3\),
so there exist at most 42 triples. The table applies these two formulas
to the 42 pairs and contains no repetitions. The enumeration of the
104\,976 seeds finds all the triples in the table and none outside
it; the bound is attained.
\end{proof}

\section{Factorization of the \(2\) observables}

Let \(\mathcal S=\mathbb Z_9^2\times\{N,E,S,O\}\), with
\(|\mathcal S|=324\). To each additive seed \(P\in\mathcal S\) is
associated the six-window signature
\[
s(P)=(S_1,\ldots,S_6)\pmod{10},
\]
and to each multiplicative seed \(T\in\mathcal S\),
\[
e(T)=(E_1,\ldots,E_6)\pmod{10}.
\]
The two signatures are combined coordinate by coordinate by
\[
G(s,e)_k=
100s_k+10(s_k+e_k)_{10}+(3s_k+5e_k+1)_{10}.
\]
The map \(G\) is injective: from a triple one recovers
\(s_k=d_1\) and \(e_k=d_2-d_1\pmod{10}\), while the third digit verifies
membership in the image.

\begin{lemma}[Exact classification of the additive channel]
\label{pf84:E02-18}
\[
 |\operatorname{im}s|=18,
\]
and each additive signature has exactly \(18\) preimages.
\end{lemma}
\begin{proof}
The signature depends only on
\((i+j\bmod9,\varepsilon)\), where \(\varepsilon=+1\) for \(E,S\) and
\(-1\) for \(N,O\). The \(9\cdot2\) pairs produce distinct signatures; each
has nine positions on its diagonal and two directions, that is,
\(18\) preimages. A different number of signatures or a class of size
different from \(18\) would refute the classification.
\end{proof}

\begin{lemma}[Exact classification of the multiplicative channel]
\label{pf84:E02-19}
\[
 |\operatorname{im}e|=26.
\]
There are \(24\) classes of size \(8\), one of size \(24\), and one of size
\(108\).
\end{lemma}
\begin{proof}
Enumerating the \(9^2\cdot4=324\) seeds using the explicit table of \(\ell_\times\) gives \(26\) sextets
\((E_1,\ldots,E_6)\) and the histogram
\[
 24\cdot8+24+108=324.
\]
An additional, absent, or differently sized class would refute the census.
\end{proof}

\begin{theorem}[Exact injective factorization of the catalog]
\label{pf84:E02-20}
\[
|\operatorname{im}s|=18,\qquad
|\operatorname{im}e|=26,\qquad
|\operatorname{im}G|=468=18\cdot26.
\]
Each additive signature has 18 preimages. The multiplicative signatures possess
the histogram
\[
24\text{ size classes }8,\qquad
1\text{ of size }24,\qquad
1\text{ of size }108.
\]
\end{theorem}
\begin{proof}
The two preceding lemmas give the cardinalities of the signatures. In each emitted
component, the hundreds digit is \(\bar S_r\), while the tens digit is
\(\bar S_r+\bar E_r\bmod10\). Therefore,
\[
 \bar S_r=\left\lfloor\frac{U_r}{100}\right\rfloor,\qquad
 \bar E_r=
 \left(\left\lfloor\frac{U_r}{10}\right\rfloor-\bar S_r\right)\bmod10.
\]
The units digit verifies \(3S_r+5E_r+1\bmod10\); it introduces no other
identification. Thus, \(G\) is injective, and its image has
\(18\cdot26=468\) elements. The \(24\) multiplicative classes of size
\(8\), combined with the \(18\) additive signatures, produce \(432\) sextets
of multiplicity \(144\); the other two classes produce \(18\) sextets of
multiplicity \(432\) and \(18\) of multiplicity \(1944\). In particular,
\[
 432\cdot144+18\cdot432+18\cdot1944=104\,976.
\]
Two pairs of signatures with the same \(U_6\), or a histogram that
failed to recover the complete domain, would refute the theorem.
\end{proof}

The complete composition factors, therefore, as
\[
\mathcal S^2\xrightarrow{s\times e}
\operatorname{im}s\times\operatorname{im}e
\xrightarrow{G}\mathcal U_6
\xrightarrow{\rho_3}\mathcal W_6\subset\mathbb F_3^6.
\]
The final projection has 243 values. The first descent preserves the
sum--product separation; the second preserves 468 decimal states; the
third removes the fiber coordinate and retains a ternary word.

\begin{corollary}[Census of the observable projection]
\label{cor:censo-proyeccion-observable}
\label{pf84:E02-21}
The complete enumeration determines the chain of cardinalities
\[
 104\,976\longrightarrow468\longrightarrow243\subset729.
\]
The \(468\) intermediate elements are sextuples of integers between
zero and \(999\); only their final image belongs to \(\mathbb F_3^6\). The
fiber multiplicities of the first map satisfy
\[
 432\cdot144+18\cdot432+18\cdot1944=104\,976.
\]
\end{corollary}

\begin{proof}
The factorization theorem gives \(18\cdot26=468\), and the exhaustive enumeration
of the \(104\,976\) seeds produces exactly the three
indicated multiplicities. Componentwise reduction traverses the
\(243\) words recorded in the ternary catalogue. The appended catalogues of
signatures and routes reproduce both enumerations.
\end{proof}

\begin{proposition}[Quintuple microfiber of the word \(010211\)]
\label{pf84:E02-22}
For the restricted reduction
\[
 \Psi:\mathcal U_6^{\rm cat}\longrightarrow\mathbb F_3^6,
 \qquad \Psi(U_6)=(-U_6)\bmod3,
\]
the fiber of \(010211\) consists of exactly the following five sextets:
\begingroup
\scriptsize
\[
\begin{array}{@{}cclc@{}}
\toprule
U_6&\sigma_+&\sigma_\times&\text{multiplicity}\\
\midrule
(501,614,498,169,272,272)&(5,6,4,1,2,2)&(5,5,5,5,5,5)&432\\
(810,923,870,169,272,272)&(8,9,8,1,2,2)&(3,3,9,5,5,5)&144\\
(810,983,810,169,272,272)&(8,9,8,1,2,2)&(3,9,3,5,5,5)&144\\
(870,923,810,169,272,272)&(8,9,8,1,2,2)&(9,3,3,5,5,5)&144\\
(870,923,810,418,674,521)&(8,9,8,4,6,5)&(9,3,3,7,1,7)&144\\
\bottomrule
\end{array}
\]
\endgroup
They are five preimages of the catalog \(U_6\), not five complete TPK
histories. They preserve distinct signatures and multiplicities although they share the same
visible word.
\end{proposition}

\begin{proof}
The exhaustive scan of the catalogue of \(468\) sextuples selects exactly
the five rows of the table. The signature-recovery function reproduces
the two sextuples of every row, and the sum of multiplicities is
\(432+4\cdot144=1008\). Enumeration additionally verifies that every cell is the
Cartesian product of its channel seeds. A cardinality other than five
or identification of two rows despite their enriched data would refute the
statement.
\end{proof}
 \end{HMTScientificOwnerSubordinated}

\begin{HMTScientificOwnerSubordinated}
\chapter{Inverse systems, Weyl--Heisenberg representation and predictive quotient}
\label{chap:numero-heisenberg-d108}
\index{HMT number}
\index{Heisenberg!finite group}

The construction distinguishes three levels: the finite seed catalogue, the
state space that preserves carry and orientation, and the inverse
limit of compatible prolongations. This stratification makes it possible to
formalize contraction of cylindrical intervals to a real point,
persistence of trajectory data in a single evaluation fiber, and
the finite Weyl--Heisenberg representation associated with the alternating form
of the nonadic torus supporting the additive and multiplicative observables.

The symbols \(V_t\), \(X_t\), and \(\rho_{t+1,t}\) used here are the
same sets of local states, admissible prefixes, and truncations
introduced in \cref{chap:numero-medicion}, with the purely
notational change \(t=m\) and \(\rho_{n,m}=p_{n,m}\). This section makes their
coordinates explicit and introduces no second inverse system.

\section{Observable projection and complete state space of the TPK}

The observable projection of the phase transport system (TPK) is the
finite encoding application
\[
104\,976\longrightarrow468\longrightarrow243.
\]
which classifies seeds, signatures, and words \(w_6\). The complete TPK
state space further preserves
\[
\begin{gathered}
\text{bidirectional orientation},\quad
\text{Hensel state},\quad
\text{cylinder and carry},\\
\text{joint boundary signature},\quad
\text{visible nonadic phase }g_{\rm vis}(m).
\end{gathered}
\]
Two candidates with the same local signature may belong to distinct classes
when the following boundary is conserved. We do not introduce a second tuple
of state. For
\[
x_m=(q_m,o_m,h_m,c_m,\sigma_m,C_m,b_m,\lambda_m)\in\XTPK^{(m)}
\]
we use the coordinate chart
\[
\chi_m(x_m)=
\bigl(
B_m,h_m,(C_m,b_m),\sigma_m,
o_m,\mu(q_m),g_{\rm vis}(m),\lambda_m
\bigr),
\]
with
\[
B_m=(u_m^\pi,u_m^e,u_m^\varphi)\in(\mathbb F_3^6)^3.
\]
Here \(h_m\) is the Hensel extension; \((C_m,b_m)\), the cylinder and its
boundary; \(o_m,\mu(q_m)\), the orientation and the bidirectional observable; and
\[
g_{\rm vis}(m)=1+((m-1)\bmod9)\in\{1,\ldots,9\},
\]
the visible nonadic phase label. The corresponding class coordinate
is \(g_0(\theta_m)\in\ZZ/9\ZZ\), with \(\theta_m=[m-1]_{108}\). The joint signature
is the coordinate already typed by the phase kernel,
\[
\sigma_m=\operatorname{Sig}_{q_m}(h_m,c_m,\lambda_m)
=(q_\partial,a_\partial,c_\partial,
\operatorname{colw},\rho_W,r_\partial).
\]
The joint signature does not reduce to three independent channel signatures, because the
selection associated with the involution \(\pi/e\) and the axis \(\varphi\) depends on
correlated data. Two nodes with the same \(B_m\) represent distinct
states when they differ in extension, cylinder, carry, signature, orientation,
or transport record.

\subsection{Exact compatibility of the cylinders under truncation and nonadic prolongation}

Let \(N\) be the integer of a ternary prefix of length \(L\), and let \(D\)
be the integer obtained by concatenating its \(b\) decimal triads. We define
\[
A=N1000^b-D3^L,
\qquad
B=(N+1)1000^b-D3^L.
\]
The cylinders intersect if and only if
\[
\boxed{A<3^L,\qquad B>0.}
\]
Adding a ternary block \(u\in\{0,\ldots,728\}\) gives
\[
N'=729N+u.
\]
If a new decimal triad does not appear,
\[
A'=729A+u1000^b.
\]
If \(\delta\) occurs, so that \(D'=1000D+\delta\), then
\[
A'=729000A+u1000^{b+1}-\delta\,729\,3^L.
\]
If \(\Delta b\in\{0,1\}\) records the appearance of that new triad, in
both cases
\[
B'=A'+1000^{b+\Delta b}.
\]
The transition, once the admissible pair \((u,\delta)\) has been fixed, does not consult a
target value. The joint filtration by intersection of cylinders and nonadic
return acts at a distinct level: it selects the boundary pairs that
admit compatible prolongation.

\subsection{Cyclic registration of 108 stages and its reduced projection}

We call the sequence of \emph{108-stage cyclic record} the sequence of
\(108\) steps grouped into twelve sectors of nine steps, with the index
convention fixed in \cref{chap:tpk-definicion}.

Let \(\XTPK^{[108]}\subseteq\XTPK\) be the domain of complete states whose
register \(\lambda\) contains the declared \(108\) steps. Let us denote by
\(\mathcal C_{108}^{\rm red}\) the set of registers that result by
conserving only the reduced fields, and by \(\mathcal U_{12}^{\rm sgn}\) the
set of admissible signed dodecaphases. The instance distinguishes two declared maps:
\[
F:\XTPK^{[108]}\longrightarrow\mathcal C_{108}^{\rm red},
\qquad
R:\XTPK^{[108]}\longrightarrow\mathcal U_{12}^{\rm sgn}.
\]
The map \(F\) retains only the fields of the reduced register, whereas \(R\) evaluates the signed dodecaphase. The existence of a map \(\Xi\) such that \(R=\Xi\circ F\) is equivalent to the factorization of \(R\) through a projection that may eliminate the observable coordinate, orientation, carry, or boundary. This factorization constitutes an additional proposition and is not an implicit part of the reconstruction on the complete state space.

The reduced record is, by definition, a projection of the complete state,
not the initial state. Subsequent results use \(R\) on the
complete state and do not presuppose a factorization through \(F\).

\section{Inverse systems of survivors}

The TPK transition need not be a function. To type this fact correctly,
let \(V_t\) be the finite set of possible complete states at
depth \(t\), and let
\[
\mathcal R_t\subseteq V_t\times V_{t+1}
\]
the admissible transition relation, with all fields of the integral holonomic
state preserved. We then define \(X_t\), not as the set of
isolated terminal states, but as the set of admissible prefixes, that is,
of complete admissible histories:
\[
X_t=
\Bigl\{
(v_0,\ldots,v_t)\in\prod_{j=0}^{t}V_j:
(v_j,v_{j+1})\in\mathcal R_j\ (0\le j<t),
\ \text{and the closures are satisfied up to }t
\Bigr\}.
\]
On these objects, truncation is indeed a canonical application:
\[
\begin{aligned}
\rho_{t+1,t}:X_{t+1}&\longrightarrow X_t,\\
(v_0,\ldots,v_t,v_{t+1})&\longmapsto(v_0,\ldots,v_t).
\end{aligned}
\]
It is well defined precisely because every prefix of an admissible history remains an admissible history for the closures already imposed. If a filter were to use future information and could retroactively invalidate a prefix, \(X_t\) would first have to be replaced by the subset of prolongable prefixes; under that convention, the same truncation is again well defined.

A global HMT history is a compatible family of prefixes.
\[
x=(x_t)_{t\ge0},
\qquad
x_t\in X_t,
\qquad
\rho_{t+1,t}(x_{t+1})=x_t.
\]
It is equivalent to an infinite sequence \((v_0,v_1,\ldots)\) whose consecutive
pairs belong to the relations \(\mathcal R_t\) and whose prefixes
satisfy all closures. The history space is the inverse limit
\[
X_\infty=\varprojlim X_t.
\]
This formulation avoids imposing a single-valued local transition \(B_t\mapsto B_{t+1}\). The transition may be a finite relation, several branches may share the same visible terminal state, and uniqueness may appear only after future compatibility is imposed. By preserving the complete prefix, the information that makes it possible to distinguish them is not lost prematurely.

\begin{theorem}[Unique coinductive section]
Suppose that, for every \(t\), the survivor set \(S_t\subseteq
X_t\) is nonempty, the restrictions
\(\rho_{t+1,t}:S_{t+1}\to S_t\) are compatible, and every \(S_t\) contains a
unique prefix. Then \(\varprojlim S_t\) contains a unique global
history.
\end{theorem}

\begin{proof}
Let \(x_t\) be the unique prefix of \(S_t\). Compatibility forces
\(\rho_{t+1,t}(x_{t+1})=x_t\), so that \((x_t)_t\) belongs to the inverse
limit. Every other history would have a member of \(S_t\) at each level
and, by uniqueness, would coincide prefix by prefix with \((x_t)_t\).
\end{proof}

The theorem formalizes the passage from level-by-level compatibility to a global
history. It does not turn a finite window into a coinductive conclusion: its
application to a concrete family requires verifying at every depth non-
emptiness, compatibility of the restrictions, and uniqueness of the
surviving prefixes.

\section{Nonadic Survival Connection}
\label{sec:conexion-nonadica}

Cylinder contraction does not proceed block by block as a deterministic substitution. At certain boundaries, there are several local allocations of the same aggregate datum; prolongability determines which belongs to the published history. The resulting structure is a connection on the nine phase classes: upon completing one turn, the phase class returns, but the state of the fiber has changed. This section constructs the complete finite window in which that return is observed.

The domain of the calculation is the attached calibrated transport register. Its initial states,
its integer propagation matrix, and its boundary relations are published
premises, not outputs of an autonomous enumeration of TPK at arbitrary
depth. All statements in this section are exact relative to
those premises. Extension to every depth is governed by the preceding
coinductive theorem and requires its hypotheses to be verified at every level.

\subsection{Prefix of \(30\) trits and 1st boundary}

We write \(u_k^\chi\in\FF_3^6\), with
\(\chi\in\{\pi,e,\varphi\}\), for the visible block of channel \(\chi\) at
depth \(k\). The first five blocks form three prefixes of thirty
trits:
\begin{align*}
w_{30}^{\pi}&=010211|012222|010211|002111|110221,\\
w_{30}^{e}&=201101|121221|102011|012222|102011,\\
w_{30}^{\varphi}&=121200|112202|121020|010210|010200.
\end{align*}
The sixth block is the first boundary following those prefixes:
\[
R_{36}=
\begin{pmatrix}
u_5^\pi\\ u_5^e\\ u_5^\varphi
\end{pmatrix}
=
\begin{pmatrix}
222220\\021222\\102011
\end{pmatrix}.
\]
The notation \(R_{36}\) indicates that thirty-six trits per
channel have been displayed. It does not denote a terminal state: the record reconstructed below
contains twenty blocks, and the recurrence admits further continuations
when the corresponding Hensel states are supplied.

\subsection{Self-scale axis and residual fiber}

\ifdefined\HMTAxialTheoremDefined
The axial decomposition of the signature and the identification of the axis
\(\varphi\) against the residual fiber \(\pi/e\) were already proved in the
\cref{thm:eje-espejo-nonadico}. Here its consequence for the
nonadic window is conserved without duplicating the theorem or the proof.
\else
For
\[
B_k=(u_k^\pi,u_k^e,u_k^\varphi)\in(\FF_3^6)^3
\]
we define, coordinate by coordinate,
\[
q_k=u_k^\pi+u_k^e+u_k^\varphi,
\qquad
a_k=u_k^\pi+u_k^e-u_k^\varphi.
\]

\begin{theorem}[Axial decomposition of the boundary signature]
\label{thm:eje-espejo-nonadico}
The pair \((q_k,a_k)\) uniquely determines the axial channel and the aggregate of
the other two channels:
\[
u_k^\varphi=2(q_k-a_k),
\qquad
u_k^\pi+u_k^e=q_k-u_k^\varphi.
\]
Consequently, a fiber of the application
\[
q_{\rm ax}(u^\pi,u^e,u^\varphi)
=(u^\pi+u^e,u^\varphi)
\]
can preserve ambiguity only in the internal allocation between \(u^\pi\) and
\(u^e\).
\end{theorem}

\begin{proof}
Subtracting the two defining equations gives
\(q_k-a_k=2u_k^\varphi\). The inverse of \(2\) in \(\FF_3\) is again
\(2\), whence
\(u_k^\varphi=2(q_k-a_k)\). Subtracting this vector from \(q_k\) produces
\(u_k^\pi+u_k^e\). The operation does not recover each summand separately; for
this reason the unique residual fiber is the set of distributions of the fixed aggregate.
\end{proof}

In the model's geometric terminology, \(\varphi\) constitutes the axis
fixed by the signature and the pair \(\pi/e\) the residual fiber. The expression does not
postulate an analytic symmetry of the real constants: it names the
exact linear decomposition just proved in \(\FF_3^6\).
\fi

\subsection{Orbita complete of the \(9\) fases consecutivas of the Topological Phase Kernel}

We fix the absolute phase by
\[
g(k)=
\begin{cases}
k\bmod9,&k\not\equiv0\pmod9,\\
9,&k\equiv0\pmod9.
\end{cases}
\]
Since \(w_{30}\) occupies \(k=0,\ldots,4\), the first subsequent complete circuit
after the prefix traverses \(k=5,\ldots,13\), with phases
\(5,6,7,8,9,1,2,3,4\). Enumeration of the \(729=3^6\) candidates of
aggregate signature at each level produces the following window:
\begin{center}
\small
\begin{tabular}{@{}ccrrlll@{}}
\toprule
\(k\)&\(g(k)\)&locals&selected&\(u_k^\pi\)&\(u_k^e\)&\(u_k^\varphi\)\\
\midrule
5&5&1&1&222220&021222&102011\\
6&6&1&1&111201&201202&221021\\
7&7&3&1&212121&222210&200221\\
8&8&3&1&200121&212212&110110\\
9&9&2&1&100100&020112&010122\\
10&1&1&1&101222&112221&112002\\
11&2&1&1&022212&110001&100021\\
12&3&1&1&012012&202222&000120\\
13&4&1&1&111210&112102&211122\\
\bottomrule
\end{tabular}
\end{center}
The “local” column counts the admissible input states at the indicated
level; it does not count the outputs of the boundary relation. This distinction
is essential in \(k=9\): the two local input states precede a
three-output fiber of the selected state.

\subsection{Return 9→1 of the nonadic phase with state change and memory advance}

Let
\[
 \pi_B:V_k\longrightarrow\mathcal B:=(\FF_3^6)^3
\]
the projection onto the visible triple. Once the fiber data of the record are fixed,
the transport relation on complete states induces a visible
relation \(\MonNine^B\). The selected state entering the ninth phase is
\[
B_9=(100100\mid020112\mid010122).
\]
Its relational image consists exactly of
\[
\begin{aligned}
B_{10}^{(0)}&=(101222\mid112221\mid112002),\\
B_{10}^{(1)}&=(102221\mid111222\mid112002),\\
B_{10}^{(2)}&=(111221\mid102222\mid112002).
\end{aligned}
\]
The three outputs have
\[
q_{10}=022112,
\qquad a_{10}=101111,
\]
and, by \cref{thm:eje-espejo-nonadico}, share
\[
u_{10}^\varphi=112002,
\qquad u_{10}^\pi+u_{10}^e=210110.
\]
They differ exclusively in the three published distributions of that aggregate
between \(\pi\) and \(e\).

\begin{theorem}[Two-level extendibility selection]
\label{thm:seleccion-novena-fase}
In the calibrated window, the three preceding outputs admit a prolongation
of depth one. Upon also imposing the next block
\[
B_{11}=(022212\mid110001\mid100021)
\]
and the decimal boundary \((502,662,638)\), only \(B_{10}^{(0)}\) admits a
prolongation of the same complete state. Therefore
\[
\#\operatorname{Surv}^{\rm cal}_{1}(B_9)=3,
\qquad
\#\operatorname{Surv}^{\rm cal}_{2}(B_9)=1.
\]
\end{theorem}

\begin{proof}
The decimal triples associated with the three outputs are, respectively,
\[
(279,352,365),\quad(280,351,365),\quad(282,350,365).
\]
For a ternary word of length \(L\), with associated integer \(N\), and a
prefix of \(b\) triples, with associated integer \(D\), compatibility of the
two cylinders is equivalent to the two integer inequalities
\[
N1000^b<(D+1)3^L,
\qquad
(N+1)1000^b>D3^L.
\]
The six inequalities for each output are satisfied at the current level. Upon
adjoining \(B_{11}\) and \((502,662,638)\), they are again satisfied for the three
channels of \(B_{10}^{(0)}\). For \(B_{10}^{(1)}\) and \(B_{10}^{(2)}\), the
inequality of channel \(\varphi\) is preserved, but those of \(\pi\) and
\(e\) fail. The evaluation is performed with integers in the corresponding finite certificate; no
numerical tolerance intervenes. Consequently, the candidate set
is reduced from three to one when passing from depth one to depth two.
\end{proof}

The other two words can precede the same ternary tail only if the
deep Hensel extension is changed. They are therefore not prolongations of
the same complete state of the register: the visible residue partially coincides, but
the genealogy of the fiber is distinct.

\subsection{Phase return with state change}

The phase satisfies \(g(k+9)=g(k)\). The visible state, by contrast, is not
reset. In particular,
\[
B_1=(012222\mid121221\mid112202),
\qquad
B_{10}=(101222\mid112221\mid112002),
\]
so that \(g(1)=g(10)=1\) but \(B_1\ne B_{10}\). The verification of the
nine pairs \(k\mapsto k+9\), for \(k=1,\ldots,9\), also gives
\[
u_{k+9}^\varphi\ne u_k^\varphi,
\qquad
u_{k+9}^\pi+u_{k+9}^e\ne u_k^\pi+u_k^e.
\]
Therefore, the window satisfies exactly
\[
\boxed{g(k+9)=g(k),\qquad B_{k+9}\ne B_k
\quad(1\le k\le9).}
\]
This is a phase return with nontrivial visible holonomy. The designation
\emph{nonadic monodromy} refers to this return action on the
enriched fiber; it does not introduce a tenth phase. To promote the finite remark
to a global monodromy, the action at every depth must be constructed
and the coinductive compatibility required at the beginning of the section must be proved.

\subsection{The branch of \(20\) blocks}

The preceding return does not close the emission. The twenty blocks that the
certificate reconstructs from the three published integer states are
\begin{align*}
\pi:\;&010211|012222|010211|002111|110221|222220|111201|\\
&212121|200121|100100|101222|022212|012012|111210|\\
&121011|200220|120210|000101|022010|020111,\\[2pt]
e:\;&201101|121221|102011|012222|102011|021222|201202|\\
&222210|212212|020112|112221|110001|202222|112102|\\
&102010|022020|222002|012010|001112|022212,\\[2pt]
\varphi:\;&121200|112202|121020|010210|010200|102011|221021|\\
&200221|110110|010122|112002|100021|000120|211122|\\
&202200|202110|221000|100102|111122|020010.
\end{align*}
Thus, \(w_{30}\) is the prefix and \(R_{36}\) the first boundary of an
explicitly continued branch; neither object replaces the integral holonomic
state that transports that continuation.

\subsection{Exact decoding of \(12\) decimal triples by means of the nonadic reader}

Let \(N_{\chi,k}\) be the integer represented in base three by the first \(k\)
blocks of channel \(\chi\). With \(k=13\), the word has \(78\) trits and
defines
\[
D_{\chi,12}
=\left\lfloor\frac{1000^{12}N_{\chi,13}}{3^{78}}\right\rfloor.
\]
The escritura of \(D_{\chi,12}\) in twelve blocks of three cifras is
\[
\begin{array}{c|rrrrrrrrrrrr}
\pi&141&592&653&589&793&238&462&643&383&279&502&884\\
e&718&281&828&459&045&235&360&287&471&352&662&497\\
\varphi&618&033&988&749&894&848&204&586&834&365&638&117.
\end{array}
\]
The identity is entire: it is equivalent to
\[
D_{\chi,12}3^{78}
\le 1000^{12}N_{\chi,13}
<(D_{\chi,12}+1)3^{78}.
\]
Thus, the two cylinders intersect in every channel, and the reading
operation produces exactly the thirty-six published digits. This
result is a consequence of the calibrated register. The initial states in
the attached chart were historically constructed with those constants as
reference data; therefore, this identity does not establish the independent
provenance of the initial states from the primitives of APP\@.

\subsection{Schedule for cylinder reopening and periodicity of the compatible prolongation}

The conversion of a six-trit block to decimal triples has average slope
\[
\lambda=6\log_{1000}3\in(0,1),
\qquad k(t)=\lfloor\lambda(t+1)\rfloor.
\]
An index \(t\) is a reopening without a new triple when
\(k(t+1)=k(t)\). Let \(\delta=1-\lambda\).

\begin{theorem}[Beatty Reading Calendar Ternary-Decimal]
\label{thm:calendario-beatty}
The set of reopenings is
\[
\mathcal T=
\left\{\left\lceil\frac{m}{\delta}\right\rceil-2:m\ge1\right\}.
\]
Its first elements are
\[
20,42,64,86,108,130,151,173,195,217,\ldots,
\]
and two reaperturas consecutivas distan \(21\) or \(22\) blocks.
\end{theorem}

\begin{proof}
The quantity \(\lambda=\log_{1000}729\) is irrational: an equality
\(\lambda=p/q\) would imply \(729^q=1000^p\), impossible by prime
factorization. Also \(\delta\) is irrational, so that
\[
k(t)=(t+1)-\lceil\delta(t+1)\rceil.
\]
The equality \(k(t+1)=k(t)\) occurs exactly when the ceiling function of the
second term increases by one. The \(m\)-th occurrence is obtained
by solving
\(\delta(t+1)<m<\delta(t+2)\), which gives
\(t=\lceil m/\delta\rceil-2\). The differences of a Beatty sequence
belong to
\[
\left\{\left\lfloor\delta^{-1}\right\rfloor,
\left\lceil\delta^{-1}\right\rceil\right\}=\{21,22\}.\qedhere
\]
\end{proof}

\subsection{Bockstein--Hensel Recurrence of Residue and Quotient}

The memory carried by the branch proceeds from the exact sequence.
\[
0\longrightarrow\ZZ^6\xrightarrow{\,\times3\,}\ZZ^6
\longrightarrow\FF_3^6\longrightarrow0
\]
and of the unimodular matrix
\[
A=
\begin{pmatrix}
2&2&2&1&2&1\\
2&1&2&2&1&1\\
1&0&0&1&0&2\\
0&0&1&1&1&0\\
2&2&2&0&2&1\\
2&1&2&1&0&0
\end{pmatrix},
\qquad\det A=1.
\]
For a row state \(x_k^\chi\in\ZZ^6\), it is defined
\[
y_k^\chi=x_k^\chi A,
\qquad
u_k^\chi=y_k^\chi\bmod3\in\{0,1,2\}^6,
\qquad
x_{k+1}^\chi=\frac{y_k^\chi-u_k^\chi}{3}.
\]

\begin{proposition}[Exact conservation of quotient memory]
\label{prop:memoria-bockstein-hensel}
With \(x_k^\chi\) and \(A\) fixed, the pair
\((u_k^\chi,x_{k+1}^\chi)\) is unique and satisfies
\[
x_k^\chi A=u_k^\chi+3x_{k+1}^\chi.
\]
Conversely, since \(A\in\operatorname{GL}_6(\ZZ)\), the pair uniquely determines
\(x_k^\chi\).
\end{proposition}

\begin{proof}
Euclidean division of each coordinate of \(x_k^\chi A\) by three produces
a unique residue in \(\{0,1,2\}\) and a unique integer quotient. This proves
the first assertion and the identity. Since \(\det A=1\), the inverse
\(A^{-1}\) has integer entries; therefore
\[
x_k^\chi=(u_k^\chi+3x_{k+1}^\chi)A^{-1}
\]
recovers the previous state unambiguously.
\end{proof}

Residue \(u_k^\chi\) is the visible part, and quotient
\(x_{k+1}^\chi\) the transported 3-adic memory. The certificate reproduces
sixty coordinated steps ---twenty per channel---, verifies \(AA^{-1}=I\) with
integer arithmetic, and checks all blocks against the register. The designation
Bockstein--Hensel refers here to this residue--quotient separation and its
3-adic iteration; it is not identified without additional data with a
cohomological Bockstein operator.

\section{From the prefix to a real point}

Let us associate with each compatible prefix \(x_t\in X_t\) a closed interval
\(I_t(x_t)\subset\mathbb R\). The geometric interpretation is that of a
cylinder: all futures sharing the same prefix project into
\(I_t\).

\begin{theorem}[Archimedean convergence of nested intervals]
Let \((x_t)_{t\ge0}\) be a compatible history. If
\[
I_{t+1}(x_{t+1})\subseteq I_t(x_t)
\]
and
\[
\operatorname{diam} I_t(x_t)\longrightarrow0,
\]
then there exists a unique \(r\in\mathbb R\) such that
\[
\bigcap_{t\ge0}I_t(x_t)=\{r\}.
\]
\end{theorem}

\begin{proof}
This is the principle of nested intervals. The left endpoints form a
bounded increasing sequence and the right endpoints a bounded decreasing one. Their limits
exist; the difference between them is bounded by the diameter, which tends to
zero. Both limits coincide at a single point contained in all the
intervals.
\end{proof}

Let \(X_\infty^{\rm Arch}\subseteq X_\infty\) be the domain of histories for
which the cylinders are nested and their diameters tend to zero. If those
two properties are part of global admissibility, then
\(X_\infty^{\rm Arch}=X_\infty\); if not, this domain restriction is
necessary. We define the evaluation
\[
\operatorname{ev}_{\mathbb R}:X_\infty^{\rm Arch}\longrightarrow\mathbb R,
\qquad
\operatorname{ev}_{\mathbb R}(x)=
\text{the unique point of }\bigcap_{t\ge0}I_t(x_t).
\]
Uniqueness of the value does not entail uniqueness of the history. Two
genealogies \(x\ne y\) may satisfy
\[
\operatorname{ev}_{\mathbb R}(x)
=\operatorname{ev}_{\mathbb R}(y).
\]
We define reading equivalence
\[
x\sim_{\rm ev}y
\quad\Longleftrightarrow\quad
\operatorname{ev}_{\mathbb R}(x)
=\operatorname{ev}_{\mathbb R}(y).
\]

\begin{theorem}[Archimedean evaluation quotient]
The map
\[
\begin{aligned}
\overline{\operatorname{ev}}_{\mathbb R}:
X_\infty^{\rm Arch}/\!\sim_{\rm ev}
&\longrightarrow
\operatorname{Im}(\operatorname{ev}_{\mathbb R})\subseteq\mathbb R,\\
[x]&\longmapsto\operatorname{ev}_{\mathbb R}(x)
\end{aligned}
\]
is a bijection. In particular, the quotient is identified with all of
\(\mathbb R\) if and only if the evaluation
\(\operatorname{ev}_{\mathbb R}\) is surjective.
\end{theorem}

\begin{proof}
The definition of \(\sim_{\rm ev}\) makes the value independent of the
representative, so the map is well defined. It is injective because
two classes with the same value are, by definition, a single class; it is
surjective onto the image by the very definition of image. The last
assertion is equivalent to
\(\operatorname{Im}(\operatorname{ev}_{\mathbb R})=\mathbb R\).
\end{proof}

The covering theorem for the Archimedean reader constructs precisely those
prefixes: for every compact interval \(K\subset\mathbb R\),
\(\operatorname{ev}_{\mathbb R}(\Omega_{\infty,K})=K\), and the intervals
\([-m,m]\) form a directed family whose union is \(\mathbb R\).
Consequently, every real number possesses at least one compatible HMT history, and the
proved assertion is
\[
X_\infty^{\rm Arch}/\!\sim_{\rm ev}
\cong\operatorname{Im}(\operatorname{ev}_{\mathbb R})
=\mathbb R.
\]
The Archimedean value is, therefore, the quotient that identifies the fibers of
the evaluation; the enriched HMT number also conserves the class of
trajectory, the carry, and the orientation. The identification does not follow from the
mere abstract formation of the quotient: it follows from the covering theorem already
proved in the discrete construction of the continuum.

Algebraic descent is another problem and must be solved operation by
operation. Suppose that two candidate operations
\(\oplus\) and \(\odot\) have been constructed on \(X_\infty^{\rm Arch}\) and are closed on that
domain. Addition descends to the quotient if and only if
\[
x\sim_{\rm ev}x',\ y\sim_{\rm ev}y'
\Longrightarrow
x\oplus y\sim_{\rm ev}x'\oplus y',
\]
while the product descends, independently, if and only if
\[
x\sim_{\rm ev}x',\ y\sim_{\rm ev}y'
\Longrightarrow
x\odot y\sim_{\rm ev}x'\odot y'.
\]
The first congruence does not imply the second. Moreover, for the induced operations
to be precisely the ordinary sum and product of the Archimedean quotient, it is necessary to prove separately
\[
\operatorname{ev}_{\mathbb R}(x\oplus y)
=\operatorname{ev}_{\mathbb R}(x)+\operatorname{ev}_{\mathbb R}(y),
\]
\[
\operatorname{ev}_{\mathbb R}(x\odot y)
=\operatorname{ev}_{\mathbb R}(x)\operatorname{ev}_{\mathbb R}(y),
\]
and that the image of the evaluation be closed under the corresponding operation.
Only after these certificates does the quotient inherit the structure of a
subring ---or of a subfield, if additive opposites and multiplicative inverses are also constructed---
of \(\mathbb R\). The set-theoretic bijection of the preceding theorem is formal;
the content specific to HMT lies in constructing the evaluation, characterizing its
fibers, proving coverage, and establishing each descent law.

\section{Weyl--Heisenberg Representation Associated with the Nonadic Torus}

The oriented area form of the nonadic torus is
\[
\sigma((a,b),(c,d))=ad-bc\pmod9.
\]
This alternating form is obtained, with the chosen orientation, by
integrating the mixed curvature of APP over parallelograms of the torus. Let
\(\omega=\exp(2\pi i/9)\) and
\(\mathcal H_9=\mathbb C^9\), with basis \(\{|n\rangle:n\in\mathbb Z_9\}\).
We define
\[
X|n\rangle=|n+1\rangle,
\qquad
Z|n\rangle=\omega^n|n\rangle.
\]
Then
\[
ZX=\omega XZ,
\]
and, more generally,
\[
(X^aZ^b)(X^cZ^d)
=\omega^{bc}X^{a+c}Z^{b+d}.
\]
Therefore
\[
 c\bigl((a,b),(c,d)\bigr)=bc\pmod9
\]
is the polarized \(2\)-cocycle that appears in this section of the central
extension. We explicitly define
\[
 \operatorname{Heis}_9
 =\{\omega^kX^aZ^b:a,b,k\in\ZZ/9\ZZ\}.
\]
It is a group of order \(9^3=729\). It is a central extension associated
with the translation group of the APP support \(X_9\), not a subgroup of the
graph or the APP path category\@.
The commutator of the two transports is
\[
(X^aZ^b)(X^cZ^d)(X^aZ^b)^{-1}(X^cZ^d)^{-1}
=\omega^{bc-ad}I.
\]
Consequently, its exponent is \(-\sigma((a,b),(c,d))\). The polarized cocycle
\(c\) and the form \(\sigma\) are not the same object: the second is,
under this convention, the opposite of the antisymmetrization
\(c(u,v)-c(v,u)\). The commutator phase is thus the inverse of the
Wilson phase for the fixed orientation.

To construct observables, we take
\[
Q_{\rm H}=\operatorname{diag}(0,1,\ldots,8),
\qquad
F_{jk}=\frac13\omega^{jk},
\qquad
P_{\rm H}=F^\dagger Q_{\rm H}F.
\]
The operators \(Q_{\rm H},P_{\rm H}\) are self-adjoint. Their eigenbases are mutually
unbiased because \(|F_{jk}|^2=1/9\). For every unit vector
\(\psi\in\mathcal H_9\), the Robertson inequality gives
\[
\operatorname{Var}_\psi(Q_{\rm H})\,
\operatorname{Var}_\psi(P_{\rm H})
\ge
\frac14\left|\langle\psi,[Q_{\rm H},P_{\rm H}]\psi\rangle\right|^2,
\]
and the entropy bound for mutually unbiased bases produces
\[
H_{Q_{\rm H}}(\psi)+H_{P_{\rm H}}(\psi)\ge\log9.
\]
Here \(H\) is the Shannon entropy, with natural logarithm, of the
corresponding spectral distributions. The constructed relations are
summarized, without attributing any additional physical causality to them, by
\[
\begin{gathered}
\text{alternating form }(X_9,\sigma)
\longrightarrow
\operatorname{Heis}_9,\\
\operatorname{Heis}_9\longrightarrow(X,Z)
\longrightarrow(Q_{\rm H},P_{\rm H})
\longrightarrow\text{finite uncertainty}.
\end{gathered}
\]
This is a finite algebraic model of operators; by itself it does not specify
physical dynamics or units of measurement.

\section{Cyclic dynamics on \texorpdfstring{\(\mathbb Z_9\times\mathbb Z_{12}\)}{Z9 x Z12} and minimal predictive quotient}

The finite dynamical system \(\mathfrak C_{9,12}\) makes it possible to determine the minimal
memory required for the observed evolution to be deterministic. Let
\[
X=\{(b,m):b\in\mathbb Z/9\mathbb Z, 0\le m<12\}
\]
with successor
\[
S(b,m)=
\begin{cases}
(b,m+1),&m<11,\\
(b+1\bmod9,0),&m=11.
\end{cases}
\]
The marginal labels are
\[
\beta(b,m)=4b\pmod{12},
\qquad
d(b,m)=1+(m+\beta(b,m)\pmod{12}).
\]
For a remark application \(q\), the inclusive future word of horizon \(h\) is
\[
Q_h^q(x)=\bigl(q(x),q(Sx),\ldots,q(S^hx)\bigr).
\]

\begin{theorem}[Minimum predictive quotient]
Let \(X\) finite, \(T:X\to X\) and \(q:X\to V\). Definamos
\[
\operatorname{Eq}(f):=\{(x,y)\in X^2:f(x)=f(y)\},
\qquad
E_h=\operatorname{Eq}(Q_h^q),
\qquad
E_\infty=\bigcap_{k\ge0}\operatorname{Eq}(q\circ T^k).
\]
Then
\(E_{h+1}=E_h\cap(T\times T)^{-1}(E_h)\), the sequence stabilizes in finite
time, and \(X/E_\infty\) is the quotient with the smallest number of states that
refines \(q\) and admits an induced deterministic dynamics.
\end{theorem}

\begin{proof}
\(E_h\) requires equality of observations at the times \(0,\ldots,h\), whereas
\((T\times T)^{-1}(E_h)\) requires equality at \(1,\ldots,h+1\). Their intersection is
\(E_{h+1}\). Because \(X\) is finite, the number of classes cannot grow
indefinitely. At the first stable level, \(T\) respects the classes and
descends. Every \(T\)-compatible equivalence contained in
\(\operatorname{Eq}(q)\)
identifies only states with the same complete future and is therefore
contained in \(E_\infty\); hence the minimality property.
\end{proof}

The evaluation exhaustiva of \(\mathfrak C_{9,12}\) da:
\begin{center}
\small
\begin{tabular}{@{}lrrrr@{}}
\toprule
observable & instantaneous values & index \(h\) & observations & stable quotient\\
\midrule
\(\beta\) & 3 & 11 & 12 & \(C_{36}\), fibra 3\\
\(d\) & 12 & 8 & 9 & \(C_{36}\), fibra 3\\
\((\beta,d)\) & 36 & 0 & 1 & \(C_{36}\), fibra 3\\
\bottomrule
\end{tabular}
\end{center}
In the two marginal observables,
\[
\operatorname{Eq}(Q_{11}^{\beta})
=\operatorname{Eq}(Q_8^d)
=\operatorname{Eq}(\pi_{36}),
\qquad
\pi_{36}(b,m)=(b\bmod3,m).
\]
The factor \(36\) is not introduced as a target. It appears as the minimal predictive quotient of each marginal. The phase \(m\) is a state coordinate required to predict those futures; this mathematical assertion does not presuppose a physical causal interpretation. The transformation \(b\mapsto b+3\) is a symmetry invisible to those observables at every time.

As an abstract dynamical system, the quotient is a cycle of thirty-six states. Writing \(C_{36}\) does not, however, choose a distinguished origin: it must be fixed to obtain a coordinate \(\mathbb Z/36\mathbb Z\). Nor does this automatically turn those states into the thirty-six pairs of a nonad. Indeed, \(\operatorname{Aut}J(9,2)\cong S_9\). If the two points of a pair lie in distinct cycles of lengths \(r,s\) of a permutation, the orbit of the pair has length \(\operatorname{mcm}(r,s)\), with \(r+s\leq9\), and hence length at most \(20\). If they lie in the same cycle of length \(r\leq9\), the orbit has length \(r\), except in the antipodal case, when it has length \(r/2\). Thus no automorphism induces an orbit of length \(36\) on the pairs. Therefore, the passage
\[
C_{36}\longleftrightarrow \binom{[9]}2
\]
remains a calibrated identification, not a natural equivariant identification.
Even a Hamiltonian ordering of the pairs would prove local preservation
of incidence step by step, not a global automorphism or the canonicity of the
dictionary.

\section{HMT number as a compatible history in \(X_\infty\): genealogical identity,
reading structure, and Archimedean evaluation}

The preceding results determine the following instance of the governing definition.
In the notation of this section, an HMT number is a compatible
history
\[
x\in X_\infty=\varprojlim X_t.
\]
A reader \(L\) is additional structure: the pair \((x,L)\) constitutes a read presentation or a measurement, but does not alter the identity type of \(x\). When \(x\) belongs to \(X_\infty^{\rm Arch}\), its Archimedean value is the point intersection of the cylinders; the set of values so constructed is \(\operatorname{Im}(\operatorname{ev}_{\mathbb R})\subseteq\mathbb R\), and its identification with the entire line requires the additional coverage theorem. Its minimal memory relative to an observable is the predictive quotient of futures; and the alternating form of the sum--product support admits the finite Weyl--Heisenberg representation constructed above. Value, trajectory class, and uncertainty are defined by three operations on specified domains of the same transport architecture.
 \end{HMTScientificOwnerSubordinated}
  \part{Discrete structure of the continuum and double holographic projection}

\chapter{Inverse system of compatible trajectories}
\label{chap:83-03}

\section{Genealogical construction of the continuum}

\subsection{APP, TRIT and TPK: governing definitions}

The construction begins with three irreducible and causally ordered objects. The object of \emph{All Possible Paths} (hereafter, \APP) constructs the common nonadic support, its graph of all finite paths, and the two sum--product evaluation sheets. The tritic unit (hereafter, \TRIT) is the minimal local unit of topological, geometric, informational, and, once ordered, temporal relation: its visible signature has three states, and its lift preserves the carry that a binary or residual reading would lose. The \emph{Topological Phase Kernel} (hereafter, \TPK) is the transport category that composes these units into routes and preserves, within the same state, sheet, orientation, boundary, genealogical ledger, memory, phase, and survival.

\begin{center}
\begin{tabularx}{.94\textwidth}{@{}p{.10\textwidth}Y Y@{}}
\toprule
Object & Construction & Invariantes transportados\\
\midrule
APP & positive chart, finite torus, Cayley graph, paths, and leaves
\(\Sigma,\Pi\) & marker, cell, neighborhood, sum, product, quotient and closure\\
TRIT & ternary signature visible and local topological-temporal lift &
residue, carry, sheet, orientation, threshold and reversal\\
TPK & category of paths, fibers, and compatible extensions & cylinder,  
boundary, genealogical record, memory, phase, holonomy and survival\\
\bottomrule
\end{tabularx}
\end{center}

The exact construction of APP is established on a marked segment and a
finite chart; the real line appears afterward as a limiting publication. The
reference interval, its marks, and the oriented path that traverses them are:
\begin{equation}
I_{10}=[0,10],\qquad
\mathcal M_{10}=I_{10}\cap\mathbb Z=\{0,1,\ldots,10\},
\qquad
\mathsf P_{10}:0\longrightarrow1\longrightarrow\cdots
\longrightarrow9\longrightarrow10,
\label{p12-83-c3-s1:eq:app-marked-segment}
\end{equation}
The endpoints fix the edge and the nine interior marks
\(I_{10}^{\circ}\cap\mathbb Z\) form the positive alphabet. The exact
mathematical object of departure is the chart
\begin{equation}
\mathcal D_9=\{1,\ldots,9\},
\qquad \mathcal D_8=\{1,\ldots,8\},
\qquad
\lambda_9:\mathcal D_9\xrightarrow{\sim}\mathbb Z/9\mathbb Z,
\qquad \lambda_9(9)=[0]_9.
\label{p12-83-c3-s1:eq:app-positive-chart}
\end{equation}
The mark \(9\) is simultaneously the last positive representative and the
publication of the null class. For every \(n\geq1\), the positive digital root and
quotient are separated by means of
\begin{equation}
\operatorname{dr}_9(n)=1+((n-1)\bmod9),
\qquad q_9(n)=\left\lfloor\frac{n-1}{9}\right\rfloor,
\qquad n=9q_9(n)+\operatorname{dr}_9(n).
\label{p12-83-c3-s1:eq:positive-digital-root-carry}
\end{equation}
To the path \(1\to2\to\cdots\to9\) we add the unique closing edge \(9\to1\). The successor and its memory are separated:
\begin{equation}
u_+(a)=\begin{cases}a+1,&a<9,\\1,&a=9,\end{cases}
\qquad
\kappa_+(a)=\begin{cases}0,&a<9,\\1,&a=9,\end{cases}
\qquad a+1=u_+(a)+9\kappa_+(a).
\label{p12-83-c3-s1:eq:app-successor-carry}
\end{equation}
For that reason \(9\to1\) closes the phase, increases the memory, and prolongs the
genealogical state.

The square of the chart produces the finite torus and its movement graph:
\begin{equation}
X_9=(\mathbb Z/9\mathbb Z)^2,
\qquad
\Gamma_{\rm APP}=\operatorname{Cay}
\bigl(X_9,\{\pm e_1,\pm e_2\}\bigr),
\qquad
\operatorname{Path}(\Gamma_{\rm APP}).
\label{p12-83-c3-s1:eq:app-cayley-torus}
\end{equation}
Identifying opposite sides preserves the cardinality of the chart and adds the continuation law. The core of \(8^2=64\) interstices is completed by the minimal gnomon that publishes the zero class in both coordinates:
\begin{equation}
\mathcal D_9^2=\mathcal D_8^2\sqcup\mathcal G_9,
\qquad
81=64+(9+8)=64+17.
\label{p12-83-c3-s1:eq:app-gnomon-public}
\end{equation}
The same law is prolonged to charts of greater depth; the scale changes,
but the closure, the shared edge, and the memory of the turn preserve their type.

The associated volumetric completion is expressed through the strata identity.
\begin{equation}
10^3=(9+1)^3=729+243+27+1.
\label{p12-83-c3-s1:eq:completion-unit}
\end{equation}
\begin{figure}[tbp]
\centering
\begin{minipage}[t]{.53\textwidth}
\centering
\resizebox{.96\linewidth}{!}{%
\begin{tikzpicture}[
  font=\scriptsize,
  cell/.style={minimum width=4.35mm,minimum height=4.35mm,
    inner sep=0pt,draw=black!28}]
  \node[font=\bfseries] at (1.72,.88)
    {Multiplicative leaf of the APP chart: $\Pi(i,j)=\operatorname{dr}_9(ij)$};
  \foreach \i in {1,...,9}{
    \node[font=\bfseries] at (-.35,{-(\i-1)*.44}) {\i};
    \foreach \j in {1,...,9}{
      \pgfmathtruncatemacro{\v}{mod(\i*\j-1,9)+1}
      \pgfmathtruncatemacro{\border}{(\i==9)||(\j==9)}
      \ifnum\border=1
        \node[cell,fill=black!16,font=\bfseries]
          at ({(\j-1)*.44},{-(\i-1)*.44}) {\v};
      \else
        \node[cell,fill=white]
          at ({(\j-1)*.44},{-(\i-1)*.44}) {\v};
      \fi
    }
  }
  \foreach \j in {1,...,9}{
    \node[font=\bfseries] at ({(\j-1)*.44},.38) {\j};
  }
  \draw[very thick] (3.74,.22)--(3.74,-3.74)--(-.22,-3.74);
  \draw[densely dashed] (-.22,.22) rectangle (3.30,-3.30);
  \node[align=center,anchor=north,font=\scriptsize] at (1.76,-3.98)
    {$8^2=64$ cells interior; gnomon of closure
     $9+8=17$\\corona positive of nueves: the mark $9$ publishes $[0]_9$};
\end{tikzpicture}}
\end{minipage}\hfill
\begin{minipage}[t]{.44\textwidth}
\centering
\resizebox{.98\linewidth}{!}{%
\begin{tikzpicture}[font=\scriptsize]
  \begin{scope}[x={(.78cm,0cm)},y={(.36cm,.26cm)},z={(0cm,.78cm)}]
    \coordinate (O) at (0,0,0);
    \coordinate (X) at (3.1,0,0);
    \coordinate (Y) at (0,3.1,0);
    \coordinate (XY) at (3.1,3.1,0);
    \coordinate (Z) at (0,0,3.1);
    \coordinate (XZ) at (3.1,0,3.1);
    \coordinate (YZ) at (0,3.1,3.1);
    \coordinate (XYZ) at (3.1,3.1,3.1);
    \path[fill=black!5] (O)--(X)--(XZ)--(Z)--cycle;
    \path[fill=black!12] (X)--(XY)--(XYZ)--(XZ)--cycle;
    \path[fill=black!20] (Z)--(XZ)--(XYZ)--(YZ)--cycle;
    \draw[densely dashed] (O)--(Y)--(XY) (Y)--(YZ);
    \draw[thick] (O)--(X)--(XZ)--(Z)--cycle
      (X)--(XY)--(XYZ)--(XZ) (Z)--(YZ)--(XYZ);
    \node[font=\bfseries\large,fill=white,inner sep=1.5pt]
      at (1.55,.52,1.55) {$9^3=729$};
  \end{scope}
  \node[font=\bfseries] at (3.62,3.68) {Cube of Extreme and Mean Holographic Ratio};
  \draw[fill=black!8] (0.2,-.78) rectangle (3.15,-.05);
  \node[align=center,font=\tiny] at (1.68,-.415)
    {inner region\\$9^3=729$};
  \draw[fill=black!16] (3.15,-.78) rectangle (5.15,-.05);
  \node[align=center,font=\tiny] at (4.15,-.415)
    {$3$ caras nuevas\\$3\cdot9^2=243$};
  \draw[fill=black!24] (5.15,-.78) rectangle (6.35,-.05);
  \node[align=center,font=\tiny] at (5.75,-.415)
    {$3$ edges\\$3\cdot9=27$};
  \draw[fill=black!34] (6.35,-.78) rectangle (7.02,-.05);
  \node[align=center,font=\tiny] at (6.685,-.415)
    {$1$ vertex};
  \node[font=\bfseries] at (3.61,-1.30)
    {$10^3=729+243+27+1=729+271$};
  \node[align=center] at (3.61,-1.88)
    {growth corona $9^3\subset10^3$: $271$ states\\
     internal boundary of the cube of side $9$: $386$ points};
\end{tikzpicture}}
\end{minipage}
\caption{Two finite publications of the support. Left: the product table with positive digital root and its closure gnomon. Right: the EMRH cubic completion, which distinguishes interior, faces, edges, and vertex. The growth corona $271$ and the internal boundary $386$ are two typed counts: the former measures the enlargement $9^3\subset10^3$ and the latter the boundary of the cube of side $9$.}
\label{fig:app-emrh-mini}
\end{figure}
 The interior volume, faces, edges and point closure are thus
typed as distinct layers of the same completed unit. On the 81
cells of the positive chart act the evaluations
\[
\Sigma,\Pi:\mathcal D_9^2\longrightarrow\mathcal D_9,
\qquad
\Sigma(i,j)=\operatorname{dr}_9(i+j),
\qquad
\Pi(i,j)=\operatorname{dr}_9(ij),
\]
Their residues and quotients raised are defined by
\[
\begin{aligned}
R^+(i,j)&:=\Sigma(i,j),&
Q^+(i,j)&:=\frac{i+j-R^+(i,j)}9,\\
R^\times(i,j)&:=\Pi(i,j),&
Q^\times(i,j)&:=\frac{ij-R^\times(i,j)}9,
\end{aligned}
\]
and, therefore,
\[
\begin{aligned}
\widehat\Sigma(i,j)&:=\bigl(R^+(i,j),Q^+(i,j)\bigr),\\
\widehat\Pi(i,j)&:=\bigl(R^\times(i,j),Q^\times(i,j)\bigr),
\end{aligned}
\qquad
\widehat\Sigma,\widehat\Pi:
\mathcal D_9^2\longrightarrow\mathcal D_9\times\mathbb Z_{\geq0}.
\]
The bijection \(\lambda_9^{\times2}:\mathcal D_9^2\to X_9\) transports
these operations and their lifts to the torus. Therefore, APP is the tuple
\begin{equation}
\APP=\bigl(\mathcal D_9,\lambda_9,X_9,\Gamma_{\rm APP},
\operatorname{Path}(\Gamma_{\rm APP}),
\Sigma,\Pi,\widehat\Sigma,\widehat\Pi\bigr),
\label{p12-83-c3-s1:eq:app-full-definition-public}
\end{equation}
The table constitutes a two-dimensional representation of this complete tuple.
The additive sheet is Latin; the multiplicative sheet conserves
the ternary radical of the ring \(\mathbb Z/9\mathbb Z\)
\begin{equation}
R_9:=\mathbb Z/9\mathbb Z,
\qquad
\mathfrak m:=([3]_9)=\{[0]_9,[3]_9,[6]_9\}\triangleleft R_9,
\qquad \mathfrak m^2=(0).
\label{p12-83-c3-s1:eq:app-radical-369}
\end{equation}
The internal rotation is the automorphism
\[
\rho:\mathfrak m^3\longrightarrow\mathfrak m^3,
\qquad \rho(x_1,x_2,x_3)=(x_2,x_3,x_1),
\]
whose distinguished orbit is
\[
([3]_9,[6]_9,[0]_9)\xrightarrow{\rho}
([6]_9,[0]_9,[3]_9)\xrightarrow{\rho}
([0]_9,[3]_9,[6]_9)\xrightarrow{\rho}
([3]_9,[6]_9,[0]_9).
\]
The direct reduction sends \(3,6,9\) to \(0,0,0\). The internal coordinate
of the ideal, by contrast, conserves the extracted factor:
\begin{equation}
\eta_{\mathfrak m}:(\mathfrak m,+)
\xrightarrow{\;\sim\;}(\mathbb Z/3\mathbb Z,+),
\qquad
\eta_{\mathfrak m}([3a]_9):=[a]_3,
\qquad
([3]_9,[6]_9,[0]_9)
\xmapsto{\;\eta_{\mathfrak m}^{\times3}\;}
([1]_3,[2]_3,[0]_3),
\label{p12-83-c3-s1:eq:app-369-120}
\end{equation}
so that the \(3,6,9\) sector remains intact and its internal coordinate is
read as \(1,2,0\). This map belongs to the radical of APP; the
temporal TRIT projection constitutes the subsequent calendrical reader. The
three calendrical classes
remain separated as
\[
\mathcal C_+=\{1,4,7\},\qquad
\mathcal C_-=\{2,5,8\},\qquad
\mathcal C_0=\{3,6,9\}.
\]
The incompatibility between the two leaves becomes a measurable invariant. On
\(\mathscr A_3=\mathcal D_9^3\), with uniform measure
\(\nu_3(x)=9^{-3}\), let
\(\mathcal H_3=L^2(\mathscr A_3,\nu_3)\). We define
\[
\Sigma_3,\Pi_3:\mathscr A_3\longrightarrow\mathcal D_9,
\qquad
\Sigma_3(x)=\operatorname{dr}_9(x_1+x_2+x_3),
\qquad
\Pi_3(x)=\operatorname{dr}_9(x_1x_2x_3),
\]
and the orthogonal projectors
\[
P_{\Sigma,3}:=\mathbb E_{\nu_3}[\,\cdot\mid\sigma(\Sigma_3)],
\qquad
P_{\Pi,3}:=\mathbb E_{\nu_3}[\,\cdot\mid\sigma(\Pi_3)]
\in\mathcal B(\mathcal H_3).
\]
The joint enumeration of the \(9^3\) words then determines
\[
\bigl\|[P_{\Sigma,3},P_{\Pi,3}]\bigr\|=\frac{\sqrt3}{4},
\qquad
B_c=7I+2R=9P_++5P_-,
\qquad
459-405=54.
\]
The first number reappears as the electronic prefactor; the second
organizes the central block; the third records the global sum--product
defect.

\subsection{TRIT: local topological unit and temporal relation unit}

HMT first fixes the visible signature
\begin{equation}
\mathrm{TRIT}_{\rm vis}=\{-1,0,+1\},
\qquad \tau\in\mathrm{TRIT}_{\rm vis}.
\label{p12-83-c3-s1:eq:trit-visible-signature}
\end{equation}
In contrast to the binary unit, the TRIT preserves three local regimes. In
the uniform distribution, its information and entropy are distinguished as
\begin{equation}
I_{\rm TRIT}=\log3\ \text{nats}=\log_2 3\ \text{bits},
\qquad S_{\rm TRIT}=k_BI_{\rm TRIT}.
\label{p12-83-c3-s1:eq:trit-information-unit}
\end{equation}
The bit remains the binary unit; the TRIT is the minimal unit capable of simultaneously distinguishing return, threshold, and opening. The equality \eqref{p12-83-c3-s1:eq:trit-information-unit} quantifies the information of the triad; the cost \(k_B\Theta\log3\) belongs to the subsequent thermal reset reader. Dimensional Mechanics also preserves the arithmetic antecedent of that signature. For \((i,j)\in\mathcal D_9^2\), the associated lifted arithmetic cell is
\begin{equation}
c_{\rm MD}(i,j;\tau)=
\bigl(i,j;R^+,Q^+,R^\times,Q^\times;\tau\bigr).
\label{p12-83-c3-s1:eq:trit-complete-local-unit}
\end{equation}
The residues \(R^+,R^\times\) are the visible publication of the two sheets;
the quotients \(Q^+,Q^\times\) preserve the exact memory of Euclidean
division; and \(\tau\) orients the transition. The visible signature occupies the first
level; the cell bearing it adds the arithmetic antecedents, and the integral
holonomic state then completes boundary, path, phase and survival.

It is useful to distinguish four strata. The visible signature is \(\tau\); the
residue--carry pair is its local lift; \(c_{\rm MD}\) is the APP
arithmetic cell equipped with that signature; and the TPK state adds extension,
boundary, genealogical ledger, phase, and survival. This hierarchy preserves the distinct
identity of the observed coordinate, its lift, the carrier cell, and the
enriched history that transports it.

In the balanced ternary chart, the lifted local integer
\(n\in\{-4,-3,\ldots,3,4\}\) decomposes uniquely as
\begin{equation}
\widehat n=(r,\kappa),
\qquad n=r+3\kappa,
\qquad r,\kappa\in\{-1,0,+1\}.
\label{p12-83-c3-s1:eq:trit-balanced-lift}
\end{equation}
The visible signature is \(\tau=r\); \(\kappa\) conserves the carry that the projection
\(n\mapsto r\) would forget.
The fiber over the visible zero contains three distinct states,
\[
0^+=(0,+1),\qquad 0^0=(0,0),\qquad 0^-=(0,-1),
\]
because \(n=-3,0,3\) have the same visible residue and different carries. This
is the minimal topological lift of MD: an oriented APP transition
is endowed with the information required to invert it and to decide which
prolongations remain compatible.

The signature also has a sheet-memory rule. If \(m_k\) is the active
arithmetic mode, then
\begin{equation}
m_{k+1}=
\begin{cases}
\Sigma,&\tau_k=+1,\\
\Pi,&\tau_k=-1,\\
m_k,&\tau_k=0.
\end{cases}
\label{p12-83-c3-s1:eq:trit-zero-memory}
\end{equation}
Zero encodes retention of the previous sheet. For that reason the word
\((+1,-1,0)\) publishes the sequence \((\Sigma,\Pi,\Pi)\): the third state
conserves the active multiplicative mode.

The temporal relation is obtained by orderly parametrizing these cells.
Each signature carries a quadratic algebra.
\begin{equation}
\mathbb A_\tau=\mathbb R[X]/(X^2+\tau),
\qquad J_\tau=[X],
\label{p12-83-c3-s1:eq:trit-quadratic-algebra-local}
\end{equation}
so that
\[
\begin{array}{c@{\qquad}c@{\qquad}l}
\toprule
\tau & J_\tau^2 & \text{ordered temporal relation}\\
\midrule
+1 & -I & \text{elliptic return, memory and past},\\
0  & 0  & \text{parabolic threshold and present},\\
-1 & +I & \text{bidirectional hyperbolic opening and future}.\\
\bottomrule
\end{array}
\]
The TRIT precedes external time: it first classifies the algebraic and
topological type of the transport; an ordered parameterization then publishes it
as a temporal relation. The twofold direction is typed by the exact involution
\begin{equation}
\mathcal C_{\rm rev}(\tau_1,\ldots,\tau_n)
=(-\tau_n,\ldots,-\tau_1),
\qquad \mathcal C_{\rm rev}^2=I,
\label{p12-83-c3-s1:eq:trit-time-reversal}
\end{equation}
which exchanges return and opening and preserves the threshold. On the
hyperbolic sheet, \(P_\pm=(I\pm J_{-1})/2\) separate the two eigendirections, and
reversal exchanges their orientations. Topology and time are thus the two
coordinated readings of the TRIT.

The combinatorial involution \(\mathcal C_{\rm rev}\) and an operatorial
conjugation \(C\) belong, however, to distinct domains. In a
separate dynamic realization, endowed with \(H=H^*\),
\(CJ_EC^{-1}=-J_E\), and \(CHC^{-1}=H\), we fix an internal action scale
\(\hbar_{\rm int}>0\) and write
\(h_{\rm int}:=2\pi\hbar_{\rm int}\). Then
\begin{equation}
U(t)=\exp(-J_EHt/\hbar_{\rm int}),
\qquad CU(t)C^{-1}=U(-t).
\label{p12-83-c3-s1:eq:trit-bidirectional-propagator}
\end{equation}
The intertwiner that transports TRIT words to the dynamical realization fixes
the correspondence between \(\mathcal C_{\rm rev}\) and \(C\). The twofold
local orientation provides the domain of that construction.

\subsection{TPK: category of transport and phase evolution}

The \emph{Topological Phase Kernel} is the transport category over the
observable base of APP that composes digits, words, operators and prolongation
relations in the same integral holonomic state. Its internal architecture is distributed among
eight operator classes, differentiated by domain, codomain and function:
\begin{enumerate}[label=\textup{(\roman*)},leftmargin=2.25em]
\item observable support and integral holonomic state;
\item internal sheet selector \(M_{\rm ph}:\{1,\ldots,9\}\to\{+1,-1,0\}\);
\item cardinal translations and visible evolution \(F_{\rm obs}\);
\item readers and quotients nonadic, tritico and of block \(1000\);
\item memory, carry, cylinder, boundary, and genealogical register;
\item periodicity and return
\(C_{12}\hookrightarrow C_{108}\twoheadrightarrow C_9\),
\(\mathcal K_{27}\), and channels \(90/120\);
\item prolongations \(\operatorname{Ext}_r\), cylinders, survival, and
natural bidirectional extension;
\item Archimedean-output, calibrated-exceptional, and dimensional functors.
\end{enumerate}
The return of 27 steps is the composite morphism.
\begin{equation}
\mathcal K_{27}:=F_{\rm obs}^{27},
\qquad \mathcal K_{27}^{4}=F_{\rm obs}^{108}.
\label{p12-83-c3-s1:eq:tpk-kernel-27}
\end{equation}
Its fourth iterate restores the observable phase of the calendar, while the
memory and the complete TPK state continue along their trajectory;
\(\mathcal K_{27}\) is the internal return morphism of \(C_{108}\).
The calendar \(C_{108}\), the nonadic connection, the bilateral extension, and
the readers occupy typed places within this architecture. Let
\(\mathsf P_{\rm APP}\) be the free category
of the directed graph of APP, and let
\(\mathsf P_{\rm APP}^{(2)}=\mathsf P_{\rm APP}\times
\mathsf P_{\rm APP}\) be its simultaneous sum--product reading. Let, moreover,
\[
\mathcal Q_{\rm obs}
=X_9^2\times\{N,E,S,O\}^2\times\Theta_{108},
\qquad \Theta_{108}=\mathbb Z/108\mathbb Z,
\]
where, for the representative
\(\widetilde\theta\in\{0,\ldots,107\}\),
\[
r(\theta)=1+(\widetilde\theta\bmod9),
\qquad g_0(\theta)=[\widetilde\theta]_9,
\]
\[
\tau(\theta)=
\begin{cases}
+1,&r(\theta)\in\{1,4,7\},\\
-1,&r(\theta)\in\{2,5,8\},\\
0,&r(\theta)\in\{3,6,9\},
\end{cases}
\qquad
\mu(\theta)=
\begin{cases}
\Sigma,&\tau(\theta)=+1,\\
\Pi,&\tau(\theta)=-1,\\
\varnothing,&\tau(\theta)=0.
\end{cases}
\]
The calendar label \(\mu=\varnothing\) marks the threshold, and the dynamic memory \(m_k\) of \eqref{p12-83-c3-s1:eq:trit-zero-memory} retains the last active sheet. A configuration \(q=(p^\Sigma,p^\Pi,d^\Sigma,d^\Pi,\theta)\) preserves the two positions, the two cardinal directions, and the calendar index. The visible evolution \(F_{\rm obs}\) activates the sheet indicated by TRIT, transports the corresponding position, and advances \(\theta\). Let \(\mathsf P_{\rm obs}\) denote the free category generated by the arrows \(q\to F_{\rm obs}(q)\). Each arrow induces a pair of APP paths and defines the base functor
\begin{equation}
B:\mathsf P_{\rm obs}\longrightarrow\mathsf P_{\rm APP}^{(2)}.
\label{p12-83-c3-s1:eq:tpk-bivariant-base-functor}
\end{equation}

Over each \(q\) lies the typed fibre
\begin{equation}
\mathcal F_q=
\mathsf O_q\times\mathsf H_q\times\mathsf C_q\times\mathsf S_q
\times\mathsf B_q\times\mathsf\Lambda_q,
\label{p12-83-c3-s1:eq:tpk-fiber}
\end{equation}
Its six factors conserve, respectively, orientation and sheet; residue and
quotient; carry; signature; cylinder and boundary; and ordered register of
transport. Memory and survival are properties of the composition and of
the prolongation relations.

For a route \(r:q\to q'\), the relation
\begin{equation}
\operatorname{Ext}_r\subseteq\mathcal F_q\times\mathcal F_{q'}
\label{p12-83-c3-s1:eq:tpk-extension-relation}
\end{equation}
contains all admissible continuations and satisfies
\[
\Delta_q\subseteq\operatorname{Ext}_{\operatorname{id}_q},
\qquad
\operatorname{Ext}_{r_2}\circ\operatorname{Ext}_{r_1}
\subseteq\operatorname{Ext}_{r_2\circ r_1}.
\]
Transports respect identities and composition:
\begin{equation}
\begin{gathered}
\mathsf H(\operatorname{id}_q)=I,
\quad \mathsf C(\operatorname{id}_q)=I,
\quad \mathsf\Lambda(\operatorname{id}_q)=I,
\quad \kappa(\operatorname{id}_q)=0,\\
\mathsf H(r_2\!\circ r_1)=\mathsf H(r_2)\mathsf H(r_1),
\quad
\mathsf C(r_2\!\circ r_1)=\mathsf C(r_2)\mathsf C(r_1),
\quad
\mathsf\Lambda(r_2\!\circ r_1)=
\mathsf\Lambda(r_2)\mathsf\Lambda(r_1).
\end{gathered}
\label{p12-83-c3-s1:eq:tpk-functorial-transport}
\end{equation}
Along with these functorial laws, the transported coordinates obey
\begin{equation}
\begin{aligned}
h'&=\mathsf H(r)(h),\\
c'&=\mathsf C(r)(c)+\kappa(r),\\
\lambda'&=\mathsf\Lambda(r)(\lambda),\\
\sigma'&=\operatorname{Sig}_{q'}(h',c',\lambda'),
\end{aligned}
\qquad
\kappa(r_2\circ r_1)
=\kappa(r_2)+\mathsf C(r_2)\kappa(r_1).
\label{p12-83-c3-s1:eq:tpk-typed-transport}
\end{equation}
The last identity is the carry cocycle that makes the composition with memory associative.

The TPK is then the category \(\mathsf{TPK}\) whose objects are
\begin{equation}
x=(q,o,h,c,\sigma,C,b,\lambda),
\qquad (o,h,c,\sigma,(C,b),\lambda)\in\mathcal F_q,
\label{p12-83-c3-s1:eq:tpk-enriched-object}
\end{equation}
and whose morphisms \(x\to x'\) are the triples \((r,x,x')\) such that
the coordinates satisfy \eqref{p12-83-c3-s1:eq:tpk-typed-transport} and
\((f_x,f_{x'})\in\operatorname{Ext}_r\), where
\(f_x\in\mathcal F_q\) and \(f_{x'}\in\mathcal F_{q'}\). Composition is
\begin{equation}
(r_2,x',x'')\circ(r_1,x,x')
=(r_2\circ r_1,x,x''),
\label{p12-83-c3-s1:eq:tpk-morphism-composition}
\end{equation}
\Needspace{7\baselineskip}
the identity is \((\operatorname{id}_q,x,x)\), and the functor
\begin{equation}
p:\mathsf{TPK}\longrightarrow\mathsf P_{\rm obs},
\qquad p(x)=q,
\qquad p(r,x,x')=r,
\label{p12-83-c3-s1:eq:tpk-base-functor}
\end{equation}
publishes the visible configuration without identifying states of the same fiber.
The definition fixes a category over a base; it does not presuppose that it is a
categorical fibration nor that every arrow admits a Cartesian lift.
The local projections remain separated by type:
\[
\pi_{\rm TRIT}:\operatorname{Ob}(\mathsf{TPK})
\longrightarrow\mathrm{TRIT}_{\rm vis},
\qquad
\pi_{\rm TRIT}(x)=\tau(\theta),
\]
\[
\pi_{\rm loc}(x)=
\bigl(\pi_{\rm TRIT}(x),\mu(\theta),o,h,c,
\sigma,\lambda,g_0(\theta)\bigr).
\]
The first recovers only the visible TRIT signature. The second additionally
preserves the calendrical sheet label, the orientation, the
residue--quotient pair, the carry, the transport signature, the genealogical register and the phase,
but still forgets the cylinder, its boundary label and other coordinates of
\(q\). This is why neither visible TRIT nor its local register
replaces the integral holonomic state of the TPK.

The name encompasses three inseparable functions. \emph{Topological} refers to
the category of paths, incidence, boundary, and prolongation by
cylinders; \emph{Phase} refers to the nonadic index, the oriented calendar, and
holonomy; \emph{Kernel} names the persistent transport kernel that
remains before the state is contracted by a reader. Formally, the
object is the preceding category: an additional topology on objects or
morphisms, or its identification with the linear kernel of a map, requires
independent data.

Consequently, two routes with the same visible endpoint may remain
objects distinct by sheet, quotient, carry, boundary, cylinder, phase, or
register. The TPK does not choose a scalar output: it conserves the common antecedent
from which the Archimedean, exceptional, and dimensional publications depart.

% Estas consecuencias se publican después de la construcción correlativa y
% del paso al límite de los capítulos 4 y 5. Se definen aquí para preservar
% íntegramente su procedencia, pero no se ejecutan antes de que exista el
% continuo al que se aplican.
\long\def\HMTDeferredContinuumConsequences{%
\section{Consequences and publications of the constructed continuum}

\subsection{Terminal Assembly Theorem of the Correlative Constructions Generated by the Common State}

Each depth \(k\) has an integral holonomic state \(\mathcal E_k\) with
all its surviving emissions and a single correlative constructor
\begin{equation}
\mathcal O_k:\mathcal E_k\longrightarrow\mathcal Z_k,
\qquad
\widehat x_k\longmapsto
\mathcal O\bigl(\{\mathfrak e_\alpha:
\alpha\in\operatorname{Ch}_k(\widehat x_k)\}\bigr).
\label{p12-83-c3-s1:eq:five-features-single-constructor}
\end{equation}
The value of \(\mathcal O_k\) simultaneously preserves the spectral,
solenoidal, gauge, coherent, and determinantal coordinates,
\begin{equation}
(\mathrm{sp},\mathrm{sol},\mathrm{gau},\mathrm{coh},\mathrm{det}),
\label{p12-83-c3-s1:eq:five-features-coordinates}
\end{equation}
as five characteristics of the same state and of the same genealogical register. The
truncations
\[
u^{\mathfrak e}_{\ell,k}:\mathcal E_\ell\to\mathcal E_k,
\qquad
u^{\mathcal O}_{\ell,k}:\mathcal Z_\ell\to\mathcal Z_k
\]
satisfy the unique naturality
\begin{equation}
u^{\mathcal O}_{\ell,k}\circ\mathcal O_\ell
=\mathcal O_k\circ u^{\mathfrak e}_{\ell,k},
\qquad \ell\ge k.
\label{p12-83-c3-s1:eq:five-features-naturality}
\end{equation}
Therefore, the corresponding action passes to the limit as
\begin{equation}
\mathcal O_\infty:
\mathcal C_{\rm cont}^{\rm disc}:=\varprojlim_k\mathcal E_k
\longrightarrow
\mathcal Z_\infty:=\varprojlim_k\mathcal Z_k.
\label{p12-83-c3-s1:eq:five-features-infinite-operator}
\end{equation}
The five terminal applications are seen as intrinsic manifestations of that single value:
\begin{equation}
\boxed{
\mathfrak G_T^{\rm int}:=
\operatorname{View}_T\!\left(
\mathcal O_\infty(\mathcal C_{\rm cont}^{\rm disc})\right),
\quad
T\in\{\mathrm{sp},\mathrm{sol},\mathrm{gau},
\mathrm{coh},\mathrm{det}\}.}
\label{p12-83-c3-s1:eq:five-features-views}
\end{equation}
Simultaneity precedes the specialized resolutions: each view exposes
the required coordinate and the complete state conserves the other four.

\subsection{Survival and inverse limit}

Let \(\operatorname{Surv}^{(\chi)}_N\) be the set of histories of
depth \(N\) that simultaneously satisfy the active laws of channel
\(\chi\) and admit prolongation. We define the compatible truncations
\[
\tau_{N+1,N}:\operatorname{Surv}^{(\chi)}_{N+1}
\longrightarrow\operatorname{Surv}^{(\chi)}_N
\]
These preserve the enriched prefix. The genealogical continuum is the limit.
\begin{equation}
X_\chi=\varprojlim_N
\bigl(\operatorname{Surv}^{(\chi)}_N,\tau_{N+1,N}\bigr),
\qquad
\tau_{N+1,N}(x_{N+1})=x_N.
\label{p12-83-c3-s1:eq:inverse-survival}
\end{equation}
The global space of surviving sections and one of its enriched sections
are denoted, respectively, by
\begin{equation}
X_\infty:=\bigsqcup_\chi X_\chi,
\qquad
\widehat\Omega_\infty\in X_\infty.
\label{p12-83-c3-s1:eq:global-survival-space}
\end{equation}
The global Archimedean reader is
\begin{equation}
P_{\rm arq}:=\operatorname{ev}_{\mathbb R}:
X_\infty\longrightarrow\mathbb R;
\label{p12-83-c3-s1:eq:archimedean-reader}
\end{equation}
its restriction to each channel evaluates the unitary intersection of the compatible cylinders of the corresponding section. When the exceptional reader intervenes, \(X_\infty^{\rm cal}\subseteq
X_\infty\) denotes the subspace in which the calibration type of \eqref{p12-83-c3-s1:eq:exceptional-calibration} has been fixed. Survival determines the section by internal compatibility: every history must be prolongable in the complete system. The numerical reader intervenes afterward to certify and publish the intersection of the already compatible descendant.

\subsection{Joint nonadic connection, holonomy, and vertical memory}

At depth \(K\), the holonomy of the nonadic return is initially the typed correspondence
\begin{equation}
\Gamma_{9,K}=\operatorname{Hol}_{C_9}(\nabla_{\rm TPK}):
\widehat{\mathcal X}_K\rightrightarrows
\widehat{\mathcal X}_{K+9}.
\label{p12-83-c3-s1:eq:gamma9}
\end{equation}
Its iterates are relational compositions. Only when each state has
a unique compatible prolongation is it restricted to a map, and only when that
map is invertible does it constitute a monodromy. This distinction conserves all
surviving extensions before selecting a section.
After nine phases the local phase returns, while the complete state advances:
\[
g(K+9)=g(K),\qquad q_{{\rm mem},K+9}=q_{{\rm mem},K}+1,
\qquad \widehat x_{K+9}\ne\widehat x_K\ \text{in general}.
\]
Block, cylinder, prefix, carry, and memory continue to evolve. This
is the difference between visible periodicity and return with biography.

The transported state at level (K) can be exhibited as
\begin{equation}
\widehat x_K=
(B_K,C_K,p_K,c_K,\Sigma_K,W_K,g(K),q_{{\rm mem},K},s_K),
\label{p12-83-c3-s1:eq:enriched-state-k}
\end{equation}
where each coordinate has its own transport law. The phase return is represented on the catalog.
\[
\mathcal U_6^{\rm cat}\cong\mathfrak S_+\times\mathfrak S_\times,
\qquad |\mathfrak S_+|=18,
\qquad |\mathfrak S_\times|=26.
\]
Conditional averages
\[
(E_+f)(s,e)=\frac1{18}\sum_{s'\in\mathfrak S_+}f(s',e),
\qquad
(E_\times f)(s,e)=\frac1{26}\sum_{e'\in\mathfrak S_\times}f(s,e')
\]
commute. The phase word
\(\tau=(+1,-1,0,+1,-1,0,+1,-1,0)\) applies
\(E_+,E_\times,I\) in that order, and one revolution closes through
\begin{equation}
\mathcal K_{\rm ph}^{\,9}
=(E_+E_\times)\otimes I_{\ell^2(C_9)},
\qquad
\bigl\langle\delta_u,E_+E_\times\delta_v\bigr\rangle
=\frac1{468}\quad(u,v\in\mathcal U_6^{\rm cat}).
\label{p12-83-c3-s1:eq:joint-nine-holonomy}
\end{equation}
The sum and product sheets do not act as independent filters: they are two
components of a single holonomy that closes only after traversing the
complete cycle.

The compression--expression colligation separates what returns from what remains
as memory. If (J) embeds the history space into the phase
space, (T) is the surviving compression and \(\eta\) the defect channel,
then
\begin{equation}
\boxed{
U_{\Gamma_9}J=JT+\eta,
\qquad J^*\eta=0,
\qquad I-T^*T=\eta^*\eta.}
\label{p12-83-c3-s1:eq:compression-expression}
\end{equation}
The first equality decomposes transport into compressed return and memory
expression; the second proves their orthogonality; the third measures exactly
the information that cannot fit within the periodic shadow. The central
statement is thereby formalized: the phase returns, not the complete state.

The vertical memory is organized through the cyclic extension oriented
\begin{equation}
0\longrightarrow C_{12}\xrightarrow{\,\iota\,}C_{108}
\xrightarrow{\,p\,}C_9\longrightarrow0,
\qquad
c(r,r')=\left[\left\lfloor\frac{r+r'}9\right\rfloor\right]_{12}.
\label{p12-83-c3-s1:eq:c108-extension}
\end{equation}
Here \(\iota([b]_{12})=[9b]_{108}\) and
\(p([\theta]_{108})=[\theta]_9\). The extension does not split:
\(C_{108}\) has exponent 108, whereas
\(C_{12}\times C_9\) has exponent 36.
The cocycle records the carry from each turn of the base \(C_9\) in the
dodecaphase memory \(C_{12}\). \(C_{108}\) is therefore an oriented
phase--memory calendar. Only after the dimensional chart is that
extension realized as a finite 108-step quantum clock.

The local MD lift carrying the TRIT signature is the
topological--temporal unit of that genealogical crystal.
On
\(\mathcal H_{\rm clk}=\mathbb C^{108}\), with basis
\(\{|j\rangle:j\in\mathbb Z/108\mathbb Z\}\), let
\[
\zeta=e^{2\pi\ii/108},\qquad
|\widetilde k\rangle=\frac1{\sqrt{108}}
\sum_{j=0}^{107}\zeta^{jk}|j\rangle,
\qquad U_{\rm clk}|j\rangle=|j+1\rangle.
\]
\Needspace{12\baselineskip}
With \(\omega_0=2\pi/(108t_0)\), the Hamiltonian
\begin{equation}
H_{\rm clk}=\hbar_{\rm int}\omega_0
\sum_{k=0}^{107}k\,|\widetilde k\rangle\langle\widetilde k|
\label{p12-83-c3-s1:eq:trit-time-crystal-hamiltonian}
\end{equation}
generates exactly
\begin{equation}
U_{\rm step}=\exp\!\left(
-\frac{\ii H_{\rm clk}t_0}{\hbar_{\rm int}}\right)=U_{\rm clk},
\qquad U_{\rm step}^{108}=I,
\label{p12-83-c3-s1:eq:trit-time-crystal-period}
\end{equation}
If \(Z_{\rm clk}|j\rangle=\zeta^j|j\rangle\), then
\begin{equation}
\langle j_0|U_{\rm step}^{-n}Z_{\rm clk}U_{\rm step}^{n}|j_0\rangle
=\zeta^{j_0+n},
\label{p12-83-c3-s1:eq:trit-time-crystal-observable}
\end{equation}
whose orbit has primitive period 108. In HMT vocabulary, this is the
internal genealogical time crystal: the observable calendar closes,
\(C_{12}\) records the turn and the integral holonomic state need not
return. The subsequent realization as a macroscopic physical phase additionally
requires a Floquet protocol, primitive subharmonic response, robustness,
persistence and a protection mechanism; that chart does not alter the finite
theorem of unitary periodicity.

\subsection{Archimedean publication, genealogical number and measurement}

The Archimedean publication uses semi-open ternary cylinders.
\begin{equation}
I_{N,L}=\left[\frac{N}{3^L},\frac{N+1}{3^L}\right),
\qquad \operatorname{diam}I_{N,L}=3^{-L}\longrightarrow0.
\label{p12-83-c3-s1:eq:ternary-cylinders}
\end{equation}
A compatible section determines nested cylinders and, by completeness, a
unique real point. The HMT number is the section itself,
\begin{equation}
x=(x_N)_N\in X_\chi
=\varprojlim_N\operatorname{Surv}^{(\chi)}_N,
\label{p12-83-c3-s1:eq:hmt-number-as-section}
\end{equation}
not the tuple formed by it and its evaluations. The value appears when a
reader forgets part of the genealogy.
Therefore, precision is depth: measuring means descending to a
finer cylinder while the causal prefix remains in the antecedent.
The real line is reconstructed as the Archimedean quotient of a space of
much richer histories.

\subsection{Dimensional realization of the HMT state: screen map, observable magnitude, and spectral chart}

Dimensional Mechanics begins where the genealogical state couples to a
realization boundary. With a field \(K_n\) fixed, let \(M_n(x)\) be the module
generated by the prolongation channels of a state \(x\), let
\(\mathcal R_n(x)\) be the submodule of its relations, and let \(\beta_n(x)\) be the
conserved boundary datum. The dimensional screen is
\begin{equation}
\Sigma_n(x)=\bigl(K_n,M_n(x),\beta_n(x),\mathcal R_n(x)\bigr),
\qquad
d_{\Sigma_n}(x)=\dim_{K_n}\frac{M_n(x)}{\mathcal R_n(x)}.
\label{p12-83-c3-s1:eq:md-screen}
\end{equation}
If \(M_n(x)\simeq K_n^{g_n}\) and the columns of \(R_n(x)\) generate the
relations, then
\begin{equation}
\boxed{d_{\Sigma_n}(x)=g_n-\operatorname{rank}_{K_n}R_n(x).}
\label{p12-83-c3-s1:eq:md-rank}
\end{equation}
Dimension is not attached as a label: it is the number of independent
channels that survive after imposing the boundary relations.

A physical quantity is a section of a dimensional line \(L_D\); a unit is a chart \(\chi_u:L_D\to\mathbb R\) that publishes its coordinate. Hence
\begin{equation}
\text{TPK state}\longrightarrow\Sigma_n(x)
\longrightarrow s_D\in\Gamma(L_D)
\xrightarrow{\;\chi_u\;}[s_D]_u\in\mathbb R
\label{p12-83-c3-s1:eq:md-measurement-chain}
\end{equation}
separates structure, dimensional type, and measured number exactly. Measuring
consists in coupling a section to a reversible apparatus, imposing
edge compatibility, and applying the reader; changing units changes
\(\chi_u\), not the section or its genealogy.

\subsection{Double projection mathematical and dimensional realization}

The exceptional branch preserves gauge data that the real publication does not
need. Let \(\mathscr C_{\rm exc}\) be the space of gauges, and let
\begin{equation}
\mathfrak E:
X_\infty^{\rm cal}\times\mathscr C_{\rm exc}
\longrightarrow\mathbf{Exc}
\label{p12-83-c3-s1:eq:exceptional-calibrated-reader}
\end{equation}
the incidence reader. We fix once
\begin{equation}
\begin{aligned}
\chi_{\rm exc}=\bigl(&\text{projective order},\ \text{Paley gauge},
 \ \text{dodecad},\\
 &\text{incidence alignment},\ \text{lattice chamber}\bigr)
 \in\mathscr C_{\rm exc}
\end{aligned}
\label{p12-83-c3-s1:eq:exceptional-calibration}
\end{equation}
and abreviamos
\(\mathfrak E_{\chi_{\rm exc}}(x):=\mathfrak E(x,\chi_{\rm exc})\).
The gauge initial it transports by the history; not vuelve to elegirse in
each depth.

The same calibrated limit then admits two readers with distinct codomains:
\begin{equation}
\mathbb R
\xleftarrow{\;\operatorname{ev}_{\mathbb R}\;}
X_\infty^{\rm cal}
\xrightarrow{\;\mathfrak E_{\chi_{\rm exc}}\;}
\mathbf{Exc}.
\label{p12-83-c3-s1:eq:double-projection}
\end{equation}
The first branch contracts cylinders to a continuous value. The second
preserves incidences, orientations, gauge, and boundary. Both descend in
parallel from the same source state. This bifurcation jointly explains
metric continuity and the emergence of exceptional symmetries. The second
branch is a calibrated reconstruction: it is exact once
\(\chi_{\rm exc}\) has been fixed, but the bare state does not retrospectively select
the five data composing that gauge.

The non-serializing dimensional realization does not serialize those two branches. It is a third reader of the same antecedent:
\begin{equation}
P_{\rm dim}:X_\infty\longrightarrow
\mathcal R_{\rm MD},
\qquad
\mathcal R_{\rm MD}=
\{\text{action, mass, field, geometry, information}\}.
\label{p12-83-c3-s1:eq:dimensional-publication}
\end{equation}
In this manner, a number, a symmetry, and a physical magnitude may share a
genealogy without being confused with or reduced to one another.

The thesis of the continuum is thus formulated precisely: Archimedean
reality is the zero-radius shadow of infinitely refinable discrete histories;
what the shadow omits reappears as memory, incidence,
geometry, and physical structure.

% RESULTADO_RECUPERADO. Owner íntegro de la dinámica del cociente, el retorno
% observable y los cociclos de modo y orientación. Se subordina a la única
% construcción APP--TRIT--TPK sin abrir un origen alternativo.
\begin{HMTScientificOwnerSubordinated}
\chapter{Surviving extensions, monodromy, and dynamics of the quotient}
\label{chap:extensiones-dinamica-cociente}
\index{survival!global extension}
\index{monodromy!nonadic return}
\index{measure!emissions catalog}

The language of paths requires a decisive precision. An APP path,
a visible emission, an enriched TPK history, and a global section
are not four names for the same object. They constitute four consecutive levels
of a single construction. The path preserves positional order;
the emission preserves a finite reading; the history adds signature,
carry, orientation, boundary, and memory; the global section requires all
of its prefixes to survive simultaneously at every depth.

This hierarchy makes it possible to separate two questions that had remained intermingled. The first is logical: under what hypotheses a finite prefix truly belongs to a compatible infinite extension. The second is dynamical: what it means to transport a single history and what it means to transfer all its compatible continuations simultaneously. The distinction is essential in APP\@. A given trajectory carries one state to another; a path transfer carries a distribution of extensions to another distribution.

This section resolves the finite problem at depth six. The
published deterministic chronology has visible-phase return and cannot select
a unique probability. By contrast, transport biased by the nonadic phase,
when each active channel is interpreted as the normalized renewal of its entire
fibre of compatible continuations, has a uniform nine-step return
over the (468) emissions. The weights (1/3) then cease to be
an external choice: they follow from the three occurrences of each signature
(+1,-1,0) in the oriented nonadic cycle. Normalization within each fibre remains
a precise mathematical premise ---the uniform count of all compatible
paths--- and must not be conflated with the successor of a single
trajectory. Passage to the inverse limit over all depths preserves a
distinct statement, formulated at the end of the section.

\section{\(5\) levels of the notion of path}

Let \(\Gamma_{\rm APP}\) be the category of cardinal paths of the APP torus\@.
For each depth \(m\), let \(\mathcal H_m\) be the finite set of
admissible enriched histories
\[
 h_m=(v_0,e_1,v_1,\ldots,e_m,v_m),
\]
where each arrow \(e_j\) transports all fields declared by TPK\@. If
\(n\geq m\), the application
\[
 p_{n,m}:\mathcal H_n\longrightarrow\mathcal H_m
\]
removes the tail beyond depth \(m\). The positional shadow
\[
 \operatorname{sh}_m:\mathcal H_m\longrightarrow\Gamma_{\rm APP}
\]
forgets the memory fiber and preserves the underlying APP path.

\begin{definition}[Hierarchy of extensions]
\label{def:jerarquia-extensiones}
We distinguish the following objects.
\begin{enumerate}
  \item An \emph{APP path} is a morphism of
  \(\Gamma_{\rm APP}\).
  \item An \emph{observable route} is a reading of that morphism that
  jointly conserves the sum and product leaves.
  \item A \emph{TPK extension} is a lift of the observable route to
  a history \(h_m\in\mathcal H_m\), with its residue, quotient,
  signature, orientation, boundary, carry, and record data.
  \item A \emph{surviving extension} is a finite history that admits
  compatible prolongations to every horizon.
  \item A \emph{global section} is a compatible family
  \(x=(h_m)_{m\geq0}\) of surviving extensions.
\end{enumerate}
In particular, the route is the shadow of the extension; it is not an external object
that the extension subsequently chooses.
\end{definition}

For \(h\in\mathcal H_m\), its cylinder of prolongations is
\[
 Z_m(h)=
 \{x=(h_n)_n\in\varprojlim_n\mathcal H_n:x_m=h\}.
\]
In particular, every family of \(Z_m(h)\) also satisfies
\(p_{n,\ell}(h_n)=h_\ell\) for \(n\geq\ell\); collections of
incompatible truncations are not admitted.
Let us also define
\[
 \mathcal S_m^{(N)}=p_{N,m}(\mathcal H_N),
 \qquad
 \mathcal S_m^\infty=\bigcap_{N\geq m}\mathcal S_m^{(N)}.
\]

\begin{theorem}[Conditional extension survivor criterion]
\label{thm:criterio-extension-superviviente}
Let us fix \(h\in\mathcal H_m\) and, for \(N\geq m\), write
\[
 \mathcal P_N(h)=\{g\in\mathcal H_N:p_{N,m}(g)=h\}.
\]
Suppose that the truncations are coherent and that each
\(\mathcal P_N(h)\) is finite. The following are equivalent:
\begin{enumerate}
  \item \(\mathcal P_N(h)\ne\varnothing\) for every \(N\geq m\), that is,
  \(h\in\mathcal S_m^\infty\);
  \item the tree of finite extensions of \(h\) is finitely branching
  and has a nonempty level at every depth;
  \item \(Z_m(h)\ne\varnothing\).
\end{enumerate}
Consequently, whenever the surviving finite truncations are nonempty, the set
\[
 \Omega_{\rm H}^{\rm surv}
 =\varprojlim_m\mathcal S_m^\infty
\]
is exactly the space of surviving global sections. The theorem does not
assert that these truncations are nonempty for every APP input: that
property must be proved by means of the survival protocol.
\end{theorem}

\begin{proof}
The implication from the third condition to the first two is obtained by
truncation. Now assume the first condition and order the elements of
\(\bigsqcup_{N\geq m}\mathcal P_N(h)\) by the prefix relation. Each level
is finite and nonempty, and coherence of the truncations makes this union
a finitely branching tree with nodes at every depth; this yields the
second condition. Conversely, from the second condition König's lemma
provides an infinite branch. Its truncations form a compatible family
belonging to \(Z_m(h)\), which is the third condition. The same construction
identifies the indicated inverse limit wherever its finite truncations are
nonempty.
\end{proof}

The theorem resolves the possible discrepancy between “extendable to each
horizon separately” and “possessing a global prolongation”: finiteness of the
branches is the hypothesis that permits the choice of a single compatible chain. It does not
claim that all APP paths belong to
\(\operatorname{sh}(\Omega_{\rm H}^{\rm surv})\). That coverage requires proving
that the enriched closures do not empty the corresponding fibers.

\section{Phase return and monodromy relation}

The return of the phase--length quotient is formalized separately in
\cref{def:odometro-9-12,prop:proyeccion-fase-longitud}:
\[
(b,r)\xrightarrow{\;9\;}(b+1,r),
\qquad
(b,r)\xrightarrow{\;108\;}(b,r).
\]
These equalities belong to the odometer \(9\times12\). They do not constitute the
proof of \(F_{\rm obs}^{108}=I\), which uses the observable state and appears
below.

Let \(\phi\) be the phase coordinate with values in \(\mathbb Z/9\mathbb Z\).
A visible-phase return satisfies
\[
 \phi(h_{k+9})=\phi(h_k).
\]
The complete record may, however, satisfy
\[
 B(h_{k+9})\ne B(h_k).
\]
This inequality is not a failure of periodicity: it is the memory of the circuit.
An ordinary action of \(C_9\) on the complete state would require the
ninth power to be the identity and would erase precisely that information.

To type the return, let \(H\) be a hexad and \(P_2(H)\) its set of
fifteen pairs. Let \(\mathcal B_k\) be the record states that survive at phase
\(k\), let \(\mathcal T_k\subseteq\mathcal B_k\times\mathcal B_{k+1}\) be the published transport relation, and
suppose a fiber encoder is given
\[
 c_k:\mathcal B_k\longrightarrow P_2(H).
\]
Only after declaring those encoders is the observable relation defined.
\[
 R_k=(c_k\times c_{k+1})(\mathcal T_k)
 \subseteq P_2(H)\times P_2(H).
\]
The relational monodromy of a loop is
\[
 \mathcal M_k
 =R_{k+8}\circ R_{k+7}\circ\cdots\circ R_k.
\]

\begin{proposition}[Classification of the nonadic return given the encoder]
\label{prop:clasificacion-retorno-nonadico}
Three logically distinct levels are given.
\begin{enumerate}
  \item Without a proof of functionality, the return is the relation
  \(\mathcal M_k\).
  \item If each \(R_k\) is the graph of a map
  \(\mu_k:P_2(H)\to P_2(H)\), the transport is the skew product
  \[
   T(k,p)=(k+1,\mu_k(p)),
  \]
  and
  \[
   T^9(k,p)=(k,M_k(p)),
   \qquad
   M_k=\mu_{k+8}\circ\cdots\circ\mu_k.
  \]
  \item If the \(\mu_k\) are bijective, \(M_k\) is an automorphic holonomy.
  The transport descends to an ordinary action of
  \(C_9\) on the complete state if and only if \(M_k=I\) in all
  fibers.
\end{enumerate}
\end{proposition}

\begin{proof}
The first claim is the definition of image and relational composition once
\((c_k)_k\) have been fixed. Under
functionality, nine consecutive applications keep the first
coordinate fixed and compose their actions on the second. The composition of
bijections is a bijection. Finally, \(T^9=I\) is exactly equivalent to
\(M_k=I\) for each phase and each point of the fiber.
\end{proof}

\begin{corollary}[No return under cumulative transport]
\label{pf84:E02-29}
Let \(q:X\to V\), \(T:X\to X\), and \(x\in X\) be given. If
\[
 q(T^9x)=q(x),\qquad
 T^9x=\mathcal M_x\,x,\qquad
 \mathcal M_x\,x\ne x,
\]
then the projection returns and the state does not:
\[
 q(T^9x)=q(x),\qquad T^9x\ne x.
\]
The pointwise hypothesis \(\mathcal M_x\,x\ne x\) is indispensable; the global
inequality \(\mathcal M_x\ne I\) does not exclude \(x\) from being a fixed point.
\end{corollary}

\begin{proof}
The first equality is a hypothesis. The second and the pointwise separation give
\(T^9x=\mathcal M_x\,x\ne x\). The statement would be refuted by a case that
satisfies the three premises and had \(T^9x=x\); it does not authorize replacing the
third premise by \(\mathcal M_x\ne I\).
\end{proof}

\begin{theorem}[Functional sufficiency of the enriched tuple]
\label{pf84:E02-30}
Let
\[
 E:\mathscr X_{\rm TPK}^{\rm enr}\longrightarrow\mathscr S_E
\]
the tuple that conserves the visible block, new blocks and carry, lifting
state, cylinders, boundary signature, orientation, observable, phase, and
dodecaphasic register. Let
\[
 T:\mathscr X_{\rm TPK}^{\rm enr}\to
 \mathscr X_{\rm TPK}^{\rm enr},
 \qquad
 O:\mathscr X_{\rm TPK}^{\rm enr}\to Y.
\]
When maps exist
\[
 \widehat T:\mathscr S_E\to\mathscr X_{\rm TPK}^{\rm enr},
 \qquad
 \widehat O:\mathscr S_E\to Y
\]
such that
\[
 T=\widehat T\circ E,\qquad O=\widehat O\circ E,
\]
the equality \(E(x)=E(y)\) implies \(T(x)=T(y)\), \(O(x)=O(y)\) and, by
determinism, the same continuations and subsequent outputs.
\end{theorem}

\begin{proof}
By factorization,
\[
 T(x)=\widehat T(E(x))=\widehat T(E(y))=T(y),
 \qquad
 O(x)=\widehat O(E(x))=\widehat O(E(y))=O(y).
\]
The first equality permits iteration of the argument. This is a sufficient
condition, not an unproved equivalence. Two states with the same tuple and
different transitions or outputs, access to an absent coordinate, or
querying a target value are falsifiers. The visible five-element microfiber
of \cref{pf84:E02-22} controls only the nonfaithfulness of the
word \(w_6\); it does not prove this positive factorization.
\end{proof}

The published finite window contains nine phase returns with a change in the
integral holonomic state and a local selection from three continuations to one when
increasing by two horizons. This proves return with memory and refutes the
identification of the register with a memoryless cycle. The global extraction of
the nine encoders \(c_k\) and, through them, of the relations
\(R_k\), from all surviving histories is a distinct operation.
The preceding proposition classifies the return once the register--pairs
map is given; it does not derive that map from the integral holonomic state of TPK.

The concrete episode of the ninth phase, with its three continuations, the
block \(B_{11}\), the boundary \((502,662,638)\), and the integer inequalities
of cylinders, is proved only once in the
\cref{thm:seleccion-novena-fase}. In this residence only its
typed consequence is used:
\[
\#\operatorname{Surv}_{1}(B_9)=3,\qquad
\#\operatorname{Surv}_{2}(B_9)=1.
\]

The other two words can precede the same ternary tail only by changing
the deep Hensel extension. They share part of the visible residue,
but not the genealogy of the state.

\section{The catalog of 468 emissions as a product of channels}

The finite catalogue starts from the set
\(\Omega_{\rm seed}\) of \(104\,976\) microscopic presentations. Let
\(\mathcal U_6\) be the set of distinct sextuple emissions and let
\[
 F:\Omega_{\rm seed}\twoheadrightarrow\mathcal U_6.
\]
The exact census is
\[
 |\mathcal U_6|=468,
 \qquad
 \#\{U:|F^{-1}(U)|=144,432,1944\}=432,18,18.
\]
These counts belong to \(F\). Comparing them with TPK histories of
depth six would require an explicit map from the domain of histories
to \(\Omega_{\rm seed}\); no result in this section presupposes that
map or identifies catalog presentations with integral holonomic states.
The finite catalog provides, for each emission, an additive signature
\(s\) and a multiplicative signature \(e\). The joint map is a
bijection
\[
 \mathcal U_6\xrightarrow{\sim}
 \mathcal S_+\times\mathcal S_\times,
 \qquad
 |\mathcal S_+|=18,
 \quad
 |\mathcal S_\times|=26.
\]
Thus,
\[
 \boxed{468=18\cdot26}
\]
is a factorization of the object and not merely of its cardinality. The first
coordinate preserves how the emission sums; the second, how the
product accumulates. This separation is the explicit finite form of
sum--product incompatibility: both readings coexist in the state, but
each projection forgets the complementary signature.

The first coordinate admits an additional parametrization that will be useful when
comparing it with the nonadic boundary. If an additive seed has coordinates
\((i,j)\in(\mathbb Z/9\mathbb Z)^2\) and orientation
\(\varepsilon=+1\) for the east--south directions and \(-1\) for
north--west, its complete signature is determined bijectively by
\[
 \bigl(i+j\bmod9,\varepsilon\bigr)
 \in(\mathbb Z/9\mathbb Z)\times\{\pm1\}.
\]
Therefore,
\[
 \boxed{\mathcal S_+\simeq
 (\mathbb Z/9\mathbb Z)\times\{\pm1\}.}
\]
The product coordinate has no analogous uniformity: its \(26\) signatures
have fibers of seeds of sizes \(8\), \(24\), and \(108\), with histogram
\(24\times8+1\times24+1\times108=324\). This asymmetry is structural. The
factor \(18\) already contains phase and orientation; any subsequent descent
of the product factor to a class of pairs must preserve additional memory or
accept nonuniform fibers.

On \(\ell^2(\mathcal S_+\times\mathcal S_\times)\), let us define the
conditional expectations
\[
\begin{aligned}
 (E_+f)(s,e)&=\frac1{18}\sum_{s'\in\mathcal S_+}f(s',e),\\
 (E_\times f)(s,e)&=\frac1{26}\sum_{e'\in\mathcal S_\times}f(s,e').
\end{aligned}
\]
The first operator renews the additive signature and preserves the multiplicative
one; the second performs the dual operation. The third channel retains both.

\section{The unique path and path transfer}

Unlike the return of the odometer defined in
\cref{def:odometro-9-12,prop:proyeccion-fase-longitud}, the following proposition
proves the return of the observable operator \(F_{\rm obs}\).

The published observable law updates a single cursor at each step. In the
phases \(+1\), the additive cursor advances; in the phases \(-1\), the
multiplicative cursor advances; in the phases \(0\), both remain stationary. After the
steps \(27,54,81,108\), the two orientations are inverted. Denote this
deterministic evolution by \(F_{\rm obs}\).

\begin{proposition}[Return of the Observable Chronology]
\label{prop:retorno-fobs-identidad}
For every choice of the two initial positions and orientations, the
positions and orientations are restored after \(54\) steps. Upon
also restoring the phase, one obtains
\[
 F_{\rm obs}^{108}=I.
\]
Therefore, the Poincaré return of \(108\) steps induces the identity on
\(\mathcal U_6\) and does not select a unique stationary measure.
\end{proposition}

\begin{proof}
In every block of \(27\) steps, each cursor makes exactly nine advances in a fixed orientation. Since the position space is \(X_9=(\mathbb Z/9\mathbb Z)^2\), its displacement is \(9d=0\). At the end of the block, \(d\) is inverted. Two blocks restore position and orientation; four also restore the index in \(\Theta_{108}\). The induced return is consequently the identity. Every probability is stationary for the identity, so none is selected by that chronology.
\end{proof}

The problem arises even before the return if one attempts to forget the
route memory. The product signature
\[
 e_0=(5,5,5,5,5,5)
\]
has twenty-four representatives. After an advance in its own orientation,
those representatives separate, eight by eight, among
\[
 (5,5,5,7,1,7),\qquad
 (5,5,5,7,7,1),\qquad
 (5,5,5,1,7,7).
\]
Thus, the product advance does not descend to a map of the twenty-six
signatures: the representative that had been forgotten becomes relevant again. This
witness explains why an individual route cannot be turned, by mere
projection, into the refreshment of a complete fiber.

The APP reading is different. Instead of choosing a representative, it preserves all
compatible continuations. To formalize it in the finite catalogue,
consider the subalgebras of functions that depend only on one of the two
signatures. With respect to uniform counting, \(E_+\) and \(E_\times\) are the
conditional expectations that forget exactly the active coordinate. They are the
only Markov projectors that simultaneously satisfy:
\begin{enumerate}
  \item preservation of the coordinate inactive;
  \item complete forgetting of the active coordinate;
  \item preservation of the uniform counting state.
\end{enumerate}
Indeed, a row that has forgotten the source coordinate must be
constant on the \(18\), respectively \(26\), compatible destinations; the
normalization forces the values \(1/18\) and \(1/26\). This is the relative
canonicity of the renewal: no preference is postulated between branches that APP
declares equally possible.

Let \(C_9=\mathbb Z/9\mathbb Z\), with the phase word
\[
 \tau=(+1,-1,0,+1,-1,0,+1,-1,0).
\]
Write
\[
 P_{+1}=E_+,
 \qquad P_{-1}=E_\times,
 \qquad P_0=I.
\]

\begin{definition}[Phase-biased TPK transfer]
\label{def:transferencia-tpk-fase}
On \(\mathcal U_6\times C_9\) we define
\[
 \mathcal K_{\rm ph}\bigl((u,r),(v,r+1)\bigr)
 =P_{\tau(r)}(u,v),
\]
and set the remaining entries to zero. This definition is relative to the reader of
all compatible paths: an active phase uniformly renews the fiber
read by that phase, and the neutral phase retains the emission.
\end{definition}

\begin{theorem}[Uniform nonadic return]
\label{thm:retorno-nonadico-uniforme}
The kernel \(\mathcal K_{\rm ph}\) is doubly stochastic and irreducible,
has period exactly nine, and possesses a unique stationary state,
\[
 \pi_{\rm ph}(u,r)=\frac1{9\cdot468}.
\]
In each phase fiber, its return of nine steps is
\[
 \mathcal K_{\rm ph}^{\,9}\big|_r=E_+E_\times,
 \qquad
 (E_+E_\times)(u,v)=\frac1{468}.
\]
If the initial phase is uniformly distributed, the shadow of a step over
the emissions is
\[
 \overline K_{\rm ph}=\frac13(I+E_++E_\times).
\]
\end{theorem}

\begin{proof}
Each \(P_{\tau(r)}\) is doubly stochastic and the shift
\(r\mapsto r+1\) is a permutation; their skew product conserves both sums.
The three operators are idempotent and commute. Since each appears three
times in the nonadic word,
\[
 \prod_{r\in C_9}P_{\tau(r)}
 =E_+^3E_\times^3I^3=E_+E_\times.
\]
The first expectation forgets the additive signature and the second the multiplicative signature,
so their composition assigns to each of the \(18\cdot26=468\)
destinations the probability \(1/468\). Hence, in nine steps any two states of the same phase communicate; the
remaining steps communicate the phases, and the kernel is irreducible. Every return
must have length divisible by nine, whereas the diagonal of \(\mathcal K_{\rm ph}^9\) is positive;
the period is nine. Double stochasticity and irreducibility yield the unique
uniform state. Finally, averaging the nine phases counts three
copies of each channel and produces the final formula.
\end{proof}

Forgetting the phase is not a strong Markov quotient: from the same emission, two distinct phases apply \(I\), \(E_+\), or \(E_\times\), which have distinct rows. The operator \(\overline K_{\rm ph}\) is an averaged shadow in the stationary phase, not the chronological successor of a phaseless emission. This distinction preserves both the dynamics and the meaning of the measure.

\begin{definition}[Averaged kernel induced by transferrence]
\label{def:nucleo-tritico-refresco}
We call the phase shadow the averaged kernel of the catalog
\[
 \boxed{
 K_{\rm cat}=\frac13(I+E_++E_\times).
 }
\]
The weights are induced by the TPK calendar, which contains three occurrences
of each sign. The renewal operators are induced relative to the
uniform normalization of all compatible continuations of the active
fiber. Neither statement identifies \(K_{\rm cat}\) with
\(F_{\rm obs}\).
\end{definition}

\begin{proposition}[Convex family of stochastic kernels with a common uniform measure]
\label{prop:familia-refrescos-abc}
For \(a,b,c\geq0\), \(a+b+c=1\), the family
\[
 K_{a,b,c}=aI+bE_++cE_\times
\]
consists of symmetric, doubly stochastic kernels. If \(b,c>0\),
they are irreducible and, because they have positive diagonal, aperiodic. Their spectrum is
\[
 \operatorname{Spec}(K_{a,b,c})=
 \{1^{[1]},(a+c)^{[17]},(a+b)^{[25]},a^{[425]}\}.
\]
In particular, there exists a continuum of dynamics with the same factorization and the
same uniform measure. The choice
\(K_{\rm cat}=K_{1/3,1/3,1/3}\) is not characterized by the census alone. It is
selected jointly by the TPK phase word and the normalized transfer
of the \cref{def:transferencia-tpk-fase}.
\end{proposition}

\begin{proof}
The three summands are symmetric stochastic kernels; by convexity, the
same is true of \(K_{a,b,c}\), and symmetry converts row stochasticity
into column stochasticity. If \(b,c>0\), two steps make it possible to renew
both coordinates successively and connect any pair of states. Moreover,
\[
 K_{a,b,c}((s,e),(s,e))=a+\frac b{18}+\frac c{26}>0,
\]
so that the kernel is aperiodic. In the orthogonal decomposition
\[
 \mathbf1\oplus\ell_0^2(\mathcal S_+)
 \oplus\ell_0^2(\mathcal S_\times)
 \oplus\bigl(\ell_0^2(\mathcal S_+)\otimes
              \ell_0^2(\mathcal S_\times)\bigr),
\]
the pairs \((E_+,E_\times)\) act, respectively, as
\((1,1),(0,1),(1,0),(0,0)\). The four eigenvalues and the
multiplicities \(1,17,25,425\) are thereby obtained. For \(b,c>0\), irreducibility makes
the uniform stationary measure unique, which exhibits non-uniqueness even under
that requirement.
\end{proof}

\begin{theorem}[Exact spectrum of the balanced refreshment]
\label{thm:dinamica-cociente-468}
The kernel \(K_{\rm cat}\) is symmetric, doubly stochastic, irreducible and aperiodic. Its spectrum, with multiplicities, is
\[
 \boxed{
 \operatorname{Spec}(K_{\rm cat})
 =\left\{1^{[1]},
 \left(\frac23\right)^{[42]},
 \left(\frac13\right)^{[425]}\right\}.
 }
\]
Consequently, its spectral gap is (1/3) and its only stationary state is
\[
 \boxed{
 \tau_{468}(f)=\frac1{468}\sum_{U\in\mathcal U_6}f(U).
 }
\]
\end{theorem}

\begin{proof}
The operators \(E_+\) and \(E_\times\) are commuting orthogonal
averaging projectors. The rows and columns of each sum to one; the same
is true of their average. From \((s,e)\), \((s',e')\) is reached in two
renewals, and the contribution of \(I\) provides a self-edge. This
proves irreducibility and aperiodicity.

The space decomposes orthogonally into four sectors:
\[
 \mathbf1,
 \qquad
 \ell_0^2(\mathcal S_+),
 \qquad
 \ell_0^2(\mathcal S_\times),
 \qquad
 \ell_0^2(\mathcal S_+)\otimes
 \ell_0^2(\mathcal S_\times).
\]
Their dimensions are \(1,17,25,17\cdot25\). In the constant sector
the three summands act as the identity; in each one-coordinate sector
they act as (I+I+0); in the interaction sector, as (I+0+0).
Dividing by three gives the announced eigenvalues and multiplicities. A finite
irreducible and aperiodic kernel has a unique stationary state; since it is
doubly stochastic, that state is uniform.
\end{proof}

The theorem is exact relative to the factorization
\(\mathcal U_6\simeq\mathcal S_+\times\mathcal S_\times\), the TPK
calendar, and the uniform count of the compatible continuations of each fiber. The
family \(K_{a,b,c}\) proves that the dynamics do not follow from the census
alone. The \cref{prop:retorno-fobs-identidad} distinguishes the transfer of
all paths from the chronology of one of them.

\section{Compensated elevation constructed from the quotient}

Let us write \(m(U)=|F^{-1}(U)|\). From the kernel already chosen in the quotient
we define the kernel of presentations by
\[
 \boxed{
 K_{\rm seed}(x,y)
 =\frac{K_{\rm cat}(F(x),F(y))}{m(F(y))}
 }
\]
This kernel is stochastic. Summing over a target fiber yields
\[
 \sum_{F(y)=V}K_{\rm seed}(x,y)
 =K_{\rm cat}(F(x),V),
\]
independently of the representative \(x\). This identity shows that the
constructed lift projects exactly onto \(K_{\rm cat}\). It does not demonstrate that
a preexisting microscopic dynamics of the integral holonomic state of TPK descends
to the catalog.

\begin{theorem}[Reversible elevation constructed from the quotient]
\label{thm:elevacion-reversible-cociente}
The probability
\[
 \widetilde\tau(x)
 =\frac1{468\,m(F(x))}
\]
is estacionaria and reversible for \(K_{\rm seed}\), and
\(F_\#\widetilde\tau=\tau_{468}\). If
\[
 J_F\delta_U
 =m(U)^{-1/2}\sum_{F(x)=U}\delta_x,
 \qquad
 D=\operatorname{diag}(m(U)),
\]
then
\[
 K_{\rm seed}J_F
 =J_F\bigl(D^{1/2}K_{\rm cat}D^{-1/2}\bigr).
\]
In particular, Hilbert-space intertwining without conjugation by \(D\)
occurs if and only if \(K_{\rm cat}\) commutes with the fiber
multiplicities.
\end{theorem}

\begin{proof}
Normalization follows because each of the 468 fibres receives total mass
(1/468). Using the symmetry of \(K_{\rm cat}\), for
\(F(x)=U\) and \(F(y)=V\) we have
\[
 \widetilde\tau(x)K_{\rm seed}(x,y)
 =\frac{K_{\rm cat}(U,V)}{468m(U)m(V)}
 =\widetilde\tau(y)K_{\rm seed}(y,x).
\]
This proves detailed balance, stationarity, and the identity of the
image measure.
Applying both sides of the last identity to a basis vector
\(\delta_V\) gives, on the fiber \(U\), the coefficient
\(K_{\rm cat}(U,V)/\sqrt{m(V)}\); that is exactly the coefficient of the
conjugate operator shown.
\end{proof}

\begin{theorem}[Phase Lift Faithful to Quiescence]
\label{thm:elevacion-fase-semillas}
Let \(X=\Omega_{\rm seed}\), \(m(u)=|F^{-1}(u)|\), and let us define
\[
\begin{aligned}
 L_0(x,y)&=\mathbf1_{\{x=y\}},\\
 L_+(x,y)&=\frac{E_+(F(x),F(y))}{m(F(y))},\\
 L_\times(x,y)&=\frac{E_\times(F(x),F(y))}{m(F(y))}.
\end{aligned}
\]
On \(X\times C_9\), let
\[
 \widehat{\mathcal K}_{\rm ph}\bigl((x,r),(y,r+1)\bigr)
 =L_{\tau(r)}(x,y).
\]
Then \((F,\mathrm{id}_{C_9})\) satisfies the property of
\emph{strong aggregability} (\emph{strong lumpability}) from
\(\widehat{\mathcal K}_{\rm ph}\) to \(\mathcal K_{\rm ph}\). The microscopic
kernel is irreducible, has period nine, and its unique stationary state
is
\[
 \widetilde\pi_{\rm ph}(x,r)
 =\frac1{9\cdot468\,m(F(x))}.
\]
Its nonadic return satisfies
\[
 \widehat{\mathcal K}_{\rm ph}^{\,9}
 \bigl((x,r),(y,r)\bigr)
 =\frac1{468\,m(F(y))}.
\]
\end{theorem}

\begin{proof}
The sum of \(L_+\), respectively \(L_\times\), over a target fiber
is \(E_+\), respectively \(E_\times\), and \(L_0\) projects to the identity.
This proves aggregability in each phase. The factors \(m(F(y))\) cancel
under composition: first the active signature is selected uniformly and then a
representative of the obtained emission. Thus, \(L_+\) and \(L_\times\) are
commuting projectors and
\[
 L_+L_\times(x,y)=\frac1{468\,m(F(y))}.
\]
The positivity of this entry connects all representatives in nine
steps. The phase forces the return lengths to be multiples of nine,
and the diagonal entry of the return is positive. The stationary formula is
verified by first summing the \(m(u)\) representatives of each emission; each
of the \(9\cdot468\) phase--emission classes receives the same mass.
\end{proof}

This lift improves the reading of the averaged kernel \(K_{\rm seed}\): in the neutral phase it preserves the microscopic representative, instead of silently remixing its fiber. Upon averaging and forgetting the phase, both constructions have the same shadow \(K_{\rm cat}\), but not the same hidden dynamics. The distinction is necessary in order to respect that the zero of TRIT is memory of quiescence.

\section{Minimal congruences stable under advance, half-turn and 4th of a turn: \(28\) and \(162\) memory classes}

The witness \(e_0=(5,5,5,5,5,5)\) shows that the product signature is too
coarse to transport a route. The exact missing minimal memory can be
computed. Let
\[
 \mathcal X_\times=X_9\times\mathscr D_{\rm card}
\]
the space of \(324\) product seeds. Let \(P\) denote the advance by one
cell in the stored orientation, \(H\) the half-turn of that
orientation, and \(Q\) the quarter-turn
\[
 N\longmapsto E\longmapsto S\longmapsto O\longmapsto N.
\]

\begin{theorem}[Forward and spin congruences]
\label{thm:congruencias-avance-giro}
The minimal congruence that refines the twenty-six product signatures and is stable
under \(P\) and \(H\) has \(28\) classes: twenty-seven have size eight, and one
has size \(108\). The only new splitting is
\[
 (5,5,5,5,5,5)\longrightarrow
 \mathcal L_0\sqcup\mathcal L_1\sqcup\mathcal L_2,
 \qquad |\mathcal L_j|=8.
\]
Its orbits under \(P,H\) have sizes
\[
 9,\quad9,\quad9,\quad1.
\]
By combining this refinement with the eighteen additive signatures one obtains
\[
 18\cdot28=504=468+36.
\]

The minimal congruence stable under \(P\) and \(Q\) has \(162\) classes of two
seeds. Each class is an orbit of the involution
\[
 J(i,j,d)=\bigl(7-i,\,7-j,\,\overline d\bigr)\pmod9.
\]
The operators induced by \(P,Q\) act transitively on the
\(162\) classes.
\end{theorem}

\begin{proof}
One begins with the partition by the twenty-six signatures and iteratively refines a class whenever two of its elements have images in distinct classes under any of the generators. The process terminates because \(\mathcal X_\times\) is finite. For \((P,H)\), the exhaustive census produces the histogram \(27\times8+1\times108=324\), the four indicated orbits, and three classes over the exceptional fiber of size twenty-four. For \((P,Q)\), it produces \(162\times2=324\). Direct verification proves that \(J\) preserves every signature, commutes with both generators, and pairs exactly the two elements of each class; the graph of classes generated by \(P,Q\) is connected.
\end{proof}

The excess \(36\) and the object \(R_{36}\) have distinct types. The boundary
at depth thirty-six is
\[
 R_{36}=
 \begin{pmatrix}
 222220\\021222\\102011
 \end{pmatrix},
\]
that is, three ordered blocks of six trits, one per channel. It is not a set of thirty-six states. The term \(504-468=36\), for its part, counts additional sheet memory: the half-turn that makes product transport functional adds exactly two new sheets to a single signature and, when repeated over the eighteen additive signatures, produces \(18(3-1)=36\). The coincidence of the integer (36) expresses two gradings of transport, boundary depth and sheet multiplicity, but does not authorize identification of their objects.

The quarter-turn \(Q\) was not hidden in \(F_{\rm obs}\) either. The
published chronology contains programmed advances and half-turns; incorporating
\(Q\) is an explicit enlargement of the route alphabet. Once it is incorporated,
a lazy symmetric random walk on \(P^{\pm1},Q^{\pm1}\) is irreducible and
aperiodic on the \(162\) classes, and its uniform measure is unique. This is the
exact mathematical form of “counting the turns”: the turn is a transport
coordinate and not a subsequent decimal correction.

\section{Mode and Orientation Cocycles}
\label{sec:cociclos-modo-orientacion}

The same rule makes it possible to define the operational counters without introducing
Dimensional Mechanics. Add to the state the last active mode
\(m\in\{\Sigma,\Pi\}\). A signature \(+1\) fixes \(m=\Sigma\), a signature
\(-1\) fixes \(m=\Pi\), and a signature \(0\) preserves the previous mode. For an
arrow \(a\), define the signed cocycle
\[
 \omega_{\rm sw}(a)=
 \begin{cases}
  +1,&\Pi\longrightarrow\Sigma,\\
  -1,&\Sigma\longrightarrow\Pi,\\
  0,&\text{the mode is preserved}.
 \end{cases}
\]
Its positive variation is \(s(a)=|\omega_{\rm sw}(a)|\). We further define
\(g(a)=1\) when the ordered pair of cardinal orientations changes and
\(g(a)=0\) when it is preserved. Keeping the mode and the terminal
orientation fixed at the endpoints, the sums
\[
 \Omega_{\rm sw}(\gamma)=\sum_{a\subset\gamma}\omega_{\rm sw}(a),
 \qquad
 S(\gamma)=\sum_{a\subset\gamma}s(a),
 \qquad
 G(\gamma)=\sum_{a\subset\gamma}g(a)
\]
are additive under concatenation.

In a nonadic closed loop, the sequence of modes is
\[
 (\Sigma,\Pi,\Pi,\Sigma,\Pi,\Pi,\Sigma,\Pi,\Pi).
\]
Therefore,
\[
 \Omega_{\rm sw}(C_9)=0,\qquad S(C_9)=6:
\]
there are three changes \(\Pi\to\Sigma\), three changes
\(\Sigma\to\Pi\), and three retentions. In the cycle of \(108\) steps,
\[
 \#(+1)=\#(-1)=\#(0)=36,\qquad
 \Omega_{\rm sw}=0,\qquad S=72,\qquad G=4.
\]
The final number counts the four simultaneous inversions programmed in
\(27,54,81,108\). These cocycles proceed directly from the TPK calendar\@.

\subsection{Quotient of modes and areal-volumetric correspondence}
\label{sec:cociente-modos-fav}

The calendar sign, mode label, and geometric sheet are three typed
objects. In particular, the phase sign is not identified with the
angular observable on digits. For the geometric correspondence, we fix the
typed HMT reader
\[
 \mathfrak g_{\rm av}(\Sigma)=\mathscr H_{\rm ar},
 \qquad
 \mathfrak g_{\rm av}(\Pi)=\mathscr H_{\rm vol}.
\]
The first assignment expresses the additive or boundary reading; the second, the
multiplicative or interior reading. The map \(\mathfrak g_{\rm av}\) is
part of the sheet typing. The following result proves that, once
that typing is fixed, TPK dynamics forces the incidence between them.

\begin{theorem}[Necessary descent of leaf incidence]
\label{thm:descenso-necesario-fav}
Let \(m_r\) be the mode obtained after applying phase \(r\) of the
primitive block \((+1,-1,0)\), with the rule
\[
 +1:\ m\mapsto\Sigma,\qquad
 -1:\ m\mapsto\Pi,\qquad
 0:\ m\mapsto m.
\]
The cyclic orbit of modes is \((\Sigma,\Pi,\Pi)\). If
\[
 N_3(i,j)
 :=
 \#\{r\in\mathbb Z/3\mathbb Z:(m_r,m_{r+1})=(i,j)\},
\]
then, in the order \((\Sigma,\Pi)\),
\[
 \boxed{
 N_3=
 \begin{pmatrix}0&1\\1&1\end{pmatrix}.}
\]
By repetition of the calendar,
\[
 N_9=3N_3,\qquad N_{27}=9N_3,\qquad N_{108}=36N_3.
\]
After applying \(\mathfrak g_{\rm av}\), the primitive incidence
matrix is, in the order
\((\mathscr H_{\rm ar},\mathscr H_{\rm vol})\),
\[
 \boxed{
 F_{\rm av}=
 \begin{pmatrix}0&1\\1&1\end{pmatrix}.}
\]
Thus, the two cross transitions, the unique volumetric self-retention, the
absence of areal self-retention and the unit primitive multiplicities
are consequences of the TPK mode quotient; they are not four added
postulates.
\end{theorem}

\begin{proof}
The three consecutive pairs, with cyclic closure, are
\[
 \Sigma\longrightarrow\Pi,\qquad
 \Pi\longrightarrow\Pi,\qquad
 \Pi\longrightarrow\Sigma.
\]
Each appears once, and the pair \(\Sigma\to\Sigma\) does not appear. This determines the four entries of \(N_3\). The calendars of lengths \(9,27,108\) respectively contain \(3,9,36\) copies of the primitive block, so their incidence matrices are the indicated multiples. The greatest common divisor of the nonzero entries of each long matrix removes only that global repetition and returns \(N_3\). Finally, \(\mathfrak g_{\rm av}\) transports \(\Sigma,\Pi\) to the areal and volumetric sheets without altering the incidence.
\end{proof}

The real hyperbolic realization that uses this matrix and proves its
Lorentz invariance is formulated in
\cref{sec:tres-generadores-concretos-trit}; there \(F_{\rm av}\) is used
as an antecedent, without repeating its derivation.

This theorem must be distinguished from a stronger Markov assertion.
Forgetting the phase is not a strong Markov quotient of
\(\mathcal K_{\rm ph}\): the visible mode does not determine which of the three
operators \(I,E_+,E_\times\) acts at the next elementary iteration. What descends
without an additional hypothesis is the edge multigraph of a TPK\@ block.

There is, however, a memory-one completion uniquely determined by that descent. Let \(X_{\rm av}\) be the smallest subshift of finite type over \(\{\mathscr H_{\rm ar},\mathscr H_{\rm vol}\}\) that contains all the words obtained by forgetting the phase and retains only consecutive pairs. Every one-step subshift with that property must admit exactly the three observed pairs; the smallest such subshift excludes the sole unobserved pair. Its adjacency matrix is therefore \(F_{\rm av}\).

\begin{corollary}[Parry Measure of the One-Memory Wrapper]
\label{cor:parry-envolvente-fav}
The subshift \(X_{\rm av}\) is primitive, has topological entropy
\(\log\varphi\), and possesses a unique measure of maximal entropy. Its stationary
vertex weights satisfy
\[
 \boxed{
 \frac{\mu_{\rm ar}}{\mu_{\rm vol}}=\varphi^{-2}.}
\]
\end{corollary}

\begin{proof}
The Perron root of \(F_{\rm av}\) is \(\varphi\), and a positive Perron
vector is \((1,\varphi)\). Since the matrix is symmetric, the Parry vertex
weights are proportional to the product of the left and right vectors,
that is, to \((1,\varphi^2)\).
\end{proof}

The Parry measure belongs to the memory-one path envelope. It is not
the image of \(\tau_{468}\), nor the temporal frequency of the periodic orbit
of modes. The latter is
\[
 \nu_{\rm cyc}(\mathscr H_{\rm ar})=\frac13,\qquad
 \nu_{\rm cyc}(\mathscr H_{\rm vol})=\frac23,
\]
whereas \(\tau_{468}\) is uniform over emissions, and the Parry measure
weights renewal histories of arbitrary length. The three measures
answer distinct questions. The irrational ratio
\(\varphi^{-2}\) is thus derived canonically from the path envelope,
without attributing it to a quotient of finite cardinalities.

\subsubsection{From chronological counting to the return marked by \texorpdfstring{\(\pi\)}{π}}
\label{sec:retorno-marcado-pi-phi2}

The two formulas
\[
 \left(\frac13,\frac23\right)
 \qquad\text{and}\qquad
 \frac{\mu_{\rm ar}}{\mu_{\rm vol}}=\varphi^{-2}
\]
are not approximations of one another. The first counts the three instants of the
primitive calendar \((\Sigma,\Pi,\Pi)\); the second weights all the
admissible histories of the memory-one envelope. Hence the chronological
frequency is rational, and the Parry ratio is irrational. In what follows,
we write
\[
 \frac{\pi}{\varphi^2}
 =\pi\varphi^{-2};
\]
this is not division by \(\varphi^{-2}\).

The appearance of \(\varphi^{-2}\) can be seen directly in the linear
dynamics. If \(F_n\) denotes the \(n\)-th Fibonacci number,
\[
 F_{\rm av}^{\,n}
 =
 \begin{pmatrix}
  F_{n-1}&F_n\\
  F_n&F_{n+1}
 \end{pmatrix}
 \qquad(n\geq1).
\]
The eigenvalues of \(F_{\rm av}\) are \(\varphi\) and \(-\varphi^{-1}\). The contracted branch reverses orientation in one iteration. Upon passing to the quadratic memory, namely \(\operatorname{Sym}^2(\mathbb R^2)\), that reversal is recorded before the sign disappears. In the basis \((x^2,2xy,y^2)\), with the images of the basis vectors represented by rows, the induced operator is
\[
 \operatorname{Sym}^2(F_{\rm av})=
 \begin{pmatrix}
  0&0&1\\
  0&1&2\\
  1&1&1
 \end{pmatrix},
 \qquad
 \chi(t)=(t+1)(t^2-3t+1).
\]
Its three eigenvalues are
\[
 \varphi^2,\qquad -1,\qquad \varphi^{-2}.
\]
Thus, \(\varphi^{-2}\) is the positive stable mode of second-order
memory. This step explains why the golden ratio appears here as a return
law and not as an added statistical mean.

The closure \(\pi\) belongs to another level of the same HMT state: it is the
coinductive output of the regional reader already constructed, not a new eigenvalue of
\(F_{\rm av}\). Let
\[
 P_{\rm ar}=
 \begin{pmatrix}1&0\\0&0\end{pmatrix},
 \qquad
 P_{\rm vol}=
 \begin{pmatrix}0&0\\0&1\end{pmatrix}.
\]
Among the positive operators that preserve both sheets, there is exactly
one whose volumetric normalization is one and whose areal mark is the generated
closure \(\pi\):
\[
 D_\pi={\pi}P_{\rm ar}+P_{\rm vol}.
\]
Uniqueness is relative to these three explicit clauses: positivity,
diagonality with respect to the leaves, and the two normalizations. With
\(e_{\rm ar}=(1,0)^T\) and \(e_{\rm vol}=(0,1)^T\), the marked return of
length \(n\) is
\[
 \Theta_n
 :=
 \frac{\langle e_{\rm ar},
 D_\pi F_{\rm av}^{\,n}e_{\rm ar}\rangle}
 {\langle e_{\rm vol},
 F_{\rm av}^{\,n}e_{\rm vol}\rangle}
 =
 \pi\frac{F_{n-1}}{F_{n+1}}.
\]
Therefore,
\[
 \boxed{\displaystyle
 \lim_{n\to\infty}\Theta_n
 =\frac{\pi}{\varphi^2}.}
\]
The mark \(\pi\) is applied only once, at the boundary return;
it is not reinjected at every transition. Finite combinatorics produces the
stable mode \(\varphi^{-2}\), the coinductive reader produces the closure
\(\pi\), and \(D_\pi\) composes both data without reversing their
causal order.

The native length \(108\) provides an exact witness to that convergence:
\[
 (F_{\rm av}^{108})_{\rm ar,ar}
 =F_{107}=10284720757613717413913,
\]
\[
 (F_{\rm av}^{108})_{\rm vol,vol}
 =F_{109}=26925748508234281076009.
\]
Consequently,
\[
 \left|
 \frac{F_{107}}{F_{109}}-\varphi^{-2}
 \right|
 =
 6.1684979916614407936\ldots\times10^{-46},
\]
and, after the only closing marker,
\[
 \left|
 \pi\frac{F_{107}}{F_{109}}
 -\frac{\pi}{\varphi^2}
 \right|
 =
 1.9378907974286976068\ldots\times10^{-45}.
\]
These errors measure the finite approximation of the return; they are not adjustments of
\(\pi\) or of \(\varphi\).

\paragraph{Cylindrical reader of the ratio.}
Let \(I_n^\pi=[\underline\pi_n,\overline\pi_n]\) and
\(I_n^\varphi=[\underline\varphi_n,\overline\varphi_n]\) be the positive,
nested cylinders produced by the two HMT readers\@. The interval
operation
\[
 \mathcal Q(I,J)
 :=
 \left[
 \frac{\inf I}{(\sup J)^2},
 \frac{\sup I}{(\inf J)^2}
 \right]
\]
is the smallest interval containing \(x/y^2\) for every
\((x,y)\in I\times J\). It preserves nesting: if the input cylinders
are refined, then \(\mathcal Q(I_n^\pi,I_n^\varphi)\) is also refined. Since their
diameters tend to zero and the limit of \(I_n^\varphi\) is positive,
\[
 \bigcap_n\mathcal Q(I_n^\pi,I_n^\varphi)
 =
 \left\{\frac{\pi}{\varphi^2}\right\}.
\]
With the twelve already published base-thousand blocks, the quotient interval forces
thirty-five decimal digits:
\[
 \frac{\pi}{\varphi^2}
 =
 1.19998161486432666111577775331680692\ldots.
\]
This construction preserves the provenance of each digit: the numerator
comes from the closure cylinder and the denominator from the autoscale.

\paragraph{Algebraic separation between areal reading, volumetric reading and their respective codomains}
The complete ratio cannot arise from the rational algebra of
\(F_{\rm av}\) alone. Indeed, its rational spectral functions belong to
\(\mathbb Q(\sqrt5)\); \(\varphi^{-2}\) is found there, but not the
transcendental number \(\pi\). Consequently, any derivation purporting to
obtain \(\pi/\varphi^2\) solely from the matrix \(2\times2\) would be
ill typed. The transcendental datum must derive from the coinductive closure
reader, or from an equivalent boundary cocycle proved
independently. This negative control fixes the division of labor between
the finite correspondence and the number reader.

\subsubsection{Exchange involution of the quotient}
\label{sec:espejo-pi-phi-electron-hmt}

The correspondence admits two exact orientations of the same law. If
\[
 \mathscr R(x,y)=\frac{x}{y^2},
 \qquad
 \mathsf J(x,y)=(y,x),
\]
then \(\mathsf J^2=I\) and
\[
 \mathscr R(\pi,\varphi)=\frac{\pi}{\varphi^2},
 \qquad
 (\mathscr R\circ\mathsf J)(\pi,\varphi)
 =\frac{\varphi}{\pi^2}.
\]
The identities are also verified.
\[
 \mathscr R(\pi,\varphi)\mathscr R(\varphi,\pi)
 =\frac1{\pi\varphi},
 \qquad
 \frac{\mathscr R(\pi,\varphi)}
 {\mathscr R(\varphi,\pi)}
 =\left(\frac{\pi}{\varphi}\right)^3.
\]
Writing
\(\mathcal B_{\rm av}^{(\pi)}=\mathscr R(\pi,\varphi)\) also yields
the exact reparametrization
\[
 \frac{\varphi}{\pi^2}
 =
 \frac1{\varphi^3
 \bigl(\mathcal B_{\rm av}^{(\pi)}\bigr)^2},
 \qquad
 e^{\varphi/\pi^2}
 =
 \exp\!\left[
 \frac1{\varphi^3
 \bigl(\mathcal B_{\rm av}^{(\pi)}\bigr)^2}
 \right].
\]
Thus, the change of orientation between closure and growth transports the
boundary ratio to its exchanged reading:
\[
 \frac{\pi}{\varphi^2}
 \xleftrightarrow{\ \mathsf J\ }
 \frac{\varphi}{\pi^2}.
\]
The first orientation reads the marked areal closure over quadratic
volumetric memory; the second reads interior growth over
quadratic closure. This correspondence is a theorem of the HMT reader. Dimensional Mechanics subsequently applies the second orientation within the
electronic construction, without converting the electron into an antecedent of
\(\pi\), \(\varphi\), or of the correspondence.

\subsubsection{Ambivalence Principle of the areal and volumetric readers}
\label{sec:ambivalencia-grav-inercial-pi-phi2}

The Ambivalence Principle asserts that the areal reading and the
volumetric reading are two coordinated, but not interchangeable, functionals of one
and the same state with memory. Their coincidence on the diagonal is
\[
 \mathcal L_{\rm ar}^{\rm diag}(x)
 =\mathcal L_{\rm vol}^{\rm diag}(x).
\]
The diagonal equality does not erase the ambivalence outside of it.

On the areal--volumetric boundary, both readers factor through the same
positive memory amplitude \(\mathcal M(x)\):
\[
 \mathcal L_{\rm vol}(x)=\mu_{\rm vol}\mathcal M(x),
 \qquad
 \mathcal L_{\rm ar}(x)=
 \pi\mu_{\rm ar}\mathcal M(x).
\]
Then the Parry theorem and the return mark give
\[
 \boxed{\displaystyle
 \frac{\mathcal L_{\rm ar}(x)}
 {\mathcal L_{\rm vol}(x)}
 =\frac{\pi}{\varphi^2}.}
\]
With volumetric normalization, the relative difference of readers is
\[
 \frac{\mathcal L_{\rm ar}-\mathcal L_{\rm vol}}
 {\mathcal L_{\rm vol}}
 =
 \frac{\pi}{\varphi^2}-1
 =
 0.1999816148643266611\ldots.
\]
The physical reading of this correspondence is undertaken later, in the
Cartan--Holst chapter, after constructing holonomy, co-tetrad, connection, torsion,
and curvature. Here the result is entirely HMT: an exact ratio between
two readers of the same sheet memory.

They should be distinguished from the enriched extractor subsequently used in
Dimensional Mechanics. There a complete history preserves the action count
\(n_A\), the twelve oriented events \(e_0,\ldots,e_{11}\) and the torsional
cochain at step level; these yield
\[
 s_C=\sum_{i=0}^{5}(e_i+e_{i+6}),
 \qquad
 W_\pm=6n_A\pm s_C.
\]
The channels already generated by HMT then satisfy the exact identity
\[
 \exp\!\left(n_AA_{\rm rad}+\frac{s_C}{6}C^\ast_{\rm rad}\right)
 =\left(q_+^{-W_+}q_-^{-W_-}\right)^{1/12}.
\]
Therefore, the route--angular-character passage is not absent: it is realized on
the enriched dodecaphase record. What this phase calculation proves by itself
are \(\Omega_{\rm sw},S,G\); it neither replaces them nor permits reconstruction of the
complete record after it has been projected.

The uniform measure on the \(104\,976\) presentations descends to atoms
of masses \(1/729\), \(1/243\) and \(1/54\), but not to \(1/468\). In the
fivefold microfiber of \(\pi\), the two probabilities are
\[
 \tau_{468}(\mathscr R_\pi)=\frac5{468},
 \qquad
 F_\#\mu_{\rm seed}(\mathscr R_\pi)=\frac7{729}.
\]
Both the averaged lift and the phase lift realize the first: they do not
push forward the uniform measure on seeds, but assign the same total
mass to each observable class. The phase lift also preserves
microscopic stillness at neutral steps.

To relate phase elevation to the enriched histories of
depth six, a kernel \(\mathcal K_6\) on
\(\mathcal H_6^{\rm cat}\times C_9\) is required. It descends to
\(\widehat{\mathcal K}_{\rm ph}\) through
\(\rho_6\times\mathrm{id}_{C_9}\) if and only if it satisfies, at each phase, the
strong lumpability condition
\[
 \sum_{\rho_6(g)=y}\mathcal K_6((h,r),(g,r+1))
 =L_{\tau(r)}(\rho_6(h),y),
\]
independently of \(h\) within each fiber of \(\rho_6\). In that case,
composition with \(F\) projects to \(\mathcal K_{\rm ph}\). Constructing
a family \((\mathcal K_m)_{m\geq6}\) that satisfies this identity and
commutes with all truncations is the remaining step toward a
transfer of the inverse limit.

\begingroup
\setlength{\parskip}{1pt}
\setlength{\abovedisplayskip}{3pt}\setlength{\belowdisplayskip}{3pt}
\setlength{\abovedisplayshortskip}{2pt}\setlength{\belowdisplayshortskip}{2pt}
\section{Scale cocycle and cylindrical conservation}
\label{sec:cociclo-parry-escala}

Let \(r=(1,\varphi)\) be a Perron--Frobenius vector of
\(F_{\rm av}\). The Parry transition matrix is
\[
P_{ij}=\frac{(F_{\rm av})_{ij}r_j}{\varphi r_i}
=
\begin{pmatrix}
0&1\\
\varphi^{-2}&\varphi^{-1}
\end{pmatrix}.
\]
For an admissible route \(\gamma:i\to j\) of length \(n\),
\[
\mathbb P(\gamma\mid i)=\frac{r_j}{\varphi^nr_i},
\qquad
\boxed{\chi_P(\gamma)=\varphi^n\frac{r_i}{r_j}.}
\]
The telescope proves
\[
\chi_P(\gamma_2\circ\gamma_1)
=\chi_P(\gamma_2)\chi_P(\gamma_1).
\]
On a closed route, \(\chi_P=\varphi^n\); its square and its inverse
generate the combinations
\[
C_n=\frac{\chi_P^2+\chi_P^{-2}}2,\qquad
S_n=\frac{\chi_P^2-\chi_P^{-2}}2,
\]
which satisfy
\[
x_{n+1}=3x_n-x_{n-1}.
\]

Since \(F_{\rm av}\) is symmetric, let
\[
\ell=\frac{r}{1+\varphi^2},\qquad \ell^{\mathsf T}r=1.
\]
The measure of a cylinder is
\[
\mu[x_0\cdots x_n]=\ell_{x_0}r_{x_n}\varphi^{-n}.
\]
Therefore,
\[
\boxed{
V_n(x)
=\frac{\mu[x_0]}{\mu[x_0\cdots x_n]}
=\varphi^n\frac{r_{x_0}}{r_{x_n}}
=\chi_P(x_0\cdots x_n).}
\]
For a typed dimension scale \(D>0\),
\[
a_n^{(D)}(x)=V_n(x)^{1/D}
\]
satisfies the cylindrical conservation law
\[
\boxed{
\bigl(a_n^{(D)}(x)\bigr)^D
\mu[x_0\cdots x_n]=\mu[x_0].}
\]

In spatial rank three and on closed returns, we write
\(a_n:=a_n^{(3)}\) and obtain
\[
\boxed{
\frac{a_9}{a_0}=\varphi^3,\qquad
\frac{a_{27}}{a_0}=\varphi^9,\qquad
\frac{a_{108}}{a_0}=\varphi^{36}.}
\]
If \(A_k:=a_{9k}\) for \(k\geq1\), then, for \(k\geq2\),
\[
A_{k+1}-2A_k+A_{k-1}
=A_{k-1}(\varphi^3-1)^2>0.
\]
The inequality proves convexity in return depth; a physical
acceleration further requires a duration chart.

The primary computation of
\(\operatorname{Sym}^2(F_{\rm av})\) in
\cref{sec:retorno-marcado-pi-phi2} shows that its spectrum is
\(\{\varphi^2,-1,\varphi^{-2}\}\). Here we conserve only the corollary
for the scale cocycle: on a closed route with \(V_n=\varphi^n\), the
stable mode obeys
\[
\boxed{M_n=\varphi^{-2n}=V_n^{-2}.}
\]
Identification of these three eigenspaces with the three TRIT sheets
requires an intertwiner; coincidence of cardinality does not replace it. The
corollary would be refuted if the closed route failed to satisfy
\(V_n=\varphi^n\) or if the stable mode did not have eigenvalue
\(\varphi^{-2}\).

\section{Dynamic prolongation theorem: global hypotheses and certified domain}

The following objects have been explicitly constructed and separated from the finite problem:
\[
 F,
 \qquad
 F_{\rm obs},
 \qquad
 \mathcal K_{\rm ph},
 \qquad
 \widehat{\mathcal K}_{\rm ph},
 \qquad
 K_{\rm cat},
 \qquad
 \tau_{468},
 \qquad
 F_{\rm av},
 \qquad
 \operatorname{Sym}^2(F_{\rm av}),
 \qquad
 D_\pi,
 \qquad
 \mathcal Q.
\]
For \(F_{\rm obs}\), deterministic return has been proved and an
explicit witness of failure to descend as a product has been exhibited. For \(\mathcal K_{\rm ph}\),
double stochasticity, irreducibility, period nine, uniform
nonadic return, and uniqueness of
\(\pi_{\rm ph}\) have been proved. For \(\widehat{\mathcal K}_{\rm ph}\), strong
lumpability, microscopic conservation of rest, and the compensated
stationary measure have been proved. \(K_{\rm cat}\) is no longer an average without
provenance: it is the shadow of one step after the uniform phase is averaged.
The TPK mode quotient induces \(F_{\rm av}\); its quadratic lift
separates the modes \(\varphi^2,-1,\varphi^{-2}\). The operator \(D_\pi\)
marks the boundary return exactly once with the generated closure, and the
interval reader \(\mathcal Q\) preserves the nesting of the cylinders of
\(\pi\) and \(\varphi\). These four objects construct
\(\pi/\varphi^2\) without conflating temporal frequency,
the uniform measure, and the Parry measure.

The force of this induction is relative and concrete. The TPK word forces the
weights \(1/3\). The APP interpretation of all compatible continuations,
together with uniform counting, forces \(E_+\) and \(E_\times\). The deterministic
trajectory \(F_{\rm obs}\) does not force those operators and, indeed, does not
produce them. This is not a defect of the result, but the mathematical difference
between evolution along a path and transfer over a space of paths.

\ifdefined\HMTInternalTraceability
\begin{criterion}[Enriched lift of reversion]
\label{pf84:SYN31-RHO-LIFT}
A subdomain \(D_\rho\subset X_{\rm TPK}^{\rm enr}\) and a map are sought
\[
 \Theta_\rho:D_\rho\longrightarrow X_{\rm TPK}^{\rm enr}
\]
that lifts the integer reversal \(27\mapsto72\) of
\cref{pf84:SYN30-RHO-DESC,pf84:SYN30-CROSS-WITNESS}. It must commute with the
residual, integer, and genealogical readers, quantify the quotient and the
carry, and preserve orientation, route, return, and, when the transport is
invertible, monodromy, region, and survival.
Only if \(D_\rho\) is total and invariant under \(\Theta_\rho\) may
\(\Theta_\rho^2=I\) be required and the map be called an involution.

Neither this construction nor a proof of involutivity yet exists. The
program is falsified if the map does not exist, if one of the reading squares
fails, if memory is lost without being quantified, or if
\(\Theta_\rho^2\ne I\) on a domain where involution is claimed.
\end{criterion}

\begin{criterion}[Enriched lift of the square]
\label{pf84:SYN31-SQ-LIFT}
A subdomain \(D_s\subset X_{\rm TPK}^{\rm enr}\) and a map are sought
\[
 \Theta_s:D_s\longrightarrow X_{\rm TPK}^{\rm enr}
\]
that lifts \(s(n)=n^2\), in particular \(27\mapsto729\), and commutes with the
residual, integer, and genealogical readers of
\cref{pf84:SYN30-SQ-DESC,pf84:SYN30-CROSS-WITNESS}. It must quantify
quotient and carry and preserve orientation, route, return, and, when
transport is invertible, monodromy, region, and survival. This map is not
involutive, and
\(\Theta_s^2=I\) is not imposed on it.

The proof remains absent. The nonexistence of \(\Theta_s\), the failure of
a reading square, or the uncontrolled loss of any of the
enriched coordinates would refute the program. The ablation
\(81\mapsto72\) upon changing \(q^\times(9,9)\) is an independent control:
it is a branch of neither \(\Theta_\rho\) nor \(\Theta_s\), and it proves neither
of the two maps.
\end{criterion}
\fi

The finite conclusion is complete in its domain: the uniform return, the
stationary measure, the lumpability at depth six, and the separation
between evolution of a route and transfer of the route space remain
proved by the constructed objects. In particular,
\(\widehat{\mathcal K}_{\rm ph}\) is not identified with the chronology of
\(F_{\rm obs}\), as \cref{prop:retorno-fobs-identidad} shows.
\par\endgroup
 \end{HMTScientificOwnerSubordinated}
}
 
\chapter{Correlative and inseparable construction of the continuum}
\chaptermark{Discrete structure of the continuum}
\label{chap:83-04}

The five constructions are generated correlatively on a single
APP--TRIT--TPK state. They are consubstantial operations of the continuum,
not an exhaustive enumeration of its properties, five independent branches,
or a retrospective selection determined by subsequent applications.

\begin{HMTScientificOwnerSubordinated}
% Corrección exclusivamente editorial del segundo epígrafe de 04b1. El
% fuente teoremática canónico permanece inmutable y se conserva íntegro; aquí se
% traduce la nomenclatura interna «levantamiento TRIT» a una formulación
% pública que nombra la operación y el resultado matemático.
\newcount\HMTContinuoCommonSection
\HMTContinuoCommonSection=0
\let\HMTContinuoCommonOriginalSection\section
\renewcommand{\section}[1]{%
  \advance\HMTContinuoCommonSection by 1
  \ifnum\HMTContinuoCommonSection=2
    \HMTContinuoCommonOriginalSection{Tritic extension of the incidence datum: balanced orientation and genealogical carry cocycle}%
  \else
    \HMTContinuoCommonOriginalSection{#1}%
  \fi
}
\chapter{Common state, prolongation tree, and genealogical measure}
\label{chap:continuo-estado-arbol-medida-comun}

\begin{thesisbox}[title={A single integral holonomic state before any internal discrimination}]
The object of this chapter is the complete finite state produced by
\(\APP\), oriented by TRIT and transported by \(\TPK\).  From it are
constructed the groupoid of reversible symmetries, the tree of all
admissible prolongations, the multisection of complete histories and its canonical
measure.  No specialized domain is introduced: the domain is the
entire set of states actually generated by the same mechanism
\(\APP\)--TRIT--\(\TPK\).  Nor is a sheet, a continuation or
a subsequent reading selected.
\end{thesisbox}

\section{Finite data produced by APP}
\label{sec:continuo-comun-app-finito}

For each depth \(k\geq0\), let \(\mathsf A_k\) denote the finite set
of cellular prefixes produced by APP, and let
\begin{equation}
\label{eq:continuo-comun-app-truncamiento}
 \rho^{\APP}_{k+1,k}:\mathsf A_{k+1}\longrightarrow\mathsf A_k
\end{equation}
the erasure of the last refinement.  This erasure conserves the already
constructed prefix.  The neutral cell of APP gives a prolongation
\begin{equation}
\label{eq:continuo-comun-app-degeneracion}
 \eta^{\APP}_k:\mathsf A_k\longrightarrow\mathsf A_{k+1},
 \qquad
 \rho^{\APP}_{k+1,k}\eta^{\APP}_k=I_{\mathsf A_k}.
\end{equation}
We write
\begin{equation}
\label{eq:continuo-comun-emision-neutra-app}
 \varepsilon^0_k(a):a\longrightarrow\eta^{\APP}_k(a)
\end{equation}
for the emission that realizes that prolongation. Thus, its
destination, \(\eta^{\APP}_k(a)\), and the operation that produces it,
\(\varepsilon^0_k(a)\), are kept separate.
\(\eta^{\APP}_k\) is not interpreted as an absence of history: it records a
length-one refinement with neutral emission and materially preserves the
preceding state.

Each \(a\in\mathsf A_k\) simultaneously bears the two sheets of APP,
\begin{equation}
\label{eq:continuo-comun-dos-hojas}
 \operatorname{Sh}_k(a)
 =\bigl(\Sigma_k(a),\Pi_k(a)\bigr)
 \in\mathsf H_k^{\Sigma}\times\mathsf H_k^{\Pi},
\end{equation}
and not a chosen element of the pair.  For an elementary emission
\(\epsilon:a\to a'\), APP further produces two typed Euclidean divisions,
\begin{equation}
\label{eq:continuo-comun-residuo-cociente-app}
 \delta^{\diamond}(\epsilon)
 =\operatorname{res}^{\diamond}(\epsilon)
  +9\operatorname{quo}^{\diamond}(\epsilon),
 \qquad
 \operatorname{res}^{\diamond}(\epsilon)\in\{0,\ldots,8\},
 \quad
 \operatorname{quo}^{\diamond}(\epsilon)\in\mathbb Z,
\end{equation}
where \(\diamond\in\{\Sigma,\Pi\}\).  The label \(\diamond\) merely indicates
which leaf is being read; both divisions remain in the same ledger.

The cellular source and target of \(\epsilon\) determine its oriented
boundary
\begin{equation}
\label{eq:continuo-comun-frontera-elemental}
 \partial\epsilon:=t(\epsilon)-s(\epsilon).
\end{equation}
The neutral emission satisfies
\begin{equation}
\label{eq:continuo-comun-neutral-app-datos}
 \operatorname{res}^{\diamond}(\varepsilon^0_k(a))=0,
 \qquad
 \operatorname{quo}^{\diamond}(\varepsilon^0_k(a))=0,
 \qquad
 \partial\varepsilon^0_k(a)=0.
\end{equation}

\begin{definition}[Directed graph of APP emissions]
\label{def:continuo-comun-grafo-app}
Let \(\mathsf E_k(a)\) be the finite set of APP emissions of length one
that depart from \(a\in\mathsf A_k\).  We conserve the emissions with
multiplicity: two operations with the same target but with different
residue, quotient, or provenance are different elements of
\(\mathsf E_k(a)\).  By \eqref{eq:continuo-comun-app-degeneracion},
\begin{equation}
\label{eq:continuo-comun-out-no-vacio-app}
 \varepsilon^0_k(a)\in\mathsf E_k(a),
 \qquad
 1\leq \#\mathsf E_k(a)<\infty.
\end{equation}
\end{definition}

\section{TRIT Lift: produced genealogical orientation and carry}
\label{sec:continuo-comun-trit-lift}

The signed incidence of an APP emission is an integer
\(\Delta(\epsilon)\). The balanced representative modulo three uniquely
gives
\begin{equation}
\label{eq:continuo-comun-trit-euclid}
 \Delta(\epsilon)
 =\operatorname{or}(\epsilon)+3\kappa(\epsilon),
 \qquad
 \operatorname{or}(\epsilon)\in\{-1,0,+1\},
 \quad
 \kappa(\epsilon)\in\mathbb Z.
\end{equation}
Thus, orientation and carry are not conditions imposed on a history:
they are total functions of the APP data.

To make the composition visible, let
\(\mathsf C_{\epsilon}:\mathsf M_{s(\epsilon)}\to
\mathsf M_{t(\epsilon)}\) be the memory transport induced by the emission.
The native carry of a composition satisfies
\begin{equation}
\label{eq:continuo-comun-cociclo-carry}
 \kappa(\epsilon_2\epsilon_1)
 =\kappa(\epsilon_2)
  +\mathsf C_{\epsilon_2}\kappa(\epsilon_1),
\end{equation}
whenever \(t(\epsilon_1)=s(\epsilon_2)\). For the identity and the neutral emission,
\begin{equation}
\label{eq:continuo-comun-carry-unidad}
 \kappa(I_a)=0,
 \qquad
 \kappa(\varepsilon^0_k(a))=0,
 \qquad
 \mathsf C_{I_a}=I_{\mathsf M_a}.
\end{equation}

\begin{lemma}[Associativity of the carry ledger]
\label{lem:continuo-comun-asociatividad-carry}
For three composable emissions, the evaluation of
\(\kappa(\epsilon_3\epsilon_2\epsilon_1)\) does not depend on the parenthesization.
\end{lemma}

\begin{proof}
Applying \eqref{eq:continuo-comun-cociclo-carry} twice to the left
parenthesization gives
\[
 \kappa(\epsilon_3)
 +\mathsf C_{\epsilon_3}\kappa(\epsilon_2)
 +\mathsf C_{\epsilon_3}\mathsf C_{\epsilon_2}
  \kappa(\epsilon_1).
\]
The right parenthesization produces the same expression because the transport of
memory is functorial:
\(\mathsf C_{\epsilon_3\epsilon_2}
=\mathsf C_{\epsilon_3}\mathsf C_{\epsilon_2}\).
\end{proof}

The oriented reversal of a reversible emission is denoted by
\(\overline\epsilon\).  TRIT acts by
\begin{equation}
\label{eq:continuo-comun-reversion-trit}
 \operatorname{or}(\overline\epsilon)
 =-\operatorname{or}(\epsilon),
 \qquad
 \partial\overline\epsilon=-\partial\epsilon,
 \qquad
 \mathsf C_{\overline\epsilon}=\mathsf C_\epsilon^{-1}.
\end{equation}
The inverse carry is determined, not chosen, by
\begin{equation}
\label{eq:continuo-comun-carry-inverso}
 \kappa(\overline\epsilon)
 =-\mathsf C_\epsilon^{-1}\kappa(\epsilon),
\end{equation}
for \(\kappa(\overline\epsilon\epsilon)=0\).

\section{TPK transport of the \(2\) leaves and of memory}
\label{sec:continuo-comun-tpk-transporte}

TPK chronology is part of the type of each emission and is not
reconstructed a posteriori by counting letters. Here we recover the action of the
ometer of \cref{def:odometro-9-12,prop:proyeccion-fase-longitud}. Let
\begin{equation}
\label{eq:continuo-comun-ciclo-fase-tpk}
 C_9:=\mathbb Z/9\mathbb Z,\qquad
 \widetilde{\mathcal O}_{9}:=\mathbb Z_{\geq0}\times C_9,\qquad
 \mathcal O_{9,12}:=\mathbb Z/12\mathbb Z\times C_9.
\end{equation}
The unrolled phase-and-memory state is
\(\mathsf Q_k=(w_k,r_k)\in\widetilde{\mathcal O}_{9}\); its reduction
\((w,r)\mapsto([w]_{12},r)\) recovers
\(\mathcal O_{9,12}\). The canonical seed of the clock is
\begin{equation}
\label{eq:continuo-comun-semilla-fase-tpk}
 \mathsf Q_0(a_0):=(0,[0]_9)\qquad(a_0\in\mathsf A_0);
\end{equation}
the other choices of origin are their coherent translations by \(C_9\)
and do not change the update law. For
\(\bar r\in\{0,\ldots,8\}\), we set
\begin{equation}
\label{eq:continuo-comun-carry-fase-tpk}
 \chi_9(r):=\left\lfloor\frac{\bar r+1}{9}\right\rfloor.
\end{equation}
Every emission \(\epsilon\) that lleva depth \(k\) to depth \(k+1\)
acts in the odometro by
\begin{equation}
\label{eq:continuo-comun-accion-fase-tpk}
 \sigma_9(w,r):=(w+\chi_9(r),r+[1]_9),\qquad
 \mathsf Q_{k+1}(\epsilon_*\widehat x_k)
 =\sigma_9\mathsf Q_k(\widehat x_k).
\end{equation}
Its phase component is, by typed definition,
\begin{equation}
\label{eq:continuo-comun-fase-emision-tpk}
 \operatorname{ph}(\epsilon):r_k\longmapsto r_k+[1]_9,
 \qquad
 \operatorname{ph}(\epsilon_n\cdots\epsilon_1):
 r_k\longmapsto r_k+[n]_9.
\end{equation}
We write \(g_k:=r_k\in C_9\) for the published phase and
\(q^{(9)}_k:=w_k\) for the vertical memory, and
\(\beta_k:=[w_k]_{12}\) for its duodecaphase reduction. The iteration of
\eqref{eq:continuo-comun-accion-fase-tpk} is
\begin{equation}
\label{eq:continuo-comun-iteracion-fase-tpk}
 \sigma_9^n(w,r)=
 \left(w+\left\lfloor\frac{\bar r+n}{9}\right\rfloor,
       r+[n]_9\right).
\end{equation}
In efecto, the sumandos
\(\chi_9(r),\chi_9(r+[1]_9),\ldots,\chi_9(r+[n-1]_9)\)
cuentan exactly the cruces of the representative \(8\) to the representative
\(0\), that are
\(\lfloor(\bar r+n)/9\rfloor\). In particular,
\begin{equation}
\label{eq:continuo-comun-retorno-fase-avance-memoria}
 \sigma_9^9(w,r)=(w+1,r),\qquad
 g_{k+9}=g_k,\qquad q^{(9)}_{k+9}=q^{(9)}_k+1.
\end{equation}
Even the neutral emission of APP realizes
\(g_k\mapsto g_{k+1}\): it is neutral in residue, quotient, carry, and boundary,
but records that an extension has occurred. The phase word
\begin{equation}
\label{eq:continuo-comun-palabra-fase-tpk}
 \boldsymbol\tau_k^{(N)}:=(g_k,g_{k+1},\ldots,g_N),
 \qquad g_{j+1}=g_j+[1]_9,
\end{equation}
and the memory \(q^{(9)}\) type the phase return and the
state change as distinct objects.

\subsection{Simultaneous Renewal of the Complete Fiber}
\label{sec:continuo-comun-kph}

The preceding odometer conserves the chronology of each extension. The same
nonadic turn also has the internal Markov realization, already constructed
over the APP--TPK catalogue, which simultaneously renews the two coordinates
of the complete fiber. It is essential to conserve both realizations and not to
identify a generic edge of the tree with an emission of the catalogue without
a typing map.

\begin{proposition}[Internal Markov Realization of the Nonadic Connection]
\label{prop:continuo-comun-kph-heredada}
\label{pII-full:thm:seleccion-novena-fase}
Let the APP seed domain, its observable emission, and its multiplicities be
\[
 T_{\min}:\Omega_{\rm seed}\twoheadrightarrow\mathcal U_6,
 \qquad
 m(u):=\#T_{\min}^{-1}(u),
\]
and the factorization produced by the TPK
\[
 \mathcal U_6\simeq\mathcal S_+\times\mathcal S_\times,
 \qquad |\mathcal S_+|=18,\qquad |\mathcal S_\times|=26.
\]
On functions in \(\mathcal U_6\), the two conditional expectations act
by
\[
 (E_+f)(s,e)=\frac1{18}\sum_{s'\in\mathcal S_+}f(s',e),
 \qquad
 (E_\times f)(s,e)=\frac1{26}
 \sum_{e'\in\mathcal S_\times}f(s,e').
\]
For
\[
 \tau=(+1,-1,0,+1,-1,0,+1,-1,0),
 \qquad
 P_{+1}=E_+,\quad P_{-1}=E_\times,\quad P_0=I,
\]
the catalog transfer and its elevation to all seeds are
\begin{align}
 \mathcal K_{\rm ph}\bigl((u,r),(v,r+1)\bigr)
 &=P_{\tau(r)}(u,v),
 \label{eq:continuo-comun-kph-catalogo}\\
 L_0(\omega,\omega')&=\mathbf1_{\{\omega=\omega'\}},\nonumber\\
 L_+(\omega,\omega')
 &=\frac{E_+(T_{\min}\omega,T_{\min}\omega')}
         {m(T_{\min}\omega')},\nonumber\\
 L_\times(\omega,\omega')
 &=\frac{E_\times(T_{\min}\omega,T_{\min}\omega')}
         {m(T_{\min}\omega')},\nonumber\\
 \widehat{\mathcal K}_{\rm ph}
 \bigl((\omega,r),(\omega',r+1)\bigr)
 &=L_{\tau(r)}(\omega,\omega').
 \label{eq:continuo-comun-kph-elevada}
\end{align}
For \(u=T_{\min}\omega\), it holds phase by phase that
\begin{equation}
\label{eq:continuo-comun-kph-lumpability}
 \sum_{\omega':\,T_{\min}\omega'=v}
 \widehat{\mathcal K}_{\rm ph}
 \bigl((\omega,r),(\omega',r+1)\bigr)
 =
 \mathcal K_{\rm ph}\bigl((u,r),(v,r+1)\bigr).
\end{equation}
In particular,
\begin{equation}
\label{eq:continuo-comun-kph-novena-potencia}
 \boxed{\;
 \mathcal K_{\rm ph}^{\,9}\big|_r=E_+E_\times,
 \qquad
 (E_+E_\times)(u,v)=\frac1{468}.
 \;}
\end{equation}
\end{proposition}

\begin{proof}
In a phase \(+\) or \(\times\), summing
\eqref{eq:continuo-comun-kph-elevada} over
\(T_{\min}^{-1}(v)\) exactly cancels the denominator \(m(v)\); in a
zero phase, \(L_0\) projects onto the identity. This proves
\eqref{eq:continuo-comun-kph-lumpability}. The operators \(E_+\),
\(E_\times\), and \(I\) are idempotent and commute, and each appears three
times in the nonadic word. Therefore
\[
 \prod_{r\in C_9}P_{\tau(r)}
 =E_+^3E_\times^3I^3=E_+E_\times.
\]
The composition first renews one coordinate of cardinal \(18\) and then
the other of cardinal \(26\), so that each target receives
\(1/(18\cdot26)=1/468\).
\end{proof}

\begin{remark}[Two internal realizations, without fictitious identification]
\label{rem:continuo-comun-kph-no-identificacion-aristas}
The transfer \(\mathcal K_{\rm ph}\) is the global realization of the
complete catalog, and \(\sigma_9\) is the chronological action on each state
of the tree. One is not substituted for the other. An edge-by-edge equality
would require an additional map
\[
 \operatorname{Cat}_{k,x}:\operatorname{Out}(x)\longrightarrow\mathcal U_6
\]
that verified the corresponding sum identity on fibers. That map
is not used in this volume. The complete monodromy
\eqref{eq:continuo-comun-kph-novena-potencia} proceeds from the
APP--TPK catalogue, while the genealogy of each history remains in the tree and
its ledger.
\end{remark}

TPK assigns to each admissible emission a typed transport of the leaf pair,
\begin{equation}
\label{eq:continuo-comun-transporte-hojas}
 U_\epsilon:
 \mathsf H_{s(\epsilon)}^{\Sigma}\times
 \mathsf H_{s(\epsilon)}^{\Pi}
 \longrightarrow
 \mathsf H_{t(\epsilon)}^{\Sigma}\times
 \mathsf H_{t(\epsilon)}^{\Pi},
\end{equation}
with
\begin{equation}
\label{eq:continuo-comun-transporte-hojas-funtorial}
 U_{I_a}=I,
 \qquad
 U_{\epsilon_2\epsilon_1}=U_{\epsilon_2}U_{\epsilon_1}.
\end{equation}
The codomain of \(U_\epsilon\) remains the complete pair. No
operation in this chapter replaces
\((\Sigma,\Pi)\) by one of its coordinates.

Let \(m\in\mathsf M_{s(\epsilon)}\) be an accumulated memory.  The TPK action
over it is the affine transformation
\begin{equation}
\label{eq:continuo-comun-accion-memoria}
 \mathsf U_\epsilon(m)
 :=\mathsf C_\epsilon m+\kappa(\epsilon).
\end{equation}

\begin{lemma}[TPK Action on Memory]
\label{lem:continuo-comun-accion-memoria}
The transformations \(\mathsf U_\epsilon\) satisfy
\begin{equation}
\label{eq:continuo-comun-accion-memoria-composicion}
 \mathsf U_{I_a}=I,
 \qquad
 \mathsf U_{\epsilon_2\epsilon_1}
 =\mathsf U_{\epsilon_2}\mathsf U_{\epsilon_1}.
\end{equation}
\end{lemma}

\begin{proof}
The identity follows from \eqref{eq:continuo-comun-carry-unidad}.  For the
composition,
\begin{align*}
 \mathsf U_{\epsilon_2}\mathsf U_{\epsilon_1}(m)
 &=\mathsf C_{\epsilon_2}
   \bigl(\mathsf C_{\epsilon_1}m+\kappa(\epsilon_1)\bigr)
   +\kappa(\epsilon_2)\\
 &=\mathsf C_{\epsilon_2\epsilon_1}m
   +\kappa(\epsilon_2)
   +\mathsf C_{\epsilon_2}\kappa(\epsilon_1)\\
 &=\mathsf C_{\epsilon_2\epsilon_1}m
   +\kappa(\epsilon_2\epsilon_1),
\end{align*}
where the last equality is
\eqref{eq:continuo-comun-cociclo-carry}.
\end{proof}

A finite route is a composable word.
\begin{equation}
\label{eq:continuo-comun-ruta}
 \gamma=\epsilon_1\epsilon_2\cdots\epsilon_k.
\end{equation}
Its boundary, transport, carry, and memory are
\begin{align}
 \partial\gamma
 &:=\sum_{j=1}^{k}\partial\epsilon_j,
 \label{eq:continuo-comun-frontera-ruta}\\
 U_\gamma
 &:=U_{\epsilon_k}\cdots U_{\epsilon_1},
 \label{eq:continuo-comun-transporte-ruta}\\
 \kappa(\gamma)
 &:=\kappa(\epsilon_k)
   +\sum_{j=1}^{k-1}
    \mathsf C_{\epsilon_k}\cdots\mathsf C_{\epsilon_{j+1}}
    \kappa(\epsilon_j),
 \label{eq:continuo-comun-carry-ruta}\\
 m(\gamma;m_0)
 &:=\mathsf U_{\epsilon_k}\cdots
       \mathsf U_{\epsilon_1}(m_0).
 \label{eq:continuo-comun-memoria-ruta}
\end{align}
In \eqref{eq:continuo-comun-frontera-ruta}, the interior endpoints
cancel, leaving
\begin{equation}
\label{eq:continuo-comun-frontera-telescopica}
 \partial\gamma=t(\epsilon_k)-s(\epsilon_1).
\end{equation}

\section{The common state and its total domain}
\label{sec:continuo-comun-estado-total}

\begin{definition}[Elementary ledger]
\label{def:continuo-comun-ledger-elemental}
For an emission that appears between two consecutive enriched prefixes,
we denote by
\(\mathsf Q^-(\epsilon)=(w^-(\epsilon),r^-(\epsilon))\) the odometer of the
source prefix and by
\begin{equation}
\label{eq:continuo-comun-ledger-fase-extremos}
 \mathsf Q^+(\epsilon):=\sigma_9\mathsf Q^-(\epsilon)
 =(w^-(\epsilon)+\chi_9(r^-(\epsilon)),r^-(\epsilon)+[1]_9)
\end{equation}
that of the destination prefix. Thus the odometer is not applied to a bare APP
cell. The emission record is
\begin{equation}
\label{eq:continuo-comun-ledger-elemental}
 \lambda(\epsilon):=\bigl(
 s(\epsilon),t(\epsilon),
 \operatorname{res}^{\Sigma}(\epsilon),
 \operatorname{quo}^{\Sigma}(\epsilon),
 \operatorname{res}^{\Pi}(\epsilon),
 \operatorname{quo}^{\Pi}(\epsilon),
 \operatorname{or}(\epsilon),
 \kappa(\epsilon),
 \partial\epsilon,
 \mathsf Q^-(\epsilon),\mathsf Q^+(\epsilon)\bigr).
\end{equation}
The ledger of \(\gamma\) is the word ordered
\begin{equation}
\label{eq:continuo-comun-ledger-ruta}
 \operatorname{Led}(\gamma)
 :=\bigl(\lambda(\epsilon_1),\ldots,
          \lambda(\epsilon_k)\bigr).
\end{equation}
Preserving this word, and not only its sum, distinguishes two histories with the same terminal output.
\end{definition}

\subsection{Complete reader of the TPK state}
\label{sec:continuo-comun-lector-estado-tpk}

The depth \(K\) of TPK blocks and the depth \(k\) of the
extension tree are distinct types. In this subsection, \(K\) denotes
exclusively the depth of the TPK state already generated by the operator
mechanism; it is not identified with \(k\) or with a generic edge.

\begin{definition}[Compact nonadic memory reader]
\label{def:continuo-comun-lector-compacto-nonadico}
The integral holonomic TPK state of depth \(K\) is
\begin{equation}
\label{eq:continuo-comun-estado-tpk-completo}
 v_K=
 \bigl(
 B_K,\boldsymbol\delta_K,\boldsymbol x_K,
 \boldsymbol{\mathcal C}_K,\Sigma_K,
 \omega_K,\ell_K,g(K),\Lambda_K
 \bigr),
 \qquad
 B_K=(u_K^\pi,u_K^e,u_K^\varphi).
\end{equation}
Its compact reader does not add coordinates: it publishes them from that same state,
\begin{equation}
\label{eq:continuo-comun-lector-compacto-nonadico}
 \mathfrak r_9(v_K)
 :=\bigl(B_K,C_K,p_K,c_K,\Sigma_K,W_K,g(K)\bigr),
\end{equation}
where
\begin{equation}
\label{eq:continuo-comun-lector-compacto-componentes}
\begin{aligned}
 C_K&:=\boldsymbol{\mathcal C}_K
      =\bigl(C_K^\chi\bigr)_{\chi\in\{\pi,e,\varphi\}},\\
 C_K^\chi&:=(N_K^\chi,L_K^\chi,r_K^\chi,
             D_K^\chi,A_K^\chi,B_K^\chi),\\
 p_K&:=(B_0,\ldots,B_K),\\
 c_K&:=(\boldsymbol\delta_K,\boldsymbol x_K,
       \kappa(\gamma^{\rm TPK}_{0:K})),
\end{aligned}
\end{equation}
and the oriented shadow of Witt is preserved as
\begin{equation}
\label{eq:continuo-comun-sombra-witt}
 W_K:=
 \bigl((u_K^\chi,u_K^\chi A_W)\bigr)_
       {\chi\in\{\pi,e,\varphi\}},
 \qquad
 A_W=
 \begin{pmatrix}
 0&1&1&1&1&1\\
 1&0&1&2&2&1\\
 1&1&0&1&2&2\\
 1&2&1&0&1&2\\
 1&2&2&1&0&1\\
 1&1&2&2&1&0
 \end{pmatrix},
 \qquad A_W^2=-I_6.
\end{equation}
The route \(\gamma^{\rm TPK}_{0:K}\) is the word of extensions encoded
by \(\Lambda_K\); its prefix of length \(j\) is the one that produces
\(v_j\). The signature is the output of the explicit internal primitive
\begin{equation}
\label{eq:continuo-comun-primitiva-firma-tpk}
 \Sigma_B:
 \mathsf B_{\rm TPK}\times\mathsf C_{\rm TPK}
 \times\mathbb Z[\mathsf A]\times\Omega_{12}\times\mathsf L_{\rm TPK}
 \longrightarrow\mathsf S_{\rm TPK},
 \quad
 (B,c,\partial\gamma,\omega,\ell)\longmapsto
 \Sigma_B(B,c,\partial\gamma,\omega,\ell),
\end{equation}
with \(\Sigma_B(g\cdot-)=g\cdot\Sigma_B(-)\) for each TPK symmetry. Its
explicit starting structure is
\((B_K,c_K,\partial\gamma^{\rm TPK}_{0:K},\omega_K,\ell_K)\). It is not replaced by an
external signature nor inferred retrospectively from a numerical reader.
\end{definition}

The action of step \(K\mapsto K+1\) is written without omitting the deep state.
For each \(\chi\in\{\pi,e,\varphi\}\),
\begin{equation}
\label{eq:continuo-comun-update-hensel-tpk}
 y_K^\chi=x_K^\chi A_H,\qquad
 u_{K+1}^\chi=y_K^\chi\bmod3,\qquad
 x_{K+1}^\chi=\frac{y_K^\chi-u_{K+1}^\chi}{3},
\end{equation}
with
\[
 A_H=
 \begin{pmatrix}
 2&2&2&1&2&1\\
 2&1&2&2&1&1\\
 1&0&0&1&0&2\\
 0&0&1&1&1&0\\
 2&2&2&0&2&1\\
 2&1&2&1&0&0
 \end{pmatrix},
 \qquad \det A_H=1.
\]
If
\[
 \operatorname{ind}_3(u_1,\ldots,u_6)
 :=\sum_{j=1}^{6}\overline u_j\,3^{6-j}
 \in\{0,\ldots,728\},
\]
the update of the prefix and its \(729\) subcylinders is
\begin{align}
 N_{K+1}^\chi
 &=729N_K^\chi+\operatorname{ind}_3(u_{K+1}^\chi),
 &L_{K+1}^\chi&=L_K^\chi+6,
 \label{eq:continuo-comun-update-cilindro-indice}\\
 I_{K,u}^\chi
 &:=
 \left[
  \frac{729N_K^\chi+u}{3^{L_K^\chi+6}},
  \frac{729N_K^\chi+u+1}{3^{L_K^\chi+6}}
 \right),
 &0&\leq u<729.
 \label{eq:continuo-comun-update-cilindro-intervalo}
\end{align}
The decimal cylinder and the associated carry are updated by the two integer
cases of the same operator \(\mathcal C_{3,1000}\). If a new decimal triad is not
published,
\begin{equation}
\label{eq:continuo-comun-update-cilindro-sin-triada}
 r_{K+1}^\chi=r_K^\chi,\qquad D_{K+1}^\chi=D_K^\chi,\qquad
 A_{K+1}^\chi=729A_K^\chi+
 \operatorname{ind}_3(u_{K+1}^\chi)1000^{r_K^\chi},\qquad
 B_{K+1}^\chi=A_{K+1}^\chi+1000^{r_K^\chi}.
\end{equation}
The primitive \(\mathcal C_{3,1000}\) produces the option datum
\[
 d_{K+1,\chi}^{\rm dec}\in
 \{\bot\}\sqcup\{0,\ldots,999\};
\]
in the first case it equals \(d_{K+1,\chi}^{\rm dec}=\bot\). If a triad occurs, we write
\(d_{K+1,\chi}^{\rm dec}=\delta_{K+1}^\chi\in\{0,\ldots,999\}\) and
\begin{equation}
\label{eq:continuo-comun-update-cilindro-con-triada}
 \begin{aligned}
 r_{K+1}^\chi&=r_K^\chi+1,\qquad
 D_{K+1}^\chi=1000D_K^\chi+\delta_{K+1}^\chi,\\
 A_{K+1}^\chi&=729000A_K^\chi
 +\operatorname{ind}_3(u_{K+1}^\chi)1000^{r_K^\chi+1}
 -729\delta_{K+1}^\chi3^{L_K^\chi},\\
 B_{K+1}^\chi&=A_{K+1}^\chi+1000^{r_K^\chi+1}.
 \end{aligned}
\end{equation}
The register \(\boldsymbol\delta_{K+1}\) adds those triads and the two joint
cocycles
\begin{equation}
\label{eq:continuo-comun-update-triadas-carry-conjunto}
 \begin{gathered}
 q_{K+1}=(u_{K+1}^\pi+u_{K+1}^e+u_{K+1}^\varphi)\bmod3,
 \quad
 c_{K+1}^{\rm joint}=\frac{u_{K+1}^\pi+u_{K+1}^e+u_{K+1}^\varphi-q_{K+1}}3,\\
 a_{K+1}=(u_{K+1}^\pi+u_{K+1}^e-u_{K+1}^\varphi)\bmod3,
 \quad
 h_{K+1}^{\rm joint}=\frac{u_{K+1}^\pi+u_{K+1}^e-u_{K+1}^\varphi-a_{K+1}}3,\\
 \boldsymbol\delta_{K+1}
 =\boldsymbol\delta_K\star
 \bigl((d_{K+1,\chi}^{\rm dec})_\chi,
       c_{K+1}^{\rm joint},h_{K+1}^{\rm joint}\bigr).
 \end{gathered}
\end{equation}
Thus \(C_{K+1}\), \(\boldsymbol\delta_{K+1}\), and
\(\boldsymbol x_{K+1}\) have material actions, not merely names.
The other coordinates are simultaneously updated by
\begin{align}
 p_{K+1}&=p_K\star B_{K+1},
 \label{eq:continuo-comun-update-prefijo-tpk}\\
 \kappa(\gamma^{\rm TPK}_{0:K+1})
 &=\kappa(\varepsilon^{\rm TPK}_{K+1})
   +\mathsf C_{\varepsilon^{\rm TPK}_{K+1}}
    \kappa(\gamma^{\rm TPK}_{0:K}),
 \label{eq:continuo-comun-update-carry-lector-tpk}\\
 c_{K+1}^{\rm comp}
 &:=(\boldsymbol\delta_{K+1},\boldsymbol x_{K+1},
       \kappa(\gamma^{\rm TPK}_{0:K+1})),
 \label{eq:continuo-comun-update-carry-compacto-tpk}\\
 \Sigma_{K+1}
 &=\Sigma_B\bigl(
 B_{K+1},c_{K+1}^{\rm comp},
 \partial\gamma^{\rm TPK}_{0:K}
 +\partial\varepsilon^{\rm TPK}_{K+1},
 \omega_{K+1},\ell_{K+1}\bigr),
 \label{eq:continuo-comun-update-firma-tpk}\\
 W_{K+1}
 &=\bigl((u_{K+1}^\chi,u_{K+1}^\chi A_W)\bigr)_\chi,
 &g(K+1)&=g(K)+[1]_9.
 \label{eq:continuo-comun-update-witt-fase-tpk}
\end{align}
Since \(A_H\) is invertible modulo three and \(A_W^2=-I_6\), the residue
and shadow updates do not collapse distinct blocks. Truncation preserves
the complete sequences
\[
 (B_0,\ldots,B_K),\ (C_0,\ldots,C_K),\
 (p_0,\ldots,p_K),\ (c_0,\ldots,c_K),\
 (\Sigma_0,\ldots,\Sigma_K),\ (W_0,\ldots,W_K)
\]
and simultaneously removes its last entry. Hence the reader
\(\mathfrak r_9\) commutes with the erasure of a block.

\begin{proposition}[Universal scope and certified scope of non-rebooting]
\label{prop:continuo-comun-no-reinicio-alcance}
On every horizontal circuit of nine steps, the following are preserved
\[
 g(K+9)=g(K),\qquad
 q_{\rm mem}(K+9)=q_{\rm mem}(K)+1,
\]
and the prefix is prolonged with nine emissions, so the state is not
reset. In the certified joint section
\(\chi\in\{\pi,e,\varphi\}\), \(0\leq K<387\), it also holds that
\begin{equation}
\label{eq:continuo-comun-no-reinicio-certificado}
 (p_K,C_K)\neq(p_{K+9},C_{K+9}),\qquad
 (B_K,\mathbf N_K,W_K)\neq
 (B_{K+9},\mathbf N_{K+9},W_{K+9}),
 \qquad
 \mathbf N_K:=(N_K^\pi,N_K^e,N_K^\varphi).
\end{equation}
\end{proposition}

\begin{proof}
The two universal identities are the exact iteration of
\eqref{eq:continuo-comun-accion-fase-tpk}; the prefix preserves the nine
entries adjoined by \eqref{eq:continuo-comun-update-prefijo-tpk}. For the
joint section, the recurrence
\eqref{eq:continuo-comun-update-cilindro-indice} compares the \(387\) nonadic pairs
of each of the three channels and produces
\eqref{eq:continuo-comun-no-reinicio-certificado}; this is the assertion
certified in
\cref{thm:generacion-monodromica-conjunta-rev7}. Finally,
\(u\mapsto(u,uA_W)\) is injective, so the change of block also preserves
the change of its complete shadow. This last inequality is not extended
to an arbitrary neutral prolongation outside the certified
section.
\end{proof}

\begin{definition}[APP--TRIT--TPK state of depth \(k\)]
\label{def:continuo-comun-estado}
For a path \(\gamma=\epsilon_1\cdots\epsilon_k\) originating in an APP
cell \(a_0\), the common integral holonomic state is
\begin{equation}
\label{eq:continuo-comun-estado}
 \widehat x_k(\gamma):=\Bigl(
  a_0,\ldots,a_k;
  \operatorname{Sh}_0,\ldots,\operatorname{Sh}_k;
  \mathbf{res}_k,\mathbf{quo}_k;
  \mathbf{or}_k,\mathbf\kappa_k;
  \mathsf Q_k;g_k;q^{(9)}_k;\beta_k;m_k;\partial\gamma;\gamma;
  \operatorname{Led}(\gamma)
 \Bigr),
\end{equation}
where
\begin{align*}
 \mathbf{res}_k
 &:=\bigl(
  (\operatorname{res}^{\Sigma}(\epsilon_j),
   \operatorname{res}^{\Pi}(\epsilon_j))
  \bigr)_{1\leq j\leq k},\\
 \mathbf{quo}_k
 &:=\bigl(
  (\operatorname{quo}^{\Sigma}(\epsilon_j),
   \operatorname{quo}^{\Pi}(\epsilon_j))
  \bigr)_{1\leq j\leq k},\\
 \mathbf{or}_k
 &:=\bigl(\operatorname{or}(\epsilon_j)\bigr)_{1\leq j\leq k},
 \qquad
 \mathbf\kappa_k
 :=\bigl(\kappa(\epsilon_j)\bigr)_{1\leq j\leq k},\\
 m_k&:=m(\gamma;m_0).
\end{align*}
The intermediate sheets are obtained by
\(\operatorname{Sh}_{j}=U_{\epsilon_j}\operatorname{Sh}_{j-1}\).
\end{definition}

The simultaneous presence of the route and the ledger prevents identifying
distinct words that terminate in the same cell.  The state is not a list of
cardinals: each coordinate has the type and composition equations of the
preceding sections.

\begin{definition}[Total depth finite domain]
\label{def:continuo-comun-dominio-total}
Let
\begin{equation}
\label{eq:continuo-comun-dominio-total}
 \widehat{\mathcal X}^{\rm tot}_k
 :=\left\{
  \widehat x_k(\epsilon_1\cdots\epsilon_k):
  \begin{array}{l}
   a_0\in\mathsf A_0,\\
   \epsilon_j\in\mathsf E_{j-1}(a_{j-1}),\\
   s(\epsilon_j)=a_{j-1},\ t(\epsilon_j)=a_j
  \end{array}
 \right\}.
\end{equation}
No condition from a later publication is added to \eqref{eq:continuo-comun-dominio-total}.
Being an element of the domain means exactly
having been produced by the total actions
\eqref{eq:continuo-comun-residuo-cociente-app},
\eqref{eq:continuo-comun-trit-euclid}, and
\eqref{eq:continuo-comun-accion-fase-tpk}--
\eqref{eq:continuo-comun-accion-memoria}.
\end{definition}

The common truncation deletes the last emission and all and only the entries of the ledger that it produced:
\begin{equation}
\label{eq:continuo-comun-truncamiento-total}
 \rho_{k+1,k}:\widehat{\mathcal X}^{\rm tot}_{k+1}
 \longrightarrow\widehat{\mathcal X}^{\rm tot}_k,
 \qquad
 \rho_{k+1,k}\widehat x_{k+1}
 =\widehat x_k.
\end{equation}
Let \(\varepsilon^0_k(x)\) be the neutral emission
\eqref{eq:continuo-comun-emision-neutra-app} applied to the APP endpoint of
\(x\). The neutral prolongation defines
\begin{equation}
\label{eq:continuo-comun-seccion-total}
 \eta_k:\widehat{\mathcal X}^{\rm tot}_k
 \longrightarrow\widehat{\mathcal X}^{\rm tot}_{k+1},
 \qquad
 \eta_k(x):=(\varepsilon^0_k(x))_*x,
 \qquad
 \rho_{k+1,k}\eta_k=I.
\end{equation}

\begin{proposition}[Finiteness, non-emptiness, and surjectivity]
\label{prop:continuo-comun-total-finito}
For every \(k\), the set
\(\widehat{\mathcal X}^{\rm tot}_k\) is finite and nonempty, and
\(\rho_{k+1,k}\) is surjective.
\end{proposition}

\begin{proof}
APP produces a nonempty finite set at depth zero. At each step,
every state has a finite set of emissions through
\eqref{eq:continuo-comun-out-no-vacio-app}; by induction, there are only finitely many
words of length \(k\). The sequence of neutral prolongations
produces at least one word of each length. Finally,
\eqref{eq:continuo-comun-seccion-total} is a right inverse of
truncation, so the latter is surjective.
\end{proof}

\section{Finite reversible symmetry groupoid}
\label{sec:continuo-comun-groupoide}

The reversible operations of APP, TRIT, and TPK act on the entire state.
Let \(\mathsf U\) be the group of coherent families
\(g=(g_k)_{k\geq0}\) generated by those operations, where
\begin{equation}
\label{eq:continuo-comun-simetria-coherente}
 \rho_{k+1,k}g_{k+1}=g_k\rho_{k+1,k},
 \qquad
 g_{k+1}\eta_k=\eta_kg_k.
\end{equation}
We denote by \(\mathsf U_k\) its image in
\(\operatorname{Sym}(\widehat{\mathcal X}^{\rm tot}_k)\). This image is
finite because it acts on a finite set. The action of \(g_k\) on
\begin{equation*}
 \widehat x=(a_{\leq k};\operatorname{Sh}_{\leq k};
 \mathbf{res},\mathbf{quo};\mathbf{or},\mathbf\kappa;
 \mathsf Q_k;g_k;q_k^{(9)};\beta_k;m;\partial\gamma;\gamma;
 \operatorname{Led}(\gamma))
\end{equation*}
is
\begin{equation}
\label{eq:continuo-comun-accion-reversible}
\begin{aligned}
 g_k\cdot\widehat x:=\bigl(&
 g_k a_{\leq k};
 U_{g_k}\operatorname{Sh}_{\leq k};
 \mathbf{res}(g_k\gamma),\mathbf{quo}(g_k\gamma);
 \mathbf{or}(g_k\gamma),\mathbf\kappa(g_k\gamma);\\
 &\mathsf Q(g_k\gamma);g(g_k\gamma);q^{(9)}(g_k\gamma);
 \beta(g_k\gamma);
 \mathsf U_{g_k}(m);
 (g_k)_*\partial\gamma;
 g_k\gamma;
 \operatorname{Led}(g_k\gamma)\bigr).
\end{aligned}
\end{equation}
All the terms on the right-hand side are recalculated by means of the APP,
TRIT, and TPK rules; they are not copied as independent labels.

\begin{definition}[Common action groupoid]
\label{def:continuo-comun-groupoide}
The finite groupoid of depth \(k\) is
\begin{equation}
\label{eq:continuo-comun-groupoide}
 \mathfrak G_k
 :=\mathsf U_k\ltimes\widehat{\mathcal X}^{\rm tot}_k.
\end{equation}
Its objects are the common states.  An arrow
\((g_k,\widehat x)\) has source \(\widehat x\), target
\(g_k\cdot\widehat x\), identity \((I,\widehat x)\), and inverse
\((g_k^{-1},g_k\cdot\widehat x)\).
\end{definition}

\begin{lemma}[The action preserves all equations of state]
\label{lem:continuo-comun-groupoide-conserva}
The action \eqref{eq:continuo-comun-accion-reversible} preserves the pair of
sheets, the APP divisions, the TRIT orientation, the carry, the phase
odometer, the vertical memory, its dodecaphase reduction, the affine memory, the
boundary, the route, and the ledger, in their respective types.  Moreover,
\begin{equation}
\label{eq:continuo-comun-groupoide-accion}
 I\cdot\widehat x=\widehat x,
 \qquad
 (g_{2,k}g_{1,k})\cdot\widehat x
 =g_{2,k}\cdot(g_{1,k}\cdot\widehat x).
\end{equation}
\end{lemma}

\begin{proof}
The APP divisions are unique, so the action on an emission
redetermines its residues and quotients.  The balanced decomposition
\eqref{eq:continuo-comun-trit-euclid} is likewise unique.  Composition of the
carry is associative by
\cref{lem:continuo-comun-asociatividad-carry}; the action on memory is associative
by \cref{lem:continuo-comun-accion-memoria}; and the boundary is functorial by
\eqref{eq:continuo-comun-frontera-telescopica}.  Finally, the action on
the ordered word induces the action on its ledger entry by entry.
Coherence with the odometer recalculates \(Q\), \(g\), and \(q^{(9)}\) from
the transformed route and preserves
\eqref{eq:continuo-comun-accion-fase-tpk}. These identities prove
\eqref{eq:continuo-comun-groupoide-accion}.
\end{proof}

\section{Finite tree of all prolongations}
\label{sec:continuo-comun-arbol}

For \(x\in\widehat{\mathcal X}^{\rm tot}_k\), let
\(\operatorname{Out}(x)\) be the set of all admissible TPK emissions
of length one from the endpoint of \(x\).  An emission belongs to
\(\operatorname{Out}(x)\) only when its action jointly updates all
coordinates of \eqref{eq:continuo-comun-estado}.  In particular,
\begin{equation}
\label{eq:continuo-comun-out-finito}
 1\leq d(x):=\#\operatorname{Out}(x)<\infty,
 \qquad
 \varepsilon^0_k(x)\in\operatorname{Out}(x).
\end{equation}

\begin{definition}[Truncated prolongation tree]
\label{def:continuo-comun-arbol-finito}
Given \(N\geq k\) and
\(x\in\widehat{\mathcal X}^{\rm tot}_k\), we define
\(\mathscr T_{k,N}(x)\) as the rooted tree whose vertices at height
\(j\), \(0\leq j\leq N-k\), are the words
\begin{equation}
\label{eq:continuo-comun-vertice-arbol}
 h=(\epsilon_{k+1},\ldots,\epsilon_{k+j}),
 \qquad
 \epsilon_{n+1}\in\operatorname{Out}(x_n),
 \qquad
 x_{n+1}=(\epsilon_{n+1})_*x_n,
\end{equation}
with the empty word as root. An edge appends an emission. Two
words reaching the same terminal state remain distinct vertices
if their ledgers differ.
\end{definition}

The last condition is essential: the tree records genealogies and not the
quotient that remembers only the target.  Its terminal level is
\begin{equation}
\label{eq:continuo-comun-terminal-finito}
 \operatorname{Term}_{k,N}(x)
 :=\{h\in\mathscr T_{k,N}(x):|h|=N-k\}.
\end{equation}
Each \(h\) preserves the complete terminal state.
\begin{equation}
\label{eq:continuo-comun-registro-terminal}
 \operatorname{ter}_{N}(h)
 :=\bigl(x_N,m_N,\partial(\gamma h),
          \gamma h,\operatorname{Led}(\gamma h)\bigr).
\end{equation}

\begin{proposition}[Finiteness and non-emptiness at every horizon]
\label{prop:continuo-comun-arbol-finito-no-vacio}
For all \(x\), \(k\), and \(N\geq k\), the tree
\(\mathscr T_{k,N}(x)\) is finite and
\(\operatorname{Term}_{k,N}(x)\neq\varnothing\).
\end{proposition}

\begin{proof}
The root has a finite positive number of successors by
\eqref{eq:continuo-comun-out-finito}.  Applying the same argument at each
level yields, by induction, finiteness at every finite horizon.  The word
formed by neutral emissions
\((\varepsilon^0_k,\varepsilon^0_{k+1},\ldots,
\varepsilon^0_{N-1})\) is a terminal vertex, so
level \eqref{eq:continuo-comun-terminal-finito} is nonempty.
\end{proof}

The erasure of the last emission defines
\begin{equation}
\label{eq:continuo-comun-proyeccion-terminal}
 \pi_{N+1,N}:\operatorname{Term}_{k,N+1}(x)
 \longrightarrow\operatorname{Term}_{k,N}(x).
\end{equation}

\begin{lemma}[Terminal surjectivity]
\label{lem:continuo-comun-terminal-sobreyectivo}
The map \(\pi_{N+1,N}\) is surjective for every \(N\geq k\).
\end{lemma}

\begin{proof}
Let \(h\in\operatorname{Term}_{k,N}(x)\), with terminal state \(v_h\).
Adding the emission \(\varepsilon_N^0(v_h)\) produces
\(h'\in\operatorname{Term}_{k,N+1}(x)\) and
\(\pi_{N+1,N}(h')=h\).
\end{proof}

\section{Nonadic Connection: Phase, Holonomy, and Memory}
\label{sec:continuo-comun-gamma9-total}

A word of nine emissions is not called a holonomy merely because it has
length nine. The horizontal connection determined by
the phase functor is first preserved. For \(r\in C_9\), we define
\begin{equation}
\label{eq:continuo-comun-relacion-horizontal-fase}
 \begin{split}
 \mathscr R_{r,k}:=\bigl\{
 (x,\epsilon,\epsilon_*x):\;&
 x\in\widehat{\mathcal X}^{\rm tot}_k,\quad
 \epsilon\in\operatorname{Out}(x),\\
 &\mathsf Q_k(x)=(w,r),\quad
 \mathsf Q_{k+1}(\epsilon_*x)=\sigma_9(w,r)
 \bigr\}.
 \end{split}
\end{equation}
The last equality does not restrict the domain: it is the total action
\eqref{eq:continuo-comun-accion-fase-tpk} of every elementary TPK emission. The connection simultaneously transports sheet, memory and ledger:
\begin{equation}
\label{eq:continuo-comun-conexion-tpk-generador}
 \nabla_{\rm TPK}(\epsilon):=
 \bigl(\operatorname{ph}(\epsilon),U_\epsilon,
       \mathsf U_\epsilon,\lambda(\epsilon)\bigr).
\end{equation}
Its composition is calculated by means of
\eqref{eq:continuo-comun-transporte-hojas-funtorial},
\eqref{eq:continuo-comun-accion-memoria-composicion}, and the concatenation of the
ledger.

The apex of the nonadic turn is now the set of phase-compatible horizontal lifts.
\begin{equation}
\label{eq:continuo-comun-gamma9-apice}
 \begin{split}
 \mathsf Z^{\Gamma}_{9,k}:=\coprod_{r\in C_9}\bigl\{
 (x,\epsilon_1,\ldots,\epsilon_9):\;&x_0=x,\quad
 x_i=(\epsilon_i)_*x_{i-1},\\
 &(x_{i-1},\epsilon_i,x_i)
 \in\mathscr R_{r+i-1,k+i-1},\quad 1\leq i\leq9
 \bigr\},
 \end{split}
\end{equation}
 where the subscripts of \(r+i-1\) are read in \(C_9\). The summand is
uniquely determined by \(r=g_k(x)\), so the coproduct does not
duplicate witnesses. The source and
target maps are
\begin{equation}
\label{eq:continuo-comun-gamma9-span}
 \widehat{\mathcal X}^{\rm tot}_k
 \xleftarrow{\ s_{9,k}\ }
 \mathsf Z^{\Gamma}_{9,k}
 \xrightarrow{\ t_{9,k}\ }
 \widehat{\mathcal X}^{\rm tot}_{k+9},
 \qquad
 s_{9,k}(x,\boldsymbol\epsilon)=x,\quad
 t_{9,k}(x,\boldsymbol\epsilon)
 =(\epsilon_9\cdots\epsilon_1)_*x.
\end{equation}
The relational holonomy is
\begin{equation}
\label{eq:continuo-comun-gamma9-relacion}
 \Gamma_{9,k}:=
 \{(s_{9,k}z,t_{9,k}z):z\in\mathsf Z^{\Gamma}_{9,k}\}
 =\operatorname{Hol}_{C_9}(\nabla_{\rm TPK}).
\end{equation}

\begin{lemma}[The ninth level of the tree realizes the holonomy]
\label{lem:continuo-comun-gamma9-arbol}
The map that forgets the intermediate states already determined by the emissions induces a bijection
\begin{equation}
\label{eq:continuo-comun-gamma9-apice-arbol}
 \mathsf Z^{\Gamma}_{9,k}
 \xrightarrow{\ \cong\ }
 \coprod_{x\in\widehat{\mathcal X}^{\rm tot}_k}
 \operatorname{Term}_{k,k+9}(x).
\end{equation}
\end{lemma}

\begin{proof}
Each edge of \(\operatorname{Out}(x_i)\) satisfies
\eqref{eq:continuo-comun-accion-fase-tpk}; hence a terminal word of
length nine determines the nine elements of
\eqref{eq:continuo-comun-relacion-horizontal-fase}. Reciprocally, a horizontal lift
of \eqref{eq:continuo-comun-gamma9-apice} is a composable word
of the tree. The two operations conserve the emissions and are
inverse. Thus \eqref{eq:continuo-comun-gamma9-apice-arbol} is a theorem
derived from the connection, not its definition.
\end{proof}

\begin{proposition}[Phase return, memory advance and no restart]
\label{prop:continuo-comun-gamma9-retorno-memoria}
For every \(z=(x,\epsilon_1,\ldots,\epsilon_9)\) of the apex,
\begin{align}
 g(t_{9,k}z)&=g(s_{9,k}z),
 \label{eq:continuo-comun-gamma9-retorno-fase}\\
 q^{(9)}(t_{9,k}z)&=q^{(9)}(s_{9,k}z)+1,
 \label{eq:continuo-comun-gamma9-avance-memoria}\\
 \operatorname{Led}(t_{9,k}z)
 &=\operatorname{Led}(s_{9,k}z)\star
   \bigl(\lambda(\epsilon_1),\ldots,\lambda(\epsilon_9)\bigr).
 \label{eq:continuo-comun-gamma9-ledger}
\end{align}
In particular, the phase identification
\(g(t_{9,k}z)=g(s_{9,k}z)\) does not induce an identification of states: after
transporting both registers to the same comparison type, their vertical
memories differ by one unit. The published phase returns, not the integral
holonomic state.
\end{proposition}

\begin{proof}
The first two identities are
\eqref{eq:continuo-comun-retorno-fase-avance-memoria}. The third is the
definition of ledger concatenation. Since the vertical memory differs by
one unit, the complete states are distinct even if some additional visible
coordinate could coincide.
\end{proof}

\begin{proposition}[Isometric compression of the same holonomy]
\label{prop:continuo-comun-gamma9-compresion-isometrica}
Let \(U_{\Gamma_9}:\mathcal H_{\rm full}\to\mathcal H_{\rm full}\) be an
isometric realization of the preceding holonomy and let
\(J:\mathcal H_{\rm vis}\to\mathcal H_{\rm full}\) be the isometric inclusion
of the visible register, with \(E:=J^*\) and \(P:=JE\). We define
\begin{equation}
\label{eq:continuo-comun-gamma9-compresion-datos}
 T:=EU_{\Gamma_9}J,\qquad
 \eta:=(I-P)U_{\Gamma_9}J.
\end{equation}
Then
\begin{equation}
\label{eq:continuo-comun-gamma9-compresion}
 \boxed{\;
 U_{\Gamma_9}J=JT+\eta,\qquad
 I-T^*T=\eta^*\eta.
 \;}
\end{equation}
\end{proposition}

\begin{proof}
Inserting \(I=P+(I-P)\) before \(U_{\Gamma_9}J\) gives
\[
 U_{\Gamma_9}J
 =JEU_{\Gamma_9}J+(I-P)U_{\Gamma_9}J
 =JT+\eta.
\]
The two summands belong to orthogonal subspaces. Since
\(U_{\Gamma_9}\) and \(J\) are isometries,
\[
 I=(U_{\Gamma_9}J)^*(U_{\Gamma_9}J)
   =T^*T+\eta^*\eta,
\]
which is the second identity.
\end{proof}

\begin{proposition}[Totality and Equivariance of Holonomy]
\label{prop:continuo-comun-gamma9-total-natural}
The map \(s_{9,k}\) is surjective. The nine neutral prolongations form
the section
\begin{equation}
\label{eq:continuo-comun-gamma9-seccion-neutra}
 \begin{gathered}
 x_0:=x,\qquad
 \varepsilon^0_{k+i}:=\varepsilon^0_{k+i}(x_i),\qquad
 x_{i+1}:=(\varepsilon^0_{k+i})_*x_i\quad(0\leq i<9),\\
 \eta^{(9)}_k(x):=
 (x,\varepsilon^0_k,\ldots,\varepsilon^0_{k+8})
 \in\mathsf Z^{\Gamma}_{9,k},
 \qquad s_{9,k}\eta^{(9)}_k=I.
 \end{gathered}
\end{equation}
Every coherent symmetry \(a\) whose action on the odometer satisfies
\(\mathsf Q\,a=\bar a\,\mathsf Q\) and
\(\bar a\,\sigma_9=\sigma_9\,\bar a\) induces
\[
 a^Z(x,\epsilon_1,\ldots,\epsilon_9)
 :=(a x,a\epsilon_1,\ldots,a\epsilon_9)
\]
and satisfies
\begin{equation}
\label{eq:continuo-comun-gamma9-naturalidad}
 s_{9,k}a^Z=a_ks_{9,k},
 \qquad
 t_{9,k}a^Z=a_{k+9}t_{9,k}.
\end{equation}
\end{proposition}

\begin{proof}
The neutral word exists in every source and satisfies the phase action \eqref{eq:continuo-comun-accion-fase-tpk}; this proves the section and surjectivity. The condition \(\bar a\sigma_9=\sigma_9\bar a\) sends horizontal relations to horizontal relations. Acting entry by entry preserves composition, carry, memory, boundary, and ledger; reading the first or the final state before or after acting gives \eqref{eq:continuo-comun-gamma9-naturalidad}.
\end{proof}

\section{Total survival and complete multisection}
\label{sec:continuo-comun-supervivencia}

\begin{definition}[Produced survival]
\label{def:continuo-comun-supervivencia}
A state \(x\in\widehat{\mathcal X}^{\rm tot}_k\) survives when
\begin{equation}
\label{eq:continuo-comun-supervivencia}
 \operatorname{Term}_{k,N}(x)\neq\varnothing
 \qquad\text{for every }N\geq k.
\end{equation}
We denote by \(\operatorname{Surv}_k\) the set of such states.
\end{definition}

\begin{theorem}[The total domain is survivor]
\label{thm:continuo-comun-total-superviviente}
For every depth,
\begin{equation}
\label{eq:continuo-comun-surv-igual-total}
 \operatorname{Surv}_k=\widehat{\mathcal X}^{\rm tot}_k.
\end{equation}
In particular, survival does not define a subset supposed before constructing the state.
\end{theorem}

\begin{proof}
The inclusion \(\operatorname{Surv}_k\subseteq
\widehat{\mathcal X}^{\rm tot}_k\) belongs to the definition.  The opposite
inclusion is exactly the nonemptiness at every horizon proved in
\cref{prop:continuo-comun-arbol-finito-no-vacio}.
\end{proof}

\begin{definition}[Multisection of complete prolongations]
\label{def:continuo-comun-multiseccion}
The complete multisection over \(x\) is the inverse limit
\begin{equation}
\label{eq:continuo-comun-multiseccion}
 \operatorname{Msect}(x)
 :=\varprojlim_{N\geq k}
 \bigl(\operatorname{Term}_{k,N}(x),\pi_{N+1,N}\bigr).
\end{equation}
Its elements are all the compatible infinite histories that prolong
\(x\).  None of them receives priority within
\eqref{eq:continuo-comun-multiseccion}.
\end{definition}

\begin{theorem}[No material emptiness of the multisection]
\label{thm:continuo-comun-multiseccion-no-vacia}
For every \(x\in\widehat{\mathcal X}^{\rm tot}_k\),
\begin{equation}
\label{eq:continuo-comun-multiseccion-no-vacia}
 \operatorname{Msect}(x)\neq\varnothing.
\end{equation}
Moreover, every finite coordinate of a history in
\(\operatorname{Msect}(x)\) preserves both leaves, orientation, carry,
memory, boundary, path, and ledger of the state that it prolongs.
\end{theorem}

\begin{proof}
The terminal sets are finite and nonempty by \cref{prop:continuo-comun-arbol-finito-no-vacio}, and their bonding maps are surjective by \cref{lem:continuo-comun-terminal-sobreyectivo}. A compatible element can be constructed recursively: one begins at the root and, at every step, retains a prolongation of the preceding prefix. The neutral prolongation guarantees that the process never stops. Compatibility of the coordinates follows from the fact that each edge of the tree acts through the unique state action defined in \eqref{eq:continuo-comun-estado}.
\end{proof}

\section{Canonical measure generated by the tree}
\label{sec:continuo-comun-medida}

The measure is not introduced as an external weight on the outputs.  It is obtained
by counting the admissible emissions at each vertex before choosing a
continuation.

\begin{definition}[Local weight and horizon measure]
\label{def:continuo-comun-medida-finita}
For \(v\in\widehat{\mathcal X}^{\rm tot}_j\) and
\(\epsilon\in\operatorname{Out}(v)\), we set
\begin{equation}
\label{eq:continuo-comun-peso-local}
 p_v(\epsilon):=\frac{1}{d(v)}\in\mathbb Q_{>0}.
\end{equation}
If
\(h=(\epsilon_{k+1},\ldots,\epsilon_N)\in
\operatorname{Term}_{k,N}(x)\) and \(x_j\) is the prefix that precedes
\(\epsilon_{j+1}\), we define
\begin{equation}
\label{eq:continuo-comun-medida-horizonte}
 \mu_{x,N}(h)
 :=\prod_{j=k}^{N-1}p_{x_j}(\epsilon_{j+1})
 =\prod_{j=k}^{N-1}\frac{1}{d(x_j)}.
\end{equation}
\end{definition}

\begin{lemma}[Normalization]
\label{lem:continuo-comun-medida-normalizada}
For every horizon \(N\geq k\),
\begin{equation}
\label{eq:continuo-comun-medida-suma-uno}
 \sum_{h\in\operatorname{Term}_{k,N}(x)}\mu_{x,N}(h)=1.
\end{equation}
\end{lemma}

\begin{proof}
At height one,
\(\sum_{\epsilon\in\operatorname{Out}(x)}d(x)^{-1}=1\).  Assuming the
normalization at height \(N-k\), each terminal prefix \(h\), with final state
\(v_h\), contributes at the next level
\[
 \sum_{\epsilon\in\operatorname{Out}(v_h)}
 \mu_{x,N}(h)p_{v_h}(\epsilon)
 =\mu_{x,N}(h).
\]
Summing over all prefixes proves the identity by induction.
\end{proof}

\begin{proposition}[Projective consistency]
\label{prop:continuo-comun-medida-consistente}
Finite measures satisfy
\begin{equation}
\label{eq:continuo-comun-medida-pushforward}
 (\pi_{N+1,N})_*\mu_{x,N+1}=\mu_{x,N}.
\end{equation}
\end{proposition}

\begin{proof}
For \(h\in\operatorname{Term}_{k,N}(x)\), the mass of its fiber is
\begin{align*}
 \sum_{\pi_{N+1,N}(h')=h}\mu_{x,N+1}(h')
 &=\mu_{x,N}(h)
   \sum_{\epsilon\in\operatorname{Out}(v_h)}p_{v_h}(\epsilon)\\
 &=\mu_{x,N}(h).
\end{align*}
This is precisely the equality of point measures equivalent to
\eqref{eq:continuo-comun-medida-pushforward}.
\end{proof}

The cylinders of the multisection are
\begin{equation}
\label{eq:continuo-comun-cilindro-msect}
 [h]_{x,N}
 :=\{\omega\in\operatorname{Msect}(x):
       \omega|_N=h\}.
\end{equation}

\begin{theorem}[Canonical Measure of Extensions]
\label{thm:continuo-comun-medida-canonica}
There exists a unique measure of cilindros \(\mu_x\) on
\(\operatorname{Msect}(x)\) such that
\begin{equation}
\label{eq:continuo-comun-medida-cilindro}
 \mu_x([h]_{x,N})=\mu_{x,N}(h).
\end{equation}
It is normalized, has support on the whole multisection, and is obtained only from
the finite sets of APP--TRIT--TPK prolongations.
\end{theorem}

\begin{proof}
The consistency of \cref{prop:continuo-comun-medida-consistente} makes \eqref{eq:continuo-comun-medida-cilindro} independent of the horizon at which a cylinder is represented. The normalization follows from \cref{lem:continuo-comun-medida-normalizada}. Two cylinders are either disjoint or one contains the other; hence the finite masses uniquely define an additive measure on the cylinder algebra. The multisection is a limit of finite discrete sets, and the cylinders form its compact basis; the measure extends uniquely to the \(\sigma\)-algebra they generate. Finally, \(p_v(\epsilon)>0\) for every admissible emission, so every nonempty cylinder has positive mass and the support is the entire multisection.
\end{proof}

\section{Naturality under symmetries and truncations}
\label{sec:continuo-comun-naturalidad}

A coherent family \(g=(g_j)_{j\geq0}\in\mathsf U\) transports an
outgoing emission of depth \(j\) by means of
\begin{equation}
\label{eq:continuo-comun-accion-out}
 (g_{j+1},g_j)_*:\operatorname{Out}(x)\longrightarrow
      \operatorname{Out}(g_j\cdot x),
 \qquad
 \epsilon\longmapsto g_{j+1}\epsilon g_j^{-1}.
\end{equation}
The TPK relations make \eqref{eq:continuo-comun-accion-out} a bijection
that preserves multiplicities.  Acting on each entry yields a tree
isomorphism
\begin{equation}
\label{eq:continuo-comun-isomorfismo-arbol}
 g_{*,N}:\mathscr T_{k,N}(x)
 \xrightarrow{\ \cong\ }
 \mathscr T_{k,N}(g_k\cdot x).
\end{equation}

\begin{theorem}[Naturalness of the tree, the multisection, and the measure]
\label{thm:continuo-comun-naturalidad}
For every coherent reversible symmetry \(g=(g_j)_{j\geq0}\),
\begin{align}
 \pi_{N+1,N}g_{*,N+1}
 &=g_{*,N}\pi_{N+1,N},
 \label{eq:continuo-comun-naturalidad-terminal}\\
 g_*\operatorname{Msect}(x)
 &=\operatorname{Msect}(g_k\cdot x),
 \label{eq:continuo-comun-naturalidad-msect}\\
 (g_{*,N})_*\mu_{x,N}
 &=\mu_{g_k\cdot x,N},
 \qquad
 (g_*)_*\mu_x=\mu_{g_k\cdot x}.
 \label{eq:continuo-comun-naturalidad-medida}
\end{align}
\end{theorem}

\begin{proof}
The first equality states that entrywise action commutes with deleting the
final entry. Passing to the inverse limit yields
\eqref{eq:continuo-comun-naturalidad-msect}. Since
\eqref{eq:continuo-comun-accion-out} is bijective,
\(d(g_j\cdot v)=d(v)\); by
\eqref{eq:continuo-comun-medida-horizonte}, a word and its image have the
same weight. This proves the first equality of measures in
\eqref{eq:continuo-comun-naturalidad-medida}. The second follows on the
cylinders from \eqref{eq:continuo-comun-medida-cilindro}, and the cylinders
determine the measure.
\end{proof}

The truncations of the state and those of the tree also form the square.
\begin{equation}
\label{eq:continuo-comun-cuadrado-truncamiento}
\begin{CD}
 \operatorname{Term}_{k,N+1}(x)
 @>{\pi_{N+1,N}}>>
 \operatorname{Term}_{k,N}(x)\\
 @V{\operatorname{ter}_{N+1}}VV
 @VV{\operatorname{ter}_{N}}V\\
 \widehat{\mathcal X}^{\rm tot}_{N+1}
 @>{\rho_{N+1,N}}>>
 \widehat{\mathcal X}^{\rm tot}_{N}.
\end{CD}
\end{equation}
Commutativity is literal: both paths erase the last emission and its
last ledger entry.

\section{Inverse limit \(\widehat\Omega\) of complete histories and common APP–TRIT–TPK integral holonomic state}
\label{sec:continuo-comun-limite}

The surjections \eqref{eq:continuo-comun-truncamiento-total} produce the
object of complete histories
\begin{equation}
\label{eq:continuo-comun-omega}
 \widehat\Omega
 :=\varprojlim_k
 \bigl(\widehat{\mathcal X}^{\rm tot}_k,\rho_{k+1,k}\bigr).
\end{equation}
An element \(\widehat x=(\widehat x_k)_{k\geq0}\) contains, at each
depth, the same pair of transported sheets, the produced orientations and
carries, the accumulated memory, the boundaries, the route, and the complete
ledger.  Its fiber of continuations from any prefix is
\(\operatorname{Msect}(\widehat x_k)\), equipped with
\(\mu_{\widehat x_k}\).

\begin{theorem}[Common APP--TRIT--TPK integral holonomic state]
\label{thm:continuo-comun-fuente teoremática}
The system
\begin{equation}
\label{eq:continuo-comun-fuente teoremática}
 \mathfrak C^{\rm gen}:=\Biggl(
 \begin{aligned}
 &\{\widehat{\mathcal X}^{\rm tot}_k,\rho_{k+1,k},\eta_k\}_{k\geq0},
   \widehat\Omega,\{\mathfrak G_k\}_{k\geq0},
   \{\mathscr T_{k,N}\}_{N\geq k},\\
 &\operatorname{Msect},\operatorname{Surv},\{\mu_x\},
   \{\mathfrak r_9(v_K)\}_{K\geq0},
   \mathcal K_{\rm ph},\widehat{\mathcal K}_{\rm ph},\\
 &\operatorname{Sh},\operatorname{or},\kappa,
   \mathsf Q,g,q^{(9)},\beta,m,\partial,\gamma,\operatorname{Led},\\
 &\{\mathsf Z^\Gamma_{9,k},s_{9,k},t_{9,k},\Gamma_{9,k}\}_{k\geq0}
 \end{aligned}
 \Biggr)
\end{equation}
is defined for every state produced by APP--TRIT--TPK\@.  It is non-empty,
natural under reversible symmetries, and compatible with refinement.  In
particular:
\begin{enumerate}
 \item the two leaves are transported together;
 \item orientation and carry it calculan by the division TRIT;
 \item the memory obeys the affine TPK action;
 \item the phase obeys the TPK odometer, returns after nine emissions, and the
       vertical memory advances by one unit;
 \item the Markov realization of that turn satisfies
       \(\mathcal K_{\rm ph}^9=E_+E_\times\) and uniformly renews the complete fiber of \(468=18\cdot26\) emissions;
 \item the compact reader conserves block, cylinder, prefix, carry, signature,
       Witt shadow, and phase, with their updates and truncations;
 \item boundary, path, and ledger are composed without erasing the endpoints or the
       order of the emissions;
 \item every state survives and has a nonempty multisection of complete prolongations;
 \item the canonical measure is projective and has support on all those
       prolongations.
\end{enumerate}
\end{theorem}

\begin{proof}
Totality and finiteness follow from \cref{prop:continuo-comun-total-finito}. The sheet, carry, and memory identities are proved in \cref{lem:continuo-comun-asociatividad-carry,lem:continuo-comun-accion-memoria}; the phase, vertical memory, and nonadic holonomy, including its fiber realization and its state reader, are proved in \cref{prop:continuo-comun-gamma9-retorno-memoria,prop:continuo-comun-gamma9-total-natural,prop:continuo-comun-kph-heredada,prop:continuo-comun-gamma9-compresion-isometrica,prop:continuo-comun-no-reinicio-alcance}; the boundary preserves its endpoints by \eqref{eq:continuo-comun-frontera-telescopica}. Total survival and nonemptiness of the multisections are \cref{thm:continuo-comun-total-superviviente,thm:continuo-comun-multiseccion-no-vacia}. The existence, consistency, and support of the measure are \cref{prop:continuo-comun-medida-consistente,thm:continuo-comun-medida-canonica}. Finally, compatibility with the symmetries and truncations is proved in \cref{thm:continuo-comun-naturalidad} and \eqref{eq:continuo-comun-cuadrado-truncamiento}.
\end{proof}

\section{Structural necessity and origin}
\label{sec:continuo-comun-falsadores}

\begin{proposition}[Falsifiers of the common integral holonomic state]
\label{prop:continuo-comun-falsadores}
Any of the following operations replaces
\(\mathfrak C^{\rm gen}\) with another object and blocks the transfer of the
preceding theorems:
\begin{enumerate}
 \item elegir \(\Sigma\) or \(\Pi\) and descartar the another sheet;
 \item conserving the residue and erasing the quotient;
 \item preserving the TRIT sign and erasing its carry;
 \item replacing the ordered ledger with its terminal value;
 \item identifying two words because they reach the same cell;
 \item choosing a continuation and replacing the multisection with it;
 \item declaring survival without exhibiting prolongations to each horizon;
 \item uniformly normalizing only the last level, ignoring the degrees of
       the prefixes and losing projective consistency;
 \item using a subsequent reading to decide which states enter
       \(\widehat{\mathcal X}^{\rm tot}_k\).
\end{enumerate}
\end{proposition}

\begin{proof}
The first six points eliminate coordinates or relations that appear in the
definition of the state, tree, or multisection. The seventh eliminates the proof of
\cref{thm:continuo-comun-total-superviviente}. The eighth may violate
\eqref{eq:continuo-comun-medida-pushforward} when the terminal degrees are not
constant. The last replaces the total domain
\eqref{eq:continuo-comun-dominio-total} with a retrospective selection.
\end{proof}

The APP cells, the pair of sheets, the TRIT lift, the TPK extensions,
the carry cocycle, memory, boundary, path and conservation
of the register constitute the starting structures. On them, the finite groupoid, the word tree, the multisection and
the normalized local measure are brought together in a single domain. The theorems of total survival, nonemptiness,
consistency and naturality certify this joint construction.
 \let\section\HMTContinuoCommonOriginalSection
\chapter{Intrinsic Correlative Operations of the Same Continuum}
\label{chap:pdfv2-operaciones-intrinsecas-correlativas}

\begin{thesisbox}
The symbols \(\mathrm{sp},\mathrm{sol},\mathrm{gau},\mathrm{coh}\) and
\(\mathrm{det}\) that appear in this chapter mark five correlative operations
over a single APP--TRIT--TPK state\@. Each operation receives
the same surviving extension, the
same orientation, the same carry, the same memory, the same boundary, and the
same route. The following sections are successive mathematical stages of a
single construction and not autonomous expositions.
\end{thesisbox}

\section{Generating edge and common genealogical state}

We do not repeat here the state, tree, or measure constructed in the preceding
chapter. We take
\(\widehat x=\widehat x_k(\gamma)\in
\widehat{\mathcal X}^{\rm tot}_k\) from
\eqref{eq:continuo-comun-estado} and all its admissible emissions
\(\alpha\in\operatorname{Out}(\widehat x)\). If
\(\widehat y=\alpha_*\widehat x\), we write
\begin{equation}
\label{eq:pdfv2-cor-hijos-supervivientes}
 \operatorname{Ch}_k(\widehat x):=
 \{(\alpha,\alpha_*\widehat x):
      \alpha\in\operatorname{Out}(\widehat x)\},
 \qquad
 d_k(\widehat x):=\#\operatorname{Ch}_k(\widehat x).
\end{equation}
The finiteness and non-emptiness of this set are already
\eqref{eq:continuo-comun-out-finito}; survival at all horizons,
the complete multisection, and its measure are
\cref{thm:continuo-comun-total-superviviente,thm:continuo-comun-multiseccion-no-vacia,thm:continuo-comun-medida-canonica}. A second non-emptiness is not again assumed or proved.

Each pair \((\alpha,\widehat y)\) of
\eqref{eq:pdfv2-cor-hijos-supervivientes} preserves the unique record
\begin{equation}
\label{eq:pdfv2-cor-registro-extension}
 \begin{aligned}
 \mathfrak e_\alpha:=\bigl(&
 \alpha,\operatorname{Ext}_{\gamma_\alpha};
 \operatorname{Sh}_{s(\alpha)},\operatorname{Sh}_{t(\alpha)};\\
 &
 \operatorname{res}^{\Sigma,\Pi}(\alpha),
 \operatorname{quo}^{\Sigma,\Pi}(\alpha);\\
 &\operatorname{or}(\alpha),\kappa(\alpha),\mathsf C_\alpha;
 \mathsf Q(s(\alpha)),\mathsf Q(t(\alpha)),g(s(\alpha)),g(t(\alpha)),
 q^{(9)}(s(\alpha)),q^{(9)}(t(\alpha));\\
 &
 m(\gamma\alpha;m_0),\partial\alpha,\gamma\alpha,
 \operatorname{Led}(\gamma\alpha);\\
 &\{\operatorname{Term}_{k+1,N}(\widehat y)\}_{N\geq k+1},
 \operatorname{Msect}(\widehat y),\mu_{\widehat y}\bigr).
 \end{aligned}
\end{equation}
Here \(\operatorname{Ext}_{\gamma_\alpha}\) denotes the complete relation that
the admissible emission already satisfies; it introduces no additional filter. The subsequent
operations are coordinated applications of \(\mathfrak e_\alpha\) and do not
reconstruct the record from their outputs.

\section{Intrinsic Correlative Actions on the Common Generator}

The internal marks are fixed without changing their source: \(\mathrm{sp}\) denotes transport and grading of the two sheets; \(\mathrm{sol}\), signed increment, balance, and memory; \(\mathrm{gau}\), incidence and carry curvature; \(\mathrm{coh}\), ports, rectangular difference, and cellular complex; and \(\mathrm{det}\), based complex, correlative filling, and determinant line. The following formulas display their couplings over the same \(\mathfrak e_\alpha\); their actions literally share source, target, route, and ledger.

The double APP evaluation provides, on the same support, the labels
\(\Sigma\) and \(\Pi\). For an extension \(\alpha\), their multiplicities
are calculated from the complete ledger, without an external counter:
\begin{equation}
\label{eq:pdfv2-cor-multiplicidades-hojas}
 N_\diamond(\alpha):=
 \sum_{\epsilon\ {\rm en}\ \gamma\alpha}
 \mathbf1_{\{\operatorname{res}^{\diamond}(\epsilon)\neq0\}},
 \qquad \diamond\in\{\Sigma,\Pi\}.
\end{equation}
Let there also be the total oriented coefficient
\[
 \Delta_{\gamma_\alpha}:=
 \sum_{\epsilon\ {\rm en}\ \gamma\alpha}\operatorname{or}(\epsilon)
 \in\mathbb Z.
\]
The signed action is
\begin{equation}
\label{eq:pdfv2-cor-b-alpha}
 b_\alpha:=\Delta_{\gamma_\alpha}
 \bigl(N_\Sigma(\alpha)-N_\Pi(\alpha)\bigr),
 \qquad
 b_\alpha^\pm:=\frac{|b_\alpha|\pm b_\alpha}{2}.
\end{equation}
By construction,
\begin{equation}
\label{eq:pdfv2-cor-b-identidades}
 b_\alpha=b_\alpha^+-b_\alpha^-,\qquad
 |b_\alpha|=b_\alpha^++b_\alpha^-,\qquad
 b_\alpha^+b_\alpha^-=0.
\end{equation}

The same set of extensions simultaneously produces the ports
\begin{equation}
\label{eq:pdfv2-cor-puertos-comunes}
 I_k(\widehat x_k):=
 \operatorname{Ch}_k(\widehat x_k)\times\{\Sigma\},
 \qquad
 J_k(\widehat x_k):=
 \operatorname{Ch}_k(\widehat x_k)\times\{\Pi\}.
\end{equation}
A visible leaf is not assumed to contain more prolongations than the other:
each surviving extension is evaluated on both leaves of the same support.
By \eqref{eq:continuo-comun-out-finito}, both sets of ports are
nonempty.

For \(A_N=\mathbb Z/3^N\mathbb Z\), the cellular action on generators is
\begin{equation}
\label{eq:pdfv2-cor-kappa-celda}
 \begin{aligned}
 \kappa_{k,N}:A_N^{I_k}\oplus A_N^{J_k}
 &\longrightarrow A_N^{I_k\times J_k},\\
 \kappa_{k,N}(u,v)_{(\alpha,\Sigma),(\beta,\Pi)}
 &:=u_{(\alpha,\Sigma)}-v_{(\beta,\Pi)}.
 \end{aligned}
\end{equation}
Selected first ports according to the TPK order, the rectangular difference
\begin{equation}
\label{eq:pdfv2-cor-delta-celda}
 \delta(w)_{ij}:=w_{ij}-w_{ij_0}-w_{i_0j}+w_{i_0j_0}
\end{equation}
download the split sequence
\begin{equation}
\label{eq:pdfv2-cor-sucesion-celda}
 0\longrightarrow A_N\xrightarrow{\Delta}
 A_N^{I_k}\oplus A_N^{J_k}
 \xrightarrow{\kappa_{k,N}}A_N^{I_k\times J_k}
 \xrightarrow{\delta}A_N^{(|I_k|-1)(|J_k|-1)}
 \longrightarrow0.
\end{equation}
This is an identity of the same APP fibre: it adds no cell foreign to the
generator \(\mathfrak e_\alpha\).

The local cellular action uses the two integer coordinates accumulated by that
same ledger. For
\(\widehat x=\widehat x_k(\gamma)\) and
\(\widehat y=\alpha_*\widehat x=\widehat x_{k+1}(\gamma\alpha)\), we set
\begin{equation}
\label{eq:pdfv2-cor-c-v}
 \begin{aligned}
 c_{\widehat x}&:=\kappa(\gamma), &
 c_\alpha:=c_{\widehat y}&:=\kappa(\gamma\alpha),\\
 v_{\widehat x}&:=\sum_{\epsilon\ {\rm en}\ \gamma}
                  \operatorname{or}(\epsilon), &
 v_{\widehat y}&:=\sum_{\epsilon\ {\rm en}\ \gamma\alpha}
                  \operatorname{or}(\epsilon),\\
 v_\alpha&:=v_{\widehat y}-v_{\widehat x}
           =\operatorname{or}(\alpha).&&
 \end{aligned}
\end{equation}
On the generators \(f,r,e,b,g\) is defined
\begin{equation}
\label{eq:pdfv2-cor-complejo-local}
 \partial f=3c_\alpha e+b,\qquad
 \partial r=0,\qquad
 \partial e=g,\qquad
 \partial b=-3c_\alpha g,\qquad
 \partial g=0.
\end{equation}
The cancellation
\begin{equation}
\label{eq:pdfv2-cor-dos-cero}
 \partial^2f=3c_\alpha g-3c_\alpha g=0
\end{equation}
is literal and uses the same carry with the two signs imposed by the
orientation.

The incidence seed is obtained from the free module of the two leaves,
\begin{equation}
\label{eq:pdfv2-cor-v-tpk}
 V_{\rm TPK}:=\mathbf F_3e_\Sigma\oplus\mathbf F_3e_\Pi,
 \qquad
 \mathscr I:=\mathbf P^1(V_{\rm TPK}),
 \qquad |\mathscr I|=4,
\end{equation}
and of its unordered pairs
\begin{equation}
\label{eq:pdfv2-cor-incidencias}
 \mathscr E:=\binom{\mathscr I}{2},
 \qquad |\mathscr E|=6,
 \qquad
 \mathscr B:=\mathscr I\times\{+,-\},
 \qquad |\mathscr B|=8.
\end{equation}
The incidence application is determined on the base by
\begin{equation}
\label{eq:pdfv2-cor-b-incidencia-generadores}
 B[e_L]=\sum_{L'\neq L}e_{\{L,L'\}},
 \qquad L\in\mathscr I,
\end{equation}
and, by linearity,
\begin{equation}
\label{eq:pdfv2-cor-b-incidencia-coordenadas}
 (Ba)_{\{L,L'\}}=a_L+a_{L'}.
\end{equation}
If \(s(q)=\sum_{E\in\mathscr E}q_E\), each \(L\) appears three times and
\begin{equation}
\label{eq:pdfv2-cor-sb-cero}
 s(Ba)=0.
\end{equation}
Therefore the function
\begin{equation}
\label{eq:pdfv2-cor-wc}
 W_c(q):=\mathbf1_{\{s(q)=0\}}-\frac13
\end{equation}
satisfies \(W_c(q+Ba)=W_c(q)\). The identities
\eqref{eq:pdfv2-cor-v-tpk}--\eqref{eq:pdfv2-cor-wc} are internal and exact.
The material lift of \(q\) from the cells and common carries is treated in
the assembly layer; this incidence identity is not conflated with the
construction of that lift.

\section{Unitary transport, polarization, and balance on the common fiber}

The non-emptiness of \eqref{eq:pdfv2-cor-hijos-supervivientes} allows one to fix,
without selecting a prolongation,
\begin{equation}
\label{eq:pdfv2-cor-peso-conteo}
 p_k(\alpha,\widehat y\mid\widehat x)
 :=\frac1{d_k(\widehat x)}.
\end{equation}
Let the space of the two leaves of the same state be
\begin{equation}
\label{eq:pdfv2-cor-hilbert-hojas}
 \mathscr H_k:=
 \bigoplus_{\widehat x\in\widehat{\mathcal X}^{\rm tot}_k}
 \bigl(\mathsf H_{\widehat x}^{\Sigma}
       \oplus\mathsf H_{\widehat x}^{\Pi}\bigr),
\end{equation}
and let \(\iota_{\widehat x}\) be the inclusion of the indicated fiber. The TPK
transport of \eqref{eq:continuo-comun-transporte-hojas} gives the exact action
\begin{equation}
\label{eq:pdfv2-cor-u-comun}
 U_k\iota_{\widehat x}v:=
 \sum_{(\alpha,\widehat y)\in\operatorname{Ch}_k(\widehat x)}
 p_k(\alpha,\widehat y\mid\widehat x)^{1/2}
 \iota_{\widehat y}U_\alpha v.
\end{equation}
The children of two distinct parents are disjoint because truncation is a
function. Moreover, two distinct witnesses \(\alpha\) do not collapse into the same
direct summand: the complete route \(\gamma\), including its last step, is a
coordinate of \(\widehat y\). Consequently, the Gram identity is
\begin{equation}
\label{eq:pdfv2-cor-u-gram}
 \left.U_k^*U_k\right|_{\mathsf H_{\widehat x}^{\Sigma}
       \oplus\mathsf H_{\widehat x}^{\Pi}}
 =
 \sum_{(\alpha,\widehat y)\in\operatorname{Ch}_k(\widehat x)}
 p_k(\alpha,\widehat y\mid\widehat x)\,U_\alpha^*U_\alpha.
\end{equation}
Therefore, \(U_k^*U_k=I\) exactly when the sheet transports
\(U_\alpha\) published by TPK are isometries. Normalization of the weights
alone cannot replace that verification.
For \(\ell>k\), the complete history is preserved through
\begin{equation}
\label{eq:pdfv2-cor-u-iterado}
 U_{\ell,k}:=U_{\ell-1}\cdots U_k,
 \qquad U_{k,k}:=I.
\end{equation}

The grading acts on the two coordinates already present in each fiber of \eqref{eq:pdfv2-cor-hilbert-hojas}:
\begin{equation}
\label{eq:pdfv2-cor-sigma}
 \sigma_k\iota_{\widehat x}(v_\Sigma,v_\Pi)
 :=\iota_{\widehat x}(v_\Sigma,-v_\Pi),
\qquad
 P_k:=\frac{I+\sigma_k}{2},\quad
 Q_k:=\frac{I-\sigma_k}{2}.
\end{equation}
The TRIT threshold belongs to the orientation and to the carry of
\(\widehat x\); it erases neither of the two APP sheets and does not require a third
case in \eqref{eq:pdfv2-cor-sigma}.

The leaf-preserving part and the leaf-changing part are two components of
the same \(U_k\):
\begin{align}
 A_k&:=P_{k+1}U_kP_k+Q_{k+1}U_kQ_k,
 \label{eq:pdfv2-cor-a-comun}\\
 \eta_k&:=Q_{k+1}U_kP_k+P_{k+1}U_kQ_k.
 \label{eq:pdfv2-cor-eta-comun}
\end{align}
By orthogonality of \(P_{k+1}\) and \(Q_{k+1}\),
\begin{equation}
\label{eq:pdfv2-cor-balance-hojas}
 U_k=A_k+\eta_k,\qquad
 A_k^*A_k+\eta_k^*\eta_k=U_k^*U_k.
\end{equation}
The last member is \(I\) under, and only under, the isometric condition
identified after \eqref{eq:pdfv2-cor-u-gram}.
This identity and the increment \eqref{eq:pdfv2-cor-b-alpha} are not two
transports: \(A_k,\eta_k\) act on the base of histories and
\(b_\alpha\) reads the edge of that same base.

\subsection{Signed polarization, non-autonomous telescoping, and return without reinitialization}
No independent spectral branch is introduced here. The notations are
restricted to the same space \(\mathscr H_k\) and the same already constructed
transport \(U_k\)
\[
 \mathcal H_k^{\mathrm{sp}}:=\mathscr H_k,
 \qquad
 \Lambda(\gamma_{\widehat x})
 :=\bigl(\operatorname{or},\operatorname{Carr},\operatorname{Mem},
          \partial,\gamma\bigr)(\widehat x),
\]
and \(\Sigma_{\rm reg}\Lambda(\gamma_{\widehat x})\) denotes its family of
compatible truncations. In this way, the following signed reader, defect and
terminal are actions of the common correlative constructor, not a
retrospective selection of histories.
\begingroup
  \let\section\subsection
  % Extracción quirúrgica del fuente teoremática V2 05a: líneas 293--394.
% El resto permanece en este archivo como procedencia, pero no se publica.
\iffalse
\chapter{Spectral macrostructure--internal polar continuum}
\label{chap:continuo-macro-sp-interna}

\section{Strictly internal scope}

The construction in this chapter terminates before any external recognition. Its sole data are the two-sheet APP support, the TRIT orientation, the TPK extensions and carry, the projective system of histories, the cylindrical measure, and the signed ledger. The result is a graded macrostructure with transport, a sheet-change expression, a signed reader, and nonautonomous telescoping. No external function is used to define any of its maps.

\section{Surviving domain of the branch}

For \(x_k\in\mathcal A_k\), let
\begin{equation}
\label{eq:continuo-sp-msect-n}
 \operatorname{Msect}^{\mathrm{sp},(N)}_k(x_k)
 =\left\{(x_j)_{j=k}^{N}:
 \tau_{j+1,j}x_{j+1}=x_j,
 \ (x_j,x_{j+9})\in\Gamma_{9,j},
 \ \text{and the sheaf is preserved }\{\Sigma,\Pi\}\right\}.
\end{equation}
We define
\begin{equation}
\label{eq:continuo-sp-dominio-superviviente}
 \mathcal C_k^{\mathrm{sp}}
 =\left\{x_k:
 \operatorname{Msect}^{\mathrm{sp},(N)}_k(x_k)\neq\varnothing
 \text{ for every }N\geq k\right\},
 \qquad
 \mathcal C_{\rm cont}^{\mathrm{sp}}
 =\varprojlim_k\mathcal C_k^{\mathrm{sp}}.
\end{equation}
The TRIT threshold is not excluded from this domain. The bundle preserves the two
arithmetic sheets even if a visible chart publishes neither of them in a
specific event.

\section{Thirteen-coordinate package}

For \(x_k\in\mathcal C_k^{\mathrm{sp}}\), the dependent package is
\begin{equation}
\label{eq:continuo-sp-d13}
\begin{aligned}
 D_{\mathrm{sp},k}(x_k)=\bigl(&
 \mathcal F_{\mathrm{sp},k}(x_k),
 \operatorname{Ext}_{\Gamma,\mathrm{sp},k}(x_k),
 \operatorname{Rel}_{\mathrm{sp},k}(x_k),
 \mathcal N_{\mathrm{sp},k}(x_k),
 \operatorname{or}_{\mathrm{sp},k}(x_k),\\
 &\operatorname{Carr}_{\mathrm{sp},k}(x_k),
 \operatorname{Mem}_{\mathrm{sp},k}(x_k),
 \partial_{\mathrm{sp},k}(x_k),
 \gamma_{\mathrm{sp},k}(x_k),
 \operatorname{Msect}_{\mathrm{sp},k}(x_k),\\
 &\operatorname{Surv}_{\mathrm{sp},k}(x_k),
 \operatorname{Term}_{\mathrm{sp},k}(x_k),
 \mathcal L^{\rm sgn}_{\mathrm{sp},k}(x_k)\bigr).
\end{aligned}
\end{equation}
Their components are determined correlatively as follows.

The fiber \(\mathcal F_{\mathrm{sp},k}\) contains the TPK state and the ordered
bundle \(\{\Sigma,\Pi\}\). The extension is the complete restriction of
\(\operatorname{Ext}_\Gamma\) to histories satisfying
\eqref{eq:continuo-sp-dominio-superviviente}. The relations contain the
identity and composition laws, the residue--quotient lift, the carry
cocycle, boundary compatibility and truncation naturality. The
numerical coordinate contains depth, phase, residues, quotients and internal
signature. Orientation stores sheet and TRIT orientation separately.
Memory contains the quotient, the ledger and the turn increment. The
boundary and route remain uncontracted. The multisection is the complete
system of prolongations of \eqref{eq:continuo-sp-msect-n}.
Survival is the stable membership of
\eqref{eq:continuo-sp-dominio-superviviente}. The terminal and reader are
constructed in the following sections.

The truncations of all coordinates form a map
\(u^{\mathrm{sp}}_{\ell,k}\) satisfying
\begin{equation}
\label{eq:continuo-sp-d-naturalidad}
 u^{\mathrm{sp}}_{\ell,k}D_{\mathrm{sp},\ell}
 =D_{\mathrm{sp},k}\tau_{\ell,k}.
\end{equation}

\section{Graph of the constraint and dependent section}

We define
\begin{equation}
\label{eq:continuo-sp-g-graph}
 \mathcal G_{\mathrm{sp},k}
 =\operatorname{Graph}(D_{\mathrm{sp},k})
 =\{(x_k,D_{\mathrm{sp},k}x_k):
     x_k\in\mathcal C_k^{\mathrm{sp}}\}.
\end{equation}
The total restriction is
\begin{equation}
\label{eq:continuo-sp-res-graph}
 \operatorname{Res}_{\mathrm{sp}}=s_{\mathrm{sp}}:
 \mathcal C_{\rm cont}^{\mathrm{sp}}
 \longrightarrow\mathcal G_{\mathrm{sp}},
 \qquad
 s_{\mathrm{sp}}(x)=(x,D_{\mathrm{sp}}x).
\end{equation}
The projection \(\pi_1(x,d)=x\) verifies
\begin{equation}
\label{eq:continuo-sp-res-inversa}
 \pi_1s_{\mathrm{sp}}=I,
 \qquad
 s_{\mathrm{sp}}\pi_1=I_{\mathcal G_{\mathrm{sp}}}.
\end{equation}

Let \(r_{\mathrm{sp},k}(x_k)\) be the multisection of sheets and paths that already
belongs to \(D_{\mathrm{sp},k}(x_k)\). Define
\begin{equation}
\label{eq:continuo-sp-x-dependiente}
 X_{\chi,\mathrm{sp},k}
 =\{(r_{\mathrm{sp},k}x_k,x_k,D_{\mathrm{sp},k}x_k):
     x_k\in\mathcal C_k^{\mathrm{sp}}\},
\end{equation}
\begin{equation}
\label{eq:continuo-sp-pr-x}
 \operatorname{pr}_{X,\mathrm{sp}}(x,D_{\mathrm{sp}}x)
 =(r_{\mathrm{sp}}x,x,D_{\mathrm{sp}}x).
\end{equation}
Forgetting the first coordinate of
\eqref{eq:continuo-sp-x-dependiente} is the inverse of
\(\operatorname{pr}_{X,\mathrm{sp}}\).

\section{Bundle of leaves and projective measure}

Let
\begin{equation}
\label{eq:continuo-sp-haz-hojas}
 \widetilde{\mathcal C}_k^{\mathrm{sp}}
 =\{(x_k,h):x_k\in\mathcal C_k^{\mathrm{sp}},
              h\in\{\Sigma,\Pi\}\}.
\end{equation}
The projective measure of the continuum induces probabilities
\(\widetilde\mu_k\) on this bundle. We restrict to the positive support; for
\(z\) in that support,
\begin{equation}
\label{eq:continuo-sp-consistencia-medida}
 \widetilde\mu_k(z)
 =\sum_{\widetilde\tau_{k+1,k}(z')=z}
  \widetilde\mu_{k+1}(z').
\end{equation}
Level linearization is
\begin{equation}
\label{eq:continuo-sp-hilbert-nivel}
 \mathcal H_k^{\mathrm{sp}}
 =\ell^2(\operatorname{supp}\widetilde\mu_k),
\end{equation}
with orthonormal basis \(e_z\).

\section{Transport of all extensions and carry}

For \(z'\) extending \(z\), we define the positive conditional weight
\begin{equation}
\label{eq:continuo-sp-peso-condicional}
 w_k(z'\mid z)
 =\left(\frac{\widetilde\mu_{k+1}(z')}
               {\widetilde\mu_k(z)}\right)^{1/2}.
\end{equation}
Transport is
\begin{equation}
\label{eq:continuo-sp-u-accion}
 U_ke_z
 =\sum_{\substack{z':\widetilde\tau_{k+1,k}(z')=z\\
            (z,z')\in\operatorname{Ext}_{\Gamma,\mathrm{sp},k}}}
 w_k(z'\mid z)e_{z'}.
\end{equation}
The sum contains all admissible extensions. It does not depend on a selector.

By \eqref{eq:continuo-sp-consistencia-medida},
\[
 \sum_{z':\widetilde\tau(z')=z}|w_k(z'\mid z)|^2=1.
\]
The sets of descendants of two distinct parents are disjoint. Consequently,
\begin{equation}
\label{eq:continuo-sp-u-isometria}
 U_k^*U_k=I_{\mathcal H_k^{\mathrm{sp}}}.
\end{equation}

Each basis \(e_{z'}\) includes the carry of the corresponding extension. For
composable paths, we preserve
\begin{align}
 \kappa(\gamma_2\gamma_1)
 &=\kappa(\gamma_2)+
   \mathsf C(\gamma_2)\kappa(\gamma_1),
 \label{eq:continuo-sp-carry-nativo}\\
 \operatorname{Carr}(\gamma_2\gamma_1)
 &=\operatorname{Carr}(\gamma_2)H(\gamma_1)
  +Y(\gamma_2)\operatorname{Carr}(\gamma_1).
 \label{eq:continuo-sp-carry-operador}
\end{align}
The conditional weights multiply along the unique chain of
truncations of a descendant. Equations
\eqref{eq:continuo-sp-carry-nativo}--
\eqref{eq:continuo-sp-carry-operador} ensure that the carry labels
coincide. Therefore,
\begin{equation}
\label{eq:continuo-sp-u-composicion}
 U_{\ell,k}=U_{\ell-1}\cdots U_k.
\end{equation}

\section{Graduation, projectors and decomposition of transport}

The sheet graduation is the involution.
\begin{equation}
\label{eq:continuo-sp-sigma}
 \sigma_ke_{(x,\Sigma)}=+e_{(x,\Sigma)},
 \qquad
 \sigma_ke_{(x,\Pi)}=-e_{(x,\Pi)}.
\end{equation}
Thus,
\[
 \sigma_k^*=\sigma_k,
 \qquad
 \sigma_k^2=I.
\]
The internal projectors are
\begin{equation}
\label{eq:continuo-sp-p-q}
 P_k=\frac{I+\sigma_k}{2},
 \qquad
 Q_k=\frac{I-\sigma_k}{2}.
\end{equation}
Satisfy
\begin{equation}
\label{eq:continuo-sp-p-q-relaciones}
 P_k^2=P_k,
 \quad Q_k^2=Q_k,
 \quad P_kQ_k=0,
 \quad P_k+Q_k=I.
\end{equation}

We separate the transport part that preserves sheet from the one that changes it:
\begin{align}
 A_k&=P_{k+1}U_kP_k+Q_{k+1}U_kQ_k,
 \label{eq:continuo-sp-a}\\
 \eta_k&=Q_{k+1}U_kP_k+P_{k+1}U_kQ_k.
 \label{eq:continuo-sp-eta}
\end{align}
Then
\begin{equation}
\label{eq:continuo-sp-u-a-eta}
 U_k=A_k+\eta_k,
\end{equation}
and the graduation verifies
\begin{equation}
\label{eq:continuo-sp-paridad-a-eta}
 \sigma_{k+1}A_k=A_k\sigma_k,
 \qquad
 \sigma_{k+1}\eta_k=-\eta_k\sigma_k.
\end{equation}

\begin{proposition}[Leaf local balance]
\label{prop:continuo-sp-balance-local}
For each level,
\begin{equation}
\label{eq:continuo-sp-balance-local}
 A_k^*A_k+\eta_k^*\eta_k=I.
\end{equation}
\end{proposition}

\begin{proof}
The orthogonality of \(P_{k+1}\) and \(Q_{k+1}\) gives
\begin{align*}
 A_k^*A_k
 &=P_kU_k^*P_{k+1}U_kP_k
   +Q_kU_k^*Q_{k+1}U_kQ_k,\\
 \eta_k^*\eta_k
 &=P_kU_k^*Q_{k+1}U_kP_k
   +Q_kU_k^*P_{k+1}U_kQ_k.
\end{align*}
Adding and using \(P_{k+1}+Q_{k+1}=I\),
\[
 A_k^*A_k+\eta_k^*\eta_k
 =P_kU_k^*U_kP_k+Q_kU_k^*U_kQ_k
 =P_k+Q_k=I.
\]
\end{proof}

The operator \(\eta_k\) records departure from the sheet-preserving channel. It
is not the integer carry: the coordinates
\eqref{eq:continuo-sp-carry-nativo}--
\eqref{eq:continuo-sp-carry-operador} remain in the package
\eqref{eq:continuo-sp-d13}.

\fi
\section{Signed reader with genealogical register}

For \(v\in\mathcal H_k^{\mathrm{sp}}\), the signed polarization is
\begin{equation}
\label{eq:continuo-sp-polarizacion-firmada}
 \operatorname{pol}_k(v)
 =\langle v,\sigma_kv\rangle
 =\|P_kv\|^2-\|Q_kv\|^2.
\end{equation}
The specialized reader furthermore preserves the source:
\begin{equation}
\label{eq:continuo-sp-lector-firmado}
 \mathcal L^{\rm sgn}_{\mathrm{sp},k}(x_k;v)
 =\bigl(
   \Lambda(\gamma_{x_k}),
   \Sigma_{\rm reg}(\Lambda(\gamma_{x_k})),
   \|P_kv\|^2,
   \|Q_kv\|^2,
   \operatorname{pol}_k(v)\bigr).
\end{equation}
The TRIT orientation is preserved in \(\Lambda(\gamma_{x_k})\); it is not
identified with the sheet sign. At horizon \(N\), the record keeps the
family of readers for all levels \(k\leq N\). Truncating means
removing only subsequent entries.

\section{Non-autonomous and terminal product}

We define
\begin{equation}
\label{eq:continuo-sp-f-producto}
 F_{0,0}=I_{\mathcal H_0^{\mathrm{sp}}},
 \qquad
 F_{0,k}=A_{k-1}\cdots A_1A_0
 \quad(k\geq1).
\end{equation}
Therefore,
\begin{equation}
\label{eq:continuo-sp-f-recursion}
 F_{0,k+1}=A_kF_{0,k}.
\end{equation}
The term published upon leaving the channel at step \(k\) and the terminal at
horizon \(N\) are
\begin{equation}
\label{eq:continuo-sp-defecto-terminal}
 \mathcal D_k=F_{0,k}^*\eta_k^*\eta_kF_{0,k},
 \qquad
 \mathcal T_N=F_{0,N}^*F_{0,N}.
\end{equation}

\begin{theorem}[Spectral telescope--non-autonomous polar]
\label{thm:continuo-sp-telescopia-no-autonoma}
For every finite horizon \(N\),
\begin{equation}
\label{eq:continuo-sp-telescopia-no-autonoma}
 \boxed{
 I=\sum_{k=0}^{N-1}
 F_{0,k}^*\eta_k^*\eta_kF_{0,k}
 +F_{0,N}^*F_{0,N}.}
\end{equation}
Moreover,
\[
 0\preceq\mathcal T_{N+1}
 \preceq\mathcal T_N\preceq I,
\]
so that the strong positive limit exists
\[
 \mathcal T_\infty
 =\operatorname{s-lim}_{N\to\infty}\mathcal T_N
\]
and
\begin{equation}
\label{eq:continuo-sp-telescopia-infinita}
 I=\sum_{k=0}^{\infty}
 F_{0,k}^*\eta_k^*\eta_kF_{0,k}
 +\mathcal T_\infty
\end{equation}
in the strong topology. The theorem does not assert that
\(\mathcal T_\infty\) is zero.
\end{theorem}

\begin{proof}
Multiplying \eqref{eq:pdfv2-cor-balance-hojas} by
\(F_{0,k}^*\) and \(F_{0,k}\), and using
\eqref{eq:continuo-sp-f-recursion}, one obtains
\[
 F_{0,k}^*F_{0,k}
 =F_{0,k+1}^*F_{0,k+1}
  +F_{0,k}^*\eta_k^*\eta_kF_{0,k}.
\]
The sum from \(k=0\) to \(N-1\) cancels all intermediate
terms and leaves precisely
\eqref{eq:continuo-sp-telescopia-no-autonoma}. Since
\(A_k^*A_k\preceq I\),
\[
 \mathcal T_{N+1}
 =F_{0,N}^*A_N^*A_NF_{0,N}
 \preceq F_{0,N}^*F_{0,N}=\mathcal T_N.
\]
Monotonicity and the positive bound provide the strong limit. Passing to the
limit in the finite identity gives
\eqref{eq:continuo-sp-telescopia-infinita}.
\end{proof}

\iffalse
\section{Assembler as graph}

For
\[
 z_N=(r_{\mathrm{sp},N}x_N,x_N,D_{\mathrm{sp},N}x_N)
 \in X_{\chi,\mathrm{sp},N},
\]
the material ensemble is defined
\begin{equation}
\label{eq:continuo-sp-a-accion}
\begin{aligned}
 a_{\mathrm{sp},N}(z_N)=\bigl(&
 (\mathcal H_k^{\mathrm{sp}},\sigma_k,P_k,Q_k)_{k\leq N},
 (U_k,A_k,\eta_k)_{k<N},\\
 &(F_{0,k},\mathcal D_k,\mathcal T_k,
   \mathcal L^{\rm sgn}_{\mathrm{sp},k})_{k\leq N}\bigr).
\end{aligned}
\end{equation}
The macrostructure is
\begin{equation}
\label{eq:continuo-sp-m-graph}
 \mathfrak M_{\mathrm{sp},N}
 =\operatorname{Graph}(a_{\mathrm{sp},N}),
 \qquad
 \mathcal A_{\mathrm{sp},N}(z_N)
 =(z_N,a_{\mathrm{sp},N}z_N).
\end{equation}
Projection onto \(z_N\) is the inverse of
\(\mathcal A_{\mathrm{sp},N}\). Truncation removes the entries whose
index exceeds the new horizon and satisfies
\begin{equation}
\label{eq:continuo-sp-a-naturalidad}
 m^{\mathrm{sp}}_{\ell,k}a_{\mathrm{sp},\ell}
 =a_{\mathrm{sp},k}\xi^{\mathrm{sp}}_{\ell,k}.
\end{equation}

\section{Publisher as graph}

The publisher does not contract the macrostructure to a scalar. Return the
complete certificate.
\begin{equation}
\label{eq:continuo-sp-p-accion}
 p_{\mathrm{sp},N}(m_N)
 =\left(
 \mathcal L^{\rm sgn}_{\mathrm{sp},k},
 \mathcal D_k,
 \mathcal T_k,
 I=\sum_{j=0}^{k-1}\mathcal D_j+\mathcal T_k
 \right)_{0\leq k\leq N}.
\end{equation}
We define
\begin{equation}
\label{eq:continuo-sp-c-graph}
 \mathcal C_{\mathrm{sp},N}
 =\operatorname{Graph}(p_{\mathrm{sp},N}),
 \qquad
 \mathcal P_{\mathrm{sp},N}(m_N)
 =(m_N,p_{\mathrm{sp},N}m_N).
\end{equation}
Projection onto \(m_N\) is its inverse. Truncating the certificate removes only
subsequent identities, so
\begin{equation}
\label{eq:continuo-sp-p-naturalidad}
 c^{\mathrm{sp}}_{\ell,k}p_{\mathrm{sp},\ell}
 =p_{\mathrm{sp},k}m^{\mathrm{sp}}_{\ell,k}.
\end{equation}

\section{Theorem of the complete internal chain}

\begin{theorem}[Spectral-polar macrostructure generated by the continuum]
\label{thm:continuo-sp-cadena-interna}
The actions
\eqref{eq:continuo-sp-res-graph},
\eqref{eq:continuo-sp-pr-x},
\eqref{eq:continuo-sp-m-graph}, and
\eqref{eq:continuo-sp-c-graph} form the natural chain
\begin{equation}
\label{eq:continuo-sp-cadena-interna}
 \mathcal C_{\rm cont}^{\rm disc}
 \supset\mathcal C_{\rm cont}^{\mathrm{sp}}
 \xrightarrow[\cong]{\operatorname{Res}_{\mathrm{sp}}}
 \mathcal G_{\mathrm{sp}}
 \xrightarrow[\cong]{\operatorname{pr}_{X,\mathrm{sp}}}
 X_{\chi,\mathrm{sp}}
 \xrightarrow[\cong]{\mathcal A_{\mathrm{sp}}}
 \mathfrak M_{\mathrm{sp}}
 \xrightarrow[\cong]{\mathcal P_{\mathrm{sp}}}
 \mathcal C_{\mathrm{sp}}.
\end{equation}
Its published content consists of the two-sheet bundle, the grading
\(\sigma_k\), the projectors \(P_k,Q_k\), the total transport \(U_k\),
the separation \(U_k=A_k+\eta_k\), the carry and the memory, the signed reader,
and the telescoping structure
\eqref{eq:continuo-sp-telescopia-no-autonoma} with an explicit terminal object. Each
object preserves a canonical projection to all its antecedents.
\end{theorem}

\begin{proof}
The identities
\eqref{eq:continuo-sp-res-inversa} prove the first equivalence. The
dependent definition \eqref{eq:continuo-sp-x-dependiente} proves the
second. The graphs
\eqref{eq:continuo-sp-m-graph} and
\eqref{eq:continuo-sp-c-graph} prove the remaining two. Naturality of
the four steps is
\eqref{eq:continuo-sp-d-naturalidad},
\eqref{eq:continuo-sp-a-naturalidad}, and
\eqref{eq:continuo-sp-p-naturalidad}, together with naturality of
\(r_{\mathrm{sp},k}\). The operatorial content is proved by
\eqref{eq:continuo-sp-u-isometria}, the
\cref{prop:continuo-sp-balance-local}, and the
\cref{thm:continuo-sp-telescopia-no-autonoma}. All extensions appear
in \eqref{eq:continuo-sp-u-accion}; consequently, no choice of
route intervenes.
\end{proof}

\section{Phase return and state evolution}

The branch explicitly preserves
\begin{equation}
\label{eq:continuo-sp-no-reset}
 g(k+9)=g(k),
 \qquad
 x_{k+9}\neq x_k\quad\text{in general}.
\end{equation}
The second member records the increase in memory and possible changes
of sheet, carry, cylinder and boundary. Neither \(\sigma_k\), nor the projectors, nor
the signed reader contract this difference.

\section{Origin}

The two APP sheets, the holonomic tritic state, the TPK fibers and extensions,
carry, return without reinitialization, survival and the ledger are
\texttt{ARCHITECTURE\_AUTHORIAL\_PREEXISTING}. The obligation to preserve the
finite terminal is a \texttt{RESULT\_RECOVERED}. The explicit sheaf of sheets,
the involution \(\sigma_k\), the projectors, the linearization summing all
extensions, the decomposition \(U_k=A_k+\eta_k\), the nonautonomous
product \(F_{0,k}\) and the dependent graphs are
\texttt{FORMALIZATION\_NEW}. The
\cref{thm:continuo-sp-telescopia-no-autonoma,thm:continuo-sp-cadena-interna}
constitutes the \texttt{CERTIFICATE\_NEW} of that formalization.

\section{Falsadores internos}

\begin{proposition}[Falsifiers of the macrostructure]
\label{prop:continuo-sp-falsadores}
The conclusion of
\cref{thm:continuo-sp-cadena-interna} is not transferable if any
of the following substitutions occurs:
\begin{enumerate}
 \item choosing a prolongation in
       \eqref{eq:continuo-sp-u-accion} and deleting the remaining ones;
 \item replacing \(\mathcal G_{\mathrm{sp}}\) with the external image of
       \(D_{\mathrm{sp}}\);
 \item deleting \((x,D_{\mathrm{sp}}x)\) when forming
       \(X_{\chi,\mathrm{sp}}\);
 \item identifying the sheet with the TRIT orientation or deleting the threshold;
 \item breaking the consistency
       \eqref{eq:continuo-sp-consistencia-medida};
 \item deleting the carry or violating
       \eqref{eq:continuo-sp-carry-nativo}--
       \eqref{eq:continuo-sp-carry-operador};
 \item identifying \(\eta_k\) with the entire carry;
 \item replacing the family \((A_k)_k\) with a single operator without proving
       stationarity;
 \item removing \(F_{0,N}^*F_{0,N}\) at a finite horizon;
 \item asserting \(\mathcal T_\infty=0\) without additional proof;
 \item reconstructing a history from the signed polarization;
 \item using \(g(k+9)=g(k)\) to assert \(x_{k+9}=x_k\);
 \item introducing an external domain, function, or criterion to
define any of the internal arrows;
 \item introducing a subsequent output as an antecedent of the branch.
\end{enumerate}
\end{proposition}
\fi
   % Extracción quirúrgica del fuente teoremática V2 05a: líneas 512--523.
% El resto permanece en este archivo como procedencia, pero no se publica.
\iffalse
\chapter{Spectral macrostructure--internal polar continuum}
\label{chap:continuo-macro-sp-interna}

\section{Strictly internal scope}

The construction in this chapter terminates before any external recognition. Its sole data are the two-sheet APP support, the TRIT orientation, the TPK extensions and carry, the projective system of histories, the cylindrical measure, and the signed ledger. The result is a graded macrostructure with transport, a sheet-change expression, a signed reader, and nonautonomous telescoping. No external function is used to define any of its maps.

\section{Surviving domain of the branch}

For \(x_k\in\mathcal A_k\), let
\begin{equation}
\label{eq:continuo-sp-msect-n}
 \operatorname{Msect}^{\mathrm{sp},(N)}_k(x_k)
 =\left\{(x_j)_{j=k}^{N}:
 \tau_{j+1,j}x_{j+1}=x_j,
 \ (x_j,x_{j+9})\in\Gamma_{9,j},
 \ \text{and the sheaf is preserved }\{\Sigma,\Pi\}\right\}.
\end{equation}
We define
\begin{equation}
\label{eq:continuo-sp-dominio-superviviente}
 \mathcal C_k^{\mathrm{sp}}
 =\left\{x_k:
 \operatorname{Msect}^{\mathrm{sp},(N)}_k(x_k)\neq\varnothing
 \text{ for every }N\geq k\right\},
 \qquad
 \mathcal C_{\rm cont}^{\mathrm{sp}}
 =\varprojlim_k\mathcal C_k^{\mathrm{sp}}.
\end{equation}
The TRIT threshold is not excluded from this domain. The bundle preserves the two
arithmetic sheets even if a visible chart publishes neither of them in a
specific event.

\section{Thirteen-coordinate package}

For \(x_k\in\mathcal C_k^{\mathrm{sp}}\), the dependent package is
\begin{equation}
\label{eq:continuo-sp-d13}
\begin{aligned}
 D_{\mathrm{sp},k}(x_k)=\bigl(&
 \mathcal F_{\mathrm{sp},k}(x_k),
 \operatorname{Ext}_{\Gamma,\mathrm{sp},k}(x_k),
 \operatorname{Rel}_{\mathrm{sp},k}(x_k),
 \mathcal N_{\mathrm{sp},k}(x_k),
 \operatorname{or}_{\mathrm{sp},k}(x_k),\\
 &\operatorname{Carr}_{\mathrm{sp},k}(x_k),
 \operatorname{Mem}_{\mathrm{sp},k}(x_k),
 \partial_{\mathrm{sp},k}(x_k),
 \gamma_{\mathrm{sp},k}(x_k),
 \operatorname{Msect}_{\mathrm{sp},k}(x_k),\\
 &\operatorname{Surv}_{\mathrm{sp},k}(x_k),
 \operatorname{Term}_{\mathrm{sp},k}(x_k),
 \mathcal L^{\rm sgn}_{\mathrm{sp},k}(x_k)\bigr).
\end{aligned}
\end{equation}
Their components are determined correlatively as follows.

The fiber \(\mathcal F_{\mathrm{sp},k}\) contains the TPK state and the ordered
bundle \(\{\Sigma,\Pi\}\). The extension is the complete restriction of
\(\operatorname{Ext}_\Gamma\) to histories satisfying
\eqref{eq:continuo-sp-dominio-superviviente}. The relations contain the
identity and composition laws, the residue--quotient lift, the carry
cocycle, boundary compatibility and truncation naturality. The
numerical coordinate contains depth, phase, residues, quotients and internal
signature. Orientation stores sheet and TRIT orientation separately.
Memory contains the quotient, the ledger and the turn increment. The
boundary and route remain uncontracted. The multisection is the complete
system of prolongations of \eqref{eq:continuo-sp-msect-n}.
Survival is the stable membership of
\eqref{eq:continuo-sp-dominio-superviviente}. The terminal and reader are
constructed in the following sections.

The truncations of all coordinates form a map
\(u^{\mathrm{sp}}_{\ell,k}\) satisfying
\begin{equation}
\label{eq:continuo-sp-d-naturalidad}
 u^{\mathrm{sp}}_{\ell,k}D_{\mathrm{sp},\ell}
 =D_{\mathrm{sp},k}\tau_{\ell,k}.
\end{equation}

\section{Graph of the constraint and dependent section}

We define
\begin{equation}
\label{eq:continuo-sp-g-graph}
 \mathcal G_{\mathrm{sp},k}
 =\operatorname{Graph}(D_{\mathrm{sp},k})
 =\{(x_k,D_{\mathrm{sp},k}x_k):
     x_k\in\mathcal C_k^{\mathrm{sp}}\}.
\end{equation}
The total restriction is
\begin{equation}
\label{eq:continuo-sp-res-graph}
 \operatorname{Res}_{\mathrm{sp}}=s_{\mathrm{sp}}:
 \mathcal C_{\rm cont}^{\mathrm{sp}}
 \longrightarrow\mathcal G_{\mathrm{sp}},
 \qquad
 s_{\mathrm{sp}}(x)=(x,D_{\mathrm{sp}}x).
\end{equation}
The projection \(\pi_1(x,d)=x\) verifies
\begin{equation}
\label{eq:continuo-sp-res-inversa}
 \pi_1s_{\mathrm{sp}}=I,
 \qquad
 s_{\mathrm{sp}}\pi_1=I_{\mathcal G_{\mathrm{sp}}}.
\end{equation}

Let \(r_{\mathrm{sp},k}(x_k)\) be the multisection of sheets and paths that already
belongs to \(D_{\mathrm{sp},k}(x_k)\). Define
\begin{equation}
\label{eq:continuo-sp-x-dependiente}
 X_{\chi,\mathrm{sp},k}
 =\{(r_{\mathrm{sp},k}x_k,x_k,D_{\mathrm{sp},k}x_k):
     x_k\in\mathcal C_k^{\mathrm{sp}}\},
\end{equation}
\begin{equation}
\label{eq:continuo-sp-pr-x}
 \operatorname{pr}_{X,\mathrm{sp}}(x,D_{\mathrm{sp}}x)
 =(r_{\mathrm{sp}}x,x,D_{\mathrm{sp}}x).
\end{equation}
Forgetting the first coordinate of
\eqref{eq:continuo-sp-x-dependiente} is the inverse of
\(\operatorname{pr}_{X,\mathrm{sp}}\).

\section{Bundle of leaves and projective measure}

Let
\begin{equation}
\label{eq:continuo-sp-haz-hojas}
 \widetilde{\mathcal C}_k^{\mathrm{sp}}
 =\{(x_k,h):x_k\in\mathcal C_k^{\mathrm{sp}},
              h\in\{\Sigma,\Pi\}\}.
\end{equation}
The projective measure of the continuum induces probabilities
\(\widetilde\mu_k\) on this bundle. We restrict to the positive support; for
\(z\) in that support,
\begin{equation}
\label{eq:continuo-sp-consistencia-medida}
 \widetilde\mu_k(z)
 =\sum_{\widetilde\tau_{k+1,k}(z')=z}
  \widetilde\mu_{k+1}(z').
\end{equation}
Level linearization is
\begin{equation}
\label{eq:continuo-sp-hilbert-nivel}
 \mathcal H_k^{\mathrm{sp}}
 =\ell^2(\operatorname{supp}\widetilde\mu_k),
\end{equation}
with orthonormal basis \(e_z\).

\section{Transport of all extensions and carry}

For \(z'\) extending \(z\), we define the positive conditional weight
\begin{equation}
\label{eq:continuo-sp-peso-condicional}
 w_k(z'\mid z)
 =\left(\frac{\widetilde\mu_{k+1}(z')}
               {\widetilde\mu_k(z)}\right)^{1/2}.
\end{equation}
Transport is
\begin{equation}
\label{eq:continuo-sp-u-accion}
 U_ke_z
 =\sum_{\substack{z':\widetilde\tau_{k+1,k}(z')=z\\
            (z,z')\in\operatorname{Ext}_{\Gamma,\mathrm{sp},k}}}
 w_k(z'\mid z)e_{z'}.
\end{equation}
The sum contains all admissible extensions. It does not depend on a selector.

By \eqref{eq:continuo-sp-consistencia-medida},
\[
 \sum_{z':\widetilde\tau(z')=z}|w_k(z'\mid z)|^2=1.
\]
The sets of descendants of two distinct parents are disjoint. Consequently,
\begin{equation}
\label{eq:continuo-sp-u-isometria}
 U_k^*U_k=I_{\mathcal H_k^{\mathrm{sp}}}.
\end{equation}

Each basis \(e_{z'}\) includes the carry of the corresponding extension. For
composable paths, we preserve
\begin{align}
 \kappa(\gamma_2\gamma_1)
 &=\kappa(\gamma_2)+
   \mathsf C(\gamma_2)\kappa(\gamma_1),
 \label{eq:continuo-sp-carry-nativo}\\
 \operatorname{Carr}(\gamma_2\gamma_1)
 &=\operatorname{Carr}(\gamma_2)H(\gamma_1)
  +Y(\gamma_2)\operatorname{Carr}(\gamma_1).
 \label{eq:continuo-sp-carry-operador}
\end{align}
The conditional weights multiply along the unique chain of
truncations of a descendant. Equations
\eqref{eq:continuo-sp-carry-nativo}--
\eqref{eq:continuo-sp-carry-operador} ensure that the carry labels
coincide. Therefore,
\begin{equation}
\label{eq:continuo-sp-u-composicion}
 U_{\ell,k}=U_{\ell-1}\cdots U_k.
\end{equation}

\section{Graduation, projectors and decomposition of transport}

The sheet graduation is the involution.
\begin{equation}
\label{eq:continuo-sp-sigma}
 \sigma_ke_{(x,\Sigma)}=+e_{(x,\Sigma)},
 \qquad
 \sigma_ke_{(x,\Pi)}=-e_{(x,\Pi)}.
\end{equation}
Thus,
\[
 \sigma_k^*=\sigma_k,
 \qquad
 \sigma_k^2=I.
\]
The internal projectors are
\begin{equation}
\label{eq:continuo-sp-p-q}
 P_k=\frac{I+\sigma_k}{2},
 \qquad
 Q_k=\frac{I-\sigma_k}{2}.
\end{equation}
Satisfy
\begin{equation}
\label{eq:continuo-sp-p-q-relaciones}
 P_k^2=P_k,
 \quad Q_k^2=Q_k,
 \quad P_kQ_k=0,
 \quad P_k+Q_k=I.
\end{equation}

We separate the transport part that preserves sheet from the one that changes it:
\begin{align}
 A_k&=P_{k+1}U_kP_k+Q_{k+1}U_kQ_k,
 \label{eq:continuo-sp-a}\\
 \eta_k&=Q_{k+1}U_kP_k+P_{k+1}U_kQ_k.
 \label{eq:continuo-sp-eta}
\end{align}
Then
\begin{equation}
\label{eq:continuo-sp-u-a-eta}
 U_k=A_k+\eta_k,
\end{equation}
and the graduation verifies
\begin{equation}
\label{eq:continuo-sp-paridad-a-eta}
 \sigma_{k+1}A_k=A_k\sigma_k,
 \qquad
 \sigma_{k+1}\eta_k=-\eta_k\sigma_k.
\end{equation}

\begin{proposition}[Leaf local balance]
\label{prop:continuo-sp-balance-local}
For each level,
\begin{equation}
\label{eq:continuo-sp-balance-local}
 A_k^*A_k+\eta_k^*\eta_k=I.
\end{equation}
\end{proposition}

\begin{proof}
The orthogonality of \(P_{k+1}\) and \(Q_{k+1}\) gives
\begin{align*}
 A_k^*A_k
 &=P_kU_k^*P_{k+1}U_kP_k
   +Q_kU_k^*Q_{k+1}U_kQ_k,\\
 \eta_k^*\eta_k
 &=P_kU_k^*Q_{k+1}U_kP_k
   +Q_kU_k^*P_{k+1}U_kQ_k.
\end{align*}
Adding and using \(P_{k+1}+Q_{k+1}=I\),
\[
 A_k^*A_k+\eta_k^*\eta_k
 =P_kU_k^*U_kP_k+Q_kU_k^*U_kQ_k
 =P_k+Q_k=I.
\]
\end{proof}

The operator \(\eta_k\) records departure from the sheet-preserving channel. It
is not the integer carry: the coordinates
\eqref{eq:continuo-sp-carry-nativo}--
\eqref{eq:continuo-sp-carry-operador} remain in the package
\eqref{eq:continuo-sp-d13}.

\section{Signed reader with ledger}

For \(v\in\mathcal H_k^{\mathrm{sp}}\), the signed polarization is
\begin{equation}
\label{eq:continuo-sp-polarizacion-firmada}
 \operatorname{pol}_k(v)
 =\langle v,\sigma_kv\rangle
 =\|P_kv\|^2-\|Q_kv\|^2.
\end{equation}
The specialized reader furthermore preserves the source:
\begin{equation}
\label{eq:continuo-sp-lector-firmado}
 \mathcal L^{\rm sgn}_{\mathrm{sp},k}(x_k;v)
 =\bigl(
   \Lambda(\gamma_{x_k}),
   \Sigma_{\rm reg}(\Lambda(\gamma_{x_k})),
   \|P_kv\|^2,
   \|Q_kv\|^2,
   \operatorname{pol}_k(v)\bigr).
\end{equation}
The TRIT orientation is preserved in \(\Lambda(\gamma_{x_k})\); it is not
identified with the sheet sign. At horizon \(N\), the record keeps the
family of readers for all levels \(k\leq N\). Truncating means
removing only subsequent entries.

\section{Non-autonomous and terminal product}

We define
\begin{equation}
\label{eq:continuo-sp-f-producto}
 F_{0,0}=I_{\mathcal H_0^{\mathrm{sp}}},
 \qquad
 F_{0,k}=A_{k-1}\cdots A_1A_0
 \quad(k\geq1).
\end{equation}
Therefore,
\begin{equation}
\label{eq:continuo-sp-f-recursion}
 F_{0,k+1}=A_kF_{0,k}.
\end{equation}
The term published upon leaving the channel at step \(k\) and the terminal at
horizon \(N\) are
\begin{equation}
\label{eq:continuo-sp-defecto-terminal}
 \mathcal D_k=F_{0,k}^*\eta_k^*\eta_kF_{0,k},
 \qquad
 \mathcal T_N=F_{0,N}^*F_{0,N}.
\end{equation}

\begin{theorem}[Spectral telescope--non-autonomous polar]
\label{thm:continuo-sp-telescopia-no-autonoma}
For every finite horizon \(N\),
\begin{equation}
\label{eq:continuo-sp-telescopia-no-autonoma}
 \boxed{
 I=\sum_{k=0}^{N-1}
 F_{0,k}^*\eta_k^*\eta_kF_{0,k}
 +F_{0,N}^*F_{0,N}.}
\end{equation}
Moreover,
\[
 0\preceq\mathcal T_{N+1}
 \preceq\mathcal T_N\preceq I,
\]
so that the strong positive limit exists
\[
 \mathcal T_\infty
 =\operatorname{s-lim}_{N\to\infty}\mathcal T_N
\]
and
\begin{equation}
\label{eq:continuo-sp-telescopia-infinita}
 I=\sum_{k=0}^{\infty}
 F_{0,k}^*\eta_k^*\eta_kF_{0,k}
 +\mathcal T_\infty
\end{equation}
in the strong topology. The theorem does not assert that
\(\mathcal T_\infty\) is zero.
\end{theorem}

\begin{proof}
Multiplying \eqref{eq:continuo-sp-balance-local} by
\(F_{0,k}^*\) and \(F_{0,k}\), and using
\eqref{eq:continuo-sp-f-recursion}, one obtains
\[
 F_{0,k}^*F_{0,k}
 =F_{0,k+1}^*F_{0,k+1}
  +F_{0,k}^*\eta_k^*\eta_kF_{0,k}.
\]
The sum from \(k=0\) to \(N-1\) cancels all intermediate
terms and leaves precisely
\eqref{eq:continuo-sp-telescopia-no-autonoma}. Since
\(A_k^*A_k\preceq I\),
\[
 \mathcal T_{N+1}
 =F_{0,N}^*A_N^*A_NF_{0,N}
 \preceq F_{0,N}^*F_{0,N}=\mathcal T_N.
\]
Monotonicity and the positive bound provide the strong limit. Passing to the
limit in the finite identity gives
\eqref{eq:continuo-sp-telescopia-infinita}.
\end{proof}

\section{Assembler as graph}

For
\[
 z_N=(r_{\mathrm{sp},N}x_N,x_N,D_{\mathrm{sp},N}x_N)
 \in X_{\chi,\mathrm{sp},N},
\]
the material ensemble is defined
\begin{equation}
\label{eq:continuo-sp-a-accion}
\begin{aligned}
 a_{\mathrm{sp},N}(z_N)=\bigl(&
 (\mathcal H_k^{\mathrm{sp}},\sigma_k,P_k,Q_k)_{k\leq N},
 (U_k,A_k,\eta_k)_{k<N},\\
 &(F_{0,k},\mathcal D_k,\mathcal T_k,
   \mathcal L^{\rm sgn}_{\mathrm{sp},k})_{k\leq N}\bigr).
\end{aligned}
\end{equation}
The macrostructure is
\begin{equation}
\label{eq:continuo-sp-m-graph}
 \mathfrak M_{\mathrm{sp},N}
 =\operatorname{Graph}(a_{\mathrm{sp},N}),
 \qquad
 \mathcal A_{\mathrm{sp},N}(z_N)
 =(z_N,a_{\mathrm{sp},N}z_N).
\end{equation}
Projection onto \(z_N\) is the inverse of
\(\mathcal A_{\mathrm{sp},N}\). Truncation removes the entries whose
index exceeds the new horizon and satisfies
\begin{equation}
\label{eq:continuo-sp-a-naturalidad}
 m^{\mathrm{sp}}_{\ell,k}a_{\mathrm{sp},\ell}
 =a_{\mathrm{sp},k}\xi^{\mathrm{sp}}_{\ell,k}.
\end{equation}

\section{Publisher as graph}

The publisher does not contract the macrostructure to a scalar. Return the
complete certificate.
\begin{equation}
\label{eq:continuo-sp-p-accion}
 p_{\mathrm{sp},N}(m_N)
 =\left(
 \mathcal L^{\rm sgn}_{\mathrm{sp},k},
 \mathcal D_k,
 \mathcal T_k,
 I=\sum_{j=0}^{k-1}\mathcal D_j+\mathcal T_k
 \right)_{0\leq k\leq N}.
\end{equation}
We define
\begin{equation}
\label{eq:continuo-sp-c-graph}
 \mathcal C_{\mathrm{sp},N}
 =\operatorname{Graph}(p_{\mathrm{sp},N}),
 \qquad
 \mathcal P_{\mathrm{sp},N}(m_N)
 =(m_N,p_{\mathrm{sp},N}m_N).
\end{equation}
Projection onto \(m_N\) is its inverse. Truncating the certificate removes only
subsequent identities, so
\begin{equation}
\label{eq:continuo-sp-p-naturalidad}
 c^{\mathrm{sp}}_{\ell,k}p_{\mathrm{sp},\ell}
 =p_{\mathrm{sp},k}m^{\mathrm{sp}}_{\ell,k}.
\end{equation}

\section{Theorem of the complete internal chain}

\begin{theorem}[Spectral-polar macrostructure generated by the continuum]
\label{thm:continuo-sp-cadena-interna}
The actions
\eqref{eq:continuo-sp-res-graph},
\eqref{eq:continuo-sp-pr-x},
\eqref{eq:continuo-sp-m-graph}, and
\eqref{eq:continuo-sp-c-graph} form the natural chain
\begin{equation}
\label{eq:continuo-sp-cadena-interna}
 \mathcal C_{\rm cont}^{\rm disc}
 \supset\mathcal C_{\rm cont}^{\mathrm{sp}}
 \xrightarrow[\cong]{\operatorname{Res}_{\mathrm{sp}}}
 \mathcal G_{\mathrm{sp}}
 \xrightarrow[\cong]{\operatorname{pr}_{X,\mathrm{sp}}}
 X_{\chi,\mathrm{sp}}
 \xrightarrow[\cong]{\mathcal A_{\mathrm{sp}}}
 \mathfrak M_{\mathrm{sp}}
 \xrightarrow[\cong]{\mathcal P_{\mathrm{sp}}}
 \mathcal C_{\mathrm{sp}}.
\end{equation}
Its published content consists of the two-sheet bundle, the grading
\(\sigma_k\), the projectors \(P_k,Q_k\), the total transport \(U_k\),
the separation \(U_k=A_k+\eta_k\), the carry and the memory, the signed reader,
and the telescoping structure
\eqref{eq:continuo-sp-telescopia-no-autonoma} with an explicit terminal object. Each
object preserves a canonical projection to all its antecedents.
\end{theorem}

\begin{proof}
The identities
\eqref{eq:continuo-sp-res-inversa} prove the first equivalence. The
dependent definition \eqref{eq:continuo-sp-x-dependiente} proves the
second. The graphs
\eqref{eq:continuo-sp-m-graph} and
\eqref{eq:continuo-sp-c-graph} prove the remaining two. Naturality of
the four steps is
\eqref{eq:continuo-sp-d-naturalidad},
\eqref{eq:continuo-sp-a-naturalidad}, and
\eqref{eq:continuo-sp-p-naturalidad}, together with naturality of
\(r_{\mathrm{sp},k}\). The operatorial content is proved by
\eqref{eq:continuo-sp-u-isometria}, the
\cref{prop:continuo-sp-balance-local}, and the
\cref{thm:continuo-sp-telescopia-no-autonoma}. All extensions appear
in \eqref{eq:continuo-sp-u-accion}; consequently, no choice of
route intervenes.
\end{proof}

\fi
\section{Phase return and non-reinitialized evolution of the common state}

The branch explicitly preserves
\begin{equation}
\label{eq:continuo-sp-no-reset}
 g(k+9)=g(k),
 \qquad
 x_{k+9}\neq x_k\quad\text{in general}.
\end{equation}
The second member records the increase in memory and possible changes
of sheet, carry, cylinder and boundary. Neither \(\sigma_k\), nor the projectors, nor
the signed reader contract this difference.

\iffalse
\section{Origin}

The two APP sheets, the holonomic tritic state, the TPK fibers and extensions,
carry, return without reinitialization, survival and the ledger are
\texttt{ARCHITECTURE\_AUTHORIAL\_PREEXISTING}. The obligation to preserve the
finite terminal is a \texttt{RESULT\_RECOVERED}. The explicit sheaf of sheets,
the involution \(\sigma_k\), the projectors, the linearization summing all
extensions, the decomposition \(U_k=A_k+\eta_k\), the nonautonomous
product \(F_{0,k}\) and the dependent graphs are
\texttt{FORMALIZATION\_NEW}. The
\cref{thm:continuo-sp-telescopia-no-autonoma,thm:continuo-sp-cadena-interna}
constitutes the \texttt{CERTIFICATE\_NEW} of that formalization.

\section{Falsadores internos}

\begin{proposition}[Falsifiers of the macrostructure]
\label{prop:continuo-sp-falsadores}
The conclusion of
\cref{thm:continuo-sp-cadena-interna} is not transferable if any
of the following substitutions occurs:
\begin{enumerate}
 \item choosing a prolongation in
       \eqref{eq:continuo-sp-u-accion} and deleting the remaining ones;
 \item replacing \(\mathcal G_{\mathrm{sp}}\) with the external image of
       \(D_{\mathrm{sp}}\);
 \item deleting \((x,D_{\mathrm{sp}}x)\) when forming
       \(X_{\chi,\mathrm{sp}}\);
 \item identifying the sheet with the TRIT orientation or deleting the threshold;
 \item breaking the consistency
       \eqref{eq:continuo-sp-consistencia-medida};
 \item deleting the carry or violating
       \eqref{eq:continuo-sp-carry-nativo}--
       \eqref{eq:continuo-sp-carry-operador};
 \item identifying \(\eta_k\) with the entire carry;
 \item replacing the family \((A_k)_k\) with a single operator without proving
       stationarity;
 \item removing \(F_{0,N}^*F_{0,N}\) at a finite horizon;
 \item asserting \(\mathcal T_\infty=0\) without additional proof;
 \item reconstructing a history from the signed polarization;
 \item using \(g(k+9)=g(k)\) to assert \(x_{k+9}=x_k\);
 \item introducing an external domain, function, or criterion to
define any of the internal arrows;
 \item introducing a subsequent output as an antecedent of the branch.
\end{enumerate}
\end{proposition}
\fi
 \endgroup

The conditional expectation associated with the same weight is
\begin{equation}
\label{eq:pdfv2-cor-esperanza-comun}
 (E_kF)(\widehat x):=
 \sum_{(\alpha,\widehat y)\in\operatorname{Ch}_k(\widehat x)}
 p_k(\alpha,\widehat y\mid\widehat x)F(\alpha,\widehat y).
\end{equation}
For the fine increment \(b_f\), set
\begin{equation}
\label{eq:pdfv2-cor-kappa-jensen}
 b_c:=E_kb_f,\qquad
 \kappa_k^{\rm can}:=\frac{E_k|b_f|-|b_c|}{2}\geq0.
\end{equation}
Then
\begin{equation}
\label{eq:pdfv2-cor-jensen-hojas}
 E_kb_f^+=b_c^++\kappa_k^{\rm can},\qquad
 E_kb_f^-=b_c^-+\kappa_k^{\rm can}.
\end{equation}
The superscript \({\rm can}\) distinguishes this cancellation carry from the integer
carry \(\kappa_\alpha\); both arise from the same transport, but have
distinct types.

Signed memories require no external initial condition: for
\(\gamma=(\epsilon_1,\ldots,\epsilon_k)\) they are defined by summing the ledger itself,
\begin{equation}
\label{eq:pdfv2-cor-memorias}
 \begin{aligned}
 M_0^\pm&:=0,\\
 M_k^\pm(\gamma)&:=\sum_{j=1}^k b_{\epsilon_j}^\pm,\\
 M_{k+1}^\pm(\gamma\alpha)&:=M_k^\pm(\gamma)+b_\alpha^\pm,\\
 D_k^0&:=M_k^+-M_k^-, &
 H_k^0&:=M_k^++M_k^-.
 \end{aligned}
\end{equation}
Thus, the signed balance, the total memory, and the polarization
\(\langle v,\sigma_kv\rangle\) are three coordinated readings of the same
prefix, not three states reconstructed separately.

\subsection{Nonadic Refinement, Orthogonal Losses and Terminal Memory}
The index \(j\) of the following development is the level of the same tree of
prolongations and its thirteen-coordinate register is the restriction of the
common register of \eqref{eq:pdfv2-cor-registro-extension}. Therefore
\(D_{{\rm sol},j}:=D_j^{(13)}\) is fixed and no additional domain is postulated.
\begingroup
  \let\section\subsection
  % Extracción quirúrgica del fuente teoremática V2 05b: líneas 100--359.
% El resto permanece en este archivo como procedencia, pero no se publica.
\iffalse
% Macroestructura solenoidal estrictamente interna del continuo discreto.
% Este módulo no introduce ninguna ecuación clásica ni un problema exterior.

\chapter{Internal macrostructure of balance, memory, and telescopia}
\label{chap:pdfv2-sol-interna}

\begin{statusbox}
The only input of this chapter is the discrete structure of the continuum.
The branch \(\mathrm{sol}\) is constructed before any exterior scalar evaluation
and before any classical differential formulation.  Its complete chain is
\[
 \mathcal C_{\rm cont}^{\rm disc}
 \supset\mathcal C_{\rm cont}^{\rm sol}
 \xrightarrow{\operatorname{Res}_{\rm sol}}
 \mathcal G_{\rm sol}
 \xrightarrow{\operatorname{pr}_{X,{\rm sol}}}
 X_{\chi,{\rm sol}}
 \xrightarrow{\mathcal A_{\rm sol}}
 \mathfrak M_{\rm sol}
 \xrightarrow{\mathcal P_{\rm sol}}
 \mathcal C_{\rm sol}^{\rm HMT}.
\]
Each arrow preserves as coordinates the history and the apparatus that
produces it.  No image in this chain is a quotient that forgets its source.
\end{statusbox}

\section{Proper domain prior to any subsequent evaluation}
\label{sec:pdfv2-sol-dominio}

Let
\[
 x=(x_0\xrightarrow{e_0}x_1\xrightarrow{e_1}\cdots)
 \in\mathcal C_{\rm cont}^{\rm disc}
\]
a surviving history.  Each edge \(e_j\) preserves its two sheets,
orientation, residue, quotient, and integer carry.  Let
\(N_j^+(e_j)\) and \(N_j^-(e_j)\) be the multiplicities of the two sheets, and let
\(\Delta_{\gamma_j}\in\mathbb Z\) be the oriented coefficient of the route, with the
carry of the composition already incorporated.  The internal signed increment is
\begin{equation}\label{eq:pdfv2-sol-b}
 b_j=b(e_j):=
 \Delta_{\gamma_j}\bigl(N_j^+(e_j)-N_j^-(e_j)\bigr)\in\mathbb Z.
\end{equation}
It is not obtained by evaluating an external path: it is a coordinate of the ledger of the edge itself.

Its positive and negative parts are constructed by
\begin{equation}\label{eq:pdfv2-sol-bpm}
 b_j^+:=\frac{|b_j|+b_j}{2},
 \qquad
 b_j^-:=\frac{|b_j|-b_j}{2}.
\end{equation}
Therefore,
\begin{equation}\label{eq:pdfv2-sol-bpm-identidades}
 b_j=b_j^+-b_j^-,
 \qquad
 |b_j|=b_j^++b_j^-,
 \qquad
 b_j^+b_j^-=0.
\end{equation}

The two original memoirs and their two readers are
\begin{align}
 M_0^+=M_0^-&=0,
 &M_j^+&=\sum_{0\leq i<j}b_i^+,
 &M_j^-&=\sum_{0\leq i<j}b_i^-;
 \label{eq:pdfv2-sol-memorias}\\
 D_j^0&:=M_j^+-M_j^-,
 &H_j^0&:=M_j^++M_j^-.
 \label{eq:pdfv2-sol-dh}
\end{align}
Consequently,
\begin{equation}\label{eq:pdfv2-sol-dh-diferencias}
 D_{j+1}^0-D_j^0=b_j,
 \qquad
 H_{j+1}^0-H_j^0=|b_j|.
\end{equation}
Signed memory \(D^0\) and total memory \(H^0\) remain distinct
objects; retaining only their final difference would destroy the genealogy.

The substructure \(\mathcal C_{\rm cont}^{\rm sol}\) consists of the histories
for which, in every prefix and every nonadic refinement:
\begin{enumerate}
 \item the increments \eqref{eq:pdfv2-sol-b} and the two memories
       \eqref{eq:pdfv2-sol-memorias} exist;
 \item the fine cylinders are grouped into fibres of nine extensions with the
       same projective probability law;
 \item the depth and resolution projections defined below are
       compatible with the history extensions;
 \item each loss row and each terminal column is finite at a finite
       horizon;
 \item every prefix admits at least one compatible extension to any
       subsequent horizon.
\end{enumerate}
These are conditions of membership in the branch, not silent premises
attributed to an arbitrary history.

\fi
\section{Common nonadic refinement: inclusion, balance expectation and memory}
\label{sec:pdfv2-sol-e9j9}

Let \(Y_j(x)\) be the finite set of cylinders compatible with the prefix
\(x_{\leq j}\), endowed with the projective measure \(\mu_j\).  The truncation
\[
 \tau_{j+1,j}:Y_{j+1}(x)\longrightarrow Y_j(x)
\]
has nine admissible extensions over each cylinder. Finite spaces are considered.
\[
 \mathscr H_j(x):=\ell^2\bigl(Y_j(x),\mu_j\bigr).
\]
The inclusion and the nonadic expectation are
\begin{align}
 J_{9,j}:\mathscr H_j&\longrightarrow\mathscr H_{j+1},
 &(J_{9,j}f)(y)&=f(\tau_{j+1,j}y),
 \label{eq:pdfv2-sol-j9}\\
 E_{9,j}:\mathscr H_{j+1}&\longrightarrow\mathscr H_j,
 &(E_{9,j}g)(z)&=\frac1{9}
   \sum_{\tau_{j+1,j}y=z}g(y).
 \label{eq:pdfv2-sol-e9}
\end{align}
The projective compatibility of \(\mu_{j+1}\) and \(\mu_j\) gives
\begin{equation}\label{eq:pdfv2-sol-ej-projector}
 E_{9,j}J_{9,j}=I_{\mathscr H_j},
 \qquad
 J_{9,j}^*=E_{9,j},
 \qquad
 P_{9,j}:=J_{9,j}E_{9,j}=P_{9,j}^*=P_{9,j}^2.
\end{equation}
The complementary microscopic projector is
\begin{equation}\label{eq:pdfv2-sol-qmicro-preliminar}
 Q_{9,j}:=I-P_{9,j}.
\end{equation}

If \(b_f\) is the increment on the fine fiber and
\begin{equation}\label{eq:pdfv2-sol-bcoarse}
 b_c:=E_{9,j}b_f,
\end{equation}
the elementary convexity of \(|\cdot|\) constructs, rather than postulates, the carry of
cancellation
\begin{equation}\label{eq:pdfv2-sol-kappa}
 \kappa_j
 :=\frac{E_{9,j}|b_f|-|b_c|}{2}\geq0.
\end{equation}
Indeed, using \(t^\pm=(|t|\pm t)/2\), one obtains
\begin{align}
 E_{9,j}b_f^+
 &=\frac{E_{9,j}|b_f|+E_{9,j}b_f}{2}
   =b_c^++\kappa_j,
 \label{eq:pdfv2-sol-jensen-plus}\\
 E_{9,j}b_f^-
 &=\frac{E_{9,j}|b_f|-E_{9,j}b_f}{2}
   =b_c^-+\kappa_j.
 \label{eq:pdfv2-sol-jensen-minus}
\end{align}
Both sheets lose the same amount when compressed; therefore, the signed
difference is preserved while total memory records the cancellation:
\begin{equation}\label{eq:pdfv2-sol-jensen-dh}
 E_{9,j}(b_f^+-b_f^-)=b_c,
 \qquad
 E_{9,j}(b_f^++b_f^-)=|b_c|+2\kappa_j.
\end{equation}

\section{Orthogonal factorization of the degrees of depth and of the resolution layers in \(\mathscr H_j^{\mathrm{depth}}\widehat\otimes\mathscr H_j^{\mathrm{res}}\)}
\label{sec:pdfv2-sol-micro-shell}

To prevent the same loss from being counted both as depth and as
resolution, the space at each level is typed as
\begin{equation}\label{eq:pdfv2-sol-factor-hilbert}
 \mathscr H_j
 =\mathscr H_j^{\rm depth}\widehat\otimes
  \mathscr H_j^{\rm res}.
\end{equation}
For \(j\geq1\), the first factor carries
\[
 P_j^{\rm depth}:=J_{9,j-1}E_{9,j-1}
 \quad\text{on }\mathscr H_j^{\rm depth},
\]
and \(P_0^{\rm depth}=I\) is fixed. The second factor carries the finite
shell partition inherited from the cylinder boundary. If
\(\sigma_j:Y_j^{\rm res}\to S_j\) is that partition, its operators are
\begin{align}
 J_j^{\rm res}f&=f\circ\sigma_j,
 &E_j^{\rm res}g(s)&=
 \frac1{\#\sigma_j^{-1}(s)}
 \sum_{\sigma_j(y)=s}g(y),
 \label{eq:pdfv2-sol-shell-je}\\
 P_j^{\rm res}&:=J_j^{\rm res}E_j^{\rm res}.
 \label{eq:pdfv2-sol-shell-p}
\end{align}

The three mutually orthogonal projectors are defined.
\begin{align}
 P_j^{\rm term}
 &:=P_j^{\rm depth}\otimes P_j^{\rm res},
 \label{eq:pdfv2-sol-pterm}\\
 Q_j^{\rm micro}
 &:=(I-P_j^{\rm depth})\otimes I,
 \label{eq:pdfv2-sol-qmicro}\\
 Q_j^{\rm shell}
 &:=P_j^{\rm depth}\otimes(I-P_j^{\rm res}).
 \label{eq:pdfv2-sol-qshell}
\end{align}
Then
\begin{equation}\label{eq:pdfv2-sol-ortogonalidad}
 Q_j^{\rm micro}Q_j^{\rm shell}=0,
 \qquad
 I=P_j^{\rm term}+Q_j^{\rm micro}+Q_j^{\rm shell},
\end{equation}
and, for all \(\xi\in\mathscr H_j\),
\begin{equation}\label{eq:pdfv2-sol-pitagoras}
 \|\xi\|^2
 =\|P_j^{\rm term}\xi\|^2
  +\|Q_j^{\rm micro}\xi\|^2
  +\|Q_j^{\rm shell}\xi\|^2.
\end{equation}
The term
\((I-P_j^{\rm depth})\otimes(I-P_j^{\rm res})\) does not appear a second time: it is already
contained exactly once in \(Q_j^{\rm micro}\).

\section{Cylinder extensions and full row of losses}
\label{sec:pdfv2-sol-gamma-loss}

An extension of cylinders from level \(j\) to level \(j+1\) induces the map
\begin{equation}\label{eq:pdfv2-sol-gamma}
 \Gamma_j:\mathscr H_j\longrightarrow\mathscr H_{j+1},
 \qquad
 \Gamma_j f=f\circ\tau_{j+1,j},
\end{equation}
with all leaf, orientation, carry, memory, and path labels
transported in the cylinder basis.  For \(n>j\), we write
\begin{equation}\label{eq:pdfv2-sol-gamma-composition}
 \Gamma_{j:n}:=\Gamma_{n-1}\cdots\Gamma_j,
 \qquad
 \Gamma_{j:j}:=I.
\end{equation}

On \(\mathscr H_j\), the nonnegative increments define multiplication
operators
\begin{equation}\label{eq:pdfv2-sol-loss-multipliers}
 B_j^+:=M_{\sqrt{b_j^+}},
 \qquad
 B_j^-:=M_{\sqrt{b_j^-}},
 \qquad
 K_j:=M_{\sqrt{2\kappa_j}}.
\end{equation}
Each root is taken pointwise over the finite set of cylinders; if a
coordinate is zero, its operator is zero.  The space of losses is the external
orthogonal sum
\[
 \mathscr E_j
 :=\mathscr H_j^{(+)}\oplus\mathscr H_j^{(-)}
   \oplus\mathscr H_j^{(\kappa)}
   \oplus\mathscr H_j^{({\rm micro})}
   \oplus\mathscr H_j^{({\rm shell})},
\]
and the complete row is
\begin{equation}\label{eq:pdfv2-sol-loss-row}
 L_j:\mathscr H_j\longrightarrow\mathscr E_j,
 \qquad
 L_j\xi
 =\bigl(
 B_j^+\xi,B_j^-\xi,K_j\xi,
 Q_j^{\rm micro}\xi,Q_j^{\rm shell}\xi
 \bigr).
\end{equation}
Therefore,
\begin{equation}\label{eq:pdfv2-sol-loss-norm}
 \|L_j\xi\|^2
 =\|B_j^+\xi\|^2+\|B_j^-\xi\|^2+\|K_j\xi\|^2
  +\|Q_j^{\rm micro}\xi\|^2
  +\|Q_j^{\rm shell}\xi\|^2.
\end{equation}
The positive part, negative part, cancellation carry, microstructure, and shell
each appear once. They are neither merged into an anonymous constant nor
repeated in two summands.

\section{Complete history terminal}
\label{sec:pdfv2-sol-terminal}

Let \(\operatorname{Hist}_{J}^{\rm sol}\) be the set of histories
surviving to the horizon \(J\), each with its package of thirteen
coordinates.  The terminal space is defined as the free space on these
complete records:
\begin{equation}\label{eq:pdfv2-sol-terminal-space}
 \mathscr T_J^{\rm sol}
 :=\ell^2\!\left(
 \left\{(h,D_{{\rm sol},J}(h)):
 h\in\operatorname{Hist}_{J}^{\rm sol}\right\}
 \right).
\end{equation}
The terminal evaluation
\begin{equation}\label{eq:pdfv2-sol-terminal-evaluation}
 E_J^{\rm term}:\mathscr H_J\longrightarrow\mathscr T_J^{\rm sol}
\end{equation}
sends each basic cylinder vector to the complete record of the history that
generates it and extends linearly.  In particular, the terminal retains
fiber, relations, numbers, orientation, carry, memory, boundary, route,
multisection, survival, and reader.  It is not replaced by the last visible
value of the route.

\section{Observation column, Gram operator, and telescoping}
\label{sec:pdfv2-sol-gram}

For \(0\leq j\leq J\), the complete future column is
\begin{equation}\label{eq:pdfv2-sol-g-column}
 \begin{aligned}
 G_{j;J}:\mathscr H_j&\longrightarrow
 \mathscr T_J^{\rm sol}\oplus
 \bigoplus_{n=j}^{J-1}\mathscr E_n,\\
 G_{j;J}\xi
 &:={}
 E_J^{\rm term}\Gamma_{j:J}\xi
 \ \oplus\ 
 \bigoplus_{n=j}^{J-1}L_n\Gamma_{j:n}\xi.
 \end{aligned}
\end{equation}
For \(j=J\), the sum of perdidas is empty and
\(G_{J;J}=E_J^{\rm term}\).  The Gram terminal is
\begin{equation}\label{eq:pdfv2-sol-u-gram}
 U_{j;J}:=G_{j;J}^*G_{j;J}\succeq0.
\end{equation}

Rearranging the orthogonal sum of the codomain yields the literal identity.
\begin{equation}\label{eq:pdfv2-sol-g-recursion}
 G_{j;J}\xi
 \simeq
 \bigl(L_j\xi,\,G_{j+1;J}\Gamma_j\xi\bigr).
\end{equation}
Consequently,
\begin{equation}\label{eq:pdfv2-sol-stein}
 \boxed{U_{j;J}
 =L_j^*L_j+\Gamma_j^*U_{j+1;J}\Gamma_j}
\end{equation}
and
\begin{equation}\label{eq:pdfv2-sol-stein-difference}
 U_{j;J}-\Gamma_j^*U_{j+1;J}\Gamma_j
 =L_j^*L_j\succeq0.
\end{equation}
There is no sign mutation: the loss term is on the positive side
of the identity because it arises from a sum of squares.

Iterating \eqref{eq:pdfv2-sol-stein} yields the exact telescoping identity
\begin{equation}\label{eq:pdfv2-sol-telescopia-operatorial}
 U_{j;J}
 =\Gamma_{j:J}^*(E_J^{\rm term})^*E_J^{\rm term}\Gamma_{j:J}
  +\sum_{n=j}^{J-1}
   \Gamma_{j:n}^*L_n^*L_n\Gamma_{j:n},
\end{equation}
o, in quadratic form,
\begin{equation}\label{eq:pdfv2-sol-telescopia-normas}
 \langle U_{j;J}\xi,\xi\rangle
 =\|E_J^{\rm term}\Gamma_{j:J}\xi\|^2
  +\sum_{n=j}^{J-1}\|L_n\Gamma_{j:n}\xi\|^2.
\end{equation}
The first term preserves the complete terminal future; the sum preserves,
level by level, all the losses and all the carries that led to it.

\iffalse
\section{The thirteen coordinates of the constraint}
\label{sec:pdfv2-sol-trece}

For a prefix \(x_{\leq j}\), it is defined
\begin{equation}\label{eq:pdfv2-sol-d-package}
 \begin{aligned}
 D_{{\rm sol},j}(x):=\bigl(&
 \mathcal F_{{\rm sol},j},
 \operatorname{Ext}_{\Gamma,{\rm sol},j},
 \operatorname{Rel}_{{\rm sol},j},
 \mathcal N_{{\rm sol},j},
 \operatorname{or}_{{\rm sol},j},
 \operatorname{Carr}_{{\rm sol},j},
 \operatorname{Mem}_{{\rm sol},j},\\
 &\partial_{{\rm sol},j},
 \gamma_{{\rm sol},j},
 \operatorname{Msect}_{{\rm sol},j},
 \operatorname{Surv}_{{\rm sol},j},
 \operatorname{Term}_{{\rm sol},j},
 \mathcal L_{{\rm sol},j}\bigr).
 \end{aligned}
\end{equation}
Its components are as follows.
\begin{enumerate}
 \item \(\mathcal F_{{\rm sol},j}\) contains the cylinders \(Y_j\), their projective
       measure, \(\mathscr H_j^{\rm depth}\),
       \(\mathscr H_j^{\rm res}\), their tensor products, and the five
       spaces of losses.
 \item \(\operatorname{Ext}_{\Gamma,{\rm sol},j}\) contains
       \(J_{9,j},E_{9,j},\Gamma_j,\Gamma_{j:n}\), the shell inclusions, and
       all the truncation and refinement maps.
 \item \(\operatorname{Rel}_{{\rm sol},j}\) contains
       \eqref{eq:pdfv2-sol-bpm-identidades},
       \eqref{eq:pdfv2-sol-ej-projector},
       \eqref{eq:pdfv2-sol-jensen-plus}--
       \eqref{eq:pdfv2-sol-jensen-minus},
       \eqref{eq:pdfv2-sol-ortogonalidad} and
       \eqref{eq:pdfv2-sol-stein}.
 \item \(\mathcal N_{{\rm sol},j}\) conserves
       \(j,J,9,b_j,b_j^+,b_j^-,\kappa_j\), sheet multiplicities, fiber sizes,
       and shell dimensions.
 \item \(\operatorname{or}_{{\rm sol},j}\) preserves
       \(\Delta_{\gamma_j}\), the path order, and the inversion action
       \(b_j\mapsto-b_j\), which interchanges \(b_j^+\) and \(b_j^-\) without changing
       \(|b_j|\) or \(\kappa_j\).
 \item \(\operatorname{Carr}_{{\rm sol},j}\) preserves the integer TPK carry, the composition carry and the cancellation carry \(\kappa_j\) as distinct coordinates.
 \item \(\operatorname{Mem}_{{\rm sol},j}\) is
       \((M_j^+,M_j^-,D_j^0,H_j^0)\) together with the ordered ledger of all
       preceding increments.
 \item \(\partial_{{\rm sol},j}\) contains the oriented cylinder boundary,
       the increments \(b_j\), the row \(L_j\), and the Stein differences.
 \item \(\gamma_{{\rm sol},j}\) is the complete route \(e_0e_1\cdots e_{j-1}\), not merely
its endpoints.
 \item \(\operatorname{Msect}_{{\rm sol},j}\) contains all compatible
       nonadic and shell extensions; it does not choose a prolongation.
 \item \(\operatorname{Surv}_{{\rm sol},j}\) records the horizons
       \(J\geq j\) for which \(G_{j;J}\), \(U_{j;J}\), and
       compatible terminal objects exist.
 \item \(\operatorname{Term}_{{\rm sol},j}\) contains
       \(E_J^{\rm term}\Gamma_{j:J}\), the complete record of history and
       the trece coordinates in each horizon.
 \item \(\mathcal L_{{\rm sol},j}\) internally publishes
       \((b^\pm,M^\pm,D^0,H^0,\kappa,Q^{\rm micro},Q^{\rm shell},
       L,G,U)\) and its identities.
\end{enumerate}

For an extension \(u:x_j\to x_{j'}\), the coordinated transport
\(u_{{\rm sol},j',j}\) satisfies
\begin{equation}\label{eq:pdfv2-sol-naturality-d}
 u_{{\rm sol},j',j}D_{{\rm sol},j'}(x_{j'})
 =D_{{\rm sol},j}(\tau_{j',j}x_{j'}).
\end{equation}
The equality includes all five components of \(L\), the complete terminal, and
the composition of \(\Gamma\); it is not an equality of cardinalities alone.

\section{Constraint graph and dependent object}
\label{sec:pdfv2-sol-graphs}

We define
\begin{equation}\label{eq:pdfv2-sol-g-graph}
 \mathcal G_{\rm sol}:=\operatorname{Graph}(D_{\rm sol}),
 \qquad
 \operatorname{Res}_{\rm sol}(x):=(x,D_{\rm sol}x).
\end{equation}
The first projection is inverse over the graph:
\begin{equation}\label{eq:pdfv2-sol-res-inverse}
 \pi_1\operatorname{Res}_{\rm sol}=I,
 \qquad
 \operatorname{Res}_{\rm sol}\pi_1=I_{\mathcal G_{\rm sol}}.
\end{equation}

Let \(\chi_{\rm sol}(x)\) be the complete multisection formed by all the
cylinders, refinements, shells, rows of losses, and terminals compatible
with \(x\).  Then
\begin{equation}\label{eq:pdfv2-sol-xchi}
 X_{\chi,{\rm sol}}
 :=\{(\chi_{\rm sol}(x),x,D_{\rm sol}x):
 x\in\mathcal C_{\rm cont}^{\rm sol}\},
\end{equation}
and
\begin{equation}\label{eq:pdfv2-sol-prx}
 \operatorname{pr}_{X,{\rm sol}}(x,D_{\rm sol}x)
 :=(\chi_{\rm sol}(x),x,D_{\rm sol}x).
\end{equation}
Again, the projection that preserves \((x,D_{\rm sol}x)\) is inverse on
this dependent image.

\section{Internal assembler and publisher}
\label{sec:pdfv2-sol-assembler-publisher}

The ambient space \(\mathscr M_{\rm sol}^{\rm amb}\) consists of compatible
systems
\[
 \bigl(
 \{b,b^\pm,M^\pm,D^0,H^0,\kappa\},
 \{E_9,J_9,P^{\rm term},Q^{\rm micro},Q^{\rm shell}\},
 \{\Gamma,L,E^{\rm term},G,U\}
 \bigr)
\]
with all the previous identities. The raw assembler is
\begin{equation}\label{eq:pdfv2-sol-assembler}
 \begin{aligned}
 a_{\rm sol}:X_{\chi,{\rm sol}}&\longrightarrow
 \mathscr M_{\rm sol}^{\rm amb},\\
 a_{\rm sol}(\chi(x),x,Dx)&:=
 \bigl(
 \{b,b^\pm,M^\pm,D^0,H^0,\kappa\}_x,
 \{E_9,J_9,P,Q\}_x,
 \{\Gamma,L,E^{\rm term},G,U\}_x
 \bigr).
 \end{aligned}
\end{equation}
The macrostructure preserves its antecedent through
\begin{equation}\label{eq:pdfv2-sol-macro-graph}
 \mathfrak M_{\rm sol}:=\operatorname{Graph}(a_{\rm sol}),
 \qquad
 \mathcal A_{\rm sol}(z):=(z,a_{\rm sol}z).
\end{equation}

Let \(\mathscr Y_{\rm sol}^{\rm int}\) be the type of certificates formed by
the identities
\begin{equation}\label{eq:pdfv2-sol-certificate-list}
 \begin{gathered}
 b=b^+-b^-,\qquad |b|=b^++b^-,\\
 E_9J_9=I,\qquad J_9E_9=P_9,\\
 E_9b_f^\pm=b_c^\pm+\kappa,\qquad\kappa\geq0,\\
 I=P^{\rm term}+Q^{\rm micro}+Q^{\rm shell},
 \qquad Q^{\rm micro}Q^{\rm shell}=0,\\
 U_{j;J}=L_j^*L_j+\Gamma_j^*U_{j+1;J}\Gamma_j,\\
 \langle U_{j;J}\xi,\xi\rangle
 =\|E_J^{\rm term}\Gamma_{j:J}\xi\|^2
  +\sum_{n=j}^{J-1}\|L_n\Gamma_{j:n}\xi\|^2,
 \end{gathered}
\end{equation}
together with the naturality of all arrows. The raw publisher
\begin{equation}\label{eq:pdfv2-sol-publisher}
 p_{\rm sol}^{\rm raw}:\mathfrak M_{\rm sol}
 \longrightarrow\mathscr Y_{\rm sol}^{\rm int}
\end{equation}
publish that certificate without removing the macroobject. Finally,
\begin{equation}\label{eq:pdfv2-sol-output-graph}
 \mathcal C_{\rm sol}^{\rm HMT}
 :=\operatorname{Graph}(p_{\rm sol}^{\rm raw}),
 \qquad
 \mathcal P_{\rm sol}(m):=(m,p_{\rm sol}^{\rm raw}m).
\end{equation}

\section{Solenoidal Generation Internal Theorem}
\label{sec:pdfv2-sol-theorem}

\begin{theorem}[Internal generation of the balance macrostructure]
\label{thm:pdfv2-sol-internal-generation}
For every \(x\in\mathcal C_{\rm cont}^{\rm sol}\), the chain
\[
 \mathcal C_{\rm cont}^{\rm sol}
 \xrightarrow{\operatorname{Res}_{\rm sol}}
 \mathcal G_{\rm sol}
 \xrightarrow{\operatorname{pr}_{X,{\rm sol}}}
 X_{\chi,{\rm sol}}
 \xrightarrow{\mathcal A_{\rm sol}}
 \mathfrak M_{\rm sol}
 \xrightarrow{\mathcal P_{\rm sol}}
 \mathcal C_{\rm sol}^{\rm HMT}
\]
naturally constructs a macrostructure containing both memories, the cancellation carry, the orthogonal micro/shell separation, the complete row of losses, the future column, the full-history terminal object, the Gram matrix, and the telescoping structure. The thirteen coordinates of the continuum can be recovered by successive projections from any point in the chain.
\end{theorem}

\begin{proof}
The identities \eqref{eq:pdfv2-sol-bpm-identidades} and
\eqref{eq:pdfv2-sol-dh-diferencias} are obtained directly from the
positive and negative decomposition of the integer increment; therefore the
memories are not added data.  The formulas
\eqref{eq:pdfv2-sol-j9}--\eqref{eq:pdfv2-sol-e9}, together with the projective
measure, prove \(E_9J_9=I\), \(J_9^*=E_9\), and that \(P_9\) is an
orthogonal projector.

The triangle inequality gives
\(E_9|b_f|\geq|E_9b_f|=|b_c|\), so that
\(\kappa\geq0\).  Substituting
\(b^\pm=(|b|\pm b)/2\) proves
\eqref{eq:pdfv2-sol-jensen-plus} and
\eqref{eq:pdfv2-sol-jensen-minus} without an additional hypothesis.

The tensorial definitions
\eqref{eq:pdfv2-sol-pterm}--\eqref{eq:pdfv2-sol-qshell} prove that
\(P^{\rm term}\), \(Q^{\rm micro}\), and \(Q^{\rm shell}\) are orthogonal and
sum to the identity. Consequently, row \(L_j\) counts each sector only
once.

The definition \eqref{eq:pdfv2-sol-g-column} orthogonally separates the terminal
and the losses of each level.  Separating the first summand of losses produces
\eqref{eq:pdfv2-sol-g-recursion}.  Taking the adjoint and composing proves the
Stein identity \eqref{eq:pdfv2-sol-stein}; iterating it proves
\eqref{eq:pdfv2-sol-telescopia-operatorial} and
\eqref{eq:pdfv2-sol-telescopia-normas}.  Every term is a squared norm,
so the opposite sign cannot appear.

Compatibility of cylinder truncation, expectation, projections, and
extensions \(\Gamma\) yields the naturality
\eqref{eq:pdfv2-sol-naturality-d}. Finally,
\(\mathcal G_{\rm sol}\), \(\mathfrak M_{\rm sol}\), and
\(\mathcal C_{\rm sol}^{\rm HMT}\) are graphs, and
\(X_{\chi,{\rm sol}}\) is a dependent object that explicitly preserves
its antecedent. The first projection at each stage recovers the preceding
stage. Thus no composition erases history, memory, multisection,
terminal, or reader.
\end{proof}

\section{Probative origin and forgers}
\label{sec:pdfv2-sol-provenance-falsifiers}

\paragraph{Result recovered.}
The signed increment with carry, the separation \(b=b^+-b^-\), the memories
\(M^\pm,D^0,H^0\), the pair \(E_9,J_9\), the Jensen carry \(\kappa\), cylinder
extension, the loss row, the terminal, the Gram matrix, Stein, and telescoping
belong to the already constructed architecture.

\paragraph{New Formalization.}
Explicit formalizations of that architecture are the typing prior to every
external evaluation, the factorization
\(\mathscr H^{\rm depth}\widehat\otimes\mathscr H^{\rm res}\), the orthogonal
separation preventing double counting, the free terminal on histories with
thirteen coordinates, and the four nonablative graph packagings.

\begin{criterion}[Internal falsifiers of the branch \(\mathrm{sol}\)]
\label{crit:pdfv2-sol-falsifiers}
The construction fails if any of the following occurs:
\begin{enumerate}
 \item \(b\neq b^+-b^-\) o \(|b|\neq b^++b^-\);
 \item \(E_9J_9\neq I\), the measure is not projective or \(P_9\) is not a
       projector;
 \item \(\kappa<0\) or fail
       \(E_9b_f^\pm=b_c^\pm+\kappa\);
 \item \(Q^{\rm micro}Q^{\rm shell}\neq0\), or the same sector appears in
       both summands of row \(L\);
 \item one of \(b^+,b^-,\kappa,Q^{\rm micro},Q^{\rm shell}\) is omitted or
       counted twice;
 \item the terminal preserves only the visible value and not the
history with its thirteen coordinates;
 \item the Stein identity fails or the telescoping contains a sign
       contrary to a sum of squares;
 \item an extension does not commute with \(E_9,J_9,L,G,U\);
 \item the multisection of continuations is replaced with a retrospectively chosen prolongation;
 \item any of the thirteen coordinates is omitted.
\end{enumerate}
In the last case, the binding diagnosis is: \emph{object replaced;
conclusion not transferable}.
\end{criterion}

\begin{warningbox}
This chapter concludes only the internal construction of
\(\mathcal C_{\rm sol}^{\rm HMT}\). It contains no classical differential
equation, classical initial data, map from external data, or
conclusion concerning an external problem. Those intertwiners belong to another
document and do not govern the existence of this macrostructure.
\end{warningbox}
\fi
 \endgroup

\section{Orientation, carry, and simultaneous composition of paths}

Path inversion acts only once on
\(\mathfrak e_\alpha\). Let
\(\mathsf J(v_\Sigma,v_\Pi):=(v_\Pi,v_\Sigma)\) be the leaf exchange,
and let \(P_{\rm tr}\) be the transposition of port matrices. Its five
local expressions are
\begin{equation}
\label{eq:pdfv2-cor-orientacion-cruzada}
 \begin{gathered}
 \mathsf J\sigma_k=-\sigma_k\mathsf J,
 \qquad b_{\bar\alpha}=-b_\alpha,
 \qquad (b_{\bar\alpha}^+,b_{\bar\alpha}^-)
       =(b_\alpha^-,b_\alpha^+),\\
 \vartheta_L(K,\epsilon):=
 \begin{cases}
  (L,-\epsilon),&K=L,\\
  (K,\epsilon),&K\neq L,
 \end{cases}
 \qquad
 \kappa_{J,I}(v,u)=-P_{\rm tr}\kappa_{I,J}(u,v),\\
 \kappa(\bar\alpha)
 =-\mathsf C_\alpha^{-1}\kappa(\alpha),
 \qquad v_{\bar\alpha}=-v_\alpha.
 \end{gathered}
\end{equation}
Here \(\mathsf J\) exchanges the two sheets, \(P_{\rm tr}\) transposes the incidence matrix, and \(\vartheta_L\) exchanges exactly the two faces in the direction \(L\), fixing the other six. In the free basis of \(\mathscr B=\mathscr I\times\{+,-\}\), the eight vectors are distinct, \(\vartheta_L^2=I\) and \(\vartheta_L\neq\vartheta_{L'}\) for \(L\neq L'\). The same orientation explains all signs in \eqref{eq:pdfv2-cor-orientacion-cruzada}.

For composable extensions \(\alpha:x\to y\) and \(\beta:y\to z\), the
native carry satisfies
\begin{equation}
\label{eq:pdfv2-cor-carry-nativo}
 \kappa(\beta\alpha)=
 \kappa(\beta)+\mathsf C(\beta)\kappa(\alpha).
\end{equation}
The matrix lift is not left as a name. In the ordered basis
\((f,r,e,b,g)\), let
\(\widetilde L_N(\alpha)\in\operatorname{Mat}_5(\mathbb Z)\) be the
centered integral representative of the matrix of
\(\Psi_\alpha\) of \eqref{eq:pdfv2-cor-psi-generadores} modulo \(3^N\), and
define
\begin{equation}
\label{eq:pdfv2-cor-carry-matricial-def}
 c_N(\beta,\alpha):=
 \frac{\widetilde L_N(\beta)\widetilde L_N(\alpha)
       -\widetilde L_N(\beta\alpha)}{3^N}
 \in\operatorname{Mat}_5(\mathbb Z).
\end{equation}
Divisibility follows from equality of the two compositions modulo
\(3^N\). Writing \(L_N=[\widetilde L_N]_{3^N}\), the integer lifts
satisfy
\begin{equation}
\label{eq:pdfv2-cor-carry-matricial}
 \widetilde L_N(\beta)\widetilde L_N(\alpha)
 =\widetilde L_N(\beta\alpha)+3^Nc_N(\beta,\alpha)
\end{equation}
and associativity produces the cocycle
\begin{equation}
\label{eq:pdfv2-cor-cociclo-matricial}
 c_N(\delta,\beta\alpha)+\widetilde L_N(\delta)c_N(\beta,\alpha)
 =c_N(\delta\beta,\alpha)
  +c_N(\delta,\beta)\widetilde L_N(\alpha).
\end{equation}

The same carry increment acts upon the local complex through
\begin{equation}
\label{eq:pdfv2-cor-psi-generadores}
 \begin{gathered}
 \Psi_\alpha(r_x)=r_y,\quad
 \Psi_\alpha(e_x)=e_y,\quad
 \Psi_\alpha(g_x)=g_y,\quad
 \Psi_\alpha(f_x)=f_y,\\
 \Psi_\alpha(b_x)=b_y+3(c_y-c_x)e_y.
 \end{gathered}
\end{equation}
The oriented boundary change remains in
\begin{equation}
\label{eq:pdfv2-cor-k-alpha}
 K_\alpha:=(v_y-v_x)f_y,
 \qquad
 K_{\beta\alpha}=K_\beta+\Psi_\beta K_\alpha.
\end{equation}
The equations \eqref{eq:pdfv2-cor-carry-nativo},
\eqref{eq:pdfv2-cor-cociclo-matricial}, and
\eqref{eq:pdfv2-cor-k-alpha} are three degrees of the same ledger of composition.
None authorizes erasing the other two.

\subsection{Nerve of extensions and Alexander--Whitney diagonal}
For a history of the common register one writes
\(A_n:=\mathbf Z/3^n\mathbf Z\) and \(\mathsf T_m(x)\) for the finite
category of \emph{all} its TPK prefixes of depth \(m\), with the
extension morphisms already constructed. The diagonal and counit that follow
apply to that same nerve; they do not define a fifth autonomous branch.
\begingroup
  \let\section\subsection
  % Extracción quirúrgica del fuente teoremática V2 05e: líneas 68--110.
% El resto permanece en este archivo como procedencia, pero no se publica.
\iffalse
\chapter[Internal Determinantal Macrostructure]{Internal Determinantal Macrostructure Generated by the Discrete Continuum}
\label{chap:detint-macroestructura}

\begin{thesisbox}
The \({\rm det}\) branch starts exclusively from surviving histories of
\(\mathcal C_{\rm cont}^{\rm disc}\). The nerve of extensions, the Alexander--Whitney diagonal,
the local complex, the unit, the two publications, the filler, the terminal
reduction, and the \(3\)-adic tower are constructed before any subsequent
evaluation. The balance and acyclicity conditions are incorporated into the
typed domain and are never tacitly assumed.
\end{thesisbox}

\section{Surviving determinant domain}

Let
\begin{equation}\label{eq:detint-an}
 A_n:=\mathbf Z/3^n\mathbf Z,
 \qquad
 \rho_{n+1,n}:A_{n+1}\longrightarrow A_n.
\end{equation}
For \(x\in\mathcal C_{\rm cont}^{\rm disc}\), let
\(\mathsf T_m(x)\) be the finite category of TPK prefixes of depth \(m\)
compatible with \(x\). Each object preserves sheet, orientation, residue,
quotient, integer carry, memory, boundary, route, multisection, survival,
and terminal; its morphisms are the admissible extensions that preserve those
coordinates.

From each object \(y\in\mathsf T_m(x)\), the internal integers are extracted
\begin{equation}\label{eq:detint-c-v-lifts}
 \widetilde c_y,\widetilde v_y\in\mathbf Z,
 \qquad
 c_{y,n}:=\widetilde c_y\bmod3^n,
 \qquad
 v_{y,n}:=\widetilde v_y\bmod3^n.
\end{equation}
The first coordinate is the unoriented carry and the second, the oriented
boundary coefficient. Reversal satisfies
\begin{equation}\label{eq:detint-or-cv}
 c_{\bar y,n}=c_{y,n},
 \qquad
 v_{\bar y,n}=-v_{y,n}.
\end{equation}

We write
\begin{equation}\label{eq:detint-z3-k3}
 \mathbf Z_3:=\varprojlim_n A_n,
 \qquad
 K_3:=\operatorname{Frac}(\mathbf Z_3).
\end{equation}
The determinantial restriction is defined by
\begin{equation}\label{eq:detint-dominio}
\begin{split}
 \mathcal C_{\rm cont}^{\rm det}:=
 \bigl\{x\in\mathcal C_{\rm cont}^{\rm disc}:{}&
 x\text{ survives at every depth};\\
 &\widehat P_x^{\rm red}\text{ is finite and based};\\
 &\chi(\widehat P_x^{\rm red})=0;\\
 &H_*(\widehat P_x^{\rm red}\otimes_{\mathbf Z_3}K_3)=0;\\
 &\text{the canonical reduction produces a square block}\\
 &\text{stable under refinement}\bigr\}.
\end{split}
\end{equation}
The objects \(\widehat P_x^{\rm red}\) and the square block will be constructed
in \cref{sec:detint-terminal}. Therefore
\eqref{eq:detint-dominio} is a verifiable membership condition and not
a conclusion inserted into the definition of the maps.

\fi
\section{Nerve of routes, Alexander--Whitney diagonal and counit}

The normalized complex is defined
\begin{equation}\label{eq:detint-gmn}
 G_{m,n}(x):=
 N_*^{\rm norm}\!\left(N\mathsf T_m(x);A_n\right).
\end{equation}
For a nondegenerate simplex
\(\sigma=[y_0\to y_1\to\cdots\to y_q]\), the diagonal is
\begin{equation}\label{eq:detint-aw}
 \Delta_{\rm AW}(\sigma)
 =\sum_{j=0}^q
 [y_0\to\cdots\to y_j]\otimes
 [y_j\to\cdots\to y_q].
\end{equation}
The counit is given by
\begin{equation}\label{eq:detint-counidad-g}
 \epsilon_G([y])=1,
 \qquad
 \epsilon_G(\sigma)=0\quad(q>0).
\end{equation}
In particular,
\begin{equation}\label{eq:detint-vertice-grouplike}
 \Delta_{\rm AW}[y]=[y]\otimes[y],
 \qquad
 \epsilon_G([y])=1.
\end{equation}

If \(m'\ge m\) and \(n'\ge n\), truncation and coefficient reduction
induce
\begin{equation}\label{eq:detint-rmn}
 R_{m',m}^{n',n}:G_{m',n'}(x)\longrightarrow G_{m,n}(x),
\end{equation}
and they are verified
\begin{equation}\label{eq:detint-aw-natural}
 R\partial=\partial R,
 \qquad
 (R\otimes R)\Delta_{\rm AW}=\Delta_{\rm AW}R,
 \qquad
 \epsilon_GR=\rho_{n',n}\epsilon_G.
\end{equation}
The path is not reduced to its endpoints: the entire simplex remains in
\(G_{m,n}(x)\).

\iffalse
\section{Internal local complex}

For \(c\in A_n\), define
\begin{equation}\label{eq:detint-p0-grados}
 P^0_{n,1}(c)=A_nf,
 \qquad
 P^0_{n,0}(c)=A_nr\oplus A_ne\oplus A_nb,
 \qquad
 P^0_{n,-1}(c)=A_ng,
\end{equation}
with differential
\begin{equation}\label{eq:detint-p0-diferencial}
 \partial f=3ce+b,
 \qquad
 \partial r=0,
 \qquad
 \partial e=g,
 \qquad
 \partial b=-3cg.
\end{equation}
The chain identity is literal:
\begin{equation}\label{eq:detint-d2}
 \partial^2f=3c\,\partial e+\partial b=3cg-3cg=0,
 \qquad
 \partial^2e=\partial^2b=\partial^2r=0.
\end{equation}

The increase
\begin{equation}\label{eq:detint-epsilon-p}
 \epsilon_P:P_n^0(c)\longrightarrow A_n[0]
\end{equation}
is fixed by
\[
 \epsilon_P(r)=1,
 \qquad
 \epsilon_P(e)=\epsilon_P(b)=\epsilon_P(f)=\epsilon_P(g)=0.
\]
The complex possesses the differential coalgebra.
\begin{equation}\label{eq:detint-coalgebra-p}
 \Delta_P(r)=r\otimes r,
 \qquad
 \Delta_P(q)=q\otimes r+r\otimes q
 \quad(q\in\{e,b,f,g\}).
\end{equation}
Thus, \(r\) is group-like and the other four generators are primitive
relative to \(r\). Equations
\eqref{eq:detint-p0-diferencial} prove that \(\Delta_P\) commutes with the
tensor differential.

\section{Local extension system and memory ledger}

For an extension \(\alpha:y\to y'\), we define
\begin{equation}\label{eq:detint-psi}
 \Psi_\alpha:P_n^0(c_y)\longrightarrow P_n^0(c_{y'})
\end{equation}
through means of
\begin{equation}\label{eq:detint-psi-generadores}
\begin{gathered}
 \Psi_\alpha(r_y)=r_{y'},\quad
 \Psi_\alpha(e_y)=e_{y'},\quad
 \Psi_\alpha(g_y)=g_{y'},\quad
 \Psi_\alpha(f_y)=f_{y'},\\
 \Psi_\alpha(b_y)=b_{y'}+3(c_{y'}-c_y)e_{y'}.
\end{gathered}
\end{equation}
It is a map of complexes because
\begin{equation}\label{eq:detint-psi-b-chain}
 \partial\Psi_\alpha(b_y)
 =-3c_{y'}g+3(c_{y'}-c_y)g
 =-3c_yg
 =\Psi_\alpha(\partial b_y),
\end{equation}
and
\begin{equation}\label{eq:detint-psi-f-chain}
 \Psi_\alpha(\partial f_y)
 =3c_ye+b+3(c_{y'}-c_y)e
 =3c_{y'}e+b
 =\partial\Psi_\alpha(f_y).
\end{equation}
Moreover,
\begin{equation}\label{eq:detint-psi-functorial}
 \Psi_\beta\Psi_\alpha=\Psi_{\beta\alpha}.
\end{equation}

The oriented coordinate change does not erase. It is recorded by
\begin{equation}\label{eq:detint-k-alpha}
 K_\alpha:=(v_{y'}-v_y)f_{y'}.
\end{equation}
For composable extensions,
\begin{equation}\label{eq:detint-k-coherencia}
 K_{\beta\alpha}=K_\beta+\Psi_\beta K_\alpha.
\end{equation}
This equality is the exact ledger of the boundary increment and makes visible
the memory that a purely nominal composition would lose.

The complete local diagram is the Grothendieck construction.
\begin{equation}\label{eq:detint-grothendieck-local}
 \int_{y\in\mathsf T_m(x)}P_n^0(c_y).
\end{equation}
In it, an arrow over \(\alpha:y\to y'\) carries \(\Psi_\alpha\)
as its linear component and \(K_\alpha\) as its coherence. When applied to an
element supported at a vertex, the notation is used
\begin{equation}\label{eq:detint-psi-soporte-vertice}
 \Psi_\alpha(p,[y]):=(\Psi_\alpha p,[y']).
\end{equation}
The edge \([y\to y']\) remains in the path coordinate of the nerve. This
convention will be used in the naturality formulas for
\(z_\pm\) and \(h_*\); no global map erasing the remaining
simplices of \(G_{m,n}(x)\) is asserted.

\section{Pullback, unit, root, publications and filler}

The pullback of complexes is
\begin{equation}\label{eq:detint-pgen}
 P^{\rm gen}_{y;m,n}
 :=P_n^0(c_y)\times_{A_n[0]}G_{m,n}(x)
 =\{(p,\gamma):\epsilon_P(p)=\epsilon_G(\gamma)\}.
\end{equation}
Its differential is
\begin{equation}\label{eq:detint-pgen-d}
 \partial(p,\gamma)=(\partial p,\partial\gamma).
\end{equation}
The common unit associated with the vertex \(y\) is
\begin{equation}\label{eq:detint-unidad}
 \mathbf1_y:=(r,[y]).
\end{equation}
Its two projections are group-like by
\eqref{eq:detint-vertice-grouplike} and
\eqref{eq:detint-coalgebra-p}; that is the precise meaning of a compatible
group-like unit of the pullback.

The root projection is the map of complexes
\begin{equation}\label{eq:detint-proyeccion-root}
 \varpi_{\rm root}:P^{\rm gen}_{y;m,n}\longrightarrow
 \operatorname{Root}_{y;m,n},
 \qquad
 \varpi_{\rm root}(p,\gamma):=(\epsilon_P(p)r,\gamma).
\end{equation}
The two internal publications are defined.
\begin{equation}\label{eq:detint-zpm}
 z_-(y):=(r,[y]),
 \qquad
 z_+(y):=(r+3c_yv_ye+v_yb,[y]).
\end{equation}
Both are cycles:
\begin{equation}\label{eq:detint-zpm-ciclos}
 \partial z_-(y)=0,
 \qquad
 \partial z_+(y)=3c_yv_yg-3c_yv_yg=0.
\end{equation}
Their roots coincide:
\begin{equation}\label{eq:detint-raiz-comun}
 \varpi_{\rm root}z_-(y)
 =(r,[y])
 =\varpi_{\rm root}z_+(y).
\end{equation}

The filler is constructed directly:
\begin{equation}\label{eq:detint-hstar}
 h_*(y):=(-v_yf,0)\in P^{\rm gen}_{y;m,n,1}.
\end{equation}
Then
\begin{equation}\label{eq:detint-filler-identidad}
 \partial h_*(y)
 =(-3c_yv_ye-v_yb,0)
 =z_-(y)-z_+(y).
\end{equation}
In particular,
\begin{equation}\label{eq:detint-clases-iguales}
 [z_-(y)]=[z_+(y)]
 \quad\text{in}\quad H_0(P^{\rm gen}_{y;m,n}),
\end{equation}
and the equality preserves its witness chain.

Under an extension \(\alpha:y\to y'\), we have
\begin{equation}\label{eq:detint-z-natural}
 z_-(y')=\Psi_\alpha z_-(y),
 \qquad
 z_+(y')-\Psi_\alpha z_+(y)=\partial K_\alpha,
\end{equation}
and
\begin{equation}\label{eq:detint-h-natural}
 h_*(y')-\Psi_\alpha h_*(y)=-K_\alpha.
\end{equation}
Therefore, naturality does not identify distinct values of \(v\): it records
their difference through the coherent homotopy \(K_\alpha\).

\section{Internal Orientation}

On \(P_n^0(c)\) is defined
\begin{equation}\label{eq:detint-omega-p}
 \Omega(r)=r,
 \qquad
 \Omega(q)=-q\quad(q\in\{e,b,f,g\}).
\end{equation}
It is an automorphism of complexes and of the relative coalgebra. In the nerve,
the oriented inversion is
\begin{equation}\label{eq:detint-omega-g}
 \mathfrak o_G[y_0\to\cdots\to y_q]
 :=(-1)^{q(q+1)/2}
 [\bar y_q\to\cdots\to\bar y_0].
\end{equation}
Together with \eqref{eq:detint-or-cv}, these actions give
\begin{equation}\label{eq:detint-or-z-h}
 \Omega z_-(y)=z_-(\bar y),
 \qquad
 \Omega z_+(c_y,v_y)=z_+(c_{\bar y},v_{\bar y}),
 \qquad
 \Omega h_*(y)=h_*(\bar y).
\end{equation}
Orientation acts on path, complex, publications, filler and determinant
line; it is not a label without an action.

\section{Terminal reduction and survival condition}
\label{sec:detint-terminal}

Only the common root is eliminated:
\begin{equation}\label{eq:detint-p-reducido}
 \overline P_{y;m,n}
 :=P^{\rm gen}_{y;m,n}/A_n\mathbf1_y,
 \qquad
 \widehat P_{y;m}^{\rm red}:=\varprojlim_n\overline P_{y;m,n}.
\end{equation}
The base is first ordered by the TPK order of routes and then by
\begin{equation}\label{eq:detint-orden-base}
 f\prec e\prec b\prec g.
\end{equation}
The following internal algorithm is executed.
\begin{enumerate}
 \item Find the coefficients that are units of \(\mathbf Z_3\).
 \item Choose the first according to \eqref{eq:detint-orden-base} and the order of
 paths.
 \item Cancel the corresponding pair and record the elementary
 operations, their sign, and their contribution to the orientation.
 \item Repeat until no cancellable unit pivot remains.
\end{enumerate}
The survival condition of \eqref{eq:detint-dominio} requires that,
cofinally in \(m\), the result be a two-term block
\begin{equation}\label{eq:detint-sx}
 \mathbf Z_3^{r_x}\xrightarrow{S_x}\mathbf Z_3^{r_x},
 \qquad
 \det S_x\ne0\text{ in }K_3,
\end{equation}
and that a refinement only add unit contractible blocks and recorded
changes of basis. If these conditions fail, the history does not belong to
\(\mathcal C_{\rm cont}^{\rm det}\); a nonexistent terminal is not forced.

The identity route provides an elementary witness. After removing the root,
the local complex is
\begin{equation}\label{eq:detint-testigo-contractil}
 A_nf\xrightarrow{\binom{3c}{1}}
 A_ne\oplus A_nb\xrightarrow{(1,-3c)}A_ng.
\end{equation}
Is contractible via
\begin{equation}\label{eq:detint-contraccion-local}
 s(g)=e,
 \qquad
 s(ae+\beta b)=\beta f,
 \qquad
 s(f)=0.
\end{equation}
Thus, the terminal conditions admit at least the internal identity witness
and do not describe a formally empty subdomain.

\section{Smith, \texorpdfstring{\(3\)}{3}-adic order and initial unit}

In oriented bases, a valuation diagonalization has the form
\begin{equation}\label{eq:detint-smith}
 U_xS_xV_x
 =\operatorname{diag}(3^{a_1}u_1,\ldots,3^{a_r}u_r),
 \qquad
 a_i\in\mathbf N,quad u_i\in\mathbf Z_3^\times.
\end{equation}
We define
\begin{equation}\label{eq:detint-orden-d}
 d_x:=v_3(\det S_x)=\sum_{i=1}^r a_i
\end{equation}
and the initial unit
\begin{equation}\label{eq:detint-leading-unit}
 \lambda_x:=3^{-d_x}\det S_x\in\mathbf Z_3^\times.
\end{equation}
The orientation coordinate trivializes the determinant line. If a
transition produces
\begin{equation}\label{eq:detint-estabilidad-bloques}
 S_x\oplus I_q=U S_{x'}V,
\end{equation}
the orientation is transported by the inverse factor of \(\det U\det V\).
The oriented unit \(\lambda_x^{\rm or}\) is then natural and does not change
when contractible blocks are added.

At finite precision,
\begin{equation}\label{eq:detint-smith-finito}
 S_{x,n}=S_x\bmod3^n,
 \qquad
 a_i^{(n)}=\min(a_i,n).
\end{equation}
The terminal is, therefore, the compatible family of all precisions
and not an isolated integer.

\section{Bockstein Tower and Higher Order Carry}

Let
\[
 C_x^1\xrightarrow{S_x}C_x^0
\]
the terminal block. If a class \([\bar y]\) admits a lift
\(y\in C_x^1\) with \(S_xy\in3^qC_x^0\), we define
\begin{equation}\label{eq:detint-beta-q}
 \beta_q[\bar y]
 :=\left[3^{-q}S_xy\bmod3\right].
\end{equation}
Changing the lift alters the representative by a boundary of the corresponding
page, so \(\beta_q\) is well defined. The higher
carry is
\begin{equation}\label{eq:detint-carry-superior}
 \operatorname{Carr}_q(y):=3^{-q}S_xy.
\end{equation}
It does not replace the TPK carry of \eqref{eq:detint-c-v-lifts}; it records its
prolongation in the coefficient tower.

For a block \(3^{a_i}u_i\), the class persists until page \(a_i\),
and on that page the nonzero differential is multiplication by
\(\bar u_i\in\mathbf F_3^\times\). Consequently,
\begin{equation}\label{eq:detint-bock-longitudes}
 \{\text{survival lengths}\}=\{a_1,\ldots,a_r\},
 \qquad
 d_x=\sum_i a_i.
\end{equation}
The oriented product of the page trivializations satisfies
\begin{equation}\label{eq:detint-bock-leading}
 \Theta_{\rm Bock}(x)=\overline{\lambda_x^{\rm or}}
 \in\mathbf F_3^\times.
\end{equation}
The proof is carried out in each diagonal block and transported by the recorded
unimodular changes. Smith, Bockstein and the initial unit are thus three
readings of the same terminal.

\section{The thirteen coordinates of the determinantal constraint}

For each \((m,n)\), it is defined
\begin{equation}\label{eq:detint-d-trece}
 D_{{\rm det},m,n}(x):=
 (\mathcal F,\operatorname{Ext}_\Gamma,\operatorname{Rel},
 \mathcal N,\operatorname{or},\operatorname{Carr},\operatorname{Mem},
 \partial,\gamma,\operatorname{Msect},\operatorname{Surv},
 \operatorname{Term},\mathcal L)_{{\rm det},m,n}(x),
\end{equation}
with the following content.
\begin{enumerate}
 \item The fiber is
 \[
  \mathcal F=(A_n,\mathsf T_m(x),
  \{P_n^0(c_y)\}_y,\{P^{\rm gen}_{y;m,n}\}_y).
 \]
 \item \(\operatorname{Ext}_\Gamma\) contains all morphisms
 \(\alpha\), the maps \(\Psi_\alpha\), the homotopies \(K_\alpha\) and
 the maps \(R_{m',m}^{n',n}\).
 \item \(\operatorname{Rel}\) contains \(\partial^2=0\), the
 pullback equation, AW, counit, \(\Psi_\beta\Psi_\alpha=\Psi_{\beta\alpha}\) and
 \(K_{\beta\alpha}=K_\beta+\Psi_\beta K_\alpha\).
 \item
 \(\mathcal N=(m,n,\widetilde c_y,\widetilde v_y,c_{y,n},v_{y,n},
 r_x,a_1,\ldots,a_{r_x},d_x)\).
 \item \(\operatorname{or}\) contains \(y\mapsto\bar y\),
 \(\mathfrak o_G\), \(\Omega\) and the orientation of the determinant
 line.
 \item \(\operatorname{Carr}\) separately preserves the integer carry,
 its reduction \(c_{y,n}\) and all higher carries
 \(\operatorname{Carr}_q\).
 \item \(\operatorname{Mem}\) contiene
 \[
  y_0\to\cdots\to y_m,
  \quad\{(\widetilde c_{y_j},\widetilde v_{y_j},\mathbf1_{y_j})\}_{j=0}^m,
  \quad\{K_{\alpha_j}\}.
 \]
 \item
 \(\partial=(\partial_G,\partial_P,\partial_{P^{\rm gen}},
 \partial_{\overline P},h_*)\).
 \item \(\gamma=[y_0\to\cdots\to y_q]\) preserves the complete oriented
 simplex.
 \item \(\operatorname{Msect}\) contains all compatible
 extensions, all \(3\)-adic lifts, and all unimodular
 presentations of the same oriented terminal; it does not choose one a posteriori.
 \item \(\operatorname{Surv}\) records arbitrary depth,
 coefficient compatibility, balance, acyclicity over \(K_3\),
 stability under unit blocks and persistence of the root.
 \item
 \[
  \operatorname{Term}=(\widehat P_x^{\rm red},S_x,
  \{a_i,u_i\},d_x,\lambda_x^{\rm or},
  \{\beta_q\},\Theta_{\rm Bock}).
 \]
 \item The internal reader is
 \[
  \mathcal L_{\rm det}^{\rm int}
  =(\mathbf1_y,
  \varpi_{\rm root}z_-=\varpi_{\rm root}z_+,
  [z_-]=[z_+],d_x,\lambda_x^{\rm or},\Theta_{\rm Bock}).
 \]
\end{enumerate}

The transitions satisfy, coordinate by coordinate,
\begin{equation}\label{eq:detint-naturalidad-trece}
 u^{\rm det}_{(m',n'),(m,n)}D_{{\rm det},m',n'}(x')
 =D_{{\rm det},m,n}(\tau_{m',m}x').
\end{equation}

\section{Assembler, publisher, and non-ablative chain}

We define
\begin{equation}\label{eq:detint-graph-d}
 \mathcal G_{\rm det}:=\operatorname{Graph}(D_{\rm det}),
 \qquad
 \operatorname{Res}_{\rm det}(x):=(x,D_{\rm det}x).
\end{equation}
Let \(\chi_{\rm det}(x)\) be the complete multisection of germs
\[
 (y,\widetilde c_y,\widetilde v_y,\operatorname{or}_y,\gamma_y)
\]
compatible at all depths. The dependent object is
\begin{equation}\label{eq:detint-x-dependiente}
 X_{\chi,{\rm det}}
 :=\{(\chi_{\rm det}(x),x,D_{\rm det}x):
 x\in\mathcal C_{\rm cont}^{\rm det}\},
\end{equation}
and
\begin{equation}\label{eq:detint-prx}
 \operatorname{pr}_{X,{\rm det}}(x,D_{\rm det}x)
 :=(\chi_{\rm det}(x),x,D_{\rm det}x).
\end{equation}

The raw assembler
\begin{equation}\label{eq:detint-ensamblador-crudo}
 a_{\rm det}:X_{\chi,{\rm det}}\longrightarrow
 \mathscr M_{\rm det}^{\rm amb}
\end{equation}
yields
\begin{equation}\label{eq:detint-ensamblador-salida}
\begin{split}
 a_{\rm det}(z)=\bigl(&
 \{G_{m,n},\Delta_{\rm AW},\epsilon_G\}_{m,n};\\
 &\{P_n^0(c),\Psi_\alpha,K_\alpha,P^{\rm gen},
 \mathbf1,z_-,z_+,h_*\}_{m,n};\\
 &\{\widehat P^{\rm red},S,d,\lambda^{\rm or},
 \beta_q,\Theta_{\rm Bock}\}\bigr).
\end{split}
\end{equation}
The macroobject preserves its entry:
\begin{equation}\label{eq:detint-graph-a}
 \mathfrak M_{\rm det}:=\operatorname{Graph}(a_{\rm det}),
 \qquad
 \mathcal A_{\rm det}(z):=(z,a_{\rm det}z).
\end{equation}

The raw publisher
\begin{equation}\label{eq:detint-publicador-crudo}
 p_{\rm det}:\mathfrak M_{\rm det}\longrightarrow\mathscr Y_{\rm det}
\end{equation}
publishes the certificate jointly
\begin{equation}\label{eq:detint-certificado-publicado}
\begin{gathered}
 \partial^2=0,quad \text{AW and counit},quad
 \mathbf1_y\text{ compatible group-like},\\
 \partial z_\pm=0,quad
 \varpi_{\rm root}z_-=\varpi_{\rm root}z_+,quad
 \partial h_*=z_--z_+,\\
 \Psi_\beta\Psi_\alpha=\Psi_{\beta\alpha},quad
 K_{\beta\alpha}=K_\beta+\Psi_\beta K_\alpha,\\
 d=v_3(\det S)=\sum_i a_i,quad
 \Theta_{\rm Bock}=\overline{\lambda^{\rm or}},quad
 \text{full naturality under refinement}.
\end{gathered}
\end{equation}
Finally,
\begin{equation}\label{eq:detint-graph-p}
 \mathcal C_{\rm det}^{\rm HMT}:=\operatorname{Graph}(p_{\rm det}),
 \qquad
 \mathcal P_{\rm det}(M):=(M,p_{\rm det}M).
\end{equation}
The complete chain is
\begin{equation}\label{eq:detint-cadena-completa}
 \mathcal C_{\rm cont}^{\rm disc}\supset
 \mathcal C_{\rm cont}^{\rm det}
 \xrightarrow{\operatorname{Res}_{\rm det}}\mathcal G_{\rm det}
 \xrightarrow{\operatorname{pr}_{X,{\rm det}}}X_{\chi,{\rm det}}
 \xrightarrow{\mathcal A_{\rm det}}\mathfrak M_{\rm det}
 \xrightarrow{\mathcal P_{\rm det}}\mathcal C_{\rm det}^{\rm HMT}.
\end{equation}
Each arrow has an inverse on its image given by a coordinate
projection. History is not replaced by value, multisection by selector,
complex by determinant or proof by name.

\section{Theorem of Internal Determinantal Realization}

\begin{theorem}[Complete internal determinant realization]
\label{thm:detint-realizacion}
For every \(x\in\mathcal C_{\rm cont}^{\rm det}\), the chain
\eqref{eq:detint-cadena-completa} constructs a macro-object natural in
depth and precision. It contains two cycles with a common root, an explicit filler
that identifies them, a compatible group-like unit, and a
\(3\)-adic terminal whose order, Smith form, and Bockstein tower produce the same
oriented reader. The thirteen coordinates of \eqref{eq:detint-d-trece}
remain recoverable in the output.
\end{theorem}

\begin{proof}
Equations \eqref{eq:detint-p0-diferencial}--
\eqref{eq:detint-d2} prove that \(P_n^0(c)\) is a complex. The two
augmentations are maps of complexes, so the pullback is typed.
Alexander--Whitney is natural and makes each vertex group-like; together with
\(r\), this gives the common unit. The direct calculation
\eqref{eq:detint-zpm-ciclos}--\eqref{eq:detint-filler-identidad} proves that
the publications are cycles, have the same root and are joined by
\(h_*\). The formulas for \(\Psi_\alpha\) prove functoriality, and
\(K_\alpha\) proves homotopic naturality while preserving the memory
increment. The definition of \(\mathcal C_{\rm cont}^{\rm det}\) ensures that
terminal reduction produces a square block without postulating it outside the
domain. Smith over \(\mathbf Z_3\) gives \(d=\sum_i a_i\); the calculation of each
block \(3^{a_i}u_i\) gives
\(\Theta_{\rm Bock}=\overline{\lambda^{\rm or}}\). Truncation, reduction and
addition of unit blocks preserve these identities. Finally, the
three graph constructions retain the input as the first coordinate,
so the complete composition is non-ablative.
\end{proof}

\section{Provenance and Forgers}

\paragraph{Origin.}
The TPK nerve, Alexander--Whitney, the counit, the common unit, the normal
form \(\partial f=3ce+b\), \(\partial e=g\),
\(\partial b=-3cg\), the common root, the filler, and the tower architecture
have the status of \texttt{PREEXISTING\_AUTHORIAL\_ARCHITECTURE}. The literal computation
of the two cycles, their root, and the identity
\(\partial h_*=z_--z_+\) has the status of a \texttt{RECOVERED\_RESULT}. The
local system \(\Psi_\alpha/K_\alpha\), the balanced surviving subdomain,
the exclusively \(3\)-adic terminal, and the packaging by
graphs have the status of a \texttt{NEW\_FORMALIZATION}; they formalize without altering
the conceptual provenance of the apparatus.

\begin{ablation}[Internal falsifiers of the determinantal branch]
\label{abl:detint-falsadores}
The construction is falsified at the corresponding level if any of
the following facts occurs.
\begin{enumerate}
 \item \(\partial^2=0\) fails; then the local complex does not exist.
 \item Alexander--Whitney does not commute with refinement, the counit fails, or
 the vertex ceases to be group-like.
 \item \(z_\pm\) are not cycles, their roots differ, or
 \(\partial h_*=z_--z_+\) fails.
 \item \(\Psi_\alpha\) is not a map of complexes or
 \(K_{\beta\alpha}=K_\beta+\Psi_\beta K_\alpha\) fails; in that case
 carry or memory has been lost.
 \item The reduced complex is not balanced or is not acyclic over
 \(K_3\); the history lies outside the typed terminal domain.
 \item The Bockstein lengths do not coincide with the \(a_i\), or
 \(\Theta_{\rm Bock}\ne\overline{\lambda^{\rm or}}\).
 \item Refinement modifies \(d\) or \(\lambda^{\rm or}\) outside
 unit blocks and recorded orientation changes.
 \item One of the thirteen coordinates is omitted: \emph{object replaced;
 conclusion not transferable}.
\end{enumerate}
These forgers check the internally constructed object; they do not replace it with a subsequent control.
\end{ablation}
\fi
 \endgroup

\section{Cellular incidence and correlative filling of the same trajectory}

The sequence \eqref{eq:pdfv2-cor-sucesion-celda} assigns to each prefix the
two-term complex
\begin{equation}
\label{eq:pdfv2-cor-cono-celda}
 C_{k,N}(\widehat x):=
 \left[A_N^{I_k}\oplus A_N^{J_k}
 \xrightarrow{\kappa_{k,N}}A_N^{I_k\times J_k}\right].
\end{equation}
A refinement acts on ports by truncating the extension and on coefficients
by \(A_{N'}\twoheadrightarrow A_N\). When \(\alpha_u:I_{k'}\to I_k\) and \(\beta_u:J_{k'}\to J_k\) are those maps, their
action on generators is
\begin{equation}
\label{eq:pdfv2-cor-refinamiento-celda}
 (p,q)\longmapsto(p\circ\alpha_u,q\circ\beta_u),
 \qquad
 w\longmapsto w\circ(\alpha_u\times\beta_u),
\end{equation}
and satisfies
\begin{equation}
\label{eq:pdfv2-cor-refinamiento-cadena}
 \kappa_{k',N}\Gamma(u)_1=\Gamma(u)_0\kappa_{k,N}.
\end{equation}

The incidence is not fabricated by labeling four copies of a neutral emission.
It is generated by the free cellular completion of the module of the two sheets that
the same state already bears. For
\[
 V_{\widehat x}:=\mathbf F_3e_{\Sigma,\widehat x}
 \oplus\mathbf F_3e_{\Pi,\widehat x},
\]
the four internal lines are
\begin{equation}
\label{eq:pdfv2-cor-cuatro-rectas-tpk}
 \mathscr I_{\widehat x}:=
 \{L_0,L_\infty,L_+,L_-\}
 =\mathbf P^1(V_{\widehat x}),\qquad
 \begin{aligned}
 L_0&=\langle e_\Sigma\rangle,&
 L_\infty&=\langle e_\Pi\rangle,\\
 L_+&=\langle e_\Sigma+e_\Pi\rangle,&
 L_-&=\langle e_\Sigma-e_\Pi\rangle.
 \end{aligned}
\end{equation}
If \(v_L\) is, respectively,
\((1,0)^t,(0,1)^t,(1,1)^t,(1,-1)^t\), the action of the direction \(L\)
on the generators of sheet is the proyector
\begin{equation}
\label{eq:pdfv2-cor-proyectores-cuatro-direcciones}
 P_L:=v_L(v_L^tv_L)^{-1}v_L^t,
\end{equation}
that is, over \(\mathbf F_3\),
\[
 P_0=\begin{pmatrix}1&0\\0&0\end{pmatrix},\quad
 P_\infty=\begin{pmatrix}0&0\\0&1\end{pmatrix},\quad
 P_+=\begin{pmatrix}2&2\\2&2\end{pmatrix},\quad
 P_-=\begin{pmatrix}2&1\\1&2\end{pmatrix}.
\]
Direct calculation gives
\begin{equation}
\label{eq:pdfv2-cor-proyectores-no-colapso}
 P_L^2=P_L,\qquad \operatorname{rk}P_L=1,\qquad
 \operatorname{Im}P_L=L,\qquad
 P_L=P_{L'}\Longrightarrow L=L'.
\end{equation}
Therefore, the four directions are distinct outputs of the module
\(V_{\widehat x}\); they are not four names placed upon the same action.

\begin{definition}[Cell-free leaf TPK completion]
\label{def:pdfv2-cor-completacion-celular-tpk}
Let \(\operatorname{Cell}_{\rm TPK}(\widehat x)\) be the free complex
\[
 C_2(\widehat x)\xrightarrow{\partial_2}
 C_1(\widehat x)\xrightarrow{\partial_1}C_0(\widehat x)
\]
with
\begin{align*}
 C_0(\widehat x)
 &:=
 \mathbf F_3[v_\star]\oplus
 \bigoplus_{L\in\mathscr I}\mathbf F_3[v_L]\oplus
 \bigoplus_{(L,\epsilon)\in\mathscr B}
       \mathbf F_3[v_{L,\epsilon}]\oplus
 \bigoplus_{E\in\mathscr E}\mathbf F_3[v_E],\\
 C_1(\widehat x)
 &:=
 \bigoplus_{L\in\mathscr I}\mathbf F_3[u_L]\oplus
 \bigoplus_{(L,\epsilon)\in\mathscr B}
       \mathbf F_3[u_{L,\epsilon}]\oplus
 \bigoplus_{E=\{L,L'\}\in\mathscr E}
 \bigl(\mathbf F_3[u_{L'\mid L}]
       \oplus\mathbf F_3[u_{L\mid L'}]\bigr),\\
 C_2(\widehat x)
 &:=
 \bigoplus_{E\in\mathscr E}\mathbf F_3[\omega_E].
\end{align*}
Its action on generators is
\begin{align}
 \partial u_L&=v_L-v_\star,
 &
 \partial u_{L'\mid L}&=v_E-v_L,
 \label{eq:pdfv2-cor-cell-frontera-aristas}\\
 \partial u_{L,\epsilon}&=v_{L,\epsilon}-v_L,
 &
 \partial\omega_{\{L,L'\}}
 &=u_L+u_{L'\mid L}-u_{L'}-u_{L\mid L'}.
 \label{eq:pdfv2-cor-cell-frontera-celda}
\end{align}
\end{definition}

The eight generators \(v_{L,\epsilon}\) materially realize the screen
\(\mathscr B\). The preceding involution extends to the complex by
\begin{equation}
\label{eq:pdfv2-cor-pantalla-ocho-generadores}
 \vartheta_Lv_{K,\epsilon}:=v_{\vartheta_L(K,\epsilon)},\qquad
 \vartheta_Lu_{K,\epsilon}:=u_{\vartheta_L(K,\epsilon)},
\end{equation}
and fixes the other generators. Since the basis is free,
\begin{equation}
\label{eq:pdfv2-cor-pantalla-ocho-no-colapso}
 |\{v_{L,\epsilon}:(L,\epsilon)\in\mathscr B\}|=8,\qquad
 \vartheta_L^2=I,\qquad
 \vartheta_L\vartheta_{L'}=\vartheta_{L'}\vartheta_L.
\end{equation}
Moreover, \(\partial\vartheta_L=\vartheta_L\partial\), because both paths
send \(u_{L,\epsilon}\) to \(v_{L,-\epsilon}-v_L\) and fix the remaining
boundaries.

The two words
\[
 p_{L,L'}:=u_{L'\mid L}u_L,\qquad
 p_{L',L}:=u_{L\mid L'}u_{L'}
\]
are different paths in the free complex. Their difference is the boundary of
\(\omega_{\{L,L'\}}\), and
\begin{equation}
\label{eq:pdfv2-cor-cell-dos-cero}
 \partial^2\omega_E
 =(v_L-v_\star)+(v_E-v_L)
 -(v_{L'}-v_\star)-(v_E-v_{L'})=0.
\end{equation}
The representation of TPK on the leaves is fixed by
\[
 \varrho(u_L)=P_L,\qquad
 \varrho(u_{L'\mid L})=P_{L'},\qquad
 \varrho(p_{L,L'})=P_{L'}P_L.
\]
Even if two published matrices coincided in a specialization, the two
paths do not collapse: they remain distinct free generators of
\(C_1(\widehat x)\), joined by an explicit 2--cell.

The complete incidence coordinate is the group of 2--cochains.
\begin{equation}
\label{eq:pdfv2-cor-q-incidencia-total}
 Q_{\rm inc}(\widehat x):=
 C^2(\operatorname{Cell}_{\rm TPK}(\widehat x);\mathbf F_3)
 =\bigoplus_{E\in\mathscr E}\mathbf F_3\omega_E^*
 \simeq\mathbf F_3^6.
\end{equation}
The internal operation returns all of \(Q_{\rm inc}(\widehat x)\), not a
selected element. Each \(q\) has a unique expansion
\(q=\sum_Eq_E\omega_E^*\). For
\(a\in\mathbf F_3^{\mathscr I}\), its change of potential is
\begin{equation}
\label{eq:pdfv2-cor-collar-variacion}
 q\longmapsto q+Ba,\qquad (Ba)_{\{L,L'\}}=a_L+a_{L'}.
\end{equation}
As each line belongs to three pairs,
\[
 s(Ba)=3\sum_{L\in\mathscr I}a_L=0,
\]
and therefore
\begin{equation}
\label{eq:pdfv2-cor-wc-invariante-celular}
 W_c(q+Ba)=W_c(q)\qquad(q\in Q_{\rm inc}(\widehat x)).
\end{equation}
The unit of the completion only anchors the complex in the state:
\begin{equation}
\label{eq:pdfv2-cor-cell-unidad}
 \eta^{\rm cell}_{\widehat x}:V_{\widehat x}\longrightarrow
 V_{\widehat x}\otimes C_0(\operatorname{Cell}_{\rm TPK}(\widehat x)),\qquad
 e_{\diamond,\widehat x}\longmapsto
 e_{\diamond,\widehat x}\otimes v_\star,
\end{equation}
and marks the origin \(q=0\); it does not assert that the orbit of that origin is all of
\(Q_{\rm inc}\).

The construction is functorial under common truncation. If
\(\tau_{\ell,k}\widehat x_\ell=\widehat x_k\), we define
\begin{align*}
 \tau_V&:V_{\widehat x_\ell}\longrightarrow V_{\widehat x_k},&
 \tau_Ve_{\diamond,\widehat x_\ell}&=e_{\diamond,\widehat x_k},\\
 \tau_a&:\mathbf F_3^{\mathscr I_{\widehat x_\ell}}
          \longrightarrow\mathbf F_3^{\mathscr I_{\widehat x_k}},&
 (\tau_aa)_{\tau_{\mathscr I}L}&=a_L,\\
 \tau_Q&:Q_{\rm inc}(\widehat x_\ell)
          \longrightarrow Q_{\rm inc}(\widehat x_k),&
 (\tau_Qq)_{\tau_{\mathscr E}E}&=q_E,
\end{align*}
where \(\tau_{\mathscr I}[v]=[\tau_Vv]\) and
\(\tau_{\mathscr E}\{L,L'\}=\{\tau_{\mathscr I}L,
\tau_{\mathscr I}L'\}\). The map of cadenas
\(\operatorname{Cell}(\tau):C_\bullet(\widehat x_\ell)
\to C_\bullet(\widehat x_k)\) acts in generators by
\begin{equation}
\label{eq:pdfv2-cor-cell-truncamiento-generadores}
 v_\bullet(\widehat x_\ell)\mapsto v_\bullet(\widehat x_k),\quad
 u_\bullet(\widehat x_\ell)\mapsto u_\bullet(\widehat x_k),\quad
 \omega_E(\widehat x_\ell)\mapsto\omega_E(\widehat x_k).
\end{equation}
From \eqref{eq:pdfv2-cor-cell-frontera-aristas}--
\eqref{eq:pdfv2-cor-cell-frontera-celda} one obtains
\begin{equation}
\label{eq:pdfv2-cor-collar-naturalidad}
\begin{aligned}
 \operatorname{Cell}(\tau)\partial
 &=\partial\operatorname{Cell}(\tau),\\
 \tau_VP_L&=P_{\tau_{\mathscr I}L}\tau_V,\qquad
 \tau_QB=B\tau_a,\qquad W_c\circ\tau_Q=W_c,\\
 \operatorname{Cell}(\tau)\vartheta_L
 &=\vartheta_{\tau_{\mathscr I}L}\operatorname{Cell}(\tau),\\
 (\tau_V\otimes\operatorname{Cell}_0(\tau))
 \eta^{\rm cell}_{\widehat x_\ell}
 &=\eta^{\rm cell}_{\widehat x_k}\tau_V.
\end{aligned}
\end{equation}
Identity and composition are verified on each generator. Thus
\(\operatorname{Cell}_{\rm TPK}\) is the functorial cellular completion of
the same APP--TRIT--TPK fiber, not a carrier added from a subsequent
publication.

\subsection{Incidence measure, carry curvature, and adjoint naturality}
In the following development, they are identified exclusively as abbreviations
of the common completion,
\[
 I_{\rm gau}:=\mathscr I_{\widehat x},\qquad
 \mathscr B_{\rm gau}:=\mathscr B_{\widehat x},\qquad
 Q_{\rm gau}:=Q_{\rm inc}(\widehat x).
\]
With the uniform counting average on \(Q_{\rm gau}\), the function
\(W_c\) of \eqref{eq:pdfv2-cor-wc} satisfies
\begin{equation}
\label{eq:pdfv2-cor-wc-momentos}
 \mathbb E_QW_c=0,\qquad
 \mathbb E_QW_c^2=\frac29,\qquad
 \mathbb E_QW_c^3=\frac{2}{27}.
\end{equation}
These identities are properties of the same cell and the same carry, not
data from an external theory.
\begingroup
  \let\section\subsection
  % Extracción quirúrgica del fuente teoremática V2 05c: líneas 99--335.
% El resto permanece en este archivo como procedencia, pero no se publica.
\iffalse
\chapter[Internal Gauge Macrostructure]{Internal Gauge Macrostructure Generated by the Discrete Continuum}
\label{chap:gauint-macroestructura}

\begin{thesisbox}
The macrostructure of this chapter is constructed solely from a
surviving restriction of \(\mathcal C_{\rm cont}^{\rm disc}\). The numbers
\(4\), \(6\), and \(8\), the screen, the four involutions, the
projective measure, the transfer operator, the incidence mode, and the cellular
terminal are derived within that restriction. No subsequent evaluation
forms part of the domain or of the arrows that follow.
\end{thesisbox}

\section{The projective seed and the internal derivation of \texorpdfstring{\(4/6/8\)}{4/6/8}}

Let \(V_{\rm TPK}=\mathbf F_3^2\), with the ordered basis inherited from the two
sheets of the TPK state. We define the set of directions
\begin{equation}\label{eq:gauint-cuatro-direcciones}
 I_{\rm gau}:=\mathbf P^1(\mathbf F_3)
 =\{L_\infty,L_0,L_1,L_{-1}\},
\end{equation}
where
\[
 L_\infty=[1:0],\qquad L_0=[0:1],\qquad
 L_1=[1:1],\qquad L_{-1}=[1:-1].
\]
Therefore,
\begin{equation}\label{eq:gauint-cardinal-cuatro}
 |I_{\rm gau}|=\frac{3^2-1}{3-1}=4.
\end{equation}
A set of four indices is not introduced: it appears as the quotient of
the eight nonzero vectors by the action of \(\mathbf F_3^\times\).

The unoriented incidences are
\begin{equation}\label{eq:gauint-seis-incidencias}
 E_{\rm gau}:=\binom{I_{\rm gau}}2,
 \qquad |E_{\rm gau}|=\binom42=6,
 \qquad Q_{\rm gau}:=\mathbf F_3^{E_{\rm gau}}.
\end{equation}
The oriented screen duplicates each direction, not each incidence:
\begin{equation}\label{eq:gauint-ocho-caras}
 \mathscr B_{\rm gau}:=I_{\rm gau}\times\{+,-\},
 \qquad |\mathscr B_{\rm gau}|=2|I_{\rm gau}|=8.
\end{equation}
Thus, \(4\), \(6\), and \(8\) arise, respectively, from directions, pairs of
directions, and oriented faces of the same seed.

For each \(L\in I_{\rm gau}\), let \(\vartheta_L\) be the screen
involution that exchanges \((L,+)\) and \((L,-)\) and fixes the other six
faces. Then
\begin{equation}\label{eq:gauint-cuatro-involuciones}
 \vartheta_L^2=I,
 \qquad
 \vartheta_L\vartheta_{L'}=\vartheta_{L'}\vartheta_L
 \quad(L\ne L').
\end{equation}
The four involutions have been obtained prior to defining measure or operator.

\section{Incidence, internal mode and non-constant witness}

The incidence application is
\begin{equation}\label{eq:gauint-incidencia-b}
 B:\mathbf F_3^{I_{\rm gau}}\longrightarrow Q_{\rm gau},
 \qquad
 (Ba)_{\{L,L'\}}:=a_L+a_{L'}.
\end{equation}
If
\[
 s(q):=\sum_{e\in E_{\rm gau}}q_e\in\mathbf F_3,
\]
each direction appears exactly three times in the sum of the six
incidences, and therefore
\begin{equation}\label{eq:gauint-incidencia-invariante}
 s(Ba)=3\sum_{L\in I_{\rm gau}}a_L=0,
 \qquad
 s(q+Ba)=s(q).
\end{equation}

On \(Q_{\rm gau}\), only the normalized counting average is used.
The incidence mode is the rational function
\begin{equation}\label{eq:gauint-wc-definicion}
 W_c(q):=\mathbf 1_{\{s(q)=0\}}-\frac13.
\end{equation}
Since \(s:Q_{\rm gau}\to\mathbf F_3\) is surjective, each fiber has \(3^5\) elements.
Exact counting then yields,
\begin{equation}\label{eq:gauint-wc-momentos}
 \mathbb E_{Q}W_c=0,
 \qquad
 \mathbb E_{Q}W_c^2=\frac29,
 \qquad
 \mathbb E_{Q}W_c^3=\frac{2}{27}.
\end{equation}
Furthermore, \eqref{eq:gauint-incidencia-invariante} gives
\begin{equation}\label{eq:gauint-wc-invariancia}
 W_c(q+Ba)=W_c(q).
\end{equation}
Therefore, \(W_c\) is nonzero, centered, and invariant under the internally
generated incidence.

\fi
\section{Incidence, carry curvature, and common domain of naturalness}

For \(x\in\mathcal C_{\rm cont}^{\rm disc}\), let us denote by
\(Y_m(x)\) the finite set of gauge prefixes of depth \(m\) that
simultaneously conserve
\begin{equation*}
\begin{gathered}
 \text{fiber, extension relations, numbers, orientation, carry, memory,}\\
 \text{boundary, route, multisection, survival, terminal and reader}.
\end{gathered}
\end{equation*}
There exists a collar neighborhood application.
\begin{equation}\label{eq:gauint-collar}
 q_m:Y_m(x)\longrightarrow Q_{\rm gau}
\end{equation}
which records the six incidences without erasing the history that produces them.
For \(\ell\ge m\), the truncation is written
\[
 \tau_{\ell,m}:Y_\ell(x)\longrightarrow Y_m(x).
\]

For \(y,z\in Y_m(x)\), let
\begin{equation}\label{eq:gauint-todas-extensiones}
 \operatorname{Ext}^{\rm adm}_{m}(y,z)
\end{equation}
the set of \emph{all} admissible TPK extensions from \(y\) to \(z\).
An extension belongs to this set only if it transports the entire fiber,
the relations \(\operatorname{Ext}_\Gamma\), the native and operatorial carry,
the orientation, the accumulated memory, the boundary, the route, and the
survival condition. We define
\begin{equation}\label{eq:gauint-conteo-kernel}
 a_m(y,z):=\#\operatorname{Ext}^{\rm adm}_m(y,z),
 \qquad
 d_m(y):=\sum_{z\in Y_m(x)}a_m(y,z),
\end{equation}
and, when \(d_m(y)>0\),
\begin{equation}\label{eq:gauint-kernel-normalizado}
 k_m(y,z):=\frac{a_m(y,z)}{d_m(y)}.
\end{equation}
Thus, \(k_m\) is a normalized count of complete extensions, not a
function introduced after projecting the state.

The restriction \(\mathcal C_{\rm cont}^{\rm gau}\) is formed by the
surviving histories \(x\) for which the following conditions
hold at every depth.
\begin{enumerate}
 \item The sets \(Y_m(x)\) are nonempty, \(d_m(y)>0\), and the
 truncations \(\tau_{\ell,m}\) are surjective.
 \item There exists a nonempty multisection of strictly positive rational
 weights \(\mu_m:Y_m(x)\to\mathbf Q_{>0}\) such that
 \begin{equation}\label{eq:gauint-medida-proyectiva}
  \sum_{y\in Y_m}\mu_m(y)=1,
  \qquad
  \mu_m(y)=\sum_{\tau_{\ell,m}\widetilde y=y}\mu_\ell(\widetilde y).
 \end{equation}
 \item Each \(\mu_m\) is stationary:
 \begin{equation}\label{eq:gauint-estacionariedad}
  \sum_{y\in Y_m}\mu_m(y)k_m(y,z)=\mu_m(z).
 \end{equation}
 \item The marginal of \((q_m)_*\mu_m\) is the uniform counting measure on
 \(Q_{\rm gau}\). This condition makes the calculations
 of \eqref{eq:pdfv2-cor-wc-momentos} literally applicable to the lift \(W_c\circ q_m\).
 \item Direct lumpability holds
 \begin{equation}\label{eq:gauint-lump-directa}
  \sum_{\tau_{\ell,m}\widetilde z=z}
   k_\ell(\widetilde y,\widetilde z)
  =k_m(\tau_{\ell,m}\widetilde y,z)
 \end{equation}
 for every \(\widetilde y\in Y_\ell\) and \(z\in Y_m\).
 \item For the adjoint kernel
 \begin{equation}\label{eq:gauint-kernel-adjunto}
  k_m^*(z,y):=\frac{\mu_m(y)k_m(y,z)}{\mu_m(z)},
 \end{equation}
 the adjoint lumpability condition is satisfied
 \begin{equation}\label{eq:gauint-lump-adjunta}
  \sum_{\tau_{\ell,m}\widetilde y=y}
   k_\ell^*(\widetilde z,\widetilde y)
  =k_m^*(\tau_{\ell,m}\widetilde z,y).
 \end{equation}
 \item The categories of prefixes possess the TPK root as an initial object,
 so that their nerve has the explicit terminal contraction of
 \eqref{eq:gauint-sdr} below.
\end{enumerate}
These conditions define the typed domain; outside it, a measure or natural
operator is not published by mere declaration.

\section{Projective measure, kernel, and direct and adjoint naturality}

Having fixed a branch of the multisection \(\{\mu_m\}_m\), let
\begin{equation}\label{eq:gauint-hm}
 H_m:=\operatorname{Fun}(Y_m,\mathbf Q),
 \qquad
 \langle f,g\rangle_m:=\sum_{y\in Y_m}\mu_m(y)f(y)g(y).
\end{equation}
The operator generated by all extensions is
\begin{equation}\label{eq:gauint-km}
 (K_mf)(y):=\sum_{z\in Y_m}k_m(y,z)f(z).
\end{equation}
The Markov equations and stationarity give
\begin{equation}\label{eq:gauint-markov-estacionario}
 K_m\mathbf1=\mathbf1,
 \qquad
 K_m^*\mathbf1=\mathbf1,
 \qquad
 \|K_mf\|_m\le\|f\|_m.
\end{equation}

The lift and the truncated expectation are
\begin{equation}\label{eq:gauint-lift-esperanza}
 J_{\ell,m}f:=f\circ\tau_{\ell,m},
 \qquad
 E_{m,\ell}:=J_{\ell,m}^*.
\end{equation}
The projectivity of \(\mu\) implies
\begin{equation}\label{eq:gauint-ej-identidad}
 E_{m,\ell}J_{\ell,m}=I_{H_m}.
\end{equation}
The two lumpability conditions discharge, respectively,
\begin{equation}\label{eq:gauint-naturalidad-directa-adjunta}
 K_\ell J_{\ell,m}=J_{\ell,m}K_m,
 \qquad
 K_\ell^*J_{\ell,m}=J_{\ell,m}K_m^*.
\end{equation}
The second equality is not inferred from the first without the attached condition:
both remain registered in the domain.

The positive internal operator is defined
\begin{equation}\label{eq:gauint-transferencia-t}
 T_m:=K_m^*K_m.
\end{equation}
Then
\begin{equation}\label{eq:gauint-t-positivo-natural}
 \langle f,T_mf\rangle_m=\|K_mf\|_m^2\ge0,
 \qquad
 T_\ell J_{\ell,m}=J_{\ell,m}T_m.
\end{equation}

\section{Common law of \(8\) faces and \(4\) Fubini squares}

On \(Y_m^{\mathscr B_{\rm gau}}\), we define the rational mass
\begin{equation}\label{eq:gauint-ley-ocho-caras}
 \Omega_m^{(8)}(\mathbf y)
 :=\sum_{z\in Y_m}\mu_m(z)
 \prod_{(L,\varepsilon)\in\mathscr B_{\rm gau}}
 k_m(z,y_{L,\varepsilon}).
\end{equation}
It is a probability because each factor is Markov. All faces arise
from the same latent state \(z\) and the same kernel; they are not eight independently
added measures.

For functions \(f_{L,+},f_{L,-}\in H_m\), finite Fubini gives
\begin{equation}\label{eq:gauint-fubini-ocho}
\begin{aligned}
 &\sum_{\mathbf y}\Omega_m^{(8)}(\mathbf y)
  \prod_{L\in I_{\rm gau}}
  f_{L,+}(y_{L,+})f_{L,-}(y_{L,-})\\
 &\qquad=
 \sum_{z\in Y_m}\mu_m(z)
 \prod_{L\in I_{\rm gau}}
 (K_mf_{L,+})(z)(K_mf_{L,-})(z).
\end{aligned}
\end{equation}
Upon taking \(f_{L,+}=f_{L,-}=f_L\), the four squares appear
simultaneously
\begin{equation}\label{eq:gauint-cuatro-cuadrados}
 \sum_{z\in Y_m}\mu_m(z)
 \prod_{L\in I_{\rm gau}}(K_mf_L(z))^2.
\end{equation}
The marginal of any opposite pair satisfies
\begin{equation}\label{eq:gauint-marginal-t}
 \sum_{y_+,y_-}\Omega_{m,L}^{(2)}(y_+,y_-)
 f(y_+)g(y_-)
 =\langle K_mf,K_mg\rangle_m
 =\langle f,T_mg\rangle_m.
\end{equation}
The commutativity of the product in \eqref{eq:gauint-ley-ocho-caras} proves
\begin{equation}\label{eq:gauint-involuciones-ley}
 (\vartheta_L)_*\Omega_m^{(8)}=\Omega_m^{(8)}
 \qquad(L\in I_{\rm gau}).
\end{equation}

\section{Exact internal operator and self-test witness}

The projection onto the constant sector and its complement are
\begin{equation}\label{eq:gauint-pq}
 P_{\Omega,m}f:=\langle\mathbf1,f\rangle_m\mathbf1,
 \qquad
 Q_m:=I-P_{\Omega,m}.
\end{equation}
The internal cell operator is defined by
\begin{equation}\label{eq:gauint-d-hmt}
 D_m^{\rm HMT}:=3Q_m.
\end{equation}
The coefficient \(3\) is the cardinality of the residue field used in
\eqref{eq:pdfv2-cor-proyectores-cuatro-direcciones}; it does not represent an
added scale.
For the lift
\[
 W_{c,m}:=W_c\circ q_m
\]
one has, by \eqref{eq:pdfv2-cor-wc-momentos},
\begin{equation}\label{eq:gauint-w-eigen}
 P_{\Omega,m}W_{c,m}=0,
 \qquad
 D_m^{\rm HMT}W_{c,m}=3W_{c,m},
 \qquad
 \|W_{c,m}\|_m^2=\frac29.
\end{equation}
Thus, the complementary sector possesses a nonzero witness constructed by the
incidence of the six edges.

\section{Cellular terminal and strong retractor}

Let \(C_*(N\mathsf Y_m;A_n)\) be the normalized chain complex of the nerve of the
prefix category, and let \(y_*\) be its initial root. The inclusion and the
augmentation are
\[
 \iota_m:A_n[0]\longrightarrow C_*(N\mathsf Y_m;A_n),
 \quad \iota_m(1)=[y_*],
 \qquad
 \epsilon_m:C_*(N\mathsf Y_m;A_n)\longrightarrow A_n[0].
\]
The initial object provides the extra degeneracy
\begin{equation}\label{eq:gauint-homotopia-celular}
 H_m[y_0\to\cdots\to y_q]
 :=[y_*\to y_0\to\cdots\to y_q],
\end{equation}
with the normalized convention when a degeneration appears. It is verified
\begin{equation}\label{eq:gauint-sdr}
 \epsilon_m\iota_m=I,
 \qquad
 \partial H_m+H_m\partial=I-\iota_m\epsilon_m,
 \qquad
 H_m^2=H_m\iota_m=\epsilon_mH_m=0.
\end{equation}
This cellular terminal is distinct from the constant projection
\(P_{\Omega,m}\): the former contracts the nerve, while the latter separates the
sectors of the function space. Both are preserved.

\iffalse
\section{The thirteen coordinates of the gauge constraint}

For each \((m,n)\), it is defined
\begin{equation}\label{eq:gauint-d-trece}
 D_{{\rm gau},m,n}(x):=
 (\mathcal F,\operatorname{Ext}_\Gamma,\operatorname{Rel},
 \mathcal N,\operatorname{or},\operatorname{Carr},\operatorname{Mem},
 \partial,\gamma,\operatorname{Msect},\operatorname{Surv},
 \operatorname{Term},\mathcal L)_{{\rm gau},m,n}(x),
\end{equation}
with the following content.
\begin{enumerate}
 \item \(\mathcal F=(V_{\rm TPK},I_{\rm gau},E_{\rm gau},
 Q_{\rm gau},\mathscr B_{\rm gau},Y_m,H_m)\).
 \item \(\operatorname{Ext}_\Gamma\) contains all the sets
 \(\operatorname{Ext}^{\rm adm}_m(y,z)\), truncations, lifts,
 expectations, and their compositions.
 \item \(\operatorname{Rel}\) contains \(B\), \(sB=0\), the four
 involutions, Markov, stationarity, and both lumpabilities.
 \item \(\mathcal N=(3,4,6,8,m,n,|Y_m|,a_m,d_m)\); those numbers do not
 replace any of the remaining coordinates.
 \item \(\operatorname{or}\) preserves the faces \(+/-\), the four
 inversions \(\vartheta_L\), and the reversal of each path.
 \item \(\operatorname{Carr}\) contains the complete carry of every extension counted in \(a_m(y,z)\), including the carry already propagated.
 \item \(\operatorname{Mem}\) conserves the full prefix, the common
 state \(z\) of \(\Omega_m^{(8)}\), and the memory of its eight extensions.
 \item \(\partial=(q_m,\partial_{N\mathsf Y_m},Q_m,H_m)\) preserves
 collar boundary, cellular boundary, and complement.
 \item \(\gamma\) is the oriented simplex of the nerve that records the route
 and not only its endpoints.
 \item \(\operatorname{Msect}\) is the family of all compatible
 projective stationary measures and of all the extensions that
 enter into the counts; no retrospective branch is selected.
 \item \(\operatorname{Surv}\) records nonemptiness, arbitrary
 depth, projectivity, both lumpabilities, and the persistence of the
 uniform collar marginal.
 \item \(\operatorname{Term}=(\iota_m,\epsilon_m,H_m;
 P_{\Omega,m},Q_m)\) keeps the two terminals separate.
 \item \(\mathcal L=(\Omega_m^{(8)},K_m,T_m,D_m^{\rm HMT},W_{c,m})\)
 is the internal reader posterior to the construction of the history.
\end{enumerate}

For \(\ell\ge m\) and \(n'\ge n\), all the coordinates satisfy
\begin{equation}\label{eq:gauint-naturalidad-trece}
 u^{\rm gau}_{(\ell,n'),(m,n)}D_{{\rm gau},\ell,n'}(x')
 =D_{{\rm gau},m,n}(\tau_{\ell,m}x').
\end{equation}

\section{Assembler, publisher, and non-ablative chain}

Let
\begin{equation}\label{eq:gauint-graph-d}
 \mathcal G_{\rm gau}:=\operatorname{Graph}(D_{\rm gau}),
 \qquad
 \operatorname{Res}_{\rm gau}(x):=(x,D_{\rm gau}x).
\end{equation}
The dependent multisection \(\chi_{\rm gau}(x)\) contains all compatible branches
of collar, measure, and extension. We define
\begin{equation}\label{eq:gauint-x-dependiente}
 X_{\chi,{\rm gau}}
 :=\{(\chi_{\rm gau}(x),x,D_{\rm gau}x):
 x\in\mathcal C_{\rm cont}^{\rm gau}\},
\end{equation}
\begin{equation}\label{eq:gauint-prx}
 \operatorname{pr}_{X,{\rm gau}}(x,D_{\rm gau}x)
 :=(\chi_{\rm gau}(x),x,D_{\rm gau}x).
\end{equation}

The raw assembler
\begin{equation}\label{eq:gauint-ensamblador-crudo}
 a_{\rm gau}:X_{\chi,{\rm gau}}\longrightarrow
 \mathscr M_{\rm gau}^{\rm amb}
\end{equation}
produce the complete family
\[
 \bigl(I,E,\mathscr B,B,W_c;
 Y_m,\mu_m,k_m,J_{\ell,m},E_{m,\ell};
 K_m,T_m,\Omega_m^{(8)},P_{\Omega,m},Q_m,D_m^{\rm HMT};
 \iota_m,\epsilon_m,H_m\bigr)_{m,n}.
\]
The macroobject preserves its entry through
\begin{equation}\label{eq:gauint-graph-a}
 \mathfrak M_{\rm gau}:=\operatorname{Graph}(a_{\rm gau}),
 \qquad
 \mathcal A_{\rm gau}(z):=(z,a_{\rm gau}z).
\end{equation}

The raw publisher
\begin{equation}\label{eq:gauint-publicador-crudo}
 p_{\rm gau}:\mathfrak M_{\rm gau}\longrightarrow
 \mathscr Y_{\rm gau}
\end{equation}
publishes jointly the identities
\[
 \begin{gathered}
 |I|=4,quad |E|=6,quad |\mathscr B|=8,quad sB=0,
 \quad W_c(q+Ba)=W_c(q),\\
 K_\ell J=JK_m,quad K_\ell^*J=JK_m^*,
 \quad T_m=K_m^*K_m,\\
 (\vartheta_L)_*\Omega_m^{(8)}=\Omega_m^{(8)},
 \quad \text{the four quadratic Fubini factorizations},\\
 D_m^{\rm HMT}=3Q_m,quad D_m^{\rm HMT}W_{c,m}=3W_{c,m},
 \quad \partial H_m+H_m\partial=I-\iota_m\epsilon_m.
 \end{gathered}
\]
Finally,
\begin{equation}\label{eq:gauint-graph-p}
 \mathcal C_{\rm gau}^{\rm HMT}:=\operatorname{Graph}(p_{\rm gau}),
 \qquad
 \mathcal P_{\rm gau}(M):=(M,p_{\rm gau}M).
\end{equation}
The complete chain is
\begin{equation}\label{eq:gauint-cadena-completa}
 \mathcal C_{\rm cont}^{\rm disc}\supset
 \mathcal C_{\rm cont}^{\rm gau}
 \xrightarrow{\operatorname{Res}_{\rm gau}}\mathcal G_{\rm gau}
 \xrightarrow{\operatorname{pr}_{X,{\rm gau}}}X_{\chi,{\rm gau}}
 \xrightarrow{\mathcal A_{\rm gau}}\mathfrak M_{\rm gau}
 \xrightarrow{\mathcal P_{\rm gau}}\mathcal C_{\rm gau}^{\rm HMT}.
\end{equation}
Each arrow has an inverse on its image given by a projection of
coordinates. No output replaces its generative history.

\section{Theorem of internal gauge realization}

\begin{theorem}[Complete and natural internal realization]
\label{thm:gauint-realizacion}
For every \(x\in\mathcal C_{\rm cont}^{\rm gau}\), the chain \eqref{eq:gauint-cadena-completa} constructs, from the discrete structure of the continuum, a natural macrostructure with four directions, six incidences, eight faces, four involutions, a projective measure, a complete counting kernel, its positive transfer, a common eight-face law, a nonconstant incidence witness, and a cellular terminal. The thirteen coordinates of \eqref{eq:gauint-d-trece} remain recoverable in the output.
\end{theorem}

\begin{proof}
Equations \eqref{eq:gauint-cuatro-direcciones}--\eqref{eq:gauint-ocho-caras} derive \(4/6/8\). The incidence sum \eqref{eq:gauint-incidencia-invariante} proves the invariance of \(W_c\), and the uniform count proves \eqref{eq:gauint-wc-momentos}. The complete extensions produce \(k_m\) by \eqref{eq:gauint-kernel-normalizado}; the explicit domain conditions prove the Markov property, stationarity, and the two naturality equalities. Hence \(T_m=K_m^*K_m\) is positive and natural. Applying Fubini once to \eqref{eq:gauint-ley-ocho-caras} yields the four quadratic forms without changing the measure. The uniform marginal centers \(W_{c,m}\), whence \eqref{eq:gauint-w-eigen}. The initial object of the nerve gives the additional degeneracy and the SDR identity \eqref{eq:gauint-sdr}. Finally, the three graph constructions retain their inputs as first coordinates, so the composition forgets none of the thirteen inputs.
\end{proof}

\section{Provenance and Forgers}

\paragraph{Origin.}
The TPK seed, the oriented screen, the surviving extensions, the common measure, and the terminal contraction have the status \texttt{PREEXISTING\_AUTHORIAL\_ARCHITECTURE}. The counting transfer, the incidence mode, and the exact sector \(3Q\) are presented as a \texttt{RECOVERED\_RESULT}. The explicit derivation \(\mathbf P^1(\mathbf F_3)\to4/6/8\), the two lumpability conditions, and the non-ablative packaging by graphs have the status of a \texttt{NEW\_FORMALIZATION}; they formalize the preceding architecture without claiming a new conceptual origin.

\begin{ablation}[Internal falsifiers of the gauge branch]
\label{abl:gauint-falsadores}
The construction is falsified at the indicated level if any of the
following facts occurs.
\begin{enumerate}
 \item \(I_{\rm gau}\), \(E_{\rm gau}\), or
 \(\mathscr B_{\rm gau}\) is replaced by a list without the maps that derive them.
 \item \(sB=0\) fails or \(W_c\) ceases to be invariant and nonconstant.
 \item The count \(a_m(y,z)\) omits carry, orientation, memory, boundary, or
 route: in that case a different object has been counted.
 \item Projectivity, stationarity, or one of the two
 lumpabilities fails; then the corresponding naturality cannot be published.
 \item The eight faces do not follow from a unique \(\mu_m\) and a unique
 \(k_m\), or some \(\vartheta_L\) does not preserve \(\Omega_m^{(8)}\).
 \item \(P_{\Omega,m}W_{c,m}\ne0\) o
 \(D_m^{\rm HMT}W_{c,m}\ne3W_{c,m}\).
 \item The SDR identity \(\partial H+H\partial=I-\iota\epsilon\) fails.
 \item One of the thirteen coordinates is omitted: \emph{object replaced;
 conclusion not transferable}.
\end{enumerate}
These falsifiers control membership and naturality; they do not degrade a
history that satisfies all the identities.
\end{ablation}
\fi
 \endgroup

The same rule defines its action on an elementary emission
\(\alpha:x\to y\). Since \(U_\alpha\) transports the ordered pair \emph{}
of sheets, its label map is
\(\overline U_\alpha e_{\Sigma,x}=e_{\Sigma,y}\) and
\(\overline U_\alpha e_{\Pi,x}=e_{\Pi,y}\). Its projectivization
\(\alpha_{\mathscr I}[v]:=[\overline U_\alpha v]\) acts in
\(\mathbf P^1(V_x)\) and, therefore,
\begin{equation}
\label{eq:pdfv2-cor-cell-accion-emision}
 \begin{gathered}
 \operatorname{Cell}_{\rm TPK}(\alpha)v_\star(x)=v_\star(y),\qquad
 \operatorname{Cell}_{\rm TPK}(\alpha)v_L(x)
   =v_{\alpha_{\mathscr I}L}(y),\\
 \operatorname{Cell}_{\rm TPK}(\alpha)v_{\{L,L'\}}(x)
   =v_{\{\alpha_{\mathscr I}L,\alpha_{\mathscr I}L'\}}(y),\qquad
 \operatorname{Cell}_{\rm TPK}(\alpha)v_{L,\epsilon}(x)
   =v_{\alpha_{\mathscr I}L,\epsilon}(y),\\
 \operatorname{Cell}_{\rm TPK}(\alpha)u_L(x)
   =u_{\alpha_{\mathscr I}L}(y),\qquad
 \operatorname{Cell}_{\rm TPK}(\alpha)u_{L,\epsilon}(x)
   =u_{\alpha_{\mathscr I}L,\epsilon}(y),\\
 \operatorname{Cell}_{\rm TPK}(\alpha)u_{L'\mid L}(x)
   =u_{\alpha_{\mathscr I}L'\mid\alpha_{\mathscr I}L}(y),\qquad
 \operatorname{Cell}_{\rm TPK}(\alpha)\omega_{\{L,L'\}}(x)
   =\omega_{\{\alpha_{\mathscr I}L,
                 \alpha_{\mathscr I}L'\}}(y).
 \end{gathered}
\end{equation}
In the leaf fiber and on its screen, one also has
\begin{equation}
\label{eq:pdfv2-cor-cell-entrelaza-proyectores-pantalla}
 \overline U_\alpha P_L=P_{\alpha_{\mathscr I}L}\overline U_\alpha,
 \qquad
 \operatorname{Cell}_{\rm TPK}(\alpha)\vartheta_L
 =\vartheta_{\alpha_{\mathscr I}L}
  \operatorname{Cell}_{\rm TPK}(\alpha).
\end{equation}
Boundary formulas give
\begin{equation}
\label{eq:pdfv2-cor-cell-accion-emision-natural}
 \partial\operatorname{Cell}_{\rm TPK}(\alpha)
 =\operatorname{Cell}_{\rm TPK}(\alpha)\partial,\qquad
 \operatorname{Cell}_{\rm TPK}(\beta\alpha)
 =\operatorname{Cell}_{\rm TPK}(\beta)
  \operatorname{Cell}_{\rm TPK}(\alpha).
\end{equation}
Thus the map of complexes used in the terminal register is the free action
induced by the same emission, not an added symbol.

\subsection{Module of paths, totalization and nonadic return cone}
Let \(J_d(x)\) be the complete category of prefixes of depth \(d\) of the
common tree, \(\operatorname{Term}_d(x)\) its terminal set and
\(\Gamma_{d,N}(\nu):=C_{d,N}(\nu)\) the cellular complex of
\eqref{eq:pdfv2-cor-cono-celda}. The module, totalization and cone that
follow simultaneously preserve all continuations; they do not select
a history or separate a cohomological construction from the common continuum.
\begingroup
  \let\section\subsection
  % Extracción quirúrgica del fuente teoremática V2 05d: líneas 286--512.
% El resto permanece en este archivo como procedencia, pero no se publica.
\iffalse
% Macroestructura de coherencia estrictamente interna del continuo discreto.
% No usa geometrías, clases ni conclusiones clásicas.

\chapter{Internal macrostructure of cells, coherence and completeness}
\label{chap:pdfv2-coh-interna}

\begin{statusbox}
The sole input of this chapter is the discrete structure of the continuum. The
branch \(\mathrm{coh}\) is constructed on the rings
\(A_N=\mathbb Z/3^N\mathbb Z\), the APP cells, their histories, and their
terminals. Its nonablative chain is
\[
 \mathcal C_{\rm cont}^{\rm disc}
 \supset\mathcal C_{\rm cont}^{\rm coh}
 \xrightarrow{\operatorname{Res}_{\rm coh}}
 \mathcal G_{\rm coh}
 \xrightarrow{\operatorname{pr}_{X,{\rm coh}}}
 X_{\chi,{\rm coh}}
 \xrightarrow{\mathcal A_{\rm coh}}
 \mathfrak M_{\rm coh}
 \xrightarrow{\mathcal P_{\rm coh}}
 \mathcal C_{\rm coh}^{\rm HMT}.
\]
No classical object appears as an input or as an output. Each image literally preserves its generating history.
\end{statusbox}

\section{Typed domain and index separation}
\label{sec:pdfv2-coh-domain}

To avoid identifying depth with precision, the directed set is used.
\begin{equation}\label{eq:pdfv2-coh-lambda}
 \Lambda:=\mathbb N_{\rm prof}\times\mathbb N_{\rm prec},
 \qquad
 \lambda=(d,N),
 \qquad
 A_N:=\mathbb Z/3^N\mathbb Z.
\end{equation}
If \(\lambda\leq\lambda'=(d',N')\), the transition
\begin{equation}\label{eq:pdfv2-coh-rho-lambda}
 \rho_{\lambda',\lambda}:
 \mathcal C_{{\rm cont},\lambda'}^{\rm disc}
 \longrightarrow
 \mathcal C_{{\rm cont},\lambda}^{\rm disc}
\end{equation}
truncates the history after depth \(d\) and reduces coefficients by
\(A_{N'}\twoheadrightarrow A_N\), preserving sheet, orientation, residue,
quotient, carry, memory, boundary, path, survival, and terminal.

The substructure \(\mathcal C_{\rm cont}^{\rm coh}\) consists of the histories
\(x\in\mathcal C_{\rm cont}^{\rm disc}\) that satisfy, at every level:
\begin{enumerate}
 \item the prefix of \(x\) carries the oriented and based cells defined in
       the following section;
 \item each refinement induce a map of the complex of the celda;
 \item truncation, coefficient reduction, orientation, and nonadic
       extension commute;
 \item the surviving-extension multisections are nonempty at every
       depth;
 \item the nonadic return preserves the terminal difference between the history and its continuation.
\end{enumerate}
These conditions define the domain of the branch. They are not attributed to an
arbitrary history without verifying them.

\section{APP cell over \texorpdfstring{\(A_N\)}{AN}}
\label{sec:pdfv2-coh-cell}

An oriented cell is
\[
 \nu=(h_\nu,I_\nu,J_\nu,\epsilon_\nu),
\]
where \(h_\nu\) is a finite history, and \(I_\nu,J_\nu\) are nonempty
ordered sets of ports,
\[
 a_\nu:=|I_\nu|\geq1,
 \qquad
 b_\nu:=|J_\nu|\geq1,
\]
and \(\epsilon_\nu\in\{+1,-1\}\) is its orientation.

We define
\begin{equation}\label{eq:pdfv2-coh-kappa}
 \begin{aligned}
 \kappa_{\nu,N}:
 A_N^{I_\nu}\oplus A_N^{J_\nu}
 &\longrightarrow A_N^{I_\nu\times J_\nu},\\
 \kappa_{\nu,N}(u,v)_{ij}&:=u_i-v_j.
 \end{aligned}
\end{equation}
Once the base ports \(i_0\in I_\nu\), \(j_0\in J_\nu\) have been chosen, let
\begin{equation}\label{eq:pdfv2-coh-delta}
 \begin{aligned}
 \delta_{\nu,N}:A_N^{I_\nu\times J_\nu}
 &\longrightarrow
 A_N^{(I_\nu\setminus i_0)\times(J_\nu\setminus j_0)},\\
 \delta_{\nu,N}(w)_{ij}
 &:={}
 w_{ij}-w_{i j_0}-w_{i_0j}+w_{i_0j_0}.
 \end{aligned}
\end{equation}

\begin{proposition}[Exact sequence of the cell]
\label{prop:pdfv2-coh-cell-exact}
For every celda \(\nu\) and every precision \(N\), the sucesion
\begin{equation}\label{eq:pdfv2-coh-exact-sequence}
 0\longrightarrow A_N
 \xrightarrow{\ \Delta\ }
 A_N^{I_\nu}\oplus A_N^{J_\nu}
 \xrightarrow{\ \kappa_{\nu,N}\ }
 A_N^{I_\nu\times J_\nu}
 \xrightarrow{\ \delta_{\nu,N}\ }
 A_N^{(a_\nu-1)(b_\nu-1)}
 \longrightarrow0,
\end{equation}
where
\[
 \Delta(t)=\bigl((t)_{i\in I_\nu},(t)_{j\in J_\nu}\bigr),
\]
is exact and split over \(A_N\).
\end{proposition}

\begin{proof}
If \(\kappa_{\nu,N}(u,v)=0\), then \(u_i=v_j\) for every \(i,j\), so \((u,v)\) is diagonal. Thus,
\(\ker\kappa_{\nu,N}=\operatorname{im}\Delta\). Direct substitution gives
\(\delta_{\nu,N}\kappa_{\nu,N}=0\).

If \(\delta_{\nu,N}(w)=0\), set
\[
 u_i=w_{i j_0},
 \qquad
 v_j=w_{i_0j_0}-w_{i_0j}.
\]
Then
\[
 u_i-v_j=w_{i j_0}+w_{i_0j}-w_{i_0j_0}=w_{ij},
\]
so that \(\ker\delta_{\nu,N}=\operatorname{im}\kappa_{\nu,N}\). To prove the surjectivity of \(\delta_{\nu,N}\), the base
row and column are set to zero and the desired coordinates are placed
arbitrarily in the remaining positions. Those same formulas provide
splittings.
\end{proof}

The complex cell, in degrees \(1\to0\), is the homological cone of
\(\kappa_{\nu,N}\) with the convention that places its source in degree \(1\):
\begin{equation}\label{eq:pdfv2-coh-cell-complex}
 C_{\nu,N}:=\operatorname{Cone}(\kappa_{\nu,N})
 :=\left[
 A_N^{I_\nu}\oplus A_N^{J_\nu}
 \xrightarrow{\ \kappa_{\nu,N}\ }
 A_N^{I_\nu\times J_\nu}
 \right].
\end{equation}
The preceding proposition gives
\begin{equation}\label{eq:pdfv2-coh-cell-homology}
 H_1(C_{\nu,N})\simeq A_N,
 \qquad
 H_0(C_{\nu,N})\simeq A_N^{(a_\nu-1)(b_\nu-1)}.
\end{equation}
The cell defect is thus constructed, not added as a cardinal without
map.

\section{Orientation and carry of the cell}
\label{sec:pdfv2-coh-orientation-carry}

Let
\[
 \tau_1(u,v)=(v,u),
 \qquad
 P(w)_{ji}=w_{ij}.
\]
Oriented exchange is the complex map
\begin{equation}\label{eq:pdfv2-coh-theta}
 \Theta_{\nu,N}:=(\tau_1,-P):
 C_{a_\nu,b_\nu,N}\longrightarrow C_{b_\nu,a_\nu,N},
\end{equation}
because
\begin{equation}\label{eq:pdfv2-coh-orientation}
 \kappa_{b_\nu,a_\nu,N}\tau_1
 =-P\kappa_{a_\nu,b_\nu,N}.
\end{equation}
Moreover,
\(\Theta_{\nu^{\rm op},N}\Theta_{\nu,N}=I\). The opposite route is not
identified with the original: it changes the degree-zero sign and preserves the act
of inversion.

For \(r\in A_N\), let
\(\langle r\rangle_N\in\{0,\ldots,3^N-1\}\) be its canonical representative. The
local borrow/carry is
\begin{equation}\label{eq:pdfv2-coh-beta}
 \beta_N(r,s):=
 \frac{\langle r-s\rangle_N-\langle r\rangle_N+\langle s\rangle_N}
 {3^N}\in\{0,1\}.
\end{equation}
Thus,
\begin{equation}\label{eq:pdfv2-coh-carry-cell}
 \langle u_i-v_j\rangle_N
 =\langle u_i\rangle_N-\langle v_j\rangle_N
  +3^N\beta_N(u_i,v_j).
\end{equation}
The matrix
\[
 \operatorname{Carr}_{\nu,N}(u,v)
 :=\bigl(\beta_N(u_i,v_j)\bigr)_{ij}
\]
is preserved alongside \(\kappa_{\nu,N}(u,v)\): residue and carry are two
distinct coordinates.

When port sizes follow from the two APP leaves, the admissibility
of the cell includes the integer identity
\begin{equation}\label{eq:pdfv2-coh-app-defect}
 (a_\nu-1)(b_\nu-1)
 =\rho_{\Pi,\nu}-\rho_{\Sigma,\nu}+1
  +9(c_{\Pi,\nu}-c_{\Sigma,\nu}).
\end{equation}
The equality is required only for a cell with that genealogy; it is not asserted
for integers devoid of APP leaves.

If \(L_N(u)\) is the lift whole canonica of the matrix of a flecha
\(u\), it define the carry of composition by
\begin{equation}\label{eq:pdfv2-coh-carry-composition}
 L_N(v)L_N(u)=L_N(vu)+3^N c_N(v,u).
\end{equation}
The associativity of three arrows gives the exact cocycle.
\begin{equation}\label{eq:pdfv2-coh-carry-cocycle}
 c_N(w,vu)+L_N(w)c_N(v,u)
 =c_N(wv,u)+c_N(w,v)L_N(u).
\end{equation}

\section{Category of histories and cellular functor}
\label{sec:pdfv2-coh-category}

For a prefix \(x_{\leq d}\), let \(J_d(x)\) be the finite category whose
objects are all its cells \(\nu\). Its arrows are generated by:
\begin{enumerate}
 \item refinements \(u:\nu\to\nu'\) with surjective maps of
       ports
       \[
        \alpha_u:I_{\nu'}\twoheadrightarrow I_\nu,
        \qquad
        \beta_u:J_{\nu'}\twoheadrightarrow J_\nu;
       \]
 \item the exchanges of orientation \(\Theta_\nu\);
 \item the nonadic extensions \(\Gamma_9\nu\);
 \item identity and composition, subject to the boundary,
       path, orientation, and carry relations of the history.
\end{enumerate}
Finiteness refers to the prefix; the system of all prefixes is not declared finite.

The cellular functor is
\begin{equation}\label{eq:pdfv2-coh-gamma-functor}
 \Gamma_{d,N}:J_d(x)\longrightarrow
 \operatorname{Ch}^{[0,1]}(A_N),
 \qquad
 \Gamma_{d,N}(\nu)=C_{\nu,N}.
\end{equation}
For a refinement arrow \(u:\nu\to\nu'\), it acts by means of
\begin{align}
 \Gamma_{d,N}(u)_1(p,q)
 &:=(p\circ\alpha_u,q\circ\beta_u),
 \label{eq:pdfv2-coh-gamma-one}\\
 \Gamma_{d,N}(u)_0(w)
 &:=w\circ(\alpha_u\times\beta_u).
 \label{eq:pdfv2-coh-gamma-zero}
\end{align}
It is verified coordinate by coordinate.
\begin{equation}\label{eq:pdfv2-coh-gamma-chain-map}
 \kappa_{\nu',N}\Gamma_{d,N}(u)_1
 =\Gamma_{d,N}(u)_0\kappa_{\nu,N}.
\end{equation}
On an inversion of orientation it acts by
\(\Theta_{\nu,N}\), and on compositions it acts by composition, separately preserving
\(c_N(v,u)\).

The reduction \(r_{N+1,N}:A_{N+1}\twoheadrightarrow A_N\) satisfies
\begin{equation}\label{eq:pdfv2-coh-kappa-reduction}
 r_{N+1,N}^{I\times J}\kappa_{\nu,N+1}
 =\kappa_{\nu,N}
  \bigl(r_{N+1,N}^{I}\oplus r_{N+1,N}^{J}\bigr).
\end{equation}
Therefore, the cell diagram is naturally simultaneous in history and
precision.

\fi
\section{Module of paths, cellular totalization and nonadic terminal}
\label{sec:pdfv2-coh-route-module}

Let \(\operatorname{Term}_d(x)\) be the conjunto of terminal of memory of the
prefix. For \(\nu\in J_d(x)\), definase
\begin{equation}\label{eq:pdfv2-coh-w-module}
 W_{d,N}(\nu):=
 A_N\bigl[
 \operatorname{Path}_{J_d(x)}(\nu,\operatorname{Term}_d(x))
 \bigr].
\end{equation}
It is the free module on all terminal routes; it selects none. If
\(u:\nu\to\nu'\), the contravariant action is
\begin{equation}\label{eq:pdfv2-coh-w-action}
 W_{d,N}(u):W_{d,N}(\nu')\longrightarrow W_{d,N}(\nu),
 \qquad
 [\lambda]\longmapsto[\lambda\circ u].
\end{equation}
Each generator preserves orientation, carry, boundary, route, and memory.

\section{Bilateral bar complex and totalization of the bicolumn}
\label{sec:pdfv2-coh-bar}

For \(q\geq0\) and \(r\in\{0,1\}\), let
\begin{equation}\label{eq:pdfv2-coh-bar-bigraded}
 B_{q,r}^{d,N}(x):=
 \bigoplus_{\nu_0\xrightarrow{u_1}\cdots\xrightarrow{u_q}\nu_q}
 W_{d,N}(\nu_q)\otimes_{A_N}\Gamma_{d,N}(\nu_0)_r.
\end{equation}
The bar differential is defined without abbreviation by
\begin{align}
 d_{\rm bar}(w;u_1,\ldots,u_q;c)
 :={}&(w;u_2,\ldots,u_q;\Gamma(u_1)c)
 \nonumber\\
 &+\sum_{i=1}^{q-1}(-1)^i
 (w;u_1,\ldots,u_{i+1}u_i,\ldots,u_q;c)
 \nonumber\\
 &+(-1)^q
 (W(u_q)w;u_1,\ldots,u_{q-1};c).
 \label{eq:pdfv2-coh-bar-differential}
\end{align}
For \(q=1\), this is exactly the relation
\begin{equation}\label{eq:pdfv2-coh-coend-relation}
 d_{\rm bar}(w;u;c)
 =w\otimes\Gamma(u)c-W(u)w\otimes c.
\end{equation}
A transported incidence is thereby identified with the same incidence written
in the refined chart; the relation is not concealed beneath the name of the
totalizer.

The total complex is
\begin{equation}\label{eq:pdfv2-coh-total-complex}
 K_{d,N}(x):=\operatorname{Tot}B_{\bullet,\bullet}^{d,N}(x),
 \qquad
 D_{\rm tot}:=d_{\rm bar}+(-1)^q\partial_\Gamma
 \quad\text{on }B_{q,r}^{d,N}.
\end{equation}
The functoriality of \(W\) and \(\Gamma\), together with
\(\partial_\Gamma^2=0\), gives
\begin{equation}\label{eq:pdfv2-coh-total-square-zero}
 D_{\rm tot}^2=0.
\end{equation}
In compact notation, but now already discharged by
\eqref{eq:pdfv2-coh-bar-bigraded}--
\eqref{eq:pdfv2-coh-bar-differential},
\begin{equation}\label{eq:pdfv2-coh-derived-coend}
 K_{d,N}(x)
 \simeq
 W_{d,N}\otimes^{\mathbf L}_{A_N[J_d(x)]}\Gamma_{d,N}.
\end{equation}

\section{Depth and terminal nonadic colimit}
\label{sec:pdfv2-coh-terminal}

The prefixes form the colimit.
\begin{equation}\label{eq:pdfv2-coh-k-infinity}
 K_{\infty,N}(x):=\varinjlim_d K_{d,N}(x).
\end{equation}
The nine-phase extension carries depth \(d\) to depth \(d+9\).
Only after forming \eqref{eq:pdfv2-coh-k-infinity} does it induce a
well-typed endomorphism
\begin{equation}\label{eq:pdfv2-coh-u-gamma9}
 U_{\Gamma_9,N}:K_{\infty,N}(x)
 \longrightarrow K_{\infty,N}(x).
\end{equation}
The terminal is
\begin{equation}\label{eq:pdfv2-coh-terminal-cone}
 \mathcal T_{{\rm coh},N}(x)
 :=\operatorname{Cone}\bigl(I-U_{\Gamma_9,N}\bigr).
\end{equation}
With the homological convention
\[
 \operatorname{Cone}(f)_q
 =K_{\infty,N,q-1}\oplus K_{\infty,N,q},
\]
its differential is
\begin{equation}\label{eq:pdfv2-coh-cone-differential}
 \partial_{\rm Cone}(a,b)
 =\bigl(-\partial_Ka,\partial_Kb+f(a)\bigr).
\end{equation}

Although the visible phase of \(h\) and \(\Gamma_9h\) coincides, the generators
\([h]\) and \([\Gamma_9h]\) are different because their ledgers and depths
are different. The terminal object preserves
\begin{equation}\label{eq:pdfv2-coh-terminal-difference}
 (I-U_{\Gamma_9,N})[h]=[h]-[\Gamma_9h].
\end{equation}
Precision reductions commute with the terminal:
\begin{equation}\label{eq:pdfv2-coh-u-reduction}
 r_{N+1,N}U_{\Gamma_9,N+1}
 =U_{\Gamma_9,N}r_{N+1,N}.
\end{equation}

\section{Terminal cohomological reader: incidence cycle and class in \(H_0(\operatorname{Cone}(I-U_{\Gamma_9,N}))\)}
\label{sec:pdfv2-coh-reader}

Each cell of \(x_{\leq d}\) carries a vector of incidence
\[
 \omega_{\nu,N}(x)
 \in A_N^{I_\nu\times J_\nu}=C_{\nu,N,0}
\]
and the complete terminal route \(\lambda_{\nu,x}\). The raw cycle is
\begin{equation}\label{eq:pdfv2-coh-z-cycle}
 z_{d,N}(x):=
 \sum_{\nu\in\operatorname{Cell}_d(x)}
 \epsilon_\nu[\lambda_{\nu,x}]
 \otimes\omega_{\nu,N}(x)
 \in B_{0,0}^{d,N}(x).
\end{equation}
As the total complex is homologically non-negative,
\begin{equation}\label{eq:pdfv2-coh-z-closed}
 D_{\rm tot}z_{d,N}(x)=0.
\end{equation}
Both the cycle and its class are preserved.
\[
 [z_{d,N}(x)]\in H_0(K_{d,N}(x)).
\]
After including it in \(K_{\infty,N}\), its terminal class is
\[
 [(0,z_{d,N}(x))]
 \in H_0(\mathcal T_{{\rm coh},N}(x)).
\]
The complete reader is
\begin{equation}\label{eq:pdfv2-coh-reader-full}
 \mathcal L_{{\rm coh},d,N}(x)
 :=\bigl(
 z_{d,N}(x),[z_{d,N}(x)],
 U_{\Gamma_9,N}z_{d,N}(x),
 [(0,z_{d,N}(x))]
 \bigr).
\end{equation}
The publication is not merely a quotient class: the representative, the route that produced it, its continuation and the terminal remain.

\section{Surviving fibres and corrections without a selector}
\label{sec:pdfv2-coh-surviving-fibers}

Let \(\mathscr S_\lambda^{\rm coh}\) be the free \(A_N\)-module generated by the
admissible finite histories at level \(\lambda=(d,N)\). The reader extends
linearly to
\begin{equation}\label{eq:pdfv2-coh-e-lambda}
 e_\lambda:\mathscr S_\lambda^{\rm coh}
 \longrightarrow\mathscr O_\lambda^{\rm coh},
 \qquad
 \mathscr O_\lambda^{\rm coh}
 :=H_0(K_{d,N})\oplus H_0(\mathcal T_{{\rm coh},N}).
\end{equation}
The surviving codomain is
\begin{equation}\label{eq:pdfv2-coh-surviving-output}
 \mathscr O_\infty^{\rm surv}
 :=\operatorname{Im}\!\left(
 e_\infty:\mathcal C_{\rm cont}^{\rm coh}
 \longrightarrow\varprojlim_\lambda\mathscr O_\lambda^{\rm coh}
 \right).
\end{equation}
For \(a=(a_\lambda)_\lambda\in\mathscr O_\infty^{\rm surv}\), define
\begin{equation}\label{eq:pdfv2-coh-fiber}
 \mathscr F_\lambda(a):=
 \left\{s_\lambda\in\mathscr S_\lambda^{\rm coh}:
 \begin{array}{l}
 e_\lambda(s_\lambda)=a_\lambda,\\
 \exists y\in\mathcal C_{\rm cont}^{\rm coh}\text{ with }
 y|_\lambda=s_\lambda\text{ and }e_\infty(y)=a
 \end{array}
 \right\}.
\end{equation}
These fibers are nonempty by the definition of the surviving codomain. If
\(\lambda'\geq\lambda\), the restriction is surjective:
\begin{equation}\label{eq:pdfv2-coh-fiber-surjection}
 \rho_{\lambda',\lambda}:\mathscr F_{\lambda'}(a)
 \twoheadrightarrow\mathscr F_\lambda(a).
\end{equation}
Indeed, a complete history that witnesses
\(s_\lambda\in\mathscr F_\lambda(a)\) provides by truncation an
antecedent at any later level. That antecedent is not chosen as
part of the object.

The complete multisect is
\begin{equation}\label{eq:pdfv2-coh-extension-multisection}
 \Sigma_{\lambda',\lambda}^{\rm coh}(s_\lambda;a)
 :=\{s_{\lambda'}\in\mathscr F_{\lambda'}(a):
 \rho_{\lambda',\lambda}s_{\lambda'}=s_\lambda\}.
\end{equation}
For a provisional prolongation \(\widetilde s_{\lambda'}\), let
\begin{equation}\label{eq:pdfv2-coh-discrepancy}
 \delta_{\lambda'}:=
 a_{\lambda'}-e_{\lambda'}(\widetilde s_{\lambda'}).
\end{equation}
The relation of all corrections is
\begin{equation}\label{eq:pdfv2-coh-correction-relation}
 \operatorname{Corr}_{\lambda',\lambda}
 (\widetilde s_{\lambda'},\delta_{\lambda'})
 :=\left\{\eta\in\mathscr S_{\lambda'}^{\rm coh}:
 \begin{array}{l}
 \rho_{\lambda',\lambda}\eta=0,\\
 e_{\lambda'}(\eta)=\delta_{\lambda'},\\
 \eta\text{ is supported on new cells},\\
 \eta\text{ preserves orientation, carry, memory and terminal}
 \end{array}
 \right\}.
\end{equation}
If two lifts are comparable, their difference
\(\eta=s_{\lambda'}-\widetilde s_{\lambda'}\) belongs to this relation
exactly when it satisfies the four preceding conditions. It is not asserted
that \eqref{eq:pdfv2-coh-correction-relation} is nonempty for an
arbitrary discrepancy: the governing non-emptiness is that of
\eqref{eq:pdfv2-coh-extension-multisection} on the surviving domain.

\iffalse
\section{The thirteen coordinates of the constraint}
\label{sec:pdfv2-coh-thirteen}

For \(x_\lambda\in\mathcal C_{{\rm cont},\lambda}^{\rm coh}\), let
\begin{equation}\label{eq:pdfv2-coh-d-package}
 \begin{aligned}
 D_{{\rm coh},\lambda}(x_\lambda):=\bigl(&
 \mathcal F_{{\rm coh},\lambda},
 \operatorname{Ext}_{\Gamma,{\rm coh},\lambda},
 \operatorname{Rel}_{{\rm coh},\lambda},
 \mathcal N_{{\rm coh},\lambda},
 \operatorname{or}_{{\rm coh},\lambda},
 \operatorname{Carr}_{{\rm coh},\lambda},
 \operatorname{Mem}_{{\rm coh},\lambda},\\
 &\partial_{{\rm coh},\lambda},
 \gamma_{{\rm coh},\lambda},
 \operatorname{Msect}_{{\rm coh},\lambda},
 \operatorname{Surv}_{{\rm coh},\lambda},
 \operatorname{Term}_{{\rm coh},\lambda},
 \mathcal L_{{\rm coh},\lambda}\bigr).
 \end{aligned}
\end{equation}
Their coordinates are:
\begin{enumerate}
 \item \(\mathcal F_{{\rm coh},\lambda}\): all the modulos
       \(A_N^{I_\nu}\), \(A_N^{J_\nu}\),
       \(A_N^{I_\nu\times J_\nu}\), the complejos \(C_{\nu,N}\), the
       modulos \(W_{d,N}(\nu)\) and \(\mathscr S_\lambda^{\rm coh}\).
 \item \(\operatorname{Ext}_{\Gamma,{\rm coh},\lambda}\): all maps
       \(\Gamma(u)\), coefficient reductions, truncations, extensions
       \(\Gamma_9\), and the relations \(\Sigma_{\lambda',\lambda}\).
 \item \(\operatorname{Rel}_{{\rm coh},\lambda}\):
       \(\kappa\), \(\delta\), the exact sequence, the bar faces, the
       coend relation, the orientation, and the composition laws.
 \item \(\mathcal N_{{\rm coh},\lambda}\):
       \[
       (d,N,3^N,a_\nu,b_\nu,(a_\nu-1)(b_\nu-1),
       \rho_{\Sigma,\nu},\rho_{\Pi,\nu},c_{\Sigma,\nu},c_{\Pi,\nu}).
       \]
 \item \(\operatorname{or}_{{\rm coh},\lambda}\): the signs
       \(\epsilon_\nu\), the port orderings, and the maps
       \(\Theta_{\nu,N}\).
 \item \(\operatorname{Carr}_{{\rm coh},\lambda}\): the matrices
       \(\beta_N\), the cocycles \(c_N(v,u)\), residues, and quotients,
       conserved separately.
 \item \(\operatorname{Mem}_{{\rm coh},\lambda}\): the inherited ledger and the
       ordered word of refinements, carries, corrections, and advances
       \(\Gamma_9\), not merely its aggregate value.
 \item \(\partial_{{\rm coh},\lambda}\):
       \(\kappa_{\nu,N}\), \(D_{\rm tot}\),
       \(\partial_{\rm Cone}\) and the boundaries of the cells.
 \item \(\gamma_{{\rm coh},\lambda}\): all composable chains
       \(\nu_0\to\cdots\to\nu_q\) and their terminal continuations.
 \item \(\operatorname{Msect}_{{\rm coh},\lambda}\): all cells,
       terminal routes, extensions \(\Sigma\), and corrections
       \(\operatorname{Corr}\), without a selector.
 \item \(\operatorname{Surv}_{{\rm coh},\lambda}\): the system of fibres
       \(\mathscr F_\lambda(a)\) and surjective restrictions.
 \item \(\operatorname{Term}_{{\rm coh},\lambda}\):
       \(\operatorname{Cone}(I-U_{\Gamma_9,N})\), together with \(z\), \(Uz\)
       and their difference.
 \item \(\mathcal L_{{\rm coh},\lambda}\):
       \((z,[z],Uz,[(0,z)])\) and all its reduction compatibilities.
\end{enumerate}

For \(\lambda'\geq\lambda\), the coordinated transport satisfies
\begin{equation}\label{eq:pdfv2-coh-d-naturality}
 u_{\lambda',\lambda}^{\rm coh}
 D_{{\rm coh},\lambda'}(x_{\lambda'})
 =D_{{\rm coh},\lambda}
  (\rho_{\lambda',\lambda}x_{\lambda'}).
\end{equation}

\section{Constraint graph and dependent object}
\label{sec:pdfv2-coh-graphs}

The total restriction is the graph
\begin{equation}\label{eq:pdfv2-coh-g-graph}
 \mathcal G_{\rm coh}:=\operatorname{Graph}(D_{\rm coh}),
 \qquad
 \operatorname{Res}_{\rm coh}(x):=(x,D_{\rm coh}x).
\end{equation}
Thus,
\begin{equation}\label{eq:pdfv2-coh-res-inverse}
 \pi_1\operatorname{Res}_{\rm coh}=I,
 \qquad
 \operatorname{Res}_{\rm coh}\pi_1=I_{\mathcal G_{\rm coh}}.
\end{equation}

The full extractor is
\begin{equation}\label{eq:pdfv2-coh-chi}
 \chi_{\rm coh}(x):=
 \bigl(
 \{C_{\nu,N}\},
 \{z_\lambda,[z_\lambda]\},
 \{\mathscr F_\lambda(e_\infty x)\},
 \{\Sigma_{\lambda',\lambda}\},
 \{\operatorname{Corr}_{\lambda',\lambda}\}
 \bigr).
\end{equation}
We define
\begin{equation}\label{eq:pdfv2-coh-xchi}
 X_{\chi,{\rm coh}}
 :=\{(\chi_{\rm coh}(x),x,D_{\rm coh}x):
 x\in\mathcal C_{\rm cont}^{\rm coh}\}
\end{equation}
and
\begin{equation}\label{eq:pdfv2-coh-prx}
 \operatorname{pr}_{X,{\rm coh}}(x,D_{\rm coh}x)
 :=(\chi_{\rm coh}(x),x,D_{\rm coh}x).
\end{equation}
The projection onto \((x,D_{\rm coh}x)\) is inverse on this object
dependent.

\section{Internal assembler and publisher}
\label{sec:pdfv2-coh-assembler-publisher}

The ambient space \(\mathscr M_{\rm coh}^{\rm amb}\) consists of systems
compatible of exact sequences of cells, categories, functors, modules of
path, totalizers, terminals, readers, fibers and multisections. The
assembler raw is
\begin{equation}\label{eq:pdfv2-coh-assembler}
 \begin{aligned}
 a_{\rm coh}:X_{\chi,{\rm coh}}&\longrightarrow
 \mathscr M_{\rm coh}^{\rm amb},\\
 a_{\rm coh}(\chi(x),x,Dx)
 &:={}
 \bigl(
 \{C_{\nu,N},\kappa_{\nu,N},\delta_{\nu,N}\},
 \{J_d,\Gamma,W,K\},\\
 &\hspace{31mm}
 \{U_{\Gamma_9,N},\mathcal T_N\},
 \{e_\lambda,\mathscr F_\lambda,\Sigma,\operatorname{Corr}\}
 \bigr)_x.
 \end{aligned}
\end{equation}
The macrostructure is another graph:
\begin{equation}\label{eq:pdfv2-coh-macro-graph}
 \mathfrak M_{\rm coh}:=\operatorname{Graph}(a_{\rm coh}),
 \qquad
 \mathcal A_{\rm coh}(z):=(z,a_{\rm coh}z).
\end{equation}

Let \(\mathscr Y_{\rm coh}^{\rm int}\) be the type of certificates containing
\begin{equation}\label{eq:pdfv2-coh-certificate-list}
 \begin{gathered}
 \text{la sucesión exacta \eqref{eq:pdfv2-coh-exact-sequence}},\\
 \kappa_{b,a}\tau=-P\kappa_{a,b},
 \qquad D_{\rm tot}^2=0,
 \qquad\partial_{\rm Cone}^2=0,\\
 rU_{\Gamma_9}=U_{\Gamma_9}r,
 \qquad
 e_\lambda\rho_{\lambda',\lambda}
 =r_{\lambda',\lambda}e_{\lambda'},\\
 \rho_{\lambda',\lambda}:\mathscr F_{\lambda'}(a)
 \twoheadrightarrow\mathscr F_\lambda(a),
 \qquad
 \Sigma_{\lambda',\lambda}(s_\lambda;a)\neq\varnothing,
 \end{gathered}
\end{equation}
together with the raw cycles, their classes and their terminals. The publisher
\begin{equation}\label{eq:pdfv2-coh-publisher}
 p_{\rm coh}^{\rm raw}:\mathfrak M_{\rm coh}
 \longrightarrow\mathscr Y_{\rm coh}^{\rm int}
\end{equation}
produces that certificate without eliminating the macrostructure. Finally,
\begin{equation}\label{eq:pdfv2-coh-output-graph}
 \mathcal C_{\rm coh}^{\rm HMT}
 :=\operatorname{Graph}(p_{\rm coh}^{\rm raw}),
 \qquad
 \mathcal P_{\rm coh}(m):=(m,p_{\rm coh}^{\rm raw}m).
\end{equation}

\section{Internal Theorem of Coherent Generation}
\label{sec:pdfv2-coh-theorem}

\begin{theorem}[Internal Generation of the Coherence Macrostructure]
\label{thm:pdfv2-coh-internal-generation}
On the explicit domain \(\mathcal C_{\rm cont}^{\rm coh}\), the chain
\[
 \mathcal C_{\rm cont}^{\rm coh}
 \xrightarrow{\operatorname{Res}_{\rm coh}}
 \mathcal G_{\rm coh}
 \xrightarrow{\operatorname{pr}_{X,{\rm coh}}}
 X_{\chi,{\rm coh}}
 \xrightarrow{\mathcal A_{\rm coh}}
 \mathfrak M_{\rm coh}
 \xrightarrow{\mathcal P_{\rm coh}}
 \mathcal C_{\rm coh}^{\rm HMT}
\]
naturally constructs the cells, their cones, the diagram of histories, the
two-sided bar, the totalization, the nonadic terminal, and the
surviving fibers. It conserves the thirteen coordinates and all corrections as
multisections. Each stage admits, on its image, a projection that recovers
the preceding stage.
\end{theorem}

\begin{proof}
The \cref{prop:pdfv2-coh-cell-exact} constructs each cell over \(A_N\), with
diagonal kernel and defect \((a_\nu-1)(b_\nu-1)\). The identity
\eqref{eq:pdfv2-coh-orientation} makes the oriented exchange an involutive map of
complexes. The formulas
\eqref{eq:pdfv2-coh-gamma-one}--\eqref{eq:pdfv2-coh-gamma-zero} prove that
each refinement is a map of complexes, whereas
\eqref{eq:pdfv2-coh-kappa-reduction} proves compatibility in precision.

Precomposition of paths makes \(W\). The identities
functorial of \(W\) and \(\Gamma\) imply
\(d_{\rm bar}^2=0\); its commutation with the cellular differential gives
\(D_{\rm tot}^2=0\). Thus, \(K_{d,N}\) is constructed by the bar formulae,
not by to nominal invocation of the coend.

The extension \(\Gamma_9\) increases the depth. Therefore it induces the
endomorphism \(U_{\Gamma_9,N}\) only after passage to the colimit
\(K_{\infty,N}\). The cone
\(\operatorname{Cone}(I-U_{\Gamma_9,N})\) then conserves the difference
\([h]-[\Gamma_9h]\) and does not identify phase repetition with repetition
of state.

The cycle \(z_{d,N}(x)\) is formed with all the cells, signs, routes, and incidences of the history. Truncation and reduction commute with the reader. If \(s_\lambda\in\mathscr F_\lambda(a)\), the complete history appearing in the very definition of the fiber provides an antecedent at every level \(\lambda'\geq\lambda\); this proves the surjectivity \eqref{eq:pdfv2-coh-fiber-surjection}. The object preserves the entire set \(\Sigma_{\lambda',\lambda}\), not a choice. Differences between lifts are registered through \eqref{eq:pdfv2-coh-correction-relation}, with support, orientation, carry, memory, and terminal made explicit.

Finally, \(\mathcal G_{\rm coh}\), \(\mathfrak M_{\rm coh}\), and
\(\mathcal C_{\rm coh}^{\rm HMT}\) are graphs, whereas
\(X_{\chi,{\rm coh}}\) retains \(x\) and \(D_{\rm coh}x\) as coordinates.
The first projection of each object recovers its predecessor. Therefore, the
complete composition does not replace history with class, multisection with
selector, or terminal object with visible phase.
\end{proof}

\section{Probative origin and forgers}
\label{sec:pdfv2-coh-provenance-falsifiers}

\paragraph{Result recovered.}
Belonging to the already constructed architecture are the cell \(\kappa_{a,b}\), its
diagonal kernel and defect, the orientation
\(\kappa_{b,a}\tau=-P\kappa_{a,b}\), the APP carry, the cone of the cell, the
category of histories, the totalization by coend, the action \(\Gamma_9\), the
terminal \(\operatorname{Cone}(I-U_{\Gamma_9})\), and the architecture of
surviving extension.

\paragraph{New Formalization.}
Explicit formalizations include the exact sequence \eqref{eq:pdfv2-coh-exact-sequence} over \(A_N\), the separation between depth and precision, the complete faces of the bar, formation of the terminal object only after the colimit, packaging of the thirteen coordinates, corrections as a relation without a selector, and the four dependent objects or graphs that prevent loss of genealogy.

\begin{criterion}[Internal falsifiers of the branch \(\mathrm{coh}\)]
\label{crit:pdfv2-coh-falsifiers}
The construction fails if any of the following occurs:
\begin{enumerate}
 \item the sequence is not exact
       \eqref{eq:pdfv2-coh-exact-sequence};
 \item \(\kappa_{b,a}\tau=-P\kappa_{a,b}\) fails;
 \item some refinement does not satisfy
       \(\kappa_{\nu'}\Gamma(u)_1=\Gamma(u)_0\kappa_\nu\);
 \item the composition carry does not satisfy
       \eqref{eq:pdfv2-coh-carry-cocycle};
 \item \(D_{\rm tot}^2\neq0\);
 \item reduction, reader, or terminal do not commute with \(U_{\Gamma_9}\);
 \item \(\operatorname{Cone}(I-U_{\Gamma_9})\) is formed by previously identifying
       the depths \(d\) and \(d+9\);
 \item a fiber declared to be surviving is empty, or a restriction is not
       surjective;
 \item is replaced \(\Sigma_{\lambda',\lambda}\) by a single
       lift;
 \item only \([z]\) is published, and \(z\), route, carry, memory, or
       terminal is eliminated;
 \item any of the thirteen coordinates is omitted.
\end{enumerate}
In the last case, the binding diagnosis is: \emph{object replaced;
conclusion not transferable}.
\end{criterion}

\begin{warningbox}
This chapter concludes only the internal construction of
\(\mathcal C_{\rm coh}^{\rm HMT}\). It introduces neither a classical space,
nor a classical class, nor a classical publication. Any later intertwiner
belongs to another document and does not govern the existence of this
macrostructure.
\end{warningbox}
\fi
 \endgroup

The local complex \eqref{eq:pdfv2-cor-complejo-local} and the vertex
\(y_\alpha\), endpoint of that same extension in the nerve, form the augmented
pullback. In it are defined
\begin{equation}
\label{eq:pdfv2-cor-zpm-h}
 \mathbf1_\alpha=(r,[y_\alpha]),\qquad
 z_-(\alpha)=(r,[y_\alpha]),
\end{equation}
\begin{equation}
\label{eq:pdfv2-cor-zplus-h}
 z_+(\alpha)=
 (r+3c_\alpha v_\alpha e+v_\alpha b,[y_\alpha]),
 \qquad
 h_*(\alpha)=(-v_\alpha f,0).
\end{equation}
The action on generators gives, without an additional abstract existence,
\begin{equation}
\label{eq:pdfv2-cor-filler}
 \partial z_-=\partial z_+=0,\qquad
 \partial h_*=z_--z_+.
\end{equation}
The correlative filling, the cell, and the collar preserve the same
\((\alpha,\operatorname{or}(\alpha),\kappa(\alpha),\partial\alpha,
\gamma\alpha,m(\gamma\alpha;m_0))\); they merely apply different internal readers
to the common record.

For the same finite based complex, the total determinant line is defined, without requiring balance or acyclicity.
\begin{equation}
\label{eq:pdfv2-cor-linea-determinante-total}
 \operatorname{Det}(C_{k,N}):=
 \bigotimes_q
 \left(\bigwedge^{\rm top}C_{k,N,q}\right)^{(-1)^q}.
\end{equation}
Refinement induces the determinantal morphism of the map of complexes and
simultaneously transports homology modules, Smith normal forms, and
Bockstein pages as fields of the register. Neither the existence of
\eqref{eq:pdfv2-cor-linea-determinante-total} nor its transport requires
homology to vanish. If the cone of refinement is contractible and oriented,
the morphism is canonically trivialized; in the general case the
cone is preserved and not replaced by a fictitious isomorphism.

\subsection{Smith Normal Form of the Terminal Operator and Primary Decomposition of Its Cokernel}
For each prefix \(y\), the unoriented carry and the oriented coefficient
of boundary of the common register admit the lifts
\begin{equation}
\label{eq:detint-c-v-lifts}
 \widetilde c_y,\widetilde v_y\in\mathbf Z,\qquad
 c_{y,n}:=\widetilde c_y\bmod3^n,\qquad
 v_{y,n}:=\widetilde v_y\bmod3^n.
\end{equation}
The determinantal output is neither an input nor a separate branch: it is the sector
produced by the same terminal reduction when the typed conditions
are verified
\begin{equation}
\label{eq:detint-dominio}
\begin{split}
 \mathcal C_{\rm cont}^{\rm det}:=\bigl\{x\in\mathcal C_{\rm cont}^{\rm disc}:{}&
 x\text{ survives at every depth};\\
 &\widehat P_x^{\rm red}\text{ is finite and based};\\
 &\chi(\widehat P_x^{\rm red})=0;\\
 &H_*(\widehat P_x^{\rm red}\otimes_{\mathbf Z_3}K_3)=0;\\
 &\text{the reduction produces a square block stable under refinement}
 \bigr\}.
\end{split}
\end{equation}
\begingroup
  \let\section\subsection
  % Extracción quirúrgica del fuente teoremática V2 05e: líneas 325--447.
% El resto permanece en este archivo como procedencia, pero no se publica.
\iffalse
\chapter[Internal Determinantal Macrostructure]{Internal Determinantal Macrostructure Generated by the Discrete Continuum}
\label{chap:detint-macroestructura}

\begin{thesisbox}
The \({\rm det}\) branch starts exclusively from surviving histories of
\(\mathcal C_{\rm cont}^{\rm disc}\). The nerve of extensions, the Alexander--Whitney diagonal,
the local complex, the unit, the two publications, the filler, the terminal
reduction, and the \(3\)-adic tower are constructed before any subsequent
evaluation. The balance and acyclicity conditions are incorporated into the
typed domain and are never tacitly assumed.
\end{thesisbox}

\section{Surviving determinant domain}

Let
\begin{equation}\label{eq:detint-an}
 A_n:=\mathbf Z/3^n\mathbf Z,
 \qquad
 \rho_{n+1,n}:A_{n+1}\longrightarrow A_n.
\end{equation}
For \(x\in\mathcal C_{\rm cont}^{\rm disc}\), let
\(\mathsf T_m(x)\) be the finite category of TPK prefixes of depth \(m\)
compatible with \(x\). Each object preserves sheet, orientation, residue,
quotient, integer carry, memory, boundary, route, multisection, survival,
and terminal; its morphisms are the admissible extensions that preserve those
coordinates.

From each object \(y\in\mathsf T_m(x)\), the internal integers are extracted
\begin{equation}\label{eq:detint-c-v-lifts}
 \widetilde c_y,\widetilde v_y\in\mathbf Z,
 \qquad
 c_{y,n}:=\widetilde c_y\bmod3^n,
 \qquad
 v_{y,n}:=\widetilde v_y\bmod3^n.
\end{equation}
The first coordinate is the unoriented carry and the second, the oriented
boundary coefficient. Reversal satisfies
\begin{equation}\label{eq:detint-or-cv}
 c_{\bar y,n}=c_{y,n},
 \qquad
 v_{\bar y,n}=-v_{y,n}.
\end{equation}

We write
\begin{equation}\label{eq:detint-z3-k3}
 \mathbf Z_3:=\varprojlim_n A_n,
 \qquad
 K_3:=\operatorname{Frac}(\mathbf Z_3).
\end{equation}
The determinantial restriction is defined by
\begin{equation}\label{eq:detint-dominio}
\begin{split}
 \mathcal C_{\rm cont}^{\rm det}:=
 \bigl\{x\in\mathcal C_{\rm cont}^{\rm disc}:{}&
 x\text{ survives at every depth};\\
 &\widehat P_x^{\rm red}\text{ is finite and based};\\
 &\chi(\widehat P_x^{\rm red})=0;\\
 &H_*(\widehat P_x^{\rm red}\otimes_{\mathbf Z_3}K_3)=0;\\
 &\text{the canonical reduction produces a square block}\\
 &\text{stable under refinement}\bigr\}.
\end{split}
\end{equation}
The objects \(\widehat P_x^{\rm red}\) and the square block will be constructed
in \cref{sec:detint-terminal}. Therefore
\eqref{eq:detint-dominio} is a verifiable membership condition and not
a conclusion inserted into the definition of the maps.

\section{Nerve, Alexander--Whitney, and counit}

The normalized complex is defined
\begin{equation}\label{eq:detint-gmn}
 G_{m,n}(x):=
 N_*^{\rm norm}\!\left(N\mathsf T_m(x);A_n\right).
\end{equation}
For a nondegenerate simplex
\(\sigma=[y_0\to y_1\to\cdots\to y_q]\), the diagonal is
\begin{equation}\label{eq:detint-aw}
 \Delta_{\rm AW}(\sigma)
 =\sum_{j=0}^q
 [y_0\to\cdots\to y_j]\otimes
 [y_j\to\cdots\to y_q].
\end{equation}
The counit is given by
\begin{equation}\label{eq:detint-counidad-g}
 \epsilon_G([y])=1,
 \qquad
 \epsilon_G(\sigma)=0\quad(q>0).
\end{equation}
In particular,
\begin{equation}\label{eq:detint-vertice-grouplike}
 \Delta_{\rm AW}[y]=[y]\otimes[y],
 \qquad
 \epsilon_G([y])=1.
\end{equation}

If \(m'\ge m\) and \(n'\ge n\), truncation and coefficient reduction
induce
\begin{equation}\label{eq:detint-rmn}
 R_{m',m}^{n',n}:G_{m',n'}(x)\longrightarrow G_{m,n}(x),
\end{equation}
and they are verified
\begin{equation}\label{eq:detint-aw-natural}
 R\partial=\partial R,
 \qquad
 (R\otimes R)\Delta_{\rm AW}=\Delta_{\rm AW}R,
 \qquad
 \epsilon_GR=\rho_{n',n}\epsilon_G.
\end{equation}
The path is not reduced to its endpoints: the entire simplex remains in
\(G_{m,n}(x)\).

\section{Internal local complex}

For \(c\in A_n\), define
\begin{equation}\label{eq:detint-p0-grados}
 P^0_{n,1}(c)=A_nf,
 \qquad
 P^0_{n,0}(c)=A_nr\oplus A_ne\oplus A_nb,
 \qquad
 P^0_{n,-1}(c)=A_ng,
\end{equation}
with differential
\begin{equation}\label{eq:detint-p0-diferencial}
 \partial f=3ce+b,
 \qquad
 \partial r=0,
 \qquad
 \partial e=g,
 \qquad
 \partial b=-3cg.
\end{equation}
The chain identity is literal:
\begin{equation}\label{eq:detint-d2}
 \partial^2f=3c\,\partial e+\partial b=3cg-3cg=0,
 \qquad
 \partial^2e=\partial^2b=\partial^2r=0.
\end{equation}

The increase
\begin{equation}\label{eq:detint-epsilon-p}
 \epsilon_P:P_n^0(c)\longrightarrow A_n[0]
\end{equation}
is fixed by
\[
 \epsilon_P(r)=1,
 \qquad
 \epsilon_P(e)=\epsilon_P(b)=\epsilon_P(f)=\epsilon_P(g)=0.
\]
The complex possesses the differential coalgebra.
\begin{equation}\label{eq:detint-coalgebra-p}
 \Delta_P(r)=r\otimes r,
 \qquad
 \Delta_P(q)=q\otimes r+r\otimes q
 \quad(q\in\{e,b,f,g\}).
\end{equation}
Thus, \(r\) is group-like and the other four generators are primitive
relative to \(r\). Equations
\eqref{eq:detint-p0-diferencial} prove that \(\Delta_P\) commutes with the
tensor differential.

\section{Local extension system and memory ledger}

For an extension \(\alpha:y\to y'\), we define
\begin{equation}\label{eq:detint-psi}
 \Psi_\alpha:P_n^0(c_y)\longrightarrow P_n^0(c_{y'})
\end{equation}
through means of
\begin{equation}\label{eq:detint-psi-generadores}
\begin{gathered}
 \Psi_\alpha(r_y)=r_{y'},\quad
 \Psi_\alpha(e_y)=e_{y'},\quad
 \Psi_\alpha(g_y)=g_{y'},\quad
 \Psi_\alpha(f_y)=f_{y'},\\
 \Psi_\alpha(b_y)=b_{y'}+3(c_{y'}-c_y)e_{y'}.
\end{gathered}
\end{equation}
It is a map of complexes because
\begin{equation}\label{eq:detint-psi-b-chain}
 \partial\Psi_\alpha(b_y)
 =-3c_{y'}g+3(c_{y'}-c_y)g
 =-3c_yg
 =\Psi_\alpha(\partial b_y),
\end{equation}
and
\begin{equation}\label{eq:detint-psi-f-chain}
 \Psi_\alpha(\partial f_y)
 =3c_ye+b+3(c_{y'}-c_y)e
 =3c_{y'}e+b
 =\partial\Psi_\alpha(f_y).
\end{equation}
Moreover,
\begin{equation}\label{eq:detint-psi-functorial}
 \Psi_\beta\Psi_\alpha=\Psi_{\beta\alpha}.
\end{equation}

The oriented coordinate change does not erase. It is recorded by
\begin{equation}\label{eq:detint-k-alpha}
 K_\alpha:=(v_{y'}-v_y)f_{y'}.
\end{equation}
For composable extensions,
\begin{equation}\label{eq:detint-k-coherencia}
 K_{\beta\alpha}=K_\beta+\Psi_\beta K_\alpha.
\end{equation}
This equality is the exact ledger of the boundary increment and makes visible
the memory that a purely nominal composition would lose.

The complete local diagram is the Grothendieck construction.
\begin{equation}\label{eq:detint-grothendieck-local}
 \int_{y\in\mathsf T_m(x)}P_n^0(c_y).
\end{equation}
In it, an arrow over \(\alpha:y\to y'\) carries \(\Psi_\alpha\)
as its linear component and \(K_\alpha\) as its coherence. When applied to an
element supported at a vertex, the notation is used
\begin{equation}\label{eq:detint-psi-soporte-vertice}
 \Psi_\alpha(p,[y]):=(\Psi_\alpha p,[y']).
\end{equation}
The edge \([y\to y']\) remains in the path coordinate of the nerve. This
convention will be used in the naturality formulas for
\(z_\pm\) and \(h_*\); no global map erasing the remaining
simplices of \(G_{m,n}(x)\) is asserted.

\section{Pullback, unit, root, publications and filler}

The pullback of complexes is
\begin{equation}\label{eq:detint-pgen}
 P^{\rm gen}_{y;m,n}
 :=P_n^0(c_y)\times_{A_n[0]}G_{m,n}(x)
 =\{(p,\gamma):\epsilon_P(p)=\epsilon_G(\gamma)\}.
\end{equation}
Its differential is
\begin{equation}\label{eq:detint-pgen-d}
 \partial(p,\gamma)=(\partial p,\partial\gamma).
\end{equation}
The common unit associated with the vertex \(y\) is
\begin{equation}\label{eq:detint-unidad}
 \mathbf1_y:=(r,[y]).
\end{equation}
Its two projections are group-like by
\eqref{eq:detint-vertice-grouplike} and
\eqref{eq:detint-coalgebra-p}; that is the precise meaning of a compatible
group-like unit of the pullback.

The root projection is the map of complexes
\begin{equation}\label{eq:detint-proyeccion-root}
 \varpi_{\rm root}:P^{\rm gen}_{y;m,n}\longrightarrow
 \operatorname{Root}_{y;m,n},
 \qquad
 \varpi_{\rm root}(p,\gamma):=(\epsilon_P(p)r,\gamma).
\end{equation}
The two internal publications are defined.
\begin{equation}\label{eq:detint-zpm}
 z_-(y):=(r,[y]),
 \qquad
 z_+(y):=(r+3c_yv_ye+v_yb,[y]).
\end{equation}
Both are cycles:
\begin{equation}\label{eq:detint-zpm-ciclos}
 \partial z_-(y)=0,
 \qquad
 \partial z_+(y)=3c_yv_yg-3c_yv_yg=0.
\end{equation}
Their roots coincide:
\begin{equation}\label{eq:detint-raiz-comun}
 \varpi_{\rm root}z_-(y)
 =(r,[y])
 =\varpi_{\rm root}z_+(y).
\end{equation}

The filler is constructed directly:
\begin{equation}\label{eq:detint-hstar}
 h_*(y):=(-v_yf,0)\in P^{\rm gen}_{y;m,n,1}.
\end{equation}
Then
\begin{equation}\label{eq:detint-filler-identidad}
 \partial h_*(y)
 =(-3c_yv_ye-v_yb,0)
 =z_-(y)-z_+(y).
\end{equation}
In particular,
\begin{equation}\label{eq:detint-clases-iguales}
 [z_-(y)]=[z_+(y)]
 \quad\text{in}\quad H_0(P^{\rm gen}_{y;m,n}),
\end{equation}
and the equality preserves its witness chain.

Under an extension \(\alpha:y\to y'\), we have
\begin{equation}\label{eq:detint-z-natural}
 z_-(y')=\Psi_\alpha z_-(y),
 \qquad
 z_+(y')-\Psi_\alpha z_+(y)=\partial K_\alpha,
\end{equation}
and
\begin{equation}\label{eq:detint-h-natural}
 h_*(y')-\Psi_\alpha h_*(y)=-K_\alpha.
\end{equation}
Therefore, naturality does not identify distinct values of \(v\): it records
their difference through the coherent homotopy \(K_\alpha\).

\section{Internal Orientation}

On \(P_n^0(c)\) is defined
\begin{equation}\label{eq:detint-omega-p}
 \Omega(r)=r,
 \qquad
 \Omega(q)=-q\quad(q\in\{e,b,f,g\}).
\end{equation}
It is an automorphism of complexes and of the relative coalgebra. In the nerve,
the oriented inversion is
\begin{equation}\label{eq:detint-omega-g}
 \mathfrak o_G[y_0\to\cdots\to y_q]
 :=(-1)^{q(q+1)/2}
 [\bar y_q\to\cdots\to\bar y_0].
\end{equation}
Together with \eqref{eq:detint-or-cv}, these actions give
\begin{equation}\label{eq:detint-or-z-h}
 \Omega z_-(y)=z_-(\bar y),
 \qquad
 \Omega z_+(c_y,v_y)=z_+(c_{\bar y},v_{\bar y}),
 \qquad
 \Omega h_*(y)=h_*(\bar y).
\end{equation}
Orientation acts on path, complex, publications, filler and determinant
line; it is not a label without an action.

\fi
\section{Lifting of the cokernel through the Bockstein tower and preservation of carry}
\label{sec:detint-terminal}

Only the common root is eliminated:
\begin{equation}\label{eq:detint-p-reducido}
 \overline P_{y;m,n}
 :=P^{\rm gen}_{y;m,n}/A_n\mathbf1_y,
 \qquad
 \widehat P_{y;m}^{\rm red}:=\varprojlim_n\overline P_{y;m,n}.
\end{equation}
The base is first ordered by the TPK order of routes and then by
\begin{equation}\label{eq:detint-orden-base}
 f\prec e\prec b\prec g.
\end{equation}
The following internal algorithm is executed.
\begin{enumerate}
 \item Find the coefficients that are units of \(\mathbf Z_3\).
 \item Choose the first according to \eqref{eq:detint-orden-base} and the order of
 paths.
 \item Cancel the corresponding pair and record the elementary
 operations, their sign, and their contribution to the orientation.
 \item Repeat until no cancellable unit pivot remains.
\end{enumerate}
The survival condition of \eqref{eq:detint-dominio} requires that,
cofinally in \(m\), the result be a two-term block
\begin{equation}\label{eq:detint-sx}
 \mathbf Z_3^{r_x}\xrightarrow{S_x}\mathbf Z_3^{r_x},
 \qquad
 \det S_x\ne0\text{ in }K_3,
\end{equation}
and that a refinement only add unit contractible blocks and recorded
changes of basis. If these conditions fail, the history does not belong to
\(\mathcal C_{\rm cont}^{\rm det}\); a nonexistent terminal is not forced.

The identity route provides an elementary witness. After removing the root,
the local complex is
\begin{equation}\label{eq:detint-testigo-contractil}
 A_nf\xrightarrow{\binom{3c}{1}}
 A_ne\oplus A_nb\xrightarrow{(1,-3c)}A_ng.
\end{equation}
Is contractible via
\begin{equation}\label{eq:detint-contraccion-local}
 s(g)=e,
 \qquad
 s(ae+\beta b)=\beta f,
 \qquad
 s(f)=0.
\end{equation}
Thus, the terminal conditions admit at least the internal identity witness
and do not describe a formally empty subdomain.

\section{Smith, \texorpdfstring{\(3\)}{3}-adic order and initial unit}

In oriented bases, a valuation diagonalization has the form
\begin{equation}\label{eq:detint-smith}
 U_xS_xV_x
 =\operatorname{diag}(3^{a_1}u_1,\ldots,3^{a_r}u_r),
 \qquad
 a_i\in\mathbf N,quad u_i\in\mathbf Z_3^\times.
\end{equation}
We define
\begin{equation}\label{eq:detint-orden-d}
 d_x:=v_3(\det S_x)=\sum_{i=1}^r a_i
\end{equation}
and the initial unit
\begin{equation}\label{eq:detint-leading-unit}
 \lambda_x:=3^{-d_x}\det S_x\in\mathbf Z_3^\times.
\end{equation}
The orientation coordinate trivializes the determinant line. If a
transition produces
\begin{equation}\label{eq:detint-estabilidad-bloques}
 S_x\oplus I_q=U S_{x'}V,
\end{equation}
the orientation is transported by the inverse factor of \(\det U\det V\).
The oriented unit \(\lambda_x^{\rm or}\) is then natural and does not change
when contractible blocks are added.

At finite precision,
\begin{equation}\label{eq:detint-smith-finito}
 S_{x,n}=S_x\bmod3^n,
 \qquad
 a_i^{(n)}=\min(a_i,n).
\end{equation}
The terminal is, therefore, the compatible family of all precisions
and not an isolated integer.

\section{Bockstein Tower and Higher Order Carry}

Let
\[
 C_x^1\xrightarrow{S_x}C_x^0
\]
the terminal block. If a class \([\bar y]\) admits a lift
\(y\in C_x^1\) with \(S_xy\in3^qC_x^0\), we define
\begin{equation}\label{eq:detint-beta-q}
 \beta_q[\bar y]
 :=\left[3^{-q}S_xy\bmod3\right].
\end{equation}
Changing the lift alters the representative by a boundary of the corresponding
page, so \(\beta_q\) is well defined. The higher
carry is
\begin{equation}\label{eq:detint-carry-superior}
 \operatorname{Carr}_q(y):=3^{-q}S_xy.
\end{equation}
It does not replace the TPK carry of \eqref{eq:detint-c-v-lifts}; it records its
prolongation in the coefficient tower.

For a block \(3^{a_i}u_i\), the class persists until page \(a_i\),
and on that page the nonzero differential is multiplication by
\(\bar u_i\in\mathbf F_3^\times\). Consequently,
\begin{equation}\label{eq:detint-bock-longitudes}
 \{\text{survival lengths}\}=\{a_1,\ldots,a_r\},
 \qquad
 d_x=\sum_i a_i.
\end{equation}
The oriented product of the page trivializations satisfies
\begin{equation}\label{eq:detint-bock-leading}
 \Theta_{\rm Bock}(x)=\overline{\lambda_x^{\rm or}}
 \in\mathbf F_3^\times.
\end{equation}
The proof is carried out in each diagonal block and transported by the recorded
unimodular changes. Smith, Bockstein and the initial unit are thus three
readings of the same terminal.

\iffalse
\section{The thirteen coordinates of the determinantal constraint}

For each \((m,n)\), it is defined
\begin{equation}\label{eq:detint-d-trece}
 D_{{\rm det},m,n}(x):=
 (\mathcal F,\operatorname{Ext}_\Gamma,\operatorname{Rel},
 \mathcal N,\operatorname{or},\operatorname{Carr},\operatorname{Mem},
 \partial,\gamma,\operatorname{Msect},\operatorname{Surv},
 \operatorname{Term},\mathcal L)_{{\rm det},m,n}(x),
\end{equation}
with the following content.
\begin{enumerate}
 \item The fiber is
 \[
  \mathcal F=(A_n,\mathsf T_m(x),
  \{P_n^0(c_y)\}_y,\{P^{\rm gen}_{y;m,n}\}_y).
 \]
 \item \(\operatorname{Ext}_\Gamma\) contains all morphisms
 \(\alpha\), the maps \(\Psi_\alpha\), the homotopies \(K_\alpha\) and
 the maps \(R_{m',m}^{n',n}\).
 \item \(\operatorname{Rel}\) contains \(\partial^2=0\), the
 pullback equation, AW, counit, \(\Psi_\beta\Psi_\alpha=\Psi_{\beta\alpha}\) and
 \(K_{\beta\alpha}=K_\beta+\Psi_\beta K_\alpha\).
 \item
 \(\mathcal N=(m,n,\widetilde c_y,\widetilde v_y,c_{y,n},v_{y,n},
 r_x,a_1,\ldots,a_{r_x},d_x)\).
 \item \(\operatorname{or}\) contains \(y\mapsto\bar y\),
 \(\mathfrak o_G\), \(\Omega\) and the orientation of the determinant
 line.
 \item \(\operatorname{Carr}\) separately preserves the integer carry,
 its reduction \(c_{y,n}\) and all higher carries
 \(\operatorname{Carr}_q\).
 \item \(\operatorname{Mem}\) contiene
 \[
  y_0\to\cdots\to y_m,
  \quad\{(\widetilde c_{y_j},\widetilde v_{y_j},\mathbf1_{y_j})\}_{j=0}^m,
  \quad\{K_{\alpha_j}\}.
 \]
 \item
 \(\partial=(\partial_G,\partial_P,\partial_{P^{\rm gen}},
 \partial_{\overline P},h_*)\).
 \item \(\gamma=[y_0\to\cdots\to y_q]\) preserves the complete oriented
 simplex.
 \item \(\operatorname{Msect}\) contains all compatible
 extensions, all \(3\)-adic lifts, and all unimodular
 presentations of the same oriented terminal; it does not choose one a posteriori.
 \item \(\operatorname{Surv}\) records arbitrary depth,
 coefficient compatibility, balance, acyclicity over \(K_3\),
 stability under unit blocks and persistence of the root.
 \item
 \[
  \operatorname{Term}=(\widehat P_x^{\rm red},S_x,
  \{a_i,u_i\},d_x,\lambda_x^{\rm or},
  \{\beta_q\},\Theta_{\rm Bock}).
 \]
 \item The internal reader is
 \[
  \mathcal L_{\rm det}^{\rm int}
  =(\mathbf1_y,
  \varpi_{\rm root}z_-=\varpi_{\rm root}z_+,
  [z_-]=[z_+],d_x,\lambda_x^{\rm or},\Theta_{\rm Bock}).
 \]
\end{enumerate}

The transitions satisfy, coordinate by coordinate,
\begin{equation}\label{eq:detint-naturalidad-trece}
 u^{\rm det}_{(m',n'),(m,n)}D_{{\rm det},m',n'}(x')
 =D_{{\rm det},m,n}(\tau_{m',m}x').
\end{equation}

\section{Assembler, publisher, and non-ablative chain}

We define
\begin{equation}\label{eq:detint-graph-d}
 \mathcal G_{\rm det}:=\operatorname{Graph}(D_{\rm det}),
 \qquad
 \operatorname{Res}_{\rm det}(x):=(x,D_{\rm det}x).
\end{equation}
Let \(\chi_{\rm det}(x)\) be the complete multisection of germs
\[
 (y,\widetilde c_y,\widetilde v_y,\operatorname{or}_y,\gamma_y)
\]
compatible at all depths. The dependent object is
\begin{equation}\label{eq:detint-x-dependiente}
 X_{\chi,{\rm det}}
 :=\{(\chi_{\rm det}(x),x,D_{\rm det}x):
 x\in\mathcal C_{\rm cont}^{\rm det}\},
\end{equation}
and
\begin{equation}\label{eq:detint-prx}
 \operatorname{pr}_{X,{\rm det}}(x,D_{\rm det}x)
 :=(\chi_{\rm det}(x),x,D_{\rm det}x).
\end{equation}

The raw assembler
\begin{equation}\label{eq:detint-ensamblador-crudo}
 a_{\rm det}:X_{\chi,{\rm det}}\longrightarrow
 \mathscr M_{\rm det}^{\rm amb}
\end{equation}
yields
\begin{equation}\label{eq:detint-ensamblador-salida}
\begin{split}
 a_{\rm det}(z)=\bigl(&
 \{G_{m,n},\Delta_{\rm AW},\epsilon_G\}_{m,n};\\
 &\{P_n^0(c),\Psi_\alpha,K_\alpha,P^{\rm gen},
 \mathbf1,z_-,z_+,h_*\}_{m,n};\\
 &\{\widehat P^{\rm red},S,d,\lambda^{\rm or},
 \beta_q,\Theta_{\rm Bock}\}\bigr).
\end{split}
\end{equation}
The macroobject preserves its entry:
\begin{equation}\label{eq:detint-graph-a}
 \mathfrak M_{\rm det}:=\operatorname{Graph}(a_{\rm det}),
 \qquad
 \mathcal A_{\rm det}(z):=(z,a_{\rm det}z).
\end{equation}

The raw publisher
\begin{equation}\label{eq:detint-publicador-crudo}
 p_{\rm det}:\mathfrak M_{\rm det}\longrightarrow\mathscr Y_{\rm det}
\end{equation}
publishes the certificate jointly
\begin{equation}\label{eq:detint-certificado-publicado}
\begin{gathered}
 \partial^2=0,quad \text{AW and counit},quad
 \mathbf1_y\text{ compatible group-like},\\
 \partial z_\pm=0,quad
 \varpi_{\rm root}z_-=\varpi_{\rm root}z_+,quad
 \partial h_*=z_--z_+,\\
 \Psi_\beta\Psi_\alpha=\Psi_{\beta\alpha},quad
 K_{\beta\alpha}=K_\beta+\Psi_\beta K_\alpha,\\
 d=v_3(\det S)=\sum_i a_i,quad
 \Theta_{\rm Bock}=\overline{\lambda^{\rm or}},quad
 \text{full naturality under refinement}.
\end{gathered}
\end{equation}
Finally,
\begin{equation}\label{eq:detint-graph-p}
 \mathcal C_{\rm det}^{\rm HMT}:=\operatorname{Graph}(p_{\rm det}),
 \qquad
 \mathcal P_{\rm det}(M):=(M,p_{\rm det}M).
\end{equation}
The complete chain is
\begin{equation}\label{eq:detint-cadena-completa}
 \mathcal C_{\rm cont}^{\rm disc}\supset
 \mathcal C_{\rm cont}^{\rm det}
 \xrightarrow{\operatorname{Res}_{\rm det}}\mathcal G_{\rm det}
 \xrightarrow{\operatorname{pr}_{X,{\rm det}}}X_{\chi,{\rm det}}
 \xrightarrow{\mathcal A_{\rm det}}\mathfrak M_{\rm det}
 \xrightarrow{\mathcal P_{\rm det}}\mathcal C_{\rm det}^{\rm HMT}.
\end{equation}
Each arrow has an inverse on its image given by a coordinate
projection. History is not replaced by value, multisection by selector,
complex by determinant or proof by name.

\section{Theorem of Internal Determinantal Realization}

\begin{theorem}[Complete internal determinant realization]
\label{thm:detint-realizacion}
For every \(x\in\mathcal C_{\rm cont}^{\rm det}\), the chain
\eqref{eq:detint-cadena-completa} constructs a macro-object natural in
depth and precision. It contains two cycles with a common root, an explicit filler
that identifies them, a compatible group-like unit, and a
\(3\)-adic terminal whose order, Smith form, and Bockstein tower produce the same
oriented reader. The thirteen coordinates of \eqref{eq:detint-d-trece}
remain recoverable in the output.
\end{theorem}

\begin{proof}
Equations \eqref{eq:detint-p0-diferencial}--
\eqref{eq:detint-d2} prove that \(P_n^0(c)\) is a complex. The two
augmentations are maps of complexes, so the pullback is typed.
Alexander--Whitney is natural and makes each vertex group-like; together with
\(r\), this gives the common unit. The direct calculation
\eqref{eq:detint-zpm-ciclos}--\eqref{eq:detint-filler-identidad} proves that
the publications are cycles, have the same root and are joined by
\(h_*\). The formulas for \(\Psi_\alpha\) prove functoriality, and
\(K_\alpha\) proves homotopic naturality while preserving the memory
increment. The definition of \(\mathcal C_{\rm cont}^{\rm det}\) ensures that
terminal reduction produces a square block without postulating it outside the
domain. Smith over \(\mathbf Z_3\) gives \(d=\sum_i a_i\); the calculation of each
block \(3^{a_i}u_i\) gives
\(\Theta_{\rm Bock}=\overline{\lambda^{\rm or}}\). Truncation, reduction and
addition of unit blocks preserve these identities. Finally, the
three graph constructions retain the input as the first coordinate,
so the complete composition is non-ablative.
\end{proof}

\section{Provenance and Forgers}

\paragraph{Origin.}
The TPK nerve, Alexander--Whitney, the counit, the common unit, the normal
form \(\partial f=3ce+b\), \(\partial e=g\),
\(\partial b=-3cg\), the common root, the filler, and the tower architecture
have the status of \texttt{PREEXISTING\_AUTHORIAL\_ARCHITECTURE}. The literal computation
of the two cycles, their root, and the identity
\(\partial h_*=z_--z_+\) has the status of a \texttt{RECOVERED\_RESULT}. The
local system \(\Psi_\alpha/K_\alpha\), the balanced surviving subdomain,
the exclusively \(3\)-adic terminal, and the packaging by
graphs have the status of a \texttt{NEW\_FORMALIZATION}; they formalize without altering
the conceptual provenance of the apparatus.

\begin{ablation}[Internal falsifiers of the determinantal branch]
\label{abl:detint-falsadores}
The construction is falsified at the corresponding level if any of
the following facts occurs.
\begin{enumerate}
 \item \(\partial^2=0\) fails; then the local complex does not exist.
 \item Alexander--Whitney does not commute with refinement, the counit fails, or
 the vertex ceases to be group-like.
 \item \(z_\pm\) are not cycles, their roots differ, or
 \(\partial h_*=z_--z_+\) fails.
 \item \(\Psi_\alpha\) is not a map of complexes or
 \(K_{\beta\alpha}=K_\beta+\Psi_\beta K_\alpha\) fails; in that case
 carry or memory has been lost.
 \item The reduced complex is not balanced or is not acyclic over
 \(K_3\); the history lies outside the typed terminal domain.
 \item The Bockstein lengths do not coincide with the \(a_i\), or
 \(\Theta_{\rm Bock}\ne\overline{\lambda^{\rm or}}\).
 \item Refinement modifies \(d\) or \(\lambda^{\rm or}\) outside
 unit blocks and recorded orientation changes.
 \item One of the thirteen coordinates is omitted: \emph{object replaced;
 conclusion not transferable}.
\end{enumerate}
These forgers check the internally constructed object; they do not replace it with a subsequent control.
\end{ablation}
\fi
 \endgroup

\section{Naturalness of the ensemble and compatibility under refinement}

Before forming the horizon record, we make explicit the total action
of the five operations developed here on \emph{the same} edge
\(\mathfrak e_\alpha\). If \(P_x,Q_x\) denote the restrictions of
\eqref{eq:pdfv2-cor-sigma} to the fiber of \(x=s(\alpha)\), and analogously in
\(y=t(\alpha)\), we first set
\begin{equation}
\label{eq:pdfv2-cor-a-eta-emision}
 A_\alpha:=P_yU_\alpha P_x+Q_yU_\alpha Q_x,
 \qquad
 \eta_\alpha:=Q_yU_\alpha P_x+P_yU_\alpha Q_x.
\end{equation}
The five readings are dependent families indexed by the
\(\mathfrak e_\alpha\) itself. The index fixes source, target, route, and ledger, but
is not copied again within the value. Their types are the dependent
records
\(\operatorname{LedgerEntry}:=\operatorname{cod}(\lambda)\) and
\begin{align*}
 \mathsf Y_{\rm sp}[\alpha]&:=
 \operatorname{Rec}[U:\mathsf H_x\to\mathsf H_y;
 \sigma_x\in\operatorname{End}\mathsf H_x;
 \sigma_y\in\operatorname{End}\mathsf H_y;A,\eta:\mathsf H_x\to\mathsf H_y],\\
 \mathsf Y_{\rm sol}[\alpha]&:=
 \operatorname{Rec}[b,b^+,b^-\in\operatorname{LedgerEntry};
 \operatorname{upd}:\mathsf{Mem}_x\to\mathsf{Mem}_y],\\
 \mathsf Y_{\rm gau}[\alpha]&:=
 \mathcal P_{\rm fin}\operatorname{Rec}[
 \operatorname{Cell}(\alpha);q\in Q_{\rm inc}(y);B;W_c(q)],\\
 \mathsf Y_{\rm coh}[\alpha]&:=
 \operatorname{Rec}[I_y,J_y,\kappa_{y,N},\delta_y,C_{y,N};
 F_\alpha:C_{y,N}\to C_{x,N}],\\
 \mathsf Y_{\rm det}[\alpha]&:=
 \operatorname{Rec}[\Psi_\alpha,K_\alpha,\mathbf1_\alpha,z_\pm,h_*;
 \operatorname{Det},H_*,\operatorname{Smith},\operatorname{Boc}].
\end{align*}
Therefore, the common judgment is
\[
 \mathfrak e_\alpha:\operatorname{Ext}_k\ \vdash\
 \mathcal O_\bullet(\mathfrak e_\alpha)
 \in\mathsf Y_\bullet[\alpha],
 \qquad
 \bullet\in\{\mathrm{sp,sol,gau,coh,det}\}.
\]
\begin{align}
 \mathcal O_{\rm sp}(\mathfrak e_\alpha)
 &:=
 \bigl(U_\alpha,\sigma_x,\sigma_y,
 A_\alpha,\eta_\alpha\bigr),
 \label{eq:pdfv2-cor-accion-sp}\\
 \mathcal O_{\rm sol}(\mathfrak e_\alpha)
 &:=
 \bigl(b_\alpha,b_\alpha^+,b_\alpha^-;\,
 (M^+,M^-,D,H)\mapsto
 (M^++b_\alpha^+,M^-+b_\alpha^-,
  D+b_\alpha,H+|b_\alpha|)\bigr),
 \label{eq:pdfv2-cor-accion-sol}\\
 \mathcal O_{\rm gau}(\mathfrak e_\alpha)
 &:=
 \left\{
 \bigl(\operatorname{Cell}_{\rm TPK}(\alpha),
 q,B,W_c(q)\bigr):
 q\in Q_{\rm inc}(y)
 \right\},
 \label{eq:pdfv2-cor-accion-gau}\\
 \mathcal O_{\rm coh}(\mathfrak e_\alpha)
 &:=
 \bigl(I_y,J_y,\kappa_{y,N},\delta_y,C_{y,N},
 F_\alpha:C_{y,N}\longrightarrow C_{x,N}\bigr),
 \label{eq:pdfv2-cor-accion-coh}\\
 \mathcal O_{\rm det}(\mathfrak e_\alpha)
 &:=
 \bigl(\Psi_\alpha,K_\alpha,\mathbf1_\alpha,z_\pm(\alpha),
 h_*(\alpha),\operatorname{Det}(C_{y,N}),H_*(C_{y,N}),
 \operatorname{Smith}(C_{y,N}),\operatorname{Boc}(C_{y,N})\bigr).
 \label{eq:pdfv2-cor-accion-det}
\end{align}

The internal Archimedean reading is the first component of this simultaneous
action; it does not yet introduce any classical form:
\begin{equation}
\label{pII-full:eq:continuo-lector-arquimediano}
 \mathcal L_{\rm arq}^{\rm int}(\mathfrak e_\alpha)
 :=\mathcal O_{\rm sp}(\mathfrak e_\alpha).
\end{equation}
For every value in those families, including each element of the
multisection \(\mathcal O_{\rm gau}\), its support is the dependent index
\begin{equation}
\label{eq:pdfv2-cor-soporte-dependiente}
 \begin{aligned}
 \mathsf B_\alpha&:=
 \{(s(\alpha),t(\alpha),\gamma\alpha,
          \operatorname{Led}(\gamma\alpha))\},\\
 \operatorname{supp}_\alpha:\mathsf Y_\bullet[\alpha]&\longrightarrow
 \mathsf B_\alpha,\qquad
 o\longmapsto(s(\alpha),t(\alpha),\gamma\alpha,
          \operatorname{Led}(\gamma\alpha)).
 \end{aligned}
\end{equation}
Therefore, they remain typed, without a projection that recovers the edge,
\begin{equation}
\label{eq:pdfv2-cor-soporte-cinco-lecturas}
 \begin{gathered}
 s_\alpha(\mathcal O_\bullet):=s(\alpha),\qquad
 t_\alpha(\mathcal O_\bullet):=t(\alpha),\\
 \gamma_\alpha(\mathcal O_\bullet):=\gamma\alpha,
 \qquad
 \operatorname{Led}_\alpha(\mathcal O_\bullet)
 :=\operatorname{Led}(\gamma\alpha).
 \end{gathered}
\end{equation}
Here \(\operatorname{Cell}_{\rm TPK}(\alpha)\) is the map of complexes
\eqref{eq:pdfv2-cor-cell-accion-emision} induced by transport of the
same sheet fiber, and \(F_\alpha\) is
\eqref{eq:pdfv2-cor-refinamiento-celda} for that extension. In particular,
\(\mathcal O_{\rm gau}\) preserves the complete multisection of \(q\)'s; it does not
choose the origin or an orbit.

The five formulas above constitute a single dependent action.
\begin{equation}
\label{eq:pdfv2-cor-accion-unica-dependiente}
 \mathcal O(\mathfrak e_\alpha)
 :=\left\langle
 \mathcal O_{\rm sp},\mathcal O_{\rm sol},\mathcal O_{\rm gau},
 \mathcal O_{\rm coh},\mathcal O_{\rm det};
 \operatorname{Cpl}_\alpha
 \right\rangle,
\end{equation}
where \(\operatorname{Cpl}_\alpha\) is the certificate formed by the
identities
\begin{equation}
\label{eq:pdfv2-cor-acoplamientos-unicos}
 \begin{gathered}
 \operatorname{supp}_\alpha(\mathcal O_{\rm sp})
 =\operatorname{supp}_\alpha(\mathcal O_{\rm sol})
 =\operatorname{supp}_\alpha(\mathcal O_{\rm gau})
 =\operatorname{supp}_\alpha(\mathcal O_{\rm coh})
 =\operatorname{supp}_\alpha(\mathcal O_{\rm det}),\\
 \mathsf Q(t(\alpha))=\sigma_9\mathsf Q(s(\alpha)),\quad
 g(t(\alpha))=g(s(\alpha))+[1]_9,\quad
 q^{(9)}(t(\alpha))=q^{(9)}(s(\alpha))+\chi_9(g(s(\alpha))),\\
 A_\alpha^*A_\alpha+\eta_\alpha^*\eta_\alpha=U_\alpha^*U_\alpha,\qquad
 (M^+,M^-,D,H)'=(M^++b_\alpha^+,M^-+b_\alpha^-,
 D+b_\alpha,H+|b_\alpha|),\\
 Q_{\rm inc}=C^2(\operatorname{Cell}_{\rm TPK};\mathbf F_3),\qquad
 sB=0,\qquad W_c(q+Ba)=W_c(q),\\
 C_{y,N}\text{ is the same complex in }
 \mathcal O_{\rm coh}\text{ and }\mathcal O_{\rm det},\qquad
 \partial h_*=z_--z_+,\qquad
 K_{\beta\alpha}=K_\beta+\Psi_\beta K_\alpha.
 \end{gathered}
\end{equation}
\Needspace{18\baselineskip}
The symbol \(\bullet\) ranges over the five readings of
\eqref{eq:pdfv2-cor-accion-unica-dependiente}; all preserve the same source, destination, route, and ledger.

We collect, without separating the state, the operations constructed up to here:
\begin{equation}
\label{eq:pdfv2-cor-constructor-unico}
 \begin{aligned}
 \mathfrak D_{k,N}(\widehat x)
 &:=\bigl(\{\mathfrak e_\alpha\}_\alpha;\\
 &\quad U_k,\sigma_k,P_k,Q_k,A_k,\eta_k;\\
 &\quad b_\alpha,b_\alpha^\pm,M_k^\pm,D_k^0,H_k^0,
 \kappa_k^{\rm can};\\
 &\quad V_{\rm TPK},\mathscr I,\mathscr E,\mathscr B,B;\\
 &\quad \operatorname{Cell}_{\rm TPK}(\widehat x),Q_{\rm inc}(\widehat x),W_c;\\
 &\quad I_k,J_k,\kappa_{k,N},\delta,C_{k,N};\\
 &\quad c_\alpha,v_\alpha,\Psi_\alpha,K_\alpha,
 \mathbf1_\alpha,z_\pm(\alpha),h_*(\alpha);\\
 &\quad
 \operatorname{Det}(C_{k,N}),H_*(C_{k,N}),
 \operatorname{Smith}(C_{k,N}),\operatorname{Boc}(C_{k,N});\\
 &\quad \{\mathcal O(\mathfrak e_\alpha)\}_\alpha
 \bigr).
 \end{aligned}
\end{equation}
The semicolon merely separates levels of interpretation of the same dependent tuple; it does not define independent structures. The input state is neither retained as the first coordinate nor recovered through a projection. Each operation is proved through its action on the generators \(\mathfrak e_\alpha\), which preserve the same support, the same route, and the same ledger.

For \(\ell\geq k\) and \(N'\geq N\), transport is defined component by
component through truncation of histories, reduction of coefficients, and
precomposition of ports. The naturality assertion that must be proved is
a single one:
\begin{equation}
\label{eq:pdfv2-cor-naturalidad-unica}
 u_{(\ell,N'),(k,N)}\mathfrak D_{\ell,N'}(\widehat x_\ell)
 =\mathfrak D_{k,N}(\tau_{\ell,k}\widehat x_\ell).
\end{equation}
Naturality is formulated over the generated history, not by pretending that an
isolated fine edge must remain an edge after truncating its two
endpoints. For
\(\gamma=(\varepsilon_1,\ldots,\varepsilon_\ell)\), we define
\[
 \mathsf E_{\leq\ell}(\gamma)
 :=(\mathfrak e_{\varepsilon_i})_{1\leq i\leq\ell},
 \qquad
 \mathcal O_{\leq\ell}(\gamma)
 :=(\mathcal O(\mathfrak e_{\varepsilon_i}))_{1\leq i\leq\ell}.
\]
The two transports are the explicit prefix maps.
\begin{align}
 u^{\mathfrak e}_{\ell,k}:
 \mathsf E_{\leq\ell}(\gamma)&\longrightarrow
 \mathsf E_{\leq k}(\gamma),&
 (\mathfrak e_{\varepsilon_i})_{i\leq\ell}
 &\longmapsto(\mathfrak e_{\varepsilon_i})_{i\leq k},
 \label{eq:pdfv2-cor-transporte-arista}\\
 u^{\mathcal O}_{\ell,k}:
 \mathcal O_{\leq\ell}(\gamma)&\longrightarrow
 \mathcal O_{\leq k}(\gamma),&
 (\mathcal O(\mathfrak e_{\varepsilon_i}))_{i\leq\ell}
 &\longmapsto
 (\mathcal O(\mathfrak e_{\varepsilon_i}))_{i\leq k}.
 \label{eq:pdfv2-cor-transporte-accion}
\end{align}
In each retained entry, the component \(\mathrm{sp}\) intertwines
\(U,\sigma,P,Q,A,\eta\); the component \(\mathrm{sol}\) applies the conditional
expectation to \(b,b^\pm,\kappa\); the component \(\mathrm{gau}\) applies
\(\operatorname{Cell}(\tau)\oplus\tau_Q\); the component \(\mathrm{coh}\)
precomposes the ports with \(F\); and the component \(\mathrm{det}\) applies
\(\Psi,K,H_*,\operatorname{Smith},\operatorname{Boc}\) to the same map of
complexes. The prefix acts simultaneously on those five components.
Its components already discharged are
\begin{equation}
\label{eq:pdfv2-cor-naturalidad-componentes}
 \begin{gathered}
 U_{m,k}=U_{m,\ell}U_{\ell,k}\quad(m\geq\ell\geq k),\qquad
 E_kb_f^\pm=b_c^\pm+\kappa_k^{\rm can},\\
 r_{N',N}\kappa_{k,N'}=
 \kappa_{k,N}(r_{N',N}\oplus r_{N',N}),\\
 \Psi_\beta\Psi_\alpha=\Psi_{\beta\alpha},
 \qquad K_{\beta\alpha}=K_\beta+\Psi_\beta K_\alpha,\\
 u^{\mathcal O}_{\ell,k}\mathcal O_{\leq\ell}(\gamma)
 =\mathcal O_{\leq k}(\gamma)
 =\operatorname{map}(\mathcal O)\bigl(u^{\mathfrak e}_{\ell,k}
                    \mathsf E_{\leq\ell}(\gamma)\bigr).
 \end{gathered}
\end{equation}
The first equality uses multiplication of the conditional weights over the same
system of common cylinders; the remaining ones are the preceding explicit
formulas. No visible projection
can be used to reconstruct or select \(\widehat x\).

To complete the last component, let us denote by
\(F_{(\ell,N'),(k,N)}:C_{\ell,N'}\to C_{k,N}\) the map of complexes of
\eqref{eq:pdfv2-cor-refinamiento-celda}. Its cone is preserved as part of the
transition. Functoriality of homology and of the Bockstein sequence gives
\begin{equation}
\label{eq:pdfv2-cor-naturalidad-homologia-bockstein}
 H_*(F_{m,k})=H_*(F_{\ell,k})H_*(F_{m,\ell}),
 \qquad
 \operatorname{Boc}(F_{m,k})
 =\operatorname{Boc}(F_{\ell,k})
  \operatorname{Boc}(F_{m,\ell}).
\end{equation}
The determinant line is transported by the identity of the exact triangle.
\begin{equation}
\label{eq:pdfv2-cor-naturalidad-determinante}
 \operatorname{Det}(C_{k,N})
 \simeq
 \operatorname{Det}(C_{\ell,N'})\otimes
 \operatorname{Det}\!\left(\operatorname{Cone}
 F_{(\ell,N'),(k,N)}\right).
\end{equation}
This formula also records the refinements that add homology. The
Smith reading is conserved without feigning invariance: each transition is
assigned the total certificate
\begin{equation}
\label{eq:pdfv2-cor-transicion-smith}
 \mathcal S(F):=
 \bigl(\operatorname{Smith}(C_{\ell,N'}),
       \operatorname{Smith}(C_{k,N}),
       \operatorname{Smith}(\operatorname{Cone}F)\bigr),
\end{equation}
which literally determines which summands are conserved, born, or die.

The action of a coherent symmetry \(g\) on the common tree induces
permutation isometries \(\widetilde g_k\) on
\(\mathscr H_k\). The bijection of children and preservation of the weights of
\eqref{eq:continuo-comun-naturalidad-medida} prove, on the basis vectors,
\begin{equation}
\label{eq:pdfv2-cor-naturalidad-transporte-hojas}
 \widetilde g_{k+1}U_k=U_k\widetilde g_k,
 \qquad
 \widetilde g_k\sigma_k=\sigma_k\widetilde g_k,
 \qquad
 \widetilde g_kP_k=P_k\widetilde g_k,
 \qquad
 \widetilde g_kQ_k=Q_k\widetilde g_k.
\end{equation}

\begin{theorem}[Totality and naturality of the correlative constructor]
\label{thm:pdfv2-cor-totalidad-naturalidad}
For every total state of \eqref{eq:continuo-comun-estado}, every finite horizon,
and every coefficient level, the expression
\eqref{eq:pdfv2-cor-constructor-unico} is defined. Its five
characteristics are simultaneous operations on the same extension ledger and
satisfy the single identity \eqref{eq:pdfv2-cor-naturalidad-unica}. The collar, the complex, and the determinant
line remain nonempty even when their subsequent readers are neither acyclic nor
invertible.
\end{theorem}

\begin{proof}
Nonemptiness of the tree and its ports is
\cref{prop:continuo-comun-arbol-finito-no-vacio}; the measure and the expectation
are total by \cref{thm:continuo-comun-medida-canonica}. The formulas
  \eqref{eq:pdfv2-cor-u-comun}--\eqref{eq:pdfv2-cor-memorias} define transport,
polarity, balance, and memory for all generators. The four
  directions of \eqref{eq:pdfv2-cor-cuatro-rectas-tpk} are distinct by
  \eqref{eq:pdfv2-cor-proyectores-no-colapso}; the cellular completion
  \cref{def:pdfv2-cor-completacion-celular-tpk} contains its two path orders and
  six cells, while
  \eqref{eq:pdfv2-cor-collar-naturalidad} proves its naturality. The
cellular sequence is exact by \eqref{eq:pdfv2-cor-sucesion-celda}, the
local complex satisfies \(\partial^2=0\), and the correlated filling is explicit by
\eqref{eq:pdfv2-cor-filler}. Every finite based construction possesses the line
\eqref{eq:pdfv2-cor-linea-determinante-total}; homology, Bockstein, and the
Smith certificate are transported by
\eqref{eq:pdfv2-cor-naturalidad-homologia-bockstein}--
\eqref{eq:pdfv2-cor-transicion-smith}. Finally, the identities are
verified on each generator \(\mathfrak e_\alpha\) through
\eqref{eq:pdfv2-cor-naturalidad-componentes},
\eqref{eq:pdfv2-cor-collar-naturalidad}, and
\eqref{eq:pdfv2-cor-naturalidad-transporte-hojas}; by linearity and by the
functorial composition of routes,
\eqref{eq:pdfv2-cor-naturalidad-unica} is obtained for the integral record.
\end{proof}

\section{Causal separation of the identities, deductions, and realizations of the correlator constructor}

\paragraph{No vacuity and exact identities of the correlator constructor}
The nonemptiness imported from \eqref{eq:continuo-comun-out-finito} and the
identities
\eqref{eq:pdfv2-cor-b-identidades},
\eqref{eq:pdfv2-cor-sucesion-celda},
\eqref{eq:pdfv2-cor-sb-cero},
\eqref{eq:pdfv2-cor-u-gram},
\eqref{eq:pdfv2-cor-balance-hojas},
\eqref{eq:pdfv2-cor-jensen-hojas},
\eqref{eq:pdfv2-cor-dos-cero},
\eqref{eq:pdfv2-cor-cociclo-matricial}, and
\eqref{eq:pdfv2-cor-filler} follow from the definitions and the exhibited
APP--TRIT--TPK relations. Their status is
\texttt{EXACT\_INTERNAL}.

\paragraph{Starting hypotheses for the isometry of \(U_k\), the cellular sequence, and the composition \(\Psi/K\)}
The construction of \(U_k\) uses as its starting structure the finiteness of the fibers, survival at every horizon, functional truncation, and the transports of the two sheets already produced by TPK\@. The equality \(U_k^*U_k=I\) additionally requires verification of \(U_\alpha^*U_\alpha=I\) for every emission; without that verification, the Gram identity \eqref{eq:pdfv2-cor-u-gram}, which is unconditional, remains. The cellular succession uses the two ordered APP ports. The composition \(\Psi/K\) uses integer carry, orientation, and boundary. Under those enumerated premises, the stated conclusions are affirmative.

\paragraph{Additional data for isometric realizations, balanced retracts, and trivialization of the determinant line}
The internal constructor is completed by the positive Gram form of
\eqref{eq:pdfv2-cor-u-gram}, the total space \(Q_{\rm inc}\), the cellular sequence,
and the determinant line with its transition cone. A specialized
publication may require additional data: the isometry of each
\(U_\alpha\), a concrete numerical identification of the APP defect, a
balanced acyclic retraction, or trivialization of a determinant
line. Those data are intertwiners of the published codomain; they do not
form part of the domain or of the existence of
\(\mathfrak D_{k,N}\). Their absence precludes exclusively the promotion that
mentions them.

\begin{criterion}[Falsifiers of the intrinsic correlation]
\label{crit:pdfv2-cor-falsadores}
The joint construction fails at the first level at which any of the
following facts occurs:
\begin{enumerate}
 \item an operation uses an extension, orientation, carry, memory, boundary or route different from that recorded in \(\mathfrak e_\alpha\);
 \item a child of
 \(\operatorname{Ch}_k(\widehat x_k)\) is chosen and the remaining ones are eliminated;
 \item the weight of \eqref{eq:pdfv2-cor-peso-conteo} does not sum to one,
 \eqref{eq:pdfv2-cor-u-gram} fails, or \(U_k^*U_k=I\) is declared without proving
 \(U_\alpha^*U_\alpha=I\) for every emission;
 \item \(\kappa_k^{\rm can}\) is identified with the integer carry
 \(\kappa_\alpha\);
 \item \(sB=0\) fails, \(Q_{\rm inc}(\widehat x)\) is contracted
 to its origin, the action \(q\mapsto q+Ba\) is altered, or the two path orders are identified before
 applying the reader;
 \item a refinement breaks
 \eqref{eq:pdfv2-cor-refinamiento-cadena} or the cocycle
 \eqref{eq:pdfv2-cor-cociclo-matricial};
 \item fails \(\partial^2=0\), \(z_\pm\) are not ciclos or
 \(\partial h_*\neq z_--z_+\);
 \item a nonzero determinantal scalar is asserted without the retract, the
 balance, and the acyclicity required by that publication;
 \item any of the five operations is presented as an object
 separate from the constructor \(\mathfrak D_{k,N}\).
\end{enumerate}
In all those cases the diagnosis is: \emph{common object substituted;
conclusion not transferable}.
\end{criterion}
 \end{HMTScientificOwnerSubordinated}
 \label{chap:83-05}

\begin{HMTScientificOwnerSubordinated}
\chapter[Terminal, naturality, and limit]{Multivalued terminal, unique naturality and passage to the limit}
\label{chap:pdfv2-terminal-naturalidad-limite}

\begin{thesisbox}
The discrete continuum produces a single system of histories, extensions, memories, and terminals. The internal discriminants \(\mathrm{sp},\mathrm{sol},\mathrm{gau},\mathrm{coh},\mathrm{det}\) name five characteristics that are simultaneous within that same system: they neither fix distinct inputs nor authorize the selection of different histories. In this chapter, the terminal preserves all continuations, the connection remains a nonadic correspondence, and naturality is expressed by a single identity.
\end{thesisbox}

\section{Surviving histories with extension witnesses}

Let
\(\widehat x_k\in\widehat{\mathcal X}^{\rm tot}_k\) be the total state of
\eqref{eq:continuo-comun-dominio-total}. For \(N\geq k\), we define without
any other predicate of membership
\begin{equation}
\label{eq:pdfv2-terminal-historias-finita}
 \mathscr H_{k,N}(\widehat x_k)
 :=\operatorname{Term}_{k,N}(\widehat x_k).
\end{equation}
Each
\(h=(\varepsilon_{k+1},\ldots,\varepsilon_N)\) on the right contains through
\eqref{eq:continuo-comun-vertice-arbol} its intermediate states, the witnesses
of \(\operatorname{Ext}\), the carries, boundaries, routes, and ledger.
Each edge satisfies the phase action
\eqref{eq:continuo-comun-accion-fase-tpk}; hence every nine-step subpath
determines, with its nine emissions and not merely with its length, an
element of \(\mathsf Z^\Gamma_{9,j}\) and a witness of
\(\Gamma_{9,j}\) through
\eqref{eq:continuo-comun-gamma9-apice}--
\eqref{eq:continuo-comun-gamma9-relacion}. Therefore, Ext and nonadic holonomy
arise from the same phase-compatible word. To \eqref{eq:pdfv2-terminal-historias-finita} there is adjoined neither a survival filter nor a
second condition \(\Gamma_9\).

For \(k\leq\ell\leq N\), let
\begin{equation}
\label{eq:pdfv2-terminal-prefijos-intermedios}
 \mathscr P_{\ell\mid k}(\widehat x_k)
 :=\{h_{[k,\ell]}:h\in\mathscr H_{k,N}(\widehat x_k)
       \text{ for some }N\geq\ell\}.
\end{equation}
For \(p\in\mathscr P_{\ell\mid k}(\widehat x_k)\), we use the abbreviation
\begin{equation}
\label{eq:pdfv2-terminal-continuaciones-prefijo}
 \mathscr H_{\ell,N}(p)
 :=\{q:p\star q\in\mathscr H_{k,N}(\widehat x_k)\}.
\end{equation}
The prolongation neutral of \eqref{eq:continuo-comun-seccion-total} extends
each prefix to cualquier horizon. Consequently,
\begin{equation}
\label{eq:pdfv2-terminal-no-vacuidad-historias}
 \mathscr H_{k,N}(\widehat x_k)\neq\varnothing,
 \qquad
 \mathscr P_{\ell\mid k}(\widehat x_k)\neq\varnothing,
\end{equation}
by \cref{prop:continuo-comun-arbol-finito-no-vacio}.

The truncation
\begin{equation}
\label{eq:pdfv2-terminal-truncamiento-historias}
 \pi_{N+1,N}^{\mathscr H}:
 \mathscr H_{k,N+1}(\widehat x_k)
 \longrightarrow\mathscr H_{k,N}(\widehat x_k)
\end{equation}
is exactly \eqref{eq:continuo-comun-proyeccion-terminal} and is
surjective by \cref{lem:continuo-comun-terminal-sobreyectivo}. Survival is, by
\eqref{eq:continuo-comun-surv-igual-total}, a demonstrated output of the total
domain and not a prior selection.

\section{Generated terminal record and its \(13\) coordinates}

Let us fix a history
\[
 h=(\varepsilon_{j+1}:x_j\longrightarrow x_{j+1})_{j=k}^{N-1},
 \qquad
 \gamma_{k,k}:=I_{x_k},\qquad
 \gamma_{k,j+1}:=\gamma_{k,j}\star\varepsilon_{j+1},
\]
The operator notation
\(U_{\gamma_{k,j+1}}
=U_{\varepsilon_{j+1}}U_{\gamma_{k,j}}\) preserves the order of application,
whereas \(\star\) preserves the chronological order of the word. We further fix
a coefficient level \(r\geq1\). For brevity, we write
\[
 C_{x_j,r}:=C_{j,r}(x_j),\qquad
 F_{\varepsilon_{j+1}}:C_{x_{j+1},r}\longrightarrow C_{x_j,r}
\]
for the complex and the map of complexes of
\eqref{eq:pdfv2-cor-cono-celda} and
\eqref{eq:pdfv2-cor-refinamiento-celda}. The type of the level record
\(j\), denoted \(\mathfrak R_{j,r}\), is a dependent record with the
following thirteen fields. In each field we indicate the type, the seed and the
action of an edge \(\varepsilon:=\varepsilon_{j+1}\).

The state \(x_k\) already contains the complete word
\(\gamma_{0,k}=\varepsilon_1\star\cdots\star\varepsilon_k\). Therefore a
``seed in \(k\)'' never restarts the genealogy: it is the fold of the same
updates from the root to \(k\). We denote
\begin{equation}
\label{eq:pdfv2-terminal-prefijo-heredado}
 \begin{gathered}
 M_k^\pm:=\sum_{i=1}^{k}b_{\varepsilon_i}^\pm,qquad
 D_k:=\sum_{i=1}^{k}b_{\varepsilon_i},qquad
 H_k:=\sum_{i=1}^{k}|b_{\varepsilon_i}|,\\
 \mathbf c_{\leq k}:=
  \bigl(c_r(\varepsilon_i,\gamma_{0,i-1})\bigr)_{1\leq i\leq k},
 \qquad
 \mathbf k_{\leq k}:=
  \bigl(\kappa_{i-1}^{\rm can}\bigr)_{1\leq i\leq k},\\
 \mathbf F_{\leq k}:=
  \bigl(F_{\varepsilon_i}:C_{x_i,r}\to C_{x_{i-1},r}\bigr)_{1\leq i\leq k}.
 \end{gathered}
\end{equation}
All empty lists appearing below are understood literally only
in \(k=0\); for \(k>0\) they are replaced by prefixes inherited from
\eqref{eq:pdfv2-terminal-prefijo-heredado} and from the ledger of
\(\gamma_{0,k}\). Thus the record does not preserve an opaque copy of \(x_k\),
but neither does it lose its APP--TRIT--TPK past.

\paragraph{Fiber \(\mathsf F\): local domain and conservation under prolongation}
The elementary fiber is
\begin{equation}
\label{eq:pdfv2-terminal-tipo-fibra}
 \operatorname{Fib}_r(x):=
 \bigl(\mathsf H_x^\Sigma\oplus\mathsf H_x^\Pi,\,
       \mathsf M_x,\,
       \operatorname{Cell}_{\rm TPK}(x),\,C_{x,r}\bigr).
\end{equation}
The field \(\mathsf F_{k,j}\) is a word in the category of these fibers.
Its seed and update are
\begin{align}
 \mathsf F_{k,k}
 &:=\bigl(\operatorname{Fib}_r(x_0),I\bigr)\star
 \mathop{\star}_{i=1}^{k}
 \bigl(\operatorname{Fib}_r(x_i),U_{\varepsilon_i},
 \mathsf C_{\varepsilon_i},\operatorname{Cell}_{\rm TPK}(\varepsilon_i),
 F_{\varepsilon_i}\bigr),
 \label{eq:pdfv2-terminal-semilla-f}\\
 \operatorname{Upd}^{\mathsf F}_{\varepsilon}\mathsf F_{k,j}
 &:=
 \mathsf F_{k,j}\star
 \bigl(\operatorname{Fib}_r(x_{j+1}),U_\varepsilon,
 \mathsf C_\varepsilon,\operatorname{Cell}_{\rm TPK}(\varepsilon),
 F_\varepsilon\bigr).
 \label{eq:pdfv2-terminal-update-f}
\end{align}
It is nonempty because the two sheets, the memory, the cellular unit
\eqref{eq:pdfv2-cor-cell-unidad}, and the port complex have already been
constructed. The composition of its five maps proves naturality.

\paragraph{Extensions \(\mathsf E\): morphisms of compatible prolongation of the state}
The codomain is the category of typed relation words. We set
\begin{align}
 \mathsf E_{k,k}
 &:=(I_{x_0},\operatorname{Ext}_{I_{x_0}})\star
 \mathop{\star}_{i=1}^{k}
 \bigl(\varepsilon_i,\operatorname{Ext}_{\varepsilon_i},
 \operatorname{Ext}_{\varepsilon_i}\circ
 \operatorname{Ext}_{\gamma_{0,i-1}}
 \subseteq\operatorname{Ext}_{\gamma_{0,i}}\bigr),
 \label{eq:pdfv2-terminal-semilla-e}\\
 \operatorname{Upd}^{\mathsf E}_{\varepsilon}\mathsf E_{k,j}
 &:=
 \mathsf E_{k,j}\star
 \bigl(\varepsilon,\operatorname{Ext}_{\varepsilon},
 \operatorname{Ext}_{\varepsilon}\circ
 \operatorname{Ext}_{\gamma_{0,j}}
 \subseteq\operatorname{Ext}_{\gamma_{0,j+1}}\bigr).
 \label{eq:pdfv2-terminal-update-e}
\end{align}
Identity and neutral emission prove nonemptiness; associativity of the
composition relational proof naturality.

\paragraph{Relations \(\mathsf R\): incidence, orientation, and compatibility of composition}
Let \(\operatorname{RelBlock}_{j,r}\) be the type of blocks of equations
labeled by the edge and the prefix over which they are verified, and let
\(\operatorname{Pres}_{j,r}:=\operatorname{List}
(\operatorname{RelBlock}_{j,r})\). The unit block and the seed are
\begin{equation}
\label{eq:pdfv2-terminal-semilla-r}
 \mathsf B_I:=(I_{x_0};U_I=I,\ \mathsf U_I=I,\ \kappa(I)=0,\ \partial I=0),
 \qquad
 \mathsf R_{k,k}:=(\mathsf B_I)\star
 \mathop{\star}_{i=1}^{k}
 \mathsf B_R(\varepsilon_i;\gamma_{0,i-1}).
\end{equation}
For an edge \(\varepsilon:x_j\to x_{j+1}\), the new block is the
labeled tuple
\begin{equation}
\label{eq:pdfv2-terminal-bloque-r}
 \begin{split}
 \mathsf B_R(\varepsilon;\gamma_{0,j}):=\bigl(
 &\varepsilon,\gamma_{0,j};\ 
 \delta^\diamond(\varepsilon)
 =\operatorname{res}^\diamond(\varepsilon)
  +9\operatorname{quo}^\diamond(\varepsilon),\\
 &\Delta(\varepsilon)=\operatorname{or}(\varepsilon)
  +3\kappa(\varepsilon),\quad
 U_{\gamma_{0,j+1}}=U_\varepsilon U_{\gamma_{0,j}},\\
 &\mathsf U_{\gamma_{0,j+1}}
 =\mathsf U_\varepsilon\mathsf U_{\gamma_{0,j}},\quad
 \kappa(\gamma_{0,j+1})
 =\kappa(\varepsilon)+\mathsf C_\varepsilon\kappa(\gamma_{0,j}),\\
 &\partial\gamma_{0,j+1}
 =\partial\varepsilon+\partial\gamma_{0,j},\quad
 \operatorname{Cpl}_\varepsilon
 \text{ of }\eqref{eq:pdfv2-cor-acoplamientos-unicos}
 \bigr).
 \end{split}
\end{equation}
The update concatenates the block without losing date or multiplicity:
\begin{equation}
\label{eq:pdfv2-terminal-update-r}
 \operatorname{Upd}^{\mathsf R}_{\varepsilon}\mathsf R_{k,j}
 :=\mathsf R_{k,j}\star\mathsf B_R(\varepsilon;\gamma_{0,j}).
\end{equation}
The truncation \(\operatorname{tr}^{\mathsf R}_{j+1,j}\) removes the final
labeled block. Moreover,
\(a_*\mathsf B_R(\varepsilon;\gamma)
=\mathsf B_R(a\varepsilon;a\gamma)\), because each equation is natural over
the same edge and the same prefix.

\paragraph{Numbers \(\mathsf N\): publication of value with conservation of genealogy}
For an edge \(\varepsilon:x_j\to x_{j+1}\), we define the numerical datum
\begin{equation}
\label{eq:pdfv2-terminal-numero-arista}
 \nu(\varepsilon):=
 \bigl(j+1,\mathsf Q_{j+1},g_{j+1},q^{(9)}_{j+1},\beta_{j+1},
 \operatorname{res}^{\Sigma},\operatorname{quo}^{\Sigma},
 \operatorname{res}^{\Pi},\operatorname{quo}^{\Pi},
 \operatorname{or},\kappa,\Delta\bigr)(\varepsilon).
\end{equation}
Its type is
\[
 \mathbb N\times\widetilde{\mathcal O}_9\times C_9\times\mathbb N
 \times\mathbb Z/12\mathbb Z
 \times\{0,\ldots,8\}^2\times\mathbb Z^5.
\]
The seed contains the complete numerical list of the prefix \(0\to k\);
the update adds exactly the datum of the next edge:
\begin{equation}
\label{eq:pdfv2-terminal-update-n}
 \mathsf N_{k,k}:=
 ((0,\mathsf Q_0,g_0,q^{(9)}_0,\beta_0;0,\ldots,0))
 \star\mathop{\star}_{i=1}^{k}\nu(\varepsilon_i),\qquad
 \operatorname{Upd}^{\mathsf N}_{\varepsilon}\mathsf N_{k,j}
 :=\mathsf N_{k,j}\star\nu(\varepsilon).
\end{equation}
Totality proceeds from the divisions APP, the division TRIT, and the odometer TPK.

\paragraph{Orientation \(\mathsf O\): sign transport and leaf selection}
The type of an orientation state is
\[
 \mathbb Z\times
 (\mathsf H^\Sigma_x\oplus\mathsf H^\Pi_x)
 \times\mathbb Z[\mathsf A],
\]
with seed and update
\begin{align}
 \mathsf O_{k,k}&:=
 \left(\sum_{i=1}^{k}\operatorname{or}(\varepsilon_i),
       \operatorname{Sh}(x_k),\partial\gamma_{0,k}\right),
 \label{eq:pdfv2-terminal-semilla-o}\\
 \operatorname{Upd}^{\mathsf O}_{\varepsilon}
 (\mathfrak o,\operatorname{Sh},\mathfrak b)
 &:=
 \bigl(\mathfrak o+\operatorname{or}(\varepsilon),
 U_\varepsilon\operatorname{Sh},
 \mathfrak b+\partial\varepsilon\bigr).
 \label{eq:pdfv2-terminal-update-o}
\end{align}
The APP leaves prove nonemptiness. Functoriality of \(U\), additivity of
\(\operatorname{or}\), and telescoping of the boundary prove naturality.

\paragraph{Carry \(\mathsf K\): update of the quotient and residual memory}
We use the typed integer lifts in
\eqref{eq:pdfv2-cor-carry-matricial-def}. The field has type
\[
 \mathsf M_{x_j}\times
 \operatorname{Hom}(\mathsf M_{x_0},\mathsf M_{x_j})
 \times\operatorname{Mat}_5(\mathbb Z)
 \times\operatorname{List}(\operatorname{Mat}_5(\mathbb Z))
 \times\operatorname{List}(\mathbb Q_{\geq0}).
\]
For \(\gamma_{k,j}=\varepsilon_j\cdots\varepsilon_{k+1}\), we explicitly
fix
\begin{equation}
\label{eq:pdfv2-terminal-lift-entero-ruta}
 \widetilde L_{\gamma_{k,k}}:=I_5,\qquad
 \widetilde L_{\gamma_{k,j}}
 :=\widetilde L_r(\varepsilon_j)\cdots
   \widetilde L_r(\varepsilon_{k+1}).
\end{equation}
For the inherited prefix we also fix
\[
 \widetilde L_{\gamma_{0,k}}
 :=\widetilde L_r(\varepsilon_k)\cdots
   \widetilde L_r(\varepsilon_1).
\]
The seed and update are
\begin{align}
 \mathsf K_{k,k}&:=
 \bigl(\kappa(\gamma_{0,k}),\mathsf C_{\gamma_{0,k}},
       \widetilde L_{\gamma_{0,k}},
       \mathbf c_{\leq k},\mathbf k_{\leq k}\bigr),
 \label{eq:pdfv2-terminal-semilla-k}\\
 \operatorname{Upd}^{\mathsf K}_{\varepsilon}
 (\kappa_\gamma,\mathsf C_\gamma,\widetilde L_\gamma,
   \mathbf c,\mathbf k)
 &:=
 \bigl(\kappa(\varepsilon)+\mathsf C_\varepsilon\kappa_\gamma,\,
 \mathsf C_\varepsilon\mathsf C_\gamma,\,
 \widetilde L_r(\varepsilon)\widetilde L_\gamma,\,
 \mathbf c\star c_r(\varepsilon,\gamma),\,
 \mathbf k\star\kappa_j^{\rm can}\bigr).
 \label{eq:pdfv2-terminal-update-k}
\end{align}
The identity provides the seed; the two cocycles
\eqref{eq:continuo-comun-cociclo-carry} and
\eqref{eq:pdfv2-cor-cociclo-matricial} prove independence of
parenthesization and naturality.

\paragraph{Memory \(\mathsf M\): record of the traversal and compatibility under truncation}
Its type at the same time preserves affine memory, ledger, and signed memories:
\begin{equation}
\label{eq:pdfv2-terminal-tipo-ledger-entry}
 \operatorname{LedgerEntry}:=
 \{\lambda(\varepsilon):\varepsilon\text{ typed elementary emission}\}.
\end{equation}
\[
 \mathsf M_x\times\operatorname{List}(\operatorname{LedgerEntry})
 \times\mathbb N^2\times\mathbb Z\times\mathbb N.
\]
We set
\begin{align}
 \mathsf M_{k,k}&:=
 \bigl(m_k,\operatorname{Led}(\gamma_{0,k}),
       M_k^+,M_k^-,D_k,H_k\bigr),
 \label{eq:pdfv2-terminal-semilla-m}\\
 \operatorname{Upd}^{\mathsf M}_{\varepsilon}
 (m,\mathbf{Led},M^+,M^-,D,H)
 &:=
 \bigl(\mathsf C_\varepsilon m+\kappa(\varepsilon),\,
 \mathbf{Led}\star\lambda(\varepsilon),\,
 M^++b_\varepsilon^+,\,M^-+b_\varepsilon^-,\,
 D+b_\varepsilon,\,H+|b_\varepsilon|\bigr).
 \label{eq:pdfv2-terminal-update-m}
\end{align}
The affine action and \eqref{eq:pdfv2-cor-memorias} prove totality and
naturality; the past is not recomputed from the last value.

\paragraph{Boundary \(\mathsf B\): gluing data and survival of the extension}
The type is a boundary chain accompanied by a word of complejos
based and chain maps. Its seed and update are
\[
 \mathsf B_{k,j}\in
 \operatorname{Rec}[\mathfrak b\in\mathbb Z[\mathsf A_j];
 C_{x_j,r}\in\operatorname{Ch}_{\rm bas};
 \mathbf F\in\operatorname{Path}_{\operatorname{Ch}_{\rm bas}}
 (C_{x_j,r},C_{x_0,r})].
\]
\begin{align}
 \mathsf B_{k,k}&:=
 \bigl(\partial\gamma_{0,k},C_{x_k,r},\mathbf F_{\leq k}\bigr),
 \label{eq:pdfv2-terminal-semilla-b}\\
 \operatorname{Upd}^{\mathsf B}_{\varepsilon}
 (\mathfrak b,C_{x_j,r},\mathbf F)
 &:=
 \bigl(\mathfrak b+\partial\varepsilon,\,
 C_{x_{j+1},r},\,\mathbf F\star F_\varepsilon\bigr).
 \label{eq:pdfv2-terminal-update-b}
\end{align}
Non-emptiness proceeds from the common ports; \(\partial^2=0\),
\(F_{\beta\varepsilon}=F_\varepsilon F_\beta\), and
\eqref{eq:continuo-comun-frontera-telescopica} prove compatibility.

\paragraph{Route \(\mathsf G\): observable history of the surviving extension}
The field is the list of prefixes in the path category:
\begin{equation}
\label{eq:pdfv2-terminal-update-g}
 \mathsf G_{k,k}:=(\gamma_{0,i})_{0\leq i\leq k},\qquad
 \operatorname{Upd}^{\mathsf G}_{\varepsilon}\mathsf G_{k,j}
 :=\mathsf G_{k,j}\star\gamma_{0,j+1}.
\end{equation}
Identity, associative concatenation, and
\(a(\varepsilon\gamma)=a(\varepsilon)a(\gamma)\) prove the three
required properties.

\paragraph{Multisection \(\mathsf X\): stable finite fiber and obstruction to pointwise selection}
For \(M\geq j\), we define by action
\begin{equation}
\label{eq:pdfv2-terminal-multiseccion-interna}
 \mathsf X_{k,j}(M):=
 \{\gamma_{k,j}\star q:
 q\in\operatorname{Term}_{j,M}(x_j)\}.
\end{equation}
Thus
\[
 \mathsf X_{k,k}(M)=\operatorname{Term}_{k,M}(x_k),
\]
and
\begin{equation}
\label{eq:pdfv2-terminal-update-x}
 \operatorname{Upd}^{\mathsf X}_{\varepsilon}\mathsf X_{k,j}(M)
 :=
 \{\gamma_{k,j}\star\varepsilon\star q:
 q\in\operatorname{Term}_{j+1,M}(x_{j+1})\}.
\end{equation}
The neutral prolongation proves nonemptiness. Deleting the final step and acting
by a symmetry commute literally with concatenation.

\paragraph{Survival \(\mathsf S\): persistence under refinement and passage to the limit}
Here, a proposition is not stored, but witnesses. Let
\begin{equation}
\label{eq:pdfv2-terminal-testigo-neutro}
 s_{j,j}:=I_{x_j},\qquad
 s_{j,M+1}:=s_{j,M}\star
 \varepsilon_M^0(t(s_{j,M})).
\end{equation}
The seed is
\(\mathsf S_{k,k}:=(s_{k,M})_{M\geq k}\), and
\begin{equation}
\label{eq:pdfv2-terminal-update-s}
 \operatorname{Upd}^{\mathsf S}_{\varepsilon}\mathsf S_{k,j}
 :=(\gamma_{k,j}\star\varepsilon\star s_{j+1,M})_{M\geq j+1}.
\end{equation}
By construction
\(\pi_{M+1,M}s_{j,M+1}=s_{j,M}\); the naturality of neutral emission
proves equivariance.

\paragraph{Terminal \(\mathsf T\): universality of the limit object and factorization of readers}
The terminal is the inverse system, not a selected queue:
\begin{equation}
\label{eq:pdfv2-terminal-sistema-inverso-campo}
 \mathsf T(x_j):=
 \bigl((\operatorname{Term}_{j,M}(x_j))_{M\geq j},
       (\pi_{M+1,M})_{M\geq j}\bigr).
\end{equation}
For each \(M\geq j\), the map of prefix tipado is
\begin{equation}
\label{eq:pdfv2-terminal-update-t}
 \operatorname{pref}_{\gamma_{k,j},M}:
 \operatorname{Term}_{j,M}(x_j)\longrightarrow
 \operatorname{Term}_{k,M}(x_k),\qquad
 q\longmapsto\gamma_{k,j}\star q.
\end{equation}
We define the field state and its update by
\begin{align}
 \mathsf T_{k,j}
 &:=
 \left(\mathsf T(x_j),
   (\operatorname{pref}_{\gamma_{k,j},M})_{M\geq j}\right),
 \label{eq:pdfv2-terminal-estado-t}\\
 \operatorname{Upd}^{\mathsf T}_{\varepsilon}\mathsf T_{k,j}
 &:=
 \left(\mathsf T(x_{j+1}),
   (\operatorname{pref}_{\gamma_{k,j+1},M})_{M\geq j+1}\right),
 \qquad \gamma_{k,j+1}:=\gamma_{k,j}\star\varepsilon.
 \label{eq:pdfv2-terminal-update-t-explicita}
\end{align}
In particular, \(\mathsf T_{k,k}\) uses
\(\gamma_{k,k}=I_{x_k}\). The witnesses
\eqref{eq:pdfv2-terminal-testigo-neutro} and the surjectivity of \(\pi\)
prove that the system is nonempty; level by level,
\begin{equation}
\label{eq:pdfv2-terminal-prefijo-natural}
 \pi^{\mathscr H}_{M+1,M}
 \operatorname{pref}_{\gamma_{k,j},M+1}
 =\operatorname{pref}_{\gamma_{k,j},M}
  \pi^{\mathscr H}_{M+1,M},
\end{equation}
because both traversals only remove the last emission of \(q\).

\paragraph{13. Terminal reader \(\mathsf L\): evaluation of the compatible extension and conservation of genealogy}
The reader contains no classical object. Its type is the category of
paths of the unique dependent constructor
\(\mathfrak D_{j,r}\) of \eqref{eq:pdfv2-cor-constructor-unico}. We set
\begin{equation}
\label{eq:pdfv2-terminal-transicion-lector-tipado}
 u^{\mathfrak D}_{j+1,j}
 :=u_{(j+1,r),(j,r)}:
 \mathfrak D_{j+1,r}(x_{j+1})
 \longrightarrow\mathfrak D_{j,r}(x_j),
\end{equation}
the component of the common transport
\eqref{eq:pdfv2-cor-naturalidad-unica} that truncates the last emission, reduces
the coefficients and precomposes the ports. By composition,
\begin{equation}
\label{eq:pdfv2-terminal-transicion-lector-composicion}
 u^{\mathfrak D}_{m,k}
 =u^{\mathfrak D}_{\ell,k}u^{\mathfrak D}_{m,\ell}
 \qquad(m\geq\ell\geq k).
\end{equation}
\begin{align}
 \mathsf L_{k,k}
 &:=(\mathfrak D_{0,r}(x_0))\star
 \mathop{\star}_{i=1}^{k}
 \bigl(\mathfrak D_{i,r}(x_i),u^{\mathfrak D}_{i,i-1},
       \mathcal O(\mathfrak e_{\varepsilon_i})\bigr),
 \label{eq:pdfv2-terminal-semilla-l}\\
 \operatorname{Upd}^{\mathsf L}_{\varepsilon}\mathsf L_{k,j}
 &:=
 \mathsf L_{k,j}\star
 \bigl(\mathfrak D_{j+1,r}(x_{j+1}),
 u^{\mathfrak D}_{j+1,j},\mathcal O(\mathfrak e_\varepsilon)\bigr).
 \label{eq:pdfv2-terminal-update-l}
\end{align}
The totality of \(\mathfrak D\) and
\eqref{eq:pdfv2-terminal-transicion-lector-composicion} proves non-emptiness
and naturality.

The thirteen preceding seeds form a single element.
\begin{equation}
\label{eq:pdfv2-terminal-semilla-trece}
 \mathfrak r^0_{k,r}:=
 \langle
 \mathsf F_{k,k};\mathsf E_{k,k};\mathsf R_{k,k};\mathsf N_{k,k};
 \mathsf O_{k,k};\mathsf K_{k,k};\mathsf M_{k,k};\mathsf B_{k,k};
 \mathsf G_{k,k};\mathsf X_{k,k};\mathsf S_{k,k};\mathsf T_{k,k};
 \mathsf L_{k,k}
 \rangle\in\mathfrak R_{k,r}.
\end{equation}
The action of an edge is the total application.
\begin{equation}
\label{eq:pdfv2-terminal-actualizacion-unica}
 \operatorname{Upd}_\varepsilon:
 \mathfrak R_{j,r}(\gamma_{0,j})
 \longrightarrow\mathfrak R_{j+1,r}(\gamma_{0,j+1}),\qquad
 \operatorname{Upd}_\varepsilon:=
 \langle\operatorname{Upd}^{\mathsf F}_\varepsilon,\ldots,
 \operatorname{Upd}^{\mathsf L}_\varepsilon\rangle,
\end{equation}
with the thirteen actions
\eqref{eq:pdfv2-terminal-update-f}--
\eqref{eq:pdfv2-terminal-update-l}. Finally, the terminal register is
\emph{generated} by the fold
\begin{equation}
\label{eq:pdfv2-terminal-registro-trece}
 \boxed{
 \mathfrak r_{k,N}(h):=
 \operatorname{Upd}_{\varepsilon_N}\circ\cdots\circ
 \operatorname{Upd}_{\varepsilon_{k+1}}(\mathfrak r^0_{k,r}).}
\end{equation}
It is not the unpacking of an intact copy of the state.
On the image of this fold, the truncation is defined componentwise
\begin{equation}
\label{eq:pdfv2-terminal-truncamiento-trece-def}
 \operatorname{tr}_{13;j+1,j}
 \mathfrak r_{k,j+1}(h\star\varepsilon)
 :=\mathfrak r_{k,j}(h).
\end{equation}
It is well defined: the fields \(\mathsf E\) and \(\mathsf G\) preserve the
ordered list of emissions and prefixes, so that equality of two
records entails equality of their labeled histories. In
\(\mathsf F,\mathsf E,\mathsf R,\mathsf N,\mathsf G,\mathsf T,\mathsf L\),
the last entry or block is deleted; in \(\mathsf O,\mathsf K,\mathsf M,
\mathsf B\), the cumulative formula is applied to the recovered prefix; and in
\(\mathsf X,\mathsf S\), the cone of tails is restricted. Neither an
idempotent union is inverted nor a continuation selected.

\begin{proposition}[Totality and naturalness of the thirteen updates]
\label{prop:pdfv2-terminal-trece-total-natural}
All the seeds of \eqref{eq:pdfv2-terminal-semilla-trece} exist, each
\(\operatorname{Upd}_\varepsilon\) is total, and, for every coherent symmetry
\(a\) and every horizon truncation,
\begin{equation}
\label{eq:pdfv2-terminal-trece-naturalidad}
 a_*\operatorname{Upd}_\varepsilon
 =\operatorname{Upd}_{a\varepsilon}a_*,
 \qquad
 \operatorname{tr}_{13;j+1,j}
 \operatorname{Upd}_\varepsilon\mathfrak r_{k,j}(h)
 =\mathfrak r_{k,j}(h).
\end{equation}
\end{proposition}

\begin{proof}
The existence of each seed was exhibited in the thirteen paragraphs. In
\(\mathsf F,\mathsf E,\mathsf G,\mathsf X,\mathsf S,\mathsf T,\mathsf L\),
the formula is concatenation with an already typed map. In
\(\mathsf R,\mathsf N,\mathsf O,\mathsf K,\mathsf M,\mathsf B\), the
formula is, respectively, concatenation of a labeled block of equations,
adjunction of a numerical datum, transport of sheets, composition
of cocycles, affine action, and chain map. The naturality identities
cited after each formula prove the first equality of
\eqref{eq:pdfv2-terminal-trece-naturalidad} componentwise. The
second is exactly
\eqref{eq:pdfv2-terminal-truncamiento-trece-def}; its sound definition uses the
ordered history preserved by \(\mathsf E\) and \(\mathsf G\).
\end{proof}

\begin{proposition}[Joint Certificate of the Five Internal Developments]
\label{prop:pdfv2-terminal-cinco-desarrollos-integrales}
The updates of the thirteen coordinates jointly transport:
the spectral--polar telescoping and return without reinitialization of
\cref{thm:continuo-sp-telescopia-no-autonoma,eq:continuo-sp-no-reset};
the orthogonal balance, the remark column, and the Stein identity of
\cref{eq:pdfv2-sol-ortogonalidad,eq:pdfv2-sol-stein}; direct
and adjoint naturality and the cellular retraction of
\cref{eq:gauint-naturalidad-directa-adjunta,eq:gauint-sdr}; totalization,
the cone of the nonadic return, and the complete fibers of
\cref{eq:pdfv2-coh-total-square-zero,eq:pdfv2-coh-terminal-cone,eq:pdfv2-coh-extension-multisection}; and the Alexander--Whitney
diagonal, Smith reduction, and the Bockstein tower of
\cref{eq:detint-aw-natural,eq:detint-smith,eq:detint-bock-leading}.
All these actions have as their domain the same register generated by
\eqref{eq:pdfv2-terminal-registro-trece}; none introduces an initial
subdomain, a privileged history, or a selection motivated by a subsequent
application.
\end{proposition}

\begin{proof}
Each cited identity acts on a coordinate of the common register, and its truncation maps are restrictions of \eqref{eq:pdfv2-terminal-truncamiento-trece-def}. The componentwise naturality proved in \eqref{eq:pdfv2-terminal-trece-naturalidad} permits their assembly into a single update. Equality of domain and preservation of the ordered history preclude replacing correlative generation with five retrospective projections.
\end{proof}

\section{Multivalued terminal without selector}

The finite terminal generated by \(\widehat x_k\) is the multisection
\begin{equation}
\label{eq:pdfv2-terminal-multivaluado}
 \mathfrak T_{k,N}(\widehat x_k)
 :=\left\{\mathfrak r_{k,N}(h):
            h\in\mathscr H_{k,N}(\widehat x_k)\right\}.
\end{equation}
The expression of the right returns the set whole. No contains a
rule that distinguishes a history privileged.

The terminal level environment is obtained by gathering those multisects without
identifying records:
\begin{equation}
\label{eq:pdfv2-terminal-ambiente-nivel}
 \mathfrak T_{k,N}
 :=\bigcup_{\widehat x_k\in\widehat{\mathcal X}^{\rm tot}_k}
   \mathfrak T_{k,N}(\widehat x_k).
\end{equation}

The linear version, when contributions must be added, is defined on
the free module by
\begin{equation}
\label{eq:pdfv2-terminal-evaluacion-todas}
 \begin{aligned}
 \mathsf E_{k,N}^{\rm all}:
 \mathbb Z[\widehat{\mathcal X}^{\rm tot}_k]&\longrightarrow
 \mathbb Z[\mathfrak T_{k,N}],\\
 [\widehat x_k]&\longmapsto
 \sum_{h\in\mathscr H_{k,N}(\widehat x_k)}
 [\mathfrak r_{k,N}(h)].
 \end{aligned}
\end{equation}
Each witness appears with its multiplicity. If exists a measure projective
positive \(\mu\), the version isometric uses all the probabilities
conditional:
\begin{equation}
\label{eq:pdfv2-terminal-evaluacion-ponderada}
 \widehat{\mathsf E}_{k,N}^{\mu}e_{\widehat x_k}
 :=\sum_{h\in\mathscr H_{k,N}(\widehat x_k)}
 \mu_{\widehat x_k,N}(h)^{1/2}
 e_{\mathfrak r_{k,N}(h)}.
\end{equation}
The projectivity gives
\begin{equation}
\label{eq:pdfv2-terminal-isometria}
 \left\|\widehat{\mathsf E}_{k,N}^{\mu}e_{\widehat x_k}\right\|^2
 =\sum_{h\in\mathscr H_{k,N}(\widehat x_k)}
   \mu_{\widehat x_k,N}(h)=1.
\end{equation}
The sum is not substituted by ``the history that generates the cylinder''.

\section{The nonadic connection as correspondence}

The nonadic holonomy was generated by the odometer and the TPK connection,
without a membership hypothesis, in
--
. In this chapter\eqref{eq:continuo-comun-gamma9-apice} we use
\eqref{eq:continuo-comun-gamma9-relacion}
\begin{equation}
\label{eq:pdfv2-terminal-apice-gamma9}
 \mathsf Z_{9,k}:=\mathsf Z^{\Gamma}_{9,k}.
\end{equation}
The maps
\begin{equation}
\label{eq:pdfv2-terminal-span-gamma9}
 \widehat{\mathcal X}^{\rm tot}_k
 \xleftarrow{\ s_{9,k}\ }
 \mathsf Z^{\Gamma}_{9,k}
 \xrightarrow{\ t_{9,k}\ }
 \widehat{\mathcal X}^{\rm tot}_{k+9}
\end{equation}
are those of \eqref{eq:continuo-comun-gamma9-span}; its source map is surjective by \cref{prop:continuo-comun-gamma9-total-natural}. Every point of the apex preserves the nine phase-compatible emissions. It is not replaced by the pair of its endpoints.

For \(N\geq k+9\), define the apex of horizon before of linearizing:
\begin{equation}
\label{eq:pdfv2-terminal-apice-horizonte-gamma9}
 \mathsf Z^\Gamma_{9;k,N}:=
 \coprod_{z=(\widehat x_k,\boldsymbol\omega)
            \in\mathsf Z^\Gamma_{9,k}}
 \operatorname{Term}_{k+9,N}(t_{9,k}z).
\end{equation}
An element is written \((z,q)\), where \(\boldsymbol\omega\) is the
nonadic word of nine emissions and \(q\) its tail up to \(N\). The maps on
histories use the labeled codomain
\begin{equation}
\label{eq:pdfv2-terminal-gamma9-codominio-colas}
 \mathscr D^\Gamma_{9;k,N}:=
 \coprod_{z\in\mathsf Z^\Gamma_{9,k}}
 \operatorname{Term}_{k+9,N}(t_{9,k}z).
\end{equation}
Thus the maps are
\begin{align}
 S_{k,N}(z,q)&:=\boldsymbol\omega\star q,
 \label{eq:pdfv2-terminal-gamma9-s-horizonte}\\
 T_{k,N}(z,q)&:=(z,q)\in\mathscr D^\Gamma_{9;k,N}.
 \label{eq:pdfv2-terminal-gamma9-t-horizonte}
\end{align}
The first lands in
\(\coprod_x\mathscr H_{k,N}(x)\); the second expressly preserves the
labelled summand by \(z\).

The link map of the apex is the explicit action
\begin{equation}
\label{eq:pdfv2-terminal-gamma9-enlace-apice}
 \pi^Z_{N+1,N}(z,q\star\varepsilon):=(z,q).
\end{equation}
It is surjective: if \((z,q)\) is given, one adds to the tail the emission
neutral at the endpoint of \(q\). The codomain of tails has link
\begin{equation}
\label{eq:pdfv2-terminal-gamma9-enlace-codominio}
 \pi^{\mathscr D}_{N+1,N}(z,q\star\varepsilon):=(z,q).
\end{equation}
We denote by
\(\pi^{\mathscr H,j}_{N+1,N}\) the map
\eqref{eq:pdfv2-terminal-truncamiento-historias} with initial level \(j\).
Then
\begin{align}
 \pi^{\mathscr H,k}_{N+1,N}S_{k,N+1}
 &=S_{k,N}\pi^Z_{N+1,N},
 \label{eq:pdfv2-terminal-gamma9-cuadrado-fuente}\\
 \pi^{\mathscr D}_{N+1,N}T_{k,N+1}
 &=T_{k,N}\pi^Z_{N+1,N}.
 \label{eq:pdfv2-terminal-gamma9-cuadrado-destino}
\end{align}
In the first square, both paths send
\(\boldsymbol\omega\star q\star\varepsilon\) to
\(\boldsymbol\omega\star q\); in the second, they send
\((z,q\star\varepsilon)\) to \((z,q)\). By unique decomposition of a word
into its first nine emissions and its tail, \(S_{k,N}\) is bijective. The
second square is Cartesian by the coproduct--fibered definition of
\eqref{eq:pdfv2-terminal-apice-horizonte-gamma9}. These are the
Beck--Chevalley squares preserving witnesses and multiplicities.

The corresponding apex of registers, now generated by the thirteen updates, is
\begin{equation}
\label{eq:pdfv2-terminal-apice-terminal-gamma9}
 \begin{split}
 \mathsf Z^{\mathfrak T}_{9;k,N}:=\bigl\{
 (&z,q;
   \mathfrak r_{k,N}(\boldsymbol\omega\star q),
   \mathfrak r_{k+9,N}(q)):\\
 &z=(\widehat x_k,\boldsymbol\omega)\in
       \mathsf Z^\Gamma_{9,k},\quad
 q\in\operatorname{Term}_{k+9,N}(t_{9,k}z)
 \bigr\}.
 \end{split}
\end{equation}
Its source and target maps read, respectively, the primer or the second
record, but never remove from their domain the witness \((z,q)\):
\begin{equation}
\label{eq:pdfv2-terminal-span-terminal}
 \mathfrak T_{k,N}
 \xleftarrow{\ s^{\mathfrak T}_{9;k,N}\ }
 \mathsf Z^{\mathfrak T}_{9;k,N}
 \xrightarrow{\ t^{\mathfrak T}_{9;k,N}\ }
 \mathfrak T_{k+9,N}.
\end{equation}
If \(\zeta=(z,q;r_s,r_t)\), its action is
\begin{equation}
\label{eq:pdfv2-terminal-span-terminal-accion}
 s^{\mathfrak T}_{9;k,N}(\zeta):=r_s,
 \qquad
 t^{\mathfrak T}_{9;k,N}(\zeta):=r_t.
\end{equation}
Its link is defined by
\begin{equation}
\label{eq:pdfv2-terminal-gamma9-enlace-registros}
 \begin{split}
 \pi^{Z,\mathfrak T}_{N+1,N}
 \bigl(&z,q\star\varepsilon;
       \mathfrak r_{k,N+1}
          (\boldsymbol\omega\star q\star\varepsilon),
       \mathfrak r_{k+9,N+1}(q\star\varepsilon)\bigr)\\
 &:=(z,q;\mathfrak r_{k,N}(\boldsymbol\omega\star q),
       \mathfrak r_{k+9,N}(q)).
 \end{split}
\end{equation}
The pair \((z,q)\) is a coordinate of the apex, not a representation
chosen of a triad; therefore the link is well defined even if two
registros terminals coincide. The compatibility of
\(\operatorname{Upd}\) with the deletion of the ultima edge proof that
\eqref{eq:pdfv2-terminal-gamma9-enlace-registros} commutes with both maps of
\eqref{eq:pdfv2-terminal-span-terminal}.

If an action on free modules is desired, this
correspondence is linearized by summing all witnesses:
\begin{equation}
\label{eq:pdfv2-terminal-linealizacion-span}
 [\Gamma_9]_{k,N}[r]
 :=\sum_{\zeta\in(s^{\mathfrak T}_{9;k,N})^{-1}(r)}
 [t^{\mathfrak T}_{9;k,N}(\zeta)].
\end{equation}
The formula preserves multiplicities. Equality of phases after
nine steps does not contract the states, the witnesses, or their ledgers.

\section{Uniqueness of the prefix decomposition and naturality of the truncation system}

For \(k\leq\ell\leq N\), every history of
\(\mathscr H_{k,N}(\widehat x_k)\) has a unique prefix
\(p\in\mathscr P_{\ell\mid k}(\widehat x_k)\). Therefore there is a disjoint decomposition
by prefixes, expressed without choosing any of them:
\begin{equation}
\label{eq:pdfv2-terminal-particion-historias}
 \mathscr H_{k,N}(\widehat x_k)
 =\bigsqcup_{p\in\mathscr P_{\ell\mid k}(\widehat x_k)}
 \left\{p\star q:q\in\mathscr H_{\ell,N}(p)\right\}.
\end{equation}
Let \(\operatorname{Conc}_p\) be the operation that prepends the record of the
prefix \(p\) and successively applies
\eqref{eq:pdfv2-terminal-actualizacion-unica}. We define the transition in common over a complete family of terminals by
\begin{equation}
\label{eq:pdfv2-terminal-transicion-comun}
 \mathfrak u_{\ell,k}
 \left(
   \bigl(\mathfrak T_{\ell,N}(p)\bigr)_{p\in\mathscr P_{\ell\mid k}(\widehat x_k)}
 \right)
 :=\bigcup_{p\in\mathscr P_{\ell\mid k}(\widehat x_k)}
   \operatorname{Conc}_p\mathfrak T_{\ell,N}(p).
\end{equation}

Writing
\(\mathbf G_{k,N}(\widehat x_k):=\mathfrak T_{k,N}(\widehat x_k)\), all naturality is concentrated in the single identity
\begin{equation}
\label{eq:pdfv2-terminal-naturalidad-unica}
 \boxed{
 \mathfrak u_{\ell,k}
 \left(
   \bigl(\mathbf G_{\ell,N}(p)\bigr)_{p\in\mathscr P_{\ell\mid k}(\widehat x_k)}
 \right)
 =\mathbf G_{k,N}(\widehat x_k).}
\end{equation}
There is no distinct law for each internal discriminant. The five
characteristics obey
\eqref{eq:pdfv2-terminal-naturalidad-unica} because each is updated
by means of the same prefix, the same extension, and the same ledger.

The nonadic correspondence is compatible with this identity through the
maps already constructed, not through a condition that is implicit. For every
horizon, \eqref{eq:pdfv2-terminal-gamma9-cuadrado-fuente} and
\eqref{eq:pdfv2-terminal-gamma9-cuadrado-destino} preserve the nonadic witness with its multiplicity, and
\eqref{eq:pdfv2-terminal-gamma9-enlace-registros} does the same in the
thirteen coordinates. If \(a\) is a coherent symmetry of the common
state, its action
\[
 a^Z(z,q):=(a^Zz,aq)
\]
satisfies
\begin{equation}
\label{eq:pdfv2-terminal-gamma9-equivarianza-completa}
 S a^Z=aS,\qquad T a^Z=aT,\qquad
 \pi^Za^Z=a^Z\pi^Z,\qquad
 s^{\mathfrak T}a^Z=as^{\mathfrak T},\qquad
 t^{\mathfrak T}a^Z=at^{\mathfrak T}.
\end{equation}
These equalities are verified entry by entry using
\eqref{eq:continuo-comun-gamma9-naturalidad} and the naturality of the thirteen
updates. They are the material condition that permits the use of
\eqref{eq:pdfv2-terminal-linealizacion-span}; a mere equality of visible
values does not replace it.

\section{Passage to the limit without interchanging image and limit}

Truncating a history and each of the thirteen fields defines
\begin{equation}
\label{eq:pdfv2-terminal-transicion-horizonte}
 \pi^{\mathfrak T}_{N+1,N}:
 \mathfrak T_{k,N+1}(\widehat x_k)\longrightarrow
 \mathfrak T_{k,N}(\widehat x_k),
 \qquad
 \pi^{\mathfrak T}_{N+1,N}\mathfrak r_{k,N+1}(h)
 =\mathfrak r_{k,N}(h_{\leq N}).
\end{equation}
It is surjective by
\eqref{eq:pdfv2-terminal-truncamiento-historias}. The complete terminal is
therefore
\begin{equation}
\label{eq:pdfv2-terminal-limite-inverso}
 \mathfrak T_{k,\infty}(\widehat x_k)
 :=\left\{(r_N)_{N\geq k}:
 r_N\in\mathfrak T_{k,N}(\widehat x_k),\quad
 \pi^{\mathfrak T}_{N+1,N}r_{N+1}=r_N\right\}.
\end{equation}
The finite sets of
\eqref{eq:pdfv2-terminal-multivaluado} are nonempty and their transition maps
are surjective; therefore,
\begin{equation}
\label{eq:pdfv2-terminal-limite-no-vacio}
 \mathfrak T_{k,\infty}(\widehat x_k)\neq\varnothing.
\end{equation}

\begin{corollary}[Simultaneous persistence of the five developments]
\label{cor:pdfv2-terminal-limite-cinco-desarrollos}
In each element of
\(\mathfrak T_{k,\infty}(\widehat x_k)\) there simultaneously persist the strong
limit \(\mathcal T_\infty\) of the polar telescopy, the Gram family and
losses, the incidence measure with its two naturalities, the cone
\(\operatorname{Cone}(I-U_{\Gamma_9,N})\), and the compatible family of Smith forms and
Bockstein pages. None of these data is obtained
by reconstructing only the visible terminal value.
\end{corollary}

\begin{proof}
Surjectivity of the truncations preserves the complete witnesses. The
telescoping, naturality, and reduction equations of
\cref{prop:pdfv2-terminal-cinco-desarrollos-integrales} commute with those
truncations; by universality of the inverse limit, they define a unique compatible
family on each surviving history.
\end{proof}

The apex nonadico in the limit is defined, not inferred from an image:
\begin{equation}
\label{eq:pdfv2-terminal-gamma9-apice-limite-raw}
 \mathsf Z^\Gamma_{9;k,\infty}
 :=\varprojlim_{N\geq k+9}
 \bigl(\mathsf Z^\Gamma_{9;k,N},\pi^Z_{N+1,N}\bigr).
\end{equation}
Since the links do not modify the finite witness \(z\) and only truncate its
tail, there exists the explicit isomorphism
\begin{equation}
\label{eq:pdfv2-terminal-gamma9-apice-limite-msect}
 \mathsf Z^\Gamma_{9;k,\infty}
 \xrightarrow{\ \cong\ }
 \coprod_{z\in\mathsf Z^\Gamma_{9,k}}
 \operatorname{Msect}(t_{9,k}z),\qquad
 ((z,q_N))_N\longmapsto(z,(q_N)_N).
\end{equation}
The inverse takes a witness \(z\) and a compatible tail \((q_N)_N\), and forms
\(((z,q_N))_N\). Both compositions are identities coordinate by
coordinate. The limit maps are
\begin{equation}
\label{eq:pdfv2-terminal-gamma9-mapas-limite-raw}
 S_{k,\infty}((z_N)_N):=(S_{k,N}z_N)_N,\qquad
 T_{k,\infty}((z_N)_N):=(T_{k,N}z_N)_N;
\end{equation}
are well defined precisely by \eqref{eq:pdfv2-terminal-gamma9-cuadrado-fuente} and \eqref{eq:pdfv2-terminal-gamma9-cuadrado-destino}.

The apex of registers is defined in the same way
\begin{equation}
\label{eq:pdfv2-terminal-gamma9-apice-registros-limite}
 \mathsf Z^{\mathfrak T}_{9;k,\infty}
 :=\varprojlim_{N\geq k+9}
 \bigl(\mathsf Z^{\mathfrak T}_{9;k,N},
       \pi^{Z,\mathfrak T}_{N+1,N}\bigr).
\end{equation}
The correspondences
\eqref{eq:pdfv2-terminal-span-terminal} then form, component by
component, the span
\begin{equation}
\label{eq:pdfv2-terminal-span-limite}
 \mathfrak T_{k,\infty}
 \xleftarrow{\ s^{\mathfrak T}_{9;k,\infty}\ }
 \mathsf Z^{\mathfrak T}_{9;k,\infty}
 \xrightarrow{\ t^{\mathfrak T}_{9;k,\infty}\ }
 \mathfrak T_{k+9,\infty}.
\end{equation}
The image of a limit has not been identified with the limit of certain
images: \eqref{eq:pdfv2-terminal-limite-inverso},
\eqref{eq:pdfv2-terminal-gamma9-apice-limite-raw} and
\eqref{eq:pdfv2-terminal-gamma9-apice-registros-limite} construct
directly the three compatible families and their maps.

\section{Internal reconstruction of the integral holonomic state from its \(13\) genealogical coordinates}

Recovering does not unpack names: it repeats the fold that generated the
record. For
\(h=(\varepsilon_{k+1},\ldots,\varepsilon_N)\), we define
\begin{equation}
\label{eq:pdfv2-terminal-recuperacion-trece}
 \boxed{
 \operatorname{Rec}_{13;k,N}(h):=
 \operatorname{Upd}_{\varepsilon_N}\circ\cdots\circ
 \operatorname{Upd}_{\varepsilon_{k+1}}
 (\mathfrak r^0_{k,r}).}
\end{equation}
By \eqref{eq:pdfv2-terminal-registro-trece}, \(\operatorname{Rec}_{13;k,N}(h)=\mathfrak r_{k,N}(h)\), but the equality is now an identity between two compositions of actions, not the projection of a preserved input. If \(\operatorname{pr}_{\mathsf A}\) denotes the projector from the register to field \(\mathsf A\in\{\mathsf F,\mathsf E,\mathsf R,\mathsf N,\mathsf O,
\mathsf K,\mathsf M,\mathsf B,\mathsf G,\mathsf X,\mathsf S,\mathsf T,
\mathsf L\}\), then
\begin{equation}
\label{eq:pdfv2-terminal-recuperacion-campo-fold}
 \operatorname{pr}_{\mathsf A}\operatorname{Rec}_{13;k,N}(h)
 =\operatorname{Upd}^{\mathsf A}_{\varepsilon_N}\circ\cdots\circ
  \operatorname{Upd}^{\mathsf A}_{\varepsilon_{k+1}}
  (\mathsf A_{k,k}).
\end{equation}
The proof is induction on \(N-k\): for \(N=k\) both sides are the
seed; the inductive step applies the component \(\mathsf A\) of
\(\operatorname{Upd}_{\varepsilon_N}\).

Moreover,
\begin{equation}
\label{eq:pdfv2-terminal-recuperacion-compatible}
 \operatorname{Rec}_{13;k,N}(h_{\leq N})
 =\operatorname{tr}_{13,N+1,N}
   \bigl(\operatorname{Rec}_{13;k,N+1}(h)\bigr),
\end{equation}
where \(\operatorname{tr}_{13,N+1,N}\) removes only the last
update and restricts its cone of tails. Consequently, an element
of \(\mathfrak T_{k,\infty}(\widehat x_k)\) determines a compatible family of the
thirteen coordinates at every depth.

\section{Multivalued terminality, naturality, and compatible limit of the continuum}

\begin{theorem}[Multivalued terminal, naturality, and limit of the continuum]
\label{thm:pdfv2-terminal-conjunto}
For every \(\widehat x_k\in\widehat{\mathcal X}^{\rm tot}_k\), the rules
\eqref{eq:pdfv2-terminal-historias-finita}--
\eqref{eq:pdfv2-terminal-recuperacion-compatible} construct, from the
unique discrete continuum:
\begin{enumerate}
 \item a multisection terminal finite not empty at each horizon;
 \item an evaluation that conserves all histories and multiplicities;
 \item the connection nonadic as a correspondence with witnesses;
 \item the unique identity of naturality
       \eqref{eq:pdfv2-terminal-naturalidad-unica};
 \item a nonempty inverse terminal and a nonadic correspondence in the
       limit;
 \item the compatible recovering of the fiber, extensions, relations,
       numbers, orientation, carry, memory, boundary, route,
       multisection, survival, terminal, and reader.
\end{enumerate}
The five intrinsic characteristics remain correlated within
each record and each extension witness.
\end{theorem}

\begin{proof}
The non-emptiness of
\(\mathscr H_{k,N}(\widehat x_k)\) is nonempty by
\cref{prop:continuo-comun-arbol-finito-no-vacio}, and its truncation is
surjective by \cref{lem:continuo-comun-terminal-sobreyectivo}. The rule
\eqref{eq:pdfv2-terminal-actualizacion-unica} applies the APP,
TRIT, and TPK identities to each witness \(\varepsilon_j\); by induction on
\(N-k\), it produces all the fields of
\eqref{eq:pdfv2-terminal-registro-trece} and conserves their relations.

Taking the complete set in
\eqref{eq:pdfv2-terminal-multivaluado}, or the complete sum in
\eqref{eq:pdfv2-terminal-evaluacion-todas}, proves that the terminal does not
contain a selector. The weighted version is isometric by the projective
equality \eqref{eq:pdfv2-terminal-isometria}.

The witness \(z=(\widehat x_k,\boldsymbol\omega)\) carries the phase action of each of its nine emissions; by \cref{prop:continuo-comun-gamma9-retorno-memoria} it defines the two maps in \eqref{eq:pdfv2-terminal-span-gamma9} without converting the relation into a function or identifying phase with state. Equations \eqref{eq:pdfv2-terminal-gamma9-s-horizonte}--\eqref{eq:pdfv2-terminal-gamma9-enlace-registros} prove the two horizon squares and the same construction over registers. Summing the source space of each witness yields \eqref{eq:pdfv2-terminal-linealizacion-span} without erasing multiplicities.

The partition
\eqref{eq:pdfv2-terminal-particion-historias} is disjoint because every
history has a determined prefix of length \(\ell\). Concatenating each
prefix with all its continuations reconstructs once and only once each
history of \(\mathscr H_{k,N}(\widehat x_k)\). This directly proves
\eqref{eq:pdfv2-terminal-naturalidad-unica} for the whole register; five disconnected proofs are not
needed.

Surjectivity of
\eqref{eq:pdfv2-terminal-transicion-horizonte} and finiteness of each level
imply nonemptiness of
\eqref{eq:pdfv2-terminal-limite-inverso}. The explicit links
\eqref{eq:pdfv2-terminal-gamma9-enlace-apice} and
\eqref{eq:pdfv2-terminal-gamma9-enlace-registros}, together with their squares,
construct the limits
\eqref{eq:pdfv2-terminal-gamma9-apice-limite-raw} and
\eqref{eq:pdfv2-terminal-gamma9-apice-registros-limite}; from them follows
\eqref{eq:pdfv2-terminal-span-limite}. Finally,
\eqref{eq:pdfv2-terminal-recuperacion-trece} reruns the
successive updates, and
\eqref{eq:pdfv2-terminal-recuperacion-compatible} proves that those
readings pass jointly to the limit.
\end{proof}

\section{Obstructions to terminality of the integral holonomic state: loss of fiber, orientation, memory or reader}

\begin{proposition}[Falsifiers of the terminal set]
\label{prop:pdfv2-terminal-falsadores}
Any of the following operations replaces the constructed object and
blocks the conclusion of
\cref{thm:pdfv2-terminal-conjunto}:
\begin{enumerate}
 \item send to cylinder to to single history when it admits several;
 \item choosing a tail of the terminal field \(\mathsf T\) and eliminating the
       remaining ones;
 \item replace \(\Gamma_9\) by a function without proving uniqueness of
       each target and of its witness;
 \item linearize to correspondence while ignoring multiplicities;
 \item identifying phase return with state return;
 \item preserve the input intact as the first coordinate and use its
       recoverability as purported generation;
 \item omit any of the thirteen fields of
       \eqref{eq:pdfv2-terminal-registro-trece};
 \item allowing two intrinsic characteristics to use different histories,
       carries, boundaries, or terminals;
 \item replace
       \eqref{eq:pdfv2-terminal-naturalidad-unica} by five independent
       assertions;
 \item assert operatorial compatibility of \(\Gamma_9\) without preserving the
       space of witnesses and its multiplicities;
 \item interchanging, without proof, formation of images and passage to the limit;
 \item present a list of thirteen names without the actions
       \eqref{eq:pdfv2-terminal-actualizacion-unica} that produce them.
 \item omit the polar telescoping, the loss row, the naturality
       adjoint, the module of all continuations or the reduction
       Smith--Bockstein certified in
       \cref{prop:pdfv2-terminal-cinco-desarrollos-integrales};
 \item turning any of those five developments into an autonomous branch
       or using a later classical application to select its domain.
\end{enumerate}
\end{proposition}

\begin{proof}
The first five operations contract the multisection or the correspondence. The
next four break their common support and naturality. The final
three eliminate the generative proof or alter passage to the limit. In
each case, the resulting object ceases to satisfy at least one of
\eqref{eq:pdfv2-terminal-multivaluado},
\eqref{eq:pdfv2-terminal-span-terminal},
\eqref{eq:pdfv2-terminal-naturalidad-unica}, or
\eqref{eq:pdfv2-terminal-recuperacion-compatible}; therefore, the conclusion
does not carry over.
\end{proof}

\begin{statusbox}
The organization of the complete terminal object, unique naturality, and passage to the
limit are established jointly. The surviving multisection, the
carries, the register, the boundary, the route, and the nonadic connection are its
structural antecedents.\quad
The linearization over free modules is proved by
\eqref{eq:pdfv2-terminal-linealizacion-span} and its Beck--Chevalley
squares. A subsequent promotion to a bounded operator on an
analytic completion additionally requires that norm to be declared and the bound proved; the
correspondence \eqref{eq:pdfv2-terminal-span-limite} does not depend on that
interface.
\end{statusbox}
 \end{HMTScientificOwnerSubordinated}

\HMTDeferredContinuumConsequences
 
\chapter{Projective measure and evolutionary conservation of the unit}
\label{chap:83-06}

\section{Projective measure and evolutionary conservation of the unit}
\label{p12-83-c6-s1:chap:medida-proyectiva}

Cylinder topology provides a notion of proximity; a theory of paths also requires weights. The evolutionary conservation of the unit here takes an exact probabilistic form: every cylinder distributes its weight among its extensions, and the total sum remains equal to one. The unit is not fixed in a cell; it is transported through refinement.

\subsection{Consistent weights on finite levels}

Let \(\mu_n\) be a probability on the finite set
\(\mathcal S_n\). The family is \emph{projectively consistent} when
\begin{equation}
 \mu_n(a)
 =\sum_{\substack{b\in\mathcal S_{n+1}\\
                   \rho_{n+1,n}(b)=a}}
   \mu_{n+1}(b)
 \qquad(a\in\mathcal S_n).
 \label{p12-83-c6-s1:eq:consistencia-pesos}
\end{equation}
This equality expresses the local preservation of the unit: the weight of a
truncated history is the sum of the weights of all its children.

\begin{theorem}[Genealogical measure on cylinders]
\label{p12-83-c6-s1:thm:medida-genealogica}
Every family of finite probabilities that satisfies
\eqref{p12-83-c6-s1:eq:consistencia-pesos} determines a unique Borel probability
\(\mu_\infty\) on \(\Omega_\infty\) such that
\begin{equation}
 \mu_\infty([a]_n)=\mu_n(a).
 \label{p12-83-c6-s1:eq:medida-cilindro}
\end{equation}
\end{theorem}

\begin{proof}
\[
 \mathscr A
 :=\{\text{finite unions of cylinders}\}
\]
is an algebra: the intersection of two cylinders is either empty or a cylinder of the
deeper level. Every element of \(\mathscr A\) can be refined into a
disjoint union of cylinders of a common level. If
\(A=\bigsqcup_{a\in I}[a]_n\), we set
\[
 \mu_0(A)=\sum_{a\in I}\mu_n(a).
\]
Iterating \eqref{p12-83-c6-s1:eq:consistencia-pesos} proves that the value does not depend on the
chosen common level.

If a cylinder \([a]_n\) is empty, finite branching and Koenig's lemma
\cite{Koenig1936}
imply that there is a depth with no descendants of \(a\). Iterated
consistency then gives \(\mu_n(a)=0\). Hence \(\mu_0\) is well
defined even for representations containing empty cylinders and is
finitely additive.

The cylinders are both open and closed. If a decreasing sequence of
elements of \(\mathscr A\) had empty intersection, compactness
would imply that one of them was already empty; hence \(\mu_0\) is continuous at
the empty set and constitutes a countably additive premeasure. The Caratheodory extension theorem ---equivalently, the Kolmogorov extension for this finite
projective system--- extends it to the \(\sigma\)-algebra generated by the
cylinders \cite{Caratheodory1914,Kolmogorov1933}. This is the Borel algebra of
\(\Omega_\infty\), and finiteness of the
measure ensures uniqueness.
\end{proof}

The measure may follow from local transitions. If
\(P_n(b\mid a)\ge0\) satisfies
\[
 \sum_{\rho_{n+1,n}(b)=a}P_n(b\mid a)=1,
\]
then
\(\mu_{n+1}(b)=\mu_n(a)P_n(b\mid a)\) is consistent. Uniformity is a
possible choice, not an obligation. The weights may depend on the sheet,
carry, orientation, a nonnegative action cost, or a record, provided
they preserve the sum. For example,
\[
 P_n(b\mid a)
 =\frac{e^{-\beta a_n(b\mid a)}}
        {\sum_{b'\mapsto a}e^{-\beta a_n(b'\mid a)}},
 \qquad a_n\geq0,
\]
defines normalized positive weights. The orthogonal action phase belongs to the
kernel of \cref{chap:integral-genealogica}; it is not inserted into these
positive probabilities.

\subsection{Decomposition of the genealogical measure on fibers of equal Archimedean value}

The Archimedean evaluation transports the measure to \(K\):
\begin{equation}
 \nu:=\Val_*\mu_\infty,
 \qquad
 \nu(B)=\mu_\infty(\Val^{-1}(B)).
 \label{p12-83-c6-s1:eq:pushforward-arquimediano}
\end{equation}
The visible probability of a region sums the weight of all the genealogies that
publish it. When two sections possess the same value, their contributions are
aggregated in \(\nu\), although they remain separate in \(\mu_\infty\).

\begin{proposition}[Value Decomposition]
If \(K\) is a standard Borel space, there exists a family of conditional
measures \(\{\mu_y\}_{y\in K}\), defined for \(\nu\)-almost every \(y\),
such that \(\mu_y\) is supported on \(\mathfrak G_y=\Val^{-1}(y)\) and
\[
 \int_{\Omega_\infty}F\dd\mu_\infty
 =\int_K\left(\int_{\mathfrak G_y}F\dd\mu_y\right)\dd\nu(y)
\]
for every integrable function \(F\).
\end{proposition}

This formula separates value and genealogy at the level of the measure. Physics that
depends only on the value first integrates the entire fibre; a memory observable
can distinguish different distributions \(\mu_y\) over the same
real coordinate.

\subsection{Conditioning as pruning of futures}

Let \(A\subseteq\Omega_\infty\) be an observable event with
\(\mu_\infty(A)>0\). The conditional measure is
\begin{equation}
 \mu_\infty(B\mid A)
 =\frac{\mu_\infty(B\cap A)}{\mu_\infty(A)}.
 \label{p12-83-c6-s1:eq:condicionamiento-genealogico}
\end{equation}
When \(A\) is the reading of a register, its intersection with each cylinder
eliminates the incompatible extensions. Normalization does not create new branches:
it redistributes the unit among those that have survived.

\begin{proposition}[Compatible update]
If \(A\) is a union of cylinders at level \(m\), then for every \(n\ge m\)
the conditional distribution on \(\mathcal S_n\) is obtained by removing the
states whose truncations do not belong to \(A\) and renormalizing the remaining
weights. The conditional distributions continue to satisfy
\eqref{p12-83-c6-s1:eq:consistencia-pesos}.
\end{proposition}

\begin{proof}
The identity is Bayes' rule on finite sets. The sum of the
weights of the conditioned children coincides with the conditioned weight of the parent,
because intersection and disjoint union commute.
\end{proof}

\subsection{Combinatorial conservation and dimensional completion}

The same law appears in the binomial completion:
\[
 1=\sum_{r=0}^{d}\binom dr
 \left(\frac9{10}\right)^{d-r}
 \left(\frac1{10}\right)^r.
\]
In three dimensions,
\[
 1000=729+243+27+1,
\]
so that the interior, new faces, extreme edges, and final vertex form
a partition of decimal unity. The passage
\(729+270=999\to729+271=1000\) distinguishes the state prior to closure of the
completed unity. In the path system, the analogous identity is
\[
 \mu([a]_n)=\sum_{b\mapsto a}\mu([b]_{n+1}).
\]
The unit evolves as it is distributed among the strata and is recovered by summing
the complete genealogy.

\begin{figure}[htbp]
\centering
\begin{tikzpicture}[>=Latex,node distance=12mm and 17mm,
  box/.style={draw=hmtblue,rounded corners,fill=hmtlightblue,align=center,
              minimum width=28mm,minimum height=8mm},
  leaf/.style={draw=hmtgreen,rounded corners,fill=hmtlightgreen,align=center,
               minimum width=22mm,minimum height=7mm}]
\node[box] (a) {$[a]_n$\\weight $\mu_n(a)$};
\node[leaf,below left=of a] (b1) {$[b_1]_{n+1}$\\$\mu_{n+1}(b_1)$};
\node[leaf,below=of a] (b2) {$[b_2]_{n+1}$\\$\mu_{n+1}(b_2)$};
\node[leaf,below right=of a] (b3) {$[b_3]_{n+1}$\\$\mu_{n+1}(b_3)$};
\draw[->] (a)--(b1); \draw[->] (a)--(b2); \draw[->] (a)--(b3);
\node[below=20mm of b2,align=center,text width=.78\textwidth] (eq)
 {$\mu_n(a)=\mu_{n+1}(b_1)+\mu_{n+1}(b_2)+\mu_{n+1}(b_3)$:\\
the unit is preserved when the genealogy is refined.};
\end{tikzpicture}
\caption{Projective weight conservation on cylinders.}
\label{p12-83-c6-s1:fig:conservacion-peso-cilindros}
\end{figure}

\begin{statusbox}
The projective measure is constructed on the HMT inverse system.
The extension theorem is classical mathematics; the identification of its
cylinders with TPK states is exact once the transition weights have been fixed.
The selection of a specific physical measure belongs to the dynamical model and does not
follow solely from the census of branches.
\end{statusbox}

The projective measure is characterized by positivity, conservation of
weight between parents and children, and conditioning within the observed event. A
complete experimental realization adds preparation and the sampling
protocol; with those data, the same measure produces reading probabilities
comparable with frequencies.
 
\chapter{Archimedean and exceptional morphisms from a common domain}
\label{chap:83-17}

\section{Archimedean and exceptional morphisms from a common enriched domain}
\label{p12-83-c17-s1:chap:doble}
\index{double projection}
\index{Monstrous Moonshine}

The double projection compares two evaluations of the common enriched domain
\(X_\infty^{\rm cal}\) already constructed in
\cref{p12-83-c17-s2:thm:lector-excepcional-global-rev6}. One produces
Archimedean values; the other, together with a calibration of its codomain, preserves
the combinatorial incidence that leads to codes, designs, and
automorphism groups. Both evaluations start from the unique initial state
\(X_\infty^{\rm cal}\) and preserve their common provenance in distinct
codomains.

\subsection{Domains of the \(2\) projections}

On \(X_\infty^{\rm cal}\), the cylinders contract to a point and the
boundary admits the incidence prolongation. In this residence, we fix
\(\mathbf{Arch}:=\mathbb R\). The Archimedean evaluation and the exceptional
evaluation are, respectively,
\[
\mathfrak R:X_\infty^{\rm cal}\longrightarrow\mathbf{Arch},
\qquad
\mathfrak E:X_\infty^{\rm cal}\times\mathscr C_{\rm exc}
\longrightarrow\mathbf{Exc}.
\]
Here \(\mathscr C_{\rm exc}\) contains the projective order, Paley gauge,
dodecad, incidence alignment, and lattice chamber
specified in each construction. The first map produces values; the
second produces codes, designs, lattices, and automorphisms. The preceding
global theorem is the source of compatibility and permits comparison of
invariants obtained through distinct operations without again intersecting
domains defined separately. The finite comparison in this work is
performed on the flag
\(B_K\subset H^-_\alpha\).

\begin{figure}[H]
\centering
\resizebox{\textwidth}{!}{\begin{tikzpicture}[font=\small,
  big/.style={draw,rounded corners=3mm,align=center,minimum width=3.7cm,minimum height=1.25cm},
  arr/.style={-{Stealth[length=2.6mm]},very thick}]
  \node[big,draw=hmtgold,fill=hmtlightgold,font=\bfseries] (x) at (0,0) {integral holonomic state\\\(\XTPK\)};
  \node[big,draw=hmtblue,fill=hmtlightblue] (arq) at (5,2.6) {evaluation Archimedean\\cylinders \(\to\RR\)};
  \node[big,draw=hmtgreen,fill=hmtlightgreen] (exc) at (5,0) {exceptional projection\\Witt, 3-neighbor, and Leech};
  \draw[arr,hmtblue] (x)--node[above left]{\(\mathfrak R\)} (arq);
  \draw[arr,hmtgreen] (x)--node[above]{\(\mathfrak E\)} (exc);
  \node[big,draw=hmtgray,fill=hmtpaper] (value) at (10,2.6) {real value\\evaluation class};
  \node[big,draw=hmtgray,fill=hmtpaper] (sym) at (10,0) {automorphism\\of the incidence};
  \draw[arr,hmtblue] (arq)--(value);
  \draw[arr,hmtgreen] (exc)--(sym);
  \draw[dashed,hmtred,thick,<->] (value.south west) to[bend right=18] node[fill=white,inner sep=1pt]{invariantes comunes} (sym.north west);
\end{tikzpicture}
 }
\caption{Typed scheme of the two projections. The Archimedean evaluation
identifies states within the same fiber; the exceptional evaluation, relative
to its calibration, preserves incidence structure and automorphisms.}
\label{p12-83-c17-s1:fig:doble-proyeccion}
\end{figure}

\subsection{Topology of the space of histories}

The formulation by cylinders admits a standard realization as an
inverse limit. Let \(\mathsf S_n\) be the finite sets of states
admissible at depth \(n\), endowed with the discrete topology, and let
\[
 \tau_{n+1,n}:\mathsf S_{n+1}\longrightarrow\mathsf S_n
\]
the aplicaciones of truncation. Supondremos in this seccion that are
sobreyectivas and definiremos
\[
 \mathcal X=\varprojlim_n\mathsf S_n.
\]
If \(x\ne y\), let \(m(x,y)\) be the first depth at which their components differ. The function
\[
 d(x,y)=3^{-m(x,y)},\qquad d(x,x)=0,
\]
is an ultrametric that induces the topology of the inverse limit.

\begin{theorem}[Space of compact histories]
\label{p12-83-c17-s1:thm:compacto-historias}
\label{thm:compacto-historias}
The space \(\mathcal X\) is compact and totally disconnected. When
every history admits bifurcations at arbitrarily large depths,
\(\mathcal X\) has no isolated points and is homeomorphic to Cantor
space.
\end{theorem}

\begin{proof}
The product \(\prod_n\mathsf S_n\) is compact. The equations
\(\tau_{n+1,n}(x_{n+1})=x_n\) define a closed subset, so that
\(\mathcal X\) is compact. The sets of histories with a fixed prefix
are simultaneously open and closed and form a basis; this proves total
disconnectedness. Under the hypothesis of arbitrarily deep branching,
no cylinder contains an isolated history. The characterization of
compact, metric, perfect, and totally disconnected spaces proves
the final assertion.
\end{proof}

Associate with each \(s\in\mathsf S_n\) a closed interval \(I_n(s)\),
compatibly with truncation, and suppose that
\[
 \operatorname{diam}I_n(s)\leq C\lambda^n,
 \qquad 0<\lambda<1.
\]
The intersection of the intervals corresponding to a history then contains a unique point.

\begin{theorem}[Archimedean Evaluation]
\label{p12-83-c17-s1:thm:evaluacion-holder}
\label{thm:evaluacion-holder}
The map
\[
 \operatorname{val}:\mathcal X\longrightarrow\mathbb R,
 \qquad
 \{\operatorname{val}(x)\}=\bigcap_n I_n(x_n),
\]
is continuous and Hölder with respect to \(d\). If
\(x\sim_{\mathrm{Arch}}y\) is equivalent to
\(\operatorname{val}(x)=\operatorname{val}(y)\), the induced application
is a homeomorphism
\[
 \mathcal X/{\sim_{\mathrm{Arch}}}
 \;\cong\;\operatorname{val}(\mathcal X).
\]
\end{theorem}

\begin{proof}
Two histories with a common prefix of length \(n\) have values in the
same interval of diameter no greater than \(C\lambda^n\). Since
\(d(x,y)=3^{-n}\), this estimate is a Hölder bound. The induced
map on the quotient is a continuous bijection from a compact space onto
a Hausdorff space and therefore a homeomorphism.
\end{proof}

The theorem formalizes the Archimedean projection: the value retains the
evaluation class and eliminates the distinction between histories lying in
the same fiber. Its image is \(\operatorname{val}(\mathcal X)\); identifying it with the whole of
\(\mathbb R\) additionally requires surjectivity to be proved.

\subsection{Separation by the \(2\) projections}

\begin{theorem}[Joint separation]
\label{p12-83-c17-s1:thm:separacion-conjunta}
\label{thm:separacion-conjunta}
Let \(P_1:\mathcal X\to Y_1\) and \(P_2:\mathcal X\to Y_2\) be continuous maps
with Hausdorff codomains. If
\[
 x\ne y\quad\Longrightarrow\quad
 P_1(x)\ne P_1(y)\ \text{or}\ P_2(x)\ne P_2(y),
\]
the diagonal application
\[
 (P_1,P_2):\mathcal X\longrightarrow Y_1\times Y_2
\]
is a topological embedding.
\end{theorem}

\begin{proof}
The hypothesis expresses precisely the injectivity of the diagonal map. A
continuous injective map from a compact space to a Hausdorff space is a
homeomorphism onto its image.
\end{proof}

With a declared gauge \(c\in\mathscr C_{\rm exc}\) fixed, let
\[
 \mathfrak E_c:X_\infty^{\rm cal}\longrightarrow\mathbf{Exc},
 \qquad
 \mathfrak E_c(x):=\mathfrak E(x,c).
\]
Applied to HMT, the theorem takes
\[
 P_1=\mathfrak R,\qquad P_2=\mathfrak E_c
\]
on the common domain, or on the quotient by the coordinate changes that
have been declared. The global separation assertion requires verifying
that no pair of inequivalent histories remains in the same fiber of both
maps.

At a finite level \(S\), this condition is decidable. Given
\(p:S\to A\), \(q:S\to B\), and an equivalence relation of coordinates,
the pairs \(x\not\sim y\) that simultaneously satisfy are enumerated
\[
 p(x)=p(y),\qquad q(x)=q(y).
\]
The absence of such pairs is equivalent to joint injectivity at that
level. The calculation must perform the two maps separately: when
\(q=f\circ p\), the second has the same fibers as the first and does not
recover information eliminated by \(p\).

At the level developed here, the equality
\[
 B_K=H_e\cap P_\pi
\]
constitutes a finite witness of compatibility between the two branches. It
does not by itself establish joint separation at every depth; that
assertion corresponds exactly to the hypothesis of
\cref{p12-83-c17-s1:thm:separacion-conjunta}.

\subsection{Archimedean quotient by the evaluation fibers}

An HMT section contains a block, an orientation, a carry, an observable label,
and an extension coordinate. The
evaluation \(\val\) contracts its cylinders to a point. If two sections
\(x,y\) satisfy
\[
\val(x)=\val(y),
\]
the Archimedean evaluation identifies them although their transport is different. In
this precise sense,
\[
\boxed{\text{a real value is an equivalence class of the evaluation}.}
\]
On any subdomain on which addition, multiplication, and evaluation descend
compatibly, the image inherits those operations. This section does not
thereby construct all of \(\RR\) or prove surjectivity onto
\(\RR\).

A global identification with \(\RR\) would require existence of the limit, determination
of its image, descent of sum and product, and characterization of the
fibres. The finite windows and the conditional inverse-limit theorem
provide the structure; prolongation to arbitrary depth requires the
hypotheses of unit survival and compatibility.

\subsection{Common invariant with the Witt system}

The vector \(K\) defines the support superior
\[
B_K=\{m:K_m\ge729\}=\{4,6,9,10\}.
\]
In the Witt atlas appear the supports
\[
P_\pi=\{2,4,5,6,7,8,9,10,11\},
\]
\[
H_e=\{1,3,4,6,9,10\},
\qquad
H_\varphi=\{1,2,3,4,7,9\}.
\]
The support of the negative carries of \(\alpha\) is
\[
H_\alpha^-=\{4,5,6,7,9,10\}.
\]
The support identities are verified.
\[
\boxed{
B_K=H_e\cap P_\pi=H_e\cap H_\alpha^-,}
\]
and
\[
\boxed{H_\alpha^-=B_K\sqcup\{5,7\}.}
\]
Witt polarity sends the four points of \(B_K\) to the pairs
\[
\{1,3\},\quad\{2,11\},\quad\{8,12\},\quad\{5,7\}.
\]
Three receive the positive label; the last completes the negative support. These identities are exact in the declared gauge, but they are conditioned controls: the vector, the threshold, the words, and the Paley gauge are not statistically independent inputs. The equality follows by direct evaluation and does not constitute a probabilistic estimate.

\subsection{Residual design on a dodecada and elevation of hexads}
\label{p12-83-c17-s1:sec:w12-w24-elevacion}

The relation between the two Witt designs is obtained within the extended
binary Golay code and not through an identification of cardinalities.
Let
\[
\mathcal G_{24}\subset\mathbb F_2^{24}
\]
the extended binary Golay code, and denote by \(\mathcal O_{24}\) the family
of supports of its weight-eight words. It is a classical result that
\((\{1,\ldots,24\},\mathcal O_{24})\) is the Steiner system
\(W_{24}=S(5,8,24)\) \cite{MacWilliamsSloane1977,ConwaySloane1999}.
Fix a weight-twelve word and write \(D\) for its support. We define
\begin{equation}
 \mathcal H(D)
 :=\{O\cap D:O\in\mathcal O_{24},\ |O\cap D|=6\}.
 \label{p12-83-c17-s1:eq:hexadas-residuales-dodecada}
\end{equation}
The family \(\mathcal H(D)\) is intrinsic to the pair
\((\mathcal G_{24},D)\): it does not require choosing coordinates in \(D\).

\begin{theorem}[Residual design of a dodecada]
\label{p12-83-c17-s1:thm:dodecada-w12}
The incidence system
\[
 (D,\mathcal H(D))
\]
is a Steiner system \(S(5,6,12)\). In particular,
\(|\mathcal H(D)|=132\), and every subset of five points of \(D\) is
contained in a unique hexad of \(\mathcal H(D)\).
\end{theorem}

\begin{proof}
Let \(A\subset D\), with \(|A|=5\). Since the octads form
\(S(5,8,24)\), there is a unique octad \(O_A\) containing \(A\). If
\(j=|O_A\cap D|\), the binary sum of the characteristic words of
\(O_A\) and \(D\) belongs to \(\mathcal G_{24}\) and has weight
\[
 \operatorname{wt}(\mathbf1_{O_A}+\mathbf1_D)=8+12-2j=20-2j.
\]
The weights of the extended binary Golay code are
\(0,8,12,16,24\). Since \(j\leq8\), it follows that
\(j\in\{2,4,6\}\). The inclusion \(A\subset O_A\cap D\) gives \(j\geq5\),
hence \(j=6\). Therefore, \(O_A\cap D\in\mathcal H(D)\) is a hexad containing
\(A\).

If two members of \(\mathcal H(D)\) contained \(A\), their octads of
origin would contain the same set of five points; the Steiner property of
\(W_{24}\) would make them coincide. This proves uniqueness. The number
of blocks is obtained by counting incidences with five-point subsets:
\[
 |\mathcal H(D)|
 =\frac{\binom{12}{5}}{\binom65}=132.\qedhere
\]
\end{proof}

\begin{theorem}[Unique elevation of a hexad]
\label{p12-83-c17-s1:thm:elevacion-unica-hexada}
For each \(H\in\mathcal H(D)\), there exists a unique octad
\(O_H\in\mathcal O_{24}\) such that
\[
 O_H\cap D=H.
\]
Consequently, the application
\[
 \iota_D:\mathcal H(D)\longrightarrow\mathcal O_{24},
 \qquad H\longmapsto O_H,
\]
is well defined and is injective.
\end{theorem}

\begin{proof}
Existence forms part of the definition of \(\mathcal H(D)\). If two
octads \(O_1,O_2\) had intersection \(H\) with \(D\), both
would contain any five-point subset of \(H\). The uniqueness of
the octad incident with those five points implies \(O_1=O_2\).
Injectivity follows by intersecting again with \(D\).
\end{proof}

The lift has a local law that will be used to compare the five
regions of \(\pi\). In a system \(S(t,k,v)\), the number of blocks
containing a fixed subset of \(s\leq t\) points is
\[
 \lambda_s=\frac{\binom{v-s}{t-s}}{\binom{k-s}{t-s}}.
\]
Consequently,
\begin{equation}
 \lambda_4(W_{12})
 =\frac{\binom81}{\binom21}=4,
 \qquad
 \lambda_4(W_{24})
 =\frac{\binom{20}1}{\binom41}=5.
 \label{p12-83-c17-s1:eq:lambda4-w12-w24}
\end{equation}

\begin{corollary}[Decomposition \(4+1\) on a face]
\label{p12-83-c17-s1:cor:cuatro-mas-uno-w24}
Let \(B\subset D\), with \(|B|=4\). The five octads of \(W_{24}\) that
contain \(B\) decompose uniquely into
\[
 \{O_H:H\in\mathcal H(D),\ B\subset H\}
 \;\sqcup\;\{O_B^\perp\}.
\]
The first set contains four octads and satisfies
\(|O_H\cap D|=6\); the remaining octad satisfies
\[
 O_B^\perp\cap D=B.
\]
\end{corollary}

\begin{proof}
The first equality in \eqref{p12-83-c17-s1:eq:lambda4-w12-w24} provides four residual
hexads containing \(B\), and
\cref{p12-83-c17-s1:thm:elevacion-unica-hexada} provides four distinct octads. The
second equality in \eqref{p12-83-c17-s1:eq:lambda4-w12-w24} shows that there exists exactly
one fifth octad \(O_B^\perp\).

Since \(B\subset O_B^\perp\cap D\), the intersection spectrum used in
the proof of \cref{p12-83-c17-s1:thm:dodecada-w12} leaves only the possibilities four and
six. An intersection of size six would be a fifth residual hexad that
contains \(B\), contradicting \(\lambda_4(W_{12})=4\). Therefore, the
intersection has size four and coincides with \(B\).
\end{proof}

\subsubsection{Alineamiento with the \(12\) positions HMT}

The residual design is defined on the points of \(D\), whereas the
design obtained from the ternary code \(C_W\) is defined on the twelve
dodecaphase positions. To avoid conflating the two sets, we shall call a bijection
\emph{HMT binary alignment}
\begin{equation}
 \ell:\{1,\ldots,12\}_{\rm HMT}\xrightarrow{\sim}D
 \quad\text{such that}\quad
 \{\ell(H):H\in W_{12}^{\rm HMT}\}=\mathcal H(D).
 \label{p12-83-c17-s1:eq:alineamiento-hmt-w24}
\end{equation}
Here \(W_{12}^{\rm HMT}\) denotes the design of weight-six supports
obtained from the ternary code \(C_W\). Uniqueness of the Witt system
\(S(5,6,12)\) up to isomorphism ensures the existence of \(\ell\). Two
distinct alignments differ by automorphisms of the source and target
designs; hence \(\ell\) is coordinate data and not an additional
invariant.
The pair \((D,\ell)\) is the calibration datum required to transport
the lift \(\iota_D\) to the HMT\@ coordinate system. Once fixed, each
HMT hexad \(H\) unambiguously determines the octad
\[
 \widehat H:=\iota_D(\ell(H)).
\]
The dodecad and the alignment are data of the binary codomain; they are not identified
with a decimal output of the transport kernel.

For the distinguished face
\[
 B_K=\{4,6,9,10\},
\]
the four incident hexads are
\[
 B_K\sqcup\{1,3\},\quad
 B_K\sqcup\{2,11\},\quad
 B_K\sqcup\{5,7\},\quad
 B_K\sqcup\{8,12\}.
\]
The third agrees with
\(H_\alpha^-=B_K\sqcup\{5,7\}\). The alignment \(\ell\) lifts these
four hexads to the four residual octads over \(\ell(B_K)\);
\cref{p12-83-c17-s1:cor:cuatro-mas-uno-w24} also determines a unique transverse octad.

\subsubsection{Incidence correspondence between the microfiber of \texorpdfstring{\(\pi\)}{pi} and the residual octads: multiplicity \texorpdfstring{\(432+4\cdot144=1008\)}{432+4x144=1008}}

The exact microfiber of the word
\(w_\pi=010211\) contains five distinct decimal regions:
\[
\begin{array}{c|c|c}
\text{multiplicity}&U_6& e(U_6)\\
\hline
432&501|614|498|169|272|272&(5,5,5,5,5,5)\\
144&810|923|870|169|272|272&(3,3,9,5,5,5)\\
144&810|983|810|169|272|272&(3,9,3,5,5,5)\\
144&870|923|810|169|272|272&(9,3,3,5,5,5)\\
144&870|923|810|418|674|521&(9,3,3,7,1,7).
\end{array}
\]
Here \(e(U_6)\) is the multiplicative signature recovered from each
decimal block \(\overline{d_1d_2d_3}\) by means of
\(e=d_2-d_1\pmod {10}\). The first region is the only one of multiplicity
432 and constant signature. The remaining four have multiplicity 144; among
them, three share the return block \(169|272|272\), whereas the
fourth has return \(418|674|521\). The catalogue distinguishes the five
rows by a reproducible predicate: the three rows of type
\(\mathrm{ES}\) form the positive triple; among the two rows of type
\(\mathrm{NO}\), that of maximal multiplicity and constant signature is termed
transverse, and the remaining one is termed negative. Relative to this role
predicate, the partition is
\[
 \boxed{1_\perp+3_{++}+1_{--}.}
\]
The sign labels are a convention; what is intrinsic to the catalog is the
partition obtained from its columns, not a prior binary polarity.

The alineamiento regional complete the givesto \((D,\ell)\) with to bijection
that sends the transver region to \(O_{\ell(B_K)}^\perp\), the region
negativa to \(\widehat{H_\alpha^-}\) and the tris regionis positivas to the
elevacionis of
\[
 B_K\sqcup\{1,3\},\qquad
 B_K\sqcup\{2,11\},\qquad
 B_K\sqcup\{8,12\}.
\]
The ordered assignment between the three positive regions and those three hexads is a choice of coordinates. The two distinguished rows and the triple are reproducible through the preceding catalogue predicate; their designation as negative, transversal, and positive belongs to the alignment. Thus, the law \(4+1\) does not identify the five regions with five octads by mere cardinality: it provides the incidence space, whereas \((D,\ell)\) and the regional alignment specify the transport between the two objects.

\subsection{Lattice prolongation and dimensional screens}

The lift of hexads opens two branches that must not be confused. The first
conserves binary incidence and prolongs \(W_{12}\) toward \(W_{24}\). The
second conserves ternary transport:
\[
C_{12}^{(3)}
\longrightarrow N(A_2^{12})
\longrightarrow\Lambda_{24}
\longrightarrow V^\natural
\longrightarrow\mathbb M
\longrightarrow\text{Moonshine}.
\]
The primary proof of the gluing, the rootless \(3\) neighbor, and the corridor
to Moonshine resides in
\cref{chap:bandera-excepcional-acumulativa,thm:corredor-hmt-moonshine}.
There is no canonical arrow \(W_{24}\to N(A_2^{12})\), nor any direct
action of the Monster on TPK words: both branches start from the same
calibrated code and preserve different structures.

The reductions \(243/11\), \(12\to11\to10\), and \(26\to24\) are not
identities between numerals either. The \cref{chap:pantallas-dualidad} constructs their
spaces, quotients, and maps: the perfect ternary quotient, the screen
marked by \(K\), and the isotropic reduction of the Lorentzian lattice. This
section thereby retains a single function: to link the already proved finite
incidence with its two prolongations, without repeating the reference proofs.

\subsection{Residue–quotient involution and conservation of state information}

For \(n=9Q+r\), we define the lattice energy
\[
E_R(r,Q)=\frac{r^2}{R^2}+Q^2R^2.
\]
Then
\[
\boxed{E_R(r,Q)=E_{1/R}(Q,r).}
\]
The exchange \((r,Q,R)\leftrightarrow(Q,r,1/R)\) is an exact involution
of the lattice model: the residue and quotient exchange their functions.

\subsection{Finite compatibility and joint prolongation}

The finite witness \(B_K\), the carry pattern, and the Witt frame fix
local compatibility. Its prolongation is no longer left as an abstract
scheme: the \cref{p12-83-c17-s2:thm:lector-excepcional-global-rev6} constructs, block by
block, the exceptional reader with a unique transported initial gauge, and the
\cref{p12-83-c17-s2:thm:doble-realizacion-global-rev6} proves that this reading and the
Archimedean evaluation globalize over the same TPK section. Truncating the
history before reading is equivalent in both cases to truncating the output after
reading.

The Archimedean evaluation identifies histories by their limiting value. The
exceptional evaluation preserves the visible current and its Paley--Witt
incidence; it does not thereby recover the hidden memories that the visible emitter
does not publish. The double projection is therefore a joint construction with
two distinct codomains, not an identification between the continuum and the
exceptional branch.

\begin{remark}[Compatibility of the two projections]
The support \(B_K\), the carry pattern, and the Witt frame provide the
common invariants. The lattice chain prolongs this compatibility to a
rootless neighbor and, by classical classification, to the Leech lattice. The
binary branch toward \(W_{24}\) remains separate: its incidence morphism is
\(\iota_D\), defined after fixing the dodecad and alignment
\((D,\ell)\), and its codomain is neither a Niemeier lattice nor the
Leech lattice.
\end{remark}


\section{Exceptional projection of the integral holonomic state: compatible incidence at every depth}
\label{p12-83-c17-s2:chap:lector-excepcional-global-rev6}
\index{exceptional reader}\index{double projection}
\index{Paley--Witt}\index{inverse system}

The finite witness of the double projection acquires global scope when the
exceptional reading is applied to every visible block and the gauge is
transported, rather than chosen anew, as the depth increases. This
section constructs that family. The domain remains the enriched TPK
history; the exceptional reading is faithful on its visible stream and does not
purport to replace the memory fields that this stream does not publish.

\subsection{The local operation on the 729 words}

In the Paley--Witt chart one fixes
\[
A_W=
\begin{pmatrix}
0&1&1&1&1&1\\
1&0&1&2&2&1\\
1&1&0&1&2&2\\
1&2&1&0&1&2\\
1&2&2&1&0&1\\
1&1&2&2&1&0
\end{pmatrix}
\in M_6(\mathbb F_3).
\]
Matrix multiplication gives
\[
A_W^{\mathsf T}=A_W,\qquad A_W^2=-I_6.
\]
For every conjugated operator \(A=fA_Wf^{-1}\), define
\[
\Gamma_A:\mathbb F_3^6\longrightarrow\mathbb F_3^{12},
\qquad
\Gamma_A(w)=(w,wA).
\]
Its image is a calibrated copy \(C_W(A)\) of the ternary Witt code.
Projection onto the first six coordinates is a left inverse:
\[
\operatorname{pr}_1\Gamma_A=\operatorname{id}_{\mathbb F_3^6}.
\]
Therefore, \(\Gamma_A\) is total and injective on the \(729\) words and
does not need to distinguish in advance \(\pi\), \(e\), \(\varphi\), or any
other output.

\subsection{Compatible extensions and uniqueness of the initial scale}

Let \(X_n\) be the set of admissible TPK prefixes of depth \(n\), with
truncations
\[
p_{n,m}:X_n\longrightarrow X_m,\qquad m\leq n.
\]
Fix an initial exceptional frame \(f_0\). Since the changes \(g_j\) form
part of the transport stored by the history, a calibrated prefix is the
pair \((x,f_0)\); we write
\[
 X_n^{\rm cal}:=X_n\times\{f_0\},\qquad
 p_{n,m}^{\rm cal}(x,f_0):=(p_{n,m}x,f_0),
\]
and
\[
 X_\infty^{\rm cal}
 :=\varprojlim_n(X_n^{\rm cal},p_{n,m}^{\rm cal}).
\]
In what follows we abbreviate \(p_{n,m}^{\rm cal}\) as \(p_{n,m}\).
The compatible visible emission is
\[
\varepsilon_n:X_n^{\rm cal}\longrightarrow(\mathbb F_3^6)^n,
\qquad
\varepsilon_n(x)=(w_0,\ldots,w_{n-1}),
\]
and satisfies
\[
\operatorname{pr}_{n,m}\varepsilon_n
=\varepsilon_m p_{n,m}.
\]
This application distinguishes its three types:
\[
|\mathcal U_6|=468,\qquad
|\operatorname{im}w_6|=243,\qquad
|\mathbb F_3^6|=729.
\]
The emission \(w_j\) is not identified with the complete state that conserves
leaf, quotient, carry, orientation, phase, and survival.

The exceptional gauge consists of a Witt chart, a dodecad, an
alignment, a lattice origin, and an oriented chamber. It is fixed only
once in the initial block. For
\(x\in X_n^{\rm cal}\), transport of the history supplies
\[
 g_j(x)\in\operatorname{GL}_6(\mathbb F_3),
 \qquad 0\le j<n-1,
\]
and \(g_j(x)\) depends only on the prefix
\(p_{n,j+1}(x)\). The frame is transported by
\[
 f_{j+1}(x)=g_j(x)f_j(x),\qquad
 A_j=f_jA_Wf_j^{-1}.
\]
Thus, the family \((A_j)\) is determined by the history and the initial gauge; it is not a sequence of independent choices.

Let \(\mathcal F_{\rm cal}\subseteq\operatorname{GL}_6(\mathbb F_3)\) be the
set of frames reachable from \(f_0\). The emission preserving the
transported gauge is
\[
 \widetilde\varepsilon_n:X_n^{\rm cal}
 \longrightarrow(\mathbb F_3^6\times\mathcal F_{\rm cal})^n,
 \qquad
 \widetilde\varepsilon_n(x)
 =\bigl((w_0,f_0),\ldots,(w_{n-1},f_{n-1})\bigr).
\]
Prefix dependence implies
\(\operatorname{pr}_{n,m}\widetilde\varepsilon_n
=\widetilde\varepsilon_m p_{n,m}\). Its projection onto the words is
\(\varepsilon_n\); forgetting the frames does not suffice to reconstruct the
exceptional reading.

\subsection{Naturality and compatibility at every depth}

The local operation is covariant under the row-vector convention. For
every change of coordinates \(f\in\operatorname{GL}_6(\mathbb F_3)\),
\[
A^f=fAf^{-1}
\quad\Longrightarrow\quad
\Gamma_{A^f}(wf^{-1})
=\Gamma_A(w)(f^{-1}\oplus f^{-1}).
\]
Change of frame simultaneously transports the word, the operator, and the
atlas. It is not required to remain within the stabilizer of a single
labeled copy.

With the set of reachable frames \(\mathcal F_{\rm cal}\) already defined,
we construct the fixed calibrated atlas
\[
 \mathscr C_W^{\rm cal}
 :=
 \left\{
 \bigl(f,\Gamma_{fA_Wf^{-1}}(w)\bigr):
 f\in\mathcal F_{\rm cal},\ w\in\mathbb F_3^6
 \right\}.
\]
The label \(f\) preserves the calibrated copy in which each word resides.
Define
\[
Q_n(x)=
\bigl(
\bigl(f_0,\Gamma_{A_0}(w_0)\bigr),\ldots,
\bigl(f_{n-1},\Gamma_{A_{n-1}}(w_{n-1})\bigr)
\bigr)\in\bigl(\mathscr C_W^{\rm cal}\bigr)^n.
\]
Let
\[
 r_{n,m}:
 \bigl(\mathscr C_W^{\rm cal}\bigr)^n
 \longrightarrow
 \bigl(\mathscr C_W^{\rm cal}\bigr)^m
\]
the truncation of the last \(n-m\) blocks.

\begin{theorem}[Projective compatibility of the exceptional incidence]
\label{p12-83-c17-s2:thm:lector-excepcional-global-rev6}
\label{thm:lector-excepcional-global-rev6}
For every \(m\leq n\),
\[
\boxed{r_{n,m}\circ Q_n=Q_m\circ p_{n,m}.}
\]
Consequently, there exists a unique application
\[
\boxed{
Q_\infty:X_\infty^{\rm cal}
\longrightarrow
\bigl(\mathscr C_W^{\rm cal}\bigr)^{\mathbb N}}
\]
whose finite projections are the \(Q_n\). The map is faithful on the
visible current.
\end{theorem}

\begin{proof}
Write
\[
\varepsilon_n(x)
=(w_0,\ldots,w_{m-1},w_m,\ldots,w_{n-1}).
\]
The compatibility of \(\varepsilon\) shows that the visible current of \(p_{n,m}(x)\) is \((w_0,\ldots,w_{m-1})\). For \(j<m\), the frame \(f_j\) depends only on the initial gauge and on changes preceding \(j\), all of which are contained in that prefix. The first \(m\) operators \(A_j\) are therefore the same before and after truncation. Applying \(\Gamma_{A_j}\) block by block yields the asserted identity. The universal property of the inverse limit produces \(Q_\infty\). Finally, \(\operatorname{pr}_1\Gamma_{A_j}=I\) recovers every \(w_j\), proving faithfulness on the visible current.
\end{proof}

The last assertion conserves its type. Two histories with the same current
\((w_j)_j\) and distinct hidden memory may have the same
exceptional image. The global reader preserves exactly what each visible block
publishes and the incidence added to that block.

The compatibility of enriched emissions induces a unique application.
\[
 \widetilde\varepsilon_\infty:X_\infty^{\rm cal}
 \longrightarrow
 (\mathbb F_3^6\times\mathcal F_{\rm cal})^{\mathbb N}.
\]
Let \(\operatorname{pr}_w\) the projection onto the corriente of palabras and
\[
 \widetilde\Gamma_\infty
 \bigl((w_j,f_j)_{j\ge0}\bigr)
 :=\bigl(
   (f_j,\Gamma_{f_jA_Wf_j^{-1}}(w_j))
 \bigr)_{j\ge0}.
\]
Then
\(Q_\infty=\widetilde\Gamma_\infty
\circ\widetilde\varepsilon_\infty\), with declared domains and codomains.

\subsection{Bifurcation of a surviving prefix in Archimedean value and exceptional incidence}

Fixed \(m\in\mathbb Z\), let us denote by
\[
 \operatorname{flat}:(\mathbb F_3^6)^{\mathbb N}
 \longrightarrow\mathbb F_3^{\mathbb N}
\]
the concatenation of the blocks in the declared coordinate order, and by
\[
 \operatorname{ev}_3((d_k)_{k\ge1})
 :=\sum_{k\ge1}\frac{d_k}{3^k}\in[0,1],
 \qquad
 \tau_m(t):=m+t,
\]
the ternary evaluation and translation by the integer part. The visible prefix
\((d_1,\ldots,d_{6n})\) determines the cylinder
\[
I_n(x)=
\left[
m+\sum_{k=1}^{6n}\frac{d_k}{3^k},
m+\sum_{k=1}^{6n}\frac{d_k}{3^k}+3^{-6n}
\right].
\]
A prolongation restricts the cylinder and
\(\operatorname{diam}I_n=3^{-6n}\to0\). Therefore, every compatible
infinite history determines a unique real point
\(\mathcal P_{\rm arq}(x)\).

\begin{theorem}[Double global realization]
\label{p12-83-c17-s2:thm:doble-realizacion-global-rev6}
\label{thm:doble-realizacion-global-rev6}
Every compatible section \(x=(x_n)_n\in X_\infty^{\rm cal}\) determines,
from the same prefixes, two compatible outputs:
\[
\mathcal P_{\rm arq}(x)\in\mathbb R,
\qquad
\mathcal P_{\rm exc}(x)=Q_\infty(x)
\in\bigl(\mathscr C_W^{\rm cal}\bigr)^{\mathbb N}.
\]
Truncating before reading is equivalent to truncating each reading after reading. The
first output identifies histories by their limiting value; the second preserves
each visible block and adjoins its Paley--Witt rotation.
\end{theorem}

\begin{proof}
The nesting of the cylinders and the convergence of their diameters prove the
first assertion. The
\cref{p12-83-c17-s2:thm:lector-excepcional-global-rev6} proves the second and the
compatibility with truncations. Both factor through the same emission:
\[
\mathcal P_{\rm arq}
=\tau_m\circ\operatorname{ev}_3\circ\operatorname{flat}
 \circ\operatorname{pr}_w
 \circ\widetilde\varepsilon_\infty,
\qquad
\mathcal P_{\rm exc}
=\widetilde\Gamma_\infty\circ\widetilde\varepsilon_\infty.
\]
\end{proof}

The double projection is thus constructed on admissible histories and a
transported initial gauge. The continuum is the Archimedean evaluation of
the section; the exceptional branch is the compatible realization of its
visible current. From \(C_W\), with the already declared alignments, the
branches toward \(W_{12}\), \(W_{24}\), and toward
\(N(A_2^{12})\), Leech, \(V^\natural\), and the Monster are opened. These branches do not turn
the digits into a set on which the Monster acts directly:
they transport the incidence that evaluation by value ceases to see.

\subsection{Postgenerative coupling with Euler's operator closure}
\label{p12-83-c17-s2:sec:acoplamiento-euler-doble-proyeccion-rev3}

Let \(x_\pi\in X_\infty^{\rm cal}\) be the surviving section of the
\(\pi\) channel. The double realization of
\cref{p12-83-c17-s2:thm:doble-realizacion-global-rev6} produces
\[
 \mathcal P_{\rm arq}(x_\pi)\in\mathbb R,
 \qquad
 \mathcal P_{\rm exc}(x_\pi)=Q_\infty(x_\pi).
\]
The first branch publishes the generated angle; the second preserves the visible
current and realizes the quadratic Paley--Witt relation in each block.

\begin{theorem}[Euler--postgenerative coupling double projection]
\label{p12-83-c17-s2:thm:acoplamiento-euler-doble-proyeccion-rev3}
Suppose that the coinductive construction of the channel of \(\pi\) determines
\[
 \mathcal P_{\rm arq}(x_\pi)=\pi,
\]
and that the exceptional branch uses the gauge transported with
\(A_W^2=-I_6\). Let \(J_{\mathrm e}\) be the real realization of the integral
source of \cref{thm:dos-realizaciones-fuente-cuadratica-rev3}. Then
\begin{equation}
 \boxed{
 \operatorname{Exp}\!\bigl(
   \mathcal P_{\rm arq}(x_\pi)J_{\mathrm e}
 \bigr)+I=0.}
 \label{p12-83-c17-s2:eq:acoplamiento-euler-doble-proyeccion-rev3}
\end{equation}
\end{theorem}

\begin{proof}
By the \cref{thm:dos-realizaciones-fuente-cuadratica-rev3}, \(A_W\) and
\(J_{\mathrm e}\) realizan the relation \(u^2+1=0\) in categorias distinct.
In the branch real,
\[
 \operatorname{Exp}(\theta J_{\mathrm e})
 =\cos\theta\,I+\sin\theta\,J_{\mathrm e}.
\]
Substitution of the previously generated output
\(\theta=\mathcal P_{\rm arq}(x_\pi)=\pi\) produces \(-I\), and the
sum of \(I\) gives the stated identity.
\end{proof}

\paragraph{Causal order and absence of circularity.}
The demonstrated dependence is
\[
 \text{emitter}
 \longrightarrow x_\pi
 \longrightarrow\mathcal P_{\rm arq}(x_\pi)
 \longrightarrow\pi
 \longrightarrow
 \operatorname{Exp}({\pi}J_{\mathrm e})+I.
\]
The exceptional branch fixes the quadratic type; it selects neither the
angle nor its digits. Euler's identity subsequently recognizes the
compatibility of the two realizations and does not enter as an input to the
generator of \(\pi\).


\section{Natural globalization of the Archimedean and exceptional projections over a common inverse system}
\label{p12-83-c17-s3:chap:globalizacion-doble-proyeccion}

The double projection thesis acquires mathematical content when both
readings are defined on the same state system and commute with
depth restriction. Archimedean evaluation contracts a family of
cylinders to a value; exceptional evaluation preserves incidence
relations that the first may forget. These are not two juxtaposed
numerical expressions, but two natural transformations with a common
domain.

\subsection{Compatible systems at finite depth}

Let \((X_n,p_{n+1,n})\), \((A_n,a_{n+1,n})\), and
\((E_n,e_{n+1,n})\) be three inverse systems. The first contains TPK
histories; the second, cylinders or Archimedean approximants; the third,
exceptional incidence data. Suppose that maps are given
\[
 R_n:X_n\longrightarrow A_n,
 \qquad
 Q_n:X_n\times C_n\longrightarrow E_n,
\]
where \(C_n\) gathers the coordinate data of the exceptional reader. The
gauges are also required to form a compatible family and, for every
\(n\), the diagrams must commute
\[
 a_{n+1,n}\circ R_{n+1}=R_n\circ p_{n+1,n},
\]
\[
 e_{n+1,n}\circ Q_{n+1}
 =Q_n\circ(p_{n+1,n}\times c_{n+1,n}).
\]

\begin{theorem}[Globalization of the two evaluations]
\label{p12-83-c17-s3:thm:globalizacion-evaluaciones}
Under the preceding compatibilities, there exist canonical maps
\[
 R_\infty:\varprojlim X_n\longrightarrow\varprojlim A_n,
 \qquad
 Q_\infty:\varprojlim X_n\times\varprojlim C_n
            \longrightarrow\varprojlim E_n.
\]
Its diagonal application
\[
 (R_\infty,Q_\infty):X_\infty\times C_\infty
 \longrightarrow A_\infty\times E_\infty
\]
is unique among the applications whose coordinate projections are
\(R_n\) and \(Q_n\).
\end{theorem}

\begin{proof}
Let \(x=(x_n)_n\in\varprojlim X_n\). The first commutativity equation
implies
\[
 a_{n+1,n}(R_{n+1}(x_{n+1}))=R_n(x_n),
\]
and therefore \((R_n(x_n))_n\) belongs to \(\varprojlim A_n\). We define
\(R_\infty(x)\) by means of that family. The same argument, applied to
\((x_n,c_n)_n\), defines \(Q_\infty\). Uniqueness follows because an element of
an inverse limit is determined by all of its coordinates.
\end{proof}

The theorem identifies the concrete proof required by an HMT realization:
it is not enough to have an Archimedean result at one depth and a code at
another; both must arise from compatible states and respect the truncation
maps.

\subsection{Archimedean contraction of compatible cylinders and realization of the real line}

To each \(x_n\in X_n\) is associated a nonempty compact interval
\(I_n(x_n)\subset\mathbb R\). If
\[
 p_{n+1,n}(x_{n+1})=x_n
 \quad\Longrightarrow\quad
 I_{n+1}(x_{n+1})\subseteq I_n(x_n)
\]
and there exists \(\varepsilon_n\to0\) with
\(\operatorname{diam}I_n(x_n)\leq\varepsilon_n\), the intersection
\[
 \bigcap_{n\geq0}I_n(x_n)
\]
contains a unique real. Thus it is obtained.
\[
 \operatorname{ev}_{\rm arq}:X_\infty^{\rm arq}\longrightarrow\mathbb R.
\]
The phrase ``the cylinders tend to zero'' means here that their diameter
converges to zero, not that the evaluated number tends to zero.

The relation
\[
 x\sim_{\rm arq}y
 \quad\Longleftrightarrow\quad
 \operatorname{ev}_{\rm arq}(x)=\operatorname{ev}_{\rm arq}(y)
\]
describes exactly the information removed by the reader. The real value
is an evaluation class; the history additionally preserves orientation,
carry, sheet, and boundary signature.

\subsection{Exceptional projection by incidence and boundary conservation}

The second evaluation does not take values in \(\mathbb R\). Its codomain is a
category of incidence configurations. In the current finite realization,
the coordinate datum contains:
\[
 \begin{aligned}
 C_{\rm exc}=(\,&\text{projective order},\text{Paley gauge},
 \text{dodecad},\\
 &\text{alignment},\text{lattice origin},\text{chamber}).
 \end{aligned}
\]
After fixing that data, the chain contains two descendants that must remain separated:
\[
 C_W\longrightarrow W_{12}\longrightarrow W_{24},
\]
\[
 C_W\longrightarrow N(A_2^{12})
 \longrightarrow\Lambda_{24}
 \longrightarrow V^\natural
 \longrightarrow\mathbb M.
\]
The first transports designs; the second leads, through later
classical results, to the Leech lattice, the vertex operator algebra, and
the Monster group. The central holonomy and the
Schläfli--\(E_6\) corridor form a parallel boundary extraction; they are not obtained
by replacing one of the two preceding descents.

\subsection{Finite Compatibility of the Common Domain}

The dodecaphonic seal
\[
 K=(234,543,140,729,659,824,621,58,914,794,146,601)
\]
determines the support
\[
 B_K=\{m:K_m\geq729\}=\{4,6,9,10\}.
\]
In the published Witt chart,
\[
 P_\pi=\{2,4,5,6,7,8,9,10,11\},
 \qquad
 H_e=\{1,3,4,6,9,10\},
\]
and, by independent evaluation of the two supports,
\[
 \boxed{B_K=H_e\cap P_\pi.}
\]
Furthermore, for the support of negative carry
\[
 H_\alpha^-=\{4,5,6,7,9,10\}
\]
we have
\[
 B_K=H_e\cap H_\alpha^-,
 \qquad H_\alpha^-=B_K\sqcup\{5,7\}.
\]
These identities constitute a witness of finite compatibility: the same
subset is recovered from the arithmetic seal and from Witt
incidence. This is not an equality of cardinalities, since it identifies specific
positions.

\subsection{Typed separation of the Archimedean and exceptional publications of the common state}

\begin{theorem}[Joint Fidelity Criterion]
Let \(X\) be compact and \(Y,Z\) be Hausdorff spaces. If two
continuous maps \(R:X\to Y\) and \(Q:X\to Z\) satisfy
\[
 x\neq y\Longrightarrow R(x)\neq R(y)\ \text{or}\ Q(x)\neq Q(y),
\]
then \((R,Q):X\to Y\times Z\) is a topological embedding.
\end{theorem}

\begin{proof}
The condition is the injectivity of the diagonal application. Every continuous
injective application from a compact space into a Hausdorff space is a
homeomorphism onto its image.
\end{proof}

At each finite depth, the condition is decidable by enumerating the pairs \(x\neq y\) that remain simultaneously in a fiber of \(R_n\) and of \(Q_n\). Globalization requires compatibility of these controls in passing from \(n+1\) to \(n\). The witness \(B_K\) proves that the two readings already share a nontrivial invariant; joint faithfulness requires control of all fibers and is not replaced by that single witness.

\subsection{Precise formulation of the thesis on the real line}

The enriched object associated to the Archimedean reading is the arrow
\[
 \mathbb R_{\mathfrak H}:=
 \bigl(X_\infty^{\rm arq}
 \xrightarrow{\operatorname{ev}_{\rm arq}}
 \operatorname{Im}(\operatorname{ev}_{\rm arq})\subseteq\mathbb R\bigr).
\]
The domain preserves genealogies; the image preserves values. Coverage of
the cylinders at every scale and compatible prolongability are proved in
the construction of the Archimedean reader
\eqref{eq:continuo-lector-arquimediano}: for every compact set
\(K\subset\mathbb R\),
\(\operatorname{ev}_{\mathbb R}(\Omega_{\infty,K})=K\), and the directed union
of those intervals covers \(\mathbb R\). Therefore,
\[
 \operatorname{Im}(\operatorname{ev}_{\rm arq})=\mathbb R,
 \qquad
 X_\infty^{\rm arq}/\!\sim_{\mathbb R}\ \cong\mathbb R.
\]
The real field does not act as the starting point of the construction, but as
the Archimedean quotient of a space of sections with memory. The descent of
sum and product is verified operation by operation and does not condition this
already proved surjectivity.

The double reading is then expressed by
\[
 X_\infty^{\rm com}\times C_\infty
 \xrightarrow{\ (\operatorname{ev}_{\rm arq},Q_\infty)\ }
 \mathbb R\times E_\infty.
\]
This formula is the mathematical translation of the conceptual assertion: one
and the same history may contract to a value and, in the transverse reading
of its boundary, preserve the incidence leading to the exceptional
symmetries.


% Fuente normativa activa de la obstrucción de factorización.
\section{Dependence of the incidence reader on the integral holonomic state}
\label{p12-83-c17-s4:sec:obstruccion-lector-excepcional-minimo}

Finite compatibility of the two evaluations does not imply that every
incidence reader refines the partition produced by the visible reader. The
regional catalog makes it possible to decide this question exactly for the
minimal Paley--Witt realization. Let \(\mathcal U_{468}\) be the set of
the 468 published regional states, and let
\[
 \Psi:\mathcal U_{468}\longrightarrow\mathcal W_{243}\subset\mathbb F_3^6,
 \qquad
 \Psi(U_1,\ldots,U_6)=(-U_1,\ldots,-U_6)\pmod 3.
\]
For \(w\in\mathbb F_3^6\), the map
\[
 c(w)=(w,wA_W)\in C_W\subset\mathbb F_3^{12}
\]
is the ternary encoder fixed in the Paley--Witt section. We define
the minimal incidence reader by
\[
 Q_{\min}=\operatorname{Supp}\circ c\circ\Psi:
 \mathcal U_{468}\longrightarrow\mathcal P([12]).
\]

For a map \(f:X\to Y\), denote by
\[
 \operatorname{Eq}(f)=\{(x,x')\in X^2:f(x)=f(x')\}
\]
its fiberwise equivalence relation.

\begin{theorem}[Obstruction of the minimal exceptional reader]
\label{p12-83-c17-s4:thm:obstruccion-lector-excepcional-minimo}
In the complete catalogue of 468 states one verifies
\[
 \operatorname{Eq}(\Psi,Q_{\min})=\operatorname{Eq}(\Psi).
\]
In the usual abbreviated notation for fiber partitions,
\[
 \ker(\Psi,Q_{\min})=\ker\Psi.
\]
Consequently, \((\Psi,Q_{\min})\) is not injective and the minimal
incidence reader does not distinguish any pair of states that has already been
identified by \(\Psi\).
\end{theorem}

\begin{proof}
The factorization \(Q_{\min}=S\circ\Psi\), with
\(S=\operatorname{Supp}\circ c\), implies
\[
 \Psi(u)=\Psi(v)\Longrightarrow Q_{\min}(u)=Q_{\min}(v).
\]
The converse implication between the equivalence relations is immediate,
because \(\Psi\) is the first component of the diagonal map. Equality of the
relations therefore follows without statistical assumptions. Exhaustive
enumeration further confirms that both partitions have the same
histogram:
\[
 132\,[1]+66\,[2]+18\,[3]+3\,[4]+6\,[5]+18\,[6],
\]
where \(a\,[m]\) means that there exist \(a\) fibres of cardinality \(m\).
\end{proof}

The microfiber of \(w_\pi=010211\) provides the distinguished witness. Its
five regions are
\[
\begin{split}
 &(501,614,498,169,272,272),\quad
 (810,923,870,169,272,272),\\
 &(810,983,810,169,272,272),\quad
 (870,923,810,169,272,272),\\
 &(870,923,810,418,674,521).
\end{split}
\]
The five have the same visible image and the same exceptional support.
\[
 P_\pi=\{2,4,5,6,7,8,9,10,11\}.
\]
At the level of the 243 words there appear 213 distinct supports: 185 have
one preimage, 27 have two, and one has four. Taking support therefore introduces
a second quotient in addition to the one imposed by \(\Psi\).

The theorem delimits the minimal reader, not the enriched exceptional
evaluation. A reader capable of separating the regions of the same word
must act before the visible quotient and use at least one of the data
that it eliminates: orientation, sheet, signature, return, or boundary state. The
dodecaphase intertwiner developed in the following part belongs to that
enriched class and does not contradict the preceding obstruction.

Complete enumeration of the fibers produces the preceding histogram and
proves that taking support introduces a second quotient.


\section{Archimedean and incidental publication}

For each level, let \(\mathcal E_n\) be an incidence object and let
\begin{equation}
 e_n:S_n\longrightarrow\mathcal E_n,
 \qquad
 \sigma_{n,m}:\mathcal E_n\longrightarrow\mathcal E_m
 \label{p12-83-c17-s5:eq:datos-lector-incidencial-u024}
\end{equation}
applications that satisfy
\begin{equation}
 \sigma_{n,m}\circ e_n=e_m\circ r_{n,m}.
 \label{p12-83-c17-s5:eq:naturalidad-lector-incidencial-u024}
\end{equation}
The universal property constructs
\begin{equation}
 \operatorname{ev}_{\rm inc}:\Omega_\infty
 \longrightarrow\mathcal E_\infty,
 \qquad
 \mathcal E_\infty:=\varprojlim_n\mathcal E_n.
 \label{p12-83-c17-s5:eq:evaluacion-incidencial-limite-u024}
\end{equation}
The double publication is
\begin{equation}
 \mathcal P_\infty
 :=(\operatorname{ev}_{\mathbb R},\operatorname{ev}_{\rm inc}):
 \Omega_\infty\longrightarrow\mathbb R\times\mathcal E_\infty.
 \label{p12-83-c17-s5:eq:doble-publicacion-continuo-u024}
\end{equation}

\begin{figure}[H]
\centering
\begin{tikzcd}[column sep=10em,row sep=5em,
  cells={nodes={font=\large}}]
&\Omega_\infty
 \arrow[dl,"\operatorname{ev}_{\mathbb R}"']
 \arrow[dr,"\operatorname{ev}_{\rm inc}"]&\\
\mathbb R
 \arrow[dr,"\operatorname{coord}"']&&
\mathcal E_\infty
 \arrow[dl,"\operatorname{pub}_{\rm inc}"]\\
&\mathcal O&
\end{tikzcd}
\caption{Two publications of the same history. Commutativity in
\(\mathcal O\) does not identify the codomains \(\mathbb R\) and
\(\mathcal E_\infty\).}
\label{p12-83-c17-s5:fig:doble-publicacion-continuo-u024}
\end{figure}

If it is proven
\begin{equation}
 \operatorname{pub}_{\rm inc}\circ\operatorname{ev}_{\rm inc}
 =\operatorname{coord}\circ\operatorname{ev}_{\mathbb R},
 \label{p12-83-c17-s5:eq:compatibilidad-doble-publicacion-u024}
\end{equation}
a typed comparison of the two publications is obtained. The identity does not
assert that one branch is constructed from the other. Both start from
\(\Omega_\infty\).


\section{Gauge of the exceptional publication}

The second projection is not a unary reader devoid of incidence data. Let
\[
 \mathscr C_{\rm exc}
 =\mathscr C_{\rm ord}\times\mathscr C_{\rm Paley}
  \times\mathscr C_{\rm dod}\times\mathscr C_{\rm inc}
  \times\mathscr C_{\rm ret}
\]
the space of exceptional gauges, whose coordinates respectively fix
the projective order, the Paley gauge, the dodecad, the incidence
alignment, and the lattice chamber. One fixes
\[
 \chi_{\rm exc}
 =(\chi_{\rm ord},\chi_{\rm Paley},\chi_{\rm dod},
   \chi_{\rm inc},\chi_{\rm ret})\in\mathscr C_{\rm exc}.
\]
The typed reader has domain
\[
 \mathfrak E:
 X_\infty^{\rm cal}\times\mathscr C_{\rm exc}
 \longrightarrow\mathbf{Exc},
 \qquad
 \mathfrak E_{\chi_{\rm exc}}(u):=\mathfrak E(u,\chi_{\rm exc}).
\]
Henceforth, every abbreviation \(\mathfrak E(u)\) means
\(\mathfrak E_{\chi_{\rm exc}}(u)\) for this explicitly fixed gauge.
The dependence on \(\chi_{\rm exc}\) does not affect the existence of the
continuum or its Archimedean publication: it types the exceptional conservation
of incidence once the common state has been constructed.

\begin{falsifier}[Loss of the exceptional gauge]
The comparison becomes incomplete if the projective order, the dodecad,
the incidence alignment, or the lattice chamber is changed without transporting the remaining
gauge coordinates, or if \(\mathfrak E\) is used as a unary reader before
fixing \(\chi_{\rm exc}\). Such a failure invalidates that exceptional comparison;
it does not modify the Archimedean quotient or regenerate the continuum.
\end{falsifier}


The second reading does not arise from the real number. Let
\[
 \mathfrak E:
 \mathscr O_{\mathbb R,\leftrightarrow}^{\rm HMT}
 \longrightarrow\mathbf{Exc}
\]
the continuous incidence reader, with Hausdorff codomain, that prolongs
the Paley--Witt, Golay,
Leech, \(V^\natural\), Monster, and Moonshine chain. Both maps begin from
the same history:
\begin{equation}
 \mathfrak D(u)=(\mathfrak R(u),\mathfrak E(u)).
 \label{p12-83-c17-s7:hmtch:eq-doble-proyeccion-two-way}
\end{equation}
The first contracts the compatible fibres to their Archimedean mean; the
second preserves incidence, orientation, and boundary memory.

Fixing the origin of the dodecaphonic calendar, let it be
\[
 R_+:=\{k\in\mathbb Z:1+(k\bmod 9)\in\{1,4,7\}\}
\]
the bi-infinite set of complete returns \((\tau=+1)\). It is syndetic:
its gaps are uniformly bounded. For
\(u=(m,(x_k)_{k\in\mathbb Z})\), define
\[
 \operatorname{Ret}_+(u)
 :=\bigl(m,(x_k)_{k\in R_+}\bigr)
 \in\mathbb Z\times\prod_{k\in R_+}X,
\]
but each \(x_k\) here is the whole state
\eqref{p12-83-c19-s1:hmtch:eq-celula-trit-completa} with leaf, carry, route, and incidence,
not the isolated sign.

\begin{proposition}[Lossless Sampling of the Return Genealogical Record]
\label{p12-83-c17-s7:hmtch:prop-retorno-reconstruye}
The map \(\operatorname{Ret}_+\), with the subspace topology in the
product, is a homeomorphism onto its image.
Consequently, there is a unique continuous map
\(\mathfrak E_+\) on that image such that
\begin{equation}
 \mathfrak E=\mathfrak E_+\circ\operatorname{Ret}_+.
 \label{p12-83-c17-s7:hmtch:eq-factorizacion-excepcional-retorno}
\end{equation}
\end{proposition}

\begin{proof}
Between two consecutive returns there is a uniformly bounded number of steps
of the calendar. The complete state at a return contains the transition and
the genealogical register that determine the intermediate steps; in the other direction, the
inverse arrow of the groupoid completion is used. The entire
orbit is thereby reconstructed. The map is continuous and, when restricted to a
compact chart, a continuous injection into a Hausdorff space is a
homeomorphism onto its image. The integer coordinate is preserved
explicitly. The factorization is defined by transporting \(\mathfrak E\) through
the inverse and is continuous chart by chart.
\end{proof}

For this reason, \(+1\) participates in the exceptional corridor as a complete memory
of return: sampling the integer states does not lose the exceptional output.
This proposition proves recovery and factorization, not a generation of
the corridor from the isolated scalar \(+1\). The Archimedean output and the
exceptional output are two coordinated resultants of the same antecedent.
 
\chapter{Archimedean reconstruction through half-closed cylinders with memory}
\label{chap:83-18}

The inverse limit \(\Omega_\infty\) preserves the complete history. To
obtain a real coordinate it is necessary to add a compatible family of
intervals and prove three different facts: each history determines a
unique value; every value of the declared compactum is reached; and the finite
measures descend consistently. None of these facts follows
from the mere growth in the number of states.

The Archimedean publication likewise preserves an approximation of the
identity by cylindrical conditional expectations. The concentration of unit
mass on compatible prefixes provides the exact passage from the
nonadic cylinders to the Dirac measure, without identifying density,
support measure, and total mass.

\begin{HMTScientificOwnerSubordinated}
\chapter{Nonadic cylinders, identity and Dirac kernel}

\section{Inverse system of compatible nonadric cylinders with genealogical memory}

Let
\[
X=(\Znine)^{\mathbb N}
\]
with the product topology and the uniform measure \(\mu\). For
\(x=(x_1,x_2,\ldots)\), the cylinder of depth \(n\) is
\[
C_n(x)=\{y\in X:y_j=x_j\text{ for }1\le j\le n\}.
\]
Cylinders are simultaneously open and closed, form a base and
satisfy
\[
\mu(C_n(x))=9^{-n},\qquad
C_{n+1}(x)\subset C_n(x),\qquad
\bigcap_n C_n(x)=\{x\}.
\]
Adding a digit does not create a new number: it restricts the cylinder of
admissible prolongations. This statement gives topological content to the
genealogical thesis of number.

\section{Cylindrical concentration of the identity and construction of the Dirac kernel}

\begin{theorem}[Unitary nonadic concentration]
Define
\[
f_{n,x}=9^n\mathbf1_{C_n(x)}.
\]
Then
\[
\mu(C_n(x))\to0,\qquad f_{n,x}(x)\to\infty,\qquad
\int_Xf_{n,x}\,\dd\mu=1,
\]
and the measures \(f_{n,x}\mu\) converge weak-* to \(\delta_x\).
\end{theorem}

\begin{proof}
The first identity is immediate. The height of \(f_{n,x}\) over its support
is \(9^n\), so it diverges. The integral equals
\(9^n9^{-n}=1\). For \(g\in C(X)\), the oscillation of \(g\) in \(C_n(x)\)
tends to zero by uniform continuity; therefore
\[
\left|\int g f_{n,x}\,\dd\mu-g(x)\right|
\le \sup_{y\in C_n(x)}|g(y)-g(x)|\longrightarrow0.
\]
This is weak-* convergence to \(\delta_x\).
\end{proof}

The theorem does not assert \(1/0=\infty\). It exhibits three evaluations of the same
object: support measure, density and total mass. In blocks of three digits,
\[
\mu(C_{3m})=729^{-m},\qquad
f_{3m}=729^m\mathbf1_{C_{3m}},\qquad
729^{-m}729^m=1.
\]
The evolutionary conservation of unity here assumes an exact analytical form.

\section{Convergence of the Kronecker Kernel to the Dirac Distribution}

At the finite level \(X_n=(\Znine)^n\), let
\[
K_n(x,y)=9^n\mathbf1_{\{x=y\}}.
\]
With the uniform measure, the operator
\[
(E_nf)(x)=\sum_{y\in X_n}K_n(x,y)f(y)9^{-n}
\]
is the identity. In the cylindrical limit, the condition \(x=y\) is replaced by
“same prefix of length \(n\)” and \(E_n\) is the conditional expectation on
the cylindrical algebra. One has
\[
E_n^2=E_n,\qquad E_n^*=E_n,\qquad E_nf\to f
\quad\text{in }L^2(X,\mu).
\]

The exact bridge is, therefore,
\[
\delta_{ij}
\longrightarrow K_n(x,y)
\longrightarrow\delta_x.
\]
The Kronecker delta is a matrix kernel; the Dirac delta is a measure or
distribution. Neither is the Dirac operator
\(D=\gamma^\mu\partial_\mu\). The transition to spinors and Clifford requires
a second intertwiner and will be treated later.

\section{Scaling Law, Effective Dimension, and Density of the Cylindrical Measure}

Suppose that a contraction of ratio \(\lambda\) acts independently
in \(d\) directions and that the product measure scales as \(\lambda^d\). To
preserve unit mass, the density must scale as
\[
\rho\longmapsto\lambda^{-d}\rho.
\]
The exponents \((-1,-2,-3)\) then appear as density weights on a line,
area, and volume. This proposition is general and exact under the product
hypothesis. It does not permit \(\lambda\) to be identified automatically with
\(\alpha\); that identification requires a metric chart.

\section{Nonadic Monodromy and Computability of the Archimedean Evaluation}

The existence of an infinite branch in a finitely branching tree may
follow from König's lemma, but König's lemma does not furnish a computable selector. There
are infinite recursive binary trees with no recursive branch. The joint nonadic
connection does not rest on that abstract existence: it has a certified uniform
selector that generates to arbitrary finite depth and opens the decimal witnesses
only after generation.

\begin{falsadorBox}[title={Three prohibited confusions}]
\(1/0\) shall not be written; the Dirac delta shall not be represented as an ordinary
function with “an infinite one”; König shall not be used as an algorithm. The strong
result is the exact compatibility among cylinders, unit density, and the weak
limit, together with the constructive selector independent of decimal comparison.
\end{falsadorBox}
 \end{HMTScientificOwnerSubordinated}

\section{System of intervals associated with finite states}

Let \((S_n,r_{n,m})\) be the inverse system of the
\cref{chap:sistema-inverso-continuidad-genealogica}. For each
\(a\in S_n\), assign a nonempty compact interval
\begin{equation}
 I_n(a)=[\ell_n(a),u_n(a)]\subset\mathbb R.
 \label{p12-83-c18-s1:eq:intervalo-asociado-estado}
\end{equation}
The assignment is \emph{compatible with refinement} when
\begin{equation}
 r_{n+1,n}(b)=a
 \quad\Longrightarrow\quad
 I_{n+1}(b)\subseteq I_n(a).
 \label{p12-83-c18-s1:eq:anidamiento-intervalos-estados}
\end{equation}
It is \emph{uniformly contractive} when there exists a sequence
\(\delta_n\downarrow0\) such that
\begin{equation}
 \operatorname{diam}I_n(a)\leq\delta_n
 \qquad(a\in S_n).
 \label{p12-83-c18-s1:eq:contraccion-uniforme-intervalos}
\end{equation}

For a history \(x=(x_n)_n\), the inclusions
\eqref{p12-83-c18-s1:eq:anidamiento-intervalos-estados} produce a chain
\begin{equation}
 I_0(x_0)\supseteq I_1(x_1)\supseteq I_2(x_2)\supseteq\cdots.
 \label{p12-83-c18-s1:eq:cadena-intervalos-historia}
\end{equation}

\begin{theorem}[Evaluation of a history]
\label{p12-83-c18-s1:thm:evaluacion-unica-historia-u024}
If the assignment is compatible and uniformly contractive, then for every \(x\in\Omega_\infty\) there exists a unique real number
\begin{equation}
 \operatorname{ev}_{\mathbb R}(x)
 \in\bigcap_{n\geq0}I_n(x_n).
 \label{p12-83-c18-s1:eq:evaluacion-real-historia-u024}
\end{equation}
The map
\begin{equation}
 \operatorname{ev}_{\mathbb R}:\Omega_\infty\longrightarrow\mathbb R
 \label{p12-83-c18-s1:eq:aplicacion-evaluacion-real-u024}
\end{equation}
is continuous.
\end{theorem}

\begin{proof}
The intervals of \eqref{p12-83-c18-s1:eq:cadena-intervalos-historia} are compact,
nonempty and nested. The nested-interval theorem guarantees that
their intersection is nonempty. If \(s,t\) belong to it, then
\[
 |s-t|\leq\operatorname{diam}I_n(x_n)\leq\delta_n
\]
for every \(n\). Setting \(n\to\infty\) yields \(s=t\).

For continuity, let \(\varepsilon>0\) and choose \(n\) with
\(\delta_n<\varepsilon\). If \(y\in[x_n]_n\), then
\(I_n(y_n)=I_n(x_n)\). The two evaluations belong to that interval and,
therefore,
\[
 |\operatorname{ev}_{\mathbb R}(x)
  -\operatorname{ev}_{\mathbb R}(y)|<\varepsilon.
\]
The cylinder \([x_n]_n\) is an open neighborhood of \(x\).
\end{proof}

The construction does not identify \(x\) with
\(\operatorname{ev}_{\mathbb R}(x)\). A single coordinate may have
several histories, and a history contains more data than any point on
the line.

\section{Coherent coverage of the visible compactum}

Let \(K\subset\mathbb R\) a compacto not vacio. Moreover of anidamiento and
contraccion, exigimos
\begin{align}
 K&=\bigcup_{a\in S_0}I_0(a),
 \label{p12-83-c18-s1:eq:cobertura-inicial-compacto-u024}\\
 I_n(a)&=
 \bigcup_{\substack{b\in S_{n+1}\\r_{n+1,n}(b)=a}}
 I_{n+1}(b)
 \qquad(a\in S_n).
 \label{p12-83-c18-s1:eq:cobertura-extensiones-compatible-u024}
\end{align}
The second equality is an additional geometric condition. The
surjectivity of \(r_{n+1,n}\) only guarantees the existence of states
above \(a\); it does not guarantee that their intervals cover all of
\(I_n(a)\).

\begin{theorem}[Surjectivity of the evaluation]
\label{p12-83-c18-s1:thm:sobreyectividad-evaluacion-compacto}
If
\eqref{p12-83-c18-s1:eq:cobertura-inicial-compacto-u024} and
\eqref{p12-83-c18-s1:eq:cobertura-extensiones-compatible-u024} hold, then
\begin{equation}
 \operatorname{ev}_{\mathbb R}(\Omega_\infty)=K.
 \label{p12-83-c18-s1:eq:imagen-evaluacion-compacto}
\end{equation}
\end{theorem}

\begin{proof}
The left-to-right inclusion follows from
\(I_n(x_n)\subseteq I_0(x_0)\subseteq K\).

Let us fix \(t\in K\). Let us define
\begin{equation}
 T_n(t):=\{a\in S_n:t\in I_n(a)\}.
 \label{p12-83-c18-s1:eq:estados-que-cubren-punto}
\end{equation}
The initial cover makes \(T_0(t)\) nonempty. If \(a\in T_n(t)\), the identity \eqref{p12-83-c18-s1:eq:cobertura-extensiones-compatible-u024} provides at least one \(b\in T_{n+1}(t)\) with \(r_{n+1,n}(b)=a\). The graph of these extensions is finitely branching and has nonempty levels at every depth. There therefore exists a compatible chain \(x=(x_n)_n\) with \(t\in I_n(x_n)\) for every \(n\). Uniqueness of the intersection forces \(\operatorname{ev}_{\mathbb R}(x)=t\).
\end{proof}

\section{Archimedean genealogical quotient and homeomorphism with \texorpdfstring{\(\mathbb R\)}{R}}

We define the relation
\begin{equation}
 x\sim_{\mathbb R}y
 \quad\Longleftrightarrow\quad
 \operatorname{ev}_{\mathbb R}(x)
 =\operatorname{ev}_{\mathbb R}(y).
 \label{p12-83-c18-s1:eq:relacion-equivalencia-arquimediana}
\end{equation}
Let
\begin{equation}
 q_{\mathbb R}:\Omega_\infty
 \longrightarrow\Omega_\infty/\!\sim_{\mathbb R}
 \label{p12-83-c18-s1:eq:proyeccion-cociente-arquimediano}
\end{equation}
the canonical projection.

\begin{theorem}[Reconstruction of the compact as a quotient]
\label{p12-83-c18-s1:thm:homeomorfismo-cociente-compacto-u024}
Under the preceding hypotheses, there exists a unique homeomorphism
\begin{equation}
 \overline{\operatorname{ev}}_{\mathbb R}:
 \Omega_\infty/\!\sim_{\mathbb R}
 \xrightarrow{\ \cong\ }K
 \label{p12-83-c18-s1:eq:homeomorfismo-cociente-compacto-u024}
\end{equation}
such that
\begin{equation}
 \operatorname{ev}_{\mathbb R}
 =\overline{\operatorname{ev}}_{\mathbb R}\circ q_{\mathbb R}.
 \label{p12-83-c18-s1:eq:factorizacion-evaluacion-cociente}
\end{equation}
The fiber over \(t\in K\) is
\begin{equation}
 \mathfrak G_t
 :=\operatorname{ev}_{\mathbb R}^{-1}(t),
 \label{p12-83-c18-s1:eq:fibra-genealogias-mismo-valor}
\end{equation}
the complete set of genealogies that publish \(t\).
\end{theorem}

\begin{proof}
The map \(\operatorname{ev}_{\mathbb R}\) is continuous, surjective, and
constant exactly on the classes of \(\sim_{\mathbb R}\). By the
universal property of the quotient, it induces a continuous bijective map
\(\overline{\operatorname{ev}}_{\mathbb R}\). The domain is compact because it is a continuous
image of the compact space \(\Omega_\infty\), and \(K\) is Hausdorff.
A continuous bijection from a compact space to a Hausdorff space is a homeomorphism.
\end{proof}

The theorem proves two things and keeps them separate:
\begin{equation}
 \boxed{
 \begin{array}{c}
 \text{the history space is profinite and retains memory,}\\
 \text{its quotient by equality of evaluation is the visible compact space.}
 \end{array}}
 \label{p12-83-c18-s1:eq:dos-niveles-continuo-u024}
\end{equation}
It is not asserted that \(\Omega_\infty\) is homeomorphic to \(K\).

\section{Genealogical Atlas of the Real Line}

For \(m\geq1\), let \(K_m=[-m,m]\). Suppose that for each \(m\) a
system \(\Omega_\infty^{(m)}\) is given that reconstructs \(K_m\), together with
closed inclusions
\begin{equation}
 j_m:\Omega_\infty^{(m)}\hookrightarrow
 \Omega_\infty^{(m+1)}
 \label{p12-83-c18-s1:eq:inclusion-atlas-genealogico}
\end{equation}
which satisfy
\begin{equation}
 \operatorname{ev}_{m+1}\circ j_m
 =\iota_m\circ\operatorname{ev}_m,
 \label{p12-83-c18-s1:eq:compatibilidad-atlas-real}
\end{equation}
where \(\iota_m:K_m\hookrightarrow K_{m+1}\) is the inclusion ordinary.

\begin{proposition}[Exhaustion of the line]
\label{p12-83-c18-s1:prop:agotamiento-recta-real-u024}
The Archimedean quotient of the strict colimit
\begin{equation}
 \Omega_\infty^{\mathbb R}
 :=\varinjlim_m\Omega_\infty^{(m)}
 \label{p12-83-c18-s1:eq:colimite-historias-recta}
\end{equation}
is
\begin{equation}
 \varinjlim_m[-m,m]
 =\bigcup_{m\geq1}[-m,m]
 =\mathbb R
 \label{p12-83-c18-s1:eq:colimite-compactos-recta}
\end{equation}
with its usual topology.
\end{proposition}

\begin{proof}
Each piece has quotient \(K_m\) by the
\cref{p12-83-c18-s1:thm:homeomorfismo-cociente-compacto-u024};
\eqref{p12-83-c18-s1:eq:compatibilidad-atlas-real} intertwines the quotients. The final topology
of the exhaustion by compact sets coincides with the usual topology of
\(\mathbb R\): a subset \(U\subseteq\mathbb R\) is open if and only
if \(U\cap[-m,m]\) is open in \([-m,m]\) for every \(m\).
\end{proof}

Passage from a bounded interval to the line does not erase the fibers
\(\mathfrak G_t\). The inclusions \(j_m\) preserve the history and only
enlarge the available visible chart.

\section{Balanced ternary and block-decimal charts}

The balanced ternary representation of an integer has the form
\begin{equation}
 n=\sum_{j=0}^{J}\epsilon_j3^j,
 \qquad
 \epsilon_j\in\{-1,0,+1\},
 \label{p12-83-c18-s1:eq:carta-ternaria-balanceada-u024}
\end{equation}
with a fixed carry rule. This chart preserves the local sign and the
quotient of the reduction; it must not be confused with the nonnegative word in
\(\mathbb F_3\).

The block decimal chart utilizes
\begin{equation}
 D_N=1000D_{N-1}+u_N,
 \qquad
 u_N\in\{0,\ldots,999\},
 \qquad D_0=0.
 \label{p12-83-c18-s1:eq:concatenacion-base-mil-u024}
\end{equation}
The block \(u_N\), its class modulo nine, and the carry are distinct fields.
Since \(1000\equiv1\pmod9\),
\begin{equation}
 D_N\equiv\sum_{j=1}^{N}u_j\pmod9.
 \label{p12-83-c18-s1:eq:base-mil-compatible-mod9-u024}
\end{equation}
This congruence does not recover the order of the blocks; the sequence
\((u_1,\ldots,u_N)\) remains in the record.

For \(M\geq1\), the cells decimales half-open are
\begin{equation}
 \mathcal D_{M,q}
 :=[q10^{-M},(q+1)10^{-M}),
 \qquad q\in\mathbb Z.
 \label{p12-83-c18-s1:eq:celda-decimal-semiabierta-u024}
\end{equation}
A half-open interval \([\ell,h)\) is contained in a unique cell if and only if
\begin{equation}
 \left\lfloor10^M\ell\right\rfloor
 =\left\lceil10^Mh\right\rceil-1.
 \label{p12-83-c18-s1:eq:criterio-exacto-celda-decimal-u024}
\end{equation}
The formula with a ceiling at the right endpoint also covers the case in which
\(h\) coincides exactly with a decimal boundary.

\begin{definition}[Truncation and Rounding]
For \(x\geq0\), define
\begin{align}
 \operatorname{Tr}_M(x)
 &:=10^{-M}\left\lfloor10^Mx\right\rfloor,
 \label{p12-83-c18-s1:eq:truncamiento-decimal-u024}\\
 \operatorname{Rd}_M(x)
 &:=10^{-M}\left\lfloor10^Mx+\frac12\right\rfloor.
 \label{p12-83-c18-s1:eq:redondeo-decimal-u024}
\end{align}
\end{definition}
These operations coincide except when the discarded fraction crosses the
half-unit threshold at the final position. The distinction is semantic and
arithmetic; one cannot be substituted for the other to force prefix
compatibility.

\section{Descent of operations along evaluation fibers}

Let \(D_\star\subseteq\Omega_\infty\times\Omega_\infty\) the domain of to
operation parcial continuous
\begin{equation}
 \star:D_\star\longrightarrow\Omega_\infty.
 \label{p12-83-c18-s1:eq:operacion-historias-u024}
\end{equation}
Assume that \(D_\star\) is saturated by the fibers of
\(\operatorname{ev}_{\mathbb R}\times\operatorname{ev}_{\mathbb R}\).

\begin{theorem}[Exact Descent Criterion]
\label{p12-83-c18-s1:thm:criterio-descenso-operaciones-u024}
There exists a unique visible operation
\begin{equation}
 \overline\star:
 (\operatorname{ev}_{\mathbb R}\times\operatorname{ev}_{\mathbb R})
 (D_\star)\longrightarrow K
 \label{p12-83-c18-s1:eq:operacion-descendida-visible}
\end{equation}
which satisfies
\begin{equation}
 \operatorname{ev}_{\mathbb R}(x\star y)
 =\operatorname{ev}_{\mathbb R}(x)
  \ \overline\star\ 
  \operatorname{ev}_{\mathbb R}(y)
 \label{p12-83-c18-s1:eq:entrelazamiento-operacion-descendida}
\end{equation}
if and only if
\begin{equation}
 \begin{split}
 &\operatorname{ev}_{\mathbb R}(x)=
  \operatorname{ev}_{\mathbb R}(x'),\\
 &\operatorname{ev}_{\mathbb R}(y)=
  \operatorname{ev}_{\mathbb R}(y')
 \end{split}
 \quad\Longrightarrow\quad
 \operatorname{ev}_{\mathbb R}(x\star y)=
 \operatorname{ev}_{\mathbb R}(x'\star y')
 \label{p12-83-c18-s1:eq:constancia-operacion-sobre-fibras}
\end{equation}
for all pairs in the domain.
\end{theorem}

\begin{proof}
If \(\overline\star\) exists, the right-hand side of
\eqref{p12-83-c18-s1:eq:entrelazamiento-operacion-descendida} depends only on the two
evaluations, and
\eqref{p12-83-c18-s1:eq:constancia-operacion-sobre-fibras} follows. Conversely, this
constancy makes it possible to define
\[
 s\ \overline\star\ t
 :=\operatorname{ev}_{\mathbb R}(x\star y)
\]
for any representatives \(x,y\) with evaluations \(s,t\). The
condition proves that the definition is independent of the representatives.
Uniqueness is immediate.
\end{proof}

The additive and multiplicative sheets of APP supply operations on a
common support, but their carries and their memories are not identified. The
preceding theorem allows sum and product to descend to the visible chart
without requiring their genealogical realizations to coincide.

\section{Correspondence between discrete cardinality, cylindrical measure, and Archimedean publication}

At each depth three operations appear that must be kept separated:
\begin{align}
 \operatorname{Count}_n(A)&:=|A|,
 \label{p12-83-c18-s1:eq:operacion-conteo-u024}\\
 \operatorname{Measure}_n(A)&:=\mu_n(A),
 \label{p12-83-c18-s1:eq:operacion-medida-u024}\\
 \operatorname{Publish}_n(a)&:=L_n(a).
 \label{p12-83-c18-s1:eq:operacion-publicacion-u024}
\end{align}
The count depends on the number of states; the measure depends on compatible
weights; the publication depends on the reader. Two fibers of equal
cardinality can have distinct weights, and two states of equal publication
can belong to distinct cylinders. Equality of one of these three
outputs does not imply equality of the other two.

\section{Obstructions to Archimedean evaluation: nesting, continuity, quotient, projective measures, and descent of operations}

\begin{criterion}[Refutation of the evaluation and the measure]
The construction is refuted if any of the following facts
holds:
\begin{enumerate}
 \item a compatible extension receives an interval that is not contained
       in the interval of its restriction;
 \item the diameters do not tend to zero, but uniqueness of the
       intersection is asserted;
 \item the evaluation is not continuous in the cylinder topology;
 \item it is asserted that every point of \(K\) appears without proving the coherent
       covering \eqref{p12-83-c18-s1:eq:cobertura-extensiones-compatible-u024};
 \item the map induced by the quotient is not bijective;
 \item the profinite space is identified with its visible quotient;
 \item the inclusions of the atlas do not intertwine their evaluations;
 \item \(\lfloor10^Mh\rfloor\) is used instead of
       \(\lceil10^Mh\rceil-1\) for a semi-open interval, and the
       boundary case is lost;
 \item truncation and rounding are conflated;
 \item a descended operation depends on the representative chosen in a
       fibre;
 \item to family \((\mu_n)_n\) viola
       \eqref{eq:consistencia-proyectiva-medidas-u024};
 \item the measure of a cylinder depends on the level to which it is refined;
 \item a conditioned measure ceases to sum to one or loses consistency;
 \item the visible measure \(\nu\) is identified with the internal measures
       \(\mu_t\);
 \item an arrow between \(\mathbb R\) and
       \(\mathcal E_\infty\) is inferred solely because both publications share an
       antecedent;
 \item the count of a fibre is used as though it determined its measure or its
       published value.
\end{enumerate}
\end{criterion}


\section{From the space of histories to the real line}

Let
\[
 \Psi:\mathbb Z_3^6\xrightarrow{\ \sim\ }
       (\mathbb F_3^6)^{\mathbb N}
\]
the Hensel conjugacy between the initial memory and the complete tape of
six-trit blocks. Concatenating the six coordinates of each block yields
a ternary tape
\(
 (a_n)_{n\geq 1}\in\{0,1,2\}^{\mathbb N}
\), on which the evaluation acts
\[
 \operatorname{ev}_3((a_n))
 :=\sum_{n\geq1}\frac{a_n}{3^n}\in[0,1].
\]
The two expansions of a ternary endpoint represent the same value and
remain distinct before evaluation is applied. This does not constitute
an ambiguity of the genealogical object: it is precisely a fiber of the reader.

We define
\[
 X_{\mathrm{raw}}:=\mathbb Z\times\mathbb Z_3^6,
 \qquad
 P(m,x):=m+\operatorname{ev}_3(\operatorname{flat}(\Psi x)).
\]
The integer coordinate provides the countable atlas of the line, and the
3-adic coordinate preserves the complete history of the fractional part.

\begin{theorem}[Exact reconstruction of the Archimedean quotient]
\label{p12-83-c18-s2:hmtch:thm-cociente-real}
The map \(P:X_{\mathrm{raw}}\to\mathbb R\) is sobreyectiva. If
\[
 x\sim_P y
 \quad\Longleftrightarrow\quad
 P(x)=P(y),
\]
then
\[
 \overline P:X_{\mathrm{raw}}/\!\sim_P\;\longrightarrow\mathbb R,
 \qquad [x]\longmapsto P(x),
\]
is a bijection. Consequently,
\[
 \left|X_{\mathrm{raw}}/\!\sim_P\right|
 =|\mathbb R|
 =2^{\aleph_0}.
\]
\end{theorem}

\begin{proof}
Every ternary sequence is grouped uniquely into consecutive blocks of
six digits. The conjugation \(\Psi\) assigns a unique initial
3-adic memory to that tape. Given \(r\in\mathbb R\), take
\(m=\lfloor r\rfloor\) and a ternary expansion of
\(r-m\in[0,1)\); the preceding reconstruction produces
\(x\in\mathbb Z_3^6\) with \(P(m,x)=r\). Thus, \(P\) is surjective.

The map \(\overline P\) is well defined by the definition of
\(\sim_P\), is surjective because \(P\) is, and is injective because two
classes with the same image belong to the same fiber. Finally,
\[
 |\mathbb Z_3^6|
 =\left|(\mathbb F_3^6)^{\mathbb N}\right|
 =729^{\aleph_0}
 =2^{\aleph_0},
\]
and multiplying by the countable coordinate \(\mathbb Z\) does not alter that
cardinal. The calculation agrees with that of the quotient, already identified with
\(\mathbb R\).
\end{proof}

\begin{remark}[Survivor Version]
\label{p12-83-c18-s2:hmtch:rem-version-superviviente}
The preceding theorem is unconditional for the raw completion. For a
surviving subspace \(X_{\mathrm{surv}}\), the same result is
obtained when the covers at every level are coherent, the diameters
contract uniformly, and every visible cell has a prolongation. Under
these hypotheses, the infinite-branch lemma produces a section over every
point, and evaluation induces, compact set by compact set, a homeomorphism of the
quotient with the published interval. Exhaustion by
\([-m,m]\) then reconstructs the entire real line.
\end{remark}

\section{Archimedean mean and double projection of the same state}

Let \(u=(m,\mathbf u)\in
\mathscr O_{\mathbb R,\leftrightarrow}^{\rm HMT}\). A prefix
\(\mathbf u_n\) determines in the fractional chart an oriented
semi-open cylinder
\[
 I_n(\mathbf u_n)=[a_n,b_n),
\]
with the complementary convention at the endpoints of double expansion. To
apply compactness, its closure
\(K_n=[a_n,b_n]\) is used. The compatibilities give
\[
 K_{n+1}\subseteq K_n,
 \qquad
 \delta_n:=b_n-a_n\longrightarrow0.
\]
The global Archimedean mean of the cylinder is
\[
 \mu_n(u)=m+\frac{a_n+b_n}{2}.
\]
If \(q\geq n\), then
\[
 |\mu_q(u)-\mu_n(u)|\leq\delta_n,
\]
so that
\begin{equation}
 \mathfrak R(u):=\lim_{n\to\infty}\mu_n(u)
 \label{p12-83-c18-s2:hmtch:eq-media-arquimediana}
\end{equation}
exists. What tends to zero is the diameter of the cylinder; the real value
is the resultant of the nesting, not that which vanishes.

\begin{theorem}[Two-way quotient and real line]
\label{p12-83-c18-s2:hmtch:thm-cociente-two-way-real}
Let
\[
 u\sim_{\mathfrak R}v
 \quad\Longleftrightarrow\quad
 \mathfrak R(u)=\mathfrak R(v).
\]
Under the coherent covering of the HMT charts, the evaluation induces a homeomorphism.
\begin{equation}
 \overline{\mathfrak R}:
 \mathscr O_{\mathbb R,\leftrightarrow}^{\rm HMT}/\!
 \sim_{\mathfrak R}
 \xrightarrow{\ \sim\ }\mathbb R.
 \label{p12-83-c18-s2:hmtch:eq-homeomorfismo-two-way-real}
\end{equation}
In particular, this is a continuous HMT morphism: not the result of choosing
an isolated history from each fiber.
\end{theorem}

\begin{proof}
The cylinder bound proves the continuity of \(\mathfrak R\). The covering produces a history over every real and hence surjectivity. Upon passage to the quotient by the fibers, the induced map is bijective. On every compact chart, a continuous bijection onto a Hausdorff interval is a homeomorphism. The global topology of the quotient is the final topology of the gluing of these charts; the evaluation fiber identifies endpoint \(m+1\) of one chart with the translated endpoint \(0\) of the following chart. The resulting cover is locally finite, so the chart homeomorphisms glue into a global homeomorphism onto \(\mathbb R\).
\end{proof}
  \part{Mathematical characters of coinductive extension}

\chapter{Seeds, cubic completion and synchronization of the fundamental characters}
\section{The two mechanical projections and exact conservation of the boundary unit}

The statement that the unit is preserved during extension already admits an arithmetic formulation independent of a physical analogy. Let

\[
 \delta:=\log_{10}\!\left(\frac{10}{9}\right),
 \qquad 0<\delta<1,
\]

be the discrepancy between the triadic capacity \(3^6=729\) and the decimal completion \(10^3=1000\). The irrationality of \(\delta\) has already been proved in the construction through prime valuations; its transcendence has been proved in this notebook through Gelfond--Schneider. Consider the two one-sided codings of the same irrational line:

\[
 z_n^-:=\lfloor(n+1)\delta\rfloor-\lfloor n\delta\rfloor,
 \qquad
 z_n^+:=\lceil(n+1)\delta\rceil-\lceil n\delta\rceil,
 \qquad n\geq0.
\]

Both take values in \(\{0,1\}\). They are neither two densities nor two dynamics: they are the two closure conventions of the same irrational cut. Since \(n\delta\notin\mathbb Z\) for every \(n\geq1\),

\[
 z_0^-=0,\qquad z_0^+=1,\qquad z_n^-=z_n^+\quad(n\geq1).
\]

By telescoping, their counts in the prefix of length \(N>0\) are

\[
 Z_N^-:=\sum_{n=0}^{N-1}z_n^-=\lfloor N\delta\rfloor,
 \qquad
 Z_N^+:=\sum_{n=0}^{N-1}z_n^+=\lceil N\delta\rceil.
\]

We therefore obtain the boundary theorem

\[
 \boxed{Z_N^+-Z_N^-=1\qquad(N>0).}
\]

The unit is not created at the end of the block: it was localized in the choice of which half-boundary belongs to the cylinder. If the two polarizations are defined by

\[
 r_N^-:=N\delta-\lfloor N\delta\rfloor,
 \qquad
 r_N^+:=\lceil N\delta\rceil-N\delta,
\]

then

\[
 0<r_N^\pm<1,
 \qquad
 \boxed{r_N^-+r_N^+=1.}
\]

This is the first exact form of the \textbf{conservative evolution of the unit}: the two holographic projections distribute one boundary unit and their sum recovers it at every finite depth. Energy conservation is not yet being asserted. Thermodynamic promotion also requires the physical realization of the operator to be reversible. In the full TPK state, bijective updating already implies \(H(X\mid U(X))=0\); this logical statement is distinct both from the zero topological entropy of the Sturmian word and from a concrete thermal implementation.

The law is also bidirectional. The lower reading fixes the outgoing half-boundary and the upper reading fixes the incoming one. Reversing orientation interchanges the two conventions and transports the unit, but alters neither the slope nor the long-range order. This is the minimum content that any \emph{two-way} reader must preserve: inversion changes the location of the boundary unit rather than eliminating it.

\textbf{Refutation criterion.} An \(N>0\) for which \(Z_N^+-Z_N^-\neq1\), or wording that turns this arithmetic identity into an energy equality without a transducer.

\section{Concatenation cocycle, carry memory and exact restriction of the TPK}

For an initial phase \(k\in\mathbb Z\), the number of vacancies in a block of length \(N\) is

\[
 V_N(k):=\lfloor(k+N)\delta\rfloor-\lfloor k\delta\rfloor.
\]

The word is balanced, so that

\[
 V_N(k)\in
 \left\{\lfloor N\delta\rfloor,\lceil N\delta\rceil\right\}
\]

for every phase. More important for the enriched dynamics is the composition identity

\[
 \boxed{V_{M+N}(k)=V_M(k)+V_N(k+M).}
\]

The second summand cannot be evaluated afresh at the origin: it receives the phase transported by the first block. Forgetting that phase loses exactly the carry

\[
 \kappa_\delta(M,N)
 :=\lfloor(M+N)\delta\rfloor
   -\lfloor M\delta\rfloor-\lfloor N\delta\rfloor
 \in\{0,1\}.
\]

Thus,

\[
 \lfloor(M+N)\delta\rfloor
 =\lfloor M\delta\rfloor+\lfloor N\delta\rfloor
  +\kappa_\delta(M,N).
\]

Associativity of addition forces the cocycle equation

\[
 \kappa_\delta(L,M)+\kappa_\delta(L+M,N)
 =\kappa_\delta(M,N)+\kappa_\delta(L,M+N).
\]

Thus the carry is not a numerical correction added after completing the calculation: it is the datum making block concatenation associative when their history is preserved. This is precisely the structure concealed by projection onto mere cardinalities.

The restriction of TPK dynamics to the vacancy subsystem can be written without replacing the full fibre. On

\[
 \mathcal Y_{108}:=C_{108}\times\mathbb T\times\mathbb Z
\]

we define

\[
 \mathcal F(r,\theta,m)
 =\bigl(r+1,\{\theta+\delta\},
         m+\lfloor\theta+\delta\rfloor\bigr),
 \qquad 0\leq\theta<1.
\]

The first component is the finite clock; the second, the irrational phase; the third, the integrated memory. The map is bijective, with inverse

\[
 \mathcal F^{-1}(r,\theta,m)
 =\bigl(r-1,\{\theta-\delta\},
 m-\lfloor\{\theta-\delta\}+\delta\rfloor\bigr).
\]

In the governing operator language, this formula is a \textbf{factor} of \(U_t=\operatorname{Upd}_t\circ\operatorname{Tra}_t\circ
\operatorname{Sel}_t\): the selector distinguishes event and non-event; transport advances the nonadic and irrational phases; updating incorporates the pulse into the memory record. It does not replace the TPK fibre, which also preserves sheet, orientation, residue, quotient, signature, cylinder, boundary, route and survival.

After \(108\) iterations,

\[
 \mathcal F^{108}(r,\theta,m)
 =\bigl(r,\{\theta+108\delta\},m+V_{108}(\theta)\bigr),
\]

with \(V_{108}(\theta)\in\{4,5\}\). The first clock returns; neither the irrational phase nor the memory returns. This equality types the phrase “the phase returns, not the state” exactly within the aperiodic subsystem.

\textbf{Refutation criterion.} Failure of the cocycle equation, non-invertibility of \(\mathcal F\), or omission by the factor projection of a field subsequently used to distinguish states.

\section{The renormalized spectrum of cardinalities and its two exact self-references}

Define the balanced length reader

\[
 \mathcal W_\delta(N)
 :=\{\lfloor N\delta\rfloor,\lceil N\delta\rceil\}.
\]

This operator assigns to each integer scale the only two counts permitted by the irrational cut. It consults neither \(\pi\), \(\varphi\), \(e\), \(\alpha\), a mass nor an experimental constant. Applying it to cardinalities already produced by APP--TRIT--TPK yields the following exact equalities:

\[
 \begin{aligned}
 \mathcal W_\delta(54)&=\{2,3\},&
 \mathcal W_\delta(81)&=\{3,4\},\\
 \mathcal W_\delta(90)&=\{4,5\},&
 \mathcal W_\delta(120)&=\{5,6\},\\
 \mathcal W_\delta(137)&=\{6,7\},&
 \mathcal W_\delta(202)&=\{9,10\},\\
 \mathcal W_\delta(210)&=\{9,10\},&
 \mathcal W_\delta(243)&=\{11,12\},\\
 \mathcal W_\delta(271)&=\{12,13\},&
 \mathcal W_\delta(324)&=\{14,15\},\\
 \mathcal W_\delta(369)&=\{16,17\},&
 \mathcal W_\delta(468)&=\{21,22\},\\
 \mathcal W_\delta(729)&=\{33,34\},&
 \mathcal W_\delta(1000)&=\{45,46\}.
 \end{aligned}
\]

Two rows are self-referential in a precise sense. The integer \(137\), obtained previously as the trace of the synchronization matrix, reproduces the secondary returns when read by \(\mathcal W_\delta\):

\[
 \boxed{\mathcal W_\delta(137)=\{6,7\}.}
\]

The proved count \(468\) similarly reproduces the primary gaps:

\[
 \boxed{\mathcal W_\delta(468)=\{21,22\}.}
\]

This does not round \(\alpha^{-1}\) to \(137\). The trace \(137\) was constructed beforehand from integer synchronization, and the vacancy reader subsequently recovers \(6/7\). Nor is \(21/22\) introduced to select \(468\): the count belongs to nonadic extension and the same reader returns it as a window spectrum. These two identities constitute a renormalization closure within the genealogy, not a retrospective definition.

There is a third particularly clear structural identity. Since

\[
 210=90+120,
\]

the ordered signature of the interval is

\[
 \bigl(V_{90}(0),V_{120}(90)\bigr)=(4,5),
\]

and therefore

\[
 Z^-_{210}=4+5=9,
 \qquad
 Z^+_{210}=1+4+5=10.
\]

The channels \(90/120\), generated in the TPK before their parangular evaluators, thus compress the transition \(9\to10\) through the same boundary unit. The equality does not say that \(90\) and \(120\) are “the numbers” nine and ten; it says that the capacity reader \(729/1000\), applied to their causal concatenation, produces exactly those two coordinates.

\textbf{Refutation criterion.} A single discrepancy in any of the exact floors or ceilings, or the use of an externally recognized value to choose the lengths.

\section{Cubic completion jointly contains scales 34 and 11}

The completion

\[
 10^3=(9+1)^3=729+243+27+1
\]

now has an exact image under the vacancy cocycle. Preserving the order of the four summands, the lower mechanical word produces the signature

\[
 \boxed{(33,11,1,0),}
\]

because

\[
 \begin{aligned}
 V_{729}(0)&=33,\\
 V_{243}(729)&=11,\\
 V_{27}(972)&=1,\\
 V_1(999)&=0.
 \end{aligned}
\]

The sum is \(45\), whereas the upper reading places the boundary unit in the first block and gives

\[
 \boxed{(34,11,1,0),\qquad 34+11+1+0=46.}
\]

This yields a refined form of the law \(45/46\): it not only counts events in a window of one thousand positions, but also preserves the provenance of each event within the completion \(729+243+27+1\). The order of the summands is part of the data. A permutation may redistribute the carry between blocks even though the total sum remains fixed; the enriched signature therefore contains more information than the cardinality \(45\) or \(46\).

The scale \(729\) in turn decomposes into three visible blocks of length \(243\):

\[
 \bigl(V_{243}(0),V_{243}(243),V_{243}(486)\bigr)
 =(11,11,11).
\]

Consequently,

\[
 Z^-_{729}=11+11+11=33,
 \qquad
 \boxed{Z^+_{729}=1+11+11+11=34.}
\]

This result places \(11\) and \(34\) in a single aperiodic hierarchy before any dimensional interpretation. Moreover,

\[
 \begin{aligned}
 324&=243+81,& (V_{243}(0),V_{81}(243))&=(11,3),\\
 369&=243+126,&(V_{243}(0),V_{126}(243))&=(11,5),
 \end{aligned}
\]

from which

\[
 Z^-_{324}=14,\quad Z^+_{324}=15,
 \qquad
 Z^-_{369}=16,\quad Z^+_{369}=17.
\]

Two connections with previously proved objects are exact at the cardinal level:

\[
 \boxed{Z^+_{324}=15
 =\det\begin{pmatrix}21&22\\6&7\end{pmatrix},}
\]

and

\[
 \boxed{Z^-_{369}=16=|\det H_4|=|\operatorname{coker}H_4|.}
\]

Combined with the active determinantal theorem,

\[
 \operatorname{ord}_{10}
 \bigl(\varphi\alpha^{16}\bigr)=-34,
\]

they produce the scalar factorization

\[
 \boxed{
 -\operatorname{ord}_{10}
 \bigl(\varphi\alpha^{Z^-_{369}}\bigr)
 =Z^+_{729}.}
\]

The preliminary formulation of this result described it as a square reduced to cardinalities. Sections 95--96 complete its typing: the relevant domain is not an arbitrary enumeration of sixteen vacancies, but the ordered signature of cubic completion and its depth reader. The action branch factors exactly as

\[
 \boxed{
 \mathcal V^-_{369}
 \xrightarrow{\;Z^-\;}16
 =|\operatorname{coker}(H_4)|
 \xrightarrow{\;\mathcal N_\alpha\;}
 \alpha^{16}
 \xrightarrow{\;\varphi\cdot\;}
 \Sigma_H
 \xrightarrow{\;-\operatorname{ord}_{10}\;}34
 =Z^+_{729}.}
\]

Here \(\mathcal N_\alpha(n)=\alpha^{n}\) is the norm reader of the determinantal exponent. The Smith normal form of \(H_4\) produces the cokernel of order \(16\); the section \(\alpha\) transports it to the norm \(\alpha^{16}\); the autoscale \(\varphi\) forms \(\Sigma_H\); and decimal order determines the depth \(34\). The equality \(Z^-_{369}=|\operatorname{coker}(H_4)|\) is therefore not used as an elementwise bijection: it acts as an equality of invariants in the composition proving the action depth. This is precisely the type of statement required by the theorem.

Likewise,

\[
 Z^-_{243}=11
\]

is the second coordinate of the same cubic reader. Its typed factorization is

\[
 \boxed{
 \mathcal V^-_{243}
 \xrightarrow{\;Z^-\;}11
 =\operatorname{pr}_{243}\!\left(
   \mathcal D_{\rm cub}(\sigma^-,\sigma^+)\right)
 \xrightarrow{\;\operatorname{Pub}_G\;}
 \mathcal D_{10}(G)=11.}
\]

Assembly with the action branch gives, without consulting target values,

\[
 \boxed{
 \bigl(\mathcal D_{10}(\hbar),\mathcal D_{10}(G)\bigr)=(34,11),
 \qquad
 \mathcal D_{10}(G\otimes\hbar)=34+11=45,
 \qquad 45+1=46.}
\]

The identity \(3^6\cdot243=3^{11}\) preserves its own domain as a dimension reader; the completion coordinate \(11\) and gravitational depth are connected through \(\mathcal D_{\rm cub}\) and \(\operatorname{Pub}_G\), not through tacit identification of their elements. The product \(G_a\hbar_a\), proved in Section 95, closes the dynamic square: two-way return reciprocally inverts both branches and preserves their product.

\textbf{Refutation criterion.} Failure of the ordered signatures to telescope, a determinant or cokernel differing from the declared values, failure of \(\mathcal D_{\rm cub}\) to produce \((34,11,45,46)\), or a dimensional realization depending on \(\hbar\), \(G\), CODATA or the very decimal position it must generate.

\section{The electronic splitting of the ambient space and the internal emergence of 9, 10 and 24}

At the first joint reopening, the electronic channel contains \(527\) extensions of an ambient space of \(729\). Its deficit is

\[
 202=729-527=6+28\cdot7.
\]

The ordered decomposition of the ambient space

\[
 729=202+527
\]

is exactly compatible with the aperiodic reader:

\[
 \bigl(V_{202}(0),V_{527}(202)\bigr)=(9,24).
\]

Therefore,

\[
 \boxed{Z^-_{729}=9+24=33,\qquad
        Z^+_{729}=10+24=34.}
\]

The same deficit that singles out the electronic window reproduces, under \(\mathcal W_\delta\), the nonadic--decimal boundary:

\[
 \boxed{\mathcal W_\delta(202)=\{9,10\}.}
\]

The boundary unit acts here on the deficit sector; the surviving sector preserves twenty-four events in the canonical phase. This decomposition refines the preceding equality \(34=1+3\cdot11\): the same coordinate is now read as \(34=10+24\). Both factorizations coexist because they record different cuts of the same cylinder, not because one replaces the other.

The integer \(24\) coincides with the rank of the Leech lattice, but cardinal equality does not identify the \(527\) electronic extensions with lattice vectors. Its exact content is a rank connection in the second projection. Exceptional intertwining is constructed, with its own incidence and quadratic form, through the Paley--Witt--Golay--Leech chain; no action of the Monster group is assigned to digits or vacancies.

\textbf{Refutation criterion.} An electronic count other than \(527\), a complementary block not containing twenty-four events, or absence of any transport preserving the required exceptional incidence.

\section{The two finite clocks and aperiodic helical suspension}

The clock \(108\) divides into two half-cycles of \(54\) steps. The lower word has signature

\[
 \bigl(V_{54}(0),V_{54}(54)\bigr)=(2,2),
\]

so that

\[
 Z^-_{108}=2+2=4,
 \qquad Z^+_{108}=1+2+2=5.
\]

The return \(54+54\) closes the finite phase and preserves a choice of half-boundary. Triple repetition of the clock does not, however, repeat the same signature:

\[
 \bigl(V_{108}(0),V_{108}(108),V_{108}(216)\bigr)=(4,5,5),
\]

and therefore

\[
 Z^-_{324}=4+5+5=14,
 \qquad Z^+_{324}=1+4+5+5=15.
\]

This directly manifests aperiodic order: three blocks of equal length have balanced but nonidentical counts. The sequence preserves long-range order and discrepancy less than one without acquiring a translational period.

The cardinality \(324\) admits two internal readings already present in the construction: \(324=81\cdot4\) for oriented routes and \(324=108\cdot3\) for a tritic lift of the clock. The second factorization defines the minimal cyclic model

\[
 0\longrightarrow C_3\longrightarrow C_{324}
 \longrightarrow C_{108}\longrightarrow0.
\]

The extension is non-split: \(C_{324}\) contains an element of order \(324\), whereas \(C_3\times C_{108}\) has exponent \(108\). In the lifted clock, a translation of \(108\) steps returns to the same coordinate of \(C_{108}\) but advances to the next sheet of the kernel \(C_3\); only after \(324\) steps does the lifted coordinate return. This is the smallest algebraic model of a three-sheeted helix over the clock \(108\).

Decorating this clock with the irrational rotation of slope \(\delta\) yields a Sturmian helical suspension: the finite base may return, the lifted sheet may return with a longer period, and the irrational phase never returns exactly. The accumulated memory distinguishes these three facts. The ascending helix uses \(\mathcal F\); the descending helix uses \(\mathcal F^{-1}\). They are conjugate readings of the same reversible dynamics, not two independent series. The mirror \(\pi/e\) remains another involution: it acts on the output windows and fixes \(\varphi\); it does not by itself reverse temporal orientation.

\textbf{Refutation criterion.} A group section of \(C_{324}\to C_{108}\), a repetition \((4,4,4)\) or a conjugation that erases memory.

\section{The first supercell that closes clock 108}

The continued fraction of \(\delta\) begins

\[
 [0;21,1,5,1,6,2,3,1,1,8,1,2,\ldots].
\]

The first convergent whose denominator is divisible by \(108\) is

\[
 \boxed{\frac{7813}{170748}.}
\]

No earlier convergent denominator has this property. Moreover, the rational certificate proves that

\[
 \left\lfloor n\delta\right\rfloor
 =\left\lfloor\frac{7813n}{170748}\right\rfloor
 \qquad(0\leq n<170748),
\]

whereas at the final endpoint

\[
 \lfloor170748\delta\rfloor=7812,
 \qquad
 \frac{7813\cdot170748}{170748}=7813.
\]

The aperiodic word and its periodic Christoffel approximant coincide on every proper prefix of the supercell and differ only at closure. The periodic approximant does not correct a cloud of errors: it transports the boundary unit exactly from the lower projection to the upper one.

The denominator has the factorization

\[
 \boxed{
 170748=108\cdot1581
       =324\cdot527
       =2\cdot54\cdot3\cdot527.}
\]

Here \(527\) is exactly the electronic-channel count at the first reopening and \(324\) the cardinality of oriented routes. Since \(\gcd(324,527)=1\), there is the canonical clock decomposition

\[
 C_{170748}\cong C_{324}\times C_{527}.
\]

The Chinese remainder theorem justifies this group decomposition. The second factor has the cardinality of the electronic fibre and the first that of the oriented atlas; this coincidence of cardinalities neither identifies the objects nor automatically transfers their actions.

The pair of clocks provides a second unexpected datum. In the lower branch, after \(170748\) blocks,

\[
 7812\equiv36\pmod{108},
 \qquad
 170748-7812=162936\equiv72\pmod{108}.
\]

Therefore, although the block clock closes,

\[
 \boxed{(r_{\rm vac},r_{\rm cap})=(36,72),\qquad36+72=108.}
\]

The upper branch transfers the boundary unit:

\[
 (r_{\rm vac}^+,r_{\rm cap}^+)=(37,71),
 \qquad37+71=108.
\]

The appearance of \(36\) in the lower branch coincides with the index of stage \(R_{36}\) of coinductive extension. This congruence does not by itself identify the vacancy residue with the coordinate of \(R_{36}\); it proves that the supercell closes the block clock, preserves the distribution of the unit and leaves nontrivial residual memory in the capacity clock.

The numerator has another internal factorization:

\[
 7813=13\cdot601,
 \qquad
 7812=36\cdot217,
 \qquad217=2\cdot108+1.
\]

This yields the unimodular frame

\[
 \boxed{
 M_\star=
 \begin{pmatrix}
 601&217\\
 36&13
 \end{pmatrix}\in\operatorname{SL}_2(\mathbb Z),
 \qquad
 601\cdot13-217\cdot36=1.}
\]

The four integers are localized in the architecture: \(601\) is the twelfth coordinate of the dodecaphasic seal \(K\); \(13\) is the number of atlas multisections; \(36\) is stage \(R_{36}\); and \(216\) is the extended kinematic return. The equality is exact and uses no physical value. It defines a unimodular frame of their cardinalities; it does not on that basis alone postulate a categorical equivalence between the four objects.

\textbf{Refutation criterion.} An earlier convergent with denominator a multiple of \(108\), a discrepancy in a proper prefix, a false factorization or a transport that does not conjugate the actions.

\section{The second derivation of \(\alpha\) as renormalization of a single word}

The polynomial \(P_9\) of the second route to \(\alpha\) contains the pairs \(21/22\), \(6/7\) and \(45/46\). Their coefficients had previously been reconstructed through the synchronization matrix, Beatty returns and the boundary \(7{:}9\). The new reader proves that all three pairs are also images of \textbf{the same} mechanical word at three internal scales:

\[
 \boxed{
 \mathcal W_\delta(468)=\{21,22\},\qquad
 \mathcal W_\delta(137)=\{6,7\},\qquad
 \mathcal W_\delta(1000)=\{45,46\}.}
\]

The first length is the decimal-emission count, the second is the integer synchronization trace and the third is cubic completion. The genealogy of the second route can therefore be written without introducing its coefficients as an independent list:

\[
 \begin{aligned}
 &\mathrm{APP}\longrightarrow\mathrm{TRIT}\longrightarrow\mathrm{TPK}
 \longrightarrow X^{\rm enr},\\
 &729\longrightarrow1000\longrightarrow\delta
 \longrightarrow\mathcal W_\delta,\\
 &(468,137,1000)
 \longmapsto
 \bigl(\{21,22\},\{6,7\},\{45,46\}\bigr),\\
 &\text{synchronization and boundary signature }7{:}9
 \longrightarrow P_9
 \longrightarrow\alpha.
 \end{aligned}
\]

This chain does not replace the first dodecaphasic route. Both start from the coinductive state that also determines \(\pi\), \(\varphi\) and \(e\); they subsequently preserve their own readers. The root of \(P_9\) is calculated at the end, once the signature has been constructed. The value of \(\alpha\) never selects \(468\), \(137\), \(1000\) or their windows.

The numerator \(7813=13\cdot601\) also locates a possible closure between the aperiodic supercell and the dodecaphasic seal: \(601=K_{12}\) belongs to the first route, whereas \(13\) belongs to the subsequent atlas. The unimodular identity of the preceding subsection allows the question to be formulated correctly: one must determine whether the memory module on which \(D_3K\), \(R_{36}\), the return \(216\) and the thirteen multisections act admits the change of basis \(M_\star\). Coincidence of factors must not be turned into proof, but neither should it be discarded as a mere curiosity: it is localized in the first supercell selected by the clock \(108\), not in an integer sought afterwards.

\textbf{Refutation criterion.} Any pair not proceeding from \(\mathcal W_\delta\), involvement of \(P_9\) in selecting its own inputs, or failure of the module proposed for \(M_\star\) to preserve the memory record.

\section{Seed deepening and rigorous separation of physical scales}

The space of an oriented seed has cardinality

\[
 |\mathcal S|=9^2\cdot4=324,
\]

and the pair sweep has

\[
 |\mathcal S\times\mathcal S|=324^2=104976.
\]

Applying the aperiodic reader to the total cardinality gives

\[
 \mathcal W_\delta(104976)=\{4803,4804\}.
\]

This equality is a count of the mechanical line, not yet a seed count. To turn it into a property of the APP sweep, one must fix a canonical ordering

\[
 \operatorname{ord}_{\mathcal S^2}:
 \mathcal S\times\mathcal S\overset\sim\longrightarrow
 \{0,\ldots,104975\}
\]

compatible with the two sheets, the dihedral action, the TRIT selector and TPK updating. Only then can one ask which seeds occupy the vacancies and how the count changes under inward refinement or outward expansion. Cardinality alone must not select this ordering.

The new analysis nevertheless already determines where to seek the scales \(10^{-11}\) and \(10^{-34}\). They are not obtained by arbitrarily raising a seed to a decimal power. The crystal first produces the integers

\[
 Z^-_{243}=11,
 \qquad
 Z^+_{729}=34,
\]

within the same completion signature. Two distinct transducers can subsequently realize those orders on dimensional lines:

\[
 \begin{aligned}
 \mathscr T_G&:\bigl(X^{\rm enr},\mathcal V^-_{243}\bigr)
 \longrightarrow 10^{-11}\mathcal L_G,\\
 \mathscr T_\hbar&:\bigl(X^{\rm enr},\mathcal V^+_{729},G_H,\alpha,\varphi\bigr)
 \longrightarrow10^{-34}\mathcal L_S.
 \end{aligned}
\]

The second transducer already has a proved route through \(|G_H|=16\) and \(\varphi\alpha^{16}\). The first retains a decimal chart and selection criteria in the authoritative source and must be reconciled with the vacancy reader without using CODATA. The phrase “fractalize inwards down to \(10^{-11}\) and \(10^{-34}\)” thus acquires a precise mathematical meaning: refine the same state to two window cardinalities, and only then apply the corresponding dimensional transduction. Powers of ten belong to the chart of the dimensional line, not to the bare seed.

The unit remains conserved during this step if the transducer simultaneously preserves both projections and the cocycle. The requirement is a naturality square

\[
 \mathscr T\circ\operatorname{Upd}
 =\operatorname{Upd}_{\mathcal L}\circ\mathscr T
\]

in which the left-hand side first updates the discrete memory and the right-hand side first transports the dimensional section. This equation, rather than an isolated decimal equality, is the requisite proof that a physical constant emerges “at the appropriate scale”.

The discrete structure of the continuum must receive all these fibres together. Seed refinement, extension solvability, measure construction, coherence and the determinantal reader are not five separable objects: they are five inseparable aspects of the same passage to the limit. The aperiodic crystal supplies the balanced calendar; APP preserves sum, product, residue, quotient and carry; TRIT types orientation and the three curvatures; TPK preserves history throughout selection, transport and updating. Only their assembly produces the real line as an output. In this genealogy, the line \(\mathbb R\) is not the support presupposed to construct the protocol: it is one of its terminal realizations.

\textbf{Refutation criterion.} An ordering not natural with respect to APP, TRIT and TPK; a decimal power chosen from the physical value; or a transducer not preserving the boundary unit and memory record.

\section{The cellular corona of the continuum contains clock 108 and the TRIT signature}

Review of the cellular reference construction supplies a component linking completion, the clock and double projection more rigidly than a numerical coincidence. For the cubical pair

\[
 Q_{9,3}\subset Q_{10,3},
\]

the relative complex has the cardinality vector

\[
 \mathbf c=(c_0,c_1,c_2,c_3)=(271,756,702,217).
\]

Its Euler characteristic vanishes:

\[
 \chi(\mathbf c)=271-756+702-217=0.
\]

This vanishing, however, is not an opaque cancellation. Each coordinate admits a decomposition into cardinalities already belonging to the architecture:

\[
 \begin{aligned}
 271&=243+27+1,\\
 756&=729+27,\\
 702&=729-27,\\
 217&=216+1=2\cdot108+1.
 \end{aligned}
\]

Pairing the outer and inner degrees yields two equal defects:

\[
 \boxed{c_0-c_3=271-217=54,}
 \qquad
 \boxed{c_1-c_2=756-702=54.}
\]

Consequently,

\[
 \boxed{(c_0-c_3)+(c_1-c_2)=54+54=108.}
\]

The clock \(108\) has not been imposed on the corona from an external calendar: its double half-cycle is contained in the difference between the boundary degrees and in the difference between the inner degrees. The vanishing characteristic is rewritten as

\[
 \chi(\mathbf c)=(c_0-c_3)-(c_1-c_2)=54-54=0.
\]

This identity distinguishes two senses. The first \(54\) measures separation between the extremities of the complex; the second measures internal separation. Their equality makes closure compatible: what exits through one pairing returns through the other without erasing cellular degree. The sum \(54+54\) measures the length of the bidirectional path; the difference \(54-54\) certifies its Euler closure. These are two distinct operations on the same pair and must not be flattened into the single number \(108\).

The aperiodic reader is now applied to each degree

\[
 \mathcal W_\delta(N)=
 \{\lfloor N\delta\rfloor,\lceil N\delta\rceil\}.
\]

The two vacancy projections of the corona are

\[
 \mathbf v^-=(12,34,32,9),
 \qquad
 \mathbf v^+=(13,35,33,10).
\]

Their capacity complements are

\[
 \mathbf b^+=\mathbf c-\mathbf v^-
 =(259,722,670,208),
\]

\[
 \mathbf b^-=\mathbf c-\mathbf v^+
 =(258,721,669,207).
\]

In both one-sided closures, component by component,

\[
 \mathbf c=\mathbf v^-+\mathbf b^+
           =\mathbf v^++\mathbf b^-.
\]

The decisive information appears upon taking the alternating characteristic:

\[
 \boxed{\chi(\mathbf v^-)=\chi(\mathbf v^+)=+1,}
\]

\[
 \boxed{\chi(\mathbf b^-)=\chi(\mathbf b^+)=-1.}
\]

The corona of zero total charge therefore polarizes exactly as

\[
 \boxed{0=(+1)+(-1).}
\]

Independence from the half-boundary convention is structural:

\[
 \mathbf v^+-\mathbf v^-=(1,1,1,1),
 \qquad
 \chi(1,1,1,1)=0.
\]

Moving the boundary unit between the two closures does not alter the Euler charge of the branch. This yields an internal realization of the TRIT alphabet:

\[
 \boxed{
 \text{vacancy projection}\mapsto+1,
 \quad
 \text{complete crown}\mapsto0,
 \quad
 \text{capacity projection}\mapsto-1.}
\]

This does not replace the general definition of the TRIT or by itself identify the three charges with physical fields. It proves something more elementary and prior: double projection of a single cellular object generates the three oriented classes (+1,0,-1), without adding symbols. The algebras

\[
 \mathbb A_\tau=\mathbb R[J_\tau]/(J_\tau^2+\tau I)
\]

can subsequently act on these classes. Elliptic, parabolic or hyperbolic type is no longer selected from a presupposed continuous geometry: the signature indexing it has emerged from the discrete polarization of the complex.

\textbf{Refutation criterion.} A change in any of the four cardinalities, a defect no longer equal to \(54\), projections not having charges (+1/-1), or use of classical geometry to select the sign retrospectively.

\section{Polarized frame, boundary root and conservative evolution of the unit}

The preceding result is a case of an exact law valid for every length. Let

\[
 a_N=\lfloor N\delta\rfloor,
 \qquad N>0.
\]

The two closure conventions determine the vectors

\[
 u_N^-=(a_N,N-a_N),
 \qquad
 u_N^+=(a_N+1,N-a_N-1).
\]

Each row preserves the block exactly:

\[
 (1,1)u_N^-=(1,1)u_N^+=N.
\]

The difference is independent of (N):

\[
 \boxed{u_N^+-u_N^-=(1,-1).}
\]

The unit neither appears nor disappears; it is transported from the capacity coordinate to the vacancy coordinate. The vector ((1,-1)) generates the kernel of the augmentation

\[
 \mathbb Z^2\xrightarrow{\ (x,y)\mapsto x+y\ }\mathbb Z,
\]

that is, the root \((A_1)\). Double projection therefore has a canonical boundary direction. No physical metric need be chosen to define it.

The frame of the two polarizations has determinant

\[
 \boxed{
 \det\begin{pmatrix}
 a_N&N-a_N\\
 a_N+1&N-a_N-1
 \end{pmatrix}=-N.}
\]

The proof is immediate:

\[
 a_N(N-a_N-1)-(N-a_N)(a_N+1)=-N.
\]

Thus, transferring a single boundary unit sweeps out an entire parallelogram of oriented area (N). After normalization by (N), the two rows lie in the simplex (x+y=1), are separated by ((1,-1)/N) and converge to

\[
 (\delta,1-\delta)=(\delta,\rho).
\]

This is a precise form of the conservative evolution of the dimensional unit: at each scale there is one normalized unit; refinement does not multiply it, but reduces the size of its boundary quantum while preserving the integral direction that allows it to be recovered.

The complete trajectory is written in the lattice as

\[
 \gamma(N)=
 \bigl(\lfloor N\delta\rfloor,
       N-\lfloor N\delta\rfloor\bigr).
\]

Each step increments exactly one coordinate:

\[
 \gamma(N+1)-\gamma(N)\in\{(1,0),(0,1)\}.
\]

Moreover,

\[
 \left\|\gamma(N)-N(\delta,\rho)\right\|_\infty<1.
\]

In its lattice section, the aperiodic helix is a digital straight line with uniformly bounded discrepancy. It has long-range order because it never departs by more than one cell from the limiting direction, but has no translational period because \(\delta\) is irrational. The Sturmian word is not an embellishment of this trajectory: it is the sequence of decisions specifying which of the two coordinates receives the next unit.

The same augmentation structure reappears in the three channels. The reopening vector

\[
 v_\Omega=(3,-7,4)
\]

satisfies (3-7+4=0), and therefore belongs to the lattice

\[
 A_2=\ker\bigl(\mathbb Z^3\xrightarrow{x+y+z}\mathbb Z\bigr).
\]

The root \((A_1)\) governs transfer between the two projections; the defect \((A_2)\) governs redistribution among the three correlated windows. They are not the same vector and do not have the same rank, but they are two successive degrees of a single conservation law through augmentation. This explains why double reading and ternary orientation can coexist without turning the theory into binary coding: conservation fixes the diagonal, whereas the HMT datum lies in its oriented complement.

\textbf{Refutation criterion.} A scale whose frame does not have determinant (-N), a step changing both coordinates at once, or an identification of \((A_1)\) and \((A_2)\) that does not declare the rank-changing map.

\section{Topological--temporal boundary charge and non-destructive curvature selection}

The tritic derivative of the vacancy word,

\[
 \tau_n=z_{n-1}-z_n\in\{+1,0,-1\},
\]

has a conservation law completing its interpretation as a unit of topological--temporal information. For every finite interval \((a\leq n\leq b)\),

\[
 \boxed{
 Q_{[a,b]}:=\sum_{n=a}^{b}\tau_n
 =z_{a-1}-z_b\in\{-1,0,+1\}.}
\]

The oriented interior charge telescopes completely and is supported at the two endpoints. The visible TRIT does not accumulate an arbitrary extensive charge: it records how the block opens, persists or closes relative to its boundary. The TPK memory record preserves the interior terms concealed by the alternating sum, so telescoping does not amount to loss of history.

Since \((0<\delta<1/2)\), two vacancies are not consecutive. The nonzero symbols of \(((\tau_n))\) therefore alternate strictly:

\[
 -1,+1,-1,+1,\ldots
\]

when intermediate zeros are removed. An opening must close before another opening can occur. Zeros do not mean absence of dynamics: they are the plateaux during which the regime is preserved while the block clock continues advancing. Upon inducing the dynamics on events, these plateaux have lengths \(6\) or \(7\), as already proved.

Each symbol selects one of the three exponentials

\[
 \exp(tJ_\tau)=C_\tau(t)I+S_\tau(t)J_\tau,
 \qquad J_\tau^2=-\tau I.
\]

The word thus orders elliptic rotations, nilpotent transports and hyperbolic openings. Selection does not destroy aperiodicity because the differential map is invertible once one boundary bit is preserved:

\[
 z_n=z_0-\sum_{k=1}^{n}\tau_k.
\]

This is the exact reason why the TRIT can read three curvatures without replacing the word generating them. The tritic unit is the visible change; the TPK preserves the boundary datum, APP sheet, carry and operator order needed to invert the reading.

Time reversal

\[
 C_{\rm rev}(\tau_1,\ldots,\tau_n)
 =(-\tau_n,\ldots,-\tau_1)
\]

interchanges opening and closure, preserves the threshold and satisfies \((C_{\rm rev}^2=I)\). In the matrix fibre, it also reverses the order of the product

\[
 \mathcal A_N=exp(t_NJ_{\tau_N})\cdots
               \exp(t_1J_{\tau_1}).
\]

The double temporal direction is therefore not the reading of a list backwards. It is conjugation of a word, its signs and the noncommutative order of its transports, retaining the boundary that allows the state to be reconstructed.

\textbf{Refutation criterion.} Two consecutive openings in the induced word, an interval charge not supported on the boundary, or reconstruction dispensing with \((z_0)\).

\section{Nonstationary \(S\)-adic tower of the two clocks}

The continued fraction of the internal slope begins

\[
 \delta=[0;21,1,5,1,6,2,3,1,1,8,1,2,\ldots].
\]

For each positive coefficient (a), define the continuant operator

\[
 A(a)=\begin{pmatrix}a&1\\1&0\end{pmatrix},
 \qquad\det A(a)=-1.
\]

If \((p_n/q_n)\) is the (n)th convergent, then

\[
 A(a_1)\cdots A(a_n)
 =\begin{pmatrix}q_n&q_{n-1}\\p_n&p_{n-1}\end{pmatrix}.
\]

The first stages are

\[
 \begin{aligned}
 A(21)&=\begin{pmatrix}21&1\\1&0\end{pmatrix},\\
 A(21)A(1)&=\begin{pmatrix}22&21\\1&1\end{pmatrix},\\
 A(21)A(1)A(5)&=\begin{pmatrix}131&22\\6&1\end{pmatrix},\\
 A(21)A(1)A(5)A(1)&=\begin{pmatrix}153&131\\7&6\end{pmatrix}.
 \end{aligned}
\]

This proves that \(21/22\) and \(6/7\) are not two regularities placed side by side. The integers (21,22) are the first denominators of the Diophantine skeleton; (6,7) are the numerators appearing two inductions later:

\[
 \begin{aligned}
 131&=5\cdot22+21,&6&=5\cdot1+1,\\
 153&=1\cdot131+22,&7&=1\cdot6+1.
 \end{aligned}
\]

The primary gaps \(21/22\) and secondary returns \(6/7\) therefore belong to the same renormalization. The next convergents are

\[
 \frac{48}{1049},\quad
 \frac{103}{2251},\quad
 \frac{357}{7802},\quad
 \frac{460}{10053},\quad
 \frac{817}{17855},\quad
 \frac{6996}{152893},\quad
 \frac{7813}{170748}.
\]

Each consecutive pair preserves a primitive lattice cell:

\[
 \boxed{p_nq_{n-1}-p_{n-1}q_n=(-1)^{n-1}.}
\]

The first denominator in this tower divisible by \(108\) is \(170748\), and its appearance preserves the complete recurrence:

\[
 \boxed{
 170748=152893+17855,
 \qquad
 7813=6996+817.}
\]

At that stage,

\[
 \begin{pmatrix}
 170748&152893\\
 7813&6996
 \end{pmatrix}
\in GL_2(\mathbb Z),
 \qquad\det=-1.
\]

The tower therefore preserves the lattice unit even when the supercell grows large enough to close the finite clock. Closure of \(C_{108}\) does not make the irrational direction periodic: the rational approximant coincides with the word on every proper prefix and transfers one unit only at the final half-boundary.

Renormalization is also genuinely ordered. For \(a\ne b\),

\[
 \boxed{
 A(a)A(b)-A(b)A(a)
 =(a-b)\begin{pmatrix}0&1\\-1&0\end{pmatrix}.}
\]

For example, \(a=21\) and \(b=5\) produce the oriented defect

\[
 16\begin{pmatrix}0&1\\-1&0\end{pmatrix}.
\]

The rotation \((x\mapsto x+\delta)\) on the circle remains abelian. Noncommutativity appears in ordering the renormalization operators and, in the full fibre, in ordering the APP sheets and curvature generators. This distinction removes two possible errors: calling the base rotation nonabelian, or reducing all dynamics to an abelian rotation while forgetting the cocycle.

Since \(\delta\) is transcendental, it is not a quadratic irrational. By Lagrange’s theorem, its continued fraction is not eventually periodic. The continuant directive does not stabilize into a single periodic substitution: the crystal is \(S\)-adic and nonstationary. It nevertheless preserves linear complexity, zero entropy and unimodular determinants. This is a much stronger definition of “aperiodic crystal” than mere absence of repetition: renormalization exists at every scale, but there is no final renormalization cell repeating indefinitely.

\textbf{Refutation criterion.} A discrepancy in the coefficients, a determinant other than \(\pm1\), an earlier convergent denominator divisible by \(108\), or an eventually periodic directive.

\section{Cardinal genealogy of the corona, the boundary \(369\) and the dodecaphasic seal}

The tower allows the structural cardinalities to be resolved without selecting them by decimal resemblance. The following pairs proceed from the same reader and have independent derivations:

\[
 \begin{aligned}
 \mathcal W_\delta(217)&=\{9,10\},\\
 \mathcal W_\delta(271)&=\{12,13\},\\
 \mathcal W_\delta(601)&=\{27,28\},\\
 \mathcal W_\delta(973)&=\{44,45\}.
 \end{aligned}
\]

The integer \(217\) is the cardinality of the \(3\)-cells of the corona and also \(2\cdot108+1\). Its reading determines the boundary between the nonadic base and decimal completion. The integer \(271\) is the cubic corona \(1000-729\); its reading determines \(12/13\), precisely the pair formed by the dodecaphase and the number of atlas multisections. The integer \(601\) is the twelfth coordinate of the seal \(K\); its reading determines \(27/28\), the cubic boundary and the cardinality of the microfibre. Finally,

\[
 973=271+702=756+217
\]

is the common sum of the even and odd degrees of the corona, and

\[
 \mathcal W_\delta(973)=\{44,45\}.
\]

Here \(44=396/9\) is the reduced APP count and \(45\) is the determinant of the boundary matrix and the count of the radical region. The Euler balance of the continuum thus projects onto two cardinalities already situated in APP and the radical branch.

These equalities do not declare that a multisection is literally a vacancy or that \((K_{12})\) is a set of twenty-eight elements. The proved content is more precise: a single cardinal cut functor, constructed before these realizations, sends four typed objects to four internal pairs with known provenance. The next step is not to seek further resemblances, but to lift this cardinal functor to the actions preserved by each representation.

The boundary \(369\) admits an even more rigid genealogy. The construction distinguishes the cardinality \(386\) of the internal boundary from the \(387\) pairs \((K,K+9)\) per channel. The reader gives

\[
 \mathcal W_\delta(386)=\mathcal W_\delta(387)=\{17,18\}.
\]

Using opposite closures,

\[
 \boxed{386=17+369,\qquad387=18+369.}
\]

Thus, \(369\) is the common capacity surviving transfer of the boundary unit between two consecutive cardinalities. The transition

\[
 (386,17)\longleftrightarrow(387,18)
\]

increments both the ambient space and the vacancy, but leaves the radical capacity \(369\) fixed. This is an exact definition of stability under change of half-boundary. It complements, without replacing, the already proved genealogies

\[
 369=396-27,
 \qquad
 369=3(132-6)-9,
 \qquad
 369+126=495.
\]

The convergence of four independent routes—adjacent boundary, APP square, induced returns and conservative transfer—makes \(369\) an internal node of the architecture, not a chosen decimal string.

Other cardinal projections retain a distinct status:

\[
 \mathcal W_\delta(756)=\{34,35\},
 \qquad
 \mathcal W_\delta(702)=\{32,33\}.
\]

The \(34\) coincides with the action depth already derived and the \(33\) with the lower count of \(729\). The physical derivation of the decade \(-34\) belongs to the cokernel of \(H_4\) and to \(\varphi\alpha^{16}\); reading the \(756\) edges does not replace it. The same corona producing the Euler trit thus contains an additional cardinal projection of that depth, without identifying the edge reader with the action line.

The unimodular frame of the supercell also acquires a further strong algebraic interpretation:

\[
 M_\star=
 \begin{pmatrix}601&217\\36&13\end{pmatrix},
 \qquad\det M_\star=1.
\]

Its four entries are now the dodecaphasic coordinate \((K_{12}=601)\), the cardinality of the \(3\)-cells of the corona \(217\), the stage \(R_{36}\) and the thirteen multisections. The identity \(217=216+1\) remains compatible with this reading. The unit determinant proves lattice conservation of the frame formed by the four cardinalities, without identifying their domains.

\textbf{Refutation criterion.} A change in any window, failure of \(369\) to be the common complement, or a proposed map preserving cardinalities while destroying incidence, action or memory.

\section{The crystal kernel and the continuum as output}

The preceding components allow the governing thesis to be specified without confusing the result with a cosmological metaphor. APP supplies a finite support with two inseparable evaluations. The additive sheet accumulates; the multiplicative sheet folds; residue, quotient and carry preserve the difference between their histories. TRIT extracts the local orientation of that difference and realizes it in three quadratic types. TPK turns local signatures into dynamics through selection, transport and updating, preserving phase, sheet, boundary and memory record. The completion \(729\to1000\) then induces a digital line of transcendental slope; its vacancy word is aperiodic, balanced and of zero entropy; its derivative is tritic; its renormalization tower is \(S\)-adic; and its suspension over \(C_{108}\) is helical.

The continuum does not enter this chain as a previously available \(\mathbb R\). Compatible finite cylinders, their two one-sided closures, the boundary unit and the refinement maps exist first. The cellular corona proves that double projection separates an object of zero charge into branches of charges \(+1\) and \(-1\) without altering their sum. The inverse limit preserves compatible histories; the Archimedean reader subsequently determines a real coordinate. In this sense, \(\mathbb R\) is a terminal realization of an enriched discrete protocol, not the ontological basis from which APP is reconstructed.

The two clocks are now unambiguous. The chronological clock records the number of applications of \(U_t\) and has the finite quotient \(C_{108}\). The capacity clock records how often refinement has opened an effective decimal pruning. Their difference is the integrated vacancy memory. A turn of \(108\) steps can close the first clock while leaving the second, the irrational phase and memory open. The supercell of length \(170748\) is the first convergent that closes the first clock exactly; the boundary unit remaining at the endpoint proves that this closure does not make the crystal periodic.

The correlated system \(\pi,\varphi,e,\alpha\) resides on this same coinductive state. The windows of \(\pi,e,\varphi\) synchronize and reopen; the oriented defect lies in \(A_2\); the second route to \(\alpha\) receives the pairs \(21/22\), \(6/7\) and \(45/46\) from three scales of the same word; the first route preserves the dodecaphasic seal. No conventional coordinate generates the others. Dodecaphasic closure determines \(\alpha\) within the correlated output, and the Machin, Perron--Fibonacci or exponential-semigroup characterizations enter only afterwards as recognitions.

The electronic vacancy also acquires a precise provenance. On the APP support, the transition from the centre \((5,5)\) to the branch \((8,2)\) preserves the additive sheet, product residue and TRIT regime while losing one unit of multiplicative carry. In the temporal word, a vacancy preserves the capacity phase while the block and memory advance. The HMT--MD correspondence jointly preserves this unit loss, the fixed centre, the mode \(m_{\rm ph}\), the incidence \(90/120\) and band conjugation. Electricity determines the oriented vacancy and magnetism its conjugate rotational transport; the common and oriented constitutive responses are obtained from the two sheets before being written in a dimensional chart.

The depths \(11\) and \(34\) appear in a common completion signature. The action line has its determinantal derivation and the gravitational branch first receives the length \(L=(5759/23040)\alpha^{16}L_*\). Marked inscription and the positive radial metric prove \(R_X\Lambda_X=L^2I\), \(H_X\Lambda_X=\hbar_{\rm ret}cI\) and \(G=c^3L^2/\hbar_{\rm ret}\), with the rules and domains of \cref{g26:sec:rigidez-radial-es}. The functional restitution of \cref{cr28:sec:restitucion-constitutiva} links the Catalan moment to the responses of the two sheets and preserves their common propagation. Each constant is a section or evaluation of the same predefined state and preserves the reader determining its dimension, normalization and scale. The Barbero--Immirzi parameter, Catalan constant, Feigenbaum constants, permittivity and permeability are distinct readers of a common genealogy; they are therefore coordinated without becoming a single function or a list of juxtaposed values.

The second projection preserves boundary incidence. The connection \(24\) in the splitting \(729=202+527\), the pairs \(12/13\) and \(27/28\), the unimodular frame \(M_\star\), the Witt design and the Golay and Leech gluings compose the typed passage to \(V^\natural\) and the Monster group. Exceptional symmetry is therefore the boundary realization of the same discrete structure that produces the Archimedean coordinate through the other projection; form, incidence, orientation and carry remain visible throughout transport.

The kernel reached can be summarized in a single mathematical statement without losing its architecture: the unit propagates through two complementary projections; the boundary difference generates orientation; the derivative of that orientation selects three regimes; the order of their transports generates a noncommutative fibre; and unlimited refinement preserves the unit, memory and aperiodicity while producing its continuous realizations. This is the precise content of the conservative evolution of the dimensional unit through double holographic projection and fractal refinement.

The formulation gathers a single constructive chain. Its arithmetic, physical and exceptional realizations preserve different domains and readers, but none is placed on a lower scale of certainty: each is the exact image of the integral holonomic state in its declared codomain. External comparison neither generates nor selects these outputs.

\section{Exhaustive morphology of the \(104\,976\) seeds}

The cardinal correction is not minor: \(104\,976\) and \(170\,748\) belong to distinct mathematical objects. The first is the exhaustive domain of initial conditions of the APP--TPK emitter; the second is the denominator of the first continuant approximation closing the chronological clock \(C_{108}\). Neither can replace the other.

Let

\[
 I_9=\{0,1,\ldots,8\},\qquad
 D=\{N,E,S,O\},\qquad
 \mathcal S=I_9^2\times D.
\]

Here \(I_9\) is exclusively the positional chart of the nine places, indexed from zero. It does not replace the canonical arithmetic alphabet \(\mathcal D_9=\{1,\ldots,9\}\) on which APP evaluates sum and product; passage between the two charts is a relabelling of positions, not an additional physical operator.

APP causally precedes this construction. An element \(s=(i,j,d)\in\mathcal S\) does not generate APP: it fixes a cursor position and orientation on the APP support already constructed. Since additive and multiplicative evaluation must be preserved simultaneously, a complete seed is the ordered pair

\[
 \omega=(s_+,s_\times)\in
 \Omega_{\rm seed}:=\mathcal S_+\times\mathcal S_\times.
\]

Therefore,

\[
 |\mathcal S|=9^2\cdot4=324,
 \qquad
 \boxed{|\Omega_{\rm seed}|=324^2=104\,976=2^4 3^8.}
\]

The seed is not a bare integer. It is a six-field oriented cell,

\[
 (i_+,j_+,d_+;i_\times,j_\times,d_\times),
\]

which keeps the two arithmetic sheets and their directions separate. This is the finite unit on which TRIT types the local regime and TPK applies selection, transport and updating.

The exhaustive sweep reveals a factorization stronger than the count alone. The additive signature induces exactly \(18\) classes, all of cardinality \(18\). The multiplicative signature induces \(26\) classes, but its fibres stratify into

\[
 24\text{ size classes }8,\qquad
 1\text{ size class }24,\qquad
 1\text{ size class }108.
\]

If \(\mathcal A_{18}\) and \(\mathcal M_{26}\) denote the respective signature sets, the emitter factors as

\[
 \Omega_{\rm seed}
 \longrightarrow \mathcal A_{18}\times\mathcal M_{26}
 \xrightarrow{\;\mathrm{Comb}\;}\mathcal U_6^{\rm cat},
 \qquad
 |\mathcal U_6^{\rm cat}|=18\cdot26=468.
\]

The first arrow preserves the provenance of each sheet; the second combines the two signatures into a decimal emission. The fibre multiplicities are the exact products

\[
 18\cdot8=144,\qquad
 18\cdot24=432,\qquad
 18\cdot108=1944,
\]

and the global closure is

\[
 \boxed{432\cdot144+18\cdot432+18\cdot1944=104\,976.}
\]

This asymmetry is structural. The additive sheet has a uniform partition; the multiplicative sheet contains two singular strata. Vacancy and torsion are not introduced afterwards as metaphors: the seed domain already distinguishes uniform accumulation from folding with fibre degeneracies. The TPK preserves precisely which of these fibres was traversed, together with residue, quotient, carry, orientation and boundary.

After this factorization, the previously certified reductions are applied

\[
 104,976\longrightarrow468\ U_6
 \longrightarrow243\ w_6\subset\mathbb F_3^6,
 \qquad |\mathbb F_3^6|=729.
\]

Four types are thus distinguished and must not be confused again: seed pairs, decimal emissions, visible words and tritic ambient space. The crystal structure is none of these cardinalities in isolation, but the system of fibres and morphisms linking them.

\textbf{Refutation criterion.} An additive class of size other than \(18\), a multiplicative histogram other than \(24\cdot8+1\cdot24+1\cdot108=324\), or loss of the ordered sheet pair in passage to \(U_6\).

\section{Bipartite geometry of the \(\pi\) microfibre and electronic spectral connection}

The visible output \(w_\pi=010211\) has \(1008\) antecedents in \(\Omega_{\rm seed}\). To recover its form and not merely its mass, define the bipartite graph

\[
 \mathcal G_\pi=(P_\pi\sqcup T_\pi,E_\pi),
\]

where \(P_\pi\) contains the additive-sheet seeds participating in the microfibre, \(T_\pi\) those of the multiplicative sheet, and \((p,t)\in E_\pi\) if the pair produces one of the five emissions reducing to \(w_\pi\). The exhaustive sweep proves

\[
 |E_\pi|=1008,\qquad |P_\pi|=54,\qquad |T_\pi|=56.
\]

The structure is not an arbitrary graph. Forgetting only the emission colour, its connected components are exactly

\[
 \boxed{
 \mathcal G_\pi\simeq
 K_{18,24}\sqcup K_{18,24}\sqcup K_{18,8}.}
\]

One of the two components \(K_{18,24}\) corresponds to the region of multiplicity \(432\). The other refines into three complete rectangles of multiplicity \(144\), since the three multiplicative classes of size eight share the same additive class of size eighteen. The last component is the fourth rectangle of multiplicity \(144\). The region colour therefore recovers

\[
 \boxed{1008=432+4\cdot144}
\]

as a decomposition of incidence operators, not merely as a sum of integers.

Let \(B_R\) be the biadjacency matrix of a rectangular region \(R\). For a complete graph \(K_{m,n}\), \(B_R\) has rank one and its only nonzero singular value is

\[
 \|B_R\|=\sqrt{mn}.
\]

The region of \(432=18\cdot24\) routes has norm \(12\sqrt3\); each of the four regions of \(144=18\cdot8\) routes has norm \(12\). The normalization treating the four smaller regions symmetrically is determined by the regional partition itself:

\[
 \Xi_\pi:=
 \frac{\|B_{432}\|}{\sum_{r=1}^{4}\|B_{144}^{(r)}\|}
 =\frac{12\sqrt3}{4\cdot12}
 =\boxed{\frac{\sqrt3}{4}}.
\]

Its square is \(3/16\). This is exactly the invariant that appears, in the previously certified APP electronic mode, as

\[
 \|[P_{\Sigma,3},P_{\Pi,3}]\|^2
 =s^2(1-s^2)=\frac14\frac34=\frac3{16}.
\]

This yields an \textbf{exact internal spectral connection} between two constructions arising from APP: the regional contrast of the \(\pi\) microfibre and the noncommutativity of the two electronic projections. Neither the electron mass nor any experimental parameter has been used. Nor is it asserted that the matrices \(B_R\) and the commutator act on the same space. The strong statement already proved is equality of their canonical scalars; a subsequent operator equivalence must transport the regional colouring, the mode \(m_{\rm ph}\), the sheet and the memory record.

This geometry also specifies the place of vacancy. On the additive side, the classes are uniform; on the multiplicative side, the degree changes from \(24\) to \(8\). Porosity of the support is not an inserted zero: it is the incidence stratification surviving the common reduction \(w_\pi\). The electric realization can read that oriented loss and the magnetic realization its conjugate transport, but both must preserve this original incidence.

\textbf{Refutation criterion.} A component ceasing to be complete, the five colours not having sizes \((432,144,144,144,144)\), or a spectral ratio other than \(\sqrt3/4\). Operator equality, which is not asserted here, would be falsified by any transport not preserving colour, sheet and inner product.

\section{Three-channel morphism, tritic gluing class and diagonal memory}

The count of the nine nonadic positions defines a vector

\[
 N_K=(N_\pi(K),N_e(K),N_\varphi(K))\in\mathbb Z^3.
\]

The relational part is obtained through the morphism

\[
 \partial:\mathbb Z^3\longrightarrow A_2,
 \qquad
 \partial(x,y,z)=(x-y,y-z,z-x),
\]

where

\[
 A_2=\{(a,b,c)\in\mathbb Z^3:a+b+c=0\}.
\]

The kernel is the diagonal \((\Delta\mathbb Z=\mathbb Z(1,1,1))\), and \(\partial\) is surjective. Consequently,

\[
 \boxed{0\longrightarrow\mathbb Z
 \xrightarrow{\Delta}\mathbb Z^3
 \xrightarrow{\partial}A_2\longrightarrow0.}
\]

Synchronization of the three channels is exactly the vanishing of \(\partial N_K\). The quotient removes the common quantity, but preserves which channel exceeds which and with what orientation.

Integer reconstruction contains additional information. Define

\[
 \mathcal J(x,y,z)=(s,a,b)
 =(x+y+z,x-y,y-z).
\]

Over \(\mathbb Q\), the three data reconstruct \((x,y,z)\) through

\[
 x=\frac{s+2a+b}{3},\qquad
 y=\frac{s-a+b}{3},\qquad
 z=\frac{s-a-2b}{3}.
\]

Over \(\mathbb Z\), the image is not all of \(\mathbb Z\oplus\mathbb Z^2\): it satisfies the congruence

\[
 \boxed{s-a+b\equiv0\pmod3.}
\]

The determinant of \(\mathcal J\) has absolute value \(3\). There is therefore an exact sequence

\[
 0\longrightarrow\mathbb Z^3
 \xrightarrow{\mathcal J}\mathbb Z^3
 \xrightarrow{\chi_3}C_3\longrightarrow0,
 \qquad
 \chi_3(s,a,b)=[s-a+b]_3.
\]

This result is particularly important for the meaning of TRIT. Separating the common content from the relative defect leaves a gluing class with exactly three values. A trit is the minimum amount of information needed to record that class and select the integral correction recomposing the three-channel state. The topological--temporal interpretation is not added by analogy: the class types the local gluing and TPK transports it in time while preserving its orientation and memory. This quotient \(C_3\) is a precise internal realization of that function; it does not purport to exhaust all coordinates of the governing TRIT.

The passage from \(K=8\) to \(K=9\) shows why synchronization does not mean state return:

\[
 N_8=(59,59,59),\qquad N_9=(43,43,43),
\]

\[
 \partial N_8=\partial N_9=0,
 \qquad
 \boxed{N_9-N_8=-16(1,1,1).}
\]

The state remains at the relational origin \(A_2\), but the diagonal memory record changes by sixteen units per channel, forty-eight in total. The synchronization phase returns; memory does not. This equality is the entirely integral three-channel analogue of the TPK rule that phase return does not reset the integral holonomic state.

\textbf{Refutation criterion.} A count violating the congruence, a kernel of \(\partial\) larger than the diagonal, or an equality \(N_8=N_9\).

\section{Wall itinerary and finite quotient of the synchronization matrix}

Between \(K=4\) and \(K=9\), the preceding morphism turns the complete table into a geometric itinerary:

\[
\begin{array}{c|c|c}
K&N_K&\partial N_K\\
\hline
4&(207,207,122)&(0,85,-85)\\
5&(39,151,151)&(-112,0,112)\\
6&(110,110,69)&(0,41,-41)\\
7&(81,29,81)&(52,-52,0)\\
8&(59,59,59)&(0,0,0)\\
9&(43,43,43)&(0,0,0).
\end{array}
\]

The successive equalities traverse the walls

\[
 H_{\pi e}\longrightarrow H_{e\varphi}
 \longrightarrow H_{\pi e}\longrightarrow H_{\varphi\pi}
 \longrightarrow0\longrightarrow0.
\]

The wall word \((H_{12},H_{23},H_{12},H_{13})\) has the shadow of the Coxeter pattern \(A_2\), since \(s_{12}s_{23}s_{12}=s_{13}\). This does not automatically turn every TPK transition into a reflection: the step also changes the diagonal component. What is proved is that the relational trajectory crosses the three wall systems in the order of the braid relation, while its lift preserves memory invisible to the quotient.

The amplitudes of the four walls preceding closure form

\[
 D_{\rm syn}=\begin{pmatrix}85&112\\41&52\end{pmatrix}.
\]

In addition to the identities already gathered

\[
 \operatorname{tr}D_{\rm syn}=137,
 \qquad
 \det D_{\rm syn}=-172=-4\cdot43,
\]

the Smith normal form adds quotient structure. Since the greatest common divisor of the four entries is one,

\[
 \operatorname{SNF}(D_{\rm syn})=\operatorname{diag}(1,172),
\]

and, therefore,

\[
 \boxed{\operatorname{coker}D_{\rm syn}\simeq C_{172}
 \simeq C_4\times C_{43}.}
\]

The factorization \((4\cdot43)\) no longer appears only in the determinant: it is the primary decomposition of the lattice quotient of the wall amplitudes. The right null directions can be chosen as

\[
 (0,1)\pmod4,
 \qquad
 (1,5)\pmod{43}.
\]

The factor \(43\) coincides with the count of the second triple closure, and the factor \(4\) with the oriented coefficient obtained internally as \((-\det(D_{\rm syn})/43)\). This proves that both are canonical divisors of the same quotient. A subsequent physical representation must specify which subspaces realize these two factors; the joint arithmetic structure is already closed.

The itinerary thus supplies a complete mathematical picture: a path in the wall arrangement of \((A_2)\), lifted by a diagonal memory coordinate and governed by a finite quotient \(C_4\times C_{43}\). It is more informative than a table of coincidences because it distinguishes relation, amplitude, orientation and the state invisible to the quotient.

\textbf{Refutation criterion.} A determinant other than \(-172\), a Smith form other than \((1,172)\), or a transition erasing the diagonal component needed to distinguish \(K=8\) from \(K=9\).

\section{Nonadic Pascal tower and infinite concatenation law with carry}

Cubic completion belongs to an exact system of arbitrary depth. For \(n\geq1\), write

\[
 10^n=(9+1)^n=\sum_{j=0}^{n}B_{n,j},
 \qquad
 B_{n,j}:=\binom nj9^{n-j}.
\]

The blocks obey the weighted Pascal recurrence

\[
 \boxed{B_{n+1,j}=9B_{n,j}+B_{n,j-1}},
\]

with \(B_{n,j}=0\) outside \(0\leq j\leq n\). This is the length-refinement law; the vacancy law must also preserve the initial phase. If

\[
 s_{n,j}=\sum_{i<j}B_{n,i},
 \qquad
 v^-_{n,j}=V_{B_{n,j}}(s_{n,j}),
\]

then

\[
 \sum_{j=0}^{n}v^-_{n,j}=\lfloor10^n\delta\rfloor.
\]

The upper reading transports the half-boundary unit to the first block:

\[
 v^+_{n,0}=v^-_{n,0}+1,
 \qquad
 v^+_{n,j}=v^-_{n,j}\quad(j>0),
\]

and therefore

\[
 \boxed{\sum_jv^+_{n,j}=\lceil10^n\delta\rceil
 =\lfloor10^n\delta\rfloor+1.}
\]

The transition \(n\mapsto n+1\) cannot be linearized by forgetting the starting point. For any block \(B\),

\[
 V_{9B+B'}(s)=
 \sum_{r=0}^{8}V_B(s+rB)+V_{B'}(s+9B).
\]

The shift of \(s\) in each summand is the concatenation memory. Projecting all blocks to the origin reveals the binary cocycle

\[
 \kappa_\delta(M,N)=
 \lfloor(M+N)\delta\rfloor-
 \lfloor M\delta\rfloor-
 \lfloor N\delta\rfloor,
\]

which satisfies

\[
 \kappa_\delta(L,M)+\kappa_\delta(L+M,N)
 =\kappa_\delta(M,N)+\kappa_\delta(L,M+N).
\]

This is the exact law preventing refinement from depending on parenthesization. The Pascal tower supplies the lengths; the cocycle supplies the carry; TPK preserves the phase and memory record that allow their composition.

The first levels are

\[
\begin{array}{c|c|c|c}
n&(B_{n,0},\ldots,B_{n,n})&v_n^-&v_n^+\\
\hline
2&(81,18,1)&(3,1,0)&(4,1,0)\\
3&(729,243,27,1)&(33,11,1,0)&(34,11,1,0)\\
4&(6561,2916,486,36,1)&(300,133,22,2,0)&(301,133,22,2,0).
\end{array}
\]

The cubic level is therefore not an isolated coincidence. It is the first level at which the component \(729=3^6\), the subscale \(243=3^5\), the boundary \(27=3^3\) and the unit appear simultaneously in a decimal completion. Its lower signature contains \(33=3\cdot11\), and the upper one preserves the unit to produce \(34\). The next level proves that the rule continues without resetting the state: the unit again modifies only the first component, while the remaining entries preserve the history of the ordering.

The concatenation \(210=90+120\) is a section of the same law:

\[
 (V_{90}(0),V_{120}(90))=(4,5),
\]

whence the lower reading produces \(9\) and the upper one \(10\). The nonadic--decimal transition is thus realized by a reading morphism with carry, not by the graphic observation \(4+5=9\).

Define the decimal carry digit

\[
 d_{n+1}:=\lfloor10^{n+1}\delta\rfloor
          -10\lfloor10^n\delta\rfloor.
\]

Since \(\delta\) is transcendental, the sequence \((d_n)_{n\geq1}\) cannot be eventually periodic: such periodicity would make the decimal expansion of \(\delta\) rational. Since the mechanical word is Sturmian, every level retains linear complexity and zero topological entropy. The tower therefore combines unlimited refinement, the boundary unit and exact aperiodicity. This is the mathematical sense in which the crystal simultaneously grows outwards and refines inwards.

One must distinguish \(\kappa_\delta\) from the governing nonadic cocycle \(\kappa_9\). The former is the mechanical restriction recording the carry of the irrational cut; the latter preserves the full APP carry under composition. A restriction square may relate them, but they are not identified merely because both use carry values.

\textbf{Refutation criterion.} Failure of the Pascal recurrence, a nontelescoping signature, loss of the upper unit, or concatenation whose result depends on parenthesization.

\section{Mathematical type of the helical aperiodic crystal}

The results allow the object to be named without reducing it to an image. Its finite layer of initial conditions is \(\Omega_{\rm seed}\); its emission quotient is \(\mathcal U_6^{\rm cat}\); its visible alphabet lies in \(\mathbb F_3^6\); and its compatible extensions form the enriched coinductive state. Acting on this state, in the governing order, are

\[
 U_t=\operatorname{Upd}_t\circ\operatorname{Tra}_t
 \circ\operatorname{Sel}_t,
 \qquad
 U_t\longrightarrow\Gamma_9\longrightarrow K_{\rm ph}.
\]

The holonomy \(\Gamma_9\) belongs to the enriched dynamics; the reduced memory \(K_{\rm ph}\) is its subsequent realization. The aperiodic factor can be written on

\[
 \mathcal Y_{108}=C_{108}^{(t)}\times\mathbb T_\delta\times\mathbb Z,
\]

through means of

\[
 \mathcal F(r,\theta,m)=
 \bigl(r+1,\{\theta+\delta\},
 m+\lfloor\theta+\delta\rfloor\bigr).
\]

The first coordinate is the block clock; the second, the irrational phase; the third, the integrated vacancy memory. The capacity clock is obtained as

\[
 K_t=t-Z_t,
 \qquad
 g_t=[K_t]_9,
\]

and does not coincide with the chronological clock. The inverse dynamics \(\mathcal
F^{-1}\) constitutes the descending reading of the same helix. Ascent and descent are conjugate; they are not two disconnected generators.

The complete object is therefore an invertible affine extension of an irrational rotation over a finite clock, fibred by the integral holonomic APP--TRIT--TPK state. Its symbolic shadow is Sturmian; its suspension is helical; its fibre cocycle may be noncommutative even though the base rotation is abelian. The designation “helical aperiodic crystal” is justified by four proved properties:

\[
 \begin{gathered}
 \text{long-range order and repetitivity},\\
 p(m)=m+1,\qquad h_{\rm top}=0,\\
 \text{absence of translational period},\\
 \text{finite phase return with memory advance}.
 \end{gathered}
\]

In HMT, the expression \textbf{aperiodic time crystal} designates precisely the preceding skew product: there is a finite temporal support, phase return is ordered, memory advances in each holonomy and the occupation word has no translational period despite preserving long-range order and zero topological entropy. Temporal dynamics is not added to the crystal afterwards: it is contained in the coordination of clock, orientation, return and memory defining its state.

The correlated constant output is situated after the extension

\[
 w_6\to w_{12}\to w_{18}\to w_{24}\to w_{30}
 \to R_{36}\to G_9,
\]

without introducing conventional values. \(\pi,\varphi,e\) and the dodecaphasic coordinate \(\alpha\) are realizations of the same coinductive state; their synchronization windows, mirror and own readers are not generator inputs. The completion tower supplies the capacity calendar and the carry law by which these realizations can be extended to arbitrary precision.

Finally, \(\mathbb R\) appears through the compatible-cylinder reader of the joint continuum. The crystal is not “inside \(\mathbb R\)” as an embellishment of a line given beforehand. Compatible histories, double projection, boundaries and their limits are constructed first; one of the realizations subsequently determines the Archimedean coordinate. This reversal of perspective is the strong thesis: the real line is the terminal shadow of a finite oriented unit that is conserved, refined and accumulates memory.

The aperiodic dynamics, its complexity, clocks and cocycle constitute the mathematical kernel of the crystal. The electromagnetic, gravitational and exceptional realizations are typed readers of that kernel: they preserve the APP--TRIT--TPK genealogy and express the same conservative evolution of the unit in different codomains.

\section{Coinductive predefinition: the output precedes its evaluation}

The decisive distinction separates generation from evaluation. HMT--MD does not take an external measurement to infer the law that might have produced it. It first constructs the complete section, fixes its compatible extensions and only then allows a reader to express a coordinate of that section. This is the exact mathematical meaning of \textbf{predefinition}.

Let

\[
 (X_n,\rho_{n+1,n})_{n\geq0}
\]

the inverse system of APP--TRIT--TPK states truncated at depth \(n\), where each \(X_n\) preserves sheet, orientation, residue, quotient, carry, boundary, route and memory up to that level. The coinductive extension already constructed defines

\[
 \mathcal G_{\rm HMT}:\Omega_{\rm seed}\longrightarrow
 X_\infty^{\rm enr}:=\varprojlim_n X_n,
 \qquad
 p_n:X_\infty^{\rm enr}\longrightarrow X_n,
\]

with \(\rho_{n+1,n}p_{n+1}=p_n\). For a system of outputs \((Y_n,\sigma_{n+1,n})\) and readers \(L_n:X_n\to Y_n\), the condition expressing that the reader does not destroy genealogy is the square

\[
\begin{CD}
X_{n+1}@>{L_{n+1}}>>Y_{n+1}\\
@V{\rho_{n+1,n}}VV @VV{\sigma_{n+1,n}}V\\
X_n@>>{L_n}>Y_n .
\end{CD}
\]

Therefore,

\[
 \operatorname{Pred}_{Y,n}
 :=L_n\circ p_n\circ\mathcal G_{\rm HMT}
 :\Omega_{\rm seed}\longrightarrow Y_n
\]

satisfies

\[
 \sigma_{n+1,n}\circ\operatorname{Pred}_{Y,n+1}
 =\operatorname{Pred}_{Y,n}.
\]

The system \(\bigl(\operatorname{Pred}_{Y,n}(\omega)\bigr)_{n\geq0}\) is then a single element of \(\varprojlim Y_n\). The output is not decided anew each time depth increases: the depth-\(n\) output is the restriction of the depth-\(n+1\) output. Precision no longer generates the object; it merely decides how much of the predefined object is determined.

This provides an internal anticipation theorem. If, at stage \(n\), the coordinate \(n+k\) of the same compatible section is consulted, a future output is obtained without extrapolating a trend or introducing the value one wishes to find. The anticipation is

\[
 \operatorname{Ant}_{n,k}(\omega)
 :=L_{n+k}p_{n+k}\mathcal G_{\rm HMT}(\omega),
\]

and its subsequent comparison must restrict to \(L_np_n\mathcal G_{\rm HMT}(\omega)\) through \(\sigma_{n+k,n}\). The compatibility identity proves that both belong to the same section. Consequently, the correct conceptual transition is not

\[
 \text{more measure}\longrightarrow\text{more precision}
 \longrightarrow\text{best prediction},
\]

but

\[
 \boxed{
 \begin{gathered}
 \text{oriented seed}\longrightarrow
 \text{predefined coinductive section}\longrightarrow\\
 \text{reading at the chosen depth}\longrightarrow
 \text{testable anticipation}
 \end{gathered}}
\]

Indefinite digit generation is an instance of this theorem. For the readers of \(\pi\), \(\varphi\), \(e\) and \(\alpha\), each decimal block is a coordinate of a compatible system; extending from one thousand to ten thousand digits neither changes the generator nor resets selection. The new window restricts to the preceding one and continues the same history. The correlated system is a single correlated output of the coinductive state, even though its four coordinates have internal readers and different realization stages.

TPK reversibility strengthens the result. Since every update \(\mathcal U_t\) is bijective on the integral holonomic state, for every finite distribution transported by it,

\[
 H(X_{t+1}\mid X_t)=H(X_t\mid X_{t+1})=0,
 \qquad H(X_{t+1})=H(X_t).
\]

This is the zero Landauer cost of the complete state: no information is erased in updating. It does not assert that every reduced entropy is zero; it asserts that the complete history can be recovered in both directions. Apparent loss arises only when a reader contracts the fibre and forgets coordinates preserved by the integral holonomic state.

\section{TRIT as the topological--temporal unit combining regime and orientation}

TRIT is neither a ternary digit added to APP nor a temporal label superimposed on TPK. It is the local unit allowing a difference between the additive and multiplicative sheets to acquire geometric type, temporal orientation and lifting memory simultaneously. Let

\[
 \mathcal T=\{-1,0,+1\},\qquad
 x=(t,\mu,o,h,c,\sigma,\lambda,g)\in X_{\rm TPK}^{\rm enr}.
\]

Two coordinated readers act on the same state,

\[
 \tau_{\rm g}:X_{\rm TPK}^{\rm enr}\to\mathcal T,
 \qquad
 \varepsilon_{\rm t}:X_{\rm TPK}^{\rm enr}\to\mathcal T.
\]

The first types the local quadratic regime through

\[
 \mathcal A_\tau=\mathbb R[J_\tau]/(J_\tau^2+\tau I),
\]

so that

\[
\begin{array}{c|c|c}
\tau_{\rm g}&J_{\tau_{\rm g}}^2&\exp(sJ_{\tau_{\rm g}})\\
\hline
+1&-I&\cos s\,I+\sin s\,J_{+}\\
0&0&I+sJ_0\\
-1&+I&\cosh s\,I+\sinh s\,J_{-}.
\end{array}
\]

The second types the active temporal orientation:

\[
 \begin{aligned}
 +1&\longleftrightarrow\text{return, memory and past},\\
 0&\longleftrightarrow\text{threshold and present},\\
 -1&\longleftrightarrow\text{bidirectional opening and future}.
 \end{aligned}
\]

The conceptual depth lies in both readers arising from the same state, not in their having to be identified term by term. The correct object is the fibred coupling

\[
 \Theta_{\rm ITT}(x)
 =\bigl(\tau_{\rm g}(x),\varepsilon_{\rm t}(x),
        h(x),c(x),\lambda(x),g(x)\bigr),
\]

which may be called the \textbf{unit of topological--temporal information}. Geometry indicates how the unit transforms locally; temporal orientation indicates the direction in which that transformation is read; carry and record prevent two returns with the same phase from being confused. Thus, the past is the closure remaining in memory, the present is the threshold at which continuation is decided, and the future is the already predefined system of compatible extensions.

TPK makes this unit dynamic through

\[
 \mathcal U_t=\operatorname{Upd}_t\circ
 \operatorname{Tra}_t\circ\operatorname{Sel}_t.
\]

\(\operatorname{Sel}_t\) selects the regime and sheet without erasing the other APP evaluation; \(\operatorname{Tra}_t\) transports direction, phase and boundary; \(\operatorname{Upd}_t\) updates residue, quotient, carry, signature and memory. A chronological turn can therefore return \(g(t+9)=g(t)\) while simultaneously producing

\[
 q_{\rm mem}(\Gamma_{9,t}x)=q_{\rm mem}(x)+1.
\]

Temporal conjugation is represented by the oriented inversion

\[
 \mathcal C_{\rm rev}(\tau_1,\ldots,\tau_n)
 =(-\tau_n,\ldots,-\tau_1),
 \qquad \mathcal C_{\rm rev}^2=I.
\]

In the helical suspension, this involution interchanges ascending and descending dynamics:

\[
 \mathcal C_{\rm rev}\,\mathcal F\,
 \mathcal C_{\rm rev}^{-1}=\mathcal F^{-1}.
\]

History does not disappear: its orientation changes while the record enabling its reconstruction is preserved. The HMT time crystal is consequently the dynamics of an ITT returning in the finite clock and progressing in memory, while the Sturmian decoration prevents translational repetition of the full state. The clock \(C_{108}\) supplies finite temporal order; the holonomy \(\Gamma_9\), unlimited lifting; the Sturmian word, exact aperiodicity; TRIT, the joint discrimination of geometry and orientation; and TPK, the law preserving all of this under evolution.

\section{The singularity of \(\alpha\) as the seam of the two clocks}

The function of \(\alpha\) becomes transparent once it is no longer treated as an isolated decimal. \(\alpha\) is the seam coordinate making finite dodecaphasic closure compatible with aperiodic capacity extension. The first clock supplies a signed word of twelve sectors; the second supplies a unique root in the mechanical dynamics governed by

\[
 \delta=\log_{10}\frac{10}{9}.
\]

In the first route, the dodecaphasic state produces

\[
 K=\frac14\mathcal H_{12}U_{12}^{\rm sgn},
 \qquad \mathcal H_{12}^2=4I_{12},
\]

and the base carry \(1000\) satisfies, block by block,

\[
 P_m+E_m-\Phi_m-K_m+c_{m+1}=A_m+1000c_m.
\]

Telescoping concatenation gives

\[
 \boxed{\iota(P)+\iota(E)-\iota(\Phi)-\iota(K)=\iota(A).}
\]

The coordinates \(P,E,\Phi\) are the HMT realizations of \(\pi,e,\varphi\); they are not conventional values introduced from outside. The vector \(K\) is the signed memory of closure, and \(A\) is the finite section determining \(\alpha\) in the declared chart. The term distinguishing \(\alpha\) from a bare combination of \(\pi,e,\varphi\) is precisely \(K\): the dodecaphasic memory of how concatenation was performed.

In the second route, the same vacancy arithmetic produces the jet

\[
\begin{aligned}
P_6(x)={}&2\pi x-\frac74x^2+\frac{x^3}{2\pi}
+\frac{x^4}{20}-\frac{2x^5}{21}-\frac{x^6}{46},\\
T_{7:9}(x)={}&-\frac{x^7}{120}+\frac{x^8}{45}
+\frac{2x^9}{495},\\
P_9(x)={}&P_6(x)+T_{7:9}(x),
\end{aligned}
\]

and determines, through a simple positive root,

\[
 P_9(x)=\delta.
\]

The coefficients are not interpolated parameters. They proceed from the half-corona \(120\), the central APP block of determinant \(45\), the dodecaphasic transverse rank \(11\) and \(495=45\cdot11\). The nilpotent operator

\[
 Ne_7=-\frac83e_8,\qquad Ne_8=\frac2{11}e_9,
 \qquad Ne_9=0
\]

produces the tail exactly through

\[
 (I-N)^{-1}\!\left(-\frac{e_7}{120}\right)
 =-\frac{e_7}{120}+\frac{e_8}{45}+\frac{2e_9}{495}.
\]

The relation between the two routes is expressed without confusing a finite section with an infinite root. Let \(d_m\) be the reader of the certified common resolution. The object \(\alpha\) is then the fibre product

\[
 \mathfrak A_\alpha=
 \left\{(x_{12},x):
 d_m\!\left(\Gamma_{12}^{\alpha}(x_{12})\right)=d_m(x),\
 P_9(x)=\delta\right\}.
\]

The first coordinate fixes dodecaphasic closure; the second fixes unique extension; equality in the common codomain fixes their synchronization. This is why \(\alpha\) is singular within the correlated system. \(\pi\), \(\varphi\) and \(e\) are regions and coordinates that lifting makes coincide, separate and coincide again; \(\alpha\) records the seam preserving these coincidences when twelve phases close without stopping infinite extension. It does not arrive from outside afterwards: it is the holonomy coordinate through which joint generation becomes stable under return.

\section{Dodecaphasic fan: \(\alpha\), electron, Witt and gravitation}

The wheel of twelve positions does not produce an accidental succession of objects; it carries several irreducible readings of the same state. On

\[
 V_{12}=\mathbb R[A_5/C_5]
\]

the exact decomposition is

\[
 \boxed{V_{12}\simeq\mathbf1\oplus V_3\oplus V_{3'}\oplus V_5.}
\]

If \(P_1,P_3,P_{3'},P_5\) are the central projectors, every dodecaphasic linear form \(\ell\) decomposes canonically as

\[
 \ell(K)=\ell(P_1K)+\ell(P_3K)+\ell(P_{3'}K)+\ell(P_5K).
\]

The scalar reader determining \(\alpha\) acts on the full vector \(K\); it therefore integrates information from all summands without confusing them. The reader \(P_3\) leads, through the Witt incidence isometry, to the rank-three exceptional projector. The reader \(P_5\) leads through

\[
 P_5:V_{12}\to V_5,
 \qquad
 \Gamma_5:V_5\overset{\simeq}{\longrightarrow}
 \operatorname{Sym}_0^2(\mathbb R^3),
\]

to the five-component tensor carrier, and the projection \(\Lambda(n)\) determines its transverse traceless sector. These are not three juxtaposed theories: they are three resolutions of the same dodecaphasic carrier—a scalar seam resolution, an exceptional incidence resolution and a tensor gravitation resolution.

The electronic intersection makes this unity visible through an identity of supports. If

\[
 B_K=\{4,6,9,10\},\quad
 P_\pi=\{2,4,5,6,7,8,9,10,11\},
\]

\[
 H_e=\{1,3,4,6,9,10\},\quad
 H_\alpha^-=\{4,5,6,7,9,10\},
\]

then

\[
 \boxed{B_K=H_e\cap P_\pi=H_e\cap H_\alpha^-},
 \qquad
 \boxed{H_\alpha^-=B_K\sqcup\{5,7\}.}
\]

This equality answers the question of the “connection” of \(\alpha\). The electron does not join an external constant; its internal support intersects both the \(\pi\) microfibre and the negative-carry support of \(\alpha\). The additional pair \(\{5,7\}\) completes the negative hexad, and Witt polarity lifts the four pairs associated with \(B_K\) to their incidence blocks. Electronic vacancy is the local singularity; \(\alpha\) is the global seam preserving its orientation throughout the dodecaphase.

The second realization of \(\alpha\) is not scalar. From the same state one obtains

\[
 v_\alpha=4\rho_1+\rho_2+\cdots+\rho_{12},
 \qquad \|v_\alpha\|^2=54,
\]

and the neighbour

\[
 N_{\alpha,0}=
 \{x\in N(C_W):\langle x,v_\alpha\rangle\equiv0\pmod3\},
 \qquad
 N_\alpha=N_{\alpha,0}+\mathbb Z\frac{v_\alpha}{3}
 \cong\Lambda_{24}.
\]

The order-three action preserves this neighbour, and the corresponding orbifold determines \(V^\natural\) and the class \(3B\) of the Monster group. The scalar \(\alpha\) and vector \(v_\alpha\) are not the same object; they are two coordinates of the same dodecaphasic closure. This plurality of codomains is precisely what explains how \(\alpha\) links the Archimedean branch, electronic vacancy and exceptional symmetry without functioning as a numerical adjustment.

\section{From electronic vacancy to electromagnetic response}

The exact definition of vacancy already appears in the arithmetic of the support. Normalizing the electronic pair by nine produces

\[
 \frac1{9}(369,126)=(41,14)\longmapsto(5,5),
\]

whereas conservative transport satisfies

\[
 369+126=495=396+99,
 \qquad
 (396,99)/9=(44,11)\longmapsto(8,2).
\]

The two routes preserve the additive sheet, multiplicative residue and TRIT regime; the second differs by one unit in the multiplicative quotient. Vacancy is therefore the minimum loss of carry in one sheet while preserving the remaining state coordinates. It is not an added zero: it is a relative defect visible only because APP preserves sum and product together.

The tritic phase mode

\[
 m_{\rm ph}=(1,-1,0)^3
\]

is also the canonical singular electronic mode of APP:

\[
 C^{(3)}m_{\rm ph}=C^{(3)\,*}m_{\rm ph}
 =-\frac12m_{\rm ph}.
\]

From here, one obtains

\[
 s=\frac12,qquad
 1-s^2=\frac34=\frac{90}{120},
 \qquad
 \|[P_{\Sigma,3},P_{\Pi,3}]\|=\frac{\sqrt3}{4},
\]

and, on the principal block, normalization of the commutator satisfies

\[
 J_{\rm comm}^2=-I.
\]

Electromagnetic polarity is then determined by four layers of the same genealogy. APP supplies the sheets whose noncommutativity produces orientation; TRIT determines the elliptic generator \((J^2=-I)\); TPK preserves vacancy and its return; the channels \(90/120\) determine the common and differential responses. If

\[
\begin{aligned}
\mathcal E={}&
\log[(1-q_+^{90})(1-q_-^{90})]
-\log[(1-q_+^{120})(1-q_-^{120})],\\
\mathcal B={}&
\log\frac{1-q_-^{90}}{1-q_+^{90}}
-\log\frac{1-q_-^{120}}{1-q_+^{120}},
\end{aligned}
\]

then \(\mathcal E\) is even under sheet interchange and \(\mathcal B\) is odd. The internal responses

\[
 \widehat\mu=e^{\mathcal B+\mathcal E},
 \qquad
 \widehat\varepsilon=e^{-\mathcal B+\mathcal E}
\]

satisfy exactly

\[
 \sqrt{\frac{\widehat\mu}{\widehat\varepsilon}}
 =e^{\mathcal B},
 \qquad
 \frac1{\sqrt{\widehat\mu\widehat\varepsilon}}
 =e^{-\mathcal E}.
\]

Permeability expresses rotational retention; permittivity, electric opening; their quotient preserves orientation and their product preserves common propagation. The electromagnetic wave is not added to the Sturmian word as an analogy: it is the realization in two conjugate components of oriented vacancy and its elliptic transport. The electric microfibre can therefore be aperiodic while maintaining global coherence, stable polarity and zero transport entropy in the full state.

\section{Action, area, information and gravitation as predefined depths}

Once the correlated system has been determined, the determinantal branch fixes

\[
 \Sigma_H=\varphi\alpha^{16},
 \qquad
 d_H=\lfloor\log_{10}\Sigma_H\rfloor=-34.
\]

The regional reader determines the mantissas \(\eta_{\rm pre}=H_5\) and

\[
 \eta_{\rm ret}=H_5-\frac{90}{\pi}\mathscr C_\pi,
\]

so that

\[
 \hbar_a=\eta_a10^{-34}\mathcal U_S,
 \qquad h_a=2\pi\hbar_a.
\]

The scale \(10^{-34}\) does not proceed from a Planck measurement: it is the predefined depth at which the action reader finds its section. The phase \(\theta\) subsequently acts through

\[
 S(\theta)=\hbar_{\rm ret}\theta_{\rm rad}
 =h_{\rm ret}\frac{\theta_{\rm deg}}{360}.
\]

The factor \(2\pi\) does not correct a chosen unit: it relates action per radian to action per turn. The angular chart and the action line determine two coordinates of the same oriented motion.

The hexad--octad transition subsequently evaluates the channels and produces the Barbero--Immirzi parameter. On the boundary, if

\[
 N=N_{90}+N_{120},\qquad D=90N_{90}+120N_{120},
\]

then

\[
 N_{90}=4N-\frac{D}{30},
 \qquad
 N_{120}=\frac{D}{30}-3N.
\]

The joint reader \((N,D)\) preserves sectorial provenance in full:

\[
 H((N_{90},N_{120})\mid(N,D))=0.
\]

This is the exact formulation of zero Landauer cost in the area--information transition. The entropy of a reduced state may be positive; what vanishes is the conditional loss when the full reconstructive pair is retained. The area spectrum

\[
 \frac{\mathcal A_{\rm phys}}{\ell_P^2}
 =\gamma_{\rm HMT}a_0(N_{90}+N_{120})
\]

and the modular Hamiltonian

\[
 K_A=cI+\Delta_{\rm mod}(90N_{90}+120N_{120})
\]

are thus two complementary readers: the first counts total area; the second preserves how it was distributed between the two sectors.

Gravitation occupies another resolution of the same dodecaphasic closure. With

\[
 \omega_{108}=\frac{2\pi}{108t_0},
 \qquad R_{108}=\frac{54}{\pi}\ell_0,
\]

and

\[
 G_a(\vartheta)=\vartheta\frac{c_{\rm int}^5t_0^2}{\hbar_a},
\]

the coincidence condition for the three lengths

\[
 R_{108}=\bar\lambda_{C,a}=r_{g,a}
\]

uniquely selects

\[
 \vartheta_{108}=\left(\frac{54}{\pi}\right)^2
\]

and produces

\[
 \ell_P=\sqrt{\frac{G_a\hbar_a}{c_{\rm int}^3}}=R_{108}.
\]

The depth \(10^{-11}\) of the gravitational branch and the depth \(10^{-34}\) of the action branch are not observed scales introduced into the crystal. They are locations in the refinement where distinct readers of the same predefined section determine gravitation and action. The identity \(G_a\hbar_a=\vartheta_{108}c_{\rm int}^5t_0^2\) shows that the two action orientations change gravitation reciprocally and preserve the Planck length. Clock, action and gravitation are therefore synchronized coordinates of a single dimensional evolution.

\section{Governing law of the article: conservative evolution of the dimensional unit}

The preceding results allow the central thesis to be stated without turning it into a metaphor or a catalogue of constants.

\textbf{Theorem of conservative evolution and holographic predefinition.} Let \(\omega\in\Omega_{\rm seed}\) be an oriented APP seed and let \(\mathcal G_{\rm HMT}(\omega)\in X_\infty^{\rm enr}\) be its extension through TRIT and TPK. Then:

\begin{enumerate}
\item updating invertibly preserves sheet, residue, quotient, carry, orientation, boundary, route and memory;
\item the finite clock returns and holonomy advances, so the orbit is helical;
\item the irrational mechanical decoration has zero topological entropy, linear complexity and no translational period;
\item every compatible reader defines a predefined system in an inverse limit, so deepening reveals rather than recalibrates;
\item double projection determines, from the same section, the real line and the exceptional Witt--Golay--Leech--\(V^\natural\)--Monster branch;
\item the correlated system \((\pi,\varphi,e,\alpha)\) is a correlated output and \(\alpha\) is its dodecaphasic seam with the aperiodic clock;
\item electronic vacancy, Planck action, the Barbero--Immirzi parameter, the vacuum constitutive quantities and gravitation are realizations of typed readers on that same genealogy.
\end{enumerate}

The proof is the composition of the squares already written. Compatibility of truncations proves predefinition; inversion of \(\mathcal U_t\) proves conservation; the Sturmian cocycle proves zero-entropy aperiodicity; the identity of supports proves the electron--\(\pi\)--\(\alpha\) connection; the decomposition of \(V_{12}\) proves the coordinated separation of the exceptional and tensor readers; the action, area and information readers prove persistence of the unit under change of dimension.

The compact form of the theorem is

\[
\boxed{
\begin{CD}
\Omega_{\rm seed}@>{\mathrm{APP\to TRIT\to TPK}}>>
X_\infty^{\rm enr}\\
@. @V{(\mathcal P_{\rm arq},\mathcal P_{\rm exc},
\mathcal P_{\rm dim})}VV\\
@. \mathbb R\times\mathscr E_{\rm exc}\times\mathscr D_{\rm phys},
\end{CD}}
\]

where the three realizations preserve the same provenance despite having distinct codomains. The unit does not remain motionless: it is conserved while being refined, changing chart, acquiring memory and increasing its dimension. The continuum is not the prior support of this process but its terminal projection. This is why the helical aperiodic crystal constitutes the kernel of HMT--MD: in a single reversible dynamics, it predefines the genealogy of the numbers and constants subsequently appearing as mathematics and physics.
 
\chapter{Correlated generation of \(\pi\), \(e\) and \(\varphi\)}
\label{chap:83-09}

% Unidad U012. Biografía estructural completa de pi.

% La biografía de pi que seguía en esta residencia es exactamente la ya
% publicada en el capítulo 8, salvo el namespace de sus etiquetas. Se conserva
% en la fuente por su procedencia, pero no se imprime dos veces.
\iffalse

\section{Construction of the closure character and structural genealogical trajectories of \(\pi\)}
\label{p12-83-c9-s1:chap:biografia-pi-autonoma}

\subsection{Causal separation between seed, character, and evaluation}

The construction is divided into three strata:
\[
 \text{finite selection}
 \longrightarrow
 \text{construction of the closure character}
 \longrightarrow
 \text{Archimedean evaluation}.
\]
The decimal expansion does not intervene in the first stratum. The type chain
is
\begin{equation}
 104\,976
 \longrightarrow468\,\mathcal U_6
 \xrightarrow{\Psi}243\,\mathcal W_6
 \lhook\joinrel\longrightarrow\mathbb F_3^6,
 \qquad|\mathbb F_3^6|=729.
 \label{p12-83-c9-s1:eq:cadena-tipica-pi-autonoma}
\end{equation}
The cardinalities count, respectively, enriched initial conditions,
distinct emissions, visible ternary words, and elements
of the complete ambient fiber.

For \(U_6=(u_1,\ldots,u_6)\), the projection is
\[
 \Psi(U_6)=((-u_1)\bmod3,\ldots,(-u_6)\bmod3).
\]
The census of the \(43\) orbits of \(D_3\) proves that the closure
orbit is the only one whose six orientations possess at once:
\[
 \text{five lifts},\qquad
 \text{mass }1008,\qquad
 1008=432+4\cdot144.
\]
The gauge
\[
 \operatorname{diag}_c=6,
 \qquad
 \operatorname{card}=ES
\]
select a single representative:
\begin{equation}
 \boxed{w_6^\pi=010211.}
 \label{p12-83-c9-s1:eq:semilla-pi-autonoma}
\end{equation}
The selection usis catalogue, multiplicitiis, orbital action, and orientation;
it dois not consult to prefix of \(\pi\).

\subsection{Liftings \(L_0,L_1\) and 1st coinductive boundary}

On vectors row of \(\mathbb F_3^6\), let
\[
L_0=
\begin{pmatrix}
2&2&2&1&2&1\\
2&1&2&2&1&1\\
1&0&0&1&0&2\\
0&0&1&1&1&0\\
2&2&2&0&2&1\\
2&1&2&1&0&0
\end{pmatrix},
\quad
L_1=
\begin{pmatrix}
2&1&2&0&2&2\\
1&2&1&1&2&0\\
1&2&1&1&1&2\\
0&1&0&0&2&1\\
2&0&2&1&0&1\\
0&2&2&2&1&1
\end{pmatrix}.
\]
If \(b_1=w_6^\pi\), the recurrence of blocks is
\[
 b_{k+1}=b_kL_{\varepsilon_k},
 \qquad
 (\varepsilon_1,\varepsilon_2,\varepsilon_3,\varepsilon_4)=(0,0,1,1),
\]
and the accumulated word is formed by concatenation:
\[
 w_{6(k+1)}=w_{6k}\mathbin{\|}b_{k+1}.
\]
A word of length \(6k\) is not multiplied by a matrix \(6\times6\).

For the closing channel:
\[
\begin{array}{c|c}
\text{depth}&\text{new block}\\ \hline
6&010211\\
12&012222\\
18&010211\\
24&002111\\
30&110221
\end{array}
\]
and
\begin{equation}
 \boxed{
 w_{30}^{\pi}
 =010211|012222|010211|002111|110221.}
 \label{p12-83-c9-s1:eq:w30-pi-autonomo}
\end{equation}
The following boundary is
\[
 u_5^\pi=222220.
\]
Along with the other channels,
\[
 R_{36}=
 \begin{pmatrix}
 222220\\021222\\102011
 \end{pmatrix}.
\]
The symbol \(R_{36}\) records the sixth block of three histories and opens a
boundary; it is not a terminal.

The certified branch of twenty blocks begins
\begin{align*}
\pi:\;&010211|012222|010211|002111|110221|222220|111201|\\
&212121|200121|100100|101222|022212|012012|111210|\\
&121011|200220|120210|000101|022010|020111.
\end{align*}
For this reason, ni \(w_{30}\) ni \(R_{36}\) terminate the genealogy.

\subsection{Oriented decomposition of the microfiber of \texorpdfstring{\(\pi\)}{pi}: central region, \(4\) lateral regions, and multiplicity \texorpdfstring{\(432+4\cdot144=1008\)}{432+4x144=1008}}

The microfiber
\[
 \mathscr R_\pi^{\rm micro}
 :=\{U\in\mathcal U_6:\Psi(U)=010211\}
\]
contains exactly five regions:
\[
\begin{array}{c|c|c|c}
U_6&E_{\rm sig}&\operatorname{mult}&\text{role}\\ \hline
501|614|498|169|272|272&555555&432&\perp\\
810|923|870|169|272|272&339555&144&++\\
810|983|810|169|272|272&393555&144&++\\
870|923|810|169|272|272&933555&144&++\\
870|923|810|418|674|521&933717&144&--
\end{array}
\]
Consequently,
\begin{equation}
 \boxed{
 5=1_\perp+3_{++}+1_{--},
 \qquad
 1008=432+4\cdot144.}
 \label{p12-83-c9-s1:eq:descomposiciones-cinco-regiones-pi}
\end{equation}
The first identity classifies orientation functions. The second counts
provenance multiplicities. They do not count five indistinguishable copies.

The transversal region is
\[
 501|614|498|169|272|272,
 \qquad E_{\rm sig}=555555,
 \qquad\operatorname{mult}=432.
\]
The mutated return is
\[
 870|923|810|418|674|521,
 \qquad E_{\rm sig}=933717,
 \qquad\operatorname{mult}=144.
\]
Interchanging them would preserve the census \(3/1/1\) but falsify the complete
regional structure.

\subsection{Descent of a complete nonadic block to the phase transfer multigraph}

For \(U_6=(U_1,\ldots,U_6)\), let
\[
 A(U_6)=(U_1,U_2,U_3),
 \qquad
 B(U_6)=(U_4,U_5,U_6).
\]
Each emission defines an edge between \(A(U_6)\) and \(B(U_6)\). The raw
multigraph has \(96\) vertices, distributed into components of sizes
\[
 28,28,28,4,4,4.
\]
The internal phase is quotiented by
\[
 (a,b,c)\sim(b,c,a)\sim(c,a,b).
\]
The representative is the lexicographically minimal rotation. The quotient
graph has \(32\) vertices and two components of sizes \(28\) and \(4\).

The anchors are
\[
\begin{aligned}
U_{\rm APP}&=903|850|850|252|109|252,\\
U_e&=850|963|890|149|252|212,\\
U_\varphi&=830|943|830|109|252|252.
\end{aligned}
\]
We number the regions.
\[
\begin{aligned}
U_\pi^{(1)}&=810|923|870|169|272|272,\\
U_\pi^{(2)}&=810|983|810|169|272|272,\\
U_\pi^{(3)}&=870|923|810|169|272|272,\\
U_\pi^{(4)}&=870|923|810|418|674|521,\\
U_\pi^{(5)}&=501|614|498|169|272|272.
\end{aligned}
\]

\subsection{The \(5\) oriented closure cycles}

\subsubsection{Closure of \texorpdfstring{\(U_\pi^{(1)}\)}{U-pi 1}}
\begin{align*}
&252|109|252|903|850|850\to
903|850|850|252|109|252\to\\
&252|109|252|923|870|810\to
810|923|870|169|272|272\to\\
&272|169|272|963|890|850\to
850|963|890|149|252|212\to\\
&212|149|252|943|830|830\to
830|943|830|109|252|252\to
252|109|252|903|850|850.
\end{align*}

\subsubsection{Closure of \texorpdfstring{\(U_\pi^{(2)}\)}{U-pi 2}}
\begin{align*}
&903|850|850|252|109|252\to
272|169|272|903|850|850\to\\
&272|169|272|983|810|810\to
810|983|810|169|272|272\to\\
&272|169|272|963|890|850\to
850|963|890|149|252|212\to\\
&212|149|252|943|830|830\to
830|943|830|109|252|252\to
903|850|850|252|109|252.
\end{align*}

\subsubsection{Closure of \texorpdfstring{\(U_\pi^{(3)}\)}{U-pi 3}}
\begin{align*}
&903|850|850|252|109|252\to
292|189|232|903|850|850\to\\
&292|189|232|923|810|870\to
870|923|810|169|272|272\to\\
&272|169|272|963|890|850\to
850|963|890|149|252|212\to\\
&212|149|252|943|830|830\to
830|943|830|109|252|252\to
903|850|850|252|109|252.
\end{align*}

\subsubsection{Closure of \texorpdfstring{\(U_\pi^{(4)}\)}{U-pi 4}}
\begin{align*}
&903|850|850|252|109|252\to
272|169|272|903|850|850\to\\
&272|169|272|963|890|850\to
850|963|890|149|252|212\to\\
&212|149|252|923|810|870\to
870|923|810|418|674|521\to\\
&943|830|830|521|418|674\to
830|943|830|109|252|252\to
903|850|850|252|109|252.
\end{align*}

\subsubsection{Closure of \texorpdfstring{\(U_\pi^{(5)}\)}{U-pi 5}}
\begin{align*}
&252|109|252|903|850|850\to
903|850|850|252|109|252\to\\
&614|498|501|252|109|252\to
501|614|498|169|272|272\to\\
&272|169|272|963|890|850\to
850|963|890|149|252|212\to\\
&212|149|252|943|830|830\to
830|943|830|109|252|252\to
252|109|252|903|850|850.
\end{align*}

\begin{theorem}[Minimality of the Five Closures]
\label{p12-83-c9-s1:thm:minimalidad-cierres-pi-autonoma}
For each \(i\in\{1,\ldots,5\}\), the minimum length of a closed walk
that simultaneously traverses
\[
 U_\pi^{(i)},\quad U_e,\quad U_\varphi,\quad U_{\rm APP}
\]
in the quotient multigraph is eight.
\end{theorem}

\begin{proof}
Given \(i\), let
\[
 E_*=\{U_\pi^{(i)},U_e,U_\varphi,U_{\rm APP}\}.
\]
The finite space is considered.
\[
 \mathscr S_i=\{(v,M):v\in V(G),\ M\subseteq E_*\}.
\]
The first component records the vertex and the second the target edges
already traversed. If an edge \(e\) takes \(v\) to \(v'\), the transition is
\[
 (v,M)\longmapsto
 \bigl(v',M\cup(E_*\cap\{e\})\bigr).
\]
Breadth-first search enumerates lengths in increasing order. The first
return to \((v,E_*)\) occurs at length eight for all five regions.
The printed walks attain the bound; the optimality of the search excludes
a shorter length.
\end{proof}

The closure function is encoded by
\[
 \mathfrak C_\pi=
 \bigl(\mathscr R_\pi^{\rm micro},
 w_\pi,\rho_\pi,m_\pi,\ell_{\rm cl}\bigr),
\]
\[
 \rho_\pi=(\perp,++ ,++ ,++ ,--),
 \qquad
 m_\pi=(432,144,144,144,144),
 \qquad
 \ell_{\rm cl}=8.
\]

\subsection{Incidence operator and oriented half-turn}

The incidence shadow is
\[
 A_W=
 \begin{pmatrix}
 0&1&1&1&1&1\\
 1&0&1&2&2&1\\
 1&1&0&1&2&2\\
 1&2&1&0&1&2\\
 1&2&2&1&0&1\\
 1&1&2&2&1&0
 \end{pmatrix}.
\]
Multiplication in \(\mathbb F_3\) gives
\begin{equation}
 \boxed{A_W^2=2I_6=-I_6.}
 \label{p12-83-c9-s1:eq:aw-cuadrado-menos-identidad-pi}
\end{equation}
Products between distinct rows and columns vanish modulo three,
whereas each diagonal product equals two. One application effects a ternary
quarter-turn; two applications effect the oriented inversion.

\subsection{The slope 1/5 forced by the \(5\) regions}

Let
\[
 \lambda=(\lambda_1,\ldots,\lambda_5)^{\mathsf T}
 \in\mathbb Q_{\geq0}^5,
 \qquad
 \mathbf1^{\mathsf T}\lambda=1.
\]
Forgetting the regional label applies the symmetrizer
\[
 P_5=\frac15\mathbf1\mathbf1^{\mathsf T}.
\]
For every normalized vector,
\[
 P_5\lambda=\frac15\mathbf1.
\]
The invariance \(P_5\lambda=\lambda\) therefore has a unique solution:
\begin{equation}
 \boxed{\lambda=\frac15\mathbf1,\qquad t_1=\frac15.}
 \label{p12-83-c9-s1:eq:pendiente-primitiva-pi}
\end{equation}
The integer five comes from the regional microfiber; it is not recognized in a
formula for \(\pi\).

\subsection{Rational composition and compensator \texorpdfstring{\(239\)}{239}}

The law of slope composition is
\[
 t\star s:=\frac{t+s}{1-ts}.
\]
The first duplication gives
\[
 t_2=t_1\star t_1
 =\frac{2/5}{1-1/25}
 =\frac5{12}.
\]
The second gives
\[
 t_4=t_2\star t_2
 =\frac{10/12}{1-25/144}
 =\frac{120}{119}.
\]
The return to the unit diagonal requires \(q>0\) such that
\[
 \frac{120/119-1/q}{1+(120/119)/q}=1.
\]
Multiplying by \(119q\) produces
\[
 120q-119=119q+120,
\]
and forces
\begin{equation}
 \boxed{q=239.}
 \label{p12-83-c9-s1:eq:compensador-239-pi}
\end{equation}

\subsection{Gaussian identity and 1st half-turn}

We define
\begin{equation}
 \Theta_\pi
 :=16\arctan\frac15-4\arctan\frac1{239}.
 \label{p12-83-c9-s1:eq:theta-pi-autonoma}
\end{equation}
The successive squares are
\[
\begin{array}{c|r@{}l}
\text{power}&\multicolumn{2}{c}{\text{Gaussian integer}}\\ \hline
(5+i)^2&24&+10i\\
(5+i)^4&476&+480i\\
(5+i)^8&-3824&+456960i\\
(5+i)^{16}&-208797818624&-3494830080i\\
(239-i)^2&57120&-478i\\
(239-i)^4&3262465916&-54606720i
\end{array}
\]
and the final multiplication gives
\begin{equation}
 \boxed{
 (5+i)^{16}(239-i)^4
 =-681386607803576157184.}
 \label{p12-83-c9-s1:eq:identidad-gaussiana-pi-autonoma}
\end{equation}
The imaginary part is zero and the real part is negative. Therefore,
\[
 \Theta_\pi\equiv\pi\pmod{2\pi}.
\]
Para \(0<x<1\),
\[
 x-\frac{x^3}{3}<\arctan x<x.
\]
From here
\[
 \Theta_\pi>
 16\left(\frac15-\frac1{3\cdot5^3}\right)-\frac4{239}>3
\]
and
\[
 \Theta_\pi<\frac{16}{5}<4.
\]
The only negative orientation in that interval is the first positive half-turn:
\begin{equation}
 \boxed{\Theta_\pi=\pi.}
 \label{p12-83-c9-s1:eq:reconocimiento-pi-exacto}
\end{equation}

The order of subsequent Archimedean recognition is
\[
 \text{five regions}
 \to\frac15
 \to\frac5{12}
 \to\frac{120}{119}
 \to239
 \to\Theta_\pi
 \to\pi.
\]

\subsection{Arbitrary precision rational horocycles}

For \(0<x<1\), let
\[
 S_n(x)=\sum_{r=0}^{n}(-1)^r\frac{x^{2r+1}}{2r+1}.
\]
Leibniz's criterion gives
\[
 S_{2m+1}(x)<\arctan x<S_{2m}(x).
\]
We define
\[
 a_m^-=S_{2m+1}(1/5),\quad a_m^+=S_{2m}(1/5),
\]
\[
 b_m^-=S_{2m+1}(1/239),\quad b_m^+=S_{2m}(1/239),
\]
and
\begin{equation}
 \ell_{\pi,m}=16a_m^- -4b_m^+,
 \qquad
 h_{\pi,m}=16a_m^+ -4b_m^-.
 \label{p12-83-c9-s1:eq:horquilla-pi-autonoma}
\end{equation}
Then
\[
 \ell_{\pi,m}<\pi<h_{\pi,m}
\]
and
\begin{equation}
 h_{\pi,m}-\ell_{\pi,m}
 =
 \frac{16}{(4m+3)5^{4m+3}}
 +\frac{4}{(4m+3)239^{4m+3}}
 \longrightarrow0.
 \label{p12-83-c9-s1:eq:anchura-horquilla-pi}
\end{equation}
The lower bounds grow and the upper bounds decrease.

Since \(3<\pi<4\), the ternary reader acts on
\[
 r_\pi:=\pi-3\in(0,1),
\]
with bracket
\[
 [\ell_{\pi,m}-3,h_{\pi,m}-3].
\]

\subsection{Subsequent Archimedean certification of the already generated fiber}

Suppose the internal relation \(\operatorname{Ext}\) and coinductive
survival have already produced, without using Machin or a tabulated expansion,
the compatible extension \(N_{K+1}\) of the prefix \(N_K\), of length
\(L_K=6K\). The \(729\) subcylinders of its ambient fiber are
\[
 I_{K,u}=
 \left[
 \frac{729N_K+u}{3^{L_K+6}},
 \frac{729N_K+u+1}{3^{L_K+6}}
 \right),
 \qquad0\leq u<729.
\]
The order \(m=m(K)\) is then increased until the Leibniz bracket
locates the already produced subcylinder. Let
\[
 \ell_{\pi,K}:=\ell_{\pi,m(K)}-3,
 \qquad
 h_{\pi,K}:=h_{\pi,m(K)}-3.
\]
The compatible integer endpoints are
\[
 j_{\min}=
 \max\!\left(
 729N_K,\left\lfloor
 \ell_{\pi,K}3^{L_K+6}\right\rfloor
 \right),
\]
\[
 j_{\max}=
 \min\!\left(
 729N_K+728,\left\lceil
 h_{\pi,K}3^{L_K+6}\right\rceil-1
 \right).
\]
The condition of unitary localization is
\begin{equation}
 \boxed{j_{\min}=j_{\max},\qquad N_{K+1}=j_{\min}.}
 \label{p12-83-c9-s1:eq:seleccion-unitaria-subcilindro-pi}
\end{equation}
The first equality certifies unit localization; the second verifies
that this localization agrees with the internally generated extension. Neither
defines the TPK transition or selects among the \(729\) states.
Existence and uniqueness come from \(\operatorname{Ext}\) and
survival; Machin and the Leibniz brackets are subsequent
recognitions.

\subsection{Coinductive prolongation of the section of \(\pi\) through the inverse system of surviving histories}

Let \(\mathcal H_m^\pi\) be the set of enriched histories at depth
\(m\), and let
\[
 p_{n,m}:\mathcal H_n^\pi\longrightarrow\mathcal H_m^\pi
\]
the truncation. We define
\[
 \operatorname{Surv}_m^{(N),\pi}
 :=p_{N,m}(\mathcal H_N^\pi),
\]
\[
 \operatorname{Surv}_m^\pi
 :=\bigcap_{N\geq m}\operatorname{Surv}_m^{(N),\pi}.
\]
When the relation \(\operatorname{Ext}\) is nonempty, single-valued and
compatible at every level, the sets are singletons and satisfy
\[
 p_{m+1,m}(\operatorname{Surv}_{m+1}^{\pi})
 =\operatorname{Surv}_{m}^{\pi}.
\]
In that domain of extendability, there exists, therefore, a unique section.
\[
 x_\pi\in
 X_\pi:=\varprojlim_m\operatorname{Surv}_m^\pi.
\]
The cylinders are nested and their diameters are \(3^{-6K}\), which tends to
zero. Their intersection contains a unique point, identified by the angular
character with \(r_\pi=\pi-3\).

For every decimal depth \(M\), one may choose a level \(K(M)\) whose
cylinder lies within a single decimal cell of diameter \(10^{-M}\).
The deep prefixes truncate exactly to the preceding ones.

\subsection{Characterization of the coinductive channel of \(\pi\): selection of \(w_6^\pi\), \(5\) regions, oriented inversion, and uniqueness of the limit section}

\begin{theorem}[Characterization of the closure channel]
\label{p12-83-c9-s1:thm:caracterizacion-final-pi-autonoma}
Starting from the catalog, the action \(D_3\), the oriented gauge, the
operators \(L_0,L_1,A_W\), the slope law and the rational filtration:
\begin{enumerate}
 \item \(w_6^\pi=010211\) is selected;
 \item five regions are constructed with the identities of
       \eqref{p12-83-c9-s1:eq:descomposiciones-cinco-regiones-pi};
 \item each region belongs to a minimum closure of length eight;
 \item \(A_W^2=-I_6\) fixes the oriented inversion;
 \item regional symmetrization forces \(1/5\);
 \item the composition forces \(120/119\);
 \item the unit return forces \(239\);
 \item the Gaussian identity recognizes the first half-turn;
 \item the brackets and the fiber of \(729\) elements produce a section
       compatible at every finite depth.
\end{enumerate}
Its evaluation is
\[
 \boxed{\pi=16\arctan\frac15-4\arctan\frac1{239}.}
\]
\end{theorem}

The complete chain is
\[
 \mathrm{APP}\to\mathrm{TRIT}\to\mathrm{TPK}
 \to104\,976\to468\to243\to w_6^\pi
 \to w_{30}^\pi\to R_{36}\to\Gamma_9
 \to X_\pi\to\pi.
\]

\begin{criterion}[Channel Closure Falsification]
The derivation is refuted if:
\begin{enumerate}
 \item the closure orbit ceases to be unique;
 \item one of the five regions disappears or its complete identity changes;
 \item one of the five closures has length less than eight;
 \item \(A_W^2\ne2I_6\);
 \item the symmetrizer admits another normalized distribution;
 \item the compositions do not produce \(5/12\), \(120/119\), and \(239\);
 \item the Gaussian identity fails;
 \item a bracket ceases to contain the angular character;
 \item a level does not traverse the \(729\) elements of the ambient fiber or does not
       locate exactly one extension;
 \item a target expansion intervenes before recognition.
\end{enumerate}
\end{criterion}
\fi

The complete construction of the closure character, the five regions, and the
structural biography of \(\pi\) was established in
\cref{chap:83-08}. Starting from that same already constructed character, the present
chapter now prolongs the refinement toward \(e\) and \(\varphi\), without
restarting the genealogy or repeating its first stage.


% Unidad U013. Biografía estructural completa del carácter exponencial.

\section{Propagation, memory return and construction of the exponential character}
\label{p12-83-c9-s2:chap:biografia-completa-e}

\subsection{Finite orbital selector: existence and uniqueness of compatible emission}

Let \(\mathcal W=\mathbb F_3^6\). For an emission
\(U=(u_1,\ldots,u_6)\), the visible word is
\[
 \Psi(U)=((-u_1)\bmod3,\ldots,(-u_6)\bmod3).
\]
On \(\mathcal W\), the dihedral action is given by
\[
\begin{aligned}
 r(u_1,u_2,u_3,u_4,u_5,u_6)
 &=(u_3,u_1,u_2,u_5,u_6,u_4),\\
 s(u_1,u_2,u_3,u_4,u_5,u_6)
 &=(u_4,u_5,u_6,u_1,u_2,u_3).
\end{aligned}
\]
For \(w\in\mathcal W\), define
\[
 n(w)=\#\Psi^{-1}(w),
 \qquad
 M(w)=\sum_{U\in\Psi^{-1}(w)}\operatorname{mult}(U),
\]
and
\[
 \nu(w)=
 \bigl(
 \#\{j:w_j=0\},
 \#\{j:w_j=1\},
 \#\{j:w_j=2\}
 \bigr).
\]

\begin{theorem}[Structural selection of the exponential channel]
\label{p12-83-c9-s2:thm:seleccion-estructural-e}
Among the \(43\) orbits of the visible catalog, there is a unique orbit whose
fibers simultaneously satisfy:
\[
 n(w)=1,\qquad M(w)=144,\qquad
 \mathfrak D(w)=\{0,3,6\},\qquad
 \nu(w)=(2,3,1).
\]
The gauge
\[
 \operatorname{diag}_c=6,
 \qquad \operatorname{card}=ES
\]
select in it
\[
 \boxed{w_6^e=201101.}
\]
\end{theorem}

\begin{proof}
The four conditions are evaluated over the \(243\) words in the catalog, not over a sample. A single orbit satisfies the joint predicate. In that orbit, exactly one emission has both diagonal \(6\) and orientation \(ES\). The enumeration uses only fields from the APP--TRIT catalog and does not consult an expansion of \(e\).
\end{proof}

\subsection{Visible lift and chain of blocks}

The recurrence \(L_0,L_0,L_1,L_1\) produces
\[
 w_{30}^{e}
 =201101|121221|102011|012222|102011,
\]
and the joint boundary contributes
\[
 u_5^e=021222.
\]
The record of thirty-six trits remains
\[
 w_{36}^{e}
 =201101|121221|102011|012222|102011|021222.
\]

\subsection{Decimal signature and positional square}

The initial reading of twelve digits is
\[
 \mathbf d_e^{(12)}
 =718281828459
 =7\,\underbrace{1828\,1828}_{\text{local square}}\,459.
\]
The central equality is the concatenation
\[
 1828\mathbin{\|}1828,
\]
not an arithmetic product nor the supposed beginning of a period.

\begin{definition}[Positional square]
A word \(d_1\cdots d_{12}\) contains a square of length \(\ell\)
at position \(p\) when
\[
 d_p\cdots d_{p+\ell-1}
 =d_{p+\ell}\cdots d_{p+2\ell-1}.
\]
The position and the lateral contexts belong to the signature.
\end{definition}

\begin{table}[htbp]
\centering
\caption{Comprehensive census of squares in twelve-digit readings.}
\label{p12-83-c9-s2:tab:censo-cuadrados-e}
\begin{tabular}{c*{6}{r}}
\toprule
Length \(\ell\)&1&2&3&4&5&6\\
\midrule
Occurrences&240&21&3&2&0&0\\
Words that contain them&153&19&3&2&0&0\\
\bottomrule
\end{tabular}
\end{table}

The other word with a square of length four is
\[
 575157518472
 =\underbrace{5751\,5751}_{\text{initial position }1}\,8472.
\]

\begin{proposition}[Positional uniqueness]
\label{p12-83-c9-s2:prop:cuadrado-posicional-e}
The reading of channel \(e\) is the only one that has a square of length
four after a prefix of length one and with contexts \(7\) and \(459\).
\end{proposition}

\begin{proof}
The census leaves only two words with squares of length four. One starts
it at the first position; the other after the prefix \(7\). The
position and contexts separate the two possibilities.
\end{proof}

\subsection{Visible return and memory rupture}

In the elevation
\[
 201101|121221|102011|012222|102011
\]
the third and the fifth blocks coincide:
\[
 B_3=B_5=102011.
\]
Their prestates modulo \(27\) are, however,
\[
 p_3=(25,11,7,17,8,16),
 \qquad
 p_5=(7,8,4,23,26,7),
\]
and the subsequent quotients:
\[
 q_3=(6,13,9,2,13,1),
 \qquad
 q_5=(15,0,3,19,23,25).
\]
\(p_3\ne p_5\) and \(q_3\ne q_5\) hold.

\begin{figure}[htbp]
\centering
\begin{tikzpicture}[
 node distance=11mm and 7mm,
 every node/.style={font=\small},
 block/.style={draw,rounded corners,minimum width=17mm,minimum height=8mm},
 state/.style={draw, dashed, rounded corners, align=center, inner sep=3pt},
 >=stealth
]
\node[block] (b1) {\(201101\)};
\node[block,right=of b1] (b2) {\(121221\)};
\node[block,right=of b2] (b3) {\(102011\)};
\node[block,right=of b3] (b4) {\(012222\)};
\node[block,right=of b4] (b5) {\(102011\)};
\draw[->] (b1)--(b2);
\draw[->] (b2)--(b3);
\draw[->] (b3)--(b4);
\draw[->] (b4)--(b5);
\node[state,below=of b3] (p3)
 {\(p_3=(25,11,7,17,8,16)\)\\
  \(q_3=(6,13,9,2,13,1)\)};
\node[state,below=of b5] (p5)
 {\(p_5=(7,8,4,23,26,7)\)\\
  \(q_5=(15,0,3,19,23,25)\)};
\draw[->] (p3)--(b3);
\draw[->] (p5)--(b5);
\draw[<->,thick] (b3) to[bend left=28]
 node[above]{same word visible} (b5);
\end{tikzpicture}
\caption{Visible return of \(102011\) with renewal of the prestate and the
quotient.}
\label{p12-83-c9-s2:fig:ruptura-memoria-e}
\end{figure}

\begin{theorem}[Visible-phase return without periodicity of the state]
\label{p12-83-c9-s2:thm:retorno-visible-no-periodico-e}
\(B_3=B_5\) is a return of the visible projection, not a periodic orbit of the integral holonomic state.
\end{theorem}

\begin{proof}
The return of the complete state would require equality of all its coordinates.
The prestate and quotient inequalities exclude it. Nor does
\(1828|1828\) begin a decimal period: after the second block,
\(4\) appears, whereas a periodic continuation would require \(1\).
\end{proof}

\subsection{Subsequent factorial certification of the coinductive exponential section}

Let \(\mathcal A\) be a unital commutative algebra over \(\mathbb Q\), and
let \(P_t\in\mathcal A[[t]]\) be the formal propagator. Temporal
concatenation imposes
\[
 P_{s+t}=P_sP_t,
 \qquad
 P_0=1.
\]
We write
\[
 P_t=\sum_{n=0}^{\infty}a_nt^n.
\]
The unit fixes \(a_0=1\), and normalized elementary transport fixes
\(a_1=1\).

\begin{theorem}[Factorial recurrence]
\label{p12-83-c9-s2:thm:recurrencia-factorial-e}
\[
 (n+1)a_{n+1}=a_n
 \qquad(n\geq0),
\]
and, therefore,
\[
 a_n=\frac1{n!},
 \qquad
 P_t=\sum_{n=0}^{\infty}\frac{t^n}{n!}.
\]
En \(t=1\),
\[
 \boxed{P_1=e.}
\]
\end{theorem}

\begin{proof}
The coefficient of \(s^nt\) in \(P_{s+t}\) is
\[
 \binom{n+1}{n}a_{n+1}=(n+1)a_{n+1}.
\]
In \(P_sP_t\) it is \(a_na_1=a_n\). Equality of series produces the
recurrence. Induction from \(a_0=1\) gives \(a_n=1/n!\).
\end{proof}

\begin{corollary}[Differential equation]
\[
 \frac{d}{dt}P_t=P_t,
 \qquad P_0=1.
\]
\end{corollary}

\begin{proof}
\[
 P_t'=\sum_{n\geq0}(n+1)a_{n+1}t^n
 =\sum_{n\geq0}a_nt^n=P_t.
\]
\end{proof}

\subsection{Cofinal Forking of Rational Tangles and Convergence to the Published Value}

Let
\[
 S_N=\sum_{n=0}^{N}\frac1{n!}.
\]

\begin{theorem}[Explicit factorial bound]
\label{p12-83-c9-s2:thm:cota-factorial-e}
For \(N\geq1\),
\begin{equation}
 \boxed{
 S_N<e<
 S_N+\frac{N+2}{(N+1)(N+1)!}.}
 \label{p12-83-c9-s2:eq:horquilla-factorial-e}
\end{equation}
\end{theorem}

\begin{proof}
The rest is
\[
 R_N=\sum_{k=0}^{\infty}\frac1{(N+1+k)!}>0.
\]
Para \(k\geq0\),
\[
 (N+1+k)!
 \geq(N+1)!(N+2)^k.
\]
Therefore,
\[
 R_N
 \leq\frac1{(N+1)!}
 \sum_{k=0}^{\infty}\frac1{(N+2)^k}
 =\frac{N+2}{(N+1)(N+1)!}.
\]
The inequality is strict from \(k=2\) onward.
\end{proof}

We define
\[
 J_N^{(e)}=
 \left[
 S_N,\,
 S_N+\frac{N+2}{(N+1)(N+1)!}
 \right]
\]
and, for the fractional part,
\[
 \widetilde J_N^{(e)}=J_N^{(e)}-2.
\]
The width converges to zero. For each \(M\), there exists \(N\) such that
\[
 \operatorname{diam}J_N^{(e)}<10^{-M}.
\]
If the endpoints belong to the same decimal cell, the first \(M\)
digits are certified by rational arithmetic.

\subsection{Coinductive prolongation of the exponential section through ternary subcylinders of the Topological Phase Kernel}

At level \(K\), TPK and the relation \(\operatorname{Ext}\) have already generated
a unique compatible subcylinder. The smallest \(N=N_e(K)\) for which
the factorial bracket is contained in that subcylinder is taken. \(N_e(K)\) is an
order of analytic certification, not a generating index or a rule for
selecting the fiber.

The TPK cylinders are compatible under truncation independently of the
semigroup. The semigroup theorem subsequently identifies their Archimedean
evaluation with \(e-2\); the integer part recovers \(e\). If the orbit,
\(\operatorname{Ext}\) or truncation fails, generation fails; if the
recurrence or the factorial bound fails, this analytic certificate fails, not the
APP--TRIT--TPK genealogy.

\subsection{Adversarial certificate for the generation of \(e\): orbital uniqueness, factorial recurrence, rational brackets and cylindrical prolongation}

\begin{criterion}[Exponential Channel Falsification]
The construction is refuted if:
\begin{enumerate}
 \item another orbit satisfies the complete finite predicate;
 \item the signature \(7[1828\,1828]459\) is not positionally unique;
 \item the prestates or quotients of the two blocks \(102011\) coincide;
 \item the semigroup does not force \((n+1)a_{n+1}=a_n\);
 \item the bound \eqref{p12-83-c9-s2:eq:horquilla-factorial-e} fails;
 \item a level does not locate a compatible extension uniquely;
 \item a tabulated expansion intervenes before recognition.
\end{enumerate}
\end{criterion}


% Unidad U014. Biografía estructural completa del carácter de autoescala.

\section{Modal incidence, Perron ray, and construction of the self-similar character}
\label{p12-83-c9-s3:chap:biografia-completa-phi}

\subsection{Finite Selection of the Self-Scaling Seed}

Let \(\mathcal U_6\subset\{0,\ldots,999\}^6\) be the catalog of the
\(468\) distinct emissions constructed in
\cref{chap:semillas-emisor-54}, and let
\[
 \Psi:\mathcal U_6\longrightarrow\mathbb F_3^6,
 \qquad
 \Psi(U)=((-U_j)\bmod3)_{j=1}^{6}.
\]
For \(w\in\Psi(\mathcal U_6)\), we recall the invariants
\[
 n(w)=\#\Psi^{-1}(w),
 \qquad
 M(w)=\sum_{U\in\Psi^{-1}(w)}\operatorname{mult}(U),
\]
and the occupation vector
\[
 \nu(w)=
 \bigl(\#\{j:w_j=0\},\#\{j:w_j=1\},\#\{j:w_j=2\}\bigr).
\]

\begin{theorem}[Orbital selection of the self-scale]
\label{p12-83-c9-s3:thm:seleccion-orbital-phi}
The dihedral action on \(\Psi(\mathcal U_6)\) contains a unique orbit
that satisfiis the self-scale predicate
\[
 n(w)=1,
 \qquad
 M(w)=144,
 \qquad
 \nu(w)=(2,2,2).
\]
The internal caliber
\[
 \operatorname{diag}_c=6,
 \qquad
 \operatorname{card}=ES
\]
select in that orbit the word
\[
 \boxed{w_6^{\varphi}=121200.}
\]
Its only emission is
\[
 \boxed{U_{\varphi}=(830,943,830,109,252,252),}
\]
of multiplicity \(144\), multiplicative signature \(555933\), additive
orientation \(ES\), multiplicative support of cardinality six and residual type
\(A\).
\end{theorem}

\begin{proof}
The direct projection gives
\[
 U_{\varphi}\bmod3=(2,1,2,1,0,0),
 \qquad
 -U_{\varphi}\bmod3=(1,2,1,2,0,0).
\]
The enumeration is carried out over the \(243\) visible words and their
\(43\) orbits, not over a selection of examples. The joint predicate
\(n=1\), \(M=144\), \(\nu=(2,2,2)\) leaves one orbit. Within it, the
declared diagonal and orientation leave one emission. At this stage neither
the equation \(x^2-x-1=0\), nor a decimal table, nor the
conventional name of the constant has been introduced.
\end{proof}

\subsection{Blockwise lift and 1st coinductive boundary}

Let \(b_1=w_6^{\varphi}\), and let
\(\varepsilon=(0,0,1,1)\). With the matrices \(L_0,L_1\) constructed in
\cref{hmtctx:p12:chap:orbitas-elevacion-r36}, we define
\[
 b_{k+1}=b_kL_{\varepsilon_k}\pmod3
 \qquad(1\leq k\leq4),
\]
and concatenate
\[
 w_{6(k+1)}^{\varphi}=w_{6k}^{\varphi}\mathbin{\|}b_{k+1}.
\]
Matrix calculus produces, block by block,
\[
\begin{aligned}
 b_1&=121200,\\
 b_2&=112202,\\
 b_3&=121020,\\
 b_4&=010210,\\
 b_5&=010200,
\end{aligned}
\]
and, therefore,
\begin{equation}
 \boxed{
 w_{30}^{\varphi}
 =121200|112202|121020|010210|010200.}
 \label{p12-83-c9-s3:eq:w30-phi}
\end{equation}
The dual neutrality of the boundary adds
\[
 u_5^{\varphi}=102011,
\]
so that
\begin{equation}
 \boxed{
 w_{36}^{\varphi}
 =121200|112202|121020|010210|010200|102011.}
 \label{p12-83-c9-s3:eq:w36-phi}
\end{equation}

\begin{remark}[Two distinct uses of a visible word]
The word \(102011\) is also visible in the history of \(e\). Equality of one projection does not identify the integral holonomic states: channel, prefix, quotient, carry, phase and memory belong to the object’s type. This distinction is mandatory before constructing any real character.
\end{remark}

\subsection{From the tritic calendar to the incidence multigraph}

\subsubsection{Primitive modal incidence \(F_{\rm av}\) and memory dynamics of order \(1\)}

The primitive triticale calendar is
\[
 (+1,-1,0).
\]
We fix two visible modis, \(\Sigma\) and \(\Pi\), and the rule
\begin{equation}
 +1:m\longmapsto\Sigma,
 \qquad
 -1:m\longmapsto\Pi,
 \qquad
 0:m\longmapsto m.
 \label{p12-83-c9-s3:eq:regla-modal-primitiva-phi}
\end{equation}
Under the cyclic closure condition, the modal orbit produced by
\eqref{p12-83-c9-s3:eq:regla-modal-primitiva-phi} is
\[
 \boxed{(\Sigma,\Pi,\Pi).}
\]
The three consecutive pairs are
\[
 \Sigma\longrightarrow\Pi,
 \qquad
 \Pi\longrightarrow\Pi,
 \qquad
 \Pi\longrightarrow\Sigma.
\]

\begin{definition}[Chronological incidence matrix]
For \(i,j\in\{\Sigma,\Pi\}\), let
\[
 N_3(i,j)=
 \#\{r\in\mathbb Z/3\mathbb Z:(m_r,m_{r+1})=(i,j)\}.
\]
The subscript \(3\) counts instants of the primitive calendar; it does not denote a
matrix power.
\end{definition}

\begin{theorem}[Necessary descent of modal incidence]
\label{p12-83-c9-s3:thm:descenso-incidencia-modal-phi}
In the order \((\Sigma,\Pi)\),
\begin{equation}
 \boxed{
 N_3=
 \begin{pmatrix}
 0&1\\
 1&1
 \end{pmatrix}.}
 \label{p12-83-c9-s3:eq:N3-phi}
\end{equation}
By exact repetition of the calendar,
\begin{equation}
 \boxed{
 N_9=3N_3,
 \qquad
 N_{27}=9N_3,
 \qquad
 N_{108}=36N_3.}
 \label{p12-83-c9-s3:eq:incidencias-9-27-108-phi}
\end{equation}
These three equalities are equalities of edge counts; they do not assert
\(N_9=N_3^3\), \(N_{27}=N_3^9\) or \(N_{108}=N_3^{36}\).
\end{theorem}

\begin{proof}
\(\Sigma\to\Sigma\) does not occur; each of the other three pairs occurs
once. This determines the four entries of \(N_3\). The calendars of
lengths \(9\), \(27\), and \(108\) contain respectively \(3\), \(9\),
and \(36\) copies of the primitive block. Each repetition multiplies all the
counts by the same integer, which proves
\eqref{p12-83-c9-s3:eq:incidencias-9-27-108-phi}.
\end{proof}

\subsubsection{Publication areal–volumetrica and razon estable \(\varphi^{-2}\) under the measure of Parry}

The modal label and the geometric sheet are distinct objects. The publication chart is
\[
 \mathfrak g_{\rm av}(\Sigma)=\mathscr H_{\rm ar},
 \qquad
 \mathfrak g_{\rm av}(\Pi)=\mathscr H_{\rm vol}.
\]
Transporting \eqref{p12-83-c9-s3:eq:N3-phi} through \(\mathfrak g_{\rm av}\), we obtain
the incidence operator
\begin{equation}
 \boxed{
 F_{\rm av}=
 \begin{pmatrix}
 0&1\\
 1&1
 \end{pmatrix}.}
 \label{p12-83-c9-s3:eq:Fav-derivada-phi}
\end{equation}
Thus, the two cross-transitions, the unique volumetric self-retention, and the
absence of areal self-retention are derived from the calendar; they are not four
independent inputs.

\begin{remark}[Scope of the descent]
Forgetting the elementary phase does not by itself produce a strong Markov
quotient of the Topological Phase Kernel, because the visible mode does not determine
which of the three internal operators will act at the next step. What
does descend without an additional hypothesis is the edge multigraph of a
complete block. Its minimal memory-one envelope has adjacency
matrix \(F_{\rm av}\).
\end{remark}

\subsection{Invariant positive ray and uniqueness of self-scaling}

\begin{theorem}[Perron ray derived]
\label{p12-83-c9-s3:thm:rayo-perron-phi}
The positive cone of \(F_{\rm av}\) contains a unique eigenray. If we
normalize it as \(v=(1,x)^{\mathsf T}\), then
\[
 x^2-x-1=0,
 \qquad
 x>0,
\]
and, therefore,
\begin{equation}
 \boxed{x=\frac{1+\sqrt5}{2}=\varphi.}
 \label{p12-83-c9-s3:eq:phi-perron}
\end{equation}
\end{theorem}

\begin{proof}
The equation \(F_{\rm av}v=\lambda v\) gives, coordinate by coordinate,
\[
 x=\lambda,
 \qquad
 1+x=\lambda x.
\]
Eliminating \(\lambda\) yields \(x^2-x-1=0\). Its roots are \((1\pm\sqrt5)/2\), and only the positive root lies in the interior of the cone. The matrix \(F_{\rm av}\) is primitive, since \(F_{\rm av}^2\) has all its entries positive; the Perron--Frobenius theorem ensures that this positive ray is simple and unique. The name \(\varphi\) is assigned only after this construction.
\end{proof}

\begin{corollary}[Scalar Self-Similarity Equation]
\[
 \varphi=1+\frac1\varphi,
 \qquad
 \varphi^2=\varphi+1,
 \qquad
 \varphi^{-1}=\varphi-1.
\]
Thi identitiis are consequencis of \eqref{p12-83-c9-s3:eq:phi-perron}; they select neither
the seed nor the matrix.
\end{corollary}

\subsection{Fibonacci Recurrence of the Perron Ray and Diophantine Embedding of Rational Gaps}

Let \(F_0=0\), \(F_1=1\), and
\[
 F_{n+1}=F_n+F_{n-1}\qquad(n\geq1).
\]

\begin{theorem}[Exact Powers of Incidence]
\label{p12-83-c9-s3:thm:potencias-fibonacci-phi}
For every \(n\geq1\),
\begin{equation}
 \boxed{
 F_{\rm av}^{\,n}=
 \begin{pmatrix}
 F_{n-1}&F_n\\
 F_n&F_{n+1}
 \end{pmatrix}.}
 \label{p12-83-c9-s3:eq:potencias-Fav-fibonacci}
\end{equation}
\end{theorem}

\begin{proof}
For \(n=1\), the identity is the definition of \(F_{\rm av}\). If it holds for \(n\), right multiplication by \(F_{\rm av}\) gives
\[
 \begin{pmatrix}
 F_n&F_{n-1}+F_n\\
 F_{n+1}&F_n+F_{n+1}
 \end{pmatrix}
 =
 \begin{pmatrix}
 F_n&F_{n+1}\\
 F_{n+1}&F_{n+2}
 \end{pmatrix},
\]
which is the case \(n+1\).
\end{proof}

\begin{theorem}[Cassini's Identity]
\label{p12-83-c9-s3:thm:cassini-phi}
For every \(n\geq1\),
\begin{equation}
 \boxed{F_{n+1}F_{n-1}-F_n^2=(-1)^n.}
 \label{p12-83-c9-s3:eq:cassini-phi}
\end{equation}
\end{theorem}

\begin{proof}
The determinant of \(F_{\rm av}\) is \(-1\). Taking determinants in
\eqref{p12-83-c9-s3:eq:potencias-Fav-fibonacci},
\[
 F_{n-1}F_{n+1}-F_n^2
 =\det(F_{\rm av}^{\,n})=(-1)^n.
\]
\end{proof}

We define \(r_n=F_{n+1}/F_n\) for \(n\geq1\). From Cassini it follows that
\[
 r_n-r_{n-1}
 =\frac{F_{n+1}F_{n-1}-F_n^2}{F_nF_{n-1}}
 =\frac{(-1)^n}{F_nF_{n-1}}.
\]
Therefore, the approximations alternate around the Perron ray.

\begin{proposition}[Perron Rational Gaps]
\label{p12-83-c9-s3:prop:horquillas-perron-phi}
For \(m\geq1\),
\begin{equation}
 \frac{F_{2m+2}}{F_{2m+1}}
 <\varphi<
 \frac{F_{2m+1}}{F_{2m}},
 \label{p12-83-c9-s3:eq:horquilla-fibonacci-phi}
\end{equation}
and the width is
\begin{equation}
 \boxed{
 \frac{F_{2m+1}}{F_{2m}}
 -\frac{F_{2m+2}}{F_{2m+1}}
 =\frac{1}{F_{2m}F_{2m+1}}.}
 \label{p12-83-c9-s3:eq:anchura-fibonacci-phi}
\end{equation}
\end{proposition}

\begin{proof}
The preceding alternation places the even quotients below and the odd ones
above \(\varphi\). The formula for the width is
\eqref{p12-83-c9-s3:eq:cassini-phi} with \(n=2m+1\), after putting both fractions over a
common denominator. Since \(F_n\to\infty\), the width tends to zero.
\end{proof}

\subsection{Memory incidence of order \(1\) and Parry measure of the symbolic envelope}

Let \(X_{\rm av}\) be the smallest subshift of finite type over
\(\{\mathscr H_{\rm ar},\mathscr H_{\rm vol}\}\) that admits exactly
the three pairs observed in the modal calendar. Its adjacency matrix is
\(F_{\rm av}\).

\begin{theorem}[Transition and Parry Measure]
\label{p12-83-c9-s3:thm:parry-phi}
Let \(r=(1,\varphi)^{\mathsf T}\). The Parry stochastic matrix is
\begin{equation}
 \boxed{
 P_{ij}=
 \frac{(F_{\rm av})_{ij}r_j}{\varphi r_i}
 =
 \begin{pmatrix}
 0&1\\
 \varphi^{-2}&\varphi^{-1}
 \end{pmatrix}.}
 \label{p12-83-c9-s3:eq:parry-matriz-phi}
\end{equation}
The measure estacionaria of vertices is proporcional to
\((1,\varphi^2)\), hence
\begin{equation}
 \boxed{
 \frac{\mu_{\rm ar}}{\mu_{\rm vol}}=\varphi^{-2}.}
 \label{p12-83-c9-s3:eq:parry-razon-phi}
\end{equation}
\end{theorem}

\begin{proof}
The sum of row \(i\) of \(P\) is
\[
 \frac{(F_{\rm av}r)_i}{\varphi r_i}=1.
\]
Since \(F_{\rm av}\) is symmetric, the left and right eigenvectors
coincide. The stationary weights are proportional to their coordinatewise
products, that is, to \((1,\varphi^2)\). The ratio between the
two weights is \(\varphi^{-2}\).
\end{proof}

\begin{remark}[Three measures that must not be confused]
The chronological frequency of the primitive orbit is
\[
 \nu_{\rm cyc}(\mathscr H_{\rm ar})=\frac13,
 \qquad
 \nu_{\rm cyc}(\mathscr H_{\rm vol})=\frac23.
\]
The uniform measure \(\tau_{468}\) lives on the emission catalog. The
Parry measure lives on admissible histories of arbitrary length. The
first counts instants, the second counts emissions, and the third weights
dynamic cylinders. None is an approximation of the others.
\end{remark}

\subsection{Multiplicative character and cylindrical conservation}

For an admissible route \(\gamma:i\to j\) of length \(n\), we define
\begin{equation}
 \chi_P(\gamma)=\varphi^n\frac{r_i}{r_j}.
 \label{p12-83-c9-s3:eq:caracter-parry-phi}
\end{equation}

\begin{proposition}[Concatenation multiplicativity]
\label{p12-83-c9-s3:prop:multiplicatividad-parry-phi}
If \(\gamma_1:i\to j\) and \(\gamma_2:j\to k\), then
\[
 \chi_P(\gamma_2\circ\gamma_1)
 =\chi_P(\gamma_2)\chi_P(\gamma_1).
\]
On to closed route of length \(n\),
\[
 \chi_P(\gamma)=\varphi^n.
\]
\end{proposition}

\begin{proof}
Length adds and the intermediate coordinate telescopes:
\[
 \varphi^{n_1+n_2}\frac{r_i}{r_k}
 =\left(\varphi^{n_1}\frac{r_i}{r_j}\right)
  \left(\varphi^{n_2}\frac{r_j}{r_k}\right).
\]
If \(i=k\), the quotient terminal is one.
\end{proof}

We normalize the left vector as
\[
 \ell=\frac{r}{1+\varphi^2},
 \qquad
 \ell^{\mathsf T}r=1.
\]
The measure of a cylinder is
\[
 \mu[x_0\cdots x_n]
 =\ell_{x_0}r_{x_n}\varphi^{-n}.
\]
Therefore,
\begin{equation}
 V_n(x)
 :=\frac{\mu[x_0]}{\mu[x_0\cdots x_n]}
 =\varphi^n\frac{r_{x_0}}{r_{x_n}}
 =\chi_P(x_0\cdots x_n).
 \label{p12-83-c9-s3:eq:volumen-cilindrico-phi}
\end{equation}
For to typed dimension \(D>0\), the scale
\[
 a_n^{(D)}(x)=V_n(x)^{1/D}
\]
satisfies the exact law
\begin{equation}
 \boxed{
 \bigl(a_n^{(D)}(x)\bigr)^D
 \mu[x_0\cdots x_n]=\mu[x_0].}
 \label{p12-83-c9-s3:eq:conservacion-cilindrica-phi}
\end{equation}

In three dimensions and over closed returns,
\[
 \frac{a_9}{a_0}=\varphi^3,
 \qquad
 \frac{a_{27}}{a_0}=\varphi^9,
 \qquad
 \frac{a_{108}}{a_0}=\varphi^{36}.
\]
If \(A_k=a_{9k}\), then
\begin{equation}
 A_{k+1}-2A_k+A_{k-1}
 =A_{k-1}(\varphi^3-1)^2>0.
 \label{p12-83-c9-s3:eq:convexidad-retornos-phi}
\end{equation}
This is convexity in return depth; it is not interpreted as a
kinematic quantity without an additional duration chart.

\subsection{Quadratic memory and contracted branch}

The eigenvalues of \(F_{\rm av}\) are
\[
 \varphi,
 \qquad
 -\varphi^{-1}.
\]
The contracted branch changes orientation at each iteration. To preserve
that sign before squaring, we pass to
\(\operatorname{Sym}^2(\mathbb R^2)\). In the basis
\((x^2,2xy,y^2)\),
\begin{equation}
 \operatorname{Sym}^2(F_{\rm av})=
 \begin{pmatrix}
 0&0&1\\
 0&1&2\\
 1&1&1
 \end{pmatrix}.
 \label{p12-83-c9-s3:eq:sym2-Fav-phi}
\end{equation}

\begin{theorem}[Spectrum of the second-order memory]
\label{p12-83-c9-s3:thm:espectro-sym2-phi}
\[
 \det\bigl(tI-\operatorname{Sym}^2(F_{\rm av})\bigr)
 =(t+1)(t^2-3t+1),
\]
and
\begin{equation}
 \boxed{
 \operatorname{Spec}\bigl(\operatorname{Sym}^2(F_{\rm av})\bigr)
 =\{\varphi^2,-1,\varphi^{-2}\}.}
 \label{p12-83-c9-s3:eq:espectro-sym2-phi}
\end{equation}
The only positive stable mode is \(\varphi^{-2}\).
\end{theorem}

\begin{proof}
The eigenvalues of a degree-two symmetric power are the products
\(\lambda_1^2\), \(\lambda_1\lambda_2\), and \(\lambda_2^2\). With
\(\lambda_1=\varphi\) and \(\lambda_2=-\varphi^{-1}\), one obtains
\(\varphi^2\), \(-1\), and \(\varphi^{-2}\). The first exceeds one, the
second has modulus one, and the third belongs to \((0,1)\).
\end{proof}

On a closed route with \(V_n=\varphi^n\), the stable mode satisfies
\begin{equation}
 \boxed{M_n=\varphi^{-2n}=V_n^{-2}.}
 \label{p12-83-c9-s3:eq:modo-estable-phi}
\end{equation}

\subsection{Normalized hyperbolic realization of the self-scale: \(\exp(2\log\varphi\,J_h)=F_{\mathrm{av}}^2\)}

Let
\[
 A=F_{\rm av}^{\,2}=
 \begin{pmatrix}1&1\\1&2\end{pmatrix},
 \qquad
 J_h=\frac{2A-3I}{\sqrt5}.
\]
Then
\[
 J_h^2=I,
 \qquad
 A=\frac32I+\frac{\sqrt5}{2}J_h.
\]
Since
\[
 \cosh(2\log\varphi)=\frac32,
 \qquad
 \sinh(2\log\varphi)=\frac{\sqrt5}{2},
\]
we obtain
\begin{equation}
 \boxed{
 \exp(2\log\varphi\,J_h)=F_{\rm av}^{\,2}.}
 \label{p12-83-c9-s3:eq:realizacion-hiperbolica-phi}
\end{equation}
The self-scaling constant is therefore simultaneously the Perron
eigenvalue of the incidence and the exponential parameter of its normalized
hyperbolic realization.

\subsection{Characterization, brackets and prolongation}

The finite construction determines the positive ray of \(F_{\rm av}\). The
subsequent analytic characterization is
\[
 x^2-x-1=0,
 \qquad x>0.
\]
The brackets \eqref{p12-83-c9-s3:eq:horquilla-fibonacci-phi} have rational endpoints
and explicit width \eqref{p12-83-c9-s3:eq:anchura-fibonacci-phi}. For each
nonadic level \(K\), we choo the least \(m=m_{\varphi}(K)\) who bracket remains
contained in a single subcylinder among the \(729\) elements of the
ambient fiber. Minimality converts localization into a function, not into to
tacit choice.

The selected cylinders are compatible with the truncations, and their
diameters tend to zero. Their intersection contains a unique real number; the
\cref{p12-83-c9-s3:thm:rayo-perron-phi} identifies it with \(\varphi\). The Fibonacci
quotients are therefore rational certificates of localization, not decimal data
used to decide the seed.

\subsection{Internal certificate of uniqueness, compatibility, and convergence of self-scaling}

\begin{criterion}[Refutation of the self-scale biography]
The construction remains refuted if any of the following facts
occurs:
\begin{enumerate}
 \item the finite predicate selects more than one orbit or more than one emission;
 \item the recurrence \(L_0,L_0,L_1,L_1\) does not produce
       \eqref{p12-83-c9-s3:eq:w30-phi};
 \item the calendar does not produce exactly the three pairs used in
       \eqref{p12-83-c9-s3:eq:N3-phi};
 \item a count multiple \(N_{9},N_{27},N_{108}\) is confused with a
       power of \(N_3\);
 \item \(F_{\rm av}\) has another positive eigenray;
 \item \eqref{p12-83-c9-s3:eq:potencias-Fav-fibonacci} or Cassini’s identity fails;
 \item the rational brackets are not nested or do not contract their width;
 \item the chronological frequency, the uniform measure on the
       catalog or the Parry measure are identified;
 \item a decimal table is involved before the Perron ray is constructed;
 \item a selected cylinder does not truncate to the previous cylinder.
\end{enumerate}
\end{criterion}


% Unidad U015. Acoplamiento de horquillas racionales, refinamiento TPK y
% certificación a profundidad arbitraria.

\section{Diagonal certification, inverse limit and generation to arbitrary precision}
\label{p12-83-c9-s4:chap:precision-arbitraria-constantes}

\subsection{Typed separation of the \(4\) indices of nonadic prolongation}

This unit has the status of subsequent certification. The prefixes and their
compatible extensions come from the internal relation
\(\operatorname{Ext}\), coinductive survival and the APP--TRIT--TPK
transition laws. The rational brackets that follow do not choose
among the \(729\) elements of a fiber: they locate in the Archimedean chart
the already generated extension and certify its precision at arbitrary depth.

Four parameters of different kinds enter each character \(\chi\in\{\pi,e,\varphi\}\):
\begin{enumerate}
 \item \(n\), order of the rational analytic bracket;
 \item \(K\), refinement level of the Topological Phase Kernel;
 \item \(L_K=6K\), ternary length of the published prefix;
 \item \(M\), the requested number of decimal positions.
\end{enumerate}
Neither \(n=K\), \(n=L_K\), \(K=M\) nor a fixed proportionality
between them is postulated. The coupling is constructed through interval inclusions
decided by integer arithmetic.

The integer parts are
\[
 m_\pi=3,
 \qquad
 m_e=2,
 \qquad
 m_\varphi=1,
\]
and the fractional characters are
\[
 \rho_\pi=\pi-3,
 \qquad
 \rho_e=e-2,
 \qquad
 \rho_\varphi=\varphi-1.
\]
The three families of semiopen forks are written
\[
 J_{\chi,n}=[\ell_{\chi,n},h_{\chi,n}).
\]
Their endpoints are rational, contain \(\rho_\chi\) and satisfy
\[
 \ell_{\chi,n}\nearrow\rho_\chi,
 \qquad
 h_{\chi,n}\searrow\rho_\chi,
 \qquad
 h_{\chi,n}-\ell_{\chi,n}\longrightarrow0.
\]
For \(\pi\), the endpoints come from the alternating sums of Machin’s
identity proved in \cref{chap:biografia-pi-autonoma}; for \(e\),
from the factorial sums and their rational remainder in
\cref{p12-83-c9-s2:chap:biografia-completa-e}; for \(\varphi\), from the Fibonacci
and Cassini quotients in \cref{p12-83-c9-s3:chap:biografia-completa-phi}.

\subsection{Ternary cylinders and the immediate ambient fibre}

A ternary prefix of length \(L_K=6K\) is represented by an integer
\(N_K\) with
\[
 0\leq N_K<3^{L_K}.
\]
His semi-open cylinder is
\begin{equation}
 C(N_K,L_K)=
 \left[
 \frac{N_K}{3^{L_K}},
 \frac{N_K+1}{3^{L_K}}
 \right).
 \label{p12-83-c9-s4:eq:cilindro-nivel-K}
\end{equation}
The next emission adds six trits. The immediate ambient fiber contains
\(3^6=729\) elements, indexed by \(u\in\{0,\ldots,728\}\):
\begin{equation}
 C_{K,u}=
 \left[
 \frac{729N_K+u}{3^{L_K+6}},
 \frac{729N_K+u+1}{3^{L_K+6}}
 \right).
 \label{p12-83-c9-s4:eq:729-subcilindros}
\end{equation}
The \(C_{K,u}\) are disjoint and
\[
 C(N_K,L_K)=\bigsqcup_{u=0}^{728}C_{K,u}.
\]
The number \(729\) counts coordinates of this ambient fiber; it does not count
seeds, distinct emissions, orbits, or surviving states.

\subsubsection{Diophantine comparison of the rational endpoints of the gap}

Let \(J=[\ell,h)\subset C(N_K,L_K)\), with
\(\ell,h\in\mathbb Q\). Escribimos \(S_K=3^{L_K+6}\) and define
\begin{align}
 j_{\min}(J,K)
 &=\max\!\left(729N_K,\left\lfloor S_K\ell\right\rfloor\right),
 \label{p12-83-c9-s4:eq:jmin-diagonal}\\
 j_{\max}(J,K)
 &=\min\!\left(729N_K+728,\left\lceil S_Kh\right\rceil-1\right).
 \label{p12-83-c9-s4:eq:jmax-diagonal}
\end{align}

\begin{lemma}[Unitary localization criterion]
\label{p12-83-c9-s4:lem:criterio-localizacion-unitaria}
The bracket \(J\) is contained in a unique element of the fiber
\eqref{p12-83-c9-s4:eq:729-subcilindros} if and only if
\begin{equation}
 \boxed{j_{\min}(J,K)=j_{\max}(J,K).}
 \label{p12-83-c9-s4:eq:jmin-igual-jmax}
\end{equation}
In that case, if \(j\) is the common value,
\[
 u=j-729N_K,
 \qquad
 N_{K+1}=j.
\]
\end{lemma}

\begin{proof}
The integers \(j\) whose elementary intervals of length \(S_K^{-1}\)
intersect \(J\) form the integer interval
\[
 \left[
 \lfloor S_K\ell\rfloor,
 \lceil S_Kh\rceil-1
 \right]\cap
 [729N_K,729N_K+728].
\]
This contains exactly one element if and only if its endpoints coincide.
The equality identifies the numerator of the subcylinder and, by truncating its last
six trits, \(N_K\) is recovered.
\end{proof}

\subsection{Minimal diagonal certification}

Generation does not fix the same analytic order for all
levels in advance. Once the compatible extension has been produced, the order of the
bracket is increased only until it is certified uniquely.

\begin{definition}[Diagonal certification recurrence]
\label{p12-83-c9-s4:def:recurrencia-diagonal-constantes}
Given a family of compatible cylinders fixed by \(\operatorname{Ext}\),
\(C(N_K,L_K)\), whose publication contains \(\rho_\chi\), an initial order
\(n_\chi(K_0-1)\geq1\) is chosen and we recursively define
\begin{equation}
 n_\chi(K)=
 \min\left\{
 n\geq n_\chi(K-1):
 j_{\min}(J_{\chi,n},K)=j_{\max}(J_{\chi,n},K)
 \right\},
 \label{p12-83-c9-s4:eq:nchi-K-minimo}
\end{equation}
and we verify
\begin{equation}
 N_{\chi,K+1}
 =j_{\min}(J_{\chi,n_\chi(K)},K).
 \label{p12-83-c9-s4:eq:Nchi-K-mas-uno}
\end{equation}
The set
\[
 \Delta_\chi=
 \{(n_\chi(K),K):K\geq K_0\}
\]
is the analytic--nonadic diagonal of the character \(\chi\).
\end{definition}

\begin{theorem}[Existence and uniqueness of the certification diagonal]
\label{p12-83-c9-s4:thm:existencia-unicidad-diagonal}
For \(\chi\in\{\pi,e,\varphi\}\), the recurrence
\eqref{p12-83-c9-s4:eq:nchi-K-minimo} is defined at every finite level. The index
\(n_\chi(K)\) and the coordinate \(N_{\chi,K+1}\) are unique.
\end{theorem}

\begin{proof}
Suppose the cylindrical extension of level \(K\) containing
\(\rho_\chi\) has been generated internally. Since \(\rho_\chi\) is irrational, it belongs to no
ternary boundary \(3^{-L}\mathbb Z\). It therefore lies in the interior of a
unique subcylinder \(C_{K,u_*}\), at a positive distance from both endpoints.
The brackets \(J_{\chi,n}\) contain \(\rho_\chi\) and their diameter
tends to zero; hence there exists an \(n\) for which the bracket is
contained in \(C_{K,u_*}\). The set over which the minimum is taken in
\eqref{p12-83-c9-s4:eq:nchi-K-minimo} is nonempty and, by the well-ordering of
\(\mathbb N\), has a least element.
\cref{p12-83-c9-s4:lem:criterio-localizacion-unitaria} proves that the bracket
locates it uniquely and that the located numerator agrees with the one already
generated. Induction on \(K\) completes the argument.
\end{proof}

\begin{corollary}[Absence of decimal parameters]
The diagonal is computed using sums, products, powers, Euclidean divisions,
floors and ceilings of integers. It receives no decimal digit of \(\pi\), \(e\) or
\(\varphi\) as certification data.
\end{corollary}

\subsection{Compatibility with the joint nonadic connection}

Let \(\widehat x_{\chi,K}\) be the integral holonomic state of channel \(\chi\) at level \(K\). Its coordinates include the prefix, cylinder, phase, quotient, carry, unbounded vertical memory, terminal term and provenance record. The residual phase is
\[
 g_0(K)=[K]_9\in\mathbb Z/9\mathbb Z,
\]
while the public label is
\[
 g(K)=1+((K-1)\bmod9)\in\{1,\ldots,9\}.
\]
The vertical memory \(q_{\rm mem}(K)\in\mathbb Z_{\geq0}\) is not reduced
modulo \(12\); only its dodecaphasic projection is
\[
 \beta(K)=q_{\rm mem}(K)\bmod12
\]
does so.

For each \(K\), let
\[
 \Gamma_{9,K}:\widehat X_K\longrightarrow\widehat X_{K+9}
\]
the holonomy of a phase loop. Then
\[
 g_0(K+9)=g_0(K),
 \qquad
 q_{\rm mem}(K+9)=q_{\rm mem}(K)+1,
 \qquad
 \beta(K+9)=\beta(K)+1\pmod{12}.
\]

\begin{theorem}[Phase return with genealogical conservation]
\label{p12-83-c9-s4:thm:retorno-fase-conservacion-genealogica}
Phase equality does not reset the state:
\[
 g_0(K+9)=g_0(K),
 \qquad
 \widehat x_{\chi,K+9}\ne\widehat x_{\chi,K}.
\]
The depths \(27\), \(54\), and \(108\) are the typed compositions
\begin{align*}
 \Gamma_{27,K}
 &=\Gamma_{9,K+18}\circ\Gamma_{9,K+9}\circ\Gamma_{9,K},\\
 \Gamma_{54,K}
 &=\Gamma_{9,K+45}\circ\cdots\circ\Gamma_{9,K},\\
 \Gamma_{108,K}
 &=\Gamma_{9,K+99}\circ\cdots\circ\Gamma_{9,K}.
\end{align*}
In no case are they restarts of the prefix, the carry, the terminal, or the memory.
\end{theorem}

\begin{proof}
The first coordinate is calculated in \(\mathbb Z/9\mathbb Z\), whereas \(q_{\rm mem}\) belongs to an unbounded monoid. Each application of \(\Gamma_{9,K}\) adds one memory unit, six blocks per intermediate level and the corresponding composition of cocycles. The phase may therefore coincide even when the integral holonomic state does not. The depth formulas show the indices of each map and prevent composition as though they were all endomorphisms of the same space.
\end{proof}

\subsubsection{Compression, expression, and terminal term}

The operational realization of a loop satisfies
\begin{equation}
 U_\Gamma J=JT+\eta,
 \qquad
 J^*\eta=0,
 \qquad
 I-T^*T=\eta^*\eta.
 \label{p12-83-c9-s4:eq:compresion-expresion-precision}
\end{equation}
The visible strict contraction is reserved for
\[
 M_9=\frac59V_9;
\]
is not identified with the compression operator \(T\). Telescoping at
depth \(L\) preserves the terminal term:
\begin{equation}
 S_0=
 \sum_{r=0}^{L-1}M_9^{*r}D_9^*D_9M_9^r
 +M_9^{*L}S_0M_9^L.
 \label{p12-83-c9-s4:eq:telescopia-terminal-precision}
\end{equation}
Omitting the last summand before proving its disappearance would erase
boundary information and would turn phase return into a false
restart.

\subsection{Inverse system of compatible histories}

Let \(\mathcal H_K^\chi\) be the set of admissible enriched histories
of channel \(\chi\) at depth \(K\), and let
\[
 p_{N,K}:\mathcal H_N^\chi\longrightarrow\mathcal H_K^\chi
 \qquad(N\geq K)
\]
the truncation. We define
\[
 \operatorname{Surv}_K^{(N)}(\chi)
 =p_{N,K}(\mathcal H_N^\chi)
\]
and
\begin{equation}
 \operatorname{Surv}_K(\chi)
 =\bigcap_{N\geq K}\operatorname{Surv}_K^{(N)}(\chi).
 \label{p12-83-c9-s4:eq:supervivencia-todas-profundidades}
\end{equation}
The coinductive object is
\begin{equation}
 X_\chi=
 \varprojlim_K
 \bigl(\operatorname{Surv}_K(\chi),p_{K+1,K}\bigr).
 \label{p12-83-c9-s4:eq:limite-inverso-caracteres}
\end{equation}

\begin{theorem}[Section compatible under enhanced extendibility]
\label{p12-83-c9-s4:thm:seccion-compatible-caracteres}
Suppose further that, over each generated numerical subcylinder, the
relation \(\operatorname{Ext}\) is nonempty, single-valued and compatible with the
truncations. Then \(\operatorname{Ext}\) and survival produce a family
\[
 x_\chi=(\widehat x_{\chi,K})_{K\geq K_0}
\]
which satisfies
\[
 p_{K+1,K}(\widehat x_{\chi,K+1})
 =\widehat x_{\chi,K}.
\]
Therefore, \(x_\chi\in X_\chi\).
\end{theorem}

\begin{proof}
The subcylinder of level \(K+1\) has numerator
\(N_{\chi,K+1}=729N_{\chi,K}+u\). Truncation removes its last six
trits and recovers \(N_{\chi,K}\). The other coordinates are obtained through the
functorial laws of carry, memory, terminal and phase. The single-valuedness
hypothesis for \(\operatorname{Ext}\) fixes the subcylinder and the enriched
extension; the local uniqueness of
\cref{p12-83-c9-s4:lem:criterio-localizacion-unitaria} subsequently certifies its
Archimedean localization, and the compatibility of \(\operatorname{Ext}\)
provides truncation equality.
\end{proof}

\subsection{Decimal cells and prescribed precision}

For \(M\geq1\) and \(q\in\{0,\ldots,10^M-1\}\), let
\[
 D_{M,q}=
 \left[
 \frac{q}{10^M},
 \frac{q+1}{10^M}
 \right).
\]
A bracket \(J=[\ell,h)\) is contained in a unique decimal cell if
and only if
\begin{equation}
 \boxed{
 \left\lfloor10^M\ell\right\rfloor
 =\left\lceil10^Mh\right\rceil-1=q.}
 \label{p12-83-c9-s4:eq:criterio-celda-decimal}
\end{equation}

\begin{definition}[Minimum level of decimal certification]
For \(\chi\in\{\pi,e,\varphi\}\), we defines \(K_\chi(M)\) as the least
level for which there exist an analytic order \(n\) and an integer \(q\) such
that
\begin{enumerate}
 \item \(J_{\chi,n}\) satisfies
       \eqref{p12-83-c9-s4:eq:criterio-celda-decimal};
 \item \(J_{\chi,n}\) is contained in the cylinder generated by
       state \(\widehat x_{\chi,K}\);
 \item the transition to the next level satisfies
       \eqref{p12-83-c9-s4:eq:jmin-igual-jmax}.
\end{enumerate}
The pair \((K,n)\) is minimized lexicographically.
\end{definition}

\begin{theorem}[Certification of coinductive outputs to arbitrary decimal precision]
\label{p12-83-c9-s4:thm:precision-decimal-arbitraria}
For each \(\chi\in\{\pi,e,\varphi\}\) and each \(M\geq1\), there exists
\(K_\chi(M)<\infty\). The generated section determines a unique decimal prefix
of \(M\) positions, and the bracket recognizes and certifies that prefix. If
\(M'<M\), the prefix of order \(M\) truncates to the one of
order \(M'\), and the corresponding cells are nested.
\end{theorem}

\begin{proof}
The three characters are irrational: \(\varphi\) is irrational by its quadratic
polynomial with nonsquare discriminant; the irrationality of \(e\) and
\(\pi\) is taken from their classical characterization theorems, subsequent to
the finite construction. Thus none belongs to a rational decimal or
ternary boundary. For a fixed \(M\), the distance from the character to the finite
set of relevant decimal boundaries is positive. The rational
brackets contract their diameter until they lie in a unique decimal cell.
\cref{p12-83-c9-s4:thm:existencia-unicidad-diagonal} couples that bracket to a
TPK cylinder and a unique extension. The lexicographic minimum exists by
well-ordering. Compatibility of the section proves nesting as
\(M\) increases.
\end{proof}

\begin{corollary}[Real uniqueness]
The nested cells have diameters \(10^{-M}\to0\). Their intersection is
a singleton, identified by the preceding characterizations with
\(\pi\), \(e\) or \(\varphi\), respectively.
\end{corollary}

\subsection{Equivalence of \(2\) finite realizations of the uniform law}

The witnesses of \(1,000\) and \(10,000\) positions do not define the
characters or limit the prolongation. They are executions of the same uniform
recurrence at two different depths.

\begin{table}[htbp]
\centering
\caption{Exact dimensions of certified runs, by channel.}
\label{p12-83-c9-s4:tab:dimensiones-ejecuciones-1000-10000}
\begin{tabular}{lrr}
\toprule
Counted object & \(1,000\) positions & \(10,000\) positions\\
\midrule
Partitions exhaustive of the fiber ambient & 396 & 3501\\
Trits transportados & 2376 & 21006\\
Pairs of equal phase \((K,K+9)\) & 387 & 3492\\
Complete nonadic turns & 43 & 388\\
Elements examined by partition & 729 & 729\\
Compatible extensions preserved by partition & 1 & 1\\
\bottomrule
\end{tabular}
\end{table}

Para \(10,000\) posiciones,
\[
 B=3501=389\cdot9,
 \qquad
 6B=21006,
 \qquad
 3^{6B}>10^{10000}.
\]
The first \(396\) blocks of the long run coincide with the
run of \(1,000\) positions. In each of the \(3492\) pairs of
equal phase, the phase coordinate is preserved and the prefix, depth,
cylinder, and memory are extended. The certificate also records any
repetitions of local projections without conflating them with the return of
the complete state.

\subsection{Separation between generation and verification}

\begin{definition}[Structural generator]
The generator receives the finite catalog, \(L_0,L_1\), the initial state, the
connection laws and the internal relation \(\operatorname{Ext}\). At each level, it:
\begin{enumerate}
 \item enumerates the \(729\) elements of the ambient fiber;
 \item applies \(\operatorname{Ext}\) and retains the unique compatible extension;
 \item updates phase, quotient, carry, memory, and terminal;
 \item writes the ternary block, the selection record and the genealogical record.
\end{enumerate}
\end{definition}

\begin{definition}[Independent post hoc verification]
The independent subsequent check receives the closed outputs of the
generator and the three families of rational brackets; it computes
\(j_{\min}\) and \(j_{\max}\), certifies their equality with the extension already
produced, reconstructs the
partitions, checks the truncation squares, the pairs
\((K,K+9)\), the partition sums
\[
 \operatorname{killed}_{\mathrm{izq}}+1+\operatorname{killed}_{\mathrm{der}}=729,
\]
the integrity of the outputs and, finally, the agreement with reference expansions.
\end{definition}

\begin{theorem}[Isolation of the target data]
\label{p12-83-c9-s4:thm:aislamiento-datos-objetivo-precision}
Reference decimal expansions and metrological data are not
involved in selecting a compatible extension. The comparison is
performed after the three prefixes have been written.
\end{theorem}

\begin{proof}
The mathematical input to each generative decision consists of the finite
catalog, \(L_0,L_1\), the inherited state and the internal relation
\(\operatorname{Ext}\). The rational endpoints, powers of \(3\), floors and
ceilings belong to the subsequent check. The generator finishes before
that check opens the brackets or runs the decimal controls. The direction
of the data flow is
\[
 \text{structure}\longrightarrow
 \text{prefixes and ledgers}\longrightarrow
 \text{independent subsequent check};
\]
there is no arrow in the opposite direction.
\end{proof}

\subsection{Differentiated obstructions to generation and its certification at arbitrary precision}

\begin{criterion}[Refutation of generation to arbitrary precision]
Failures of orbit, extension, truncation or return refute the generation;
failures of brackets, order or decimal control refute the corresponding
analytic certificate; and a discrepancy between the two refutes the
joint identification of the published output. With this separation of types,
the complete assertion is refuted if:
\begin{enumerate}
 \item any of the indices \(n,K,L_K,M\) are identified without a
       linking theorem;
 \item a bracket ceases to contain its character or does not contract its diameter;
 \item a stage satisfies \(j_{\min}\ne j_{\max}\) after the order
       declared to be minimal;
 \item an extension does not truncate to the cylinder of the previous level;
 \item the phase return resets the prefix, carry, terminal or memory;
 \item the terminal term in
       \eqref{p12-83-c9-s4:eq:telescopia-terminal-precision} is omitted without proving that it vanishes;
 \item the first \(396\) blocks of the long execution differ from the
       short execution;
 \item a decimal or metrological control is read before the output is closed;
 \item a printed digit differs from the certified mathematical section;
 \item the proof depends on the requested precision being \(1,000\) or
       \(10,000\), rather than an arbitrary \(M\).
\end{enumerate}
\end{criterion}
 
\chapter{Dodecaphasic closure and two internal derivations of \(\alpha\)}
\label{chap:83-10}

% Unidad U016. Registro firmado, lectura terminal y reconstruccion de K.
% Redaccion autonoma desde la autoridad material 1063 y su sucesor V2.

\section{Initial generation of the signed dodecaphase state and reconstruction of the seal \(K\)}
\label{p12-83-c10-s1:chap:registro-dodecafasico-hadamard-k}

Construction of the integer seal does not begin with \(\alpha\), its inverse or a metrological table. Its domain is the terminal holonomic state of the Topological Phase Kernel obtained after the chain already constructed
\begin{equation}
 \boxed{
 \mathrm{APP}\longrightarrow\mathrm{TRIT}\longrightarrow\mathrm{TPK}
 \longrightarrow\Omega_{\rm seed}
 \xrightarrow{T_{\min}}\mathcal U_6^{\rm cat}
 \xrightarrow{\rho_3}W_{243}
 \longrightarrow w_{30}\longrightarrow R_{36}
 \longrightarrow G_9\longrightarrow x_{\rm term}.}
 \label{p12-83-c10-s1:eq:cadena-upstream-registro-dodecafasico}
\end{equation}
The foregoing arrows were defined with their respective domains:
pairs of seeds, signatures, emissions, ternary words, boundaries,
compatible histories, and nonadic holonomy. This unit does not reselect
any of those data. It constructs the integral charts of the same
terminal state and proves their equivalence.

\subsection{Oriented tracis over windows of length \(108\)}

Let
\[
 \Theta_{108}=\mathbb Z/108\mathbb Z
\]
the observable calendar. Its twelve nonadic windows are
\[
 E_m=\{9m+1,\ldots,9m+9\},
 \qquad 0\leq m<12.
\]
The partition is disjoint and exhaustive:
\[
 \{1,\ldots,108\}=\bigsqcup_{m=0}^{11}E_m.
\]
The set of address codes that publish an event is fixed as
\[
 D_{\rm evt}=\{2,4,6,8\}.
\]
The orientation of the twelve windows is
\begin{equation}
 \Sigma_{12}=(+,+,+,-,-,-,+,+,+,-,-,-).
 \label{p12-83-c10-s1:eq:signatura-doce-ventanas}
\end{equation}
If \(d_t\) is the direction encoded at step \(t\), the signed event
of window \(m\) is
\begin{equation}
 e_m=\Sigma_{12}(m)\,
 \#\{t\in E_m:d_t\in D_{\rm evt}\}.
 \label{p12-83-c10-s1:eq:evento-firmado-ventana}
\end{equation}
The formula distinguishes the nonnegative number of events from the sign of the
orientation. It does not replace a cardinal direction with a sign: both
coordinates remain present in the TPK state.

\begin{definition}[Enhanced dodecaphonic registration]
\label{p12-83-c10-s1:def:registro-enriquecido-dodecafasico}
The record transported by the trace is
\[
 \mathcal R_{12}(x)=
 \bigl((E_m,\tau_m,\sigma_m,\kappa_m,\Lambda_m)\bigr)_{m=0}^{11},
\]
where \(\tau_m\) is the TRIT orientation, \(\sigma_m\) the boundary
signature, \(\kappa_m\) the carry, and \(\Lambda_m\) the genealogical register segment
that preserves the route. There is a chain of projections
\[
 \mathcal X_{\rm TPK}^{\rm enr}
 \longrightarrow\mathcal R_{12}
 \longrightarrow\Theta_{108},
\]
but neither of the two arrows is an identification.
\end{definition}

In particular, \(\Theta_{108}\) preserves only the position in the calendar.
The record \(\mathcal R_{12}\) additionally preserves orientation, signature,
carry, and provenance. The complete state further preserves cylinder,
boundary, quotient, terminal, and vertical memory.

\subsection{Integral chart of the dodecaphasic opposition}

For \(0\leq i<6\), the opposite positions determine
\begin{equation}
 p_i=e_i+e_{i+6},
 \qquad
 a_i=e_i-e_{i+6}.
 \label{p12-83-c10-s1:eq:pares-opuestos-registro}
\end{equation}
The central charge and the five relative margins are
\begin{equation}
 s_C=\sum_{i=0}^{5}p_i,
 \qquad
 M_i=p_i-p_5\quad(0\leq i<5).
 \label{p12-83-c10-s1:eq:carga-margenes-registro}
\end{equation}

\begin{definition}[Integral chart \(1+5+6\)]
\label{p12-83-c10-s1:def:carta-integral-156}
It is defined
\begin{equation}
 \mathcal L_{12}(e_0,\ldots,e_{11})
 =\bigl(s_C,M_0,\ldots,M_4,a_0,\ldots,a_5\bigr).
 \label{p12-83-c10-s1:eq:carta-integral-156}
\end{equation}
The coordinate partition is
\[
 1+5+6=12.
\]
\end{definition}

\begin{theorem}[Exact inverse of the integral chart]
\label{p12-83-c10-s1:thm:inversion-carta-integral-156}
A tuple
\[
 \ell=(s_C,M_0,\ldots,M_4,a_0,\ldots,a_5)\in\mathbb Z^{12}
\]
belongs to the image of \(\mathcal L_{12}\) if and only if
\begin{align}
 s_C-\sum_{i=0}^{4}M_i&\equiv0\pmod6,
 \label{p12-83-c10-s1:eq:congruencia-seis-carta}\\
 p_i&\equiv a_i\pmod2
 \quad(0\leq i<6),
 \label{p12-83-c10-s1:eq:congruencia-paridad-carta}
\end{align}
where
\begin{align}
 p_5&=\frac{s_C-\sum_{i=0}^{4}M_i}{6},
 \label{p12-83-c10-s1:eq:inversa-p5-carta}\\
 p_i&=p_5+M_i\quad(0\leq i<5).
 \label{p12-83-c10-s1:eq:inversa-pi-carta}
\end{align}
In that case, the only entire preimage is
\begin{equation}
 e_i=\frac{p_i+a_i}{2},
 \qquad
 e_{i+6}=\frac{p_i-a_i}{2}.
 \label{p12-83-c10-s1:eq:inversa-eventos-carta}
\end{equation}
\end{theorem}

\begin{proof}
Summing \(M_i=p_i-p_5\) for \(0\leq i<5\), we obtain
\[
 \sum_{i=0}^{4}M_i
 =\sum_{i=0}^{4}p_i-5p_5
 =s_C-6p_5,
\]
which forces \eqref{p12-83-c10-s1:eq:inversa-p5-carta} and the congruence modulo six.
Once the six \(p_i\) have been reconstructed, the system
\[
 p_i=e_i+e_{i+6},
 \qquad a_i=e_i-e_{i+6}
\]
has the unique solution \eqref{p12-83-c10-s1:eq:inversa-eventos-carta}. This solution is
integral exactly when \(p_i\) and \(a_i\) have the same parity.
The two families of congruences are therefore necessary and sufficient.
\end{proof}

This theorem precedes the numerical closure of \(\alpha\). It explains how the
oriented trace produces a reversible integral chart, and why preserving
the reduced calendar is not enough.

\subsection{Selection of the common terminal state}

The three compatible sections reach terminal depth with catalogs
of cardinalities
\[
 |\mathcal T_\pi|=196,
 \qquad
 |\mathcal T_e|=197,
 \qquad
 |\mathcal T_\varphi|=196.
\]
The product contains
\begin{equation}
 196\cdot197\cdot196=7\,567\,952
 \label{p12-83-c10-s1:eq:cardinal-producto-terminal}
\end{equation}
ordered triples.

The column margin of the terminal reader is
\[
 q_\partial=(1,2,2,1,2,0)\in\mathbb F_3^6,
\]
and the row orientation of the combined channel is
\[
 \sigma=(1,1,-1)=(1,1,2)\in\mathbb F_3^3.
\]
The first datum acts on columns; the second acts on rows. Neither is deducible from the other.

The local signatures of Gram matrix, norm, weight, and distance reduce
\eqref{p12-83-c10-s1:eq:cardinal-producto-terminal} to \(331\) triples. The margin
\(q_\partial\) leaves two indices:
\[
 (120,197,157),
 \qquad
 (129,197,148).
\]
Their matrices are
\[
 B_0=
 \begin{pmatrix}
 0&2&0&2&2&0\\
 0&0&1&2&2&0\\
 1&0&1&0&1&0
 \end{pmatrix},
 \qquad
 B_1=
 \begin{pmatrix}
 0&2&1&0&2&0\\
 0&0&1&2&2&0\\
 1&0&0&2&1&0
 \end{pmatrix}.
\]
The row sums are
\[
 (0,2,0),
 \qquad
 (2,2,1).
\]
Since
\[
 \langle\sigma\rangle
 =\{(0,0,0),(1,1,2),(2,2,1)\},
\]
only \(B_1\) satisfies the axial condition.

\begin{theorem}[Uniqueness of the Visible Terminal Publication]
\label{p12-83-c10-s1:thm:unicidad-publicacion-terminal-visible}
The joint selection determines
\begin{equation}
 B_{\rm term}=
 \begin{pmatrix}
 021020\\001220\\100210
 \end{pmatrix},
 \label{p12-83-c10-s1:eq:matriz-terminal-visible}
\end{equation}
corresponding to the triple \((129,197,148)\). No decimal evaluation of
\(\alpha\) enters the census or the two selection conditions.
\end{theorem}

\begin{proof}
The census and margin leave exactly \(B_0,B_1\). The charge of \(B_0\)
does not belong to \(\langle\sigma\rangle\), whereas that of \(B_1\) is
\(2\sigma\). The second candidate is therefore the only one that simultaneously satisfies
both typed requirements.
\end{proof}

\subsection{Archimedean and exceptional publications of the common state: inequivalent codomains and shared genealogy}

Let
\[
 x_{\rm term}\in
 \mathcal X_{\rm TPK}^{\rm term,enr}
 \subseteq\mathcal X_{\rm TPK}^{[108]}.
\]
The visible reading is
\[
 \pi_{\rm term}(x_{\rm term})=B_{\rm term}.
\]
The signed chart preserves another coordinate of the same state.

\begin{definition}[Signed terminal subspace]
\label{p12-83-c10-s1:def:subespacio-firmado-terminal}
Let \(\mathcal H_{12}\) be the integral transformation constructed in the next section. Define
\begin{equation}
 \mathcal X_{\rm TPK}^{[108],\rm sgn}
 :=\{(x_{\rm core},K,U):U=\mathcal H_{12}K\}
 \label{p12-83-c10-s1:eq:subespacio-firmado-terminal}
\end{equation}
and the publication module
\begin{equation}
 \mathcal U_{12}^{\rm sgn}
 :=\operatorname{im}
 \bigl(\mathcal H_{12}:\mathbb Z^{12}\to\mathbb Z^{12}\bigr).
 \label{p12-83-c10-s1:eq:modulo-publicacion-firmada}
\end{equation}
The signed reading is
\begin{equation}
 R_{\rm sgn}:
 \mathcal X_{\rm TPK}^{[108],\rm sgn}
 \longrightarrow\mathcal U_{12}^{\rm sgn},
 \qquad
 R_{\rm sgn}(x_{\rm core},K,U)=U.
 \label{p12-83-c10-s1:eq:lector-firmado-terminal}
\end{equation}
\end{definition}

The correct joint writing is
\begin{equation}
 x_{\rm term}
 \xrightarrow{(\pi_{\rm term},R_{\rm sgn})}
 (B_{\rm term},U_{12}^{\rm sgn}).
 \label{p12-83-c10-s1:eq:doble-publicacion-terminal}
\end{equation}
No arrow \(B_{\rm term}\to U_{12}^{\rm sgn}\) is interposed: the visible matrix has forgotten precisely the orientation, opposition, quotient, carry and genealogical-record coordinates used by the signed chart. The source of both expressions is the integral holonomic state produced by \eqref{p12-83-c10-s1:eq:cadena-upstream-registro-dodecafasico}.

The evaluation of the signed record gives
\begin{equation}
 \boxed{
 U=(2378,1406,2479,-452,998,-551,
 -668,-204,-371,-322,-28,-997).}
 \label{p12-83-c10-s1:eq:vector-u-firmado-autonomo}
\end{equation}

\subsection{Naturality with respect to depth}

For \(N'\geq N\geq108\), the truncations of the signed system are
\begin{equation}
 \tau_{N',N}^{\rm sgn}
 (x_{\rm core}^{[N']},K,U)
 =\bigl(\tau_{N',N}x_{\rm core}^{[N']},K,U\bigr).
 \label{p12-83-c10-s1:eq:truncamiento-preserva-ku-autonomo}
\end{equation}

\begin{theorem}[Square of naturality of the signed reading]
\label{p12-83-c10-s1:thm:naturalidad-lector-firmado-autonomo}
We have
\begin{equation}
 \boxed{
 R_{{\rm sgn},N}\circ\tau_{N',N}^{\rm sgn}
 =\operatorname{id}_{\mathcal U_{12}^{\rm sgn}}
  \circ R_{{\rm sgn},N'}.}
 \label{p12-83-c10-s1:eq:cuadrado-naturalidad-lector-firmado}
\end{equation}
\end{theorem}

\begin{proof}
Both sides, applied to
\((x_{\rm core}^{[N']},K,U)\), literally return \(U\). The
truncation modifies the subsequent history but preserves the already fixed
terminal coordinate. Commutativity is an equality of maps, not
an analogy between diagrams.
\end{proof}

\subsection{Orbits of step \(3\) and transform of Hadamard}

The cyclic order of the seal is distributed over the three orbits.
\begin{align}
 K^{(1)}&=(K_1,K_4,K_7,K_{10}),
 \label{p12-83-c10-s1:eq:orbita-k-1}\\
 K^{(2)}&=(K_2,K_5,K_8,K_{11}),
 \label{p12-83-c10-s1:eq:orbita-k-2}\\
 K^{(3)}&=(K_3,K_6,K_9,K_{12}).
 \label{p12-83-c10-s1:eq:orbita-k-3}
\end{align}
Let
\begin{equation}
 H_4=
 \begin{pmatrix}
 1&1&1&1\\
 1&1&-1&-1\\
 1&-1&1&-1\\
 1&-1&-1&1
 \end{pmatrix}.
 \label{p12-83-c10-s1:eq:matriz-hadamard-cuatro-autonoma}
\end{equation}
If \(P_{\rm orb}\) reorders the twelve coordinates according to
\eqref{p12-83-c10-s1:eq:orbita-k-1}--\eqref{p12-83-c10-s1:eq:orbita-k-3}, the global transformation is
\begin{equation}
 \mathcal H_{12}
 =P_{\rm orb}^{-1}
 (H_4\oplus H_4\oplus H_4)P_{\rm orb}.
 \label{p12-83-c10-s1:eq:hadamard-global-doce}
\end{equation}

The three blocks of \eqref{p12-83-c10-s1:eq:vector-u-firmado-autonomo}, read by
orbits, are
\begin{align*}
 U^{(1)}&=(2378,-452,-668,-322),\\
 U^{(2)}&=(1406,998,-204,-28),\\
 U^{(3)}&=(2479,-551,-371,-997).
\end{align*}

\begin{lemma}[Algebraic Hadamard quarter-turn]
\label{p12-83-c10-s1:lem:cuarto-inversa-hadamard}
The matrix \(H_4\) satisfies
\begin{equation}
 H_4^2=4I_4,
 \qquad
 H_4^{-1}=\frac14H_4,
 \qquad
 \det H_4=-16.
 \label{p12-83-c10-s1:eq:identidades-hadamard-cuatro}
\end{equation}
\end{lemma}

\begin{proof}
The rows of \(H_4\) are orthogonal, and each has squared norm four.
Since the matrix is symmetric, \(H_4^{\mathsf T}H_4=H_4^2=4I_4\). The
formula for the inverse and the determinant follow immediately.
\end{proof}

\begin{theorem}[Unique full reconstruction of the seal]
\label{p12-83-c10-s1:thm:reconstruccion-entera-k-hadamard}
The inverse
\begin{equation}
 K^{(r)}=\frac14H_4U^{(r)}
 \qquad(r=1,2,3)
 \label{p12-83-c10-s1:eq:inversion-bloques-hadamard}
\end{equation}
yields
\begin{align*}
 K^{(1)}&=(234,729,621,794),\\
 K^{(2)}&=(543,659,058,146),\\
 K^{(3)}&=(140,824,914,601).
\end{align*}
By undoing the orbital order one obtains
\begin{equation}
 \boxed{
 K=(234,543,140,729,659,824,621,058,914,794,146,601).}
 \label{p12-83-c10-s1:eq:sello-k-doce-autonomo}
\end{equation}
It is the unique rational preimage of \(U\), and it belongs to
\(\mathbb Z^{12}\).
\end{theorem}

\begin{proof}
By \eqref{p12-83-c10-s1:eq:identidades-hadamard-cuatro}, the rational preimage is unique. The explicit multiplications are
\begin{align*}
 H_4U^{(1)}&=(936,2916,2484,3176),\\
 H_4U^{(2)}&=(2172,2636,232,584),\\
 H_4U^{(3)}&=(560,3296,3656,2404).
\end{align*}
All coordinates are divisible by four. Dividing and reinterleaving
yields \eqref{p12-83-c10-s1:eq:sello-k-doce-autonomo}.
\end{proof}

The factor \(1/4\) is not chosen to fit an output: it is the inverse
forced by \(H_4^2=4I_4\). Nor is it the square root of \(-1\). The
internal complex structure \(A_W^2=-I\) will appear after constructing the
Paley matrix; these are two facts of different types.

\subsection{Integral structure and Smith normal form}

\begin{theorem}[Integral image of \(H_4\)]
\label{p12-83-c10-s1:thm:smith-hadamard-autonomo}
The Smith normal form is
\begin{equation}
 \operatorname{SNF}(H_4)=\operatorname{diag}(1,2,2,4).
 \label{p12-83-c10-s1:eq:snf-hadamard-autonoma}
\end{equation}
Consequently,
\begin{equation}
 \operatorname{coker}\mathcal H_{12}
 \cong(\mathbb Z/2\mathbb Z)^6
 \oplus(\mathbb Z/4\mathbb Z)^3,
 \label{p12-83-c10-s1:eq:cociente-hadamard-doce}
\end{equation}
and the image of \(\mathcal H_{12}\) has index
\[
 16^3=4096
\]
in \(\mathbb Z^{12}\). For \(u\in\mathbb Z^4\),
\begin{equation}
 u\in H_4\mathbb Z^4
 \quad\Longleftrightarrow\quad
 H_4u\equiv0\pmod4.
 \label{p12-83-c10-s1:eq:criterio-imagen-hadamard}
\end{equation}
\end{theorem}

\begin{proof}
The greatest common divisor of the entries is one. The minors of order two are
even, and one has absolute value two; those of order three are multiples of
four, and one has absolute value four; the determinant has modulus
sixteen. The determinantal divisors are \(1,2,4,16\), yielding
the invariant factors \(1,2,2,4\). The triple direct sum repeats the
factors and multiplies the indices. Finally,
\(H_4^{-1}=H_4/4\), so the preimage is integral exactly under the
congruence \eqref{p12-83-c10-s1:eq:criterio-imagen-hadamard}.
\end{proof}

\subsection{2nd coordinate system: step differences \(3\) and \(4\)}

With cyclic indices modulo twelve, we define
\begin{equation}
 (D_3K)_m=K_m-K_{m+3},
 \qquad
 (D_4K)_m=K_m-K_{m+4},
 \qquad
 Q(K)=\sum_{m=1}^{12}K_m.
 \label{p12-83-c10-s1:eq:extractor-diferencias-k}
\end{equation}
For \eqref{p12-83-c10-s1:eq:sello-k-doce-autonomo},
\begin{align*}
 D_3K={}&(-495,-116,-684,108,601,-90,\\
         &\qquad-173,-88,313,560,-397,461),\\
 D_4K={}&(-425,-281,-481,671,-255,30,\\
         &\qquad475,-543,680,251,6,-128),\\
 Q(K)={}&6263.
\end{align*}

\begin{theorem}[Completeness of the transversal extractor]
\label{p12-83-c10-s1:thm:completitud-extractor-transversal-k}
The operator
\[
 \mathcal E_{\rm dod}=(D_3,D_4,Q):
 \mathbb Z^{12}\longrightarrow\mathbb Z^{25}
\]
has rank twelve and Smith normal form
\begin{equation}
 \operatorname{SNF}(\mathcal E_{\rm dod})
 =\operatorname{diag}(\underbrace{1,\ldots,1}_{11},12).
 \label{p12-83-c10-s1:eq:snf-extractor-transversal-k}
\end{equation}
In particular, \((D_3K,D_4K,Q(K))\) determina a single \(K\).
\end{theorem}

\begin{proof}
If \(D_3v=D_4v=0\), then \(v\) is constant along steps
three and four. Since \(\gcd(3,4)=1\), those steps generate all of
\(\mathbb Z/12\mathbb Z\), so that
\[
 \ker D_3\cap\ker D_4=\langle\mathbf1\rangle.
\]
The total charge does not vanish on that line, since \(Q(\mathbf1)=12\). The rank
is twelve. The rows of \((D_3,D_4)\) are oriented incidences of a connected
Cayley graph; a spanning tree supplies eleven unit
invariant factors. Every maximal minor must contain the row \(Q\), and the
unimodular operations reduce its final contribution to twelve. There is a maximal minor of
modulus twelve; therefore, the last factor is exactly twelve.
\end{proof}

The reconstruction of \(U\) from these coordinates does not need to
multiply again by \(H_4\). For \(r=1,2,3\), let
\begin{equation}
 B_r=\sum_{j=0}^{3}(D_4K)_{r+3j}.
 \label{p12-83-c10-s1:eq:br-diferencias-k}
\end{equation}
Then
\[
 (B_1,B_2,B_3)=(972,-1073,101).
\]
Orbital sums are
\begin{align}
 A_1&=\frac{Q+2B_1+B_2}{3}=2378,
 \label{p12-83-c10-s1:eq:a1-reconstruccion-u}\\
 A_2&=A_1-B_1=1406,
 \label{p12-83-c10-s1:eq:a2-reconstruccion-u}\\
 A_3&=A_2-B_2=2479.
 \label{p12-83-c10-s1:eq:a3-reconstruccion-u}
\end{align}
If \(d_{r,j}=(D_3K)_{r+3j}\), the signed block of orbit \(r\) is
\begin{equation}
 \bigl(A_r,
 d_{r,1}-d_{r,3},
 d_{r,0}+d_{r,2},
 d_{r,0}-d_{r,2}\bigr),
 \label{p12-83-c10-s1:eq:bloque-firmado-desde-diferencias}
\end{equation}
that reproduces the three blocks \(U^{(r)}\).

\subsection{Threshold Tetrad of the Seal}

The local decimal completion satisfies
\[
 1000=729+271.
\]
The support of the seal coordinates that reach the threshold \(729\) is
\begin{equation}
 \boxed{
 B_K:=\{m:K_m\geq729\}=\{4,6,9,10\}.}
 \label{p12-83-c10-s1:eq:tetrada-umbral-k}
\end{equation}
This tetrad is obtained after constructing \(K\). In this unit it is not yet
called a Witt star: that object will be defined only after
constructing the ternary code, its hexads, and the Steiner system.

\subsection{Independence of the genealogical construction with respect to external metrological or physical data}

\begin{theorem}[Isolation of \(U\) and \(K\)]
\label{p12-83-c10-s1:thm:aislamiento-u-k-alpha-codata}
The composition
\begin{equation}
 \Omega_{\rm seed}\longrightarrow x_{\rm term}
 \xrightarrow{R_{\rm sgn}}U
 \xrightarrow{\mathcal H_{12}^{-1}}K
 \label{p12-83-c10-s1:eq:composicion-upstream-u-k}
\end{equation}
does not receive \(\alpha\), \(\alpha^{-1}\), CODATA or a target decimal
string as input. The seal \(K\) is determined before defining the
recurrence that will produce \(\alpha\).
\end{theorem}

\begin{proof}
The first part of the composition uses only seeds, emissions, ternary readers, operators \(L_0,L_1\), survival relations, phase transport, and the charts of the TPK state. The reading \(R_{\rm sgn}\) is a projection of the terminal register, and the final arrow is the linear transformation \(\mathcal H_{12}^{-1}=\mathcal H_{12}/4\) on its integral image. None of the formulas contains a physical variable or a metrological table. The base-thousand recurrence is introduced only in the following unit, after \(K\) has been fixed.
\end{proof}

\subsection{Obstructions of the dodecaphase record and the seal \(K\): naturality, Hadamard integrity, and metrological isolation}

\begin{criterion}[Refutation of the register and the seal]
The construction is refuted if any of the following facts
is established:
\begin{enumerate}
 \item the windows \(E_m\) do not partition the one hundred eight steps;
 \item the chart \(\mathcal L_{12}\) fails its inverse or its congruences;
 \item the terminal census does not produce \(331\) triples, two matrices and a unique
       axial charge;
 \item an attempt is made to obtain \(U\) from \(B_{\rm term}\) after forgetting the
       signed coordinates;
 \item the naturality square
       \eqref{p12-83-c10-s1:eq:cuadrado-naturalidad-lector-firmado} fails;
 \item \(H_4^2\neq4I_4\), some preimage is not integral or the
       reconstructed vector differs from \eqref{p12-83-c10-s1:eq:sello-k-doce-autonomo};
 \item the Smith normal form differs from
       \eqref{p12-83-c10-s1:eq:snf-hadamard-autonoma};
 \item extractor \((D_3,D_4,Q)\) loses rank or fails to reproduce \(U\);
 \item \(B_K\) is introduced before constructing \(K\), or is assigned a
       Witt structure before defining \(W_{12}\);
 \item a digit of \(\alpha\) or a metrological reference intervenes in
       any of the arrows of \eqref{p12-83-c10-s1:eq:composicion-upstream-u-k}.
\end{enumerate}
\end{criterion}


% Unidad U017. Cierre afín, prolongación y semántica de frontera de alpha.
% Redacción autónoma desde la autoridad material 1063 y su sucesor V2.

\section{Initial construction of the boundary \texorpdfstring{\(662/663\)}{662/663} through the dodecaphase carry cocycle}
\label{p12-83-c10-s2:chap:acarreo-alpha-frontera-662-663}

\subsection{Aritmetic closure domain}

The construction receives four dodecaphasic coordinates obtained in the preceding units:
\[
\begin{aligned}
 P={}&(141,592,653,589,793,238,462,643,383,279,502,884),\\
 E={}&(718,281,828,459,045,235,360,287,471,352,662,497),\\
 \Phi={}&(618,033,988,749,894,848,204,586,834,365,638,117),\\
 K={}&(234,543,140,729,659,824,621,058,914,794,146,601).
\end{aligned}
\]
The first three vectors are the first twelve triads of the fractional Archimedean readings of \(\pi\), \(e\) and \(\varphi\). The fourth is the dodecaphasic coordinate of the terminal holonomic state of the Topological Phase Kernel, constructed before opening the carry operation.

We fix the base.
\[
 B:=10^3=1000
\]
and the concatenation
\begin{equation}
 \iota(v_1,\ldots,v_{12})
 :=\sum_{m=1}^{12}v_mB^{12-m}.
 \label{p12-83-c10-s2:eq:concatenacion-base-mil-alpha}
\end{equation}
\Needspace{8\baselineskip}
In particular,
\[
\begin{aligned}
 \iota(P)&=141592653589793238462643383279502884,\\
 \iota(E)&=718281828459045235360287471352662497,\\
 \iota(\Phi)&=618033988749894848204586834365638117,\\
 \iota(K)&=234543140729659824621058914794146601.
\end{aligned}
\]

The operation studied in this unit is the map
\begin{equation}
 \mathfrak C_{12}:
 (P,E,\Phi,K,c_{13})
 \longmapsto(A,C),
 \label{p12-83-c10-s2:eq:mapa-cierre-dodecafasico-alpha}
\end{equation}
where \(A=(A_1,\ldots,A_{12})\) is the output word and
\[
 C=(c_1,\ldots,c_{13})
\]
is the complete system of thirteen Euclidean quotients: twelve internal
quotients and one right boundary condition.

\subsection{Left-to-right oriented Euclidean division}

For \(m=12,11,\ldots,1\), define
\begin{equation}
\begin{aligned}
 s_m&:=P_m+E_m-\Phi_m-K_m+c_{m+1},\\
 A_m&:=s_m\bmod B,
       \qquad 0\leq A_m<B,\\
 c_m&:=\frac{s_m-A_m}{B}
      =\left\lfloor\frac{s_m}{B}\right\rfloor .
\end{aligned}
\label{p12-83-c10-s2:eq:recurrencia-corta-alpha}
\end{equation}
The second equality for \(c_m\) is also valid when \(s_m<0\),
because Euclidean division with remainder in
\(\{0,\ldots,B-1\}\) is used. Therefore,
\[
 -431=569-1000,
 \qquad
 -717=283-1000,
 \qquad
 -1200=800-2\cdot1000,
\]
and negative loans do not introduce ambiguity.

\begin{theorem}[Existence and Uniqueness of the Oriented Carry]
\label{p12-83-c10-s2:thm:unicidad-acarreo-orientado-alpha}
Given \(P,E,\Phi,K\) and the terminal carry \(c_{13}\), the recurrence
\eqref{p12-83-c10-s2:eq:recurrencia-corta-alpha} uniquely determines
\[
 (A_1,\ldots,A_{12};c_1,\ldots,c_{12}).
\]
\end{theorem}

\begin{proof}
Row \(m=12\) is determined by the Euclidean division of
\[
 P_{12}+E_{12}-\Phi_{12}-K_{12}+c_{13}
\]
by \(B\). Once \(c_{12}\) is known, row \(m=11\) is determined by
the same property. Descending induction reaches \(m=1\). At each
step, uniqueness of the Euclidean quotient and remainder prevents two distinct
outputs.
\end{proof}

\subsection{Dodecaphasic window of the signed state and reversible closure of the cycle}

The finite window is defined by the terminal condition
\[
 c_{13}^{\rm fin}=0.
\]
The complete evaluation is as follows. The table is printed in ascending
order, although it must be read causally from the last row to the
first.

\begingroup
\scriptsize
\setlength{\tabcolsep}{3.2pt}
\begin{longtable}{@{}r|rrrr|r|r|r|r@{}}
\caption{The twelve Euclidean divisions and thirteen carries of the finite
dodecaphasic window.}
\label{p12-83-c10-s2:tab:acarreo-alpha-finito-autonomo}\\
\toprule
\(m\)&\(P_m\)&\(E_m\)&\(\Phi_m\)&\(K_m\)&
\(c_{m+1}\)&\(s_m\)&\(A_m\)&\(c_m\)\\
\midrule
\endfirsthead
\toprule
\(m\)&\(P_m\)&\(E_m\)&\(\Phi_m\)&\(K_m\)&
\(c_{m+1}\)&\(s_m\)&\(A_m\)&\(c_m\)\\
\midrule
\endhead
1  &141&718&618&234& \(0\)& \(7\)&007& \(0\)\\
2  &592&281&033&543& \(0\)& \(297\)&297& \(0\)\\
3  &653&828&988&140& \(-1\)& \(352\)&352& \(0\)\\
4  &589&459&749&729& \(-1\)& \(-431\)&569& \(-1\)\\
5  &793&045&894&659& \(-2\)& \(-717\)&283& \(-1\)\\
6  &238&235&848&824& \(-1\)& \(-1200\)&800& \(-2\)\\
7  &462&360&204&621& \(0\)& \(-3\)&997& \(-1\)\\
8  &643&287&586&058& \(-1\)& \(285\)&285& \(0\)\\
9  &383&471&834&914& \(-1\)& \(-895\)&105& \(-1\)\\
10 &279&352&365&794& \(0\)& \(-528\)&472& \(-1\)\\
11 &502&662&638&146& \(0\)& \(380\)&380& \(0\)\\
12 &884&497&117&601& \(0\)& \(663\)&663& \(0\)\\
\bottomrule
\end{longtable}
\endgroup

The complete word of thirteen carries is
\[
 \boxed{
 C^{\rm fin}
 =(c_1,\ldots,c_{13})
 =(0,0,0,-1,-1,-2,-1,0,-1,-1,0,0,0).}
\]
The dodecaphasic output is
\[
 \boxed{
 A^{\rm fin}
 =(007,297,352,569,283,800,997,285,105,472,380,663).}
\]
We define the finite rational
\begin{equation}
 \boxed{
 \alpha_{12}^{\rm fin}
 :=10^{-36}\iota(A^{\rm fin})
 =\frac{7297352569283800997285105472380663}{10^{36}}.}
 \label{p12-83-c10-s2:eq:alpha-finita-autonoma}
\end{equation}
Therefore,
\[
 \alpha_{12}^{\rm fin}
 =0.007297352569283800997285105472380663.
\]

\subsection{Posterior metrological recognition: CODATA 2018}

The metrological comparison is performed only after closing the recurrence,
the right boundary, and the preceding rational. The integer identity gives
\begin{equation}
 \boxed{
 \alpha_{12}^{\rm fin}
 =10^{-36}
 \left\lfloor\frac{10^{36}}{137.035999084}\right\rfloor.}
 \label{p12-83-c10-s2:eq:alpha-finita-codata-2018}
\end{equation}
Consequently,
\begin{equation}
 (\alpha_{12}^{\rm fin})^{-1}
 =137.03599908400000000000000000000001066588\ldots,
 \label{p12-83-c10-s2:eq:alpha-inversa-codata-2018}
\end{equation}
and exactly reproduces the central chain published by CODATA 2018
\cite{Tiesinga2021CODATA}:
\[
 \boxed{137.035999084.}
\]
The word \emph{exactly} refers to the published central chain, not to a purported experimental measurement of thirty-six digits. The subsequent digits in \eqref{p12-83-c10-s2:eq:alpha-inversa-codata-2018} are consequences of the HMT rational already fixed. CODATA does not select the state, does not determine (K), does not decide the direction of the carry, and does not intervene in the dodecaphase table.

\subsection{Complete affine identity}

Each row of the table satisfies exactly
\begin{equation}
 P_m+E_m-\Phi_m-K_m+c_{m+1}
 =A_m+Bc_m.
 \label{p12-83-c10-s2:eq:identidad-local-acarreo-alpha}
\end{equation}
Introducing
\[
 C^-:=(c_1,\ldots,c_{12}),
 \qquad
 C^+:=(c_2,\ldots,c_{13}),
\]
the twelve local identities come together in
\begin{equation}
 \boxed{
 P+E-\Phi-K+C^+-BC^-=A.}
 \label{p12-83-c10-s2:eq:identidad-afin-vectorial-alpha}
\end{equation}
The term
\[
 \partial_BC:=C^+-BC^-
\]
is the carry cocycle of the oriented chain. Consequently,
\(P+E-\Phi-K=A\) is not a coordinatewise equality in
\(\mathbb Z^{12}\): it becomes correct only upon incorporating
\(\partial_BC\).

\begin{theorem}[Telescoping of the Thirteen Carries]
\label{p12-83-c10-s2:thm:telescopado-acarreos-alpha}
For any terminal condition \(c_{13}\), the recurrence satisfies
\begin{equation}
 \boxed{
 \iota(P)+\iota(E)-\iota(\Phi)-\iota(K)-\iota(A)
 =B^{12}c_1-c_{13}.}
 \label{p12-83-c10-s2:eq:telescopado-general-alpha}
\end{equation}
In the finite window, \(c_1=c_{13}=0\), and therefore
\begin{equation}
 \boxed{
 \iota(P)+\iota(E)-\iota(\Phi)-\iota(K)
 =\iota(A^{\rm fin}).}
 \label{p12-83-c10-s2:eq:identidad-entera-alpha-finita}
\end{equation}
\end{theorem}

\begin{proof}
We multiply \eqref{p12-83-c10-s2:eq:identidad-local-acarreo-alpha} by
\(B^{12-m}\) and sum over \(m=1,\ldots,12\). The contribution of the carries is
\[
 \sum_{m=1}^{12}
 (Bc_m-c_{m+1})B^{12-m}.
\]
The internal components cancel:
\begin{align*}
 \sum_{m=1}^{12}(Bc_m-c_{m+1})B^{12-m}
 &=\sum_{m=1}^{12}c_mB^{13-m}
   -\sum_{m=1}^{12}c_{m+1}B^{12-m}\\
 &=B^{12}c_1-c_{13}.
\end{align*}
This proves \eqref{p12-83-c10-s2:eq:telescopado-general-alpha}. The specialization
\(c_1=c_{13}=0\) gives \eqref{p12-83-c10-s2:eq:identidad-entera-alpha-finita}.
\end{proof}

\subsection{Exact inversion of the seal \texorpdfstring{\(K\)}{K}}

The local recurrence is reversible. Solving for \(K_m\) in
\eqref{p12-83-c10-s2:eq:identidad-local-acarreo-alpha},
\begin{equation}
 \boxed{
 K_m=P_m+E_m-\Phi_m+c_{m+1}-A_m-Bc_m.}
 \label{p12-83-c10-s2:eq:inversion-local-k-desde-alpha}
\end{equation}
By concatenation,
\begin{equation}
 \boxed{
 \iota(K)
 =\iota(P)+\iota(E)-\iota(\Phi)-\iota(A)
  -B^{12}c_1+c_{13}.}
 \label{p12-83-c10-s2:eq:inversion-concatenada-k-desde-alpha}
\end{equation}
On the finite boundary,
\[
 \iota(K)
 =\iota(P)+\iota(E)-\iota(\Phi)-\iota(A^{\rm fin}).
\]

Algebraic reversibility does not invert the causal order. The effective order
is
\[
 x_{\rm term}
 \longrightarrow K
 \longrightarrow(P,E,\Phi,K,c_{13})
 \longrightarrow(A,C).
\]
Formula \eqref{p12-83-c10-s2:eq:inversion-local-k-desde-alpha} is a commutativity
check: it proves that the output and the carries recover the seal
that was already fixed. It is not a rule for retrospectively selecting
\(K\) from a desired value of \(\alpha\).

\subsection{Coinductive prolongation of the dodecaphasic seal and inverse carry recurrence in \(\mathcal C=\{-2,-1,0,1,2\}\)}

Let
\[
 (P_k)_{k\geq1},
 \qquad
 (E_k)_{k\geq1},
 \qquad
 (\Phi_k)_{k\geq1}
\]
the infinite sequences of triads produced by the three coinductive branches, and let
\[
 \kappa(k):=1+((k-1)\bmod12).
\]
The prolongation of the stamp is periodic:
\[
 K_k^{\rm per}:=K_{\kappa(k)}.
\]
For a finite depth \(N\) and a terminal condition
\(t=c_{N+1}\), the recurrence is
\begin{equation}
\begin{aligned}
 s_k^{(N,t)}
 &=P_k+E_k-\Phi_k-K_{\kappa(k)}+c_{k+1}^{(N,t)},\\
 A_k^{(N,t)}
 &=s_k^{(N,t)}\bmod1000,\\
 c_k^{(N,t)}
 &=\left\lfloor\frac{s_k^{(N,t)}}{1000}\right\rfloor,
 \qquad k=N,\ldots,1.
\end{aligned}
\label{p12-83-c10-s2:eq:recurrencia-larga-alpha}
\end{equation}

The set of terminal carries
\[
 \mathcal C:=\{-2,-1,0,1,2\}
\]
is invariant for the certified inputs: for each triad and each
\(c_{k+1}\in\mathcal C\), the quotient \(c_k\) again belongs to
\(\mathcal C\). No favorable terminal condition is chosen. All
five conditions are propagated, and only the prefix common to all of them is published.

The execution of certified depth employs
\[
\begin{array}{c|r}
\text{quantity}&\text{value}\\ \hline
\text{ternary blocks per channel}&3501\\
\text{trits per channel}&3501\cdot6=21006\\
\text{guard decimal digits}&10020\\
\text{guard triads}&3340\\
\text{propagated terminal conditions}&5\\
\text{common triads}&3339\\
\text{common digits}&10017\\
\text{published digits}&10000
\end{array}
\]
The five terminal branches coalesce before the published prefix. In
particular, no terminal condition extracted from a target value of
\(\alpha\) intervenes in the output.

\begin{theorem}[Stabilized Output Compatibility]
\label{p12-83-c10-s2:thm:compatibilidad-prolongacion-alpha}
Let \(A^{[N]}\) and \(A^{[N']}\), with \(N'>N\), be two runs whose
five terminal conditions have coalesced before position \(N\).
Then the truncation of \(A^{[N']}\) to its first \(N\) blocks
coincides with \(A^{[N]}\). The stabilized prefixes define a unique
section in the inverse limit of finite words.
\end{theorem}

\begin{proof}
The recurrence is deterministic from right to left once the carry entering a
position is known. Coalescence asserts that all remote
conditions produce the same carry before the boundary of the
prefix. From that point, the same triads \(P_k,E_k,\Phi_k\) and the same
periodic seal produce, by Euclidean uniqueness, the same remainders and
quotients at all preceding positions. This is exactly
commutativity with truncation.
\end{proof}

\subsection{Boundary originating from the remote tail}

The prolongation determines on the visible boundary.
\[
 c_{13}^{\infty}=-1.
\]
The twelve internal carries remain the same as in the finite window:
\[
 \boxed{
 C_{\leq13}^{\infty}
 =(0,0,0,-1,-1,-2,-1,0,-1,-1,0,0,-1).}
\]
Only the information entering from the thirteenth triad changes. The
two relevant divisions are
\[
\begin{aligned}
 s_{12}^{\infty}
 &=884+497-117-601-1=662,\\
 A_{12}^{\infty}&=662,
 \qquad c_{12}^{\infty}=0,
\end{aligned}
\]
and
\[
\begin{aligned}
 s_{13}^{\infty}
 &=197+757-720-234-1=-1,\\
 A_{13}^{\infty}&=999,
 \qquad c_{13}^{\infty}=-1.
\end{aligned}
\]
Therefore, the beginning of the prolonged word is
\[
 \boxed{
 A^{\infty}
 =(007,297,352,569,283,800,997,285,105,472,380,662,999,\ldots).}
\]
The corresponding expansion begins
\begin{equation}
 \boxed{
 \alpha_{\infty}
 =0.007297352569283800997285105472380662999564172539642\ldots.}
 \label{p12-83-c10-s2:eq:alpha-infinita-autonoma}
\end{equation}

The difference between both boundaries is local and exact:
\[
 c_{13}^{\rm fin}-c_{13}^{\infty}=1,
 \qquad
 A_{12}^{\rm fin}-A_{12}^{\infty}=1,
\]
while
\[
 A_m^{\rm fin}=A_m^\infty,
 \qquad 1\leq m\leq11.
\]

Applying the telescoping theorem to the first twelve blocks of the
prolongation,
\[
 \iota(P)+\iota(E)-\iota(\Phi)-\iota(K)
 -\iota(A^\infty_{\leq12})
 =-c_{13}^{\infty}=1.
\]
Then
\begin{equation}
 \boxed{
 \iota(A^{\rm fin})
 =\iota(A^\infty_{\leq12})+1.}
 \label{p12-83-c10-s2:eq:relacion-unidad-fronteras-alpha}
\end{equation}
The ending \(663\) does not contradict the ending \(662\): the former
is the finite representative obtained by closing the window with
\(c_{13}=0\); the latter retains the carry \(-1\) from the remote
tail.

\begin{figure}[htbp]
\centering
\begin{tikzpicture}[
  >=Stealth,
  node distance=12mm and 16mm,
  box/.style={draw,rounded corners=1.5mm,align=center,
    inner xsep=3.2mm,inner ysep=2.5mm},
  arr/.style={->,thick}
]
\node[box] (shared) {shared data\\
 \(P,E,\Phi,K\)};
\node[box,below left=of shared] (finite)
 {finite boundary\\\(c_{13}=0\)};
\node[box,below right=of shared] (remote)
 {remote tail\\\(c_{13}=-1\)};
\node[box,below=of finite] (a663)
 {\(A_{12}=663\)\\clausura finita};
\node[box,below=of remote] (a662)
 {\(A_{12}=662\), \(A_{13}=999\)\\prolonged prefix};
\draw[arr] (shared)--(finite);
\draw[arr] (shared)--(remote);
\draw[arr] (finite)--(a663);
\draw[arr] (remote)--(a662);
\draw[<->,densely dashed] (a663)--node[below,align=center]
 {exact integer difference \(1\)} (a662);
\end{tikzpicture}
\caption{Boundary bifurcation of a single recurrence. Two distinct generators
are not being compared: only the datum entering from the
right changes.}
\label{p12-83-c10-s2:fig:bifurcacion-frontera-alpha}
\end{figure}

\subsection{Truncation and rounding to \(36\) digits}

For \(x\geq0\), define
\[
 \operatorname{Tr}_{36}(x)
 :=10^{-36}\left\lfloor10^{36}x\right\rfloor
\]
and
\[
 \operatorname{Rd}_{36}(x)
 :=10^{-36}\left\lfloor10^{36}x+\frac12\right\rfloor.
\]
Let
\[
 a:=7297352569283800997285105472380662.
\]
The expansion \eqref{p12-83-c10-s2:eq:alpha-infinita-autonoma} can be written as
\[
 \alpha_\infty=\frac{a+\rho}{10^{36}},
 \qquad
 \rho=0.999564172539642709\ldots,
\]
so that \(0.999<\rho<1\).

\begin{theorem}[Exact resolution of the boundary \(662/663\)]
\label{p12-83-c10-s2:thm:frontera-alpha-662-663}
The coinductive prolongation satisfies
\[
 \boxed{
 \operatorname{Tr}_{36}(\alpha_\infty)
 =0.007297352569283800997285105472380662,}
\]
while its rounding to thirty-six digits satisfies
\[
 \boxed{
 \operatorname{Rd}_{36}(\alpha_\infty)
 =0.007297352569283800997285105472380663
 =\alpha_{12}^{\rm fin}.}
\]
\end{theorem}

\begin{proof}
Since \(0\leq\rho<1\),
\[
 \left\lfloor10^{36}\alpha_\infty\right\rfloor=a.
\]
This gives the truncation ending in \(662\). Since \(\rho>1/2\),
\[
 \left\lfloor10^{36}\alpha_\infty+\frac12\right\rfloor=a+1,
\]
which ends in \(663\). By \eqref{p12-83-c10-s2:eq:alpha-finita-autonoma},
\((a+1)10^{-36}\) is exactly \(\alpha_{12}^{\rm fin}\).
\end{proof}

Moreover,
\[
 0<\alpha_{12}^{\rm fin}-\alpha_\infty
 =\frac{1-\rho}{10^{36}}<10^{-39}.
\]
The difference is not an arithmetic error in the table: it quantifies the
difference between two distinct boundary operators.

\subsection{Causal isolation from \texorpdfstring{\(\alpha\)}{alpha} and metrological data}

At depth \(N\), the full domain of the operation is
\[
 \mathscr D_N
 =\bigl((P_k)_{k=1}^{N},(E_k)_{k=1}^{N},
        (\Phi_k)_{k=1}^{N},K,c_{N+1}\bigr).
\]
The carry map
\[
 \mathfrak C_N:\mathscr D_N
 \longrightarrow
 \{0,\ldots,999\}^{N}\times\mathcal C^{N}
\]
is defined exclusively by \eqref{p12-83-c10-s2:eq:recurrencia-larga-alpha}. Its
domain contains no target value of \(\alpha\), its inverse, a metrological
table, a target decimal string or a terminal carry chosen by
comparison.

\begin{theorem}[No interference from the metrological target]
\label{p12-83-c10-s2:thm:no-interferencia-metrologica-alpha}
Let \(\mathscr Y\) be any space of external data containing metrological
values or comparison tables. The formal extension
\[
 \widetilde{\mathfrak C}_N:
 \mathscr D_N\times\mathscr Y
 \longrightarrow
 \{0,\ldots,999\}^{N}\times\mathcal C^{N},
 \qquad
 \widetilde{\mathfrak C}_N(d,y):=\mathfrak C_N(d),
\]
factorizes by the projection
\[
 \operatorname{pr}_{\mathscr D_N}:
 \mathscr D_N\times\mathscr Y\longrightarrow\mathscr D_N.
\]
Therefore, for any \(y,y'\in\mathscr Y\),
\[
 \widetilde{\mathfrak C}_N(d,y)
 =\widetilde{\mathfrak C}_N(d,y').
\]
No intervention on a value of \(\alpha\) or a metrological
table alters \(K\), the quotients, the remainders or the prefix produced.
\end{theorem}

\begin{proof}
The statement results from
\[
 \widetilde{\mathfrak C}_N
 =\mathfrak C_N\circ\operatorname{pr}_{\mathscr D_N}.
\]
All operations of \(\mathfrak C_N\) are additions, subtractions and Euclidean
divisions of the elements of \(\mathscr D_N\). No coordinate of
\(\mathscr Y\) appears in them.
\end{proof}

The complete causal order is
\[
 \text{enriched branches of }\pi,e,\varphi
 \longrightarrow(P,E,\Phi),
\]
\[
 x_{\rm term}\longrightarrow K,
\]
\[
 (P,E,\Phi,K,\text{boundary})
 \longrightarrow(A,C)
 \longrightarrow\alpha_\infty,
\]
\[
 \alpha_\infty
 \longrightarrow\text{external comparison}.
\]
The last arrow cannot be reversed to select any of the previous ones.

\subsection{Obstructions of the affine prolongation of \(\alpha\): telescoping, inversion of \(K\), boundary \(662/663\) and isolation of target data}

\begin{criterion}[Refutation of the affine closure and its prolongation]
The block is refuted in its domain if any of the following facts
occurs:
\begin{enumerate}
 \item some row of the \cref{p12-83-c10-s2:tab:acarreo-alpha-finito-autonomo}
       fails to satisfy \(s_m=A_m+1000c_m\);
 \item the published vector of thirteen carries does not reproduce the twelve rows;
 \item the telescoping identity fails
       \[
       \iota(P)+\iota(E)-\iota(\Phi)-\iota(K)-\iota(A)
       =1000^{12}c_1-c_{13};
       \]
 \item the inversion \eqref{p12-83-c10-s2:eq:inversion-local-k-desde-alpha} does not recover
       the twelve coordinates of \(K\);
 \item remote propagation produces \(c_{13}\neq-1\) for the certified
       prefix;
 \item the five terminal conditions of
       \(\mathcal C=\{-2,-1,0,1,2\}\) do not coalesce before the published
       digits;
 \item the triad following \(662\) is not \(999\), in which case the
       identification of the rounding with \(663\) is no longer proved;
 \item the generator opens a target expansion of \(\alpha\), its inverse,
       a metrological table, or a terminal condition chosen to force
       the coincidence.
\end{enumerate}
\end{criterion}

Thus, two distinct boundary identities are thereby demonstrated:
\[
 \boxed{
 663=\text{finite closure with }c_{13}=0
     =\text{rounding to 36 digits of }\alpha_\infty,}
\]
\[
 \boxed{
 662=\text{truncation to 36 digits of }\alpha_\infty
 \quad\text{with}\quad c_{13}=-1,\ A_{13}=999.}
\]
The central metrological result is finally published in the same residence as the derivation:
\begin{equation}
 \boxed{
 \begin{gathered}
  (\alpha_{12}^{\rm fin})^{-1}
  =137.03599908400000000000000000000001066588\ldots,\\
  \text{central CODATA 2018 string:}\qquad137.035999084.
 \end{gathered}}
 \label{p12-83-c10-s2:eq:resultado-final-alpha-codata-2018}
\end{equation}
The implication expresses a subsequent comparison with the central string published; it does not introduce that string into the constructor.
 % Insercion editor-ready para el capitulo 31.

\section{Graded nonadic prolongation of the kernel \texorpdfstring{\(E_{21/22}\)}{E21/22}: operator \texorpdfstring{\(T_{7:9}\)}{T7:9}, genealogy \texorpdfstring{\(120/45/11/495\)}{120/45/11/495}, and residue of order \texorpdfstring{\(9\)}{9}}
\label{sec:l9s102:alpha-dos-vias}

This section develops two internal constructions relating to the coordinate \(\alpha\) of the
\hyperref[sec:tesis-inaugural-helicoidal-hmt-md]{tesis inaugural}. Route A operates through reversible dodecaphasic closure of the signed state and expresses \(\alpha_{12}\). Route B, once the analytic kernel \(E_{21/22}\) and its coefficients have been constructed internally, expresses in its analytic chart a simple unique positive root, denoted \(\alpha_E:=x_{21/22}\). Both outputs share a genealogy; their proved agreement belongs to the typed decimal codomain specified below. \(\alpha_E\) is not identified with \(\alpha_{12}\) by definition. Their domains, operations and certificates remain distinct:
\begin{equation}
 X_{12}^{\mathrm{sig}}
 \xrightarrow{\ \Gamma_{12}^{\alpha}\ }
 \alpha_{12},
 \qquad
 E_{21/22}:I_\alpha\longrightarrow\mathbb R,
 \quad E_{21/22}(x_{21/22})=0,
 \quad E_{21/22}'(x_{21/22})\ne0.
 \label{eq:l9s102:alpha-dos-mapas}
\end{equation}
Integer closure expresses \(\alpha_{12}\) through carry, whereas the analytic equation determines \(\alpha_E=x_{21/22}\) through existence and uniqueness of the root in the certified interval. The currently proved agreement of both outputs resides exactly in the codomain of the morphism
\begin{equation}
 \operatorname{Tr}_9(y):=10^{-9}\lfloor10^9y\rfloor,
 \qquad
 \operatorname{Tr}_9(x_{21/22}^{-1})
 =137.035999084
 =\operatorname{Tr}_9(\alpha_{12}^{-1}).
 \label{eq:l9s102:concordancia-tipica-alpha}
\end{equation}
The exact analytic identity corresponds to \(E_{21/22}(\alpha_E)=0\), whereas \eqref{eq:l9s102:concordancia-tipica-alpha} is the exact equality of the two expressions after applying \(\operatorname{Tr}_9\). The stronger promotion \(\alpha_E=\alpha_{12}\) in \(\mathbb R_{\rm HMT}\) requires proving \(E_{21/22}(\alpha_{12})=0\), or an injective morphism lifting the truncated equality; the nonzero residue of \cref{tab:l9s102:alpha-residuos} does not supply that identity. Neither route introduces the target value as an input; metrological identification belongs solely to subsequent conventional recognition.

\subsection{Nucleo analytic \texorpdfstring{\(P_6\)}{P6} and root unique of the equation \texorpdfstring{\(E_{21/22}\)}{E21/22}}
\label{sec:l9s102:alpha-p6}

The characterization of vacancies employs
\begin{equation}
 \begin{aligned}
 P_6(x)={}&2\pi x-\frac74x^2
 +\frac{x^3}{2\pi}
 +\frac{x^4}{20}-\frac{2x^5}{21}-\frac{x^6}{46},\\
 P_6(x)={}&\log_{10}\frac{10}{9}.
 \end{aligned}
 \label{eq:l9s102:p6}
\end{equation}
The certified positive root of this equation satisfies
\begin{equation}
 x_6^{-1}=137.03599908400002702009\ldots .
 \label{eq:l9s102:raíz-p6}
\end{equation}
The polynomial \(P_6\) provides the initial jet of the analytic
characterization; the nonadic holonomy now prolongs its degrees \(7,8,9\).

\subsection{APP--TRIT--TPK genealogy of the invariants \texorpdfstring{\(120,45,11\)}{120, 45, 11} and \texorpdfstring{\(495\)}{495}}
\label{sec:l9s102:alpha-genealogia-7-9}

The phase semicorona of the TPK constructs
\begin{equation}
 |\mathcal F_8^{\mathrm{ph}}|=120,
 \qquad |L_{\mathrm{ph}}|=45.
 \label{eq:l9s102:alpha-120-45}
\end{equation}
APP provides the central block
\begin{equation}
 B_c=\begin{pmatrix}7&2\\2&7\end{pmatrix},
 \qquad
 \det B_c=45,
 \qquad
 \operatorname{Spec}(B_c)=\{9,5\}.
 \label{eq:l9s102:alpha-bc}
\end{equation}
The dodecaphase transverse extractor provides
\begin{equation}
 r_{\mathrm{dod}}=\operatorname{rank}(D_3,D_4)=11,
 \qquad
 (D_3K)_1=-495,
 \qquad
 495=45\cdot11=\det(B_c)\,r_{\mathrm{dod}}.
 \label{eq:l9s102:alpha-495}
\end{equation}
The coefficients of the prolonged jet therefore proceed from three
prior invariants of APP--TRIT--TPK:
\begin{equation}
 -\frac1{120}\longmapsto\frac1{45}
 \longmapsto\frac2{45\cdot11}.
 \label{eq:l9s102:alpha-coeficientes}
\end{equation}

\subsection{Nilpotent resolvent \texorpdfstring{\((I-N)^{-1}\)}{(I-N)^-1} and exact construction of the tail \texorpdfstring{\(T_{7:9}\)}{T7:9}}
\label{sec:l9s102:alpha-resolvente}

On the basis \(e_7,e_8,e_9\), we define
\begin{equation}
 Ne_7=-\frac83e_8,
 \qquad
 Ne_8=\frac2{11}e_9,
 \qquad
 Ne_9=0,
 \qquad
 v_7=-\frac1{120}e_7.
 \label{eq:l9s102:alpha-nilpotente}
\end{equation}
Since \(N^3=0\), its finite resolvent acts exactly by
\begin{equation}
 (I-N)^{-1}v_7
 =v_7+Nv_7+N^2v_7
 =-\frac{e_7}{120}+\frac{e_8}{45}+\frac{2e_9}{495}.
 \label{eq:l9s102:alpha-resolvente-finito}
\end{equation}
The evaluation \(e_n\mapsto x^n\) produces
\begin{equation}
 \boxed{
 T_{7:9}(x)=-\frac{x^7}{120}+\frac{x^8}{45}
 +\frac{2x^9}{495}.}
 \label{eq:l9s102:alpha-t79}
\end{equation}

\subsection{Nonadic jet \texorpdfstring{\(P_9\)}{P9}, prolonged root, and scale of certified residues}
\label{sec:l9s102:alpha-jet-p9}

The nonadic jet of order \(9\) is defined by
\begin{equation}
 P_9(x)=P_6(x)+T_{7:9}(x),
 \qquad
 P_9(x)=\log_{10}\frac{10}{9}.
 \label{eq:l9s102:alpha-p9}
\end{equation}
Its positive root satisfies
\begin{equation}
 x_9^{-1}
 =137.0359990840000000000000111926291900\ldots .
 \label{eq:l9s102:alpha-raíz-p9}
\end{equation}

\begin{table}[htbp]
\centering
\caption{Certified descent of the residue of the analytical characterization}
\label{tab:l9s102:alpha-residuos}
\begin{tabular}{@{}lc@{}}
\toprule
Nucleo evaluado in \(\alpha_{12}\) & Residue \\
\midrule
\(P_6\) & \(9.0038886\times10^{-18}\) \\
\(P_6-x^7/120\) & \(-1.7893115\times10^{-19}\) \\
\(P_6-x^7/120+x^8/45\) & \(-2.3708601\times10^{-22}\) \\
\(P_9\) & \(3.7297132\times10^{-27}\) \\
\bottomrule
\end{tabular}
\end{table}

\subsection{Operador dodecafasico alternativo \texorpdfstring{\(A_{\mathrm{dod}}\)}{A_dod} and proof of the not uniqueness of the selection}
\label{sec:l9s102:alpha-control}

The control operator
\begin{equation}
 A_{\mathrm{dod}}
 =I-\frac14(D_3^*D_3+D_4^*D_4)
 \label{eq:l9s102:alpha-control-alternativo}
\end{equation}
produces a different coefficient sequence. This control delimits the selective power of the resolvent \cref{eq:l9s102:alpha-resolvente-finito,eq:l9s102:alpha-t79}: a generic dodecaphasic operator does not by itself determine the extension \(T_{7:9}\). It does not weaken the vacancy route, whose proof condition is existence and uniqueness of the root of the declared kernel. The control morphism \(\Gamma_{E21}^{(9)}\) is reserved exclusively for canonical uniqueness of the nonadic extension \(7{:}9\); it neither certifies nor conditions the status of route B.

\subsection{Scope of the nilpotent operator \texorpdfstring{\(N\)}{N} on the nonadic jet \texorpdfstring{\(P_9\)}{P9} and uniqueness control of the prolongation \texorpdfstring{\(7{:}9\)}{7:9}}
\label{sec:l9s102:alpha-alcance}

The arithmetic of \(120,45,11,495\), the nilpotent operator, the resolvent and the identity \cref{eq:l9s102:alpha-t79} are internally exact. The simple unique positive root of the kernel \(E_{21/22}\) constitutes a second exact internal construction, \(\alpha_E\), for the declared domain and coefficients. Its identification with the dodecaphasic output is proved exactly after \(\operatorname{Tr}_9\); raising this to equality in \(\mathbb R_{\rm HMT}\) requires the identity or injective lifting stated at the beginning of the section. Independently, canonical uniqueness of the extension \(7{:}9\) is typed by the morphism
\begin{equation}
 \Gamma_{E21}^{(9)}:
 X_{\mathrm{APP\text{-}TRIT\text{-}TPK}}
 \longrightarrow \mathbb Q(\pi)[x]/(x^{10}),
 \label{eq:l9s102:alpha-gamma-e21}
\end{equation}
whose action jointly determines the oriented seed, both transfers, the grading \(7,8,9\) and the jet quotient \(F^7/F^{10}\). The falsifier therefore separates two statements: it would invalidate route B if monotonicity, the sign change or root isolation failed; it would invalidate canonical promotion of \(T_{7:9}\) if another reader on the same domain produced the same grading without factoring through \(\Gamma_{E21}^{(9)}\).
 
% Selección íntegra de los cuerpos matemáticos de la Parte III rectora. Las
% cartas dimensionales se reservan para la parte siguiente, de modo que una
% constante física no pueda retroactuar como generador de estos caracteres.
% Punto de montaje redefinible por el compilador único.
% La ruta por defecto es válida cuando el módulo se consume desde
% ``manuscrito/``; el editor único puede redefinirla antes del \input.
\providecommand{\HMTParteIIIRoot}{../colaboracion/parte_iii}
\providecommand{\HMTParteIIISource}{\HMTParteIIIRoot/source/public_final}
\providecommand{\HMTParteIIICapitulos}{\HMTParteIIIRoot/tex/capitulos_finales}
\providecommand{\APP}{\mathrm{APP}}
\providecommand{\TRIT}{\mathrm{TRIT}}
\providecommand{\TPK}{\mathrm{TPK}}
\providecommand{\etaRet}{\eta_{\mathrm{ret}}^{(\mathrm{deg})}}
\providecommand{\etaRetrad}{\eta_{\mathrm{ret}}^{(\mathrm{rad})}}
\providecommand{\alphaAngle}{\alpha_{\angle}}
 
% \newcommand{\HMTMonographChapterInput}[2]{%
%   \chapter{#1}%
%   \begingroup
%   \renewcommand{\chapter}[2][]{}%
%   \input{#2}%
%   \endgroup}

\chapter{Algebraic and genealogical classification of the fundamental constants}
\begingroup
\renewcommand{\chapter}[2][]{}%
\chapter{Algebra and genealogical classification of fundamental constants}
\label{chap:p3-final:27:algebra_clasificacion}
\label{chap:p3-83-23-algebra_clasificacion}


\label{int:chap:producto-nativo}
\index{native product}

The Euclidean residue--quotient extension of APP contains two local maps: the additive map normalizes collisions, and the multiplicative map produces partial products. In this section, the global multiplication of words is constructed using exclusively both maps. The correspondence with integer multiplication will be a consequence of the intertwiner, not the definition of the operation.

\section{Euclidean coordinates of the local cell}

\begin{definition}[Additive and multiplicative applications on the same support]
Let $\Dnine=\{1,\ldots,9\}$. For $a,b\in\Dnine$ they are defined by Euclidean division
with remainder in $\Dnine$:
\begin{align}
 a+b&=9q^+(a,b)+R^+(a,b),\label{int:eq:local-plus}\\
 ab&=9q^\times(a,b)+R^\times(a,b).\label{int:eq:local-times}
\end{align}
In particular, the zero remainder is represented by $9$ and the quotient is reduced by
one unit. The quadruple
\[
 \CellRQ(a,b)=
 (R^+(a,b),q^+(a,b);R^\times(a,b),q^\times(a,b))
\]
is the local arithmetic extension by residue and quotient.
\end{definition}

Aggregation by the three classes modulo three of the visible table
\(R^\times\) defines a matrix distinct from the local maps:
\[
(M_\Pi)_{rs}
=\frac19
\sum_{\substack{a\equiv r\;(3)\\b\equiv s\;(3)}}R^\times(a,b),
\qquad r,s\in\{1,2,0\},
\]
and, in that order of classes,
\[
M_\Pi=
\begin{pmatrix}4&5&6\\5&4&6\\6&6&9\end{pmatrix}.
\]
Thus, \(q^+\) and \(q^\times\) are quotient coordinates of a cell,
whereas \(M_\Pi\) is an aggregate operator on residual classes.

The $81$ cells satisfy exactly
\[
\begin{array}{c|rrrr}
 &\sum R&\sum q&\sum(9q+R)\\\hline
 +&405&45&810\\
 \times&459&174&2025
\end{array}
\]
and therefore
\[
 2025-810=(459-405)+9(174-45)=54+9\cdot129=1215.
\]
The quotient $q^\times$ is necessary for reconstruction. For example,
\[
 2\cdot5=9\cdot1+1.
\]
If only $R^\times$ is preserved, the product $10$ is conflated with the
visible value $1$.

\subsection{Cocycles of associativity and distributivity}

The coherence of the transducer can already be decided in the cell. For every
$a,b,c\in\Dnine$ the three identities are verified
\begin{align}
q^+(a,b)+q^+(R^+(a,b),c)
&=q^+(b,c)+q^+(a,R^+(b,c)),\label{int:eq:cocycle-plus}\\
q^\times(R^\times(a,b),c)+cq^\times(a,b)
&=q^\times(a,R^\times(b,c))+aq^\times(b,c),
\label{int:eq:cocycle-times}\\
q^\times(a,R^+(b,c))+aq^+(b,c)
&=q^+(R^\times(a,b),R^\times(a,c))\notag\\
&\quad+q^\times(a,b)+q^\times(a,c).
\label{int:eq:cocycle-distributive}
\end{align}

\begin{theorem}[Local coherence of the Euclidean extension]
The identities \eqref{int:eq:cocycle-plus}--\eqref{int:eq:cocycle-distributive},
together with the corresponding equalities of residues, express
respectively additive associativity, multiplicative associativity, and
distributivity in the residue--quotient extension. The three families have
zero discrepancies in the $9^3=729$ triples.
\end{theorem}

\begin{proof}
Expanding $(a+b)+c$ through \eqref{int:eq:local-plus}, the coefficient of
nine is the left-hand side of \eqref{int:eq:cocycle-plus}; expanding
$a+(b+c)$ yields the right-hand side. Uniqueness of division with remainder in
$\Dnine$ equates both residues and quotients. The same argument applied to
$(ab)c=a(bc)$ gives \eqref{int:eq:cocycle-times}. Finally, the two expansions of
$a(b+c)=ab+ac$ produce \eqref{int:eq:cocycle-distributive}. Direct enumeration
of the \(729\) cases of each identity confirms the algebraic
calculation.
\end{proof}

\section{Bijectivity of numbering}

The symbol $9$ is not identified with zero. In $\CellRQ$, $9$ represents a
complete crossing of the nonadic boundary. Replacing it with the digits
$0,\ldots,8$ would eliminate the Euclidean quotient that intervenes in the
intertwining.

\begin{definition}[Canonical words]
Let
\[
 \Wnine^+=\bigsqcup_{\ell\geq1}\Dnine^\ell,
 \qquad \Wnine=\{\varnothing\}\sqcup\Wnine^+.
\]
Words are written from the least significant digit. We define
\[
 \operatorname{ev}_9(u_0,\ldots,u_{\ell-1})
 =\sum_{j=0}^{\ell-1}u_j9^j,
 \qquad \operatorname{ev}_9(\varnothing)=0.
\]
\end{definition}

\begin{theorem}[Uniqueness of the normal form]
The evaluation $\operatorname{ev}_9$ is a bijection from $\Wnine^+$ onto
$\N_{\geq1}$, and from $\Wnine$ onto $\N_0$.
\end{theorem}

\begin{proof}
The words of length $\ell$ cubren the intervalo
\[
 \frac{9^\ell-1}{8}
 \leq \operatorname{ev}_9(u)
 \leq 9\frac{9^\ell-1}{8}.
\]
The minimum of the next interval is one unit greater than the maximum of the
preceding one. Within each interval, division of $n-1$ by nine uniquely determines
the digit
\[
 d=(n-1)\bmod9+1
\]
and the quotient that is encoded recursively. The process terminates because the
quotient strictly decreases.
\end{proof}

Examples:
\[
 9\longleftrightarrow(9),\qquad
 10\longleftrightarrow(1,1),\qquad
 81\longleftrightarrow(9,8),\qquad
 1000\longleftrightarrow(1,3,3,1).
\]

\section{Fibers of the residual projections on the integral liftings}

For this block, we distinguish two alphabets. The normal form
\(\Wnine\) uses the digits \(\{1,\ldots,9\}\), written from the least
significant. By contrast, \(\mathscr B_9\) denotes the conventional positional
representations in base nine, with digits \(\{0,\ldots,8\}\) and the most
significant on the left. The two types are not identified.

Let
\[
 \mathscr W_{10}^{(2)}=\{0,\ldots,9\}^2,\qquad
 \rho_{10}(a,b)=(b,a),\qquad
 \operatorname{ev}_{10}(a,b)=10a+b,
\]
and let \(q_9:\mathbb N\to\mathbb Z/9\mathbb Z\) be the residual reduction.

\begin{lemma}[Residual descent of decimal reversal]
\label{int:pf84:SYN30-RHO-DESC}
Reversion on decimal words of length two is involutive and its
shadow modulo nine is the identity:
\[
 \rho_{10}^2=I,\qquad
 q_9\!\left(\operatorname{ev}_{10}(\rho_{10}(a,b))\right)
 =
 q_9\!\left(\operatorname{ev}_{10}(a,b)\right).
\]
Fixed length preserves the initial zeros.
\end{lemma}

\begin{proof}
The doble permutation recovers \((a,b)\). Ifnce \(10\equiv1\pmod9\),
\[
 10b+a\equiv a+b\equiv10a+b\pmod9.
\]
A word whose double reversal were not itself, or whose digit sum
changed modulo nine, would refute the lemma.
\end{proof}

\begin{lemma}[Residual Square Descent]
\label{int:pf84:SYN30-SQ-DESC}
The map \(s(n)=n^2\) induces the residual endomap
\[
 [n]_9\longmapsto[n]_9^2,
 \qquad q_9(s(n))=q_9(n)^2.
\]
\end{lemma}

\begin{proof}
The compatibility of a congruence with the product gives
\(n\equiv m\pmod9\Rightarrow n^2\equiv m^2\pmod9\). An integer for which
\(q_9(n^2)\ne q_9(n)^2\) would refute the statement.
\end{proof}

\begin{proposition}[Residual fibers of the full reversal]
\label{int:pf84:SYN30-RHO-NONFAITH}
The residual shadow does not reconstruct the integral lift
\(\widehat\rho_{10}=\operatorname{ev}_{10}\circ\rho_{10}\). Indeed,
\[
 18\equiv27\pmod9,\qquad
 \widehat\rho_{10}(1,8)=81\ne72=\widehat\rho_{10}(2,7).
\]
\end{proposition}

\begin{proof}
Both integers belong to residue class zero, but their reversals are
distinct. Therefore, there is no
\(\overline\rho:\mathbb Z/9\mathbb Z\to\mathbb N\) that simultaneously factors
those two values.
\end{proof}

\begin{proposition}[Residual fibres of the entire square]
\label{int:pf84:SYN30-SQ-NONFAITH}
The same shadow likewise does not reconstruct the integer square:
\[
 18\equiv27\pmod9,\qquad
 18^2=324\ne729=27^2.
\]
\end{proposition}

\begin{proof}
A function \(\overline s:\mathbb Z/9\mathbb Z\to\mathbb N\) that
factorized \(s\) would have to assign the two values
\(324\) and \(729\) to the same zero class, which is impossible.
\end{proof}

\begin{corollary}[Cross-witness \(27\)--\(72\)--\(729\)]
\label{int:pf84:SYN30-CROSS-WITNESS}
\[
 \widehat\rho_{10}(2,7)=72,\qquad
 27^2=729=9^3=3^6,
\]
and, in the alfabeto convencional \(\mathscr B_9\),
\[
 27=30_9,\qquad72=80_9,\qquad729=1000_9.
\]
The three integers have residual shadow zero and positive digital root nine.
\end{corollary}

\begin{proof}
The two operations are computed directly; successive divisions by
nine produce the three expressions. The result would fail if
\(\mathscr B_9\) were confused with \(\Wnine\), if any of the expressions were incorrect,
or if it were claimed that faithful residual descent reconstructs the operation and its
complete genealogy.
\end{proof}

\section{Native adder and quotient propagation}

\begin{algorithm}[Word addition]
Two words are traversed by positions, with initial carry \(c_0=0\). If there are
two digits at a position, \eqref{int:eq:local-plus} is applied; if there is only one,
that digit is taken as the provisional residue and the provisional quotient is zero.
When \(c_j\ne0\), a second additive cell is applied to the provisional residue
and to \(c_j\). The new residue is published, and the sum of the
quotients is transported to the next position. If a nonzero carry remains upon
completion, it is published as the final digit. In this manner, the local cells never receive the
absent symbol \(0\), which does not belong to \(\Dnine\).
\end{algorithm}

\begin{lemma}[Adder Invariant]
If $u\boxplus_\#v$ is the word produced by the preceding algorithm,
\[
 \operatorname{ev}_9(u\boxplus_\#v)
 =\operatorname{ev}_9(u)+\operatorname{ev}_9(v).
\]
The intermediate carry lies in $\{0,1,2\}$.
\end{lemma}

\begin{proof}
After processing the positions $0,\ldots,k-1$, let $c_k$ be the pending quotient.
The local identities imply
\[
 \sum_{j<k}(u_j+v_j)9^j
 =\sum_{j<k}w_j9^j+c_k9^k,
\]
including the absent digits as zero. Each new cell preserves this
equality. The local maximum is $9+9+2=20$, so that $c_{k+1}\leq2$. At the
end the last quotient is published, and the global identity remains.
\end{proof}

\section{Multiplication by a digit}

\begin{algorithm}[Scalar transducer]
With $d\in\Dnine$ fixed, each digit $u_j$ produces
\[
 u_jd=9q_j+r_j,
 \qquad
 (r_j,q_j)=\bigl(R^\times(u_j,d),q^\times(u_j,d)\bigr).
\]
Let $c_j\in\{0,\ldots,9\}$ be the incoming quotient. Since $0\notin\Dnine$,
the transition is defined by branches:
\[
 (w_j,c_{j+1})=
 \begin{cases}
 (r_j,q_j),&c_j=0,\\[1mm]
 \bigl(R^+(r_j,c_j),q_j+q^+(r_j,c_j)\bigr),&c_j\in\Dnine.
 \end{cases}
\]
Thus, no local map is evaluated outside its domain. The first
branch expresses the absence of a pending quotient; the second normalizes the
collision by means of the additive map.
\end{algorithm}

\begin{lemma}[Scalar invariant]
The word $d\odot_\#u$ produced by the transducer satisfies
\[
 \operatorname{ev}_9(d\odot_\#u)
 =d\operatorname{ev}_9(u).
\]
The multiplicative quotient belongs to $\{0,\ldots,9\}$.
\end{lemma}

\begin{proof}
The argument is that of the adder, replacing each local summand by
$u_jd$. If $c_j=0$, the identity
$u_jd=9q_j+r_j$ directly preserves the partial equality. If
$c_j\in\Dnine$, the second Euclidean division gives
\[
 r_j+c_j=9q^+(r_j,c_j)+R^+(r_j,c_j),
\]
and preserves the same equality after transporting
$q_j+q^+(r_j,c_j)$. The maximum $9\cdot9+9=90$ and the remainder form in
$\Dnine$ keep the quotient within the declared bound. The last quotient is
published as a digit.
\end{proof}

\section{Global product by Horner}

\begin{definition}[Native product]
Let $v=(v_0,\ldots,v_{m-1})$. Partiendo of $a_m=\varnothing$, is defined for
$j=m-1,\ldots,0$
\[
 a_j=(9\odot_\#a_{j+1})\boxplus_\#(v_j\odot_\#u).
\]
The result is
\[
 u\boxtimes_\#v:=a_0.
\]
\end{definition}

This definition does not contain
\[
 \operatorname{dec}_9(\operatorname{ev}_9(u)\operatorname{ev}_9(v)).
\]
The evaluation will appear only in the proof of entanglement.

\begin{theorem}[Native multiplicativity]
For all words $u,v\in\Wnine$,
\begin{equation}
 \label{int:eq:ev-multiplicative}
 \boxed{
 \operatorname{ev}_9(u\boxtimes_\#v)
 =\operatorname{ev}_9(u)\operatorname{ev}_9(v).}
\end{equation}
The empty word is absorbing, and $(1)$ is the unit.
\end{theorem}

\begin{proof}
By the two previous lemmas,
\[
 \operatorname{ev}_9(a_j)
 =9\operatorname{ev}_9(a_{j+1})
  +v_j\operatorname{ev}_9(u).
\]
Horner's iteration gives
\[
 \operatorname{ev}_9(a_0)
 =\operatorname{ev}_9(u)
  \sum_{j=0}^{m-1}v_j9^j.
\]
The last factor is $\operatorname{ev}_9(v)$. The uniqueness of the normal form
identifies the resulting word.
\end{proof}

\begin{corollary}[Multiplicative monoid of words]
The structure
\[
\WordMonoid=(\Wnine,\boxtimes_\#,(1))
\]
is a commutative monoid; the empty word is its absorbing element. Its
construction does not use global multiplication of integers.
\end{corollary}

The shortest boundary example is
\[
 (9)\boxtimes_\#(9)=(9,8),
 \qquad 9+8\cdot9=81.
\]
The \(17\) local cells of multiplication produce
\[
 (9,8)\boxtimes_\#(9,8)=(9,8,8,8),
 \qquad \operatorname{ev}_9(9,8,8,8)=6561.
\]

\section{Matrix linearization of the genealogical state and construction of the completely positive entangler}

\begin{definition}[Hilbert Sectors]
Let
\[
 \Krad=\ell^2(\Wnine^+),
 \qquad
 \Hplus=\ell^2(\N_{\geq1}).
\]
We define on the canonical bases
\[
 W_\times e_u=e_{\operatorname{ev}_9(u)}.
\]
In \(\Hplus\), the number operator is denoted by
\[
 \mathsf N e_n=n e_n,
 \qquad \Spec(\mathsf N)=\N_{\geq1}.
\]
The vacuum may be added separately as a vacuum vector $e_0$; it is not
confused with $\Hplus$, whose spectrum begins at one.
\end{definition}

\begin{theorem}[Unitarity and multiplicative square]
$W_\times$ is unitary. In the algebraic kernel
$\mathbb C[\Wnine^+\times\Wnine^+]$, let
\[
 \mu_\#(e_u\otimes e_v)=e_{u\boxtimes_\#v},
 \qquad
 \mu_{\N}(e_a\otimes e_b)=e_{ab}.
\]
Then the square holds.
\[
\begin{CD}
\Krad\otimes_{\mathrm{alg}}\Krad @>{W_\times\otimes W_\times}>>
\Hplus\otimes_{\mathrm{alg}}\Hplus\\
@V{\mu_\#}VV @VV{\mu_{\N}}V\\
\Krad @>>{W_\times}> \Hplus
\end{CD}
\]
and therefore, the square commutes.
\end{theorem}

\begin{proof}
The bijection of normal forms sends one orthonormal basis to another and proves
unitarity. On a basic tensor, both paths produce
$e_{\operatorname{ev}_9(u)\operatorname{ev}_9(v)}$ by
\eqref{int:eq:ev-multiplicative}. Linearity gives the equality on the kernel.
\end{proof}

The maps $\mu_\#$ and $\mu_{\N}$, initially defined on the respective
algebraic kernels, are densely defined and unbounded. For
\(c=(c_{u,v})\in\ell^2(\Wnine^+\times\Wnine^+)\), each sum over a product fiber
is finite. The adjoint of \(\mu_\#\) contains all finite-support vectors
in its domain, so that \(\mu_\#\) is closable. The domain of its
closure is
\[
 \Dom(\overline{\mu_\#})=
 \left\{c\in\ell^2(\Wnine^+\times\Wnine^+):
 \sum_{w\in\Wnine^+}
 \left|\sum_{u\boxtimes_\#v=w}c_{u,v}\right|^2<\infty
 \right\}.
\]
Each fiber is finite because $\operatorname{ev}_9$ identifies it with the ordered
pairs of divisors of an integer. The same description holds for
\(\overline{\mu_{\N}}\) after applying \(W_\times\otimes W_\times\), and
$W_\times$ intertwines the graphs of both closures.

\section{Native coproduct and native antipode}

The multiplicative opening is now defined before looking at the codomain:
\[
 \Delta_\#e_w
 =\sum_{u\boxtimes_\#v=w}e_u\otimes e_v.
\]
The sum is obtained by enumerating products of the transducer. One does not ask whether
one ordinary integer divides another.
In the integer sector define, on the corresponding algebraic kernel,
\[
 \Delta_\times e_n
 =\sum_{\substack{a,b\in\N_{\geq1}\\ab=n}}e_a\otimes e_b.
\]
Both coproducts enumerate ordered pairs.

\begin{theorem}[Native coassociativity]
In the algebraic core,
\[
 (\Delta_\#\otimes I)\Delta_\#
 =(I\otimes\Delta_\#)\Delta_\#.
\]
Moreover,
\[
 (W_\times\otimes W_\times)\Delta_\#
 =\Delta_\times W_\times.
\]
\end{theorem}

\begin{proof}
Both sides of the first equality enumerate every native triple
$a\boxtimes_\#b\boxtimes_\#c=w$ exactly once. The second equality follows from the bijection
between fibers induced by \eqref{int:eq:ev-multiplicative}.
\end{proof}

\begin{theorem}[Complete arithmetic coproduct]
\label{int:pf84:C06}
In the algebraic kernel
\(c_{00}(\mathbb N_{\geq1})\subset\ell^2(\mathbb N_{\geq1})\), let
\[
 \Delta_\times e_n=\sum_{ab=n}e_a\otimes e_b.
\]
Its closure has domain
\[
 \operatorname{Dom}(\overline{\Delta_\times})
 =
 \left\{x=(x_n):
 \sum_{n\geq1}d(n)|x_n|^2<\infty\right\},
\]
is coassociative on the kernel and, for \(Qe_n=ne_n\), satisfies
\[
 (\log Q\otimes I+I\otimes\log Q)\Delta_\times
 =\Delta_\times\log Q.
\]
\end{theorem}

\begin{proof}
The supports of \(\Delta_\times e_n\) are orthogonal for distinct integers,
and their squared norm is the number \(d(n)\) of divisors, which yields the weighted
domain of the closure. The two coassociativity relations reindex the same sum
over \(abc=n\). Finally, \(\log a+\log b=\log n\) in each term
proves the intertwining. A triple factorization counted differently, or failure of
this identity, would falsify the theorem.
\end{proof}

\begin{theorem}[Reduced coproduct and irreducible selector]
\label{int:pf84:C07}
In the sector integer, the coproducto reducido
\[
 \Delta_\times^{\rm red}e_n
 =\sum_{\substack{ab=n\\a,b>1}}e_a\otimes e_b
\]
is closable. The operator
\[
 N_\times^{\rm red}
 =\bigl(\overline{\Delta_\times^{\rm red}}\bigr)^*
  \overline{\Delta_\times^{\rm red}}
\]
is positive self-adjoint and diagonal:
\[
 N_\times^{\rm red}e_n=d_{\rm red}(n)e_n.
\]
Its kernel is generated by \(e_1\) and by the \(e_p\) with \(p\) irreducible;
therefore,
\[
 P_{\Irr}
 =\mathbf1_{\{0\}}(N_\times^{\rm red})
  -\lvert e_1\rangle\langle e_1\rvert
\]
select the irreducibles and exclude the unit.

In the native sector, the corresponding coproduct is
\[
 \widetilde\Delta_\#e_w
 =\sum_{\substack{u\boxtimes_\#v=w\\u,v\ne(1)}}e_u\otimes e_v.
\]
The entrelazador \(W_\times\) transports
\(\Delta_\times^{\rm red}\), \(N_\times^{\rm red}\) and \(P_{\Irr}\) toward
\(\widetilde\Delta_\#\),
\[
 A_\#:=
 \bigl(\overline{\widetilde\Delta_\#}\bigr)^*
 \overline{\widetilde\Delta_\#},
 \qquad
 P_{\Irr,\#}
 =(I-P_{(1)})\mathbf1_{\{0\}}(A_\#),
 \qquad
 P_{(1)}=\lvert e_{(1)}\rangle\langle e_{(1)}\rvert.
\]
\end{theorem}

\begin{proof}
The supports of the images of two distinct basis vectors are orthogonal.
The squared norm of
\(\Delta_\times^{\rm red}e_n\) is \(d_{\rm red}(n)\), which vanishes
exactly for \(n=1\) or \(n\) irreducible. This proves the diagonalization
and the declared kernel; subtracting the projector of \(e_1\) removes the unit.
The multiplicativity of \(W_\times\) identifies the integer fibers with the
native fibers and transports the three operators. A composite in the kernel,
an irreducible outside it, or inclusion of the unit in \(P_{\Irr}\)
would refute the theorem.
\end{proof}

\begin{proposition}[Nonadic coincidence of primes and composites]
\label{int:pf84:B05}
Let \(\gcd(n,90)=1\),
\(\tau_n=\operatorname{ord}_n(10)\) and
\[
 M_n=\frac{10^{\tau_n}-1}{n}.
\]
Then \(M_n\equiv0\pmod9\), both for primes and for composites. In
particular, the nonadic closure and the isolated decimal period are not a
selector of irreducibles.
\end{proposition}

\begin{proof}
Because \(9\mid10^{\tau_n}-1\) and because \(n\) is invertible modulo nine, divisibility by nine survives passage to the quotient. The primes \(7,13\) and the composites \(77,91,143\) share period six; the Carmichael numbers provide additional controls. The statement would be falsified only if the congruence or the period separated all primes from all composites in the declared domain.
\end{proof}

This statement is strictly arithmetic: it identifies factorization
fibers and the operators they induce on the sector of normal
forms.

\ifdefined\HMTInternalTraceability
\section{Monoidal extension of the native product to compatible histories}
\label{int:pf84:SYN-05D-HIST}

The result distinguishes the space \(\Krad\) of canonical normal forms and
the space \(\mathcal K_{\#,\mathrm{hist}}\) of complete histories of
cells and propagated quotients.
Many histories may produce the same normal form. Therefore,
$W_\times$ is unitary on $\Krad$, but the projection of a history onto
its value prevents it from being injective on all transport classes.

Let \(\mathscr X_n^{\rm hist}\) be the set of enriched histories of
depth \(n\) and let
\(\mathcal H_n=\ell^2(\mathscr X_n^{\rm hist})\). Denote by \(P_n\) the
orthogonal projection onto the subspace generated by surviving
histories, and suppose a multiplication \(\mu_n\) is defined on an
algebraic domain of
\(\mathcal H_n\otimes_{\mathrm{alg}}\mathcal H_n\). A monoidal extension
of the word product to the inverse system consists of a family of
maps \(\iota_n:\Krad\to\mathcal H_n\) that commutes with the
depth restrictions and satisfies, for every \(n\),
\[
 P_n\,\iota_n\mu_\#
 =P_n\,\mu_n
 (\iota_n\otimes\iota_n)
\]
in the algebraic kernel. This equality is the compatibility condition that
defines the extension. The preceding theorems establish neither the existence nor
the uniqueness of a family \((\iota_n)_n\); they prove the product in the
monoid \(\WordMonoid\) and determine the equation that an extension to the history
space must satisfy. The transfer therefore retains the status of an
\emph{open program}: it does not inherit the exactness of the product in
\(\WordMonoid\). The nonexistence of a compatible family, failure to
commute with truncations, loss of survival, quotient,
carry, orientation, or memory, or failure to satisfy the preceding identity
are its falsifiers.
\fi

\section{Exhaustive enumeration and necessity of structural coordinates}

The finite enumeration comprises:
\[
\begin{array}{@{}lr@{}}
\toprule
\text{Test}&\text{Result}\\\midrule
\text{Cells per local map}&81\\
\text{Reversible numbering paths}&59050\\
\text{Intertwiner pairs}&532900\\
\text{Intertwiner errors}&0\\
\text{Associative triples}&531441\\
\text{Associative errors}&0\\
\text{Coproduct window}&1\text{ to }500\\
\text{Irreducibles obtained}&95\\
\bottomrule
\end{array}
\]

\begin{definition}[Multiplicative quotient]
Preserving only $R^\times$ transforms $2\cdot5=10$ into the visible value $1$.
Perturbing $q^\times(9,9)$ from $8$ to $7$ transforms $81$ into $72$. The local
quotient is necessary.
\end{definition}

\begin{definition}[Additive application]
The multiplicative map produces partial products, but does not resolve their collisions.
Removing $q^+$ breaks the global normalization; for example, it no longer reproduces
the intermediate quotients of $10\cdot18=180$.
\end{definition}

\begin{definition}[Coordinates and boundary symbol]
Ordinary zero is not a digit of the construction. The empty word is
$0$, and the symbol $9$ preserves its value and quotient. Mixing
zero-indexed positional coordinates with the labels $1,\ldots,9$ changes the
cells.
\end{definition}

\begin{proposition}[Exact scope of the ablations and of the finite census]
\label{int:pf84:D10}
The preceding ablations prove, within the declared conventions:
\begin{enumerate}
  \item without \(q^\times\), \(2\cdot5=10\) colapsa to the residue visible \(1\);
  \item cambiar \(q^\times(9,9)=8\) by \(7\) transforms \(81\) in \(72\);
  \item without \(q^+\), normalization of the additive collisions fails.
\end{enumerate}
The exhaustive enumeration comprises the \(2\cdot81=162\) cells of the two
maps, the \(729\) local triples, \(532\,900\) word pairs, and
the \(95\) irreducibles no greater than \(500\). The indicated identities are
proved on these finite domains, but the coinductive survival
of every TPK history is not.
\end{proposition}

\begin{proof}
The three counterexamples are computed directly with
\eqref{int:eq:local-plus} and \eqref{int:eq:local-times}. The enumeration of the four
finite domains produces the cardinals written in the statement and does not add
any coinductive identity to those already proved in the monoid
of words.
\end{proof}

\section{Algebraic structure of the word monoid}

The construction proves, in the monoid \(\WordMonoid\):
\begin{enumerate}
 \item a unique normal form for each nonnegative integer;
 \item global sum and product constructed solely from the two local maps;
 \item symbolic conservation of the value under quotient propagation;
 \item a unitary $W_\times$ and the exact multiplicative square;
 \item a divisor coproduct and an irreducible selector defined natively;
 \item the need to preserve the transport class when leaving the
normal-form sector.
\end{enumerate}

Passage to holonomic space now reduces to an existence problem: constructing a system \((\iota_n)_n\) compatible with the depth restrictions and proving that it satisfies the equation of the preceding section. This question is distinct from the exact intertwiner on \(\WordMonoid\), already determined by the normal form and local maps.
 

\label{p3-83-c23-s1:ch:foundation}

\section{Problem of determining the constants from the genealogical state}

A physical constant is represented as a section of a magnitude line equipped with a rule of evaluation, orientation, and normalization. Its numerical coordinate depends on the chosen chart, whereas its genealogy remains invariant. Consequently, the determination problem comprises:

\begin{enumerate}
\item the identification of the antecedent HMT state;
\item the specification of the operator or character of selection;
\item determination of the conserved invariant;
\item the corresponding magnitude line;
\item the chart in which the terminal value is expressed.
\end{enumerate}

The causal chain of the work is

\begin{equation}
\begin{aligned}
\APP&\longrightarrow\TRIT\longrightarrow\TPK
\longrightarrow\text{internal realization}
\longrightarrow\text{dimensionless invariant}\\
&\longrightarrow\text{magnitude section}
\longrightarrow\R.
\end{aligned}
\label{p3-83-c23-s1:eq:master-chain}
\end{equation}

The term \(\R\) in the chain represents the terminal evaluation of the
section; it does not constitute the original domain from which the
route is retrospectively reconstructed.

\section{Antecedent chain APP--TRIT--TPK and domain of the extraction operators}

APP supplies the geometry of positions, residues, quotients, and
boundaries. TRIT incorporates orientation \((-1,0,+1)\) and distinguishes
opening, threshold, and return. TPK completes the state through carry, route,
sheet, memory, and survival. In particular, transition \(9\to1\)
preserves phase and modifies the memory state; it therefore does not constitute
a reinitialization devoid of genealogy.

The realizations developed below preserve that antecedent without replacing it. The vacuum is obtained through a spectral function of the bilayer operator; gravity selects a section through the chronological structure \(108\); criticality uses a dodecaphasic family; the Gaussian constants follow from the evaluation of cylinders; and the modular Hamiltonian is defined on a boundary that preserves the incidence \(6\to8\). The production of scalar values does not establish functional identity among these realizations.

\begin{proposition}[Insufficiency of the Scalar Coincidence]
Let \(F_1:X_1\to L_1\) and \(F_2:X_2\to L_2\) be two typed functionals. Equality of
coordinates \(c_1(F_1(x_1))=c_2(F_2(x_2))\) does not by itself induce either an
identification \(X_1\simeq X_2\) or an intertwiner
\(F_2I=JF_1\).
\end{proposition}

In particular, a numerical proximity such as
\(I_{\rm or}\approx\gammaH a_0\) cannot be promoted to an identity without
constructing the corresponding intertwiner.

\section{Lines of magnitude, oriented letters and dimensional coordinate systems}

Let \(L_Q\) be an oriented real line representing the magnitude \(Q\).
A unit \(u_Q\in L_Q\setminus\{0\}\) determines a
provisional coordinate \(L_Q\simeq\R\), without replacing the intrinsic structure of
the line. The constitutive relation

\[
L_\mu=L_S L_C^{-1}L_Q^{-2},\qquad
L_\varepsilon=L_S^{-1}L_C^{-1}L_Q^2
\]

precedes the SI choice. In the same way,

\[
L_Z=L_SL_Q^{-2},\qquad L_G=L_C^5L_T^2L_S^{-1}
\]

type impedance and the gravitational constant before expressing their magnitudes in
ohms or SI units. In this work we shall write

\[
Q_{\rm int}=\widehat Q\,u_Q
\]

to distinguish the dimensionless character \(\widehat Q\) from the reference
section \(u_Q\).

\section{Torsor of frames and covariance of dimensional evaluations}

The admissible charts form a torsor and therefore none has a character that is canonical relative to the others. The frame change \(u_Q\mapsto\lambda_Q u_Q\) induces the contragredient transformation of the coordinate. Physical identities must be homogeneous with respect to this action. For example,

\[
\mu\varepsilon=c^{-2},\qquad \frac\mu\varepsilon=Z^2
\]

are valid in every compatible frame. Mathematical unification and
dimensional unification do not eliminate units; rather, they reconstruct them through the
covariance of the lines and allow the invariant to precede its coordinate
expression.

\section{Final metrological evaluation and prohibition of retrospective selection}

Each entry of the final atlas adopts the form

\begin{equation}
\boxed{\text{antecedent}\to\text{operator}\to\text{invariant}
\to\text{meaning}\to\text{terminal value}.}
\label{p3-83-c23-s1:eq:atlas-rule}
\end{equation}

This order expresses a causal dependency: the numerical value is only
determined after fixing the antecedent, operator, invariant, and
interpretation. Omission of any of these elements reduces the
result to a reconstruction that is numerical, not a derivation that is
genealogical.

\section{Exact domain of the genealogical generation and its terminal metrical chart}

The generation of
\(\pi,e,\varphi\), the dodecaphasic closure of \(\alpha\),
the action section, and the Barbero--Immirzi output \(\gammaH\) are adopted as initial results. The
present development does not recalibrate those constructions; it determines the constants
and invariants resulting from their application to the constitutive,
thermal, gravitational, modular, and spectral realizations.

The electron mass is likewise not included as a governing problem, since it
belongs to the General Mass Law. Every occurrence of \(m_e\) in an
atomic identity is treated as a datum proceeding from that domain and not as
a new selection of the present development.

\subsection*{Logical Scope and Domain of Validity}
The APP--TRIT--TPK architecture and the dimensional lines constitute the
starting structure. The uniform rule of five fields of the atlas makes
explicit its causal discipline in each entry.



\section{Functional taxonomy of \(114\) numerical expressions and \(106\) operator readings}
\label{p3-83-c23-s2:sec:atlas-taxonomico-114-constantes-rev11}
\index{constants!taxonomic atlas of 114 entries}

Let the finite indexed set be
\[
 E_{114}=\{\varepsilon_1,\ldots,\varepsilon_{114}\}
\]
of expresiones numericas and let \(\mathcal O\) the space of readings of the
reader triadico. The map
\[
 T:E_{114}\longrightarrow\mathcal O,
 \qquad \varepsilon_i\longmapsto T_{\rm tri}(\varepsilon_i),
 \qquad |\operatorname{im}T|=106,
\]
classifies expressions according to their operatorial role. Each
\(\varepsilon_i\), including the compared targets, is fixed as
input of \(T\); therefore, this map does not generate the values that
it evaluates. The functional partition is
\[
\begin{array}{l|r}
\text{functional family}&\text{entries}\\ \hline
\text{structural kernel}&5\\
\text{coupling and closure}&1\\
\text{bridges between modes}&7\\
\text{angular and closure}&22\\
\text{exponential--logarithmic}&12\\
\text{spectral and regularization}&14\\
\text{hyperbolic and modular}&2\\
\text{algebraic and metric}&50\\
\text{typed miscellany}&1
\end{array}
\qquad
\sum n_i=114.
\]
Network centrality, shortest paths, and communities computed on
the graph induced by \(T\) describe connectivity within the classification;
they do not establish mathematical causality or HMT genealogy.

Seven fibers of \(T\) gather fifteen expressions with a common reading. A
shared fiber does not identify their values either. In particular,
\(\varphi^{-1}\), \(\varphi\), and \(\varphi^2\) are distinct real numbers that the
map \(T\) groups into the same self-scaling orbit; the common
projection preserves dynamic function and forgets the Archimedean coordinate that
separates them.

The exhaustive enumeration traversed the \(114\) entries of \(E_{114}\), four hundred sixty-eight direct candidates \(U_{12}\), and one hundred four thousand nine hundred seventy-six weighted seeds without finding, within that finite domain, a rival that reproduced the oriented dodecaphasic correction of \(\alpha\). The search domain is thereby bounded by those three sets. The census proves that no rival appears in the enumerated class; it does not prove uniqueness among all possible functionals and does not enter as a premise of theorem \(\alpha\).

For each \(\varepsilon_i\), the taxonomy assigns the quadruple
\((\text{object},\text{reader},\text{domain},\text{structural function})\)
used in the enumeration.
 \section{Spectral, informational, and thermodynamic confluence of \(\log3\)}

The identity

\[
\log3=C_1+C_2-2C_0
=S_{\TRIT}/k_B
=E_P/(k_B\Theta_{\rm clk}),
\]

relates three distinct realizations. The first member is a residual
defect, the second an entropy of decision, and the third a ratio
between energy and temperature. Their equality proves the conservation of the
scalar under a change of realization, without identifying the state
spaces.

The first equality is deduced from the three residual channels.
\[
Z_r(s)=\sum_{\substack{n\ge1\\n\equiv r\ ({\rm mod}\ 3)}}n^{-s},
\qquad r=0,1,2.
\]
Then
\[
Z_0(s)=3^{-s}\zeta(s),\quad
Z_1(s)+Z_2(s)=(1-3^{-s})\zeta(s),\quad
Z_1(s)-Z_2(s)=L(s,\chi_{-3}).
\]
If \(C_r\) denotes the part finite of \(Z_r\) in \(s=1\), the expansion
of \(3^{-s}\) gives
\begin{equation}
C_0+C_1+C_2=\gamma,qquad
-2C_0+C_1+C_2=\log3,qquad
C_1-C_2=\frac{\pi}{3\sqrt3}.
\label{eq:residual-three-covectors}
\end{equation}
The three covectors are linearly independent. From a single residual
decomposition, they recover the total finite term, the dilation of the
zero class, and the orientation between the two nonzero classes. \(\log3\) is
therefore a scale-change invariant before being interpreted as
entropy or temperature.

\section{Epistemological classification of the derived constants}

The chapter's contributions are classified into three categories:

\begin{enumerate}
\item assignment of an internal operator to a known constant, as in
the Catalan case;
\item exact composed evaluation of two internal characters, such as
\(L(1,\chi_{12})\);
\item derivation of a constant from a field or a family
constructed internally, such as \(\gamma_{\mathbb Q(\sqrt3)}\) or \(A_k\).
\end{enumerate}

Mathematical provenance determines the status of recovery or
novelty; decimal proximity plays no role in that classification.

\begin{remark}[Negative control]
The character \(\chi_{12}\) must have exactly the pattern
\((+,-,-,+)\) in \((1,5,7,11)\); any other pattern invalidates
\cref{eq:L1chi12,eq:L2chi12}. The identity
\cref{eq:odd-zeta-bendersky} and the factorization
\cref{eq:dedekind-sqrt3} are independent controls.
\end{remark}

\subsection*{Logical Scope and Domain of Validity}
\(J^2=-I\), Catalan, \(\chi_{12}\), their values \(L\), Euler--Kronecker, and the Bendersky--Glaisher tower: \emph{Exact internal theorem}. The gathering in a lattice \(3/4/12\) is \emph{New formalization} that preserves the mathematical provenances.

 All these results are joined because they turn the terminal value into the quantitative expression of a previously constructed structure. Numerical exactness accompanies a different conception of what a fundamental constant is. The nonadic connection does not first generate three constants and then add a fourth: the prolongation \(w_6\to w_{12}\to w_{18}\to w_{24}\to w_{30}\to R_{36}\to G_9\) stabilizes a common coinductive state from which \(\pi\), \(e\), \(\varphi\) and \(\alpha\) are published in a correlated manner. Their readers distinguish closure, propagation, self-scaling and coupling; within this joint generation, coordinate \(\alpha\) retains its signed, dodecaphasic realization and its two internal derivation routes.\par

In the usual description, a constant is commonly presented as a number placed before an equation: a quantity that the theory requires but whose value must be received from outside. In Triadic Modular Holography, the opposite occurs. The constant appears at the end, after the same state has been traversed by several readers and it has been determined which relation has survived all of them.\par

This is why the constants in the document are magnitudes endowed with genealogy: invariant remnants of a transformation. Each states what has been preserved when something has changed:\par

\begin{itemize}[leftmargin=*,itemsep=0.35em]
\item to sheet has been inverted;
\item a cycle has returned to its initial phase;
\item a memory has crossed the seam;
\item an interior description has been published on the boundary;
\item a dimensionless magnitude has received a physical realization;
\item a system has changed scale while fully preserving its genealogy.
\end{itemize}

Dimensional Mechanics subsequently adds the magnitude lines and the units. But the structure that selects the constant is prior. The real value appears at the end because \(\mathbb R\) functions here as a terminal evaluation: it is the quantitative shadow of a construction that was already complete before becoming a decimal.\par

With that view, the nine developments considered form a single narrative on orientation, memory, geometry, information, and scale.\par
  \endgroup
% Rectificación quirúrgica: generación conjunta, profundidad arbitraria y
% estatuto tipado de los transportes ternarios.  Este bloque pertenece al
% capítulo 27 y precede a la separación expositiva de los cuatro lectores.

\section{Joint theorem of coinductive generation and arbitrary depth}
\label{sec:c27-generacion-coinductiva-profundidad-arbitraria}

The starting structure of this theorem is not the bare set of nine
symbols. It is the typed system
\[
 \mathfrak H_9=
 \bigl(\mathrm{APP},\mathrm{TRIT},\mathrm{TPK};
       \widehat{\mathcal X}^{\mathrm{enr}},\mathcal U,
       \Gamma_9,\operatorname{Ext},\operatorname{tr}\bigr),
\]
constructed on the double positional and arithmetic realization of
\(\mathcal D_9=\{1,\ldots,9\}\). APP accumulates and folds through its sheets
\(\Sigma\) and \(\Pi\), preserving residue and quotient; TRIT orients each
transition in one of its three regimes; and TPK selects, transports,
memorizes and returns on the enriched fiber. In particular,
\[
 \mathcal U_t
 =\operatorname{Rec}_t\circ\operatorname{Mem}_t
  \circ T_{d(t)}\circ M_{\mathrm{ph}}(g(t)).
\]
This composition is the effective dynamic datum: the nine-symbol chart is its finite support, not a substitute for the dynamics.

For \(r\in\{6,12,18,24,30,36\}\), let
\(\widehat{\mathcal X}^{\mathrm{enr}}_r\) be the truncated states with sheet,
orientation, residue, quotient, carry, route, boundary and memory, and let
\[
 \operatorname{Ext}_{r+6,r}:
 \widehat{\mathcal X}^{\mathrm{enr}}_r
 \longrightarrow
 \widehat{\mathcal X}^{\mathrm{enr}}_{r+6},
 \qquad
 \operatorname{tr}_{r+6,r}:
 \widehat{\mathcal X}^{\mathrm{enr}}_{r+6}
 \longrightarrow
 \widehat{\mathcal X}^{\mathrm{enr}}_r
\]
admissible extension and truncation. The compatibility \(\operatorname{tr}_{r+6,r}\operatorname{Ext}_{r+6,r}=\mathrm{id}\) extends to every level. The notation
\[
 w_6\longrightarrow w_{12}\longrightarrow w_{18}\longrightarrow w_{24}
 \longrightarrow w_{30}\longrightarrow R_{36}\longrightarrow G_9
\]
designates these maps and their type changes, not a list of labels.
The holonomy \(\Gamma_9\) satisfies
\[
 g(k+9)=g(k),
 \qquad q^{(9)}(k+9)=q^{(9)}(k)+1,
\]
so that the phase returns while the quotient, the route and the
memory of the state advance.

\begin{theorem}[Inaugural construction from the seeds]
\label{thm:c27-construccion-inaugural-semillas}
Let
\[
 I_9=\{0,\ldots,8\},\qquad D=\{N,E,S,O\},\qquad
 \mathcal S=I_9^2\times D,
\]
and let \(\Omega_{\mathrm{seed}}=\mathcal S_+\times\mathcal S_\times\) be the
domain of the two cursors realizing the additive and multiplicative sheets
of APP. Then
\[
 |\Omega_{\mathrm{seed}}|=(9^2\cdot4)^2=104\,976,
\]
and the \(54\)-step recurrence and ternary reduction define
\[
 \Omega_{\mathrm{seed}}
 \xrightarrow{\ T_{54}^{\rm seed}\ }
 \mathcal U_6^{\mathrm{cat}}
 \xrightarrow{\ \rho_3\ }
 \mathcal W_6
 \hookrightarrow\mathbb F_3^6,
\]
with
\[
 |\mathcal U_6^{\mathrm{cat}}|=468,
 \qquad |\mathcal W_6|=243,
 \qquad |\mathbb F_3^6|=729.
\]
The fibers satisfy exactly
\[
 432\cdot144+18\cdot432+18\cdot1944=104\,976.
\]
The TPK extension relation extends the surviving words to a compatible section \(x_\infty\) of infinite depth. On that section, the Archimedean shadow and the exceptional realization of \cref{thm:c27-tesis-rectora-doble-proyeccion} are constructed jointly.
\end{theorem}

\begin{proof}
The preceding Cartesian product exhaustively enumerates the positions and orientations of both cursors. The map \(T_{54}^{\rm seed}\) is obtained by composing, at each step, sheet selection, cardinal transport, quotient updating and carry memory. Enumeration of its image contains \(468\) distinct sextuples; the sum of its fibre cardinalities is the stated identity. The componentwise map \(\rho_3\) contains \(243\) distinct words. Finally, nonemptiness, single-valuedness on the surviving subdomain and compatibility of \(\operatorname{Ext}\) with truncations produce the inverse-limit section. Each statement is a finite equality or a consequence of projective compatibility; none uses a real number, a recognized constant or a decimal table as an input.
\end{proof}

\begin{theorem}[Governing thesis: one continuum and two projections]
\label{thm:c27-tesis-rectora-doble-proyeccion}
Let \(\Omega_\Gamma^{\mathrm{cal}}\subset
X_\infty^{\mathrm{cal}}\) be the space of surviving sections of the
enriched turns of the TPK. The carry of the three turns
\(\gamma_-,\gamma_0,\gamma_+\) defines, without prior Archimedean evaluation,
the homeomorphism
\[
 \beta:\Omega_\Gamma^{\mathrm{cal}}\xrightarrow{\ \sim\ }
 (\mathbb F_3^6)^{\mathbb N}=:\mathcal W_\infty.
\]
The normalized balanced forms
\[
 a_k+c_k=\rho_k+3c_{k+1},\qquad
 \rho_k\in\{-1,0,+1\},
\]
construct the ordered ring \(\mathbb Z_{\rm HMT}\), its field of
fractions \(\mathbb Q_{\rm HMT}\) and the completion
\[
 \mathbb R_{\rm HMT}:=
 \operatorname{Cau}(\mathbb Q_{\rm HMT})/{\asymp}.
\]
The joint discrete structure is the fiber product
\[
\mathcal C_{\mathrm{cont}}^{\mathrm{disc}}
=\bigl(\mathbb Z_{\rm HMT}\times
\Omega_\Gamma^{\mathrm{cal}}\bigr)
\times_{\mathcal W_\infty}X_\infty^{\mathrm{cal}}.
\]
Its two natural maps are
\[
 P_{\mathrm{ar}}:\mathbb Z_{\rm HMT}\times
 \Omega_\Gamma^{\mathrm{cal}}\longrightarrow\mathbb R_{\rm HMT},
 \qquad
 Q_\infty:X_\infty^{\mathrm{cal}}\longrightarrow
 (\mathscr C_W^{\mathrm{cal}})^{\mathbb N}.
\]
If \(\beta(\xi)=(b_q)_{q\geq0}\),
\(\nu(b)=\sum_{i=1}^{6}b_i3^{6-i}\) and
\[
 q_n(m,\xi)=m+\sum_{q=0}^{n-1}
 \frac{\nu(b_q)}{729^{q+1}},
\]
then \((q_n)\) is Cauchy and \(P_{\mathrm{ar}}(m,\xi)\) is its class
in \(\mathbb R_{\rm HMT}\). The rational cylinders
\[
 I_n(m;b_0,\ldots,b_{n-1})=
 \left[q_n,q_n+729^{-n}\right]
\]
are compatible, have diameter tending to zero and the recurrent partition
into \(729\) subcylinders constructs, for every Cauchy class, a tape
\((b_q)\) representing it. Consequently,
\[
 P_{\mathrm{ar}}\ \text{is surjective},\qquad
 (\mathbb Z_{\rm HMT}\times\Omega_\Gamma^{\mathrm{cal}})
 /{\sim_{\rm ar}}\ \cong\ \mathbb R_{\rm HMT}.
\]
Only after this internal construction does uniqueness of the complete Archimedean ordered
field recognize \(\mathbb R_{\rm HMT}\cong\mathbb R\).

The second preserves the calibrated incidence of each level; in coordinates \(x=((w_j,f_j))_{j\geq1}\), it acts by
\[
 Q_\infty(x)
 =\bigl(f_j,\Gamma_{f_jA_Wf_j^{-1}}(w_j)\bigr)_{j\geq1}.
\]
Thus, \(\mathbb R_{\rm HMT}\), subsequently recognized as \(\mathbb R\), is the Archimedean shadow obtained by contracting the complete genealogy of cylinders. The exceptional projection preserves that genealogy as boundary incidence and admits the successive Paley--Witt--Golay--Leech--\(V^\natural\)--Monster realization. Both branches proceed from the same state, and neither is reconstructed retrospectively from the other.
\end{theorem}

\begin{proof}
The three enriched turns exist after each prefix and their carry
distinguishes them by \(-1,0,+1\). Concatenation and truncation are
inverses, proving bijectivity and continuity of \(\beta\). For a
Cauchy class of \(\mathbb Q_{\rm HMT}\), one first chooses its ordered unit
cylinder and then, recursively, the unique half-open subcylinder
of the internal partition into \(729\) pieces containing the class; at a
common boundary, the two writings are identified by \(\sim_{\rm ar}\).
The sequence of left endpoints is precisely \((q_n)\), converges to the
given class and proves surjectivity without presupposing an expansion of
\([0,1]\) or a point of \(\mathbb R\). Identification with the real field
is the uniqueness theorem applied afterwards to \(\mathbb R_{\rm HMT}\).

The subspaces satisfying closure of a particular character are survival filters within the global capacity. Their function is to individuate the sections of \(\pi,e,\varphi,\alpha\); surjectivity of the continuum proceeds from completion of \(\mathbb Q_{\rm HMT}\). This separation of domains prevents a constant-selection condition from becoming a premise for coverage of \(\mathbb R\).

In the second branch, compatibility
\(r_{n,m}Q_n=Q_m p_{n,m}\) allows passage to the limit and defines \(Q_\infty\).
Conjugation by \(f_j\) transports the gauge without erasing incidence.
The subsequent code, lattice and vertex-algebra functors realize
the exceptional chain. Equality of the tapes in the fiber product
constructs the joint source and preserves complementary information: Archimedean
value in the first; incidence, boundary and genealogy in the second.
\end{proof}

\begin{corollary}[Infinite expansion as a mathematical object]
\label{cor:c27-expansion-infinita-objeto}
For each numerical character \(\chi\) published by the coinductive state, the
decimal reading is a complete sequence
\[
 \operatorname{Dec}_\chi(x_\infty)
 =(d_{\chi,1},d_{\chi,2},\ldots)
 \in\{0,\ldots,9\}^{\mathbb N}.
\]
Its prefix of length \(N\) is the projection
\(\rho_N\operatorname{Dec}_\chi(x_\infty)\). Thus a thousand, ten thousand or a
million digits are cuts of the same infinite object; generation does not
consist in successively raising an experimental bound.
\end{corollary}

\begin{theorem}[Correlated four-coordinate system of the integral coinductive holonomic state]
\label{thm:c27-cuadrirrelacion-generada}
The unique APP--TRIT--TPK composition of the system \(\mathfrak H_9\) determines a compatible state
\[
 x_\infty\in
 \widehat\Omega:=
 \varprojlim_r\widehat{\mathcal X}^{\mathrm{enr}}_r
\]
and four correlated characters
\[
 \bigl(\Pi_{\mathrm{cl}},\Pi_{\mathrm{prop}},
       \Pi_{\mathrm{aut}},\Pi_{\mathrm{dod}}\bigr)(x_\infty)
 =\bigl(\pi,e,
        \varphi,\alpha\bigr).
\]
The first three characters proceed respectively from the unique closure orbit, the unique propagation orbit and the unique autoscale orbit of the enriched catalogue. The fourth proceeds from the primitive signed coordinate of the same terminal state and its dodecaphasic closure. The four coordinates are therefore outputs of a single coinductive extension, despite their distinct expression maps.
\end{theorem}

\begin{proof}
The exhaustive census
\(104,976\to468\to243\subset\mathbb F_3^6\) decomposes the realized image
into \(43\) dihedral orbits. The fiber invariants leave a single orbit
in each of the closure, propagation and self-scaling classes; the internal
gauge orients in them, without decimal consultation, the representatives
\[
 w_6^{\mathrm{cl}}=010211,
 \qquad w_6^{\mathrm{prop}}=201101,
 \qquad w_6^{\mathrm{aut}}=121200.
\]
Extension preserves the predicates defining the three classes. In the closure channel, pentaregional composition with unit return produces the ratios \(1/5\), \(1/239\) and the relation \(120/119\); Machin’s identity recognizes the half-turn exactly. In the propagation channel, the normalized semigroup law forces \((n+1)a_{n+1}=a_n\), hence \(a_n=1/n!\) and \(P_1=e\). In the autoscale channel, the incidence
\[
 F_{\mathrm{av}}=\begin{pmatrix}0&1\\1&1\end{pmatrix}
\]
has a unique positive ray and its eigenvalue equation forces
\(x^2-x-1=0\), \(x>0\), that is, \(x=\varphi\).

At the terminal phase of the unique APP--TRIT--TPK structure, the reading \(x_{\mathrm{term}}\mapsto(B_{\mathrm{term}},U_{12}^{\mathrm{sgn}})\) simultaneously expresses the three-channel chart and the primitive signed coordinate. The latter is reversibly identified with the seal \(K\) and its dodecaphasic differences. The Euclidean recurrence with carry
\[
 s_m=P_m+E_m-\Phi_m-K_m+c_{m+1},\qquad
 A_m=s_m\bmod1000,\qquad
 c_m=(s_m-A_m)/1000
\]
produces \(\alpha\) without receiving it as an input. Independent analytic characterization through the kernel \(E_{21/22}\) determines a simple root in \((0,10^{-2})\), whose reciprocal coincides with that of the first route after applying \(\operatorname{Tr}_9\). Thus, separating the readers does not alter the common origin of the four coordinates.
\end{proof}

\begin{proposition}[Dodecaphase register of the complete state prior to \(\alpha\)]
\label{prop:c27-registro-u-k-producido}
The terminal holonomic state has the double expression
\[
x_{\mathrm{term}}
\xrightarrow{(\pi_{\mathrm{term}},R_{\mathrm{sgn}})}
(B_{\mathrm{term}},U),
\]
where
\[
U=(2378,1406,2479,-452,998,-551,
-668,-204,-371,-322,-28,-997).
\]
Upon ordering the coordinates by the three step-three orbits, the transform
\[
\mathcal H_{12}
=P_{\mathrm{orb}}^{-1}(H_4\oplus H_4\oplus H_4)P_{\mathrm{orb}},
\qquad H_4^2=4I_4,
\]
has inverse \(\mathcal H_{12}/4\) on its integral image and determines
uniquely
\[
K=(234,543,140,729,659,824,621,058,914,794,146,601).
\]
Therefore, within the integral holonomic state,
\[
x_{\mathrm{term}}\xrightarrow{R_{\mathrm{sgn}}}U
\xrightarrow{\mathcal H_{12}^{-1}}K
\]
is an exact composition preceding the carry that expresses \(\alpha\). The visible matrix \(B_{\mathrm{term}}\) and signed record \(U\) are two coordinates of the same state; the latter preserves orientation, opposition, quotient and record fields contracted by the former. There is no second APP or competing reduced profile here: APP produces the sheets and arithmetic record, TRIT orients them, and TPK accumulates in its ledger the signed window projected by \(R_{\mathrm{sgn}}\). Requiring isolated APP to express a coordinate belonging to the subsequent TPK codomain is a type error. The complete composition does not receive \(U\), \(K\), their differences, \(\alpha\) or target expansions as inputs.
\end{proposition}

\begin{theorem}[Subsequent bidirectional closure of the character \(\alpha\)]
\label{thm:c27-alpha-segunda-via-bidireccional}
Let \(\alpha=\alpha\) be the coordinate already expressed by dodecaphasic carry. Evaluating the vacancy jet of degrees \(7{:}9\) at that output determines the functional
\[
\begin{aligned}
B_9(\alpha):={}&\log_{10}\!\left(\frac{10}{9}\right)
+\frac74\alpha^2-\frac{\alpha^4}{20}
+\frac{2\alpha^5}{21}+\frac{\alpha^6}{46}\\
&+\frac{\alpha^7}{120}-\frac{\alpha^8}{45}
-\frac{2\alpha^9}{495}.
\end{aligned}
\]
Upon lifting the rotation to the positive unknown \(T\), the resolvent is equivalent to
\[
\boxed{\alpha T^2-B_9(\alpha)T+\alpha^3=0.}
\]
With
\(\xi=\operatorname{arcosh}(B_9(\alpha)/(2\alpha^2))\), the solutions are
\[
T_\pm=\alpha e^{\pm\xi},
\qquad T_+T_-=\alpha^2,
\qquad
J_\alpha(T)=\frac{\alpha^2}{T},
\qquad J_\alpha(T_\pm)=T_\mp.
\]
The chamber \(6<T<7\) selects \(T_+\); the finite closure \(\pi_9=T_+/2\) is compatible, at the declared resolution, with the closure channel independently extended by \(\Gamma_9\) to \(\pi\). The identity \(\alpha=\sqrt{T_+T_-}\) fixes the geometric mean of the conjugate branches. The jet \(7{:}9\), both roots and the involution constitute an exact bidirectional closure \emph{on} the already generated character. This theorem receives no conventional constant, but neither is it a second independent derivation of \(\alpha\): its domain contains \(\alpha\), and its result is a subsequent functional certificate.
\end{theorem}

\begin{theorem}[Exceptional expression correlated with the sector \(\alpha\)]
\label{thm:c27-alpha-rama-excepcional}
In the incidence chart of the same dodecaphase register, the radial roots
of \(A_2^{12}\) determine
\[
v_\alpha=4\rho_1+\rho_2+\cdots+\rho_{12},
\qquad v_\alpha^2=54,
\]
and the index-three neighbor
\[
N_{\alpha,0}
=\{x\in N(C_W):\langle x,v_\alpha\rangle\equiv0\pmod3\},
\qquad
N_\alpha=N_{\alpha,0}+\mathbb Z\frac{v_\alpha}{3}.
\]
The lattice \(N_\alpha\) is even, unimodular, of rank twenty-four and
rootless, so \(N_\alpha\cong\Lambda_{24}\). The chain
\[
C_W\longrightarrow N(A_2^{12})
\xrightarrow{v_\alpha}\Lambda_{24}
\longrightarrow V^\natural
\xrightarrow{\operatorname{Aut}}\mathbb M
\]
exposes the exceptional hinge. The scalar \(\alpha\) and vector \(v_\alpha\) belong to different codomains and share the originating dodecaphasic state. The exact result of this branch is the construction \(N_\alpha\cong\Lambda_{24}\to V^\natural\) and the subsequent action of \(\operatorname{Aut}(V^\natural)\cong\mathbb M\) on \(V^\natural\). The scalar is not identified with the vector, nor is an action of the Monster on digits or routes deduced from correlation; such a promotion would require an equivariant morphism with explicit domain and codomain.
\end{theorem}

\begin{proposition}[Status of the transports \(L_0,L_1\)]
\label{prop:c27-estatuto-l0-l1}
The endomorphisms \(L_0,L_1\in\operatorname{End}(\mathbb F_3^6)\) and calendar \((L_0,L_0,L_1,L_1)\) belong to the transport law of the integral holonomic state. The six expressed full-rank transitions uniquely fix each matrix through \(L=X^{-1}Y\) in \(\mathbb F_3\), and the calendar is the only one among the \(16\) binary calendars simultaneously preserving the three structural histories. Their function is to extend the already selected representatives; Archimedean evaluation and conventional nomenclature are applied afterwards.
\end{proposition}

The domain of this proposition is the unique structure \(\mathfrak H_9\). Its forgetful projections are subsequent readers and do not constitute alternative generating profiles. This typing prevents confusion of a modal shadow with the state producing it.

\begin{theorem}[Arbitrary depth and determination of the complete expansions]
\label{thm:c27-profundidad-arbitraria}
For every \(\chi\in\{\pi,e,
\varphi\}\) and every \(N\ge1\), there is a unique decimal prefix \(p_{\chi,N}\) of length \(N\), produced by a finite truncation of the integral holonomic state, such that
\[
 \rho_{N+1,N}(p_{\chi,N+1})=p_{\chi,N}.
\]
The compatible family is the family of projections of a unique element
\[
 p_{\chi,\infty}
 \in\varprojlim_N\{0,\ldots,9\}^{N}
 \cong\{0,\ldots,9\}^{\mathbb N},
\]
and the intersection of its Archimedean cylinders consists of a unique real number.
Consequently, the connection generates the complete expansion; each finite
precision is a projection of that infinite output.
\end{theorem}

\begin{proof}
For a precision \(N\), the rational closure and propagation brackets and consecutive Perron quotients are refined until contained in a single decimal cylinder of diameter \(10^{-N}\). The map \(\operatorname{Ext}\), the carry cocycle and TPK truncation make the cylinder of order \(N+1\) project onto that of order \(N\). The cylinders are nested and their diameters tend to zero; their inverse limit determines a unique point. Since \(N\) is arbitrary and the extension relation has no terminal depth, the conclusion holds for every finite length and, jointly, for the infinite expansion.
\end{proof}

\begingroup\fontsize{10}{11.4}\selectfont
\setlength{\parskip}{1.5pt}\setlength{\abovedisplayskip}{4pt}\setlength{\belowdisplayskip}{4pt}
The run of \(396\) levels, \(2376\) trits, \(43\) turns, \(387\) pairs \((K,K+9)\) per channel and \(1000\) digits verifies an extensive material instance of the theorem. The run of \(10,000\) digits verifies another instance. The quantifier \(\forall N\) and compatibility \(\rho_{N+1,N}p_{N+1}=p_N\), together with existence of the inverse-limit element \(p_\infty\), prove the infinite decimal output; the finite calculations are reproducible certificates of its projections.

The coordinate \(\alpha\) belongs to the same limit state. Its integer route extends through \(\mathsf T_{\alpha,\infty}\), whereas the simple root of \(E_{21/22}\) supplies a second internal reader. Stability checked to \(10,000\) digits certifies remote carry propagation; the mathematical object expressed by both readers is fixed before any metrological comparison.

\begin{proposition}[Complete prolongation of the integer route of \(\alpha\)]
\label{prop:c27-alpha-prolongacion-completa}
Let \(B=1000\), let \((P_k),(E_k),(\Phi_k)\) be the three block sequences
published by the coinductive state and let \((K_k)\) be the prolonged
dodecaphase sequence. The four series
\[
 p=\sum_{k\geq1}P_kB^{-k},\qquad
 e_0=\sum_{k\geq1}E_kB^{-k},\qquad
 \phi_0=\sum_{k\geq1}\Phi_kB^{-k},\qquad
 \kappa=\sum_{k\geq1}K_kB^{-k}
\]
converge absolutely and determine
\[
 \alpha_A=p+e_0-\phi_0-\kappa.
\]
The carry recurrence is the base-\(B\) writing of this equality.
For each \(M\), a sufficiently deep truncation determines the
first \(M\) digits of the canonical expansion of \(\alpha_A\); the tail
error is bounded by a constant multiplied by \(B^{-N}\).
Consequently, \(\mathsf T_{\alpha,\infty}\) publishes a complete decimal
sequence, and its runs at \(3000\) or \(10\,000\) digits are finite
projections of that output.
\end{proposition}

\begin{proof}
Each block belongs to a bounded integer set, so the four
series are dominated by a geometric series. The partial sum through \(N\)
differs from the complete sum by \(O(B^{-N})\). Euclidean division from
right to left transports that tail exactly through the carry;
by uniqueness of the canonical expansion, every prefix separated from
a cylinder boundary stabilizes as \(N\) increases, and at a boundary
the canonical nonterminal convention applies. This simultaneously defines
all digits, not a sequence of independent numerical targets.
\end{proof}

\begingroup\fontsize{10.2}{11.7}\selectfont
\setlength{\parskip}{2pt}\setlength{\abovedisplayskip}{5pt}\setlength{\belowdisplayskip}{5pt}
\begin{corollary}[Separation between generation and recognition]
\label{cor:c27-generacion-reconocimiento}
Machin’s identities, the exponential series, the Perron ray and the decimal chart are exact realizations of characters previously selected by \(\mathfrak H_9\). Conventional names \(\pi,e,\varphi,\alpha\) and metrological tables recognize the outputs; they do not determine orbits, transports, carries or prefixes.
\end{corollary}
 

\par\endgroup
\par\endgroup
\chapter{Construction of the elliptic generator \texorpdfstring{\(i=\sqrt{-1}\)}{i=sqrt(-1)} and phase geometry}
\begingroup
\renewcommand{\chapter}[2][]{}%
\chapter{Internal construction of \(\sqrt{-1}\): rotation of order \(4\), complex operator, and phase geometry}
\label{chap:p3-final:32:raiz_menos_uno}
\label{chap:p3-83-26-raiz_menos_uno}

\section*{The internal root of \texorpdfstring{\(-1\)}{-1} and the tritic exponential family}

The root of \(-1\) is not introduced as a formal symbol without an
antecedent. The six units modulo nine are coordinatized over
\(\mathbb P^1(\mathbb F_5)\), and the quadratic character of \(\mathbb F_5\) determines the
Paley conference matrix. Its ternary reduction is
\[
 A_W=
 \begin{pmatrix}
 0&1&1&1&1&1\\
 1&0&1&2&2&1\\
 1&1&0&1&2&2\\
 1&2&1&0&1&2\\
 1&2&2&1&0&1\\
 1&1&2&2&1&0
 \end{pmatrix}\in M_6(\mathbb F_3).
\]
The conference identity reduces to
\begin{equation}\label{p3-83-c26-s1:eq:constantes-raiz-interna-menos-uno}
 \boxed{A_W^2=2I_6=-I_6.}
\end{equation}
Therefore, \(A_W\) is a quarter-turn operator in the ternary space.
After extending scalars to \(\mathbb F_9\), it decomposes into the conjugate
eigenspaces of the roots of \(X^2+1\). The common integral source
\[
 \mathcal Q_{\mathbb Z}=\mathbb Z[u]/(u^2+1)
\]
has a finite realization \(u\mapsto A_W\) and a real realization
\(u\mapsto J_{\rm e}\). The source coordinates the quadratic relation; it does not
identify fields of different characteristics.

The three TRIT regimes are carried out by means of operators
\[
 J_{+1}^2=-I,\qquad J_0^2=0,\qquad J_{-1}^2=I.
\]
Separating the even and odd powers of the exponential yields a single family:
\begin{equation}\label{p3-83-c26-s1:eq:constantes-euler-extendido}
 \boxed{
 \exp(tJ_\tau)=C_\tau(t)I+S_\tau(t)J_\tau,
 \qquad \tau\in\{+1,0,-1\},}
\end{equation}
where
\[
 (C_{+1},S_{+1})=(\cos,\sin),\qquad
 (C_0,S_0)=(1,t),\qquad
 (C_{-1},S_{-1})=(\cosh,\sinh).
\]
Its three governing specializations are
\begin{equation}\label{p3-83-c26-s1:eq:constantes-euler-tres-ramas}
 \exp(\pi J_{+1})+I=0,\qquad
 \exp(J_0)=I+J_0,\qquad
 \exp\!\bigl(2(\log\varphi)J_{-1}\bigr)=F_{\rm av}^{\,2}.
\end{equation}
The elliptic branch contains the complex circular closure; the parabolic branch,
threshold transport; the hyperbolic branch, golden self-scaling. Euler’s equation
thus appears as a specialization of a family that coordinates rotation,
boundary and hyperbolic expansion.


\section{Elliptic realization of \(J^2=-I\) through the Hopf fibration and Möbius holonomy}

\subsection{Complex structure and Hopf fibration}

Let \(S^3\subset\mathbb C^2\) and
\[
h(z_0,z_1)=
\bigl(2\Re(z_0\overline z_1),
2\Im(z_0\overline z_1),|z_0|^2-|z_1|^2\bigr)
\in S^2.
\]
On the Hopf torus
\[
T_\eta=\{(z_0,z_1):|z_0|=\cos\eta,
|z_1|=\sin\eta\},
\]
the curves
\[
\gamma_+(t)=(\cos\eta\,e^{it},\sin\eta\,e^{it}),
\]
\[
\gamma_-(t)=(\cos\eta\,e^{it},\sin\eta\,e^{-it})
\]
have distinct behaviors: \(\gamma_+\) is a Hopf fiber and
\(\gamma_-\) projects with degree two. Under stereographic projection, the
homology classes \((1,1)\) and \((1,-1)\) produce the two diagonal
Villarceau families.

The total complex conjugation
\[
K(z_0,z_1)=(\overline z_0,\overline z_1)
\]
preserves each unoriented family and reverses its orientation. It does not interchange
the two families. Partial conjugation can interchange them, but it is neither a
complex-linear map nor the standard global antiunitary. This correction
is essential to avoid manufacturing a particle--antiparticle symmetry from the
toroidal drawing.

\subsection{Double covering and spinor return}

The map \(SU(2)\to SO(3)\) is a double cover. A loop of
\(2\pi\) in \(SO(3)\) can lift to \(-I\) in \(SU(2)\); the loop of
\(4\pi\) recovers \(I\). The sign structure does not turn a circle into a
Möbius band. Obtaining a band requires a real line bundle
whose holonomy character is \(-1\).

Let \(L_\chi\to S^1\) be the associated bundle to
\[
\chi:\pi_1(S^1)\longrightarrow\{\pm1\},
\qquad \chi(1)=-1.
\]
Its total space is a Möbius band. The route \(C_M=-\operatorname{rev}\)
provides the sign candidate; the Villarceau circle provides the loop; the
physical intertwiner must prove that the former is the holonomy of the latter.

\subsection{Toroidal realization of the operator \texorpdfstring{\(B_{\mathrm{mem}}^{\varepsilon}\)}{Bmem} and determination of the Villarceau angle}

With \(B_c=7I+2S\), define the oriented memory operator
\[
B_{\rm mem}^{\varepsilon}=5I+2\varepsilon S,
\qquad\varepsilon=\pm1.
\]
For \(\varepsilon=+1\),
\[
\Spec(B_{\rm mem}^{+})=\{7,3\},
\qquad\det B_{\rm mem}^{+}=21.
\]
Interpreting the values \((7,3)\) as the outer and inner equatorial boundaries
in a common positive chart yields
\[
R+r=7\ell,\qquad R-r=3\ell,
\]
and therefore
\[
R=5\ell,\qquad r=2\ell,\qquad
\sin\beta=\frac rR=\frac25.
\]
The Villarceau angle is
\[
\beta=\arcsin\frac25
=23.5781784782\ldots^\circ.
\]

The algebraic theorem \(\{7,3\}\leftrightarrow(5,2)\) is exact. The geometric
premise interpreting the spectrum as the two boundaries of a torus is a
declared interface. It must not be concealed, nor may the angle be used as a target.

The realization of the full block \(B_c\) retains another reading:
\(R:r=7:2\) and \(\beta=\arcsin(2/7)\). These are not rival values; they belong to
two functors: total block and memory mode. The physical priority of the latter
must be decided by the intertwiner with the dynamics, not by decimal proximity to
another constant.

\subsection{Action of the operator \texorpdfstring{\(J^2=-I\)}{J2=-I} on the conjugate angular coordinates}

HMT’s native angle is a phase \(u\in\mathbb R/\mathbb Z\). The readers
publish
\[
\theta_{\rm rad}=2\pi u,
\qquad
\theta_{\rm deg}=360u.
\]
Radians and degrees are distinct charts of the same circular element. The radian is not more “ontological” by virtue of being dimensionless; its naturality follows from arc length. Degrees provide another absolute coordinate of the full turn. Any relation with \(\alpha^{-1}\), \(\varphi\), or CKM must first be formulated as an invariant quotient or holonomy and only thereafter evaluated in an angular chart.

The exact identity already available
\[
\tanh^2(2\log\varphi)=\frac59
\]
connects the Perron--Fibonacci channel to the spectral quotient of \(B_c\). Its
complement is \(4/9\), the normalized gap. This relation is structural.
It does not imply that every angle constructed with \(\varphi\) is a physical angle.

\subsection{Brouwer, Möbius, and degree}

Brouwer's fixed-point theorem, fractional linear Mobius transformations,
and the Mobius band are three distinct objects. They are related only
after fixing domain, action, and bundle. The elliptic advance and
retardation orbits may be modeled as two sections of a bundle with sign
holonomy; the Mobius transformation describes the conformal action in the chart;
Brouwer's theorem guarantees a fixed point for continuous maps of the disk.
None of these statements replaces the realization diagram.

\begin{remark}[Result and boundary]
The double degree of the Hopf curve, the orientation of total conjugation, and the algebraic construction \(7,3\mapsto5,2\) are proved. The toroidal reading (5{:}2), the Möbius holonomy, and the particle--antiparticle interpretation form a chain of interfaces that must share a single intertwiner. Volume does not force that intertwiner through an angular coincidence.
\end{remark}
  \endgroup

\chapter{Hensel--Witt extension of \texorpdfstring{\(-1/12\)}{-1/12} and modular memory}
\begingroup
\renewcommand{\chapter}[2][]{}%
\chapter{Hensel--Witt derivation of \(-1/12\): compatible prolongation, regularization, and modular memory}
\label{chap:p3-final:33:menos_un_doceavo}
\label{chap:p3-83-27-menos_un_doceavo}

\section{Hensel--Witt domain, Teichmüller maps, and rational target \(-1/12\)}

The NTHW is the data and maps package
\[
\NTHW=
\bigl(
\mathscr R_\pi^{\rm micro},
\Gnine,
\Cdoce,
K,
\mathcal W_{12\to24},
\mathcal I_{6\to8}
\bigr),
\]
with the following compatibilities:
\begin{enumerate}
  \item the Hensel regions share a quotient word and retain
  distinct memory;
  \item \(\Gnine\) transports orientation and produces monodromy;
  \item \(\Cdoce\) organizes the signed observable and the seal \(K\);
  \item the Witt reading converts supports of the same boundary into hexads and
  stars;
  \item the extension \(W_{12}\to W_{24}\) produces the
  hexad--octad interface;
  \item the Archimedean, exceptional and dimensional readers receive
  invariants of the same state.
\end{enumerate}

\begin{figure}[!htbp]
\centering
\resizebox{\textwidth}{!}{\begin{tikzpicture}[font=\small,node distance=1.2cm and 1cm,
  nodebox/.style={draw,rounded corners=2mm,align=center,minimum width=3cm,minimum height=1.05cm},
  arr/.style={-{Stealth[length=2.4mm]},thick}]
  \node[nodebox,draw=hmtblue,fill=hmtlightblue] (hensel) at (-5,2.3) {Henselian cylinders\\carry cocycle and \(\MonNine\)};
  \node[nodebox,draw=hmtblue,fill=hmtlightblue] (regions) at (-5,-.1) {five regions of \(\pi\)\\regions of \(e,\varphi\)};
  \node[nodebox,draw=hmtgold,fill=hmtlightgold,minimum width=3.8cm,minimum height=2.05cm,font=\bfseries] (core) at (0,1.1) {\(\mathscr T_{\rm HW}\)\\Correspondences\\Hensel--Witt};
  \node[nodebox,draw=hmtviolet,fill=hmtlightviolet] (c12) at (0,-1.8) {\(\Cdoce\), \(\Usgn\), \(K\)\\differences \(D_3,D_4,Q\)};
  \node[nodebox,draw=hmtgreen,fill=hmtlightgreen] (witt) at (5,2.3) {Witt \(W_{12}\)\\supports and \(M_{12}\)};
  \node[nodebox,draw=hmtgreen,fill=hmtlightgreen] (oct) at (5,-.1) {incidence inclusion \(6\to8\)\\channels \(90/120\), normalization \(-1/12\)};
  \node[nodebox,draw=hmtred,fill=hmtlightred] (readers) at (5,-2.5) {coordinated readers\\value and symmetry};
  \draw[arr,hmtblue] (hensel)--(core);
  \draw[arr,hmtblue] (regions)--(core);
  \draw[arr,hmtgold] (core)--(c12);
  \draw[arr,hmtgreen] (core)--(witt);
  \draw[arr,hmtgreen] (core)--(oct);
  \draw[arr,hmtred] (c12)--(readers);
  \draw[arr,hmtred] (witt.east)--++(1.9,0)|-(readers.east);
  \draw[arr,hmtred] (oct)--(readers);
  \node[align=center,hmtgray,font=\footnotesize] at (0,-3.45) {Compatibility of cylinders, dodecaphase and incidences on the same integral holonomic state.};
\end{tikzpicture}
 }
\caption{The NTHW as a typed interface. The central mathematical task is to make
the diagram commute, not to accumulate coincidences around a name.}
\label{p3-83-c27-s1:fig:nthw}
\end{figure}

NTHW gives a formal name to the place previously described as the “Rosetta
Stone.” The substitution is conceptual: a stone translates through resemblance; a
transduction specifies input, output, preserved datum and loss of
information.

\section{Teichmüller map from the ternary code to the Witt ring}

The reversible APP boundary selects, with declared gauges, a matrix
\(A_W\) over \(\FF_3\) satisfying
\[
A_WA_W^T=2I\pmod3.
\]
The graphic code
\[
C_W=\{(q,qA_W):q\in\FF_3^6\}
\]
is the extended ternary Golay code. Its weight enumerator is
\[
1+264y^6+440y^9+24y^{12}.
\]
The \(264\) words of weight six, identified up to sign, yield \(132\)
supports. Every five-point subset belongs to a unique hexad:
\[
W_{12}=S(5,6,12).
\]

The field of \(495\) star involutions generates a group of order
\(95\,040\), recognized as \(M_{12}\). An isolated star does not generate the
group. The ordered HMT frame has trivial stabilizer and therefore fixes a
gauge within the action.

\section{Hexad--octad incidence and modular channels \(90/120\)}

Let \(H\) be a hexad and \(O_H\) an octad containing it. Marking one point
and a pair of points of the hexad yields
\[
F_6(H)=H\times\binom H2,
\qquad
|F_6|=6\binom62=90.
\]
If the marked point can traverse the octad:
\[
F_8(H)=O_H\times\binom H2,
\qquad
|F_8|=8\binom62=120.
\]
The inclusion \(H\hookrightarrow O_H\) induces \(F_6\hookrightarrow F_8\).
Its normalized defect is
\[
\boxed{
\frac{90-120}{24\binom62}
=\frac{6-8}{24}
=-\frac1{12}.}
\]
This combinatorial theorem precedes Barbero’s angular reading. The
two objects share the interface \(6\to8\), but are not identical.

\section{4 compatible routes to \texorpdfstring{\(-1/12\)}{-1/12}}

The construction contains four realizations that must be presented together, not
counted as four independent proofs of a single provenance.

\subsection{Angular realization through the elementary increment dodecaphasic}

One tooth of \(\Cdoce\) measures \(30^\circ\). Then
\[
\frac{30}{120}-\frac{30}{90}
=\frac14-\frac13
=-\frac1{12}.
\]
It is the angular shadow of the pattern \(90/120\).

\subsection{Oriented pseudoinverse of the Gram operator on the surviving subspace}

In the certified jump \(60\to81\), two fibers have sizes \(0\) and
\(12\). The Gram operator is
\[
K_{60,81}=\operatorname{diag}(0,12).
\]
On the live subspace, its pseudoinverse is \(1/12\); with the negative
boundary orientation,
\[
E_{\rm surv}=-K_{60,81}^{+}
\]
yields \(-1/12\).

\subsection{Scale extraction by the 2nd triadic difference}

The realization via a triadic scale extractor is the second-difference construction
developed in full in
\cref{sec:c27-segunda-diferencia-triadica}. This first occurrence establishes its
role within the inventory of realizations; the indicated residence preserves
the definitions of \(T_L\), \(a(L)\), \(A_1(L)\), \(C(L)\), the independence
of the residue from \(L\), and the exact status of the normalization.

\subsection{Modular realization through the quotient \(\eta(\tau)/\eta(3\tau)\) and the function \(t_3\)}

The eta function satisfies
\[
\frac{\eta(\tau)}{\eta(3\tau)}
=q^{-1/12}\prod_{n\ge1}(1+q^n+q^{2n})^{-1}.
\]
The exponent comes from \(1/24-3/24=-1/12\). Taking twelve copies:
\[
t_3(q)=\left(\frac{\eta(\tau)}{\eta(3\tau)}\right)^{12}
=q^{-1}-12+54q-76q^2-243q^3+\cdots.
\]
The coefficient \(54\) coincides with the APP defect. The corridor
\[
j=\frac{(t_3+27)(t_3+243)^3}{t_3^3}
\]
connects afterwards with the classical modular branch.

\begin{remark}[What Has Been Proved in the Corridor]
The APP defect \(54\), the eta exponent \(-1/12\), the
hexad--octad defect, and the modular identities once \(t_3\) has been fixed are exact. The
strong assertion is that all of them are natural images of a single TPK operator.
That naturality must be proved by means of the NTHW diagram; it does not
follow merely from repeating the same integers.
\end{remark}


\section{5 algebraic realizations of \texorpdfstring{\(-1/12\)}{-1/12} and separation of their domains}

In addition to the normalized defect of the hexad--octad inclusion,
five operator realizations of the same rational number appear. Their domains are
specified separately; equality of values does not identify the
operators without an additional intertwiner.

\subsection{Difference of dodecaphasic periodicities}

One tooth of \(\Cdoce\) measures \(30^\circ\). Then
\[
\frac{30}{120}-\frac{30}{90}
=\frac14-\frac13
=-\frac1{12}.
\]
It is the oriented difference between the periods of \(90^\circ\) and
\(120^\circ\) in this realization.

\subsection{Pseudoinverse of the Gram Operator on the Live Component}

In the exact jump \(60\to81\), two fibers have sizes \(0\) and
\(12\). The Gram operator is
\[
K_{60,81}=\operatorname{diag}(0,12).
\]
On the nonzero component of its range, the pseudoinverse is \(1/12\); with the negative
boundary orientation,
\[
E_{\rm surv}=-K_{60,81}^{+}
\]
yields \(-1/12\).

\subsection{2nd difference between triadic scales}
\label{sec:c27-segunda-diferencia-triadica}

Define
\[
T_L=\sum_{n=1}^{L-1}n(L-n)=\frac{L(L^2-1)}6,
\qquad
a(L)=-\frac{T_L}{L^2}=-\frac L6+\frac1{6L}.
\]
The first scale subtraction is
\[
A_1(L)=a(L)-\frac{a(3L)}3=\frac4{27L},
\]
and the second
\[
C(L)=LA_1(L)-\frac{3L}{3}A_1(3L)=\frac8{81}.
\]
The residue \(8/81\) is independent of \(L\). To obtain
\(-1/12\), the declared normalization is applied
\[
R(L)=-\frac{27}{32}C(L)=-\frac1{12}.
\]
The calculation derives \(8/81\); it does not by itself derive the factor \(-27/32\).

\subsection{Modular chain}

The eta function satisfies
\[
\frac{\eta(\tau)}{\eta(3\tau)}
=q^{-1/12}\prod_{n\ge1}(1+q^n+q^{2n})^{-1}.
\]
The exponent comes from \(1/24-3/24=-1/12\). Taking twelve copies:
\[
t_3(q)=\left(\frac{\eta(\tau)}{\eta(3\tau)}\right)^{12}
=q^{-1}-12+54q-76q^2-243q^3+\cdots.
\]
The coefficient \(54\) coincides with the APP defect\@. The rational identity
\[
j=\frac{(t_3+27)(t_3+243)^3}{t_3^3}
\]
establishes the connection with the classical modular branch.

\Needspace{32\baselineskip}
\subsection{Defect of cyclic difference projectors}

Let \(S\) be the cyclic shift on \(\QQ^{12}\) and define
\[
 D_r=I-S^r.
\]
The kernel of \(D_r\) consists of the vectors constant on each orbit
of \(S^r\); therefore,
\[
 \dim\ker D_r=\gcd(12,r).
\]
Let \(\Pi_{\ker D_3}\) and \(\Pi_{\ker D_4}\) be the rational projectors onto
\(\ker D_3\) and \(\ker D_4\), and let
\(\tau(T)=\operatorname{Tr}(T)/12\) be the normalized trace. Then
\[
 \boxed{
 \tau(\Pi_{\ker D_3}-\Pi_{\ker D_4})
 =\frac{\operatorname{rank}\Pi_{\ker D_3}
       -\operatorname{rank}\Pi_{\ker D_4}}{12}
 =\frac{3-4}{12}
 =-\frac1{12}.}
\]
This realization expresses the rational as the transverse defect of steps
three and four in the twelve-phase cycle.

\begin{remark}[Domains of the realizations]
The APP defect \(54\), the eta exponent \(-1/12\), the combinatorial
defect of the hexad--octad lift, and the
modular identities once \(t_3\) has been fixed are exact. The value
of the pseudoinverse, the normalized extractor, and the projector defect are likewise exact.
Each equality preserves its domain and its corresponding map.
\end{remark}
  \endgroup

\chapter{Quadratic characters and spectral constants: Catalan, Apéry, Euler--Mascheroni and Stieltjes}
\begingroup
\renewcommand{\chapter}[2][]{}%
\chapter{Quadratic characters and spectral constants: Catalan, Apéry and the Euler--Stieltjes tower}
\label{chap:p3-final:41:catalan_apery_euler}
\label{ch:spectral}

\section{Operatorial representation of Dirichlet characters}

The constants in this chapter are grouped because they share a spectral
construction independent of their decimal expansions and free of target
values. Let

\[
Ne_n=ne_n,\qquad n\in\N,
\]

in \(\ell^2(\N)\). A residual character \(\chi\) defines the diagonal operator \(Q_\chi e_n=\chi(n)e_n\); whenever the power is trace class,

\begin{equation}
\Tr(N^{-s}Q_\chi)=L(s,\chi).
\label{eq:L-trace}
\end{equation}

The equation preserves the modulus, orientation, and provenance of each
character; therefore, it does not identify distinct residual objects.

\section{HMT construction of the residue characters}

The preceding diagonal operator is not the initial algebraic datum. The
selection chain is

\[
\begin{aligned}
\APP&\longrightarrow\TRIT\longrightarrow\TPK
\longrightarrow\{\text{residue modulo }3,\ C_P\}\\
&\longrightarrow\{\chi_{-3},\ J=C_P\bmod3\}
\longrightarrow\{\chi_{-3},\ J^2=-I,\chi_{-4}\}\\
&\longrightarrow Q_3,Q_4.
\end{aligned}
\]

The triadic residual functional determines the difference between the two nonzero
oriented classes. The incidence sector uses the symmetric Paley conference
matrix of order six,
\[
C_P^{\mathsf T}=C_P,\qquad C_PC_P^{\mathsf T}=5I_6.
\]
Its reduction \(J=C_P\bmod3\) satisfies
\[
J^2=2I_6=-I_6
\quad\text{in }\operatorname{End}_{\mathbb F_3}(\mathbb F_3^6).
\]
The torsional quarter then determines the cyclic phase
\(1,\ii,-1,-\ii\). Applying the Chinese remainder theorem preserves
both residual data and produces \(Q_{12}=Q_3Q_4\). The trace
\(\Tr(N^{-s}Q)\) is evaluated at the terminal stage; neither the Dirichlet series
nor its decimal expression enters the selection of the character.

The construction is made explicit by the tables
\begin{equation}
\chi_{-3}(n)=
\begin{cases}
0,&3\mid n,\\
1,&n\equiv1\pmod3,\\
-1,&n\equiv2\pmod3,
\end{cases}
\qquad
\chi_{-4}(n)=
\begin{cases}
0,&2\mid n,\\
1,&n\equiv1\pmod4,\\
-1,&n\equiv3\pmod4.
\end{cases}
\label{eq:characters34-tables}
\end{equation}
En \(\ell^2(\N)\),
\[
Q_3e_n=\chi_{-3}(n)e_n,
\qquad Q_4e_n=\chi_{-4}(n)e_n.
\]
As they are diagonals, they commute; the Chinese remainder theorem gives
\begin{equation}
(Q_3Q_4)e_n=\chi_{-3}(n)\chi_{-4}(n)e_n
=\chi_{12}(n)e_n.
\label{eq:q12-product-proof}
\end{equation}
The product simultaneously preserves the triadic residue and the quarter of
torsion. Consequently, the character modulo twelve is determined
before performing any evaluation of its values \(L\).

\begin{remark}[Negative control]
The chain is invalidated if \(Q_3\) does not distinguish the two oriented classes,
if \(J^2\ne-I\), if \(Q_3Q_4\) does not have the declared CRT table, or if the
operator is selected from the value of the constant.
\end{remark}

\section{Internal complex structure and Catalan's constant}

The internal complex structure satisfies

\begin{equation}
J^2=-I.
\label{eq:J-square}
\end{equation}

The equality \(J^2=-I\) defines a representation
\(\rho:C_4\to\operatorname{Aut}_{\mathbb F_3}(\mathbb F_3^6)\),
\(\rho(j)=J\). To realize it spectrally, take the unitary character
\(\upsilon:C_4\to U(1)\), \(\upsilon(j)=\ii\), and the phase map
\(n\mapsto j^n\). On \(\ell^2(\N)\), the induced diagonal representation
is
\[
U_4e_n=\upsilon(j^n)e_n=\ii^ne_n.
\]
Its oriented part is
\[
Q_4=\frac{U_4-U_4^*}{2\ii},\qquad
Q_4e_n=\sin\!\left(\frac{\pi n}{2}\right)e_n
=\chi_{-4}(n)e_n.
\]
This is the explicit lift between the finite generator of order four and the
residual operator; \(\mathbb F_9\) is not identified with \(\mathbb C\).

On odd integers, the same character can be written

\[
\chi_J(n)=\ii^{n-1}=\chi_{-4}(n).
\]

Equivalently,

\begin{equation}
Q_4=\frac{U_4-U_4^*}{2\ii},
\qquad Q_4e_n=\chi_{-4}(n)e_n.
\label{eq:Q4}
\end{equation}

The Catalan constant appears as

\begin{equation}
G=\Tr(N^{-2}Q_4)=L(2,\chi_{-4})
=\sum_{n\ge0}\frac{(-1)^n}{(2n+1)^2}.
\label{eq:catalan}
\end{equation}

The unitary lift of the torsion quarter generates exactly the
character whose quadratic moment is Catalan's constant.

\begin{proof}[Operational proof]
The operator \(N^{-2}Q_4\) is trace-class because
\(\sum_{n\ge1}|\chi_{-4}(n)|n^{-2}<\infty\). Its trace is therefore the
sum of the diagonal elements. The table \cref{eq:characters34-tables}
annuls the even pairs and alternates the odd ones:
\[
\Tr(N^{-2}Q_4)
=1-\frac1{3^2}+\frac1{5^2}-\frac1{7^2}+\cdots.
\]
The identity with \(L(2,\chi_{-4})\) is the definition of the
Dirichlet series; the equality with \(\operatorname{Im}\operatorname{Li}_2(\ii)\)
is obtained by taking the imaginary part of
\(\sum_{n\ge1}\ii^n/n^2\).
\end{proof}

The terminal value is
\[
G=0.9159655941772190150546035149\ldots.
\]
The determining structure is the identity \(J^2=-I\), which induces
exactly the alternating sign of the odd classes; the decimal
expansion is its terminal evaluation.

\section{Residual composition of the triadic--quaternary type}

The residual functional modulo three provides \(\chi_{-3}\). Its CRT
product with \(\chi_{-4}\) defines

\begin{equation}
\chi_{12}=\chi_{-3}\chi_{-4}.
\label{eq:chi12}
\end{equation}

On the units modulo twelve,

\[
\chi_{12}(1)=1,\quad
\chi_{12}(5)=-1,\quad
\chi_{12}(7)=-1,\quad
\chi_{12}(11)=1,
\]

and vanishes outside of them. Two exact evaluations are

\begin{align}
L(1,\chi_{12})&=\frac{\log(2+\sqrt3)}{\sqrt3},
\label{eq:L1chi12}\\
L(2,\chi_{12})&=\frac{\pi^2}{6\sqrt3}.
\label{eq:L2chi12}
\end{align}

The first contains the hyperbolic unit

\begin{equation}
\sqrt3L(1,\chi_{12})=\arcosh2=\log(2+\sqrt3).
\label{eq:pell2}
\end{equation}

Together with \(\log(3+2\sqrt2)=2\log(1+\sqrt2)\) and \(\log(30+\sqrt{899})\), this identity determines a
family of Pell scales whose members retain their respective algebraic
provenances.

Certified terminal evaluations are

\[
L(1,\chi_{12})=0.7603459963009463475310942549\ldots,
\]

\[
L(2,\chi_{12})=0.9497031262940093952634984917\ldots.
\]

\begin{proof}[Evaluations dodecaphasic]
Grouping by residue classes,
\begin{align*}
L(1,\chi_{12})
&=\sum_{n\ge0}\left(
\frac1{12n+1}-\frac1{12n+5}
-\frac1{12n+7}+\frac1{12n+11}\right)\\
&=\int_0^1
\frac{1-x^4-x^6+x^{10}}{1-x^{12}}\,dx
=\int_0^1\frac{1-x^4}{1+x^6}\,dx.
\end{align*}
A real decomposition into the quadratic factors of \(1+x^6\)
produces
\(\log(2+\sqrt3)/\sqrt3\). For the second moment, if
\(\psi_1\) is the trigamma function,
\begin{align*}
L(2,\chi_{12})
&=\frac1{144}\bigl[
\psi_1(1/12)-\psi_1(5/12)
-\psi_1(7/12)+\psi_1(11/12)\bigr]\\
&=\frac{\pi^2}{144}\left(
\csc^2\frac\pi{12}-\csc^2\frac{5\pi}{12}\right)
=\frac{\pi^2}{6\sqrt3},
\end{align*}
where has been used
\(\psi_1(z)+\psi_1(1-z)=\pi^2\csc^2(\pi z)\). The two proofs use
the CRT table and not a target decimal value.
\end{proof}

\section{Euler and Stieltjes Constants and the Theta--Mellin Transform}

Let
\[
\vartheta(t)=\sum_{n\in\Z}e^{-\pi n^2t},\qquad t>0.
\]
For \(\Re s>1\), the theta--Mellin functional recovers the convergent identity

\[
\pi^{-s/2}\Gamma\!\left(\frac{s}{2}\right)\zeta(s)
=\frac12\int_0^\infty
\bigl(\vartheta(t)-1\bigr)t^{s/2-1}\,dt.
\]

The meromorphic continuation is obtained by splitting the integral at \(t=1\), using
\(\vartheta(t)=t^{-1/2}\vartheta(1/t)\), and separating the polar terms. Its expansion at \(s=1\),

\begin{equation}
\zeta(s)=\frac1{s-1}+\sum_{j\ge0}
\frac{(-1)^j\gamma_j}{j!}(s-1)^j,
\label{eq:stieltjes}
\end{equation}

situates \(\gamma=\gamma_0\) and the Stieltjes constants in a single
family of renormalized moments. This organization expresses an
analytical common provenance, without attributing an interpretation that is physically identical to all its levels.

 The Catalan constant often appears as an alternating series:\par

\[
G
=
1-\frac1{3^2}+\frac1{5^2}-\frac1{7^2}+\cdots.
\]

That formula gives its value, and the HMT genealogy explains why this specific alternation exists.\par

In HMT, the origin is in the torsion chamber.\par

There exists an internal operator \(J\) such that\par

\[
J^2=-I.
\]

The root of \(-1\) emerges from two successive applications of a
quarter turn. The complex structure is the algebraic memory of an internal
rotation.\par

That turn generates a character of order four. When it is published on the positive integers, it produces exactly the residual pattern that alternates among \(+1\), \(0\), \(-1\), \(0\).\par

The Catalan constant then appears as a spectral trace:\par

\[
G
=
L(2,\chi_{-4})
=
\operatorname{Tr}(N^{-2}Q_4).
\]

This changes its meaning. Catalan's constant is simultaneously the sum of a particular series and the spectral response of the integers to a quarter of torsion.\par

Each term of the series records how the internal rotation acts on a residue class. The alternation is the arithmetic shadow of the rotation.\par

The triadic structure has its own character modulo \(3\). Combining it with the character of order four, the Chinese remainder theorem produces a character modulo \(12\):\par

\[
\chi_{12}
=
\chi_{-3}\chi_{-4}.
\]

The module \(12\) appears as the minimal common space in which triadic orientation and the quarter-torsion can coexist.\par

The dodecaphase is, in that sense, a reconciliation of \(3\) and \(4\).\par

Of the combined character emerges\par

\[
L(1,\chi_{12})
=
\frac{\log(2+\sqrt3)}{\sqrt3},
\]

y\par

\[
L(2,\chi_{12})
=
\frac{\pi^2}{6\sqrt3}.
\]

La cantidad\par

\[
2+\sqrt3
\]

is the fundamental unit of the real field \(\mathbb Q(\sqrt3)\). Its logarithm represents a hyperbolic shift. Once again, a finite residue structure opens an infinite multiplicative dynamics.\par

Here rotation, torsion, arithmetic character, Pell unit, and hyperbolic displacement meet.\par

The Euler–Kronecker constant of the field \(\mathbb Q(\sqrt3)\) prolongs this structure. It measures the regular term that remains after separating the principal pole of the field's zeta function. In the HMT genealogy, this term constitutes the regular invariant of fine arithmetic memory that remains after removing the universal divergence.\par

The Bendersky–Glaisher family extends the mechanism into a tower. The spectral derivatives of zeta generate that tower. Apery's constant and other constants associated with odd values are organized as members of the same functional family.\par

The novel contribution lies in the genealogy that unifies the classical constants:\par

\begin{itemize}[leftmargin=*,itemsep=0.35em]
\item the quarter of torsion produces the complex structure;
\item the residue triadico produces the structure modulo \(3\);
\item their composition produces the character modulo \(12\);
\item the character generates values \(L\);
\item the values \(L\) open the quadratic field and its spectral constants.
\end{itemize}

Geometric torsion has become arithmetic while preserving entirely its orientation.\par
  \endgroup

\chapter{Euler--Kronecker, Meissel--Mertens and prime distributions in the double reading of the sheets}
\begingroup
\renewcommand{\chapter}[2][]{}%
\chapter{Euler--Kronecker constants and prime distributions: dodecaphasic characters, Meissel--Mertens, and twin-prime series}
\label{chap:p3-final:42:euler_kronecker_primos}
\section{Euler--Kronecker constant of the character of dodecaphase type}

The factorization of Dedekind

\begin{equation}
\zeta_{\mathbb Q(\sqrt3)}(s)=\zeta(s)L(s,\chi_{12})
\label{eq:dedekind-sqrt3}
\end{equation}

gives, upon comparing the constant terms in \(s=1\),

\begin{equation}
\gamma_{\mathbb Q(\sqrt3)}
=\gamma+\frac{L'(1,\chi_{12})}{L(1,\chi_{12})}
=1.0539656082484825800\ldots.
\label{eq:euler-kronecker}
\end{equation}

This constant represents the normalized finite term at the pole of the
zeta function of the real dodecaphase field. It is not a variant of
Euler's \(\gamma\), but the arithmetic correction induced by
\(\chi_{12}\).

In fact, writing \(u=s-1\),
\[
\zeta(s)=u^{-1}+\gamma+O(u),
\qquad
L(s,\chi_{12})=L_1+L_1'u+O(u^2),
\]
we obtain
\begin{equation}
\zeta_{\mathbb Q(\sqrt3)}(s)
=\frac{L_1}{u}+\bigl(\gamma L_1+L_1'\bigr)+O(u).
\label{eq:ek-laurent-derivation}
\end{equation}
The Euler--Kronecker constant is the quotient of the finite term by
the residue; dividing \cref{eq:ek-laurent-derivation} by \(L_1\) proves
\cref{eq:euler-kronecker}. This derivation simultaneously explains the
residue
\[
\operatorname*{Res}_{s=1}\zeta_{\mathbb Q(\sqrt3)}(s)
=\frac{\log(2+\sqrt3)}{\sqrt3}
\]
and the finite correction. The numerical value is certified by evaluating
\(L(1,\chi_{12})\) and \(L'(1,\chi_{12})\) with the same residue table and a
tail bound; it is not imported from a table of quadratic fields.

\section{Arithmetic constants associated with prime numbers}

Once the irreducible arithmetic cycles have been selected, two
distinct families are obtained:

\begin{align}
B_1&=\lim_{x\to\infty}\left(
\sum_{p\le x}\frac1p-\log\log x
\right),
\label{eq:meissel-mertens}\\
C_2&=\prod_{p>2}\frac{p(p-2)}{(p-1)^2}.
\label{eq:twin-prime-constant}
\end{align}

\(B_1\) regularizes a local sum; \(C_2\) is a multiplicative density of twin-prime pairs. Both require tail control. They are not deduced from one another by sharing the primes.

  \endgroup

\chapter{Regular products and the Bendersky--Glaisher--Kinkelin system}
\begingroup
\renewcommand{\chapter}[2][]{}%
\chapter{The Bendersky--Glaisher--Kinkelin family: regularized products, derivatives of the zeta function, and normalization hierarchies}
\label{chap:p3-final:43:bendersky_glaisher}
\section{Generation of the Bendersky--Glaisher--Kinkelin family via zeta regularization}

Let \(\mathcal Z_3(s):=\zeta(s)\) be the scalar realization of the triadic
zeta functional in the normalization of this volume,
\(H_k=\sum_{j=1}^k j^{-1}\) (with \(H_0=0\)), and let \(B_{k+1}\) be the Bernoulli number
under convention \(B_1=-1/2\). We define

\begin{equation}
A_k=\exp\!\left(
\frac{H_kB_{k+1}}{k+1}-\mathcal Z_3'(-k)
\right),\qquad k\ge0.
\label{eq:bendersky}
\end{equation}

The first levels include

\[
A_0=\sqrt{2\pi},
\]

and the Glaisher--Kinkelin constant in \(A_1\). For \(n\ge1\),

\begin{equation}
\zeta(2n+1)=(-1)^{n+1}
\frac{2(2\pi)^{2n}}{(2n)!}\log A_{2n}.
\label{eq:odd-zeta-bendersky}
\end{equation}

Apéry's constant, \(\zeta(3)\), is the first nontrivial even level.
Its presence in the photon gas of \cref{ch:metrology} constitutes a
subsequent physical realization of this spectral value.

The first part of the tower can be seen in a synoptic form:
\begin{center}
\small
\begin{tabular}{@{}cll@{}}
\toprule
\(k\) & identidad estructural & terminal value\\
\midrule
0 & \(A_0=\exp[-\zeta'(0)]=\sqrt{2\pi}\)
  & \(2.5066282746310005024\ldots\)\\
1 & \(\log A_1=1/12-\zeta'(-1)\)
  & \(1.2824271291006226369\ldots\)\\
2 & \(\log A_2=-\zeta'(-2)=\zeta(3)/(4\pi^2)\)
  & \(1.0309167521973921\ldots\)\\
3 & \(\log A_3=-11/720-\zeta'(-3)\)
  & \(0.9795555269428446\ldots\)\\
4 & \(\log A_4=-\zeta'(-4)\)
  & \(0.9920479745250403\ldots\)\\
\bottomrule
\end{tabular}
\end{center}

\begin{proof}[Relation to the odd values]
The functional equation of \(\zeta\), developed in \(s=-2n\), gives
\begin{equation}
\zeta'(-2n)=(-1)^n
\frac{(2n)!}{2(2\pi)^{2n}}\zeta(2n+1).
\label{eq:zeta-derivative-even-negative}
\end{equation}
Since \(B_{2n+1}=0\) for \(n\ge1\), the definition
\cref{eq:bendersky} reduces \(\log A_{2n}=-\zeta'(-2n)\). Solving for
\(\zeta(2n+1)\) proves \cref{eq:odd-zeta-bendersky}. The index \(2n\)
preserves the derivation level: it is not an editorial label.
\end{proof}

  \endgroup

\chapter{Gauss dynamics and the Khinchin, Lévy and Lochs constants}
\begingroup
\renewcommand{\chapter}[2][]{}%
\chapter{Gauss dynamics and the Khinchin, Lévy, and Lochs constants: cofinality of cylinders, transfer operator, and Gauss--Kuzmin--Wirsing spectrum}
\label{chap:p3-final:44:gauss_khinchin_levy_lochs}
\label{ch:gauss}

\section{Cofinality between HMT cylinders and continued fraction cylinders}

An irrational HMT section determines a nested family of compact cylinders. Its Archimedean evaluation is a point \(x\in(0,1)\setminus\mathbb Q\). The continued fraction of \(x\) determines classical cylinders

\[
I(a_1,\ldots,a_n).
\]

\begin{proposition}[Cofinality]
For every finite cylinder \(I(a_1,\ldots,a_n)\) containing the evaluation \(x\), there is an HMT depth \(m(n)\) such that the HMT cylinder at level \(m\ge m(n)\) is contained in \(I(a_1,\ldots,a_n)\).
\end{proposition}

\begin{proof}
Given \(n\), the continued-fraction cylinder \(I(a_1,\ldots,a_n)\) is an open neighborhood of the irrational \(x\). There then exists \(\varepsilon_n>0\) with \((x-\varepsilon_n,x+\varepsilon_n)\subset I(a_1,\ldots,a_n)\). The HMT cylinders \(C_m\) contain \(x\), are nested, and satisfy \(\operatorname{diam}C_m\to0\). Choose \(m(n)\) with \(\operatorname{diam}C_{m(n)}<\varepsilon_n\). Since \(x\in C_{m(n)}\), every point of that cylinder lies from \(x\) at a distance of less than \(\varepsilon_n\), hence \(C_{m(n)}\subset I(a_1,\ldots,a_n)\). Nesting gives the inclusion for every subsequent level.
\end{proof}

The proof uses compactness, nesting, and uniqueness of
evaluation. The result establishes the transport of approximations, but
does not imply conjugation of the TPK memory shift with the Gauss
map on all deep fibers.

The complete genealogy of this realization is
\[
\APP\to\TRIT\to\TPK
\to (C_m,\rho_{m+1,m})_{m\ge0}
\to x=\bigcap_m C_m
\to \bigl(I(a_1,\ldots,a_n)\bigr)_{n\ge1}.
\]
The HMT cylinders preserve route, phase, and carry; the terminal map projects
these coordinates onto continued-fraction prefixes. Cofinality
proves that this projection converges to the same point. A dynamical
intertwiner would also require a lift \(\widetilde T_G\) with
commutative squares over the valuation fibers, a property not
deduced from cofinality.

\section{Gauss measure and Khinchin and Levy observables}

The Gauss map \cite{Khinchin,IosifescuKraaikamp}

\[
T_G(x)=\frac1x-\left\lfloor\frac1x\right\rfloor
\]

preserve

\begin{equation}
d\mu_G(x)=\frac{dx}{\log2(1+x)}.
\label{eq:gauss-measure}
\end{equation}

Khinchin's constant is defined by

\begin{equation}
\log K=\int_0^1\log a_1(x)\,d\mu_G(x).
\label{eq:khinchin}
\end{equation}

The distribution of the first digit is computed directly:
\begin{equation}
\mu_G(a_1=k)
=\frac1{\log2}\int_{1/(k+1)}^{1/k}\frac{dx}{1+x}
=\log_2\!\left(\frac{(k+1)^2}{k(k+2)}\right).
\label{eq:gauss-digit-law}
\end{equation}
By the ergodic theorem,
\begin{equation}
K=\prod_{k\ge1}
k^{\,\log_2((k+1)^2/[k(k+2)])}
=2.685452001065306445309714835\ldots
\label{eq:khinchin-product}
\end{equation}
for \(\mu_G\)-almost every section. The product is not applied to each point:
the section \(\varphi^{-1}\), for example, has all its digits equal
to one and is an immediate falsification criterion for a pointwise reading. A
certification of \cref{eq:khinchin-product} truncates the sum of logarithms and
bounds the tail by the estimate
\[
\mu_G(a_1=k)=O(k^{-2}),
\qquad
\sum_{k>N}\mu_G(a_1=k)\log k
=O\!\left(\frac{\log N+1}{N}\right).
\]

Lévy's constant is

\begin{equation}
L=-\int_0^1\log x\,d\mu_G(x)
=\frac{\pi^2}{12\log2}.
\label{eq:levy}
\end{equation}

Indeed,

\[
\frac1{1+x}=\sum_{n\ge0}(-1)^nx^n,\qquad
\int_0^1(-\log x)x^n\,dx=\frac1{(n+1)^2},
\]

so that the integral is \(\eta(2)/\log2=\pi^2/(12\log2)\).

Khinchin's constant determines the logarithmic mean of the
partial quotients, whereas Lévy's constant determines the
exponential growth rate of the denominators. They are distinct observables
of the same measured system.

\section{Entropy (metric), Lyapunov exponent, and Lochs's constant}

Since

\[
\log|T_G'(x)|=-2\log x,
\]

the Lyapunov exponent and the entropy metric are

\begin{equation}
\chi_G=h_G=2L=\frac{\pi^2}{6\log2}.
\label{eq:gauss-entropy}
\end{equation}

The Lochs ratio between continued fractions and a base-\(b\) expansion is

\begin{equation}
\Lambda_b=\frac{\log b}{h_G}
=\frac{6\log2\log b}{\pi^2}.
\label{eq:lochs}
\end{equation}

In particular,

\begin{align*}
\Lambda_{10}&=0.9702701143920339257402560192\ldots,\\
\Lambda_{729}&=\frac{36\log2\log3}{\pi^2}
=2.777618966374305777705782976\ldots.
\end{align*}

The basis \(729=3^6\) is relevant to the ambient space \(\mathbb F_3^6\), but \cref{eq:lochs} remains an almost-sure theorem of \(\mu_G\), not a pointwise equality for every HMT route.

The value \(\Lambda_{729}\) represents the asymptotic typical rate of
partial quotients obtained per complete block of six trits. It does not identify the
ambient set of \(729\) words with the other objects of cardinality \(729\) in the
corpus. The equality is an information rate between partitions, not a
bijection between states.

\section{Gauss--Kuzmin--Wirsing Transfer Operator}

The Gauss--Kuzmin--Wirsing operator acts by

\begin{equation}
(\mathcal L f)(x)=\sum_{n\ge1}\frac1{(x+n)^2}
f\!\left(\frac1{x+n}\right).
\label{eq:gkw-operator}
\end{equation}

Its dominant eigenvalue corresponds to the invariant density. The GKW constant is the modulus of the subdominant eigenvalue on the zero-mean subspace. An HMT-compatible scheme uses partitions of \(9^m\) cylinders, transfer matrices with intervals, and a tail bound.

To promote the calculation, the following must be certified: isolation of the eigenvalue, invariance of the subspace, truncation error, convergence of spectral projectors, and absence of another eigenvalue of equal modulus. The known decimal value cannot be introduced as the center of the initial interval.

A joint deformation contains Khinchin and Lévy in a single object:
\begin{equation}
(\mathcal L_{s,t}f)(x)
=\sum_{n\ge1}n^t(n+x)^{-2s}
 f\!\left(\frac1{n+x}\right),
\qquad
P(s,t)=\log r(\mathcal L_{s,t}).
\label{eq:gauss-pressure-operator}
\end{equation}
At the point \((s,t)=(1,0)\), the spectral perturbation gives
\begin{equation}
\partial_tP(1,0)=\log K,
\qquad
-\partial_sP(1,0)=\chi_G=2L.
\label{eq:pressure-derivatives}
\end{equation}
The Hessian of \(P\) determines the asymptotic variances and covariances of
\(\log a_1\) and \(-2\log x\). Therefore, the Khinchin and
Levy constants are two transverse derivatives of the same pressure function.

The numerical value of the subdominant eigenvalue is
\[
\lambda_2\simeq-0.3036630028987326586,\qquad
\kappa_{\rm GKW}=|\lambda_2|,\qquad
-\log\kappa_{\rm GKW}\simeq1.19183673556055.
\]
In this volume, those decimals constitute reference values for
comparison and not a spectral certificate.
The proposed HMT proof uses Galerkin or Ulam on the cofinal cylinders of
depth \(9^m\), a Lasota--Yorke inequality, interval arithmetic to
isolate \(\lambda_2\), and a rigorous bound on the tail. An
interval containing two eigenvalues of the same modulus refutes the
promotion, although it does not affect Khinchin, Levy, or Lochs.

 \section{Cofinal certificate of the Gauss--Kuzmin--Wirsing spectrum over the nonadic hierarchy}
\label{sec:p3-final:cierre-gkw}

For Gauss--Kuzmin--Wirsing, let \(P_N\) be the certified family of
truncations on the declared function space, and define
\[
 \Pi_m=P_{9^m}.
\]
The hierarchy \((\Pi_m)\) is a cofinal subsequence; uniform bounds, projector convergence and absence of spectral pollution are inherited from the system \((P_N)\). A different cylindrical discretization must supply its intertwiner with \((P_{9^m})\).
\endgroup

\chapter{Unimodal criticality, Feigenbaum constants and hyperbolic renormalization}
\begingroup
\renewcommand{\chapter}[2][]{}%
\chapter{Unimodal criticality of the dodecaphasic family: exact analytic profile, Feigenbaum approximants and reduction of universality to the hyperbolic certificate}
\label{chap:p3-final:45:feigenbaum}
\label{mab:rm:ch:chaos}

\section{Intrinsic construction of the dodecaphasic family}

The Feigenbaum constants characterize the universality of period
doubling. Their derivation in HMT first requires constructing
an analytic unimodal family, proving its quadratic criticality, and
then studying the action of the renormalization operator.

The initial map does not follow from an external harmonic choice. The return
of \(\TPK\) determines the oriented dodecaphase wheel
\(U_{12}^{\rm sign}\). In the uniform sector \(\eta=0\), the normalized
trace of its twelve sectors fixes \(w_m=1/12\), and the phase functional
\begin{equation}
\mathcal P_{\rm crit}:U_{12}^{\rm sign}\longrightarrow\R/2\pi\Z,
\qquad
m\longmapsto \frac{(30m-10)\pi}{180}
=\pi\frac{3m-1}{18}
\label{mab:rm:eq:critical-phase-reader}
\end{equation}
preserves order and orientation. The causal composition is
\[
\APP\to\TRIT\to\TPK\to U_{12}^{\rm sign}
\xrightarrow{\mathcal P_{\rm crit}}
\{(\phi_m,w_m)\}_{m=1}^{12}
\to u_0\to f_\lambda.
\]
Every modification of \(\mathcal P_{\rm crit}\), of the weights, or of the order of
the phases alters the family and constitutes a criterion of refutation. These
choices must be fixed before the contrast with \(\delta_F\).

The resulting dodecaphasic family starts from

\[
a_m=\frac{3m-1}{18},\qquad
\phi_m=\pi a_m,\qquad m=1,\ldots,12,
\]

and of the second moment

\[
D_0=\frac1{12}\sum_{m=1}^{12}a_m^2.
\]

The identity

\begin{equation}
\sum_{m=1}^{12}(3m-1)^2=5394=6\cdot899
\label{mab:rm:eq:chaos-sum2}
\end{equation}

yields

\[
D_0=\frac{899}{648},\qquad
s_0=\frac{36}{\pi\sqrt{899}}.
\]

Therefore,

\begin{equation}
s_0\phi_m=\frac{2(3m-1)}{\sqrt{899}}.
\label{mab:rm:eq:pi-cancellation}
\end{equation}

The factor \(\pi\) cancels before the iterations are performed. The family does not
depend on \(\delta_F\) or \(\alpha_F\) as input data.

Cancellation is a structural property of
normalization. Before normalization, the second angular moment
is
\[
\frac1{12}\sum_{m=1}^{12}\phi_m^2
=\pi^2D_0.
\]
The condition of quadratic criticality fixes
\(s_0^2\pi^2D_0=2\), and therefore each normalized phase depends only on the
integer arithmetic of \(3m-1\). The functional preserves the angular provenance of
the wheel, but the iteration receives an intrinsic real family. This is
precisely what makes it possible to distinguish a derivation from a calibration
by \(\pi\).

\section{Exact analytic representation and closed trigonometric expression}

We define

\begin{equation}
u_0(x)=1-\frac1{12}\sum_{m=1}^{12}
\cos\!\left(\frac{2(3m-1)}{\sqrt{899}}x\right).
\label{mab:rm:eq:chaos-profile}
\end{equation}

The development at the origin is

\begin{equation}
u_0(x)=x^2-\frac{715769}{2424603}x^4+O(x^6).
\label{mab:rm:eq:chaos-taylor}
\end{equation}

The finite sum admits the following closed expression. If \(y=x/\sqrt{899}\),

\begin{equation}
u_0(x)=1-\frac{\cos(37y)\sin(36y)}{12\sin(3y)}
=1-\frac{\sin(73y)-\sin y}{24\sin(3y)}.
\label{mab:rm:eq:chaos-closed}
\end{equation}

The apparent singularity at \(y=0\) is removable and reproduces \cref{mab:rm:eq:chaos-taylor}.

\begin{proof}[Derivation of the closed form]
The twelve arguments form an arithmetic progression:
\[
2(3m-1)y=(6m-2)y,\qquad m=1,\ldots,12.
\]
The sum of cosines over an arithmetic progression satisfies
\[
\sum_{m=1}^{N}\cos(a+md)
=\frac{\sin(Nd/2)}{\sin(d/2)}
 \cos\!\left(a+\frac{N+1}{2}d\right).
\]
Taking \(N=12\), \(d=6y\) and \(a=-2y\) yields
\[
\sum_{m=1}^{12}\cos((6m-2)y)
=\frac{\sin36y}{\sin3y}\cos37y.
\]
The product--to--sum identity
\(2\cos37y\sin36y=\sin73y-\sin y\) gives the second expression for
\cref{mab:rm:eq:chaos-closed}.
\end{proof}

The quartic coefficient is not fitted either. The integer sums
\begin{equation}
\sum_{m=1}^{12}(3m-1)^2=5394,\qquad
\sum_{m=1}^{12}(3m-1)^4=4294614
\label{mab:rm:eq:chaos-moments24}
\end{equation}
inserted in the development of \(\cos z\) produce exactly
\[
\frac1{24\cdot12}\sum_m
\left(\frac{2(3m-1)}{\sqrt{899}}\right)^4
=\frac{715769}{2424603}.
\]
Thus, both the quadratic normalization and the first nonlinear
correction are finite moments of the same wheel.

\section{Membership in the quadratic unimodal analytic class}

Consider

\[
f_\lambda(x)=1-\lambda u_0(x),\qquad \lambda>0.
\]

For \(0<x\le1\), all arguments

\[
\frac{2(3m-1)}{\sqrt{899}}x
\]

belong to \((0,\pi)\). Therefore,

\[
u_0'(x)>0,\qquad u_0'''(x)<0.
\]

Since \(u_0\) is even, \(x=0\) is its unique critical point and \(u_0''(0)=2\). The Schwarzian of \(f_\lambda\) is

\begin{equation}
S(f_\lambda)
=\frac{u_0'''}{u_0'}
-\frac32\left(\frac{u_0''}{u_0'}\right)^2<0
\qquad(0<|x|\le1).
\label{mab:rm:eq:negative-schwarzian}
\end{equation}

\begin{theorem}[Exact membership in the unimodal class]
For
\[
0<\lambda\le \frac{2}{u_0(1)},
\]
the family \(f_\lambda=1-\lambda u_0\) is a self-map that is real--analytic on \([-1,1]\) and even, with a single critical point of quadratic type at the origin and negative Schwarzian away from that point. Consequently, without using Feigenbaum values as a target value, it belongs to the \(S\)-unimodal quadratic class subject to the action of the period-doubling renormalizer \cite{rm:Feigenbaum1978,rm:Lanford1982}. For \(\lambda\) outside that interval, the analytic germ and its local form are retained, but one does not automatically assert a self-map of the interval.
\end{theorem}

The theorem establishes the contribution intrinsic to HMT prior to the
universality theorem: the analytic germ and its type of criticality.

\begin{proof}[Details of the sign and the self-map property]
Differentiating \cref{mab:rm:eq:chaos-profile},
\begin{align*}
u_0'(x)&=\frac1{12}\sum_m k_m\sin(k_mx),\\
u_0'''(x)&=-\frac1{12}\sum_m k_m^3\sin(k_mx),
\qquad k_m=\frac{2(3m-1)}{\sqrt{899}}.
\end{align*}
For \(0<x\le1\), \(0<k_mx<\pi\), so all summands in the first
line are positive and those in the second are negative. Parity extends
monotonicity to both sides of the origin. Since \(u_0(0)=0\) and
\(0\le u_0(x)\le u_0(1)\), the bound
\(0<\lambda\le2/u_0(1)\) implies
\(-1\le1-\lambda u_0(x)\le1\). Finally,
\(u_0''(0)=1/12\sum k_m^2=2\), so the critical point is
quadratic.
\end{proof}

\section{Finite sequences of superstable parameters}

The superstable sequence certified through level eight gives

\[
\delta_8=4.66921218915\ldots,
\qquad
\alpha_8=2.50296847988\ldots.
\]

Its status is \emph{finite certificate free of target values}: each
approximant is obtained from the previously defined family, without introducing
the universal limit as a datum. Decimal agreement provides a
convergence indicator but does not replace the hyperbolicity theorem.

An independent numerical extension of the same sequence, with a step size
adapted to the separation between roots and control of the exact period,
provisionally reaches:

\begin{center}
\small
\begin{tabular}{@{}ccll@{}}
\toprule
Level & \(\lambda_n\) & \(\delta_n\) & \(\alpha_n\)\\
\midrule
11 & 1.87403828195439908045 & 4.66920170694424982745 & 2.50290847517165181764\\
12 & 1.87403836411023460451 & 4.66920163036308158848 & 2.50290800379001546915\\
13 & 1.87403838170549870372 & 4.66920161361839219913 & 2.50290790268748193108\\
14 & 1.87403838547386545431 & 4.66920161007472591023 & 2.50290788100969754274\\
\bottomrule
\end{tabular}
\end{center}

Levels 11--14 retain their status as numerical evaluations beyond the persistent certificate through level eight. Aitken extrapolation gives $\lambda_\infty^{\HMT}\simeq1.8740383865008918080$ and $\alpha_F\simeq2.50290787509$. The theorem in \ref{sec:p3-final:cierre-feigenbaum} establishes parametric scaling of the first roots; the spatial evaluation retains its own scope.

To avoid an unspecified numerical procedure, the method is formulated as
follows. Define
\begin{equation}
F_n(\lambda)=f_\lambda^{\,2^n}(0).
\label{mab:rm:eq:superstable-equation}
\end{equation}
The superstable root \(\lambda_n\) is the first root of \(F_n\) to the
right of \(\lambda_{n-1}\) that satisfies simultaneously
\[
f_{\lambda_n}^{\,2^j}(0)\ne0
\quad(0\le j<n).
\]
This inequality excludes divisor periods. With disjoint intervals
\(I_n\ni\lambda_n\), one computes
\begin{equation}
\delta_n=
\frac{\lambda_{n-1}-\lambda_{n-2}}
     {\lambda_n-\lambda_{n-1}},
\qquad
\alpha_n=-\frac{d_{n-1}}{d_n},
\label{mab:rm:eq:feigenbaum-approximants}
\end{equation}
where \(d_n=f_{\lambda_n}^{\,2^{n-1}}(0)\) is the oriented distance from the
last new point to the critical point. A persistent certificate must retain
the root intervals, the exclusion of smaller periods, and the propagation
of error from \cref{mab:rm:eq:feigenbaum-approximants}; a mere printed decimal does not
satisfy those three conditions.

The discretization rule also constitutes a negative control. A
fixed lower bound greater than the separation
\(|\lambda_n-\lambda_{n-1}|\) can omit the first root and nevertheless produce a numerically plausible
succession. For this reason the increment
must scale with the previous separation and each candidate must be verified
through its exact period.

\section{Conditions for certifying the renormalization operator}

On a space of even analytic functions, we define

\begin{equation}
\mathcal R(f)(x)=-a(f)^{-1}f\bigl(f(-a(f)x)\bigr),
\qquad a(f)=-f(1).
\label{mab:rm:eq:renormalization}
\end{equation}

The complete closure of \(\delta_F\) and \(\alpha_F\) requires jointly certifying:

\begin{enumerate}
\item existence of a fixed point \(g=\mathcal R(g)\) in a declared analytic ball;
\item local uniqueness of that fixed point;
\item a single unstable real eigenvalue of \(D\mathcal R_g\), equal to \(\delta_F\);
\item spectral radius strictly less than one on the stable complement;
\item transversality of the curve \(\lambda\mapsto f_\lambda\) with respect to the stable manifold of \(g\).
\end{enumerate}

A natural HMT implementation employs Chebyshev truncations of order
\(9^m\), interval arithmetic, and radii polynomials. The proof must
provide a ball, an approximate inverse, a defect bound, and a
Lipschitz bound satisfying \(Y+Z(r)<r\). Thus, the nonadic
depth parameterizes a fixed-point certificate and not a mere sequence
of decimal approximations.

In an even Chebyshev basis,
\[
g(x)=\sum_{j=0}^{N}g_{2j}T_{2j}(x),\qquad N=9^m,
\]
the finite residue is \(F_N(g)=\Pi_N(\mathcal Rg-g)\). If \(A_N\) approximates
\((DF_N(\bar g))^{-1}\), the quantities that must be certified are
\begin{align*}
Y&=\|A_NF_N(\bar g)\|,\\
Z_0&=\|I-A_NDF_N(\bar g)\|,\\
Z_1(r)&=\sup_{\|h\|\le r}
\|A_N(DF_N(\bar g+h)-DF_N(\bar g))\|.
\end{align*}
An inequality
\(Y+(Z_0+Z_1(r))r<r\) encloses a unique fixed point. The later spectral block
must isolate an unstable direction and contract the complement. The
use of \(9^m\) thus enters as a refinement law of the certificate, not as
numerological decoration.

\begin{corollary}[Exact convergence of mode thirty]
Mode \(30\) admits three simultaneous and typologically
distinct realizations:
\begin{align*}
30&\longmapsto(3\cdot30,4\cdot30)=(90,120)
&&\text{in the constitutive completion},\\
30&\longmapsto30^2-1=899=29\cdot31
&&\text{in the critical normalization},\\
30&\longmapsto
M_{30}=\begin{pmatrix}30&899\\1&30\end{pmatrix}
&&\text{in the hyperbolic Pell transformation}.
\end{align*}
The three identities are exact and are obtained without prescribing target values. The conserved invariant is the central integer \(30\) together with its orientation: extension \(3\to4\), quadratic defect \(-1\), and a unit of norm one. What is not asserted is the equality between the vacuum operator and the duplication renormalizer.
\end{corollary}

\begin{remark}[Negative control]
A second eigenvalue outside the unit disk, a ball of radii without solution, loss of uniqueness, or tangency of the HMT family to the stable variety refutes the universal promotion. None refutes the exact germ or the Schwarzian theorem.
\end{remark}

\begin{proposition}[Main Result]
The structural significance of the doubling constant lies in the
cancellation in dodecaphase form of \(\pi\), in the identity
\(899=30^2-1\), in the closed trigonometric form, and in its proved membership
in the analytic class on which the renormalizer acts. Its
decimal expansion constitutes exclusively the terminal evaluation.
\end{proposition}

\subsection*{Logical Scope and Domain of Validity}
The profile, normalization, Pell unit, unimodality and Schwarzian retain the preceding proofs. The eight certified levels and the extrapolations at levels 11--14 retain their respective numerical scope. The development in \ref{sec:p3-final:cierre-feigenbaum} proves parametric scaling for the original first-root sequence in the uniform sector, through identification of centres and remainder control. This conclusion does not transfer a proof about the parametric ratio to the spatial ratio.

\subsection{Weighted harmonic deformation and selection of the superstable branch}
\label{mab:sec:extension-dodecafasica-ponderada-armonica}

The uniform profile has a parametric extension that preserves the twelve
phases and separates the harmonic deformation from the dynamic parameter. For
\(1\le m\le12\), let

\begin{equation}
 \phi_m=\frac{(30m-10)\pi}{180},
 \qquad
 p(\phi)=12\cos(3\phi)-\cos(4\phi).
 \label{mab:eq:fases-y-deformacion-armonica-armonizada}
\end{equation}

For every real \(\eta\) such that the weights remain positive, we define

\begin{equation}
 w_m(\eta)=\frac{1+\eta p(\phi_m)}{12}.
 \label{mab:eq:pesos-dodecafasicos-armonizados}
\end{equation}

The sums of the third and fourth harmonics over the twelve phases vanish;
therefore,

\[
 \sum_{m=1}^{12}w_m(\eta)=1.
\]

Let

\begin{equation}
 s_\eta^2=\frac{2}{\sum_{m=1}^{12}w_m(\eta)\phi_m^2},
 \qquad
 u_\eta(x)=1-\sum_{m=1}^{12}w_m(\eta)
 \cos(s_\eta\phi_m x).
 \label{mab:eq:familia-ponderada-armonizada}
\end{equation}

The cosine expansion and the definition of \(s_\eta\) give

\[
 u_\eta(x)=x^2+O(x^4)\qquad(x\to0).
\]

Consequently, \(u_\eta\) is real--analytic, even, and retains a quadratic
critical point at the origin. We introduce the dynamical parameter through

\begin{equation}
 f_{\lambda,\eta}(x)=1-\lambda u_\eta(x).
 \label{mab:eq:familia-feigenbaum-ponderada-armonizada}
\end{equation}

A superstable parameter of exact period \(2^n\) satisfies

\[
 f_{\lambda,\eta}^{\,2^n}(0)=0,
 \qquad
 f_{\lambda,\eta}^{\,2^k}(0)\ne0\quad(0\le k<n).
\]

The equation may have several branches. A branch
\(\mathscr B\) is fixed by traversing an adaptive mesh immediately after the parameter of the
preceding level; the first interval with a sign change is taken
and bisection is applied. The notation
\(\lambda_n^{\mathscr B}(\eta)\) incorporates this selector and does not attribute
global uniqueness to the superstable equation. The first level is determined
exactly:

\[
 \lambda_1^{\mathscr B}(\eta)=\frac1{u_\eta(1)},
\]

if \(u_\eta(1)>0\), because
\(f_{\lambda,\eta}^{\,2}(0)=f_{\lambda,\eta}(1)=0\).
Given three consecutive levels and

\[
 x_n^{\mathscr B}(\eta)
 :=f_{\lambda_n^{\mathscr B}(\eta),\eta}^{\,2^{n-1}}(0),
\]

we form

\begin{equation}
 \delta_n^{\mathscr B}(\eta)
 =\frac{\lambda_{n-1}^{\mathscr B}(\eta)-
         \lambda_{n-2}^{\mathscr B}(\eta)}
        {\lambda_n^{\mathscr B}(\eta)-
         \lambda_{n-1}^{\mathscr B}(\eta)},
 \qquad
 \alpha_n^{\mathscr B}(\eta)
 =\frac{|x_{n-1}^{\mathscr B}(\eta)|}
        {|x_n^{\mathscr B}(\eta)|}.
 \label{mab:eq:aproximantes-ponderados-armonizados}
\end{equation}

For \(\eta=0\), the uniform profile is recovered exactly. The weighted
extension introduces a transverse parameter: it permits study of the
transversality of the family with respect to the stable manifold of the
renormalizer without converting a finite coincidence into a proof of
universality.

\begin{remark}[Negative control]
The extension is refuted if the weights lose positivity on the declared
domain, if normalization ceases to produce a unit quadratic
coefficient, if the selector does not exclude divisor periods, or if a branch is chosen
using the terminal universal values as input.
\end{remark}
 \section{First-root selection and Feigenbaum parametric scaling}
\label{sec:p3-final:cierre-feigenbaum}

The uniform dodecaphasic reader also determines a sequence of first roots of critical return. Its separation ratio is proved to converge to the Feigenbaum constant, preserving the original selection, orientation and memory of the construction.

% Focal insertion; hierarchy inherited from the host.
\begingroup
\let\HMTFGHostSection\section
\edef\HMTFGSavedSecnumdepth{\number\value{secnumdepth}}
\setcounter{secnumdepth}{4}
\let\HMTFGHostSubsection\subsection
\let\HMTFGHostSubsubsection\subsubsection
\let\section\HMTFGHostSubsection
\let\subsection\HMTFGHostSubsubsection
\let\subsubsection\paragraph
% Fragmento maestro: incluir desde el anfitrion, con \HMTFGRoot fijada a la raiz del paquete.
% Requiere amsmath, amssymb/mathrsfs o unicode-math, amsthm, booktabs, longtable, hyperref, needspace.
% El anfitrion conserva sus entornos theorem, proposition, lemma y proof.
\providecommand{\HMTFGRoot}{.}
\begingroup
\providecommand{\ZZ}{}\renewcommand{\ZZ}{\mathbb Z}
\providecommand{\RR}{}\renewcommand{\RR}{\mathbb R}
\providecommand{\FF}{}\renewcommand{\FF}{\mathbb F}
\providecommand{\APP}{}\renewcommand{\APP}{\mathrm{APP}}
\providecommand{\TPK}{}\renewcommand{\TPK}{\mathrm{TPK}}
\providecommand{\TRIT}{}\renewcommand{\TRIT}{\{-1,0,+1\}}
% unicode-math usa lBrack/rBrack; newtxmath ya proporciona llbracket/rrbracket.
\providecommand{\llbracket}{\lBrack}
\providecommand{\rrbracket}{\rBrack}
% Fragmento de inserción; se incluye desde la raíz de este paquete editorial.
% Las fuentes completas y los cortes materiales constan en ANTECEDENTES_Y_DEPENDENCIAS.md.
% Procedencia: XI REV11 ES, sections/nucleo.tex, tpk_desarrollo_integrado.tex,
% ch_tres_vueltas_ext.tex, ch_elevacion_catalogal.tex, ch_cierre_cardinal.tex,
% ch_descenso_por_fibras.tex; perfil integral public_final/feigenbaum/c37_feigenbaum.tex;
% Grassmannianos REV07 ES, sections/jacobi_spectral.tex;
% TRANSVERSALIDAD_Y_ESCALADO_LOCAL_FEIGENBAUM.md, sección 14.
% Ampliación expositiva de demostraciones existentes; no introduce resultados nuevos.
\providecommand{\ZZ}{\mathbb Z}
\providecommand{\RR}{\mathbb R}
\providecommand{\FF}{\mathbb F}
\providecommand{\APP}{\mathrm{APP}}
\providecommand{\TPK}{\mathrm{TPK}}
\providecommand{\TRIT}{\{-1,0,+1\}}

\section{Operator antecedents of the criticality family}
\label{hmtfg:antecedentes-operatorios}

The criticality family is constructed from a readout of the TPK calendar. The two arithmetic sheets, tritic orientation and transport with memory first determine the states and their prolongations. The ordered reading of the continuum then allows analytic functions to be evaluated, parameters to be compared and roots to be selected. This development preserves the genealogical state and its two readouts: the Archimedean chart provides the evaluation space, while the incidence readout preserves boundary relations. The joint discrete structure of the continuum retains its five consubstantial constructions. The present exposition uses its ordered reader and preserves the common antecedent from which it originates.

\subsection{Arithmetic cells, orientation and transport}
\label{hmtfg:app-trit-transporte}

% XI REV11, sections/nucleo.tex: definiciones de APP y TRIT;
% tpk_desarrollo_integrado.tex:22--165, transporte levantado y calendario.
Let \(D_9=\{1,\ldots,9\}\) and \(X_9=(\ZZ/9\ZZ)^2\). For
\(n\in\ZZ\), division with a positive representative is
\[
 \rho_9(n)=1+((n-1)\bmod9),\qquad
 q_9(n)=\frac{n-\rho_9(n)}9,\qquad
 n=\rho_9(n)+9q_9(n).
\]
The APP cell associated with \((p,q)\in D_9^2\) is
\[
 a_{p,q}=(p,q;R^\Sigma,q^\Sigma,R^\Pi,q^\Pi),
 \quad
 p+q=R^\Sigma+9q^\Sigma,\qquad
 pq=R^\Pi+9q^\Pi.
\]
The two pairs \((R^\Sigma,q^\Sigma)\) and \((R^\Pi,q^\Pi)\) retain the additive and multiplicative evaluations, respectively. The division formula reconstructs both evaluations before any subsequent reduction. Selection of an active sheet is recorded in a state coordinate and preserves the complete cell.

The balanced tritic representative \(r_3(n)\in\{-1,0,+1\}\) satisfies
\(n\equiv r_3(n)\pmod3\). Its carry is
\[
 \chi(a,b)=\frac{a+b-r_3(a+b)}3,\qquad
 a,b\in\{-1,0,+1\}.
\]
The equality
\[
 \chi(a,b)+\chi(r_3(a+b),c)
 =\chi(b,c)+\chi(a,r_3(b+c))
\]
is obtained by writing both sides as
\((a+b+c-r_3(a+b+c))/3\). This identity preserves the carry when regrouping a tritic sum. The operative selector is
\[
 M_{\rm ph}(r)=
 \begin{cases}
 +1,&r\in\{1,4,7\},\\
 -1,&r\in\{2,5,8\},\\
 0,&r\in\{3,6,9\}.
 \end{cases}
\]
The value \(+1\) activates the additive sheet, the value \(-1\) the multiplicative sheet, and the value \(0\) preserves the position of both cursors during the neutral event. All three events belong to the complete chronology and are recorded in the route.

Let \(\mathcal D=\{N,E,S,O\}\), with vectors
\(\nu_N=(-1,0)\), \(\nu_E=(0,1)\), \(\nu_S=(1,0)\) and
\(\nu_O=(0,-1)\), and \(\Theta_{108}=\ZZ/108\ZZ\). The observable configuration is
\[
 q=(x_+,d_+,x_\times,d_\times,\theta)
 \in\mathcal Q_{\rm obs}=(X_9\times\mathcal D)^2\times\Theta_{108}.
\]
For the representative \(\bar\theta\in\{0,\ldots,107\}\), define
\[
 r(\theta)=1+(\bar\theta\bmod9),\qquad
 \beta(\theta)=\left[\left\lfloor\frac{\bar\theta}{9}\right\rfloor\right]_{12},
\]
\[
 a_+(\theta)=\mathbf1_{\{1,4,7\}}(r(\theta)),\qquad
 a_\times(\theta)=\mathbf1_{\{2,5,8\}}(r(\theta)).
\]
The position transition is
\[
 x_+'=x_++a_+(\theta)\nu_{d_+},\qquad
 x_\times'=x_\times+a_\times(\theta)\nu_{d_\times}.
\]
The involution \(\jmath\) exchanges \(N,S\) and \(E,O\). With
\[
 h(\theta)=\mathbf1_{\{0,27,54,81\}}\bigl((\bar\theta+1)\bmod108\bigr),
\]
the complete observable operator is fixed by
\[
 F_{\rm obs}(q)=
 (x_+',\jmath^{h(\theta)}d_+,x_\times',
   \jmath^{h(\theta)}d_\times,\theta+1).
\]

\begin{proposition}[Conservation of lifted transport and observable return]
\label{hmtfg:transporte-calendario}
Transport preserves the integer displacement through the record of boundary crossings. The operator \(F_{\rm obs}\) is invertible and satisfies
\(F_{\rm obs}^{108}=I\).
\end{proposition}
\begin{proof}
Let \(s:X_9\to D_9^2\) be the section of positive representatives. For
\(x'=x+a\nu_d\), \(a\in\{0,1\}\), define
\[
 b_d(x,a)=\frac{s(x)+a\nu_d-s(x')}{9}\in\ZZ^2.
\]
If \(z'=z+b_d(x,a)\), then
\[
 s(x')+9z'=s(x)+9z+a\nu_d.
\]
Summing this equality along a route gives
\[
 s(x_n)+9z_n=s(x_0)+9z_0+\sum_{t=0}^{n-1}a_t\nu_{d_t}.
\]
Partitioning the sum preserves the identity under concatenation. To invert \(F_{\rm obs}\), first subtract one unit from the phase; that phase determines the reversal of directions and the displaced cursor. The previous directions are recovered and the corresponding vector is subtracted. Each block of \(27\) steps contains nine activations of each cursor. The next block reverses directions and cancels their integer displacements. In \(54\) steps, positions and directions are restored; in \(108\), the calendar coordinate is also restored.
\end{proof}

In addition to this configuration, the complete state preserves both arithmetic evaluations, carries, orientation, trajectory, prefix, compatible cylinder, its boundary, memory degree and survival record. Its update factors as
\[
 U_t=\mathsf{Mem}_t\circ\mathsf{Tra}_t\circ\mathsf{Sel}_t.
\]
The holonomy \(\Gamma_9\) composes nine complete transitions. Phase return and memory advance are different coordinates of this composition. The following construction fixes the transitions and their update field by field.

% Recuperación íntegra del propietario de 855 líneas, con etiquetas prefijadas.
\input{inserciones_20260930/feigenbaum/00a_transiciones_nonadicas.tex}

\subsection{Cylinder refinement and ordered readout}
\label{hmtfg:refinamiento-ordenado}

% XI REV11, ch_elevacion_catalogal.tex: bloques de 54 aristas, beta_blk,
% partición efectiva de 729 subcilindros; ch_cierre_cardinal.tex:375--497.
For \(\mathbf b=(b_1,\ldots,b_6)\in\FF_3^6\), let
\(\upsilon:\FF_3\to\{-1,0,+1\}\) be the balanced representative. The three prolongations constructed edge by edge define
\[
 \Gamma^{[54]}_{\mathbf b,k}
 =\gamma_{\upsilon(b_6),k+45}\circ\cdots\circ
   \gamma_{\upsilon(b_1),k}.
\]
Starting from \(Y_0=\{x_\star\}\), put
\[
 Y_{N+1}
 =\{\Gamma^{[54]}_{\mathbf b,1+54N}(x):
          x\in Y_N,\ \mathbf b\in\FF_3^6\}.
\]
Truncation \(\rho^{\rm blk}_{N+1,N}\) removes the last fifty-four edges, recovering the previous state through the transition record.

\begin{proposition}[Effective coding of prolongation blocks]
\label{hmtfg:bloques-729}
Reading the six tritic pairs of each block defines bijections
\[
 \beta_N:Y_N\xrightarrow{\ \cong\ }(\FF_3^6)^N
\]
that commute with truncation. Each state of \(Y_N\) has exactly \(729\) block prolongations in \(Y_{N+1}\).
\end{proposition}
\begin{proof}
On each turn, the record retains the tritic pair and its carry. The three pairs produce the three distinct outputs \(\chi(\varepsilon,\varepsilon)=\varepsilon\), and the record allows recovery of the branch used, including the zero-carry branch. The six concatenations therefore produce all words of \(\FF_3^6\). Two different words are distinguished at the first recorded pair in which they differ. Recovery of the prefix proves compatibility with truncation. By induction on \(N\), every word of length \(N\) has a unique antecedent state; the next word can be any of the \(3^6=729\) possibilities.
\end{proof}

Let \(\Omega_{\Gamma}^{\rm cal}=\varprojlim_NY_N\). The preceding bijections induce
\[
 \Omega_{\Gamma}^{\rm cal}\cong(\FF_3^6)^{\mathbb N}.
\]
The correspondence and its inverse are continuous: each finite coordinate depends on a finite prefix. The cylinders are open and closed. Each cylinder of depth \(N\) is partitioned into its \(729\) subcylinders of depth \(N+1\); each contains an infinite prolongation, obtained, for example, by continuing with the zero branch in all subsequent blocks.

% Antecedentes efectivos del continuo conjunto y del universo interno.
% Se imprimen en ambos destinos; la etiqueta de universo no reemplaza su regla.
\input{inserciones_20260930/feigenbaum/00b_continuo_universo.tex}

\subsection{Ordered readout of block cylinders}
\label{hmtfg:publicacion-cilindros-bloque}

For the representative \([b]_3\in\{0,1,2\}\), fix
\[
 \nu(\mathbf b)=\sum_{i=1}^6[b_i]_3\,3^{6-i}\in\{0,\ldots,728\}.
\]
In the HMT ordered completion, a prefix
\(s=(m;\mathbf b_0,\ldots,\mathbf b_{N-1})\), \(m\in\ZZ\), determines
\[
 a_N(s)=m+\sum_{j=0}^{N-1}\nu(\mathbf b_j)\,729^{-j-1},\qquad
 I_N(s)=[a_N(s),a_N(s)+729^{-N})_{\rm HMT}.
\]
The carry preserves the joining of the two representatives of a boundary. If \(C_N(s)=\overline{I_N(s)}\), the preceding relations give
\[
 I_N(s)=\bigsqcup_{\mathbf b\in\FF_3^6}I_{N+1}(s;\mathbf b),
 \qquad \operatorname{diam}C_N(s)=729^{-N}\longrightarrow0.
\]

\begin{proposition}[Coverage of the Archimedean readout]
\label{hmtfg:cobertura-arquimediana}
The map
\[
 P_{\rm ar}:\ZZ\times\Omega_\Gamma^{\rm cal}
       \longrightarrow\mathbb R_{\rm HMT},\qquad
 \{P_{\rm ar}(m,\xi)\}
    =\bigcap_N C_N(m;\beta_N(\rho_{\infty,N}\xi))
\]
is well defined, is surjective and determines
\[
 (\ZZ\times\Omega_\Gamma^{\rm cal})/\!\sim_{P_{\rm ar}}
       \ \cong\ \mathbb R_{\rm HMT}.
\]
\end{proposition}
\begin{proof}
The left endpoints form a Cauchy sequence in the HMT completion, since the intervals are nested and their diameters tend to zero. Their limit belongs to all the closed intervals and is unique. For any element of the completion, first choose the integer interval containing it and then the unique half-open subinterval of each partition containing it. The preceding proposition realizes each of those choices through an effective TPK prolongation. The compatible prefixes define \(\xi\), whose readout is the chosen element. At boundaries, the carry identifies equivalent representations after their construction. By definition and surjectivity, the quotient by equality of readout is in bijection with the completion.
\end{proof}

The construction is also carried out in the internal genealogical universe \(\mathfrak G=\mathfrak G_{\rm HMT}(h,m_\infty)\). Here \(h\) encodes the operator rules and \(m_\infty\) the compatible history; the universe is obtained through initial transitive closure, definability at each successor and union at limit stages:
\[
 G_0=\operatorname{TC}\{\omega,u_{h,m},h,m_\infty\},\qquad
 G_{\beta+1}=\operatorname{Def}_{\Sigma_{\rm HMT}}(G_\beta),\qquad
 G_\lambda=\bigcup_{\beta<\lambda}G_\beta .
\]
The relations of \(\Sigma_{\rm HMT}\) are those of the operations already constructed and their codes. The internal domain is therefore retained in all coverage claims. Denote its completion by \(\mathbb R_{\rm HMT}^{\mathfrak G}\) and the ordered realization within the same universe by \(\mathbb R^{\mathfrak G}\).

\begin{proposition}[Internal ordered chart]
\label{hmtfg:carta-interna}
There exists a unique isomorphism of complete Archimedean ordered fields
\[
 \iota:\mathbb R_{\rm HMT}^{\mathfrak G}
             \xrightarrow{\ \cong\ }\mathbb R^{\mathfrak G}
\]
that fixes the unit. It preserves rational endpoints, positive square roots and limits of convergent sequences.
\end{proposition}
\begin{proof}
The unit fixes the integers and rationals. Each element of the HMT completion is associated with the cut of rationals less than it; completeness of the ordered realization provides the unique element with that cut. Rational density proves injectivity and strict order preservation. The same procedure in the reverse direction proves surjectivity. Rational approximations from above and below preserve addition and multiplication and thereby the field structure. The positive element whose square is \(a>0\) is transported to the positive element whose square is \(\iota(a)\). If a sequence converges, for every positive rational tolerance its final terms lie in the rational neighborhood of the limit; order preservation transports that condition. All these operations are performed within \(\mathfrak G\). Uniqueness follows from preservation of rational cuts.
\end{proof}

\subsection{Descent along fibers and preservation of the selector}
\label{hmtfg:descenso-selector}

% XI REV11, ch_descenso_por_fibras.tex:10--39; prueba completa pertinente.
\begin{proposition}[Descent criterion for an operation]
\label{hmtfg:criterio-descenso}
Let \(q:\Omega\twoheadrightarrow J\) and
\(\star:D\subseteq\Omega^2\to\Omega\), with \(D\) saturated by the fibers of \(q\times q\). There exists a unique operation \(\bar\star\) on \((q\times q)(D)\) such that
\[
 q(x\star y)=q(x)\,\bar\star\,q(y)
\]
if and only if
\[
 q(x)=q(x'),\quad q(y)=q(y')
 \ \Longrightarrow\ q(x\star y)=q(x'\star y')
\]
for pairs in the domain.
\end{proposition}
\begin{proof}
If \(\bar\star\) exists, both results are the evaluation of the same pair of arguments. Conversely, choose antecedents of the visible arguments and define the result as the readout of their operation. Saturation fixes the domain; the condition on fibers makes the result independent of the chosen antecedents. Surjectivity proves uniqueness.
\end{proof}

In an application through \(P_{\rm ar}\), this criterion identifies the precise condition for transporting an operation from complete states to their readout. In the chart \(\iota\), the field structure and limits already satisfy that condition. Next, the analytic profile is constructed from the calendar and this compatibility is used to transport its iteration and roots.

\subsection{Dodecaphase readout and quadratic normalization}
\label{hmtfg:perfil-dodecafasico}

The profile, its phases and its normalization are those already constructed in \eqref{mab:rm:eq:critical-phase-reader}--\eqref{mab:rm:eq:chaos-profile}. Its closed form and moments are proved in \eqref{mab:rm:eq:chaos-closed} and \eqref{mab:rm:eq:chaos-moments24}. The finite cosine sum determines the removable extension at every zero of the closed-form denominator. The angular moment retains oriented real representatives: changing them modulo $2\pi$ changes the moment. The reader and its uniform sector are the explicit starting structure of this realization.

\subsection{Transport of iterates, roots and minima}
\label{hmtfg:transporte-raices}

% Reunión existente: TRANSVERSALIDAD_Y_ESCALADO_LOCAL_FEIGENBAUM.md, sección 14.1.
Within \(\mathbb R_{\rm HMT}^{\mathfrak G}\), the convergent series
\[
 \cos_{\rm HMT}(t)=\sum_{j=0}^{\infty}
                   \frac{(-1)^j t^{2j}}{(2j)!}
\]
defines the corresponding analytic realization. Let
\[
 u_{\rm HMT}(x)=\frac1{12}\sum_{m=1}^{12}
  \left[1-\cos_{\rm HMT}
   \left(\frac{2(3m-1)x}{\sqrt{899}_{\rm HMT}}\right)\right],
 \qquad F_a(x)=1-a\,u_{\rm HMT}(x),
\]
\[
 S_n^{\rm HMT}(a)=F_a^{\,2^n}(0),\qquad
 S_n(\lambda)=f_\lambda^{\,2^n}(0).
\]

\begin{proposition}[Ordered conjugacy of the family and the root selector]
\label{hmtfg:conjugacion-selector}
For the internal arguments of the preceding expressions,
\[
 \iota\circ F_a=f_{\iota(a)}\circ\iota,\qquad
 \iota(S_n^{\rm HMT}(a))=S_n(\iota(a)).
\]
The chart transports exactly the conditions of return, exact period and parameter order. If a set of admissible roots has a minimum, the image of that minimum is the minimum of the image set.
\end{proposition}
\begin{proof}
The chart fixes the rationals, transports the positive square root and preserves limits of the cosine partial sums. Hence, \(\iota(u_{\rm HMT}(x))=u_0(\iota(x))\). Compatibility with addition and multiplication proves the first equality; induction on the number of iterations proves the second. A root is characterized by equality to zero, which is preserved in both directions. The exact period \(2^n\) adds the inequalities \(F_a^{2^j}(0)\ne0\), \(0\le j<n\), also preserved by the injectivity of \(\iota\). If the parameter is required to exceed a preceding center, that inequality is transported by order preservation. Finally, \(a\le b\) for every \(b\) in the set is equivalent to \(\iota(a)\le\iota(b)\) for every element of its image. Surjectivity onto the sets thus defined completes the assertion about the minimum.
\end{proof}

The domain of this transport is the indicated internal continuum. The existence of each minimum and its identification with a specific dynamical branch are results of the criticality construction that follows; the chart preserves those properties once they have been established.

\subsection{Finite spectral representation of the uniform reader}
\label{hmtfg:jacobi-finito}

% Grassmannianos REV07, sections/jacobi_spectral.tex:35--136,
% mecanismo de Hankel, Gram--Schmidt, recurrencia y representación espectral.
% Especialización atómica ya reunida en la sección 14.2 del desarrollo Feigenbaum.
The squared frequencies of the profile define the positive measure
\[
 s_m=k_m^2=\frac{4(3m-1)^2}{899},\qquad
 \nu_0=\frac1{12}\sum_{m=1}^{12}\delta_{s_m},\qquad
 M_r=\int s^r\,d\nu_0(s).
\]
Its twelve nodes are distinct and positive. For \(1\le r\le12\), the Hankel matrix \(H_r=(M_{i+j})_{0\le i,j<r}\) is positive definite. Indeed, for \(p\) of degree less than \(r\),
\[
 \mathbf p^{\mathsf T}H_r\mathbf p
   =\frac1{12}\sum_{m=1}^{12}p(s_m)^2>0
\]
unless \(p=0\), because a polynomial of degree at most eleven that vanishes at twelve distinct nodes is zero. At degree twelve appears the polynomial \(\prod_m(s-s_m)\), with zero norm in this measure. The corresponding spectral representation has dimension twelve.

Let \(p_0,\ldots,p_{11}\) be the monic orthogonal polynomials,
\(h_j=\int p_j^2\,d\nu_0>0\) and
\(\psi_j=p_j/\sqrt{h_j}\). Since \(\nu_0\) is a probability measure,
\(\psi_0=1\). Multiplication by \(s\) in \(L^2(\nu_0)\) is self-adjoint and has a tridiagonal matrix \(J_{12}\) in this basis.

\begin{proposition}[Exact recurrence and spectral intertwining]
\label{hmtfg:jacobi-entrelazamiento}
With
\[
 a_j=\frac{\int s p_j(s)^2\,d\nu_0(s)}{h_j},\qquad
 \beta_j=\frac{h_j}{h_{j-1}}>0\quad(1\le j\le11),
\]
the matrix \(J_{12}\) has diagonal \(a_j\) and subdiagonal
\(\sqrt{\beta_j}\). If
\[
 Q_{mj}=\frac{\psi_j(s_m)}{\sqrt{12}}
 \quad(1\le m\le12,\ 0\le j\le11),\qquad
 D=\operatorname{diag}(s_1,\ldots,s_{12}),
\]
then
\[
 Q^{\mathsf T}Q=I,\qquad DQ=QJ_{12},\qquad
 J_{12}=Q^{\mathsf T}DQ,
\]
and the profile satisfies
\begin{equation}
 u_0(x)=e_0^{\mathsf T}
          \bigl[I-\cos(x\sqrt{J_{12}})\bigr]e_0.
 \label{hmtfg:perfil-espectral}
\end{equation}
\end{proposition}
\begin{proof}
Positivity of the Hankel systems uniquely provides the monic orthogonal polynomials. On expanding \(s p_j\) in the orthogonal basis, the coefficients of \(p_i\), \(i\le j-2\), vanish because
\(\langle sp_j,p_i\rangle=\langle p_j,sp_i\rangle=0\).
The coefficient of \(p_{j+1}\) is one by monicity, that of \(p_j\) is \(a_j\) and that of \(p_{j-1}\) is \(h_j/h_{j-1}\). At the last degree, the same equality is understood in \(L^2(\nu_0)\), where the monic polynomial of degree twelve vanishing at the nodes represents the zero vector. Normalization gives the stated positive subdiagonal.

Orthonormality is equivalent to \(Q^{\mathsf T}Q=I\). The recurrence equality at each node gives \(DQ=QJ_{12}\). Since \(Q\) is square, we also have \(QQ^{\mathsf T}=I\), and obtain the diagonalization of \(J_{12}\). Its positive eigenvalues allow its positive square root to be defined. For every function defined at those twelve nodes,
\[
 e_0^{\mathsf T}\Phi(J_{12})e_0
   =\frac1{12}\sum_{m=1}^{12}\Phi(s_m),
\]
since \(Qe_0=(1,\ldots,1)^{\mathsf T}/\sqrt{12}\).
The choice \(\Phi(s)=1-\cos(x\sqrt{s})\) proves
\eqref{hmtfg:perfil-espectral}.
\end{proof}

The representation preserves exactly the uniform functional, its moments and its profile. Its first coefficients are
\[
 a_0=M_1=2,\qquad
 \beta_1=M_2-M_1^2=\frac{2493348}{808201}.
\]
The complete genealogy remains in the TPK state from which the reader originates; \(J_{12}\) represents its specific spectral readout. The Gram--Schmidt construction and the Jacobi recurrence are applied to this atomic measure, terminating in dimension twelve.

The operator chain used in the following sections is thus fixed: APP cells and TRIT orientation feed the TPK transition; its prolongations produce cylinders and their ordered readout; the twelve oriented windows and the uniform functional produce \(u_0\); the family \(f_\lambda\) determines returns and roots; the internal chart preserves the order of their selectors. The spectral realization describes the same profile through a finite self-adjoint matrix and a cyclic vector.

% Fragmento integrable. Preambulo anfitrion: amsmath, amssymb, longtable, hyperref.
\section{Dodecaphase profile and quadratic return}
\label{hmtfg:fragmento:perfil-dodecafasico-y-retorno-cuadratico}

\subsection{Genealogy and object of the proof}
\label{hmtfg:desarrollo:1}

The initial object is the critical readout of the oriented dodecaphase wheel already constructed in the corpus. APP evaluates addition and multiplication on the two sheets of the digit lattice, preserving quotient and remainder; TRIT determines regime and orientation; TPK transport selects, transports and updates the state with its carry, boundary and memory records. The nonadic connection prolongs that state, with phase return and memory advance. Its readers act on the same discrete structure of the continuum.

The focal reader of the source is
\[
\mathcal P_{\mathrm{crit}}:U_{12}^{\mathrm{sign}}\longrightarrow
\mathbb R/2\pi\mathbb Z,\qquad
m\longmapsto\phi_m=\frac{(3m-1)\pi}{18},\quad 1\le m\le12.
\]
To evaluate the angular moment, the oriented representative of that chart is retained: replacing it with another representative modulo \(2\pi\) would change the moment. The uniform trace fixes \(w_m=1/12\). The focal chain is
\[\begin{aligned}
\text{APP–TRIT–TPK}&\longrightarrow
\text{discrete structure of the continuum}\\ &\longrightarrow U_{12}^{\mathrm{sign}}
\xrightarrow{\mathcal P_{\mathrm{crit}}}(\phi_m,w_m)
\\ &\longrightarrow u_0\longrightarrow f_\lambda
\longrightarrow\text{returns and renormalization}.
\end{aligned}\]

The abbreviated notation retains the complete state and its prolongations as antecedent; the scalar family is a readout of that antecedent. The already fixed reader, with its order, orientation and uniform sector, is taken as an explicit starting structure. The conclusions concern that particular reader; uniqueness of that reader among all possible TPK readouts lies outside these proofs.

The coordinate \(\pi_{\mathrm{HMT}}\) used by the chart belongs to the internal outputs of the corpus. Its factor cancels exactly when normalizing the second moment. The Feigenbaum values enter exclusively in subsequent recognition, never in the definition of the family or the certification of roots.


\subsection{Exact profile and quadratic criticality}
\label{hmtfg:desarrollo:2}

We retain the notation $s_m=3m-1$, $\kappa_0=s_0$, $k_m=2s_m/\sqrt{899}$ and $f_{\lambda,0}=f_\lambda$ of the profile already proved. Its entire extension has the rational series
\[
u_0(z)=\sum_{j\ge1}(-1)^{j+1}
\frac{4^j\sum_m s_m^{2j}}{12\,899^j(2j)!}z^{2j}
=z^2-\frac{715769}{2424603}z^4
+\frac{1353832978}{32695771455}z^6+\cdots.
\]
Thus the finite record of twelve sectors produces an entire profile, with quadratic criticality and scale fixed before iteration.


\subsection{Structural robustness with uniform bounds}
\label{hmtfg:desarrollo:3}

The deformation is defined in \eqref{mab:eq:pesos-dodecafasicos-armonizados}--\eqref{mab:eq:familia-feigenbaum-ponderada-armonizada}. We write $M_{2,\eta}=\sum_mw_m\phi_m^2$ and $\kappa_\eta=s_\eta=\sqrt{2/M_{2,\eta}}$, retaining the same weights, phases and profiles.
The notation is $p_m=p(\phi_m)=12\cos(3\phi_m)-\cos(4\phi_m)$ and $k_{m,\eta}=\kappa_\eta\phi_m$; in the following weighted formulae, $k_m$ abbreviates $k_{m,\eta}$, whereas $k_{m,0}$ denotes the uniform sector.

\textbf{Structural robustness theorem.} For \(|\eta|\le1/40\), the weights are positive and sum to one. The profile is even and entire, satisfies \(u_\eta''(0)=2\) and is strictly increasing on \((0,1]\). For \(0<\lambda\le2/u_\eta(1)\), the map sends \([-1,1]\) into itself, has a unique critical point in that interval and satisfies
\[
S f_{\lambda,\eta}(x)\le-\frac{640}{47647}<0,\qquad0<|x|\le1.
\]

\textbf{Proof.} The sums of cosines \(3\phi_m\) and \(4\phi_m\) vanish through cycles of lengths four and three. Since \(|p_m|\le13\),
\[
\frac9{160}\le w_m\le\frac{53}{480},\qquad
\frac{27}{40}M_{2,0}\le M_{2,\eta}\le\frac{53}{40}M_{2,0}.
\]
It follows that
\[
k_{\max,\eta}^2\le\frac{196000}{24273}<9<\pi^2,\qquad
k_{\min,\eta}^2\ge\frac{640}{47647}.
\]
For \(0<x\le1\), all sines \(\sin(k_{m,\eta}x)\) are positive, whence
\[
u_\eta'(x)=\sum_mw_mk_m\sin(k_mx)>0,\qquad
u_\eta'''(x)=-\sum_mw_mk_m^3\sin(k_mx)<0.
\]
The quotient \(u_\eta'''/u_\eta'\) is the negative of a positive mean of the \(k_m^2\). Invariance of the Schwarzian under affine changes of the value gives
\[
S f_{\lambda,\eta}
=\frac{u_\eta'''}{u_\eta'}-\frac32
\left(\frac{u_\eta''}{u_\eta'}\right)^2
\le-k_{\min,\eta}^2.
\]
Evenness covers \(x<0\), and the image of the interval is
\([1-\lambda u_\eta(1),1]\). \(\square\)

The width \(1/40\) is a choice made to obtain uniform bounds, distinct from a fundamental constant or an optimal threshold. The original sector remains \(\eta=0\).


\subsection{Global quadratic holomorphic coordinate on a strip}
\label{hmtfg:desarrollo:4}

\textbf{Holomorphic factorization theorem.} For \(|\eta|\le1/40\), in
\[
\mathcal S=\{z\in\mathbb C:|\operatorname{Re}z|<\pi/3\}
\]
there exists a holomorphic, odd and univalent function \(b_\eta\), with \(b_\eta'(0)=1\), such that
\[
\boxed{u_\eta(z)=b_\eta(z)^2,\qquad
f_{\lambda,\eta}(z)=1-\lambda b_\eta(z)^2.}
\]

\textbf{Proof.} For \(z=x+iy\) in the right half-strip,
\[
\operatorname{Re}u_\eta'(z)
=\sum_mw_mk_m\sin(k_mx)\cosh(k_my)>0.
\]
The integral of \(u_\eta'\) over the segment between two points, divided by the difference of the endpoints, has positive real part. Convexity of the half-strip proves injectivity there. Evenness gives the corresponding result on the left.

Moreover,
\[
\operatorname{Im}u_\eta(x+iy)
=\sum_mw_m\sin(k_mx)\sinh(k_my)
\]
has the sign of \(y\) for \(x>0\). On the positive real axis \(u_\eta>0\); on the imaginary axis,
\[
u_\eta(iy)=\sum_mw_m(1-\cosh(k_my))
\]
is strictly negative away from zero and decreases strictly with \(|y|\). These separations, injectivity on each half-strip and evenness prove that the fibers of \(u_\eta\) in \(\mathcal S\) are exactly \(\{z,-z\}\), except for the fiber of the origin.

The function \(u_\eta(z)/z^2\), extended with value one at zero, is holomorphic and has no zeros in the simply connected strip. It has a holomorphic logarithm \(L_\eta\) with \(L_\eta(0)=0\), which is even by uniqueness of the normalization. Define
\[
b_\eta(z)=z\exp\!\left(\frac12L_\eta(z)\right).
\]
Then \(b_\eta^2=u_\eta\), \(b_\eta\) is odd and \(b_\eta'(0)=1\). The equality \(b_\eta(z)=b_\eta(w)\) forces \(w=z\) or \(w=-z\). The second case gives equality only at the origin. This proves its univalence. \(\square\)

The twelve phases thus determine an explicit conformal deformation of the quadratic fold. The complete dynamical conjugacy preserves that deformation:
\[
b_\eta\circ f_{\lambda,\eta}\circ b_\eta^{-1}(w)
=b_\eta(1-\lambda w^2),
\]
on their composition domains. The factorization proved and a conjugacy to the entire quadratic family are distinct assertions.


\subsection{Return, rescaling and memory}
\label{hmtfg:desarrollo:5}

Let \(a=-f(1)>0\), \(H_a(x)=-ax\), \(J=H_a([-1,1])\). When \(J\subset[-1,1]\) is admissible for the return, \(f^2(J)\subset J\) and the compositions are defined,
\[
\mathcal Rf=H_a^{-1}f^2H_a,\qquad
\mathcal Rf(x)=-a^{-1}f(f(-ax)).
\]
If in addition \(f(J)\subset(0,1]\), a unique central quadratic maximum is preserved:
\[
(\mathcal Rf)(0)=1,\qquad
(\mathcal Rf)''(0)=-a f'(1)f''(0)<0.
\]
Composition of Schwarzians gives, at regular points,
\[
S(\mathcal Rf)(x)=a^2
\left[(Sf)(f(-ax))f'(-ax)^2+(Sf)(-ax)\right]<0.
\]
Domain inclusions are part of the hypotheses and are checked at every scale.

Let \(v_\eta=(u_\eta|_{[0,1]})^{-1}\). The inverse branches are
\[
\beta_\sigma(y)=\sigma v_\eta((1-y)/\lambda),\quad
\sigma\in\{-1,+1\},
\]
with the critical record \(\beta_0(1)=0\). The step
\[
(x,w)\longmapsto(f(x),w\mathbin{\|}\sigma(x))
\]
is inverted on its image by reading the last sheet and applying \(\beta_\sigma\). The word preserves branch decisions; the previous value is reconstructed through the map equation.

The return groups two transitions. From the endpoint \(x'\) and its sheets \((\sigma_0,\sigma_1)\) one recovers
\[
x=H_a^{-1}\beta_{\sigma_0}
\!\left(\beta_{\sigma_1}(H_ax')\right).
\]
The intermediate points are also reconstructed. The APP–TRIT–TPK records of sheet, carry, route, boundary and memory are concatenated while preserving their identity. The sign of the analytic branch is an additional record: it does not replace those data with a binary label.

If \(C_N\) groups histories of length \(2N\) into pairs, the restrictions satisfy
\[
\pi^{\mathrm{ret}}_{N,M}C_N
=C_M\pi^f_{2N,2M},\qquad M\le N.
\]
Both sides remove the same terminal pairs and preserve the same inverse branches. Compatible transport of finite histories is thus proved.


\subsection{Accumulated scales and derivative of the operator}
\label{hmtfg:desarrollo:6}

Let \(c_n=f^{2^n}(0)\), \(c_0=1\) and \(H_n(x)=c_nx\). As long as the compositions are admissible and \(c_n\ne0\),
\[
\boxed{\mathcal R^nf=H_n^{-1}f^{2^n}H_n},\qquad
(\mathcal R^nf)(1)=\frac{c_{n+1}}{c_n}=-a_n.
\]
Induction follows from \(H_nH_{a_n}=H_{n+1}\). The convention \(a_n>0\) requires sign alternation of \(c_n\). At a superstable root of period \(2^N\), \(c_N=0\): that normalization becomes singular and the preceding levels or a controlled limit must be used.

For \(z=-ax\), \(w=f(z)\) and a perturbation \(h\), the complete derivative is
\[
\boxed{
D\mathcal R_f[h](x)=
-\frac{h(w)+f'(w)h(z)}a+
\frac{h(1)}a
\left[\mathcal Rf(x)-xf'(w)f'(z)\right].
}
\]
It is differentiated using \(\dot a=-h(1)\), \(\dot z=xh(1)\),
\(\dot w=h(z)+xf'(z)h(1)\) and the variation of the denominator. If \(h(0)=0\), this gives \(D\mathcal R_f[h](0)=0\). The initial HMT parameter direction is \(h=-u_\eta\).

This formula includes scale variation; freezing \(a\) would give a different linearized operator.


\subsection{Certified roots and itineraries through period 256}
\label{hmtfg:desarrollo:7}

Let \(F_n(\lambda)=f_{\lambda,0}^{2^n}(0)\). Compute
\[
x_0=p_0=0,\qquad x_{j+1}=1-\lambda u_0(x_j),\qquad
p_{j+1}=-u_0(x_j)-\lambda u_0'(x_j)p_j.
\]
Then \(F_n=x_{2^n}\) and \(F_n'=p_{2^n}\).

The attached program uses rational intervals with outward integer rounding, square roots enclosed by integer square roots, and sine/cosine with an explicit Taylor remainder. The exact bound \(|u_0''|\le2\) allows centered evaluation with a quadratic remainder.

The historical approximations only propose boxes of radius \(10^{-42}\). Opposite signs at their endpoints, separation of \(F_n'\) from zero throughout the box, and exclusion of returns with periods that properly divide \(2^n\) are certified. The intermediate value theorem and monotonicity prove local existence and uniqueness; exclusion of divisors proves the least period.

\begingroup\small
\begin{longtable}{p{.12\linewidth}p{.17\linewidth}p{.62\linewidth}}
\hline
Level & Least period & Box center, rounded \\
\hline
1 & 2 & 1.345913990471832792210949 \\
2 & 4 & 1.763379787846910127290794 \\
3 & 8 & 1.850413796950136464530567 \\
4 & 16 & 1.868979756606798125150574 \\
5 & 32 & 1.872955034435659740004205 \\
6 & 64 & 1.873806367467030119412262 \\
7 & 128 & 1.873988695221365362307169 \\
8 & 256 & 1.874027744154450672897008 \\
\hline\end{longtable}
\endgroup

The first root has the exact expression \(\lambda_1=1/u_0(1)\). The complete boxes and their checks are in the JSON certificate.

The 502 preclosure states of these eight levels, all separated from zero, and their complete sign words have been enclosed. In particular,
\[
\operatorname{sign}f_{\lambda_n,0}^{2^j}(0)=(-1)^j,\qquad0\le j<n.
\]
The finite orientations of the rescalings preceding closure are established. Inclusions of restrictive intervals at every scale constitute an additional check.

For \(\delta_{\mathrm F,n}=(\lambda_{n-1}-\lambda_{n-2})/(\lambda_n-\lambda_{n-1})\), an abbreviated outer decimal interval is
\[
\boxed{\begin{gathered}
4.66921218915020415455609621588966392804<\delta_{\mathrm F,8},\\
\delta_{\mathrm F,8}<4.66921218915020415455609621588966392863.
\end{gathered}}
\]
The width of the complete interval is less than \(5.808\cdot10^{-37}\). It measures the numerical error of the finite ratio \(\delta_{\mathrm F,8}\), distinct from the error in approximating the universal limit.


\subsection{Analytic persistence of the deformed roots}
\label{hmtfg:desarrollo:8}

For each certified level, \(F_n'(\lambda_n)\ne0\). The implicit function theorem produces a local analytic branch \(\lambda_n(\eta)\) with
\[
f_{\lambda_n(\eta),\eta}^{2^n}(0)=0.
\]
By continuity, the proper returns remain separated from zero for \(\eta\) sufficiently small; the least period and itinerary are preserved. The existence of those local neighborhoods is proved. Their common width is not yet identified with the entire structural strip \(1/40\).

With \(M=M_{2,\eta}\), the profile response is
\[
v_\eta(x)=\partial_\eta u_\eta(x)
=\sum_m\frac{p_m}{12}(1-\cos(k_{m,\eta}x))
-\frac{M'}{2M}\,x u_\eta'(x).
\]
The second term preserves normalization of the second moment. If \(q_j=\partial_\eta x_j\) at fixed \(\lambda\),
\[
q_0=0,\qquad q_{j+1}=-\lambda v_\eta(x_j)-\lambda u_\eta'(x_j)q_j,
\qquad
\boxed{\lambda_n'(\eta)=-q_{2^n}/p_{2^n}}.
\]
The weights–scale–response–superstable-parameter chain is thus explicit. This finite nondegeneracy is distinct from asymptotic transversality to the stable leaf of the universal fixed point.


\subsection{Analytic tails and error propagation}
\label{hmtfg:desarrollo:9}

For \(\eta=0\), let
\[
\mu_{2j}=\frac1{12}\sum_mk_m^{2j},\quad
p_N(z)=\sum_{j=1}^N(-1)^{j+1}\frac{\mu_{2j}}{(2j)!}z^{2j},
\quad\Omega=\frac{70}{\sqrt{899}}.
\]
If \(q_N=\Omega^2r^2/[(2N+3)(2N+4)]<1\),
\[
\|u_0-p_N\|_r\le
\frac{\mu_{2N+2}r^{2N+2}}{(2N+2)!(1-q_N)}.
\]
The inequality \(\mu_{2j+2}\le\Omega^2\mu_{2j}\) bounds the tail by a geometric series. For \(0\le s\le2N+2\),
\[
\|(u_0-p_N)^{(s)}\|_r\le
\frac{\mu_{2N+2}r^{2N+2-s}}{(2N+2-s)!}
\left(1-\frac{\Omega^2r^2}{(2N+3-s)(2N+4-s)}\right)^{-1},
\]
when the denominator is positive.

To transport errors through a return, its domains are also controlled: if \(\|f-g\|\le\varepsilon\), \(a_f,a_g\ge a_*>0\), the compositions lie in a common convex domain and \(\|f'\|\le L,\ \|g\|\le M\) there, then
\[
\|\mathcal Rf-\mathcal Rg\|_r
\le\varepsilon\left[\frac{1+L+L^2r}{a_*}+\frac{M}{a_*^2}\right].
\]
The proof adds the outer perturbation, the inner perturbation, the change of argument due to \(|a_f-a_g|\le\varepsilon\) and the variation of the denominator. These bounds prevent reuse of the initial cosine sum as if all renormalized profiles retained that same form.

% Fragmento integrable. Preambulo anfitrion: amsmath, amssymb, longtable, hyperref.
\section{Nonadic depth and persistence of the itinerary}
\label{hmtfg:fragmento:profundidad-nonadica-y-persistencia-del-itinerario}

\subsection{Outer growth and inner resolution}
\label{hmtfg:prolongacion:2}

The refinement law of article XI, with integer base \(B\ge2\), is
\[
L_c(x,y)=(Bx-c,By+c),\qquad 0\le c<B.
\]
For a word \(c_1\ldots c_m\), its record is
\[
C_m=\sum_{j=1}^m c_jB^{m-j},\qquad
(x_m,y_m)=(B^mx_0-C_m,B^my_0+C_m).
\]
The outer sum is \(B^mM_0\), with \(M_0=x_0+y_0\). The normalized reading of the record belongs to the cylinder
\[
I_m(C_m)=\left[\frac{C_m}{B^m},\frac{C_m+1}{B^m}\right],
\qquad \operatorname{diam}I_m=B^{-m}.
\]
The integer pair and the scale recover the quotients, remainders and leading zeros of the word. Endpoints with two expansions retain their boundary mark. Support growth and readout refinement are thus coordinated without identifying distinct histories solely by their limit image.

The composition is exact:
\[
L_D^{[b]}\circ L_C^{[a]}=L_{B^bC+D}^{[a+b]}.
\]
By associativity, grouping a history of eighteen events into nine blocks of two or two blocks of nine produces the same record. At return level \(2^N\), the common horizon is \(9\,2^N\). The proof is an equality of composition on the same ordered history; it retains its carries and its recoverable subblocks.

The restriction of states and the prolongation of their readers are expressed through the following limits:
\[
X_\infty=\varprojlim X_m,\qquad
r^*a=a\circ r,\qquad
r^*e_x=\sum_{r(y)=x}e_y,
\qquad \mathfrak H_{\rm cyl}=\varinjlim\mathfrak H_m.
\]
States restrict; readers extend. Pullback preserves products, the unit and idempotents. An exactly compatible family induces a reader on the direct limit. Approximants that are compatible only up to an error additionally require a convergence estimate in their completion. The next section provides that estimate for the critical return.

\subsubsection{Capacity and vacancies of the record}

The capacity reader of XI compares blocks of \(729\) and \(1000\) possibilities. With \(\rho=\log_{1000}729=\log_{10}9\), it retains
\[
\mathcal C_k=\lfloor(k+1)\rho\rfloor,\qquad
\mathcal C_k-\mathcal C_{k-1}=1-z_k,
\]
\[
\mathcal C_b-\mathcal C_a+\sum_{k=a+1}^b z_k=b-a,
\qquad T_{\rm cap}(a)=\lceil a/\rho\rceil-1\quad(a\ge1).
\]
Each vacancy records a transition during which capacity is maintained; chronological length is conserved in the balance. Capacity phase and chronological phase are distinct readers.

To store a prefix of \(m\) base-nine digits, a sufficient decimal capacity is
\[
a(m)=\left\lceil\frac{m\rho}{3}\right\rceil.
\]
The quotient and remainder allow the integer \(C_m\) to be converted between bases, keeping the depth \(m\) separate. The following hold
\[
9^m\le1000^{a(m)}\le729^{T_{\rm cap}(a(m))+1}.
\]
These inequalities are also decided by integer powers, without rounding logarithms. The readout budget of the next section is thus translated into record capacity while preserving units. This count sizes the storage; the survival rules additionally determine which native histories are admissible.


\subsection{Explicit depth budget for each return}
\label{hmtfg:prolongacion:3}

\textbf{Theorem.} With \(|\eta|\le1/40\) fixed, normalize
\[
t=\lambda u_\eta(1)/2\in[0,1],\qquad z=(x+1)/2,
\]
\[
F_{t,\eta}(z)=1-t\frac{u_\eta(2z-1)}{u_\eta(1)},\qquad z\in[0,1].
\]
For every integer \(q\ge1\) and two parameters \(t,s\),
\[
\left|F_{t,\eta}^{q}(1/2)-F_{s,\eta}^{q}(1/2)\right|
\le S_q|t-s|,
\qquad S_q=\sum_{j=0}^{q-1}(6\sqrt2)^j<9^q.
\]
Consequently, a base-nine parameter cylinder of depth \(m=2^N+r\) encloses the readout of return \(q=2^N\) in an interval of diameter less than \(9^{-r}\).

\textbf{Proof.} The function \((1-\cos v)/v^2\) decreases for \(0<v<\pi\): the sign of its derivative is that of \(v\sin v-2(1-\cos v)\), which is negative because \(\tan(v/2)>v/2\). By the second moment,
\[
u_\eta(1)\ge\frac{1-\cos3}{9}\sum_jw_jk_j^2
=\frac{2(1-\cos3)}9>\frac13.
\]
The last inequality is verified, for example, by bounding \(\cos3\) by its alternating Taylor sum through degree eight, which is less than \(-1/2\). Moreover, Cauchy–Schwarz gives \(|u_\eta'(x)|\le\sqrt2\) on the real line. It follows that
\[
\operatorname{Lip}_zF_{t,\eta}\le6\sqrt2<9,\qquad
\operatorname{Lip}_tF_{t,\eta}\le1.
\]
If \(d_j\) is the difference between the two orbits, then \(d_0=0\) and \(d_{j+1}\le6\sqrt2\,d_j+|t-s|\). Induction proves the geometric sum. For \(|t-s|\le9^{-m}\), the diameter is less than \(9^{q-m}\). Substituting \(q=2^N\) and \(m=q+r\) completes the proof. \(\square\)

The bound controls parameter uncertainty with \(\eta\) fixed. An implementation separately adds and propagates cosine-evaluation and rounding errors; the preceding rational certificate provides these operations. The bound is sufficient and conservative, with a cost that grows with the horizon.

\textbf{Limit reader.} Realize \(t\) from \((x_0,y_0)=(1,0)\), \(B=9\), so that
\[
(x_m,y_m)=(9^m-C_m,C_m),\qquad t_m=C_m/9^m.
\]
The readers \(E_{q,m}=F_{t_m,\eta}^{q}(1/2)\), brought to a common level, satisfy
\[
\|E_{q,m+a}-r^*E_{q,m}\|_\infty\le S_q9^{-m}.
\]
They converge uniformly to a continuous reader \(E_q\) in the completion of the cylinder algebra. The set-theoretic version
\[
J_{q,m}(C)=\{F_{t,\eta}^{q}(1/2):t\in I_m(C)\}
\]
is exactly nested and its diameters tend to zero. This result gives the passage to the limit of the \textbf{fixed-horizon readout}, with depth control at any requested horizon.


\subsection{Doubling word and nonadic reader of chronology}
\label{hmtfg:prolongacion:4}

The eight itineraries certified in the preceding work obey
\[
W_1=+,\qquad W_{n+1}=W_n\,(-1)^n\,W_n.
\]
Its unique prefix continuation is
\[
\tau_j=(-1)^{v_2(j)},\qquad j\ge1;
\qquad \tau_{2j}=-\tau_j,\quad\tau_{2j+1}=+.
\]
The proof uses \(v_2(2^n+r)=v_2(r)\) for \(0<r<2^n\). It is a recurrence at every depth, and agreement with the eight computed returns constitutes its finite check in the family.

In the parity-lexicographic order of a map with a critical maximum, the orientation factor preceding position \(k\) is \(\prod_{j<k}(-\tau_j)\). The word \(\tau\) is strictly greater than any of its positive shifts. Indeed, the first disagreement with an odd shift occurs at \(k=1\) or \(2\); an even shift halves the comparison and doubles \(k\). The first disagreement is therefore at \(k=2^r\). There
\[
\prod_{j<2^r}(-\tau_j)=(-1)^r=\tau_{2^r},
\]
and the oriented difference with the opposite sign is positive. This argument also proves aperiodicity.

The chronological record retains
\[
K=\rho_9(K)+9q_9(K).
\]
Its dyadic reader is
\[
\ell_n(B_K)=[\rho_9(K)+9q_9(K)]_{2^n}=[K]_{2^n}.
\]
On the chart of histories in which \(\Gamma_9 B_K=B_{K+9}\),
\[
\ell_n(\Gamma_9B_K)=\ell_n(B_K)+9\pmod{2^n},\qquad
\pi_{n+1,n}\ell_{n+1}=\ell_n.
\]
Since nine is odd, advancement by nine visits the \(2^n\) positions. Its conjugacy with unit advancement is \(j\mapsto9j\). For \(0<j<2^n\),
\[
v_2(9j\bmod2^n)=v_2(j),
\]
so that nonadic resampling also preserves the oriented word in each certified cycle.

This reader uses remainder \textbf{and} quotient. The phases of \(K=1\) and \(K=10\) agree modulo nine, while their dyadic signs are opposite. Memory distinguishes the two histories. The arithmetic property of resampling holds for every odd multiplier; the HMT calendar fixes the choice of nine here.


\subsection{Compatible bilateral memory at infinite depth}
\label{hmtfg:prolongacion:5}

Let \(\mathcal H_n=\ell^2(\mathbb Z/2^n\mathbb Z)\),
\[
C_ne_j=e_{j+1},\qquad
J_ne_j=(e_j+e_{j+2^n})/\sqrt2.
\]
A check on the basis gives \(J_n^*J_n=I\) and \(C_{n+1}J_n=J_nC_n\), including the boundary terms.

Article XI provides the projectors
\[
P_0=\frac19\begin{pmatrix}8&-\sqrt8\\-\sqrt8&1\end{pmatrix},\qquad
P_1=\frac19\begin{pmatrix}1&\sqrt8\\\sqrt8&8\end{pmatrix}.
\]
They are orthogonal and complementary. Therefore
\[
\mathbb U_n=P_0\otimes I+P_1\otimes C_n
=\begin{pmatrix}T_n&D_n\\D_n&R_n\end{pmatrix}
\]
is unitary, with
\[
T_n=(8I+C_n)/9,\quad D_n=\sqrt8(C_n-I)/9,\quad R_n=(I+8C_n)/9.
\]
The naturality already proved implies
\[
\mathbb U_{n+1}(I_2\otimes J_n)=(I_2\otimes J_n)\mathbb U_n,
\quad \mathbb U_n^a=P_0\otimes I+P_1\otimes C_n^a
\quad(a\in\mathbb Z).
\]
The Hilbert limit is identified with \(L^2(\mathbb Z_2,\mu_{\rm Haar})\), taking \(e_j^{(n)}\) to \(2^{n/2}\mathbf1_{j+2^n\mathbb Z_2}\). The identities extend by density and boundedness. For every vector \(v\) and every integer time \(a\),
\[
\boxed{\|T_{\infty,a}v\|^2+\|D_{\infty,a}v\|^2=\|v\|^2,}
\]
\[
\boxed{v=T_{\infty,a}^*(T_{\infty,a}v)+D_{\infty,a}^*(D_{\infty,a}v).}
\]
Here \(T_{\infty,a}=(8I+C_\infty^a)/9\). Bilateral composition preserves the archive between steps; composing only \(T_n\) describes a different operation.

The permutation \(S_ne_j=e_{9j\bmod2^n}\) satisfies
\[
S_{n+1}J_n=J_nS_n,\qquad S_nC_nS_n^{-1}=C_n^9.
\]
Thus one obtains in the limit
\[
S_\infty C_\infty S_\infty^{-1}=C_\infty^9,
\quad (I_2\otimes S_\infty)\mathbb U_\infty
(I_2\otimes S_\infty^{-1})=\mathbb U_\infty^9.
\]
The temporal representation is faithful: \(C_\infty^a=C_\infty^b\) forces \(2^n\mid a-b\) for every \(n\), whence \(a=b\). Particular preparations may preserve temporal symmetries. The operator identity preserves this distinction.

The projective point \(([K]_{2^n})_n\) and a vector of \(L^2\) have different types: a point history induces a Dirac measure, whereas normalizable amplitudes belong to the Hilbert space. The result constructs an equivariant reader of chronology; the remaining native coordinates of the TPK accompany the complete record.


\subsection{Existence and separation of asymptotic obligations}
\label{hmtfg:prolongacion:6}

\subsubsection{Realization of the infinite word throughout the strip}
\label{hmtfg:prolongacion:6.1}

\textbf{Existence theorem.} For each \(|\eta|\le1/40\) there exists at least one
\[
\lambda_\infty(\eta)\in
\left[\frac1{2u_\eta(1)},\frac2{u_\eta(1)}\right]
\]
such that
\[
\boxed{\operatorname{sign}f_{\lambda_\infty(\eta),\eta}^{j}(0)
=(-1)^{v_2(j)}\quad\text{for every }j\ge1.}
\]
The conclusion realizes the infinite doubling itinerary within the fixed family. The existence quantifier remains separate from a unique selection and from the law of distances between superstable parameters.

\textbf{Proof.} We fix \(\eta\). We write \(K_\lambda\) for the sign word of the critical orbit; when it first returns to zero, the word ends at that zero. We use the unimodal order with a maximum, whose factor before comparing symbol \(j\) is the product of the signs \(-K_i\), \(i<j\).

At the left endpoint the image of the entire interval is in \([1/2,1]\); hence \(K_-=+++\cdots\). At the right endpoint the orbit is \(0\mapsto1\mapsto-1\mapsto-1\), whence \(K_+=+---\cdots\). Comparing the second and third symbols, respectively, gives
\[
K_-<\tau<K_+.
\]

If \(K_{\lambda_0}\) is infinite and distinct from \(\tau\), the first difference occurs at a finite position. Continuity of the corresponding iterates makes the side of the comparison locally constant. It remains to examine a superstable parameter of first period \(p\), with word \(W0\), where \(|W|=p-1\).

For \(\lambda\) near \(\lambda_0\), let \(g_\lambda=f_{\lambda,\eta}^p\) and \(a=g_\lambda(0)\). We have \(g_\lambda'(0)=0\), and a uniform local bound on the second derivative gives
\[
g_\lambda(x)=a+O(x^2).
\]
For \(|a|\) sufficiently small and nonzero, \(g_\lambda\) maps \([-2|a|,2|a|]\) into \([a-|a|/2,a+|a|/2]\). The returns retain the sign of \(a\), and the \(p-1\) intermediate positions retain \(W\). The neighboring words are therefore \((W+)^\infty\) or \((W-)^\infty\); parameters with \(a=0\) retain \(W0\).

If \(\tau\) differs from \(W\) before position \(p\), that difference fixes the same side for all three words. If it shares \(W\), put
\[
s=\tau_p,\qquad e=\prod_{j<p}(-\tau_j),\qquad A=Ws.
\]
The sign of the comparison of \(\tau\) with \(W0\) and with \((W,-s)^\infty\) is \(es\). To compare it with \(A^\infty\), let \(q\) be the first position at which \(\tau_{p+q}\ne\tau_q\). It exists by aperiodicity. Up to position \(p+q-1\), both words agree, and the orientation sign factors as \(O_pO_{q-1}\), with \(O_p=-es\). The strict maximality already proved gives
\[
O_{q-1}(\tau_q-\tau_{p+q})>0.
\]
Consequently, the comparison of \(\tau\) with \(A^\infty\) has sign \(-O_p=es\). The three local possibilities lie on the same side of \(\tau\).

If no parameter realized \(\tau\), the sets \(\{\lambda:K_\lambda<\tau\}\) and \(\{\lambda:K_\lambda>\tau\}\) would be open, disjoint, nonempty and would cover the parameter interval. Connectedness of the interval excludes that partition. The stated parameter therefore exists. \(\square\)

This itinerary-continuity argument applies after the HMT generator and its analytic family have been fixed; its hypotheses are verified directly here. The construction preserves the infinite word without introducing a parameter of the quadratic family or the universal value as an input.

\subsubsection{Margins preventing loss of the itinerary when taking limits}
\label{hmtfg:prolongacion:6.2}

Write \(c_p=f_{\lambda,\eta}^p(0)\). A sufficient uniform rational bound is
\[
|f'|,|f''|\le16,\qquad
B_p:=16^p\frac{16^p-1}{15}\ge\|(f^p)''\|_{[-1,1]}.
\]
Indeed, \(1-\cos v\ge v^2/2-v^4/24\ge v^2/8\) for \(|v|\le3\), so that \(u_\eta(1)\ge1/4\), \(\lambda\le8\) and both derivatives are bounded by \(16\), using the moments. The chain rule proves the bound \(B_p\) by induction.

The word \(\tau\) satisfies \(\tau_{2p}=-\tau_p\) for every \(p\). For any of its realizations, \(c_{2p}\) and \(c_p\) have opposite signs. Since \((f^p)'(0)=0\), Taylor gives
\[
|c_{2p}-c_p|\le\tfrac12B_p|c_p|^2.
\]
The left-hand side is greater than \(|c_p|\), and therefore
\[
\boxed{|c_p|>2/B_p>0.}
\]
For each fixed horizon there is thus a uniform margin preventing a sequence of realizations from losing its sign upon convergence. In the compact coordinates \((\eta,b)\), \(b=\lambda u_\eta(1)\in[1/2,2]\), the set of realizations of \(\tau\) is closed and compact, and its projection onto the entire \(\eta\) strip is surjective.

The same argument applies to a prefix tree: if all prefixes of \(\tau\) are realized, a subsequence in the compact set preserves every sign, because sufficiently long prefixes simultaneously include positions \(p\) and \(2p\). The existence theorem now provides the nonemptiness that this argument requires.

\subsubsection{Invariant return intervals at every level}
\label{hmtfg:prolongacion:6.3}

\textbf{Nested returns theorem.} For any of the realizations of the existence theorem in section \ref{hmtfg:prolongacion:6.1}, let \(c_j=f^j(0)\) and
\[
J_n=\operatorname{conv}\{c_{2^n},c_{2^{n+1}}\},\qquad n\ge1.
\]
Then \(0\in\operatorname{int}J_n\), \(J_{n+1}\subset J_n\),
\[
f^{2^n}(J_n)=J_n,
\]
and the \(2^n\) images \(J_n,f(J_n),\ldots,f^{2^n-1}(J_n)\) are disjoint. The normalized returns are unimodal with a maximum, with critical value one and the same word \(\tau\). These intervals are invariant return cores; a normalization requiring periodic endpoints may enlarge the core to the corresponding maximal restrictive interval and must specify that step.

\textbf{Proof of the first return.} Let \(F\) be a unimodal map with a maximum, with \(F(0)=1\) and word \(\tau\), and \(d_j=F^j(0)\). Its first signs are \(+,-,+,+,+,-,+\). The interval \(I=[d_2,1]\) is invariant: \(F(d_2)=d_3>0\), \(F(1)=d_2<0\), and the maximum is one. Put \(J=[d_2,d_4]\). Since \(F\) decreases on positive arguments and
\[
F(d_3)=d_4>0> d_6=F(d_5),
\]
we have \(d_3<d_5\), and therefore \(F(J)=[d_3,1]\).

Moreover, \(d_3>d_4\). To check this, \(F^2\) is strictly increasing between the two positive points \(d_3,d_4\), because their images \(d_4,d_5\) are positive. Its values are \(F^2(d_3)=d_5>0>d_6=F^2(d_4)\), which forces that order. Thus \(J\) and \(F(J)\) are separated by an open interval, and
\[
F^2(J)=F([d_3,1])=[d_2,d_4]=J.
\]
The map \(F^2|_J\) has a unique critical minimum: the image \(F(J)\) remains positive. The conjugacy \(h(x)=d_2x\) reverses orientation and normalizes the return as a maximum with value one. Its itinerary is
\[
\operatorname{sign}\frac{d_{2j}}{d_2}
=-\tau_{2j}=\tau_j.
\]
The same argument can be repeated at any depth.

\textbf{Inductive disjointness.} Suppose the \(p=2^n\) images of the parent \(J_n\) are disjoint, and let \(g=f^p\). The first return applied to the normalized map produces \(B=J_{n+1}\), \(A=g(B)\), with \(A\cap B=\varnothing\), \(g(B)=A\), \(g(A)=B\). For \(0\le j<p\), an intersection of \(f^j(A)\) and \(f^j(B)\) would, upon applying \(f^{p-j}\), produce an intersection of \(B\) and \(A\). Those two images are therefore disjoint. Images corresponding to distinct \(j\) belong to disjoint images of the parent. This proves the \(2p\) disjoint images and the return \(f^{2p}(B)=B\). The endpoint formula results from composing the normalizations \(x\mapsto c_{2^n}x\). \(\square\)

The compact sets
\[
\Lambda_n=\bigcup_{0\le j<2^n}f^j(J_n)
\]
are nonempty and nested. Their intersection \(\Lambda_\infty\) has a continuous reader onto the dyadic odometer: each point is assigned the label of the unique component it occupies at each level. Every compatible sequence of labels has a preimage by compactness, and advancement by \(f\) adds one to the labels. This map is a surjective semiconjugacy. Promoting it to a pointwise conjugacy requires proving that the component diameters tend to zero; that metric promotion remains separate.

% Fragmento integrable. Preambulo anfitrion: amsmath, amssymb, longtable, hyperref.
\section{Compatibility of readers and recovery of the parameter}
\label{hmtfg:fragmento:compatibilidad-de-lectores-y-restitucion-del-parametro}

\subsection{Compatibility of readers, minima and recovery of history}
\label{hmtfg:escalado:13}

The prolongation operations and ordered readers preserve the information required for root selection. The following identities specify their domains and composition.

\subsubsection{Recovered forward rule}
\label{hmtfg:escalado:13.1}

The deterministic update starts from the R36 boundary state with character \(\chi\) and internal remainder. The update is
\begin{equation}
\mathcal U_K^\chi=
\operatorname{Upd}_K^\chi\circ
\operatorname{Tra}_K^\chi\circ
\operatorname{Sel}_K^\chi,
\qquad
N_{K+1}=729N_K+d_K^\chi.
\label{hmtfg:escalado:eq:26}
\end{equation}
The digit and the update depend on the preceding state. Two forward sections with the same complete R36 state coincide by induction. The subrelation \(\operatorname{Ext}^{\chi,\rightarrow}\) is the graph of that update; the total relation \(\operatorname{Ext}\) admits multiple continuations depending on the data and characters.

Compatible prolongations transport the regional readers and the joint signature over the same prefixes.

These rules transport orientation, remainder, quotient, carry, boundary and memory. Their application to the cascade must also preserve ordered selection of the critical reader’s parameter.

\subsubsection{Independence of the antecedent constants from the free parameter}
\label{hmtfg:escalado:13.2}

Let \(B\) be the HMT state providing the wheel and its antecedent constants. Write \(\mathcal C(B)\) for the set of those readouts and \(u_B\) for the critical profile. In the domain of the critical reader, the subsequent construction is
\begin{equation}
(B,\lambda)\longmapsto f_{B,\lambda}(x)=1-\lambda u_B(x),
\qquad 0<\lambda\le 2/u_B(1).
\label{hmtfg:escalado:eq:27}
\end{equation}
For the same \(B\), the reading of the preceding constants factors through projection onto the first component:
\begin{equation}
\widetilde{\mathcal C}(B,\lambda)
=\mathcal C\circ\operatorname{pr}_B(B,\lambda)
=\mathcal C(B).
\label{hmtfg:escalado:eq:28}
\end{equation}
Consequently, two distinct values of \(\lambda\) preserve the same antecedent readouts. The equality is a property of this causal construction: the profile is fixed before \(\lambda\) is varied, and the factor \(\pi_{\mathrm{HMT}}\) cancels in its quadratic normalization.

\textbf{Lemma.} Suppose a condition \(A\) is evaluated solely on those antecedent readouts. For every fixed \(B\),
\begin{equation}
\{\lambda\in\Lambda:A(\widetilde{\mathcal C}(B,\lambda))\}
=
\begin{cases}
\Lambda,&A(\mathcal C(B))\text{ holds},\\
\varnothing,&A(\mathcal C(B))\text{ fails}.
\end{cases}
\label{hmtfg:escalado:eq:29}
\end{equation}
\textbf{Proof.} Substituting \eqref{hmtfg:escalado:eq:28} makes the predicate independent of \(\lambda\). Therefore all parameters in the domain satisfy it, or none do. This includes comparison of the antecedent constants with metrological intervals.

Uniqueness of the section \(B\) and the subsequent restriction \(f_{B,\lambda}^{2^n}(0)=0\) produce a solution fiber in \(\lambda\). The first-root rule acts on this fiber through the return and its itinerary. Selection and its prolongation are proved in the following results.

This result maintains the HMT order: metrology compares previously constructed outputs, and root selection is proved within the reader that produces it.

\subsubsection{Lemma on ordered transport of the selector}
\label{hmtfg:escalado:13.3}

Let \(P_n\) be an ordered set of genealogical candidates with minimum \(p_n^*\). Let \(Z_n\) be the corpus’s original set of candidate roots, with the real order, and let \(L_n:P_n\to Z_n\) be a monotone reader. Its image is said to be \textbf{coinitial} when
\begin{equation}
\forall z\in Z_n\quad\exists p\in P_n:
L_n(p)\le z.
\label{hmtfg:escalado:eq:30}
\end{equation}
Then
\begin{equation}
\boxed{L_n(p_n^*)=\min Z_n
\quad\Longleftrightarrow\quad
L_n(P_n)\text{ is coinitial in }Z_n.}
\label{hmtfg:escalado:eq:31}
\end{equation}
\textbf{Proof.} If the image of the minimum is the global minimum, take \(p=p_n^*\) in \eqref{hmtfg:escalado:eq:30}. Conversely, given \(z\in Z_n\), coinitiality provides \(p\) with \(L_n(p)\le z\). Monotonicity gives \(L_n(p_n^*)\le L_n(p)\le z\). Since \(L_n(p_n^*)\in Z_n\), it is the global minimum. Surjectivity of \(L_n\) is a sufficient condition, stronger than \eqref{hmtfg:escalado:eq:30}.

The ordered isomorphism of the continuum chart satisfies reader surjectivity. The concrete application to the return and its complete set of zeros is proved in the ordered-transport section; classification and isolation then determine its dynamical continuation.

\subsubsection{Minimal selection through survival at arbitrary depth}
\label{hmtfg:escalado:13.6}

Composition with survival also produces a precise global selector. The uniform family \(\eta=0\), constructed in section~\ref{hmtfg:escalado:1}, is retained, and the theorems on realization of the infinite itinerary and preservation of its signs are used. Let
\[
I=\left[\frac1{2u_0(1)},\frac2{u_0(1)}\right],\qquad
c_j(\lambda)=f_\lambda^j(0),\qquad
\tau_j=(-1)^{v_2(j)}.
\]
The set
\[
\Lambda_\tau=\{\lambda\in I:\operatorname{sign}c_j(\lambda)=\tau_j
\text{ for every }j\ge1\}
\]
is compact and nonempty. The derivative control in that source provides
\begin{equation}
B_j=\frac{16^j(16^j-1)}{15},\qquad
|c_j(\lambda)|>\frac2{B_j}\quad(\lambda\in\Lambda_\tau).
\label{hmtfg:escalado:eq:32}
\end{equation}
Define
\begin{equation}
\varepsilon_j=\frac1{B_j},\qquad
A_N=\{\lambda\in I:\tau_jc_j(\lambda)\ge\varepsilon_j,\
1\le j\le N\}.
\label{hmtfg:escalado:eq:33}
\end{equation}
Each \(A_N\) is compact by continuity of the iterates and nonempty because it contains \(\Lambda_\tau\). Moreover,
\begin{equation}
A_{N+1}\subseteq A_N,\qquad
\bigcap_{N\ge1}A_N=\Lambda_\tau.
\label{hmtfg:escalado:eq:34}
\end{equation}
The right-to-left inclusion follows from \eqref{hmtfg:escalado:eq:32}. The reverse inclusion follows from the strictly positive margins in \eqref{hmtfg:escalado:eq:33}, which fix each sign.

\textbf{Surviving-minimum theorem.} The minima \(a_N=\min A_N\) satisfy
\begin{equation}
\boxed{a_N\uparrow a_*=\min\Lambda_\tau.}
\label{hmtfg:escalado:eq:35}
\end{equation}
\textbf{Proof.} The inclusion \eqref{hmtfg:escalado:eq:34} makes the sequence of minima increasing and bounded within \(I\); it therefore has a limit \(a\). For each \(r\), all terms \(a_N\), \(N\ge r\), belong to the closed set \(A_r\); thus \(a\in A_r\). By \eqref{hmtfg:escalado:eq:34}, \(a\in\Lambda_\tau\). For any \(\lambda\in\Lambda_\tau\), we have \(a_N\le\lambda\) at every level, and on passing to the limit, \(a\le\lambda\). This proves \eqref{hmtfg:escalado:eq:35}.

In the HMT cylindrical reader, the object \(a_*\) has compatible prefixes at every depth, with orientation, remainder and memory preserved. Its choice uses the order of the parameter reading and complete survival. Compactness proves the existence and convergence of these minima.

The parameters \(a_N\) solve sign constraints with a margin. Their boundaries may satisfy \(\tau_jc_j=\varepsilon_j\), whereas the superstable roots satisfy \(c_{2^n}=0\). Therefore \eqref{hmtfg:escalado:eq:35} establishes a global selector for infinite realization and preserves, as a distinct object, the corpus’s first-root selector.

\subsubsection{Recovery of the parameter from the complete critical history}
\label{hmtfg:escalado:13.7}

The second term of the critical history provides an exact inverse reader:
\begin{equation}
c_1(\lambda)=1,\qquad
c_2(\lambda)=1-\lambda u_0(1),\qquad
\boxed{\lambda=\frac{1-c_2(\lambda)}{u_0(1)}.}
\label{hmtfg:escalado:eq:36}
\end{equation}
The positivity \(u_0(1)>0\), proved for the profile, makes division valid. Consequently, two identical complete critical histories have the same parameter. The sign word \(\tau\) is a reduced readout of that history; it preserves its combinatorics and omits the magnitude of \(c_2\).

Formulas \eqref{hmtfg:escalado:eq:35}–\eqref{hmtfg:escalado:eq:36} unite survival-based selection and recovery of the parameter from the critical history. Equality of its limit with the stable crossing is proved by first exit, and identification of the centers by analytic transport of their hybrid classes.

% Fragmento integrable. Preambulo anfitrion: amsmath, amssymb, longtable, hyperref.
\section{Ordered transport and Jacobi representation}
\label{hmtfg:fragmento:transporte-ordenado-y-representacion-de-jacobi}

\subsection{Complete transport from the continuum and spectral realization of the profile}
\label{hmtfg:escalado:14}

The ordered transport of the continuum and the spectral representation of the reader are composed before parameter selection.

\subsubsection{Coverage of all roots and transport of their minimum}
\label{hmtfg:escalado:14.1}

The ordered construction of the continuum developed in the preceding works produces the surjective Archimedean readout from compatible cylinders and their ordered quotient. Within its declared universe \(\mathfrak G\), it proves the isomorphism of complete ordered fields
\begin{equation}
\iota:\mathbb R_{\rm HMT}^{\mathfrak G}
   \xrightarrow{\cong}\mathbb R^{\mathfrak G}.
\label{hmtfg:escalado:eq:37}
\end{equation}
The name \(\eta\) used in the source work is changed here to \(\iota\), to distinguish it from the weighting parameter, fixed at zero. The second projection preserves incidence, boundary and memory over the same antecedent; the evaluation quotient and the complete history retain their respective domains.

Let us construct in the HMT field, before selecting a root,
\[
\cos_{\rm H}(t)=\sum_{k=0}^{\infty}\frac{(-1)^kt^{2k}}{(2k)!},
\qquad
u_{\rm H}(x)=\frac1{12}\sum_{m=1}^{12}
\left[1-\cos_{\rm H}\!\left(\frac{2(3m-1)x}{\sqrt[\rm H]{899}}\right)\right],
\]
\begin{equation}
F_a(x)=1-a\,u_{\rm H}(x),\qquad S_n^{\rm H}(a)=F_a^{\,2^n}(0).
\label{hmtfg:escalado:eq:38}
\end{equation}
The integers and the second moment are those already derived from the dodecaphase reader. The square root is the positive one in the ordered field. The convergent series analytically realizes that reader; it preserves the profile and the normalizing coefficient of section~\ref{hmtfg:escalado:1}.

\textbf{Return transport theorem.} For every parameter in the internal domain of \eqref{hmtfg:escalado:eq:38},
\begin{equation}
\iota(F_a(x))=f_{\iota(a)}(\iota(x)),\qquad
\iota(S_n^{\rm H}(a))=S_n(\iota(a)).
\label{hmtfg:escalado:eq:39}
\end{equation}
\textbf{Proof.} The ordered isomorphism fixes the rationals and preserves the positive square root. It preserves convergent limits: every neighborhood of positive rational radius is transported to another of the same radius. Applied to the partial sums of \eqref{hmtfg:escalado:eq:38}, it transports the cosine and \(u_{\rm H}\). The first equality follows from preservation of addition and multiplication; the second follows by induction on the number of iterations.

With a previous center \(a_{n-1}\) fixed, define \(Z_n^{\rm H}\) by zero return, exact critical period \(2^n\) and \(a>a_{n-1}\), within the declared parameter domain. Let \(Z_n\) be the set defined by the same conditions for \(f_\lambda\), to the right of \(\iota(a_{n-1})\). Then
\begin{equation}
\boxed{\iota[Z_n^{\rm H}]=Z_n,\qquad
\iota(\min Z_n^{\rm H})=\min Z_n.}
\label{hmtfg:escalado:eq:40}
\end{equation}
The second equality is asserted when the minimum exists; the first also preserves the nonexistence of candidates.

\textbf{Proof.} \eqref{hmtfg:escalado:eq:39} preserves the zero return and the earlier returns determining the exact period. Order preserves position relative to the previous center. Surjectivity of \eqref{hmtfg:escalado:eq:37}, applied to any root parameter, proves reverse coverage. If \(a_*\) is the minimum, every \(z\in Z_n\) has the form \(\iota(a)\), with \(a\in Z_n^{\rm H}\), whence \(\iota(a_*)\le z\). The converse assertion uses \(\iota^{-1}\).

Surjectivity covers every parameter in the chart, including those whose roots have not been isolated by a finite computation. The isomorphism preserves the declared universe of interpretation and the selection order.

\subsubsection{Jacobi representation of the complete dodecaphase profile}
\label{hmtfg:escalado:14.2}

Define the atomic measure of the reader,
\begin{equation}
s_m=\frac{4(3m-1)^2}{899},\qquad
\nu_0=\frac1{12}\sum_{m=1}^{12}\delta_{s_m},\qquad
M_r=\int s^r\,d\nu_0(s).
\label{hmtfg:escalado:eq:41}
\end{equation}
The twelve nodes are distinct and positive. For a polynomial \(p\ne0\) of degree less than \(r\le12\),
\begin{equation}
\sum_{i,j=0}^{r-1}p_iM_{i+j}p_j
=\frac1{12}\sum_{m=1}^{12}p(s_m)^2>0.
\label{hmtfg:escalado:eq:42}
\end{equation}
A polynomial of degree less than twelve does not vanish at all twelve nodes. At degree twelve, \(\prod_m(s-s_m)\) appears, with zero norm.

Orthogonalization with respect to this measure produces orthonormal polynomials \(\psi_0,\ldots,\psi_{11}\), with \(\psi_0=1\). Let
\[
Q_{mj}=\psi_j(s_m)/\sqrt{12},\quad
D=\operatorname{diag}(s_m),\quad J_{12}=Q^{\mathsf T}DQ.
\]
The following hold
\begin{equation}
Q^{\mathsf T}Q=I,\quad DQ=QJ_{12},\qquad
\boxed{u_0(x)=e_0^{\mathsf T}
\bigl[I-\cos(x\sqrt{J_{12}})\bigr]e_0.}
\label{hmtfg:escalado:eq:43}
\end{equation}
\textbf{Proof.} Orthogonality follows from \eqref{hmtfg:escalado:eq:42}. Multiplication by \(s\) produces a three-term recurrence: entries separated by more than one position vanish by orthogonality and degree. Choosing positive leading coefficients makes the subdiagonals positive. By functional calculus,
\(\cos(x\sqrt{J_{12}})=Q^{\mathsf T}\cos(x\sqrt D)Q\).
The vector \(Qe_0\) is constant, with value \(1/\sqrt{12}\); substituting it recovers the sum in section~\ref{hmtfg:escalado:1}.

The article on Grassmannians develops this mechanism for marked measures and refinements. Here it is applied to \eqref{hmtfg:escalado:eq:41}; rank twelve belongs to this reader. The genealogical marks are retained alongside the reader, in accordance with the marked faithfulness of the source work. \(J_{12}\) represents the analytic profile; the complete state additionally preserves routes, orientation and memory.

The verifier for this extension checks, with exact fractions, the twelve degrees of positive norm, vanishing at degree twelve, the intertwiner \(VT=DV\) in the monic basis and weighted self-adjointness \(HT=T^{\mathsf T}H\). It checks moments of degree zero through thirty; the identity for every degree follows from the intertwiner. Its first coefficients are
\begin{equation}
a_0=2,\qquad \beta_1=\frac{2493348}{808201}.
\label{hmtfg:escalado:eq:44}
\end{equation}

\subsubsection{Oriented sensitivity of the return}
\label{hmtfg:escalado:14.3}

For an exact return of period \(p=2^n\), let
\[
c_j=f_\lambda^j(0),\quad q_j=\partial_\lambda c_j,\quad
D_j=\prod_{k=1}^j f_\lambda'(c_k),\quad D_0=1.
\]
Before the critical return, the factors of \(D_{p-1}\) are nonzero. Differentiation and telescoping give
\begin{equation}
q_{j+1}=-u_0(c_j)+f_\lambda'(c_j)q_j,\qquad
\boxed{\frac{q_p}{D_{p-1}}
=-\sum_{j=1}^{p-1}\frac{u_0(c_j)}{D_j}.}
\label{hmtfg:escalado:eq:45}
\end{equation}
In the canonical word, \(\operatorname{sign}c_j=(-1)^{v_2(j)}\). The sign of \(f_\lambda'(c_j)\) is opposite; therefore,
\begin{equation}
\operatorname{sign}D_j=(-1)^{j+\sum_{k=1}^jv_2(k)}
=(-1)^{s_2(j)},
\label{hmtfg:escalado:eq:46}
\end{equation}
using \(\sum_{k=1}^jv_2(k)=j-s_2(j)\), where \(s_2(j)\) counts the binary ones. The normalized sensitivity is
\begin{equation}
\mathcal T_n=-\sum_{j=1}^{2^n-1}(-1)^{s_2(j)}t_j,\qquad
t_j=\frac{u_0(c_j)}{|D_j|}>0.
\label{hmtfg:escalado:eq:47}
\end{equation}
Nodal positivity of \eqref{hmtfg:escalado:eq:42} determines \(u_0\); \eqref{hmtfg:escalado:eq:47} additionally retains the orbital products and their orientations.

The additional moment representation would require a positive measure \(\rho\) on \([0,1]\) with \(t_j=\int z^{j-1}\,d\rho\). Under that condition,
\begin{equation}
\mathcal T_n=\int_0^1
\frac{1-\prod_{r=0}^{n-1}(1-z^{2^r})}{z}\,d\rho(z)>0
\label{hmtfg:escalado:eq:48}
\end{equation}
for \(\rho\ne0\), extending the integrand by its value \(1\) at zero. The identity is obtained by expanding the binary product.

The rational check at the previously certified period-four root obtains
\[
t_1\approx0.40425224267600139730,\quad
t_2\approx0.04925249167432646403,\quad
t_3\approx0.15740870098294833737
\]
and the rigorous outer interval
\begin{equation}
\begin{gathered}
0.296096033367379523967764444238832345360107<\mathcal T_2,\\
\mathcal T_2<0.296096033367379523967764444238832345360112.
\end{gathered}
\label{hmtfg:escalado:eq:49}
\end{equation}
It also certifies \(t_3>t_2\). A positive measure on \([0,1]\) would impose \(t_3\le t_2\); that specific lifting of orbital weights to positive moments is ruled out. The positive sign of \eqref{hmtfg:escalado:eq:49} remains certified. This check distinguishes the failure of a proposed proof from the failure of the property it sought to prove.

% Fragmento integrable. Preambulo anfitrion: amsmath, amssymb, longtable, hyperref.
\section{Classification and continuation of the first-root selector}
\label{hmtfg:fragmento:clasificacion-y-continuacion-del-selector-de-primeras-raices}

\subsection{Object and provenance}
\label{hmtfg:clasificacion:1}

This appendix preserves the family rendered by the HMT dodecaphase reader:

\[
u_0(x)=\frac1{12}\sum_{m=1}^{12}\left(1-\cos\frac{2(3m-1)x}{\sqrt{899}}\right),
\qquad f_\lambda(x)=1-\lambda u_0(x).
\]

The family originates in the operator and analytic antecedents presented in this development. The existence and compactness of its realizations of the infinite itinerary were proved in the subsections on realization of the word and preservation of its signs.

The following argument is original and uses only those properties and the return operations of the fixed family. It does not use a parameter from another family or the value of a universal constant as an input. The sheets of the itineraries used here are the two monotone branches of the critical reader; the remaining coordinates and HMT memory retain their residence in the state from which the reader originates.

We write
\[
c_j(\lambda)=f_\lambda^j(0),\qquad
\tau_j=(-1)^{v_2(j)},\qquad
W_n=\tau_1\cdots\tau_{2^n-1}.
\]
The itinerary of a critical orbit that returns to the origin for the first time ends in the zero symbol. The itinerary order is the unimodal order with a maximum: before comparing the symbol of index \(j\), the orientation factor is
\[
O_{j-1}(K)=\prod_{i=1}^{j-1}(-K_i).
\]
The ordinary order of the symbols is \(-1<0<1\).


\subsection{Classification before the infinite word}
\label{hmtfg:clasificacion:2}

\textbf{Classification theorem.} Let \(F:I\to I\), with \(I=[\ell,1]\), \(\ell<0\), be a continuous map, strictly increasing to the left of zero and strictly decreasing to its right, whose maximum satisfies \(F(0)=1\). Suppose zero has exact period \(2^n\), with \(n\ge1\), and its critical itinerary \(K\) satisfies \(K<\tau\). Then
\[
\boxed{K=W_n0.}
\]
The conclusion concerns the itinerary and allows distinct parameters in a family to have that same itinerary.

\textbf{Proof.} Write \(d_j=F^j(0)\). For period two the itinerary is \(+0=W_10\). Consider an exact period \(p=2^n\ge4\).

We have \(d_2<0\). Indeed, \(d_2=0\) would give period two. If \(d_2>0\), monotonicity of the positive branch gives
\[
F([d_2,1])=[d_2,d_3]\subset[d_2,1],
\]
and the critical orbit remains in a strictly positive interval, incompatible with a return to zero.

The initial signs are then \(+,-\). The orientation factor for comparing the third sign is \(-1\); hence \(K<\tau\) forces \(d_3>0\). A zero at that position would produce period three.

Nor can \(d_4<0\) occur. Under this hypothesis,
\[
d_2<d_4<0,
\qquad F([d_2,d_4])=[d_3,d_5]\subset(0,1],
\]
since \(d_4=F(d_3)>F(1)=d_2\) and \(d_5=F(d_4)>F(d_2)=d_3\). The positive branch is decreasing, so that
\[
F^2([d_2,d_4])=[d_6,d_4]\subset[d_2,d_4].
\]
Here \(d_6\ge d_2\), because \(d_5\le1\), and \(d_6<d_4\), because \(d_5>d_3\). This strictly negative interval traps the even returns of the critical orbit and excludes a return to zero. Therefore \(d_4\ge0\). If \(p=4\), one obtains precisely \(+-+0=W_20\).

Now suppose \(p\ge8\). Then \(d_4>0\). The orientation factors for the fifth and sixth signs are, respectively, \(-1\) and \(+1\). Since \(\tau_5=+1\) and \(\tau_6=-1\), the inequality \(K<\tau\) forces
\begin{equation}
\operatorname{sign}(d_1,\ldots,d_6)=(+,-,+,+,+,-).
\label{hmtfg:clasificacion:eq:1}
\end{equation}
Zero signs at those positions are excluded by the exact period.

We construct the return interval
\[
J=[d_2,d_4].
\]
The inequality \(F(d_3)=d_4>0>d_6=F(d_5)\), together with \(d_3,d_5>0\), implies \(d_3<d_5\); therefore
\[
F(J)=[d_3,1].
\]
Moreover, \(d_4<d_3\). To check this, the image under \(F\) of the interval with endpoints \(d_3,d_4\) is strictly positive, since its endpoint images are \(d_4,d_5>0\). On that interval \(F^2\) is strictly increasing. Its endpoint values satisfy
\[
F^2(d_3)=d_5>0>d_6=F^2(d_4),
\]
which forces \(d_3>d_4\). Consequently,
\begin{equation}
J\cap F(J)=\varnothing,
\qquad F^2(J)=F([d_3,1])=[d_2,d_4]=J.
\label{hmtfg:clasificacion:eq:2}
\end{equation}

The map \(F^2|_J\) has a unique minimum at zero. Conjugacy by \(h(x)=d_2x\), which reverses orientation, defines
\begin{equation}
G=h^{-1}\circ F^2\circ h:
\left[\frac{d_4}{d_2},1\right]\longrightarrow
\left[\frac{d_4}{d_2},1\right].
\label{hmtfg:clasificacion:eq:3}
\end{equation}
This map satisfies the same maximum hypotheses as \(F\), with \(G(0)=1\), and its critical point has exact period \(p/2\). Its itinerary is
\begin{equation}
K'_j=-K_{2j}.
\label{hmtfg:clasificacion:eq:4}
\end{equation}
All odd-indexed symbols of \(K\) are positive, because the even positions belong to \(J\) and the odd positions to \(F(J)\subset(0,1]\).

The transformation \eqref{hmtfg:clasificacion:eq:4} preserves the comparison order with \(\tau\). Indeed, \(\tau_{2j-1}=+1\) and \(\tau_{2j}=-\tau_j\). If \(q\) is the first position at which \(K'\) and \(\tau\) differ, the first position of difference between \(K\) and \(\tau\) is \(2q\), and
\begin{equation}
O_{2q-1}(K)=-O_{q-1}(K'),
\qquad K_{2q}-\tau_{2q}=-(K'_q-\tau_q).
\label{hmtfg:clasificacion:eq:5}
\end{equation}
The product determining the order is identical. The identities remain valid when the first distinct symbol is the final zero. Thus \(K'<\tau\).

Induction applied to \(G\) gives \(K'=W_{n-1}0\). The positive odd positions and \eqref{hmtfg:clasificacion:eq:4} reconstruct exactly \(K=W_n0\). The theorem is proved. \(\square\)

The theorem also proves that each classified map has the restrictive doubling returns preceding its superstable return, with the domains of \eqref{hmtfg:clasificacion:eq:2}–\eqref{hmtfg:clasificacion:eq:3}. It does not presuppose monotonicity of a family parameter.

\textbf{Symmetric normalization of the fixed family.} When \(F\) is even and defined on \([-1,1]\), the preceding inequalities additionally give
\[
F(d_4)=d_5>d_3=F(d_2)=F(-d_2).
\]
Since both arguments \(d_4,-d_2\) are positive and \(F\) decreases on that branch,
\begin{equation}
0<d_4<-d_2<1.
\label{hmtfg:clasificacion:eq:5a}
\end{equation}
Therefore \(F^2([d_2,-d_2])=[d_2,d_4]\), and the same conjugacy \(h(x)=d_2x\) defines on all of \([-1,1]\) an even map with a maximum whose image is \([d_4/d_2,1]\subset[-1,1]\). It agrees exactly with the operator \(RF=-F(F(-ax))/a\), \(a=-F(1)\), used in the corpus. This symmetric extension preserves the critical return, its period and its itinerary.


\subsection{First parameter with infinite itinerary}
\label{hmtfg:clasificacion:3}

We work on the compact interval already used in the existence proof,
\[
L=\left[\frac1{2u_0(1)},\frac2{u_0(1)}\right].
\]
Let
\[
\Lambda_\tau=\{\lambda\in L:
\operatorname{sign}c_j(\lambda)=\tau_j\text{ for every }j\ge1\}.
\]
By the existence and sign-separation proofs of the prolongation, this set is nonempty and compact. Therefore there exists
\begin{equation}
a=\min\Lambda_\tau.
\label{hmtfg:clasificacion:eq:6}
\end{equation}

\textbf{Comparison lemma.} For every \(\lambda\in[1/(2u_0(1)),a)\), we have \(K_\lambda<\tau\).

\textbf{Proof.} At the left endpoint, the itinerary is \(+^\infty<\tau\). The separation argument for realization of the infinite word proves that the strict comparison sets with \(\tau\) are open: in an infinite itinerary distinct from \(\tau\), the first finite difference persists by continuity; at a superstable parameter, the two neighboring periodic words and the word with a final zero lie on the same side of \(\tau\), by its strict maximality. By definition, the interval preceding \(a\) contains no parameters with itinerary \(\tau\). Its connectedness maintains the initial comparison throughout that interval. \(\square\)


\subsection{Existence of each first new root}
\label{hmtfg:clasificacion:4}

We retain the original selector. The first root is
\[
\lambda_1=\frac1{u_0(1)},
\]
and, for \(n\ge2\), \(\lambda_n\) is the first root of \(c_{2^n}\) to the right of \(\lambda_{n-1}\) satisfying
\begin{equation}
c_{2^j}(\lambda_n)\ne0\quad(1\le j<n).
\label{hmtfg:clasificacion:eq:7}
\end{equation}
Since every first period dividing \(2^n\) is a power of two, \eqref{hmtfg:clasificacion:eq:7} is equivalent to requiring exact critical period \(2^n\).

\textbf{Existence lemma.} All roots of this selector exist and satisfy
\begin{equation}
\lambda_1<\lambda_2<\cdots<\lambda_n<a.
\label{hmtfg:clasificacion:eq:8}
\end{equation}

\textbf{Proof.} The sign \(c_2(a)<0\) implies \(\lambda_1<a\). Suppose \(\lambda_{n-1}<a\) has been constructed, and put \(p=2^{n-1}\). The analytic function \(c_p(\lambda)\) is not identically zero, since \(c_p(a)\ne0\). Its zero set in \([\lambda_{n-1},a]\) is finite, nonempty and does not contain \(a\). Let \(z<a\) be its last zero.

On \((z,a]\), the sign of \(c_p\) is constant and equal to \(s=\tau_p\). For \(g_\lambda=f_\lambda^p\), we have
\[
g_\lambda(0)=c_p(\lambda),\qquad g'_\lambda(0)=0.
\]
A uniform local bound on the second derivative gives
\begin{equation}
c_{2p}(\lambda)
=g_\lambda(c_p(\lambda))
=c_p(\lambda)+O(c_p(\lambda)^2).
\label{hmtfg:clasificacion:eq:9}
\end{equation}
As \(\lambda\) approaches \(z\) from the right, the first-order term dominates; hence \(c_{2p}(\lambda)\) has sign \(s\). At \(a\), its sign is
\[
\tau_{2p}=-\tau_p=-s.
\]
The intermediate value theorem provides a zero \(\beta\in(z,a)\) of \(c_{2p}\). Since \(c_p\ne0\) on that interval, the critical period at \(\beta\) is exactly \(2p\).

Analyticity and \(c_{2p}(a)\ne0\) imply that the zeros of \(c_{2p}\) in \([\lambda_{n-1},a]\) are finite. Among its zeros of exact period \(2p\) to the right of \(\lambda_{n-1}\), there is therefore a first one. It is the original selector \(\lambda_n\), and it satisfies \(\lambda_n\le\beta<a\). \(\square\)

Use of the last auxiliary zero \(z\) belongs solely to the existence proof. The selected parameter remains the \textbf{first new root} \(\lambda_n\).


\subsection{Convergence of the global selector to the first infinite itinerary}
\label{hmtfg:clasificacion:5}

\textbf{Global continuation theorem.} For the fixed HMT family and the original selector, all selected itineraries are canonical and
\begin{equation}
\boxed{K_{\lambda_n}=W_n0,\qquad
\lambda_n\nearrow\min\Lambda_\tau.}
\label{hmtfg:clasificacion:eq:10}
\end{equation}

\textbf{Proof.} The comparison and existence lemmas place each selected parameter below \(a\), with an itinerary less than \(\tau\). The classification theorem identifies that itinerary with \(W_n0\). The sequence is increasing and bounded by \(a\); let \(b\le a\) be its limit.

Fix \(p\ge1\). For \(n\) sufficiently large, \(2p<2^n\). Then the signs of \(c_p(\lambda_n)\) and \(c_{2p}(\lambda_n)\) agree with \(\tau_p\) and \(\tau_{2p}=-\tau_p\), respectively. The uniform bound on second derivatives,
\[
\|(f_\lambda^p)''\|\le B_p,
\qquad B_p=16^p\frac{16^p-1}{15},
\]
and the identity \((f_\lambda^p)'(0)=0\) imply
\[
|c_{2p}(\lambda_n)-c_p(\lambda_n)|
\le\tfrac12B_p|c_p(\lambda_n)|^2.
\]
Since the two terms on the left-hand side have opposite signs,
\begin{equation}
|c_p(\lambda_n)|>2/B_p.
\label{hmtfg:clasificacion:eq:11}
\end{equation}
By continuity, \(c_p(b)\) retains the sign \(\tau_p\) and satisfies \(|c_p(b)|\ge2/B_p>0\). This holds for each \(p\), so that \(b\in\Lambda_\tau\). Minimality of \(a\) gives \(b\ge a\), and therefore \(b=a\). \(\square\)

The complete proof preserves the global selector and produces its itineraries and its limit. The information from the first eight levels is compatible with the result, but the proof of \eqref{hmtfg:clasificacion:eq:10} does not extrapolate those eight levels: it uses induction, invariant intervals and uniform separation of every sign.


\subsection{Logical checks and falsifiers}
\label{hmtfg:clasificacion:7}

\par\noindent\textbf{1.}\enspace The classification theorem requires a strict unimodal maximum and an invariant interval. If either branch loses monotonicity, the intervals \eqref{hmtfg:clasificacion:eq:2} may cease to be restrictive and the classification must be checked again.
\par\noindent\textbf{2.}\enspace The condition \(K<\tau\) is essential. Beyond the limit word, centers of period \(2^n\) with other combinatorics may exist; the theorem does not identify them with \(W_n0\).
\par\noindent\textbf{3.}\enspace Analyticity in the parameter provides isolation and finiteness of the zeros on the compact set. For a merely continuous family, the existence of a strictly subsequent first root requires an additional argument.
\par\noindent\textbf{4.}\enspace A limit of critical roots may, in general, lose signs upon reaching a return of lower period. Inequality \eqref{hmtfg:clasificacion:eq:11} excludes precisely that phenomenon for this sequence and each fixed horizon.
\par\noindent\textbf{5.}\enspace Convergence of the first roots to the first parameter with itinerary \(\tau\) is not equivalent to uniqueness of all parameters with that itinerary and does not by itself determine a metric rate. Transverse crossing and branch identification provide those subsequent conclusions.

% Fragmento integrable. Preambulo anfitrion: amsmath, amssymb, longtable, hyperref.
\section{Transverse entry and local scaling}
\label{hmtfg:fragmento:entrada-transversal-y-escalado-local}

\subsection{Provenance and result}
\label{hmtfg:escalado:1}

APP evaluates addition and multiplication on the two sheets of the lattice, preserving quotient and remainder. TRIT determines regime and orientation. TPK transport selects, transports and updates the state with its carry, boundary and memory records. Its readers act on the same discrete structure of the continuum. The prolongation \texttt{w6→w12→w18→w24→w30→R36→G9} preserves this provenance and distinguishes phase return from memory increment.

The critical reader of the dodecaphase wheel, constructed in the operator antecedents and transported by compatible refinement, produces
\[
\phi_j=\frac{(3j-1)\pi_{\mathrm{HMT}}}{18},\qquad
w_j=\frac1{12},\qquad
k_j=\sqrt{\frac2{\sum_{\ell=1}^{12}w_\ell\phi_\ell^2}}\,\phi_j
=\frac{2(3j-1)}{\sqrt{899}}.
\]
The coordinate \(\pi_{\mathrm{HMT}}\) retains its APP–TRIT–TPK origin; its common factor cancels in this second moment. The resulting reader is
\begin{equation}
u_0(x)=\frac1{12}\sum_{j=1}^{12}(1-\cos k_jx),\qquad
f_\lambda(x)=1-\lambda u_0(x).
\label{hmtfg:escalado:eq:1}
\end{equation}
The focal sector is the uniform one, \(\eta=0\). The preceding extension separately retains the weighted family \(|\eta|\le1/40\). The new numerical bounds correspond to \eqref{hmtfg:escalado:eq:1}.

The family is even and entire, with \(f_\lambda(0)=1\), a quadratic critical point and second moment \(\sum_jw_jk_j^2=2\). The parameter \(\lambda\) ranges over a family downstream of the HMT generator. The target Feigenbaum constant remains outside the inputs of \eqref{hmtfg:escalado:eq:1}.

\textbf{Main result.} There exists a unique
\begin{equation}
\lambda_*\in[1.874038,1.874039]
\label{hmtfg:escalado:eq:2}
\end{equation}
for which the fourth normalized return of \eqref{hmtfg:escalado:eq:1} belongs to the local stable leaf of the real doubling fixed point. The crossing of the family with that leaf is transverse, with quantified separation. For every \(n\ge0\),
\begin{equation}
\boxed{\|R^{4+n}f_{\lambda_*}-g\|_{\mathcal A}
<0.001666\,(0.83996)^n.}
\label{hmtfg:escalado:eq:3}
\end{equation}
The local superstable cascade associated with this crossing has parameters \(\mu_j\) and constants \(C\ne0\), \(0<\theta<1\), such that
\begin{equation}
\mu_j-\lambda_*=C\delta_{\mathrm F}^{-j}\bigl[1+O(\theta^j)\bigr],
\qquad
\frac{\mu_{j-1}-\mu_{j-2}}{\mu_j-\mu_{j-1}}\longrightarrow\delta_{\mathrm F}.
\label{hmtfg:escalado:eq:4}
\end{equation}
Here \(\delta_{\mathrm F}>1\) is the unstable eigenvalue of the doubling operator. Selection of \(\mu_j\) is performed through the local returns of section~\ref{hmtfg:escalado:8}. Identification with the global selector is proved in sections~\ref{hmtfg:escalado:15}--\ref{hmtfg:escalado:17}: all global terms retain the combinatorics and the centers agree at every sufficiently high level when their periods are matched.


\subsection{Analytic chart and effective return}
\label{hmtfg:escalado:2}

The chart of even functions is used
\[
s=x^2,\qquad z=\frac{s-1}{5/2},\qquad
f(x)=1-x^2V(z),\qquad V(z)=\sum_{j\ge0}v_jz^j.
\]
The coordinates of the Banach space are
\begin{equation}
u=10v_0,\qquad \nu=(v_1,v_2,\ldots),\qquad
\|(\Delta u,\Delta\nu)\|_{\mathcal A}
=|\Delta u|+\sum_{j\ge1}|\Delta v_j|.
\label{hmtfg:escalado:eq:5}
\end{equation}
In particular, the constant coefficient of \(V\) has weight ten. This normalization agrees with the Hertling–Spandl chart \(u/10\), section 2 of the source.

Let \(a=-f(1)=v_0-1\). For admissible real returns,
\begin{equation}
(Rf)(x)=-\frac{f(f(-ax))}{a}.
\label{hmtfg:escalado:eq:6}
\end{equation}
The program evaluates an exact factorization of \eqref{hmtfg:escalado:eq:6}. Defining
\[
S=\frac{a^2s-1}{5/2},\qquad b=V(S),\qquad
h=1-a^2sb,\qquad t=\frac{h^2-1}{5/2},
\]
\[
\operatorname{shift}V(z)=\sum_{j\ge0}v_{j+1}z^j,
\]
one obtains
\begin{equation}
\boxed{V_R=a\,b\,(h+1)\left[V(t)+\frac25\operatorname{shift}V(t)\right].}
\label{hmtfg:escalado:eq:7}
\end{equation}
To verify it, one uses \(v_0=1+a\),
\(h^2-1=-a^2sb(h+1)\) and
\(V(t)-v_0=t\operatorname{shift}V(t)\). Thus,
\(a+1-h^2V(t)=a^2sb(h+1)[V(t)+(2/5)\operatorname{shift}V(t)]\),
which is precisely \(asV_R\).

The initial moments are rational:
\begin{equation}
\frac1{12}\sum_jk_j^{2m}
=\frac1{12}\left(\frac4{899}\right)^m
\sum_{j=1}^{12}(3j-1)^{2m}.
\label{hmtfg:escalado:eq:8}
\end{equation}
Therefore both \eqref{hmtfg:escalado:eq:7} and the initial construction use rational intervals; the initial trigonometric functions are represented by their series with an outer remainder bound.


\subsection{External bounds used as subsequent recognition}
\label{hmtfg:escalado:3}

Once \eqref{hmtfg:escalado:eq:1} has been constructed, Lanford’s hyperbolicity theorem is used as an analytic certificate for the operator \eqref{hmtfg:escalado:eq:6}. Let \(g_0\) be the published central polynomial and \(g\) the fixed point. The estimates used are
\[
\|g-g_0\|_{\mathcal A}<\varepsilon_0=4\,10^{-5},
\]
and, in the ball \(\|f-g_0\|_{\mathcal A}<0.01\), the blocks of \(DR\), written \(R=(U,V)\), satisfy
\begin{equation}
U_u\ge a_0=4.521,\quad
\|U_\nu\|\le b_0=0.560,\quad
\|V_u\|\le c_0=0.756,\quad
\|V_\nu\|\le d_0=0.719.
\label{hmtfg:escalado:eq:9}
\end{equation}
The fixed point has a unique unstable direction, with positive real eigenvalue \(\delta_{\mathrm F}>1\); the rest of the spectrum lies inside the unit disk. The figures in \eqref{hmtfg:escalado:eq:9} come from the published estimates, not from a new reproduction of Lanford’s proof in this dossier. The coefficients of \(g_0\) are explicitly incorporated into the program’s central reference polynomial, in accordance with table 2 of Hertling–Spandl.

Sources: \href{https://repo-archives.ihes.fr/A_GARDER/20180409/Test/I_Prepublications/LANFORD/1968-2013/P_81_17/P_81_17_web.pdf}{Lanford, hyperbolicity estimates}; \href{https://arxiv.org/pdf/1410.3277}{Hertling–Spandl, analytic chart and table 2}.


\subsection{Rational entry and direction certificate}
\label{hmtfg:escalado:4}

The renormalization entry certificate divides \eqref{hmtfg:escalado:eq:2} into sixteen equal rational intervals. It propagates four returns of \eqref{hmtfg:escalado:eq:7} and the exact derivative with respect to \(\lambda\), using differentiation of the composition and outward rounding to one hundred decimal places.

Each series consists of a polynomial of degree 32 with interval coefficients and an arbitrary analytic error \(E\), \(\|E\|_1\le\epsilon\). This error may also modify low-order coefficients. For \(\|q\|_1<1\),
\begin{equation}
\|E\circ q\|_1\le\epsilon,\qquad
\|E'\circ q\|_1\le\frac{\epsilon}{(1-\|q\|_1)^2},\qquad
\|\operatorname{shift}E\|_1\le\epsilon.
\label{hmtfg:escalado:eq:10}
\end{equation}
The products include polynomial–error and error–error terms. The initial tail is bounded by its first absolute term and a geometric ratio, using \(\|1+(5/2)z\|_1=7/2\). Each composed argument remains strictly inside the unit ball of this algebra.

For \(\Gamma(\lambda)=R^4f_\lambda=(u_\lambda,\nu_\lambda)\), the certificate gives the following conservative bounds, uniform in \eqref{hmtfg:escalado:eq:2}:

\begingroup\small
\begin{longtable}{p{0.465\linewidth}p{0.465\linewidth}}
\hline
Quantity & Certified outer bound \\
\hline
\(\|\Gamma(\lambda)-g_0\|_{\mathcal A}\) & \(<0.004993128\) \\
\(\|\nu_\lambda-\nu_{g_0}\|_1\) & \(<0.001396077\) \\
\(u'_\lambda\) & \(>3821.289115\) \\
\(\|\nu'_\lambda\|_1\) & \(<551.757000\) \\
\(\|\nu'_\lambda\|_1/u'_\lambda\) & \(<0.144386<1/4\) \\
Analytic error of the function & \(<8.482\,10^{-19}\) \\
Analytic error of the tangent & \(<4.569\,10^{-15}\) \\
\hline\end{longtable}
\endgroup

The complete decimals, the sixteen subintervals and every intermediate return are retained in the rational entry certificate.

The cone \(\|\dot\nu\|_1\le|\dot u|/4\) is invariant under \eqref{hmtfg:escalado:eq:9}:
\[
|\dot u_{\rm nuevo}|\ge(4.521-0.560/4)|\dot u|
=4.381|\dot u|,
\]
\begin{equation}
\|\dot\nu_{\rm nueva}\|_1\le(0.756+0.719/4)|\dot u|
=0.93575|\dot u|<\tfrac14(4.381)|\dot u|.
\label{hmtfg:escalado:eq:11}
\end{equation}
The certified parameter direction therefore enters a uniformly expanding cone of the operator.


\subsection{Quantitative construction of the stable leaf}
\label{hmtfg:escalado:5}

Translate \(g\) to the origin and write the coordinates as \((x,y)=(u-u_g,\nu-\nu_g)\). Fix
\begin{equation}
L=0.16,\qquad r=0.004,\qquad
A=a_0-Lc_0=4.40004,\qquad q=d_0+c_0L=0.83996.
\label{hmtfg:escalado:eq:12}
\end{equation}
The cylinder \(|x|\le Lr\), \(\|y\|_1\le r\) lies inside the ball of \eqref{hmtfg:escalado:eq:9}, since
\((1+L)r+\varepsilon_0=0.00468<0.01\).

Consider the complete space of functions \(\sigma:B_r\to\mathbb R\) with \(\sigma(0)=0\) and Lipschitz constant at most \(L\), equipped with the uniform distance. For each \(y\), define \((\mathcal G\sigma)(y)\) as the root in \([-L\|y\|_1,L\|y\|_1]\) of
\begin{equation}
F_\sigma(x,y)=U(x,y)-\sigma(V(x,y)).
\label{hmtfg:escalado:eq:13}
\end{equation}
On that interval, \(\|V(x,y)\|_1\le q\|y\|_1<r\). As \(x\) increases, \eqref{hmtfg:escalado:eq:9} shows that \(F_\sigma\) grows with slope at least \(A\). At its endpoints,
\[
F_\sigma(L\|y\|_1,y)
\ge[AL-b_0-Ld_0]\|y\|_1,
\]
and the opposite inequality holds at the negative endpoint. The margin is
\begin{equation}
AL-b_0-Ld_0=0.0289664>0.
\label{hmtfg:escalado:eq:14}
\end{equation}
A unique root is obtained. Comparing \eqref{hmtfg:escalado:eq:13} for two values of \(y\),
\[
\operatorname{Lip}(\mathcal G\sigma)
\le\frac{b_0+Ld_0}{A}
=\frac{0.67504}{4.40004}<0.16.
\]
Comparing two graphs,
\begin{equation}
\|\mathcal G\sigma-\mathcal G\widetilde\sigma\|_\infty
\le A^{-1}\|\sigma-\widetilde\sigma\|_\infty,
\qquad A^{-1}<0.228.
\label{hmtfg:escalado:eq:15}
\end{equation}
The contraction theorem produces a unique fixed graph \(x=\sigma_*(y)\). On it,
\begin{equation}
\|y_n\|_1\le q^n\|y_0\|_1,
\qquad
\|(x_n,y_n)\|_{\mathcal A}\le(1+L)q^n\|y_0\|_1.
\label{hmtfg:escalado:eq:16}
\end{equation}
Moreover, the vertical distance \(d(x,y)=x-\sigma_*(y)\) satisfies
\[
|d(R(x,y))|\ge A|d(x,y)|
\]
as long as the orbit remains in the cylinder. Every orbit that remains there indefinitely belongs to the graph. This graph identifies exactly the local stable leaf and provides estimate \eqref{hmtfg:escalado:eq:16}. Its local regularity agrees with the analytic stable manifold of the hyperbolic fixed point.


\subsection{Existence, uniqueness and transversality of the parameter}
\label{hmtfg:escalado:6}

For \(\lambda\) in \eqref{hmtfg:escalado:eq:2}, define
\[
H(\lambda)=u_\lambda-u_g-
\sigma_*(\nu_\lambda-\nu_g).
\]
The certificate ensures \(\|\nu_\lambda-\nu_g\|_1<0.001436077<r\). For \(\lambda_2>\lambda_1\),
\begin{equation}
H(\lambda_2)-H(\lambda_1)
\ge\int_{\lambda_1}^{\lambda_2}
\bigl[u'_t-L\|\nu'_t\|_1\bigr]dt
>3733.007995\,(\lambda_2-\lambda_1).
\label{hmtfg:escalado:eq:17}
\end{equation}
To determine the endpoint signs, \(\|g-g_0\|_{\mathcal A}<\varepsilon_0\) and \(|\sigma_*(y)|\le L\|y\|_1\) are used. Rational point evaluation produces
\begin{equation}
H(1.874038)<-0.0011856110945,
\qquad H(1.874039)>0.0021866053399.
\label{hmtfg:escalado:eq:18}
\end{equation}
By continuity, a root exists; \eqref{hmtfg:escalado:eq:17} proves its uniqueness and the transversality of the crossing. At that root, \eqref{hmtfg:escalado:eq:16} and
\((1+L)(0.001396077+0.00004)<0.001666\)
prove \eqref{hmtfg:escalado:eq:3}.

The bound is uniform in the number of returns: it converts increasing depth into an explicit geometric operator error.


\subsection{Real admissibility of all returns}
\label{hmtfg:escalado:7}

The four initial returns are accompanied, on each rational subinterval, by the inequalities
\begin{equation}
0<a=-f(1)<1,\qquad b=f(a)>a,\qquad f(b)<a.
\label{hmtfg:escalado:eq:19}
\end{equation}
The program stores \(a,b,f(b)\) before each application of \eqref{hmtfg:escalado:eq:7}. Unimodality, evenness and the quadratic critical point are preserved under these normalized returns.

For subsequent iterations, a uniform check in the ball of \eqref{hmtfg:escalado:eq:9} suffices. If \(E=V-V_{g_0}\), the norm \eqref{hmtfg:escalado:eq:5} gives, for \(-0.4\le z\le0\),
\[
|E(z)|\le0.004,\qquad |E'(z)|\le0.01.
\]
The second inequality uses \(\sup_{j\ge1}j(0.4)^{j-1}=1\). Rational evaluation of the central coefficients provides
\begin{equation}
\frac65<V(z)<\frac85,\qquad |V'(z)|<\frac12,
\qquad 0.39<a<0.41.
\label{hmtfg:escalado:eq:20}
\end{equation}
Consequently,
\[
f'(x)=-2x\left[V(z)+\frac25x^2V'(z)\right],
\]
and the bracketed expression exceeds one for \(0\le x\le1\). The map is unimodal and preserves \([-1,1]\). Moreover,
\[
b=f(a)>1-(0.41)^2\frac85=\frac{4569}{6250}>a,
\]
\begin{equation}
f(b)<1-\left(\frac{4569}{6250}\right)^2\frac65
=\frac{35028967}{97656250}<0.39<a.
\label{hmtfg:escalado:eq:21}
\end{equation}
The orbits of the stable leaf satisfy \eqref{hmtfg:escalado:eq:19} at every depth. Thus the analytic iterations of the operator correspond to real doubling returns, with their domains and orientations preserved.


\subsection{Scaling of the local superstable cascade}
\label{hmtfg:escalado:8}

The metric conclusion requires two crossings: that of the family with the stable leaf, proved in section~\ref{hmtfg:escalado:6}, and that of the unstable manifold with the superstable target. The second is proved by Eckmann–Wittwer for
\(\Sigma_2=\{\psi:\psi(1)=0\}\), corresponding to the period-two superstable cycle. Normalization \eqref{hmtfg:escalado:eq:6} agrees with that of the even real operator used there.

The general theorem of Collet–Eckmann–Lanford, section 6 of the cited source, proposition 6.1 and theorems 6.2–6.3, applies to a Banach-space map \(C^2\), a fixed point with a unique unstable eigenvalue \(\delta_{\mathrm F}>1\), a curve transverse to its stable leaf and a target transverse to its unstable manifold. The local preimages of the target intersect the curve at points satisfying
\begin{equation}
\delta_{\mathrm F}^j(\mu_j-\lambda_*)\longrightarrow C\ne0.
\label{hmtfg:escalado:eq:22}
\end{equation}
The index \(j\) counts local returns. Since the initial curve here is \(R^4f_\lambda\), the physical period is \(2^{j+5}\) when the target is taken in \(\Sigma_2\); any fixed number of steps used to bring the target into the local neighborhood is incorporated into a fixed index shift and into \(C\).

The real conditions of section~\ref{hmtfg:escalado:7} and the target’s real doubling branch fix the exact period and itinerary. The sequence is selected by this oriented local continuation. Taking differences in \eqref{hmtfg:escalado:eq:22} gives the ratio limit of \eqref{hmtfg:escalado:eq:4}.

Sources of the two results used: \href{https://link.springer.com/article/10.1007/BF01013368}{Eckmann–Wittwer, \emph{A complete proof of the Feigenbaum conjectures}}; \href{https://people.math.harvard.edu/~knill/history/lanford/papers/ColletEckmannLanford.pdf}{Collet–Eckmann–Lanford, section 6}.

\subsubsection{Power-law asymptotic remainder}

The analyticity of our family allows \eqref{hmtfg:escalado:eq:22} to be made more precise. In the normal-coordinate construction of Collet–Eckmann–Lanford, only the stable manifold, the unstable manifold and the target hypersurface are flattened. These three submanifolds admit local \(C^2\) changes of coordinates with bounded derivatives. The smooth cutoff used is applied in the scalar unstable coordinate; the complete foliation is then obtained through the constructed change of coordinates.

In those coordinates, the correction is written as a series of increments \(w_j\) whose norm \(C^1\) satisfies
\[
\|w_j\|_{C^1}\le C_1\rho^j,\qquad 0<\rho<1.
\]
The chain rule for the cut-off map, with uniformly bounded first- and second-order derivatives, provides
\[
\|w_j\|_{C^2}\le C_2B^j
\]
for some \(B\ge1\). The weighted norm of the construction controls the ordinary \(C^1\) norm on the domain considered. By interpolation,
\[
[Dw_j]_\beta\le C_3
\left(\rho^{1-\beta}B^\beta\right)^j.
\]
Choose
\begin{equation}
0<\beta<\min\left\{1,
\frac{-\log\rho}{\log B-\log\rho}\right\}.
\label{hmtfg:escalado:eq:23}
\end{equation}
The series converges in \(C^{1,\beta}\). Let \(z\) be the resulting normal coordinate, normalized so that the preimages of the target satisfy \(z=\delta_{\mathrm F}^{-j}\), absorbing the fixed index shift. Since our curve is analytic and transverse,
\[
z(\Gamma(\lambda))=c(\lambda-\lambda_*)
+O(|\lambda-\lambda_*|^{1+\beta}),\qquad c\ne0.
\]
Local inversion gives
\begin{equation}
\mu_j-\lambda_*=C\delta_{\mathrm F}^{-j}
+O\left(\delta_{\mathrm F}^{-(1+\beta)j}\right).
\label{hmtfg:escalado:eq:24}
\end{equation}
This proves the remainder in \eqref{hmtfg:escalado:eq:4}, with \(\theta=\delta_{\mathrm F}^{-\beta}\). The statement asserts the existence of the constants in the parameter remainder. The explicit operator rate \(0.83996\) of \eqref{hmtfg:escalado:eq:3}, for its part, controls convergence of the renormalized maps.

\subsubsection{Conversion of the remainder into ratio error}

If \(\mu_j-\lambda_*=C\delta_{\mathrm F}^{-j}(1+e_j)\), with \(|e_j|\le M\theta^j\), write
\[
\mu_{j-1}-\mu_j=C(\delta_{\mathrm F}-1)\delta_{\mathrm F}^{-j}(1+r_j),
\quad
r_j=\frac{\delta_{\mathrm F} e_{j-1}-e_j}{\delta_{\mathrm F}-1}.
\]
With \(K=M(\delta_{\mathrm F}+\theta)/(\delta_{\mathrm F}-1)\), we have \(|r_j|\le K\theta^{j-1}\). Therefore,
\begin{equation}
\boxed{
\left|\frac{\mu_{j-2}-\mu_{j-1}}{\mu_{j-1}-\mu_j}-\delta_{\mathrm F}\right|
\le\frac{\delta_{\mathrm F} K(1+\theta)\theta^{j-2}}
{1-K\theta^{j-1}}}
\label{hmtfg:escalado:eq:25}
\end{equation}
when \(K\theta^{j-1}<1\). This formula separates the evaluation precision of each root from the truncation error of the cascade.


\subsection{Certified continuation of the eight initial roots}
\label{hmtfg:escalado:9}

The finite continuation certificate preserves the family \eqref{hmtfg:escalado:eq:1} and the previous eight-root certificate. For \(2\le n\le8\), it studies
\[
S_n(\lambda)=f_\lambda^{2^n}(0)
\]
from the certified box of \(\lambda_{n-1}\) to \(1.874039\). The complete rational covering contains exactly the inherited root \(\lambda_{n-1}\) and the new root \(\lambda_n\). In particular, \(\lambda_n\) is the unique new root in \((\lambda_{n-1},1.874039]\).

Boxes are excluded by the sign of the return or by a local derivative of certified sign. Anchor boxes retain existence and uniqueness through endpoint signs and local monotonicity. Their endpoints cover each interval exactly; adjacencies are checked. Evaluation uses the moments \eqref{hmtfg:escalado:eq:8}, order-32 Horner evaluation, outer geometric tails and \(|u_0''|\le2\). The uniform tail of the potential is less than \(1.294\,10^{-58}\) on \(|x|\le3/2\).

The final certificate contains 650 processed boxes and 332 covering leaves. The first level has the exact expression \(\lambda_1=1/u_0(1)\). The following approximations are shown solely for guidance; their complete boxes are retained in the certificates:

\begingroup\small
\begin{longtable}{p{0.465\linewidth}p{0.465\linewidth}}
\hline
Level & Superstable parameter \\
\hline
1 & 1.345913990471832792210949095… \\
2 & 1.763379787846910127290794030… \\
3 & 1.850413796950136464530566778… \\
4 & 1.868979756606798125150574309… \\
5 & 1.872955034435659740004205128… \\
6 & 1.873806367467030119412262311… \\
7 & 1.873988695221365362307168780… \\
8 & 1.874027744154450672897007573… \\
\hline\end{longtable}
\endgroup

% Fragmento integrable. Preambulo anfitrion: amsmath, amssymb, longtable, hyperref.
\section{Isolation of the limit and scaling of the global selector}
\label{hmtfg:fragmento:aislamiento-del-limite-y-escalado-del-selector-global}

\subsection{Indefinite continuation of the first root and exclusion of the intermediate interval}
\label{hmtfg:escalado:15}

The complete classification proved above preserves the family and the first-new-root selector. Ordered transport preserves the zeros and their minimum; the return word preserves the nonadic chronological reading.

\subsubsection{Classification and existence at every depth}
\label{hmtfg:escalado:15.1}

Write \(W_n=\tau_1\cdots\tau_{2^n-1}\), with \(\tau_j=(-1)^{v_2(j)}\), and \(a_*=\min\Lambda_\tau\). The existence and compactness of \(\Lambda_\tau\) are proved in the prolongation development, \ref{hmtfg:prolongacion:6.1}–\ref{hmtfg:prolongacion:6.2}; its minimum exists by compactness and order.

\textbf{Classification.} Every strictly unimodal map with a maximum, on an invariant interval, whose critical point has exact period \(2^n\) and whose itinerary \(K\) satisfies \(K<\tau\), has itinerary \(W_n0\).

The induction proceeds through the effective return. For period at least eight, comparison with \(\tau\), together with the intervals that would trap orbits with incompatible signs, forces the prefix \((+,-,+,+,+,-)\). This yields
\[
J_1=[d_2,d_4],\quad F(J_1)=[d_3,1],\quad
d_4<d_3,\quad F^2(J_1)=J_1.
\]
The return normalized by \(h(x)=d_2x\) has maximum one, half the period, and itinerary \(K'_j=-K_{2j}<\tau\). Order is preserved because
\[
O_{2q-1}(K)=-O_{q-1}(K'),\qquad
K_{2q}-\tau_{2q}=-(K'_q-\tau_q).
\]
The initial cases are \(+0\) and \(+-+0\); induction reconstructs \(W_n0\). For even maps, also \(0<d_4<-d_2<1\), so the return extends to \([-1,1]\) and coincides with \eqref{hmtfg:escalado:eq:6}.

Before \(a_*\), all itineraries are smaller than \(\tau\): the strict-comparison sets are open by the separation proved in \ref{hmtfg:prolongacion:6.1}; the interval preceding \(a_*\) is connected and begins with \(+^\infty<\tau\).

\textbf{Existence of each first root.} Suppose \(\lambda_{n-1}<a_*\) has been constructed, and let \(p=2^{n-1}\). The zeros of \(c_p\) in \([\lambda_{n-1},a_*]\) form a nonempty finite set: the function is analytic and \(c_p(a_*)\ne0\). Let \(z<a_*\) be its last zero. In \((z,a_*]\), \(c_p\) has sign \(\tau_p\). For \(g_\lambda=f_\lambda^p\),
\[
g_\lambda'(0)=0,\qquad
c_{2p}(\lambda)=c_p(\lambda)+O(c_p(\lambda)^2).
\]
Near \(z\), on the right, \(c_{2p}\) has sign \(\tau_p\); at \(a_*\), it has the opposite sign. There is a zero of \(c_{2p}\) in \((z,a_*)\), and its period is exactly \(2p\), since there \(c_p\ne0\). Finiteness of the zeros provides the first new root in \((\lambda_{n-1},a_*)\). The last auxiliary zero proves existence; the effective choice remains the first root. The base case is \(\lambda_1=1/u_0(1)<a_*\).

\subsubsection{Limit of the original sequence}
\label{hmtfg:escalado:15.2}

\textbf{Theorem.} The original selector exists at every level and satisfies
\begin{equation}
\boxed{K_{\lambda_n}=W_n0,\qquad
\lambda_n\nearrow a_*=\min\Lambda_\tau.}
\label{hmtfg:escalado:eq:50}
\end{equation}
\textbf{Proof.} The preceding existence result gives an increasing sequence bounded by \(a_*\). The classification fixes all its itineraries. Let \(b\le a_*\) be its limit. For each \(p\), sufficiently late terms preserve the opposite signs \(\tau_p,\tau_{2p}\). Taylor’s formula and the uniform bound \(B_p=16^p(16^p-1)/15\) give
\[
|c_{2p}(\lambda_n)-c_p(\lambda_n)|
\le\tfrac12 B_p|c_p(\lambda_n)|^2,\qquad
|c_p(\lambda_n)|>2/B_p.
\]
Passing to the limit preserves \(\operatorname{sign}c_p(b)=\tau_p\), with absolute value at least \(2/B_p>0\). For every \(p\), this implies \(b\in\Lambda_\tau\); minimality gives \(b=a_*\).

The result establishes indefinite existence, preservation of the symbolic cascade, and selection of its first limiting parameter for the original global rule. The HMT completion transports this sequence and its limit while preserving order and reading.

\subsubsection{Certificate excluding the intermediate interval}
\label{hmtfg:escalado:15.3}

The bridge program certifies
\begin{equation}
\Lambda_\tau\cap[\lambda_{8,\rm inf},1.874038]=\varnothing,
\label{hmtfg:escalado:eq:51}
\end{equation}
where
\[
\lambda_{8,\rm inf}
=1.874027744154450672897007573595092408435133414523352839687903472138975
\]
is the rational lower endpoint of the certified interval for \(\lambda_8\).

The cover contains 374 exactly adjacent intervals and no unresolved region: 319 with \(c_{512}>0\), 10 with \(c_{1024}<0\), two with \(c_{2048}>0\), all opposite to the sign of \(\tau\), and 43 near centers of periods 256, 512 or 1024, excluded by their second return.

In the latter case, \(g=f_\lambda^p\), \(d=g(0)\) and a uniform bound \(|g''|\le M\) between zero and \(d\) satisfy \(M|d|<2\). Then
\[
|g(d)-d|\le\tfrac12M|d|^2<|d|\quad(d\ne0).
\]
Both returns have the same sign; for \(d=0\), both are zero. Both situations exclude \(\tau_{2p}=-\tau_p\) with strict signs.

The evaluation uses rational moments, geometric tails of the potential and its second derivative, and \(|u_0'''|<6\). Parameter centering intersects the natural enclosure with
\[
c_j(\lambda_{\rm medio})+
\partial_\lambda c_j(I)\,(I-\lambda_{\rm medio});
\]
the derivative is propagated over the full box. The self-map property of \([-1,1]\) is proved before restricting the enclosures to it.

The complete certificate retains the cover and checks its exact tiling, together with every sign or Taylor inequality.

\subsubsection{Joint localization}
\label{hmtfg:escalado:15.4}

The local crossing \(\lambda_*\) realizes \(\tau\) through its real returns at every depth, so that \(a_*\le\lambda_*\). By \eqref{hmtfg:escalado:eq:50}, \(\lambda_8<a_*\); together with \eqref{hmtfg:escalado:eq:51},
\begin{equation}
\boxed{1.874038<a_*\le\lambda_*<1.874039.}
\label{hmtfg:escalado:eq:52}
\end{equation}
The localization preserves the canonical itinerary and places its first limiting parameter in the same interval as the stable crossing.

\subsection{First-exit isolation and coincidence of the limits}
\label{hmtfg:escalado:16}

This section completes the isolation identified in section~\ref{hmtfg:escalado:15.4}. It preserves the HMT family \eqref{hmtfg:escalado:eq:1}, the minimal selection \eqref{hmtfg:escalado:eq:50}, the chart \eqref{hmtfg:escalado:eq:5} and the stable sheet of sections~\ref{hmtfg:escalado:5}--\ref{hmtfg:escalado:7}. The proof uses the order \(a_*\le\lambda_*\): it suffices to exclude a separation toward the negative side of the expanding cone.

\subsubsection{Upper and lower bounds for transverse transport}
\label{hmtfg:escalado:16.1}

Besides \eqref{hmtfg:escalado:eq:9}, Lemma 2.1 of \href{https://arxiv.org/pdf/1410.3277}{Hertling–Spandl, section 2} provides \(\|D\Phi\|<0.9\) in the same ball of radius \(0.01\), where
\[
\Phi_u(u,\nu)=u-\frac{U(u,\nu)-u}{3.669}.
\]
Therefore,
\begin{equation}
\left|1-\frac{U_u-1}{3.669}\right|<0.9,\qquad
U_u<1+1.9(3.669)=7.9711.
\label{hmtfg:escalado:eq:53}
\end{equation}
The number \(3.669\) belongs to the published preconditioner of this analytic certificate; the generated family \eqref{hmtfg:escalado:eq:1} retains its prior coefficients.

Take \(\kappa=0.21\). For two points in the ball joined by a secant with
\(\Delta u<0\), \(\|\Delta\nu\|_1\le\kappa|\Delta u|\), integration of the blocks of \(DR\) along the segment gives
\begin{equation}
4.4034|\Delta u|
\le|\Delta u'|\le8.0887|\Delta u|,\qquad \Delta u'<0,
\label{hmtfg:escalado:eq:54}
\end{equation}
\begin{equation}
\|\Delta\nu'\|_1
\le(0.756+0.719\kappa)|\Delta u|
=0.90699|\Delta u|
<\kappa(4.4034)|\Delta u|.
\label{hmtfg:escalado:eq:55}
\end{equation}
Here \(4.4034=4.521-0.560\kappa\) and \(8.0887=7.9711+0.560\kappa\). Thus, the secant cone preserves orientation and expands while the segment remains in the ball.

\subsubsection{Reduction of an infinite separation to a finite region}
\label{hmtfg:escalado:16.2}

Suppose \(a_*<\lambda_*\), and set
\[
f_n=R^{4+n}f_{a_*},\qquad
s_n=R^{4+n}f_{\lambda_*},\qquad
(\Delta u_n,\Delta\nu_n)=f_n-s_n.
\]
The positive bound for \(u'\) and \(\|\nu'\|_1/u'<0.144386<\kappa\) on \(\Gamma(J)\) place the initial secant in the negative cone. The stable sheet is written relative to \(g\), with
\begin{equation}
s_n=g+e_n,\quad
\|e_n\|_{\mathcal A}<0.001666(0.83996)^n,\quad
|(e_n)_u|\le0.16\|(e_n)_\nu\|_1.
\label{hmtfg:escalado:eq:56}
\end{equation}
The entry checks imply
\begin{equation}
|\Delta u_0|
<0.004993128+0.00004
 +0.16(0.001396077+0.00004)
=0.00526290032< h,\qquad h=0.006.
\label{hmtfg:escalado:eq:57}
\end{equation}
While \(|\Delta u_n|<h\), \eqref{hmtfg:escalado:eq:55}–\eqref{hmtfg:escalado:eq:56} give
\begin{equation}
\|f_n-g_0\|_{\mathcal A}
<(1+\kappa)h+0.001666+0.00004
=0.008966<0.01.
\label{hmtfg:escalado:eq:58}
\end{equation}
The point \(s_n\) and the segment between the two also remain in that convex ball. The lower expansion \eqref{hmtfg:escalado:eq:54} forces a finite first exit \(N\), and the upper bound applied to the preceding step provides
\begin{equation}
0.006\le-\Delta u_N<0.0485322,\qquad
\|\Delta\nu_N\|_1\le0.21|\Delta u_N|.
\label{hmtfg:escalado:eq:59}
\end{equation}
The inequality \(8.0887h=0.0485322\) is used; the full jump, not only the first-exit surface, is included.

Consequently, \(f_N\) belongs to the following region:
\begin{equation}
\begin{split}
\mathcal C_-=\{\,g+(d_u,d_\nu)+e:\;&
-0.0485322\le d_u\le-0.006,\\
&\|d_\nu\|_1\le0.21|d_u|,\quad
|e_u|+\|e_\nu\|_1\le0.001666,\\
&|e_u|\le0.16\|e_\nu\|_1\,\}.
\end{split}
\label{hmtfg:escalado:eq:60}
\end{equation}
Since \(a_*\) realizes \(\tau\), all its normalized returns are unimodal self-maps of \([-1,1]\) with the same itinerary. This property follows from return induction, independently of the distance from \(f_N\) to \(g_0\). This is why the following certificate treats \(\mathcal C_-\) intersected with that class of self-maps.

\subsubsection{Certified exclusion of the first-exit region}
\label{hmtfg:escalado:16.3}

The rational first-exit exclusion certificate, reproduced in the computational supplement, covers the interval of \eqref{hmtfg:escalado:eq:59} exactly with 64 adjacent rational bands. All are excluded; no additional subdivision or unresolved region enters the result.

Of the 64 bands, 54 display an opposite sign and ten a contracting return of period 16 or 32. The largest Lipschitz bound for those returns is less than \(0.475146491\); the maximum exclusion horizon is iterate 64.

The evaluation retains the first forty perturbation coefficients as shared variables, together with an explicit bound on their infinite tail. For each band, the center is \(g_0\) shifted in coordinate \(u\). If \(a,b\ge0\) bound the dual sensitivities in \(u,\nu\), the support of the uncertainty in the fixed point and the stable residual is
\begin{equation}
0.00004\max(a,b)
+0.001666\max\left(b,\frac{0.16a+b}{1.16}\right).
\label{hmtfg:escalado:eq:61}
\end{equation}
The second expression is obtained by maximizing \(a|e_u|+b\|e_\nu\|_1\) under the two constraints of \eqref{hmtfg:escalado:eq:56}. This uses the orientation of the stable sheet proved previously, as well as its norm.

Let \(x_j\) be the center’s orbit and \(L_j\) its linear sensitivity to all shared coefficients. The Taylor model has the form
\[
c_j=x_j+L_j(\Delta v)+r_j,\qquad |r_j|\le E_j.
\]
If \(\rho_j\) bounds \(|L_j(\Delta v)|+E_j\), the outward recurrence used is
\begin{equation}
E_{j+1}\le |f_0'(x_j)|E_j
+\tfrac12\sup|f_0''|\rho_j^2
+\sup|\Delta f'|\rho_j.
\label{hmtfg:escalado:eq:62}
\end{equation}
The derivatives are bounded along the segment between the two states. The central orbit is verified within \([-1,1]\); the other orbits belong to that interval by the self-map property included in the certified domain. The calculations use outward rational rounding, including for the geometric tails.

Each band is excluded by one of two inequalities:

\par\noindent\textbf{1.}\enspace A critical value has a sign strictly opposite to \(\tau_j\).
\par\noindent\textbf{2.}\enspace For \(p=2^k\), the return \(H=f^p\) has Lipschitz constant \(L_p<1\) on the segment between \(0\) and \(H(0)\). Then
\[
|H(H(0))-H(0)|\le L_p|H(0)|.
\]
If \(H(0)\ne0\), the two values have the same sign; if \(H(0)=0\), both vanish. In both cases the strict itinerary \(\tau_{2p}=-\tau_p\) is excluded.

The constant \(L_p\) is obtained by propagating the initial segment \([-r,r]\), \(r\ge|H(0)|\), alongside the critical-orbit boxes and multiplying the outward derivative bounds. The maps of the full band are bounded, retaining \eqref{hmtfg:escalado:eq:61}.

\subsubsection{Coincidence theorem for the accumulation parameter}
\label{hmtfg:escalado:16.4}

\textbf{Theorem.} For the uniform HMT family \eqref{hmtfg:escalado:eq:1}, the limit of the original first-root selection coincides with the certified stable crossing:
\begin{equation}
\boxed{\lambda_n\nearrow a_*=\lambda_*,
\qquad 1.874038<\lambda_*<1.874039.}
\label{hmtfg:escalado:eq:63}
\end{equation}

\textbf{Proof.} \eqref{hmtfg:escalado:eq:50}–\eqref{hmtfg:escalado:eq:52} give \(\lambda_n\uparrow a_*\le\lambda_*\), with both parameters in \(J\). If the inequality were strict, \eqref{hmtfg:escalado:eq:54}–\eqref{hmtfg:escalado:eq:59} would produce a return \(f_N\in\mathcal C_-\) realizing \(\tau\). The complete cover of section~\ref{hmtfg:escalado:16.3} excludes precisely that possibility. Hence \(a_*=\lambda_*\).

The proof controls every depth through a finite first exit and a uniform cover of its region. The identity \eqref{hmtfg:escalado:eq:63} links the global selector’s limit to the transverse crossing.


\subsection{Identification of the original centers and global scaling}
\label{hmtfg:escalado:17}

The final step preserves the selection of \eqref{hmtfg:escalado:eq:50}. The local continuation and the original parameters are compared through their exact period and their hybrid class; coincidence is obtained through uniqueness in a fixed parameter neighborhood.

\subsubsection{Degree-two analytic realization and transversality}
\label{hmtfg:escalado:17.1}

The fixed point \(g\) of section~\ref{hmtfg:escalado:3} is even, analytic, with a quadratic critical point and stationary doubling combinatorics. Theorem 2.1 of \href{https://annals.math.princeton.edu/wp-content/uploads/annals-v164-n3-p01.pdf}{de Faria–de Melo–Pinto}, applied to that single combinatorics, gives a one-element limit set, with a proper holomorphic extension of degree two between nested topological disks. Its attraction on the interval identifies that element with \(g\), since \(Rg=g\).

Lemma 3.2 of the same source extends a fixed number of returns to a full analytic neighborhood of \(g\), with degree-two images and positive modulus. The norm \eqref{hmtfg:escalado:eq:5} controls the uniform norm on a sufficiently small complex neighborhood of \([-1,1]\), because \(|(x^2-1)/2.5|<1\) there. The convergence \eqref{hmtfg:escalado:eq:3} enters that neighborhood; analytic dependence of the returns provides, for some fixed integer \(d\), a family
\begin{equation}
\mathcal F_\lambda=R^d f_\lambda
\label{hmtfg:escalado:eq:64}
\end{equation}
of degree-two maps, analytic in a complex neighborhood of \(\lambda_*\). The disks and the number of returns are fixed before varying \(\lambda\) in that neighborhood.

The parameter tangent belongs to the expanding cone, whereas the stable sheet satisfies the slope bound of section~\ref{hmtfg:escalado:5}. The returns preserve that separation. The degree-two realization locally identifies the stable sheet with the hybrid class of the infinitely renormalizable map; hence \eqref{hmtfg:escalado:eq:64} is transverse to that class.

The change of analytic space is justified in detail in \ref{hmtfg:clasificacion:8.2}, devoted to the transport of transversality. A direction tangent to the hybrid class would produce decreasing renormalized tangents, by uniform convergence on a hybrid disk and Cauchy’s estimate. The certified direction, by contrast, satisfies \(|\dot f(1)|=|\dot u|/10\) and grows under the returns by the expanding cone. This common evaluation of both representatives proves transversality without assuming equivalence of their norms.

\subsubsection{Uniqueness in a fixed neighborhood}
\label{hmtfg:escalado:17.2}

For \(n>d\), the classification of section~\ref{hmtfg:escalado:15} gives
\begin{equation}
K(\mathcal F_{\lambda_n})=W_{n-d}0.
\label{hmtfg:escalado:eq:65}
\end{equation}
Straightening a degree-two map preserves the critical orbit; \eqref{hmtfg:escalado:eq:65} identifies its hybrid class with that of the canonical quadratic center of period \(2^{n-d}\). The a priori bounds for the real doubling sequence are those used in Lyubich’s universality theorem.

\href{https://www.maths.tcd.ie/EMIS/journals/Annals/149_2/lyubich.pdf\#page=69}{Lyubich’s Theorem 7.4, pp. 387–388} gives a single intersection with each of those classes, at sufficiently high levels, on the analytic transversal and within a fixed neighborhood of the accumulation parameter.

By \eqref{hmtfg:escalado:eq:63}, \(\lambda_n\) eventually enters that neighborhood. The local continuation of section~\ref{hmtfg:escalado:8}, indexed by the same period, belongs to the same hybrid class. Uniqueness implies
\begin{equation}
\boxed{\lambda_n=\widehat\mu_n\quad(n\ge n_0),}
\label{hmtfg:escalado:eq:66}
\end{equation}
where \(\widehat\mu_n\) denotes the local center of physical period \(2^n\). This is a reindexing of \(\mu_j\) by a fixed shift determined by the returns and the target of section~\ref{hmtfg:escalado:8}. The original selection of \(\lambda_n\) remains intact.

The integer \(n_0\) is existential in this proof. Continuation at every level and the word \(W_n0\) are already proved in \eqref{hmtfg:escalado:eq:50}; \eqref{hmtfg:escalado:eq:66} provides the metric identity needed for the limit.

\subsubsection{Scaling theorem for the first HMT roots}
\label{hmtfg:escalado:17.3}

\textbf{Theorem.} For the uniform family \eqref{hmtfg:escalado:eq:1}, let \(\lambda_n\) be the successive first new roots of exact critical period \(2^n\). There exist \(C>0\), \(0<\theta<1\) and \(M<\infty\) such that, for every sufficiently large \(n\),
\begin{equation}
\lambda_*-\lambda_n=C\delta_{\mathrm F}^{-n}(1+e_n),
\qquad |e_n|\le M\theta^n,
\label{hmtfg:escalado:eq:67}
\end{equation}
and
\begin{equation}
\boxed{\lim_{n\to\infty}
\frac{\lambda_{n-1}-\lambda_{n-2}}
{\lambda_n-\lambda_{n-1}}=\delta_{\mathrm F}.}
\label{hmtfg:escalado:eq:68}
\end{equation}

\textbf{Proof.} \eqref{hmtfg:escalado:eq:66} transports the local scaling \eqref{hmtfg:escalado:eq:24} to the original roots; the finite index shift is absorbed into \(C,M\). Its sign is positive because \(\lambda_n<\lambda_*\).

The power-law remainder also admits a direct justification through the holonomy of hybrid classes. Lyubich’s Lemma 7.3 gives \(C^{1+\beta}\) regularity, with nonzero derivative, on the transversal at the accumulation point. The one-dimensional unstable dynamics is analytic and has multiplier \(\delta_{\mathrm F}>1\); its linearizing coordinate places successive centers at distances proportional to \(\delta_{\mathrm F}^{-n}\). Composing with the inverse holonomy gives a nonzero linear term and a remainder \(O(\delta_{\mathrm F}^{-(1+\beta)n})\). Thus one may take \(\theta=\delta_{\mathrm F}^{-\beta}\).

Writing
\[
r_n=\frac{\delta_{\mathrm F} e_{n-1}-e_n}{\delta_{\mathrm F}-1},
\qquad K=\frac{M(\delta_{\mathrm F}+\theta)}{\delta_{\mathrm F}-1},
\]
one obtains
\[
\lambda_n-\lambda_{n-1}
=C(\delta_{\mathrm F}-1)\delta_{\mathrm F}^{-n}(1+r_n),\qquad
|r_n|\le K\theta^{n-1}.
\]
Therefore,
\begin{equation}
\left|
\frac{\lambda_{n-1}-\lambda_{n-2}}
{\lambda_n-\lambda_{n-1}}-\delta_{\mathrm F}
\right|
\le
\frac{\delta_{\mathrm F} K(1+\theta)\theta^{n-2}}
{1-K\theta^{n-1}},
\quad K\theta^{n-1}<1.
\label{hmtfg:escalado:eq:69}
\end{equation}
The right-hand side tends to zero, proving \eqref{hmtfg:escalado:eq:68}.

\subsubsection{Chain of results and scope of the closure}
\label{hmtfg:escalado:17.4}

The effective composition is: APP–TRIT–TPK dodecaphase reader; profile \eqref{hmtfg:escalado:eq:1}; ordered transport of returns \eqref{hmtfg:escalado:eq:38}–\eqref{hmtfg:escalado:eq:40}; classification and existence of the first roots \eqref{hmtfg:escalado:eq:50}; localization \eqref{hmtfg:escalado:eq:52}; first-exit isolation \eqref{hmtfg:escalado:eq:63}; eventual identity of centers \eqref{hmtfg:escalado:eq:66}; scaling \eqref{hmtfg:escalado:eq:67}–\eqref{hmtfg:escalado:eq:69}. The Jacobi representation \eqref{hmtfg:escalado:eq:43} preserves this reader’s full profile and its marks remain in the antecedent.

The spatial expansion of the HMT reader and nonadic refinement provide the genealogy and evaluation at arbitrary depth. Scaling between parameters is proved through the return dynamics, transversality and isolation that we have just composed. The cited analytic proofs act on the family already constructed; their comparison values do not select its phases or coefficients.

The closure concerns the original global selector of the uniform sector \(\eta=0\), with the listed theorem dependencies. The asymptotic law retains its existential quantifiers. The width of a finite-root box measures evaluation precision; the estimate \eqref{hmtfg:escalado:eq:69} expresses convergence toward the limit.


The programs and certificates in the reproducibility supplement make it possible to reproduce the finite parts of the proof. Continuation at every depth rests on the inductions, the uniform first-exit argument and the identification theorems, with their explicit domains.

% Fragmento integrable. Preambulo anfitrion: amsmath, amssymb, longtable, hyperref.
\section{Analytic identification of the original centers}
\label{hmtfg:fragmento:identificacion-analitica-de-los-centros-originales}

\subsection{Degree-two realization and metric identification of the centers}
\label{hmtfg:clasificacion:8}

The equality between the first parameter of the infinite itinerary and the stable crossing, proved by first exit, provides
\begin{equation}
\boxed{a_*:=\min\Lambda_\tau=\lambda_*.}
\label{hmtfg:clasificacion:eq:H}
\end{equation}

The family, dodecaphase profile, returns and selector remain exactly as in sections~\ref{hmtfg:clasificacion:1}--\ref{hmtfg:clasificacion:5}. The external results used below recognize the dynamics of that already constructed object. Their universal values do not enter the selection of coefficients or parameters of the HMT family.

\textbf{Metric transfer theorem.} The identification \eqref{hmtfg:clasificacion:eq:H}, together with the previously proved entry, admissibility and transversality bounds, implies that the original sequence of first new roots satisfies, for certain constants \(C>0\) and \(0<\theta<1\),
\begin{equation}
\boxed{\lambda_*-\lambda_n=C\delta_{\mathrm F}^{-n}
\bigl(1+O(\theta^n)\bigr),}
\label{hmtfg:clasificacion:eq:14}
\end{equation}
where \(\delta_{\mathrm F}>1\) is the unstable eigenvalue of the doubling operator at the recognized fixed point. In particular,
\begin{equation}
\frac{\lambda_{n-1}-\lambda_{n-2}}
     {\lambda_n-\lambda_{n-1}}
=\delta_{\mathrm F}+O(\theta^n).
\label{hmtfg:clasificacion:eq:15}
\end{equation}
The sequence eventually coincides with the centers of the local continuation, once both indices refer to the same physical period. The proof retains the first-root selector of section~\ref{hmtfg:clasificacion:4}.

\subsubsection{Obtaining a common quadratic-like domain}
\label{hmtfg:clasificacion:8.1}

Write \(g\) for the renormalization fixed point and
\(\Gamma(\lambda)=R^4f_\lambda\). The analytic chart used in the certificate is
\[
f(x)=1-x^2V\!\left(\frac{x^2-1}{2.5}\right),\qquad
V(z)=\frac{u}{10}+\sum_{k\ge1}\nu_kz^k,
\]
with difference norm \(|\Delta u|+\|\Delta\nu\|_1\). The admissibility inequalities prove that \(g\) is even, analytic, unimodal, quadratic and of doubling type; moreover, \(g(x)=\varphi(x^2)\) with \(\varphi'<0\) in \([0,1]\).

Theorem 2.1 of de Faria–de Melo–Pinto, applied to the combinatorial set consisting solely of doubling, provides a one-point limit set with a quadratic-like extension. Since \(Rg=g\), the real convergence in that theorem identifies this point with \(g\). Lemma 3.2 provides, in an analytic neighborhood of \(g\), a finite renormalization iterate with a quadratic-like domain and uniformly positive modulus. The locations are sections 2.1–2.3, pp. 736–738, and section 3.1, Lemma 3.2, p. 745, of \href{https://annals.math.princeton.edu/wp-content/uploads/annals-v164-n3-p01.pdf}{de Faria–de Melo–Pinto, \emph{Global hyperbolicity of renormalization for \(C^r\) unimodal mappings}}.

To check membership in the neighborhood of that lemma, let \(\Omega_\epsilon\) be a sufficiently small complex neighborhood of \([-1,1]\), with
\[
\rho_\epsilon:=\sup_{x\in\Omega_\epsilon}
\left|\frac{x^2-1}{2.5}\right|<1,
\qquad M_\epsilon:=\sup_{x\in\Omega_\epsilon}|x|^2<\infty.
\]
The restriction of the certified chart satisfies
\begin{equation}
\|\Delta f\|_{\Omega_\epsilon}
\le M_\epsilon\left(\frac{|\Delta u|}{10}
          +\rho_\epsilon\|\Delta\nu\|_1\right)
\le C_\epsilon\|\Delta f\|_{\mathcal A}.
\label{hmtfg:clasificacion:eq:16}
\end{equation}
The certificate gives
\[
\|R^m\Gamma(\lambda_*)-g\|_{\mathcal A}
<0.001666(0.83996)^m.
\]
A fixed integer \(m\) is therefore chosen to place this image in the neighborhood of Lemma 3.2. Continuity of the finite composition provides an interval \(I_*\) around \(\lambda_*\) where the same operations are defined. If \(N\) is the iterate supplied by that lemma, we set
\begin{equation}
d=4+m+N,\qquad
\mathcal F_\lambda=R^df_\lambda,\quad \lambda\in I_*.
\label{hmtfg:clasificacion:eq:17}
\end{equation}
The family \(\mathcal F\) has quadratic-like representatives of degree two, on a common domain and with uniformly positive modulus. Its real-analytic dependence admits a local complexification of the parameter. The number \(d\) is fixed and counts the doublings performed before applying the parameter theorem.

This construction fixes a common domain and a finite number of returns \(d\), sufficient to apply the parameter theorem.

\subsubsection{Transport of certified transversality}
\label{hmtfg:clasificacion:8.2}

The direction of \(\Gamma\) satisfies the cone condition
\[
\|\dot\nu\|_1\le\tfrac14|\dot u|,\qquad \dot u\ne0.
\]
The certified differential bounds imply invariance of this cone and
\begin{equation}
|\dot u_{k+1}|\ge4.381|\dot u_k|
\label{hmtfg:clasificacion:eq:18}
\end{equation}
along the orbit of \(\Gamma(\lambda_*)\). All its points remain in the certified domain. In particular, the direction transported through the finite steps of \eqref{hmtfg:clasificacion:eq:17} retains a nonzero expanding component.

The hybrid class of the fixed point coincides with its stable manifold in the space of quadratic-like germs: \href{https://www.maths.tcd.ie/EMIS/journals/Annals/149_2/lyubich.pdf}{Lyubich, \emph{Feigenbaum–Coullet–Tresser universality and Milnor's hairiness conjecture}, Theorem 6.1, pp. 371–372}. Transversality remains valid on passing to that space, as can be checked without identifying different Banach norms. Indeed, if \(\partial_\lambda\mathcal F_{\lambda_*}\) were tangent to the hybrid class, it would be the tangent of an analytic disk contained in it. Uniform exponential convergence of renormalization on a compact subdisk, followed by Cauchy’s estimate in its parameter, would make the evaluation at \(x=1\) of its renormalized tangents converge to zero. However, in the preceding chart,
\[
\dot f(1)=-\dot u/10,
\]
and \eqref{hmtfg:clasificacion:eq:18} makes its absolute value grow. Both derivatives represent the same germ and have the same real evaluation. The contradiction proves
\begin{equation}
\partial_\lambda\mathcal F_{\lambda_*}
\notin T_{\mathcal F_{\lambda_*}}\mathcal H_{c_\infty}.
\label{hmtfg:clasificacion:eq:19}
\end{equation}
Thus, \(S=\{\mathcal F_\lambda\}\) is a complex-analytic transversal to the hybrid class of the doubling limit.

\subsubsection{Eventual identification of the first roots by their period}
\label{hmtfg:clasificacion:8.3}

Denote by \(c_r\) the canonical quadratic center of period \(2^r\), obtained through \(r\) successive doublings of the period-one center. This symbol denotes the parameter of the quadratic recognition family; it is distinct from the orbital functions \(c_j(\lambda)=f_\lambda^j(0)\) of the preceding sections. We write \(c_r^{\rm quad}\) to preserve the distinction.

The global continuation theorem and \eqref{hmtfg:clasificacion:eq:H} give \(\lambda_n\to\lambda_*\), hence \(\lambda_n\in I_*\) for sufficiently large \(n\). The classification and restrictive intervals of section~\ref{hmtfg:clasificacion:2} show that, when \(n>d\),
\begin{equation}
\mathcal F_{\lambda_n}=R^df_{\lambda_n}
\quad\text{has exact critical period }2^{n-d}
\quad\text{and itinerary }W_{n-d}0.
\label{hmtfg:clasificacion:eq:20}
\end{equation}
Its critical orbit remains in the invariant real interval, contained in the quadratic-like domain; its filled Julia set is therefore connected. Straightening preserves the postcritical orbit and the real branches. Following its restrictive returns gives \(n-d\) canonical doublings and, finally, a fixed critical point. The latter straightens to the center \(0\); the successive inverses of the straightening maps of the doubling copies give the unique center \(c_{n-d}^{\rm quad}\). This is the concrete application of the combinatorial parameterization and the straightening homeomorphism of the copies described in Lyubich, sections 5.1–5.3, pp. 362–364. We conclude
\begin{equation}
\mathcal F_{\lambda_n}\in\mathcal H_{c_{n-d}^{\rm quad}}.
\label{hmtfg:clasificacion:eq:21}
\end{equation}

Lyubich’s Theorem 7.4, pp. 387–388, gives a unique local intersection of \(S\) with each class \(\mathcal H_{c_r^{\rm quad}}\) for sufficiently large \(r\). Its assumptions of a priori complex bounds hold for the real quadratic doubling cascade, according to Theorem 5.6, p. 367, of the same work. Write \(\zeta_r\) for the parameter of that intersection. Equations \eqref{hmtfg:clasificacion:eq:20}–\eqref{hmtfg:clasificacion:eq:21}, convergence to the crossing and local uniqueness imply
\begin{equation}
\boxed{\lambda_n=\zeta_{n-d}\quad\text{for every }n\text{ sufficiently large}.}
\label{hmtfg:clasificacion:eq:22}
\end{equation}
The uniqueness used here concerns \textbf{each canonical hybrid class} in a fixed neighborhood of the crossing. It is not inferred merely from the existence of a local cascade or from uniqueness of the crossing with the stable sheet. Equation \eqref{hmtfg:clasificacion:eq:22} identifies the original selector, already defined globally, with the local intersections; it does not introduce a replacement selector.

In particular, let \(\mu_j\) be the local cascade of section \ref{hmtfg:escalado:8}, and let \(s\) be the fixed integer determined by its exact physical period:
\[
\operatorname{per}_{\rm crit}(f_{\mu_j})=2^{j+s}.
\]
In the presentation with curve \(R^4f_\lambda\) and period-two target, \(s=5\); the additional fixed steps used to localize the target modify \(s\). By \eqref{hmtfg:clasificacion:eq:21} and the same uniqueness,
\begin{equation}
\mu_j=\zeta_{j+s-d},\qquad
\boxed{\lambda_n=\mu_{n-s}\quad(n\text{ sufficiently large}).}
\label{hmtfg:clasificacion:eq:23}
\end{equation}
The shift \(d\) in the quadratic-like preparation and the shift \(s\) in the target periods are distinct objects. Identification is made through the physical period, not by equating those two integers.

\subsubsection{Power-law remainder through transverse holonomy}
\label{hmtfg:clasificacion:8.4}

Lyubich’s Lemma 7.3, p. 384, with proof on pp. 384–387, provides \(C^{1+\beta}\)-conformal regularity, for some \(\beta>0\), of the holonomy between transversals to the hybrid class of the Feigenbaum point. Its definition on p. 384 concerns the connectedness loci on the transversals; these contain, in particular, the centers considered here. We may reduce the exponent so that \(0<\beta\le1\).

Let \(h\) be the holonomy from \(S\) to the unstable manifold, and let \(z\) be an analytic coordinate linearizing the one-dimensional restriction of renormalization:
\begin{equation}
z(g)=0,\qquad z(Rq)=\delta_{\mathrm F} z(q).
\label{hmtfg:clasificacion:eq:24}
\end{equation}
The existence of this coordinate follows by applying the linearization of a repelling holomorphic germ, with multiplier \(\delta_{\mathrm F}>1\), to the analytic unstable manifold. The expanding direction of \eqref{hmtfg:clasificacion:eq:18}, cone invariance and the stationary return identify that multiplier with the positive unstable eigenvalue used in the certificate.

If \(q_r\) is the center of the class \(\mathcal H_{c_r^{\rm quad}}\) lying on the unstable manifold, compatibility of straightening and renormalization gives \(Rq_r=q_{r-1}\). This identity is the one used in the proof of Theorem 7.4, p. 388. By \eqref{hmtfg:clasificacion:eq:24}, there exists \(b\ne0\), absorbing any fixed initial index, such that
\begin{equation}
z(q_r)=b\delta_{\mathrm F}^{-r}
\label{hmtfg:clasificacion:eq:25}
\end{equation}
for every sufficiently large \(r\).

The regularity of \(h\) and transversality \eqref{hmtfg:clasificacion:eq:19} give, on the connectedness locus near \(\lambda_*\),
\begin{equation}
H(\lambda):=z\bigl(h(\mathcal F_\lambda)\bigr)
=A(\lambda-\lambda_*)
+O\!\left(|\lambda-\lambda_*|^{1+\beta}\right),
\qquad A\ne0.
\label{hmtfg:clasificacion:eq:26}
\end{equation}
Since \(H(\zeta_r)=b\delta_{\mathrm F}^{-r}\), the lower estimate
\(|H(\lambda)|\ge |A|\,|\lambda-\lambda_*|/2\) in a sufficiently small neighborhood first implies
\(|\zeta_r-\lambda_*|=O(\delta_{\mathrm F}^{-r})\). Substituting this bound into the remainder of \eqref{hmtfg:clasificacion:eq:26} yields
\begin{equation}
\zeta_r-\lambda_*=\frac bA\delta_{\mathrm F}^{-r}
+O\!\left(\delta_{\mathrm F}^{-(1+\beta)r}\right).
\label{hmtfg:clasificacion:eq:27}
\end{equation}
We apply \eqref{hmtfg:clasificacion:eq:22} with \(r=n-d\). The finite shift modifies the leading coefficient and the remainder constant. Since \(\lambda_n<\lambda_*\), we obtain \eqref{hmtfg:clasificacion:eq:14}, with
\[
C=-\frac bA\delta_{\mathrm F}^d>0,
\qquad \theta=\delta_{\mathrm F}^{-\beta}\in(0,1).
\]
Finally,
\[
\lambda_n-\lambda_{n-1}
=C(\delta_{\mathrm F}-1)\delta_{\mathrm F}^{-n}\bigl(1+O(\theta^n)\bigr),
\]
which proves \eqref{hmtfg:clasificacion:eq:15}. \(\square\)


\subsubsection{Quantitative scope of the theorem}
\label{hmtfg:clasificacion:8.5}

The identification of limits proved by first exit activates metric transfer for the global first roots of the uniform sector. The operator rate \(0.83996\) controls the renormalized maps on the stable sheet. The parameter rate \(\theta=\delta_{\mathrm F}^{-\beta}\) follows from the regularity of transverse holonomy. Both estimates retain their variable and domain.

The theorem establishes the existence of \(\beta,\theta,C\), the neighborhood \(I_*\), the finite number of returns \(d\) and an index \(n_0\) beyond which the centers coincide. Their individual numerical values are not assigned in this asymptotic statement. This formulation completes the limit and its power-law remainder; numerical evaluation of those witnesses is a distinct quantitative specialization.

\section{Return, memory and universality of the dodecaphase cascade}
\label{sec:hmt-feig-interpretacion}

\subsection{The generating structure and the critical reading}

The Feigenbaum construction in Triadic Modular Holography links an
oriented arithmetic organization to an asymptotic scaling law.
Its starting point is the effective composition of APP, TRIT and TPK. APP evaluates
sum and product on the two sheets of the digit lattice and preserves the
quotients and remainders of those operations. TRIT determines the local regime and
the orientation of its reading. TPK selects and transports the states,
updates their records and prolongs their routes with carry, boundary and memory.
The joint discrete structure of the continuum provides the common residence
of those states and of the readers acting on them.

The prolongation
\[
w_6\longrightarrow w_{12}\longrightarrow w_{18}
\longrightarrow w_{24}\longrightarrow w_{30}
\longrightarrow R_{36}\longrightarrow G_9
\]
organizes this persistence. The initial seed is lifted while preserving its
prefixes; successive lifts change the regime without losing orientation
or memory; the coinductive boundary allows the construction to continue, and the
nine charts bring together its local readings. Nonadic holonomy expresses
phase return accompanied by an increase in memory. Thus, a visible
turn preserves its position within the complete history. From the outset,
the return operation acquires two inseparable components:
reiteration of a rule and preservation of its provenance.

The critical reader fixed on the oriented dodecaphase wheel turns this
structure into twelve phases with a uniform trace. Their representatives are
\[
\phi_m=\frac{(3m-1)\pi_{\mathrm{HMT}}}{18},
\qquad w_m=\frac1{12},\qquad 1\le m\le12.
\]
The choice of oriented representative matters because the second moment
uses its magnitude. The angular quotient records the phase, while the
representative retains the information needed to normalize it. Coordinate
\(\pi_{\mathrm{HMT}}\) comes from the coinductive readers of the same
APP--TRIT--TPK core. Within its correlated image, coordinates
\(\pi,\varphi,e\) and the dodecaphase closure of \(\alpha_{\mathrm{HMT}}\)
retain their constructive order and their own domains. The critical reader
considered here uses the angular coordinate already generated.

Normalization is determined through an integer sum:
\[
\sum_{m=1}^{12}(3m-1)^2=5394=6\cdot899,
\qquad k_m=\frac{2(3m-1)}{\sqrt{899}}.
\]
The common factor \(\pi_{\mathrm{HMT}}\) cancels when the second moment is fixed.
This cancellation preserves the provenance of the phases and produces a
particularly precise representation of their normalized reading. We obtain
\[
u_0(x)=\frac1{12}\sum_{m=1}^{12}\bigl(1-\cos(k_mx)\bigr),
\qquad f_\lambda(x)=1-\lambda u_0(x),
\qquad \sum_mw_mk_m^2=2.
\]
The profile is even and entire; it satisfies \(u_0(0)=u_0'(0)=0\) and
\(u_0''(0)=2\). For \(0<\lambda\le2/u_0(1)\), the family maps
\([-1,1]\) into itself and has a quadratic critical maximum with
\(f_\lambda''(0)=-2\lambda\). Parameter \(\lambda\) ranges over this family
after the reader, orientation, weights and scale have been fixed. The universal
number recognized at the end is thus separated from the selection
of the coefficients.

\subsection{Twelve phases, moments and spectral representation}

The profile admits a second expression bringing its phases together in a finite
operator. The nodes
\[
s_m=k_m^2=\frac{4(3m-1)^2}{899}
\]
determine the positive measure \(\mu_0=12^{-1}\sum_m\delta_{s_m}\). Its
moments \(M_r=\int s^r\,d\mu_0(s)\) define strictly positive Hankel forms
up to rank twelve. This positivity follows from evaluating
a polynomial at twelve distinct nodes: a polynomial of degree less than
twelve retains at least one nonzero evaluation. At degree twelve, the product
of the twelve nodal factors annihilates the norm and determines the reader’s
rank exactly.

Orthogonalization of these moments produces a Jacobi matrix
\(J_{12}\), real symmetric, with positive subdiagonal entries. Its functional calculus
recovers the complete profile:
\[
u_0(x)=e_0^{\mathsf T}
\bigl[I-\cos(x\sqrt{J_{12}})\bigr]e_0.
\]
The cosine sum and the Jacobi operator are two realizations of the same
reader. The moments retain all frequencies and their weights; the
orientation, route and memory marks remain in the state that
generates them. Rank twelve identifies the spectral dimension of this
specific reading. The depth of its prolongations belongs to another dimension
of the construction: the length of compatible histories and the resolution
at which they are evaluated.

The analytic geometry is also made explicit. For the weighted
deformation \(u_\eta\), in the strip \(|\eta|\le1/40\), the weights remain
positive and the second moment remains normalized. In
\(|\operatorname{Re}z|<\pi/3\) there exists an odd, univalent holomorphic function
\(b_\eta\), normalized by \(b_\eta'(0)=1\), such that
\(u_\eta=b_\eta^2\). The wheel therefore produces a conformal deformation
of the quadratic fold. Its dynamical conjugacy preserves that deformation
through \(b_\eta(1-\lambda w^2)\). The parameter-scaling theorem
developed here concerns the uniform member \(\eta=0\); the
structural properties of the weighted strip and the local persistence of
its roots retain their specific quantifiers.

\subsection{Reading depth and recovery of histories}

Nonadic depth refines the reading of a parameter, whereas dynamical
doubling changes the return horizon. Refinement with
carry expresses the first operation through
\[
L_c(x,y)=(Bx-c,By+c),\qquad
C_m=\sum_{j=1}^m c_jB^{m-j},
\]
\[
(x_m,y_m)=(B^mx_0-C_m,B^my_0+C_m).
\]
The outer support grows as \(B^m\), and the normalized record cylinder
has diameter \(B^{-m}\). The integer pair, depth and boundary
mark preserve reconstruction of quotients, remainders and prefixes.
Outer growth and inner resolution thus describe compatible
operations on the same history.

The exact composition
\(L_D^{[b]}L_C^{[a]}=L_{B^bC+D}^{[a+b]}\) allows a history to be grouped into
blocks without altering its data. A doubling return combines two
transitions; the nonadic calendar combines nine. Both groupings are
compared over a common horizon of \(9\,2^N\) events, preserving
the order of the subblocks and their carries.

The contribution of memory can be observed directly in the profile’s
inverse branches. If \(v=(u_0|_{[0,1]})^{-1}\), the two branches are
\[
\beta_\sigma(y)=\sigma v\!\left(\frac{1-y}{\lambda}\right),
\qquad \sigma\in\{-1,+1\}.
\]
The final value, together with the last sheet, reconstructs the preceding value. A
sequence of sheets allows all intermediate steps to be recovered. The branch
word accompanies the APP--TRIT--TPK records of route, boundary, carry and
memory; each record retains its function. This concrete recovery
explains what information is required to invert the critical reading.

The resolution of each return also has an explicit bound. With
\(t=\lambda u_\eta(1)/2\), \(z=(x+1)/2\) and
\[
F_{t,\eta}(z)=1-t\frac{u_\eta(2z-1)}{u_\eta(1)},
\]
one obtains
\[
\left|F_{t,\eta}^{q}(1/2)-F_{s,\eta}^{q}(1/2)\right|
\le S_q|t-s|,\qquad S_q=\sum_{j=0}^{q-1}(6\sqrt2)^j<9^q.
\]
A base-nine cylinder of depth \(m=2^N+r\) therefore allows the
return of horizon \(2^N\) to be read with diameter less than \(9^{-r}\). This
estimate gives quantitative content to arbitrary depth: for every
requested horizon and precision, it determines a sufficient depth.
Evaluation of the analytic functions and rounding have their own
remainders, which propagate together with the parameter uncertainty.

Vacancies enter the capacity balance of the record.
Comparison between blocks of 729 and 1000 possibilities distinguishes accumulated
capacity from chronological length, preserving the steps at which the former
remains constant. This accounting determines the space required to
preserve a reading and its depth. The admissibility of histories
follows, in turn, from the survival rules. Both operations
contribute to the construction and retain distinct mathematical objects.

\subsection{Nonadic chronology and doubling combinatorics}

The critical doubling word has an expression at every depth:
\[
\tau_j=(-1)^{v_2(j)},\qquad
\tau_{2j}=-\tau_j,\qquad \tau_{2j+1}=+1.
\]
Its prefixes satisfy \(W_{n+1}=W_n(-1)^nW_n\). The rule determines how
the two copies of the return are oriented and where their separation appears.
The nonadic reading of chronology preserves the complete integer
\(K=\rho_9(K)+9q_9(K)\). Reducing it modulo \(2^n\), the holonomy
that advances nine events induces translation by nine on that cycle. Since
nine is odd, this translation visits all its positions. The
permutation \(j\mapsto9j\pmod{2^n}\), which conjugates it to unit advance,
preserves the dyadic valuation of each nonzero position. The chronological remainder and quotient
jointly support that compatibility.

Bilateral memory also has an operator realization compatible
with refinement. On cycles of length \(2^n\), the shift operator
\(C_n\) is incorporated into
\[
\mathbb U_n=P_0\otimes I+P_1\otimes C_n,
\]
with the complementary projectors of the nonadic reader. In the limit, its
visible and complementary components satisfy
\[
\|T_{\infty,a}v\|^2+\|D_{\infty,a}v\|^2=\|v\|^2,
\qquad
v=T_{\infty,a}^*T_{\infty,a}v+D_{\infty,a}^*D_{\infty,a}v.
\]
These identities describe preservation and recovery within the declared space
and representation. Contraction of the renormalized maps,
which will appear later, measures another operation: the approach of their profiles
to a stable form. The history preserved by the lift and the
convergence of a functional reading are articulated at different levels.

\subsection{The ordered chart and global root selection}

The Archimedean readout of the HMT continuum transports the operations
through an isomorphism of complete ordered fields, within the declared
universe \(\mathfrak G\):
\[
\iota:\mathbb R_{\mathrm{HMT}}^{\mathfrak G}
\longrightarrow\mathbb R^{\mathfrak G}.
\]
The internal construction of the cosine through its series, the positive square root of 899,
the profile and its iterates precedes zero selection. Preservation
of sum, product, order and convergent limits gives
\[
\iota\bigl(S_n^{\mathrm H}(a)\bigr)=S_n\bigl(\iota(a)\bigr),
\qquad S_n(\lambda)=f_\lambda^{2^n}(0).
\]
Surjectivity covers all roots in the domain of that chart; order
preserves preceding returns, exact period and the condition of
lying to the right of the preceding root. In this way,
the minimum of each candidate set is also transported. The genealogy reaches
the effective parameter selection, with the second projection and its
marks preserved on the same antecedent.

The selector is fixed by \(\lambda_1=1/u_0(1)\) and by the first
new root of exact critical period \(2^n\) lying to the right of
\(\lambda_{n-1}\). The itinerary classification proves that each selected
root has word \(W_n0\). Restrictive returns halve its
period while preserving unimodal order, and induction restores
the complete word. Analyticity of returns that are not identically zero
isolates their zeros in each compact set;
comparison of signs before the first infinite-itinerary parameter
guarantees a new root at every level.

Let \(\Lambda_\tau\) be the nonempty compact set of parameters realizing
the infinite word, and let \(a_*\) be its minimum. It is proved that
\[
\lambda_n\nearrow a_*.
\]
Passage to the limit preserves each sign through a uniform bound at a fixed
horizon. If \(c_p=f_\lambda^p(0)\), the opposite signs of \(c_p\) and
\(c_{2p}\), together with \((f_\lambda^p)'(0)=0\), imply
\[
|c_p|>\frac2{B_p},\qquad
B_p=16^p\frac{16^p-1}{15}.
\]
Thus, the prefixes of increasing periods preserve the infinite word as
they accumulate. The proof uses the original sequence of first roots from
its global definition.

\subsection{The stable sheet and identification of the selected limit}

The normalized return
\[
(Rf)(x)=-\frac{f(f(-ax))}{a},\qquad a=-f(1)>0,
\]
combines two transitions, changes the scale and preserves the orientation needed
to restore a maximum of value one. The restrictive domains establish
its interpretation as a real return. The derivative of this operator includes
the variation of scale \(a\); that contribution determines the parameter
sensitivity of the complete transformation.

In the chart \(f(x)=1-x^2V((x^2-1)/(5/2))\), with coordinates \(u=10v_0\)
and \(\nu=(v_1,v_2,\ldots)\), the fourth return of the family enters a
quantified neighborhood of the fixed point \(g\). The curve crosses
its stable sheet transversely at a unique local parameter
\[
1.874038<\lambda_*<1.874039,
\]
and satisfies
\[
\|R^{4+n}f_{\lambda_*}-g\|_{\mathcal A}
<0.001666(0.83996)^n.
\]
The stable sheet is a graph of slope at most \(0.16\). The
parameter secants belong to a transverse expanding cone. Both
properties distinguish the decrease in stable deformations from the
growth of separation between parameters.

Identifying \(a_*\) with \(\lambda_*\) brings this local geometry together
with the global minimum. First, the interval between the eighth root
and the lower endpoint of the local box is excluded. Then the order
\(a_*\le\lambda_*\) allows a hypothetical separation to be studied through
the negative cone. While its transverse coordinate has absolute value less
than \(0.006\), the secant remains in the analytic domain and grows by a
factor between \(4.4034\) and \(8.0887\). A persistent separation
would therefore produce a first exit with amplitude between
\(0.006\) and \(0.0485322\).

The stable-orbit residual retains information more precise than its
norm: its components satisfy \(|e_u|\le0.16\|e_\nu\|_1\). This
relation allows their effects on the critical orbit to be bounded jointly.
The cover of the entire first-exit region shows in each case a
sign incompatible with \(\tau\), or a contracting return forcing two
successive dyadic critical values to have the same sign. The infinite
word requires opposite signs. The contradiction identifies the parameters:
\[
\lambda_n\nearrow a_*=\lambda_*.
\]
The proof combines an induction valid at every depth with a
uniform finite exclusion of the region that any separation
would reach. Its infinite scope rests on that composition.

\subsection{Universal scaling and the meaning of precision}

The degree-two holomorphic realization of the returns and the established
transversality allow their hybrid classes to be compared. For each sufficiently
high period, the canonical class has a unique local intersection
with the family. The global first roots already converge to the crossing and
preserve the doubling combinatorics; hence they eventually coincide
with the local centers when their physical periods are matched. The proof
transports the original selection to the universal metric law.

There exist \(C>0\), \(M<\infty\) and \(0<\theta<1\) such that
\[
\lambda_*-\lambda_n=C\delta_{\mathrm F}^{-n}(1+e_n),
\qquad |e_n|\le M\theta^n,
\]
and, consequently,
\[
\lim_{n\to\infty}
\frac{\lambda_{n-1}-\lambda_{n-2}}
{\lambda_n-\lambda_{n-1}}=\delta_{\mathrm F}.
\]
The factor \(\delta_{\mathrm F}\) is the expanding eigenvalue of the doubling operator
recognized at the end of the construction. The specific family preserves its
twelve phases, moments and provenance history; the limit of its
parameter ratios belongs to the same universality class of quadratic
criticality. This relation jointly explains the specificity of the
generator and the universality of the output.

Precision has three residences here. The root intervals
control their finite evaluation. The estimate
\(0.001666(0.83996)^n\) controls operator convergence on the stable
sheet. The term \(M\theta^n\) controls the transfer of that geometry to
the parameter scale through transverse holonomy. Each estimate
corresponds to its variable and operation. The eighth quotient is enclosed
around \(4.669212189150204154556096215889663928\), with width less
than \(5.808\cdot10^{-37}\); that width measures the evaluation of \(\delta_{\mathrm F,8}\).
The distance from \(\delta_{\mathrm F,8}\) to the limit requires quantifying the constants
of the parameter remainder, whose existence is already proved.

The notation keeps separate \(\delta_{\mathrm F}\), the Feigenbaum parameter-scaling
factor; \(\alpha_{\mathrm F}\), the spatial constant associated with its
renormalization problem; and \(\alpha_{\mathrm{HMT}}\), the corpus’s fine-structure
constant. The present result identifies \(\delta_{\mathrm F}\)
through the first roots of the uniform sector \(\eta=0\). The weighted
strip retains its structural and existence theorems, and its quantified
metric extension constitutes a distinct specific statement.

The contribution of the chain lies in continuity between generation,
reading, memory, selection and limit. Oriented arithmetic determines the
profile; the spectral representation preserves its twelve components; the
cylinders allow it to be resolved with arbitrary precision at each horizon;
history transport restores its paths; the ordered chart
preserves all roots and their minima; classification prolongs the
selector; and transverse expansion, together with stable contraction,
determines its asymptotic law. Universality emerges as a proved
property of this complete construction and its analytic realization.

\endgroup

% Fragmento compartido por VII y CRISTAL; incluir despues del cuerpo Feigenbaum.
% Procedencia: resultados recuperados, con demostraciones elementales reunidas.
% Perfil, Jacobi y escalado: 00_antecedentes_operatorios.tex, 03, 06 y 07
% del presente paquete; no se cambia su selector ni su dominio uniforme.
% Constitutiva: CRISTAL_ES/main_autosuficiente.tex:58738--58840,58864--58894.
% Inversion: CRISTAL_ES/sections/cr28_restitucion_constitutiva.tex:18--59.
% Catalan: el mismo propietario, 277--373; identidad integral, unicidad y serie.
% Corte material: CUMPLIMIENTO_EDITORIAL_20260928/fuentes/CRISTAL_ES.
% No contiene una nueva calibracion ni una nueva campana de calculo.
\section{Articulation of the critical, spectral and constitutive readers}
\label{hmtfg:articulacion:seccion}

The APP--TRIT--TPK composition precedes the readers in this section.
APP preserves the additive and multiplicative evaluations in each cell,
with their remainders and quotients; tritic reduction determines the regime,
orientation and carry of the transition. The effective composition
of selection, transport and memory update produces the enriched
state and its prolongations, developed in
\ref{hmtfg:app-trit-transporte} and
\ref{hmtfg:subsubsec:tres-vueltas-ext-ch}.
Nonadic holonomy composes those transitions: phase return
advances memory and preserves history. The readers act on the
same joint discrete structure of the continuum, with its five
consubstantial constructions and with the two readouts of the same
antecedent: ordered evaluation and boundary incidence.

The oriented wheel and angular channels are outputs of this
composition. Their already generated HMT coordinates are used in subsequent
operations, with their marks and domains; they are not metrological numbers
introduced to select an answer. The ordered readout
allows the following formulas to be evaluated, while the provenance state
preserves sheet, route, remainder, quotient, carry, boundary and memory.
The recoveries that are proved recover the object of their declared
domain; the evaluated representation remains accompanied by those
records.

\subsection{Critical separation as an intrinsic asymptotic reading}
\label{hmtfg:articulacion:separacion}

In the uniform sector of the dodecaphase wheel, the oriented
representatives and their weights are
\[
 \phi_m=\frac{(3m-1)\pi_{\mathrm{HMT}}}{18},\qquad
 w_m=\frac1{12},\qquad 1\le m\le12.
\]
The second moment is normalized through
\(\sum_{m=1}^{12}(3m-1)^2=5394=6\cdot899\).
Thus the common angular factor cancels, leaving
\begin{equation}
 k_m=\frac{2(3m-1)}{\sqrt{899}},\qquad
 u_0(z)=\frac1{12}\sum_{m=1}^{12}\bigl(1-\cos(k_mz)\bigr),
 \qquad f_\lambda(z)=1-\lambda u_0(z).
 \label{hmtfg:articulacion:perfil}
\end{equation}
Variable \(\lambda\) ranges over the family after the
reader, its weights and normalization have been fixed. In
\(0<\lambda\le2/u_0(1)\), the dynamical domain is \([-1,1]\).
The first root is \(\lambda_1=1/u_0(1)\); for \(n\ge2\),
\(\lambda_n\) is the smallest root greater than \(\lambda_{n-1}\)
whose critical point has exact period \(2^n\).
The classification, existence of each minimum and metric
identification of these centers have been proved in
\ref{hmtfg:escalado:15}--\ref{hmtfg:escalado:17}.

\begin{proposition}[Structural reading of the scaling constant]
\label{hmtfg:articulacion:delta}
For the uniform reader \eqref{hmtfg:articulacion:perfil} and its
first-root selector, the critical-separation constant is
\begin{equation}
 \delta_{\mathrm{HMT,crit}}
 :=\lim_{n\to\infty}
 \frac{\lambda_n-\lambda_{n-1}}{\lambda_{n+1}-\lambda_n}
 =\delta_{\mathrm F}.
 \label{hmtfg:articulacion:limite}
\end{equation}
It is the asymptotic parameter-separation invariant of the critical
returns of that reader, with its selection order preserved.
\end{proposition}
\begin{proof}
The scaling theorem \eqref{hmtfg:escalado:eq:67} provides
\(\lambda_*-\lambda_n=C\delta_{\mathrm F}^{-n}(1+e_n)\),
with \(C>0\) and \(e_n\to0\). Subtracting two consecutive terms,
\[
 \lambda_n-\lambda_{n-1}
 =C\delta_{\mathrm F}^{-n}
   \bigl[\delta_{\mathrm F}-1+
          \delta_{\mathrm F}e_{n-1}-e_n\bigr].
\]
The quotient of this equality and the one at index \(n+1\) converges to
\(\delta_{\mathrm F}\). The eventual identity of the local centers
with the first roots, proved in
\eqref{hmtfg:escalado:eq:66}, allows the scaling to be applied precisely
to the selected sequence and not to another cascade chosen afterward.
\end{proof}

This formulation preserves more information than the terminal numeral:
it identifies the wheel, its reading, the family and the selector whose limit
is evaluated. Its interpretation is intrinsic to that HMT composition;
recognition of the universal constant subsequently uses the analytic
theorem applied to the already constructed reader.

\subsection{The finite measure and spectral recovery of the profile}
\label{hmtfg:articulacion:jacobi}

The squared frequencies produce the atomic probability measure
\begin{equation}
 s_m=k_m^2=\frac{4(3m-1)^2}{899},\qquad
 \nu_0=\frac1{12}\sum_{m=1}^{12}\delta_{s_m},\qquad
 M_j=\int s^j\,d\nu_0(s).
 \label{hmtfg:articulacion:medida-finita}
\end{equation}
Here \(\delta_{s_m}\) is the Dirac mass at a node, distinct
from the constant \(\delta_{\mathrm F}\). The measure preserves the
twelve frequencies and their weights. For a nonzero polynomial of degree less
than \(r\le12\),
\[
 \int p(s)^2\,d\nu_0(s)
 =\frac1{12}\sum_{m=1}^{12}p(s_m)^2>0,
\]
since a polynomial of degree less than twelve cannot vanish at the
twelve distinct nodes. Thus the Hankel matrices
\((M_{i+j})_{0\le i,j<r}\) are positive definite.
The polynomial \(\prod_m(s-s_m)\) has zero norm; the space
\(L^2(\nu_0)\) has dimension twelve.

\begin{proposition}[Exact representation of the critical reader]
\label{hmtfg:articulacion:representacion}
Let \(\psi_0,\ldots,\psi_{11}\) be the orthonormal polynomials
with positive leading coefficient for \(\nu_0\), with
\(\psi_0=1\). Multiplication by \(s\) has Jacobi
matrix \(J_{12}\), real symmetric with positive subdiagonal entries.
With \(e_0=(1,0,\ldots,0)^{\mathsf T}\),
\begin{equation}
 u_0(z)=\int\bigl(1-\cos(z\sqrt{s})\bigr)\,d\nu_0(s)
 =e_0^{\mathsf T}\bigl[I-\cos(z\sqrt{J_{12}})\bigr]e_0.
 \label{hmtfg:articulacion:perfil-jacobi}
\end{equation}
\end{proposition}
\begin{proof}
The preceding positivity allows orthogonalization through degree eleven.
If \(i<j-1\), self-adjointness of multiplication gives
\(\langle s\psi_j,\psi_i\rangle
=\langle\psi_j,s\psi_i\rangle=0\).
By symmetry, only the diagonal and its two neighbors remain; the positive
ratio of the leading coefficients establishes positivity of the
subdiagonal. Define
\[
 Q_{mj}=\frac{\psi_j(s_m)}{\sqrt{12}},\qquad
 D=\operatorname{diag}(s_1,\ldots,s_{12}).
\]
Orthonormality gives \(Q^{\mathsf T}Q=I\), and evaluating the
recurrence at the nodes gives \(DQ=QJ_{12}\). Since \(Q\) is
square, \(J_{12}=Q^{\mathsf T}DQ\). Its eigenvalues are
positive, so the positive square root and cosine are calculated in
this diagonalization. Finally,
\(Qe_0=(1,\ldots,1)^{\mathsf T}/\sqrt{12}\), which proves
\eqref{hmtfg:articulacion:perfil-jacobi}.
\end{proof}

The realization through \((J_{12},e_0)\) recovers the same profile;
therefore, preserving the family and selector also recovers
its returns and first roots. Spectral positivity
justifies this representation. Parameter transversality
also uses the derivative of the return and its orbital products,
developed in \ref{hmtfg:escalado:14.3} and
\ref{hmtfg:escalado:6}. These are properties of different operations:
the scaling proof preserves both and does not replace
tangent control with positivity of the moments.

\subsection{Constitutive channels and their labeled inversion}
\label{hmtfg:articulacion:constitutiva}

The constitutive reader acts on another typed object on the same
support: two-sheet angular transport. Its already generated
coordinates are \(x=A_{\rm rad}\) and \(y=C^*_{\rm rad}\), in the
contracting chamber \(x>|y|\). The sheet involution
\(R=R^*=R^{-1}\) determines the marked projectors
\(P_\pm=(I\pm R)/2\). The transport representation is
\begin{equation}
 q_+=\exp[-(x+y)],\qquad q_-=\exp[-(x-y)],\qquad
 T=q_+P_++q_-P_-,\qquad 0<q_\pm<1.
 \label{hmtfg:articulacion:canales}
\end{equation}
Exponential evaluation realizes those channels after their
angular construction. The transport incidences \(90\) and \(120\)
fix the response
\[
 \mathcal V=(I-T^{90})(I-T^{120})^{-1}.
\]
The inverse exists because \(1-q_\pm^{120}>0\).
Writing \(s=q^{30}\), with \(30=\gcd(90,120)\), gives
\begin{equation}
 f(s)=\frac{1-s^3}{1-s^4}
     =\frac{1+s+s^2}{1+s+s^2+s^3},\qquad
 r_\pm=f(q_\pm^{30}),\qquad
 \mathcal V=r_+P_++r_-P_-.
 \label{hmtfg:articulacion:respuesta}
\end{equation}
In this subsection \(f\) denotes the constitutive response;
\(f_\lambda\) still denotes the dynamical map of
\eqref{hmtfg:articulacion:perfil}. Their arguments and domains
are distinct.

\begin{proposition}[Recovery of the channels and orientation]
\label{hmtfg:articulacion:inversion}
The constitutive function \(f:(0,1)\to(3/4,1)\) is a decreasing
bijection. Its inverse \(\psi\) assigns to \(r\) the unique root
\(s\in(0,1)\) of
\begin{equation}
 r s^3+(r-1)(1+s+s^2)=0.
 \label{hmtfg:articulacion:cubica}
\end{equation}
The normalized constitutive readings, with their labeled sheets,
\[
 \widehat\varepsilon=r_+^2,\qquad
 \widehat\mu=r_-^2,\qquad
 \widehat Z=\frac{r_-}{r_+},\qquad
 \widehat c=\frac1{r_+r_-},
\]
with \(\widehat\varepsilon,\widehat\mu\in(9/16,1)\), allow recovery of
\begin{equation}
 \begin{aligned}
 r_+&=\sqrt{\widehat\varepsilon},&
 r_-&=\sqrt{\widehat\mu},\\
 s_\pm&=\psi(r_\pm),& q_\pm&=s_\pm^{1/30}.
 \end{aligned}
 \label{hmtfg:articulacion:recuperacion-canales}
\end{equation}
\begin{equation}
 x=-\frac12\log(q_+q_-),\qquad
 y=\frac12\log\frac{q_-}{q_+}.
 \label{hmtfg:articulacion:recuperacion-angular}
\end{equation}
\end{proposition}
\begin{proof}
Differentiating the rational quotient gives
\begin{equation}
 f'(s)=-\frac{s^2(3+2s+s^2)}{(1+s+s^2+s^3)^2}<0
 \qquad(0<s<1).
 \label{hmtfg:articulacion:monotonia}
\end{equation}
The limits \(f(0^+)=1\) and \(f(1^-)=3/4\), together with
continuity, prove bijectivity. Multiplying \(f(s)=r\) by
the positive denominator produces \eqref{hmtfg:articulacion:cubica}.
Positivity of \(r_\pm\) determines the square roots
and that of \(s_\pm\) determines the thirtieth roots.
By \eqref{hmtfg:articulacion:canales},
\(\log q_+=-x-y\) and \(\log q_-=-x+y\);
adding and subtracting recovers \(x,y\).
Conversely, the recovered roots belong to \((0,1)\),
so that \(x\pm y=-\log q_\pm>0\); the resulting pair
remains in the chamber \(x>|y|\).
\end{proof}

In particular,
\(\widehat\mu\widehat\varepsilon=\widehat c^{-2}\) and
\(\widehat\mu/\widehat\varepsilon=\widehat Z^2\).
Exchanging the labels permutes \(q_+\) and \(q_-\),
preserves \(x\) and changes the sign of \(y\). Recovery
requires that orientation: the isolated product \(r_+r_-\)
preserves common propagation, but does not determine the two labeled
responses. The positive roots and cubic inversion recover
the channels within the fixed reader; the complete record of their
generation remains in the enriched state.

\subsection{The Catalan moment and functional recovery}
\label{hmtfg:articulacion:catalan}

The same incidences determine a relation between functions,
before either channel is evaluated. For \(0<s<1\), let
\(\mathcal D=s\,d/ds\), and define the oriented moment
\begin{equation}
 \mathcal C_2(s)
 =\int_0^s\log\frac{s}{t}\,
     \frac{1+t}{1+t+t^2}f(t)\,dt
 =\int_0^s\frac{\log(s/t)}{1+t^2}\,dt.
 \label{hmtfg:articulacion:integral}
\end{equation}
The second equality follows from the factorization
\(1+t+t^2+t^3=(1+t)(1+t^2)\). The continuous positive weight
\(dt/(1+t^2)\) is thus produced by the rational composition
of the constitutive reader, not by an identification with
\(\nu_0\).

\begin{proposition}[Differential inversion, uniqueness and endpoint reading]
\label{hmtfg:articulacion:restitucion-catalan}
The moment \(\mathcal C_2\) is the unique function
\(F\in C^2((0,1))\) satisfying
\begin{equation}
 \mathcal D^2F(s)=\frac{s(1+s)}{1+s+s^2}f(s),\qquad
 \lim_{s\downarrow0}F(s)=0.
 \label{hmtfg:articulacion:problema-diferencial}
\end{equation}
The constitutive response is recovered through
\begin{equation}
 \mathcal D^2\mathcal C_2(s)=\frac{s}{1+s^2},\qquad
 f(s)=\frac{1+s+s^2}{s(1+s)}\,
          \mathcal D^2\mathcal C_2(s).
 \label{hmtfg:articulacion:inversa-diferencial}
\end{equation}
Its continuous extension to the closed interval has the series
\begin{equation}
 \mathcal C_2(s)=\sum_{j=0}^{\infty}
       \frac{(-1)^j s^{2j+1}}{(2j+1)^2},\qquad 0\le s\le1,
 \qquad
 \mathcal C_2(1)=\sum_{j=0}^{\infty}\frac{(-1)^j}{(2j+1)^2}
 =G_{\rm Cat}.
 \label{hmtfg:articulacion:serie-catalan}
\end{equation}
\end{proposition}
\begin{proof}
The integrand of \eqref{hmtfg:articulacion:integral} is nonnegative
and
\[
 0\le\mathcal C_2(s)\le\int_0^s\log(s/t)\,dt=s.
\]
This proves integrability at zero and the required limit.
The identity \(\log(s/t)=\int_t^s du/u\), followed by
interchanging nonnegative integrals, provides
\[
 \mathcal C_2(s)
 =\int_0^s\frac1u\left(\int_0^u\frac{dt}{1+t^2}\right)du.
\]
On every compact interval of the open domain, the
fundamental theorem of calculus applies. Therefore,
\[
 \mathcal D\mathcal C_2(s)=\int_0^s\frac{dt}{1+t^2},
 \qquad
 \mathcal D^2\mathcal C_2(s)=\frac{s}{1+s^2}.
\]
The rational factorization used in
\eqref{hmtfg:articulacion:integral} gives
\eqref{hmtfg:articulacion:problema-diferencial}; solving for \(f\)
gives \eqref{hmtfg:articulacion:inversa-diferencial}, whose
denominators are positive in \((0,1)\).

If two solutions have the required zero limit, their difference
\(H\) satisfies \(\mathcal D^2H=0\). The substitution \(z=\log s\)
transforms this equation into \(d^2H(e^z)/dz^2=0\), whence
\(H(s)=a\log s+b\). Existence of the zero limit forces
\(a=b=0\), proving uniqueness.

For \(s<1\), the geometric expansion of \((1+t^2)^{-1}\)
can be integrated term by term with the logarithmic weight, since
\[
 \int_0^s t^{2j}\log(s/t)\,dt
 =\frac{s^{2j+1}}{(2j+1)^2},\qquad
 \sum_{j\ge0}\frac{s^{2j+1}}{(2j+1)^2}<\infty.
\]
This proves the series in the interior. The bound
\(1/(2j+1)^2\) gives absolute and uniform convergence of that
series on \([0,1]\). In the integral, the integrand extended
by zero for \(t>s\) is dominated by \(\log(1/t)\)
on \((0,1)\). Dominated convergence allows us to reach
\(s=1\) and agrees with the extension of the series. Its endpoint
value is precisely the stated Catalan series.
\end{proof}

Differential inversion uses the family \(\mathcal C_2(s)\),
not only its endpoint value \(G_{\rm Cat}\). In particular,
\(\mathcal D^2\mathcal C_2(s)\to1/2\) when \(s\uparrow1\).
The series differentiated twice with \(\mathcal D\) has
a radial Abel limit there; the series of the moment itself is
absolutely convergent at the boundary. This distinction preserves
the precise conditions of each passage to the limit.

\subsection{Composition of the recoveries and joint scope}
\label{hmtfg:articulacion:conclusion}

The preceding interrelations have explicit operations and
domains. The moments of the atomic measure \(\nu_0\)
produce \(J_{12}\), and its functional calculus recovers the profile
\(u_0\). The profile and first-root selector determine
the sequence whose scaling yields the readout \(\delta_{\mathrm F}\).
On the other hand, the constitutive function \(f\) produces
\(\mathcal C_2\) through the integral moment, and
\eqref{hmtfg:articulacion:inversa-diferencial} recovers \(f\)
from that family. Its labeled evaluations at
\(s_\pm=q_\pm^{30}\) determine the constitutive responses;
cubic inversion recovers both channels and their angular
coordinates. The endpoint reading of the moment gives \(G_{\rm Cat}\).

The measure \(\nu_0\) is positive, atomic and of rank twelve.
The Catalan moment integrates a continuous positive weight and has
an alternating expansion. The signs of that expansion are
coefficients of the geometric series; they are not negative weights of
\(\nu_0\), nor do they turn the positive integral into the same
spectral representation. Each reader retains its support,
orientation, normalization and evaluation operation.

Unity lies in the APP--TRIT--TPK structure that generates and
transports their antecedents; the functional relations are
established through the proved compositions. Equality of
origin does not identify operators with different domains.
In particular, this articulation does not postulate a scalar formula
determining \(\delta_{\mathrm F}\) from Catalan’s constant and
other physical constants. It places critical scaling alongside
effective recoveries between readers on the same support,
preserving the information required by each direction.

\setcounter{secnumdepth}{\HMTFGSavedSecnumdepth}
\endgroup

\endgroup

\chapter{Rogers--Ramanujan identities and nonadic Pell lifting}
\begingroup
\renewcommand{\chapter}[2][]{}%
\chapter{Arithmetic of mode \(30\): Rogers--Ramanujan identities, normalizer \(899\) and nonadic Pell lift}
\label{chap:p3-final:46:modo_30_pell}
\section{Normalizer \(899\), Pell unit, and arithmetic structure of the mode \(30\)}

The normalizer has the factorization arithmetic

\begin{equation}
899=30^2-1=29\cdot31.
\label{mab:rm:eq:899}
\end{equation}

The same mode \(30\) that generates sectors \(90\) and \(120\) simultaneously
determines the twin-prime pair \(29,31\). Moreover,

\begin{equation}
\varepsilon_{30}=30+\sqrt{899},\qquad
\varepsilon_{30}^{-1}=30-\sqrt{899},\qquad
\log\varepsilon_{30}=\arcosh30.
\label{mab:rm:eq:pell30}
\end{equation}

This Pell unit determines the hyperbolic scale associated with the
normalizer and relates the dodecaphase phase, mode \(30\), twin
primes, and hyperbolic geometry. This relation does not identify the
constitutive operator of the vacuum with the doubling renormalizer; such an
identification would require an explicit intertwiner.

The unit acts integrally through
\begin{equation}
M_{30}=\begin{pmatrix}30&899\\1&30\end{pmatrix},
\qquad \det M_{30}=1.
\label{mab:rm:eq:pell-matrix30}
\end{equation}
On \(\Z[\sqrt{899}]\), this matrix represents multiplication by
\(30+\sqrt{899}\). Its eigenvalues are \(\varepsilon_{30}^{\pm1}\); therefore
\(\log\varepsilon_{30}\) is the translation length of the associated hyperbolic
transformation. The same central integer \(30\) reappears in the vacuum
because \(90=3\cdot30\) and \(120=4\cdot30\), but the two operators continue to
have distinct domains. The demonstrated confluence concerns the mode and its
completions; identity of the dynamics is not postulated.

 \section{Lift nonadic of the unit of Pell}

For \(k\ge1\), let the unramified quadratic ring be
\begin{equation}
\mathcal R_k=(\Z/3^k\Z)[r]/(r^2-2),
\label{eq:pell-ring}
\end{equation}
and let

\[
u=3+2r.
\]

The conjugation \(r\mapsto-r\) gives
\[
N_{\mathcal R_k/(\Z/3^k\Z)}(u)
=(3+2r)(3-2r)=9-8=1,
\]
so that \(u\) is a unit of norm one. Moreover,
\[
u^2=17+12r,\qquad
u^4-1=576+408r=3(192+136r).
\]
The second factor is a unit modulo three because it reduces to \(r\ne0\). Hence
the valuation \(3\)-adic satisfies
\(v_3(u^4-1)=1\). The lifting lemma for a unit congruent to one
modulo three yields, for every \(j\ge0\),
\[
v_3\!\left(u^{4\cdot3^j}-1\right)=j+1.
\]
Since \(u\bmod3=2r\) and \((2r)^2=-1\), its order modulo three is four.
Consequently,

\begin{equation}
\operatorname{ord}(u)=4\cdot3^{k-1},
\label{eq:pell-order}
\end{equation}

and the profinite closure of the compatible cyclic groups is
\begin{equation}
\varprojlim_k\langle u\bmod3^k\rangle
\simeq C_4\times\mathbb Z_3.
\label{eq:pell-profinite}
\end{equation}
Its Archimedean realization is

\[
3+2\sqrt2=(1+\sqrt2)^2.
\]

This construction proves the conservation of nonadic depth
by a hyperbolic unit. The units \(3+2\sqrt2\),
\(2+\sqrt3\), and \(30+\sqrt{899}\) form a Pell family with
distinct provenances; any identification among them requires
explicit maps between the corresponding extensions.

 The mode \(30\) is one of the points at which the architecture begins to exhibit a unit determined by the operatorial and arithmetic genealogy.\par

It appears first in the vacuum because\par

\[
90=3\cdot30,
\qquad
120=4\cdot30.
\]

It is the elementary cadence that permits comparison of the sectors \(90\) and \(120\). The vacuum reads that cadence as the completion \(3\to4\).\par

But the same \(30\) reappears in the construction of the criticality germ. There the central identity is\par

\[
899=30^2-1=29\cdot31.
\]

Normalization of the twelve phases produces \(\sqrt{899}\), and the analytic profile can be written in closed form. What is truly striking is that \(\pi\) disappears before the iteration begins.\par

This must be read carefully. The angular geometry of \(\pi\) is absorbed into the organization of the phases. Once the dodecaphase configuration has been normalized, the dynamic germ is governed by the internal arithmetic of \(899\).\par

Therefore, the construction begins by generating internally a dynamic object of its own:\par

\begin{itemize}[leftmargin=*,itemsep=0.35em]
\item is even;
\item is analytic;
\item has a quadratic critical point;
\item belongs to the corresponding unimodal class;
\item has negative Schwarzian;
\item can be subjected to the operator of renormalization.
\end{itemize}

The superstable ratios then begin to approximate the universal constants. The germ was produced first, and universality appears as the terminal behavior of its dynamics.\par

This is a profound epistemological difference. A numerical coincidence would state: “our calculation resembles Feigenbaum.” The HMT construction states: “this is the exact object that enters the universality class, and these are the quotients that emerge upon iterating it.”.\par

The \(30\) reappears a third time in modular memory. If\par

\[
N=N_{90}+N_{120}
\]

measures the total occupancy and\par

\[
D=90N_{90}+120N_{120}
\]

measures the occupation weighted by provenance, the matrix that transforms \((N_{90},N_{120})\) into \((N,D)\) has determinant \(30\).\par

This allows the process to be inverted:\par

\[
N_{90}=4N-\frac D{30},
\qquad
N_{120}=\frac D{30}-3N.
\]

Thus, the same \(30\) performs three functions:\par

\begin{itemize}[leftmargin=*,itemsep=0.35em]
\item in the vacuum, makes the propagations \(90\) and \(120\) commensurable;
\item at criticality, determines the normalization \(899=30^2-1\);
\item in modular memory, is the determinant that allows the origin of each occupation to be recovered.
\end{itemize}

The vacuum operator, the chaos germ, and the modular Hamiltonian remain typologically distinct. The confluence is subtler: all three preserve a common partition of the genealogy.\par

We might say that the mode \(30\) is a gear. In one machine it measures propagation; in another it normalizes criticality; in another it makes the apparent loss of provenance reversible. The gear is the same, but each machine performs a different function.\par

The emergence of the Pell unit\par

\[
30+\sqrt{899}
\]

reinforces this interpretation. The \(30\) simultaneously organizes a finite sum and hyperbolic growth. The same arithmetic that closes the germ opens a scale dynamics.\par

That is perhaps the profound reason why the result proves so beautiful: the mode that makes the finite completion of the vacuum possible is also the mode that prepares an infinite universality of criticality.\par

Finite and infinite are articulated through the coherent repetition of a finite structure that preserves memory; from it the dynamic infinite is born.\par

  \endgroup
\begin{HMTSpiralChapterBody}
% Corte mecánico íntegro: parte_ii.tex:323--419.
% Residencia material del ensamblaje sucesor de 102 capítulos.
% Residencia causal: capítulo 46.

\chapter{Spectral Pell screw and logarithmic suspension}

\section{Screw law}

The corpus contains the exact identity
\begin{equation}
 V^{-1}\mathscr S V=(2+\sqrt3)\ii\,\mathscr S.
 \label{espiral:eq:tornillo-pell}
\end{equation}
Let \(\varepsilon=2+\sqrt3\). The orbit transversal
\[
 z_n=z_0(\varepsilon\ii)^n
\]
satisfies
\[
 r_n=r_0\varepsilon^n,
 \qquad
 \theta_n=\theta_0+\frac{\pi n}{2}.
\]
Eliminating \(n\),
\begin{equation}
 r(\theta)=r_0
 \exp\!\left[
 \frac{2\log(2+\sqrt3)}{\pi}
 (\theta-\theta_0)
 \right].
 \label{espiral:eq:espiral-log-pell}
\end{equation}
The logarithmic spiral is an output of the operator: its equation was not chosen
to fit \(\varepsilon\) afterward.

\begin{proof}[Derivation of \cref{espiral:eq:espiral-log-pell}]
From \(\theta_n-\theta_0=n\pi/2\) it follows that
\(n=2(\theta_n-\theta_0)/\pi\). Substituting into
\(r_n=r_0\varepsilon^n\),
\[
 r_n=r_0\exp\left(
 \frac{2\log\varepsilon}{\pi}(\theta_n-\theta_0)
 \right).
\]
Continuous interpolation is the indicated equation.
\end{proof}

\section{Conjugate contraction involution}

The bilateral extension \(n\in\Z\) produces
\[
 n\to+\infty:\ r_n\to\infty,
 \qquad
 n\to-\infty:\ r_n\to0.
\]
The inversion \(n\mapsto-n\) interchanges expansion and contraction. Since
\(\varepsilon^{-1}=2-\sqrt3\), both branches belong to the same Pell unit.
They are not two independent spirals, but two orientations of the same
spectral screw.

\begin{figure}[htbp]
\centering
\begin{tikzpicture}[scale=.82]
 \draw[->,HMTGray] (-5.1,0)--(5.1,0) node[right] {$x$};
 \draw[->,HMTGray] (0,-3.1)--(0,3.1) node[above] {$y$};
 \draw[HMTBlue,very thick,domain=-7:5.2,samples=240,smooth,variable=\t]
   plot ({.27*exp(.22*\t)*cos(deg(\t))},{.27*exp(.22*\t)*sin(deg(\t))});
 \draw[HMTRed,thick,domain=-7:5.2,samples=240,smooth,variable=\t]
   plot ({.27*exp(.22*\t)*cos(deg(-\t))},{.27*exp(.22*\t)*sin(deg(-\t))});
 \foreach \k in {-4,-3,...,4}{
   \pgfmathsetmacro{\ang}{1.57079632679*\k}
   \fill[HMTNavy]
     ({.27*exp(.22*\ang)*cos(deg(\ang))},
      {.27*exp(.22*\ang)*sin(deg(\ang))}) circle (.8pt);
 }
 \node[HMTBlue,font=\small\sffamily] at (3.5,2.4) {expansion};
 \node[HMTRed,font=\small\sffamily] at (3.5,-2.4) {conjugate orientation};
\end{tikzpicture}
\caption{Transverse projection of the Pell screw. The marks represent quarter turns.}
\label{espiral:fig:tornillo-pell}
\end{figure}

\section{Scale covariance and energy conservation}

The dilation \(\varepsilon\) must not be interpreted as gratuitous growth
of the energetic amplitude of a wave in vacuum. Its coherent realization
acts on the spatial or spectral scale:
\[
 (r,\varphi,\lambda)
 \longmapsto
 (\varepsilon r,\varphi+\pi/2,\varepsilon\lambda).
\]
This distinction will be essential in the electromagnetic realization of
\cref{espiral:chap:realizacion-em}.

\begin{resultbox}[title={Exact compatibility between the Pell spectral screw and the phase--memory suspension}]
The minimal phase--memory lift is a circular helix. The spectral screw
combines a quarter-turn and a Pell dilation; its projection is a logarithmic
spiral. Both arise from the HMT ledger, but they are neither the same
periodicity nor obtained from one another by forgetting types.
\end{resultbox}
 \end{HMTSpiralChapterBody}

\chapter{Arithmetic theory of nonadic vacancies and boundary torsion}
\begingroup
\renewcommand{\chapter}[2][]{}%
\chapter{Arithmetic theory of nonadic vacancies: defect \(10^3-3^6\), radical network \(3\to6\to9\), power--vacancy law, and boundary torsion}
\label{chap:p3-final:47:vacancias_nonadicas}
\section{Classification of the irreducible arithmetic cycles of the nonadic radical lattice}

In the arithmetic realization, a prime corresponds
to an irreducible cycle that cannot be factored into smaller positive
cycles. Nonadic periodicity may select candidates and organize
residue classes; primality remains defined by the exact multiplicative
property

\[
p=ab\Longrightarrow a=1\ \text{or}\ b=1.
\]

Once the coproduct of prime arithmetic cycles has been selected,
\cref{eq:meissel-mertens,eq:twin-prime-constant} are obtained. The helical structure and the digital root
parametrize phase and carry, but do not replace the factorization criterion.

\section{Density, regularized term, and vacancy holonomy}

The basic positive density is

\begin{equation}
\varepsilon_{\rm vac}=\log_{10}\frac{10}{9}.
\label{eq:vacancy-density}
\end{equation}

The sequence is deduced from the completion
\(1000=729+271\) determines
\begin{equation}
\varepsilon_{\rm vac}
=1-\log_{1000}(729)
=\log_{10}\frac{10}{9},
\label{eq:vacancy-density-two-forms}
\end{equation}
and the exact mechanical word
\begin{equation}
z_n=\lfloor(n+1)\varepsilon_{\rm vac}\rfloor
-\lfloor n\varepsilon_{\rm vac}\rfloor,
\qquad n\ge1.
\label{eq:vacancy-mechanical-word}
\end{equation}
Since \(\varepsilon_{\rm vac}\) is irrational, the word is Sturmian:
it has complexity \(p(N)=N+1\), balance one and topological entropy topological zero.
This dynamic differs from the full shift on the triadic
alphabet triadic: the vacancy determines an ordered boundary of zero entropy.

The suma telescopica produces
\begin{equation}
\sum_{n=1}^{N}z_n
=\lfloor(N+1)\varepsilon_{\rm vac}\rfloor,
\qquad
\sum_{n=1}^{N}(z_n-\varepsilon_{\rm vac})
=\varepsilon_{\rm vac}
-\{(N+1)\varepsilon_{\rm vac}\}.
\label{eq:vacancy-discrepancy}
\end{equation}
The partial sums of the discrepancy have modulus less than one. This
is the estimation supporting the subsequent harmonic constants.

The fraction continuous begins
\[
\varepsilon_{\rm vac}
=[0;21,1,5,1,6,2,3,\ldots].
\]
Thus the ones in \(z_n\) are separated by \(21\) or \(22\) steps.
If
\[
\sigma=22-\varepsilon_{\rm vac}^{-1}
=1-T_G(\varepsilon_{\rm vac}),
\]
then \(a_1(\sigma)=6\) and the returns of the short separation have
length \(6\) or \(7\). The pairs \(21/22\) and \(6/7\) are two scales of
the same rotation, not coincidences of integers.

For the exact word in \cref{eq:vacancy-mechanical-word}, the regularized constant is

\begin{equation}
\gamma_{\rm vac}^{+}
=\lim_{N\to\infty}\left(
\sum_{n\le N}\frac{z_n}{n}
-\varepsilon_{\rm vac}\log N
\right).
\label{eq:vacancy-gamma}
\end{equation}

The associated holonomy is

\begin{equation}
\rho_\varepsilon(k)=\exp(2\pi\ii k\varepsilon_{\rm vac}).
\label{eq:vacancy-holonomy}
\end{equation}

The three quantities have distinct types: density, finite term, and phase. The vacancy Stieltjes tower prolongs \cref{eq:vacancy-gamma}:

\begin{equation}
\gamma_{\rm vac,j}^{+}
=\varepsilon_{\rm vac}\gamma_j
+\sum_{n\ge1}\frac{(z_n-\varepsilon_{\rm vac})(\log n)^j}{n}.
\label{eq:vacancy-stieltjes}
\end{equation}

For \(j=0\), one recovers \(\gamma_{\rm vac}^{+}\). Convergence requires the
discrepancy bound for the integral development of vacancies.

Convergence is proved here. For \(j\ge0\), the sequence
\(b_n=(\log n)^j/n\) tends to zero and is decreasing from
\(n>e^j\) onward. The uniform bound \cref{eq:vacancy-discrepancy} and Dirichlet's criterion prove the convergence of
\[
\sum_{n\ge1}(z_n-\varepsilon_{\rm vac})b_n.
\]
For \(j=0\), separating
\(\varepsilon_{\rm vac}\sum_{n\le N}n^{-1}\) gives exactly
\cref{eq:vacancy-gamma}. Abel summation conservatively bounds the
tail after \(N\) by \(2/(N+1)\): one term comes from the boundary and
another from the total variation of \(1/n\). The declared exhaustive execution uses
\(N=10^8\), enumerates the Beatty positions of the ones, checks the
gap sets \(\{21,22\}\) and returns \(\{6,7\}\), and obtains as its
finite part
\[
-0.11207107505876366224207105242526397929\ldots.
\]
Combining the analytic bound with the recorded rounding and boundary
guards, the finite certificate encloses
\begin{equation}
-0.112071094143613634
<\gamma_{\rm vac}^{+}
<-0.112071055516338732.
\label{eq:vacancy-gamma-interval}
\end{equation}
This interval is a terminal value with controlled error. It does not define the
word; the word and its discrepancy precede it.

A second evaluation using multiprecision arithmetic and directed
rounding control provides, to eighteen decimal places, the certified interval
\[
-0.112071086058763562
<\gamma_{\rm vac}^{+}
<-0.112071064058763762,
\]
contained strictly within \cref{eq:vacancy-gamma-interval}. The inclusion of
both intervals verifies the stability of the finite part with respect to
truncation, the boundary guards, and the precision regime. The
result depends solely on the recurrences, the tail bounds, and the
interval arithmetic set out in this chapter.

\section{Operative content of the helicoidal representation}

The helical representation encodes a periodic phase together with a
memory increment. It has operatorial content when there exists a cocycle

\[
c(\gamma\eta)=c(\gamma)+c(\eta)
\]

that transports the carry and whose monodromy is interlinked with the
arithmetic operator. In the absence of that map, the helix constitutes only a
phase parametrization; with the interlinker, it distinguishes phase return
from state return.

\begin{remark}[Negative control]
Cofinality of cylinders does not imply dynamical
conjugacy. A counterexample to commutation between
the TPK shift and \(T_G\) precludes the intertwiner, but does not affect
\cref{eq:levy,eq:gauss-entropy}. In the vacancy sector, the divergence of
\cref{eq:vacancy-stieltjes} under the declared bound refutes the family of
regularized constants.
\end{remark}

\subsection*{Logical Scope and Domain of Validity}
Cofinality, Gauss identities, Levy, entropy, and Lochs: \emph{Exact internal theorem in the transported system}. GKW and dynamic intertwiner: \emph{Open program with falsification criterion}. Density, constant, and vacancy holonomy: \emph{Exact internal theorem with discrepancy control}.
 \section{Governing theorem of the radical lattice and the power--vacancy law}

The appearance of \(369\) and \(693\) is not reducible to a graphic similarity. There are two exact mechanisms that converge:

\begin{enumerate}
\item \textbf{nonadic algebraic mechanism:} \(3,6,9\) are the visual representatives of the nilpotent radical \(\mathfrak m=(3)\subset\mathbb Z/9\mathbb Z\); the words \(369\), \(693\), and \(936\) are the three phases of its multiplicative orbit;
\item \textbf{Archimedean vacancy mechanism:} the difference between the decimal length of \(9^n\) and the boundary \(10^n\) is governed exactly by the same irrational rotation \(\delta_{\rm vac}=\log_{10}(10/9)\) that already organizes the vacancy calendar.
\end{enumerate}

For \(n=9^9\), both mechanisms meet in the identity

\[
9^9-\left\lceil 9^9\log_{10}\frac{10}{9}\right\rceil
=369\,693\,099,
\]

so that

\[
\#\operatorname{cifras}\bigl(9^{9^9}\bigr)=369\,693\,100.
\]

The block \(369\mid693\) is, therefore, an exact output of the logarithmic reader of the nonadic power; it is not the decimal prefix of \(9^{9^9}\). The final unit proceeds from the universal operator that transforms decimal order into decimal length.

This identity opens a broader structural reading: addition, product, power, subtraction, division, root, and logarithm can be organized as a hierarchy of transports and quotients. Each ascent of operation compactifies histories; TPK preserves the coordinates that make their reconstruction possible.


\section{Radical structure of the APP on \(\mathbb Z/9\mathbb Z\)}

\subsection{Nonadic local ring \(\mathbb Z/9\mathbb Z\) and nilpotent radical}

Let

\[
R=\mathbb Z/9\mathbb Z.
\]

Its group of units and its maximal ideal are

\[
R^\times=\{1,2,4,5,7,8\},
\qquad
\mathfrak m=(3)=\{0,3,6\}=\{9,3,6\}_{\rm vis}.
\]

Moreover,

\[
\mathfrak m^2=0.
\]

The preceding equality is an equality of ideals: every product of two elements of \(\mathfrak m\) vanishes modulo nine. It must not be confused with the Cartesian product \(\mathfrak m\times\mathfrak m\), which contains nine APP positions.

\subsection{Exact sequence of the multiplier \(3\)}

Multiplication by three,

\[
M_3:R\longrightarrow R,
\qquad x\longmapsto3x,
\]

has

\[
\ker M_3=\mathfrak m,
\qquad
\operatorname{im}M_3=\mathfrak m.
\]

It therefore induces an isomorphism of spaces over \(\mathbb F_3\),

\[
\overline M_3:R/\mathfrak m\xrightarrow{\sim}\mathfrak m,
\qquad
\overline x\longmapsto3x.
\]

In visual representatives,

\[
\overline0\longmapsto9,
\qquad
\overline1\longmapsto3,
\qquad
\overline2\longmapsto6.
\]

Thus, the band \(3,6,9\) is not a decimal ornament. It is an internal copy of the ternary quotient lodged in the nonadic radical.

\subsection{Cyclic orbit \(369\to693\to936\to369\)}

The row of the multiplication table corresponding to \(3\) is

\[
3,6,9,3,6,9,3,6,9.
\]

The cambio of phase \(j\mapsto j+1\) produces

\[
369\longmapsto693\longmapsto936\longmapsto369.
\]

These three words are the three cyclic readings of the application.

\[
\mathbb F_3\longrightarrow\mathfrak m,
\qquad
\bar j\longmapsto3j.
\]

The row \(6\), by satisfying \(6\equiv-3\pmod9\), reverses the orientation
of the same cycle. The mark \(9\) fulfills two typed functions:

\[
9\bmod9=0,
\qquad
\operatorname{dr}_9(9)=9.
\]

The module publishes the null class; the positive digital root preserves the mark of closure. This difference of readers is precisely what allows \(9\) to be at once residual zero and visible closure.

\subsection{Radical cross of APP and coordinate reduction to \(\mathbb F_3^2\)}

On

\[
X_9=R^2,
\]

the union

\[
\mathscr C_{369}
=(\mathfrak m\times R)\cup(R\times\mathfrak m)
\]

is the cross formed by rows and columns \(3,6,9\). It contains

\[
3\cdot9+3\cdot9-3\cdot3=45
\]

positions. Their central intersection \(\mathfrak m\times\mathfrak m\) contains nine positions. In the multiplicative observable, the sum of the cross is

\[
297=216+81=3(72+27).
\]

The coordinate reduction produces

\[
X_9/(\mathfrak m\times\mathfrak m)\cong\mathbb F_3^2
\]

as a quotient set typed by the classes modulo three. The inverse lift requires preserving the coordinate within each three-element fiber.

\subsection{Kernel \(7227\), spectral decomposition and link to the radical cross}

The central block of the APP is

\[
B_c=
\begin{pmatrix}
7&2\\
2&7
\end{pmatrix}
=7I+2R
=9P_++5P_-.
\]

Its rowwise reading produces \(72\mid27\); its diagonal reading produces \(77\mid22\). One verifies

\[
72+27=77+22=99,
\]

\[
72-27=45=\det B_c,
\]

\[
77-22=55=54+1.
\]

The relation with the radical cross is not nominal:

\[
3\cdot72=216,
\qquad
3\cdot27=81,
\qquad
3\cdot99=297.
\]

The central unitary block and the non-unitary network \(3,6,9\) are thereby linked by an exact local--global identity. The hierarchy of the sum--product difference preserves its types:

\[
4\quad\text{(local spectral jump)},\qquad
2\quad\text{(off-diagonal coupling)},
\]

\[
6\quad\text{(compressed difference modulo three)},\qquad
54\quad\text{(two-dimensional unfolding)}.
\]


\section{Power--vacancy law for the scales \texorpdfstring{\(9^n\)}{9 to the power n} and \texorpdfstring{\(10^n\)}{10 to the power n}}

\subsection{Decimal Length Functions and Accumulated Vacancy}

For \(n\ge1\), let

\[
D_9(n)=\#\operatorname{cifras}(9^n)
=\left\lfloor n\log_{10}9\right\rfloor+1.
\]

The decimal boundary \(10^n\) has \(n+1\) digits. We define the accumulated vacancy of the nonadic power by

\[
V_9(n)=(n+1)-D_9(n).
\]

We introduce the two complementary numbers

\[
\rho_{\rm vac}=\log_{10}9,
\qquad
\delta_{\rm vac}=1-\rho_{\rm vac}
=\log_{10}\frac{10}{9}.
\]

\subsection{Exact identity between decimal length and vacancy rotation}

\textbf{Theorem.} For every integer \(n\ge1\),

\[
\boxed{V_9(n)=\left\lceil n\delta_{\rm vac}\right\rceil},
\]

and, therefore,

\[
\boxed{D_9(n)=n-\left\lceil n\delta_{\rm vac}\right\rceil+1}.
\]

\textbf{Proof.} Since \(\delta_{\rm vac}\) is irrational, \(n\delta_{\rm vac}\notin\mathbb Z\). Then

\[
\begin{aligned}
D_9(n)
&=\lfloor n(1-\delta_{\rm vac})\rfloor+1\\
&=\lfloor n-n\delta_{\rm vac}\rfloor+1\\
&=n-\lceil n\delta_{\rm vac}\rceil+1.
\end{aligned}
\]

Subtracting from \(n+1\) yields the first identity. \(\square\)

\subsection{Mechanical coding of vacancies through nonadic powers}

The increments of the vacancy are

\[
Z_n=V_9(n+1)-V_9(n)
=\lceil(n+1)\delta_{\rm vac}\rceil
-\lceil n\delta_{\rm vac}\rceil
\in\{0,1\}.
\]

This is exactly the mechanical word of slope \(\delta_{\rm vac}\) used by the HMT vacancy calendar. Therefore, that calendar admits an additional, completely exact reading:

\begin{quote}
\(Z_n=1\) if and only if, in passing from \(9^n\) to \(9^{n+1}\), a new decimal range is lost relative to the boundary \(10^n\to10^{n+1}\).
\end{quote}

The power does not add a coincidence to the theory of vacancies: it provides its natural realization as a counter of decimal rank loss.

\subsection{Law \(21/22\) as a case of the power--vacancy theorem}

Since

\[
\delta_{\rm vac}^{-1}
=21.85434532678283256\ldots,
\]

the indices for which \(Z_n=1\) are separated by \(21\) or \(22\). The first twenty-one exponents satisfy

\[
D_9(n)=n,\qquad1\le n\le21,
\]

whereas

\[
D_9(22)=21.
\]

In particular,

\[
9^9=387\,420\,489
\]

has exactly nine digits. This initial preservation of rank is not an indefinite rule: the vacancy calendar determines exactly when it ceases to hold.

\subsection{Exact evaluation at \(n=9^9\) and length of \(9^{9^9}\)}

Let

\[
n_*=9^9=387\,420\,489.
\]

Then

\[
n_*\delta_{\rm vac}
=17\,727\,389.3684296412564569\ldots,
\]

and, therefore,

\[
\boxed{V_9(n_*)=17\,727\,390}.
\]

The conserved coordinate after discounting vacancies is

\[
n_*-V_9(n_*)
=387\,420\,489-17\,727\,390
=369\,693\,099.
\]

Therefore,

\[
\boxed{\operatorname{ord}_{10}\bigl(9^{9^9}\bigr)=369\,693\,099},
\]

\[
\boxed{D_9(9^9)=369\,693\,100}.
\]

Furthermore, the congruences are obtained.

\[
9^9\equiv0\pmod9,\qquad
V_9(9^9)\equiv0\pmod9,
\]

\[
369\,693\,099\equiv0\pmod9,\qquad
369\,693\,100\equiv1\pmod9.
\]

The corresponding digit sums are

\[
s_{10}(9^9)=45,\qquad
s_{10}(V_9(9^9))=36,
\]

\[
s_{10}(369\,693\,099)=54,\qquad
s_{10}(369\,693\,100)=37.
\]

The return \(9\to1\) appears here with precise types:

\begin{enumerate}
\item \(9^9\) and the accumulated vacancy belong to the null class modulo nine;
\item its difference, the decimal order, remains in the null class and has digit sum \(54\);
\item the universal operator \(D=\operatorname{ord}+1\) adds the closure unit and publishes the positive class \(1\).
\end{enumerate}

\subsection{Domain of validity of the power--vacancy law}

It is proved:

\begin{itemize}
\item the theorem potencia–vacancia for every \(n\);
\item the identity with the mechanical vacancy word;
\item the exact evaluation in \(n=9^9\);
\item the decimal order and length of \(9^{9^9}\);
\item the preceding congruences and digit sums.
\end{itemize}

It is not claimed that \(9\) is the only basis with an analogous property in every domain, that each external occurrence of \(369\) arises from this tower, or that the tower replaces the construction of the nine phases or the generation certificate for \(\pi,e,\varphi\).


\section{Operator Hierarchy, Reversibility, and Genealogical Memory}

\subsection{Translation, Dilation, and Exponentiation as Operational Levels}

Elementary operations can be seen as actions:

\[
T_a(x)=x+a,\qquad
M_a(x)=ax,\qquad
P_k(x)=x^k.
\]

Their partial inverses are, respectively, subtraction, division, and extraction of roots. The passage from one operation to the next modifies the structure of the fibers.

\subsection{Stratified reversibility in \(\mathbb Z/9\mathbb Z\)}

\textbf{proposition.} En the anillo nonadico:

\begin{enumerate}
\item every translation \(T_a\) is bijective;
\item \(M_a\) is bijective if and only if \(a\in R^\times\);
\item if \(a\in\mathfrak m\), multiplication has nontrivial fibers and
its image lies within \(\mathfrak m\);
\item for \(x\in\mathfrak m\), \(x^k=0\) for every \(k\ge2\);
\item over \(R^\times\cong C_6\), the map \(P_k\) is an automorphism if and only if \(\gcd(k,6)=1\).
\end{enumerate}

\textbf{Proof.} The first assertion uses the additive group structure. The second is the characterization of the units. The third follows from \(aR\subseteq\mathfrak m\). The fourth follows from \(\mathfrak m^2=0\). The fifth is the condition for multiplication by \(k\) to be an automorphism of the cyclic group of order six. \(\square\)

The result produces an exact hierarchy:

\[
\text{sum: global reversibility},
\]

\[
\text{product: reversibility restricted to units},
\]

\[
\text{power: cyclic dynamics on units and nilpotent collapse in }\mathfrak m.
\]

Roots are functions only when the corresponding power is injective on the chosen domain; on the total domain they are multibranch correspondences or empty relations.

\subsection{Reading of the Topological Phase Kernel on the operator hierarchy}

This stratification suggests a functional definition:

\begin{quote}
TPK does not add an external operation to sum, product, and power. It preserves the sheet, quotient, carry, orientation, route, and boundary that each ordinary evaluation eliminates when identifying several histories with the same output.
\end{quote}

Sum and product are the first two evaluations on the APP. A power is repeated composition of the product; a root selects or reconstructs branches of that composition. The record distinguishes the published value from the operatorial history producing it.

\subsection{Logarithmic descent between consecutive levels of the hierarchy}

In a positive domain,

\[
\log(xy)=\log x+\log y,\qquad
\log(x^k)=k\log x,\qquad
\log(a^x)=x\log a.
\]

The logarithm converts product into sum, power into product, and exponential into dilation. The defect appears when the subsequent reader discretizes:

\[
\mathcal L_{10}(a^n)
=\lfloor n\log_{10}a\rfloor+1.
\]

The integer part introduces the boundary-crossing calendar. For \(a=9\), that calendar is precisely the vacancy word.

\subsection{Affine conjugation of translations and dilations}

Translations and dilations satisfy

\[
M_aT_bM_a^{-1}=T_{ab}
\]

when \(a\) is invertible. This identity expresses that the product reorganizes the units of addition. When \(a\) ceases to be invertible, conjugation breaks because branches disappear. The radical \(3,6,9\) localizes precisely that failure.


\section{Path Algebra: Multiplicative Composition, Additive Action and Superposition}

In a path algebra, addition aggregates alternatives and multiplication concatenates trajectories. For local weights \(w(e)\),

\[
K(x,y)
=\sum_{\gamma:x\to y}
\prod_{e\in\gamma}w(e).
\]

If a coherent logarithmic branch is fixed,

\[
\prod_{e\in\gamma}w(e)
=\exp\!\left(\sum_{e\in\gamma}\log w(e)\right),
\]

and, with the appropriate discrete action,

\[
K(x,y)=\sum_{\gamma:x\to y}e^{-S(\gamma)/\hbar}.
\]

The equation combines the operations:

\begin{itemize}
\item the product composes successive steps;
\item the logarithm accumulates the local action;
\item the exponential reconstructs the route weight;
\item the suma agrega all the alternativas.
\end{itemize}

APP is the minimal support on which sum and product act on the same histories. TPK adds the data that prevent the sum over paths from becoming a sum over values without genealogy.

In this framework, a digit is the output of a reader; a value is the Archimedean projection of a section; an HMT number preserves the compatible family of paths that supports that value; a particle preserves that genealogy under a persistent spectral realization.


\section{Internal root of \(-1\), exponential structure, and elliptic, parabolic, and hyperbolic geometries}

\subsection{Internal root of \(-1\) in the finite branch}

The Paley incidence matrix provides an operator \(A_W\) with

\[
A_W^2=-I.
\]

This turns \(\mathbb F_3^6\) into a three-dimensional space over

\[
\mathbb F_9=\mathbb F_3[\iota]/(\iota^2+1).
\]

The root of \(-1\) is not adjoined from outside as complex notation. It emerges as an endomorphism of an already constructed ternary support.

\subsection{Continuous realization of the induced quadratic type}

In the elliptic real branch, the generator \(J_{\rm e}\) satisfies

\[
J_{\rm e}^2=-I,
\]

and

\[
e^{tJ_{\rm e}}=\cos t\,I+\sin t\,J_{\rm e}.
\]

In particular,

\[
e^{\frac\pi2J_{\rm e}}=J_{\rm e},\qquad
e^{\pi J_{\rm e}}=-I,\qquad
e^{2\pi J_{\rm e}}=I.
\]

Thus, \(J_{\rm e}\) is the internal quarter-turn and Euler's equation is its half-turn closure.

\subsection{Common operator support of \(3\) geometric readings}

Let

\[
\mathbb D=\{z\in\mathbb C:|z|<1\}.
\]

This same set contains the real diameter \((-1,1)\), the complex
interior, the trigonometric boundary \(S^1\), and the Poincaré
hyperbolic metric

\[
ds_{\rm P}^2=\frac{4\ell^2|dz|^2}{(1-|z|^2)^2}.
\]

The map

\[
R_{\pi/2}(z)=iz
\]

preserves \(|z|\) and \(|dz|\). By consiguiente, is simultaneamente:

\begin{itemize}
\item complex multiplication by \(\sqrt{-1}\);
\item trigonometric rotation of \(90^\circ\);
\item Euclidean isometry of the disk;
\item Poincaré hyperbolic isometry;
\item automorphism of its circular boundary.
\end{itemize}

The coincidence is literal: the same map acts on the same carrier; the metric reader changes.

\subsection{Elliptic, parabolic, and hyperbolic geometries of the extended Euler equation}

The generators

\[
J_{+1}^2=-I,\qquad
J_0^2=0,\qquad
J_{-1}^2=I
\]

produce

\[
e^{tJ_{+1}}=\cos t\,I+\sin t\,J_{+1},
\]

\[
e^{tJ_0}=I+tJ_0,
\]

\[
e^{tJ_{-1}}=\cosh t\,I+\sinh t\,J_{-1}.
\]

The same exponential series separates into three regimes according to the
square of the generator: elliptic, parabolic, and hyperbolic. The
specialization

\[
e^{2\log\varphi\,J_{-1}}=F_{\rm av}^{\,2}
\]

integrates the Fibonacci self-scale. Euler and Fibonacci appear as
specializations of a common quadratic exponential family.

\subsection{Common carrier, common operator, and quadratic exponential family}

Geometric fusion is formulated at three levels:

\begin{enumerate}
\item \textbf{common carrier:} the unit disk contains the real diameter, the complex interior, and the circular boundary;
\item \textbf{common operator:} \(z\mapsto iz\) is a quarter-turn, complex multiplication, and a hyperbolic isometry;
\item \textbf{common family:} the tritic exponential equation brings together the elliptic, parabolic, and hyperbolic branches.
\end{enumerate}

The result does not identify the Euclidean and hyperbolic metrics. It states
that they are distinct realizations on the same support and share the
automorphism originating in the internal root of \(-I\).


\section{Restriction of internal rotation, edge holonomy and minimal torsion}

The construction contains two positive and differentiated objects:

\begin{enumerate}
\item the internal generator \(J_{\rm e}\), with \(J_{\rm e}^2=-I\);
\item the global invariant
\end{enumerate}
   \[
   \Delta_4=
   \frac{\pi-e\log\pi}{270}
   \left(1+\frac{\pi}{729}\right),
   \]
   denoted minimal torsion and used prior to species selection.

The authorial thesis adds that internal rotational dynamics, when published
at the boundary, is absorbed as minimal torsion. The relation must not be
written as a bare equality \(J_{\rm e}=\Delta_4\): the objects have different types.
Its natural form requires

\[
\text{interior rotation}
\xrightarrow{\operatorname{Res}_{\partial}}
\text{boundary holonomy}
\xrightarrow{\operatorname{Def}}
\text{scalar defect}.
\]

Let \(q_\partial\) be the boundary reader, \(\nabla_\partial\) the induced connection, and \(e_\partial\) the discrete frame. The edge torsion has type

\[
T_\partial=d_{\nabla_\partial}e_\partial.
\]

The author's claim acquires probative form through a functional

\[
\Lambda_\partial:
\operatorname{Tors}(\partial X)
\longrightarrow\mathbb R
\]

such that

\[
\Lambda_\partial(T_\partial)=\Delta_4
\]

and a theorem deriving \(T_\partial\) from the restriction of \(J_{\rm e}\) through the boundary operators of the TPK.

The internal root of \(-I\), the quarter-turn, the extended Euler equation, and the formula for \(\Delta_4\) are proved in their respective domains. The passage from interior rotation to minimal boundary torsion requires the declared morphism \(\operatorname{Res}_{\partial}\) and functional \(\Lambda_{\partial}\); it is not identified by a scalar equality.


\section{Generational triplicity, conjugation and mixing on the integral holonomic state}

The phrase of Isidor Isaac Rabi may serve as a historical reference, but it must not govern the title or constitute an autonomous section. The scientific question is:

\begin{quote}
What internal mechanism distinguishes three stable orders of fermionic realization and preserves that distinction under color, conjugation, and mixing?
\end{quote}

The HMT response is not "because TRIT has three values". The structure separates

\[
\operatorname{flav}(X)
=\bigl(\sigma_{ud}(X),g(X),\chi_{\rm fina}(X)\bigr),
\]

where \(\sigma_{ud}\) distinguishes branches, \(g\) distinguishes physical generations, and \(\chi_{\rm fina}\) preserves orientation and memory within a branch. For quarks, the color orbit is added after flavor and before the character.

Triplicity is obtained as a consequence of the enriched structure; its
effects on color, conjugation, and mixing are developed below:

\begin{enumerate}
\item colour, flavour, generation and mass are distinct readers of the same integral holonomic state;
\item particle–antiparticle conjugation acts on the orientation without erasing the generational class;
\item CKM and PMNS transport already typed state bases; they do not create triplicity;
\item the neutrino classification must be derived from representation, occupation, and mixing, not by renaming depths \(G0,\ldots,G4\).
\end{enumerate}

That neutrinos are half-integer-spin leptons obeying fermionic statistics is a known physical classification. The structural contribution of HMT consists in obtaining internally why the neutral sector occupies that representation, why it admits three physical realizations, and how it couples to the mixing operator without reading its external angles.


\section{Operator memory of the TPK: preservation of residue, quotient, carry, sheet, orientation, route, and boundary under projection}

\subsection{Loss of reversibility and genealogical lifting}

In APP, different histories may produce the same residue. The product folds more histories than the sum; the power folds more than the product; the root attempts to reconstruct branches already identified. TPK preserves the datum that makes it possible to invert that loss:

\[
\text{residue},\quad
\text{quotient},\quad
\text{carry},\quad
\text{sheet},\quad
\text{orientation},\quad
\text{route},\quad
\text{boundary},\quad
\text{memory}.
\]

An algebraic operation may be characterized not only by its value, but also by the type of genealogy that it preserves or destroys.

\subsection{Projection of Compatible Histories onto the Archimedean Continuum}

The real line appears when a family of compatible cylinders contracts under Archimedean evaluation. The real value is the shadow of the section, not its complete genealogy. The exceptional projection preserves another part as incidence and symmetry.

The path integral expresses the same principle:

\[
\text{alternatives}\xrightarrow{\sum}\text{global amplitude},
\qquad
\text{successive stages}\xrightarrow{\prod}\text{route weight}.
\]

The logarithm converts multiplicative composition into additive action; the exponential returns it to amplitude. The discrete structure of the continuum can be formulated as genealogical conservation of all routes before their projection onto \(\mathbb R\).

\subsection{Distinction between numerical section and persistent realization of particle}

A number and a particle share minimal ontological invariants: memory, orientation, persistence, closure, signature, and position in a family of prolongations.

The number publishes an Archimedean coordinate. The particle adds spectral realization, occupation, propagation, and stability. The passage HMT \(\to\) MD changes the reader applied to the same enriched genealogy; it does not arbitrarily change the object.


\section{Joint closure of the radical network, the power--vacancy law, Euler geometry, and the operative hierarchy}

\subsection{Theorem of the Nonadic Radical Network}

The APP net \(3,6,9\) is the visual representation of the nilpotent radical \((3)\subset\mathbb Z/9\mathbb Z\). Its phases \(369\), \(693\), and \(936\) realize the isomorphism \(R/(3)\cong(3)\) induced by multiplication by three. The bands locate exactly the failure of invertibility of the product.

\subsection{Power-void Theorem}

For every \(n\ge1\), the number of decimal ranks lost by \(9^n\) relative to \(10^n\) is \(\lceil n\log_{10}(10/9)\rceil\). Their increments form the mechanical vacancy word. In \(n=9^9\), the retained coordinate is \(369\,693\,099\) and the resulting length is \(369\,693\,100\).

\subsection{Exponential Family Euler's Geometric Theorem}

The internal root of \(-I\) produces a quarter-turn. On the unit disk, the same map \(z\mapsto iz\) is simultaneously complex multiplication, trigonometric rotation, and a Poincaré isometry. The tritic exponential equation extends this regime to the parabolic and hyperbolic regimes, incorporating the Fibonacci self-scaling in the latter.

\subsection{Operatorial conservation principle between levels}

The sum--product--power hierarchy is a hierarchy of fibers and loss of
reversibility. The TPK is the lift that preserves the coordinates required
to distinguish histories identified by the ordinary readers.


\section{Typing hypotheses and negative controls}

\begin{enumerate}
\item Do not identify \(369\,693\,100\) with a TPK word without declaring the reader between domains.
\item No confundir the ideal \(\mathfrak m^2=0\) With \(\mathfrak m\times\mathfrak m\).
\item Do not confuse modulo nine with digital root: \(9\mapsto0\) in the former and \(9\mapsto9\) in the latter.
\item Do not use the threefold character of TRIT as a sufficient proof of three physical generations.
\item Do not turn \(J^2=-I\) into \(\Delta_4\) without the boundary map and the torsion functional.
\item Do not present \(369\) as a prefix of \(9^{9^9}\): it pertains to its order and its decimal length.
\item Do not reduce the ten thousand digits to an order-of-magnitude test: they certify extensive truncations of a parametrizable process whose thesis is indefinite prolongation.
\end{enumerate}

\section{Joint consequences for the radical net, the operational memory, and the edge torsion}

The principal finding is not that a tower visually contains the digits \(3,6,9\). The nonadic power and the vacancy calendar are two presentations of the same logarithmic defect, while APP localizes \(3,6,9\) as the radical where multiplication loses invertibility. The tower connects both extremes:

\[
\text{nonadic radical}
\longrightarrow
\text{power}
\longrightarrow
\text{logarithm}
\longrightarrow
\text{vacancy}
\longrightarrow
\text{decimal order}
\longrightarrow
\text{closure }9\to1.
\]

The extended Euler equation adds the second direction: the internal root of \(-I\) turns the discrete structure into rotational dynamics and unites, on a common support, the real line, trigonometric circle and hyperbolic disk. Subsequent passage to minimal torsion must preserve the boundary map already developed in the corpus, not replace it with a nominal identification.
\endgroup

\chapter{Radion, angular phase and equivalence between cyclic and linear realizations}
\begingroup
\renewcommand{\chapter}[2][]{}%
% Integración científica del capítulo del radión. La numeración y los rótulos
% se adaptan a la monografía; los cuerpos matemáticos permanecen íntegros.
% Adaptador de compatibilidad del fuente teoremática íntegro del radión.
%
% El tratado rector ya carga tipografía, geometría, hipervínculos, teoremas,
% tablas, TikZ y tcolorbox.  Este archivo aporta únicamente los nombres y
% entornos exclusivos del fuente teoremática, sin importar su preámbulo autónomo ni
% alterar la jerarquía tipográfica general.

\makeatletter
\@ifundefined{color@RadNavy}{\definecolor{RadNavy}{HTML}{102A43}}{}
\@ifundefined{color@RadBlue}{\definecolor{RadBlue}{HTML}{1F5E7A}}{}
\@ifundefined{color@RadTeal}{\definecolor{RadTeal}{HTML}{168C8C}}{}
\@ifundefined{color@RadCopper}{\definecolor{RadCopper}{HTML}{B66A3C}}{}
\@ifundefined{color@RadGold}{\definecolor{RadGold}{HTML}{D4A73A}}{}
\@ifundefined{color@RadInk}{\definecolor{RadInk}{HTML}{1A202C}}{}
\@ifundefined{color@RadMist}{\definecolor{RadMist}{HTML}{EDF4F7}}{}
\@ifundefined{color@RadCream}{\definecolor{RadCream}{HTML}{FAF7F0}}{}
\@ifundefined{color@RadRose}{\definecolor{RadRose}{HTML}{8D3F55}}{}
\@ifundefined{color@RadGreen}{\definecolor{RadGreen}{HTML}{2F7D5D}}{}
\@ifundefined{color@RadGray}{\definecolor{RadGray}{HTML}{66788A}}{}

\providecommand{\TRIT}{\ensuremath{\mathrm{TRIT}}}
\providecommand{\HMT}{\ensuremath{\mathrm{HMT}}}
\providecommand{\Rstate}{\mathcal R_{12}}
\providecommand{\epsR}{\varepsilon_R}
\providecommand{\epsone}{\varepsilon_R^{(1)}}
\providecommand{\pih}{\pi}
\providecommand{\alphah}{\alpha}
\providecommand{\hret}{\hbar_{\mathrm{ret}}}
\providecommand{\hfull}{h_{\mathrm{ret}}}
\providecommand{\diff}{\mathop{}\!\mathrm d}
\providecommand{\e}{\mathrm e}
\providecommand{\R}{\mathbb R}
\providecommand{\C}{\mathbb C}
\providecommand{\F}{\mathbb F}
\providecommand{\Z}{\mathbb Z}
\providecommand{\dd}{\,\mathrm d}
\providecommand{\status}[1]{\textsf{\bfseries #1}}
\providecommand{\sourcepath}[1]{\nolinkurl{#1}}

\@ifundefined{typebox}{%
  \newtcolorbox{typebox}[1][]{enhanced,breakable,colback=white,
    colframe=RadBlue,boxrule=.7pt,arc=1.5mm,left=4mm,right=4mm,
    top=3mm,bottom=3mm,title={Type receipt},fonttitle=\bfseries,
    coltitle=RadNavy,#1}%
}{}
\@ifundefined{narrativebox}{%
  \newtcolorbox{narrativebox}[1][]{enhanced,breakable,colback=white,
    colframe=RadGold,borderline west={2.4pt}{0pt}{RadGold},
    boxrule=.25pt,arc=0mm,left=5mm,right=4mm,top=3mm,bottom=3mm,#1}%
}{}

% Las once portadillas internas del fuente teoremática autónomo se convierten en
% divisiones científicas ordinarias dentro del capítulo 52.  Se evita así
% introducir partes autónomas, gigantismo tipográfico o páginas vacías.
\providecommand{\partopener}[3]{\section{#2}\noindent\emph{#3}\par\medskip}

\providecommand{\BigRadionEquation}{%
  \begin{equation*}
  \boxed{\displaystyle
  \frac{180^{\circ}}{\pih}
  =50^{\circ}+A^{\circ}-\frac{20}{81}\,\Delta_4^{\circ}+\epsR^{\circ}}
  \end{equation*}}
\makeatother
 
\chapter{The angular radion: carry-mediated joining of direct and torsional sections and exact closure of the phase, action and memory unit}
\label{chap:sucesor-102-radion}

\section{Composition APP--TRIT--TPK of the phase state: generation of the memory-oriented unit}
% Fragmento público generado desde fuentes teoremáticas íntegros.
% Residencia: 52.1.
% Fuente: /Users/ruben/Documents/HMT2/EL_CIERRE_HOLOGRAFICO_DEL_INFINITO_SUCESOR_102_CAPITULOS_2026-08-28_EN_TRABAJO/incoming/radion_monografia_20260827/manuscrito/sections/01_resumen_y_guia.tex:1-128
\paragraph{Mathematical Synthesis of the Object and Its Causal Chain}

This monograph reconstructs the radian from the generative chain
\[
\APP\longrightarrow\TRIT\longrightarrow\TPK
\longrightarrow X_{\mathrm{enr}}
\longrightarrow\text{discrete structure of the continuum}
\longrightarrow\MD,
\]
without taking as input either the metrological value of \(180/\pi\) or the small angular residue to be explained. The result is an enriched unit, the \emph{radion}, defined as the oriented unit that traverses one radian of phase and sweeps exactly an action \(\hret\) in the positive ellipse associated with the same state.

The governing closure is
\BigRadionEquation
with
\[
S_E=2\pih\left(\frac5{36}+\frac{25}{9}\alphah\right),
\qquad
\Delta_4=
\frac{\pih-e\log\pih}{270}
\left(1+\frac{\pih}{729}\right),
\]
\[
\rho_{\mathrm{tw}}=\frac{2\cdot90}{729}=\frac{20}{81},
\qquad
\epsone=\frac{20}{81}\Delta_4-(S_E-1).
\]
APP division proves \(1<S_E<2\), so that the quotient is exactly
one and the remainder is \(D_E=S_E-1\). The equality
\[
1=S_E-\frac{20}{81}\Delta_4+\epsone
\]
is then an internal identity. The angular chart produces
\[
\epsR^\circ
=5.13830167585752820716851646226459474899085744\ldots
\times10^{-8}\,{}^\circ.
\]

The text also demonstrates that the action ellipse
\[
Q=a_S\cos\theta,\qquad P=b_S\sin\theta,\qquad
a_Sb_S=2\hret
\]
satisfies
\[
\operatorname{Area}(0,\theta)=\hret\theta.
\]
Therefore, a radion sweeps \(\hret\) and the turn \(2\pi\) encloses
\(h_{\mathrm{ret}}=2\pi\hret\). The same ellipse has a finite symplectic
boundary and an interior of infinite area when equipped with the
Hilbert--Beltrami--Klein metric. This is the exact form of the physical thesis guiding
the work: the visible surface closes the action; the interior depth
preserves a development that cannot be exhausted by the planar projection.

The exposition maintains typed four realizations of the same state:

\begin{enumerate}[label=\textbf{\arabic*.}]
\item the circle \(S^1\), which publishes the phase;
\item the positive ellipse, which publishes the action and the anisotropy;
\item the interior of Klein, that miof the depth by means of reason doble;
\item the nonadic helix, which preserves return, memory, and contraction.
\end{enumerate}

These are not four geometries superimposed without a criterion. They are four readers with different domains, codomains and preserved data, applied to a single integral holonomic state. The sector \(90/120\), the dodecaphase, radion carry, the Lambert tail, the Barbero kernel and global towers are likewise presented as related but noninterchangeable objects.

\begin{thesisbox}
The radian is the oriented unit that sweeps \(\hret\) in the action ellipse.
The circle publishes its phase, Klein measures its depth, the helix preserves its
memory, and the coinductive carry exactly splices its direct and
torsional sections.
\end{thesisbox}

\paragraph{Expository order and reading dependencies}

The order is not ornamental. Parts I and II produce the object and its
closure from APP--TRIT--TPK. Only thereafter are realizations introduced in
the languages of Hilbert, LC, Maxwell, Minkowski, Hodge, Lorentz,
modularity, and orbits. Each chapter distinguishes three levels:

\begin{description}[style=nextline,leftmargin=4.2cm,labelwidth=3.9cm]
\item[Internal result.] Equality or construction proved in the
HMT domain with its declared primitives.
\item[Consolidated formalization.] A new map articulating already
existing results without using the target as a selector.
\item[Physical interface.] Subsequent dimensional translation that recognizes a
conventional structure without making it the origin of the selection.
\end{description}

The \emph{status and falsifier} boxes indicate which concrete remark would invalidate each strong promotion. This discipline does not weaken the thesis; it makes the thesis auditable. The purpose of the extended narrative is to ensure that no word such as \emph{circle}, \emph{ellipse}, \emph{torsion}, \emph{charge}, or \emph{memory} conceals a change of type.

\paragraph{Pivotal results and proof scope}

\begin{enumerate}[label=\textbf{C\arabic*.}]
\item Closure of the radion independent of its final value through two publications of the same state and
APP subtraction with carry.
\item Derivation of the coefficient \(20/81\) as two sheets of the active sector
\(90\), normalized by the fiber \(3^6=729\).
\item Typed equalizer between the seam \(\Re s=1/2\), the hundred-mark APP
face, and the dodecaphasic phase \(\theta_2=50^\circ\).
\item Radian area theorem: \(\operatorname{Area}(\theta)=\hret\theta\).
\item Exact focal study of the ellipse and of its Euclidean,
hyperbolic, symplectic, and modular invariants.
\item Boundary--interior theorem: finite total action \(h\) and Klein interior
of infinite hyperbolic area.
\item Exact realization through minimal covariance and an LC oscillator, followed by
the electromagnetic publication \(90/120\).
\item Atlas of non-identifications for all objects \(90/120\), including
the separation among \(\epsR\), the spectral tail, and the Barbero kernel \(12:-1\).
\item Structural characterization of \(\pi\) in the radion chart, distinct
from the primary coinductive generator.
\item Ontological synthesis: dimension as the minimal rank of memories that a
phase projection can no longer preserve.
\end{enumerate}
% Fuente: /Users/ruben/Documents/HMT2/EL_CIERRE_HOLOGRAFICO_DEL_INFINITO_SUCESOR_102_CAPITULOS_2026-08-28_EN_TRABAJO/incoming/radion_monografia_20260827/manuscrito/sections/01b_mapa_resultados.tex:1-133
\paragraph{Correspondence between results, causal types, and causal residences}

The following correspondence orders the governing assertions and their proofs,
while keeping visible the hierarchy that prevents
physical richness from erasing the mathematical closure supporting it.

% HMT_RESULT_BEGIN:RDN-01-GENEALOGIA:theorem
\begin{exactbox}[title={APP--TRIT--TPK genealogy of the radion}]
The radion arises from the effective composition
\(\APP\to\TRIT\to\TPK\to X_{\mathrm{enr}}\to\MD\); it is not obtained
by retroactively redefining the conventional radian.
\end{exactbox}
% HMT_RESULT_END:RDN-01-GENEALOGIA:theorem

% HMT_RESULT_BEGIN:RDN-02-CIERRE:theorem
\begin{exactbox}[title={Exact closure and prospective independence}]
APP carry internally produces
\[
1=S_E-\frac{20}{81}\Delta_4+\varepsilon_R^{(1)},
\qquad
\varepsilon_R^{(1)}=\frac{20}{81}\Delta_4-(S_E-1),
\]
and the angular chart transports this same identity to
\(180^\circ/\pi\) without introducing the target residue.
\end{exactbox}
% HMT_RESULT_END:RDN-02-CIERRE:theorem

% HMT_RESULT_BEGIN:RDN-03-BISAGRA-50:theorem
\begin{exactbox}[title={Critical coordinate of \(50^\circ\)}]
The critical coordinate \(\Re s=1/2\), published on the APP face of one hundred
marks and in the dodecaphase wheel, selects the phase
\(\theta_2=50^\circ\). It is an equality of typed readers of the same state.
\end{exactbox}
% HMT_RESULT_END:RDN-03-BISAGRA-50:theorem

% HMT_RESULT_BEGIN:RDN-04-COEFICIENTE-2081:theorem
\begin{exactbox}[title={Nonadic Bidirectional Density \(20/81\)}]
The torsional coefficient is not fitted on the lattice \(1/81\): it is constructed as
\(\rho_{\mathrm{tw}}=2N_6/3^6=2\cdot90/729=20/81\), preserving the sheet,
orientation, and normalization of the nonadic fibre.
\end{exactbox}
% HMT_RESULT_END:RDN-04-COEFICIENTE-2081:theorem

% HMT_RESULT_BEGIN:RDN-05-AREA-ACCION:theorem
\begin{exactbox}[title={Areal law of the radion}]
For the positive ellipse
\(Q=a_S\cos\theta\), \(P=b_S\sin\theta\),
\(a_Sb_S=2\hbar_{\mathrm{ret}}\), the oriented area swept is
\(\mathcal A(\theta)=\hbar_{\mathrm{ret}}\theta\). One radion sweeps exactly
\(\hbar_{\mathrm{ret}}\), and the revolution \(2\pi\) encloses
\(h_{\mathrm{ret}}=2\pi\hbar_{\mathrm{ret}}\).
\end{exactbox}
% HMT_RESULT_END:RDN-05-AREA-ACCION:theorem

% HMT_RESULT_BEGIN:RDN-06-INVARIANTES-FOCALES:theorem
\begin{exactbox}[title={Focal invariants of the positive operator}]
The published pair \(A\pm C^\ast\) fixes the semiaxes, eccentricity,
focal separation, and symplectic compression parameter. In particular,
\((d_1/d_0)^2=C^\ast/2\), a dimensionless identity separating shape from scale.
\end{exactbox}
% HMT_RESULT_END:RDN-06-INVARIANTES-FOCALES:theorem

% HMT_RESULT_BEGIN:RDN-07-BORDE-INTERIOR:theorem
\begin{exactbox}[title={Symplectic finiteness of the boundary and divergence of the interior hyperbolic area}]
The symplectic boundary of the ellipse encloses the finite action
\(h_{\mathrm{ret}}\), whereas the Hilbert--Beltrami--Klein area of the interior
diverges upon reaching the boundary. Visible phase and depth are not the same
measure.
\end{exactbox}
% HMT_RESULT_END:RDN-07-BORDE-INTERIOR:theorem

% HMT_RESULT_BEGIN:RDN-08-HELICE-MEMORIA:theorem
\begin{exactbox}[title={Nonadic Helical Lift with Phase Return and Memory Advance}]
The nonadic monodromy returns the phase without resetting the state. The helix
lifts the circle while preserving carry and memory; the closure \(2\pi/4\pi\)
distinguishes projective return, reversal of orientation, and complete return.
\end{exactbox}
% HMT_RESULT_END:RDN-08-HELICE-MEMORIA:theorem

% HMT_RESULT_BEGIN:RDN-09-OSCILADOR-EM:development
\begin{narrativebox}[title={Oscillatory and constitutive electromagnetic realization}]
The same anisotropy is realized through a minimal covariance and an LC oscillator:
one coordinate stores the electric--elastic aspect and its conjugate the
magnetic--rotational aspect. Frequency and impedance remain separately typed.
\end{narrativebox}
% HMT_RESULT_END:RDN-09-OSCILADOR-EM:development

% HMT_RESULT_BEGIN:RDN-10-TIPOS-90120:theorem
\begin{exactbox}[title={Typed separation of the operators \(90/120\)}]
Dodecaphase extractor, radion carry, Lambert tail,
hexad--octad functional, Barbero kernel, and global tower semigroup are
distinct objects. The weight \(12:-1\) is forced only upon fixing the pole \(1/24\) and
the minimal integer counterterm.
\end{exactbox}
% HMT_RESULT_END:RDN-10-TIPOS-90120:theorem

% HMT_RESULT_BEGIN:RDN-11-GEOMETRIAS-CAMPO:development
\begin{narrativebox}[title={Complex, Lorentzian, and Hodge-dual realizations}]
The elliptic, parabolic, and hyperbolic regimes of TRIT are published as
rotation, threshold, and hyperbolic transformation. The Hodge star, the constitutive relations, and the Lorentz
force act subsequently on the already oriented state; they do not select its
constants.
\end{narrativebox}
% HMT_RESULT_END:RDN-11-GEOMETRIAS-CAMPO:development

% HMT_RESULT_BEGIN:RDN-12-MODULAR-ORBITAL:development
\begin{narrativebox}[title={Modular, orbital, and arithmetic boundary characters}]
The modular involution, elliptic orbits, advance--delay, and
3--6--9 arithmetic recognize distinct memories of the same state. The modular function and
vacancies are used as posterior readers, never as substitutes for the
APP--TRIT--TPK generator.
\end{narrativebox}
% HMT_RESULT_END:RDN-12-MODULAR-ORBITAL:development

% HMT_RESULT_BEGIN:RDN-13-SINTESIS-ONTOLOGICA:conclusion
\begin{thesisbox}[title={Operatorial synthesis of phase, action, depth, and memory}]
The circle is the phase projection; the ellipse is the action section; Klein
measures depth; the helix preserves history. The radion is the oriented unit
that keeps these publications coupled without erasing the residue that
distinguishes them.
\end{thesisbox}
% HMT_RESULT_END:RDN-13-SINTESIS-ONTOLOGICA:conclusion

% HMT_RESULT_BEGIN:RDN-14-CONTROL-ALCANCE:control
\begin{statusbox}[title={Delimitation of domains and causal continuity}] The complete mass laws, electron refinement, CKM/CP, and the exceptional corridor retain their own governing proofs. This monograph cites them only where they share an antecedent or chart; it neither transfers their specific selectors to the radion nor confuses their probative residences.
\end{statusbox}
% HMT_RESULT_END:RDN-14-CONTROL-ALCANCE:control
 
\subsection{Bivariate APP support, ternary TRIT orientation and updating of the integral holonomic state through TPK}
% Fragmento público generado desde fuentes teoremáticas íntegros.
% Residencia: 52.1.1.
% Fuente: /Users/ruben/Documents/HMT2/EL_CIERRE_HOLOGRAFICO_DEL_INFINITO_SUCESOR_102_CAPITULOS_2026-08-28_EN_TRABAJO/incoming/radion_monografia_20260827/manuscrito/sections/02_genealogia.tex:1-180
\noindent The operative genealogy precedes angular measurement and determines the variables preserved by each reader.\par\medskip

\subsubsection{APP--TRIT--TPK support and construction of the integral holonomic state}

\paragraph{Priority of the integral holonomic state over the circular expression}

The elementary definition of radian begins with an already given circle, an already
given center, and an already measurable arc length. If the radius and arc have the
same length, the subtended angle measures one radian. The definition is correct
within Euclidean geometry, but does not answer a prior question:
what structure makes it possible for a turn to possess phase, orientation, action, and
memory, and why is the natural unit related to \(2\pi\)?

HMT reverses the causal order. It first constructs the finite state, then its
coinductive prolongation, then the revolution, and only at the end its angular
publication. The radion does not correct the conventional radian; it lifts it to its
genealogical domain.

\begin{typebox}
\[
\APP\to\TRIT\to\TPK\to X_{\mathrm{enr}}
\to\mathcal C^{\mathrm{disc}}_{\mathrm{cont}}
\to\bigl(\pih,e,\varphi,\alphah,A,C^*,\Delta_4,\hret\bigr)
\to\mathfrak r_\sigma.
\]
Conventional mathematics appears only afterward as a language of realization:
\(S^1\), symplectic ellipse, Klein metric, Hilbert space, LC
circuit, Minkowski screen, and field equations.
\end{typebox}

\paragraph{Bivariate support of APP: leaves, residues, quotients, and carry}

The Positional Arithmetic Platform (APP) simultaneously preserves two
realizations of the symbols \(1,\dots,9\): an additive sheet and a
multiplicative sheet. A number is not reduced to its published value. Alongside the value
remain sheet, orientation, quotient, remainder, carry, boundary, route, and
survival.

The basic census has three strata that must never be identified by mere cardinal coincidence:
\[
104,976\longrightarrow468\;U_6\longrightarrow243\;w_6
\subset\F_3^6,\qquad |\F_3^6|=729.
\]
The \(468\) decimal emissions, the \(243\) visible words, and the ambient set
of \(729\) words are distinct objects. The ambient set will be decisive for the
normalization \(20/81\).

In cubic base \(B=1000\), the oriented subtraction of two triad words is
performed from right to left:
\begin{equation}
u_i=t_i-v_i+c_{i+1},\qquad
r_i=u_i\bmod1000,\qquad
c_i=\frac{u_i-r_i}{1000}.
\label{eq:subacarreo}
\end{equation}
Once the operands and the remote carry are fixed, quotient and remainder are unique. This is
the operator that will subsequently publish the digits of \(\epsone\). No
decimal is chosen after seeing the radian.

The completion
\[
10^3=729+243+27+1
\]
is not a decimal curiosity. It expresses conservation of the unit in passing
from the ternary ambient space to a chamber of one thousand positions. The difference
\(1000-729=271\) is the completing corona; the correct identity is
\(999=729+270\), not \(729+271\).

\paragraph{Ternary orientation and classification of the 3 operational regimes of the TRIT}

TRIT assigns one of three oriented regimes to each local state. Under the
active temporal convention,
\[
\begin{aligned}
+1&\longleftrightarrow\text{return, memory, past},\\
0&\longleftrightarrow\text{threshold, present},\\
-1&\longleftrightarrow\text{bidirectional opening, future}.
\end{aligned}
\]
The quadratic realization distinguishes:
\begin{equation}
J^2=-I,\qquad N^2=0,\qquad R^2=I.
\label{eq:tres-regimenes}
\end{equation}
Thus arise, respectively, elliptic rotation, the parabolic threshold, and
hyperbolic separation. Three metrics are not being mixed. Three typed
actions on a common antecedent are preserved.

The internal root of \(-1\) is not imported by presupposing the complex
plane. In the APP block, two projections \(P,Q\) with ratio
\(r=3/4=90/120\) produce
\[
J_{\mathrm{comm}}=\frac{[P,Q]}{\sqrt{r(1-r)}},\qquad J_{\mathrm{comm}}^2=-I,
\]
while
\[
R_{\mathrm{comm}}=2P-I,\qquad R_{\mathrm{comm}}^2=I,\qquad
R_{\mathrm{comm}}J_{\mathrm{comm}}=-J_{\mathrm{comm}}R_{\mathrm{comm}}.
\]
The pair generates the minimal form of \(\mathrm{Cl}_{1,1}\simeq M_2(\R)\):
rotation and opening emerge together, but remain distinct operations.

\paragraph{Dynamic composition of the TPK: selection, transport, memory, and return}

The Topological Phase Kernel is not a name for a box of objects. It is the
composition that selects phase, transports cardinality, lifts sheets,
records memory, and executes return. The fundamental dependence is
\[
U_t\longrightarrow\Gamma_9\longrightarrow K_{\mathrm{ph}}.
\]
The nonadic holonomy \(\Gamma_9\) acts on the integral holonomic state; the expression \(K_{\mathrm{ph}}\) is subsequent and preserves reduced memory. The identity \(K_{\mathrm{ph}}^9|_r=E_+E_\times\) does not turn \(K_{\mathrm{ph}}\) into a generator.

The nonadic certificate preserves the complete chain
\[
w_6\to w_{12}\to w_{18}\to w_{24}\to w_{30}
\to R_{36}\to G_9,
\]
with \(396\) levels, \(2376\) trits, \(43\) circuits, \(387\) pairs
\((K,K+9)\) per channel, and five regions of \(\pi\). The phase satisfies
\[
g(k+9)=g(k),
\]
but the complete state does not return:
\[
B_{k+9}\ne B_k.
\]
This is the mathematical seed of the helix: the angular projection closes; the
memory advances.

\paragraph{Non-split temporal extension \(C_{12}\to C_{108}\to C_9\)}

The twelve-step memory is organized via
\begin{equation}
0\longrightarrow C_{12}\longrightarrow C_{108}
\longrightarrow C_9\longrightarrow0.
\label{eq:extension-108}
\end{equation}
The cocycle
\[
c(r,r')=\left[\left\lfloor\frac{r+r'}9\right\rfloor\right]_{12}
\]
measures the carry between phase return and memory advance. Nine steps close
the observable phase \(C_9\); they do not reset the dodecaphasic
record.

It is convenient to separate four objects:
\begin{center}
\begin{tabularx}{.95\textwidth}{>{\bfseries}l X}
\toprule
\(C_{12}\) & cyclic group of twelve positions;\\
\(\Rstate\) & enhanced record \((E_m,\tau_m,\sigma_m,\kappa_m,\Lambda_m)_{m=1}^{12}\);\\
\(\Theta_{108}\) & raw publication on the 108 calendar;\\
\(\Gamma_{108}\) & history or lifted holonomy.\\
\bottomrule
\end{tabularx}
\end{center}
Accordingly, the vague expression
\enquote{post-\(C_{12}\)} shall be replaced by
\enquote{subsequent to the dodecaphasic record
\(\Rstate\), in its hexadic} reading when that is the effective
map.

\paragraph{Correlative emergence of the 5 constructions of the continuum from the unique state}

The discrete structure of the continuum is not obtained by assembling five
independent theories. Its five internal aspects---spectral, solenoidal, gauge,
cohomological, and determinant---emerge correlatively from the same
APP--TRIT--TPK state. This work does not prove the integral treatise anew, but
preserves the causal fact required by the radion: the coinductive section that
generates \(\pi,e,\varphi\) and \(\alpha\) does not live on a bare real line; it lives
in an object with fibers, relations, orientation, carry, memory, terminal object, and
reader.

\begin{statusbox}
\textbf{Genealogy falsifier.} If a later formula receives
\(180/\pi\), \(\epsR\), a mass, or a metrological table in order to choose a phase,
a coefficient, or a branch, it ceases to be an HMT-first generation. The external
contrast may recognize an output; it cannot select it.
\end{statusbox}
 
\subsection{Prospective independence of the closure with respect to the final value}
% Fragmento público generado desde fuentes teoremáticas íntegros.
% Residencia: 52.1.2.
% Fuente: /Users/ruben/Documents/HMT2/EL_CIERRE_HOLOGRAFICO_DEL_INFINITO_SUCESOR_102_CAPITULOS_2026-08-28_EN_TRABAJO/incoming/radion_monografia_20260827/manuscrito/sections/02_genealogia.tex:181-272

\subsubsection{Genealogy of the constants and construction of the angular pair}

\paragraph{Coinductive section and uniqueness of the generated value of \(\pi\)}

Be \(x_\infty\) the section global of the inverse system. The lectores
coinductivos produce
\[
\operatorname{ev}_\pi(x_\infty)=\pih,\qquad
\operatorname{ev}_e(x_\infty)=e,\qquad
\operatorname{ev}_\varphi(x_\infty)=\varphi.
\]
These evaluations determine the universal constants themselves, not variants dependent on the procedure constructing them. In particular,
\[
\pih=3.14159265358979323846264338327950288419716939937510\ldots.
\]
The positive loop is
\begin{equation}
T_+(x_\infty)=2\pih.
\label{eq:vuelta-plus}
\end{equation}
Its angular publication of a unit is
\begin{equation}
R_{\HMT}^{\circ}=\frac{360^\circ}{T_+}=\frac{180^\circ}{\pih}.
\label{eq:radian-publicado}
\end{equation}
The equation does not yet define the enriched radion; it defines the visible
chart of the radian once the turn has been generated.

\paragraph{Bidirectionality of conjugate routes and prolongations}

Bidirectionality is not enunciated as a slogan. The concrete involution
\begin{equation}
J_\alpha(T)=\frac{\alphah^2}{T},\qquad
T_-=J_\alpha(T_+),\qquad T_+T_-=\alphah^2
\label{eq:bidireccional-turn}
\end{equation}
preserves a product and reverses the sheet. Expansion and contraction are two
conjugate prolongations of the same state. The finite root of a
quadratic jet can approximate \(T_+\), but it does not replace the
coinductive section: the certified difference is
\[
T_{P_9}-T_+=-5.1110566290335\ldots\times10^{-25}.
\]
Finite control is a witness; the foundation is the limit.

\paragraph{Causal dependence among \(\alpha\), action, and angular publication}

The dodecaphase closure generates \(\alphah\) from the signed state and its
reversible charts. The local determinantal reader then fixes the action decade
through
\[
\Sigma_H=\varphi\alphah^{16},
\qquad \operatorname{ord}_{10}\Sigma_H=-34.
\]
The return section produces \(\hret\). Neither \(\alpha\) nor \(\hbar\) is
\enquote{measured in degrees}. They are subsequently published through the angular pair
\begin{equation}
P(\alpha,\eta_{\mathrm{ret}})=
\left([10^3\alpha]_{360},[2(\eta_{\mathrm{ret}}+\alpha)]_{360}\right)
=(A,C^*).
\label{eq:parangular}
\end{equation}
In the canonical branch,
\[
A=7.297352569283800997285105472380663\ldots{}^\circ,
\]
\[
C^*=2.1237383389686457242619513803933607937\ldots{}^\circ.
\]
Only then the analytic chart
\[
x=\frac{\pi A}{180},\qquad y=\frac{\pi C^*}{180}
\]
open the channels \(q_\pm=e^{-x\mp y}\).

The complete causal order is
\[
\alpha\to(\eta\parallel-34)\to\hbar
\to P\to\kappa_{360}\to q_\pm
\to\text{MD realizations}.
\]
The relation \(C^*/2-\alpha\) may serve as a reverse verification. It is not used
to generate \(\hret\) retroactively.

\begin{narrativebox}
The angular chart does not turn a dimensionless constant or an action into
an angle by fiat. It publishes two state coordinates on a finite wheel. The
degree is the language of that publication; it is not the ontological
substance of \(\alpha\) or \(\hbar\).
\end{narrativebox}
% Fuente: /Users/ruben/Documents/HMT2/EL_CIERRE_HOLOGRAFICO_DEL_INFINITO_SUCESOR_102_CAPITULOS_2026-08-28_EN_TRABAJO/incoming/radion_monografia_20260827/manuscrito/sections/03_cierre_exacto.tex:254-285
\paragraph{Prospective independence with respect to the final value}

The depth operator verifier explicitly prohibits three inputs:
\[
\frac{180}{\pi},\qquad \epsR,\qquad
\text{the target decimal}.
\]
First construct \(S_E\) and \(S_T\), then execute the APP subtraction, and only at
the end compare the closure. The controls give
\[
\texttt{target\_epsilon\_used=False},\qquad
\texttt{target\_180\_over\_pi\_used=False}.
\]
Also:
\begin{itemize}
\item a single sheet gives a defect of order \(10^{-5}\);
\item replacing \(\theta_2\) by \(\theta_1\) or \(\theta_3\) shifts the
result by approximately \(\pm\pi/6\);
\item removing channel \(e\) leaves a defect of order \(10^{-3}\);
\item the variante historica \(\Delta_4(A,C^*)\) produce other number;
\item the spectral tail following \(\Rstate\) and the four-term harmonic operator do not coincide with \(\epsone\) either.
\end{itemize}

\begin{statusbox}
\textbf{Status.} The closure is a \status{new formalization} proved
with an explicit initial structure and a computational certificate independent of the final value.
The equality is exact in the coinductive limit. The short 36-digit profile
of \(\alpha\) generates a variant that diverges only after many digits;
it is retained as a historical witness, not as the canonical value.
\end{statusbox}

 
\section{Direct and torsional sections: APP carry and exact closure of the radion}
\subsection{Independent construction of the sections \(S_E\) and \(S_T\) before defining the cocycle \(\varepsilon_R=S_T-S_E\)}
% Fragmento público generado desde fuentes teoremáticas íntegros.
% Residencia: 52.2.1.
% Fuente: /Users/ruben/Documents/HMT2/EL_CIERRE_HOLOGRAFICO_DEL_INFINITO_SUCESOR_102_CAPITULOS_2026-08-28_EN_TRABAJO/incoming/radion_monografia_20260827/manuscrito/sections/03_cierre_exacto.tex:1-50
\noindent Exact closure is obtained through two sections constructed before the cocycle that relates them.\par\medskip

\subsubsection{Independent construction of direct and torsional sections}

\paragraph{Direct section of the dodecaphasic temporal system}

The dodecaphase wheel has centers
\begin{equation}
\theta_m=30m-10^\circ,\qquad m=1,\ldots,12.
\label{eq:rueda-dodecafasica}
\end{equation}
The second phase is \(\theta_2=50^\circ\). Together with the publication
\(A=10^3\alphah\), it forms the direct section
\begin{align}
S_E
&=\frac{T_+}{360}\,(50+10^3\alphah)\\
&=2\pih\left(\frac5{36}+\frac{25}{9}\alphah\right).
\label{eq:seccion-electrica}
\end{align}
The subscript \(E\) here means \emph{direct or common}; its electrical
realization will come after the constitutive reader. One does not literally
postulate \(A=E\).

The coinductive cylinders prove
\begin{equation}
1<S_E<2.
\label{eq:SE-intervalo}
\end{equation}
The APP division then enforces
\begin{equation}
q_E=\lfloor S_E\rfloor=1,
\qquad
D_E=\operatorname{rem}_1(S_E)=S_E-1.
\label{eq:DE}
\end{equation}
The integer one is not the number of turns of the helix. It is the unique quotient of the direct section within the unit cell.

\paragraph{Torsional Section Determined by the Operator \(\Delta_4\)}

The second section uses a previous coinductive output:
\begin{equation}
\Delta_4=
\frac{\pih-e\log\pih}{270}
\left(1+\frac{\pih}{729}\right).
\label{eq:delta4}
\end{equation}
The word \emph{torsion} here denotes the
global quantum of the HMT chart. It does not automatically identify it
with Cartan torsion \(T^I=D_\omega e^I\). The bridge to Cartan requires another
typed arrow.
 \subsection{Oriented density of the active sector and unique determination of \(20/81\)}
% Fragmento público generado desde fuentes teoremáticas íntegros.
% Residencia: 52.2.2.
% Fuente: /Users/ruben/Documents/HMT2/EL_CIERRE_HOLOGRAFICO_DEL_INFINITO_SUCESOR_102_CAPITULOS_2026-08-28_EN_TRABAJO/incoming/radion_monografia_20260827/manuscrito/sections/03_cierre_exacto.tex:51-103

\paragraph{Oriented density \(20/81\) of the active sector \(90/120\)}

The triadic phase separator selects
\[
H_+=\{1,4,7\},\qquad H_-=\{2,5,8\},\qquad H_0=\{3,6,9\}.
\]
There are six active phases and fifteen unordered pairs:
\[
H_{\mathrm{act}}=H_+\cup H_-,\qquad
\left|\binom{H_{\mathrm{act}}}2\right|=\binom62=15.
\]
The active mode and its extension to eight positions are
\begin{equation}
N_6=6\binom62=90,\qquad
N_8=8\binom62=120.
\label{eq:90-120-incidencia}
\end{equation}
Two conjugate sheets of the active sector, normalized by the ambient fiber,
give
\begin{equation}
\rho_{\mathrm{tw}}
=\frac{2N_6}{|\F_3^6|}
=\frac{2\cdot90}{729}
=\frac{180}{729}
=\boxed{\frac{20}{81}}.
\label{eq:rho-twoway}
\end{equation}

\begin{proposition}[Relative uniqueness of the density]
With the active sector \(N_6=90\), the two conjugate sheets, and the
normalization by \(|\F_3^6|=729\) fixed, the oriented density is necessarily
\(20/81\).
\end{proposition}
\begin{proof}
The transported cardinality is \(2N_6=180\). Normalization by the ambient space
gives \(180/729\); dividing the numerator and denominator by \(9\) yields \(20/81\).
No comparison with \(180/\pi\) or with \(\epsR\) is involved.
\end{proof}

The historical reading as a projection onto a \(1/81\) lattice is a useful rigidity control, but it is no longer the foundation: the coefficient is obtained before closure. Ablating one sheet produces \(10/81\), and the output changes at order \(10^{-5}\); hence, the factor of two is not decorative.

The torsional section is
\begin{equation}
\Phi_T=\rho_{\mathrm{tw}}\Delta_4,
\qquad
S_T=1+\Phi_T.
\label{eq:seccion-torsional}
\end{equation}

 \subsection{The carry \(\varepsilon_R\) as a transition cocycle between sections}
% Fragmento público generado desde fuentes teoremáticas íntegros.
% Residencia: 52.2.3.
% Fuente: /Users/ruben/Documents/HMT2/EL_CIERRE_HOLOGRAFICO_DEL_INFINITO_SUCESOR_102_CAPITULOS_2026-08-28_EN_TRABAJO/incoming/radion_monografia_20260827/manuscrito/sections/03_cierre_exacto.tex:104-253
\subsubsection{Transition cocycle between the direct and torsional sections}

\paragraph{definition operatoria of the carry \(\varepsilon_R\)}

\begin{definition}[Real carry of the radion]
The unitary output of the radion is the transition cocycle
\begin{equation}
\epsone:=S_T-S_E
=\Phi_T-D_E
=\frac{20}{81}\Delta_4-operatorname{rem}_1(S_E).
\label{eq:epsilon-def}
\end{equation}
\end{definition}

The equation contains a subtraction, but it is not a fit to the target. Its two operands \(S_T\) and \(S_E\) are generated separately. The independence demonstrated is independence from \(180/\pi\) and from a target \(\varepsilon\); algebraic independence between a difference and its summands is not asserted.

\begin{theorem}[Exact unitary closure]
With \(S_E\), \(\Phi_T\), and \(\epsone\) defined by
\eqref{eq:seccion-electrica}, \eqref{eq:delta4}, \eqref{eq:rho-twoway}, and
\eqref{eq:epsilon-def}, one has exactly
\begin{equation}
\boxed{1=S_E-\Phi_T+\epsone.}
\label{eq:cierre-unitario}
\end{equation}
\end{theorem}
\begin{proof}
By \eqref{eq:SE-intervalo}, \(D_E=S_E-1\). By definition,
\(\epsone=\Phi_T-D_E\). Then
\[
S_E-\Phi_T+\epsone
=S_E-\Phi_T+\Phi_T-(S_E-1)=1.
\]
The strength of the proof lies not in this elementary cancellation, but in the fact that the
prior genealogy fixes both operands without using the closure as input.
\end{proof}

\paragraph{Publication of the closure on the angular chart}

The full-turn chart multiplies a unit fraction by \(360/T_+=180/\pih\). In particular,
\[
\epsR^\circ=\frac{360^\circ}{T_+}\epsone.
\]
Multiplying \eqref{eq:cierre-unitario} by \(360^\circ/T_+\) and using
\[
\frac{360}{T_+}S_E
=360\left(\frac5{36}+\frac{25}{9}\alphah\right)
=50+1000\alphah=50+A
\]
produce the angular closure.

\begin{theorem}[Exact equation of the radion]
\BigRadionEquation
\end{theorem}

The equation says something more precise than \enquote{the radian is
\(57.2957\ldots{}^\circ\)}. It decomposes that publication into four operations
with a reference proof:
\begin{center}
\begin{tabularx}{.96\textwidth}{>{\bfseries}l X}
\toprule
\(50^\circ\) & second phase of the wheel; critical hinge on a declared chart;\\
\(A^\circ\) & common coordinate of the parangular publication;\\
\(-\frac{20}{81}\Delta_4^\circ\) & orientation-transport of global torsion;\\
\(+\epsR^\circ\) & real carry that mates the two lifts.\\
\bottomrule
\end{tabularx}
\end{center}

\paragraph{Convergence by depth of refinement}

At depth \(N\), the generator delivers rational cylinders
\(p_N\to\pih\) and \(e_N\to e\). Define
\[
S_{E,N}=2p_N\left(\frac5{36}+\frac{25}{9}\alphah\right),
\]
\[
\Delta_{4,N}=\frac{p_N-e_N\log p_N}{270}
\left(1+\frac{p_N}{729}\right),\qquad
S_{T,N}=1+\frac{20}{81}\Delta_{4,N},
\]
and
\[
\mathcal E_N=S_{T,N}-S_{E,N}.
\]
Continuity of the functional on \(3<p<4\), \(2<e<3\), together with the
monotone derivatives, transports each pair of cylinders to a certified
interval. At a depth of \(45\) blocks, its width is less than
\(4.81\times10^{-130}\).

Each nine-level return contains six trits per level. Therefore, the refinement rate is
\begin{equation}
729^{-9}=3^{-54}.
\label{eq:contraccion-54}
\end{equation}
This quantity controls the width of the intervals; it is not the value of
\(\epsR\).

The subtraction of triads stabilizes
\begin{center}
\ttfamily
000|000|000|896|802|822|044|562|981|999|050|695|020|651|286|413
\end{center}
and the carry prefix
\begin{center}
\ttfamily
0,0,0,0,0,-1,0,-1,-1,-1,0,-1,0.
\end{center}
This is the concrete mechanism of decimal memory: not a chosen tail, but
the recurrence \eqref{eq:subacarreo} applied to previously
constructed operands.

\paragraph{Evaluation of the canonical values}

The coinductive output is
\begin{align}
\epsone
&=8.9680282204456298199905069502065128641372736205009\ldots
\times10^{-10},\\
\epsR^\circ
&=5.1383016758575282071685164622645947489908574400054\ldots
\times10^{-8}\,{}^\circ.
\label{eq:epsilon-valores}
\end{align}
The complete publication results
\[
\frac{180^\circ}{\pih}
=57.2957795130823208767981548141051703324054724665643\ldots{}^\circ.
\]

\begin{exactbox}
\begin{align*}
50^\circ+A^\circ
&=57.29735256928380099728510547238066299956\ldots{}^\circ,\\
\frac{20}{81}\Delta_4^\circ
&=0.00157310758449687906223272996065728980465\ldots{}^\circ,\\
50^\circ+A^\circ-\frac{20}{81}\Delta_4^\circ
&=57.29577946169930411822287274242000570976\ldots{}^\circ,\\
\epsR^\circ
&=0.00000005138301675857528207168516462264594749\ldots{}^\circ,\\
\text{final sum}
&=57.29577951308232087679815481410517033241\ldots{}^\circ.
\end{align*}
\end{exactbox}

 
\section{Critical coordinate and oriented dodecaphasic temporal system: affine determination of the phase of \(50^\circ\)}
\subsection{Affine map \(j_{100}\) of the APP chart and uniqueness of the phase \(m=2\) in \(j_{100}(1/2)=\theta_m=50^\circ\)}
% Fragmento público generado desde fuentes teoremáticas íntegros.
% Residencia: 52.3.1.
% Fuente: /Users/ruben/Documents/HMT2/EL_CIERRE_HOLOGRAFICO_DEL_INFINITO_SUCESOR_102_CAPITULOS_2026-08-28_EN_TRABAJO/incoming/radion_monografia_20260827/manuscrito/sections/04_cincuenta_critica.tex:1-105
\noindent The critical coordinate links the Cayley chart, the APP record, and the dodecaphasic phase through typed maps.\par\medskip

\subsubsection{Two exact realizations of the angular coordinate \(50^\circ\)}

\paragraph{Second phase of the dodecaphasic temporal system}

By \eqref{eq:rueda-dodecafasica}, the sequence of centros begins
\[
20^\circ,\ 50^\circ,\ 80^\circ,\ 110^\circ,\ldots.
\]
Solve
\[
30m-10=50
\]
gives only \(m=2\). Therefore, \(50^\circ\) is not a rounding of \(A\),
a half-turn, or the centre of a circle: it is a typed phase of the
dodecaphase calendar. The fraction of a turn is
\[
\frac{50}{360}=\frac5{36},
\]
exactly the first summand of \(S_E/T_+\).

\paragraph{Cayley Critical Coordinate Matching}

The functional involution
\[
\rho\longmapsto1-\bar\rho
\]
fixed
\[
\operatorname{Fix}(\rho\mapsto1-\bar\rho)
=\{\rho:\Re\rho=1/2\}.
\]
The number \(1/2\) is the self-dual real coordinate of the functional interval \(0\leq\Re s\leq1\). In the longitudinal sense, it is not \enquote{half of a line}, and it does not contain the spectral height \(\gamma\). For
\[
\rho_\gamma=\frac12+i\gamma,
\]
The Cayley transformation can be written
\begin{equation}
\tau_\gamma
=\frac{\rho_\gamma-1}{\rho_\gamma}
=\frac{\gamma+i/2}{\gamma-i/2}
\in S^1,
\qquad
\gamma=\frac12\cot\frac{\theta_\gamma}{2}.
\label{eq:cayley-critica}
\end{equation}
Thus, \(1/2\) fixes the seam, \(\gamma\) determines the angular character, and
\(\Gamma_9\) preserves the depth of its iterations.

\paragraph{Affine injection \(j_{100}\) of the APP record into the angular chart}

Cubic completion arranges a chamber of ten marks per coordinate. One
face contains \(10^2=100\) positions. We declare the affine chart
\begin{equation}
j_{100}:[0,1]\longrightarrow[0,100^\circ],
\qquad j_{100}(\sigma)=100\sigma^\circ.
\label{eq:j100}
\end{equation}
Then
\[
j_{100}\left(\frac12\right)=50^\circ.
\]
Requiring \(j_{100}(1/2)=\theta_m\) simultaneously yields
\[
100\cdot\frac12=30m-10
\quad\Longleftrightarrow\quad m=2.
\]

\begin{theorem}[Critical Equalizer--Dodecaphasic]
In the hundred-mark APP chart, the fixed coordinate of the critical involution
is published exactly in the second phase of the wheel:
\begin{equation}
\boxed{j_{100}(1/2)=\theta_2=50^\circ.}
\label{eq:igualador-50}
\end{equation}
The phase \(m=2\) is unique.
\end{theorem}

The theorem is exact once \eqref{eq:j100} has been declared. Its status is
\status{new formalization}: the critical midpoint, the \(10^3\) completion, the hundred-mark face, and \(\theta_2\) had already been constructed; choosing that face as the affine critical chart brings their domains together for the first time.

\paragraph{Compatibility of the critical value with the Cayley chart}

The Cayley transform does not send \(1/2\) to \(50^\circ\). If
\(\gamma=0\), \eqref{eq:cayley-critica} gives
\[
\tau_0=-1=e^{i\pi},
\]
that is, \(180^\circ\). The \(50^\circ\) belong to the APP chart of one hundred
marks and to the dodecaphasic calendar; the Cayley angle depends on \(\gamma\).
This distinction prevents all the zeros of the zeta function from being collapsed into a single
phase.

\begin{falsifier}
The identification \(1/2\leftrightarrow50^\circ\) fails if the
chart \eqref{eq:j100} is removed. The scalar equality \(100(1/2)=50\) is not sufficient
by itself to identify two objects. The theorem depends on the distinguished
APP face and on the dodecaphase wheel.
\end{falsifier}

 \subsection{Compatibility of the critical phase with the nonadic return and memory advance}
% Fragmento público generado desde fuentes teoremáticas íntegros.
% Residencia: 52.3.2.
% Fuente: /Users/ruben/Documents/HMT2/EL_CIERRE_HOLOGRAFICO_DEL_INFINITO_SUCESOR_102_CAPITULOS_2026-08-28_EN_TRABAJO/incoming/radion_monografia_20260827/manuscrito/sections/04_cincuenta_critica.tex:106-171
\subsubsection{Typed separation of the 2 objects with value \(1/2\)}

APP contains, in addition, a singular mode with
\[
s=\frac12,\qquad s^2=\frac14,\qquad
1-s^2=\frac34=\frac{90}{120},\qquad
s\sqrt{1-s^2}=\frac{\sqrt3}{4}.
\]
This \(s\) arises from the principal angles of the sum and product sheets. The
critical seam \(\Re\rho=1/2\) arises from a functional involution. That both
coordinates equal \(1/2\) is a scalar coincidence with possible content,
but it does not authorize their identification without a typed square.

The correct way to preserve the finding is an equalizer diagram:
\[
\begin{tikzpicture}[baseline=(current bounding box.center),node distance=14mm and 23mm,
 every node/.style={font=\small}]
\node[draw=RadBlue,rounded corners,fill=RadMist] (crit) {costura crítica \(\Re s=1/2\)};
\node[draw=RadTeal,rounded corners,fill=RadCream,right=of crit] (face) {carta \(j_{100}\)};
\node[draw=RadCopper,rounded corners,fill=RadCream,right=of face] (theta) {fase \(\theta_2=50^\circ\)};
\node[draw=RadRose,rounded corners,fill=white,below=of face] (mode) {modo APP \(s=1/2\), \(1-s^2=90/120\)};
\draw[-{Latex},thick,RadBlue] (crit)--node[above]{\(j_{100}\)}(face);
\draw[-{Latex},thick,RadTeal] (face)--node[above]{igualador}(theta);
\draw[-{Latex},dashed,RadRose] (mode)--node[right,align=left]{misma cifra;\\dominio distinto}(face);
\end{tikzpicture}
\]
The diagram does not feign the missing arrow. It shows what has been proved, what has
been formalized, and what equality still requires additional transport if
a global operator identification is intended.

\subsubsection{Spectral Circle and Helical Memory Compatibility}

The Cayley--Pontryagin reader transports a spectral fiber via
\[
\widetilde U_{\Gamma_9^n}J\psi_\gamma
=\tau_\gamma^nJ\psi_\gamma.
\]
The circular publication is \(\tau_\gamma^n\in S^1\). The lift
\[
(\tau_\gamma^n,n)\in S^1\times\Z
\]
preserves the number of returns. On the contractive sheet appears
\[
\widetilde M_9^nJ\psi_\gamma
=\left(\frac59\right)^n\tau_\gamma^nJ\psi_\gamma.
\]
We therefore have three typed images:
\[
\text{phase circle},\qquad
\text{memory helix},\qquad
\text{contractive spiral}.
\]

A zero is not a revolution, and a nine is not a particular zero of the zeta
function. The zero fixes a character \(\tau_\gamma\) that persists through the
revolutions; nonadic return increments memory. Nine functions as residual zero
in \(\Z/9\Z\), an arithmetic operation distinct from the enumeration
of spectral zeros.

\begin{narrativebox}
The critical line becomes a circle when we forget the linear height and
retain the unitary character. It becomes a helix when we reinsert
the memory of how many times the character has been transported. Depth
is not a third decorative coordinate: it is the information that the circular
projection cannot retain.
\end{narrativebox}
 
\section{Angular pair and reduced action: construction of the positive ellipse and area per radion}
\subsection{Areal law of the radion: \(\mathcal A(\theta)=\hbar_{\mathrm{ret}}\theta\) and \(\mathcal A(2\pi)=h_{\mathrm{ret}}\)}
% Fragmento público generado desde fuentes teoremáticas íntegros.
% Residencia: 52.4.1.
% Fuente: /Users/ruben/Documents/HMT2/EL_CIERRE_HOLOGRAFICO_DEL_INFINITO_SUCESOR_102_CAPITULOS_2026-08-28_EN_TRABAJO/incoming/radion_monografia_20260827/manuscrito/sections/05_geometria_accion.tex:1-217
\noindent Spectral geometry determines the form, and the areal law determines the action transported by the phase.\par\medskip

\subsubsection{Equivalence between the spectral and geometric charts}

\paragraph{Spectral positive operator and conjugate eigenvalues}

The angular momentum produces the two-layer operator.
\begin{equation}
B_1=AI+C^*R,\qquad R^2=I,\qquad
P_\pm=\frac{I\pm R}{2}.
\label{eq:B1}
\end{equation}
Its eigenvalues are \(A\pm C^*\). Since \(0<C^*<A\), both are positive.
In the chart \(\kappa=\pi/180\), the spectral ellipse has semiaxes
\begin{equation}
a_1^2=\kappa(A+C^*),\qquad
b_1^2=\kappa(A-C^*).
\label{eq:semiejes-espectrales}
\end{equation}
The decomposition is reversible:
\begin{equation}
A=\frac{a_1^2+b_1^2}{2\kappa},\qquad
C^*=\frac{a_1^2-b_1^2}{2\kappa}.
\label{eq:reconstruccion-AC}
\end{equation}
The ellipse is not drawn to illustrate a number. It exactly reconstructs the pair that generated it.

Define
\begin{equation}
\rho_1=\sqrt{A^2-C^{*2}},\qquad
\eta=\operatorname{artanh}\frac{C^*}{A}
=\frac12\log\frac{A+C^*}{A-C^*}.
\label{eq:eta}
\end{equation}
Then
\[
A\pm C^*=\rho_1e^{\pm\eta},
\]
and
\begin{equation}
a_1=\sqrt{\kappa\rho_1}\,e^{\eta/2},
\qquad
b_1=\sqrt{\kappa\rho_1}\,e^{-\eta/2}.
\label{eq:semiejes-eta}
\end{equation}
The product preserves scale; the quotient preserves form:
\[
a_1b_1=\kappa\rho_1,
\qquad
\frac{b_1}{a_1}=e^{-\eta}
=\sqrt{\frac{A-C^*}{A+C^*}}.
\]

\paragraph{Focal Invariants of the Positive Operator}

The semifocal distance is
\begin{equation}
c_1^2=a_1^2-b_1^2=2\kappa C^*,
\qquad
F_\pm=(\pm c_1,0).
\label{eq:focos}
\end{equation}
The total distance is equal to
\begin{equation}
d_1=2c_1=2\sqrt{2\kappa C^*}
=0.544545509325626623406299779866023\ldots.
\label{eq:d1}
\end{equation}
The isolated number is not the thesis. The identity \(c_1^2=2\kappa C^*\) states that
\(C^*\) is literally the quadratic focal splitting in the
distinguished chart. The foci are the conjugate endpoints of the spectral aperture;
they are not, by themselves, the electric and magnetic poles. That reading requires
the constitutive transducer.

The eccentricity is
\begin{equation}
e_f=\frac{c_1}{a_1},\qquad
e_f^2=1-e^{-2\eta}=\frac{2C^*}{A+C^*}.
\label{eq:excentricidad}
\end{equation}
Numerically,
\[
\begin{aligned}
\eta&=0.299689683921630799684926562863927\ldots,\\
e_f&=0.671451895550191986916277055864332\ldots.
\end{aligned}
\]

\paragraph{Reference APP Block and Normalized Transport}

The central block
\[
B_0=7I+2R=9P_++5P_-
\]
possesses
\[
\rho_0=\sqrt{45},\qquad
\eta_0=\frac12\log\frac95,\qquad
e_0=\frac23.
\]
Its focal length is
\[
d_0=4\sqrt\kappa
=0.528443639680801468446003835349358\ldots.
\]
The exact comparison is
\begin{equation}
\boxed{\left(\frac{d_1}{d_0}\right)^2=\frac{C^*}{2}},
\qquad
d_1=d_0\sqrt{\frac{C^*}{2}}.
\label{eq:invariante-focal-relativo}
\end{equation}
This is the defensible genealogy of \(0.544545\ldots\): a focal APP scale
of reference multiplied by a relative invariant. Its visual proximity to
\(0.54\) is not an identity with the APP defect \(54\).

Transport from \(B_0\) separates into scale and anisotropy:
\begin{equation}
T_{\mathrm{rel}}=\frac{\rho_1}{\sqrt{45}}
\exp\bigl((\eta-\eta_0)R\bigr).
\label{eq:transporte-rel}
\end{equation}
On the leaves,
\[
t_+=\frac{A+C^*}{9}=1.0467878786947163\ldots,
\qquad
t_-=\frac{A-C^*}{5}=1.0347228460630311\ldots.
\]
The deformation is not a single decimal. Its isotropic component is \(\sqrt{t_+t_-}\); its hyperbolic component is \(\tfrac12\log(t_+/t_-)\).

\subsubsection{Positive ellipse associated with the angular pair}

\paragraph{Construction via Conjugate Semi-axes}

The preceding spectral scale has no units of action. Dimensional
Mechanics publishes the same form on the internal action line:
\begin{equation}
a_S=\sqrt{2\hret e^\eta},
\qquad
b_S=\sqrt{2\hret e^{-\eta}}.
\label{eq:semiejes-accion}
\end{equation}
Immediately,
\begin{equation}
a_Sb_S=2\hret,\qquad
\frac{a_S}{b_S}=e^\eta.
\label{eq:producto-forma}
\end{equation}
The product fixes the action; the quotient fixes the anisotropy. The ellipse is
parametrized by
\begin{equation}
Q(\theta)=a_S\cos\theta,\qquad
P(\theta)=b_S\sin\theta.
\label{eq:elipse-param}
\end{equation}

The spectral and action ellipses are homothetic:
\begin{equation}
\lambda_{\mathrm{act}}^2=\frac{2\hret}{\kappa\rho_1},
\qquad
(a_S,b_S,c_S)=\lambda_{\mathrm{act}}(a_1,b_1,c_1).
\label{eq:homotecia}
\end{equation}
They share the form \(e^{-\eta}\), not the physical dimension.

\paragraph{Areal law of action per radion}

The forma simplectica canonica is \(\omega=\diff Q\wedge\diff P\). The area
oriented barrida from \(0\) until \(\theta\) is
\begin{align}
\mathcal A(\theta)
&=\frac12\int_0^\theta
\bigl(Q(t)P'(t)-P(t)Q'(t)\bigr)\,\diff t\\
&=\frac12a_Sb_S\int_0^\theta
(\cos^2t+\sin^2t)\,\diff t.
\end{align}
Using \eqref{eq:producto-forma}, we obtain the central result.

\begin{theorem}[Action per radion]
For the ellipse \eqref{eq:elipse-param},
\begin{equation}
\boxed{\mathcal A(\theta)=\hret\theta.}
\label{eq:area-theorem}
\end{equation}
In particular,
\begin{equation}
\boxed{\mathcal A(1)=\hret},
\qquad
\boxed{\mathcal A(2\pi)=2\pi\hret=\hfull}.
\label{eq:area-radion-vuelta}
\end{equation}
\end{theorem}

This equality allows the definition of the radius without first appealing to the arc length.

\begin{definition}[Oriented radion]
\begin{equation}
\mathfrak r_\sigma=(\theta=1,\sigma),
\qquad \sigma\in\{+1,-1\},
\qquad
\mathcal A(\mathfrak r_\sigma)=\sigma\hret.
\label{eq:radion-oriented}
\end{equation}
\end{definition}

The radian is the projection that forgets \(\sigma\), the sheet, the
memory, and the area. The radion preserves that information. Hence the name
\emph{rad--ion}: orientation precedes
physical charge, and charge appears later as
\[
Q(x)=n_Q(x)e_{\mathrm{int}},\qquad
e_{\mathrm{int}}=\sqrt{\frac{2\alphah\hfull}{Z_{\mathrm{int}}}}.
\]

 \subsection{Conjugate semi-axes, focal invariants, and orientation of anisotropy}
% Fragmento público generado desde fuentes teoremáticas íntegros.
% Residencia: 52.4.2.
% Fuente: /Users/ruben/Documents/HMT2/EL_CIERRE_HOLOGRAFICO_DEL_INFINITO_SUCESOR_102_CAPITULOS_2026-08-28_EN_TRABAJO/incoming/radion_monografia_20260827/manuscrito/sections/05_geometria_accion.tex:218-290
\paragraph{Correspondence between phase, action, area, and orientation}

\begin{figure}[H]
\centering
\begin{tikzpicture}[scale=1.12]
  \def\aa{4.0}
  \def\bb{2.95}
  \def\cc{2.70}
  \draw[->,RadGray] (-4.7,0)--(4.7,0) node[right]{\(Q\)};
  \draw[->,RadGray] (0,-3.5)--(0,3.5) node[above]{\(P\)};
  \fill[RadTeal!8] (0,0) ellipse[x radius=\aa,y radius=\bb];
  \draw[RadTeal,line width=1.35pt] (0,0) ellipse[x radius=\aa,y radius=\bb];
  \fill[RadCopper] (-\cc,0) circle (2pt) node[below=2mm]{\(F_-\)};
  \fill[RadCopper] (\cc,0) circle (2pt) node[below=2mm]{\(F_+\)};
  \draw[RadGold,line width=1.1pt,-{Latex}] (\aa,0)
     arc[start angle=0,end angle=57.3,x radius=\aa,y radius=\bb];
  \node[RadGold] at (3.7,2.3) {a radion};
  \draw[dashed,RadBlue] (0,0)--(\aa,0) node[midway,below]{\(a_S\)};
  \draw[dashed,RadBlue] (0,0)--(0,\bb) node[midway,left]{\(b_S\)};
  \node[align=center,RadNavy] at (0,-3.25)
    {total area \(\pi a_Sb_S=2\pi\hret=\hfull\)};
\end{tikzpicture}
\caption{The action ellipse. A parametric arc of one radian sweeps
\(\hret\); one turn encloses \(\hfull\). The foci encode the aperture
\(C^*\) in the distinguished chart.}
\label{fig:elipse-accion}
\end{figure}

\paragraph{Relationship between the reduced action, uncertainty, and area}

The relation \(\hbar=h/(2\pi)\) is often presented as a notational
convenience: when describing oscillations, rotations, and Fourier transforms,
the factors \(2\pi\) disappear if \(\hbar\) is used. The action ellipse
reveals a stronger reading. Division by \(2\pi\) is exactly the
passage from the action of one revolution to the action per unit phase:
\[
h=\mathcal A(2\pi),\qquad \hbar=\mathcal A(1).
\]
The reduced constant is not merely \(h\) in more convenient notation; it is the
areal density of action with respect to phase.

Heisenberg, Dirac, and Jordan established the language of conjugate variables,
commutators, and canonical transformations. No literal
\enquote{ontological ellipse} lacking a source is attributed to them here. The HMT--MD contribution is
to identify the positive figure, construct its semiaxes from \((A,C^*)\), and
prove \eqref{eq:area-theorem}. The historical genealogy opens the language;
the present theorem fixes the object.

\begin{narrativebox}
Action is not appended to motion afterward. In the radion ellipse, action is the area generated by motion. A full turn does not transport a prefabricated constant: it encloses it.
\end{narrativebox}

\Needspace{10\baselineskip}
\paragraph{Algebraic controls of the focal invariants}

Decimal similarities are subject to control:
\begin{center}
\begin{tabularx}{.96\textwidth}{l r X}
\toprule
Comparison & Residue & Verdict\\
\midrule
\(d_1\) versus \(0.54\) & \(4.5455\times10^{-3}\) & not equality; \(54\) is not deduced;\\
\(d_1^2\) versus \(8/27\) & \(2.3352\times10^{-4}\) & proximity without map;\\
\(\eta\) frente a \(3/10\) & \(-3.1032\times10^{-4}\) & no igualdad;\\
\(e_f\) versus \(2/3\) & \(4.7852\times10^{-3}\) & \(2/3\) belongs exactly to \(B_0\);\\
\(e^{-\eta}\) versus \(20/27\) & \(3.0740\times10^{-4}\) & no igualdad.\\
\bottomrule
\end{tabularx}
\end{center}
The robust invariants are \eqref{eq:eta}, \eqref{eq:excentricidad}, and
\eqref{eq:invariante-focal-relativo}. The HMT criterion does not consist in favouring
a decimal resemblance, but in requiring an object, an operation, and conservation.
 
\section{Nonadic Return and Internal Projective: Klein Disk and Helical Depth of Memory}
\subsection{Normalization to the Hilbert--Beltrami--Klein disk and hyperbolic interior}
% Fragmento público generado desde fuentes teoremáticas íntegros.
% Residencia: 52.5.1.
% Fuente: /Users/ruben/Documents/HMT2/EL_CIERRE_HOLOGRAFICO_DEL_INFINITO_SUCESOR_102_CAPITULOS_2026-08-28_EN_TRABAJO/incoming/radion_monografia_20260827/manuscrito/sections/06_interior_helice.tex:1-108
\noindent Projective geometry distinguishes the finite action from hyperbolic depth and return memory.\par\medskip

\subsubsection{Hyperbolic Interior of Hilbert--Beltrami--Klein}

\paragraph{Projective normalization on the Klein disk}

The affine application
\[
u=\frac{Q}{a_S},\qquad v=\frac{P}{b_S}
\]
carries the action ellipse to the unit disk
\[
D=\{(u,v):u^2+v^2<1\}.
\]
On each chord, the Hilbert distance is defined through the cross-ratio
of its boundary endpoints. The resulting infinitesimal metric is the
Beltrami--Klein form
\begin{equation}
\diff s_K^2=
\frac{(1-r^2)(\diff u^2+\diff v^2)+(u\diff u+v\diff v)^2}
{(1-r^2)^2},
\qquad r^2=u^2+v^2.
\label{eq:klein-metric}
\end{equation}
With this normalization, the Gaussian curvature is constant:
\begin{equation}
K=-1.
\label{eq:klein-curvature}
\end{equation}
The geodesics are Euclidean chords, but their hyperbolic length diverges as
they approach the boundary.

\paragraph{Finiteness of the boundary and infinitude of the hyperbolic depth}

The metric determinant of \eqref{eq:klein-metric} produce
\begin{equation}
\diff A_K=\frac{\diff u\,\diff v}{(1-r^2)^{3/2}}.
\label{eq:klein-area-density}
\end{equation}
The area of the subdisk \(r<R<1\) is
\begin{align}
A_K(R)
&=\int_0^{2\pi}\int_0^R
\frac{r\,\diff r\,\diff\theta}{(1-r^2)^{3/2}}\\
&=2\pi\left(\frac1{\sqrt{1-R^2}}-1\right).
\label{eq:klein-area-R}
\end{align}
\Needspace{4\baselineskip}
Therefore,
\[
\lim_{R\to1^-}A_K(R)=+\infty.
\]
Simultaneously, the same boundary encloses the symplectic area.
\[
\pi a_Sb_S=2\pi\hret=\hfull<\infty.
\]

\begin{theorem}[Finiteness of Action and Infinitude of Depth]
The radion ellipse has a finite turn action \(\hfull\), whereas
its interior, equipped with the Hilbert--Beltrami--Klein metric, has infinite
hyperbolic area:
\begin{equation}
\boxed{
\operatorname{Area}_{\mathrm{simp}}(\partial\mathbb E_S)=\hfull<\infty,
\qquad
\operatorname{Area}_{K}(\mathbb E_S)=+\infty.}
\label{eq:finite-infinite}
\end{equation}
\end{theorem}

There is no contradiction: \(\operatorname{Area}_{\mathrm{simp}}\) and
\(\operatorname{Area}_K\) are different measures. The first counts oriented action
in phase; the second measures interior projective depth. The thesis
\enquote{finite boundary, infinite interior} is now an equality and a limit,
not a metaphor.

\paragraph{Focal Coordinates Inside Klein}

In the normalized disk, the foci are located at \((\pm e_f,0)\). The distance from the center to a focus is
\begin{equation}
u_f=\operatorname{artanh}e_f
=\operatorname{arcosh}(e^\eta)
=0.813382331175342333760889731088228\ldots.
\label{eq:klein-focal}
\end{equation}
The focal--focal distance is
\[
d_K(F_-,F_+)=2u_f
=1.62676466235068466752177946217646\ldots.
\]
We have
\begin{equation}
e_f=\tanh u_f,\qquad
e^{-\eta}=\operatorname{sech}u_f,\qquad
\eta=\log\cosh u_f.
\label{eq:focal-hyperbolic-chain}
\end{equation}
The region up to the focal radius has area
\[
A_K(e_f)=2\pi(e^\eta-1)
=2.19559620884061869541206115980360\ldots.
\]

All open ellipses are projectively isometric as Hilbert domains. Curvature
\(-1\) alone therefore does not retain eccentricity. The HMT form
resides in the joint structure: the distinguished affine chart, symplectic
form, sheet labeling, and boundary parametrization.
 \subsection{Difference between phase return and return of the integral holonomic state}
% Fragmento público generado desde fuentes teoremáticas íntegros.
% Residencia: 52.5.2.
% Fuente: /Users/ruben/Documents/HMT2/EL_CIERRE_HOLOGRAFICO_DEL_INFINITO_SUCESOR_102_CAPITULOS_2026-08-28_EN_TRABAJO/incoming/radion_monografia_20260827/manuscrito/sections/06_interior_helice.tex:109-201
\subsubsection{Helical elevation of the nonadic return}

\paragraph{Nonadic decomposition and memory quotient}

The elementary operation
\begin{equation}
n=9q_9(n)+\rho_9(n),
\qquad
H_9(n)=(\rho_9(n),q_9(n))
\label{eq:H9}
\end{equation}
satisfies
\begin{equation}
H_9(n+9)=H_9(n)+(0,1).
\label{eq:H9-helix}
\end{equation}
The first coordinate closes; the second ascends. This identity is the
minimal arithmetic version of a screw.

For a spectral fiber,
\[
n\longmapsto(\tau_\gamma^n,n)
\]
is a helix over \(S^1\). Adding the contraction \((5/9)^n\) produces an
inner spiral. The radion participates in both readings: a unit of phase
on the boundary and an oriented unit of action in the lift.

\paragraph{Memory range as dimensional reader}

Suppose that a visible projection preserves the phase value but loses \(d\) linearly independent memory cocycles. Any faithful lift requires at least \(d\) additional coordinates to reconstruct them. This motivates a precise reading of Dimensional Mechanics:
\begin{equation}
\begin{aligned}
\text{emergent dimension}
&=\text{minimum rank of independent memory}\\
&\phantom{=}\text{not publishable on the visible sheet}.
\end{aligned}
\label{eq:dimension-memory}
\end{equation}
It is not asserted that every physical dimension is a counter. Rather, in the
TPK domain, every independent memory that must survive forces a
coordinate of the lift.

\begin{narrativebox}
Dimension appears when the projection can no longer preserve the entire
history. The circle suffices to state where the phase is; it does not
suffice to state how many times it returned, which sheet it transported,
which interval contracted, or which carry survived.
\end{narrativebox}

\paragraph{Cylinder Assembly, Spirals, and Helical Trajectories}

The coinductive cylinders are not physical tubes added to a spiral. They are
nested intervals that preserve prefixes and narrow the value. The spiral
describes transverse contraction; the helix describes longitudinal memory;
the cylinder describes the set of compatible values at a finite
depth. In the limit, the intersection of the cylinders selects a unique
section.

The spatial image brings together, without confusion:
\[
\begin{array}{ccl}
\text{angle} &\leftrightarrow& \text{visible phase},\\
\text{height} &\leftrightarrow& \text{return/memory},\\
\text{cylinder radius} &\leftrightarrow& \text{refinement uncertainty},\\
\text{sheet} &\leftrightarrow& \text{bidirectional orientation}.
\end{array}
\]
The rate \(3^{-54}\) contracts the cylinder by return; it is not added to the angle and
does not replace the actual carry \(\epsR\).

\subsubsection{Phase closures of orders \(2\pi\) and \(4\pi\)}

In the circular projection, \(2\pi\) restores the phase. In an oriented lift
with two sheets, one turn may close on the conjugate sheet and require a
second turn to restore the orientation:
\[
U_{2\pi}=-I,\qquad U_{4\pi}=I.
\]
The structure agrees with spinorial periodicity when the
functor is applied to a representation \(SU(2)\); the TPK lift is not identified with
\(SU(2)\) without that map.

The areal reading is consistent:
\[
\mathcal A(2\pi)=\hfull,\qquad
\mathcal A(4\pi)=2\hfull.
\]
On a Möbius strip, one visible circuit reverses the sheet and two circuits restore orientation. The \(2\pi/4\pi\) is not an accidental doubling; it is the difference between phase return and return of the integral holonomic state.
 
\section{Oriented ellipse and conjugate variables: minimum uncertainty and electromagnetic oscillation}
\subsection{Conjugate symplectic variables and determination of the minimum uncertainty}
% Fragmento público generado desde fuentes teoremáticas íntegros.
% Residencia: 52.6.1.
% Fuente: /Users/ruben/Documents/HMT2/EL_CIERRE_HOLOGRAFICO_DEL_INFINITO_SUCESOR_102_CAPITULOS_2026-08-28_EN_TRABAJO/incoming/radion_monografia_20260827/manuscrito/sections/07_incertidumbre_oscilador.tex:1-111
\noindent The symplectic form coordinates action, uncertainty, and oscillatory dynamics without altering the genealogy of the state.\par\medskip

\subsubsection{Realization of the state in a Hilbert space}

\paragraph{Minimum Covariance Matrix}

Once the action ellipse has been constructed, it can be performed on a pair of
conjugate operators
\[
[\widehat Q,\widehat P]=i\hret.
\]
We define
\begin{equation}
V_\eta=\frac{\hret}{2}
\begin{pmatrix}
e^\eta&0\\[2pt]0&e^{-\eta}
\end{pmatrix}.
\label{eq:covariance}
\end{equation}
Its determinant is
\begin{equation}
\det V_\eta=\frac{\hret^2}{4}.
\label{eq:cov-det}
\end{equation}
Consequently,
\begin{equation}
(\Delta Q)^2=\frac{\hret}{2}e^\eta,
\qquad
(\Delta P)^2=\frac{\hret}{2}e^{-\eta},
\qquad
\Delta Q\,\Delta P=\frac{\hret}{2}.
\label{eq:uncertainty}
\end{equation}
The Robertson--Schrödinger inequality is saturated.

The compressed annihilator
\begin{equation}
\widehat a_\eta=
\frac{e^{-\eta/2}\widehat Q+i e^{\eta/2}\widehat P}
{\sqrt{2\hret}}
\label{eq:a-eta}
\end{equation}
satisfies \([\widehat a_\eta,\widehat a_\eta^\dagger]=1\). Its vacuum has
wavefunction
\begin{equation}
\psi_\eta(Q)=
(\pi\hret e^\eta)^{-1/4}
\exp\left(-\frac{Q^2}{2\hret e^\eta}\right).
\label{eq:gaussian}
\end{equation}

\paragraph{Ellipse as level set of the Mahalanobis functional}

The set
\begin{equation}
\mathbb E_{\mathrm{act}}
=\{x\in\R^2:x^{\mathsf T}V_\eta^{-1}x\le4\}
\label{eq:mahalanobis}
\end{equation}
has semiaxes
\begin{equation}
a_S=2\Delta Q,\qquad b_S=2\Delta P.
\label{eq:two-sigma}
\end{equation}
Its area is
\[
4\pi\sqrt{\det V_\eta}=2\pi\hret=\hfull.
\]
The previously constructed HMT ellipse coincides exactly with the contour of two
deviations from the minimal Gaussian state.

\begin{proposition}[Unique chain of the quotient form]
After declaring the spectral, symplectic, and Hilbert interfaces,
\begin{equation}
\boxed{
e^{-\eta}
=\sqrt{\frac{A-C^*}{A+C^*}}
=\frac{b_1}{a_1}
=\frac{b_S}{a_S}
=\frac{\Delta P}{\Delta Q}.}
\label{eq:shape-chain-hilbert}
\end{equation}
\end{proposition}

The action ellipse corresponds to the classical level \(H=\hret\omega\) and to the
contour \(2\sigma\) of the compressed vacuum. It must not be conflated with the mean
vacuum energy, \(\hret\omega/2\). This factor-of-two difference is a
mandatory physical control.

\paragraph{Invariance Conditions of the Ellipse as an Operational Object}

In the conventional formulation, an uncertainty ellipse may be viewed as
a graphical device for representing a covariance matrix. Here the causal
order is reversed: the positive ellipse already exists as a publication of
\((A,C^*,\hret)\); Hilbert space recognizes it through
\eqref{eq:covariance}. The claimed ontological reality does not mean that a
classical particle physically traverses an oval in the plane \((Q,P)\). It means
that the state preserves a positive form, a total action, and an orientation
quotient before a probabilistic theory interprets them.

The wave function does not create the form; it represents it. Uncertainty does not create
\(\hbar\); it publishes the impossibility of simultaneously compressing the two
directions without altering the symplectic product.

\begin{narrativebox}
The probabilistic reading is one possible realization of the ellipse, not its cause.
The primary object is the joint preservation of product and orientation.
Probability appears when that geometry is published as the covariance of a
state on a Hilbert space.
\end{narrativebox}

 \subsection{Exchange \(\mu\leftrightarrow\varepsilon\), inversion \(Z\leftrightarrow Z^{-1}\), and conservation of causal velocity}
% Fragmento público generado desde fuentes teoremáticas íntegros.
% Residencia: 52.6.2.
% Fuente: /Users/ruben/Documents/HMT2/EL_CIERRE_HOLOGRAFICO_DEL_INFINITO_SUCESOR_102_CAPITULOS_2026-08-28_EN_TRABAJO/incoming/radion_monografia_20260827/manuscrito/sections/07_incertidumbre_oscilador.tex:156-197
\paragraph{Energy exchange between the electric and magnetic sectors}

The two terms of \eqref{eq:HLC} become
\begin{equation}
U_E(\theta)=\hret\omega\cos^2\theta,
\qquad
U_B(\theta)=\hret\omega\sin^2\theta.
\label{eq:energy-exchange}
\end{equation}
Therefore,
\[
U_E+U_B=\hret\omega.
\]
Electric--magnetic polarity does not consist in arbitrarily labeling
the foci. It consists of two conjugate reserves that are exchanged during the
orbit while keeping the frequency and total action constant.

The common quotient extends:
\begin{equation}
\boxed{
e^{-\eta}
=\frac{b_S}{a_S}
=\frac{\Delta P}{\Delta Q}
=\frac{Z_\eta}{Z_*}
=\sqrt{\frac{A-C^*}{A+C^*}}.}
\label{eq:shape-chain-LC}
\end{equation}
The conserved products are
\[
a_Sb_S=2\hret,\qquad
\det V_\eta=\frac{\hret^2}{4},
\qquad
L_\eta C_\eta=\omega^{-2}.
\]

\begin{exactbox}
The LC oscillator is the minimal physical realization of the double reading: the
product preserves action and frequency; the quotient preserves orientation and
impedance. A complete cycle exchanges electricity and magnetism without
losing the area \(\oint P\,\diff Q=\hfull\).
\end{exactbox}

 \subsection{Inductive--capacitive pair \((L_\eta,C_\eta)\): invariance of \(L_\eta C_\eta=\omega^{-2}\) and orientation of impedance}
% Fragmento público generado desde fuentes teoremáticas íntegros.
% Residencia: 52.6.3.
% Fuente: /Users/ruben/Documents/HMT2/EL_CIERRE_HOLOGRAFICO_DEL_INFINITO_SUCESOR_102_CAPITULOS_2026-08-28_EN_TRABAJO/incoming/radion_monografia_20260827/manuscrito/sections/07_incertidumbre_oscilador.tex:112-155
\subsubsection{Inductive--capacitive oscillator associated with the angular pair}

\paragraph{Independent construction of the 2 conjugated reservoirs}

Given \(\omega>0\) and a reference impedance \(Z_*>0\), we define
\begin{equation}
L_\eta=\frac{Z_*}{\omega}e^{-\eta},
\qquad
C_\eta=\frac{1}{Z_*\omega}e^\eta.
\label{eq:LC-def}
\end{equation}
Then
\begin{equation}
L_\eta C_\eta=\omega^{-2},
\qquad
\sqrt{\frac{L_\eta}{C_\eta}}=Z_*e^{-\eta}.
\label{eq:LC-invariants}
\end{equation}
The frequency remains fixed; the impedance preserves the orientation of the form.

The Hamiltonian is
\begin{equation}
H_{LC}=\frac{q^2}{2C_\eta}+\frac{\Phi^2}{2L_\eta}.
\label{eq:HLC}
\end{equation}
With the normalized variables
\begin{equation}
Q=\sqrt{Z_*}\,q,\qquad
P=\frac{\Phi}{\sqrt{Z_*}},
\label{eq:LC-map}
\end{equation}
we obtain
\begin{equation}
H_{LC}=\frac\omega2
\left(e^{-\eta}Q^2+e^\eta P^2\right).
\label{eq:HLC-QP}
\end{equation}
On \eqref{eq:elipse-param},
\begin{equation}
H_{LC}=\hret\omega.
\label{eq:LC-level}
\end{equation}

% Fuente: /Users/ruben/Documents/HMT2/EL_CIERRE_HOLOGRAFICO_DEL_INFINITO_SUCESOR_102_CAPITULOS_2026-08-28_EN_TRABAJO/incoming/radion_monografia_20260827/manuscrito/sections/07_incertidumbre_oscilador.tex:198-239
\paragraph{LC Oscillator as an Elementary Dynamic Unit}

The harmonic oscillator pervades nearly all of physics because it linearizes the neighborhood of an equilibrium, decomposes fields into modes, organizes quanta, and allows propagation to be read as a superposition of phases. In the radion, this universality receives a concrete genealogy: the mode does not begin in an external differential equation, but in a positive form with a fixed return action. Conventional dynamics appears upon choosing a clock \(\omega\) and a transducer \(Z_*\).

Frequency does not alter the action ellipse: it changes the level energy
through \(E=\hret\omega\). This separation is crucial. Geometry fixes
action; the clock converts action into energy.

\subsubsection{Time reversal and conjugate propagations}

Let \(H=H^*\) be given, and let \(J_E\) be a complex structure with \(J_E^2=-I\). Suppose there is an involution \(C\) such that
\[
CJ_EC^{-1}=-J_E,\qquad CHC^{-1}=H.
\]
The evolution
\begin{equation}
U(t)=\exp\left(-\frac{J_EHt}{\hret}\right)
\label{eq:evolution-J}
\end{equation}
satisfies exactly
\begin{equation}
CU(t)C^{-1}=U(-t).
\label{eq:time-reversal}
\end{equation}
This is the minimal operator form of bidirectionality: conjugation
interchanges the temporally opposed prolongations.

Wheeler--Feynman absorber theory and the reading of antiparticles
as temporally conjugate trajectories offer a subsequent historical realization.
Identity \eqref{eq:time-reversal} is demonstrated; converting it
into a complete theory of advanced and retarded propagators requires declaring
action, boundary conditions, causality, and absorber interactions.
The monograph preserves the physical intuition without converting a correspondence
into an equality of theories.

 
\section{The three TRIT regimes: electric--magnetic and elastic--rotational polarities}
\subsection{Coordination of the elliptic, parabolic, and hyperbolic realizations}
% Fragmento público generado desde fuentes teoremáticas íntegros.
% Residencia: 52.7.1.
% Fuente: /Users/ruben/Documents/HMT2/EL_CIERRE_HOLOGRAFICO_DEL_INFINITO_SUCESOR_102_CAPITULOS_2026-08-28_EN_TRABAJO/incoming/radion_monografia_20260827/manuscrito/sections/08_complejo_electromagnetismo.tex:1-58
\noindent The 3 TRIT regimes admit complex, electromagnetic, and mechanical realizations with distinct types.\par\medskip

\subsubsection{Operator quadratic of root internal of \(-1\)}

\paragraph{Exponential representation of the elliptic regime \(J^2=-I\)}

With \(J^2=-I\),
\begin{equation}
T_{\mathrm{ell}}(x,y)=e^{-xI-yJ}
=e^{-x}(\cos y\,I-\sin y\,J).
\label{eq:elliptic-exp}
\end{equation}
Under the identification generated by \(J\),
\[
T_{\mathrm{ell}}(x,y)\longleftrightarrow e^{-x-iy}.
\]
The scalar \(x\) controls dilation; \(y\) controls oriented
rotation. The form \(a+bi\) is recognized as elasticity plus rotation,
not as two arbitrarily adjoined quantities.

\paragraph{Exponential representation of the parabolic regime \(J^2=0\)}

With \(N^2=0\), the exponential is truncated:
\begin{equation}
T_{\mathrm{par}}(x,y)=e^{-x}(I-yN).
\label{eq:parabolic-exp}
\end{equation}
The regime describes the threshold at which there is neither periodic rotation nor exponential opening. In the active temporal reading it is the Euclidean present, but it retains memory of the branches that bound it.

\paragraph{Exponential Representation of the Hyperbolic Regime \(J^2=+I\)}

With \(R^2=I\) and \(P_\pm=(I\pm R)/2\),
\begin{equation}
T_{\mathrm{hyp}}(x,y)=e^{-xI-yR}
=q_+P_++q_-P_-,
\qquad q_\pm=e^{-x\mp y}.
\label{eq:hyperbolic-exp}
\end{equation}
The elementary invariants are
\begin{equation}
q_+q_-=e^{-2x},
\qquad
\frac{q_-}{q_+}=e^{2y}.
\label{eq:q-product-ratio}
\end{equation}
The product preserves the common mode; the quotient preserves the oriented separation.

Applied to the angular pair,
\[
x=\frac{\pi A}{180},\qquad y=\frac{\pi C^*}{180}.
\]
Therefore, \(A\) is the common coordinate of the publication and \(C^*\) its oriented contrast. Only after a physical reader has been applied are they recognized as electric and magnetic poles.

 \subsection{Typed separation of the constitutive and mechanical polarities}
% Fragmento público generado desde fuentes teoremáticas íntegros.
% Residencia: 52.7.2.
% Fuente: /Users/ruben/Documents/HMT2/EL_CIERRE_HOLOGRAFICO_DEL_INFINITO_SUCESOR_102_CAPITULOS_2026-08-28_EN_TRABAJO/incoming/radion_monografia_20260827/manuscrito/sections/08_complejo_electromagnetismo.tex:59-158
\subsubsection{Constitutive Operator on the Channels \(90/120\)}

\paragraph{Common and Oriented Component Decomposition}

We define
\begin{align}
\mathcal E
&=\log\bigl[(1-q_+^{90})(1-q_-^{90})\bigr]
-\log\bigl[(1-q_+^{120})(1-q_-^{120})\bigr],
\label{eq:Ecommon}\\
\mathcal B
&=\log\frac{1-q_-^{90}}{1-q_+^{90}}
-\log\frac{1-q_-^{120}}{1-q_+^{120}}.
\label{eq:Borient}
\end{align}
Under \(C^*\mapsto-C^*\), \(q_+\leftrightarrow q_-\) are interchanged. Therefore,
\(\mathcal E\) is even and \(\mathcal B\) is odd.

The normalized constitutive equations are
\begin{equation}
\widehat\mu=e^{\mathcal E+\mathcal B},
\qquad
\widehat\varepsilon=e^{\mathcal E-\mathcal B},
\qquad
\widehat Z=e^{\mathcal B},
\qquad
\widehat c=e^{-\mathcal E}.
\label{eq:constitutives}
\end{equation}
They are verified exactly.
\begin{equation}
\widehat Z=\sqrt{\frac{\widehat\mu}{\widehat\varepsilon}},
\qquad
\widehat c=(\widehat\mu\widehat\varepsilon)^{-1/2}.
\label{eq:constitutive-identities}
\end{equation}
Conjugation produces
\begin{equation}
C^*\mapsto-C^*:\qquad
\widehat\mu\leftrightarrow\widehat\varepsilon,
\quad
\widehat Z\leftrightarrow\widehat Z^{-1},
\quad
\widehat c\mapsto\widehat c.
\label{eq:EM-conjugation}
\end{equation}

\begin{theorem}[Constitutive polarity]
The reader \(90/120\) converts the common mode and oriented contrast of the
pair \((A,C^*)\) into two conjugate constitutive quantities. Sheet inversion
exchanges permeability and permittivity, reverses impedance, and preserves normalized
velocity.
\end{theorem}

This is the precise formulation of the intuition \enquote{electrical part and
magnetic part}. One writes neither \(A=E\) nor \(C^*=B\). Instead, one writes the map
\[
(A,C^*)\to(x,y)\to(q_+,q_-)
\to(\mathcal E,\mathcal B)
\to(\widehat\varepsilon,\widehat\mu,\widehat Z,\widehat c).
\]

\paragraph{Elastic-rotational realization of the constitutive polarity}

The same information may be published in the complex regime by
\begin{equation}
Z_{\mathrm{em}}=e^{\mathcal E/2}
\left(\cos\frac{\mathcal B}{2}\,I
-\sin\frac{\mathcal B}{2}\,J\right).
\label{eq:Zem}
\end{equation}
The factor \(e^{\mathcal E/2}\) is an elastic dilation; the factor generated by \(J\) is a rotation. Complex notation does not add a mysterious imaginary dimension to a previously given real world: it records two operational modes already separated by the TRIT.

\paragraph{Transduction of Orientation into Effective Load}

The radion sheet \(\sigma\) fixes an orientation before a
unit of charge is introduced. Dimensional Mechanics then provides the transducer
\[
e_{\mathrm{int}}=\sqrt{\frac{2\alphah\hfull}{Z_{\mathrm{int}}}},
\qquad
Q=n_Qe_{\mathrm{int}}.
\]
The causal chain is
\[
\text{sheet}\to\text{orientation}\to\text{constitutive}
\to\text{dimensional transducer}\to\text{physical charge}.
\]
Calling the unit \emph{radion} does not prematurely identify the orientation
with the coulomb; it makes visible the locus where the charge sign may reside.

\begin{narrativebox}
Electricity and magnetism are not affixed to the circle from outside. They emerge
when the same state is read through product and quotient: the product
retains the common mode, while the quotient retains the orientation. The LC oscillator
realizes the exchange; the reader \(90/120\) publishes the constitutive relations; the
dimensional chart assigns units.
\end{narrativebox}
 
\section{Incidence of the sectors \(90/120\): record, spectrum, Barbero--Immirzi kernel and towers}
\subsection{Typed distinction between the incidence realizations of the sector \(90/120\)}
% Fragmento público generado desde fuentes teoremáticas íntegros.
% Residencia: 52.8.1.
% Fuente: /Users/ruben/Documents/HMT2/EL_CIERRE_HOLOGRAFICO_DEL_INFINITO_SUCESOR_102_CAPITULOS_2026-08-28_EN_TRABAJO/incoming/radion_monografia_20260827/manuscrito/sections/09_sector_90120.tex:176-244
\subsubsection{Global harmonic hierarchy and semigroup \(\langle90,120\rangle\)}

To avoid a symbol collision, we fix
\[
R_{\mathrm{tow}}:=-R_{\mathrm{ch}}.
\]
Then
\begin{equation}
T=e^{-xI+yR_{\mathrm{tow}}}
=q_-P_+^{\mathrm{tow}}+q_+P_-^{\mathrm{tow}}.
\label{eq:tower-operator}
\end{equation}
The even and odd towers are
\begin{align}
c_n&=\operatorname{Tr}(T^n)
=q_-^n+q_+^n
=2e^{-nx}\cosh(ny),
\label{eq:cn}\\
s_n&=\operatorname{Tr}(R_{\mathrm{tow}}T^n)
=q_-^n-q_+^n
=2e^{-nx}\sinh(ny).
\label{eq:sn}
\end{align}
Under \(C^*\mapsto-C^*\), \(c_n\) is even and \(s_n\) is odd. It holds that
\begin{equation}
c_n^2-s_n^2=4e^{-2nx},
\label{eq:tower-invariant}
\end{equation}
and both sequences satisfy the order-two recurrence determined by the
eigenvalues \(q_\pm\).

The semigroup
\begin{equation}
\langle90,120\rangle=30\langle3,4\rangle
\label{eq:semigroup}
\end{equation}
is the harmonic hierarchy of a single bilayer operator. It is not a list of eight
lateral corrections. Its completeness means that every positive ratio
admits a convergent expansion; it does not mean that an expansion obtained
from the target is a prospective physical selector.

\subsubsection{Typed classification and separation of the realizations \(90/120\)}

\begin{longtable}{>{\bfseries}p{.19\textwidth} p{.24\textwidth} p{.45\textwidth}}
\toprule
Object & Operation & What is conserved / why it is not another object\\
\midrule
\endfirsthead
\toprule
Object & Operation & What is conserved / why it is not another object\\
\midrule
\endhead
Incidence \(N_6,N_8\) & finite counting over active phases & cardinality and orientation; it produces \(90,120\), not a series;\\
\(-1/12\) & diferencia normalized \eqref{eq:minus-one-twelve} & incidence scalar; no Barbero weight;\\
\(E_{\mathrm{dod}}\) & harmonic extractor of the wheel & dodecaphase phase; no real carry;\\
\(h=C^*/6\) & reconstruction of six pairs & opening by pair; no metrological normalization;\\
\(F_{\mathrm{spec}}\) & cola logaritmica \(j\ge2\) & complete spectral publication; differs from \(\epsR^\circ\);\\
\(\epsone\) & subtraction APP between \(S_T\) and \(S_E\) & real join of two lifts; not a Lambert series;\\
\(K_{6\to8}\) & \(12\Delta S_{90}-\Delta S_{120}\) & pole \(1/24\) and marked jump; Barbero output;\\
\(12:-1\) & solution Diophantine of \eqref{eq:pole-condition} & peso of the kernel; no \(-1/12\);\\
\(\langle90,120\rangle\) & semigroup of powers of \(T\) & global harmonic hierarchy; no residue per species.\\
\bottomrule
\end{longtable}

\begin{thesisbox}
All these objects descend from the same APP--TRIT--TPK state. That common genealogy explains why they share numbers and symmetries. Their typing explains why they cannot replace one another.
\end{thesisbox}
 \subsection{Conservation of Incidence and Orientation Between the Tritic Kernel and Its Subsequent Realizations}
% Fragmento público generado desde fuentes teoremáticas íntegros.
% Residencia: 52.8.2.
% Fuente: /Users/ruben/Documents/HMT2/EL_CIERRE_HOLOGRAFICO_DEL_INFINITO_SUCESOR_102_CAPITULOS_2026-08-28_EN_TRABAJO/incoming/radion_monografia_20260827/manuscrito/sections/09_sector_90120.tex:1-30
\noindent The typed hierarchy separates incidence, record, spectral evaluation, Barbero functional, and harmonic towers.\par\medskip

\subsubsection{Tritic incidence, dodecaphasic recording and hexadic reading}

\paragraph{Primary combinatorial derivation of \(90\) and \(120\)}

The counts
\[
90=6\binom62,
\qquad
120=8\binom62
\]
arise from the phase selector and its active incidences. The normalized orientation satisfies
\begin{equation}
\frac{30}{120}-\frac{30}{90}=-\frac1{12}.
\label{eq:minus-one-twelve}
\end{equation}
This scalar is an incidence result. It is neither the weight \(12:-1\) of the
Barbero kernel nor a dodecaphase carry.

The dodecaphonic extractor can be published through a combination of
harmonics as
\[
E_{\mathrm{dod}}(\phi)=12\cos(3\phi)-\cos(4\phi).
\]
Its domain is the phase of the wheel; its output is a harmonic selector. Although
the numbers \(12\) and \(-1\) appear, it does not act on Lambert series or
produce \(\epsR\).

 \subsection{Relative Diophantine uniqueness of the weight \(12:-1\) under the pole condition \(1/24\) and the minimal integer counterterm}
% Fragmento público generado desde fuentes teoremáticas íntegros.
% Residencia: 52.8.3.
% Fuente: /Users/ruben/Documents/HMT2/EL_CIERRE_HOLOGRAFICO_DEL_INFINITO_SUCESOR_102_CAPITULOS_2026-08-28_EN_TRABAJO/incoming/radion_monografia_20260827/manuscrito/sections/09_sector_90120.tex:115-175
\subsubsection{Barbero--Immirzi hexad--octad functional}

\paragraph{Oriented series and incidence jump}

We define
\begin{equation}
S_{90}(q)=\frac{q^{90}}{1-q^{270}},
\qquad
S_{120}(q)=\frac{q^{120}}{1-q^{360}},
\label{eq:barbero-series}
\end{equation}
and
\[
\Delta S_m=S_m(q_-)-S_m(q_+).
\]
The marked inclusion functional is
\begin{equation}
K_{6\to8}=12\Delta S_{90}-\Delta S_{120}.
\label{eq:barbero-núcleo}
\end{equation}

\paragraph{Diophantine condition for the weight \(12:-1\)}

Consider \(aS_{90}-bS_{120}\). The declared pole condition is
\begin{equation}
\frac{a}{270}-\frac{b}{360}=\frac1{24}.
\label{eq:pole-condition}
\end{equation}
Multiplying by \(1080\) gives
\[
4a-3b=45.
\]
If the contratermino is integer minimo, \(|b|=1\). For \(b=1\), \(a=12\); for
\(b=-1\), \(a=21/2\) is not integer. therefore,
\begin{equation}
\boxed{a:b=12:1,\qquad aS_{90}-bS_{120}=12S_{90}-S_{120}.}
\label{eq:12minus1}
\end{equation}

\begin{proposition}[Relative Diophantine uniqueness]
The signed weight \(12:-1\) is unique among integer coefficients when the
pole condition \(1/24\), positivity of the principal channel, and the
minimal counterterm \(|b|=1\) are fixed.
\end{proposition}

The hexad--octad incidence supplies the channels \(90/120\). By itself, it does not force the Lambert series or the condition \eqref{eq:pole-condition}. The strength of \(12:-1\) is relative to these explicit analytic premises.

The Barbero-Deformable chain is subsequent:
\[
(A,C^*)\to(x,y)\to q_\pm
\to K_{6\to8}\to\gamma_{\mathrm{BI}},
\]
with evaluation
\[
\gamma_{\mathrm{BI}}=0.274007672694459\ldots.
\]
Barbero--Immirzi does not retroactively generate \(C^*\), \(\Delta_4\), or
\(\epsR\).

 \subsection{Extraction of the finite component subsequent to the oriented dodecaphasic system \(R_{12}\): \(C^*\mapsto C^*/6\), remainder \(j\ge2\) and reconstruction of six sections}
% Fragmento público generado desde fuentes teoremáticas íntegros.
% Residencia: 52.8.4.
% Fuente: /Users/ruben/Documents/HMT2/EL_CIERRE_HOLOGRAFICO_DEL_INFINITO_SUCESOR_102_CAPITULOS_2026-08-28_EN_TRABAJO/incoming/radion_monografia_20260827/manuscrito/sections/09_sector_90120.tex:31-54
\paragraph{Extraction of component \(C^*/6\) from the dodecaphase record}

The twelve sectors group into six opposite pairs:
\begin{equation}
p_i=e_i+e_{i+6},\qquad
a_i=e_i-e_{i+6},\qquad i=1,\ldots,6.
\label{eq:sexadic-pairs}
\end{equation}
The reconstruction has profile \(1+5+6=12\). The opening per pair is
\begin{equation}
h=\frac{C^*}{6}.
\label{eq:h-C6}
\end{equation}
It is neither an external normalization nor one sixth of a target: it is the
internal index of the hexadic reconstruction of the record \(\Rstate\).

We define
\begin{equation}
q_{h,-}=\exp\left[-\frac\pi{180}(A-h)\right],
\qquad
q_{h,+}=\exp\left[-\frac\pi{180}(A+h)\right].
\label{eq:q-h}
\end{equation}

 \subsection{Quantitative decomposition of the carry and spectral realizations: \(F_{\mathrm{spec}}=\varepsilon_R^\circ+\chi_R\)}
% Fragmento público generado desde fuentes teoremáticas íntegros.
% Residencia: 52.8.5.
% Fuente: /Users/ruben/Documents/HMT2/EL_CIERRE_HOLOGRAFICO_DEL_INFINITO_SUCESOR_102_CAPITULOS_2026-08-28_EN_TRABAJO/incoming/radion_monografia_20260827/manuscrito/sections/09_sector_90120.tex:55-114
\subsubsection{Spectral publication of the tail of orders \(j\ge2\)}

\paragraph{Definition of the spectral functional}

For \(m\in\{90,120\}\), let
\begin{equation}
L_m^{(2)}(q)=\sum_{j\ge2}\frac{q^{mj}}{j}
=-\log(1-q^m)-q^m.
\label{eq:L2}
\end{equation}
The reading of a pair is
\begin{equation}
T_h^{(2)}=\frac{180}{\pi}
\Bigl[
L_{90}^{(2)}(q_{h,-})-L_{90}^{(2)}(q_{h,+})
-L_{120}^{(2)}(q_{h,-})+L_{120}^{(2)}(q_{h,+})
\Bigr].
\label{eq:T-h}
\end{equation}
The reconstruction of the six pairs gives
\begin{equation}
F_{\mathrm{spec}}=6T_h^{(2)}
=5.1499235153094574410\ldots\times10^{-8}\,{}^\circ.
\label{eq:Fspec}
\end{equation}

The radion carry is
\[
F_{\mathrm{acarreo}}=\epsR^\circ
=5.1383016758575282072\ldots\times10^{-8}\,{}^\circ.
\]
They do not coincide:
\begin{equation}
\chi_R:=F_{\mathrm{spec}}-F_{\mathrm{acarreo}}
=1.16218394519292338\ldots\times10^{-10}\,{}^\circ.
\label{eq:chiR}
\end{equation}

The equality
\[
F_{\mathrm{acarreo}}=F_{\mathrm{spec}}-\chi_R
\]
is a valid change of chart once both outputs have been constructed. It must not
be presented as an independent generation of \(\epsR^\circ\) if \(\chi_R\) is
defined by the subtraction itself. The tail is closed and meaningful; what is false
is its direct identification with carry.

\begin{falsifier}
After converting the tail expressed in degrees to the internal angular unit,
\[
\frac{\pi}{180^\circ}F_{\mathrm{spec}}-\epsone
=2.0283936357433839\ldots\times10^{-12}\ne0.
\]
Therefore, neither the complete logarithm nor the tail \(j\ge2\), even after applying the
chart of units, is \(\epsone\). A
second autonomous spectral factorization would need to generate \(\chi_R\) without
using the difference as a definition; the APP--TPK closure of the radion does not depend
on that promotion.
\end{falsifier}

 
\section{Oriented holonomy and physical geometry: Minkowski, Hodge, Lorentz and Cartan--Holst realizations}
% Fragmento público generado desde fuentes teoremáticas íntegros.
% Residencia: 52.9.
% Fuente: /Users/ruben/Documents/HMT2/EL_CIERRE_HOLOGRAFICO_DEL_INFINITO_SUCESOR_102_CAPITULOS_2026-08-28_EN_TRABAJO/incoming/radion_monografia_20260827/manuscrito/sections/10_minkowski_hodge_lorentz.tex:1-55
\noindent Oriented holonomy admits subsequent causal and geometric realizations in Minkowski, Hodge, Lorentz, and Cartan--Holst.\par\medskip

\paragraph{Construction of the Lorentzian metric from the APP block}

\paragraph{APP matrix of parameters \(7\) and \(2\)}

The central matrix
\begin{equation}
B_c=\begin{pmatrix}7&2\\2&7\end{pmatrix}
=9P_++5P_-
\label{eq:Bc}
\end{equation}
it is normalized by its determinant:
\begin{equation}
M_c=\frac{B_c}{\sqrt{45}}
=\exp\left(\frac12\log\frac95\,R\right)
\in SO^+(1,1).
\label{eq:Mc}
\end{equation}
The internal speed is
\[
\eta_c=\frac12\log\frac95.
\]
The associated projective publication satisfies
\begin{equation}
P_B^n=P_++\left(\frac59\right)^nP_-.
\label{eq:PB-contract}
\end{equation}
The same ratio \(5/9\) that appears as a spectral contraction is recognized here
as the separation of one stable mode and one contractive mode. The identity is
operatorial; it does not turn every occurrence of \(5/9\) into the same map.

\paragraph{Incidence \(6\to8\) and causal-signature quotient}

Let \(M\) be the typed incidence between six faces and eight
vertices. The recovered spectral relation is
\begin{equation}
MM^{\mathsf T}=12P_u+4P_-,
\label{eq:MMt}
\end{equation}
where \(P_u\) projects onto the uniform mode and \(P_-\) onto the
oriented subspace. The surviving quotient of dimension four decomposes as
\begin{equation}
Q_4\simeq\R u\oplus E_-.
\label{eq:radion-Q4}
\end{equation}
Declaring the uniform mode negative and the three oriented modes positive,
one obtains
\begin{equation}
B_{\mathrm{caus}}=\operatorname{diag}(-1,1,1,1).
\label{eq:minkowski}
\end{equation}
The Minkowski screen does not select the HMT state. It is the subsequent causal realization of the quotient already constructed by the incidence.

 \subsection{Hodge duality on \(2\)-Lorentzian forms}
% Fragmento público generado desde fuentes teoremáticas íntegros.
% Residencia: 52.9.1.
% Fuente: /Users/ruben/Documents/HMT2/EL_CIERRE_HOLOGRAFICO_DEL_INFINITO_SUCESOR_102_CAPITULOS_2026-08-28_EN_TRABAJO/incoming/radion_monografia_20260827/manuscrito/sections/10_minkowski_hodge_lorentz.tex:56-100
\subsubsection{Hodge Duality and Quarter-turn Operator}

With the signature \((-+++ )\) and an orientation fixed, the Hodge star on
two-forms satisfies
\begin{equation}
\star:\Lambda^2Q_4^*\to\Lambda^2Q_4^*,
\qquad \star^2=-I.
\label{eq:hodge-star}
\end{equation}
Thus, the quarter-turn expressed in the phase plane by
\(J^2=-I\) reappears in field space as a complex structure on
\(\Lambda^2\). The domains differ, but the genealogy of orientation
is preserved by the functor.

Let \(A_{\mathrm{em}}\) be a potential and
\begin{equation}
F=\diff A_{\mathrm{em}}.
\label{eq:F-dA}
\end{equation}
The decomposition of \(F\) relative to a temporal observer produces electric
and magnetic fields; \(\star F\) exchanges their roles with the sign fixed
by the orientation and signature. The constitutive polarity of
\eqref{eq:EM-conjugation} thus acquires a differential realization.

\paragraph{Electromagnetic and causal transduction chain}

The transport that must remain visible is
\begin{equation}
\Gamma_9
\longrightarrow\text{memory helix}
\longrightarrow(Q_4,B_{\mathrm{caus}},\star)
\longrightarrow A_{\mathrm{em}}
\longrightarrow F=\diff A_{\mathrm{em}}
\longrightarrow
m\nabla_u u=q(\iota_uF)^\sharp.
\label{eq:lorentz-chain}
\end{equation}
Each arrow adds structure: the holonomy preserves return; the quotient fixes
causality; Hodge fixes duality; the potential fixes the field; coupling
to charge produces dynamics.

The final equation is the Lorentz force in geometric form. The TPK helix
is its memory antecedent, not a physical orbit already moved by a field.
Dynamics appears only upon introducing \(F\), \(q\), and the connection \(\nabla\).

 \subsection{Separation between internal torsional memory and Cartan--Holst geometric realization}
% Fragmento público generado desde fuentes teoremáticas íntegros.
% Residencia: 52.9.2.
% Fuente: /Users/ruben/Documents/HMT2/EL_CIERRE_HOLOGRAFICO_DEL_INFINITO_SUCESOR_102_CAPITULOS_2026-08-28_EN_TRABAJO/incoming/radion_monografia_20260827/manuscrito/sections/10_minkowski_hodge_lorentz.tex:101-186
\subsubsection{Lorentzian Realization of Helical Trajectories}

In a uniform magnetic field, a particle of charge \(q\), mass \(m\), and relativistic factor \(\gamma\) has
\begin{equation}
\omega_c=\frac{|q|B}{\gamma m},
\qquad
r_L=\frac{p_\perp}{|q|B},
\qquad
\Delta z=\frac{2\pi p_\parallel}{|q|B}.
\label{eq:lorentz-helix}
\end{equation}
The cyclotron phase closes while the parallel motion advances: the orbit is helical. The parallel with \eqref{eq:H9-helix} is structural and can be formalized by a transducer that preserves phase, orientation, and step. The quotient \(q_9\) is not identified with the spatial coordinate \(z\); rather, the same architecture of return plus advance is recognized.

The Lorentz force is not derived from the image of a spiral by similarity.
It is derived by completing \eqref{eq:lorentz-chain}. The image helps explain
why circular phase and longitudinal memory are compatible; the field
form determines the dynamics.

\subsubsection{Confluence between action, area, information, and gravity}

\paragraph{Cartan--Holst Realization of Holonomy}

The gravitational promotion is expressed via a holonomy morphism.
\begin{equation}
\operatorname{Hol}_{\mathcal A}
\bigl(\mathfrak R_{\mathrm{grav}}(\gamma)\bigr)
=\rho_C\bigl(\operatorname{Hol}_{\TPK}(\gamma)\bigr).
\label{eq:holonomy-gravity}
\end{equation}
The right-hand side contains route memory; \(\rho_C\) realizes it in a gravitational connection. Cartan torsion
\[
T^I=D_\omega e^I
\]
appears in the codomain. The global torsion \(\Delta_4\) is an upstream
coordinate that can feed the map, but it is not renamed as \(T^I\) without
\eqref{eq:holonomy-gravity}.

\paragraph{Transduction of the action area into informational area}

The radion fixes one quantum of action per phase. The Barbero hexad--octad
reader subsequently uses the channels \(q_\pm\), the incidence
\(90/120\), and the functional \eqref{eq:barbero-núcleo}. The causal direction is
\[
(A,C^*,\hret)\to q_\pm\to K_{6\to8}
\to\gamma_{\mathrm{BI}}\to\text{area/information}.
\]
The tritic alphabet provides the unit of entropy \(\log3\), and a cut of
\(81\) channels provides \(81\log3\). This entropy is not to be confused with
\(\log2\), with \(\log\varphi\), with a microcanonical entropy, or with a
von Neumann entropy unless the state is specified.

The physical synthesis is affirmative but typed: the radion supplies the oriented unit of action that a holonomy can transport; Barbero publishes the area normalization; the causal screen and Hodge make it possible to realize fields; Cartan--Holst receives the holonomy in gravitational geometry.

\begin{statusbox}
The scalar radion equation is not by itself an Einstein equation. Its
contribution to gravity is to provide an exact section of the architecture of
action, orientation, and holonomy consumed by the gravitational realization.
\end{statusbox}

\subsubsection{Wick Rotation as Analytic Reader Change}

In the conventional formulation, Wick rotation relates a Lorentzian time
variable to a Euclidean coordinate through complex continuation. In HMT,
the root \(J^2=-I\), the involutive map \(R^2=I\), and the threshold
\(N^2=0\) already exist. The change of reader
\[
J\longleftrightarrow R
\]
transports a common/oriented pair between elliptic publication and
hyperbolic publication. Complex continuation subsequently recognizes that
operation in analytic coordinates.

This explains why the circle, the action ellipse, and the hyperbolic interior can
belong to the same state without being the same metric. The circle uses \(J\) for
the phase; the shape of the ellipse uses a symplectic compression generated by \(R\); Klein measures
the interior by means of the cross-ratio; the Minkowski screen realizes causality
in another codomain.
 
\section{Ellipse, Lattice, and Path Phase: Modular Character, Orbits, and Propagation}
\subsection{Rectangular Lattice of Periods, Involution \(\tau\mapsto-1/\tau\), and Invariance of \(j\)}
% Fragmento público generado desde fuentes teoremáticas íntegros.
% Residencia: 52.10.1.
% Fuente: /Users/ruben/Documents/HMT2/EL_CIERRE_HOLOGRAFICO_DEL_INFINITO_SUCESOR_102_CAPITULOS_2026-08-28_EN_TRABAJO/incoming/radion_monografia_20260827/manuscrito/sections/11_modularidad_orbitas.tex:1-84
\noindent Modular, orbital, and path characters preserve distinct invariants of the same oriented form.\par\medskip

\subsubsection{Complex torus associated with the action ellipse}

\paragraph{Rectangular lattice of periods}

The ellipse by itself is not a torus. We declare the lattice
\begin{equation}
\Lambda_\eta=a_S\Z+i,b_S\Z
\label{eq:lattice}
\end{equation}
and the quotient \(\C/\Lambda_\eta\). Its modulus is
\begin{equation}
\tau_\eta=i\frac{b_S}{a_S}=ie^{-\eta}
=i\,0.741048144159375160795483663369300\ldots.
\label{eq:tau-eta}
\end{equation}
The bidirectional inversion \(\eta\mapsto-\eta\) exchanges the axes:
\begin{equation}
\tau_{-\eta}=ie^\eta=-\frac1{\tau_\eta}.
\label{eq:modular-S}
\end{equation}
It is exactly the modular transformation \(S:\tau\mapsto-1/\tau\).

\paragraph{Invariance of the Modular Function \(j\)}

By modular invariance,
\begin{equation}
\boxed{j(\tau_{-\eta})=j(\tau_\eta).}
\label{eq:j-invariance}
\end{equation}
In the normalization \(j(i)=1728\),
\[
j(\tau_\eta)
=5597.43843932408729830097089254768\ldots.
\]
The normalized moonshine function \(J=j-744\) gives
\[
J(\tau_\eta)
=4853.43843932408729830097089254768\ldots.
\]
For the basal APP block,
\[
\tau_0=i\frac{\sqrt5}{3},
\qquad
j(\tau_0)=5369.48832024674306021916882626352\ldots.
\]

\begin{proposition}[What the modular class preserves]
The sheet involution reverses \(\eta\), but \(j\) preserves the class of the
torus. The foci and the axis labelling retain the orientation that \(j\)
forgets.
\end{proposition}

This result connects the two-way structure exactly with modularity once
the quotient has been declared. It does not authorize making the Monster act directly
on the charge, digits, or foci. The exceptional corridor
\[
\mathrm{APP}\to\mathrm{Paley}\to\mathrm{Witt}\to\mathrm{Golay}
\to A_2^{12}\to\mathrm{Leech}\to V^\natural
\to\mathrm{Monster}\to\mathrm{Moonshine}
\]
is another projection of the same integral holonomic state. Identifying a concrete action on the radion requires an intertwiner preserving its relevant coordinates.

\paragraph{Chain of character form invariants}

Gathering the declared realizations,
\begin{equation}
\boxed{
r_\eta=e^{-\eta}
=\sqrt{\frac{A-C^*}{A+C^*}}
=\frac{b_1}{a_1}
=\frac{b_S}{a_S}
=\frac{\Delta P}{\Delta Q}
=\frac{Z_\eta}{Z_*}
=\Im\tau_\eta.}
\label{eq:shape-master}
\end{equation}
A single ratio runs through spectral geometry, action, covariance, impedance, and modulus. Each equality rests on an explicit map; none depends on decimal resemblance.

 \subsection{Compatibility between the modular character and propagation over oriented paths}
% Fragmento público generado desde fuentes teoremáticas íntegros.
% Residencia: 52.10.2.
% Fuente: /Users/ruben/Documents/HMT2/EL_CIERRE_HOLOGRAFICO_DEL_INFINITO_SUCESOR_102_CAPITULOS_2026-08-28_EN_TRABAJO/incoming/radion_monografia_20260827/manuscrito/sections/11_modularidad_orbitas.tex:85-187
\subsubsection{Areal Conservation in Kepler and Sommerfeld Orbits}

\paragraph{Areal Equivalence of the Orbital and Oscillatory Realizations}

Kepler's ellipse lives in configuration space and responds to a central
potential. The radion ellipse lives in phase space and responds to an action
form. They are different objects. They share an areal law because both
realizations possess a conserved two-form.

For a Keplerian orbit,
\begin{equation}
r(\phi)=\frac{p}{1+e\cos(\phi-\phi_0)},
\qquad
\frac{\diff A}{\diff t}=\frac{\ell}{2m}.
\label{eq:kepler}
\end{equation}
The third law takes the form
\[
T^2=\frac{4\pi^2\mu}{K}a^3.
\]
For the phase oscillator,
\begin{equation}
J=\oint p\,\diff q=\frac{2\pi E}{\omega}.
\label{eq:osc-action}
\end{equation}
If the elementary cell is \(J=\hfull\), then
\begin{equation}
E=\hret\omega,\qquad \frac{J}{2\pi}=\hret.
\label{eq:E-hbarw}
\end{equation}
This is the clean convergence between Planck action, angular frequency, and radion.

Sommerfeld quantizes action integrals over cycles. The HMT reading adds memory of which cycle, sheet and return survive. On a Möbius strip, one visible loop can carry \(\hfull/2\) and two loops restore \(\hfull\). This topological factor is not added to the Maslov index of EBK quantization; they belong to corrections with different theorem-based sources.

\paragraph{Orbital Holonomy and Advanced and Retarded Propagations}

An orbit that accumulates holonomy satisfies a condition of the type
\begin{equation}
\epsilon_\gamma\exp\left(\frac{iJ_\gamma}{\hret}\right)=1.
\label{eq:holonomy-action}
\end{equation}
The sign \(\epsilon_\gamma\) preserves the sheet; \(J_\gamma\) preserves the action. Apsidal advance or retardation is interpreted as a decoupling between visible angular closure and complete return. The radion supplies the oriented unit with which that difference is measured.

\subsubsection{Composition of Schrödinger amplitudes and path summation}

\paragraph{Four realizations of the same phase}

Once \(\hret\) is fixed, distinct conventional equations recognize
aspects of the same transport:
\begin{enumerate}
\item Schrödinger publishes the phase through
\(i\hret\partial_t\psi=H\psi\).
\item Pauli adds the double spin leaf and its magnetic coupling.
\item Dirac linearizes the energy--momentum relation and organizes particle and antiparticle within a causal representation.
\item The path sum assigns to each history the phase
\(\exp(iS[\gamma]/\hret)\).
\end{enumerate}
These theories do not retrospectively generate \(\hret\). They receive the action unit and realize different contents of the integral holonomic state.

\paragraph{Phase functional over the path space}

In Feynman's formulation,
\begin{equation}
\mathcal K(b,a)=\int_{\gamma:a\to b}
\exp\left(\frac{i}{\hret}S[\gamma]\right)\mathcal D\gamma.
\label{eq:path-integral}
\end{equation}
An action increment \(\hret\) changes the phase by one radian. An increment
\(\hfull\) changes it by one revolution. The theorem
\(\mathcal A(\theta)=\hret\theta\) converts this relation into area geometry:
the phase of a path can be read as swept action in units of the
radion.

The inversion \(C_M=-\operatorname{rev}\) on dodecaphase words and the identity \eqref{eq:time-reversal} organize conjugate routes. In an absorber realization, the following are distinguished
\[
G_{\mathrm{sym}}=\tfrac12(G_{\mathrm{ret}}+G_{\mathrm{adv}}),
\qquad
G_{\mathrm{odd}}=G_{\mathrm{ret}}-G_{\mathrm{adv}}.
\]
The first sector is even under retarded--advanced interchange; the second is
odd. This reproduces the separation between even energy/form and odd
oriented action of the radion when the transducer preserves the boundary and the
symplectic form.

\begin{statusbox}
Classical absorber theory and the Feynman--Stueckelberg interpretation of
antiparticles are distinct historical realizations. The former uses a
temporally symmetric interaction; the latter reorganizes propagation in
relativistic theory. HMT supplies a common antecedent of reversal, but does not
merge them without a map of propagators.
\end{statusbox}
 
\section{Radical nonadic and decimal boundary: conservation of the radion carry}
\subsection{The carry of the radion in the decimal completion}
% Fragmento público generado desde fuentes teoremáticas íntegros.
% Residencia: 52.11.1.
% Fuente: /Users/ruben/Documents/HMT2/EL_CIERRE_HOLOGRAFICO_DEL_INFINITO_SUCESOR_102_CAPITULOS_2026-08-28_EN_TRABAJO/incoming/radion_monografia_20260827/manuscrito/sections/12_aritmetica_borde.tex:1-159
\noindent The nonadic boundary preserves the sum--product defect, vacancy, and completion carry.\par\medskip

\subsubsection{Radical \(3\! -\! 6\! -\! 9\) of the boundary}

\paragraph{Generation and nilpotency of the nonadic radical}

In the ring \(\Z/9\Z\), the ideal
\begin{equation}
\mathfrak m=(3)=\{0,3,6\}
\label{eq:radical-369}
\end{equation}
satisfies
\begin{equation}
\mathfrak m^2=(9)=0.
\label{eq:m-square-zero}
\end{equation}
It is the nilpotent sector of the nonadic boundary. The visible orbit of its three
marks is published as
\[
369\longrightarrow693\longrightarrow936\longrightarrow369.
\]
The return does not imply that all states are equal: the visible
permutation closes while the route cocycles may advance.

The double convention
\[
9\equiv0\pmod9,
\qquad
\operatorname{dr}_9(9)=9
\]
expresses two readers: algebraic residual zero and visible digital-root
closure. Thus a nine may absorb boundary torsion without becoming
a specific zero of the zeta function.

\paragraph{Sum--product defect \(459-405=54\)}

The APP sheets produce complementary totals
\[
S_+=405,\qquad S_\times=459,\qquad S_\times-S_+=54.
\]
The \(54\) is an exact defect of the tables and reappears as a contraction
exponent through the identity \(9\cdot6=54\). Each occurrence nevertheless
requires its own map. The focal distance \(d_1=0.544545\ldots\) is not equal to
\(0.54\), and invariant analysis has produced no arrow identifying it with
the \(54\) defect.

The power
\[
9^9=387,420,489
\]
generates a scale where
\begin{equation}
\operatorname{ord}_{10}\bigl(9^{9^9}\bigr)=369,693,099,
\qquad
\#\,\text{digits}=369,693,100.
\label{eq:369693}
\end{equation}
The block \(369|693\) appears in the order--length pair; it is not asserted to be
the decimal prefix of the power.

\subsubsection{Arithmetic Theory of Decimal Vacuums}

\paragraph{Asymptotic vacancy slope}

The separation between the decade and the nonary base is
\begin{equation}
\delta_{\mathrm{vac}}=\log_{10}\frac{10}{9}.
\label{eq:delta-vac}
\end{equation}
The accumulated counter
\begin{equation}
V_9(n)=\left\lceil n\delta_{\mathrm{vac}}\right\rceil
\label{eq:vacancy-count}
\end{equation}
records how many correction positions are required by the decimal publication of
a nonadic evolution. Its jumps form a binary word
\[
z_n=V_9(n)-V_9(n-1)\in\{0,1\}.
\]
The gaps are distributed with lengths \(21/22\), the returns with
lengths \(6/7\), and the combined determinant is \(15\). These numbers
arise from the vacancy reader, not from the Barbero kernel, although \(15\)
coincides with \(\binom62\).

\paragraph{Polarization and Arithmetic Current of the Nonadic Grid}

We define
\begin{equation}
P_N=\sum_{n=1}^N(z_n-\delta_{\mathrm{vac}}),
\qquad
J_N=P_N-P_{N-1}=z_N-\delta_{\mathrm{vac}}.
\label{eq:arith-current}
\end{equation}
The balanced sequence satisfies \(|P_N|<1\). The \enquote{current}
\(J_N\) is here an arithmetic current: it measures the creation and absorption of
vacancies relative to the mean slope. Only a subsequent dimensional transducer
can publish it as an electromagnetic current density.

The physical intuition is powerful because the same algebraic form separates an
even average from an oriented fluctuation. Rigor requires retaining the type:
\[
J_N\in\R\quad\text{dimensionless},
\qquad
j^\mu\in\text{physical current line}.
\]

\subsubsection{Completion \(729\to1000\) and preservation of carry}

\paragraph{Decomposition \(1000=729+243+27+1\)}

The completion
\[
1000=729+243+27+1
\]
preserves the unit through four ternary scales. The total crown is
\[
1000-729=271,
\]
while the last integer before the closure satisfies
\[
999=729+270.
\]
The difference between \(270\) and \(271\) distinguishes an occupied boundary from a
completing unit. The same care prevents identifying a vacancy with zero or a
carry with an actual residue.

\paragraph{Residence of the carry of the radion in the decimal crown}

The radion's operands begin, in blocks of three digits, with
\[
\Phi_T:\quad000|027|455|906|837|565|\cdots,
\]
\[
D_E:\quad000|027|455|010|034|743|\cdots.
\]
The subtraction APP gives
\[
\epsone:\quad000|000|000|896|802|822|\cdots.
\]
The small size of the result is not obtained by truncating the low-order
summands. It arises because two distinct publications share a long prefix,
and the operator \(\operatorname{SubAcarreo}_{1000}\) preserves the borrows required to subtract
them exactly.

The word \emph{memory} has two levels here:
\begin{enumerate}
\item \(q_{\mathrm{mem}}\) records the depth of nonadic returns;
\item \(c_i\) records the local borrowings of decimal subtraction.
\end{enumerate}
The first organizes the helix; the second publishes the decimal. \(\epsR\) uses
both because its operands follow from cylinders at depth and its subtraction is
performed with carry, but it is equal to neither counter.

\begin{thesisbox}
The boundary is not a place where the theory loses information. It is the
place where information changes representation: phase that returns, memory
that ascends, cylinder that contracts, vacancy that is recorded, and carry
that transports the difference.
\end{thesisbox}
 \subsection{Compatibility of the direct and torsional sections: structural characterization of \(\pi\)}
% Fragmento público generado desde fuentes teoremáticas íntegros.
% Residencia: 52.11.2.
% Fuente: /Users/ruben/Documents/HMT2/EL_CIERRE_HOLOGRAFICO_DEL_INFINITO_SUCESOR_102_CAPITULOS_2026-08-28_EN_TRABAJO/incoming/radion_monografia_20260827/manuscrito/sections/03_cierre_exacto.tex:286-309
\subsubsection{Structural characterization of \(\pi\) by compatibility of sections}

Reorder \eqref{eq:cierre-unitario} gives
\begin{equation}
\boxed{
2\pih=
\frac{1+\Phi_T-\epsone}
{\frac5{36}+\frac{25}{9}\alphah}.}
\label{eq:pi-radion}
\end{equation}
This equality may be called a new definition of \(\pi\) within the
radion chart if \emph{structural definition} is understood as an
exact characterization by compatibility of sections. It is not a second
independent generator: \(\Delta_4\) and the carry operator already receive the
coinductive section of \(\pi\). The novelty is that the turn is
characterized by the gluing among direct publication, torsion, and memory.

\begin{narrativebox}
The equation \eqref{eq:pi-radion} does not say that we first do not know \(\pi\) and then guess it by closing an angle. It says that the \(\pi\) generated by the nonadic connection is exactly the one that makes two internal readings of the same state compatible. It is an identity of intrinsic recognition, not a metrological calibration.
\end{narrativebox}
 
\section{Evolutionary invariance of the holographic unit across dimensions}
% Conversión TeX fiel del cuerpo científico del delta narrativo.
% Fuente congelada: /Users/ruben/Documents/HMT2/CONTINUIDAD_EDITORIAL/RECONSTRUCCION_ARMONICA_MACROTRATADO_20260822/REVISION_CONJUNTA_PARTES_I_II/auditar_pdf_rigor_academico/docs/DELTA_NARRATIVO_UNIDAD_RADION_AUTOCONSERVACION_AUTOINERCIA_GRAVEDAD_PENTADICA_20260828.md:12-366.
\subsection{Typed Separation among the Curvature Regime \(\tau\in\{+1,0,-1\}\), the Radial Response \(p\), and the Global Spin Holonomy}

The radion's biography can be represented by a raised curve.

\[
\gamma(\theta)
=
\bigl(r(\theta)\cos\theta,\,
r(\theta)\sin\theta,\,
m(\theta)\bigr),
\]

where the first two coordinates belong to the Euclidean plane \(\mathbb R^2\), and \(m(\theta)\) records the memory or depth of the lift. This equation permits a joint reading of circle, helix, and spiral.

When the angle evolves while the radius and memory remain constant, the planar publication is a circle. When angle and memory evolve while the radius retains its value, the complete trajectory is a helix of constant radius. When angle and radius evolve over a fixed memory, a spiral appears in \(\mathbb R^2\). When all three coordinates evolve, the trajectory is a radial helix: a biography in which phase, amplitude, and memory are jointly transformed.

The radial character \(p\) classifies the response law of the radius. One useful family is

\[
r_p(\theta)
=
r_0\bigl[1+p\lambda(\theta-\theta_0)\bigr]^{1/p},
\qquad p\neq 0,
\]

with logarithmic limit

\[
r_0\exp\bigl(\lambda(\theta-\theta_0)\bigr)
\qquad\text{when }p\to0.
\]

In this parameterization, \(p=2\) publishes the Fermat regime, \(p=1\) the Archimedean regime, \(p=0\) the logarithmic regime, \(p=-1\) the hyperbolic regime, and \(p=-2\) the lituus regime, with the corresponding choices of domain and orientation. Fractional values describe intermediate regimes. Thus, \(p=\tfrac12\), \(p=\tfrac32\), or \(p=-\tfrac12\) express radial responses that interpolate among the integer families and extend the available geometric classification.

The character can also change when crossing a threshold. The differential writing

\[
\frac{dx}{d\theta}
=
\beta(m)\,x^{1-p(m)},
\qquad x=\frac r{r_0},
\]

describes an evolution through regimes: \(p(m)\) remains stable within a phase and changes when the memory reaches a publication boundary. The resulting trajectory connects circular, helical, and spiral sectors through a common law. Continuity of the state and the threshold ledger make those sectors phases of a single evolution.

Three invariants occupy here distinct and coordinated levels. The curvature sign

\[
\tau\in\{+1,0,-1\}
\]

classifies the elliptic, parabolic, and hyperbolic regimes. The character \(p\) classifies the radial response. Spin records the orientation holonomy accumulated by the lift. Curvature, radial response, and spin form a reading system: curvature determines the local type of space, \(p\) determines how the trajectory occupies that space, and spin preserves the orientation acquired while traversing it. Closed and open curves then appear as distinct geometric biographies of the same oriented unit.

\subsection{Particle as a Persistent Section of an Enriched Orbit: Phase, Route, Fibre, Multisection, and Memory}

Dimensional Mechanics can formulate the particle as a persistent section of that geometric and dynamical extension. A compact conceptual notation is

\[
\mathfrak P
=
\operatorname{Sec}_{\mathrm{pers}}
\bigl[
\tau,p,\operatorname{Hol};
\text{sheet},\text{fiber},\text{route},\text{memory}
\bigr].
\]

The particle preserves a recognizable route through changes of phase, curvature, and dimension. Its spin publishes orientation holonomy; its charge publishes phase polarity; its color publishes an internal orientation of the gauge fiber; its flavor distinguishes families of routes or multisections; its mass expresses the dimensional-transduction character by which areal action acquires a volumetric reading. Each property is a reader of the persistent state, and all refer to a common genealogy.

The electron occupies the basal threshold of this construction because it offers the first stable reader in which phase, charge, spin and mass are simultaneously published. The rector electronic mass

\[
m_e^\beta=0.510998950690000808608\ldots\ \mathrm{MeV}
\]

fixes a transduction reference; the preceding basal states retain their function as genealogical stages. Starting from the electron, the quotients \(K_D/K_\Omega\), towers \(90/120\) and domain \(81/56/13/324\) organize the expression of massive systems. The mass law therefore reads how each persistent section traverses holonomic space and converts areal action density into volumetric presence.

This interpretation connects spirals to particles precisely. The index \(p\) supplies the radial response; \(\tau\) supplies the curvature regime; holonomy supplies spin; the route and the fiber supply flavor and color; area--volume transduction supplies the mass character. The particle atlas thereby becomes a classification of stable modes of dimensional traversal.

\subsection{Transduction of the oriented phase in electric polarization, magnetic response, and boundary radiative emission}

Electricity is presented as directed transport of phase and action through a vacancy structure. A vacancy is an admissible transition of the support: it opens a channel, fixes a threshold, and permits the update to advance. The electric component publishes transport tangential to the route. The magnetic component publishes the transverse curvature induced by that transport. Their perpendicularity expresses the geometric relation between advance and rotation.

When two sinusoidal transverse components maintain the appropriate phase difference, the field vector rotates. Longitudinal propagation lifts that rotation to a helix. Variation of amplitude transforms its frontal projection into a spiral. The geometric reading is then fully coordinated: angle carries phase; angular velocity carries frequency; radius carries amplitude and action density; orientation carries polarization and direction; the vertical coordinate carries propagation and memory; \(p\) carries the radial regime; the threshold carries emission, absorption, or change of state.

Electromagnetism appears as a phase–action–memory dynamics with complementary publications. Electricity records polar advance; magnetism records the transverse response; the wave records periodicity; the helix records advance with memory; the spiral records radial evolution; radiation records the publication of the state when accumulation reaches a threshold.

This architecture also orders the role of vacancies and cycles. Vacancy supplies the local possibility of transition; the cycle preserves phase; memory distinguishes turns; the threshold determines the next residence of the action. Propagation may remain confined, change radial regime, or be published as controlled radiation. Field forces thereby acquire a continuous genealogical description from the discrete support to the observed wave.

\subsection{Reversibility of the complete dynamics and thermodynamic cost of information-reducing projections}
The Landauer zero represents a limiting case of this law. For a bijective update of the complete state,

\[
H(X\mid U(X))=0,
\qquad
\inf Q_L=0.
\]

The information remains available in the genealogy. When a reader forms a quotient \(q\), entropy decomposes as

\[
H(X)=H(q(X))+H(X\mid q(X)).
\]

The conditional term measures the information retained outside the projection. A materialized tritic reset satisfies the threshold

\[
Q_{\mathrm{reset}}\ge k_B\Theta\log 3.
\]

The Boltzmann constant translates the multiplicity grouped by the thermal reader; the Nyquist threshold separates faithful sampling from aliasing; null Landauer characterizes the reversible transport of the integral state. Information, temperature and action are articulated by the type of reading applied.

\subsection{Self-conservative forces as transport of action among the observable,
radial, memory, vacancy, and radiation sectors}

Evolutionary conservation is better expressed through exchange than through scalar immobility. A complete state distributes action and momentum among visible motion, memory, radial deformation, vacancies, radiation, and the conjugate return branch. The enriched balance may be written as

\[
\Delta J_{\mathrm{visible}}
+\Delta J_{\mathrm{memoria}}
+\Delta J_{\mathrm{radial}}
+\Delta J_{\mathrm{vacancias}}
+\Delta J_{\text{radiation}}
+\Delta J_{\mathrm{retorno}}
=0.
\]

A force is \textbf{self-conserving} when its own update law contains the channels that transport, store, publish, and reabsorb action. Conservation belongs to the complete genealogical balance of the interaction. Each reader observes part of the exchange; the genealogy reconstructs its totality.

The Earth--Moon system provides a direct physical image. Earth's rotation, the lunar orbit, tides, deformation, and radiation exchange angular momentum. The net consumption of the complete system closes in the genealogical balance of those transfers. The magnitude called energy functions as a reading derived from the action and momentum published in each sector. Fundamental conservation resides in compatibility among all readers of the exchange.

This formulation dissolves the former boundary between conservative forces and so-called nonconservative forces through a deeper category. Friction, apparent dissipation, and radiation occupy explicit channels of the extended state. A loss in the visible mechanical projection corresponds to a transfer toward memory, deformation, heat, vacancies, or radiation. Self-conservation follows the residence of action through those changes.

\subsection{Self-inertial systems: publication of acceleration through orientation,
curvature, sheet, and memory}

The appropriate expression is \textbf{relational self-inertia}. It denotes the compatibility between the total state's own connection and the response that appears when it is projected within the global distribution. The endogenous and exogenous aspects are analytic components of a single dynamic relation.

Let \(\Gamma\) be a horizontal trajectory or geodesic of the complete state. Its intrinsic law satisfies

\[
\nabla_{\dot\Gamma}\dot\Gamma=0.
\]

When it is projected onto a reader \(S\), the observed acceleration is determined by the curvature of the projection:

\[
\nabla^S_{\dot\Gamma_S}\dot\Gamma_S
=
\mathcal K_{\pi_S}(\dot\Gamma,\dot\Gamma).
\]

The published acceleration expresses the relation between total motion and the geometry of the reader. Mach's principle completes this mechanism through the global boundary trace

\[
b:\mathcal X_\infty\longrightarrow B_\infty
\]

and the factorization of the local inertial response

\[
I_v=\bar I_v\circ b.
\]

The global distribution enters the local response through a typed map.
Relational self-inertia therefore combines the endogenous connection of
motion, the global condition transmitted by \(b\), and the
observer's concrete projection. A synthetic expression for this
coordination is

\[
\nabla^{\mathrm{auto}}
=
\mathcal M\bigl(
\nabla^{\mathrm{end}},
b[\mathcal X_\infty],
\pi_S
\bigr),
\]

where \(\mathcal M\) represents the Machian coupling and preserves the type of each component.

The Ambivalence Principle adds the double orientation of reading. Every direct publication preserves a conjugate return branch; emission and absorption form related endpoints of the same biography; temporal orientation is read from both directions of transport. Absorber theory finds a natural formulation here: source and receiver jointly close the genealogical balance of phase and action. The torsional present is the interface at which both orientations are rendered compatible.

TRIT acts at this level as a time crystal: it organizes the local periodicity of updating and, through memory, preserves the difference between successive returns. The direct branch and the conjugate branch relate emission to absorption, and the observer--observable coupling acquires a dynamical residence within the same state. Published time is the reading of that orientation; genealogical time is the complete sequence of phase, return, and memory.

Mach and Ambivalence perform complementary functions. Mach relates the local response to the global distribution. Ambivalence preserves the forward and conjugate directions of evolution. Relational self-inertia articulates both and replaces the division between inertial and non-inertial systems with a single law of connection, projection, and memory.

\subsection{Area--volume duality, hexad--octad transition, and the role of the Barbero--Immirzi parameter}

The passage from area to volume constitutes a central dimensional
transduction. Area publishes boundary action; volume publishes interior
depth. The HMT derivation of the Barbero–Immirzi parameter organizes the
compatibility between the two measures through hexad–octad incidence and
the channels \(90/120\):

\[
90=6\binom62,
\qquad
120=8\binom62.
\]

The cardinals \(6\) and \(8\) admit a cubic realization by means of faces and vertices, whereas the Steiner hexad and octad are distinct incidence objects linked only by the typed morphism \(6\to8\). The \(12\) edges of the cube retain their own combinatorial function and do not by themselves define Steiner incidence. The difference in incidence expresses a productive dimensional incompatibility: surface information requires a transduction rule in order to acquire volumetric residence.

The same principle appears in the carry \(999+1=1000\). The nines reach the boundary of the decimal register; the increment simultaneously publishes the closure of the preceding level and the origin \(0/1\) of the next. The holographic ratio among \(729\), \(243\), \(27\), and \(1\) unfolds that transition at triadic scales. The Barbero–Immirzi transduction is inserted into this family of reader changes: boundary and interior preserve the unit while redistributing its measure.

The normalized defect \(-1/12\) and the oriented weight \(12:-1\) perform distinct functions within this geometry. The former quantifies an incidence correction; the latter records a signed orientation in the transduction. Their coordination with the minimal area, the modular Hamiltonian, and memory provides the bridge to dimensional dynamics.

\subsection{Hyperbolic interior contraction, exterior expansion, and holographic conservation of the boundary}

The thesis may be formulated from a single idea: unity preserves its identity while changing scale, reader, curvature, memory, and dimension. This permanence has an evolutionary character because unity unfolds, acquires new coordinates, and publishes different properties; it has a holographic character because each finite publication retains the rule for reconstructing the state from which it proceeds. The Holographic Closure of Infinity designates precisely this compatibility between identity and unfolding: each stage contains the law that permits the same genealogy to be deepened inward and prolonged outward.

The governing designation is \textbf{principle of evolutionary conservation of the unit through dimensional transduction and holographic projection}. The expression places each term at its proper level. Conservation pertains to genealogical identity; evolution pertains to updating; transduction relates dimensional realizations; holographic projection preserves the reconstructive correspondence between its readers.

Elementary arithmetic offers the first image of this architecture. The decimal completion

\[
1000=10^3=729+243+27+1=3^6+3^5+3^3+3^0
\]

exposes four triadic strata within a decimal cubic unit. Upon normalization,

\[
1=\frac{729}{1000}+\frac{243}{1000}+\frac{27}{1000}+\frac{1}{1000},
\]

the unit is recognized simultaneously as totality and as graded unfolding. The reciprocal reading

\[
3^{-6},\quad 3^{-5},\quad 3^{-3},\quad 3^0
\]

describes the inner depth of those strata; its normalization produces the corresponding partition. Relation \(999+1=1000\) expresses the transport mechanism: the last unit completes a register and acts simultaneously as the carry inaugurating the next. Zero and one thus form a continuity interface. The intervals \([0,1]\), \([0,10]\), and the prolongation toward \([0,\infty)\) are scalar readers of a single completion law.

The first splitting \(1000=729+271\), followed by \(271=243+27+1\), shapes the extreme and mean holographic ratio of the deployment: a dominant triadic core preserves the same resolution rule within the remainder. The decimal cubic unit thus contains a nested triadic depth. Movement toward \(729\), \(243\), \(27\), and \(1\) reads the inward contraction; the carry toward one thousand reads the outward expansion.

This double orientation organizes HMT infinity. Inward, refinement grows hyperbolically: each cell admits a new depth without losing its filiation. Outward, unity expands through successive dimensional publications. Closure consists in the commutativity of both movements. If \(U\) represents the update of the state, \(R_d\) the reader on stratum \(d\), and \(T_{d\to d'}\) the transduction between strata, evolutionary conservation is expressed by means of

\[
R_{d'}\circ U
=
T_{d\to d'}\circ R_d.
\]

The same event may provide an angular measure, an action density, a published mass, or a gravitational response. Their common identity resides in the genealogy that makes those readings commute.

\subsubsection{Nonadic Monodromy and Helical Lifting with Memory Advance}

APP, TRIT, and TPK provide the formal mechanism for this unfolding. APP constructs the support and separates the evaluations that must be preserved. TRIT organizes the local regimes and their ambivalence. TPK transports phase and memory over the enriched support. The causal dependence is fixed by

\[
U_t\longrightarrow \Gamma_9\longrightarrow K_{\mathrm{ph}}.
\]

Phase return and memory advance form coordinated operations. After nine steps, the same angular class may reappear,

\[
g(K+9)=g(K),
\]

as the memory counter progresses,

\[
q_{\mathrm{mem}}\Gamma_9=q_{\mathrm{mem}}+1,
\qquad
B_{K+9}\neq B_K.
\]

The precise geometric figure is a helical lift. A closed orbit in the angular projection lifts to an open trajectory when each phase return increments the memory. The circle preserves the cycle; the helix preserves the cycle and additionally records the depth reached. The coinductive monodromy thereby generates infinite depth through a sequence of locally closed and globally ascending returns.

This image has more than merely illustrative value. The angular wheel constitutes a section of the helix; the helix constitutes the wheel's complete biography when the state transports memory. Each turn recovers a phase position and simultaneously initiates a new sheet. Infinitude arises from compatibility between return and lifting: iterating the cycle produces depth without dissolving orbital identity.

The nonadic certificate gives finite and verifiable content to that architecture. The chain

\[
w_6\longrightarrow w_{12}\longrightarrow w_{18}\longrightarrow
w_{24}\longrightarrow w_{30}\longrightarrow R_{36}\longrightarrow G_9
\]

orders the lift; the \(396\) levels, the \(2376\) trits, the \(43\) turns, and the \(387\) pairs \((K,K+9)\) per channel quantify the transit; the five regions of \(\pi\) and the ambient fiber of \(729\) elements publish subsequent internal readings. This sequence makes a productive monodromy visible: the return recovers the phase and coinduction preserves the growth of memory.

\subsubsection{Rational unit and coordination of the circular, elliptic, projective, and helical readers}

The radion concentrates within a metrological unit the structure that the radian publishes angularly. Its natural definition is that of an \textbf{oriented section of phase, action, and memory}. One unit of phase \(\theta=1\) transports one unit of action \(\sigma\hbar_{\mathrm{ret}}\), belongs to a complete turn \(T_+=2\pi\), and preserves sheet, orientation, nonadic return, and position in the refinement cylinder.

The radian expresses the visible angular coordinate. The radion preserves, together with that coordinate, the transported action, orientation, carry, and memory of the revolution. Angular metrology thereby acquires a genealogy. Measuring an arc means identifying how much rotation has been published and which complete state has realized that rotation.

The relationship between phase and area takes on a particularly clear form:

\[
\mathcal A(\theta)=\hbar_{\mathrm{ret}}\,\theta.
\]

One radional unit sweeps an action area \(\hbar_{\mathrm{ret}}\); one complete turn encloses \(h_{\mathrm{ret}}\). The passage from phase to area is incorporated into the reading unit itself. The angular metric, action, and memory cease to appear as juxtaposed quantities and become coordinates of the same oriented section.

The radion admits several coordinated geometric publications. The circular reader publishes phase; the elliptical reader publishes action and form; the Klein interior reader publishes depth; the nonadic helical reader publishes return and memory. The unit remains genealogically identical across those realizations. The radion therefore constitutes the appropriate metrological link between HMT and Dimensional Mechanics: it offers a unit that can traverse changes of reader while preserving the state presupposed by each partial measurement.

\subsection{Five-Component Gravitational Polarity and Joint Transduction of Action, Area, Volume, Time, and Mass}

Dimensional Mechanics makes it possible to define dimension as a stable publication regime characterized by the minimum rank of independent memory required by the lift. To traverse a dimension is to apply a transducer that preserves genealogy and changes the type of reading. Gravity occupies the level of connection between these regimes.

In this formulation, \textbf{interdimensional gravitation with five correlative polarizations}, abbreviated as the \textbf{pentapolar structure of gravitation}, is the connection that transports genealogical unity and leaves minimal torsion at each interface. A compact notation for this trace is

\[
\Theta_d^{\min}
=
\nabla_{d+1}\circ T_{d\to d+1}
-
T_{d\to d+1}\circ\nabla_d.
\]

The expression measures the difference between transporting after differentiating and differentiating after transporting. That difference is the minimal geometric memory of the dimensional crossing. The composition of interfaces produces a gravitational holonomy

\[
\mathfrak G_{d\to D}
=
\operatorname{Hol}
\bigl(
\Theta_d^{\min},
\Theta_{d+1}^{\min},
\ldots,
\Theta_{D-1}^{\min}
\bigr).
\]

The fundamental object is this interdimensional connection. A local graviton-like mode occupies the derived rank of the quantized publication of torsion; gravitation retains its primary residence in compatibility among dimensions.

The term pentapolar encompasses the five polarizations inherent to the continuous state:

\[
\mathbf G_d^{\pm}
=
\bigl(
G_{\mathrm{sp},d}^{\pm},
G_{\mathrm{sol},d}^{\pm},
G_{\mathrm{gau},d}^{\pm},
G_{\mathrm{coh},d}^{\pm},
G_{\mathrm{det},d}^{\pm}
\bigr).
\]

The five components emerge correlatively from the same integral holonomic state. The polarity \(\pm\) records direct concentration and conjugate expression. Gravitation crosses the dimensional strata by transporting these five readings simultaneously; minimal torsion preserves the trace of each crossing.

Geometry thereby acquires a unified image. The elliptic regime organizes closures; the parabolic regime organizes thresholds; the hyperbolic regime organizes depth. Dimensional transduction connects these regimes with the radial responses \(p\), spin holonomies, and mass publications. The pentapolar structure of gravitation expresses the global coherence of that connection.

M-theory supplies subsequent physical screens for this connection: additional dimensions, strings, branes and T-duality realize concrete modes of transduction. HMT supplies the APP–TRIT–TPK genealogy, the integral holonomic state, memory and bidirectional reading; the M-theory interface expresses these structures in typed physical geometries. T-duality interchanges compatible dimensional readers and preserves the genealogical class of the state.

\subsubsection{Compatibility of the classical, relativistic, quantum, and HMT-MD realizations}

Classical mechanics publishes trajectory and momentum. Relativity publishes causal structure, metric, and curvature. Quantum mechanics publishes phase, holonomy, vacancies, and amplitudes. Dimensional Mechanics provides the transducers that relate those publications, and HMT preserves the discrete genealogy that sustains them.

In the classical reading, an orbit describes the visible projection of motion. In the relativistic reading, that orbit is incorporated into the geometry of the reader. In the quantum reading, phase and spin record the holonomy of the section. In the HMT--MD reading, the three expressions belong to the same lifted biography: the radion quantifies phase, action, and memory; \(p\) classifies the radial response; \(\tau\) classifies curvature; holonomy publishes spin; the dimensional transducer publishes mass; the pentapolar connection publishes gravitation.

Apparent classical dissipation acquires a residence in the expanded genealogical balance. Relativistic acceleration is expressed through connection and projection. Quantum indeterminacy is organized by readers, vacancies, and branches of Ambivalence. Mach's principle relates the local response to the global distribution. Zero Landauer cost guarantees the informational reversibility of the complete state. All these elements converge in a deeper conservation: the identity of the unit across changes of publication.

\subsubsection{Genealogical Conservation and Holographic Prolongation Narrative Theorem}

The closure can now be narrated as a single sequence. The unit decomposes without losing identity. The oriented phase acquires measure in the radion. The nonadic return converts the cycle into a helix through memory. The character \(p\) transforms the helix into radial families and classifies their regimes. The persistence of those trajectories produces sections with spin, charge, flavour, colour, and mass. Vacancies open channels; electrical polarity transports phase; magnetic response curves transversely; thresholds publish radiation. Self-conservation closes the genealogical balance of action and momentum. Relational self-inertia links the proper connection, projection, and Machian totality. Ambivalence preserves the direct and conjugate directions. Area--volume transduction leads to Barbero--Immirzi and the mass law. Minimal torsion links dimensions and publishes the pentapolar structure of gravitation.

Inward growth and outward expansion are conjugate orientations of the same process. Hyperbolic refinement generates depth; dimensional transduction generates extension. Carry \(0/1\) links scales; monodromy links turns; memory links states; holonomy links orientations; gravitation links dimensions. Each linkage preserves an evolving unity, and every finite publication contains the rule of its prolongation.

The fundamental thesis is formulated as follows:

\begin{quote}
\textbf{The holographic closure of infinity realizes the principle of the evolutionary conservation of unity through dimensional transduction and holographic projection: unity maintains its genealogical identity as it unfolds as phase, action, curvature, memory, particle, field, and gravitation; every finite reading preserves the rule of reconstruction and prolongation of the complete state.}
\end{quote}

This formulation provides a joint physical and mathematical image. Infinity acquires a law of local closure; the unit acquires depth; dimension acquires memory; conservation acquires relational character; force acquires self-preservation; inertia acquires Machian structure; gravity acquires a pentapolar connection; and holography acquires the precise meaning of reconstruction among compatible readers.
 
\section{Operator definition of the angular radion through the integral holonomic state and four typed readers}
\subsection{Operator definition and fourfold typed realization}
% Fragmento público generado desde fuentes teoremáticas íntegros.
% Residencia: 52.13.
% Fuente: /Users/ruben/Documents/HMT2/EL_CIERRE_HOLOGRAFICO_DEL_INFINITO_SUCESOR_102_CAPITULOS_2026-08-28_EN_TRABAJO/incoming/radion_monografia_20260827/manuscrito/sections/13_sintesis.tex:1-80
\noindent The operator synthesis brings together the phase, action, depth, and memory readers without identifying them with one another.\par\medskip

\subsubsection{Quadruple realization of the radional state}

Let \(x_\infty\in X_{\mathrm{enr}}\) be the coinductive section. The complete radion is not
an isolated point, but rather the bundle
\begin{equation}
\mathfrak R_\sigma(x_\infty)
=\bigl(
\theta=1,\sigma,
\pih,\alphah,A,C^*,\Delta_4,
\Phi_T,D_E,\epsone,\hret,
q_{\mathrm{mem}},\mathcal I_\infty
\bigr).
\label{eq:full-radion}
\end{equation}
Four readers extract different objects from it:

\begin{figure}[H]
\centering
\begin{tikzpicture}[
  node distance=17mm and 18mm,
  state/.style={draw=RadNavy,very thick,rounded corners,fill=RadCream,
    minimum width=39mm,minimum height=13mm,align=center},
  readout/.style={draw=RadTeal,rounded corners,fill=RadMist,
    minimum width=39mm,minimum height=12mm,align=center},
  arr/.style={-{Latex},thick,RadCopper}]
\node[state] (x) {integral holonomic state\\\(x_\infty\)};
\node[readout,above left=of x] (circ) {circle \(S^1\)\\phase};
\node[readout,above right=of x] (ell) {ellipse positive\\action and form};
\node[readout,below left=of x] (klein) {interior Klein\\profundidad};
\node[readout,below right=of x] (helix) {nonadic helix\\memory and return};
\draw[arr] (x)--(circ) node[midway,below left]{\small \(P_{\mathrm{ph}}\)};
\draw[arr] (x)--(ell) node[midway,below right]{\small \(P_{\mathrm{act}}\)};
\draw[arr] (x)--(klein) node[midway,above left]{\small \(P_K\)};
\draw[arr] (x)--(helix) node[midway,above right]{\small \(P_{\mathrm{mem}}\)};
\end{tikzpicture}
\caption{Four typed publications of the same antecedent. The unity of the
state does not eliminate the difference among codomains.}
\label{fig:four-readers}
\end{figure}

The circle preserves the phase class and forgets the lift. The ellipse preserves
action, anisotropy, and symplectic orientation. Klein endows the interior with a
projective distance. The helix preserves the number of returns and refinement.
Dimensional Mechanics does not consist in mixing them, but in constructing the
transducers that make them compatible.

\subsubsection{Operator interpretation of the cocycle \(\varepsilon_R\)}

\(\varepsilon_R\) is not a universal constant independent of every context. It is the
coordinate of a transition cocycle associated with the radion chart. Its
domain contains two sections of the same state:
\[
S_E=1+D_E,
\qquad
S_T=1+\Phi_T.
\]
His rule is
\begin{equation}
\epsone=S_T-S_E=\Phi_T-D_E.
\label{eq:epsilon-final-meaning}
\end{equation}
Its codomain is the dimensionless line of angular unity; the chart
\(360/T_+\) publishes it in degrees.

\begin{exactbox}
\textbf{In plain language:} the direct section and torsional section reach
the same unit cell by two different internal routes. Their integer
parts coincide; their remainders do not. \(\epsR\) is the oriented difference of those
remainders, computed with APP carry and carried to the limit by the TPK
cylinders. It is not introduced to make the radian close: it appears before the
radian is consulted and, when published, closes it.
\end{exactbox}

This definition explains why the number is small. It is not because a
small correction was sought, but because \(S_E\) and \(S_T\) share a
deep prefix. Smallness measures the structural proximity of two
publications; the sign measures their orientation; the digits record the carry.

 % Fragmento público generado desde fuentes teoremáticas íntegros.
% Residencia: 52.13.1.
% Fuente: /Users/ruben/Documents/HMT2/EL_CIERRE_HOLOGRAFICO_DEL_INFINITO_SUCESOR_102_CAPITULOS_2026-08-28_EN_TRABAJO/incoming/radion_monografia_20260827/manuscrito/sections/13_sintesis.tex:81-134
\subsubsection{Operational definition of the radion}

\begin{definition}[HMT--MD radion]
The radion is the phase-oriented section of action and memory that simultaneously satisfies:
\begin{enumerate}
\item traverses a phase unit \(\theta=1\);
\item sweeps signed action \(\sigma\hret\);
\item belongs to a coinductively generated turn \(T_+=2\pih\);
\item preserves the \(\sigma\) sheet, the nonadic return, and the
refinement cylinder;
\item satisfies the exact junction
\(1=S_E-\Phi_T+\epsone\).
\end{enumerate}
\end{definition}

The conventional radian is the image of the radon under the forgetful functor.
\[
\mathfrak R_\sigma
\longmapsto
\theta=1\in\R/2\pi\Z.
\]
That forgetting is useful and legitimate for Euclidean trigonometry. It becomes
insufficient when physics must distinguish action, sign, return,
field, and history.

\paragraph{Genealogy of the Natural Angular Unit}

The naturalness of the radian is usually justified because the derivative of
\(e^{i\theta}\), \(\sin\theta\) and \(\cos\theta\) assumes its simplest form
when \(\theta\) is measured as the arc/radius quotient. HMT preserves that advantage,
but explains it from an earlier layer:
\[
\theta\quad\leftrightarrow\quad
\frac{\mathcal A(\theta)}{\hret}.
\]
The phase is natural because it measures action in units of \(\hret\); the turn is
natural because it encloses \(\hfull\). The arc-to-radius quotient is the
Euclidean publication of a unit that already possesses symplectic structure and memory.

\paragraph{Compatibility between action ellipse and dimensional growth}

The circle is the form that results when anisotropy is forgotten:
\(C^*=0\Rightarrow\eta=0\Rightarrow a=b\). The ellipse is more informative because
it preserves two eigenvalues, two axes and an exchange orientation at fixed
product. In this sense, the ellipse is a prolongation of the circle, not a circle
deformed by accident.

Klein depth and the helical lift add two types of information that the planar
contour does not contain: interior distance and return memory. This is the
ontological image of Dimensional Mechanics that emerges from the radion:
a new dimension appears when one requires an independent datum, which the
preceding projection had to forget, to be preserved.

 \subsection{Structural properties of the radional unit: genealogy, critical hinge, density, carry, action, boundary--interior compatibility, form, polarity, and modularity}
% Fragmento público generado desde fuentes teoremáticas íntegros.
% Residencia: 52.13.2.
% Fuente: /Users/ruben/Documents/HMT2/EL_CIERRE_HOLOGRAFICO_DEL_INFINITO_SUCESOR_102_CAPITULOS_2026-08-28_EN_TRABAJO/incoming/radion_monografia_20260827/manuscrito/sections/13_sintesis.tex:135-184
\subsubsection{System of component theorems of the radion}

\begin{theorem}[Genealogy]
The object \(\mathfrak R_\sigma\) descends from
\(\APP\to\TRIT\to\TPK\to X_{\mathrm{enr}}\) and does not receive the metrological target
as a selector.
\end{theorem}

\begin{theorem}[Hinge]
In the hundred-mark APP chart, the critical seam and the second
dodecaphasic phase satisfy \(j_{100}(1/2)=\theta_2=50^\circ\), with unique \(m=2\).
\end{theorem}

\begin{theorem}[Density]
The two sheets of the active sector \(N_6=90\), normalized by \(3^6=729\),
produce \(\rho_{\mathrm{tw}}=20/81\).
\end{theorem}

\begin{theorem}[Carry]
The APP subtraction of the sections \(S_T\) and \(S_E\) converges to
\(\epsone\) at interval rate \(3^{-54}\) per return and exactly closes
the unit.
\end{theorem}

\begin{theorem}[Action]
The positive ellipse satisfies \(\mathcal A(\theta)=\hret\theta\); a radion
sweeps \(\hret\), and one turn encloses \(\hfull\).
\end{theorem}

\begin{theorem}[Boundary--interior]
The same boundary has finite action \(\hfull\), whereas the Klein
interior has infinite hyperbolic area.
\end{theorem}

\begin{theorem}[Form]
The quotient \(e^{-\eta}\) is preserved across spectral semiaxes,
action semiaxes, uncertainties, impedance, and the modulus of the torus.
\end{theorem}

\begin{theorem}[Polarity]
The inversion \(C^*\mapsto-C^*\) interchanges constitutive relations, inverts
impedance, and preserves the normalized velocity.
\end{theorem}

\begin{theorem}[Modularity]
The bidirectional inversion acts as \(S:\tau\mapsto-1/\tau\) and preserves
\(j(\tau)\), while the focal labeling preserves the orientation that
\(j\) forgets.
\end{theorem}

% Fuente: /Users/ruben/Documents/HMT2/EL_CIERRE_HOLOGRAFICO_DEL_INFINITO_SUCESOR_102_CAPITULOS_2026-08-28_EN_TRABAJO/incoming/radion_monografia_20260827/manuscrito/sections/A_demostraciones.tex:1-179
\subsubsection{Algebraic and Geometric Proofs United}

\paragraph{Linear proof of the unit closure}

To facilitate independent verification, we gather the argument without narrative interruption.

\begin{enumerate}[label=\textbf{Step \arabic*.},leftmargin=2.2cm]
\item The nonadic connection generates in a correlated manner
\(\pih,e,\varphi,\alphah\).
\item The turn is \(T_+=2\pih\).
\item The wheel fixes \(\theta_2=50^\circ\), and the parangular chart fixes
\(A=1000\alphah\).
\item The direct section is
\[
S_E=T_+\left(\frac{50}{360}+\frac{1000\alphah}{360}\right)
=2\pih\left(\frac5{36}+\frac{25}{9}\alphah\right).
\]
\item The certified intervals yield \(1<S_E<2\); then
\(D_E=S_E-1\).
\item The incidence produces \(N_6=90\); two sheets and the ambient space \(729\)
produce \(\rho=2N_6/729=20/81\).
\item Global torsion is
\[
\Delta_4=\frac{\pih-e\log\pih}{270}
\left(1+\frac{\pih}{729}\right).
\]
\item The section torsional is \(S_T=1+(20/81)\Delta_4\).
\item The subtraction APP defines
\[
\epsone=S_T-S_E=\frac{20}{81}\Delta_4-D_E.
\]
\item By substitution,
\[
S_E-\frac{20}{81}\Delta_4+\epsone=1.
\]
\item Multiplying by \(360/T_+=180/\pih\) produces
\[
\frac{180^\circ}{\pih}
=50^\circ+A^\circ-\frac{20}{81}\Delta_4^\circ+\epsR^\circ.
\]
\end{enumerate}

No step receives the left-hand side of the final equality. The
metrological result is evaluated only at the end as a recognition.

\paragraph{Uniqueness of the carry-subtraction}

Let \(B=1000\). For each position, every integer \(u_i\) admits a unique
Euclidean division
\[
u_i=Bc_i+r_i,\qquad 0\le r_i<B.
\]
Therefore,
\[
r_i=u_i\bmod B,\qquad c_i=(u_i-r_i)/B
\]
are unique. Traversing the positions from the remote end toward the
front, the output carry of one position fixes the input of the next.
By induction, the entire word of residues and carries is unique.

In pseudocode:
\begin{verbatim}
acarreo = remote_acarreo
for i from last_block downto first_block:
    u = torsion[i] - direct[i] + acarreo
    remainder[i] = u mod 1000
    acarreo = (u - remainder[i]) / 1000
return remainder, acarreo_word
\end{verbatim}
The procedure does not consult the expected value of the difference.

\paragraph{Convergence of the depth functional}

Let
\[
F(p,e)=1+\frac{20}{81}
\frac{p-e\log p}{270}\left(1+\frac{p}{729}\right)
-2p\left(\frac5{36}+\frac{25}{9}\alphah\right).
\]
On the compact rectangle \(3\le p\le4\), \(2\le e\le3\), \(F\) is
continuous and continuously differentiable. If \(I_N^\pi\searrow\{\pih\}\) and
\(I_N^e\searrow\{e\}\), the interval image
\[
F(I_N^\pi,I_N^e)
\]
forms a nested family whose width tends to zero. The intersection is a
single real number, \(\epsone\). Monotonicity of the derivatives permits the
endpoints to be evaluated without interior sampling.

The rate \(3^{-54}\) per return follows from the refinement of six trits over
nine levels. At depth \(45\) triads, the width bound
\(<4.81\times10^{-130}\) certifies more than 129 digits of the value.

\paragraph{Derivation of the metric and hyperbolic area of Klein}

On the unit disk, the matrix form of the Klein metric is
\[
g(u,v)=\frac1{1-r^2}I
+\frac1{(1-r^2)^2}
\begin{pmatrix}u\\v\end{pmatrix}
\begin{pmatrix}u&v\end{pmatrix}.
\]
The lemma for the determinant of a rank one perturbation gives
\[
\det g=(1-r^2)^{-3}.
\]
Henceforth
\[
\sqrt{\det g}\,\diff u\diff v
=(1-r^2)^{-3/2}\diff u\diff v.
\]
Integration in polar coordinates yields \eqref{eq:klein-area-R}. The infinite limit is
due to the metric singularity at the boundary; the boundary continues to have
finite symplectic area because that measure uses \(\diff Q\wedge\diff P\),
not \(\sqrt{\det g}\).

\paragraph{Focal Distance on the Klein Disk}

In a diametro, the distance from the center to \(r\) is
\[
d_K(0,r)=\operatorname{artanh}r.
\]
For \(r=e_f\), \eqref{eq:excentricidad} implies
\[
1-e_f^2=e^{-2\eta}.
\]
If \(u_f=\operatorname{artanh}e_f\), then
\[
\operatorname{sech}u_f=\sqrt{1-e_f^2}=e^{-\eta},
\]
and by therefore \(\eta=\log\cosh u_f\), as in
\eqref{eq:focal-hyperbolic-chain}.

\paragraph{Symplectic Compression Transformation}

Define
\begin{equation}
S_\eta=
\begin{pmatrix}e^{\eta/2}&0\\0&e^{-\eta/2}\end{pmatrix},
\qquad \det S_\eta=1.
\label{eq:compresión simpléctica}
\end{equation}
The transported complex structure is
\begin{equation}
J_\eta=S_\eta
\begin{pmatrix}0&-1\\1&0\end{pmatrix}
S_\eta^{-1}
=\begin{pmatrix}0&-e^\eta\\e^{-\eta}&0\end{pmatrix},
\qquad J_\eta^2=-I.
\label{eq:Jeta}
\end{equation}
The curve
\[
\gamma(\theta)=\sqrt{2\hret}\,
S_\eta(\cos\theta,\sin\theta)
=e^{\theta J_\eta}\gamma(0)
\]
is exactly the action ellipse. The hyperbolic symplectic compression preserves area and
transports complex rotation instead of destroying it.

For orientation \(\sigma\),
\[
\gamma_\sigma(\theta)
=(a_S\cos\theta,\sigma b_S\sin\theta),
\qquad
\mathcal A_\sigma(\theta)=\sigma\hret\theta.
\]
The two sheets share foci and action modulus; they differ in signed action.

\paragraph{Relative Diophantine uniqueness of the Barbero weight}

The equation \(4a-3b=45\) has integer solutions
\[
(a,b)=(12+3t,1+4t),\qquad t\in\Z.
\]
The condition \(|b|=1\) leaves \(b=1\), since \(b=-1\) does not belong to the class
\(1\pmod4\). Hence \(a=12\). The proof clearly separates the
incidence datum \(90/120\) from the pole and minimality premises.
 \subsection{Structural necessity, stability and closure of the radional unit}

The construction admits negative controls at each of its strata. If the direct section satisfied \(S_E\notin(1,2)\), the APP quotient and closure cell would change simultaneously; if the active incidence did not produce \(N_6=90\), or if one sheet were removed, the density \(20/81\) would no longer be derived. Generation would be circular if the chart received \(180/\pi\) or \(\varepsilon_R\) before constructing its operands. Nor would the coinductive section exist if subtraction by triads failed to stabilize compatible prefixes. In the elliptic realization, the conditions \(|C^*|<A\) and \(a_Sb_S=2\hbar_{\mathrm{ret}}\) are necessary for positivity and the area law. In the modular fibre, \(C^*=0\) forces \(d=\eta=0\), \(\tau=i\) and \(j=1728\). In the projective interior, the depth thesis requires the hyperbolic area law; in the spectral sector, the equality \(F_{\mathrm{spec}}=\varepsilon_R^\circ\) is excluded by the counterterm \(\chi_R\). Finally, the weight \(12:-1\) jointly requires incidence \(90/120\) and the pole condition \(1/24\), whereas invariance of \(j\) prevents its use as a charge-sign selector.

The physical realization is not a perturbation of a presupposed circle. The arithmetic support preserves unit, residue, quotient, carry and fibre; ternary orientation distinguishes regime and sheet; phase transport preserves the difference between return and memory; the ellipse determines action and anisotropy; the Hilbert--Beltrami--Klein metric measures interior depth; and helical lifting preserves the history of each turn. These realizations are articulated through
\[
 \mathcal A(1)=\hbar_{\mathrm{ret}},\qquad
 \mathcal A(2\pi)=h_{\mathrm{ret}},\qquad
 K=-1,\qquad
 g(k+9)=g(k),\quad B_{k+9}\ne B_k.
\]
The circle is the angular projection, the ellipse the action realization, the Klein metric resolves the interior and the helix preserves return history. The radion is the oriented unit allowing the four realizations to be read without identifying them with one another.

The generated invariants satisfy, in the coinductive branch considered,
\[
\begin{aligned}
 \pi&=3.1415926535897932384626433832795028841\ldots,\\
 \alpha&=0.0072973525692838009972851054723806629\ldots,\\
 A&=7.2973525692838009972851054723806629995\ldots{}^\circ,\\
 C^*&=2.1237383389686457242619513803933607937\ldots{}^\circ,\\
 \eta&=0.299689683921630799684926562863927\ldots,\\
 \hbar_{\mathrm{ret}}&=1.05457181691503906113369058472430\ldots
 \times10^{-34}U_S .
\end{aligned}
\]
The coordinate \(e\) entering this system proceeds from the coinductive Euler reader, not from an external table. Returns at depths \(9,18,27,36,45\) preserve the phase \(8\) and lift memory according to \(0,1,2,3,4\); contraction of the intervals tends to
\[
 3^{-54}=1.7196982334549856127610399316004871\ldots\times10^{-26}.
\]
Controlled removal of one sheet, phase displacement, elimination of the \(e\) channel, substitution of global torsion or alteration of the hexadic tail produce nonzero residues of different orders. This scale separation proves that closure jointly requires both sheets, phase orientation, the \(\pi/e\) channels, global torsion and carry.

Mathematical exactness proceeds from the unit-closure identity; evaluation at increasing depth only certifies its convergence. Residues arising when branches evaluated at different decimal precisions are combined are rounding errors, not new physical terms. This distinction separates the structural identity from its numerical realization without altering the genealogy of the operands.

The physical genealogy of the construction can be expressed precisely. Planck fixes the action dimension of \(h\); Bohr and Sommerfeld relate \(h/2\pi\) to angular quantization; Born and Jordan introduce matrix structure and the oscillator; Dirac formulates \(E=\hbar\omega\); Robertson and Schrödinger establish the covariant form of uncertainty; Feynman relates action to path phase; Wheeler and Feynman develop symmetry between retarded and advanced propagation. These frameworks recognize subsequent realizations of the object: they do not select the sheets, incidence, carry, phase or memory generating it.

The Hodge star on Lorentzian two-forms satisfies \(\star^2=-I\) and supplies the relevant geometric duality. Hochschild cohomology belongs to another domain and enters neither the ellipse, the carry nor the incidence sector \(90/120\). The distinction follows from the domains and codomains of the operations, not from terminological resemblance.

The synthesis closes in the integral holonomic state. Together with the value, APP preserves sheet, quotient, residue, carry and boundary; TRIT determines regime and orientation; TPK returns phase without resetting the state. The nonadic connection constructs the coinductive section of \(\pi\), and the dodecaphasic system determines the hinge of \(50^\circ\). The angular pair supplies the common mode \(A\) and oriented opening \(C^*\); incidence of six active phases produces \(90\), their two sheets on the ambient space of cardinality \(729\) produce \(20/81\), and global torsion supplies the second section. The resulting cocycle is
\[
 \varepsilon_R=\frac{20}{81}\Delta_4-(S_E-1).
\]
Its limit in the angular chart is
\[
 5.13830167585752820716851646226459474899\ldots
 \times10^{-8}\,{}^\circ .
\]
The pair \((A,C^*)\) determines a positive ellipse: the product of its semiaxes fixes \(2\hbar_{\mathrm{ret}}\) and their quotient fixes the shape. One phase unit sweeps out \(\hbar\), whereas the full turn encloses \(h\). Division by \(2\pi\) thus ceases to be an external quotient: it distinguishes action per turn from the action transported by one phase unit.

The circular projection preserves phase and omits anisotropy, interior and memory; the ellipse preserves action and shape; Klein geometry preserves depth; the helix preserves return and history. The symplectic realization recognizes the ellipse as a minimum-covariance contour, the \(LC\) oscillator realizes it as exchange of electric and magnetic stores, and the incidence \(90/120\) determines permeability, permittivity, impedance and speed. Hodge turns orientation into field duality; Lorentz turns field and charge into dynamics; Cartan--Holst receives holonomy in gravitational geometry. The causal order remains unchanged: the discrete construction determines the object, and physical formulations subsequently realize its characters.

The radion is therefore an oriented section of phase, action and memory. Its two orientations satisfy
\[
 \mathcal A(\mathfrak r_+)=+\hbar_{\mathrm{ret}},\qquad
 \mathcal A(\mathfrak r_-)=-\hbar_{\mathrm{ret}}.
\]
Charge is obtained by applying the dimensional transducer; its sign was already preserved in sheet orientation. The incidences \(90/120\), the dodecaphasic extractor, the spectral tail, the Barbero--Immirzi parameter, towers and carry can be compared because they share a common origin, but remain distinct because they act through different operators. The fourfold unity of the radion is summarized by
\[
\boxed{
\begin{aligned}
 \text{phase:}\quad&\theta=1,\\
 \text{action:}\quad&\mathcal A=\sigma\hbar_{\mathrm{ret}},\\
 \text{memory:}\quad&H_9(n+9)=H_9(n)+(0,1),\\
 \text{compatibility:}\quad&1=S_E-\Phi_T+\varepsilon_R.
\end{aligned}}
\]
The radian is its angular coordinate; the radion also preserves action, orientation, carry and memory of the turn.
\endgroup

\let\HMTMonographChapterInput\relax
  \part{Dimensional realizations: electron, action, gravitation, area and information}

\chapter{Trigeometric predefinition, boundary potential and dimensional depths}
\section{Theorem of trigeometric temporal predefinition}

The word \emph{predefinition} names a mathematical property of the generator, not a weakened form of prediction. Specifying it requires gathering the three objects previously presented separately: the bundle of TRIT quadratic algebras, the two TPK clocks and the coinductive section of arbitrary depth.

Let

\[
 X^{\rm enr}=\{(\tau,\mu,o,h,c,\sigma,\lambda,g)\}
\]

the space of integral holonomic states. The coordinate \(\tau\in\{-1,0,+1\}\) is the TRIT regime; \(\mu\in\{\Sigma,\Pi,\varnothing\}\) preserves the sheet; \(o\) preserves orientation; \(h\), helical height; \(c\), carry; \(\sigma\), boundary; \(\lambda\), the route memory record; and \(g\in C_9\), the nonadic phase. Acting on each regime fibre is the algebra

\[
 \mathbb A_\tau=\mathbb R[J_\tau]/(J_\tau^2+\tau I),
\]

so that

\[
 \exp(sJ_\tau)=
 \begin{cases}
  \cos s+J_{+1}\sin s,&\tau=+1,\\
  I+sJ_0,&\tau=0,\\
  \cosh s+J_{-1}\sinh s,&\tau=-1.
 \end{cases}
\]

This is the exact form of coordination among the three geometries. Three metrics are not averaged, nor are three previously given spaces superimposed. At each step, TRIT chooses which of the three exponential laws governs local action; TPK transports the state between fibres without erasing sheet or history. With explicit intermediate domains,

\[
 \operatorname{Sel}_t:X^{\rm enr}\longrightarrow X_t^{\rm sel},
 \qquad
 \operatorname{Tra}_t:X_t^{\rm sel}\longrightarrow X_t^{\rm tra},
 \qquad
 \operatorname{Upd}_t:X_t^{\rm tra}\longrightarrow X^{\rm enr},
\]

and

\[
 \boxed{U_t=\operatorname{Upd}_t\circ
              \operatorname{Tra}_t\circ
              \operatorname{Sel}_t.}
\]

Selection determines sheet and regime; transport applies the selected direction and geometric law; updating incorporates quotient, carry, boundary and memory record. The apparent geometric “mixture” is therefore an ordered word of morphisms between elliptic, parabolic and hyperbolic fibres. Since composition depends on order, its operator part is noncommutative even though the base rotation is abelian.

The same TRIT coordinate has a temporal reading, on the same state but through a second reader:

\[
 \begin{aligned}
 +1&\longleftrightarrow\text{return, memory and past},\\
 0&\longleftrightarrow\text{threshold and present},\\
 -1&\longleftrightarrow\text{bidirectional opening and future}.
 \end{aligned}
\]

Geometry and time are not two names for the same sign. They are two coordinated evaluations of a single unit of topological--temporal information. The past is the datum already accumulated in the memory record; the present is the boundary cut at which continuation is selected; the future is the bundle of compatible extensions opened by the hyperbolic regime. The involution

\[
 C_{\rm rev}(\tau_1,\ldots,\tau_n)
 =(-\tau_n,\ldots,-\tau_1)
\]

satisfies \(C_{\rm rev}^2=I\), preserves the threshold and interchanges opening and return. Lifted to the integral holonomic state, it conjugates the direct dynamics with its inverse:

\[
 \widehat C_{\rm rev}\,\mathcal U\,\widehat C_{\rm rev}
 =\mathcal U^{-1}.
\]

Past and future are thus conjugated around the present without eliminating the memory distinguishing the two orientations.

The first clock is the finite non-split extension

\[
 0\longrightarrow C_{12}\longrightarrow C_{108}
 \longrightarrow C_9\longrightarrow0,
\]

whose cocycle

\[
 c(r,r')=\left[\left\lfloor\frac{r+r'}9\right\rfloor\right]_{12}
\]

records the carry that a mere turn in \(C_9\) would lose. The second clock is the capacity rotation

\[
 R_\delta:\theta\longmapsto\theta+\delta\pmod1,
 \qquad \delta=\log_{10}(10/9),
\]

whose mechanical word is

\[
 z_t=\lfloor(t+1)\delta\rfloor-\lfloor t\delta\rfloor.
\]

They are coupled by

\[
 K_t-K_{t-1}=1-z_t.
\]

This equality proves that aperiodicity does not embellish the finite clock: it governs its capacity. The total system can be written as the skew product

\[
 \mathfrak C_{\rm HMT}=
 \bigl(C_{108}\times\mathbb T\times X^{\rm enr},\mathcal U_{\rm HMT}\bigr),
\]

\[
 \mathcal U_{\rm HMT}(r,\theta,x)=
 \bigl(r+1,\theta+\delta,
 U_{\tau(r,\theta,x)}x\bigr).
\]

If \(\mathcal U_{\rm HMT}^n=I\), the first factor forces \(108\mid n\), whereas the second would require \(n\delta\in\mathbb Z\). Irrationality of \(\delta\) excludes every \(n>0\). Moreover, Sturmian coding has complexity \(p(m)=m+1\), and its nonadic and tritic enrichments satisfy

\[
 p_9(m)=9(m+1),\qquad p_3(m)=3(m+1).
\]

Consequently, the crystal is aperiodic and has long-range order and zero topological entropy. Its additive spectrum contains the module

\[
 \mathcal M_{108,\delta}
 =\frac1{108}\mathbb Z+\delta\mathbb Z\pmod1.
\]

The factor \(1/108\) preserves the finite temporal clock; the irrational generator preserves unlimited ascent. Closure of the first never implies closure of the second or return of the state.

Now let \((X_n,r_{n+1,n})\) be the inverse system of refinements and

\[
 \mathcal G_{\rm HMT}:\Omega_{\rm seed}\longrightarrow
 X_\infty^{\rm enr}=\varprojlim X_n
\]

the APP--TRIT--TPK extension of a seed. For every compatible reader \(L_{Y,n}:X_n\to Y_n\), define

\[
 \operatorname{Predef}_{Y,n}
 =L_{Y,n}\,p_n\,\mathcal G_{\rm HMT}.
\]

Naturality

\[
 s_{n+1,n}L_{Y,n+1}=L_{Y,n}r_{n+1,n}
\]

implies

\[
 \boxed{
 s_{n+1,n}\operatorname{Predef}_{Y,n+1}
 =\operatorname{Predef}_{Y,n}.}
\]

The outputs at all depths are therefore coordinates of a single section of \(\varprojlim Y_n\). Increasing precision does not readjust a constant: it opens a deeper coordinate of the same object. On the complete state, bijectivity of updating also gives

\[
 H(X_{t+1}\mid X_t)=H(X_t\mid X_{t+1})=0,
\]

so the appearance of new information in a reduced shadow does not correspond to creation or erasure in the generator. This is the exact meaning of zero Landauer cost within the conservative evolution of the unit.

Predefinition results from composing the generator with its readers:

\[
 \mathcal L_{\rm HMT--MD}=
 (L_\pi,L_\varphi,L_e,L_\alpha,L_{\rm el},L_\hbar,
 L_{\rm vac},L_{\rm BI},L_G,\ldots).
\]

The coordinates \(\pi,\varphi,e,\alpha\) belong to the same coinductive output; \(\alpha\) occupies the dodecaphasic seam linking the finite clock with the capacity clock. Electronic vacancy reads the local carry defect on that seam; \(\hbar\) reads action depth; Barbero--Immirzi reads areal--volumetric transduction; gravitation reads the rank-five tensor sector. None of these operations receives an external measurement to select its value. The system does not obtain its objects through metrology: it predefines them as typed images of the coinductive section.

Finally, an exact factorization is not a lower-rank result. If

\[
 L_Y=H\circ F,
\]

the factorization exhibits the intermediate object and proves what information each arrow transports. If two routes satisfy

\[
 H_1F_1=H_2F_2=L_Y,
\]

the corresponding square is precisely the certificate that two distinct readings determine the same coordinate without circularity. In this development, “theorem”, “factorization” and “transduction” name complementary forms of the same proof: the theorem states the law, factorization shows its genealogy and transduction determines the type of its output.

The governing formulation is thus established:

\[
 \boxed{
 \begin{gathered}
 \text{APP}\to\text{TRIT}\to\text{TPK}\to
 \text{aperiodic time crystal}\\
 \to\text{coinductive section}\to
 \text{HMT--MD predefinition}
 \end{gathered}}
\]

The ordered combination of the three geometric regimes and conjugation of the three temporal modes are not analogies added afterwards. They are, respectively, the local and global laws of the dynamics preserving the unit while evolving it through all its depths.

\section{APP potential triangle and local realization of the defect \(A_2\)}

The discrepancy between the sum and product sheets can be turned into a mathematical notion of potential without introducing an external physical field. Let \(D_9=\{1,\ldots,9\}\), with positive representatives of residue classes modulo nine, and write the two Euclidean divisions preserved by APP:

\[
 i+j=R^+_{ij}+9q^+_{ij},
 \qquad
 ij=R^\times_{ij}+9q^\times_{ij}.
\]

The joint lift retains not only the visible residues \((R^+_{ij},R^\times_{ij})\). It also retains the heights \((q^+_{ij},q^\times_{ij})\), so that the integer difference between the two evaluations decomposes canonically into a residual tension and a carry tension:

\[
 \begin{aligned}
 \rho_{\Sigma\Pi}(i,j)
   &:=R^\times_{ij}-R^+_{ij},\\
 \kappa_{\Sigma\Pi}(i,j)
   &:=9\bigl(q^\times_{ij}-q^+_{ij}\bigr),\\
 \Delta_{\Sigma\Pi}(i,j)
   &:=ij-(i+j)
    =\rho_{\Sigma\Pi}(i,j)+\kappa_{\Sigma\Pi}(i,j).
 \end{aligned}
\]

Thus,

\[
 \mathscr P_{\Sigma\Pi}:D_9^2\longrightarrow\mathbb Z^2,
 \qquad
 (i,j)\longmapsto
 \bigl(\rho_{\Sigma\Pi}(i,j),\kappa_{\Sigma\Pi}(i,j)\bigr)
\]

is the lifted APP potential, whereas \(\Delta_{\Sigma\Pi}=(1,1)\circ\mathscr P_{\Sigma\Pi}\) is its scalar shadow. The term “potential” has a precise content here: it measures the arithmetic energy needed to pass from cumulative to folded evaluation while preserving the full Euclidean division. A projection modulo nine sees only \(\rho_{\Sigma\Pi}\); a projection forgetting the residue and retaining height sees only \(\kappa_{\Sigma\Pi}\). APP preserves both components simultaneously.

Three cells already selected by electronic genealogy form an exact triangle:

\[
 V=(8,2),\qquad E=(5,5),\qquad Q=(8,5).
\]

Their two sheets and potentials are

\[
\begin{array}{c|c|c|c|c}
x&(R^+,q^+)&(R^\times,q^\times)
 &\mathscr P_{\Sigma\Pi}(x)&\Delta_{\Sigma\Pi}(x)\\
\hline
V=(8,2)&(1,1)&(7,1)&(6,0)&6\\
E=(5,5)&(1,1)&(7,2)&(6,9)&15\\
Q=(8,5)&(4,1)&(4,4)&(0,27)&27.
\end{array}
\]

The cell \(V\) is a pure residual tension: the two sheets have the same height and differ in residue. The cell \(Q\) is a pure memory tension: the residues coincide and the entire discrepancy is stored in three nonadic carry units. The central cell \(E\) combines a six-unit residual tension and one carry unit, so that

\[
 \Delta_{\Sigma\Pi}(E)=6+9=15.
\]

This decomposition gives a finer definition of electronic vacancy. On the additive fibre \(i+j=10\), parametrized by

\[
 Q_t=(5+t,5-t),\qquad -4\leq t\leq4,
\]

the product is

\[
 \Pi(Q_t)=25-t^2.
\]

Preserving the multiplicative residue of the central cell requires \(t^2\equiv0\pmod9\). In the permitted interval this occurs only for \(t=0,\pm3\). Therefore, the only two noncentral cells simultaneously preserving the additive sheet, product residue and tritic regime are

\[
 Q_{+3}=(8,2),\qquad Q_{-3}=(2,8).
\]

In them the product decreases from \(25=7+2\cdot9\) to \(16=7+1\cdot9\). Vacancy is exactly the loss of one multiplicative carry unit:

\[
 q^\times_{5,5}-q^\times_{8,2}=1.
\]

No zero has been placed in a cell. The unique noncentral displacement of the fibre preserving all visible data and reducing folded memory by one unit has been constructed.

The triangle \(V,E,Q\) also contains the three-channel reopening vector. To avoid concealing an orientation choice, fix the coboundary operator

\[
 d_\triangle f
 :=\bigl(f(E)-f(V),\,f(V)-f(Q),\,f(Q)-f(E)\bigr).
\]

Applied to the scalar potential,

\[
 \begin{aligned}
 d_\triangle\Delta_{\Sigma\Pi}
 &=(15-6,\,6-27,\,27-15)\\
 &=(9,-21,12)\\
 &=3(3,-7,4).
 \end{aligned}
\]

Since the three coordinates sum to zero,

\[
 d_\triangle\Delta_{\Sigma\Pi}\in
 A_2:=\ker\!\left(
 \mathbb Z^3\xrightarrow{x+y+z}\mathbb Z\right).
\]

The vector \((3,-7,4)\) is not called a root of \(A_2\): its Euclidean squared norm is \(74\). It is an integral vector of the lattice \(A_2\), and the electronic triangle realizes exactly its tritic multiple. The same vector had appeared in the global reopening of the seven tetrads of the \(\pi\) microfibre. The equality

\[
 \boxed{
 d_\triangle\Delta_{\Sigma\Pi}=3v_\Omega,
 \qquad v_\Omega=(3,-7,4)}
\]

proves that the global defect does not float above APP: it has a local representative in the three cells generating the electronic centre, conservative vacancy and rational branch \(22/7\). The factor three is the tritic lifting between the local potential and the primitive reopening coordinate.

The connection with the aperiodic crystal is equally exact. Let

\[
 \delta=\log_{10}\frac{10}{9},
 \qquad
 V_{\rm vac}=
 \begin{pmatrix}
 21&22\\
 6&7
 \end{pmatrix}.
\]

The first induced scale of the mechanical word satisfies

\[
 \left\lceil\delta^{-1}\right\rceil=22,
 \qquad
 \det V_{\rm vac}=21\cdot7-22\cdot6=15.
\]

Consequently,

\[
 \boxed{
 \Delta_{\Sigma\Pi}(E)=\det V_{\rm vac}=15.}
\]

The electronic centre and renormalization frame share more than a number: the former measures the sum--product tension of the central cell and the latter measures the oriented area surviving between the primary clock \(21/22\) and secondary clock \(6/7\). Both evaluations proceed from the same carry cocycle.

This allows the basal stage of the electronic equation to be rewritten without any opaque coefficient:

\[
 \boxed{
 \mathfrak e_0=
 \frac{\sqrt3}{4}
 \left[
  \exp\!\left(\frac{\varphi}{\pi^2}\right)
  -\left\lceil\delta^{-1}\right\rceil
   \alpha^3
 \right]
 \left(1+\det(V_{\rm vac})\Delta_4\right).}
\]

Indeed,

\[
 \left\lceil\delta^{-1}\right\rceil
 =22=\frac{396}{18},
 \qquad
 \det(V_{\rm vac})=15=\frac{270}{18}.
\]

The remaining prefactor is also internal:

\[
 s=\frac12,\qquad
 \|[P_{\Sigma,3},P_{\Pi,3}]\|
 =s\sqrt{1-s^2}=\frac{\sqrt3}{4},
 \qquad
 1-s^2=\frac34=\frac{90}{120}.
\]

Thus, the three factors of the basal stage have an explicit genealogy: noncommutativity of the APP sheets, the primary gap of aperiodic renormalization and the oriented area of its induced return. Mass selects none of them. The electronic equation is the subsequent dimensional evaluation of a previously constructed potential and survival operator.

\textbf{Refutation criterion.} The result fails if any of the three Euclidean divisions in the table is incorrect, if another noncentral point exists in \(i+j=10\) with \(t^2\equiv0\pmod9\), if the coboundary is not \(3(3,-7,4)\), or if the return frame does not have determinant \(15\).

\section{Tritic curvature of the potential and coinductive selection of futures}

The preceding APP potential is a spatial coordinate of the support. The vacancy word supplies its temporal evolution. Both layers are assembled through TRIT without improperly identifying space and time.

Let the lower mechanical word be

\[
 z_n(\beta)=
 \lfloor(n+1)\delta+\beta\rfloor
 -\lfloor n\delta+\beta\rfloor,
 \qquad z_n\in\{0,1\}.
\]

Define the accumulated polarization, increment field and tritic curvature by

\[
 \begin{aligned}
 P_N(\beta)
  &:=\sum_{n=0}^{N-1}(z_n-\delta),\\
 J_N(\beta)
  &:=P_{N+1}-P_N=z_N-\delta,\\
 \tau_N(\beta)
  &:=z_{N-1}-z_N\in\{+1,0,-1\}.
 \end{aligned}
\]

Polarization remains uniformly bounded:

\[
 P_N(\beta)
 =\lfloor N\delta+\beta\rfloor-\lfloor\beta\rfloor-N\delta,
 \qquad |P_N(\beta)|<1.
\]

More importantly, TRIT is exactly the discrete curvature of this potential:

\[
 \begin{aligned}
 J_N-J_{N-1}
 &=z_N-z_{N-1}=-\tau_N,\\
 \boxed{\tau_N}
 &\boxed{=-\Delta J_{N-1}=-\Delta^2P_{N-1}.}
 \end{aligned}
\]

This identity gives operational content to the three signs. If \(\tau_N=-1\), the code moves from \(0\) to \(1\): a vacancy opens and the fibre enables the bidirectional branching associated with the future. If \(\tau_N=+1\), the code moves from \(1\) to \(0\): the vacancy returns and its effect is incorporated into the memory of the past. If \(\tau_N=0\), the polarization maintains its local slope: the state crosses the present threshold without changing regime. Since \(0<\delta<1/2\), there are no two consecutive ones; hence the case \(\tau_N=0\) is a genuine plateau and the nonzero signs alternate.

The charge of any block is purely a boundary quantity:

\[
 \sum_{n=a}^{b}\tau_n=z_{a-1}-z_b.
\]

The interior telescopes in the tritic shadow, but does not disappear from the TPK state: the memory record preserves the ordered sequence of pulses, their quotients, their carries and the helical height. This is the difference between a local Gauss law and erasing history. TRIT determines the net boundary flux; TPK preserves the chain whose boundary has been calculated.

The relation between spatial potential and temporal curvature can be assembled in the boundary state

\[
 \mathcal V_n=
 \bigl(
 \mathscr P_{\Sigma\Pi}(i_n,j_n),
 P_n,J_n,\tau_n;
 g_n,h_n,c_n,\lambda_n
 \bigr).
\]

Here \(\mathscr P_{\Sigma\Pi}\) measures the tension between the sheets in the occupied cell; \(P_n\) accumulates the aperiodic discrepancy without extensive drift; \(J_n\) is the potential difference per step; \(\tau_n\) is its oriented curvature; and \((g,h,c,\lambda)\) preserve phase, height, carry and boundary. The transition

\[
 U_n=\operatorname{Upd}_n\circ
      \operatorname{Tra}_n\circ
      \operatorname{Sel}_n
\]

acts on the entire state: \(\operatorname{Sel}_n\) determines which curvature regime corresponds to the pulse; \(\operatorname{Tra}_n\) transports sheet, orientation and position; \(\operatorname{Upd}_n\) incorporates the coboundary into the memory record. Thus the potential difference is not added to TPK: it is the differential reading of one of its preserved coordinates.

The selection of compatible futures also admits an exact formulation. Let \(\mathcal L_m\) be the set of factors of length \(m\) of the Sturmian word, and let

\[
 d^+(u)=
 \#\{a\in\{0,1\}:ua\in\mathcal L_{m+1}\}
\]

be the number of right extensions of \(u\). Since \(p(m)=|\mathcal L_m|=m+1\),

\[
 \sum_{u\in\mathcal L_m}d^+(u)=p(m+1)
\]

and, therefore,

\[
 \boxed{
 \sum_{u\in\mathcal L_m}\bigl(d^+(u)-1\bigr)
 =p(m+1)-p(m)=1.}
\]

Every factor of a recurrent word has at least one continuation, and in a binary alphabet it has at most two. The preceding equality proves that, at each length \(m\), there is exactly one right-special factor with two possible futures; all the others have a unique future. The survival choice is therefore not an arbitrary selection among \(2^m\) words. It is a single topological branching at each scale.

Upon enrichment by the nonadic phase,

\[
 p_9(m)=9(m+1),
\]

exactly nine special states appear, one per phase origin. The tritic quotient

\[
 p_3(m)=3(m+1)
\]

organizes them into three curvature classes without destroying aperiodicity. The irrational phase \(\theta_n=\{\theta_0+n\delta\}\), together with the nonadic origin and the memory record, decides which of the two extensions occurs. Therefore:

\[
 \begin{array}{ccl}
 \text{length shadow }m
 &\longrightarrow&
 \text{a special factor with two continuations},\\
 \text{complete TPK state}
 &\longrightarrow&
 \text{a determined and invertible continuation}.
 \end{array}
\]

This is the predefinition of possible futures in mathematical form. The future is not guessed through an external probability: the reduced cylinder contains two compatible extensions, and the deep coordinate of the coinductive state selects one. Nor is the past erased, because the complete transition is bijective:

\[
 H(X_{n+1}\mid X_n)=0,
 \qquad
 H(X_n\mid X_{n+1})=0.
\]

The shadow can have positive conditional entropy precisely at the special factor; the integral state has zero net logical cost. This difference between branching in a projection and reversibility of the generator is the informational core of the crystal.

The same law contains its two renormalization clocks. Write

\[
 \delta^{-1}=21+\beta,
 \qquad 0<\beta<1,
\]

and define the induced slope

\[
 \mathcal R(\delta)
 :=1-\{\delta^{-1}\}
 =1-\beta
 =22-\delta^{-1}.
\]

The first Diophantine inequality produces the gaps \(\{21,22\}\). Since

\[
 6<\mathcal R(\delta)^{-1}<7,
\]

the induced dynamics on the short gaps produces returns \(\{6,7\}\). Renormalization does not introduce a second code: it takes the return section of the first word. Its continuant matrices

\[
 A(a)=
 \begin{pmatrix}a&1\\1&0\end{pmatrix},
 \qquad \det A(a)=-1,
\]

preserve the lattice cell, while the order of the products preserves orientation:

\[
 A(a)A(b)-A(b)A(a)
 =(a-b)
 \begin{pmatrix}0&1\\-1&0\end{pmatrix}.
\]

The basal rotation remains Abelian. Noncommutativity lies in the ordered renormalization cocycle, in the APP sheets and in curvature transports. The crystal can have zero topological entropy and simultaneously noncommutative fibred dynamics: the first property measures language growth; the second, the dependence of the result on operator order.

The three bridges are thus brought together:

\[
 \boxed{
 \begin{aligned}
 \text{APP discrepancy}
 &\xrightarrow{\ d\ }\text{boundary potential},\\
 \text{aperiodic potential}
 &\xrightarrow{\ -\Delta^2\ }\text{TRIT curvature},\\
 \text{oriented curvature}
 &\xrightarrow{\ \operatorname{Sel},\operatorname{Tra},
 \operatorname{Upd}\ }\text{TPK survival}.
 \end{aligned}}
\]

\textbf{Refutation criterion.} The construction fails if \(|P_N|\) ceases to be bounded by one, if \(\tau_N\neq-\Delta^2P_{N-1}\), if there is more than one right-special factor at some length, if the nonadic lift does not multiply complexity by nine, or if two distinct TPK states become indistinguishable after all the fields of the memory record have been preserved.

\section{Common bidirectional holonomy of the electron, action and gravity}

The action scale and gravitational scale do not interrupt the investigation of the crystal. They are two contragredient representations of the same return holonomy already acting on the electronic vacancy.

The correlated system \((\pi,\varphi,e,
\alpha)\) is first constructed as a coinductive output. On this output, the determinantal action reader produces

\[
 \Sigma_H=\varphi\alpha^{16},
 \qquad
 d_H=\left\lfloor\log_{10}\Sigma_H\right\rfloor=-34.
\]

The integer \(16\) is not the number of digits required of a measurement. It is the order of the local cokernel:

\[
 \operatorname{SNF}(H_4)=\operatorname{diag}(1,2,2,4),
 \qquad
 |\operatorname{coker}H_4|=16.
\]

The mantissa before return is

\[
 H_5=
 \frac12\sqrt{\frac{1000\alpha}{\varphi}}
 -\alpha
 +\frac9{16}\alpha^2
 -\frac59\alpha^3
 +\frac7{48}\alpha^4
 -\frac1{54}\alpha^5.
\]

The regional memory constructs

\[
 D_A=
 \exp\!\left(
 -\frac{100\pi}9\alpha\right),
\]

\[
 \mathscr C_\pi(\alpha)
 =169\alpha^6+
 \frac{D_A\alpha^7}
 {1-D_A\alpha},
\]

and the returned section is

\[
 \eta_{\rm ret}
 =H_5-\frac{90}{\pi}\mathscr C_\pi.
\]

It is useful to isolate the bidirectional action potential:

\[
 \boxed{
 \delta_{2w}
 :=H_5-\eta_{\rm ret}
 =\frac{90}{\pi}\mathscr C_\pi>0.}
\]

The two orientations belong to a single line:

\[
 \hbar_{\rm pre}=H_5\,10^{-34}\mathcal U_S,
 \qquad
 \hbar_{\rm ret}=\eta_{\rm ret}\,10^{-34}\mathcal U_S,
\]

\[
 h_a=2\pi\hbar_a.
\]

The decade fixes the common depth; the difference \(\delta_{2w}\) preserves the additive shift; and the ratio

\[
 \boxed{
 R_{\rm act}
 :=\frac{\hbar_{\rm pre}}{\hbar_{\rm ret}}
 =\frac{H_5}{\eta_{\rm ret}}
 =\left(1-\frac{\delta_{2w}}{H_5}\right)^{-1}}
\]

is the multiplicative holonomy of the outward path and return. The passage from \(\hbar\) to \(h\) preserves that ratio:

\[
 \frac{h_{\rm pre}}{h_{\rm ret}}=R_{\rm act}.
\]

The complete electronic equation uses exactly the same holonomy. With

\[
 \kappa_e=
 80-54\alpha+6\Delta_4,
\]

we have

\[
 \mathfrak e_e^\beta
 =\mathfrak e_0R_{\rm act}^{\kappa_e}.
\]

Incorporating the aperiodic factorization of the preceding section:

\[
 \boxed{
 \begin{aligned}
 \mathfrak e_e^\beta
 ={}&
 \frac{\sqrt3}{4}
 \left[
 \exp\!\left(
 \frac{\varphi}{\pi^2}\right)
 -\left\lceil\delta^{-1}\right\rceil
  \alpha^3
 \right]\\
 &\times
 \left(1+\det(V_{\rm vac})\Delta_4\right)
 R_{\rm act}^{\,80-54\alpha+6\Delta_4}.
 \end{aligned}}
\]

The basal stage measures the vacancy and sheet incompatibility; the power of \(R_{\rm act}\) transports their memory through the return. Equivalently,

\[
 \log\frac{\mathfrak e_e^\beta}{\mathfrak e_0}
 =\kappa_e
 \left(\log H_5-\log\eta_{\rm ret}\right)
 =-\kappa_e
 \log\!\left(1-\frac{\delta_{2w}}{H_5}\right).
\]

Gravity reads the same pair contragrediently. With \(t_0>0\), \(c_{\rm int}>0\) and \(\ell_0=c_{\rm int}t_0\) fixed, the clock \(108\) constructs

\[
 \omega_{108}=\frac{2\pi}{108t_0},
 \qquad
 R_{108}=\frac{54}{\pi}\ell_0,
\]

\[
 m_{108,a}=
 \frac{\hbar_a\omega_{108}}{c_{\rm int}^2}.
\]

The gravitational transducer is

\[
 G_a(\vartheta)
 =\vartheta\frac{c_{\rm int}^5t_0^2}{\hbar_a}.
\]

The closure condition

\[
 r_{g,a}(\vartheta)=R_{108}
\]

uniquely selects

\[
 \vartheta_{108}
 =\left(\frac{54}{\pi}\right)^2.
\]

Therefore,

\[
 G_a=
 \left(\frac{54}{\pi}\right)^2
 \frac{c_{\rm int}^5t_0^2}{\hbar_a},
\]

and the two sheets satisfy the synchronization theorem:

\[
 \boxed{
 R_{\rm act}
 =\frac{\hbar_{\rm pre}}{\hbar_{\rm ret}}
 =\frac{h_{\rm pre}}{h_{\rm ret}}
 =\frac{m_{108,\rm pre}}{m_{108,\rm ret}}
 =\frac{G_{\rm ret}}{G_{\rm pre}}
 =\left(
 \frac{\mathfrak e_e^\beta}{\mathfrak e_0}
 \right)^{1/\kappa_e}.}
\]

The proof consists solely in substituting the definitions. Chronological mass is linear in \(\hbar_a\); gravity is inversely linear in \(\hbar_a\); and the electronic correction is the weight-\(\kappa_e\) representation of the same ratio. Three independent correctors do not appear. A single positive element \(R_{\rm act}\) appears, with three representations:

\[
 \rho_{\rm act}(R)=R,\qquad
 \rho_{\rm grav}(R)=R^{-1},\qquad
 \rho_e(R)=R^{\kappa_e}.
\]

The action--gravity product is invariant:

\[
 \boxed{
 G_a\hbar_a
 =\vartheta_{108}c_{\rm int}^5t_0^2.}
\]

Consequently,

\[
 \ell_P^2
 =\frac{G_a\hbar_a}{c_{\rm int}^3}
 =\vartheta_{108}\ell_0^2,
\]

\[
 \boxed{
 \ell_P=R_{108}=\bar\lambda_{C,a}=r_{g,a}}
\]

for both orientations. Action preserves the sheet orientation; gravity compensates for that orientation contragrediently; the Planck area is the common invariant. This is the exact form of the assertion that something situated between past and future continuously calibrates the realization: the return changes the coordinates \(\hbar_a\), \(G_a\) and \(\mathfrak e_a\), while leaving the geometric combination \(G_a\hbar_a\) fixed.

There is also an exact filtration by order in \(\alpha\). At the basal electronic stage, the vacancy enters as

\[
 22\alpha^3.
\]

In the regional action memory,

\[
 \mathscr C_\pi
 =169\alpha^6+O(\alpha^7).
\]

Finally, the absolute depth uses \(\alpha^{16}\). Hence,

\[
 \boxed{
 3\longrightarrow6\longrightarrow16}
\]

is the hierarchy of orders of the same coordinate: local vacancy amplitude, quadratic return of that amplitude, and determinantal cokernel fixing the action depth. The step \(3\to6\) is literally a doubling of order; the step \(6\to16\) changes the operation and belongs to the determinantal reader. This distinction prevents the hierarchy from becoming a numerological sequence: each exponent has a generating operator.

The result answers the question of how a single configuration can be realized as number, constant and particle. The APP configuration provides the potential; TRIT assigns curvature and orientation to it; TPK integrates its return; the action reader determines \(\hbar\); the electronic representation determines \(\mathfrak e_e^\beta\); and the contragredient representation determines \(G\). The three outputs preserve the same \(R_{\rm act}\).

\textbf{Refutation criterion.} It suffices for one of the equalities in the chain of \(R_{\rm act}\) to fail, for \(\delta_{2w}\neq(90/\pi)\mathscr C_\pi\), for \(G_a\hbar_a\) to depend on the sheet, or for the first nonzero power of \(\mathscr C_\pi\) to differ from six.

\section{Joint depth theorem for \(34\), \(11\), \(45\) and \(46\)}

The cubic completion already contained the two aperiodic signatures

\[
 \sigma^-=(33,11,1,0),
 \qquad
 \sigma^+=(34,11,1,0),
\]

associated with the canonical order

\[
 1000=729+243+27+1.
\]

The second differs from the first by a single half-boundary unit:

\[
 \sigma^+-\sigma^-=(1,0,0,0).
\]

The essential information lies not only in the sums \(45\) and \(46\), but also in the position of each contribution. Define the cubic depth reader

\[
 \mathcal D_{\rm cub}(\sigma^-,\sigma^+)
 :=
 \left(
 \operatorname{pr}_{729}\sigma^+,\,
 \operatorname{pr}_{243}\sigma^-,\,
 |\sigma^-|_1,\,
 |\sigma^+|_1
 \right).
\]

The evaluation is

\[
 \boxed{
 \mathcal D_{\rm cub}(\sigma^-,\sigma^+)
 =(34,11,45,46).}
\]

This map does not consult \(\hbar\), \(G\), an experimental table or its digits. It reads exclusively the mechanical word on the four strata of the completion. Coordinate \(34\) is sensitive to the boundary unit; coordinate \(11\) is common to both orientations; \(45\) is the total interior charge; \(46\) is that charge after preserving the half-boundary.

We now define the decimal depth of a positive section of a dimensional line. If

\[
 x=m_x10^{-d_x}\mathcal U_X,
 \qquad 1\leq m_x<10,
\]

we set

\[
 D_{10}(x)=d_x.
\]

For two normalized sections \(x,y\), the depth of the tensor product satisfies

\[
 D_{10}(x\otimes y)
 =D_{10}(x)+D_{10}(y)-\chi(m_xm_y\geq10).
\]

In the HMT action and gravity sections,

\[
 \hbar_{\rm ret}
 =1.0545718169\ldots\times10^{-34}\mathcal U_S,
\]

\[
 G_{\rm ret}
 =6.674299659414022706367776189\ldots\times10^{-11}\mathcal U_G,
\]

and the product of mantissas belongs to the interval \((7,8)\). The gravitational coordinate corresponds to the longitudinal and radial composition of \cref{g26:eq:g-rigido-es,cr28:eq:g-evaluation}; the central value \(6.67430\times10^{-11}\) remains a metrological reference. No additional decimal carry is generated. Hence:

\[
 D_{10}(\hbar_{\rm ret})=34,\qquad
 D_{10}(G_{\rm ret})=11,
\]

\[
 \boxed{
 D_{10}(G_{\rm ret}\otimes\hbar_{\rm ret})
 =34+11=45.}
\]

The holonomy theorem of the preceding section proves that \(G_a\hbar_a\) is independent of the sheet; therefore this depth \(45\) is also a bidirectional invariant. The half-boundary unit of the upper coding then completes

\[
 \boxed{
 34+11+1=46.}
\]

The four digits of the cubic reader are thus identified with four distinct functions:

\[
\begin{array}{c|l}
34&\text{depth of the oriented action section},\\
11&\text{depth of the contragredient gravitational section},\\
45&\text{depth of the invariant }G\hbar,\\
46&\text{interior depth plus the conserved half-boundary unit}.
\end{array}
\]

The coincidence has a second internal certification. The vacancy reader had already proved

\[
 Z^-_{369}=16=|\operatorname{coker}H_4|,
\]

\[
 Z^+_{729}=34
 =-\operatorname{ord}_{10}
 \bigl(\varphi\alpha^{16}\bigr),
\]

and

\[
 Z^-_{243}=11.
\]

Thus the action decade occupies the boundary component constructed from the same determinantal index appearing in the block \(369\); gravity occupies the interior component \(243\); and their product exhausts the lower signature:

\[
 \boxed{
 Z^+_{729}+Z^-_{243}
 =|\sigma^-|_1
 =45.}
\]

The second route to \(\alpha\) preserves precisely the same integers. In its nilpotent operator,

\[
 T_{7:9}(x)
 =-\frac{x^7}{120}
  +\frac{x^8}{45}
  +\frac{2x^9}{495},
\]

with

\[
 495=45\cdot11.
\]

We may therefore write

\[
 \boxed{
 T_{7:9}(x)
 =-\frac{x^7}{120}
 +\frac{x^8}{D_{10}(G\hbar)}
 +\frac{2x^9}
 {D_{10}(G\hbar)\,D_{10}(G)}.}
\]

The equality is an exact re-expression of the integers already generated: the coefficient of order eight uses the preserved depth of the action--gravity product, and the coefficient of order nine adds the stable gravitational depth. The denominator \(46\) of the preceding extension stage corresponds, in turn, to the complete upper signature. Thus \(34\), \(11\), \(45\), \(46\) and \(495\) cease to be isolated occurrences: they form the depth spectrum of a single completion with boundary.

This theorem completes the link previously formulated only as a cardinal coincidence. The effective map is the composition

\[
 \begin{CD}
 \mathcal V_{\delta}(729,243,27,1)
 @>{(\sigma^-,\sigma^+)}>>
 \mathbb Z^4\times\mathbb Z^4\\
 @V{\operatorname{Pub}_{\rm dim}}VV
 @VV{\mathcal D_{\rm cub}}V\\
 \mathcal L_S^+\times\mathcal L_G^+
 @>{D_{10}\times D_{10}}>>
 \mathbb Z^2,
 \end{CD}
\]

where the first coordinate is determined on the action line and the second on the gravitational line; the additive character of \(D_{10}\) determines the interior charge \(45\), and the half-boundary determines \(46\). The choice of coordinates is structural: action distinguishes the two orientations \(\mathrm{pre}/\mathrm{ret}\) and therefore reads the boundary-sensitive entry; gravity compensates for that orientation and reads the stable entry.

\textbf{Refutation criterion.} The theorem fails if the lower signature is not \((33,11,1,0)\), if the upper signature does not differ by a single unit in block \(729\), if the mantissas of \(G\hbar\) generate a decimal carry, if \(G\hbar\) varies between sheets, or if the denominators \(45\), \(46\) and \(495=45\cdot11\) cease to appear in the extension \(7{:}9\).

\section{Electromagnetic polarity, neutrality and rank-five gravitational carrier}

The word ‘polarity’ can be typed through the parity of the two channel characters. Let

\[
 \mathcal E=
 \log[(1-q_+^{90})(1-q_-^{90})]
 -\log[(1-q_+^{120})(1-q_-^{120})],
\]

\[
 \mathcal B=
 \log\frac{1-q_-^{90}}{1-q_+^{90}}
 -\log\frac{1-q_-^{120}}{1-q_+^{120}}.
\]

Under the exchange involution \(\mathsf C(q_+,q_-)=(q_-,q_+)\),

\[
 \mathcal E\circ\mathsf C=\mathcal E,
 \qquad
 \mathcal B\circ\mathsf C=-\mathcal B.
\]

The character \(\mathcal E\) is common; the character \(\mathcal B\) is oriented. The constitutive responses

\[
 \widehat\mu=e^{\mathcal B+\mathcal E},
 \qquad
 \widehat\varepsilon=e^{-\mathcal B+\mathcal E}
\]

separate product and quotient exactly:

\[
 \widehat c
 =(\widehat\mu\widehat\varepsilon)^{-1/2}
 =e^{-\mathcal E},
\]

\[
 \widehat Z
 =(\widehat\mu/\widehat\varepsilon)^{1/2}
 =e^{\mathcal B}.
\]

The oriented potential difference lies in \(\mathcal B\); the common propagation capacity lies in \(\mathcal E\). Permittivity and permeability are the two conjugate recombinations. This explains the dual function of the vacancy: as an electrical opening, it changes the sign of the oriented coordinate; as magnetic transport, it preserves its circulation through \(J_{\rm comm}^2=-I\).

Neutrality then has an internal definition:

\[
 X_{\rm neu}:=\ker\mathcal B.
\]

In that subspace,

\[
 \widehat Z=1,
 \qquad
 \widehat\mu=\widehat\varepsilon=e^{\mathcal E}.
\]

Every configuration fixed by conjugation belongs to the neutral sector. Indeed, if \(\mathsf Cx=x\), then

\[
 \mathcal B(x)=\mathcal B(\mathsf Cx)
 =-\mathcal B(x),
\]

and therefore \(\mathcal B(x)=0\). The inclusion

\[
 \boxed{
 \operatorname{Fix}(\mathsf C)\subseteq X_{\rm neu}}
\]

is the algebraic content of a Majorana-type sector: a self-conjugate section cannot carry odd oriented charge. The neutral sector can be larger than the fixed locus; equality would require every configuration with \(\mathcal B=0\) to be self-conjugate. The distinction separates neutrality from self-conjugacy without leaving the operator.

Gravity receives only the common character in its scalar factor. In the internal chart

\[
 c_{\rm int}=c_{\rm ref}\widehat c
 =c_{\rm ref}e^{-\mathcal E},
\]

the transducer becomes

\[
 \boxed{
 G_a=
 \vartheta_{108}
 \frac{c_{\rm ref}^5t_0^2}{\hbar_a}
 e^{-5\mathcal E}.}
\]

The oriented character \(\mathcal B\) cancels exactly. The scalar gravitational response is even under sheet exchange and reads the common propagation coordinate five times. This fifth power is not confused with the five tensor components; the two layers combine in the operator

\[
 \mathbb G_{{\rm TT},a}(n)
 =G_a\Lambda(n)\Gamma_5P_5:
 V_{12}\longrightarrow\mathcal H_{\rm TT}(n),
\]

where

\[
 P_5:V_{12}\to V_5,
 \qquad
 \Gamma_5:V_5\to\operatorname{Sym}_0^2(\mathbb R^3).
\]

The complete composition can be written as

\[
 \boxed{
 (\mathcal E,\mathcal B,x)
 \longmapsto
 \vartheta_{108}
 \frac{c_{\rm ref}^5t_0^2}{\hbar_a}
 e^{-5\mathcal E}
 \Lambda(n)\Gamma_5P_5x.}
\]

The scalar eliminates odd dependence on \(\mathcal B\); the projector \(P_5\) preserves the five-component carrier; and \(\Gamma_5\) realizes transverse traceless polarization. This gives a precise account of the relation between electromagnetic polarity and gravitational pentapolarity. The former lies in the even/odd pair \((\mathcal E,\mathcal B)\). The latter lies in the rank of the tensor carrier. They are linked because the same state feeds both and because the gravitational scalar uses the fifth power of the common channel.

The role of symmetry is also clarified. The involution is not the object preserved first; it is the visible effect of preserving the boundary unit under two readings. The state retains

\[
 \mathcal E,\quad |\mathcal B|,\quad
 G_a\hbar_a,\quad
 \text{and the orbit of }P_5x,
\]

while conjugation reverses \(\mathcal B\). Symmetry appears as equivalence between two orientations of the same preserved potential. It is therefore a boundary property, not an interior decoration.

\textbf{Refutation criterion.} The link fails if \(\mathcal E\) is not even, if \(\mathcal B\) is not odd, if a self-conjugate section has \(\mathcal B\neq0\), if the scalar transducer depends on \(\mathcal B\), or if \(P_5\) ceases to have rank five.

\section{Integral holonomic configuration law}

The configuration assembling the preserved invariants and the operation making them dynamic is called the \textbf{integral APP--TRIT--TPK holonomic state}. A state of depth \(n\) is written

\[
 X_n^{\rm hol}=
 \bigl(
 \Sigma_n,\Pi_n;
 R_n^+,q_n^+,R_n^\times,q_n^\times;
 \tau_n,\varepsilon_n;
 g_n,h_n,c_n,\partial_n,\lambda_n,\mathfrak m_n
 \bigr).
\]

The first coordinates preserve the two APP sheets; the next, residue, quotient and carry; \((\tau,\varepsilon)\) distinguish geometric regime and temporal orientation; and \((g,h,c,\partial,\lambda,\mathfrak m)\) preserve phase, height, carry, boundary, survival and memory. The TPK transition

\[
 U_n=
 \operatorname{Upd}_n\circ
 \operatorname{Tra}_n\circ
 \operatorname{Sel}_n
\]

is an automorphism of the complete fibre on its reversibility domain. The restriction maps

\[
 r_{n+1,n}:X_{n+1}^{\rm hol}\to X_n^{\rm hol}
\]

satisfy

\[
 r_{n+1,n}U_{n+1}=U_nr_{n+1,n}.
\]

The infinite configuration is not postulated as a decimal sequence. It is the section

\[
 X_\infty^{\rm hol}
 =\varprojlim X_n^{\rm hol}.
\]

Three systems of realizations act on it:

\[
 \mathcal L_{\rm num}:
 X_\infty^{\rm hol}\to\mathbb R,
\]

\[
 \mathcal L_{\rm const}:
 X_\infty^{\rm hol}\to
 \prod_{\xi}\mathcal L_\xi^+,
\]

\[
 \mathcal L_{\rm part}:
 X_\infty^{\rm hol}\to
 \operatorname{Sect}(\mathscr A_{\rm masas}).
\]

The first determines numbers and expansions; the second, sections of dimensional lines; the third, fibres, orbits and multisections of the matter atlas. All three commute with restriction:

\[
 s_{n+1,n}^\bullet L_{\bullet,n+1}
 =L_{\bullet,n}r_{n+1,n}.
\]

Thus a number, a constant and a particle are not objects identified by having the same scalar value. They are three images of the same configuration when their readers preserve the same genealogy. The number is the Archimedean coordinate; the constant is that coordinate on a dimensional line; the particle is a persistent section preserving the route that produced the coordinate.

The single law can now be stated.

\textbf{Integral holonomic configuration theorem.} Every surviving APP seed determines, through TRIT and TPK, a section \(X_\infty^{\rm hol}\) such that:

\begin{enumerate}
\item the sum--product discrepancy defines a lifted integer potential;
\item the aperiodic word converts its temporal difference into tritic curvature;
\item TPK selects the unique extension compatible with phase and memory;
\item phase return acts through a positive holonomy on action, matter and gravity;
\item the numerical, dimensional and material realizations preserve the same section under changes in depth;
\item the dual Archimedean and exceptional projection preserves a common origin and distinct codomains;
\item the half-boundary unit is preserved at every refinement and the logical entropy of the complete state remains zero.
\end{enumerate}

The compact proof chain is

\[
 \boxed{
 \begin{aligned}
 \Omega_{\rm seed}
 &\xrightarrow{\rm APP}
 \mathscr P_{\Sigma\Pi}
 \xrightarrow{-\Delta^2}
 \tau\\
 &\xrightarrow{\rm TRIT}
 (\mathbb A_+,\mathbb A_0,\mathbb A_-)
 \xrightarrow{\rm TPK}
 X_\infty^{\rm hol}\\
 &\xrightarrow{(\mathcal L_{\rm num},
 \mathcal L_{\rm const},\mathcal L_{\rm part})}
 (\text{number},\text{constant},\text{particle}).
 \end{aligned}}
\]

The joint nonadic connection does not choose ‘the most probable future’. It chooses the unique extension compatible with the integral section when the reduced shadow admits two continuations. This operation is a predefinition: the future coordinate already belongs to the inverse limit, although it has not yet been determined at the visible depth. The net Landauer cost is zero because selection does not erase the unchosen branch from the universe of cylinders; it records in the memory register which cylinder contains the state.

The crystal is evolutionary in a precise sense. Its language does not grow exponentially, since \(p(m)=m+1\), but neither does it freeze into a period. Each depth adds exactly one factor class and each nonadic phase multiplies it by nine. The system preserves the lattice unit, boundary and determinant while modifying memory and helical height. Evolution and conservation are not opposing terms: they describe, respectively, the advance of the state and the invariance of its augmentation class.

The chain

\[
 \pi,\varphi,e
 \longrightarrow\alpha
 \longrightarrow(\mathfrak e_e^\beta,\hbar,G,\gamma_{\rm HMT},\ldots)
\]

must be read within the complete correlated generation. The first three coordinates are terminal readers of the same nonadic state; \(\alpha\) is the dodecaphasic seam of that output; and the physical constants are subsequent representations of the already closed state. The mirror \(\pi/e\), self-scaling \(\varphi\), expansion \(e+\pi\) and coincidences of the three channels are internal synchronization operations, not conventional values introduced to manufacture \(\alpha\).

The final sequence now proved is

\[
 \boxed{
 \begin{gathered}
 \text{APP potential }(6,15,27)
 \longrightarrow
 d_\triangle\Delta=3(3,-7,4),\\
 P_N\longrightarrow
 \tau_N=-\Delta^2P_{N-1}
 \longrightarrow
 \text{unique future bifurcation per scale},\\
 (21/22)\longrightarrow(6/7)
 \longrightarrow
 V_{\rm vac},\quad\det V_{\rm vac}=15,\\
 \alpha^3\longrightarrow\alpha^6
 \longrightarrow\alpha^{16}
 \longrightarrow10^{-34},\\
 R_{\rm act}\longrightarrow
 (\mathfrak e_e^\beta,\hbar,G)
 \longrightarrow G\hbar\ \text{invariant},\\
 (33,11,1,0)\leftrightarrow(34,11,1,0)
 \longrightarrow(34,11,45,46).
 \end{gathered}}
\]

This is the mathematical core that allows the work to be carried into the article: a dynamics of potential, curvature, survival, return and realization that preserves unity while converting it into number, constant and matter.

\section{Intrinsic arithmetic of the holonomic state}

The question of whether the coordinates of the holonomic state have an arithmetic of their own admits an affirmative and precise answer. The primary arithmetic object is not a field of scalars, because one history can be multiplied by another only when the terminal state of the first coincides with the initial state of the second. The natural object is therefore an \textbf{enriched path algebra} over the category of TPK transports.

Let \(\mathcal C_{\rm TPK}\) be the category whose objects are holonomic states

\[
 x=(\tau,\ell,o,h,R,q,\kappa,\partial,\gamma,M,S,b,\chi,g)
\]

and whose arrows are admissible transitions

\[
 u:x\longrightarrow y,
 \qquad
 U_t={\rm Upd}_t\circ {\rm Tra}_t\circ {\rm Sel}_t.
\]

Here \(\tau\) is the TRIT regime; \(\ell\), the APP sheet; \(o\), orientation; \(h\), helical height; \((R,q)\), residue and quotient; \(\kappa\), carry; \(\partial\), boundary; \(\gamma\), route; \(M\), the memory record; \(S\), survival; \(b\), block; \(\chi\), exceptional shadow; and \(g\), nonadic phase. Each object contains the local quadratic algebra

\[
 \mathbb A_{\tau(x)}
 =\mathbb R[J_x]/(J_x^2+\tau(x)I),
 \qquad \tau(x)\in\{+1,0,-1\}.
\]

To record the increments of carry, boundary and memory without loss, let \(\Lambda=\mathbb R[\mathcal M]\) be the algebra of the memory-record monoid. Each arrow \(u\) is associated with a symbol \([u]\). The sum

\[
 a[u]+b[v]
\]

is the formal superposition of two histories compatible with the same reading framework; the product exists causally only through concatenation. If \(u:x\to y\) and \(v:y\to z\), define

\[
 (a[v])\star(b[u])
 =a\,\rho_v(b)\,\Omega(v,u)[v\circ u],
 \tag{99.1}
\]

and set it to zero when the arrows are not composable. The transport \(\rho_v\) carries the local coefficient from the fibre of \(y\) to that of \(z\); the factor \(\Omega(v,u)\in\Lambda\) preserves the increment of carry, torsion, boundary and memory produced by concatenation. The law

\[
 \Omega(w,v)\,\Omega(w\!\circ\!v,u)
 =\rho_w\!\left(\Omega(v,u)\right)
  \Omega(w,v\!\circ\!u)
 \tag{99.2}
\]

is exactly the cocycle identity making the product associative:

\[
 ([w]\star[v])\star[u]=[w]\star([v]\star[u]).
\]

This equality is not an added convention. It is the common algebraic form of the APP carry identity, the route-concatenation cocycle and the nonsplit extension of the clock \(108\). Ordinary arithmetic forgets the factor \(\Omega\) when realizing only the value; holonomic arithmetic also preserves the procedure that produced it.

The constant states \(e_x=[{\rm id}_x]\) are local idempotents:

\[
 e_x\star e_x=e_x,
 \qquad e_y\star e_x=0\quad(x\ne y).
\]

Thus unity is not merely the scalar \(1\), but the compatible system of idempotents identifying each resolution of the same state. In the inverse limit one obtains

\[
 \mathbf 1_{\rm hol}=(e_{x_n})_{n\ge0},
 \qquad r_{n+1,n}(e_{x_{n+1}})=e_{x_n}.
 \tag{99.3}
\]

This is the dimensional unit that can be refined inward and realized outward without losing its identity. Refinement increases resolution; representation changes the codomain; neither act creates the unit anew.

Nor is noncommutativity postulated from an external structure. Two arrows \(u\) and \(v\) may be composable in one order and not in the reverse; when both orders exist, their cocycles and sheets may differ:

\[
 [v]\star[u]-[u]\star[v]
 =\Omega(v,u)[v\circ u]-\Omega(u,v)[u\circ v].
 \tag{99.4}
\]

The difference measures the effective order of accumulation and folding. The base phase system can be Abelian while the total algebra is not: noncommutativity lies in sheet composition, orientation transport and the affine memory cocycle.

\section{Complex numbers, Clifford algebras, quaternions and octonions}

The comparison with hypercomplex systems is resolved by distinguishing three levels. At the local level, TRIT already contains the three fundamental real quadratic algebras:

\[
 \begin{array}{c|c|c}
 \tau & J_\tau^2 & \mathbb A_\tau\\ \hline
 +1&-I&\mathbb C\quad\hbox{(elliptic regime)},\\
 0&0&\mathbb R[\varepsilon]/(\varepsilon^2)
       \quad\hbox{(parabolic regime)},\\
 -1&+I&\mathbb R\oplus\mathbb R
       \quad\hbox{(hyperbolic regime)}.
 \end{array}
 \tag{100.1}
\]

The complex number is therefore not the universal language of the system: it is one of its three curvature fibres. The parabolic fibre records an infinitesimal deformation without quadratic rotation; the hyperbolic fibre separates two conjugate real directions; the elliptic fibre produces internal rotation. TRIT does not choose among three geometric metaphors: it types which of these three arithmetics governs each transition.

An even stronger structure appears in the principal electronic block. Let

\[
 P=\begin{pmatrix}1&0\\0&0\end{pmatrix},\qquad
 Q=\begin{pmatrix}
 3/4&\sqrt3/4\\[2pt]\sqrt3/4&1/4
 \end{pmatrix}.
\]

Both matrices are projectors and

\[
 [P,Q]={\sqrt3\over4}
 \begin{pmatrix}0&1\\-1&0\end{pmatrix}.
\]

Define

\[
 J={4\over\sqrt3}[P,Q]
   =\begin{pmatrix}0&1\\-1&0\end{pmatrix},
 \qquad
 S=2P-I=\begin{pmatrix}1&0\\0&-1\end{pmatrix},
 \qquad K=SJ.
\]

Then

\[
 J^2=-I,\qquad S^2=I,\qquad SJ=-JS,
 \qquad K^2=I.
 \tag{100.2}
\]

Consequently, the block generates exactly

\[
 \langle I,J,S,K\rangle
 \cong {\rm Cl}_{1,1}(\mathbb R)
 \cong M_2(\mathbb R),
 \tag{100.3}
\]

the split-quaternion algebra. Furthermore,

\[
 2Q-I={1\over2}S+{\sqrt3\over2}K,
 \tag{100.4}
\]

so that the product projector occupies a direction at angle \(\pi/3\) in the plane generated by the two split involutions. The electronic prefactor \(\sqrt3/4\), the internal root of \(-1\), the fraction \(3/4=90/120\) and the separation of the APP sheets are brought together in a single exact algebraic representation.

This identification clarifies the proposed analogy. The first object to arise is not Hamilton's quaternion division algebra \(\mathbb H\), but its split form \({\rm Cl}_{1,1}\). The difference is substantive: the HMT system needs idempotents, thresholds and null directions, whereas a division algebra would eliminate precisely those vacancies. Hamilton quaternions may appear in a representation possessing two anticommuting roots of \(-I\); octonions, in a subsequent exceptional representation whose composition law reproduces the Fano plane. Neither should replace the path algebra preceding them.

Associativity of the complete TPK is preserved by the cocycle (99.2). If a subsequent exceptional projection has a nonzero associator,

\[
 [a,b,c]=(a\star b)\star c-a\star(b\star c),
\]

that associator does not describe an inconsistency of the generator: it quantifies the memory left outside the projection's codomain. This precisely separates intrinsic, associative arithmetic with a memory record from its potentially nonassociative hypercomplex shadows.

\section{Ternary synchronization and the dodecaphasic singularity of alpha}

The coordinated generation of \(\pi\), \(e\) and \(\varphi\) is organized in the difference lattice of type \(A_2\). For

\[
 N_K=(N_\pi(K),N_e(K),N_\varphi(K))\in\mathbb Z^3
\]

define the relative boundary

\[
 \partial_{A_2}(x,y,z)=(x-y,\,y-z,\,z-x).
 \tag{101.1}
\]

Its kernel is the diagonal \(\mathbb Z(1,1,1)\). Thus each triple has two distinct components: a common coordinate measuring the advance of the memory record, and a relative coordinate measuring desynchronization of the three channels. Over \(\mathbb Q\),

\[
 \mathbb Q^3=\mathbb Q(1,1,1)\oplus A_{2,\mathbb Q};
\]

over \(\mathbb Z\), reconstruction preserves a cyclic defect of order three. This is the algebraic reason why coordination is ternary rather than a binary comparison of constants.

The coincidence table follows an itinerary of walls:

\[
 H_{\pi e}\longrightarrow H_{e\varphi}longrightarrow H_{\pi e}
 \longrightarrow H_{\varphi\pi}\longrightarrow 0\longrightarrow0.
 \tag{101.2}
\]

At \(K=8\) and \(K=9\) the relative component vanishes:

\[
 N_8=59(1,1,1),\qquad N_9=43(1,1,1).
\]

However,

\[
 N_9-N_8=-16(1,1,1).
 \tag{101.3}
\]

The system returns to phase coincidence, but not to the same state: desynchronization returns to zero while diagonal memory advances sixteen units per channel. The visible equality \(\pi=e=\varphi\) expresses not immobility but relative closure with holonomic translation.

The coordinate \(\alpha\) occupies precisely the seam allowing this return to be realized without confusing it with a restart. Dodecaphasic closure takes the already synchronized relative class, preserves the diagonal memory record and lifts it to a signed section. Its two derivations are not two fits of an external quantity: they are two representations of the same holonomic class, one through integer closure with carry and the other through the simple root of the nonadic extension. Equality of the two routes expresses commutativity of the realization square.

In this way, \(\alpha\) does not merely follow \(\pi,\varphi,e\) in a list. It is the coupling coordinate transforming relative closure of the correlated system into dimensional transduction capacity. The twelve-phase wheel does not generate the three regions again: it certifies that their common return can cross a boundary while preserving the memory subsequently feeding action, the electron, the determinantal scale and the exceptional branch.

\section{Conservative evolution of unity as an algebraic theorem}

The expression ‘conservative evolution of the dimensional unit through holographic projection’ can be formulated without relying on an image. Let \(X_n\) be the holonomic states at depth \(n\), \(r_{n+1,n}:X_{n+1}\to X_n\) their restrictions, and

\[
 X_\infty^{\rm hol}=\varprojlim_n X_n
\]

their inverse limit. A dimensional unit is a compatible section \(\mathbf 1_{\rm hol}=(e_{x_n})_n\) of local idempotents. The TPK transition induces algebraic transports \(U_{t,*}\) and satisfies

\[
 U_{t,*}(e_x)=e_{U_t x},\qquad
 U_{t,*}(e_x)^2=U_{t,*}(e_x).
 \tag{102.1}
\]

There is also an augmentation morphism

\[
 \epsilon:\mathscr A_{\rm hol}\longrightarrow\mathbb R
\]

that forgets history but preserves the normalized multiplicity of the unit. Conservation and evolution are expressed simultaneously by

\[
 \epsilon(U_{t,*}a)=\epsilon(a),
 \qquad
 d_M(\Gamma_9 a)=d_M(a)+1,
 \tag{102.2}
\]

where \(d_M\) is the memory degree. The first equality states that unity is neither created nor destroyed; the second states that the state does not remain motionless. The holonomy \(\Gamma_9\) returns to the same phase and raises the memory record by one unit.

The completion

\[
 1000=729+243+27+1
\]

realizes this law finitely. The lower word distributes vacancies as \((33,11,1,0)\), while the upper reading preserves the boundary unit and produces \((34,11,1,0)\). The increment does not alter the total; it changes the location of the unit within the cubic filtration. The vector

\[
 (34,11,45,46),
 \qquad 34+11=45,
 \qquad 45+1=46,
 \tag{102.3}
\]

is therefore a signature of depth and boundary, not a juxtaposition of decimal scales.

Now let

\[
 \rho_{\mathbb R},\quad \rho_{\rm exc},\quad \rho_{\rm MD}
\]

the Archimedean, exceptional and dimensional representations of the holonomic algebra. Their holographic character consists in their being unital homomorphisms from the same source:

\[
 \rho_j(\mathbf1_{\rm hol})=I_j,
 \qquad
 \rho_j(a\star b)=\rho_j(a)\rho_j(b).
 \tag{102.4}
\]

Dual projection does not divide the origin into two universes; it preserves the same element and changes its representation. Fractal refinement operates in the opposite direction: rather than projecting to a codomain, it extends the section to a greater depth. The two movements kept separate in the preceding exposition are thus coordinated:

\[
 \boxed{
 \begin{array}{rcl}
 \text{interior refinement}&:&X_n\longleftarrow X_{n+1}
 \longleftarrow X_{n+2}\longleftarrow\cdots,\\[2pt]
 \text{holographic realization}&:&
 \mathscr A_{\rm hol}\longrightarrow
 \operatorname{End}(V_j).
 \end{array}}
 \tag{102.5}
\]

Unity is preserved in both directions: through compatibility under restriction and preservation of the identity element under representation. This dual invariance is the algebraic content of conservative evolution of the dimensional unit.

\section{The two clocks and the Sturmian helical suspension}

The temporality of the crystal arises from two mutually irreducible clocks. The first is finite and exact: the extension

\[
 0\longrightarrow C_{12}\longrightarrow C_{108}
 \longrightarrow C_9\longrightarrow0
 \tag{103.1}
\]

coordinates nonadic phase and dodecaphasic memory. It is not the direct product \(C_9\times C_{12}\): its extension class preserves the carry appearing when the phase boundary is crossed. The second clock is the irrational rotation

\[
 R_\delta:\theta\longmapsto\theta+\delta\pmod1,
 \qquad \delta=\log_{10}(10/9),
 \tag{103.2}
\]

whose mechanical coding produces the Sturmian word of vacancies.

The joint dynamics can be written as a skew product

\[
 \mathcal F(r,\theta,M)
 =\bigl(r+1,\theta+\delta,
 M+c_{108}(r)+c_{\rm St}(\theta)\bigr),
 \tag{103.3}
\]

where \(c_{108}\) is the carry of the finite clock and \(c_{\rm St}\) records crossing of the irrational window. The base \(C_9\times\mathbb T\) is Abelian; the affine extension ceases to be so when the APP sheets and the order of \({\rm Sel}\), \({\rm Tra}\) and \({\rm Upd}\) are incorporated. The helix is therefore not a circle drawn in three dimensions, but the lift of a phase rotation whose vertical coordinate is the memory record.

The ascending and descending readings are conjugate. If \(\iota\) reverses orientation and route, then

\[
 \iota\circ\mathcal F\circ\iota^{-1}=\mathcal F^{-1},
 \tag{103.4}
\]

with reversal of the oriented increments and preservation of the common responses. The mirror \(\pi/e\) is another involution: it exchanges two regional readers at the same height, whereas (103.4) reverses the temporal direction of the helix. Their distinction explains why a round-trip reading can preserve the common value without returning to the same memory.

The crystal has long-range order and zero topological entropy. For the Sturmian language,

\[
 p(m)=m+1,
\]

and the typed extensions satisfy

\[
 p_3(m)=3(m+1),\qquad p_9(m)=9(m+1).
 \tag{103.5}
\]

Phase multiplies the diversity of factors but does not change their linear growth. Hence the system is simultaneously nonperiodic, ordered and capable of unlimited extension. The gaps \(21/22\), the returns \(6/7\) and the thousand-position windows with \(45/46\) events are three observations of this same suspension, not three ad hoc rhythms.

This construction makes precise the designation aperiodic topological-temporal crystal. It is topological because information resides in route, boundary and holonomy classes; temporal because the transition distinguishes orientation, return and memory; crystalline because it possesses stable combinatorial order; and aperiodic because no nonzero translation fixes its infinite word.

\section{Vacancy, potential difference and electromagnetic polarity}

APP contains an arithmetic notion of potential before any physical chart. For each cell \((i,j)\),

\[
 \Delta_{\Sigma\Pi}(i,j)
 =ij-(i+j)
 =(R^\times-R^+)+9(q^\times-q^+).
 \tag{104.1}
\]

The first term is the visible difference of residues; the second preserves the carry difference. In the antidiagonal fibre through the centre, the electronic vacancy is the minimal loss of one multiplicative-carry unit that leaves the additive sheet, product residue and TRIT regime invariant. It is thus defined without introducing an external zero.

The potential boundary is the coboundary of (104.1). In the central triangle it produces

\[
 d_\triangle\Delta=3(3,-7,4),
 \qquad 3-7+4=0.
 \tag{104.2}
\]

The zero sum expresses conservation; the nonzero components express polarization. The potential difference is not a quantity subsequently added to the graph: it is the impossibility of making accumulation and folding coincide locally without transporting the memory record.

The representation (100.3) separates two constitutive actions. The generator \(J\) is rotational and changes sign under orientation reversal; the generator \(S\) is a sheet involution and preserves sign. Thus an even and an odd character of the same algebra produce two readers

\[
 \chi_E(SaS)=\chi_E(a),
 \qquad
 \chi_B(SaS)=-\chi_B(a),
 \tag{104.3}
\]

that realize, respectively, the electric and magnetic sectors. Electricity reads potential separation between sheets; magnetism reads the oriented circulation produced by their commutator. Their perpendicularity is not imposed through a prior Euclidean plane: it follows from \(SJ=-JS\) and \(J^2=-I\).

The regime \(\tau=0\) occupies the threshold between the two orientations. Its nilpotent element preserves a first-order perturbation and annihilates its square; it provides the algebraic type of a neutral channel without erasing the history passing through it. This explains why neutrality, vacancy and polarity can arise from the same support without being identified with one another.

The action scale and gravity extend the same holonomic character. If

\[
 R_{\rm act}={H_5\over\eta_{\rm ret}},
\]

then the common transport satisfies

\[
 R_{\rm act}
 ={\hbar_{\rm pre}\over\hbar_{\rm ret}}
 ={m_{108,{\rm pre}}\over m_{108,{\rm ret}}}
 ={G_{\rm ret}\over G_{\rm pre}}
 =\left({\mathfrak e_e^\beta\over e_0}\right)^{1/\kappa_e},
 \qquad G_a\hbar_a=\text{constant}.
 \tag{104.4}
\]

Electron, action and gravity are not three independent calibrations here: they are representations of the same return quotient at different scales and in different codomains.

\section{Exact sequence reading and the biological aperiodic crystal}

The relation to DNA can be formulated as a concrete mathematical reader, without turning a historical analogy into a premise. Encode the four bases using two binary characters:

\[
 A\mapsto(0,0),\quad G\mapsto(0,1),\quad
 T\mapsto(1,0),\quad C\mapsto(1,1),
 \tag{105.1}
\]

where the first coordinate distinguishes purine from pyrimidine and the second the weak or strong pairing type. Watson--Crick complementarity is the involution

\[
 \mathcal C_{\rm DNA}(p,h)=(1-p,h),
 \qquad \mathcal C_{\rm DNA}^2=I.
 \tag{105.2}
\]

An oriented chain is lifted to

\[
 \widetilde s_n=(p_n,h_n,n\bmod3),
\]

so that reading phase is not confused with the symbol. Route reversal and complementarity are combined in reverse complementation, which acts as an involution on the path algebra. APP can evaluate both character accumulation and concatenation in each window; TRIT types the second difference of the resulting potential; TPK preserves phase, boundary, returns, carry and return words.

This representation allows one to ask exactly whether a sequence region belongs to the minimal aperiodic class: for its factors of length \(m\), one calculates complexity \(p_s(m)\), the number of right extensions and their return words. The minimal HMT signature is

\[
 p_s(m)=m+1,
 \qquad
 \sum_{u\in L_m}(d^+(u)-1)=1,
 \tag{105.3}
\]

or, when a finite phase fibre is preserved, an extension \(q(m+1)\). The equality selects a unique combinatorial branching per length; TPK determines the compatible continuation without erasing the unrealized alternative from the cylinder space.

The relevance of Schrödinger's expression ‘aperiodic crystal’ thus becomes structural. A periodic crystal would repeat a block and lose informational capacity; a random chain would have exponential complexity and lose long-range order. The Sturmian regime occupies the exact minimum between the two: it adds a single new class at each scale. The TRIT topological-temporal unit can read this aperiodicity because it separates curvature, orientation and phase without destroying the sequence it classifies.

The criterion is directly refutable: it fails if reverse complementation does not intertwine the route involution, if observed complexity does not belong to the declared linear class, if several special words appear where the model requires one, or if compatible extension requires erasing memory. Biology does not select the HMT generator; it receives from it a sequence reader capable of distinguishing repetition, minimal aperiodicity and disorder.

\section{Pentadicity, the Witt star and gravitational polarizations}

The recurrent appearance of the number five must be typed through two distinct representations. The five regions of \(\pi\) form a regional fibre on which the symmetrizer acts

\[
 S_{\pi,5}={1\over5}\mathbf1\mathbf1^{\mathsf T},
 \qquad \operatorname{rank}S_{\pi,5}=1.
 \tag{106.1}
\]

Its function is to extract the common component of the pentafibre. In the gravitational sector, by contrast, the dodecaphasic decomposition

\[
 V_{12}\cong\mathbf1\oplus V_3\oplus V_3'\oplus V_5
 \tag{106.2}
\]

contains a projector \(P_{G,5}:V_{12}\to V_5\) of rank five. The two matrices are not the same operator: one has rank one on the regional fibre and the other rank five on the dodecaphasic representation. Their correct relation is that both are images of the same holonomic arithmetic:

\[
 (\rho_{\pi,5},\rho_{G,5}):
 \mathscr A_{\rm hol}\longrightarrow
 \operatorname{End}(\mathbb Q^5)
 \times\operatorname{End}(V_{12}).
 \tag{106.3}
\]

The Witt star organizes exceptional pentadic incidence; the module \(V_5\) organizes the gravitational tensor carrier. Coordination occurs in the APP--TRIT--TPK source, not by identifying their codomains.

The gravitational realization continues through

\[
 V_{12}\xrightarrow{P_{G,5}}V_5
 \xrightarrow{\Gamma_5}{\rm STF}_3
 \xrightarrow{\Pi_{\rm TT}}V_{\rm TT},
 \tag{106.4}
\]

with dimensions \(12\to5\to5\to2\). The five components constitute the structural carrier before propagation; the two transverse traceless components are the propagating output. There is no contradiction between original pentapolarity and radiative bipolarity: they belong to distinct stages of the same transduction.

In the joint discrete structure of the continuum, pentapolarity is expressed by the inseparable vector

\[
 \mathcal G_5=(\mathcal G_{\rm sp},\mathcal G_{\rm sol},
 \mathcal G_{\rm gau},\mathcal G_{\rm coh},\mathcal G_{\rm det}).
 \tag{106.5}
\]

Its five coordinates are not five theories added afterwards, but five inherent aspects of the same state: spectral, solenoidal, gauge, coherent and determinantal. The gravitational tensor

\[
 \mathbb G_{\rm TT}
 =G\,\Lambda\,\Gamma_5\,P_{G,5}
 \tag{106.6}
\]

jointly transports these aspects to the physical sector. This formula explains why regional pentadicity, exceptional structure and the five polarizations can resonate without being reduced to a coincidence of cardinalities.

\section{Variance duality between histories and readers}

Retrospective reading reveals a categorical structure uniting inward fractalization and outward holographic realization. Let \((X_n)\) be the set of holonomic states at depth (n), and let

\[
 r_{n+1,n}:X_{n+1}\longrightarrow X_n
 \tag{108.1}
\]

be the restriction removing the last layer while preserving the prefix, sheet, accumulated carry and corresponding boundary. The complete history is a state of the inverse limit

\[
 \boxed{X_\infty=\varprojlim_n X_n.}
 \tag{108.2}
\]

The direction of observations is opposite. If (f) is a cylindrical reader defined on \((X_n)\), its extension to \((X_{n+1})\) is

\[
 r_{n+1,n}^*f=f\circ r_{n+1,n}.
 \tag{108.3}
\]

Therefore the reader algebras form a direct system:

\[
 \boxed{
 \mathfrak H_{\rm cyl}
 =\varinjlim_n(\mathfrak H_n,r_{n+1,n}^*).}
 \tag{108.4}
\]

This dual variance is the exact mathematical formulation of the two movements described by the author's intuition. Deepening inward means specifying a compatible section of the inverse system. Expanding outward means extending the repertoire of readers through the direct system. Neither direction replaces the other: one constructs the history and the other constructs the capacity to realize it.

Let \((e_x)\) be the indicator idempotent of the coarse cylinder \((x\in X_n)\). The lift of that cylinder is

\[
 r^*e_x=\sum_{y:r(y)=x}e_y.
 \tag{108.5}
\]

The summands are orthogonal, so \((r^*e_x)\) remains idempotent. Furthermore,

\[
 r^*\!\left(\sum_{x\in X_n}e_x\right)
 =\sum_{y\in X_{n+1}}e_y.
 \tag{108.6}
\]

Global unity is preserved while the number of refinements increases. An individual trajectory in turn determines a compatible system \(((e_{x_n})_n)\). The former is the algebra identity; the latter is the genealogical unit of a history. This distinction makes conservative evolution precise: the total possibility space preserves its unity and each route preserves its membership through compatible restrictions.

A continuous reader (L) of the deep state is equivalent to a system of finite readers \((L_n)\) satisfying

\[
 L_{n+1}\circ r_{n+1,n}^*=L_n.
 \tag{108.7}
\]

Consequently, each digit or coordinate determined at finite depth is predefined by a cylinder; increasing depth does not recalibrate that result, but determines new compatible cylinders. Predefinition is the naturality of (108.7), not the advance insertion of an experimental value.

\textbf{Refutation criterion.} The construction fails if a restriction does not recover the prefix, if the \emph{pullback} is not multiplicative, if unity is not preserved, or if a deep reader modifies a coordinate already fixed in an earlier cylinder.

\section{Universal presentation of the holonomic algebra}

Let \((mathcal C_n)\) be the category of states at depth (n) and admissible histories generated by

\[
 U_t={\rm Upd}_t\circ{\rm Tra}_t\circ{\rm Sel}_t.
 \tag{109.1}
\]

For each object (x), let \((e_x)\) be its local identity and \((J_x)\) the generator of the TRIT fibre:

\[
 J_x^2=-\tau(x)e_x,
 \qquad \tau(x)\in\{+1,0,-1\}.
 \tag{109.2}
\]

For each arrow \((u:x\to y)\), let \((T_u)\) be the transport. The memory record defines a central multiplier (Omega(v,u)) on each composable pair. The finite presentation is

\[
 \boxed{
 \mathfrak H_n
 =\mathbb R\langle e_x,J_x,T_u\rangle/\mathcal R_{\rm HMT},}
 \tag{109.3}
\]

where

\[
 \begin{gathered}
 e_xe_y=\delta_{xy}e_x,
 \qquad
 T_u=e_{t(u)}T_ue_{s(u)},\\
 T_u a=\rho_u(a)T_u,
 \qquad
 T_vT_u=\Omega(v,u)T_{v\circ u}.
 \end{gathered}
 \tag{109.4}
\]

The product is declared zero when the arrows are not composable. Associativity is equivalent to

\[
 \rho_w(\Omega(v,u))\Omega(w,v\circ u)
 =\Omega(w,v)\Omega(w\circ v,u).
 \tag{109.5}
\]

Indeed, the two products \((((T_wT_v)T_u))\) and \((T_w(T_vT_u))\) differ only in the two sides of (109.5). The cocycle identity proves that every history has a memory independent of the chosen parenthesization.

The sum in (109.3) superposes alternative histories. Its product concatenates causally composable histories. This distinction turns the APP sum--product difference into dynamic arithmetic: the sum does not erase the alternative and the product does not erase provenance.

The nonadic connection appears in the elementary section of the quotient. For \((a,b\in\{0,\ldots,8\})\),

\[
 c_9(a,b)=\left\lfloor{a+b\over9}\right\rfloor
 \tag{109.6}
\]

satisfies

\[
 c_9(a,b)+c_9(a+b\bmod9,c)
 =c_9(b,c)+c_9(a,b+c\bmod9).
 \tag{109.7}
\]

If (X) represents a step and (lambda) the memory of one turn,

\[
 X^9=\lambda I,
 \qquad
 X^aX^b=\lambda^{c_9(a,b)}X^{a+b\bmod9}.
 \tag{109.8}
\]

The helix is condensed in (109.8): the exponent returns; the coefficient remembers how many times the boundary was crossed.

\section{Universal property and spectrum of realizations}

Let \(({E_x})\) be a system of modules, let \((pi_x:A_x\to\operatorname{End}(E_x))\) be representations of the TRIT fibres, and let \((V_u:E_x\to E_y)\) be transports satisfying

\[
 V_u\pi_x(a)=\pi_y(\rho_u(a))V_u,
 \tag{110.1}
\]

\[
 V_vV_u=\pi_z(\Omega(v,u))V_{v\circ u}.
 \tag{110.2}
\]

The assignment of generators defines a homomorphism from the free algebra. Equations (110.1)--(110.2) make that homomorphism annihilate the ideal \((mathcal R_{\rm HMT})\). There is then a unique representation

\[
 \Pi:\mathfrak H_n\longrightarrow
 \operatorname{End}\!\left(\bigoplus_xE_x\right)
 \tag{110.3}
\]

realizing the prescribed data. Uniqueness follows because \((e_x,J_x,T_u)\) generate the algebra.

This theorem provides the principle of ultrasimplification. Physics, number theory, the continuum and biology need not add generators to the kernel when their readers satisfy the same relations. Each is a representation of (109.3) in a different codomain.

Not all readers are scalar characters. A character \((mathfrak H\to\mathbb R)\) would annihilate commutators and could not preserve the electronic block. We therefore define the typed spectrum

\[
 \operatorname{Read}(\mathfrak H)
 =\coprod_B\operatorname{Rep}_{\rm comp}(\mathfrak H,B).
 \tag{110.4}
\]

A scalar output takes the form

\[
 c=\sigma_B(\rho(x_\infty)),
 \tag{110.5}
\]

where \(\rho\) preserves genealogy and \(\sigma_B\) is the invariant specific to the codomain. Five concepts that previously seemed heterogeneous are thus typed:

\[
 \begin{array}{c|l}
 \text{number}&\text{coordinate of an Archimedean representation},\\
 \text{constant}&\text{stable invariant of return, scale or depth},\\
 \text{particle}&\text{persistent section of a module},\\
 \text{field}&\text{representation transporting sections},\\
 \text{symmetry}&\text{boundary stabilizer and preserved relations}.
 \end{array}
 \tag{110.6}
\]

\section{Exponential identity of the three curvatures}

We introduce the total algebra

\[
 \mathbb A_{\rm tri}=A_+\oplus A_0\oplus A_-
 \tag{111.1}
\]

with central idempotents \((p_+,p_0,p_-)\) and the generator

\[
 \mathbb J=p_+J_++p_0J_0+p_-J_-.
 \tag{111.2}
\]

The exponential is defined within the algebra by

\[
 \operatorname{Exp}_\star(X)
 =\sum_{n\ge0}{X^{\star n}\over n!}.
 \tag{111.3}
\]

The three quadratic relations split the series into

\[
 \boxed{
 \begin{aligned}
 \operatorname{Exp}_\star(t\mathbb J)
 ={}&p_+(\cos t+J_+\sin t)\\
 &+p_0(1+tJ_0)\\
 &+p_-(\cosh t+J_-\sinh t).
 \end{aligned}}
 \tag{111.4}
\]

The first row is rotation; the second is infinitesimal threshold transport; the third is dilation. A single power law contains the three curvatures.

Euler's identity occupies the elliptic face. In the closure produced by \((pi_{\rm HMT})\),

\[
 \boxed{
 p_+\operatorname{Exp}_\star(\pi\mathbb J)+p_+=0.}
 \tag{111.5}
\]

The complex projection preserves (111.5) and forgets the parabolic and hyperbolic components. HMT jointly preserves closure, jet and scale. Euler's identity thus appears as an exact face of the tritic identity, not as an external model of the generator.

The number \((e)\) has a different function:

\[
 \operatorname{Exp}_\star(I)=e\,I.
 \tag{111.6}
\]

The equation defines the function converting additive transport into multiplicative transport. The nonadic generation of its digits retains its own certificate; (111.6) identifies its role in the algebra.

Self-scaling \(\varphi\) occupies the hyperbolic face. For

\[
 F_{\rm av}=\begin{pmatrix}0&1\\1&1\end{pmatrix},
 \qquad
 H_\varphi={2F_{\rm av}^2-3I\over\sqrt5},
 \tag{111.7}
\]

we verify

\[
 H_\varphi^2=I,
 \tag{111.8}
\]

and

\[
 \boxed{
 F_{\rm av}^2
 =\operatorname{Exp}_\star(2\log\varphi\,H_\varphi).}
 \tag{111.9}
\]

Thus \(\pi\) measures elliptic closure, \(e\) realizes the exponential step, and \(\varphi\) measures stable hyperbolic dilation. Equalizations of the three channels are encounters of these functions within the same orbit, not universal equalities between their real values.

\section{\(\alpha\) as dodecaphasic transgression}

The nonsplit extension produces the cocycle \(c_9\) of (109.6). A phase turn determines an element of the kernel \(C_{12}\), and the signed word \(K\) preserves how its sectors were crossed. Define the seam functional

\[
 \mathcal T_\alpha:
 Z^2(C_9;C_{12})\times X_{12}^{\rm sgn}
 \longrightarrow\mathcal L_\alpha.
 \tag{112.1}
\]

The determined coordinate is

\[
 \boxed{
 \alpha=\mathcal T_\alpha(c_9,K).}
 \tag{112.2}
\]

The formula does not identify the number \(\alpha\) with a cohomology class. It identifies its type: it is the scalar obtained by evaluating the splitting obstruction together with dodecaphasic memory.

The first route realizes (112.2) through

\[
 \iota(P)+\iota(E)-\iota(\Phi)-\iota(K)=\iota(A).
 \tag{112.3}
\]

The second route realizes the same gluing on the aperiodic rotation:

\[
 P_9(\alpha)=\delta,
 \qquad
 \delta=\log_{10}{10\over9}.
 \tag{112.4}
\]

Their compatibility is expressed in the fibre product

\[
 \mathfrak A_\alpha
 =\{(x_{12},x):
 d_m(\Gamma_{12}^\alpha x_{12})=d_m(x),
 \ P_9(x)=\delta\}.
 \tag{112.5}
\]

A kernel matrix of the correlated system can now be formulated:

\[
 \boxed{
 \begin{pmatrix}
 \operatorname{Exp}_\star(\pi J_+)&0&0&0\\
 0&\operatorname{Exp}_\star(I)&0&0\\
 0&0&\operatorname{Exp}_\star(2\log\varphi H_\varphi)&0\\
 0&0&0&\operatorname{Exp}_\star(P_9(\alpha)J_0)
 \end{pmatrix}
 =
 \begin{pmatrix}
 -I&0&0&0\\
 0&e\,I&0&0\\
 0&0&F_{\rm av}^2&0\\
 0&0&0&I+\delta J_0
 \end{pmatrix}.}
 \tag{112.6}
\]

The matrix does not linearize the genealogy. It displays four functions of a single algebra: closure, exponentiation, self-scaling and seam. Signed memory remains in (112.3)--(112.5).

In the lattice \(A_2\), the same seam appears as

\[
 \partial_{A_2}N_8=\partial_{A_2}N_9=0,
 \qquad
 N_9-N_8=-16\mathbf1.
 \tag{112.7}
\]

\(\alpha\) determines closure of the relative part without erasing diagonal translation. This is the exact meaning of its singularity: it is the coordinate distinguishing synchronization from restart.

\section{The functional spectrum of constants}

The uniform dodecaphasic reader also determines a sequence of first roots of critical return. Its separation ratio is proved to converge to the Feigenbaum constant, preserving the original selection, orientation and memory of the construction. \eqref{hmtfg:articulacion:limite}.


The preceding structure reorganizes the constants already derived according to their function in the algebra, not their historical affiliation with mathematics or physics.

Catalan's constant is the second-order moment of the elliptic character:

\[
 G_{\rm Cat}=\operatorname{Tr}(N^{-2}Q_4)
 =L(2,\chi_{-4}).
 \tag{113.1}
\]

Its functional setting preserves the entire system \(\mathcal C_2(s)\), not only the endpoint \(s=1\). In \cref{cr28:thm:catalan-recovery} it is proved that the constitutive response \(f(s)=(1-s^3)/(1-s^4)\), together with the condition \(\mathcal C_2(0)=0\), uniquely determines the moment through \(\mathcal D^2\mathcal C_2=s(1+s)f(s)/(1+s+s^2)\). Cubic inversion in turn recovers the two oriented arguments.

The quarter-turn character and the triadic character are brought together in

\[
 \chi_{12}=\chi_{-3}\chi_{-4}.
 \tag{113.2}
\]

Feigenbaum occupies the nonlinear reader. The dodecaphasic wheel generates

\[
 u_0(x)=1-{1\over12}\sum_{m=1}^{12}
 \cos\!\left({2(3m-1)\over\sqrt{899}}x\right),
 \tag{113.3}
\]

with

\[
 899=30^2-1=29\cdot31,
 \qquad
 M_{30}=\begin{pmatrix}30&899\\1&30\end{pmatrix},
 \qquad \det M_{30}=1.
 \tag{113.4}
\]

The constructed system is analytic, even, unimodal, quadratically critical and has negative Schwarzian derivative on the proved domain. The doubling renormalizer reads the accumulation rate and spatial rescaling. \(\alpha_{\rm F}\) is a coordinate of that nonlinear reader; \(\alpha\) is the seam of (112.2). The distinct subscripts express a difference of operation.

Boltzmann occupies the transducer between the thermal and energy lines:

\[
 k_B:\mathcal L_\Theta\longrightarrow\mathcal L_E,
 \tag{113.5}
\]

\[
 k_B\Theta_{\rm clk}\log3
 =\hbar_{\rm int}{2\pi\over108t_0}.
 \tag{113.6}
\]

\(\log 3\) quantifies a tritic unit; \(\log\varphi\) quantifies hyperbolic self-scaling. They are not interchangeable. For complete transport,

\[
 H(X\mid U(X))=0,
 \qquad \inf Q_L=0.
 \tag{113.7}
\]

For a quotient erasing one trit,

\[
 Q_{\rm reset}\ge k_B\Theta\log3.
 \tag{113.8}
\]

The cost belongs to the reducing projection, not to the advance of holonomy.

Barbero--Immirzi occupies the dual of the incidence module \(90/120\). The equation

\[
 {a\over270}-{b\over360}={1\over24}
 \tag{113.9}
\]

selects the minimal integer covector \((a,-b)=(12,-1)\), and therefore

\[
 K_{6\to8}=12\Delta S_{90}-\Delta S_{120}.
 \tag{113.10}
\]

\(\alpha\) closes the clocks; \(\gamma_{\rm BI}\) evaluates the seam on the area--volume defect. This functional difference explains why both objects are central without performing the same operation.

\section{Symmetry, pentadicity, continuum and sequences}

The definition of symmetry compatible with the kernel is

\[
 \operatorname{Sym}(\rho,b_\partial)
 =\{g\in\operatorname{Aut}(\rho(\mathfrak H)):
 g\rho(\mathcal R_{\rm HMT})=\rho(\mathcal R_{\rm HMT}),
 \ gb_\partial=b_\partial\}.
 \tag{114.1}
\]

Symmetries are not the primary preserved quantity. They are the automorphisms remaining available after preservation of relations and boundary is required. The exceptional branch realizes (114.1) through Witt--Golay--Leech--\(V^\natural\)--Monster incidence.

Pentadicity has two representations:

\[
 S_{\pi,5}={1\over5}\mathbf1\mathbf1^{\mathsf T},
 \qquad \operatorname{rank}S_{\pi,5}=1,
 \tag{114.2}
\]

and

\[
 P_{G,5}:V_{12}\to V_5,
 \qquad \operatorname{rank}P_{G,5}=5.
 \tag{114.3}
\]

Both receive the same holonomic source, but one extracts the common component of five regions and the other preserves the five-dimensional gravitational carrier. The chain

\[
 V_{12}\to V_5\to {\rm STF}_3\to V_{\rm TT},
 \qquad 12\to5\to5\to2,
 \tag{114.4}
\]

proves how five structural modes determine two propagating modes.

Dual projection of the continuum is expressed as

\[
 \rho_{\rm arq}:\mathfrak H_{\rm cyl}\to\mathcal A_{\mathbb R},
 \qquad
 \rho_{\rm exc}:\mathfrak H_{\rm cyl}\to\mathcal A_{\rm exc}.
 \tag{114.5}
\]

The first contracts compatible cylinders; the second preserves boundary incidence. The five spectral, solenoidal, gauge, coherent and determinantal aspects belong to the same domain of (114.5). \(\mathbb R\) is the output of a commutative representation; holonomic arithmetic additionally preserves commutators and the memory record invisible to that output.

Sequence reading constitutes another representation. With

\[
 A\mapsto(0,0),\quad G\mapsto(0,1),\quad
 T\mapsto(1,0),\quad C\mapsto(1,1),
 \tag{114.6}
\]

and phase \((n\bmod3)\), reverse complementation defines an involution on the category of words. APP evaluates accumulation and concatenation; TRIT types curvature; TPK preserves returns and boundary. The Sturmian condition

\[
 p_s(m)=m+1
 \tag{114.7}
\]

provides the exact threshold between periodicity and exponential complexity. Thus the idea of an aperiodic crystal applied to biological sequences becomes a falsifiable reader of factors, extensions and return words.

\section{Kernel theorem of HMT holonomic arithmetic}

The preceding sections allow genealogy to be condensed without reducing it. Define

 \[
 \boxed{
 \mathfrak H_{\rm HMT}^{\rm alg}
 =\left[\varinjlim_n
 \left(
 \mathbb R\langle e_x,J_x,T_u\rangle/
 \mathcal R_{\rm HMT}
 \right)\right]
 \rtimes_{\Gamma_9}\mathbb N.}
 \tag{115.1}
 \]

In the reversible sector, the last factor is \(\mathbb Z\). The state space is \(X_\infty=\varprojlim X_n\). For a compatible bounded representation \(\rho\), the completion \(\mathfrak H_{\rm HMT}^{\rho}
=\overline{\rho(\mathfrak H_{\rm HMT}^{\rm alg})}^{\|\cdot\|}\) belongs to the reader; the universal generator does not presuppose a norm.

\begin{quote}
\textbf{Kernel theorem.} Dual APP evaluation, TRIT trigeometric orientation and TPK dynamics generate an associative noncommutative arithmetic of histories, twisted by the carry cocycle and graded by phase and memory. Refinement preserves unity, holonomy returns the phase while advancing memory, and the typed spectrum of representations determines the correlated system \((\pi,\varphi,e,\alpha)\), the joint discrete structure of the continuum, the electron, action and gravity scales, geometric parameters, derived constants and boundary symmetries.
\end{quote}

The proof is obtained by concatenating five established facts:

\begin{enumerate}
\item the APP equations preserve both sheets, residue and quotient;
\item \(J_\tau^2=-\tau I\) preserves the three curvatures within a single TRIT;
\item the cocycle identity proves associativity of the history product;
\item the nonsplit extension proves that phase return and state return are distinct;
\item the universal property proves that every compatible reader factors through (115.1).
\end{enumerate}

The first consequence is that \(\mathbb R\) is not the fundamental support. It is a terminal representation of a structure preserving more information than its real coordinate. The second is that a constant is not a stored numeral: it is an invariant of a reader. The third is that a particle is not a number: it is a persistent section of a module. The fourth is that a symmetry is not postulated: it appears as the stabilizer of a preserved boundary.

The shortest form of conservative evolution is

\[
 \boxed{
 \underbrace{\varprojlim X_n}_{\text{history that deepens}}
 \quad\longleftrightarrow\quad
 \underbrace{\varinjlim\mathfrak H_n}_{\text{expanding readers}}.}
 \tag{115.2}
\]

The helical aperiodic crystal coordinates both limits. The trigeometric identity (111.4) provides its local calculus; \(\alpha\) provides the memory seam; the block \({\rm Cl}_{1,1}\) provides its electronic corner; the invariant \((G\hbar)\) provides its dimensional conservation; Barbero--Immirzi provides its boundary covector; and dual projection provides its Archimedean and exceptional outputs. This is the minimal structure from which the standalone article can be developed without expository dependence on the (2\,099) pages and without losing any essential operator.

\textbf{Refutation criterion.} The theorem is refuted if the cocycle identity fails, if refinement ceases to be unital, if holonomy restarts memory, if a TRIT fibre disappears, if a realization fails to satisfy the relations it claims to represent, or if an output uses its own conventional value to select the history that was to produce it.

\section{Triadic irreducibility and potential difference in sequence readers}

The APP--TRIT--TPK composition can be proved irreducible through three forgetful quotients. Let \(q_{\rm APP}\) be the quotient identifying the additive and multiplicative sheets. Then

\[
q_{\rm APP}(\Delta_{\Sigma\Pi})=0,
\qquad
q_{\rm APP}([P_\Sigma,P_\Pi])=0.
\tag{116.1}
\]

Let \(q_{\rm TRIT}\) be the quotient identifying the three regime idempotents. The triple of squares is then lost

\[
(J_+^2,J_0^2,J_-^2)=(-I,0,I).
\tag{116.2}
\]

Let \(q_{\rm TPK}\) be the quotient trivializing the carry cocycle and memory grading. In that quotient,

\[
q_{\rm TPK}\!\left(d_M(\Gamma_9x)-d_M(x)\right)=0,
\tag{116.3}
\]

although in the complete holonomic state the difference equals \(1\). The three invariants (116.1)--(116.3) have distinct types, and none can be reconstructed from the other two after its quotient is applied. APP provides discrepancy and sheet order; TRIT, curvature signature; TPK, composition, holonomy and memory. This is the exact algebraic form of the ternary irreducibility of the kernel.

The same separation clarifies the application to sequences. In the alphabet \(\{A,G,T,C\}\), define

\[
\chi_p(A,G,T,C)=(0,0,1,1),
\qquad
\chi_h(A,G,T,C)=(0,1,0,1),
\tag{116.4}
\]

and the complement involution \(\mathcal C_{\rm DNA}(A,G,T,C)=(T,C,A,G)\). One verifies

\[
\chi_p(\mathcal C_{\rm DNA}b)=1-\chi_p(b),
\qquad
\chi_h(\mathcal C_{\rm DNA}b)=\chi_h(b).
\tag{116.5}
\]

Consequently,

\[
\mathcal V_n=\chi_p(b_{n+1})-\chi_p(b_n)
\tag{116.6}
\]

is an oriented symbolic potential difference, whereas \(\chi_h\) preserves the transverse bond class. The lift

\[
b_n\longmapsto
\bigl(\chi_p(b_n),\chi_h(b_n),n\bmod3,\lambda_n,c_n,m_n\bigr)
\tag{116.7}
\]

adds TRIT regime, boundary, carry and TPK memory. This formulates the structural reason why polarity, potential and aperiodic word can be read within the same arithmetic without turning biology into the generator of the formalism.

\textbf{Refutation criterion.} Irreducibility fails if one of the three invariants survives the quotient that should erase it, or can be reconstructed from the other two without reintroducing the eliminated datum. The sequence reader fails if complementation does not reverse \(\chi_p\), does not preserve \(\chi_h\), or if the lift does not intertwine ternary phase with TPK dynamics.

\section{The electronic vacancy as a Witt null pair}

The arithmetic of the electronic corner allows the three regimes to be completed locally without introducing another algebra. With

\[
S=2P-I,\qquad
J=\frac{[P,Q]}{\sqrt3/4},
\tag{117.1}
\]

we have

\[
S^2=I,\qquad J^2=-I,\qquad SJ=-JS.
\tag{117.2}
\]

The combinations

\[
n_+=\frac{S+J}{2},\qquad
n_-=\frac{S-J}{2}
\tag{117.3}
\]

satisfy

\[
n_+^2=n_-^2=0,\qquad
n_+n_-+n_-n_+=I.
\tag{117.4}
\]

Therefore,

\[
\{J,n_\pm,S\}
\quad\longmapsto\quad
\{\text{elliptic},\text{parabolic},\text{hyperbolic}\}
\tag{117.5}
\]

is a complete local realization of TRIT typing within \({\rm Cl}_{1,1}\cong M_2(\mathbb R)\). The vacancy is characterized by two compatible descriptions: in APP it is the unit loss of the product quotient that preserves the additive sheet, product residue and regime; in the Clifford corner it is a parabolic null direction. The pair \((n_+,n_-)\) is the local Witt pair of the two orientations. This decomposition provides the precise algebraic domain from which subsequent Paley--Witt incidence can realize the same genealogy in its own combinatorial module.

The assertion is stronger than the quaternion analogy. Hamilton quaternions have three generators squaring to \(-I\) and form a division algebra. The block (117.2) is split: it contains zero divisors, idempotents and null directions, precisely the objects required by sheet, vacancy and threshold. Octonions enlarge dimension at the cost of associativity; holonomic arithmetic maintains associativity through the cocycle and enlarges the structure through routes, memory and refinement.

\textbf{Refutation criterion.} It suffices for one of the four identities (117.2)--(117.4) to fail, for the vacancy not to preserve its declared APP data, or for TPK transport not to preserve the regime idempotent.

\section{Weyl relations of the finite and irrational clocks}

Let \(T\) be the orbit step, let \(\delta=\log_{10}(10/9)\), let \(\zeta_9\) and \(\omega_\delta\) be the internal characters of nonadic phase and capacity, and define on \(\lvert t,m\rangle\)

\[
V_9\lvert t,m\rangle=\zeta_9^t\lvert t,m\rangle,
\qquad
W_\delta\lvert t,m\rangle=\omega_\delta^t\lvert t,m\rangle.
\tag{118.1}
\]

Temporal advance satisfies

\[
V_9T=\zeta_9TV_9,\qquad
W_\delta T=\omega_\delta TW_\delta.
\tag{118.2}
\]

The turn \(\Gamma_9=T^9\) produces

\[
V_9\Gamma_9=\Gamma_9V_9,\qquad
W_\delta\Gamma_9=\omega_\delta^9\Gamma_9W_\delta.
\tag{118.3}
\]

The two clocks are thus distinguished algebraically. The first has period nine; the second preserves an irrational phase. The mechanical vacancy word is the threshold coding of the second clock's orbit. If \(D_M\) is the memory-degree operator, the TPK lift satisfies

\[
[D_M,\Gamma_9]=\Gamma_9.
\tag{118.4}
\]

The subsequent elliptic representation recognizes these characters through

\[
\rho_+(\zeta_9)
=\operatorname{Exp}_+\!\left(\frac{2\pi}9J_+\right),
\qquad
\rho_+(\omega_\delta)
=\operatorname{Exp}_+\!\left(2\pi\delta J_+\right).
\tag{118.5}
\]

The characters, not their recognized complex expressions, are the clock's operational data.

Equation (118.4) states exactly that \(\Gamma_9\) is a raising operator: phase returns and the state ascends one unit along the helix. The translation action is Abelian; the crossed product of that action with the observables \(V_9,W_\delta,D_M\) is noncommutative by (118.2)--(118.4). The temporal structure therefore does not require choosing between ‘finite clock’ and ‘aperiodic crystal’: both are conjugate characters of a single skew system.

Zero topological entropy follows from irrational rotation and complexity \(p(m)=m+1\); memory growth follows from (118.4). Zero entropy and an infinite history coexist because the latter measures reconstructible depth whereas the former measures exponential growth of factors.

\textbf{Refutation criterion.} The model fails if \(V_9\) does not commute with \(\Gamma_9\), if \(W_\delta\) acquires a finite period, if the memory commutator differs from \(\Gamma_9\), or if the threshold word ceases to have Sturmian complexity.

\section{Minimal equation of the trigeometric generator}

Let

\[
\mathbb J=p_+J_++p_0J_0+p_-J_-,
\qquad p_++p_0+p_-=I.
\tag{119.1}
\]

The three curvature relations are equivalent to

\[
\boxed{\mathbb J^6=\mathbb J^2.}
\tag{119.2}
\]

Indeed, \(K=\mathbb J^2\) satisfies \(K^3=K\), and the three idempotents are recovered without external data:

\[
p_+=\frac{K^2-K}{2},\qquad
p_0=I-K^2,\qquad
p_-=\frac{K^2+K}{2}.
\tag{119.3}
\]

The intrinsic exponential identity is

\[
\operatorname{Exp}_\star(t\mathbb J)=
p_+(\cos t+\mathbb J\sin t)+
p_0(I+t\mathbb J)+
p_-(\cosh t+\mathbb J\sinh t).
\tag{119.4}
\]

Equations (118.4) and (119.4) are combined in the holonomic propagator

\[
\boxed{
\mathcal E_{\rm HMT}(t)
=\Gamma_9\operatorname{Exp}_\star(t\mathbb J),
\qquad
[D_M,\Gamma_9]=\Gamma_9.}
\tag{119.5}
\]

This is the sought operational compression. The first coordinate of (119.5) governs local response in the three curvatures; the second governs memory raising. APP enters through the sheet idempotents, their projectors and the cocycle defining \(\Gamma_9\); TRIT enters through \(\mathbb J\); TPK enters through composition, the lift and \(D_M\). The correlated system, electron, action, gravity, Barbero--Immirzi, scale constants, symmetries and sequence readers are obtained as typed representations of the universal algebra containing (119.5).

The function of \(\alpha\) becomes transparent in this presentation. \(\operatorname{Exp}_\star\) contains local dynamics and \(\Gamma_9\) the turn with memory; dodecaphasic closure evaluates their compatibility and determines \(\alpha\) as seam. Thus \(\alpha\) links synchronization of \((\pi,e,\varphi)\) to dimensional readers because it measures a closure class of the holonomic propagator, not because it is added to three pre-existing constants.

\textbf{Refutation criterion.} The minimal equation fails if \(\mathbb J^6\ne\mathbb J^2\), if the polynomials (119.3) cease to be orthogonal idempotents, if (119.4) does not restrict to the three certified exponentials, or if \(\alpha\) does not coincide along its two typed routes within the existing compatibility square.
 
% Realizaciones dimensionales completas tomadas de la autoridad material viva.
% El orden respeta la precedencia electrón estructural -> acción -> constitución
% del vacío -> gravitación -> área e información.
% \newcommand{\HMTMonographChapterInput}[2]{%
%   \chapter{#1}%
%   \begingroup
%   \renewcommand{\chapter}[2][]{}%
%   \input{#2}%
%   \endgroup}

\chapter{The electron as a persistent section of the oriented vacancy}
\begingroup
\renewcommand{\chapter}[2][]{}%
\chapter{The electron as the 1st material realization: basal stage, beta closure and the atomistic structure of electricity}
\label{chap:parte-iv-48}

The genealogical state here acquires its first complete material realization. The central localization, the electron--positron band, and the bidirectional reading are constructed before any metrological comparison; the basal stage remains distinct from the governing beta closure.

\section{Closure \(\beta\) of the electronic section and dimensional publication of its rest energy}
\label{chap:biografia-electron-rev6}
\index{electron!holographic construction}
\index{positron!conjugate band}
\index{mass!electronic law}

The electron appears in HMT as the point at which six
preceding and mutually typed constructions are assembled: the incompatibility between
the two arithmetic sheets of APP, the orientation induced by the involution
of the areal--volumetric hinge, two integer quotients from the APP--TPK census,
the dodecaphasic
closure of \(\alpha\), the global defect
\(\Delta_4\), and the unique center of the nonadic atlas. The route band then adds
electron--positron conjugation. The MeV chart and the comparison
with CODATA are subsequent to all of them.

This precedence is part of the mathematical content: the factors, their
types and their causal order are determined before opening the tabulated chart, which
intervenes only in the subsequent metrological comparison.

\subsection{Electronic scalar \(\mathfrak e\) and causal graph of its factors APP–TPK, \(\alpha\), \(\Delta_4\) and nonadic atlas}

The internal output is the dimensionless scalar.
\begin{equation}
 \label{eq:electron-holografico-rev6}
 \boxed{
 \mathfrak e
 =
 \frac{\sqrt3}{4}
 \left[
 \exp\!\left(\frac{\varphi}{\pi^2}\right)
 -\frac{396}{18}\,\alpha^3
 \right]
 \left[
 1+\frac{270}{18}\,\Delta_4
 \right].}
\end{equation}
The identities
\[
 \frac{396}{18}=22=27-5,
 \qquad
 \frac{270}{18}=15=\binom62=3\cdot5
 \]
allow the abbreviated notation to be recovered
\[
 \mathfrak e
 =
 \frac{\sqrt3}{4}
 \left(e^{\varphi/\pi^2}-22\alpha^3\right)
 (1+15\Delta_4).
\]
Here the letter \(e\) inside the exponential is Euler's number; the
subscript of \(\mathfrak e\) denotes the electronic section.

The causal graph of the equation is
\[
\begin{array}{ccl}
\text{APP projectors in dimension three}
&\longrightarrow& \sqrt3/4,\\[2mm]
\text{TPK calendar and areal--volumetric sheets}
&\longrightarrow&
F_{\mathrm{av}},\ \pi/\varphi^2,\
\varphi/\pi^2,\\[2mm]
\text{dodecaphasic closure}
&\longrightarrow&\alpha,\\[2mm]
\text{APP--TPK census and electronic family}
&\longrightarrow&396,\ 18,\ 270,\ 22,\ 15,\ 5,\\[2mm]
\text{global readings of }\pi,e\text{ and register}
&\longrightarrow&\Delta_4,\\[2mm]
\text{nonadic atlas and anti-section}
&\longrightarrow&(5,5),\
\text{electron--positron band}.
\end{array}
\]
No arrow in this diagram receives as an argument the reference electronic mass.

\subsection{1st factor: exact incompatibility of the APP leaves}

For \(d\geq2\), let
\[
 \mathscr A_d=\{1,\ldots,9\}^d
\]
with uniform measure, and let
\[
 \Sigma_d(x)=\operatorname{dr}(x_1+\cdots+x_d),
 \qquad
 \Pi_d(x)=\operatorname{dr}(x_1\cdots x_d)
\]
the additive and multiplicative observables. Their orthogonal conditional expectations
on \(\ell^2(\mathscr A_d)\) are
\(P_{\Sigma,d}\) and \(P_{\Pi,d}\). The
\cref{prop:app-no-conmutatividad} proves by exact enumeration that
\[
 \left\|[P_{\Sigma,3},P_{\Pi,3}]\right\|
 =\frac{\sqrt3}{4}.
\]

\begin{proposition}[Electronic prefactor APP content]
\label{prop:prefactor-app-electron-rev6}
The prefactor \(\sqrt3/4\) of
\cref{eq:electron-holografico-rev6} is the exact norm of the
noncommutativity between the sum and product sheets in dimension three.
It is not a geometric factor introduced from the electronic formula.
\end{proposition}

\begin{proof}
The joint table of \(\Sigma_3\) and \(\Pi_3\) determines the principal
angles between
\(\operatorname{Ran}P_{\Sigma,3}\) and
\(\operatorname{Ran}P_{\Pi,3}\). If \(s\) is one of its singular
values, the corresponding contribution to the commutator norm is
\(s\sqrt{1-s^2}\). Enumeration of the \(9^3\) words gives at most
\[
 s^2(1-s^2)=\frac3{16},
 \qquad s^2=\frac14.
\]
Therefore the maximum norm is \(\sqrt{3/16}=\sqrt3/4\), before
introducing \(\alpha\), \(\Delta_4\), a unit of energy
or a reference mass.
\end{proof}

The physical function of the factor is thereby fixed: it measures the maximal incompatibility of the two arithmetic readings on the same ternary support. The electron does not inherit two independent numbers, one additive and the other multiplicative; it inherits the exact obstruction to rendering the two sheets simultaneously compatible. The prefactor preserves, within the mass law, the local memory of the APP disagreement from which the route arises.

\subsection{2nd factor: the exchange involution between \(\pi\) and \(\varphi\)}

The primitive calendar of modes
\((\Sigma,\Pi,\Pi)\) induces, on the areal and volumetric sheets, the matrix
\[
 F_{\mathrm{av}}
 =
 \begin{pmatrix}
 0&1\\
 1&1
 \end{pmatrix}.
\]
\cref{thm:descenso-necesario-fav} obtains it by counting the transitions
\(\Sigma\to\Pi\), \(\Pi\to\Pi\) and \(\Pi\to\Sigma\), each once.
The matrix is not postulated from the golden ratio: its eigenvalues
\(\varphi\) and \(-\varphi^{-1}\) appear after the incidence has been
constructed.

In the second-order memory,
\(\operatorname{Sym}^2(F_{\mathrm{av}})\) has spectrum
\[
 \varphi^2,\qquad -1,\qquad\varphi^{-2}.
\]
Thus, \(\varphi^{-2}\) is the positive stable mode of the hinge. The closure
mark \(\pi\), applied just once to the areal return, produces
\[
 \frac{\pi}{\varphi^2}.
\]
This reason preserves the closure--self-scaling order. The electronic law uses its opposite orientation.

\begin{theorem}[Exact Inversion of the Electronic Hinge]
\label{thm:espejo-electron-rev6}
Let
\[
 \mathscr R(x,y)=\frac{x}{y^2},
 \qquad
 \mathsf J(x,y)=(y,x).
\]
Then \(\mathsf J^2=I\) and
\[
 \mathscr R(\pi,\varphi)=\frac{\pi}{\varphi^2},
 \qquad
 (\mathscr R\circ\mathsf J)(\pi,\varphi)
 =\frac{\varphi}{\pi^2}.
\]
Consequently,
\[
 \frac{\varphi}{\pi^2}
 =
 \frac{1}{
 \varphi^3\bigl(\pi/\varphi^2\bigr)^2},
 \qquad
 e^{\varphi/\pi^2}
 =
 \exp\!\left[
 \frac{1}{
 \varphi^3\bigl(\pi/\varphi^2\bigr)^2}
 \right].
\]
\end{theorem}

\begin{proof}
The first pair of identities is obtained by applying \(\mathscr R\) before
and after the involution \(\mathsf J\). For the reparametrization,
\[
 \varphi^3
 \left(\frac{\pi}{\varphi^2}\right)^2
 =\frac{\pi^2}{\varphi},
\]
and its inverse is \(\varphi/\pi^2\). The last equality follows by applying
the exponential.
\end{proof}

The simultaneous presence of \(\pi\) and \(\varphi\) therefore has a precise structural meaning. \(\pi/\varphi^2\) reads areal closure over quadratic volumetric memory; \(\varphi/\pi^2\) reads interior self-growth over quadratic closure. They are not two similar quotients chosen for convenience. They are the two orientations of the same reader, exchanged by an exact involution. The term \(e^{\varphi/\pi^2}\) publishes the interior orientation of that hinge in the electron section.

\subsection{3rd and 4th factors: the integers 22, 15 and the center 5}

The census of the APP chart employed by the electronic family provides
three integer sums:
\[
 \Sigma(\mathrm{APP})=396,
 \qquad
 \Sigma(\mathrm{centro}\ 2\times2)=18,
 \qquad
 \Sigma(\mathrm{canal})=270.
\]
The first admits the ring decomposition.
\[
 396=18+72+108+198.
\]
The kernel \(2\times2\) is therefore an explicit substructure of the
complete chart, and the integer 270 belongs to the global channel of the same
APP--TPK census.
The sum of the even rings is
\[
 18+108=126,
 \qquad
 \frac{396}{126}=\frac{22}{7}.
\]
This last quotient is the rational closure seal of the chart, not an
identification of \(22/7\) with the number \(\pi\) generated by the coinductive
reader. It does show that the numerator \(22\) already serves an
internal function before appearing in the electronic law.

\begin{theorem}[Internal Origin of the Electronic Coefficients]
\label{thm:coeficientes-electron-rev6}
In the electronic family fixed by the APP chart and its global channel,
\[
 \boxed{
 22=\frac{\Sigma(\mathrm{APP})}
 {\Sigma(\mathrm{centro})}
 =\frac{396}{18}=27-5,}
\]
\[
 \boxed{
 15=\frac{\Sigma(\mathrm{canal})}
 {\Sigma(\mathrm{centro})}
 =\frac{270}{18}=3\cdot5=\binom62.}
\]
In particular, the coefficients are determined without consulting
\(Z_0\), the electron mass or a chart of units.
\end{theorem}

\begin{proof}
The first two equalities are exact divisions of the finite census:
\[
 396=22\cdot18,
 \qquad
 270=15\cdot18.
\]
By defining
\[
 p:=27-22,
\]
one obtains \(p=5\), and then
\[
 15=3p=3\cdot5=\binom62.
\]
All quantities belong to the integer domain of the chart and the
family before metrological transduction.
\end{proof}

\subsubsection{Persistence selector in the central return}

The same pair has a dynamic construction that also fixes its sign without
consulting a mass. The four orientations of the central seed
\((5,5)\) return to that cell for the first time after twenty-seven steps. The
typed return space is
\[
 \mathcal K_e=\mathbb Z/9\mathbb Z\times\mathbb F_3,
 \qquad |\mathcal K_e|=27.
\]
Moreover, the microfiber of \(\pi\) contains the five genealogically
distinct regions
\[
 \mathscr R_\pi=3_{++}+1_{--}+1_{\perp}.
\]
The regional gauge \(\operatorname{diag}_c\) and, within a repeated phase,
the internal order of \(U_6\) define the injection
\[
 \iota:\mathscr R_\pi\hookrightarrow\mathcal K_e
\]
through the five distinct places
\[
 (4,0),\quad(5,0),\quad(6,0),\quad(6,1),\quad(6,2).
\]
The first two correspond, respectively, to the transverse region and
the mutated return; the last three separate the positive regions by their
tritic sheet. The common visible word does not intervene as the inverse of the map.

\begin{theorem}[Internal selector of the electronic coefficients]
\label{thm:selector-persistencia-electron-rev7}
With the central return and the five typed regions fixed, the coefficients of
the electronic equation are
\[
 \boxed{
 n_e=-\bigl|\mathcal K_e\setminus\iota(\mathscr R_\pi)\bigr|=-22,
 \qquad
 m_e^{\rm tor}=\bigl|\mathscr R_\pi\times\mathbb F_3\bigr|=+15.}
\]
The first symbol records the boundary deficit; the second, the retained torsional memory.
\end{theorem}

\begin{proof}
The five places above are distinct, so \(\iota\) is
injective and
\[
 |\mathcal K_e\setminus\iota(\mathscr R_\pi)|=27-5=22.
\]
Each of the five regions preserves three triticale leaves, hence
\[
 |\mathscr R_\pi\times\mathbb F_3|=5\cdot3=15.
\]
The orientation rule assigns a negative sign to states removed from the
shell and a positive sign to retained memory. The construction opens only the
center, the return of twenty-seven states, and the regional identities; it does not
open an observed mass, an impedance, or a calibrated signature.
\end{proof}

The number \(22\) measures how many central kernels of weight \(18\) are recomposed by the
APP sum \(396\). The number \(15\) measures how many of those same kernels
are recomposed by the channel \(270\), and simultaneously counts the pairs of a
hexad. The number \(5\) is not an additional parameter:
\[
 5=27-22,\qquad 15=3\cdot5.
\]
The three expressions display the same closure from depth 27,
from the central normalization, and from hexadic combinatorics.

\subsubsection{Conservation of coefficients 22 and 15 in the material realization of the electron}

The historical exploration first located \(p=5\) through a finite
comparison with \(Z_0\), and from there wrote
\((27-p,3p)=(22,15)\). That procedure retains value as a history of
discovery and as an external check. It is not, however, the proof
that governs the law: the subsequent count
\[
 \frac{396}{18}=22,
 \qquad
 \frac{270}{18}=15
\]
constructs both integers within APP--TPK and recovers \(p=5\) as a
consequence. The history explains how the pair was found; the census
explains why it belongs to the equation.

\subsection{5.º factor: the cubo of the closure \(\alpha\)}

The closure of \(\alpha\) precedes Dimensional Mechanics.
As established in \cref{chap:alpha}, the twelve base-thousand blocks
of \(\pi\), \(e\), \(\varphi\) and the seal \(K\) combine through the
reversible carry
\[
 P_m+E_m-\Phi_m-K_m+c_{m+1}
 =A_m+1000c_m,
 \qquad c_{13}=c_1=0.
\]
The output is
\[
 A_\alpha=
 (007,297,352,569,283,800,997,285,105,472,380,663)
\]
and, by concatenation, the exact dodecaphasic window
\[
 \alpha_{12}=\pi_{12}(\alpha)
 =
 \frac{
 7297352569283800997285105472380663
 }{10^{36}}.
\]
Telescoping the carries proves the complete integer identity; the
compatible prolongation constructs
\(\alpha=\varprojlim_N A^{(N)}\), and \(\alpha_{12}\) is its
thirty-six-digit rational projection. The
comparison with the fine-structure constant is performed subsequently and does not
intervene in its construction.

The electronic law takes the third power
\[
 \alpha^3.
\]
Mathematically, this is the multiplicative evaluation of the third
symmetric power of the same scalar already closed; neither three
metrological copies nor three different parameters are introduced. Its contribution is
\[
 22\alpha^3
 =
 0.00000854906599200551727014979807598228\ldots.
\]

\begin{proposition}[Statute of the Electronic Cubic Channel]
\label{prop:alpha-cubica-electron-rev6}
The term \(22\alpha^3\) of the electronic law depends
only on the prior closure of \(\alpha\), on internal
multiplication, and on the integer \(22\) of
\cref{thm:coeficientes-electron-rev6}. It does not depend on a reading of the
electron mass.
\end{proposition}

\begin{proof}
\(\alpha_{12}\) is the rational number fixed by the concatenation and carry
identity, and the compatible family of windows determines the section
\(\alpha\). The projection \(\alpha_{12}^3\) certifies the
printed digits of \(\alpha^3\).
The coefficient \(22\) is determined by \(396/18\). The composition of
both operations contains only data constructed before the
energy chart.
\end{proof}

This proposition fixes exactly what is used for the electron. The third power is the cubic degree of this channel in the published law. There is no need here to postulate a universal rule according to which every occurrence of \(\alpha^d\) automatically represents a dimension \(d\); a universal statement of that kind would require its own theorem of multiplicativity among degrees. The electron genealogy is complete without turning that generalization into a premise.

\subsection{6th factor: minimal torsion and global defect \(\Delta_4\)}
\index{minimum torsion!\(\Delta_4\)}

The terms \emph{minimal torsion} and \emph{global defect}
denote a single invariant here, not two distinct corrections. The object
used by the mass law is
\begin{equation}
 \label{eq:delta4-electron-rev6}
 \boxed{
 \Delta_4
 =
 \frac{\pi-e\log\pi}{270}
 \left(1+\frac{\pi}{729}\right).}
\end{equation}
Its inputs are the already fixed HMT readings of \(\pi\) and \(e\), the global
channel \(270\), and the interior cell \(729\). Consequently,
\[
 \Delta_4
 =
 0.00011119642269214005405113478612067909\ldots.
\]
The first ratio measures the signed defect \(\pi-e\log\pi\) in units of the 270 channel; the second factor couples that ratio to the interior cell through \(\pi/729\). This reading adds no parameters to the formula: it separates its two constituent operations. The \cref{p3-20260826:chap:medida-accion-positiva} shows that it is a pre-dimensional mathematical input and distinguishes it from the operator that eventually realizes it as a change of sheet.

\begin{proposition}[Globality of the Electronic Defect]
\label{prop:delta4-global-electron-rev6}
\(\Delta_4\) is a global invariant preceding the species. In the electronic
law it enters through
\[
 1+15\Delta_4
 =
 1.00166794634038210081076702179181019\ldots.
\]
It is not a corrector chosen after knowing the residue of the electron mass.
\end{proposition}

\begin{proof}
The \cref{eq:delta4-electron-rev6} contains neither a species label, a
mass, nor a unit. The integer \(15\) has already been constructed by \(270/18\).
Therefore, the composition \(1+15\Delta_4\) is determined by the
global data and the internal census.
\end{proof}

\(\Delta_4\) must be distinguished from the other blocks of the mass character and from the other notions of torsion. In particular, it is neither \(s_{270}\), nor the quadratic invariant \(135\Delta_4^2\), nor the angular residue \(\varepsilon_{\rm res}\), nor the geometric Cartan--Holst torsion. Substituting one for another would repeat or alter the level of the correction. For the electron, \(\Delta_4\) preserves the global memory of refinement; the factor \(15\) states how many times the electron family reads it. The qualifier \emph{minimal} preserves the historical name of this invariant and is not used as a substitute for a variational theorem of minimality.

\subsection{Material realization of the fixed point (5,5) in the electronic closure}

Let
\[
 \mathcal C_{81}=\{1,\ldots,9\}^2
\]
the nonadic atlas and
\[
 \rho(r,c)=(10-r,10-c)
\]
its central inversion. \cref{thm:censo-atlas-81} proves that the atlas
contains a unique generation-zero cell and that this cell bears the positional
label \(e_{\mathrm{core}}\).

\begin{proposition}[Uniqueness of the Electronic Center]
\label{prop:centro-electron-rev6}
The involution \(\rho\) has a unique fixed point:
\[
 \boxed{(r,c)=(5,5).}
\]
Its uniqueness is combinatorial and does not depend on the value of
\(\mathfrak e\).
\end{proposition}

\begin{proof}
The equation \((r,c)=\rho(r,c)\) is equivalent to
\[
 2r=10,\qquad2c=10.
\]
In \(\{1,\ldots,9\}^2\) it has the unique solution \(r=c=5\). The evaluation
of the mass does not appear in the argument.
\end{proof}

\begin{corollary}[Universality of the Central Section]
\label{cor:universalidad-electron-central-rev6}
Every admissible automorphism of the atlas that intertwines the central inversion
preserves the cell \((5,5)\).
\end{corollary}

\begin{proof}
If \(f\rho=\rho f\) and \(\rho(5,5)=(5,5)\), then
\(\rho f(5,5)=f\rho(5,5)=f(5,5)\). Therefore \(f(5,5)\) is another fixed point
of \(\rho\); uniqueness forces \(f(5,5)=(5,5)\).
\end{proof}

This is the structural reason why the electronic class is the same in
every realization preserving the atlas and its inversion: a local decimal does not
recur by chance; rather, every admissible chart must
transport the same fixed point and the same chain of invariants. The
physical reading of this uniqueness identifies each electron with a realization of that
central class; the proposition proves the internal universality of the atlas.

In the electronic family fixed by the construction, this point is the geometric
anchor of the section. Four objects should not be confused:
\[
 (5,5),\qquad
 e_{\mathrm{core}},\qquad
 v_e=(0,0,0,0,0),\qquad
 \text{the raw signature of the central cell}.
\]
The first is a position; the second, its internal label; the third,
the chosen origin in the fine action lattice; the fourth preserves the
raw atlas data and need not be the zero vector. Their
coordination explains why the electron occupies the center without erasing the
genealogy of the representations describing it.

\paragraph{Oriented electronic routes \(e_{++},e_{--}\), envelope \(e_{\mathrm{env}}\), action origin and signatures of the central cell}
The central cell carries two oriented histories of twelve events. Writing
\(v^2\) for the concatenation of two copies of a sextuple, the active
route of orientation \({++}\), its opposite history, and the historical
envelope are, respectively,
\[
\begin{aligned}
 e_{\rm act}=e_{++}
   &=(2,2,5,-4,-3,-2)^2,\\
 e_{--}
   &=(4,3,2,-2,-2,-5)^2,\\
 e_{\rm env}
   &=(4,3,5,-4,-3,-5)^2.
\end{aligned}
\tag{E-centro-1}\label{eq:electron-rutas-finas-rev8}
\]
The two stories publish the same word
\[
 \tau_{12}=\texttt{+++---+++---},
\]
but they are not the same route. The envelope takes the absolute maximum
coordinate by coordinate while preserving the sign; it records amplitudes, but
flattens the genealogical order and does not replace \(e_{\rm act}\).

The other readers of the center live in distinct codomains:
\[
\begin{array}{rcll}
 v_e&=&(0,0,0,0,0),
   &\text{affine origin of the action lattice},\\
 R(5,5)&=&(24,0,12,0,-12),
   &\text{raw APP signature of the cell},\\
 \operatorname{Sig}^{\Gamma_{108}}_{\rm masa}
   &=&(-12,-4,12,-6,12),
   &\text{even mass projection},\\
 \operatorname{Sig}^{\Gamma_{108}}_{\rm carga}
   &=&(0,0,0,0),
   &\text{odd charge projection}.
\end{array}
\tag{E-centro-2}\label{eq:electron-firmas-tipadas-rev8}
\]
The fact that the charge word has zero odd signature does not turn the route into
zero, nor does it identify the raw signature with the affine origin. Position, active
route, opposite route, envelope, origin and both signatures are coordinated
objects, not six notations for the same vector.

\subsection{Particle–antiparticle conjugation and persistence of the electronic band}

Electronic identity does not end at the central point or at the mass
value. Material conjugation is realized on dodecaphasic words. Let
\[
 C_M(\mathbf t)=-\operatorname{rev}(\mathbf t)
\]
be the oriented anti-section and let
\(\chi\in\{-1,+1\}\) be the charge orientation. The band is the class
\[
 [(\mathbf t,\chi)]_M
 =
 \{(\mathbf t,\chi),(C_M\mathbf t,-\chi)\}
\]
of \cref{def:banda-dodecafasica-ruta}. The band preserves its even
components and exhibits two sections with opposite odd components.

\begin{theorem}[Electron and positron as sections of a band]
\label{thm:electron-positron-biografia-rev6}
For the band character of
\cref{thm:principio-conjugacion-masa-banda},
\[
 m(C_M\mathbf t,-\chi)=m(\mathbf t,\chi).
\]
In the electronic channel,
\[
 (n,\chi)=(-22,+1)
 \quad\longleftrightarrow\quad
 (+22,-1)
\]
and both pairs produce
\[
 \chi n=-22.
\]
Therefore, electron and positron have the same mass and opposite charge orientations.
\end{theorem}

\begin{proof}
If \(E_+\) and \(E_-\) are, respectively, the even and odd
components of the character with respect to \(C_M\), then
\[
 E_+(C_M\mathbf t)=E_+(\mathbf t),
 \qquad
 E_-(C_M\mathbf t)=-E_-(\mathbf t).
\]
From here,
\[
 E_{\mathrm{band}}(C_M\mathbf t,-\chi)
 =
 E_{\mathrm{band}}(\mathbf t,\chi),
\]
and the equality of mass is obtained by applying the exponential. In the electronic term,
\[
 (+1)(-22)=(-1)(+22)=-22.
\]
\end{proof}

The finite closure of \cref{prop:cierre-antiseccion-ct108} reinforces this
interpretation: over the eighty-one CT--108 routes, neither pure sign nor
pure reversal preserves the taxonomy, whereas
\(-\operatorname{rev}\) does preserve it, with
\[
 81=54+9+9+9.
\]
The antiparticle is therefore neither a negative mass nor an independent
copy. It is the other oriented section of the same band: it preserves the
even invariants that determine mass and reverses the odd invariants that
determine charge and orientation.

\subsection{Electronic Closure Theorem}

\begin{theorem}[Holographic composition of the electronic section]
\label{thm:cierre-electron-holografico-rev6}
With the objects constructed in the preceding sections fixed, the
dimensionless electronic section is
\[
 \mathfrak e
 =
 \left\|[P_{\Sigma,3},P_{\Pi,3}]\right\|
 \left[
 \exp\bigl((\mathscr R\circ\mathsf J)(\pi,\varphi)\bigr)
 -\frac{\Sigma(\mathrm{APP})}
 {\Sigma(\mathrm{centro})}\,
 \alpha^3
 \right]
 \left[
 1+
 \frac{\Sigma(\mathrm{canal})}
 {\Sigma(\mathrm{centro})}\,
 \Delta_4
 \right].
\]
Equivalently,
\[
 \boxed{
 \mathfrak e
 =
 \frac{\sqrt3}{4}
 \left(e^{\varphi/\pi^2}
 -22\alpha^3\right)
 (1+15\Delta_4).}
\]
Its evaluation is
\[
 \boxed{
 \mathfrak e
 =
 0.51099892248806607250987087719134943763\ldots.}
\]
\end{theorem}

\begin{proof}
The \cref{prop:prefactor-app-electron-rev6} contributes
\[
 \left\|[P_{\Sigma,3},P_{\Pi,3}]\right\|=\sqrt3/4.
\]
The \cref{thm:espejo-electron-rev6} contributes
\[
 (\mathscr R\circ\mathsf J)(\pi,\varphi)=\varphi/\pi^2.
\]
The \cref{thm:coeficientes-electron-rev6} contributes
\[
 \frac{\Sigma(\mathrm{APP})}{\Sigma(\mathrm{centro})}=22,
 \qquad
 \frac{\Sigma(\mathrm{canal})}{\Sigma(\mathrm{centro})}=15.
\]
The dodecaphasic closure contributes the rational window \(\alpha_{12}\) of the
coinductive section \(\alpha\), and
\cref{eq:delta4-electron-rev6} contributes the global defect. Substitution
yields the abbreviated formula. Successive evaluation of
\[
\begin{aligned}
 e^{\varphi/\pi^2}
 &=1.17814494259630678081160095680888910\ldots,\\
 22\alpha^3
 &=0.00000854906599200551727014979807598228\ldots,\\
 1+15\Delta_4
 &=1.00166794634038210081076702179181019\ldots
\end{aligned}
\]
gives the declared value.
\end{proof}

The theorem is holographic in a verifiable sense: each local factor
retains a distinct layer of the complete system. The commutator retains sheet
duality; the exponential retains the interior orientation of the
hinge \(\pi\)--\(\varphi\); \(22\) and \(15\) retain the global census and
its central kernel; \(\alpha^3\) retains fine closure; and
\(\Delta_4\) retains global refinement. The center determines where the
section is anchored, and the band determines how it survives under conjugation.
The value does not replace this structure: it is its common evaluation.

\subsubsection{Anterior and posterior action splitting upon return}
\label{sec:electron-desdoblamiento-two-way-rev9}

The regional realization of the conjugate angular coordinate \(C^*\) distinguishes two coordinates on the
same internal action line. The coordinate before the return is \(H_5\) and
the subsequent coordinate is
\[
\etaRet
 =H_5-\frac{90}{\pi}\mathscr C_\pi(\alpha).
\]
The notation \(\etaRet\) is the canonical sexagesimal coordinate defined on the
action branch by \(\etaRet=\Cstar/2-\alphaAngle\). The preceding equality is
evaluated in that chart; it does not introduce a second return coordinate.
Therefore, the additive and multiplicative charts of the oriented splitting
are
\begin{equation}
 \boxed{
 \delta_{2w}:=H_5-\etaRet
 =\frac{90}{\pi}\mathscr C_\pi(\alpha),
 \qquad
 R_{\rm act}:=\frac{\hbar^{\rm pre}_{\rm HMT}}
 {\hbar^{\rm ret}_{\rm HMT}}
 =\frac{H_5}{\etaRet}.}
 \label{eq:electron-desdoblamiento-two-way-rev9}
\end{equation}
Both coordinates follow from a single pre-return/post-return transport.
They are not two Planck constants: the common factor
\(10^{-34}\mathcal U_S\) fixes the dimensional line and cancels exactly in
\(R_{\rm act}\). The action ellipse subsequently realizes
\(\operatorname{Area}(\theta)=\hbar_{\rm HMT}^{\rm ret}\theta\), so that
the splitting precedes the holonomy reading.

The coupling
\begin{equation}
 \kappa_e^{\beta}
 =80-54\alpha+6\Delta_4,
 \qquad
 \mathfrak e_e^{\beta}
 =\mathfrak eR_{\rm act}^{\kappa_e^{\beta}}
 =0.5109989506900008086\ldots
 \label{eq:electron-beta-two-way-rev9}
\end{equation}
is the bidirectional quantitative closure of the electron. The two independent
readers of the electronic register constructed in the following chapter
select, without consulting the comparison mass,
\[
 K_D(\Gamma_e)=K_\Omega(\Gamma_e)=(80,54,6),
\]
and derive exactly \(\kappa_e^\beta\). The integer \(80\) comes from the
first coordinate of those readers; the word \(w_e=201101\) retains its
own residence in the channel of Euler's constant.

\begin{statusbox}[title={Internal Closure of Bidirectional Refinement}]
The identity \eqref{eq:electron-desdoblamiento-two-way-rev9}, the reduction
\((80,54,6)\), and the equation
\eqref{eq:electron-beta-two-way-rev9} constitute internally exact,
target--free results on the electronic record. Propagation to fibers of
higher rank is formulated by means of the multisection operators.
\end{statusbox}

\subsection{Energy chart and contrast of the bidirectional closure}

Let \(\mathcal U_E\) a real line of energy dimensions and let
\(u_E\) a chosen basis. The metrological transduction is
\[
 \varsigma_{u_E}:\mathbb R\longrightarrow\mathcal U_E,
 \qquad
 \varsigma_{u_E}(x)=x\,u_E.
\]
The chart preserves two stages of the same genealogy. The basal stage is
\[
 E_e^{(0)}
 =
 \varsigma_{u_E}(\mathfrak e).
\]
With \(u_E=1\,\mathrm{MeV}\),
\[
 E_e^{(0)}
 =
 0.5109989224880660725\ldots\ \mathrm{MeV}.
\]
The action ratio and the derived exponent complete the leading stage,
\[
 E_e^\beta=E_e^{(0)}R_{\rm act}^{\kappa_e^\beta}
 =0.5109989506900008086082571648\ldots\ \mathrm{MeV},
 \qquad
 m_e^\beta=E_e^\beta/c^2.
\]
The unit does not appear in \cref{eq:electron-holografico-rev6}; it is the basis of a dimensional
line applied to an already constructed output.

The CODATA 2022 reference for the electron energy equivalent is
\cite{CODATA2022}
\[
 E_e^{\mathrm{CODATA}}
 =
 0.51099895069(16)\ \mathrm{MeV}.
\]
In that same chart, the bidirectional output satisfies
\[
 E_e^\beta-E_e^{\mathrm{CODATA}}
 =
 8.086082571648\times10^{-16}\ \mathrm{MeV},
\]
with relative difference
\[
 1.582406883758\times10^{-15}.
\]
The comparison is carried out after freezing the internal evaluation. The identities of the basal selector remain
\[
 22=\frac{396}{18},
 \qquad
 15=\frac{270}{18},
 \qquad
 p=27-22=5.
\]
and additionally receive the combinatorial proof
\(-22=-|K_e\setminus\iota(R_\pi)|\),
\(+15=|R_\pi\times\mathbb F_3|\). The metrological comparison verifies the
output; the route, the readers, and the action ratio generate it.

\subsection{Structural Consequences of Electronic Closure}

The preceding structural reading imposes several distinctions:

\begin{enumerate}
\item \(\sqrt3/4\) is the norm of the APP commutator in dimension three. Its
agreement with other geometric formulas does not supersede this provenance.

\item \(\pi/\varphi^2\) and \(\varphi/\pi^2\) are not interchangeable. The
first ratio is the closure orientation; the second, obtained through
\(\mathsf J\), is the electronic orientation of self-growth.

\item \(22\) and \(15\) follow from \(396/18\) and \(270/18\).
Presenting them as values logically selected by \(Z_0\) would reverse
the demonstrated dependence.

\item \(\alpha^3\) uses the prior closure of
\(\alpha\). The cubic power belongs to the
electron law; by itself, it does not authorize a universal rule
\(\alpha^d\leftrightarrow d\) without an additional theorem.

\item \(\Delta_4\) is global. No must reemplazarse by
\(s_{270}\), by \(135\Delta_4^2\) neither by a residue particular of the
especie.

\item The position \((5,5)\), the label \(e_{\mathrm{core}}\), the fine
origin \(v_e=0\), and the central raw signature are coordinated objects, not
synonyms.

\item The electron and positron are not distinguished by two mass laws: they are the
two oriented sections of a single band, with even mass and odd charge.

\item \(0.5109989224\ldots\) is first a dimensionless scalar. MeV,
\(\mathrm{MeV}/c^2\), CODATA, and \(Z_0\) belong to the subsequent
comparison.
\end{enumerate}

These distinctions explain why the electronic equation is a
miniature of the HMT--MD system. In a single section, support,
sheets, orientation, census, fine closure, refinement, center, persistence,
and conjugation reappear. This is the content of the word \emph{holographic}: not
that a digit poetically represents the whole, but that each factor
preserves a verifiable dependence on the whole that produces it.
 
\section{Bidirectional derivation of the electronic exponent from the enriched route}
\label{sec:cierre-electron-two-way-rev11}
\index{electron!bidirectional reader}
\index{calendar of order 108! electronic exponent}

The enriched route of the central point publishes two complementary representations
of its memory: the complete tritic word
\(\gamma_{108}\in\mathbb F_3^{108}\) and the signed dodecaphase record
\(\tau_{12}\in\{-1,0,+1\}^{12}\). The two readers defined below
are computed exclusively from those representations, the calendar, and the
HMT constants already constructed.

\subsection{Reader of displacements on the cyclic calendar of order 108}

Let \(S\) be the cyclic shift on \(\mathbb Z/108\mathbb Z\). For
\(k\in\mathbb Z/108\mathbb Z\), define
\begin{equation}
 D_k(\Gamma):=
 \#\{t\in\mathbb Z/108\mathbb Z:\gamma_t\ne\gamma_{t+k}\}
 =\lVert(I-S^k)\gamma\rVert_0.
 \label{eq:lector-dk-ct108-rev11}
\end{equation}
The integer triple of the reader is
\begin{equation}
 K_D(\Gamma)=
 \left(D_{27}+D_{36},\ D_1,\ \frac{D_4-D_3}{2}\right)
 =:(A_D,B_D,C_D),
 \label{eq:KD-rev11}
\end{equation}
and its scalar evaluation is
\begin{equation}
 \kappa_D(\Gamma)=A_D-{\alpha}B_D+\Delta_4 C_D.
 \label{eq:kappaD-rev11}
\end{equation}
The shifts (27) and (36) are the chronological lifts of
\(90^\circ\) and \(120^\circ\) because nine discrete steps represent
\(30^\circ\). The term \(D_1\) measures the immediate local boundary, and the
difference \(D_4-D_3\) compares the two genealogies before their height
quotient. In the domain of the 324 routes,
\(\gamma_{t+54}=\gamma_t\) holds; all the \(D_k\) are even and \(C_D\) is an integer.

\subsection{Dodecaphasic correlation difference reader}

For the word \(\tau=(\tau_r)_{r\in\mathbb Z/12\mathbb Z}\), let
\[
 C_k^{\rm corr}(\tau)=\sum_{r=0}^{11}\tau_r\tau_{r+k},
 \qquad
 A_\ell(\tau)=\sum_{r=0}^{11}
 \left(\sum_{j=0}^{\ell-1}\tau_{r+j}\right)^2.
\]
The primitive genealogical obstruction and the antipodal memory are
\begin{align}
 \Omega_{90\mid120}(\tau)
 &=4A_3-3A_4
 =-2C_1^{\rm corr}-4C_2^{\rm corr}-6C_3^{\rm corr},
 \label{eq:omega-two-way-rev11}\\
 P_6(\tau)&=\sum_{r=0}^{5}\tau_r\tau_{r+6}.
 \label{eq:p6-two-way-rev11}
\end{align}
The symbol \(P_6\) prevents the colision with the short cycle of the corredor of
constantes. The second triple and its evaluation are
\begin{equation}
 K_\Omega(\Gamma)=\bigl(\Omega_{90\mid120}(\tau_\Gamma),54,
 P_6(\tau_\Gamma)\bigr),
 \qquad
 \kappa_\Omega=\Omega_{90\mid120}-54\alpha+P_6\Delta_4.
 \label{eq:kappa-omega-rev11}
\end{equation}

\HMTClaim{MD-R-ELECTRON-OMEGA-PRIMITIVE}
\begin{lemma}[Linear integral primitive uniqueness]
\label{lem:unicidad-omega-rev11}
Let \(L=uA_3+vA_4\), with \(u,v\in\mathbb Z\), be a contrast that vanishes
when the genealogies descend to the same height,
\(A_3=3a\), \(A_4=4a\). Then
\((u,v)=n(4,-3)\). Once the sign \(90-120\) has been fixed,
\(\Omega_{90\mid120}=4A_3-3A_4\) is the primitive generator.
\end{lemma}

\begin{proof}
Vanishing for every \(a\) requires \(3u+4v=0\). Since
\(\gcd(3,4)=1\), the integer solutions are
\((u,v)=n(4,-3)\). Primitivity corresponds to \(|n|=1\).
\end{proof}

The lemma selects \(\Omega\) within the class of integral linear contrasts
of the energies \(A_3,A_4\). Its uniqueness domain is that functional
class; nonlinear invariants belong to a distinct selection problem.

\subsection{Evaluation of the central route}

Before metrological comparison, the electronic route \(\Gamma_e=(5,5,++)\) provides,
\[
 (D_1,D_3,D_4,D_{27},D_{36})=(54,56,68,48,32),
\]
and its signed dodecaphonic register is
\[
 \tau_e=(+,+,+,-,-,-,+,+,+,-,-,-).
\]
Therefore,
\[
 K_D(\Gamma_e)=(48+32,54,(68-56)/2)=(80,54,6),
\]
\[
 K_\Omega(\Gamma_e)=(80,54,6).
\]

\HMTClaim{MD-R-ELECTRON-TWOWAY-REDUCTION}
\begin{theorem}[Common Electronic Reduction of Bidirectional Readers]
\label{thm:reduccion-electron-two-way-rev11}
The readers \(K_D\) and \(K_\Omega\) coincide on the central route and
determine
\begin{equation}
 \kappa_e=80-54\alpha+6\Delta_4.
 \label{eq:kappa-electron-demostrado-rev11}
\end{equation}
\end{theorem}

\begin{proof}
The two preceding evaluations yield the same integer triple. Substitution
of its components into
\(A-{\alpha}B+\Delta_4C\) proves
\eqref{eq:kappa-electron-demostrado-rev11}. Every quantity used in this
reduction belongs to the route, to its two records, or to the previously
constructed HMT constants.
\end{proof}

\paragraph{Subsequent realization on the energy chart.}
With \(R_{\rm act}=H_5/\etaRet\), the positive action on the baseline
\(\mathfrak e_0\) yields
\begin{equation}
 \mathfrak e_e^\beta
 =\mathfrak e_0R_{\rm act}^{\kappa_e}.
 \label{eq:masa-electron-adimensional-two-way-rev11}
\end{equation}
After applying the current chart to MeV is obtained
\begin{equation}
 \varsigma_{u_E}(\mathfrak e_e^\beta)\big|_{u_E=1\,\mathrm{MeV}}
 =0.5109989506900008086082571648\ldots\ \mathrm{MeV}.
 \label{eq:masa-electron-two-way-rev11}
\end{equation}
The first equality belongs to the internal theorem conditioned by the basal and
the action line; the second is its dimensional realization. The external
value is reserved for metrological comparison.

The integer (80) follows from \(D_{27}+D_{36}\) and, in the dodecaphase reader,
from \(\Omega_{90\mid120}\). Its coincidence with other internal cardinals is
a compatibility control and does not replace this operator provenance. The
orientation route \((--)\) is a temporal translation of route \((++)\)
and preserves the same profile, in accordance with mass parity under conjugation.

\subsection{Separation of readers and return scales}

Numerical coincidence of internal components is subordinate to their
type. At the central point, the objects involved include
\[
 R_{\rm atlas}(5,5)=(24,0,12,0,-12),
 \qquad
 L_{\kappa_0}(\Gamma_e^{++})=(24,0,0,0,-12),
\]
\[
 \begin{gathered}
 k_{\Delta_4}^{\rm dod},k_{\Delta_4}^{\rm pro}:
 \mathscr R_{324}\longrightarrow\mathbb Z,\\
 k_{\Delta_4}^{\rm dod}(\Gamma_e^{++})=12,
 \qquad
 k_{\Delta_4}^{\rm pro}(\Gamma_e^{++})=4,\\
 K_D(\Gamma_e^{++})=K_\Omega(\Gamma_e^{++})=(80,54,6).
 \end{gathered}
\]
Here \(\mathscr R_{324}\) denotes the domain of the 324 oriented routes.
The first is a static family signature; the second, a genealogy
extractor; the third and fourth terms are torsional readers defined,
respectively, by dodecaphasic reading and prospective prolongation;
the final object is the pair of bidirectional readers. No partial equality
identifies those morphisms or permits transferring a coefficient between
readers.

Four temporal scales are also distinguished: \(27\) is the first return of class, orbit, and cell; \(54=D_1(\Gamma_e)\) is the observable kinematic return within the cycle; \(108\) closes the calendar; and \(216\) closes the oriented kinematic state in the extension. This typing keeps positional return, remark, and oriented holonomy separate.

\begingroup\fontsize{10.2}{11.7}\selectfont
\setlength{\parskip}{2pt}\setlength{\abovedisplayskip}{5pt}\setlength{\belowdisplayskip}{5pt}
\paragraph{Causal chain of the electronic outcome}
The route and its two registers determine \(K_D\), \(K_\Omega\), their common
reduction, and the evaluation of \(\kappa_e\). The realization in MeV subsequently applies
the energy chart; external mass enters only in the comparison. The pair
\(\mathbf K=(K_D,K_\Omega)\) is the complete operator output of HMT on
each route. In a noncentral fiber, its representation by a single physical
scalar requires an additional morphism between the fiber operator and the
dimensional chart; this requirement neither modifies the internal law nor authorizes
the selection of a cell by metrological proximity.
 
\paragraph{From the Cartan–Holst realization to horizon–area–information confluence through conservation of action} The electronic closure fixes a material realization; the next step identifies the grammar that converts action, clock, and mass into physical magnitudes without using that realization as input data.

\par\endgroup
 \endgroup

\chapter{Determinantal derivation of the absolute action scale and valuation \texorpdfstring{\(10^{-34}\)}{10^-34}}
\begingroup
\renewcommand{\chapter}[2][]{}%
\chapter{Determinantal derivation of the action scale: Smith normal form, cokernel of order \(16\) and valuation \(10^{-34}\)}
\label{chap:p3-final:34:escala_accion}
\label{ch:vacuum}

\section{Derivation of the action scale from the dodecaphasic state}

The action decade is determined intrinsically, without adopting it as an
inherited unit. The local determinantal functional is applied to the previously
constructed HMT output \(\alpha\) and defines
\begin{equation}
\Sigma_H=\varphi\alpha^{16}.
\label{eq:action-determinantal-reader}
\end{equation}
The exponent \(16\) is the order of the local cokernel of the signed closure; its decimal valuation fixes
\begin{equation}
\operatorname{ord}_{10}(\Sigma_H)=-34.
\label{eq:action-decade}
\end{equation}

The origin of the exponent can be verified within this volume. On the
local orbit of four leaves, the Hadamard block acts
\begin{equation}
H_4=
\begin{pmatrix}
1&1&1&1\\
1&1&-1&-1\\
1&-1&1&-1\\
1&-1&-1&1
\end{pmatrix}.
\label{eq:action-hadamard}
\end{equation}

\begin{lemma}[Local index of the action line]
The form normal of Smith of \(H_4\) is
\begin{equation}
\operatorname{SNF}(H_4)=\operatorname{diag}(1,2,2,4).
\label{eq:action-smith}
\end{equation}
Consequently,
\[
\operatorname{coker}(H_4)\simeq
\Z/2\Z\oplus\Z/2\Z\oplus\Z/4\Z,
\qquad
|\operatorname{coker}(H_4)|=16.
\]
\end{lemma}

\begin{proof}
The greatest common divisor of the entries is one; that of the nonzero
\(2\times2\) minors is two; that of the \(3\times3\) minors is four; and
\(|\det H_4|=16\). The successive quotients of these determinantal divisors
are \(1,2,2,4\), which prove \cref{eq:action-smith}. The product of the
invariants gives the order of the cokernel.
\end{proof}

The norm of the section \(\alpha\) over those sixteen sheets
produces \(\alpha^{16}\); the self-scaling \(\varphi\) acts
only once on the basis, not once per sheet. Hence
\cref{eq:action-determinantal-reader}. The certified comparisons
\begin{equation}
\alpha^{16}<10^{-34}
<\varphi\alpha^{16}<10^{-33}
\label{eq:action-decade-bounds}
\end{equation}
prove, without introducing the SI unit, that the decimal order is
exactly \(-34\). The terminal evaluation begins with
\begin{equation}
\Sigma_H
=1.046243838104756580184229208303089\ldots\times10^{-34}.
\label{eq:action-sigma-value}
\end{equation}
The substitution of the kernel by three copies with exponent \(16^3\), or the repeated application of \(\varphi\) on each sheet, would alter the genealogy. Both operations constitute falsification criteria rather than equivalent conventions. The causal chain that is preserved is therefore,
\[
\APP\to\TRIT\to\TPK\to U_{12}^{\rm sign}
\to\alpha\to\Sigma_H
\to10^{-34}\U_S\to
\{\hbar_{\rm pre},\hbar_{\rm ret}\}.
\]
This volume does not rederive \(\alpha\), but develops the metrological geometry
subsequent to the selection of the type and decade by
\(\Sigma_H\). The integer \(-34\) is neither a mantissa nor an adjustment to the International
System, but the terminal valuation of the local functional.

\section{Oriented sections of the action line}

The realizations before and after the return constitute two sections
of the same action line:

\begin{equation}
\hbar_{\rm pre}=H_5\,10^{-34}\U_S,\qquad
\hbar_{\rm ret}=\eta_{\rm ret}\,10^{-34}\U_S.
\label{eq:action-twoway}
\end{equation}

The common decade determines the line, while the oriented
information is expressed through

\begin{equation}
R_{\rm act}=\frac{\hbar_{\rm pre}}{\hbar_{\rm ret}}
=\frac{H_5}{\eta_{\rm ret}},
\qquad
\xi_{\rm act}=\frac12\log R_{\rm act}.
\label{eq:action-rapidity}
\end{equation}

The first quantity provides a multiplicative coordinate and the second
an additive coordinate of rapidity type. They do not represent two independent
Planck constants, but the two orientations of the same section
with respect to the return.

If \(h_a=2\pi\hbar_a\), \(a\in\{\mathrm{pre},\mathrm{ret}\}\), a quantity proportional to \(h_a^{1/2}\) changes as \(R_{\rm act}^{1/2}\), one proportional to \(h_a^{-1/2}\) changes as \(R_{\rm act}^{-1/2}\), and a quotient in which the action cancels is a bidirectional invariant.

  \endgroup
% Distribución causal exacta del delta Ley 9; contenido conservado sin síntesis.
\Needspace{28\baselineskip}
\section{Causal bifurcation subsequent to \texorpdfstring{\(\alpha\)}{alpha HMT}}
\label{sec:l9s102:bifurcacion-accion-angular}

The genealogy of the action and the angular pair forms a bifurcated diagram.
The regional branch produces the reduced mantissa \(\eta_{\mathrm{ret}}\); the
determinantal branch produces the decimal order \(\varepsilon_H=-34\). Both
converge on the action line. The angular pair is constructed in a sister
branch from \(\alpha\) and
\(\eta_{\mathrm{ret}}\):
\begin{equation}
\begin{tikzcd}[column sep=small,row sep=normal]
\alpha,\pi,\varphi
  \arrow[r] \arrow[d]
& H_5,\mathscr C_\pi
  \arrow[d]
&& \alpha,\varphi,H_4
  \arrow[d] & {}\\
x_A,D_A,\rho_\pi
  \arrow[r]
& \eta_{\mathrm{ret}}
  \arrow[r] \arrow[d]
& \hbar_{\mathrm{ret}}\in\mathcal L_S
& \Sigma_H,\varepsilon_H=-34
  \arrow[l] & {}\\
& (A_{\deg},C^*_{\deg})
  \arrow[r]
& (A_{\mathrm{rad}},C^*_{\mathrm{rad}})
  \arrow[r]
& (q_+,q_-)
  \arrow[r]
& V_{90,120}.
\end{tikzcd}
\label{diag:l9s102:dag-accion-angular}
\end{equation}
The expository order places the Planck derivation before the complete
angular publication; the diagram simultaneously preserves the exact
mathematical dependence of each output.

\section{Local cokernel and derivation of the absolute order \texorpdfstring{\(10^{-34}\)}{10^-34}}
\label{sec:l9s102:orden-accion}

The Smith normal form of the local dodecaphase block is
\begin{equation}
 \operatorname{SNF}(H_4)=\operatorname{diag}(1,2,2,4),
 \qquad
 \operatorname{coker}(H_4)
 \cong\mathbb Z/2\mathbb Z\oplus\mathbb Z/2\mathbb Z
 \oplus\mathbb Z/4\mathbb Z.
 \label{eq:l9s102:snf-h4}
\end{equation}
The local cokernel has order \(16\). For the constant section
\(s_\alpha(g)=\alpha\), the local norm and an action of
\(\varphi\) produce
\begin{equation}
 \boxed{
 \Sigma_H=\varphi\mathcal N_{G_H}(s_\alpha)
 =\varphi\alpha^{16}.}
 \label{eq:l9s102:sigma-h}
\end{equation}
The exact inequalities
\begin{equation}
 \alpha^{16}<10^{-34}
 <\varphi\alpha^{16}<10^{-33}
 \label{eq:l9s102:orden-34}
\end{equation}
determine
\begin{equation}
 \boxed{
 \varepsilon_H=\left\lfloor\log_{10}\Sigma_H\right\rfloor=-34.}
 \label{eq:l9s102:epsilon-h}
\end{equation}
 
\chapter{Degree-five regional character and positive action measure}
\begingroup
\renewcommand{\chapter}[2][]{}%
\chapter{Quotient measure and positive character of action: canonical Hilbertization, regional--dodecaphasic defect and functorial law on paths}
\label{chap:p3-final:35:medida_accion_positiva}
\label{p3-20260826:chap:medida-accion-positiva}
\index{measure!emissions catalog}
\index{action!positive cocycle}
\index{normalized trace}

The preceding section isometrically identifies a rank-three positive sector in the oriented interface \(t=7\), the dodecaphase, and the Witt incidence module. The following construction derives the other coefficient of the angular law, \(5/468\), and proves why the two terms combine through a positive sum. The question requires three operations to be distinguished:
\begin{enumerate}
  \item count microscopic presentations;
  \item count distinct emissions after forming the TPK quotient;
  \item assign a action additive to a path.
\end{enumerate}
The first two operations produce different measures. The third is
characterized by a universal property, not by the subsequent proximity of
a decimal value.

We shall denote by \(\Gtr\) the directed graph of \(1944\) states of the
enriched-path reader. Its free path category will be written
\(\mathsf P(\Gtr)\). The local definition of its edges and their action is
developed in the corresponding section.

\section{Seed space, quotient of emissions and regional fibers}

Let
\[
 \Omega_{\rm seed}
 =\bigl((\mathbb Z/9\mathbb Z)^2\times D_4\bigr)^2,
 \qquad
 |\Omega_{\rm seed}|=104\,976,
\]
the executable set of seed pairs of the finite sender, and let
\[
 F:\Omega_{\rm seed}\twoheadrightarrow\mathcal U_6,
 \qquad
 |\mathcal U_6|=468,
\]
the application that preserves the distinct sextuple emission. The exact count of
fibers is
\[
\begin{array}{c|c|c}
|F^{-1}(U)|&\#\{U\text{ of that size}\}&\text{supplied seeds}\\
\hline
144&432&62\,208\\
432&18&7\,776\\
1\,944&18&34\,992.
\end{array}
\]
The three contributions sum to \(104\,976\). Therefore, \(F\) is not a
regular covering: different emissions have different numbers of
microscopic presentations.

The microfiber of \(\pi\) contains five distinct emissions. Four have
multiplicity \(144\), and one has multiplicity \(432\). This datum suffices
to prove that the uniform measure on seeds and the uniform measure
on emissions are not the same.

\section{\(2\) exact measures on the quotient of emissions and words}

If each seed receives weight \(1/104\,976\), the image measure under \(F\) is
\[
 \nu_{\rm seed}(U)
 =\frac{|F^{-1}(U)|}{104\,976}.
\]
Its possible atoms are
\[
 \frac1{729},\qquad\frac1{243},\qquad\frac1{54}.
\]
\Needspace{5\baselineskip}
In particular,
\[
 \nu_{\rm seed}(\mathscr R_\pi^{\rm micro})
 =\frac{4\cdot144+432}{104\,976}
 =\boxed{\frac7{729}}.
\]

If the elementary state is instead a distinct emission
\(U\in\mathcal U_6\), the diagonal algebra of observables is
\[
 \mathcal D_{468}=\ell^\infty(\mathcal U_6)
\]
and its normalized count state is
\[
 \tau_{468}(f)
 =\frac1{468}\sum_{U\in\mathcal U_6}f(U).
\]
Let \(\Pi_\pi^{(468)}\) be the characteristic function of the five regions of the
microfiber. Then
\[
 \boxed{\tau_{468}(\Pi_\pi^{(468)})=\frac5{468}.}
\]
The exact difference between both readings is
\[
 \frac5{468}-\frac7{729}
 =\frac{41}{37\,908}.
\]
They are not two approaches. The first measure counts presentations; the second counts emission classes.

\begin{theorem}[Uniqueness of the uniform measure of the catalog]
\label{p3-20260826:thm:unicidad-medida-468}
The unique probability measure on \(\mathcal U_6\) invariant under every
relabeling of its \(468\) atoms is the uniform measure. Equivalently,
the unique positive unital trace on \(M_{468}(\mathbb C)\) is
\[
 \tau_{468}(A)=\frac1{468}\operatorname{Tr}A.
\]
\end{theorem}

\begin{proof}
Transpositions take any atom to any other. Invariance
forces all of them to have the same mass, and normalization fixes it at
\(1/468\). In the matrix version, unitary invariance equates the value
of all rank-one projectors. Spectral decomposition and
linearity reconstruct the normalized trace.
\end{proof}

The decisive hypothesis is explicit: the reader works with the
\emph{quotient catalogue}. The theorem does not claim that the uniform measure on
seeds becomes uniform when the fibers are forgotten. The symmetry used
is the naturality of the finite unlabeled object, not an ergodic property
attributed without proof to the microscopic dynamics.

\section{Canonical Hilbertization of the quotient}

The relationship between seeds and emissions can be formulated unambiguously.
Let
\[
 H_{\rm seed}=\ell^2(\Omega_{\rm seed}),
 \qquad
 H_U=\ell^2(\mathcal U_6).
\]
The emitter determines
\[
 \boxed{
 J_F\delta_U
 =|F^{-1}(U)|^{-1/2}
 \sum_{\omega\in F^{-1}(U)}\delta_\omega.
 }
\]
Distinct fibers are disjoint and each column has norm one; therefore,
\[
 J_F^*J_F=I_{H_U}.
\]
Moreover, for every observable \(f\) of the catalogue,
\[
 M_{f\circ F}J_F=J_FM_f.
\]
The image \(J_FH_U\) is exactly the subspace of functions that are constant
on each fiber.

If \(F^\sharp:H_U\to H_{\rm seed}\) denotes the map induced by
precomposition, without normalization,
then
\[
 J_F
 =F^\sharp\bigl((F^\sharp)^*F^\sharp\bigr)^{-1/2}.
\]
It is the isometric factor of the polar decomposition of that map. This
expression proves its naturality under isomorphisms and relabelings of the
emitter. It does not imply a false functoriality under arbitrary composition of
nonregular surjections: the normalization depends on the fibers of the
specific map.

In the language of measures, the elevation of the quotient trace that is
uniform within each fiber equals
\[
 \widetilde\tau(\omega)
 =\frac1{468\,|F^{-1}(F\omega)|}.
\]
Each fiber receives total mass \(1/468\). The inverse compensation by
multiplicity, and not an informal invocation of orbits, is the exact proof.

\begin{corollary}[Canonicity of \(5/468\)]
\label{p3-20260826:cor:canonicidad-cinco-468}
Once the regional reader is defined on distinct emissions, the mass
of the microfiber of \(\pi\) is forced:
\[
 \tau_{468}(\Pi_\pi^{(468)})
 =\frac{\operatorname{rank}\Pi_\pi^{(468)}}{468}
 =\frac5{468}.
\]
\end{corollary}

\section{Positive defect between the regional and dodecaphasic measures}

Write
\[
 H_{\rm ang}=H_U\otimes V_{11}.
\]
The internal scalar
\[
 \Delta_4
 =\frac{\pi-e\log\pi}{270}
 \left(1+\frac\pi{729}\right)
 =0.0001111964226921402\ldots
\]
is constructed with the already fixed HMT readings of \(\pi\) and \(e\), and with the
integers \(270\) and \(729\) of the ledger. Here it is a pre-dimensional
mathematical input: neither units nor physical observables intervene.

The internal construction of the scalar does not by itself determine which operator
must multiply. To keep this premise visible, we declare:
\begin{description}
  \item[\HDef --- sheet-change defect.] The dodecaphase coordinate of the
  sheet-change defect is \(\Delta_4P_3\).
\end{description}
\HAdd will fix the additive form of the action but does not select \HDef: for any
\(c\ge0\), the same argument would admit the pair
\((\Pi_\pi^{(468)},cP_3)\). Proving \HDef directly from a total transition
of the enriched topological phase kernel is an additional condition.

We define the Kronecker sum
\[
 \boxed{
 D_{\rm switch}^{\rm dod}
 =\Pi_\pi^{(468)}\otimes I_{11}
 +\Delta_4 I_{468}\otimes P_3.
 }
\]
The two summands are positive, commute, and act on different factors.
Therefore,
\[
 D_{\rm switch}^{\rm dod}\ge0.
\]
The intertwiner \(U_W\) of the
\cref{thm:entrelazador-dodecafase-witt} produces the Witt chart:
\[
\begin{aligned}
 D_{\rm switch}^{\rm Witt}
 &=(I_{468}\otimes U_W)
 D_{\rm switch}^{\rm dod}
 (I_{468}\otimes U_W)^*\\
 &=\Pi_\pi^{(468)}\otimes I_{E_{30}}
 +\Delta_4I_{468}\otimes Q_3.
\end{aligned}
\]
Thus, the same positive operator is read in the dodecaphase and in the
incidence module.

The oriented isometric chain also identifies the active sectors.
\[
 I_{L_{\rm or}}\otimes P_G,\qquad P_3,\qquad Q_3.
\]
No equivalence of the complete complements of the three
modules is asserted: \(W_{GK}\) is a partial equivalence of the selected
ranges. This is exactly the information required to identify
the second contribution to the trace.

\begin{theorem}[Exact Scope of Entanglement and Action Induction]
\label{p3-20260826:thm:alcance-induccion-accion}
Let \(W_{GK}\) and \(U_W\) be the partial isometries of the
\cref{thm:identificacion-isometrica-tres-proyectores}. Then:
\begin{enumerate}
  \item for every \(c\geq0\),
  \[
   W_{GK}^{*}(cP_3)W_{GK}
   =c(I_{L_{\rm or}}\otimes P_G),
   \qquad
   U_W(cP_3)U_W^{*}=cQ_3;
  \]
  thus, entanglement carries the coefficient, but does not select it;
  \item if \(P_{\rm sgn}\) and \(P_{\rm std}\) are the central projectors
  of the decomposition
  \[
   P_3V_{11}\big|_{H_*}
   \simeq \operatorname{sgn}_{H_*}\oplus\operatorname{std},
   \qquad
   \operatorname{rank}P_{\rm sgn}=1,
   \quad
   \operatorname{rank}P_{\rm std}=2,
  \]
  every positive self-adjoint operator supported on \(P_3V_{11}\) and
  \(H_*\)-equivariant has the form
  \[
   D_{a,b}=aP_{\rm sgn}+bP_{\rm std},
   \qquad a,b\geq0.
  \]
  Even impose the trace of \(\Delta_4P_3\) leaves the family
  \[
   a+2b=3\Delta_4.
  \]
  The isotropy \(a=b\) is not deduced from the residual symmetry \(S_3\);
  \item the uniform probability of the catalog,
  \(\mu_U(U)=1/468\), is the canonical trace of the finite quotient, and its
  compensated lift is determined by \(J_F\). It is not the
  image measure of the uniform probability on seeds, whose atoms are
  \(1/729,1/243,1/54\). In particular,
  \[
   \mu_U(\mathscr R_\pi^{\rm micro})=\frac5{468},
   \qquad
   F_\#\mu_{\rm seed}(\mathscr R_\pi^{\rm micro})=\frac7{729};
  \]
  \item the complete law
  \[
   \delta_*=\frac5{468}+\frac3{11}\Delta_4,
   \qquad
   \lambda_*=1-\delta_*,
   \qquad
   a_*(e)=g(e)+\lambda_*s(e)
  \]
  is a theorem relative to the five clauses \HTr--\HAdd--\HDef--\HRes--\HLoc. The
  executable engine of the reduced topological phase kernel does not induce, by directed cover, the branched operator
  \Gtr: the associated dynamics returns to the identity after \(54\) iterations,
  while each state of
  \Gtr has four distinct successors and no self-edge.
\end{enumerate}
Consequently, the preceding results prove the canonicity of the isometric
identifications, of the trace on the quotient catalogue, and of the additive character
within their premises. The stronger assertion that a total transition
of the enriched Topological Phase Kernel necessarily induces \(\Delta_4P_3\), the uniform measure, and the
local law still requires the additional dynamical theorem specified in the
induction analysis.
\end{theorem}

\begin{proof}
The identities in the first assertion are obtained by multiplying
\(W_{GK}^{*}W_{GK}=I_{L_{\rm or}}\otimes P_G\),
\(W_{GK}W_{GK}^{*}=P_3\), and \(U_WP_3U_W^{*}=Q_3\) by the scalar \(c\).
For the second assertion, the restriction to \(H_*\simeq S_3\) is the multiplicity-free
sum of two nonisomorphic irreducible representations. Schur's lemma
diagonalizes every equivariant self-adjoint operator on those two blocks;
positivity is equivalent to \(a,b\geq0\). Its normalized trace is
\((a+2b)/11\), which proves the residual freedom even after the
trace has been fixed. The full action of \(A_5\) on the irreducible triplet would reduce
the family to \(cP_3\), but it would not fix the value of \(c\) either.

The third point follows from the fibre census: \(F\) has multiplicities \(144,432,1944\), so the image measure weights an emission by its number of presentations. By contrast, the trace \(\tau_{468}\) assigns the same weight to every class; \(J_F\) realizes exactly its normalized lift. The fourth point is the logical concatenation already proved: \HTr fixes the quotient trace, \HAdd the sum of traces, \HDef declares \(\Delta_4P_3\), \HRes defines the residual capacity, and \HLoc carries it to the edges. The finite enumeration proves the obstruction between the reduced kernel and \Gtr: a directed cover would preserve the outgoing star, and a semiconjugacy of the identity return would require self-edges in the image, contradicting both censuses of \Gtr.
\end{proof}

The complete spectrum of \(D_{\rm switch}^{\rm dod}\) is
\[
\begin{array}{c|c}
\text{eigenvalue}&\text{multiplicity}\\ \hline
0&(468-5)(11-3)=3704\\
\Delta_4&(468-5)3=1389\\
1&5(11-3)=40\\
1+\Delta_4&5\cdot3=15.
\end{array}
\]
Consequently,
\[
\begin{aligned}
 \tau_{468\cdot11}(D_{\rm switch})
 &=\tau_{468}(\Pi_\pi^{(468)})+\Delta_4\tau_{11}(P_3)\\
 &=\boxed{\frac5{468}+\frac3{11}\Delta_4}.
\end{aligned}
\]
The same identity holds in the Witt chart. In the oriented-sector chart, the term
\(3/11\) is the trace of \(P_G\) on the homogeneous space of eleven branches.

\section{Universal property of the positive action functional}

The word \emph{action} is used in an additive sense: the concatenation of two
routes receives the sum of their increments. To separate this content from a
chosen formula, we first define the category in which the law will be unique.

\begin{definition}[Positive Action Character HMT]
\label{p3-20260826:def:caracter-accion-hmt}
A character local of action on the product
regional--dodecafasico is one map
\[
 \delta:
 M_{468}(\mathbb C)_+\oplus M_{11}(\mathbb C)_+
 \longrightarrow\mathbb R_{\ge0}
\]
which satisfies:
\begin{enumerate}
  \item exact additivity on the product cone;
  \item homogeneidad positiva;
  \item invariance under independent unitary conjugation in the two
  factors;
  \item bifactorial normalization,
  \[
  \delta(I_{468},0)=\delta(0,I_{11})=1.
  \]
\end{enumerate}
\end{definition}

Additivity translates concatenation; invariance eliminates the
coordinates; bifactorial normalization fixes a scale in each summand.
This is not a unital functional on the direct-sum algebra: its unit
is \((I_{468},I_{11})\) and additivity gives
\(\delta(I_{468},I_{11})=2\). No clause contains the value of an
angle.

\begin{theorem}[Uniqueness of the positive character and relative residual capacity]
\label{p3-20260826:thm:ley-accion-hmt}
En the class of the definition preceding there exists a unique character positive of
action HMT\@. It is
\[
 \boxed{\delta(A,B)=\tau_{468}(A)+\tau_{11}(B).}
\]
Applied to the defect declared by \HDef,
\((\Pi_\pi^{(468)},\Delta_4P_3)\), it yields
\[
 \delta_*=\frac5{468}+\frac3{11}\Delta_4.
\]
If, in addition, \HRes is declared, according to which the self-dual capacity for sheet
change is normalized to one and the defect reduces it through
\(\lambda=1-\delta\), the residual capacity within \HTr--\HAdd--\HDef--\HRes is
\[
 \boxed{
 \lambda_*
 =1-\delta_*
 =1-\frac5{468}-\frac3{11}\Delta_4.
 }
\]
\end{theorem}

\begin{proof}
The separability of additivity separates the factors:
\[
 \delta(A,B)=\delta(A,0)+\delta(0,B).
\]
Each restriction is positive, homogeneous, and invariant under conjugation. For
a rank-one projector, unitary invariance gives a common value.
Additivity over a spectral decomposition forces each restriction
to be a multiple of the trace. The bifactorial normalization fixes those multiples at
\(1/468\) and \(1/11\). The sum of traces follows. The last formula
applies the declared normalization of residual capacity.
\end{proof}

Numerically,
\[
 \boxed{\lambda_*=0.9892859130191415\ldots},
 \qquad
 0<\lambda_*<1.
\]
Positivity already follows from
\(0<\Delta_4<10^{-3}\) and \(5/468<1\). It does not depend on a
metrological comparison.

\section{Positive Action Functor on the Category of Compatible Paths}

Recall the directed graph \(\Gtr\) and its free category
\(\mathsf P(\Gtr)\). For an edge
\(e:x\to y\), define
\[
 g(e)=\mathbf1_{\{d_y\ne d_x\}},
 \qquad
 s(e)=\mathbf1_{\{m_y\ne m_x\}},
\]
where \(g\) records a turn and \(s\) a change of sheet. The arrow from the scalar \(\lambda_*\) to the edges does not follow from the trace theorem. It is declared through the clause \HLoc: a turn has unit cost, a change of sheet has cost \(\lambda_*\), and the two indicators are added when they occur on the same edge. Under \HLoc, the local action is
\[
 \boxed{a_*(e)=g(e)+\lambda_*s(e).}
\]
For a route \(\gamma=e_1\cdots e_n\),
\[
 \mathcal A_*(\gamma)
 =\sum_{j=1}^na_*(e_j)
 =\#\operatorname{giros}(\gamma)
 +\lambda_*\#\operatorname{cambios\ de\ hoja}(\gamma).
\]
For the empty path \(\mathbf1_x\), one fixes
\(\mathcal A_*(\mathbf1_x)=0\). Then
\[
 \boxed{
 \mathcal A_*(\gamma_1\gamma_2)
 =\mathcal A_*(\gamma_1)+\mathcal A_*(\gamma_2).
 }
\]
Therefore,
\[
 \mathcal A_*:\mathsf P(\Gtr)\longrightarrow(\mathbb R_{\ge0},+)
\]
is a monoidal functor, or additive path cocycle.

The indicators \(g\) and \(s\) are insensitive to the direction
of traversal. Orientation is not lost: it resides in the line
\(L_{\rm or}\) and in the star involution, whereas positive cost resides in
the projectors. Conflating the two data would introduce negative signs into
the action.

For \(\beta>0\), the associated transfer operator is
\[
 (\mathcal L_{\beta,*}f)(y)
 =\sum_{e:x\to y}e^{-\beta a_*(e)}f(x).
\]
All its weights are strictly positive. Since the potential depends only
on the two local indicators, it commutes with the coarse action
\((\mathbb Z/3\mathbb Z)^3\) and respects the decomposition into \(27\)
Bloch sectors used in the following section.

A finite unweighted deterministic dynamics acts through permutation matrices and has singular values equal to one. The cocycle \(\mathcal A_*\) turns the inventory of paths into a transfer operator that preserves the positive cone, that is, into a matrix with strictly positive weights on the permitted edges. That property does not by itself prove contraction, simplicity of the dominant eigenvalue, or spectral separation; these conclusions require a normalization or an additional spectral argument.

\ifdefined\HMTInternalTraceability
\section{Induction by the quotient and dynamic extension}

The frase \emph{inducida by the kernel topological of phase enriquecido} has two sentidos that must mantenerse
separate.

\begin{enumerate}
  \item \textbf{Induction through the quotient object.} The executable emitter determines \(F\), its polar factor \(J_F\), the space \(H_U\) and the normalized trace. The integral holonomic state supplies \(\Pi_\pi^{(468)}\), \(P_3\), \(K\), \(B_K\), the Witt frame and \(\Delta_4\). In this categorical sense, \(D_{\rm switch}\) is constructed solely from internal APP--TRIT--TPK data\@.
  \item \textbf{Induction by a total transition.} This reading would require
  a finite transition executable on all coordinates of
  \(\mathcal X_{\mathrm{TPK}}^{\mathrm{enr}}\) and a proof that its dynamic
  pushforward measure
  preserves \(J_FH_U\) and produces the potential \(a_*\). That total transition
  does not belong to the executable domain constructed in this section.
\end{enumerate}

The \cref{p3-20260826:thm:ley-accion-hmt} and the construction of caminos cierran the first
reading by means of five clausulas explicit of the reader:
\begin{description}
  \item[\HTr --- quotient trace.] When the presentation multiplicity
  remains in the forgotten fiber, a distinct emission is an atom and the
  catalog state is the normalized trace of its finite algebra.
  \item[\HAdd --- action cocycle.] The defect of a path is the positive,
  bifactorially normalized, additive, and invariant character of the
  regional--dodecaphasic product.
  \item[\HDef --- defecto of cambio of sheet.] The second argumento of the character is
  \(\Delta_4P_3\).
  \item[\HRes --- residual capacity.] The self-dual capacity is normalized to
  one and the defect reduces it according to \(\lambda=1-\delta\).
  \item[\HLoc --- local action.] On the edges of \Gtr, fix
  \(a(e)=g(e)+\lambda s(e)\), and the action of a route is the sum over its
  edges.
\end{description}
\HTr forces \(5/468\). \HAdd forces the sum of traces and excludes mixed
terms. \HDef fixes the scalar occupying the switch coordinate. \HRes transforms
the defect into residual capacity, and \HLoc transports that capacity to the edges.
Thus, \(\delta_*\) is a consequence of \HTr--\HAdd--\HDef, \(\lambda_*\) is a consequence of
\HTr--\HAdd--\HDef--\HRes, and \(\mathcal A_*\) is a consequence of
\HTr--\HAdd--\HDef--\HRes--\HLoc. \HTr and \HAdd possess universal characterizations; \HDef, \HRes, and
\HLoc are internal constitutive identifications; they are not deduced from the
preceding finite transition.

\section{Families Compatible with Partial Symmetries}

\subsection{Symmetry of the \(1/2\) cycle and compatible family of measures}

Interchanging the two halves of an emission decomposes the \(468\)
emissions into \(234\) pairs. Both the uniform measure on the catalogue and the
image measures of seeds are invariant under that interchange. So too is
any assignment constant on each pair. Invariance of the return
alone does not prove uniformity.

\subsection{Ifmetria of \texorpdfstring{\Gtr}{Gamma tr} and family of costis}

The \(7776\) edges of \Gtr split into \(288\) orbits of size
\(27\) under \((\mathbb Z/3\mathbb Z)^3\). The \(2160\) sheet-change
edges form \(80\) orbits:
\[
 80=5\cdot16.
\]
Symmetry allows different costs to be assigned to those orbits. The common
scalar cost comes from \HAdd, not from cardinality.

\subsection{Scalar freedom of the \(2^{\mathrm a}\) coordinate}

For every \(c\ge0\), the operator
\[
D_c=\Pi_\pi^{(468)}\otimes I_{11}
+cI_{468}\otimes P_3
\]
is positive and the unique character of \HAdd gives
\[
\tau(D_c)=\frac5{468}+\frac3{11}c.
\]
The formula internal of \(\Delta_4\) constructs a distinguished candidate, but
the universal property of the trace no demonstrates by itself
\(c=\Delta_4\). That equality is precisely the content of \HDef.

\subsection{Positivity, linearity, and terms of order \(2\)}

The functions
\[
 1-u-v+cuv,\qquad(1-u)(1-v),\qquad e^{-(u+v)}
\]
may share normalization and first order. The absence of \(uv\)
follows from exact additivity. Without \HAdd, the linear law would
not be unique.

\section{Existence Hypothesis, Positivity, Naturalness, and Dynamic Extension}

The result depends on the following verifiable conditions:
\begin{enumerate}
  \item the catalogue does not contain \(468\) emissions or the microfibre does not have
  rank five;
  \item the fiber census is not
  \(432\times144,\ 18\times432,\ 18\times1944\);
  \item \(J_F\) is not isometric or does not intertwine quotient observables;
  \item one of \(P_G,P_3,Q_3\) ceases to be a projector of rank three;
  \item there exists another positive, bifactorially normalized, additive and
  invariant character other than the sum of traces;
  \item a total transition of the enriched topological phase kernel, once constructed, induces a
  switch coefficient distinct from \(\Delta_4\);
  \item \(D_{\rm switch}\) adquiere to negative eigenvalue;
  \item the accumulated action ceases of be additive under concatenation.
\end{enumerate}

The certificates reconstruct the fibres, the two measures, the polar factor, the defect spectrum, the edge orbits and the ablations. The value of \(\lambda_*\) appears at the end as a consequence of these objects. \fi \endgroup
% Distribución causal exacta del delta Ley 9; contenido conservado sin síntesis.
\section{Regional character of degree \texorpdfstring{\(5\)}{5} and reduced action mantissa}
\label{sec:l9s102:eta-ret}

The central block
\begin{equation}
 \mathcal B_c=\begin{pmatrix}7&2\\2&7\end{pmatrix}
 \label{eq:l9s102:bc-accion}
\end{equation}
has eigenvalues \(9\) and \(5\). The length \(4\) of the step
orbits \(3\), the rank \(12\) of the dodecaphasic cycle, and
\(104976/1944=54\) determine the character
\begin{equation}
 \boxed{
 H_5(\alpha,\varphi)
 =\frac12\sqrt{\frac{1000\alpha}{\varphi}}-\alpha
 +\frac9{16}\alpha^2-\frac59\alpha^3
 +\frac7{48}\alpha^4-\frac1{54}\alpha^5.}
 \label{eq:l9s102:h5}
\end{equation}
The nonadic completion constructs the analytic coordinate and the regional determinant.
\begin{equation}
 x_A=\frac{50\pi}{9}\alpha,
 \qquad
 D_A=e^{-2x_A}
 =\exp\!\left(-\frac{100\pi}{9}\alpha\right).
 \label{eq:l9s102:xa-da}
\end{equation}
The microfiber supplies \(\rho_\pi=169\), and its absolutely
convergent series is
\begin{equation}
 \mathscr C_\pi(\alpha)
 =169\alpha^6+\frac{D_A\alpha^7}{1-D_A\alpha}.
 \label{eq:l9s102:cpi}
\end{equation}
The reduced mantissa after return is then constructed by
\begin{equation}
 \boxed{
 \eta_{\mathrm{ret}}
 =H_5-\frac{90}{\pi}\mathscr C_\pi(\alpha)
 \in\mathbb R_{>0}.}
 \label{eq:l9s102:eta-ret}
\end{equation}
Its right-hand side contains exclusively objects generated previously
by APP--TRIT--TPK and by the closure of \(\alpha\).

 
\chapter{Self-dual derivation of the Planck constant and its action line}
\begingroup
\renewcommand{\chapter}[2][]{}%
\chapter{Self-dual derivation of the Planck constant: temporal periodicity of \(108\) phases, action and Compton wavelength}
\label{chap:p3-final:36:planck_autodual}
\label{chap:planck-autodual}

\section{Compton realization associated with the temporal periodicity of \(108\) phases}

On the cycle of \cref{eq:realizacion-circular-ct108}, fix a positive
action section \(\hbar\). The associated angular quantum is
\begin{equation}
E_{108}=\hbar\omega_{108},
\qquad
m_{108}=\frac{E_{108}}{c^2}
=\frac{\hbar\omega_{108}}{c^2}.
\label{eq:cuanto-ct108}
\end{equation}
\(G\) has not yet been used. The geometry of the cycle and angular
quantization suffice to fix its two Compton readers.

\begin{theorem}[Identity between the temporal periodicity of \(108\) phases and the Compton length]
\label{thm:ct108-compton}
For the quantum \eqref{eq:cuanto-ct108},
\begin{equation}
\lambdabarC(m_{108})=\frac{\hbar}{m_{108}c}=R_{108},
\qquad
\lambdaC(m_{108})=\frac{h}{m_{108}c}=C_{108}.
\label{eq:identidad-ct108-compton}
\end{equation}
Therefore, the cycle radius is its reduced Compton wavelength, and the
length of one revolution is its full Compton wavelength.
\end{theorem}

\begin{proof}
Substituting \(m_{108}=\hbar\omega_{108}/c^2\),
\[
\frac{\hbar}{m_{108}c}
=\frac{\hbar c^2}{\hbar\omega_{108}c}
=\frac{c}{\omega_{108}}=R_{108}.
\]
Since \(h=2\pi\hbar\), the second equality is \(2\pi R_{108}=C_{108}\).
\end{proof}

The coincidence is not a numerical fit. It is a homogeneous identity: it holds
for all \(\hbar,c,t_0>0\) and depends only on the same period determining
the frequency, the cycle length, and the quantum of action.

\section{Gravitational Transducer and Self-Dual Selection Equation}

The gravitational line HMT--MD contains the transducer
\begin{equation}
G(\theta)=\theta\frac{c^5t_0^2}{\hbar},
\qquad \theta>0.
\label{eq:transductor-g-planck-anexo}
\end{equation}
Homogeneity is exact, but by itself does not select \(\theta\). The
CT108 cycle supplies an additional internal condition: the gravitational reader
must return the same radial scale as the geometric and Compton
readers.

We define the reduced gravitational radius
\begin{equation}
\rg(\theta):=\frac{G(\theta)m_{108}}{c^2}.
\label{eq:radio-gravitatorio-reducido}
\end{equation}
Substituting \eqref{eq:cuanto-ct108} and
\eqref{eq:transductor-g-planck-anexo},
\begin{equation}
\rg(\theta)
=\theta\frac{\pi}{54}\ell_0.
\label{eq:rg-theta-lineal}
\end{equation}

\begin{theorem}[Self-dual selector of the temporal periodicity of \(108\) phases]
\label{thm:selector-autodual-ct108}
The equation
\[
\rg(\theta)=R_{108}
\]
has a unique positive solution:
\begin{equation}
\boxed{\thetaStar=\left(\frac{54}{\pi}\right)^2.}
\label{eq:theta-star}
\end{equation}
\end{theorem}

\begin{proof}
By \eqref{eq:realizacion-circular-ct108} and
\eqref{eq:rg-theta-lineal}, the equation is
\[
\theta\frac{\pi}{54}\ell_0
=\frac{54}{\pi}\ell_0.
\]
Since \(\ell_0>0\), simplification yields
\(\theta=(54/\pi)^2\). The function
\(\theta\mapsto\rg(\theta)\) is strictly increasing on
\(\RR_{>0}\), so there is no other positive solution.
\end{proof}

\begin{corollary}[Internal Planck Unit]
\label{cor:unidad-planck-interna}
With \(G_*=G(\thetaStar)\), the Planck sections satisfy
\begin{equation}
\ellP:=\sqrt{\frac{\hbar G_*}{c^3}}=R_{108},
\qquad
\tP:=\sqrt{\frac{\hbar G_*}{c^5}}=\frac1{\omega_{108}},
\label{eq:planck-long-time-ct108}
\end{equation}
\begin{equation}
\mP:=\sqrt{\frac{\hbar c}{G_*}}=m_{108},
\qquad
\EP:=\mP c^2=E_{108}.
\label{eq:planck-mass-energy-ct108}
\end{equation}
In particular,
\begin{equation}
\boxed{
R_{108}=\lambdabarC=\rg=\ellP,
\qquad
C_{108}=\lambdaC=2\pi\rg=2\pi\ellP.}
\label{eq:tres-lectores-planck}
\end{equation}
\end{corollary}

\begin{proof}
\(G_*=(54/\pi)^2c^5t_0^2/\hbar\) is substituted. Then
\[
\ellP=\frac{54}{\pi}ct_0=R_{108},
\qquad
\tP=\frac{54}{\pi}t_0=\frac1{\omega_{108}},
\]
and
\[
\mP=\frac{\pi}{54}\frac{\hbar}{c^2t_0}
=\frac{\hbar\omega_{108}}{c^2}=m_{108}.
\]
The rest follows from \cref{thm:ct108-compton}.
\end{proof}

\begin{remark}[Three readers, one section]
Circular geometry publishes \(R_{108}\) and \(C_{108}\); Compton quantization
publishes \(\lambdabarC\) and \(\lambdaC\); gravitational realization
publishes \(\rg\) and \(2\pi\rg\). In \(\thetaStar\) those three pairs are
coordinates of a single section. This is the precise reason why the
Planck unit ceases to appear as an accidental dimensional combination
within this formalization.
\end{remark}

\section{Equivalence of normalizations on the self-dual identity}

The transducer \eqref{eq:transductor-g-planck-anexo} admits two ways of
fixing the elementary scale, which must not be confused.

\begin{enumerate}[label=\textup{(\roman*)}]
\item If \(\theta=1\) is fixed, then
\(\sqrt{\hbar G/c^3}=\ell_0\). The elementary transition \(ct_0\) is identified
with the Planck length of that chart.
\item If the \emph{complete orbit} CT108 is required to have Planck radius,
the condition selects \(\thetaStar=(54/\pi)^2\) and
\(R_{108}=\ellP\).
\end{enumerate}

They are not rival values of \(G\). They are two gauges of the same
homogeneous equation: one assigns the Planck unit to the elementary step; the other,
to the radius of the complete cycle. The annex uses the latter because its thesis concerns
the identity of the CT108 closure.

\section{Gravitational conventions: reduced radius and Schwarzschild radius}

The gravitational notation is fixed in an unambiguous manner:
\[
\rg=\frac{Gm}{c^2},
\qquad
\rs=\frac{2Gm}{c^2}=2\rg.
\]
The preceding theorem uses \(\rg\), not \(\rs\). In the usual convention of
Planck mass, the equality \(\hbar/(mc)=Gm/c^2\) is produced in \(m=\mP\),
whereas
\begin{equation}
\frac{\hbar}{mc}=\frac{2Gm}{c^2}
\quad\Longleftrightarrow\quad
m=\frac{\mP}{\sqrt2}.
\label{eq:cruce-schwarzschild-factor-dos}
\end{equation}
For this reason the phrase “gravitational scale” designates, in the rector thesis, the reduced
radius. The conventional horizon radius remains a distinct reader.

\section{Bifurcation of the action line}

The internal line of action possesses two oriented sections:
\begin{equation}
\hbarpre=H_5\,10^{-34}\mathcal U_S,
\qquad
\hbarret=\eta_{\mathrm{ret}}\,10^{-34}\mathcal U_S,
\qquad
\Ract=\frac{H_5}{\eta_{\mathrm{ret}}}.
\label{eq:hbar-pre-ret-anexo}
\end{equation}
The decade \(10^{-34}\) fixes the absolute scale; the quotient \(\Ract\)
transports the oriented difference. This separation makes it possible to ask whether
Planck self-duality survives in both coordinates.

\begin{theorem}[Reciprocal Covariance of Action and Gravity]
\label{thm:accion-gravedad-reciproca}
For \(a\in\{\mathrm{pre},\mathrm{ret}\}\), let
\[
G_a(\theta)=\theta\frac{c^5t_0^2}{\hbar_a},
\qquad
m_{108,a}=\frac{\hbar_a\omega_{108}}{c^2}.
\]
Then
\begin{equation}
\frac{G_a(\theta)m_{108,a}}{c^2}
=\theta\frac{\pi}{54}\ell_0
\label{eq:invariante-pre-ret-planck}
\end{equation}
is independent of \(a\). Moreover,
\begin{equation}
\frac{m_{108,\mathrm{pre}}}{m_{108,\mathrm{ret}}}=\Ract,
\qquad
\frac{G_{\mathrm{pre}}}{G_{\mathrm{ret}}}=\Ract^{-1},
\label{eq:reciprocidad-pre-ret}
\end{equation}
and the Planck length
\(\sqrt{\hbar_aG_a/c^3}=\sqrt\theta\,\ell_0\) is the same in both
orientations.
\end{theorem}

\begin{proof}
The factors \(\hbar_a\) cancel in the product \(G_am_{108,a}\).
The two ratios in \eqref{eq:reciprocidad-pre-ret} follow directly
from the definitions. The invariance of the Planck length follows from
\(\hbar_aG_a=\theta c^5t_0^2\).
\end{proof}

A symmetric parametrization makes the geometry visible:
\begin{equation}
\hbar_{\mathrm{geom}}:=\sqrt{\hbarpre\hbarret},
\qquad
\xi:=\frac12\log\Ract,
\qquad
\hbarpre=\hbar_{\mathrm{geom}}e^\xi,
\quad
\hbarret=\hbar_{\mathrm{geom}}e^{-\xi}.
\label{eq:rapidez-accion-pre-ret}
\end{equation}
Bidirectional memory is lodged in \(\xi\); Planck geometry is
lodged in the invariant product. This allows a real difference between the
two action coordinates without requiring a difference in
propagation speed.

\begin{remark}[Physical scope of the theorem]
The reciprocity of \cref{thm:accion-gravedad-reciproca} is exact within the transducer. To identify \(G_{\mathrm{pre}}\) and \(G_{\mathrm{ret}}\) with two oriented gravitational measurements, the protocol, the observable, and the symmetry relating them must be specified. The theorem proves structural compatibility; by itself, it neither anticipates a Lorentz violation nor an observable variation of \(G\).
\end{remark}
 \label{mdg:sec:identidad-autodual-planck-ct108-rev3}
\index{Planck!self-dual identity}
\index{calendar of order 108! circular realization}
\index{Compton!reduced and complete readings}

The action line constructed in the preceding section admits a
circular realization compatible with temporal periodicity of order 108.
This realization coordinates three objects that already possess their own types: the
action section, the geometry of one revolution, and the gravitational
transduction. The resulting coordination provides an internal Planck
identity and prepares the positive line on which the general mass
operator will subsequently act.

\subsection{Oriented action letters and declared circular realization}

Let
\[
 \mathcal L_S^+,\quad \mathcal L_T^+,\quad \mathcal L_C^+,
 \quad \mathcal L_L^+,\quad \mathcal L_M^+,\quad \mathcal L_G^+
\]
the positive half-lines of action, time, velocity, length, and mass.
The pre-return and post-return coordinates are written
\begin{equation}
 \hbar_{\rm pre}=H_5\,10^{-34}\mathcal U_S,
 \qquad
 \hbar_{\rm ret}=\eta_{\rm ret}\,10^{-34}\mathcal U_S,
 \qquad
 R_{\rm act}=\frac{\hbar_{\rm pre}}{\hbar_{\rm ret}}
 =\frac{H_5}{\eta_{\rm ret}}.
 \label{mdg:eq:planck-ct108-cartas-pre-ret-rev3}
\end{equation}
The subscript \(a\in\{\mathrm{pre},\mathrm{ret}\}\) shall designate either
of the two charts. In each one the exact distinction is preserved
\begin{equation}
 h_a=2\pi\hbar_a .
 \label{mdg:eq:planck-ct108-h-hbar-rev3}
\end{equation}

Let us fix an elementary duration \(t_0\in\mathcal L_T^+\), an internal
velocity \(c_{\rm int}\in\mathcal L_C^+\), and
\(\ell_0=c_{\rm int}t_0\). The declared circular realization of the
temporal periodicity of order 108 is determined by
\begin{equation}
 T_{108}=108t_0,
 \qquad
 \omega_{108}=\frac{2\pi}{T_{108}},
 \qquad
 R_{108}=\frac{c_{\rm int}}{\omega_{108}}
 =\frac{54}{\pi}\ell_0,
 \label{mdg:eq:planck-ct108-realizacion-circular-rev3}
\end{equation}
and by the length of the full turn
\begin{equation}
 C_{108}=2\pi R_{108}=c_{\rm int}T_{108}=108\ell_0.
 \label{mdg:eq:planck-ct108-perimetro-rev3}
\end{equation}
Thus, \(R_{108}\) is the radius of the realization and \(C_{108}\) its
perimeter. The angular frequency, the cyclic frequency and the two
action constants are related through
\(\omega_{108}=2\pi/T_{108}\) and \(h_a=2\pi\hbar_a\).

\subsection{Correspondence between the order-108 periodicity and Compton readings}

In the chart \(a\), the circular quantum and its mass section are
\begin{equation}
 E_{108,a}=\hbar_a\omega_{108},
 \qquad
 m_{108,a}=\frac{E_{108,a}}{c_{\rm int}^{2}}
 =\frac{\hbar_a\omega_{108}}{c_{\rm int}^{2}}.
 \label{mdg:eq:planck-ct108-cuanto-circular-rev3}
\end{equation}

\begin{theorem}[Identity between circular periodicity and Compton readings]
\label{mdg:thm:planck-ct108-compton-rev3}
For each chart \(a\in\{\mathrm{pre},\mathrm{ret}\}\), the reduced and complete
Compton readings satisfy
\begin{equation}
 \boxed{
 \bar\lambda_{\rm C}(m_{108,a})
 =\frac{\hbar_a}{m_{108,a}c_{\rm int}}=R_{108},
 \qquad
 \lambda_{\rm C}(m_{108,a})
 =\frac{h_a}{m_{108,a}c_{\rm int}}=C_{108}.}
 \label{mdg:eq:planck-ct108-identidad-compton-rev3}
\end{equation}
\end{theorem}

\begin{proof}
Substituting
\eqref{mdg:eq:planck-ct108-cuanto-circular-rev3} produces
\[
 \frac{\hbar_a}{m_{108,a}c_{\rm int}}
 =\frac{c_{\rm int}}{\omega_{108}}=R_{108}.
\]
The second equality follows by applying simultaneously
\eqref{mdg:eq:planck-ct108-h-hbar-rev3} and
\eqref{mdg:eq:planck-ct108-perimetro-rev3}.
\end{proof}

The theorem is homogeneous in
\((\hbar_a,c_{\rm int},t_0)\): the same periodicity determines the radius,
the perimeter, and the two Compton charts. The circular realization is the
declared geometric primitive; the preceding equalities are its
exact consequences.

\subsection{Circular selector and internal Planck unit}

The gravitational line is coordinated with the chart \(a\) through the homogeneous
transducer
\begin{equation}
 G_a(\theta)=\theta\,
 \frac{c_{\rm int}^{5}t_0^2}{\hbar_a}
 \in\mathcal L_G^+,
 \qquad \theta\in\mathbb R_{>0}.
 \label{mdg:eq:planck-ct108-transductor-rev3}
\end{equation}
Its reduced gravitational radius, evaluated on the circular quantum, is
\begin{equation}
 r_{{\rm g},a}(\theta)
 :=\frac{G_a(\theta)m_{108,a}}{c_{\rm int}^{2}}
 =\theta\frac{\pi}{54}\ell_0.
 \label{mdg:eq:planck-ct108-radio-gravitatorio-rev3}
\end{equation}

\begin{proposition}[Planck Circular Selector]
\label{mdg:prop:selector-circular-planck-ct108-rev3}
Within the circular realization
\eqref{mdg:eq:planck-ct108-realizacion-circular-rev3}, the condition
\(r_{{\rm g},a}(\theta)=R_{108}\) selects a unique positive coordinate,
\begin{equation}
 \boxed{\theta_{108}^{\circ}=\left(\frac{54}{\pi}\right)^2.}
 \label{mdg:eq:theta-circular-planck-ct108-rev3}
\end{equation}
The coordinate is independent of the chart oriented \(a\).
\end{proposition}

\begin{proof}
The equations
\eqref{mdg:eq:planck-ct108-realizacion-circular-rev3} and
\eqref{mdg:eq:planck-ct108-radio-gravitatorio-rev3} reduce the condition to
\[
 \theta\frac{\pi}{54}\ell_0=\frac{54}{\pi}\ell_0.
\]
The positivity of \(\ell_0\) and the strict monotonicity of the map
\(\theta\mapsto r_{{\rm g},a}(\theta)\) provide the solution and its
uniqueness.
\end{proof}

Let \(G_{108,a}^{\circ}:=G_a(\theta_{108}^{\circ})\). The internal units
associated with this chart satisfy
\begin{align}
 \ell_{{\rm P},a}^{\circ}
 &:=\sqrt{\frac{\hbar_a G_{108,a}^{\circ}}{c_{\rm int}^{3}}}
 =R_{108},
 &
 t_{{\rm P},a}^{\circ}
 &:=\sqrt{\frac{\hbar_a G_{108,a}^{\circ}}{c_{\rm int}^{5}}}
 =\omega_{108}^{-1},
 \label{mdg:eq:planck-ct108-longitud-tiempo-rev3}\\
 m_{{\rm P},a}^{\circ}
 &:=\sqrt{\frac{\hbar_a c_{\rm int}}{G_{108,a}^{\circ}}}
 =m_{108,a},
 &
 E_{{\rm P},a}^{\circ}
 &:=m_{{\rm P},a}^{\circ}c_{\rm int}^{2}=E_{108,a}.
 \label{mdg:eq:planck-ct108-masa-energia-rev3}
\end{align}
The equality of radii is therefore published as
\begin{equation}
 R_{108}=\bar\lambda_{\rm C}=r_{\rm g}=\ell_{\rm P}^{\circ},
 \qquad
 C_{108}=\lambda_{\rm C}=2\pi r_{\rm g}=2\pi\ell_{\rm P}^{\circ}.
 \label{mdg:eq:planck-ct108-tres-lectores-rev3}
\end{equation}

The convention of this theorem is the reduced gravitational radius
\(r_{\rm g}=Gm/c_{\rm int}^{2}\). The Schwarzschild radius satisfies
\(r_{\rm s}=2r_{\rm g}\); if this second convention is adopted, its crossing with
\(\bar\lambda_{\rm C}\) occurs at
\(m=m_{{\rm P},a}^{\circ}/\sqrt2\). The declaration of the radius therefore forms
part of the chart and prevents a factor of two from being inadvertently transferred
between the two identities.

The coordinate \(\theta_{108}^{\circ}\) belongs to the declared circular
specialization. The general gravitational selector, formulated on the
enriched holonomy and its spectral response, retains its residence in
\cref{mdg:sec:selector-gravitatorio-rev11}. Thus, the circular
identity fixes an internal Planck normalization and the general criterion
fixes the conditions that the physical realization of the gravitational
section must satisfy.

\subsection{Pre-return/post-return covariance}

The product of the reciprocal sections of action and gravity preserves the
circular geometry:
\begin{equation}
 \frac{m_{108,\rm pre}}{m_{108,\rm ret}}=R_{\rm act},
 \qquad
 \frac{G_{108,\rm pre}^{\circ}}{G_{108,\rm ret}^{\circ}}
 =R_{\rm act}^{-1},
 \qquad
 \hbar_a G_{108,a}^{\circ}
 =\theta_{108}^{\circ}c_{\rm int}^{5}t_0^2.
 \label{mdg:eq:planck-ct108-covariancia-pre-ret-rev3}
\end{equation}
In particular, \(\ell_{{\rm P},\rm pre}^{\circ}\) and
\(\ell_{{\rm P},\rm ret}^{\circ}\) coincide with \(R_{108}\). The
bidirectional memory resides in the quotient of the masses and in its
gravitational reciprocal; the geometric scale resides in the invariant product.

Reciprocity admits a symmetric parametrization that separates the common
scale from the orientation of the return. Define
\begin{equation}
 \begin{aligned}
 \hbar_{\rm geom}
 &:=\sqrt{\hbar_{\rm pre}\hbar_{\rm ret}},
 &
 G_{108,\rm geom}^{\circ}
 &:=\sqrt{G_{108,\rm pre}^{\circ}G_{108,\rm ret}^{\circ}},
 &
 \xi_{\rm act}&:=\frac12\log R_{\rm act}.
 \end{aligned}
 \label{mdg:eq:planck-ct108-parametrizacion-simetrica-rev3}
\end{equation}
Then
\[
 \hbar_{\rm pre}=\hbar_{\rm geom}e^{\xi_{\rm act}},
 \qquad
 \hbar_{\rm ret}=\hbar_{\rm geom}e^{-\xi_{\rm act}},
\]
whereas the conjugate gravitational sections satisfy
\[
 G_{108,\rm pre}^{\circ}
 =G_{108,\rm geom}^{\circ}e^{-\xi_{\rm act}},
 \qquad
 G_{108,\rm ret}^{\circ}
 =G_{108,\rm geom}^{\circ}e^{\xi_{\rm act}}.
\]
The coordinate \(\xi_{\rm act}\) preserves the pre-return/post-return
orientation; the product \(\hbar_aG_{108,a}^{\circ}\)
preserves geometric scale. This parametrization of the action charts
does not by itself introduce distinct propagation speeds.

\subsection{Split self-dual mass line}

Given a chart \(a\), let the split algebra be
\[
 \mathbb D_{\rm sp}:=\mathbb R[j]/(j^2-1),
 \qquad
 P_\pm:=\frac{I\pm j}{2},
\]
and consider its spectral decomposition
\[
 \mathbb D_{\rm sp}
 =\mathbb R P_+\oplus\mathbb R P_-,
 \qquad
 P_\pm^2=P_\pm,\quad
 P_+P_-=0,\quad
 P_++P_-=I.
\]
The broken conjugation swaps the projectors,
\[
 \overline{uP_++vP_-}=vP_++uP_-,
\]
and its norm is
\[
 N_{\rm sp}(uP_++vP_-):=uv.
\]

For \(m\in\mathcal L_M^+\), define the section
\begin{equation}
 Z_a(m)
 :=\frac{m_{{\rm P},a}^{\circ}}{m}P_+
   +\frac{m}{m_{{\rm P},a}^{\circ}}P_- .
 \label{mdg:eq:planck-ct108-seccion-partida-masa-rev3}
\end{equation}
The unit \(m_{{\rm P},a}^{\circ}\) also determines the involution
\begin{equation}
 \iota_{{\rm P},a}(m)
 =\frac{(m_{{\rm P},a}^{\circ})^2}{m}.
 \label{mdg:eq:planck-ct108-involucion-masas-rev3}
\end{equation}

\begin{theorem}[Split self-dual mass line]
\label{mdg:thm:planck-ct108-recta-partida-masa-rev3}
For every \(m\in\mathcal L_M^+\),
\begin{equation}
 \boxed{
 N_{\rm sp}(Z_a(m))=1,
 \qquad
 \overline{Z_a(m)}
 =Z_a\!\left(\iota_{{\rm P},a}(m)\right).}
 \label{mdg:eq:planck-ct108-norma-conjugacion-rev3}
\end{equation}
Conjugation has a unique positive fixed point,
\[
 m=m_{{\rm P},a}^{\circ},
 \qquad
 Z_a(m_{{\rm P},a}^{\circ})=I.
\]
\end{theorem}

\begin{proof}
The norm is the product of the two reciprocal coordinates:
\[
 \frac{m_{{\rm P},a}^{\circ}}m
 \frac m{m_{{\rm P},a}^{\circ}}=1.
\]
Conjugation interchanges both coordinates, which is equivalent to replacing
\(m\) by \((m_{{\rm P},a}^{\circ})^2/m\). The positive fixed point satisfies
\(m^2=(m_{{\rm P},a}^{\circ})^2\).
\end{proof}

With the speed
\begin{equation}
 x_a(m)=\log\frac{m}{m_{{\rm P},a}^{\circ}},
 \label{mdg:eq:planck-ct108-rapidez-masa-rev3}
\end{equation}
the section and its conjugate are written
\[
 Z_a(m)=e^{-x_a}P_++e^{x_a}P_-,
 \qquad
 \overline{Z_a(m)}
 =e^{x_a}P_++e^{-x_a}P_-.
\]
Their even and odd readers are
\begin{equation}
 Q_a(m):=\frac{e^{x_a}+e^{-x_a}}2=\cosh x_a,
 \qquad
 A_a(m):=\frac{e^{x_a}-e^{-x_a}}2=\sinh x_a.
 \label{mdg:eq:planck-ct108-lectores-par-impar-rev3}
\end{equation}
The first preserves the distance to the self-dual point and the second preserves its orientation.

The reciprocal lattice vectors associated with the same coordinate are
\[
 r_{{\rm g},a}(m)=\frac{G_{108,a}^{\circ}m}{c_{\rm int}^{2}},
 \qquad
 \bar\lambda_{{\rm C},a}(m)=\frac{\hbar_a}{mc_{\rm int}},
\]
and satisfy
\begin{equation}
 \frac{r_{{\rm g},a}(m)}{\bar\lambda_{{\rm C},a}(m)}
 =e^{2x_a(m)},
 \qquad
 \frac{\bar\lambda_{{\rm C},a}(m)}{r_{{\rm g},a}(m)}
 =e^{-2x_a(m)}.
 \label{mdg:eq:planck-ct108-lectores-exponenciales-rev3}
\end{equation}
Conjugation interchanges the two readers; Planck's identity is their unitary equilibrium.

\begin{remark}[Planck identity and electronic closure]
Identity \(m=m_{{\rm P},a}^{\circ}\) fixes the self-dual origin of the
positive mass line. The electronic section fixes a distinct position on
that same line through its own genealogy. The
\cref{mdg:sec:ley-general-masa-accion-autonoma} constructs the universal operator
that transports absolute scale, character, towers, and bidirectional reading;
the \cref{mdg:sec:cierre-electron-two-way-rev11} subsequently realizes its central
rank-one reduction. Thus, the Planck unit supplies the identity of
the coordinate, and the electron supplies the first complete material transition
on it.
\end{remark}

\begin{proposition}[Genealogical realization of the electronic center]
\label{mdg:prop:planck-ct108-realizacion-centro-electron-rev3}
Let
\[
 F_e^+=\mathbb R_{>0}(P_++P_-)
\]
the positive half-line of the rank-one central fiber, let
\(\mathcal G_e\) be the space of enriched electronic
genealogies and let \(\mathcal L_M^+\) be the positive mass
line. A realization morphism
\[
 \mathfrak R_e:
 F_e^+\times\mathcal G_e\longrightarrow\mathcal L_M^+
\]
can send the identity \(I=P_++P_-\) to a mass different from
\(m_{{\rm P},a}^{\circ}\) without altering the
\cref{mdg:thm:planck-ct108-recta-partida-masa-rev3}: the Planck identity fixes the
autodual origin of the chart, while the electronic genealogy fixes a
displacement on that chart.
\end{proposition}

\begin{proof}
If \(\chi_e:\mathcal G_e\to\mathbb R_{>0}\) is the positive character constructed
by the electronic record, the rule
\[
 \mathfrak R_e(\lambda I,g)
 =\lambda\,m_{{\rm P},a}^{\circ}\chi_e(g)
\]
is positive and multiplicative in the genealogical coordinate. For the neutral
genealogy, \(\chi_e=1\), the Planck origin is recovered; for a genealogy with
\(\chi_e(g)\ne1\), the same internal identity occupies a different position of
the line. The concrete electronic closure preserves its antecedent development in the
\cref{mdg:sec:cierre-electron-two-way-rev11}.
\end{proof}
 \section{Autoduality of the Planck internal unit under orientation inversion}

The CT108 selector constructs

\begin{equation}
\theta_{108}=\left(\frac{54}{\pi}\right)^2,
\qquad
\ell_P=\frac{54}{\pi}\ell_0,
\qquad
t_P=\frac{54}{\pi}t_0,
\label{eq:planck-length-time}
\end{equation}

and

\begin{equation}
E_P=\frac{\hbar_{\rm int}}{t_P}
=\frac{\pi\hbar_{\rm int}}{54t_0}.
\label{eq:planck-energy}
\end{equation}

Comparing \cref{eq:boltzmann-clock,eq:planck-energy} gives

\begin{equation}
\boxed{E_P=k_B\Theta_{\rm clk}\log3.}
\label{eq:planck-trit}
\end{equation}

The relation does not depend on a coincidence between units: both sides
follow from the same chronological structure \(108\).

When the gravitational constant of \cref{ch:gravity} is introduced, the Planck family satisfies

\begin{align}
\ell_P^2&=\frac{G\hbar}{c^3}=\theta_{108}\ell_0^2,
\label{eq:planck-area}\\
F_P&=\frac{c^4}{G},
\qquad
P_P=\frac{c^5}{G},
\label{eq:planck-force-power}\\
\rho_P&=\frac{\hbar}{\theta_{108}^2c^5t_0^4}.
\label{eq:planck-density}
\end{align}

Although \(G\) and \(\hbar\) transform reciprocally between the sections before and after return, their product in \(\ell_P^2\) preserves geometry. This invariance determines area as the interface between the geometric, informational and entanglement realizations. \endgroup
% Distribución causal exacta del delta Ley 9; contenido conservado sin síntesis.
\section{Action line and Planck constant}
\label{sec:l9s102:recta-planck}

Let \(\mathcal L_S\) be the real action line with positive basis
\(\mathcal U_S\). The transduction
\begin{equation}
 \tau_S:\mathbb R_{>0}\longrightarrow\mathcal L_S,
 \qquad u\longmapsto u\,10^{\varepsilon_H}\mathcal U_S
 \label{eq:l9s102:transductor-accion}
\end{equation}
publish the sections preceding and following the return:
\begin{equation}
 \hbar_{\mathrm{pre}}=H_5\,10^{-34}\mathcal U_S,
 \qquad
 \boxed{
 \hbar_{\mathrm{ret}}=\eta_{\mathrm{ret}}\,10^{-34}\mathcal U_S.}
 \label{eq:l9s102:hbar-pre-ret}
\end{equation}
The full cycle actions are
\begin{equation}
 h_{\mathrm{pre}}=2\pi\hbar_{\mathrm{pre}},
 \qquad
 h_{\mathrm{ret}}=2\pi\hbar_{\mathrm{ret}}.
 \label{eq:l9s102:h-pre-ret}
\end{equation}
For a phase \(\theta\), the angular and action charts commute:
\begin{equation}
 S(\theta)
 =\hbar_{\mathrm{ret}}\theta_{\mathrm{rad}}
 =\hbar_{\mathrm{ret}}\frac{\pi}{180}\theta_{\deg}
 =h_{\mathrm{ret}}\frac{\theta_{\deg}}{360}.
 \label{eq:l9s102:accion-fase}
\end{equation}

 
\chapter{Phase operator, spectral radius and factorization of dual projection}
\begingroup
\renewcommand{\chapter}[2][]{}%
\chapter{Predimensional angular operator: decomposition into \(27\) sectors, spectral radius and effective composition of the double projection}
\label{chap:p3-final:37:operador_angular}
\label{p3-20260826:chap:operador-angular-dictamen}
\index{angular operator}
\index{angular phase!spectral output}
\index{closure!pre-dimensional}

The two preceding sections have gathered three pieces that previously appeared
separately: an isometric identification between projectors, a quotient measure,
and a positive action cocycle. The present section studies the spectral consequence
of that law and places the result within the complete architecture
of HMT\@.

The separation of statuses is precise. The triple identification is an exact theorem
relative to the new internal variational rule and an oriented anchor.
The measure and the positive character are exact theorems of \HTr--\HAdd; the defect
is relative, in addition, to \HDef, the capacity to \HRes, and the local cost to \HLoc. The
angular value is an interval-certified output of the operator \Gtr. Its
proximity to the canonical conjugate angular coordinate is extraordinary, but the interval for
the difference excludes equality.

\section{Decomposition of the angular operator into \(27\) sectors of \(72\) states}

The all-paths reader \Gtr has \(1944\) microscopic states. The
coarse translation action of \((\mathbb Z/3\mathbb Z)^3\) permits the
decomposition
\[
 1944=27\cdot72.
\]
Each block of \(72\) states is parametrized by
\[
 (r,s,m,d)
 \in(\mathbb Z/3\mathbb Z)^2\times\{0,1\}\times D_4,
\]
where \(m\) is the sheet and \(d\) the preceding direction. From each state
four cardinal continuations depart. The TRIT of the arrival cell
decides the following sheet.

With the action of the \cref{p3-20260826:chap:medida-accion-positiva}, the weight of an edge
is
\[
 \exp[-g(e)-\lambda_*s(e)].
\]
The character transform reduces the operator to complex matrices
\(L_{k_a,k_b}\) of size \(72\times72\). The program enumerates the
edges of each block; it does not receive a precomputed spectral matrix. Power
iteration is retained as an exploration, but the result that
follows no longer rests on it: the spectral radii are isolated by means of a
complete rational certification.

Let \(\rho_{01}\) and \(\rho_{12}\) be the spectral radii of the oriented
sectors \((0,1)\) and \((1,2)\). We define the \emph{band phase}
\[
 \boxed{y_{\rm band}=\log\rho_{12}-\log\rho_{01}.}
\]
For
\[
 \lambda_*
 =1-\frac5{468}-\frac3{11}\Delta_4,
\]
the certification gives
\[
\begin{aligned}
1.460655120862011111
&<\rho_{01}<1.460655120862404024,\\
1.515812008270944328
&<\rho_{12}<1.515812008271195416,
\end{aligned}
\]
and, by interval propagation,
\[
 0.037066226434427529
 <y_{\rm band}<
 0.037066226434862172,
\]
\[
 \boxed{
 2.123738337168943^\circ
 <C_{\rm band}:=\frac{180}{\pi}y_{\rm band}
 <2.123738337193847^\circ.}
\]
The primary quantity is \(y_{\rm band}\), a dimensionless logarithm of spectral
quotient. Degrees appear later when the chart
\(180/\pi\) is applied. The operator first produces an internal phase and only later
an Archimedean angular representation.

\subsection{Exact calculation of the spectral radius of the angular operator}

The blocks are nonnormal. Therefore, a small residual
\(\|Lv-\mu v\|\) does not suffice to prove that \(\mu\) is the eigenvalue of
greatest modulus. The proof used here consists of four steps.

First, twelve structurally zero rows occur in each block. Expanding the
characteristic polynomial along these rows yields a factor
\(z^{12}\) and a principal minor of order sixty. A further
eight structurally zero rows occur in this minor. Therefore,
\[
 \det(zI_{72}-L)=z^{20}\det(zI_{52}-A),
\]
and every nonzero eigenvalue is contained, with multiplicity, in a
reduced block \(A\) of order fifty-two.

Second, the certificate provides rational Gaussian matrices
\(S,R\), with denominator \(10^{16}\). The verifier calculates
\(E=I-RS\) exactly and obtains
\[
 \|E\|_\infty\le8.5664\times10^{-14}
 \quad\hbox{and}\quad
 \|E\|_\infty\le3.7628\times10^{-14}
\]
for the sectors \((0,1)\) and \((1,2)\), respectively. The Neumann
series then proves that \(S\) is invertible.

Third, if \(A_0\) is the rational midpoint of the interval block, the
identity
\[
 S^{-1}AS=(I-E)^{-1}RAS
\]
bounds the distance to \(C_0=RA_0S\) by
\[
 1.25125\times10^{-13}
 \quad\hbox{and}\quad
 5.7037\times10^{-14}.
\]
The Gershgorin discs of \(C_0\), enlarged by those bounds, contain
the discs of the exact similarity. In each block, one disc remains separated
from the other fifty-one; the second Gershgorin theorem assigns it
exactly one eigenvalue. Its lower modulus exceeds the upper modulus of
every other disc. This proves that the isolated eigenvalue is the
spectral radius.

Fourth, \(\pi,e,\log\pi,e^{-1},e^{-\lambda_*}\) and \(\sqrt3\) are not accepted as library decimals. The verifier encloses them by means of Machin's formula, the factorial and logarithmic series with remainder, the alternating series for the exponential, and a quadratic rational bracket for \(\sqrt3\). Verification employs only integers and fractions; the argument is closed by the rational reduction, the Neumann series, and the Gershgorin separation presented in this section. The reproducible decimal expansions to a depth of ten thousand are published, as independent finite witnesses, in the
\hyperref[app:desarrollos-10000-rev9]{appendix of expansions}.

\section{Effective band operator \texorpdfstring{\(T_{\rm band}\)}{Tband} and isolation of the angular phase}

Let
\[
 x=\frac{50\pi}{9}\alpha
 =0.127362829012869966\ldots
\]
and let
\[
 R=\begin{pmatrix}1&0\\0&-1\end{pmatrix}.
\]
The effective bilayer operator associated exclusively with the band phase is
\[
 \boxed{T_{\rm band}=\exp(-xI+y_{\rm band}R).}
\]
At the base of oriented sectors,
\[
 T_{\rm band}
 =\begin{pmatrix}
 e^{-x+y_{\rm band}}&0\\
 0&e^{-x-y_{\rm band}}
 \end{pmatrix}.
\]
We verify
\[
 \det T_{\rm band}=e^{-2x},
 \qquad
 \frac{(T_{\rm band})_{11}}{(T_{\rm band})_{22}}
 =e^{2y_{\rm band}}.
\]
Therefore,
\[
 -\frac9{100\pi}\log\det T_{\rm band}=\alpha,
 \qquad
 \frac{90}{\pi}
 \log\frac{(T_{\rm band})_{11}}{(T_{\rm band})_{22}}
 =C_{\rm band}.
\]

The two recoveries have distinct statuses. The second returns the output
produced by the gap. The first rereads the \(\alpha\) that entered
\(x\); it proves reversibility of the chart, not a second independent
derivation of \(\alpha\).

The exponential form separates two functions. The common factor \(e^{-x}\)
describes the shared contraction of the two sheets. The factors
\(e^{\pm y_{\rm band}}\) distinguish orientations while preserving the determinant. Since
\(I\) and \(R\) commute, the pair
\[
 \left(\det T_{\rm band},
 \frac{(T_{\rm band})_{11}}{(T_{\rm band})_{22}}\right)
\]
uniquely reconstructs \((x,y_{\rm band})\).

\section{Posterior comparison and precision scale}

The following section constructs another phase, the \emph{regularized regional
phase} \(y_{\rm reg}\) ---also denoted there by \(y_*\)---. To
formulate a comparison here with that subsequent result, we anticipate the
notation
\[
 C_{\rm reg}:=\frac{180}{\pi}y_{\rm reg}
 =2.123738338968646^\circ.
\]
This line does not constitute a second definition or an entry of the
spectral operator: it cites the conclusion of
\cref{p3-20260826:thm:contraangulo-regional}, whose construction is developed in the
following section.
The two pairs of objects are not identified:
\[
 (y_{\rm band},C_{\rm band})
 \ne
 (y_{\rm reg},C_{\rm reg}).
\]
The certified interval difference of the linear law is
\[
 -1.799703\times10^{-9}\ {}^\circ
 <C_{\rm band}-C_{\rm reg}
 <-1.774799\times10^{-9}\ {}^\circ,
\]
with relative difference enclosed by
\[
 -8.475\times10^{-10}
 <\frac{C_{\rm band}-C_{\rm reg}}{C_{\rm reg}}
 <-8.356\times10^{-10}.
\]
The law that produces \(C_{\rm band}\) does not receive \(C_{\rm reg}\), CKM, Barbero,
masses, or SI data\@. This target independence makes the proximity a
candidate for strong, falsifiable recognition. The spectral certification
eliminates the former numerical gap; the interval difference, which excludes
zero, now fixes the exact boundary: \(C_{\rm band}\) is a predimensional spectral
output distinct from the regional realization
\(C_{\rm reg}\).

The additional term
\[
 -\frac{\Delta_4^2}{27}
\]
en \(\lambda\) produce
\[
 C_{\rm num}\approx2.12373833894104980{}^\circ,
\]
at \(2.76\times10^{-11}\) degrees of \(C_{\rm reg}\). The integer \(27\) is internal,
since it counts the characters of \((\mathbb Z/3\mathbb Z)^3\), but that
fact does not force the coefficient of the second moment. The exact additivity of
\HAdd excludes quadratic terms. The variant is preserved as an ablation and not
as the principal law.

\section{Functional dependence of the spectral radius with respect to the components of the operator}

All subsequent rows use the same operator; only one input
of the law changes. They are numerical ablation diagnostics; the preceding certified
interval pertains exclusively to the complete law
\HTr--\HAdd--\HDef--\HRes--\HLoc (and, within it, the positive defect is fixed in the
\HTr--\HAdd--\HDef stratum).
\[
\begin{array}{l|c|c}
\text{rule}&\lambda&C_{\rm num}\text{ in degrees}\\ \hline
\text{self-dual}&1&2.082755794784418\\
\text{microfiber only}&1-5/468&2.123621809570142\\
\text{trace }5/468+(3/11)\Delta_4
&0.989285913019141&2.123738337181414\\
4/468\text{ instead of }5/468
&0.991422665155894&2.115535228138553\\
3/12\text{ instead of }3/11
&0.989288440210566&2.123728626433667\\
5/243\text{ instead of }5/468
&0.979393542015659&2.161907060464779\\
5/729\text{ instead of }5/468
&0.993110963140488&2.109064222037112\\
7/729\text{, seed measure}
&0.990367478915522&2.119584299852700.
\end{array}
\]
The ablations demonstrate that the output distinguishes the complete microfiber,
the reduced dimension eleven, and the emission measure. By themselves, they do
not assign a chance probability: such a probability would first require defining a
space of alternative models and a measure on it.

\section{Typed composition of the morphisms Archimedean, exceptional, and dimensional}

The complete construction can be written as a sequence of applications
with visible domains:
\[
\begin{aligned}
\Omega_{\rm seed}
&\xrightarrow{F}\mathcal U_6
\xrightarrow{J_F}H_U,\\
\mathscr R_\pi^{++}
&\xrightarrow{\beta_\pi}P_GH_G,\\
L_{\rm or}\otimes P_GH_G
&\xrightarrow{W_{GK}}P_3V_{11}
\xrightarrow{U_W}Q_3E_{30},\\
(\Pi_\pi^{(468)},\Delta_4P_3)
&\longmapsto D_{\rm switch}
\xrightarrow{\tau}\delta_*
\xrightarrow{\rm \HRes}\lambda_*,\\
\lambda_*
&\xrightarrow{\rm \HLoc}a_*
\xrightarrow{\mathcal L_{1,*}}y_{\rm band}
\xrightarrow{180/\pi}C_{\rm band},\\
(\alpha,y_{\rm band})
&\xrightarrow{\exp(-xI+y_{\rm band}R)}T_{\rm band}.
\end{aligned}
\]
Each arrow declares what it preserves:
\begin{itemize}
  \item \(F\) preserves the emission class and leaves its presentation
  multiplicity in the fibre;
  \item \(J_F\) compensates that multiplicity in the norm;
  \item \(\beta_\pi\) preserves the exceptional position;
  \item \(W_{GK}\) preserves the oriented representation;
  \item \(U_W\) preserves incidence and metric;
  \item the trace forgets coordinates and preserves ranks;
  \item the transfer operator sums paths with the action weight;
  \item the chart \(180/\pi\) represents an internal phase in degrees.
\end{itemize}

\section{Dependency diagram of the angular operator and its publications}

\begin{longtable}{>{\raggedright\arraybackslash}p{.31\textwidth}>{\raggedright\arraybackslash}p{.22\textwidth}>{\raggedright\arraybackslash}p{.37\textwidth}}
\toprule
Object&Mathematical Nature&Contenido demostrado\\
\midrule
\endfirsthead
\toprule
Object&Mathematical Nature&Contenido demostrado\\
\midrule
\endhead
\(U_W=I/6:V_{11}\to E_{30}\)
&exact theorem
&isometry, surjectivity, equivariant incidence, and positive uniqueness\\
\(P_3\leftrightarrow Q_3\)
&exact theorem
&conjugate projectors, rank three, and trace \(3/11\)\\
star \(R_{B_K}\)
&typed exact theorem
&spectrum \(1^{[7]},(-1)^{[4]}\), index \(3/11\), and conjugation obstruction\\
\(P_G\) in \(t=7\)
&finite exact theorem
&filtration \(11\to3\to1\), positive projector, and action \(S_3\)\\
\(\mathscr R_\pi^{++}\leftrightarrow P_G\)
&exact theorem
&\(S_3\)-equivariant bijection by exceptional position\\
oriented isometric identification
&new relative theorem
&unique selector through \(B_K,u_K\), Reynolds, polar, and reduction \(240\to1\)\\
measure \(1/468\)
&exact theorem of \HTr
&unique normalized trace and polar application induced by fibers\\
\(D_{\rm switch}\) y \(\delta_*\)
&theorem exact relativo to \HTr--\HAdd--\HDef
&positivity, spectrum, and trace \(5/468+(3/11)\Delta_4\)\\
\(\lambda_*\)
&theorem internal relativo to \HTr--\HAdd--\HDef--\HRes
&residual capacity obtained from the defect through the rule \HRes\\
\(\mathcal A_*\)
&internal theorem relative to \HTr--\HAdd--\HDef--\HRes--\HLoc
&cocycle additive of paths with cost local \(g+\lambda_*s\)\\
\(T_{\rm band}\) y \(C_{\rm band}\)
&intervalar computational theorem
&reduction \(72\to60\to52\), rational similarity, Gershgorin
isolation and output whose certified distance lies between
\(1.774799\times10^{-9}\) and \(1.799703\times10^{-9}\) degrees of
\(C_{\rm reg}\)\\
\(C_{\rm band}=C_{\rm reg}\)
&refuted statement for the linear law
&the certified interval of the difference excludes zero\\
total transition
\(\mathcal X_{\mathrm{TPK}}^{\mathrm{enr}}\to \Gtr\)
&required global realizer
&\HTr--\HAdd--\HDef--\HRes--\HLoc
defines the relative reader; the complete transition additionally requires
a dynamical pushforward measure on \(\Gtr\)\\
\bottomrule
\end{longtable}

\section{Predimensional action at the center of the double projection of the genealogical state}

Dual projection starts from the same integral holonomic state and produces two readings. In the Archimedean branch, compatible cylinders are refined and their diameters decrease; evaluation fixes a limiting value and forgets part of the genealogy. In the exceptional branch, incidences extend towards Witt, Golay, Niemeier--Leech, \(V^\natural\), the Monster and moonshine.

The new result adds a common object stronger than a coincidence of cardinalities. The eleven-mode module, its rank-three positive sector, and the trace of the action pass from dodecaphase to Witt through \(U_W\). The oriented sector \(t=7\) is incorporated through \(W_{GK}\). The microfiber of \(\pi\) provides the regional projector \(\Pi_\pi^{(468)}\), and the uniform measure on the catalogue fixes its weight.

Positive action thus stands at the center of the two readings:
\[
 \text{region}
 \quad\oplus\quad
 \text{oriented polarization}
 \quad\longmapsto\quad
 \text{positive defect}.
\]
The Archimedean evaluation and the exceptional evaluation remain
different maps; the intertwiner identifies a common sector, it does not
collapse their codomains.

\section{Compatibility of the angular operator with the cylinders of the discrete continuum}

An HMT history determines nested cylinders
\[
 C_1\supset C_2\supset\cdots
\]
with compatible prefixes. When the Archimedean reader satisfies
\[
 \operatorname{diam}(C_n)\longrightarrow0,
\]
the completeness of \(\mathbb R\) fixes at most one value in the intersection.
The new action adds structure before that limit: it distinguishes the cost
of the route, the emission class, and the multiplicity of presentations that the
final value no longer displays.

This is the precise mathematical meaning of the assertion that the continuum is read as a
projection of discrete genealogies. The system of cylinders and its
evaluations is prolonged at every scale; the global coverage theorem proves
\(\operatorname{ev}_{\mathbb R}(\Omega_{\infty,K})=K\) for every compact set
\(K\subset\mathbb R\) and, upon directing the intervals, proves that the image is all of
\(\mathbb R\). The finite theorem in this section controls the angular action
at each level and is integrated into that global coverage; it neither reopens nor
conditions it. The algebraic descent of each operation retains its own
verification.

\section{Compatibility with the exceptional rider, Moonshine, and M-theory screens}

The intertwiner \(U_W\) lands in the incidence module of \(W_{12}\).
From the ternary code, the two exceptional continuations remain
separate:
\[
 C_W\longrightarrow W_{12}
 \xrightarrow[\text{relative to }(D,\ell)]{}W_{24},
\]
and
\[
 C_W\xrightarrow{\ \phi\ } N(A_2^{12})
\xrightarrow{\ v_\alpha\ }\Lambda_{24}
\longrightarrow V_\Lambda
\longrightarrow V^\natural
\longrightarrow\mathbb M
\longrightarrow\text{moonshine}.
\]
Positive action contributes to this branch a transported sector and an
orientation line. The arrows from Leech to moonshine are the classical
constructions applied to the finite object obtained. The Monster action is
realized on \(V^\natural\); the HMT words are not identified here with a
direct module for that action.

The screens \(12\to11\to10\), \(26\to24\), and the T-dual involution form
an exact mathematical interface with the dimensions and dualities used in
M-theory. The full physical identification belongs to Dimensional Mechanics
and is typed in the second part by the additional map to
\(g_{MN}\), \(C_{MNP}\), \(\psi_M\), action, and supersymmetry. It is neither
deferred to another work nor deduced from the Monster action.

\section{Spectral factorization and isolation of the dominant sectors}

The pre-dimensional line is organized into four strata.
\begin{enumerate}
  \item \textbf{Finite kernel.} APP, TRIT, the TPK emitter, the catalog
  \(104\,976\to468\to243\), the five regions of \(\pi\), the
  dodecaphase, \(K\), \(\alpha\), Witt, and the exceptional branch have
  explicit constructions and certificates.
  \item \textbf{Operatorial bridge.} Incidence constructs \(U_W\);
the oriented sector \(t=7\) constructs \(P_G\); \(B_K\) and \(u_K\) select
the oriented isometric identification; the three rank-three projectors
are identified relative to that rule.
  \item \textbf{Action extension.}
\HTr fixes the measure \(1/468\), \HAdd fixes the additive
character, and \HDef declares the coupling \(\Delta_4P_3\).
\HRes transforms the defect into capacity, and \HLoc carries
capacity to the edges. The five clauses produce the local law.
  \item \textbf{Spectral consequence.} \Gtr produces
  \(T_{\rm band}\) and a certified predimensional angular phase, with
  relative difference between \(8.356\times10^{-10}\) and
  \(8.475\times10^{-10}\) with respect to \(C_{\rm reg}\).
\end{enumerate}

This permits a strong and typed closure of this route: the finite emitter of the
infinite coinductive section, the incidence, the identification of
projectors, the measure, and the action are closed
relative to the declared primitives, the variational rule
\(\mathcal E_B\), and the action clauses \HTr--\HAdd--\HDef--\HRes--\HLoc. \HAdd
selects the linear form; \HDef selects its second coordinate; \HRes and \HLoc
specify how it is read as a cost. The global theorem covering all of
\(\mathbb R\), induction from a total transition of
\(\mathcal X_{\mathrm{TPK}}^{\mathrm{enr}}\), and physical equivalence with
M-theory are formulated as extensions of the preceding results.

\section{Compatibility between the matrix, angular, and lattice realizations}

The constructed representations simultaneously satisfy the following compatibility conditions:
\begin{enumerate}
  \item five regions of \(\pi\), with three positivas, one negativa and one
  transversal;
  \item \(P_G,P_3,Q_3\) as positive projectors and \(R_{B_K}\) as an
  orientation involution;
  \item the oriented homogeneous sector \(t=7\), not the heterogeneous quotient of \(t=8\);
  \item the line \(L_{\rm or}\) in the triple isometric identification;
  \item the relative character of the variational rule that selects
  \(H_*,c_*,r_*\);
  \item \(5/468\) on emissions and \(7/729\) on seeds, without
  interchanging them;
  \item the linearity of the character as a consequence of \HAdd and the coefficient
  \(\Delta_4\) as content separated from \HDef;
  \item \((y_{\rm band},C_{\rm band})\) as a certified spectral output, separate from
\((y_{\rm reg},C_{\rm reg})\), the canonical regional realization of the following section;
  \item \(\alpha\) as input to \(x\), so that the determinant
  reading is reversible and not independent;
  \item the branches \(W_{12}\to W_{24}\) and
  \(A_2^{12}\to\Lambda_{24}\) as separate constructions;
  \item the relative reader \HTr--\HAdd--\HDef--\HRes--\HLoc distinguished from
  a total transition, whose dynamic image measure constitutes an
  additional construction.
\end{enumerate}

These conditions are not editorial notes. They are part of the logical structure of the result. \endgroup

\chapter{Equivalence of the global and analytic charts of the angular pair}
The positive measure after return is obtained before the counterangle. The determinantal branch determines the action valuation and both constructions converge in the section of the dimensional line. Only then are the angle, its oriented complement and the characters \(q_{+}\) and \(q_{-}\) formed, evaluating the incidences already generated by the kernel.
% Adaptador único del capítulo 38 del sucesor 102.
% Sustituye exclusivamente la jerarquía de encabezados. Los cuerpos de los
% dos fuentes teoremáticas científicos y del delta Ley 9 se conservan línea a línea,
% reordenados conforme al DAG causal aprobado. Véase MAPA_SEGMENTOS_CAPITULO_38.tsv.

% HMT_SOURCE_SEGMENT_BEGIN id=B00 source=colaboracion/parte_iii/source/public_final/base_83/c30_body.tex body_lines=1-2
\label{chap:p3-83-30-angulo_contraangulo}

% HMT_SOURCE_SEGMENT_END id=B00
% HMT_SOURCE_SEGMENT_BEGIN id=A00 source=colaboracion/parte_iii/source/public_final/ascensos_20260826/16f_realizacion_analitica_contraangulo.tex body_lines=1-4
\label{p3-20260826:chap:realizacion-analitica-contraangulo}
\index{counter-angle}
\index{microfiber of pi@microfiber of \(\pi\)!analytic character}
\index{determinantal character}
% HMT_SOURCE_SEGMENT_END id=A00
% HMT_SOURCE_SEGMENT_BEGIN id=D00 source=manuscrito/sucesor_102/deltas_ley9/c38_par_angular_cartas_y_canales.tex body_lines=1
% Distribución causal exacta del delta Ley 9; contenido conservado sin síntesis.
% HMT_SOURCE_SEGMENT_END id=D00

\section{Intrinsic angular phase \texorpdfstring{\(\alpha\)}{alpha HMT}, nonadic lift \texorpdfstring{\(10^3\)}{10^3}, and precursor \texorpdfstring{\(x_A\)}{x_A}}
\label{sec:s102:c38:1}

\subsection{Construction of the analytic precursor \texorpdfstring{\(x_A\)}{x_A}, the contraction \texorpdfstring{\(D_A\)}{D_A}, and the regional correction \texorpdfstring{\(\mathscr C_\pi(\alpha)\)}{C_pi(alpha)}}
\label{sec:s102:c38:1:1}

% HMT_SOURCE_SEGMENT_BEGIN id=B01 source=colaboracion/parte_iii/source/public_final/base_83/c30_body.tex body_lines=4-32
\label{p3-83-c30-s1:chap:realizacion-analitica-contraangulo}
\index{conjugate angular coordinate \(C^*\)}
\index{microfiber of pi@microfiber of \(\pi\)!analytic character}
\index{determinantal character}

The finite catalog and an analytic evaluation answer questions of different
types. The catalog determines which regions exist, which share a
return and what information each preserves. An analytic evaluation must
also decide how the length of a word is transformed into degree, how
a branch is prolonged after its return and with which scalar character that
prolongation is read. This chapter explicitly constructs that second
operation.

The construction begins with four data already obtained in the preceding
chapters: the five-region microfiber of \(\pi\), the length six of
each regional word, the dodecaphasic closure of \(\alpha\), and the bilayer
transport associated with the common coordinate
\(A=10^3\alpha\). No value of the conjugate angular coordinate \(C_*^{\rm HMT}\), of the reduced action
component, of the Barbero--Immirzi functional, or of a metrological constant is
used in the evaluation. The result is a convergent regional character,
an exact recurrence, and an oriented phase whose sexagesimal chart reproduces
the published conjugate angular coordinate.

The separation between the finite part and the analytic part is essential. The
integer \(169\) will be a theorem of the catalogue. The series that prolongs it will be the
theorem of an analytic reader whose rules are declared before it is evaluated.
Thus, the additional premises are not concealed within a decimal
coincidence.

% HMT_SOURCE_SEGMENT_END id=B01
\section{Preangular construction of \texorpdfstring{\(\eta_{\mathrm{ret}}\)}{eta_ret} through \texorpdfstring{\(H_5\)}{H5}, \texorpdfstring{\(D_A\)}{D_A} and \texorpdfstring{\(\mathscr C_\pi(\alpha)\)}{C_pi(alpha)}}
\label{sec:s102:c38:2}

\subsection{Independent derivation of \texorpdfstring{\(\eta_{\mathrm{ret}}=H_5-(90/\pi)\mathscr C_\pi(\alpha)\)}{eta_ret=H5-(90/pi)C_pi(alpha)} through the regional analytic reader}
\label{sec:s102:c38:2:1}

% HMT_SOURCE_SEGMENT_BEGIN id=A01 source=colaboracion/parte_iii/source/public_final/ascensos_20260826/16f_realizacion_analitica_contraangulo.tex body_lines=5-38,40-47

The finite catalogue and an analytic evaluation answer questions of different types. The catalogue determines which regions exist, which share a return, and what information each preserves. An analytic evaluation must further determine how the length of a word is transformed into degree, how a branch is prolonged after its return, and by which scalar character that prolongation is read. This section constructs that second operation explicitly.

The construction begins with the five-region microfiber of \(\pi\), the
length six of each regional word, the dodecaphasic closure of \(\alpha\),
and the declared angular insertion
\[
 \iota_\angle(x):=10^3x\,u_\circ,\qquad
 \Adeg:=10^3\alpha .
\]
Here \(u_\circ\) is the sexagesimal base whose complete turn has
coordinate \(360\).
The factor \(10^3\) belongs to that chart normalization and is not deduced from the
arithmetic closure of \(\alpha\). The two-sheet transport is subsequently constructed
from \(\Adeg\). The function evaluated by the reader does not receive as
an argument the counter-angle, the post-return coordinate, the
hexad--octad functional, or a metrological constant. The result is a
convergent regional character, an exact recurrence, and an oriented phase in the
declared sexagesimal chart.

The separation between the finite part and the analytic part is essential. The
integer \(169\) is a theorem of the catalogue. The series that prolongs it belongs
to an analytic reader whose rules are declared before evaluating it.

The phase is denoted by \(y_*\) and, when its regional origin is emphasized, by
\(y_{\rm reg}\); both symbols denote the same reader constructed
below.


The evaluation receives the catalogue of \(468\) emissions, the exact value of
\(\alpha\), \(\varphi=(1+\sqrt5)/2\), the fixed formula of \(H_5\),
the angular chart, and the five rules of the regional reader. With those data, it
computes \(169\), \(D_A\), the recurrence, and \(C^*\). The theorem is exact in
this declared class of readers; it does not use \(C^*\) as an input or assert that
the rules of \(H_5\) are the only possible natural class.

% HMT_SOURCE_SEGMENT_END id=A01
\subsection{Persistence of the return in the \texorpdfstring{\(5\)}{5} regions of the microfiber}
\label{sec:s102:c38:2:2}

% HMT_SOURCE_SEGMENT_BEGIN id=B02 source=colaboracion/parte_iii/source/public_final/base_83/c30_body.tex body_lines=34-155

Let
\[
\mathcal F_\pi=\Psi^{-1}(010211)\subset\mathcal U_6
\]
the microfiber relative to the catalog. Recall its five elements, now
separated into initial block and terminal block:
\[
\begin{aligned}
 &(501,614,498\mid169,272,272),\\
 &(810,923,870\mid169,272,272),\\
 &(810,983,810\mid169,272,272),\\
 &(870,923,810\mid169,272,272),\\
 &(870,923,810\mid418,674,521).
\end{aligned}
\]
We define the return projection.
\[
p_{\rm ret}:\mathbb Z^6\longrightarrow\mathbb Z^3,
\qquad
p_{\rm ret}(u_1,\ldots,u_6)=(u_4,u_5,u_6).
\]

\begin{proposition}[Persistent terminal block]
\label{p3-83-c30-s1:prop:bloque-terminal-persistente}
The multiset \(p_{\rm ret}(\mathcal F_\pi)\) has a unique element of
maximum multiplicity:
\[
b_\pi=(169,272,272),
\qquad
\operatorname{mult}_{\mathcal F_\pi}(b_\pi)=4.
\]
The only remaining terminal block is
\[
b_\pi^-=(418,674,521)
\]
and has multiplicity one.
\end{proposition}

\begin{proof}
The assertion follows by exact enumeration of the five rows of the
catalog that satisfy \(\Psi(U)=010211\). Four rows have terminal
block \((169,272,272)\); the fifth, which carries the mutated return, has
\((418,674,521)\). The maximum multiplicity is four and is attained only
once.
\end{proof}

The oriented difference between the mutated return and the persistent is
\[
 \omega_\pi=b_\pi^- - b_\pi=(249,402,249),
 \qquad
 \langle(1,1,1),\omega_\pi\rangle=900.
\]
This identity separates two objects that must not be conflated. The integer
\(169\) is a coordinate of a terminal datum, that is, a datum of degree
zero. The vector \(\omega_\pi\) can be read as a one-cochain
after orienting the edge joining the two blocks. Calling it a cocycle would additionally
require fixing a cochain complex and verifying that its coboundary vanishes.
The analytic prolongation constructed below does not presuppose that
identification.

The selection of \(169\) does not require privileging the first position of the block. The symmetric group \(\mathfrak S_3\) acts by permuting its three coordinates. The stabilizer of \(b_\pi\) exchanges the two occurrences of \(272\) and fixes the unique coordinate of multiplicity one. Denote that coordinate by \(\operatorname{sing}(b)\) when the block has type \((r,s,s)\), with \(r\ne s\). Equivalently,
\[
\operatorname{sing}(b)
=\sum_{j=1}^3b_j-2\operatorname{moda}(b).
\]
In particular,
\[
\boxed{\operatorname{sing}(b_\pi)=169.}
\]

\begin{definition}[Residual fiber selector]
\label{p3-83-c30-s1:def:selector-residual-pi}
The selector \(\rho_\pi\) successively applies two operations:
\begin{enumerate}
  \item selects the terminal block of maximum multiplicity within
  \(p_{\rm ret}(\mathcal F_\pi)\);
  \item retains the coordinate fixed by the stabilizer of that block.
\end{enumerate}
By \cref{p3-83-c30-s1:prop:bloque-terminal-persistente},
\[
\boxed{\rho_\pi=169.}
\]
\end{definition}

\begin{theorem}[Equivariant uniqueness of return]
\label{p3-83-c30-s1:thm:unicidad-retorno-169}
Among the selectors that
\begin{enumerate}
  \item are invariant under permutations of the five regions;
  \item are equivariant under permutations of the three terminal
  coordinates;
  \item select the block of maximal multiplicity;
  \item retain a coordinate fixed by its stabilizer,
\end{enumerate}
the return value is uniquely determined and is \(169\).
\end{theorem}

\begin{proof}
Invariance under \(\mathfrak S_5\) prevents using the order of the rows. The
third condition selects \(b_\pi\), which is the unique mode by
\cref{p3-83-c30-s1:prop:bloque-terminal-persistente}. Its stabilizer in
\(\mathfrak S_3\) has two orbits on the coordinates: the pair of
values \(272\) and the singleton value \(169\). The fourth condition selects the
second orbit. Equivariance preserves its value when its position is shifted.
\end{proof}

This argument improves the positional reading of the return. The \(169\) is not
«the first coordinate» of a selected row: it is an invariant of the
regional multiset and of the action of its reordering symmetries.
Nor is it selected by global rarity. Among the \(468\) emissions in the
catalogue, \(169\) appears \(84\) times, uniformly distributed with
\(14\) occurrences in each of the six positions. The uniqueness of
\cref{p3-83-c30-s1:thm:unicidad-retorno-169} is therefore relational and regional: it resides
in the combination of microfiber, modal persistence, and stabilizer, not in the
isolated frequency of the integer.

% HMT_SOURCE_SEGMENT_END id=B02
\subsection{Construction of the regional analytic reader and its convergence domain}
\label{sec:s102:c38:2:3}

% HMT_SOURCE_SEGMENT_BEGIN id=B05 source=colaboracion/parte_iii/source/public_final/base_83/c30_body.tex body_lines=261-301

Analytic continuation is fixed by the following rules.

\begin{definition}[Regional Determinantal Reader]
\label{p3-83-c30-s1:def:lector-regional-determinantal}
The determinantal regional reader satisfies:
\begin{enumerate}
  \item \emph{degree compatibility}: cylindrical length is the analytic
  degree;
  \item \emph{renewal decomposition}: the finite return impulse and
  strict continuation belong to distinct additive summands;
  \item \emph{continuation normalization}: after recording the residue
  \(\rho_\pi\) just once, the surviving branch begins at the unit of
  the determinant line;
  \item \emph{concatenation covariance}: each elementary prolongation
  acts through \(\bigwedge^2\mathcal T=D_A\);
  \item \emph{orientation}: in the cardinal convention fixed for APP, the
  regional return belongs to the compensating sector and is subtracted from the fifth-degree
action phase.
\end{enumerate}
\end{definition}

The third rule has mathematical content. The value \(169\) is recorded once as the amplitude of the return; it is not multiplied again at each prolongation. The strict continuation counts the sole line that continues and normalizes it by its unit. This is the distinction between a renewal recurrence and homogeneous iteration of the initial impulse.

The need to declare that normalization can be proved. For every
\(\lambda\in\mathbb R\), the family
\[
F_\lambda(z)
=169z^6+\lambda\frac{D_Az^7}{1-D_Az}
\]
preserves the same finite return and the same determinantal ratio between
consecutive coefficients. The catalogue and the ratio alone do not distinguish
\(\lambda\). The unit of the determinant line fixes
\(\lambda=1\) before \(z=\alpha\) is evaluated. Thus, the normalization
is not obtained by fitting the final value; it forms part of the definition of the
character.

% HMT_SOURCE_SEGMENT_END id=B05
\subsection{Analytical solution of the return recurrence}
\label{sec:s102:c38:2:4}

% HMT_SOURCE_SEGMENT_BEGIN id=B06 source=colaboracion/parte_iii/source/public_final/base_83/c30_body.tex body_lines=303-364

We define the coefficients \(\kappa_n\) by
\[
\kappa_6=169,
\qquad
\kappa_7=D_A,
\qquad
\kappa_{n+1}=D_A\kappa_n\quad(n\ge7).
\]
The first equality is the regional impulse; the next two describe the normalized continuation.

\begin{theorem}[Graduated analytic realization]
\label{p3-83-c30-s1:thm:realizacion-analitica-graduada}
The analytic germ compatible with the rules of
\cref{p3-83-c30-s1:def:lector-regional-determinantal} is unique and equals
\[
\boxed{
\mathscr C_\pi(z)
=169z^6+\frac{D_Az^7}{1-D_Az}.}
\]
Their coefficients satisfy
\[
\kappa_n=D_A^{n-6}\qquad(n\ge7).
\]
\end{theorem}

\begin{proof}
Write
\[
F(z)=169z^6+\sum_{k\ge1}c_kz^{6+k}.
\]
The normalization of continuation and the determinantal covariance give
\[
c_1=D_A,
\qquad
c_{k+1}=D_Ac_k.
\]
By induction, \(c_k=D_A^k\). The geometric sum yields
\[
\sum_{k\ge1}D_A^kz^{6+k}
=\frac{D_Az^7}{1-D_Az}.
\]
No free coefficients remain.
\end{proof}

Since \(0<\alpha<1\) and
\(0<D_A<1\),
\[
0<D_A\alpha<1.
\]
Consequently, \(\mathscr C_\pi(\alpha)\) converges absolutely and its
radius of convergence is \(D_A^{-1}>1\). Numerically,
\[
D_A=0.7751291192471564301888840964802\ldots,
\]
\[
D_A\alpha=0.00565639046986492673885079095469\ldots.
\]
The speed of this contraction explains why the complete sum and its
degree-seven truncation have the same first fifteen digits in the
final chart.

% HMT_SOURCE_SEGMENT_END id=B06
\subsection{Persistence of return in the
\texorpdfstring{\(5\)}{5}
regions and compatibility of the
\texorpdfstring{\(\eta_{\mathrm{ret}}\)}{eta_ret}
component}
\label{sec:s102:c38:2:5}

% HMT_SOURCE_SEGMENT_BEGIN id=A02 source=colaboracion/parte_iii/source/public_final/ascensos_20260826/16f_realizacion_analitica_contraangulo.tex body_lines=49-170

Let
\[
\mathcal F_\pi=\Psi^{-1}(010211)\subset\mathcal U_6
\]
the microfiber relative to the catalog. Recall its five elements, now
separated into initial block and terminal block:
\[
\begin{aligned}
 &(501,614,498\mid169,272,272),\\
 &(810,923,870\mid169,272,272),\\
 &(810,983,810\mid169,272,272),\\
 &(870,923,810\mid169,272,272),\\
 &(870,923,810\mid418,674,521).
\end{aligned}
\]
We define the return projection.
\[
p_{\rm ret}:\mathbb Z^6\longrightarrow\mathbb Z^3,
\qquad
p_{\rm ret}(u_1,\ldots,u_6)=(u_4,u_5,u_6).
\]

\begin{proposition}[Persistent terminal block]
\label{p3-20260826:prop:bloque-terminal-persistente}
The multiset \(p_{\rm ret}(\mathcal F_\pi)\) has a unique element of
maximum multiplicity:
\[
b_\pi=(169,272,272),
\qquad
\operatorname{mult}_{\mathcal F_\pi}(b_\pi)=4.
\]
The only remaining terminal block is
\[
b_\pi^-=(418,674,521)
\]
and has multiplicity one.
\end{proposition}

\begin{proof}
The assertion follows by exact enumeration of the five rows of the
catalog that satisfy \(\Psi(U)=010211\). Four rows have terminal
block \((169,272,272)\); the fifth, which carries the mutated return, has
\((418,674,521)\). The maximum multiplicity is four and is attained only
once.
\end{proof}

The oriented difference between the mutated return and the persistent is
\[
 \omega_\pi=b_\pi^- - b_\pi=(249,402,249),
 \qquad
 \langle(1,1,1),\omega_\pi\rangle=900.
\]
This identity separates two objects that must not be conflated. The integer
\(169\) is a coordinate of a terminal datum, that is, a datum of degree
zero. The vector \(\omega_\pi\) can be read as a one-cochain
after orienting the edge joining the two blocks. Calling it a cocycle would additionally
require fixing a cochain complex and verifying that its coboundary vanishes.
The analytic prolongation constructed below does not presuppose that
identification.

The selection of \(169\) does not require privileging the first position of the block. The symmetric group \(\mathfrak S_3\) acts by permuting its three coordinates. The stabilizer of \(b_\pi\) exchanges the two occurrences of \(272\) and fixes the unique coordinate of multiplicity one. Denote that coordinate by \(\operatorname{sing}(b)\) when the block has type \((r,s,s)\), with \(r\ne s\). Equivalently,
\[
\operatorname{sing}(b)
=\sum_{j=1}^3b_j-2\operatorname{moda}(b).
\]
In particular,
\[
\boxed{\operatorname{sing}(b_\pi)=169.}
\]

\begin{definition}[Residual fiber selector]
\label{p3-20260826:def:selector-residual-pi}
The selector \(\rho_\pi\) successively applies two operations:
\begin{enumerate}
  \item selects the terminal block of maximum multiplicity within
  \(p_{\rm ret}(\mathcal F_\pi)\);
  \item retains the coordinate fixed by the stabilizer of that block.
\end{enumerate}
By \cref{p3-20260826:prop:bloque-terminal-persistente},
\[
\boxed{\rho_\pi=169.}
\]
\end{definition}

\begin{theorem}[Equivariant uniqueness of return]
\label{p3-20260826:thm:unicidad-retorno-169}
Among the selectors that
\begin{enumerate}
  \item are invariant under permutations of the five regions;
  \item are equivariant under permutations of the three terminal
  coordinates;
  \item select the block of maximal multiplicity;
  \item retain a coordinate fixed by its stabilizer,
\end{enumerate}
the return value is uniquely determined and is \(169\).
\end{theorem}

\begin{proof}
Invariance under \(\mathfrak S_5\) prevents using the order of the rows. The
third condition selects \(b_\pi\), which is the unique mode by
\cref{p3-20260826:prop:bloque-terminal-persistente}. Its stabilizer in
\(\mathfrak S_3\) has two orbits on the coordinates: the pair of
values \(272\) and the singleton value \(169\). The fourth condition selects the
second orbit. Equivariance preserves its value when its position is shifted.
\end{proof}

This argument improves the positional reading of the return. The \(169\) is not
«the first coordinate» of a selected row: it is an invariant of the
regional multiset and of the action of its reordering symmetries.
Nor is it selected by global rarity. Among the \(468\) emissions in the
catalogue, \(169\) appears \(84\) times, uniformly distributed with
\(14\) occurrences in each of the six positions. The uniqueness of
\cref{p3-20260826:thm:unicidad-retorno-169} is therefore relational and regional: it resides
in the combination of microfiber, modal persistence, and stabilizer, not in the
isolated frequency of the integer.

% HMT_SOURCE_SEGMENT_END id=A02
\subsection{Application of the regional reader to the conjugate angular momentum pair}
\label{sec:s102:c38:2:6}

% HMT_SOURCE_SEGMENT_BEGIN id=A05 source=colaboracion/parte_iii/source/public_final/ascensos_20260826/16f_realizacion_analitica_contraangulo.tex body_lines=277-317

Analytic continuation is fixed by the following rules.

\begin{definition}[Regional Determinantal Reader]
\label{p3-20260826:def:lector-regional-determinantal}
The determinantal regional reader satisfies:
\begin{enumerate}
  \item \emph{degree compatibility}: cylindrical length is the analytic
  degree;
  \item \emph{renewal decomposition}: the finite return impulse and
  strict continuation belong to distinct additive summands;
  \item \emph{continuation normalization}: after recording the residue
  \(\rho_\pi\) just once, the surviving branch begins at the unit of
  the determinant line;
  \item \emph{concatenation covariance}: each elementary prolongation
  acts through \(\bigwedge^2\mathcal T=D_A\);
  \item \emph{orientation}: in the cardinal convention fixed for APP, the
  regional return belongs to the compensating sector and is subtracted from the fifth-degree
action phase.
\end{enumerate}
\end{definition}

The third rule has mathematical content. The value \(169\) is recorded once as the amplitude of the return; it is not multiplied again at each prolongation. The strict continuation counts the sole line that continues and normalizes it by its unit. This is the distinction between a renewal recurrence and homogeneous iteration of the initial impulse.

The need to declare that normalization can be proved. For every
\(\lambda\in\mathbb R\), the family
\[
F_\lambda(z)
=169z^6+\lambda\frac{D_Az^7}{1-D_Az}
\]
preserves the same finite return and the same determinantal ratio between
consecutive coefficients. The catalogue and the ratio alone do not distinguish
\(\lambda\). The unit of the determinant line fixes
\(\lambda=1\) before \(z=\alpha\) is evaluated. Thus, the normalization
is not obtained by fitting the final value; it forms part of the definition of the
character.

% HMT_SOURCE_SEGMENT_END id=A05
\subsection{Closed sum and publication of the posterior component upon return}
\label{sec:s102:c38:2:7}

% HMT_SOURCE_SEGMENT_BEGIN id=A06 source=colaboracion/parte_iii/source/public_final/ascensos_20260826/16f_realizacion_analitica_contraangulo.tex body_lines=319-380

We define the coefficients \(\kappa_n\) by
\[
\kappa_6=169,
\qquad
\kappa_7=D_A,
\qquad
\kappa_{n+1}=D_A\kappa_n\quad(n\ge7).
\]
The first equality is the regional impulse; the next two describe the normalized continuation.

\begin{theorem}[Graduated analytic realization]
\label{p3-20260826:thm:realizacion-analitica-graduada}
The analytic germ compatible with the rules of
\cref{p3-20260826:def:lector-regional-determinantal} is unique and equals
\[
\boxed{
\mathscr C_\pi(z)
=169z^6+\frac{D_Az^7}{1-D_Az}.}
\]
Their coefficients satisfy
\[
\kappa_n=D_A^{n-6}\qquad(n\ge7).
\]
\end{theorem}

\begin{proof}
Write
\[
F(z)=169z^6+\sum_{k\ge1}c_kz^{6+k}.
\]
The normalization of continuation and the determinantal covariance give
\[
c_1=D_A,
\qquad
c_{k+1}=D_Ac_k.
\]
By induction, \(c_k=D_A^k\). The geometric sum yields
\[
\sum_{k\ge1}D_A^kz^{6+k}
=\frac{D_Az^7}{1-D_Az}.
\]
No free coefficients remain.
\end{proof}

Since \(0<\alpha<1\) and
\(0<D_A<1\),
\[
0<D_A\alpha<1.
\]
Consequently, \(\mathscr C_\pi(\alpha)\) converges absolutely and its
radius of convergence is \(D_A^{-1}>1\). Numerically,
\[
D_A=0.7751291192471564301888840964802\ldots,
\]
\[
D_A\alpha=0.00565639046986492673885079095469\ldots.
\]
The speed of this contraction explains why the complete sum and its
degree-seven truncation have the same first fifteen digits in the
final chart.

% HMT_SOURCE_SEGMENT_END id=A06
\section{Confluence of \texorpdfstring{\(\eta_{\mathrm{ret}}\)}{eta_ret} With the valuation \texorpdfstring{\(\varepsilon_H=-34\)}{epsilon_H=-34} in the magnitude \texorpdfstring{\(\hbar_{\mathrm{ret}}\)}{hbar_ret}}
\label{sec:s102:c38:3}

\subsection{Derivation of the determinantal character of parangular transport}
\label{sec:s102:c38:3:1}

% HMT_SOURCE_SEGMENT_BEGIN id=B04 source=colaboracion/parte_iii/source/public_final/base_83/c30_body.tex body_lines=196-259

The dodecaphasic closure provides \(\alpha\). The completion of the three
nonadic pairs provides the scale \(10^3\), so that
\[
A=10^3\alpha.
\]
Having chosen the Archimedean angular chart, the common coordinate in radians is
\[
x=\frac{\pi A}{180}=\frac{50\pi}{9}\alpha.
\]
Let \(R\) be a real involution of zero trace in dimension two,
\[
R^2=I_2,
\qquad
\operatorname{tr}R=0,
\]
and consider the formal two-layer transport
\[
\mathcal T(x,y)=\exp(-xI_2-yR).
\]
The oriented variable \(y\) has not yet been evaluated. Nevertheless, the action
of \(\mathcal T\) on the determinant line is independent of it:
\[
\begin{aligned}
\det\mathcal T(x,y)
&=\exp\!\left(\operatorname{tr}(-xI_2-yR)\right)\\
&=e^{-2x}.
\end{aligned}
\]
We define
\[
\boxed{
D_A=e^{-2x}=e^{-100\pi\alpha/9}.}
\]
The subscript recalls that this is the determinant of the common contraction
associated with \(A\), not a function of \(C_*^{\rm HMT}\).

\begin{proposition}[Character of the determinant line]
\label{p3-83-c30-s1:prop:caracter-determinante}
The map
\[
\delta_A:(\mathbb N_0,+)\longrightarrow(\mathbb R_{>0},\cdot),
\qquad
\delta_A(k)=D_A^k,
\]
is the unique multiplicative character whose value on an elementary
prolongation is \(D_A\).
\end{proposition}

\begin{proof}
The multiplicativity gives
\[
\delta_A(k)
=\delta_A(\underbrace{1+\cdots+1}_{k\text{ times}})
=\delta_A(1)^k=D_A^k.
\]
The same identity proves existence and uniqueness.
\end{proof}

The determinant remains invariant under conjugation of
\(\mathcal T\) and under the exchange of sheets \(R\mapsto-R\). Therefore, the
tail constructed from \(D_A\) belongs to the common mode. The orientation
will appear in the sign with which the regional character corrects the phase.

% HMT_SOURCE_SEGMENT_END id=B04
\subsection{action character of degree \texorpdfstring{\(5\)}{5} in the construction regional}
\label{sec:s102:c38:3:2}

% HMT_SOURCE_SEGMENT_BEGIN id=B07 source=colaboracion/parte_iii/source/public_final/base_83/c30_body.tex body_lines=366-405

The central block
\[
\mathcal B_c=\begin{pmatrix}7&2\\2&7\end{pmatrix}
\]
has eigenvalues \(9\) and \(5\). The dodecaphasic cycle contributes rank
\(12\); its step-three orbits have length four; and the complete census
contributes the exact quotient
\[
\frac{104976}{1944}=54.
\]
From these invariants the fifth-degree action character is fixed.
\[
\boxed{
H_5(\alpha,\varphi)
=\frac12\sqrt{\frac{1000\alpha}{\varphi}}-\alpha
+\frac9{16}\alpha^2-\frac59\alpha^3
+\frac7{48}\alpha^4-\frac1{54}\alpha^5.}
\]
The four rational coefficients can be written as
\[
\frac9{16}=\frac{\lambda_+}{4^2},
\qquad
\frac59=\frac{\lambda_-}{\lambda_+},
\qquad
\frac7{48}=\frac{\tfrac12\operatorname{tr}\mathcal B_c}{4\cdot12},
\qquad
\frac1{54}=\frac{1944}{104976},
\]
with \((\lambda_+,\lambda_-)=(9,5)\). These identities certify the
provenance of the numbers used. The successive assignment of those
invariants to degrees two, three, four, and five forms part of the analytic
character; it does not follow from the mere existence of the finite catalog.

This precision removes a common ambiguity. A combination of
invariants of the construction is an internal and reproducible definition; its
naturality requires the reader rules specifying order, signs, and
normalization. Here those rules are visible and are subject to the same
negative controls as the rest of the construction.

% HMT_SOURCE_SEGMENT_END id=B07
\subsection{Exact decomposition of the action line into pre-return and post-return sections}
\label{sec:s102:c38:3:3}

% HMT_SOURCE_SEGMENT_BEGIN id=B12 source=colaboracion/parte_iii/source/public_final/base_83/c30_body.tex body_lines=571-617

The fifth-degree character and the post-action component after return
are two distinct values:
\[
H_5
=1.05457181764615446962582182943306\ldots,
\]
\[
\bar h_{\rm HMT}^{\rm ret}
=1.05457181691503906113369058472430\ldots.
\]
The first is the local action before incorporating the regional character; the
second is the oriented action remaining after the return. The correction
that generates the conjugate angular coordinate prevents identifying them.

The metrological comparison must respect this distinction. The first value gives
the mantissa
\[
2\pi H_5
=6.62607014999998792600111034989947\ldots,
\]
while the subsequent action gives
\[
2\pi\bar h_{\rm HMT}^{\rm ret}
=6.62607014540625433351074984759288\ldots.
\]
In the International System, the value
\[
h=6.62607015\times10^{-34}\;\mathrm{J\,s}
\]
is exact by definition \cite{BIPMSI}. Therefore,
\[
2\pi H_5-6.62607015
=-1.2073998889650101\ldots\times10^{-14},
\]
and it is not an exact equality. The proximity of the value before the return is
a remarkable numerical check; it does not authorize replacing the defining constant
of the SI with the subsequent component or confusing the two HMT actions\@.

This result precisely determines the mathematical situation. The regional
reader generates \(C_*^{\rm HMT}\) and, from it, the functional
\(\Gamma_{6\to8}\). The same calculation separates the Planck approximation
produced by \(H_5\) from the reduced action intervening in
the conjugate angular coordinate. The separation does not weaken the construction: it eliminates an
identification incompatible with the exact figures and converts every
comparison into a falsifiable assertion of a well-defined type.

% HMT_SOURCE_SEGMENT_END id=B12
\subsection{Factorization of the determinantal character in the \texorpdfstring{\(2\)}{2} phase charts}
\label{sec:s102:c38:3:4}

% HMT_SOURCE_SEGMENT_BEGIN id=A04 source=colaboracion/parte_iii/source/public_final/ascensos_20260826/16f_realizacion_analitica_contraangulo.tex body_lines=211-275

Dodecaphase closure supplies \(\alpha\). The declared angular insertion
uses the scale \(10^3\), compatible with the grouping of three nonadic
pairs, so that
\[
\Adeg=10^3\alpha\,{}^\circ.
\]
Having chosen the Archimedean angular chart, the common coordinate in radians is
\[
x=: \Arad=\frac{\pi\Adeg}{180}=\frac{50\pi}{9}\alpha.
\]
Let \(R\) be a real involution of zero trace in dimension two,
\[
R^2=I_2,
\qquad
\operatorname{tr}R=0,
\]
and consider the formal two-layer transport
\[
\mathcal T(x,y)=\exp(-xI_2-yR).
\]
The oriented variable \(y\) has not yet been evaluated. Nevertheless, the action
of \(\mathcal T\) on the determinant line is independent of it:
\[
\begin{aligned}
\det\mathcal T(x,y)
&=\exp\!\left(\operatorname{tr}(-xI_2-yR)\right)\\
&=e^{-2x}.
\end{aligned}
\]
We define
\[
\boxed{
D_A=e^{-2x}=e^{-100\pi\alpha/9}.}
\]
The subscript recalls that this is the determinant of the common contraction
associated with \(A\), not a function of the counter-angle.

\begin{proposition}[Character of the determinant line]
\label{p3-20260826:prop:caracter-determinante}
The map
\[
\delta_A:(\mathbb N_0,+)\longrightarrow(\mathbb R_{>0},\cdot),
\qquad
\delta_A(k)=D_A^k,
\]
is the unique multiplicative character whose value on an elementary
prolongation is \(D_A\).
\end{proposition}

\begin{proof}
The multiplicativity gives
\[
\delta_A(k)
=\delta_A(\underbrace{1+\cdots+1}_{k\text{ times}})
=\delta_A(1)^k=D_A^k.
\]
The same identity proves existence and uniqueness.
\end{proof}

The determinant remains invariant under conjugation of
\(\mathcal T\) and under the exchange of sheets \(R\mapsto-R\). Therefore, the
tail constructed from \(D_A\) belongs to the common mode. The orientation
will appear in the sign with which the regional character corrects the phase.

% HMT_SOURCE_SEGMENT_END id=A04
\subsection{Publication of the action character of degree \texorpdfstring{\(5\)}{5} in the parangular chart}
\label{sec:s102:c38:3:5}

% HMT_SOURCE_SEGMENT_BEGIN id=A07 source=colaboracion/parte_iii/source/public_final/ascensos_20260826/16f_realizacion_analitica_contraangulo.tex body_lines=382-418

The central block
\[
\mathcal B_c=\begin{pmatrix}7&2\\2&7\end{pmatrix}
\]
has eigenvalues \(9\) and \(5\). The dodecaphasic cycle contributes rank
\(12\); its step-three orbits have length four; and the complete census
contributes the exact quotient
\[
\frac{104976}{1944}=54.
\]
The current reader declares the nature of fifth-degree action
\[
\boxed{
H_5(\alpha,\varphi)
=\frac12\sqrt{\frac{1000\alpha}{\varphi}}-\alpha
+\frac9{16}\alpha^2-\frac59\alpha^3
+\frac7{48}\alpha^4-\frac1{54}\alpha^5.}
\]
The four rational coefficients can be written as
\[
\frac9{16}=\frac{\lambda_+}{4^2},
\qquad
\frac59=\frac{\lambda_-}{\lambda_+},
\qquad
\frac7{48}=\frac{\tfrac12\operatorname{tr}\mathcal B_c}{4\cdot12},
\qquad
\frac1{54}=\frac{1944}{104976},
\]
with \((\lambda_+,\lambda_-)=(9,5)\). These identities fix the
internal arithmetic representability of the character coefficients. Their
order, signs, and normalization belong to the graded rule of the regional
reader and are applied before any chart of units. The ablations
certify that every term effectively intervenes in the resulting phase.
In particular, neither the SI mantissa of action nor the sexagesimal value of the
counter-angle belongs to the domain of the generating function.

% HMT_SOURCE_SEGMENT_END id=A07
\section{Parangular application \texorpdfstring{\(P(\alpha,\eta_{\mathrm{ret}})\)}{P(alpha,eta_ret)} and bilayer closure \texorpdfstring{\(C^*=2(\eta_{\mathrm{ret}}+\alpha)\)}{C*=2(eta_ret+alpha)}}
\label{sec:s102:c38:4}

\subsection{Oriented phase and conjugate angular coordinate}
\label{sec:s102:c38:4:1}

% HMT_SOURCE_SEGMENT_BEGIN id=B08 source=colaboracion/parte_iii/source/public_final/base_83/c30_body.tex body_lines=407-472

The fifth-degree phase prior to return is
\[
y_5=\frac\pi{90}\bigl(H_5+\alpha\bigr).
\]
The complete regional character produces the correction
\(\mathscr C_\pi(\alpha)\). With the orientation fixed in
\cref{p3-83-c30-s1:def:lector-regional-determinantal}, we define
\[
\boxed{
y_*=\frac\pi{90}\bigl(H_5+\alpha\bigr)
-169\alpha^6
-\frac{D_A\alpha^7}{1-D_A\alpha}.}
\]
This is the primary formula. The sexagesimal unit does not come into play until applying the chart.
\[
C_*^{\rm HMT}=\frac{180}{\pi}y_*.
\]

\begin{theorem}[Conjugate angular coordinate of the regional reader determinantal]
\label{p3-83-c30-s1:thm:contraangulo-regional}
With the dodecaphasic closure of \(\alpha\),
\(\varphi=(1+\sqrt5)/2\) and the rules of the determinantal regional reader fixed, the
phase \(y_*\) and its conjugate angular coordinate are determined without introducing the sought
value. Their evaluation is
\[
\boxed{
y_*=0.03706622646583826398493779409230638\ldots,}
\]
\[
\boxed{
C_*^{\rm HMT}
=2.12373833896864572426195138039336\ldots{}^\circ.}
\]
Therefore, to fifteen decimal places,
\[
\boxed{C_*^{\rm HMT}=2.123738338968646^\circ.}
\]
\end{theorem}

\begin{proof}
\cref{p3-83-c30-s1:thm:unicidad-retorno-169} determines \(169\). The closure of
\(\alpha\) determines \(A\), \(x\), and
\(D_A=e^{-100\pi\alpha/9}\). \cref{p3-83-c30-s1:thm:realizacion-analitica-graduada}
determines the complete series. The formula for \(H_5\) contains only
\(\alpha\), \(\varphi\), and declared structural invariants. Substituting
these data into the expression for \(y_*\) gives the printed value. The decimal value
of \(C_*^{\rm HMT}\) occurs in none of the right-hand sides.
\end{proof}

The reduced action component after return is now as follows
output:
\[
\boxed{
\bar h_{\rm HMT}^{\rm ret}
=\frac{C_*^{\rm HMT}}2-\alpha
=H_5-\frac{90}{\pi}\mathscr C_\pi(\alpha).}
\]
Numerically,
\[
\bar h_{\rm HMT}^{\rm ret}
=1.05457181691503906113369058472430\ldots.
\]
The inverse identity has not been used to introduce \(C_*^{\rm HMT}\); it is
obtained after generating \(y_*\).

% HMT_SOURCE_SEGMENT_END id=B08
\subsection{Conjugate coordinate \texorpdfstring{\(C^*\)}{C*} and reversal of its orientation under the bilayer involution}
\label{sec:s102:c38:4:2}

% HMT_SOURCE_SEGMENT_BEGIN id=A08 source=colaboracion/parte_iii/source/public_final/ascensos_20260826/16f_realizacion_analitica_contraangulo.tex body_lines=420-503

The fifth-degree phase prior to return is
\[
y_5=\frac\pi{90}\bigl(H_5+\alpha\bigr).
\]
The complete regional character produces the correction
\(\mathscr C_\pi(\alpha)\). With the orientation fixed in
\cref{p3-20260826:def:lector-regional-determinantal}, we define
\[
\boxed{
\Crad:=y_*=\frac\pi{90}\bigl(H_5+\alpha\bigr)
-169\alpha^6
-\frac{D_A\alpha^7}{1-D_A\alpha}.}
\]
From now on we also write
\[
 \boxed{y_{\rm reg}:=y_*}
\]
when it is necessary to distinguish this realization from the spectral phase of
\Gtr. The two notations here designate the same canonical regional object.
This is the primary formula. The sexagesimal unit does not intervene until the chart is applied
\[
\boxed{\Cdeg:=\frac{180}{\pi}\Crad.}
\]
The symbols \(\Crad\) and \(\Cdeg\) do not denote two counterangles: they are the
radial and sexagesimal coordinates of the same angular element.

\begin{theorem}[Counter-angle of the regional determinantal reader]
\label{p3-20260826:thm:contraangulo-regional}
Given the dodecaphase closure of \(\alpha\),
\(\varphi=(1+\sqrt5)/2\), the germ \(H_5\) and the rules of the regional
determinantal reader, the phase \(\Crad\) and its counter-angle are determined without
introducing the sought value into the evaluation function. Its evaluation is
\[
\boxed{
\Crad=0.03706622646583826398493779409230638\ldots,}
\]
\[
\boxed{
\Cdeg
=2.12373833896864572426195138039336\ldots{}^\circ.}
\]
Therefore, to fifteen decimal places,
\[
\boxed{\Cdeg=2.123738338968646^\circ.}
\]
\end{theorem}

\begin{proof}
The \cref{p3-20260826:thm:unicidad-retorno-169} determines \(169\). The closure of \(\alpha\) determines \(\Adeg\), \(\Arad\), and \(D_A=e^{-100\pi\alpha/9}\). The \cref{p3-20260826:thm:realizacion-analitica-graduada} determines the complete series. The fixed formula for \(H_5\) contains only \(\alpha\), \(\varphi\), and declared structural invariants. Substituting these data into the expression for \(\Crad\) yields the printed value. The decimal value of the counter-angle appears in none of the right-hand sides of this evaluation. Comparison with a metrological action mantissa is undertaken only after the phase has been obtained and does not alter the generating function.
\end{proof}

Let \(u_\circ\) be the sexagesimal base whose complete turn has coordinate
\(360\), and let us introduce the angular coordinate
\(\alpha_{\angle,{\rm deg}}:=\alpha\). The output posterior to
return is then typed by
\[
\boxed{
\etaRet
=\frac{\Cdeg}2-\alpha_{\angle,{\rm deg}}
=H_5-\frac{90}{\pi}\mathscr C_\pi(\alpha).}
\]
Its radial coordinate is
\[
 \etaRetrad=\frac{\pi}{180}\etaRet
 =\frac{\Crad}{2}-\frac{\pi}{180}\alpha.
\]
Numerically,
\[
\etaRet
=1.05457181691503906113369058472430\ldots.
\]
The inverse identity has not been used to introduce \(\Cdeg\); it is obtained
after generating \(\Crad\). The symbol \(\hbar\) is reserved for the
element of the action line constructed in the following section.

% HMT_SOURCE_SEGMENT_END id=A08
\subsection{Direct theorem \texorpdfstring{\(C^*=2(\eta_{\mathrm{ret}}+\alpha)\)}{C*=2(eta_ret+alpha)} and inverse character of the identity \texorpdfstring{\(\eta_{\mathrm{ret}}=C^*/2-\alpha\)}{eta_ret=C*/2-alpha}}
\label{sec:s102:c38:4:3}

% HMT_SOURCE_SEGMENT_BEGIN id=A13 source=colaboracion/parte_iii/source/public_final/ascensos_20260826/16f_realizacion_analitica_contraangulo.tex body_lines=750-797

The fifth-degree character and the coordinate after return
are two distinct values:
\[
H_5
=1.05457181764615446962582182943306\ldots,
\]
\[
\etaRet
=1.05457181691503906113369058472430\ldots.
\]
The former is the local character before incorporation of the regional character; the
latter is the oriented coordinate remaining after the return. The correction
generating the contra-angle prevents their identification.

The metrological comparison must respect this distinction. The first value gives
the mantissa
\[
2\pi H_5
=6.62607014999998792600111034989947\ldots,
\]
while the posterior coordinate gives
\[
2\pi\etaRet
=6.62607014540625433351074984759288\ldots.
\]
As an external control only, we retain the mantissa \(6.62607015\) of the
defining value of the International System \cite{BIPMSI}. Neither its exponent nor its unit is yet incorporated into the HMT line: the decade is derived in the
\cref{chap:orden-determinantal-accion} and the unit belongs to a later
metrological chart. The comparison of mantissas gives
\[
2\pi H_5-6.62607015
=-1.2073998889650101\ldots\times10^{-14},
\]
and is not an exact equality. The proximity of the pre-return value is
a noteworthy numerical control; it does not authorize replacing the defining SI
constant by the subsequent coordinate or confusing an angular coordinate with
the action element that will be constructed upon incorporating the decade.

This result determines the mathematical situation precisely. The regional
reader generates \((\Crad,\Cdeg)\) and, from that single angular element, the functional
\(\Gamma_{6\to8}\). The same calculation separates the external proximity of the
mantissa produced by \(H_5\) from the return coordinate that enters the
counter-angle. The separation does not weaken the construction: it eliminates an
identification incompatible with the exact digits and turns each
comparison into a falsifiable assertion of a well-defined type.

% HMT_SOURCE_SEGMENT_END id=A13
% HMT_SOURCE_SEGMENT_BEGIN id=D01 source=manuscrito/sucesor_102/deltas_ley9/c38_par_angular_cartas_y_canales.tex body_lines=3-20
\label{sec:l9s102:par-angular}

The cubic completion
\begin{equation}
 10^3=(9+1)^3=9^3+3\cdot9^2+3\cdot9+1
 =729+243+27+1
 \label{eq:l9s102:base-1000}
\end{equation}
determines the chart global of the coordinate direct. The closure bilayer contributes
the factor \(2\) of the coordinate conjugate:
\begin{equation}
 \boxed{
 A_{\deg}=10^3\alpha,
 \qquad
 C^*_{\deg}=2(\eta_{\mathrm{ret}}+\alpha)
 \quad\text{in }\mathbb R/(360\mathbb Z).}
 \label{eq:l9s102:a-cstar-deg}
\end{equation}
% HMT_SOURCE_SEGMENT_END id=D01
% HMT_SOURCE_SEGMENT_BEGIN id=D04 source=manuscrito/sucesor_102/deltas_ley9/c38_par_angular_cartas_y_canales.tex body_lines=60-66
The inverse identity
\begin{equation}
 \eta_{\mathrm{ret}}=\frac12C^*_{\deg}-\alpha
 \label{eq:l9s102:eta-corolario-inverso}
\end{equation}
is established as an algebraic corollary of the direct construction.

% HMT_SOURCE_SEGMENT_END id=D04
\section{Exact commutativity of the charts in turns, degrees, and radians}
\label{sec:s102:c38:5}

\subsection{Derivation of the cylindrical length and the analytic degree of the parangular precursor}
\label{sec:s102:c38:5:1}

% HMT_SOURCE_SEGMENT_BEGIN id=B03 source=colaboracion/parte_iii/source/public_final/base_83/c30_body.tex body_lines=157-194

Let \(\mathcal P\) be the free module generated by the regional words and
let \(F^n\mathcal P\) be its length filtration. The associated grading
records the first level at which each word appears. For a parameter
\(z\), the ordinary length character is
\[
\chi_z([w])=z^{|w|}.
\]
Multiplicativity
\[
\chi_z([uv])=\chi_z([u])\chi_z([v])
\]
expresses that concatenation adds lengths. When evaluated at
\(z=\alpha\), a word of length \(n\) receives degree \(\alpha^n\).

\begin{definition}[Degree compatibility]
\label{p3-83-c30-s1:def:compatibilidad-grado}
The regional analytic reader is compatible with the cylindrical filtration
when it sends the homogeneous component of length \(n\) to analytic degree
\(n\):
\[
\operatorname{gr}_n\mathcal P\longrightarrow
\alpha^n\mathbb R.
\]
\end{definition}

The five regions belong to \(\mathcal U_6\); therefore, the return datum
appears in degree six:
\[
\rho_\pi[U_6]\longmapsto169\alpha^6.
\]
This six is the length of the regional word. It is not the length of the
closed walks of \cref{chap:cinco-ciclos-pi}. Those walks traverse
the four anchors
\(U_\pi,U_e,U_\varphi,U_{\rm APP}\) simultaneously and have minimum length eight. A
word, an edge of the block graph, and a closed walk are objects of
distinct types; the grading used here belongs to the first.

% HMT_SOURCE_SEGMENT_END id=B03
\subsection{Publication of the analytic degree on the global phase chart}
\label{sec:s102:c38:5:2}

% HMT_SOURCE_SEGMENT_BEGIN id=A03 source=colaboracion/parte_iii/source/public_final/ascensos_20260826/16f_realizacion_analitica_contraangulo.tex body_lines=172-209

Let \(\mathcal P\) be the free module generated by the regional words and
let \(F^n\mathcal P\) be its length filtration. The associated grading
records the first level at which each word appears. For a parameter
\(z\), the ordinary length character is
\[
\chi_z([w])=z^{|w|}.
\]
Multiplicativity
\[
\chi_z([uv])=\chi_z([u])\chi_z([v])
\]
expresses that concatenation adds lengths. When evaluated at
\(z=\alpha\), a word of length \(n\) receives degree \(\alpha^n\).

\begin{definition}[Degree compatibility]
\label{p3-20260826:def:compatibilidad-grado}
The regional analytic reader is compatible with the cylindrical filtration
when it sends the homogeneous component of length \(n\) to analytic degree
\(n\):
\[
\operatorname{gr}_n\mathcal P\longrightarrow
\alpha^n\mathbb R.
\]
\end{definition}

The five regions belong to \(\mathcal U_6\); therefore, the return datum
appears in degree six:
\[
\rho_\pi[U_6]\longmapsto169\alpha^6.
\]
This six is the length of the regional word. It is not the length of the
closed walks of \cref{chap:cinco-ciclos-pi}. Those walks traverse
the four anchors
\(U_\pi,U_e,U_\varphi,U_{\rm APP}\) simultaneously and have minimum length eight. A
word, an edge of the block graph, and a closed walk are objects of
distinct types; the grading used here belongs to the first.

% HMT_SOURCE_SEGMENT_END id=A03
% HMT_SOURCE_SEGMENT_BEGIN id=D02 source=manuscrito/sucesor_102/deltas_ley9/c38_par_angular_cartas_y_canales.tex body_lines=21-34
The global calendar exhibits the three factorizations
\begin{equation}
 360=12\cdot30=4\cdot90=3\cdot120.
 \label{eq:l9s102:360-factorizaciones}
\end{equation}
The analytical chart is obtained through
\begin{equation}
 A_{\mathrm{rad}}=\frac\pi{180}A_{\deg}
 =\frac{50\pi}{9}\alpha,
 \qquad
 C^*_{\mathrm{rad}}=\frac\pi{180}C^*_{\deg}
 =\frac\pi{90}(\eta_{\mathrm{ret}}+\alpha).
 \label{eq:l9s102:a-cstar-rad}
\end{equation}
% HMT_SOURCE_SEGMENT_END id=D02
% HMT_SOURCE_SEGMENT_BEGIN id=D03 source=manuscrito/sucesor_102/deltas_ley9/c38_par_angular_cartas_y_canales.tex body_lines=35-59

\begin{theorem}[Commutativity of Regional and Action Factorizations]
\label{thm:l9s102:conmutatividad-cstar}
The two routes toward the conjugate phase produce the same output:
\begin{equation}
 \boxed{
 \frac\pi{180}C^*_{\deg}
 =\frac\pi{90}(\eta_{\mathrm{ret}}+\alpha)
 =\frac\pi{90}(H_5+\alpha)
 -\mathscr C_\pi(\alpha)=y_*.}
 \label{eq:l9s102:cuadrado-cstar}
\end{equation}
\end{theorem}

\begin{proof}
Substituting \cref{eq:l9s102:eta-ret} in the second member yields
\[
 \frac\pi{90}
 \left(H_5-\frac{90}{\pi}\mathscr C_\pi+\alpha\right)
 =\frac\pi{90}(H_5+\alpha)-\mathscr C_\pi.
\]
The final expression is the regional definition of \(y_*\), whereas the
first equality follows from \cref{eq:l9s102:a-cstar-deg}.
\end{proof}

% HMT_SOURCE_SEGMENT_END id=D03
\section{spectral decomposition \texorpdfstring{\(q_\pm\)}{q+/-} and calculo functional on the channels \texorpdfstring{\(90/120\)}{90/120} previously generated}
\label{sec:s102:c38:6}

\subsection{Spectral reconstruction of the angular pair and hexad-octad transition}
\label{sec:s102:c38:6:1}

% HMT_SOURCE_SEGMENT_BEGIN id=B11 source=colaboracion/parte_iii/source/public_final/base_83/c30_body.tex body_lines=523-569

The common phase and the oriented phase determine the two channels.
\[
q_-=e^{-x+y_*},
\qquad
q_+=e^{-x-y_*}.
\]
The invariants are recovered exactly.
\[
q_-q_+=D_A,
\qquad
\alpha=-\frac9{100\pi}\log D_A,
\qquad
C_*^{\rm HMT}=\frac{90}{\pi}\log\frac{q_-}{q_+}.
\]
The first equality reads the common contraction; the third, the oriented
separation of the sheets. The exchange \(R\mapsto-R\) permutes
\(q_+\) and \(q_-\), preserves \(D_A\) and reverses the sign of \(y_*\).

For
\[
S_{90}(q)=\frac{q^{90}}{1-q^{270}},
\qquad
S_{120}(q)=\frac{q^{120}}{1-q^{360}},
\]
the hexad--octad transition functional is written entirely in the
phase chart:
\[
\Gamma_{6\to8}
=\frac{180}{\pi}xy_*
+12\bigl(S_{90}(q_-)-S_{90}(q_+)\bigr)
-\bigl(S_{120}(q_-)-S_{120}(q_+)\bigr).
\]
The evaluation produced by the coordinate \(C_*^{\rm HMT}\) of
\cref{p3-83-c30-s1:thm:contraangulo-regional} is
\[
\boxed{
\Gamma_{6\to8}
=0.274007672694459274747255260774\ldots.}
\]
Thus, the internal value recognized in Dimensional Mechanics as the Barbero--Immirzi functional ceases to receive \(C_*^{\rm HMT}\) as an independent literal: both descend from the same pair of phases generated in this chapter. The geometric and physical interpretation of \(\Gamma_{6\to8}\) belongs to the work on Dimensional Mechanics; its mathematical evaluation already belongs to the previously constructed interface.

% HMT_SOURCE_SEGMENT_END id=B11
\subsection{Relative transporter and spectral disjunction of the blocks}
\label{sec:s102:c38:6:2}

% HMT_SOURCE_SEGMENT_BEGIN id=B13 source=colaboracion/parte_iii/source/public_final/base_83/c30_body.tex body_lines=619-675
\label{p3-83-c30-s1:sec:transportador-relativo-split}

In the ordered basis of channels, let
\[
 R=\begin{pmatrix}0&1\\1&0\end{pmatrix},
 \qquad P_\pm=\frac{I\pm R}{2},
 \qquad
 B_0=7I+2R,
 \qquad B_1=\Adeg I+\Cdeg R.
\]
The split decomposition
\[
 xI+yR=\rho(x,y)e^{\eta(x,y)R},
 \quad \rho(x,y)=\sqrt{x^2-y^2},
 \quad \eta(x,y)=\operatorname{artanh}(y/x)
\]
constructs the unique transporter
\[
 T_{\rm rel}
 =\frac{\sqrt{(\Adeg)^2-(\Cdeg)^2}}{\sqrt{45}}
   \exp\!\left[
   \left(\operatorname{artanh}\frac{\Cdeg}{\Adeg}
   -\operatorname{artanh}\frac27\right)R\right],
 \qquad T_{\rm rel}B_0=B_1.
\]
In the spectral basis it acts by
\(9\mapsto\Adeg+\Cdeg\) and
\(5\mapsto\Adeg-\Cdeg\).

\begin{theorem}[Spectral Disjunction of the Local and Angular Realizations]
\label{p3-83-c30-s1:thm:no-go-entrelazador-bloques-split}
The spectra
\[
 \operatorname{Spec}(B_0)=\{9,5\},\qquad
 \operatorname{Spec}(B_1)=\{\Adeg+\Cdeg,\Adeg-\Cdeg\}
\]
are disjoint. Consequently,
\[
 \operatorname{Hom}_{B_0,B_1}
 :=\{\Phi:\Phi B_0=B_1\Phi\}=\{0\}.
\]
\end{theorem}

\begin{proof}
The rational bounds constructed for the angular coordinates give
\(7.29<\Adeg<7.30\) and \(2.12<\Cdeg<2.13\); therefore,
\(9.41<\Adeg+\Cdeg<9.43\) and
\(5.16<\Adeg-\Cdeg<5.18\). In the basis that diagonalizes \(R\), each entry
of an intertwiner is multiplied by the difference between an eigenvalue
of \(B_0\) and one of \(B_1\); disjointness forces all of them to vanish.
\end{proof}

The relative transporter compares two blocks within the split algebra;
the dodecaphasic--Witt and incidence intertwiners act between
equivalent representations in their own spaces of residence. This separation
positively types both operations and preserves their domains.

% HMT_SOURCE_SEGMENT_END id=B13
\subsection{Spectral evaluation \texorpdfstring{\(V_{90,120}\)}{V_90,120} of the incidence channels previously generated by the TPK}
\label{sec:s102:c38:6:3}

% HMT_SOURCE_SEGMENT_BEGIN id=A11 source=colaboracion/parte_iii/source/public_final/ascensos_20260826/16f_realizacion_analitica_contraangulo.tex body_lines=554-572

The common phase and the oriented phase determine the two channels.
\[
q_-=e^{-\Arad+\Crad},
\qquad
q_+=e^{-\Arad-\Crad}.
\]
The invariants are recovered exactly.
\[
q_-q_+=D_A,
\qquad
\alpha=-\frac9{100\pi}\log D_A,
\qquad
\Cdeg=\frac{90}{\pi}\log\frac{q_-}{q_+}.
\]
The first equality reads the common contraction; the third, the oriented
separation of the sheets. The exchange \(R\mapsto-R\) permutes
\(q_+\) and \(q_-\), preserves \(D_A\) and reverses the sign of \(\Crad\).

% HMT_SOURCE_SEGMENT_END id=A11
% HMT_SOURCE_SEGMENT_BEGIN id=D05 source=manuscrito/sucesor_102/deltas_ley9/c38_par_angular_cartas_y_canales.tex body_lines=68-97
\label{sec:l9s102:qpm-posteriores}

Let \(H^2=I\) and \(P_\pm=(I\pm H)/2\). The angular block
\begin{equation}
 B=A_{\mathrm{rad}}I+C^*_{\mathrm{rad}}H
 =(A_{\mathrm{rad}}+C^*_{\mathrm{rad}})P_+
 +(A_{\mathrm{rad}}-C^*_{\mathrm{rad}})P_-
 \label{eq:l9s102:b-angular}
\end{equation}
is published multiplicatively by
\begin{equation}
 e^{-B}=q_+P_++q_-P_-,
 \qquad
 q_\pm=\exp[-(A_{\mathrm{rad}}\pm C^*_{\mathrm{rad}})].
 \label{eq:l9s102:qpm-angular}
\end{equation}
The involution \(C^*\mapsto-C^*\) permutes \(q_+\) and \(q_-\), and preserves
\begin{equation}
 q_+q_-=e^{-2A_{\mathrm{rad}}}=D_A.
 \label{eq:l9s102:det-angular}
\end{equation}
The cardinalities \(90/120\), previously constructed through TPK
incidence, now receive their spectral evaluation:
\begin{equation}
 V_{90,120}
 =(I-T^{90})(I-T^{120})^{-1}
 =\sum_{\sigma\in\{+,-\}}
 \frac{1-q_\sigma^{90}}{1-q_\sigma^{120}}P_\sigma.
 \label{eq:l9s102:v90120-posterior}
\end{equation}
% HMT_SOURCE_SEGMENT_END id=D05
\subsubsection{Split relative transporter}
\label{sec:s102:c38:6:3:1}

% HMT_SOURCE_SEGMENT_BEGIN id=A12 source=colaboracion/parte_iii/source/public_final/ascensos_20260826/16f_realizacion_analitica_contraangulo.tex body_lines=574-748
\label{p3-20260826:sec:transportador-relativo-split}

The comparison With the central block of APP can be carried out without identifying
representations distinct. Fijemos the ordered basis of channels and the
involution
\[
 R=\begin{pmatrix}0&1\\1&0\end{pmatrix},
 \qquad R^2=I,
 \qquad
 P_\pm=\frac{I\pm R}{2}.
\]
In commutative algebra
\[
 \mathscr B_R:=\{xI+yR:x,y\in\mathbb R\}
\]
consider a, within the already declared sexagesimal chart,
\[
 B_0:=7I+2R,
 \qquad
 B_1:=\Adeg I+\Cdeg R.
\]
The orthogonal change of basis
\[
 S=\frac1{\sqrt2}
 \begin{pmatrix}1&1\\1&-1\end{pmatrix}
\]
simultaneously diagonalizes the involution and both blocks. For \(x>|y|\),
we define
\[
 \rho(x,y):=\sqrt{x^2-y^2},
 \qquad
 \eta(x,y):=\operatorname{artanh}\!\left(\frac yx\right).
\]
Then
\[
 xI+yR=\rho(x,y)e^{\eta(x,y)R}.
\]
If
\[
 \rho_0=\sqrt{45},\quad
 \eta_0=\operatorname{artanh}(2/7),\quad
 \rho_1=\rho(\Adeg,\Cdeg),\quad
 \eta_1=\eta(\Adeg,\Cdeg),
\]
the unique left multiplication that carries \(B_0\) to \(B_1\) is
\[
 \boxed{
 T_{\rm rel}
 :=\frac{\rho_1}{\sqrt{45}}
 e^{(\eta_1-\eta_0)R},
 \qquad
 T_{\rm rel}B_0=B_1.}
\]

\begin{proposition}[Relative carrier of the two blocks]
\label{p3-20260826:prop:transportador-relativo-split}
Once \(R\), the channel order, the positive branch, and the sexagesimal
chart have been fixed, \(T_{\rm rel}\) is the unique element of
\(\mathscr B_R\) satisfying \(TB_0=B_1\). In the spectral basis it acts by
\[
 9\longmapsto\Adeg+\Cdeg,
 \qquad
 5\longmapsto\Adeg-\Cdeg.
\]
\end{proposition}

\begin{proof}
The split polar decomposition yields the printed formula and proves existence.
Since \(\det B_0=45\ne0\), every matrix solution of \(TB_0=B_1\) satisfies
\(T=B_1B_0^{-1}\); it is therefore unique. Diagonalization by means of \(S\)
produces the two spectral quotients
\((\Adeg+\Cdeg)/9\) and \((\Adeg-\Cdeg)/5\), equivalent to the exponential
expression of \(T_{\rm rel}\).
\end{proof}

This transporter is not a representational intertwiner. Indeed, an
intertwiner \(\Phi\) between the endomorphisms \(B_0\) and \(B_1\) should
satisfy
\[
 \Phi B_0=B_1\Phi.
\]

\begin{theorem}[Spectral Obstruction to Entanglement]
\label{p3-20260826:thm:no-go-entrelazador-bloques-split}
For the values of \(\Adeg\) and \(\Cdeg\) constructed in this section,
\[
 \boxed{\Phi B_0=B_1\Phi\quad\Longrightarrow\quad\Phi=0.}
\]
\end{theorem}

\begin{proof}
The spectra are
\[
 \operatorname{Spec}(B_0)=\{9,5\},
 \qquad
 \operatorname{Spec}(B_1)
 =\{\Adeg+\Cdeg,\Adeg-\Cdeg\}.
\]
The previously demonstrated rational bounds provide the disjoint intervals.
\[
 7.29<\Adeg<7.30,
 \qquad
 2.12<\Cdeg<2.13,
\]
and, consequently,
\[
 9.41<\Adeg+\Cdeg<9.43,
 \qquad
 5.16<\Adeg-\Cdeg<5.18.
\]
In the basis that diagonalizes \(R\), each entry \(\phi_{ij}\) of \(\Phi\)
is multiplied by the difference between an eigenvalue of \(B_0\) and one of
\(B_1\). All these differences are nonzero; hence all entries of
\(\Phi\) vanish.
\end{proof}

Numerically,
\[
 \rho_1=6.9814819335172378048\ldots,
 \qquad
 \frac{\rho_1}{\sqrt{45}}=1.0407378791354140779\ldots,
\]
\[
 \eta_1-\eta_0
 =0.005796351470571295590\ldots.
\]
The formula preserves the split algebra, the projectors \(P_\pm\), and the
Lorentzian form up to the conformal factor \(\rho_1^2/45\). It nevertheless depends
on \(B_1\) already constructed: it does not constitute a prior derivation of
\(\Adeg\) or \(\Cdeg\), nor should it be confused with the
equivariant intertwiners between realizations of the same module.

For
\[
S_{90}(q)=\frac{q^{90}}{1-q^{270}},
\qquad
S_{120}(q)=\frac{q^{120}}{1-q^{360}},
\]
the hexad--octad transition functional is written entirely in the
phase chart:
\[
\Gamma_{6\to8}
=\frac{180}{\pi}\Arad\Crad
+12\bigl(S_{90}(q_-)-S_{90}(q_+)\bigr)
-\bigl(S_{120}(q_-)-S_{120}(q_+)\bigr).
\]
The evaluation produced by the counter-angle of the
\cref{p3-20260826:thm:contraangulo-regional} is
\[
\boxed{
\Gamma_{6\to8}
=0.274007672694459274747255260774\ldots.}
\]
The mathematically equivalent forms of the internal functional are
\[
\boxed{
\Gamma_{6\to8}
=\Arad\Cdeg+K_{6\to8}
=\frac{180}{\pi}\Arad\Crad+K_{6\to8}
=\frac{\pi\Adeg\Cdeg}{180}+K_{6\to8}.}
\]
These identities evaluate \(\Gamma_{6\to8}\) after generating
\((\Crad,\Cdeg)\); it is not used in the inverse sense to select the
counter-angle. The physical reading in the Ashtekar--Barbero connection is
then represented by means of
\[
 \Gamma_{6\to8}\dashrightarrow\gamma_{\rm BI}.
\]
The arrow introduces no third functional and does not act retroactively on \(C^*\). The
series \(S_{90}\), \(S_{120}\), the coefficients \(12\) and \(-1\), and their
assignment to the \(6\to8\) passage belong to the declared mathematical
construction. Physical identification additionally requires the connection's own
normalization.

% HMT_SOURCE_SEGMENT_END id=A12
\section{Naturalness under Refinement and Separation with Respect to Metrological Recognition}
\label{sec:s102:c38:7}

\subsection{Exact estimation of the remainder of the analytic prolongation}
\label{sec:s102:c38:7:1}

% HMT_SOURCE_SEGMENT_BEGIN id=B09 source=colaboracion/parte_iii/source/public_final/base_83/c30_body.tex body_lines=474-499

The expression initially found,
\[
y_7=\frac\pi{90}(H_5+\alpha)-169\alpha^6-D_A\alpha^7,
\]
is the seventh-degree truncation of the full character. Produce
\[
C_7=2.12373833896864600265403827103919\ldots{}^\circ.
\]
The omitted remainder is not an indeterminate term; it is the exact geometric tail.
\[
\mathscr C_\pi(\alpha)
-169\alpha^6-D_A\alpha^7
=\frac{D_A^2\alpha^8}{1-D_A\alpha}.
\]
Therefore,
\[
\left|C_7-C_*^{\rm HMT}\right|
=\frac{180}{\pi}\frac{D_A^2\alpha^8}{1-D_A\alpha}
=2.7839208689064583\ldots\times10^{-16}\;{}^\circ.
\]
The recurrence thus transforms a seventh-degree approximation into a
complete analytic realization and provides a closed-form bound for each
subsequent truncation.

% HMT_SOURCE_SEGMENT_END id=B09
\subsection{Functional dependence of the regional selector and the analytic continuation}
\label{sec:s102:c38:7:2}

% HMT_SOURCE_SEGMENT_BEGIN id=B10 source=colaboracion/parte_iii/source/public_final/base_83/c30_body.tex body_lines=501-521

The construction changes when any of its active components is removed.
The following values are obtained without recalibrating the other terms:
\[
\begin{array}{@{}lc@{}}
\toprule
\text{variant}&\text{conjugate angular coordinate in degrees}\\
\midrule
\rho_\pi=168&2.1237383389772976863904487\ldots\\
\rho_\pi=169&2.1237383389686457242619514\ldots\\
\rho_\pi=170&2.1237383389599937621334541\ldots\\
\text{without determinantal tail}&2.1237383389686949415301675\ldots\\
D_A\text{ replaced by }1&2.1237383389686313409965826\ldots\\
\bottomrule
\end{array}
\]
These ablations do not in themselves constitute a probabilistic estimate.
They demonstrate that the equivariant selector and the determinantal character are
effective components of the output and cannot be removed while maintaining the
same realization.

% HMT_SOURCE_SEGMENT_END id=B10
\subsection{Operatorial Interpretation of the Conjugate Angular Coordinate}
\label{sec:s102:c38:7:3}

% HMT_SOURCE_SEGMENT_BEGIN id=B14 source=colaboracion/parte_iii/source/public_final/base_83/c30_body.tex body_lines=677-715

The entire procedure can be summarized as a chain of applications:
\[
\begin{aligned}
\mathcal F_\pi
&\longmapsto b_\pi
\longmapsto\rho_\pi=169,\\
\alpha
&\longmapsto A
\longmapsto x
\longmapsto D_A,\\
(\rho_\pi,D_A)
&\longmapsto\mathscr C_\pi(\alpha),\\
(H_5,\mathscr C_\pi(\alpha))
&\longmapsto y_*
\longmapsto C_*^{\rm HMT},\\
(x,y_*)
&\longmapsto(q_+,q_-)
\longmapsto\Gamma_{6\to8}.
\end{aligned}
\]
The first line belongs to the finite catalog; the second, to the common closure of
\(\alpha\); the third, to the determinantal prolongation; the fourth produces the
oriented coordinate; the fifth transports it toward the dimensional interface.

In this organization, \(\alpha\) measures the common contraction of the transport, and \(C_*^{\rm HMT}\) measures the oriented separation of its two sheets. The regional return adds no external constant: it selects a finite amplitude and prolongs it through the determinant already fixed by \(\alpha\). The chart in degrees is subsequent. The primary object is the pair of phases \((x,y_*)\).

An exhaustive verification reconstructs the five regions directly
from the catalogue, traverses the actions of
\(\mathfrak S_5\) and \(\mathfrak S_3\), checks the recurrence, and sums the
series. The input of \(\alpha\) is the exact output of its
dodecaphasic closure, not a transcribed metrological digit. The verification keeps
the catalogue theorems, reader rules, and external comparison
separate; therefore, no experimental comparison intervenes in the generating
function.
% HMT_SOURCE_SEGMENT_END id=B14
\subsection{Angular Truncation Convergence Certification}
\label{sec:s102:c38:7:4}

% HMT_SOURCE_SEGMENT_BEGIN id=A09 source=colaboracion/parte_iii/source/public_final/ascensos_20260826/16f_realizacion_analitica_contraangulo.tex body_lines=505-530

The expression initially found,
\[
y_7=\frac\pi{90}(H_5+\alpha)-169\alpha^6-D_A\alpha^7,
\]
is the seventh-degree truncation of the full character. Produce
\[
C_7=2.12373833896864600265403827103919\ldots{}^\circ.
\]
The omitted remainder is not an indeterminate term; it is the exact geometric tail.
\[
\mathscr C_\pi(\alpha)
-169\alpha^6-D_A\alpha^7
=\frac{D_A^2\alpha^8}{1-D_A\alpha}.
\]
Therefore,
\[
\left|C_7-\Cdeg\right|
=\frac{180}{\pi}\frac{D_A^2\alpha^8}{1-D_A\alpha}
=2.7839208689064583\ldots\times10^{-16}\;{}^\circ.
\]
The recurrence thus transforms a seventh-degree approximation into a
complete analytic realization and provides a closed-form bound for each
subsequent truncation.

% HMT_SOURCE_SEGMENT_END id=A09
\subsection{Naturality of the regional selector under refinement and chart change}
\label{sec:s102:c38:7:5}

% HMT_SOURCE_SEGMENT_BEGIN id=A10 source=colaboracion/parte_iii/source/public_final/ascensos_20260826/16f_realizacion_analitica_contraangulo.tex body_lines=532-552

The construction changes when any of its active components is removed.
The following values are obtained without recalibrating the other terms:
\[
\begin{array}{@{}lc@{}}
\toprule
\text{variant}&\text{counter-angle in degrees}\\
\midrule
\rho_\pi=168&2.1237383389772976863904487\ldots\\
\rho_\pi=169&2.1237383389686457242619514\ldots\\
\rho_\pi=170&2.1237383389599937621334541\ldots\\
\text{without determinantal tail}&2.1237383389686949415301675\ldots\\
D_A\text{ replaced by }1&2.1237383389686313409965826\ldots\\
\bottomrule
\end{array}
\]
These ablations do not in themselves constitute a probabilistic estimate.
They demonstrate that the equivariant selector and the determinantal character are
effective components of the output and cannot be removed while maintaining the
same realization.

% HMT_SOURCE_SEGMENT_END id=A10
\subsection{Metric independence of the parangular construction}
\label{sec:s102:c38:7:6}

% HMT_SOURCE_SEGMENT_BEGIN id=A14 source=colaboracion/parte_iii/source/public_final/ascensos_20260826/16f_realizacion_analitica_contraangulo.tex body_lines=799-831

The entire procedure can be summarized as a chain of applications:
\[
\begin{aligned}
\mathcal F_\pi
&\longmapsto b_\pi
\longmapsto\rho_\pi=169,\\
\alpha
&\longmapsto \Adeg
\longmapsto \Arad
\longmapsto D_A,\\
(\rho_\pi,D_A)
&\longmapsto\mathscr C_\pi(\alpha),\\
(H_5,\mathscr C_\pi(\alpha))
&\longmapsto \Crad
\longmapsto \Cdeg,\\
(\Arad,\Crad)
&\longmapsto(q_+,q_-)
\longmapsto\Gamma_{6\to8}.
\end{aligned}
\]
The first line belongs to the finite catalog; the second, to the common closure of
\(\alpha\); the third, to the determinantal prolongation; the fourth produces the
oriented coordinate; the fifth transports it toward the dimensional interface.

In this organization, \(\alpha\) measures the common contraction of transport, and
the counter-angle measures the oriented separation of its two sheets. In the
current evaluation, the regional return selects a finite amplitude and
prolongs it through the determinant already fixed by \(\alpha\), without inserting an
additional external constant. The degree chart is subsequent. The primary
object is the phase pair \((\Arad,\Crad)\). The recurrence, regional sum,
and ablations depend on the catalogue and on
\(\alpha\), but receive no additional metrological constant.
% HMT_SOURCE_SEGMENT_END id=A14

 
\chapter{Positive constitutive vacuum operator and joint derivation of \texorpdfstring{\(\varepsilon,\mu,Z,c\)}{epsilon, mu, Z, c}}
\begingroup
\renewcommand{\chapter}[2][]{}%
\chapter{Positive constitutive operator of the vacuum: completion \(3\to4\), spectral reconstruction, and joint derivation of \(\varepsilon\), \(\mu\), \(Z\), and \(c\)}
\label{chap:p3-final:39:vacio_constitutivo}
\section{Bilayer spectral decomposition of the angular operator}

Let \(R=R^*=R^{-1}\) be the active involution and

\[
P_\pm=\frac{I\pm R}{2},
\qquad
T=q_+P_+ +q_-P_-,
\qquad 0<q_+<q_-<1.
\]

With
\[
x=A_{\rm rad}=\frac{\pi A}{180},\qquad
y=C^*_{\rm rad}=\frac{\pi C^*}{180},
\]
the two channels are constructed before evaluating the vacuum:
\begin{equation}
q_+=e^{-x-y},\qquad q_-=e^{-x+y}.
\label{eq:vacuum-angular-channels}
\end{equation}
The exchange of sheets reverses \(C^*\), that is, \(y\mapsto-y\), and permutes \(q_+\leftrightarrow q_-\). The pair is thus the spectral reading of the same angular antecedent, not two parameters fitted separately.

\begin{definition}[Constitutive response]
\begin{equation}
\mathcal V_{90,120}=(I-T^{90})(I-T^{120})^{-1}.
\label{eq:vacuum-operator}
\end{equation}
\end{definition}

Since \(\|T\|<1\), the denominator is positive and invertible. Functional calculus gives

\begin{equation}
\mathcal V_{90,120}=r_+P_+ +r_-P_-,
\qquad
r_\pm=\frac{1-q_\pm^{90}}{1-q_\pm^{120}}.
\label{eq:vacuum-eigenvalues}
\end{equation}

The assertion of invertibility deserves to be made explicit. For each leaf,
\(0<q_\pm^{120}<1\); therefore
\[
I-T^{120}=(1-q_+^{120})P_++(1-q_-^{120})P_-
\]
is strictly positive and
\[
(I-T^{120})^{-1}
=\frac{P_+}{1-q_+^{120}}+
 \frac{P_-}{1-q_-^{120}}.
\]
The response \(\mathcal V_{90,120}\) is not obtained by solving four
separate scalar equations: it is the result of the functional calculus of a
single positive contraction. This observation prevents permeability,
permittivity, impedance and propagation from losing their shared antecedent.

\begin{theorem}[Joint spectral determination of the four constitutive relations]
The four constituent characters
\begin{equation}
\widehat\varepsilon=r_+^2,\qquad
\widehat\mu=r_-^2,\qquad
\widehat Z=\frac{r_-}{r_+},\qquad
\widehat c=\frac1{r_+r_-}
\label{eq:vacuum-four}
\end{equation}
are spectral functions of \(\mathcal V_{90,120}\). They satisfy
\begin{equation}
\widehat\mu\widehat\varepsilon=\widehat c^{-2},
\qquad
\frac{\widehat\mu}{\widehat\varepsilon}=\widehat Z^2.
\label{eq:vacuum-relations}
\end{equation}
\end{theorem}

\begin{proof}
The identities are obtained by substituting \cref{eq:vacuum-four}. No dimensional chart nor any exterior value intervenes.
\end{proof}

\begin{corollary}[Reconstruction of the two sheets]
The four characters do not contain four degrees of freedom. Any of
the pairs \((\widehat\mu,\widehat\varepsilon)\),
\((\widehat Z,\widehat c)\), or \((\mathcal E,\mathcal B)\) reconstructs the
two positive eigenvalues:
\begin{align}
r_-&=\sqrt{\widehat\mu}
=\sqrt{\widehat Z/\widehat c},
&r_+&=\sqrt{\widehat\varepsilon}
=\frac1{\sqrt{\widehat Z\widehat c}},
\label{eq:vacuum-reconstruct-r}\\
\log r_-&=\frac{\mathcal E+\mathcal B}{2},
&\log r_+&=\frac{\mathcal E-\mathcal B}{2}.
\label{eq:vacuum-reconstruct-log}
\end{align}
In particular, orientation is lost if only the product is retained,
but is recovered by adding the quotient.
\end{corollary}

\begin{proof}
The first line is the positive root of the identities
\cref{eq:vacuum-four,eq:vacuum-relations}. The second inverts the linear
system
\(\mathcal E=\log r_-+\log r_+\),
\(\mathcal B=\log r_--\log r_+\).
\end{proof}

This corollary expresses a central property of HMT: product and quotient are
complementary functionals of the same state. The product determines the
propagation in common, whereas the quotient preserves the relative orientation
of the inner and outer sheets.

In the V2 certificate,

\begin{align*}
r_+&=0.9999996285482493631786952281\ldots,\\
r_-&=0.9997241371895183641315165853\ldots,\\
\widehat\varepsilon&=0.9999992570966367027604416155\ldots,\\
\widehat\mu&=0.9994483504793269350899815559\ldots,\\
\widehat Z&=0.9997245085389372154555543466\ldots,\\
\widehat c&=1.0002763104861575261063862689\ldots.
\end{align*}

These numerals are the terminal recognition of \cref{eq:vacuum-operator}; they do not define its exponents \(90,120\).

\section{Harmonic factorization of the \(90\) and \(120\) sectors}

The harmonic semigroup is

\[
\langle90,120\rangle=30\langle3,4\rangle.
\]

If \(s=q^{30}\), then

\begin{equation}
\frac{1-q^{90}}{1-q^{120}}
=\frac{1-s^3}{1-s^4}
=\frac{1+s+s^2}{1+s+s^2+s^3}.
\label{eq:three-four-completion}
\end{equation}

The identity also separates the response from its defect. If
\[
F_{3\to4}(s)=\frac{1+s+s^2}{1+s+s^2+s^3},
\qquad 0<s<1,
\]
then
\begin{equation}
d_{3\to4}(s):=1-F_{3\to4}(s)
=\frac{s^3}{1+s+s^2+s^3}>0.
\label{eq:vacuum-defect34}
\end{equation}
The additional term is preserved exactly as a positive
defect. Moreover,
\begin{equation}
F_{3\to4}'(s)
=-\frac{s^2(3+2s+s^2)}{(1+s+s^2+s^3)^2}<0.
\label{eq:vacuum-monotonic34}
\end{equation}
Consequently, the order of the angular channels uniquely determines
the order of the constitutive responses. The assignment of
\(\widehat\mu\) and \(\widehat\varepsilon\) to the respective sheets is obtained
analytically, without resorting to their decimal expansions.

Each sheet of the vacuum constitutes the exact completion from three to four
terms of the same elementary mode \(30\). Consequently, permeability
and permittivity are interpreted as two quadratic functionals of the same
completion evaluated on conjugate sheets, and not as independent
constitutive parameters.

Since the function \((1-q^{90})/(1-q^{120})\) decreases in \(0<q<1\) and \(q_+<q_-\), we obtain

\[
0<r_-<r_+<1,\qquad
0<\widehat\mu<\widehat\varepsilon<1,\qquad
0<\widehat Z<1<\widehat c.
\]

The limits of the functional determine its geometry. When \(q\to0^+\),
\(F(q)\to1\): the four readings approach unity and the memory
defect disappears. When \(q\to1^-\),
\[
F(q)\longrightarrow\frac{90}{120}=\frac34.
\]
The response therefore takes values in the positive interval \((3/4,1)\). The
endpoint \(3/4\) is not a fitted constant: it is the quotient of the two
channel lengths and the continuous realization of the
three--to--four completion.

\section{Sheet involution and preservation of orientation}

The sheet exchange \(C^*\mapsto-C^*\) acts by

\begin{equation}
\widehat\mu\longleftrightarrow\widehat\varepsilon,\qquad
\widehat Z\longmapsto\widehat Z^{-1},
\qquad
\widehat c\longmapsto\widehat c.
\label{eq:vacuum-sheet}
\end{equation}

The common propagation is even; impedance preserves orientation.
This distinction will reappear in Barbero and CKM: not all functionals
associated with the change of sheet have the same parity.

\section{Logarithmic characters and tracial compatibility}

Define

\[
\mathcal E=\log(r_+r_-),\qquad
\mathcal B=\log(r_-/r_+).
\]

With the normalized trace \(\tau=\Tr/2\),

\begin{equation}
2\tau(\log\mathcal V)=\mathcal E,
\qquad
-2\tau(R\log\mathcal V)=\mathcal B.
\label{eq:vacuum-traces}
\end{equation}

The pair \((\mathcal E,\mathcal B)\) constitutes a linear coordinate system for the operator
logarithm:
\begin{equation}
\log\mathcal V
=\frac{\mathcal E}{2}I-\frac{\mathcal B}{2}R.
\label{eq:vacuum-log-decomposition}
\end{equation}
The scalar component is invariant under sheet exchange; the
oriented component changes sign. This decomposition directly
demonstrates the parity of propagation and the oriented character of
impedance.

The nonadic amplification

\[
\mathcal V_m=\mathcal V\otimes I_{9^m}
\]

is compatible with the expectations \(E_m=\mathrm{id}\otimes\tau_9\):

\[
E_m(\mathcal V_{m+1})=\mathcal V_m.
\]

Therefore, \(\mathcal E\) and \(\mathcal B\) are not peculiarities of a \(2\times2\) matrix; they persist in the UHF tower and in its tracial closure \(II_1\).

More precisely, for every polynomial \(p\) and afterwards, by continuity,
for every continuous function on the spectrum,
\begin{equation}
E_m\bigl(p(\mathcal V_{m+1})\bigr)=p(\mathcal V_m).
\label{eq:vacuum-functional-refinement}
\end{equation}
Compatibility does not preserve only two numbers: it preserves the entire
functional algebra generated by the response. In particular, it preserves its
spectral projectors, powers, logarithm, and quadratic forms
that measure the defects \cref{eq:vacuum-defect34}.

\section{Dimensional realization of the constitutive relations}

Dimensional Mechanics determines

\begin{align}
Z_{\rm vac}&=\widehat Z\,u_Su_Q^{-2},\\
c_{\rm vac}&=\widehat c\,u_C,\\
\mu_{\rm vac}&=\widehat\mu\,u_Su_C^{-1}u_Q^{-2},\\
\varepsilon_{\rm vac}&=\widehat\varepsilon\,u_S^{-1}u_C^{-1}u_Q^2.
\label{eq:vacuum-sections}
\end{align}

The constitutive identities hold on the straight lines:

\begin{equation}
\mu_{\rm vac}\varepsilon_{\rm vac}=c_{\rm vac}^{-2},
\qquad
\mu_{\rm vac}/\varepsilon_{\rm vac}=Z_{\rm vac}^2.
\end{equation}

The SI chart comes later. It is not necessary to declare \(c\) primitive and then solve \(\mu\varepsilon=c^{-2}\): the three objects descend jointly from \(\mathcal V\) and from the declared lines.

The separation of levels is expressed by the diagram
\[
\begin{array}{c}
(A,C^*)\longrightarrow T\longrightarrow\mathcal V
\longrightarrow
(\widehat\mu,\widehat\varepsilon,\widehat Z,\widehat c)
\\[2mm]
\Big\downarrow\text{ line charts}\hspace{36mm}
\Big\downarrow\text{ final evaluation}
\\[2mm]
(u_S,u_Q,u_C)\longrightarrow
(\mu_{\rm vac},\varepsilon_{\rm vac},Z_{\rm vac},c_{\rm vac})
\longrightarrow\text{ SI coordinates}.
\end{array}
\]
The first line is dimensionless and operatorial; the second determines types of
quantity; the third selects a metrological chart. Altering the final chart
changes neither the eigenvalues nor the constitutive identities. In this
precise sense, mathematical--physical unification consists in distinguishing
the genealogical invariant from its coordinate system, not in removing the
dimensional structure.

\begin{proposition}[Spectral minimality]
Let \(W\) be a positive two-sheet operator with simple spectrum
\(\{r_+,r_-\}\). If four positive characters satisfy
\(\mu\varepsilon=c^{-2}\) and \(\mu/\varepsilon=Z^2\), if their propagation in common is determined by \(c^{-1}=\det W=r_-r_+\), and if its orientation
fixes \(Z=r_-/r_+\), then necessarily
\[
(\mu,\varepsilon,Z,c)
=\left(r_-^2,r_+^2,\frac{r_-}{r_+},\frac1{r_-r_+}\right).
\]
Therefore, \cref{eq:vacuum-four} is not one parametrization among many:
it is the minimal positive realization compatible with product, quotient and
orientation.
\end{proposition}

\begin{proof}
From \(Z=r_-/r_+\) and \(c^{-1}=\det W=r_-r_+\) the positive roots are obtained uniquely
roots. The two constitutive identities then force afterwards
\(\mu=r_-^2\) and \(\varepsilon=r_+^2\).
\end{proof}

 The first conceptual inversion consists in conceiving the vacuum as a positive relational response.\par

In HMT, the vacuum is the minimal dynamical response of the system when its
two sheets must coexist, propagate, and close their difference while
preserving the orientation that distinguishes them.\par

The two sheets constitute two conjugate readings of the same angular antecedent. One preserves an orientation, and the other preserves the opposite orientation. Hence two channels, \(q_+\) and \(q_-\), appear as the two correlated spectral faces of a single state.\par

The operator\par

\[
\mathcal V_{90,120}
=
(I-T^{90})(I-T^{120})^{-1}
\]

compares two propagation depths: \(90\) and \(120\). Those depths are determined by the same elementary mode:\par

\[
90=3\cdot30,
\qquad
120=4\cdot30.
\]

Therefore, the response of each sheet can be rewritten as\par

\[
\frac{1-q^{90}}{1-q^{120}}
=
\frac{1+s+s^2}{1+s+s^2+s^3},
\qquad
s=q^{30}.
\]

And here is where the interpretation becomes genuinely powerful. The numerator contains three states of the mode \(30\); the denominator contains those same three states and one more. The vacuum is, therefore, an exact completion \(3\to4\).\par

That fourth term integrates and preserves the preceding three. It completes them. What is added is recorded as a positive completion residue:\par

\[
1-\frac{1+s+s^2}{1+s+s^2+s^3}
=
\frac{s^3}{1+s+s^2+s^3}>0.
\]

Vacuum memory is precisely the exact increment that allows three terms to become four. It is the positive and necessary trace of the completion.\par

From the two eigenvalues of that single operator arise the four constitutive relations:\par

\[
\widehat\varepsilon=r_+^2,
\qquad
\widehat\mu=r_-^2,
\qquad
\widehat Z=\frac{r_-}{r_+},
\qquad
\widehat c=\frac1{r_+r_-}.
\]

This completely changes the interpretation of permeability and permittivity.\par

Permittivity is the quadratic response of one sheet. Permeability is the quadratic response of the conjugate sheet. Impedance preserves the oriented relation between them. Propagation speed depends on their common product.\par

Thus, product and quotient perform two different tasks:\par

\begin{itemize}[leftmargin=*,itemsep=0.35em]
\item the product forgets which sheet was which and preserves the common propagation;
\item the quotient preserves what the relative orientation between them was.
\end{itemize}

The product answers the question: “what do the two leaves share?”. The quotient answers the question: “how are they distinguished?”.\par

Por eso\par

\[
\widehat\mu\,\widehat\varepsilon=\widehat c^{-2}
\]

y\par

\[
\frac{\widehat\mu}{\widehat\varepsilon}
=
\widehat Z^2
\]

are the two complementary ways of reading the same spectral pair.\par

The holographic structure appears here in a particularly transparent form: four visible magnitudes actually contain only two internal degrees of freedom. Any of certain suitable pairs allows the two sheets to be reconstructed. The constitutive publication preserves genealogy when product and quotient remain jointly available.\par

This is the deep repair of the ablation of \(\mu\) and \(\varepsilon\). The repair reunites their values with their common operator-theoretic origin.\par

The International System chart may subsequently transform these intrinsic sections into \(\mu_0\), \(\varepsilon_0\), \(Z_0\), and \(c\). Its function is to express, within a metrological convention, a structure whose provenance has already been selected by the operator.\par

The vacuum thereby ceases to be a repository of four constantes. becomes a unique spectral operation With two conjugate memories.\par

And that view has an important philosophical consequence: the vacuum is a compatibility relation between two oriented ways of realizing one and the same state.\par

% BEGIN CR28_CONSTITUTIVE_INVERSE
\section{Functional recovery, composition and the Catalan moment}
\label{cr28:sec:restitucion-constitutiva}

The incidence response \(90/120\) preserves more information than an isolated scalar value. In the same notation as \eqref{eq:three-four-completion}, write
\[
 f(s)=\frac{1+s+s^2}{1+s+s^2+s^3},\qquad
 \mathcal D=s\,\frac{d}{ds},\qquad
 \mathcal V=f(T^{30}).
\]
The marks \(P_\pm\) maintain channel orientation. Angular transport is the antecedent of the response; its functional reconstruction recovers this representation without replacing the memory archive from which it arises.

\begin{theorem}[Inversion of the constitutive response]
\label{cr28:thm:inversion-constitutiva} The function \(f\) is a decreasing bijection from \((0,1)\) onto \((3/4,1)\). Its inverse \(\psi\) assigns to \(r\) the unique root \(s\in(0,1)\) of
\begin{equation}
 r s^3+(r-1)(1+s+s^2)=0.
 \label{cr28:eq:cubic}
\end{equation}
The two labelled responses \(r_\pm\) determine
\begin{equation}
 s_\pm=\psi(r_\pm),\quad q_\pm=s_\pm^{1/30},\quad
 x=-\frac12\log(q_+q_-),\quad
 y=\frac12\log\frac{q_-}{q_+}.
 \label{cr28:eq:angular-recovery}
\end{equation}
For the complete operator \(S=\mathcal V^2\), positive with spectrum in \((9/16,1)\), functional calculus gives
\begin{equation}
 T=\bigl[\psi(S^{1/2})\bigr]^{1/30},\qquad
 \det T=\det\bigl[\psi(S^{1/2})\bigr]^{1/30}.
 \label{cr28:eq:operator-recovery}
\end{equation}
\end{theorem}
\begin{proof}
The negative derivative of \eqref{eq:vacuum-monotonic34} and the limits \(f(0^+)=1,\ f(1^-)=3/4\) prove bijectivity. Multiplying \(f(s)=r\) by its positive denominator gives \eqref{cr28:eq:cubic}. The positive root of order \(30\) recovers \(q_\pm\); the sum and difference of \(\log q_\pm=-x\mp y\) give the two angular coordinates. By positivity, \(S^{1/2}=\mathcal V\). Applying \(\psi\) to each spectral projector recovers \(T^{30}\) and the positive root gives \(T\). Taking its determinant completes the proof.
\end{proof}

This recovery of the complete operator does not extend to an isolated determinant. Indeed,
\[
 T_1=\operatorname{diag}(1/5,1/3),\qquad
 T_2=\operatorname{diag}(1/6,2/5)
\]
have determinant \(1/15\) but distinct spectra in the same contractive chamber. The proved injectivity implies \(f(T_1^{30})^2\ne f(T_2^{30})^2\). Both channels and their marks preserve recovery; reducing them to a product loses information.
% END CR28_CONSTITUTIVE_INVERSE

% BEGIN CR28_TRANSPORT_COMPOSITION
% Recuperación expositiva de INTERRELACIONES_ESTRUCTURALES_VACIO.md, §§6--8.
% Insertar después de la inversión cúbica y antes de las coordenadas logarítmicas.
\subsection{Transported composition of constitutive responses}
\label{cr28:comp-sec}

The angular transport and the incidences $90$ and $120$ originate in the
same APP--TRIT--TPK structure: the sheets retain evaluations and their
carries; the tritic regime retains orientation; the TPK calendar transports
phase and memory to the angular channels. The constitutive response
applies the function $f$ of \eqref{eq:vacuum-monotonic34} to those channels.
Its inversion expresses, in response coordinates, the composition of
transports sharing a common frame. The constructive order is thereby
preserved: transport, composition and spectral evaluation. The resulting
operation acts on this representation of the enriched state; its inverse
recovers the channels, while the complete memory history remains in the
underlying structure.

Let $s=q^{30}$ and $\psi=f^{-1}$. The exponent $30$ is the greatest
common divisor of the incidences $90$ and $120$. Cubic inversion determines
$\psi:(3/4,1)\longrightarrow(0,1)$. Adjoining the boundary response
$f(1)=3/4$ extends the bijection to
\[
 \psi:D\longrightarrow(0,1],\qquad D=[3/4,1).
\]
At this boundary the rational function $f$ is used; the quotient
$(1-q^{90})/(1-q^{120})$ has a removable singularity there.

\begin{proposition}[Response monoid and additive coordinate]
\label{cr28:comp-monoid}
The operation
\begin{equation}
 r\odot t=f\bigl(\psi(r)\psi(t)\bigr),\qquad r,t\in D,
 \label{cr28:comp-law}
\end{equation}
makes $D$ a commutative cancellative monoid with identity $3/4$.
The map
\begin{equation}
 \ell(r)=-\frac1{30}\log\psi(r)
 \label{cr28:comp-character}
\end{equation}
is an isomorphism from this monoid to $(\mathbb R_{\geq0},+)$.
The interval $(3/4,1)$ is closed under $\odot$; the identity belongs
to the boundary extension. The identity is the only invertible element
within $D$.
\end{proposition}
\begin{proof}
The product of two elements of $(0,1]$ belongs to the same interval, and
\[
 \psi(r\odot t)=\psi(r)\psi(t).
\]
For three responses, either association has image
$\psi(r)\psi(t)\psi(u)$; injectivity of $\psi$ proves associativity.
Commutativity follows from scalar multiplication. If
$r\odot t=r\odot u$, positivity of $\psi(r)$ permits cancellation,
giving $\psi(t)=\psi(u)$ and hence $t=u$. The parameter $1$ corresponds
to $3/4$ and provides the identity. The logarithm transforms products into
sums, and its range on $(0,1]$ shows that $\ell$ is a bijection onto
$\mathbb R_{\geq0}$. In particular, the response of a channel $q$
satisfies $\ell(r)=-\log q$. If both parameters are strictly less than
$1$, so is their product. Finally, two numbers in $(0,1]$ have product
$1$ only when both are $1$.
\end{proof}

\begin{proposition}[Admissible cancellation and inverse extension]
\label{cr28:comp-inverse}
For $r,u\in D$, the equation $r\odot t=u$ has a solution $t\in D$
exactly when $u\geq r$. This solution is unique and is given by
\begin{equation}
 t=f\!\left(\frac{\psi(u)}{\psi(r)}\right).
 \label{cr28:comp-cancellation}
\end{equation}
The same function $f$ is a bijection from $(0,\infty)$ onto $(0,1)$.
With this extension of $\psi$, the law~\eqref{cr28:comp-law} transports
the multiplicative group of positive parameters to the interval $(0,1)$.
Group inversion satisfies
\begin{equation}
 r^{\ominus}=\psi(r)r,\qquad
 r\odot r^{\ominus}=\frac34.
 \label{cr28:comp-group-inverse}
\end{equation}
\end{proposition}
\begin{proof}
The equation is equivalent to
$\psi(t)=\psi(u)/\psi(r)$. This quotient belongs to $(0,1]$ exactly
when $\psi(u)\leq\psi(r)$, which, by decreasing monotonicity, is
equivalent to $u\geq r$. The bijection proves existence and uniqueness.
The derivative in \eqref{eq:vacuum-monotonic34} is negative for every $s>0$;
the limits of $f$ at $0$ and infinity are $1$ and $0$, respectively.
This establishes the extended bijection. Multiplying the numerator and
denominator by $s^3$ gives
\[
 f(s^{-1})=\frac{s^3+s^2+s}{s^3+s^2+s+1}=s f(s).
\]
The multiplicative inverse of $s=\psi(r)$ therefore has response
$\psi(r)r$. Its composition with $r$ corresponds to the parameter $1$.
\end{proof}

The extension to $s>1$, equivalently to $q>1$, is the algebraic extension
of the reader. The chapter's physical contractive sector retains its
declared domain. The coordinate $\ell$ extends to an isomorphism with
$(\mathbb R,+)$ and changes sign under group inversion.

\subsubsection{Operator realization in a common sheet frame}

\begin{proposition}[Composition with marked projectors]
\label{cr28:comp-operator}
Let $\mathcal R=\mathcal R^*$, $\mathcal R^2=I$ and
$P_\pm=(I\pm\mathcal R)/2$. For $j=a,b$, let
\[
 T_j=q_{j,+}P_++q_{j,-}P_-,\qquad
 \mathcal V_j=f(T_j^{30}),\qquad 0<q_{j,\pm}<1.
\]
The composition $T_{a\circ b}=T_aT_b$ satisfies
\begin{equation}
 \mathcal V_{a\circ b}
 =f\bigl(\psi(\mathcal V_a)\psi(\mathcal V_b)\bigr).
 \label{cr28:comp-operator-law}
\end{equation}
Its two responses are $r_{a,\pm}\odot r_{b,\pm}$. If
$T_j=\exp[-(x_jI+y_j\mathcal R)]$ with $x_j>|y_j|$, then
\begin{equation}
 T_aT_b=
 \exp[-((x_a+x_b)I+(y_a+y_b)\mathcal R)],
 \label{cr28:comp-angular-addition}
\end{equation}
y $x_a+x_b>|y_a+y_b|$.
\end{proposition}
\begin{proof}
Orthogonality $P_+P_-=0$ gives
\[
 T_aT_b=q_{a,+}q_{b,+}P_++q_{a,-}q_{b,-}P_-.
\]
In particular, $T_a$ and $T_b$ commute. Functional calculus on these
projectors yields
$\psi(\mathcal V_j)=T_j^{30}$ and
$(T_aT_b)^{30}=T_a^{30}T_b^{30}$; applying $f$ proves
\eqref{cr28:comp-operator-law}. More explicitly, for every function $g$
defined on the two eigenvalues,
\[
 P_\pm g(T_j)=g(q_{j,\pm})P_\pm.
\]
Since $q_{j,\pm}=\exp[-(x_j\pm y_j)]$, their products prove
\eqref{cr28:comp-angular-addition}. The triangle inequality concludes
$x_a+x_b>|y_a|+|y_b|\geq|y_a+y_b|$.
\end{proof}

The recovered coordinates can be written directly as
\[
 x=\frac{\ell(r_+)+\ell(r_-)}2,\qquad
 y=\frac{\ell(r_+)-\ell(r_-)}2.
\]
Composition adds these pairs. Simultaneous sheet exchange permutes the
two channels of each transport, acts as $(x,y)\mapsto(x,-y)$ and is an
automorphism of the componentwise law. Inversion of both transports in
the algebraic extension acts as $(x,y)\mapsto(-x,-y)$. The two
involutions commute and retain distinct functions: one exchanges
orientation markings; the other inverts the transport.

To represent exchange by conjugation, use an operator $L$ satisfying
$L\mathcal R L^{-1}=-\mathcal R$. In the representation
\[
 \mathcal R=\begin{pmatrix}0&1\\1&0\end{pmatrix},\qquad
 L=\begin{pmatrix}1&0\\0&-1\end{pmatrix},
\]
direct multiplication gives $LP_\pm L^{-1}=P_\mp$, and hence
$LT_jL^{-1}=q_{j,-}P_++q_{j,+}P_-$. Conjugation by $\mathcal R$
preserves its own projectors. A trace combining the channels likewise
cannot replace the projectors in this functional identity: their labels
contribute to the composition.

The common-frame hypothesis is substantive. For example, the positive
definite operators
\[
 A_0=\begin{pmatrix}1/2&0\\0&1/3\end{pmatrix},\qquad
 B_0=\begin{pmatrix}1/2&1/6\\1/6&1/2\end{pmatrix}
\]
have positive eigenvalues, but
\[
 A_0B_0=\begin{pmatrix}1/4&1/12\\1/18&1/6\end{pmatrix}
\]
is not self-adjoint. The proposition specifies transport composition on
common projectors; its domain does not encompass an arbitrary mixture of
material media. Realizing a concatenation of HMT routes additionally
retains their memory data: algebraic closure of the parameter cone
describes the representation and does not by itself select new primary
routes.

\subsubsection{Channel product and transport composition}

The normalized speed of one state remains determined by
\[
 \widehat c=(r_+r_-)^{-1}.
\]
Here the ordinary product combines the two responses of that same state.
The law $\odot$ evaluates a different operation: composition of two
transports on each sheet. The following exact calculation distinguishes
their arguments:
\begin{equation}
 f(1/2)=\frac{14}{15},\qquad
 \frac{14}{15}\odot\frac{14}{15}=f(1/4)=\frac{84}{85}
 \ne\frac{196}{225}=\left(\frac{14}{15}\right)^2.
 \label{cr28:comp-rational-example}
\end{equation}
The fractions in this example are algebraic test parameters. The
expression for $\widehat c$ and the preceding constitutive identities
remain unchanged. The added content is an explicit law transporting
composition and recovering its additive angular coordinate.

% END CR28_TRANSPORT_COMPOSITION

% BEGIN CR28_CATALAN_RECOVERY
\subsection{Functional reconstruction of the oriented moment}
\label{cr28:sec:catalan-recovery}

The differential relation~\eqref{cr28:eq:catalan-problem} admits an inverse preserving the complete system of moments. Its domain is the constitutive function produced by the cycle sectors, with the orientation of~\eqref{eq:vacuum-angular-channels}. Incidences $90=3\cdot30$ and $120=4\cdot30$ fix that function before its evaluation in channels $s_\pm=q_\pm^{30}$. The condition at the origin corresponds to disappearance of the moment when its contractive argument vanishes. This condition also determines the integration constants.

\begin{theorem}[Integral reconstruction and uniqueness]
\label{cr28:thm:catalan-recovery}
Let $f(s)=(1-s^3)/(1-s^4)$, for $0<s<1$, be the response
of~\eqref{eq:three-four-completion}, and let $\mathcal D=s\,d/ds$.
There is a unique function $F\in C^2((0,1))$ satisfying
\begin{equation}
 \mathcal D^2F(s)=\frac{s(1+s)}{1+s+s^2}f(s),
 \qquad \lim_{s\downarrow0}F(s)=0.
 \label{cr28:eq:catalan-problem}
\end{equation}
It is given by
\begin{align}
 F(s)&=\int_0^s\log\frac{s}{t}\,
          \frac{1+t}{1+t+t^2}f(t)\,dt
       =\int_0^s\frac{\log(s/t)}{1+t^2}\,dt
       =\mathcal C_2(s),
 \label{cr28:eq:catalan-integral}\\
 f(s)&=\frac{1+s+s^2}{s(1+s)}\mathcal D^2F(s).
 \label{cr28:eq:catalan-inverse}
\end{align}
Its continuous extension to $s=1$ recovers $F(1)=G_{\rm Cat}$.
\end{theorem}
\begin{proof}
The factorization $1+t+t^2+t^3=(1+t)(1+t^2)$ gives
\[
 \frac{1+t}{1+t+t^2}f(t)=\frac1{1+t^2}.
\]
The integrand of~\eqref{cr28:eq:catalan-integral} is
integrable at zero, since
\[
 0\leq F(s)\leq\int_0^s\log(s/t)\,dt=s.
\]
Hence $F(s)\to0$. Differentiation on the open interval,
first on a truncated interval and then by dominated convergence,
gives
\[
 \mathcal DF(s)=\int_0^s\frac{dt}{1+t^2},
 \qquad \mathcal D^2F(s)=\frac{s}{1+s^2}
     =\frac{s(1+s)}{1+s+s^2}f(s).
\]
In particular, $F$ is of class $C^2$ and satisfies the equation.
If $F_1,F_2$ are two solutions, their difference $H=F_1-F_2$
satisfies $\mathcal D^2H=0$. Setting $z=\log s$ gives
$d^2H(e^z)/dz^2=0$, so $H(s)=a\log s+b$.
The existence of a zero limit as $s\downarrow0$ first forces
$a=0$ and then $b=0$. The condition on $F$ therefore determines
the solution and also the limit $\mathcal DF(s)\to0$.

For $s<1$, the geometric expansion of $(1+t^2)^{-1}$ and
absolute integrability of its weighted terms allow termwise integration:
\[
 \int_0^s t^{2k}\log(s/t)\,dt
     =\frac{s^{2k+1}}{(2k+1)^2},\qquad
 F(s)=\sum_{k\geq0}\frac{(-1)^ks^{2k+1}}{(2k+1)^2}.
\]
The series represents the integral system~\eqref{cr28:eq:catalan-integral}. Absolute and uniform convergence of this last series on $[0,1]$ permits passage to the endpoint $s=1$. In the integral, this is also permitted by the integrable bound $\log(1/t)$, extending the integrand by zero for $t>s$. Finally, solving the differential equation for $f$ gives \eqref{cr28:eq:catalan-inverse}, with positive denominators on $(0,1)$.
\end{proof}

The limit of the moment and the limit of its second derivative are
taken in their respective convergence classes. In particular,
\[
 \lim_{s\uparrow1}\mathcal D^2\mathcal C_2(s)=\frac12
\]
is the radial Abel limit of the differentiated series. Evaluation
of $\mathcal C_2(1)$ directly uses its absolutely convergent
series. This distinction preserves the same depth operator in
the interior and in its boundary readout.

\subsection{Compatibility of reconstruction with channels and speed}
\label{cr28:sec:channels-velocity}

The theorem restores a function from another previously constructed function. Cubic inversion~\eqref{cr28:eq:cubic} instead acts on evaluations $r_\pm$: it recovers arguments $s_\pm$ within that fixed system. Composition of the two operations preserves quarter-turn orientation, labels $P_\pm$ and transport incidences. Thus $G_{\rm Cat}=\mathcal C_2(1)$ is the endpoint reading of a functional object also participating in the interior responses $r_\pm=f(s_\pm)$.

For $\mathcal V_m=\mathcal V\otimes I_{9^m}$ and
$\tau_m=\operatorname{Tr}/(2\cdot9^m)$, the speed readout
retains exactly the normalization of
\eqref{eq:vacuum-traces}:
\begin{equation}
 \exp[-2\tau_m(\log\mathcal V_m)]
   =\frac1{r_+r_-}
   =(\det\mathcal V)^{-1}=\widehat c.
 \label{cr28:eq:velocity-normalized}
\end{equation}
Indeed, each eigenvalue $r_\pm$ is repeated $9^m$ times, so
$2\tau_m(\log\mathcal V_m)=\log r_++\log r_-$.
The identity uses the determinant of the two-sheet realization
and the normalized trace of its amplification. The previously fixed
dimensional sections then give
\[
 c_{\rm int}=c_{\rm vac}=\widehat c\,u_C,
 \qquad \ell_0=c_{\rm vac}t_0,
\]
with the same state, chart and elementary duration as~\eqref{cr28:eq:same-velocity}. Functional restoration preserves this velocity section when composing it with action and length; spectral multiplicity is absorbed by the declared normalization.

% END CR28_CATALAN_RECOVERY

% BEGIN CR28_CONSTITUTIVE_RADIAL_LINK
\subsection{A common constitutive section in action, length and gravity}
\label{cr28:sec:constitutive-radial} Relations \eqref{eq:vacuum-sections} and the preceding restoration maintain a unique velocity section:
\begin{equation}
 c_{\rm int}=c_{\rm vac}=\widehat c\,u_C,\qquad
 \ell_0=c_{\rm int}t_0,\qquad
 c_{\rm int}=(\mu_{\rm vac}\varepsilon_{\rm vac})^{-1/2}.
 \label{cr28:eq:same-velocity}
\end{equation}
The final equality is an identity of sections in compatible dimensional bases. The radial proof of \cref{g26:sec:rigidez-radial-es} first constructs \(L=(5759/23040)\alpha^{16}L_*\) and then determines
\begin{equation}
 G_{\rm ret}=\frac{c_{\rm int}^3L^2}{\hbar_{\rm ret}}
 =\frac{L^2}
 {\hbar_{\rm ret}(\mu_{\rm vac}\varepsilon_{\rm vac})^{3/2}}.
 \label{cr28:eq:constitutive-gravity}
\end{equation}
The second expression results from substituting the constitutive identity into the first; it introduces no new calibration of either channel. In coordinates, if \(c_{\rm int}=c_*\widehat c\), then \(G_{\rm ret}=c_*^3\widehat c^{\,3}L^2/\hbar_{\rm ret}\). The propagation factor enters only once. Catalan recovery, composed transport and radial response thus preserve the same angular origin, its marks and its normalization.
% END CR28_CONSTITUTIVE_RADIAL_LINK



  \endgroup

\chapter{Charge, von Klitzing resistance, conductance, Josephson constant and magnetic flux}
\begingroup
\renewcommand{\chapter}[2][]{}%
\chapter{Derivation of charge and of the electric quanta: von Klitzing resistance, conductance, Josephson constant, and magnetic flux}
\label{chap:p3-final:40:cuantos_electricos}
\label{ch:metrology}

\section{Derivation of the charge section from the vacuum}

With \(Z_{\rm vac}\), \(h_a=2\pi\hbar_a\) and \(\alpha\) fixed, the charge of sheet \(a\in\{\mathrm{pre},\mathrm{ret}\}\) is

\begin{equation}
e_a=\sqrt{\frac{2\alphah_a}{Z_{\rm vac}}}.
\label{eq:charge}
\end{equation}

The section \(e_a\) is not introduced through an International System
coordinate intended to reconstruct \(\alpha\). The causal order proceeds from
\(\alpha\), the action, and the vacuum toward the charge section.

\section{Von Klitzing and Josephson constants, conductance and magnetic flux}

From \cref{eq:charge} we deduce

\begin{align}
R_K&=\frac{h_a}{e_a^2}=\frac{Z_{\rm vac}}{2\alpha},
\label{eq:rk}\\
G_0&=\frac{2e_a^2}{h_a}=\frac{4\alpha}{Z_{\rm vac}},
\label{eq:g0}\\
\Phi_{0,a}&=\frac{h_a}{2e_a}
=\sqrt{\frac{h_aZ_{\rm vac}}{8\alpha}},
\label{eq:phi0}\\
K_{J,a}&=\frac{2e_a}{h_a}=\Phi_{0,a}^{-1}.
\label{eq:kj}
\end{align}

Consequently,

\begin{equation}
R_KG_0=2,\qquad G_0Z_{\rm vac}=4\alpha,
\qquad K_{J,a}\Phi_{0,a}=1.
\label{eq:quantum-electrical-identities}
\end{equation}

The bidirectional covariance takes the form

\[
R_K^{\rm pre}=R_K^{\rm ret},\qquad
G_0^{\rm pre}=G_0^{\rm ret},
\]

\[
\frac{\Phi_{0,\rm pre}}{\Phi_{0,\rm ret}}=R_{\rm act}^{1/2},
\qquad
\frac{K_{J,\rm pre}}{K_{J,\rm ret}}=R_{\rm act}^{-1/2}.
\]

The resistance and conductance relations are independent of the
orientation of action. The flux--frequency pair preserves that
orientation and transforms its two components reciprocally.

  \endgroup

\chapter{Thermal constants and radiation: Boltzmann, Wien and Stefan--Boltzmann}
\begingroup
\renewcommand{\chapter}[2][]{}%
\chapter{Thermal Constants and Radiation: Boltzmann, Wien, Stefan--Boltzmann, and the Spectral Thermodynamics of the Photon Gas}
\label{chap:p3-final:48:constantes_termicas}
\section{Relation termica intrinseca of the escala \(108\)}

The chronological structure \(108\) fixes

\begin{equation}
k_B\Theta_{\rm clk}\log3
=\frac{h_{\rm int}}{108t_0}
=\frac{2\pi\hbar_{\rm int}}{108t_0}.
\label{eq:boltzmann-clock}
\end{equation}

This equality does not separately determine a numerical coordinate of \(k_B\) and another of \(\Theta_{\rm clk}\) without an additional normalization. It does determine their product as a section of energy. The factor \(\log3\) has three compatible antecedents:

\begin{equation}
\log3=C_1+C_2-2C_0
=\frac{S_{\TRIT}}{k_B}
=\frac{E_P}{k_B\Theta_{\rm clk}}.
\label{eq:log3-triple}
\end{equation}

The three equalities correspond to residual, informational and
thermodynamic realizations of the same scalar. The coincidence of value does not identify
their respective domains.

\section{Spectral formulations of Wien's displacement law}

For the spectral density per wavelength, the maximum satisfies

\[
5(1-\ee^{-x_5})=x_5,
\qquad
x_5=5+W_0(-5\ee^{-5}).
\]

With \cref{eq:boltzmann-clock},

\begin{equation}
\lambda_{\max}(\Theta_{\rm clk})
=\frac{108\log3}{x_5}\,\ell_0.
\label{eq:wien-lambda}
\end{equation}

For the frequency density,

\[
3(1-\ee^{-x_3})=x_3,
\qquad
x_3=3+W_0(-3\ee^{-3}),
\]

and

\begin{equation}
\nu_{\max}(\Theta_{\rm clk})
=\frac{x_3}{108t_0\log3}.
\label{eq:wien-frequency}
\end{equation}

The maxima are not reciprocals, since the spectral densities
transform with the corresponding Jacobian. The quotient \(x_5/x_3\)
belongs to the structure of the parameterization and does not represent an
inconsistency.

\section{Stefan--Boltzmann law and thermodynamics of the photon gas}

Evaluating the radiation law at \(\Theta_{\rm clk}\) produces the intrinsic flux

\begin{equation}
j_\star(\Theta_{\rm clk})
=\frac{2\pi^5}{15\,108^4(\log3)^4}\,
\frac{h_{\rm int}}{\ell_0^2t_0^2},
\label{eq:stefan-hmt}
\end{equation}

in the declared normalization. The Stefan--Boltzmann constant is not inserted as a primitive: it results from integrating the bosonic spectrum once action, propagation, and temperature have been fixed.
The spectral laws employed are those of Planck and Wien \cite{Planck,Wien}; HMT fixes here the section on which they are evaluated, not their constants by retrospective comparison.

In units of the Planck scale,

\begin{align}
n_\gamma\ell_P^3
&=\frac{2\zeta(3)}{\pi^2\log^3 3},
\label{eq:photon-number}\\
\frac{u_\gamma\ell_P^3}{E_P}
&=\frac{\pi^2}{15\log^4 3},
\label{eq:photon-energy}\\
\frac{s_\gamma\ell_P^3}{k_B}
&=\frac{4\pi^2}{45\log^3 3}.
\label{eq:photon-entropy}
\end{align}

\(\zeta(3)\) represents the third-order bosonic moment, whereas
\(\log3\) determines the energy associated with the TRIT decision. The equation preserves the
distinction between the two invariants.

The quantum thermal conductance is expressed as

\begin{equation}
\frac{\kappa_Qt_0}{k_B}=\frac{\pi^2}{324\log3}.
\label{eq:thermal-conductance}
\end{equation}

\section{Metrological Confluence of Thermal and Radiative Invariants}

\begin{longtable}{@{}p{0.18\textwidth}p{0.27\textwidth}p{0.45\textwidth}@{}}
\toprule
Object & Equation internal & Significado estructural\\
\midrule
\endhead
\(R_K\) & \(Z_{\rm vac}/(2\alpha)\) & bidirectionally invariant quantum resistance\\
\(G_0\) & \(4\alpha/Z_{\rm vac}\) & conductancia complementaria\\
\(\Phi_0\) & \(\sqrt{hZ/(8\alpha)}\) & flow preserving orientation of action\\
\(K_J\) & \(\Phi_0^{-1}\) & reciprocal relation between frequency and flux\\
\(E_P\) & \(k_B\Theta\log3\) & energy of the cycle and TRIT informational content\\
Wien's law & \(108\log3/x_5\) & endpoint spectral in the scale chronological\\
Photonic Apéry & \(2\zeta(3)/(\pi^2\log^3 3)\) & number density of the bosonic field\\
\bottomrule
\end{longtable}

\begin{remark}[Negative control]
The block is falsified if a change of spectral variable improperly preserves
the same Wien extremum, if \(\zeta(3)\) replaces
\(\log3\), or if separate values are assigned to \(k_B\) and \(\Theta\) without
specifying a chart. The identities
\cref{eq:quantum-electrical-identities,eq:planck-trit} provide direct
algebraic controls.
\end{remark}

\subsection*{Logical Scope and Domain of Validity}
Electric quanta and Planck--TRIT confluence: \emph{Exact internal theorem}. Wien, Stefan--Boltzmann, and the photonic gas: \emph{Exact mathematical composition} with the declared bosonic continuum. The SI chart remains an \emph{External metrological comparison}.

  \endgroup

\chapter{Electroweak constants, \texorpdfstring{\(W/Z\)}{W/Z} relation and angular mixing}
\begingroup
\renewcommand{\chapter}[2][]{}%
\chapter{Electroweak constants and mixing: angular determinant, W/Z ratio, Weinberg angle, and rational CKM reflection}
\label{chap:p3-final:49:sector_electrodebil}
\section{Angular determinant and Weinberg electroweak relation}

The determinant of \(T\) is

\begin{equation}
\det T=q_+q_-=D_A.
\label{eq:detT}
\end{equation}

The constitutive sector of the vacuum evaluates \(F(T)^2\), where
\(F(z)=(1-z^{90})/(1-z^{120})\), whereas the electroweak angular
sector evaluates \(\det T\). These are two functions of a
common antecedent, not a horizontal factorization between constants.
This distinction is developed in the \cref{ch:ew}.

\begin{remark}[Negative control]
The block is refuted if any of these exact identities is violated:
positivity of \(I-T^{120}\), equality
\cref{eq:three-four-completion}, relations \cref{eq:vacuum-relations}, or
compatibility \(E_m(\mathcal V_{m+1})=\mathcal V_m\). An SI coincidence
cannot compensate for an operator failure.
\end{remark}

\subsection*{Logical Scope and Domain of Validity}
Operator, spectrum, completion \(3\to4\), parity, amplification, and constitutive sections: \emph{Exact internal theorem}. The formulation that unites them in a single operator is \emph{New formalization}; the angular data and lines are an \emph{Antecedent demonstrated result} of the preceding HMT development and are declared here as explicit initial structure.

 \label{ch:ew}

\section{Constitutive and Determinantal Functions of the Angular Operator}

Recall the two-sheet operator of the \cref{ch:vacuum},
\begin{equation}
\label{eq:ew-angular-operator}
T=q_+P_+ +q_-P_-,
\qquad
q_-=\ee^{-(A_{\rm rad}-C^*_{\rm rad})},
\qquad
q_+=\ee^{-(A_{\rm rad}+C^*_{\rm rad})}.
\end{equation}
The electromagnetic vacuum applies the function
\begin{equation}
\label{eq:ew-vacuum-reader}
F(T)^2,
\qquad
F(z)=\frac{1-z^{90}}{1-z^{120}},
\end{equation}
while the electroweak angular sector applies the determinant
\begin{equation}
\label{eq:ew-determinant-reader}
\boxed{
\det T=q_+q_-=\ee^{-2A_{\rm rad}}=:D_A.}
\end{equation}
They are two functional realizations of the common antecedent \((A,C^*,T)\). Generation starts from \(T\); recovery of that representation from the complete constitutive operator is also possible in the contractive domain. By \cref{cr28:thm:inversion-constitutiva},
\[
 T=\bigl[\psi((F(T)^2)^{1/2})\bigr]^{1/30},
\]
and therefore \(\det T\) is also recovered. The isolated determinant does not reconstruct the two channels: the positive counterexample in the same theorem preserves their product and modifies their responses. Structural unity resides in \(T\), its projectors and provenance memory; functional recovery preserves that order and does not identify the different readings of permeability, permittivity and \(\theta_W\).

Under the sheet involution \(C^*\mapsto-C^*\), the eigenvalues
\(q_+,q_-\) are exchanged. Therefore \(D_A\) and the constitutive speed
are even, while \((\widehat\mu,\widehat\varepsilon)\) is exchanged and
\(\widehat Z\) is inverted. The following electroweak relation uses the
even invariant.

\section{Exact relation between sectors \(W\) and \(Z\)}

The angular coordinate of duration--perduration is
\begin{equation}
\label{eq:ew-angular-character}
\chi_{\rm dur}(n_A,s_C)
=\exp\!\left(n_AA_{\rm rad}
+\frac{s_C}{6}C^*_{\rm rad}\right).
\end{equation}
This coordinate is not introduced as an isolated label. Its origin is
\begin{equation}
\label{eq:ew-signature-provenance}
\begin{aligned}
\APP&\longrightarrow\TRIT\longrightarrow\TPK
\longrightarrow\text{surviving extension}
\longrightarrow\text{route}\\
&\longrightarrow\text{ledger}
\longrightarrow(n_A,s_C,\ldots)
\longrightarrow\chi_{\rm dur}.
\end{aligned}
\end{equation}
The signatures of \(W\), \(Z\) and \(H\) are typed quantities obtained through
that chain; this chapter adopts them without selecting them from
their terminal masses. In the established angular realization, the
charged and neutral sectors satisfy
\begin{equation}
\label{eq:ew-WZ-signatures}
(n_A,s_C)_W=(94,-1),
\qquad
(n_A,s_C)_Z=(95,-1).
\end{equation}
This assertion refers to the common angular component. The
additional refinement coordinates of the General Mass Law are not
eliminated unless the physical transducer proves their cancellation.

\begin{theorem}[Weinberg defect in the physical-mass scheme of the angular chart]
\label{thm:ew-weinberg}
If \(W\) and \(Z\) are evaluated on the same scale section and share
the refinement that accompanies \cref{eq:ew-WZ-signatures}, then
\begin{equation}
\label{eq:ew-WZ-ratio}
\boxed{
\frac{m_W^{\angle}}{m_Z^{\angle}}
=\ee^{-A_{\rm rad}}=\sqrt{D_A},
\qquad
\left(\frac{m_W^{\angle}}{m_Z^{\angle}}\right)^2=D_A.}
\end{equation}
In the physical-mass scheme,
\begin{equation}
\label{eq:ew-sin2thetaW}
\boxed{
\sin^2\theta_W^{\rm OS,HMT}=1-D_A
=0.2248708807528435698111159035\ldots .}
\end{equation}
Moreover,
\begin{equation}
\label{eq:ew-WZ-terminal-ratio}
\frac{m_W^{\angle}}{m_Z^{\angle}}
=0.8804141748331613670128885459\ldots .
\end{equation}
\end{theorem}

\begin{proof}
By \cref{eq:ew-angular-character,eq:ew-WZ-signatures},
\[
\frac{\chi_{\rm dur}(94,-1)}{\chi_{\rm dur}(95,-1)}
=\exp(-A_{\rm rad});
\]
the coordinate \(C^*\) cancels because both sectors share
\(s_C=-1\). Squaring and using
\cref{eq:ew-determinant-reader} yields \(D_A\). The definition
of physical masses \(s_W^2=1-m_W^2/m_Z^2\) yields
\cref{eq:ew-sin2thetaW}.
\end{proof}

The principal content is the identity between two internal constructions: the determinant of the angular operator and the one-step difference in the realization \(W/Z\). Promotion to a physical gauge selector requires an intertwiner that carries the sheets of \(T\) to the neutral basis \((W^3,B)\) and demonstrates which refinements are preserved or cancel.

\section{Gauge couplings in rationalized normalization}

We define the dimensionless gauge coupling
\begin{equation}
\label{eq:ew-egauge}
e_g:=\sqrt{4\pi\alpha}
=0.3028221208717526427662442689\ldots .
\end{equation}
\(e_g\) must not be confused with a dimensional charge section
\(e_{\rm int}\): the former is a coupling in natural units; the
latter lives in the charge line of Dimensional Mechanics.

Let
\begin{equation}
\label{eq:ew-sw-cw}
s_W^2=1-D_A,
\qquad
c_W^2=D_A.
\end{equation}
The convention electrodebil racionalizada to level of arbol gives
\begin{align}
g&=\frac{e_g}{s_W}
=2\sqrt{\frac{\pi\alpha}{1-D_A}}
=0.6385883427334574283647003809\ldots,
\label{eq:ew-g}\\
g'&=\frac{e_g}{c_W}
=2\sqrt{\frac{\pi\alpha}{D_A}}
=0.3439541633108502429362319257\ldots,
\label{eq:ew-gprime}\\
\frac{g'}{g}
&=\sqrt{D_A^{-1}-1}
=0.5386164141966093594996610933\ldots .
\label{eq:ew-gprime-over-g}
\end{align}

\begin{proposition}[Tree-level custodial relation]
\label{prop:ew-rho}
At the same interface,
\begin{equation}
\label{eq:ew-rho}
\rho_{\rm EW}^{(0)}
:=\frac{(m_W^{\angle})^2}
{(m_Z^{\angle})^2c_W^2}=1.
\end{equation}
\end{proposition}

\begin{proof}
By \cref{eq:ew-WZ-ratio,eq:ew-sw-cw}, both the relative numerator and
\(c_W^2\) are \(D_A\).
\end{proof}

The equations \cref{eq:ew-g,eq:ew-gprime} are exact compositions once the
physical mass scheme, rationalization, and scale have been declared. They do not
assert that a change of renormalization scheme leaves their
numerical values unchanged.

\section{Fermi Constant in the Approximation at Tree Level}

In natural units, integration of the charged propagator at low
energy yields
\begin{equation}
\label{eq:ew-GF-def}
\frac{G_F^{(0)}}{\sqrt2}=\frac{g^2}{8m_W^2}.
\end{equation}
Substituting \cref{eq:ew-g} gives the composition
\begin{equation}
\label{eq:ew-GF}
\boxed{
G_F^{(0)}
=\frac{\pi\alpha}
{\sqrt2(1-D_A)m_W^2}.}
\end{equation}
By adopting the previously obtained mass
\begin{equation}
\label{eq:ew-W-output}
m_W^{\HMT}=80.381592550221\ldots\ \mathrm{GeV},
\end{equation}
the terminal evaluation is
\begin{equation}
\label{eq:ew-GF-value}
\boxed{
G_F^{(0)}
=1.1157162817659203186621784\times10^{-5}\ \mathrm{GeV}^{-2}.}
\end{equation}

The difference with respect to a renormalized Fermi constant is not used
to choose any coefficient. Under the convention
\begin{equation}
\label{eq:ew-Delta-r}
G_F
=\frac{\pi\alpha}
{\sqrt2(1-D_A)m_W^2(1-\Delta r)},
\end{equation}
the remaining obligation is to derive \(\Delta r\) from the internal sector of
loops, states, and scale. Opening the observed value of \(G_F\) to define
\(\Delta r\) would be a calibrated reconstruction, not a prediction.

\section{Higgs Coupling Oriented Dependence}

In the angular chart, the relevant routes are
\begin{equation}
\label{eq:ew-HW-signatures}
(n_A,s_C)_H=(97,9),
\qquad
(n_A,s_C)_W=(94,-1).
\end{equation}
Therefore,
\begin{equation}
\label{eq:ew-HW-ratio}
\boxed{
\frac{m_H^{\angle}}{m_W^{\angle}}
=\frac{\chi_{\rm dur}(97,9)}{\chi_{\rm dur}(94,-1)}
=\exp\!\left(3A_{\rm rad}+\frac53C^*_{\rm rad}\right)
=1.55872087212325548564\ldots .}
\end{equation}
Here, \(C^*\) does not cancel: the scalar sector preserves the orientation
that the determinant \(D_A=\ee^{-2A_{\rm rad}}\) does not distinguish.

Let us take the convention of the potential
\begin{equation}
\label{eq:ew-Higgs-potential}
V(H)=-\mu_H^2H^\dagger H
+\lambda_H(H^\dagger H)^2,
\end{equation}
for which
\begin{equation}
\label{eq:ew-Higgs-tree-relations}
m_H^2=2\lambda_Hv^2,
\qquad
m_W^2=\frac{g^2v^2}{4}.
\end{equation}

\begin{corollary}[Oriented Higgs self-coupling]
\label{cor:ew-Higgs-lambda}
At the declared tree-level interface,
\begin{equation}
\label{eq:ew-Higgs-lambda}
\boxed{
\lambda_H
=\frac{\pi\alpha}{2(1-D_A)}
\exp\!\left(6A_{\rm rad}+\frac{10}{3}C^*_{\rm rad}\right)
=0.1238479115482466804662\ldots .}
\end{equation}
Moreover,
\begin{equation}
\label{eq:ew-GF-Higgs}
G_F^{(0)}m_H^2=\sqrt2\,\lambda_H.
\end{equation}
\end{corollary}

\begin{proof}
Dividing the two equations in
\cref{eq:ew-Higgs-tree-relations} gives
\(m_H^2/m_W^2=8\lambda_H/g^2\). With
\(g^2=4\pi\alpha/(1-D_A)\) and the square of
\cref{eq:ew-HW-ratio}, we obtain \cref{eq:ew-Higgs-lambda}. On the other hand,
\(G_F^{(0)}=1/(\sqrt2v^2)\) and
\(m_H^2=2\lambda_Hv^2\) imply
\cref{eq:ew-GF-Higgs}.
\end{proof}

The value of \(\lambda_H\) precedes the incorporation of radiative corrections and depends on the declared convention and scale. Its structural interest is that it separates two memories: the common gauge sector uses the even invariant \(D_A\), whereas the Higgs route explicitly retains \(C^*\).

\section{Determinantal closure of the electroweak sector: mixing, couplings, and constitutive spectrum}

The complete chain of this chapter is
\begin{equation}
\label{eq:ew-causal-diagram}
\boxed{
\begin{aligned}
(A,C^*)&\longrightarrow T
\longrightarrow
\begin{cases}
F(T)^2&\longrightarrow
(\widehat\mu,\widehat\varepsilon,\widehat Z,\widehat c),\\
\det T=D_A&\longrightarrow
(s_W^2,c_W^2),
\end{cases}\\
(D_A,\alpha)&\longrightarrow(e_g,g,g'),\\
(D_A,\alpha,m_W)&\longrightarrow G_F^{(0)},\\
(A,C^*,D_A,\alpha)&\longrightarrow
\left(\frac{m_H}{m_W},\lambda_H\right).
\end{aligned}}
\end{equation}
The figure always appears at the end. Neither \(\sin^2\theta_W\), nor
\(G_F\), nor \(\lambda_H\) is introduced as a target value for selecting
\(A,C^*\) or the route signatures.

\begin{proposition}[Main Result]
The same angular operator determines the vacuum through a rational function
of its sheets and the electroweak defect through its determinant. The
one-step difference between \(W\) and \(Z\) turns that determinant into
\(\sin^2\theta_W^{\rm OS}=1-D_A\); the Higgs route, in contrast, preserves
the counterangle and preserves the orientation eliminated by the determinant.
\end{proposition}

\begin{remark}[Negative control]
The Weinberg relation is refuted in its domain if the angular
signatures \(W/Z\) do not differ only by \(n_A:94\to95\), if a refinement
coordinate does not cancel at the declared interface, or if
\((m_W^{\angle}/m_Z^{\angle})^2\neq D_A\). Gauge promotion is
blocked if no intertwiner with the basis \((W^3,B)\) exists. The figures
of \(g,g',G_F,\lambda_H\) change if the scheme or scale is altered without
recording it. Using observed \(G_F\) or \(m_H\) to select
\(\Delta r\), the signatures, or \(C^*\) introduces feedback from the
target value and invalidates the derivation. Conflating \(e_g\) with the dimensional
charge constitutes a type error.
\end{remark}

\subsection*{Logical Scope and Domain of Validity}
Angular determinant, quotient \(W/Z\) and complementary Weinberg defect:
\emph{Exact internal theorem in the conditional angular chart}. Couplings \(g,g'\),
\(\rho_{\rm EW}^{(0)}\), Fermi and Higgs: \emph{Exact composition at the
tree-level electroweak interface}. The gauge-basis intertwiner and the
internal correction \(\Delta r\): \emph{Open frontier with a refutation criterion}.
Fermi and Higgs are not independent constants selected without target
values: they are exact compositions using the signatures, the quantity
\(m_W\) and the
declared tree-level convention. Subsequent conventional values
are \emph{External metrological comparison}.

 \section{Involution of leaves in \(4\) representations}

The involution \(C^*\mapsto-C^*\) organizes four distinct responses:
\begin{equation}
\label{eq:modular-sheet-parities}
\begin{aligned}
S_{\rm bi}(A,-C^*)&=S_{\rm bi}(A,C^*),\\
\Gamma_{6\to8}(A,-C^*)&=-\Gamma_{6\to8}(A,C^*),\\
(\widehat\mu,\widehat\varepsilon)&\longleftrightarrow
(\widehat\varepsilon,\widehat\mu),\\
\widehat c&\longmapsto\widehat c,
\qquad D_A\longmapsto D_A.
\end{aligned}
\end{equation}
The bilayer information is even; the areal orientation is odd; the two
constitutive quantities form a doublet, and the determinants remain fixed. This parity
table precedes any interpretation as a physical symmetry.

\section{Involutive rational transformation in the CKM parameter space}

Let
\begin{equation}
\label{eq:modular-Mckm}
M_{\rm CKM}=
\begin{pmatrix}
2&-5&28\\
0&7&-25/2\\
1&-20&-2/3
\end{pmatrix},
\qquad
\det M_{\rm CKM}=-\frac{3857}{6}.
\end{equation}
The inverse exists on \(\mathbb Q\) and is
\begin{equation}
\label{eq:modular-Mckm-inverse}
M_{\rm CKM}^{-1}=
\begin{pmatrix}
1528/3857&3380/3857&801/3857\\
75/3857&176/3857&-150/3857\\
6/551&-30/551&-12/551
\end{pmatrix}.
\end{equation}
With
\begin{equation}
\label{eq:modular-ckm-coordinates}
h=\frac{C^*}{6},
\qquad
\Delta=\frac{180}{\pi}\Delta_4,
\qquad
\boldsymbol\theta=M_{\rm CKM}(A,h,\Delta)^T,
\end{equation}
the odd coordinate enters in the direction
\begin{equation}
\label{eq:modular-ch}
c_h=\begin{pmatrix}-5\\7\\-20\end{pmatrix}.
\end{equation}

If \(c_A=(2,0,1)^T\) and
\(c_\Delta=(28,-25/2,-2/3)^T\), the three columns
\((c_A,c_h,c_\Delta)\) are precisely the basis transported by
\(M_{\rm CKM}\). The second row of
\cref{eq:modular-Mckm-inverse} is the covector \(\ell_h\) that extracts the
odd coordinate:
\begin{equation}
\label{eq:modular-dual-coordinate-check}
\ell_hc_A=0,
\qquad
\ell_hc_h=1,
\qquad
\ell_hc_\Delta=0.
\end{equation}

\begin{theorem}[Reflexion CKM inducida by the sheet]
\label{thm:modular-ckm-reflection}
Let
\begin{equation}
\label{eq:modular-Dh-lh}
D_h=\operatorname{diag}(1,-1,1),
\qquad
\ell_h=\frac1{3857}\begin{pmatrix}75&176&-150\end{pmatrix}.
\end{equation}
The involution induced on the space of \(\boldsymbol\theta\) is
\begin{equation}
\label{eq:modular-Rckm}
\boxed{
\mathcal R_{\rm CKM}
=M_{\rm CKM}D_hM_{\rm CKM}^{-1}
=I-2c_h\ell_h.}
\end{equation}
Satisfies
\begin{equation}
\label{eq:modular-Rckm-properties}
\mathcal R_{\rm CKM}^2=I,
\qquad
\det\mathcal R_{\rm CKM}=-1,
\qquad
\mathcal R_{\rm CKM}M_{\rm CKM}=M_{\rm CKM}D_h.
\end{equation}
\end{theorem}

\begin{proof}
Multiplication of
\cref{eq:modular-Mckm,eq:modular-Mckm-inverse} gives the identity; in
particular, \cref{eq:modular-dual-coordinate-check} proves that
\(c_h\ell_h\) is the rank-one projector onto \(\mathbb Qc_h\) along
\(\operatorname{span}\{c_A,c_\Delta\}\). The exact product
\(\ell_hc_h=1\) implies
\[
(I-2c_h\ell_h)^2
=I-4c_h\ell_h+4c_h(\ell_hc_h)\ell_h=I.
\]
The operator fixes the hyperplane
\(\ker\ell_h=\operatorname{span}\{c_A,c_\Delta\}\) pointwise and changes the sign of
\(c_h\). Its eigenvalues are therefore \(1,1,-1\); hence
\(\det\mathcal R_{\rm CKM}=-1\) and
\(\Tr\mathcal R_{\rm CKM}=1\). Finally,
\[
\mathcal R_{\rm CKM}M_{\rm CKM}
=M_{\rm CKM}D_hM_{\rm CKM}^{-1}M_{\rm CKM}
=M_{\rm CKM}D_h,
\]
which proves the entanglement identity.
\end{proof}

The explicit matrix is
\begin{equation}
\label{eq:modular-Rckm-matrix}
\mathcal R_{\rm CKM}=
\begin{pmatrix}
4607/3857&1760/3857&-1500/3857\\
-150/551&199/551&300/551\\
3000/3857&7040/3857&-2143/3857
\end{pmatrix}.
\end{equation}
Although \(\mathcal R_{\rm CKM}\) is not an orthogonal reflection for the
Euclidean product of the three angular coordinates, it is for the
transported rational metric
\begin{equation}
\label{eq:modular-ckm-metric}
G_{\rm CKM}=(M_{\rm CKM}^{-1})^TM_{\rm CKM}^{-1},
\qquad
\mathcal R_{\rm CKM}^TG_{\rm CKM}\mathcal R_{\rm CKM}=G_{\rm CKM}.
\end{equation}
Indeed, in the basis \((c_A,c_h,c_\Delta)\), the metric is the identity
and the involution is \(D_h\). Hence the term ``reflection'' has
precise metric content and does not merely describe that
\(\mathcal R_{\rm CKM}^2=I\).
Phase \(\delta_{\rm CP}=9A\) remains fixed under this reflection. Therefore,
\(\mathcal R_{\rm CKM}\) is not physical CP conjugation, but rather the
transport of HMT sheet inversion to mixing space.

The rational inverse of \(M_{\rm CKM}\) recovers
\begin{align}
A&=\frac{1528\theta_{12}+3380\theta_{23}+801\theta_{13}}{3857},
\label{eq:modular-A-from-ckm}\\
C^*&=\frac{6(75\theta_{12}+176\theta_{23}-150\theta_{13})}{3857}.
\label{eq:modular-C-from-ckm}
\end{align}
The third coordinate is also recovered without loss:
\begin{equation}
\label{eq:modular-Delta-from-ckm}
\Delta
=\frac{6\theta_{12}-30\theta_{23}-12\theta_{13}}{551}.
\end{equation}
These equations are coordinate changes. They do not turn CKM into a
retrospective cause of \(A,C^*\) or of \(\gammaH\).

As a terminal recognition, the previous constructor determines
\begin{equation}
\label{eq:modular-ckm-values}
(\theta_{12},\theta_{23},\theta_{13},\delta_{\rm CP})
=(13.003313589509^\circ,
2.398056157332^\circ,
0.213977382244^\circ,
65.676173123554^\circ).
\end{equation}
The rational reflection is proved before using these decimals.

\begin{proposition}[Main Result]
The transition \(6\to8\) is first formalized through the inclusion
\(\iota_{6,8}:\mathcal F_6\hookrightarrow\mathcal F_8\), from which
\(90,120\) and the defect \(-1/12\) are
obtained. The six-trit word then provides the bijection \(729\leftrightarrow729\), and
\(\mathcal J_{6\to8}\) transports the two collective modes to the boundary. On
that antecedent, area and the modular Hamiltonian exactly reconstruct
the two families; thermofield-double purification determines the entanglement entropy,
and the finite first law
preserves the relative-entropy defect. In CKM, sheet inversion
becomes a rational rank-one odd reflection.
\end{proposition}

\begin{remark}[Negative control]
The initial genealogy fails if \(\iota_{6,8}\) is not injective, if the
counts are not \(90/120\), if \(\delta_{6,8}^{\rm inc}\neq-1/12\), if
\(\beta_{729}\) is not bijective or if it is falsely presented as an
additive isomorphism. Collective transport fails if
\(\mathcal J_{6\to8}D_{\rm inc}\neq
D_{\rm area}\mathcal J_{6\to8}\). Reconstruction fails if the
determinant of the boundary observables is not thirty or if
\cref{eq:modular-inverse-channels} does not recover the occupations.
The TFD fails if its partial trace is not \(\rho_A\). The first law does not
extend to finite changes without the relative entropy of
\cref{eq:modular-finite-first-law}. Type
\(III_{\lambda_\partial}\) is invalidated if a channel does not persist, if there is no
compatible product or if the ratio group changes. The CKM reflection
is refuted by any of
\(\mathcal R^2\neq I\), \(\det\mathcal R\neq-1\) or
\(\mathcal RM\neq MD_h\). A formula
\(\boldsymbol\theta=f(\gammaH)\) without an intertwiner violates the typed
independence of the two functional realizations of the common antecedent.
\end{remark}

\subsection*{Logical Scope and Domain of Validity}
Inclusion \(H\subset O\), counts \(90/120\), two normalizations of the
defect \(-1/12\), bijection of six trits \(729\leftrightarrow729\),
collective intertwiner, scale \(a_0\), finite operators, area spectrum,
TFD, oriented defect, local and finite first law, and CKM reflection:
\emph{Internally exact theorem with its explicit starting structures}.
Classification
\(III_{\lambda_\partial}\): \emph{Theorem with explicit
persistence hypotheses}. Interpretation of the thermofield-double copy as a physical exterior:
\emph{Conditional physical realization}. The magnitude \(\gammaH\) is adopted
as \emph{Proved antecedent result}; it is not rederived in this volume.

 The electroweak link starts from the same angular operator that produced the vacuum constitutive relations, but uses another functional.\par

To obtain \(\mu\), \(\varepsilon\), \(Z\), and \(c\) it was necessary to preserve the full spectrum. It was necessary to distinguish the two sheets.\par

The determinant does the opposite:\par

\[
D_A=\det T=q_+q_-.
\]

By multiplying the eigenvalues, the determinant integrates the two orientations into a common invariant angular magnitude. It publishes the even part and reserves the leaf identity in the oriented memory.\par

By eso \(D_A\) is even under the inversion\par

\[
C^*\longmapsto-C^*.
\]

The sector \(W/Z\) has precisely that structure. The routes of \(W\) and \(Z\) are adjacent in the principal angular coordinate and share the same oriented coordinate. Upon forming the quotient, the contribution of \(C^*\) cancels.\par

Remains\par

\[
\left(\frac{m_W}{m_Z}\right)^2
=
D_A,
\]

and therefore\par

\[
\sin^2\theta_W
=
1-D_A.
\]

The identity reveals that the Weinberg angle is the complementary defect of the angular determinant.\par

The determinant measures the common part that survives when the two sheets are collapsed into their product. \(1-D_A\) measures the complementary part required for mixing.\par

It may be stated as follows: the vacuum preserves the entire response; the electroweak sector uses the even compression of that response.\par

The same angular antecedent produces two distinct functions:\par

\begin{itemize}[leftmargin=*,itemsep=0.35em]
\item the complete functional calculus generates the constitutive relations;
\item the determinant generates the mixing invariant.
\end{itemize}

This explains why the vacuum and the Weinberg angle are two distinct functional realizations of the same angular antecedent.\par

From that point onward, the couplings \(g\) and \(g'\) are, at the stated level, the two projections of the same rationalized electromagnetic coupling onto the sine and cosine directions of the mixing:\par

\[
g^2
=
\frac{4\pi\alpha}{1-D_A},
\qquad
g'^2
=
\frac{4\pi\alpha}{D_A}.
\]

The tree custodial identity\par

\[
\rho_{\rm EW}^{(0)}=1
\]

states that the same angular partition organizes mass and coupling coherently.\par

The Fermi constant appears later as the effective low-energy response of the charged channel:\par

\[
G_F^{(0)}
=
\frac{\pi\alpha}
{\sqrt2(1-D_A)m_W^2}.
\]

The Fermi constant appears as a composition of the fine structure, the angular defect, and the carrier scale \(W\).\par

The Higgs sector introduces a decisive difference. Its route has an oriented coordinate distinct from the route of \(W\), so \(C^*\) persists. The scalar self-coupling publishes the orientation reserved by the determinant.\par

The gauge sector is even: it uses \(D_A\).\par
The scalar sector is oriented: it retains \(C^*\).\par

This division is physically suggestive. Gauge fields organize the common part of the transformation; the scalar sector preserves the memory of the broken orientation.\par

Radiative corrections belong to a later level and must be constructed prospectively from the internal loops, states, and scales, with \(G_F\) as output and terminal evaluation. The causal skeleton is already present: vacuum, mixing, couplings, and Fermi response follow from different readings of the same enriched angular state.\par

 \section{Electroweak Angular Closure and Compatibility Conditions of the Gauge Transducer}
\label{sec:p3-final:cierre-electrodebil}

Let
\[
 c_W=\sqrt{D_A},
 \qquad
 s_W=\sqrt{1-D_A}.
\]
The rotation
\[
 \begin{pmatrix}Z\\A\end{pmatrix}
 =
 \begin{pmatrix}c_W&-s_W\\s_W&c_W\end{pmatrix}
 \begin{pmatrix}W^3\\B\end{pmatrix}
\]
is orthogonal and preserves orientation. The identities
\[
 \left(\frac{m_W}{m_Z}\right)^2=D_A,
 \qquad
 \sin^2\theta_W=1-D_A
\]
are closed in the angular chart. Gauge promotion requires that the
transducer preserve the kinetic form, the charges, and the Ward identities;
the correction \(\Delta r\) belongs to radiative dynamics and is not obtained
by fitting it to \(G_F\).

  \endgroup

\chapter{Intrinsic derivation of gravitation and contragredient reciprocity \texorpdfstring{\(G\leftrightarrow\hbar\)}{G<->hbar}}
\begingroup
\renewcommand{\chapter}[2][]{}%
\chapter{Intrinsic Derivation of \(G\) and \(G\leftrightarrow\hbar\) Reciprocity: Phase Closure, Compton and Planck Lengths, and Areal Invariance}
\label{chap:p3-final:50:gravitacion_accion}
\label{ch:gravity}
% BEGIN CR28_RADIAL_FULL
% Sección incorporable; requiere amsmath y amssymb.
% Fuente principal: RIGIDEZ_SIMULTANEA_DEL_LECTOR_RADIAL_20260926.md.
% Complementos demostrativos: COMPOSICION_METRICA_MEMORIA_Y_ACOPLAMIENTO.md
% y ACCION_ACOPLADA_FASE_Y_MEMORIA_DE_FRONTERA.md.
% La integración, la compilación y la revisión visual corresponden al editor.
\section{Simultaneous rigidity of the longitudinal reader and radial coupling}
\label{g26:sec:rigidez-radial-es}

The Triadic Modular Holography construction transmits to the radial realization an enriched state with its coordinates, orientation, carry, memory and incidences. The sequence
\[
 \begin{gathered}
 \mathrm{APP}\longrightarrow\mathrm{TRIT}\longrightarrow\mathrm{TPK}
 \longrightarrow\text{enriched state}\\
 \longrightarrow\text{joint discrete structure of the continuum}
 \end{gathered}
\]
preserves the Archimedean projection and incidence projection, together with the five inherent constructions of the continuum. The extension
\[
 w_6\longrightarrow w_{12}\longrightarrow w_{18}\longrightarrow
 w_{24}\longrightarrow w_{30}\longrightarrow R_{36}\longrightarrow G_9
\]
is articulated with the operator order \(U_t\longrightarrow\Gamma_9\longrightarrow K_{\rm ph}\): phase returns and the memory record increases. The regional coordinates, selected record \(K\), coordinate \(\alpha\) and return action are received with this common provenance and with their previously constructed readers.

The result of this section is a classification of realizations simultaneously preserving the generated data and the longitudinal, metric and temporal rules specified below. For each positive character the radial response is determined; the coefficient relating that response to mass is common to all admissible characters. The gravitational interpretation is formulated after both readings have been constructed on the same space.

\subsection{Incidence record and constitutive constraints}

Consider complex Hilbert spaces \(V_{\rm in}\) and \(V_{\rm out}\) of dimension \(N\) with fixed inner products. The marked cardinality specified in R1 is positive; in particular, both fibres are nonzero. Adjoints are taken with respect to these products. The symbol \(\mathsf U\) will denote normalized archive transport, and \(U_t\) will retain its meaning as evolution of the enriched state.

\paragraph{R1. Marked record.}
The incidence set is
\begin{equation}
 \Omega=\left\{(j,k,q,u,v):
 \begin{array}{l}
 j,k\in\mathbb Z/12\mathbb Z,\quad k-j\in\{0,3,6,9\},\\
 1\leq q\leq8,\quad u<v,\quad
 u,v\in\{1,2,4,5,7,8\}
 \end{array}\right\}.
 \label{g26:eq:omega-rigidez-es}
\end{equation}
The twelve phases, four relative positions, eight positions of index \(q\) and fifteen pairs in the last set give
\begin{equation}
 N=|\Omega|=12\cdot4\cdot8\binom62
   =12\cdot4\cdot120=5760.
 \label{g26:eq:cardinal-rigidez-es}
\end{equation}
Each incidence has its individual mark. In the corresponding basis \((e_i)_{i\in\Omega}\) of \(V_{\rm out}\), let \(P_i=e_ie_i^*\). The common mode \(\boldsymbol u\) is normalized by \(|u_i|^2=1/N\); write \(P=\boldsymbol u\boldsymbol u^*\). The marked algebra is \(\mathcal A_\Omega=\operatorname{span}\{P_i:i\in\Omega\}\).

\paragraph{R2. Normalization and complete archive.}
The generated inverter \(B:V_{\rm in}\to V_{\rm out}\) satisfies
\begin{equation}
 B^*B=\tfrac14 I_{\rm in},\qquad
 BB^*=\tfrac14 I_{\rm out}.
 \label{g26:eq:normalizacion-b-es}
\end{equation}
The two components are preserved
\begin{equation}
 C_0=(I-P)B,\qquad M_0=PB,\qquad
 C_0+M_0=B,\qquad \mathsf U=2B.
 \label{g26:eq:archivo-completo-es}
\end{equation}
The first records variation relative to the common mode; the second preserves that mode. Their images are orthogonal and their sum recovers the complete transport. By \eqref{g26:eq:normalizacion-b-es}, \(\mathsf U\) is unitary.

\paragraph{R3. Individual inscription.}
The map \(\mathcal E:\operatorname{End}(V_{\rm out})\to\mathcal A_\Omega\) is a linear conditional expectation fixing each \(P_i\) and satisfying
\begin{equation}
 \mathcal E(A_1TA_2)=A_1\mathcal E(T)A_2,
 \qquad A_1,A_2\in\mathcal A_\Omega.
 \label{g26:eq:bimodularidad-es}
\end{equation}
Inscription of the result preserves the complete construction data in a complementary record. The expectation determines the inscribed observable; archive preservation determines its subsequent reconstruction.

\paragraph{R4. Longitudinal composition.}
Fix the generated coordinate \(\alpha>0\) and the longitudinal reference section \(L_*>0\). Longitudinal transport is defined by
\begin{equation}
 D=L_*\alpha^{16}\mathcal E(C_0C_0^*)\mathsf U:
 V_{\rm in}\longrightarrow V_{\rm out}.
 \label{g26:eq:regla-longitudinal-es}
\end{equation}
This rule applies the inscribed coefficient linearly to the longitudinal section. Its type fixes the position of the power \(\alpha^{16}\) and return coefficient within the composition. Reading an intensity root is a different operation with its own transformation rule.

\paragraph{R5. Positive radial metric.}
Let \(X=X^*>0\) be the complete character on \(V_{\rm in}\). Define the positive radial reading \(R_X\) on \(V_{\rm out}\) by
\begin{equation}
 \|R_Xv\|^2=\|XD^*v\|^2
 \qquad(v\in V_{\rm out}).
 \label{g26:eq:metrica-radial-es}
\end{equation}
This equality fixes the metric on every vector of the fibre. In polar notation, \(R_X\) is the arrival modulus \(\lvert(DX)^*\rvert=\sqrt{DX^2D^*}\).

\paragraph{R6. Action and clock.}
The generated action section \(\hbar_{\rm ret}>0\) and the real phase lift satisfy
\begin{equation}
 S(\theta)=\hbar_{\rm ret}\theta.
 \label{g26:eq:accion-fase-es}
\end{equation}
The constitutive velocity \(c>0\) is preserved. Once the intensive length \(L\) of \eqref{g26:eq:regla-longitudinal-es} is obtained, the conjugate clock satisfies \(\omega L=c\). For \(\theta(t)=\omega t\), the action rate determines the energy scale
\begin{equation}
 E_{L}=\frac{dS}{dt}=\hbar_{\rm ret}\omega
   =\frac{\hbar_{\rm ret}c}{L}.
 \label{g26:eq:escala-energetica-es}
\end{equation}
The lift and its refinements remain associated with route memory.

\paragraph{R7. Readings of the same character.}
On \(Y=\mathsf U X\mathsf U^*>0\), define
\begin{equation}
 H_X=E_{L}Y,\qquad M_X=\frac{H_X}{c^2},\qquad
 \Lambda_X=LY^{-1}.
 \label{g26:eq:lecturas-conjugadas-es}
\end{equation}
Energy, mass and conjugate length preserve the same character and transport. The inverse acts in the positive sector; the null sector retains its own treatment.

\paragraph{R8. Reduced radius and dimensional sections.}
The physical identification of this realization is
\begin{equation}
 R_X=\frac{G}{c^2}M_X.
 \label{g26:eq:radio-reducido-es}
\end{equation}
Thus \(R_X\) represents the reduced gravitational radius. The corresponding Schwarzschild radius is \(2R_X\). \(L_*\), the action section and the constitutive velocity section are jointly maintained; a change of units transforms their coordinates according to their dimensions.

\paragraph{Provenance of the constraints.}
Hadamard transport, archive preservation, inscription, phase action and reciprocity have antecedents in the corpus. The expanded record \eqref{g26:eq:omega-rigidez-es} and compositions \eqref{g26:eq:regla-longitudinal-es}--\eqref{g26:eq:metrica-radial-es} are explicitly incorporated in the assembled construction. R4 and R5 are constitutive rules of this realization; the following proof establishes their joint consequence with the remaining antecedents. R8 specifies the physical identification fixing the coefficient's meaning.

\subsection{Simultaneous determination theorem}

\paragraph{Theorem.}
For the same generated data \(B\), \((P_i)\), \(P\), \(\alpha\), \(L_*\), \(\hbar_{\rm ret}\), \(c\) and the same character \(X\), constraints R1--R8 uniquely determine inscription, longitudinal transport and radial response. For all admissible positive characters, the coefficient \(G\) is common. One has
\begin{align}
 \mathcal E&=\Delta_\Omega,
 &\Delta_\Omega(T)&=\sum_{i\in\Omega}P_iTP_i,\label{g26:eq:delta-unica-es}\\
 r_N&=\frac{N-1}{4N}=\frac{5759}{23040},
 &L&=L_*\alpha^{16}r_N,\label{g26:eq:longitud-unica-es}\\
 D&=L\mathsf U,
 &R_X&=L\mathsf U X\mathsf U^*,\label{g26:eq:radial-unica-es}\\
 R_X\Lambda_X&=L^2I,
 &H_X\Lambda_X&=\hbar_{\rm ret}cI.\label{g26:eq:productos-radiales-es}
\end{align}
The coupling is determined by
\begin{equation}
 \boxed{G=\frac{c^3L^2}{\hbar_{\rm ret}}
 =\frac{c^3L_*^2}{\hbar_{\rm ret}}
   \left(\frac{5759}{23040}\right)^2\alpha^{32}.}
 \label{g26:eq:g-rigido-es}
\end{equation}

\paragraph{Proof: uniqueness of inscription.}
Let \(E_{ij}=e_ie_j^*\) be the matrix units. By bimodularity,
\[
 \mathcal E(E_{ij})=P_i\mathcal E(E_{ij})P_j.
\]
For \(i\ne j\), membership of \(\mathcal E(E_{ij})\) in the diagonal algebra implies \(P_i\mathcal E(E_{ij})P_j=0\). For \(i=j\), fixing \(P_i\) gives \(\mathcal E(E_{ii})=P_i\). Linearity then determines \eqref{g26:eq:delta-unica-es} on all matrices. Marked inscription is thus fixed by its algebraic properties.

\paragraph{Proof: return coefficient and archive.}
Normalization of \(B\) and \(P^2=P=P^*\) give
\[
 C_0C_0^*=\tfrac14(I-P),\qquad M_0M_0^*=\tfrac14P.
\]
Since \(P_iPP_i=|u_i|^2P_i=P_i/N\), one obtains
\begin{equation}
 \Delta_\Omega(C_0C_0^*)=\frac{N-1}{4N}I,
 \qquad
 \Delta_\Omega(M_0M_0^*)=\frac{1}{4N}I.
 \label{g26:eq:retorno-y-memoria-es}
\end{equation}
Both readings sum to \(I/4\). The first is scalar, so any normalized distribution \((p_i)\) on incidences, with \(p_i\geq0\) and \(\sum_i p_i=1\), produces \(\sum_i p_i r_N=r_N\). The same coefficient is therefore preserved under arbitrary diagonal weighting.

\paragraph{Proof: length and metric.}
Substituting \eqref{g26:eq:retorno-y-memoria-es} into R4 gives \(D=L\mathsf U\), with \(L\) as in \eqref{g26:eq:longitud-unica-es}. Unitarity of \(\mathsf U\) implies \(DD^*=L^2I\). By R5, the quadratic forms of \(R_X^2\) and \(DX^2D^*\) coincide. The polarization identity determines
\[
 R_X^2=DX^2D^*=L^2\mathsf U X^2\mathsf U^*.
\]
The operator \(L\mathsf U X\mathsf U^*\) is positive and its square is the operator on the right-hand side. Uniqueness of the positive square root proves \eqref{g26:eq:radial-unica-es}. In particular, \(R_I=LI\).

\paragraph{Proof: action, reciprocity and coupling.}
R6 and R7 give
\[
 H_X=\frac{\hbar_{\rm ret}c}{L}Y,
 \qquad M_X=\frac{\hbar_{\rm ret}}{cL}Y,
 \qquad \Lambda_X=LY^{-1}.
\]
Multiplication produces both products of \eqref{g26:eq:productos-radiales-es}. Next, R8 is equivalent to
\[
 LY=\frac{G\hbar_{\rm ret}}{c^3L}Y.
\]
Invertibility of \(Y\) determines \eqref{g26:eq:g-rigido-es}. The same substitution verifies the equality for each positive character. If \(G_1\) and \(G_2\) satisfy the identity with the same data, \((G_1-G_2)Y=0\), and therefore \(G_1=G_2\). This conclusion distinguishes the state-dependent response \(R_X\) from the common coefficient \(G\). \hfill\(\square\)

\subsection{Reciprocity and phase rigidity}

The power--logarithmic antecedent uses, for \(x>0\),
\[
 \ell_p(x)=\frac{x^p-1}{p}\quad(p\ne0),
 \qquad \ell_0(x)=\log x.
\]
Direct substitution proves \(\ell_{-p}(x^{-1})=-\ell_p(x)\). On the positive domain of their inverses, this identity is equivalent to
\[
 \exp_p(s)\exp_{-p}(-s)=1.
\]
Spectral composition on a positive character transmits this product to the pair \(LY\), \(LY^{-1}\). In the present realization, R5 can equivalently be expressed, once R1--R4 and R7 are fixed, by \(R_X\Lambda_X=L^2I\): multiplication by \(\Lambda_X^{-1}\) determines \(R_X=L\mathsf U X\mathsf U^*\), and this expression satisfies R5.

The two products of \eqref{g26:eq:productos-radiales-es} also give a direct form of rigidity:
\begin{equation}
 R_X=\frac{L^2}{\hbar_{\rm ret}c}H_X.
 \label{g26:eq:proporcion-operadores-es}
\end{equation}
For a positive rescaling \(s\), the pair \((sR_X,\Lambda_X/s)\) preserves the radial product; the product \(H_X(\Lambda_X/s)=\hbar_{\rm ret}cI/s\), with energy and action fixed, requires \(s=1\). Joint preservation determines the scale. Likewise, any positive radial operator with the same metric R5 equals \(R_X\) by uniqueness of the positive root.

The lifted action is also exactly determined. For \(\theta\ne0\), the equality \(S/\hbar_{\rm ret}=\theta\) fixes \(S\). Even when rotations of all subdivisions \(\theta/9^n\) are preserved, a fixed real rescaling \(a\) satisfying
\[
 \exp(-ia\theta/9^n)=\exp(-i\theta/9^n)
 \qquad(n\geq0)
\]
satisfies \((a-1)\theta/9^n\in2\pi\mathbb Z\). This sequence tends to zero; for sufficiently large \(n\) its only possible value is zero. From \(\theta\ne0\) one deduces \(a=1\).

\subsection{Planck scales and specialization of the unit character}

Write \(L_\alpha=L\); on the return section, \(E_{P,\rm ret}=E_L\).

The constructed length $L_\alpha>0$, action section $\hbar_a>0$ and velocity $c>0$ determine the scales

\[
t_P=\frac{L_\alpha}{c},\qquad
m_{P,a}=\frac{\hbar_a}{cL_\alpha},\qquad
E_{P,a}=\frac{\hbar_a c}{L_\alpha}.
\]

Indeed, substituting $G_a=c^3L_\alpha^2/\hbar_a$ into the expressions $\sqrt{\hbar_aG_a/c^5}$, $\sqrt{\hbar_ac/G_a}$ and $\sqrt{\hbar_ac^5/G_a}$ gives those three identities respectively. Positivity fixes the root in each case. Length precedes recognition of these scales; their conventional expressions are identities of the already constructed realization.

For a mode of character $y>0$, the mass and longitudinal readings are

\[
m(y)=m_{P,a}y,\qquad r_g(y)=L_\alpha y,\qquad
\bar\lambda(y)=\frac{L_\alpha}{y}.
\]

Thus $r_g(y)\bar\lambda(y)=L_\alpha^2$ for all positive modes. Equality of the two lengths characterizes the mode $y=1$, since $y=y^{-1}$ and $y>0$ imply $y=1$. In this mode, mass is $m_{P,a}$ and both lengths equal $L_\alpha$. General reciprocity is thus distinguished from its specialization to the unit character.

The chronological step and Planck time satisfy $t_0=(\pi_{\rm HMT}/54)t_P$. The complete 108-step phase period is $108t_0=2\pi_{\rm HMT}t_P$ and its angular frequency is $\omega_{108}=1/t_P=c/L_\alpha$. A horizon of radius $R_h=\zeta_hL_\alpha$, with $\zeta_h>0$, preserves its own geometric variable. The Schwarzschild radius associated with a mass is $2r_g$.

If the two action sections are compared with $L_\alpha$ and $c$ fixed, the ratio $\rho=\hbar_b/\hbar_a$ gives $G_b=G_a/\rho$, $m_{P,b}=\rho m_{P,a}$ and $E_{P,b}=\rho E_{P,a}$; $L_\alpha$ and $t_P$ remain equal. This comparison between sections remains distinct from temporal evolution of the constants.

\subsection{Archive covariance and transformation of units}

Let \(T:V_{\rm in}\to V'_{\rm in}\) and \(S:V_{\rm out}\to V'_{\rm out}\) be unitary. Joint transport of the data is
\[
 B'=SBT^*,\quad P'=SPS^*,\quad P_i'=SP_iS^*,\quad X'=TXT^*.
\]
The transported expectation \(\mathcal E'(A')=S\mathcal E(S^*A'S)S^*\) fixes \(P_i'\) and is bimodular in the transformed marked algebra. By substitution,
\begin{align*}
 \mathsf U'&=S\mathsf UT^*,&D'&=SDT^*,\\
 R'_{X'}&=SR_XS^*,&H'_{X'}&=SH_XS^*,&
 \Lambda'_{X'}&=S\Lambda_XS^*.
\end{align*}
Conjugation preserves the R1--R8 equalities and the same coefficient \(G\). Marks, orientations and records are transported with the bases. Incidence permutations are a particular case.

Incorporating an auxiliary memory space \(\mathcal F\) transports \(P\) to \(P\otimes I_{\mathcal F}\), \(\mathsf U\) to \(\mathsf U\otimes I_{\mathcal F}\) and \(\Delta_\Omega\) to \(\Delta_\Omega\otimes\operatorname{id}\). The inscription of \(C_0C_0^*\otimes I_{\mathcal F}\) is \(r_N I\otimes I_{\mathcal F}\); the coefficient preserves the original marked cardinality.

To specify dimensional covariance, let \(\ell\), \(\tau\) and \(m\) be the positive factors of the new length, time and mass units relative to the previous units. Numerical coordinates satisfy
\[
 L'=L/\ell,\qquad c'=c\tau/\ell,\qquad
 \hbar'_{\rm ret}=\hbar_{\rm ret}\tau/(m\ell^2).
\]
Thus,
\[
 \frac{(c')^3(L')^2}{\hbar'_{\rm ret}}
 =G\frac{m\tau^2}{\ell^3},
\]
which is the transformation of a quantity of dimension \(\mathrm{longitud}^3\mathrm{masa}^{-1}\mathrm{tiempo}^{-2}\). The section \(L_*\) transforms as \(L\), while \(\alpha\) and \(r_N\) remain dimensionless.

\subsection{Exact elimination of coupled memory}

Let \(\mathcal K\) and \(\mathcal H\) be finite Hilbert spaces, \(W=W^*>0\) a global weight, and \(\mathcal C:\mathcal K\to\mathcal H\) surjective. The action of residuals \(r\), with boundary defect \(z\), is
\[
 \mathcal S(r)=\mathfrak a\,r^*Wr,
 \qquad \mathcal Cr=z,\qquad \mathfrak a>0.
\]
The global weight preserves its cross blocks between incidences. Define \(\mathcal R=\mathcal CW^{-1}\mathcal C^*\). For \(v\ne0\), injectivity of \(\mathcal C^*\) gives
\[
 v^*\mathcal Rv=\|W^{-1/2}\mathcal C^*v\|^2>0.
\]
Thus \(\mathcal R\) is invertible. The least-action residual is
\[
 r_*(z)=W^{-1}\mathcal C^*\mathcal R^{-1}z.
\]
One verifies \(\mathcal Cr_*=z\). Each admissible residual is written \(r=r_*+k\), with \(\mathcal Ck=0\) and \(k^*Wr_*=k^*\mathcal C^*\mathcal R^{-1}z=0\). Consequently,
\begin{equation}
 \mathcal S(r)=\mathfrak a\,z^*\mathcal R^{-1}z
             +\mathfrak a\,k^*Wk.
 \label{g26:eq:accion-residuo-es}
\end{equation}
Positivity of \(W\) proves uniqueness of the minimum. The pair \((z,k)\) preserves and reconstructs the complete residual.

For a route with transports \(U_j\), the residuals \(r_j=f_j-U_jf_{j-1}\) are reconstructed through \(f_j=U_jf_{j-1}+r_j\). The map \(\mathcal C\) collects the subsequent transports of each residual towards the endpoint. If its blocks are \(\mathcal C_j\), then
\[
 \mathcal R=\sum_{j,k}\mathcal C_j(W^{-1})_{jk}\mathcal C_k^*.
\]
This expression preserves correlations between route segments. For an invertible longitudinal transport \(D\), the boundary \(y=Dz\) has response
\[
 \mathcal R_L=D\mathcal RD^*,\qquad
 \mathcal S_{{\rm ef},L}(y)
 =\mathfrak a\,y^*(D\mathcal RD^*)^{-1}y
 =\mathfrak a\,z^*\mathcal R^{-1}z.
\]
The same action is obtained whether memory is eliminated first or the boundary is transported first. The prefactor preserves the normalization of the source action; when that action carries \(\mathfrak a=\hbar_{\rm ret}\), reduction transmits the same section.

\subsection{Joint reduction of radius, energy and conjugate length}

Write the same positive character in a boundary and memory decomposition:
\[
 Y=\begin{pmatrix}A&F\\F^*&C\end{pmatrix},\qquad C>0,
 \qquad Y_{\rm ef}=A-FC^{-1}F^*.
\]
The identity
\[
 \begin{pmatrix}b\\y\end{pmatrix}^{\!*}Y
 \begin{pmatrix}b\\y\end{pmatrix}
 =b^*Y_{\rm ef}b+\eta^*C\eta,
 \qquad \eta=y+C^{-1}F^*b,
\]
proves \(Y_{\rm ef}>0\), characterizes the least-energy extension and preserves the reconstruction \(y=\eta-C^{-1}F^*b\). For \(s>0\), substitution into the Schur complement gives \(\operatorname{Schur}(sY)=sY_{\rm ef}\). Consequently,
\[
 R_{\rm ef}=LY_{\rm ef},\qquad
 H_{\rm ef}=E_{L}Y_{\rm ef}.
\]
To calculate the compressed conjugate length, solve \(Y(x,y)=(f,0)\). The second equation gives \(y=-C^{-1}F^*x\), and the first gives \(x=Y_{\rm ef}^{-1}f\). The boundary block of \(Y^{-1}\) is therefore \(Y_{\rm ef}^{-1}\). Thus,
\begin{equation}
 \Lambda_{\rm b}=LY_{\rm ef}^{-1},\qquad
 R_{\rm ef}\Lambda_{\rm b}=L^2I,\qquad
 R_{\rm ef}=\frac{L}{E_{L}}H_{\rm ef}.
 \label{g26:eq:schur-reciproco-es}
\end{equation}
Reduction acts on the same blocks in both readings and preserves the residual \(\eta\). The effective character may depend on memory in matrix form, while the coefficient \(L/E_{L}\) remains fixed.

\subsection{Dynamic memory and induced temporal norm}

Dynamic elimination preserves an additional spectral dependence. In a stationary representation, let
\[
 H=\begin{pmatrix}H_{\rm bb}&V\\V^*&H_{\rm ii}\end{pmatrix},
 \qquad H_{\rm ii}>0.
\]
For \(z\) in the relevant resolvent domain, elimination of the interior component in \((H-zI)(b,y)=(f,0)\) gives
\begin{equation}
 \mathcal K_H(z)=H_{\rm bb}-zI
       -V(H_{\rm ii}-zI)^{-1}V^*.
 \label{g26:eq:resolvente-memoria-es}
\end{equation}
When both inverses exist, the boundary block of \((H-zI)^{-1}\) is \(\mathcal K_H(z)^{-1}\). For \(|z|<\lambda_{\min}(H_{\rm ii})\), the convergent resolvent series gives
\begin{align}
 \mathcal K_H(z)&=H_{\rm est}-zZ-z^2VH_{\rm ii}^{-3}V^*-\cdots,
 \label{g26:eq:expansion-resolvente-es}\\
 H_{\rm est}&=H_{\rm bb}-VH_{\rm ii}^{-1}V^*,\qquad
 Z=I+VH_{\rm ii}^{-2}V^*\geq I.
 \label{g26:eq:norma-z-es}
\end{align}
Positivity follows from \(Z-I=(H_{\rm ii}^{-1}V^*)^*(H_{\rm ii}^{-1}V^*)\). The static extension \((b,-H_{\rm ii}^{-1}V^*b)\) has squared norm \(b^*Zb\); energy and temporal norm arise from the same extension. The physical parameter is \(z=\hbar_{\rm ret}\omega\).

To first temporal order, \(T=Z^{1/2}\) and \(\psi=Tb\) transform the pair \((H_{\rm est},Z)\) into \((T^{-*}H_{\rm est}T^{-1},I)\). The section \(\hbar_{\rm ret}\) remains fixed. If \(T\) depends on time, the transformation also preserves the term \(\dot T b\) of \(\dot\psi=\dot T b+T\dot b\).

Radial proportionality is preserved in the complete resolvent. Writing \(\chi_{\rm rad}=L/E_{L}\) and \(R=\chi_{\rm rad} H\), substitution of each block proves
\begin{equation}
 \mathcal K_R(\chi_{\rm rad} z)=\chi_{\rm rad}\mathcal K_H(z).
 \label{g26:eq:resolvente-proporcional-es}
\end{equation}
This identity transports the operator and spectral parameter simultaneously and preserves every order of dynamic memory. For the effective forms of \eqref{g26:eq:schur-reciproco-es}, canonical normalization requires dual transport of the inverse length:
\begin{align*}
 H_{\rm can}&=T^{-*}H_{\rm ef}T^{-1},&
 R_{\rm can}&=T^{-*}R_{\rm ef}T^{-1},&
 \Lambda_{\rm can}&=T\Lambda_{\rm b}T^*.
\end{align*}
Multiplying these expressions gives
\[
 R_{\rm can}=\frac{L}{E_{L}}H_{\rm can},\qquad
 R_{\rm can}\Lambda_{\rm can}=L^2I.
\]
The proof allows noncommuting \(T\), \(R_{\rm ef}\) and \(\Lambda_{\rm b}\). The induced temporal norm and dual length jointly preserve the same \(G\).

\subsection{Refinement, spectral domains and scope}

Let \(J\) and \(K\) be isometric inscriptions between two levels satisfying \(X'J=JX\) and \(\mathsf U'J=K\mathsf U\), with the same intensive section \(L\). Then
\[
 R'K=L\mathsf U'X'\mathsf U'^*K
     =KL\mathsf U X\mathsf U^*=KR,
\]
and the same calculation proves \(H'K=KH\). Intertwining of the inverses is formulated on the corresponding domains. When the compatible system has a positive injective self-adjoint limit \(X\), spectral calculus defines \(Y=\mathsf U X\mathsf U^*\), \(R=LY\), \(H=E_{L}Y\) and \(\Lambda=LY^{-1}\). Proportionality \(R=(L/E_{L})H\) holds on \(\operatorname{Dom}(Y)\), whereas \(R\Lambda=L^2I\) and \(H\Lambda=\hbar_{\rm ret}cI\) hold on \(\operatorname{Dom}(Y^{-1})\), with bounded extensions of their products. Spectral cuts \(\mathbf1_{[\varepsilon,M]}(Y)\) verify the identities before passage to the limit on these domains. A refined edge length \(L/9\) incorporates its own scale factor; the intensive section is preserved according to the declared intertwining.

The classification distinguishes transformations preserving the realization from those modifying a constitutive rule. Diagonal weightings, unitary equivalences, the auxiliary archive and joint reduction preserve \(G\). Changing the marked resolution modifies R1; replacing \(r_N\) with \(\sqrt{r_N}\), \(\alpha^{16}\) with its square, or the inscribed component with the total norm modifies R4. Isolated radius rescaling changes R5, and isolated action rescaling changes R6. R8 maintains the distinction between reduced radius, Schwarzschild radius and change of units.

One thus obtains mathematical uniqueness of the R1--R8 realization: the generated data and character fix the response, and the common composition fixes \(G\). The longitudinal and metric rules appear explicitly in this conclusion as incorporated compositions; gravitational identification retains its physical status. The proof covers the declared class with its constitutive definitions and compatible transformations. The rational checks of the development verify finite cases of the stated identities; the general proof resides in the preceding equations and arguments.
% END CR28_RADIAL_FULL


% BEGIN CR28_RADIAL_CONSTITUTIVE_EVALUATION
\section{Constitutive continuity and evaluation of the gravitational section}
\label{cr28:sec:radial-evaluation}

The longitudinal construction determines \(L\) before using gravitational identification. The composition of rules R1--R8 preserves the velocity section of \eqref{cr28:eq:same-velocity} and the previously generated return action. With
\[
 L=\frac{5759}{23040}\alpha^{16}L_*,
 \qquad \hbar_{\rm ret}=\eta_{\rm ret}S_*,
 \qquad c_{\rm int}=c_*\widehat c,
\]
the result of \eqref{g26:eq:g-rigido-es} is expressed as
\begin{equation}
 G_{\rm ret}
 =\frac{c_*^3L_*^2}{S_*}
   \left(\frac{5759}{23040}\right)^2
   \frac{\alpha^{32}\widehat c^{\,3}}{\eta_{\rm ret}}
 =\frac{L^2}
 {\hbar_{\rm ret}(\mu_{\rm vac}\varepsilon_{\rm vac})^{3/2}}.
 \label{cr28:eq:g-constitutive-evaluation}
\end{equation}
The channels are the same as those of \eqref{eq:vacuum-angular-channels}:
\[
 q_\pm=\exp\!\left[-\frac{\pi}{180}(A_{\rm deg}\pm C^*_{\rm deg})\right],
 \qquad r_\pm=\frac{1-q_\pm^{90}}{1-q_\pm^{120}},
 \qquad \widehat c=(r_+r_-)^{-1}.
\]
The constitutive bases satisfy \(\varepsilon=\varepsilon_*r_+^2\), \(\mu=\mu_*r_-^2\) and \(c_*=(\mu_*\varepsilon_*)^{-1/2}\). Their product proves \(c_{\rm int}=(\mu\varepsilon)^{-1/2}\). If the dimensional coordinate of \(c_{\rm int}\) is declared directly, its basis is recovered as \(c_*=c_{\rm int}/\widehat c\). The constitutive coefficient is not multiplied again: both expressions represent the same section.

In the representation \(L_*=1\,\mathrm m\), \(S_*=10^{-34}\,\mathrm{J\,s}\) and \(c_{\rm int}=299792458\,\mathrm{m\,s^{-1}}\), evaluation of the readers and radial composition gives
\begin{equation}
 \boxed{
 G_{\rm ret}
 =6.674299659414022706367776189\ldots\times10^{-11}
 \,\mathrm{m^3\,kg^{-1}\,s^{-2}}.}
 \label{cr28:eq:g-evaluation}
\end{equation}
The comparison value \(G_{\rm CODATA\,2022}=6.67430(15)\times10^{-11}\,
\mathrm{m^3\,kg^{-1}\,s^{-2}}\) retains its experimental uncertainty. In the same units, the difference between the evaluation and the central value is \(-3.40585977\ldots\times10^{-18}\), approximately \(-0.00227057\) standard uncertainties. The evaluation digits correspond to the declared readers and bases; uncertainty belongs to the measured result. No reference digit selects \(N\), \(r_N\), the longitudinal rule or the radial norm.

\subsection{Chronological realization of the same constructed radius}
\label{cr28:sec:radial-clock} With the elementary duration
\begin{equation}
 t_0=\frac{\pi L}{54c_{\rm int}},
 \qquad \ell_0=c_{\rm int}t_0=\frac{\pi L}{54},
 \label{cr28:eq:clock-from-length}
\end{equation}
the calendar of \(108\) positions satisfies
\[
 T_{108}=108t_0=\frac{2\pi L}{c_{\rm int}},\qquad
 \omega_{108}=\frac{c_{\rm int}}{L},\qquad
 R_{108}=\frac{c_{\rm int}}{\omega_{108}}=L.
\]
By substitution,
\begin{equation}
 \left(\frac{54}{\pi}\right)^2
 \frac{c_{\rm int}^5t_0^2}{\hbar_a}
 =\frac{c_{\rm int}^3L^2}{\hbar_a}.
 \label{cr28:eq:circular-radial-agreement}
\end{equation}
This identity places the circular selector retained below in the same longitudinal normalization. Its uniqueness theorem preserves its stated geometric premise; the preceding construction additionally specifies the length used in that realization.

The positive character \(Y=I\) gives simultaneously \(R_X=\Lambda_X=L\) and \(M_X=\hbar_a/(c_{\rm int}L)\). For a general character \(Y>0\), the lengths are reciprocal: \(R_X=LY\), \(\Lambda_X=LY^{-1}\). Their product remains \(L^2I\) without their individual values necessarily coinciding. Thus identification of the four lengths in the circular case corresponds to a precise sector of the radial proof and does not eliminate dependence on the character.

The proportionality \(R_X=(G/c_{\rm int}^{2})M_X\) is common to all characters in the domain. Its preservation under change of basis, refinement, minimization with memory, Schur complementation and dynamic normalization has been proved in \cref{g26:sec:rigidez-radial-es}. The five-component carrier and its reduction to two helicities, developed subsequently, describe the tensor representation; they do not replace the proof of the common radial coefficient.
% END CR28_RADIAL_CONSTITUTIVE_EVALUATION


\section{Constitutive chain of the gravitational sector: chronological scale \(108\), tensor reduction, holonomy and metrological evaluation of \(G\)}

Gravity is organized here into four strata that must not be merged:

\begin{equation}
\label{eq:gravity-causal-chain}
\begin{aligned}
\text{inscription and radial norm R1--R8}
&\longrightarrow L
\longrightarrow G=c_{\rm int}^3L^2/\hbar_a,\\
\text{icosahedral incidence}
&\longrightarrow V_{12}\xrightarrow{P_5}V_5
\xrightarrow{\Gamma_5}\operatorname{STF}_3
\xrightarrow{\Lambda(n)}\mathcal H_{\rm TT}(n),\\
\text{CT108 calendar with }t_0=\pi L/(54c_{\rm int})
&\longrightarrow R_{108}=L,\quad
\theta_{108}^{\rm circ}=(54/\pi)^2,\\
\text{metrological chart}
&\longrightarrow 6.6742996594140227\ldots\times10^{-11}\,
\mathrm{m^3\,kg^{-1}\,s^{-2}}.
\end{aligned}
\end{equation}

The first line is the determination proved in \cref{g26:sec:rigidez-radial-es}: length precedes coupling. The second constructs the carrier and its transverse traceless reduction. The third preserves the calendar as a realization of that same length, according to \eqref{cr28:eq:circular-radial-agreement}. The fourth evaluates the section in the bases of \cref{cr28:sec:radial-evaluation}. The alternative holonomy branches and the decimal reader are preserved with their own domains; they do not replace the radial proof.

\section{Gravitational magnitude line and chronological scale \(108\)}

Let \(c_{\rm int}>0\), \(t_0>0\),
\(\ell_0=c_{\rm int}t_0\) and a positive section
\(\hbar_a\) of the action line, with
\(a\in\{\mathrm{pre},\mathrm{ret}\}\). The dimensional transducer is

\begin{equation}
\label{eq:gravity-transducer}
G_a(\theta)
=\theta\frac{c_{\rm int}^{5}t_0^2}{\hbar_a}
=\theta\frac{c_{\rm int}^{3}\ell_0^2}{\hbar_a},
\qquad \theta>0.
\end{equation}

Equality of the two expressions uses only \(\ell_0=c_{\rm int}t_0\) and \([G]=L^3M^{-1}T^{-2}\). Dimension is fixed by the line; the dimensionless character \(\theta\) distinguishes the realizations of that system. In the longitudinal composition already proved, \(t_0=\pi L/(54c_{\rm int})\) and \(\theta=(54/\pi)^2\) reproduce exactly \(G_a=c_{\rm int}^3L^2/\hbar_a\).

The CT108 calendar defines

\begin{equation}
\label{eq:gravity-clock108}
T_{108}=108t_0,
\qquad
\omega_{108}=\frac{2\pi}{108t_0},
\qquad
R_{108}=\frac{c_{\rm int}}{\omega_{108}}
=\frac{54}{\pi}\ell_0.
\end{equation}

For each action sheet, let

\begin{equation}
\label{eq:gravity-clock-energy-mass}
E_{108,a}=\hbar_a\omega_{108},
\qquad
m_{108,a}=\frac{E_{108,a}}{c_{\rm int}^2}.
\end{equation}

Then, the reduced Compton wavelength coincides with the radius of the cycle:

\begin{equation}
\label{eq:gravity-compton-clock}
\bar\lambda_{C,a}
=\frac{\hbar_a}{m_{108,a}c_{\rm int}}
=R_{108}.
\end{equation}

\begin{theorem}[Uniqueness of the selector under the HMT periodicity condition]
\label{thm:gravity-circular-selector}
With the reduced convention
\(r_g=Gm/c_{\rm int}^2\) and the HMT periodicity premise that identifies the gravitational
radius of the CT108 quantum with its phase radius, the closure condition
\begin{equation}
\label{eq:gravity-closing-condition}
r_{g,a}(\theta)=R_{108}
\end{equation}
selects a unique positive solution,
\begin{equation}
\label{eq:gravity-theta108}
\boxed{
\theta_{108}^{\rm circ}=\left(\frac{54}{\pi}\right)^2
=295.4525715010569415303525\ldots .}
\end{equation}
In that embodiment,
\begin{equation}
\label{eq:gravity-four-lengths}
R_{108}=\bar\lambda_{C,a}=r_{g,a}=\ell_P^{\rm circ}.
\end{equation}
\end{theorem}

\begin{proof}
By \cref{eq:gravity-transducer,eq:gravity-clock-energy-mass},
\[
r_{g,a}(\theta)
=\frac{G_a(\theta)m_{108,a}}{c_{\rm int}^2}
=\theta c_{\rm int}t_0^2\omega_{108}
=\theta\frac{\pi}{54}\ell_0.
\]
Comparing it with
\(R_{108}=(54/\pi)\ell_0\) necessarily gives
\(\theta=(54/\pi)^2\). The equation is linear in \(\theta\) and all the
factors are positive, so the solution is unique.
\end{proof}

The equality \(r_g=R_{108}\) does not follow from the dimensionality of \cref{eq:gravity-transducer}. It is the explicit geometric condition of this HMT realization: action, Compton, and periodic return must determine the same length. The theorem proves that, once this structural condition is accepted, the character \(\theta\) is fixed without using the SI value of \(G\).

\begin{remark}[Radius convention]
The theorem uses the reduced gravitational radius \(Gm/c^2\). If it is replaced
by the Schwarzschild radius \(2Gm/c^2\), the selector changes by a factor of
two. The two expressions correspond to distinct geometric conditions,
so the convention must be declared before the calculation.
\end{remark}

\section{Reciprocity \(G\leftrightarrow\hbar\) and conservation of the Planck area}

\begin{corollary}[Planck family associated with the temporal periodicity of \(108\) phases]
\label{cor:gravity-planck-family}
For \(G_a=G_a(\theta_{108}^{\rm circ})\),
\begin{align}
\ell_P&=\frac{54}{\pi}\ell_0,
&t_P&=\frac{54}{\pi}t_0,
\label{eq:gravity-planck-length-time}\\
E_{P,a}&=\frac{\hbar_a}{t_P}
=\frac{\pi\hbar_a}{54t_0}
=\frac{h_a}{108t_0},
&\ell_P^2&=\frac{G_a\hbar_a}{c_{\rm int}^3}
=\theta_{108}^{\rm circ}\ell_0^2.
\label{eq:gravity-planck-energy-area}
\end{align}
Moreover,
\begin{equation}
\label{eq:gravity-planck-force-power-density}
F_{P,a}=\frac{c_{\rm int}^4}{G_a},
\qquad
P_{P,a}=\frac{c_{\rm int}^5}{G_a},
\qquad
\rho_{P,a}=\frac{\hbar_a}
{(\theta_{108}^{\rm circ})^2c_{\rm int}^5t_0^4}.
\end{equation}
\end{corollary}

\begin{proof}
The first three identities follow from
\(\ell_P=R_{108}\), \(t_P=\ell_P/c_{\rm int}\), and
\(h_a=2\pi\hbar_a\). For the area,
\[
\frac{G_a\hbar_a}{c_{\rm int}^3}
=\theta_{108}^{\rm circ}c_{\rm int}^2t_0^2
=\theta_{108}^{\rm circ}\ell_0^2.
\]
The remaining expressions are the dimensional compositions of force,
power, and mass per volume with the same section.
\end{proof}

Let
\begin{equation}
\label{eq:gravity-Ract}
R_{\rm act}:=\frac{\hbar_{\rm pre}}{\hbar_{\rm ret}}.
\end{equation}
Since the action appears in the numerator of the chronological mass and in the
denominator of \(G\),
\begin{equation}
\label{eq:gravity-twoway-covariance}
\frac{m_{108,\rm pre}}{m_{108,\rm ret}}=R_{\rm act},
\qquad
\frac{G_{\rm pre}}{G_{\rm ret}}=R_{\rm act}^{-1},
\qquad
G_a\hbar_a
=\theta_{108}^{\rm circ}c_{\rm int}^5t_0^2.
\end{equation}
Therefore \(\ell_P^2=G_a\hbar_a/c^3\) does not change sheet. Gravity and action are oriented reciprocally, while area preserves the common geometry. This cancellation is the convergence interface towards the area operator of \cref{ch:modular}. The second fundamental relation appears when action and gravity are compared.\par

In HMT, action measures how much occurs while simultaneously preserving the orientation of the realization. There is one section before the return and another after it. The phase may return accompanied by a different memory history. That difference is stored in the bidirectional ratio of the two action sections.\par

Gravity responds in a reciprocal manner.\par

The gravitational transducer has the form\par

\[
G_a(\theta)
=
\theta\frac{c_{\rm int}^{5}t_0^2}{\hbar_a}.
\]

Action appears in the denominator. Therefore, if \(\hbar\) is multiplied by a ratio \(R_{\rm act}\) when the sheet changes, the gravitational constant is multiplied by the inverse ratio:\par

\[
\frac{\hbar_{\rm pre}}{\hbar_{\rm ret}}
=
R_{\rm act},
\qquad
\frac{G_{\rm pre}}{G_{\rm ret}}
=
R_{\rm act}^{-1}.
\]

Action and gravity form a contragredient pair here: when one advances in one direction, the other exactly compensates for that displacement.\par

What remains invariant is\par

\[
G_a\hbar_a.
\]

And that is precisely the combination that determines the Planck area:\par

\[
\ell_P^2=\frac{G_a\hbar_a}{c_{\rm int}^3}.
\]

Action remembers the sheet. Gravity responds to that memory. Area remains
invariant.\par

This is an extraordinarily fertile idea. The Planck area emerges as the object that survives when action and gravity transform into one another reciprocally.\par

It may be stated as follows: \(\hbar\) preserves the dynamical orientation; \(G\) preserves the geometric compensation; \(\ell_P^2\) preserves that on which both realizations agree.\par

The area is, therefore, a currency of exchange between quantum physics and geometry, because the bidirectional transformation of \(\hbar\) and \(G\) selects exactly that combination as invariant.\par

This reciprocity affords another understanding of the relation between information and gravity. Information requires a distinction among states; action measures the cost of transition between them. Yet if that transition must acquire a geometry, the gravitational response changes inversely so as to keep the quantum of area stable.\par

Geometry incorporates phase information through an exact compensation.\par

From a Machian reading, this is already a highly suggestive result. A local magnitude of geometry is determined relationally by the orientation and the return of the complete state. The invariant area restricts that relational dependence and fixes its consistency.\par

The HMT version of the relational principle consists in something more precise than the maxim "matter determines geometry": the orientation of action and the gravitational response co-determine one another in such a way that elementary geometry remains invariant.\par

The action contains the history. Gravity adjusts the response. The area preserves the agreement.\par

 The derivation of \(G\) takes that reciprocity to a much more concrete geometric condition.\par

Dimensional Mechanics first determines what type of object \(G\) can be. It fixes its magnitude line and dimensional homogeneity. Chronological closure then selects its value: dimension determines the class of magnitude, and closure determines its dimensionless coefficient.\par

The selection proceeds from the temporal periodicity of (108) phases.\par

The cycle \(108\) determines a frequency, a phase radius, and an energy. From that energy arises a chronological mass. And constructing the reduced Compton wavelength of that mass yields exactly the same radius as that of the cycle:\par

\[
\bar\lambda_C=R_{108}.
\]

Phase and matter are already stating the same length before \(G\) is introduced.\par

Then gravity comes into play. The reduced gravitational radius is constructed.\par

\[
r_g=\frac{Gm}{c^2},
\]

and gravitational closure is required to identify the phase radius, Compton
length, and gravitational radius as a single structural scale:\par

\[
r_g=R_{108}=\bar\lambda_C.
\]

That condition uniquely selects\par

\[
\theta_{108}
=
\left(\frac{54}{\pi}\right)^2.
\]

In this way,\par

\[
R_{108}
=
\bar\lambda_C
=
r_g
=
\ell_P.
\]

The structural function of \(\theta_{108}\) determines the significance of its terminal decimal: it shows what \(G\) is doing.\par

\(G\) appears as the coefficient required to make the four lengths published by these languages coincide:\par

\begin{itemize}[leftmargin=*,itemsep=0.35em]
\item the lenguaje of the fase produce \(R_{108}\);
\item the quantum language produces \(\bar\lambda_C\);
\item the gravitational language produces \(r_g\);
\item the lenguaje of Planck produce \(\ell_P\).
\end{itemize}

The gravitational constant is the mediator that prevents those four readers from separating.\par

Stated in more physical terms: \(G\) expresses how strongly geometry must respond for a quantum whose mass follows from the clock \(108\) to close upon its own phase.\par

It is the relation that makes possible for matter, phase, and geometry to be three compatible publications of the same event.\par

Here the Machian reading becomes stronger. The local gravitational scale is selected by a global cycle-closure condition. The local radius of a quantum acquires meaning when integrated into a complete periodicity. Inertia, phase, and geometry are determined relationally.\par

The demonstrated scope provides a concrete mathematical mechanism for a Machian reading: a local property—the gravitational coupling—is selected by the consistency of the global state. The other historical formulations of Mach's principle preserve their own domain.\par

Then comes the fifth dimension.\par

The scalar \(G\) fixes the gravitational intensity, and the carrier specifies its representational structure. HMT constructs a symmetric traceless tensor space of dimension five. That carrier admits five helicoidal weights:\par

\[
+2,\quad +1,\quad 0,\quad -1,\quad -2.
\]

They are the five natural degrees of freedom of a symmetric traceless tensor in three dimensions before the dynamical constraints of a specific theory are imposed.\par

That clarifies a point that is often confused. “Five components” designates the complete carrier, whereas the two helicities designate the transverse massless reduction in general relativity.\par

It means that the complete carrier has five irreducible coordinates. The dynamics then decides which part of that carrier is physical:\par

\begin{itemize}[leftmargin=*,itemsep=0.35em]
\item in a massive spin-two realization, all five weights may survive;
\item in a massless realization with gauge invariance, the transverse traceless projection eliminates the sectors \(\pm1\) and \(0\);
\item the helicities \(\pm2\) remain.
\end{itemize}

Thus, HMT provides a space prior to the dynamical decision. It constructs the complete five-component carrier and contains, through the corresponding projection, its physical reduction to two helicities.\par

Beauty resides in the ordered coexistence of two statements:\par

\begin{itemize}[leftmargin=*,itemsep=0.35em]
\item the spin-two geometry has a five-component carrier;
\item massless gravitational radiation occupies a two-dimensional subspace.
\end{itemize}

Both statements form a coherent and precise reduction.\par

And that reduction preserves another central lesson: the scalar \(G\) fixes the common intensity, while the orientation information resides in the tensor carrier. The intensity resides in the constant, the helicities reside in the representation, and the complete genealogy brings magnitude and carrier together.\par

  \endgroup
% Secciones 1--3 del delta gravitatorio íntegro.
\section{Genealogical domain of gravitational realization}
\label{sec:l9s102:gravedad-dominio-genealogico}

The gravitational realization receives as its initial object a section of the
discrete continuum already constructed by the composition
\begin{equation}
 \mathrm{APP}\longrightarrow\mathrm{TRIT}\longrightarrow\mathrm{TPK}
 \longrightarrow X_{\mathrm{TPK}}^{\mathrm{enr}}
 \longrightarrow\mathcal C_{\mathrm{cont}}^{\mathrm{disc}}.
 \label{eq:l9s102:gravedad-cadena-genealogica}
\end{equation}
APP preserves the additive and multiplicative evaluations through residue,
quotient and carry; TRIT determines the regime and orientation; the TPK
incorporates the tritic phase selector without eliminating either of the two APP
sheets, performs cardinal transport, preserves memory and realizes phase
return without resetting the state. The dimensional
transducer then acts on the surviving section of the continuum:
\begin{equation}
 \mathfrak G_{108}:\mathcal C_{\mathrm{cont}}^{\mathrm{disc}}
 \times\mathcal L_S^+
 \longrightarrow\mathcal G_{\mathrm{HMT\text{-}MD}},
 \qquad
 (x,\hbar_a)\longmapsto
 \bigl(G_a,R_{108},m_{108,a},\ell_{P,a}^{,2}\bigr).
 \label{eq:l9s102:gravedad-mapa-realizacion}
\end{equation}
The cardinals \(90/120\), the periodicity \(108\), the action, and the angular
pair enter with their genealogy preserved. No subsequent gravitational
identification retrospectively selects these antecedents.

\section{Gravitational Transducer and Chronological Selector of Periodicity}
\label{sec:l9s102:gravedad-transductor}

Let \(c_{\mathrm{int}}>0\), \(t_0>0\),
\(\ell_0=c_{\mathrm{int}}t_0\), and
\(\hbar_a\in\mathcal L_S^+\), with
\(a\in\{\mathrm{pre},\mathrm{ret}\}\). The dimensional transducer is
\begin{equation}
 G_a(\theta)
 =\theta\frac{c_{\mathrm{int}}^5t_0^2}{\hbar_a}
 =\theta\frac{c_{\mathrm{int}}^3\ell_0^2}{\hbar_a},
 \qquad \theta>0.
 \label{eq:l9s102:g-transductor}
\end{equation}
The chronology cyclic of \(108\) states determina
\begin{equation}
 T_{108}=108t_0,
 \qquad
 \omega_{108}=\frac{2\pi}{108t_0},
 \qquad
 R_{108}=\frac{c_{\mathrm{int}}}{\omega_{108}}
 =\frac{54}{\pi}\ell_0.
 \label{eq:l9s102:g-reloj108}
\end{equation}
The chronological quantum and its associated mass are
\begin{equation}
 E_{108,a}=\hbar_a\omega_{108},
 \qquad
 m_{108,a}=\frac{E_{108,a}}{c_{\mathrm{int}}^2}.
 \label{eq:l9s102:g-cuanto-masa}
\end{equation}
The reduced Compton wavelength satisfies
\begin{equation}
 \bar\lambda_{C,a}
 =\frac{\hbar_a}{m_{108,a}c_{\mathrm{int}}}=R_{108}.
 \label{eq:l9s102:g-compton}
\end{equation}

\begin{theorem}[Uniqueness of the gravitational periodicity selector]
\label{thm:l9s102:g-selector}
Under the reduced convention \(r_g=Gm/c_{\mathrm{int}}^2\), the condition
\(r_{g,a}(\theta)=R_{108}\) has the unique positive solution
\begin{equation}
 \boxed{
 \theta_{108}^{\mathrm{circ}}=\left(\frac{54}{\pi}\right)^2.}
 \label{eq:l9s102:g-theta}
\end{equation}
The corresponding realization identifies structurally
\begin{equation}
 R_{108}=\bar\lambda_{C,a}=r_{g,a}=\ell_P^{\mathrm{circ}}.
 \label{eq:l9s102:g-cuatro-longitudes}
\end{equation}
\end{theorem}

\begin{proof}
By \cref{eq:l9s102:g-transductor,eq:l9s102:g-cuanto-masa},
\[
 r_{g,a}(\theta)
 =\frac{G_a(\theta)m_{108,a}}{c_{\mathrm{int}}^2}
 =\theta\frac{\pi}{54}\ell_0.
\]
Equality with \(R_{108}=(54/\pi)\ell_0\) imposes
\(\theta=(54/\pi)^2\). Linearity in \(\theta\) and positivity of
the factors determine uniqueness.
\end{proof}

In the chart in which \(c_{\mathrm{int}}=(\mu\varepsilon)^{-1/2}\), the
transducer assumes the form
\begin{equation}
 \boxed{
 G_a=\left(\frac{54}{\pi}\right)^2
 \frac{t_0^2}{\hbar_a(\mu\varepsilon)^{5/2}}.}
 \label{eq:l9s102:g-mu-epsilon}
\end{equation}
The constitutive product \(\mu\varepsilon\) governs the causal cone and the
gravitational scale; the quotient \(Z^2=\mu/\varepsilon\) preserves the
impedance orientation.

\section{Bidirectional Covariance of Gravity and Action}
\label{sec:l9s102:g-hbar}

We define
\begin{equation}
 R_{\mathrm{act}}=\frac{\hbar_{\mathrm{pre}}}
 {\hbar_{\mathrm{ret}}}.
 \label{eq:l9s102:r-act}
\end{equation}
The chronological mass and the gravitational coupling transform contragrediently:
\begin{equation}
 \frac{m_{108,\mathrm{pre}}}{m_{108,\mathrm{ret}}}=R_{\mathrm{act}},
 \qquad
 \frac{G_{\mathrm{pre}}}{G_{\mathrm{ret}}}=R_{\mathrm{act}}^{-1}.
 \label{eq:l9s102:g-hbar-reciprocidad}
\end{equation}
Its product remains fixed:
\begin{equation}
 \boxed{
 G_a\hbar_a
 =\theta_{108}^{\mathrm{circ}}c_{\mathrm{int}}^5t_0^2,
 \qquad
 \ell_P^2=\frac{G_a\hbar_a}{c_{\mathrm{int}}^3}
 =\theta_{108}^{\mathrm{circ}}\ell_0^2.}
 \label{eq:l9s102:g-hbar-invariante}
\end{equation}
The action preserves the orientation of the sheet; the area of Planck preserves the
geometry common of both orientationes.

\begin{theorem}[Chart independence of the sections \(h\) and \(G\)]
\label{thm:l9s102:covariancia-unidades-g}
Let \(u_L,u_T,u_S\) be positive bases of the dimensional lines of
length, time and action, and replace them by
\[
 u_L'=\lambda_Lu_L,\qquad
 u_T'=\lambda_Tu_T,\qquad
 u_S'=\lambda_Su_S,
 \qquad \lambda_L,\lambda_T,\lambda_S>0.
\]
The coordinates of the same sections satisfy
\begin{equation}
 c_{\rm int}'=\frac{\lambda_T}{\lambda_L}c_{\rm int},\qquad
 t_0'=\frac{t_0}{\lambda_T},\qquad
 \hbar_a'=\frac{\hbar_a}{\lambda_S},\qquad
 G_a'=\frac{\lambda_S\lambda_T^3}{\lambda_L^5}G_a.
 \label{eq:l9s102:cambio-carta-g}
\end{equation}
For the induced basis
\(u_G=u_L^5u_T^{-3}u_S^{-1}\) one obtains
\begin{equation}
 \hbar_a'u_S'=\hbar_au_S,
 \qquad G_a'u_G'=G_au_G.
 \label{eq:l9s102:secciones-h-g-invariantes}
\end{equation}
Consequently,
\(h_a=2\pi\hbar_a\),
\(G_a=\theta_{108}^{\rm circ}c_{\rm int}^5t_0^2/\hbar_a\) and
\(G_a\hbar_a=\theta_{108}^{\rm circ}c_{\rm int}^5t_0^2\)
are identities of sections, not equalities dependent on a choice of
units. The metrological chart publishes their numerical coordinates after
the HMT--MD construction.
\end{theorem}

\begin{proof}
The first three transformations of
\eqref{eq:l9s102:cambio-carta-g} are the contragredient law of coordinates
in \(L\otimes T^{-1}\), \(T\) and \(S\). Substituting them into
\eqref{eq:l9s102:g-transductor} produces the fourth. The basis of the gravitational
line transforms by the inverse factor, proving
\eqref{eq:l9s102:secciones-h-g-invariantes}. The scalars
\(\pi\) and \(\theta_{108}^{\rm circ}\) are dimensionless and
remain invariant.
\end{proof}
 
\chapter{Barbero--Immirzi parameter, hexadic--octadic incidence and spectrum of the area operator}
\begingroup
\renewcommand{\chapter}[2][]{}%
\chapter{HMT derivation of the Barbero--Immirzi parameter and minimal area:
hexad--octad incidence, channels \(90/120\), and spectrum of the area
operator}
\label{chap:p3-final:51:barbero_area}
\section{Hexad--Octad Incidence Morphism and Conservation of the Oriented Channels \(90/120\)}

With the dodecad and the alignment of the
\cref{sec:w12-w24-elevacion} fixed, let \(H\) be a hexad in the HMT coordinates,
let \(\ell(H)\subset D\) be its image, and let
\(O_H=\iota_D(\ell(H))\) be its unique lift. Marking a point
and a pair of points of the residual hexad yields
\[
F_6(H)=\ell(H)\times\binom{\ell(H)}2,
\qquad
|F_6|=6\binom62=90.
\]
If the marked point can traverse the octad:
\[
F_8(H)=O_H\times\binom{\ell(H)}2,
\qquad
|F_8|=8\binom62=120.
\]
The inclusion \(\ell(H)\hookrightarrow O_H\) induces the incidence morphism
\[
 \mathcal I_{6\to8}:F_6(H)\hookrightarrow F_8(H).
\]
Normalizing the difference by the chosen
factor \(24\binom62\) yields
\[
\boxed{
\frac{90-120}{24\binom62}
=\frac{6-8}{24}
=-\frac1{12}.}
\]
The equality is a normalized defect of the combinatorial interface \(6\to8\). The factor \(24\binom62\) is part of the definition of this normalization; the inclusion alone determines the difference \(-30\). The operator \(\mathcal I_{6\to8}\) has marked incidence configurations as its domain; it is neither the transverse extractor of the dodecaphasic state nor the system of angular moments. \label{mdg:chap:barbero}

The transition from a hexad to an octad combines three structures of distinct
types: a finite incidence, two Lambert series associated with the
channels \(q_\pm\), and an identification with the real parameter of the Holst
action. The three are articulated through the incidence, the angular functional, and
the areal dictionary into the geometry of the
Ashtekar--Barbero connection developed in this same residence.

\section{Incidence cardinalities 90 and 120}

Let \(H\) be a set of six points and let \(O\) be a set of eight points
containing a distinguished copy of \(H\). Consider
\[
 \mathcal F_6=H\times\binom H2,
 \qquad
 \mathcal F_8=O\times\binom H2.
\]
The inclusion distinguished \(H\hookrightarrow O\) induce
\[
 \iota_{6,8}:\mathcal F_6\hookrightarrow\mathcal F_8,
 \qquad
 (h,\{h_1,h_2\})\longmapsto(h,\{h_1,h_2\}).
\]
This incidence morphism is the combinatorial antecedent of the
cardinals \(90/120\); it is neither the dodecaphasic extractor nor an angular series.

\begin{theorem}[Cardinalities of the Inclusion]
\[
 |\mathcal F_6|=6\binom62=90,
 \qquad
 |\mathcal F_8|=8\binom62=120.
\]
The normalized difference by twenty-four coordinates is
\[
 \frac{90-120}{24\binom62}=-\frac1{12}.
\]
\end{theorem}
\begin{proof}
Each element of the first coordinate can be combined with any of the
fifteen pairs of \(H\). The multiplication principle gives the two cardinalities;
the last identity follows by substituting \(\binom62=15\).
\end{proof}

\begin{figure}[htbp]
\centering
\begin{tikzpicture}[
  font=\small,
  setnode/.style={circle,draw,minimum size=6mm,inner sep=0pt,
    font=\scriptsize\bfseries},
  box/.style={draw,rounded corners=1.5mm,align=center,
    text width=5.0cm,inner xsep=6pt,inner ysep=5pt},
  arr/.style={-{Stealth[length=2.2mm]},line width=.7pt}]

  % Inclusión geométrica. Las flechas parten del contorno y no atraviesan nodos.
  \node[setnode,draw=hmtblue,fill=hmtlightblue] (h1) at (-2.3,1.25) {1};
  \node[setnode,draw=hmtblue,fill=hmtlightblue] (h2) at (-1.1,1.9) {2};
  \node[setnode,draw=hmtblue,fill=hmtlightblue] (h3) at (.1,1.25) {3};
  \node[setnode,draw=hmtblue,fill=hmtlightblue] (h4) at (.1,-.15) {4};
  \node[setnode,draw=hmtblue,fill=hmtlightblue] (h5) at (-1.1,-.8) {5};
  \node[setnode,draw=hmtblue,fill=hmtlightblue] (h6) at (-2.3,-.15) {6};
  \node[setnode,draw=hmtgreen,fill=hmtlightgreen] (o7) at (-3.2,.55) {7};
  \node[setnode,draw=hmtgreen,fill=hmtlightgreen] (o8) at (1.0,.55) {8};
  \node[draw=hmtblue,dashed,rounded corners=4mm,
    fit=(h1)(h2)(h3)(h4)(h5)(h6),inner sep=4mm,
    label={[hmtblue,font=\bfseries]above:{hexada \(H\), \(|H|=6\)}}] (H) {};
  \node[draw=hmtgreen,line width=.8pt,rounded corners=6mm,
    fit=(H)(o7)(o8),inner sep=5mm,
    label={[hmtgreen,font=\bfseries]below:{octada \(O\), \(|O|=8\), \(H\subset O\)}}] (O) {};

  \node[box,draw=hmtblue,fill=hmtlightblue] (f6) at (5.3,1.35)
    {Incidences defined on \(H\)\\[1mm] \(\mathcal F_6=H\times\binom H2\),\quad \(|\mathcal F_6|=6\binom62=90\)};
  \node[box,draw=hmtgreen,fill=hmtlightgreen] (f8) at (5.3,-1.0)
    {Extensión de incidencias a \(O\)\\[1mm]
     \(\mathcal F_8=O\times\binom H2\),\quad
     \(|\mathcal F_8|=8\binom62=120\)};
  \draw[arr,hmtblue] (O.east) -- ++(.45,0) |- (f6.west);
  \draw[arr,hmtgreen] (O.east) -- ++(.45,0) |- (f8.west);
  \draw[arr,hmtgold] (f6.south) -- node[right,align=left,font=\footnotesize]
    {inclusión \(\iota_{6,8}\)} (f8.north);

  \node[box,draw=hmtgold,fill=hmtlightgold,text width=10.9cm]
    (def) at (1.45,-3.80)
    {Defecto combinatorio orientado\\[-.2mm]
     con la normalización declarada:\\[1mm]
     \(\displaystyle
       \frac{|\mathcal F_6|-|\mathcal F_8|}{24\binom62}
       =\frac{90-120}{24\cdot15}=-\frac1{12}.\)};
  \draw[arr,hmtgold] (f8.south) -- (f8.south |- def.north);

  \node[align=center,text=hmtgray,font=\footnotesize,text width=11.2cm]
    at (1.45,-5.50)
    {Esta identidad pertenece al mapa de incidencias. No identifica el
     extractor dodecafásico, el semigrupo de torres ni el funcional angular.};
\end{tikzpicture}
 \caption{Inclusión de un subconjunto distinguido de seis elementos en uno de
ocho elementos y conteos de incidencia. Los cardinales \(90\) y \(120\) no fusionan el extractor
dodecafásico, el funcional angular ni el semigrupo de torres.}
\label{mdg:fig:incidencia-seis-ocho-rev9}
\end{figure}

Los índices iniciales de la jerarquía angular quedan así asociados a dos
espacios de incidencia: 90 configuraciones sobre la hexada y 120 sobre la
octada ambiente. El cociente \(-1/12\) es una identidad combinatoria de esta
normalización.

\section{Series de Lambert asociadas a los 2 sectores de incidencia}

Para \(0<q<1\), definimos simultáneamente
\[
 S_{90}(q)=\frac{q^{90}}{1-q^{270}},
 \qquad
 S_{120}(q)=\frac{q^{120}}{1-q^{360}},
 \qquad
 \Delta S_m=S_m(q_-)-S_m(q_+).
\]
Geometric expansions
\[
 \Delta S_{90}=\sum_{j\geq0}s_{90+270j},
 \qquad
 \Delta S_{120}=\sum_{j\geq0}s_{120+360j}
\]
place both terms within the global semigroup of
\cref{mdg:chap:torres}.

The angular functional associated with the transition \(6\to8\) is
\[
 K_{6\to8}=12\Delta S_{90}-\Delta S_{120},
\]
and does not follow from the incidence cardinalities alone. Its weights
are determined only after fixing the two series, the pole
residue \(1/24\), integer positivity, and the minimality criterion proved
below.
The resulting angular coordinate is
\[
\Gamma_{6\to8}
=\frac{\pi A\Cstar}{180}+K_{6\to8}.
\]
At the point HMT,
\[
\begin{aligned}
 \Delta S_{90}&=2.95169439297457621139847987861\times10^{-4},\\
 \Delta S_{120}&=1.96835112502951720093738706217\times10^{-5},\\
 \frac{\pi A\Cstar}{180}&=0.270485322934140078465586458790\ldots,
\end{aligned}
\]
and, therefore,
\[
 \boxed{\Gamma_{6\to8}
 =0.27400767269445927474725526077398041940\ldots.}
\]
The tables computed with the conjugate angular coordinate \(C^*\) rounded to fifteen decimal places give
\(0.2740076726944593\) in double-precision arithmetic; the preceding value
is the multiprecision evaluation of the complete regional phase.

\section{Diophantic equation of the weights}

The identity
\[
 \lim_{q\to1^-}(1-q)\frac{q^m}{1-q^n}=\frac1n
\]
associates to the combination
\(a\Delta S_{90}-b\Delta S_{120}\) the normalized residue
\[
 \frac a{270}-\frac b{360}.
\]

\begin{theorem}[Integer Weights with Residue \texorpdfstring{\(1/24\)}{1/24}]
The positive integer solutions of
\[
 \frac a{270}-\frac b{360}=\frac1{24}
\]
are
\[
 (a,b)=(12+3t,1+4t),\qquad t\in\ZZ_{\geq0}.
\]
The pair \((12,1)\) is the unique coordinatewise minimal solution and the
unique one of minimal \(\ell^1\) norm.
\end{theorem}
\begin{proof}
Multiplying by \(1080\) produces \(4a-3b=45\). Modulo three, one obtains
\(a\equiv0\pmod3\); writing \(a=3u\) yields \(4u-b=15\).
Positivity implies \(u=4+t\), \(b=1+4t\), with \(t\geq0\). Therefore,
\(a+b=13+7t\), which proves minimality.
\end{proof}

Primitivity \(\gcd(a,b)=1\) does not by itself select the minimal pair, since
other primitive solutions exist. The weight \(12:-1\) results from combining the
residue condition \(1/24\) with minimality. The choice of that residue is
a normalization condition of the construction; it is not deduced from the
decimal evaluation of \(\Gamma_{6\to8}\).

\newpage
\section{Monotonicity of the functional with respect to the conjugate angular coordinate}

Para \(m\in\{90,120\}\),
\[
 S_m'(q)=\frac{m q^{m-1}(1+2q^{3m})}{(1-q^{3m})^2}>0.
\]
Moreover,
\[
 \frac{\dd q_-}{\dd C}=\frac\pi{180}q_-,
 \qquad
 \frac{\dd q_+}{\dd C}=-\frac\pi{180}q_+.
\]

\begin{proposition}[Strict Monotonicity]
For \(1^\circ\leq C\leq3^\circ\), the function
\(C\mapsto\Gamma_{6\to8}(A,C)\) is strictly increasing. Consequently,
every value in its image has at most one preimage in that interval.
\end{proposition}
\begin{proof}
The derivative of each difference \(\Delta S_m\) is positive. On the
specified interval one obtains
\(\Delta S'_{120}<5.35\times10^{-4}\), whereas
\(\pi A/180>0.1273\). Therefore,
\(\Gamma'(C)>0.1267\). At \(C=\Cstar\),
\[
 \Gamma'(\Cstar)=0.1328995105618907\ldots.
\]
\end{proof}

Local injectivity excludes a second nearby angular coordinate with the
same functional evaluation. This control neither constructs nor selects
\(C^*\): the conjugate angular coordinate has already been obtained in the preceding regional branch.
The proposition only quantifies sensitivity of the downstream evaluation
once \(A\) and \(C^*\) are fixed; it does not authorize inversion of
\(\Gamma_{6\to8}\) or a physical value of \(\gamma_{\mathrm{BI}}\) to
redefine them.

 \label{ch:modular}

\section{Hexad--octad inclusion and associated functional realizations}

Let us fix a set \(H\) of six points and a set \(O\) of eight
points with a distinguished inclusion \(j:H\hookrightarrow O\). In the
Witt residence, \(H\) is a hexad of a dodecad and \(O\) its lifted octad;
in this chapter only the already fixed inclusion is used, not the
terminal value of Barbero--Immirzi. We write
\begin{equation}
\label{eq:modular-pairs-H}
\binom H2:=\{P\subset H:|P|=2\},
\qquad \left|\binom H2\right|=\binom62=15,
\end{equation}
and we define two incidence spaces with distinct types:
\begin{equation}
\label{eq:modular-F6-F8}
\mathcal F_6:=H\times\binom H2,
\qquad
\mathcal F_8:=O\times\binom H2.
\end{equation}
The inclusion of incidences is
\begin{equation}
\label{eq:modular-iota68}
\iota_{6,8}:\mathcal F_6\hookrightarrow\mathcal F_8,
\qquad
(h,P)\longmapsto(j(h),P).
\end{equation}
It is injective because \(j\) is injective and preserves the marked pair \(P\).

\begin{theorem}[Incidence and defect counts oriented]
\label{thm:modular-incidence-90-120}
The two spaces of \cref{eq:modular-F6-F8} satisfy
\begin{equation}
\label{eq:modular-counts-90-120}
|\mathcal F_6|=6\binom62=90,
\qquad
|\mathcal F_8|=8\binom62=120,
\qquad
|\mathcal F_8\setminus\iota_{6,8}(\mathcal F_6)|=30.
\end{equation}
Normalization by the twenty-four coordinates of the
dodecad--twenty-four pair and by the fifteen pairs of \(H\) produces
\begin{equation}
\label{eq:modular-incidence-defect}
\boxed{
\delta_{6,8}^{\rm inc}
=\frac{|\mathcal F_6|-|\mathcal F_8|}
{24\binom62}=-\frac1{12}.}
\end{equation}
The same fraction is obtained through the reciprocal functional associated
with the incorporated memory,
\begin{equation}
\label{eq:modular-reciprocal-defect}
\boxed{
\delta_{6,8}^{\rm rec}
=\frac{30}{120}-\frac{30}{90}=-\frac1{12}.}
\end{equation}
They are two exact realizations of the same oriented defect, but have
distinct domains and are not identified with \(\zeta(-1)\) without an
additional intertwiner.
\end{theorem}

\begin{proof}
The multiplicative principle applied to
\cref{eq:modular-F6-F8} gives \(6\cdot15=90\) and \(8\cdot15=120\).
The image of \(\iota_{6,8}\) has ninety elements, so the
complement has thirty. Substitution into
\cref{eq:modular-incidence-defect} gives
\(-30/(24\cdot15)=-1/12\), while
\cref{eq:modular-reciprocal-defect} is
\(1/4-1/3=-1/12\).
\end{proof}

The hexad--octad transition admits two functional realizations that
are not deduced from one another. The algebraic diagram, formulated on the spaces
that justify the counts, is
\begin{equation}
\label{eq:modular-68-two-readers}
\begin{aligned}
(H\xhookrightarrow{\,j\,}O)
&\longmapsto
\left(|\mathcal F_6|,|\mathcal F_8|\right)=(90,120)
\longmapsto\Gamma_{6\to8}=\gammaH,\\
(h_6,o_8,t)=(6,8,3)
&\longmapsto-\frac{o_8-h_6}{t}=-\frac23,
\end{aligned}
\end{equation}
where \(-2/3\) is the coefficient of the third coordinate of row \(13\) of the CKM constructor. The first realization transforms the incidence into area and modular memory; the second transforms the same transition into angular-mixing parameters. \(\boldsymbol\theta_{\rm CKM}=f(\gammaH)\) is not postulated: both outputs preserve a common antecedent.

The morphism determining the sectors \(90/120\) is
\(\iota_{6,8}\). Magnitude \(\gammaH\) is adopted as a theorem previously
demonstrated in the hexad--octad construction and is not reconstructed from
its decimal representation. The
operatorial consequences begin only after fixing the elementary
scale \(a_0\), defined below without using entropy or the
terminal product \(\gammaH a_0\).

\section{Correspondence between the Nonadic Interior and the Boundary Degrees of Freedom}

Let \(\mathbb T_3=\{0,1,2\}\). Every \(x\in\mathbb Z/9\mathbb Z\) has
a unique digit representation
\(x=x_0+3x_1\), with \(x_i\in\mathbb T_3\), and every
\(y\in\mathbb Z/27\mathbb Z\) a unique representation
\(y=y_0+3y_1+9y_2\). Therefore
\begin{equation}
\label{eq:modular-bulk-boundary-729}
\mathcal B=(\mathbb Z/9\mathbb Z)^3,
\qquad
\mathcal B_\partial=(\mathbb Z/27\mathbb Z)^2,
\qquad
|\mathcal B|=|\mathcal B_\partial|=3^6=729.
\end{equation}
Equality of cardinalities is accompanied by the explicit
map. If \(x_j=r_{2j-1}+3r_{2j}\) for \(j=1,2,3\), we define
\begin{equation}
\label{eq:modular-beta729}
\beta_{729}(x_1,x_2,x_3)
=\bigl(r_1+3r_2+9r_3,\;r_4+3r_5+9r_6\bigr).
\end{equation}

\begin{proposition}[Genealogical bijection \(729\leftrightarrow729\)]
\label{prop:modular-beta729}
The map \(\beta_{729}:\mathcal B\to\mathcal B_\partial\) is bijective and
preserves the ordered six-trit word. In general, it is not an
isomorphism of additive groups: regrouping the digits changes where the
carry is recorded.
\end{proposition}

\begin{proof}
Extracting the six digits from three nonadic coordinates and grouping them into
two triples is a permutation of the coordinates of \(\mathbb T_3^6\). The
positional decomposition is unique on both sides, so the
operation has an inverse. The additive caveat is verified, for example,
by adding words that generate a carry between \(r_3\) and \(r_4\): that
boundary separates the two coordinates of order twenty-seven, but crosses
two distinct coordinates in the nonadic packing.
\end{proof}

On the boundary torus \(\mathcal B_\partial\), we take the block of vertices
\begin{equation}
\label{eq:modular-block-A}
A=\{0,1,\ldots,8\}\times\{0,1,\ldots,8\}.
\end{equation}
Its oriented links crossing the boundary are separated by direction:
\begin{align}
\partial A_{90}
&=\{((-1,j),(0,j)),((8,j),(9,j)):0\leq j\leq8\},
\label{eq:modular-boundary90}\\
\partial A_{120}
&=\{((i,-1),(i,0)),((i,8),(i,9)):0\leq i\leq8\},
\label{eq:modular-boundary120}
\end{align}
where all coordinates are read modulo twenty-seven. Thus
\begin{equation}
\label{eq:modular-boundary-counts}
\partial A=\partial A_{90}\sqcup\partial A_{120},
\qquad
|\partial A_{90}|=|\partial A_{120}|=18.
\end{equation}
The subscripts \(90,120\) label the channel provenance; they do not say that
the boundary has ninety or one hundred and twenty links.

\begin{theorem}[Law areal of the cut triadic]
\label{thm:modular-triadic-cut}
In the cubic graph \(N\times N\times N\) with unit capacity, every
cut between two opposite faces has capacity at least \(N^2\), and a
transverse layer attains that bound. For \(N=9\),
\begin{equation}
\label{eq:modular-mincut-entropy}
\mathcal A_{\min}=81,
\qquad
S_{\rm tri}=81\log3.
\end{equation}
\end{theorem}

\begin{proof}
There are \(N^2\) lines of edges parallel to the normal direction, pairwise
disjoint. Every cut must intersect each one, so its capacity
is at least \(N^2\). A single transverse layer contains exactly
\(N^2\) edges and attains the minimum. If each cut link carries a
maximally mixed trit, its entropy is \(\log3\); additivity over the
eighty-one links gives \cref{eq:modular-mincut-entropy}.
\end{proof}

The entropy \(S_{\rm tri}\) measures the capacity of the cubic cut. The
modular entropy of the thirty-six links of
\cref{eq:modular-boundary-counts} measures another state and must not be equated with
it.

\section{Collective incidence entangler between interior and boundary}

In the sum
\(\ell^2(\mathcal F_6)\oplus\ell^2(\mathcal F_8)\), we define
\begin{equation}
\label{eq:modular-u90-u120}
u_{90}=\frac1{\sqrt{90}}\sum_{f\in\mathcal F_6}(\delta_f,0),
\qquad
u_{120}=\frac1{\sqrt{120}}\sum_{f\in\mathcal F_8}(0,\delta_f).
\end{equation}
In
\(\ell^2(\partial A_{90})\oplus\ell^2(\partial A_{120})\), let
\begin{equation}
\label{eq:modular-v90-v120}
v_{90}=\frac1{\sqrt{18}}\sum_{e\in\partial A_{90}}(\delta_e,0),
\qquad
v_{120}=\frac1{\sqrt{18}}\sum_{e\in\partial A_{120}}(0,\delta_e).
\end{equation}
The pairs \((u_{90},u_{120})\) and \((v_{90},v_{120})\) are orthonormal.
Therefore there exists a unique isometry between the collective subspaces
that satisfies
\begin{equation}
\label{eq:modular-J68}
\mathcal J_{6\to8}u_{90}=v_{90},
\qquad
\mathcal J_{6\to8}u_{120}=v_{120}.
\end{equation}
Define
\begin{align}
D_{\rm inc}
&=90|u_{90}\rangle\langle u_{90}|
+120|u_{120}\rangle\langle u_{120}|,
\label{eq:modular-Dinc}\\
D_{\rm area}
&=90|v_{90}\rangle\langle v_{90}|
+120|v_{120}\rangle\langle v_{120}|.
\label{eq:modular-Darea}
\end{align}

\begin{theorem}[Exact transport of the two labels]
\label{thm:modular-J68}
The map \(\mathcal J_{6\to8}\) is unitary between the two collective
planes and intertwines the incidence and area operators:
\begin{equation}
\label{eq:modular-J-intertwining}
\boxed{
\mathcal J_{6\to8}D_{\rm inc}
=D_{\rm area}\mathcal J_{6\to8}.}
\end{equation}
Therefore, every polynomial combination of the labels \(90,120\), in
particular the primitive combination \(12:-1\), is transported with the
same coefficients.
\end{theorem}

\begin{proof}
Orthonormality shows that \(\mathcal J_{6\to8}\) maps one
orthonormal basis to another. In \(u_{90}\), both sides of
\cref{eq:modular-J-intertwining} equal \(90v_{90}\); in \(u_{120}\),
both equal \(120v_{120}\). Equality on a basis proves the identity.
The functional calculus, in its polynomial form, preserves the intertwining.
\end{proof}

This map does not purport to inject ninety individual incidences into
eighteen links: it transports exactly the two uniform modes and their
spectral labels. Extending it to non-uniform fluctuations would require
a new incidence map; the modular operator defined here only uses
the collective subspace that has just been constructed.

\section{Determination of the Elementary Scale of the Spectrum of the Area Operator}

Let \(C^*\) the angular coordinate conjugate, expresada in grados, and fijemos
\begin{equation}
\label{eq:modular-theta-clock-a0}
\theta_{\rm clk}=\frac{C^*}{6}\frac{\pi}{180},
\qquad
\boxed{
a_0=\frac{1+2\cos(10^\circ)}{3}\,\theta_{\rm clk}.}
\end{equation}
The factor preceding \(\theta_{\rm clk}\) is not a decimal
correction. It is the normalized real character of the phase triplet
\begin{equation}
\label{eq:modular-c10-character}
\frac13\Tr\operatorname{diag}
\left(1,\ee^{\ii\pi/18},\ee^{-\ii\pi/18}\right)
=\frac{1+2\cos(10^\circ)}3.
\end{equation}
Thus, \(a_0\) is the spectral mean of the angular structure on the oriented
trit. It is positive, is fixed before the area operator, and is not
selected from entropy. The subsequent identification
\(\gammaH=\Gamma_{6\to8}\) realizes the HMT output in the Ashtekar--Barbero connection
\cite{Barbero,Immirzi}; it does not reopen its derivation.

\begin{remark}[Alcance of the simbolo \(a_0\)]
In this chapter, \(a_0\equiv a_0^{\rm area}\) denotes exclusively the
normalization of \cref{eq:modular-theta-clock-a0}. It is neither the Bohr radius
nor the initial cosmological scale factor, although those objects are traditionally written
with the same symbol. No identity in this chapter
allows one to substitute one for the other.
\end{remark}
 \section{Areal Law of the Triadic Cut: Capacity \(81\log 3\), Area Operator, Spectrum, and Collective Intertwiner of the Channels \(90/120\)}
\label{hap:sec:area-entrelazamiento-barbero-rev11}
\index{Barbero--Immirzi!area operator}
\index{modular entropy!modular Hamiltonian}

The HMT identification of the functional \(\Gamma_{6\to8}\) with the
Barbero--Immirzi coordinate is made on the same triadic boundary that supports the
incidence (90/120). The interior and its boundary sheet have equal cardinality,
\[
 \mathcal B=(\mathbb Z/9\mathbb Z)^3,
 \quad |\mathcal B|=729,
 \qquad
 \partial\mathcal B=(\mathbb Z/27\mathbb Z)^2,
 \quad |\partial\mathcal B|=729,
\]
by regrouping six trits into three nonadic coordinates or two
coordinates of order twenty-seven.

\begin{theorem}[Areal Law of Triadic Cutoff]
\label{hap:thm:ley-areal-corte-triadico-rev11}
In the cubic graph \(N\times N\times N\) with unit capacity, the minimum
cut between two opposite faces has capacity \(N^2\). For \(N=9\),
\[
 \mathcal A_{\min}=81,
 \qquad S_{\rm tri}=\mathcal A_{\min}\log3=81\log3.
\]
\end{theorem}
\begin{proof}
The \(N^2\) lines parallel to the normal direction are edge-disjoint
paths, so every cut intercepts them. A transverse layer contains
exactly \(N^2\) edges and attains the bound. Each triadic link contributes
\(\log3\) to the entropy of a maximally mixed state.
\end{proof}

On the boundary torus \(27\times27\), let \(A\) be the canonical block
\(9\times9\). Its boundary decomposes exactly into
\begin{equation}
 \partial A=\partial A_{90}\sqcup\partial A_{120},
 \qquad |\partial A_{90}|=|\partial A_{120}|=18.
 \label{hap:eq:corte-90-120-rev11}
\end{equation}
The horizontal links represent channel (90), and the vertical ones
channel (120). This decomposition turns the pair of incidences into an
operator decomposition of the boundary.

The areal entropy \(S_{\rm tri}=81\log3\) belongs to the maximally
mixed state of the cut links. The bilayer entropy
\(S_{\rm bi}(y)\) of the \cref{hap:def:entropia-bicapa-rev6} belongs to the
normalized distribution between the two sheets; the two objects retain
distinct domains.

Let \(N=\operatorname{diag}(0,1,2)\) be the number operator of a tritic link.
The areal realization declares the positive scale
\[
 \theta_{\rm clk}=\frac{C^*}{6}\frac{\pi}{180},
 \qquad
 a_0=\frac{1+2\cos(10^\circ)}{3}\,\theta_{\rm clk},
 \qquad \gamma_{\rm HMT}=\Gamma_{6\to8}.
\]
The formula for \(a_0\) constitutes the starting structure of this
realization; once that scale is fixed, the operator and its spectrum follow
exactly.

\begin{definition}[Area Operator HMT]
On the boundary sheet
\(\mathcal H_{\partial}=\bigotimes_{e\in\partial A}\mathbb C^3\), define
\begin{equation}
 \widehat{\mathcal A}_{\rm HMT}
 =\gamma_{\rm HMT}a_0\sum_{e\in\partial A}N_e.
 \label{hap:eq:operador-area-hmt-rev11}
\end{equation}
Its spectrum in the canonical cut is
\[
 \operatorname{Spec}\widehat{\mathcal A}_{\rm HMT}
 =\{n\gamma_{\rm HMT}a_0:0\leq n\leq72\},
\]
with multiplicities given by the coefficients of
\((1+z+z^2)^{36}\). Prolongation through compatible cuts produces the
protected family \(\mathcal A_n^{\rm prot}=n\gamma_{\rm HMT}a_0\).
\end{definition}

The intertwiner that fixes normalization is constructed in the subspace of
collective channels. Let \(u_{90},u_{120}\) be the uniform unit vectors
of the hexadic and octadic incidences, and let \(v_{90},v_{120}\) be the
uniform unit vectors of the two sets of links of
\eqref{hap:eq:corte-90-120-rev11}. The map
\begin{equation}
 \begin{aligned}
 \mathcal J_{6\to8}:{}
 &\operatorname{span}\{u_{90},u_{120}\}
 \longrightarrow
 \operatorname{span}\{v_{90},v_{120}\},\\
 &\mathcal J_{6\to8}u_{90}=v_{90},
 \qquad
 \mathcal J_{6\to8}u_{120}=v_{120}
 \end{aligned}
 \label{hap:eq:entrelazador-area-barbero-rev11}
\end{equation}
is an isometry between these collective subspaces. If
\[
 D_{\rm inc}=90|u_{90}\rangle\langle u_{90}|
 +120|u_{120}\rangle\langle u_{120}|,
 \qquad
 D_{\rm area}=90|v_{90}\rangle\langle v_{90}|
 +120|v_{120}\rangle\langle v_{120}|,
\]
then
\[
 \boxed{\mathcal J_{6\to8}D_{\rm inc}
 =D_{\rm area}\mathcal J_{6\to8}.}
\]
In particular, the primitive combination \(12:-1\) acts
with the same coefficients before and after \(\mathcal J_{6\to8}\). This
equality fixes the normalization dictionary that inserts
\(\gamma_{\rm HMT}=\Gamma_{6\to8}\) into the
Ashtekar--Barbero connection.

 The passage from hexad to octad is an incidence structure transformation.\par

In passing from \(6\) to \(8\), the available pairs, the internal relations, and the way in which orientation can be preserved change. That transition generates an oriented defect. The cardinal variation realizes a turn in the mode of organization.\par

The same sheet inversion is expressed by\par

\[
C^*\longmapsto -C^*.
\]

But different readers respond in different ways to that inversion.\par

The hexad–octad functional leading to the Barbero–Immirzi parameter is odd:
when the sheet is inverted, it changes sign. Geometric orientation is
reversed.\par

Bilayer entropy, by contrast, is even: it assigns the same entropy to \(C^*\) and \(-C^*\) when the probability distribution is observed. The total information may remain unchanged even though the orientation has been reversed.\par

CKM space realizes a third response. The inversion induces an exact rational reflection on the three angles:\par

\[
R_{\rm CKM}
=
M_{\rm CKM}
\operatorname{diag}(1,-1,1)
M_{\rm CKM}^{-1},
\]

with\par

\[
R_{\rm CKM}^2=I,
\qquad
\det R_{\rm CKM}=-1.
\]

This means that the change of sheet is transported to the mixing space as a true reflection. There exists an odd direction of rank one and a plane that remains fixed.\par

The counter-angle \(C^*\) is precisely the coordinate that changes. The phase constructed from \(A\) remains invariant. The transformation thus separates two memories:\par

\begin{itemize}[leftmargin=*,itemsep=0.35em]
\item an orientation memory;
\item a phase memory.
\end{itemize}

This provides a profound reading of CKM. The mixing angles are coordinates
of a space acted upon by the same involution that already organized the
hexad–octad transition.\par

The turn \(6\to8\) supplies the orientation operation that CKM realizes in its own codomain; physical causality belongs to the corresponding transducer.\par

We then have three sister publications of the same transit:\par

\begin{enumerate}[leftmargin=*,itemsep=0.65em]
\item \textbf{Geometric publication.}
The Barbero–Immirzi parameter orients the quantum of area.\par

\item \textbf{Informational publication.}
The reduced state and its modular Hamiltonian measure what information from that orientation remains on the boundary.\par

\item \textbf{Mixing publication.}
The CKM reflection transforms the mixing coordinates and isolates their odd component.\par
\end{enumerate}

\(\gamma_{\rm BI}\), entropy and the CKM angles remain distinct realizations of a common genealogical hinge. Unity resides in the antecedent prior to ramification. They share the hinge, but each preserves a distinct part of it.\par

This is one of the most beautiful expressions of the HMT methodology: the fundamental unit is recovered by tracing the genealogies backward until the common operation that generates distinct terminal results is found.\par

Gravity conserves the spin as the orientation of the area.\par
Information conserves the spin as memory of the state.\par
The mixture preserves the twist as a reflection of coordinates.\par

The same transformation «breathes» in three ways.\par

 \section{Causal separation between the area spectrum, modular memory and the Cartan--Holst realization}
The derivation of $\gamma_{\mathrm{BI}}$, the minimal area and the spectrum of the area operator concludes in this chapter. Joint reconstruction of occupation and provenance through the boundary modular Hamiltonian, and the Palatini--Cartan--Holst realization, retain their own domains, operators and proofs in Part IV. \endgroup
% Secciones 4--7 del delta gravitatorio íntegro.
\section{Hexad--octad incidence, Barbero--Immirzi parameter, and area spectrum}
\label{sec:l9s102:barbero-area}

The areal derivation receives as antecedents the incidence inclusion
\(6\to8\), the TPK cardinals \(90/120\), the normalized defect
\(-1/12\), and the spectral channels \(q_\pm\). The
hexad--octad functional determines the Barbero--Immirzi parameter and the
normalization of the area operator. This residence contains the HMT derivation
and the spectrum; the Palatini--Cartan--Holst realization occupies a subsequent
residence in Dimensional Mechanics.

Let \(H\subset O\), with \(|H|=6\) and \(|O|=8\), and preserve the marked
pair of the hexad in the inclusion
\begin{equation}
 \iota_{6,8}:\mathcal F_6=H\times\binom H2
 \lhook\joinrel\longrightarrow
 \mathcal F_8=O\times\binom H2.
 \label{eq:l9s102:barbero-inclusion-incidencia}
\end{equation}
The internal count gives
\begin{equation}
 |\mathcal F_6|=6\binom62=90,
 \qquad |\mathcal F_8|=8\binom62=120,
 \qquad
 \frac{90-120}{24\binom62}=-\frac1{12}.
 \label{eq:l9s102:barbero-conteos}
\end{equation}
Once \(\pi\), the angle \(A\) and the
unique counterangle \(C^*\) have been published, the two channels are
\begin{equation}
 q_\pm=\exp\!\left[-\frac{\pi}{180}(A\pm C^*)\right],
 \qquad 0<q_+<q_-<1,
 \label{eq:l9s102:barbero-qpm}
\end{equation}
and their oriented-return series are
\begin{equation}
 S_{90}(q)=\frac{q^{90}}{1-q^{270}},\qquad
 S_{120}(q)=\frac{q^{120}}{1-q^{360}},\qquad
 \Delta S_m=S_m(q_-)-S_m(q_+).
 \label{eq:l9s102:barbero-series}
\end{equation}

\begin{theorem}[Diophantine determination of the Barbero--Immirzi functional]
\label{thm:l9s102:barbero-funcional} Among positive combinations \(a\Delta S_{90}-b\Delta S_{120}\), the residue condition \(1/24\) and coordinatewise minimality determine the unique primitive pair \((a,b)=(12,1)\). Thus the dimensionless output is
\begin{equation}
 \boxed{
 \gamma_{\rm HMT}=\Gamma_{6\to8}
 =\frac{{\pi}AC^*}{180}
   +12\Delta S_{90}-\Delta S_{120}.}
 \label{eq:l9s102:barbero-gamma}
\end{equation}
This formula receives the already generated angular pair and does not use
 the conventional value of the Barbero--Immirzi parameter as input.
\end{theorem}

\begin{proof}
The residue of \(S_{90}\) at \(q=1\) is \(1/270\), and that of \(S_{120}\) is \(1/360\). The required condition is equivalent to
\[
 \frac a{270}-\frac b{360}=\frac1{24}
 \quad\Longleftrightarrow\quad 4a-3b=45.
\]
Its positive integer solutions are \((a,b)=(12+3t,1+4t)\), \(t\ge0\); the unique minimal solution is \((12,1)\). Substituting this pair into the incidence functional proves \eqref{eq:l9s102:barbero-gamma}.
\end{proof}

The causal sequence is fixed without exchanging roles:
\begin{equation}
 (\mathcal F_6\hookrightarrow\mathcal F_8)
 \longrightarrow (90,120,-1/12)
 \longrightarrow (A,C^*,q_\pm)
 \longrightarrow \Gamma_{6\to8}=\gamma_{\rm HMT}
 \longrightarrow \widehat{\mathcal A}_{\rm phys}.
 \label{eq:l9s102:barbero-cadena-completa}
\end{equation}
The subsequent realization in the Holst action recognizes the same dimensionless section; it does not participate in its selection.

For the two boundary sectors, let
\begin{equation}
 N_m=\sum_{e\in\partial A_m}N_e,
 \qquad m\in\{90,120\},
 \qquad \operatorname{Spec}(N_e)=\{0,1,2\}.
 \label{eq:l9s102:n90-n120}
\end{equation}
The physical area operator is written
\begin{equation}
 \boxed{
 \frac{\widehat{\mathcal A}_{\mathrm{phys}}}{\ell_P^2}
 =\gamma_{\mathrm{HMT}}a_0(N_{90}+N_{120}).}
 \label{eq:l9s102:area-operador}
\end{equation}
On \(36\) links qutrit,
\begin{equation}
 \operatorname{Spec}(N_{90}+N_{120})=\{0,1,\ldots,72\},
 \label{eq:l9s102:area-numero-spectrum}
\end{equation}
and the multiplicity of the eigenvalue \(n\) is the coefficient of \(z^n\) in
\((1+z+z^2)^{36}\). From here it follows
\begin{equation}
 \operatorname{Spec}\!\left(
 \frac{\widehat{\mathcal A}_{\mathrm{phys}}}{\ell_P^2}\right)
 =\{n\gamma_{\mathrm{HMT}}a_0:0\le n\le72\}.
 \label{eq:l9s102:area-spectrum}
\end{equation}

\section{Total Geometry and Origin Memory}
\label{sec:l9s102:area-hamiltoniano}

The modular Hamiltonian of the boundary is
\begin{equation}
 \boxed{
 K_A=-\log\rho_A
 =cI+\Delta_{\mathrm{mod}}(90N_{90}+120N_{120}).}
 \label{eq:l9s102:hamiltoniano-modular}
\end{equation}
Let
\begin{equation}
 N=N_{90}+N_{120},
 \qquad
 D=90N_{90}+120N_{120}.
 \label{eq:l9s102:n-d}
\end{equation}
The transition matrix
\(\left(\begin{smallmatrix}1&1\\90&120\end{smallmatrix}\right)\)
has determinant \(30\), and its inverse reconstructs the sectors:
\begin{equation}
 \boxed{
 N_{90}=4N-\frac{D}{30},
 \qquad
 N_{120}=\frac{D}{30}-3N.}
 \label{eq:l9s102:reconstruccion-sectores}
\end{equation}
The area publishes the total number of excitations; the modular Hamiltonian
publishes their provenance sector. The joint pair reconstructs the complete
occupation of the boundary.

\section{Bilayer state and oriented information}
\label{sec:l9s102:estado-bicapa}

Let \(R=R^*=R^{-1}\) be the sheet involution and
\(y=C^*_{\mathrm{rad}}\). The normalized state is
\begin{equation}
 \boxed{
 \rho_y=\frac{e^{-yR}}{2\cosh y}.}
 \label{eq:l9s102:rho-y}
\end{equation}
Its modular Hamiltonian and its entropy are
\begin{equation}
 K_y=-\log\rho_y=\log(2\cosh y)I+yR,
 \label{eq:l9s102:ky}
\end{equation}
\begin{equation}
 S(\rho_y)=\log(2\cosh y)-y\tanh y.
 \label{eq:l9s102:entropia-y}
\end{equation}
The inversion \(y\mapsto-y\) leaves entropy fixed and produces the oriented
informational distance
\begin{equation}
 \boxed{
 D(\rho_y\Vert\rho_{-y})=2y\tanh y.}
 \label{eq:l9s102:entropia-relativa-hojas}
\end{equation}
The even component measures the spectral distribution; the odd component measures
the informational cost of reversing the orientation.

\section{How much informational energy and gravitational cell of medium action}
\label{sec:l9s102:horizonte-media-accion}

The periodicity of \(108\) states determines
\begin{equation}
 t_P=\frac{54}{\pi}t_0,
 \qquad
 E_P=\frac{\hbar}{t_P}=\frac{h}{108t_0}.
 \label{eq:l9s102:ep-tp}
\end{equation}
The thermal realization of the tritico clock satisfies
\begin{equation}
 \boxed{
 E_P=k_B\Theta_{\mathrm{clk}}\log3.}
 \label{eq:l9s102:planck-trit}
\end{equation}
For the horizon of radius \(R\),
\begin{equation}
 E_\Lambda(R)=\frac{c^4R}{2G},
 \qquad
 E_\Lambda(R)=T_H(R)S_H(R).
 \label{eq:l9s102:energia-horizonte}
\end{equation}
In the Planck cell,
\begin{equation}
 \boxed{
 E_\Lambda
 =T_HS_H
 =\frac{E_P}{2}
 =\frac{k_B\Theta_{\mathrm{clk}}\log3}{2}.}
 \label{eq:l9s102:confluencia-horizonte}
\end{equation}
Therefore,
\begin{equation}
 \boxed{
 E_\Lambda t_P=\frac{\hbar}{2},
 \qquad
 E_\Lambda T_{108}=\frac h2.}
 \label{eq:l9s102:media-accion}
\end{equation}
Vacuum energy, area entropy, chronological temperature, and
action share a single confluence cell with typed domains.

When \(S_H/k_B=\pi\), the multiplicity efectiva satisfies
\begin{equation}
 \Omega_{\mathrm{eff}}=e^{S_H/k_B}=e^\pi.
 \label{eq:l9s102:omega-e-pi}
\end{equation}
The hyperbolic Euler operator possesses
\begin{equation}
 \operatorname{Spec}(e^{\pi H})=\{e^\pi,e^{-\pi}\}.
 \label{eq:l9s102:e-pi-spectrum}
\end{equation}
Equality of the scalars establishes spectral confluence; the
intertwiner between microscopic horizon multiplicity and Perron
dynamics belongs to its specific physical residence.
 
\chapter{Exact reversibility, erasure cost and zero Landauer limit}
\begingroup
\renewcommand{\chapter}[2][]{}%
\chapter{Reversible update, deletion cost, and exact domain of the Landauer null limit}
\label{chap:parte-iv-57}

The full transport preserves fibre information and may have zero logical cost; the projection or reinitialization that discards the genealogy is subject to Landauer's cost. The two domains are separated explicitly.

\section{Logical reversibility and Landauer's principle in TPK fibers}
\label{sec:informacion-fibra-landauer-rev7}
\index{Landauer's principle}\index{information!fiber}

Let \((\Omega,2^\Omega,\mu)\) be a finite probability space, let
\(X:\Omega\to\mathcal X\) be a random variable with values in the space
\(\mathcal X\) of complete states, and let \(q:\mathcal X\to Y\) be a reader.
The chain rule gives
\[
 \mathsf H(X)
 =
 \mathsf H(q(X))+\mathsf H(X\mid q(X)).
\]
The second summand measures the information of the fiber that the reader does not
publish.

\begin{theorem}[Landauer zero for full transport and projection cost]
\label{thm:landauer-relativo-lector-rev7}
If \(U:\mathcal X\to\mathcal X\) is a bijective update of the complete
HMT state,
then
\[
 \boxed{\mathsf H(X\mid U(X))=0.}
\]
For a projected output \(q(U(X))\), the discarded information is
\[
 \mathsf H(X\mid q(U(X))).
\]
A physical reset that identifies distinct memory states must
account for that information in the environment.
\end{theorem}

\begin{proof}
Bijectivity implies
\(X=U^{-1}(U(X))\); therefore \(X\) is determined by \(U(X)\) and the
conditional entropy is zero. For the reader \(q\), the chain rule
separates published entropy from fiber entropy. If the reset operation is
not injective, the output no longer permits reconstruction of the memory
state. Landauer's thermodynamic principle bounds the cost of implementation in
terms of that erased entropy; the reversible update does not by itself
contribute an erasure term.
\end{proof}

In nats and bits, respectively,
\[
 Q\ge k_B\Theta\,\mathsf H_{\rm nat},
 \qquad
 Q\ge k_B\Theta\log 2\,\mathsf H_{\rm bit}.
\]
For bijective transport of the complete state, the erasure entropy is
zero and, therefore, the minimum Landauer term associated exclusively with
erasure satisfies
\[
\boxed{Q_{\rm L}^{\min}[U]=k_B\Theta\,
 \mathsf H(X\mid U(X))=0.}
\]
For the restart of a uniform TRIT, on the other hand,
\[
 \boxed{Q_{\rm reset}\ge k_B\Theta\log3.}
\]
These two equations locate the cost in information loss: the cost appears when projection or resetting forgets the fiber. Preparation, projection, and resetting may produce entropy; reversible transport preserves genealogical information. The null value is the logical Landauer limit for that complete update, whereas the dynamical costs of control, finite speed, and coupling to the environment belong to the physical realization of the device.

\begin{statusbox}[title={Scope of zero Landauer cost}] The equality \(\mathsf H(X\mid U(X))=0\) is exact for a bijective update of the integral holonomic state. The bound \(Q_{\rm reset}\ge k_B\Theta\log3\) specializes Landauer's principle to erasure of a uniform triadic decision. Realization of both formulas in joules uses the Boltzmann transducer and a declared physical temperature.
\end{statusbox}
 
\paragraph{Edge Information Transduction in Minimum Triadic Cut Capacity.} The information preserved by the boundary reappears geometrically in the minimal triadic cut and its capacity. \endgroup

\chapter{Minimal area of the triadic cut and capacity \texorpdfstring{\(81\log 3\)}{81 log 3}}
\begingroup
\renewcommand{\chapter}[2][]{}%
\chapter{Minimum area of the triadic cut, capacity \texorpdfstring{\(81\log 3\)}{81 log 3}, and collective incidence transport}
\label{chap:parte-iv-58}

Areal quantization begins with the nonadic interior, its boundary, and the collective transport of incidence. The minimal cut contains \(81\) units and its qutrit capacity is \(81\log 3\); the result precedes its gravitational realization.

\section{Areal quantization, modular dynamics, and rational CKM transformation}
\label{rm:ch:modular}

\subsection{Hexad--octad inclusion and associated functional realizations}

Let us fix a set \(H\) of six points and a set \(O\) of eight
points with a distinguished inclusion \(j:H\hookrightarrow O\). In the
Witt residence, \(H\) is a hexad of a dodecad and \(O\) its lifted octad;
in this chapter only the already fixed inclusion is used, not the
terminal value of Barbero--Immirzi. We write
\begin{equation}
\label{rm:eq:modular-pairs-H}
\binom H2:=\{P\subset H:|P|=2\},
\qquad \left|\binom H2\right|=\binom62=15,
\end{equation}
and we define two incidence spaces with distinct types:
\begin{equation}
\label{rm:eq:modular-F6-F8}
\mathcal F_6:=H\times\binom H2,
\qquad
\mathcal F_8:=O\times\binom H2.
\end{equation}
The inclusion of incidences is
\begin{equation}
\label{rm:eq:modular-iota68}
\iota_{6,8}:\mathcal F_6\hookrightarrow\mathcal F_8,
\qquad
(h,P)\longmapsto(j(h),P).
\end{equation}
It is injective because \(j\) is injective and preserves the marked pair \(P\).

\begin{theorem}[Incidence and defect counts oriented]
\label{rm:thm:modular-incidence-90-120}
The two spaces of \cref{rm:eq:modular-F6-F8} satisfy
\begin{equation}
\label{rm:eq:modular-counts-90-120}
|\mathcal F_6|=6\binom62=90,
\qquad
|\mathcal F_8|=8\binom62=120,
\qquad
|\mathcal F_8\setminus\iota_{6,8}(\mathcal F_6)|=30.
\end{equation}
Normalization by the twenty-four coordinates of the
dodecad--twenty-four pair and by the fifteen pairs of \(H\) produces
\begin{equation}
\label{rm:eq:modular-incidence-defect}
\boxed{
\delta_{6,8}^{\rm inc}
=\frac{|\mathcal F_6|-|\mathcal F_8|}
{24\binom62}=-\frac1{12}.}
\end{equation}
The same fraction is obtained through the reciprocal functional associated
with the incorporated memory,
\begin{equation}
\label{rm:eq:modular-reciprocal-defect}
\boxed{
\delta_{6,8}^{\rm rec}
=\frac{30}{120}-\frac{30}{90}=-\frac1{12}.}
\end{equation}
They are two exact realizations of the same oriented defect, but have
distinct domains and are not identified with \(\zeta(-1)\) without an
additional intertwiner.
\end{theorem}

\begin{proof}
The multiplicative principle applied to
\cref{rm:eq:modular-F6-F8} gives \(6\cdot15=90\) and \(8\cdot15=120\).
The image of \(\iota_{6,8}\) has ninety elements, so the
complement has thirty. Substitution into
\cref{rm:eq:modular-incidence-defect} gives
\(-30/(24\cdot15)=-1/12\), while
\cref{rm:eq:modular-reciprocal-defect} is
\(1/4-1/3=-1/12\).
\end{proof}

The hexad--octad transition admits two functional realizations that
are not deduced from one another. The algebraic diagram, formulated on the spaces
that justify the counts, is
\begin{equation}
\label{rm:eq:modular-68-two-readers}
\begin{aligned}
(H\xhookrightarrow{\,j\,}O)
&\longmapsto
\left(|\mathcal F_6|,|\mathcal F_8|\right)=(90,120)
\longmapsto\Gamma_{6\to8}=\gammaH,\\
(h_6,o_8,t)=(6,8,3)
&\longmapsto-\frac{o_8-h_6}{t}=-\frac23,
\end{aligned}
\end{equation}
where \(-2/3\) is the coefficient of the third coordinate of row \(13\) of the CKM constructor. The first realization transforms the incidence into area and modular memory; the second transforms the same transition into angular-mixing parameters. \(\boldsymbol\theta_{\rm CKM}=f(\gammaH)\) is not postulated: both outputs preserve a common antecedent.

The morphism determining the sectors \(90/120\) is
\(\iota_{6,8}\). Magnitude \(\gammaH\) is adopted as a theorem previously
demonstrated in the hexad--octad construction and is not reconstructed from
its decimal representation. The
operatorial consequences begin only after fixing the elementary
scale \(a_0\), defined below without using entropy or the
terminal product \(\gammaH a_0\).

\subsection{Correspondence between the Nonadic Interior and the Boundary Degrees of Freedom}

Let \(\mathbb T_3=\{0,1,2\}\). Every \(x\in\mathbb Z/9\mathbb Z\) has
a unique digit representation
\(x=x_0+3x_1\), with \(x_i\in\mathbb T_3\), and every
\(y\in\mathbb Z/27\mathbb Z\) a unique representation
\(y=y_0+3y_1+9y_2\). Therefore
\begin{equation}
\label{rm:eq:modular-bulk-boundary-729}
\mathcal B=(\mathbb Z/9\mathbb Z)^3,
\qquad
\mathcal B_\partial=(\mathbb Z/27\mathbb Z)^2,
\qquad
|\mathcal B|=|\mathcal B_\partial|=3^6=729.
\end{equation}
Equality of cardinalities is accompanied by the explicit
map. If \(x_j=r_{2j-1}+3r_{2j}\) for \(j=1,2,3\), we define
\begin{equation}
\label{rm:eq:modular-beta729}
\beta_{729}(x_1,x_2,x_3)
=\bigl(r_1+3r_2+9r_3,\;r_4+3r_5+9r_6\bigr).
\end{equation}

\begin{proposition}[Genealogical bijection \(729\leftrightarrow729\)]
\label{rm:prop:modular-beta729}
The map \(\beta_{729}:\mathcal B\to\mathcal B_\partial\) is bijective and
preserves the ordered six-trit word. In general, it is not an
isomorphism of additive groups: regrouping the digits changes where the
carry is recorded.
\end{proposition}

\begin{proof}
Extracting the six digits from three nonadic coordinates and grouping them into
two triples is a permutation of the coordinates of \(\mathbb T_3^6\). The
positional decomposition is unique on both sides, so the
operation has an inverse. The additive caveat is verified, for example,
by adding words that generate a carry between \(r_3\) and \(r_4\): that
boundary separates the two coordinates of order twenty-seven, but crosses
two distinct coordinates in the nonadic packing.
\end{proof}

On the boundary torus \(\mathcal B_\partial\), we take the block of vertices
\begin{equation}
\label{rm:eq:modular-block-A}
A=\{0,1,\ldots,8\}\times\{0,1,\ldots,8\}.
\end{equation}
Its oriented links crossing the boundary are separated by direction:
\begin{align}
\partial A_{90}
&=\{((-1,j),(0,j)),((8,j),(9,j)):0\leq j\leq8\},
\label{rm:eq:modular-boundary90}\\
\partial A_{120}
&=\{((i,-1),(i,0)),((i,8),(i,9)):0\leq i\leq8\},
\label{rm:eq:modular-boundary120}
\end{align}
where all coordinates are read modulo twenty-seven. Thus
\begin{equation}
\label{rm:eq:modular-boundary-counts}
\partial A=\partial A_{90}\sqcup\partial A_{120},
\qquad
|\partial A_{90}|=|\partial A_{120}|=18.
\end{equation}
The subscripts \(90,120\) label the channel provenance; they do not say that
the boundary has ninety or one hundred and twenty links.

\begin{theorem}[Law areal of the cut triadic]
\label{rm:thm:modular-triadic-cut}
In the cubic graph \(N\times N\times N\) with unit capacity, every
cut between two opposite faces has capacity at least \(N^2\), and a
transverse layer attains that bound. For \(N=9\),
\begin{equation}
\label{rm:eq:modular-mincut-entropy}
\mathcal A_{\min}=81,
\qquad
S_{\rm tri}=81\log3.
\end{equation}
\end{theorem}

\begin{proof}
There are \(N^2\) lines of edges parallel to the normal direction, pairwise
disjoint. Every cut must intersect each one, so its capacity
is at least \(N^2\). A single transverse layer contains exactly
\(N^2\) edges and attains the minimum. If each cut link carries a
maximally mixed trit, its entropy is \(\log3\); additivity over the
eighty-one links gives \cref{rm:eq:modular-mincut-entropy}.
\end{proof}

The entropy \(S_{\rm tri}\) measures the capacity of the cubic cut. The
modular entropy of the thirty-six links of
\cref{rm:eq:modular-boundary-counts} measures another state and must not be equated with
it.

\subsection{Collective incidence entangler between interior and boundary}

In the sum
\(\ell^2(\mathcal F_6)\oplus\ell^2(\mathcal F_8)\), we define
\begin{equation}
\label{rm:eq:modular-u90-u120}
u_{90}=\frac1{\sqrt{90}}\sum_{f\in\mathcal F_6}(\delta_f,0),
\qquad
u_{120}=\frac1{\sqrt{120}}\sum_{f\in\mathcal F_8}(0,\delta_f).
\end{equation}
In
\(\ell^2(\partial A_{90})\oplus\ell^2(\partial A_{120})\), let
\begin{equation}
\label{rm:eq:modular-v90-v120}
v_{90}=\frac1{\sqrt{18}}\sum_{e\in\partial A_{90}}(\delta_e,0),
\qquad
v_{120}=\frac1{\sqrt{18}}\sum_{e\in\partial A_{120}}(0,\delta_e).
\end{equation}
The pairs \((u_{90},u_{120})\) and \((v_{90},v_{120})\) are orthonormal.
Therefore there exists a unique isometry between the collective subspaces
that satisfies
\begin{equation}
\label{rm:eq:modular-J68}
\mathcal J_{6\to8}u_{90}=v_{90},
\qquad
\mathcal J_{6\to8}u_{120}=v_{120}.
\end{equation}
Define
\begin{align}
D_{\rm inc}
&=90|u_{90}\rangle\langle u_{90}|
+120|u_{120}\rangle\langle u_{120}|,
\label{rm:eq:modular-Dinc}\\
D_{\rm area}
&=90|v_{90}\rangle\langle v_{90}|
+120|v_{120}\rangle\langle v_{120}|.
\label{rm:eq:modular-Darea}
\end{align}

\begin{theorem}[Exact transport of the two labels]
\label{rm:thm:modular-J68}
The map \(\mathcal J_{6\to8}\) is unitary between the two collective
planes and intertwines the incidence and area operators:
\begin{equation}
\label{rm:eq:modular-J-intertwining}
\boxed{
\mathcal J_{6\to8}D_{\rm inc}
=D_{\rm area}\mathcal J_{6\to8}.}
\end{equation}
Therefore, every polynomial combination of the labels \(90,120\), in
particular the primitive combination \(12:-1\), is transported with the
same coefficients.
\end{theorem}

\begin{proof}
Orthonormality shows that \(\mathcal J_{6\to8}\) maps one
orthonormal basis to another. In \(u_{90}\), both sides of
\cref{rm:eq:modular-J-intertwining} equal \(90v_{90}\); in \(u_{120}\),
both equal \(120v_{120}\). Equality on a basis proves the identity.
The functional calculus, in its polynomial form, preserves the intertwining.
\end{proof}

This map does not purport to inject ninety individual incidences into
eighteen links: it transports exactly the two uniform modes and their
spectral labels. Extending it to non-uniform fluctuations would require
a new incidence map; the modular operator defined here only uses
the collective subspace that has just been constructed.

\subsection{What is the elementary and minimal scale of the area operator}

Let \(C^*\) the angular coordinate conjugate, expresada in grados, and fijemos
\begin{equation}
\label{rm:eq:modular-theta-clock-a0}
\theta_{\rm clk}=\frac{C^*}{6}\frac{\pi}{180},
\qquad
\boxed{
a_0=\frac{1+2\cos(10^\circ)}{3}\,\theta_{\rm clk}.}
\end{equation}
The factor preceding \(\theta_{\rm clk}\) is not a decimal
correction. It is the normalized real character of the phase triplet
\begin{equation}
\label{rm:eq:modular-c10-character}
\frac13\Tr\operatorname{diag}
\left(1,\ee^{\ii\pi/18},\ee^{-\ii\pi/18}\right)
=\frac{1+2\cos(10^\circ)}3.
\end{equation}
Thus, \(a_0\) is the spectral mean of the angular structure on the oriented
trit. It is positive, is fixed before the area operator, and is not
selected from entropy. The subsequent identification
\(\gammaH=\Gamma_{6\to8}\) realizes the HMT output in the Ashtekar--Barbero connection
\cite{Barbero1995,Immirzi1997}; it does not reopen its derivation.

\begin{remark}[Alcance of the simbolo \(a_0\)]
In this chapter, \(a_0\equiv a_0^{\rm area}\) denotes exclusively the
normalization of \cref{rm:eq:modular-theta-clock-a0}. It is neither the Bohr radius
nor the initial cosmological scale factor, although those objects are traditionally written
with the same symbol. No identity in this chapter
allows one to substitute one for the other.
\end{remark}
 
\paragraph{Construction of the area operator from the triad cut and channel separation.} The cut and its capacity determine an area operator; its spectrum must preserve total occupation and channel provenance separately. \endgroup

\chapter{Spectrum of the area operator and reconstruction of channel provenance}
\begingroup
\renewcommand{\chapter}[2][]{}%
\chapter{Spectrum of the area operator and joint reconstruction of occupation and origin}
\label{chap:parte-iv-59}

The sectors \(90\) and \(120\) define compatible occupancy and
memory operators. Inverting the linear system reconstructs both channels
and prevents total area from erasing provenance.

\section{Area and modular operators in the sectors \texorpdfstring{\(90\)}{90} and \texorpdfstring{\(120\)}{120}}

In each of the two families of
\cref{rm:eq:modular-boundary-counts} there are eighteen qutrit channels. For
\(m\in\{90,120\}\), define
\begin{equation}
\label{rm:eq:modular-number-operators}
N_m=\sum_{e\in\partial A_m}N_e,
\qquad \Spec(N_e)=\{0,1,2\}.
\end{equation}
The Barbero--Immirzi output is used here as a prior result, without
rederiving it. With the area normalization constructed in
\cref{rm:eq:modular-theta-clock-a0}, one obtains
\begin{equation}
\label{rm:eq:modular-area-operator}
\boxed{
\frac{\widehat{\mathcal A}_{\rm phys}}{\ell_P^2}
=\gammaH a_0(N_{90}+N_{120}),}
\qquad
\gammaH a_0=0.0016755940749224991.
\end{equation}

\begin{theorem}[Finite spectrum of the area quantum]
\label{rm:thm:modular-area-spectrum}
On the thirty-six links of the cut,
\begin{equation}
\label{rm:eq:modular-total-number-spectrum}
\Spec(N_{90}+N_{120})=\{0,1,\ldots,72\}.
\end{equation}
The multiplicity of the eigenvalue \(n\) is the coefficient of \(z^n\) in
\((1+z+z^2)^{36}\). Consequently,
\begin{equation}
\label{rm:eq:modular-area-spectrum}
\boxed{
\Spec\!\left(\frac{\widehat{\mathcal A}_{\rm phys}}{\ell_P^2}\right)
=\{n\gammaH a_0:0\leq n\leq72\}.}
\end{equation}
The compatible cuts prolong the family
\(\mathcal A_n^{\rm prot}=n\gammaH a_0\ell_P^2\).
\end{theorem}

\begin{proof}
Each operator \(N_e\) has a spectral generating function
\(1+z+z^2\). The thirty-six factors commute and act on distinct
tensor factors; therefore, the generating function of the total is
\((1+z+z^2)^{36}\), whose degrees range over all integers from zero to
seventy-two. Multiplication by the positive scalar
\(\gammaH a_0\ell_P^2\) gives \cref{rm:eq:modular-area-spectrum}.
\end{proof}

Two quantization procedures must be distinguished. The cut law \cref{rm:thm:modular-triadic-cut} counts eighty-one unit capacities and determines \(81\log3\); the operator \cref{rm:eq:modular-area-operator} quantizes occupations of thirty-six response links, up to seventy-two number trits. They share the same boundary sheet, but not the same observable.

The torsion modular produces
\begin{equation}
\label{rm:eq:modular-delta-value}
\Dmod
=0.0001111970500599773249414566131933856\ldots,
\end{equation}
and the boundary Hamiltonian
\begin{equation}
\label{rm:eq:modular-hamiltonian}
\boxed{
\HHM^{(A)}=-\log\rho_A
=cI+\Dmod(90N_{90}+120N_{120}).}
\end{equation}
The qualifier \emph{holographic} does not designate an analogy: the
incidence \(6\to8\) is transported to the boundary sets
\(\partial A_{90}\) and \(\partial A_{120}\), and the reduced state resides
in the algebra generated by its number operators.

\begin{theorem}[Joint reconstruction of geometry and memory]
\label{rm:thm:modular-area-memory}
Let
\begin{equation}
\label{rm:eq:modular-ND}
N=N_{90}+N_{120},
\qquad
D=90N_{90}+120N_{120}.
\end{equation}
Then
\begin{equation}
\label{rm:eq:modular-inverse-channels}
\boxed{
N_{90}=4N-\frac{D}{30},
\qquad
N_{120}=\frac{D}{30}-3N.}
\end{equation}
In particular,
\([\widehat{\mathcal A}_{\rm phys},\HHM^{(A)}]=0\), and the pair
\((\widehat{\mathcal A}_{\rm phys},\HHM^{(A)}-cI)\) reconstructs the two
occupation sectors.
\end{theorem}

\begin{proof}
The operators \(N_{90}\) and \(N_{120}\) act on distinct factors
and commute. The linear map carrying \((N_{90},N_{120})\) to \((N,D)\)
has matrix
\[
\begin{pmatrix}1&1\\90&120\end{pmatrix},
\qquad \det=30.
\]
Its inverse is exactly \cref{rm:eq:modular-inverse-channels}. Since
\cref{rm:eq:modular-area-operator,rm:eq:modular-hamiltonian} are linear
combinations of the same number operators, the commutator vanishes.
\end{proof}

The operator interpretation is precise: the area observable preserves the total number of excitations, while the modular Hamiltonian preserves their sector of origin. Neither observable alone retains both coordinates; the joint pair does determine them.

\section{Confluence of area, information and energy}

The angular scale \(\theta_{\rm clk}\) in
\cref{rm:eq:modular-theta-clock-a0} and the chronological temperature
\(\Theta_{\rm clk}\) are distinct objects; the use of an uppercase letter is part
of the typing. The thermodynamic realization in \cref{rm:ch:metrology} proves
\begin{equation}
\label{rm:eq:modular-planck-trit-bridge}
E_P=k_B\Theta_{\rm clk}\log3,
\end{equation}
whereas the areal realization constructs the elementary increment
\begin{equation}
\label{rm:eq:modular-area-quantum}
\delta\mathcal A
=\gammaH a_0\ell_P^2.
\end{equation}
They are not equated: \cref{rm:eq:modular-planck-trit-bridge} fixes the informational quantum of energy, and \cref{rm:eq:modular-area-quantum} the geometric quantum of area. Their quotient defines the typed tension of conversion
\begin{equation}
\label{rm:eq:modular-information-area-tension}
\boxed{
\mathcal T_{I/A}
=\frac{E_P}{\delta\mathcal A}
=\frac{k_B\Theta_{\rm clk}\log3}
{\gammaH a_0\ell_P^2}.}
\end{equation}
The expression introduces no additional constant: it determines
the normalizations that must be preserved in the transformation
of information into geometry.

\enlargethispage{7\baselineskip}
At the minimal cut of eighty-one links, the physical informational
capacity and the energy of one trit per link are
\begin{equation}
\label{rm:eq:modular-cut-info-energy}
S_{\rm cut}^{\rm phys}=81k_B\log3,
\qquad
E_{\rm cut}^{(1)}=81E_P
=\Theta_{\rm clk}S_{\rm cut}^{\rm phys}.
\end{equation}
The final equality is an equation of state of this realization, not a universal
equipartition law. In particular, the Gibbs state of the thirty-six
boundary links has a different entropy and is treated next.

This separation determines the structural content of the confluence:
Barbero--Immirzi
orients the area spectrum; \(\log3\) measures the irreducible decision of the
TRIT; \(\Theta_{\rm clk}\) converts that information into energy; and
\(\ell_P^2\) connects the area quantum to the gravitational family.
Gravity is described as a confluence of functional realizations of a
common antecedent, and not as an imposed equality among their values.
 
\paragraph{Transport of the channel memory to the modular Hamiltonian and to the double thermal field purification.} Channel memory is represented below through the modular Hamiltonian and a thermofield-double purification. \endgroup

\chapter{Boundary modular Hamiltonian, thermofield-double purification and oriented information}
\begingroup
\renewcommand{\chapter}[2][]{}%
\chapter{Modular boundary Hamiltonian, double thermal field purification, and oriented information deficit}
\label{chap:parte-iv-60}

The faithful qutrit state admits a double-thermofield purification. The modular Hamiltonian records the orientation of the channels, and the informational deficit measures the finite separation between states without being identified with area.

\section{Qutrit state, thermofield-double purification, and
informational deficit}

For a link with weight \(w>0\), let
\begin{equation}
\label{rm:eq:modular-link-state}
Z(w)=1+\ee^{-w}+\ee^{-2w},
\qquad
\rho(w)=\frac1{Z(w)}\sum_{n=0}^2\ee^{-wn}|n\rangle\langle n|.
\end{equation}
The weights of the two families are
\begin{align}
w_{90}&=90\Dmod
=0.0100077345053979592447310951874\ldots,
\label{rm:eq:modular-w90}\\
w_{120}&=120\Dmod
=0.0133436460071972789929747935832\ldots,
\qquad \frac{w_{120}}{w_{90}}=\frac43.
\label{rm:eq:modular-w120}
\end{align}
The finite boundary state is
\begin{equation}
\label{rm:eq:modular-product-state}
\rho_A=\rho(w_{90})^{\otimes18}
\otimes\rho(w_{120})^{\otimes18}.
\end{equation}
Its logaritmo reproduce \cref{rm:eq:modular-hamiltonian}, with
\(c=18\log Z(w_{90})+18\log Z(w_{120})\).

\begin{theorem}[Double-thermofield purification]
\label{rm:thm:modular-tfd}
The vector
\begin{equation}
\label{rm:eq:modular-tfd-link}
|\operatorname{TFD}(w)\rangle
=\frac1{\sqrt{Z(w)}}\sum_{n=0}^2
\ee^{-wn/2}|n\rangle_A\otimes|n\rangle_{\bar A}
\end{equation}
purify \(\rho(w)\). The product of eighteen copies for each of the
weights \(w_{90},w_{120}\) purifies \(\rho_A\); in the doubled system,
\(S(\rho_A)\) is exactly entanglement entropy.
\end{theorem}

\begin{proof}
Upon tracing out the second factor, the orthogonality of \(|n\rangle_{\bar A}\)
eliminates the cross terms and leaves
\[
\operatorname{Tr}_{\bar A}
|\operatorname{TFD}(w)\rangle\langle\operatorname{TFD}(w)|
=\rho(w).
\]
The partial trace commutes with tensor products, which proves the
assertion for the thirty-six links.
\end{proof}

For a link,
\begin{equation}
\label{rm:eq:modular-link-entropy}
S(w)=\log Z(w)
+w\frac{\ee^{-w}+2\ee^{-2w}}{Z(w)}.
\end{equation}
Consequently,
\begin{equation}
\label{rm:eq:modular-entropy-value}
S(\rho_A)=18S(w_{90})+18S(w_{120})
=39.54837320881807994957319730\ldots\ \text{nats}.
\end{equation}
The maximally mixed state of thirty-six qutrits has entropy
\(36\log3\). We define the oriented defect
\begin{equation}
\label{rm:eq:modular-oriented-information}
\boxed{
I_{\rm or}=36\log3-S(\rho_A)
=0.001669183233868940655631225913\ldots\ \text{nats}.}
\end{equation}
Its positivity follows from the entropic maximality of the
uniform state. The proximity of \(I_{\rm or}\) to \(\gammaH a_0\) does not constitute
an identity and is not used to select either value.
 
\section{Even information and odd geometric transport in the bilayer}

The normalized distribution between sheets separates even entropy from odd
hexad--octad transport. This bilayer parity preserves the boundary information that
the modular Hamiltonian records and avoids identifying entropy, area, and
orientation.

\begingroup
\let\chapter\section
\let\section\subsection
\let\subsection\subsubsection
\let\subsubsection\paragraph
\let\clearpage\relax
\let\cleardoublepage\relax
\chapter{The bicover: even information and odd geometric transport}
\label{chap:paridad-bicapa-barbero-rev6}
\index{entropy!bilayer}
\index{Barbero--Immirzi!parity}
\index{sheet!exchange}

Barbero--Immirzi and the entropy of the two sheets are not two juxtaposed
quantities. They arise from the same channels \(q_\pm\), but preserve
different data. Probabilistic normalization forgets the common scale and
measures the distribution between orientations; the hexad--octad functional preserves the
sign of that orientation. The first reading is even and the second is odd.

\section{Normalized distribution between leaves}

Write
\[
 x=A_{\rm rad},
 \qquad
 y=C^*_{\rm rad},
 \qquad
 q_-=e^{-x+y},
 \qquad
 q_+=e^{-x-y}.
\]
The sum \(q_-+q_+=2e^{-x}\cosh y\) allows defining
\[
 p_-=\frac{q_-}{q_-+q_+}
     =\frac{e^y}{2\cosh y},
 \qquad
 p_+=\frac{q_+}{q_-+q_+}
     =\frac{e^{-y}}{2\cosh y}.
\]
The common factor \(e^{-x}\) disappears. The distribution preserves the asymmetry
between sheets, but not the common angular scale.

\begin{definition}[Bilayer entropy]
\label{def:entropia-bicapa-rev6}
The entropy of the oriented distribution is
\[
 S_{\rm bi}(y)
 :=
 -p_-\log p_--p_+\log p_+
 =
 \log(2\cosh y)-y\tanh y.
\]
\end{definition}

\begin{proposition}[Parity and maximum of the bicover entropy]
\label{prop:paridad-entropia-bicapa-rev6}
The function \(S_{\rm bi}\) is even and satisfies
\[
 S_{\rm bi}(-y)=S_{\rm bi}(y),
 \qquad
 S_{\rm bi}'(y)=-y\,\operatorname{sech}^2y,
 \qquad
 0<S_{\rm bi}(y)\leq\log2.
\]
The maximum \(\log2\) is attained only at \(y=0\), when both sheets
have equal weight.
\end{proposition}

\begin{proof}
Parity follows because \(\cosh\) is even and \(y\tanh y\) is also even.
The direct derivation produces
\[
 \frac{\dd}{\dd y}\log(2\cosh y)=\tanh y,
 \qquad
 \frac{\dd}{\dd y}(y\tanh y)
 =\tanh y+y\operatorname{sech}^2y.
\]
Therefore \(S_{\rm bi}'(y)=-y\operatorname{sech}^2y\): the function increases on
\((-\infty,0)\), decreases on \((0,\infty)\), and takes the value
\(\log2\) at the origin. Positivity is that of the entropy of a distribution with two
strictly positive weights.
\end{proof}

\section{Parity of the hexad-octad transport}

Let
\[
 s_n(q_-,q_+)=q_-^n-q_+^n.
\]
The leaf exchange is the involution
\[
 \mathsf J_{\rm sh}:(x,y)\longmapsto(x,-y),
\]
which permutes \(q_-\leftrightarrow q_+\). Hence,
\[
 s_n\circ\mathsf J_{\rm sh}=-s_n.
\]
The functional constructed in \cref{mdg:chap:barbero} is
\[
 \Gamma_{6\to8}(A,C^*)
 =
 \frac{\pi A C^*}{180}
{}+12\bigl(S_{90}(q_-)-S_{90}(q_+)\bigr)
-\bigl(S_{120}(q_-)-S_{120}(q_+)\bigr).
\]

\begin{theorem}[Theorem of double-layer parity]
\label{thm:paridad-bicapa-barbero-rev6}
Under the exchange \(C^*\mapsto-C^*\),
\[
 \boxed{S_{\rm bi}(-C^*_{\rm rad})
        =S_{\rm bi}(C^*_{\rm rad}),}
\]
\[
 \boxed{\Gamma_{6\to8}(A,-C^*)
        =-\Gamma_{6\to8}(A,C^*).}
\]
The same bilayer therefore produces an even informational reading and an odd geometric transport.
\end{theorem}

\begin{proof}
The first equality is
\cref{prop:paridad-entropia-bicapa-rev6}. In the second, the term
\(\pi AC^*/180\) changes sign, and each Lambert difference changes
sign because its two arguments are interchanged. The complete combination is
odd.
\end{proof}

\section{Interpretation of the constitutive operator as coupling between vacuum, field, and response}

Bilayer entropy and \(\Gamma_{6\to8}\) are not two names for the same
number. The former quantifies how much weight the two orientations share
after cancellation of the common scale. The latter preserves the orientation and
incidence of the step \(6\to8\). At the HMT point,
\[
 S_{\rm bi}
 =0.69246069960358548\ldots\ {\rm nats},
 \qquad
 \Gamma_{6\to8}
 =0.27400767269445927\ldots.
\]
The link between information and geometry is that both readings
factor through the same channels \(q_\pm\), not that their values are equated.

The particle--antiparticle band of \cref{chap:banda-particula-virtual} reproduces the same parity distinction: mass and the amount of correlation descend to the band quotient, whereas charge, orientation and transport change sign. In the Cartan--Holst action, \(\Gamma_{6\to8}\) occupies the oriented coordinate; in the informational reading, \(S_{\rm bi}\) counts the distribution between sheets. This explains how a single structure can preserve information and reverse orientation without confusing the two operations. \endgroup

\paragraph{Tangential variation of the modular state and derivation of the areal law.} Tangent variation of the modular state produces the area first law; finite variations preserve the relative-entropy term. \endgroup

\chapter{Modular area first law and finite extension through relative entropy}
\begingroup
\renewcommand{\chapter}[2][]{}%
\chapter{1st modular law areal and finite extension via relative entropy}
\label{chap:parte-iv-61}

The first modular law relates variations of entropy and area within the declared finite system. Its nonlinear extension retains the positive relative-entropy defect and therefore does not conflate tangent equality with global identity.

\section{1st modular law for area variations}

Let us define the two normalized areas
\begin{equation}
\label{rm:eq:modular-normalized-area-channels}
\mathcal A_m=\gammaH a_0N_m,
\qquad m\in\{90,120\}.
\end{equation}
Then
\begin{equation}
\label{rm:eq:modular-beta-definition}
\HHM^{(A)}-cI
=\beta_{90}\mathcal A_{90}+\beta_{120}\mathcal A_{120},
\qquad
\beta_m=\frac{m\Dmod}{\gammaH a_0}.
\end{equation}
The terminal values are
\begin{align}
\beta_{90}&=5.9726485401070929914651798889\ldots,
\label{rm:eq:modular-beta90}\\
\beta_{120}&=7.9635313868094573219535731852\ldots,
\qquad
\frac{\beta_{120}}{\beta_{90}}=\frac43.
\label{rm:eq:modular-beta120}
\end{align}

\begin{theorem}[First modular law in areal variables]
\label{rm:thm:modular-first-law}
For a differentiable family of normalized states
\(\rho(\varepsilon)\) passing through \(\rho_A\), while keeping
\(\Dmod,\gammaH,a_0\) fixed,
\begin{equation}
\label{rm:eq:modular-first-law}
\boxed{
\delta S
=\beta_{90}\,\delta\langle\mathcal A_{90}\rangle
+\beta_{120}\,\delta\langle\mathcal A_{120}\rangle.}
\end{equation}
\end{theorem}

\begin{proof}
The first entanglement law in the reference state is \(\delta S=\delta\langle-\log\rho_A\rangle\). The
scalar term \(cI\) does not contribute because
\(\delta\Tr\rho=0\). Substituting \cref{rm:eq:modular-beta-definition} proves the identity.
\end{proof}

The formula is local in state space. For finite changes,
relative entropy appears; omitting it would turn a tangent identity into a
false claim of global linearity.

\Needspace{22\baselineskip}
\begin{theorem}[Finite Identity with a Defect Term]
\label{rm:thm:modular-finite-first-law}
For every state \(\sigma\) whose support is contained in that of
\(\rho_A\),
\begin{equation}
\label{rm:eq:modular-finite-first-law}
\boxed{
S(\sigma)-S(\rho_A)
=\beta_{90}\Delta_\sigma\langle\mathcal A_{90}\rangle
+\beta_{120}\Delta_\sigma\langle\mathcal A_{120}\rangle
-D(\sigma\Vert\rho_A),}
\end{equation}
where
\(D(\sigma\Vert\rho_A)=\Tr\sigma(\log\sigma-\log\rho_A)\geq0\).
In particular, the linear variation of
\cref{rm:thm:modular-first-law} is the tangent term of this identity.
\end{theorem}

\begin{proof}
By definition and by \(\HHM^{(A)}=-\log\rho_A\),
\[
D(\sigma\Vert\rho_A)
=-S(\sigma)+\Tr(\sigma\HHM^{(A)}).
\]
For \(\sigma=\rho_A\), the left-hand side is zero and
\(S(\rho_A)=\Tr(\rho_A\HHM^{(A)})\). Subtracting the two identities gives
\[
S(\sigma)-S(\rho_A)
=\Delta_\sigma\langle\HHM^{(A)}\rangle
-D(\sigma\Vert\rho_A).
\]
The scalar \(cI\) cancels between normalized states; substituting
\cref{rm:eq:modular-beta-definition} yields
\cref{rm:eq:modular-finite-first-law}. The positivity of relative entropy
is Klein's inequality.
\end{proof}

The correction \(D(\sigma\Vert\rho_A)\) measures the informational distance
between the actual state and the reference modular state. Thus, the first
law loses no content upon leaving the infinitesimal regime: it becomes
an exact identity with a controlled positive defect.
 
\paragraph{Inductive persistence of modular weights and factorial classification of the limit.} Inductive persistence of modular weights leads to infinite operator closure and its factorial classification. \endgroup

\chapter{Persistence of the incidences \texorpdfstring{\(90/120\)}{90/120} and hyperfinite factorial limit}
\begingroup
\renewcommand{\chapter}[2][]{}%
\chapter{Persistence of the channels \texorpdfstring{\(90/120\)}{90/120} and classification of the limit as the hyperfinite factor \texorpdfstring{\(III_{\lambda_\partial}\)}{III(partial lambda)}}
\label{chap:parte-iv-62}

Faithful local products preserve the ratios of channels \(90\) and \(120\) under refinement. The GNS representation and the ITPFI limit determine the hyperfinite type III factor with explicit boundary parameter.

\section{Local Limit and Modular Classification of the Infinite System}

In an infinite UHF algebra, there is in general no global trace-class
density matrix within the algebra. The correct object is the
net
\begin{equation}
\label{rm:eq:modular-local-net}
\HHM^{(m)}=-\log\rho_m
\end{equation}
with compatible-state restrictions, or the modular generator in the
GNS representation. The tracial expectation \(\mathrm{id}\otimes\tau_9\) does
not automatically preserve a nontracial Gibbs state and does not
replace state compatibility.

\begin{theorem}[Modular type under faithful persistence of both channels]
\label{rm:thm:modular-typeIII}
Suppose that:
\begin{enumerate}[label=(\roman*)]
\item each cut is a product of faithful qutrit states with weights
\(w_{90}\) and \(w_{120}\);
\item the inclusions preserve the state and the channel label;
\item infinitely many factors of each type occur, both with positive
asymptotic density;
\item no additional spectral quotients are introduced, and the
product representation is factorial.
\end{enumerate}
Then the additive group of logarithms of ratios of eigenvalues is
\begin{equation}
\label{rm:eq:modular-ratio-group}
90\Dmod\Z+120\Dmod\Z=30\Dmod\Z.
\end{equation}
The GNS closure ITPFI is the hyperinfinitesimal factor
\begin{equation}
\label{rm:eq:modular-typeIII-lambda}
\boxed{
III_{\lambda_\partial},
\qquad
\lambda_\partial=\ee^{-30\Dmod}
=0.9966696464689573786108299769\ldots .}
\end{equation}
\end{theorem}

\begin{proof}
At each link, the eigenvalue quotients are powers of \(\ee^{-w_m}\). Infinite persistence of both types generates the subgroup \(w_{90}\Z+w_{120}\Z\). Since \(\gcd(90,120)=30\), it follows that \cref{rm:eq:modular-ratio-group}. The Araki--Woods asymptotic ratio criterion for the faithful product yields the hyperfinite ITPFI factor with parameter \(\ee^{-30\Dmod}\) \cite{rm:ArakiWoods1968,rm:Takesaki}.
\end{proof}

The theorem states its actual hypotheses. If one of the channels disappears,
if the product is incompatible, or if the ratio group changes, the
type must be recalculated.
 
\paragraph{Gravitational publication of the modular closure via the pentacomponent carrier.} Modular closure prepares the gravitational expression: first the five-component carrier is constructed from icosahedral incidence. \endgroup

\chapter{Icosahedral descent and five-component gravitational carrier}
\begingroup
\renewcommand{\chapter}[2][]{}%
\chapter{Icosahedral descent \texorpdfstring{\(A_5/C_5\)}{A₅/C₅} and selection of the five-component gravitational carrier}
\label{chap:parte-iv-63}

The gravitational line preserves the constructed length and its chronological realization of \cref{g26:sec:rigidez-radial-es,cr28:sec:radial-clock}; its carrier then descends to the icosahedral quotient. The irreducible projector of rank \(5\) selects the tensor carrier without yet assuming a massive or radiative realization.

\section{Intrinsic gravitational constant and tensor representation}
\label{rm:ch:gravity}

\subsection{Constitutive chain of the gravitational sector: chronological scale \texorpdfstring{\(108\)}{108}, tensor reduction, holonomy and metrological evaluation of \texorpdfstring{\(G\)}{G}}

Gravity is organized here into four strata that must not be merged:

\begin{equation}
\label{rm:eq:gravity-causal-chain}
\begin{aligned}
\text{inscription and radial norm R1--R8}
&\longrightarrow L
\longrightarrow G=c_{\rm int}^3L^2/\hbar_a,\\
\text{icosahedral incidence}
&\longrightarrow V_{12}\xrightarrow{P_5}V_5
\xrightarrow{\Gamma_5}\operatorname{STF}_3
\xrightarrow{\Lambda(n)}\mathcal H_{\rm TT}(n),\\
\text{CT108 calendar with }t_0=\pi L/(54c_{\rm int})
&\longrightarrow R_{108}=L,\quad
\theta_{108}^{\rm circ}=(54/\pi)^2,\\
\text{metrological chart}
&\longrightarrow 6.6742996594140227\ldots\times10^{-11}\,
\mathrm{m^3\,kg^{-1}\,s^{-2}}.
\end{aligned}
\end{equation}

The first line preserves determination of length and coupling through rules R1--R8 of \cref{g26:sec:rigidez-radial-es}. The second constructs the carrier and its transverse traceless reduction. The third is the chronological realization of the already constructed length, and the fourth evaluates the same dimensional section. The reduction \(12\to5\to5\to2\) does not alter the common radial coefficient: intensity and representation preserve their provenance and respective functions.

\subsection{Gravitational magnitude line and chronological scale \texorpdfstring{\(108\)}{108}}

Let \(c_{\rm int}>0\), \(t_0>0\),
\(\ell_0=c_{\rm int}t_0\) and a positive section
\(\hbar_a\) of the action line, with
\(a\in\{\mathrm{pre},\mathrm{ret}\}\). The dimensional transducer is

\begin{equation}
\label{rm:eq:gravity-transducer}
G_a(\theta)
=\theta\frac{c_{\rm int}^{5}t_0^2}{\hbar_a}
=\theta\frac{c_{\rm int}^{3}\ell_0^2}{\hbar_a},
\qquad \theta>0.
\end{equation}

Equality of the two expressions uses only \(\ell_0=c_{\rm int}t_0\) and \([G]=L^3M^{-1}T^{-2}\). Dimension is fixed by the line; the dimensionless character \(\theta\) distinguishes the realizations of that system. In the longitudinal composition already proved, \(t_0=\pi L/(54c_{\rm int})\) and \(\theta=(54/\pi)^2\) reproduce exactly \(G_a=c_{\rm int}^3L^2/\hbar_a\).

The CT108 calendar defines

\begin{equation}
\label{rm:eq:gravity-clock108}
T_{108}=108t_0,
\qquad
\omega_{108}=\frac{2\pi}{108t_0},
\qquad
R_{108}=\frac{c_{\rm int}}{\omega_{108}}
=\frac{54}{\pi}\ell_0.
\end{equation}

For each action sheet, let

\begin{equation}
\label{rm:eq:gravity-clock-energy-mass}
E_{108,a}=\hbar_a\omega_{108},
\qquad
m_{108,a}=\frac{E_{108,a}}{c_{\rm int}^2}.
\end{equation}

Then, the reduced Compton wavelength coincides with the radius of the cycle:

\begin{equation}
\label{rm:eq:gravity-compton-clock}
\bar\lambda_{C,a}
=\frac{\hbar_a}{m_{108,a}c_{\rm int}}
=R_{108}.
\end{equation}

\begin{theorem}[Uniqueness of the selector under the HMT periodicity condition]
\label{rm:thm:gravity-circular-selector}
With the reduced convention
\(r_g=Gm/c_{\rm int}^2\) and the HMT periodicity premise that identifies the gravitational
radius of the CT108 quantum with its phase radius, the closure condition
\begin{equation}
\label{rm:eq:gravity-closing-condition}
r_{g,a}(\theta)=R_{108}
\end{equation}
selects a unique positive solution,
\begin{equation}
\label{rm:eq:gravity-theta108}
\boxed{
\theta_{108}^{\rm circ}=\left(\frac{54}{\pi}\right)^2
=295.4525715010569415303525\ldots .}
\end{equation}
In that embodiment,
\begin{equation}
\label{rm:eq:gravity-four-lengths}
R_{108}=\bar\lambda_{C,a}=r_{g,a}=\ell_P^{\rm circ}.
\end{equation}
\end{theorem}

\enlargethispage{7\baselineskip}
\begin{proof}
By \cref{rm:eq:gravity-transducer,rm:eq:gravity-clock-energy-mass},
\[
r_{g,a}(\theta)
=\frac{G_a(\theta)m_{108,a}}{c_{\rm int}^2}
=\theta c_{\rm int}t_0^2\omega_{108}
=\theta\frac{\pi}{54}\ell_0.
\]
Comparing it with
\(R_{108}=(54/\pi)\ell_0\) necessarily gives
\(\theta=(54/\pi)^2\). The equation is linear in \(\theta\) and all the
factors are positive, so the solution is unique.
\end{proof}

The equality \(r_g=R_{108}\) does not follow from the dimensionality of \cref{rm:eq:gravity-transducer}. It is the explicit geometric condition of this HMT realization: action, Compton, and periodic return must determine the same length. The theorem proves that, once this structural condition is accepted, the character \(\theta\) is fixed without using the SI value of \(G\).

\begin{remark}[Radius convention]
The theorem uses the reduced gravitational radius \(Gm/c^2\). If it is replaced
by the Schwarzschild radius \(2Gm/c^2\), the selector changes by a factor of
two. The two expressions correspond to distinct geometric conditions,
so the convention must be declared before the calculation.
\end{remark}

\subsection{Planck System and Bidirectional Covariance}

\begin{corollary}[Planck Family Associated with the CT108 Chronological Structure]
\label{rm:cor:gravity-planck-family}
For \(G_a=G_a(\theta_{108}^{\rm circ})\),
\begin{align}
\ell_P&=\frac{54}{\pi}\ell_0,
&t_P&=\frac{54}{\pi}t_0,
\label{rm:eq:gravity-planck-length-time}\\
E_{P,a}&=\frac{\hbar_a}{t_P}
=\frac{\pi\hbar_a}{54t_0}
=\frac{h_a}{108t_0},
&\ell_P^2&=\frac{G_a\hbar_a}{c_{\rm int}^3}
=\theta_{108}^{\rm circ}\ell_0^2.
\label{rm:eq:gravity-planck-energy-area}
\end{align}
Moreover,
\begin{equation}
\label{rm:eq:gravity-planck-force-power-density}
F_{P,a}=\frac{c_{\rm int}^4}{G_a},
\qquad
P_{P,a}=\frac{c_{\rm int}^5}{G_a},
\qquad
\rho_{P,a}=\frac{\hbar_a}
{(\theta_{108}^{\rm circ})^2c_{\rm int}^5t_0^4}.
\end{equation}
\end{corollary}

\begin{proof}
The first three identities follow from
\(\ell_P=R_{108}\), \(t_P=\ell_P/c_{\rm int}\), and
\(h_a=2\pi\hbar_a\). For the area,
\[
\frac{G_a\hbar_a}{c_{\rm int}^3}
=\theta_{108}^{\rm circ}c_{\rm int}^2t_0^2
=\theta_{108}^{\rm circ}\ell_0^2.
\]
The remaining expressions are the dimensional compositions of force,
power, and mass per volume with the same section.
\end{proof}

Let
\begin{equation}
\label{rm:eq:gravity-Ract}
R_{\rm act}:=\frac{\hbar_{\rm pre}}{\hbar_{\rm ret}}.
\end{equation}
Since the action appears in the numerator of the chronological mass and in the
denominator of \(G\),
\begin{equation}
\label{rm:eq:gravity-twoway-covariance}
\frac{m_{108,\rm pre}}{m_{108,\rm ret}}=R_{\rm act},
\qquad
\frac{G_{\rm pre}}{G_{\rm ret}}=R_{\rm act}^{-1},
\qquad
G_a\hbar_a
=\theta_{108}^{\rm circ}c_{\rm int}^5t_0^2.
\end{equation}
Therefore, \(\ell_P^2=G_a\hbar_a/c^3\) does not change sheets. Gravity and
action are reciprocally oriented, while area preserves the
common geometry. This cancellation is the interface of confluence toward the area
operator of the \cref{rm:ch:modular}.

\subsection{Icosahedral Component of Dimension \(5\)}

The antecedent is neither \(\R\) nor a subsequent selection of five
polarizations. The object responsible for transporting the incidence is defined
explicitly before introducing the exceptional realization.

Let \(X_{\rm TPK}^{\rm exc}\) be the subspace of the enriched state \(\TPK\)
that preserves, in addition to the arithmetic output, the orientation, route,
carry, and exceptional incidence. To make the geometric target space explicit,
let \(\varphi_g=(1+\sqrt5)/2\), \(\nu=\sqrt{1+\varphi_g^2}\), and define
\begin{equation}
\label{rm:eq:gravity-icosahedron-vertices}
\Omega_{12}^{\rm ico}
=\frac1\nu\bigl(
\{0\}\times\{\pm1\}\times\{\pm\varphi_g\}
\;\cup\;
\{\pm1\}\times\{\pm\varphi_g\}\times\{0\}
\;\cup\;
\{\pm\varphi_g\}\times\{0\}\times\{\pm1\}
\bigr)\subset S^2.
\end{equation}
Its twelve elements are the oriented vertex classes of an icosahedron.
The incidence is not left verbal: it is fixed by
\begin{equation}
\label{rm:eq:gravity-icosahedron-incidence}
\mathcal I_{12}^{\rm ico}
=\{(v,w)\in\Omega_{12}^{\rm ico}\times\Omega_{12}^{\rm ico}:
v\neq w,\ \langle v,w\rangle=1/\sqrt5\}.
\end{equation}
Each vertex has five neighbors in this relation. The reading bridge
is typed as a map
\begin{equation}
\label{rm:eq:gravity-exceptional-reader-map}
\pi_{\rm ico}:X_{\rm TPK}^{\rm exc}\longrightarrow\Omega_{12}^{\rm ico}.
\end{equation}
The cardinality twelve of the image of \(\pi_{\rm ico}\) is not sufficient.
To define an incidence realization, the mapping must satisfy
simultaneously:
\begin{enumerate}[label=(\roman*)]
\item being surjective and constant exactly on the fibers that
discard the arithmetic chart but preserve the exceptional route;
\item intertwine the action of the oriented automorphisms of the antecedent
with a transitive action on \(\Omega_{12}^{\rm ico}\);
\item transporting the declared TPK incidence to
\(\mathcal I_{12}^{\rm ico}\), not
merely its cardinalities;
\item preserve the leaf involution and the nonadic return in the
stabilizer of the corresponding fiber.
\end{enumerate}

These conditions separate three assertions: the existence of the map
\(\pi_{\rm ico}\), the identification of the group that acts, and the
linear representation constructed thereafter. The second and third
cannot be used to define the first retrospectively.

\begin{theorem}[Incidence Descent to the Icosahedral Quotient]
\label{rm:thm:gravity-exceptional-descent}
Suppose that \(\pi_{\rm ico}\) satisfies (i)--(iv) and that the oriented action
induced on \(\Omega_{12}^{\rm ico}\) has group
\(G_{\rm ico}\simeq A_5\). Then, for every class
\(v_0\in\Omega_{12}^{\rm ico}\),
\begin{equation}
\label{rm:eq:gravity-icosahedral-quotient}
H_{v_0}=\operatorname{Stab}_{G_{\rm ico}}(v_0)\simeq C_5,
\qquad
\Omega_{12}^{\rm ico}\simeq G_{\rm ico}/H_{v_0}\simeq A_5/C_5.
\end{equation}
The application
\begin{equation}
\label{rm:eq:gravity-coset-map}
\kappa_{v_0}:G_{\rm ico}/H_{v_0}\longrightarrow\Omega_{12}^{\rm ico},
\qquad gH_{v_0}\longmapsto g\cdot v_0,
\end{equation}
is an equivariant bijection and preserves incidence if, in the quotient,
one declares
\begin{equation}
\label{rm:eq:gravity-coset-incidence}
gH_{v_0}\sim g'H_{v_0}
\quad\Longleftrightarrow\quad
(g\cdot v_0,g'\cdot v_0)\in\mathcal I_{12}^{\rm ico}.
\end{equation}
Therefore, the composite \(\kappa_{v_0}^{-1}\pi_{\rm ico}\) is the typed map from the TPK
antecedent to the quotient \(A_5/C_5\).
\end{theorem}

\begin{proof}
The list of \cref{rm:eq:gravity-icosahedron-vertices} has twelve elements. By transitivity, \(|G_{\rm ico}\cdot v_0|=12\), and the orbit--stabilizer theorem gives \(|H_{v_0}|=60/12=5\). Every group of prime order five is cyclic; hence \(H_{v_0}\simeq C_5\). The formula for \(\kappa_{v_0}\) does not depend on the representative: if \(gH_{v_0}=g'H_{v_0}\), then \(g^{-1}g'\in H_{v_0}\) and \(g'v_0=gv_0\). Surjectivity follows from transitivity, and equality of the cardinalities gives injectivity. Equivariance is immediate: \(\kappa_{v_0}(agH_{v_0})=a\kappa_{v_0}(gH_{v_0})\). Finally, \cref{rm:eq:gravity-coset-incidence} is precisely the transport of structure through that bijection; hypothesis (iii) makes this transport commute with \(\pi_{\rm ico}\).
\end{proof}

The theorem does not replace the construction of \(\pi_{\rm ico}\): it proves
that it is obtained once its four requirements have been verified. In the complete
HMT profile, the incidence corridor reaching
Paley--Witt--Golay--Leech--\(V^\natural\)--Monster--Moonshine is the
domain of definition of that map. In this autonomous volume, the exact input is
the map \cref{rm:eq:gravity-exceptional-reader-map} with its conditions; the entire
segment from \(A_5/C_5\) to TT reduction is proved below
without again consulting a gravitational value.

The chain of provenance is expressed through all its morphisms:
\begin{equation}
\label{rm:eq:gravity-exceptional-provenance}
\APP\longrightarrow\TRIT\longrightarrow\TPK
\longrightarrow X_{\rm TPK}^{\rm exc}
\xrightarrow{\ \pi_{\rm ico}\ }\Omega_{12}^{\rm ico}
\xleftarrow[\simeq]{\ \kappa_{v_0}\ }A_5/C_5.
\end{equation}
The complete exceptional corridor is proved in its own dedicated development. The hypothesis that this chapter cannot replace with a count is the existence of \(\pi_{\rm ico}\) satisfying (i)--(iv). With that input fixed, descent to the quotient is \cref{rm:thm:gravity-exceptional-descent}, and \(P_5\), \(\Gamma_5\) and \(\Lambda(n)\) are constructed from it. Realizing the resulting carrier as a gravitational field remains a typed physical interface, not a consequence of cardinality.

Let \(G=A_5\), let \(H\simeq C_5\) be the stabilizer of an oriented vertex
of the icosahedron, and consider
\begin{equation}
\label{rm:eq:gravity-v12}
V_{12}=\R[G/H],
\qquad |G/H|=\frac{60}{5}=12.
\end{equation}
The permutation representation \(\rho_{12}\) decomposes as
\begin{equation}
\label{rm:eq:gravity-v12-decomposition}
V_{12}\simeq\mathbf1\oplus V_3\oplus V_{3'}\oplus V_5.
\end{equation}
The central idempotent of the character summand
\(\chi_5=(5,1,-1,0,0)\) is
\begin{equation}
\label{rm:eq:gravity-p5}
\boxed{
P_5=\frac5{60}\sum_{g\in A_5}\chi_5(g^{-1})\rho_{12}(g)
=\frac1{12}\left(
5I+\sum_{g\in2A}\rho_{12}(g)
-\sum_{g\in3A}\rho_{12}(g)
\right).}
\end{equation}

\begin{theorem}[Rank-five gravitational projector and coupler]
\label{rm:thm:gravity-p5-coupling}
The operator \(P_5\) satisfies
\begin{equation}
\label{rm:eq:gravity-p5-properties}
P_5^2=P_5=P_5^*,
\qquad \Tr P_5=\rank P_5=5.
\end{equation}
For
\begin{equation}
\label{rm:eq:gravity-G5}
\mathbb G_{5,a}:=G_a(\theta_{108}^{\rm circ})P_5
\end{equation}
we have
\begin{equation}
\label{rm:eq:gravity-G5-spectrum}
\Spec(\mathbb G_{5,a})
=\{G_a(\theta_{108}^{\rm circ})\text{ with multiplicity }5,
\;0\text{ with multiplicity }7\}.
\end{equation}
\end{theorem}

\begin{proof}
Character orthogonality converts \cref{rm:eq:gravity-p5} into the
central idempotent that acts as the identity on the unique copy of
\(V_5\) and as zero on the other three summands of
\cref{rm:eq:gravity-v12-decomposition}. Reality of the character and
orthogonality of \(\rho_{12}\) give self-adjointness. Its trace and rank
are five. Multiplying a rank-five projector by the positive scalar
\(G_a\) immediately produces the indicated spectrum.
\end{proof}

\begingroup\fontsize{10.2}{11.7}\selectfont \setlength{\parskip}{2pt}\setlength{\abovedisplayskip}{5pt}\setlength{\belowdisplayskip}{5pt} The expression \(\mathbb G_{5,a}\) does not add five different values of the gravitational constant. It is the same section of \(G\) acting on the selected tensor summand and annihilating the other seven modes of the permutation module.
 
\paragraph{Equivariant reduction of the carrier to the traceless tensor sector.} The carrier intertwines with the symmetric traceless tensors and is reduced equivariantly before transversality is imposed.

\par\endgroup
 \endgroup

\chapter{Tensor intertwiner and transverse reduction \texorpdfstring{\(12\to5\to5\to2\)}{12->5->5->2}}
\begingroup
\renewcommand{\chapter}[2][]{}%
\chapter{Intertwiner with symmetric traceless tensors and transverse reduction \texorpdfstring{\(12\to5\to5\to2\)}{12→5→5→2}}
\label{chap:parte-iv-64}

The icosahedral framework constructs an equivariant isometry into the space of traceless symmetric tensors. The transverse projection realizes \(12\to5\to5\to2\), preserves the five-component carrier, and localizes the radiative subspace of dimension \(2\).

\section{Intertwiner with the space of symmetric traceless tensors}

Let \(\mathcal V\subset S^2\) be the set of the twelve unit vertices
of the icosahedron. For each \(v\in\mathcal V\), define
\begin{equation}
\label{rm:eq:gravity-Qv}
Q_v=vv^T-\frac13I_3\in
\operatorname{STF}_3:=\operatorname{Sym}^2_0(\R^3).
\end{equation}
The frame map
\begin{equation}
\label{rm:eq:gravity-Gamma0}
\Gamma_0\!\left(\sum_{v\in\mathcal V}x_ve_v\right)
=\sum_{v\in\mathcal V}x_vQ_v
\end{equation}
is \(A_5\)-equivariant and satisfies
\begin{equation}
\label{rm:eq:gravity-tight-frame}
\Gamma_0\Gamma_0^*=\frac85I_{\operatorname{STF}_3},
\qquad
\Gamma_0=\Gamma_0P_5.
\end{equation}

\begin{theorem}[Normalized tensor entangler]
\label{rm:thm:gravity-gamma5}
The map
\begin{equation}
\label{rm:eq:gravity-gamma5}
\boxed{
\Gamma_5:=\sqrt{\frac58}\,\Gamma_0|_{V_5}:
V_5\xrightarrow{\ \simeq\ }\operatorname{STF}_3}
\end{equation}
is a \(A_5\)-equivariant isometry. Its adjoint is
\begin{equation}
\label{rm:eq:gravity-gamma5-adjoint}
\Gamma_5^*Q=\sqrt{\frac58}
\sum_{v\in\mathcal V}(v^TQv)e_v.
\end{equation}
\end{theorem}

\begin{proof}
Equivariance follows from \(Q_{gv}=gQ_vg^T\). The operator
\(\Gamma_0\Gamma_0^*\) commutes with the irreducible representation of
\(A_5\) on \(\operatorname{STF}_3\), and is therefore scalar. Its trace
is
\[
\sum_{v\in\mathcal V}\|Q_v\|_F^2
=12\cdot\frac23=8;
\]
dividing by dimension five yields \(8/5\). Hence,
\(\Gamma_5\Gamma_5^*=I\). Since domain and codomain have dimension
five, also \(\Gamma_5^*\Gamma_5=I\). The formula for the adjoint
follows from \(\langle Q,Q_v\rangle_F=v^TQv\).
\end{proof}

Normalization through the icosahedral frame
eliminates a purely cardinal coincidence: it constructs the map that
transports the finite summand to the continuous spin-two
carrier.

\section{Transverse traceless projection}

For \(n\in S^2\), let \(\Pi(n)=I-nn^T\). We define
\begin{equation}
\label{rm:eq:gravity-lambda-tt}
\boxed{
\Lambda(n)Q
=\Pi Q\Pi-\frac12\Tr(\Pi Q\Pi)\Pi.}
\end{equation}

\begin{theorem}[Reduction from five components to two]
\label{rm:thm:gravity-tt}
The restriction of \(\Lambda(n)\) to \(\operatorname{STF}_3\) is the
orthogonal projector onto
\begin{equation}
\label{rm:eq:gravity-htt}
\mathcal H_{\rm TT}(n)
=\{Q\in\operatorname{STF}_3:Qn=0\}.
\end{equation}
In particular,
\begin{equation}
\label{rm:eq:gravity-lambda-properties}
\Lambda(n)^2=\Lambda(n)=\Lambda(n)^*,
\qquad \rank\Lambda(n)=2.
\end{equation}
For \(n=e_3\), the image is generated by
\begin{equation}
\label{rm:eq:gravity-plus-cross}
E_+=\frac1{\sqrt2}\operatorname{diag}(1,-1,0),
\qquad
E_\times=\frac1{\sqrt2}
\begin{pmatrix}0&1&0\\1&0&0\\0&0&0\end{pmatrix}.
\end{equation}
\end{theorem}

\begin{proof}
Since \(\Pi n=0\), the output is transverse. Its trace is
\[
\Tr(\Pi Q\Pi)-\frac12\Tr(\Pi)\Tr(\Pi Q\Pi)=0,
\]
since \(\Tr\Pi=2\). If \(Qn=0\) and \(\Tr Q=0\), then
\(\Pi Q\Pi=Q\) and \(\Tr(\Pi Q\Pi)=0\), so that
\(\Lambda(n)Q=Q\). The formula is self-adjoint with respect to the
Frobenius product. In the frame \(n=e_3\), its matrix in an adapted STF basis is
\(\operatorname{diag}(1,1,0,0,0)\), which proves the rank.
\end{proof}

The complete chain thus remains,
\begin{equation}
\label{rm:eq:gravity-12-5-5-2}
\boxed{
V_{12}\xrightarrow{P_5}V_5
\xrightarrow{\Gamma_5}\operatorname{STF}_3
\xrightarrow{\Lambda(n)}\mathcal H_{\rm TT}(n),
\qquad 12\longrightarrow5\longrightarrow5\longrightarrow2.}
\end{equation}
The reduced coupling
\begin{equation}
\label{rm:eq:gravity-effective-map}
\mathbb G_{{\rm TT},a}(n)
=G_a(\theta_{108}^{\rm circ})\Lambda(n)\Gamma_5P_5
\end{equation}
has rank two.
 
\paragraph{Helical decomposition of the spectrum and separation of massive and non-massive modes.} The reduction distinguishes the complete helicity spectrum of a massive realization from its two massless helicities. \endgroup

\chapter{Helicity decomposition of the spin-two carrier}
\begingroup
\renewcommand{\chapter}[2][]{}%
\chapter{Helicity Decomposition of the Spin-\texorpdfstring{\(2\)}{2} Carrier and Distinction between Massive and Massless Realizations}
\label{chap:parte-iv-65}

The real carrier of dimension \(5\) decomposes into \(5\)
helicity weights. The massive realization retains all
\(5\) components; the massless transverse reduction retains only the
radiative helicities \(\pm2\), without converting a collective
possibility into a postulated fundamental particle.

\section{Decomposition into \(5\) polarization modes}

Dimension five acquires tensorial content upon choosing a unit direction
of propagation \(n\) and an oriented orthonormal frame
\((e_1,e_2,n)\). Define the five real tensors
\begin{align}
E_{+}&=\frac1{\sqrt2}(e_1e_1^T-e_2e_2^T),
&E_{\times}&=\frac1{\sqrt2}(e_1e_2^T+e_2e_1^T),
\label{rm:eq:gravity-helicity2-real}\\
E_{1}&=\frac1{\sqrt2}(e_1n^T+ne_1^T),
&E_{2}&=\frac1{\sqrt2}(e_2n^T+ne_2^T),
\label{rm:eq:gravity-helicity1-real}\\
E_{0}&=\frac1{\sqrt6}(2nn^T-e_1e_1^T-e_2e_2^T).
\label{rm:eq:gravity-helicity0-real}
\end{align}
All are symmetric and traceless. The first two are transverse; the
next two mix the transverse plane with the direction of
propagation, and the last is the traceless longitudinal scalar mode.

\begin{theorem}[Complete helical decomposition of the carrier]
\label{rm:thm:gravity-five-helicities}
The tensors of
\cref{rm:eq:gravity-helicity2-real,rm:eq:gravity-helicity1-real,rm:eq:gravity-helicity0-real}
form an orthonormal basis of \(\operatorname{STF}_3\). In the
complexification, the combinations
\begin{equation}
\label{rm:eq:gravity-complex-helicities}
E_{\pm2}=\frac1{\sqrt2}(E_+\mp\ii E_\times),
\qquad
E_{\pm1}=\frac1{\sqrt2}(E_1\mp\ii E_2),
\qquad E_0=E_0
\end{equation}
are weight-\(\pm2,\pm1,0\) vectors, respectively, under rotations
about \(n\). Moreover,
\begin{equation}
\label{rm:eq:gravity-tt-on-five}
\Lambda(n)E_{\pm2}=E_{\pm2},
\qquad
\Lambda(n)E_{\pm1}=0,
\qquad
\Lambda(n)E_0=0.
\end{equation}
\end{theorem}

\begin{proof}
In the frame \((e_1,e_2,n)\), each tensor has Frobenius norm one.
Their off-diagonal entry patterns are disjoint, and the three diagonals
of \(E_+,E_0\) have zero inner product; therefore, the five tensors
are orthonormal. Since \(\dim\operatorname{STF}_3=6-1=5\), they form a basis.

A rotation through angle \(\psi\) about \(n\) sends
\((e_1,e_2)\) to
\((\cos\psi\,e_1+\sin\psi\,e_2,
-\sin\psi\,e_1+\cos\psi\,e_2)\). Substituting into the formulas gives
a rotation through angle \(2\psi\) on
\(\operatorname{span}\{E_+,E_\times\}\), one through angle \(\psi\) on
\(\operatorname{span}\{E_1,E_2\}\), and \(E_0\) fixed. Hence the weights of
\cref{rm:eq:gravity-complex-helicities}. Finally,
\(\Pi E_{+}\Pi=E_+\) and \(\Pi E_\times\Pi=E_\times\), whereas
\(\Pi E_1\Pi=\Pi E_2\Pi=0\). For \(E_0\),
\(\Pi E_0\Pi=-(e_1e_1^T+e_2e_2^T)/\sqrt6\), which is pure trace in
the plane, and subtracting \cref{rm:eq:gravity-lambda-tt} annihilates it.
\end{proof}

The theorem specifies the meaning of ``five polarizations'' without assigning all of them the same dynamic status. Before field equations are imposed, the spin-two carrier has five irreducible coordinates. In a massive Fierz--Pauli realization, the constraints leave precisely the five weights of \cref{rm:eq:gravity-complex-helicities}; in a massless realization with gauge invariance, the sectors \(\pm1\) and \(0\) are eliminated, leaving \(\pm2\). Here HMT supplies the carrier, the basis, and the projector; the choice of massive or massless dynamics belongs to the physical interface. Thus five does not mean ``five observed waves,'' but rather five tensorial degrees available before reduction.

The orientation also remains visible: replacing
\((e_1,e_2,n)\) with \((e_1,-e_2,-n)\) preserves \(E_+,E_0\) and exchanges
the phases of the helical pairs. That information is not in the
scalar \(G_a\); it resides in the carrier on which
\(\mathbb G_{5,a}\) acts.

\section{Representational chain \(12\to5\to5\to2\): carrier \(V_5\), trace-free tensor, Fierz–Pauli polarizations, and transverse helicities}

Dimension five appears in distinct objects, and only
\cref{rm:eq:gravity-12-5-5-2} authorizes transport between them:

\begin{enumerate}
\item \(V_5\) is the irreducible summand of the representation on the
twelve vertices. This is a finite result in representation theory.
\item \(\operatorname{STF}_3\) is the real five-dimensional carrier of
the \(m=-2,-1,0,1,2\) weights of spin two.
\item A massive Fierz--Pauli theory can realize these five weights
as five physical states, after fixing the action and constraints.
\item Massless general relativity retains two radiative helicities,
\(\pm2\), represented by \(E_+\) and \(E_\times\).
\item The \(E(2)\) classes of metric waves, such as \(III_5\), classify
possible amplitudes under other hypotheses; they are not automatically
\(V_5\), Fierz--Pauli, or general relativity.
\end{enumerate}

In particular, the correct statement is: \emph{HMT constructs a five-component
tensor carrier before the restrictions; the TT reduction of the massless
realization has rank two}. It is not asserted that
general relativity has five propagating polarizations.
 
\paragraph{Existence conditions of the mode through natural holonomic selection.} Physical existence of the mode still requires a natural, falsifiable holonomic selector. \endgroup

\chapter{Natural gravitational selector and dimensional transducer}
\section[Structural Criteria and Metrological Chart of \(G\)]{Criteria for selecting the gravitational coordinate: enriched holonomy, spectral
response, and metrological chart}
\label{sec:selector-gravitatorio-rev11}
\index{gravitational constant!transducer and selector}
\index{gravity!metric coordinate}
\index{holonomy!gravitational selector}

The Palatini--Cartan--Holst action and its constrained Hamiltonian first determine
the dynamics. Upon it, the gravitational constant is organized into three typed
levels: the dimensional line that contains \(G\), the transducer that converts
a dimensionless coordinate into a section of that line, and the structural
selector of that coordinate. The decimal chart and experimental uncertainty
are presented at the end as representation and comparison.

\subsection{Dimensional line and internal transducer}

In the dimensional lattice of Dimensional Mechanics,
\[
 \mathbf G=-\mathbf S+5\mathbf C+2\mathbf T,
 \qquad [G]=L^3M^{-1}T^{-2}.
\]
Let
\[
 c_{\rm int}\in\mathcal L_C^+,
 \qquad t_0\in\mathcal L_T^+,
 \qquad \ell_0=c_{\rm int}t_0,
 \qquad \hbar_{\rm int}\in\mathcal L_S^+.
\]

\HMTClaim{MD-R-G-TRANSDUCER}
\begin{definition}[Gravitational transducer]
\label{def:transductor-gravitatorio-rev11}
For \(\theta_G>0\), define
\begin{equation}
 \boxed{
 G_{\rm int}(\theta_G)
 =\theta_G\frac{c_{\rm int}^{5}t_0^2}{\hbar_{\rm int}}
 =\theta_G\frac{c_{\rm int}^{3}\ell_0^2}{\hbar_{\rm int}}
 \in\mathcal L_G^+.}
 \label{eq:transductor-gravitatorio-rev11}
\end{equation}
\end{definition}

The equality of the two expressions and the type of
\eqref{eq:transductor-gravitatorio-rev11} are exact. If
\[
 L_G^2:=\frac{\hbar_{\rm int}G_{\rm int}}{c_{\rm int}^3},
\]
then
\begin{equation}
 \theta_G=\left(\frac{L_G}{\ell_0}\right)^2.
 \label{eq:thetaG-razon-longitudes-rev11}
\end{equation}
The gravitational coordinate is therefore a quadratic ratio of lengths
or a normalized holonomy response. The decade
\(10^{-34}\mathcal U_S\) already belongs to \(\hbar_{\rm int}\): it fixes the absolute
action section and does not constitute a relative factor added to \(G\).

\subsection{Observable torsion and positive characters}

The observable return of \(27\) steps
\[
 \mathcal K_{27}:=F_{\rm obs}^{27}
\]
satisfies
\[
 \mathcal K_{27}^{4}=F_{\rm obs}^{108}=I.
\]
In coordinate \(\theta\in\mathbb Z/108\mathbb Z\), the order is
exactly four, since
\(\mathcal K_{27}(\theta)=\theta+27\),
\(\mathcal K_{27}^{2}(\theta)=\theta+54\neq\theta\), and
\(\mathcal K_{27}^{4}(\theta)=\theta\).

\HMTClaim{MD-R-G-POSITIVE-CHARACTER-NOGO}
\begin{theorem}[Triviality of the Observable Positive Character]
\label{thm:caracter-positivo-G-rev11}
Every multiplicative homomorphism
\[
 \chi:\langle\mathcal K_{27}\rangle
 \longrightarrow\mathbb R_{>0}^{\times}
\]
cumple \(\chi(\mathcal K_{27})=1\).
\end{theorem}

\begin{proof}
Multiplicativity implies
\(\chi(\mathcal K_{27})^4=\chi(I)=1\).
The only positive real number whose fourth power is one is \(1\).
\end{proof}

The same cycle does carry the unitary characters.
\[
 \chi_j(\mathcal K_{27})=\exp(2\pi ij/4),
 \qquad j=0,1,2,3.
\]
The observable quotient admits phase, whereas its torsion eliminates a
nontrivial positive multiplicative scale. A multiplicative continuation must therefore
act on a lift \(\widetilde{\mathcal K}_{27}\) whose class in
the abelianization of the isotropy group is nontorsional, and must naturally select
both the character and its normalization.

\subsection{Chronological memory and spectral response}
\label{sec:respuesta-espectral-G-rev11}

The directed chronological covering realizes the return of \(27\) steps as the
translation
\[
 (G_{\rm or},S)\longmapsto(G_{\rm or},S)+(1,18),
\]
of infinite semigroup order. This route preserves accumulation, but its
spectral response requires interference among histories.

Let \(\Gamma=(V,E)\) be a finite directed multigraph with positive measure,
weights \(a_\ast(e)\geq0\), and cocycle \(c(e)\in\mathbb Z^2\). The twisted
operator
\[
 (\mathcal L_{\beta,\vartheta}f)(y)
 =\sum_{e:x\to y}
 e^{-\beta a_\ast(e)}e^{i\langle\vartheta,c(e)\rangle}f(x)
\]
induce the operator positive
\(\mathcal H_{\beta,\vartheta}
=\mathcal L_{\beta,\vartheta}^{*}\mathcal L_{\beta,\vartheta}\).

\HMTClaim{MD-R-G-PHASE-INTERFERENCE}
\begin{proposition}[Interference criterion for the phase response]
\label{prop:interferencia-selector-G-rev11}
In a deterministic orbit with a unique incoming edge at each vertex,
\(\mathcal H_{\beta,\vartheta}\) is independent of \(\vartheta\). Every nonzero-phase spectral
response requires multiple coherent histories or a magnetic differential
whose holonomy survives in the enriched cycles.
\end{proposition}

\begin{proof}
On the deterministic orbit, the sole phase of each term is multiplied by its
conjugate in
\(\mathcal L_{\beta,\vartheta}^{*}\mathcal L_{\beta,\vartheta}\)
and cancels. The absence of cross terms removes every dependence on
\(\vartheta\).
\end{proof}

If a simple band \(\lambda_0(\vartheta)\) has a minimum at the origin,
its Hessian \(Q\) is positive semidefinite. With
\[
 v_{27}=(1,18),
 \qquad
 \mathcal R_{27}=v_{27}^{\mathsf T}Qv_{27},
\]
the condition \(\mathcal R_{27}\neq0\) constitutes the minimal control of
sensitivity to accumulated memory. A natural normalization
\(\mathfrak N\) would allow one to define
\[
 \Lambda_G=\ell_0\mathfrak N(Q,v_{27}),
 \qquad
 G_{\rm spec}
 =\frac{c_{\rm int}^{3}}{\hbar_{\rm int}}
  \bigl(\delta_\theta\Lambda_G\bigr)^2,
 \qquad
 \delta_\theta=\frac{A_{\rm deg}-C_{\rm deg}^{\ast}}{360}.
\]
The homogeneity of this template is exact. Its physical realization is
characterized by four positive conditions: an enriched multigraph,
a nonzero response, stability under refinement, and independence under rescaling,
together with a normalization \(\mathfrak N\) fixed before the
metrological contrast.

\subsection{Geometric Functional on Enriched Holonomy}

The second continuation defines directly a length functional.
\[
 \mathscr L:
 \operatorname{Loop}(\mathcal X_{\rm TPK}^{\rm enr})
 \longrightarrow\mathcal L_L^+\cup\{0\},
\]
invariant under cyclic change of basepoint, conjugation, admissible
relabeling, and subdivision. For an enriched lift, write
\[
 \Lambda_{27}:=\mathscr L(\widetilde{\mathcal K}_{27}),
 \qquad
 \rho_{27}:=\frac{\Lambda_{27}}{c_{\rm int}t_0},
 \qquad
 \theta_G^{\rm hol}:=(\delta_\theta\rho_{27})^2.
\]
Transduction produces the homogeneous template
\begin{equation}
 \boxed{
 G_{\rm hol}
 =\frac{c_{\rm int}^{3}}{\hbar_{\rm int}}
  \bigl(\delta_\theta\Lambda_{27}\bigr)^2.}
 \label{eq:plantilla-holonomia-G-rev11}
\end{equation}
The concatenation law of this functional on the torsional cycle forms
part of its definition and its proof. Physical selection is completed
by constructing \(\widetilde{\mathcal K}_{27}\) and \(\mathscr L\), fixing a
unique normalization, and verifying the stability of the output. The value
\(G_{\rm CODATA\,2022}\) enters exclusively afterward, in the metrological
comparison.

These conditions describe this holonomy template. The longitudinal and radial realization already specified in \cref{g26:sec:rigidez-radial-es} uses the marked record, its conditional expectation and the positive norm; its normalization and the common value of \(G\) are proved there. \(\delta_\theta\Lambda_{27}\) is not identified with \(L\) by formal resemblance: the alternative template retains the obligation to construct its functional before making that identification.

\begin{proposition}[Separation of the two gravitational selectors]
\label{prop:separacion-selectores-G}
The circular selector
\(\theta_{108}^{\rm circ}=(54/\pi)^2\), constructed by action--Compton--gravity
closure, and the holonomic selector
\(\theta_G^{\rm hol}=(\delta_\theta\rho_{27})^2\) are two distinct typed realizations
of the transducer \(G_{\rm int}\).  They are not identified with
one another, and the decimal chart \(\mathscr R_G\) does not constitute an evaluation of
\(G_{\rm hol}\).
\end{proposition}
\begin{proof}
The first selector depends solely on the circular closure of order 108. The second additionally contains the spectral length \(\Lambda_{27}=\mathscr L(\widetilde{\mathcal K}_{27})\). An equality between them would require computing \(\Lambda_{27}\) and proving the agreement of normalizations; no operation of the decimal reader performs that computation. Their domains and input data therefore distinguish them before any numerical comparison.
\end{proof}

\subsection{Structural selection criterion}

\HMTClaim{MD-R-G-SELECTOR-CRITERION}
\begin{criterion}[Structural selector of the gravitational constant]
\label{crit:selector-G-rev11}
An HMT selector of \(G\) is identified when it jointly satisfies:
\begin{enumerate}
 \item enriched domain, composition, and orientation declared;
 \item descent to a nontorsional Abelian class, or a geometric functional
 invariant under basepoint, conjugation, relabelling, and subdivision;
 \item nontrivial positive response stable under refinement;
 \item unique normalization independent of the operator's overall scale;
 \item internal selection of the coordinate and its decimal position;
 \item freezing of the code and entries before revealing the
 comparison value;
 \item final realization as a section of \(\mathcal L_G\) in a common
 dimensional chart.
\end{enumerate}
\end{criterion}

The dimensionless ratios
\[
 \frac{\mathcal G_{\rm tor}}{\mathcal I_{\rm tor}}=1,
 \qquad
 \frac{\mathcal G_{\rm av}}{\mathcal I_{\rm av}}
 =\frac{\pi}{\varphi^2}
\]
describe distinct gravitational--inertial readers and remain
compatible with a single dimensional section \(G\). The construction of that
section also uses
\((\hbar_{\rm int},c_{\rm int},\ell_0)\) and the selector \(\theta_G\).

\subsection{Evaluation of the decimal gravitational chart and subsequent metrological contrast}

The decimal representation preserved in the corpus is
\begin{equation}
 \mathscr R_G([n,k,d]_3;t,e)
 :=10^{-n}\left(k+\frac{t}{10^3}+\frac{e}{10^d}\right).
 \label{eq:lector-decimal-G-rev11}
\end{equation}
For the chart
\([n,k,d]_3=[11,6,5]_3\), \(t=674\) and \(e=30\), the evaluation racional is
\begin{equation}
 \boxed{
 \mathscr R_G
 =10^{-11}\left(6+\frac{674}{1000}+\frac{30}{10^5}\right)
 =6.67430\times10^{-11}.}
 \label{eq:evaluacion-decimal-G-rev11}
\end{equation}
The triple, the two blocks and the decimal position constitute the declared data of this chart. Its evaluation is exact; its structural selection is governed by \cref{crit:selector-G-rev11}. This decomposition is retained as a historical decimal witness. The radial branch of \eqref{g26:eq:g-rigido-es} obtains its coordinate through length and action and gives \eqref{cr28:eq:g-evaluation}; it does not use blocks \(674\) and \(30\) to define the coupling.

After choosing the SI basis of \(\mathcal L_G\), the current contrast is
\begin{equation}
 \boxed{
 G_{\rm CODATA\,2022}=6.67430(15)\times10^{-11}
 \ \mathrm{m^3\,kg^{-1}\,s^{-2}}.}
 \label{eq:G-CODATA-2022-rev11}
\end{equation}

\begin{center}
\begin{tabularx}{0.94\textwidth}{@{}L{0.27\textwidth}Y@{}}
\toprule
Field & Metrological Statement\\
\midrule
Central value &
\(6.67430\times10^{-11}\ \mathrm{m^3\,kg^{-1}\,s^{-2}}\)\\
Standard uncertainty &
\(u(G)=0.00015\times10^{-11}
=1.5\times10^{-15}\ \mathrm{m^3\,kg^{-1}\,s^{-2}}\)\\
Incertidumbre relativa & \(u_r(G)\simeq2.2\times10^{-5}\)\\
Source and temporality &
2022 CODATA recommended values, published in 2024
\cite{CODATA2022}\\
Proof function &
Subsequent external contrast; neither the unit nor the uncertainty enters into
the rational evaluation of \eqref{eq:evaluacion-decimal-G-rev11}.\\
\bottomrule
\end{tabularx}
\end{center}

\paragraph{Reproducibility of the gravitational reader.}
An independent enumeration jointly verifies the dimensional type, the
decimal reader, and the monodromy identities on the declared finite
domain.

\begin{proposition}[Causal chain of the gravitational constant]
The covariant action, the restricted Hamiltonian, and the
Einstein--Schrödinger equation govern the gravitational dynamics. The dimensional type
and the transducer \eqref{eq:transductor-gravitatorio-rev11} are exact; the
evaluation \eqref{eq:evaluacion-decimal-G-rev11} is exact for its declared
chart. The value, unit, and uncertainty of
\eqref{eq:G-CODATA-2022-rev11} constitute an external comparison with CODATA 2022. The
autonomous selection of \(\theta_G\) is typed by the enriched
lift, the interfering response, and the holonomy functional subject
to the \cref{crit:selector-G-rev11}.
\end{proposition}
In the radial realization under rules R1--R8, selection is determined by \eqref{g26:eq:longitud-unica-es} and \eqref{g26:eq:g-rigido-es}; compatibility with the calendar is proved in \eqref{cr28:eq:circular-radial-agreement}. The holonomy template separately retains the conditions of its own domain.
 % Integración por delta de la adenda de realizaciones físicas.
% Residencia de procedencia: tipo dimensional de G, selector de holonomía y carta decimal.

\chapter{Dimensional Transduction and Holonomic Selection of the Gravitational Constant}
\label{chap:transductor-gravitatorio-G}
\index{gravitational constant!dimensional transducer}
\index{naturalness!gravitational selector}
\index{decimal reader!control of G}

\section{Gravitational line and dimensional form}
\label{sec:recta-gravitatoria-interna}

In the lattice of \cref{chap:funtor-dimensional-hmt}, the dimension of the
gravitational constant is
\[
 \mathbf G=-\mathbf S+5\mathbf C+2\mathbf T,
 \qquad
 [G]=L^3M^{-1}T^{-2}.
\]
Let
\[
 c_{\rm int}\in\mathcal L_C^+,\qquad
 t_0\in\mathcal L_T^+,\qquad
 \ell_0:=c_{\rm int}t_0\in\mathcal L_L^+,
 \qquad
 \hbar_{\rm int}\in\mathcal L_S^+.
\]
The section \(\hbar_{\rm int}\) is obtained only after the decade
determinantal reader \(-34\) and the choice of a basis for the action line.
It is not the dimensionless angular character preceding the return.

For every dimensionless character \(\theta_G>0\), define
\[
 \boxed{
 G_{\rm int}(\theta_G)
 =\theta_G
 \frac{c_{\rm int}^{5}t_0^2}{\hbar_{\rm int}}
 =\theta_G
 \frac{c_{\rm int}^{3}\ell_0^2}{\hbar_{\rm int}}
 \in\mathcal L_G^+.}
\]
The equality of the two expressions and their type are
\textsc{internally exact} within the dimensional lattice. They do not select
\(\theta_G\). If
\[
 L_G^2:=\frac{\hbar_{\rm int}G_{\rm int}}{c_{\rm int}^3},
\]
then
\[
 \theta_G=\left(\frac{L_G}{\ell_0}\right)^2.
\]
This final identity locates the selecting datum: the length ratio must be produced by a holonomy or spectral-response structure prior to the metrological chart. The realization of \cref{g26:sec:rigidez-radial-es} satisfies this requirement through inscription and the radial norm, with \(L_G=L\). What follows examines other return and holonomy templates, preserving the obstructions specific to their quotients.

\section{Obstruction of the Positive Character on the Observable Return}
\label{sec:selector-holonomia-G}

The \cref{sec:jerarquia-retornos-27-54-108} defines
\[
 \mathcal K_{27}:=F_{\rm obs}^{27}.
\]
The primary proposition \cref{prop:retorno-fobs-identidad} proves
\[
 F_{\rm obs}^{108}=I;
\]
Therefore, in the observable quotient,
\[
 \mathcal K_{27}^{4}=I.
\]

\begin{theorem}[Torsional Obstruction for the Observable Positive Selector]
\label{thm:no-go-chi-G-observable}
Every multiplicative homomorphism
\[
 \chi:\langle\mathcal K_{27}\rangle
 \longrightarrow\mathbb R_{>0}^{\times}
\]
is trivial. In particular,
\[
 \chi(\mathcal K_{27})=1.
\]
\end{theorem}

\begin{proof}
Multiplicativity and \(\mathcal K_{27}^{4}=I\) imply
\[
 \chi(\mathcal K_{27})^4=\chi(I)=1.
\]
The sole positive real element whose fourth power is one is \(1\).
\end{proof}

The observable order is exactly four, and not merely a divisor of four:
in the coordinate \(\theta\in\mathbb Z/108\mathbb Z\),
\[
 \mathcal K_{27}(\theta)=\theta+27,
 \qquad
 \mathcal K_{27}^{2}(\theta)=\theta+54\neq\theta,
 \qquad
 \mathcal K_{27}^{4}(\theta)=\theta.
\]
The same cyclic group admits unitary characters.
\[
 \chi_k(\mathcal K_{27})
 =\exp\!\left(\frac{2\pi i k}{4}\right),
 \qquad k=0,1,2,3,
\]
and, in particular, the phase values \(\pm i\). The typed conclusion is
that the observable cycle can transport phase but not a nontrivial positive
scale. These phase characters are not identified with the Paley operator
\(J_W\): they satisfy an analogous quadratic relation on distinct
domains.

The obstruction does not depend on a metrological comparison. It affects the very
form of the selector: a positive character on the observable cycle cannot
produce a nontrivial \(\theta_G\). Moreover, a mere set of conjugacy
classes is not, in general, a group and cannot be declared the domain
of a character without first defining its composition.

\subsection{Unique admissible multiplicative continuation}

A multiplicative route can reopen only on an enriched object that
does not inherit the observable torsion. For an object \(x\) of the memory groupoid,
the correctly typed domain would be the isotropy group
\[
 H_x:=\operatorname{Aut}_{\mathcal G_{\rm mem}}(x)
\]
and, by multiplicativity towards an abelian group, its abelianization
\[
 H_x^{\rm ab}\longrightarrow\mathbb R_{>0}^{\times}.
\]
Every positive character annihilates the torsion part of \(H_x^{\rm ab}\). Therefore,
the proposal must construct a lift
\[
 \widetilde{\mathcal K}_{27}\in H_x,
 \qquad
 p(\widetilde{\mathcal K}_{27})=\mathcal K_{27},
\]
prove that its Abelian class is non-torsion and naturally select
a character \(\widetilde\chi_G\). The positive route over the observable cycle
is thus closed by obstruction: without that non-torsion datum there is no
nontrivial multiplicative selector. If every admissible lift
satisfied
\(\widetilde{\mathcal K}_{27}^{\,4}=I\), the multiplicative route would also be
ruled out in the enriched domain.

\section{Chronological path of length, action, and spectrum}
\label{sec:via-cronologica-espectral-G}

The \cref{thm:cobertura-acumulativa-dirigida} provides a third route,
distinct both from the obstructed observable character and from a possible
lift in an isotropy group. On
\[
 \widetilde{\mathcal O}_{108}^{+}
 =\mathcal O_{108}\times\mathbb N^2,
\]
the twenty-seven-step return adds \((1,18)\) and its fourth power adds \((4,72)\). This action belongs to a directed chronological category. It is neither a lift of the groupoid nor a field proved native to the complete integral holonomic state.

The already constructed positive functional,
\[
 \mathcal A_\ast(\gamma)
 =G_{\rm or}(\gamma)+\lambda_\ast S(\gamma),
\]
is additive under concatenation of directed histories. It is not a character of the
isotropy group, because its indicators do not change sign under a
formal inversion. The completion
\(\mathbb N^2\to\mathbb Z^2\) permits unitary characters to be introduced for
harmonic analysis, but does not alter this fact.

\begin{equation}
\label{eq:fraccion-angular-orientada-G}
 \delta_\theta
 :=\frac{A_{\rm deg}-C_{\rm deg}^{\ast}}{360}.
\end{equation}
This oriented angular fraction is defined once the unique angle and
counter-angle have been constructed. It will be used in the two subsequent
dimensional templates and does not intervene retroactively in obtaining
\(A_{\rm deg}\) or \(C_{\rm deg}^{\ast}\).

\subsection{Twisted operator and deterministic negative control}

For the exact control, let \(\Gamma=(V,E)\) be an enriched directed
\emph{finite} multigraph, with measure
\(\mu:V\to\mathbb R_{>0}\), weights \(a_\ast(e)\geq0\), parameter
\(\beta\geq0\), and cocycle
\(c(e)\in\mathbb Z^2\). For a fixed normalization, consider
\[
 (\mathcal L_{\beta,\vartheta}f)(y)
 =
 \sum_{e:x\to y}
 e^{-\beta a_\ast(e)}
 e^{i\langle\vartheta,c(e)\rangle}f(x),
 \qquad
 \vartheta\in\mathbb T^2,
\]
and take the entire finite Hilbert space
\(\ell^2(V,\mu)\) as the domain. Then \(\mathcal L_{\beta,\vartheta}\) is bounded and
\[
 \mathcal H_{\beta,\vartheta}
 :=\mathcal L_{\beta,\vartheta}^{*}
   \mathcal L_{\beta,\vartheta}\geq0
\]
is self-adjoint. An extension to an infinite graph would additionally require
bounded operators, or densely defined closed operators such
that \(\mathcal L_{\beta,\vartheta}^{*}
\mathcal L_{\beta,\vartheta}\) were self-adjoint on a common domain;
that extension is not assumed here.

\begin{proposition}[Phase cancellation in the deterministic orbit]
\label{prop:control-negativo-espectral-determinista}
If \(\Gamma\) is solely the observable deterministic orbit and every vertex
has a single incoming edge, then
\(\mathcal H_{\beta,\vartheta}\) is independent of \(\vartheta\). In
particular, the Hessian of any of its bands with respect to
\(\vartheta\) is null.
\end{proposition}

\begin{proof}
In each component, \(\mathcal L_{\beta,\vartheta}\) multiplies the sole
incoming term by a unitary phase. In
\(\mathcal L_{\beta,\vartheta}^{*}\mathcal L_{\beta,\vartheta}\), that phase
is multiplied by its conjugate and cancels. No cross-terms remain
between distinct histories.
\end{proof}

This control prevents spectral rigidity from being attributed to the mere deterministic cycle.
A nontrivial phase response requires at least one of the following
constructions:
\begin{enumerate}
 \item an enriched multigraph in which two or more coherent histories
 contribute to the same amplitude; or
 \item to magnetic differential
 \[
  (d_\vartheta f)(e)
  =
  e^{i\langle\vartheta,c(e)\rangle/2}f(t(e))
  -
  e^{-i\langle\vartheta,c(e)\rangle/2}f(s(e)),
 \]
 with Laplacian
 \(\Delta_\vartheta=d_\vartheta^*d_\vartheta\), whose phase depends on the
 holonomy of enriched cycles.
\end{enumerate}
In both cases, the Hilbert space, domain, measure,
weights, and normalization must be declared, and it must be proved that the family is not
unitarily equivalent to the family without phase.

Suppose that
\(\vartheta\mapsto\mathcal H_{\beta,\vartheta}\) is an analytic self-adjoint family on a common domain and that
\(\lambda_0(\vartheta)\) is a simple isolated band in a neighborhood of
\(\vartheta=0\). If \(0\) is a critical point and a local minimum,
\[
 Q_{ij}
 =
 \left.
 \frac{\partial^2\lambda_0(\vartheta)}
      {\partial\vartheta_i\partial\vartheta_j}
 \right|_{\vartheta=0}
\]
is positive semidefinite. Without the minimum hypothesis, \(Q\) represents only
the quadratic response. For
\[
 v_{27}=(1,18),
 \qquad
 \mathcal R_{27}=v_{27}^{\mathsf T}Qv_{27},
\]
the condition \(\mathcal R_{27}\neq0\) is a minimal test of spectral
sensitivity to the accumulated memory.

\subsection{Typed normalization of the spectral response and construction of
\texorpdfstring{\(\theta_G^{\mathrm{spec}}\)}{theta G
spec}}

If a natural and unique normalization \(\mathfrak N\) converts the preceding
response into a positive ratio independent of the operator's global
scale, one may define
\[
 \Lambda_G
 :=\ell_0\,\mathfrak N(Q,v_{27}),
 \qquad
 \theta_G^{\rm spec}
 :=
 \left(
 \delta_\theta\frac{\Lambda_G}{\ell_0}
 \right)^2,
\]
and, by the dimensional transducer,
\[
 \boxed{
 G_{\rm spec}
 =
 \frac{c_{\rm int}^{3}}{\hbar_{\rm int}}
 \bigl(\delta_\theta\Lambda_G\bigr)^2.}
\]
The homogeneity of the template is exact. The class of its evaluations is
characterized by the conditions of canonical operator selection,
nonzero response, naturality and stability under refinement, independence
of rescaling, uniqueness of \(\mathfrak N\), and sealing prior to any query of
the metrological value of \(G\). Without those conditions, rescaling
invariance proves the nonuniqueness of the selector; with them, the evaluation is
determined on the declared normalized domain.

\section{Nonhomomorphic functional over enriched paths}
\label{sec:funcional-longitud-K27-G}

The second possible continuation abandons the multiplicative character and uses
a geometric, spectral, or action functional. Let
\[
 \mathscr L:
 \operatorname{Loop}(\mathcal X_{\rm TPK}^{\rm enr})
 \longrightarrow\mathcal L_L^+\cup\{0\}
\]
invariant under cyclic change of base point, admissible relabeling,
conjugation, and subdivision. Its concatenation law must be stated
explicitly---for example, as a subadditive norm or a spectral functional---;
it is not presumed to be a homomorphism of the torsion group. For a
candidate enriched lift, we write
\[
 \Lambda_{27}:=\mathscr L(\widetilde{\mathcal K}_{27}),
 \qquad
 \rho_{27}:=\frac{\Lambda_{27}}{c_{\rm int}t_0}.
\]
The oriented angular fraction of
\eqref{eq:fraccion-angular-orientada-G} produces the dimensionless candidate
\[
 \theta_G^{\rm hol}:=(\delta_\theta\rho_{27})^2.
\]
His specialization in the transducer gives
\[
 \boxed{
 G_{\rm hol}
 =\frac{c_{\rm int}^{3}}{\hbar_{\rm int}}
  \bigl(\delta_\theta\Lambda_{27}\bigr)^2.}
\]
The temporal form is the same identity:
\[
 T_{27}:=\frac{\Lambda_{27}}{c_{\rm int}},
 \qquad
 t_\ast:=\delta_\theta T_{27}.
\]
In particular, it is not appropriate to replace \(t_\ast\) with
\(T_{27}/\delta_\theta\).

These formulas are exactly homogeneous. This does not promote their status to a derivation of \(G\) by this template: selection remains open until \(\widetilde{\mathcal K}_{27}\) and \(\mathscr L\) are constructed without consulting the external value, their naturality is proved and a unique normalization is fixed. Downstream use of \(A\) and \(C^\ast\) in \(\delta_\theta\) does not authorize using \(G\) to redefine them. The radial proof of \cref{g26:sec:rigidez-radial-es} is a different specified realization: its data, norm, length and coupling are explicitly constructed and its reduction preserves memory. The limitation of the preceding template does not transfer to that proof.

\section{Transduction of the gravitational decimal chart onto the magnitude line and causal separation of metrological recognition}
\label{sec:lector-decimal-independiente-G}

Let
\[
 \mathscr R_G([n,k,d]_3;t,e)
 :=10^{-n}\left(k+\frac{t}{10^3}+\frac{e}{10^d}\right),
\]
where the subscript of \([n,k,d]_3\) records the reader triad and does not denote
a ternary numeral. For the declared data
\[
 [n,k,d]_3=[11,6,5]_3,\qquad t=674,\qquad e=30,
\]
rational evaluation is
\[
 \boxed{
 \mathscr R_G
 =10^{-11}\left(6+\frac{674}{1000}+\frac{30}{10^5}\right)
 =6.67430\times10^{-11}.}
\]
The arithmetic equality is exact once the parameters have been declared. The triad,
the blocks \(674\), \(30\), and their decimal position have the status of chart
data, so this branch is preserved as a parametrized control and subsequent
recognition. The number is dimensionless until a basis
\(u_G\in\mathcal L_G^+\) is chosen; only then is the section defined
\[
 G_{\rm dec}:=\mathscr R_G\,u_G.
\]
The agreement of its coordinate with a metrological publication is \textsc{external recognition}. A future promotion would require the integers \(674\) and \(30\), the triple \([11,6,5]\), and the decimal placement to be uniquely forced by published symmetries and fixed before the external value is consulted.

\section{Hierarchy of Evidence for Gravitational Constructions}
\label{sec:cuadrado-contraste-G}

The longitudinal and radial realization determines the coupling in the explicit R1--R8 domain. Alternative templates and preceding controls retain their own scopes; they do not by themselves constitute other proved genealogies of that realization. The organization is:
\begin{center}
\begin{tabularx}{0.96\textwidth}{@{}L{0.25\textwidth}Y L{0.26\textwidth}@{}}
\toprule
Object & Result available & Verdict\\
\midrule Inscription and radial norm R1--R8 &
\(L=(5759/23040)\alpha^{16}L_*\),
\(R_X\Lambda_X=L^2I\), \(H_X\Lambda_X=\hbar_{\rm ret}cI\),
\(G=c^3L^2/\hbar_{\rm ret}\). &
Unique determination in the declared domain; memory, reduction and refinement preserved.\\
CT108 circular realization &
\(t_0=\pi L/(54c)\), \(R_{108}=L\) y
\(\theta_{108}^{\rm circ}=(54/\pi)^2\). &
Exact compatibility with the constructed length.\\
Transductor dimensional &
Types \(G_{\rm int}(\theta_G)\) exactly, without fixing \(\theta_G\). &
Dimensional theorem; no selector.\\
Character of the observable \(\mathcal K_{27}\) &
\(\mathcal K_{27}^4=I\) force
\(\chi(\mathcal K_{27})=1\). &
Obstructed formulation.\\
Levantamiento enriquecido &
Possible only with non-torsion abelian class, natural character and unique normalization. &
Obstruction of canonicity demonstrated without those data.\\
Chronological coverage and spectral response &
The cumulative translation has infinite semigroup order; along the deterministic
orbit, the phase cancels. &
Null response proved on the deterministic orbit; nontrivial class
characterized by interference and natural normalization.\\
Path functional &
The formula for \(G_{\rm hol}\) is homogeneous once
\(\Lambda_{27}\) has been fixed. &
Exact parametrized family; no unique selector exists without specifying
\(\mathscr L\) and its normalization.\\
Decimal reader &
Reproduce exactly the declared coordinate, but receive its blocks and
position. &
Control parametrizado.\\
\bottomrule
\end{tabularx}
\end{center}

\begin{criterion}[Gravitational constant selector] \label{crit:selector-G} For the alternative lift, spectral-response and decimal-reader templates examined in this chapter, promotion is rejected if:
\begin{enumerate}
 \item the enriched lift preserves only torsion or does not descend
 to a well-defined Abelian class;
 \item the chronological operator does not preserve interference, produces
 \(\mathcal R_{27}=0\), lacks a domain or measure, or its response changes under
 an admissible refinement;
 \item \(\mathscr L\) depends on the basepoint, relabelling, the
 chosen orientation, a subdivision, or a free normalization;
 \item the normalization \(\mathfrak N\) depends on the global scale of
the operator, is nonunique, or is fixed by consulting \(G\);
 \item in the decimal reader, \(674\), \(30\) or the decimal position admit alternatives with equal structural entitlement;
 \item the candidate consults the metrological value before being sealed;
 \item the value is not stable under refinement or does not define a
section of \(\mathcal L_G\) in a common dimensional chart.
\end{enumerate}
\end{criterion}

\section{Gravitational--inertial ratios and the dimensional constant}
\label{sec:razones-no-rivales-G}

The local normalization of the torsion chart,
\[
 \frac{\mathcal G_{\rm tor}}{\mathcal I_{\rm tor}}=1,
\]
and the interface areal--volumetric of the
\cref{sec:ambivalencia-grav-inercial-pi-phi2},
\[
 \frac{\mathcal G_{\rm av}}{\mathcal I_{\rm av}}
 =\frac{\pi}{\varphi^{2}},
\]
are dimensionless reader ratios. They are not rival values of the
dimensional constant \(G\). The first expresses a local diagonal
constraint; the second is an interface conditional upon factorization through a
common memory. Constructing a section of
\(\mathcal L_G\) additionally requires the
\((\hbar_{\rm int},c_{\rm int},\ell_0)\) framework and a
\(\theta_G\) selector.
 
\chapter{Perron dynamics, torsional memory and Einstein--Cartan realization}
\begingroup
\renewcommand{\chapter}[2][]{}%
\chapter{Perron dynamics, \texorpdfstring{\(a^{-6}\)}{a⁻⁶} dilution of memory, and Einstein–Cartan realization}
\label{chap:parte-iv-67}

The Perron operator fixes the stable memory mode; when published as a spin density, its contribution decays as the inverse sixth power of the scale factor. The Einstein--Cartan realization preserves this provenance and its conditions of validity.

\section{Perron dynamics and the stable memory mode}

On a closed route, the cylindrical Perron cocycle satisfies

\begin{equation}
V_n=\varphi^n.
\label{rm:eq:cosmo-perron-V}
\end{equation}

In spatial dimension \(D=3\), the scale determined by the cocycle is

\begin{equation}
\frac{a_n}{a_0}=V_n^{1/3}=\varphi^{n/3},
\label{rm:eq:cosmo-perron-a}
\end{equation}

whereas the stable mode of the quadratic action is

\begin{equation}
M_n=\varphi^{-2n}=V_n^{-2}.
\label{rm:eq:cosmo-perron-M}
\end{equation}

Eliminating \(n\) between \cref{rm:eq:cosmo-perron-a,rm:eq:cosmo-perron-M} gives the
central result:

\begin{theorem}[Law cylindrical of sixth order for the memory]
For every closed level of the Perron realization in \(D=3\),
\begin{equation}
\boxed{
M_n=\left(\frac{a_n}{a_0}\right)^{-6}.}
\label{rm:eq:cosmo-minus-six}
\end{equation}
The exponent \(-6\) is not fitted: it is the quotient of the stable exponent
\(-2\) by the dimensional root \(1/3\).
\end{theorem}

In the distinguished returns one obtains

\[
\frac{a_9}{a_0}=\varphi^3,
\qquad
\frac{a_{27}}{a_0}=\varphi^9,
\qquad
\frac{a_{108}}{a_0}=\varphi^{36},
\]

and the same levels satisfy exactly
\(M_n=(a_n/a_0)^{-6}\). If the additive temporal parameter \(t_n=nt_0\) is declared, the
logarithmic rate is constant:

\begin{equation}
H_\varphi
:=\frac{\Delta\log a_n}{\Delta t}
=\frac{\log\varphi}{3t_0},
\qquad
H_\varphi t_0
=0.1604039416865344824992529711\ldots.
\label{rm:eq:cosmo-Hphi}
\end{equation}

The symbol \(H_\varphi\) names the rate of this discrete scale. Reading it
as a Hubble parameter requires the cosmographic transducer; the internal
identity does not depend on that promotion.

\section{Realization in the system of Einstein--Cartan}

In an Einstein--Cartan spin fluid, the quadratic spin density
obeys

\begin{equation}
s^2\propto a^{-6}.
\label{rm:eq:cosmo-EC-scaling}
\end{equation}

The \cref{rm:eq:cosmo-minus-six} provides the exponent and the stable mode without
consulting a bounce. To make the physical reading, one must declare a
positive section \(s_*^2\) and a map

\begin{equation}
\mathcal T_{\rm EC}:M_n\longmapsto
s_n^2:=s_*^2M_n
=s_*^2\left(\frac{a_n}{a_0}\right)^{-6}.
\label{rm:eq:cosmo-EC-transducer}
\end{equation}

The map preserves the scaling law, but does not by itself determine the amplitude
\(s_*^2\), the spin current, the sign of the term of contorsion, or the
equation of state. With mass density, an effective equation has the
form

\begin{equation}
H^2=\frac{8\pi G}{3}
\bigl(\rho_{\rm m}-\rho_{\rm spin,m}\bigr)+\cdots,
\label{rm:eq:cosmo-EC-mass}
\end{equation}

whereas with energy density it must be written

\begin{equation}
H^2=\frac{8\pi G}{3c^2}
\bigl(\rho_{\rm E}-\rho_{\rm spin,E}\bigr)+\cdots.
\label{rm:eq:cosmo-EC-energy}
\end{equation}

The two conventions are not mixed. A bounce candidate requires, in addition to
equality of the densities, regularity, matter, and boundary data.

\begin{falsifier}
The Einstein--Cartan transduction fails if the scale functional produces an
exponent other than \(-6\), if \(V_n\ne\varphi^n\) along the declared route, or
if the physical map does not preserve positivity and dimensions. Retrospectively choosing
\(s_*^2\) to force a bounce threshold does not establish the
transducer.
\end{falsifier}

\section{Structural comparison between the 108-state cycle and Perron dynamics}

\begin{center}
\begin{tabularx}{\textwidth}{>{\bfseries}p{29mm}XX}
\toprule
& CT108--de Sitter framework & Perron--Einstein--Cartan framework\\
\midrule
Antecedent & circular return of order \(108\) & cylindrical cocycle and stable mode\\
Selection mechanism & \(r_g=R_{108}\) & \(V_n=\varphi^n\), \(D=3\)\\
Magnitudes determinadas & radius, curvature, temperature, and entropy & expansion, rate, and dilution of memory\\
Invariant & \(G_a\hbar_a\) and circular radius & \(M_na_n^6\)\\
Status & de Sitter conditioned interface & exact internal law; conditioned EC reading\\
Falsification criterion & incompatible radius or background equation & perdida of the exponente \(-6\) or of the transductor\\
\bottomrule
\end{tabularx}
\end{center}

\begin{statusbox}
The CT108 identities, the selector \(\theta_{108}\), and the law
\(M_n=(a_n/a_0)^{-6}\) are \status{Internal exact theorem}. The reading of the same radius
as de Sitter and the identification of \(M_n\) with a spin density are
\status{Conditional physical realizations}, each with its own falsification criterion.
\end{statusbox}
 
\paragraph{Coordination of torsion and area in the Palatini--Cartan--Holst action.} Torsion and area are then coordinated in the first-order Palatini--Cartan--Holst action. \endgroup

\chapter{Palatini--Cartan--Holst realization of the Barbero--Immirzi parameter}
\begingroup
\renewcommand{\chapter}[2][]{}%
\chapter{Palatini–Cartan–Holst realization of the Barbero–Immirzi parameter and physical transduction of the area spectrum}
\label{chap:parte-iv-68}

The structural derivation of the parameter and the area spectrum precedes this chapter. Here the tetrad, connection, curvature, and first-order action are constructed, and only then is the Holst parameter specialized as a physical realization of the areal reader.

\section{Curvature, torsion, and Einstein--Cartan--Holst action}
\label{chap:gravedad-cartan-holst}

The HMT transport structure distinguishes the memory variable,
holonomy, and closure defect. Its gravitational realization is formulated
by means of a tetrad and a connection. This section establishes the first-order
geometry and specifies the type of the map between it and the TPK state space
\@.

\subsection{Discrete Gauss–Stokes identity and incidence transport to the boundary}

Let \(K\) be a cellular complex, let \(W\) be an abelian
group, and let \(\Omega\in C^1(K;W)\) be a cochain. We denote by
\(\delta:C^1(K;W)\to C^2(K;W)\) the cellular coboundary and by \(\langle\,\cdot\,,\,\cdot\,\rangle\) the pairing
between cochains with values in \(W\) and integral cellular chains.

\begin{theorem}[Universal Stokes Cell Identity]
For every \(2\)-cadena finite \(z\in C_2(K;\mathbb Z)\),
\[
 \boxed{\langle\delta\Omega,z\rangle
 =\langle\Omega,\partial z\rangle.}
\]
\end{theorem}
\begin{proof}
This is the identity that defines the coboundary as the adjoint operator of the cellular
boundary. The finiteness of \(z\) ensures that both pairings are finite
sums in \(W\).
\end{proof}

\begin{corollary}[Gauss--Stokes on a cellular pseudomanifold]
Let \(K\) be a finite, compact, regular two-dimensional cellular pseudomanifold,
coherently oriented and possibly with boundary. If \(z_K\) is the sum of
its \(2\)-cells with the coherent orientation, then
\[
 \sum_{f\in K^{(2)}}(\delta\Omega)(f)
 =\sum_e[\partial z_K:e]\,\Omega(e).
\]
In particular, regularity allows the uncancelled incidences of
\(\partial z_K\) to be identified with the edges of the geometric boundary.
\end{corollary}
\begin{proof}
The theorem is applied to \(z_K\). Coherent orientation makes each
interior edge appear twice in \(\partial z_K\) with opposite coefficients;
the uncancelled incidences form the boundary chain.
\end{proof}

The non-Abelian version is not obtained merely by replacing sums with products.
In this work it is used only when transports to a common
base, an ordering of the holonomies, and the corresponding conjugation have been fixed. Outside
that subdomain, no non-Abelian Stokes identity is asserted.

\subsection{Area capacity of the qutrit cell: \(81\log 3\) and collective incidence transport}

\begin{theorem}[Minimum cut of a discrete cube]
In the cubic graph \(N\times N\times N\) with unit capacity, every cut
separating two opposite faces has capacity at least \(N^2\), and one exists
with capacity \(N^2\).
\end{theorem}
\begin{proof}
There are \(N^2\) edge-disjoint paths crossing the cube in the direction
normal to the two faces. Every cut must intersect them, so its capacity is
at least \(N^2\). A transverse layer contains exactly \(N^2\) edges and
attains the bound.
\end{proof}

For \(N=9\), the cut equals 81 and
\[
 9^3=729=27^2.
\]
The last equality permits six trits to be regrouped as a volume \(9^3\) or a
grid of cardinality \(27^2\). It is a bijection of sets; it is not an
isomorphism between
\((\ZZ/9\ZZ)^3\) and \((\ZZ/27\ZZ)^2\), whose exponents are distinct.

\subsection{Tetrad, connection, and \(2\) defects}

Let \(e^I\) be a tetrad and \(\omega^{IJ}=-\omega^{JI}\) a Lorentz
connection. Define
\[
 T_{\rm tor}^I=D_\omega e^I
 =\dd e^I+\omega^I{}_J\wedge e^J,
\]
\[
 F^{IJ}=\dd\omega^{IJ}+\omega^I{}_K\wedge\omega^{KJ},
 \qquad
 g=\eta_{IJ}e^I\otimes e^J.
\]
The decomposition
\[
 \omega^{IJ}=\mathring\omega^{IJ}(e)+K^{IJ}
\]
separates the Levi--Civita connection and the contorsion. Then
\[
 T_{\rm tor}^I=K^I{}_J\wedge e^J.
\]
Curvature and torsion are distinct defects: \(F\) measures closure of
vector transport; \(T_{\rm tor}\), closure of the tetrad. The
\(T\) memory of the Ambivalence Principle is likewise not geometric torsion: only a
realization functor can relate them.

\subsection{1st-order action}

Let \(M\) be a smooth oriented four-dimensional manifold. If
\(\partial M\ne\varnothing\), we assume compactly supported variations in the
interior or, equivalently for the variational problem considered, a
boundary term and boundary conditions that annul the boundary
contribution. When \(\Psi\) contains spinorial fields, we also assume a
spin structure and fields compatible with it. Let
\[
 \Sigma^{IJ}=e^I\wedge e^J,
 \qquad
 (\star X)^{IJ}=\frac12\epsilon^{IJ}{}_{KL}X^{KL},
 \qquad
 \kappa=\frac{8\pi G}{c^4},
 \qquad
 \gamma\in\mathbb R\setminus\{0\}.
\]
The Palatini--Cartan--Holst action is
\[
\boxed{
 S[e,\omega,\Psi]
 =\frac1{2\kappa c}\int_M
 \Sigma_{IJ}\wedge\left(\star F^{IJ}+\frac1\gamma F^{IJ}\right)
 -\frac\Lambda{\kappa c}\int_M\operatorname{vol}_e
 +S_{\rm m}[e,\omega,\Psi].}
\]
The variation of this action determines the torsional corrections
\cite{Sciama1964,Kibble1961,Hehl1976,Holst1996}.

We define the spin current by
\[
 \delta_\omega S_{\rm m}
 =-\frac12\int_M\delta\omega^{IJ}\wedge\sigma_{IJ}.
\]

\begin{theorem}[Cartan--Holst Equation]
Under the preceding geometric and boundary hypotheses, variation with respect
to \(\omega\) gives
\[
 D_\omega\!\left[(\star+\gamma^{-1}I)\Sigma^{IJ}\right]
 =\kappa c\,\sigma^{IJ},
\]
with the global normalization fixed by the preceding definition of
\(\sigma^{IJ}\).
\end{theorem}
\begin{proof}
One uses \(\delta F^{IJ}=D_\omega\delta\omega^{IJ}\), integrates by parts, and
sets the coefficient of the arbitrary variation
\(\delta\omega^{IJ}\) equal to zero. With the sign convention fixed in the definition of
\(\sigma^{IJ}\), the matter term moves to the right-hand side with a
positive sign. The factor \(1/(2\kappa c)\) produces the coefficient
\(\kappa c\). The cosmological term does not depend on \(\omega\), and the
boundary conditions eliminate the term generated by integration by
parts.
\end{proof}

In signature lorentziana \(\star^2=-I\) on bivectores and, for
\(\gamma\in\RR\setminus\{0\}\),
\[
 (\star+\gamma^{-1}I)^{-1}
 =\frac{\gamma^2}{1+\gamma^2}(-\star+\gamma^{-1}I).
\]

\begin{corollary}[Torsional decoupling]
For a co-tetrada no degenerada and
\(\gamma\in\mathbb R\setminus\{0\}\),
\[
 \sigma^{IJ}=0\Longrightarrow T_{\rm tor}^I=0
 \Longrightarrow\omega=\mathring\omega(e).
\]
\end{corollary}

To make the normalization explicit, we define the normalized sources by
the action through
\[
 \delta_g S_{\rm m}
 =-\frac12\int_M\operatorname{vol}_e\,
 \mathcal T^{S}_{\mu\nu}\,\delta g^{\mu\nu}.
\]
When the tetrad coordinates have the dimension of length, the conventional
physical energy--momentum tensor is
\(T^{\rm phys}_{\mu\nu}:=c\,\mathcal T^{S}_{\mu\nu}\). We likewise denote by
\(\mathcal U^{S,\rm spin}\) the contribution normalized by the
action and by \(U^{\rm phys,\rm spin}:=c\,\mathcal U^{S,\rm spin}\) its physical
convention. Upon algebraically eliminating the connection, the metric
equation then assumes the two equivalent forms
\[
 G_{\mu\nu}+\Lambda g_{\mu\nu}
 =\kappa c\,
 \bigl(\mathcal T_{\mu\nu}^{S,\rm LC}
       +\mathcal U_{\mu\nu}^{S,\rm spin}\bigr)
 =\kappa\,
 \bigl(T_{\mu\nu}^{\rm phys,LC}
       +U_{\mu\nu}^{\rm phys,spin}\bigr),
\]
where the spin contribution is quadratic in the corresponding current
and depends on the matter coupling. The coefficient of that contribution is not
fixed without specifying the fermionic action and the duality
conventions.

\subsection{Bianchi Identity with Torsion and Flow Balance in the Cartan Connection}

For every connection,
\[
 D_\omega T_{\rm tor}^I=F^I{}_J\wedge e^J,
 \qquad
 D_\omega F^{IJ}=0.
\]
After removing torsion,
\[
\nabla^\mu(T_{\mu\nu}^{\rm phys,LC}
          +U_{\mu\nu}^{\rm phys,spin})=0.
\]
These identities are mandatory consistency checks: every HMT
gravitational term must arise from a compatible action or satisfy the corresponding
Noether balance.

\subsection{Map of realization and specialization of the Holst parameter}

The insertion defined in \cref{mdg:chap:barbero} takes
\[
 \gamma=\Gamma_{6\to8}.
\]
Once this substitution has been made, all the preceding variational consequences
are theorems of the specialized action. The physical selection of
this action by HMT transport is expressed through a realization map
\[
 \mathfrak R_{\rm grav}:\mathcal X_{\rm TPK}
 \dashrightarrow\{(e,\omega,\Psi)\}
\]
endowed with three properties: intertwining TPK transport with that of
\(\omega\), identifying the holonomy classes, and furnishing the scales of
\(G\) and \(\Lambda\). These properties fix the precise domain of a
gravitational realization of the formalism.

For a nondegenerate co-tetrad, nonzero real \(\gamma\), and matter without spin current, the torsional sector reduces to the first-order formulation of general relativity with the same cosmological constant, subject to the already declared boundary conditions. In a spin fluid, the scaling \(s^2\propto a^{-6}\) can produce an Einstein--Cartan bounce; obtaining its density and initial conditions from the TPK state is a question distinct from the variational consistency of the action.
 
\section{Nonadic Return, Memory Groupoid and Cartan Entanglement Condition}

The hierarchy \(27\to54\to108\), the directed cumulative coverage, and the
memory groupoid supply the exact condition under which the
TPK return admits a compatible Cartan connection. This bridge precedes the
Palatini--Cartan--Holst specialization and preserves its domain and
refutation criteria intact.

\begingroup
\let\chapter\section
\let\section\subsection
\let\subsection\subsubsection
\let\subsubsection\paragraph
\let\clearpage\relax
\let\cleardoublepage\relax
\chapter{Nonadic Return, Memory, and Cartan Holonomy}
\label{chap:puente-retorno-holonomia-cartan}
\index{return!27--54--108}\index{Holonomy!TPK--Cartan}
\index{memory!distinction with respect to torsion}

The gravitational realization does not begin by identifying a discrete variable
with a geometric magnitude. It begins by distinguishing three levels: return
of observable transport, memory distinguishing histories with equal
projection, and holonomy of a connection. The bridge formulated here
receives the returns and cocycles already demonstrated in HMT; it does not reconstruct them
from the Holst action, from a mass, or from a metrological constant.

\section{Exact hierarchy of returns}
\label{sec:jerarquia-retornos-27-54-108}

Let \(F_{\rm obs}\) be the observable chronology defined in the HMT kernel, and let
\[
 \mathcal K_{27}:=F_{\rm obs}^{27}.
\]
The symbol \(\mathcal K_{27}\) denotes exclusively the composition of twenty-seven updates of the positional and oriented loop; it does not denote a generic window, a trace of \(108\) steps, or the complete integral holonomic state. Denote by \(p\) the projection onto the positions of the two cursors, and by \(\mathsf J_{\rm or}\) simultaneous reversal of their orientations. The primary proof is in \cref{prop:retorno-fobs-identidad}; its typed consequences are
\[
 p(\mathcal K_{27}x)=p(x),\qquad
 o(\mathcal K_{27}x)=\mathsf J_{\rm or}\,o(x),
 \qquad \mathsf J_{\rm or}^{\,2}=I,
\]
\[
 F_{\rm obs}^{54}:
 \text{restores positions and orientations},
 \qquad
 F_{\rm obs}^{108}=I:
 \text{also restores the observable phase}.
\]
Thus, position return, orientation return, and phase return are
distinct assertions. None of them by itself implies the return of
all fields of an enriched extension.

The mode quotient and orientation cocycles are established in
\cref{sec:cociclos-modo-orientacion,thm:descenso-necesario-fav}. In areal--volumetric
order,
\[
 N_{27}=9F_{\rm av},
 \qquad
 F_{\rm av}=
 \begin{pmatrix}0&1\\1&1\end{pmatrix}.
\]
If
\[
 \Omega_{\rm sw}(\gamma)=\sum_{a\subset\gamma}\omega_{\rm sw}(a),
 \qquad
 S(\gamma)=\sum_{a\subset\gamma}|\omega_{\rm sw}(a)|,
 \qquad
 G_{\rm or}(\gamma)=\sum_{a\subset\gamma}g(a),
\]
then a block of twenty-seven steps contains three nonadic turns and one
orientation reversal:
\[
 \Omega_{\rm sw}(\mathcal K_{27})=0,\qquad
 S(\mathcal K_{27})=18,\qquad
 G_{\rm or}(\mathcal K_{27})=1.
\]
In the observable return of \(108\) steps,
\[
 \Omega_{\rm sw}(C_{108})=0,\qquad
 S(C_{108})=72,\qquad
 G_{\rm or}(C_{108})=4.
\]
Thus, for the additive functional declared
\[
 \mathcal A_\ast(\gamma)
 :=G_{\rm or}(\gamma)+\lambda_\ast S(\gamma),
\]
are obtained, without introducing a dimensional scale,
\[
 \mathcal A_\ast(\mathcal K_{27})=1+18\lambda_\ast,
 \qquad
 \mathcal A_\ast(C_{108})=4+72\lambda_\ast.
\]
These values are record counts. Their subsequent thermal or gravitational
use does not alter their provenance.

\subsection{Directed system of coverings, refinement, and limit of the Cartan connection}
\label{sec:cobertura-acumulativa-dirigida}

Let \(\mathcal O_{108}\) be the observable orbit and, for each chronological edge
\(e\), let
\[
 c(e)=\bigl(g(e),s(e)\bigr)\in\mathbb N^2
\]
the pair formed by the indicators of orientation inversion and effective
sheet change. For a directed history
\(\gamma=e_1\cdots e_n\), define
\[
 c(\gamma)=\sum_{j=1}^{n}c(e_j).
\]
The concatenation law takes the form
\[
 c_{m+n}(x)
 =c_m(x)+c_n\!\left(F_{\rm obs}^{m}x\right).
\]
The previous counts express, uniformly over the orbit,
\[
 c_{27}(x)=(1,18),
 \qquad
 c_{108}(x)=(4,72).
\]

\begin{theorem}[Cumulative coverage of the observable chronology]
\label{thm:cobertura-acumulativa-dirigida}
On
\[
 \widetilde{\mathcal O}_{108}^{+}
 :=\mathcal O_{108}\times\mathbb N^2
\]
define
\[
 \widetilde F^{+}(x,m)
 :=\bigl(F_{\rm obs}x,m+c(x\to F_{\rm obs}x)\bigr).
\]
Then the transport of twenty-seven steps satisfies
\[
 \widetilde{\mathcal K}_{27}^{+}(x,m)
 =\bigl(\mathcal K_{27}x,m+(1,18)\bigr),
\]
and
\[
 \left(\widetilde{\mathcal K}_{27}^{+}\right)^4(x,m)
 =\bigl(x,m+(4,72)\bigr)\neq(x,m).
\]
Therefore, its observable projection has order four, whereas the
cumulative action of the chronological monoid has infinite semigroup order.
\end{theorem}

\begin{proof}
The cocyclic law and the constancy of \(c_{27}\) give
\[
 c_{108}(x)
 =\sum_{j=0}^{3}c_{27}(\mathcal K_{27}^{j}x)
 =4(1,18)=(4,72).
\]
The first coordinate returns through
\(\mathcal K_{27}^{4}=I\); the second does not return. After \(q>0\) iterations,
the increment is \(q(4,72)\), which is nonzero.
\end{proof}

\begin{remark}[Type of the cumulative extension]
The preceding construction extends directed histories through positive
counters. It is not an automorphic lift of groupoid
\(\mathcal G_{\rm mem}\), nor does it prove that the unbounded coordinate already belongs
to state \(\mathcal X_{\mathrm{TPK}}^{\mathrm{enr}}\). Its Grothendieck completion
\(\mathbb Z^2\) may be used for harmonic analysis, but it does not transform
the positive cost
\(\mathcal A_\ast=G_{\rm or}+\lambda_\ast S\) into a character of the isotropy
group: upon formally inverting a route, its indicators do not acquire
opposite signs.
\end{remark}

\paragraph{Domain of validity of the cumulative extension.}
The hierarchy \(27\to54\to108\), the incidences of \(F_{\rm av}\) and the preceding cocycles are exact identities in the declared chronology and readers. The assertion is restricted to the observable state those readers preserve; extending it to the complete integral holonomic state would require proving return of each forgotten field.

\section{Memory Groupoid}
\label{sec:grupoide-memoria-tpk}

Let \(\mathsf P_{\rm TPK}^{\rm enr}\) be the enriched path groupoid:
its objects are surviving TPK states, and its arrows are admissible routes
with declared source and target, composition by concatenation, and inversion
by route reversal. Let \(\sim_{\rm mem}\) be an equivalence relation that
identifies two arrows only when it preserves the memory fields required by
the reader under consideration and is congruent with identities, inversion, and
concatenation. We define
\[
 \mathcal G_{\rm mem}
 :=\mathsf P_{\rm TPK}^{\rm enr}/\!\sim_{\rm mem}.
\]
This definition does not transform memory into torsion. Memory belongs to
the discrete state and distinguishes histories; torsion belongs to a
geometric connection and measures \(D_\omega e\). Only a map between the two
domains can relate them.

A geometric realization requires, as additional data, a length
\(\ell_0>0\), a representation of the groupoid
\[
 \rho_{\rm C}:
 \mathcal G_{\rm mem}
 \longrightarrow
 \operatorname{Spin}(1,3)\ltimes\RR^{1,3},
\]
and a field map
\[
 \mathfrak R_{\rm grav}:
 \mathsf P_{\rm TPK}^{\rm enr}
 \dashrightarrow
 \mathsf{Cartan}_{1,3}.
\]
To define the second map, one must specify its action on objects and
routes, its compatibility with concatenation and inversion, and the regularity of the
resulting connection. None of these properties follows from the mere count
of returns.

\section{Joint connection and entanglement condition}
\label{sec:conexion-cartan-conjunta}

Under the geometric hypotheses of \cref{chap:gravedad} ---variedad
oriented of dimension four, co-tetrad \(e^I\), Lorentz connection
\(\omega^{IJ}\) and boundary conditions declared there---, the connection of
Cartan joint is written
\[
 \mathcal A_{\ell_0}
 =
 \begin{pmatrix}
  \omega&\ell_0^{-1}e\\
  0&0
 \end{pmatrix}.
\]
Its curvature is the algebraic identity.
\[
 \mathcal F_{\ell_0}
 =\dd\mathcal A_{\ell_0}
  +\mathcal A_{\ell_0}\wedge\mathcal A_{\ell_0}
 =
 \begin{pmatrix}
  F_\omega&\ell_0^{-1}T_{\rm tor}\\
  0&0
 \end{pmatrix},
\]
where
\[
 F_\omega^{IJ}
 =\dd\omega^{IJ}+\omega^I{}_K\wedge\omega^{KJ},
 \qquad
 T_{\rm tor}^{I}=D_\omega e^I.
\]
Curvature and torsion appear as distinct components of the same
transport defect. The matrix equality follows algebraically from the
explicit data \(e,\omega,\ell_0\), but does not construct these inputs from TPK.

The strong content of the bridge is the commutativity of holonomies:
\[
 \boxed{
 \Hol_{\mathcal A_{\ell_0}}
 \bigl(\mathfrak R_{\rm grav}(\gamma)\bigr)
 =
 \rho_{\rm C}\!\left(
 \Hol_{\rm TPK}(\gamma)
 \right)}
 \qquad
 (\gamma\in\mathsf P_{\rm TPK}^{\rm enr}).
\]
The equality must be compatible with identities, inversion, concatenation,
subdivision, change of base point, and conjugation. Only then is discrete
memory represented in a connection; it is not merely renamed as
torsion or curvature.

\begin{criterion}[Existence conditions for the bridge TPK--Cartan]
\label{crit:puente-tpk-cartan}
There is not realization with the declared commutativity and compatibility
if any of the following facts occurs:
\begin{enumerate}
 \item two equivalent representatives of the same route produce
 nonconjugate geometric holonomies;
 \item the image does not preserve identities, inverses, or concatenation;
 \item an admissible subdivision alters the holonomy;
 \item the equality of holonomies fails for \(\mathcal K_{27}\) or for their
 concatenations of depths \(54\) and \(108\);
 \item the map identifies memory and torsion without exhibiting a morphism that
 preserves their types;
 \item some descending choice ---the Holst action, \(G\), a mass or
 an SI chart--- is used to redefine \(A\), \(C^\ast\) or the original
 TPK return.
\end{enumerate}
\end{criterion}

\paragraph{Antecedence of the functional \(\Gamma_{6\to8}\) with respect to the Palatini--Cartan--Holst realization.}
\cref{mdg:chap:barbero} fixes, once \(A,C^\ast,q_\pm\) have been constructed, the internal functional \(\Gamma_{6\to8}\). \cref{chap:gravedad} then studies the Palatini--Cartan--Holst action and its variational equations. The present bridge introduces no inverse route between the two stages. \endgroup

% Owner íntegro de la realización discreta del retorno, la holonomía y la
% respuesta colectiva. Se sitúa después del puente que declara sus hipótesis.
\begin{HMTScientificOwnerSubordinated}
\chapter{Gravity: Holonomy, Torsion, and Collective Response}
\label{chap:gravedad}
\index{gravity}\index{Cartan--Holst}\index{holonomy}
\index{torsion}\index{graviton!non-primitive}

Gravity is not introduced here through a decimal formula for \(G\).
Its construction assembles two already typed branches. In the first, a route preserves
orientation and memory; a connection compares frames at distinct points; its
holonomy measures return and its curvature the closure defect. In the second,
the dimensional chart constructs \(G_{\rm int}(\theta_G)\) and
\(\kappa_{\rm int}\) from the internal action scale, velocity, and
elementary duration. Only after both branches are available is the
Palatini--Cartan--Holst action written. No external metrological figure participates in
that precedence.

\section{Returns: position, orientation, phase, and memory}

Let \(F_{\rm obs}\) be the observable chronology of the transport and
\[
 \mathcal K_{27}=F_{\rm obs}^{27}.
\]
The exact hierarchy distinguishes three returns:
\[
 \begin{array}{c|c}
 27\ \text{steps}&
 \text{position return and orientation reversal},\\
 54\ \text{steps}&
 \text{position and orientation return},\\
 108\ \text{steps}&
 \text{complete return of the observable phase}.
 \end{array}
\]
The observable return does not erase the accumulated record. If \(g(e)\) counts
an orientation reversal and \(s(e)\) an effective sheet change, then
\[
 c(\gamma)=\sum_{e\subset\gamma}(g(e),s(e))
\]
satisfies the cocyclic law of concatenation. In the canonical blocks,
\[
 c_{27}=(1,18),\qquad c_{108}=(4,72).
\]
Therefore, the observable projection of \(\mathcal K_{27}\) has order
four, whereas the cumulative coordinate continues to grow.  Returning to the
same phase does not mean returning to the same history.

\begin{theorem}[Cumulative coverage of return]
\label{thm:cobertura-retorno-grav-rev6}
Over \(\mathcal O_{108}\times\mathbb N^2\), the lift
\[
 \widetilde F(x,m)=
 \bigl(F_{\rm obs}x,m+c(x\to F_{\rm obs}x)\bigr)
\]
satisfies
\[
 \widetilde{\mathcal K}_{27}(x,m)
 =\bigl(\mathcal K_{27}x,m+(1,18)\bigr)
\]
and
\[
 \widetilde{\mathcal K}_{27}^{\,4}(x,m)
 =\bigl(x,m+(4,72)\bigr)\neq(x,m).
\]
\end{theorem}

\begin{proof}
The constancy of \(c_{27}\) along the orbit and the concatenation law give
\[
 c_{108}=\sum_{j=0}^{3}c_{27}(\mathcal K_{27}^jx)
 =4(1,18).
\]
The first coordinate returns; the second preserves the accumulated memory.
\end{proof}

This is the datum that to geometric realization must preserve.  One does not
identify the discrete memory with torsion: one constructs to representation
that allows them to is compared.

\section{From the path groupoid to the Cartan group}

Let \(\mathcal G_{\rm mem}\) be the groupoid of enriched routes, with
composition by concatenation and inverse by reversal. Fix an element
\(R_4\in\operatorname{Spin}^+(1,3)\) of order four and a unit temporal vector
\(u\in\mathbb R^{1,3}\) fixed by \(R_4\). If
\[
c(\gamma)=\bigl(c_g(\gamma),c_s(\gamma)\bigr)\in\mathbb Z^2
\]
is the additive cocycle of inversion and sheet change, we define
\[
\rho_{\rm C}^{\rm disc}(\gamma)
=
\left(
R_4^{\,c_g(\gamma)},
\ell_0c_s(\gamma)u
\right)
\in\operatorname{Spin}^+(1,3)\ltimes\mathbb R^{1,3}.
\]

\begin{theorem}[Discrete realization of the return cocycle]
\label{thm:realizacion-discreta-cartan-rev6}
\(\rho_{\rm C}^{\rm disc}\) is a functor from the route groupoid to the
Cartan group. In particular, it preserves identity, concatenation, and inversion, and
transports the return \(27\to54\to108\) without erasing the cumulative coordinate.
\end{theorem}

\begin{proof}
The semidirect product law is
\[
(R,a)(S,b)=(RS,a+Rb).
\]
Since \(R_4u=u\) and \(c\) is additive,
\[
\begin{aligned}
\rho_{\rm C}^{\rm disc}(\gamma_2)
\rho_{\rm C}^{\rm disc}(\gamma_1)
&=
\left(
R_4^{c_g(\gamma_2)+c_g(\gamma_1)},
\ell_0\bigl(c_s(\gamma_2)+c_s(\gamma_1)\bigr)u
\right)\\
&=\rho_{\rm C}^{\rm disc}(\gamma_2\gamma_1).
\end{aligned}
\]
The identity route has zero cocycle, and reversal changes its sign, so
identity and inverse are preserved. The return formulas follow from
\(c_{27}=(1,18)\) and \(c_{108}=4c_{27}\).
\end{proof}

This theorem constructs the discrete bridge. A subsequent smooth realization consists of a representation
\[
 \rho_{\rm C}:\mathcal G_{\rm mem}
 \longrightarrow\operatorname{Spin}(1,3)\ltimes\mathbb R^{1,3}
\]
and in fields \(e^I,\omega^{IJ}\) whose holonomy intertwines that representation.
The diagram fixing the type is
\[
 \operatorname{Hol}_{\mathcal A_{\ell_0}}
   \bigl(\mathfrak R_{\rm C}(\gamma)\bigr)
 =
 \rho_{\rm C}\bigl(\operatorname{Hol}_{\rm TPK}(\gamma)\bigr).
\]
Identities, inversion, concatenation, and subdivision must be preserved.  This
equation does not rename the carry as curvature: it requires the already represented
discrete transport and the geometric transport to realize the same composition.

Once to positive scale is fixed \(\ell_0\), the joint connection is
\[
 \mathcal A_{\ell_0}
 =
 \begin{pmatrix}
 \omega&\ell_0^{-1}e\\
 0&0
 \end{pmatrix}.
\]

\begin{proposition}[Cartan Joint Curvature]
\label{prop:curvatura-conjunta-cartan-rev6}
The curvature of \(\mathcal A_{\ell_0}\) is
\[
 \mathcal F_{\ell_0}
 =\dd\mathcal A_{\ell_0}
  +\mathcal A_{\ell_0}\wedge\mathcal A_{\ell_0}
 =
 \begin{pmatrix}
 F_\omega&\ell_0^{-1}T_{\rm tor}\\
 0&0
 \end{pmatrix},
\]
where
\[
 F_\omega^{IJ}
 =\dd\omega^{IJ}+\omega^I{}_K\wedge\omega^{KJ},
 \qquad
 T_{\rm tor}^{I}=D_\omega e^I.
\]
\end{proposition}

\begin{proof}
The matrices of forms are multiplied using the exterior product.  The
diagonal block produces \(F_\omega\); the translation block produces
\(\dd e+\omega\wedge e=D_\omega e\).
\end{proof}

Curvature and torsion are two components of the same transport defect,
but they perform distinct functions.  Curvature measures frame return;
torsion measures closure of the co-tetrad.

\section{Gauss--Stokes and areal capacity}

In a finite cellular complex, for every \(1\)-cochain \(\Omega\) and every
\(2\)-chain \(z\),
\[
 \boxed{\quad
 \langle\delta\Omega,z\rangle
 =\langle\Omega,\partial z\rangle.
 \quad}
\]
The identity is the definition of the coboundary as the adjoint of the boundary.
When \(z\) sums the coherently oriented faces of a
pseudomanifold, the interior edges cancel, and only the geometric
boundary remains.  The same equality has already been read as
vorticity--circulation and as curvature--holonomy.

The areal scale also has an exact model.  In the cubic graph
\(N\times N\times N\), every cut between two opposite faces contains at least
\(N^2\) edges, and to transverse layer attains the bound.  For \(N=9\),
\[
 9^3=729=27^2.
\]
This cardinal equality relates volume and minimal cut; it does not identify the
groups \((\mathbb Z/9)^3\) and \((\mathbb Z/27)^2\), which have distinct
exponents.

\section{The Palatini--Cartan--Holst Action}

Let \(M\) be a smooth oriented manifold of dimension four, with nondegenerate
co-tetrad \(e^I\), Lorentz connection
\(\omega^{IJ}=-\omega^{JI}\), and boundary conditions that make the boundary
variational term vanish.  Let us write
\[
 \Sigma^{IJ}=e^I\wedge e^J,\qquad
 (\star X)^{IJ}=\frac12\epsilon^{IJ}{}_{KL}X^{KL},
 \qquad
 G_{\rm int}=G_{\rm int}(\theta_G),\qquad
 \kappa_{\rm int}=\frac{8\pi G_{\rm int}}{c_{\rm int}^{4}}.
\]
These last two quantities are not defined through the action: they follow from
the constitutive construction of \cref{eq:G-interno}. The
\cref{thm:homogeneidad-pch} first proves that the following combination
belongs to the action line.
The first-order action is
\[
\boxed{
 S[e,\omega,\Psi]
 =\frac1{2\kappa_{\rm int}c_{\rm int}}\int_M
 \Sigma_{IJ}\wedge
 \left(\star F^{IJ}+\frac1\gamma F^{IJ}\right)
 -\frac{\Lambda_{\rm int}}{\kappa_{\rm int}c_{\rm int}}
  \int_M\operatorname{vol}_e
 +S_{\rm m}[e,\omega,\Psi].}
\]
The internal parameter constructed in the previous chapter specializes
\[
 \gamma=\Gamma_{6\to8}.
\]
The substitution is performed after obtaining
\(\Gamma_{6\to8}\) from \(A,C^\ast,q_\pm\) and the hexad--octad inclusion; the action is not used to redefine those inputs.

\begin{theorem}[Cartan--Holst Equation]
\label{thm:cartan-holst-rev6}
If the spin current is normalized through
\[
 \delta_\omega S_{\rm m}
 =-\frac12\int_M\delta\omega^{IJ}\wedge\sigma_{IJ},
\]
the variation with respect to \(\omega\) produces
\[
 \boxed{\quad
 D_\omega
 \left[(\star+\gamma^{-1}I)\Sigma^{IJ}\right]
 =\kappa_{\rm int}c_{\rm int}\,\sigma^{IJ}.
 \quad}
\]
\end{theorem}

\begin{proof}
One uses \(\delta F^{IJ}=D_\omega\delta\omega^{IJ}\), integrates by parts, and
sets the coefficient of the arbitrary variation equal to zero.  The cosmological
term does not depend on \(\omega\).
\end{proof}

In Lorentzian signature, \(\star^2=-I\) on bivectors.  For
\(\gamma\in\mathbb R\setminus\{0\}\),
\[
 (\star+\gamma^{-1}I)^{-1}
 =\frac{\gamma^2}{1+\gamma^2}
  (-\star+\gamma^{-1}I).
\]
Therefore, if the spin current vanishes,
\[
 T_{\rm tor}^I=0,\qquad
 \omega=\mathring\omega(e).
\]
With spinorial matter, torsion is determined algebraically by the
current.  In the minimal action, it does not constitute an independent
propagating degree of freedom.

\section{Metric equation, spin, and geometric balance}

To fix the normalization unambiguously, we defines the sources measured on
the action line by
\[
 \delta_g S_{\rm m}
 =-\frac12\int_M\operatorname{vol}_e\,
 \mathcal T^{S}_{\mu\nu}\,\delta g^{\mu\nu}.
\]
When the tetrad coordinates have the dimension of length, we write
\[
 T^{\rm phys}_{\mu\nu}:=c_{\rm int}\mathcal T^{S}_{\mu\nu},
 \qquad
 U^{\rm phys,spin}_{\mu\nu}
 :=c_{\rm int}\mathcal U^{S,\rm spin}_{\mu\nu}.
\]
After algebraically eliminating the connection, the metric equation adopts the two equivalent forms
\[
 \boxed{
 G_{\mu\nu}+\Lambda_{\rm int}g_{\mu\nu}
 =\kappa_{\rm int}c_{\rm int}
 \bigl(\mathcal T_{\mu\nu}^{S,\rm LC}
       +\mathcal U_{\mu\nu}^{S,\rm spin}\bigr)
 =\kappa_{\rm int}
 \bigl(T_{\mu\nu}^{\rm phys,LC}
       +U_{\mu\nu}^{\rm phys,spin}\bigr).}
\]
The spin contribution is quadratic in the corresponding current and
depends on the matter coupling; its coefficient is not fixed without declaring
the fermionic action and the duality conventions.

The geometric identities
\[
 D_\omega T_{\rm tor}^{I}=F^{I}{}_{J}\wedge e^{J},
 \qquad
 D_\omega F^{IJ}=0
\]
imply, after removing the torsion,
\[
 \nabla^\mu\bigl(T_{\mu\nu}^{\rm phys,LC}
             +U_{\mu\nu}^{\rm phys,spin}\bigr)=0.
\]
This Bianchi--Noether balance is to consistency condition for the
gravitational transducer: not additional term can is incorporated without
proceeding from to compatible action or satisfying the same balance.

\section{Area--volumetric correspondence \(\pi/\varphi^2\)}

The calendar \((\Sigma,\Pi,\Pi)\) produces the incidence matrix
\[
 F_{\rm av}=\begin{pmatrix}0&1\\1&1\end{pmatrix}.
\]
Its Parry measure gives
\(\mu_{\rm ar}/\mu_{\rm vol}=\varphi^{-2}\), and the regional reader marks the
closure exactly once \(\pi\).  The marked return satisfies
\[
 \Theta_n
 =\pi\frac{F_{n-1}}{F_{n+1}}
 \longrightarrow
 \frac{\pi}{\varphi^2}.
\]
This ratio is not another version of \(G\). It is the dimensionless correspondence between
the areal and volumetric readings. Under the same positive amplitude
\(\mathcal M\),
\[
 \mathcal I_{\rm av}=\mu_{\rm vol}\mathcal M,\qquad
 \mathcal G_{\rm av}=\pi\mu_{\rm ar}\mathcal M
\]
give exactly
\[
 \frac{\mathcal G_{\rm av}}{\mathcal I_{\rm av}}
 =\frac{\pi}{\varphi^2}.
\]
On the flat sheet, the diagonal restriction is
\(\mathcal G_{\rm tor}=\mathcal I_{\rm tor}\).  Local equivalence and
boundary ambivalence are therefore two regimes of the reader, not two
incompatible values of a mass.

\section{Typing of \(G\) as gravitational transducer and causal order of its readers}

The type and homogeneity of \(G\) belong to the prior constitutive
construction. The \cref{eq:G-interno} fixes the covariant family
\(G_{\rm int}(\theta_G)\) once and for all, and the
\cref{thm:homogeneidad-pch} proves that its insertion into the action has
the dimension of action. This chapter does not redefine that family: it uses its
section together with the holonomy already constructed. Causality is not linear,
but convergent:
\[
\begin{aligned}
 \mathcal B_{\rm geom}&:
 \text{return and cocycle}\to\text{Cartan representation}
 \to(e,\omega),\\
 \mathcal B_{\rm dim}&:
 (\hbar_{\rm int},c_{\rm int},t_0,\theta_G)
 \to G_{\rm int}\to\kappa_{\rm int},
\end{aligned}
\]
\[
 (\mathcal B_{\rm geom},\mathcal B_{\rm dim})
 \to S_{\rm PCH}
 \to\text{algebraic torsion, metric equation and balance}.
\]
The metrological comparison belongs to the end and does not select any of
the two branches.

\Needspace{22\baselineskip}
\section{Collective origin of the gravitational mode in the holonomy of paths}

The basic objects of the preceding action are the co-tetrad and the connection.
Cartan torsion is algebraic; curvature has the propagating sector that
contains gravitational waves. Accordingly, the formalism does not need to
postulate an elementary graviton in order to define the interaction. A possible
quantization of the perturbations may produce a collective mode of
helicities \(\pm2\); that mode would be an excitation of geometric
transport, not a datum added at the level of APP, TRIT, or TPK.

This distinction also excludes a frequent confusion. The existence of
classical waves proves propagation of curvature; by itself it does not prove
a fundamental quantum. Conversely, describing the field geometrically does not
eliminate its radiative degrees of freedom. HMT--MD locates the question properly:
memory and its holonomy are constructed first; the spectrum of their excitations
is then studied.
 \end{HMTScientificOwnerSubordinated}

\paragraph{Typed comparison of the area quantum with energy, entropy, and horizon information.} The first-order action allows the area quantum to be compared with energy, entropy and horizon information without merging these observables. \endgroup

\chapter{Convergence of horizon, area, information, entropy and action}
\begingroup
\renewcommand{\chapter}[2][]{}%
\chapter{Horizon–area–information–entropy–action confluence}
\label{chap:parte-iv-69}

The horizon patch coordinates energy, temperature, entropy, and action through typed identities. The half-action cell and chronological closure are presented with their Schwarzschild and de Sitter controls, without turning an interface into a retroactive definition.

\section{Horizon energy and half-action cell}
\label{rm:chap:horizonte-media-accion}

\subsection{Geometric conventions of the de Sitter patch}

Consider a four-dimensional de Sitter space with areal radius
\(R>0\), causal velocity \(c\), and gravitational constant \(G\). Adopt
the convention
\begin{equation}
 H_R=\frac cR,
 \qquad
 \Lambda_R=\frac3{R^2},
 \qquad
 \ell_P^2=\frac{G\hbar}{c^3}.
 \label{rm:eq:horizon-ds-conventions}
\end{equation}
The density of energy associated to the term cosmological is
\begin{equation}
 \varepsilon_\Lambda
 :=\frac{\Lambda_Rc^4}{8\pi G}.
 \label{rm:eq:horizon-vacuum-density}
\end{equation}
The spatially flat section of the patch contains, up to radius \(R\), the
volume
\begin{equation}
 V_R=\frac{4\pi}{3}R^3.
 \label{rm:eq:horizon-flat-volume}
\end{equation}
These conventions positively fix the quasilocal energy used in the
chapter. Choosing another sign for the de Sitter Killing charge
belongs to a distinct construction and must be translated before formulas
are compared.

\subsection{energy of vacuum and energy of Misner--Sharp}

In a spherically symmetric geometry, let
\begin{equation}
 \dd s^2=g_{ab}(x)\dd x^a\dd x^b
 +\mathcal R(x)^2\dd\Omega_2^2.
 \label{rm:eq:horizon-spherical-metric}
\end{equation}
The Misner--Sharp energy \cite{rm:MisnerSharp1964}, under the adopted positive convention, is
defined by
\begin{equation}
 E_{\rm MS}(\mathcal R)
 :=\frac{c^4\mathcal R}{2G}
 \left(1-g^{ab}\nabla_a\mathcal R\nabla_b\mathcal R\right).
 \label{rm:eq:horizon-Misner-Sharp}
\end{equation}
A marginal sphere satisfies
\begin{equation}
 g^{ab}\nabla_a\mathcal R\nabla_b\mathcal R=0.
 \label{rm:eq:horizon-marginal-sphere}
\end{equation}

\begin{theorem}[Confluence volumetrica and cuasilocal]
\label{rm:thm:horizon-vacuum-MS}
On the de Sitter horizon \(\mathcal R=R\), the integrated vacuum
energy and the Misner--Sharp energy coincide exactly:
\begin{equation}
 \boxed{
 E_\Lambda
 :=\varepsilon_\Lambda V_R
 =E_{\rm MS}(R)
 =\frac{c^4R}{2G}.}
 \label{rm:eq:horizon-E-Lambda}
\end{equation}
\end{theorem}

\begin{proof}
The equations
\cref{rm:eq:horizon-ds-conventions,rm:eq:horizon-vacuum-density,rm:eq:horizon-flat-volume}
give
\[
 \varepsilon_\Lambda V_R
 =\frac{3c^4}{8\pi GR^2}\frac{4\pi R^3}{3}
 =\frac{c^4R}{2G}.
\]
The marginal condition \cref{rm:eq:horizon-marginal-sphere} applied to
\cref{rm:eq:horizon-Misner-Sharp} yields the same result.
\end{proof}

The equality brings together two geometric publications. The first integrates a
volumetric density; the second evaluates the trapping defect of the
areal sphere. Both preserve the radius \(R\) and gravitational compliance
\(G/c^4\).

\subsection{Temperature, entropy and product of horizon}

The de Sitter surface gravity has magnitude \(c^2/R\). The Gibbons--Hawking
temperature \cite{rm:GibbonsHawking1977} and the Bekenstein--Hawking
entropy \cite{rm:Bekenstein1973,rm:Hawking1975} are
\begin{align}
 k_BT_H&=\frac{\hbar c}{2\pi R},
 \label{rm:eq:horizon-temperature}\\
 \frac{S_H}{k_B}
 &=\frac{A_H}{4\ell_P^2}
 =\frac{4\pi R^2}{4\ell_P^2}
 =\frac{\pi R^2}{\ell_P^2}.
 \label{rm:eq:horizon-entropy}
\end{align}

\begin{theorem}[de Sitter thermodynamic identity]
\label{rm:thm:horizon-TS}
For every radius \(R\) of the declared family,
\begin{equation}
 \boxed{
 T_HS_H=\frac{c^4R}{2G}=E_\Lambda.}
 \label{rm:eq:horizon-TS-E}
\end{equation}
\end{theorem}

\begin{proof}
Multiplying \cref{rm:eq:horizon-temperature,rm:eq:horizon-entropy} gives
\[
 T_HS_H
 =\frac{\hbar c}{2\pi Rk_B}
  \frac{k_B\pi R^2}{\ell_P^2}
 =\frac{\hbar cR}{2\ell_P^2}.
\]
The identity \(\ell_P^2=G\hbar/c^3\) turns the last member into
\(c^4R/(2G)\), equal to \cref{rm:eq:horizon-E-Lambda}.
\end{proof}

The factor \(1/2\) has a decomposition geometric transparent. The
areal factor \(4\pi\), the normalization holographic by four and the
periodicity thermal \((2\pi)^{-1}\) satisfy
\begin{equation}
 \frac{4\pi}{4}\frac1{2\pi}=\frac12.
 \label{rm:eq:horizon-half-factor}
\end{equation}
The half is therefore the joint contraction between areal quantization
and Euclidean periodicity.

\subsection{How complete of the horizon and Planck force}

We define the tension or fuerza of Planck and the complete quantum associated with the
radius by
\begin{equation}
 \mathscr F_P:=\frac{c^4}{G},
 \qquad
 \boxed{
 \mathcal Q_H(R):=2T_HS_H
 =\frac{c^4R}{G}=\mathscr F_P\,R.}
 \label{rm:eq:horizon-full-quantum}
\end{equation}
Quasilocal energy occupies half a cell of the complete quantum:
\begin{equation}
 \boxed{E_\Lambda=T_HS_H=\frac{\mathcal Q_H(R)}2.}
 \label{rm:eq:horizon-half-quantum-general}
\end{equation}

\begin{proposition}[Radial compliance]
\label{rm:prop:horizon-compliance}
The derivative of the complete quantum with respect to radius is the Planck force,
and the derivative of the quasilocal energy is half of it:
\begin{equation}
 \frac{\dd\mathcal Q_H}{\dd R}=\mathscr F_P,
 \qquad
 \frac{\dd E_\Lambda}{\dd R}=\frac{\mathscr F_P}{2}.
 \label{rm:eq:horizon-radial-derivatives}
\end{equation}
The compliance gravitatoria \(\mathscr C_P=G/c^4\) satisfies
\begin{equation}
 R=\mathscr C_P\mathcal Q_H.
 \label{rm:eq:horizon-compliance-law}
\end{equation}
\end{proposition}

\begin{proof}
The three identities are direct derivatives or inversions of
\cref{rm:eq:horizon-full-quantum}.
\end{proof}

The equation \cref{rm:eq:horizon-compliance-law} expresses the radial response:
the complete quantum applies tension, and geometry publishes a radius through
universal compliance. The product \(T_HS_H\) occupies the quasilocal half
of that response.

\subsection{Specialization to the closure chronological of 108 states}

The chronological closure of \cref{rm:ch:cosmo} determines
\begin{equation}
 R_{108}=\ell_P,
 \qquad
 t_P=\frac{\ell_P}{c},
 \qquad
 T_{108}=2\pi t_P,
 \label{rm:eq:horizon-CT108-scales}
\end{equation}
and the Planck quantum
\begin{equation}
 E_P=\frac{\hbar}{t_P}
 =k_B\Theta_{\rm clk}\log3.
 \label{rm:eq:horizon-Planck-trit}
\end{equation}

\begin{theorem}[Cell CT108 of half action]
\label{rm:thm:horizon-half-action}
In the specialization \(R=R_{108}=\ell_P\),
\begin{align}
 \mathcal Q_H(\ell_P)
 &=\frac{c^4\ell_P}{G}=E_P,
 \label{rm:eq:horizon-QH-EP}\\
 \boxed{
 E_\Lambda=T_HS_H
 =\frac{E_P}{2}
 =\frac{k_B\Theta_{\rm clk}\log3}{2},}
 \label{rm:eq:horizon-CT108-half-energy}\\
 \boxed{E_\Lambda t_P=\frac{\hbar}{2},}
 \qquad
 \boxed{E_\Lambda T_{108}=\frac h2.}
 \label{rm:eq:horizon-half-actions}
\end{align}
\end{theorem}

\begin{proof}
The first equality is obtained by substituting
\(\ell_P=\sqrt{G\hbar/c^3}\) and
\(E_P=\sqrt{\hbar c^5/G}\). The
\cref{rm:eq:horizon-half-quantum-general} produces
\(E_\Lambda=E_P/2\), and
\cref{rm:eq:horizon-Planck-trit} supplies the informational chart. Finally,
\[
 E_\Lambda t_P=\frac{E_Pt_P}{2}=\frac\hbar2
\]
and, as \(T_{108}=2\pi t_P\) and \(h=2\pi\hbar\),
\[
 E_\Lambda T_{108}
 =\frac{E_P}{2}(2\pi t_P)=\pi\hbar=\frac h2.
\]
\end{proof}

The first equality in \cref{rm:eq:horizon-half-actions} publishes half a reduced action
over Planck time; the second publishes half a complete action
over the CT108 period. Both express the same cell in the angular
and circular charts.

The temperature and entropy are also determined:
\begin{equation}
 k_BT_H=\frac{E_P}{2\pi},
 \qquad
 \boxed{\frac{S_H}{k_B}=\pi.}
 \label{rm:eq:horizon-CT108-TS}
\end{equation}
The temperature converts the frequency of horizon in energy; the
entropy converts the area of the sphere of Planck in information. Their
product recovers the same half-quantum.

\subsection{Effective horizon weight and hyperbolic exponential}

We define the effective thermodynamic weight
\begin{equation}
 \Omega_{\rm eff}(R):=\exp\!\left(\frac{S_H(R)}{k_B}\right).
 \label{rm:eq:horizon-effective-weight}
\end{equation}
In the CT108 set,
\begin{equation}
 \boxed{\Omega_{\rm eff}(\ell_P)=\ee^\pi.}
 \label{rm:eq:horizon-effective-e-pi}
\end{equation}
The weight \(\Omega_{\rm eff}\) is a positive quantity defined by
entropy. Its real value expresses effective multiplicity; an interpretation
as an integer cardinality of microstates requires an additional discretization.

The areal--volumetric incidence operator is
\begin{equation}
 F_{\rm av}=\begin{pmatrix}0&1\\1&1\end{pmatrix},
 \qquad
 J_\varphi:=\frac{2F_{\rm av}^2-3I}{\sqrt5}.
 \label{rm:eq:horizon-J-phi}
\end{equation}
A direct calculation gives
\begin{equation}
 J_\varphi^2=I,
 \qquad
 \Spec(J_\varphi)=\{+1,-1\}.
 \label{rm:eq:horizon-J-involution}
\end{equation}
The hyperbolic flow at closure time \(\pi\) is
\begin{equation}
 \mathcal E_\pi
 :=\ee^{\pi J_\varphi}
 =\cosh\pi\,I+\sinh\pi\,J_\varphi,
 \label{rm:eq:horizon-hyperbolic-flow}
\end{equation}
and possesses spectrum
\begin{equation}
 \boxed{
 \Spec(\mathcal E_\pi)=\{\ee^\pi,\ee^{-\pi}\},
 \qquad
 \rho(\mathcal E_\pi)=\ee^\pi.}
 \label{rm:eq:horizon-hyperbolic-spectrum}
\end{equation}

\begin{corollary}[Horizon--Perron scalar confluence]
\label{rm:cor:horizon-Perron-scalar}
In the CT108 closure,
\begin{equation}
 \boxed{
 \frac{S_H}{k_B}
 =\log\Omega_{\rm eff}
 =\log\rho(\mathcal E_\pi)
 =\pi.}
 \label{rm:eq:horizon-Perron-log}
\end{equation}
\end{corollary}

\begin{proof}
The equations
\cref{rm:eq:horizon-CT108-TS,rm:eq:horizon-effective-e-pi,rm:eq:horizon-hyperbolic-spectrum}
give the three members.
\end{proof}

The scalar identity is exact and preserves two distinct genealogies. The horizon produces \(\ee^\pi\) through exponentiation of area entropy; the previously constructed transfer dynamics recognizes the same scalar as the expanding eigenvalue of a hyperbolic involution. Terminal equality fixes a boundary condition for the intertwiner and preserves the distinction of domains.

\subsection{Consistency of the Gravitational Transducer with the Schwarzschild Horizon}

A Schwarzschild horizon with the same area radius satisfies
\begin{equation}
 R_H=\frac{2GE_H}{c^4},
 \qquad
 E_H=\frac{c^4R_H}{2G}.
 \label{rm:eq:horizon-Schwarzschild-energy}
\end{equation}
The entropy areal coincides with \cref{rm:eq:horizon-entropy}, whereas the
temperature is
\begin{equation}
 k_BT_{\rm Sch}=\frac{\hbar c}{4\pi R_H}
 =\frac12k_BT_H\big|_{R=R_H}.
 \label{rm:eq:horizon-Schwarzschild-temperature}
\end{equation}
Consequently,
\begin{equation}
 E_H=2T_{\rm Sch}S_H,
 \qquad
 T_H^{\rm dS}=2T_{\rm Sch}
 \label{rm:eq:horizon-Schwarzschild-Smarr}
\end{equation}
for the same radius. In \(R_H=\ell_P\),
\begin{equation}
 E_H=\frac{E_P}{2},
 \qquad
 T_{\rm Sch}S_H=\frac{E_P}{4}.
 \label{rm:eq:horizon-Schwarzschild-Planck}
\end{equation}
The contrast separates the two geometries: they share quasilocal energy,
area, and entropy, and their surface gravities produce temperatures
in the ratio two.

With \(\beta=(k_BT_{\rm Sch})^{-1}\), the action Euclidean satisfies
\begin{equation}
 \frac{I_E}{\hbar}
 =\frac{E_H}{k_BT_{\rm Sch}}-\frac{S_H}{k_B}
 =2\pi-\pi=\pi
 \label{rm:eq:horizon-Schwarzschild-action}
\end{equation}
at the Planck radius. The same \(\pi\) appears as reduced entropy,
reduced Euclidean action, hyperbolic-flow time, and the logarithm of the
effective weight.

\subsection{Horizontal Horizon–Perron Entangler: Spectral compatibility between boundary energy and return memory}

The coincidence
\(\Omega_{\rm eff}=\rho(\mathcal E_\pi)\) fixes a spectral datum and leaves
open the relation between the complete objects. Let
\(\mathcal A_H^{(m)}\) denote the algebra of horizon observables at nonadic level
\(m\), let \(\omega_H^{(m)}\) denote its state, and let
\(\mathcal P_\varphi^{(m)}\) denote the Perron transfer system.

\begin{definition}[Thermodynamic--dynamic intertwiner]
\label{rm:def:horizon-Perron-intertwiner}
An admissible intertwiner is a family of unital completely positive
linear maps (UCP)
\begin{equation}
 \mathfrak I_m:\mathcal A_H^{(m)}\longrightarrow
 \mathcal A_\varphi^{(m)}
 \label{rm:eq:horizon-intertwiner}
\end{equation}
which satisfies:
\begin{enumerate}[label=\roman*)]
 \item preservation of the state, in
addition to unitality and complete positivity already
incorporated into the UCP condition;
 \item compatibility with nonadic refinement;
 \item transport of the involution interior--exterior toward
       \(J_\varphi\mapsto-J_\varphi\);
 \item commutation between the horizon modular semigroup and the transfer
       semigroup on the declared subspace;
 \item equality of the spectral radius at time \(\pi\), given by
       \cref{rm:eq:horizon-Perron-log}.
\end{enumerate}
\end{definition}

The fifth condition is necessary and scalar; the first four provide the
algebraic content that a numerical coincidence does not contain. An
intertwiner constructed in this manner would identify the effective
thermodynamic multiplicity with a dynamics of expansion and contraction
while preserving orientation, state, and refinement.

\begin{proofstatusbox}
The identities
\(E_\Lambda=E_{\rm MS}=T_HS_H=c^4R/(2G)\), the definition
\(\mathcal Q_H=\mathscr F_P\,R\), the CT108 specialization, the two
half-action cells, and the equality
\(S_H/k_B=\log\rho(\ee^{\pi J_\varphi})=\pi\) are exact results under
the declared conventions and de Sitter realization. The identification
of \(\Omega_{\rm eff}\) with a microscopic cardinal and the intertwiner
\(\mathfrak I_m\) belong to a delimited physical interface.
\end{proofstatusbox}

\enlargethispage{4\baselineskip}
\begin{falsifierbox}
Geometric closure fails if \(\Lambda_R\ne3/R^2\), if the sphere of
radius \(R\) fails to satisfy the marginal condition, or if the integrated energy differs from
\(E_{\rm MS}(R)\). Thermodynamic closure fails if the temperature or the
area law produces \(T_HS_H\ne c^4R/(2G)\). The CT108 specialization
fails if \(R_{108}\ne\ell_P\), \(T_{108}\ne2\pi t_P\), or
\(E_P\ne k_B\Theta_{\rm clk}\log3\). The Perron interface is refuted by
an obstruction of positivity, state, orientation, or refinement, even
when the eigenvalue \(\ee^\pi\) agrees scalarly.
\end{falsifierbox}
 
\paragraph{Machian Factorization of the Transport Memory from the Global Closure and the Local Response.} The relation between global closure and local response leads to Mach's principle formulated on transport memory. \endgroup

\let\HMTMonographChapterInput\relax
  \part{Boundary incidence and exceptional stabilizers}

\chapter{Dodecaphasic--Witt intertwiner}
\section{Construction of the point--hexad incidence intertwiner and transport of the residual action of \texorpdfstring{\(S_3\)}{S3}}
\label{chap:entrelazador-accion}
\index{intertwiner!dodecaphase--Witt}
\index{Witt!incidence module}
\index{survival!rank-three projector}

The repeated appearance of quotient \(3/11\) on the dodecaphasic screen,
in the polarity of a Witt star, and in nonadic filtration raises
a question stronger than the arithmetic equality of three fractions. The
correct question is operatorial: to determine whether those three appearances are
traces of the same positive sector viewed in different realizations and, if they
are, to construct the maps intertwining them without introducing a subsequent
metrological datum.

The answer has four levels. First, point--hexad incidence
constructs a canonical isometry between the eleven-dimensional
dodecaphase module and an eleven-dimensional eigenspace of the Witt design. Second, the
Witt star acts on that same module, but does so as a signed involution;
its \(3/11\) is an index, not the trace of a positive projector.
Third, the oriented survival sector appears in \(t=7\), where a
single eleven-branch space contains a rank-three sector. Fourth, the
orientation line of that triplet, together with the internal witnesses
\(B_K\) and \(u_K\), variationally selects a unique normalizer
\(S_3\), fixes its orientation, and, through Reynolds averaging and
polar decomposition, produces the sector--dodecaphase isometric intertwining. The
dodecaphase--Witt passage is a second incidence isometry; it is not presupposed
that both arrows realize a single representation of \(S_3\).
The Reynolds averages, the decompositions into representations, and the
polar factor used in this section are the standard ones
\cite{Serre1977Representations,HornJohnson2013}.

\begin{figure}[htbp]
\centering
\begin{tikzpicture}[
  font=\footnotesize,
  box/.style={draw,rounded corners=2mm,align=center,minimum width=2.35cm,minimum height=.98cm},
  arr/.style={-{Stealth[length=2.2mm]},thick},
  link/.style={<->,thick,dashed}
]
  \node[box,draw=hmtblue,fill=hmtlightblue] (seed) at (0,2.2)
    {semillas\\\(104\,976\) translated};
  \node[box,draw=hmtblue,fill=hmtlightblue] (u) at (3,2.2)
    {emissions\\\(\mathcal U_6,\ |\mathcal U_6|=468\)};
  \node[box,draw=hmtgold,fill=hmtlightgold] (ppi) at (6,2.2)
    {microfibra\\\(\Pi_\pi^{(468)},\ \operatorname{rank}=5\)};
  \node[box,draw=hmtred,fill=hmtlightred] (d) at (9,2.2)
    {positive defect\\\(D_{\rm switch}\)};
  \node[box,draw=hmtviolet,fill=hmtlightviolet] (lam) at (12,2.2)
    {capacidad residual\\\(\lambda_*=1-\delta_*\)};

  \draw[arr,hmtblue] (seed)--node[above=16pt,font=\scriptsize,fill=white,inner sep=1pt]{\(F\)} (u);
  \draw[arr,hmtgold] (u)--node[above=16pt,font=\scriptsize,fill=white,inner sep=1pt]{\(\tau_{468}\)} (ppi);
  \draw[arr,hmtred] (ppi)--node[above=16pt,font=\scriptsize,fill=white,inner sep=1pt]{\(D_{\rm switch}\)} (d);
  \draw[arr,hmtviolet] (d)--node[above=16pt,font=\scriptsize,fill=white,inner sep=1pt]{\(\tau+\mathrm{CR}\)} (lam);

  \node[box,draw=hmtgold,fill=hmtlightgold] (g) at (1.5,-.2)
    {sector oriented \(t=7\)\\\(L_{\rm or}\otimes(H_G,P_G)\)};
  \node[box,draw=hmtgreen,fill=hmtlightgreen] (v) at (5,-.2)
    {dodecaphase\\\(V_{11},P_3\)};
  \node[box,draw=hmtgreen,fill=hmtlightgreen] (e) at (8.5,-.2)
    {Witt incidence\\\(E_{30},Q_3\)};
  \node[box,draw=hmtgray,fill=hmtpaper] (t) at (12,-.2)
    {reader spectral\\\(\mathcal L_{\beta,*}\to T_{\rm band}\)};

  \draw[arr,hmtgold] (g)--node[above=16pt,font=\scriptsize,fill=white,inner sep=1pt]{\(W_{GK}\)} (v);
  \draw[arr,hmtgreen] (v)--node[above=16pt,font=\scriptsize,fill=white,inner sep=1pt]{\(U_W=I/6\)} (e);
  \draw[arr,hmtgray] (lam)--(t);
  \draw[arr,hmtred] (v) to[bend left=20]
    node[below=3pt,font=\scriptsize,fill=white,inner sep=1pt]{\((3/11)\Delta_4\)} (d);
\end{tikzpicture}
 \caption{Architecture of the predimensional closure. The upper row constructs
the measure and the action from the quotient of emissions; the lower row transports
the oriented sector \(t=7\) to the dodecaphase and then the dodecaphase to
Witt. The solid arrows are proved projector isometries. The
first is fixed by the variational selection of
\cref{thm:identificacion-isometrica-tres-proyectores}; for the second, no
equivariance with respect to the selected \(H_*\) is asserted.}
\label{fig:entrelazador-accion}
\end{figure}

\section{Typing of the objects associated with \texorpdfstring{\(3/11\)}{3/11}}

Let
\[
 V_{12}=\mathbb R^{12},
 \qquad
 V_{11}=\mathbf 1^\perp
 =\left\{v\in\mathbb R^{12}:\sum_{p=1}^{12}v_p=0\right\}.
\]
The icosahedral action declared on the twelve positions decomposes as
\[
 V_{12}=V_1\oplus V_3\oplus V_{3'}\oplus V_5.
\]
Denotemos by \(P_3\) the orthogonal projector on the first summand
three-dimensional. Then
\[
 P_3^2=P_3=P_3^*,
 \qquad
 \operatorname{rank}P_3=3,
 \qquad
 \tau_{11}(P_3)
 =\frac1{11}\operatorname{Tr}_{V_{11}}P_3
 =\frac3{11}.
\]

There are three nearby objects and a fourth object of a different type:
\begin{enumerate}
  \item the point--hexad incidence transports \(P_3\) to a rank-three
  projector \(Q_3\) on a Witt module;
  \item survival in \(t=7\) distinguishes three branches among eleven and
  produces a projector \(P_G\) with normalized trace \(3/11\);
  \item an oriented intertwiner identifies
  \(P_3V_{11}\) with the oriented realization of \(P_GH_G\);
  \item a Witt star has seven even modes and four odd ones, so
  its normalized trace also equals \(3/11\), but it contains negative spectrum
  and is not a projector.
\end{enumerate}
The distinction between projector and involution will be essential to the positivity of the action law in the following section.

\section{The point--hexad incidence matrix}

Let \(\mathcal H\) be the set of the \(132\) hexads of the design
\[
 W_{12}=S(5,6,12).
\]
We define
\[
 I:\mathbb R^{12}\longrightarrow\mathbb R^{\mathcal H},
 \qquad
 (Iv)(H)=\sum_{p\in H}v_p,
\]
or, equivalently,
\[
 I_{H,p}=\mathbf 1_{\{p\in H\}}.
\]
This operator records only incidence. It contains no \(\alpha\), units,
masses, CKM, Barbero, or any physical target.

\begin{lemma}[Metric Identity of Incidence]
\label{lem:incidencia-metrica-witt}
We have
\[
 \boxed{I^*I=36I_{12}+30J_{12}.}
\]
Therefore, for every \(v\in V_{11}\),
\[
 \|Iv\|^2=36\|v\|^2.
\]
\end{lemma}

\begin{proof}
In an \(S(5,6,12)\) system, each point belongs to
\[
 r=\frac{\binom{11}{4}}{\binom54}=66
\]
hexads, and each pair of points belongs to
\[
 \lambda_2=\frac{\binom{10}{3}}{\binom43}=30
\]
hexads. The diagonal entries of \(I^*I\) have value \(66\) and the
nondiagonal entries have value \(30\). Hence
\[
 I^*I=(66-30)I_{12}+30J_{12}.
\]
Since \(J_{12}\) vanishes on \(\mathbf1^\perp\), the second
identity is obtained.
\end{proof}

To identify the spectral branch selected by \(I\), we introduce
\[
 G_{H,H'}=\binom{|H\cap H'|}{4}.
\]

\begin{lemma}[Spectral identity of incidence]
\label{lem:incidencia-espectral-witt}
With \(J_{132,12}\) denoting the rectangular all-ones matrix,
\[
 \boxed{GI=30I+15J_{132,12}.}
\]
In particular, if \(v\in V_{11}\), then
\[
 G(Iv)=30Iv.
\]
\end{lemma}

\begin{proof}
The entry \((H,p)\) of \(GI\) is
\[
 \sum_{H'\ni p}\binom{|H\cap H'|}{4}.
\]
The exhaustive census of \(W_{12}\) gives \(45\) when \(p\in H\) and \(15\)
when \(p\notin H\). Consequently, the entry equals
\(30I_{H,p}+15\). Reconstruction of the \(132\) hexads from the extended ternary
Golay code checks the \(132\cdot12\) entries as
integers.
\end{proof}

Define
\[
 E_{30}^{\rm inc}:=I(V_{11})\subseteq\ker(G-30I).
\]
Its dimension is eleven by the \cref{lem:incidencia-metrica-witt}. The calculation
modulo \(1009\) obtains
\[
 \operatorname{rank}_{\mathbb F_{1009}}(G-30I)=121.
\]
Since reduction modulo \(1009\) gives a lower bound for the rational
rank, it follows that
\(\dim_{\mathbb Q}\ker(G-30I)\le11\). The preceding inclusion provides
eleven linearly independent vectors in that kernel. Therefore,
\[
 E_{30}^{\rm inc}=\ker(G-30I).
\]
From now on we will write \(E_{30}\).

\begin{theorem}[Canonical dodecaphase entangler--Witt]
\label{thm:entrelazador-dodecafase-witt}
The operator
\[
 \boxed{
 U_W=\frac16 I\big|_{V_{11}}:V_{11}\longrightarrow E_{30}
 }
\]
is a isometry surjective. For every
\(g\in\operatorname{Aut}(W_{12})\cong M_{12}\),
\[
 U_W\rho_{12}(g)=\rho_{\mathcal H}(g)U_W.
\]
Among the equivariant point--hexad kernels, the isometry condition and
the positive sign on incidences uniquely determine \(U_W\).
\end{theorem}

\begin{proof}
The isometry follows from \(I^*I=36I\) over \(V_{11}\). Its image has
dimension eleven and is contained in \(E_{30}\), which also has dimension eleven;
it is therefore surjective. Equivariance translates
\[
 p\in H\quad\Longleftrightarrow\quad g(p)\in g(H).
\]

For uniqueness, \(M_{12}\) is transitive on the flags
\((p,H)\), with \(p\in H\), and on the antiflags. Every equivariant kernel
has two values, \(a\) on flags and \(b\) on
antiflags:
\[
 T=aI+b(J-I)=(a-b)I+bJ.
\]
The term \(J\) vanishes on \(V_{11}\), so there \(T\) is a
multiple of \(I\). The isometry forces \(|a-b|=1/6\) and the positive
normalization selects \(a-b=1/6\).
\end{proof}

\section{Positive transport of the tritic projector}

The entangler allows defining
\[
 \boxed{Q_3=U_WP_3U_W^*.}
\]
We obtain
\[
 Q_3^2=Q_3=Q_3^*,
 \qquad
 \operatorname{rank}Q_3=3,
\]
and the diagram
\[
\begin{CD}
V_{11} @>{U_W}>> E_{30}\\
@V{P_3}VV @VV{Q_3}V\\
V_{11} @>{U_W}>> E_{30}
\end{CD}
\]
commutes. In particular,
\[
 \boxed{\tau_{V_{11}}(P_3)=\tau_{E_{30}}(Q_3)=\frac3{11}.}
\]

The canonicity of \(U_W\) belongs to the labeled design and its
automorphism action. It must be distinguished from the screen \(A_5/C_5\) used to
define \(P_3\): in the constructed labeled realization, only the identity
of that concrete action of \(A_5\) preserves the fixed set of hexads. Hence
\(U_W\) canonically transports the already given \(P_3\) toward
\(Q_3\), but does not by itself convert the action \(A_5/C_5\) into a
subaction of \(M_{12}\). The bridge to the oriented sector will be constructed separately
from the marked data \(K\) and \(B_K\).

\section{Signed index of the Witt star}

Let \(B\) be a tetrad. The four hexads that contain it are written
\[
 H_i=B\sqcup P_i,\qquad |P_i|=2,
\]
where the four petals \(P_i\) partition the complement of \(B\).
The star involution \(R_B\) fixes \(B\) and exchanges the two points
of each petal. Its cycle type on twelve points is \(1^4\,2^4\). In
\(V_{11}\), and by means of \(U_W\) also in \(E_{30}\), its spectrum is
\[
 \operatorname{Spec}(R_B)=\{1^{[7]},(-1)^{[4]}\}.
\]
Therefore,
\[
 \boxed{\tau_{11}(R_B)=\frac{7-4}{11}=\frac3{11}.}
\]

\begin{theorem}[Obstruction projector--star]
\label{thm:no-go-proyector-estrella}
There is no isomorphism \(F:V_{11}\to E_{30}\) such that
\[
 FP_3=R_BF.
\]
\end{theorem}

\begin{proof}
If it existed, \(P_3\) and \(R_B\) would be similar. However,
\[
 \operatorname{Spec}(P_3)=\{1^{[3]},0^{[8]}\},
 \qquad
 \operatorname{Spec}(R_B)=\{1^{[7]},(-1)^{[4]}\}.
\]
Their minimal polynomials are \(x(x-1)\) and
\((x-1)(x+1)\), respectively. Similarity is impossible.
\end{proof}

The equality \(3/11\) thus has two readings on the same module. For
\(Q_3\) it is the mass of a positive rank-three sector. For \(R_B\) it is
a parity index. A positive action must involve \(Q_3\); the
star preserves orientation and distinguishes sheets but does not replace the
projector.

\section{homogeneous reduction \texorpdfstring{\(11\to3\to1\)}{11-3-1} in \texorpdfstring{\(t=7\)}{t=7}}

The preceding reading of \(3/11\) in \(t=8\) compared three local roots
with eleven continuations of a single root. The numerator and denominator did not
belong to the same space. The detailed table locates the correct object
in \(t=7\): there are eleven local branches; all eleven survive one horizon, three
survive two, and one survives three,
\[
 \boxed{11\longrightarrow3\longrightarrow1.}
\]

\begin{center}
\small
\begin{tabularx}{\textwidth}{cYcc}
\toprule
index&triple of words&\(h=2\)&\(h=3\)\\
\midrule
0&\texttt{200112|212221|110110}&3&0\\
1&\texttt{200121|212212|110110}&3&1\\
2&\texttt{200211|212122|110110}&3&0\\
3&\texttt{202122|210211|110110}&0&0\\
4&\texttt{202212|210121|110110}&0&0\\
5&\texttt{202221|210112|110110}&0&0\\
6&\texttt{210111|202222|110110}&0&0\\
7&\texttt{210222|202111|110110}&0&0\\
8&\texttt{212112|200221|110110}&0&0\\
9&\texttt{212121|200212|110110}&0&0\\
10&\texttt{212211|200122|110110}&0&0\\
\bottomrule
\end{tabularx}
\end{center}

Let \(\mathcal B_7\) be the set of these eleven branches and
\[
 H_G=\ell^2(\mathcal B_7).
\]
We define
\[
 P_G\delta_b=
 \begin{cases}
 \delta_b,&b\text{ admits continuation to horizon two},\\
 0,&\text{otherwise}.
 \end{cases}
\]
Then
\[
 P_G^2=P_G=P_G^*,
 \qquad
 \operatorname{rank}P_G=3,
 \qquad
 \boxed{\tau_{H_G}(P_G)=\frac3{11}.}
\]
This is the third positive realization.

\section{The residual action of \texorpdfstring{\(S_3\)}{S3} and the positive regions of \texorpdfstring{\(\pi\)}{pi}}

A permutation of the three final positions, applied simultaneously to
the first two words of each branch and leaving the third fixed, preserves
\(\mathcal B_7\). This rule defines an action of \(S_3\). Its orbits
have sizes
\[
 3+3+1+1+3=11.
\]
The support of \(P_G\), formado by the indices \(0,1,2\), is one orbit
complete, so that \(P_G\) conmuta with \(S_3\).

The character of the representation on the eleven branches, evaluated at the identity,
a transposition, and a cycle of order three, is
\[
 \chi_{\mathcal B_7}=(11,5,2)
 =5\chi_{\rm triv}+3\chi_{\rm std}.
\]
On the image of \(P_G\),
\[
 \chi_G=(3,1,0)
 =\chi_{\rm triv}+\chi_{\rm std}.
\]

The three positive regions of the microfibre of \(\pi\) have signatures
\[
 339555,\qquad393555,\qquad933555.
\]
The tails of the first word of the three surviving branches are
\[
 112,\qquad121,\qquad211.
\]
The coordinate substitution \(1\mapsto3\), \(2\mapsto9\) gives
\[
\begin{array}{c|c|c}
\text{branch}&\text{tail}&\text{positive region}\\ \hline
0&112&339555\\
1&121&393555\\
2&211&933555.
\end{array}
\]

\begin{theorem}[Horizontal--vertical bridge of the positive microfiber]
\label{thm:puente-pi-puerta-s3}
The preceding rule is an \(S_3\)-equivariant bijection between the three
positive regions of \(\pi\) and the three branches selected by \(P_G\).
The map preserves the position of the unique exceptional symbol.
\end{theorem}

\begin{proof}
Both sets are the three words obtained by placing a single exceptional
symbol in one of three positions. The substitution is applied coordinate by
coordinate and therefore commutes with every permutation of positions. The
exceptional position uniquely recovers the source element.
\end{proof}

\section{Dependence of orientation in the dodecaphase entangler--Witt}

For any normalizer \(H=N_{A_5}(C_3)\simeq S_3\) of a subgroup of
order three in \(A_5\), we define the abstract character of the quotient
\[
 \operatorname{sgn}_H:H\longrightarrow H/C_3\simeq C_2
 \longrightarrow\{+1,-1\},
 \qquad
 \operatorname{sgn}_H(h)=
 \begin{cases}
 +1,&h\in C_3,\\
 -1,&h\in H\setminus C_3.
 \end{cases}
\]
This character is not the ambient sign of a permutation of five letters:
every element of \(A_5\) has ambient sign \(+1\). It is the
intrinsic character of the reflection in the quotient \(H/C_3\). With this convention, the
character of the icosahedral sector restricted to \(H\) is
\[
 \chi_{P_3V_{11}|H}=(3,-1,0)
 =\chi_{\operatorname{sgn}_H}+\chi_{\rm std}.
\]
By contrast,
\[
 \chi_{\operatorname{im}P_G}=(3,1,0)
 =\chi_{\rm triv}+\chi_{\rm std}.
\]

\begin{corollary}[Obstruction to the unoriented braider]
\label{cor:obstruccion-entrelazador-incidencia-p3}
There is no \(S_3\)-equivariant isomorphism
\[
 \operatorname{im}P_G\longrightarrow P_3V_{11}
\]
for the non-twisted actions.
\end{corollary}

\begin{proof}
Isomorphic representations have the same character. Their values on a
transposition are \(1\) and \(-1\).
\end{proof}

Let \((e_0,e_1,e_2)\) the basis posicional of
\(\mathbb R^3_{\rm perm}\). We defines
\[
 L_{\rm or}
 =\bigwedge\nolimits^3\mathbb R^3_{\rm perm},
 \qquad
 \ell_{\rm or}=e_0\wedge e_1\wedge e_2
\]
the orientation line of the triplet, equipped with the positive anchor
\(\ell_{\rm or}\). The action of \(S_3\) on this line is precisely its
abstract sign character. Upon tensoring by it,
\[
 \chi_{L_{\rm or}\otimes\operatorname{im}P_G}
 =(3,-1,0),
\]
of mode that the character obstruction disappears. The equality of
caracteres proves the existence of one class of isomorphisms, but aun not
selects one. The selection is obtained of two witnesses already present in the
dodecaphase.

\section{Variational selection of the normalizer \texorpdfstring{\(S_3\)}{S3}}

Let
\[
 K=(234,543,140,729,659,824,621,58,914,794,146,601)
\]
the decimal chart of the seal. The common minimal witness of the two readings is
\[
 B_K=H_e\cap P_\pi
 =\{p:K_p\ge729\}
 =\{4,6,9,10\},
\]
in the mathematical numbering \(1,\ldots,12\). En the equivalent convention
with zero-based indices, the same set is \(\{3,5,8,9\}\). We define two
vectors of the sector \(P_3\):
\[
 b_K=P_3\left(\mathbf1_{B_K}-\frac{|B_K|}{12}\mathbf1\right),
 \qquad
 u_K=P_3\left(K-\overline K\,\mathbf1\right).
\]
The first satisfies exactly.
\[
 \|b_K\|^2=1.
\]

The group \(A_5\) contains ten subgroups \(C_3\). For each one, let
\[
 H(C_3)=N_{A_5}(C_3)\simeq S_3
\]
and let be the sign projector
\[
 e_{{\rm sgn},H}
 =\frac16\sum_{h\in H}\operatorname{sgn}_H(h)\rho(h).
\]
Here \(\operatorname{sgn}_H\) is the abstract character of \(H/C_3\)
defined above, not the —trivial— ambient sign of \(A_5\).
The orientation energy associated with the common witness is
\[
 \mathcal E_B(H)=\|e_{{\rm sgn},H}b_K\|^2.
\]
The exact census in \(\mathbb Q(\sqrt5)\) produces a unique maximum:
\[
 \mathcal E_B(H_*)
 =\frac23+\frac{2}{15}\sqrt5.
\]
The second level is
\[
 \frac13+\frac{2}{15}\sqrt5,
\]
so that the exact gap is
\[
 \boxed{\mathcal E_B(H_*)-\mathcal E_B(H_{\rm segundo})=\frac13.}
\]
The same \(H_*\) is the unique maximum when \(b_K\) is replaced by \(u_K\).
Therefore, the selector does not depend on a single encoding of the witness.

In this labeling,
\[
 H_*=N_{A_5}\bigl(\langle c_*\rangle\bigr),
\qquad
 c_*=(2\,3\,4).
\]
The rule is natural under simultaneous relabeling by \(A_5\): when
\(B_K\), \(K\), and \(P_3\) are transported, the maximum is conjugated by the same
element.

\section{Selection of orientation and presentation of \texorpdfstring{\(S_3\)}{S3}}

Within of \(H_*\), consider
\[
 s(g)=\langle b_K,\rho(g)u_K\rangle.
\]
The maximum among the two order-three elements uniquely selects
\[
 c_*=(0,1,3,4,2),
\]
in image notation of the permutation of five points, and the maximum among
the three involutions selects
\[
 r_*=(1,0,4,3,2).
\]
It is verified exactly.
\[
 r_*c_*r_*=c_*^{-1}.
\]
The presentation of the oriented sector by the cycle of positions
\(c=(0\,1\,2)\) and the reflection \(r=(0\,2)\), which fixes the central position,
therefore determines a unique isomorphism
\[
 \theta:S_3^{\rm or}\xrightarrow{\sim}H_*,
 \qquad
 \theta(c)=c_*,
 \qquad
 \theta(r)=r_*.
\]
Here \(S_3^{\rm or}\) denotes the group of relabelings of the oriented sector
\(t=7\), generated by \(c\) and \(r\); it introduces no additional operatorial
class.
The only branch that reaches horizon three is precisely the central branch
\(1\); this datum fixes which of the three positions remains fixed under
reflection.

\section{Reynolds average and polar factor}

Let
\[
 W_{\rm or}=L_{\rm or}\otimes\mathbb R^3_{\rm perm}
\]
the oriented sector triple. With the anchor already fixed, we take
\[
 \boxed{w=\ell_{\rm or}\otimes e_1,}
\]
where \(e_1\) represents the central position, the only one reaching horizon
three. Define
\[
 A
 =\sum_{g\in S_3}
 \rho_{W_{\rm or}}(g)\,
 |w\rangle\langle b_K|\,
 \rho_{V_{11}}(\theta(g))^{-1}.
\]
By construction,
\[
 A\rho_{V_{11}}(\theta(g))
 =\rho_{W_{\rm or}}(g)A,
\qquad
 A=AP_3.
\]
The exact computation in \(\mathbb Q(\sqrt5)\) obtains
\[
 \operatorname{rank}A=3.
\]
Moreover, \(G_A=AA^*\) has all its diagonal entries equal to
\[
 3+\frac25\sqrt5
\]
and all its off-diagonal entries equal to
\[
 \frac52+\frac35\sqrt5.
\]
Its two isotype eigenvalues are
\[
 \lambda_{\rm sgn}=8+\frac85\sqrt5,
 \qquad
 \lambda_{\rm std}=\frac12-\frac15\sqrt5.
\]
Both are strictly positive. In particular, \(G_A^{-1/2}\) exists, and
the polar decomposition
\[
 \boxed{U_{GK}=G_A^{-1/2}A}
\]
is one equivariant isometry of \(P_3V_{11}\) on \(W_{\rm or}\):
\[
 U_{GK}^*U_{GK}=P_3,
 \qquad
 U_{GK}U_{GK}^*=I_{W_{\rm or}}.
\]

Let
\[
 J_G:W_{\rm or}\longrightarrow L_{\rm or}\otimes H_G
\]
the isometric inclusion in the three selected branches. Then
\[
 J_GJ_G^*=I_{L_{\rm or}}\otimes P_G.
\]
The sector--dodecaphase partial isometry is denoted
\[
 \boxed{
 W_{GK}=U_{GK}^*J_G^*
 :L_{\rm or}\otimes H_G\longrightarrow V_{11}
 }
\]
satisfies
\[
 W_{GK}^*W_{GK}
 =I_{L_{\rm or}}\otimes P_G,
\qquad
 W_{GK}W_{GK}^*=P_3,
\]
and, therefore,
\[
 P_3W_{GK}
 =W_{GK}
 =W_{GK}(I_{L_{\rm or}}\otimes P_G).
\]

\begin{theorem}[Relative Isometric Identification of Three Projectors]
\label{thm:identificacion-isometrica-tres-proyectores}
Relative to the action \(A_5/C_5\), the projector \(P_3\), the seal \(K\), the doubly constructed witness \(B_K\), the oriented finite sector \(t=7\), the abstract character \(\operatorname{sgn}_{H_*}\), the anchor \(\ell_{\rm or}\), and the variational rule \(\mathcal E_B\), there exists a unique normalized partial isometry \(W_{GK}\) that identifies \(P_G\) with \(P_3\). Independently, the incidence isometry \(U_W\) identifies \(P_3\) with \(Q_3\). Consequently, the following chain of isometric identifications of projectors is obtained:
\[
\begin{CD}
L_{\rm or}\otimes H_G
 @>{W_{GK}}>>
 V_{11}
 @>{U_W}>>
 E_{30}\\
@V{I_{L_{\rm or}}\otimes P_G}VV
 @VV{P_3}V
 @VV{Q_3}V\\
L_{\rm or}\otimes H_G
 @>{W_{GK}}>>
 V_{11}
 @>{U_W}>>
 E_{30}.
\end{CD}
\]
Typed traces satisfy
\[
 \tau_{L_{\rm or}\otimes H_G}(I_{L_{\rm or}}\otimes P_G)
 =\tau_{V_{11}}(P_3)
 =\tau_{E_{30}}(Q_3)
 =\frac3{11}.
\]
The statement identifies
three projectors isometrically; it does not assert that their three spaces are
the same representation of a common group.
\end{theorem}

\begin{proof}
The energy \(\mathcal E_B\) selects a unique normalizer \(H_*\); the scoring criterion \(s(g)\) selects a unique pair \((c_*,r_*)\), and with it a unique isomorphism \(\theta\). The Reynolds average is then equivariant for the actions of \(S_3^{\rm or}\) and \(H_*\) related by \(\theta\), and has rank three. The positive part of its polar decomposition fixes the normalization without introducing a basis. The preceding partial identities prove that \(W_{GK}\) identifies \(P_G\) and \(P_3\). On the other hand, the \cref{thm:entrelazador-dodecafase-witt} identifies \(P_3\) and \(Q_3\) as projectors through \(U_W\). These two operatorial identities make the diagram commute, without extending the equivariance of the first arrow to the second.
\end{proof}

\begin{proposition}[Obstruction to triple equivariance in the fixed gauge]
\label{prop:no-go-equivarianza-hstar-witt}
The preceding construction does not prove that \(U_W\) is
\(H_*\)-equivariant. Under the declared labeling, the fixed action of
\(A_5/C_5\) on the twelve points preserves the set of hexads of
\(W_{12}\) only for the identity element. Therefore, the
\(H_*\)-equivariant intertwining of \(W_{GK}\) does not extend, in that
gauge, to a common representation on \(E_{30}\).
\end{proposition}

\begin{proof}
The first assertion separates equivariance domains:
\(W_{GK}\) uses \(H_*\subset A_5\), whereas the incidence theorem
proves equivariance of \(U_W\) for
\(\operatorname{Aut}(W_{12})\). The exact reproducible census applies the sixty
elements of action \(A_5/C_5\) to the labeled hexad system; only the
identity preserves that system. Since \(H_*\) contains six elements, its
action cannot enter by restriction of that fixed action in
\(\operatorname{Aut}(W_{12})\).
\end{proof}

To promote the chain of projectors to a common representation, it is
necessary to construct explicitly a monomorphism
\[
 \boxed{
 \iota:H_*\hookrightarrow\operatorname{Aut}(W_{12})
 }
\]
such that, for every \(h\in H_*\),
\[
 \boxed{
 U_W\,\rho_{V_{11}}(h)
 =\rho_{E_{30}}\!\bigl(\iota(h)\bigr)\,U_W.
 }
\]
This equation is also a sufficient condition for the equivariance of the
second arrow. No \(\iota\) with this property has been constructed in the
current gauge; therefore the proved equivariance terminates in \(P_3\) and is not
prolonged by mere equality of ranks.

Before using \(B_K\), \(u_K\), the orientation, and the polar factor, there were \(10\cdot6\cdot4=240\) labeled maps, or \(40\) after internal relabelings of the sector had been identified. The preceding rules reduce that set to one. This is the verifiable content of the word \emph{canonical} for \(W_{GK}\), within the declared rules. It does not extend canonicity to a common action of \(H_*\) on the Witt module.

The variational rule \(\mathcal E_B\) uses only prior HMT data
and fixes the class in which uniqueness is obtained. It does not identify the star
involution with \(P_3\), since their spectra remain distinct.

\section{Comparison of the readings of \texorpdfstring{\(t=7\) and \(t=8\)}{t=7 and t=8}}

In \(t=8\) there are three local roots and eleven extensions to horizon nine of
the unique root that continues. The sizes of the three fibers are
\[
 (11,0,0).
\]
The natural branch--chain matrix has one row of eleven ones and two
zero rows, so its rank is one. In \(t=9\), moreover, the visible quotient
changes from \(3/11\) to \(3/6\). That fraction was not the trace of a
stable projector. The correction does not eliminate the sector: it shifts the reading
to the homogeneous object \(P_G\) of \(t=7\), where rank and dimension
belong to the same space.

\section{Consequence for a positive action}

The law of action may u now a unique class isometric of projectors
positive in tris spaces of support:
\[
 P_G
 \xleftrightarrow{\ W_{GK}\ }
 P_3
 \xleftrightarrow{\ U_W\ }
 Q_3.
\]
The star \(R_{B_K}\) is reserved for parity and orientation. The projector \(P_G\) records survival; \(P_3\) records dodecaphasic polarization; \(Q_3\) records its incidence realization in Witt. This separation of functions is the mathematical premise of the measure and the positive law in the following section. The chain expresses isometric equivalence of projectors, not a common representation: the latter would require the monomorphism \(\iota\) and the equivariance identity formulated above.
 
\chapter{From the spectral gap to the Leech lattice}
\chaptermark{Exceptional publication of the continuum}
\label{chap:83-20}

\section{From the spectral gap \texorpdfstring{\(4\to2\to6\to54\)}{4 to 2 to 6 to 54}: local coupling, modular compression, and two-dimensional integration}
\label{sec:sucesor83-salto-espectral-4-2-6-54}

For two consecutive representatives \((k,k+1)\), consider
\[
B_k=
\begin{pmatrix}
k^2&k(k+1)\\
k(k+1)&(k+1)^2
\end{pmatrix}\pmod 9.
\]

\begin{theorem}[Uniqueness of the Adjacent Diagonal Block]
\label{thm:sucesor83-app-bloque-central}
There is a unique block \(B_k\) whose diagonal entries coincide. It is
determined by \(k=4\) and is
\[
B_c=\begin{pmatrix}7&2\\2&7\end{pmatrix}.
\]
\end{theorem}

\begin{proof}
The equality of the diagonals is equivalent to
\[
k^2\equiv(k+1)^2\pmod 9,
\]
that is, \(2k+1\equiv0\pmod 9\). Since \(2\) is a unit of
\(\mathbb Z/9\mathbb Z\), the congruence has the unique solution
\(k\equiv4\pmod 9\).
\end{proof}

Therefore,
\[
B_c=7I+2R,
\qquad
R=\begin{pmatrix}0&1\\1&0\end{pmatrix},
\qquad R^2=I.
\]
With \(P_\pm=(I\pm R)/2\),
\[
B_c=9P_++5P_-.
\]

\begin{theorem}[Hierarchy of the separation between sum and product]
\label{thm:sucesor83-app-jerarquia-4-2-6-54}
The central block has local spectral jump \(9-5=4\). Its off-diagonal
coefficient is \(2\). Compression by the \(3\) classes modulo
\(3\) produces the aggregate difference \(51-45=6\), and expansion
over the \(9\) positions produces the global defect
\[
9\cdot6=459-405=54.
\]
The integers \(4,2,6,54\) pertain, respectively, to the spectral gap,
the local coupling, modular compression, and two-dimensional
integration.
\end{theorem}

\begin{proof}
The decomposition \(B_c=7I+2R=9P_++5P_-\) gives the jump and the local
coefficient. The aggregated matrices \(M_\Pi,M_\Sigma\) satisfy
\(\sum M_\Pi=51\) and \(\sum M_\Sigma=45\). Each aggregated entry represents
\(9\) positions, so unfolding multiplies the difference by
\(9\) and recovers the totals \(459\) and \(405\).
\end{proof}

Their ordered concatenations produce the blocks \(72\) and \(27\):
\[
72+27=99,
\qquad
72-27=45=\det B_c.
\]
The concatenation of the diagonal and antidiagonal entries produces,
respectively, \(77\) and \(22\). Therefore,
\[
72+27=77+22=99,
\qquad
77-22=55=54+1.
\]
These equalities express two exact partitions of the same block; they do not
by themselves identify an external constant.
 
% Unidad U018. Paley, extension cuadratica finita y codigo ternario.
% Redaccion autonoma desde la autoridad material 1063.

\section{Paley conference form, finite quadratic extension and self-dual ternary code}
\label{p12-83-c20-s1:chap:paley-extension-cuadratica-codigo-ternario}

The tetrad \(B_K\) constructed in the preceding unit does not yet have a
Witt star: to define it, one must first construct the incidence system that
contains its hexads. This unit performs that step from the
projective gauge of the six units modulo nine, obtains an internal square
root of \(-I_6\), and constructs the ternary graph code of
length twelve. No identification follows from a mere equality of
cardinalities.

\subsection{Projective gauge of the \(6\) units modulo \(9\)}

Let
\[
 \mathcal U_9=(\mathbb Z/9\mathbb Z)^\times
 =\{1,2,4,8,7,5\},
\]
ordered by the successive powers of \(2\) modulo \(9\). Fijamos the
coordinate bijection
\begin{equation}
 \theta:\mathcal U_9\longrightarrow\mathbb P^1(\mathbb F_5),
 \qquad
 \begin{array}{c|cccccc}
 u&1&2&4&8&7&5\\ \hline
 \theta(u)&\infty&0&1&2&3&4
 \end{array}.
 \label{p12-83-c20-s1:eq:calibre-unidades-paley}
\end{equation}
The map \(\theta\) is a coordinate gauge. It is not asserted to be a homomorphism between multiplication in \(\mathcal U_9\) and an operation on \(\mathbb P^1(\mathbb F_5)\).

Let
\[
 \chi:\mathbb F_5\longrightarrow\{-1,0,1\}
\]
the quadratic character:
\[
 \chi(0)=0,
 \qquad
 \chi(1)=\chi(4)=1,
 \qquad
 \chi(2)=\chi(3)=-1.
\]
In the order \((\infty,0,1,2,3,4)\), we define
\begin{align}
 (C_P)_{\infty,x}&=(C_P)_{x,\infty}=1,
 &&x\in\mathbb F_5,\label{p12-83-c20-s1:eq:paley-fila-infinito}\\
 (C_P)_{x,y}&=\chi(x-y),
 &&x,y\in\mathbb F_5, x\neq y,\label{p12-83-c20-s1:eq:paley-entradas-finitas}\\
 (C_P)_{x,x}&=0.
 \label{p12-83-c20-s1:eq:paley-diagonal}
\end{align}
Thus,
\begin{equation}
 C_P=
 \begin{pmatrix}
 0& 1& 1& 1& 1& 1\\
 1& 0& 1&-1&-1& 1\\
 1& 1& 0& 1&-1&-1\\
 1&-1& 1& 0& 1&-1\\
 1&-1&-1& 1& 0& 1\\
 1& 1&-1&-1& 1& 0
 \end{pmatrix}.
 \label{p12-83-c20-s1:eq:matriz-conferencia-paley-autonoma}
\end{equation}

\begin{lemma}[Quadratic correlation sum]
\label{p12-83-c20-s1:lem:correlacion-cuadratica-f5}
For \(a\neq b\) in \(\mathbb F_5\),
\begin{equation}
 \sum_{t\in\mathbb F_5}
 \chi(t-a)\chi(t-b)=-1.
 \label{p12-83-c20-s1:eq:correlacion-cuadratica-f5}
\end{equation}
\end{lemma}

\begin{proof}
The substitution \(u=(t-a)/(b-a)\) is to bijection of \(\mathbb F_5\).
Ifnce \(\chi((b-a)^2)=1\), the sum  reduce to
\[
 \sum_{u\in\mathbb F_5}\chi(u(u-1)).
\]
Evaluation of \(u=0,1,2,3,4\) gives, respectively,
\(0,0,1,-1,-1\), whose sum is \(-1\).
\end{proof}

\begin{proposition}[Conference Identity]
\label{p12-83-c20-s1:prop:paley-identidad-conferencia-autonoma}
The matrix \(C_P\) is symmetric and satisfies
\begin{equation}
 \boxed{C_PC_P^{\mathsf T}=5I_6.}
 \label{p12-83-c20-s1:eq:paley-identidad-conferencia-autonoma}
\end{equation}
\end{proposition}

\begin{proof}
Each row contains five entries of modulus one and one zero entry; its
squared norm is five. For two distinct finite rows, the contribution
of the coordinate \(\infty\) is \(+1\), whereas
\cref{p12-83-c20-s1:lem:correlacion-cuadratica-f5} gives \(-1\) for the five finite
coordinates. Their inner product is zero. The sum of a nontrivial character
over \(\mathbb F_5\) is zero, so the row corresponding to
\(\infty\) is also orthogonal to every finite row.
\end{proof}

\subsection{ternary reduction and internal quadratic root of \texorpdfstring{\(-I_6\)}{-I₆}}

We reduce \(C_P\) modulo tris, with \(-1\equiv2\pmod3\):
\begin{equation}
 A_W=
 \begin{pmatrix}
 0&1&1&1&1&1\\
 1&0&1&2&2&1\\
 1&1&0&1&2&2\\
 1&2&1&0&1&2\\
 1&2&2&1&0&1\\
 1&1&2&2&1&0
 \end{pmatrix}
 \in M_6(\mathbb F_3).
 \label{p12-83-c20-s1:eq:matriz-paley-witt-ternaria}
\end{equation}

\begin{theorem}[Internal square root of \(-I_6\)]
\label{p12-83-c20-s1:thm:aw-cuadrado-menos-identidad-autonomo}
We have
\begin{equation}
 \boxed{A_W^2=-I_6.}
 \label{p12-83-c20-s1:eq:aw-cuadrado-menos-identidad-autonomo}
\end{equation}
In particular,
\[
 A_W^{-1}=-A_W.
\]
\end{theorem}

\begin{proof}
The \cref{p12-83-c20-s1:prop:paley-identidad-conferencia-autonoma} gives
\[
 A_WA_W^{\mathsf T}
 \equiv5I_6
 \equiv2I_6
 =-I_6
 \pmod3.
\]
Since \(A_W=A_W^{\mathsf T}\), \(A_W^2=-I_6\) is obtained.
\end{proof}

The preceding identity is the algebraic residence of the root of \(-1\)
in this corridor. No proceeds of the factor \(1/4\) of the Hadamard inversion:
the former is a rational scalar forced by \(H_4^2=4I_4\); this is one
quadratic structure on characteristic three.

\subsection{Extension quadratic to \texorpdfstring{\(\mathbb F_9\)}{F₉}}

We define
\begin{equation}
 \mathbb F_9=\mathbb F_3[\iota]/(\iota^2+1),
 \qquad \iota^2=-1.
 \label{p12-83-c20-s1:eq:definicion-f9-autonoma}
\end{equation}
The polynomial \(X^2+1\) has no roots in \(\mathbb F_3\), so that the
quotient is a field of nine elements. For
\(V=\mathbb F_3^6\), let
\[
 V_9=V\otimes_{\mathbb F_3}\mathbb F_9.
\]
The operators
\begin{equation}
 P_\iota=\frac12(I-\iota A_W),
 \qquad
 P_{-\iota}=\frac12(I+\iota A_W)
 \label{p12-83-c20-s1:eq:proyectores-f9-paley}
\end{equation}
are well defined because \(2\) is invertible in \(\mathbb F_3\).

\begin{proposition}[Conjugate projectors]
\label{p12-83-c20-s1:prop:proyectores-conjugados-paley}
The following hold
\[
 P_\iota^2=P_\iota,
 \qquad
 P_{-\iota}^2=P_{-\iota},
 \qquad
 P_\iota P_{-\iota}=0,
 \qquad
 P_\iota+P_{-\iota}=I.
\]
\end{proposition}

\begin{proof}
Usando \(\iota^2=-1\) and \(A_W^2=-I\),
\[
 (\iota A_W)^2=I.
\]
The four identities are then the elementary identities for the projectors
\((I\mp\iota A_W)/2\) of an involution.
\end{proof}

\begin{theorem}[Conjugate spectral decomposition]
\label{p12-83-c20-s1:thm:descomposicion-espectral-f9-paley}
Let
\[
 E_{\pm\iota}=\ker(A_W\mp\iota I).
\]
Then
\begin{equation}
 V_9=E_\iota\oplus E_{-\iota},
 \qquad
 \dim_{\mathbb F_9}E_\iota
 =\dim_{\mathbb F_9}E_{-\iota}=3.
 \label{p12-83-c20-s1:eq:descomposicion-espectral-f9-paley}
\end{equation}
\end{theorem}

\begin{proof}
The minimal polynomial of \(A_W\) divides \(X^2+1\), which splits over
\(\mathbb F_9\) with simple roots \(\pm\iota\). The projectors of
\cref{p12-83-c20-s1:prop:proyectores-conjugados-paley} produce the direct sum. Since
\(A_W\) has coefficients in \(\mathbb F_3\), the Frobenius automorphism
exchanges the two eigenspaces; their dimensions are equal and
sum to six.
\end{proof}

This construction dota \(\mathbb F_3^6\) of istructura of modulo of rango
tris over \(\mathbb F_9\), by means of
\begin{equation}
 (a+b\iota)\cdot v:=av+b(vA_W).
 \label{p12-83-c20-s1:eq:accion-f9-sobre-f3-seis}
\end{equation}
The equality \(3^6=9^3\) is thus accompanied by an explicit action; it is not
used as a cardinal argument.

\subsection{Ternary graph code of length \(12\)}

We define
\begin{equation}
 c:\mathbb F_3^6\longrightarrow\mathbb F_3^{12},
 \qquad
 c(q)=(q,qA_W),
 \label{p12-83-c20-s1:eq:aplicacion-codigo-grafico-witt}
\end{equation}
and
\begin{equation}
 C_W:=\operatorname{im}c.
 \label{p12-83-c20-s1:eq:definicion-codigo-witt-ternario}
\end{equation}
Its generating matrix is
\[
 G_W=(I_6\mid A_W).
\]

\begin{theorem}[Injectivity and self-duality]
\label{p12-83-c20-s1:thm:cw-inyectivo-autodual}
The map \(c\) is injective, \(\dim C_W=6\) and
\begin{equation}
 \boxed{C_W=C_W^\perp.}
 \label{p12-83-c20-s1:eq:cw-autodual}
\end{equation}
\end{theorem}

\begin{proof}
The first half of \(c(q)\) is \(q\), so
\(\operatorname{pr}_1\circ c=\operatorname{id}\); this proves
injectivity and the dimension. Moreover,
\[
 G_WG_W^{\mathsf T}
 =I_6+A_WA_W^{\mathsf T}
 =I_6+2I_6=0
\]
in \(\mathbb F_3\). Therefore, \(C_W\subseteq C_W^\perp\). Both spaces
have dimension six, so they are equal.
\end{proof}

For \(r=0,\ldots,12\), let
\[
 N_r:=\#\{q\in\mathbb F_3^6:
            \operatorname{wt}(q,qA_W)=r\}.
\]
The enumeration of the \(3^6=729\) entries produces
\begin{equation}
 N_0=1,
 \qquad
 N_6=264,
 \qquad
 N_9=440,
 \qquad
 N_{12}=24,
 \label{p12-83-c20-s1:eq:censo-pesos-cw}
\end{equation}
and \(N_r=0\) in the remaining weights.

\begin{theorem}[Enumerator and minimum distance]
\label{p12-83-c20-s1:thm:enumerador-cw-autonomo}
The code \(C_W\) has parameters
\[
 \boxed{C_W=[12,6,6]_3}
\]
and weight enumerator
\begin{equation}
 \boxed{W_{C_W}(y)=1+264y^6+440y^9+24y^{12}.}
 \label{p12-83-c20-s1:eq:enumerador-pesos-cw}
\end{equation}
\end{theorem}

\begin{proof}
The dimension and self-duality have already been proved. The census
\eqref{p12-83-c20-s1:eq:censo-pesos-cw} traverses the \(729\) words without sampling. Its
least nonzero weight is six, which is the minimum distance of a linear code.
The sum \(1+264+440+24=729\) checks that no word was omitted.
\end{proof}

\subsection{Nature of the structure obtained}

Three distinct objects and three typed arrows have been constructed:
\[
 \mathcal U_9
 \xrightarrow{\theta}\mathbb P^1(\mathbb F_5)
 \xrightarrow{\chi}C_P
 \xrightarrow{\bmod3}A_W,
\]
\[
 A_W^2=-I_6
 \quad\Longrightarrow\quad
 \mathbb F_3^6\cong\mathbb F_9^3
 \quad\Longrightarrow\quad
 C_W=\{(q,qA_W):q\in\mathbb F_3^6\}.
\]
The first arrow is a gauge, the second uses the quadratic character, and the third is coefficient reduction. The structure over
\(\mathbb F_9\) is finite; it is not identified with a real or analytic complex structure.

\subsection{Obstructions in the Paley–\(C_W\) corridor: bijection, \(A_W^2=-I_6\), self-duality, and census of \(729\) words}

\begin{criterion}[Falsification of the Paley--\(C_W\) corridor]
The construction is refuted if any of the following occurs:
\begin{enumerate}
 \item \(\theta\) is not bijective or is used as a homomorphism without proof;
 \item the sum \eqref{p12-83-c20-s1:eq:correlacion-cuadratica-f5} differs from \(-1\);
 \item \(C_PC_P^{\mathsf T}\neq5I_6\);
 \item \(A_W^2\neq-I_6\);
 \item the projectors \eqref{p12-83-c20-s1:eq:proyectores-f9-paley} are not
       complementary;
 \item the map \(q\mapsto(q,qA_W)\) loses injectivity or
       self-duality;
 \item the census does not sum to \(729\), differs from
       \eqref{p12-83-c20-s1:eq:censo-pesos-cw}, or presents a word of weight between one
       and five;
 \item \(\mathbb F_3^6\cong\mathbb F_9^3\) is inferred only from
       \(3^6=9^3\), omitting the action \eqref{p12-83-c20-s1:eq:accion-f9-sobre-f3-seis}.
\end{enumerate}
\end{criterion}


% Unidad U019. Sistema de Witt, campo de estrellas y grupo M12.
% Redaccion autonoma desde la autoridad material 1063.

\section{Ternary Steiner system, Witt-star involutions and generation of \texorpdfstring{\(M_{12}\)}{M12}}
\label{p12-83-c20-s2:chap:sistema-witt-estrellas-m12}

The preceding unit constructed the self-dual code
\(C_W=[12,6,6]_3\). Its words of weight six are now projectivized,
the Steiner property is proved, and only then is the star
associated with the tetrad \(B_K\) defined. This order prevents an incidence
figure from being turned into an object added by visual similarity.

\subsection{Supports projective of weight \(6\)}

Let
\begin{equation}
 \mathcal H_W
 :=\{\operatorname{supp}(c):c\in C_W,
                              \operatorname{wt}(c)=6\}.
 \label{p12-83-c20-s2:eq:familia-soportes-peso-seis}
\end{equation}
Each support occurs for the word pair \(\{c,-c\}\). Indeed, if two
weight-six words with the same support were not proportional, one could
scale one of them to cancel a coordinate upon subtraction; this would yield
a nonzero word of \(C_W\) of weight less than six, contrary to the minimum
distance. Since neither is there a nonzero word equal to its opposite in
characteristic three, every projective orbit has exactly two elements.
Therefore,
\begin{equation}
 |\mathcal H_W|=\frac{264}{2}=132.
 \label{p12-83-c20-s2:eq:numero-hexadas-cw}
\end{equation}

For \(T\in\binom{[12]}5\), define
\begin{equation}
 \nu(T):=\#\{H\in\mathcal H_W:T\subseteq H\}.
 \label{p12-83-c20-s2:eq:incidencia-cinco-hexada}
\end{equation}
The exact enumeration of the
\[
 \binom{12}{5}=792
\]
subsets of five points produce
\begin{equation}
 \nu(T)=1
 \qquad
 \text{for every }T\in\binom{[12]}5.
 \label{p12-83-c20-s2:eq:unicidad-cinco-en-hexada}
\end{equation}

\begin{theorem}[Ternary Witt System]
\label{p12-83-c20-s2:thm:w12-autonomo}
The incidence system
\[
 W_{12}:=([12],\mathcal H_W)
\]
is the Steiner system
\begin{equation}
 \boxed{W_{12}=S(5,6,12).}
 \label{p12-83-c20-s2:eq:w12-steiner-autonomo}
\end{equation}
\end{theorem}

\begin{proof}
The equality \eqref{p12-83-c20-s2:eq:unicidad-cinco-en-hexada} is precisely the
Steiner property. The number of blocks is also recovered from the double
count of pairs \((T,H)\) with \(T\subset H\):
\[
 |\mathcal H_W|\binom65=\binom{12}{5},
 \qquad
 |\mathcal H_W|=\frac{792}{6}=132.
\]
\end{proof}

\subsection{Residue of a tetrad and decomposition into \(4\) incident components}

Let \(B\in\binom{[12]}4\). In an \(S(5,6,12)\) system, the number of
hexads containing \(B\) is
\begin{equation}
 \lambda_4
 =\frac{\binom{12-4}{5-4}}{\binom{6-4}{5-4}}
 =\frac82=4.
 \label{p12-83-c20-s2:eq:lambda-cuatro-w12}
\end{equation}
We write these four hexads as
\[
 H_i=B\sqcup P_i,
 \qquad |P_i|=2,
 \qquad i=1,\ldots,4.
\]

\begin{lemma}[Residual partition of a tetrad]
\label{p12-83-c20-s2:lem:particion-residual-tetrada-witt}
The four pairs \(P_1,\ldots,P_4\) are disjoint and
\begin{equation}
 [12]\setminus B
 =P_1\sqcup P_2\sqcup P_3\sqcup P_4.
 \label{p12-83-c20-s2:eq:particion-residual-tetrada-witt}
\end{equation}
\end{lemma}

\begin{proof}
If two residual pairs shared a point \(p\), the corresponding
hexads would contain the same subset \(B\cup\{p\}\) of five points,
in contradiction with the Steiner uniqueness. The four pairs are, therefore,
disjoint; contain eight points and agotan the complement of \(B\).
\end{proof}

\begin{definition}[Star and Witt star involution]
\label{p12-83-c20-s2:def:estrella-witt-involucion}
The star centered at \(B\) is
\begin{equation}
 \mathscr S_B
 :=\{H\in\mathcal H_W:B\subseteq H\}
 =\{B\sqcup P_1,\ldots,B\sqcup P_4\}.
 \label{p12-83-c20-s2:eq:estrella-witt-general}
\end{equation}
If \(P_i=\{p_i^-,p_i^+\}\), its involution is
\begin{equation}
 \sigma_B:=\prod_{i=1}^{4}(p_i^-\ p_i^+).
 \label{p12-83-c20-s2:eq:involucion-estrella-witt-general}
\end{equation}
The permutation \(\sigma_B\) fixes \(B\) pointwise and exchanges the
ends of each residual petal.
\end{definition}

The mathematical term used henceforth is \emph{Witt star} or
\emph{Witt-star involution}. It is not a Euclidean
star.

\subsection{Star selected by the tetrad of the seal}

The dodecaphasic unit produced, without using Steiner's design,
\[
 B_K=\{4,6,9,10\}.
\]
Now that \(W_{12}\) has been constructed, its star is well defined. The
four incident hexads are
\begin{align}
 H_{13}&=B_K\sqcup\{1,3\}
       =\{1,3,4,6,9,10\},
 \label{p12-83-c20-s2:eq:hexada-petalo-13}\\
 H_{2,11}&=B_K\sqcup\{2,11\}
       =\{2,4,6,9,10,11\},
 \label{p12-83-c20-s2:eq:hexada-petalo-211}\\
 H_{5,7}&=B_K\sqcup\{5,7\}
       =\{4,5,6,7,9,10\}
       =H_\alpha^-,
 \label{p12-83-c20-s2:eq:hexada-petalo-57}\\
 H_{8,12}&=B_K\sqcup\{8,12\}
       =\{4,6,8,9,10,12\}.
 \label{p12-83-c20-s2:eq:hexada-petalo-812}
\end{align}
Consequently,
\begin{equation}
 \boxed{
 \sigma_{B_K}=(1\ 3)(2\ 11)(5\ 7)(8\ 12).}
 \label{p12-83-c20-s2:eq:involucion-estrella-witt-bk}
\end{equation}

\begin{figure}[htbp]
\centering
\begin{tikzpicture}[
  >=Stealth,
  font=\small,
  core/.style={draw,double,rounded corners=2mm,align=center,
    minimum width=4.0cm,minimum height=1.3cm,fill=black!4},
  petal/.style={draw,rounded corners=2mm,align=center,
    minimum width=3.55cm,minimum height=1.08cm,fill=black!2},
  negative/.style={petal,line width=1.1pt,double,fill=black!10},
  arr/.style={->,thick}
]
\node[core] (B) at (0,0)
 {\(\displaystyle B_K=\{4,6,9,10\}\)\\
  points fixed by \(\sigma_{B_K}\)};
\node[petal] (P13) at (0,3.0)
 {\(P_1=\{1,3\}\), \((1\ 3)\)\\
  \(H_{13}=B_K\sqcup P_1\)};
\node[petal] (P211) at (5.25,0)
 {\(P_2=\{2,11\}\), \((2\ 11)\)\\
  \(H_{2,11}=B_K\sqcup P_2\)};
\node[negative] (P57) at (0,-3.0)
 {\(P_3=\{5,7\}\), \((5\ 7)\)\\
  \(H_{5,7}=H_\alpha^-\)};
\node[petal] (P812) at (-5.25,0)
 {\(P_4=\{8,12\}\), \((8\ 12)\)\\
  \(H_{8,12}=B_K\sqcup P_4\)};
\draw[arr] (B)--(P13);
\draw[arr] (B)--(P211);
\draw[arr] (B)--(P57);
\draw[arr] (B)--(P812);
\node[align=center,font=\footnotesize] at (0,-4.35)
 {\(\mathscr S_{B_K}
   =\{H_{13},H_{2,11},H_{5,7},H_{8,12}\}\).\\
  The diagram represents incidence, not geometric distance.};
\end{tikzpicture}
\caption{Witt star centered at \(B_K\). The center is the fixed tetrad;
the four petals are the residual pairs that determine the
involution.}
\label{p12-83-c20-s2:fig:estrella-witt-bk-autonoma}
\end{figure}

\subsection{Global field of the 495 stars}

There exists a star for each tetrad, so
\begin{equation}
 \#\{\mathscr S_B:B\in\tbinom{[12]}4\}
 =\binom{12}{4}=495.
 \label{p12-83-c20-s2:eq:numero-estrellas-witt}
\end{equation}
The exact census verifies
\begin{equation}
 \sigma_B(\mathcal H_W)=\mathcal H_W
 \qquad
 \text{for every }B\in\binom{[12]}4.
 \label{p12-83-c20-s2:eq:estrellas-conservan-hexadas}
\end{equation}
We define
\begin{equation}
 \Gamma_\star
 :=\langle\sigma_B:B\in\tbinom{[12]}4\rangle
 \leq\operatorname{Aut}(W_{12}).
 \label{p12-83-c20-s2:eq:grupo-generado-estrellas-witt}
\end{equation}

A generating set of six tetrads is
\[
 1234,
 \quad1235,
 \quad1236,
 \quad1245,
 \quad1345,
 \quad2345.
\]
The orders of the subgroups obtained by adding their involutions successively are
\begin{equation}
 1, 2, 6, 18, 360, 7920, 95040.
 \label{p12-83-c20-s2:eq:cadena-ordenes-generadores-m12}
\end{equation}

\begin{theorem}[Generation of the Mathieu group \(M_{12}\)]
\label{p12-83-c20-s2:thm:495-estrellas-m12-autonomo}
We have
\begin{equation}
 \boxed{
 \langle\sigma_B:B\in\tbinom{[12]}4\rangle
 =\operatorname{Aut}(W_{12})
 \cong M_{12}.}
 \label{p12-83-c20-s2:eq:estrellas-generan-m12}
\end{equation}
\end{theorem}

\begin{proof}
Each \(\sigma_B\) preserves the \(132\) hexads, hence
\(\Gamma_\star\leq\operatorname{Aut}(W_{12})\). The Schreier enumeration of
the six indicated generators yields order \(95040\). The automorphism group
of \(S(5,6,12)\) is \(M_{12}\) and has that same order; therefore
the inclusion is an equality.
\end{proof}

An isolated star generates only
\[
 \langle\sigma_B\rangle\cong C_2.
\]
The conclusion about \(M_{12}\) belongs to the global field of the
\(495\) involutions, not to \(\sigma_{B_K}\) considered in isolation.

\subsection{Action on the augmentation module}

Let
\[
 V_{11}=\mathbf1^\perp\subset\mathbb R^{12}.
\]
The involution \(\sigma_B\) fixes the four coordinates of \(B\) and acts by four transpositions on its complement. In \(\mathbb R^{12}\) it has eight directions with eigenvalue \(+1\) and four with eigenvalue \(-1\). The constant direction belongs to the first eigenspace. Therefore,
\begin{equation}
 \operatorname{Spec}(\sigma_B|_{V_{11}})
 =\{1^{[7]},(-1)^{[4]}\},
 \qquad
 \frac1{11}\operatorname{Tr}(\sigma_B|_{V_{11}})=\frac3{11}.
 \label{p12-83-c20-s2:eq:espectro-involucion-estrella-aumentacion}
\end{equation}
It is a signed involution; it is not a positive projector of rank three.

\subsection{Obstructions concerning \(W_{12}\), the field of \(495\) star involutions and the generation of \(M_{12}\)}

\begin{criterion}[Refutation of \(W_{12}\), its stars, and \(M_{12}\)]
The construction is refuted if any of the following facts occurs:
\begin{enumerate}
 \item the projectivization of the \(264\) words of weight six does not produce
       \(132\) supports;
 \item there exists a subset of five points contained in zero or in more of
       one hexad;
 \item a tetrad does not belong to exactly four hexads;
 \item the four residual pairs do not partition the complement of the
       tetrad;
 \item the petals of \(B_K\) are not
       \(\{1,3\},\{2,11\},\{5,7\},\{8,12\}\);
 \item any of the \(495\) involutions fails to preserve
       \(\mathcal H_W\);
 \item the generated group does not have order \(95040\);
 \item \(M_{12}\) is attributed to a single star instead of the global field of involutions;
 \item  introduces the istrella before of constructing \(W_{12}\).
\end{enumerate}
\end{criterion}


% Unidad U020. Prolongacion binaria W12--W24 y M24.
% Redaccion autonoma desde la autoridad material 1063.

\section{Binary extension relative to a dodecad, system \texorpdfstring{\(W_{24}\)}{W24} and group \texorpdfstring{\(M_{24}\)}{M24}}
\label{p12-83-c20-s3:chap:dodecada-w24-m24}

The system \(W_{12}\) constructed over twelve phases admits a binary prolongation
when an extended Golay code, a dodecad, and an incidence
isomorphism are declared. These data are relative inputs of this branch:
they are not deduced solely from \(C_W\), and they do not intervene in the ternary lattice
branch that will lead to the Leech lattice.

\subsection{Golay binary code and octads}

Let
\[
 \mathcal G_{24}\subset\mathbb F_2^{24}
\]
the extended binary Golay code. Its weights are
\begin{equation}
 0, 8, 12, 16, 24.
 \label{p12-83-c20-s3:eq:pesos-golay-binario-extendido}
\end{equation}
The supports of the weight eight words form the Steiner system
\begin{equation}
 W_{24}=S(5,8,24).
 \label{p12-83-c20-s3:eq:w24-steiner-octadas}
\end{equation}
Its blocks are called \emph{octads}. Fix a dodecad \(D\), that is,
the support of a word of weight twelve, and let \(\mathcal O_{24}\)
be the family of octads.

We define
\begin{equation}
 \mathcal H(D)
 :=\{O\cap D:O\in\mathcal O_{24},\ |O\cap D|=6\}.
 \label{p12-83-c20-s3:eq:hexadas-residuales-dodecada}
\end{equation}

\begin{theorem}[Residual design of a dodecada]
\label{p12-83-c20-s3:thm:diseno-residual-dodecada-autonomo}
The residual system is
\begin{equation}
 \boxed{(D,\mathcal H(D))\cong S(5,6,12).}
 \label{p12-83-c20-s3:eq:diseno-residual-dodecada-w12}
\end{equation}
Moreover, each \(H\in\mathcal H(D)\) has one unique octad \(O_H\) such that
\begin{equation}
 O_H\cap D=H.
 \label{p12-83-c20-s3:eq:elevacion-unica-hexada-octada}
\end{equation}
\end{theorem}

\begin{proof}
Let \(A\subset D\), with \(|A|=5\). In \(W_{24}=S(5,8,24)\) there is a
unique octad \(O_A\) containing \(A\). If \(j=|O_A\cap D|\), the binary
sum of the characteristic words of \(O_A\) and \(D\) belongs to
\(\mathcal G_{24}\) and has weight
\[
 8+12-2j=20-2j.
\]
Since \(A\subseteq O_A\cap D\), one has \(5\leq j\leq8\). The four
possible weights \(10,8,6,4\), combined with
\eqref{p12-83-c20-s3:eq:pesos-golay-binario-extendido}, force \(j=6\). Thus every
five-point subset of \(D\) belongs to a unique residual hexad.

If two octads had the same hexadic intersection with \(D\), both
would contain the same five points of that hexad, contradicting
uniqueness in \(S(5,8,24)\). This also proves uniqueness of the
lift.
\end{proof}

\subsection{Relative scale between the HMT design and the dodecagon}

To relate the preceding construction to the design obtained from
\(C_W\), one must explicitly declare
\begin{equation}
 (D,\ell),
 \qquad
 \ell:W_{12}^{\rm HMT}
 \xrightarrow{\sim}(D,\mathcal H(D)),
 \label{p12-83-c20-s3:eq:dato-relativo-dodecada-alineamiento}
\end{equation}
where \(\ell\) is an incidence isomorphism. If \(\ell_*\) denotes its
action on blocks, we define
\begin{equation}
 \widehat\iota_{D,\ell}
 :=\iota_D\circ\ell_*:
 W_{12}^{\rm HMT}\longrightarrow W_{24},
 \qquad
 \widehat\iota_{D,\ell}(H)=O_{\ell(H)}.
 \label{p12-83-c20-s3:eq:elevacion-tipificada-w12-w24}
\end{equation}
The composition \(\widehat\iota_{D,\ell}\) avoids applying
\(\iota_D\), whose domain is \(\mathcal H(D)\), directly to a hexad
that still lies in \([12]_{\rm HMT}\).

\begin{proposition}[Injectivity of the Relative Lift]
\label{p12-83-c20-s3:prop:inyectividad-elevacion-w12-w24}
The map \(\widehat\iota_{D,\ell}\) is injective and preserves the
incidence between points and hexads after applying \(\ell\).
\end{proposition}

\begin{proof}
The isomorphism \(\ell\) is bijective on blocks. By the
\cref{p12-83-c20-s3:thm:diseno-residual-dodecada-autonomo}, each residual hexad has a
unique lifting octad. Two HMT hexads with the same image would have the same
residual hexad and, by bijectivity of \(\ell\), would be equal.
\end{proof}

There is no absolute choice of \((D,\ell)\) prior to this gauge. The
invariant object is the class of the construction under simultaneous change of
dodecad and alignment.

\subsection{Incidence law \texorpdfstring{\(4+1\)}{4+1} on a tetrad}

For a tetrad \(B\subset D\), the parameters derived from the two
Steiner systems are
\begin{equation}
 \lambda_4(W_{12})
 =\frac{\binom81}{\binom21}=4,
 \qquad
 \lambda_4(W_{24})
 =\frac{\binom{20}1}{\binom41}=5.
 \label{p12-83-c20-s3:eq:ley-cuatro-mas-uno-witt}
\end{equation}
Therefore, the five octads containing \(B\) decompose uniquely
as
\begin{equation}
 \{O_H:H\in\mathcal H(D),\ B\subset H\}
 \sqcup\{O_B^\perp\}.
 \label{p12-83-c20-s3:eq:descomposicion-cinco-octadas}
\end{equation}
The four first satisfy \(|O_H\cap D|=6\); the quinta satisfies
\begin{equation}
 O_B^\perp\cap D=B.
 \label{p12-83-c20-s3:eq:octada-transversal-tetrada}
\end{equation}

For \(B=B_K\), the four residual lifts correspond to the
petals
\[
 \{1,3\},
 \qquad
 \{2,11\},
 \qquad
 \{5,7\},
 \qquad
 \{8,12\},
\]
and the fifth octad is the transversal \(O_{B_K}^\perp\). This law explains
the transition from four hexads to five octads without identifying a face
of a cube with a hexad or a vertex with an octad.

\subsection{Strata Correspondence and Positivity of the Triple Freedom}

The five-region structure of channel \(\pi\) has types
\[
 1_\perp+1_{--}+3_{++}.
\]
The incidence on \(B_K\) has types
\[
 1_{\rm transversal}+1_{H^-}+3_{\rm restantes}
\]
after marking the negative hexad. There is therefore a natural
correspondence between
\emph{role} strata:
\begin{equation}
 1_\perp\longleftrightarrow O_{B_K}^\perp,
 \qquad
 1_{--}\longleftrightarrow
 \widehat\iota_{D,\ell}(H_\alpha^-),
 \qquad
 3_{++}\longleftrightarrow
 \mathcal O_{B_K}^{+,3}.
 \label{p12-83-c20-s3:eq:correspondencia-estratos-cinco-cinco}
\end{equation}
Here \(\mathcal O_{B_K}^{+,3}\) is the unordered set of the other
three residual octads. An individual bijection between the three positive
regions and those three octads requires declaring an additional gauge of
\(S_3\). The coincidence of the pattern \(1+1+3\) does not by itself determine that
ordering.

\subsection{Automorphism group of the octad system}

\begin{theorem}[Mathieu group \(M_{24}\)]
\label{p12-83-c20-s3:thm:automorfismos-w24-m24}
The automorphism group of the octad Steiner system is
\begin{equation}
 \boxed{
 \operatorname{Aut}(W_{24})\cong M_{24},
 \qquad
 |M_{24}|=244823040.}
 \label{p12-83-c20-s3:eq:automorfismos-w24-m24}
\end{equation}
\end{theorem}

The theorem identifies the classical automorphism group of the input
binary data. It does not assert that \(M_{24}\) is generated from a single
dodecad, or that the ternary code \(C_W\) determines
\(\mathcal G_{24}\) by itself.

\subsection{Categorical separation from the lattice branch}

The branch constructed in this unit is
\begin{equation}
 (W_{12}^{\rm HMT};\mathcal G_{24},D,\ell)
 \longrightarrow W_{24}
 \longrightarrow\operatorname{Aut}(W_{24})\cong M_{24}.
 \label{p12-83-c20-s3:eq:rama-binaria-w24-m24}
\end{equation}
The lattice branch to be constructed subsequently is
\begin{equation}
 C_W\longrightarrow N(A_2^{12})
 \xrightarrow{\text{index-three neighbor}}\Lambda_{24}.
 \label{p12-83-c20-s3:eq:rama-ternaria-niemeier-leech-anunciada}
\end{equation}
Their constructors, domains, and codomains are distinct. In particular,
\begin{equation}
 \boxed{
 \text{no arrow is introduced }
 W_{24}\longrightarrow\Lambda_{24}.}
 \label{p12-83-c20-s3:eq:no-flecha-w24-leech}
\end{equation}

\begin{figure}[htbp]
\centering
\begin{tikzpicture}[
  >=Stealth,
  node distance=13mm and 16mm,
  box/.style={draw,rounded corners=1.5mm,align=center,
    inner xsep=3mm,inner ysep=2.2mm},
  arr/.style={->,thick}
]
\node[box] (cw) {\(C_W\)};
\node[box,right=of cw] (w12) {\(W_{12}\)};
\node[box,right=of w12] (w24) {\(W_{24}\)};
\node[box,right=of w24] (m24) {\(M_{24}\)};
\node[box,below=of cw] (cw2) {\(C_W\)};
\node[box,right=of cw2] (n) {\(N(A_2^{12})\)};
\node[box,right=of n] (leech) {\(\Lambda_{24}\)};
\draw[arr] (cw)--(w12);
\draw[arr] (w12)--node[above,font=\scriptsize]
 {\((\mathcal G_{24},D,\ell)\)} (w24);
\draw[arr] (w24)--(m24);
\draw[arr] (cw2)--node[below,font=\scriptsize]
 {pegado ternario} (n);
\draw[arr] (n)--node[below,font=\scriptsize]
 {\(3\)-neighbor} (leech);
\node[below=5mm of w24,font=\footnotesize,align=center]
 {Vertical alignment is not an arrow.};
\end{tikzpicture}
\caption{Typed bifurcation of the binary and ternary extensions.}
\label{p12-83-c20-s3:fig:bifurcacion-w24-leech-autonoma}
\end{figure}

\subsection{Obstructions to the binary prolongation \(W_{12}\to W_{24}\): Golay code, octads, alignment and separation of the Leech branch}

\begin{criterion}[Refutation of the binary prolongation]
The construction is refuted if any of the following facts occurs:
\begin{enumerate}
 \item the weight spectrum of \(\mathcal G_{24}\) contains a weight not
       listed in \eqref{p12-83-c20-s3:eq:pesos-golay-binario-extendido};
 \item the supports of weight eight not form \(S(5,8,24)\);
 \item a residual intersection of an octad with \(D\) that contains five
       points of \(D\) has cardinality other than six;
 \item a residual hexad admits either no lifting octad or more than one;
 \item \(\iota_D\) is applied directly to an HMT hexad without the alignment \(\ell\);
 \item the incidence law on a tetrad is not \(4+1\);
 \item the positive triple is ordered without declaring the gauge \(S_3\);
 \item \(\operatorname{Aut}(W_{24})\) has no order \(244823040\);
 \item it is asserted that \(W_{24}\) arises from \(W_{12}\) without the binary datum
       \(\mathcal G_{24}\), the dodecad, and the alignment;
 \item the two branches are serialized by means of an arrow
       \(W_{24}\to\Lambda_{24}\).
\end{enumerate}
\end{criterion}


% Unidad U021. Pegado discriminante ternario y Niemeier A2^12.
% Redaccion autonoma desde la autoridad material 1063.

\section{Ternary discriminant gluing and construction of the Niemeier lattice of type \texorpdfstring{\(A_2^{12}\)}{A₂¹²}}
\label{p12-83-c20-s4:chap:pegado-niemeier-a2-doce}

This unit begins the lattice branch directly from the self-dual ternary code
\(C_W\). The binary system \(W_{24}\), the dodecad, and the group
\(M_{24}\) are not inputs of this construction. The typed sequence is
\[
 C_W
 \longrightarrow
 (A_2^*/A_2)^{12}
 \longrightarrow
 N(C_W)
 \cong N(A_2^{12}).
\]

\subsection{Reticulo root \texorpdfstring{\(A_2\)}{A₂}}

Let
\[
 A_2=\mathbb Z\alpha_1\oplus\mathbb Z\alpha_2
\]
the root lattice whose Gram matrix, in the ordered basis
\((\alpha_1,\alpha_2)\), is
\begin{equation}
 G_{A_2}=
 \begin{pmatrix}
 2&-1\\
 -1&2
 \end{pmatrix}.
 \label{p12-83-c20-s4:eq:gram-a2-autonomo}
\end{equation}
For \(x=x_1\alpha_1+x_2\alpha_2\),
\begin{equation}
 \langle x,x\rangle
 =2(x_1^2-x_1x_2+x_2^2).
 \label{p12-83-c20-s4:eq:norma-a2-autonoma}
\end{equation}
In particular, \(A_2\) is even, positive defined and
\[
 \det A_2=3.
\]

We define
\[
 \rho:=\alpha_1+\alpha_2.
\]
Then
\begin{equation}
 \rho^2=2,
 \qquad
 \langle\rho,\alpha_1\rangle
 =\langle\rho,\alpha_2\rangle=1.
 \label{p12-83-c20-s4:eq:identidades-rho-a2}
\end{equation}

\subsection{Discriminant module and form}

The dual lattice is
\[
 A_2^*
 :=\{x\in A_2\otimes_{\mathbb Z}\mathbb Q:
       \langle x,A_2\rangle\subseteq\mathbb Z\}.
\]
We choose the representatives
\begin{equation}
 \mu_0=0,
 \qquad
 \mu_1=\frac{2\alpha_1+\alpha_2}{3},
 \qquad
 \mu_2=\frac{\alpha_1+2\alpha_2}{3}.
 \label{p12-83-c20-s4:eq:representantes-discriminantes-a2}
\end{equation}
Their pairings with the simple base are
\[
\begin{array}{c|cc}
 &\alpha_1&\alpha_2\\ \hline
 \mu_1&1&0\\
 \mu_2&0&1
\end{array},
\]
thereby \(\mu_1,\mu_2\in A_2^*\). Furthermore,
\[
 3\mu_1=2\alpha_1+\alpha_2\in A_2,
 \qquad
 3\mu_2=\alpha_1+2\alpha_2\in A_2,
\]
while \(\mu_1,\mu_2\notin A_2\). It is obtained
\begin{equation}
 A_2^*/A_2
 =\{\overline\mu_0,\overline\mu_1,\overline\mu_2\}
 \cong\mathbb Z/3\mathbb Z
 \cong\mathbb F_3.
 \label{p12-83-c20-s4:eq:discriminante-a2-f3}
\end{equation}

The sum of representatives is explicitly reduced by means of
\[
 \mu_1+\mu_1-\mu_2=\alpha_1,
 \qquad
 \mu_1+\mu_2=\rho,
 \qquad
 \mu_2+\mu_2-\mu_1=\alpha_2.
\]
Consequently,
\begin{equation}
 \mu_r+\mu_s-\mu_{r+s\,({\rm mod}\,3)}\in A_2
 \qquad(r,s\in\mathbb F_3).
 \label{p12-83-c20-s4:eq:adicion-representantes-discriminantes-a2}
\end{equation}

\begin{proposition}[Discriminant Form of \(A_2\)]
\label{p12-83-c20-s4:prop:forma-discriminante-a2}
Under \eqref{p12-83-c20-s4:eq:discriminante-a2-f3}, the discriminant bilinear and quadratic
forms are
\begin{align}
 b(\overline\mu_r,\overline\mu_s)
 &\equiv\frac{2rs}{3}\pmod{\mathbb Z},
 \label{p12-83-c20-s4:eq:bilineal-discriminante-a2}\\
 q(\overline\mu_r)
 &\equiv\frac{2r^2}{3}\pmod{2\mathbb Z}.
 \label{p12-83-c20-s4:eq:cuadratica-discriminante-a2}
\end{align}
\end{proposition}

\begin{proof}
The formula \eqref{p12-83-c20-s4:eq:norma-a2-autonoma} gives
\[
 \mu_1^2=\mu_2^2=\frac23,
 \qquad
 \langle\mu_1,\mu_2\rangle=\frac13.
\]
For \(r,s\in\{0,1,2\}\), these values coincide, modulo
\(\mathbb Z\), with \(2rs/3\). The quadratic formula results by taking
\(r=s\) and recalling that the lattice is even.
\end{proof}

\subsection{\(12\) factors and gluing map}

Let
\[
 L_0:=A_2^{12}.
\]
Then
\[
 L_0^*=(A_2^*)^{12},
 \qquad
 L_0^*/L_0\cong\mathbb F_3^{12}.
\]
For \(c=(c_1,\ldots,c_{12})\in\mathbb F_3^{12}\), define
\begin{equation}
 \phi(c):=(\mu_{c_1},\ldots,\mu_{c_{12}})\in L_0^*.
 \label{p12-83-c20-s4:eq:aplicacion-pegado-cw}
\end{equation}
The extension determined by \(C_W\) is
\begin{equation}
 \begin{aligned}
 N(C_W)
 &:=\bigcup_{c\in C_W}(L_0+\phi(c))\\
 &=\{x\in L_0^*:x\bmod L_0\in C_W\}.
 \end{aligned}
 \label{p12-83-c20-s4:eq:definicion-niemeier-cw}
\end{equation}

\begin{lemma}[Additive closure of the gluing]
\label{p12-83-c20-s4:lem:clausura-pegado-cw}
The set \(N(C_W)\) is a discrete subgroup of \(L_0^*\) of rank
twenty-four.
\end{lemma}

\begin{proof}
Let \(x=\ell+\phi(c)\) and \(y=m+\phi(d)\), with \(\ell,m\in L_0\) and
\(c,d\in C_W\). By \eqref{p12-83-c20-s4:eq:adicion-representantes-discriminantes-a2},
\[
 \phi(c)+\phi(d)-\phi(c+d)\in L_0.
\]
Since \(C_W\) is linear, \(c+d\in C_W\), and then
\(x+y\in L_0+\phi(c+d)\). The same identity applied to the opposite gives
inverses. The inclusions
\[
 L_0\subseteq N(C_W)\subseteq L_0^*
\]
They test discretion and range twenty-four.
\end{proof}

\subsection{Integrality induced by self-duality}

The discriminant bilinear form of \(L_0\) is the orthogonal sum of twelve
copies of \eqref{p12-83-c20-s4:eq:bilineal-discriminante-a2}. For
\(c,d\in\mathbb F_3^{12}\),
\begin{equation}
 b(c,d)
 \equiv\frac23\sum_{i=1}^{12}c_id_i
 \pmod{\mathbb Z}.
 \label{p12-83-c20-s4:eq:bilineal-discriminante-a2-doce}
\end{equation}

\begin{lemma}[Integrity of the gluing]
\label{p12-83-c20-s4:lem:integralidad-pegado-cw}
The lattice \(N(C_W)\) is integral.
\end{lemma}

\begin{proof}
The self-duality \(C_W=C_W^\perp\) implica
\[
 \sum_{i=1}^{12}c_id_i\equiv0\pmod3
 \qquad(c,d\in C_W).
\]
By \eqref{p12-83-c20-s4:eq:bilineal-discriminante-a2-doce}, the discriminant form
vanishes on \(C_W\times C_W\). Therefore, the pairing of any
two glued classes is integral.
\end{proof}

\subsection{Parity Induced by the Code Weights}

At a nonzero coordinate of the code, the discriminant quadratic form
equals \(2/3\) modulo \(2\mathbb Z\). Therefore,
\begin{equation}
 q(c)\equiv\frac23\operatorname{wt}(c)\pmod{2\mathbb Z}.
 \label{p12-83-c20-s4:eq:cuadratica-discriminante-cw}
\end{equation}

\begin{lemma}[Parity of the gluing]
\label{p12-83-c20-s4:lem:paridad-pegado-cw}
The lattice \(N(C_W)\) is even.
\end{lemma}

\begin{proof}
The only weights of \(C_W\) are \(0,6,9,12\), and
\[
 \frac23\{0,6,9,12\}=\{0,4,6,8\}\subseteq2\mathbb Z.
\]
Every glue class is isotropic for the discriminant quadratic form, so all norms of \(N(C_W)\) are even.
\end{proof}

Let \(g_1,\ldots,g_6\) be the rows of \(G_W=(I_6\mid A_W)\) and
\(\gamma_i=\phi(g_i)\). Exact multiplication in \(A_2^*\) gives
\begin{equation}
 (\langle\gamma_i,\gamma_j\rangle)_{i,j=1}^{6}
 =
 \begin{pmatrix}
 4&2&2&2&2&2\\
 2&4&2&2&2&2\\
 2&2&4&2&2&2\\
 2&2&2&4&2&2\\
 2&2&2&2&4&2\\
 2&2&2&2&2&4
 \end{pmatrix}.
 \label{p12-83-c20-s4:eq:gram-generadores-pegado-cw}
\end{equation}
This matrix provides an additional direct control of integrality and parity for the six gluing generators.

\subsection{Index, determinant, and unimodularity}

The cosets of \(L_0\) that appear in
\eqref{p12-83-c20-s4:eq:definicion-niemeier-cw} are indexed by the
\(|C_W|=729=3^6\) elements of the code and are distinct. Consequently,
\begin{equation}
 [N(C_W):L_0]=3^6.
 \label{p12-83-c20-s4:eq:indice-niemeier-sobre-a2-doce}
\end{equation}
Since
\[
 \det L_0=(\det A_2)^{12}=3^{12},
\]
the determinant formula for a finite extension gives
\begin{equation}
 \det N(C_W)
 =\frac{\det L_0}{[N(C_W):L_0]^2}
 =\frac{3^{12}}{(3^6)^2}=1.
 \label{p12-83-c20-s4:eq:determinante-niemeier-cw}
\end{equation}

\subsection{Root system induced by Witt incidence and classification of its orbits}

En each one of the two classes discriminantes not nulas of
\(A_2^*/A_2\), the norm minima is
\[
 \mu_1^2=\mu_2^2=\frac23.
\]
If \(c\in C_W\setminus\{0\}\), then
\(\operatorname{wt}(c)\geq6\). By orthogonality of the twelve factors,
any vector of \(L_0+\phi(c)\) has norm at least
\begin{equation}
 6\cdot\frac23=4.
 \label{p12-83-c20-s4:eq:cota-clases-no-triviales-niemeier}
\end{equation}
Therefore, no nontrivial gluing class contains vectors of norm
two. The roots of \(N(C_W)\) are exactly the roots already present in
\(L_0=A_2^{12}\).

\begin{theorem}[Realization of the Niemeier of type \(A_2^{12}\)]
\label{p12-83-c20-s4:thm:niemeier-a2-doce-autonomo}
The lattice \(N(C_W)\) is positive definite, even, and unimodular, and its root
system is \(A_2^{12}\). Consequently,
\begin{equation}
 \boxed{N(C_W)\cong N(A_2^{12}).}
 \label{p12-83-c20-s4:eq:identificacion-niemeier-a2-doce}
\end{equation}
\end{theorem}

\begin{proof}
Positivity follows from \(L_0^*\). The
\cref{p12-83-c20-s4:lem:integralidad-pegado-cw,p12-83-c20-s4:lem:paridad-pegado-cw} proves
integrality and parity; \eqref{p12-83-c20-s4:eq:determinante-niemeier-cw} proves
unimodularity. The bound \eqref{p12-83-c20-s4:eq:cota-clases-no-triviales-niemeier}
identifies the root system. The Niemeier classification uniquely associates
the indicated lattice with that root system.
\end{proof}

\subsection{Obstructions to the ternary gluing \(C_W\to N(A_2^{12})\): parity, unimodularity, and root system}

\begin{criterion}[Refutation of ternary gluing]
The construction is refuted if any of the following facts occurs:
\begin{enumerate}
 \item any of the representatives \(\mu_0,\mu_1,\mu_2\) does not belong to
       \(A_2^*\), two of its classes coincide, or
       \eqref{p12-83-c20-s4:eq:adicion-representantes-discriminantes-a2} fails;
 \item the discriminant form is not the one given in
       \cref{p12-83-c20-s4:prop:forma-discriminante-a2};
 \item \(N(C_W)\) is not closed under sum or inversos;
 \item there exists \(c,d\in C_W\) with
       \(\frac23\sum_i c_id_i\notin\mathbb Z\);
 \item one palabra of \(C_W\) produces
       \(\frac23\operatorname{wt}(c)\notin2\mathbb Z\);
 \item the index is not \(3^6\) or the determinant is not one;
 \item a nontrivial gluing class contains a vector of norm two;
 \item \(W_{24}\), a dodecad, or a binary datum is used in any of
the preceding gluing operations.
\end{enumerate}
\end{criterion}


% Unidad U022. Origen incidencial, vecino ternario y Leech.
% Redaccion autonoma desde la autoridad material 1063.

\section{Incidental origin, typed radial problem and construction of the rootless unimodular neighbor}
\label{p12-83-c20-s5:chap:origen-incidencial-vecino-leech}

The lattice \(N(A_2^{12})\) has already been constructed. This unit determines a
base point from an ordered incidence frame, solves a
variational problem on a declared domain, verifies the three conditions
for the index-three neighbor, and proves that the result is the Leech
lattice. The neighbor vector is denoted \(v_{\mathcal F_*}\): it depends on the
incidence frame and does not enter the numerical carry of \(\alpha\).

\subsection{Ordered set of incidences}

Consider the four hexads
\begin{align}
 H^-&=\{4,5,6,7,9,10\},
 \label{p12-83-c20-s5:eq:hexada-marco-negativa}\\
 H_e&=\{1,3,4,6,9,10\},
 \label{p12-83-c20-s5:eq:hexada-marco-e}\\
 H_\varphi&=\{1,2,3,4,7,9\},
 \label{p12-83-c20-s5:eq:hexada-marco-phi}\\
 H_{\rm APP}&=\{2,3,5,10,11,12\}.
 \label{p12-83-c20-s5:eq:hexada-marco-app}
\end{align}
The datum is the ordered framework.
\begin{equation}
 \mathcal F:=(H^-,H_e,H_\varphi,H_{\rm APP}).
 \label{p12-83-c20-s5:eq:marco-incidencial-ordenado}
\end{equation}
We define
\begin{equation}
 o(\mathcal F)
 :=(H_e\cap H_\varphi)\setminus(H^-\cup H_{\rm APP}).
 \label{p12-83-c20-s5:eq:operador-origen-incidencial}
\end{equation}

\begin{theorem}[Equivariant Reconstruction of the Origin]
\label{p12-83-c20-s5:thm:origen-incidencial-equivariante}
For the frame canonical \(\mathcal F_*\),
\begin{equation}
 \boxed{o(\mathcal F_*)=\{1\}.}
 \label{p12-83-c20-s5:eq:origen-incidencial-singleton}
\end{equation}
Moreover, for every simultaneous permutation \(g\in M_{12}\),
\begin{equation}
 o(g\mathcal F)=g\,o(\mathcal F).
 \label{p12-83-c20-s5:eq:equivarianza-origen-incidencial}
\end{equation}
\end{theorem}

\begin{proof}
We have
\[
 H_e\cap H_\varphi=\{1,3,4,9\}
\]
and
\[
 (H^-\cup H_{\rm APP})\cap\{1,3,4,9\}=\{3,4,9\}.
\]
The difference is \(\{1\}\). For equivariance one applies the
identities
\[
 g(A\cap B)=gA\cap gB,
 \quad
 g(A\cup B)=gA\cup gB,
 \quad
 g(A\setminus B)=gA\setminus gB,
\]
valid for all bijections.
\end{proof}

The tetrad \(B_K=\{4,6,9,10\}\) does not by itself determine the origin. The
singleton is obtained from the complete ordered frame; it is not a
coordinate chosen after the result has been observed.

\subsection{Positive radial chamber of the root system and selection of the dominant cone}

Let \(\rho_i\) the copia of \(\rho=\alpha_1+\alpha_2\) in the
\(i\)-esimo factor of \(A_2^{12}\). Consideramos the dominio
\begin{equation}
 \mathscr R_+
 :=\left\{
 v(a)=\sum_{i=1}^{12}a_i\rho_i:
 a_i=1+3b_i, b_i\in\mathbb Z_{\geq0}
 \right\}.
 \label{p12-83-c20-s5:eq:camara-radial-positiva}
\end{equation}
Since the factors are orthogonal and \(\rho_i^2=2\), if
\(S=\sum_i b_i\), then
\begin{align}
 v(a)^2
 &=2\sum_{i=1}^{12}(1+3b_i)^2\notag\\
 &=24+12S+18\sum_{i=1}^{12}b_i^2.
 \label{p12-83-c20-s5:eq:norma-radial-expandida}
\end{align}
Reducing modulo eighteen,
\[
 v(a)^2\equiv6+12S\pmod{18},
\]
and therefore
\begin{equation}
 v(a)^2\equiv0\pmod{18}
 \quad\Longleftrightarrow\quad
 S\equiv1\pmod3.
 \label{p12-83-c20-s5:eq:congruencia-radial-suma-b}
\end{equation}

\begin{proposition}[Minima in the declared radial chamber]
\label{p12-83-c20-s5:prop:minimos-camara-radial}
On the subset of \(\mathscr R_+\) defined by
\(v(a)^2\equiv0\pmod{18}\), the minimum norm is \(54\). It is attained
exactly at twelve vectors: for each \(j\in\{1,\ldots,12\}\), the
coefficient \(a_j\) equals four, and the remaining eleven equal one.
\end{proposition}

\begin{proof}
Condition \eqref{p12-83-c20-s5:eq:congruencia-radial-suma-b} and the nonnegativity of
the \(b_i\) imply that the smallest possible value of \(S\) is one. This
forces a unique index \(j\) with \(b_j=1\) and \(b_i=0\) for
\(i\neq j\). Substituting into \eqref{p12-83-c20-s5:eq:norma-radial-expandida},
\[
 v(a)^2=24+12+18=54.
\]
There are twelve choices of \(j\). The next admissible value of \(S\) is
four and produces a strictly greater norm.
\end{proof}

The origin \(o(\mathcal F_*)=\{1\}\) selects
\begin{equation}
 \boxed{
 v_{\mathcal F_*}
 =4\rho_1+\rho_2+\cdots+\rho_{12},
 \qquad
 v_{\mathcal F_*}^2=54.}
 \label{p12-83-c20-s5:eq:vector-marco-incidencial}
\end{equation}

\begin{remark}[Exact scope of minimality]
The norm \(54\) is minimal in the positive radial chamber
\eqref{p12-83-c20-s5:eq:camara-radial-positiva}, not in the entire residue class. The vector
\[
 u=(\alpha_1-2\alpha_2,\rho,\ldots,\rho)
\]
satisfies
\[
 \alpha_1-2\alpha_2-\rho=-3\alpha_2,
\]
so that
\[
 u\equiv(\rho,\ldots,\rho)
 \equiv v_{\mathcal F_*}\pmod{3A_2^{12}}.
\]
However,
\[
 (\alpha_1-2\alpha_2)^2
 =2(1+2+4)=14,
 \qquad
 u^2=14+11\cdot2=36.
\]
This witness prevents extending the minimality assertion outside the domain
of \cref{p12-83-c20-s5:prop:minimos-camara-radial}.
\end{remark}

\subsection{Character and congruence kernel}

We write
\[
 N:=N(A_2^{12}),
 \qquad
 v:=v_{\mathcal F_*}.
\]
We define
\begin{equation}
 \chi_v:N\longrightarrow\mathbb Z/3\mathbb Z,
 \qquad
 \chi_v(x)=\langle x,v\rangle\bmod3,
 \label{p12-83-c20-s5:eq:caracter-vecino-leech}
\end{equation}
its nucleus
\begin{equation}
 N_0:=\ker\chi_v
 =\{x\in N:\langle x,v\rangle\equiv0\pmod3\},
 \label{p12-83-c20-s5:eq:nucleo-congruencia-vecino}
\end{equation}
and the extension
\begin{equation}
 N_v:=N_0+\mathbb Z\frac v3.
 \label{p12-83-c20-s5:eq:definicion-vecino-indice-tres}
\end{equation}

\subsection{Condition I1: surjectivity and exclusion of roots}

By construction, \(v\in A_2^{12}\subseteq N\). The roots of a factor
\(A_2\) are
\[
 \pm\alpha_1,
 \quad
 \pm\alpha_2,
 \quad
 \pm\rho.
\]
Their pairings with \(\rho\) are \(\pm1\) or \(\pm2\). All the
coefficients of \(v\) are congruent to one modulo three:
\[
 4\equiv1\pmod3,
 \qquad
 1\equiv1\pmod3.
\]
Therefore, no root of \(A_2^{12}\) has zero pairing with
\(v\) modulo three. In particular,
\[
 \chi_v(\alpha_{1,1})=1,
\]
so that \(\chi_v\) is surjective and
\begin{equation}
 [N:N_0]=3.
 \label{p12-83-c20-s5:eq:indice-nucleo-vecino}
\end{equation}
Moreover, no root of \(N\) belongs to \(N_0\).

\subsection{Condition I2: norm divisibility}

The identity \eqref{p12-83-c20-s5:eq:vector-marco-incidencial} gives
\begin{equation}
 v^2=54\equiv0\pmod{18},
 \qquad
 \left(\frac v3\right)^2=6.
 \label{p12-83-c20-s5:eq:condicion-i2-vecino}
\end{equation}
In particular, \(\langle v,v\rangle\equiv0\pmod3\), hence
\(v\in N_0\).

\subsection{Condition I3: class dual of order \(3\)}

If \(x\in N_0\), there exists \(k\in\mathbb Z\) with
\(\langle x,v\rangle=3k\). Therefore,
\[
 \left\langle x,\frac v3\right\rangle=k\in\mathbb Z,
\]
and
\begin{equation}
 \frac v3\in N_0^*.
 \label{p12-83-c20-s5:eq:v-tercio-en-dual-nucleo}
\end{equation}
On the other hand,
\[
 \left\langle\frac v3,\alpha_{1,1}\right\rangle=\frac43.
\]
Since \(\alpha_{1,1}\in N\) and \(N\) is integral, this proves
\begin{equation}
 \frac v3\notin N.
 \label{p12-83-c20-s5:eq:v-tercio-no-en-niemeier}
\end{equation}
Finally, \(3(v/3)=v\in N_0\). The class of \(v/3\) in
\(N_0^*/N_0\) is nontrivial and has order exactly three. Consequently,
\begin{equation}
 [N_v:N_0]=3.
 \label{p12-83-c20-s5:eq:indice-extension-vecino}
\end{equation}

\subsection{Parity, integrality and determinant of the neighbor}

\begin{proposition}[Unimodular structure of the neighbor]
\label{p12-83-c20-s5:prop:estructura-unimodular-vecino}
The lattice \(N_v\) is integral, even, and unimodular.
\end{proposition}

\begin{proof}
Since \(N\) is unimodular and \([N:N_0]=3\),
\[
 \det N_0=[N:N_0]^2\det N=9.
\]
Since \([N_v:N_0]=3\),
\[
 \det N_v=\frac{\det N_0}{[N_v:N_0]^2}=1.
\]

Let \(y=x+m(v/3)\), with \(x\in N_0\), and write
\(\langle x,v\rangle=3k\). Then
\begin{align*}
 y^2
 &=x^2+2m\left\langle x,\frac v3\right\rangle
       +m^2\left(\frac v3\right)^2\\
 &=x^2+2mk+6m^2.
\end{align*}
The three summands form an even integer. The polarized form proves
integrality, and determinant one proves unimodularity.
\end{proof}

\subsection{\(2\) clasis globalis and \(6\) clasis local}

Globally, \eqref{p12-83-c20-s5:eq:definicion-vecino-indice-tres} adds only two nontrivial
cosets:
\[
 N_0+\frac v3,
 \qquad
 N_0-\frac v3.
\]
Locally, in each factor \(A_2\), a coordinate of those classes belongs to one of the six classes
\begin{equation}
 \mathscr C_{j,\pm}
 :=A_2+\mu_j\pm\frac\rho3,
 \qquad j\in\{0,1,2\}.
 \label{p12-83-c20-s5:eq:seis-clases-locales-vecino}
\end{equation}
In the first coordinate, \(4\rho/3\equiv\rho/3\pmod{A_2}\), so
the same list holds for the twelve factors.

\begin{table}[htbp]
\centering
\caption{Minimal representatives of the six local cosets.}
\label{p12-83-c20-s5:tab:seis-clases-locales-vecino}
\begin{tabular}{cclc}
\toprule
class&minimal representative&\((u,v)\)&
\(2(u^2-uv+v^2)/9\)\\
\midrule
\(\mathscr C_{0,-}\)&\(-\rho/3\)&\((-1,-1)\)&\(2/9\)\\
\(\mathscr C_{0,+}\)&\( \rho/3\)&\((1,1)\)&\(2/9\)\\
\(\mathscr C_{1,-}\)&\( \mu_1-\rho/3\)&\((1,0)\)&\(2/9\)\\
\(\mathscr C_{1,+}\)&\( \mu_1+\rho/3-\rho\)&\((0,-1)\)&\(2/9\)\\
\(\mathscr C_{2,-}\)&\( \mu_2-\rho/3\)&\((0,1)\)&\(2/9\)\\
\(\mathscr C_{2,+}\)&\( \mu_2+\rho/3-\rho\)&\((-1,0)\)&\(2/9\)\\
\bottomrule
\end{tabular}
\end{table}

Here \((u,v)\) means the representative
\((u\alpha_1+v\alpha_2)/3\).

\begin{lemma}[Exact minimum of the six classes]
\label{p12-83-c20-s5:lem:minimo-seis-clases-locales}
Each class in \eqref{p12-83-c20-s5:eq:seis-clases-locales-vecino} has minimum norm
exactly \(2/9\).
\end{lemma}

\begin{proof}
The norm of \((u\alpha_1+v\alpha_2)/3\) is
\[
 \frac29(u^2-uv+v^2).
\]
The integral form \(u^2-uv+v^2\) is positive definite. On a nonzero
residual class it does not take the value zero, and hence is at least one. The six
representatives in \cref{p12-83-c20-s5:tab:seis-clases-locales-vecino} attain one.
\end{proof}

\subsection{Criterion for the absence of roots in the Leech lattice and consequence for the minimum distance}

\begin{theorem}[Absence of vectors of norm two]
\label{p12-83-c20-s5:thm:vecino-sin-raices}
The lattice \(N_v\) contains no vectors of norm two.
\end{theorem}

\begin{proof}
All roots of \(N\) are those of \(A_2^{12}\). Condition I1 proves
that none belongs to \(N_0\). Hence class \(N_0\) contains no
vectors of norm two.

In a new class \(N_0\pm v/3\), each of the twelve coordinates
belongs to one of the six local classes. By
\cref{p12-83-c20-s5:lem:minimo-seis-clases-locales}, each coordinate contributes at least
\(2/9\). Orthogonality gives
\[
 \left\|x\pm\frac v3\right\|^2
 \geq12\cdot\frac29=\frac83>2.
\]
Nor do the new classes contain roots.
\end{proof}

\subsection{Identification with the Leech lattice}

\begin{theorem}[Construction of the Leech lattice]
\label{p12-83-c20-s5:thm:identificacion-vecino-leech}
The lattice
\begin{equation}
 N_v
 =\ker(\langle\,\cdot\,,v_{\mathcal F_*}\rangle\bmod3)
  +\mathbb Z\frac{v_{\mathcal F_*}}3
 \label{p12-83-c20-s5:eq:formula-final-vecino-leech}
\end{equation}
is isometric to the Leech lattice:
\begin{equation}
 \boxed{N_v\cong\Lambda_{24}.}
 \label{p12-83-c20-s5:eq:identificacion-leech-autonoma}
\end{equation}
\end{theorem}

\begin{proof}
The lattice \(N_v\) has rank twenty-four and is positive definite.
\cref{p12-83-c20-s5:prop:estructura-unimodular-vecino} proves that it is even and unimodular;
\cref{p12-83-c20-s5:thm:vecino-sin-raices} proves that it has no roots. By the
Niemeier classification, there exists a unique isometry class of positive-definite,
even, unimodular lattices of rank twenty-four without roots:
the Leech lattice.
\end{proof}

\subsection{The 243 syndromes of the ternary punctured code}

Let \(G_{11}\) be the ternary Golay code with parameters
\([11,6,5]_3\). Its syndrome space has cardinality
\begin{equation}
 3^{11-6}=3^5=243.
 \label{p12-83-c20-s5:eq:numero-sindromes-golay-ternario}
\end{equation}
The number of weight errors of at most two is
\begin{equation}
 1+2\binom{11}{1}+2^2\binom{11}{2}
 =1+22+220=243.
 \label{p12-83-c20-s5:eq:bola-radio-dos-golay-ternario}
\end{equation}

\begin{proposition}[Perfect correspondence error--syndrome]
\label{p12-83-c20-s5:prop:correspondencia-error-sindrome-golay}
Each syndrome of \(G_{11}\) has a unique error representative of weight at most two.
\end{proposition}

\begin{proof}
Minimum distance five implies that two distinct errors of weight at
most two cannot differ by a codeword: their difference would have
weight at most four. The error--syndrome map is therefore
injective on the ball of radius two. Both sets have cardinality
\(243\) by \eqref{p12-83-c20-s5:eq:numero-sindromes-golay-ternario} and
\eqref{p12-83-c20-s5:eq:bola-radio-dos-golay-ternario}; the map is bijective.
\end{proof}

These \(243\) syndromes are not the \(243\) visible words obtained from the
TPK catalogue. They share a cardinal, but they possess distinct domains, equivalences and
functions.

\subsection{Obstructions of the ternary neighbor \(N_v\cong\Lambda_{24}\): class of order \(3\), absence of roots, and separation of the cardinals \(243\)}

\begin{criterion}[Refutation of the neighbor and of the Leech identification]
The construction is refuted if any of the following facts occurs:
\begin{enumerate}
 \item the operator \(o(\mathcal F)\) not produces \(\{1\}\) or is not
       equivariant;
 \item there exists in the radial chamber an admissible vector of norm less than
       \(54\), or the number of minimizers is not twelve;
 \item global minimality is asserted despite the witness of norm \(36\);
 \item \(\chi_v\) is not sobreyectivo or one root belongs to \(N_0\);
 \item \(v^2\not\equiv0\pmod{18}\) o \((v/3)^2\neq6\);
 \item \(v/3\notin N_0^*\), \(v/3\in N\), or its class does not have order three;
 \item the determinant of \(N_v\) is not one, or the neighbor ceases to be even;
 \item some of the six clasis local has minimum least that \(2/9\);
 \item one of the two new global classes contains a root;
 \item the two global classes are confused with the six local classes;
 \item \(W_{24}\) is used as an antecedent of the gluing or of the neighbor;
 \item the \(243\) syndromes are identified with the \(243\) words
       visible by sharing cardinality.
\end{enumerate}
\end{criterion}


\section{Finite holonomy, Schläfli configuration and symmetry \texorpdfstring{\(E_6\)}{E6}}
\label{p12-83-c20-s6:chap:holonomia-schlafli-e6-nuevo}
\index{finite holonomy}
\index{Schläfli configuration}
\index{E6@\(E_6\)}

The Witt branch does not exhaust the exceptional structures accessible from
APP. The parallel branch of this chapter starts from the two mixed
sum--product and product--sum polarizations, together with a plaquette
orientation and a declared central section. The calculated curvature determines
the alternating form; this defines the central extension. An
extension is not inferred from cardinality 27 alone.

\subsection{Mixed APP curvatures}

On \(T_9=(\mathbb Z/9\mathbb Z)^2\), with
\(S(x,y)=x+y\) and \(P(x,y)=xy\), let
\[
 \Delta_xf=f(x+1,y)-f(x,y),\qquad
 \Delta_yf=f(x,y+1)-f(x,y).
\]
Mixed connections are
\[
 A^{SP}=(\Delta_xS,\Delta_yP)=(1,x),\qquad
 A^{PS}=(\Delta_xP,\Delta_yS)=(y,1).
\]
Its discrete curvature is
\[
 F_A=\Delta_xA_y-\Delta_yA_x.
\]

\begin{proposition}[Oriented opposite curvatures]
\label{p12-83-c20-s6:prop:curvaturas-app-e6-nuevo}
\[
 F_{A^{SP}}=1,\qquad F_{A^{PS}}=-1\pmod9.
\]
The circulation over an oriented rectangle of sides \(r,s\) is,
respectively, \(rs\) and \(-rs\).
\end{proposition}

\begin{proof}
\(\Delta_xx-\Delta_y1=1\) and
\(\Delta_x1-\Delta_yy=-1\). A rectangle contains \(rs\) plaquettes; upon
summing, the interior edges cancel and the boundary circulation remains.
\end{proof}

For displazamientos \(u=(a,b)\), \(v=(c,d)\), the class alternante is
\begin{equation}
 \sigma(u,v)=ad-bc\pmod9.
 \label{p12-83-c20-s6:eq:forma-alternante-e6-nueva}
\end{equation}
The orientation fixes the sign of \(\sigma\). The central section used
below belongs to the gauge of this realization; replacing it with a coboundary
produces an isomorphic extension, not a new invariant.

\subsection{Central extension and ternary reduction}

We define the nonadic Heisenberg group
\begin{equation}
 \mathcal H_9=\{(a,b,z):a,b,z\in\mathbb Z/9\mathbb Z\},
 \quad
 (a,b,z)(c,d,w)=(a+c,b+d,z+w+bc).
 \label{p12-83-c20-s6:eq:heisenberg-9-nuevo}
\end{equation}
The commutator of two horizontal shifts is
\[
 [(a,b,0),(c,d,0)]=(0,0,bc-ad),
\]
which records \(-\sigma(u,v)\). Reducing each coordinate modulo three
gives
\begin{equation}
 H_3=\{(a,b,z):a,b,z\in\mathbb F_3\},\qquad |H_3|=27,
 \label{p12-83-c20-s6:eq:heisenberg-3-nuevo}
\end{equation}
with the same law. The nucleo of
\(\mathcal H_9\to H_3\) is
\((3\mathbb Z/9\mathbb Z)^3\), of cardinal 27.

\subsection{Weyl--Heisenberg representation and central phase}

The nonadic extension admits an explicit unitary realization. Let
\[
 \zeta_9:=\exp(2\pi i/9),
 \qquad
 \mathscr H_9^{\rm Sch}:=\mathbb C^9,
\]
based on \(\{|n\rangle:n\in\mathbb Z/9\mathbb Z\}\). We define
\begin{equation}
 X|n\rangle:=|n+1\rangle,
 \qquad
 Z|n\rangle:=\zeta_9^n|n\rangle.
 \label{p12-83-c20-s6:eq:weyl-desplazamiento-fase-u029}
\end{equation}
We have
\[
 ZX=\zeta_9XZ
\]
and, by multiplication,
\begin{equation}
 (X^aZ^b)(X^cZ^d)
 =\zeta_9^{bc}X^{a+c}Z^{b+d}.
 \label{p12-83-c20-s6:eq:cociclo-weyl-nonadico-u029}
\end{equation}

\begin{theorem}[Faithful Weyl--Heisenberg Representation]
\label{p12-83-c20-s6:thm:representacion-weyl-heisenberg-u029}
The map
\begin{equation}
 \rho_{\rm WH}:\mathcal H_9\longrightarrow
 U(\mathscr H_9^{\rm Sch}),
 \qquad
 \rho_{\rm WH}(a,b,z):=\zeta_9^zX^aZ^b,
 \label{p12-83-c20-s6:eq:representacion-weyl-heisenberg-u029}
\end{equation}
is a faithful unitary representation. For
\(u=(a,b)\) and \(v=(c,d)\), its commutator satisfies
\begin{equation}
 [\rho_{\rm WH}(u,0),\rho_{\rm WH}(v,0)]
 =\zeta_9^{\,bc-ad}I
 =\zeta_9^{-\sigma(u,v)}I.
 \label{p12-83-c20-s6:eq:conmutador-weyl-holonomia-u029}
\end{equation}
\end{theorem}

\begin{proof}
The identity \eqref{p12-83-c20-s6:eq:cociclo-weyl-nonadico-u029} reproduces the term
\(bc\) of \eqref{p12-83-c20-s6:eq:heisenberg-9-nuevo}; hence the
homomorphism property. If \(\zeta_9^zX^aZ^b=I\), the action on the basis successively forces
\(a=0\), \(b=0\), and \(z=0\). This proves faithfulness. Comparing
the products in the two orders gives
\(\zeta_9^{bc-ad}\), which is the alternating phase of
\eqref{p12-83-c20-s6:eq:forma-alternante-e6-nueva}.
\end{proof}

The representation distinguishes two conjugate observables. Let
\[
 Q_{\rm H}:=\operatorname{diag}(0,1,\ldots,8),
 \qquad
 F_{jk}:=\frac13\zeta_9^{jk},
 \qquad
 P_{\rm H}:=F^\dagger Q_{\rm H}F.
\]
The eigenbases of \(Q_{\rm H}\) and \(P_{\rm H}\) are mutually
unbiased, since \(|F_{jk}|^2=1/9\). For every unit vector
\(\psi\in\mathscr H_9^{\rm Sch}\),
\begin{align}
 \operatorname{Var}_\psi(Q_{\rm H})
 \operatorname{Var}_\psi(P_{\rm H})
 &\ge
 \frac14\left|
 \langle\psi,[Q_{\rm H},P_{\rm H}]\psi\rangle
 \right|^2,
 \label{p12-83-c20-s6:eq:incertidumbre-weyl-u029}\\
 H_{Q_{\rm H}}(\psi)+H_{P_{\rm H}}(\psi)
 &\ge\log9.
 \label{p12-83-c20-s6:eq:incertidumbre-entropica-weyl-u029}
\end{align}
These operators are dimensionless. A physical unit chart is a
later realization and does not intervene in the central extension.

\subsection{Action of \(10\) displazamientos over the grafo of Cayley}

For a linear functional
\(\ell(a,b)=pa+qb\), set
\[
 D_\ell(a,b)=\left(a,b,\frac{ab}{2}+\ell(a,b)\right),
 \qquad Z=(0,0,1),
\]
where \(2^{-1}=2\) in \(\mathbb F_3\). The inverse-closed set
\begin{equation}
 S_\ell=
 \{D_\ell(v):v\in\mathbb F_3^2\setminus\{0\}\}
 \cup\{Z,Z^{-1}\}
 \label{p12-83-c20-s6:eq:conjunto-conexion-e6}
\end{equation}
has ten elements and does not contain the identity.

\begin{theorem}[Independence of the linear scale]
\label{p12-83-c20-s6:thm:independencia-calibre-e6}
The nine functionals \(\ell\) produce isomorphic Cayley graphs. The
nine sets \(S_\ell\) form a single orbit under
\(\operatorname{Aut}(H_3)\).
\end{theorem}

\begin{proof}
The map
\[
 \phi_\ell(a,b,z)=(a,b,z+\ell(a,b))
\]
is a central automorphism and satisfies
\(\phi_\ell(S_0)=S_\ell\). To write the complete family in the
polarized coordinates of \cref{p12-83-c20-s6:eq:heisenberg-9-nuevo}, let
\(c(v,w)=b c\) for \(v=(a,b)\), \(w=(c,d)\). Given
\(M\in\operatorname{GL}(2,3)\), the defect
\[
 B_M(v,w)=c(Mv,Mw)-\det(M)c(v,w)
\]
is symmetric, because the alternating parts of the two terms coincide. We define the quadratic correction
\begin{equation}
 q_M(v)=2B_M(v,v),
 \qquad
 q_M(v+w)-q_M(v)-q_M(w)=B_M(v,w),
 \label{p12-83-c20-s6:eq:correccion-cuadratica-automorfismos-h3}
\end{equation}
where \(2=2^{-1}\) in \(\mathbb F_3\). The complete family of
automorphisms is then written
\begin{equation}
 (v,t)\longmapsto
 \bigl(Mv,\det(M)t+q_M(v)+\ell(v)\bigr),
 \qquad M\in\operatorname{GL}(2,3),\
 \ell\in(\mathbb F_3^2)^*.
 \label{p12-83-c20-s6:eq:automorfismos-h3-completos}
\end{equation}
The polar identity of \cref{p12-83-c20-s6:eq:correccion-cuadratica-automorfismos-h3}
directly proves that the central product is preserved. The family has
\(48\cdot9=432\) elements. Exhaustive enumeration of the
\(\binom{13}{5}=1287\) inverse-closed subsets of cardinality ten
finds exactly those nine with the parameters derived
below.
\end{proof}

Write \(\underline S=\sum_{s\in S}s\) in
\(\mathbb Z[H_3]\). Counting ordered products gives
\begin{equation}
 \underline S^{\,2}
 =10e+\underline S+5(\underline{H_3}-e-\underline S)
 =5\underline{H_3}+5e-4\underline S.
 \label{p12-83-c20-s6:eq:identidad-grupo-schlafli-nueva}
\end{equation}
Indeed, the identity receives ten factorizations; an element of
\(S\) receives one; each of the other sixteen receives five.

\begin{theorem}[Graph of the twenty-seven lines]
\label{p12-83-c20-s6:thm:grafo-27-rectas-nuevo}
The matrix of adjacency \(A_\ell\) of
\(\Gamma_\ell=\operatorname{Cay}(H_3,S_\ell)\) satisfies
\begin{equation}
 A_\ell^2=5I-4A_\ell+5J.
 \label{p12-83-c20-s6:eq:identidad-schlafli-nueva}
\end{equation}
Therefore,
\[
 \Gamma_\ell=\operatorname{srg}(27,10,1,5),
 \qquad
 \operatorname{Spec}(A_\ell)=\{10^1,1^{20},(-5)^6\}.
\]
It is isomorphic to the intersection graph of the twenty-seven lines of a
smooth cubic surface.
\end{theorem}

\begin{proof}
The regular representation of
\cref{p12-83-c20-s6:eq:identidad-grupo-schlafli-nueva} produces
\cref{p12-83-c20-s6:eq:identidad-schlafli-nueva}. Its entries state that every vertex
has ten neighbors, that an adjacent pair has one common neighbor, and a
nonadjacent pair has five. On the constant line, \(A_\ell\) acts by 10;
on its complement, \(J=0\) and
\((A_\ell-I)(A_\ell+5I)=0\). Zero trace fixes the multiplicities 20 and 6.

In the model of a cubic obtained by blowing up six points of
\(\mathbb P^2\), the 27 classes are
\[
 E_i,\qquad L_{ij}=H-E_i-E_j,\qquad
 Q_i=2H-\sum_{j\ne i}E_j.
\]
With \(H^2=1\), \(E_i^2=-1\), \(H\cdot E_i=0\), and
\(E_i\cdot E_j=0\) for \(i\ne j\), its intersection matrix has the
same parameters. Identification with the classical configuration uses the
uniqueness theorem for this realization of the 27 lines
\cite{Manivel2005}; it is not attributed to a nonexistent certificate.
\end{proof}

The isomorphism assertion can be made fully effective by means of
a computed bijection. The labeling is not absolute: another linear section
or an automorphism of the surface produces another table in the same class.
\begin{table}[!htbp]
\centering\fontsize{9.4}{10.8}\selectfont
\setlength{\tabcolsep}{4pt}\renewcommand{\arraystretch}{1.04}
\caption{Explicit bijection between \(H_3\) and the twenty-seven classes of
lines.}
\label{p12-83-c20-s6:tab:h3-27-rectas-u029}
\begin{tabular}{@{}c@{\,}c@{\,}c@{\qquad}l@{}}

\toprule
\(a\)&\(b\)&\(z\)&class of line\\
\midrule

0&0&0&\(E_1\)\\
0&0&1&\(L_{12}\)\\
0&0&2&\(Q_2\)\\
0&1&0&\(L_{13}\)\\
0&1&1&\(Q_1\)\\
0&1&2&\(E_3\)\\
0&2&0&\(Q_3\)\\
0&2&1&\(E_2\)\\
0&2&2&\(L_{23}\)\\
1&0&0&\(L_{14}\)\\
1&0&1&\(L_{35}\)\\
1&0&2&\(L_{26}\)\\
1&1&0&\(L_{36}\)\\
1&1&1&\(L_{24}\)\\
1&1&2&\(L_{15}\)\\
1&2&0&\(L_{25}\)\\
1&2&1&\(L_{16}\)\\
1&2&2&\(L_{34}\)\\
2&0&0&\(Q_4\)\\
2&0&1&\(L_{46}\)\\
2&0&2&\(E_6\)\\
2&1&0&\(E_5\)\\
2&1&1&\(Q_6\)\\
2&1&2&\(L_{56}\)\\
2&2&0&\(L_{45}\)\\
2&2&1&\(E_4\)\\
2&2&2&\(Q_5\)\\
\bottomrule
\end{tabular}
\end{table}

As a local control, the ten neighbors of \(e=(0,0,0)\) are the elements of
\(S_0\). The table transforms them into the ten classes that intersect \(E_1\):
the five \(L_{1j}\) and the five \(Q_j\), with \(2\le j\le6\). The
complete verification compares all 729 entries of both adjacency
matrices.

The complement \(B_\ell=J-I-A_\ell\) satisfies
\begin{equation}
 B_\ell^2=8I+2B_\ell+8J,
 \qquad
 \operatorname{srg}(27,16,10,8),
 \qquad
 \operatorname{Spec}(B_\ell)=\{16^1,4^6,(-2)^{20}\}.
 \label{p12-83-c20-s6:eq:complemento-schlafli}
\end{equation}
This is the Schläfli graph in the convention of complementary adjacency;
\(A_\ell\) is the intersection graph with parameters
\((27,10,1,5)\).

Each edge belongs to a unique triangle, since \(\lambda=1\). The number
of triangles is
\[
 \frac{27\cdot10}{6}=45.
\]
The full automorphism group has order
\[
 27\cdot1920=51840
\]
and, through its action on the line configuration, is identified with
\(W(E_6)\), not merely by equality of orders \cite{Manivel2005}.

\subsection{The root lattice realizing \texorpdfstring{\(E_6\)}{E6}}

To make explicit the object on which that group acts, let
\begin{equation}
 \operatorname{Pic}(X)
 =\mathbb ZH\oplus\bigoplus_{i=1}^6\mathbb ZE_i,
 \qquad
 H^2=1,\qquad E_i^2=-1,\qquad H\cdot E_i=E_i\cdot E_j=0\ (i\ne j),
 \label{p12-83-c20-s6:eq:picard-cubica-e6}
\end{equation}
the Picard lattice of the blow-up of six general points of
\(\mathbb P^2\). Its canonical class is
\[
 K_X=-3H+E_1+\cdots+E_6.
\]
The orthogonal complement
\begin{equation}
 L_{E_6}=K_X^\perp\subset\operatorname{Pic}(X),
 \qquad (x,y)_{E_6}:=-x\cdot y,
 \label{p12-83-c20-s6:eq:reticulo-raices-e6}
\end{equation}
is positive definite of rank six. A simple root basis is
\begin{equation}
 \begin{aligned}
 \alpha_1&=E_1-E_2,& \alpha_2&=E_2-E_3,&
 \alpha_3&=E_3-E_4,\\
 \alpha_4&=E_4-E_5,& \alpha_5&=E_5-E_6,&
 \alpha_6&=H-E_1-E_2-E_3.
 \end{aligned}
 \label{p12-83-c20-s6:eq:raices-simples-e6}
\end{equation}
Each \(\alpha_i\) is orthogonal to \(K_X\), has norm two, and its products
are \(-1\) exactly for adjacent pairs of the diagram
\(E_6\): the chain \(1-2-3-4-5\), with node \(6\) joined to node \(3\).
Therefore, its Gram matrix is the Cartan matrix of \(E_6\).

\begin{proposition}[The seventy-two roots and their reflections]
\label{p12-83-c20-s6:prop:setenta-dos-raices-e6}
The roots of \(L_{E_6}\) are
\begin{equation}
 \begin{aligned}
 &E_i-E_j &&(i\ne j),\\
 &\pm(H-E_i-E_j-E_k) &&(1\le i<j<k\le6),\\
 &\pm(2H-E_1-\cdots-E_6),
 \end{aligned}
 \label{p12-83-c20-s6:eq:enumeracion-raices-e6}
\end{equation}
in total \(30+40+2=72\). For each root \(\alpha\),
\begin{equation}
 s_\alpha(x)=x+(x\cdot\alpha)\alpha
 \label{p12-83-c20-s6:eq:reflexion-raiz-e6}
\end{equation}
preserves the intersection and \(K_X\). The six simple reflections generate
\(W(E_6)\) and act faithfully by permutations on the 27 classes
\(E_i,L_{ij},Q_i\) in \cref{p12-83-c20-s6:thm:grafo-27-rectas-nuevo}.
\end{proposition}

\begin{proof}
Substituting each class of \cref{p12-83-c20-s6:eq:enumeracion-raices-e6} into
\cref{p12-83-c20-s6:eq:picard-cubica-e6} gives \(K_X\cdot\alpha=0\) and
\(\alpha^2=-2\). The class count is as stated. The reflection formula
is the orthogonal reflection for the form \(-(\cdot)\); it preserves \(K_X\) because
\(K_X\cdot\alpha=0\). The simple system of
\cref{p12-83-c20-s6:eq:raices-simples-e6} generates all roots, and its Coxeter group is
\(W(E_6)\). Its action on the 27 classes of exceptional curves preserves
intersection, and therefore coincides with the action on the already
constructed graph \cite{Manivel2005}.
\end{proof}

\subsection{Stratification of the 729 Words}

We write a six-trit word as
\[
 w=(\ell_1,h_1,\ell_2,h_2,\ell_3,h_3),
\]
and we define two elements of \(H_3\),
\[
 L(w)=(\ell_1,\ell_2,\ell_3),\qquad
 H(w)=(h_1,h_2,h_3),
\]
using the central law in their products. According to
\(L(w)^{-1}H(w)\), the 729 words are divided into
\begin{align*}
 \mathcal D_\ell&=\{w:L(w)^{-1}H(w)=e\},\\
 \mathcal I_\ell&=\{w:L(w)^{-1}H(w)\in S_\ell\},\\
 \mathcal S_\ell&=\{w:L(w)^{-1}H(w)\notin\{e\}\cup S_\ell\}.
\end{align*}

\begin{theorem}[Hensel--Schläfli Partition]
\label{p12-83-c20-s6:thm:particion-hensel-schlafli-nueva}
\begin{equation}
 729=27+270+432,
 \label{p12-83-c20-s6:eq:particion-729-e6-nueva}
\end{equation}
corresponding to diagonal, incidence, and warping.
\end{theorem}

\begin{proof}
For each one of the 27 values of \(L(w)\), there is a unique value of
\(H(w)\) with identity difference, ten with difference in \(S_\ell\) and
sixteen remaining ones. \(27\) is multiplied by \(1,10,16\).
\end{proof}

The cardinal \(432\) of the warped stratum numerically coincides with the
multiplicity of the transverse region of \(\pi\) in
\cref{eq:descomposiciones-cinco-regiones-pi}. The two objects are not thereby identified: here
one counts words classified by \(L(w)^{-1}H(w)\), whereas there one
counts preimages of a register \(U_6\). Without an explicit and
equivariant intertwiner between those domains, \(432=432\) is only a cardinal control.

\subsection{Necessity of the central phase}

Enumerating all inverse-closed subsets of cardinality ten in the five groups of order 27 gives
\[
\begin{array}{c|c|c}
\text{group}&\text{abelian}&
\#\operatorname{srg}(27,10,1,5)\\ \hline
C_{27}&\text{yes}&0\\
C_9\times C_3&\text{yes}&0\\
C_3^3&\text{yes}&0\\
C_9\rtimes C_3&\text{no}&9\\
H_3&\text{no}&9.
\end{array}
\]
Thus, retaining 27 labels while abelianizing composition destroys the
configuration. Neither cardinality nor degree suffices.

If \(\mathcal T\) is the conjunto of 45 triangulos, the cubico
\[
 \mathcal C_\Gamma(x_1,\ldots,x_{27})
 =\sum_{\{i,j,k\}\in\mathcal T}x_ix_jx_k
\]
recovers adjacency, because every edge lies in a unique triangle. Its
permutation stabilizer is exactly
\[
 \operatorname{Aut}(\Gamma_\ell)\cong W(E_6).
\]

\begin{figure}[htbp]
\centering
\[
 \text{APP curvatures }\pm1
 \longrightarrow\sigma
 \longrightarrow\mathcal H_9
 \longrightarrow H_3
 \longrightarrow\Gamma_{27}
 \longrightarrow\mathcal T_{45}
 \longrightarrow W(E_6).
\]
\caption{Causal chain of the mixed-holonomy branch. It is parallel to
\(C_W\to W_{12}\to M_{12}\) and not a link of it.}
\label{p12-83-c20-s6:fig:cadena-holonomia-e6-u029}
\end{figure}

\begin{criterion}[Falsification criteria for branch \(E_6\)]
\label{p12-83-c20-s6:crit:refutacion-rama-e6-u029}
The branch is refuted if the curvatures are not opposite, if the central cocycle is erased, if the quadratic correction \eqref{p12-83-c20-s6:eq:correccion-cuadratica-automorfismos-h3} does not preserve the product, if \eqref{p12-83-c20-s6:eq:identidad-grupo-schlafli-nueva} does not produce the coefficients \((10,1,5)\), if an abelian group produces the same configuration, if the partition does not sum to 729, if the table \ref{p12-83-c20-s6:tab:h3-27-rectas-u029} does not intertwine the adjacency matrices, or if \(W(E_6)\) is identified only by its order and not by its action on the twenty-seven classes and the 72 roots.
\end{criterion}


\section{Oriented APP–Fock coupling and transverse rigidity of \(12\) fibres}
\label{p12-83-c20-s7:chap:acoplamiento-app-fock}
\index{APP--Fock!oriented coupling}
\index{Fock!degree two}
\index{sheet orientation}

The ordinary cyclotomic determinant does not distinguish \(C\) from \(C^{-1}\).
Sum--product information is preserved by explicitly introducing
sheet parity and an orientation line. The resulting coupling
produces index \(54\) from the APP aggregate. Once the four
coupling conditions have been fixed, \(54\), \(J\), and \(196884\) are not inputs to
the evaluation. This evaluative independence asserts neither universal uniqueness
of the ansatz nor historical independence of its design.

\subsection{Triangular mark on each orbit}

Let
\[
M=\mathbb Z[\mathbb Z/9\mathbb Z]
\]
and let \(M_{\mathrm{bal}}^{C_3}\) be the submodule of invariant elements
\(\nu\) whose six unit coefficients coincide:
\[
\nu_u=c(\nu)
\qquad
\bigl(u\in(\mathbb Z/9\mathbb Z)^\times\bigr).
\]
The coefficients of \(e_0,e_3,e_6\) permanecen libres.

For a free orbit
\(\mathcal O=\{u,4u,7u\}\), set
\[
b_v=e_v-e_{4v}\in A(\mathcal O).
\]
If \(b_v^\flat\) is the covector defined by the form of \(A_2\), the three
roots of the triangular frame satisfy
\[
\boxed{
\sum_{v\in\mathcal O}b_v\otimes b_v^\flat
=3I_{A(\mathcal O)}.}
\]
Consequently,
\[
\sum_{v\in\mathcal O}\nu_v\,b_v\otimes b_v^\flat
=3c(\nu)I.
\]
The factor three proceeds from the resolution of the identity by the root framework.

\subsection{Oriented coupling \(\mathcal H_{\Sigma,\Pi}^{(m)}\) between the balance \(C_3\), leaf parity, and degree \(2\) of Fock space}

For \(m\in2\mathbb N_{>0}\), let
\[
\begin{aligned}
\mathcal V_m&=\mathcal V_{\Sigma,m}\oplus\mathcal V_{\Pi,m},\\
\mathcal V_{\Sigma,m}&=(A_2\otimes\mathbb C)^{\oplus m/2},\qquad
\mathcal V_{\Pi,m}=(A_2\otimes\mathbb C)^{\oplus m/2},\\
g_m&=C^{\oplus m/2}\oplus(C^{-1})^{\oplus m/2},\\
\mathsf E_m&=I_{\mathcal V_{\Sigma,m}}
\oplus(-I_{\mathcal V_{\Pi,m}}).
\end{aligned}
\]
The specialization \(m=12\), with the labels of pair, sheet and orbit,
is \(\mathcal V_{12}=W_{\mathrm{APP}}\otimes\mathbb C\). This notation
prevents confusing the representational space with the Witt system
\(W_{12}=S(5,6,12)\). Degree two of the Fock space and the induced action are
\[
\mathcal F_2(\mathcal V_m)
=\mathcal V_m[2]\oplus\operatorname{Sym}^2(\mathcal V_m[1]),
\]
\[
\Gamma_2(g_m)
=g_m|_{\mathcal V_m[2]}
\oplus\operatorname{Sym}^2(g_m|_{\mathcal V_m[1]}).
\]

Let \(o_{\Sigma\Pi}\) be the signed orientation line
\(\Pi-\Sigma\), and let \(M_{\mathrm{bal}}^{C_3,-}\) be the oriented copy in
which the exchange of leaves acts by \(\nu\mapsto-\nu\). We define
\[
\boxed{
\begin{aligned}
\mathcal H_{\Sigma,\Pi}^{(m)}&:
M_{\mathrm{bal}}^{C_3,-}\longrightarrow
\operatorname{End}(\mathcal F_2(\mathcal V_m))\otimes o_{\Sigma\Pi},\\
\mathcal H_{\Sigma,\Pi}^{(m)}(\nu)
&=3c(\nu)
\left(0_{\mathcal V_m[2]}\oplus\operatorname{Sym}^2\mathsf E_m\right)
\otimes o_{\Sigma\Pi}.
\end{aligned}}
\]
The transformation is linear in \(c\), is supported on
\(\operatorname{Sym}^2(\mathcal V_m[1])\), uses leaf parity, and adopts the
normalization of the triangular frame.

Let \(\sigma\) the exchange of sheets. Then
\[
\nu\longmapsto-\nu,\qquad
o_{\Sigma\Pi}\longmapsto-o_{\Sigma\Pi},\qquad
\mathsf E_m\longmapsto-\mathsf E_m,
\]
and
\[
\operatorname{Sym}^2(-\mathsf E_m)
=\operatorname{Sym}^2(\mathsf E_m).
\]
Therefore,
\[
\mathcal H_{\Sigma,\Pi}^{(m)}(\sigma\nu)
=\sigma\mathcal H_{\Sigma,\Pi}^{(m)}(\nu).
\]
When the orientation is reversed, the dual covector also changes sign; thus
\((-o_{\Sigma\Pi}^{\vee})(-o_{\Sigma\Pi})=1\), and the oriented evaluation
remains normalized.

Fixing
\(o_{\Sigma\Pi}^{\vee}(o_{\Sigma\Pi})=1\), the oriented trace is
\[
\operatorname{Tr}_{o,\mathcal F_2}
(T\otimes o_{\Sigma\Pi})
=\operatorname{Tr}_{\mathcal F_2}(T).
\]

\begin{theorem}[APP-oriented Index--Fock]
\label{p12-83-c20-s7:thm:indice-orientado-app-fock}
Suppose that \(m\) fibers \(A_2\), with \(m\) even, are distributed equally
between the orientations \(C\) and \(C^{-1}\). Then
\[
\boxed{
\operatorname{Tr}_{o,\mathcal F_2}
\left(
\mathcal H_{\Sigma,\Pi}^{(m)}(\nu)\Gamma_2(g_m)
\right)
=-\frac{3m\,c(\nu)}2.}
\]
For the aggregate produced by APP,
\[
\delta=\nu_\Pi-\nu_\Sigma
=12e_0+3e_3+3e_6
-3\sum_{u\in(\mathbb Z/9\mathbb Z)^\times}e_u,
\]
one has \(c(\delta)=-3\) and
\[
m=3\ \text{pairs}\cdot2\ \text{sheets}\cdot2\ \text{orbits}=12.
\]
Therefore,
\[
\boxed{
\operatorname{Tr}_{o,\mathcal F_2}
\left(
\mathcal H_{\Sigma,\Pi}^{(12)}(\delta)\Gamma_2(g_{12})
\right)=54.}
\]
\end{theorem}

\begin{proof}
Let \(A=\mathsf E_mg_m\). The contributions of the two leaves cancel:
\(\operatorname{tr}A=0\). Since \(\mathsf E_m\) commutes with \(g_m\) and
\(\mathsf E_m^2=I\),
\[
\operatorname{tr}(A^2)=\operatorname{tr}(g_m^2)=-m.
\]
The identity
\[
\operatorname{tr}(\operatorname{Sym}^2A)
=\frac{(\operatorname{tr}A)^2+\operatorname{tr}(A^2)}2
\]
gives \(-m/2\); the frame supplies \(3c(\nu)\).
\end{proof}

\begin{corollary}[Minimality of the Bosonic Degree Detector]
\label{p12-83-c20-s7:cor:minimalidad-grado-dos-app-fock-u030}
In the declared class of couplings, degree one does not detect the oriented
aggregate, and degree two is the first positive bosonic degree that
detects it:
\[
 \operatorname{tr}(\mathsf E_mg_m)=0,
 \qquad
 \operatorname{tr}\!\left(
 \operatorname{Sym}^2(\mathsf E_mg_m)
 \right)=-\frac m2\neq0.
\]
\end{corollary}

\begin{proof}
The first equality is the cancellation of the two sheets. The second was
obtained in the proof of \cref{p12-83-c20-s7:thm:indice-orientado-app-fock}. Since degree
zero contains only the unit and degree one has zero trace, degree two is
the first with a nonzero response.
\end{proof}

\subsection{Transverse rigidity of the number of fibers}

For an integer \(m\ge1\), we define the formal series
\begin{equation}
 T_m(q):=
 q^{-1}\prod_{n\ge1}(1+q^n+q^{2n})^{-m}
 \in q^{-1}\mathbb Z[[q]]
 \label{p12-83-c20-s7:eq:familia-fock-m-fibras-u030}
\end{equation}
and, in the orthogonal sum \(A_2^m\), the radial vector
\begin{equation}
 v_m:=4\rho_1+\rho_2+\cdots+\rho_m,
 \qquad
 \|\rho_j\|^2=2.
 \label{p12-83-c20-s7:eq:familia-radial-m-fibras-u030}
\end{equation}
We construct two entire functions without using \(m=12\) beforehand:
\[
 F(m):=[q]T_m(q),
 \qquad
 R(m):=\|v_m\|^2.
\]

\begin{proposition}[Fock coefficient and radial norm]
\label{p12-83-c20-s7:prop:dos-funciones-rigidez-doce-u030}
For every \(m\ge1\),
\begin{equation}
 F(m)=\frac{m(m-3)}2,
 \qquad
 R(m)=2(m+15).
 \label{p12-83-c20-s7:eq:funciones-rigidez-doce-u030}
\end{equation}
\end{proposition}

\begin{proof}
Up to second order,
\[
 (1+q+q^2)^{-m}
 =1-mq+\frac{m(m-1)}2q^2+O(q^3),
\]
\[
 (1+q^2+q^4)^{-m}=1-mq^2+O(q^4).
\]
The factors with \(n\ge3\) do not contribute to the coefficient of \(q^2\).
Since the exterior factor of \(T_m\) is \(q^{-1}\),
\[
 F(m)=\frac{m(m-1)}2-m=\frac{m(m-3)}2.
\]
By orthogonality of the \(m\) copies of \(A_2\),
\[
 R(m)=4^2\|\rho_1\|^2+
 \sum_{j=2}^{m}\|\rho_j\|^2
 =32+2(m-1)=2(m+15).
\]
\end{proof}

\begin{theorem}[Cyclotomic--radial rigidity]
\label{p12-83-c20-s7:thm:rigidez-transversal-doce-u030}
The equality \(F(m)=R(m)\) is equivalent to
\begin{equation}
 (m-12)(m+5)=0.
 \label{p12-83-c20-s7:eq:ecuacion-rigidez-doce-u030}
\end{equation}
In the domain \(m\in\mathbb Z_{>0}\), the unique solution is
\[
 \boxed{m=12,\qquad F(12)=R(12)=54.}
\]
\end{theorem}

\begin{proof}
Equating the two expressions of
\eqref{p12-83-c20-s7:eq:funciones-rigidez-doce-u030} produce
\[
 m(m-3)=4(m+15),
\]
that is,
\[
 m^2-7m-60=(m-12)(m+5)=0.
\]
The second root is negative and lies outside the domain.
\end{proof}

This rigidity belongs to the declared cyclotomic--radial family. It does not use the
value \(54\) to choose \(m\): it first constructs \(F(m)\) and \(R(m)\), and
then resolves their equality.

\subsection{Factorization of the APP–Fock coupling: hypotheses, invariants, necessary
conditions for validity, and obstructions}

The coupling factorizes by
\[
M_{\mathrm{bal}}^{C_3,-}/\ker c
\cong\mathbb Z\otimes o_{\Sigma\Pi},
\]
and the induced map is injective because
\(\operatorname{Sym}^2\mathsf E_m\neq0\).
The formula is selected by the four declared conditions:
linearity in \(c\), support in
\(\operatorname{Sym}^2(\mathcal V_m[1])\), sheet parity,
and normalization by the triangular frame. No uniqueness is asserted
among all natural couplings, nor historical independence of the
ansatz. Once those conditions have been fixed, \(54\), \(J\), and \(196884\) are not
inputs to the evaluation.

Three controls show that the result is not determined by a retroactive normalization:
\[
(m,c)=(12,-2)\longmapsto36,\qquad
(m,c)=(10,-3)\longmapsto45,
\]
and
\[
\delta_k=\delta+ke_0
\]
leaves \(\mathcal H_{\Sigma,\Pi}^{(12)}\) invariant while the digital weight
changes from \(54\) to \(54+9k\). In particular, for \(k=1\),
\[
\operatorname{Tr}_{o,\mathcal F_2}
\left(
\mathcal H_{\Sigma,\Pi}^{(12)}(\delta+e_0)\Gamma_2(g_{12})
\right)=54,
\qquad
w(\delta+e_0)=63.
\]

\begin{proposition}[Need for orientation data]
\label{p12-83-c20-s7:prop:orientacion-app-fock}
If \(o_{\Sigma\Pi}\) is forgotten, every natural scalar that is linear in
\(\nu_\Pi-\nu_\Sigma\) and depends only on the class of the representation
vanishes.
\end{proposition}

\begin{proof}
\(C\) and \(C^{-1}\) are conjugate, so the representational class is
even under sheet exchange. The aggregate
\(\nu_\Pi-\nu_\Sigma\) is odd. Without the orientation line, a
transformation that is simultaneously even and odd is zero.
\end{proof}

The theorem thus establishes the exact transport of the aggregate to Fock
degree two. The local state-by-state defect preserves the descent obstruction
of \cref{p12-83-c1-s5:cor:obstruccion-isotipica-app-autonoma}; the two statements belong to
distinct domains and are compatible.


% Unidad U031. Alineamiento APP--Witt, automorfia de orden tres,
% construccion VOA y Moonshine.
 \section{From spectral jump 4 to Monster symmetries: APP unfolding
\texorpdfstring{$2\to6\to54$}{2--6--54}, the Leech lattice,
\texorpdfstring{$V^\natural$}{V natural}, and the Moonshine corridor}
\label{sec:cierre-equivalencias-corredor-moonshine-20260825}

The corridor is not a sequence of numerical coincidences. It begins at the
central block of APP, preserves the sum--product orientation, branches into the
APP--Fock and APP--Witt couplings, and is coordinated anew through the
grading. The binding causal order is
\[\begin{aligned}
 B_c=9P_++5P_-
 &\longrightarrow 4\longrightarrow2\longrightarrow6\longrightarrow54\\
 &\longrightarrow
 \left\{
 \begin{array}{l}
  \text{Fock degree two and }[q]t_3=54,\\
  \text{APP--Witt alignment and }N_\alpha\cong\Lambda_{24}
 \end{array}
 \right.\\
 &\longrightarrow V^\natural
 \longrightarrow
 \bigl(\operatorname{Aut},\operatorname{ch},g\mapsto T_g\bigr)
 \longrightarrow\text{Moonshine}.
\end{aligned}\]
The two intermediate branches have distinct codomains. The common key \(54\)
is a graded invariant transported by the corridor; it is not a
linear map from a scalar into a vertex operator algebra.

\subsection{Typed hierarchy and \(2\) realizations of index \(54\)}

\begin{theorem}[Spectral Hierarchy and APP Unfolding]
\label{thm:cierre-vtm-jerarquia-app}
Let \(B_c=7I+2R=9P_++5P_-\) be the central block of APP.\@ If
\(M_\Pi,M_\Sigma\) are the defined modular compressions, then
\[
 9-5=4,
 \qquad (B_c)_{i\ne j}=2,
 \qquad \sum M_\Pi-\sum M_\Sigma=51-45=6,
\]
and the deployment over the nine positions gives
\[
 9\cdot6=459-405=54.
\]
The four integers pertain, respectively, to the spectral gap, the
local coupling, modular compression, and two-dimensional integration.
\end{theorem}

\begin{proof}
The spectral decomposition of \(B_c\) gives \(9\) and \(5\), whereas
\(7I+2R\) fixes the nondiagonal coefficient. The two aggregate sums are \(51\)
and \(45\). Each compressed entry represents nine positions, so that the
unfolding multiplies the difference by nine. None of the four integers
is obtained by retrospectively identifying its types.
\end{proof}

Once the orientation line \(o_{\Sigma\Pi}\), leaf parity, and the
triangular normalization have been fixed, the oriented coupling
\[
 \mathcal H_{\Sigma,\Pi}^{(m)}:
 M_{\mathrm{bal}}^{C_3,-}
 \longrightarrow
 \operatorname{End}\!\bigl(\mathcal F_2(\mathcal V_m)\bigr)
 \otimes o_{\Sigma\Pi}
\]
satisfies
\[
 \operatorname{Tr}_{o,\mathcal F_2}
 \left(\mathcal H_{\Sigma,\Pi}^{(m)}(\nu)\Gamma_2(g_m)\right)
 =-\frac{3m\,c(\nu)}2.
\]
For \(m=12\) and \(c(\delta)=-3\), the trace equals \(54\). On the other hand,
\[
 t_3(q)=q^{-1}\prod_{n\ge1}(1+q^n+q^{2n})^{-12}
       =q^{-1}-12+54q-76q^2-243q^3+\cdots,
\]
from which \([q]t_3=54\).

\begin{theorem}[Common Index Compatibility Cone]
\label{thm:cierre-vtm-indice-comun-54}
With the starting structures of each branch preserved, it holds that
\[
 \boxed{
 D_{\mathrm{APP}}
 =\operatorname{Tr}_{o,\mathcal F_2}
   \left(\mathcal H_{\Sigma,\Pi}^{(12)}(\delta)\Gamma_2(g_{12})\right)
 =[q]t_3
 =v_\alpha^2
 =\frac{196884}{5\cdot729+1}
 =54.}
\]
This equality constitutes a cone of evaluations toward the same integer. It does not
identify \(W_{\mathrm{APP}}\), the Fock space, the neighboring lattice, and
\(V^\natural\).
\end{theorem}

\begin{proof}
The APP census yields \(459-405=54\); the oriented trace yields
\(54\); expansion of the cyclotomic product yields \([q]t_3=54\); and
the radial vector of the neighbor satisfies \(v_\alpha^2=4^2\cdot2+11\cdot2=54\). Finally,
\(5\cdot729+1=3646\) and \(54\cdot3646=196884\). The equalities are compared only after each
object has been constructed and therefore do not numerically select the
corridor.
\end{proof}

\paragraph{Exact status of the hierarchy \(4\to2\to6\to54\) and of the lattice
identification.}
The hierarchy \(4\to2\to6\to54\), the APP--Fock trace, and the extraction
\([q]t_3\) are exact internal identities. Identification of the neighbor with
Leech uses the declared chamber, neighboring vector, and lattice
classification.
%
 
\chapter{Vertex operator algebra and monstrous Moonshine}
\section{Alineamiento ciclotomico, orbifold of order \(3\) and Monstrous Moonshine}
\label{p12-83-c20-s8:chap:alineamiento-orbifold-moonshine}

The Leech lattice constructed in
\cref{p12-83-c20-s5:thm:identificacion-vecino-leech} makes it possible to pass from finite incidence
to an infinite graded structure. This transition will not be effected
through a chain of names. The following will be constructed, in typed order:
\begin{enumerate}
 \item the intertwining of the twelve fibres \(A_2\);
 \item a fixed-point-free lattice isometry of order three;
 \item its graded trace and its type-zero lift;
 \item the cyclic orbifold and its dual automorphism of class \(3B\);
 \item the moonshine vertex operator algebra;
 \item through distinct publications, its automorphism group, its graded
       character, and the McKay--Thompson series.
\end{enumerate}
Moonshine sera the theorem that coordina the action and the series. is not usara
the atajo not tipado
\(V^\natural\to\mathbb M\to\text{Moonshine}\).

\subsection{Alignment of the \(12\) APP and Witt fibers}

The twelve fibers of \(W_{\mathrm{APP}}\cong A_2^{12}\) carry three
independent labels:
\[
 i\in\mathbb Z/3\mathbb Z
 \quad\text{(coordinate pair)},\qquad
 \mu\in\{0,1\}
 \quad\text{(sheet)},\qquad
 \varepsilon\in\{0,1\}
 \quad\text{(cyclotomic orbit)}.
\]
Fixing the dodecaphasic origin constructed on the TPK, we define
\begin{equation}
 \lambda(i,\mu,\varepsilon)
 :=4i+2\mu+\varepsilon\pmod{12}.
 \label{p12-83-c20-s8:eq:alineamiento-doce-fibras-u031}
\end{equation}

\begin{lemma}[Bijectivity of the alignment]
\label{p12-83-c20-s8:lem:biyectividad-alineamiento-u031}
The map \(\lambda\) from
\eqref{p12-83-c20-s8:eq:alineamiento-doce-fibras-u031} is a bijection onto
\(\mathbb Z/12\mathbb Z\).
\end{lemma}

\begin{proof}
For each \(i\), the four values \(4i+2\mu+\varepsilon\) form the block
\(\{4i,4i+1,4i+2,4i+3\}\). The three blocks for \(i=0,1,2\) are disjoint
and cover the twelve classes.
\end{proof}

Let \(C\) be the Coxeter element of \(A_2\) and let \(\mathsf R\) be the reflection that
satisfies
\begin{equation}
 \mathsf R C\mathsf R=C^{-1}.
 \label{p12-83-c20-s8:eq:reflexion-conjuga-coxeter-u031}
\end{equation}
Once the trivializations \(\iota_b\) and the oriented frames \(P_b\) have been fixed,
we set
\begin{align}
 g_b&:=\iota_b^{-1}P_bCP_b^{-1}\iota_b,
 \label{p12-83-c20-s8:eq:coxeter-calibrado-fibra-u031}\\
 T_{i\mu\varepsilon}
 &:=\iota_{\lambda(i,\mu,\varepsilon)}^{-1}
 P_{\lambda(i,\mu,\varepsilon)}\mathsf R^\mu.
 \label{p12-83-c20-s8:eq:entrelazador-fibra-u031}
\end{align}

\begin{theorem}[APP Entanglement--Witt]
\label{p12-83-c20-s8:thm:alineamiento-app-witt-u031}
For each terna \((i,\mu,\varepsilon)\),
\begin{equation}
 g_{\lambda(i,\mu,\varepsilon)}T_{i\mu\varepsilon}
 =T_{i\mu\varepsilon}C^{(-1)^\mu}.
 \label{p12-83-c20-s8:eq:identidad-entrelazamiento-app-witt-u031}
\end{equation}
The direct sum defines a \(C_3\)-equivariant isomorphism
\begin{equation}
 T_\lambda:
 W_{\mathrm{APP}}\otimes\mathbb C
 \xrightarrow{\;\cong\;}
 W_c:=\bigoplus_{b=0}^{11}
 (A_{2,b}\otimes\mathbb C).
 \label{p12-83-c20-s8:eq:isomorfismo-app-witt-u031}
\end{equation}
\end{theorem}

\begin{proof}
For \(\mu=0\), the definition of \(g_b\) directly gives conjugation
by \(C\). For \(\mu=1\), one inserts
\(\mathsf RC\mathsf R=C^{-1}\). The
\cref{p12-83-c20-s8:lem:biyectividad-alineamiento-u031} guarantees that the direct sum
uses every Witt position exactly once and preserves the pair, sheet, and orbit.
\end{proof}

The theorem is relative to the simultaneous transport of the order of pairs, leaves, and
orbits, of the Witt origin, of the frames, and of the orientation. Changing a
single one of these coordinates does not define the same functor.

\subsection{Graded trace of the ternary generator}

Let
\[
 \rho_W:=\bigoplus_{b=0}^{11}g_b
\]
and consider the bosonic envelope
\begin{equation}
 \mathcal H_{\rm bos}(W_c)
 :=
 \operatorname{Sym}\!\left(
 \bigoplus_{n\ge1}W_c[n]
 \right).
 \label{p12-83-c20-s8:eq:fock-bosonico-doce-fibras-u031}
\end{equation}
If \(N\) is the degree operator, we define
\begin{equation}
 Z_{\rm Fock}(q)
 :=
 q^{-1}\operatorname{Tr}_{\mathcal H_{\rm bos}(W_c)}
 \left(\Gamma_{\rm bos}(\rho_W)q^N\right).
 \label{p12-83-c20-s8:eq:traza-graduada-fock-u031}
\end{equation}

\begin{proposition}[Cyclotomic polynomial of level three]
\label{p12-83-c20-s8:prop:producto-tres-doce-fibras-u031}
As a formal Laurent series,
\begin{equation}
 \boxed{
 Z_{\rm Fock}(q)
 =q^{-1}\prod_{n\ge1}(1+q^n+q^{2n})^{-12}
 =:t_3(q).}
 \label{p12-83-c20-s8:eq:t3-producto-u031}
\end{equation}
Its expansion begins with
\begin{equation}
 t_3(q)
 =q^{-1}-12+54q-76q^2-243q^3+1188q^4+O(q^5).
 \label{p12-83-c20-s8:eq:t3-expansion-u031}
\end{equation}
\end{proposition}

\begin{proof}
Each \(g_b\) is conjugate to \(C\), whence
\[
 \det(I-g_bq^n)=1+q^n+q^{2n}.
\]
The identity
\[
 \sum_{k\ge0}
 \operatorname{Tr}(\operatorname{Sym}^k\rho_W)t^k
 =\det(I-\rho_Wt)^{-1}
\]
applied in each degree produces the product. Formal multiplication through
degree five gives the expansion. In particular, the coefficient of \(q\) is
obtained as
\[
 [q^2]\,
 (1+q+q^2)^{-12}(1+q^2+q^4)^{-12}
 =66-12=54.
\]
\end{proof}

The value \(54\) appears here as the coefficient of a trace constructed from
twelve fibers. Unit U030 also obtained it as the oriented index of the
APP aggregate and proved that twelve is the unique positive solution of
cyclotomic--radial rigidity. The equality of the two values is a compatibility
between two realizations; neither uses the coefficient of \(J\) as input.

\subsection{Full-support word and fixed-point-free isometry}

Let
\[
 q_*=(1,1,1,1,1,1)\in\mathbb F_3^6.
\]
Multiplication by the Paley--Witt matrix gives
\begin{equation}
 q_*A_W=(2,1,1,1,1,1).
 \label{p12-83-c20-s8:eq:producto-qstar-aw-u031}
\end{equation}
Therefore,
\begin{equation}
 c_*:=(q_*,q_*A_W)
 =(1,1,1,1,1,1,2,1,1,1,1,1)
 \in C_W
 \label{p12-83-c20-s8:eq:palabra-soporte-total-u031}
\end{equation}
and all its coordinates are non-zero.

For each position \(b\), we fix
\begin{equation}
 P_b(c_*)=
 \begin{cases}
  \mathsf R,&(c_*)_b=1,\\
  I,&(c_*)_b=2,
 \end{cases}
 \qquad
 g_c:=\bigoplus_{b=0}^{11}P_b(c_*)CP_b(c_*)^{-1}.
 \label{p12-83-c20-s8:eq:isometria-orden-tres-u031}
\end{equation}

\begin{proposition}[Order, absence of fixed points and gluing]
\label{p12-83-c20-s8:prop:isometria-fijo-libre-pegado-u031}
The operator \(g_c\) has order three, has no nonzero fixed vectors, and
preserves the Niemeier lattice \(N(A_2^{12})\) constructed in U021.
\end{proposition}

\begin{proof}
Each block is conjugate to the Coxeter \(C\), whose minimal polynomial is
\(X^2+X+1\). Hence it has order three and has no eigenvalue one. The direct
sum preserves both properties. The Coxeter acts trivially on
\(A_2^*/A_2\); therefore each block preserves the discriminant class and the
union of lateral classes that defines the gluing.
\end{proof}

\subsection{Preservation of the Leech Neighbor}

Write
\[
 v:=v_{\mathcal F_*}
 =4\rho_1+\rho_2+\cdots+\rho_{12},
 \qquad
 v^2=54,
\]
for the vector selected by the incidence origin of
\eqref{p12-83-c20-s5:eq:vector-marco-incidencial}. Let
\[
 N:=N(A_2^{12}),
 \qquad
 N_0:=\ker(\langle\,\cdot\,,v\rangle\bmod3),
 \qquad
 N_v:=N_0+\mathbb Z\frac v3\cong\Lambda_{24}.
\]

\begin{theorem}[Neighbor-Oriented Preservation]
\label{p12-83-c20-s8:thm:preservacion-vecino-orden-tres-u031}
The displacements
\begin{equation}
 y_*:=\frac{g_cv-v}{3},
 \qquad
 z_*:=\frac{g_c^{-1}v-v}{3}
 \label{p12-83-c20-s8:eq:desplazamientos-vecino-u031}
\end{equation}
belong to \(N_0\). Consequently,
\[
 g_cN_0=N_0,
 \qquad
 g_cN_v=N_v.
\]
\end{theorem}

\begin{proof}
The discriminant words of \(y_*\) and \(z_*\) are
\(c_*\) and \(-c_*\), both in \(C_W\); hence
\(y_*,z_*\in N\). Computation in each factor \(A_2\) gives
\begin{equation}
 \langle y_*,v\rangle
 =\langle z_*,v\rangle
 =-(4^2+11)=-27\equiv0\pmod3.
 \label{p12-83-c20-s8:eq:emparejamientos-desplazamientos-u031}
\end{equation}
Thus both lie in \(N_0\). For \(x\in N_0\),
\[
 \langle g_cx,v\rangle
 =\langle x,g_c^{-1}v\rangle
 =\langle x,v\rangle+3\langle x,z_*\rangle
 \equiv0\pmod3.
\]
The same count with \(g_c^{-1}\) gives equality of the kernel. Finally,
\[
 g_c(v/3)=v/3+y_*,
\]
so that the class generated by the neighbor is preserved.
\end{proof}

\subsection{\(24\) orientations and lattice naturality}

The weight enumerator of \eqref{p12-83-c20-s1:eq:enumerador-pesos-cw} contains
\(24y^{12}\). Therefore, the set
\[
 C_W^\times:=
 \{u\in C_W:u_b\in\{1,2\}\text{ for every }b\}
\]
has 24 elements. For each \(u\in C_W^\times\), we define
\[
 P_b(u)=
 \begin{cases}
  \mathsf R,&u_b=1,\\
  I,&u_b=2,
 \end{cases}
 \qquad
 g_u:=\bigoplus_{b=0}^{11}P_b(u)CP_b(u)^{-1}.
\]

\begin{proposition}[Robustness of Full Support Orientations]
\label{p12-83-c20-s8:prop:robustez-orientaciones-u031}
For every \(u\in C_W^\times\),
\[
 \frac{g_uv-v}{3},
 \quad
 \frac{g_u^{-1}v-v}{3}
 \in N_0.
\]
In particular, each \(g_u\) preserves the neighbor constructed with the
transported orientation.
\end{proposition}

\begin{proof}
In simple root coordinates,
\[
 C\rho-\rho=(-2,-1),
 \qquad
 C^{-1}\rho-\rho=(-1,-2).
\]
The discriminant classes of the two displacements are \(u\) and \(-u\).
For the radial coefficient \(a_b\in\{4,1\}\),
\[
 \left\langle
 \frac{(C^{\pm1}-I)a_b\rho}{3},a_b\rho
 \right\rangle=-a_b^2.
\]
The sum over the twelve factors again gives \(-27\). The proof of
\cref{p12-83-c20-s8:thm:preservacion-vecino-orden-tres-u031} is applied component by
component.
\end{proof}

To type the recalibration, let
\(h\in\operatorname{Aut}(C_W)\) be monomial: a permutation \(\sigma\) of
the twelve coordinates followed by multipliers
\(\epsilon_b\in\{1,2\}\). We define its lattice lift
\begin{equation}
 \widetilde h
 :=P_\sigma\circ
 \bigoplus_{b=0}^{11}\mathsf R^{\nu_b},
 \qquad
 \nu_b=
 \begin{cases}
 0,&\epsilon_b=1,\\
 1,&\epsilon_b=2.
 \end{cases}
 \label{p12-83-c20-s8:eq:elevacion-monomial-u031}
\end{equation}
Its action on the discriminant is exactly \(h\). If the change simultaneously
transports the code, flag, incidence origin, vector \(v\), neighbor, and
frames, then
\begin{equation}
 g_{hu}=\widetilde h\,g_u\,\widetilde h^{-1}.
 \label{p12-83-c20-s8:eq:naturalidad-isometrias-u031}
\end{equation}
The ternary matrix \(h\) does not act without a lift on \(A_2^{12}\) or on
the Leech lattice.

\subsection{Lattice trace and computation of type \(0\)}

Let \(\widehat g_c\) be the standard lift of \(g_c\) to the lattice VOA
\[
 V_{N_v}=M(1)\otimes\mathbb C_\varepsilon[N_v].
\]
For \(j=1,2\), the oscillator--lattice decomposition produces
\begin{equation}
 \begin{split}
 Z_{0,j}(\tau)
 &:=
 \operatorname{Tr}_{V_{N_v}}
 \left(\widehat g_c^{\,j}q^{L_0-1}\right)\\
 &=q^{-1}
 \left(
 \sum_{\lambda\in N_v^{g_c^j}}
 \varepsilon_j(\lambda)
 q^{\langle\lambda,\lambda\rangle/2}
 \right)
 \prod_{n\ge1}\det(I-g_c^jq^n)^{-1}.
 \end{split}
 \label{p12-83-c20-s8:eq:factorizacion-traza-reticular-u031}
\end{equation}

\begin{proposition}[Trace of the nontrivial sectors]
\label{p12-83-c20-s8:prop:traza-reticular-t3-u031}
We have
\[
 Z_{0,1}(\tau)=Z_{0,2}(\tau)=t_3(\tau).
\]
\end{proposition}

\begin{proof}
The absence of fixed points implies
\[
 N_v^{g_c}=N_v^{g_c^2}=\{0\},
\]
so that the lattice sum in
\eqref{p12-83-c20-s8:eq:factorizacion-traza-reticular-u031} equals one. In each factor,
\(g_c^j\) has eigenvalues \(\zeta_3,\zeta_3^2\); therefore,
\[
 \det(I-g_c^jq^n)=1+q^n+q^{2n}.
\]
The twelve factors reproduce \eqref{p12-83-c20-s8:eq:t3-producto-u031}.
\end{proof}

\begin{theorem}[Twisted weight and type of the lift]
\label{p12-83-c20-s8:thm:peso-torcido-tipo-cero-u031}
The standard lift has order three, twisted weight
\begin{equation}
 \rho_{g_c}=\frac43
 \label{p12-83-c20-s8:eq:peso-torcido-u031}
\end{equation}
and type zero.
\end{theorem}

\begin{proof}
In
\(\mathfrak h=N_v\otimes_{\mathbb Z}\mathbb C\), each of the twelve
factors contributes one eigenvector with eigenvalue \(\zeta_3\) and another with eigenvalue
\(\zeta_3^2\). Therefore,
\[
 \dim\mathfrak h_{\zeta_3}
 =\dim\mathfrak h_{\zeta_3^2}=12.
\]
The formula for the weight of the twisted vacuum of order three gives
\[
 \rho_{g_c}
 =\frac1{4\cdot3^2}
 \bigl(1\cdot2\cdot12+2\cdot1\cdot12\bigr)
 =\frac43.
\]
For a lattice isometry of odd order, the standard elevation preserves
the order. The type is the class of
\[
 3^2\rho_{g_c}=12\equiv0\pmod3.
\]
Thus, the two initially independent hypotheses are calculated here.
\end{proof}

\subsection{Cyclic orbifold and dual class \texorpdfstring{\(3B\)}{3B}}

\begin{theorem}[Triadic orbifold of the Leech neighbor]
\label{p12-83-c20-s8:thm:orbifold-triadico-vnatural-u031}
The classical cyclic-orbifold theorems, applied to the standard lift
in \cref{p12-83-c20-s8:thm:peso-torcido-tipo-cero-u031}, yield
\begin{equation}
 \operatorname{Orb}_{\mathbb Z/3}
 (V_{N_v},\widehat g_c)
 \cong V^\natural.
 \label{p12-83-c20-s8:eq:orbifold-triadico-vnatural-u031}
\end{equation}
The dual automorphism belongs to the class \(3B\) and its series is
\begin{equation}
 T_{3B}(\tau)=t_3(\tau)+12.
 \label{p12-83-c20-s8:eq:serie-3b-u031}
\end{equation}
\end{theorem}

\begin{proof}
The isometry is fixed-point-free of order three; the twisted weight and type zero were calculated, not assumed. The cyclic-orbifold theorems identify the holomorphic extension with the Moonshine module and classify the dual automorphism as \(3B\) \cite{ChenLamShimakura2016,AbeLamYamada2017,Carnahan2017}. The identification uses these theorems after constructing the lattice object; it is not deduced solely from an equality of series.
\end{proof}

If \(V^{(i)}\) is the \(\widehat g_c^{\,i}\)-twisted module and
\[
 Z_{i,j}(\tau):=
 \operatorname{Tr}_{V^{(i)}}\!
 \left(
 \phi_i(\widehat g_c)^j q^{L_0-1}
 \right),
\]
the character of the orbifold is the average of the nine sectors:
\begin{equation}
 \operatorname{ch}_{V^\natural}(\tau)
 =\frac13\sum_{i,j\in\mathbb Z/3}Z_{i,j}(\tau)
 =J(\tau).
 \label{p12-83-c20-s8:eq:promedio-nueve-sectores-u031}
\end{equation}
The level three rational identity is
\begin{equation}
 J
 =t_3+12+\frac{10\,3^9}{t_3}
 +\frac{4\,3^{14}}{t_3^2}
 +\frac{3^{18}}{t_3^3}.
 \label{p12-83-c20-s8:eq:j-desde-t3-u031}
\end{equation}
Since \(t_3^{-1}=q+O(q^2)\),
\begin{equation}
 [q]J=54+10\cdot3^9=196884
 =54(5\cdot729+1).
 \label{p12-83-c20-s8:eq:coeficiente-196884-u031}
\end{equation}
This is a graded compatibility within the constructed corridor, not a
scalar morphism \(54\to V^\natural\).

\subsection{Lattice construction of order \(2\)}

The preceding route of order three and the classical construction of
Frenkel--Lepowsky--Meurman produce two presentations of the same moonshine VOA.
To declare the second, we fix the central extension
\begin{equation}
 1\longrightarrow\{\pm1\}
 \longrightarrow\widehat\Lambda_{24}
 \longrightarrow\Lambda_{24}
 \longrightarrow1
 \label{p12-83-c20-s8:eq:extension-central-leech-u031}
\end{equation}
with commutator
\((-1)^{\langle\lambda,\mu\rangle}\), and a cocycle
\(\varepsilon(\lambda,\mu)\) that represents it. The twisted group algebra is \(\mathbb C_\varepsilon[\Lambda_{24}]\), and
\begin{equation}
 V_\Lambda
 :=M(1)\otimes\mathbb C_\varepsilon[\Lambda_{24}]
 \label{p12-83-c20-s8:eq:voa-reticular-leech-u031}
\end{equation}
is the lattice VOA of central charge 24.

The negation \(\lambda\mapsto-\lambda\) has a lift
\(\theta\) compatible with \eqref{p12-83-c20-s8:eq:extension-central-leech-u031}. Its trace
is
\begin{equation}
 t_2(\tau)
 :=
 \operatorname{Tr}_{V_\Lambda}
 (\theta q^{L_0-1})
 =q^{-1}\prod_{n\ge1}(1+q^n)^{-24}.
 \label{p12-83-c20-s8:eq:t2-leech-u031}
\end{equation}
Indeed, negation pairs \(e^\lambda\) with \(e^{-\lambda}\), and the
absence of torsion leaves only \(\lambda=0\) in the lattice trace; each
of the 24 oscillators contributes \((1+q^n)^{-1}\).

Let \(V_\Lambda^T\) the modulo irreducible torcido by \(\theta\), and
\((V_\Lambda^T)^+\) its sector of signo positive. The construction FLM is
\begin{equation}
 V^\natural
 =(V_\Lambda)^+\oplus(V_\Lambda^T)^+.
 \label{p12-83-c20-s8:eq:vnatural-flm-u031}
\end{equation}
The route \(\mathbb Z/2\) of
\eqref{p12-83-c20-s8:eq:vnatural-flm-u031} and the route \(\mathbb Z/3\) of
\eqref{p12-83-c20-s8:eq:orbifold-triadico-vnatural-u031} are parallel constructions
identified by classical theorems; neither replaces the other.

\subsection{Automorphism group, character, and graded series}

eleven constructed \(V^\natural\), appear two publications distinct:
\begin{align}
 \operatorname{Aut}(V^\natural)&\cong\mathbb M,
 \label{p12-83-c20-s8:eq:aut-vnatural-monstruo-u031}\\
 \operatorname{ch}_{V^\natural}(\tau)
 &=J(\tau)=j(\tau)-744
 =q^{-1}+196884q+21493760q^2+\cdots.
 \label{p12-83-c20-s8:eq:caracter-vnatural-j-u031}
\end{align}
The former identifies the symmetry group; the latter counts graded
dimensions. The vacuum line is
\[
 (V^\natural)_0=\mathbb C\mathbf1,
\]
and the absence of roots in \(\Lambda_{24}\) gives
\[
 (V^\natural)_1=0.
\]
In degree two,
\[
 196884=1+196883=196560+324=54(5\cdot729+1).
\]
Only the first equality is a decomposition into dimensions of
representations of the Monster; the others are arithmetic controls.

For \(g\in\mathbb M\), the McKay--Thompson series is defined
\begin{equation}
 T_g(\tau):=
 \operatorname{Tr}_{V^\natural}
 \left(gq^{L_0-1}\right).
 \label{p12-83-c20-s8:eq:serie-mckay-thompson-u031}
\end{equation}

\begin{theorem}[Monstrous Moonshine]
\label{p12-83-c20-s8:thm:moonshine-monstruoso-u031}
The graded action of \(\mathbb M\) on \(V^\natural\) produces the family
\(\{T_g\}_{g\in\mathbb M}\). The Moonshine theorems identify these
series with genus-zero modular functions and prove their replicability
identities
\cite{ConwayNorton1979,Borcherds1992,FLM1988}.
\end{theorem}

This theorem has as input the pair
\[
 \bigl(V^\natural,\operatorname{Aut}(V^\natural)\bigr)
\]
and its grading. It is not the consequence of an arrow whose domain is the
abstract group \(\mathbb M\).

For \(J=T_{1A}\), let \(P_n\) be the Faber polynomials normalized by
\[
 P_n(J)=q^{-n}+O(q).
\]
The Hecke--Faber identity is
\begin{equation}
 P_n(J)=n\,T_nJ.
 \label{p12-83-c20-s8:eq:hecke-faber-u031}
\end{equation}
It controls all degrees; the coefficient 196884 is only the first
nontrivial manifestation.

\subsection{Paley–Witt–Golay–Leech–Moonshine Runner Dependency Graph}

\begin{figure}[htbp]
\centering
\begin{tikzcd}[column sep=huge,row sep=large]
 &C_W
   \arrow[dl,"W_{12}"']
   \arrow[d,"N(A_2^{12})"',font=\scriptsize]
   \arrow[dr,phantom,"\substack{\text{distinct}\\\text{branches}}",
          description,pos=.72,font=\scriptsize]&\\
 M_{12}
 &\Lambda_{24}
   \arrow[d,"V_\Lambda"]
 &(\mathcal G_{24},D,\ell)
   \arrow[d,"W_{24}"]\\
 &V^\natural
   \arrow[dl,"\operatorname{Aut}"']
   \arrow[d,"\operatorname{ch}"]
   \arrow[dr,"g\mapsto T_g"]
 &M_{24}\\
 \mathbb M
 &J
 &\{T_g\}_{g\in\mathbb M}
\end{tikzcd}
\caption{Exceptional bifurcation and three publications of
\(V^\natural\). The binary input \((\mathcal G_{24},D,\ell)\) is
independent; there is no arrow \(W_{24}\to\Lambda_{24}\).}
\label{p12-83-c20-s8:fig:grafo-excepcional-moonshine-u031}
\end{figure}

\begin{criterion}[Refutation criteria for the moonshine corridor]
\label{p12-83-c20-s8:crit:refutacion-corredor-lunar-u031}
The construction is refuted if any of the following facts occurs:
\begin{enumerate}
 \item \(\lambda\) is not bijective or
       \eqref{p12-83-c20-s8:eq:identidad-entrelazamiento-app-witt-u031} falla;
 \item \(c_*\notin C_W\), \(g_c\) has fixed points or does not preserve the
       neighbor;
 \item the displacements of
       \eqref{p12-83-c20-s8:eq:desplazamientos-vecino-u031} do not belong to \(N_0\);
 \item the lattice lift \(\widetilde h\) does not realize the monomial
       automorphism \(h\) on the discriminant;
 \item the twisted weight is not \(4/3\) or the lift is not of type zero;
 \item the orbifold is identified with \(V^\natural\) without verifying the
       hypotheses of the classical theorem;
 \item \(W_{24}\) is used as a constructor of the Leech lattice;
 \item an equality of dimensions is presented as a construction of
       \(\mathbb M\);
 \item the publications
       \(\operatorname{Aut}\), \(\operatorname{ch}\), and \(g\mapsto T_g\) are improperly serialized.
\end{enumerate}
\end{criterion}
 \section{Theta control of the Leech lattice: identity \texorpdfstring{\(\Theta_{\Lambda}=E_4^3-720\Delta\)}{Theta Lambda = E4^3 - 720 Delta} and separation between arithmetic coincidence and structural morphism}
\label{sec:control-theta-leech}

From \(\Lambda_{24}\), the string
\[
\Lambda_{24}
\longrightarrow V^\natural
\longrightarrow\mathbb M
\longrightarrow\text{Moonshine}
\]
is classical mathematics: vertex operator algebra, the Monster group, and
McKay--Thompson functions. The contribution verified here lies in the finite
input, the selection of the boundary, and the rootless \(3\)-neighbor, not in claiming
those theorems anew.

Three numerical identities accompany this chain:
\[
\Theta_{\Lambda}=E_4^3-720\Delta,
\]
\[
196884=196560+324,
\]
\[
196884=54(3645+1)=54(5\cdot729+1).
\]
They are exact arithmetic identities, but without a morphism or a character
decomposition, they do not by themselves constitute mathematical signals of
compatibility. None replaces the construction of \(V^\natural\) or
provides a direct action of the Monster on HMT decimal words.
 \section{Uniformization of order \texorpdfstring{\(2\)}{2} in Moonshine: functions \texorpdfstring{\(T_{2A}\)}{T2A} and \texorpdfstring{\(j\)}{j}, modular structure and control of the micro--macro corridor}
\label{sec:voa-moonshine-orden-dos-rev11}
\index{VOA!Leech lattice}
\index{Moonshine!class 2A}

The incidence branch first constructs the Leech lattice
\(\Lambda_{24}\). From that datum, the classical theory of vertex operator
algebras provides a canonical prolongation. If
\(\mathfrak h=\Lambda_{24}\otimes_{\mathbb Z}\mathbb C\) and \(M(1)\) is
the Fock module of its Heisenberg algebra, the lattice VOA is
\begin{equation}
 V_{\Lambda}=M(1)\otimes\mathbb C[\Lambda_{24}],
 \label{eq:voa-reticular-leech-rev11}
\end{equation}
with central charge \(24\). For \(\alpha\in\Lambda_{24}\), its vertex
operators are written
\[
 Y(e^\alpha,z)=E^-(-\alpha,z)E^+(-\alpha,z)e^\alpha z^{\alpha_0}.
\]
The absence of roots in \(\Lambda_{24}\) implies the vanishing of the
weight-one sector in the Moonshine orbifold.

The involution
\[
 \theta:\Lambda_{24}\longrightarrow\Lambda_{24},
 \qquad \theta(\lambda)=-\lambda,
\]
is intrinsic and prolongs to \(V_{\Lambda}\). Its graded trace is obtained
by direct counting of the oscillator modes.

\begin{theorem}[Trace of the negation involution]
\label{thm:traza-theta-leech-rev11}
\begin{equation}
 t_2(\tau)
 :=\operatorname{Tr}_{V_\Lambda}
 \bigl(\theta q^{L_0-1}\bigr)
 =q^{-1}\prod_{n\geq1}(1+q^n)^{-24}.
 \label{eq:t2-producto-leech-rev11}
\end{equation}
\end{theorem}
\begin{proof}
The involution pairs \(e^\lambda\) with \(e^{-\lambda}\); since the lattice
is torsion-free, only \(\lambda=0\) contributes to the trace of the
lattice sector. Each one of the twenty-four energy oscillators \(n\) contributes
\(\sum_{k\geq0}(-1)^kq^{nk}=(1+q^n)^{-1}\). The vacuum displacement
produces the factor \(q^{-1}\).
\end{proof}

The Fricke completion associated with this trace is
\begin{equation}
 T_{2A}(\tau)=t_2(\tau)+24+\frac{2^{12}}{t_2(\tau)}
 =q^{-1}+4372q+96256q^2+1240002q^3+\cdots.
 \label{eq:T2A-hmt-rev11}
\end{equation}
The term \(24\) cancels the constant term of \(t_2\); the factor
\(2^{12}\) is the degeneracy of the twisted sector of twenty-four bosons.
The same boundary series determines the modular invariant through
\begin{equation}
 \boxed{j(\tau)=\frac{(t_2(\tau)+256)^3}{t_2(\tau)^2}}.
 \label{eq:j-desde-t2-rev11}
\end{equation}
Equations \eqref{eq:t2-producto-leech-rev11}--
\eqref{eq:j-desde-t2-rev11} construct the canonical branch of order two without
continuous parameters.

The Frenkel--Lepowsky--Meurman orbifold is defined by
\[
 V^\natural=(V_\Lambda)^+\oplus(V_\Lambda^T)^+.
\]
It is a holomorphic VOA of central charge \(24\), satisfies
\((V^\natural)_1=0\), and has character
\[
 \operatorname{ch}_{V^\natural}(\tau)=j(\tau)-744
\]
and its automorphism group is the Monster group
\cite{FLM1988,ConwayNorton1979,Borcherds1992}. In this chain, HMT fixes the
Leech lattice and the involution \(\theta\); the classical VOA and
FLM theorems identify the resulting prolongation with \(V^\natural\) and its
automorphism group. The assertion does not require an action of the Monster on the
digits: the Archimedean evaluation and the exceptional branch descend from a
common HMT preimage and preserve different invariants.

The neighboring function
\(T_{3A}=t_3+12+729/t_3\) belongs to the classical theory of Hauptmoduln of
order three. The active HMT realization in the tritic corridor is the class
\(3B\); an identification with \(3A\) requires declaration of the specific lift
that changes the class.
 \subsection{APP--Witt intertwining, Leech neighbor, and orbifold of order \(3\)}

The twelve fibers bear the labels
\[
 (i,\mu,\varepsilon)\in
 \mathbb Z/3\mathbb Z\times\{0,1\}\times\{0,1\},
 \qquad
 \lambda(i,\mu,\varepsilon)=4i+2\mu+\varepsilon\pmod{12}.
\]
The map \(\lambda\) is bijective and, once trivializations and frames are fixed,
the intertwiners satisfy
\[
 g_{\lambda(i,\mu,\varepsilon)}T_{i\mu\varepsilon}
 =T_{i\mu\varepsilon}C^{(-1)^\mu}.
\]
Its direct sum is an \(C_3\)-equivariant isomorphism
\[
 T_\lambda:W_{\mathrm{APP}}\otimes\mathbb C
 \xrightarrow{\ \cong\ }
 \bigoplus_{b=0}^{11}(A_{2,b}\otimes\mathbb C).
\]

Let \(N=N(A_2^{12})\), let \(v=v_\alpha\) the vector radial of norm \(54\),
and definamos
\[
 N_0=\ker\bigl(\langle\,\cdot\,,v\rangle\bmod3\bigr),
 \qquad N_v=N_0+\mathbb Z\frac v3.
\]
The lattice calculation verifies that \(N_v\cong\Lambda_{24}\), that the
isometry \(g_c\) has order three and has no nonzero fixed points, and that
\[
 \frac{g_cv-v}{3},\quad\frac{g_c^{-1}v-v}{3}\in N_0.
\]
Hence \(g_cN_v=N_v\). Its standard lift to the lattice VOA has twisted
weight \(4/3\) and type zero. The classical cyclic-orbifold theorems give
\[
 V_{(3)}:=\operatorname{Orb}_{\mathbb Z/3}
       (V_{N_v},\widehat g_c)\cong V^\natural,
 \qquad T_{3B}=t_3+12.
\]
The character is recovered by the average of the nine sectors and by
the level three uniformization:
\[
 \operatorname{ch}_{V_{(3)}}
 =\frac13\sum_{i,j\in\mathbb Z/3}Z_{i,j}=J,
\]
\[
 J=t_3+12+\frac{10\,3^9}{t_3}
       +\frac{4\,3^{14}}{t_3^2}+\frac{3^{18}}{t_3^3}.
\]
In particular,
\[
 [q]J=54+10\cdot3^9=196884=54(5\cdot729+1).
\]

\subsection{Constructions of orders \(2\) and \(3\) of the vertex operator algebra \texorpdfstring{\(V^\natural\)}{V natural}}

Fix the central extension and the Leech cocycle.
\[
 1\longrightarrow\{\pm1\}\longrightarrow\widehat\Lambda_{24}
 \longrightarrow\Lambda_{24}\longrightarrow1,
 \qquad
 V_\Lambda=M(1)\otimes\mathbb C_\varepsilon[\Lambda_{24}].
\]
The negation \(\lambda\mapsto-\lambda\) lifts to \(\theta\), and
\[
 t_2(\tau)=\operatorname{Tr}_{V_\Lambda}
  \left(\theta q^{L_0-1}\right)
 =q^{-1}\prod_{n\ge1}(1+q^n)^{-24}.
\]
The FLM presentation is
\[
 V_{(2)}:=(V_\Lambda)^+\oplus(V_\Lambda^T)^+\cong V^\natural.
\]

\begin{theorem}[Typed equivalence of the presentations of orders \(3\) and \(2\)]
\label{thm:cierre-vtm-equivalencia-voa-23}
Let
\[
 \phi_3:V_{(3)}\xrightarrow{\ \cong\ }V^\natural,
 \qquad
 \phi_2:V_{(2)}\xrightarrow{\ \cong\ }V^\natural
\]
the graded isomorphisms provided by the orbifold theorems
applied to the preceding data. Then
\[
 \boxed{\Phi_{23}:=\phi_2^{-1}\circ\phi_3:
 V_{(3)}\xrightarrow{\ \cong\ }V_{(2)}}
\]
is a graded isomorphism of vertex operator algebras. Moreover,
\[
 a\longmapsto\Phi_{23}a\Phi_{23}^{-1}
\]
identifies its automorphism groups and preserves the graded traces.
\end{theorem}

\begin{proof}
The composition of VOA isomorphisms is a VOA isomorphism and preserves the
vacuum, the conformal vector, the vertex operators, and the grading by
\(L_0\). Conjugation transports automorphisms. For every automorphism
\(a\), invariance of the trace under conjugation gives equality of the
corresponding graded series.
\end{proof}

The isomorphism \(\Phi_{23}\) depends on the identifications \(\phi_2\) and
\(\phi_3\); absolute canonicity is not asserted. The order-two construction
also preserves the exclusive content
\[
 T_{2A}=t_2+24+\frac{2^{12}}{t_2},
 \qquad
 \boxed{j=\frac{(t_2+256)^3}{t_2^2}},
\]
which completes the modular uniformization of order two.

Once \(V^\natural\) has been constructed, the three publications are stated without
unduly serializing them:
\[
 \operatorname{Aut}(V^\natural)\cong\mathbb M,
 \qquad
 \operatorname{ch}_{V^\natural}=J,
 \qquad
 g\longmapsto T_g(\tau)
 =\operatorname{Tr}_{V^\natural}\!\left(gq^{L_0-1}\right).
\]
The Moonshine theorem coordinates the graded action with genus-zero functions
and their replicability identities. The Monster action resides
in \(V^\natural\); the fact that it is not a pointwise permutation of digits is a final type
control, not the thesis of the corridor.

 \section{Continuous evidentiary chain of the exceptional publication: from the APP unfolding to Paley--Witt,
\texorpdfstring{\(V^\natural\)}{V natural} and Moonshine}
\label{sec:prueba-corredor-excepcional}

The exceptional publication starts from the incidence reader
\(\mathfrak E_{\rm exc}\) of \eqref{eq:frontal-doble-proyeccion}. It does not take an already given
code, lattice or exceptional group as input: it takes the
boundary, orientation, sheet and memory of an enriched history and
expresses them in the projective gauge fixed by nonadic holonomy. Only
afterwards are the classical theorems of codes, lattices and vertex operator
algebras applied to the produced objects.

\subsection{\texorpdfstring{APP unfolding and graded cone of the index \(54\)}{APP unfolding and graded cone of the index 54}}

Let \(R\) be the involution exchanging the APP sheets and let
\(P_+=(I+R)/2\), \(P_-=(I-R)/2\) be its projectors. The central block
\begin{equation}
 B_c=7I+2R=9P_++5P_-
 \label{eq:exc-bloque-central-app}
\end{equation}
produces four objects of distinct types:
\begin{equation}
 9-5=4,\qquad (B_c)_{i\ne j}=2,\qquad
 \sum M_\Pi-\sum M_\Sigma=51-45=6,\qquad
 9\cdot6=459-405=54.
 \label{eq:exc-despliegue-42654}
\end{equation}
The \(4\) is the spectral jump between sheets, the \(2\) the local coupling,
the \(6\) the difference of the modular compressions and the \(54\) the
integration of that difference over the nine positions. Therefore,
\(4\to2\to6\to54\) designates a hierarchy of morphisms and evaluations, not
a chain of equalities between objects.

With the orientation \(o_{\Sigma\Pi}\), sheet parity and triangular
normalization fixed, the Fock branch has the coupling
\begin{equation}
 \mathcal H_{\Sigma,\Pi}^{(m)}:
 M_{\rm bal}^{C_3,-}\longrightarrow
 \operatorname{End}\!\bigl(\mathcal F_2(\mathcal V_m)\bigr)
 \otimes o_{\Sigma\Pi},
 \label{eq:exc-acoplamiento-fock}
\end{equation}
and its oriented trace satisfies
\begin{equation}
 \operatorname{Tr}_{o,\mathcal F_2}
 \bigl(\mathcal H_{\Sigma,\Pi}^{(m)}(\nu)\Gamma_2(g_m)\bigr)
 =-\frac{3m\,c(\nu)}2.
 \label{eq:exc-traza-fock}
\end{equation}
For \(m=12\) and \(c(\delta)=-3\), \(54\) holds. Independently, the
cyclotomic product
\begin{equation}
 t_3(q)=q^{-1}\prod_{n\ge1}(1+q^n+q^{2n})^{-12}
 =q^{-1}-12+54q-76q^2-243q^3+\cdots
 \label{eq:exc-t3}
\end{equation}
gives \([q]t_3=54\).

\begin{theorem}[Typed cone of the common index]
\label{thm:exc-cono-54}
Once the APP census, the Fock branch, the cyclotomic
product, the lattice neighbor of
\cref{eq:exc-vecino-indice-tres} and the lunar module of
\cref{eq:exc-flm} have been constructed separately, their readers satisfy
\begin{equation}
 D_{\rm APP}
 =\operatorname{Tr}_{o,\mathcal F_2}
  \bigl(\mathcal H_{\Sigma,\Pi}^{(12)}(\delta)\Gamma_2(g_{12})\bigr)
 =[q]t_3=v^2=\frac{196884}{5\cdot729+1}=54.
 \label{eq:exc-cono-54}
\end{equation}
The equality gathers evaluations in a cone toward \(\mathbb Z\); it does not
identify their domains or use \(196884\) to select the corridor.
\end{theorem}

\begin{proof}
The first three equalities are
\eqref{eq:exc-despliegue-42654}, \eqref{eq:exc-traza-fock} and
\eqref{eq:exc-t3}. For the radial vector
\(v=4\rho_1+\rho_2+\cdots+\rho_{12}\) of
\eqref{eq:exc-vector-vecino}, orthogonality of the factors \(A_2\) gives
\(v^2=4^2\cdot2+11\cdot2=54\). Finally,
\(5\cdot729+1=3646\) and \(54\cdot3646=196884\). Each identity is proved in
its own codomain before they are compared.
\end{proof}

\subsection{\texorpdfstring{Paley form and internal quadratic root of \(-I_6\)}{Paley form and internal quadratic root of -I6}}

Let
\[
 \mathcal U_9=(\mathbb Z/9\mathbb Z)^\times=\{1,2,4,8,7,5\},
\]
in the order of the powers of \(2\), and let
\begin{equation}
 \theta:\mathcal U_9\xrightarrow{\sim}\mathbb P^1(\mathbb F_5),
 \qquad
 \begin{array}{c|cccccc}
 u&1&2&4&8&7&5\\ \hline
 \theta(u)&\infty&0&1&2&3&4
 \end{array}.
 \label{eq:exc-calibre-proyectivo}
\end{equation}
This bijection is a coordinate gauge of the reader; it is not used as a
homomorphism. For the quadratic character
\(\chi:\mathbb F_5\to\{-1,0,1\}\), define
\begin{equation}
 C_P=
 \begin{pmatrix}
 0& 1& 1& 1& 1& 1\\
 1& 0& 1&-1&-1& 1\\
 1& 1& 0& 1&-1&-1\\
 1&-1& 1& 0& 1&-1\\
 1&-1&-1& 1& 0& 1\\
 1& 1&-1&-1& 1& 0
 \end{pmatrix},
 \label{eq:exc-matriz-paley}
\end{equation}
where the finite off-diagonal entries are \(\chi(x-y)\) and the
entries incident with \(\infty\) are one.

\begin{lemma}[Quadratic correlation]
\label{lem:exc-correlacion-cuadratica}
If \(a\ne b\) in \(\mathbb F_5\), then
\[
 \sum_{t\in\mathbb F_5}\chi(t-a)\chi(t-b)=-1.
\]
\end{lemma}

\begin{proof}
The substitution \(u=(t-a)/(b-a)\) reduces the sum to
\(\sum_u\chi(u(u-1))\). For \(u=0,1,2,3,4\), the summands are
\(0,0,1,-1,-1\).
\end{proof}

\begin{proposition}[Conference identity and ternary reduction]
\label{prop:exc-paley-witt}
The matrix \(C_P\) is symmetric and satisfies
\[
 C_PC_P^{\mathsf T}=5I_6.
\]
Its reduction \(A_W=C_P\bmod3\) satisfies
\begin{equation}
 \boxed{A_W^2=-I_6\quad\text{in }M_6(\mathbb F_3).}
 \label{eq:exc-aw-cuadrado}
\end{equation}
\end{proposition}

\begin{proof}
Each row of \(C_P\) has squared norm five. Two distinct finite
rows have inner product \(1-1=0\) by the
\cref{lem:exc-correlacion-cuadratica}; the sum of a nontrivial character
also proves orthogonality with the row of \(\infty\). Reducing modulo
three yields
\(A_WA_W^{\mathsf T}=5I_6=2I_6=-I_6\). Since \(A_W\) is symmetric,
\eqref{eq:exc-aw-cuadrado} follows.
\end{proof}

The identity \eqref{eq:exc-aw-cuadrado} constructs the finite complex
structure of the corridor. Indeed,
\[
 \mathbb F_9=\mathbb F_3[\iota]/(\iota^2+1),
 \qquad
 P_{\pm\iota}=\frac12(I\mp\iota A_W)
\]
are complementary rank-three projectors and endow
\(\mathbb F_3^6\) with a rank-three \(\mathbb F_9\)-module structure.
The action is explicit:
\begin{equation}
 (a+b\iota)\cdot v=av+b(vA_W).
 \label{eq:exc-accion-f9}
\end{equation}
Therefore, \(3^6=9^3\) is not used as a cardinal substitute for an action.

\subsection{Self-dual ternary code and Witt system}

Define
\begin{equation}
 c:\mathbb F_3^6\longrightarrow\mathbb F_3^{12},
 \qquad c(q)=(q,qA_W),
 \qquad C_W:=\operatorname{im}c.
 \label{eq:exc-codigo-grafico}
\end{equation}

\begin{theorem}[Ternary code produced by incidence]
\label{thm:exc-codigo-ternario}
The map \(c\) is injective and
\begin{equation}
 C_W=C_W^\perp,
 \qquad
 W_{C_W}(y)=1+264y^6+440y^9+24y^{12}.
 \label{eq:exc-enumerador-cw}
\end{equation}
In particular, \(C_W\) has parameters \([12,6,6]_3\).
\end{theorem}

\begin{proof}
The first projection of \(c(q)\) is \(q\), so \(c\) is injective.
With \(G_W=(I_6\mid A_W)\),
\[
 G_WG_W^{\mathsf T}=I_6+A_WA_W^{\mathsf T}=0
\]
in \(\mathbb F_3\). Thus \(C_W\subseteq C_W^\perp\), and equality of
dimensions gives self-duality. The enumerator is the exact finite sum
\[
 \sum_{q\in\mathbb F_3^6}
 y^{\operatorname{wt}(q)+\operatorname{wt}(qA_W)}.
\]
Evaluating the \(3^6=729\) rows of the explicit matrix
\eqref{eq:exc-matriz-paley} gives, respectively, the counts
\(1,264,440,24\) at weights \(0,6,9,12\), and zero at the others. The sum
of the four counts is \(729\), so no word remains outside
the finite certificate.
\end{proof}

The \(264\) weight-six vectors occur in pairs \(\{x,-x\}\); their
\(132\) supports form the Steiner system
\begin{equation}
 W_{12}=S(5,6,12).
 \label{eq:exc-w12}
\end{equation}
This is the recognition, after constructing \(C_W\), of the extended ternary
Golay code and its Witt design
\citep{MacWilliamsSloane1977,ConwaySloane1999}. The design is not introduced
to select the matrix \(A_W\).

The binary prolongation remains separate. If
\(\mathcal G_{24}\subset\mathbb F_2^{24}\) is the extended binary
Golay code, \(D\) a dodecad and
\(\ell:\{1,\ldots,12\}_{\rm HMT}\to D\) a design isomorphism, the
weight-eight supports form \(W_{24}=S(5,8,24)\). For a fixed tetrad,
\[
 \lambda_4(W_{12})=4,
 \qquad
 \lambda_4(W_{24})=5,
\]
so that four octads lift the four incident hexads and a fifth
is transversal. The datum \((\mathcal G_{24},D,\ell)\) belongs to the binary
codomain; there is no canonical arrow \(W_{24}\to\Lambda_{24}\).

\subsection{Ternary gluing and Leech neighbor}

Let \(Q=A_2^{12}\), with \(\det Q=3^{12}\). The discriminant of each
factor \(A_2\) is \(\mathbb F_3\), and the gluing defined by
\eqref{eq:exc-codigo-grafico} is
\begin{equation}
 N:=\{x\in Q^*:x\bmod Q\in C_W\}.
 \label{eq:exc-niemeier-pegado}
\end{equation}

\begin{theorem}[Niemeier lattice produced by the code]
\label{thm:exc-niemeier}
The lattice \(N\) is positive definite, even and unimodular of rank
twenty-four, and its root system is \(A_2^{12}\). Consequently,
\(N\cong N(A_2^{12})\).
\end{theorem}

\begin{proof}
Self-duality of \(C_W\) makes the gluing integral. All its weights are
multiples of three, which makes the extension even. Moreover,
\([N:Q]=|C_W|=3^6\), and therefore
\[
 \det N=\frac{\det Q}{[N:Q]^2}
 =\frac{3^{12}}{3^{12}}=1.
\]
A nonzero discriminant class of \(A_2\) has minimum norm \(2/3\).
Since every nonzero word of \(C_W\) has weight at least six, every gluing
vector outside \(Q\) has norm at least four. The norm-two
vectors are exactly the roots of the twelve factors \(A_2\). The
classification of even unimodular lattices of rank twenty-four
identifies the indicated type \citep{ConwaySloane1999}.
\end{proof}

Let \(\rho_i\) be the highest root of the factor \(A_2^{(i)}\), normalized by
\(\rho_i^2=2\), and define
\begin{equation}
 v=4\rho_1+\rho_2+\cdots+\rho_{12},
 \qquad v^2=54.
 \label{eq:exc-vector-vecino}
\end{equation}
The ordered frame of the exceptional reader fixes which of the twelve fibers occupies
the first position; no Archimedean value intervenes in this selection.
Set
\begin{equation}
 \chi_v(x)=\langle x,v\rangle\bmod3,
 \quad N_0=\ker\chi_v,
 \quad N_v=N_0+\mathbb Z\frac v3.
 \label{eq:exc-vecino-indice-tres}
\end{equation}

\begin{lemma}[Structure of the neighbor]
\label{lem:exc-estructura-vecino}
The character \(\chi_v\) is surjective, \([N:N_0]=3\),
\(v/3\in N_0^*\), and \(N_v\) is even and unimodular.
\end{lemma}

\begin{proof}
For every root of a factor \(A_2\), the product with the highest root is
\(\pm1\) or \(\pm2\); the coefficients of \(v\) are congruent to one
modulo three. Therefore, no root belongs to \(N_0\) and \(\chi_v\) is
nontrivial, hence surjective. The definition of the kernel gives index three and
\(\det N_0=9\). For \(x\in N_0\),
\(\langle x,v/3\rangle\in\mathbb Z\), while
\((v/3)^2=6\). Thus the extension by \(v/3\) is integral and even. Since
\([N_v:N_0]=3\),
\(\det N_v=9/3^2=1\).
\end{proof}

\begin{lemma}[Minimum of the new classes]
\label{lem:exc-minimo-clases}
Each coordinate of a local class appearing in
\(N_0\pm v/3\) has minimum norm \(2/9\). Consequently,
\[
 \|x\pm v/3\|^2\ge12\cdot\frac29=\frac83
 \qquad(x\in N_0).
\]
\end{lemma}

\begin{proof}
Writing a coordinate of \(A_2^*\) in the simple-root basis, the
six relevant oriented classes are the shifts of
\(\pm\rho/3\) by the three discriminant classes. Completing squares in the
Gram form
\(\bigl(\begin{smallmatrix}2&-1\\-1&2\end{smallmatrix}\bigr)\)
gives minimum \(2/9\) in all six classes. The orthogonal sum of the twelve
coordinates gives the inequality.
\end{proof}

\begin{theorem}[Construction of the Leech lattice]
\label{thm:exc-leech}
The neighbor \(N_v\) of \eqref{eq:exc-vecino-indice-tres} is isometric to the
Leech lattice:
\begin{equation}
 \boxed{N_v\cong\Lambda_{24}.}
 \label{eq:exc-leech}
\end{equation}
\end{theorem}

\begin{proof}
The roots of \(N\) do not belong to \(N_0\), and the two new classes
contain no norm-two vectors by the
\cref{lem:exc-minimo-clases}. The \cref{lem:exc-estructura-vecino} proves that
\(N_v\) is positive definite, even, unimodular and of rank twenty-four. The
uniqueness of the lattice with these properties and no roots identifies it with
\(\Lambda_{24}\) \citep{ConwaySloane1999}.
\end{proof}

\subsection{Order-three isometry, orbifold and lunar module}

The Coxeter rotation of each factor \(A_2\) induces an order-three isometry
\(g_c\). Stability of the code \(C_W\) proves that
\(g_cN=N\). For the oriented frame fixing \(v\), the shifts
\begin{equation}
 y_*:=\frac{g_cv-v}{3},
 \qquad
 z_*:=\frac{g_c^{-1}v-v}{3}
 \label{eq:exc-desplazamientos-vecino}
\end{equation}
belong to \(N_0\): their discriminant words are a pair
\(c_*,-c_*\in C_W\), and
\(\langle y_*,v\rangle=\langle z_*,v\rangle=-27\). Therefore,
\begin{equation}
 g_cN_v=N_v.
 \label{eq:exc-preservacion-vecino}
\end{equation}
The action has no fixed points in \(N_v\otimes\mathbb R\), since each
factor \(A_2\) has the two primitive eigenvalues of order three.

Let \(\widehat g_c\) be the standard lift to the lattice VOA
\(V_{\Lambda_{24}}\). Its two nontrivial eigenspaces have dimension
twelve. The vacuum energy of the twisted module is therefore
\begin{equation}
 \rho_{g_c}
 =\frac1{4\cdot3^2}
 \sum_{k=1}^{2}k(3-k)\dim\mathfrak h_{(k)}
 =\frac43.
 \label{eq:exc-peso-torcido}
\end{equation}
The congruence \(3^2\rho_{g_c}=12\equiv0\pmod3\) proves that the lift
is of type zero.

\begin{theorem}[Triadic orbifold]
\label{thm:exc-orbifold}
The cyclic orbifold of the preceding lift satisfies
\begin{equation}
 \operatorname{Orb}_{\mathbb Z/3}
 (V_{\Lambda_{24}},\widehat g_c)\cong V^\natural.
 \label{eq:exc-orbifold}
\end{equation}
The dual automorphism belongs to the class \(3B\).
\end{theorem}

\begin{proof}
The three hypotheses that cannot be replaced by an equality of series
have been verified: fixed-point-free order-three isometry, preservation of the
neighbor and type zero. The cyclic orbifold theorems then identify
the holomorphic extension with the lunar module and classify the dual automorphism
\citep{ChenLamShimakura2016,AbeLamYamada2017,Carnahan2017}.
\end{proof}

The order-two construction is a parallel realization. For the lift
\(\theta\) of negation on \(\Lambda_{24}\),
\begin{equation}
 V^\natural=(V_{\Lambda_{24}})^+
 \oplus(V_{\Lambda_{24}}^T)^+.
 \label{eq:exc-flm}
\end{equation}
The order-two and order-three routes are identified by the VOA theorems; they are not
causally serialized one inside the other \citep{FLM1988}.

\begin{theorem}[Publications of the lunar module]
\label{thm:exc-moonshine}
Once \(V^\natural\) has been constructed, one obtains
\begin{align}
 \operatorname{Aut}(V^\natural)&\cong\mathbb M,
 \label{eq:exc-aut-monster}\\
 \operatorname{ch}_{V^\natural}(\tau)
 &=J(\tau)=j(\tau)-744
 =q^{-1}+196884q+21493760q^2+\cdots,
 \label{eq:exc-caracter-j}\\
 T_g(\tau)&=\operatorname{Tr}_{V^\natural}
 \bigl(gq^{L_0-1}\bigr),\qquad g\in\mathbb M.
 \label{eq:exc-mckay-thompson}
\end{align}
The series \(T_g\) are the genus-zero modular functions of the Moonshine
theorem and satisfy the replicability identities.
\end{theorem}

\begin{proof}
The first identification and the graded construction come from the
FLM realization. The coordination between the graded action, the McKay--Thompson
series and their modular properties is the Moonshine theorem of
Conway--Norton and Borcherds
\citep{FLM1988,ConwayNorton1979,Borcherds1992}. In particular,
\((V^\natural)_0=\mathbb C\mathbf1\) and
\((V^\natural)_1=0\); the coefficient \(196884=1+196883\) is the first
nontrivial decomposition into dimensions of Monster representations.
\end{proof}

\begin{falsifier}[Exceptional corridor]
\label{fals:exc-corredor}
The publication is refuted if any of the effective steps fails:
\eqref{eq:exc-despliegue-42654}, the trace
\eqref{eq:exc-traza-fock}, the coefficient \([q]t_3=54\),
\(C_PC_P^{\mathsf T}=5I_6\), \(A_W^2=-I_6\), self-duality or the
enumerator of \(C_W\), evenness or unimodularity of the gluing, absence
of roots of the neighbor, its preservation by \(g_c\), the twisted weight
\(4/3\) or type zero. A narrative is also refuted if it uses
\(W_{24}\) to construct \(\Lambda_{24}\), confuses the two sets
of \(243\) on account of their cardinality, or makes the Monster act directly
on digits and routes. The action of \(\mathbb M\) resides in \(V^\natural\).
\end{falsifier}
 
\chapter{Spectral--polar image and exact Riemann--Weil reduction}
\begingroup
\renewcommand{\chapter}[2][]{}%
\chapter{Spectral--polar image and exact Riemann--Weil reduction}

\section{Starting point: the complete spectral--polar structure}

The input to the resolution is not the isolated Weil form or a list of
zeros. It is the complete spectral--polar structure reproduced as data in this
volume. The causal chain specific to this chapter is
\begin{equation}\label{eq:pdf2-rh-cadena-procedencia}
 \mathfrak G_{\rm sp}
 \xrightarrow{\ \operatorname{pr}_{X,{\rm sp}}\ }X_{\chi_{\rm sp}}
 \xrightarrow{\ \mathcal A_{\rm sp}\ }\mathfrak M_{\rm sp}
 \xrightarrow{\ \mathcal P_{\rm sp}\ }\mathcal C_{\rm GW}.
\end{equation}
The classical codomain appears only in the last arrow; it does not
retrospectively select the state, the charts, the projection, or the terminal.

The starting structure contains the object of thirteen coordinates
\begin{equation}\label{eq:pdf2-rh-res-sp-completa}
\begin{aligned}
 \mathfrak G_{\rm sp}^{(13)}=\bigl(&
 \mathcal F_{\rm sp},
 \operatorname{Ext}_{\Gamma,{\rm sp}},
 \operatorname{Rel}_{\rm sp},
 \mathcal N_{\rm sp},
 \operatorname{or}_{\rm sp},
 \operatorname{Carr}_{\rm sp},
 \operatorname{Mem}_{\rm sp},\\
 &\partial_{\rm sp},
 \gamma_{\rm sp},
 \Sigma_{\rm sp}^{\rm mult},
 \operatorname{Surv}_{\rm sp},
 \tau_{\rm sp}^{\rm term},
 \mathcal L_{\rm sp}\bigr).
\end{aligned}
\end{equation}
Each coordinate enters substantively:
\begin{enumerate}
 \item \(\mathcal F_{\rm sp}\) preserves the two sheets, their orientation, the
 residue--quotient pair, and the signature of the double publication.
 \item \(\operatorname{Ext}_{\Gamma,{\rm sp}}\) transports that fibre through the
 directed class of cutoffs and nonadic levels.
    \item \(\operatorname{Rel}_{\rm sp}\) contains the carry cocycle, the
    orthogonality \(P=P^*=P^2\), the two publications, and the structural
    complement \(B_{\rm sp}^{\rm str}:=\Xi^*(I-P)\Xi\). The identification
    of this complement with the Guinand--Weil form is a distinct arrow
    and is audited separately in this chapter.
 \item \(\mathcal N_{\rm sp}\) records the level \(m\), the window \(R\),
 the quotient \(q=5/9\), the moments of both sheets, and the indices of the
 prime, gamma, and polar channels.
 \item \(\operatorname{or}_{\rm sp}=\widehat R=(R,K)\) keeps the two oriented
 charts distinct.
 \item \(\operatorname{Carr}_{\rm sp}=\operatorname{Carr}_R\) is the
 operatorial carry and is not annulled when a translation is published.
 \item \(\operatorname{Mem}_{\rm sp}=(S_0,M,\mathcal O_9)\) preserves the
 finite Stein identity and its terminal object.
 \item \(\partial_{\rm sp}=(I-P)\Xi\) is the boundary absent from the
 \(P\Xi\) compression and whose energy defines the structural complement. Its
 identification with the Weil form is not incorporated into this definition.
 \item \(\gamma_{\rm sp}\) is the cofinal route of windows, with a single
 regulator for all channels.
 \item \(\Sigma_{\rm sp}^{\rm mult}\) is the multisection of the two
 publications and their prime--gamma--polar components; it is not the signature of a
 single state.
 \item \(\operatorname{Surv}_{\rm sp}\) preserves the family under
 \(m\mapsto m+1\) and \(\mathcal D_R\hookrightarrow\mathcal D_{R'}\).
 \item \(\tau_{\rm sp}^{\rm term}\) preserves the terminal
 \(M^{*N}S_0M^N\), which in the mirror chart is published through \(E_R\).
    \item \(\mathcal L_{\rm sp}=(\Xi,P,B_{\rm sp}^{\rm str})\) is the final reader; it is applied
 after antecedent, charts, memory, route, and terminal object.
\end{enumerate}

The assembler and the publisher are
\begin{equation}\label{eq:pdf2-rh-ensamblador-sp}
\begin{aligned}
 \mathcal A_{\rm sp}(\operatorname{pr}_{X,{\rm sp}}\widehat x)
  ={}&(\widehat X_\infty,\Gamma_9,\widehat R,
      \operatorname{Carr}_R,S_0,M,\mathcal O_9,
      \Xi,P,\mathcal D_R^{\rm cof}),\\
 \mathfrak M_{\rm sp}={}&\operatorname{Im}\mathcal A_{\rm sp},\\
 \mathcal P_{\rm sp}(\mathfrak M_{\rm sp})
  ={}&(B_{\rm sp}^{\rm str},\Xi^*(I-P)\Xi,
       \mathscr I_{\rm GW},\Delta_{\rm bind},
       \text{cofinal prime--gamma--polar family}),\\
 \Delta_{\rm bind}:={}&B_{\rm sp}^{\rm str}-\mathscr I_{\rm GW}.
\end{aligned}
\end{equation}
The finite memory of the starting object satisfies
\begin{equation}\label{eq:pdf2-rh-stein-finito}
 S_0=\sum_{r=0}^{N-1}M^{*r}\mathcal O_9^*\mathcal O_9M^r
     +M^{*N}S_0M^N.
\end{equation}
The last summand remains visible until its extinction is proved.
Therefore, replacing \(\mathfrak G_{\rm sp}\) with the published shadow
\(B_{\rm sp}^{\rm str}\) changes the starting object; it is not an abbreviation
of the same theorem. The residual \(\Delta_{\rm bind}\) is preserved until the
source--first colligation materially proves its vanishing in
\eqref{eq:pdf2-rh-binding-source-first}.

\section{Cofinal domain and target form}

For $R>0$ is fixed
\[
 \mathcal D_R=C_c^\infty((-R/2,R/2)),\qquad
 \widehat\psi(t)=\int_{\mathbb R}\psi(x)e^{-itx}\,dx,
\]
and, for $\psi,\phi\in\mathcal D_R$,
\[
 h_{\psi,\phi}(z)=\widehat\psi(z)
 \overline{\widehat\phi(\overline z)},\qquad
 \check h_{\psi,\phi}(a)=\langle\psi,T_a\phi\rangle.
\]
The union $\bigcup_R\mathcal D_R=C_c^\infty(\mathbb R)$ is exact. For every
finite set of distinct complex points, the evaluations
$\psi\mapsto(\widehat\psi(z_1),\ldots,\widehat\psi(z_n))$ are jointly
surjective when $R$ is sufficiently large. Therefore, this cofinal class
separates the evaluations required by the Weil criterion.

\begin{lemma}[Proof of the joint separation]
\label{lem:pdf2-rh-separacion}
The map
\[
 \mathcal D_R\longrightarrow\mathbb C^n,
 \qquad
 \psi\longmapsto
 (\widehat\psi(z_1),\ldots,\widehat\psi(z_n))
\]
is surjective for every set of distinct complex points
\(z_1,\ldots,z_n\), once \(R\) contains a nonempty interval.
\end{lemma}

\begin{proof}
If the evaluation functionals were linearly dependent, there would exist
\(c_1,\ldots,c_n\), not all zero, such that
\[
 \int_{\mathbb R}\psi(x)
 \left(\sum_{j=1}^nc_je^{-iz_jx}\right)\,dx=0
 \qquad(\psi\in\mathcal D_R).
\]
The analytic function \(\sum_jc_je^{-iz_jx}\) would vanish at
\((-R/2,R/2)\), hence throughout the plane. The independence of exponentials
with distinct exponents forces \(c_j=0\) for every \(j\), a contradiction.
The functionals are independent; since the codomain is finite-dimensional,
the map has rank \(\mathbb C^n\).
\end{proof}

\section{The two columns and the three channels computed before the sign}

Be $\vartheta_L$ a regulador comun. With
\[
 \ell_{p,k}=k\log p,\qquad w_{p,k}=(\log p)p^{-k/2},\qquad
 \mu_\Gamma(a)=\frac{e^{-a/2}}{1-e^{-2a}},
\]
the uniform inclusion $j_m$ and the sewing channel
$\widehat\delta_{a,m}$, which satisfy, are used
\[
 j_m^*j_m=I,\qquad
 \widehat\delta_{a,m}^*\widehat\delta_{b,m}
 =(I-T_a)^*(I-T_b).
\]
The regulated columns are
\begin{align*}
 J_{+,m,R,L}\psi={}&
 \Bigl(
  (\sqrt{\vartheta_L(\ell_{p,k})w_{p,k}}\,
       \widehat\delta_{\ell_{p,k},m}\psi)_{\ell_{p,k}\leq R},\\
 &\hspace{9mm}a\mapsto
  \sqrt{\vartheta_L(a)\mu_\Gamma(a)}\,
       \widehat\delta_{a,m}\psi,\ \zeta_+b_+(\psi)
 \Bigr),\\
 J_{-,m,R,L}\psi={}&
 \Bigl(
  (\sqrt{2\vartheta_L(\ell_{p,k})w_{p,k}}\,j_m\psi)_{\ell_{p,k}\leq R},
  \sqrt{-c_\Gamma}\,j_m\psi,\ \zeta_-b_-(\psi)
 \Bigr),
\end{align*}
where $c_\Gamma=\psi_\Gamma(1/4)-\log\pi<0$ and
$b_\pm=(\widehat\psi(i/2)\pm\widehat\psi(-i/2))/\sqrt2$.

A term-by-term expansion, without using zeros or the sought sign, gives
\begin{align*}
 &\langle J_{+,m,R,L}\psi,J_{+,m,R,L}\phi\rangle
 -\langle J_{-,m,R,L}\psi,J_{-,m,R,L}\phi\rangle
 =\mathscr I_{{\rm GW},L}(\psi,\phi),\\
 \mathscr I_{{\rm GW},L}(\psi,\phi)
 ={}&h_{\psi,\phi}(i/2)+h_{\psi,\phi}(-i/2)
 +c_\Gamma\check h_{\psi,\phi}(0)\\
 &+\int_0^\infty\vartheta_L(a)\mu_\Gamma(a)
 [2\check h(0)-\check h(a)-\check h(-a)]\,da\\
 &-\sum_{p,k}\vartheta_L(\ell_{p,k})w_{p,k}
 [\check h(\ell_{p,k})+\check h(-\ell_{p,k})].
\end{align*}
The truncation of both columns at \(\ell_{p,k}\leq R\) is mandatory, not
notational. If \(\psi,\phi\in\mathcal D_R\), then
\(\check h_{\psi,\phi}(\pm\ell)=0\) for \(\ell>R\); the contribution of each
omitted pair of coordinates to the \emph{difference} of Grams is therefore zero.
The untruncated negative column, by contrast, would have infinite norm because
\(\sum_p(\log p)p^{-1/2}=\infty\). In the finite set
\(\ell_{p,k}\leq R\) the prime coordinates converge individually when
\(L\to\infty\). In the gamma channel, the integrand is \(O(a)\) at zero and
\(\mu_\Gamma(a)=O(e^{-a/2})\) at infinity, so convergence is also
in norm. We therefore define
\begin{equation}\label{eq:pdf2-rh-columnas-reducidas}
 J_{\pm,m,R}:=\lim_{L\to\infty}J_{\pm,m,R,L}
 \quad\text{in their truncated-coordinate Hilbert spaces},
\end{equation}
and we suppress \(m\) when it remains fixed. By dominated convergence,
\[
 \mathscr I_{{\rm GW},L}\longrightarrow\mathscr I_{\rm GW}
 \qquad(L\to\infty),
\]
and the identities are natural under $m\mapsto m+1$ and $R\le R'$.

The reduced columns \(J_{+,R}\) and \(J_{-,R}\) are now material
operators from \(\mathcal D_R\) to Hilbert. The explicit expansion proves, without
consulting zeros or sign,
\begin{equation}\label{eq:pdf2-rh-gram-difference-gw}
 \langle J_{+,R}f,J_{+,R}g\rangle
 -\langle J_{-,R}f,J_{-,R}g\rangle
 =\mathscr I_{\rm pr}(f,g)+\mathscr I_\Gamma(f,g)
  +\mathscr I_{\rm pol}(f,g)
 =\mathscr I_{\rm GW}(f,g).
\end{equation}
The clocks \(\ell_{p,k}\) produce the prime summand; the integral with
\(\mu_\Gamma\) and \(c_\Gamma\check h(0)\), the gamma summand; and
\[
 \langle b_+(f),b_+(g)\rangle-
 \langle b_-(f),b_-(g)\rangle
 =h_{f,g}(i/2)+h_{f,g}(-i/2)
\]
produce the polar term. This part of the calculation is material and independent
of the structural bridge.

\section{Source--first colligation of the spectral--polar image}

We first fix a nonadic level \(m\); this index is suppressed in
\(J_{\pm,R}\), \(E_R\) and \(\Sigma_R\) until the encoder is constructed. The
decomposition of the explicit codomains of
\eqref{eq:pdf2-rh-columnas-reducidas} is
\[
 I_+=\{\mathrm{pr},\Gamma,\mathrm{pol}\},\qquad
 I_-(R)=\{\Gamma0,\mathrm{pol}\}
 \sqcup\{(p,k):\ell_{p,k}\leq R\},
\]
\[
 \mathscr H_{+,R}^{\rm amb}=\bigoplus_{i\in I_+}H_{+,i,R},
 \qquad
 \mathscr H_{-,R}^{\rm amb}=\bigoplus_{j\in I_-(R)}H_{-,j,R}.
\]
We do not identify a reachable subspace with a sum of its coordinate
projections. With the ambient inclusions understood, we define
\[
 H_{+,R}:=\overline{\operatorname{Ran}J_{+,R}},\qquad
 H_{-,R}:=\overline{\operatorname{Ran}J_{-,R}},
\]
as closed subspaces of the two ambient Hilbert spaces. The structure \(\mathfrak G_{\rm sp}\) contains, as part of its mathematical data and before the classical expression, an antispecular state space \(\mathcal H_R^{\rm sh}\) and the source--first isometric row
\begin{equation}\label{eq:pdf2-rh-coligacion-source-first}
 \Sigma_R=
 \begin{pmatrix}\mathcal K_R^{\rm src}\\ \mathcal F_R^{\rm src}\end{pmatrix}:
 H_{+,R}\longrightarrow
 H_{-,R}\oplus\mathcal H_R^{\rm sh},
 \qquad
 \Sigma_R^*\Sigma_R
 = (\mathcal K_R^{\rm src})^*\mathcal K_R^{\rm src}
  +(\mathcal F_R^{\rm src})^*\mathcal F_R^{\rm src}=I_{H_{+,R}}.
\end{equation}
Its reachable action is fixed on the entire positive column by
the relations of the two publications that form part of
\(\operatorname{Rel}_{\rm sp}\). To prevent any substitution of the reader,
we denote by
\((\mathcal X_R^{\rm str},\Xi_R^{\rm str},P_R^{\rm str})\) the restriction to
\(\mathcal D_R\) of the structural reader of
\eqref{eq:pdf2-rh-res-sp-completa}. Its two moments are, literally,
\begin{equation}\label{eq:pdf2-rh-momentos-lector-estructural}
\begin{aligned}
 \langle\Xi_R^{\rm str}f,\Xi_R^{\rm str}g\rangle
  &=\langle J_{+,R}f,J_{+,R}g\rangle,\\
 \langle P_R^{\rm str}\Xi_R^{\rm str}f,
          P_R^{\rm str}\Xi_R^{\rm str}g\rangle
  &=\langle J_{-,R}f,J_{-,R}g\rangle,
 \qquad (P_R^{\rm str})^*= (P_R^{\rm str})^2=P_R^{\rm str}.
\end{aligned}
\end{equation}
The boundary coordinate comes with the structural isometry.
\[
 \jmath_R^{\rm sh}:
 \overline{\operatorname{Ran}((I-P_R^{\rm str})\Xi_R^{\rm str})}
 \longrightarrow\mathcal H_R^{\rm sh}.
\]
By definition of the source--first row, not by a subsequent identification,
\begin{equation}\label{eq:pdf2-rh-source-state-output}
 E_R:=\mathcal F_R^{\rm src}J_{+,R}
     =\jmath_R^{\rm sh}(I-P_R^{\rm str})\Xi_R^{\rm str}:
 \mathcal D_R\longrightarrow\mathcal H_R^{\rm sh},
 \qquad
 \Omega_{0,R}:=\mathcal K_R^{\rm src}J_{+,R}-J_{-,R}.
\end{equation}
There is no colligation by channel: the same row simultaneously receives the
prime--power clocks, the gamma integral, the scalar mode, and the two polar
lines. In particular, it preserves all cross-products that arise upon
polarizing \eqref{eq:pdf2-rh-gram-difference-gw}.
With the ambient projections \(\pi_j^-:\mathscr H_{-,R}^{\rm amb}
\to H_{-,j,R}\), the correctly typed output rows are
\begin{equation}\label{eq:pdf2-rh-bloques-coligacion}
 \mathcal K_{j,R}^{\rm src}
 :=\pi_j^-\mathcal K_R^{\rm src}:
 H_{+,R}\longrightarrow H_{-,j,R},
 \qquad
 \mathcal F_R^{\rm src}:H_{+,R}\longrightarrow\mathcal H_R^{\rm sh}.
\end{equation}
Thus, \(\mathcal K_R^{\rm src}\) is not applied to a positive coordinate that
might not by itself belong to \(H_{+,R}\). All cross-terms between
channels remain within
\((\mathcal K_R^{\rm src})^*\mathcal K_R^{\rm src}\) and
\((\mathcal F_R^{\rm src})^*\mathcal F_R^{\rm src}\); they are not orthogonalized
before the energy identity.

The survival coordinate of
\(\mathfrak G_{\rm sp}\) provides residual spaces
\(\mathcal R_{n,R}^{\rm sh}\), continuations
\(\Omega_{n,R}:\mathcal D_R\to\mathcal R_{n,R}^{\rm sh}\), and antispectral
isometries
\(V_{n,R}:\mathcal R_{n+1,R}^{\rm sh}\to\mathcal R_{n,R}^{\rm sh}\). Its
first space and its first operator are, by typed definition,
\[
 \mathcal R_{0,R}^{\rm sh}:=H_{-,R},\qquad
 \Omega_{0,R}:=\mathcal K_R^{\rm src}J_{+,R}-J_{-,R}.
\]
Its
full cycle recurrence is
\begin{equation}\label{eq:pdf2-rh-recurrencia-residual}
 q=\frac59,\qquad
 \Omega_{n,R}=qV_{n,R}\Omega_{n+1,R},
 \qquad V_{n,R}^*V_{n,R}=I,
 \qquad
 \sup_{n\geq0}\|\Omega_{n,R}f\|<\infty
 \quad(f\in\mathcal D_R).
\end{equation}
This recurrence belongs to the enriched colligation: it transports the
output difference, not the Weil form or its sign.

\begin{proposition}[Coinductive Cancellation of Output Difference]
\label{prop:pdf2-rh-omega-cero}
For every \(R>0\) and every \(f\in\mathcal D_R\),
\begin{equation}\label{eq:pdf2-rh-output-binding-source}
 \Omega_{n,R}f=0\quad(n\geq0),
 \qquad
 \boxed{\mathcal K_R^{\rm src}J_{+,R}=J_{-,R}}.
\end{equation}
\end{proposition}

\begin{proof}
Iterating \eqref{eq:pdf2-rh-recurrencia-residual} \(N\) times and using that each
\(V_{n,R}\) is isometric gives
\[
 \|\Omega_{n,R}f\|
 =q^N\|\Omega_{n+N,R}f\|
 \leq q^N\sup_{j\geq n}\|\Omega_{j,R}f\|.
\]
The right-hand member converges to zero because \(0<q=5/9<1\). Therefore
\(\Omega_{n,R}f=0\); the case \(n=0\) and
\eqref{eq:pdf2-rh-source-state-output} produce the framed identity.
\end{proof}

The row-by-row action of the output is no longer implicit. The projections are
the restrictions to \(H_{-,R}\) of the projections in the ambient Hilbert space, and
\begin{equation}\label{eq:pdf2-rh-salidas-generador-a-generador}
\begin{aligned}
 \mathcal K_{\Gamma0,R}^{\rm src}J_{+,R}f
  &=\sqrt{-c_\Gamma}\,j_mf,\\
 \mathcal K_{(p,k),R}^{\rm src}J_{+,R}f
  &=\sqrt{2w_{p,k}}\,j_mf
  &&(\ell_{p,k}\leq R),\\
 \mathcal K_{{\rm pol},R}^{\rm src}J_{+,R}f
  &=\zeta_-b_-(f).
\end{aligned}
\end{equation}
These rows exhaust \(J_{-,R}\), while
\begin{equation}\label{eq:pdf2-rh-frontera-generador-a-generador}
 \mathcal F_R^{\rm src}J_{+,R}f=E_Rf
\end{equation}
preserves, with all its crossings, the boundary coordinate.
Evaluating \eqref{eq:pdf2-rh-coligacion-source-first} at \(J_{+,R}f\) and
\(J_{+,R}g\), and using \eqref{eq:pdf2-rh-output-binding-source}, produces the
energy identity prior to the sign
\begin{equation}\label{eq:pdf2-rh-energia-source-first}
 \boxed{
 \langle J_{+,R}f,J_{+,R}g\rangle
 =\langle J_{-,R}f,J_{-,R}g\rangle
  +\langle E_Rf,E_Rg\rangle.}
\end{equation}
Comparison with the independent calculation
\eqref{eq:pdf2-rh-gram-difference-gw} yields, without introducing a root of
the target form,
\begin{equation}\label{eq:pdf2-rh-binding-source-first}
 \boxed{E_R^*E_R=\mathscr I_{\rm GW}\big|_{\mathcal D_R}.}
\end{equation}
In particular,
\(\Delta_{{\rm bind},R}
:=E_R^*E_R-\mathscr I_{\rm GW}|_{\mathcal D_R}=0\); this annihilation is the
output of \eqref{eq:pdf2-rh-coligacion-source-first}--
\eqref{eq:pdf2-rh-energia-source-first}, not a premise inserted into the
publisher.

\section{Structural complement, mirror sheet and terminal telescopia}

The state \(E_R=\mathcal F_R^{\rm src}J_{+,R}\) constructed by the
source--first row is the boundary coordinate of
\(\mathfrak G_{\rm sp}\). We define its structural form by
\begin{equation}\label{eq:pdf2-rh-bw-estructural}
 \mathcal B_{{\rm sp},R}^{\rm str}
 :=E_R^*E_R
 =\mathscr I_{\rm GW}\big|_{\mathcal D_R},
\end{equation}
where the second equality is
\eqref{eq:pdf2-rh-binding-source-first}, already deduced from the colligation and from
the generator-by-generator expansion.

Let \(P_+^{\rm sh}\) and \(P_-^{\rm sh}\) be the complementary orthogonal
projectors of the two sheets. Fix
\begin{equation}\label{eq:pdf2-rh-q-a-d}
 q=\frac59,\qquad
 A=P_+^{\rm sh}+qP_-^{\rm sh},\qquad
 D=\sqrt{1-q^2}\,P_-^{\rm sh}.
\end{equation}
Orthogonality gives the exact identity
\begin{equation}\label{eq:pdf2-rh-conservacion-local}
 A^*A+D^*D
 =P_+^{\rm sh}+q^2P_-^{\rm sh}
 +(1-q^2)P_-^{\rm sh}=I.
\end{equation}
The nine phases coordinate a single memory-enabled rotation:
\begin{equation}\label{eq:pdf2-rh-monodromia}
 M_9=\Phi A_8\cdots A_0=qV_9,
 \qquad V_9^*V_9=I.
\end{equation}
The contraction per turn is \(q\), not \(q^9\); the internal phases preserve
order and carry.

The colligation types
\(E_R:\mathcal D_R\to\operatorname{Ran}P_-^{\rm sh}\). Therefore,
\begin{equation}\label{eq:pdf2-rh-binding-espejo}
 P_+^{\rm sh}E_R=0,
 \qquad AE_R=qE_R,
 \qquad
 E_R^*E_R=\mathcal B_{{\rm sp},R}^{\rm str}
 =\mathscr I_{\rm GW}\big|_{\mathcal D_R}.
\end{equation}

\begin{proposition}[Finite Telescopy with Conserved Terminal]
\label{prop:pdf2-rh-telescopia}
for each \(N\geq1\),
\begin{equation}\label{eq:pdf2-rh-telescopia-finita}
 \mathcal B_{{\rm sp},R}^{\rm str}
 =\sum_{n=0}^{N-1}(DA^nE_R)^*(DA^nE_R)
 +(A^NE_R)^*(A^NE_R).
\end{equation}
Equivalently, evaluated at \(f,g\in\mathcal D_R\), the finite identity of
the colligation is
\begin{equation}\label{eq:pdf2-rh-balance-finito-columnas}
\boxed{
\begin{aligned}
 &\langle J_{+,R}f,J_{+,R}g\rangle
  -\langle J_{-,R}f,J_{-,R}g\rangle\\
 &\quad=\sum_{n=0}^{N-1}
  \langle DA^nE_Rf,DA^nE_Rg\rangle
  +\langle A^NE_Rf,A^NE_Rg\rangle.
\end{aligned}}
\end{equation}
The last term satisfies
\begin{equation}\label{eq:pdf2-rh-terminal-extincion}
 (A^NE_R)^*(A^NE_R)
 =q^{2N}\mathcal B_{{\rm sp},R}^{\rm str}
 \xrightarrow[N\to\infty]{}0.
\end{equation}
\end{proposition}

\begin{proof}
Applying \eqref{eq:pdf2-rh-conservacion-local} to \(A^nE_R\),
\[
 (A^nE_R)^*(A^nE_R)
 =(DA^nE_R)^*(DA^nE_R)
 +(A^{n+1}E_R)^*(A^{n+1}E_R).
\]
The sum from \(n=0\) to \(N-1\) cancels only the intermediate terms and
preserves \((A^NE_R)^*(A^NE_R)\). By
\eqref{eq:pdf2-rh-binding-espejo}, \(A^NE_R=q^NE_R\); hence the terminal object
is \(q^{2N}E_R^*E_R=q^{2N}\mathcal B_{{\rm sp},R}^{\rm str}\). Since
\(0<q=5/9<1\), it converges geometrically to zero.
\end{proof}

Only after proving \eqref{eq:pdf2-rh-terminal-extincion} is defined
\begin{equation}\label{eq:pdf2-rh-l-r}
 \mathscr L_Rf=(DA^nE_Rf)_{n\geq0}.
\end{equation}
The limit of finite telescoping gives
\begin{equation}\label{eq:pdf2-rh-lstar-l}
 \boxed{
 \mathscr L_R^*\mathscr L_R
 =\mathcal B_{{\rm sp},R}^{\rm str}
 =\mathscr I_{\rm GW}\big|_{\mathcal D_R}\succeq0.}
\end{equation}

\section{Common encoder, projection, and explicit double publication}

Fix \(m\geq1\) and two orthonormal vectors
\(u_m,\eta_m\in\mathbb C^{9^m}\), where \(u_m\) is the uniform mode and
\(\eta_m\) the carry mode. The coder space and projection are
\begin{equation}\label{eq:pdf2-rh-hilbert-comun}
 \mathcal X_{m,R}^{\rm sp}
 :=\mathbb C^{9^m}\widehat\otimes
   (H_{-,R}\oplus\mathcal H_R^{\rm sh}),
 \qquad
 P_{m,R}:=|u_m\rangle\langle u_m|\otimes
 I_{H_{-,R}\oplus\mathcal H_R^{\rm sh}},
 \qquad Q_{m,R}:=I-P_{m,R}.
\end{equation}
The complete action of the five-channel encoder is
\begin{equation}\label{eq:pdf2-rh-xi-comun-explicita}
 \boxed{
 \begin{aligned}
  \mathfrak C_{m,R}f
  &:=u_m\otimes(J_{-,R}f,0)
    +\eta_m\otimes(0,E_Rf),\\
  \Xi_{m,R}f&:=\mathfrak C_{m,R}f.
 \end{aligned}}
\end{equation}
No root of \(\mathscr I_{\rm GW}\) is taken: the two inputs are the
output and state of the source--first colligation already defined in
\eqref{eq:pdf2-rh-source-state-output}.

\begin{proposition}[The two moments arise from the action of the encoder]
\label{prop:pdf2-rh-vmas-isometria}
For every \(f,g\in\mathcal D_R\),
\begin{equation}\label{eq:pdf2-rh-dos-momentos-construidos}
\begin{aligned}
 \langle\mathfrak C_{m,R}f,\mathfrak C_{m,R}g\rangle
  &=\langle J_{+,R}f,J_{+,R}g\rangle,\\
 \langle P_{m,R}\mathfrak C_{m,R}f,
         P_{m,R}\mathfrak C_{m,R}g\rangle
  &=\langle J_{-,R}f,J_{-,R}g\rangle,\\
 \langle Q_{m,R}\Xi_{m,R}f,Q_{m,R}\Xi_{m,R}g\rangle
  &=\mathscr I_{\rm GW}(f,g).
\end{aligned}
\end{equation}
\end{proposition}

\begin{proof}
The orthogonality \(u_m\perp\eta_m\) and
\eqref{eq:pdf2-rh-energia-source-first} give
\[
 \langle\mathfrak C_{m,R}f,\mathfrak C_{m,R}g\rangle
 =\langle J_{-,R}f,J_{-,R}g\rangle
  +\langle E_Rf,E_Rg\rangle
 =\langle J_{+,R}f,J_{+,R}g\rangle.
\]
The projection \(P_{m,R}\) preserves the uniform summand and eliminates the
carry summand; \(Q_{m,R}\) does the reverse. The other two identities follow from
\eqref{eq:pdf2-rh-binding-source-first}.
\end{proof}

The newly written encoder is an explicit representative of the initial
structural reader, not a substitute reader. In the reachable subspaces one
defines
\begin{equation}\label{eq:pdf2-rh-isomorfismo-lector-estructural}
\begin{aligned}
 \mathcal U_{m,R}^{\rm str}
  (P_R^{\rm str}\Xi_R^{\rm str}f)
   &:=u_m\otimes(J_{-,R}f,0),\\
 \mathcal U_{m,R}^{\rm str}
  ((I-P_R^{\rm str})\Xi_R^{\rm str}f)
   &:=\eta_m\otimes(0,E_Rf).
\end{aligned}
\end{equation}
The two sources are orthogonal. The first Gram equality of each source
is \eqref{eq:pdf2-rh-momentos-lector-estructural}; for the second one uses
\(E_R=\jmath_R^{\rm sh}(I-P_R^{\rm str})\Xi_R^{\rm str}\). Hence the two
rules are well defined, are isometric, and extend to the orthogonal sum
of their closures. In that reachable closure they satisfy
\begin{equation}\label{eq:pdf2-rh-entrelazamiento-lector-estructural}
 \mathcal U_{m,R}^{\rm str}\Xi_R^{\rm str}=\Xi_{m,R},
 \qquad
 \mathcal U_{m,R}^{\rm str}P_R^{\rm str}
 =P_{m,R}\mathcal U_{m,R}^{\rm str},
 \qquad
 (\Xi_R^{\rm str})^*(I-P_R^{\rm str})\Xi_R^{\rm str}=E_R^*E_R.
\end{equation}
Thus, \(B_{{\rm sp},R}^{\rm str}\), the complement of the starting
structure, is exactly the complement of the constructed coder; the object has not
been changed in passing from \eqref{eq:pdf2-rh-res-sp-completa} to
\eqref{eq:pdf2-rh-xi-comun-explicita}.

The rule
\begin{equation}\label{eq:pdf2-rh-vmas-rango}
 V_{+,m,R}^{0}(J_{+,R}f):=\Xi_{m,R}f
\end{equation}
is well defined and is isometric by the first equality of
\eqref{eq:pdf2-rh-dos-momentos-construidos}; it extends uniquely to
an isometry
\(V_{+,m,R}:H_{+,R}\to\mathcal X_{m,R}^{\rm sp}\). The negative publication
is the isometry
\begin{equation}\label{eq:pdf2-rh-vmenos}
 V_{-,m,R}:H_{-,R}\longrightarrow\mathcal X_{m,R}^{\rm sp},
 \qquad V_{-,m,R}y=u_m\otimes(y,0).
\end{equation}
By construction,
\begin{equation}\label{eq:pdf2-rh-doble-publicacion-material}
 \boxed{
 V_{+,m,R}J_{+,R}=\Xi_{m,R},
 \qquad
 V_{-,m,R}J_{-,R}=P_{m,R}\Xi_{m,R}.}
\end{equation}
The complement preserves exactly the antiespecular state and its column of
remark:
\begin{equation}\label{eq:pdf2-rh-complemento-comun}
 \boxed{
 \Xi_{m,R}^*Q_{m,R}\Xi_{m,R}
 =E_R^*E_R
 =\mathscr L_R^*\mathscr L_R
 =\mathscr I_{\rm GW}\big|_{\mathcal D_R}.}
\end{equation}

\begin{proposition}[Induced connector and source--first coincidence]
\label{prop:pdf2-rh-conector}
The operator
\begin{equation}\label{eq:pdf2-rh-k-r}
 \mathcal K_R
 :=V_{-,m,R}^*P_{m,R}V_{+,m,R}:
 H_{+,R}\longrightarrow H_{-,R}
\end{equation}
is contractive, coincides with \(\mathcal K_R^{\rm src}\) throughout
\(H_{+,R}\), and satisfies
\begin{equation}\label{eq:pdf2-rh-k-binding}
 \|\mathcal K_R\|\leq1,
 \qquad
 \mathcal K_RJ_{+,R}=J_{-,R}.
\end{equation}
\end{proposition}

\begin{proof}
The two charts are isometries and \(P_{m,R}\) is an orthogonal projection,
hence \(\|\mathcal K_R\|\leq1\). Over the dense range of \(J_{+,R}\),
\[
 \mathcal K_RJ_{+,R}f
 =V_{-,m,R}^*P_{m,R}\Xi_{m,R}f
 =V_{-,m,R}^*V_{-,m,R}J_{-,R}f
 =J_{-,R}f.
\]
The same identity holds for \(\mathcal K_R^{\rm src}\) by
\eqref{eq:pdf2-rh-output-binding-source}; continuity and density
prove that the two operators coincide.
\end{proof}

The Gram difference is thus published in all its coincident realizations:
\begin{equation}\label{eq:pdf2-rh-tres-publicaciones}
\boxed{
\begin{aligned}
 J_{+,R}^*J_{+,R}-J_{-,R}^*J_{-,R}
 &=\Xi_{m,R}^*Q_{m,R}\Xi_{m,R}\\
 &=E_R^*E_R\\
 &=\mathscr L_R^*\mathscr L_R\\
 &=\mathcal B_{{\rm sp},R}^{\rm str}=\mathscr I_{\rm GW}.
\end{aligned}}
\end{equation}

\section{Cofinal Naturality and Reversible Weil Criterion}

For \(R\leq R'\), let
\(i_{R,R'}:\mathcal D_R\hookrightarrow\mathcal D_{R'}\). The individual
columns do not form an isometric system: when a new length
\(R<\ell_{p,k}\leq R'\) appears, both acquire the same positive energy. For
\(f,g\in\mathcal D_R\), separation of the supports gives exactly
\begin{equation}\label{eq:pdf2-rh-naturalidad}
 \left\langle(I-T_{\ell_{p,k}})f,
                    (I-T_{\ell_{p,k}})g\right\rangle
 =2\langle f,g\rangle
 =\left\langle\sqrt2j_mf,\sqrt2j_mg\right\rangle .
\end{equation}
Therefore the prime increment of
\(J_{+,R'}^*J_{+,R'}\) coincides with that of
\(J_{-,R'}^*J_{-,R'}\) and cancels in the difference, but no false
isometry between the columns is asserted. The correct cofinal compatibility is that
of the forms:
\begin{equation}\label{eq:pdf2-rh-naturalidad-p}
 \boxed{
 \mathscr I_{{\rm GW},R}
 =i_{R,R'}^*\mathscr I_{{\rm GW},R'}i_{R,R'},
 \qquad
 \mathcal B_{{\rm sp},R}^{\rm str}
 =i_{R,R'}^*\mathcal B_{{\rm sp},R'}^{\rm str}i_{R,R'}.}
\end{equation}
Since \(E_R^*E_R=\mathscr L_R^*\mathscr L_R=\mathscr I_{{\rm GW},R}\), the
minimal Kolmogorov realizations do admit unique isometries over their
reachable ranges, defined by
\begin{equation}\label{eq:pdf2-rh-naturalidad-bw}
 \iota_{R,R'}^E(E_Rf):=E_{R'}f,
 \qquad
 \iota_{R,R'}^L(\mathscr L_Rf):=\mathscr L_{R'}f
 \quad(f\in\mathcal D_R).
\end{equation}
These two rules are well defined precisely by
\eqref{eq:pdf2-rh-naturalidad-p}; they are not extended by fiat to
\(J_{\pm,R}\), \(\Xi_{m,R}\), or \(\mathcal K_R\). At fixed window, the nonadic refinement
\(m\mapsto m+1\) remains within the structural isometries
that transport \(j_m\), \(\widehat\delta_{a,m}\), \(u_m\), and \(\eta_m\).

To avoid compressing the last step, we also fix its objects and
normalizations. Let
\[
 \xi(s):=\frac12s(s-1)\pi^{-s/2}\Gamma(s/2)\zeta(s)
\]
and let \(\mathscr Z\) be the multiset of its zeros, each zero repeated according to
its multiplicity. The functional equation and the reality of the coefficients
preserve \(\mathscr Z\) under
\(\rho\mapsto1-\rho\), \(\rho\mapsto\overline\rho\) and
\(\rho\mapsto1-\overline\rho\). With our Fourier convention we define
\begin{equation}\label{eq:pdf2-rh-transformada-weil}
 \Phi_f(s):=\widehat f\!\left(i(s-\tfrac12)\right),
 \qquad
 \widetilde f(x):=\overline{f(-x)}.
\end{equation}
The complete explicit formula computed in
\eqref{eq:pdf2-rh-gram-difference-gw}, including the gamma and polar terms
that remove poles and trivial zeros, has the spectral form
\begin{equation}\label{eq:pdf2-rh-formula-espectral-weil}
 \mathscr I_{\rm GW}(f,g)
 =\sum_{\rho\in\mathscr Z}
   \Phi_f(\rho)\,
   \overline{\Phi_g(1-\overline\rho)}.
\end{equation}
The sum counts multiplicities. It converges absolutely: for
\(f,g\in C_c^\infty(\mathbb R)\), Paley--Wiener gives decay faster than any power in vertical
strips, whereas the number of zeros with \(|\operatorname{Im}\rho|\leq T\) is
\(O(T\log T)\).

\begin{proposition}[Reversible Weil Criterion with Full Quantifiers]
\label{prop:pdf2-rh-criterio-weil-desplegado}
Under the normalization
\eqref{eq:pdf2-rh-transformada-weil} and counting all nontrivial zeros
with their multiplicity,
\begin{equation}\label{eq:pdf2-rh-criterio-weil}
 \left[
  \mathscr I_{\rm GW}(f,f)\geq0
  \text{ for every }f\in C_c^\infty(\mathbb R;\mathbb C)
 \right]
 \Longleftrightarrow
 \left[
  \operatorname{Re}\rho=\frac12
  \text{ for every }\rho\in\mathscr Z
 \right].
\end{equation}
\end{proposition}

The converse implication is not obtained by interpolating an increasing number of
zeros: that operation would provide no uniform control over the tail. The
complete classical Weil criterion in Bombieri's formulation is used; its test
space and normalization are transferred here without abbreviation.\footnote{%
E.~Bombieri, \emph{Remarks on Weil's quadratic functional in the theory of
prime numbers, I}, Rendiconti Lincei, Matematica e Applicazioni, ser.~9,
vol.~11 (2000), pp.~183--233. The criterion established there states that the
quadratic functional of the explicit formula is positive semidefinite on all
\(C_c^\infty(0,\infty)\) if and only if the Riemann Hypothesis holds.}

\begin{proof}
We materially write out the change of coordinates identifying both
criteria. For \(f\in C_c^\infty(\mathbb R)\), let
\begin{equation}\label{eq:pdf2-rh-cambio-mellin}
 (\mathcal M f)(x):=x^{-1/2}f(\log x),
 \qquad x>0.
\end{equation}
The map
\[
 \mathcal M:C_c^\infty(\mathbb R)
 \longrightarrow C_c^\infty(0,\infty)
\]
is bijective, with inverse
\begin{equation}\label{eq:pdf2-rh-cambio-mellin-inverso}
 (\mathcal M^{-1}g)(u)=e^{u/2}g(e^u).
\end{equation}
If
\(\widehat g_{\rm Mel}(s):=\int_0^\infty g(x)x^{s-1}\,dx\), the change
\(x=e^u\) gives, with no residual factor,
\begin{equation}\label{eq:pdf2-rh-mellin-fourier}
 \begin{aligned}
 \widehat{\mathcal M f}_{\rm Mel}(s)
 &=\int_{\mathbb R}f(u)e^{(s-1/2)u}\,du\\
 &=\widehat f\!\left(i(s-\tfrac12)\right)
 =\Phi_f(s).
 \end{aligned}
\end{equation}
Moreover,
\begin{equation}\label{eq:pdf2-rh-mellin-involucion}
 \widehat{\overline{\mathcal M f}}_{\rm Mel}(1-\rho)
 =\overline{
   \widehat{\mathcal M f}_{\rm Mel}(1-\overline\rho)}
 =\overline{\Phi_f(1-\overline\rho)}.
\end{equation}
Therefore the Weil--Bombieri criterion functional is exactly
\begin{equation}\label{eq:pdf2-rh-bombieri-identificado}
 \sum_{\rho\in\mathscr Z}
 \widehat{\mathcal M f}_{\rm Mel}(\rho)\,
 \widehat{\overline{\mathcal M f}}_{\rm Mel}(1-\rho)
 =
 \sum_{\rho\in\mathscr Z}
 \Phi_f(\rho)\overline{\Phi_f(1-\overline\rho)}
 =\mathscr I_{\rm GW}(f,f).
\end{equation}
The class of tests is not changed, no multiplicity is lost, and the
set of zeros is not truncated, because \(\mathcal M\) is a bijection. The
cited classical criterion applied to
\eqref{eq:pdf2-rh-bombieri-identificado} gives precisely the equivalence
\eqref{eq:pdf2-rh-criterio-weil}. In the forward direction, if
\(\operatorname{Re}\rho=1/2\), the same identity reduces explicitly to
\[
 \mathscr I_{\rm GW}(f,f)
 =\sum_{\rho\in\mathscr Z}|\Phi_f(\rho)|^2\geq0.
\]
\end{proof}

The equality \(\bigcup_R\mathcal D_R=C_c^\infty(\mathbb R)\) is exact; each
function employed in the argument belongs to some window, and cofinal
naturality transports there the positivity already proved.

\begin{theorem}[Riemann--Weil Resolution from the Spectral--Polar Image]
\label{thm:pdf2-rh-resolucion}
Having fixed the complete structure
\(\mathfrak G_{\rm sp}\), the complete
Guinand--Weil form is positive on the cofinal separating class. Consequently,
every nontrivial zero \(\rho\) of the Riemann zeta function
satisfies
\[
 \boxed{\operatorname{Re}\rho=\frac12.}
\]
\end{theorem}

\begin{proof}
The prime--gamma--polar expansion identifies the difference between the two columns
with \(\mathscr I_{\rm GW}\) using a single regulator and removes it by
dominated convergence. The antispecular recurrence
\eqref{eq:pdf2-rh-recurrencia-residual} annihilates the output difference before
forming an inequality; therefore the isometry \(\Sigma_R\) gives
\eqref{eq:pdf2-rh-energia-source-first}. Comparing that identity with the
calculated expansion produces
\(E_R^*E_R=\mathscr I_{\rm GW}\). The local identity
\(A^*A+D^*D=I\) is iterated without erasing the terminal and gives
\eqref{eq:pdf2-rh-telescopia-finita}; the terminal is exactly
\(q^{2N}E_R^*E_R\) and vanishes only in the limit. Hence
\(\mathscr L_R^*\mathscr L_R=\mathscr I_{\rm GW}\succeq0\).
The action
\eqref{eq:pdf2-rh-xi-comun-explicita} then constructs the encoder, its
two moments, the projection, and the two publications, and
\eqref{eq:pdf2-rh-tres-publicaciones} identifies all realizations of the
same square. Naturality transports that identity to the entire cofinal
union. Finally,
\eqref{eq:pdf2-rh-criterio-weil} concludes the assertion.
\end{proof}

\section{Subsequent, non-governing local Redheffer control}

The following development verifies that local passivity can be obtained from
the carry without using the Weil form. It is a noncircularity control and not a
premise of \cref{thm:pdf2-rh-resolucion}. This truncated transfer is not
the complete colligation \(\Sigma_R\) of
\eqref{eq:pdf2-rh-coligacion-source-first}.

For \(N=9^m\), be \(U_{a,N}\) the odometro unitario
\[
 U_{a,N}(e_j\otimes f)=e_{j+1}\otimes f\ (j<N-1),
 \qquad U_{a,N}(e_{N-1}\otimes f)=e_0\otimes T_af.
\]
With \(P_N=|u_N\rangle\langle u_N|\otimes I\), \(Q_N=I-P_N\),
\(A_{a,N}=P_NU_{a,N}P_N\), and \(C_{a,N}=Q_NU_{a,N}P_N\),
\begin{equation}\label{eq:pdf2-rh-source-energia}
 A_{a,N}^*A_{a,N}+C_{a,N}^*C_{a,N}=P_N,
 \qquad
 C_{a,N}^*C_{a,N}
 =\frac{N-1}{N^2}(2I-T_a-T_a^*).
\end{equation}
For \(0<s<1\), the transfer
\begin{equation}\label{eq:pdf2-rh-source-redheffer}
 \mathscr R_{a,N}(s)=(sI-A_{a,N})(I-sA_{a,N})^{-1}
\end{equation}
satisfies
\begin{equation}\label{eq:pdf2-rh-source-defecto}
\begin{aligned}
 I-\mathscr R_{a,N}(s)^*\mathscr R_{a,N}(s)
 ={}&(1-s^2)(I-sA_{a,N}^*)^{-1}\\
 &\quad\cdot C_{a,N}^*C_{a,N}(I-sA_{a,N})^{-1}\succeq0.
\end{aligned}
\end{equation}
With
\[
 s_{p,N}=\frac{Np^{-1/2}}{1+(N-1)p^{-1/2}},
 \qquad s_\Gamma(a)=e^{-a/2},
\]
the sum over prime clocks and the direct gamma integral have norm at
most one. Explicitly,
\begin{equation}\label{eq:pdf2-rh-source-relojes}
 \mathscr R_{N,R}
 =\bigoplus_{\log p\leq R}
   \mathscr R_{\log p,N}(s_{p,N})
 \ \oplus\!
 \int_0^\infty{}^\oplus
   \mathscr R_{a,N}(s_\Gamma(a))\,da,
 \qquad \|\mathscr R_{N,R}\|\leq1.
\end{equation}
The finite chart of roles
\[
 \mathcal H_+=\operatorname{span}\{f_3,f_6,f_9\},\quad
 \mathcal H_-=\operatorname{span}\{f_4,f_8\},\quad
 C_\angle=|f_4\rangle\langle f_3|+|f_8\rangle\langle f_9|
\]
satisfies \(I-C_\angle^*C_\angle=P_6\). The two-way energy charts fix
\begin{equation}\label{eq:pdf2-rh-source-cartas}
\begin{aligned}
 \mathcal K_q&=K_2\otimes(\mathbb C^9)^{\otimes q},\\
 \mathsf B_{\pm,q}:\mathcal M_{\pm,q,R}
 &\longrightarrow
 \mathcal K_q\widehat\otimes\mathcal H_\pm
 \widehat\otimes\mathcal H_{\mathfrak C_R},\\
 \mathsf B_{\pm,q}^*\mathsf B_{\pm,q}&=\mathsf G_{\pm,q},
 \qquad
 \mathsf B_{\pm,q}^{\sharp}
 =\mathsf G_{\pm,q}^{-1}\mathsf B_{\pm,q}^*.
\end{aligned}
\end{equation}
The transfer
\begin{equation}\label{eq:pdf2-rh-source-transferencia}
 \mathscr C_{q,R}^{\rm src}
 =\mathsf B_{-,q}^{\sharp}
  (I\otimes C_\angle\otimes\mathscr R_{N,R})\mathsf B_{+,q}
\end{equation}
has defect
\begin{align}
 \Delta_{q,R}^{\rm src}
 ={}&I\otimes\bigl[P_6\otimes I
 +(P_3+P_9)\otimes
 (I-\mathscr R_{N,R}^*\mathscr R_{N,R})\bigr]\succeq0,
 \label{eq:pdf2-rh-source-delta}\\
 I-(\mathscr C_{q,R}^{\rm src})^\sharp\mathscr C_{q,R}^{\rm src}
 ={}&\mathsf B_{+,q}^{\sharp}\Delta_{q,R}^{\rm src}
 \mathsf B_{+,q}\succeq0.
 \label{eq:pdf2-rh-source-contractividad}
\end{align}
This control preserves route, clock, role, carry, and channels before forming
a difference. It is not identified with \(\mathcal K_R\) and does not replace the
explicit common publication of
\eqref{eq:pdf2-rh-doble-publicacion-material}; its local scope therefore does not
reopen the cofinal conclusion.

\section{Falsifiers of an ablated shadow}

\begin{falsifier}[Substitutions that change the starting object]
The preceding argument does not transfer to a shadow that eliminates one of the
thirteen coordinates of \eqref{eq:pdf2-rh-res-sp-completa}, constructs \(J_+\) and
\(J_-\) from different antecedents, replaces an orthogonal projection with
a non-self-adjoint idempotent, omits a channel, uses incompatible regulators,
chooses \(P_R\) after the sign, erases the terminal object before the limit, or reduces
the class to a noncofinal family.
\end{falsifier}

\begin{proof}
Each operation respectively violates the domain of
\(\mathfrak G_{\rm sp}\), the proper definition of
\eqref{eq:pdf2-rh-vmas-rango}, the orthogonal square, the functional identity,
the causal direction, finite telescoping, or the reversibility criterion. The
falsifier affects the replaced shadow; it does not dislodge the complete
structural image used in \cref{thm:pdf2-rh-resolucion}.
\end{proof}

\paragraph{Provenance and evidentiary force.}
The spectral--polar restriction, source--first colligation, its complement and the mirror sheet are \texttt{PRE-EXISTING\_AUTHORIAL\_ARCHITECTURE}. The coinductive vanishing of \(\Omega\), the finite identity with terminal and the action encoder \eqref{eq:pdf2-rh-xi-comun-explicita} are a \texttt{NEW\_FORMALIZATION} of that complete structure. The Guinand--Weil identification and cofinal closure are a \texttt{RECOVERED\_RESULT}. Equations \eqref{eq:pdf2-rh-coligacion-source-first}--\eqref{eq:pdf2-rh-tres-publicaciones} exhibit the object, its domains, its action, the terminal and the limit without introducing an abstract root of the form; their strength is \emph{proved with explicit starting structure}. \endgroup
 \part{Theorem-based synthesis}

\chapter{Arithmetic before Archimedean realization}
\section{VI. The discrete structure of the continuum}

\subsection{The continuum is a realization of compatible cylinders}

The inverse limit \(X_\infty\) preserves all compatible histories. Its half-closed cylinders determine a topology; refinement determines shape; the cocycle determines holonomy; orientation determines boundaries; and the inductive algebra determines readers. The five internal spectral, solenoidal, gauge, coherent and determinantal aspects emerge jointly from this construction. Separating them would destroy the object giving them their common naturality.

Dual projection is a pair of representations of the same algebra:

\[
\rho_{\rm arq}:\mathfrak H_{\rm cyl}\to\mathcal A_{\mathbb R},
\qquad
\rho_{\rm exc}:\mathfrak H_{\rm cyl}\to\mathcal A_{\rm exc}.
\]

The first contracts compatible cylinders towards the Archimedean reading. The second preserves boundary incidence towards Witt--Golay--Leech--\(V^\natural\)--Monster. \(\mathbb R\) and the exceptional branch are not two universes of origin: they are two representational quotients of a richer source.

Abelianization necessarily forgets the electronic commutator and part of the memory record. The real line can therefore be complete as a scalar continuum while remaining insufficient as a description of genealogy. Holonomic arithmetic precedes \(\mathbb R\) just as a history precedes one of its coordinates.

\subsection{Predefinition of compatible futures}

Let

\[
\mathcal S_n(x)=\{y\in X_{n+1}:r_{n+1,n}(y)=x
\text{ and }y\text{ is admissible}\}.
\]

The coinductive tree contains possible extensions. The Sturmian law guarantees a unique right-special word per length; TRIT types its curvature and TPK selects the continuation while preserving the complementary branch at the boundary. Predefinition does not mean that an external numeral has been stored. It means that every finite output is fixed by a cylinder and that the system of cylinders is compatible at every depth.

This makes precise the expression ‘choice of survival futures’. The future is not guessed from a measurement: it is selected within the already typed set of extensions preserving the relations of the state. Past, present and future are functions of the return operator, threshold and opening, not mere verbal labels.

\section{VII. Kernel theorem and consequences}

\subsection{Minimal statement}

We define the kernel object

\[
\boxed{
\mathfrak H_{\rm HMT}^{\rm alg}
=\left[\varinjlim_n
\left(
\mathbb R\langle e_x,J_x,T_u\rangle/
\mathcal R_{\rm HMT}
\right)\right]
\rtimes_{\Gamma_9}\mathbb N.}
\]

In the reversible sector, \(\mathbb N\) is replaced by \(\mathbb Z\). Its space of infinite states is

\[
X_\infty=\varprojlim_nX_n.
\]

The expression is algebraic. A bounded representation \(\rho\) determines its own completion \(\mathfrak H_{\rm HMT}^{\rho}
=\overline{\rho(\mathfrak H_{\rm HMT}^{\rm alg})}^{\|\cdot\|}\); the topology of the codomain therefore belongs to the reader and is not placed before the generator.

\textbf{Kernel theorem of conservative evolution of unity.} APP arithmetic, oriented by the fibres \(J_\tau^2=-\tau I\) and transported by the TPK cocycle, generates an associative noncommutative algebra whose refinement preserves unity and whose returns raise memory. The typed spectrum of representations of this algebra correlatively determines the correlated system \((\pi,\varphi,e,\alpha)\), the joint structure of the continuum, the electronic vacancy, action and gravity scales, geometric parameters, catalogue constants and boundary symmetries. Each realization preserves its domain, reader and degree.

The proof divides into six consecutive identities:

\begin{enumerate}
\item the APP equations prove preservation of residue and carry;
\item the tritic relations prove coexistence of the three curvatures;
\item the cocycle identity proves associativity of concatenation;
\item nonsplitting of \(C_{108}\) proves that phase return raises memory;
\item the inverse-limit/direct-limit pair proves preservation of unity under refinement and reading;
\item the universal property proves that every typed output is a representation of the same source, not a disguised conventional input.
\end{enumerate}

\subsection{Conceptual consequence: scalar reality is a faithful but incomplete shadow}

The Archimedean representation can be faithful to the determined value yet unfaithful to genealogy. Two histories can have the same real coordinate and different carry, height or memory. The discrete structure of the continuum resolves this multiplicity by preserving fibres before contracting them. The phrase ‘one point becomes two’ thus translates as follows: a scalar coordinate can admit two distinct holonomic lifts, and bidirectionality resides in the fibre, not in an arbitrary duplication of the point.

\subsection{Physical consequence: unity does not first preserve symmetries; it produces them}

What is primarily preserved is the possibility of reconstruction: sheet, residue, quotient, carry, orientation, boundary and memory. A symmetry appears when a representation has automorphisms preserving these relations. Exceptional symmetries, electromagnetic polarity and gravitational pentapolarity are thus different ways of stabilizing an already generated boundary. Symmetry conservation is a consequence of genealogical conservation.

\subsection{Informational consequence: infinite growth at zero internal cost}

Linear complexity \(p(m)=m+1\) allows unlimited nonperiodic information to be produced without exponential entropy growth. Refinement adds memory, but the integral update is reversible. Cost appears when a reader decides to erase the fibre that would allow the step to be reconstructed. This combination—aperiodicity, zero topological entropy, growing memory and preserved unity—is the operational mathematical definition of the HMT topological-temporal crystal.

\section{VIII. Independent contributions arising from retrospective reading}

\subsection{First contribution: inverse-limit--direct-limit duality}

The distinction between the inverse limit of states and the direct limit of observables had not been formulated at kernel level in the earlier developments. It explains why deepening and projection are conjugate operations without being trivial inverses; corrects the identification of global unity with an individual section; and provides the natural framework for predefinition.

\subsection{Second contribution: universal history algebra}

The presentation by generators and relations brings together the TRIT fibres, APP order, TPK transports and the memory-record cocycle. Its universal property proves that subsequent readers need no new generators if they preserve the relations.

\subsection{Third contribution: trigeometric exponential identity}

The common exponential of \(J_+,J_0,J_-\) shows that complex rotation, the nilpotent jet and hyperbolic dilation are three components of a single law. Euler's identity is its elliptic projection; \(\varphi\) is its hyperbolic scale; \(e\) is the exponential law; and \(\alpha\) is the seam gluing the jet to dodecaphasic holonomy.

\subsection{Fourth contribution: constants as a typed spectrum of representations}

The notion of scalar character is extended to a spectrum of representations and functionals. This brings constants, particles and fields together without confusing their codomains. The same definition locates Catalan as a moment, Feigenbaum as nonlinear rescaling, Boltzmann as transducer, Barbero--Immirzi as incidence covector, Planck as action section and gravity as contragredient character.

Functional recovery makes an additional relation explicit: the constitutive incidence response \(90/120\), with its regularity at the origin, determines the complete Catalan moment; its inversion recovers the labelled angular channels. Transport composition is then expressed through the law \(r\odot t=f(\psi(r)\psi(t))\), preserving product, orientation and normalized trace. On that same velocity section, inscription and the radial norm first determine the length \(L\), then the products \(R_X\Lambda_X=L^2I\) and \(H_X\Lambda_X=\hbar_{\rm ret}cI\), and finally the common coefficient \(G=c^3L^2/\hbar_{\rm ret}\). The proofs of \cref{cr28:sec:restitucion-constitutiva,g26:sec:rigidez-radial-es} preserve the memory and constitutive hypotheses of each realization.

The Feigenbaum constant acquires a precise structural definition: it is the asymptotic separation invariant of the parameters selected by the critical returns of the uniform dodecaphasic reader. The construction preserves its APP--TRIT--TPK origin, the joint continuum and the first-root order. Identifying the selected centres with the analytic centres makes the scaling theorem applicable to that same sequence; the proofs in \ref{hmtfg:escalado:15}--\ref{hmtfg:escalado:17} and \eqref{hmtfg:articulacion:limite} establish the parametric result. The Jacobi representation recovers the profile from its twelve-node measure, while constitutive inversion recovers the oriented channels and the moment whose endpoint is Catalan's constant. These operations exhibit relations between readers of the same antecedent, preserving the domain and information of each.

\subsection{Fifth contribution: symmetry as reconstruction stabilizer}

The intuition that symmetries result from conservation is converted into a definition of the stabilizer of boundary and relations. This formulation links dual projection, the exceptional corridor and physical representations without postulating symmetry as the source.

\subsection{Sixth contribution: triadic irreducibility of the kernel}

The author's image of the Borromean link admits an exact formulation through three forgetful quotients. If the two APP sheets are identified, both the potential \(\Delta_{\Sigma\Pi}\) and the commutator \([P_\Sigma,P_\Pi]\) disappear. If the three TRIT regimes are identified, the sign spectrum disappears

\[
J_+^2=-I,\qquad J_0^2=0,\qquad J_-^2=I.
\]

If the TPK cocycle is trivialized, the relation \(d_M(\Gamma_9x)-d_M(x)=1\) becomes zero and phase return becomes indistinguishable from state return. Each quotient destroys an invariant the other two factors cannot reconstruct. APP, TRIT and TPK are therefore three functionally irreducible generators of the kernel object: APP supplies discrepancy and sheet order; TRIT supplies the local curvature signature; TPK supplies composition, holonomy and memory.

This irreducibility explains why the kernel cannot be reduced to an arithmetic table, a quadratic algebra or phase dynamics. The structure appears only in their composition.

\begingroup\fontsize{10.2}{11.7}\selectfont
\setlength{\parskip}{2pt}\setlength{\abovedisplayskip}{5pt}\setlength{\belowdisplayskip}{5pt}
\fontsize{9.5}{10.7}\selectfont
\setlength{\parskip}{1.5pt}\setlength{\abovedisplayskip}{2pt}\setlength{\belowdisplayskip}{2pt}
\subsection{Seventh contribution: the electron contains the Witt null pair}

Retrospective comparison of the electronic block with the three TRIT regimes reveals a local closure that had remained implicit. Starting from

\[
S^2=I,\qquad J^2=-I,\qquad SJ=-JS
\]

we form

\[
n_\pm=\frac{S\pm J}{2}.
\]

Then

\[
n_\pm^2=0,\qquad n_+n_-+n_-n_+=I.
\]

The corner \({\rm Cl}_{1,1}\cong M_2(\mathbb R)\) therefore contains the three local classes: \(J\) is elliptic, \(S\) is hyperbolic and \(n_\pm\) are parabolic. The electronic vacancy is simultaneously a minimal carry defect in APP and a null direction of the split algebra. TRIT does not introduce the threshold as an external correction to the electron: it recognizes its type and preserves it under TPK transport. The designation Witt refers here to the local null decomposition; the Paley--Witt incidence structure is obtained subsequently as a combinatorial boundary representation of the same genealogy.

\subsection{Eighth contribution: Weyl relations of the two clocks}

The phrase ‘the helix contains two clocks’ admits an exact operator presentation. Let \(T\) be the step, \(V_9\) the nonadic-phase character, \(W_\delta\) the irrational-capacity character, and \(\delta=\log_{10}(10/9)\). If \(\zeta_9\) and \(\omega_\delta\) denote their internal characters, then

\[
V_9T=\zeta_9TV_9,\qquad
W_\delta T=\omega_\delta TW_\delta.
\]

For \(\Gamma_9=T^9\),

\[
V_9\Gamma_9=\Gamma_9V_9,\qquad
W_\delta\Gamma_9=\omega_\delta^9\Gamma_9W_\delta.
\]

The phase clock closes and the capacity clock continues. If \(D_M\) measures helical height,

\[
[D_M,\Gamma_9]=\Gamma_9.
\]

This final identity is the shortest algebraic form of the difference between phase return and state return. It also locates noncommutativity precisely: basal translation belongs to an Abelian group, but its crossed product with the phase and capacity observables satisfies noncommutative Weyl relations. The HMT aperiodic crystal is thus an ordered suspension of a finite clock by an irrational character, lifted by a memory operator.

In the subsequent elliptic representation, \(\zeta_9\) and \(\omega_\delta\) are recognized as \(\operatorname{Exp}_+(2\pi{}J_+/9)\) and \(\operatorname{Exp}_+(2\pi{}\delta J_+)\). The HMT character precedes that complex notation.

\subsection{Ninth contribution: a polynomial equation for the three geometries}

The three local generators can be combined in

\[
\mathbb J=p_+J_++p_0J_0+p_-J_-.
\]

The identity

\[
\boxed{\mathbb J^6=\mathbb J^2}
\]

contains exactly the spectrum of squares \(\{-I,0,I\}\). If \(K=\mathbb J^2\), the regimes are recovered from the operator itself:

\[
p_+=\frac{K^2-K}{2},\qquad
p_0=I-K^2,\qquad
p_-=\frac{K^2+K}{2}.
\]

This formula dispels the appearance that TRIT is merely a ternary label. The regime is an intrinsic spectral decomposition. The corresponding exponential identity

\[
\operatorname{Exp}_\star(t\mathbb J)=
p_+(\cos t+\mathbb J\sin t)+
p_0(I+t\mathbb J)+
p_-(\cosh t+\mathbb J\sinh t)
\]

contains as faces Euler's elliptic closure, the nilpotent threshold and hyperbolic self-scaling. Its composition with \(\Gamma_9\), together with \([D_M,\Gamma_9]=\Gamma_9\), constitutes the minimal operator identity from which conservative evolution of the dimensional unit can be expounded.
  

\par\endgroup
\appendix
\label{app:desarrollos-10000-rev9}
% Copia mecánica íntegra de apendices.tex:1--761.
% Residencia material del ensamblaje sucesor de 102 capítulos.
\part{Proof appendices and classification tables: algebraic, arithmetic and coinductive certificates}

\chapter{Ordered affine calculation of the APP leaves}

\section{Affine semidirect product and route composition}

Let \(U_{\APP}\) be the radial module prior to its Archimedean evaluation.
The correct envelope is the monoid
\[
 \AffEnd(U_{\APP})
 =\operatorname{End}(U_{\APP})\ltimes U_{\APP},
\]
because the threshold or folding leaves may contain noninvertible
endomorphisms. The action of \((\Lambda,C)\) on \(u\in U_{\APP}\) is
\[
 (\Lambda,C)\cdot u=\Lambda u+C.
\]
The composition is calculated without commuting factors:
\begin{align}
 (\Lambda_2,C_2)(\Lambda_1,C_1)\cdot u
 &=\Lambda_2(\Lambda_1u+C_1)+C_2 \notag\\
 &=(\Lambda_2\Lambda_1)u+(C_2+\Lambda_2C_1).
 \label{espiral:eq:ap-comp-afin}
\end{align}
Therefore,
\[
 (\Lambda_2,C_2)(\Lambda_1,C_1)
 =(\Lambda_2\Lambda_1,C_2+\Lambda_2C_1).
\]
The second coordinate is not an ordinary sum of shifts: it is a
cocycle twisted by the accumulated linear action.

\begin{proposition}[Prefix Formula]
For a route with \(N\) factors \(h_j=(\lambda_j,C_j)\), ordered from
\(j=0\) through \(j=N-1\), the prefix holonomy is
\[
 H_N=h_{N-1}\cdots h_1h_0=(\Lambda_N,\mathcal C_N),
\]
where
\[
 \Lambda_N=\prod_{j=0}^{N-1}\lambda_j
\]
and
\[
 \mathcal C_N
 =C_{N-1}+\lambda_{N-1}C_{N-2}
 +\lambda_{N-1}\lambda_{N-2}C_{N-3}
 +\cdots+
 \lambda_{N-1}\cdots\lambda_1C_0.
\]
\end{proposition}
\begin{proof}
For \(N=1\) the formula is immediate. If it holds for \(N\), then
\[
 h_NH_N
 =(\lambda_N\Lambda_N,C_N+\lambda_N\mathcal C_N),
\]
which yields exactly the expression of length \(N+1\). The induction
preserves the order of the route and does not identify paths with the same
endpoints.
\end{proof}

\section{Difference sum--product as curvature of order}

Let \(A_a=(I,a)\) and \(M_\lambda=(\lambda,0)\) be given. In the invertible sector,
\[
 A_aM_\lambda A_a^{-1}M_\lambda^{-1}
 =A_{a(1-\lambda)}.
\]
The direct proof for \(u\) is
\[
 u\xmapsto{M_\lambda^{-1}}\lambda^{-1}u
 \xmapsto{A_a^{-1}}\lambda^{-1}u-a
 \xmapsto{M_\lambda}u-\lambda a
 \xmapsto{A_a}u+a(1-\lambda).
\]
Thus, the distinction between accumulation and folding is published as a
translation. When \(\lambda=1\), the commutator vanishes; when
\(\lambda\neq1\), the route retains a holonomy that cannot be reconstructed
from the endpoint alone.

\begin{example}[Two routes with the same visible data]
Fix \(u_0=0\), \(a=1\), \(\lambda=2\). The routes
\[
 A_1M_2:u_0\longmapsto1,
 \qquad
 M_2A_{1/2}:u_0\longmapsto1
\]
publish the same final value. However, their ledgers contain distinct factors,
shifts, and prefixes. Pointwise evaluation identifies the
routes; enriched affine holonomy distinguishes them.
\end{example}

\section{Exhaustive table of the local chart of 2 orientations}

The following table does not posit that every global radial degree arises from a
single pair. It records the elementary chart obtained by composing two TRIT
orientations. The degree of a general route is the integer cocycle of
\cref{espiral:eq:cociclos-radiales}; this chart provides its nine local
generators.

\begin{longtable}{rrrrrrp{.27\textwidth}}
\caption{The \(9\) oriented pairs of the local TRIT chart.}
\label{espiral:tab:ap-nueve-pares}\\
\toprule
\(\tau_+\)&\(\tau_-\)&\(p\)&\(a\)&\(r_3\)&\(c_3\)&Lectura local\\
\midrule
\endfirsthead
\toprule
\(\tau_+\)&\(\tau_-\)&\(p\)&\(a\)&\(r_3\)&\(c_3\)&Lectura local\\
\midrule
\endhead
\( 1\)&\( 1\)&\( 2\)&\( 0\)&\(-1\)&\( 1\)&double radial opening with positive carry\\
\( 1\)&\( 0\)&\( 1\)&\( 1\)&\( 1\)&\( 0\)&opening and threshold\\
\( 1\)&\(-1\)&\( 0\)&\( 2\)&\( 0\)&\( 0\)&positive orientation equilibrium\\
\( 0\)&\( 1\)&\( 1\)&\(-1\)&\( 1\)&\( 0\)&threshold and aperture\\
\( 0\)&\( 0\)&\( 0\)&\( 0\)&\( 0\)&\( 0\)&umbral doble\\
\( 0\)&\(-1\)&\(-1\)&\( 1\)&\(-1\)&\( 0\)&threshold and contraction\\
\(-1\)&\( 1\)&\( 0\)&\(-2\)&\( 0\)&\( 0\)&negative orientation equilibrium\\
\(-1\)&\( 0\)&\(-1\)&\(-1\)&\(-1\)&\( 0\)&contraction and threshold\\
\(-1\)&\(-1\)&\(-2\)&\( 0\)&\( 1\)&\(-1\)&double contraction with negative carry\\
\bottomrule
\end{longtable}

Here \(p=r_3+3c_3\), with balanced residue
\(r_3\in\{-1,0,1\}\). The carry distinguishes \(p=2\) from its residue \(-1\) and
\(p=-2\) from its residue \(1\). The pair
\((p,a)\) likewise distinguishes the three local genealogies of
degree \(0\).

\chapter{Nonadic Certificate and Geometry of Memory}

\section{Typed distinction between decimal catalog, ternary words and ambient fiber}

\begin{longtable}{p{.18\textwidth}p{.18\textwidth}p{.52\textwidth}}
\caption{Strata of the certificate of the joint nonadic connection.}
\label{espiral:tab:certificado-nonadico}\\
\toprule
Object&Cardinality&Mathematical function\\
\midrule
\endfirsthead
\toprule
Object&Cardinality&Mathematical function\\
\midrule
\endhead
Catalog of source triples&\(104\,976\)&Finite domain prior to the
joint selection; it is not a decimal table.\\
Decimal emissions \(U_6\)&\(468\)&Compatible emissions of length \(6\)
in the decimal publication.\\
Visible words \(w_6\)&\(243\)&Published ternary words after typing; they are not the complete environment.\\
Ambient space \(\F_3^6\)&\(729\)&Space of all ternary words of
length \(6\); it contains the \(243\) visible words.\\
Certified levels&\(396\)&Finite prolongation up to the declared boundaries.\\
Trits recorridos&\(2\,376\)&Longitud total \(396\cdot6\).\\
Vueltas completas&\(43\)&Nonadic cycles in the certified segment.\\
Pairs \((K,K+9)\)&\(387\)&Pairs of same-phase by channel; the ledger may differ.\\
\bottomrule
\end{longtable}

The prolongation chain is preserved in the order
\[
 w_6\longrightarrow w_{12}\longrightarrow w_{18}
 \longrightarrow w_{24}\longrightarrow w_{30}
 \longrightarrow R_{36}\longrightarrow G_9.
\]
None of these objects can be replaced by the subsequent modal quotient.
The holonomy \(\Gamma_9\) acts on the enriched dynamics; the
reduced-memory publication preserves only the information specified
by its codomain.

\section{Phase return, memory advance, and complete state}

With
\[
 \rho_9(N)=1+((N-1)\bmod9),
 \qquad
 q_9(N)=\left\lfloor\frac{N-1}{9}\right\rfloor,
\]
we have
\[
 \rho_9(N+9)=\rho_9(N),
 \qquad q_9(N+9)=q_9(N)+1.
\]
Therefore,
\[
 H_9(N+9)=H_9(N)+(0,0,1).
\]
The equality of the first phase coordinate does not imply equality of the state:
\[
 g(K+9)=g(K),
 \qquad B_{K+9}\neq B_K
\]
in general. The helical formula is precisely the minimal representation
of this distinction.

\section{Non-Split Dodecaphase Extension}

The sequence
\[
 0\longrightarrow C_{12}\longrightarrow C_{108}
 \longrightarrow C_9\longrightarrow0
\]
couples the dodecaphasic calendar to the nonadic return. If a
global homomorphic section \(s:C_9\to C_{108}\) existed, the generator of \(C_9\)
would have to lift to an element of order dividing \(9\) that projected onto the
generator. The active temporal lift, however, preserves the accumulated
dodecaphasic displacement and does not identify the state after one phase turn.
The seam is recorded as carry; geometrically, it manifests as
axial translation of the helix.

\section{5 regions of \texorpdfstring{\(\pi\)}{pi} and non-faithfulness of evaluation}

The microfiber of \(\pi\) contains five oriented regions with
total multiplicity
\[
 432+4\cdot144=1008.
\]
The regions publish the same real value, but retain distinct genealogies.
This result provides the internal precedent for spiral classification: a
visible curve does not by itself determine the region, the sheet, the carry,
or the memory that produced it.

\chapter{Power--logarithmic algebra and 5 canonical publications}

\section[Power--logarithmic and inverse coordinate]{The coordinate \(L_p\) and its inverse}

For \(x>0\), define
\[
 L_p(x)=
 \begin{cases}
 (x^p-1)/p,&p\neq0,\\
 \log x,&p=0.
 \end{cases}
\]
Its inverse is
\[
 \exp_p(u)=
 \begin{cases}
 (1+pu)^{1/p},&p\neq0,\ 1+pu>0,\\
 \e^u,&p=0.
 \end{cases}
\]
The derivative
\[
 \frac{\dd}{\dd x}L_p(x)=x^{p-1}>0
\]
proves that \(L_p\) is strictly increasing on its domain and is therefore
invertible.

\begin{proposition}[Sum--Product Law]
For \(x,y>0\),
\[
 L_p(xy)=L_p(x)+L_p(y)+pL_p(x)L_p(y).
\]
\end{proposition}
\begin{proof}
Si \(p\neq0\),
\begin{align*}
 L_p(x)+L_p(y)+pL_p(x)L_p(y)
 &=\frac{x^p-1}{p}+\frac{y^p-1}{p}
   +\frac{(x^p-1)(y^p-1)}p\\
 &=\frac{x^py^p-1}{p}=L_p(xy).
\end{align*}
For \(p=0\), the equality reduces to
\(\log(xy)=\log x+\log y\).
\end{proof}

\begin{corollary}[Deformed group in the positive domain]
The operation
\[
 u\boxplus_pv=u+v+puv
\]
satisfies
\[
 1+p(u\boxplus_pv)=(1+pu)(1+pv).
\]
Its neutral element is \(0\), and the inverse of \(u\) is
\[
 \ominus_pu=-\frac{u}{1+pu},
\]
si \(1+pu\neq0\).
\end{corollary}

\section{Derivation of the 5 canonical curves}

Let \(m\) be the published depth, and suppose
\[
 L_p(r/r_0)=\beta m,
 \qquad
 \theta=\theta_0+\omega m,
 \qquad\omega\neq0.
\]
Eliminating \(m=(\theta-\theta_0)/\omega\),
\[
 r(\theta)=r_0
 \left[1+p\frac{\beta}{\omega}
 (\theta-\theta_0)\right]^{1/p},
 \qquad p\neq0,
\]
and
\[
 r(\theta)=r_0
 \exp\!\left[\frac{\beta}{\omega}
 (\theta-\theta_0)\right]
\]
for \(p=0\). Positive changes of origin and rescalings yield the
conventional forms:

\begin{longtable}{rllp{.31\textwidth}}
\caption{Canonical publications in the translational sector.}
\label{espiral:tab:cinco-publicaciones-detalle}\\
\toprule
\(p\)&Normal equation&Conventional name&Domain and feature\\
\midrule
\endfirsthead
\toprule
\(p\)&Normal equation&Conventional name&Domain and feature\\
\midrule
\endhead
\(2\)&\(r^2=a^2\theta\)&Fermat spiral&\(\theta\ge0\); linear increment of the quadratic
radial area.\\
\(1\)&\(r=a\theta\)&Archimedes&\(\theta\ge0\); constant radial
increment per angular unit.\\
\(0\)&\(r=r_0\e^{b\theta}\)&logarithmic&\(\theta\in\R\); constant dilation
per angular unit.\\
\(-1\)&\(r=a/\theta\)&hyperbolic&\(\theta>0\); the product \(r\theta\) is
constant.\\
\(-2\)&\(r^2=a^2/\theta\)&lituus spiral&\(\theta>0\); product
\(r^2\theta\) is constant.\\
\bottomrule
\end{longtable}

\section{Sheet inversion}

Para \(p\neq0\),
\[
 L_{-p}(x^{-1})
 =\frac{x^p-1}{-p}=-L_p(x).
\]
The identity also holds for \(p=0\). Consequently, radial inversion
interchanges \(p\) with \(-p\) and reverses the sign of the
linearized coordinate. This relation pairs the Fermat spiral with the lituus and the Archimedean
spiral with the hyperbolic spiral as conjugate orientations of the reader, without identifying them
as identical curves.

\section{General affine sector}

If the radial coordinate satisfies
\[
 u_{n+1}=\Lambda u_n+C,
\]
the exact internal iteration is
\[
 u_n=\Lambda^n u_0+\sum_{k=0}^{n-1}\Lambda^kC.
\]
If \(I-\Lambda\) is invertible,
\(u_*=(I-\Lambda)^{-1}C\) and
\(u_n-u_*=\Lambda^n(u_0-u_*)\). Only in a positive scalar realization,
\(\chi\Lambda=\lambda\chi\) with \(\lambda>0\), does one define
\(u_*=C/(1-\lambda)\). With \(\theta_n=\theta_0+n\omega\),
\(\omega\neq0\), the interpolation published in that realization is
\[
 r(\theta)=r_0\exp_p\!\left[
 u_*+\lambda^{(\theta-\theta_0)/\omega}(u_0-u_*)
 \right].
\]
The complete taxonomy is therefore not a list of five names. It contains at
least the degree \(p\), the orientation \(a\), the affine class
\([\Lambda,C]\), the carry cocycle, the boundary, the angular increment, and
the memory ledger.

\chapter{Naturality, equivariance and non-faithfulness}

\section{The prefix reader}

For a truncated history
\[
 \widehat\gamma_N=\widehat e_N\cdots\widehat e_1,
\]
The enriched reader records
\[
 \mathcal P_{\rad,N}^{\enr}(\widehat\gamma_N)
 =
 \left(
 (H_0,\Led_0);
 \bigl(p(\widehat e_j),a(\widehat e_j),P_j,A_j,
 h_{e_j},H_j,\Led_j\bigr)_{j=1}^{N}
 \right).
\]
The term \(H_j=h_{e_j}\cdots h_{e_1}\) preserves the prefixes. The ledger
\(\Led_j\) preserves, according to the route, sheet, residue, quotient, carry,
boundary, orientation, cylinder, signature, survival, and memory. The
real-valued coordinates are obtained subsequently:
\[
 \mathcal P_{\arq}^{\mathrm{esp}}
 =\ev_{\R^2}\circ\mathcal P_{\rad}^{\enr}.
\]

\section{Complete proof of naturality under truncation}

Let \(\pi_{N,M}\) denote deletion of the suffix of length \(N-M\), with \(M<N\).
Let \(\operatorname{pr}_{N,M}\) denote the projection that preserves the initial record
and the first \(M\) edge records. Then
\[
 \operatorname{pr}_{N,M}\,
 \mathcal P_{\rad,N}^{\enr}(\widehat\gamma_N)
 =
 \left(
 (H_0,\Led_0);
 \bigl(p(\widehat e_j),a(\widehat e_j),P_j,A_j,
 h_{e_j},H_j,\Led_j\bigr)_{j=1}^{M}
 \right).
\]
By definition of truncation,
\[
 \pi_{N,M}(\widehat\gamma_N)=\widehat e_M\cdots\widehat e_1,
\]
and upon reapplying the reader the same sequence is obtained. Therefore,
\[
 \operatorname{pr}_{N,M}\circ
 \mathcal P_{\rad,N}^{\enr}
 =
 \mathcal P_{\rad,M}^{\enr}\circ\pi_{N,M}.
\]
The proof does not use the value of the curve. It uses prefix identity,
which is the information that the reduced seven-coordinate reader would have
lost.

\section{Non-trivial refinements}

For a refinement \(r:\gamma\to\gamma'\), naturality requires a
map \(\widehat r\) on ledgers and the local equality
\[
 h_e
 =
 h_{e'_k}\cdots h_{e'_1}
\]
for each edge \(e\) replaced by the refined chain
\((e'_1,\ldots,e'_k)\). If, in addition, the carry, boundary, and memory of the
refined chain compose the records of the original step, then
\[
 \widehat r\circ\mathcal P_{\rad}^{\enr}
 =
 \mathcal P_{\rad}^{\enr}\circ r.
\]
This condition is a verifiable commutative square; it must not be assumed for an
arbitrary refinement without exhibiting the factors.

\section{Equivariance, not invariance}

The action of \(\Gamma_9\) returns the phase and prolongs the state. In the
codomain of the reader,
\(\widehat\Gamma_9^{\rad}\) acts by shifting the records and updating the
last prefix. The correct identity is
\[
 \mathcal P_{\rad}^{\enr}(\Gamma_9\gamma)
 =
 \widehat\Gamma_9^{\rad}
 \mathcal P_{\rad}^{\enr}(\gamma).
\]
Requiring invariance would erase precisely the memory that the connection
preserves. In the reversible sector, route inversion satisfies
\[
 H_{C_{\mathrm{rev}}\gamma}=H_\gamma^{-1}
\]
when all factors belong to the invertible affine subgroup. This
typed assertion does not automatically extend to a global involution
on an arbitrary subsequent quotient.

\section{3 proofs of non-faithfulness}

\begin{proposition}[Intrinsic circle and projected circle]
There exist two histories with the same visible curve.
\[
 c(\theta)=(\cos\theta,\sin\theta)
\]
and different ledgers.
\end{proposition}
\begin{proof}
The first history has constant axial memory \(m=0\). The second is
the helix
\[
 h(\theta)=(\cos\theta,\sin\theta,\theta/2\pi),
\]
whose ledger records \(m(\theta)=\theta/2\pi\). Both satisfy
\(\pi_{xy}\circ h=c\). The projection is therefore non-injective.
\end{proof}

\begin{proposition}[Same endpoint, different affine order]
If \(a(1-\lambda)\neq0\), the paths \(A_aM_\lambda\) and
\(M_\lambda A_a\) are not equivalent under the enriched reader, even though
they may be calibrated to publish the same endpoint.
\end{proposition}
\begin{proof}
Its commutator is \(A_{a(1-\lambda)}\neq I\). The commutator datum
belongs to the prefix ledger and disappears when only the endpoint is retained.
\end{proof}

\begin{proposition}[Same curve, different genealogical region]
Equality of the published real value does not identify the five regions of
the microfiber of \(\pi\); consequently, it is likewise insufficient to identify
the spiral class when the reader preserves region, boundary, or carry.
\end{proposition}

\chapter{Stratified Taxonomy of Spiral Correspondences}

\section{Classification by mathematical level}

\begin{longtable}{p{.20\textwidth}p{.30\textwidth}p{.39\textwidth}}
\caption{Levels that must not collapse into a single notion of curvature.}
\label{espiral:tab:niveles-curvatura}\\
\toprule
Level&Invariant&Meaning\\
\midrule
\endfirsthead
\toprule
Level&Invariant&Meaning\\
\midrule
\endhead
APP&\(F(S,P)=1,\ F(P,S)=-1\)&Discrete curvature of the order of accumulation
and folding.\\
TRIT&\(\tau\in\{-1,0,1\}\), \(J_\tau^2=-\tau I\)&Local elliptic,
parabolic, or hyperbolic regime of the transport.\\
TPK&\(\Gamma_9\), carry, boundary, route, and ledger&Phase holonomy and memory on the integral holonomic state.\\
Radial level&\(p,a,[\Lambda,C]\)&Characters that determine the radial law before recognizing a curve.\\
Archimedean&\(r(\theta)\)&Curve published in \(\R^2\); it may forget
genealogy.\\
Frenet level&\(\kappa_{\mathrm{F}},\mathcal T_{\mathrm{F}}\)&Differential curvature and torsion
of the visible curve or suspension.\\
Dimensional level&fields, scales and representations&Posterior physical realization,
with its own premises.\\
\bottomrule
\end{longtable}

\section{Classification by rotation, radius and memory}

\begin{longtable}{ccccl}
\caption{Complete kinematic taxonomy of the minimal degenerate cases.}
\label{espiral:tab:taxonomia-cinematica-ap}\\
\toprule
\(\Delta\theta\)&\(\Delta u_{\rad}\)&\(\Delta m\)&Salida visible&
Genealogical Information\\
\midrule
\endfirsthead
\toprule
\(\Delta\theta\)&\(\Delta u_{\rad}\)&\(\Delta m\)&Salida visible&
Genealogical Information\\
\midrule
\endhead
\(0\)&\(0\)&\(0\)&point&published steady state\\
\(0\)&\(\neq0\)&\(0\)&recta radial&scale change\\
\(0\)&\(0\)&\(\neq0\)&fibra axial&non-transverse displacement memory\\
\(0\)&\(\neq0\)&\(\neq0\)&radial sheet suspended&scale and memory\\
\(\neq0\)&\(0\)&\(0\)&circle&intrinsic angular closure\\
\(\neq0\)&\(0\)&\(\neq0\)&cylindrical helix&closed phase, high memory\\
\(\neq0\)&\(\neq0\)&\(0\)&espiral plana&radial character visible\\
\(\neq0\)&\(\neq0\)&\(\neq0\)&radial helical suspension&turn, radius and
depth preserved\\
\bottomrule
\end{longtable}

\section{Affine holonomy classification}

The following table belongs to a real scalar realization of the
linear part, \(\chi\Lambda=\lambda\chi\); for ease of reading, we write
\(\Lambda\) in the conditions column. It does not classify abstract
endomorphisms by numerical inequalities.

\begin{longtable}{p{.18\textwidth}p{.22\textwidth}p{.48\textwidth}}
\caption{Types of holonomy over the radial coordinate.}
\label{espiral:tab:tipos-holonomia}\\
\toprule
Condition&Orbit of \(u\)&Interpretation\\
\midrule
\endfirsthead
\toprule
Condition&Orbit of \(u\)&Interpretation\\
\midrule
\endhead
\(\Lambda=1,\ C=0\)&\(u_n=u_0\)&constant radius; the memory may be
zero or hidden by projection.\\
\(\Lambda=1,\ C\neq0\)&\(u_n=u_0+nC\)&translational sector that publishes the
family \(L_p\).\\
\(0<\Lambda<1\)&\(u_n\to u_*\)&contraction towards a fixed section.\\
\(\Lambda>1\)&\(|u_n-u_*|\) grows&expansion from a fixed section.\\
\(\Lambda<0\)&orientation alternation&requires control of the domain of
\(\exp_p\) at each step.\\
\(\Lambda=0\)&\(u_{n+1}=C\)&noninvertible folding; it belongs to the
monoid, not to the affine group.\\
\bottomrule
\end{longtable}

\section{Genealogical Equivalence Criterion}

Two routes are not identified merely because they publish the same curve. A
typed isomorphism preserving the following term by term is required:
\[
 \bigl(
 (x_0,H_0,\Led_0);
 (\widehat e_j,p(\widehat e_j),a(\widehat e_j),P_j,A_j,
 h_{e_j},H_j,\Led_j)_{j=1}^{N}
 \bigr)
\]
and that commutes with the admissible truncations and refinements. The resulting
classification is finer than the conventional taxonomy because it incorporates the
history preceding the emergence of \(\R^2\).

\chapter{Subsequent realizations and type checks}

\section{Circular electromagnetic mode}

Let
\[
 \zeta=\frac{z-c_{\mathrm{vac}}t}{\lambda},
\]
and
\[
 \mathbf E_\sigma
 =E_0(\mathbf e_1\cos2\pi\zeta
 +\sigma\mathbf e_2\sin2\pi\zeta),
 \qquad
 \mathbf B_\sigma
 =c_{\mathrm{vac}}^{-1}\widehat{\mathbf k}\times\mathbf E_\sigma.
\]
Then
\[
 \mathbf E_\sigma\perp\mathbf B_\sigma\perp\widehat{\mathbf k},
 \quad |E|=c_{\mathrm{vac}}|B|,
 \quad\omega=c_{\mathrm{vac}}k.
\]
At fixed time, the normalized endpoint of the field in the phase--position
bundle describes
\[
 (\cos2\pi\zeta,\sin2\pi\zeta,\zeta).
\]
This statement does not assert that every spatial field line is a helix:
it describes the trajectory of the tip of the polarization vector together with the
marked event.

\section{Loxodromic compatibility}

For
\[
 \Phi(r,\varphi,z,t)
 =\ell\varphi+\nu\log(r/r_0)+k_zz-\omega t,
\]
the transformation
\[
 (r,\varphi)\longmapsto
 ((2+\sqrt3)r,\varphi+\pi/2)
\]
preserves the phase modulo \(2\pi\) exactly when
\[
 \nu\log(2+\sqrt3)+\ell\frac{\pi}{2}\in2\pi\Z.
\]
Dilation is interpreted as a change of scale between solutions; not
as spontaneous energy production.

\section{Electronic centre and subsequent reader}

The electron occupies the unique fixed point \((5,5)\) of the material atlas. This
centrality is not derived anew in the present work. The prior operator
compatibility includes
\[
 C^{(3)}M_{\ph}=-\frac12M_{\ph},
 \qquad
 1-\left(\frac12\right)^2=\frac34=\frac{90}{120},
\]
and
\[
 \|[P_{\Sigma,3},P_{\Pi,3}]\|=\frac{\sqrt3}{4}.
\]
The theory of spiral correspondences adds a route reader; it does not
replace the General Mass Law or use the electron to select
\(p\), \(\Lambda\), or \(C\) retrospectively.

\section{Closed matrix of negative controls}

\begin{longtable}{p{.38\textwidth}p{.50\textwidth}}
\caption{Falsifiers and type errors excluded by the construction.}
\label{espiral:tab:falsadores-ap}\\
\toprule
Failure&Consequence\\
\midrule
\endfirsthead
\toprule
Failure&Consequence\\
\midrule
\endhead
The classical curve selects \(p,\Lambda,C\).&Circularity: the reader ceases
to be constructive--generative.\\
Only the final holonomy is preserved.&Loss of prefixes, carry, and boundary;
naturality is not established.\\
Invariance under \(\Gamma_9\) is required.&The memory advance is erased; the
correct relation is equivariance.\\
\(U_6\), \(w_6\), and \(\F_3^6\) are conflated.&Collapse of objects with
cardinalities \(468\), \(243\), and \(729\).\\
The term phase return is identified with state return.&Contradiction with
\(B_{K+9}\neq B_K\) in general.\\
\(\tau\), \(p\), and \(\kappa_{\mathrm{F}}\) are identified.&Regime collapse,
radial character and visible curvature.\\
A general helix is called any variable radius.&Differential statement
unverified Lancret condition.\\
The \(\varphi\) axis is identified with the electron.&Mixture of the fiber of
constants and the material realization.\\
Pell is interpreted as energy amplification.&Violation of scale reading and physical conservation.\\
Electromagnetic helicity is identified with spin \(1/2\).&Confusion between
distinct representations.\\
A point cell is imposed on every species.&Contradiction with the character
of a fibre or multisection of the noncentral species.\\
The five constructions of the continuum are separated.&Breakdown of the previous unique genealogy prior to any geometric publication.\\
\bottomrule
\end{longtable}

\chapter{Mathematical Glossary}

\begin{longtable}{p{.22\textwidth}p{.66\textwidth}}
\caption{Symbols and Key Terms.}
\label{espiral:tab:glosario}\\
\toprule
Symbol or term&Definition\\
\midrule
\endfirsthead
\toprule
Symbol or term&Definition\\
\midrule
\endhead
\(\Sigma,\Pi\)&Additive and multiplicative evaluations on the common support APP.\\
\(\tau\)&Local TRIT orientation regime; it is neither the radial degree nor
Frenet torsion.\\
\(M_{\ph}\)&Tritic operational-phase selector.\\
\(\Gamma_9\)&Nonadic holonomy of enriched dynamics.\\
\(H_N\)&Accumulated affine holonomy up to the prefix \(N\).\\
\(\Lambda_N,\mathcal C_N\)&Linear part and translational cocycle of \(H_N\).\\
\(p\)&Integer cocycle of the radial character.\\
\(a\)&Local orientation character that separates genealogies of the same
local degree.\\
\(L_p,\exp_p\)&Power-logarithmic coordinate and its positive inverse.\\
\(\boxplus_p\)&Deformed law that publishes the sum-product interaction.\\
\(\Led_N\)&Enriched ledger of sheet, carry, boundary, route, and memory.\\
\(\mathcal P_{\rad}^{\enr}\)&Radial reader enriched prior to the curve.\\
\(\mathcal P_{\arq}^{\mathrm{esp}}\)&Subsequent evaluation on a curve in
\(\R^2\).\\
\(\kappa_{\mathrm{F}},\mathcal T_{\mathrm{F}}\)&Curvature and torsion of Frenet of the
visible curve.\\
Radial helical suspension&Variable-radius surface or sheet that
preserves phase, depth, and path.\\
Unfaithful publication&Assessment that can identify distinct histories.\\
Constructivo--generativo&Method that constructs the object and its genealogy
before recognizing its conventional appearance.\\
\bottomrule
\end{longtable}
  \chapter{Dictionary of domains, operators and representations}
\chapter*{Notation, types, and symbol disambiguation}
\addtocontents{toc}{\protect\enlargethispage{2\baselineskip}}
\addcontentsline{toc}{chapter}{Notation, types, and symbol disambiguation}

\begin{longtable}{@{}p{.22\textwidth}p{.70\textwidth}@{}}
\toprule
Symbol & Object \\ \midrule
\endhead
\(\mathcal A_9\)&Nonadic arithmetic torus of paths with two internal evaluations.\\
\(\mathrm{TRIT}\)&Oriented ternary coordinate of the transport; the visible
output does not exhaust the enriched lift with carry.\\
\(\mathrm{TPK}\)&Topological phase kernel with integral holonomic state, transitions, routes, extensions and memory.\\
\(M_{\rm ph}\)&Tritic operational-phase selector, distinct from cardinal
transport and stroboscopic sampling.\\
\(\tau_3\)&Primitive phase block \((+1,-1,0)\).\\
\(\mathcal N,\mathcal E,\mathcal S,\mathcal O\)&Generators of transport
cardinal.\\
\(\operatorname{Obs}_3\)&Stroboscopic sampling of every third step.\\
\(\mathcal R_{12}\)&Internal TPK-oriented dodecaphasic temporal system.\\
\(\Theta_{108}=\mathbb Z/108\mathbb Z\)&Observable calendar
\(12\times9\); it is not the complete TPK state.\\
\(C_{12}\)&Rotation of the twelve temporal sectors.\\
\(\Gamma_{108}\)&Oriented transport orbit of \(108\) steps.\\
\(\mathcal G_9\)&Unitary pruning within the extension and
boundary relation; it does not by itself constitute the monodromy.\\
\(\mathfrak M_9\)&Nonadic monodromy with memory: it returns the phase and changes
the state.\\
\(\pi_{\rm term}\)&Projection of the terminal holonomic state onto the three selected visible blocks.\\
\(R_{\rm sgn}\)&Reading of the signed integer coordinate of the same terminal holonomic state.\\
\(U_{12}^{\rm cat}\)&Visible concatenation \(U_6\Vert U_6\).\\
\(U_{12}^{\rm sgn}\)&Signed coordinate of the integral holonomic state, reversible with \(K\) through \(\mathcal H_{12}\); it is not deduced from \(U_{12}^{\rm cat}\), the calendar projection or the three isolated terminal blocks.\\
\(\mathcal H_{12}\)&Reversible dodecaphase transformation between
\(U_{12}^{\rm sgn}\) and \(K\).\\
\(\mathcal E_{\rm dod}\)&Integral extractor
\((D_3,D_4,Q):\mathbb Z^{12}\to\mathbb Z^{25}\) that reconstructs \(K\).\\
\(\overline{\mathcal C}_9^{\,3}\)&Nonadic cubic completion.\\
\(A,C^*\)&Intrinsic angular and antiangular coordinates.\\
\(q_\pm\)&Oriented channels of the double-layer reader.\\
\(D_{\rm raw}^{90,120}\)&Raw diagnostic difference of series in a
fixed chart; it is not a canonical extractor.\\
\(\iota_{6,8}\)&Inclusion
\(\mathcal F_6\hookrightarrow\mathcal F_8\) of hexad--octad incidences.\\
\(K_{6\to8}\)&Angular functional of the two Lambert subtails, subsequent to
\(q_\pm\), the series, the pole, positivity, and minimality.\\
\(\langle90,120\rangle\)&Spectral semigroup governing the hierarchy of
towers.\\
\(\varepsilon_{\rm res}^{\circ}\)&Residue of the exact angular identity;
it is not a finite extraction from the preceding 90/120 objects.\\
\(\Delta_4\)&Minimal torsion or global torsional invariant prior to the
species; it is not Cartan torsion.\\
\(\mathsf T_{\alpha,\infty}\)&Remote prolongation of the carry along the entire
path of \(\alpha\), typed separately from the orbital generator
\(T_1^{\rm orb}=L_1\), the neutral counter, and the calendar.\\
\(N_1^{\rm neu}\)&Number of neutral positions in the first window of the
transport example; it equals $3$ and is not an operator.\\
\(T_1^{\rm orb}:=L_1\)&Second generator in the local orbital gauge; in its
reference section it is abbreviated \(T_1\).\\
\(\varepsilon_{\rm clk}\)&Energy quantum of the clock,
\(\hbar_{\rm int}2\pi/(108t_0)\); it is not a canonical energy basis unless the
temporal frame is specified. The symbol \(\varepsilon_0\) is reserved for
the vacuum permittivity at the SI constitutive interface.\\
\(\mathrm{CLK}1\text{--}\mathrm{CLK}6\)&Labels of the six equations of the
finite clock, from the Hamiltonian to the Schrödinger equation.\\
\(\mathcal K_{27}\)&Compound return of twenty-seven updates.\\
\(X_\infty^{\rm cal}\)&Common enriched domain of global evaluations.\\
\bottomrule
\end{longtable}

\section*{Repeated cardinalities with different sources}

Equality of cardinalities does not automatically transport domain,
operation, memory, or evidentiary force. The following table sets out the
collisions that most directly affect interpretation.

\begin{longtable}{@{}p{.12\textwidth}p{.27\textwidth}p{.51\textwidth}@{}}
\toprule
Index & Objeto tipado & Function and separation \\ \midrule
\endhead
\(108\)&\(\Theta_{108}\)&The \(12\times9\) calendar closes the observable
clock.\\
\(108\)&\(\Pi_{108}^{\rm cal}\)&Calendrical projection of a cursor with \(324\)
seeds; it forgets the second cursor, orientation, carry, boundary, enriched
chronological record, and
survival.\\
\(108\)&\(w_{108}\)&Word of \(108=18\times6\) trits; it is not the calendar.\\
\(104,976\)&\(X_{\rm TPK}^{\rm seed}\)&Two-cursor engine,
\(104,976=324^2\); it preserves a domain distinct from that of the calendar.\\
\(1080\)&\(w_{1080}\)&Length of a word.\\
\(1080\)&\(\Gamma_0(1080)\)&Level of a modular subgroup; requires its
own realization.\\
\(729\)&\(\mathbb F_3^6\)&Ternary word environment.\\
\(729\)&\(9^3\)&Cubic cell/interior EMRH.\\
\(729\)&\(\operatorname{Ext}(X_K)\)&Number of subcylinders examined at each
prolongation stage.\\
\(729\)&bulk EMRH&Inner layer of the base completion mil.\\
\(27\)&\(\mathcal K_{27}\)&Composite return operator; not the name of an
autonomous theory.\\
\(27\)&cubic stratum&Contribution of two boundary coordinates in
\(999=729+243+27\).\\
\(27\)&module or classifier&Each occurrence requires its map before
identifying itself with another.\\
\bottomrule
\end{longtable}

\section*{Six distinct sources of 90/120}

\begin{longtable}{@{}
>{\raggedright\arraybackslash}p{.21\textwidth}
>{\raggedright\arraybackslash}p{.42\textwidth}
>{\raggedright\arraybackslash}p{.27\textwidth}@{}}
\toprule
Object & Domain and operation & Exit and type prohibition \\ \midrule
\endhead
\(\mathcal E_{\rm dod}\)&\(K\in\mathbb Z^{12}\); stepwise differences
\(3,4\) and sum \(Q\)&Transverse dodecaphasic coordinates; no incidence
nor angular series.\\
\(D_{\rm raw}^{90,120}\)&A \(q\) chart; raw difference
\(S_{90}(q)-S_{120}(q)\)&Scale diagnostic, distinct from the
Barbero--Immirzi functional and the incidence operators.\\
\(\iota_{6,8}\)&\(H\subset O\); inclusion of
\(H\times\binom{H}{2}\) in \(O\times\binom{H}{2}\)&Cardinals \(90,120\) and
defect \(-1/12\); neither an extractor nor a series.\\
\(K_{6\to8}\)&Channels \(q_\pm\) and Lambert subtails; combination
\(12\Delta S_{90}-\Delta S_{120}\)&Angular functional relative to its
premises; it is not transported as a mass block.\\
\(\mathfrak S=\langle90,120\rangle\)&Moment heights
\(s_n=q_-^n-q_+^n\)&Global tower semigroup; does not end at eight heights.\\
\bottomrule
\end{longtable}

\section*{Four realizations of bidirectional reversibility}

\begin{longtable}{@{}
  >{\raggedright\arraybackslash}p{.26\textwidth}
  >{\raggedright\arraybackslash}p{.26\textwidth}
  >{\raggedright\arraybackslash}p{.38\textwidth}@{}}
\toprule
Operation & Domain & Action \\ \midrule
\endhead
\(P_\pm=(I\pm R)/2\)&Split Algebra&Project the two spectral sectors.\\
U-turn&Route and composite calendar&Invert the orientation of the
transport without erasing its word.\\
\(\operatorname{Ext}^{\pm}\)&Extension Tree&It compares future
and past valences; the balance may be read through
\(\log(\#\operatorname{Ext}^+/\#\operatorname{Ext}^-)\).\\
\(C_M(\tau)=-\operatorname{rev}(\tau)\)&Dodecaphonic Word&It conjugates order
and sign; it is not the cellular involution \(\rho_c\).\\
\bottomrule
\end{longtable}

The common function of these four realizations is to preserve reversible
orientation. It does not authorize the identification of phase return, route
return, word return, and state return.
  
\backmatter
\begingroup
\sloppy
\small
\setlength{\parskip}{0pt}
\setlength{\emergencystretch}{3em}
\bibliographystyle{alphaurl}
\begin{thebibliography}{HvdHKN76}

\bibitem[ALY17]{AbeLamYamada2017}
Toshiyuki Abe, Ching~Hung Lam, and Hiromichi Yamada.
\newblock A remark on {$\mathbb Z_p$}-orbifold constructions of the moonshine vertex operator algebra.
\newblock {\em arXiv preprint arXiv:1705.09022}, 2017.
\newblock \href {https://arxiv.org/abs/1705.09022} {\path{arXiv:1705.09022}}.

\bibitem[AW68]{rm:ArakiWoods1968}
H.~Araki and E.~J. Woods.
\newblock A classification of factors.
\newblock {\em Publ. Res. Inst. Math. Sci. Ser. A}, 4:51--130, 1968.
\newblock URL: \url{https://ems.press/content/serial-article-files/41616}.

\bibitem[Bek73]{rm:Bekenstein1973}
J.~D. Bekenstein.
\newblock Black holes and entropy.
\newblock {\em Physical Review D}, 7:2333--2346, 1973.

\bibitem[BG95a]{Barbero}
J.~Fernando Barbero~G.
\newblock Real {Ashtekar} variables for {Lorentzian} signature space-times.
\newblock {\em Physical Review D}, 51:5507--5510, 1995.

\bibitem[BG95b]{Barbero1995}
J.~Fernando Barbero~G.
\newblock Real {Ashtekar} variables for {Lorentzian} signature space-times.
\newblock {\em Physical Review D}, 51:5507--5510, 1995.

\bibitem[Bor92]{Borcherds1992}
Richard~E. Borcherds.
\newblock Monstrous moonshine and monstrous lie superalgebras.
\newblock {\em Inventiones Mathematicae}, 109:405--444, 1992.

\bibitem[{Bur}19]{BIPMSI}
{Bureau International des Poids et Mesures}.
\newblock {\em The International System of Units ({SI} Brochure)}.
\newblock BIPM, S\`evres, 9 edition, 2019.
\newblock Updated online edition.
\newblock \href {https://doi.org/10.59161/AUEZ1291}
  {\path{doi:10.59161/AUEZ1291}}.

\bibitem[Car14]{Caratheodory1914}
Constantin Carath{\'e}odory.
\newblock {"U}ber das lineare ma{\ss} von punktmengen---eine verallgemeinerung
  des l{"a}ngenbegriffs.
\newblock {\em Nachrichten von der Gesellschaft der Wissenschaften zu
  G{"o}ttingen, Mathematisch-Physikalische Klasse}, pages 404--426, 1914.

\bibitem[Car17]{Carnahan2017}
Scott Carnahan.
\newblock 51 constructions of the moonshine module.
\newblock {\em arXiv preprint arXiv:1707.02954}, 2017.
\newblock \href {https://arxiv.org/abs/1707.02954} {\path{arXiv:1707.02954}}.

\bibitem[CLS16]{ChenLamShimakura2016}
Hsian-Yang Chen, Ching~Hung Lam, and Hiroki Shimakura.
\newblock A {$\mathbb Z_3$}-orbifold construction of the moonshine vertex
  operator algebra and some maximal 3-local subgroups of the monster.
\newblock {\em arXiv preprint arXiv:1606.05961}, 2016.
\newblock \href {https://arxiv.org/abs/1606.05961} {\path{arXiv:1606.05961}}.

\bibitem[CN79]{ConwayNorton1979}
John~H. Conway and Simon~P. Norton.
\newblock Monstrous moonshine.
\newblock {\em Bulletin of the London Mathematical Society}, 11:308--339, 1979.

\bibitem[{COD}24]{CODATA2022}
{CODATA Task Group on Fundamental Constants}.
\newblock The 2022 {CODATA} recommended values of the fundamental physical
  constants, 2024.
\newblock URL: \url{https://physics.nist.gov/constants}.

\bibitem[CS99]{ConwaySloane1999}
John~H. Conway and Neil J.~A. Sloane.
\newblock {\em Sphere Packings, Lattices and Groups}.
\newblock Springer, 3 edition, 1999.

\bibitem[Fei78]{rm:Feigenbaum1978}
Mitchell~J. Feigenbaum.
\newblock Quantitative universality for a class of nonlinear transformations.
\newblock {\em Journal of Statistical Physics}, 19(1):25--52, 1978.
\newblock \href {https://doi.org/10.1007/BF01020332}
  {\path{doi:10.1007/BF01020332}}.

\bibitem[FLM88]{FLM1988}
Igor Frenkel, James Lepowsky, and Arne Meurman.
\newblock {\em Vertex Operator Algebras and the Monster}.
\newblock Academic Press, 1988.

\bibitem[GH77]{rm:GibbonsHawking1977}
G.~W. Gibbons and S.~W. Hawking.
\newblock Cosmological event horizons, thermodynamics, and particle creation.
\newblock {\em Physical Review D}, 15:2738--2751, 1977.

\bibitem[Haw75]{rm:Hawking1975}
S.~W. Hawking.
\newblock Particle creation by black holes.
\newblock {\em Communications in Mathematical Physics}, 43:199--220, 1975.

\bibitem[HJ13]{HornJohnson2013}
Roger~A. Horn and Charles~R. Johnson.
\newblock {\em Matrix Analysis}.
\newblock Cambridge University Press, 2 edition, 2013.

\bibitem[Hol96]{Holst1996}
Sören Holst.
\newblock Barbero's hamiltonian derived from a generalized hilbert--palatini
  action.
\newblock {\em Physical Review D}, 53:5966--5969, 1996.

\bibitem[HvdHKN76]{Hehl1976}
Friedrich~W. Hehl, Paul von~der Heyde, G.~David Kerlick, and James~M. Nester.
\newblock General relativity with spin and torsion: Foundations and prospects.
\newblock {\em Reviews of Modern Physics}, 48:393--416, 1976.

\bibitem[IK02]{IosifescuKraaikamp}
M.~Iosifescu and C.~Kraaikamp.
\newblock {\em Metrical Theory of Continued Fractions}.
\newblock Kluwer, 2002.

\bibitem[Imm97a]{Immirzi}
Giorgio Immirzi.
\newblock Real and complex connections for canonical gravity.
\newblock {\em Classical and Quantum Gravity}, 14:L177--L181, 1997.

\bibitem[Imm97b]{Immirzi1997}
Giorgio Immirzi.
\newblock Real and complex connections for canonical gravity.
\newblock {\em Classical and Quantum Gravity}, 14:L177--L181, 1997.

\bibitem[Khi64]{Khinchin}
A.~Ya. Khinchin.
\newblock {\em Continued Fractions}.
\newblock University of Chicago Press, 1964.

\bibitem[Kib61]{Kibble1961}
T.~W.~B. Kibble.
\newblock Lorentz invariance and the gravitational field.
\newblock {\em Journal of Mathematical Physics}, 2:212--221, 1961.

\bibitem[K{"o}36]{Koenig1936}
D{\'e}nes K{"o}nig.
\newblock {\em Theorie der endlichen und unendlichen Graphen}.
\newblock Akademische Verlagsgesellschaft, Leipzig, 1936.

\bibitem[Kol33]{Kolmogorov1933}
Andrei~N. Kolmogorov.
\newblock {\em Grundbegriffe der Wahrscheinlichkeitsrechnung}, volume 2, Heft 3
  of {\em Ergebnisse der Mathematik und ihrer Grenzgebiete}.
\newblock Julius Springer, Berlin, 1933.

\bibitem[Lan82]{rm:Lanford1982}
III Lanford, O.~E.
\newblock A computer-assisted proof of the {Feigenbaum} conjectures.
\newblock {\em Bulletin of the American Mathematical Society}, 6:427--434,
  1982.

\bibitem[Man05]{Manivel2005}
Laurent Manivel.
\newblock Configurations of lines and models of lie algebras.
\newblock arXiv:math/0507118, 2005.
\newblock \href {https://arxiv.org/abs/math/0507118}
  {\path{arXiv:math/0507118}}.

\bibitem[MS64]{rm:MisnerSharp1964}
C.~W. Misner and D.~H. Sharp.
\newblock Relativistic equations for adiabatic, spherically symmetric
  gravitational collapse.
\newblock {\em Physical Review}, 136:B571--B576, 1964.

\bibitem[MS77]{MacWilliamsSloane1977}
Florence~Jessie MacWilliams and Neil J.~A. Sloane.
\newblock {\em The Theory of Error-Correcting Codes}.
\newblock North-Holland, 1977.

\bibitem[Pla01]{Planck}
M.~Planck.
\newblock Ueber das gesetz der energieverteilung im normalspectrum.
\newblock {\em Annalen der Physik}, 309:553--563, 1901.

\bibitem[Sci64]{Sciama1964}
Dennis~W. Sciama.
\newblock The physical structure of general relativity.
\newblock {\em Reviews of Modern Physics}, 36:463--469, 1964.

\bibitem[Ser77]{Serre1977Representations}
Jean-Pierre Serre.
\newblock {\em Linear Representations of Finite Groups}.
\newblock Springer, 1977.

\bibitem[Tak03]{rm:Takesaki}
M.~Takesaki.
\newblock {\em Theory of Operator Algebras II}.
\newblock Springer, 2003.

\bibitem[TMNT21]{Tiesinga2021CODATA}
Eite Tiesinga, Peter~J. Mohr, David~B. Newell, and Barry~N. Taylor.
\newblock {CODATA} recommended values of the fundamental physical constants:
  2018.
\newblock {\em Journal of Physical and Chemical Reference Data}, 50(3):033105,
  2021.
\newblock \href {https://doi.org/10.1063/5.0064853}
  {\path{doi:10.1063/5.0064853}}.

\bibitem[Wie93]{Wien}
W.~Wien.
\newblock Eine neue beziehung der strahlung schwarzer k{\"o}rper zum zweiten
  hauptsatz der w{\"a}rmetheorie.
\newblock {\em Sitzungsberichte der K{\"o}niglich Preussischen Akademie der
  Wissenschaften}, pages 55--62, 1893.

\end{thebibliography}
 
\endgroup
\begingroup
\setlength{\columnsep}{2.2em}\fontsize{9}{10.4}\selectfont
\renewcommand{\indexspace}{\par\vspace{3pt}}
\printindex
\endgroup
\end{document}
